\documentclass[11pt,letterpaper]{article}

\usepackage[T1]{fontenc}
\usepackage[utf8]{inputenc}
\usepackage{lmodern}
\usepackage[margin=0.79in]{geometry}
\usepackage{amsmath,amssymb,amsthm,amscd,mathtools,mathrsfs}
\usepackage{tikz-cd}
\usepackage{booktabs,longtable,array}
\usepackage{bm}
\usepackage{enumitem}
\usepackage{xcolor}
\usepackage{microtype}
\usepackage{xurl}
\usepackage[hidelinks,hypertexnames=false]{hyperref}
\usepackage[nameinlink,noabbrev]{cleveref}
\setlength{\emergencystretch}{2em}
\makeatletter
\renewcommand*\l@subsection{\@dottedtocline{2}{1.5em}{3.4em}}
\renewcommand*\l@subsubsection{\@dottedtocline{3}{3.8em}{4.4em}}
\makeatother

\definecolor{Verified}{HTML}{176B3A}
\definecolor{Conditional}{HTML}{9A6700}
\definecolor{Open}{HTML}{B42318}
\definecolor{Refuted}{HTML}{7A271A}
\definecolor{Provenance}{HTML}{475467}

\newcommand{\Status}[2]{\textcolor{#1}{\textbf{[#2]}}}
\newcommand{\Verified}{\Status{Verified}{VERIFIED}}
\newcommand{\Conditional}{\Status{Conditional}{CONDITIONAL}}
\newcommand{\Open}{\Status{Open}{OPEN}}
\newcommand{\Refuted}{\Status{Refuted}{REFUTED}}
\newcommand{\Source}[1]{\textcolor{Provenance}{\nolinkurl{#1}}}

\newtheorem{theorem}{Theorem}[section]
\newtheorem{proposition}[theorem]{Proposition}
\newtheorem{lemma}[theorem]{Lemma}
\newtheorem{corollary}[theorem]{Corollary}
\newtheorem{example}[theorem]{Example}
\newtheorem{counterexample}[theorem]{Counterexample}
\theoremstyle{definition}
\newtheorem{definition}[theorem]{Definition}
\newtheorem{auditrecord}[theorem]{Audit record}
\newtheorem{test}[theorem]{Validation, repair, or disproof test}
\theoremstyle{remark}
\newtheorem{remark}[theorem]{Remark}

\crefname{auditrecord}{audit record}{audit records}
\Crefname{auditrecord}{Audit record}{Audit records}
\crefname{test}{validation/repair/disproof test}{validation/repair/disproof tests}
\Crefname{test}{Validation/repair/disproof test}{Validation/repair/disproof tests}

\title{The Claimed Complex Structure on $S^6$\\
Validate, Repair, or Disprove}
\author{Working mathematical notebook following the three-stream proof audit}
\date{6 September 2026}

\begin{document}
\ifdefined\SIXWorkbenchEmbedded
\section*{Part II. Complete validate/repair/disprove paper}
\addcontentsline{toc}{section}{Part II. Complete validate/repair/disprove paper}
\else
\maketitle
\fi

\begin{abstract}
\noindent\textbf{6 September: explicit spectral and umbral maps, and the original family's parameter identifications.}
For the literal common lattice $\mathcal I=\{\sqrt2(a,a,a):a\in D_4\}$,
write $a=a_0+v$, $t^2=|v|^2$, $Q=a_0^2+t^2$ and
$P=a_0(a_0^2-3t^2)$. Its complete cubic image is
\[
 \tau(\mathcal I)=2\sqrt2\bigl(18\mathbb Z\cup(6\mathbb Z+2)
                                      \cup(6\mathbb Z-2)\bigr).
\]
Three explicit integral affine lines give a section of this value map.
The scalar-retaining Albert matrix $X=\mu I+Y(a)$, with cyclic
off-diagonal entries $\sqrt2a$, has invariants and labelled spectrum
\[
\begin{aligned}
 (\operatorname{tr}X,\operatorname{tr}X^2,\det X)
   &=(3\mu,3\mu^2+12Q,\mu^3-6\mu Q+4\sqrt2P),\\
 \operatorname{Spec}Y
   &=(2\sqrt2a_0,-\sqrt2a_0-\sqrt6t,-\sqrt2a_0+\sqrt6t).
\end{aligned}
\]
For $t>0$, the explicit labelled eigenprojector satisfies
$(\Pi_1)_{12}=-1/6-\sqrt3(v/t)/6$; thus the frame and spectrum
recover $a$, not only its norm and cubic. The full discriminant is
$864t^2(3a_0^2-t^2)^2$. The full Cartan family retains the scalar
$\mu$, radius $r$, and angle $\theta$, with $r^2=12Q$ and
$r^3\cos3\theta=24\sqrt3P$ on this subfamily.

For the same marked root system $A_5^4D_4$, the
Cheng--Duncan--Harvey prescription gives
$S^X=(5I_{12}+J)S_6$, where $J_{rs}=\mathbf1_{s=5r\ (12)}$ and
$S_{6,r}=\sum_{k\equiv r\ (12)}kq^{k^2/24}$.
The coefficient map to the retained Fricke function is
\[
 T_6=\frac{\eta(\tau)^5\eta(3\tau)}{\eta(2\tau)\eta(6\tau)^5},
 \qquad f^X=-q\frac{d}{dq}\log T_6,
 \qquad T_6\!\left(-\frac1{6\tau}\right)=\frac{72}{T_6(\tau)}.
\]
All twelve theta indices, the multiplier determinant, the original
glue-group action and its twisted shadows are computed in the proofs.
This is an explicit use and comparison of established umbral machinery,
not a claim to a new moonshine module.

In the original analytic family $\beta_c=\beta^{\rm part}+c$,
$\operatorname{Im}c<-M$, the actual global identifications are exactly
$X_c\simeq X_{c'}\Longleftrightarrow c'-c\in2\mathbb Z$.
The middle lattices agree for every integer shift; the multiplicity-four
filling retains an order-two obstruction to odd shifts. Even shifts have
explicit finite-filling and toric cusp maps and a proved global additive
gluing correction. The resulting parameter quotient is
$w=e^{-\pi ic}$, $0<|w|<e^{-\pi M}$.
This classifies the retained analytic family without using the claimed
identification of its total space with $S^6$.


\noindent\textbf{6 September: cubic arithmetic, integral lattice maps, and the canonical ring.}
For the retained $A_5^4\perp D_4$ Niemeier lattice $N$ and its
original index-six Leech neighbour $\Lambda$, fix the Cayley--Dickson
frame $(1,i,j,k,\ell,i\ell,j\ell,k\ell)$ used in the complete cubic
proof. Put $\tau(q)=\operatorname{Re}((q_1q_2)q_3)$ and
$F=\tau\circ\mathcal L$ on the original Niemeier span,
and $\mathcal V_F(L)=\operatorname{span}_{\mathbb Z}\{F(y):y\in L\}$.
The exact additive groups are
\[
\begin{aligned}
 \mathcal V_F(N)&=\tfrac1{27}\mathbb Z\oplus
       \tfrac{\sqrt2}{54}\mathbb Z\oplus\tfrac{\sqrt3}{18}\mathbb Z,\\
 \mathcal V_F(\Lambda)&=\tfrac1{27}\mathbb Z\oplus
       \tfrac{\sqrt2}{216}\mathbb Z\oplus\tfrac{\sqrt3}{54}\mathbb Z,
 \qquad \mathcal V_F(\Lambda)/\mathcal V_F(N)\cong\mathbb Z/12.
\end{aligned}
\]
The Albert determinant is $\lambda^3-\lambda(y,y)+2F(y)$.
The quotient concerns the displayed value-generated groups in this fixed
frame, not an inclusion $N\subset\Lambda$. The actual residue
$r=\rho(361)$ cycles three Leech neighbours $\Lambda_j=r^j\Lambda$.
Their pairwise intersection is one integral lattice $J$, with
$\det J=81$ and $[\Lambda_j:J]=9$. The fixed fourth isotropic plane
in $J^\vee/J$ gives $M$ with exactly $72$ roots of type $A_2^{12}$,
$M\cap\Lambda_j=J$, and $[M:J]=9$.
\[
 \mathcal V_F(J)=\mathcal V_F(M)
 =\tfrac1{27}\mathbb Z\oplus\tfrac{\sqrt2}{216}\mathbb Z
                         \oplus\tfrac{\sqrt3}{18}\mathbb Z.
\]
The successive value-group inclusions have matrices
$\operatorname{diag}(1,4,1)$ and $\operatorname{diag}(1,1,3)$,
giving quotients $\mathbb Z/4$ and $\mathbb Z/3$.
Both containments and their reverse generator certificates are calculated
from the full cubic polarization, not inferred from lattice indices.

Wilson's construction \cite{esn:Wilson2009}, transported by the explicit
isometry $\iota(x)=x/\sqrt2$, gives $\Lambda_{\rm cyc}$ in the
original triality metric. Its cyclic fixed Gram matrix is $3G_{E_8}$;
its orthogonal complement is explicitly isometric to $A_2\otimes E_8$,
and their sum has index $3^8$, with quotient $\mathbb F_3^8$.
The directions remain $c(q_1,q_2,q_3)=(q_3,q_1,q_2)$ and
$c_{\rm W}=c^{-1}$. Its literal intersections are
\[
 \Lambda_{\rm cyc}\cap\mathcal L(N)
 =\Lambda_{\rm cyc}\cap\mathcal L(\Lambda)
 =\Delta:=\{\sqrt2(a,a,a):a\in D_{\mathbb H}\},
\]
where $D_{\mathbb H}$ is the even-coordinate-sum integer lattice in
the first four stated Cayley coordinates. Thus $\Delta$ has rank $4$,
Gram matrix $6G_{D_4}$, determinant $5184$, minimum squared norm $12$,
and value-generated subgroup $4\sqrt2\mathbb Z$ for $\tau|_\Delta$, or
$8\sqrt2\mathbb Z$ for $2\tau|_\Delta$. This is not the actual image:
$12\sqrt2$ is not attained by $\tau$ on $\Delta$.
The preceding proofs retain the triality products and spectral boundary
graph, the flag's degree-$+1$ collapse, and the retained cusp class's
nontrivial three-plane reduction; these exact maps do not supply an
interacting quantum mass-gap result.

For the original analytic fibration $f:X\to\mathbb P^1$, the full
canonical-divisor proof gives
$\omega_X\simeq\mathcal O_X(-2S_2)$ and the graded anticanonical ring
$\bigoplus_{m\ge0}H^0(X,\omega_X^{-m})\cong\mathbb C[U,V]$,
with $\deg U=1$, $\deg V=2$. The complete system $|-2K_X|$
recovers $f$ with the specified projective coordinates. This calculation
uses the original analytic charts and finite fillings, not an
identification of $X$ with $S^6$, and draws no sphere-topology conclusion.

The retained constant $\beta_c=\beta^{\rm part}+c$ now gives a proper
holomorphic deformation over $\operatorname{Im}c<-M$, locally on
relatively compact parameter discs. On a fixed regular fibre set
$T=\operatorname{Im}\tau$, $m=\operatorname{Im}\mu$,
$D=\operatorname{Im}\beta_{c_*}-6m^2/T<0$ and $\delta=c-c_*$.
In the plus-graph convention its exact Beltrami matrix is
\[
 \nu_\delta=\frac{\delta}{2iD+\delta}
                 \begin{pmatrix}0&0\\-m/T&1\end{pmatrix}.
\]
The connecting Kodaira--Spencer class is
$-[\partial_{\zeta_2}\otimes
(d\bar\zeta_2-(m/T)d\bar\zeta_1)/(2iD)]\ne0$.
The intrinsic pencil recovers $f$ also over the dual numbers; its three
fixed critical values prohibit an infinitesimal projective base change.
Consequently this deformation of the whole $X_c$ is nontrivial,
not only its marked regular-fibre deformation. The proof uses no sphere
recognition or CDP conclusion.



The 5 September continuation adds the higher-torus-wrap and octonionic
comparison in Section~\ref{sec:h24-continuation}. Its equal-diagonal
positive octonionic cone has compactified pair $(B^{25},S^{24})$.
The flag $F_4/\operatorname{Spin}(8)$ has tangent module
$8_v\oplus8_s^+\oplus8_s^-$ and integral Poincar\'e polynomial
$1+2t^8+2t^{16}+t^{24}$. The peripheral character has values
$(-1,4,-3)$ in $\mathbb Z/12$, and sixfold suspension identifies
$\pi_6(S^3)\cong\mathbb Z/12$ with the even subgroup of
$\pi_{12}(S^9)\cong\mathbb Z/24$. The particular cusp map now has
coefficient zero: its entire regular inverse image lies in one original
toric chart and its actual normal framing bounds an explicit
$D^2\times T^2$ in that chart times an interval. This proves
$g:X\to S^3$ based null-homotopic before any sphere identification;
see Theorem~\ref{h24:thm:cusp-g-zero}.
Retaining the angle in the same conical quotient gives
$H=p_3q_2:X\to S^4$ with $[H]=\eta_4\eta_5\ne0$.
Theorem~\ref{h24:ret:main} computes its entire inverse torus, its
literal normal four-frame, a periodic $SU(3)$ contraction, and Arf
invariant one. Hence $q_2$ is nonnull. The exact Hopf--cusp map
$\eta_2\wedge H:S^9\to S^6$ represents $12\nu_6\ne0$ in
$\mathbb Z/24$; after three suspensions it represents $12\nu_9$
in $\pi_{12}(S^9)$. Its peripheral composite has values $(12,0,12)$
and kernel $\chi^{-1}(2\mathbb Z/12)$, retaining parity rather than
lifting $\chi$ under reduction modulo twelve.
Section~\ref{sec:fifteen-selector-jordan}
also gives the full map from the fifteen marked $p=1201$ selector rays
to a three-dimensional real Jordan subalgebra of
$\operatorname{Herm}_3(\mathbb H)$, with exact census certificates.

Section~\ref{esn:sec:shared} constructs the modulo-$840$ lattice
$R_{\mathrm{ES}}=A_5^4\perp D_4\subset N$ with index $72$,
$\det N=1$, and exactly the original $144$ roots. Its coordinate/code
maps give $O(N)=W(R_{\mathrm{ES}})\rtimes L(G)$,
$G\cong GL_2(\mathbb F_3)$, realizing the Weyl quotient recorded by
Cheng--Duncan--Harvey \cite[Table~3]{esn:CDH2014}.
The residue six-cycle lies in $W(R_{\mathrm{ES}})$ and centralizes this
chosen $L(G)$ (Theorems~\ref{esn:fa:thm:full-quotient}
and~\ref{esn:fa:thm:residue-weyl}). The marked $D_4$ dual classes
recover the vector and two half-spin weight sets of the flag tangent
modules. Exact shell counts retain the prime and shell: at $R=3$ the
same residue-one primes $2521,3361$ have raw counts $0,6$ and unordered
counts $0,3$, respectively.

The smooth literal scalar--Gauss continuation gives
$\mathcal S_{\rm sg}\cong\{1,-1\}^2\times(0,\infty)
\times\mathbb R^2\times\mathbb T^3$, with the original coupling
$t=g^2$ as the positive coordinate. Its four components are embedded
and six-dimensional; every fixed positive coupling has four
five-dimensional components $\mathbb R^2\times\mathbb T^3$
(Corollary~\ref{sgd:smooth-addon:levels}). Both spatial equations
remain outside this domain, and this refinement gives no quantum gap.
For the cusp-regular periods, Proposition~\ref{prop:period-exact-admissible-constants}
gives the exact admissible constants
$\operatorname{Im}c_0<-\max_{B\setminus\{p_0\}}D_{\rm part}$.
The cusp-free affine solution $\widetilde\beta=\beta+i\,t\circ\pi$
has $\widetilde D=D+\operatorname{Re}t$ unbounded above.
The cusp edge maps retain the multipliers
$e^{-2\pi i\mu_0}$, $e^{12\pi i\mu_0}$ and
$e^{-2\pi i(7\mu_0+b_0)}$ in their stated character conventions
(equation~\eqref{eq:cusp-exact-opposite-edge}). Restricted chart domains
and closed-polydisc inequalities are explicit in
Lemma~\ref{lem:toric-chart-control}. These are period and local-quotient
refinements, with no additional sphere or CDP conclusion inferred.

This paper is the consolidated successor to the former 56-page workbench: it
incorporates the nonduplicate proofs from the subsequent 18-page repair branch
and the later Reviewer~B/Codex validation work.  It
validates, repairs, or attempts to disprove a reconstructed
manuscript claiming a compact complex threefold diffeomorphic to $S^6$ and
 of algebraic dimension one.  In its present form the workbench gives complete
 proofs of the simultaneous elliptic-orientation normalization.  Let
 \(\Delta\simeq(\mathbb Z/3)*(\mathbb Z/4)\) act on \(\mathfrak h_z\), let
 \(g_1,g_2\) have fixed points \(z_1,z_2\), put
 \(g_0=(g_1g_2)^{-1}\), and let
 \(\tau:\mathfrak h_z\to\mathfrak h\) be the normalized period lift.  Define
 \[
  \varphi:\Delta\longrightarrow\operatorname{PSL}_2(\mathbb Z),\qquad
  \tau\circ g=\varphi(g)\circ\tau,
 \]
 and
 \[
  \mathfrak S_1(w)=\frac{w-1}{w},\qquad
  \mathfrak S(w)=-\frac1w,\qquad \mathfrak P(w)=w-1.
 \]
 Then
 \[
  g_1'(z_1)=e^{-2\pi i/3},\qquad g_2'(z_2)=-i,\qquad
  g_0\ \text{is parabolic},
 \]
 and
 \[
  \varphi(g_1)=\mathfrak S_1,\qquad
  \varphi(g_2)=\mathfrak S,\qquad
  \varphi(g_0)=\mathfrak P,\qquad
  \mathfrak S_1\mathfrak S\mathfrak P=1.
 \]
 The sign at \(z_2\) is fixed by the simultaneous parabolic normalization;
 the quadratic ramification calculation alone gives only
 \((g_2'(z_2))^2=-1\).  The resulting nondegenerate period family descends,
 and the middle quotient is a proper holomorphic submersion with compact torus
 fibres.  The toric cusp filling, both finite logarithmic fillings, all three
 punctured-overlap biholomorphisms, and the global analytic gluing yield a
 connected compact Hausdorff smooth complex threefold and a proper surjective
 holomorphic map \(f:X\to\mathbb P^1\) with connected fibres.  Its full cusp
 fibre \(W\) is reduced, irreducible, nonnormal, and normal crossings, with
 normalization \(\mathrm{dP}_6\); the other two full scheme-theoretic fibres
 have multiplicities three and four.  The workbench also gives complete proofs
 of integral nearby-cycle specialization, the global Leray page and three
 primitive transgressions, the equivariant toric collapse, its integral
 cellular map and its equality with specialization, the based van Kampen
 presentation, sphere recognition, and the very-general-fibre N\'eron--Severi
 and meromorphic-descent calculation.  For the constructed threefold it
 therefore verifies
\[
X\cong_{\mathrm{diff}}S^6,\qquad
\langle c_3(T_X),[X]\rangle=2,\qquad
\mathcal M(X)=f^*\mathbb C(t),\qquad a(X)=1.
\]
The same verified object satisfies the stated hypotheses of
Campana--Demailly--Peternell Theorem~2.2 and contradicts its conclusions, so
the workbench records a counterexample to that theorem as stated as well as
two exact defects in its printed dependency chain.  It also proves the
strongest corrected conclusion that survives those defects: existence of one
\(L_0\) with \(H^2(X,T_X\otimes L_0)=0\) implies generic
off-degree-one vanishing and \(c_3(X)\le0\); for reduced full fibres the
invalid singular-fibre shortcut is replaced by injectivity of the exact
conductor connecting map.  Because \(c_3(X)=2\), contraposition gives
\(H^2(X,T_X\otimes L_0)\ne0\) for every line bundle \(L_0\); independently,
the conductor calculation gives \(R^2f_*(T_X\otimes L)\ne0\) for every
\(L\), so the stronger relative sufficient condition also fails.  The
independent alternate
certificate is now complete as well: the actual finite-band maps are
$\alpha_q=(\pi_{1*},-\pi_{2*})$ in primitive integral bases, the cusp
attachment reduces to blocks with determinant $\pm p$, and the resulting
obstruction squares commute with all three Leray transgressions.  The
candidate's Proposition~10.8 is also verified: the
structure-sheaf direct images, degree-one base change, Picard line, exact
finite normal characters, and uniform period/countability argument are
reconstructed in full without using the alternate collapse/Mayer--Vietoris
route.
Independently of the global topology calculation, the conductor construction
also proves the required direct-image nonvanishing for every
$L\in\operatorname{Pic}^0(X)$.  On the cusp model the invariant form
\[
 \Omega=\frac{dx_1\wedge dx_2\wedge dt}{x_1x_2}
\]
first proves intrinsically that
\(\omega_{N_0}\simeq\mathcal O_{N_0}\).  After global gluing, the open
embedding \(N_0\hookrightarrow X\) gives
\[
 \omega_X|_{N_0}\simeq\omega_{N_0}\simeq\mathcal O_{N_0}.
\]
Write $\nu:\widetilde W\simeq\mathrm{dP}_6\to W$ for the normalization.
For arbitrary line bundles the exact
obstruction is the restriction homomorphism
$\nu^*i^*:\operatorname{Pic}(X)\to\operatorname{Pic}(\mathrm{dP}_6)$.  A direct
Leray-edge calculation gives
$e(c_1(L))\in E_\infty^{0,2}\subseteq\ker d_2^{0,2}=0$; proper-base-change
naturality then makes $\nu^*i^*L$ trivial for every $L$.  This avoids
invoking the packaged theorem $H^2(X;\mathbb Z)=0$ but reuses its decisive
integral Leray transgression.  A second all-Pic proof now avoids every global
Leray differential.  Put $q:=uw+6\gamma\delta\in\bigwedge^2V=H^2(F;\mathbb Z)$.
If $E=c_1(L|_F)$, the order-four filling forces
$E=2nq$.  The pair-of-pants defect classes satisfy
\[
 \bar d_1+\bar d_2+\bar d_0=0,\qquad
 \bar d_1=4n[\delta],\quad \bar d_2=-3n[\delta],\quad \bar d_0=0,
\]
where the first two values come from the exact logarithmic-cover norm maps
and the last from a literal cusp solid torus.  Hence $n=0$; local integral
specialization then gives $\nu^*i^*L\simeq\mathcal O$.
The cusp-local map is nevertheless nonzero: a descended line bundle on the
cusp neighbourhood has generic Chern class $-2q$ and normalized central
degree vector $(-12,2,0,-12,2,0)$.  The third global certificate therefore
incorporates both finite fillings analytically.  For $c_1(L|_F)=kq$, its
$\mathcal O^*$-factor defects satisfy
\[
 (\bar d_1,\bar d_2,\bar d_0)=(2k,-\tfrac32k,0)[\delta],
 \qquad \bar d_1+\bar d_2+\bar d_0=0,
\]
The finite norm identity is computed in the typed affine deck group
\[
 \Gamma_j^{\log}
 =\langle\tau_\lambda,Q_{j,\log}:\lambda\in\Lambda\rangle
\]
Its extended defect satisfies
\[
 \widetilde d_j(Q_{j,\log})=d_j+(T_j-I)a_j,\qquad
 \widetilde d_j(\tau_v)=\iota_vE,
\]
and therefore
\[
 N_{T_j}d_j
 =\widetilde d_j((Q_{j,\log})^{m_j})
 =\widetilde d_j(\tau_{v_j})=\iota_{v_j}E.
\]
Thus the displayed balance gives \(k=0\); the cusp equality comes from the
argument principle in a Stein chart.  The comparison uses the actual
toric-open direction
 $U_{\rho_v}=D(z_wz_u)\subset U_\sigma$ and a contractible Stein chart
 $S_\sigma$, so adjacent central frames have zero horizontal winding.  If
 $\mu^+_\lambda$ is multiplication by $\chi^{m_\lambda}$ on
 $\mathcal O(D_a)$, the required map on the negative divisor is the typed
 inverse-dual
 \[
  \mu^-_\lambda=((\mu^+_\lambda)^{-1})^\vee:
  \Psi_\lambda^*\mathcal O(-D_a)\longrightarrow\mathcal O(-D_a),
  \qquad s\longmapsto\chi^{-m_\lambda}s.
 \]
Put \(M=L|_{N_0}\).  If $(d_1,d_2,d_3)$ are the independent normalized
central degrees, then
\[
 B_0=\begin{pmatrix}0&1\\-1&0\end{pmatrix},\qquad
 B=\begin{pmatrix}d_2+d_3&d_3\\d_3&d_1+d_3\end{pmatrix},
\]
\[
 E_M((\lambda,0),(0,m))=-B(B_0\lambda,m).
\]
The preceding balance gives \(k=0\), hence \(E_M=c_1(L|_F)=0\).
Unimodularity of $B_0$ now forces $B=0$ and all six degrees to vanish.  This
third proof avoids both global integral Leray and local integral
specialization/bounding-chain input.
The workbench also completes the Hodge calculation left open in the
manuscript.  The torsion-free absolute cotangent direct images are
\[
R^0f_*\Omega_X^1\simeq\mathcal O(-2),\quad
R^1f_*\Omega_X^1\simeq\mathcal O\oplus\mathcal O(-1)^{\oplus2},\quad
R^2f_*\Omega_X^1\simeq\mathcal O\oplus\mathcal O(-1).
\]
Hence $h^{1,2}=h^{2,1}=1$, $h^{1,1}=h^{2,2}=2$, and every Hodge number is
known.  The Fr\"olicher spectral sequence does not degenerate at $E_1$: its
two middle $d_1=\partial$ rows are exact
\[
0\longrightarrow\mathbb C\longrightarrow\mathbb C^2
\longrightarrow\mathbb C\longrightarrow0.
\]
Thus $\dim_{\mathbb C}E_1=10$, while
$E_2=E_\infty$ and $\dim_{\mathbb C}E_2=2$.
Put
\[
 \Sigma=p_0+p_1+p_2,\quad D=f^{-1}(\Sigma)_{\mathrm{red}},\quad
 \mathcal L=\Omega^1_{\mathbb P^1}(\log\Sigma),\quad
 \mathcal E=\Omega_X^1(\log D)/f^*\mathcal L,
 \quad \mathcal G=R^1f_*\mathcal O_X.
\]
Then
\[
 f_*\mathcal E\simeq\mathcal O\oplus\mathcal O(-1),\quad
 R^1f_*\mathcal E\simeq\mathcal O^{\oplus2}\oplus
 \mathcal O(-1)^{\oplus2},\quad
 R^2f_*\mathcal E\simeq\mathcal O\oplus\mathcal O(-1),
\]
and all three sheaves are locally free and commute with base change.  The
global calculation uses the actual cyclic invariant lattices, the simple divisor
  $p_1+p_2$ of the logarithmic Kodaira--Spencer determinant, and the normalized
  global scalar at the order-four point.  At a finite fibre write \(t=s^m\),
  with \((m,a)=(3,2)\) or \((4,1)\), and, after removing the common \(ds\)
  from \(D_{\mathrm{KS}}=(d_{rs})\), put
 \[
 A=\frac{s}{m}d_{00},\quad B=\frac{s^{1-a}}m d_{01},\quad
 C=\frac{s^{1-a}}m d_{10},\quad D_0=\frac{s^{m+1-2a}}m d_{11}.
 \]
  In the invariant source basis \((e_0,s^{m-a}e_1)\), target basis
  \((b_0,s^ab_1)\ell\), and ordered degree-one source basis
  \(e_0b_0,s^ae_0b_1,s^{m-a}e_1b_0,e_1b_1\),
 \[
  [\kappa_0]=\begin{pmatrix}A&tB\\ C&D_0\end{pmatrix},\qquad
  [\kappa_1]=(-C,A,-D_0,B).
 \]
 The exact finite residue cancellation is
 \[
  \operatorname{res}_{1,j}(xb_0+ys^ab_1)=m_jx[b_0],\qquad
  \operatorname{res}_{1,j}\kappa_0=0,
 \]
 \[
  \ker\operatorname{res}_{1,j}=\mathbb C\,s^ab_1
  =\operatorname{im}(\kappa_0\bmod t).
 \]
 At the cusp put
 \[
  Q_{\mathrm{cusp}}=\nu_*\mathcal O_{\mathrm{dP}_6}/\mathcal O_W,\qquad
  \eta_i=d\log x_i,\qquad q_1=(-1,0,1),\qquad q_2=(0,-1,-1),
 \]
 and, for the normalization connecting maps, set
 \[
  \alpha_i=\delta_0(q_i),\qquad
  \theta=-\delta_1^{-1}(\alpha_1\smile\alpha_2).
 \]
 The standard \v Cech sign gives
 \[
  \delta_1\bigl(\rho_*(\alpha\otimes\eta_i)\bigr)
  =-\alpha\smile\alpha_i,
 \]
 hence the exact row \((0,1,-1,0)\) with target \(\mathbb C\theta\).
 Let
 \[
  \xi=\frac{\partial}{\partial\zeta_2},\qquad
  h=(a(z),1)\binom{d\zeta_1}{d\zeta_2},\qquad h(\xi)=1
 \]
 be the unique global lift.  At the order-four point, in the genuine
 invariant frame
 \(\omega_2=s\,u_2(s)(\gamma\wedge u)\ell\), \(u_2(0)=1\), one has
 \(a(z_2)=\tfrac12\), \(\dot\tau(z_2)=0\),
 \(K(z_2)=6\dot\mu(z_2)^2\ne0\), and
 \[
  \kappa_1(h\otimes\gamma)(p_2)
  =-\frac14\bigl(a(z_2)\dot\tau(z_2)+\dot\mu(z_2)\bigr)
  =-\tfrac14\dot\mu(z_2)\ne0.
 \]
 The stalks at \(t=xy\) and \(t=xyz\) show that
 \(\Omega^1_{X/B}\) is torsion-free but non-locally-free at the cusp.  On a
 finite cover \(t=s^{m_j}\), its exact local module is
 \[
  R\,dy_1\oplus R\,dy_2\oplus
  \bigl(R/(s^{m_j-1})\bigr)ds,
 \]
 so its torsion is supported exactly at the two multiple fibres.  No step
 assumes that direct image commutes with duality across a multiple fibre.
McHugh's compact-threefold identities then determine the remaining
Bott--Chern and Aeppli invariants.  In total degrees $0$ through $6$,
\[
(h^k_{BC})=(1,0,2,2,2,2,1),\qquad
(h^k_A)=(1,2,2,2,2,0,1),
\]
and the non-$\partial\bar\partial$ degrees are
$(0,2,4,4,4,2,0)$.  Thus no Bott--Chern or Aeppli dimension remains open.
On a marked smooth fibre,
\[
 \eta=u\wedge w+6\gamma\wedge\delta,\qquad
 \eta^2=12\,\gamma\wedge u\wedge w\wedge\delta.
\]
If \(Q_F\) is its displayed Dolbeault representative, then
\[
 \int_{(F_b,\mathrm{or}_{\rm mark})}Q_F^2=12,\qquad
 \int_{(F_b,\mathrm{or}_{\mathbb C})}Q_F^2=-12;
\]
the marked and complex orientations are opposite because
\(\det_{\mathbb R}\Pi<0\).

Let \(\mathfrak g\in H_{\bar\partial}^{0,1}(X)\) be the distinguished
degree-one class and let \(\mathfrak q\in H_{\bar\partial}^{1,1}(X)\) be a
chosen coherent-Leray lift of \(\eta\).  The latter is determined only modulo
\(\mathbb C\,d_1\mathfrak g\).  Choose global \(\bar\partial\)-closed smooth
representatives
\[
 G\in\mathcal A^{0,1}(X),\qquad Q\in\mathcal A^{1,1}(X)
\]
of \(\mathfrak g,\mathfrak q\).  Their marked smooth-fibre cohomology classes
have the exact period-coordinate formulas proved below; no literal equality
of an arbitrary global \(Q\) with the displayed coefficient formula is used.
With \(\omega=c_1(\mathcal O_{\mathbb P^1}(1))\),
\[
 \exp(12\mathfrak g)=f^*\mathcal O_{\mathbb P^1}(1),\qquad
 c_{1,BC}(\exp(z\mathfrak g))=z\mathsf A+\bar z\,\mathsf B,
\]
\[
 c_{1,BC}(f^*\mathcal O_{\mathbb P^1}(1))
 =12(\mathsf A+\mathsf B).
\]
Put \(\mathbf1=[1]_{BC}\), and choose a smooth top form \(\vartheta\) with
\(\int_X\vartheta=1\), so \(\mathsf T=[\vartheta]_{BC}\).  Then
\[
 \mathsf A=[\partial G],\ \mathsf B=[\bar\partial\bar G],\
 \mathsf C=[\partial Q],\ \overline{\mathsf C},\
 \mathsf D=[\partial(G\wedge Q)],\ \overline{\mathsf D},\
 \mathsf E=[\partial(Q^2)],\ \overline{\mathsf E},\ \mathsf T
\]
are bases of all nonzero positive-degree Bott--Chern groups; together with
\(\mathbf1\), they give every nonzero group.  The two load-bearing fibre
pairings in the complex orientation are
 \[
  \int_{F_b}Q_F^2=-12,\qquad
  \int_{F_b}Q_F\wedge\Gamma\wedge d\zeta_2=-1.
 \]
They prove the required Dolbeault nonvanishing, while the constant fibre
coefficients give the typed identity
\[
 i_b^*\!\left(d_1[G\wedge Q]_{\bar\partial}\right)
 =[\partial_F(\Gamma\wedge Q_F)]_{\bar\partial}=0.
\]
If \(N_{\mathrm{BC}}\) is the span of the nine displayed positive-degree
classes, then
\[
 H^{*,*}_{BC}(X)=\mathbb C\mathbf1\oplus N_{\mathrm{BC}},\qquad
 N_{\mathrm{BC}}^2=0.
\]
Thus every product of two positive-degree Bott--Chern classes vanishes; see
\cref{prop:bcpr-generators-are-bases,thm:bcpr-square-zero-algebra}.  The
 result does not select a preferred literal global coefficient atlas for
 \(G,Q\) through the finite fillings and cusp; that presentation choice is not
 an open mathematical premise.
 The marked positive radial cusp also carries an exact auxiliary flat-metric
 calculation.  With
 \[
  L_T=T+\operatorname{Im}h_T,\qquad
  q_T=L_T-\operatorname{Im}b_T,\qquad
  m_T=\operatorname{Im}\mu_T,\qquad
  D_T=L_Tq_T+6m_T^2,
 \]
 where
 \(H_{\mathrm{cusp}}=\operatorname{Im}h(0)\) and
 \(B_{\mathrm{cusp}}=\operatorname{Im}b(0)\),
 the complete dual lattice is \(P_T^{-\top}\mathbb Z^4\), and its
 shortest nonzero vector gives
 \[
  \lambda_1(F_T,g_T)
  =\frac{4\pi^2}{T^2}
    \left(1+\frac{-2H_{\mathrm{cusp}}
    +2\min\{B_{\mathrm{cusp}},0\}}T+O(T^{-2})\right),
 \]
 while the unit-volume metric
 \(\widehat g_T=D_T^{-1/2}g_T\) gives
 \[
 \lambda_1(F_T,\widehat g_T)
  =\frac{4\pi^2}{T}\bigl(1+O(T^{-1})\bigr)\longrightarrow0.
\]
 For fixed \(E\geq0\), the scalar eigenvalue count with Fourier
 multiplicity is
 \[
  N_T(E):=\#\{k\in\mathbb Z^4:
  4\pi^2\|P_T^{-\top}k\|^2\leq E\}
  =\pi D_T\sum_{n\in\mathbb Z^2}
   \left(\frac{E}{4\pi^2}-\|n\|^2\right)_+ +O_E(T).
 \]
 In particular, for \(0<E<4\pi^2\),
 \(N_T(E)=ED_T/(4\pi)+O_E(T)\): the number of soft modes grows
 quadratically, rather than only one eigenvalue tending to zero.
 The mean-zero vertical direct integral is a single injective self-adjoint
 operator with \(0\) in its essential spectrum.  The unmodified regular
 four-torus fibres \(F_T\) have native pointed limit exactly
 \(\mathbb T^2\times\mathbb R^2\).  A separate ancillary construction changes
 both the fibre lattice and the cusp base: for an integer \(N\geq1\), after
 the added operations
 \[
  t_c=v^N,\qquad
  \Lambda^{(N)}
  =\mathbb Z\langle\widehat\gamma,\widehat u,
     N\widehat w,N\widehat\delta\rangle,
 \]
 the punctured degree-\(N^2\) fibre covers have pointed local limit
 \(\mathbb R^4\) and momentum lattices becoming dense in all four directions.
 The extension across the central fibre is explicit.  With source residual
 term \(C_N(v)=C(v^N)/N\), put
 \[
  f_N:N_0^{\langle N\rangle}\longrightarrow\Delta_v,\qquad
  W_N=(f_N^{-1}(0))_{\mathrm{red}}.
 \]
 Toric multiplication descends to
 \[
  \kappa_N:N_0^{\langle N\rangle}\longrightarrow N_0,
  \qquad \deg\kappa_N=N^3,\qquad
  \kappa_N^*W=NW_N,\qquad
  \operatorname{Ram}(\kappa_N)=(N-1)W_N.
 \]
 Factoring through the normal base change
 \(N_0[N]=N_0\times_{\Delta_t}\Delta_v\) gives a generic degree-\(N^2\)
 map \(Q_N\); for \(N>1\) it is quasi-\'etale but nonflat exactly over the
 double and triple strata.  On the normalized central fibre it is the toric
 map \([N]:\mathrm{dP}_6\to\mathrm{dP}_6\).  This is a completed
 \emph{local} cusp morphism, not a cover of the entire base or of \(X\).

 Independently, the compact cusp lattice
 \(K_\Lambda=\mathbb Z\langle\widehat w,\widehat\delta\rangle\) is not
 preserved by the finite monodromies, but its invariant closure is
 \(L_\Lambda=\mathbb Z\langle\widehat u,\widehat w,\widehat\delta\rangle\).
 The invariant character lattice is \(\mathbb Z\varepsilon_\gamma\); fix the
 displayed primitive generator.  Its real group quotient has exact circle
 circumference
 \(\ell_T=D_T/\sqrt{L_T^2+m_T^2}
 =T+H_{\mathrm{cusp}}-B_{\mathrm{cusp}}+O(T^{-1})\).  Four synchronized copies
 therefore give a genuine \(\mathbb T^4\)-quotient over the full regular base
 whose pointed cusp limit is \(\mathbb R^4\), and four copies are minimal for
 this invariant-kernel mechanism.  The two finite fillings extend to explicit
 Seifert targets after simultaneous-root normalization, with central degrees
 \(27\) and \(64\).  At the cusp, two paths with the same limit and different
 character values still prove that the \emph{literal fixed-target} quotient
 cannot extend; no finite base cover repairs the native two-copy kernel.

 A different, exact compactification changes the codomain rather than the
 character values.  After choosing a base point, collapsing the inner and
 outer cusp cross-sections gives a based quotient (denoted \(q_2\) in the
 detailed construction, and written \(\rho_{2,\mathrm{cone}}\) here to avoid collision with
 an earlier \v Cech vector)
 \[
  \rho_{2,\mathrm{cone}}:(X,x_0)\longrightarrow
  (\mathcal A_{\mathrm{cone}},*)
  =\bigl(\operatorname{Susp}(S^1\times\overline{\mathbb T}^{\,2}),*\bigr).
 \]
 Its Cartesian square restricts on synchronized radial fibres to
 \(\pi_T^{\times2}:F_T^{\times2}\to Q_T^{\times2}\), and basedness induces
 \(\rho_{2,\mathrm{cone}}^{\wedge2}:X^{\wedge2}\to
 \mathcal A_{\mathrm{cone}}^{\wedge2}\).
 Likewise the regular and finite-filled circle targets glue to a compact
 three-manifold \(\mathcal M_{\mathrm{Seif}}\) with boundary
 \(\mathbb T^2\); coning that boundary gives a based quotient (denoted
 \(q_1\) in the detailed construction)
 \(\rho_{1,\mathrm{cone}}:(X,x_0)\to(\mathcal C_{\mathrm{cone}},*)\) whose fourth Cartesian
 power restricts to
 \(\chi_{\gamma,T}^{\times4}:F_T^{\times4}\to C_T^{\times4}\), and hence
 induces a fourth smash power.

 Separately, composing \(\rho_{2,\mathrm{cone}}\) and the cusp-local circle quotient
 \(q_{1,\mathrm{cusp}}\) with their respective base-angle-forgetting maps
 gives continuous surjective quotient maps
 \[
  \widehat\rho_3^{\wedge2}:S^{12}\twoheadrightarrow S^6,
  \qquad
  \widehat\rho_\gamma^{\wedge4}:S^{24}\twoheadrightarrow S^8,
 \]
 after a based choice \(X\cong_{\mathrm{diff}}S^6\).  Both full sphere
 carriers are based null-homotopic, by two distinct calculations.  In the
 detailed notation \(\widehat\rho_\gamma=\widehat q_\gamma\) and
 \(\widehat\rho_3=\widehat q_3\).  For any based map
 \(f:S^6\to S^2\), the smash square has the exact factorization
 \[
  S^{12}=S^6\wedge S^6
  \xrightarrow{\ f\wedge\operatorname{id}_{S^6}=\Sigma^6f\ }
  S^2\wedge S^6=S^8
  \xrightarrow{\ \operatorname{id}_{S^2}\wedge f\ }
  S^2\wedge S^2=S^4.
 \]
 Freudenthal stability and the vanishing fourth stable stem give
 \(\pi_{12}(S^8)\cong\pi_4^{\mathrm S}=0\), so \(f\wedge f\) is based
 null-homotopic.  Associativity then writes
 \(f^{\wedge r}=(f\wedge f)\wedge f^{\wedge(r-2)}\), and hence
 \([f^{\wedge r}]=0\in\pi_{6r}(S^{2r})\) for every \(r\ge2\).  In
 particular, the entire quotient map
 \(\widehat\rho_\gamma^{\wedge4}:S^{24}\to S^8\), not merely its moving
 four-sphere restriction, is based null-homotopic.

 For the other carrier, let \(g:S^6\to S^3\) be any based map.  The exact
 smash factorization is
 \[
  g\wedge g=(\Sigma^3g)\circ(\Sigma^6g):
  S^{12}\longrightarrow S^9\longrightarrow S^6.
 \]
 Write \(\operatorname{st}\) for stabilization.  The exact low-stem data are
 \[
  \pi_6(S^3)\cong\mathbb Z/12,
  \qquad
  \pi_3^{\mathrm S}=\mathbb Z/24\langle\nu\rangle,
  \qquad
  \pi_6^{\mathrm S}=\mathbb Z/2\langle\nu^2\rangle,
 \]
 \[
  \pi_{12}(S^6)
  =\mathbb Z/2\langle\nu_6\circ\nu_9\rangle,
  \qquad
  \operatorname{st}:\pi_{12}(S^6)
  \xrightarrow{\ \cong\ }\pi_6^{\mathrm S}.
 \]
 The full groups are tabulated in \cite[p.~339]{Hatcher2002}; the signed
 smash formula and the unstable stabilization calculation are
 \cite[Propositions~3.1 and~5.11]{Toda1962}.
 Since \(12[g]=0\), its stable image has the form
 \(\operatorname{st}([g])=2a\nu\) for some \(a\in\mathbb Z/12\).  Stable
 smash multiplication therefore gives
 \[
  \operatorname{st}([g\wedge g])
  =\operatorname{st}([g])^2
  =4a^2\nu^2=0\in\pi_6^{\mathrm S}.
 \]
 The displayed stabilization isomorphism now forces
 \([g\wedge g]=0\in\pi_{12}(S^6)\).
 Hence every \(g^{\wedge r}\) with \(r\ge2\) is based null-homotopic as
 well.  In particular,
 \[
  [\widehat\rho_3^{\wedge2}]
  =[\widehat q_3^{\wedge2}]=0\in\pi_{12}(S^6),
  \qquad
  [\widehat\rho_\gamma^{\wedge4}]
  =[\widehat q_\gamma^{\wedge4}]=0\in\pi_{24}(S^8).
 \]
 The one-copy class \([\widehat q_3]\in\pi_6(S^3)\) is zero by
 the explicit framed filling in Theorem~\ref{h24:thm:cusp-g-zero}.
 The one-copy class \([\widehat q_\gamma]\in\pi_6(S^2)\) remains open. The unstable
 classes of the conical-target maps
 \(\rho_{2,\mathrm{cone}}^{\wedge2}\),
 \(q_{1,\mathrm{cusp}}^{\wedge4}\), and
 \(\rho_{1,\mathrm{cone}}^{\wedge4}\) also remain \Open; nullity after
 postcomposition to a sphere does not prove nullity before that
 postcomposition.  See
 \cref{cusp-cone-S12-sphere-map-null,cusp-cone-S24-sphere-map-null,%
 cusp-cone-sphere-map-homotopy-status,%
 cusp-cone-target-homotopy-types}.

 Their targets contain
 moving topologically embedded four-spheres, equipped separately with their
 intrinsic quotient length metrics,
 \(S_T^2\wedge S_T^2\) and \(C_T^{\wedge4}\).  With their quotient length
 metrics and midpoint base points, every radius-\(R\) ball is exactly
 Euclidean once \(T-C_{\mathrm{sys}}>4R\), respectively once
 \(\ell_T>4R\); thus both moving four-spheres converge pointed
 local-isometrically and Gromov--Hausdorff to \(\mathbb R^4\), with smooth
 convergence on those Euclidean neighborhoods.  Deleting the common smash
 basepoint gives literal nonclosed topological four-cells
 \[
  (Q_T\setminus K_T^{(1)})^2\cong\mathbb R^4\hookrightarrow S^6,
  \qquad
  (C_T\setminus\{0\})^4\cong\mathbb R^4\hookrightarrow S^8.
 \]
 The based inclusion maps
 \[
  S_T^2\wedge S_T^2\hookrightarrow S^6,
  \qquad
  C_T^{\wedge4}\hookrightarrow S^8
 \]
 are null-homotopic.  These slice-inclusion null-homotopies are separate from
 the global equalities
 \([\widehat\rho_3^{\wedge2}]=0\in\pi_{12}(S^6)\) and
 \([\widehat\rho_\gamma^{\wedge4}]=0\in\pi_{24}(S^8)\).  Two copies are
 rank-minimal for the
 native cusp-local quotient, while four are rank-minimal for a primitive
 full-monodromy-invariant quotient with four expanding directions.  Here
 \(C_{\mathrm{sys}}>0\) is a fixed bound for the holomorphic error in the
 radial period lattice.

Exact five-manifolds with boundary interpolate the positive four-torus
slices with a boundary fibre \(\mathbb R^4\).  There are continuous
surjective comparison maps to the compact cone carriers which are
homeomorphisms away from that boundary and collapse it to the cone vertex;
however, these maps are nonproper and nonquotient.  Thus they are not
topological blow-downs or resolutions, and the compact carriers are not the
quotients of the five-manifolds by their \(\mathbb R^4\) boundary fibres via
these maps.  For every unrescaled fixed compatible
target metric the radial slices collapse to the cone point.  A fixed
 cusp-adapted metric can recover the slice limits only after the blow-up
 factors \(T/r_T\) and \(\ell_T/r_T\); the corresponding full ambient
 blow-ups are \(\mathbb R^8\) and \(\mathbb R^{12}\), not \(\mathbb R^4\).
 Thus the sphere constructions yield exact carrier slices with intrinsic
 \(\mathbb R^4\) limits and a decompactification boundary, but not
 \(\mathbb R^4\) as the whole
 compact sphere or compact target.  Moreover
 \(X^{\times4}\cong(S^6)^{\times4}\not\simeq S^{24}\), and
 Borel--Serre still excludes an almost-complex structure on either
 \(S^{12}\) or \(S^{24}\) \cite{BorelSerre1953}.
 The distinct first-coordinate holomorphic quotient is now classified exactly:
 \(\mathcal A_z\to\mathbb C/(\mathbb Z\tau+\mathbb Z)\) exists precisely when
 \(6\mu=a\tau+b\) for \((a,b)\in\mathbb Z^2\), has connected kernel
 \(\mathbb C/(\mathbb Z+\mathbb Z(a\mu-\beta))\), and its marked labels obey
 \(g_0(a,b)=(a,a+b)\), \(g_2(a,b)=(-b,6+a)\).  This locus has finite image in
 the completed orbifold base and misses a punctured cusp neighborhood, so it
 cannot be the cusp gap-closing mechanism; the theorem does not exclude other
 elliptic quotients or smooth integral magnetic fluxes.
 The reduced classical Yang--Mills Hessian has the same scalar factor
 \(g_{\mathrm{YM}}^{-2}\), and the finite cubical model is diagonalized in
 every mesh regime.  On the degree-\(M_{\mathrm{cov}}^2\) cover, an
 independent \(N_{\mathrm{mesh}}^4\)-vertex cubical subdivision gives a
 correctly typed holonomy map to \(SU(2)^{E_{N_{\mathrm{mesh}}}}\).  For
 \(H_{\mathrm{link}}=\operatorname{diag}(i,-i)\), the explicit unit-flux
 link field has
 \[
  U_{12}(n)=\exp(2\pi H_{\mathrm{link}}/N_{\mathrm{mesh}}^2),\qquad
  S^{\log}_{T,M_{\mathrm{cov}},N_{\mathrm{mesh}}}
  =\frac{2\pi^2M_{\mathrm{cov}}^2}
         {g_{\mathrm{YM}}^2D_T}.
 \]
 On the reduced finite Coulomb space, the Gaussian oscillator has gap
 \[
  \sqrt{\min_{k\ne0}\frac1{g_{\mathrm{YM}}^2}
   \left\lVert P_T^{-\top}
    S_{M_{\mathrm{cov}}}^{-1}
    \mathfrak d_{N_{\mathrm{mesh}}}(k)\right\rVert^2}.
 \]
 The joint sequence
 \(T_j^{\mathrm{rad}}=j^2,\ M_{\mathrm{cov},j}=j,\
 N_{\mathrm{mesh},j}=j^3\)
 simultaneously gives the pointed cover limit \(\mathbb R^4\), mesh
 diameter \(O(j^{-1})\), Gaussian gap \(O(j^{-2})\), and logarithmic action
 \(O(j^{-2})\).  This oscillator lives on a linear Coulomb space with
 Lebesgue measure; it is not the compact-link Haar/Gauss-law Hilbert space.
 For a nondegenerate integral magnetic form
 \(\mathsf E_{\mathrm{mag}}\), the exact oscillator frequencies satisfy
 \[
  b_1b_2=\frac{4\pi^2|\operatorname{Pf}(\mathsf E_{\mathrm{mag}})|}{D_T},
  \qquad
  \operatorname{gap}(\Delta_{\mathsf E_{\mathrm{mag}},T})
  =2\min\{b_1,b_2\}
  \leq4\pi\sqrt{\frac{|\operatorname{Pf}(\mathsf E_{\mathrm{mag}})|}{D_T}}
  =O(T^{-1}),
 \]
 for the Friedrichs magnetic Laplacian on the corresponding Hermitian line
 bundle over the marked four-torus.  Independently, the explicit reducible
 \(SU(2)\) critical point has
 \[
  \operatorname{YM}_T=\frac{2\pi^2}{g_{\mathrm{YM}}^2D_T},
  \qquad
  \mathcal J_Ta=-\frac{4\pi}{D_T}a.
 \]
 These are geometric and classical operator theorems.  No positive-metric
 physical Yang--Mills Hilbert space, reconstructed Hamiltonian, or
 scale-controlled state-and-energy maps have been constructed, so the
 interacting four-dimensional mass-gap implication is \Open.

 The same native full-flux construction extends exactly to every nonzero
 trace-zero integral weight vector
 \(\mathbf m=(m_1,\ldots,m_r)\) in a split Cartan of \(SU(r)\).  Writing
 \(s_2(\mathbf m)=\sum_i m_i^2\), the resulting reducible bundle and
 connection satisfy
 \[
  c_2(\mathscr E_{\mathbf m,T})
   =-\frac{s_2(\mathbf m)}2[\eta_T]^2,
  \qquad
  \left\langle c_2,[F_T]_{x\text{-}\mathrm{complex}}\right\rangle
   =6s_2(\mathbf m),
 \]
 and
 \[
  \operatorname{YM}_T(\nabla^{\mathbf m,T})
  =\frac{12\pi^2s_2(\mathbf m)}{g_{\mathrm{YM}}^2}
   +\frac{\pi^2s_2(\mathbf m)S_T^2}{g_{\mathrm{YM}}^2D_T}.
 \]
 Its charged Coulomb--Jacobi bands, coincident-band multiplicities, wall
 nullities, and Morse index are computed root by root; on the deep cusp tail
 the real index is \(36r s_2(\mathbf m)\).  These are fibrewise classical
 Yang--Mills theorems.  The nonzero \(c_2\) also proves that the bundles do
 not extend over any proposed carrier with \(H^4=0\), including spheres of
 dimensions (6), (12), and (24); changing the gauge group therefore
 does not create a four-dimensional interacting quantum theory.

 Finally, for \(r=(r_1,r_2,r_3)\in\mathbb R^3\), put
 \[
  Q_K(r)=\frac23(r_1+r_2+r_3)^2-(r_1^2+r_2^2+r_3^2)
        =h^2-x^2-y^2,
 \]
 where
 \[
  h=\frac{r_1+r_2+r_3}{\sqrt3},\qquad
  x=\frac{r_1-r_3}{\sqrt2},\qquad
  y=\frac{r_1-2r_2+r_3}{\sqrt6}.
 \]
 The real-linear coefficient isomorphism
 \[
  (h,x,y)\longmapsto(t,m,\varrho)
  =\left(h+y,\frac{x}{\sqrt6},h-y\right)
 \]
 identifies the Koide cone \(Q_K=0\) with
 \(t\varrho-6m^2=0\) inside the three-dimensional imaginary-coefficient
 space of the canonical period ansatz.  It does not identify that cone with
 the full real period discriminant, whose three real coefficients remain
 free.  For every choice of those real parts,
 \(\det P=t\varrho-6m^2=Q_K(r)\).  Hence each Koide-cone point gives
 \(\det P=0\), whereas the actual positive radial cusp has
 \(\det P_T=-D_T<0\) and \(\det P_T\sim-T^2\); the Koide cone does not
 parametrize regular fibres.

 Independently, write
 \[
  \Delta(x,y,z)=4(xy+xz+yz)-(x^2+y^2+z^2)
 \]
 and, for \(Y\geq0\), define
 \[
  R_{\mathrm K}(Y)=\#\{(x,y,z)\in\mathbb Z_{>0}^3:
  x\leq y\leq z\leq Y,\ \Delta(x,y,z)=-1,\ z\geq4(x+y)\}.
 \]
 The integral coordinates
 \[
  n=z-4(x+y),\qquad q=y-x,\qquad a=x+y-4n
 \]
 give a bijection onto
 \[
  \{(a,q,n)\in\mathbb Z^3:a>0,\ q,n\geq0,\ a+4n>q,
    \ a^2-3q^2-18n^2=-2\},
 \]
 with height \(z=4a+17n\).  Calibrated homogeneous counting gives
 unconditionally
 \[
  R_{\mathrm K}(Y)\sim C_0Y,\qquad
  C_0=\frac9{2\pi^2}\log(2+\sqrt3)
      \log\!\left(\frac3{2\sqrt2}\right).
 \]
 The stronger estimate
 \(R_{\mathrm K}(Y)-C_0Y=O_\delta(Y^{1/2+\delta})\) for every
 \(\delta>0\) remains \Open.

 There is also an exact arithmetic bridge to the earlier
 Erd\H{o}s--Straus divisor-shell construction.  Fix a prime
 \(p\equiv1\pmod{12}\) and \(\varepsilon\in\{0,1\}\), and put
 \[
  L=\frac{p-1}{4},\qquad 0\leq j<L,\qquad R_j=4j+3,
  \qquad a_j^{\mathrm{ES}}=\frac{p+R_j}{4}.
 \]
 For \(u\in\mathbb Z_{>0}\) satisfying
 \(u\mid(a_j^{\mathrm{ES}})^2\) and
 \(R_j\mid(4u+p^\varepsilon)\), set
 \[
  T_{\mathrm{sel}}=\frac{4u+p^\varepsilon}{R_j},\qquad
  \ell_0=\frac{T_{\mathrm{sel}}+1}{4},\qquad
  \ell_1=u-jT_{\mathrm{sel}},\qquad \ell_2=1.
 \]
 Then \(T_{\mathrm{sel}}\in\mathbb Z_{>0}\),
 \(12\ell_0-4\ell_1-3=p^\varepsilon\), and
 \[
  \begin{pmatrix}1&-j\\0&1\end{pmatrix}
  \begin{pmatrix}R_j&-u\\4&-T_{\mathrm{sel}}\end{pmatrix}
  =\begin{pmatrix}3&-\ell_1\\4&1-4\ell_0\end{pmatrix}.
 \]
 For fixed \((p,\varepsilon)\), the marked assignment \((j,u)\mapsto
 (j,\ell_0,\ell_1)\) is a bijection onto the triples satisfying
 \[
  0\leq j<L,\quad12\ell_0-4\ell_1-3=p^\varepsilon,\quad
  T_{\mathrm{sel}}=4\ell_0-1>0,\quad
  u=\ell_1+jT_{\mathrm{sel}}>0,\quad
  u\mid(a_j^{\mathrm{ES}})^2.
 \]
 This is an integral-relator equivalence, not a geometric map between the
 reciprocal cubic and \(X\).
 The distinct residue-lattice and count maps are proved in
 Section~\ref{esn:sec:shared}: the fixed-shell divisor bijection and
 finite-sample Fourier pushforward retain the actual prime and shell,
 rather than replacing the divisor gate by a residue modulo $840$.

 The higher-rung calculation now has several exact, sharply separated outputs.
 For two complex columns, left Gram reduction gives
 \[
  \operatorname{Mat}_{3\times2}(\mathbb C)/SU(3)
  \cong\operatorname{Herm}_2(\mathbb C)_{\geq0},\qquad
  S^{12}_{\rm col}/SU(3)\cong B^4,
 \]
 whose orbit-type boundary is \(S^3\).  Here \(S^{12}_{\rm col}\) is the
 one-point compactification of the column representation; it is not the
 cusp-smash \(S^{12}\) constructed elsewhere in the workbench.
 For \(\mathscr V_3=\operatorname{Mat}_{3\times3}(\mathbb C)\), left Gram
 reduction and the alternating determinant give
 \[
  \mathscr V_3/U(3)\cong\operatorname{Herm}_3(\mathbb C)_{\geq0},
  \qquad
  \mathscr V_3/SU(3)
  \cong\{(H,z):H\geq0,\ |z|^2=\det H\}\cong\mathbb R^{10}.
 \]
 Consequently
 \[
  S^{17}/U(3)\cong B^8,\quad
  \mathscr V_3^+/U(3)\cong B^9,\quad
  S^{17}/SU(3)\cong S^9,\quad
  \mathscr V_3^+/SU(3)\cong S^{10}.
 \]
 The compactified rank-deficient locus is the standard locally flat
 \(S^8\subset S^{10}\), and
 \(S^{10}\setminus S^8\cong\mathbb R^8\times\mathbb C^*\simeq S^1\),
 with generator \(\det A/|\det A|\).  The full rank-three cubic Jordan cone
 is not itself the rank-two quadratic Lorentz cone; the marked angular
 quotient nevertheless supplies a split positive observation onto that cone
 while retaining all five complementary directions.

 For
 \[
  P=\begin{pmatrix}1&-1&0\\0&1&-1\end{pmatrix},\quad
  u=(1,1,1)^{\mathsf T},\quad
  R=P^*(PP^*)^{-1},\quad \Pi=uu^*/3,
 \]
 define
 \[
  \mathcal O_P(H)=PHP^*,\qquad \jmath_P(G)=RGR^*.
 \]
 Then
 \[
  \mathcal O_P\jmath_P=\operatorname{id},\qquad
  \jmath_P\mathcal O_P(H)=(I_3-\Pi)H(I_3-\Pi),
 \]
 and the positive zero fibre is exactly
 \(\{\lambda uu^*: \lambda\geq0\}\).  With \(X=(R\;u)\), every
 Hermitian \(H\) has unique coordinates
 \[
  H=X\begin{pmatrix}G&\psi\\\psi^*&\lambda\end{pmatrix}X^*,\qquad
  \det H=\lambda\det G-\psi^*\operatorname{adj}(G)\psi,
 \]
 where
 \[
  H\succeq0\Longleftrightarrow
  G\succeq0,\quad\psi\in\operatorname{im}G,\quad
  \lambda\geq\psi^*G^+\psi.
 \]
 Thus
 \[
  \operatorname{Herm}_3(\mathbb C)
  \cong\operatorname{Herm}_2(\mathbb C)\oplus\mathbb C^2\oplus\mathbb R,
  \qquad9=4+4+1.
 \]
 The faithful embedding
 \(T_P(g)=RgP+\Pi=X(g\oplus1)X^{-1}\) acts by
 \((G,\psi,\lambda)\mapsto(gGg^*,g\psi,\lambda)\).  The three summands
 are the Lorentz four-vector, the defining complex spin representation, and
 a real scalar; the vector action kills \(\{\pm I_2\}\), whereas the full
 action retains the spin double cover because \(-I_2\) negates \(\psi\).
 On positive cones, \(\jmath_P\) properly compactifies to
 \(B^4\hookrightarrow S^8\), while the reverse boundary observation is
 nonproper.  The determinant-phase equation is
 \[
  |\omega|^2=\lambda\det G-\psi^*\operatorname{adj}(G)\psi,
 \]
 and \(\omega=0\) is the rank-loss locus compactifying to
 \(S^8\subset S^{10}\).  The four-branch chart supplies the exact seed
 \[
  K(z)=\operatorname{diag}(z_0z_3,z_1z_2),\qquad
  \det K=z_0z_1z_2z_3,
 \]
 whose positive \(SL_2(\mathbb C)\)-orbit fills the future causal cone.
 Finally, with retained \(\alpha,c>0\),
 \[
  p_\alpha(Z)=\alpha Z^*Z,\qquad
  m_{\alpha,c}^2(Z)=\frac{\alpha^2}{c^2}
  \sum_{a<b}|Z_{a1}Z_{b2}-Z_{a2}Z_{b1}|^2
 \]
 is the complete \(SU(3)\)-orbit momentum, and
 \(m_{\alpha,c}^2(rZ)=r^4m_{\alpha,c}^2(Z)\).  Hence the unrestricted
 pointwise carrier is gapless.  The complete presymplectic reduction of the
 displayed worldline action is now proved.  For \(m_0>0\), the rank-two
 colour action is free, \(DP_{\alpha,Z}\) is surjective with kernel
 \(\mathfrak{su}(3)\cdot Z\), and the characteristic quotient is
 \[
  \bigl(\mathcal J_2\times\mathcal H^+_{m_0,c}\bigr)/
  \bigl((x,p)\sim(x+sp,p)\bigr)
  \cong_{\mathrm{symp}}T^*\mathcal H^+_{m_0,c};
 \]
 thus this stratum is exactly the future-timelike spinless relativistic
 particle.  For \(m_0=0\), the intrinsic nonzero rank-one quotient similarly
 gives \(T^*\mathcal N_+^\times\), but the written multiplier action is
 irregular because
 \[
  \det(\alpha Z^*Z)=\alpha^2\sum_{a<b}|\Delta_{ab}(Z)|^2,
  \qquad d(\det\circ P_\alpha)_Z=0\quad
  (\operatorname{rank}_{\mathbb C}Z\leq1).
 \]
 Its rank-one equation is \(\dot x=\lambda p\), with no equation
 \(\lambda=e\); four real transverse rank-opening directions are
 first-order null but are not colour gauge, and the apex \(Z=0\) admits
 arbitrary \(x(\tau)\).  On the later hard-constraint carrier with
 \(\omega_0\ne0\), the momentum projection has rank four for \(w\ne0\) and
 rank three at \(w=0\); direct variation forces \(\dot x=0\) strictly above
 the sharp threshold and gives \(\dot x\in\mathbb Rp\) only at the threshold.
 Its determinant bound and multiplication/Fock spectra therefore remain exact
 algebraic and free-operator results, but the multiplication operator is not
 the Hamiltonian generating the displayed worldline action and the threshold
 is not a dynamically generated interacting gap.  See
 \cref{hlob:thm:split-observation,hlob:thm:block-coordinates,%
  hlob:thm:lorentz-action,hlob:thm:cone-section,hlob:cor:Y3-block-form,%
  hlob:thm:colour-momentum,wpr:thm:massive,wpr:thm:massless,%
  wpr:thm:zero,wpr:thm:packet-action}.

 The same marked data now carry an exact field-theoretic continuation on the
 affine normal slice
 \[
  \mathcal N_P=\{Y\in\operatorname{Herm}_3(\mathbb C):Yu=0\}
  =\jmath_P\bigl(\operatorname{Herm}_2(\mathbb C)\bigr).
 \]
 Polarization of
 \(\det(uu^*+\varepsilon\jmath_P(G))=\varepsilon^2\det G\)
 gives a constant metric of signature \((1,3)\), and the congruence action
 induces the spin cover
 \(SL_2(\mathbb C)\to SO^+(1,3)\).  The standard Moore--Penrose metrics select
 \[
  W_P=R^*R=C^{-1},\qquad
  K_P=\{a\in SL_2(\mathbb C):a^*W_Pa=W_P\}\cong SU(2),
 \]
 acting on the real normal field as \(\boldsymbol1\oplus\boldsymbol3\).
 Gauging this pointwise action yields an explicit nonnegative classical
 gauge--scalar action; the centre acts trivially on the normal field, so it
 cannot screen a fundamental centre charge without the retained
 centre-nontrivial spinor or another added field.  The candidate inverse-dual
 matrix remains outside the positive Gram momentum cone because
 \[
  \det G_{\eta,T}=-\frac23D_T<0,
  \qquad \operatorname{im}p_\alpha=
  \operatorname{Herm}_2(\mathbb C)_{\geq0}.
 \]
 On this same Lorentzian host, the Klein--Gordon/Fock construction has exact
 gap \(m_{\mathrm{KG}}\) for an independently supplied
 \(m_{\mathrm{KG}}>0\), is gapless at \(m_{\mathrm{KG}}=0\), and the massless
 compact gauge Hessian has lowest nonzero Gaussian frequency \(2\pi/L\to0\).
 Thus neither the determinant identity, the Gram cone, nor entanglement alone
 chooses an interacting mass gap.  A full-domain comparison satisfying a
 uniform norm lower bound and energy-form upper bound would instead imply the
 quantitative estimate \(\operatorname{gap}(A_{\mathrm{phys}})\geq
 \gamma c_R^2/C_R\); an intertwiner on only one gapped subspace does not.
 See \cref{hnsp-lorentz-theorem,hnsp-intrinsic-compact-representation,%
 hnsp-intrinsic-gauge-scalar-theorem,hnsp-candidate-spacelike,%
 hnsp-free-field-gap,hnsp-gauge-gaussian-gap,hnsp-coercive-transfer}.

 The selectively corrected Jordan determinant/Euler audit identifies an exact family of
 Lorentz sections inside the marked Hermitian cone.  With
 \([a]_{\times}v=a\times v\), one has for every real symmetric \(S\)
 \[
  \det\bigl(S+i[a]_{\times}\bigr)=\det S-a^{\mathsf T}Sa.
 \]
 Hence every marked positive monodromy form \(B\) defines
 \[
  \ell_B(\lambda,a)=\lambda B+i[a]_{\times},\qquad
  \ell_B^{-1}\!\left(\operatorname{Herm}_3(\mathbb C)_{\geq0}\right)
  =\{\lambda\geq0:a^{\mathsf T}Ba\leq\lambda^2\det B\}.
 \]
 There is also a real-linear order-cone bijection from
 \(\operatorname{Herm}_2(\mathbb C)_{\geq0}\) onto the positive part of
 \(\mathbb RI_3\oplus i\mathfrak{so}(3)\) for which
 \[
  \det_3\Theta(H)=\frac{\operatorname{tr}H}{2}\det_2H;
 \]
 it sends positive-semidefinite ranks \(2,1,0\) to \(3,2,0\).  If \(P\) is a
 rank-one projection, then \(\Theta(P)\) has eigenvalues \(0,1/2,1\) and is
 not idempotent, which separately proves that \(\Theta\) is not a Jordan
 homomorphism.  These are exact algebraic and order-theoretic maps, not holomorphic
 toric, period-lattice, or filling maps.  In particular, the toric limit
 \(z_0z_1z_2z_3=t\to0\) keeps \(B=I_3\) fixed and is not the operation
 \(\det B\to0\).

 The corresponding rank-three toroidal degeneration is now constructed as a
 separate analytic object.  Its smooth fibre is \(E_\tau^3\), its deepest
 local equation is \(t=z_0z_1z_2z_3\), and its irreducible central fibre has
 stratum counts
 \[
  (f_0,f_1,f_2,f_3)=(1,7,12,6),\qquad \chi(W_3)=6.
 \]
 Its normalization has orbit vector \((1,14,36,24)\); the strata of positive
 codimension map with degrees \(2,3,4\) onto the \(7,12,6\) quotient strata.
 Thus any real \(S^8\)-total-space fibration retaining \(W_3\) must have its
 remaining singular fibres contribute Euler sum \(-4\).
 More generally, an injective integral shear
 \(\mathcal B:\bar L_3\hookrightarrow L_3\) gives
 \[
  \chi(W_{3,\mathcal B})=3!\,[L_3:\mathcal B\bar L_3]
  =3!\,|\det\mathcal B|.
 \]
 The toroidal object is not merely analogous to the reconstructed candidate
 cusp: the explicit epimorphisms
 \[
  P=\begin{pmatrix}1&-1&0\\0&1&-1\end{pmatrix},
  \qquad
  Q=\begin{pmatrix}0&-1&1\\1&-1&0\end{pmatrix}
 \]
 give \(\Phi=P\oplus Q:H_1(T^6)\twoheadrightarrow H_1(T^4)\) with
 \[
  \Phi M_3=M_F\Phi,\qquad c_F\Phi=Qc_3,
  \qquad PZ_3(t)=Z_2(t)Q.
 \]
 This extends to a holomorphic map of local cusp families with connected
 elliptic kernel.  An explicit unimodular change of basis identifies the
 rank-three cusp matrix with \(M_F\oplus
 \bigl(\begin{smallmatrix}1&1\\0&1\end{smallmatrix}\bigr)\), and compatible
 order-three and order-four affine peripheral actions exist on the added
 elliptic factor.  Under the marked \(A_2\)-Cartan identification, the six
 Kuhn tetrahedron orbits map equivariantly to the six diagonal Weyl chambers.
  This last correspondence is fan-theoretic: it neither identifies the
  four-branch toroidal germ with the quadratic Gram-rank germ nor supplies the
  still-open analytic gluing over both finite special fibres.

  The free rank-six object native to this construction is now identified
  without conflating it with a cyclic quotient or with the different defect
  group on the original regular locus:
  \[
   \Lambda_{\rm tor}=H_1(T^6;\mathbb Z)
   =\mathbb Z^3_{\rm angular}\oplus\mathbb Z^3_{\rm deck}.
  \]
  With diagonal coordinate permutations it fits into the exact sequence of
  integral \(S_3\)-modules
  \[
   0\longrightarrow\mathbb Z(k,0)\oplus\mathbb Z(0,k)
   \longrightarrow\Lambda_{\rm tor}
   \xrightarrow{\,P\oplus Q\,}A_F\oplus B_F\longrightarrow0,
   \qquad k=(1,1,1).
  \]
  The sequence splits integrally with monodromy but not over
  \(\mathbb Z[S_3]\); its quotient copies are \(A_2\) weight lattices, with
  the root lattices sitting at index three.  The rank identity
  \(\operatorname{rk}\Lambda_{\rm tor}=6\), the chamber count
  \(|S_3|=3!=6\), and
  \(\chi(W_{3,\mathcal B})=3!|\det\mathcal B|\) are thereby separated: the
  six homology basis vectors form two three-element \(S_3\)-orbits, whereas
  the six Kuhn/Weyl chambers form one regular orbit.  That regular orbit does
  occur exactly in the associated rank-nine carrier:
  \[
   \operatorname{Hom}(\overline L_3,L_3)
   =\bigoplus_i\mathbb ZE_{ii}\oplus
    \bigoplus_{i\ne j}\mathbb ZE_{ij},
  \]
  where the off-diagonal matrix units form \(\mathbb Z[S_3]\), and the
  chamber \(\pi=(\pi_1\pi_2\pi_3)\) corresponds to
  \(E_{\pi_1,\pi_3}\).  Equivalently, the chambers are the six mixed
  angular--deck basis pairs with zero alternating pairing; the three
  excluded diagonal pairs have pairing one.  The real-linear
  \(S_3\)-isomorphism
  \[
   X\longmapsto\frac{X+X^{\mathsf T}}2
          +i\frac{X-X^{\mathsf T}}2
  \]
  identifies the ambient matrix representation with
  \(\operatorname{Herm}_3(\mathbb C)\), but has integral index \(2^3\) in
  the symmetric/skew coordinates and sends every off-diagonal \(E_{ij}\) to
  an indefinite Hermitian matrix with eigenvalues
  \(2^{-1/2},-2^{-1/2},0\).  It is therefore an exact representation and
  chamber-incidence bridge, not a positive, Jordan, determinant, or toric
  morphism.  The separate map into the Gram boundary is only the nonadditive
  piecewise-linear chamber map
  \(\Phi_{\rm G}\circ\jmath\circ P\) on the angular cocharacters; it ignores
  the deck half and does not descend to the torus.  See
  \cref{rsl:thm:s3-sequence,rsl:thm:three-sixes,rsl:prop:gram-scope,%
  rnmt:thm:regular-torsor,rnmt:prop:representation-bridge}.
  This successful lift is distinct from the diagonal local-system candidate
 \[
  M_1=I_3\otimes S,\qquad M_2=I_3\otimes Q,\qquad M_0=I_3\otimes T,
  \qquad M_1M_2M_0=I.
 \]
 Although its finite factors have orders \(4\) and \(6\), each
 \(I-M_j:T^6\to T^6\) is an isogeny for \(j=1,2\); consequently every affine
 lift of either finite action has a fixed point.

 A literal complex-fourfold analogue on \(S^8\) is impossible: an almost
 complex tangent bundle would have
 \(c_1=c_2=c_3=0\), \(c_4=2u\), hence \(p_2=4u\), while the signature
 formula forces \(p_2=0\).  A nonsingular \(T^6\)-bundle over \(S^2\) is
 likewise impossible; any real singular replacement must derive its actual
 filling maps and satisfy the displayed Euler, van Kampen, Wang, and
 Mayer--Vietoris tests.  For its regular-piece monodromy \(\rho\), the
 corrected homological Serre filtration is
 \[
  0\longrightarrow(\textstyle\bigwedge^q\Lambda)_\rho
  \longrightarrow H_q(E;\mathbb Z)
  \longrightarrow
  H_1(B^\circ;\textstyle\bigwedge^{q-1}\Lambda)
  \longrightarrow0,
 \]
 while the dual edge map is
 \(H^q(E;\mathbb Z)\twoheadrightarrow
 (\bigwedge^q\Lambda^\vee)^\rho\).  Thus the filling maps must preserve the
 coinvariant subobject and induce the required map on the base-one-cycle
 quotient; using the reverse filtration invalidates the map-level test.
 Every effective torus action on \(S^8\) has rank at most four, and the bound
 is sharp for the coordinate action
 \[
  S^8/T^4\cong\operatorname{Susp}(\Delta^3)\cong B^4.
 \]
 This is a real two-pole singular compactification with generic \(T^4\)
 orbit, not a three-special-fibre \(T^6\)-fibration over \(S^2\).

 The exact bundle-theoretic alternative occurs on
 the regular locus
 \(U=X\setminus(W\sqcup S_1\sqcup S_2)\):
 \[
  H^4(U;\mathbb Z)
  \cong H_1(W;\mathbb Z)\oplus H_1(S_1;\mathbb Z)
       \oplus H_1(S_2;\mathbb Z)
  \cong\mathbb Z^2\oplus\mathbb Z^2\oplus\mathbb Z^2.
 \]
 Principal \(SU(3)\)-bundles on \(U\) are classified by this complete
 \(c_2\)-class and extend across the three deleted fibres exactly when it
 vanishes.  This is a free rank-six defect lattice, not a cyclic group of
 order six.  The formulas from this point through the normal-holonomy and
 common-centre comparison are exact results of the dated fibre-restriction,
 defect-lattice, and common-centre auxiliary audits.  Their full theorem--proof
 source has not yet been integrated into a substantive research fragment of
 this workbench, so they are recorded here without a \Verified\ status claim.
 In the marked basis
 \[
  (a_0,b_0,a_1,b_1,a_2,b_2)
  =(\bar\gamma,\bar u,g_1,h_1,g_2,h_2),
 \]
 the actual collar matrix gives the primitive fibre-degree functional
 \[
  [\deg_F]=(12,0,4,0,-3,0),\qquad
  \ker\deg_F=\mathbb Z\langle
  b_0,b_1,b_2,a_0-3a_1,3a_1+4a_2\rangle.
 \]
 Define
 \[
  q_{\rm col}:L\longrightarrow H_1(X^\circ;\mathbb Z),\qquad
  a_0\mapsto c,\quad a_1\mapsto\bar g_1,\quad a_2\mapsto\bar g_2,
  \quad b_0,b_1,b_2\mapsto0,
 \]
 and let
 \(\Phi_1:H_1(X^\circ;\mathbb Z)\to\mathbb Z\) have row
 \((12,4\ell_1,3\ell_2)\) in the basis \((c,\bar g_1,\bar g_2)\).
 Then the typed identity is
 \(\deg_F=\Phi_1\circ q_{\rm col}\), and
 \(\Phi_1(j_*[\widetilde a_0])=p\).  The order-three and order-four normal
 holonomies combine, with the candidate's relative sign, to characters
 \[
  \chi_{12}=\rho_{12}\circ\deg_F,\qquad
  \chi_6=\chi_{12}^2=\rho_6\circ\deg_F,\qquad
  \rho_m(k)=e^{-2\pi ik/m}.
 \]
 Thus the natural finite quotients are
 \(H^4(U)/\ker\chi_{12}\cong\mathbb Z/12\) and
 \(H^4(U)/\ker\chi_6\cong\mathbb Z/6\); the latter is not inferred from
 the rank.  Ordinary coefficient reduction instead gives
 \(H^4(U;\mathbb Z/6)\cong(\mathbb Z/6)^6\); this is not the single
 distinguished cyclic quotient selected by \(\chi_6\).  Equivalently, for a
 principal \(SU(3)\)-bundle \(P\to U\),
 \[
  \chi_6(c_2(P))
  =\exp\!\left(-\frac{2\pi i}{6}
  \langle c_2(P)|_F,[F]\rangle\right).
 \]

 Tensor product supplies a separate exact centre statement:
 \[
  SU(2)\times SU(3)\hookrightarrow SU(6),\qquad
  (A,B)\longmapsto A\otimes B,\qquad
  \mu_2\times\mu_3\xrightarrow{\sim}Z(SU(6))=\mu_6.
 \]
 On the five-dimensional CW model of \(U\), \(c_2\) classifies
 \(SU(n)\)-bundles for \(n\ge3\), and
 \(c_2(E_2\otimes E_3)=3c_2(E_2)+2c_2(E_3)\).  Hence every chosen
 \(h\in H^4(U;\mathbb Z)\) is represented by an \(SU(6)\)-bundle admitting
 a tensor-subgroup reduction; its \(c_2\)-factorizations are
 \[
 (e,f)=(h+2t,-h-3t),\qquad t\in H^4(U;\mathbb Z),
 \]
 so the geometry selects neither a nonzero \(h\) nor a preferred reduction.
 Separately,
 \[
  H^5(U;\mathbb Z)
  \cong\ker\{\mathbb Z^3\xrightarrow{(1,1,1)}\mathbb Z\}
  \cong Q(A_2)
 \]
 for the combinatorial form making the three support components
 orthonormal; this is not a cup-product form and does not identify
 \(\ker\deg_F\) with \(A_5\).
 The finite-group identification
 \(H^4(U)/\ker\chi_6\cong Z(SU(6))\) is not an \(SU(6)\)-connection or a
 gauge-transition cocycle.  The tensor map is injective and does not factor
 through \((SU(2)\times SU(3))/\mu_6\), because its central \(\mu_6\) maps
 to nontrivial scalars.  The distinct common-centre map is
 \[
  \Gamma_6=\{((\lambda\mu)^{-1},\lambda I_2,\mu I_3):
              \lambda\in\mu_2,\ \mu\in\mu_3\},
  \qquad
  \begin{aligned}
  \frac{U(1)\times SU(2)\times SU(3)}{\Gamma_6}
  &\hookrightarrow U(6),\\
  [z,A,B]&\longmapsto z(A\otimes B),
  \end{aligned}
 \]
 where \(\Gamma_6\cong\mu_6\).  Its determinant is \(z^6\), and its image
 meets \(SU(6)\) in exactly the tensor subgroup
 \(\iota(SU(2)\times SU(3))\); it is not generally contained in
 \(SU(6)\).

 Fibre evaluation therefore gives a canonical surjection
 \[
  \bar\kappa_F:H^4(U;\mathbb Z)\twoheadrightarrow
  H^4(T^4;\mathbb Z)\twoheadrightarrow\mathbb Z/6\mathbb Z.
 \]
 Identifying this degree-four quotient with a degree-two common-centre bundle
 obstruction remains unproved.  Global \(SU(3)\)-bundles on \(S^6\) instead carry the separate
 clutching class \(c_3=\pm2k u\).  The Lorentz-observation module and its
 exact presymplectic reduction identify the massive particle, the punctured
 massless quotient, the irregular massless multiplier/apex branches, and the
 strict-interior versus threshold dynamics of the hard-constraint packet.
 These carrier, reduction, and bundle theorems do not supply a Yang--Mills
 mass-gap conclusion: the required field connection,
 interacting operator and domain, gauge-invariant spectral sector, and
 continuum comparison morphism remain to be constructed.

 Reviewer~B independently reviewed the reconstruction existing up to its
evidentiary cutoff and provisionally agreed on the non-collapse route; the
later Codex proofs here are separate.  In a 26 August follow-up, Reviewer~B
also independently supplied the finite-complex/UCT proof of top-degree base
change and confirmed the conductor--duality nonvanishing from the three named
candidate inputs; see \cite{ReviewerBValidation2026}.  Throughout, \Verified\ means that a
complete auditable proof is written in this document, not that specialist
replication or community acceptance has already occurred.  The candidate
construction and its original claims remain attributed to their reported
source.  This paper is an AI-produced verification aid: its own contribution
is to expose the relevant morphisms, signs, lattices, comparison maps, and
singular-fibre arguments as complete checkable proofs, thereby reducing the
reconstruction work required of human experts.  It makes no claim of priority
for the candidate and does not replace independent specialist adjudication.
The 4 September divisor-channel audit verifies the oriented composite-shell
bijection and the diagonal-ray density-one argument. It also constructs a
separate rational map from $ac-b^2=-\delta$ to the defect-one Koide surface,
with an explicit inverse on stated open sets and exact image denominators;
this does not replace the supplied orthogonal map, whose rational image is
empty. See Section~\ref{sec:esv3-audit}.
The quaternionic continuation audit identifies the two proposed Lorentz
$4+4+4+1$ slices as the same embedded space, with
$w=-\varepsilon\bar\chi$, and proves the positive-cone sphere pair
$(S^{12},S^8)$. The retained marking transports the Jordan product to
$(MDN+NDM)/2$, $D=\operatorname{diag}(C^{-1},3)$, and its closure is the
full fifteen-dimensional Hermitian quaternionic algebra. For the supplied
scalar mechanics the omitted coordinate $\rho=u^*b$ instead obeys
$\ddot\rho=-(m^2+\Gamma s)\rho-\kappa u^*\Phi b$; explicit initial data
give $\ddot\rho=-2\kappa$. The exact constrained alternative and the
gauge-sign correction are in Section~\ref{sec:q444-reconciliation}.
Within the unrestricted equations, fixed complex $W\subset u^\perp$
with $b\in W$, $\Phi W\subset W$, and tangent velocities yield global
nonzero invariant sectors of configuration dimensions six and eight,
without reaction forces. Other initial data have zero normal derivatives
through order three but $\rho^{(4)}(0)=-12\kappa$.
The fixed-subspace sectors are not exhaustive: an exact periodic
moving-kernel family remains in $\rho=0$ outside all of them. At
$(m^2,\Gamma,\kappa)=(1,1,6)$ its norm is $s=1$, energy $7/4$, and
conserved moment $-\Omega/2\ne0$, so its zero-connection gauge lift fails
Gauss' equation. The scalar global-minimum transition at
$m^2=\kappa^2/(27\Gamma)$ and full Hessian multiplicities $1,8,5$ are
independently verified; no gauge-reduced quantum spectrum follows.
The full positive quaternionic carrier admits a split positive observation
to $\operatorname{Herm}_2(\mathbb C)$ whose exact permitted pairs satisfy
$G\succeq0$, $|\rho|^2\leq\det G$. Every such pair lifts, and explicit
positive lifts approach determinant zero with $\rho\ne0$; no uniform
positive gap follows from this cone. For trace-free scalar matrices the
separate map $Q\mapsto Q-\lambda_{\min}(Q)I_3$ gives a positive-boundary
observation retaining $PP^*$, not an assumed positive scalar trajectory.
The retained rank-six period matrix has complete positive real metric
fibres of dimensions twelve over its Hermitian Gram matrix and seventeen
over its Lorentz observation. A rank-losing observation is realized by
an explicit equivariant metric with diverging relative condition number;
bounded relative condition number excludes that rank loss.
For the periodic scalar orbit, every homogeneous temporal-gauge lift
fails Gauss' equation, even allowing arbitrary spatial matrices.
A different interacting radial family solves all gauge and scalar
equations exactly. On the positive rank-one boundary, colour-group
averaging converts the marked determinant to the gauge invariant
$\operatorname{tr}\operatorname{adj}(A)=\|b\|^2$, which varies continuously
along an explicit family of scalar minima. These are classical
observation and field-equation results, not an interacting quantum gap.
The exact continuation now supplies a current-free three-real-coordinate
sector, nonlinear null waves, and periodic full-gauge motions with
arbitrarily small positive frequency and vacuum energy excess at fixed
volume. The periodic amplitude stays finite. The full transverse
gauge-inertia spectrum and conserved invariant current identify the branch
$\mu_{\rm d}^2=(2g^2/3)\|b\|^2$ with the colour-averaged Lorentz determinant.
The literal earlier rotating orbit has no arbitrary-time real-connection
lift and no constant-spatial-field complex lift in its constant-matter
frame. A further genuinely complex rotating-pair block class and its
constant spatial rotations are excluded. Without that block restriction,
Gauss gives $C=\omega T$ with
$T=4(g^2\mathcal I+2\mathcal K)^{-1}[Z_1,Z_2]$ and the full scalar
equation fixes $\omega^2=(1/4+\sum_j\|Z_jq\|^2)/\|Tq\|^2$.
All remaining matrix/vector equations and both frequency signs are
retained in a proved inverse correspondence; its zero set and more
general complex time-dependent lifts remain open.
The full scalar equations force $s_C^2<s_{Z_1}^2+s_{Z_2}^2$ for
$s_X=-i\operatorname{tr}(HX)$ unless all three vanish; the two spatial
equations give an additional positive stress matrix and its exact
sum-of-squares discriminant, without a block ansatz.

For each fixed unrestricted spatial pair, we now prove an explicit
bounded coupling interval and a nonnegative residual polynomial of
degree exactly twenty-eight whenever its domain polynomial is nonzero.
Its admissible positive roots are precisely the allowed original $g^2$:
there are at most fourteen, with both frequency signs reconstructed.
Three vector tests decide the identically empty Gauss domain.
A nonzero, explicitly defined $24\times24$ Sylvester determinant
polynomial further confines the entire spatial solution projection to
a closed, nowhere-dense, measure-zero subset of the sixteen-coordinate
carrier. An exact unrestricted witness and rational and finite-field
Bezout certificates prove that polynomial is nonzero. The eight-dimensional linear class in which both spatial matrices
exchange the matter line with its orthogonal complement lies inside this
exceptional set, but a full-equation positivity identity excludes every
positive coupling and both frequency signs there. A denominator-free
residual identity also gives a strictly positive spatial-error bound on
its scalar--Gauss zero set, uniform for bounded spatial norm and coupling
bounded away from zero. This scalar--Gauss zero set is now completely
parameterized: the classified domain is
$\mathcal D\times\{1,-1\}^2\times U(1)^2$, and the literal domain
retains the frame phase subject only to
$(\zeta,p_+,p_-)\sim(-\zeta,-p_+,-p_-)$.
An explicit family attains every original $t=g^2>0$, with all scalar
and Gauss entries zero. The two spatial residuals cannot vanish simultaneously, and
existence on the remaining exceptional set is open.

In the period bridge, Hermitian compatibility makes
$W$ a lattice cometric with $h_3=2H_3^{-1}$.
The native Euclidean choice gives $H_3=2I_3$, $G=2PP^*$ and the exact
source/quotient Laplace asymptotics $4\pi^2/y^2$, $8\pi^2/y^2$.
Constant observed cometric therefore coexists with expanding periods and
vanishing positive Laplace gaps.
For the full homogeneous gauge--scalar mechanics, retaining all
thirty-eight configuration coordinates and imposing its constant-$SU(3)$
Gauss law, a new direct form estimate proves compact resolvent and a
unique invariant quantum ground state. Its first physical excitation
gap is strictly positive at each fixed positive volume and coupling.
This finite-model theorem does not identify the homogeneous quantum
mechanics with an invariant subsector of the full torus field theory:
an explicit mixed Fourier-mode quartic term obstructs the automatic
fixed-vacuum product construction. No continuum or volume-uniform
quantum gap is inferred.
For that exact two-coordinate Fourier mechanics, the actual positive
quartic fibre ground state gives an isometry $Jf=f\Omega_q$ and
$H_2J-JH_{\rm eff}=-\hbar^2(2f'\Omega'_q+fQ_q\Omega''_q)/(2\mathcal V_3)$.
The effective Hamiltonian retains its ground energy and the nonzero
Born--Huang potential. Explicit global derivative bounds prove a
finite-energy physical-time error of order $\mathcal V_3^{-1/2}$,
with the other parameters fixed and a stated sufficient volume regime.
This is a controlled finite-model correspondence, not a continuum
or thermodynamic mass-gap result.
\end{abstract}

\tableofcontents

\section{Results and evidence boundary}

\noindent\textbf{6 September: explicit spectral and umbral maps, and the original family's parameter identifications.}
For the literal common lattice $\mathcal I=\{\sqrt2(a,a,a):a\in D_4\}$,
write $a=a_0+v$, $t^2=|v|^2$, $Q=a_0^2+t^2$ and
$P=a_0(a_0^2-3t^2)$. Its complete cubic image is
\[
 \tau(\mathcal I)=2\sqrt2\bigl(18\mathbb Z\cup(6\mathbb Z+2)
                                      \cup(6\mathbb Z-2)\bigr).
\]
Three explicit integral affine lines give a section of this value map.
The scalar-retaining Albert matrix $X=\mu I+Y(a)$, with cyclic
off-diagonal entries $\sqrt2a$, has invariants and labelled spectrum
\[
\begin{aligned}
 (\operatorname{tr}X,\operatorname{tr}X^2,\det X)
   &=(3\mu,3\mu^2+12Q,\mu^3-6\mu Q+4\sqrt2P),\\
 \operatorname{Spec}Y
   &=(2\sqrt2a_0,-\sqrt2a_0-\sqrt6t,-\sqrt2a_0+\sqrt6t).
\end{aligned}
\]
For $t>0$, the explicit labelled eigenprojector satisfies
$(\Pi_1)_{12}=-1/6-\sqrt3(v/t)/6$; thus the frame and spectrum
recover $a$, not only its norm and cubic. The full discriminant is
$864t^2(3a_0^2-t^2)^2$. The full Cartan family retains the scalar
$\mu$, radius $r$, and angle $\theta$, with $r^2=12Q$ and
$r^3\cos3\theta=24\sqrt3P$ on this subfamily.

For the same marked root system $A_5^4D_4$, the
Cheng--Duncan--Harvey prescription gives
$S^X=(5I_{12}+J)S_6$, where $J_{rs}=\mathbf1_{s=5r\ (12)}$ and
$S_{6,r}=\sum_{k\equiv r\ (12)}kq^{k^2/24}$.
The coefficient map to the retained Fricke function is
\[
 T_6=\frac{\eta(\tau)^5\eta(3\tau)}{\eta(2\tau)\eta(6\tau)^5},
 \qquad f^X=-q\frac{d}{dq}\log T_6,
 \qquad T_6\!\left(-\frac1{6\tau}\right)=\frac{72}{T_6(\tau)}.
\]
All twelve theta indices, the multiplier determinant, the original
glue-group action and its twisted shadows are computed in the proofs.
This is an explicit use and comparison of established umbral machinery,
not a claim to a new moonshine module.

In the original analytic family $\beta_c=\beta^{\rm part}+c$,
$\operatorname{Im}c<-M$, the actual global identifications are exactly
$X_c\simeq X_{c'}\Longleftrightarrow c'-c\in2\mathbb Z$.
The middle lattices agree for every integer shift; the multiplicity-four
filling retains an order-two obstruction to odd shifts. Even shifts have
explicit finite-filling and toric cusp maps and a proved global additive
gluing correction. The resulting parameter quotient is
$w=e^{-\pi ic}$, $0<|w|<e^{-\pi M}$.
This classifies the retained analytic family without using the claimed
identification of its total space with $S^6$.


\paragraph{6 September: cubic arithmetic, integral lattice maps, and the canonical ring.}
For the retained $A_5^4\perp D_4$ Niemeier lattice $N$ and its
original index-six Leech neighbour $\Lambda$, fix the Cayley--Dickson
frame $(1,i,j,k,\ell,i\ell,j\ell,k\ell)$ used in the complete cubic
proof. Put $\tau(q)=\operatorname{Re}((q_1q_2)q_3)$ and
$F=\tau\circ\mathcal L$ on the original Niemeier span,
and $\mathcal V_F(L)=\operatorname{span}_{\mathbb Z}\{F(y):y\in L\}$.
The exact additive groups are
\[
\begin{aligned}
 \mathcal V_F(N)&=\tfrac1{27}\mathbb Z\oplus
       \tfrac{\sqrt2}{54}\mathbb Z\oplus\tfrac{\sqrt3}{18}\mathbb Z,\\
 \mathcal V_F(\Lambda)&=\tfrac1{27}\mathbb Z\oplus
       \tfrac{\sqrt2}{216}\mathbb Z\oplus\tfrac{\sqrt3}{54}\mathbb Z,
 \qquad \mathcal V_F(\Lambda)/\mathcal V_F(N)\cong\mathbb Z/12.
\end{aligned}
\]
The Albert determinant is $\lambda^3-\lambda(y,y)+2F(y)$.
The quotient concerns the displayed value-generated groups in this fixed
frame, not an inclusion $N\subset\Lambda$. The actual residue
$r=\rho(361)$ cycles three Leech neighbours $\Lambda_j=r^j\Lambda$.
Their pairwise intersection is one integral lattice $J$, with
$\det J=81$ and $[\Lambda_j:J]=9$. The fixed fourth isotropic plane
in $J^\vee/J$ gives $M$ with exactly $72$ roots of type $A_2^{12}$,
$M\cap\Lambda_j=J$, and $[M:J]=9$.
\[
 \mathcal V_F(J)=\mathcal V_F(M)
 =\tfrac1{27}\mathbb Z\oplus\tfrac{\sqrt2}{216}\mathbb Z
                         \oplus\tfrac{\sqrt3}{18}\mathbb Z.
\]
The successive value-group inclusions have matrices
$\operatorname{diag}(1,4,1)$ and $\operatorname{diag}(1,1,3)$,
giving quotients $\mathbb Z/4$ and $\mathbb Z/3$.
Both containments and their reverse generator certificates are calculated
from the full cubic polarization, not inferred from lattice indices.

Wilson's construction \cite{esn:Wilson2009}, transported by the explicit
isometry $\iota(x)=x/\sqrt2$, gives $\Lambda_{\rm cyc}$ in the
original triality metric. Its cyclic fixed Gram matrix is $3G_{E_8}$;
its orthogonal complement is explicitly isometric to $A_2\otimes E_8$,
and their sum has index $3^8$, with quotient $\mathbb F_3^8$.
The directions remain $c(q_1,q_2,q_3)=(q_3,q_1,q_2)$ and
$c_{\rm W}=c^{-1}$. Its literal intersections are
\[
 \Lambda_{\rm cyc}\cap\mathcal L(N)
 =\Lambda_{\rm cyc}\cap\mathcal L(\Lambda)
 =\Delta:=\{\sqrt2(a,a,a):a\in D_{\mathbb H}\},
\]
where $D_{\mathbb H}$ is the even-coordinate-sum integer lattice in
the first four stated Cayley coordinates. Thus $\Delta$ has rank $4$,
Gram matrix $6G_{D_4}$, determinant $5184$, minimum squared norm $12$,
and value-generated subgroup $4\sqrt2\mathbb Z$ for $\tau|_\Delta$, or
$8\sqrt2\mathbb Z$ for $2\tau|_\Delta$. This is not the actual image:
$12\sqrt2$ is not attained by $\tau$ on $\Delta$.
The preceding proofs retain the triality products and spectral boundary
graph, the flag's degree-$+1$ collapse, and the retained cusp class's
nontrivial three-plane reduction; these exact maps do not supply an
interacting quantum mass-gap result.

For the original analytic fibration $f:X\to\mathbb P^1$, the full
canonical-divisor proof gives
$\omega_X\simeq\mathcal O_X(-2S_2)$ and the graded anticanonical ring
$\bigoplus_{m\ge0}H^0(X,\omega_X^{-m})\cong\mathbb C[U,V]$,
with $\deg U=1$, $\deg V=2$. The complete system $|-2K_X|$
recovers $f$ with the specified projective coordinates. This calculation
uses the original analytic charts and finite fillings, not an
identification of $X$ with $S^6$, and draws no sphere-topology conclusion.

The retained constant $\beta_c=\beta^{\rm part}+c$ now gives a proper
holomorphic deformation over $\operatorname{Im}c<-M$, locally on
relatively compact parameter discs. On a fixed regular fibre set
$T=\operatorname{Im}\tau$, $m=\operatorname{Im}\mu$,
$D=\operatorname{Im}\beta_{c_*}-6m^2/T<0$ and $\delta=c-c_*$.
In the plus-graph convention its exact Beltrami matrix is
\[
 \nu_\delta=\frac{\delta}{2iD+\delta}
                 \begin{pmatrix}0&0\\-m/T&1\end{pmatrix}.
\]
The connecting Kodaira--Spencer class is
$-[\partial_{\zeta_2}\otimes
(d\bar\zeta_2-(m/T)d\bar\zeta_1)/(2iD)]\ne0$.
The intrinsic pencil recovers $f$ also over the dual numbers; its three
fixed critical values prohibit an infinitesimal projective base change.
Consequently this deformation of the whole $X_c$ is nontrivial,
not only its marked regular-fibre deformation. The proof uses no sphere
recognition or CDP conclusion.



The target is a reconstructed manuscript claiming a compact complex
threefold \(X\) and holomorphic map \(f:X\to\mathbb P^1\) with
\[
X\cong_{\mathrm{diff}}S^6,\qquad a(X)=1,\qquad
\langle c_3(T_X),[X]\rangle=2.
\]
As of 6 September 2026, the workbench records the following mathematical
status.  Here \Verified\ means that a complete proof is written in this
workbench; it is not a proxy for a model verdict or for community acceptance.

The 5 September direct-image recheck corrects two written proof steps.
Theorem~\ref{thm:cartier-duality} now constructs its duality morphism
from the analytic trace, not from matching outer terms. The logarithmic
gluing is the real translation
$a_j(s)=\operatorname{Re}\alpha\,v_j+\operatorname{Im}\alpha\,J_sv_j$,
where $J_s=\Pi(s)^{-1}i\Pi(s)$ and $\alpha=\log s_j/(2\pi i)$.
Its full clockwise loop has increment $-v_j$; its mixed sweep integral
is $-E(v_j,\lambda)$. The corrected proof of
Lemma~\ref{lem:picard-finite-norm} and the based calculation in
Proposition~\ref{prop:geometric-pi1} retain their original signs and
coefficients. The bounded recheck verifies these maps, not every
global candidate input. Its before/after sources and primary-source
audit are retained in the direct-image evidence folder.

The period continuation now identifies the exact admissible constants:
for the cusp-regular $\beta=\beta^{\rm part}+c_0$, set
$M=\max_{B\setminus\{p_0\}}D_{\rm part}$, an attained finite maximum.
The marked real period map is everywhere an isomorphism exactly when
$\operatorname{Im}c_0<-M$
(Proposition~\ref{prop:period-exact-admissible-constants}). Omitting the
cusp extension permits $\widetilde\beta=\beta+i\,t\circ\pi$ with
$\widetilde D=D+\operatorname{Re}t$ unbounded above.
The cusp proof uses charts restricted to $Y_\epsilon$, a bound on a
fixed closed subdisc, and the weak inequalities for closed polydiscs
(Lemma~\ref{lem:toric-chart-control}). Its exact opposite-edge maps
retain the three multipliers $e^{-2\pi i\mu_0}$,
$e^{12\pi i\mu_0}$, $e^{-2\pi i(7\mu_0+b_0)}$
with the same-character conventions of
equation~\eqref{eq:cusp-exact-opposite-edge}. These bounded period,
cusp and finite-filling rechecks add no sphere or CDP theorem.

The new higher-wrap comparison is fully written in
Section~\ref{sec:h24-continuation} and in the root draft
\path{s6_higher_rung_24d_preprint.tex}. It gives the octonionic
positive-cone compactification $(B^{25},S^{24})$, the triality flag
$F_4/\operatorname{Spin}(8)$ with integral cohomology ranks $1,2,2,1$,
and the typed peripheral and suspension maps involving
$\mathbb Z/12\hookrightarrow\mathbb Z/24$, $a\mapsto2a$.
The flag facts are attributed to Mare--Willems; the comparison preserves
all coordinate, monodromy, and coefficient choices. The particular
one-copy cusp homotopy coefficient is now zero: the full original
period-corrected derivative frame extends over an explicit collared
$D^2\times T^2$ in one cusp chart times an interval
(Theorem~\ref{h24:thm:cusp-g-zero}). The retained-angle projection of
the same $q_2$ instead gives $H=p_3q_2:X\to S^4$ with
$[H]=\eta_4\eta_5\ne0$ (Theorem~\ref{h24:ret:main}). Its entire
regular inverse image is the original two-torus with the explicitly
calculated normal four-frame and Arf invariant one. Thus $q_2$ is
nonnull and detects a quotient $\pi_6(\mathcal A)/\ker(p_3)_*
\cong\mathbb Z/2$. Smashing $H$ with the specified Hopf map gives
$12\nu_6\ne0$ in $\pi_9(S^6)\cong\mathbb Z/24$, and the
peripheral composite has values $(12,0,12)$ with kernel
$\chi^{-1}(2\mathbb Z/12)$. The full class and exact order of
$q_2$ in its conical target, and its smash powers, are not inferred
from this detected quotient. A completed global higher torus
wrap remains open. The independent arithmetic continuation in
Section~\ref{sec:fifteen-selector-jordan} proves that the fifteen
$p=1201$ projective selector rays span and Jordan-generate the real
three-dimensional symmetric corner, with exact same-residue countercounts.

Section~\ref{esn:sec:shared} adds the explicit even unimodular lattice
$A_5^4\perp D_4\subset N$ of index $72$, with its original $144$ roots,
and the full splitting $O(N)=W(A_5^4D_4)\rtimes L(G)$,
$G\cong GL_2(\mathbb F_3)$
(Theorem~\ref{esn:fa:thm:full-quotient}). This realizes in fixed
coordinate/code markings the abstract Weyl quotient already recorded in
\cite[Table~3]{esn:CDH2014}. The residue six-cycle belongs to the Weyl
group and centralizes this chosen complement
(Theorem~\ref{esn:fa:thm:residue-weyl}); the $D_4$ discriminant action
has kernel $Q_8$ in $G$ and image the full triality $S_3$.
The fixed-shell divisor bijection and finite-sample Fourier map retain
$p,R,P$: $A_3(2521)=0$, $A_3(3361)=6$, with unordered counts $0,3$,
despite equal residues modulo $840$.

The later quaternionic extension is treated separately in
Section~\ref{sec:q444-reconciliation}: the two Lorentz thirteen-slices
coincide as embedded matrix spaces, but the supplied unconstrained scalar
model can leave their trace-free slice. The explicit witness gives
$(u^*b)^{\cdot\cdot}=-2\kappa$ at $u^*b=(u^*b)^{\cdot}=0$.
There are also global interacting invariant families within the original
equations: fixing $W\subset u^\perp$, retaining $b\in W$ and
$\Phi W\subset W$, with tangent velocities, gives configuration dimensions
six and eight and identically zero reaction. A separate fourth-order
escape example has $\rho^{(j)}(0)=0$ for $0\leq j\leq3$ but
$\rho^{(4)}(0)=-12\kappa$. These results distinguish exact invariant
families from instantaneous tangency tests. The families do not exhaust
the invariant motions: Theorem~\ref{q444:moving-family} constructs
periodic moving-kernel solutions outside every fixed-$W$ family.
Their conserved nonzero colour moment gives an explicit Gauss residual
for the zero-connection gauge lift, not an obstruction to every possible
gauge-field lift. The complete scalar vacuum diagram and full metric
Hessian are also proved in Theorems~\ref{q444:vacua} and
\ref{q444:vacuum-hessian}; their gauge-orbit and quantum counterparts are
not inferred from this calculation.
The full equations and the different constrained equations are both given;
the original candidate's complex-structure claims do not follow from this
quaternionic cone calculation.

The full positive fifteen-space also has a split positive observation to
$\operatorname{Herm}_2(\mathbb C)$ with exact image condition
$|\rho|^2\leq\det G$ (Proposition~\ref{q444:full-observation}).
Every permitted $(G,\rho)$ occurs; null $G$ forces $\rho=0$, while
$G=\epsilon I_2$, $\rho=\epsilon/2$ permits positive lifts for every
$\epsilon>0$. For trace-free scalar matrices the distinct spectral-floor
map \eqref{q444:spectral-floor} gives this positive observation with the
original matrix $C=PP^*$ retained. These exact maps do not impose a
positive uniform spectral threshold or identify the coordinates as
gauge-invariant physical momenta.

The complete retained-period metric fibres are now proved in
Theorem~\ref{q444:metric-fibre}: they have dimensions twelve over $H_3$
and seventeen over $G$. Every smooth positive observed trajectory has
an equivariant lift with the original periods unchanged. An explicit
$G=\operatorname{diag}(\alpha_{\rm obs},\beta_{\rm obs}y^{-p})$
approaches a nonzero null block; its relative metric condition number
is asymptotic to
$2\max\{2\alpha_{\rm obs}/3,3\lambda\}y^{p+2}/\beta_{\rm obs}$.
Propositions~\ref{q444:metric-condition}--\ref{q444:rank-loss-metric}
give the exact inverse and bounds. Positivity and period monodromy
alone consequently do not select an observed trajectory.

Theorem~\ref{q444:temporal-no-lift} strengthens the rotating-orbit
obstruction to every homogeneous spatial connection in temporal gauge.
Proposition~\ref{q444:gauss-inertia} retains the full arbitrary-temporal
Gauss equation without inferring the remaining equations from it.
Theorem~\ref{q444:radial-gauge} constructs a different interacting
two-variable family satisfying every field equation.
Finally, Proposition~\ref{q444:vacuum-average} identifies
$\int\det(PgAg^*P^*)\,d\mu
=V_\mu\operatorname{tr}\operatorname{adj}(A)=V_\mu\|b\|^2$
on the positive rank-one boundary and exhibits gauge-inequivalent
scalar minima on which $\|b\|^2$ ranges continuously.
The fixed-mark determinant itself varies within a gauge orbit.

Theorems~\ref{q444:real-no-lift} and \ref{q444:static-complex-no-lift}
now exclude all homogeneous real connections with arbitrary time
dependence, and all complex connections constant in the spatial
constant-matter coordinates, respectively, for the literal rotating orbit.
Theorem~\ref{q444:cx:thm:block} additionally excludes the genuinely
complex rotating-pair class $C,Z_1\in\mathfrak p$, $Z_2\in\mathfrak k$
and its constant $SO(2)$ spatial saturation, where the block splitting
is defined by the fixed matter line $\mathbb Cq$. No unrestricted
pair is assumed to admit that splitting. For arbitrary spatial colour
matrices in the rotating-pair ansatz, Lemma~\ref{q444:complex-gauss-inverse}
proves $g^2\mathcal I+2\mathcal K>0$ whenever $[Z_1,Z_2]\ne0$, a
condition forced by every solution. Theorem~\ref{q444:complex-exact-zero-set}
gives its exact Gauss/frequency inverse and all four residuals
\eqref{q444:complex-full-residuals}. Whether their zero set is empty,
and the existence of still more general time-dependent complex lifts,
remain open; neither is decided by a scalar contraction alone.
The complete scalar equivalence in
Proposition~\ref{q444:complex-full-scalar-equivalence} and the
strict stabilizer signature constraint of
Theorem~\ref{q444:complex-stabilizer-spacelike} hold in every
complex coordinate, not only in the excluded block class.
Proposition~\ref{q444:complex-stress-obstruction} adds the exact
two-by-two spatial stress equation and its sum-of-squares discriminant.

Theorem~\ref{q444:finite-coupling-fibre} now computes the complete
coupling fibre over every unrestricted ordered spatial pair as the
admissible zeros of a nonnegative degree-twenty-eight polynomial.
There are at most fourteen distinct positive $g^2$ values, and its
explicit inverse retains both signs of $\omega$ and every temporal
coordinate. Proposition~\ref{q444:finite-coupling-domain-test} gives
three exact vector tests for an identically empty denominator domain.
Theorem~\ref{q444:generic-projection-obstruction} additionally places
the full spatial projection inside the zero set of a specific nonzero
Sylvester determinant polynomial on all sixteen coordinates. This
set is closed, nowhere dense and measure zero; an explicit full pair
certifies the determinant is nonzero. Proposition~\ref{q444:twocross:prop:exceptional-inclusion} places the
entire eight-dimensional two-cross spatial subspace inside that zero set,
and Theorem~\ref{q444:twocross:thm:exclusion} excludes all full solutions
there. Proposition~\ref{q444:twocross:quantitative} gives a polynomial
residual certificate and a positive uniform lower bound for spatial error
on its bounded scalar--Gauss zero locus at coupling bounded away from zero.
The full scalar--Gauss zero locus is now completely parameterized by
Theorems~\ref{sgd:parametrization} and~\ref{sgd:literal-addon:quotient}:
\[
 \mathcal D=\{(u,v_0,y):u>0,\ d=uv_0-\delta^2-y^2>0,
                       \ c=4v_0-5u>0\},
\]
\[
 \begin{gathered}
 \mathcal S_{\rm sg}\cong
 \bigl(\mathcal D\times\{1,-1\}^2\times U(1)^2\times U(1)\bigr)
 /\langle\iota\rangle,\\
 \iota:(p_+,p_-,\zeta)\mapsto(-p_+,-p_-,-\zeta),
 \qquad \delta=\sqrt{2/3}.
 \end{gathered}
\]
where $\iota$ fixes the other coordinates and $t=72d/c$.
The smooth refinement in Corollary~\ref{sgd:smooth-addon:levels} gives
\[
 \mathcal S_{\rm sg}\cong\{1,-1\}^2\times(0,\infty)
                  \times\mathbb R^2\times\mathbb T^3.
\]
These are four embedded six-dimensional literal components. The original
coupling is the displayed positive coordinate and is a surjective
submersion; every fixed $t=\tau>0$ has exactly four embedded
five-dimensional components $\mathbb R^2\times\mathbb T^3$.
For every prescribed original $\tau=g^2>0$,
\[
 \rho=\frac{\tau+\sqrt{\tau^2+4608\delta^2}}{96\delta},\qquad
 M_\rho=\delta\begin{pmatrix}\rho&1\\1&2\rho\end{pmatrix}
\]
gives $t=\tau$ by Corollary~\ref{sgd:every-coupling}, with every
scalar and Gauss entry zero and both original frequency signs retained.
The two spatial residuals cannot vanish simultaneously in this class.
Existence on the remaining exceptional locus is \Open. None of these finite-ansatz results
is a theorem about arbitrary time-dependent gauge connections or
the interacting continuum quantum spectrum.

The different sector in Proposition~\ref{q444:vk-sector} supplies
exact full-gauge solutions. Its periodic family
(Theorem~\ref{q444:vk-periodic}) has $\omega\to0$ and
$E-V_0=9h_0^2\omega^2+O(\omega^4)$ at fixed spatial volume;
its amplitude tends to $3h_0$, not zero.
Proposition~\ref{q444:vk-wave} proves nonlinear null waves and their
finite-volume energy. The entire transverse gauge spectrum is given in
Theorem~\ref{q444:gauge-spectrum}. The gauge-invariant conserved vector
current in Proposition~\ref{q444:current-mass} detects a branch with
$\mu_{\rm d}^2=(2g^2/3)\|b\|^2$, exactly proportional to the preceding
colour average. Its linear current coefficient vanishes at $b=0$.
These are proved classical equations and observables, not quantum
spectral conclusions.

There is now a separate quantum result for the complete homogeneous
gauge--scalar mechanics, not inferred from those classical families.
Section~\ref{qh:chapter} retains all thirty-eight real configuration
coordinates. The exact form bound \eqref{qh:full-form-bound}, with
its original $\hbar,g^2,\mathcal V_3,\Gamma,m^2,\kappa$, proves compact
resolvent. Theorems~\ref{qh:hq-physical} and \ref{qh:hq-positive}
impose the spatially constant Gauss constraints and prove a unique
invariant ground state; Corollary~\ref{qh:hq-ground} gives
$E^{\rm phys}_1-E_0>0$ for fixed positive volume and couplings.
This is an exact finite-dimensional interacting quantum result,
not a continuum Yang--Mills gap or a uniform infinite-volume estimate.
The classical embedding and the failure of the automatic product
quantum-subsector map are proved explicitly in
Proposition~\ref{qh:hq-embedding} and the following Fourier calculation.
The bosonic commutator-matrix discreteness mechanism is attributed
to Simon's classical work, not claimed as new.
Section~\ref{fq:section} constructs a further exact correspondence for
the retained two-Fourier-coordinate mechanics. Its true positive
quartic ground state $\Omega_q$ gives $Jf=f\Omega_q$, zero Berry
connection and the nonzero Born--Huang term in \eqref{fq:effective}.
The full residual \eqref{fq:intertwining} prevents this moving
subspace from being reducing, but Theorem~\ref{fq:dynamics} bounds
the difference of its physical-time evolutions by the explicit
$|t|\mathcal E(E)$ in \eqref{fq:time-error}. The sufficient regime
\eqref{fq:regime} is proved by min--max and a variational bound,
not an assumed spectral gap. No Gaussian is substituted for the
quartic ground state. The error is of order $\mathcal V_3^{-1/2}$
at fixed $\hbar,\Gamma,L_1,m^2,E,t$; it does not control a simultaneous
limit in which $\omega_*\to0$, nor the full continuum theory.

Theorem~\ref{q444:hermitian-cometric} identifies the original weight
$W$ as a lattice cometric. Complex-structure compatibility removes
exactly its twelve anti-complex-structure fibre directions and gives
$h_3=2H_3^{-1}$. The quotient in
Theorem~\ref{q444:metric-submersion} has $h_2=2G^{-1}$ and intertwines
the Laplacians on the exact pulled-back character subspace.
Proposition~\ref{q444:native-period-spectrum} retains the native
Euclidean metric and all inherited period corrections: $H_3=2I_3$,
$G=2PP^*$, $\lambda_1^{(3)}\sim4\pi^2/y^2$ and
$\lambda_1^{(2)}\sim8\pi^2/y^2$.
The full marked spectrum recovers the lattice cometric even though
the fixed observed matrix does not recover the changing periods.

\begin{enumerate}[leftmargin=1.8em]
\item \Verified\ Let
\(\Delta\simeq(\mathbb Z/3)*(\mathbb Z/4)\) act on
\(\mathfrak h_z\), let \(g_1,g_2\) be its elliptic generators with chosen
fixed points \(z_1,z_2\), and put \(g_0=(g_1g_2)^{-1}\).  For the normalized
period lift \(\tau:\mathfrak h_z\to\mathfrak h\), define
\[
 \varphi:\Delta\longrightarrow\operatorname{PSL}_2(\mathbb Z),
 \qquad \tau\circ g=\varphi(g)\circ\tau.
\]
To avoid confusing the modular transformations with the finite reduced
fibres \(S_j\) below, write
\[
 \mathfrak s_1(w)=\frac{w-1}{w},\qquad
 \mathfrak s_2(w)=-\frac1w,\qquad
 \mathfrak p(w)=w-1.
\]
The finite generators are normalized simultaneously by
\[
 g_1'(z_1)=e^{-2\pi i/3},\qquad g_2'(z_2)=-i,\qquad
 g_0\ \text{parabolic},
\]
and the exact modular deck certificate is
\[
 \varphi(g_1)=\mathfrak s_1,\qquad
 \varphi(g_2)=\mathfrak s_2,\qquad
 \varphi(g_0)=\mathfrak p,\qquad
 \mathfrak s_1\mathfrak s_2\mathfrak p=1.
\]
The sign at \(z_2\) is forced by the simultaneous parabolic normalization;
the quadratic ramification calculation alone gives only
\((g_2'(z_2))^2=-1\).  The cusp normalization proof now records its
covering levels explicitly: it chooses a \(g_0\)-stable component
\(U_0\subset\pi^{-1}(D_0^*)\) and a component
\(K\subset j^{-1}(\{|w|>R\})\) containing \(\tau(U_0)\).  For
\(P_{\rm cusp}:=\varphi(g_0)\), the auxiliary conjugation
\[
 \tau'=\gamma_0^{-1}\tau|_{U_0},\qquad
 P'=\gamma_0^{-1}P_{\rm cusp}\gamma_0,\qquad
 \gamma_0(\infty)=c,
\]
where \(c\) is the rational boundary cusp of \(K\), proves that
\(P'\ne1\) and preserves the selected cusp at infinity, hence
is parabolic.  The proof then discards \(\gamma_0\), returns to the original
\(\mathfrak s_1\)-normalized frame, performs only the allowed final
\(\mathfrak s_1\)-power conjugation, and obtains the displayed modular
certificate.  Only after that return does it define anew the cusp
exponential
\[
 q_{\rm cusp}:=e^{2\pi i\tau}=t_cu_1(t_c),\qquad u_1(0)\ne0;
\]
this \(q_{\rm cusp}\) is distinct from the integral class \(q\) below.
The resulting nondegenerate period family descends, and the middle quotient
\[
 J=(\mathcal T|_{\mathfrak h_z^\circ})/\Delta
 \longrightarrow B_{\rm mid}:=\mathfrak h_z^\circ/\Delta
\]
is a proper holomorphic submersion with compact torus fibres.  The toric
cusp filling, the two finite logarithmic fillings, all three
punctured-overlap biholomorphisms, and the global gluing produce a connected
compact Hausdorff smooth complex threefold and a proper surjective
holomorphic map with connected fibres
\[
 f:X\longrightarrow B\simeq\mathbb P^1.
\]
Its full scheme-theoretic cusp fibre \(X_{p_0}=W\) is reduced, irreducible,
nonnormal, and normal crossings, with local equations
\[
 z_0,\qquad z_0z_1,\qquad z_0z_1z_2,
\]
and normalization \(\nu:E_0\simeq\mathrm{dP}_6\to W\); it is not a
simple-normal-crossings divisor.  For the finite fibres, let
\[
 A_0^{(j)}=(\Lambda\otimes_{\mathbb Z}\mathbb R)/\Lambda\simeq F
\]
be the marked central torus on the covering family, and put
\[
 v_1=\widehat\gamma+2\widehat u-4\widehat w,\qquad
 v_2=-(\widehat\gamma+3\widehat u-3\widehat w),
\]
so that
\[
 S_j=A_0^{(j)}/\langle x\mapsto A_jx+v_j/m_j\rangle,\qquad
 f^*(p_j)=m_jS_j,\qquad (m_1,m_2)=(3,4).
\]
The quotient covering and its codomain are, explicitly,
\[
 \pi_j:A_0^{(j)}\longrightarrow S_j,\qquad x\longmapsto[x].
\]
The radial retraction \(N_j\to S_j\) identifies finite-fibre specialization
exactly with
\[
 \operatorname{sp}_{p_j}^r=\pi_j^*:
 H^r(S_j;\mathbb Z)\longrightarrow H^r(A_0^{(j)};\mathbb Z)
 \xrightarrow[\text{marking}]{\ \sim\ }H^r(F;\mathbb Z),
\]
with no unrecorded conjugacy or index.  This stalk identification is
\cref{lem:finite-filling-stalk-map}; see also
\cref{lem:triangle-generator-orientations,prop:period-input-audit,%
thm:candidate-cusp-geometry,prop:candidate-cusp-overlap,%
lem:middle-period-proper,prop:candidate-finite-fillings,%
thm:candidate-global-geometry}.

\item \Verified\ The monodromy, cusp invariant lattices, global
coinvariants, finite-filling stalk maps, terminal group presentation, and
terminal degree-two lattice are proved over \(\mathbb Z\) in
\cref{prop:monodromy-certificate,thm:cusp-lattices,thm:coinvariants,%
prop:terminal-group,prop:terminal-degree-two,prop:terminal-all-degree,prop:finite-indices}.

\item \Verified\ The intrinsic integral nearby-cycle specialization at the
cusp is injective with saturated image
\[
 \operatorname{im}(\operatorname{sp}^q)
   =H^q(F;\mathbb Z)^{T_0}\qquad(0\le q\le4),
\]
where the degree-two invariant class is
$q:=uw+6\gamma\delta\in\bigwedge^2V=H^2(F;\mathbb Z)$,
and the finite stalks and parabolic sheaf calculation give the complete
global Leray page.  The canonical truncation triangle identifies the
explicit Cech discrepancy cocycle with Leray \(d_2\), and multiplicativity
then gives
\[
 d_2(12\gamma)=\omega,\qquad
 d_2(2q)=\omega[\delta],\qquad
 d_2(2\gamma uw)=\omega[\gamma\delta].
\]
Consequently
\[
 H^1(X;\mathbb Z)=H^2(X;\mathbb Z)=H^3(X;\mathbb Z)=0,
 \qquad H_2(X;\mathbb Z)=0;
\]
see
\cref{thm:p0-integral-specialization,prop:leray-e2,lem:first-transgression,lem:higher-transgressions,thm:h2-vanishing}.

\item \Verified\ The candidate's independent toric-collapse construction is
continuous across every stratum and descends equivariantly to a strong
deformation retraction $r:N'_0\to W$.  Its fixed-fibre restriction
$c=r|_F$ has the exact primitive integral sequence
\[
0\longrightarrow\operatorname{im}(\textstyle\bigwedge^qM_0-I)
\longrightarrow\textstyle\bigwedge^q\Lambda
\xrightarrow{c_*}H_q(W;\mathbb Z)\longrightarrow0,
\]
and proper-base-change functoriality gives, in every degree,
\[
 c^*=i_F^*r^*=\operatorname{sp}^q.
\]
Thus the former C1--C3 obligations are closed; see
\cref{prop:verified-cellular-collapse,prop:verified-collapse-homology,%
thm:collapse-equals-specialization}.

\item \Verified\ With zero-section whiskers and the fixed clockwise
meridians, the actual overlap inclusions give
\[
 \pi_1(X)=
 \langle c,x,y\mid c\text{ central},\ xy=1,\ x^3=c,\ y^4=c^{-1}\rangle
 \simeq1.
\]
Together with the preceding integral cohomology this proves
\[
 X\cong_{\mathrm{diff}}S^6,\qquad
 \langle c_3(T_X),[X]\rangle=2;
\]
see \cref{prop:geometric-pi1,cor:candidate-sphere}.

\item \Verified\ The very-general fibre has
\[
 \eta:=q=u\wedge w+6\gamma\wedge\delta,\qquad
 \operatorname{NS}(F_b)=\mathbb Z\eta,\qquad a(F_b)=0.
\]
The square is the oriented wedge identity
\[
 \eta\wedge\eta
 =12\,\gamma\wedge u\wedge w\wedge\delta,
\]
so, for the orientation determined by the ordered marking
\((\gamma,u,w,\delta)\) and for the complex orientation respectively,
\[
 \int_{(F_b,\operatorname{or}_{\rm mark})}\eta\wedge\eta=12,
 \qquad
 \int_{(F_b,\operatorname{or}_{\mathbb C})}\eta\wedge\eta=-12.
\]
The repaired meromorphic-descent argument gives
\[
 \mathcal M(X)=f^*\mathbb C(t),\qquad a(X)=1,
\]
with \(f\) the algebraic reduction; see
\cref{prop:very-general-torus-audit,thm:algebraic-dimension-audit}.

\item \Refuted\ The p.~692 sheaf-theoretic inference of
Campana--Demailly--Peternell is false in the stated generality.  If
\(
\mathscr K_S=\ker(\Omega_X^1|_S\to\widetilde\Omega_S^1)
\), then
\[
0\longrightarrow N^*_{S/X}\longrightarrow\mathscr K_S
\longrightarrow\operatorname{tors}\Omega_S^1\longrightarrow0,
\]
and the two printed vanishings do not control sections mapping to the
conductor-supported quotient.  The proper nodal-cubic-times-
\(\mathbb P^1\) family in \cref{prop:nodal-counterexample} is an explicit
counterexample to that inference.

\item \Verified\ The conductor section, nonnormal Cartier-divisor duality,
flatness, and top-degree analytic base change are complete in
\cref{prop:conductor,thm:cartier-duality,prop:cartier-duality-absolute,%
lem:curve-flatness,prop:analytic-base-change,cor:surface-top-base-change}.
On the cusp model the invariant nowhere-vanishing form
\[
 \Omega=\frac{dx_1\wedge dx_2\wedge dt}{x_1x_2}
\]
first proves intrinsically that
\(\omega_{N_0}\simeq\mathcal O_{N_0}\).  After the global gluing, the open
embedding \(N_0\hookrightarrow X\) gives the exact restriction chain
\[
 \omega_X|_{N_0}\simeq\omega_{N_0}\simeq\mathcal O_{N_0}.
\]
Let \(\nu:E_0\simeq\mathrm{dP}_6\to W\) be the normalization.
Consequently, for every \(L\in\operatorname{Pic}^0(X)\),
\[
 \nu^*((L^\vee\otimes\omega_X)|_W)
 \simeq\mathcal O_{\mathrm{dP}_6},
\]
and the conductor--duality--base-change chain gives
\[
 R^2f_*(T_X\otimes L)\ne0
\]
without using global topology; see
\cref{thm:candidate-cusp-geometry,thm:candidate-global-geometry,%
prop:pic0-pullback,cor:candidate-pic0-route}.

\item \Verified\ For arbitrary line bundles the exact trivial-pullback
obstruction is
\[
 r=\nu^*i^*:\operatorname{Pic}(X)\longrightarrow
 \operatorname{Pic}(\mathrm{dP}_6).
\]
The conductor argument applies uniformly to all line bundles exactly when
$r=0$; see \cref{prop:all-pic-restriction}.  For the constructed candidate,
the Leray edge image of every $c_1(L)$ lies in
$E_\infty^{0,2}\subseteq\ker d_2^{0,2}$, while
$d_2^{0,2}(2q)=\omega[\delta]$ is an isomorphism.  Proper-base-change
naturality therefore gives $c_1(L|_W)=0$, and injectivity of $c_1$ on
$\operatorname{Pic}(\mathrm{dP}_6)$ forces $r=0$; see
\cref{prop:leray-edge-restriction}.  Combining this with the analytic inputs
gives
\[
 R^2f_*(T_X\otimes L)\ne0
 \qquad\text{for every }L\in\operatorname{Pic}(X).
\]
Thus the topology-free Pic$^0$ route and the all-Pic route are compatible
rather than competing statuses.  There are now three all-Pic certificates.
The first is the preceding Leray-edge proof.  The second uses no global
Leray differential: the order-four specialization lattice writes
$c_1(L|_F)=2nq$, and the pair-of-pants crossed defect has boundary values
\[
 \bar d_1=4n[\delta],\qquad
 \bar d_2=-3n[\delta],\qquad
 \bar d_0=0.
\]
The exact boundary relation forces $n=0$, and local integral cusp
specialization then forces $r(L)=0$; see
\cref{lem:picard-crossed-defect,lem:picard-finite-norm,%
lem:picard-cusp-defect,prop:direct-boundary-picard}.  This second proof is
independent of the global transgressions, but it still uses one local cusp
bounding chain and the local specialization theorem.

\item \Verified\ The local cusp Picard map itself is nonzero.  The explicit
equivariant Cartier divisor
\[
 \mathcal D_Q=\sum_{v\in\mathbb Z^2}(v_1^2-6v_2^2)E_v
\]
descends to $\mathcal L_Q\in\operatorname{Pic}(N_0)$ with
\[
 c_1(\mathcal L_Q|_{F_t})=-2q,\qquad
 \bigl(\deg_{C_i}\nu^*(\mathcal L_Q|_W)\bigr)_{i=1}^6
 =(-12,2,0,-12,2,0).
\]
Thus neither the cusp geometry nor analytic relative-Picard
representability alone can force normalized central restriction to vanish;
see \cref{prop:local-cusp-picard-obstruction}.

\item \Verified\ After the global coinvariant calculation writes
\[
 E:=c_1(L|_F)=kq\in\operatorname{Alt}^2(\Lambda,\mathbb Z)
       =\bigwedge^2V,
\]
the third all-Pic certificate completes the argument by holomorphic factors
and toric support functions.  Put
\(V_{\mathbb R}=\operatorname{Hom}(\Lambda,\mathbb R)\), let
\(P\subset B^\circ\) be the compact pair-of-pants core with clockwise
boundary meridians \(\rho_1,\rho_2,\rho_0\), and let
\(d:\pi_1(P)\to V_{\mathbb R}\) be the factor-of-automorphy crossed
homomorphism, put \(d_i=d(\rho_i)\), and let a bar denote passage to
\(V_G\otimes\mathbb R=\mathbb R[\delta]\).  Logarithms of the actual
\(\mathcal O^*\)-factors give
\[
 \bar d_1+\bar d_2+\bar d_0=0,\qquad
 \bar d_1=2k[\delta],\quad
 \bar d_2=-\tfrac32k[\delta],\quad
 \bar d_0=0.
\]
The finite norm identities now include the fully typed factor-extension
repair (5.17-A14c)--(5.17-A15), rather than identifying maps living on
different covering levels.  Put
\[
 \widetilde J_j=\Delta_j\times\mathbb C^2,
 \qquad r_j:\widetilde J_j\longrightarrow J'_j,
\]
where \(r_j\) is the lattice quotient with deck translations
\(\tau_\lambda\).  On this vector cover, and only there, define
\[
 Q_{j,\log}=\widetilde g_j^{\,v_j/m_j},\qquad
 Q_{j,\mathrm{lin}}=\widetilde g_j^{\,0}.
\]
Their exact relations are
\[
 Q_{j,\log}\tau_\lambda Q_{j,\log}^{-1}=\tau_{A_j\lambda},
 \quad (Q_{j,\log})^{m_j}=\tau_{v_j},
 \qquad
 Q_{j,\mathrm{lin}}\tau_\lambda Q_{j,\mathrm{lin}}^{-1}
 =\tau_{A_j\lambda},
 \quad (Q_{j,\mathrm{lin}})^{m_j}=1.
\]
Let \(q_{j,\log},q_{j,\mathrm{lin}}\) be their descended order-\(m_j\)
automorphisms of \(J'_j\), put
\(J_j^{\prime *}=J'_j|_{\Delta_j^*}\), and define
\[
 p_{j,\log}:J_j^{\prime *}\to
 N_j^*:=J_j^{\prime *}/\langle q_{j,\log}\rangle,
 \qquad
 p_{j,\mathrm{lin}}:J_j^{\prime *}\to
 J^*:=J_j^{\prime *}/\langle q_{j,\mathrm{lin}}\rangle.
\]
The branch-independent fibre translation
\(h_j:J_j^{\prime *}\to J_j^{\prime *}\) and
its quotient biholomorphism \(\bar h_j:N_j^*\to J^*\) satisfy
\[
 h_jq_{j,\log}h_j^{-1}=q_{j,\mathrm{lin}},
 \qquad p_{j,\mathrm{lin}}\circ h_j=\bar h_j\circ p_{j,\log},
\]
and the gluing gives the typed line-bundle comparison
\[
 p_{j,\log}^*(L|_{N_j^*})
 \simeq h_j^*p_{j,\mathrm{lin}}^*(L|_{J^*}).
\]

For the filled factor extension return to the vector cover.  With
\[
 P_{j,\log}:J'_j\to N_j,\qquad
 \pi_j^{\rm vec}:=P_{j,\log}\circ r_j:\widetilde J_j\to N_j,
 \qquad
 \Gamma_j^{\log}
 =\langle\tau_\lambda,Q_{j,\log}:\lambda\in\Lambda\rangle,
\]
where \(P_{j,\log}\) is the full quotient and
\(P_{j,\log}|_{J_j^{\prime *}}=p_{j,\log}\), pull back
\(L|_{N_j}\) by \(\pi_j^{\rm vec}\), push its integral factor
extension out along \(\mathbb Z\hookrightarrow\mathbb R\), and choose
normalized vertical lifts \(U_\lambda^{(j)}\).  For
\(\alpha\in\Gamma_j^{\log}\), define
\(A_\alpha\in\operatorname{Aut}(\Lambda)\) and
\(T_\alpha\in\operatorname{Aut}(V_{\mathbb R})\) by
\[
 \alpha\tau_\lambda\alpha^{-1}=\tau_{A_\alpha\lambda},
 \qquad T_\alpha=(A_\alpha^{-1})^{\mathsf t}.
\]
For a factor-extension lift \(\widetilde G_\alpha\), define
\(\widetilde d_j(\alpha)\in V_{\mathbb R}\) by
\[
 \widetilde G_\alpha U_{A_\alpha^{-1}\lambda}^{(j)}
 \widetilde G_\alpha^{-1}
 =\mathbf c\!\left(\langle\widetilde d_j(\alpha),\lambda\rangle\right)
 U_\lambda^{(j)}.
\]
The normalized product gives
\[
 \widetilde d_j(\alpha\beta)
 =\widetilde d_j(\alpha)+T_\alpha\widetilde d_j(\beta).
\]
Compatible basepoints in the quotient square identify the clockwise
boundary lift successively with \(q_{j,\mathrm{lin}}\), \(q_{j,\log}\), and
finally the vector-cover deck transformation \(Q_{j,\log}\).  The induced
extension isomorphism then proves, for some \(a_j\in V_{\mathbb R}\),
\[
 \widetilde d_j(Q_{j,\log})=d_j+(T_j-I)a_j,
 \qquad
 \widetilde d_j(\tau_v)=\iota_vE,\qquad
 (\iota_vE)(\lambda)=E(v,\lambda).
\]
Consequently, with \(N_{T_j}=I+T_j+\cdots+T_j^{m_j-1}\),
\[
 N_{T_j}d_j
 =N_{T_j}\widetilde d_j(Q_{j,\log})
 =\widetilde d_j((Q_{j,\log})^{m_j})
 =\widetilde d_j(\tau_{v_j})=\iota_{v_j}E,
\]
because \(N_{T_j}(T_j-I)=0\).  This proves the finite boundary values,
including their signs, inside the actual filled logarithmic deck group and
closes the factor-extension gap; the certificate therefore remains
{\Verified}.  Changes of line-bundle frame, central lift, normalized
vertical system, or logarithm branch alter only the displayed principal
character term and leave \(h_j,\bar h_j\), the quotient square, and
(5.17-A15) unchanged.
The cusp equality follows from the argument principle in a Stein chart, not
from a Chern evaluation on a bounding chain; hence \(k=0\).  For adjacent
rays in a maximal cone
\(\sigma=\{v,w,u\}\), the actual toric-open inclusions are
\[
 U_{\rho_v}=D(z_wz_u)\subset U_\sigma,\qquad
 U_{\rho_w}=D(z_vz_u)\subset U_\sigma,\qquad
 U_{\rho_v}\cap U_{\rho_w}=U_0.
\]
The contractible Stein comparison chart
\[
 S_\sigma=\{(z_v,z_w,z_u)\in\mathbb C^3:
               |z_vz_wz_u|<\epsilon\}
\]
has \(\operatorname{Pic}(S_\sigma)=0\), so its frames give zero horizontal
winding along the connected one-skeleton.  For the support-function Cartier
divisor \(D_a\), the exact deck-translate identity
\[
 \Psi_\lambda^*D_a-D_a=\operatorname{div}(\chi^{m_\lambda})
\]
shows that multiplication by \(\chi^{m_\lambda}\) extends from the dense
torus to an isomorphism
\[
 \mu^+_\lambda:\Psi_\lambda^*\mathcal O(D_a)
 \xrightarrow{\ \sim\ }\mathcal O(D_a).
\]
The map on the negative divisor required by the descent tensor product is
the inverse-dual
\[
 \mu^-_\lambda=((\mu^+_\lambda)^{-1})^\vee:
 \Psi_\lambda^*\mathcal O(-D_a)\longrightarrow\mathcal O(-D_a),\qquad
 s\longmapsto\chi^{-m_\lambda}s.
\]
Here the cusp shear is the typed lattice isomorphism
\[
 \Lambda_{\rm tor}:=\ker(M_0-I)=\operatorname{im}(M_0-I)
   =\mathbb Z\langle\widehat w,\widehat\delta\rangle,
 \qquad \overline\Lambda:=\Lambda/\Lambda_{\rm tor},
\]
\[
 B_0:\overline\Lambda\xrightarrow{\ \sim\ }\Lambda_{\rm tor},
 \qquad
 B_0\overline{\widehat\gamma}=-\widehat\delta,
 \qquad B_0\overline{\widehat u}=\widehat w.
\]
Thus its matrix is
\(\left(\begin{smallmatrix}0&1\\-1&0\end{smallmatrix}\right)\)
from the ordered basis
\((\overline{\widehat\gamma},\overline{\widehat u})\) to
\((\widehat w,\widehat\delta)\); in particular \(B_0\) is a map between
two named lattices, not an untyped endomorphism, and \(\det B_0=1\).
For \(M=L|_{N_0}\), if \((d_1,d_2,d_3)\) are the independent degrees of
\(\nu^*(M|_W)\), then
\[
 \bigl(\deg_{C_i}\nu^*(M|_W)\bigr)_{i=1}^6
 =(d_1,d_2,d_3,d_1,d_2,d_3),
\]
and their support function has polarization
\[
 B\in\operatorname{Sym}^2(\Lambda_{\rm tor}^*),\qquad
 [B]_{(\widehat w,\widehat\delta)}
 =\begin{pmatrix}d_2+d_3&d_3\\d_3&d_1+d_3\end{pmatrix}.
\]
For the Appell--Humbert form
\(E_M:=c_1(M|_F)\in\operatorname{Alt}^2(\Lambda,\mathbb Z)\), one has
\[
 E_M((\lambda,0),(0,m))=-B(B_0\lambda,m)
 \quad
 (\lambda\in\overline\Lambda,\ m\in\Lambda_{\rm tor}).
\]
The already proved antecedent \(k=0\) gives
\(E_M=c_1(L|_F)=0\); only then does the unimodularity of \(B_0\)
force \(B=0\), then \(d_1=d_2=d_3=0\), and hence
\(\nu^*i^*L\simeq\mathcal O_{\mathrm{dP}_6}\).  Beyond the preceding
coinvariant input, this proof uses neither a global integral Leray
transgression, local integral nearby-cycle specialization, nor Chern-class
evaluation on a cusp bounding chain; see
\cref{prop:analytic-ah-boundary,lem:analytic-support-specialization,%
thm:analytic-all-picard}.

\item \Refuted\ The separate split-extension calculation imported from
CDP98 replaces monodromy coinvariants by the entire fibre homology.  The
correct formula is
\[
 (K\rtimes Q)^{\mathrm{ab}}
 \simeq Q^{\mathrm{ab}}\oplus(K^{\mathrm{ab}})_Q.
\]

\item \Verified\ A corrected replacement theorem now survives both printed
defects.  If one \(L_0\in\operatorname{Pic}(X)\) satisfies
\[
 H^2(X,T_X\otimes L_0)=0,
\]
then generic endpoint vanishing, semicontinuity, and
Hirzebruch--Riemann--Roch give
\[
 H^i(X,T_X\otimes M)=0\ (i\ne1),\qquad
 \chi(X,T_X\otimes M)=\tfrac12c_3(X)\le0
\]
for general \(M\).  It is enough that
\(R^2f_*(T_X\otimes L)=0\) for one \(L\).  For reduced full fibres, the
invalid p.~692 shortcut is replaced exactly by injectivity of
\[
 \delta_c:H^0(W_c,\operatorname{tors}\Omega^1_{W_c}\otimes A_c)
 \longrightarrow H^1(W_c,N^*_{W_c/X}\otimes A_c)
\]
together with the two endpoint vanishings.  See
\cref{thm:corrected-cdp-global,cor:corrected-cdp-relative,%
 thm:corrected-cdp-fibrewise}.  The constructed candidate fails the global and
 relative conditions for separate reasons.  Since \(c_3(X)=2\), the
 contrapositive of \cref{thm:corrected-cdp-global} gives
 \(H^2(X,T_X\otimes L_0)\ne0\) for every \(L_0\).  Independently, its
 conductor makes \(R^2f_*(T_X\otimes L)\ne0\) for every \(L\), so the stronger
 relative sufficient condition also fails.  The original CDP20 theorem
 therefore remains refuted.

\item \Verified\ The constructed \(X\) satisfies the stated hypotheses of
CDP20 Theorem~2.2 but has both \(c_3(X)=2>0\) and \(a(X)=1\).  Thus it is a
counterexample to that theorem and its ensuing \(S^6\) corollary as stated.
This theorem-level conclusion uses the complete construction, topology, and
algebraic-reduction proofs above; it is not inferred merely from the two
printed proof-route defects.

\item \Verified\ The candidate's Proposition~10.8 conclusion holds.  The
workbench proves
\[
R^1f_*\mathcal O_X\simeq\mathcal O\oplus\mathcal O(-1),\qquad
\operatorname{Pic}(X)\simeq\mathbb C,
\]
and the finite-filling Hodge-line calculation now distinguishes the descended
canonical section from the direct-image frame.  In the notation of
\cref{lem:hodge-line-audit},
\[
 \operatorname{div}(\overline\sigma_j)=k_jS_j,
 \qquad
 f_{j*}\omega_{N_j/D_j}
 =\mathcal O_{D_j}\cdot s_j^{k_j}\vartheta'_j,
 \qquad k_j=m_j-1-a_j.
\]
Thus \(\overline\sigma_2\) has divisor \(2S_2\) and is not itself a
nowhere-vanishing frame, while the displayed invariant module still has the
claimed generator, a unit times \(s_j^{-a_j}e\), over the punctured cover.
This removes a local wording defect without changing the Hodge line or any
subsequent direct-image calculation.  The workbench also proves
degree-one base change at every scheme-theoretic fibre, and
\(\operatorname{ord}(\mathcal O_X(S_j)|_{S_j})=m_j\).  The smooth-fibre
period condition is exactly \(c\tau+d\mu-b=0\); hence every parameter
outside \(\mathbb R\) and a countable set gives one line bundle non-torsion
on every full fibre and on both \(S_j\).  The cusp extension is proved
intrinsically through nearby cycles and does not depend on the separately
reconstructed cellular collapse; see
\cref{thm:structure-direct-images-audit,thm:uniform-nontorsion-audit}.

\item \Verified\ The alternate Mayer--Vietoris route is now closed.  The
finite-band inclusions are exactly
$\alpha_q=(\pi_{1*},-\pi_{2*})$ in named primitive bases; the second cusp
attachment has degree-one determinant $-p$ and reduced degree-two block
$\operatorname{diag}(-p,1)$; and the dual obstruction squares commute with
the three verified Leray transgressions.  For $p=-1$ both global blocks are
 unimodular, giving an independent integral sphere certificate; see
 \cref{prop:actual-finite-band-maps,prop:cusp-mv-blocks,%
 thm:mv-leray-comparison,cor:candidate-mv-sphere}.

\item \Verified\ The Hodge diamond and Fr\"olicher sequence are now exact.
Here are the typed logarithmic objects used in the calculation:
\[
 B=\mathbb P^1,\qquad \Sigma=p_0+p_1+p_2,\qquad
 D=f^{-1}(\Sigma)_{\rm red}=W+S_1+S_2,
\]
\[
 \mathcal L=\Omega_B^1(\log\Sigma)\simeq\mathcal O_B(1),\qquad
 \mathcal A=\Omega_X^1(\log D),\qquad
 \mathcal E=\mathcal A/f^*\mathcal L,\qquad
 \mathcal G=R^1f_*\mathcal O_X.
\]
They occur in the locally free logarithmic conormal sequence
\[
 0\longrightarrow f^*\mathcal L\longrightarrow\mathcal A
 \longrightarrow\mathcal E\longrightarrow0,
\]
whose connecting maps have the exact types
\[
 \kappa_0:f_*\mathcal E\longrightarrow\mathcal L\otimes\mathcal G,
 \qquad
 \kappa_1:R^1f_*\mathcal E\longrightarrow
 \mathcal L\otimes R^2f_*\mathcal O_X\simeq\mathcal O_B.
\]
The logarithmic relative direct images are
\[
 f_*\mathcal E\simeq\mathcal O\oplus\mathcal O(-1),\qquad
 R^1f_*\mathcal E\simeq
 \mathcal O^{\oplus2}\oplus\mathcal O(-1)^{\oplus2},\qquad
 R^2f_*\mathcal E\simeq\mathcal O\oplus\mathcal O(-1),
\]
and are locally free and compatible with base change.

On a finite cover \(t=s^m\), with \((m,a)=(3,2)\) or \((4,1)\), choose
equivariant relative-cotangent frames \((e_0,e_1)\) and structure-sheaf
frames \((b_0,b_1)\) with semilinear characters
\((1,\zeta^a)\) and \((1,\zeta^{-a})\), respectively.  The actual
\(\mathbb C\{t\}\)-bases are
\[
\begin{aligned}
 f_*\mathcal E:&\quad e_0,\ s^{m-a}e_1,\\
 R^1f_*\mathcal E:&\quad
 e_0b_0,\ s^ae_0b_1,\ s^{m-a}e_1b_0,\ e_1b_1,\\
 R^2f_*\mathcal E:&\quad
 s^ae_0(b_0\wedge b_1),\ e_1(b_0\wedge b_1).
\end{aligned}
\]
Let \(\ell=dt/t\), and let
\(D_{\mathrm{KS}}=(d_{rs})_{0\le r,s\le1}\) be the marked
Kodaira--Spencer matrix after its common factor \(ds\) is removed.  Put
\[
 A=\frac{s}{m}d_{00},\qquad B_1=\frac{s^{1-a}}m d_{01},\qquad
 C=\frac{s^{1-a}}m d_{10},\qquad
 D_0=\frac{s^{m+1-2a}}m d_{11}.
\]
In source basis \((e_0,\allowbreak s^{m-a}e_1)\), target basis
\((b_0,\allowbreak s^ab_1)\ell\), ordered degree-one source basis
\((e_0b_0,\allowbreak s^ae_0b_1,\allowbreak
s^{m-a}e_1b_0,\allowbreak e_1b_1)\), and target generator
\(s^a(b_0\wedge b_1)\ell\), the coarse maps are exactly
\[
 [\kappa_0]=\begin{pmatrix}A&tB_1\\C&D_0\end{pmatrix},\qquad
 [\kappa_1]=(-C,A,-D_0,B_1).
\]
The finite residue is the typed map
\[
 \operatorname{res}_{1,j}:
 (\mathcal L\otimes\mathcal G)|_{p_j}
 \longrightarrow H^1(S_j,\mathcal O_{S_j}),\qquad
 \operatorname{res}_{1,j}\bigl((xb_0+ys^ab_1)\ell\bigr)=m_jx[b_0]
 \quad(x,y\in\mathbb C),
\]
and the derived residue square gives
\[
 \operatorname{res}_{1,j}\kappa_0=0,\qquad
 \operatorname{im}(\kappa_0\bmod t)
 =\ker\operatorname{res}_{1,j}=\mathbb C\,s^ab_1.
\]

At the cusp put
\[
 \mathcal Q_{\rm cusp}:=
 \nu_*\mathcal O_{\mathrm{dP}_6}/\mathcal O_W,
\]
and let
\(\rho:\mathcal E|_W\to\mathcal Q_{\rm cusp}\) be the relative
normalization--residue quotient map.  Retain the connecting isomorphisms
\[
 \delta_0:H^0(\mathcal Q_{\rm cusp})\xrightarrow{\ \sim\ }
 H^1(W,\mathcal O_W),\qquad
 \delta_1:H^1(\mathcal Q_{\rm cusp})\xrightarrow{\ \sim\ }
 H^2(W,\mathcal O_W).
\]
For
\[
 q_1=(-1,0,1),\qquad q_2=(0,-1,-1),\qquad
 \eta_i=d\log x_i,\qquad \alpha_i=\delta_0(q_i),
\]
let \(\theta\in H^1(\mathcal Q_{\rm cusp})\) be the unique class satisfying
\[
 \delta_1(\theta)=-\alpha_1\smile\alpha_2.
\]
The standard
\v Cech differential
\((\delta c)_{i_0\ldots i_{r+1}}
=\sum_k(-1)^kc_{i_0\ldots\widehat{i_k}\ldots i_{r+1}}\)
fixes the sign
\[
 \delta_1\bigl(\rho_*(\alpha\otimes\eta_i)\bigr)
 =-\alpha\smile\alpha_i.
\]
Consequently the two cusp maps, with their domains and codomains, are
\[
 H^0(\mathcal E|_W)\xrightarrow{\ \sim\ }
 H^0(\mathcal Q_{\rm cusp}),
\]
\[
 H^1(\mathcal E|_W)=
 \mathbb C\langle\alpha_1\eta_1,\alpha_1\eta_2,
 \alpha_2\eta_1,\alpha_2\eta_2\rangle
 \xrightarrow{\ (0,1,-1,0)\ }
 H^1(\mathcal Q_{\rm cusp})=\mathbb C\theta;
\]
the second is surjective.

Finally, the invariant section
\(\xi=\partial/\partial\zeta_2\), nowhere zero in the locally free relative
sheaf $T_{X/B}=\mathcal E^\vee$, gives the split evaluation sequence
\[
 0\longrightarrow\mathcal O(-1)\longrightarrow f_*\mathcal E
 \xrightarrow{\operatorname{ev}_\xi}\mathcal O\longrightarrow0.
\]
There is a unique lift \(h\in H^0(B,f_*\mathcal E)\) of
\(1\in H^0(B,\mathcal O)\), because
\(H^0(B,\mathcal O(-1))=0\).  On the marked cover it has the uniquely
determined form
\[
 h=(a(z),1)\binom{d\zeta_1}{d\zeta_2},\qquad a(z_2)=\frac12.
\]
Use the genuine invariant target frame
\[
 \omega_2=s\,u_2(s)(\gamma\wedge u)\ell,\qquad u_2(0)=1.
\]
The determinant and final global scalar are
\[
 \operatorname{div}(\det\kappa_0)=p_1+p_2,
 \qquad
 \kappa_1(h\otimes\gamma)|_{p_2}
 =-\frac14\dot\mu(z_2)\,\omega_2|_{p_2}\ne0
 \quad\text{in }(\mathcal L\otimes R^2f_*\mathcal O_X)|_{p_2}.
\]
This uses
\(\dot\tau(z_2)=0\) and \(K(z_2)=6\dot\mu(z_2)^2\ne0\).

Thus, with \(\mathcal K=\ker\kappa_1\),
\[
 \mathcal K\simeq\mathcal O\oplus\mathcal O(-1)^{\oplus2},
\]
and the torsion-free absolute cotangent direct images are
\[
 R^0f_*\Omega_X^1\simeq\mathcal O(-2),\quad
 R^1f_*\Omega_X^1\simeq\mathcal K,
 \quad R^2f_*\Omega_X^1\simeq\mathcal O\oplus\mathcal O(-1).
\]
Therefore
\[
 h^{1,2}=h^{2,1}=1,\qquad h^{1,1}=h^{2,2}=2.
\]
In full, with row index \(q\) and column index \(p\), the Hodge table is
\[
\begin{array}{c|cccc}
q\backslash p&0&1&2&3\\ \hline
3&0&0&0&1\\
2&0&1&2&1\\
1&1&2&1&0\\
0&1&0&0&0.
\end{array}
\]
The two middle \(d_1\)-rows are exact
\(0\to\mathbb C\to\mathbb C^2\to\mathbb C\to0\), and
\[
 \dim_{\mathbb C}E_1=10,\qquad
 \dim_{\mathbb C}E_2=2,\qquad E_2=E_\infty.
\]
See
\cref{prop:log-relative-direct-images,prop:log-ks-divisor,%
prop:absolute-cotangent-direct-images,thm:h12-scalar,cor:exact-hodge-diamond}.

\item \Refuted\ Remark~1.3's assertion that the relative differential
sheaves have torsion along all three singular fibres is false already for
$\Omega^1_{X/B}$.  The $t=xy$ and $t=xyz$ modules inject into the free
logarithmic relative module, so $\Omega^1_{X/B}$ is torsion-free but
non-locally-free near the cusp.  At a point of the smooth finite cover
\(q_j:J'_j\to N_j\), with
\(R_j=\mathcal O_{J'_j,x}\simeq\mathbb C\{s,y_1,y_2\}\), and
\(t=s^{m_j}\), the exact module is
\[
 (q_j^*\Omega^1_{N_j/B})_x
 \simeq\Omega^1_{J'_j/B,x}
 \simeq R_j\,dy_1\oplus R_j\,dy_2
 \oplus\bigl(R_j/(s^{m_j-1})\bigr)ds,
\]
and hence
\[
 (q_j^*\operatorname{tors}\Omega^1_{N_j/B})_x
 =\bigl(R_j/(s^{m_j-1})\bigr)ds.
\]
Faithful flatness descends this module and its support.  Thus the torsion
support is exactly the two multiple fibres and is absent at the cusp; see
the local calculation (11.17f).

\item \Verified\ The Bott--Chern and Aeppli dimensions are now exact as
well.  McHugh's threefold formula gives
\[
h^{1,2}_{BC}=h^{1,2}+\dim E_2^{3,1}-k^{2,0}=1+0-0=1,
\]
because the de Rham image of the pure doubly closed forms surjects onto
$K^{2,0}$, so $k^{2,0}\le b_2=0$.  The complete Bott--Chern rows
$q=3,2,1,0$ are
\[
(0,0,1,1),\quad(0,1,2,1),\quad(0,2,1,0),\quad(1,0,0,0),
\]
and Aeppli duality gives
\[
(0,0,0,1),\quad(0,1,2,0),\quad(1,2,1,0),\quad(1,1,0,0).
\]
See \cref{thm:exact-bott-chern-aeppli}; no dimension remains open.

\item \Verified\ Global representatives and the complete Bott--Chern
product structure are now known.  The two Dolbeault inputs must not be
conflated.  The distinguished class
\[
 \mathfrak g\in H^{0,1}_{\bar\partial}(X)=H^1(X,\mathcal O_X)
\]
is the unique coherent-Leray class mapping to the invariant section
\(\gamma^{0,1}\in H^0(B,R^1f_*\mathcal O_X)\).  In contrast, the marked
integral fibre class is
\[
 q=\eta=u\wedge w+6\gamma\wedge\delta\in H^2(F_b;\mathbb Z),
\]
whose holomorphic extension is a section
\(\boldsymbol\eta\in H^0(B,R^1f_*\Omega_X^1)\).  Choose a coherent-Leray
lift
\[
 \mathfrak q\in H^{1,1}_{\bar\partial}(X)
 \quad\text{with edge value}\quad \boldsymbol\eta.
\]
This lift is not distinguished: its exact ambiguity is
\[
 \mathfrak q\longmapsto\mathfrak q+c\,d_1\mathfrak g,
 \qquad c\in\mathbb C,
\]
because the kernel of the edge map is
\(H^1(B,f_*\Omega_X^1)=\mathbb C\,d_1\mathfrak g\).

Choose global \(\bar\partial\)-closed smooth representatives
\[
 G\in\mathcal A^{0,1}(X),\quad [G]_{\bar\partial}=\mathfrak g,
 \qquad
 Q\in\mathcal A^{1,1}(X),\quad [Q]_{\bar\partial}=\mathfrak q.
\]
For a marked smooth fibre put
\[
 y=\operatorname{Im}\tau,\qquad m_\mu=\operatorname{Im}\mu,\qquad
 r=\operatorname{Im}\beta,\qquad
 \Delta_F=r-\frac{6m_\mu^2}{y}<0,
\]
and define the translation-invariant forms
\[
 \Gamma=\frac{i}{2\Delta_F}
 \left(d\bar\zeta_2-\frac{m_\mu}{y}d\bar\zeta_1\right),
 \qquad
 U=\frac{i}{2y\Delta_F}
 \left(r\,d\bar\zeta_1-6m_\mu\,d\bar\zeta_2\right),
\]
\[
 Q_F=-d\zeta_1\wedge U-6d\zeta_2\wedge\Gamma.
\]
The restriction statements are equalities of typed Dolbeault classes,
not claims about a preferred literal coefficient gauge:
\[
 i_b^*\mathfrak g=[\Gamma]_{\bar\partial}
   \in H^{0,1}_{\bar\partial}(F_b),
 \qquad
 i_b^*\mathfrak q=[Q_F]_{\bar\partial}
   \in H^{1,1}_{\bar\partial}(F_b).
\]
After fixing \(\mathfrak g\) and one \(\mathfrak q\), changing smooth
representatives by
\(G\mapsto G+\bar\partial\phi\) and
\(Q\mapsto Q+\bar\partial\Psi\) changes the derived representatives only
by \(\partial\bar\partial\)-exact terms.  Thus no preferred literal global
coefficient atlas through the singular fillings is asserted or needed.
Put
\[
 \mathsf A:=[\partial G]_{BC},\qquad
 \mathsf B:=[\bar\partial\bar G]_{BC}
 \quad\text{in }H^{1,1}_{BC}(X).
\]

Let
\[
 \exp:H^1(X,\mathcal O_X)\longrightarrow\operatorname{Pic}(X)
\]
be induced by \(z\mapsto e^{2\pi iz}\), and orient
\(\omega\in H^2(B;\mathbb Z)\) by
\(\omega=c_1(\mathcal O_B(1))\).  The integral Picard normalization and
Bott--Chern first Chern class are
\[
 \exp(12\mathfrak g)=f^*\mathcal O_B(1),\qquad
 c_{1,BC}(\exp(z\mathfrak g))=z\mathsf A+\bar z\,\mathsf B,
\]
\[
 c_{1,BC}(f^*\mathcal O_B(1))=12(\mathsf A+\mathsf B).
\]
The absolute factor is \(12\); choosing the opposite primitive orientation
replaces \(\mathfrak g\) by \(-\mathfrak g\) in the pullback identity.
The exact fibre pairings, with the orientation stated, are
\[
 \int_{(F_b,\operatorname{or}_{\rm mark})}Q_F^2=12,
 \qquad
 \int_{(F_b,\operatorname{or}_{\mathbb C})}Q_F^2=-12,
\]
\[
 \int_{(F_b,\operatorname{or}_{\mathbb C})}
 Q_F\wedge\Gamma\wedge d\zeta_2=-1.
\]
They prove the needed Dolbeault nonvanishing, while the constant fibre
coefficients give the typed vanishing
\[
 i_b^*d_1[G\wedge Q]_{\bar\partial}
 =i_b^*[\partial(G\wedge Q)]_{\bar\partial}=0
 \quad\text{in }H^{2,2}_{\bar\partial}(F_b).
\]
Choose \(\upsilon\in\mathcal A^{3,3}(X)\) with \(\int_X\upsilon=1\).
Together with the unit
\(\mathbf1=[1]_{BC}\in H^{0,0}_{BC}(X)\), define
\[
\begin{array}{c|c}
\text{class}&\text{bidegree}\\ \hline
\mathsf A=[\partial G]_{BC},\quad
 \mathsf B=[\bar\partial\bar G]_{BC}&(1,1)\\
\mathsf C=[\partial Q]_{BC},\quad
 \overline{\mathsf C}=[\bar\partial\bar Q]_{BC}&(2,1),(1,2)\\
\mathsf D=[\partial(G\wedge Q)]_{BC},\quad
 \overline{\mathsf D}=[\bar\partial(\bar G\wedge\bar Q)]_{BC}
 &(2,2),(2,2)\\
\mathsf E=[\partial(Q^2)]_{BC},\quad
 \overline{\mathsf E}=[\bar\partial(\bar Q^2)]_{BC}
 &(3,2),(2,3)\\
\mathsf T=[\upsilon]_{BC},\quad\int_X\mathsf T=1&(3,3).
\end{array}
\]
These are bases of every nonzero positive-degree Bott--Chern group.  For
their span \(N_{\mathrm{BC}}\),
\[
 H^{*,*}_{BC}(X)=\mathbb C\mathbf1\oplus N_{\mathrm{BC}},\qquad
 N_{\mathrm{BC}}^2=0.
\]
See
\cref{prop:bcpr-generators-are-bases,thm:bcpr-square-zero-algebra}.  The
lift and smooth-gauge ambiguities above change neither the verified
dimensions nor the square-zero product algebra.

\item \Verified\ The positive radial cusp now has an explicitly declared
flat quotient metric and an exact scalar eigenvalue calculation.  Write
\[
 \mu_T=\alpha_T+i m_T,\qquad h_T=\eta_T+i H_T,\qquad
 b_T=\xi_T+i B_T,\qquad c_T=\xi_T-\eta_T,
\]
\[
 L_T=T+H_T,\qquad q_T=L_T-B_T,\qquad
 D_T=L_Tq_T+6m_T^2.
\]
Set also
\[
 m=\operatorname{Im}\mu(0),\qquad
 H=\operatorname{Im}h(0),\qquad
 B=\operatorname{Im}b(0).
\]
In the ordered lattice basis
\((\widehat\gamma,\widehat u,\widehat w,\widehat\delta)\)
and real coordinate order \((x_1,y_1,x_2,y_2)\), the real period
matrix, its oriented determinant, and its positive covolume are
\[
 P_T=
 \begin{pmatrix}
  6\alpha_T&\eta_T&1&0\\
  6m_T&L_T&0&0\\
  c_T&\alpha_T&0&1\\
  -q_T&m_T&0&0
 \end{pmatrix},\qquad
 \det P_T=-D_T,\qquad
 \operatorname{covol}(P_T\mathbb Z^4)=D_T>0,\qquad
 D_{\mathrm{can},T}=-\frac{D_T}{L_T}<0.
\]
The full inverse-dual certificate is
\[
 P_T^{-\mathsf T}=
 \begin{pmatrix}
 0&0&1&0\\[2pt]
 \dfrac{m_T}{D_T}&\dfrac{q_T}{D_T}&
 \dfrac{-6m_T\alpha_T-q_T\eta_T}{D_T}&
 \dfrac{-m_Tc_T-q_T\alpha_T}{D_T}\\[7pt]
 0&0&0&1\\[2pt]
 -\dfrac{L_T}{D_T}&\dfrac{6m_T}{D_T}&
 \dfrac{6(L_T\alpha_T-m_T\eta_T)}{D_T}&
 \dfrac{L_Tc_T-6m_T\alpha_T}{D_T}
 \end{pmatrix}.
\]
Thus \((P_T\mathbb Z^4)^*=P_T^{-\mathsf T}\mathbb Z^4\), and for
\(k=(r,s,n_1,n_2)^{\mathsf t}\), with
\[
 R_T(k)=r-6\alpha_Tn_1-c_Tn_2,\qquad
 S_T(k)=s-\eta_Tn_1-\alpha_Tn_2,
\]
its member indexed by \(k\) is exactly
\[
 \left(n_1,
 \frac{m_TR_T(k)+q_TS_T(k)}{D_T},n_2,
 \frac{-L_TR_T(k)+6m_TS_T(k)}{D_T}\right).
\]
For the declared Euclidean quotient metric \(g_T\), the first positive
scalar eigenvalue is, for all large \(T\),
\[
 \lambda_1(F_T,g_T)
 =4\pi^2\min\left\{
 \frac{L_T^2+m_T^2}{D_T^2},
 \frac{q_T^2+36m_T^2}{D_T^2}\right\}
 =\frac{4\pi^2}{T^2}\left(
 1+\frac{-2H+2\min\{B,0\}}{T}+O(T^{-2})\right).
\]
The metric normalization is recorded separately and gives
\[
 \begin{aligned}
 \widehat g_T&=D_T^{-1/2}g_T,&
 \operatorname{Vol}(F_T,\widehat g_T)&=1,\\
 \lambda_1(F_T,\widehat g_T)
 &=\frac{4\pi^2}{T}\left(
  1+\frac{-H-B/2+2\min\{B,0\}}{T}+O(T^{-2})\right),\\
 \lambda_1(F_T,\widehat g_T)&\longrightarrow0.
 \end{aligned}
\]
Normalized Haar measure on the fixed angular torus is only the unitary
measure convention used below; it is not the metric rescaling
\(g_T\mapsto\widehat g_T\).  See
 \cref{cusp-spec-exact-dual-lattice,cusp-spec-shortest-modes,%
 cusp-spec-unit-volume}.

\item \Verified\ The degree-
\(N^2\) regular-fibre cover extends across the completed local toric cusp,
with the base-change and total degrees kept distinct.  Retain
\[
 B_0=\begin{pmatrix}0&1\\-1&0\end{pmatrix},\qquad
 Z(s)=sB_0+C(t),\qquad t=e^{2\pi i s},
\]
let \(Y=Y_\Sigma\) be the toric threefold of the height-one
translation-invariant \(A_2\) fan, and let the target action be
\[
 \Psi_{\bar\lambda}(x,t)
 =\bigl(e^{2\pi iC(t)\bar\lambda}
        t^{B_0\bar\lambda}x,t\bigr),\qquad
 f_0:N_0:=Y/\bar\Lambda\longrightarrow\Delta_t.
\]
For every integer \(N\ge1\), after shrinking the disc, put
\[
 \rho_N(v)=v^N,\qquad
 C_N(v)=\frac1N C(v^N),\qquad
 \Psi^{(N)}_{\bar\lambda}(\xi,v)
 =\bigl(e^{2\pi iC_N(v)\bar\lambda}
        v^{B_0\bar\lambda}\xi,v\bigr).
\]
This is a free properly discontinuous action on a copy \(Y^{(N)}\) of
\(Y\), and
\[
 f_N:N_0^{\langle N\rangle}:=Y^{(N)}/\bar\Lambda
 \longrightarrow\Delta_v
\]
is a proper holomorphic map from a smooth threefold, with reduced
irreducible nonnormal central fibre \(W_N\) normalized by
\(\mathrm{dP}_6\).  Multiplication by
\(N\) on the three-dimensional cocharacter lattice gives
\[
 P_N:Y^{(N)}\longrightarrow Y,\qquad
 (\xi_1,\xi_2,v)\longmapsto(\xi_1^N,\xi_2^N,v^N),
\]
whose maximal-chart ring map is
\[
 \mathbb C\{z_0,z_1,z_2\}\longrightarrow
 \mathbb C\{w_0,w_1,w_2\},\qquad z_i\longmapsto w_i^N.
\]
It is finite flat of degree \(N^3\), is equivariant for the two displayed
actions, and descends to a finite flat degree-
\(N^3\) map
\[
 \kappa_N:N_0^{\langle N\rangle}\longrightarrow N_0,\qquad
 f_0\circ\kappa_N=\rho_N\circ f_N,\qquad
 \kappa_N^*W=NW_N,\qquad
 \operatorname{Ram}(\kappa_N)=(N-1)W_N.
\]
For the normal base change
\(N_0[N]=N_0\times_{\Delta_t}\Delta_v\), the unique factorization
\[
 N_0^{\langle N\rangle}\xrightarrow{Q_N}N_0[N]\longrightarrow N_0
\]
has generic degree \(N^2\).  If \(N>1\), \(Q_N\) is quasi-etale and
nonflat: its non-etale and nonflat loci are precisely the inverse images of
the double and triple strata, where the closed-fibre lengths are respectively
\(2N-1>N\) and \(3N^2-3N+1>N^2\).  On reduced central fibres the generic
degree is \(N^2\), and the normalization square is
\[
\begin{CD}
 \mathrm{dP}_6 @>{[N]}>> \mathrm{dP}_6\\
 @V{\nu_N}VV @VV{\nu}V\\
 W_N @>>> W,
\end{CD}
\]
where \([N]\) has degree \(N^2\) and ramification index \(N\) on each of
the six toric boundary curves.  Over \(v\ne0\), \(Q_N\) is the cover for
\[
 \mathbb Z\langle\widehat\gamma,\widehat u,
 N\widehat w,N\widehat\delta\rangle\subset\Lambda,
 \qquad \operatorname{Deck}(Q_N)\cong(\mathbb Z/N\mathbb Z)^2,
\]
and, with \(S_N=\operatorname{diag}(1,1,N,N)\),
\[
 S_N^{-1}M_0^NS_N=M_0.
\]
This is a local completed-cusp morphism, not a degree-
\(N^2\) self-map over the unchanged disc, not a normalization, and not a
cover over the entire base; see
\cref{cusp-cover-extension-theorem,cusp-cover-naive-action-failure,%
cusp-cover-global-scope}.
Its use in four-dimensional finite-cover unspooling is an ancillary
construction: it changes both the fibre lattice and the cusp base.  It does
not answer the native \(S^{12}\)/\(S^{24}\) sphere-route problem, whose
solution begins from the unaltered quotient fibres \(Q_T\) and \(C_T\).

\item \Verified\ The native cusp quotient and the invariant-circle quotient
give distinct exact conical smash-sphere carriers.  With
\[
 K_\Lambda=\mathbb Z\langle\widehat w,\widehat\delta\rangle
 =\operatorname{im}(M_0-I),\qquad
 \overline\Lambda=\Lambda/K_\Lambda,\qquad
 \overline T^2=(\overline\Lambda\otimes\mathbb R)/\overline\Lambda,
\]
there is a fibrewise real-analytic group epimorphism
\[
 \overline\pi:N_0\setminus W\longrightarrow D_0^*\times\overline T^2
\]
with connected kernel and trivial target local system.  On
\(F_T=f^{-1}(e^{-2\pi T})\), it is \(\pi_T\) after the isometry
\[
 \iota_T:\overline T^2\xrightarrow{\sim}Q_T,\qquad
 [(a,b)]\longmapsto[Y_T(a,b)^{\mathsf T}],\qquad
 Y_T=\begin{pmatrix}6m_T&L_T\\-q_T&m_T\end{pmatrix},
\]
 and, for a fixed constant \(C_{\mathrm{sys}}>0\),
 \(\operatorname{sys}(Q_T)\ge T-C_{\mathrm{sys}}\).  For
\[
 \mathcal A=
 \frac{[0,\epsilon]\times S^1\times\overline T^2}
 {\{0\}\times S^1\times\overline T^2\sim v_0,
  \ \{\epsilon\}\times S^1\times\overline T^2\sim v_1}
 \cong\operatorname{Susp}(T^3),
\]
the map
\[
 q_2:(X,x_0)\longrightarrow(\mathcal A,v_0),\qquad
 q_2(x)=[r,e^{i\phi},\overline\pi_f(x)]
\]
for \(t(f(x))=re^{i\phi}\), with
\(\overline\pi_f=\operatorname{pr}_{\overline T^2}\circ\overline\pi\),
\(q_2(W)=v_0\) and
\(q_2(X\setminus\operatorname{int}N_0)=v_1\), is a continuous surjective
based quotient.  Hence
\[
 q_2^{\wedge2}:X\wedge X\cong S^{12}\longrightarrow
 \mathcal A\wedge\mathcal A
\]
is global, and its synchronized radial restriction is exactly
\[
 \pi_T\times\pi_T:F_T\times F_T\longrightarrow Q_T\times Q_T,
 \qquad
 (Q_T\times Q_T,0)
 \xrightarrow[T\to\infty]{\mathrm{pointed}}(\mathbb R^4,0)
\]
in pointed local-isometric, pointed smooth, and pointed
Gromov--Hausdorff topology.

The cusp-supported one-circle route uses
\[
 \mathcal B=
 \frac{[0,\epsilon]\times S^1\times S^1}
 {\{0\}\times S^1\times S^1\sim b_0,\ 
  \{\epsilon\}\times S^1\times S^1\sim b_1}
 \cong\operatorname{Susp}(T^2)
\]
and the continuous surjective based quotient
\[
 q_{1,\mathrm{cusp}}:(X,x_0)\longrightarrow(\mathcal B,b_0),\qquad
 q_{1,\mathrm{cusp}}(x)=[r,e^{i\phi},\chi_\gamma(x)]
\]
on \(N_0\setminus W\), with \(q_{1,\mathrm{cusp}}(W)=b_0\) and the
exterior collapsed to \(b_1\).  Its fourth smash power is
\[
 q_{1,\mathrm{cusp}}^{\wedge4}:X^{\wedge4}\cong S^{24}
 \longrightarrow\mathcal B^{\wedge4},
\]
and its synchronized radial restriction is
\(\chi_{\gamma,T}^{\times4}:F_T^4\to C_T^4\).

For the full-base circle route, let
\[
 P=B\setminus\bigcup_{j=0}^2\operatorname{int}D_j,\qquad
 \mathcal M=(P\times S^1)\underset{h_1^{S^1}}\cup C_1
                         \underset{h_2^{S^1}}\cup C_2,\qquad
 \partial\mathcal M=\partial D_0\times S^1,
\]
and
\(\mathcal C=\mathcal M\cup_{\partial\mathcal M}
\operatorname{Cone}(\partial\mathcal M)\), with cone vertex \(c_0\).
The continuous surjective based quotient
\[
 q_1:(X,x_0)\longrightarrow(\mathcal C,c_0)
\]
 equals \((f,\theta_\gamma)\) over \(P\), equivalently
 \(\chi_\gamma|_{f^{-1}(P)}\), and equals the finite-filling maps
\(N_j\to C_j\) for \(j=1,2\), sends
\(x\in N_0\setminus W\) to
 \([r,(e^{i\phi},\theta_\gamma(x))]\), and sends \(W\) to \(c_0\).  Its
fourth smash power is
\[
 q_1^{\wedge4}:X^{\wedge4}\cong S^{24}\longrightarrow
 \mathcal C^{\wedge4},
\]
whose synchronized radial restriction is exactly
\[
\begin{aligned}
 \chi_{\gamma,T}^{\times4}&:F_T^4\longrightarrow C_T^4,\\
 C_T&=\mathbb R/\ell_T\mathbb Z,\\
 \ell_T&=\frac{D_T}{\sqrt{L_T^2+m_T^2}}
       =T+H-B+O(T^{-1}),
\end{aligned}
\]
so
\[
 (C_T^4,0)
 \xrightarrow[T\to\infty]{\mathrm{pointed}}(\mathbb R^4,0)
\]
in the same three pointed topologies.

The exponents \(2\) and \(4\), and hence the source-sphere dimensions
\(12\) and \(24\), are optimal for these two native quotient problems with
the unscaled quotient of the marked Euclidean fibre metric.  Indeed, if a
primitive kernel \(N_{\mathrm{ker}}\subset\Lambda^{\oplus r}\) defines a
connected torus
quotient of \(r\) cusp fibres whose quotient systole tends to infinity, then
every uniformly bounded period must lie in \(N_{\mathrm{ker}}\).  Thus
\[
 K_\Lambda^{\oplus r}\subset N_{\mathrm{ker}},\qquad
 \operatorname{rank}(\Lambda^{\oplus r}/N_{\mathrm{ker}})\le 2r.
\]
For a quotient over the full regular base, or over any finite cover of that
base, monodromy forces the additional sublattice
\(\mathbb Z\langle\widehat u\rangle^{\oplus r}\) into
\(N_{\mathrm{ker}}\): one uses
\(A_1\widehat w=\widehat u-\widehat w\), and on a finite cover the exact
formula
\[
 (A_2A_1)^m\widehat w
 =\widehat w+m(\widehat u-\widehat\delta)
\]
together with primitivity of \(N_{\mathrm{ker}}\).  Consequently
\[
 \operatorname{rank}(\Lambda^{\oplus r}/N_{\mathrm{ker}})\le r.
\]
A rank-four expanding quotient therefore needs at least two copies
cusp-locally and at least four copies globally; the displayed
\(Q_T^2\) and \(C_T^4\) constructions attain the two bounds.

After collapsing the coordinate one-skeleton
\[
 K^{(1)}=
 ((\mathbb R/\mathbb Z)\times\{0\})
 \cup(\{0\}\times(\mathbb R/\mathbb Z))
 \subset\overline T^2,
\]
the two native routes also give continuous surjective based quotient maps
\[
 \widehat q_3:X\twoheadrightarrow S^3,\qquad
 \widehat q_\gamma:X\twoheadrightarrow S^2,
\]
and hence, using the fixed based choice
\(X\cong_{\mathrm{diff}}S^6\),
\[
 \widehat q_3^{\wedge2}:S^{12}\twoheadrightarrow S^6,\qquad
 \widehat q_\gamma^{\wedge4}:S^{24}\twoheadrightarrow S^8.
\]
Both full sphere carriers are based null-homotopic.  First let
\(f=\widehat q_\gamma:S^6\to S^2\) after the chosen based identification.
There is an exact factorization
\[
 S^{12}=S^6\wedge S^6
 \xrightarrow{\ f\wedge\operatorname{id}_{S^6}=\Sigma^6f\ }
 S^2\wedge S^6=S^8
 \xrightarrow{\ \operatorname{id}_{S^2}\wedge f\ }
 S^2\wedge S^2=S^4.
\]
Freudenthal stability gives
\(\pi_{12}(S^8)\cong\pi_4^{\mathrm S}=0\), so \(f\wedge f\) is based
null-homotopic.  Since
\(f^{\wedge r}=(f\wedge f)\wedge f^{\wedge(r-2)}\), it follows that
\[
 [f^{\wedge r}]=0\in\pi_{6r}(S^{2r})
 \qquad\text{for every }r\ge2.
\]
In particular, the displayed surjective quotient
\(\widehat q_\gamma^{\wedge4}:S^{24}\twoheadrightarrow S^8\) is itself
based null-homotopic.

For the other carrier, let \(g:S^6\to S^3\) be any based map.  Its exact
smash factorization is
\[
 g\wedge g=(\Sigma^3g)\circ(\Sigma^6g):
 S^{12}\longrightarrow S^9\longrightarrow S^6.
\]
Write \(\operatorname{st}\) for stabilization.  The exact low-stem data are
\[
 \pi_6(S^3)\cong\mathbb Z/12,
 \quad
 \pi_3^{\mathrm S}=\mathbb Z/24\langle\nu\rangle,
 \quad
 \pi_6^{\mathrm S}=\mathbb Z/2\langle\nu^2\rangle,
\]
\[
 \pi_{12}(S^6)=\mathbb Z/2\langle\nu_6\circ\nu_9\rangle,
 \qquad
 \operatorname{st}:\pi_{12}(S^6)
 \xrightarrow{\ \cong\ }\pi_6^{\mathrm S}.
\]
The full groups are tabulated in \cite[p.~339]{Hatcher2002}; the signed smash
formula and the unstable stabilization calculation are
\cite[Propositions~3.1 and~5.11]{Toda1962}.
Since \(12[g]=0\), one has
\(\operatorname{st}([g])=2a\nu\) for some \(a\in\mathbb Z/12\).  By
multiplicativity,
\[
 \operatorname{st}([g\wedge g])
 =\operatorname{st}([g])^2
 =4a^2\nu^2=0\in\pi_6^{\mathrm S}.
\]
The stabilization isomorphism forces
\([g\wedge g]=0\in\pi_{12}(S^6)\).
Thus \(g^{\wedge r}\) is based null-homotopic for every \(r\ge2\).  Taking
\(g=\widehat q_3\) proves
\[
 [\widehat q_3^{\wedge2}]=0\in\pi_{12}(S^6),
 \qquad
 [\widehat q_\gamma^{\wedge4}]=0\in\pi_{24}(S^8).
\]
Both displayed maps remain surjective quotient maps.  The one-copy classes
\([\widehat q_3]\in\pi_6(S^3)\) and
\([\widehat q_\gamma]\in\pi_6(S^2)\) have not been computed.
Their synchronized radial restrictions contain based topologically embedded
four-spheres, equipped separately with the intrinsic quotient length metrics
\[
 \mathbb S_{12,T}^4
   =(Q_T/K_T^{(1)})^{\wedge2}\hookrightarrow S^6,\qquad
 \mathbb S_{24,T}^4=C_T^{\wedge4}\hookrightarrow S^8,
\]
obtained from, respectively,
\[
 F_T^2\xrightarrow{\ \pi_T^{\times2}\ }Q_T^2
 \xrightarrow{\ c_T^{\times2}\ }(S_T^2)^2
 \longrightarrow\mathbb S_{12,T}^4,\qquad
 F_T^4\xrightarrow{\ \chi_{\gamma,T}^{\times4}\ }C_T^4
 \longrightarrow\mathbb S_{24,T}^4,
\]
where the last arrows collapse the fat wedges.  The full maps from
\(F_T^2\) and \(F_T^4\) retain the nontrivial kernels of their first arrows;
only the quotient-stage maps are injective off their explicitly collapsed
skeleton/fat-wedge loci.  Equip the displayed products
with their Euclidean \(\ell^2\)-product metrics and give each four-sphere the
induced quotient length metric.  If
\(z_T\) is the image of \([(1/2,1/2)]\) and
\[
 \zeta_{12,T}=[z_T,z_T],\qquad
 \zeta_{24,T}=
 [\ell_T/2,\ell_T/2,\ell_T/2,\ell_T/2],
\]
then every radius-\(R\) ball about these points is exactly a Euclidean
four-ball once, respectively,
\[
 T-C_{\mathrm{sys}}>4R,\qquad \ell_T>4R.
\]
Therefore
\[
 (\mathbb S_{12,T}^4,\zeta_{12,T})\longrightarrow(\mathbb R^4,0),\qquad
 (\mathbb S_{24,T}^4,\zeta_{24,T})\longrightarrow(\mathbb R^4,0)
\]
in pointed local-isometric and pointed Gromov--Hausdorff topology, with smooth
convergence on the eventual exact Euclidean neighborhoods of the displayed
base points.  Both based inclusion maps
\[
 i_{12,T}:\mathbb S_{12,T}^4\hookrightarrow S^6,
 \qquad
 i_{24,T}:\mathbb S_{24,T}^4\hookrightarrow S^8
\]
are null-homotopic: sliding their fixed interior suspension coordinate to a
suspension endpoint contracts them through based maps.  Thus
\((i_{12,T})_*[S^4]=0\in H_4(S^6;\mathbb Z)\) and
\((i_{24,T})_*[S^4]=0\in H_4(S^8;\mathbb Z)\).  These slice-inclusion
null-homotopies are distinct from the global equalities
\([\widehat q_3^{\wedge2}]=0\in\pi_{12}(S^6)\) and
\([\widehat q_\gamma^{\wedge4}]=0\in\pi_{24}(S^8)\).

Deleting the smash basepoint nevertheless gives literal nonclosed
topologically embedded copies of \(\mathbb R^4\):
\[
 (Q_T\setminus K_T^{(1)})^2
 \cong\mathbb S_{12,T}^4\setminus\{*\}
 \cong\mathbb R^4\hookrightarrow S^6,
\qquad
 (C_T\setminus\{0\})^4
 \cong\mathbb S_{24,T}^4\setminus\{*\}
 \cong\mathbb R^4\hookrightarrow S^8.
\]
They are open only in the embedded four-spheres and are not globally
Euclidean-isometric under a fixed compatible compact-target metric.

There are also Hausdorff, second-countable topological five-manifolds with
boundary \(\mathfrak E_{12}\) and \(\mathfrak E_{24}\) whose positive fibres
are \(Q_{1/s}^{\times2}\) and \(C_{1/s}^{\times4}\), respectively, and whose
zero fibres are literally \(\mathbb R^4\).  Continuous surjections
\(\beta_\nu:\mathfrak E_\nu\to\mathcal D_\nu\) identify every positive fibre
with its radial carrier slice and collapse the zero fibre to the cone vertex.
They are homeomorphisms off the zero fibre but are nonproper and nonquotient;
thus they are comparison/collapse maps, not topological blow-downs or
resolutions, and \(\mathcal D_\nu\) is not the quotient
\(\mathfrak E_\nu/\mathbb R^4\) via \(\beta_\nu\).

The metric quantifiers are essential.  For every fixed metric compatible
with the compact conical target topology, the target slices
\(j_T(Q_T^2)\), \(k_T^{\mathcal B}(C_T^4)\), and
\(k_T^{\mathcal C}(C_T^4)\) converge in Hausdorff, hence
Gromov--Hausdorff, distance to the smash basepoint.  A fixed cusp-adapted
metric can realize the intrinsic limits only after a moving-point blow-up:
with \(r_T=e^{-2\pi T}\), the factors \(T/r_T\) and
\(\ell_T/r_T\) give the exact induced slice metrics
\[
 \left(\frac{T}{r_T}\right)^2j_T^*g_{Z_{12}}
   =g_{Q_T}\oplus g_{Q_T},\qquad
 \left(\frac{\ell_T}{r_T}\right)^2
   (k_T^{\mathcal D})^*g_{Z_{24}^{\mathcal D}}
   =g_{C_T}^{\oplus4},\quad
 \mathcal D\in\{\mathcal B,\mathcal C\}.
\]
The slice bubbles are \(\mathbb R^4\), but the corresponding full ambient
bubbles are \(\mathbb R^8\) for the \(S^{12}\) route and
\(\mathbb R^{12}\) for either \(S^{24}\) route.  More generally, if
\(\lambda_T>0\) and \(\beta_T=\lambda_Tr_T\), the slice limit is a point
when \(\beta_T\to0\), a square flat four-torus of side \(L\) when
\(\beta_T\to L\in(0,\infty)\), and \(\mathbb R^4\) when
\(\beta_T\to\infty\); in the last case the cone vertex is at scaled
distance \(\beta_T\to\infty\).  Hence the \(\mathbb R^4\) limits are
neither vertex tangent cones nor ambient limits of fixed smooth metrics on
the source spheres, whose ordinary tangent dimensions are \(12\) and \(24\).

The original maps \(q_2^{\wedge2}\), \(q_{1,\mathrm{cusp}}^{\wedge4}\),
and \(q_1^{\wedge4}\) have the displayed conically stratified codomains.
The repeated-smash null theorems for the sphere-valued composites
\(\widehat q_3\) and \(\widehat q_\gamma\) do not supply null-homotopies
before the final one-skeleton/base-angle quotients.  The unstable classes of
\[
 q_2^{\wedge2}:S^{12}\to\mathcal A^{\wedge2},
 \qquad
 q_1^{\wedge4}:S^{24}\to\mathcal C^{\wedge4}
\]
remain \Open: their zero ordinary Hurewicz images and absence of homotopy
sections do not decide nullity.  The class of
\(q_{1,\mathrm{cusp}}^{\wedge4}:S^{24}\to\mathcal B^{\wedge4}\) has
likewise not been computed.
The sphere codomains occur only after the additional one-skeleton and
base-angle collapses defining \(\widehat q_3\) and
\(\widehat q_\gamma\).  All identifications of the sources with
\(S^{12}\) and \(S^{24}\) use the chosen based diffeomorphism
\(X\cong_{\mathrm{diff}}S^6\).  See
\cref{cusp-cone-native-descent,cusp-cone-S12-route,%
cusp-cone-S24-local-route,cusp-cone-S24-route,cusp-cone-optimal-exponents,%
cusp-cone-angle-forgotten-sphere-carriers,cusp-cone-S12-sphere-map-null,%
cusp-cone-S24-sphere-map-null,%
cusp-cone-sphere-map-homotopy-status,cusp-cone-moving-S4-null-homotopy,%
cusp-cone-literal-R4-open-cells,%
cusp-cone-exact-decompactification-space,%
sphere-metric-uniform-collapse,sphere-metric-fixed-realization,%
sphere-metric-scaling-trichotomy,sphere-metric-tangent-scope,%
cusp-cone-metric-scope,cusp-cone-target-homotopy-types,%
cusp-cone-obstructions}.

\item \Refuted\ The assertion that the literal fixed-target map
\[
 \chi_\gamma:J\longrightarrow B^\circ\times(\mathbb R/\mathbb Z)
\]
extends continuously across \(W\), or that
\[
 \chi_\gamma^{[4]}:J^{[4]}\longrightarrow
 B^\circ\times(\mathbb R/\mathbb Z)^4
\]
extends with the stated regular-fibre restriction to the compact fourfold
fibre product or to \(X^4\), is false: two explicit paths have the same
toric limit in \(W\) and constant character values \(3/4\) and \(1/2\).
The cone map above changes the codomain by collapsing
\(\partial D_0\times S^1\) to \(c_0\), so it does not contradict this
nonextension.  Moreover
\[
 X^2\simeq S^6\times S^6\not\simeq S^{12},\qquad
 X^4\simeq(S^6)^4\not\simeq S^{24};
\]
only the smash products satisfy
\(X^{\wedge2}\cong S^{12}\) and \(X^{\wedge4}\cong S^{24}\).
Neither smash sphere admits an almost-complex structure, and neither route
constructs a compact complex sphere; see
\cref{cusp-spec-globalization-obstruction,cusp-cone-obstructions}.

\item \Verified\ The exceptional first-coordinate elliptic quotient is now
classified on the actual period image, rather than on an arbitrary ambient
path.  For
\[
 \mathcal A_z=\mathbb C^2/
 \mathbb Z\langle(6\mu,\beta),(\tau,\mu),(1,0),(0,1)\rangle,
\]
the first coordinate induces a surjective holomorphic homomorphism
\[
 \mathcal A_z\longrightarrow\mathbb C/(\mathbb Z\tau+\mathbb Z)
\]
if and only if \(6\mu=a\tau+b\) for a unique
\((a,b)\in\mathbb Z^2\).  Its connected kernel has period
\(\tau_f=a\mu-\beta\) and
\(\operatorname{Im}\tau_f=-D_{\mathrm{can}}>0\).  The exact marked-label
action is
\[
 g_0:(a,b)\mapsto(a,a+b),\qquad
 g_2:(a,b)\mapsto(-b,6+a).
\]
The union of these loci has finite image in the completed orbifold base and
is absent from a punctured cusp neighborhood.  Consequently this particular
elliptic/Landau sector does not produce the cusp soft modes.  This statement
does not classify all elliptic quotients, all N\'eron--Severi classes, or all
smooth integral magnetic fluxes.  See
\cref{cusp-spec-first-coordinate-descent,%
cusp-spec-elliptic-descent-finiteness,%
cusp-spec-elliptic-descent-scope}.

\item \Verified\ The fibrewise small eigenvalues are assembled into one
specified gapless operator.  On
\[
 \mathscr H_{\mathrm{aux}}
 =L^2\bigl([T_0,\infty),dT;
 L^2_0(\mathbb R^4/\mathbb Z^4,dy)\bigr)
 \cong\bigoplus_{k\in\mathbb Z^4\setminus\{0\}}
 L^2([T_0,\infty),dT),
\]
put
\[
 \lambda_k(T)=4\pi^2\lvert P_T^{-\mathsf T}k\rvert^2,\qquad
 (\mathscr L_{\mathrm{aux}}u)_k(T)=\lambda_k(T)u_k(T),
\]
\[
 D(\mathscr L_{\mathrm{aux}})=\left\{u:
 \sum_{k\ne0}\int_{T_0}^{\infty}
 \lambda_k(T)^2\lvert u_k(T)\rvert^2\,dT<\infty\right\}.
\]
This multiplication operator is nonnegative and self-adjoint,
\(\ker\mathscr L_{\mathrm{aux}}=0\), and the escaping-support Weyl
sequence in the \(k=(1,0,0,0)\) channel proves
\[
 0\in\sigma_{\mathrm{ess}}(\mathscr L_{\mathrm{aux}}).
\]
This is a spectral-gap theorem for the displayed auxiliary direct integral,
not for the complex threefold without a metric and not for a quantum gauge
Hamiltonian; see \cref{cusp-interpretation-direct-integral-gapless}.

\item \Verified\ The cover index and the cubical mesh are independent.
For integers \(M_{\mathrm{cov}}\ge1\) and
\(N_{\mathrm{mesh}}\ge2\), put
\[
 S_{M_{\mathrm{cov}}}=\operatorname{diag}(1,1,M_{\mathrm{cov}},
 M_{\mathrm{cov}}),\qquad
 P_{T,M_{\mathrm{cov}}}=P_TS_{M_{\mathrm{cov}}},
\]
\[
 \widetilde F_{T,M_{\mathrm{cov}}}
 =\mathbb R^4/P_{T,M_{\mathrm{cov}}}\mathbb Z^4
 \longrightarrow F_T,\qquad
 \operatorname{Deck}\cong(\mathbb Z/M_{\mathrm{cov}}\mathbb Z)^2.
\]
This cover has degree \(M_{\mathrm{cov}}^2\), Gram matrix and covolume
\[
 G_{T,M_{\mathrm{cov}}}
 =P_{T,M_{\mathrm{cov}}}^{\mathsf T}P_{T,M_{\mathrm{cov}}},
 \qquad
 \mathcal D_{T,M_{\mathrm{cov}}}=M_{\mathrm{cov}}^2D_T.
\]
For the cubical complex write, once and for all,
\[
 V_{N_{\mathrm{mesh}}}
   =(\mathbb Z/N_{\mathrm{mesh}}\mathbb Z)^4,\qquad
 E_{N_{\mathrm{mesh}}}
   =V_{N_{\mathrm{mesh}}}\times\{1,2,3,4\}.
\]
The diffeomorphism
\[
 \Phi_{T,M_{\mathrm{cov}}}:\mathbb R^4/\mathbb Z^4
 \xrightarrow{\ \sim\ }\widetilde F_{T,M_{\mathrm{cov}}},
 \qquad[y]\longmapsto[P_{T,M_{\mathrm{cov}}}y]
\]
transports the whole standard \(N_{\mathrm{mesh}}\)-cubical complex; its
vertex map is precisely
\[
 (\mathbb Z/N_{\mathrm{mesh}}\mathbb Z)^4
 \xrightarrow{\ \sim\ }
 \widetilde F_{T,M_{\mathrm{cov}}}[N_{\mathrm{mesh}}],
 \qquad[n]\longmapsto
 \left[\frac{P_{T,M_{\mathrm{cov}}}n}{N_{\mathrm{mesh}}}\right],
\]
so the finite set is the torsion subgroup, not the entire torus.  Its
oriented edge is
\[
 \gamma_{n,j}(t)=
 \left[\frac{P_{T,M_{\mathrm{cov}}}(n+t\mathbf e_j)}
 {N_{\mathrm{mesh}}}\right],\qquad0\le t\le1.
\]
For a compact matrix Lie group \(G\), the parallel-transport convention
\[
 U'(t)=-A_{\gamma(t)}(\dot\gamma(t))U(t),\qquad U(0)=1
\]
defines the typed map
\[
 \operatorname{Hol}_{T,M_{\mathrm{cov}},N_{\mathrm{mesh}}}:
 \Omega^1(\widetilde F_{T,M_{\mathrm{cov}}};\mathfrak g)
 \longrightarrow G^{E_{N_{\mathrm{mesh}}}},
 \qquad A\longmapsto(U_A(\gamma_{n,j}))_{n,j},
\]
and for \(A^g=g^{-1}Ag+g^{-1}dg\) its exact endpoint law is
\[
 \operatorname{Hol}(A^g)_{n,j}
 =g(\gamma_{n,j}(1))^{-1}\operatorname{Hol}(A)_{n,j}
  g(\gamma_{n,j}(0)).
\]
See \cref{cusp-link-cover-metric,cusp-link-cubical-realization,%
cusp-link-holonomy-map}.

\item \Verified\ The compact-link \(\mathrm{SU}(2)\) flux uses the
inverse-transport convention and has an exact metric-weighted action.  With
\[
 H_{\mathrm{link}}=\begin{pmatrix}i&0\\0&-i\end{pmatrix},\qquad
 \lVert H_{\mathrm{link}}\rVert^2=1,\qquad
 \exp(2\pi H_{\mathrm{link}})=1,\qquad
 N=N_{\mathrm{mesh}},
\]
set
\[
\begin{aligned}
 U_1(n)&=\exp(-2\pi n_2H_{\mathrm{link}}/N^2),\\
 U_2(n)&=
 \begin{cases}
  1,&0\le n_2\le N-2,\\
  \exp(2\pi n_1H_{\mathrm{link}}/N),&n_2=N-1,
 \end{cases}\\
 U_3(n)&=U_4(n)=1.
\end{aligned}
\]
For
\[
 U_{ij}(n)=U_i(n)U_j(n+\mathbf e_i)
 U_i(n+\mathbf e_j)^{-1}U_j(n)^{-1},
\]
including all wrap-around edges,
\begin{align*}
 U_{12}(n)&=\exp(2\pi H_{\mathrm{link}}/N^2),\\
 U_{ij}(n)&=1\quad(i<j,\ (i,j)\ne(1,2)).
\end{align*}
The principal logarithm is unambiguous because
\(0<2\pi/N^2\le\pi/2\), and
\[
 \mathcal F_N^{\log}(n)
 =N^2\operatorname{Log}(U_{12}(n))\,dy^1\wedge dy^2
 =2\pi H_{\mathrm{link}}\,dy^1\wedge dy^2.
\]
With the full inverse Gram metric,
\[
 S^{\log}_{T,M_{\mathrm{cov}},N}(U)
 =\frac{\mathcal D_{T,M_{\mathrm{cov}}}}{2g_{\mathrm{YM}}^2N^4}
 \sum_{n\in V_{N_{\mathrm{mesh}}}}
 \lVert\mathcal F_N^{\log}(n)\rVert^2_{
 G_{T,M_{\mathrm{cov}}}^{-1}\otimes\mathfrak{su}(2)}
 =\frac{2\pi^2M_{\mathrm{cov}}^2}{g_{\mathrm{YM}}^2D_T}.
\]
There is no canonical scalar Wilson action for a general nonorthogonal
metric obtained by independent scalar plaquette weights.  For this
configuration the explicitly defined \((1,2)\)-sector functional has
\[
 \kappa^{12}_{T,M_{\mathrm{cov}}}
 =(G_{T,M_{\mathrm{cov}}}^{-1})_{11}
  (G_{T,M_{\mathrm{cov}}}^{-1})_{22}
 -(G_{T,M_{\mathrm{cov}}}^{-1})_{12}^2
 =D_T^{-2},
\]
\[
 S^{W,12}_{T,M_{\mathrm{cov}},N}(U)
 =\frac{\mathcal D_{T,M_{\mathrm{cov}}}
        \kappa^{12}_{T,M_{\mathrm{cov}}}}{2g_{\mathrm{YM}}^2}
  \sum_{n\in V_{N_{\mathrm{mesh}}}}
  \left(2-\operatorname{Re}\operatorname{tr}U_{12}(n)\right)
 =\frac{M_{\mathrm{cov}}^2N^4}{2g_{\mathrm{YM}}^2D_T}
  \left(2-2\cos\frac{2\pi}{N^2}\right)
\]
Its exact mesh expansion is
\[
 S^{W,12}_{T,M_{\mathrm{cov}},N}(U)
 =\frac{M_{\mathrm{cov}}^2}{g_{\mathrm{YM}}^2D_T}
 \left(2\pi^2-\frac{2\pi^4}{3N^4}
       +\frac{4\pi^6}{45N^8}+O(N^{-12})\right),
\]
so its \(N\to\infty\) limit is
\(2\pi^2M_{\mathrm{cov}}^2/(g_{\mathrm{YM}}^2D_T)\).
See \cref{cusp-link-flux-theorem}.

\item \Verified\ The finite Gaussian oscillator is defined on the real
reduced Coulomb space with the physical inverse-Gram exterior metric.  For
\[
 \vartheta^j_{T,M_{\mathrm{cov}}}
 =P_{T,M_{\mathrm{cov}}}^{-\mathsf T}\mathbf e_j,\qquad
 \mathfrak d_{N_{\mathrm{mesh}}}(k)_j
 =N_{\mathrm{mesh}}
  (e^{2\pi ik_j/N_{\mathrm{mesh}}}-1),
\]
let
\[
 C^p_{N_{\mathrm{mesh}}}(T,M_{\mathrm{cov}})
 =\operatorname{Map}\!\left(
 (\mathbb Z/N_{\mathrm{mesh}}\mathbb Z)^4,
 \bigwedge^p(\mathbb R^4)^*\otimes\mathfrak{su}(2)\right).
\]
The exterior metric is
\[
 \left\langle
 \vartheta^{i_1}\wedge\cdots\wedge\vartheta^{i_p},
 \vartheta^{j_1}\wedge\cdots\wedge\vartheta^{j_p}
 \right\rangle
 =\det\left((G_{T,M_{\mathrm{cov}}}^{-1})_{i_a j_b}\right)_{a,b=1}^p,
\]
and the cochain inner product is
\[
 \langle a,b\rangle_p
 =\frac{\mathcal D_{T,M_{\mathrm{cov}}}}{N_{\mathrm{mesh}}^4}
 \sum_{n\in V_{N_{\mathrm{mesh}}}}\langle a(n),b(n)\rangle.
\]
Writing \(\mathsf T_ja(n)=a(n+\mathbf e_j)\), put
\[
 D_j^+=N_{\mathrm{mesh}}(\mathsf T_j-1),\qquad
 D_j^-a(n)=N_{\mathrm{mesh}}(a(n)-a(n-\mathbf e_j)),
\]
and then
\[
 d_N=\sum_{j=1}^4\varepsilon(\vartheta^j)D_j^+,\qquad
 d_N^*=-\sum_{j=1}^4\iota_{(\vartheta^j)^\sharp}D_j^-,
\]
\[
 \mathscr C_{T,M_{\mathrm{cov}},N_{\mathrm{mesh}}}
 =\ker(d_N^*:C^1\to C^0)\cap
  (\ker\Delta_N^{(1)})^\perp,\qquad
 \mathcal L=g_{\mathrm{YM}}^{-2}d_N^*d_N.
\]
The positive eigenvalue on the Fourier fibre
\(k\in(\mathbb Z/N_{\mathrm{mesh}}\mathbb Z)^4\setminus\{0\}\) is
\[
 \lambda_{T,M_{\mathrm{cov}},N_{\mathrm{mesh}}}(k)
 =\frac1{g_{\mathrm{YM}}^2}
 \left\lVert P_T^{-\mathsf T}S_{M_{\mathrm{cov}}}^{-1}
 \mathfrak d_{N_{\mathrm{mesh}}}(k)\right\rVert^2,
\]
with complex transverse multiplicity \(9\), and
\[
 \lambda(\mathbf e_1)
 =\frac{4N_{\mathrm{mesh}}^2\sin^2(\pi/N_{\mathrm{mesh}})}
 {g_{\mathrm{YM}}^2}
 \frac{L_T^2+m_T^2}{D_T^2},
\]
which is independent of \(M_{\mathrm{cov}}\).  On
\(L^2(\mathscr C_{T,M_{\mathrm{cov}},N_{\mathrm{mesh}}},da)\), the
Friedrichs realization of
\[
 \widehat H^{\mathrm{Gauss}}
 =-\frac12\Delta_{\mathscr C_{T,M_{\mathrm{cov}},N_{\mathrm{mesh}}}}
  +\frac12\langle a,\mathcal La\rangle
\]
has a unique Gaussian ground state and
\[
 \operatorname{gap}(\widehat H^{\mathrm{Gauss}})
 =\sqrt{\min_{k\ne0}\lambda(k)}
 \le\frac{2N_{\mathrm{mesh}}\sin(\pi/N_{\mathrm{mesh}})}
 {g_{\mathrm{YM}}}\frac{\sqrt{L_T^2+m_T^2}}{D_T}.
\]
For the explicit sequence
\[
 T_j=j^2,\qquad M_{\mathrm{cov},j}=j,\qquad
 N_{\mathrm{mesh},j}=j^3,
\]
one has simultaneously
\[
 \operatorname{sys}(\widetilde F_{T_j,M_{\mathrm{cov},j}})\to\infty,
 \quad
 \max_r\frac{\lVert P_{T_j,M_{\mathrm{cov},j}}\mathbf e_r\rVert}
 {N_{\mathrm{mesh},j}}\to0,
 \quad
 \operatorname{gap}(\widehat H^{\mathrm{Gauss}})\to0,
 \quad
 S^{\log}_{T_j,M_{\mathrm{cov},j},N_{\mathrm{mesh},j}}(U)\to0.
\]
See \cref{cusp-link-finite-spectrum,cusp-link-gaussian-theorem}.

\item \Verified\ The componentwise exponential
\[
 \operatorname{Exp}_{N_{\mathrm{mesh}}}:C^1
 \longrightarrow\mathrm{SU}(2)^{E_{N_{\mathrm{mesh}}}},
 \qquad a\longmapsto(\exp(a_e/N_{\mathrm{mesh}}))_e
\]
is only a local diffeomorphism at zero.  Its derivative is
\[
 D\operatorname{Exp}_{N_{\mathrm{mesh}}}(0)a
 =N_{\mathrm{mesh}}^{-1}a,
\]
and, if a smooth compact-link action \(S_W\) has the identity as a critical
point, then exactly
\[
 D^2(S_W\circ\operatorname{Exp}_{N_{\mathrm{mesh}}})_0[a,b]
 =N_{\mathrm{mesh}}^{-2}D^2(S_W)_{\mathbf1}[a,b].
\]
It is not globally injective: a cochain supported on one edge with value
\(2\pi N_{\mathrm{mesh}}H_{\mathrm{link}}\) has the same image as zero.  Thus
\[
 \mathcal H^{\mathrm{Gauss}}
 =L^2(\mathscr C_{T,M_{\mathrm{cov}},N_{\mathrm{mesh}}},da)
 \quad\text{and}\quad
 \mathcal H^{\mathrm{inv}}_{N_{\mathrm{mesh}}}
 =L^2(\mathrm{SU}(2)^{E_{N_{\mathrm{mesh}}}},dU)^{
   \mathrm{SU}(2)^{V_{N_{\mathrm{mesh}}}}}
\]
are not identified by a global coordinate chart or a constructed unitary
map; moreover \(\delta_U\notin L^2\) for product Haar measure.  No nonlinear
compact-link Hamiltonian, Gauss-law spectral problem, time direction,
reflection-positive transfer matrix, renormalized-scale limit, or
Osterwalder--Schrader reconstruction is supplied by these results.  See
\cref{cusp-link-exponential-local-only,cusp-link-physical-hilbert-space,%
cusp-link-dirac-not-l2}.

\item \Verified\ The reducible direct-sum connection on
\(\mathscr E_T=\mathscr L_T\oplus\mathscr L_T^{-1}\) is the non-flat
Yang--Mills critical point which, in the lifted splitting over
\(\mathbb R^4\), has the connection form
\[
 \nabla^T=\mathrm d+2\pi u_1\,\mathrm du_2\,H_{\mathrm{mag}},\qquad
 H_{\mathrm{mag}}=\begin{pmatrix}i&0\\0&-i\end{pmatrix},
\]
\[
 F_{\nabla^T}
 =\frac{2\pi}{D_T}H_{\mathrm{mag}}\,\mathrm dy_1\wedge\mathrm dy_2,
 \qquad \mathrm d_{\nabla^T}^*F_{\nabla^T}=0.
\]
The underlying principal \(SU(2)\)-bundle is smoothly trivial; after
choosing any of the proved global \(SU(2)\)-trivializations, the displayed
lifted-splitting form is carried to its gauge transform on that trivial
bundle.  It is not asserted to remain literally
\(\mathrm d+2\pi u_1\,\mathrm du_2\,H_{\mathrm{mag}}\) in an arbitrary
global trivialization.
For
\(\operatorname{YM}_T(\nabla)
=(2g_{\mathrm{YM}}^2)^{-1}\int_{F_T}|F_\nabla|^2\,
\mathrm d\operatorname{vol}_{F_T}\), its action is
\[
 \operatorname{YM}_T(\nabla^T)
 =\frac{2\pi^2}{g_{\mathrm{YM}}^2D_T}
 =\frac{2\pi^2}{g_{\mathrm{YM}}^2T^2}
 \left(1+\frac{B-2H}{T}+O(T^{-2})\right).
\]
On the background-gauge slice its classical Jacobi operator is
\[
 (\mathcal J_Ta)_\nu
 =-\sum_{\mu=1}^4\nabla_{e_\mu}\nabla_{e_\mu}a_\nu
  -2\sum_{\mu=1}^4[F_{\nu\mu},a_\mu],
\]
and an explicit smooth real unit vector satisfies
\[
 \mathrm d_{\nabla^T}^*a=0,\qquad
 \mathcal J_Ta=-\frac{4\pi}{D_T}a,\qquad
 \operatorname{YM}_T''|_{\nabla^T}(a,a)
 =-\frac{4\pi}{g_{\mathrm{YM}}^2D_T}.
\]
Thus the displayed reducible classical critical point is a saddle for every
finite \(T\); this is not a statement about a quantum Yang--Mills
Hamiltonian.  See
\cref{cusp-magnetic-su2-triviality,cusp-magnetic-ym-critical-point,%
cusp-magnetic-negative-mode}.

\item \Verified\ The monodromy-invariant full-flux class gives an exact
split-Cartan \(SU(r)\) theorem for every nonzero
\[
 \mathbf m=(m_1,\ldots,m_r)\in\mathbb Z^r,
 \qquad \sum_i m_i=0,
 \qquad s_2(\mathbf m)=\sum_i m_i^2.
\]
On the regular native fibre define
\[
 \mathscr E_{\mathbf m,T}
 =\bigoplus_{i=1}^r\mathscr L_{\eta,T}^{\otimes m_i},
 \qquad H_{\mathbf m}=i\operatorname{diag}(m_1,\ldots,m_r).
\]
The ordered splitting canonically and parallelly trivializes the determinant,
and, retaining the two opposite orientations used by the native construction,
\[
 c_1(\mathscr E_{\mathbf m,T})=0,
 \qquad
 c_2(\mathscr E_{\mathbf m,T})
 =-\frac{s_2(\mathbf m)}2[\eta_T]^2,
\]
\[
 \left\langle c_2,[F_T]_{\mathrm{mark}}\right\rangle=-6s_2(\mathbf m),
 \qquad
 k_{\mathbf m}:=
 \left\langle c_2,[F_T]_{x\text{-}\mathrm{complex}}\right\rangle
 =6s_2(\mathbf m).
\]
Here \(s_2(\mathbf m)\) is even, so this split family realizes precisely a
subfamily with \(k_{\mathbf m}\in12\mathbb Z_{>0}\).  The diagonal connection
\[
 \nabla^{\mathbf m,T}
 =\mathrm d-2\pi(6y^1\,\mathrm dy^4+y^2\,\mathrm dy^3)H_{\mathbf m},
 \qquad F_{\mathbf m,T}=-2\pi H_{\mathbf m}\eta_T
\]
is a reducible non-flat Yang--Mills critical point and has exact action
\[
 \operatorname{YM}_T(\nabla^{\mathbf m,T})
 =\frac{2\pi^2k_{\mathbf m}}{g_{\mathrm{YM}}^2}
  +\frac{\pi^2s_2(\mathbf m)S_T^2}{g_{\mathrm{YM}}^2D_T}.
\]
For the positive root \(i<j\), put \(q_{ij}=|m_i-m_j|\) and
\[
 \rho_T^*=\frac{\max\{\mathcal B_{T,+},\mathcal B_{T,-}\}}
                  {\min\{\mathcal B_{T,+},\mathcal B_{T,-}\}},
 \qquad
 N_-(T)=\#\{n\ge0:2n+1<\rho_T^*\}.
\]
Its complete negative normalized Coulomb spectrum is
\[
 \nu_{ij,n}(T)=q_{ij}
 \bigl((2n+1)\mathcal B_{T,\min}-\mathcal B_{T,\max}\bigr),
 \qquad 0\le n<N_-(T),
\]
with real multiplicity \(12q_{ij}^2\) before summing coincident numerical
bands.  Consequently
\[
 \operatorname{ind}_{\mathbb R}\mathcal J_{\mathbf m,T}
 =12rN_-(T)s_2(\mathbf m),
\]
and on the deep cusp tail, where \(N_-(T)=3\), this is
\(36r s_2(\mathbf m)=6rk_{\mathbf m}\).  The module also computes the
centralizer harmonic kernel and every odd-frequency wall nullity; negative
eigenspaces are descent directions, whereas only the zero-mode walls are
possible Lyapunov--Schmidt bifurcation parameters.  Finally, naturality of
\(c_2\) proves that this nonzero-flux bundle cannot extend over a carrier
with \(H^4(-;\mathbb Z)=0\), in particular not over a space diffeomorphic to
\(S^6\), \(S^{12}\), or \(S^{24}\).  See
\cref{cusp-native-full-flux-suN-topology,%
cusp-native-full-flux-suN-action-theorem,%
cusp-native-full-flux-suN-Morse-theorem,%
cusp-native-full-flux-suN-nullity-theorem,%
cusp-native-full-flux-suN-stability-boundary}.

\item \Verified\ The Koide relation has an exact ambient
quadratic-form interpretation, and it is rigorously separated from the
regular canonical cusp family.  For
\[
 Q_K(r)=\frac23(r_1+r_2+r_3)^2-(r_1^2+r_2^2+r_3^2),\qquad
 \begin{pmatrix}h\\x\\y\end{pmatrix}
 =\begin{pmatrix}
 1/\sqrt3&1/\sqrt3&1/\sqrt3\\
 1/\sqrt2&0&-1/\sqrt2\\
 1/\sqrt6&-2/\sqrt6&1/\sqrt6
 \end{pmatrix}
 \begin{pmatrix}r_1\\r_2\\r_3\end{pmatrix},
\]
one has \(Q_K(r)=h^2-x^2-y^2\).  The real-linear coefficient-space
isomorphism
\[
 (h,x,y)\longmapsto(t,m,\varrho)
 =\left(h+y,\frac{x}{\sqrt6},h-y\right)
\]
satisfies, with
\[
 (t,m,\varrho)
 =(\operatorname{Im}\tau,\operatorname{Im}\mu,
   \operatorname{Im}\beta),
\]
for the canonical realification of
\(Z=\bigl(\begin{smallmatrix}6\mu&\tau\\\beta&\mu\end{smallmatrix}\bigr)\),
\[
 \det P(\tau,\mu,\beta)
 =t\varrho-6m^2
 =\det\begin{pmatrix}t&\sqrt6\,m\\\sqrt6\,m&\varrho\end{pmatrix}
 =Q_K(r).
\]
For \(r\in\mathbb R_{\ge0}^3\setminus\{0\}\), the Koide equation
\(Q_K(r)=0\) makes this symmetric matrix positive semidefinite of rank
one and makes \(P\) singular.  Its columns therefore do not form a full
lattice and do not define a regular compact complex two-torus.  In contrast,
the actual positive radial cusp has
\[
 \det P_T=-D_T<0,\qquad
 D_T=T^2+(2H-B)T+O(1),\qquad
 \det P_T\sim-T^2.
\]
Thus the result is an isomorphism of ambient coefficient spaces with their
quadratic discriminants, not a map from Koide triples to fibres of the
constructed \(S^6\) family; see
\cref{koide-s6-thm-ambient-discriminant,koide-s6-cor-ambient-only,%
koide-s6-cor-actual-cusp}.

\item \Verified\ Independently of that discriminant comparison, the positive
defect-one shell has an exact integral parametrization and a power-saving
count with its exact constant.  Put
\[
 \Delta(x,y,z)=4(xy+xz+yz)-(x^2+y^2+z^2)
\]
and let \(R(X)\) count triples
\(0<x\le y\le z\le X\) with \(\Delta=-1\) on the terminal Vieta wall
\(z\ge4(x+y)\).  The integral change of variables
\[
 n=z-4(x+y),\qquad q=y-x,\qquad a=x+y-4n
\]
is a bijection onto
\[
 a^2-3q^2-18n^2=-2,\qquad
 a>0,\ q,n\ge0,\ a+4n>q,\qquad z=4a+17n.
\]
Write
\[
 \mathscr X=\{(a,q,n)\in\mathbb A^3:
                 a^2-3q^2-18n^2=-2\}.
\]
The exact local densities are
\[
 \sigma_2=\frac32,\qquad \sigma_3=2,\qquad
 \sigma_p=\frac{1-p^{-2}}{1-(\frac3p)p^{-1}}\qquad(p\ge5),
\]
and the Poincar\'e-residue measure of the real chamber has leading factor
\[
 I_\infty=\frac1{2\sqrt3}\log\!\left(\frac3{2\sqrt2}\right).
\]
After verifying the Gorodnik--Oh--Shah sector hypotheses and Borovoi's
adelic sign coefficient \(c_{\mathscr X}=1\), one obtains unconditionally
\[
 R(X)\sim C_0X,\qquad
 C_0=\frac9{2\pi^2}\log(2+\sqrt3)
       \log\!\left(\frac3{2\sqrt2}\right).
\]
The exact ideal parametrization now proves more.  For \(n\geq1\), put

\[
\begin{aligned}
 m_n&=18n^2-2=2(3n-1)(3n+1),\\
 \mathcal P(m_n)&=\bigl\{(C,D,E)\in\mathbb Z_{>0}^2\times\mathbb Z:\\
 &\hspace{5em}C^2D=m_n,\quad 0\leq E<D,\quad
 E^2\equiv3\pmod D\bigr\}.
\end{aligned}
\]

The map

\[
 (C,D,E)\longmapsto
 C\bigl(D\mathbb Z+(E+\sqrt3)\mathbb Z\bigr)
\]

is a bijection to the ideals of
\(\mathbb Z[\sqrt3]\) of norm \(m_n\).  Each ideal has a unique seed
\(a_0+q_0\sqrt3\) with \(a_0>2q_0\geq0\), and every solution in its orbit is

\[
 a_k+q_k\sqrt3=(a_0+q_0\sqrt3)(2+\sqrt3)^k,\qquad k\geq0.
\]

The inverse is finite and integral: apply
\((a,q)\mapsto(2a-3q,2q-a)\) while \(a\leq2q\), then recover the unique HNF
triple.  Consequently every point has multiplicity one and the number of
unit orbits at level \(n\) is

\[
 \#\mathcal P(m_n)=\rho(m_n)
 =\sum_{d\mid m_n}\chi_{12}(d)
 =\rho(3n-1)\rho(3n+1).
\]
The omitted \(n=0\) face is separate and exact:
\[
 a_k+q_k\sqrt3=(1+\sqrt3)(2+\sqrt3)^k,\qquad k\geq1,
\]
subject to the unchanged strict wall.  Its contribution is the floor count
\(R_0(X)\) of \cref{koide-param-eq:R0-exact}, and
\(R_0(X)=O(\log(2+X))\).

The exact sharp-height floor formula and its Hecke-phase Fourier expansion
are written in
\cref{koide-param-thm:integral-parametrization,koide-param-thm:exact-floor}.
A direct verification of Benoist--Oh effective well-roundedness for the
literal five-wall family, followed by their
Theorem~13.2~\cite[pp.~1931--1932]{BenoistOh_2012} on each of the
finitely many integral spin-group orbits, strengthens the asymptotic to

\[
 \boxed{R(X)=C_0X+O(X^{1-\eta})}
 \qquad\text{for some fixed }\eta>0.
\]

Neither Benoist--Oh's theorem nor its proof supplies
\(\eta\geq\tfrac12\).  The primary-source comparison is exact:
Liu--Sarnak give the smooth standard-radial exponent \(39/64\) for this
congruence arithmetic lattice, while Good and Chatzakos--Petridis give the
sharp standard-radial \(2/3\) frontier after the relevant small-spectrum
terms are displayed.  Neither theorem counts the unchanged five-wall region
with height \(4a+17n\), and neither proves its pointwise square-root
discrepancy; see \cref{koide-lit-determination}.

The exact census gives
\(R(10^8)=3{,}536{,}002\), after an independent six-coordinate replay of
all \(3{,}536{,}002\) rows.  Its largest
\(\lvert R(X)-C_0X\rvert/\sqrt X\) on
\(10^4\leq X\leq10^8\) is
\(0.1714001928088145597\ldots\), at \(X=867{,}240\); this is finite evidence
only.
The stronger proposed estimate
\(R(X)-C_0X=O_\varepsilon(X^{1/2+\varepsilon})\) is neither asserted nor
used in this status item; the proved conclusion is the displayed linear
power-saving asymptotic.  See
\cref{koide-count-thm-coordinate-bijection,koide-count-thm-local-densities,%
koide-count-thm-archimedean-factor,koide-count-thm-linear-asymptotic,%
koide-count-sec-census-1e8,koide-param-thm:power-saving,%
koide-lit-determination,%
koide-count-open-square-root-discrepancy}.

\item \Verified\ The same parity-restricted shell has an exact quaternionic
realization.  In

\[
 B=(18,-6)_{\mathbb Q}
 =\mathbb Q\langle I,J\rangle/(I^2-18,J^2+6,IJ+JI),\qquad K=IJ,
\]

put

\[
 \mathcal O=\mathbb Z1\oplus
 \mathbb Z\frac{I+J}{2}\oplus
 \mathbb Z\frac{I-J}{2}\oplus
 \mathbb Z\frac K2.
\]

The complete multiplication table proves that this is an order.  Its
reduced discriminant is \(108\), so it has index \(18\) in a maximal
overorder; the integer \(18\) is not an Eichler level, since the order is
nonmaximal at the ramified primes \(2\) and \(3\).  The algebra is a division
algebra over \(\mathbb Q\), split over \(\mathbb R\), and ramified exactly at
\(\{2,3\}\).  The bijection

\[
 (a,q,n)\longmapsto
 \gamma(a,q,n)=2+\frac{qI+aJ+nK}{2}
\]

identifies \(a^2-3q^2-18n^2=-2\), with its forced parity
\(a\equiv q\pmod2\), with the reduced-trace-\(4\), reduced-norm-\(1\) locus
of \(\mathcal O\).  Every such \(\gamma\) is primitive hyperbolic, with

\[
 \mathcal O\cap\mathbb Q(\gamma)
 =\mathbb Z[\gamma-2]\cong\mathbb Z[\sqrt3],\qquad
 C_{\mathcal O^1}(\gamma)=\{\pm\gamma^k:k\in\mathbb Z\},
\]

eigenvalues \(2\pm\sqrt3\), and translation length
\(2\log(2+\sqrt3)\).  The full shell is a finite union of
\(\mathcal O^1\)-conjugacy classes.  For the displayed real splitting, the
standard radial Huber parameter is exactly

\[
 X_i(\gamma)^2=1+\frac{(7a+15\sqrt2\,n)^2}{48},
\]

whereas the Koide height is \(4a+17n\).  Thus the hyperbolic reduction is
exact, but the radial theorem does not transfer a discrepancy estimate to
the sharp Koide polyhedron; see
\cref{koide-quat-exact-quaternion-bridge-order-shell-and-hyperbolic-conjugacy-classes,%
koide-quat-exact-reformulation-of-the-koide-chamber-count,%
koide-quat-consequence-for-the-open-square-root-discrepancy}.

\item \Verified\ The fixed-prime Erd\H{o}s--Straus divisor selector is an
exact marked slice of the candidate's terminal integral presentations.  For
\(p\equiv1\pmod{12}\) prime, put
\[
 L=\frac{p-1}{4},\qquad R_j=4j+3,\qquad
 a_j=\frac{p+R_j}{4},\qquad 0\le j<L.
\]
If \(u\in\mathbb Z_{>0}\), \(u\mid a_j^2\),
\(\varepsilon\in\{0,1\}\), and
\(R_j\mid4u+p^\varepsilon\), then
\[
 T=\frac{4u+p^\varepsilon}{R_j},\qquad
 (\ell_0,\ell_1,\ell_2)
 =\left(\frac{T+1}{4},u-jT,1\right)
\]
is integral and satisfies
\[
 12\ell_0-4\ell_1-3\ell_2=p^\varepsilon,
\]
with the exact relation-matrix identity
\[
 \begin{pmatrix}1&-j\\0&1\end{pmatrix}
 \begin{pmatrix}R_j&-u\\4&-T\end{pmatrix}
 =\begin{pmatrix}3&-\ell_1\\4&1-4\ell_0\end{pmatrix}.
\]
The inverse is \(T=4\ell_0-1\), \(u=\ell_1+jT\).  With \(p,j\), the
channel \(\varepsilon\), positive orientation, and the divisor gate retained,
this is a bijection between the complete two-channel selector and the
normalized terminal slice.  It proves an exact integral-relator crosswalk,
not a geometric morphism from the reciprocal cubic to \(X\); see
\cref{sec:es-s6-crosswalk,thm:es-s6-selector-terminal-bijection}.

\item \Verified\ The modulo-$840$ refinement has
$[N:R_{\mathrm{ES}}]=72$, $\det N=1$, and root system $A_5^4D_4$.
Theorems~\ref{esn:thm:niemeier}, \ref{esn:thm:d4-f4-triality},
\ref{esn:fa:thm:full-quotient}, and~\ref{esn:fa:thm:residue-weyl}
prove the glue, marked flag-weight comparison, full Weyl quotient and
position of the commuting residue six-cycle. For a prime
$p\equiv1\pmod{24}$ and $R>0$, $R\equiv3\pmod4$, set
$a=(p+R)/4$, $g=\gcd(R,pa)$, $R_0=R/g$, $N_0=pa/g$.
The exact divisor-to-ordered-pair map is
\[
 d>0,\quad d\mid N_0^2,\quad d\equiv-N_0\pmod{R_0}
 \quad\longmapsto\quad
 (b,c)=\left(\frac{N_0+d}{R_0},
                  \frac{N_0+N_0^2/d}{R_0}\right).
\]
It is a bijection onto the positive remaining pairs solving
$4/p=1/a+1/b+1/c$; the unordered count is
$U_R(p)=(A_R(p)+\mathbf1_{R_0\mid2})/2$
(Theorem~\ref{esn:thm:shell-bijection}).
Theorem~\ref{esn:thm:sample-fourier} gives
$Q_{R,P}=6W_{R,P}-M_{R,P}u\in A(S_{840})\subset N$ for the
retained finite sample $P$ of primes reducing into
$S_{840}=\langle289\rangle$. This is not a residue-only existence test.

\item \Verified\ The two-column \(SU(3)\) Gram quotient is the
four-dimensional positive Hermitian cone:
\[
 \operatorname{Mat}_{3\times2}(\mathbb C)/SU(3)
 \cong\operatorname{Herm}_2(\mathbb C)_{\geq0}.
\]
Its one-point compactified representation sphere satisfies
\[
 S^{12}_{\rm col}/SU(3)\cong B^4,
\]
with orbit-type boundary \(S^3\).  This \(S^{12}_{\rm col}\) is the column
representation sphere, not the cusp-smash \(S^{12}\); see
\cref{thm:hrgc-gram-quotient,thm:hrgc-s12-ball-quotient}.

\item \Verified\ The three-column orbit calculation produces an exact
rank-three carrier, not a conjectural dimensional analogy.  For
\(\mathscr V_3=\operatorname{Mat}_{3\times3}(\mathbb C)\),
\[
 \mathscr V_3/U(3)\cong\operatorname{Herm}_3(\mathbb C)_{\geq0}
\]
by the complete Gram invariant \(A^\dagger A\), whereas strict \(SU(3)\)
retains the alternating invariant:
\[
 \mathscr V_3/SU(3)
 \cong\{(H,z):H\geq0,\ |z|^2=\det H\}\cong\mathbb R^{10}.
\]
Thus
\[
 S^{17}/U(3)\cong B^8,\quad
 \mathscr V_3^+/U(3)\cong B^9,\quad
 S^{17}/SU(3)\cong S^9,\quad
 \mathscr V_3^+/SU(3)\cong S^{10}.
\]
The compactified rank-loss locus is a standard topologically locally flat
\(S^8\subset S^{10}\), and its complement is
\[
 S^{10}\setminus S^8\cong\mathbb R^8\times\mathbb C^*\simeq S^1,
\]
whose positive winding generator is \(\det A/|\det A|\).  The rank-one
trace-one stratum inside the rank-loss sphere is \(\mathbb CP^2\).  The
rank-three determinant is cubic; this cone is not the quadratic Lorentz cone
arising from two columns.  See
 \cref{thm:hr3c-u3-compactification,thm:hr3c-su3-invariants,%
 thm:hr3c-su-full-cone,thm:hr3c-rank-loss-pair,thm:hr3c-two-ladders}.

\item \Verified\ The selectively corrected Jordan determinant/Euler module supplies an exact
algebraic refinement, but no new geometric globalization.  In the fixed
orientation, every Hermitian matrix has a unique form
\(S+i[a]_{\times}\), and
\[
 \det\bigl(S+i[a]_{\times}\bigr)=\det S-a^{\mathsf T}Sa.
\]
For every marked positive-definite symmetric form \(B\), the linear section
\[
 \ell_B(\lambda,a)=\lambda B+i[a]_{\times}
\]
has the exact positive inverse image
\[
 \ell_B^{-1}\!\left(\operatorname{Herm}_3(\mathbb C)_{\geq0}\right)
 =\{\lambda\geq0:a^{\mathsf T}Ba\leq\lambda^2\det B\},
\]
a future \((3+1)\)-Lorentz cone.  The explicit order-cone lift
\(\Theta:\operatorname{Herm}_2(\mathbb C)\to
\mathbb RI_3\oplus i\mathfrak{so}(3)\) satisfies
\[
 \det_3\Theta(H)=\frac{\operatorname{tr}H}{2}\det_2H
\]
and sends positive-semidefinite ranks \(2,1,0\) to \(3,2,0\).
For a rank-one projection \(P\), the matrix \(\Theta(P)\) has eigenvalues
\(0,1/2,1\) and is not idempotent, which separately proves that \(\Theta\)
is not a Jordan homomorphism.  Neither \(\Theta\) nor \(\ell_B\) is a
toric, period-lattice, or filling map.  The claimed index-three translation
compression is also rejected: the proved six-to-two operation is the
rank-lowering map \(P:\mathbb Z^3\twoheadrightarrow\mathbb Z^2\), with the
three even and three odd Kuhn labels mapping to the two target triangle
orbits.  See
\cref{thm:teb-hermitian-determinant-splitting,%
prop:teb-theta-order-lift,prop:teb-fan-morphism}.

\item \Verified\ The rank-three toroidal successor and its exact quotient to
the reconstructed cusp are now constructed.  The local analytic family
\[
 \pi_3:N_3\longrightarrow\Delta
\]
has smooth fibre \(E_\tau^3\), deepest local equation
\(t=z_0z_1z_2z_3\), cusp monodromy
\[
 M_3=\begin{pmatrix}I_3&I_3\\0&I_3\end{pmatrix},
\]
and filling homomorphism \(a_i\mapsto0,\ b_i\mapsto e_i\).  Its irreducible
central fibre has
\[
 (f_0,f_1,f_2,f_3)=(1,7,12,6),\qquad \chi(W_3)=6.
\]
Its normalization is the permutohedral threefold with orbit vector
\((1,14,36,24)\); the positive-codimension strata map with degrees
\(2,3,4\) onto the \(7,12,6\) quotient strata.  Hence a real
\(S^8\)-total-space fibration retaining \(W_3\) requires its remaining
singular fibres to contribute Euler sum \(-4\).
For a full-rank integral shear
\(\mathcal B:\bar L_3\hookrightarrow L_3\), the exact generalization is
\[
 \chi(W_{3,\mathcal B})
 =3!\,[L_3:\mathcal B\bar L_3]
 =3!\,|\det\mathcal B|.
\]

In the candidate decomposition
\(\Lambda_F=A_F\oplus B_F\), define
\[
 P=\begin{pmatrix}1&-1&0\\0&1&-1\end{pmatrix},
 \qquad
 Q=\begin{pmatrix}0&-1&1\\1&-1&0\end{pmatrix}.
\]
Then
\[
 \Phi=P\oplus Q:H_1(T^6;\mathbb Z)
 \twoheadrightarrow H_1(T^4;\mathbb Z)
\]
obeys the literal monodromy, filling, and period identities
\[
 \Phi M_3=M_F\Phi,\qquad c_F\Phi=Qc_3,\qquad
 PZ_3(t)=Z_2(t)Q.
\]
It extends over the toric strata to a holomorphic map of cusp families with
connected elliptic kernel.  Moreover an explicit \(J\in SL_6(\mathbb Z)\)
gives
\[
 JM_3J^{-1}\cong
 M_F\oplus\begin{pmatrix}1&1\\0&1\end{pmatrix},
\]
 and the candidate's free order-three and order-four peripheral actions lift
 on the added elliptic factor.

The latter successful lifts are distinct from the diagonal candidate
\[
 M_1=I_3\otimes S,\qquad M_2=I_3\otimes Q,\qquad M_0=I_3\otimes T,
 \qquad M_1M_2M_0=I.
\]
Although its finite factors have orders \(4\) and \(6\), the maps
\(I-M_j:T^6\to T^6\) are isogenies for \(j=1,2\), so every affine lift of
either finite action has a fixed point.

The six source tetrahedron orbits also map \(S_3\)-equivariantly to the six
marked diagonal \(A_2\) Weyl chambers: with
\(\nu_i=E_{ii}-I_3/3\), the \(\pi\)-th cell maps to the cone
\[
 \mathbb R_{\geq0}\nu_{\pi_1}
 +\mathbb R_{\geq0}(-\nu_{\pi_3})
 =\{\lambda_{\pi_1}\geq\lambda_{\pi_2}
                    \geq\lambda_{\pi_3}\}.
\]
This is an exact fan/Cartan correspondence, not a geometric identification
with the Gram determinant singularity.  The global period-and-overlap gluing
over both finite-order discs remains \Open.  See
\cref{hrftd:thm:toroidal-cusp,hrftd:thm:strata,%
hrftd:cor:euler-shear-determinant,thm:teb-holomorphic-cusp-bridge,%
 prop:teb-explicit-conjugacy,prop:teb-finite-monodromy-lift,%
 prop:teb-weyl-chambers}.

\item \Verified\ The \(\mathbb Z^6\) native to the higher-rung cusp is
\[
 \Lambda_{\rm tor}=H_1(T^6;\mathbb Z)
 =L_3\oplus\overline L_3,
\]
not the cyclic group \(\mathbb Z/6\) and not the separate group
\(H^4(U;\mathbb Z)\).  For \(k=(1,1,1)\), the exact sequence
\[
 0\longrightarrow \mathbb Z(k,0)\oplus\mathbb Z(0,k)
 \longrightarrow\Lambda_{\rm tor}
 \xrightarrow{\,P\oplus Q\,}A_F\oplus B_F\longrightarrow0
\]
is \(S_3\)-equivariant and compatible with monodromy.  It splits as an
integral monodromy lattice but not as a \(\mathbb Z[S_3]\)-module.  Each
quotient copy is the \(A_2\) weight lattice; its root lattice has index three.
The six standard generators form two \(S_3\)-orbits of size three, whereas
the six Kuhn/Weyl chambers form the regular \(S_3\)-orbit.  Thus
\[
 \operatorname{rk}\Lambda_{\rm tor}=6,\qquad |S_3|=3!=6,
 \qquad\chi(W_{3,\mathcal B})=3!|\det\mathcal B|
\]
are three precisely related but nonidentical statements.  The actual map to
the Gram boundary is
\(\Phi_{\rm G}\circ\jmath\circ P\circ\operatorname{pr}_{\rm ang}\): it is
an \(S_3\)-equivariant, piecewise-linear chamber map on the real
cocharacter space, not an additive map, a toric fan morphism, or a map of
torus fibres.

The regular chamber orbit itself has an exact rank-nine realization.  With
\(\mathcal M_{\mathbb Z}=\operatorname{Hom}(\overline L_3,L_3)\),
\[
 \mathcal M_{\mathbb Z}
 =\bigoplus_i\mathbb ZE_{ii}\oplus
  \bigoplus_{i\ne j}\mathbb ZE_{ij},
\]
the second summand is \(\mathbb Z[S_3]\), and
\(c_{(\pi_1\pi_2\pi_3)}\mapsto E_{\pi_1,\pi_3}\) identifies its basis
with the six chambers.  These are precisely the six mixed basis pairs with
zero alternating pairing; the three diagonal matrix units encode the
excluded pairs with pairing one.  The real matrix carrier is
\(S_3\)-equivariantly isomorphic to \(\operatorname{Herm}_3(\mathbb C)\),
but the integral symmetric/skew splitting has index \(8\), and every
off-diagonal matrix unit maps to an indefinite Hermitian matrix.  Hence this
is an exact representation and incidence theorem inside the
nine-dimensional carrier, not a positive-semidefinite or toroidal map.  See
\cref{rsl:thm:s3-sequence,rsl:thm:three-sixes,rsl:prop:gram-scope,%
rnmt:thm:regular-torsor,rnmt:prop:representation-bridge}.

\item \Verified\ The marked angular epimorphism now gives an exact positive
observation morphism from the full rank-nine carrier to Minkowski space,
while retaining and classifying every fibre direction.  Put
\[
 P=\begin{pmatrix}1&-1&0\\0&1&-1\end{pmatrix},\quad
 u=(1,1,1)^{\mathsf T},\quad C=PP^*,\quad
 R=P^*C^{-1},\quad \Pi=uu^*/3,\quad X=(R\;u).
\]
Then \(PR=I_2\), \(RP=I_3-\Pi\), and \(\det X=1\).  The maps
\[
 \mathcal O_P(H)=PHP^*,\qquad \jmath_P(G)=RGR^*
\]
satisfy
\[
 \mathcal O_P\jmath_P=\operatorname{id},\qquad
 \jmath_P\mathcal O_P(H)=(I_3-\Pi)H(I_3-\Pi),
\]
and every \(H\in\operatorname{Herm}_3(\mathbb C)\) has unique coordinates
\[
 H=X\begin{pmatrix}G&\psi\\\psi^*&\lambda\end{pmatrix}X^*,\qquad
 \det H=\lambda\det G-\psi^*\operatorname{adj}(G)\psi.
\]
The complete positive fibre over \(G\succeq0\) is
\[
 \psi\in\operatorname{im}G,\qquad
 \lambda\geq\psi^*G^+\psi,
\]
and the zero fibre is exactly \(\{\lambda uu^*: \lambda\geq0\}\).
For \(g\in SL_2(\mathbb C)\),
\[
 T_P(g)=RgP+\Pi=X(g\oplus1)X^{-1}
\]
acts by \((G,\psi,\lambda)\mapsto(gGg^*,g\psi,\lambda)\).  Thus
\[
 \operatorname{Herm}_3(\mathbb C)
 \cong \text{Lorentz four-vector}\oplus
 \text{defining complex spin representation}\oplus\text{real scalar}.
\]
The vector action has kernel \(\{\pm I_2\}\), but the full action is
faithful because \(-I_2\) sends \(\psi\) to \(-\psi\).  Moreover
\[
 e_2(\jmath_P(G))=\tfrac13\det G,\qquad
 \det(uu^*+\varepsilon\jmath_P(G))=\varepsilon^2\det G.
\]
Consequently \(\jmath_P\) extends properly to an embedding
\(B^4\hookrightarrow S^8\); the reverse observation on the rank-loss
boundary is nonproper and has no infinity-preserving extension.  In the
strict \(SU(3)\)-quotient its determinant phase obeys
\[
 |\omega|^2=\lambda\det G-
 \psi^*\operatorname{adj}(G)\psi.
\]
Finally, for \(Z\in\operatorname{Mat}_{3\times2}(\mathbb C)\) and retained
\(\alpha,c>0\), the exact colour-invariant Lorentz momentum and mass are
\[
 p_\alpha(Z)=\alpha Z^*Z,\qquad
 m_{\alpha,c}^2(Z)=\frac{\alpha^2}{c^2}
 \sum_{a<b}|Z_{a1}Z_{b2}-Z_{a2}Z_{b1}|^2.
\]
Rank two is equivalent to positive mass, rank at most one to zero mass, and
\(m_{\alpha,c}^2(rZ)=r^4m_{\alpha,c}^2(Z)\); hence this unrestricted
kinematic carrier is gapless.  Its exact presymplectic reduction is now
complete.  For \(m_0>0\),
\[
 (\mathcal J_2\times\mathcal H^+_{m_0,c})/
 ((x,p)\sim(x+sp,p))\cong_{\rm symp}T^*\mathcal H^+_{m_0,c},
\]
so the rank-two system is exactly the future-timelike spinless particle.  For
\(m_0=0\), the intrinsic punctured rank-one reduction is
\(T^*\mathcal N_+^\times\), but the composite multiplier action is irregular:
\[
 d(\det\circ p_\alpha)_Z=0\quad
 (\operatorname{rank}_{\mathbb C}Z\leq1),
\]
so \(\dot x=\lambda p\) does not imply \(\lambda=e\); four transverse real
rank-opening directions are first-order null without being colour gauge, and
at \(Z=0\) every \(x(\tau)\) is stationary.  On the hard
\((\lambda_0,\omega_0)\)-constraint with \(\omega_0\ne0\), the momentum map
has rank four at \(w\ne0\) and rank three at \(w=0\).  Thus its own action
forces \(\dot x=0\) above the sharp threshold and permits
\(\dot x\in\mathbb Rp\) only at the threshold.  Its determinant inequality
and multiplication/Fock spectrum are exact, but the latter is not the
Hamiltonian of that action and gives a constraint-selected free spectral gap,
not a dynamically generated interacting gap.  See
\cref{hlob:thm:split-observation,hlob:thm:block-coordinates,%
 hlob:thm:lorentz-action,hlob:thm:cone-section,hlob:cor:Y3-block-form,%
 hlob:thm:colour-momentum,wpr:thm:massive,wpr:thm:massless,%
 wpr:thm:zero,wpr:thm:packet-action}.

\item \Verified\ The determinant normal slice has now been promoted from an
exact Lorentzian analogy to a completely typed host for explicit field
models, while the additional choices and the spectral limitations remain
visible.  With
\[
 \mathcal N_P=\{Y\in\operatorname{Herm}_3(\mathbb C):Yu=0\},\qquad
 \jmath_P(G)=RGR^*,
\]
the map \(\jmath_P:\operatorname{Herm}_2(\mathbb C)\to\mathcal N_P\) is a
real-linear isomorphism with inverse \(Y\mapsto PYP^*\), and
\[
 \det(uu^*+\varepsilon\jmath_P(G))=\varepsilon^2\det G.
\]
The polarized determinant is exactly the Minkowski form
\(tt'-xx'-yy'-zz'\), so the affine slice carries a constant metric of
signature \((1,3)\), its positive-semidefinite part is the future cone, and
the lifted congruence action realizes
\[
 PSL_2(\mathbb C)\cong SO^+(1,3).
\]
The standard Hermitian metrics used to define the Moore--Penrose section
select
\[
 W_P=R^*R=C^{-1},\qquad
 K_P=\{a\in SL_2(\mathbb C):a^*W_Pa=W_P\}\cong SU(2),
\]
and the positive invariant form
\(E_P(G,K)=\operatorname{tr}(GW_PKW_P)\).  The real normal representation is
\(\boldsymbol1\oplus\boldsymbol3\); its centre acts trivially.  Promoting
\(K_P\) to a local gauge symmetry gives the explicit nonnegative action
\[
 \int_M\left(
 \frac{|F_{\mathcal A}|^2}{2g_{\rm YM}^2}
 +\frac{Z_G}{2}|D_{\mathcal A}G|_{E_P}^2
 +\frac{m_G^2}{2}|G|_{E_P}^2
 +\frac{\lambda_4}{4}(\det G)^2\right)d\operatorname{vol}_g,
\]
and the lift \(\widehat{\mathcal A}=R\mathcal A P\) satisfies
\(D_{\widehat{\mathcal A}}\jmath_P(G)=
\jmath_P(D_{\mathcal A}G)\).  This gauging is an additional physical
construction; the normal field cannot screen a fundamental centre charge
without the centre-nontrivial retained spinor \(\psi\) or another added
matter field.

The candidate-specific separation is now exact rather than verbal.  The
retained inverse-dual cusp matrix obeys
\[
 G_{\eta,T}=\begin{pmatrix}2L_T&2m_T\\2m_T&-q_T/3\end{pmatrix},
 \qquad \det G_{\eta,T}=-\frac23D_T<0,
\]
whereas \(p_\alpha(Z)=\alpha Z^*Z\) ranges over the entire
positive-semidefinite cone.  Positive rescaling and complex congruence cannot
change this determinant sign/inertia, so a further morphism is required to
place these two matrices in one physical model.  On the Lorentzian host, the
free scalar Fock Hamiltonian has exact gap \(m_{\mathrm{KG}}\) when the
independently supplied coefficient is positive and is gapless when it is
zero.  For the massless compact gauge Gaussian, the lowest nonzero transverse
frequency is \(2\pi/L\), which closes as \(L\to\infty\).  Entangled vacua and
sectors occur in both explicitly gapped and explicitly gapless Hamiltonians.
Hence none of the determinant identity, the Gram carrier, or entanglement
alone proves a nonlinear four-dimensional mass gap.  The full-domain
comparison theorem identifies a sufficient exact bridge: if
\(\|Ru\|\geq c_R\|u\|\) and
\(\mathfrak b[Ru]\leq C_R(\mathfrak a_{\rm phys}[u]-
E_{\rm vac}\|u\|^2)\) throughout the physical excitation form domain and
\(\mathfrak b\geq\gamma\), then
\[
 \operatorname{gap}(A_{\rm phys})\geq\frac{\gamma c_R^2}{C_R}.
\]
An exact counterexample proves that an intertwiner defined on only one gapped
subspace is insufficient.  See
\cref{hnsp-normal-identification,hnsp-lorentz-theorem,hnsp-spin-theorem,%
hnsp-intrinsic-compact-representation,hnsp-intrinsic-gauge-scalar-theorem,%
hnsp-centre-screening-corollary,hnsp-candidate-spacelike,%
hnsp-free-field-gap,hnsp-gauge-gaussian-gap,hnsp-coercive-transfer}.

\item \Refuted\ A literal Fabel--Alp\"oge-style complex-fourfold
construction on a homotopy \(S^8\) cannot exist.  If
\(E=(TS^8,J)\) were complex of rank four, then
\[
 c_1(E)=c_2(E)=c_3(E)=0,\qquad c_4(E)=2u,
\]
and therefore \(p_2(TS^8)=2c_4(E)=4u\).  The eight-dimensional signature
formula, with \(p_1=0\) and signature zero, forces \(p_2=0\), a
contradiction.  A nonsingular bundle
\(T^6\to S^8\to S^2\) is also impossible: its homotopy sequence would
require a surjection \(\mathbb Z\to\mathbb Z^6\), and its Euler
characteristic would give \(2=0\).  A merely real singular fibration is not
refuted, but it must satisfy
\(\sum_i\chi(F_i)=2\) and the exact van Kampen, Wang, and integral
Mayer--Vietoris maps proved in the workbench.  For regular-piece monodromy
\(\rho\), the corrected homological Serre edge sequence is
\[
 0\longrightarrow(\textstyle\bigwedge^q\Lambda)_\rho
 \longrightarrow H_q(E;\mathbb Z)
 \longrightarrow
 H_1(B^\circ;\textstyle\bigwedge^{q-1}\Lambda)
 \longrightarrow0,
\]
whereas dually \(H^q(E;\mathbb Z)\) surjects onto the invariant group
\((\bigwedge^q\Lambda^\vee)^\rho\).  Thus every filling map must respect
the coinvariant subobject and the base-one-cycle quotient; the reversed
edge maps cannot be used in the Mayer--Vietoris test.  Every effective torus
action on \(S^8\) has rank at most four, and the sharp coordinate model is
\[
 S^8/T^4\cong\operatorname{Susp}(\Delta^3)\cong B^4.
\]
It is a real two-pole singular compactification with generic \(T^4\) orbit,
not a three-special-fibre \(T^6\)-fibration over \(S^2\).  See
\cref{hrs8:subsec:no-almost-complex,hrs8:eq:euler-fibres,%
 hrs8:eq:van-kampen-fillings,hrs8:eq:mv-tests,%
 hrs8:eq:serre-regular-part,hrs8:eq:serre-cohomology-invariants,%
 hrs8:eq:torus-rank,hrs8:eq:t4-quotient}.

\item \Verified\ The original reconstructed \(S^6\) carries a separate,
exact six-coordinate bundle problem on its regular locus.  With
\[
 D=W\sqcup S_1\sqcup S_2,
 \qquad U=X\setminus D,
\]
Alexander duality and the computed homology of the three special fibres give
\[
 H^4(U;\mathbb Z)
 \cong H_1(W;\mathbb Z)\oplus H_1(S_1;\mathbb Z)
       \oplus H_1(S_2;\mathbb Z)
 \cong\mathbb Z^2\oplus\mathbb Z^2\oplus\mathbb Z^2.
\]
Because \(U\) has CW dimension at most five, principal \(SU(3)\)-bundles
on \(U\) are classified exactly by their complete
\(c_2\in H^4(U;\mathbb Z)\).  The connecting map
\[
 \delta:H^4(U;\mathbb Z)\xrightarrow{\sim}H^5(X,U;\mathbb Z)
\]
shows that such a bundle extends across all three deleted fibres if and only
if \(c_2=0\).  This is the free lattice \(\mathbb Z^6\), not the cyclic
group \(\mathbb Z/6\mathbb Z\).

\item \Verified\ The primitive fibre-degree map, its saturated kernel, the
finite normal-holonomy quotients, and the common-centre matrix maps now have a
complete theorem--proof module in
\cref{zdc:sec:bridge,zdc:thm:fibre-degree,zdc:prop:finite-reductions,%
zdc:thm:tensor-centres,zdc:thm:common-centre}.  The result is exact but stops
before any gauge-bundle, connection, confinement, or spectral assertion.

In the marked basis
\[
 (a_0,b_0,a_1,b_1,a_2,b_2)
 =(\bar\gamma,\bar u,g_1,h_1,g_2,h_2),
\]
solving the transpose of the verified one-shot collar matrix gives
\[
 \operatorname{res}_F\!\left(\sum_{i=0}^2(n_i a_i+r_i b_i)\right)
 =(12n_0+4n_1-3n_2)\omega_F,
 \qquad [\deg_F]=(12,0,4,0,-3,0),
\]
and
\[
 \ker\deg_F
 =\mathbb Z\langle
 b_0,b_1,b_2,a_0-3a_1,3a_1+4a_2\rangle.
\]
The map is primitive, so restriction to the oriented fibre is surjective.
Define
\[
 q_{\rm col}:L\longrightarrow H_1(X^\circ;\mathbb Z),\qquad
 a_0\mapsto c,\quad a_1\mapsto\bar g_1,\quad a_2\mapsto\bar g_2,\quad
 b_0,b_1,b_2\mapsto0.
\]
If \(\Phi_1:H_1(X^\circ;\mathbb Z)\to\mathbb Z\) has row
\((12,4\ell_1,3\ell_2)\) in the basis \((c,\bar g_1,\bar g_2)\), then the
factorization through the actual collar quotient is
\[
 \deg_F=\Phi_1\circ q_{\rm col},\qquad
 \Phi_1(j_*[\widetilde a_0])=p.
\]

Let \(\vartheta_1(a_1)=e^{-2\pi i/3}\) and
\(\vartheta_2(a_2)=e^{-2\pi i/4}\) be the normal-line holonomy characters,
trivial on the other marked generators.  The candidate twists select
\[
 \chi_{12}\!\left(\sum_{i=0}^2(n_i a_i+r_i b_i)\right)
 =\vartheta_1(n_1a_1+r_1b_1)
  \vartheta_2(n_2a_2+r_2b_2)^{-1}
 =e^{-2\pi i(4n_1-3n_2)/12}.
\]
For \(\rho_m(k)=e^{-2\pi ik/m}\), the exact typed identities are
\[
 \boxed{\chi_{12}=\rho_{12}\circ\deg_F,\qquad
        \chi_6:=\chi_{12}^2=\rho_6\circ\deg_F.}
\]
Consequently
\[
 H^4(U;\mathbb Z)/\ker\chi_{12}\cong\mathbb Z/12\mathbb Z,
 \qquad
 H^4(U;\mathbb Z)/\ker\chi_6\cong\mathbb Z/6\mathbb Z,
\]
whereas ordinary coefficient reduction gives
\[
 H^4(U;\mathbb Z/6)\cong(\mathbb Z/6)^6.
\]
Thus this six-coordinate coefficient group is not the single distinguished
cyclic quotient selected by \(\chi_6\).  A principal \(SU(3)\)-bundle
\(P\to U\) satisfies
\[
 \chi_6(c_2(P))
 =\exp\!\left(-\frac{2\pi i}{6}
 \langle c_2(P)|_F,[F]\rangle\right).
\]
The order-six quotient is therefore the reduction of a geometrically
computed degree and of the combined local normal holonomy, not a torsion
subgroup inferred from the rank-six notation.

Tensor product gives an independent group-theoretic match:
\[
 \iota:SU(2)\times SU(3)\hookrightarrow SU(6),\qquad
 (A,B)\longmapsto A\otimes B,
 \qquad
 \mu_2\times\mu_3\xrightarrow{\sim}Z(SU(6))=\mu_6.
\]
Since \(U\) has homotopy dimension at most five, \(c_2\) classifies
\(SU(n)\)-bundles for \(n\ge3\).  For determinant-trivial bundles of ranks
two and three,
\[
 c_2(E_2\otimes E_3)=3c_2(E_2)+2c_2(E_3).
\]
Thus every chosen \(SU(6)\)-bundle class
\(h\in H^4(U;\mathbb Z)\) admits a tensor-subgroup reduction, with all
\(c_2\)-factorizations
\[
 (e,f)=(h+2t,-h-3t),\qquad t\in H^4(U;\mathbb Z).
\]
Neither \(h\) nor a reduction is preferred.  The separate support group is
\[
 H^5(U;\mathbb Z)
 \cong\ker\{\mathbb Z^3\xrightarrow{(1,1,1)}\mathbb Z\}
 \cong Q(A_2),
\]
for the combinatorial form making the three distinguished support basis
vectors orthonormal; this is not a cup-product or intersection form.  It is
not an intrinsic \(A_2\otimes\mathbb Z^2\) inside \(H^4(U)\), and the
rank-five kernel of \(\deg_F\) is not thereby an \(A_5\) root lattice.

The finite-group isomorphism
\(H^4(U)/\ker\chi_6\cong Z(SU(6))\) supplies neither transition functions
nor a connection.  The injective tensor map does not factor through
\((SU(2)\times SU(3))/\mu_6\), because its central \(\mu_6\) maps to
nontrivial scalar matrices.  The separate common-centre construction is
\[
\begin{aligned}
 \Gamma_6
 &=\bigl\{((\lambda\mu)^{-1},\lambda I_2,\mu I_3):
             \lambda\in\mu_2,\ \mu\in\mu_3\bigr\},\\
 \frac{U(1)\times SU(2)\times SU(3)}{\Gamma_6}
 &\lhook\joinrel\longrightarrow U(6),\\
 [z,A,B]&\longmapsto z(A\otimes B),
\end{aligned}
\]
where \(\Gamma_6\cong\mu_6\).  Its determinant is \(z^6\), and its image
meets \(SU(6)\) in exactly \(\iota(SU(2)\times SU(3))\); the image is not
generally contained in \(SU(6)\).  No present morphism identifies this degree-four quotient
with the degree-two lifting obstruction of a common-centre gauge quotient.
Global \(SU(3)\)-bundles on \(S^6\) are instead classified by
\(k\in\pi_5(SU(3))\cong\mathbb Z\), with \(c_3=\pm2ku\), and all restrict
trivially to \(U\).  See
\cref{hrs8:eq:h4-regular-locus,hrs8:eq:su3-classification,%
hrs8:eq:relative-delta,hrs8:eq:extension-criterion,%
hrs8:eq:c3-clutching}.

\item \Verified\ The exact operator-theoretic bridge required for any
physical no-gap inference is proved as a typed implication.  Let \(H\ge0\)
be self-adjoint on \(\mathcal H_{\mathrm{phys}}\), put
\[
 P_0=\mathbf1_{\{0\}}(H),\qquad
 \mathcal H_{\mathrm{exc}}=(I-P_0)\mathcal H_{\mathrm{phys}},\qquad
 \mathfrak h[\psi]=\lVert H^{1/2}\psi\rVert^2,\qquad
 \mathcal Q(H)=D(H^{1/2}).
\]
For a nonnegative quadratic form \(\mathfrak q_n\) and nonzero
\(u_n\in D(\mathfrak q_n)\subset\mathcal K_n\), and for
\(J_n:\operatorname{span}\{u_n\}\to\mathcal Q(H)\), define
\(\psi_n=(I-P_0)J_nu_n\).  If, for some \(c>0\),
\(C_n\ge0\), and \(\varepsilon_n\ge0\),
\[
 \lVert\psi_n\rVert\ge c\lVert u_n\rVert,\qquad
 \mathfrak h[\psi_n]\le C_n\mathfrak q_n[u_n]
 +\varepsilon_n\lVert u_n\rVert^2,
\]
\[
 C_n\frac{\mathfrak q_n[u_n]}{\lVert u_n\rVert^2}
 +\varepsilon_n\longrightarrow0,
\]
then the self-adjoint restriction
\[
 H_{\mathrm{exc}}
 =H\big|_{D(H)\cap\mathcal H_{\mathrm{exc}}}
\]
satisfies
\[
 \inf\sigma(H_{\mathrm{exc}})=0,\qquad
 0\in\sigma_{\mathrm{ess}}(H_{\mathrm{exc}}).
\]
The full spectral projection \(P_0\), rather than orthogonality to one
chosen vacuum, is essential unless vacuum uniqueness is independently
proved; see \cref{cusp-interpretation-typed-physical-transfer}.

\item \Open\ The hypotheses needed to apply that bridge to
four-dimensional quantum Yang--Mills theory have not been constructed here.
In particular, the workbench has no common gauge-invariant positive-metric
physical Hilbert space \(\mathcal H_{\mathrm{phys}}\), no reconstructed quantum
Hamiltonian \(H\), and no maps \(J_n\) satisfying the displayed norm and
energy estimates with control through gauge reduction, reflection-positive
reconstruction, renormalization, and the continuum and infinite-volume
 limits.  The exact radial scalar, continuum Hessian, compact-link flux,
 finite cochain, Gaussian-oscillator, magnetic, and split-Cartan
 \(SU(r)\) Jacobi calculations are
\Verified\ auxiliary geometric, finite-dimensional, or classical statements.
They neither prove nor disprove a mass gap for a
physical four-dimensional quantum Yang--Mills Hamiltonian; see
\cref{cusp-gauge-missing-physical-bridge,%
cusp-interpretation-present-boundary}.

\end{enumerate}

\paragraph{External validation note.}
A further context-isolated AI derivation independently recovered the cyclic
invariant lattices (11.4), determinant divisor (11.6), finite residue
cancellation (11.17), scalar test (11.19)--(11.20), all three absolute
cotangent direct images, and \(h^{1,2}=1\), \(h^{1,1}=2\).
Independent human specialist replication has not been recorded.  This is an
external-validation boundary, not a mathematical premise or an unresolved
step of the Hodge theorem; see \cref{record:h12-direct-images}.

\paragraph{Independent review checkpoint.}
Reviewer~B independently examined the reconstruction existing up to that
report's evidentiary cutoff and agreed that the intrinsic nearby-cycle,
Leray, and based van Kampen route does not depend on the then-open collapse
route.
Its disposition was provisionally favourable with focused verification
requests.  The later Codex proofs recorded here are separate work: they are
not retroactively attributed to Reviewer~B, and its report remains useful as
an independent comparison against the argument it actually saw.  Reviewer~B's
26 August analytic follow-up also independently proves the uniform
top-degree base-change corollary and the conductor--duality nonvanishing from
CAND-X, CAND-W, and TOP-H2 \cite{ReviewerBValidation2026}; those inputs are
not certified by that follow-up itself.

\paragraph{Status language.}
\Verified\ means proved here or reproduced by an explicit certificate.
\Conditional\ means a conclusion depending on an explicitly named
established or assumed hypothesis not proved in the present argument; it
never means doubtful or politely false.  A completely proved implication
whose antecedents are part of its exact statement is \Verified.  The current
candidate-specific ledger has no \Conditional\ row.  \Open\ means that no
complete proof or counterexample is recorded.  \Refuted\ means an exact
assertion has a complete contradiction or counterexample.  The status applies
to the immediately following assertion, not to an entire paper by contagion.

\section{Mathematical results so far}
\label{sec:mathematical-summary}

\noindent\textbf{6 September: explicit spectral and umbral maps, and the original family's parameter identifications.}
For the literal common lattice $\mathcal I=\{\sqrt2(a,a,a):a\in D_4\}$,
write $a=a_0+v$, $t^2=|v|^2$, $Q=a_0^2+t^2$ and
$P=a_0(a_0^2-3t^2)$. Its complete cubic image is
\[
 \tau(\mathcal I)=2\sqrt2\bigl(18\mathbb Z\cup(6\mathbb Z+2)
                                      \cup(6\mathbb Z-2)\bigr).
\]
Three explicit integral affine lines give a section of this value map.
The scalar-retaining Albert matrix $X=\mu I+Y(a)$, with cyclic
off-diagonal entries $\sqrt2a$, has invariants and labelled spectrum
\[
\begin{aligned}
 (\operatorname{tr}X,\operatorname{tr}X^2,\det X)
   &=(3\mu,3\mu^2+12Q,\mu^3-6\mu Q+4\sqrt2P),\\
 \operatorname{Spec}Y
   &=(2\sqrt2a_0,-\sqrt2a_0-\sqrt6t,-\sqrt2a_0+\sqrt6t).
\end{aligned}
\]
For $t>0$, the explicit labelled eigenprojector satisfies
$(\Pi_1)_{12}=-1/6-\sqrt3(v/t)/6$; thus the frame and spectrum
recover $a$, not only its norm and cubic. The full discriminant is
$864t^2(3a_0^2-t^2)^2$. The full Cartan family retains the scalar
$\mu$, radius $r$, and angle $\theta$, with $r^2=12Q$ and
$r^3\cos3\theta=24\sqrt3P$ on this subfamily.

For the same marked root system $A_5^4D_4$, the
Cheng--Duncan--Harvey prescription gives
$S^X=(5I_{12}+J)S_6$, where $J_{rs}=\mathbf1_{s=5r\ (12)}$ and
$S_{6,r}=\sum_{k\equiv r\ (12)}kq^{k^2/24}$.
The coefficient map to the retained Fricke function is
\[
 T_6=\frac{\eta(\tau)^5\eta(3\tau)}{\eta(2\tau)\eta(6\tau)^5},
 \qquad f^X=-q\frac{d}{dq}\log T_6,
 \qquad T_6\!\left(-\frac1{6\tau}\right)=\frac{72}{T_6(\tau)}.
\]
All twelve theta indices, the multiplier determinant, the original
glue-group action and its twisted shadows are computed in the proofs.
This is an explicit use and comparison of established umbral machinery,
not a claim to a new moonshine module.

In the original analytic family $\beta_c=\beta^{\rm part}+c$,
$\operatorname{Im}c<-M$, the actual global identifications are exactly
$X_c\simeq X_{c'}\Longleftrightarrow c'-c\in2\mathbb Z$.
The middle lattices agree for every integer shift; the multiplicity-four
filling retains an order-two obstruction to odd shifts. Even shifts have
explicit finite-filling and toric cusp maps and a proved global additive
gluing correction. The resulting parameter quotient is
$w=e^{-\pi ic}$, $0<|w|<e^{-\pi M}$.
This classifies the retained analytic family without using the claimed
identification of its total space with $S^6$.


\paragraph{6 September: cubic arithmetic, integral lattice maps, and the canonical ring.}
For the retained $A_5^4\perp D_4$ Niemeier lattice $N$ and its
original index-six Leech neighbour $\Lambda$, fix the Cayley--Dickson
frame $(1,i,j,k,\ell,i\ell,j\ell,k\ell)$ used in the complete cubic
proof. Put $\tau(q)=\operatorname{Re}((q_1q_2)q_3)$ and
$F=\tau\circ\mathcal L$ on the original Niemeier span,
and $\mathcal V_F(L)=\operatorname{span}_{\mathbb Z}\{F(y):y\in L\}$.
The exact additive groups are
\[
\begin{aligned}
 \mathcal V_F(N)&=\tfrac1{27}\mathbb Z\oplus
       \tfrac{\sqrt2}{54}\mathbb Z\oplus\tfrac{\sqrt3}{18}\mathbb Z,\\
 \mathcal V_F(\Lambda)&=\tfrac1{27}\mathbb Z\oplus
       \tfrac{\sqrt2}{216}\mathbb Z\oplus\tfrac{\sqrt3}{54}\mathbb Z,
 \qquad \mathcal V_F(\Lambda)/\mathcal V_F(N)\cong\mathbb Z/12.
\end{aligned}
\]
The Albert determinant is $\lambda^3-\lambda(y,y)+2F(y)$.
The quotient concerns the displayed value-generated groups in this fixed
frame, not an inclusion $N\subset\Lambda$. The actual residue
$r=\rho(361)$ cycles three Leech neighbours $\Lambda_j=r^j\Lambda$.
Their pairwise intersection is one integral lattice $J$, with
$\det J=81$ and $[\Lambda_j:J]=9$. The fixed fourth isotropic plane
in $J^\vee/J$ gives $M$ with exactly $72$ roots of type $A_2^{12}$,
$M\cap\Lambda_j=J$, and $[M:J]=9$.
\[
 \mathcal V_F(J)=\mathcal V_F(M)
 =\tfrac1{27}\mathbb Z\oplus\tfrac{\sqrt2}{216}\mathbb Z
                         \oplus\tfrac{\sqrt3}{18}\mathbb Z.
\]
The successive value-group inclusions have matrices
$\operatorname{diag}(1,4,1)$ and $\operatorname{diag}(1,1,3)$,
giving quotients $\mathbb Z/4$ and $\mathbb Z/3$.
Both containments and their reverse generator certificates are calculated
from the full cubic polarization, not inferred from lattice indices.

Wilson's construction \cite{esn:Wilson2009}, transported by the explicit
isometry $\iota(x)=x/\sqrt2$, gives $\Lambda_{\rm cyc}$ in the
original triality metric. Its cyclic fixed Gram matrix is $3G_{E_8}$;
its orthogonal complement is explicitly isometric to $A_2\otimes E_8$,
and their sum has index $3^8$, with quotient $\mathbb F_3^8$.
The directions remain $c(q_1,q_2,q_3)=(q_3,q_1,q_2)$ and
$c_{\rm W}=c^{-1}$. Its literal intersections are
\[
 \Lambda_{\rm cyc}\cap\mathcal L(N)
 =\Lambda_{\rm cyc}\cap\mathcal L(\Lambda)
 =\Delta:=\{\sqrt2(a,a,a):a\in D_{\mathbb H}\},
\]
where $D_{\mathbb H}$ is the even-coordinate-sum integer lattice in
the first four stated Cayley coordinates. Thus $\Delta$ has rank $4$,
Gram matrix $6G_{D_4}$, determinant $5184$, minimum squared norm $12$,
and value-generated subgroup $4\sqrt2\mathbb Z$ for $\tau|_\Delta$, or
$8\sqrt2\mathbb Z$ for $2\tau|_\Delta$. This is not the actual image:
$12\sqrt2$ is not attained by $\tau$ on $\Delta$.
The preceding proofs retain the triality products and spectral boundary
graph, the flag's degree-$+1$ collapse, and the retained cusp class's
nontrivial three-plane reduction; these exact maps do not supply an
interacting quantum mass-gap result.

For the original analytic fibration $f:X\to\mathbb P^1$, the full
canonical-divisor proof gives
$\omega_X\simeq\mathcal O_X(-2S_2)$ and the graded anticanonical ring
$\bigoplus_{m\ge0}H^0(X,\omega_X^{-m})\cong\mathbb C[U,V]$,
with $\deg U=1$, $\deg V=2$. The complete system $|-2K_X|$
recovers $f$ with the specified projective coordinates. This calculation
uses the original analytic charts and finite fillings, not an
identification of $X$ with $S^6$, and draws no sphere-topology conclusion.

The retained constant $\beta_c=\beta^{\rm part}+c$ now gives a proper
holomorphic deformation over $\operatorname{Im}c<-M$, locally on
relatively compact parameter discs. On a fixed regular fibre set
$T=\operatorname{Im}\tau$, $m=\operatorname{Im}\mu$,
$D=\operatorname{Im}\beta_{c_*}-6m^2/T<0$ and $\delta=c-c_*$.
In the plus-graph convention its exact Beltrami matrix is
\[
 \nu_\delta=\frac{\delta}{2iD+\delta}
                 \begin{pmatrix}0&0\\-m/T&1\end{pmatrix}.
\]
The connecting Kodaira--Spencer class is
$-[\partial_{\zeta_2}\otimes
(d\bar\zeta_2-(m/T)d\bar\zeta_1)/(2iD)]\ne0$.
The intrinsic pencil recovers $f$ also over the dual numbers; its three
fixed critical values prohibit an infinitesimal projective base change.
Consequently this deformation of the whole $X_c$ is nontrivial,
not only its marked regular-fibre deformation. The proof uses no sphere
recognition or CDP conclusion.



\paragraph{5 September: exact corrections to two S6 proof steps.}
The Cartier duality isomorphism in Theorem~\ref{thm:cartier-duality}
is induced by the analytic trace with $K_W^\bullet=\omega_W[2]$;
outer identities of five-term sequences alone do not define it.
For the holomorphic overlap translation by $\alpha(s)\Pi(s)v_j$,
the real marking gives
$a_j=\operatorname{Re}\alpha\,v_j+\operatorname{Im}\alpha\,J_sv_j$,
not multiplication of a real vector by the complex scalar $\alpha$.
The proved covariance $J_{g_js}A_j=A_jJ_s$ and $A_jv_j=v_j$
give $a_j(g_js)=A_ja_j(s)-v_j/m_j$. On the full clockwise
covering circle, $\Delta a_j=-v_j$ and the mixed class of
$h_j^*\Omega_E-\Omega_E$ is $-\iota_{v_j}E$.
Thus Lemma~\ref{lem:picard-finite-norm} still gives
$N_{T_j}d_j=\iota_{v_j}E$, and
Proposition~\ref{prop:geometric-pi1} retains its based meridian relation.
The complete proofs now include the base two-form term and its zero
restriction to the circle; no coefficient or orientation is suppressed.

\paragraph{5 September: exact period constants and cusp edge maps.}
For the cusp-regular particular solution define
\[
 D_{\rm part}=\operatorname{Im}\beta^{\rm part}
       -\frac{6(\operatorname{Im}\mu)^2}{\operatorname{Im}\tau},
 \qquad M=\max_{B\setminus\{p_0\}}D_{\rm part}.
\]
Proposition~\ref{prop:period-exact-admissible-constants} proves that
$M$ is finite and attained, and that $\beta=\beta^{\rm part}+c_0$
gives an everywhere invertible marked real period map exactly for
$\operatorname{Im}c_0<-M$. The cusp-free affine solution
$\widetilde\beta=\beta+i\,t\circ\pi$ has
$\widetilde D=D+\operatorname{Re}t\to+\infty$ along the positive
real base ray. The actual cusp extension excludes it.
For $C_v=E_0\cap E_v$, the cusp identification is
$\iota_v=\Psi_{B_0v}|_{C_v}:C_v\to C_{-v}$.
Using the same character $w=\chi^{(a,0)}$ with $av=0$ on both curves,
\[
 \iota_v^*w_{-v}=e^{2\pi i\,aC(0)B_0v}w_v.
\]
For $(v,a)=(e_1,(0,1)),(e_2,(1,0)),(e_1-e_2,(1,1))$, the multipliers
are respectively $e^{-2\pi i\mu_0}$, $e^{12\pi i\mu_0}$ and
$e^{-2\pi i(7\mu_0+b_0)}$, with $\mu_0=\mu(0)$ and $b_0=b(0)$.
Equation~\eqref{eq:cusp-exact-opposite-edge} retains the full $C(0)$;
Lemma~\ref{lem:toric-chart-control} uses restricted charts and the
closed-polydisc condition $\theta_i^\sigma(y)\ge0$ for torus points.
These refinements and the bounded finite-filling audit do not supply
an additional sphere-identification or CDP theorem.

\paragraph{5 September: the complete two-cross scalar--Gauss domain.}
Section~\ref{q444:sgd:section} retains the literal matrices, constants,
and all three full scalar and Gauss residuals. Put $\delta=\sqrt{2/3}$ and
\[
 \mathcal D=\{(u,v_0,y)\in\mathbb R^3:
 u>0,\ d=uv_0-\delta^2-y^2>0,\ c=4v_0-5u>0\}.
\]
Theorem~\ref{sgd:parametrization} proves a bijection from
$\mathcal D\times\{1,-1\}^2\times U(1)^2$ to all classified-frame
solutions, with $t=72d/c$, $e=\epsilon\sqrt{2+u+v_0}$,
$k=\sigma\sqrt{d/(-\det T)}$ and $\omega=e/k$, where
\[
 \begin{gathered}
 M=\begin{pmatrix}u&\delta+iy\\\delta-iy&v_0\end{pmatrix},\quad
 Q=\operatorname{diag}(1,2),\\
 T=\frac t4\operatorname{diag}(2/3,-5/6)
   -\frac12(QM+MQ).
 \end{gathered}
\]
The rank-one matrices $A_\pm=M\pm kT$ have $h_\pm=(A_\pm)_{11}>0$;
the exact vectors and phase inverse are
\[
 \begin{gathered}
 a_\pm=\begin{pmatrix}\sqrt{h_\pm}\\(A_\pm)_{21}/\sqrt{h_\pm}\end{pmatrix},\qquad
 p_\pm=\frac{(v\mp iw)_1}{\sqrt{h_\pm}},\\
 v=\frac{p_+a_++p_-a_-}{2},\quad
 w=\frac{i(p_+a_+-p_-a_-)}2.
 \end{gathered}
\]
Theorem~\ref{sgd:literal-addon:quotient} proves that the full literal
domain is the quotient of these parameters and $\zeta\in U(1)$ by
the sole free involution
$(p_+,p_-,\zeta)\mapsto(-p_+,-p_-,-\zeta)$.
The exact inverse frame recovers every literal temporal and spatial entry.
The smooth phase coordinates are
$(\eta,\alpha,\beta)=(\zeta^2,p_+/\zeta,p_-/\zeta)$.
Theorems~\ref{sgd:smooth-addon:literal-diffeomorphism}
and~\ref{sgd:smooth-addon:coupling} give
\[
 \mathcal S_{\rm sg}\cong\{1,-1\}^2\times(0,\infty)
                         \times\mathbb R^2\times\mathbb T^3,
 \qquad t=\tau=g^2,
\]
with coupling coordinates and inverse
\[
 v_0=\frac{\sqrt5}{2}\sqrt{\delta^2+y^2}+e^s,\qquad
 u=\frac{4\tau v_0+72(\delta^2+y^2)}{72v_0+5\tau},
 \quad (\tau,y,s)\in(0,\infty)\times\mathbb R^2.
\]
The literal manifold has four embedded six-dimensional components;
each positive-coupling level has four embedded five-dimensional
components $\mathbb R^2\times\mathbb T^3$
(Corollary~\ref{sgd:smooth-addon:levels}).
For each original $\tau=g^2>0$ the choice
\[
 \begin{gathered}
 \rho=\frac{\tau+\sqrt{\tau^2+4608\delta^2}}{96\delta}
       >\frac1{\sqrt2},\\
 M_\rho=\delta\begin{pmatrix}\rho&1\\1&2\rho\end{pmatrix},\qquad
 t=\frac{24\delta(2\rho^2-1)}\rho=\tau
 \end{gathered}
\]
proves nonemptiness with both temporal and frequency signs
(Corollary~\ref{sgd:every-coupling}). Both scalar equations and every
Gauss entry vanish, while Theorem~\ref{q444:twocross:thm:exclusion}
proves that the two spatial residuals cannot vanish simultaneously
on this entire class.
The quantitative spatial bound in
Proposition~\ref{q444:twocross:quantitative} therefore has an explicitly
nonempty domain. Existence on the unrestricted remaining exceptional
set stays \Open; no quantum gap claim follows.

\paragraph{5 September higher-rung and selector continuation.}
Section~\ref{sec:h24-continuation} gives the complete higher-wrap portion
of the expanded root draft.
The explicit octonionic slice has determinant
$\lambda^3-\lambda\sum_i|q_i|^2+2\operatorname{Re}((q_1q_2)q_3)$
and positive-cone compactification $(B^{25},S^{24})$.
The flag $F_4/\operatorname{Spin}(8)$ has Poincar\'e polynomial
$1+2t^8+2t^{16}+t^{24}$ and Euler class $6\beta_1^2\beta_2$.
For the retained peripheral character,
$(\chi(C_0),\chi(C_1),\chi(C'_2))=(-1,4,-3)$ in $\mathbb Z/12$.
The exact suspension is $E^6(a\alpha)=2a\nu_9$ in
$\pi_{12}(S^9)\cong\mathbb Z/24$; the particular cusp coefficient is
now $a=0$. Theorem~\ref{h24:thm:cusp-g-zero} constructs the complete
framed filling inside one original chart times an interval, retaining
the full period functions. Thus $g\simeq *$ and $\Theta_g\chi=0$.
The retained-angle map is different: Theorem~\ref{h24:ret:main}
proves $[p_3q_2]=\eta_4\eta_5\ne0$ in $\pi_6(S^4)$.
Its full normally four-framed inverse torus has spin form
$q(x,y)=x+xy+y$ over $\mathbb F_2$ and Arf invariant one,
so the original conical quotient $q_2$ is nonnull. The exact
homomorphism $[x]\mapsto[\eta_2\wedge x]$ sends this value to
$12\nu_6\ne0$ in $\pi_9(S^6)\cong\mathbb Z/24$; its
threefold suspension is $12\nu_9$ in the same group receiving
$E^6[g]=0$. Composing with the original $\chi$ gives peripheral
values $(12,0,12)$ and kernel $\chi^{-1}(2\mathbb Z/12)$.
The full class and exact order of $q_2$ are not computed by this
order-two quotient alone.
The arithmetic addition in
Section~\ref{sec:fifteen-selector-jordan} maps the complete fifteen
$p=1201$ selector rays into the real symmetric corner of
$\operatorname{Herm}_3(\mathbb H)$ and proves both linear span and
Jordan closure have dimension three. These additions retain the global
higher-completion question explicitly.

\paragraph{The modulo-$840$ lattice, full Weyl quotient, and retained counts.}
Section~\ref{esn:sec:shared} proves
\[
 R_{\mathrm{ES}}=A_5^4\perp D_4\subset N,\qquad
 [N:R_{\mathrm{ES}}]=72,\qquad \det N=1,\qquad
 \#\operatorname{Roots}(N)=144,
\]
\[
 O(N)=W(R_{\mathrm{ES}})\rtimes L(G),\qquad
 G\cong GL_2(\mathbb F_3)\cong Q_8\rtimes S_3.
\]
The full quotient and residue-position theorems
\ref{esn:fa:thm:full-quotient} and~\ref{esn:fa:thm:residue-weyl}
give $\rho(S_{840})\subset W(R_{\mathrm{ES}})$ and
$L_g\rho(r)=\rho(r)L_g$ in the original six-coordinate fibre order.
The abstract quotient is credited to \cite[Table~3]{esn:CDH2014};
the explicit code isomorphism, lattice lift and marked $D_4$ weight
comparison are proved here. The shell maps retain the actual prime:
$A_3(2521)=0$, $A_3(3361)=6$, with unordered counts $0,3$.
For a finite set $P$ of primes reducing into $S_{840}=\langle289\rangle$,
the count pushforward gives
$Q_{R,P}=6W_{R,P}-M_{R,P}u\in A(S_{840})\subset N$;
Theorem~\ref{esn:thm:sample-fourier} gives the exact Fourier inverse
for the residue-count function $w_{R,P}$.

This is the short route through the current mathematics.  It states what is
proved in this workbench and exhibits the load-bearing formulas.  Both the
intrinsic nearby-cycle/Leray certificate and the independent
collapse/Mayer--Vietoris certificate are now complete for the candidate.

\begin{enumerate}[leftmargin=1.8em,itemsep=.65em]
\item \Verified\ \textbf{The analytic threefold exists in the reconstructed
coordinate model.}
The descended modular periods give
\[
 \Pi=[Z\mid I_2],\qquad
 Z=\begin{pmatrix}6\mu&\tau\\\beta&\mu\end{pmatrix},
\]
with
\[
 \det_{\mathbb R}\Pi
 =-\det\operatorname{Im}Z
 =\operatorname{Im}\tau
 \left(\operatorname{Im}\beta-
 \frac{6(\operatorname{Im}\mu)^2}{\operatorname{Im}\tau}\right)<0.
\]
The Cayley coordinates
\(s_j(z)=(z-z_j)/(z-\bar z_j)\), the two elliptic orientations, and the
modular deck normalization are fixed simultaneously:
\[
 g_1'(z_1)=e^{-2\pi i/3},\qquad g_2'(z_2)=e^{-2\pi i/4}=-i,
 \qquad
 s_1\circ g_1=e^{-2\pi i/3}s_1,\quad s_2\circ g_2=-is_2.
\]
The period-deck homomorphism $\varphi$ is defined by
$\tau\circ g=\varphi(g)\circ\tau$.  Its cusp image is proved in a fixed
covering component, not by an untyped single-valuedness assertion: choose
\[
 U_0\subset\pi^{-1}(D_0^*),\qquad g_0U_0=U_0,\qquad
 \tau(U_0)\subset K\subset j^{-1}(\{|w|>R\}).
\]
For $P=\varphi(g_0)$ one has $P(K)=K$.  If $c$ is the boundary cusp of
$K$, choose $\gamma_0(\infty)=c$ and put
\[
 \tau'=\gamma_0^{-1}\tau|_{U_0},\qquad
 P'=\gamma_0^{-1}P\gamma_0,\qquad
 \tau'\circ g_0=P'\circ\tau'.
\]
If $P'=1$, then $\tau'$ descends to
$\bar\tau':D_0^*\to\mathfrak h$ and
$q'=\exp(2\pi i\bar\tau')$ simultaneously has a positive-order zero at the
cusp and a single-valued holomorphic logarithm, an impossible winding
combination.  Thus $P'$ and hence $P$ are nontrivial parabolic.

The auxiliary $\gamma_0$-frame is then discarded and the proof returns to
$\tau(z_1)=\rho$, $\varphi(g_1)=S_1$.  For the trace-$+2$ lift, with
primitive fixed vector $(p,r)^{\mathsf t}$, $\gcd(p,r)=1$, and $k\ne0$,
\[
 \widetilde P=I+k
 \begin{pmatrix}-pr&p^2\\-r^2&pr\end{pmatrix},
\]
and
\[
 \widetilde S_1=\begin{pmatrix}1&-1\\1&0\end{pmatrix},\qquad
 \operatorname{tr}(\widetilde S_1^{-1}\widetilde P^{-1})
 =1+k(p^2-pr+r^2)=0.
\]
Thus projective order two of $S_1^{-1}P^{-1}$ gives
$-k(p^2-pr+r^2)=1$, so $k=-1$ and
$p^2-pr+r^2=1$; conjugating the whole deck representation by a power of
$S_1$ leaves $S_1$ fixed and gives the final normalization below.  Rename
the resulting period map $\tau$ and define anew, in this cusp paragraph,
\[
 q=e^{2\pi i\tau}=t_cu_1(t_c),\qquad u_1(0)\ne0.
\]
For
\(S_1(w)=(w-1)/w\), \(S(w)=-1/w\), and \(P(w)=w-1\),
\[
 \varphi(g_1)=S_1,\qquad \varphi(g_2)=S,\qquad
 \varphi(g_0)=P,\qquad S_1SP=1
 \quad\text{in }\operatorname{PSL}_2(\mathbb Z).
\]
At $z_2$, comparison of quadratic leading terms gives
\[
 \bigl(g_2'(z_2)\bigr)^2=S'(i)=-1;
\]
the simultaneous parabolic normalization selects $g_2'(z_2)=-i$.
The periodic toric cusp action is free and properly discontinuous; its
bounded-chart argument uses
\(\lvert B_t\bar\lambda\rvert\ge\frac12\lvert\bar\lambda\rvert\).
Exponentiation identifies the punctured toric quotient with the family of
compact two-dimensional complex tori, not with an abelian-surface family.
For the finite logarithmic actions the opposite translation signs are
retained:
\[
 v_1=\varepsilon,\qquad v_2=-\varepsilon',\qquad
 \gamma(v_1)=1,\qquad\gamma(v_2)=-1,
\]
and
\[
 (z,x)\longmapsto(g_jz,A_jx+v_j/m_j).
\]
These actions are free by the invariant covector criterion.  On an overlap,
the compatible logarithm is
\[
 \ell'(w)=\ell(g_j^{-1}w)-\frac{2\pi i}{m_j},
 \qquad
 h_j\widetilde g_j^{\log}h_j^{-1}=\widetilde g_j^0,
\]
so no branch sign is suppressed.  Their radial stalk maps are the printed
quotient maps,
and the three fillings glue to a compact connected smooth complex
threefold \(X\) with a proper surjective holomorphic map
\[
 f:X\longrightarrow\mathbb P^1.
\]
On the cusp cover the toric form
\[
 \Omega=\frac{dx_1\wedge dx_2\wedge dt}{x_1x_2}
       =\pm dz_0\wedge dz_1\wedge dz_2
\]
extends nowhere vanishingly and satisfies
\(\Psi_{\bar\lambda}^*\Omega=\Omega\).  This first proves
\(\omega_{N_0}\simeq\mathcal O_{N_0}\); only then does the open embedding
\(N_0\hookrightarrow X\) give the typed two-stage chain
\[
 \omega_X|_{N_0}\simeq\omega_{N_0}\simeq\mathcal O_{N_0}.
\]
The complete chain is
\cref{prop:period-input-audit,thm:candidate-cusp-geometry,prop:candidate-cusp-overlap,prop:candidate-finite-fillings,lem:finite-filling-stalk-map,thm:candidate-global-geometry}.

\item \Verified\ \textbf{The cusp fibre has the exact geometry needed in the
topology and conductor arguments.}
Its scheme-theoretic local equations are
\[
 z_0,\qquad z_0z_1,\qquad z_0z_1z_2.
\]
It is reduced, irreducible, nonnormal, and normal crossings (not simple
normal crossings).  Its normalization is \(\mathrm{dP}_6\), with opposite
sides of the anticanonical hexagon identified, and its double locus consists
of three rational curves through two triple points.

\item \Verified\ \textbf{The algebraic reduction is computed directly and
does not use the topology or CDP20.}
For a very general smooth fibre,
\[
 \operatorname{NS}(F_b)=\mathbb Z\eta,\qquad
 \eta=u\wedge w+6\gamma\wedge\delta,\qquad
 \eta\wedge\eta=12\gamma\wedge u\wedge w\wedge\delta.
\]
The inverse-dual matrix \(G_\eta\) associated with the Appell--Humbert form
has determinant
\[
 \det G_\eta=
 \frac23\operatorname{Im}\tau
 \left(\operatorname{Im}\beta-
 \frac{6(\operatorname{Im}\mu)^2}{\operatorname{Im}\tau}\right)<0,
\]
and the same inertia as that form, hence signature \((1,1)\) and
\(a(F_b)=0\).  The repaired
meromorphic-descent argument then gives
\[
 \mathcal M(X)=f^*\mathbb C(t),\qquad a(X)=1,
\]
and \(f\) is the algebraic reduction; see
\cref{prop:very-general-torus-audit,thm:algebraic-dimension-audit}.

\item \Verified\ \textbf{The intrinsic nearby-cycle/Leray route proves the
integral topology without using the cellular collapse.}
For the cusp filling, the integral nearby-cycle map satisfies
\[
 \operatorname{sp}^r:H^r(W;\mathbb Z)
   \xrightarrow{\ \sim\ }H^r(F;\mathbb Z)^{T_0}
 \qquad(0\le r\le4)
\]
as saturated lattices.  With
$q:=uw+6\gamma\delta\in\bigwedge^2V=H^2(F;\mathbb Z)$, the global sheaf calculation gives
\[
H^0(R^1)=\mathbb Z(12\gamma),\quad
H^0(R^2)=\mathbb Z(2q),\quad
H^0(R^3)=\mathbb Z(2\gamma uw).
\]
The canonical truncation triangle
\[
 R^0f_*\mathbb Z_X\longrightarrow\tau_{\le1}Rf_*\mathbb Z_X
 \longrightarrow R^1f_*\mathbb Z_X[-1]
 \xrightarrow{\kappa}R^0f_*\mathbb Z_X[1]
\]
identifies the explicitly computed Cech discrepancy with Leray \(d_2\).
For \((\ell_0,\ell_1,\ell_2)=(0,1,-1)\), hence \(p=-1\), orientation of
\(\omega\) may be chosen so that
\[
d_2(12\gamma)=\omega,\qquad
d_2(2q)=\omega[\delta],\qquad
d_2(2\gamma uw)=\omega[\gamma\delta].
\]
Therefore
\[
 H^1(X;\mathbb Z)=H^2(X;\mathbb Z)=H^3(X;\mathbb Z)=0,
 \qquad H_2(X;\mathbb Z)=0.
\]

\item \Verified\ \textbf{The based van Kampen maps and sphere recognition
close the topology theorem.}
Zero-section whiskers fix the common basepoint, the finite logarithmic
overlaps impose the primitive relations with one common sign, the cusp
inclusion is \(\Lambda\to\bar\Lambda\), and the cusp meridian bounds the
explicit toric zero-section disc with no translation drift.  Hence
\[
\pi_1(X)=
\langle c,x,y\mid c\text{ central},\ xy=1,\ x^3=c,\ y^4=c^{-1}\rangle
\simeq1.
\]
Poincar\'e duality, universal coefficients, Hurewicz, and the smooth
homotopy-sphere recognition theorem now give
\[
 X\cong_{\mathrm{diff}}S^6,\qquad
 \langle c_3(T_X),[X]\rangle=\chi_{\mathrm{top}}(X)=2.
\]
See
\cref{thm:p0-integral-specialization,prop:leray-e2,lem:first-transgression,lem:higher-transgressions,thm:h2-vanishing,prop:geometric-pi1,cor:candidate-sphere}.

\item \Refuted\ \textbf{The printed CDP dependency has two exact failures.}
For a reduced nonnormal fibre,
\[
0\longrightarrow N^*_{S/X}\longrightarrow
\ker(\Omega_X^1|_S\to\widetilde\Omega_S^1)
\longrightarrow\operatorname{tors}\Omega_S^1\longrightarrow0,
\]
so the vanishings used on p.~692 of \cite{CDP2020} do not control the
conductor-supported middle term.  A proper
nodal-cubic-times-\(\mathbb P^1\) family is a counterexample to that
inference.  Separately, a split extension satisfies
\[
 (K\rtimes Q)^{\mathrm{ab}}
 \simeq Q^{\mathrm{ab}}\oplus(K^{\mathrm{ab}})_Q,
\]
not the direct sum with all of \(K^{\mathrm{ab}}\) unless monodromy acts
trivially.

\item \Verified\ \textbf{The singular-fibre nonvanishing has a
topology-free $\operatorname{Pic}^0$ proof and three exact all-Pic
certificates.}
The invariant toric three-form first trivializes $\omega_{N_0}$, and the
open embedding then gives
$\omega_X|_{N_0}\simeq\omega_{N_0}\simeq\mathcal O_{N_0}$, as recorded
above.  Let
$\nu:\widetilde W\simeq\mathrm{dP}_6\to W$ be the normalization.  For
$L\in\operatorname{Pic}^0(X)$, naturality of $c_1$ and injectivity of
\[
 c_1:\operatorname{Pic}(\mathrm{dP}_6)\hookrightarrow
 H^2(\mathrm{dP}_6;\mathbb Z)
\]
give
$\nu^*((L^\vee\otimes\omega_X)|_W)\simeq\mathcal O_{\mathrm{dP}_6}$.
The conductor section, Ramis--Ruget duality on the compact Cartier fibre,
and top-degree analytic base change then prove
\[
 R^2f_*(T_X\otimes L)\ne0
 \qquad(L\in\operatorname{Pic}^0(X))
\]
without any global homology calculation.  For arbitrary $L$, the exact
trivial-pullback obstruction is the homomorphism
\[
 r:\operatorname{Pic}(X)\longrightarrow\operatorname{Pic}(\mathrm{dP}_6),
 \qquad r([L])=[\nu^*i^*L].
\]
For the candidate, the Leray edge image of $c_1(L)$ lies in
$E_\infty^{0,2}\subseteq\ker d_2^{0,2}$, but
$d_2^{0,2}(2q)=\omega[\delta]$ is an isomorphism.  Hence the edge section is
zero; proper base change and naturality give $c_1(L|_W)=0$, and injectivity
of $c_1$ on $\mathrm{dP}_6$ forces
\[
 \nu^*i^*L\simeq\mathcal O_{\mathrm{dP}_6},
 \qquad r([L])=0.
\]
This proves the conclusion for
every $L\in\operatorname{Pic}(X)$ without citing the packaged
$H^2(X;\mathbb Z)=0$ theorem, although it reuses the same integral Leray
transgression.

The second certificate uses no global Leray differential.  The exact
order-four filling lattice forces $c_1(L|_F)=2nq$.  Subtracting the invariant
vertical form on the pair-of-pants bundle gives a crossed homomorphism with
global-coinvariant boundary values.  The logarithmic translations and their
cyclic norm maps give
\[
 \bar d_1=4n[\delta],\qquad \bar d_2=-3n[\delta],
\]
while the flat zero-section cusp sweep bounds a literal toric solid torus,
so $\bar d_0=0$.  The based relation gives
$\bar d_1+\bar d_2+\bar d_0=0$, hence $n=0$.  Naturality and injectivity of
the local integral cusp specialization then give $c_1(L|_W)=0$ and
\[
 \nu^*i^*L\simeq\mathcal O_{\mathrm{dP}_6},
 \qquad r([L])=0.
\]
This is independent of global $d_2$, though not of local cusp
topology.  See
\cref{prop:pic0-pullback,prop:all-pic-restriction,%
prop:leray-edge-restriction,prop:direct-boundary-picard,%
cor:candidate-pic0-route,cor:candidate-conductor-actual}.

The precise limit of a cusp-only analytic replacement is also computed.
For $Q(v)=v_1^2-6v_2^2$, the locally finite toric divisor
$\mathcal D_Q=\sum_vQ(v)E_v$ has an explicit equivariant descent
$\mathcal L_Q\in\operatorname{Pic}(N_0)$ satisfying
\[
 c_1(\mathcal L_Q|_{F_t})=-2q,\qquad
 \bigl(\deg_{C_i}\nu^*(\mathcal L_Q|_W)\bigr)
 =(-12,2,0,-12,2,0).
\]
Thus $\operatorname{Pic}(N_0)\to\operatorname{Pic}(\mathrm{dP}_6)$ is
nonzero, so no cusp-local representability argument can prove the global
restriction theorem by itself.  The completed third certificate adds the two
finite fillings analytically and retains the prior coinvariant input
\[
 (\textstyle\bigwedge^2V)^G=\mathbb Zq,\qquad
 q=u\wedge w+6\gamma\wedge\delta,\qquad
 V_G\otimes\mathbb R=\mathbb R[\delta].
\]
Thus, for a global line bundle $L$,
$E=c_1(L|_F)=kq$.  The finite norm proof keeps every covering level
distinct.  On the vector cover
\[
 \widetilde J_j=\Delta_j\times\mathbb C^2
 \xrightarrow{\ r_j\ }J'_j
\]
define the vector-cover automorphisms
\[
 Q_{j,\log},Q_{j,\mathrm{lin}}\in\operatorname{Aut}(\widetilde J_j),
 \qquad
 Q_{j,\log}=\widetilde g_j^{\,v_j/m_j},\quad
 Q_{j,\mathrm{lin}}=\widetilde g_j^{\,0}.
\]
These are not the quotient automorphisms: on the vector cover,
\[
\begin{aligned}
 Q_{j,\log}\tau_\lambda Q_{j,\log}^{-1}
 &=\tau_{A_j\lambda},&
 (Q_{j,\log})^{m_j}&=\tau_{v_j},\\
 Q_{j,\mathrm{lin}}\tau_\lambda Q_{j,\mathrm{lin}}^{-1}
 &=\tau_{A_j\lambda},&
 (Q_{j,\mathrm{lin}})^{m_j}&=1.
\end{aligned}
\]
The sign in $(Q_{j,\log})^{m_j}=\tau_{v_j}$ is positive because
$m_j$ clockwise iterations of
$s_j\circ g_j=e^{-2\pi i/m_j}s_j$ produce translation by $+v_j$.
Let
\[
 q_{j,\log},q_{j,\mathrm{lin}}\in\operatorname{Aut}(J'_j)
\]
be their descents.  They have order dividing $m_j$ by these relations and
exactly $m_j$ because they cover $g_j$.  On
$J_j^{\prime *}=J'_j|_{\Delta_j^*}$ define
\[
\begin{aligned}
 p_{j,\log}&:J_j^{\prime *}\longrightarrow
 N_j^*=J_j^{\prime *}/\langle q_{j,\log}\rangle,\\
 p_{j,\mathrm{lin}}&:J_j^{\prime *}\longrightarrow
 J_j^*=J_j^{\prime *}/\langle q_{j,\mathrm{lin}}\rangle .
\end{aligned}
\]
The branch-independent translation and its descended biholomorphism are
typed by
\[
 h_j:J_j^{\prime *}\longrightarrow J_j^{\prime *},\qquad
 \bar h_j:N_j^*\longrightarrow J_j^*,
\]
and satisfy
\[
 h_j\circ q_{j,\log}\circ h_j^{-1}=q_{j,\mathrm{lin}},\qquad
 p_{j,\mathrm{lin}}\circ h_j=\bar h_j\circ p_{j,\log},
\]
and the filled line bundle gives the typed pullback comparison
\[
 p_{j,\log}^*(L|_{N_j^*})
 \simeq h_j^*p_{j,\mathrm{lin}}^*(L|_{J_j^*}).
\]

For the factor extension return to the filled vector cover:
\[
 P_{j,\log}:J'_j\longrightarrow N_j,\qquad
 \pi_j^{\rm vec}:=P_{j,\log}\circ r_j:
 \widetilde J_j\longrightarrow N_j,\qquad
 P_{j,\log}|_{J_j^{\prime *}}=p_{j,\log},
\]
\[
 \Gamma_j^{\log}
 =\langle\tau_\lambda,Q_{j,\log}:\lambda\in\Lambda\rangle.
\]
The pullback $(\pi_j^{\rm vec})^*(L|_{N_j})$ defines the filled factor
extension.  Over $\Delta_j^*$ the quotient square identifies its
commutator form with $E$; that integer-valued form is locally constant on
the connected disc, so it remains $E$ on the filled cover.  After
trivializing this pullback, choose normalized vertical lifts
\[
 U_\lambda^{(j)}U_\mu^{(j)}
 =\mathbf c\!\left(\frac12E(\lambda,\mu)\right)
 U_{\lambda+\mu}^{(j)},
\]
write
$\alpha\tau_\lambda\alpha^{-1}=\tau_{A_\alpha\lambda}$ and
$T_\alpha=(A_\alpha^{-1})^{\mathsf t}$, and define
$\widetilde d_j(\alpha)\in V_{\mathbb R}$ by
\[
 \widetilde G_\alpha U_{A_\alpha^{-1}\lambda}^{(j)}
 \widetilde G_\alpha^{-1}
 =\mathbf c\!\left(
 \langle\widetilde d_j(\alpha),\lambda\rangle\right)
 U_\lambda^{(j)}.
\]
It is a crossed homomorphism:
\[
 \widetilde d_j(\alpha\beta)
 =\widetilde d_j(\alpha)+T_\alpha\widetilde d_j(\beta).
\]
To compare the boundary factors, use the actual common cover and its
restricted lattice quotient
\[
 \widetilde J_j^*:=\Delta_j^*\times\mathbb C^2,
 \qquad
 r_j^*:=r_j|_{\widetilde J_j^*}:
 \widetilde J_j^*\longrightarrow J_j^{\prime *}.
\]
The two quotient maps from this common domain are
\[
\begin{aligned}
 \Phi_{j,\log}:\widetilde J_j^*&\longrightarrow N_j^*,&
 \Phi_{j,\log}&:=p_{j,\log}\circ r_j^*,\\
 \Phi_{j,\mathrm{glob}}:\widetilde J_j^*&\longrightarrow J_j^*,&
 \Phi_{j,\mathrm{glob}}&:=p_{j,\mathrm{lin}}\circ h_j\circ r_j^*
 =\bar h_j\circ\Phi_{j,\log}.
\end{aligned}
\]
Both maps are $\Gamma_j^{\log}$-invariant.  Translation invariance follows
from $r_j^*\circ\tau_\lambda=r_j^*$.  For every
$\widetilde x\in\widetilde J_j^*$, invariance under the nontrivial generator
is the point-valued calculation
\[
\begin{aligned}
 \Phi_{j,\log}\bigl(Q_{j,\log}(\widetilde x)\bigr)
 &=p_{j,\log}\bigl(q_{j,\log}(r_j^*(\widetilde x))\bigr)
 =\Phi_{j,\log}(\widetilde x),\\
 \Phi_{j,\mathrm{glob}}\bigl(Q_{j,\log}(\widetilde x)\bigr)
 &=p_{j,\mathrm{lin}}\bigl(
 h_j(q_{j,\log}(r_j^*(\widetilde x)))\bigr)\\
 &=p_{j,\mathrm{lin}}\bigl(
 q_{j,\mathrm{lin}}(h_j(r_j^*(\widetilde x)))\bigr)
 =\Phi_{j,\mathrm{glob}}(\widetilde x).
\end{aligned}
\]
Thus $\widetilde J_j^*$ is the same $\Gamma_j^{\log}$-cover of the two
biholomorphic quotients $N_j^*$ and $J_j^*$.  Pulling back the gluing
isomorphism gives
\[
 \Theta_j:
 \Phi_{j,\log}^*(L|_{N_j^*})
 \xrightarrow{\ \sim\ }
 \Phi_{j,\mathrm{glob}}^*(L|_{J_j^*}).
\]
It intertwines descent and hence identifies the pushed-out filled and global
$\Gamma_j^{\log}$-factor extensions while fixing every central
$\mathbf c(r)$, $r\in\mathbb R$.

The loop and its three successive lifts are fixed pointwise as follows.  Let
\[
 \gamma_j^J:[0,1]\longrightarrow J_j^*
\]
be the zero-section image of the unwhiskered clockwise boundary meridian,
so $\gamma_j^J(1)=\gamma_j^J(0)$.  Choose the global zero-section basepoint
$\widehat b\in J|_P$ and the fixed zero-section whisker
\[
 \omega_j:[0,1]\longrightarrow J|_P,\qquad
 \omega_j(0)=\widehat b,\qquad
 \omega_j(1)=\gamma_j^J(0),
\]
where the boundary copy of $J_j^*$ is identified with its image in $J|_P$.
The retained based total-space loop is the class
\[
 \widehat\rho_j
 =[\omega_j*\gamma_j^J*\omega_j^{-1}]
 \in\pi_1(J|_P,\widehat b).
\]
Choose $x_j\in J_j^{\prime *}$ such that
$p_{j,\mathrm{lin}}(h_j(x_j))=\gamma_j^J(0)$, put
$x_{j,\mathrm{lin}}=h_j(x_j)$, and let
$\widetilde\gamma_{j,\mathrm{lin}}:[0,1]\to J_j^{\prime *}$ be the
$p_{j,\mathrm{lin}}$-lift of $\gamma_j^J$ beginning at
$x_{j,\mathrm{lin}}$.  For the chosen clockwise generator,
\[
\begin{aligned}
 p_{j,\mathrm{lin}}\circ\widetilde\gamma_{j,\mathrm{lin}}
 &=\gamma_j^J,\\
 \widetilde\gamma_{j,\mathrm{lin}}(0)&=x_{j,\mathrm{lin}},\\
 \widetilde\gamma_{j,\mathrm{lin}}(1)
 &=q_{j,\mathrm{lin}}(x_{j,\mathrm{lin}}).
\end{aligned}
\]
Set
\[
 \gamma_j^N:=\bar h_j^{-1}\circ\gamma_j^J:[0,1]\longrightarrow N_j^*,
 \qquad
 \widetilde\gamma_{j,\log}:=h_j^{-1}\circ
 \widetilde\gamma_{j,\mathrm{lin}}:[0,1]\longrightarrow J_j^{\prime *}.
\]
Then the quotient square gives the point-valued identities
\[
\begin{aligned}
 p_{j,\log}\circ\widetilde\gamma_{j,\log}&=\gamma_j^N,\\
 \widetilde\gamma_{j,\log}(0)&=x_j,\\
 \widetilde\gamma_{j,\log}(1)&=q_{j,\log}(x_j).
\end{aligned}
\]
Choose $\widetilde x_j\in\widetilde J_j^*$ with
$r_j^*(\widetilde x_j)=x_j$, and let
$\widetilde\gamma_{j,\mathrm{vec}}:[0,1]\to\widetilde J_j^*$ be the
$r_j^*$-lift of $\widetilde\gamma_{j,\log}$ beginning at
$\widetilde x_j$.  On adjacent sheets choose compatible logarithm branches
\[
 \ell'(w)=\ell(g_j^{-1}w)-\frac{2\pi i}{m_j}.
\]
The point-valued vector-cover endpoint is
\[
\begin{aligned}
 r_j^*\circ\widetilde\gamma_{j,\mathrm{vec}}
 &=\widetilde\gamma_{j,\log},\\
 \widetilde\gamma_{j,\mathrm{vec}}(0)&=\widetilde x_j,\\
 \widetilde\gamma_{j,\mathrm{vec}}(1)
 &=Q_{j,\log}(\widetilde x_j).
\end{aligned}
\]
excluding both $Q_{j,\log}^{-1}(\widetilde x_j)$ and every
$\tau_\lambda Q_{j,\log}(\widetilde x_j)$ with $\lambda\ne0$.  Fibre loops
remain the pointwise deck paths from $\widetilde x_j$ to
$\tau_\lambda(\widetilde x_j)$, since $h_j$ is a fibre translation.

The factor of $L$ along the fixed lift of $\omega_j$ identifies the target
factor over the local boundary segment with the globally based factor.  Let
$G_{\rho_j}^{(j)}$ be the target lift of the local boundary deck element
whose whiskered image is the chosen global lift $G_{\rho_j}$.  It satisfies
\[
 G_{\rho_j}^{(j)}U_{A_j^{-1}\lambda}
 (G_{\rho_j}^{(j)})^{-1}
 =\mathbf c(\langle d_j,\lambda\rangle)U_\lambda,
 \qquad d_j=d(\rho_j).
\]
Consequently, for some $a_j\in V_{\mathbb R}$ and $b_j\in\mathbb R$,
\[
 \Theta_{j,*}(U_\lambda^{(j)})
 =\mathbf c(\langle a_j,\lambda\rangle)U_\lambda,\qquad
 \Theta_{j,*}(\widetilde G_{Q_{j,\log}})
 =\mathbf c(b_j)G_{\rho_j}^{(j)}.
\]
The central factor $\mathbf c(b_j)$ disappears under conjugation.  Therefore
\[
 \widetilde d_j(Q_{j,\log})
 =d_j+(T_j-I)a_j,\qquad
 \widetilde d_j(\tau_v)=\iota_vE.
\]
Here
$A_{Q_{j,\log}}=A_j$ and $T_{Q_{j,\log}}=T_j$.  Consequently the exact,
sign-sensitive norm chain is
\[
\begin{aligned}
 N_{T_j}d_j
 &=N_{T_j}\widetilde d_j(Q_{j,\log})
 =\widetilde d_j((Q_{j,\log})^{m_j})\\
 &=\widetilde d_j(\tau_{v_j})
 =\iota_{v_j}E.
\end{aligned}
\]
Finally, on a simply connected $U\subset\Delta_j^*$ the logarithm branch
$\ell$ gives the point-typed vector-cover lift
\[
 H_j^\ell:U\times\mathbb C^2\longrightarrow U\times\mathbb C^2,
 \qquad
 H_j^\ell(z,\zeta)
 =\left(z,\zeta+\frac{\ell(z)}{2\pi i}\Pi(z)v_j\right)
\]
of $h_j$.  A branch change $\ell\mapsto\ell+2\pi i n$ gives
\[
 H_j^{\ell+2\pi i n}=\tau_{nv_j}\circ H_j^\ell.
\]
At factor-extension level this gives
\[
 \Theta'_{j,*}
 =\operatorname{Ad}(U_{nv_j})\circ\Theta_{j,*}.
\]
Thus the comparison character changes by
$a_j\mapsto a_j+n\iota_{v_j}E$.  Since
\[
 T_j(\iota_{v_j}E)=\iota_{v_j}E,
\]
both displayed $\widetilde d_j$ and norm identities are unchanged; the
descended $h_j,\bar h_j$ and quotient square are branch independent.
The exact contractions are
\[
 \iota_{v_1}q=2\eta_1,\quad \iota_{v_2}q=-3\eta_2,\quad
 [\eta_1]=3[\delta],\quad [\eta_2]=2[\delta],
\]
and logarithms of the actual $\mathcal O^*$ factors give
\[
 \bar d_1=2k[\delta],\qquad
 \bar d_2=-\tfrac32k[\delta],\qquad
 \bar d_0=0,\qquad
 \bar d_1+\bar d_2+\bar d_0=0,
\]
where the cusp equality is an argument-principle calculation in a Stein
chart.  Hence $k=0$ and $E=0$.

The toric support-function calculation uses the explicitly typed unimodular
map
\[
 B_0:\bar\Lambda\simeq\mathbb Z^2
 \longrightarrow\Lambda_{\rm tor}\simeq\mathbb Z^2,
 \qquad
 [B_0]=\begin{pmatrix}0&1\\-1&0\end{pmatrix},
 \qquad \det B_0=1.
\]
For adjacent rays $\rho_v,\rho_w<\sigma=\{v,w,u\}$ the affine-open
directions are, exactly,
\[
 U_{\rho_v}=D(z_wz_u)\subset U_\sigma,\qquad
 U_{\rho_w}=D(z_vz_u)\subset U_\sigma,\qquad
 U_{\rho_v}\cap U_{\rho_w}=D(z_vz_wz_u)=U_0.
\]
If
$\mu^+_\lambda:\Psi_\lambda^*\mathcal O(D_a)\to\mathcal O(D_a)$
is multiplication by $\chi^{m_\lambda}$, the map needed on the negative
divisor is the inverse-dual, not the ordinary dual:
\[
 \mu^-_\lambda:=((\mu^+_\lambda)^{-1})^\vee:
 \Psi_\lambda^*\mathcal O(-D_a)\longrightarrow\mathcal O(-D_a),
 \qquad s\longmapsto\chi^{-m_\lambda}s.
\]
Now set, with its type retained,
\[
 M=L|_{N_0}\in\operatorname{Pic}(N_0),\qquad
 E_M=c_1(M|_{F_t})=E=0.
\]
For
$A=\nu^*(M|_W)$ and
$(\deg_{C_1}A,\ldots,\deg_{C_6}A)=(d_1,d_2,d_3,d_1,d_2,d_3)$,
the normalized degrees determine
\[
 [B]=\begin{pmatrix}d_2+d_3&d_3\\d_3&d_1+d_3\end{pmatrix},
 \qquad
 E_M((\lambda,0),(0,m))=-B(B_0\lambda,m).
\]
The antecedent $E_M=0$ and unimodularity of $B_0$ force $B=0$, hence all
six degrees vanish and
\[
 \nu^*(L|_W)=\nu^*(M|_W)\simeq\mathcal O_{E_0};
 \qquad
 r([L])=[\mathcal O_{E_0}]=0
 \quad\text{in }\operatorname{Pic}(E_0)\simeq
 \operatorname{Pic}(\mathrm{dP}_6).
\]
This proof uses neither global Leray
nor local integral specialization.  See
\cref{prop:local-cusp-picard-obstruction,prop:analytic-ah-boundary,%
lem:analytic-support-specialization,thm:analytic-all-picard}.

\item \Verified\ \textbf{The constructed threefold is a theorem-level
counterexample to CDP20 Theorem~2.2 as stated.}
It has
\[
H^1(X;\mathbb Z)=H^2(X;\mathbb Z)=0,\qquad
a(X)=1,\qquad f:X\to\mathbb P^1
\text{ a holomorphic algebraic reduction},
\]
while
\[
c_3(X)=2>0,\qquad X\cong_{\mathrm{diff}}S^6.
\]
These are precisely the relevant hypotheses and the opposite of the
published conclusions \(c_3(X)\le0\) and \(a(X)=0\).  This conclusion uses
the verified candidate proof, not merely the fact that defects were found in
the published proof.

\item \Verified\ \textbf{The strongest corrected CDP conclusion is now
proved, with its exact extra hypothesis.}
For any compact complex threefold in the CDP holomorphic-reduction setting,
existence of a single \(L_0\) with
\[
 H^2(X,T_X\otimes L_0)=0
\]
implies, for general \(M\),
\[
 H^i(X,T_X\otimes M)=0\quad(i\ne1),\qquad
 \chi(X,T_X\otimes M)=\frac12c_3(X)\le0.
\]
The relative condition \(R^2f_*(T_X\otimes L)=0\) for one \(L\) implies
this global hypothesis by projection formula, Serre vanishing on the base,
and Leray.  On a reduced full fibre \(W_c\), after the two endpoint
vanishings, the missing ambient-cotangent vanishing is equivalent to
injectivity of
\[
 H^0(W_c,\operatorname{tors}\Omega^1_{W_c}\otimes A_c)
 \xrightarrow{\delta_c}
 H^1(W_c,N^*_{W_c/X}\otimes A_c).
\]
This repairs the proof route without asserting the false original theorem.
Because \(c_3(X)=2\), the contrapositive of the corrected global theorem gives
\(H^2(X,T_X\otimes L_0)\ne0\) for every \(L_0\).  Independently, the candidate
has nonzero \(\ker\delta_c\) at its cusp and
\(R^2f_*(T_X\otimes L)\ne0\) for every \(L\).  See
\cref{thm:corrected-cdp-global,cor:corrected-cdp-relative,%
thm:corrected-cdp-fibrewise}.

\item \Verified\ \textbf{The Proposition~10.8 uniform non-torsion branch
is complete.}
The reconstructed direct-image calculation gives
\[
R^1f_*\mathcal O_X\simeq\mathcal O\oplus\mathcal O(-1),
\qquad \operatorname{Pic}(X)\simeq\mathbb C,
\]
with degree-one base change on all scheme-theoretic fibres.  At a finite
filling, the invariant canonical section and the direct-image generator are
not conflated: if \(k_j=m_j-1-a_j\), then
\[
 \operatorname{div}(\overline\sigma_j)=k_jS_j,
 \qquad
 f_{j*}\omega_{N_j/D_j}
 =\mathcal O_{D_j}\cdot s_j^{k_j}\vartheta'_j,
 \qquad
 s_j^{k_j}\vartheta'_j
 =\frac{e^{\varphi(s_j)}}{m_j}s_j^{-a_j}e
 \quad\text{on }\Delta_j^*.
\]
In particular \(k_2=2\), so the canonical section is not a frame there;
the invariant-module identity is the exact statement that proves the Hodge
line used below.  The finite
character calculation gives
\(\operatorname{ord}(\mathcal O_X(S_j)|_{S_j})=m_j\).  For
\(\lambda=a\gamma+bu+cw+d\delta\), membership of its Dolbeault image in
the global Picard line is exactly
\[
c\tau+d\mu-b=0.
\]
Linear independence of \(1,\tau,\mu\) makes the rank-two exceptional
fibres countable, while rank-one fibres contribute bad parameters only in
\(\mathbb R\).  The three full special fibres and the reductions \(S_j\)
each add only countably many parameters.  Thus every
\(z\in\mathbb C\setminus(\mathbb R\cup\text{a countable set})\) supplies
one \(L_z\) non-torsion on every full fibre and on both \(S_j\).
The cusp step was first repaired using intrinsic nearby cycles and therefore
does not depend on the later cellular-collapse verification.

\item \Verified\ \textbf{The equivariant cellular collapse and its exact
comparison with nearby cycles are now complete.}
The positive toric quotient is a compact $3$-manifold cobordism between two
tori, its central boundary is a strong deformation retract, and the lifted
retraction spreads by compact phases and conjugates back to the original
twist.  On the fixed fibre, the resulting map $c=r|_F$ has
\[
0\longrightarrow\operatorname{im}(\textstyle\bigwedge^qM_0-I)
\longrightarrow\textstyle\bigwedge^q\Lambda
\xrightarrow{c_*}H_q(W;\mathbb Z)\longrightarrow0.
\]
In degree two, with the primitive orientations fixed in
\cref{prop:verified-collapse-homology},
\[
[c_*]_2=
\begin{pmatrix}
1&0&0&0&0&0\\
0&0&-1&0&0&0\\
0&-1&-1&1&1&0\\
0&0&0&1&0&0
\end{pmatrix},
\]
and its kernel is
$\mathbb Z\langle\widehat w\widehat\delta,
\widehat\gamma\widehat w+\widehat u\widehat\delta\rangle$.
The neighborhood realization of the nearby-cycle unit then gives, for all
$q$,
\[
\operatorname{sp}^q=i_F^*(i_W^*)^{-1}=i_F^*r^*=c^*.
\]
Thus C1--C3 are verified integrally, including saturation and the
degree-three orientation.

\item \Verified\ \textbf{The global Mayer--Vietoris realization and its
Leray comparison are complete.}
The earlier audit conflated the two stages.  On
$X^\circ=X_1\cup X_2$, the verified radial markings give the actual maps
\[
 \alpha_q=(\pi_{1*},-\pi_{2*}).
\]
Their degree-one and degree-two matrices have saturated rank three.  On the
second attachment $X=N'_0\cup X^\circ$, primitive splittings reduce the
global maps to
\[
 D_1=\begin{pmatrix}1&0&0\\0&1&0\\-12&0&-p\end{pmatrix},
 \qquad
 D_2^{\rm red}=\begin{pmatrix}-p&0\\0&1\end{pmatrix}.
\]
For $p=-1$ they are unimodular.  Dualizing the same oriented collar maps
gives commuting obstruction squares and
\[
 d_2(12\gamma)=-p\omega,\quad
 d_2(2q)=-p\omega[\delta],\quad
 d_2(2\gamma uw)=-p\omega[\gamma\delta].
\]
Thus the alternate route independently gives the integral homology of
$S^6$; see
\cref{prop:actual-finite-band-maps,prop:cusp-mv-blocks,%
thm:mv-leray-comparison,cor:candidate-mv-sphere}.

\item \Verified\ \textbf{The Hodge diamond and Fr\"olicher sequence are
computed completely.}
Put
\[
 \Sigma=p_0+p_1+p_2,\qquad
 D=f^{-1}(\Sigma)_{\rm red}=W+S_1+S_2,\qquad
 \mathcal L=\Omega_B^1(\log\Sigma)\simeq\mathcal O_B(1).
\]
The logarithmic relative bundle
$\mathcal E=\Omega_X^1(\log D)/f^*\Omega_B^1(\log\Sigma)$ has
\[
f_*\mathcal E\simeq\mathcal O\oplus\mathcal O(-1),\qquad
R^1f_*\mathcal E\simeq\mathcal O^{\oplus2}\oplus\mathcal O(-1)^{\oplus2},
\qquad R^2f_*\mathcal E\simeq\mathcal O\oplus\mathcal O(-1).
\]
All three are locally free and commute with base change.
At a finite point write
$t=s^m$ and $\ell=dt/t=m\,ds/s$, with $(m,a)=(3,2)$ at $p_1$ and
$(m,a)=(4,1)$ at $p_2$, and let
$D_{\rm KS}=(d_{rs})$ be the marked Kodaira--Spencer matrix after removing
$ds$.  Dots in the raw period determinant mean $d/ds$, and
\[
 K=\det(\dot Z)=6\dot\mu^{\,2}-\dot\tau\dot\beta.
\]
With $\zeta_m=e^{-2\pi i/m}$, choose relative-cotangent eigenframes
$(e_0,e_1)$ with characters
$(1,\zeta_m^a)$ and $R^1\mathcal O$ eigenframes $(b_0,b_1)$ with characters
$(1,\zeta_m^{-a})$.  The connecting maps are typed as
\[
 \kappa_0:f_*\mathcal E\longrightarrow
 \mathcal L\otimes R^1f_*\mathcal O_X,\qquad
 \kappa_1:R^1f_*\mathcal E\longrightarrow
 \mathcal L\otimes R^2f_*\mathcal O_X\simeq\mathcal O.
\]
In the actual invariant source basis $(e_0,s^{m-a}e_1)$ and target basis
$(b_0,s^ab_1)\ell$, put
\[
 A=\frac{s}{m}d_{00},\quad B=\frac{s^{1-a}}m d_{01},\quad
 C=\frac{s^{1-a}}m d_{10},\quad
 D_0=\frac{s^{m+1-2a}}m d_{11}.
\]
The two connecting maps then have the exact matrices
\[
 [\kappa_0]=\begin{pmatrix}A&tB\\ C&D_0\end{pmatrix},
 \qquad
 [\kappa_1]=(-C,A,-D_0,B),
\]
where the ordered source of the row is
$(e_0b_0,s^ae_0b_1,s^{m-a}e_1b_0,e_1b_1)$ and its target generator is
$s^a(b_0\wedge b_1)\ell$.  Thus
\[
 \operatorname{div}(\det\kappa_0)=p_1+p_2,\qquad
 \operatorname{ord}_{z_1}K=2,\quad \operatorname{ord}_{z_2}K=0,\qquad
 \kappa_1\text{ is sheaf-surjective}.
\]
The residue normalizations are also exact:
\[
 \mathcal O(1)\longrightarrow
 \mathbb C_{p_0}\oplus\mathbb C_{p_1}\oplus\mathbb C_{p_2},
 \qquad
 \sigma\longmapsto(\sigma(p_0),3\sigma(p_1),4\sigma(p_2)),
\]
\[
 \operatorname{res}_{1,j}(xb_0+ys^ab_1)=m_jx[b_0],
\]
and
\[
 \operatorname{im}(\kappa_0\bmod t)
 =\ker\operatorname{res}_{1,j}=\mathbb C\,s^ab_1.
\]
Hence the two finite cokernel skyscrapers map isomorphically to
$H^1(S_1,\mathcal O)\oplus H^1(S_2,\mathcal O)$ and cancel in the absolute
residue sequence.

At the cusp, the quotient from ordinary to logarithmic relative
differentials is checked on the $xy$ and $xyz$ modules.  In the fixed global
branch ordering its two triple-point maps are
\[
 d_P(a)=(a_0-a_1,a_0-a_2,a_1-a_2),\qquad
 d_Q(a)=(a_2-a_1,a_0-a_1,a_0-a_2),
\]
Put $\eta_1=d\log x_1$ and $\eta_2=d\log x_2$.  These cusp logarithmic
forms lie in the logarithmic-form module, whereas the finite-filling
covectors denoted $\eta_j$ above lie in the filling-dual lattice.  The exact
morphism used below is the branch-difference map $\rho$ from the former
module to the branch-difference lattice; no identification of the two source
modules is made.
Both branch maps give
\[
 \rho(\eta_1)=q_1=(-1,0,1),\qquad
 \rho(\eta_2)=q_2=(0,-1,-1),
\]
which is the degree-zero branch-difference isomorphism.  With
$\alpha_i=\delta_0(q_i)$ and the standard \v Cech differential, the
total-complex sign is
\[
 \delta_1\bigl(\rho_*(\alpha\otimes\eta_i)\bigr)
 =-\alpha\smile\delta_0(q_i)=-\alpha\smile\alpha_i.
\]
Thus, for
$\theta=-\delta_1^{-1}(\alpha_1\smile\alpha_2)$, the induced row is exactly
\[
 \rho_*=(0,1,-1,0):
 \mathbb C\langle
 \alpha_1\eta_1,\alpha_1\eta_2,\alpha_2\eta_1,\alpha_2\eta_2
 \rangle\longrightarrow\mathbb C\theta.
\]
The same local modules show that $\Omega^1_{X/B}$ is torsion-free but
non-locally-free at the cusp; its torsion is supported precisely on the two
multiple fibres.

For the remaining scalar, rescale the Cayley coordinate to
$s=(z_2-\bar z_2)s_2$, so that
$s\circ g_2=-is$ and $ds/dz(z_2)=1$; write a dot for $d/ds$, and put
$\ell=dt/t=4\,ds/s$.  In the marked fibre coordinate
$\zeta=(\zeta_1,\zeta_2)$, the invariant section
$\xi=\partial/\partial\zeta_2$ is nowhere zero in the locally free relative
sheaf $T_{X/B}=\mathcal E^\vee$ (although its ambient tangent-space value can
vanish at a toric fixed point), and the unique lift of its dual unit is
$h=(a(z),1)(d\zeta_1,d\zeta_2)^{\mathsf t}$ with $a(z_2)=1/2$.
At the order-four lift the genuine invariant target frame is
\[
 \omega_2=s\,u_2(s)(\gamma\wedge u)\ell,\qquad u_2(0)=1;
\]
here $b_1$ is normalized by
$\gamma\wedge b_1=u_2(s)(\gamma\wedge u)$.
$s(\gamma\wedge u)\ell$ alone is not invariant away from $s=0$.
In this frame the signed coarse value is
\[
 \kappa_1(h\otimes\gamma)(p_2)
 =-\frac14\bigl(a(z_2)\dot\tau(z_2)+\dot\mu(z_2)\bigr).
\]
Since $\operatorname{ord}_{z_2}(\tau-i)=2$, one has
$\dot\tau(z_2)=0$ and
$K(z_2)=6\dot\mu(z_2)^2\ne0$; hence the scalar is exactly
$-\frac14\dot\mu(z_2)\ne0$.  Therefore
\[
f_*\Omega_X^1\simeq\mathcal O(-2),\quad
R^1f_*\Omega_X^1\simeq\mathcal O\oplus\mathcal O(-1)^2,\quad
R^2f_*\Omega_X^1\simeq\mathcal O\oplus\mathcal O(-1),
\]
and the exact Hodge table is
\[
\begin{array}{c|cccc}
q\backslash p&0&1&2&3\\ \hline
3&0&0&0&1\\
2&0&1&2&1\\
1&1&2&1&0\\
0&1&0&0&0.
\end{array}
\]
The two middle Fr\"olicher rows are exact copies of
$0\to\mathbb C\to\mathbb C^2\to\mathbb C\to0$.  Thus the sequence does
not degenerate at $E_1$, while
$\dim E_1=10$, $E_2=E_\infty$, and $\dim E_2=2$.  See
 \cref{prop:log-relative-direct-images,prop:log-ks-divisor,%
 prop:absolute-cotangent-direct-images,thm:h12-scalar,cor:exact-hodge-diamond}.
The internal calculation is \Verified, and a further context-isolated AI
derivation has independently recovered the invariant lattices, determinant
divisor, residue cancellation, scalar test, direct-image splittings, and the
same values $h^{1,2}=1$, $h^{1,1}=2$.  The separate H2 audit record remains
\Open\ only because independent human specialist replication is still
outstanding; that external replication boundary is not a premise of the
Hodge theorem.

\item \Verified\ \textbf{Every Bott--Chern and Aeppli dimension is fixed.}
The auxiliary spaces in McHugh's compact-threefold formulas satisfy
$k^{p,q}\le b_{p+q}$: the de Rham image of the pure, doubly closed forms
surjects onto $K^{p,q}$.  Hence they vanish in all middle degrees for this
$X$; also
$E_2^{3,1}=0$.  Thus
\[
h^{1,1}_{BC}=2,\quad h^{1,2}_{BC}=h^{2,1}_{BC}=1,\quad
h^{2,2}_{BC}=2,\quad h^{2,3}_{BC}=h^{3,2}_{BC}=1,
\]
besides the two endpoint copies of $\mathbb C$, and all other Bott--Chern
dimensions vanish.  Schweitzer duality gives the complete Aeppli diamond.
Equivalently,
\[
(h^k_{BC})=(1,0,2,2,2,2,1),\quad
(h^k_A)=(1,2,2,2,2,0,1),\quad
(\Delta^k)=(0,2,4,4,4,2,0).
\]
See \cref{thm:exact-bott-chern-aeppli}.  Coordinate-level representatives
and the product structure are certified separately in the next item.

\item \Verified\ \textbf{The Bott--Chern algebra has explicit generators
and square-zero positive ideal.}
Choose global $\bar\partial$-closed smooth representatives
$G\in\mathcal A^{0,1}(X)$ and $Q\in\mathcal A^{1,1}(X)$ of the two
candidate-specific Dolbeault classes $\mathfrak g$ and $\mathfrak q$.
Here $\mathfrak g$ is distinguished, while $\mathfrak q$ is a chosen
coherent-Leray lift, unique modulo $\operatorname{im}d_1$.
Set
\[
 \mathsf A=[\partial G]_{BC},\qquad
 \mathsf B=[\bar\partial\bar G]_{BC}.
\]
On a marked smooth fibre,
\[
 \sigma_1=6\mu\gamma+\tau u+w,\qquad
 \sigma_2=\beta\gamma+\mu u+\delta,
\]
put
\[
 y=\operatorname{Im}\tau,\qquad m=\operatorname{Im}\mu,\qquad
 r=\operatorname{Im}\beta,\qquad
 \Delta=r-\frac{6m^2}{y}<0,
\]
and define
\[
 \Gamma=\frac{i}{2\Delta}
 \left(d\bar\zeta_2-\frac{m}{y}d\bar\zeta_1\right),\qquad
 U=\frac{i}{2y\Delta}
 \left(r\,d\bar\zeta_1-6m\,d\bar\zeta_2\right).
\]
These represent the fibre images of \(\gamma\) and \(u\), and the global
classes satisfy
\[
 i_b^*\mathfrak g=[\Gamma]_{\bar\partial},\qquad
 i_b^*\mathfrak q=[Q_F]_{\bar\partial}.
\]
The Appell--Humbert matrix itself is
\[
 H_\eta=\frac1{y\Delta}
 \begin{pmatrix}r&-6m\\-6m&6y\end{pmatrix},
 \qquad H_\eta=2G_\eta^{-1},\qquad
 \det H_\eta=\frac6{y\Delta}<0.
\]
The load-bearing wedge identity is
\[
 \eta=u\wedge w+6\gamma\wedge\delta
 =u\wedge\sigma_1+6\gamma\wedge\sigma_2
 =-\sigma_1\wedge u-6\sigma_2\wedge\gamma,
\]
which gives the exact fibre representative
\[
 Q_F=-d\zeta_1\wedge U-6d\zeta_2\wedge\Gamma,
 \qquad
 \eta\wedge\eta=12\gamma\wedge u\wedge w\wedge\delta.
\]
The marked and complex orientations are opposite because the real
determinant of the marked lattice columns is
$y\Delta<0$.  Accordingly,
\[
 \int_{(F_b,{\rm or}_{\rm mark})}Q_F^2=12,
 \qquad
 \int_{(F_b,{\rm or}_{\mathbb C})}Q_F^2=-12.
\]
The fibre Dolbeault--Serre pairing is normalized by
\[
 \int_{F_b}Q_F\wedge\Gamma\wedge d\zeta_2=-1
 \quad\text{for the complex orientation},
\]
so $[\Gamma\wedge Q_F]_{\bar\partial}\ne0$.

The classes $d_1\mathfrak g$ and $\mathfrak q$ are independent in
$H^{1,1}_{\bar\partial}(X)$: the first restricts to zero on every smooth
fibre, whereas the second restricts to the nonzero class of $Q_F$.
Exactness of the first middle Fr\"olicher row therefore gives
\[
 d_1\mathfrak q=[\partial Q]_{\bar\partial}\ne0,
\]
so $\mathsf C=[\partial Q]_{BC}$ and its conjugate are nonzero.  Moreover,
putting
\[
 v=[G\wedge Q]_{\bar\partial}\in H^{1,2}_{\bar\partial}(X)
\]
gives
$i_b^*v=[\Gamma\wedge Q_F]_{\bar\partial}\ne0$.  Thus $v$ spans
$H^{1,2}_{\bar\partial}(X)$, and injectivity in the second middle row gives
\[
 d_1v=[\partial(G\wedge Q)]_{\bar\partial}\ne0.
\]
Consequently $\mathsf D=[\partial(G\wedge Q)]_{BC}$ and its conjugate are
nonzero.

The global integral normalization is equally explicit.  If the base
orientation is chosen so that
$d_2(12\gamma)=c_1(\mathcal O_B(1))$ and
$\exp(z)=e^{2\pi iz}$, the holomorphic structure
$\bar\partial_zs=-2\pi izG\otimes s$ on the smooth trivial line has
unit-metric Chern connection and curvature
\[
 \nabla_z=d-2\pi i(zG+\bar z\,\bar G),\qquad
 F_{\nabla_z}=-2\pi i
 (z\,\partial G+\bar z\,\bar\partial\bar G).
\]
With the convention
$c_{1,BC}(L,h)=[\,iF_{\nabla^h}/(2\pi)\,]_{BC}$, one obtains
\[
 \exp(12\mathfrak g)=f^*\mathcal O_B(1),\qquad
 c_{1,BC}(\exp(z\mathfrak g))
 =z\mathsf A+\bar z\,\mathsf B,
\]
and hence
\[
 c_{1,BC}(f^*\mathcal O_B(1))=12(\mathsf A+\mathsf B).
\]
With the opposite primitive orientation both occurrences of $12$ acquire
the common sign $-1$:
\[
 \exp(-12\mathfrak g)=f^*\mathcal O_B(1),\qquad
 c_{1,BC}(f^*\mathcal O_B(1))=-12(\mathsf A+\mathsf B).
\]
Thus the canonical integral Picard--Leray datum fixes the absolute factor
$12$; it does not assert that $\mathsf A,\mathsf B$ are integral
Bott--Chern classes.

Put
\[
 \mathbf1=[1]_{BC},\qquad
 \mathsf T=[\nu_X]_{BC},
 \quad \nu_X\text{ a smooth positive }(3,3)\text{-form with }
 \int_X\nu_X=1.
\]
Together with the unit, the following nine positive-degree classes form
bases of all positive-degree Bott--Chern groups, with $\mathsf T$ the
normalized top class:
\[
 \mathsf A=[\partial G]_{BC},\ \mathsf B=[\bar\partial\bar G]_{BC},\
 \mathsf C=[\partial Q]_{BC},\ \overline{\mathsf C},\
 \mathsf D=[\partial(G\wedge Q)]_{BC},\ \overline{\mathsf D},\
 \mathsf E=[\partial(Q^2)]_{BC},\ \overline{\mathsf E},\ \mathsf T.
\]
If $N$ denotes their span, then
\[
 H^{*,*}_{BC}(X)=\mathbb C\mathbf1\oplus N,
 \qquad N^2=0.
\]
For the last one-dimensional generator,
$[Q^2]_{\bar\partial}$ has nonzero fibre restriction, whereas every class
in
$\operatorname{im}(d_1:H^{1,2}_{\bar\partial}\to
H^{2,2}_{\bar\partial})$ restricts to zero.  Exactness therefore gives
\[
 d_1[Q^2]_{\bar\partial}
 =[\partial(Q^2)]_{\bar\partial}\ne0,
\]
which certifies $\mathsf E$ (and conjugation certifies
$\overline{\mathsf E}$).
Equivalently every product of two positive-degree classes is zero.  Put
$a=\partial G$, $a^\dagger=\bar\partial\bar G$, and $c=\partial Q$.
The exact primitive
\[
 \partial\bar\partial(G\wedge\bar G)=-a\wedge a^\dagger
\]
gives $\mathsf A\mathsf B=0$.  Represent
$c_{1,BC}(f^*\mathcal O_B(1))$ by the pullback Fubini--Study Chern form;
its square vanishes pointwise because $B$ is a curve.  Squaring the
factor-$12$ identity and using $\mathsf A\mathsf B=0$ gives
$\mathsf A^2+\mathsf B^2=0$.  Moreover
\[
 [(a^\dagger)^2]_{\bar\partial}=0,\qquad
 [a^2]_{\bar\partial}=d_1[G\wedge a]_{\bar\partial}=0.
\]
Injectivity of
$d_1:H^{1,2}_{\bar\partial}\to H^{2,2}_{\bar\partial}$ then gives a form
$\Phi$ with
\[
 G\wedge a=\bar\partial\Phi,\qquad
 a^2=\partial\bar\partial\Phi,
\]
so $\mathsf A^2=\mathsf B^2=0$.  The last below-top product requiring a
pairing is disposed of by
\[
 \int_Xa\wedge c\wedge G
 =-\int_XG\wedge a\wedge c
 =-\int_X\bar\partial\Phi\wedge c=0,
\]
together with the perfect Dolbeault--Serre pairing and the isomorphism
$H^{3,2}_{BC}\to H^{3,2}_{\bar\partial}$; the
Bott--Chern--Aeppli perfect pairing alone would not prove this vanishing.
Every top-degree product of two positive-degree generators contains a de
Rham exact factor, while
\[
 H^{3,3}_{BC}(X)\xrightarrow{\ \int_X\ }\mathbb C
\]
is an isomorphism, so those products vanish as well.  See
\cref{prop:bcpr-generators-are-bases,thm:bcpr-square-zero-algebra}.  This
asserts global smooth representatives, but not a preferred literal
coefficient atlas through the two finite fillings and the cusp:
$[i_b^*Q]_{\bar\partial}=[Q_F]_{\bar\partial}$ does not require
$i_b^*Q=Q_F$ as forms.  The chosen $Q$ cannot be real, although its fibre
class is real.  After fixing the distinguished $\mathfrak g$ and one
coherent-Leray lift $\mathfrak q$, changing the smooth representatives and
horizontal corrections changes $G,Q$ by Dolbeault coboundaries and changes
none of the displayed classes, pairings, or products.  This is a
candidate-specific rank-one line-bundle normalization, not a statement
about arbitrary gauge theories.
\item \Verified\ \textbf{The positive radial cusp has an exact flat
period, inverse-dual, and scalar-gap certificate.}
Put
\[
 \mu_T=\alpha_T+i m_T,\qquad h_T=\eta_T+i H_T,\qquad
 b_T=\xi_T+i B_T,\qquad c_T=\xi_T-\eta_T,
\]
\[
 L_T=T+H_T,\qquad q_T=L_T-B_T,\qquad
 D_T=L_Tq_T+6m_T^2.
\]
Set also
\[
 m=\operatorname{Im}\mu(0),\qquad
 H=\operatorname{Im}h(0),\qquad
 B=\operatorname{Im}b(0).
\]
In the ordered marked lattice and interleaved real coordinate bases,
\[
 P_T=
 \begin{pmatrix}
  6\alpha_T&\eta_T&1&0\\
  6m_T&L_T&0&0\\
  c_T&\alpha_T&0&1\\
  -q_T&m_T&0&0
 \end{pmatrix},\qquad
 \det P_T=-D_T,\qquad
 \operatorname{covol}(P_T\mathbb Z^4)=D_T>0,\qquad
 D_{\mathrm{can},T}=-\frac{D_T}{L_T}<0,
\]
and the complete inverse transpose is
\[
 P_T^{-\mathsf T}=
 \begin{pmatrix}
 0&0&1&0\\[2pt]
 \dfrac{m_T}{D_T}&\dfrac{q_T}{D_T}&
 \dfrac{-6m_T\alpha_T-q_T\eta_T}{D_T}&
 \dfrac{-m_Tc_T-q_T\alpha_T}{D_T}\\[7pt]
 0&0&0&1\\[2pt]
 -\dfrac{L_T}{D_T}&\dfrac{6m_T}{D_T}&
 \dfrac{6(L_T\alpha_T-m_T\eta_T)}{D_T}&
 \dfrac{L_Tc_T-6m_T\alpha_T}{D_T}
 \end{pmatrix}.
\]
Consequently
\((P_T\mathbb Z^4)^*=P_T^{-\mathsf T}\mathbb Z^4\).  More explicitly,
for \(k=(r,s,n_1,n_2)^{\mathsf t}\), define
\[
 R_T(k)=r-6\alpha_Tn_1-c_Tn_2,\qquad
 S_T(k)=s-\eta_Tn_1-\alpha_Tn_2.
\]
The corresponding dual vector is
\[
 \left(n_1,
 \frac{m_TR_T(k)+q_TS_T(k)}{D_T},n_2,
 \frac{-L_TR_T(k)+6m_TS_T(k)}{D_T}\right),
\]
and the scalar eigenvalue is \(4\pi^2\) times its squared norm.
For every sufficiently large \(T\),
\[
 \lambda_1(F_T,g_T)
 =4\pi^2\min\left\{
 \frac{L_T^2+m_T^2}{D_T^2},
 \frac{q_T^2+36m_T^2}{D_T^2}\right\}
 =\frac{4\pi^2}{T^2}\left(
 1+\frac{-2H+2\min\{B,0\}}T+O(T^{-2})\right).
\]
This is the gap of the explicitly declared Euclidean quotient metric.  Its
unit-volume rescaling is a different operator family and obeys
\[
 \widehat g_T=D_T^{-1/2}g_T,\qquad
 \lambda_1(F_T,\widehat g_T)
 =\frac{4\pi^2}{T}\left(
 1+\frac{-H-B/2+2\min\{B,0\}}T+O(T^{-2})\right)
 \longrightarrow0.
\]
The degeneration has an exact fixed-energy density, not merely a lowest-mode
estimate.  Counting Fourier multiplicities by
\[
 N_T(E)=\#\{k\in\mathbb Z^4:
 4\pi^2\lVert P_T^{-\mathsf T}k\rVert^2\le E\},
 \qquad R_E^2=\frac{E}{4\pi^2},
\]
one has, for every fixed \(E\ge0\),
\[
 N_T(E)=\pi D_T\sum_{n\in\mathbb Z^2}
 (R_E^2-\lVert n\rVert^2)_+ +O_E(T+1),
 \qquad (x)_+=\max\{x,0\}.
\]
The sum is finite.  In particular, below the first transverse threshold,
\[
 0<E<4\pi^2
 \quad\Longrightarrow\quad
 N_T(E)=\frac{E}{4\pi}D_T+O_E(T+1)
       =\frac{E}{4\pi}T^2+O_E(T).
\]
The proof writes the admissible \((r,s)\)-indices for each fixed transverse
\(n\in\mathbb Z^2\) as the lattice points in a translated ellipse of area
\(\pi D_T(R_E^2-\lVert n\rVert^2)\) and perimeter \(O_E(T+1)\), then applies
the elementary planar area--perimeter lattice bound.  Thus the remainder is
uniform in the otherwise unrestricted real parts of the periods.  See
\cref{cusp-spec-counting}.

Normalized Haar measure in the angular-coordinate direct integral is a
unitary measure convention, not this metric rescaling.  See
 \cref{cusp-spec-exact-dual-lattice,cusp-spec-shortest-modes,%
 cusp-spec-unit-volume,cusp-spec-counting}.

\item \Verified\ \textbf{The native
\(\mathbb T^2\times\mathbb R^2\) limit and the ancillary covered
\(\mathbb R^4\) limit have different exact constructions.}
The real projection \(\rho(x_1,y_1,x_2,y_2)=(y_1,y_2)\) gives
\[
 0\longrightarrow\mathbb T^2_{\mathrm{re}}
 \longrightarrow F_T
 \longrightarrow Q_T:=\mathbb R^2/Y_T\mathbb Z^2
 \longrightarrow0,
 \qquad
 Y_T=\begin{pmatrix}6m_T&L_T\\-q_T&m_T\end{pmatrix},
 \qquad \det Y_T=D_T.
\]
The quotient systole is at least \(T-C\), whereas the kernel retains its
two unit periods; this is the native pointed
\(\mathbb T^2\times\mathbb R^2\) degeneration.  Adding
\[
 \mathsf S_M=\operatorname{diag}(1,1,M,M),\qquad
 t_c=v^M
\]
produces a degree-\(M^2\) relative fibre cover over \(v\ne0\), because
\[
 \mathsf S_M^{-1}M_0^M\mathsf S_M=M_0.
\]
The cover character indexed by \(k\in\mathbb Z^4\) has momentum
\[
 P_T^{-\mathsf T}\mathsf S_M^{-1}k,
\]
and it is pulled back from the original fibre exactly when
\(k_3,k_4\) are divisible by \(M\).  The covering radius of this momentum
 lattice is \(O(T^{-1}+M^{-1})\); hence \(T,M\to\infty\) fills all four
 momentum directions and gives eventual exact Euclidean balls.  This argument
 by itself proves only a pointed local \(\mathbb R^4\) limit over the
 punctured base; the modified-action completion across the singular fibre is
 the distinct local Kummer construction recorded below, and neither argument
 alone gives a global cover of the completed threefold or an S12/S24-type
 globalization.

There is, however, an exact four-copy regular-locus construction.  The
monodromy-invariant closure of the compact cusp lattice is
\[
 L_\Lambda=\mathbb Z\langle\widehat u,\widehat w,\widehat\delta\rangle,
\qquad
 \Lambda/L_\Lambda\cong\mathbb Z\langle\widehat\gamma\rangle.
\]
The induced real circle quotient has exact circumference
\[
 \ell_T=\frac{D_T}{\sqrt{L_T^2+m_T^2}}
     =T+H-B+O(T^{-1})\longrightarrow\infty.
\]
The fourfold fibre product of \(X^\circ=f^{-1}(B^\circ)\) therefore has a
real fibrewise-group \(\mathbb T^4\) quotient over \(B^\circ\) whose
synchronized radial fibres converge pointed-smoothly to \(\mathbb R^4\).
Two copies are insufficient for this monodromy-invariant-kernel mechanism;
four are minimal, and no finite cover of the base repairs the native two-copy
quotient.  The two finite fillings admit explicit simultaneous-root
normalizations over Seifert targets, with degrees \(27\) and \(64\) on the
reduced central fibres.  The toric cusp does not: two paths converging to the
same triple point retain different character values \(3/4\) and \(1/2\), so
the literal quotient and its fourfold product have no continuous cusp
 extension.  The Cartesian source is not \(S^{24}\):
 \(X^4\simeq(S^6)^4\), while that fixed-target fourfold-product map does not
 factor through the topological smash sphere
\(X^{\wedge4}\cong S^{24}\); Borel--Serre also rules out almost-complex
structures on \(S^{12}\) and \(S^{24}\).  Thus this exact regular-locus
construction answers the globalization question as far as this mechanism can
go and proves its compact-extension obstruction.  See
 \cref{cusp-spec-native-quotient,cusp-spec-globalization-obstruction,%
 cusp-spec-cover-momentum-theorem}.

\item \Verified\ \textbf{The degree-\(N^2\) regular-fibre cover has an exact
completed local toric Kummer model, with direct and relative degrees kept
distinct.}
Write
\[
 Z(s)=sB_0+C(t),\qquad
 B_0=\begin{pmatrix}0&1\\-1&0\end{pmatrix},\qquad
 t=e^{2\pi i s},
\]
and let the target lattice action on the toric variety \(Y=Y_\Sigma\) be
\[
 \Psi_{\bar\lambda}(x,t)
 =\bigl(e^{2\pi iC(t)\bar\lambda}
        t^{B_0\bar\lambda}x,t\bigr),
 \qquad N_0=Y/\bar\Lambda.
\]
After the base change \(t=v^N\), the correct source factor and action are
\[
 C_N(v)=\frac1N C(v^N),\qquad
 \Psi^{(N)}_{\bar\lambda}(\xi,v)
 =\bigl(e^{2\pi iC_N(v)\bar\lambda}
        v^{B_0\bar\lambda}\xi,v\bigr).
\]
They extend freely and properly discontinuously to a copy \(Y^{(N)}\) of
the same toric variety.  Thus
\[
 N_0^{\langle N\rangle}=Y^{(N)}/\bar\Lambda
 \longrightarrow\Delta_v
\]
is a proper holomorphic map from a smooth threefold, with reduced irreducible
nonnormal central fibre \(W_N\), whose normalization is \(\mathrm{dP}_6\).  The toric
power map
\[
 P_N(\xi_1,\xi_2,v)=(\xi_1^N,\xi_2^N,v^N)
\]
is finite flat of degree \(N^3\), is equivariant for the two displayed
actions, and descends to
\[
 \kappa_N:N_0^{\langle N\rangle}\longrightarrow N_0,
 \qquad \deg\kappa_N=N^3,\qquad
 \kappa_N^*W=NW_N,\qquad
 \operatorname{Ram}(\kappa_N)=(N-1)W_N.
\]
If \(N_0[N]=N_0\times_{\Delta_t}\Delta_v\), then \(N_0[N]\) is normal and
\[
 N_0^{\langle N\rangle}\xrightarrow{Q_N}N_0[N]\longrightarrow N_0
\]
is the unique relative factorization.  The generic degree of \(Q_N\) is
\(N^2\); for \(N>1\) it is quasi-etale but not flat, with non-etale and
nonflat loci exactly over the double and triple strata of the base-changed
central fibre.  It is not a normalization.  On reduced central fibres the
map has generic degree \(N^2\), and on normalizations it is the toric map
\[
 [N]:\mathrm{dP}_6\longrightarrow\mathrm{dP}_6,
\]
of degree \(N^2\), ramified with index \(N\) along all six toric boundary
curves and respecting the paired conductor curves and the two triple-point
classes.  Over \(v\ne0\), it is exactly the lattice cover
\[
 \mathbb Z\langle\widehat\gamma,\widehat u,
 N\widehat w,N\widehat\delta\rangle\subset\Lambda,
 \qquad \operatorname{Deck}(Q_N)\cong(\mathbb Z/N\mathbb Z)^2,
 \qquad S_N^{-1}M_0^NS_N=M_0,
\]
where \(S_N=\operatorname{diag}(1,1,N,N)\).  The factor
\(C(v^N)/N\) is essential: reusing the target action would require the
generally false equality
\(e^{2\pi iC(v)\bar\lambda N}=e^{2\pi iC(v^N)\bar\lambda}\).
This is a morphism of completed \emph{local} cusp families.  For \(N>1\)
the cover lattice is not preserved by the full monodromy, so it gives
neither a cover over all of \(B\) nor a product- or smash-sphere map.  See
\cref{cusp-cover-extension-theorem,cusp-cover-naive-action-failure,%
cusp-cover-global-scope}.

\item \Verified\ \textbf{The maps \(q_2\), \(q_{1,\mathrm{cusp}}\), and
\(q_1\) give exact cone-smash carriers with sources \(S^{12}\) and
\(S^{24}\), and their intrinsic metrics, ambient metrics, homology, and
limitations are all explicit.}
Let
\[
 K_\Lambda=\mathbb Z\langle\widehat w,\widehat\delta\rangle,
 \qquad \overline\Lambda=\Lambda/K_\Lambda,
 \qquad \overline T^2=(\overline\Lambda\otimes\mathbb R)/\overline\Lambda.
\]
Fix \(C_{\mathrm{sys}}>0\) so that, for all sufficiently large \(T\),
\[
 Y_T=T\begin{pmatrix}0&1\\-1&0\end{pmatrix}+E_T,
 \qquad \lVert E_T\rVert_{\mathrm{op}}\le C_{\mathrm{sys}}.
\]
Since
\[
 \operatorname{im}(M_0-I)=K_\Lambda=\ker(M_0-I),
\]
the punctured cusp has a fibrewise group epimorphism
\[
 \overline\pi:N_0\setminus W\longrightarrow D_0^*\times\overline T^2
\]
with connected two-torus kernel and trivial target local system.  On the
radial fibre, \(\overline T^2\) is isometric to
\(Q_T=\mathbb R^2/Y_T\mathbb Z^2\) and has systole at least
\(T-C_{\mathrm{sys}}\).

Put
\[
 \mathcal A=\operatorname{Susp}(S^1\times\overline T^2)
 \cong\operatorname{Susp}(T^3).
\]
For \(t(f(x))=re^{i\phi}\), the continuous surjective quotient
\[
 q_2:(X,x_0)\longrightarrow(\mathcal A,v_0),\qquad
 q_2(x)=[r,e^{i\phi},\overline\pi_f(x)]
\]
on \(N_0\setminus W\), where
\(\overline\pi_f=\operatorname{pr}_{\overline T^2}\circ\overline\pi\), sends
\(W\) to \(v_0\) and
\(X\setminus\operatorname{int}N_0\) to the other suspension vertex
\(v_1\).  Its smash square
\[
 q_2^{\wedge2}:X\wedge X\cong S^{12}\longrightarrow
 \mathcal A\wedge\mathcal A
\]
restricts on synchronized radial fibres to
\[
 \pi_T\times\pi_T:F_T^2\longrightarrow Q_T^2,
 \qquad (Q_T^2,0)\longrightarrow(\mathbb R^4,0),
\]
where the convergence is pointed local-isometric, pointed smooth, and
pointed Gromov--Hausdorff, and every fixed radius ball is eventually exactly
Euclidean.  The nonzero reduced integral homology groups are
\[
 \widetilde H_k(\mathcal A)=
 \begin{cases}\mathbb Z^3,&k=2,3,\\ \mathbb Z,&k=4,\end{cases}
\quad
 \widetilde H_k(\mathcal A^{\wedge2})=
 \begin{cases}
 \mathbb Z^9,&k=4,\\ \mathbb Z^{18},&k=5,\\
 \mathbb Z^{15},&k=6,\\ \mathbb Z^6,&k=7,\\ \mathbb Z,&k=8.
 \end{cases}
\]
More precisely,
\[
 \mathcal A^{\wedge2}\simeq
 \bigvee^9S^4\vee\bigvee^{18}S^5\vee\bigvee^{15}S^6
 \vee\bigvee^6S^7\vee S^8.
\]

The cusp-supported circle version uses
\[
 \mathcal B=\operatorname{Susp}(S^1\times S^1)=\operatorname{Susp}(T^2)
\]
and the continuous surjective based quotient
\[
 q_{1,\mathrm{cusp}}:(X,x_0)\longrightarrow(\mathcal B,b_0),
 \qquad
 q_{1,\mathrm{cusp}}(x)=[r,e^{i\phi},\chi_\gamma(x)],
 \qquad q_{1,\mathrm{cusp}}(W)=b_0,
\]
with the exterior collapsed to \(b_1\).  Its fourth smash power has source
\[
 q_{1,\mathrm{cusp}}^{\wedge4}:X^{\wedge4}\cong S^{24}
 \longrightarrow\mathcal B^{\wedge4}
\]
and restricts to
\(\chi_{\gamma,T}^{\times4}:F_T^4\to C_T^4\).  Here
\[
 \widetilde H_k(\mathcal B)=
 \begin{cases}\mathbb Z^2,&k=2,\\ \mathbb Z,&k=3,\end{cases}
\quad
 \widetilde H_k(\mathcal B^{\wedge4})=
 \begin{cases}
 \mathbb Z^{16},&k=8,\\ \mathbb Z^{32},&k=9,\\
 \mathbb Z^{24},&k=10,\\ \mathbb Z^8,&k=11,\\ \mathbb Z,&k=12.
 \end{cases}
\]

For the full-base version, glue the regular circle target and the two finite
Seifert targets to obtain the compact three-manifold with boundary
\[
 \mathcal M=(P\times S^1)\cup C_1\cup C_2,
 \qquad \partial\mathcal M=\partial D_0\times S^1,
\]
and cone its boundary:
\[
 \mathcal C=\mathcal M\cup_{\partial\mathcal M}
 \operatorname{Cone}(\partial\mathcal M).
\]
The map \(q_1:X\to\mathcal C\) is \((f,\theta_\gamma)\) over \(P\),
equivalently \(\chi_\gamma|_{f^{-1}(P)}\), is the
finite-filling map \(N_j\to C_j\) for \(j=1,2\), and on the cusp is
\[
 q_1(x)=[r,(e^{i\phi},\theta_\gamma(x))],\qquad q_1(W)=c_0.
\]
It is continuous, surjective, and based.  Its fourth smash power is
\[
 q_1^{\wedge4}:X^{\wedge4}\cong S^{24}\longrightarrow\mathcal C^{\wedge4},
\]
and its synchronized radial restriction is again
\(\chi_{\gamma,T}^{\times4}:F_T^4\to C_T^4\), with
\[
 \operatorname{sys}(C_T^4)=\ell_T=T+H-B+O(T^{-1}),
 \qquad (C_T^4,0)\longrightarrow(\mathbb R^4,0).
\]
The target homology is
\[
 \widetilde H_k(\mathcal C)=
 \begin{cases}\mathbb Z,&k=2,3,\end{cases}
\quad
 \widetilde H_k(\mathcal C^{\wedge4})=
 \begin{cases}
 \mathbb Z,&k=8,12,\\ \mathbb Z^4,&k=9,11,\\
 \mathbb Z^6,&k=10.
 \end{cases}
\]
In fact,
\[
 \mathcal C^{\wedge4}\simeq
 S^8\vee\bigvee^4S^9\vee\bigvee^6S^{10}
 \vee\bigvee^4S^{11}\vee S^{12}.
\]
All omitted degrees in these three homology displays are zero.
The conical-target maps \(q_2^{\wedge2}\) and \(q_1^{\wedge4}\) induce zero on
all positive-degree integral homology and cohomology and have zero ordinary
Hurewicz image; neither has a continuous homotopy section.  Their unstable
homotopy classes remain \Open, because those facts do not decide nullity.  The
class of the cusp-supported conical-target map
\(q_{1,\mathrm{cusp}}^{\wedge4}:S^{24}\to\mathcal B^{\wedge4}\) has likewise
not been computed.

The two-copy and four-copy exponents are sharp for the unscaled quotient
norm induced by the marked Euclidean fibre metric.  If a primitive lattice
\(N_{\mathrm{quot}}\subset\Lambda^{\oplus r}\) defines a connected quotient
over the punctured cusp whose radial systole tends to infinity, then every
bounded-length class must vanish, so
\[
 K_\Lambda^{\oplus r}\subset N_{\mathrm{quot}},\qquad
 \operatorname{rank}(\Lambda^{\oplus r}/N_{\mathrm{quot}})\le2r.
\]
Thus a rank-four quotient needs at least two copies, and equality is attained
by \((\Lambda/K_\Lambda)^{\oplus2}\).  Over the full regular base, including
after a finite base cover, monodromy invariance and primitivity additionally
force \(\widehat u\in N_{\mathrm{quot}}\): indeed
\(A_1\widehat w=\widehat u-\widehat w\), while a finite-index subgroup contains
some power \(G_*^m\) satisfying
\(G_*^m\widehat w=\widehat w+m(\widehat u-\widehat\delta)\).
Consequently
\[
 \operatorname{rank}(\Lambda^{\oplus r}/N_{\mathrm{quot}})\le r,
\]
so four copies are necessary and sufficient there.  This proves the exact
minimality of the \(S^{12}\) cusp-local and \(S^{24}\) full-base smash
sources within this native connected torus-quotient class; it asserts no
minimality among arbitrary continuous quotient maps or after changing to a
\(T\)-dependent rescaling.  See \cref{cusp-cone-optimal-exponents}.

There are also sphere targets, not merely conical targets.  Let
\[
 K^{(1)}=((\mathbb R/\mathbb Z)\times\{0\})
 \cup(\{0\}\times(\mathbb R/\mathbb Z)),\qquad
 c:\overline T^2\twoheadrightarrow\overline T^2/K^{(1)}\cong S^2,
\]
where \(c\) collapses the coordinate one-skeleton.  Forgetting the cusp base
angle in the reduced suspensions gives continuous surjections
\[
 \widehat p_2:\mathcal A\twoheadrightarrow S^3,
 \quad[r,e^{i\phi},u]\longmapsto[r,c(u)],
 \qquad
 \widehat p_1:\mathcal B\twoheadrightarrow S^2,
 \quad[r,e^{i\phi},\theta]\longmapsto[r,\theta].
\]
Thus \(\widehat q_3=\widehat p_2\circ q_2\) and
\(\widehat q_\gamma=\widehat p_1\circ q_{1,\mathrm{cusp}}\) induce, after the
chosen based diffeomorphism \(X\cong_{\mathrm{diff}}S^6\), quotient maps
\[
 \widehat q_3^{\wedge2}:S^{12}\twoheadrightarrow S^6,
 \qquad
 \widehat q_\gamma^{\wedge4}:S^{24}\twoheadrightarrow S^8.
\]
Both full sphere carriers are based null-homotopic, by two distinct exact
calculations.  First let \(f:S^6\to S^2\) be any based map.  The factorization
\[
 S^{12}=S^6\wedge S^6
 \xrightarrow{\ f\wedge\operatorname{id}_{S^6}=\Sigma^6f\ }
 S^2\wedge S^6=S^8
 \xrightarrow{\ \operatorname{id}_{S^2}\wedge f\ }
 S^2\wedge S^2=S^4
\]
and
\(\pi_{12}(S^8)\cong\pi_4^{\mathrm S}=0\) show that \(f\wedge f\) is
based null-homotopic.  Associativity gives
\(f^{\wedge r}=(f\wedge f)\wedge f^{\wedge(r-2)}\), so
\[
 [f^{\wedge r}]=0\in\pi_{6r}(S^{2r})
 \qquad(r\ge2).
\]
Taking \(f=\widehat q_\gamma\) proves, in particular,
\[
 [\widehat q_\gamma^{\wedge2}]=0\in\pi_{12}(S^4),
 \qquad
 [\widehat q_\gamma^{\wedge4}]=0\in\pi_{24}(S^8).
\]
Thus the full surjective quotient \(S^{24}\twoheadrightarrow S^8\) is based
null-homotopic.

For the other carrier, let \(g:S^6\to S^3\) be any based map.  Its exact
smash factorization is
\[
 g\wedge g=(\Sigma^3g)\circ(\Sigma^6g):
 S^{12}\longrightarrow S^9\longrightarrow S^6.
\]
Writing \(\operatorname{st}\) for stabilization, the low-stem data are
\[
 \pi_6(S^3)\cong\mathbb Z/12,
 \quad
 \pi_3^{\mathrm S}=\mathbb Z/24\langle\nu\rangle,
 \quad
 \pi_6^{\mathrm S}=\mathbb Z/2\langle\nu^2\rangle,
\]
\[
 \pi_{12}(S^6)=\mathbb Z/2\langle\nu_6\circ\nu_9\rangle,
 \qquad
 \operatorname{st}:\pi_{12}(S^6)
 \xrightarrow{\ \cong\ }\pi_6^{\mathrm S}.
\]
The full groups are tabulated in \cite[p.~339]{Hatcher2002}; the signed smash
formula and the unstable stabilization calculation are
\cite[Propositions~3.1 and~5.11]{Toda1962}.
Because \(12[g]=0\), its stable image is
\(\operatorname{st}([g])=2a\nu\) for some \(a\in\mathbb Z/12\).  Stable
smash multiplication gives
\[
 \operatorname{st}([g\wedge g])
 =\operatorname{st}([g])^2
 =4a^2\nu^2=0\in\pi_6^{\mathrm S}.
\]
The displayed stabilization isomorphism forces
\([g\wedge g]=0\in\pi_{12}(S^6)\).
Associativity again shows that \(g^{\wedge r}\) is based null-homotopic for
every \(r\ge2\).  Taking \(g=\widehat q_3\) therefore proves
\[
 [\widehat q_3^{\wedge2}]=0\in\pi_{12}(S^6),
 \qquad
 [\widehat q_\gamma^{\wedge4}]=0\in\pi_{24}(S^8).
\]
Both maps retain their point-set surjectivity and quotient-map status.  The
one-copy class \([\widehat q_3]\in\pi_6(S^3)\) is now zero by
Theorem~\ref{h24:thm:cusp-g-zero}, while
\([\widehat q_\gamma]\in\pi_6(S^2)\) remains open. Neither full-map
null theorem settles any of the conical-target classes named above; see
\cref{cusp-cone-S12-sphere-map-null,cusp-cone-S24-sphere-map-null,%
cusp-cone-sphere-map-homotopy-status,cusp-cone-target-homotopy-types}.
Give \(S_T^2=\overline T^2/K^{(1)}\) the quotient length metric transported
from \(Q_T\), and give both smash products the quotient length metrics from
the Riemannian \(\ell^2\)-product metrics on \(Q_T^2\) and \(C_T^4\).
These moving-slice metrics are assigned intrinsically and are not restrictions
of one fixed ambient sphere metric.  The targets contain moving topologically
embedded four-spheres
\[
 \mathbb S_{12,T}^4=S_T^2\wedge S_T^2\subset S^6,
 \qquad
 \mathbb S_{24,T}^4=C_T^{\wedge4}\subset S^8,
\]
and the synchronized radial restrictions are exactly
\[
 F_T^2\xrightarrow{\pi_T^{\times2}}Q_T^2
 \xrightarrow{c_T^{\times2}}(S_T^2)^2
 \longrightarrow\mathbb S_{12,T}^4,
 \qquad
 F_T^4\xrightarrow{\chi_{\gamma,T}^{\times4}}C_T^4
 \longrightarrow\mathbb S_{24,T}^4,
\]
where the last arrows collapse the corresponding fat wedges.  The full maps
from \(F_T^2\) and \(F_T^4\) retain the nontrivial kernels of their first
arrows; at the quotient stage, precisely the displayed skeleton/fat-wedge
loci are collapsed.  At the
midpoint basepoints, every radius-\(R\) ball is exactly a Euclidean four-ball
once \(T-C_{\mathrm{sys}}>4R\), respectively \(\ell_T>4R\).  Hence both
moving four-spheres converge to \(\mathbb R^4\) in pointed local-isometric and
pointed Gromov--Hausdorff topology; on those eventual Euclidean basepoint
neighborhoods the convergence is also smooth.  Each moving four-sphere is
null-homotopic in its ambient sphere by sliding its suspension coordinate to
an endpoint, so its fundamental class dies in the target.  See
\cref{cusp-cone-angle-forgotten-sphere-carriers,%
cusp-cone-moving-S4-null-homotopy}.
Removing the smash basepoint gives exact homeomorphisms
\[
 (Q_T\setminus K_T^{(1)})^2
 \cong\mathbb S_{12,T}^4\setminus\{*\}
 \cong\mathbb R^4\hookrightarrow S^6,
\qquad
 (C_T\setminus\{0\})^4
 \cong\mathbb S_{24,T}^4\setminus\{*\}
 \cong\mathbb R^4\hookrightarrow S^8.
\]
These are nonclosed topological open cells, not globally Euclidean-isometric
subspaces for a fixed compatible compact-target metric; see
\cref{cusp-cone-literal-R4-open-cells}.
Both sphere-to-sphere quotient maps induce zero on every positive-degree
integral homology group, and both full maps are based null-homotopic:
\[
 [\widehat q_3^{\wedge2}]=0\in\pi_{12}(S^6),
 \qquad
 [\widehat q_\gamma^{\wedge4}]=0\in\pi_{24}(S^8).
\]
The one-copy map $\widehat q_3$ is now null-homotopic by
Theorem~\ref{h24:thm:cusp-g-zero}; the one-copy class of
$\widehat q_\gamma$ remains open. These full-map statements are
distinct from the two explicit null-homotopies of the moving four-sphere
inclusions.

For completeness, the radial target-slice embeddings used below are typed as
follows.  Put \(r_T=e^{-2\pi T}\) and
\[
 \vartheta_T:\mathbb R/\mathbb Z\xrightarrow{\ \cong\ }C_T,
 \qquad[a]\longmapsto[\ell_Ta].
\]
Then
\[
 j_T:Q_T^2\hookrightarrow\mathcal A^{\wedge2},\qquad
 j_T(y_1,y_2)=
 \bigl[[r_T,1,\iota_T^{-1}(y_1)],
       [r_T,1,\iota_T^{-1}(y_2)]\bigr]_{\wedge},
\]
\[
 k_T^{\mathcal B}:C_T^4\hookrightarrow\mathcal B^{\wedge4},\qquad
 k_T^{\mathcal B}(c_1,\ldots,c_4)=
 \bigl[[r_T,1,\vartheta_T^{-1}(c_1)],\ldots,
       [r_T,1,\vartheta_T^{-1}(c_4)]\bigr]_{\wedge},
\]
and, in the cusp-cone piece of \(\mathcal C\),
\[
 k_T^{\mathcal C}:C_T^4\hookrightarrow\mathcal C^{\wedge4},\qquad
 k_T^{\mathcal C}(c_1,\ldots,c_4)=
 \bigwedge_{j=1}^4[r_T,(1,\vartheta_T^{-1}(c_j))].
\]
These are the embeddings of
\cref{cusp-cone-target-slice-embeddings}.

The intrinsic decompactification can be made into an actual boundary family.
Fix \(T_0\) sufficiently large, put \(r_0=e^{-2\pi T_0}\), and define the
compact radial carriers
\[
 \mathcal D_{12}
 =\frac{[0,r_0]\times(\overline T^2)^2}
 {\{0\}\times(\overline T^2)^2\sim d_{12}},
 \qquad
 \mathcal D_{24}
 =\frac{[0,r_0]\times(S^1)^4}
 {\{0\}\times(S^1)^4\sim d_{24}}.
\]
Both are homeomorphic to \(\operatorname{Cone}(\mathbb T^4)\).  At either
cone vertex \(d_\nu\),
\[
 H_k(\mathcal D_\nu,\mathcal D_\nu\setminus\{d_\nu\};\mathbb Z)
 \cong
 \begin{cases}
  \mathbb Z^4,&k=2,\\
  \mathbb Z^6,&k=3,\\
  \mathbb Z^4,&k=4,\\
  \mathbb Z,&k=5,\\
  0,&\text{otherwise},
 \end{cases}
\]
so \(d_\nu\) is not a topological-manifold point, with or without boundary,
of any dimension.  Set
\[
 A_T^{(12)}=\operatorname{diag}(Y_T,Y_T),\qquad
 A_T^{(24)}=\ell_T I_4,\qquad s_0=T_0^{-1},
\]
and, for \(\nu\in\{12,24\}\), quotient
\([0,s_0]\times\mathbb R^4\) fibrewise by
\[
 (0,x)\sim_\nu(0,y)\iff x=y,\qquad
 (s,x)\sim_\nu(s,y)\iff
 x-y\in A_{1/s}^{(\nu)}\mathbb Z^4\quad(s>0).
\]
The resulting \(\mathfrak E_\nu\) is a Hausdorff, second-countable topological
five-manifold with boundary, with
\[
 p_{12}^{-1}(s)\cong Q_{1/s}^2,\qquad
 p_{24}^{-1}(s)\cong C_{1/s}^4\quad(s>0),\qquad
 p_\nu^{-1}(0)\cong\mathbb R^4.
\]
Its other boundary component, over \(s=s_0\), is a compact four-torus.
Every fixed Euclidean ball embeds isometrically in all sufficiently small
positive fibres.  There is a continuous surjective comparison map
\(\beta_\nu:\mathfrak E_\nu\to\mathcal D_\nu\) which is a homeomorphism away
from the zero fibre, identifies each positive fibre with its radial torus
slice, and collapses the entire zero-boundary component
\(p_\nu^{-1}(0)\cong\mathbb R^4\) to the vertex of the
compact
\(\mathcal D_\nu\cong\operatorname{Cone}(\mathbb T^4)\).  Thus the Euclidean
space is an exact decompactification boundary, not the whole compact source or
target and not a continuous quotient target of a compact source.  The map
\(\beta_\nu\) is neither proper nor a quotient map, so this is not a
topological blow-down or a quotient identification
\(\mathfrak E_\nu/\mathbb R^4\cong\mathcal D_\nu\).  Explicitly, for
\(a=(1/2,1/2,1/2,1/2)\) and \(s_j\downarrow0\), the points
\[
 z_j=[s_j,A_{1/s_j}^{(\nu)}a]
\]
form a closed subset of \(\mathfrak E_\nu\).  Indeed, in the \(S^{12}\) case
\[
 \inf_{n\in\mathbb Z^4}\lVert A_{1/s_j}^{(12)}(a+n)\rVert
 \ge(1/s_j-C_{\mathrm{sys}})\operatorname{dist}(a,\mathbb Z^4)\to\infty,
\]
and in the \(S^{24}\) case the same infimum equals
\(\ell_{1/s_j}\operatorname{dist}(a,\mathbb Z^4)\to\infty\); hence the
sequence has no accumulation in a positive fibre and eventually misses every
half-space chart about every zero-fibre point.  On the other hand,
\(\beta_\nu(z_j)=[e^{-2\pi/s_j},a]\to d_\nu\).  Hence the image set is not
closed although its inverse image is closed, which proves that
\(\beta_\nu\) is not quotient;
see \cref{cusp-cone-exact-decompactification-space,%
cusp-cone-sharp-global-R4-scope}.

Unrescaled ambient metrics impose a different conclusion.  Put
\[
 Z_{12}:=\mathcal A^{\wedge2},\qquad
 Z_{24}^{\mathcal B}:=\mathcal B^{\wedge4},\qquad
 Z_{24}^{\mathcal C}:=\mathcal C^{\wedge4}.
\]
For every fixed compatible metric on these compact targets,
\[
 \operatorname{diam}_{\mathrm{amb}}j_T(Q_T^2)\to0,
 \qquad
 \operatorname{diam}_{\mathrm{amb}}k_T^{\mathcal B}(C_T^4)\to0,
 \qquad
 \operatorname{diam}_{\mathrm{amb}}k_T^{\mathcal C}(C_T^4)\to0,
\]
whereas intrinsically
\[
 \operatorname{diam}(Q_T^2)\ge\tfrac12(T-C_{\mathrm{sys}})\to\infty,
 \qquad \operatorname{diam}(C_T^4)\ge\tfrac12\ell_T\to\infty.
\]
There are fixed cusp-adapted compatible length metrics with regular-stratum
metric tensors \(g_{Z_{12}}\) and \(g_{Z_{24}^{\mathcal Q}}\),
\(\mathcal Q\in\{\mathcal B,\mathcal C\}\), for which the exact rescaled first
fundamental forms are
\[
 \left(\frac{T}{r_T}\right)^2j_T^*g_{Z_{12}}
 =g_{Q_T}\oplus g_{Q_T},
 \qquad
 \left(\frac{\ell_T}{r_T}\right)^2
 (k_T^{\mathcal Q})^*g_{Z_{24}^{\mathcal Q}}
 =g_{C_T}^{\oplus4},
 \quad\mathcal Q\in\{\mathcal B,\mathcal C\}.
\]
With their induced Riemannian path metrics (not the restricted ambient
distance functions), the rescaled slices converge to \(\mathbb R^4\), but the corresponding full
ambient pointed blow-ups converge on their eventual regular-stratum balls to
\(\mathbb R^8\) and \(\mathbb R^{12}\), respectively.  In every rescaling
that unrolls a slice to \(\mathbb R^4\), the cone/fat-wedge basepoint recedes
to infinite scaled distance; vertex-based conic tangents have dimensions
\(8\) and \(12\) for exact conic target metrics, and unit-volume normalization
gives a flat four-torus, not
\(\mathbb R^4\).  See
\cref{sphere-metric-uniform-collapse,sphere-metric-fixed-realization,%
sphere-metric-scaling-trichotomy,sphere-metric-unit-volume,%
sphere-metric-tangent-scope}.

Thus coning changes the codomain and does not contradict the proved
nonextension of the literal character map to the fixed target
\(B\times S^1\).  Nor do Cartesian powers give spheres:
\[
 X^2\simeq S^6\times S^6\not\simeq S^{12},\qquad
 X^4\simeq(S^6)^4\not\simeq S^{24};
\]
only the based smash powers are spheres.  There is no full-base two-copy
four-torus quotient, no fixed \(T^4\) target retaining the four character
classes, and no fixed \(S^4\) target extending a degree-one collapse of the
radial quotient.  No nontrivial finite base cover preserves
\(K_\Lambda\), and the local unspooling cover is not a nontrivial connected
unbranched cover of the simply connected compact sources.  Finally,
\(S^{12}\) and \(S^{24}\) admit no almost-complex structure, so the smash
sources are not compact complex spheres.  See
\cref{cusp-cone-S12-route,cusp-cone-S24-local-route,cusp-cone-S24-route,%
cusp-cone-metric-scope,cusp-cone-obstructions,%
cusp-cone-explicit-carriers,cusp-cone-target-homotopy-types}.
Every construction in this sphere-route item uses the native quotients of the
original, unscaled cusp fibres; none uses the ancillary finite-cover
unspooling from the preceding item.

\item \Verified\ \textbf{The repaired compact-link construction and its
finite Gaussian quantization have exact actions, spectra, and a resolved
zero-gap diagonal sequence.}
Keep the cover index and the mesh index independent:
\[
 M_{\mathrm{cov}}\in\mathbb Z_{\ge1},\qquad
 S_{M_{\mathrm{cov}}}=\operatorname{diag}(1,1,M_{\mathrm{cov}},M_{\mathrm{cov}}),
 \qquad P_{T,M_{\mathrm{cov}}}=P_TS_{M_{\mathrm{cov}}},
 \qquad N_{\mathrm{mesh}}\in\mathbb Z_{\ge2}.
\]
Then
\[
 \widetilde F_{T,M_{\mathrm{cov}}}
 =\mathbb R^4/P_{T,M_{\mathrm{cov}}}\mathbb Z^4
 \longrightarrow F_T
\]
has degree \(M_{\mathrm{cov}}^2\), Gram matrix and covolume
\[
 G_{T,M_{\mathrm{cov}}}
 =P_{T,M_{\mathrm{cov}}}^{\mathsf T}P_{T,M_{\mathrm{cov}}},
 \qquad
 \mathcal D_{T,M_{\mathrm{cov}}}
 =\sqrt{\det G_{T,M_{\mathrm{cov}}}}
 =M_{\mathrm{cov}}^2D_T.
\]
Put
\[
 V_{N_{\mathrm{mesh}}}
 =(\mathbb Z/N_{\mathrm{mesh}}\mathbb Z)^4,
 \qquad
 E_{N_{\mathrm{mesh}}}
 =V_{N_{\mathrm{mesh}}}\times\{1,2,3,4\}.
\]
Transport of the standard
\(N_{\mathrm{mesh}}\)-cubulation by \([y]\mapsto[P_{T,M_{\mathrm{cov}}}y]\)
is an honest cubulation of the whole cover torus; its vertex set is exactly
the \(N_{\mathrm{mesh}}\)-torsion subgroup, via
\[
 [n]\longmapsto
 \left[P_{T,M_{\mathrm{cov}}}n/N_{\mathrm{mesh}}\right].
\]
Parallel transport along its oriented edges gives the typed gauge-equivariant
holonomy map to \(\mathrm{SU}(2)^{E_{N_{\mathrm{mesh}}}}\).

Write \(M=M_{\mathrm{cov}}\) and \(N=N_{\mathrm{mesh}}\) only within the
following compact-link formulas.  Normalize
\[
 \langle X,Y\rangle_{\mathfrak{su}(2)}
 =-\tfrac12\operatorname{tr}(XY),
 \qquad H_{\mathrm{link}}=\operatorname{diag}(i,-i),
 \qquad \lVert H_{\mathrm{link}}\rVert^2=1,
 \qquad \exp(2\pi H_{\mathrm{link}})=I_2.
\]
The periodic
inverse-transport links
\[
 \begin{aligned}
 U_1(n)&=\exp(-2\pi n_2H_{\mathrm{link}}/N^2),\\
 U_2(n)&=
 \begin{cases}I_2,&0\le n_2\le N-2,\\
 \exp(2\pi n_1H_{\mathrm{link}}/N),&n_2=N-1,\end{cases}\\
 U_3(n)&=U_4(n)=I_2.
 \end{aligned}
\]
have, including every wrap plaquette,
\[
 U_{12}(n)=\exp(2\pi H_{\mathrm{link}}/N^2),\qquad U_{ij}(n)=I_2
 \quad((i,j)\ne(1,2),\ i<j),
\]
The principal-logarithm action and the
explicit continuum-matched \((1,2)\)-Wilson sector are exactly
\[
 S^{\log}_{T,M,N}(U)=\frac{2\pi^2M^2}{g_{\mathrm{YM}}^2D_T},
\]
\[
 S^{W,12}_{T,M,N}(U)
 =\frac{M^2N^4}{2g_{\mathrm{YM}}^2D_T}
   \left(2-2\cos\frac{2\pi}{N^2}\right)
 =\frac{M^2}{g_{\mathrm{YM}}^2D_T}
  \left(2\pi^2-\frac{2\pi^4}{3N^4}
  +\frac{4\pi^6}{45N^8}+O(N^{-12})\right).
\]
The inverse-Gram two-form minor producing these formulas is exactly
\[
 (G_{T,M_{\mathrm{cov}}}^{-1})_{11}
 (G_{T,M_{\mathrm{cov}}}^{-1})_{22}
 -(G_{T,M_{\mathrm{cov}}}^{-1})_{12}^2=D_T^{-2}.
\]

For \(0\le p\le4\), define the real cochain spaces
\[
 C_N^p
 =\operatorname{Map}\!\left(
 V_{N_{\mathrm{mesh}}},
 \bigwedge\nolimits^p(\mathbb R^4)^*\otimes\mathfrak{su}(2)\right).
\]
For \(\vartheta^j_{T,M}=P_{T,M}^{-\mathsf T}\mathbf e_j\), let
\(\mathsf T_j a(n)=a(n+\mathbf e_j)\) and set
\[
 D^+_{N,j}=N(\mathsf T_j-I),
 \qquad D^-_{N,j}a(n)=N(a(n)-a(n-\mathbf e_j)),
\]
\[
 d_Na=\sum_{j=1}^4
 \varepsilon(\vartheta^j_{T,M})D^+_{N,j}a,
 \qquad
 d_N^*=-\sum_{j=1}^4
 \iota_{(\vartheta^j_{T,M})^\sharp}D^-_{N,j},
\]
and
\[
 \Delta_N^{(1)}=d_Nd_N^*+d_N^*d_N,
 \qquad
 \mathscr C_{T,M,N}
 =\ker(d_N^*:C_N^1\to C_N^0)
 \cap(\ker\Delta_N^{(1)})^\perp.
\]
This last orthogonal complement removes the zero Fourier harmonic cochains.
On this real reduced Coulomb space put
\[
 \mathcal L_{T,M,N}=g_{\mathrm{YM}}^{-2}d_N^*d_N,
 \qquad \mathfrak d_N(k)_j=N(e^{2\pi ik_j/N}-1).
\]
For every \(k\ne0\), its eigenvalue and complex transverse multiplicity are
\[
 \lambda_{T,M,N}(k)=\frac1{g_{\mathrm{YM}}^2}
 \left\lVert P_T^{-\mathsf T}S_{M_{\mathrm{cov}}}^{-1}
 \mathfrak d_N(k)\right\rVert^2,
 \qquad 3\dim\mathfrak{su}(2)=9.
\]
In particular,
\[
 \lambda_{T,M,N}(\mathbf e_1)
 =\frac{4N^2\sin^2(\pi/N)}{g_{\mathrm{YM}}^2}
  \frac{L_T^2+m_T^2}{D_T^2},
\]
independently of \(M\).  The Friedrichs realization on
\(L^2(\mathscr C_{T,M,N},da)\) of
\[
 \widehat H^{\mathrm{Gauss}}_{T,M,N}
 =-\tfrac12\Delta_{\mathscr C_{T,M,N}}
  +\tfrac12\langle a,\mathcal L_{T,M,N}a\rangle
\]
has a unique Gaussian ground state and exact gap
\[
 \operatorname{gap}(\widehat H^{\mathrm{Gauss}}_{T,M,N})
 =\sqrt{\lambda_{\min}(\mathcal L_{T,M,N})}
 \le\frac{2N\sin(\pi/N)}{g_{\mathrm{YM}}}
       \frac{\sqrt{L_T^2+m_T^2}}{D_T}.
\]
For
\[
 T_j=j^2,\qquad M_j=j,\qquad N_j=j^3,
\]
one has simultaneously
\[
 \operatorname{sys}(\widetilde F_{T_j,M_j})\to\infty,
 \qquad
 \max_r\frac{\lVert P_{T_j,M_j}\mathbf e_r\rVert}{N_j}\to0,
 \qquad
 \operatorname{gap}(\widehat H^{\mathrm{Gauss}}_{T_j,M_j,N_j})\to0,
 \qquad
 S^{\log}_{T_j,M_j,N_j}(U)\to0.
\]
The exponential map \(a\mapsto(\exp(a_e/N))_e\) is only a local
diffeomorphism at zero, has derivative \(N^{-1}\operatorname{id}\), is not
globally injective, and, if the compact-link action satisfies
\(DS_W(\mathbf1)=0\), transports its Hessian with the exact factor
\[
 D^2(S_W\circ\operatorname{Exp}_N)_0
 =N^{-2}D^2(S_W)_{\mathbf1}.
\]
Consequently the displayed Gaussian Hilbert space is
not the Haar \(L^2\)-space on compact links or its nonlinear gauge-invariant
subspace; no unitary identification, nonlinear compact-link Hamiltonian,
Gauss-law spectral problem, transfer matrix, renormalized continuum limit,
or Osterwalder--Schrader reconstruction has been constructed.  This theorem
therefore remains strictly on the auxiliary finite Gaussian side of the
physical Yang--Mills boundary.  See
 \cref{cusp-link-cubical-realization,cusp-link-holonomy-map,%
 cusp-link-flux-theorem,cusp-link-finite-spectrum,cusp-link-gaussian-theorem,%
 cusp-link-exponential-local-only,cusp-link-physical-hilbert-space,%
 cusp-link-dirac-not-l2}.

\item \Verified\ \textbf{The exceptional first-coordinate elliptic quotient
is finite and does not approach the cusp.}
For the marked complex torus
\[
 \mathcal A_z=\mathbb C^2/
 \mathbb Z\langle(6\mu,\beta),(\tau,\mu),(1,0),(0,1)\rangle,
\]
the first coordinate descends onto
\(E_\tau=\mathbb C/(\mathbb Z\tau+\mathbb Z)\) exactly on
\[
 6\mu=a\tau+b,\qquad(a,b)\in\mathbb Z^2.
\]
The kernel is connected and equals
\(\mathbb C/(\mathbb Z+\mathbb Z(a\mu-\beta))\), with
\(\operatorname{Im}(a\mu-\beta)=-D_{\mathrm{can}}>0\).  Under the two
generators the labels transform by
\[
 (a,b)\xmapsto{g_0}(a,a+b),\qquad
 (a,b)\xmapsto{g_2}(-b,6+a).
\]
The marked union is therefore monodromy invariant; its image in the completed
orbifold base is finite and avoids a punctured cusp neighborhood.  Thus this
specific elliptic quotient cannot realize the cusp gap collapse.  It remains
logically possible for a fibre to have a different elliptic quotient or an
integral magnetic class not induced by this projection.  See
\cref{cusp-spec-first-coordinate-descent,%
cusp-spec-elliptic-descent-finiteness,%
cusp-spec-elliptic-descent-scope}.

\item \Verified\ \textbf{A single specified auxiliary operator has
zero essential spectral threshold.}
Fourier expansion identifies
\[
 \mathscr H_{\mathrm{aux}}
 =L^2\bigl([T_0,\infty),dT;
 L^2_0(\mathbb R^4/\mathbb Z^4,dy)\bigr)
 \cong\bigoplus_{k\in\mathbb Z^4\setminus\{0\}}
 L^2([T_0,\infty),dT).
\]
With
\[
 \lambda_k(T)=4\pi^2\lvert P_T^{-\mathsf T}k\rvert^2,\qquad
 (\mathscr L_{\mathrm{aux}}u)_k(T)=\lambda_k(T)u_k(T),
\]
the exact domain is
\[
 D(\mathscr L_{\mathrm{aux}})=\left\{u:
 \sum_{k\ne0}\int_{T_0}^{\infty}
 \lambda_k(T)^2\lvert u_k(T)\rvert^2\,dT<\infty\right\}.
\]
The operator is nonnegative and self-adjoint, has zero kernel, and the
\(k=(1,0,0,0)\) singular sequence supported on \([n,n+1]\) proves
\[
 0\in\sigma_{\mathrm{ess}}(\mathscr L_{\mathrm{aux}}).
\]
The word ``gapless'' here belongs to this auxiliary Hilbert space and
operator; it is not an invariant of the complex threefold alone.  See
\cref{cusp-interpretation-direct-integral-gapless,%
cusp-interpretation-fixed-complex-counterexample}.

\item \Verified\ \textbf{The continuum and finite classical gauge
fluctuation operators have exact spectra and cusp gaps.}
For a compact Lie group with \(d_G=\dim_{\mathbb R}\mathfrak g\), the
product-connection Hessian on the reduced Coulomb space is
\[
 \mathcal L_T=g_{\mathrm{YM}}^{-2}\mathrm d^*\mathrm d,\qquad
 \lambda_T^{\mathrm{YM}}(k)
 =\frac{4\pi^2}{g_{\mathrm{YM}}^2}
  \lvert P_T^{-\mathsf T}k\rvert^2
\]
on every \(k\ne0\) transverse Fourier fibre, with complex multiplicity
\(3d_G\).  Hence
\[
 \gamma_T=\frac{4\pi^2}{g_{\mathrm{YM}}^2}
 \min\left\{\frac{L_T^2+m_T^2}{D_T^2},
                  \frac{q_T^2+36m_T^2}{D_T^2}\right\}
 =\frac{4\pi^2}{g_{\mathrm{YM}}^2T^2}
 \left(1+\frac{-2H+2\min\{B,0\}}T+O(T^{-2})\right).
\]
For the finite cubical complex,
\[
 \mathfrak d_N(k)_j=N(e^{2\pi i k_j/N}-1),\qquad
 \zeta_{N,T}(k)=P_T^{-\mathsf T}\mathfrak d_N(k),\qquad
 \mathcal L_{N,T}=g_{\mathrm{YM}}^{-2}\mathrm d_N^*\mathrm d_N,
\]
and the reduced \(k\ne0\) eigenvalue is exactly
\[
 \lambda_{N,T}^{\mathrm{YM}}(k)
 =g_{\mathrm{YM}}^{-2}\lvert P_T^{-\mathsf T}\mathfrak d_N(k)\rvert^2
\]
with the same complex multiplicity.  Writing
\(\varrho_N^{\mathrm{spec}}=4N^2\sin^2(\pi/N)\), there is a \(T_1\), independent of
\(N\ge2\), such that for every \(T\ge T_1\) its gap is
\[
 \gamma_{N,T}=\frac{\varrho_N^{\mathrm{spec}}}{g_{\mathrm{YM}}^2}
 \min\left\{\frac{L_T^2+m_T^2}{D_T^2},
                  \frac{q_T^2+36m_T^2}{D_T^2}\right\}
 =\frac{\varrho_N^{\mathrm{spec}}}{g_{\mathrm{YM}}^2T^2}
 \left(1+\frac{-2H+2\min\{B,0\}}T+O(T^{-2})\right)
\]
uniformly for \(N\ge2\).  Thus each fixed finite model is positively
gapped, but \(\gamma_{N(T),T}\to0\) for every \(N(T)\ge2\); resolving
the varying continuum geometry is the stronger condition
\(N(T)/T\to\infty\).  The finite operator is a linearized
Lie-algebra-valued cochain model, not a nonlinear compact-link Wilson
action.  See
 \cref{cusp-gauge-finite-spectrum,cusp-gauge-finite-gap-limits,%
 cusp-gauge-continuum-symbol,cusp-gauge-continuum-gap}.

\item \Verified\ \textbf{Every fixed nondegenerate integral flux on the
marked four-torus has an exact magnetic-oscillator spectrum whose first
spacing tends to zero along the cusp.}
In angular coordinates \(u\in\mathbb R^4/\mathbb Z^4\), let
\[
 G_T=P_T^{\mathsf T}P_T,
 \qquad \sqrt{\det G_T}=D_T,
 \qquad
 \mathsf E=\sum_{j<k}E_{jk}\,\mathrm du_j\wedge\mathrm du_k
 \in\bigwedge^2(\mathbb Z^4)^*.
\]
The factor of automorphy and connection
\[
 J_n^{\mathsf E}(u)
 =\exp\!\left(-2\pi i\sum_{j<k}E_{jk}n_ju_k\right),
 \qquad
 \nabla^{\mathsf E}
 =\mathrm d+2\pi i\sum_{j<k}E_{jk}u_j\,\mathrm du_k
\]
define a Hermitian line bundle \(\mathscr L_{\mathsf E,T}\) with curvature
\(F_{\nabla^{\mathsf E}}=2\pi i\mathsf E\).  The nonnegative closed form
\[
 q_{\mathsf E,T}[f]=\int_{[0,1]^4}
 (G_T^{-1})^{jk}\overline{\nabla_j^{\mathsf E}f}
 \nabla_k^{\mathsf E}f\,D_T\,\mathrm d^4u
\]
has form domain \(H^1(F_T;\mathscr L_{\mathsf E,T})\); its Friedrichs
operator \(\Delta_{\mathsf E,T}=(\nabla^{\mathsf E})^*\nabla^{\mathsf E}\)
has domain \(H^2(F_T;\mathscr L_{\mathsf E,T})\) and compact resolvent.
For nondegenerate \(\mathsf E\), define \(b_1,b_2>0\) by
\[
 \operatorname{Spec}\!\left(
 G_T^{-1/2}(2\pi\mathsf E)G_T^{-1/2}\right)
 =\{ib_1,-ib_1,ib_2,-ib_2\}.
\]
The ladder commutators give the complete set of spectral values
\[
 \operatorname{Spec}(\Delta_{\mathsf E,T})
 =\{(2n_1+1)b_1+(2n_2+1)b_2:
 n_1,n_2\in\mathbb Z_{\ge0}\},
 \qquad
 \operatorname{gap}(\Delta_{\mathsf E,T})=2\min\{b_1,b_2\}.
\]
Writing
\[
 d=|\operatorname{Pf}(\mathsf E)|
 =|E_{12}E_{34}-E_{13}E_{24}+E_{14}E_{23}|>0,
\]
the determinant identity is
\[
 b_1b_2=\frac{4\pi^2d}{D_T},
\]
and therefore
\[
 \operatorname{gap}(\Delta_{\mathsf E,T})
 \le4\pi\sqrt{\frac d{D_T}}
 =\frac{4\pi\sqrt d}{T}
 \left(1+\frac{B-2H}{2T}+O(T^{-2})\right)\longrightarrow0.
\]
This is a fibrewise scalar magnetic Laplacian on the marked cusp tori.  It
does not assert that \([\mathsf E]\) is fixed by full monodromy or identify
this level spacing with a physical Yang--Mills particle mass.  See
\cref{cusp-magnetic-oscillator-theorem,cusp-magnetic-frequency-product}.

\item \Verified\ \textbf{A reducible non-flat \(SU(2)\) critical
connection has exact action and an exact negative Jacobi eigenvalue.}
The rank-two sum
\(\mathscr E_T=\mathscr L_T\oplus\mathscr L_T^{-1}\) has parallel trivial
determinant, \(c_1=c_2=0\), and an explicitly nullhomotopic
\(SU(2)\) clutching map.  Thus its principal \(SU(2)\)-bundle is
smoothly trivial.  In the lifted line-bundle splitting
\(\mathscr L_T\oplus\mathscr L_T^{-1}\), the direct-sum connection and its
curvature are
\[
 H_{\mathrm{mag}}=\begin{pmatrix}i&0\\0&-i\end{pmatrix},\qquad
 \nabla^T=\mathrm d+2\pi u_1\,\mathrm du_2\,H_{\mathrm{mag}},\qquad
 F_{\nabla^T}=\frac{2\pi}{D_T}H_{\mathrm{mag}}\,
 \mathrm dy_1\wedge\mathrm dy_2,
\]
where the Lie-algebra norm is the same one fixed above,
\(\langle X,Y\rangle=-\tfrac12\operatorname{tr}(XY)\), so that
\(\lVert H_{\mathrm{mag}}\rVert^2=1\).
In any chosen global \(SU(2)\)-trivialization of the smoothly trivial
principal bundle this matrix is replaced by its corresponding gauge
transform.  One has \(\mathrm d_{\nabla^T}^*F_{\nabla^T}=0\) and
\[
 \operatorname{YM}_T(\nabla)
 =\frac1{2g_{\mathrm{YM}}^2}
  \int_{F_T}|F_\nabla|^2\,\mathrm d\operatorname{vol}_{F_T},
\]
so
\[
 \operatorname{YM}_T(\nabla^T)
 =\frac{2\pi^2}{g_{\mathrm{YM}}^2D_T}
 =\frac{2\pi^2}{g_{\mathrm{YM}}^2T^2}
 \left(1+\frac{B-2H}{T}+O(T^{-2})\right).
\]
Let \((e_\mu)_{\mu=1}^4\) be a constant orthonormal frame for the descended
flat metric.  The background-gauge Jacobi expression
\[
 (\mathcal J_Ta)_\nu
 =-\sum_{\mu=1}^4\nabla_{e_\mu}\nabla_{e_\mu}a_\nu
  -2\sum_{\mu=1}^4[F_{\nu\mu},a_\mu]
\]
has an explicit smooth real unit eigenvector satisfying
\[
 \mathrm d_{\nabla^T}^*a=0,\qquad
 \mathcal J_Ta=-\frac{4\pi}{D_T}a,\qquad
 \operatorname{YM}_T''|_{\nabla^T}(a,a)
 =-\frac{4\pi}{g_{\mathrm{YM}}^2D_T}.
\]
This proves a classical saddle direction at every finite \(T\), not a
spectral assertion for a quantum Hamiltonian.  See
\cref{cusp-magnetic-su2-triviality,cusp-magnetic-ym-critical-point,%
cusp-magnetic-negative-mode}.

\item \Verified\ \textbf{The native full-flux background extends root by
root from split \(SU(2)\) to split-Cartan \(SU(r)\), with exact topology,
action, spectrum, index, and wall nullities.}
For every nonzero trace-zero \(\mathbf m\in\mathbb Z^r\), put
\[
 \mathscr E_{\mathbf m,T}
 =\bigoplus_i\mathscr L_{\eta,T}^{\otimes m_i},
 \qquad s_2(\mathbf m)=\sum_i m_i^2,
 \qquad k_{\mathbf m}=6s_2(\mathbf m).
\]
Then
\[
 c_2(\mathscr E_{\mathbf m,T})
 =-\frac{s_2(\mathbf m)}2[\eta_T]^2,
 \quad
 \langle c_2,[F_T]_{\mathrm{mark}}\rangle=-6s_2(\mathbf m),
 \quad
 \langle c_2,[F_T]_{x\text{-}\mathrm{complex}}\rangle=k_{\mathbf m},
\]
and the corresponding reducible Yang--Mills connection satisfies
\[
 \operatorname{YM}_T
 =\frac{2\pi^2k_{\mathbf m}}{g_{\mathrm{YM}}^2}
  +\frac{\pi^2s_2(\mathbf m)S_T^2}{g_{\mathrm{YM}}^2D_T}.
\]
For each positive root \(i<j\), the charge is
\(q_{ij}=|m_i-m_j|\); the complete negative normalized Coulomb bands are
\[
 q_{ij}\bigl((2n+1)\mathcal B_{T,\min}-\mathcal B_{T,\max}\bigr),
 \qquad 0\le n<N_-(T),
\]
each contributing real dimension \(12q_{ij}^2\), with coincident bands
summed.  Hence
\[
 \operatorname{ind}_{\mathbb R}\mathcal J_{\mathbf m,T}
 =12rN_-(T)s_2(\mathbf m)
 =36r s_2(\mathbf m)
\]
on the deep cusp tail.  At \(S_T=0\) the negative spectrum is empty but the
root zero modes double, so no strict-stability conclusion is made.  These are
classical fibrewise theorems, not an interacting quantum spectrum; moreover
their nonzero \(c_2\) obstructs extension over every proposed sphere carrier
with vanishing fourth integral cohomology.  See
\cref{cusp-native-full-flux-suN-topology,%
cusp-native-full-flux-suN-action-theorem,%
cusp-native-full-flux-suN-Morse-theorem,%
cusp-native-full-flux-suN-nullity-theorem,%
cusp-native-full-flux-suN-stability-boundary}.

\item \Verified\ \textbf{The Koide cone is isomorphic only to an
ambient period-coefficient discriminant, not to the regular cusp family.}
For
\[
 Q_K(r)=\frac23(r_1+r_2+r_3)^2-(r_1^2+r_2^2+r_3^2),\qquad
 \begin{pmatrix}h\\x\\y\end{pmatrix}
 =\begin{pmatrix}
 1/\sqrt3&1/\sqrt3&1/\sqrt3\\
 1/\sqrt2&0&-1/\sqrt2\\
 1/\sqrt6&-2/\sqrt6&1/\sqrt6
 \end{pmatrix}
 \begin{pmatrix}r_1\\r_2\\r_3\end{pmatrix},
\]
the coordinate change is orthogonal and
\[
 Q_K(r)=h^2-x^2-y^2.
\]
The real-linear map
\[
 (h,x,y)\longmapsto(t,m,\varrho)
 =\left(h+y,\frac{x}{\sqrt6},h-y\right)
\]
is an ambient coefficient-space isomorphism.  For the canonical convention
\[
 Z=\begin{pmatrix}6\mu&\tau\\\beta&\mu\end{pmatrix},\qquad
 P=\begin{pmatrix}
 6\operatorname{Re}\mu&\operatorname{Re}\tau&1&0\\
 6m&t&0&0\\
 \operatorname{Re}\beta&\operatorname{Re}\mu&0&1\\
 \varrho&m&0&0
 \end{pmatrix},
\]
one has, independently of the real parts,
\[
 \det P=t\varrho-6m^2
 =\det\begin{pmatrix}t&\sqrt6\,m\\\sqrt6\,m&\varrho\end{pmatrix}
 =Q_K(r).
\]
For \(r\in\mathbb R_{\ge0}^3\setminus\{0\}\), the equation
\(Q_K(r)=0\) is equivalent to the displayed symmetric matrix being
positive semidefinite of rank one.  It also forces \(\det P=0\), so
the columns of \(P\) do not form a full lattice and define no regular
compact complex two-torus.  Actual radial cusp fibres instead satisfy
\[
 \det P_T=-D_T<0,\qquad
 D_T=T^2+(2H-B)T+O(1),\qquad
 \det P_T\sim-T^2.
\]
See
\cref{koide-s6-thm-ambient-discriminant,koide-s6-cor-ambient-only,%
koide-s6-cor-actual-cusp}.

\item \Verified\ \textbf{The positive defect-one Koide shell has an exact
ideal-and-unit parametrization and an unconditional power-saving counting
law.}
Define
\[
 \Delta(x,y,z)=4(xy+xz+yz)-(x^2+y^2+z^2),
\]
and put
\[
 \mathcal R=\{(x,y,z)\in\mathbb Z_{>0}^3:
 x\le y\le z,\ \Delta(x,y,z)=-1,\ z\ge4(x+y)\},
\]
\[
 R_{\mathrm K}(Y)
 =\#\{(x,y,z)\in\mathcal R:z\le Y\}.
\]
The terminal Vieta wall is exact: below it,
the other root \(4(x+y)-z\) is positive and smaller than \(y\).  On the wall,
the coordinates
\[
n=z-4(x+y),\qquad q=y-x,\qquad a=x+y-4n
\]
give a bijection with the integral chamber
\[
a^2-3q^2-18n^2=-2,\qquad
a>0, q,n\ge0, a+4n>q,\qquad z=4a+17n,
\]
whose inverse is
\[
 (x,y,z)=
 \left(\frac{a+4n-q}{2},\frac{a+4n+q}{2},4a+17n\right).
\]
For
\[
 F_{\mathrm K}(a,q,n)=a^2-3q^2-18n^2,\qquad
 \mathscr X=\{F_{\mathrm K}=-2\},
\]
direct
Hensel counting and the good-prime formula give
\[
\sigma_2=\frac32,\qquad \sigma_3=2,\qquad
 \sigma_p=\frac{1-p^{-2}}
 {1-\left(\frac3p\right)p^{-1}}\qquad(p\ge5),
\]
and hence
\[
 \mathfrak S(\mathscr X)
 =\frac92\frac{\log(2+\sqrt3)}{\sqrt3\,\zeta(2)}.
\]
The exact Poincar\'e-residue integral over the reduced chamber is
\[
 I_\infty=\frac1{2\sqrt3}\log\!\left(\frac3{2\sqrt2}\right).
\]
The Gorodnik--Oh--Shah sector theorem applies to the genuine norm
\(\max\{|4a+17n|,|q|,|a|\}\).  An explicit element of
\(\operatorname{SO}(F_{\mathrm K},\mathbb Z_2)\) with spinor norm \(-17\) exchanges the
two Borovoi sign orbits and proves \(c_{\mathscr X}=1\).  Calibrating both
counting theorems against the same Euclidean balls therefore gives
\[
 C_0:=\frac9{2\pi^2}\log(2+\sqrt3)
 \log\!\left(\frac3{2\sqrt2}\right)
\]
and
\[
 \boxed{R_{\mathrm K}(Y)\sim C_0Y.}
\]
For \(n\geq1\), the exact ideal-and-unit parametrization puts
\[
\begin{aligned}
 m_n&=18n^2-2=2(3n-1)(3n+1),\\
 \mathcal P(m_n)&=\bigl\{(C,D,E)\in\mathbb Z_{>0}^2\times\mathbb Z:\\
 &\hspace{5em}C^2D=m_n,\quad 0\leq E<D,\quad
 E^2\equiv3\pmod D\bigr\}.
\end{aligned}
\]
The map
\[
 (C,D,E)\longmapsto
 C\bigl(D\mathbb Z+(E+\sqrt3)\mathbb Z\bigr)
\]
is a bijection to norm-\(m_n\) ideals of \(\mathbb Z[\sqrt3]\).  Each ideal
has exactly one generator \(a_0+q_0\sqrt3\) with \(a_0>2q_0\geq0\), and its
positive solutions are
\[
 a_k+q_k\sqrt3=(a_0+q_0\sqrt3)(2+\sqrt3)^k,\qquad k\geq0.
\]
The inverse descent is
\((a,q)\mapsto(2a-3q,2q-a)\), iterated while \(a\leq2q\), followed by the
unique HNF recovery.  Hence multiplicity is one and
\[
 \#\mathcal P(m_n)=\rho(m_n)
 =\sum_{d\mid m_n}\chi_{12}(d)
 =\rho(3n-1)\rho(3n+1).
\]
The \(n=0\) face is not absorbed into this positive-norm ideal family: it is
exactly
\[
 a_k+q_k\sqrt3=(1+\sqrt3)(2+\sqrt3)^k,\qquad k\geq1,
\]
with the original strict wall, and contributes the explicit floor count
\(R_0(Y)=O(\log(2+Y))\) in \cref{koide-param-eq:R0-exact}.
For every real \(Y\geq0\), the unchanged weak inequality
\(4a+17n\leq Y\) is retained in the exact floor formula of
\cref{koide-param-thm:exact-floor}; its discrepancy splits into the radial
shifted divisor term and the exact phase discrepancies.  If
\(S_\ell(m)\) denotes the norm-\(m\) coefficient of the unitary Hecke
character defined in \cref{koide-param-eq:Hecke-character}, then the angular
data obey the exact shifted identity
\[
 S_\ell(18n^2-2)
 =(-1)^\ell S_\ell(3n-1)S_\ell(3n+1),
\]
with the shifted factorization proved in
\cref{koide-param-eq:shifted-Hecke-factorization}; the endpoint half-weights
and all five sharp walls are retained in
\cref{koide-param-eq:Fejer-exact}.
A direct proof of effective well-roundedness for the unchanged sharp family
allows Benoist--Oh, Theorem~13.2~\cite[pp.~1931--1932]{BenoistOh_2012},
to be applied orbit by orbit and strengthens
the count to
\[
 \boxed{R_{\mathrm K}(Y)=C_0Y+O(Y^{1-\eta})}
 \qquad\text{for some fixed, presently unspecified }\eta>0.
\]
The primary-source comparison gives no exact-domain square-root theorem:
Liu--Sarnak's \(39/64\) exponent is for a smooth standard-radial count on
the exact congruence lattice, and the sharp standard-radial results of Good
and Chatzakos--Petridis have the \(2/3\) frontier after small-spectrum terms.
Neither transfers to the five-wall height \(4a+17n\); see
\cref{koide-lit-determination}.
The exact finite census gives
\(R_{\mathrm K}(10^8)=3{,}536{,}002\), with all rows independently replayed;
the largest \(\lvert R_{\mathrm K}(Y)-C_0Y\rvert/\sqrt Y\) for
\(10^4\leq Y\leq10^8\) is
\(0.1714001928088145597\ldots\) at \(Y=867{,}240\).

There is also an exact hyperbolic bridge.  For
\[
 B=(18,-6)_{\mathbb Q},\qquad
 \mathcal O=\mathbb Z1\oplus
 \mathbb Z\frac{I+J}{2}\oplus
 \mathbb Z\frac{I-J}{2}\oplus
 \mathbb Z\frac{IJ}{2},
\]
the map
\[
 (a,q,n)\longmapsto2+\frac{qI+aJ+nIJ}{2}
\]
is a bijection from the parity-restricted ternary shell to the
reduced-trace-\(4\), reduced-norm-\(1\) locus.  The order has reduced
discriminant \(108\), index \(18\) in a maximal overorder, and is not an
Eichler order; \(B\) is ramified exactly at \(2,3\).  Each image is primitive
hyperbolic with centralizer
\(\{\pm\gamma^k:k\in\mathbb Z\}\), eigenvalues \(2\pm\sqrt3\), and length
\(2\log(2+\sqrt3)\).  The shell is a finite union of conjugacy classes, but
its canonical height \(4a+17n\) differs from the exact radial Huber height
\(7a+15\sqrt2\,n\), since
\[
 X_i(\gamma)^2=1+\frac{(7a+15\sqrt2\,n)^2}{48}.
\]
This theorem is independent of the existence of the proposed complex
structure on \(S^6\).  The stronger square-root discrepancy
\(R_{\mathrm K}(Y)-C_0Y=O_\varepsilon(Y^{1/2+\varepsilon})\) is an explicitly \Open
problem, not a hypothesis used in the asymptotic; see
\cref{koide-count-thm-coordinate-bijection,koide-count-thm-local-densities,%
koide-count-thm-archimedean-factor,koide-count-thm-linear-asymptotic,%
koide-param-thm:integral-parametrization,koide-param-thm:power-saving,%
koide-count-sec-census-1e8,%
koide-quat-exact-reformulation-of-the-koide-chamber-count,%
koide-lit-determination,%
koide-count-open-square-root-discrepancy}.

\item \Verified\ \textbf{The earlier Erd\H{o}s--Straus divisor selector and
the \(S^6\) terminal presentation are the same marked integral matrix slice.}
For \(p\equiv1\pmod{12}\) prime, define
\[
 L=\frac{p-1}{4},\qquad 0\le j<L,\qquad
 R_j=4j+3,\qquad a_j=\frac{p+R_j}{4}.
\]
For a selector state
\[
 u\in\mathbb Z_{>0},\qquad
 u\mid a_j^2,\qquad R_j\mid4u+p^\varepsilon,\qquad
 \varepsilon\in\{0,1\},
\]
set
\[
 T=\frac{4u+p^\varepsilon}{R_j},\qquad
 \ell_0=\frac{T+1}{4},\qquad
 \ell_1=u-jT,\qquad \ell_2=1.
\]
Then \(T\equiv3\pmod4\),
\[
 12\ell_0-4\ell_1-3\ell_2=p^\varepsilon,
\]
and
\[
 \boxed{
 \begin{pmatrix}1&-j\\0&1\end{pmatrix}
 \begin{pmatrix}R_j&-u\\4&-T\end{pmatrix}
 =\begin{pmatrix}3&-\ell_1\\4&1-4\ell_0\end{pmatrix}.}
\]
The inverse \(T=4\ell_0-1\), \(u=\ell_1+jT\) proves bijectivity once the
prime, shell label, channel, orientation, and divisor support are retained.
The \(\varepsilon=0\) channel lies over terminal determinant parameter \(1\),
and the \(\varepsilon=1\) channel lies over \(p\).  The proof also derives the
complete selector from the factorization
\((R_jy-pa_j)(R_jz-pa_j)=p^2a_j^2\).  This is a structural equality of
integer presentations, not a map of the underlying geometric spaces; see
\cref{sec:es-s6-crosswalk}.

\item \Verified\ \textbf{The residue lattice and shell-count maps retain
their different domains.}
The lattice has the exact index-$72$ glue and the full coordinate
splitting in Theorems~\ref{esn:thm:niemeier}
and~\ref{esn:fa:thm:full-quotient}. Its $D_4$ shortest dual-class
sets are the vector, even half-spin and odd half-spin weights of the
flag tangent modules, in that marking
(Theorem~\ref{esn:thm:d4-f4-triality}).
For $p\equiv1\pmod{24}$ prime and $R>0$, $R\equiv3\pmod4$, put
$a=(p+R)/4$, $g=\gcd(R,pa)$, $R_0=R/g$, $N_0=pa/g$. Then
\[
 \left\{d>0:d\mid N_0^2,\ d\equiv-N_0\pmod{R_0}\right\}
 \xrightarrow{\ \cong\ }
 \left\{(b,c)\in\mathbb Z_{>0}^2:
                    \frac4p=\frac1a+\frac1b+\frac1c\right\},
\]
\[
 d\longmapsto
 \left(\frac{N_0+d}{R_0},\frac{N_0+N_0^2/d}{R_0}\right),
 \qquad U_R(p)=\frac{A_R(p)+\mathbf1_{R_0\mid2}}2.
\]
The inverse is $d=R_0b-N_0$. When $\gcd(R,pa)=1$ and $p\nmid a$,
the exact raw count is
$A_R(p)=c_R(a;-1)+2c_R(a;-p^{-1})$ in the exponent-box notation
of Theorem~\ref{esn:thm:raw-channels}. Equal prime residues do not
determine these counts (Theorem~\ref{esn:thm:non-descent}).
The finite-sample map in Theorem~\ref{esn:thm:sample-fourier}
sends these counts to an integral augmentation vector. Its Fourier
inversion recovers the aggregate residue-count function; the centering
$Q_{R,P}=6W_{R,P}-M_{R,P}u$ omits its constant Fourier component.
The sample $P$ remains in the domain, and no universal selector
nonemptiness is inferred.

\item \Verified\ \textbf{Two complex columns give the exact
four-dimensional \(SU(3)\) Gram cone.}
\[
 \operatorname{Mat}_{3\times2}(\mathbb C)/SU(3)
 \cong\operatorname{Herm}_2(\mathbb C)_{\geq0}.
\]
For the one-point compactified column representation,
\[
 S^{12}_{\rm col}/SU(3)\cong B^4,
\]
and the rank-one, zero, and infinity orbit types form its boundary \(S^3\).
This representation sphere is not the cusp-smash \(S^{12}\); see
\cref{thm:hrgc-gram-quotient,thm:hrgc-s12-ball-quotient}.

\item \Verified\ \textbf{Three complex columns give an exact cubic
rank-three cone with a distinguished rank-loss \(S^8\).}
For \(\mathscr V_3=\operatorname{Mat}_{3\times3}(\mathbb C)\), the complete
orbit invariants are
\[
 \mathscr V_3/U(3)\cong\operatorname{Herm}_3(\mathbb C)_{\geq0},
 \qquad
 \mathscr V_3/SU(3)
 \cong\{(H,z):H\geq0,\ |z|^2=\det H\}.
\]
The second space is \(\mathbb R^{10}\), so
\[
 S^{17}/SU(3)\cong S^9,\qquad
 \mathscr V_3^+/SU(3)\cong S^{10}.
\]
Inside that \(S^{10}\), the compactified rank-deficient locus is the
standard topologically locally flat \(S^8\), with
\[
 S^{10}\setminus S^8\cong\mathbb R^8\times\mathbb C^*\simeq S^1.
\]
The generator is the determinant phase \(\det A/|\det A|\), and the
trace-one rank-one stratum is \(\mathbb CP^2\).  Quotienting the determinant
phase gives \(\mathscr V_3^+/U(3)\cong B^9\) with boundary \(S^8\).  As a
whole this is a cubic Jordan/Gram carrier rather than a rank-two quadratic
Lorentz cone.  The marked angular quotient nevertheless gives an exact split
Lorentz observation with every complementary direction retained, as recorded
in the next item; see \cref{sec:hr3c-main}.

\item \Verified\ \textbf{The marked angular epimorphism gives a split
positive Lorentz observation, and its full fibre is retained.}
For
\[
 P=\begin{pmatrix}1&-1&0\\0&1&-1\end{pmatrix},\qquad
 u=(1,1,1)^{\mathsf T},\qquad
 R=P^*(PP^*)^{-1}
 =\frac13\begin{pmatrix}2&1\\-1&1\\-1&-2\end{pmatrix},
\]
put \(\Pi=uu^*/3\), \(Q_u=I_3-\Pi\), and
\[
 \mathcal O_P(H)=PHP^*,\qquad \jmath_P(G)=RGR^*.
\]
These positive maps satisfy
\[
 \mathcal O_P\jmath_P=\operatorname{id},\qquad
 \jmath_P\mathcal O_P(H)=Q_uHQ_u,
\]
with
\[
 \ker\mathcal O_P
 =\{a\,uu^*+uv^*+vu^*:a\in\mathbb R,\ v\in u^\perp\},\qquad
 \mathcal O_P^{-1}(0)\cap\mathcal P_3
 =\{\lambda uu^*: \lambda\geq0\}.
\]
Writing \(X=(R\;u)\), every Hermitian matrix is uniquely
\[
 H=X\begin{pmatrix}G&\psi\\\psi^*&\lambda\end{pmatrix}X^*,\qquad
 (G,\psi,\lambda)=
 \left(PHP^*,\frac13PHu,\frac19u^*Hu\right),
\]
so
\[
 \operatorname{Herm}_3(\mathbb C)
 \cong\operatorname{Herm}_2(\mathbb C)\oplus\mathbb C^2\oplus\mathbb R,
 \qquad 9=4+4+1,
\]
\[
 \det H=\lambda\det G-\psi^*\operatorname{adj}(G)\psi,\qquad
 H\succeq0\Longleftrightarrow
 G\succeq0, \psi\in\operatorname{im}G,
 \ \lambda\geq\psi^*G^+\psi.
\]
The faithful embedding
\[
 T_P(g)=RgP+\Pi=X(g\oplus1)X^{-1}
\]
acts by
\((G,\psi,\lambda)\mapsto(gGg^*,g\psi,\lambda)\).  For
\(g\in SL_2(\mathbb C)\), the summands are a Lorentz four-vector, the
defining complex spin representation, and a real scalar.  The vector action
factors through \(SO^+(1,3)\), while \(-I_2\) negates \(\psi\), so the full
carrier retains the spin double cover.

The section is proper and compactifies to
\[
 \jmath_P^+:B^4\cong\mathcal P_2^+
 \lhook\joinrel\longrightarrow(\partial\mathcal P_3)^+\cong S^8,\qquad
 e_2(\jmath_P(G))=\frac13\det G,
\]
and
\[
 \det(uu^*+\varepsilon\jmath_P(G))=\varepsilon^2\det G.
\]
The reverse boundary observation is nonproper.  In the strict \(SU(3)\)
quotient,
\[
 |\omega|^2=\lambda\det G-\psi^*\operatorname{adj}(G)\psi,
\]
and \(\omega=0\) is the rank-loss cone compactifying to
\(S^8\subset S^{10}\).  The four-branch chart supplies the exact toric seed
\[
 K(z)=\operatorname{diag}(z_0z_3,z_1z_2),\qquad
 \det K=z_0z_1z_2z_3,
\]
whose positive \(SL_2(\mathbb C)\)-orbit fills the future causal cone.
For \(Z\in\operatorname{Mat}_{3\times2}(\mathbb C)\) and retained
\(\alpha,c>0\),
\[
 p_\alpha(Z)=\alpha Z^*Z,\qquad
 m_{\alpha,c}^2(Z)=\frac{\alpha^2}{c^2}
 \sum_{a<b}|Z_{a1}Z_{b2}-Z_{a2}Z_{b1}|^2.
\]
This is the complete \(SU(3)\)-orbit momentum and is Lorentz covariant.
Its positive-mass locus is exactly rank two, and
\(m_{\alpha,c}^2(rZ)=r^4m_{\alpha,c}^2(Z)\), proving that the unrestricted
kinematic carrier has no positive threshold.  The exact reduction separates
four statements which the weak \(Z\)-equation alone does not.  First, for
\(m_0>0\) the colour quotient and reparametrization characteristic give
\[
 T^*\mathcal H^+_{m_0,c}
\]
with its canonical symplectic form, exactly the future spinless massive
particle.  Second, the intrinsic nonzero rank-one quotient gives
\(T^*\mathcal N_+^\times\).  Third, the full massless multiplier action is not
the standard massless variational system: the determinant constraint has zero
differential on rank one, its equation is \(\dot x=\lambda p\) without
\(\lambda=e\), four transverse rank-opening directions are non-gauge
first-order null directions, and \(Z=0\) admits arbitrary \(x(\tau)\).
Fourth, on the hard \(\omega_0\ne0\) packet the momentum differential has
rank four above threshold and rank three at threshold, so its displayed
action has \(\dot x=0\) in the strict interior and
\(\dot x\in\mathbb Rp\) only at the endpoint.  The packet's multiplication
operator is not the Hamiltonian generating this action; its positive endpoint
is a constraint-selected algebraic/free spectral threshold rather than a
dynamically generated interacting gap.  See
\cref{hlob:thm:split-observation,hlob:thm:block-coordinates,%
hlob:thm:lorentz-action,hlob:thm:cone-section,hlob:cor:Y3-block-form,%
hlob:thm:colour-momentum,wpr:thm:massive,wpr:thm:massless,%
wpr:thm:zero,wpr:thm:packet-action}.

\item \Verified\ \textbf{The determinant normal slice is an exact Lorentzian
host for explicit field models, but it does not itself select a mass gap.}
The real-linear isomorphism
\[
 \jmath_P:\operatorname{Herm}_2(\mathbb C)\xrightarrow{\ \sim\ }
 \mathcal N_P:=\{Y\in\operatorname{Herm}_3(\mathbb C):Yu=0\},
 \qquad G\longmapsto RGR^*,
\]
has inverse \(Y\mapsto PYP^*\), and
\[
 \det(uu^*+\varepsilon\jmath_P(G))=\varepsilon^2\det G.
\]
Writing
\(G(t,x,y,z)=
\bigl(\begin{smallmatrix}t+y&x+iz\\x-iz&t-y\end{smallmatrix}\bigr)\),
polarization gives
\[
 \eta_P(\jmath_P(G),\jmath_P(K))=tt'-xx'-yy'-zz'.
\]
The positive-semidefinite matrices are the future cone, the nonzero rank-one
matrices its null boundary, and the lifted congruence action induces
\(PSL_2(\mathbb C)\cong SO^+(1,3)\).  The future timelike determinant shell,
future-null locus, and candidate spacelike determinant shell are respectively
\[
 SL_2(\mathbb C)/SU(2),\qquad
 SL_2(\mathbb C)\cdot\operatorname{diag}(1,0),\qquad
 SL_2(\mathbb C)/SU(1,1).
\]

The standard metrics already retained by the Moore--Penrose construction
select
\[
 W_P=R^*R=C^{-1},\qquad
 K_P=\{a:a^*W_Pa=W_P,\ \det a=1\}\cong SU(2),
\]
and the positive form
\(E_P(G,K)=\operatorname{tr}(GW_PKW_P)\).  Under \(K_P\),
\[
 \operatorname{Herm}_2(\mathbb C)
 \cong\boldsymbol1\oplus\boldsymbol3,
\]
with trivial central action.  If this pointwise symmetry is promoted to a
gauge symmetry, the action
\[
 S[\mathcal A,G]=\int_M\left(
 \frac{|F_{\mathcal A}|^2}{2g_{\rm YM}^2}
 +\frac{Z_G}{2}|D_{\mathcal A}G|_{E_P}^2
 +\frac{m_G^2}{2}|G|_{E_P}^2
 +\frac{\lambda_4}{4}(\det G)^2\right)d\operatorname{vol}_g
\]
is nonnegative and gauge invariant, and
\(\widehat{\mathcal A}=R\mathcal A P\) intertwines the covariant
derivatives.  The normal field alone cannot screen a fundamental centre
charge; the full \((G,\psi,\lambda)\) carrier contains the centre-nontrivial
spinor \(\psi\), but its kinetic and interaction theory is not inserted by
this construction.

The candidate inverse-dual matrix and the positive Gram momentum are also
separated by an exact invariant:
\[
 \det G_{\eta,T}=-\frac23D_T<0,
 \qquad
 \operatorname{im}p_\alpha=
 \operatorname{Herm}_2(\mathbb C)_{\geq0}.
\]
Neither positive rescaling nor complex congruence can identify their inertia
types; a further morphism is required.  After a separate length/time
dictionary is supplied, the same affine host supports the ordinary
Klein--Gordon theory.  Its Fock Hamiltonian has excitation gap exactly
\(m_{\mathrm{KG}}\) for \(m_{\mathrm{KG}}>0\) and gap zero for
\(m_{\mathrm{KG}}=0\).  On a periodic spatial gauge regulator, the massless
transverse Gaussian frequency is \(2\pi/L\to0\).  Explicit countermodels also
show that entanglement occurs in both gapped and gapless systems.  A positive
auxiliary gap transfers to the physical Hamiltonian only through a map on the
entire excitation form domain satisfying
\[
 \|Ru\|\geq c_R\|u\|,\qquad
 \mathfrak b[Ru]\leq C_R
 (\mathfrak a_{\rm phys}[u]-E_{\rm vac}\|u\|^2),
\]
in which case
\(\operatorname{gap}(A_{\rm phys})\geq\gamma c_R^2/C_R\).
An intertwiner on a single gapped subspace is not enough.  See
\cref{hnsp-lorentz-theorem,hnsp-spin-theorem,%
hnsp-intrinsic-compact-representation,hnsp-intrinsic-gauge-scalar-theorem,%
hnsp-centre-screening-corollary,hnsp-candidate-spacelike,%
hnsp-determinant-orbit-geometry,hnsp-free-field-gap,%
hnsp-gauge-gaussian-gap,hnsp-coercive-transfer}.

\item \Verified\ \textbf{The marked cubic cone contains exact Lorentz
sections, but these do not identify it with the toroidal degeneration.}
With the oriented cross-product matrix \([a]_{\times}\),
\[
 \det\bigl(S+i[a]_{\times}\bigr)=\det S-a^{\mathsf T}Sa.
\]
Consequently each marked positive monodromy form \(B\) determines the exact
linear section
\[
 \ell_B(\lambda,a)=\lambda B+i[a]_{\times},\qquad
 \ell_B^{-1}\!\left(\operatorname{Herm}_3(\mathbb C)_{\geq0}\right)
 =\{\lambda\geq0:a^{\mathsf T}Ba\leq\lambda^2\det B\}.
\]
This is a future \((3+1)\)-Lorentz cone.  An explicit real-linear order-cone
map \(\Theta\) from \(2\times2\) Hermitian matrices satisfies
\[
 \det_3\Theta(H)=\frac{\operatorname{tr}H}{2}\det_2H
\]
and sends positive-semidefinite ranks \(2,1,0\) to \(3,2,0\).  It is not a
Jordan homomorphism: for a rank-one projection \(P\), the matrix
\(\Theta(P)\) has eigenvalues \(0,1/2,1\) and is not idempotent.  Neither map carries a toric period lattice, monodromy, or
singular-fibre filling.  In particular, the toric limit at fixed \(B=I_3\)
 is not the rank-drop operation \(\det B\to0\).  The rejected packet's claimed
index-three translation compression is likewise not a constructed map: the
actual six-to-two bridge is the saturated rank-lowering quotient
\(P:\mathbb Z^3\twoheadrightarrow\mathbb Z^2\), which sends the even and odd
Kuhn labels to the two unimodular target-triangle orbits.  See
\cref{thm:teb-hermitian-determinant-splitting,%
prop:teb-theta-order-lift,prop:teb-fan-morphism}.

\item \Verified\ \textbf{The rank-three toroidal successor maps onto the
candidate cusp with a connected elliptic kernel.}
The Kuhn fan gives a proper analytic family
\[
 \pi_3:N_3\longrightarrow\Delta,\qquad
 N_{3,t}\cong E_\tau^3\quad(t\ne0),
\]
whose deepest local equation is \(t=z_0z_1z_2z_3\), whose cusp matrix is
\(M_3=\bigl(\begin{smallmatrix}I_3&I_3\\0&I_3\end{smallmatrix}\bigr)\),
and whose filling map kills the three angular cycles and retains the three
deck cycles.  The irreducible central fibre has
\[
 (f_0,f_1,f_2,f_3)=(1,7,12,6),\qquad \chi(W_3)=6.
\]
Its normalization is the permutohedral threefold with orbit vector
\((1,14,36,24)\); the positive-codimension strata map with degrees
\(2,3,4\) onto the \(7,12,6\) quotient strata.  Hence any real
\(S^8\)-total-space fibration retaining \(W_3\) requires the remaining
singular fibres to contribute Euler sum \(-4\).
For an injective integral shear
\(\mathcal B:\bar L_3\hookrightarrow L_3\), this becomes the exact
determinant/index formula
\[
 \chi(W_{3,\mathcal B})
 =3!\,[L_3:\mathcal B\bar L_3]
 =3!\,|\det\mathcal B|.
\]

The explicit matrices
\[
 P=\begin{pmatrix}1&-1&0\\0&1&-1\end{pmatrix},
 \qquad
 Q=\begin{pmatrix}0&-1&1\\1&-1&0\end{pmatrix}
\]
define a surjection
\(\Phi=P\oplus Q:H_1(T^6)\twoheadrightarrow H_1(T^4)\) satisfying
\[
 \Phi M_3=M_F\Phi,\qquad c_F\Phi=Qc_3,\qquad
 PZ_3(t)=Z_2(t)Q.
\]
It extends over all cusp toric strata to a holomorphic map with connected
elliptic kernel.  Integrally,
\[
 M_3\cong M_F\oplus
 \begin{pmatrix}1&1\\0&1\end{pmatrix},
\]
and the order-three and order-four free affine peripheral actions have
compatible rank-six lifts.  The six Kuhn tetrahedron orbits map
\(S_3\)-equivariantly to the six marked diagonal Weyl chambers, but only as
fan/Cartan data; the four-branch toroidal germ and the quadratic Gram-rank
germ remain inequivalent.  What remains open is the single holomorphic
period-and-overlap family across both finite-order discs.  See
\cref{hrftd:cor:euler-shear-determinant,%
thm:teb-holomorphic-cusp-bridge,prop:teb-explicit-conjugacy,%
 prop:teb-finite-monodromy-lift,prop:teb-weyl-chambers}.

\item \Verified\ \textbf{The free rank-six cusp lattice, the six chambers,
and the Euler factor six are now separated exactly.}
The lattice native to the higher-rung construction is
\[
 \Lambda_{\rm tor}=H_1(T^6;\mathbb Z)
 =\mathbb Z^3_{\rm angular}\oplus\mathbb Z^3_{\rm deck}.
\]
For \(k=(1,1,1)\), the map \(\Phi=P\oplus Q\) gives an exact
\(S_3\)- and monodromy-equivariant sequence
\[
 0\longrightarrow\mathbb Z(k,0)\oplus\mathbb Z(0,k)
 \longrightarrow\Lambda_{\rm tor}
 \longrightarrow A_F\oplus B_F\longrightarrow0.
\]
It splits integrally as a monodromy module but not over
\(\mathbb Z[S_3]\).  The quotient copies are \(A_2\) weight lattices; their
root sublattices have index three.  The six standard homology generators are
two permutation orbits of size three, whereas the six Kuhn cells and six
marked Weyl chambers form one regular \(S_3\)-orbit.  Consequently
\[
 \operatorname{rk}\Lambda_{\rm tor}=6,\qquad
 \#\{\text{chambers}\}=3!=6,\qquad
 \chi(W_{3,\mathcal B})=3!|\det\mathcal B|
\]
do not identify six lattice directions with six strata.  The regular
six-element chamber set instead occurs in the rank-nine bilinear carrier
\[
 \operatorname{Hom}(\overline L_3,L_3)
 =\bigoplus_i\mathbb ZE_{ii}\oplus
  \bigoplus_{i\ne j}\mathbb ZE_{ij}:
\]
the off-diagonal units are a regular \(S_3\)-set, and
\(c_{(\pi_1\pi_2\pi_3)}\mapsto E_{\pi_1,\pi_3}\) is an equivariant
bijection.  These are exactly the mixed angular--deck basis pairs with zero
alternating pairing; the three excluded diagonal pairs have pairing one.
The ambient real matrix representation is equivariantly isomorphic to
\(\operatorname{Herm}_3(\mathbb C)\), but the symmetric/skew integral
splitting has index \(2^3\), and the images of the off-diagonal units are
indefinite.  This is a representation/incidence bridge, not a positive,
Jordan, determinant, or toric morphism.  The actual
nine-dimensional-boundary map is only the piecewise-linear angular map
\(\Phi_{\rm G}\circ\jmath\circ P\); the deck lattice is instead used by the
period and filling squares.  This map does not descend to a torus fibre.
The separate bundle-defect lattice \(H^4(U;\mathbb Z)\cong(\mathbb Z^2)^3\)
has a different domain and degree.  See
\cref{rsl:sec:rank-six,rsl:thm:three-sixes,rsl:prop:gram-scope,%
rnmt:thm:regular-torsor,rnmt:prop:representation-bridge}.

The successful lifted candidate actions are distinct from the diagonal
local-system candidate
\[
 M_1=I_3\otimes S,\qquad M_2=I_3\otimes Q,\qquad M_0=I_3\otimes T,
 \qquad M_1M_2M_0=I.
\]
Its finite factors have orders \(4\) and \(6\), but
\(I-M_j:T^6\to T^6\) is an isogeny for \(j=1,2\); every affine lift of
either finite action therefore has a fixed point.

\item \Refuted\ \textbf{There is no literal smooth complex-fourfold
analogue on a homotopy \(S^8\).}
An almost-complex tangent bundle would have
\[
 c_1=c_2=c_3=0,\qquad c_4=2u,\qquad p_2=2c_4=4u,
\]
whereas the signature formula forces \(p_2=0\).  A nonsingular
\(T^6\)-bundle over \(S^2\) with total space \(S^8\) is independently
excluded by both the homotopy sequence and Euler characteristic.  A real
singular \(T^6\)-over-\(S^2\) construction remains a legitimate search
problem, but its special fibres must obey
\(\sum_i\chi(F_i)=2\).  For monodromy \(\rho\) on the regular piece, its
actual filling inclusions must use the corrected Serre filtration
\[
 0\longrightarrow(\textstyle\bigwedge^q\Lambda)_\rho
 \longrightarrow H_q(E;\mathbb Z)
 \longrightarrow
 H_1(B^\circ;\textstyle\bigwedge^{q-1}\Lambda)
 \longrightarrow0,
\]
and the dual quotient
\(H^q(E;\mathbb Z)\twoheadrightarrow
(\bigwedge^q\Lambda^\vee)^\rho\).  Thus the integral Mayer--Vietoris test
uses coinvariants as the homological subobject and the base-one-cycle term
as quotient, not conversely; see
\cref{hrs8:sec:constraints,hrs8:eq:serre-regular-part,%
 hrs8:eq:serre-cohomology-invariants}.  Every effective torus action on
\(S^8\) has rank at most four, and the sharp coordinate model is
\[
 S^8/T^4\cong\operatorname{Susp}(\Delta^3)\cong B^4.
\]
This is a real two-pole singular compactification with generic \(T^4\)
orbit, not a three-special-fibre \(T^6\)-fibration over \(S^2\); see
\cref{hrs8:eq:torus-rank,hrs8:eq:t4-quotient}.

\item \Verified\ \textbf{The regular locus of the reconstructed \(S^6\)
has a six-coordinate, nonextendable defect lattice.}
For \(D=W\sqcup S_1\sqcup S_2\) and \(U=X\setminus D\),
\[
 \boxed{H^4(U;\mathbb Z)
 \cong H_1(W)\oplus H_1(S_1)\oplus H_1(S_2)
 \cong\mathbb Z^2\oplus\mathbb Z^2\oplus\mathbb Z^2.}
\]
Principal \(SU(3)\)-bundles on \(U\) are classified exactly by this
\(c_2\)-class, and the pair connecting morphism
\(H^4(U)\xrightarrow{\sim}H^5(X,U)\) proves that extension across all
three special fibres occurs exactly for \(c_2=0\).  The displayed group is
the free rank-six lattice \(\mathbb Z^6\), not a cyclic \(\mathbb Z/6\).

\item \Verified\ \textbf{The primitive fibre degree has exact order-twelve
and order-six normal-holonomy quotients.}  Complete proofs, including the
one-shot collar matrix, transpose system, saturated kernel, coefficient
reduction, tensor-centre map, and common-centre quotient, are given in
\cref{zdc:sec:bridge,zdc:thm:fibre-degree,zdc:prop:finite-reductions,%
zdc:thm:tensor-centres,zdc:thm:common-centre}.  Those group and lattice
theorems do not themselves define a gauge field or an energy operator.

In the marked basis
\[
 (a_0,b_0,a_1,b_1,a_2,b_2)
 =(\bar\gamma,\bar u,g_1,h_1,g_2,h_2),
\]
the transpose of the actual one-shot collar matrix gives
\[
 \boxed{[\deg_F]=(12,0,4,0,-3,0)},
\]
meaning
\[
 \operatorname{res}_F\!\left(\sum_{i=0}^2(n_i a_i+r_i b_i)\right)
 =(12n_0+4n_1-3n_2)\omega_F.
\]
It is primitive, with saturated kernel
\[
 \boxed{\ker\deg_F
 =\mathbb Z\langle
 b_0,b_1,b_2,a_0-3a_1,3a_1+4a_2\rangle,}
\]
Define
\[
 q_{\rm col}:L\longrightarrow H_1(X^\circ;\mathbb Z),\qquad
 a_0\mapsto c,\quad a_1\mapsto\bar g_1,\quad a_2\mapsto\bar g_2,\quad
 b_0,b_1,b_2\mapsto0.
\]
If \(\Phi_1:H_1(X^\circ;\mathbb Z)\to\mathbb Z\) has row
\((12,4\ell_1,3\ell_2)\) in the basis \((c,\bar g_1,\bar g_2)\), then the
factorization through the same collar functional that kills the final cusp is
\[
 \deg_F=\Phi_1\circ q_{\rm col},\qquad
 \Phi_1(j_*[\widetilde a_0])=p.
\]

Let \(\vartheta_1(a_1)=e^{-2\pi i/3}\) and
\(\vartheta_2(a_2)=e^{-2\pi i/4}\) be the local normal-line holonomies,
trivial on the other marked generators.  The actual relative twist gives
\[
 \chi_{12}\!\left(\sum_{i=0}^2(n_i a_i+r_i b_i)\right)
 =\vartheta_1(n_1a_1+r_1b_1)
  \vartheta_2(n_2a_2+r_2b_2)^{-1}
 =e^{-2\pi i(4n_1-3n_2)/12}.
\]
If \(\rho_m(k)=e^{-2\pi ik/m}\), then
\[
 \boxed{\chi_{12}=\rho_{12}\circ\deg_F,\qquad
        \chi_6:=\chi_{12}^2=\rho_6\circ\deg_F.}
\]
Therefore
\[
 H^4(U;\mathbb Z)/\ker\chi_{12}\cong\mathbb Z/12\mathbb Z,
 \qquad
 H^4(U;\mathbb Z)/\ker\chi_6\cong\mathbb Z/6\mathbb Z.
\]
Ordinary coefficient reduction instead gives
\[
 H^4(U;\mathbb Z/6)\cong(\mathbb Z/6)^6;
\]
it is not the single distinguished cyclic quotient selected by \(\chi_6\).
For a principal \(SU(3)\)-bundle \(P\to U\), the order-six character is
exactly
\[
 \chi_6(c_2(P))
 =\exp\!\left(-\frac{2\pi i}{6}
 \langle c_2(P)|_F,[F]\rangle\right).
\]
This is a topological congruence for the fibre second-Chern number, not a
line-operator or spectral assertion.

The tensor embedding
\[
 SU(2)\times SU(3)\hookrightarrow SU(6),\qquad
 (A,B)\longmapsto A\otimes B
\]
restricts to an isomorphism
\(\mu_2\times\mu_3\xrightarrow{\sim}Z(SU(6))=\mu_6\).  Since \(U\) has
homotopy dimension at most five, \(c_2\) classifies \(SU(n)\)-bundles for
\(n\ge3\), while
\[
 c_2(E_2\otimes E_3)=3c_2(E_2)+2c_2(E_3).
\]
Every chosen \(SU(6)\)-bundle class \(h\in H^4(U;\mathbb Z)\) therefore
admits a tensor-subgroup reduction, and all its \(c_2\)-factorizations are
\[
 (e,f)=(h+2t,-h-3t),\qquad t\in H^4(U;\mathbb Z).
\]
No nonzero \(h\) or reduction is canonically selected.  Separately,
\[
 H^5(U;\mathbb Z)
 \cong\ker\{\mathbb Z^3\xrightarrow{(1,1,1)}\mathbb Z\}
 \cong Q(A_2),
\]
for the combinatorial form making the three distinguished support basis
vectors orthonormal.  This is not a cup-product or intersection form, and
the support-redistribution lattice does not identify the rank-five kernel of
\(\deg_F\) with \(A_5\).

The finite-group isomorphism
\(H^4(U)/\ker\chi_6\cong Z(SU(6))\) supplies no \(SU(6)\) transition
functions or connection.  The injective tensor map does not factor through
\((SU(2)\times SU(3))/\mu_6\), because the central \(\mu_6\) maps to
nontrivial scalar matrices.  The distinct common-centre map is
\[
\begin{aligned}
 \Gamma_6
 &=\bigl\{((\lambda\mu)^{-1},\lambda I_2,\mu I_3):
             \lambda\in\mu_2,\ \mu\in\mu_3\bigr\},\\
 \frac{U(1)\times SU(2)\times SU(3)}{\Gamma_6}
 &\lhook\joinrel\longrightarrow U(6),\\
 [z,A,B]&\longmapsto z(A\otimes B),
\end{aligned}
\]
where \(\Gamma_6\cong\mu_6\).  Its determinant is \(z^6\), and its image
meets \(SU(6)\) in exactly \(\iota(SU(2)\times SU(3))\); the image is not
generally contained in \(SU(6)\).  The missing step is a degree-changing
identification with a degree-two common-centre lifting obstruction and its
line-operator pairing.  The separate global clutching classes on \(S^6\)
have \(c_3=\pm2ku\), \(k\in\pi_5(SU(3))\cong\mathbb Z\), and restrict
trivially to \(U\); see \cref{hrs8:subsec:su3-regular-locus}.

\item \Verified\ \textbf{The exact typed transfer from auxiliary
Rayleigh quotients to a physical excitation spectrum is proved.}
Let \(H\ge0\) be self-adjoint on \(\mathcal H_{\mathrm{phys}}\), and set
\[
 P_0=\mathbf1_{\{0\}}(H),\qquad
 \mathcal H_{\mathrm{exc}}=(I-P_0)\mathcal H_{\mathrm{phys}},\qquad
 \mathfrak h[\psi]=\lVert H^{1/2}\psi\rVert^2,\qquad
 \mathcal Q(H)=D(H^{1/2}).
\]
For a nonnegative quadratic form \(\mathfrak q_n\) and nonzero
\(u_n\in D(\mathfrak q_n)\subset\mathcal K_n\), suppose
\[
 J_n:\operatorname{span}\{u_n\}\longrightarrow\mathcal Q(H),\qquad
 \psi_n=(I-P_0)J_nu_n,
\]
and, for \(c>0\), \(C_n,\varepsilon_n\ge0\),
\[
 \lVert\psi_n\rVert\ge c\lVert u_n\rVert,\qquad
 \mathfrak h[\psi_n]\le C_n\mathfrak q_n[u_n]
 +\varepsilon_n\lVert u_n\rVert^2,\qquad
 C_n\frac{\mathfrak q_n[u_n]}{\lVert u_n\rVert^2}
 +\varepsilon_n\longrightarrow0.
\]
Then the self-adjoint excitation restriction
\[
 H_{\mathrm{exc}}
 =H\big|_{D(H)\cap\mathcal H_{\mathrm{exc}}}
\]
satisfies
\[
 \inf\sigma(H_{\mathrm{exc}})=0,\qquad
 0\in\sigma_{\mathrm{ess}}(H_{\mathrm{exc}}).
\]
The projection removes the full kernel of \(H\); orthogonality to a
single chosen vacuum is equivalent only after vacuum uniqueness is separately
proved.  This is a proved implication with fully stated antecedents, not a
claim that the present geometry supplies them.  See
\cref{cusp-interpretation-typed-physical-transfer}.

\item \Open\ \textbf{The physical Yang--Mills bridge data are absent.}
The workbench has not constructed a common gauge-invariant positive-metric
\(\mathcal H_{\mathrm{phys}}\), a four-dimensional quantum Yang--Mills
Hamiltonian \(H\), or maps \(J_n\) obeying the norm and energy inequalities
through gauge reduction, reflection-positive reconstruction,
renormalization, and the continuum and infinite-volume limits.  Therefore
the verified radial scalar spectra, auxiliary direct integral, classical
gauge Hessians, finite cochain operators, compact-link flux configurations,
 finite Gaussian oscillators, magnetic Laplacians, and reducible
 split-Cartan \(SU(r)\) Jacobi operators do not prove or disprove a mass gap for a physical quantum
Yang--Mills Hamiltonian.  See
\cref{cusp-gauge-missing-physical-bridge,%
cusp-interpretation-present-boundary}.

\end{enumerate}

\paragraph{Exact-computation certificates.}
The retained exact-arithmetic harnesses passed all \(28\) full-manuscript
checks, all \(38\) independently implemented core checks, and all \(13\)
cusp-spectral checks, with no failures; the \(13\)-check harness also passed
unchanged under \texttt{python -O}.  The first two harnesses cover finite and
cusp monodromy, invariant and coinvariant
lattices in exterior degrees \(0\) through \(4\), Smith forms and
saturation, finite-quotient norm operators, the relation-group determinant,
specialization indices, three twisting parameters, period-form identities,
and the corrected existential form of Remark~7.28.  The \(13\)-check
cusp-spectral harness is scoped to the radial period determinant and full
inverse transpose, the restricted dual maps and pairings, the integral
coframe wedge identity, the ambient Koide discriminant identity, the finite
coordinate eigenvalues, and the radial volume and first-correction
expansions.  The geometric collapse, finite-band and cusp inclusion maps,
Leray comparison, Kummer morphisms, cone carriers, compact-link construction,
and Koide counting theorem are proved separately rather than inferred from
those executions.

\paragraph{Independent review checkpoint.}
Reviewer~B independently reviewed the reconstruction existing up to its
evidentiary cutoff.  It withdrew the earlier topology-based rejection,
distinguished the intrinsic specialization \(\operatorname{sp}^q\) from the
then-open collapse pullback \(c^*\), and found no load-bearing defect in the
non-collapse nearby-cycle/Leray and based van Kampen route.  Its
provisionally favourable disposition and focused verification requests are
useful independent evidence about the argument it saw.  The later Codex
proofs in this workbench are separate and should not be attributed to that
review.  In its 26 August analytic follow-up, Reviewer~B independently proved
the finite-complex/UCT top-degree base-change statement on every
scheme-theoretic surface fibre and, from CAND-X, CAND-W, and TOP-H2, confirmed
the conductor--duality implication
$R^2f_*(T_X\otimes L)\ne0$ for every $L$ \cite{ReviewerBValidation2026}.
Those three inputs remain hypotheses in that review record; its cutoff labels
do not replace their later workbench classifications.

\paragraph{Current evidence boundary.}
No mathematical gap is presently recorded in the construction, either
topology certificate, sphere recognition, algebraic dimension, or Hodge
and Bott--Chern/Aeppli calculations.  Exact marked-fibre period-coordinate
representatives, global Bott--Chern generators, and their complete product
structure are now included.  The canonical data do not single out a literal
global coefficient atlas for the chosen smooth representatives through the
singular fillings; no such gauge is needed for the class or algebra theorem.
The new cusp spectral statements apply to the explicitly declared quotient
metrics and auxiliary or classical operators.  They now include the complete
fixed-energy lattice-point law
\(N_T(E)=\pi D_T\sum_n(E/(4\pi^2)-\lVert n\rVert^2)_++O_E(T+1)\)
and the full-torus integral magnetic bound
\(\operatorname{gap}(\Delta_{\mathsf E,T})
\le4\pi\sqrt{|\operatorname{Pf}(\mathsf E)|/D_T}\).
The modified-action toric Kummer map now completes the local
degree-\(N^2\) regular-fibre cover, while
retaining its direct degree \(N^3\), ramification, nonflat relative map, and
strictly local scope.  The exact maps \(q_2^{\wedge2}\),
\(q_{1,\mathrm{cusp}}^{\wedge4}\), and \(q_1^{\wedge4}\) now provide
topological \(S^{12}\)- and \(S^{24}\)-source carriers only by changing the
targets through suspension, collapse, and coning.  Two copies are rank-minimal
cusp-locally and four are rank-minimal over the full monodromy base.  Further
quotients \(\widehat q_3^{\wedge2}:S^{12}\twoheadrightarrow S^6\) and
\(\widehat q_\gamma^{\wedge4}:S^{24}\twoheadrightarrow S^8\) contain moving,
null-homotopic metric four-spheres with eventual exact Euclidean basepoint
balls.  Both full maps are themselves based null-homotopic.  More generally,
\(\widehat q_\gamma^{\wedge r}\) and
\(\widehat q_3^{\wedge r}\) are based null-homotopic for every \(r\ge2\).
The one-copy class of \(\widehat q_3\) is zero by the explicit framed
filling of Theorem~\ref{h24:thm:cusp-g-zero}; that of
\(\widehat q_\gamma\) remains open.
The conical-target classes of \(q_2^{\wedge2}\),
\(q_{1,\mathrm{cusp}}^{\wedge4}\), and \(q_1^{\wedge4}\) remain \Open.  The
topological five-manifolds \(\mathfrak E_{12},\mathfrak E_{24}\) realize
\(\mathbb R^4\) as a decompactification boundary fibre.  Unrescaled slices
collapse in every fixed compatible ambient metric; for the fixed
cusp-adapted metrics, the rescaled slice bubbles are \(\mathbb R^4\) while
the full ambient bubbles are \(\mathbb R^8\) and \(\mathbb R^{12}\).
This leaves intact the
nonextension to the fixed product target and supplies no compact complex
structure on either smash sphere.  The repaired compact-link calculation
keeps cover degree and mesh refinement independent and proves the exact
finite Gaussian gap and its zero-gap diagonal sequence, but not a
nonperturbative compact-link spectrum.  The physical Hilbert space, quantum
Hamiltonian, and comparison morphisms required by the typed transfer theorem
remain \Open; accordingly, these calculations neither prove nor disprove a
four-dimensional quantum Yang--Mills mass gap.  Separately, the exact
three-column quotient is now a rank-three cubic carrier with compactified
rank-loss pair \((S^{10},S^8)\), while a literal almost-complex \(S^8\) is
refuted.  The toroidal angular map gives that cubic carrier a split positive
observation onto the ordinary future Lorentz cone, with exact
\(9=4+4+1\) four-vector--spin-representation--scalar coordinates and a
proper \(B^4\hookrightarrow S^8\) section.  Its determinant-phase equation,
toric seed, and \(SU(3)\)-invariant momentum are explicit, and
\(m_{\alpha,c}^2(rZ)=r^4m_{\alpha,c}^2(Z)\) proves that the unrestricted
kinematic carrier has no positive threshold.  On the candidate's own regular
locus, the complete \(SU(3)\)-bundle
 classification is the free defect lattice
 \(H^4(U;\mathbb Z)\cong\mathbb Z^2\oplus\mathbb Z^2\oplus\mathbb Z^2\),
 and every nonzero class fails to extend across the three special fibres.
 The one-shot collar matrix computes the primitive fibre-degree row
 \((12,0,4,0,-3,0)\); the candidate's local normal holonomies are exactly its
 reductions \(\chi_{12}=\rho_{12}\circ\deg_F\) and
 \(\chi_6=\rho_6\circ\deg_F\).  The latter quotient is abstractly the same
 \(\mu_6\) as the centre of the tensor subgroup
 \(SU(2)\times SU(3)\hookrightarrow SU(6)\), but no degree-changing map to a
 common-centre lifting obstruction, gauge cocycle, connection, or physical
 line pairing has been constructed.
 The carrier theorem advances the physical bridge through an exact
colour-orbit momentum and a now fully reduced worldline action: its massive
stratum is the standard spinless particle, its punctured massless
characteristic quotient is standard while its multiplier action and apex are
irregular, and the hard-constraint packet propagates only on its threshold
stratum.  Its normal-slice
continuation now supplies a constant affine Minkowski metric, the spin and
Poincar\'e actions, an exact Klein--Gordon/Fock theory and finite regulator,
an independently chosen classical compact-group Yang--Mills theory, and the
metric-selected \(K_P\cong SU(2)\) gauge--scalar model.  These constructions
show both possible gap values: the free scalar gap is the externally supplied
\(m_{\mathrm{KG}}\), while the massless scalar and thermodynamic Gaussian
gauge limits are gapless.  They also isolate the candidate cusp matrix as
spacelike and outside the positive Gram-momentum image.  Dynamics on the
retained \(\psi,\lambda,\omega\) directions, a dynamically selected
gauge-invariant interacting spectral sector, a continuum quantum
reconstruction, and a full-domain comparison map to its physical excitation
space remain unconstructed.  The labels above mean that
complete arguments are present here for audit.  Specialist
replication and external adjudication are
separate validation boundaries, neither hypotheses nor gates on these
proofs; they remain welcome.  The corrected CDP sufficient theorem is now
proved in
\cref{thm:corrected-cdp-global,thm:corrected-cdp-fibrewise}.  Recovery of
the original theorem under its original hypotheses would instead require
overturning a verified premise of the constructed counterexample.

\section{Methodology}

The unit of audit is an exact mathematical assertion, not a paragraph's
plausibility.  Each assertion is recorded with its domain, codomain,
coefficient ring, variance convention, and dependencies.  The four outcomes
used in this paper are literal:

\begin{description}[leftmargin=8.2em,style=nextline]
\item[VERIFIED] A complete proof from the stated hypotheses is given or a
precisely cited theorem is applied after every hypothesis is checked.
\item[CONDITIONAL] The implication is proved, but its conclusion depends on
an explicitly named established or assumed hypothesis that is not proved in
the present argument.  This label is never a politeness substitute for
\textsc{open} or \textsc{refuted}.
\item[OPEN] No complete proof or counterexample is presently recorded.  This
is not a euphemism for false.
\item[REFUTED] A complete contradiction or counterexample is given to the
exact assertion.  A defect in a proof is not, by itself, a disproof of its
theorem.
\end{description}

Repair is an operation, not a fifth status: when a printed assertion is
false but a corrected statement is proved, the former is recorded as
\textsc{refuted} and the latter separately as \textsc{verified}, with every
downstream use recomputed.

Computations are converted into checkable certificates: displayed matrices,
integer row or column relations, determinant ideals, primitive kernels, or
explicit substitutions.  A software pass is corroboration, not a proof.
Likewise, an algebraically correct matrix calculation does not identify that
matrix with a geometric specialization, attaching, or boundary map; that
identification receives its own theorem and complete proof.

A theorem whose antecedents are explicitly part of its statement and whose
implication is completely proved is \textsc{verified} as that exact
implication.  \textsc{Conditional} is reserved for a conclusion that imports
an explicitly named established or assumed hypothesis not proved in the
present argument.  The current candidate-specific ledger has no such row.

Relevant human literature is searched proactively, not only when a remembered
fact seems doubtful.  Results are cited at the point of use with theorem,
section, or page locators and with their hypotheses intact.  Public working
\LaTeX{} editions such as the maintained SGA English source
\cite{SGAEnglish2026} may be used as searchable reading layers, but the
published SGA volumes \cite{SGA7I1972,SGA7II1973} remain canonical.  Every
reference in the compiled paper is resolvable without access to the editor's
computer.  The citation and tool-disclosure policy follows the Leiden
Declaration \cite{Leiden2026}; neither a model-generated review nor an
internal audit note substitutes for a human mathematical source.

\section{Relation to existing literature and exact scope}

The theorem-level conflict is with Campana--Demailly--Peternell
\cite[Theorem~2.2 and Corollary~2.3, p.~680]{CDP2020}.  Their hypotheses match
the now workbench-verified identities
$H^1(X;\mathbb Z)=H^2(X;\mathbb Z)=0$, algebraic dimension one, and
holomorphic algebraic reduction.  Their conclusion gives $c_3(X)\leq0$ and,
for a threefold homeomorphic to $S^6$, algebraic dimension zero.  The same
constructed threefold instead has $c_3(X)=2$ and $a(X)=1$.  It is therefore a
counterexample to Theorem~2.2 and Corollary~2.3 as stated; this conclusion uses
the complete candidate reconstruction, not merely the two printed proof-route
defects below, and is not resolved by observing that the candidate proof does
not invoke \cite{CDP2020}.

The first local issue isolated below occurs in the opening reduction of
\cite[\S7, p.~692]{CDP2020}.  The explicit nodal-cubic family in
\cref{prop:nodal-counterexample} refutes that reduction as a formal
sheaf-theoretic inference.  A second issue is the replacement of monodromy
coinvariants by the full fibre homology in
\cite[Lemma~3.2, pp.~84--85]{CDP1998}; CDP20 later imports the resulting
rank argument.  Either defect might in principle admit a different repair
under stronger hypotheses.  The original theorem cannot be repaired under
its printed hypotheses because the constructed threefold is a verified
counterexample.  The strongest intact replacement is now proved in
\cref{thm:corrected-cdp-global}: one global \(H^2\)-vanishing implies the
stated generic vanishing and \(c_3\le0\).  The stronger relative condition in
\cref{cor:corrected-cdp-relative} is sufficient for such a global vanishing
after a sufficiently positive twist from the base.  For reduced full fibres,
\cref{thm:corrected-cdp-fibrewise} replaces the p.~692 shortcut by the exact
connecting-map injectivity condition.

For the singular fibre, analytic duality is checked in two independent ways.
The proof in \cref{thm:cartier-duality} derives the result from smooth Serre
duality \cite[Theorem~9.37]{Lee2024}; Ramis--Ruget's dualizing-complex theorem
provides the singular analytic comparison
\cite[Theorems~1--2, pp.~81--84, and Proposition~2, p.~91]{RamisRuget1970}.
Relative duality for a proper analytic morphism is also cited in its precise
form as Ramis--Ruget--Verdier, Invent. Math. 13 (1971), 261--283
\cite{RamisRugetVerdier1971}.  Analytic cohomology and base change is
applied in the form of B\u{a}nic\u{a}--St\u{a}n\u{a}\c{s}il\u{a},
Chapter~III, \S3, Theorem~3.4 and Corollary~3.10, together with Grauert's
local-freeness/continuity theorem in Chapter~III, \S4, Theorem~4.12
\cite{BanicaStanasila1976}.  These locators cover all degrees used by the
structure-sheaf direct-image proof, not only its top-degree specialization.
The algebraic top-degree analogue in the Stacks Project, Tag~0EX2
\cite{StacksProject}, is a useful cross-check but is not substituted for the
analytic theorem.

The algebraic nearby-cycle formalism in SGA~7~II, Expos\'e XIII
\cite{SGA7II1973}, and the semistable/Raynaud framework in SGA~7~I,
Expos\'e IX \cite{SGA7I1972}, do not by themselves cover every step for the
present proper holomorphic map, which is not assumed algebraic.  The needed
analytic formalism is nevertheless available.  Dimca's Definition~4.2.1
constructs nearby cycles for complex analytic varieties; Proposition~4.2.2
identifies their stalks with Milnor-fibre cohomology; Proposition~4.2.11
commutes nearby cycles with proper analytic direct image; and
Theorem~2.3.26 gives proper base change, all for coefficient rings including
\(\mathbb Z\) \cite{Dimca2004}.  Kashiwara--Schapira
\cite[\S8.6]{KashiwaraSchapira1990} is an independent formal comparison.
Thus the general integral analytic framework is verified.  None of these
references by itself identifies the candidate's specialization image with
the full monodromy-invariant lattice.  That candidate-specific conclusion,
including saturation, is proved in
\cref{thm:p0-integral-specialization} for the reconstructed proper cusp
model, using the explicit local equations and a supported Gysin calculation.
The direct comparison with the cellular collapse is now proved in
\cref{thm:collapse-equals-specialization}: the neighborhood realization of
the nearby-cycle unit is
\[
 \operatorname{sp}^q=i_F^*(i_W^*)^{-1},
\]
and the candidate retraction gives $(i_W^*)^{-1}=r^*$, hence
$\operatorname{sp}^q=c^*$ over $\mathbb Z$ in every degree.

The topological input to the retraction is Brown's collaring theorem for the
boundary of a topological manifold \cite{Brown1962}, followed by Whitehead and
the cofibration argument in Hatcher \cite{Hatcher2002}.  Clemens's classical
collapse for semistable K\"ahler degenerations \cite{Clemens1977} is useful
context, but it is not used as a substitute for the candidate-specific proof:
the central fibre here is irreducible and nonnormal, so
\cref{lem:collapse-straightening,lem:positive-collapse,%
lem:phase-spread-collapse,prop:verified-cellular-collapse} constructs and
checks the map directly.

For compact-support Poincar\'e duality with local coefficients and its
projection formula we use Bredon \cite[Chapter~VI]{Bredon1997}, with
Hatcher \cite[\S3.H]{Hatcher2002} as a concise topological overview.  This
is the precise input behind the multiplicative identification
\(H^2(B,j_*\mathcal L)\simeq M_G\) used in the Leray calculation.

The completed global comparison uses the chain-level Mayer--Vietoris
convention of Hatcher \cite[\S2.2, pp.~149--150]{Hatcher2002}: the inclusion
of the intersection complex is \(z\mapsto(i_{A*}z,-i_{B*}z)\).  This fixes
the only global minus sign.  The candidate-specific content is not supplied
by that citation: the radial markings in the reconstructed finite fillings
prove that the two band inclusions retract to \(\pi_1\) and \(\pi_2\), while
the explicit toric disc and solid tori give the primitive cusp kernels.
Those geometric facts produce the matrices and the commuting Leray squares
in \cref{prop:actual-finite-band-maps,prop:cusp-mv-blocks,%
thm:mv-leray-comparison}.

For the Hodge-number question, Ugarte's general restrictions on a
hypothetical complex structure on $S^6$ \cite{Ugarte2000} and Angella's
later synthesis \cite{Angella2018} provide the correct general comparison
class.  The candidate-specific calculation is sharper: its reconstructed
structure-sheaf direct images and explicit vanishing of holomorphic forms
first determine the Hodge diamond up to $m=h^{1,2}$ and force
$E_2=E_\infty$.  Neither reference determines $m$ for this family.  The
workbench now does: the cyclic invariant lattices, logarithmic
Kodaira--Spencer divisor, normalized residue sequence, and order-four scalar
give $m=1$, hence $h^{1,1}=2$.  This is a candidate-specific calculation,
not a consequence imported from the general literature.

The spectral interpretation has also been audited against the primary
physical framework.  Jaffe--Witten formulate the four-dimensional
Yang--Mills mass gap on \(\mathbb R^4\) only after construction of the quantum
theory: on its physical Hilbert space, the Hamiltonian \(H\geq0\) has the
vacuum at energy zero and a mass gap precisely when there is a \(\Delta>0\)
such that
\[
  \sigma(H)\cap(0,\Delta)=\varnothing
  \qquad\text{(equivalently, }\sigma(H)\subseteq\{0\}\cup[\Delta,\infty)\text{)}.
\]
This is the physical mass-gap condition in \cite{JaffeWitten2000}.
Euclidean correlation data yield such an
\(\mathcal H_{\mathrm{phys}}\), vacuum, and positive time-translation
generator only through Osterwalder--Schrader reconstruction; the required
input includes Euclidean covariance and symmetry, distributional
regularity/growth, reflection positivity, and clustering, with the precise
reconstruction qualifications supplied by the two primary papers
\cite{OsterwalderSchrader1973,OsterwalderSchrader1975}.  Hodge decompositions,
periods, Leray terms, and coherent direct images are not Schwinger functions
and supply neither reflection positivity nor a physical Hamiltonian.  They
also do not determine even a geometric Laplace spectrum: the explicit
fixed-complex-torus calculation in
\cref{cusp-interpretation-fixed-complex-counterexample} keeps the complex
structure, integral cohomology, Hodge decomposition, coherent direct-image
data, and total volume fixed while the K\"ahler metrics
\[
 g_R=R^2(dx_1^2+dy_1^2)+R^{-2}(dx_2^2+dy_2^2)
\]
have \(\lambda_1(\Delta_{g_R})=4\pi^2/R^2\to0\).  Thus spectra depend on the
chosen metric and operator, and no spectral or physical conclusion follows
from the cohomological/direct-image data alone.  The auxiliary spectral
statements in the workbench are direct operator calculations; none is claimed
to be established by these references or to constitute an
Osterwalder--Schrader reconstruction of quantum Yang--Mills theory.

For the compact-link appendix, the historical scope has now been made
explicit.  Wilson supplies the plaquette/character lattice action
\cite{Wilson_1974}; L\"uscher supplies a selfadjoint strictly positive
transfer-matrix construction in the standard Wilson lattice theory
\cite{Luscher_1977}; and Menotti--Pelissetto establish
Osterwalder--Schrader positivity for the Wilson action in the corresponding
reflection setup \cite{MenottiPelissetto_1987}.  These citations identify the
standard constructions and the physical reconstruction boundary; they do
not prove the appendix's metric-dependent sector normalization, exact flux
action, finite Hessian spectrum, or Gaussian gap, all of which remain the
displayed candidate-specific calculations.  Nor do they turn that finite
Gaussian model into an Osterwalder--Schrader reconstruction, whose axiomatic
requirements remain those of
\cite{OsterwalderSchrader1973,OsterwalderSchrader1975}.

The general mechanisms used inside that candidate-specific calculation were
also checked against the acquired primary archival PDFs (including Deligne's
published erratum); no author- or publisher-supplied primary TeX was found for
these historical works.  Cartan's analytic quotient theorem
identifies the local quotient rings with invariants
\cite[\S1, Theorem~1, pp.~1--2]{Cartan1954Quotient}; Reynolds averaging then
gives exactness of finite-group invariants over $\mathbb C$.  Deligne's
logarithmic conormal and normalized-stratum residue formalism covers the
$d\log$ bases, relative sequence, and self-crossing residue target
\cite[Chapter~II, \S3.2--3.7, pp.~72--78]{Deligne1970}.  Katz--Oda supplies
the connecting-map and exterior-product formalism for the Gauss--Manin and
Kodaira--Spencer maps
\cite[\S2, equations~(5)--(10), pp.~202--204]{KatzOda1968}; Grothendieck and
the cited Stacks tags supply the splitting and elementary-modification
background on $\mathbb P^1$ \cite[Theorem~2.1, p.~126]{Grothendieck1957}.
Finally, Fr\"olicher's original lemmas identify the first and limiting pages
with Dolbeault and associated-graded de Rham cohomology
\cite[p.~643]{Frolicher1955}.  None of these sources supplies the printed
period derivatives or cusp branch maps; those remain the explicit
candidate-specific calculations in the workbench.  The subsequent
line-by-line audit therefore kept their normalization data explicit: the
finite determinant formula is defined up to a nowhere-zero frame unit, and
the constant cusp Kodaira--Spencer matrices are special-fibre reductions
modulo $t$.  These distinctions leave the divisor, residue, and rank
conclusions unchanged.

Once those candidate-specific Hodge and Fr\"olicher data are known, the
Bott--Chern and Aeppli dimensions are no longer open.  McHugh's general
compact-threefold formulas
\cite[Theorem~6.3, Proposition~6.4, equation~(7), and \S7.1]{McHugh2019}
express them in terms
of Dolbeault dimensions, $E_2^{3,1}$, the de Rham subgroup represented by
real $(1,1)$-forms, and the auxiliary spaces $K^{p,q}$.  For this candidate
$E_2^{3,1}=0$, $H^2_{\mathrm{dR}}(X)=0$, and
$\dim K^{p,q}\le b_{p+q}=0$ in every middle degree.  The last inequality is
not an injection $K^{p,q}\hookrightarrow H^{p+q}_{\mathrm{dR}}$: the image in
de Rham cohomology of the pure, $\partial$- and $\bar\partial$-closed forms
surjects onto $K^{p,q}$ because every pure $d$-exact form lies in McHugh's
denominator.  Direct substitution therefore yields the Bott--Chern diamond,
and Schweitzer's Hodge-star isomorphism
\cite[\S2.c, pp.~9--10]{Schweitzer2007} yields the Aeppli diamond, in
\cref{thm:exact-bott-chern-aeppli}.  Angella's published synthesis
\cite[\S4]{Angella2018} independently confirms the endpoint dimensions,
$h^{1,1}_{BC}=2h^{0,1}$, the relation
$h^{1,1}_{BC}+h^{2,2}_{BC}=2h^{2,1}_{BC}+2$, and Schweitzer duality.  The
literature supplies the general identities; the numerical specialization
depends on the workbench's candidate-specific proof that $h^{1,2}=1$ and
$E_2=E_\infty$.

The multiplicative refinement is also now complete.  Schweitzer's
definitions give the Bott--Chern algebra, its natural map to Dolbeault
cohomology, and the Chern-class convention
\cite[\S2 and \S8.a]{Schweitzer2007}, but they do not determine this
candidate's generators or products.  The workbench constructs global
$\bar\partial$-closed representatives $G,Q$, proves their exact marked-fibre
period-coordinate classes, and uses the two exact Fr\"olicher rows, the
factor-$12$ exponential--Leray comparison, Dolbeault--Serre duality, and top
integration to obtain all nine positive-degree Bott--Chern generators and
\[
 H^{*,*}_{BC}(X)=\mathbb C\mathbf1\oplus N,
 \qquad N^2=0.
\]
See
\cref{prop:bcpr-generators-are-bases,thm:bcpr-square-zero-algebra}.  This is
a candidate-specific calculation, not a consequence of the general
dimension identities.  It determines the cohomology classes and every
product, while making no claim of a preferred literal global coefficient
atlas for the noncanonical smooth representatives through the singular
fillings.

The algebraic reduction of a compact complex torus is used in the exact form
given by Birkenhake--Lange \cite[Chapter~2, \S6]{BirkenhakeLange1999}, with
the independent statement in Pareschi--Popa--Schnell
\cite[\S8, p.~2437]{PareschiPopaSchnell2017}.  The toric chart facts are
compared with Fulton \cite{Fulton1993}.  The exact classical-topology formula
\(\pi_1(X_\Sigma)\cong N/N_\Sigma\) is cited from Cox--Little--Schenck
\cite[Theorem~12.1.10]{CoxLittleSchenck2011}, with Danilov's original toric
treatment \cite[\S9]{Danilov1978} as an independent comparison.  All
assertions specific to the locally finite infinite fan and the candidate's
proper-discontinuity argument are proved directly.

Finally, the sphere-recognition endpoint uses Hurewicz and homological
Whitehead exactly as stated by Hatcher
\cite[Theorem~4.32 and Corollary~4.33]{Hatcher2002}, Smale's generalized
Poincar\'e theorem \cite[Theorem~A]{Smale1961}, and the computation of the
smooth endpoint by Smale
\cite[Theorem~1.1 and Corollary~1.2]{Smale1962}.  Kervaire--Milnor provide
an independent equivalent check: their opening table gives
$|\Theta_6|=1$, while the following Remark identifies $\Theta_6$ with
diffeomorphism classes in dimension six
\cite[pp.~504--505]{KervaireMilnor1963}.  Those results close the recognition
step after the candidate's fundamental group and integral homology are
established in \cref{thm:h2-vanishing,prop:geometric-pi1}; they are not used
to establish those inputs.

For the normalized cusp-restriction calculation, the public Stacks Project
references were checked at the formal spectral-sequence level rather than
relying on an unstated edge-map convention.  Leray is Tag~01F2; the passage
$E_r\mapsto E_{r+1}$ and the stable term are Tags~011N and~011O;
functoriality, identification with the associated graded, and convergence
are Tags~012O, 012U, and~012W \cite{StacksProject}.  Together with analytic
proper base change \cite[Theorem~2.3.26]{Dimca2004}, these results make the
evaluation square for
$H^2(X;\mathbb Z)\to H^0(B,R^2f_*\mathbb Z)$ commute with restriction to
every fibre.  Naturality of the exponential-sequence boundary then sends
$c_1(L)$ to $c_1(L|_{X_b})$.  The literature supplies this functorial
formalism; the injective differential
$d_2^{0,2}(2q)=\omega[\delta]$ remains the candidate-specific calculation
proved in \cref{lem:higher-transgressions}.

The independent direct boundary-Picard certificate uses Brown's standard
crossed-homomorphism model for cohomology of a group with coefficients in a
module \cite[Chapter~I, \S5]{Brown1982}.  That reference supplies only the
identification of $H^1$ with crossed homomorphisms modulo principal ones.  The
candidate-specific content is proved directly: the finite-cover norm, the
direction of each logarithmic conjugacy, the contraction formula for a
translated invariant two-form, the order-four parity lattice, the explicit
cusp solid torus, and Chern-class naturality under the local nearby-cycle
specialization.  These give
\[
 E=2nq,\qquad
 (\bar d_1,\bar d_2,\bar d_0)=(4n,-3n,0)[\delta],\qquad n=0.
\]
The exact Appell--Humbert source is Birkenhake--Lange, Chapter~1, \S2,
pp.~5--8, especially Theorem~2.1 on p.~7
\cite{BirkenhakeLange1999}; Chapter~2, \S6 is instead the separate source for
algebraic dimension used above.  Schnell's public notes record the precise
factor/coboundary and antisymmetrized connecting-cocycle formulas
\cite[equations~(3.3)--(3.10), pp.~3--7]{SchnellLecture3}, together with the
translation contraction $w\mapsto\iota_wE$
\cite[equation~(6.2), p.~1]{SchnellLecture6}.  These sources certify the
standard fibrewise formalism, not the candidate's chosen lifts, finite norms,
signs, or cusp winding; those are proved directly.

The completed third certificate uses only holomorphic $\mathcal O^*$ factors,
both finite fillings, the cusp argument principle, and the toric support
function.  Its exact boundary equations are
\[
 E=kq,\qquad
 (\bar d_1,\bar d_2,\bar d_0)
   =(2k,-\tfrac32k,0)[\delta],\qquad k=0.
\]
For normalized central degrees $(d_1,d_2,d_3)$, it then proves
\[
 B=\begin{pmatrix}d_2+d_3&d_3\\d_3&d_1+d_3\end{pmatrix},
 \qquad E_L((\lambda,0),(0,m))=-B(B_0\lambda,m).
\]
Fulton's chart, character-divisor, and support-function results
\cite[\S2.1, pp.~28--29; \S3.3, pp.~60--63; \S3.4,
pp.~63--72]{Fulton1993} and the parallel theorem-numbered statements in
Cox--Little--Schenck
\cite[Example~1.2.21, p.~31; Propositions~4.1.1--4.1.2, p.~171;
Theorems~4.2.8, p.~179 and~4.2.12, p.~182]{CoxLittleSchenck2011}
fix the convention: for $D=\sum a_\rho D_\rho$ the standard support function
has value $-a_\rho$, while
$\operatorname{div}(\chi^m)=\sum\langle m,u_\rho\rangle D_\rho$ has the
positive sign used in the proof.  Cartan theorem~B
\cite[\S2, fundamental theorem, p.~81]{Cartan1957} and the exponential
sequence \cite[Proposition~6.28 and equations~(6.35)--(6.36),
pp.~180--182]{Lee2024} trivialize only the residual line bundles whose
$c_1$ has already been proved zero.  All infinite-fan local finiteness,
descent factors, frame comparisons, and the final cross-block identity remain
candidate-specific proofs in \cref{subsec:analytic-all-picard}.

Bingener's analytic representability theorem was checked against the public
primary text, not inferred from the algebraic Picard literature.  For a
proper flat holomorphic map cohomologically flat in dimension zero, the
relative Picard functor is represented by a complex space which need not be
separated \cite[Introduction, p.~318; Satz~6.2, pp.~336--337]{Bingener1980}.
The same introduction states that the direct exponential-sequence quotient
works under topological local triviality and cohomological flatness in
dimension one, and explains why that quotient argument fails in general.
These statements provide representability for the cusp family but do not
provide the component-specialization or separatedness assertion required to
annihilate the normalized central restriction.  The candidate-specific
calculation in \cref{prop:local-cusp-picard-obstruction} makes this limitation
concrete: it constructs a descended local line bundle with generic Chern
class \(-2q\) and nonzero normalized central degree vector
\((-12,2,0,-12,2,0)\).  Thus all three proved global restriction theorems
genuinely use data absent from a cusp-local representability argument; no
general relative-Picard theorem has silently replaced the candidate-specific
finite-filling and common-boundary calculation.

\subsection{Targeted corpus cross-check (29 August 2026)}

Before extending the validate/repair/disprove chain, the relevant primary
texts were reread in the local LaTeX corpus and compared with their public
editions.  This check adds no unproved hypothesis; it records exactly which
formal tools the sources do, and do not, supply.

First, the primary transcription of Campana--Demailly--Peternell (2020),
Theorem~2.2 and Corollary~2.3, confirms the printed implication
\[
 H^1(X,\mathbb Z)=H^2(X,\mathbb Z)=0,\quad a(X)=1
 \Longrightarrow c_3(X)\leq 0,
 \qquad X\simeq S^6\Longrightarrow a(X)=0,
\]
under the stated compact-threefold and algebraic-reduction hypotheses.  Its
Proposition~2.4 identifies the obstruction with the vanishing of
$R^2f_*(T_X\otimes L)$, equivalently the fibrewise dual spaces involving
$H^0(\Omega_X^1\otimes L^*)$.  The workbench's conductor section,
absolute-fibre duality, and analytic base-change propositions prove a
nonzero singular-fibre contribution for the reconstructed candidate; this is
the exact point at which the earlier conditional objection is replaced by a
candidate-specific calculation.

The same primary source confirms that Proposition~2.4 is logically separate
from the defective attempts to establish its premise.  The workbench now
proves that premise-to-conclusion implication without importing either
defect.  If one \(L\) has
\(R^2f_*(T_X\otimes L)=0\), projection formula and Serre vanishing on the
base produce \(L_0\) with \(H^2(X,T_X\otimes L_0)=0\); generic endpoint
vanishing and semicontinuity then give
\[
 H^i(X,T_X\otimes M)=0\quad(i\ne1),\qquad
 \chi(X,T_X\otimes M)=\frac12c_3(X)\le0.
\]
This is a corrected sufficient theorem, not a reinstatement of Theorem~2.2:
the candidate violates this relative sufficient condition for every \(L\).
Separately, \(c_3(X)=2\) and the corrected global theorem imply by
contraposition that \(H^2(X,T_X\otimes L_0)\ne0\) for every \(L_0\).

The reread of CDP (1998), Lemma~3.2, confirms a separate printed rank defect:
for a split extension with fibre module $H_1(F)$ and monodromy group $Q$, the
abelianization contributes the coinvariants $(H_1(F))_Q$, not the whole
$H_1(F)$.  The elementary example
$\mathbb Z\rtimes_{-1}\mathbb Z$ has abelianization
$\mathbb Z\oplus\mathbb Z/2$, which is the exact correction used in
\cref{prop:monodromy-certificate}.

The SGA~7~II nearby-cycle transcription supplies the specialization--variation
sequence and the identity $\operatorname{Var}(\sigma)\alpha=\sigma-1$;
the Stacks/EGA transcriptions supply Leray functoriality and the precise
proper/flat base-change hypotheses.  Both are algebraic formal references,
so neither is silently applied as an analytic existence theorem.  The
candidate's integral specialization map, saturation, and stalk values remain
the explicit local calculations in \cref{thm:p0-integral-specialization}.

Finally, the independent Angella and McHugh transcriptions reproduce the
known six-sphere Hodge inequalities and the Bott--Chern dimension identities.
Substitution of the workbench values
\[
\begin{aligned}
 (h^{0,1},h^{0,2},h^{1,0},h^{2,0},h^{1,1},h^{1,2})&=(1,0,0,0,2,1),\\
 (h^{1,1}_{BC},h^{1,2}_{BC},h^{2,2}_{BC})&=(2,1,2)
\end{aligned}
\]
agrees with the explicit nine-generator square-zero algebra.  This is an external
consistency check, not a substitute for the candidate-specific quotient,
gluing, and period proofs.

\subsection{Stable-stem audit of the sphere-valued carriers
(3 September 2026)}

Let
\[
 g=\widehat q_3:(S^6,*)\longrightarrow(S^3,*),
 \qquad
 f=\widehat q_\gamma:(S^6,*)\longrightarrow(S^2,*)
\]
be the based sphere-valued quotients obtained after the chosen based
identification of \(X\) with \(S^6\).

Hatcher's full low-dimensional homotopy-group table
\cite[p.~339]{Hatcher2002} gives
\[
 \pi_6(S^3)=\mathbb Z/12,
 \qquad
 \pi_{12}(S^6)=\mathbb Z/2.
\]
In particular \(12g=0\).  Toda's Proposition~5.11
\cite[Proposition~5.11]{Toda1962} identifies the latter group as
\[
 \pi_{12}(S^6)=\mathbb Z/2\{\nu_6\circ\nu_9\}.
\]
The same proposition gives
\(\pi_{n+6}(S^n)=\mathbb Z/2\{\nu_n\circ\nu_{n+3}\}\) for \(n\geq5\).
Suspension carries each displayed generator to the next displayed
generator.  Hence the stabilization homomorphism
\[
 \operatorname{st}:\pi_{12}(S^6)\longrightarrow\pi_6^{\mathrm S}
\]
is an isomorphism.  Isaksen--Wang--Xu's stable-stem table
\cite[Table~1, p.~113]{IsaksenWangXu2023} identifies
\[
 \pi_3^{\mathrm S}=\mathbb Z/24\{\nu\},
 \qquad
 \pi_6^{\mathrm S}=\mathbb Z/2\{\nu^2\}.
\]
Since \(12\operatorname{st}(g)=0\) in
\(\mathbb Z/24\{\nu\}\), there is an \(a\in\mathbb Z/12\) such that
\[
 \operatorname{st}(g)=2a\nu.
\]
Stabilization is multiplicative for smash products, so
\[
 \operatorname{st}(g^{\wedge2})
 =\operatorname{st}(g)^2
 =4a^2\nu^2=0\in\pi_6^{\mathrm S}.
\]
Injectivity of the preceding stabilization isomorphism now gives
\([g^{\wedge2}]=0\in\pi_{12}(S^6)\).  Thus every based map
\(S^6\to S^3\), and in particular \(\widehat q_3\), has based
null-homotopic smash square.  Smashing this based null-homotopy with
further copies of \(g\) also gives
\[
 [g^{\wedge r}]=0\in\pi_{6r}(S^{3r})\qquad(r\geq2).
\]

For the second quotient, Hatcher's stable-range corollary
\cite[Corollary~4.24, p.~360]{Hatcher2002} says that
\(\pi_{n+k}(S^n)\) has reached the stable group \(\pi_k^{\mathrm S}\)
when \(n>k+1\).  Applying this with \(n=8\) and \(k=4\), and then using
the empty \(k=4\) row of Isaksen--Wang--Xu's complete stable-stem table
\cite[Table~1, p.~113]{IsaksenWangXu2023}, gives the exact equality
\[
 \pi_{12}(S^8)\cong\pi_4^{\mathrm S}=0.
\]
Consequently the sixfold reduced suspension
\[
 \Sigma^6f=f\wedge\operatorname{id}_{S^6}:
 S^6\wedge S^6\longrightarrow S^2\wedge S^6
\]
is based null-homotopic.  Under the canonical associativity and symmetry
identifications of based sphere smashes, the smash square factors as
\[
 f^{\wedge2}
 =\bigl(\operatorname{id}_{S^2}\wedge f\bigr)
  \circ\bigl(f\wedge\operatorname{id}_{S^6}\bigr):
 S^6\wedge S^6\longrightarrow S^2\wedge S^2.
\]
Its right-hand factor is \(\Sigma^6f\), so this factorization supplies a
based null-homotopy of \(f^{\wedge2}\).  For every \(r>2\), associativity
gives
\[
 f^{\wedge r}=f^{\wedge2}\wedge f^{\wedge(r-2)},
\]
and smashing the preceding based null-homotopy with the based map
\(f^{\wedge(r-2)}\) proves
\[
 [f^{\wedge r}]=0\in\pi_{6r}(S^{2r})
 \qquad(r\geq2).
\]
Hence both already constructed continuous surjective quotients
\[
 \widehat q_3^{\wedge2}:S^{12}\twoheadrightarrow S^6,
 \qquad
 \widehat q_\gamma^{\wedge4}:S^{24}\twoheadrightarrow S^8
\]
are based null-homotopic.  Neither calculation determines the one-copy
class \([g]\in\pi_6(S^3)\) or \([f]\in\pi_6(S^2)\).  Both arguments use
the sphere-valued postcompositions essentially, so they give no
null-homotopy of any map into a conically stratified target; those based
homotopy classes remain open.  Point-set surjectivity, quotient topology,
the embedded moving four-spheres, and their intrinsic metric limits are
unaffected by nullity of the global sphere-valued maps.

\subsection{Effective counting for the exact Koide sector}

The effective assertion used in the Koide chapter is matched to its literal
domain rather than transferred from a smoother or radial problem.  After the
five defining walls and the height \(H(a,q,n)=4a+17n\) are checked directly
against Definition~13.1, Benoist--Oh, Theorem~13.2
\cite{BenoistOh_2012}, applies to each of the finitely many integral orbits
and proves
\[
 R(X)=C_0X+O(X^{1-\eta})
 \qquad(\eta>0\text{ fixed but unspecified}).
\]
Neither the statement nor the proof of that theorem supplies
\(\eta\geq\tfrac12\).

The hyperbolic literature gives two stronger exponents for different
domains.  Liu--Sarnak, Proposition~3.2 \cite{LiuSarnak_2010}, gives the
smooth standard-radial exponent \(1/2+7/64=39/64\).  The exact order group
is a congruence arithmetic Fuchsian group because, for a maximal overorder
\(\mathcal M\),
\[
 \mathcal M^1(18\mathcal M)\subset\mathcal O^1\subset\mathcal M^1;
\]
Jacquet--Langlands and Kim--Sarnak give
\(\lambda_1\geq975/4096\) for that input
\cite{KimSarnak_2003,ChuLi_2016}.  Good's sharp radial theorem and its
Chatzakos--Petridis formulation \cite{Good_1983,ChatzakosPetridis_2016}
instead have a \(2/3\) error after the small-spectrum terms are displayed.
The exact radial Huber variable is \(7a+15\sqrt2\,n\), whereas the Koide
count uses \(4a+17n\) together with five sharp walls.  Consequently neither
radial theorem transfers its exponent to the target count.  A pointwise
square-root remainder for the unchanged domain remains open; the complete
hypothesis-by-hypothesis audit is \cref{koide-lit-determination}.

\subsection{Constrained dynamics and symplectic reduction}

The worldline calculation in
\cref{wpr:sec:worldline-reduction} uses three primary sources for its
general reduction framework.  Dirac's original constrained-Hamiltonian
paper \cite{Dirac1950} supplies the distinction between regular and singular
constraint dynamics and motivates retaining the multiplier and all primary
constraint equations.  Gotay--Nester--Hinds
\cite{GotayNesterHinds1978} supplies the presymplectic constraint algorithm:
one first restricts the closed two-form to the constraint locus, determines
its characteristic distribution, and checks whether successive consistency
conditions or rank jumps occur.  Marsden--Weinstein
\cite{MarsdenWeinstein1974} supplies the regular symplectic quotient by a
Hamiltonian symmetry action.

Their role here is structural, not candidate-specific evidence.  None of
these sources computes the Gram morphism
$P_\alpha(Z)=\alpha Z^*Z$, its $SU(3)$-orbits and stabilizers, the
rank-one zero of $\mathrm d(\det\circ P_\alpha)$, the apex branch, or the
hard-constrained packet's threshold rank jump.  Those data and the exact
cotangent sign are established directly in the cited workbench module.  In
particular, regular Marsden--Weinstein reduction applies to the massive
rank-two stratum and to the intrinsic punctured rank-one stratum after their
characteristic quotients have been proved regular; it does not authorize an
extension of that quotient through the irregular zero-momentum apex.

\subsection*{Quaternionic linear algebra for the 4 September continuation}
The local primary corpus was searched for quaternionic determinant and
Jordan-cone sources before promoting the two new Pro packets. The available
Alesker paper, \emph{Non-commutative linear algebra and plurisubharmonic
functions of quaternionic variables}, arXiv:math/0104209, Theorem 1.1.9,
gives the Moore determinant restriction and congruence law.
Wei Wang, \emph{The linear algebra in the quaternionic pluripotential theory},
arXiv:1805.07759v1, Propositions 2.1--2.3, supplies the corresponding
spectral theorem and complex matrix representation. Its primary LaTeX
source was retrieved and inspected. The sign conversion from its
representation to the present one is conjugation by
$\operatorname{diag}(I,-I)$, stated explicitly in
Section~\ref{sec:q444-reconciliation}.
These classical results are cited as background. The marked thirteen-slice
identity, nonunitary product coefficient, and scalar normal-force witness
are calculated explicitly here rather than attributed to those papers.

For the scalar-to-gauge comparison, David Tong's \emph{Lectures on Gauge
Theory}, Section 2.2.1, equation (2.27), identifies the $A_0$ equation
that must be retained after choosing temporal gauge:
\url{https://www.damtp.cam.ac.uk/user/tong/gaugetheory/gt.pdf}.
The primary PDF was obtained and the surrounding derivation read; the
local corpus was first checked for the corresponding moment-map and
constraint references. This workbench independently varies its own
anti-Hermitian $A_0$ and derives the matter moment
$[\Phi,\dot\Phi]+b\dot b^*-\dot b b^*$, so no coefficient or sign is
transferred from Tong's distinct field convention. The exact moving-frame
transformation retains the resulting nonzero Gauss residual.

For the vacuum-orbit comparison, the primary LaTeX of Andrea Gambioli,
\emph{Eight-dimensional $SU(3)$-manifolds of cohomogeneity one},
arXiv:math/0611796v1 (27 November 2006), was retrieved and its
$\mathbb{HP}^2$ example read:
\url{https://arxiv.org/abs/math/0611796}.
This is classical cohomogeneity-one geometry, not a new classification
claimed here. The displayed family and its invariant $\|b\|^2$ are
verified directly in Proposition~\ref{q444:vacuum-average}; the exact
bridge proved there is the Haar average of the marked Lorentz determinant.
The source archive and extracted TeX SHA-256 values are retained in the
supplement's \path{GAUGE_METRIC_SOURCE_RECEIPT.md}, together with the
separate Tong source used to check retention of Gauss' equation.

\subsection*{4 September: source roles in the integrated kinetic and metric continuation}
The canonical quaternionic proof now also uses the classical
gauge-invariant vector-operator context in Tong, \emph{Lectures on Gauge
Theory}, Section~2.7.3, printed pages~107--108,
\url{https://www.damtp.cam.ac.uk/user/tong/gaugetheory/gt.pdf}.
The model-specific eight-eigenvalue inertia calculation, all transverse
field equations including Gauss, conserved current, coefficient
$-4\|b\|^2$, and relation $\mu_{\mathrm d}^2=(2g^2/3)\|b\|^2$
are proved directly in Theorem~\ref{q444:gauge-spectrum} and
Proposition~\ref{q444:current-mass}; they are not imported from Tong's
different scalar example. Gambioli's Section~4 remains background for the
orbit types only. The original kinetic norm, nonzero normal acceleration,
exact current-free periodic continuation, and null waves are independently
proved in Propositions~\ref{q444:vk-horizontal}, \ref{q444:vk-sector},
Theorem~\ref{q444:vk-periodic}, and Proposition~\ref{q444:vk-wave}.

The flat-torus conventions recorded in the integrated metric section are
A.~Gogolev, \emph{Spectra of flat tori}, Section~2 (27 November 2017),
\url{https://andreygogolev.wordpress.com/2017/11/27/spectra-of-flat-tori/},
and V.~Ivrii, \emph{Partial Differential Equations}, Appendix~4.B,
\url{https://www.math.utoronto.ca/courses/apm346h1/20181/PDE-textbook/Chapter4/S4.B.html}.
The present integer-valued dual is $\mathcal N^{-\mathsf T}\mathbb Z^6$;
Ivrii incorporates $2\pi$ in the dual instead. U.~Lang,
\emph{Lecture Notes on Riemannian Geometry}, 16 June 2020,
Definition~4.1 and Proposition~4.4, printed pages~45 and~47,
\url{https://metaphor.ethz.ch/x/2020/fs/401-3532-08L/sc/DG2_16June2020.pdf},
supplies the standard horizontal-isometry definition of a Riemannian
submersion. The specialized proofs remain in
Theorems~\ref{q444:hermitian-cometric}--\ref{q444:metric-submersion}
and Proposition~\ref{q444:native-period-spectrum}: $W$ is a lattice
cometric, $h=2H^{-1}$, native $h=I$ gives $H=2I$, and the original
source/quotient lattice gaps have leading constants $4\pi^2$ and
$8\pi^2$ over $y^2$. These are source-convention pointers to the
integrated proofs, not claims of a new independent source retrieval in
this metadata pass or of a physical quantum spectral conclusion.

\paragraph{4 September: finite homogeneous quantization.}
The complete proof in Section~\ref{qh:chapter} uses the classical
bosonic-matrix mechanism of B.~Simon,
\emph{Some quantum operators with discrete spectrum but classically
continuous spectrum}, Annals of Physics \textbf{146} (1983), 209--220,
\url{https://doi.org/10.1016/0003-4916(83)90057-X}.
The original author-hosted scan was read, especially Section~2 and
Section~7, equation~(3) and Corollary~4. The present proof gives a direct
slice-oscillator and finite-cell argument in the original
$-\operatorname{tr}$ metric, rather than invoking a matrix-model
discreteness theorem. The root-frequency constant is $2\sqrt6$,
including eigenvalue walls.

The bounded-below Brownian-bridge formula is Theorem~1.1 of
B.~Simon, \emph{A Feynman--Kac Formula for Unbounded Semigroups},
arXiv:math-ph/9907022v1 (1999),
\url{https://arxiv.org/abs/math-ph/9907022}.
Its original LaTeX source was retrieved and read. The proof here
retains the anisotropic diffusion matrix and mass-volume measure.
D.~Tong, \emph{Lectures on Gauge Theory}, Section~2.2.1,
equations~(2.27)--(2.31),
\url{https://www.damtp.cam.ac.uk/user/tong/gaugetheory/gt.pdf},
supplies the Gauss-after-temporal-gauge convention. All triplet,
adjoint and volume factors are derived here, not copied from a
different gauge normalization. The corresponding infinite-volume
and continuum interacting spectral question remains open.

\paragraph{4 September: moving Fourier-state comparison and retained Gauss law.}
The original TeX of G.~Panati, H.~Spohn and S.~Teufel,
\emph{The time-dependent Born--Oppenheimer approximation},
ESAIM: M2AN \textbf{41} (2007), 297--314,
\url{https://arxiv.org/abs/0712.4369},
\url{https://doi.org/10.1051/m2an:2007023}, was retrieved and read
for the projected Hamiltonian, Berry and Born--Huang terms, and the
Duhamel/invariance distinction in Sections 1--2.1.
In particular compression alone does not justify a dynamical
approximation or higher-order expansion. Section~\ref{fq:section}
therefore proves an explicit residual and fixed-physical-time estimate
directly for the retained quartic two-mode operator; no general
Born--Oppenheimer theorem is asserted as new or invoked with unchecked
hypotheses. The positive ground-state and parameter regularity
arguments are included. Tong's cited Section 2.2.1 was also reread:
fixing temporal gauge does not remove Gauss. The complex rotating-pair
proof above retains it with the original anti-Hermitian conventions.

\paragraph{5 September: exact coupling elimination and spatial projection.}
The Sylvester map in
\eqref{q444:generic-sylvester-map} is the classical polynomial map
$(A,B)\mapsto AF+BG$; the primary Stacks Project text,
\url{https://stacks.math.columbia.edu/tag/00UA}, was checked for its
matrix description and coprimality criterion. The proof here includes
the degree-drop argument, the explicit retained coefficient ring,
the nonzero determinant witness and the elementary measure-zero proof.
No resultant theorem is presented as new. The new applications are
the exact degree-twenty-eight residual polynomial, its complete
coupling-fibre inverse, and the nonzero spatial-projection polynomial
for the specified rotating gauge ansatz. All metric, Gauss and scalar
conventions are derived in Section~\ref{sec:q444-reconciliation};
they are not inferred from a different model or an assumed quantum
spectral comparison.

\paragraph{5 September: two-cross residual certificate.}
The local primary corpus was checked, including David Tong's
\emph{Gauge Theory}, printed p.~40, equations (2.27)--(2.28),
\url{https://www.damtp.cam.ac.uk/user/tong/gaugetheory/gt.pdf},
on retaining the temporal Gauss constraint. The full matter-current
equations used here are derived in the retained-system subsection, not
imported from the pure-gauge example. The primary Stacks Project
Example 10.143.12, Tag 00UA, supplies the classical Sylvester-map
context. The degree-drop kernel, block factor, positive residual identity
and Cauchy--Schwarz constants in Section~\ref{q444:twocross:section}
are proved directly for the literal data. Neither standard linear algebra
nor Gauss law is represented as a new discovery or a quantum conclusion.

\paragraph{5 September: the modulo-$840$ lattice and its full Weyl quotient.}
Cheng--Duncan--Harvey \cite[Table~3, arXiv p.~21]{esn:CDH2014}
already record the group $GL_2(3)$ for the $A_5^4D_4$ Niemeier lattice;
their equation (2.30) defines the umbral group as the automorphism group
modulo the Weyl group. This abstract identification is not claimed as new.
Section~\ref{esn:sec:shared} proves the particular residue-indexed
index-$72$ glue, its exact minimum table, the transporter from the earlier
glue, and the signed code isomorphism with $GL_2(\mathbb F_3)$.
Theorems~\ref{esn:fa:thm:full-quotient}
and~\ref{esn:fa:thm:residue-weyl} exhaust all lattice isometries as
$wL_g$ and place the literal residue six-cycle in the Weyl group,
centralizing this chosen complement, with every fibre origin retained.

Mare--Willems \cite[Definition~2.2, Propositions~2.4 and~2.7,
and equations (5.4)--(5.6)]{esn:MareWillems2013} supply the flag
$F_4/\operatorname{Spin}(8)$ and its vector/even-half-spin/odd-half-spin
tangent modules. The marked dual classes $v,s,c$ are matched to their
shortest weight sets by the coordinate proof in
Theorem~\ref{esn:thm:d4-f4-triality}. The subsequent exact isometry
in Theorem~\ref{n24:thm:full-isometry} maps the entire rank-$24$
real span into the original three octonion sectors, with inverse,
glue test and the specified cyclic intertwining equation.
The retained residual and Foundation sources
\cite{esn:Residual2026,esn:ESMain2026,esn:Provenance} provide the earlier
shell and glue data. Their exact excerpts and snapshot digests accompany
the source package. The raw count and complementary-divisor quotient
are separately proved, with $A_3(2521)=0$, $A_3(3361)=6$ despite equal
prime residues modulo $840$. The Fourier map retains the actual finite
prime sample and its counts; neither universal selector positivity
nor a physical model is asserted by these lattice and counting maps.

\paragraph{6 September: the Leech and full triality comparisons.}
Conway--Sloane \cite[pp.~277--280]{esn:ConwaySloane1982} give the
holy coefficient construction. Borcherds \cite[Sections~6--7]{esn:Borcherds1985}
proves uniqueness and the uniform construction theorem.
Theorem~\ref{esnl:thm:leech-neighbor} now writes that construction for
the actual $72$ marked glue words: $\rho^2=84$, the two index-six
inclusions, all shifted-class minima and every norm-two exclusion are
proved in the original Euclidean coordinates. The original
$F_4$ Coxeter matrix also records its full $D_4$ discriminant action,
including $Cv_D=c_D+(1,-1,0,0)$, not equality of the chosen representatives.
Yokota \cite[Sections~1.10, 1.14, 1.16; Theorem~2.7.1]{h24:Yokota2009}
fixes the related-triple and Albert conventions. The workbench proves
the complete Peirce table, the determinant coefficients, exact Clifford
maps, and global associated-bundle descent in those conventions.
Mare--Willems \cite[Appendix~B, p.~519; orientation convention pp.~495--496]{esn:MareWillems2013}
supplies the real-flag cell theorem. The explicit finite exponential,
inverse chart, orientation sign, quotient continuity and homology maps
are all computed here. The Hopf pullback and triality reduction use
the previously stated retained $H$ and based source identification;
they do not silently recertify the global complex-threefold construction.

\section{Exact integral algebra of the printed monodromy data}

This section replaces opaque computational status reports by certificates
that can be checked by hand.  Let
\[
 V=\mathbb Z\langle\gamma,u,w,\delta\rangle,\qquad
 \Lambda=V^*
 =\mathbb Z\langle\widehat\gamma,\widehat u,
                         \widehat w,\widehat\delta\rangle .
\]
Vectors are columns in the displayed ordered bases.  The printed monodromy on
$V$ is
\[
T_1=
\begin{pmatrix}
1&0&-6&2\\0&-1&1&1\\0&-1&0&1\\0&0&0&1
\end{pmatrix},
\qquad
T_2=
\begin{pmatrix}
1&6&0&-3\\0&0&-1&1\\0&1&0&0\\0&0&0&1
\end{pmatrix}.
\tag{3.1}
\]
Put $T_0=(T_1T_2)^{-1}$ and let
$A_j=(T_j^{-1})^t$, $M_0=(T_0^{-1})^t$ act on $\Lambda$.

\begin{proposition}[monodromy certificate]\label{prop:monodromy-certificate}
\Verified\ Direct integral multiplication gives
\[
T_1^3=I,\quad T_2^4=I,\quad
T_0=
\begin{pmatrix}
1&0&0&1\\0&1&-1&0\\0&0&1&0\\0&0&0&1
\end{pmatrix},
\quad (T_0-I)^2=0,
\]
and
\[
A_1=
\begin{pmatrix}
1&0&0&0\\6&0&1&0\\-6&-1&-1&0\\-2&1&0&1
\end{pmatrix},
\quad
A_2=
\begin{pmatrix}
1&0&0&0\\0&0&-1&0\\-6&1&0&0\\3&0&1&1
\end{pmatrix},
\quad
M_0=
\begin{pmatrix}
1&0&0&0\\0&1&0&0\\0&1&1&0\\-1&0&0&1
\end{pmatrix},
\]
with $A_1A_2M_0=I$.
\end{proposition}

\begin{proof}
Multiplying the first two matrices in (3.1) gives
\[
T_1T_2=
\begin{pmatrix}
1&0&0&-1\\0&1&1&0\\0&0&1&0\\0&0&0&1
\end{pmatrix}.
\]
Its inverse is the displayed $T_0$.  Thus
$T_0-I$ sends $(a,b,c,d)^t$ to $(d,-c,0,0)^t$, and applying it twice
gives zero.  The complete order certificate is
\[
T_1^2=
\begin{pmatrix}
1&6&-6&-2\\0&0&-1&1\\0&1&-1&0\\0&0&0&1
\end{pmatrix},
\quad T_1^2T_1=I,
\]
\[
T_2^2=
\begin{pmatrix}
1&6&-6&0\\0&-1&0&1\\0&0&-1&1\\0&0&0&1
\end{pmatrix},
\quad
T_2^3=
\begin{pmatrix}
1&0&-6&3\\0&0&1&0\\0&-1&0&1\\0&0&0&1
\end{pmatrix},
\quad T_2^3T_2=I.
\]
In particular $T_1\ne I$ and $T_2^2\ne I$.  Transposing the integral inverses gives
the three displayed contragredient matrices; finally
$A_1A_2=(M_0)^{-1}$ by the definitions, hence $A_1A_2M_0=I$.
\end{proof}

\subsection{The cusp invariant, image, and coinvariant lattices}

\begin{theorem}[all exterior-power cusp lattices]
\label{thm:cusp-lattices}
\Verified\ The equality
\[
 \ker(M_0-I)=\operatorname{im}(M_0-I)
 =\mathbb Z\langle\widehat w,\widehat\delta\rangle
\tag{3.2}
\]
holds as an equality of saturated sublattices of $\Lambda$.  On $V$,
\[
V^{T_0}=\mathbb Z\langle\gamma,u\rangle,
\]
\[
(\textstyle\bigwedge^2V)^{T_0}
=\mathbb Z\langle
\gamma u,\gamma\delta,uw,\gamma w-u\delta\rangle ,
\qquad
(\textstyle\bigwedge^3V)^{T_0}
=\mathbb Z\langle\gamma uw,\gamma u\delta\rangle .
\tag{3.3}
\]
Every displayed invariant lattice is saturated.  Consequently the invariant
ranks in degrees $0,1,2,3,4$ are $(1,2,4,2,1)$.
\end{theorem}

\begin{proof}
From the displayed matrix,
\[
(M_0-I)(a,b,c,d)^t=(0,0,b,-a)^t.
\]
Its kernel is given by $a=b=0$, and its image consists of every vector
$(0,0,b,-a)^t$; this proves (3.2).  The two generators are members of the
chosen basis, so the lattice is saturated.

The action of $T_0$ on basis vectors is
\[
T_0\gamma=\gamma,\quad T_0u=u,\quad
T_0w=w-u,\quad T_0\delta=\delta+\gamma.
\tag{3.4}
\]
Thus a vector $a\gamma+bu+cw+d\delta$ is fixed exactly when $c=d=0$.
For degree two, write
\[
x=a\gamma u+b\gamma w+c\gamma\delta
  +d\,uw+e\,u\delta+f\,w\delta.
\]
Expanding (3.4) and collecting coefficients gives
\[
(\textstyle\bigwedge^2T_0-I)x
=(-b-e+f)\gamma u-f\gamma w-fu\delta .
\]
Hence $f=0$ and $e=-b$, with $a,b,c,d$ arbitrary, which is exactly the
middle lattice in (3.3).  These are integral equations with a unit
coefficient, so their solution lattice is saturated.

In the ordered basis
$(\gamma uw,\gamma u\delta,\gamma w\delta,uw\delta)$,
\[
(\textstyle\bigwedge^3T_0-I)(a,b,c,d)^t
=(d,-c,0,0)^t.
\]
The fixed lattice is therefore generated by the first two basis vectors and
is saturated.  Degrees zero and four are fixed because $\det T_0=1$.
\end{proof}

\subsection{Global coinvariants}

Let $G=\langle T_1,T_2\rangle$ and, for a $G$-lattice $L$, write
$L_G=L/((T_1-I)L+(T_2-I)L)$.

\begin{theorem}[coinvariant certificates]\label{thm:coinvariants}
\Verified\ The common invariant lattices are
\[
 V^G=\mathbb Z\gamma,\qquad
 (\textstyle\bigwedge^2V)^G=\mathbb Zq,\qquad
 (\textstyle\bigwedge^3V)^G=\mathbb Z\gamma uw,\qquad
 (\textstyle\bigwedge^4V)^G=\mathbb Z\operatorname{vol},
\]
where $q=uw+6\gamma\delta$ and
$\operatorname{vol}=\gamma uw\delta$.  There are torsion-free
identifications
\[
V_G=\mathbb Z[\delta],\qquad
(\textstyle\bigwedge^2V)_G=\mathbb Z[\gamma\delta],
\qquad
(\textstyle\bigwedge^3V)_G=\mathbb Z[uw\delta],
\tag{3.5}
\]
and $[uw]=6[\gamma\delta]$ in degree two.
\end{theorem}

\begin{proof}
Solving $(T_1-I)v=0$ and $(T_2-I)v=0$ in (3.1) gives
\[
V^{T_1}=\langle\gamma,2u+w+3\delta\rangle,\qquad
V^{T_2}=\langle\gamma,u+w+2\delta\rangle.
\]
If a vector belongs to both lattices, comparison of its $u$- and
$w$-coordinates gives simultaneously $2r=t$ and $r=t$ for the two
non-$\gamma$ coefficients; hence $r=t=0$.  Thus $V^G=\mathbb Z\gamma$.

The higher common invariants follow from the same integral nullspace
calculation.  In the ordered bases
\[
 (\gamma u,\gamma w,\gamma\delta,uw,u\delta,w\delta),
 \qquad
 (\gamma uw,\gamma u\delta,\gamma w\delta,uw\delta),
\]
the simultaneous equations
$(\bigwedge^kT_1-I)x=(\bigwedge^kT_2-I)x=0$ are respectively equivalent to
\[
 x_1=x_2=x_5=x_6=0,\qquad x_3=6x_4,
\]
and
\[
 x_2=x_3=x_4=0.
\]
Thus the integral solution lattices are generated by
$uw+6\gamma\delta=q$ and $\gamma uw$.  Both generators are primitive.
Finally $\det T_1=\det T_2=1$, so the common degree-four invariant lattice
is $\mathbb Z\operatorname{vol}$.

In $V_G$, columns of $T_1-I$ and $T_2-I$ give, among others,
\[
u+w=0,\qquad u-3\gamma=0,\qquad 2\gamma+u+w=0,
\qquad 2u+w=0.
\]
The first two relations give $w=-u$ and $u=3\gamma$; the last two then give
$2\gamma=0$ and $3\gamma=0$.  Hence $\gamma=u=w=0$.  Every column has
zero $\delta$-coordinate, so no relation involves $\delta$.  This proves
$V_G=\mathbb Z[\delta]$, including absence of torsion.

For degree two use the ordered basis
\[
(\gamma u,\gamma w,\gamma\delta,uw,u\delta,w\delta).
\]
In this basis the twelve columns of the relation map are
\[
B_2=
\left[
\begin{array}{rrrrrr|rrrrrr}
-2&1&1&-6&2&-8&-1&-1&1&-6&6&-3\\
-1&-1&1&-6&2&-6&1&-1&0&0&3&0\\
0&0&0&0&0&-6&0&0&0&0&6&0\\
0&0&0&0&0&1&0&0&0&0&-1&0\\
0&0&0&0&-2&1&0&0&0&0&-1&-1\\
0&0&0&0&-1&-1&0&0&0&0&1&-1
\end{array}
\right],
\tag{3.6a}
\]
where the two blocks are $\bigwedge^2T_1-I$ and
$\bigwedge^2T_2-I$.  Its column lattice is
\[
\operatorname{im}
[\,\textstyle\bigwedge^2T_1-I\mid\bigwedge^2T_2-I\,]
=\langle\gamma u,\gamma w,u\delta,w\delta,
          uw-6\gamma\delta\rangle .
\tag{3.6}
\]
For completeness, the third columns of the two blocks yield
$\gamma u$ and $\gamma u+\gamma w$; their difference yields $\gamma w$.
After these are removed, the fifth and sixth columns give
$2u\delta+w\delta$ and $u\delta+w\delta$, hence both
$u\delta$ and $w\delta$; the remaining nonzero relation is
$uw-6\gamma\delta$.  Conversely, substituting the twelve displayed columns
of the two exterior powers shows that they all lie in the right side.
After ordering the ambient basis as
$(\gamma u,\gamma w,u\delta,w\delta,uw,\gamma\delta)$,
the five generators have matrix $I_4\oplus(1,-6)^t$.
Its maximal minors have gcd $1$, so (3.6) is saturated.  The quotient is
therefore freely generated by $[\gamma\delta]$, with
$[uw]=6[\gamma\delta]$.

For degree three, let
$(a,b,c,d)=(\gamma uw,\gamma u\delta,\gamma w\delta,uw\delta)$.
The full relation matrix is
\[
B_3=
\left[
\begin{array}{rrrr|rrrr}
0&0&1&2&0&-1&0&-3\\
0&-2&1&-6&0&-1&-1&-6\\
0&-1&-1&-6&0&1&-1&0\\
0&0&0&0&0&0&0&0
\end{array}
\right].
\tag{3.6b}
\]
If $c_{i,r}$ denotes column $r$ of the $T_i$ block, then
$c_{1,2}-c_{2,3}=-b$, $c_{2,3}=-b-c$, and
$c_{1,3}=a+b-c$.  Hence $a,b,c$ belong to the column lattice.
Conversely every displayed column lies in $\langle a,b,c\rangle$, and the
fourth row is zero.  The relation lattice is therefore exactly the primitive
summand $\langle a,b,c\rangle$, so the quotient is
$\mathbb Z[uw\delta]$ with no hidden torsion.
\end{proof}

\subsection{Finite quotient fibres}

Put
\[
\eta_1=2u+w+3\delta,\qquad
\eta_2=u+w+2\delta,\qquad
\mathbf q=uw+6\gamma\delta,
\]
$m_1=3$, $m_2=4$, and
$v_1=(1,2,-4,0)^t$, $v_2=-(1,3,-3,0)^t$ in $\Lambda$.

\begin{proposition}[finite-fibre pullback indices]
\label{prop:finite-indices}
\Verified\ For the free affine actions and finite fibres geometrically
realized in \cref{prop:candidate-finite-fillings}, one has
\[
S_j=A_0^{(j)}/\langle x\mapsto A_jx+v_j/m_j\rangle
\]
with $H^*(A_0^{(j)};\mathbb Z)=\bigwedge^*V$, the
degree-$m_j$ covering $\pi_j:A_0^{(j)}\to S_j$ has injective pullback with
indices, in degrees $0$ through $4$,
\[
(1,3,1,1,3),\qquad (1,4,2,2,4).
\tag{3.7}
\]
Moreover
\[
\pi_1^*H^1(S_1)=\langle3\gamma,\eta_1\rangle,\qquad
\pi_2^*H^1(S_2)=\langle4\gamma,\eta_2\rangle,
\]
\[
\pi_1^*H^2(S_1)=(\textstyle\bigwedge^2V)^{T_1},\qquad
\pi_2^*H^2(S_2)
=\{a\gamma\eta_2+b\mathbf q:a\equiv b\pmod2\}.
\tag{3.8}
\]
The identification of these abstract lattices with the actual finite-fibre
stalks is \cref{lem:finite-filling-stalk-map}; hence no geometric hypothesis
remains open in this proposition.
\end{proposition}

\begin{proof}
The fundamental group is the extension
\[
\Gamma_j=\langle\Lambda,g\mid
g\lambda g^{-1}=A_j\lambda,\ g^{m_j}=v_j\rangle .
\]
After abelianization,
\[
H_1(S_j)=
(\Lambda_{A_j}\oplus\mathbb Zg)/
\mathbb Z([v_j],-m_j).
\tag{3.9}
\]
The matrices $A_j-I$ have Smith form $\operatorname{diag}(1,1,0,0)$,
so $\Lambda_{A_j}$ is free of rank two.  The relation in (3.9) is
primitive: if a prime divides $m_j$, its image $[v_j]$ modulo that prime is
nonzero because $\gamma((A_j-I)\Lambda)=0$ and
$\gamma(v_1)=1$, $\gamma(v_2)=-1$.  Thus $H_1(S_j)$ is free of rank two.
The quotient is a closed oriented four-manifold (the linear parts have
determinant one) and
$e(S_j)=e(A_0^{(j)})/m_j=0$.  Poincar\'e duality and the universal
coefficient theorem therefore give
\[
H_k(S_j;\mathbb Z)=(\mathbb Z,\mathbb Z^2,
\mathbb Z^2,\mathbb Z^2,\mathbb Z)
\]
with no torsion.
Indeed, Poincar\'e duality gives
$H_3(S_j)\simeq H^1(S_j)\simeq\mathbb Z^2$ and
$H^3(S_j)\simeq H_1(S_j)\simeq\mathbb Z^2$.  The universal coefficient
sequence for $H^3$ then forces
$\operatorname{Ext}(H_2(S_j),\mathbb Z)=0$, so $H_2(S_j)$ is torsion-free;
Euler characteristic zero gives its rank as two.

Transfer gives $\pi_{j*}\pi_j^*=m_j$; hence $\pi_j^*$ is injective and
its image lies in the $T_j$-invariants.  The
Lyndon--Hochschild--Serre spectral sequence for
$1\to\Lambda\to\Gamma_j\to\mathbb Z/m_j\to1$
\cite[Chapter~VII, \S6]{Brown1982} has
\[
E_2^{p,q}=H^p(\mathbb Z/m_j;\textstyle\bigwedge^qV).
\]
It computes $H^*(S_j;\mathbb Z)$ because the universal cover of this affine
torus quotient is $\mathbb R^4$, hence contractible.
Its transgression on an invariant $\xi\in V$ evaluates the extension
class: $d_2\xi=\xi(v_j)\bmod m_j$.  Since
$V^{T_j}=\langle\gamma,\eta_j\rangle$ and
$\eta_j(v_j)=0$, the degree-one images and indices in (3.8) follow.

In degree two, the integral invariant lattices are
\[
(\textstyle\bigwedge^2V)^{T_j}
=\langle\gamma\eta_j,\mathbf q\rangle
\]
with no omitted finite-index overlattice.  Indeed, writing an invariant as
$(a,b,c,d,e,f)$ in the ordered basis
$(\gamma u,\gamma w,\gamma\delta,uw,u\delta,w\delta)$,
the $T_1$ equations are
$e=f=0$, $a=2b$, $c=3b+6d$, while the $T_2$ equations are
$e=f=0$, $a=b$, $c=2b+6d$.  Using $\gamma u w\delta$ as orientation,
\[
(\gamma\eta_j)^2=0,\quad
(\gamma\eta_1)\mathbf q=3,\quad
(\gamma\eta_2)\mathbf q=2,\quad
\mathbf q^2=12.
\tag{3.10}
\]
Let $L_j=\pi_j^*H^2(S_j)$.  The intersection lattice of $S_j$ is
unimodular, and pullback multiplies its form by $m_j$.  Hence
$|\det L_j|=m_j^2$.  The two invariant Gram determinants from (3.10) are
$9$ and $4$.  If $k_j$ is the index of $L_j$ in the invariant lattice,
$k_j^2|\det(\bigwedge^2V)^{T_j}|=m_j^2$, so $k_1=1$ and $k_2=2$.

For $j=2$, multiplicativity of the spectral sequence gives
$d_2(\gamma\eta_2)=-[\eta_2]\ne0$ in
$\langle\gamma,\eta_2\rangle/
\langle4\gamma,2\eta_2\rangle$.  Thus
$\gamma\eta_2\notin L_2$.  The remaining index-two possibilities are
$L_2=\langle2\gamma\eta_2,a\gamma\eta_2+\mathbf q\rangle$,
$a\in\{0,1\}$.  The covering is \'{e}tale, so
$\pi_2^*K_{S_2}$ is trivial and $K_{S_2}^{\otimes4}$ is trivial.
Because $H^2(S_2;\mathbb Z)$ is torsion-free, $c_1(S_2)=0$ and
$w_2(S_2)=0$; the intersection form is even.  Dividing the pullback
pairings by four gives
\[
\begin{pmatrix}0&1\\1&a+3\end{pmatrix},
\]
so $a+3$ is even and $a=1$.  This proves (3.8).

For clarity, the displayed differential contains no suppressed coefficient
calculation.  For the cyclic group generated by $T_2$,
\[
H^2(\mathbb Z/4;V)=V^{T_2}/(I+T_2+T_2^2+T_2^3)V
=\langle\gamma,\eta_2\rangle/
  \langle4\gamma,2\eta_2\rangle.
\]
The extension class gives $d_2\gamma=\gamma(v_2)=-1$ and
$d_2\eta_2=\eta_2(v_2)=0$; the derivation rule therefore gives
$d_2(\gamma\eta_2)=-[\eta_2]$, which is nonzero because
$\eta_2\notin\langle4\gamma,2\eta_2\rangle$.

In degree four the index is the covering degree.  For degree three, direct
exterior calculation gives the primitive invariant bases
\[
(\textstyle\bigwedge^3V)^{T_1}
=\langle\gamma uw,\,uw\delta-4\gamma u\delta-2\gamma w\delta\rangle,
\]
\[
(\textstyle\bigwedge^3V)^{T_2}
=\langle\gamma uw,\,uw\delta-3\gamma u\delta-3\gamma w\delta\rangle.
\]
Against $(\gamma,\eta_j)$ their pairing matrices are
\[
\begin{pmatrix}0&1\\-3&0\end{pmatrix},
\qquad
\begin{pmatrix}0&1\\-2&0\end{pmatrix},
\]
so the absolute determinants are $3$ and $2$.  If the degree-three
index is $r_j$, the degree-one indices already proved are $3$ and $4$.
Pullback multiplies the perfect $H^1(S_j)\times H^3(S_j)$ pairing by
$m_j$, whose rank-two determinant is $m_j^2$.  Hence
$3\cdot3r_1=3^2$ and $2\cdot4r_2=4^2$; hence $r_1=1$, $r_2=2$.
Together with degrees zero through two this proves (3.7).
\end{proof}

\subsection{The terminal presentation and degree-two map}

\begin{proposition}[terminal group, including the nonabelian step]
\label{prop:terminal-group}
\Verified\ Consider
\[
G(\ell_0,\ell_1,\ell_2)=
\langle c,x,y\mid c\text{ central},\
xy=c^{\ell_0},\ x^3=c^{\ell_1},\ y^4=c^{\ell_2}\rangle
\]
and put $p=12\ell_0-4\ell_1-3\ell_2$.  If $p\ne0$, then
$G\cong\mathbb Z/|p|\mathbb Z$; if $p=0$, then $G\cong\mathbb Z$.
For $(\ell_0,\ell_1,\ell_2)=(0,1,-1)$ one has $p=-1$ and $G=1$.
\end{proposition}

\begin{proof}
The relation $xy=c^{\ell_0}$ gives $y=x^{-1}c^{\ell_0}$.  Since $c$ is
central, $x$ and $y$ commute, so the apparently nonabelian presentation is
already abelian.  In the ordered additive generators $(x,c)$ the remaining
relations form
\[
R=\begin{pmatrix}3&-\ell_1\\4&\ell_2-4\ell_0\end{pmatrix},
\qquad \det R=-p.
\]
The greatest common divisor of all entries of $R$ is $1$ because it divides
both $3$ and $4$.  Thus its Smith form is
$\operatorname{diag}(1,|p|)$ when $p\ne0$ and
$\operatorname{diag}(1,0)$ when $p=0$.  The stated groups follow.
\end{proof}

\begin{proposition}[terminal degree-two lattice]
\label{prop:terminal-degree-two}
\Verified\ Let $A:\mathbb Z^6\to\mathbb Z^4$ have columns
\[
\begin{aligned}
c_1&=(2,-4,-2,3)^t,&c_2&=(1,-2,-2,3)^t,\\
c_3&=(3,0,-4,0)^t,&c_4&=(0,1,0,-1)^t,
\qquad c_5=c_6=0.
\end{aligned}
\]
Then $\operatorname{im}A$ is saturated of rank three,
\[
\operatorname{im}A=\ker(4,2,3,2),
\qquad
\operatorname{coker}A\cong\mathbb Z,
\tag{3.11}
\]
and $(4,2,3,2)$ is the primitive left-kernel generator.
\end{proposition}

\begin{proof}
The row $\Phi=(4,2,3,2)$ annihilates the four nonzero columns:
\[
\Phi c_1=8-8-6+6=0,\quad
\Phi c_2=4-4-6+6=0,\quad
\Phi c_3=12-12=0,\quad
\Phi c_4=2-2=0.
\]
It is primitive since the gcd of its entries is $1$, so
$\ker\Phi\subset\mathbb Z^4$ is a saturated rank-three lattice.
The $3\times3$ minor on rows $(1,2,3)$ and columns $(1,2,4)$ is $2$,
whereas the minor on rows $(1,2,4)$ and the same columns is $-3$.
Consequently the ideal generated by all maximal minors contains $(2,3)=(1)$.
This proves that $A$ has rank three and saturated image.  We have
$\operatorname{im}A\subseteq\ker\Phi$ between rank-three lattices.
The finite quotient $\ker\Phi/\operatorname{im}A$ injects into the
torsion-free group $\mathbb Z^4/\operatorname{im}A$, so it is zero.
Thus equality holds, and the primitive surjection $\Phi:\mathbb Z^4\to
\mathbb Z$ induces (3.11).
\end{proof}

\begin{proposition}[all-degree terminal homology package]
\label{prop:terminal-all-degree}
\Verified\ Let $C_\bullet$ be the degree-by-degree terminal package printed
for the glued six-manifold, and assume only the following algebraic outputs
of its displayed maps:
\begin{enumerate}[label=(\roman*)]
\item $H_0(C_\bullet)=H_6(C_\bullet)=\mathbb Z$ and the degree-one
presentation is $G(\ell_0,\ell_1,\ell_2)$ of
\cref{prop:terminal-group};
\item in degree two there is an exact sequence
$0\to\mathbb Z\to H_2^\circ\to\mathbb Z\to0$, the retained boundary
lattice kills the second quotient, and its remaining generator maps to
$p$ times a primitive generator of the first copy of $\mathbb Z$;
\item the package obeys integral Poincar\'e duality and the universal
coefficient theorem in formal dimension six, and its Euler characteristic
is $2$.
\end{enumerate}
Then, for $p\ne0$,
\[
 \bigl(H_0(C_\bullet),\ldots,H_6(C_\bullet)\bigr)
 \cong
 \bigl(\mathbb Z,\mathbb Z/|p|,\mathbb Z/|p|,
 \mathbb Z/|p|,\mathbb Z/|p|,0,\mathbb Z\bigr).
\tag{3.12}
\]
In particular, the printed terminal package has the integral homology of
$S^6$ when $p=-1$.
\end{proposition}

\begin{proof}
By \cref{prop:terminal-group}, the degree-one group is
$\mathbb Z/|p|$.  In degree two, exactness and the stated retained-boundary
map first quotient $H_2^\circ$ down to the primitive first summand and then
quotient that summand by $p$; hence $H_2(C_\bullet)=\mathbb Z/|p|$.
Both groups are finite, so their Betti numbers vanish.  Universal
coefficients and duality give
\[
 H_5\cong H^1=0,
 \qquad
 H_4\cong H^2\cong
 \operatorname{Ext}(H_1,\mathbb Z)=\mathbb Z/|p|.
\]
The Euler-characteristic identity is therefore
$2=1-b_3+1$, so $b_3=0$.  Finally
\[
 H_3\cong H^3\cong
 \operatorname{Ext}(H_2,\mathbb Z)=\mathbb Z/|p|,
\]
which proves (3.12).
\end{proof}

\paragraph{Geometric boundary.}
The calculations in this section deliberately isolate the printed lattices,
presentations, and matrices.  Their geometric realization is proved
separately: the finite-band and cusp maps are identified in
\cref{prop:actual-finite-band-maps}, the based fundamental-group calculation
is \cref{record:vankampen}, and the comparison with the intrinsic Leray
transgressions is \cref{thm:mv-leray-comparison}.

\section{An exact Erd\H{o}s--Straus/S6 terminal-matrix crosswalk}
\label{sec:es-s6-crosswalk}

This section records an arithmetic relation suggested by comparing the
candidate's terminal presentation with the divisor-shell coordinates used in
the maintained Erd\H{o}s--Straus corpus~\cite{ClankersES2026}.  All maps below
are proved over the integers.  The conclusion is an exact change of relators;
it is not a geometric map from the reciprocal cubic surface to the candidate
complex threefold.

\subsection{The first-half shell and its divisor coordinates}

\begin{proposition}[first-half denominator]
\label{prop:es-first-half-denominator}
\Verified\ Let \(p\equiv1\pmod4\) be an odd prime.  If
\[
 \frac4p=\frac1x+\frac1y+\frac1z,
 \qquad x,y,z\in\mathbb Z_{>0},\qquad x\le y\le z,
\]
then
\[
 \frac p4<x\le\frac p2.
\]
\end{proposition}

\begin{proof}
The lower bound follows from \(1/x<4/p\).  Suppose \(x>p/2\).  Clearing
denominators and reducing modulo \(p\) gives
\[
 4xyz=p(xy+xz+yz),\qquad p\mid xyz.
\]
If \(p\mid x\), then \(x,y,z\ge p\), so the sum of the three reciprocals is
at most \(3/p\).  If \(p\mid y\), then \(y,z\ge p\), and
\[
 \frac1x+\frac1y+\frac1z<\frac2p+\frac2p=\frac4p.
\]
Both alternatives are impossible.  Hence \(z=kp\) for an integer \(k\ge1\),
and substitution gives
\[
 (4k-1)xy=kp(x+y).
\]
Neither \(x\) nor \(y\) is divisible by \(p\), so \(p\mid4k-1\).  Write
\(4k-1=\ell p\).  Since \(p\equiv1\pmod4\), one has
\(\ell\equiv3\pmod4\), hence \(\ell\ge3\).  Cancelling \(p\) gives
\(\ell xy=k(x+y)\).  Since \(x\le y\),
\[
 x\le\frac{2k}{\ell}
   =\frac{\ell p+1}{2\ell}
   =\frac p2+\frac1{2\ell}<\frac{p+1}{2}.
\]
But the integer \(x>p/2\) is at least \((p+1)/2\), a contradiction.
\end{proof}

Fix henceforth a prime \(p\equiv1\pmod{12}\), and put
\[
 L=\frac{p-1}{4},\qquad
 R_j=4j+3,\qquad
 a_j=\frac{p+R_j}{4},\qquad 0\le j<L.
\]
The preceding proposition shows that every solution's smallest denominator
lies in the exact candidate set \(\{a_j:0\le j<L\}\).  Notice that
\[
 4a_j=p+R_j,\qquad
 \gcd(R_j,p)=\gcd(R_j,a_j)=\gcd(p,a_j)=1.
\]

\begin{proposition}[complete two-channel divisor selector]
\label{prop:es-two-channel-selector}
\Verified\ Define
\[
 \begin{aligned}
 E_j&=\{u\in\mathbb Z_{>0}:u\mid a_j^2,\ R_j\mid4u+1\},\\
 M_j&=\{u\in\mathbb Z_{>0}:u\mid a_j^2,\ R_j\mid4u+p\}.
 \end{aligned}
\]
Then the Erd\H{o}s--Straus equation for \(p\) has a positive integral
solution if and only if at least one \(E_j\) or \(M_j\) is nonempty.
For \(u\in E_j\) set \(\varepsilon=0\), and for \(u\in M_j\) set
\(\varepsilon=1\).  With
\[
 T=\frac{4u+p^\varepsilon}{R_j},\qquad
 K=\frac{p^{1-\varepsilon}T+1}{4},
\]
the corresponding solution is
\[
 (x,y,z)=
 \left(a_j,\frac{p^\varepsilon a_jK}{u},pK\right),
\]
up to interchange of \(y\) and \(z\).
\end{proposition}

\begin{proof}
Fix \(x=a_j\), and abbreviate \(R=R_j\), \(a=a_j\), and \(N=pa\).
The remaining equation is
\[
 Ryz=N(y+z).
\]
It is equivalent to
\[
 (Ry-N)(Rz-N)=N^2.
\]
Both factors are positive.  Since \(\gcd(R,N)=1\), positive solutions are in
bijection with positive divisors \(d\mid N^2\) satisfying
\[
 d\equiv-N\pmod R;
\]
the complementary divisor \(e=N^2/d\) then also satisfies
\(e\equiv-N\pmod R\), and
\[
 y=\frac{N+d}{R},\qquad z=\frac{N+e}{R}.
\]

Because \(\gcd(p,a)=1\), the \(p\)-adic exponent of \(d\) is \(0\), \(1\),
or \(2\).  Complementation \(d\leftrightarrow e\) exchanges exponents \(0\)
and \(2\), so choose the exponent-two representative in that pair.  If
\(d=p^2u\), then \(u\mid a^2\), and
\[
 R\mid d+N
 \Longleftrightarrow R\mid pu+a
 \Longleftrightarrow R\mid4u+1,
\]
where \(4a\equiv p\pmod R\) and \(p\) is invertible modulo \(R\).  This is
the \(E_j\) channel.  If \(d=pu\), then
\[
 R\mid d+N
 \Longleftrightarrow R\mid u+a
 \Longleftrightarrow R\mid4u+p,
\]
which is the \(M_j\) channel.  This proves completeness of the two selector
sets.

For either channel, \(R_jT=4u+p^\varepsilon\).  Modulo \(4\), this says
\(-T\equiv1\), so \(T\equiv3\pmod4\); the displayed \(K\) is therefore a
positive integer.  Direct calculation using \(4a=p+R_j\) gives
\[
 R_jK=a+p^{1-\varepsilon}u,
 \qquad Ta=u+p^\varepsilon K.
\]
Since \(u\mid a^2\), the second identity implies
\(u\mid p^\varepsilon aK\).  Finally,
\[
 \frac{u}{p^\varepsilon aK}+\frac1{pK}
 =\frac{p^{1-\varepsilon}u+a}{paK}
 =\frac{R_j}{pa}
 =\frac4p-\frac1a,
\]
which proves the asserted formula and its converse.
\end{proof}

\subsection{The unimodular relation-matrix bijection}

Put
\[
 \lambda(\ell_0,\ell_1,\ell_2)
 =12\ell_0-4\ell_1-3\ell_2.
\]
By \cref{prop:terminal-group}, the candidate's ordered terminal relation
matrix in the ordered generator basis \((x,c)\), for
\(\ell=(\ell_0,\ell_1,\ell_2)\), is
\[
 \mathsf R(\ell)=
 \begin{pmatrix}
 3&-\ell_1\\
 4&\ell_2-4\ell_0
 \end{pmatrix},
 \qquad \det\mathsf R(\ell)=-\lambda(\ell).
\]

For \(\varepsilon\in\{0,1\}\), define the marked selector space
\[
\mathcal W_\varepsilon(p)=
 \left\{(j,u)\in\mathbb Z\times\mathbb Z_{>0}:
 \begin{array}{l}
 0\le j<L,\ u\in\mathbb Z_{>0},\ u\mid a_j^2,\\
 R_j\mid4u+p^\varepsilon
 \end{array}
 \right\}
\]
and the normalized divisor-supported terminal slice
\[
\mathcal H_\varepsilon(p)=
 \left\{(j,\ell_0,\ell_1)\in\mathbb Z^3:
 \begin{array}{l}
 0\le j<L,\ \lambda(\ell_0,\ell_1,1)=p^\varepsilon,\\
 T:=4\ell_0-1>0,\quad
 u:=\ell_1+jT>0,\quad u\mid a_j^2
 \end{array}
 \right\}.
\]

\begin{theorem}[exact marked selector/terminal-slice equivalence]
\label{thm:es-s6-selector-terminal-bijection}
\Verified\ The map
\[
 \Theta_{\varepsilon,p}:\mathcal W_\varepsilon(p)
 \longrightarrow\mathcal H_\varepsilon(p),
 \qquad
 (j,u)\longmapsto
 \left(j,\frac{T+1}{4},u-jT\right),
 \quad T=\frac{4u+p^\varepsilon}{R_j},
\]
is a bijection.  Under it, the two relation matrices are related by the
explicit unimodular row operation
\[
 \begin{pmatrix}1&-j\\0&1\end{pmatrix}
 \begin{pmatrix}R_j&-u\\4&-T\end{pmatrix}
 =
 \begin{pmatrix}3&-\ell_1\\4&1-4\ell_0\end{pmatrix}
 =\mathsf R(\ell_0,\ell_1,1),
\]
and both determinants equal \(-p^\varepsilon\).
\end{theorem}

\begin{proof}
For \((j,u)\in\mathcal W_\varepsilon(p)\), \(T\) is a positive integer and
the congruence computation in \cref{prop:es-two-channel-selector} gives
\(T\equiv3\pmod4\), so \(\ell_0=(T+1)/4\) is integral.  Set
\(\ell_1=u-jT\).  Since \(R_j=4j+3\),
\[
 \lambda(\ell_0,\ell_1,1)
 =3(T+1)-4(u-jT)-3
 =(4j+3)T-4u=p^\varepsilon.
\]
The positivity and divisor conditions in the definition of
\(\mathcal H_\varepsilon(p)\) follow from \(u=\ell_1+jT\).  Matrix
multiplication gives the displayed identity, and
\[
 \det\begin{pmatrix}R_j&-u\\4&-T\end{pmatrix}
 =-R_jT+4u=-p^\varepsilon.
\]

Conversely, from \((j,\ell_0,\ell_1)\in\mathcal H_\varepsilon(p)\) recover
\[
 T=4\ell_0-1,
 \qquad u=\ell_1+jT.
\]
Expanding \(\lambda(\ell_0,\ell_1,1)=p^\varepsilon\) gives
\(R_jT=4u+p^\varepsilon\).  Thus \(R_j\mid4u+p^\varepsilon\), and the
remaining positivity and divisor hypotheses are already part of
\(\mathcal H_\varepsilon(p)\).  These constructions are literal inverses.
\end{proof}

The normalization \(\ell_2=1\) automatically satisfies the candidate's odd
\(\ell_2\) admissibility condition.  It also forces \(3\nmid\ell_1\): if
\(3\mid\ell_1\), then \(\lambda(\ell_0,\ell_1,1)\) is divisible by \(3\),
whereas \(p^\varepsilon\equiv1\pmod3\).  Write
\(\operatorname{ES}(p)\) for the assertion that \(4/p\) is a sum of three
positive integral unit fractions.  Consequently
\[
 \operatorname{ES}(p)
 \Longleftrightarrow
 \mathcal H_0(p)\sqcup\mathcal H_1(p)\ne\varnothing.
\]
The \(E\)-channel is the slice with terminal parameter
\(p_{\mathrm{S6}}=1\); the \(M\)-channel is the slice with
\(p_{\mathrm{S6}}=p\).

Here the channel label is part of the datum.  The raw selector subsets need
not be disjoint: for \(p=13\), \(j=0\), and \(u=2\), one has
\(u\in E_0\cap M_0\).  Accordingly the occurrence statement above means the
tagged disjoint union
\(\{0\}\times\mathcal H_0(p)\sqcup\{1\}\times\mathcal H_1(p)\).
Moreover the map from tagged selector states to \emph{unordered} unit-fraction
solutions is not injective in the \(M\)-channel: complementation
\(u\leftrightarrow a_j^2/u\) interchanges \(y\) and \(z\).  For example,
\(p=13\), \(j=0\), and \(u=2,8\) give the same unordered solution
\(\{4,26,52\}\).  The theorem is a bijection with the marked terminal slice,
not a bijection with unordered solutions.

The tags are essential.  For every eligible prime \(q\), define the
fixed-prime forgetful map
\[
 \mathfrak f_q:\mathcal H_0(q)\longrightarrow\mathbb Z^3,
 \qquad
 \mathfrak f_q(j,\ell_0,\ell_1)=(\ell_0,\ell_1,1).
\]
For \(p=13\), the two \(E\)-channel selector states satisfy
\[
 \Theta_{0,13}(0,2)=(0,1,2),
 \qquad
 \Theta_{0,13}(1,5)=(1,1,2),
\]
so their marked images are distinct, while both are sent by
\(\mathfrak f_{13}\) to \((1,2,1)\).  Thus forgetting \(j\) destroys
injectivity.  To compare different prime fibres, define
\[
 \mathfrak F:
 \bigsqcup_{\substack{q\ \mathrm{prime}\\q\equiv1\pmod {12}}}
 \bigl(\{q\}\times\mathcal H_0(q)\bigr)
 \longrightarrow\mathbb Z^3,
 \qquad
 \mathfrak F(q;j,\ell_0,\ell_1)=(\ell_0,\ell_1,1).
\]
At \(q=73\), one has \(\Theta_{0,73}(1,5)=(1,1,2)\), so this tagged
image is also sent by \(\mathfrak F\) to \((1,2,1)\).  Hence forgetting
the external prime loses still more information.

\subsection{What the crosswalk proves, and what it does not}

The theorem identifies the complete fixed-\(p\), normalized, positively
oriented, \(j\)-marked divisor selector with an arithmetic slice of the
candidate's terminal relation matrices.  It explains a genuine structural
continuity between the earlier reciprocal-equation work and the terminal
algebra used here: the common object is a \(2\times2\) integral presentation,
and the passage between the two displayed matrices is an element of
\(\operatorname{SL}_2(\mathbb Z)\).

The orientation qualification is substantive.  The geometrically validated
candidate specialization used elsewhere in this workbench has terminal
parameter \(p_{\mathrm{S6}}=-1\), whereas the \(E\)-channel here has
\(p_{\mathrm{S6}}=+1\).  No result in this section identifies those two
oriented terminal data.  Thus the exact arithmetic equivalence does not, by
itself, place an Erd\H{o}s--Straus selector state on the particular complex
threefold reconstructed in the candidate-specific \(p=-1\) branch.

It does not prove that the selector is nonempty for every prime, so it does
not prove the Erd\H{o}s--Straus conjecture.  It also does not construct a map
from the reciprocal cubic surface to \(X\), identify selector divisors with
cycles or attaching maps in the candidate geometry, or transfer the
candidate's complex structure to the arithmetic surface.  Those would be
additional geometric assertions, not consequences of unimodular equivalence
of relation matrices.

% Proof-bearing workbench module.  The standalone root supplies its preamble.
\providecommand{\ZZ}{\mathbb Z}
\providecommand{\ZZpos}{\mathbb Z_{>0}}
\providecommand{\ES}{\operatorname{ES}}
\providecommand{\Rel}{\operatorname{Rel}}
\providecommand{\coker}{\operatorname{coker}}
\providecommand{\eps}{\varepsilon}

\section{The exact marked Erd\H{o}s--Straus--\texorpdfstring{\(S^6\)}{S6} bridge}
\label{sec:esksb-bridge}

\subsection*{Current lattice and modular continuation}
The shared continuation now proves the full cubic image
$2\sqrt2(18\mathbb Z\cup(6\mathbb Z+2)\cup(6\mathbb Z-2))$
of the literal common Niemeier--Leech lattice, with explicit integral
and local sections. The marked $A_5^4D_4$ root system maps to the
Cheng--Duncan--Harvey shadow $(5I_{12}+J)S_6$, where
$J_{rs}=\mathbf1_{s=5r\ (12)}$. The proof calculates the original
glue-group action on twisted shadows, the complete theta determinant
character, and the coefficient map to
$T_6=\eta(\tau)^5\eta(3\tau)/(\eta(2\tau)\eta(6\tau)^5)$, including
$T_6(-1/(6\tau))=72/T_6(\tau)$.
See \cref{d4i:thm:complete-image,nus:sec:shadow}.

\subsection*{Result and scope}
For each prime \(p\equiv1\pmod {12}\), the complete first-half
Erd\H{o}s--Straus divisor selector is in explicit bijection with a
normalized, positively oriented, divisor-supported, and shell-marked slice of
the terminal parameter lattice occurring in the candidate complex structure
on \(S^6\).  The coordinate change is integral in both directions.  More
strongly, the selector relation matrix and the S6 terminal relation matrix
have exactly the same row lattice: they differ by a displayed element of
\(\operatorname{SL}_2(\ZZ)\).  Composing the bijection with the shell
reconstruction gives positive denominators solving \(4/p=1/x+1/y+1/z\), and
every positive solution is detected by at least one of the two tagged slices.

The nonlinear condition \(u\mid a_j^2\), together with the external labels
\((p,j,\eps)\), is indispensable.  If it is removed, each normalized S6 fibre
is an infinite affine line and carries no Erd\H{o}s--Straus existence
information.  A particularly natural unmarked S6 point produces precisely
the already-solved congruence class \(p\equiv13\pmod {24}\) and produces no
selector state when \(p\equiv1\pmod {24}\).  Thus the bridge is an exact
arithmetic equivalence of marked presentations, not a proof of universal
selector nonemptiness and not a geometric morphism from the reciprocal cubic
to the complex threefold.

The separate continuation in Section~\ref{esn:sec:shared} constructs
the index-$72$ modulo-$840$ lattice $A_5^4\perp D_4\subset N$ and
its exact marked $D_4$ triality maps. It realizes the known Weyl quotient
by the coordinate splitting $O(N)=W(A_5^4D_4)\rtimes L(G)$,
$G\cong GL_2(\mathbb F_3)$, and places the residue six-cycle inside
the Weyl group, centralizing this chosen complement
(Theorems~\ref{esn:fa:thm:full-quotient}
and~\ref{esn:fa:thm:residue-weyl}; \cite[Table~3]{esn:CDH2014}).
Its divisor and Fourier maps retain the prime, shell, and finite sample;
in particular the shell count does not descend to the prime residue alone.

\subsection{All objects and conventions}

Fix throughout a prime
\[
 p\equiv1\pmod {12},\qquad
 h=\frac{p-1}{12}\in\ZZpos,
\]
and write
\[
 L=\frac{p-1}{4}=3h,
 \qquad J_p=\{j\in\ZZ:0\leq j<L\}.
\]
For \(j\in J_p\), retain the four integers
\begin{equation}
 R_j=4j+3,
 \qquad a_j=3h+j+1=\frac{p+R_j}{4},
 \qquad N_j=pa_j,
 \qquad 4a_j=p+R_j.                         \tag{1.1}
\end{equation}
No label in \((p,j)\) will be suppressed.  The inequalities
\[
 3\leq R_j\leq p-2,
 \qquad \frac{p+3}{4}\leq a_j\leq\frac{p-1}{2}<p
\]
give
\begin{equation}
 \gcd(p,a_j)=\gcd(p,R_j)=\gcd(a_j,R_j)=1.  \tag{1.2}
\end{equation}
Indeed \(p\) is prime and \(0<a_j<p\); then \(R_j=4a_j-p\)
gives the other two equalities.

Let
\[
 \mathscr S(p)=
 \left\{(x,y,z)\in\ZZpos^3:
 4xyz=p(xy+xz+yz)\right\}.
\]
The equation is symmetric, but whenever completeness is invoked we choose the
ordered representative \(x\leq y\leq z\).
Write \(\ES(p)\) for the assertion \(\mathscr S(p)\neq\varnothing\).

For the S6 terminal parameters put
\begin{equation}
 \lambda(\ell_0,\ell_1,\ell_2)
 =12\ell_0-4\ell_1-3\ell_2                 \tag{1.3}
\end{equation}
and define the abstract terminal group
\begin{equation}
 G(\ell_0,\ell_1,\ell_2)=
 \left\langle c,x,y\ \middle|\
 c\text{ central},\ xy=c^{\ell_0},\
 x^3=c^{\ell_1},\ y^4=c^{\ell_2}
 \right\rangle .                            \tag{1.4}
\end{equation}
After eliminating \(y\), and using the ordered generator basis \((x,c)\),
its ordered row-relation matrix is
\begin{equation}
 \mathsf R(\ell_0,\ell_1,\ell_2)=
 \begin{pmatrix}
 3&-\ell_1\\
 4&\ell_2-4\ell_0
 \end{pmatrix},
 \qquad
 \det\mathsf R=-\lambda(\ell_0,\ell_1,\ell_2).              \tag{1.5}
\end{equation}
The letter \(p\) in this note always denotes the Erd\H{o}s--Straus prime.
The signed terminal value is denoted
\[
 q_\eps=p^\eps\in\{1,p\},\qquad \eps\in\{0,1\}.
\]
Thus no identification of the external prime \(p\) with the terminal value
is made in the \(\eps=0\) channel.

For any integral \(2\times2\) matrix \(A\), define its row-relation lattice
and row-presented group by
\[
 \Rel(A)=\{v^{\mathsf T}A:v\in\ZZ^2\}\subseteq\ZZ^2,
 \qquad \mathcal G(A)=\ZZ^2/\Rel(A)=\coker(A^{\mathsf T}).    \tag{1.6}
\]
This convention makes every matrix direction in the bridge explicit.

\subsection{The complete first-half shell selector}

\begin{lemma}[first-half denominator]\label{lem:esksb-first-half}
Let \(p\equiv1\pmod4\) be an odd prime.  If
\((x,y,z)\in\mathscr S(p)\) and \(x\leq y\leq z\), then
\[
 \frac p4<x\leq\frac p2.
\]
\end{lemma}

\begin{proof}
Positivity of the two reciprocals other than \(1/x\) gives
\(1/x<4/p\), hence \(x>p/4\).  Suppose \(x>p/2\).  Clearing denominators
gives
\[
 4xyz=p(xy+xz+yz),
\]
so \(p\mid xyz\).  If \(p\mid x\), then \(x,y,z\geq p\) and the reciprocal
sum is at most \(3/p\).  If \(p\mid y\), then \(y,z\geq p\), while
\(1/x<2/p\), so the sum is strictly less than \(4/p\).  Both cases are
impossible.  Therefore \(z=kp\) for some \(k\in\ZZpos\), while neither
\(x\) nor \(y\) is divisible by \(p\).  Substitution yields
\[
 (4k-1)xy=kp(x+y).
\]
Reduction modulo \(p\) gives \(p\mid4k-1\).  Write \(4k-1=rp\).
Since \(p\equiv1\pmod4\), one has \(r\equiv3\pmod4\), hence \(r\geq3\).
After cancelling \(p\),
\[
 rxy=k(x+y).
\]
As \(x\leq y\), division by \(xy\) gives
\[
 x\leq\frac{2k}{r}
 =\frac{rp+1}{2r}
 =\frac p2+\frac1{2r}<\frac{p+1}{2}.
\]
The integer \(x>p/2\) must be at least \((p+1)/2\), a contradiction.
\end{proof}

For fixed \(j\in J_p\), define the ordered residual shell
\begin{equation}
 \mathscr P_{p,j}=
 \{(y,z)\in\ZZpos^2:R_jyz=N_j(y+z)\}                         \tag{2.1}
\end{equation}
and the oriented divisor shell
\begin{equation}
 \mathscr D_{p,j}=
 \{d\in\ZZpos:d\mid N_j^2,\ d\equiv-N_j\pmod {R_j}\}.       \tag{2.2}
\end{equation}

\begin{proposition}[oriented shell bijection]\label{prop:esksb-shell-bijection}
The maps
\begin{align*}
 \delta_{p,j}:\mathscr P_{p,j}&\longrightarrow\mathscr D_{p,j},
 &(y,z)&\longmapsto R_jy-N_j,\\
 \rho_{p,j}:\mathscr D_{p,j}&\longrightarrow\mathscr P_{p,j},
 &d&\longmapsto
 \left(\frac{N_j+d}{R_j},
       \frac{N_j+N_j^2/d}{R_j}\right)
\end{align*}
are mutually inverse.  The involution
\[
 \iota_{N_j}(d)=\frac{N_j^2}{d}
\]
corresponds exactly to \((y,z)\mapsto(z,y)\).
\end{proposition}

\begin{proof}
Abbreviate \(R=R_j\) and \(N=N_j\).  Equation~(2.1) is equivalent to
\begin{equation}
 (Ry-N)(Rz-N)=N^2.                                            \tag{2.3}
\end{equation}
Both factors are positive: for example, division of
\(Ryz=N(y+z)\) by \(z\) gives \(Ry=N(1+y/z)>N\), and similarly for
\(Rz-N\).  Therefore \(d=Ry-N\) is a positive divisor of \(N^2\) and
\(d\equiv-N\pmod R\).

Conversely, let \(d\in\mathscr D_{p,j}\) and put \(e=N^2/d\).  By~(1.2),
\(d\) and \(N\) are units modulo \(R\); hence
\[
 e\equiv N^2(-N)^{-1}\equiv-N\pmod R.
\]
Thus both components of \(\rho_{p,j}(d)\) are positive integers.  Their
 two shifted factors are \(d\) and \(e\), so~(2.3), and therefore~(2.1),
 holds.  The formulas visibly compose to the identity.  Replacing \(d\) by
\(e=N^2/d\) exchanges the two reconstructed components.
\end{proof}

For \(\eps\in\{0,1\}\), define the fixed-shell selector and the fully marked
selector by
\begin{align}
 \mathcal W_\eps(p,j)
 &=\{u\in\ZZpos:u\mid a_j^2,\ R_j\mid4u+p^\eps\},             \tag{2.4}\\
 \mathcal W_\eps(p)
 &=\{(j,u):j\in J_p,\ u\in\mathcal W_\eps(p,j)\}.           \tag{2.5}
\end{align}
The tag \(\eps=0\) is the exterior, or \(E\), channel; \(\eps=1\) is the
middle, or \(M\), channel.
Whenever a disjoint union of the two channels is written below, it means the
tagged coproduct
\[
 \{0\}\times\mathcal W_0(p)\ \sqcup\
 \{1\}\times\mathcal W_1(p).
\]
This convention is necessary because the two raw subsets can intersect.

\begin{proposition}[exact two-channel decomposition]\label{prop:esksb-two-channel}
The maps
\begin{align*}
 \iota_E:\mathcal W_0(p,j)&\longrightarrow\mathscr D_{p,j},
 &u&\longmapsto p^2u,\\
 \iota_M:\mathcal W_1(p,j)&\longrightarrow\mathscr D_{p,j},
 &u&\longmapsto pu
\end{align*}
are bijections onto the subsets on which the \(p\)-adic valuations are \(2\)
and \(1\), respectively.  The valuation-zero subset is
\(\iota_{N_j}(\iota_E(\mathcal W_0(p,j)))\).  Consequently
\begin{equation}
 \mathscr D_{p,j}=
 \iota_E(\mathcal W_0(p,j))\ \sqcup\
 \iota_M(\mathcal W_1(p,j))\ \sqcup\
 \iota_{N_j}(\iota_E(\mathcal W_0(p,j))),                    \tag{2.6}
\end{equation}
and
\begin{equation}
 |\mathscr D_{p,j}|=2|\mathcal W_0(p,j)|+|\mathcal W_1(p,j)|. \tag{2.7}
\end{equation}
\end{proposition}

\begin{proof}
Because \(p\nmid a_j\), every divisor \(d\mid p^2a_j^2\) has
\(v_p(d)\in\{0,1,2\}\).

If \(d=p^2u\), then \(u\mid a_j^2\).  The shell congruence is
\[
 R_j\mid d+N_j=p(pu+a_j).
\]
After cancelling the unit \(p\) modulo \(R_j\) and multiplying by \(4\),
the identity \(4a_j=p+R_j\) makes this equivalent to
\(R_j\mid4u+1\).  This proves the \(E\)-channel assertion in both
directions.  If \(d=pu\), the same calculation gives
\[
 R_j\mid d+N_j
 \Longleftrightarrow R_j\mid u+a_j
 \Longleftrightarrow R_j\mid4u+p,
\]
which proves the \(M\)-channel assertion.  Complementation preserves the
shell congruence by Proposition~\ref{prop:esksb-shell-bijection}, exchanges
valuations \(0\) and \(2\), and preserves valuation \(1\).  This proves the
disjoint decomposition and its cardinality formula.
\end{proof}

For \(u\in\mathcal W_\eps(p,j)\), define the two exact quotients
\begin{equation}
 T=\frac{4u+p^\eps}{R_j},
 \qquad K=\frac{p^{1-\eps}T+1}{4}.                            \tag{2.8}
\end{equation}
Since \(R_j\equiv3\pmod4\) and \(p^\eps\equiv1\pmod4\), the first equation
gives \(T\equiv3\pmod4\).  Hence \(K\in\ZZpos\).  Direct calculation,
without dropping the tag, gives
\begin{equation}
 R_jK=a_j+p^{1-\eps}u,
 \qquad Ta_j=u+p^\eps K.                                    \tag{2.9}
\end{equation}
Indeed, multiplication by \(4\) changes these identities into
\begin{align*}
 4R_jK
 &=p^{1-\eps}(4u+p^\eps)+R_j
 =4p^{1-\eps}u+p+R_j,\\
 4Ta_j
 &=T(p+R_j)
 =pT+4u+p^\eps
 =4u+4p^\eps K.
\end{align*}

\begin{proposition}[denominator reconstruction]\label{prop:esksb-reconstruct}
The map
\begin{equation}
 \Psi_{p,j,\eps}:\mathcal W_\eps(p,j)\longrightarrow\mathscr S(p),
 \qquad
 u\longmapsto
 \left(a_j,\ pK,\ \frac{p^\eps a_jK}{u}\right)             \tag{2.10}
\end{equation}
is well-defined.  Its two oriented shell factors are
\begin{equation}
 R_j(pK)-N_j=p^{2-\eps}u,
 \qquad
 R_j\left(\frac{p^\eps a_jK}{u}\right)-N_j
 =p^\eps\frac{a_j^2}{u}.                                    \tag{2.11}
\end{equation}
Moreover,
\begin{equation}
 \ES(p)\quad\Longleftrightarrow\quad
 \bigl(\{0\}\times\mathcal W_0(p)\bigr)\sqcup
 \bigl(\{1\}\times\mathcal W_1(p)\bigr)\neq\varnothing.   \tag{2.12}
\end{equation}
\end{proposition}

\begin{proof}
The last coordinate in~(2.10) is integral.  To see this without assuming it,
reduce the second identity in~(2.9) modulo \(u\), multiply by \(a_j\), and
use \(u\mid a_j^2\): this gives \(u\mid p^\eps a_jK\).
The two identities~(2.11) now follow from the first identity in~(2.9).
Their product is \(N_j^2\), so Proposition~\ref{prop:esksb-shell-bijection}
shows that~(2.10) lies in \(\mathscr S(p)\).

Conversely, suppose \(\mathscr S(p)\neq\varnothing\), and sort a solution.
Lemma~\ref{lem:esksb-first-half} shows that its smallest denominator is uniquely
of the form \(x=a_j\) with \(j\in J_p\): explicitly
\[
 R_j=4x-p\equiv3\pmod4,
 \qquad j=\frac{4x-p-3}{4}.
\]
Proposition~\ref{prop:esksb-shell-bijection} supplies a divisor
\(d\in\mathscr D_{p,j}\), and Proposition~\ref{prop:esksb-two-channel} places
either \(d\) or its complement in the image of \(\iota_E\), or places \(d\)
in the image of \(\iota_M\).  Hence one of the two marked selector sets is
nonempty.  The forward implication follows from~(2.10).
\end{proof}

\subsection{The exact selector--terminal-lattice equivalence}

Define the normalized, positively oriented, divisor-supported terminal slice
\begin{equation}
 \mathcal H_\eps(p)=
 \left\{(j,\ell_0,\ell_1)\in\ZZ^3:
 \begin{array}{l}
 j\in J_p,\quad \lambda(\ell_0,\ell_1,1)=p^\eps,\\
 T(\ell_0):=4\ell_0-1>0,\\
 u(j,\ell_0,\ell_1):=\ell_1+jT(\ell_0)>0,\\
 u(j,\ell_0,\ell_1)\mid a_j^2
 \end{array}
 \right\}.                                                   \tag{3.1}
\end{equation}
Let \(\mathcal H_\eps(p,j)\) be the fibre with the displayed marker \(j\).

For a selector state define its relation matrix
\begin{equation}
 \mathsf A_\eps(p;j,u)=
 \begin{pmatrix}
 R_j&-u\\
 4&-T
 \end{pmatrix},
 \qquad R_jT=4u+p^\eps.                                     \tag{3.2}
\end{equation}
For the marker \(j\), retain both inverse unimodular shears
\begin{equation}
 \mathsf U_j=
 \begin{pmatrix}1&-j\\0&1\end{pmatrix},
 \qquad
 \mathsf L_j=\mathsf U_j^{-1}=
 \begin{pmatrix}1&j\\0&1\end{pmatrix}.                    \tag{3.3}
\end{equation}

\begin{theorem}[strong marked bridge theorem]\label{thm:esksb-bridge}
For every \(\eps\in\{0,1\}\), the map
\begin{equation}
 \Theta_{\eps,p}:\mathcal W_\eps(p)\longrightarrow\mathcal H_\eps(p),
 \qquad
 (j,u)\longmapsto
 \left(j,\frac{T+1}{4},u-jT\right)                          \tag{3.4}
\end{equation}
is a bijection.  Its inverse is
\begin{equation}
 \Theta_{\eps,p}^{-1}(j,\ell_0,\ell_1)
 =\left(j,\ell_1+j(4\ell_0-1)\right).                       \tag{3.5}
\end{equation}
For corresponding points, with
\(\ell=(\ell_0,\ell_1,1)\), one has the exact matrix identities
\begin{equation}
 \boxed{
 \mathsf U_j\mathsf A_\eps(p;j,u)
 =\mathsf R(\ell_0,\ell_1,1),
 \qquad
 \mathsf A_\eps(p;j,u)
 =\mathsf L_j\mathsf R(\ell_0,\ell_1,1).}                  \tag{3.6}
\end{equation}
In particular,
\begin{equation}
 \Rel(\mathsf A_\eps(p;j,u))
 =\Rel(\mathsf R(\ell_0,\ell_1,1)),
 \qquad
 \det\mathsf A_\eps=\det\mathsf R=-p^\eps.                \tag{3.7}
\end{equation}
Thus the bridge intertwines the entire integral row presentation, not only
its determinant.
\end{theorem}

\begin{proof}
Take \((j,u)\in\mathcal W_\eps(p)\).  Section~2 proves that \(T\) is a
positive integer satisfying \(T\equiv3\pmod4\), so
\(\ell_0=(T+1)/4\in\ZZ\).  Put \(\ell_1=u-jT\).  Since
\(R_j=4j+3\),
\begin{align*}
 \lambda(\ell_0,\ell_1,1)
 &=3(T+1)-4(u-jT)-3\\
 &=(4j+3)T-4u=p^\eps.
\end{align*}
The positivity and divisor-support conditions in~(3.1) hold because
\(u=\ell_1+jT\).  Thus~(3.4) is well-defined.

Conversely, take \((j,\ell_0,\ell_1)\in\mathcal H_\eps(p)\), and set
\[
 T=4\ell_0-1,
 \qquad u=\ell_1+jT.
\]
Expanding \(\lambda(\ell_0,\ell_1,1)=p^\eps\) gives
\[
 (4j+3)T-4u=p^\eps,
\]
so \(R_j\mid4u+p^\eps\).  All other selector conditions occur explicitly
in~(3.1).  The two constructions recover \(T,u,\ell_0,\ell_1\) literally,
which proves bijectivity and~(3.5).

Finally, direct multiplication gives
\[
 \begin{pmatrix}1&-j\\0&1\end{pmatrix}
 \begin{pmatrix}4j+3&-u\\4&-T\end{pmatrix}
 =
 \begin{pmatrix}3&-u+jT\\4&-T\end{pmatrix}
 =
 \begin{pmatrix}3&-\ell_1\\4&1-4\ell_0\end{pmatrix}.
\]
Both shears have determinant \(1\), so their integral row operations are
invertible and the row lattices are equal.  The signed determinant is
\(-R_jT+4u=-p^\eps\), proving~(3.7).
\end{proof}

\begin{corollary}[presented groups and the Erd\H{o}s--Straus equation]
For every corresponding selector/S6 point,
\begin{equation}
 \mathcal G(\mathsf A_\eps(p;j,u))
 =\mathcal G(\mathsf R(\ell_0,\ell_1,1))
 \cong G(\ell_0,\ell_1,1)
 \cong\ZZ/p^\eps\ZZ.                                       \tag{3.8}
\end{equation}
Here \(\ZZ/1\ZZ\) is the trivial group.  Moreover the direct map
\begin{equation}
 \widehat\Psi_{\eps,p}:\mathcal H_\eps(p)\longrightarrow\mathscr S(p)
                                                                    \tag{3.9}
\end{equation}
is obtained by~(3.5), (2.8), and
\begin{equation}
 (j,\ell_0,\ell_1)\longmapsto
 \left(a_j,\ pK,\ \frac{p^\eps a_jK}{u}\right),
 \quad
 \begin{cases}
 T=4\ell_0-1,\\
 u=\ell_1+jT,\\
 K=(p^{1-\eps}T+1)/4.
 \end{cases}                                                \tag{3.10}
\end{equation}
It satisfies
\begin{equation}
 \boxed{\quad
 \ES(p)\Longleftrightarrow
 \bigl(\{0\}\times\mathcal H_0(p)\bigr)\sqcup
 \bigl(\{1\}\times\mathcal H_1(p)\bigr)\neq\varnothing.
 \quad}                                                       \tag{3.11}
\end{equation}
\end{corollary}

\begin{proof}
The relation \(xy=c^{\ell_0}\) in~(1.4) gives
\(y=x^{-1}c^{\ell_0}\); since \(c\) is central, \(x\) and \(y\) commute.
The remaining two relations are the rows of~(1.5).  The greatest common
divisor of the four entries of \(\mathsf R\) is \(1\), because it divides
both \(3\) and \(4\).  Its Smith form is therefore
\(\operatorname{diag}(1,p^\eps)\), proving~(3.8).  Formula~(3.10) is the
composition of Theorem~\ref{thm:esksb-bridge} with
Proposition~\ref{prop:esksb-reconstruct}, and~(3.11) follows from~(2.12).
\end{proof}

\begin{remark}[the selector tag and the solution map]
The channel tag is not redundant.  For \(p=13\), \(j=0\), and \(u=2\),
one has \(u\in\mathcal W_0(p,j)\cap\mathcal W_1(p,j)\), but the two tags
give different terminal values and different marked terminal points.  Nor is
the map to unordered denominator triples injective: in the \(M\)-channel at
\(p=13,j=0\), the complementary divisors \(u=2\) and \(u=8=a_0^2/2\)
both give the unordered solution \(\{4,26,52\}\).  The bijection in
Theorem~\ref{thm:esksb-bridge} is a bijection of marked selector states with marked
terminal states, not a bijection with unordered solutions.
\end{remark}

\begin{corollary}[fixed-fibre affine coordinates and the intrinsic gate]
Put
\[
 C_\eps=\frac{p^\eps+3}{4},
 \qquad D_{\eps,j}=j+C_\eps=\frac{R_j+p^\eps}{4}.
\]
Every integral point with \(\ell_2=1\) and
\(\lambda(\ell_0,\ell_1,1)=p^\eps\) is uniquely
\begin{equation}
 (\ell_0,\ell_1,\ell_2)=(n,3n-C_\eps,1),\qquad n\in\ZZ,      \tag{3.12}
\end{equation}
and, for each fixed marker \(j\in J_p\), define
\begin{equation}
 T=4n-1,
 \qquad u=R_jn-D_{\eps,j},
 \qquad K=\frac{p^{1-\eps}T+1}{4}.                          \tag{3.13}
\end{equation}
Then \(T\equiv3\pmod4\), so \(K\in\ZZ\), and the congruence equation
\(R_jT=4u+p^\eps\) holds identically.  On the locus \(T>0\) and
\(u>0\), one has
\begin{equation}
 \gcd(R_j,u)=1,
 \qquad
 u\mid a_j^2\Longleftrightarrow u\mid K^2.                  \tag{3.14}
\end{equation}
Thus the nonlinear divisor gate can be written entirely in the S6
coordinates, but it does not disappear.
\end{corollary}

\begin{proof}
Equation~(3.12) is the direct solution of
\(12\ell_0-4\ell_1-3=p^\eps\).  Substitution gives the first two
formulas in~(3.13), and \(T\equiv3\pmod4\) proves that the displayed
integer \(K\) is well defined.  The definitions give
\[
 4u=4R_jn-(R_j+p^\eps)=R_jT-p^\eps,
\]
so the congruence equation holds for every integral point in~(3.12).
If a positive integer divides both \(R_j\) and \(u\), then the congruence
equation makes it divide \(p^\eps\); (1.2) therefore gives
\(\gcd(R_j,u)=1\).  Without invoking the selector hypothesis, direct
substitution gives
\[
 4\bigl(R_jK-p^{1-\eps}u-a_j\bigr)
 =p^{1-\eps}(R_jT-4u)-p=0.
\]
Thus \(a_j\equiv R_jK\pmod u\).  Since \(R_j\) is a unit modulo \(u\),
\(a_j^2\equiv R_j^2K^2\pmod u\), which proves (3.14) in both directions.
\end{proof}

The normalization automatically lies in the admissible residue class used by
the S6 terminal presentation:
\begin{equation}
 \ell_2=1\text{ is odd},
 \qquad \ell_1\equiv2\pmod3.                                \tag{3.15}
\end{equation}
Indeed reduction of \(\lambda(\ell_0,\ell_1,1)=p^\eps\equiv1\pmod3\)
gives \(-\ell_1\equiv1\pmod3\).

\begin{remark}[orientation and the constructed S6 candidate]
The geometrically reconstructed candidate elsewhere in the audited corpus
uses the oriented terminal parameter
\(\lambda(0,1,-1)=-1\).  The two arithmetic channels in this note use
the positive terminal values \(+1\) and \(+p\).  No audited map identifies
the \(-1\) terminal datum with either positive datum.  Consequently the row
shear proved here does not, by itself, place any selector state on the
particular complex threefold constructed in the \(-1\) branch.
\end{remark}

\subsection{The exact remaining arithmetic incidence}

The bridge can be expressed without an untyped phrase such as ``some divisor
of a square.''  The next factor coordinates are canonical and retain their
orientation.

\begin{lemma}[canonical ordered factor coordinates]\label{lem:esksb-factors}
For \(a,u\in\ZZpos\), the condition \(u\mid a^2\) is equivalent to the
existence of a unique ordered triple \((r,s,t)\in\ZZpos^3\) such that
\begin{equation}
 u=r^2s,
 \qquad a=rst,
 \qquad \gcd(r,t)=1.                                        \tag{4.1}
\end{equation}
Explicitly, with \(g=\gcd(a,u)\),
\begin{equation}
 r=\frac ug,
 \qquad s=\frac{g^2}{u},
 \qquad t=\frac ag.                                         \tag{4.2}
\end{equation}
\end{lemma}

\begin{proof}
Write \(v_q(a)=\alpha\) and \(v_q(u)=\beta\leq2\alpha\) prime by prime.
Then
\[
 v_q(g)=\min(\alpha,\beta),
 \qquad
 v_q(g^2/u)=2\min(\alpha,\beta)-\beta\geq0.
\]
Thus~(4.2) consists of positive integers.  If \(\beta\leq\alpha\), then
\(q\nmid r\); if \(\beta\geq\alpha\), then \(q\nmid t\).  Hence
\(\gcd(r,t)=1\), and direct multiplication gives~(4.1).

Conversely,~(4.1) gives
\(a^2/u=s t^2\in\ZZpos\).  It also gives
\(\gcd(a,u)=rs\), because \(\gcd(r,t)=1\).  Formula (4.2) then recovers
\((r,s,t)\), proving uniqueness.
\end{proof}

\begin{proposition}[factor form of both selector channels]\label{prop:esksb-factor-congruence}
For the unique factorization
\(u=r^2s\), \(a_j=rst\), \(\gcd(r,t)=1\), one has
\begin{align}
 u\in\mathcal W_0(p,j)
 &\Longleftrightarrow R_j\mid pr+t,                           \tag{4.3}\\
 u\in\mathcal W_1(p,j)
 &\Longleftrightarrow R_j\mid r+t.                            \tag{4.4}
\end{align}
\end{proposition}

\begin{proof}
All of \(r,s,t\) are units modulo \(R_j\), because their product is \(a_j\)
and \(\gcd(a_j,R_j)=1\).  For the \(E\)-channel,
\[
 t(4u+1)=4r^2st+t=(p+R_j)r+t=(pr+t)+R_jr.
\]
Cancelling \(t\) modulo \(R_j\) proves~(4.3).  For the \(M\)-channel,
\[
 4u+p\equiv4(u+a_j)=4rs(r+t)\pmod {R_j},
\]
and cancellation of the unit \(4rs\) proves~(4.4).
\end{proof}

\begin{theorem}[necessary-and-sufficient counterexample criterion]
A prime \(p\equiv1\pmod {12}\) is a counterexample to the
Erd\H{o}s--Straus equation if and only if, for every \(j\in J_p\) and every
ordered \((r,s,t)\in\ZZpos^3\) satisfying
\[
 a_j=rst,
 \qquad \gcd(r,t)=1,
\]
one has simultaneously
\begin{equation}
 R_j\nmid pr+t,
 \qquad R_j\nmid r+t.                                       \tag{4.5}
\end{equation}
Equivalently, with the tag retained, it is necessary and sufficient that
\begin{equation}
 R_j\nmid p^{1-\eps}r+t                                     \tag{4.6}
\end{equation}
for every \(j\in J_p\), every \(\eps\in\{0,1\}\), and every such factor
triple.
\end{theorem}

\begin{proof}
Lemma~\ref{lem:esksb-factors} is a bijection between all allowed divisors in
the selector~(2.4) and all displayed ordered factor triples.  Proposition
\ref{prop:esksb-factor-congruence} translates the two selector congruences without
loss.  The assertion now follows directly from~(2.12), or equivalently
(3.11), by negating the finite existential statement.
\end{proof}

This theorem isolates the unresolved problem exactly: the S6 row shear has
not removed the bounded factor-incidence condition~(4.5).

\subsection{Why the labels and support cannot be forgotten}

\begin{lemma}[full terminal fibres are affine lattices]
The kernel of \(\lambda:\ZZ^3\to\ZZ\) is
\begin{equation}
 \ker\lambda=\ZZ(1,3,0)\oplus\ZZ(1,0,4).                    \tag{5.1}
\end{equation}
If \(q\equiv1\pmod {12}\) and \(m_q=(q+11)/12\), then
\begin{equation}
 \lambda^{-1}(q)=
 \{(m_q+r+s,\ 2+3r,\ 1+4s):r,s\in\ZZ\}.                   \tag{5.2}
\end{equation}
In the normalized slice \(\ell_2=1\), this is an infinite affine line.
\end{lemma}

\begin{proof}
If \(12a-4b-3c=0\), reduction modulo \(4\) gives \(c=4s\).  The remaining
equation gives \(b=3(a-s)\); putting \(r=a-s\) yields
\((a,b,c)=r(1,3,0)+s(1,0,4)\), uniquely.  Since
\(\lambda(m_q,2,1)=12m_q-11=q\), translation by~(5.1) proves~(5.2).
\end{proof}

The lemma proves that the mere existence of a terminal parameter with signed
value \(1\) or \(p\) is automatic.  It cannot imply selector nonemptiness.
The following examples show separately why the marker and the divisor-support
condition are both part of the theorem's domain.

\begin{example}[forgetting \(j\) and \(p\) is noninjective]\label{ex:esksb-tags}
For each eligible prime \(q\), define the fixed-prime forgetful map
\[
 \mathfrak f_q:\mathcal H_0(q)\longrightarrow\ZZ^3,
 \qquad
 \mathfrak f_q(j,\ell_0,\ell_1)=(\ell_0,\ell_1,1).
\]
At \(p=13\), the two \(E\)-states
\[
 (j,u,T)=(0,2,3),\qquad(1,5,3)
\]
satisfy
\[
 \Theta_{0,13}(0,2)=(0,1,2),
 \qquad
 \Theta_{0,13}(1,5)=(1,1,2),
\]
and both images are sent by \(\mathfrak f_{13}\) to
\((1,2,1)\).  To compare different prime fibres, form the tagged coproduct
and the global forgetful map
\[
 \mathfrak F:
 \bigsqcup_{\substack{q\ \text{prime}\\q\equiv1\pmod {12}}}
 \bigl(\{q\}\times\mathcal H_0(q)\bigr)
 \longrightarrow\ZZ^3,
 \qquad
 \mathfrak F(q;j,\ell_0,\ell_1)=(\ell_0,\ell_1,1).
\]
The state \((p,j,u)=(73,1,5)\) has
\(\Theta_{0,73}(1,5)=(1,1,2)\), so its tagged image also maps under
\(\mathfrak F\) to \((1,2,1)\).  Therefore neither \(j\) nor the external
prime \(p\) is recoverable after applying the explicitly displayed
forgetful maps.
\end{example}

\begin{example}[ambient positivity does not imply divisor support]\label{ex:esksb-ambient}
For \(p=13\), \(j=2\), one has \(R_j=11\) and \(a_j=6\).  The positive
normalized \(E\)-line has
\[
 (\ell_0,\ell_1)=(n,3n-1),
 \qquad u=11n-3,
\]
and the positive normalized \(M\)-line has
\[
 (\ell_0,\ell_1)=(n,3n-4),
 \qquad u=11n-6.
\]
Both lines are infinite.  The only positive \(u\leq a_j^2=36\) are
\[
 E:\ 8,19,30,
 \qquad M:\ 5,16,27,
\]
and none divides \(36\).  Hence
\(\mathcal H_0(13,2)=\mathcal H_1(13,2)=\varnothing\), even though both
ambient positive S6 lines are nonempty.  The prime \(13\) is solved at other
markers, so this is a fixed-shell counterexample to forgetting support, not
an Erd\H{o}s--Straus counterexample.
\end{example}

\subsection{The natural unmarked \texorpdfstring{\(S^6\)}{S6} point and its complete divisor gate}

The most direct attempted use of the S6 presentation starts from one point in
the signed \(p\)-fibre and shears it through every shell.  Its exact content
can be classified completely.

Write
\begin{equation}
 p=12m-11,
 \qquad m=h+1\geq2,
 \qquad J_p=\{0,\ldots,3m-4\}.                              \tag{6.1}
\end{equation}
Retain the unmarked S6 point
\begin{equation}
 \ell^{\mathrm{aff}}=(m,2,1),
 \qquad \lambda(\ell^{\mathrm{aff}})=p,                     \tag{6.2}
\end{equation}
and for \(j\in J_p\) put
\begin{equation}
 T_0=4m-1,
 \quad u_j=2+jT_0,
 \quad a_j=3m+j-2,
 \quad R_j=4j+3.                                             \tag{6.3}
\end{equation}
Then
\begin{equation}
 R_jT_0=4u_j+p,
 \qquad R_jm=a_j+u_j,
 \qquad T_0a_j-u_j=pm,                                      \tag{6.4}
\end{equation}
and
\begin{equation}
 \mathsf U_j
 \begin{pmatrix}R_j&-u_j\\4&-T_0\end{pmatrix}
 =
 \begin{pmatrix}3&-2\\4&1-4m\end{pmatrix}
 =\mathsf R(m,2,1).                                         \tag{6.5}
\end{equation}

\begin{theorem}[complete affine-line gate]\label{thm:esksb-affine-gate}
For prime \(p=12m-11\), \(m\geq2\), and \(0\leq j<3m-3\),
\begin{equation}
 u_j\mid a_j^2
 \Longleftrightarrow
 \bigl(j=0\text{ and }2\mid m\bigr)
 \ \text{or}\ 
 \bigl(j=2\text{ and }8\mid m\bigr).                       \tag{6.6}
\end{equation}
Consequently the single unmarked point~(6.2), with all markers restored,
produces an Erd\H{o}s--Straus witness if and only if
\(p\equiv13\pmod {24}\).  In that case \(j=0,u=2,K=m\) gives
\begin{equation}
 \frac4p=
 \frac1{3m-2}
 +\frac1{pm}
 +\frac{1}{pm(3m-2)/2}.                                    \tag{6.7}
\end{equation}
For \(p\equiv1\pmod {24}\), this affine S6 point supplies no supported
selector state at any \(j\in J_p\).
\end{theorem}

\begin{proof}
The identities~(6.4) and~(6.5) follow by direct expansion.  Assume first
that \(u_j\mid a_j^2\), and put \(g=\gcd(u_j,a_j)\).  Write
\(u_j=gu'\) and \(a_j=ga'\) with \(\gcd(u',a')=1\).  The assumed
divisibility gives \(u'\mid g(a')^2\), hence \(u'\mid g\), and therefore
\begin{equation}
 u_j\mid g^2.                                                \tag{6.8}
\end{equation}
The last identity in~(6.4) gives \(g\mid pm\).  Since \(0<a_j<p\), the
primality of \(p\) and \(g\mid a_j\) imply \(p\nmid g\); therefore
\(g\mid m\).  Reducing \(u_j=2+j(4m-1)\) modulo \(m\) now gives
\begin{equation}
 g\mid j-2.                                                  \tag{6.9}
\end{equation}

Suppose \(j\neq0,2\).  Then \(j\geq1\), and~(6.8)--(6.9) imply
\(u_j\leq g^2\leq(j-2)^2\).  But
\begin{equation}
 u_j-(j-2)^2=j(4m+3-j)-2>0                                  \tag{6.10}
\end{equation}
For \(m\geq2\) and \(1\leq j<3m-3\), indeed
\(4m+3-j\geq m+7\), so the right-hand side of~(6.10) is at least
\(m+5>0\).  This is a contradiction.  Thus only \(j=0\) or \(j=2\)
can occur.

At \(j=0\), \(u_0=2\) and \(a_0=3m-2\), so
\(u_0\mid a_0^2\) exactly when \(m\) is even.  At \(j=2\),
\(u_2=8m\) and \(a_2=3m\), so
\[
 u_2\mid a_2^2
 \Longleftrightarrow 8m\mid9m^2
 \Longleftrightarrow 8\mid m.
\]
These calculations also prove both converse implications in~(6.6).

Now \(p=12m-11\) is congruent to \(13\pmod {24}\) exactly when \(m\) is
even.  The \(j=0\) state lies in the \(M\)-channel, and~(2.8) gives
\(T=T_0\) and \(K=(T_0+1)/4=m\).  Formula~(2.10) is precisely~(6.7).
If \(m\) is odd,~(6.6) excludes every marker on this particular S6 point.
\end{proof}

The last sentence is deliberately local to~(6.2).  For example, \(p=73\)
has \(m=7\) odd, so the affine point contributes nothing; nevertheless
\((j,u,T;\ell_0,\ell_1)=(1,1,11;3,-10)\) is an \(M\)-state.  Thus failure of
the affine point is not failure of the full selector.

\subsection{Exact descendant maps transported through the bridge}

The bridge also transports a verified selector witness along multiplicative
descendant rays.  This is a map from an already nonempty domain, not a
nonemptiness theorem.

Fix a supported state at \((p,j)\), retain \(R=R_j\), and let \(b\in\ZZpos\).
Define
\begin{equation}
 a_b=ba_j,
 \qquad p_b=4a_b-R=bp+(b-1)R.                               \tag{7.1}
\end{equation}
Whenever \(p_b\) is again a prime congruent to \(1\pmod {12}\), the same
marker \(j\) belongs to \(J_{p_b}\), since \(R<p\leq p_b\), and the following
maps have the displayed domains and codomains:
\begin{align}
 \Phi_{b;p}^E:\mathcal H_0(p,j)&\longrightarrow\mathcal H_0(p_b,j),
 &(j,\ell_0,\ell_1)&\longmapsto(j,\ell_0,\ell_1),             \tag{7.2}\\
 \Phi_{b;p}^M:\mathcal H_1(p,j)&\longrightarrow\mathcal H_1(p_b,j),
 &(j,\ell_0,\ell_1)&\longmapsto
 (j,b\ell_0,b\ell_1-j(b-1)).                               \tag{7.3}
\end{align}
Write
\[
 \theta_{\eps,p,j}:\mathcal W_\eps(p,j)\longrightarrow
 \mathcal H_\eps(p,j),
 \qquad
 \theta_{\eps,p,j}(u)=\Theta_{\eps,p}(j,u)
\]
for the fixed-marker restriction of the marked bridge.

\begin{proposition}[intertwining of descendant actions]
The selector maps corresponding to~(7.2) and~(7.3) are the typed maps
\begin{align*}
 \psi_{b;p}^E:\mathcal W_0(p,j)&\longrightarrow\mathcal W_0(p_b,j),
 &u&\longmapsto u,\\
 \psi_{b;p}^M:\mathcal W_1(p,j)&\longrightarrow\mathcal W_1(p_b,j),
 &u&\longmapsto bu.
\end{align*}
They intertwine the fixed-marker bridges exactly:
\begin{equation}
 \Phi_{b;p}^\bullet\circ\theta_{\eps,p,j}
 =\theta_{\eps,p_b,j}\circ\psi_{b;p}^\bullet,
 \qquad
 (\bullet,\eps)\in\{(E,0),(M,1)\}.                         \tag{7.4}
\end{equation}
Equivalently,
\[
 \theta_{0,p_b,j}^{-1}\Phi_{b;p}^E\theta_{0,p,j}(u)=u,
 \qquad
 \theta_{1,p_b,j}^{-1}\Phi_{b;p}^M\theta_{1,p,j}(u)=bu.
\]
On the accompanying arithmetic metadata these maps replace
\((a_j,p,u)\) by \((a_b,p_b,u)\) and \((a_b,p_b,bu)\), respectively.
For the latter,
\begin{equation}
 T_b=bT+b-1,
 \qquad u_b=bu,
 \qquad K_b=bK,                                              \tag{7.5}
\end{equation}
and
\begin{equation}
 \lambda(b\ell_0,b\ell_1-j(b-1),1)=p_b.                    \tag{7.6}
\end{equation}
For either tag the maps are injective.  For \(c\in\ZZpos\), put
\[
 a_{bc}=bc\,a_j,
 \qquad
 p_{bc}=4a_{bc}-R=bc\,p+(bc-1)R.
\]
Whenever \(p_b\) and \(p_{bc}\) are primes congruent to \(1\pmod {12}\),
the displayed maps are composable and satisfy
\begin{equation}
 \Phi_{c;p_b}^\bullet\circ\Phi_{b;p}^\bullet
 =\Phi_{bc;p}^\bullet                                      \tag{7.7}
\end{equation}
\end{proposition}

\begin{proof}
In the \(E\)-channel, \(u\mid a_j^2\) implies \(u\mid a_b^2\), and the
congruence \(R\mid4u+1\) and quotient \(T\) are unchanged.  Hence
\(\ell_0,\ell_1\) are unchanged by~(3.4).

In the \(M\)-channel, \(bu\mid b^2a_j^2=a_b^2\), and
\[
 4bu+p_b=b(4u+p)+(b-1)R.
\]
Thus the descendant is a selector state and its quotient is
\(T_b=bT+b-1\).  Since \(K=(T+1)/4\) in this channel,
\(K_b=(T_b+1)/4=bK\).  Applying~(3.4) gives
\[
 \ell_{0,b}=\frac{T_b+1}{4}=b\ell_0,
 \qquad
 \ell_{1,b}=bu-jT_b=b\ell_1-j(b-1),
\]
which proves~(7.3)--(7.6), including the intertwining identity~(7.4).
Each map is visibly injective for fixed \(b\).  Substitution into
~(7.1)--(7.3), with the source prime retained, proves the typed composition
law~(7.7).
\end{proof}

Reverse descent need not preserve divisor support: the factor \(b\) can
introduce divisors absent from a proposed parent.  Therefore~(7.2)--(7.7)
propagate witnesses but cannot create the first witness on an empty ray.

\subsection{Sharp verdict}

The strongest theorem in this selector/terminal section is
Theorem~\ref{thm:esksb-bridge} together with the exhaustive equivalence~(3.11):
the complete tagged Erd\H{o}s--Straus selector for a fixed prime
\(p\equiv1\pmod {12}\) is exactly the marked divisor-supported S6 terminal
slice, and the two integral relation matrices have the same row lattice.
The exact remaining obstruction is~(4.5).

Three stronger readings are false or unsupported.

\begin{enumerate}
\item The abstract group \(G(\ell)\), its determinant, or an unmarked
terminal triple does not determine selector support.  Every \(E\)-state has
the same trivial presented group, every \(M\)-state has group \(\ZZ/p\ZZ\),
and Examples~\ref{ex:esksb-tags} and~\ref{ex:esksb-ambient} exhibit the information lost by forgetting
the marker and divisor gate.

\item The row shear is not a geometric map from the reciprocal cubic
\(4xyz=p(xy+xz+yz)\) to the candidate complex threefold \(X\).  No map of
varieties or analytic spaces, no cycle map, no attaching map, and no
cohomological correspondence is defined by left multiplication of a
relation matrix.

\item The existence of the full S6 affine fibre cannot prove universal
Erd\H{o}s--Straus positivity.  Those fibres are always nonempty by~(5.2),
whereas selector nonemptiness is exactly the additional finite nonlinear
incidence in~(4.5).  The natural affine point is completely exhausted by
Theorem~\ref{thm:esksb-affine-gate} and contributes nothing to the unresolved
\(p\equiv1\pmod {24}\) half.
\end{enumerate}

Accordingly, the bridge is a genuine theorem and is stronger than a shared
determinant or a heuristic analogy.  It is also exact enough to prove its own
boundary: without a new argument forcing one of the bounded factor
congruences in~(4.5), it neither proves nor disproves the Erd\H{o}s--Straus
conjecture.

The additional lattice and shell theorems in
Section~\ref{esn:sec:shared} supply further exact morphisms without
changing this incidence question. In particular,
Theorem~\ref{esn:thm:non-descent} gives $A_3(2521)=0$ and
$A_3(3361)=6$ although both primes are $1$ modulo $840$;
Theorem~\ref{esn:thm:sample-fourier} retains the actual counts in the
integral augmentation vector $Q_{R,P}=6W_{R,P}-M_{R,P}u\in N$.

\section{Divisor channels, density one, and two distinct arithmetic maps}
\label{sec:esv3-audit}

This dated continuation (4 September 2026) checks the supplied Pro-mode
\emph{Divisor Channels, Prime Saturation, and B\'ezout Openings}, version 3,
and extends its final arithmetic comparison. The original packet is retained
unchanged under \path{supporting_materials/workbench/evidence/koide_es_v3_audit_20260904/source/}.
The claims below concern explicit reciprocal equations, integral matrices,
and quadratic surfaces. They do not require a complex structure on a sphere.

\subsection{The complete oriented channel map}

Fix integers $n\equiv1\pmod4$, $R\equiv3\pmod4$, $0<R<n$,
$\gcd(n,R)=1$, and put $a=(n+R)/4$, $N=na$. Then
$\gcd(n,a)=\gcd(R,a)=1$. For a positive divisor $\delta\mid n^2$ define
\[
 \mathcal C_\delta(n,R)=\{(u,v)\in\mathbb Z_{>0}^2:
 uv=a^2,\ R\mid4u+\delta\}.
\]
\begin{proposition}\label{prop:esv3-channel}
The following is a bijection onto all ordered positive solutions of
$4/n=1/a+1/y+1/z$:
\[
 (\delta,u,v)\longmapsto
 \left(y,z\right)=
 \left(\frac{N+(n^2/\delta)u}{R},\frac{N+\delta v}{R}\right),
 \qquad (u,v)\in\mathcal C_\delta(n,R).
\]
Its inverse, retaining this orientation, is
\[
 d=Ry-N,\quad e=Rz-N,\quad k=\gcd(d,n^2),\quad
 \delta=n^2/k,\quad u=d/k,\quad v=e/\delta.
\]
Exchange of $y,z$ is the fixed-point-free involution
$(\delta,u,v)\mapsto(n^2/\delta,v,u)$.
\end{proposition}
\begin{proof}
The reciprocal equation is $Ryz=N(y+z)$, hence $de=N^2$.
For any positive solution $Ry=N+Ny/z>N$, so $d>0$; similarly $e>0$.
Because $n$ and $a$ are coprime, every prime dividing $n$ occurs in $de$
with precisely its exponent in $n^2$. Thus $k$ extracts the entire
$n$-primary part of $d$, $\delta$ the entire $n$-primary part of $e$,
and $u,v$ are positive integers with $uv=a^2$.
Reducing $N+e$ modulo $R$ gives
$0\equiv na+\delta v\equiv4a^2+\delta v=v(4u+\delta)$.
Since $v$ is a unit modulo $R$, this proves the channel condition.
Conversely that condition makes $N+e$ divisible by $R$; also
$\delta(N+d)=n(\delta a+nu)\equiv na(\delta+4u)\equiv0\pmod R$.
The unit $\delta$ can be cancelled modulo $R$, proving integrality of $y$.
The identity $de=N^2$ proves the reciprocal equation. These operations are
inverse. A fixed point of the involution would have $\delta=n$, $u=v=a$,
so $R\mid4a+n=2n+R$, impossible since $R\ge3$ is odd and coprime to $n$.
\end{proof}

The package therefore corrects the earlier orientation reversal: the channel
order is $n^2/\gcd(Ry-na,n^2)$, not $\gcd(Ry-na,n^2)$.
For $n=9$, $R=7$, $a=4$, $(\delta,u,v)=(3,1,16)$, the ordered
solution is $(y,z)=(9,12)$, and $\gcd(27,81)=27$.

\subsection{The integral presentation, with its exact realization boundary}

Put $T=(4u+\delta)/R$, $j=(R-3)/4$, and
$\sigma=1$ for $\delta\equiv1\pmod4$, $\sigma=-1$ otherwise. Define
\[
 \ell_0=(T+\sigma)/4,\quad \ell_1=u-jT=(3T-\delta)/4,
 \quad \ell_2=\sigma,\quad
 \mathsf R_\delta=\begin{pmatrix}3&-\ell_1\\4&-T\end{pmatrix}.
\]
All coordinates are integral because $T\equiv-\delta\pmod4$.
Direct multiplication and subtraction give
\[
 \begin{pmatrix}1&-j\\0&1\end{pmatrix}
 \begin{pmatrix}R&-u\\4&-T\end{pmatrix}=\mathsf R_\delta,
 \qquad12\ell_0-4\ell_1-3\ell_2=\delta.
\]
The entries $3,4$ are coprime, so
$\operatorname{SNF}(\mathsf R_\delta)=\operatorname{diag}(1,\delta)$;
in particular its cokernel is $\mathbb Z/\delta\mathbb Z$.
Reducing $4\ell_1=3T-\delta$ modulo three gives
$\ell_1\equiv-\delta\pmod3$. Thus the source gluing admissibility
$3\nmid\ell_1$ fails precisely when $3\mid\delta$.
For the example above $(\ell_0,\ell_1,\ell_2)=(0,0,-1)$.
The matrix and its cokernel remain valid, but assigning the source's global
manifold homology to this nonadmissible triple is not established by this
calculation. This is an explicit admissibility test, not a failure of the
reciprocal solution or the integral matrix morphism.

\subsection{An exact density-one statement and what it leaves out}

\begin{proposition}\label{prop:esv3-diagonal}
Let $p\equiv1\pmod4$ be prime and $m=(p+1)/2$. The exterior channel
states on $u=a=v$ are indexed exactly by
\[
 \{R:R\mid m,\ R\equiv3\pmod4\}.
\]
For each such $R$, set $x=(p+R)/4$ and $k=(p+1)/R$. Then
\[
 \frac4p=\frac1x+\frac1{xk}+\frac1{pxk},\qquad
 \nu(p)=\frac12\left(\tau(m)-\sum_{d\mid m}\chi_4(d)\right),
\]
where $\chi_4(d)=1$ for $d\equiv1\pmod4$ and $-1$ for $d\equiv3\pmod4$.
The count is zero exactly when every prime divisor of $m$ is $1\pmod4$.
\end{proposition}
\begin{proof}
For $\delta=1$ and $u=a$, the condition is $R\mid4a+1=p+R+1$,
equivalently $R\mid p+1$. Since $R$ is odd this is $R\mid m$.
Every such $R$ obeys $3\le R\le(p+1)/2<p$, and $k\ge2$.
Adding the final two fractions gives $(p+1)/(pxk)=R/(px)$,
and adding $1/x$ gives $(p+R)/(px)=4/p$.
Distinct $R$ give distinct $x$. Summing the indicator
$(1-\chi_4(d))/2$ over divisors of $m$ proves the count.
A divisor of residue three exists exactly when a prime factor of residue
three exists: the converse follows because products of residue-one primes
and their powers have residue one.
\end{proof}

\begin{proposition}\label{prop:esv3-density}
For fixed $4\mid M$, $\gcd(c,M)=1$, $c\equiv1\pmod4$, the preceding
diagonal states cover relative density one of primes $p\equiv c\pmod M$.
For a finite set $\mathcal Q$ of primes $q\equiv3\pmod4$, $q\nmid M$,
the subfamily with some $q\mid p+1$ has exact relative density
\[
 1-\prod_{q\in\mathcal Q}\frac{q-2}{q-1}.
\]
\end{proposition}
\begin{proof}
The Chinese remainder theorem gives $\prod(q-1)$ reduced residue classes
modulo $M\prod q$ above the specified class modulo $M$. Avoiding
$p\equiv-1\pmod q$ at every $q$ leaves exactly $\prod(q-2)$ of them.
The prime number theorem for each fixed arithmetic progression gives the
displayed ratio. For completeness, its modulus-four case and partial
summation give
\[
 \sum_{\substack{q\le Y\\q\equiv3\ (4)}}\frac1q
 =\frac{\pi(Y;4,3)}Y+\int_2^Y\frac{\pi(t;4,3)}{t^2}\,dt
 \sim\tfrac12\log\log Y.
\]
Removing finitely many primes dividing $M$ does not change divergence.
Since $\log(1-z)\le-z$ for $0<z<1$, the finite products tend to zero.
The full family's lower relative density is at least each of these finite
union densities, hence is at least one and therefore equals one.
No interchange of two unproved asymptotic limits is needed.
\end{proof}

These analytic tools are classical. Elsholtz and Tao,
\emph{Counting the number of solutions to the Erd\H{o}s--Straus equation on
unit fractions}, J. Aust. Math. Soc. 94 (2013), 50--105,
\href{https://arxiv.org/abs/1107.1010v6}{arXiv:1107.1010v6}, Appendix A,
Lemma A.3, give a stronger shifted-prime additive-function statement,
attributed there to Barban via Elliott. Choosing its prime weight to be
the indicator of residue three gives the relevant abundance of factors.
The finite-CRT proof above records the precise chosen ray and progression;
it is not a claim that almost-all solvability is new.

For $p=73$, $(p+1)/2=37$ has no residue-three prime factor. This prime
therefore misses the diagonal ray, although
\[
 \frac4{73}=\frac1{20}+\frac1{220}+\frac1{4015}.
\]
Here $R=7$, $u=5$, $v=80$, and the primitive factor ray is $[1:4]$.
Thus a prime missed by the selected ray need not be missed by the full
selector. Density one does not settle the universal Erd\H{o}s--Straus claim.

\subsection{The B\'ezout opening, with the positivity hypothesis retained}

For coprime positive integers $r,t$, choose $x,y\in\mathbb Z$ with
$ry-tx=1$, set $v=(r,t)^{\mathsf T}$, $w=(x,y)^{\mathsf T}$, and let
\[
 B=(v\ w)\in\operatorname{SL}_2(\mathbb Z),\qquad
 A=\delta vv^{\mathsf T}-ww^{\mathsf T}
 =B\operatorname{diag}(\delta,-1)B^{\mathsf T}.
\]
Then $\det A=-\delta$ for every real $\delta$. All companions are
$w+kv$ for $k\in\mathbb Z$: subtraction of two B\'ezout equations gives
$r\Delta y=t\Delta x$, so coprimality gives $(\Delta x,\Delta y)=k(r,t)$.
Choosing $|x|\le r/2$ retains an actual companion rather than quotienting
away its dependence. For a positive integer channel order,
\[
 A_{11}=\delta r^2-x^2\ge(\delta-1/4)r^2>0.
\]
The broader wording ``every $\delta>0$'' does not imply that inequality.
At $(r,t,\delta)=(2,1,1/8)$ all companions have $x$ odd, and
$A_{11}=1/2-x^2\le-1/2$. For a specified real order and companion the
exact necessary and sufficient condition is $\delta>(x/r)^2$.
This correction leaves every positive integral divisor channel intact.

With $A=\left(\begin{smallmatrix}a&b\\b&c\end{smallmatrix}\right)$ and
$a>0$, the periods $(\tau,\mu,\beta)=(ia,ib/\sqrt6,ic)$ give
\[
 P_A=\begin{pmatrix}
 0&0&1&0\\ \sqrt6 b&a&0&0\\0&0&0&1\\c&b/\sqrt6&0&0
 \end{pmatrix},\qquad
 \det P_A=ac-b^2=-\delta.
\]
This invertible real matrix sends $\mathbb Z^4$ to a discrete cocompact
lattice in $\mathbb C^2$ and therefore defines a complex two-torus.
The determinant sign and the factor $\sqrt6$ are retained.
It does not determine a point on the particular global equivariant period
curve of the proposed threefold. That incidence problem remains a distinct
map problem, not an obstruction to the torus just constructed.

There is also an explicit integral equivalence of presentations, not merely
an equality of determinants. Set
\[
 P_0=\begin{pmatrix}-1&1\\-4&3\end{pmatrix},\quad
 Q_0=\begin{pmatrix}1&T-\ell_1\\0&1\end{pmatrix},\quad
 J=\begin{pmatrix}0&1\\1&0\end{pmatrix}.
\]
Multiplication gives $P_0\mathsf R_\delta Q_0=\operatorname{diag}(1,-\delta)$
and $(-JB^{-1})A(B^{-\mathsf T}J)=\operatorname{diag}(1,-\delta)$.
Thus, for $L=P_0^{-1}(-JB^{-1})$, $V=B^{-\mathsf T}JQ_0^{-1}$,
\[
 \mathsf R_\delta=LAV.
\]
As maps $\mathbb Z^2\to\mathbb Z^2$, the commuting square has source
isomorphism $V^{-1}$ and target isomorphism $L$ because
$\mathsf R_\delta V^{-1}=LA$. Hence $[z]\mapsto[Lz]$ is an explicit
isomorphism of the two cokernels. It is not by itself a map of tori or of
global threefolds.

\subsection{The original map's rational nonintersection}

Retain the exact form and orthogonal matrix
\[
 \Delta(q)=4(q_1q_2+q_1q_3+q_2q_3)-q_1^2-q_2^2-q_3^2,
 \quad
 F=\begin{pmatrix}
 1/\sqrt3&1/\sqrt3&1/\sqrt3\\
 1/\sqrt2&0&-1/\sqrt2\\
 1/\sqrt6&-2/\sqrt6&1/\sqrt6
 \end{pmatrix}.
\]
For $\det A=-\delta$, $\delta\in\mathbb Q_{>0}$, Pro's map is
\[
 \mathfrak K_\delta(A)=F^{-1}
 \frac{((a+c)/2,b,(a-c)/2)^{\mathsf T}}{\sqrt{3\delta}}.
\]
It satisfies $\Delta(\mathfrak K_\delta(A))=-1$ but sends no rational
$A$ to a rational triple. Indeed, if $q$ were rational, put
$S=q_1+q_2+q_3$, $U=q_1-q_3$, $V=q_1+q_3-2q_2$. The rational coordinates
$H=(a+c)/2$, $X=b$, $Y=(a-c)/2$ would satisfy
\[
 H=\sqrt\delta S,\quad X=\sqrt{3\delta/2}\,U,
 \quad Y=\sqrt{\delta/2}\,V.
\]
Two nonzero coordinates would make one of $\sqrt{3/2}$, $\sqrt2$,
$\sqrt3$ rational. If only $H$ is nonzero the determinant is positive.
If only $X$ or only $Y$ is nonzero, the determinant equation forces,
respectively, $U^2=2/3$ or $V^2=2$, impossible over $\mathbb Q$.
All coordinates zero also contradicts $\delta>0$. This proves the claim.

\subsection{A different rational bridge, with exact inverse and denominators}

The preceding nonintersection concerns the specified map, not every map
between the quadratic surfaces. The following construction uses projection
from the rational point $(1,0,0)$ on $\Delta=-1$; it does not replace $F$ or
the physical meaning, if any, assigned to its coordinates.

\begin{lemma}\label{lem:esv3-anisotropic}
The equation $\Delta(q)=0$ has no nonzero rational solution.
\end{lemma}
\begin{proof}
After clearing denominators use the integral $S,U,V$ just defined.
Expansion gives $2\Delta=2S^2-3U^2-V^2$. Modulo three,
$2S^2=V^2$ forces $3\mid S,V$. Substitution then forces $3\mid U$.
Repeating gives divisibility of $S,U,V$ by every power of three, hence all
vanish. The linear map from $q$ to $(S,U,V)$ has nonzero determinant,
so $q=0$.
\end{proof}

\begin{theorem}\label{thm:esv3-rational-bridge}
For fixed $\delta\in\mathbb Q_{>0}$ write
\[
 X_\delta=\{(a,b,c):ac-b^2=-\delta\},\quad
 Y=\{q:\Delta(q)=-1\},
\]
and define
\[
 D(a,b)=-1-a^2-b^2+4a+4b+4ab,\quad K(a,b)=1-2a-2b.
\]
On $U_\delta=\{(a,b,c)\in X_\delta:aDK\ne0\}$ the map
\[
 \mathcal B_\delta(a,b,c)=
 \left(1+\frac{2K}{D},\frac{2Ka}{D},\frac{2Kb}{D}\right)
 \label{eq:esv3-birational}
\]
is an isomorphism of open varieties over $\mathbb Q$ onto the open part of
$Y$ where $q_1\ne1$, $q_2\ne0$, and
$D(q_2/(q_1-1),q_3/(q_1-1))\ne0$. Its inverse is
\[
 a=\frac{q_2}{q_1-1},\qquad b=\frac{q_3}{q_1-1},\qquad
 c=\frac{b^2-\delta}{a}.
\]
Every integral matrix on $X_\delta$ with $a>0$ belongs to $U_\delta$.
For such a matrix, the exact least common denominator of its image is
\[
 d_{\rm image}=\frac{|D(a,b)|}{\gcd(|D(a,b)|,|2K(a,b)|)}.
\]
In particular its image is integral exactly when $D(a,b)\mid2K(a,b)$.
\end{theorem}
\begin{proof}
Set $q=(1,0,0)+\lambda(1,a,b)$. Expansion, with no change of scale in
$\Delta$, gives
\[
 \Delta(q)=-1-2K\lambda+D\lambda^2.
\]
The second intersection with $Y$ is therefore $\lambda=2K/D$, nonzero
on the stated domain. The displayed image and inverse follow. Conversely,
for a point of the stated target open set, put $\lambda=q_1-1$ and recover
$a,b$. Its equation gives $\lambda D=2K$, so $K\ne0$ and the constructions
are inverse regular maps on the indicated opens. Recovering $c$ retains the
fixed determinant $-\delta$. For integral $a,b$, $K$ is odd and cannot
vanish. Moreover $D=\Delta(1,a,b)\ne0$ by
Lemma~\ref{lem:esv3-anisotropic}. Finally, $q_1-1=2K/D$ has reduced
denominator $|D|/\gcd(|D|,|2K|)$, and $q_2,q_3$ are integer multiples
of this fraction. This proves both the denominator formula and its exactness.
\end{proof}

For the packet's fixed opening
$A_*=\left(\begin{smallmatrix}1&1\\1&0\end{smallmatrix}\right)$,
$\delta=1$, $D=9$, $K=-3$, and
\[
 \mathcal B_1(A_*)=(1/3,-2/3,-2/3),\qquad
 \Delta(1/3,-2/3,-2/3)=-1.
\]
This rational image coexists with the irrationality of
$\mathfrak K_1(A_*)$, because the two maps are different and both have now
been specified. The rational bridge does not automatically preserve the
integral positive chamber, the integral orbit count, or a dynamics. The
denominator condition above is the exact next arithmetic restriction, and
the inverse retains $\delta$ as a parameter rather than erasing it.

\subsection{Independent checks and reproducibility}

The separate verifier
\path{supporting_materials/workbench/evidence/koide_es_v3_audit_20260904/verify_independent.py}
enumerates the smaller denominator $N/R<y\le2N/R$, then reconstructs
$z=Ny/(Ry-N)$, independently of the packet's divisor enumerator. It agrees
with all 1430 oriented states in 665 coprime shells through $n=161$.
It checks the 819 diagonal certificates for 329 primes below 5000, the
symbolic rational-map identity and inverse, and 820 exact denominator cases.
These finite regressions supplement, rather than replace, the proofs above.
The original packet's verifier is replayed separately without modifying its
source or its original output. The archive SHA-256 is
\begin{quote}\small\ttfamily
ab5c21cf00268f9019c95ee24b565d05c4d9\allowbreak
337c51e4a576ded9e3cc3f0f235e
\end{quote}
The correct full hash is recorded byte-for-byte in the evidence manifest;
the line breaking here is typographical only.

% Insertion-ready section.  Requires amsmath, amssymb, amsthm, booktabs, array.
\section{The fifteen $p=1201$ selector rays inside
$\operatorname{Herm}_3(\mathbb H)$}
\label{sec:fifteen-selector-jordan}

The common integer $15$ in the complete $p=1201$ selector census and in
$\dim_{\mathbb R}\operatorname{Herm}_3(\mathbb H)=15$ admits an exact test.
The test gives a nontrivial morphism, but it rules out the interpretation of
the selector states as fifteen independent Jordan-algebra directions.

\subsection{Marked selectors and their prime families}

Fix $p_0=1201$.  A marked state in the complete census is a tuple
\[
 s=(\chi,R,h,r,t),\qquad \chi\in\{E,M\},\qquad
 p_0=4rth-R,\qquad \gcd(r,t)=1,
\]
with positive integral $R,h,r,t$.  Its channel gate is
\begin{equation}
 \begin{array}{ll}
 E:&R\mid4r^2h+1,\\
 M:&R\mid r+t.
 \end{array}
 \label{eq:fifteen-gates}
\end{equation}
In the exterior channel the first gate is equivalent to $R\mid p_0r+t$,
because
\[
 p_0r+t=t(4r^2h+1)-Rr.
\]
Put
\begin{equation}
(u,a,v)=(hr^2,hrt,ht^2),\qquad uv=a^2,\qquad R=4a-p_0.
 \label{eq:fifteen-uav}
\end{equation}
If $\kappa_E=(p_0r+t)/R$ and $m=(r+t)/R$, the ordered denominator
maps are
\begin{align}
 D_E(s)&=(hrt,\ ht\kappa_E,\ p_0hr\kappa_E),\label{eq:fifteen-DE}\\
 D_M(s)&=(hrt,\ p_0hrm,\ p_0htm).\label{eq:fifteen-DM}
\end{align}
Clearing the full common denominator proves
\[
 \sum_{d\in D_\chi(s)}\frac1d=\frac4{p_0}.
\]

The fifteen states, with orientation and scale retained, are as follows.
The integer $L$ is the least prime increment that also preserves
$p\equiv361\pmod{840}$.
\begin{center}
\scriptsize
\begin{tabular}{@{}lccrcrrrrc@{}}
\toprule
tag&$\chi$&$R$&$h$&$(r,t)$&$u$&$a$&$v$&$L$&Koide $+$\\
\midrule
E23a&E&23&17&$(1,18)$&17&306&5508&57960&yes\\
M23a&M&23&3&$(6,17)$&108&306&867&14280&no\\
M23b&M&23&3&$(17,6)$&867&306&108&14280&no\\
E23b&E&23&34&$(9,1)$&2754&306&34&57960&no\\
M31a&M&31&2&$(1,154)$&2&308&47432&9240&yes\\
E31&E&31&77&$(4,1)$&1232&308&77&52080&no\\
M31b&M&31&2&$(154,1)$&47432&308&2&9240&no\\
M39a&M&39&2&$(1,155)$&2&310&48050&26040&yes\\
E39&E&39&62&$(5,1)$&1550&310&62&10920&no\\
M39b&M&39&2&$(155,1)$&48050&310&2&26040&no\\
M47a&M&47&1&$(8,39)$&64&312&1521&43680&no\\
M47b&M&47&1&$(39,8)$&1521&312&64&43680&no\\
E71&E&71&53&$(1,6)$&53&318&1908&59640&no\\
E339&E&339&11&$(35,1)$&13475&385&11&94920&no\\
E1043&E&1043&17&$(33,1)$&18513&561&17&1376760&no\\
\bottomrule
\end{tabular}
\end{center}

For an exterior row define
\[
 h_k=h_0+Rk,\qquad p_k=p_0+4rtRk;
\]
for a middle row define
\[
 h_k=h_0+k,\qquad p_k=p_0+4rtk.
\]
The corresponding gate in \eqref{eq:fifteen-gates} is preserved.  The
exterior prime class modulo $4rtR$ and the middle prime class modulo $4rt$
are reduced.  Their least intersections with $361\pmod{840}$ are
\begin{equation}
 L_s=
 \begin{cases}
  \operatorname{lcm}(4rtR,840),&\chi=E,\\
  \operatorname{lcm}(4rt,840),&\chi=M.
 \end{cases}
 \label{eq:fifteen-hard-L}
\end{equation}
Thus
\begin{equation}
 p_n=1201+L_sn,\qquad h_n=h_0+\frac{L_s}{4rt}n
 \label{eq:fifteen-hard-family}
\end{equation}
retains the same marked ray, and Dirichlet's theorem gives infinitely many
prime $p_n$ in every row.  The fifteen marked rows give eleven unordered
denominator triples at $p_0$ and ten distinct residue classes
$1201\pmod{L_s}$.  These are different quotients: middle reversal produces
the first collapse, while E23a and E23b have distinct rays but the same
$L_s=57960$.

\subsection{The exact Koide cone map}

Define
\begin{equation}
 \Gamma(s)=\begin{pmatrix}u&a\\a&v\end{pmatrix}
 =h\binom rt(r\ t).
 \label{eq:fifteen-Gamma}
\end{equation}
This matrix is positive semidefinite of rank one and has Smith normal form
$\operatorname{diag}(h,0)$.  The Lorentz-coordinate map
\begin{equation}
 \mathcal L(\Gamma)=
 \left(\frac{u+v}{2},a,\frac{u-v}{2}\right)=(H,X,Y)
 \label{eq:fifteen-Lorentz}
\end{equation}
satisfies $H^2-X^2-Y^2=uv-a^2=0$.

Let
\[
 F(q_1,q_2,q_3)=
 \left(\frac{q_1+q_2+q_3}{\sqrt3},
       \frac{q_1-q_3}{\sqrt2},
       \frac{q_1+q_3-2q_2}{\sqrt6}\right).
\]
This is orthogonal.  Its inverse at \eqref{eq:fifteen-Lorentz} is
\begin{align}
 q_1&=\frac{u+v}{2\sqrt3}+\frac a{\sqrt2}
                         +\frac{u-v}{2\sqrt6},\nonumber\\
 q_2&=\frac{u+v}{2\sqrt3}-\frac{u-v}{\sqrt6},\label{eq:fifteen-Koide}\\
 q_3&=\frac{u+v}{2\sqrt3}-\frac a{\sqrt2}
                         +\frac{u-v}{2\sqrt6}.\nonumber
\end{align}
For
\[
 Q_K(q)=\frac23(q_1+q_2+q_3)^2-(q_1^2+q_2^2+q_3^2)
\]
orthogonality gives $Q_K(q)=H^2-X^2-Y^2=0$.  Hence
\eqref{eq:fifteen-Koide} is an exact coefficient-preserving map over
$\mathbb Q(\sqrt2,\sqrt3)$, not a numerical fit.  In the displayed
orientation exactly E23a, M31a, and M39a lie in the positive Koide chamber.

\subsection{Embedding and rank obstruction in the quaternionic Jordan algebra}

Let $J=\operatorname{Herm}_3(\mathbb H)$ with
$A\circ B=(AB+BA)/2$, and put $e=\operatorname{diag}(1,1,0)$.  Define
\begin{equation}
 \iota:\operatorname{Sym}_2(\mathbb R)\longrightarrow J,
 \qquad
 \iota\begin{pmatrix}u&a\\a&v\end{pmatrix}
 =\begin{pmatrix}u&a&0\\a&v&0\\0&0&0\end{pmatrix}.
 \label{eq:fifteen-iota}
\end{equation}
Block multiplication proves
$\iota(A\circ B)=\iota(A)\circ\iota(B)$.  Thus $\iota$ is an injective
real-linear Jordan homomorphism with image
\[
 C_{\mathbb R}:=\iota(\operatorname{Sym}_2(\mathbb R))
 =eJe\cap\operatorname{Herm}_3(\mathbb R),
 \qquad \dim_{\mathbb R}C_{\mathbb R}=3.
\]
The full Peirce corner is
$eJe\cong\operatorname{Herm}_2(\mathbb H)$ of real dimension $6$;
$C_{\mathbb R}$ is its proper real-symmetric Jordan subalgebra.
It preserves trace, trace pairing, positivity, and ordinary rank.  In the
Moore convention
\[
 \begin{pmatrix}d_1&q_3&\bar q_2\\
 \bar q_3&d_2&q_1\\q_2&\bar q_1&d_3\end{pmatrix},
\]
the only nonzero coordinates of $\iota\Gamma(s)$ are
$d_1=u,d_2=v,q_3=a\in\mathbb R$.  Its Moore cubic is zero.

Projectivizing the scale gives the idempotent
\begin{equation}
 P_{r,t}=\frac1{r^2+t^2}
 \begin{pmatrix}r^2&rt&0\\rt&t^2&0\\0&0&0\end{pmatrix},
 \qquad P_{r,t}^2=P_{r,t},\quad\operatorname{tr}P_{r,t}=1.
 \label{eq:fifteen-projector}
\end{equation}
It represents the quaternionic line $[r:t:0]_{\mathbb H}$.  The fifteen
rows have pairwise distinct ratios $[r:t]$, so the composition
\begin{equation}
 \Sigma_{1201}\hookrightarrow\mathbb{RP}^1\hookrightarrow\mathbb{HP}^2,
 \qquad s\longmapsto[r:t]\longmapsto[r:t:0]_{\mathbb H}
 \label{eq:fifteen-RP1-HP2}
\end{equation}
is injective.

\begin{theorem}[Exact span and Jordan closure]
\label{thm:fifteen-span}
The fifteen matrices $\iota\Gamma(s)$ are fifteen distinct projective
rank-one points, but
\[
 \operatorname{span}_{\mathbb R}\{\iota\Gamma(s):s\in\Sigma_{1201}\}
 =\langle\iota\Gamma(\Sigma_{1201})\rangle_{\rm Jordan}=C_{\mathbb R}.
\]
In particular, both spaces have real dimension $3$, not $15$.
\end{theorem}
\begin{proof}
In coordinates
$\bigl(\begin{smallmatrix}x&y\\y&z\end{smallmatrix}\bigr)
\leftrightarrow(x,y,z)$, the projective image of $(r,t)$ is
$(r^2,rt,t^2)$.  The rows E31, E39, and E339 give
\[
 \begin{pmatrix}16&4&1\\25&5&1\\1225&35&1\end{pmatrix},
 \qquad \det=-930\ne0.
\]
Hence the images span the whole three-dimensional real-symmetric corner
$C_{\mathbb R}$.  The space $C_{\mathbb R}$ is a Jordan subalgebra by the
block-multiplication identity above, so the generated Jordan algebra is
contained in $C_{\mathbb R}$.  A Jordan subalgebra is a real vector space;
because it contains the fifteen generators, it contains their span
$C_{\mathbb R}$.  Both containments are equalities.
\end{proof}

The number fifteen is not invariant even after the hard residue class is held
fixed.  Complete divisor enumeration gives the following exact comparison:
\[
\begin{array}{c|rrrr}
 p&1201&4561&8761&9601\\ \hline
 \#E&7&26&16&17\\
 \#M&8&28&20&12\\
 \#(E\sqcup M)&15&54&36&29\\
 \#\{\text{unordered denominator triples}\}&11&40&26&23
\end{array}
\]
All four displayed primes satisfy $p\equiv361\pmod{840}$.  Thus neither the
hard residue class nor the selector construction singles out a
fifteen-element family; the equality
$\#\Sigma_{1201}=\dim_{\mathbb R}\operatorname{Herm}_3(\mathbb H)=15$ is a
feature of the individual prime $1201$, not an invariant of the arithmetic
family.  The deterministic enumeration and primality checks are included in
the accompanying verification script.

At fixed $p_0$ the unprojectivized embedding preserves the markings.  From
$\iota\Gamma$ one recovers $(u,a,v)$ and then
\[
 h=\gcd(u,a,v),\quad r=\sqrt{u/h},\quad t=\sqrt{v/h},\quad
 R=4a-1201,\quad j=(R-3)/4.
\]
The two channel gates cannot overlap in this census: overlap would imply
$R\mid1200$, and none of $23,31,39,47,71,339,1043$ divides $1200$.
Projectivization in \eqref{eq:fifteen-projector} loses $h$, and without the
fixed prime the residual formula $R=4a-p$ is unavailable.

\subsection{Symmetry and incidence obstructions}

Normalized orientation reversal fixes the seven exterior states and swaps
the four middle pairs, so its cycle type is $1^7 2^4$.  The transvection
\[
 \tau(x_1,x_2,x_3,x_4)=(x_1+x_2,x_2,x_3,x_4)
\]
on $\mathbb F_2^4\setminus\{0\}$ has the same cycle type.  For example, with
integers $1,\ldots,15$ encoding the four binary coordinates from the low bit,
the following map is $C_2$-equivariant:
\[
\begin{array}{c|rrrrrrr}
s&E23a&E23b&E31&E39&E71&E339&E1043\\\hline
\phi(s)&1&4&5&8&9&12&13
\end{array}
\]
\[
\begin{array}{c|rrrr}
s&M23a&M23b&M31a&M31b\\\hline
\phi(s)&2&3&6&7
\end{array}
\qquad
\begin{array}{c|rrrr}
s&M39a&M39b&M47a&M47b\\\hline
\phi(s)&10&11&14&15
\end{array}
\]
This proves only an equivalence of $C_2$-sets.  It is noncanonical: a set
with this cycle type has $7!2^4 4!=1{,}935{,}360$ automorphisms commuting
with the involution.

It cannot be upgraded to the natural projective incidence.  The selector
points lie on the nonsingular conic $xz-y^2=0$ in
$\mathbb P^2(\mathbb R)$, so no three distinct selector images are collinear.
By contrast, $PG(3,2)$ has $35$ three-point lines.  Nor does normalized
reversal extend linearly on the canonical real-symmetric corner: it fixes
E31, E39, and E339, which span $C_{\mathbb R}$, so such a linear map would be
the identity; but it must exchange four nonidentical middle pairs.

The other elementary count comparisons do not supply a bridge.  Type $A_3$
has $12$ oriented roots, $6$ unoriented root lines, and $4+6+4=14$ nonzero
minuscule weights.  The rank-one idempotents of $J$ form the real
eight-dimensional continuum $\mathbb{HP}^2$, while the present images occupy
the proper copy $\mathbb{RP}^1$ in \eqref{eq:fifteen-RP1-HP2}.

\subsection{Conclusion and next exact test}

The current structural statement is therefore
\[
 \boxed{
 \begin{gathered}
 \text{fifteen marked prime-carrying selector rays}\ 
 \hookrightarrow\mathbb{RP}^1\hookrightarrow\mathbb{HP}^2,\\
 \dim\operatorname{span}=\dim\operatorname{Jordan\ closure}=3.
 \end{gathered}}
\]
The equality between the number of marked rays and $\dim_{\mathbb R}J$ is
not explained by the defined morphism.

A decisive higher-rung test would construct, from the arithmetic residual and
channel markings, an explicit element $g_s$ of a named quaternionic structure
group and study
\[
 g_s\iota(\Gamma(s))g_s^*.
\]
One must then compute the span and Jordan closure over complete multi-prime
censuses and prove equivariance with the selector involutions.  If all images
remain in one common three-dimensional real-symmetric Jordan corner, the
obstruction persists; if the markings canonically move them through enough
conjugate corners and the closure is all of $J$, that would be a genuine
fifteen-dimensional bridge.  No such
marking-to-conjugation map is presently defined.

\section{The analytic threefold constructed in the candidate}
\label{sec:candidate-geometry}

This section audits the construction itself.  It does not infer existence
from the later topology.  The period family, the three local fillings, and
the gluing are proved in that order.  The notation is the candidate's:
\(B\simeq\mathbb P^1\),
\(B^\circ=B\setminus\{p_0,p_1,p_2\}\), and the local monodromy orders at
\(p_1,p_2\) are \(m_1=3,m_2=4\).

\subsection{The period family}

Let \(\Delta\simeq(\mathbb Z/3)*(\mathbb Z/4)\) be the orientation-preserving
\((3,4,\infty)\)-triangle group acting on a simply connected Riemann surface
\(\mathfrak h_z\), with elliptic generators \(g_1,g_2\) and
\(g_0=(g_1g_2)^{-1}\).  Let \(t\) be the affine coordinate on
\(B\) with values \(0,1,\infty\) at \(p_1,p_2,p_0\), respectively.
Write
\[
 \pi:\mathfrak h_z\longrightarrow B\setminus\{p_0\}
\]
for the orbifold quotient map; its stabilizers at chosen lifts
\(z_j\in\pi^{-1}(p_j)\) are \(\langle g_j\rangle\).  Let
\(\rho_V:\Delta\to\operatorname{GL}(V)\) be the printed representation,
\(\rho_V(g_j)=T_j\), and put
\[
 M_g:=\rho_V(g)^{\mathsf t},\qquad A_g:=M_g^{-1}.
\]
Thus \(M_{gh}=M_hM_g\), \(A_{g_j}=A_j\), and
\(A_{g_0}=M_0\) in the notation of \cref{prop:monodromy-certificate}.

\begin{lemma}[simultaneous elliptic-orientation normalization]
\label{lem:triangle-generator-orientations}
\Verified\ The generators may, and henceforth do, satisfy
\begin{equation}
 g_1'(z_1)=e^{-2\pi i/3},\qquad
 g_2'(z_2)=e^{-2\pi i/4}=-i,
\label{eq:triangle-generator-derivatives}
\end{equation}
while \(g_0=(g_1g_2)^{-1}\) remains parabolic.  More precisely, the Cayley
coordinates
\begin{equation}
 s_j(z)=\frac{z-z_j}{z-\bar z_j}
\label{eq:triangle-cayley-coordinates}
\end{equation}
obey
\begin{equation}
 s_1\circ g_1=e^{-2\pi i/3}s_1,
 \qquad
 s_2\circ g_2=-is_2.
\label{eq:triangle-linearizing-identities}
\end{equation}
The two signs in \cref{eq:triangle-generator-derivatives} are a simultaneous
orientation choice; they are not obtained by shrinking a coordinate disc.
\end{lemma}

\begin{proof}
Realize \(\Delta\) as the orientation-preserving subgroup of the reflection
group of a hyperbolic triangle with counterclockwise vertices \(A,B,C\),
angles \(\pi/3,\pi/4,0\), and opposite sides \(a,b,c\), respectively.  Let
\(\sigma_a,\sigma_b,\sigma_c\) denote reflection in those sides and put
\[
 x=\sigma_b\sigma_c,\qquad
 y=\sigma_c\sigma_a,\qquad
 z=\sigma_a\sigma_b.
\]
The product of reflections in two oriented geodesics meeting at angle
\(\theta\) is the rotation through \(2\theta\).  Hence \(x\) is the
counterclockwise rotation about \(A\) through \(2\pi/3\), \(y\) is the
counterclockwise rotation about \(B\) through \(\pi/2\), and \(z\) is
parabolic at the ideal vertex \(C\); this is the standard triangle-group
construction \cite[Chapters~7 and~10]{Beardon1983}.  Moreover
\[
 xyz=\sigma_b\sigma_c\sigma_c\sigma_a\sigma_a\sigma_b=1.
\]
Set
\[
 z_1=A,\qquad z_2=B,\qquad g_1=x^{-1},\qquad g_2=y^{-1}.
\]
Both finite rotations are thereby reversed, so their complex derivatives at
their fixed points are exactly those in
\cref{eq:triangle-generator-derivatives}.  Since
\(g_2g_1=y^{-1}x^{-1}=z\), one also has
\[
 g_1g_2=g_1zg_1^{-1};
\]
thus \(g_1g_2\), and therefore \(g_0=(g_1g_2)^{-1}\), is parabolic.

This simultaneous reversal is forced once the orientation of \(g_1\) and
the cusp conjugacy class are fixed.  Indeed, in
\(\langle x\rangle*\langle y\rangle\simeq C_3*C_4\) every parabolic element
is conjugate to a nonzero power of the primitive cusp word
\(z=y^{-1}x^{-1}=y^3x^2\).  The two mixed choices would instead give
\[
 (x^{-1}y)^{-1}=y^{-1}x=y^3x,
 \qquad
 (xy^{-1})^{-1}=yx^{-1}=yx^2.
\]
Both words are cyclically reduced of length two.  A nontrivial power
\(z^n\) has cyclically reduced length \(2|n|\); hence a mixed word could be
parabolic only for \(|n|=1\).  But neither \(y^3x\) nor \(yx^2\) is a cyclic
permutation of \(y^3x^2\) or its inverse \(xy\).  Thus neither mixed choice
is parabolic.  The derivative selected at \(z_1\) consequently fixes the
derivative at \(z_2\) as well.

Finally, for an elliptic M\"obius transformation \(g\) fixing
\(z_0\in\mathfrak h_z\), the Cayley map
\(s=(z-z_0)/(z-\bar z_0)\) conjugates \(g\) to a disc automorphism fixing
zero.  Such an automorphism is \(s\mapsto g'(z_0)s\).  Applying this to
\(g_1,g_2\) proves \cref{eq:triangle-linearizing-identities}.
\end{proof}

\begin{proposition}[existence and non-degeneracy of the period matrix]
\label{prop:period-input-audit}
\Verified\ There are holomorphic functions
\[
 \tau:\mathfrak h_z\to\mathfrak h,
 \qquad \mu,\beta:\mathfrak h_z\to\mathbb C
\]
satisfying
\[
\begin{array}{lll}
\tau\circ g_0=\tau-1,&\mu\circ g_0=\mu,&\beta\circ g_0=\beta+1,\\
\tau\circ g_2=-1/\tau,&\mu\circ g_2=1+\mu/\tau,
 &\beta\circ g_2=\beta-3-6\mu^2/\tau,
\end{array}
\tag{2.1}
\]
and the corresponding order-three transformation laws obtained from
\(\tau\circ g_1=(\tau-1)/\tau\).  They may be chosen so that
\[
 \operatorname{Im}\tau>0,
 \qquad
 D:=\operatorname{Im}\beta-
 \frac{6(\operatorname{Im}\mu)^2}{\operatorname{Im}\tau}<0.
\tag{2.2}
\]
Consequently
\[
 \Pi=[Z\mid I_2],\qquad
 Z=\begin{pmatrix}6\mu&\tau\\ \beta&\mu\end{pmatrix}
\tag{2.3}
\]
maps the marked rank-four lattice \(\Lambda\) isomorphically onto a lattice
in \(\mathbb C^2\) at every point.  The quotients
\(\mathbb C^2/\Pi(z)\Lambda\) form a holomorphic family of compact
two-dimensional complex tori over \(\mathfrak h_z\), equivariant for
\(\Delta\).
\end{proposition}

\begin{proof}
Normalize the modular function by
\(j(q)=q^{-1}(1+O(q))\) and put
\(J_B=1728(t\circ\pi)\); the modular facts used here are collected in
\cite[Chapter~VII]{Serre1973}.  Put
\[
 \mathfrak h_z^\dagger=
 \mathfrak h_z\setminus J_B^{-1}\{0,1728\},\qquad
 \mathfrak h_\tau^\dagger=
 \mathfrak h\setminus j^{-1}\{0,1728\}.
\]
Away from the deleted sets, the classical modular map
\(j:\mathfrak h_\tau^\dagger\to\mathbb C\setminus\{0,1728\}\) is an
unramified covering with deck group
\(\operatorname{PSL}_2(\mathbb Z)\).  The deleted set in
\(\mathfrak h_z\) is locally finite.  Since \(\mathfrak h_z\) is simply
connected, \(\pi_1(\mathfrak h_z^\dagger)\) is normally generated by small
positive meridians about its points.  At a translate of \(z_1\), \(J_B\)
has local degree three over \(0\); the corresponding meridian therefore has
deck image equal to the third power of the order-three modular generator.
At a translate of \(z_2\), \(J_B\) has local degree four over \(1728\); its
meridian has deck image equal to the fourth power of the order-two modular
generator.  Both images are the identity.  The covering-space lifting
criterion gives a global lift \(\tau\) with \(j(\tau)=J_B\).

Near a deleted point, compose the lift with a biholomorphism
\(\chi:\mathfrak h\to\mathbb D\).  The bounded map \(\chi\circ\tau\)
extends by the removable-singularity theorem.  Its value at the puncture has
modulus strictly less than one: equality would give an interior maximum and
force \(\chi\circ\tau\) to be constant, although its punctured values lie in
 \(\mathbb D\).  Thus the extension still lands in \(\mathfrak h\).

Because \(J_B\) is \(\Delta\)-invariant, \(\tau\circ g\) and \(\tau\) are
lifts of the same map through the Galois covering \(j\) on the punctured
surface.  Hence there is a unique deck transformation
\(\varphi(g)\in\operatorname{PSL}_2(\mathbb Z)\) such that
\begin{equation}
 \tau\circ g=\varphi(g)\circ\tau.
\label{eq:period-deck-homomorphism}
\end{equation}
The identity extends across the deleted points.  Uniqueness and
\((\tau\circ g)\circ h=\tau\circ(gh)\) give
\(\varphi(gh)=\varphi(g)\varphi(h)\), so \(\varphi\) is a homomorphism.

Normalize \(\tau(z_1)=\rho=e^{\pi i/3}\).  The local cubic expansion of
\(j\) at \(\rho\), together with the local degree three of \(J_B\), gives a
simple zero of \(\tau-\rho\) at \(z_1\).  At \(z_2\), choose
\(\gamma\in\operatorname{PSL}_2(\mathbb Z)\) with
\(\tau(z_2)=\gamma(i)\).  The local quadratic expansion of
\(j\circ\gamma\) at \(i\), together with the local degree four of \(J_B\),
 gives a double zero of \(\gamma^{-1}\tau-i\) at \(z_2\).  At \(z_1\),
 differentiating \cref{eq:period-deck-homomorphism} and using the simple
 zero gives
 \[
  \varphi(g_1)'(\rho)=g_1'(z_1)=e^{-2\pi i/3}.
 \]
 The stabilizer of \(\rho\) is generated by
 \(S_1(\tau)=(\tau-1)/\tau\), whose derivative at \(\rho\) is
 \(e^{-2\pi i/3}\); hence \(\varphi(g_1)=S_1\), rather than \(S_1^{-1}\).
 The order-four generator cannot have trivial modular image:
otherwise the invariant local power series \(\gamma^{-1}\tau-i\) would have
order divisible by four, contrary to its order two.  Its modular image
 therefore has order two.  If \(P=\varphi(g_0)\), the exact relation
 \(g_1g_2g_0=1\) gives
 \begin{equation}
  \varphi(g_2)=S_1^{-1}P^{-1}.
 \label{eq:period-second-generator-from-cusp}
 \end{equation}

It remains to determine the cusp image.  Choose a punctured cusp disc
\(D_0^*\subset B\) and a component
\[
 U_0\subset\pi^{-1}(D_0^*)
\]
stabilized by \(\langle g_0\rangle\).  Choose \(R\) and shrink \(D_0^*\) so
that \(\tau(U_0)\) lies in a single component
\[
 K\subset j^{-1}(\{|w|>R\}).
\]
Put \(P=\varphi(g_0)\).  Since \(g_0U_0=U_0\),
\[
 P\tau(U_0)=\tau(g_0U_0)=\tau(U_0).
\]
Hence \(P(K)=K\): the two components \(K\) and \(P(K)\) intersect in the
nonempty set \(\tau(U_0)\).  Let \(c\) be the rational boundary cusp of
\(K\), choose
\(\gamma_0\in\operatorname{PSL}_2(\mathbb Z)\) with
\(\gamma_0(\infty)=c\), and set
\[
 \tau'=\gamma_0^{-1}\tau|_{U_0},\qquad
 P'=\gamma_0^{-1}P\gamma_0.
\]
Then
\[
 \tau'\circ g_0=P'\circ\tau',
\]
and \(P'\) preserves \(K'=\gamma_0^{-1}K\), whose boundary cusp is
\(\infty\).

Suppose \(P'=1\).  Then \(\tau'\) is \(g_0\)-invariant and descends to a
holomorphic map \(\bar\tau':D_0^*\to\mathfrak h\).  Put
\[
 q'=\exp(2\pi i\bar\tau').
\]
After increasing \(R\), the \(q\)-expansion is valid on \(K'\), and
\(0<|q'|<1\).  Thus \(q'\) extends holomorphically across the puncture.
If \(|q'(0)|=1\), the maximum-modulus principle makes \(q'\), hence
\(J_B\), constant.  If \(q'(0)\ne0\), the expansion
\(j(q')=(q')^{-1}(1+O(q'))\) makes \(J_B\) bounded, contrary to its pole.
Therefore \(q'(0)=0\).  But a holomorphic zero has positive winding on a
small positive circle, while \(q'\) has the single-valued holomorphic
logarithm \(2\pi i\bar\tau'\) on \(D_0^*\), so its winding is zero.  This
contradiction proves \(P'\ne1\).  Since \(P'\) preserves the selected cusp
component \(K'\), it is parabolic; hence the original \(P\) is a nontrivial
parabolic.

We now discard the auxiliary conjugation by \(\gamma_0\) and return to the
original normalization
\[
 \tau(z_1)=\rho,\qquad \varphi(g_1)=S_1.
\]
Choose the trace-\(+2\) lift
\(\widetilde P\in\operatorname{SL}_2(\mathbb Z)\) of \(P\).  For a
primitive fixed vector \((p,r)^{\mathsf t}\), it has the unique form
\[
 \widetilde P
 =I+k\begin{pmatrix}-pr&p^2\\-r^2&pr\end{pmatrix},
 \qquad p,r\in\mathbb Z,\quad \gcd(p,r)=1,\quad k\ne0.
\]
With
\[
 \widetilde S_1=\begin{pmatrix}1&-1\\1&0\end{pmatrix},
\]
one has
\[
 \operatorname{tr}(\widetilde S_1^{-1}\widetilde P^{-1})
 =1+k(p^2-pr+r^2).
\]
Since \(S_1^{-1}P^{-1}=\varphi(g_2)\) has projective order two, this trace
is zero.  Hence
\[
 -k(p^2-pr+r^2)=1,
\]
so \(k=-1\) and \(p^2-pr+r^2=1\).  Modulo simultaneous sign, the three
solutions are cycled transitively by \(S_1\).  Conjugating the entire deck
representation by a suitable power of \(S_1\), which fixes \(\rho\) and
leaves \(S_1\) unchanged, therefore gives
\[
 P(w)=w-1,\qquad \varphi(g_2)(w)=-\frac1w.
\]
Rename the resulting period map \(\tau\).  Since \(z_2\) is fixed by the
order-four generator, the new normalization has \(\tau(z_2)=i\), and the
preceding local ramification calculation says exactly that \(\tau-i\) has
order two at \(z_2\).  Let \(t_c\) vanish at \(p_0\) and define anew
\[
 q=e^{2\pi i\tau}.
\]
Since \(\tau\circ g_0=\tau-1\), \(q\) descends to \(D_0^*\).  The simple
pole of \(J_B=j(\tau)\) gives
\[
 q=t_cu_1(t_c),\qquad u_1(0)\ne0.
\]

Writing
\[
 S_1(w)=\frac{w-1}{w},\qquad S(w)=-\frac1w,
 \qquad P(w)=w-1,
\]
the normalized deck certificate is therefore
\begin{equation}
 \varphi(g_1)=S_1,\qquad
 \varphi(g_2)=S,\qquad
 \varphi(g_0)=P,\qquad
 S_1SP=1
 \quad\text{in }\operatorname{PSL}_2(\mathbb Z).
\label{eq:period-deck-certificate}
\end{equation}
Differentiation at \(z_1\), where \(\tau-\rho\) has order one, gives
\(g_1'(z_1)=S_1'(\rho)=\rho^{-2}=e^{-2\pi i/3}\).  At \(z_2\), comparison
of the quadratic leading terms gives only
\[
 \bigl(g_2'(z_2)\bigr)^2=S'(i)=-1.
\]
The choice \(g_2'(z_2)=-i\), rather than \(+i\), is supplied by the
simultaneous parabolic normalization in
\cref{lem:triangle-generator-orientations}; it does not follow from the
quadratic ramification calculation alone.

We now solve the two affine cocycle problems, including the orbifold and
cusp regularity which the abbreviated argument had suppressed.  Define
\(\widetilde\jmath:\Delta\times\mathfrak h_z\to\mathbb C^*\) on the
generators by
\[
 \widetilde\jmath(g_1,z)=-\tau(z),\qquad
 \widetilde\jmath(g_2,z)=\tau(z),
\]
and extend it by
\(\widetilde\jmath(gh,z)=\widetilde\jmath(g,hz)
\widetilde\jmath(h,z)\).  This is well defined: the products around the
two finite cyclic subgroups are
\[
 (-1)^3\tau\frac{\tau-1}{\tau}\frac{-1}{\tau-1}=1,
 \qquad
 \tau\left(-\frac1\tau\right)\tau\left(-\frac1\tau\right)=1.
\]
Moreover \(\widetilde\jmath(g_0,\cdot)=1\).

For an open \(U\subset B\), set
\(\widetilde U=\pi^{-1}(U\setminus\{p_0\})\).  Define
\[
 \mathcal L(U)=\left\{\nu\in\mathcal O(\widetilde U):
 \nu\circ g=\nu/\widetilde\jmath(g,\cdot)\quad(g\in\Delta)\right\}
\]
and let \(\mathcal M(U)\) be the holomorphic functions on
\(\widetilde U\) satisfying
\[
 \mu\circ g_1=\frac{1-\mu}{\tau},\qquad
 \mu\circ g_2=1+\frac\mu\tau.
 \tag{2.1a}
\]
If \(p_0\in U\), both definitions additionally require that the restriction
to the distinguished \(\langle g_0\rangle\)-invariant cusp component,
viewed as a function of \(t_c\), extend holomorphically across \(t_c=0\).
Because \(\widetilde\jmath(g_0)=1\) and \(\mu\circ g_0=\mu\), this
requirement is meaningful and is stable under restriction.  Thus
\(\mathcal L\) is a line sheaf and \(\mathcal M\), if locally nonempty, is
an affine \(\mathcal L\)-torsor.

Local nonemptiness is explicit.  Over an ordinary small disc, prescribe
\(\mu=0\) on one sheet and extend using (2.1a).  At the cusp prescribe
\(\mu=0\) on the distinguished component.  At the elliptic points use
\[
 \mu_1=\frac{2-\tau}{3}\quad\hbox{on }\Delta_1,
 \qquad
 \mu_2=\frac{1-\tau}{2}\quad\hbox{on }\Delta_2.
 \tag{2.1b}
\]
Indeed
\[
 \mu_1\circ g_1=\frac{\tau+1}{3\tau}
 =\frac{1-\mu_1}{\tau},
 \qquad
 \mu_2\circ g_2=\frac{\tau+1}{2\tau}
 =1+\frac{\mu_2}{\tau}.
\]
The cyclic products above, or direct iteration of (2.1a), show that these
prescriptions close after three and four steps.

To identify \(\mathcal L\), use that \(E_6\) has a simple zero at the
order-two elliptic orbit and no other zero in \(\mathfrak h\).  Since
\(\tau-i\) has order two at every translate of \(z_2\), the function
\(E_6\circ\tau\) has only even-order zeros on the simply connected surface
\(\mathfrak h_z\).  Fix a holomorphic square root \(h\),
\(h^2=E_6\circ\tau\), and put
\[
 \Theta=\frac{(E_4\circ\tau)^2h}{\Delta_{\rm mod}\circ\tau}.
\]
For \(j=1,2\), the quotient
\(h(g_jz)/(\tau(z)^3h(z))\) has square one away from its removable zero
set, hence is a constant sign.  At \(z_1\), \(h(z_1)\ne0\) and
\(\rho^3=-1\), so the sign for \(g_1\) is \(-1\).  At \(z_2\), choose the
linearizing coordinate with \(s_2\circ g_2=-is_2\).  The function \(h\) has
a simple zero there, and therefore
\(h(g_2z)/h(z)\to-i=i^3\); the sign for \(g_2\) is \(+1\).  The modular
transformation formulae
\[
 E_4(\gamma\tau)=(c\tau+d)^4E_4(\tau),\quad
 E_6(\gamma\tau)=(c\tau+d)^6E_6(\tau),\quad
 \Delta_{\rm mod}(\gamma\tau)=(c\tau+d)^{12}\Delta_{\rm mod}(\tau)
\]
now give the exact identities
\[
 \Theta\circ g_1=-\Theta/\tau,\qquad
 \Theta\circ g_2=\Theta/\tau.
 \tag{2.1e}
\]
These are precisely the automorphy laws defining \(\mathcal L\).  Since
\(E_4\) has a simple zero at \(\rho\), whereas \(E_6\) has a simple zero
at \(i\), the orders of \(\Theta\) at \(z_1,z_2\) are respectively two
and one.  At the distinguished cusp, \(E_4,E_6\) are units and
\(\Delta_{\rm mod}=q\cdot(\text{unit})\), with
\(q=t_c\cdot(\text{unit})\), so
\(\Theta=t_c^{-1}\cdot(\text{unit})\).
Taking the stabilizer orders into account, division by \(\Theta\) identifies
the ordinary and elliptic stalks with \(\mathcal O_B\), while the cusp stalk
is \(\mathfrak m_{p_0}\Theta\).  Hence
\[
 \mathcal L\simeq\mathcal O_B(-p_0)\simeq\mathcal O_{\mathbb P^1}(-1).
\]
Consequently \(H^1(B,\mathcal L)=H^0(B,\mathcal L)=0\): the locally
nonempty torsor \(\mathcal M\) has a unique global section.  Evaluating at
the fixed points gives
\[
 \mu(z_1)=\frac{2-\rho}{3},\qquad
 \mu(z_2)=\frac{1-i}{2},
\]
and the cusp condition says precisely that its invariant branch extends
holomorphically.

For the second affine problem put
\[
 \phi_1=2-\frac{6(1-\mu)^2}{\tau},\qquad
 \phi_2=-3-\frac{6\mu^2}{\tau}.
\]
Substitution around the finite orbits gives
\[
 \sum_{k=0}^{2}\phi_1\circ g_1^k=0,
 \qquad
 \sum_{k=0}^{3}\phi_2\circ g_2^k=0.
 \tag{2.1c}
\]
For example, the three summands in the first identity are
\[
2-\frac{6(1-\mu)^2}{\tau},\quad
2-\frac{6(\tau-1+\mu)^2}{\tau(\tau-1)},\quad
2+\frac{6\mu^2}{\tau-1},
\]
 Putting them over the common denominator \(\tau(\tau-1)\), the three
 constant terms contribute \(6\tau(\tau-1)\), while the remaining terms
 contribute \(-6\tau(\tau-1)\); hence the total numerator is zero.  For the
 order-four identity the
bracket after extracting \(6/\tau\) is \(2\tau\), cancelling \(-12\).

Define \(\mathcal B(U)\) to consist of holomorphic \(\beta\) on
\(\widetilde U\) with
\(\beta\circ g_j-\beta=\phi_j\); at the cusp require
\(b=\beta+\tau\) to extend holomorphically in \(t_c\).  Differences of
sections are ordinary holomorphic functions on \(U\), so
\(\mathcal B\) is a pseudo-torsor under \(\mathcal O_B\).  It is locally
nonempty: prescribe \(\beta=0\) on one ordinary sheet; prescribe
\(\beta=-\tau\) on the distinguished cusp component; and at \(p_j\) use
\[
 \beta_j=\frac1{m_j}\sum_{k=0}^{m_j-1}
 k\,(\phi_j\circ g_j^k).
 \tag{2.1d}
\]
Equation (2.1c) makes the last expression telescope to
\(\beta_j\circ g_j-\beta_j=\phi_j\).  Therefore \(\mathcal B\) is an
\(\mathcal O_B\)-torsor.  Since \(H^1(B,\mathcal O_B)=0\), it has global
sections, and any two differ by a constant; by construction
\(b=\beta+\tau\) is invariant and holomorphic at the cusp.

The real determinant is
\[
 \det_{\mathbb R}\Pi=-\det\operatorname{Im}Z
 =\operatorname{Im}\tau\,D.
\tag{2.4}
\]
For completeness, the equivariance can be checked before any quotient is
taken.  Regard the rows of \(\Pi\) as the two covectors
\[
 \sigma_1=6\mu\gamma+\tau u+w,
 \qquad
 \sigma_2=\beta\gamma+\mu u+\delta
 \quad\text{in }V\otimes\mathbb C.
 \tag{2.3a}
\]
Multiplication by the printed matrices gives
\[
\begin{aligned}
T_1\sigma_1&=(6\mu-6)\gamma+(1-\tau)u-\tau w,\\
T_1\sigma_2&=(\beta+2)\gamma+(1-\mu)u+(1-\mu)w+\delta,\\
T_2\sigma_1&=(6\mu+6\tau)\gamma-u+\tau w,\\
T_2\sigma_2&=(\beta+6\mu-3)\gamma+u+\mu w+\delta.
\end{aligned}
\tag{2.3b}
\]
Using the four affine transformation laws in (2.1a)--(2.1d), coefficient
comparison in (2.3b) yields
\[
\begin{aligned}
 \sigma_1\circ g_1&=-\tau^{-1}T_1\sigma_1,
 &\sigma_2\circ g_1&=T_1\sigma_2+(1-\mu)\tau^{-1}T_1\sigma_1,\\
 \sigma_1\circ g_2&=\tau^{-1}T_2\sigma_1,
 &\sigma_2\circ g_2&=T_2\sigma_2-\mu\tau^{-1}T_2\sigma_1.
\end{aligned}
\tag{2.3c}
\]
For instance, dividing the first line of (2.3b) by \(-\tau\) gives the
coefficients \(6(1-\mu)/\tau\), \((\tau-1)/\tau\), and \(1\), exactly
those of \(\sigma_1(g_1z)\); subtracting \((1-\mu)\) times that normalized
row from \(T_1\sigma_2\) gives \(\sigma_2(g_1z)\).  The two \(g_2\)
identities follow by the same coefficient check.  Thus, with
\[
R_{g_1}=\begin{pmatrix}-1/\tau&0\\(1-\mu)/\tau&1\end{pmatrix},
\qquad
R_{g_2}=\begin{pmatrix}1/\tau&0\\-\mu/\tau&1\end{pmatrix},
\]
pairing (2.3c) with \(\lambda\in\Lambda\) gives
\[
 \Pi(gz)=R_g(z)\Pi(z)M_g.
 \tag{2.3d}
\]
Since \(M_{gh}=M_hM_g\), uniqueness of the left factor gives
\(R_{gh}(z)=R_g(hz)R_h(z)\); hence these formulae define an action, not
merely generator-wise maps.  Taking real determinants in (2.3d), using
\(\det M_g=1\) and
\(\operatorname{Im}\tau(g_jz)=\operatorname{Im}\tau/|\tau|^2\), shows
\(D(gz)=D(z)\).  At the cusp,
\[
 D=\operatorname{Im}b-\operatorname{Im}\tau-
 \frac{6(\operatorname{Im}\mu)^2}{\operatorname{Im}\tau}
 \longrightarrow-\infty.
\]
The descended function is bounded above on the complement of a cusp
neighbourhood.  Replacing \(\beta\) by \(\beta+c_0\) adds
\(\operatorname{Im}c_0\) to \(D\), so a sufficiently negative imaginary
constant makes (2.2) hold globally.  Equation (2.4) then makes \(\Pi\) a real
isomorphism everywhere.  The equivariance identity supplies the claimed
\(\Delta\)-action and completes the construction of the period family.
\end{proof}

\begin{proposition}[exact admissible period constants]
\label{prop:period-exact-admissible-constants}
Let \(\tau,\mu,\beta^{\mathrm{part}}\) be the cusp-regular functions
constructed in \cref{prop:period-input-audit}, with all its marked
coordinates and transformation laws retained.  The descended function
\[
 D_{\mathrm{part}}=
 \operatorname{Im}\beta^{\mathrm{part}}
 -\frac{6(\operatorname{Im}\mu)^2}{\operatorname{Im}\tau}
\]
has an attained finite maximum \(M\) on
\(B\setminus\{p_0\}\).  The exact set of constants for which
\(\beta=\beta^{\mathrm{part}}+c_0\) makes the marked real period map an
isomorphism at every point is
\[
 \{c_0\in\mathbb C:\operatorname{Im}c_0<-M\}.
\]
\end{proposition}

\begin{proof}
The preceding proposition proves continuity on
\(B\setminus\{p_0\}\), including its elliptic points, and convergence
\(D_{\mathrm{part}}\to-\infty\) at \(p_0\).  Choose a base point
\(b_*\).  Outside a compact set containing \(b_*\), the function is
smaller than \(D_{\mathrm{part}}(b_*)-1\).  Its maximum on that compact
set is therefore its maximum on the entire base.  A constant shift gives
\(D=D_{\mathrm{part}}+\operatorname{Im}c_0\), so \(D<0\) everywhere
exactly for \(\operatorname{Im}c_0<-M\).  At equality, \(D=0\) at a
maximum point.  For \(\operatorname{Im}c_0>-M\), that point has positive
\(D\), whereas the cusp still has negative \(D\).  A path toward the
cusp contains a zero.  Since the exact ordered real determinant is
\(\operatorname{Im}\tau\,D\), these and only these excluded constants
make the period matrix singular somewhere.
\end{proof}

\begin{remark}[the cusp extension is essential]
For the same \(\tau,\mu,\beta\), the function
\(\widetilde\beta=\beta+i\,t\circ\pi\) is holomorphic on
\(\mathfrak h_z\) and satisfies all the same affine transformation laws,
because \(t\circ\pi\) is invariant.  However,
\(\widetilde D=D+\operatorname{Re}t\).  On the actual ray \(t=R\to+\infty\),
\(j(\tau)=1728R\) gives
\[
 q=(1728R)^{-1}(1+O(R^{-1})),\qquad
 D(R)=-\frac{\log(1728R)}{2\pi}
       +\operatorname{Im}b(0)+o(1),
 \quad b=\beta+\tau.
\]
Consequently \(\widetilde D(R)\to+\infty\), and no additive constant
makes \(\widetilde D\) negative everywhere.  This function is excluded
by the actual construction because
\(\widetilde\beta+\tau=b+i t\) has a cusp pole.  Thus holomorphic
solvability on the open base alone does not supply the required uniform
bound; the proved cusp extension supplies it.
\end{remark}


For the cusp filling, record the local periods explicitly.  Let \(t_c\) be
the coordinate vanishing at \(p_0\).  The simple pole of \(j\) gives
\[
 q=e^{2\pi i\tau}=t_cu_1(t_c),\qquad u_1(0)\ne0.
\]
After choosing a logarithm of \(u_1\), put
\[
 h=(2\pi i)^{-1}\log u_1,\qquad s=\tau-h,
 \qquad b=\beta+\tau.
\]
Then \(e^{2\pi is}=t_c\), \(s\circ g_0=s-1\), and \(\mu,b,h\) are
holomorphic functions of \(t_c\) at the origin.  Consequently
\[
 B_0=\begin{pmatrix}0&1\\-1&0\end{pmatrix},\qquad
 Z=sB_0+C(t_c),\qquad
 C(t_c)=
 \begin{pmatrix}
 6\mu(t_c)&h(t_c)\\
 b(t_c)-h(t_c)&\mu(t_c)
 \end{pmatrix},
 \tag{2.5a}
\]
and \(C\) is holomorphic on a full cusp disc.  In particular, for
\(\bar\lambda\in\bar\Lambda\simeq\mathbb Z^2\),
\[
 \exp(2\pi iZ\bar\lambda)
 =t_c^{B_0\bar\lambda}c_{\bar\lambda}(t_c),
 \qquad
 c_{\bar\lambda}(t_c)=\exp(2\pi iC(t_c)\bar\lambda).
 \tag{2.5b}
\]

\subsection{The toric filling at the cusp}

Let \(\mathcal T_{A_2}\) be the translation-invariant triangulation of
\(\mathbb R^2\) by the two triangles
\(\{v,v+e_1,v+e_2\}\) and
\(\{v+e_1,v+e_2,v+e_1+e_2\}\).  In
\(N'=\mathbb Z^2\oplus\mathbb Z\), take the cones over all cells at height
one, and denote the associated toric threefold by \(Y\).
We use the standard orbit--cone and affine-chart constructions of
\cite[Chapters~1 and~3]{Fulton1993}; all nonstandard assertions caused by
the infinite fan are proved directly below.

For \(\bar\lambda\in\bar\Lambda\), define the integral linear map
\[
 \varphi_{\bar\lambda}(y,y_3)
 =(y+y_3B_0\bar\lambda,y_3).
\]
On the height-one slice this is translation by
\(B_0\bar\lambda\in\mathbb Z^2\), so it preserves the triangulation and
maps the cone over a cell \(P\) to the cone over
\(P+B_0\bar\lambda\).  It therefore induces a toric automorphism
\(\Phi_{\bar\lambda}\) of \(Y\), satisfying
\[
 \Phi_{\bar\lambda}(x,t)
 =(t^{B_0\bar\lambda}x,t)
\]
on the dense torus.  Because \(t\) and \(C(t)\) are holomorphic on the
full cusp model, multiplication by
\((c_{\bar\lambda}(t),1)\in(\mathbb C^*)^3\) defines a biholomorphism on
all of \(Y_\epsilon\).  Set
\[
 \Psi_{\bar\lambda}(p)
 =(c_{\bar\lambda}(t(p)),1)\cdot\Phi_{\bar\lambda}(p).
\]
The maps \(\Phi_{\bar\lambda}\) fix the first two cocharacter directions
and preserve \(t\), so they commute with this torus multiplication.  The
linearity of \(C(t)\bar\lambda\) gives
\(c_{\bar\lambda}c_{\bar\mu}=c_{\bar\lambda+\bar\mu}\); hence
\(\Psi_{\bar\lambda}\Psi_{\bar\mu}
=\Psi_{\bar\lambda+\bar\mu}\).  This proves, before taking any quotient,
both the holomorphic extension and the action law used below.

For a triangle \(\sigma=\{v_0,v_1,v_2\}\), let
\(\theta_i^\sigma\) be its barycentric coordinates and let
\(z_i=\chi^{m_i}\) be the coordinates dual to
\((v_i,1)\).  Set
\[
 \sigma^{(\eta)}=\{y:\theta_i^\sigma(y)>-\eta\text{ for }i=0,1,2\},
 \quad
 \mathcal U_\sigma(S)=\{|z_i|<S\text{ for all }i\}.
\]
We use the restricted charts
\[
 \mathcal U_\sigma(S,\epsilon)=\mathcal U_\sigma(S)\cap Y_\epsilon
\]
whenever a deck transformation acts on a chart.

\begin{lemma}[toric chart control]
\label{lem:toric-chart-control}
For a torus point \(p=(x,t)\) with \(0<|t|<1\), put
\[
 y(p)=\frac{(\log|x_1|,\log|x_2|)}{\log|t|}.
\]
Then
\[
 p\in\mathcal U_\sigma(S)
 \quad\Longleftrightarrow\quad
 y(p)\in\sigma^{(\log S/|\log|t||)}.
 \tag{2.5c}
\]
The exact closed-unit-polydisc version is
\begin{equation}
 p\in\overline{\mathcal U}_\sigma(1)
 \quad\Longleftrightarrow\quad
 \theta_i^\sigma(y(p))\ge0\ (i=0,1,2)
 \quad\Longleftrightarrow\quad y(p)\in\sigma.
\label{eq:cusp-closed-chart-identity}
\end{equation}
Every point of \(\{|t|<1\}\subset Y\) belongs to
\(\overline{\mathcal U}_\sigma(1)=\{|z_i|\le1\}\) for some \(\sigma\).
Consequently the sets \(\mathcal U_\sigma(2,\epsilon)\) cover \(Y_\epsilon\), and
every compact subset is contained in the union of finitely many members of
this cover.
\end{lemma}

\begin{proof}
Since \((y,1)=\sum_i\theta_i^\sigma(y)(v_i,1)\), duality gives
\[
 \log|z_i(p)|=\log|t(p)|\,\theta_i^\sigma(y(p)).
\]
Because \(\log|t|<0\), this is exactly (2.5c); retaining non-strict
inequalities proves the closed-unit-polydisc assertion, including triangle
edges and vertices.

For the boundary assertion let \(p\) lie in the toric orbit
\(O_{\widehat P}\) of the cone over a cell \(P\).  In a chart with
\(\sigma\supseteq P\), the coordinates corresponding to vertices of
\(P\) vanish.  On the quotient lattice
\(N'/(N'\cap\operatorname{span}\widehat P)\), the remaining coordinates
are the characters dual to the image cone \(\bar\sigma\).  Hence
\(p\in\overline{\mathcal U}_\sigma(1)\) exactly when
\(-\operatorname{Log}_{\widehat P}(p)\in\bar\sigma\).
The cones \(\bar\sigma\), for triangles \(\sigma\supseteq P\), cover the
quotient real vector space: lift a vector \(n\), choose a relative interior
point \(p_0\in P\times\{1\}\), and for sufficiently small \(u>0\) place
the height-one rescaling of \(p_0+un\) in a triangle of the locally finite
triangulation.  That triangle contains \(P\), and its quotient cone contains
\(n\).  This proves the asserted cover.  Compactness then supplies a finite
subcover.
\end{proof}

\begin{theorem}[the cusp quotient and its central fibre]
\label{thm:candidate-cusp-geometry}
\Verified\ For the period functions and holomorphic cusp matrix \(C(t)\)
constructed above, and sufficiently small \(\epsilon>0\), the maps
\[
 \Psi_{\bar\lambda}(x,t)=
 \bigl(c_{\bar\lambda}(t)t^{B_0\bar\lambda}x,t\bigr),
 \qquad
 c_{\bar\lambda}(t)=\exp(2\pi iC(t)\bar\lambda),
\tag{2.5}
\]
extend to a free properly discontinuous action of
\(\bar\Lambda\simeq\mathbb Z^2\) on
\(Y_\epsilon=\{|t|<\epsilon\}\).  The quotient
\[
 f_0:N_0=Y_\epsilon/\bar\Lambda\longrightarrow\Delta_\epsilon
\]
is a proper holomorphic map from a connected smooth complex threefold.
Its nonzero fibres are the tori of \cref{prop:period-input-audit}.  Its full
scheme-theoretic central fibre \(W\) is reduced and has local equations
\[
 z_0,\qquad z_0z_1,\qquad z_0z_1z_2.
\tag{2.6}
\]
The surface \(W\) is irreducible and nonnormal.  Its normalization is
\(\nu:E_0\simeq\mathrm{dP}_6\to W\), obtained by identifying opposite sides
of the toric hexagon; its double locus consists of three rational curves
through exactly two triple points.  Thus \(W\) is a normal-crossings surface,
but it is not a simple-normal-crossings divisor.  Moreover
\[
 \omega_{N_0}\simeq\mathcal O_{N_0}.
\tag{2.6a}
\]
\end{theorem}

\begin{proof}
For a bounded cell \(P\), the cone over \(P\times\{1\}\) is strongly
convex, and two such cones meet in the cone over the common face of the two
cells.  The locally finite triangulation therefore defines a countable
rational fan with support
\((\mathbb R^2\times\mathbb R_{>0})\cup\{0\}\).  For the determinant
calculation, order the height-one generators of the two kinds as
\[
 ((v,1),(v+e_1,1),(v+e_2,1)),\qquad
 ((v+e_1,1),(v+e_2,1),(v+e_1+e_2,1)).
\]
Their determinants are respectively \(+1\) and \(-1\), hence have absolute
value one.  Every cone is therefore unimodular.  Here is the pairwise
affine-chart separation argument, without assuming finiteness of the fan.
For two cones \(\kappa,\kappa'\), let \(\upsilon=\kappa\cap\kappa'\)
and \(S_\kappa=\kappa^\vee\cap M'\), where \(M'=\operatorname{Hom}(N',\mathbb Z)\).
The separation lemma for cones meeting in a common face gives an integral
character \(m\) nonnegative on \(\kappa\), nonpositive on \(\kappa'\),
and zero on each exactly along \(\upsilon\)
\cite[Proposition~3.1.3 and its separation lemma, pp.~106--108]{CoxLittleSchenck2011}.
For \(a\in S_\upsilon\), choose an integer \(k\ge0\) large enough
that \(a+km\) is nonnegative on every ray of \(\kappa\).  On the
common face this already holds; on each other ray \(m\) is positive.
Then \(a=(a+km)+(-km)\in S_\kappa+S_{\kappa'}\).  The reverse
inclusion follows by evaluation on \(\upsilon\).  Thus the map
\[
 \mathbb C[S_\kappa]\otimes\mathbb C[S_{\kappa'}]
 \longrightarrow\mathbb C[S_\upsilon],\qquad
 \chi^a\otimes\chi^b\longmapsto\chi^{a+b}
\]
is surjective.  It identifies the gluing graph with a closed algebraic
subset of the product of these two affine charts, hence with an
analytically closed subset.  These chart products cover \(Y\times Y\),
so its diagonal is closed.  Countability of the charts supplies a
countable topological basis.  Thus \(Y\) is a
connected Hausdorff smooth complex threefold with dense torus.

If \(m_0,m_1,m_2\) is the basis dual to the height-one generators, then the
corresponding affine chart is \(\mathbb C^3\), with coordinates
\(z_i=\chi^{m_i}\), and
\[
 t=z_0z_1z_2.
\tag{2.7}
\]
Since every primitive ray generator has height one, the divisor of \(t\) is
\(\sum_vE_v\), all with multiplicity one.  Equation (2.7) proves that the
scheme-theoretic fibre \(t^{-1}(0)\) is reduced and normal crossings.
The star of a height-one vertex has rays
\(e_1,e_2,e_2-e_1,-e_1,-e_2,e_1-e_2\), the complete smooth hexagonal fan;
hence every irreducible component \(E_v\) is \(\mathrm{dP}_6\).

The local canonical trivialization is also intrinsic to this toric model.
On the dense torus put
\[
 \Omega=\frac{dx_1\wedge dx_2\wedge dt}{x_1x_2}.
\]
If \((z_0,z_1,z_2)\) is a unimodular toric chart, then
\[
 \bigwedge_i\frac{dz_i}{z_i}
 =\pm\frac{dx_1}{x_1}\wedge\frac{dx_2}{x_2}\wedge\frac{dt}{t},
 \qquad t=z_0z_1z_2,
\]
and therefore \(\Omega=\pm dz_0\wedge dz_1\wedge dz_2\).  It extends
holomorphically and nowhere vanishingly across every boundary chart.
For a deck transformation \(\Psi_{\bar\lambda}\), formula (2.5) gives
\(\Psi_{\bar\lambda}^*(dx_i/x_i)=dx_i/x_i+\alpha_i(t)dt\) and
\(\Psi_{\bar\lambda}^*dt=dt\); every additional term in the pullback of
\(\Omega\) contains \(dt\) twice.  Hence \(\Omega\) is invariant and descends
to a nowhere-vanishing three-form on \(N_0\), proving the intrinsic
trivialization (2.6a).  No global threefold has been invoked at this stage.

Choose a radius \(r>0\) on which the full matrix \(C(t)\) is holomorphic,
and \(0<r_1<\min(r,1)\).  Put
\[
 R(t)=-2\pi\operatorname{Im}C(t),\qquad
 M=\max_{|t|\le r_1}\|R(t)\|<\infty.
\]
Choose \(0<\epsilon<r_1\) with \(|\log\epsilon|>2M\).
At a torus point fixed by a nonzero
\(\bar\lambda\), taking logarithmic absolute values in (2.5) gives
\[
 \log|t|B_0\bar\lambda+R(t)\bar\lambda=0,
\]
impossible because \(|B_0\bar\lambda|=|\bar\lambda|\) and the first term
has more than twice the norm of the second.  On \(t=0\), a cell is translated
by \(B_0\bar\lambda\ne0\); a bounded cell cannot be fixed by a nonzero
translation.  Thus the action is free.

For proper discontinuity put
\[
 B_t=B_0+\frac{R(t)}{\log|t|}.
\]
Then \(\|B_t-B_0\|\le1/2\), so
\(|B_t\bar\lambda|\ge|\bar\lambda|/2\).  In logarithmic position
coordinates, (2.5) is
\[
 y(\Psi_{\bar\lambda}p)=y(p)+B_t\bar\lambda.
\tag{2.8}
\]
Put \(\eta_0=\log2/|\log\epsilon|\).  If
\(\Psi_{\bar\lambda}\mathcal U_\sigma(2,\epsilon)\) meets
\(\mathcal U_{\sigma'}(2,\epsilon)\), the intersection is open in
\(Y_\epsilon\) and hence contains a
point of the dense invariant torus.  By \cref{lem:toric-chart-control}, for
that point \(p\) and its image \(p'\),
\[
 y(p)\in\sigma^{(\eta_0)},\qquad
 y(p')\in{\sigma'}^{(\eta_0)}.
\]
Both thickenings are bounded, so (2.8) gives
\[
 \frac12|\bar\lambda|
 \le |B_t\bar\lambda|
 =|y(p')-y(p)|
 \le \sup_{a\in{\sigma'}^{(\eta_0)},\,b\in\sigma^{(\eta_0)}}|a-b|.
\]
Thus only finitely many \(\bar\lambda\) can make this fixed pair of charts
meet.  A compact set is contained in finitely many of the charts from
\cref{lem:toric-chart-control}; taking the finite union over all chart pairs
proves proper discontinuity.  Freeness and proper discontinuity give a
Hausdorff complex-manifold quotient and a holomorphic covering map.
For completeness the restricted total space is connected: every maximal
chart intersects it in
\[
 \{(z_0,z_1,z_2)\in\mathbb C^3:|z_0z_1z_2|<\epsilon\},
\]
a set star shaped about the origin.  These connected open sets meet the
common connected dense torus
\((\mathbb C^*)^2\times\{0<|t|<\epsilon\}\).  Their union
\(Y_\epsilon\), and hence its quotient, is connected.

To prove properness, fix \(\epsilon_1<\epsilon\).  Let \(\mathcal F_0\) be
the finite set of triangles meeting the closed Euclidean ball of radius
\(\sqrt2\), and put
\[
 K_0=\bigcup_{\sigma\in\mathcal F_0}
 \left(\overline{\mathcal U}_\sigma(1)
 \cap\{|t|\le\epsilon_1\}\right).
\]
This is compact.  For a torus point \(p\) with
\(0<|t(p)|\le\epsilon_1\), put \(u=B_{t(p)}^{-1}y(p)\) and choose
\(\bar\lambda\in\mathbb Z^2\) with
\(|u+\bar\lambda|_\infty\le1/2\).  Since
\(\|B_t\|\le3/2\), equation (2.8) gives
\[
 |y(\Psi_{\bar\lambda}p)|
 =|B_t(u+\bar\lambda)|
 \le\frac{3}{2\sqrt2}<\sqrt2.
\]
Choose a triangle containing this logarithmic position.  The closed-unit-
polydisc identity of \cref{lem:toric-chart-control} places the translate
in the corresponding
\(\overline{\mathcal U}_\sigma(1)\subset K_0\).

For the central representatives, we prove the needed closed-polydisc
cover explicitly.  Denote the two ordered triangles at \(v\) by
\[
 L_v=(v,v+e_1,v+e_2),\qquad
 U_v=(v+e_1,v+e_2,v+e_1+e_2).
\]
On the dense orbit of \(E_0\), the characters \((1,0,0)\) and
\((0,1,0)\) give the original coordinates \(x_1,x_2\).  The coordinate
dual to \((0,0,1)\) vanishes there; the other two dual characters have
last exponent zero.  Pairing with the ordered height-one generators gives
the following restrictions of the chart coordinates.  Write
\(a=\log|x_1|\), \(b=\log|x_2|\) on this dense orbit.
\[
\begin{array}{c|c|c}
 \sigma&(z_0,z_1,z_2)|_{E_0^{\rm tor}}&|z_i|\le1\ (i=0,1,2)\\\hline
 L_0&(0,x_1,x_2)&a\le0,\ b\le0\\
 L_{-e_1}&((x_1x_2)^{-1},0,x_2)&a+b\ge0,\ b\le0\\
 L_{-e_2}&((x_1x_2)^{-1},x_1,0)&a+b\ge0,\ a\le0\\
 U_{-e_1}&(0,x_1^{-1},x_1x_2)&a\ge0,\ a+b\le0\\
 U_{-e_2}&(x_2^{-1},0,x_1x_2)&b\ge0,\ a+b\le0\\
 U_{-e_1-e_2}&(x_2^{-1},x_1^{-1},0)&a\ge0,\ b\ge0.
\end{array}
\]
The six regions cover \(\mathbb R^2\): the first and last cover equal
signs; for \(a\ge0\ge b\), the sign of \(a+b\) selects the second or
fourth; for \(b\ge0\ge a\), it selects the third or fifth.  Equalities
are retained throughout.  Each corresponding coordinate plane meets
\(\overline{\mathcal U}_\sigma(1)\) in a compact closed bidisc
\(Q_\sigma\subset E_0\).  Their finite union is compact and hence closed
in the already Hausdorff \(Y\).  It contains \(E_0^{\rm tor}\), so it
contains its closure \(E_0\); the reverse inclusion follows from
\(Q_\sigma\subset E_0\).  Consequently
\[
 E_0=\bigcup_{\sigma\ni0}Q_\sigma.
\]
This proves the closed-polydisc cover at every boundary point as well as
on the dense orbit, without presupposing compactness of \(E_0\).

Now let \(p\) lie in the central orbit belonging to a cell \(P\), choose
a vertex \(v\in P\), and set \(\bar\lambda=B_0v\).  The exact map
\(\Psi_{B_0v}\) translates the cell to \(P-v\), because \(B_0^2v=-v\),
and retains the full torus multiplier
\[
 c_{B_0v}(0)=\exp\bigl(2\pi iC(0)B_0v\bigr).
\]
This nonzero multiplier preserves the orbit label, so the actual
translated point belongs to \(E_0\).  The proved cover applies to it
regardless of the multiplier's modulus or argument.  Each of the six
triangles contains zero and therefore belongs to \(\mathcal F_0\), so
the translate lies in \(K_0\).  Together with the nonzero-fibre
representatives above, this proves that the image of \(K_0\) in the
quotient is exactly \(f_0^{-1}(\{|t|\le\epsilon_1\})\), which is compact.
Every compact subset of the open base disc is closed in one such closed
subdisc; its inverse image is therefore closed in that compact Hausdorff
inverse image.  Thus \(f_0\) is proper.

For \(t\ne0\), choose \(s\) with \(e^{2\pi is}=t\).  The multiplier in
(2.5) is \(\exp(2\pi iZ(s)\bar\lambda)\), so exponentiation identifies the
fibre with
\[
 \mathbb C^2/(\mathbb Z^2+Z(s)\mathbb Z^2).
\]
It is compact because \(\operatorname{Im}Z(s)\) is invertible by (2.2).

The quotient map is a local biholomorphism; therefore the local equation
(2.7), its multiplicity one, and reducedness descend to \(W\).  Every cell
has a vertex, and \(B_0\bar\Lambda=\mathbb Z^2\) acts simply transitively on
vertices, so every orbit in the central fibre meets \(E_0\).  Thus
\(E_0\to W\) is surjective.  Two interior points cannot be identified; on
the boundary, precisely opposite edges and the corresponding hexagon
vertices are identified.  The fibres have at most three points.  Since
\(E_0\) is compact, the map is finite; since it is generically injective
from a smooth irreducible surface to the reduced surface \(W\), it is the
normalization.  The two translation-orbits of triangles give the two triple
points, and the three unoriented edge directions give the three double
curves.  Irreducibility follows from the finite surjection by \(E_0\).

The actual opposite-side maps retain the entire matrix \(C(0)\).
For a neighboring vertex
\(v\in\{\pm e_1,\pm e_2,\pm(e_1-e_2)\}\), put
\[
 C_v=E_0\cap E_v,\qquad
 \bar\lambda_v=B_0v,\qquad
 \iota_v=\Psi_{B_0v}|_{C_v}:C_v\longrightarrow C_{-v}.
\]
Since \(B_0^2v=-v\), this is precisely the translation taking the edge
\(\{0,v\}\) to \(\{0,-v\}\).  For a primitive row \(a\in\mathbb Z^2\)
with \(av=0\), use the same character \(w=\chi^{(a,0)}\) on both
dense edge orbits.  Character pullback gives
\begin{equation}
 \iota_v^*w_{-v}=\exp\bigl(2\pi i\,aC(0)B_0v\bigr)w_v,
 \label{eq:cusp-exact-opposite-edge}
\end{equation}
because the exponent of \(t\) is \(-av=0\).  With
\(\mu_0=\mu(0)\) and \(b_0=b(0)\), the three maps are
\[
\begin{array}{c|c|c}
 v&a&\iota_v^*w_{-v}/w_v\\\hline
 e_1&(0,1)&e^{-2\pi i\mu_0}\\
 e_2&(1,0)&e^{12\pi i\mu_0}\\
 e_1-e_2&(1,1)&e^{-2\pi i(7\mu_0+b_0)}.
\end{array}
\]
For the third entry the full exponent coefficient is
\[
 (1,1)C(0)B_0(e_1-e_2)
 =-6\mu_0-h(0)-(b_0-h(0))-\mu_0=-7\mu_0-b_0.
\]
The cancellation of \(h(0)\) occurs in this exact character evaluation;
no entry was suppressed.  All multipliers are nonzero.  The maps extend
over both endpoints as restrictions of the ambient biholomorphisms, and
\(\iota_{-v}\) is the inverse map with reciprocal multiplier in the
same-character convention.
\end{proof}

\begin{proposition}[the cusp punctured-overlap biholomorphism]
\label{prop:candidate-cusp-overlap}
\Verified\ Let \(D_0^*=\{0<|t_c|<\epsilon\}\), and shrink it if necessary
so that its inverse image has a distinguished component \(U_0\) stabilized
exactly by \(\langle g_0\rangle\).  Exponentiation induces a
fibre-preserving biholomorphism
\[
 J|_{D_0^*}\xrightarrow{\ \sim\ }
 N_0|_{D_0^*}=N_0\setminus W.
\tag{2.8i}
\]
On the covering space it is given by
\[
 \mathcal E:U_0\times\mathbb C^2\longrightarrow
             (\mathbb C^*)^2\times D_0^*,
 \qquad
 \mathcal E(z,\zeta_1,\zeta_2)
 =\bigl(e^{2\pi i\zeta_1},e^{2\pi i\zeta_2},e^{2\pi is(z)}\bigr).
\tag{2.8ii}
\]
\end{proposition}

\begin{proof}
First identify every deck transformation.  From
\(s\circ g_0=s-1\), the equality
\(t_c\circ\pi=e^{2\pi is}\), and the fact that \(t_c\) is a base
coordinate, \(ds\) never vanishes on \(U_0\).  Thus \(s\) is a local
biholomorphism to the logarithmic cover of \(D_0^*\).  The restriction
\(U_0\to D_0^*\) is a connected regular cover with deck group
\(\langle g_0\rangle\simeq\mathbb Z\).  Under the classification of connected
covers of the punctured disc, it corresponds to a subgroup of
\(\pi_1(D_0^*)\simeq\mathbb Z\); because its deck group is infinite, that
subgroup is zero.  Hence \(U_0\) is the universal cover.  The logarithmic
cover is universal as well, and its deck generator acts by
\(s\mapsto s-1\).  Therefore \(s\) is an equivariant biholomorphism onto that
logarithmic cover.

The cusp transformation laws give
\[
 \Pi(g_0z)=[Z(s-1)\mid I_2]=[Z(s)-B_0\mid I_2].
\]
In the same block decomposition, the explicit matrix in
\cref{prop:monodromy-certificate} gives
\[
 M_{g_0}=M_0^{-1}
 =\begin{pmatrix}I_2&0\\-B_0&I_2\end{pmatrix}.
\]
Consequently
\[
 \Pi(g_0z)=\Pi(z)M_{g_0}.
\]
Comparison with (2.3d) and real invertibility of \(\Pi(g_0z)\) yields
\(R_{g_0}=I_2\).  To make the crossed action explicit, for
\(\lambda\in\Lambda\) and \(g\in\Delta\) put
\[
 T_\lambda(z,\zeta)=(z,\zeta+\Pi(z)\lambda),
 \qquad
 \widetilde g(z,\zeta)=(gz,R_g(z)\zeta).
\]
Equation (2.3d), together with \(A_g=M_g^{-1}\), gives
\[
\begin{aligned}
 \widetilde gT_\lambda(z,\zeta)
 &=\bigl(gz,R_g(z)\zeta+R_g(z)\Pi(z)\lambda\bigr)\\
 &=\bigl(gz,R_g(z)\zeta+\Pi(gz)A_g\lambda\bigr)
  =T_{A_g\lambda}\widetilde g(z,\zeta).
\end{aligned}
\]
Thus
\[
 \widetilde gT_\lambda=T_{A_g\lambda}\widetilde g.
\tag{2.8ii-a}
\]
At the cusp \(A_{g_0}=M_0\), so
\[
 \widetilde g_0(z,\zeta)=(g_0z,\zeta),
 \qquad
 \widetilde g_0T_\lambda=T_{M_0\lambda}\widetilde g_0.
\tag{2.8ii-b}
\]
In particular \(\mathcal E\) is invariant under \(\widetilde g_0\).

Let
\(\Lambda_{\rm tor}=\mathbb Z\widehat w\oplus
\mathbb Z\widehat\delta\).  By \cref{thm:cusp-lattices},
\[
 \Lambda_{\rm tor}=\ker(M_0-I)=\operatorname{im}(M_0-I).
\tag{2.8ii-c}
\]
Put
\[
 \mathcal G_0=\langle T_\lambda,\widetilde g_0:\lambda\in\Lambda\rangle
 \simeq\Lambda\rtimes_{M_0}\langle\widetilde g_0\rangle.
\]
Since \(M_0\) fixes \(\Lambda_{\rm tor}\) pointwise, (2.8ii-b) shows that
\[
 H:=\{T_\nu\widetilde g_0^{\,n}:
          \nu\in\Lambda_{\rm tor},\ n\in\mathbb Z\}
 \simeq\Lambda_{\rm tor}\times\langle\widetilde g_0\rangle.
\tag{2.8ii-d}
\]
Moreover, for \(\lambda\in\Lambda\) and
\(\nu\in\Lambda_{\rm tor}\),
\[
 \widetilde g_0T_\nu\widetilde g_0^{-1}=T_\nu,
 \qquad
 T_\lambda T_\nu T_\lambda^{-1}=T_\nu,
 \qquad
 T_\lambda\widetilde g_0T_\lambda^{-1}
 =T_{(I-M_0)\lambda}\widetilde g_0\in H.
\tag{2.8ii-e}
\]
These generator conjugations prove that \(H\triangleleft\mathcal G_0\).
Translation by \(\Lambda_{\rm tor}\) adds an integral vector to \(\zeta\),
so it lies in the deck group of \(\mathcal E\).
Conversely, if two points have the same image under \(\mathcal E\), their
\(s\)-coordinates differ by an integer and their \(\zeta\)-coordinates by
an element of \(\mathbb Z^2\).  Equivariance and injectivity of the
logarithmic coordinate then show that they differ by a unique element of \(H\).
Thus \(\mathcal E\) identifies
\((U_0\times\mathbb C^2)/H\) with
\((\mathbb C^*)^2\times D_0^*\).

It remains to identify the residual action.  Write
\(\bar\Lambda=\Lambda/\Lambda_{\rm tor}\), and let \(\lambda\) be any lift
of \(\bar\lambda\).  In the period basis, translation by \(\lambda\) adds
\(Z(s)\bar\lambda\) to \(\zeta\), up to an integral vector killed by
exponentiation.  Therefore (2.5b) gives
\[
 \mathcal E(z,\zeta+\Pi(z)\lambda)
 =\bigl(c_{\bar\lambda}(t_c)t_c^{B_0\bar\lambda}x,t_c\bigr)
 =\Psi_{\bar\lambda}\bigl(\mathcal E(z,\zeta)\bigr).
\tag{2.8iii}
\]
This is independent of the chosen lift.  Moreover \(M_0\) is the identity
on \(\bar\Lambda\), by (2.8ii-c).  The crossed multiplication law is
\[
 (T_\lambda\widetilde g_0^{\,n})
 (T_\mu\widetilde g_0^{\,m})
 =T_{\lambda+M_0^n\mu}\widetilde g_0^{\,n+m}.
\]
It follows that
\[
 \vartheta:\mathcal G_0\longrightarrow\bar\Lambda,
 \qquad
 \vartheta(T_\lambda\widetilde g_0^{\,n})=\bar\lambda,
\]
is a surjective homomorphism with kernel \(H\).  Hence
\(\mathcal G_0/H\simeq\bar\Lambda\), and (2.8iii) identifies this residual
action with the displayed \(\Psi_{\bar\lambda}\)-action.  Dividing by it
gives
\[
 (U_0\times\mathbb C^2)/\mathcal G_0
 \simeq
 ((\mathbb C^*)^2\times D_0^*)/\bar\Lambda.
\]
The left side is \(J|_{D_0^*}\), while the right side is
\(N_0|_{D_0^*}\) because \(t^{-1}(\mathbb C^*)\) is the dense torus of
\(Y\).  This proves (2.8i), including holomorphicity of the inverse.
\end{proof}

\subsection{The finite logarithmic fillings and the global gluing}

Set
\[
 \varepsilon=\widehat\gamma+2\widehat u-4\widehat w,
 \qquad
 \varepsilon'=\widehat\gamma+3\widehat u-3\widehat w.
 \tag{2.8a}
\]
Direct multiplication by the matrices in
\cref{prop:monodromy-certificate} gives
\(A_1\varepsilon=\varepsilon\) and
\(A_2\varepsilon'=\varepsilon'\).

A \emph{linearizing pair} \((\Delta_j,s_j)\) consists of a simply connected
\(g_j\)-invariant neighbourhood of \(z_j\), disjoint from all its translates
outside \(\langle g_j\rangle\), and an injective coordinate satisfying
\[
 s_j(z_j)=0,\qquad s_j\circ g_j=\zeta_js_j,
 \qquad \zeta_j=e^{-2\pi i/m_j}.
 \tag{2.8b}
\]
Choose the image of \(s_j\) to be a round disc about zero, shrinking the
linearizing neighbourhood if necessary.  Rotation invariance and
translate-disjointness are preserved by this restriction.  Put
\(D_j=\pi(\Delta_j)\),
\(\Delta_j^*=\Delta_j\setminus\{z_j\}\), and choose the base coordinate
\(t_j\) by \(t_j\circ\pi=s_j^{m_j}\).  Let
\[
 \mathcal T=(\mathfrak h_z\times\mathbb C^2)/\Lambda,
 \qquad
 \mathfrak h_z^\circ=\mathfrak h_z\setminus
 (\Delta z_1\cup\Delta z_2),
 \qquad
 J=(\mathcal T|_{\mathfrak h_z^\circ})/\Delta
 \]
\[
 J'_j=\mathcal T|_{\Delta_j},
 \qquad
 A_0^{(j)}=\mathbb C^2/\Pi(z_j)\Lambda.
 \tag{2.8c}
\]

\begin{lemma}[the middle period quotient is proper]
\label{lem:middle-period-proper}
The induced \(\Delta\)-action on
\(\mathcal T|_{\mathfrak h_z^\circ}\) is free and properly discontinuous.
Consequently \(J\) is a Hausdorff complex manifold and its natural map
\[
 J\longrightarrow B_{\rm mid}:=\mathfrak h_z^\circ/\Delta
\]
is a proper holomorphic submersion with compact torus fibres.
\end{lemma}

\begin{proof}
The deleted translates contain every elliptic fixed point, so the
\(\Delta\)-action on \(\mathfrak h_z^\circ\) is free; it is properly
discontinuous because it is the restriction of the Fuchsian action.  The
action on \(\mathcal T\) projects equivariantly to this base action.  Hence
an element fixing a point of \(\mathcal T|_{\mathfrak h_z^\circ}\) fixes
its base point and is the identity.  If \(K_1,K_2\) are compact subsets of
the total space and \(gK_1\cap K_2\ne\varnothing\), then
\(g\,\operatorname{pr}(K_1)\cap\operatorname{pr}(K_2)\ne\varnothing\).
Base proper discontinuity leaves only finitely many such \(g\), proving
proper discontinuity upstairs and hence the Hausdorff manifold assertion.

For properness, take a simply connected open set \(U\Subset B_{\rm mid}\)
evenly covered by \(\mathfrak h_z^\circ\), and one lift \(\widetilde U\).
The quotient over \(U\) is identified with
\(\mathcal T|_{\widetilde U}\).  Before quotienting, projection is a
holomorphic submersion, and the quotient maps are local biholomorphisms;
hence the displayed map is a holomorphic submersion.  The flat real-analytic trivialization
\[
 \widetilde U\times(\Lambda\otimes\mathbb R)/\Lambda
 \longrightarrow\mathcal T|_{\widetilde U},\qquad
 (z,x+\Lambda)\longmapsto[z,\Pi(z)x]
\]
shows this properness explicitly: if \(C\subset U\) is compact, its chosen
lift \(\widetilde C\subset\widetilde U\) is compact and the inverse image of
\(C\) is homeomorphic to
\(\widetilde C\times(\Lambda\otimes\mathbb R)/\Lambda\), hence compact.
Thus the map is proper over \(U\).  Properness is local on the target, so
the displayed map is proper.
\end{proof}

For \(b\in\Lambda\otimes\mathbb R\), define the varying holomorphic lift
\[
 \widetilde g_j^{b}(z,\zeta)=
 \bigl(g_jz,R_{g_j}(z)\zeta+\Pi(g_jz)b\bigr).
 \tag{2.8d}
\]
Equation (2.3d) says
\[
 R_{g_j}(z)\Pi(z)=\Pi(g_jz)A_j,
 \tag{2.8e}
\]
so (2.8d) normalizes the lattice action and descends to \(J'_j\).  In the
flat real-analytic coordinates
\[
 \Xi(z,x+\Lambda)=[z,\Pi(z)x]
\]
one has, with \(\Sigma_k=I+A_j+\cdots+A_j^{k-1}\),
\[
 \Xi^{-1}(\widetilde g_j^{b})^k\Xi(z,x)
 =\bigl(g_j^kz,A_j^kx+\Sigma_kb\bigr).
 \tag{2.8f}
\]
We call \(\widetilde g_j^{0}\) the linear lift and, for
\(v\in\Lambda^{A_j}\), call
\(\widetilde g_j^{\log}:=\widetilde g_j^{v/m_j}\) the logarithmic lift.

\begin{proposition}[finite fillings]
\label{prop:candidate-finite-fillings}
\Verified\ Let \(A_j\) be the two finite monodromies and choose
\[
 v_1=\varepsilon,\qquad v_2=-\varepsilon',
 \qquad \gamma(v_1)=1,\quad\gamma(v_2)=-1.
\]
On the local torus family over a linearizing disc \(\Delta_j\), the affine
map
\[
 (z,x)\longmapsto(g_jz,A_jx+v_j/m_j)
\tag{2.9}
\]
generates a free action of order \(m_j\).  The quotient
\(f_j:N_j\to D_j\) is a smooth complex threefold with
\[
 f_j^*(p_j)=m_jS_j,\qquad
 S_j=A_0^{(j)}/\langle x\mapsto A_jx+v_j/m_j\rangle.
\tag{2.10}
\]
Over the punctured disc it is biholomorphic to the original period family.
\end{proposition}

\begin{proof}
First prove freeness on the actual local family.  For a finite-order
integral map \(A\), the affine self-map \(x\mapsto Ax+b\) of
\((\Lambda\otimes\mathbb R)/\Lambda\) has a fixed point exactly when
\[
 \varphi(b)\in\mathbb Z
 \quad\text{for every }\varphi\in\operatorname{Hom}(\Lambda,\mathbb Z)
 \text{ with }\varphi\circ A=\varphi.
 \tag{2.9a}
\]
Necessity is immediate.  Conversely, the averaging projector
\(p=m^{-1}\sum_{r=0}^{m-1}A^r\) has image \(\ker(A-I)_{\mathbb R}\) and
kernel \(\operatorname{im}(A-I)_{\mathbb R}\); the dual lattice of
\(p(\Lambda)\) is exactly the invariant integral covector lattice.  Thus
(2.9a) puts \(pb\) in \(p(\Lambda)\), and
\(b\in(A-I)(\Lambda\otimes\mathbb R)+\Lambda\), which supplies a fixed
point.

The exact matrix calculation gives
\[
 V^{T_1^k}=\langle\gamma,\psi_1\rangle,
 \quad \psi_1=2u+w+3\delta \quad(k=1,2),
\]
and
\[
 V^{T_2^k}=\langle\gamma,\psi_2\rangle,
 \quad \psi_2=u+w+2\delta \quad(k=1,2,3).
 \tag{2.9b}
\]
For \(v_1=\varepsilon\) and \(v_2=-\varepsilon'\),
\[
 \gamma(v_1)=1,\quad\psi_1(v_1)=0,
 \qquad
 \gamma(v_2)=-1,\quad\psi_2(v_2)=0.
 \tag{2.9c}
\]
Because \(A_jv_j=v_j\), formula (2.8f) shows that the \(k\)-th power of
the logarithmic lift has translation \(kv_j/m_j\).  For \(j=1\), (2.9a)
fails for \(k=1,2\) already on \(\gamma\), since \(k/3\notin\mathbb Z\).
For \(j=2\), it fails for \(k=1,2,3\), since \(-k/4\notin\mathbb Z\).
Away from \(z_j\), the base map \(g_j^k\) has no fixed point in
\(\Delta_j\).  Hence every nonidentity power is fixed-point free on
\(J'_j\).  At \(k=m_j\), (2.8f) is translation by
\(v_j\in\Lambda\), hence the identity; the order is exactly \(m_j\)
because the map covers \(g_j\).

It follows that
\[
 N_j=J'_j/\langle\widetilde g_j^{\log}\rangle
\]
is a smooth complex threefold and the quotient map is finite \'etale of
degree \(m_j\).  The function \(s_j^{m_j}=t_j\circ\pi\) is invariant and
descends to \(f_j:N_j\to D_j\).  The projection
\(J'_j\to\Delta_j\) is proper: over a compact subset, the image of the
compact set obtained by letting \(x\) range over a fixed unit
parallelotope in \(\Lambda\otimes\mathbb R\) is the whole inverse image.
Together with finiteness of \(\Delta_j\to D_j\) and of the quotient, this
proves properness of \(f_j\).  Pulling the divisor of \(t_j\) back through
the etale quotient gives the divisor of \(s_j^{m_j}\), namely \(m_j\)
times the reduced central torus.  This proves (2.10).

It remains to identify the punctured filling with the original family and
to check that no logarithm branch has entered the answer.  On a simply
connected \(U\subset\Delta_j^*\), choose a branch \(\ell\) of
\(\log s_j\) and put
\[
 \sigma_j^\ell(z)=\frac{\ell(z)}{2\pi i}\Pi(z)v_j.
 \tag{2.9d}
\]
Changing \(\ell\) by \(2\pi in\) changes (2.9d) by the lattice vector
\(n\Pi(z)v_j\).  The classes therefore glue to a branch-independent
holomorphic section \(\sigma_j\) of \(J'_j|_{\Delta_j^*}\); translation by
that section defines a biholomorphism \(h_j\).

For the conjugacy, use \(\ell\) on \(U\) and on \(g_jU\) use the branch
\[
 \ell'(w)=\ell(g_j^{-1}w)-\frac{2\pi i}{m_j}.
\]
This is a logarithm of \(s_j(w)\) because
\(s_j\circ g_j=e^{-2\pi i/m_j}s_j\).  For \(z\in U\), equation (2.8e)
and \(A_jv_j=v_j\) give
\[
\begin{aligned}
 \widetilde g_j^{\log}h_j^{-1}[z,\zeta]
  =\Bigl[g_jz,R_{g_j}(z)\zeta
  -\Bigl(\frac{\ell(z)}{2\pi i}-\frac1{m_j}\Bigr)
    \Pi(g_jz)v_j\Bigr].
\end{aligned}
\]
But
\[
 \frac{\ell'(g_jz)}{2\pi i}
 =\frac{\ell(z)}{2\pi i}-\frac1{m_j},
\]
so applying \(h_j\) cancels the displayed translation and proves the
global identity
\[
 h_j\widetilde g_j^{\log}h_j^{-1}=\widetilde g_j^{0}.
 \tag{2.9e}
\]
Finally, by the stabilizer and translate conditions in (2.8b),
\[
 J|_{D_j^*}
 =\bigl(J'_j|_{\Delta_j^*}\bigr)/
   \langle\widetilde g_j^{0}\rangle.
\]
Thus (2.9e) descends to the asserted punctured-disc biholomorphism.
\end{proof}

\begin{lemma}[radial retraction and the finite-fibre stalk map]
\label{lem:finite-filling-stalk-map}
\Verified\ Under \cref{prop:candidate-finite-fillings}, the finite local
model \(N_j\) strongly deformation retracts onto its reduced central fibre
\(S_j\).  After identifying a nearby fibre with the marked real torus
\(A_s=(\Lambda\otimes\mathbb R)/\Lambda\) by the flat coordinates
\(\Xi\), the specialization map
\[
 H^r(S_j;\mathbb Z)\longrightarrow H^r(A_s;\mathbb Z)
\]
is the pullback \(\pi_j^*\) of the quotient covering
\[
 \pi_j:A_0^{(j)}\longrightarrow
 S_j=A_0^{(j)}/\langle x\mapsto A_jx+v_j/m_j\rangle .
\]
\end{lemma}

\begin{proof}
Use the linearizing coordinate to define the radial contraction
\(r_u:\Delta_j\to\Delta_j\) by
\[
 s_j(r_u(z))=(1-u)s_j(z),\qquad 0\le u\le1.
\]
Because \(s_j\circ g_j=\zeta_js_j\), this contraction commutes with
\(g_j\).  In the global real-analytic marking
\[
 \Xi(z,x+\Lambda)=[z,\Pi(z)x],
\]
put
\[
 H_u\bigl(\Xi(z,x+\Lambda)\bigr)
   =\Xi(r_u(z),x+\Lambda).
\]
Formula (2.8f) shows directly that \(H_u\) commutes with the logarithmic
lift: it keeps \(x\) fixed while both composites send it to
\(A_jx+v_j/m_j\).  Hence \(H_u\) descends to a strong deformation
retraction of \(N_j\) onto \(S_j\).

Fix \(t_j\ne0\) and one root \(s\) with \(s^{m_j}=t_j\).  Every orbit of
the logarithmic lift meets the fibre over \(s\) exactly once, so that fibre
maps isomorphically onto \(f_j^{-1}(t_j)\).  Following the descended radial
retraction sends the marked point \(x\in A_s\) to the orbit of the same
\(x\) in the central torus, namely \(\pi_j(x)\).  The induced map from
central-fibre cohomology to nearby-fibre cohomology is therefore exactly
\(\pi_j^*\), with no unrecorded conjugacy or index.
\end{proof}

\begin{theorem}[global analytic construction]
\label{thm:candidate-global-geometry}
\Verified\ Gluing the period family over \(B^\circ\) to
\(N_0,N_1,N_2\) by the preceding punctured-disc biholomorphisms produces a
connected compact Hausdorff complex threefold \(X\) and a proper surjective
holomorphic map
\[
 f:X\longrightarrow B\simeq\mathbb P^1
\]
with connected fibres.  The fibre over \(p_0\) is exactly the reduced
nonnormal surface \(W\) of \cref{thm:candidate-cusp-geometry}; those over
\(p_1,p_2\) have scheme multiplicities three and four.  The open embedding
\(N_0\hookrightarrow X\) identifies canonical bundles, so
\begin{equation}
 \omega_X|_{N_0}\simeq\omega_{N_0}\simeq\mathcal O_{N_0}.
\label{eq:candidate-global-canonical-on-cusp}
\end{equation}
\end{theorem}

\begin{proof}
Take the disjoint union of the four manifolds and identify the two copies of
each punctured-disc restriction by the displayed biholomorphism.  The special
discs are pairwise disjoint, so every nonsingleton equivalence class has two
points and meets each summand at most once.  The saturation of an open set is
the union of that set with its images under finitely many overlap
biholomorphisms, hence is open.  Each summand therefore embeds openly in the
quotient, and the only transition maps are biholomorphisms.  This gives a
complex threefold.

Points over distinct base values have disjoint disc preimages.  Points over
one base value lie in a common one of the four Hausdorff charts, so the
quotient is Hausdorff.  The images of countable bases of the four open
pieces give second countability.  The local maps to the base agree on
overlaps.  Moreover \(f^{-1}(D_j)\) is exactly the embedded \(N_j\),
since every middle-piece point over \(D_j^*\) was identified with its
unique representative there; similarly \(f^{-1}(B^\circ)=J\).  Each
of these full restrictions is proper, for the period-family piece by
\cref{lem:middle-period-proper} and for the three fillings by their local
proper constructions.  Explicitly, for compact \(K\subset B\), choose
finitely many open sets \(V_a\) covering \(K\), with compact closures
inside members \(U_a\) of this base cover.  Each
\(K\cap\overline V_a\) is compact in \(U_a\), so its full inverse
image is compact.  Their finite union is \(f^{-1}(K)\); hence the global
map is proper.  Since the base is compact, \(X=f^{-1}(B)\) is compact.
All four pieces are connected and every local filling meets the central
period-family piece, so \(X\) is connected.  The fibre statements follow
from \cref{thm:candidate-cusp-geometry,prop:candidate-finite-fillings}; their
fibres are connected, and surjectivity follows.
Because each summand embeds openly in the quotient, the restriction of the
intrinsic canonical bundle \(\omega_X\) to the open summand \(N_0\) is
\(\omega_{N_0}\).  Equation~(2.6a) proves
\cref{eq:candidate-global-canonical-on-cusp}.
\end{proof}

\paragraph{What has and has not been discharged.}
The formulas above reconstruct the period functions, both finite fillings,
the toric cusp quotient, and all three punctured-overlap biholomorphisms.
Thus \cref{thm:candidate-cusp-geometry,thm:candidate-global-geometry}
discharge existence, smoothness, properness,
the full central-fibre equation, and the normalization assertions formerly
labelled (H1)--(H3) in \cref{thm:candidate-conductor}.  This promotion is
local to the analytic construction: no collapse, Leray transgression, or
Mayer--Vietoris computation is used here.  Integral cohomology and sphere
recognition remain separate below.

\section{Flat spectral geometry on the positive radial cusp}
\label{cusp-spec-section}

The period family used in this section is the family already constructed and
proved nondegenerate in \cref{prop:period-input-audit}.  In particular, this
section introduces no additional period hypothesis.  We retain the proved
ordered lattices
\[
 V=\mathbb Z\langle\gamma,u,w,\delta\rangle,
 \qquad
 \Lambda=V^*
 =\mathbb Z\langle\widehat\gamma,\widehat u,
                         \widehat w,\widehat\delta\rangle
\]
and the proved period matrix
\[
 \Pi=[Z\mid I_2],
 \qquad
 Z=\begin{pmatrix}6\mu&\tau\\ \beta&\mu\end{pmatrix}.
\]
Thus the columns of \(\Pi\), in the displayed ordered basis of \(\Lambda\),
are
\[
 (6\mu,\beta)^{\mathsf t},\quad
 (\tau,\mu)^{\mathsf t},\quad
 (1,0)^{\mathsf t},\quad
 (0,1)^{\mathsf t}.
\]

The canonical cusp calculation has already produced holomorphic functions
\(\mu(t_c),b(t_c),h(t_c)\) on a full disc about \(t_c=0\), where
\[
 b=\beta+\tau,\qquad s=\tau-h,
 \qquad e^{2\pi i s}=t_c,
\]
and has proved the exact identity
\begin{equation}
 Z=s\begin{pmatrix}0&1\\-1&0\end{pmatrix}
 +\begin{pmatrix}6\mu&h\\b-h&\mu\end{pmatrix}.
 \label{cusp-spec-canonical-cusp-input}
\end{equation}
Choose \(T_0\) so that \(e^{-2\pi T}\) lies in the proved cusp disc whenever
\(T\ge T_0\).  On the positive radial ray
\[
 t_c=e^{-2\pi T},\qquad T\ge T_0,
\]
we choose the lift \(s=iT\).  The other lifts differ by the already accounted
for integral deck transformation.  Put
\begin{equation}
 \begin{aligned}
  \mu_T&:=\mu(e^{-2\pi T})=\alpha_T+i m_T,\\
  h_T&:=h(e^{-2\pi T})=\eta_T+i H_T,\\
  b_T&:=b(e^{-2\pi T})=\xi_T+i B_T,\\
  c_T&:=\xi_T-\eta_T,\qquad
  L_T:=T+H_T,\qquad q_T:=L_T-B_T,\\
  D_T&:=L_Tq_T+6m_T^2.
 \end{aligned}
 \label{cusp-spec-radial-notation}
\end{equation}
All letters on the right-hand side of
\cref{cusp-spec-radial-notation} are real.  We reserve
\begin{equation}
 m:=\operatorname{Im}\mu(0),\qquad
 H:=\operatorname{Im}h(0),\qquad
 B:=\operatorname{Im}b(0)
 \label{cusp-spec-limiting-imaginary-parts}
\end{equation}
for their limiting values.  Holomorphy at the origin gives, with the same
statement for the real parts,
\begin{equation}
 m_T=m+O(e^{-2\pi T}),\qquad
 H_T=H+O(e^{-2\pi T}),\qquad
 B_T=B+O(e^{-2\pi T}).
 \label{cusp-spec-holomorphic-errors}
\end{equation}

Identify \(\mathbb C^2\) with \(\mathbb R^4\) by the ordered coordinates
\begin{equation}
 (z_1,z_2)=(x_1+iy_1,x_2+iy_2)
 \longmapsto (x_1,y_1,x_2,y_2).
 \label{cusp-spec-real-coordinate-order}
\end{equation}
The metric in this section is, by definition, the quotient of
\(dx_1^2+dy_1^2+dx_2^2+dy_2^2\).  This makes the metric choice explicit; the
complex period construction itself neither requires nor singles out this
Euclidean metric.

\begin{proposition}[Exact radial realification and dual lattice]
\label{cusp-spec-exact-dual-lattice}
In the domain basis
\((\widehat\gamma,\widehat u,\widehat w,\widehat\delta)\) and the real
coordinate order \cref{cusp-spec-real-coordinate-order}, the positive radial
period matrix is
\begin{equation}
 P_T=
 \begin{pmatrix}
  6\alpha_T&\eta_T&1&0\\
  6m_T&L_T&0&0\\
  c_T&\alpha_T&0&1\\
  -q_T&m_T&0&0
 \end{pmatrix}.
 \label{cusp-spec-real-period-matrix}
\end{equation}
Its oriented determinant, positive covolume, and relation with the determinant
function in \cref{prop:period-input-audit} are
\begin{equation}
 \det P_T=-D_T,
 \qquad \operatorname{covol}(P_T\mathbb Z^4)=D_T>0,
 \qquad
 D_{\mathrm{can},T}=-\frac{D_T}{L_T}<0.
 \label{cusp-spec-determinant}
\end{equation}
Here
\(D_{\mathrm{can},T}=\operatorname{Im}\beta_T
-6(\operatorname{Im}\mu_T)^2/\operatorname{Im}\tau_T\) is the canonical
determinant function; in particular, \(D_T\) is not a new nondegeneracy
assumption.

The complete inverse transpose is
\begin{equation}
 P_T^{-\mathsf T}=
 \begin{pmatrix}
 0&0&1&0\\[2pt]
 \dfrac{m_T}{D_T}&\dfrac{q_T}{D_T}&
 \dfrac{-6m_T\alpha_T-q_T\eta_T}{D_T}&
 \dfrac{-m_Tc_T-q_T\alpha_T}{D_T}\\[7pt]
 0&0&0&1\\[2pt]
 -\dfrac{L_T}{D_T}&\dfrac{6m_T}{D_T}&
 \dfrac{6(L_T\alpha_T-m_T\eta_T)}{D_T}&
 \dfrac{L_Tc_T-6m_T\alpha_T}{D_T}
 \end{pmatrix}.
 \label{cusp-spec-full-inverse-transpose}
\end{equation}
Equivalently, for
\(k=(r,s,n_1,n_2)^{\mathsf t}\in\mathbb Z^4\), define
\begin{equation}
 R_T(k):=r-6\alpha_Tn_1-c_Tn_2,
 \qquad
 S_T(k):=s-\eta_Tn_1-\alpha_Tn_2.
 \label{cusp-spec-shifted-indices}
\end{equation}
Then the Euclidean dual lattice is exactly
\begin{equation}
 \begin{split}
 (P_T\mathbb Z^4)^*
 =\biggl\{\biggl(&n_1,
 \frac{m_TR_T(k)+q_TS_T(k)}{D_T},n_2,
 \frac{-L_TR_T(k)+6m_TS_T(k)}{D_T}\biggr):\\[-2pt]
 &k=(r,s,n_1,n_2)\in\mathbb Z^4\biggr\}.
 \end{split}
 \label{cusp-spec-exact-dual-lattice-formula}
\end{equation}
Consequently, if
\[
 F_T:=\mathbb R^4/P_T\mathbb Z^4
\]
 and \(\Delta_T\) is the nonnegative self-adjoint scalar Laplacian on
 \(L^2(F_T,\mathrm d\operatorname{vol}_{F_T})\), obtained from the closed
 Dirichlet form
 \[
  \mathfrak q_T[f]=\int_{F_T}|\mathrm df|^2\,
       \mathrm d\operatorname{vol}_{F_T},
  \qquad \operatorname{Dom}(\mathfrak q_T)=H^1(F_T),
 \]
 with operator domain \(\operatorname{Dom}(\Delta_T)=H^2(F_T)\), then, with
multiplicity,
\begin{equation}
 \operatorname{Spec}(\Delta_T)
 =\{\lambda_T(k):k\in\mathbb Z^4\},
 \label{cusp-spec-spectrum-set}
\end{equation}
where
\begin{equation}
 \begin{aligned}
 \lambda_T(k)=4\pi^2\biggl(&n_1^2+n_2^2\\
 &+\frac{(m_TR_T(k)+q_TS_T(k))^2
       +(-L_TR_T(k)+6m_TS_T(k))^2}{D_T^2}\biggr).
 \end{aligned}
 \label{cusp-spec-exact-eigenvalue}
\end{equation}
The eigenfunction belonging to \(k\) is the character
\begin{equation}
 \chi_{T,k}([x])
 =\exp\bigl(2\pi i\langle P_T^{-\mathsf T}k,x\rangle\bigr).
 \label{cusp-spec-characters}
\end{equation}
\end{proposition}

\begin{proof}
Substitution of \(s=iT\) into
\cref{cusp-spec-canonical-cusp-input} gives
\[
 Z_T=
 \begin{pmatrix}
  6\mu_T&h_T+iT\\ b_T-h_T-iT&\mu_T
 \end{pmatrix}.
\]
Realifying its four columns in the interlaced coordinate order
\((x_1,y_1,x_2,y_2)\) gives
\cref{cusp-spec-real-period-matrix}; in particular, the two divergent entries
have signs \(+L_T\) and \(-q_T=B_T-L_T\).

Expanding the determinant along the last two columns gives
\[
 \det P_T=-L_Tq_T-6m_T^2=-D_T.
\]
Moreover,
\[
 \operatorname{Im}\tau_T=L_T,
 \qquad \operatorname{Im}\beta_T=B_T-L_T=-q_T,
\]
so the canonical identity
\(\det_{\mathbb R}\Pi_T=\operatorname{Im}\tau_TD_{\mathrm{can},T}\)
becomes
\[
 L_TD_{\mathrm{can},T}
 =L_T\left(-q_T-\frac{6m_T^2}{L_T}\right)=-D_T.
\]
The already-proved inequalities \(L_T>0\) and
\(D_{\mathrm{can},T}<0\) imply \(D_T>0\).

To compute the inverse transpose without suppressing the real parts, solve
\(P_T^{\mathsf T}x=k\).  Its third and fourth equations give
\[
 x_1=n_1,\qquad x_3=n_2.
\]
The remaining equations are
\begin{equation}
 \begin{pmatrix}6m_T&-q_T\\L_T&m_T\end{pmatrix}
 \begin{pmatrix}x_2\\x_4\end{pmatrix}
 =\begin{pmatrix}R_T(k)\\S_T(k)\end{pmatrix}.
 \label{cusp-spec-two-by-two-system}
\end{equation}
The determinant of the matrix in
\cref{cusp-spec-two-by-two-system} is \(D_T\), and its inverse is
\[
 \frac1{D_T}
 \begin{pmatrix}m_T&q_T\\-L_T&6m_T\end{pmatrix}.
\]
This proves both \cref{cusp-spec-full-inverse-transpose} and
\cref{cusp-spec-exact-dual-lattice-formula}.

A character \(x\mapsto e^{2\pi i\langle\ell,x\rangle}\) descends to
\(F_T\) precisely when
\(\langle\ell,P_Tk\rangle\in\mathbb Z\) for every
\(k\in\mathbb Z^4\), which is precisely
\(\ell\in P_T^{-\mathsf T}\mathbb Z^4\).  These characters are the complete
Fourier basis of the compact abelian group \(F_T\), and applying the positive
Euclidean Laplacian gives \(4\pi^2\lVert\ell\rVert^2\).  Inserting
\cref{cusp-spec-exact-dual-lattice-formula} proves
\cref{cusp-spec-spectrum-set,cusp-spec-exact-eigenvalue,cusp-spec-characters}.
\end{proof}

\begin{theorem}[Shortest cusp modes and their first corrections]
\label{cusp-spec-shortest-modes}
Let
\begin{equation}
 \ell_{\gamma,T}:=P_T^{-\mathsf T}(1,0,0,0)^{\mathsf t},
 \qquad
 \ell_{u,T}:=P_T^{-\mathsf T}(0,1,0,0)^{\mathsf t}.
 \label{cusp-spec-unit-dual-modes}
\end{equation}
For all sufficiently large \(T\), every shortest nonzero vector of
\((P_T\mathbb Z^4)^*\) belongs to
\[
 \{\pm\ell_{\gamma,T},\ \pm\ell_{u,T}\}.
\]
Their exact squared lengths are
\begin{equation}
 \lVert\ell_{\gamma,T}\rVert^2
 =\frac{L_T^2+m_T^2}{D_T^2},
 \qquad
 \lVert\ell_{u,T}\rVert^2
 =\frac{q_T^2+36m_T^2}{D_T^2}.
 \label{cusp-spec-unit-mode-exact-lengths}
\end{equation}
Their exact comparison is
\begin{equation}
 D_T^2\bigl(\lVert\ell_{\gamma,T}\rVert^2
              -\lVert\ell_{u,T}\rVert^2\bigr)
 =B_T(2L_T-B_T)-35m_T^2.
 \label{cusp-spec-unit-mode-comparison}
\end{equation}
Thus the shortest Fourier indices are \(\pm(0,1,0,0)\) if the right-hand
side of \cref{cusp-spec-unit-mode-comparison} is positive,
\(\pm(1,0,0,0)\) if it is negative, and both pairs if it is zero.

In the standard basis \(e_1,e_2,e_3,e_4\) of the real coordinate space,
the vector expansions are
\begin{align}
 T\ell_{\gamma,T}
 &=-e_4+\frac{m e_2+(H-B)e_4}{T}+O(T^{-2}),
 \label{cusp-spec-gamma-vector-correction}\\
 T\ell_{u,T}
 &=e_2+\frac{-He_2+6me_4}{T}+O(T^{-2}).
 \label{cusp-spec-u-vector-correction}
\end{align}
Accordingly,
\begin{align}
 4\pi^2\lVert\ell_{\gamma,T}\rVert^2
 &=\frac{4\pi^2}{T^2}
   \left(1+\frac{2(B-H)}{T}+O(T^{-2})\right),
 \label{cusp-spec-gamma-eigenvalue-correction}\\
 4\pi^2\lVert\ell_{u,T}\rVert^2
 &=\frac{4\pi^2}{T^2}
   \left(1-\frac{2H}{T}+O(T^{-2})\right).
 \label{cusp-spec-u-eigenvalue-correction}
\end{align}
More generally, every fixed plane index has the first correction
\begin{equation}
 \frac{T^2\lambda_T(r,s,0,0)}{4\pi^2}
 =r^2+s^2+
 \frac{2(B-H)r^2-10mrs-2Hs^2}{T}+O_{r,s}(T^{-2}).
 \label{cusp-spec-general-plane-mode-correction}
\end{equation}
In particular, the first positive eigenvalue, including the coefficient after
its leading term, is
\begin{equation}
 \boxed{
 \lambda_1(F_T)
 =\frac{4\pi^2}{T^2}
  \left(1+
   \frac{-2H+2\min\{B,0\}}{T}+O(T^{-2})\right).}
 \label{cusp-spec-first-eigenvalue-correction}
\end{equation}
If \(B>0\), the bottom indices are eventually
\(\pm(0,1,0,0)\); if \(B<0\), they are eventually
\(\pm(1,0,0,0)\).  If \(B=0\), both pairs have the same displayed
first correction, and the exact decision is made by
\cref{cusp-spec-unit-mode-comparison}.  In particular, if \(B=0\) but
\(m\ne0\), the bottom indices are eventually \(\pm(1,0,0,0)\).
The first positive eigenvalue has complex multiplicity two when exactly one
pair minimizes and multiplicity four when both pairs minimize.
\end{theorem}

\begin{proof}
For \(k=(r,s,0,0)\), define
\[
 K_T:=\frac{T}{D_T}
 \begin{pmatrix}m_T&q_T\\-L_T&6m_T\end{pmatrix}.
\]
Equations \cref{cusp-spec-holomorphic-errors,cusp-spec-radial-notation}
give
\begin{equation}
 K_T=
 \begin{pmatrix}0&1\\-1&0\end{pmatrix}+O_{\mathrm{op}}(T^{-1}).
 \label{cusp-spec-scaled-dual-plane}
\end{equation}
Hence there is a constant \(C\), independent of
\((r,s)\in\mathbb R^2\), such that
\begin{equation}
 (1-C/T)\sqrt{r^2+s^2}
 \le T\lVert P_T^{-\mathsf T}(r,s,0,0)\rVert
 \le(1+C/T)\sqrt{r^2+s^2}
 \label{cusp-spec-uniform-singular-bound}
\end{equation}
for all sufficiently large \(T\).  Thus among the vectors with
\((n_1,n_2)=(0,0)\), every vector with \(r^2+s^2\ge2\) is eventually
longer than all four vectors for which \(r^2+s^2=1\).

If instead \((n_1,n_2)\ne(0,0)\),
\cref{cusp-spec-exact-dual-lattice-formula} gives
\[
 \lVert P_T^{-\mathsf T}k\rVert^2\ge n_1^2+n_2^2\ge1,
\]
whereas each of the four unit vectors just identified has length
\(T^{-1}(1+O(T^{-1}))\).  This proves the asserted exhaustion of the
shortest vectors.

Taking the first two columns of
\cref{cusp-spec-full-inverse-transpose} proves
\cref{cusp-spec-unit-mode-exact-lengths}.  Since \(q_T=L_T-B_T\),
their numerator difference is
\[
 L_T^2+m_T^2-q_T^2-36m_T^2
 =B_T(2L_T-B_T)-35m_T^2,
\]
which proves \cref{cusp-spec-unit-mode-comparison} and the exact mode
classification.

Finally,
\begin{equation}
 D_T
 =T^2+(2H-B)T+H(H-B)+6m^2+O(Te^{-2\pi T}).
 \label{cusp-spec-D-expansion}
\end{equation}
Expanding the first two columns of
\cref{cusp-spec-full-inverse-transpose} with
\cref{cusp-spec-D-expansion} gives
\cref{cusp-spec-gamma-vector-correction,cusp-spec-u-vector-correction}.
Taking their squared norms gives
\cref{cusp-spec-gamma-eigenvalue-correction,cusp-spec-u-eigenvalue-correction}.
The same multiplication for a fixed \((r,s)\) gives
\cref{cusp-spec-general-plane-mode-correction}; explicitly, the scaled
two-dimensional matrix is
\[
 \begin{pmatrix}0&1\\-1&0\end{pmatrix}
 +\frac1T\begin{pmatrix}m&-H\\H-B&6m\end{pmatrix}
 +O(T^{-2}),
\]
and twice the Euclidean inner product of its leading and first-order terms is
\(2(B-H)r^2-10mrs-2Hs^2\).
The smaller relative correction is
\[
 \min\{2(B-H),-2H\}=-2H+2\min\{B,0\},
\]
which proves \cref{cusp-spec-first-eigenvalue-correction}.
\end{proof}

\begin{corollary}[Volume and unit-volume scaling]
\label{cusp-spec-unit-volume}
The Euclidean volume is exactly
\begin{equation}
 V_T:=\operatorname{Vol}(F_T)=D_T
 =(T+H_T)(T+H_T-B_T)+6m_T^2,
 \label{cusp-spec-exact-volume}
\end{equation}
and it has the expansion \cref{cusp-spec-D-expansion}.  Define the
unit-volume metric
\begin{equation}
 \widehat g_T:=D_T^{-1/2}g_T.
 \label{cusp-spec-unit-volume-metric}
\end{equation}
Let \(\widehat\Delta_T\) denote the nonnegative self-adjoint scalar
Laplacian on
\(L^2(F_T,\mathrm d\operatorname{vol}_{\widehat g_T})\) obtained from the
closed Dirichlet form
\[
 \widehat{\mathfrak q}_T[f]
 =\int_{F_T}|\mathrm df|_{\widehat g_T}^2\,
       \mathrm d\operatorname{vol}_{\widehat g_T},
 \qquad
 \operatorname{Dom}(\widehat{\mathfrak q}_T)=H^1(F_T),
\]
with operator domain
\(\operatorname{Dom}(\widehat\Delta_T)=H^2(F_T)\).  Then the first
positive eigenvalue of \(\widehat\Delta_T\) is
\begin{equation}
 \boxed{
 \lambda_1(F_T,\widehat g_T)
 =\frac{4\pi^2}{T}
 \left(1+
  \frac{-H-B/2+2\min\{B,0\}}{T}+O(T^{-2})\right).}
 \label{cusp-spec-unit-volume-gap}
\end{equation}
Thus unit-volume normalization does not prevent the first eigenvalue from
tending to zero.  In the normalized metric the periods \(e_1,e_3\) have
length \(D_T^{-1/4}=T^{-1/2}(1+O(T^{-1}))\), while the two periods with
indices \(\widehat\gamma,\widehat u\) have length
\(T^{1/2}(1+O(T^{-1}))\).
\end{corollary}

\begin{proof}
The volume of a Euclidean lattice quotient is its covolume, so
\cref{cusp-spec-determinant} proves
\cref{cusp-spec-exact-volume}.  In real dimension four, replacing \(g\) by
\(ag\) multiplies volume by \(a^2\) and every Laplace eigenvalue by
\(a^{-1}\).  Therefore \cref{cusp-spec-unit-volume-metric} has volume one
and
\[
 \lambda_1(F_T,\widehat g_T)=D_T^{1/2}\lambda_1(F_T,g_T).
\]
Equation \cref{cusp-spec-D-expansion} gives
\[
 D_T^{1/2}
 =T\left(1+\frac{H-B/2}{T}+O(T^{-2})\right).
\]
Multiplication by \cref{cusp-spec-first-eigenvalue-correction} proves
\cref{cusp-spec-unit-volume-gap}.  The stated period-length estimates follow
directly from \cref{cusp-spec-real-period-matrix}: \(e_1,e_3\) have original
length one, whereas the \(\widehat\gamma\)- and \(\widehat u\)-columns have
original length \(T(1+O(T^{-1}))\); lengths are multiplied by
\(D_T^{-1/4}\).
\end{proof}

\begin{proposition}[Exact counting formula and its two cusp regimes]
\label{cusp-spec-counting}
For \(E\ge0\), count eigenvalues with their Fourier multiplicities by
\begin{equation}
 N_T(E):=\#\{k\in\mathbb Z^4:\lambda_T(k)\le E\}.
 \label{cusp-spec-counting-definition}
\end{equation}
Put
\[
 R_E^2:=\frac{E}{4\pi^2},\qquad
 (x)_+:=\max\{x,0\}.
\]
For every fixed \(E\),
\begin{equation}
 N_T(E)
 =\pi D_T\sum_{n\in\mathbb Z^2}
        (R_E^2-\lVert n\rVert^2)_+ +O_E(T+1).
 \label{cusp-spec-fixed-energy-count}
\end{equation}
The sum is finite.  Equivalently,
\begin{equation}
 \lim_{T\to\infty}\frac{N_T(E)}{D_T}
 =\pi\sum_{n\in\mathbb Z^2}
        (R_E^2-\lVert n\rVert^2)_+.
 \label{cusp-spec-counting-density-limit}
\end{equation}
In the first band, \(0<E<4\pi^2\), this reduces to
\begin{equation}
 N_T(E)=\frac{E}{4\pi}D_T+O_E(T+1).
 \label{cusp-spec-first-band-count}
\end{equation}

For the unit-volume metric \(\widehat g_T\), put
\begin{equation}
 \widehat N_T(E):=
 \#\{k\in\mathbb Z^4:\lambda_T(k;\widehat g_T)\le E\}.
 \label{cusp-spec-unit-volume-count-definition}
\end{equation}
For every fixed \(E>0\),
\begin{equation}
 \widehat N_T(E)
 =\frac{E}{4\pi}D_T^{1/2}+O_E(D_T^{1/4}+1),
 \qquad
 \lim_{T\to\infty}\frac{\widehat N_T(E)}{D_T^{1/2}}
 =\frac{E}{4\pi}.
 \label{cusp-spec-unit-volume-count}
\end{equation}

There is also a finite counting law at the collapsing spectral scale.  Define
\begin{equation}
 N_T^{\mathrm{coll}}(E)
 :=\#\left\{k\in\mathbb Z^4:
       \lambda_T(k)\le\frac{4\pi^2E}{T^2}\right\}.
 \label{cusp-spec-collapsing-count-definition}
\end{equation}
If \(E\) is not a sum of two integer squares, then, for all sufficiently
large \(T\),
\begin{equation}
 N_T^{\mathrm{coll}}(E)
 =\#\{(r,s)\in\mathbb Z^2:r^2+s^2\le E\}.
 \label{cusp-spec-collapsing-count}
\end{equation}
At a threshold \(E\in\{r^2+s^2:r,s\in\mathbb Z\}\), the sharp statement
available from the zeroth-order limit is
\begin{equation}
 \begin{split}
 \#\{(r,s):r^2+s^2<E\}
 &\le\liminf_{T\to\infty}N_T^{\mathrm{coll}}(E)\\
 &\le\limsup_{T\to\infty}N_T^{\mathrm{coll}}(E)
 \le\#\{(r,s):r^2+s^2\le E\}.
 \end{split}
 \label{cusp-spec-threshold-count}
\end{equation}
For a threshold vector \((r,s)\), the coefficient in
\cref{cusp-spec-general-plane-mode-correction} decides its inclusion whenever
that coefficient is nonzero.  If it vanishes, the exact eigenvalue formula
and, if necessary, higher Taylor coefficients decide the inclusion.
\end{proposition}

\begin{proof}
Set
\begin{equation}
 \mathcal A_T:=\begin{pmatrix}6m_T&-q_T\\L_T&m_T\end{pmatrix},
 \qquad
 \theta_T(n_1,n_2):=
 \begin{pmatrix}6\alpha_Tn_1+c_Tn_2\\
                 \eta_Tn_1+\alpha_Tn_2\end{pmatrix}.
 \label{cusp-spec-counting-matrices}
\end{equation}
Then \(\det\mathcal A_T=D_T\), and the exact eigenvalue formula can be
rewritten as
\begin{equation}
 \frac{\lambda_T(r,s,n_1,n_2)}{4\pi^2}
 =\lVert(n_1,n_2)\rVert^2+
  \left\lVert\mathcal A_T^{-1}
  \left(\begin{pmatrix}r\\s\end{pmatrix}
       -\theta_T(n_1,n_2)\right)\right\rVert^2.
 \label{cusp-spec-counting-eigenvalue-form}
\end{equation}
Fix \(n=(n_1,n_2)\).  If \(\lVert n\rVert^2>R_E^2\), this index contributes
nothing.  Otherwise put
\(\rho_n=(R_E^2-\lVert n\rVert^2)^{1/2}\).  The allowed pairs
\((r,s)\) are the integer points in the translated ellipse
\begin{equation}
 \theta_T(n)+\mathcal A_T\overline B_{\mathbb R^2}(0,\rho_n).
 \label{cusp-spec-counting-ellipse}
\end{equation}
Its area is \(\pi D_T\rho_n^2\).  Its boundary length is \(O_E(T+1)\),
because \(\lVert\mathcal A_T\rVert_{\mathrm{op}}=O(T+1)\).
If \(\rho_n=0\), the set in \cref{cusp-spec-counting-ellipse} is one point
and contributes either zero or one integer point, which is absorbed in the
\(O_E(1)\) remainder below.  We apply the planar area--perimeter estimate only
when \(\rho_n>0\).

For completeness, the elementary planar lattice estimate used here is
\[
 \#(\mathbb Z^2\cap K)=\operatorname{area}(K)
 +O(\operatorname{length}(\partial K)+1)
\]
for every translated compact convex set \(K\).  Indeed, attach to each
integer point its centered unit square and let \(U\) be the union of the
squares centered at points of \(K\cap\mathbb Z^2\).  A point of
\(U\mathbin\triangle K\) lies within distance \(\sqrt2\) of
\(\partial K\): this is immediate for \(U\setminus K\), while a point of
\(K\setminus U\) whose distance from the boundary exceeded \(\sqrt2\) would
have its nearest integer point inside \(K\) and would belong to the
corresponding square.  If \(P(K)\) is the perimeter, the outer parallel-set
formula for a convex polygon gives
\(\operatorname{area}(K+B_\rho)-\operatorname{area}(K)
=\rho P(K)+\pi\rho^2\): its added pieces are the edge rectangles and vertex
sectors, whose angles total \(2\pi\).  Polygonal approximation gives the same
bound for every compact convex \(K\).  Applying it to the outer and inner
boundary strips with \(\rho=\sqrt2\) gives
\(\operatorname{area}(U\mathbin\triangle K)=O(P(K)+1)\).  Since the unit
squares have disjoint interiors and total area
\(\#(K\cap\mathbb Z^2)\), the lattice estimate follows.

Only finitely many \(n\)'s occur for fixed \(E\).  Applying the estimate to
\cref{cusp-spec-counting-ellipse} and summing proves
\cref{cusp-spec-fixed-energy-count}.  Since \(D_T\sim T^2\), division by
\(D_T\) proves \cref{cusp-spec-counting-density-limit}; if
\(0<E<4\pi^2\), only \(n=0\) occurs and
\cref{cusp-spec-first-band-count} follows.

For \(\widehat g_T\), eigenvalues are multiplied by \(D_T^{1/2}\), so
\(\widehat N_T(E)=N_T(E/D_T^{1/2})\).  For large \(T\), only \(n=0\)
can occur in \cref{cusp-spec-counting-eigenvalue-form}.  Its ellipse has
radius parameter
\(\rho=(E/(4\pi^2D_T^{1/2}))^{1/2}\), hence area
\(ED_T^{1/2}/(4\pi)\) and perimeter
\(O_E(TD_T^{-1/4}+1)=O_E(D_T^{1/4}+1)\).  The same lattice estimate proves
\cref{cusp-spec-unit-volume-count}.

For the collapsing scale, any \((n_1,n_2)\ne(0,0)\) contributes at least
\(4\pi^2\), so it contributes nothing for all sufficiently large \(T\).
For \(n_1=n_2=0\),
\cref{cusp-spec-scaled-dual-plane} says uniformly that
\[
 \frac{T^2\lambda_T(r,s,0,0)}{4\pi^2}
 =\lVert K_T(r,s)^{\mathsf t}\rVert^2
 \longrightarrow r^2+s^2.
\]
The lower bound in \cref{cusp-spec-uniform-singular-bound} also shows that
only a fixed finite set of \((r,s)\)'s can enter a fixed rescaled window.
Termwise comparison on that finite set proves
\cref{cusp-spec-collapsing-count} away from thresholds and the two-sided
bound \cref{cusp-spec-threshold-count} at thresholds.
\end{proof}

\begin{theorem}[Pointed limits and expanding local covers]
\label{cusp-spec-pointed-limits}
Let \(o_T\) be the origin of \(F_T\), and set
\begin{equation}
 \mathbb T^2_{\mathrm{re}}
 :=(\mathbb Re_1\oplus\mathbb Re_3)/
   (\mathbb Ze_1\oplus\mathbb Ze_3).
 \label{cusp-spec-real-two-torus}
\end{equation}
With the product Euclidean metric,
\begin{equation}
 (F_T,o_T)\longrightarrow
 (\mathbb T^2_{\mathrm{re}}\times
  (\mathbb Re_2\oplus\mathbb Re_4),(0,0))
 \cong(\mathbb T^2\times\mathbb R^2,(0,0))
 \label{cusp-spec-pointed-torus-plane-limit}
\end{equation}
in pointed smooth, and hence pointed Gromov--Hausdorff, topology.  More
precisely, for every \(R<\infty\), the radius-\(R\) balls in the two sides
of \cref{cusp-spec-pointed-torus-plane-limit} are isometric for all
sufficiently large \(T\).

For \(M\in\mathbb Z_{\ge1}\), define the marked local sublattice
\begin{equation}
 \Lambda_{T,M}:=
 \mathbb Zv_{\gamma,T}\oplus\mathbb Zv_{u,T}
 \oplus\mathbb Z(Me_1)\oplus\mathbb Z(Me_3),
 \label{cusp-spec-cover-sublattice}
\end{equation}
where
\begin{equation}
 \begin{aligned}
 v_{\gamma,T}&=(6\alpha_T,6m_T,c_T,-q_T),\\
 v_{u,T}&=(\eta_T,L_T,\alpha_T,m_T).
 \end{aligned}
 \label{cusp-spec-long-period-vectors}
\end{equation}
Then
\begin{equation}
 \widetilde F_{T,M}:=\mathbb R^4/\Lambda_{T,M},
 \qquad
 c_{T,M}:\widetilde F_{T,M}\longrightarrow F_T
 \label{cusp-spec-degree-cover}
\end{equation}
 is a regular holomorphic local-isometric cover of degree \(M^2\), with deck
 group \((\mathbb Z/M\mathbb Z)^2\).  For every integer-valued
function \(M(T)\to\infty\),
\begin{equation}
 (\widetilde F_{T,M(T)},\widetilde o_{T,M(T)})
 \longrightarrow(\mathbb R^4,0)
 \label{cusp-spec-cover-Euclidean-limit}
\end{equation}
in pointed smooth and pointed Gromov--Hausdorff topology.  Again, every fixed
radius ball is eventually exactly isometric to its Euclidean counterpart.
If the cover is instead equipped with the pullback of the unit-volume metric
\(\widehat g_T\), then it converges in pointed smooth topology, and hence in
pointed Gromov--Hausdorff topology, to \(\mathbb R^4\) if and only if
\begin{equation}
 \frac{M(T)}{D_T^{1/4}}\longrightarrow\infty,
 \quad\text{equivalently}\quad
 \frac{M(T)}{\sqrt T}\longrightarrow\infty.
 \label{cusp-spec-unit-volume-cover-condition}
\end{equation}
By contrast, if every covering torus is itself rescaled to volume one, no
choice of \(M(T)\) gives a pointed smooth limit \(\mathbb R^4\).
\end{theorem}

\begin{proof}
The period lattice is
\begin{equation}
 P_T\mathbb Z^4
 =\mathbb Zv_{\gamma,T}\oplus\mathbb Zv_{u,T}
  \oplus\mathbb Ze_1\oplus\mathbb Ze_3.
 \label{cusp-spec-primal-lattice-decomposition}
\end{equation}
First quotient \(\mathbb R^4\) by
\(\mathbb Ze_1\oplus\mathbb Ze_3\).  This gives the proposed limit
\[
 X_\infty:=\mathbb T^2_{\mathrm{re}}
 \times(\mathbb Re_2\oplus\mathbb Re_4).
\]
The remaining quotient translations are generated by the images of
\(v_{\gamma,T},v_{u,T}\).  The projection of their integer combination
\(av_{\gamma,T}+bv_{u,T}\) to the noncompact
\((e_2,e_4)\)-plane is
\begin{equation}
 Y_T\begin{pmatrix}a\\b\end{pmatrix},
 \qquad
 Y_T:=\begin{pmatrix}6m_T&L_T\\-q_T&m_T\end{pmatrix}.
 \label{cusp-spec-long-plane-matrix}
\end{equation}
More precisely,
\begin{equation}
 Y_T
 =T\begin{pmatrix}0&1\\-1&0\end{pmatrix}
 +\begin{pmatrix}6m_T&H_T\\B_T-H_T&m_T\end{pmatrix}.
 \label{cusp-spec-long-plane-limit}
\end{equation}
Holomorphy makes the second matrix uniformly bounded for \(T\ge T_0\).
Choose \(C\) above its operator norm.  The first matrix is \(T\) times an
orthogonal matrix, so the reverse triangle inequality gives
\begin{equation}
 \left\lVert Y_T\begin{pmatrix}a\\b\end{pmatrix}\right\rVert
 \ge(T-C)\sqrt{a^2+b^2}.
 \label{cusp-spec-long-plane-lower-bound}
\end{equation}
The displacement in \(X_\infty\) of a translation is at least the norm of
its noncompact-plane projection.  Hence every nonidentity element of the
remaining deck group has displacement at least \(T-C\).
Together with the fixed lattice vectors \(e_1,e_3\), this also proves
\begin{equation}
 \operatorname{sys}(P_T\mathbb Z^4)=1,
 \qquad \operatorname{inj}(F_T,g_T)=\frac12
 \label{cusp-spec-original-injectivity-radius}
\end{equation}
for all sufficiently large \(T\).

Fix \(R\).  Once \(T-C>4R\), the quotient map
\(X_\infty\to F_T\) is injective on the radius-\(R\) ball about the origin.
It is a local isometry, and every minimizing geodesic of length at most
\(R\) in the quotient lifts from the origin, so the closed radius-\(R\)
ball maps onto the closed radius-\(R\) ball in \(F_T\).  Moreover, a
nontrivial translate cannot
shorten the distance between two points of the ball: their original distance
is at most \(2R\), while every translated distance is at least
\((T-C)-2R>2R\).  The balls are therefore isometric.  This
proves \cref{cusp-spec-pointed-torus-plane-limit}, including its pointed
smooth assertion.

Because the four vectors in
\cref{cusp-spec-primal-lattice-decomposition} are a lattice basis, the
sublattice \cref{cusp-spec-cover-sublattice} has index \(M^2\).  Thus the
identity of \(\mathbb R^4\) induces the regular holomorphic degree-\(M^2\)
local-isometric cover \cref{cusp-spec-degree-cover}.  The quotient of the two
lattices is \((\mathbb Z/M\mathbb Z)^2\), which is its deck group.  Consider
a nonzero vector of
\(\Lambda_{T,M}\):
\begin{equation}
 av_{\gamma,T}+bv_{u,T}+M(pe_1+qe_3),
 \qquad a,b,p,q\in\mathbb Z.
 \label{cusp-spec-general-cover-vector}
\end{equation}
If \((a,b)\ne(0,0)\), its \((e_2,e_4)\)-projection and
\cref{cusp-spec-long-plane-lower-bound} show that its norm is at least
\(T-C\).  If \((a,b)=(0,0)\), its norm is at least \(M\).  Therefore
\begin{equation}
 \operatorname{sys}(\Lambda_{T,M})\ge\min\{T-C,M\}.
 \label{cusp-spec-cover-systole}
\end{equation}
Conversely, \(Me_1\) has length \(M\), while either long generator has
length at most \(T+C'\) for a fixed \(C'\).  Hence
\begin{equation}
 \min\{T-C,M\}
 \le\operatorname{sys}(\Lambda_{T,M})
 \le\min\{T+C',M\},
 \qquad
 \operatorname{inj}(\widetilde F_{T,M})
 =\tfrac12\operatorname{sys}(\Lambda_{T,M}).
 \label{cusp-spec-cover-systole-two-sided}
\end{equation}
For \(M=M(T)\to\infty\), this lower bound tends to infinity.  Once it is
greater than \(4R\), the universal covering map
\(\mathbb R^4\to\widetilde F_{T,M(T)}\) identifies the radius-\(R\) balls
isometrically.  This proves \cref{cusp-spec-cover-Euclidean-limit}.

For the pullback of \(\widehat g_T\), all lengths in
\cref{cusp-spec-cover-systole} are multiplied by \(D_T^{-1/4}\).  The long
lower bound becomes
\((T-C)D_T^{-1/4}\sim\sqrt T\), which diverges automatically.  The explicit
lattice vector \(Me_1\) has normalized length \(MD_T^{-1/4}\).  Hence the
normalized systole tends to infinity exactly under
\cref{cusp-spec-unit-volume-cover-condition}.  Divergence gives Euclidean
balls by the preceding covering argument.  We now prove necessity from the
particular flat lattices rather than from any general injectivity-radius
principle.  Put
\[
 \ell_T:=M(T)D_T^{-1/4},
 \qquad A_T:=(T-C)D_T^{-1/4}.
\]
Then \(A_T\sim\sqrt T\to\infty\).  Quotienting the normalized Euclidean
cover first by
\(\ell_T\mathbb Ze_1\oplus\ell_T\mathbb Ze_3\) gives
\[
 X_{\ell_T}:=
 (\mathbb R/\ell_T\mathbb Z)^2\times\mathbb R^2.
\]
Every remaining nonidentity deck transformation has displacement at least
\(A_T\).  Consequently, for each fixed \(R\), once \(A_T>4R\), the
closed radius-\(R\) ball in the normalized covering torus is exactly the
closed radius-\(R\) ball in \(X_{\ell_T}\).  If
\(\ell_T\not\to\infty\), pass to a subsequence with
\(\ell_T\to\ell<\infty\).  When \(\ell>0\), these pointed product spaces
converge to
\((\mathbb R/\ell\mathbb Z)^2\times\mathbb R^2\); when \(\ell=0\),
the two circle factors have diameter tending to zero and the pointed
Gromov--Hausdorff limit is \(\mathbb R^2\).  Neither limit is
\(\mathbb R^4\).  Thus \(\ell_T\to\infty\), which is precisely
\cref{cusp-spec-unit-volume-cover-condition}, is necessary for both the
pointed smooth and the pointed Gromov--Hausdorff conclusion.

Finally, the volume of the covering torus with its unscaled pulled-back metric
is \(M^2D_T\).  If one instead puts
\begin{equation}
 g^{\mathrm{cov},1}_{T,M}:=(M^2D_T)^{-1/2}g_T
 \quad\text{on }\widetilde F_{T,M},
 \label{cusp-spec-cover-own-unit-volume-metric}
\end{equation}
then its volume is one and its length scale is multiplied by
\((M^2D_T)^{-1/4}\).  The explicit periods \(Me_1\) and
\(v_{\gamma,T}\) consequently have respective lengths
\begin{equation}
 \sqrt{\frac{M}{T}}\,(1+o(1)),
 \qquad
 \sqrt{\frac{T}{M}}\,(1+o(1)).
 \label{cusp-spec-cover-own-unit-volume-lengths}
\end{equation}
Their leading factors have product one, so at least one stays bounded.  The
injectivity radius therefore cannot tend to infinity, which excludes a
pointed smooth \(\mathbb R^4\) limit for the individually unit-volume covers.
\end{proof}

\begin{proposition}[Exact compact quotient of the native cusp fibre]
\label{cusp-spec-native-quotient}
Let
\begin{equation}
 \rho:\mathbb R^4\longrightarrow\mathbb R^2,
 \qquad
 \rho(x_1,x_2,x_3,x_4)=(x_2,x_4),
 \label{cusp-spec-native-projection}
\end{equation}
and let \(Y_T\) be the matrix in
\cref{cusp-spec-long-plane-matrix}.  Then
\begin{equation}
 \det Y_T=D_T>0.
 \label{cusp-spec-long-plane-determinant}
\end{equation}
Consequently
\begin{equation}
 Q_T:=\mathbb R^2/Y_T\mathbb Z^2
 \label{cusp-spec-expanding-quotient}
\end{equation}
is a compact real two-torus, and \(\rho\) induces an exact sequence of
compact real Lie groups
\begin{equation}
 0\longrightarrow\mathbb T^2_{\mathrm{re}}
 \longrightarrow F_T
 \xrightarrow{\ \pi_T\ }Q_T
 \longrightarrow0.
 \label{cusp-spec-native-exact-sequence}
\end{equation}
Moreover,
\begin{equation}
 \operatorname{sys}(Q_T)\ge T-C,
 \label{cusp-spec-expanding-quotient-systole}
\end{equation}
where \(C\) is the constant in
\cref{cusp-spec-long-plane-lower-bound}.  Thus the native
\(\mathbb T^2\times\mathbb R^2\) limit separates exactly into a fixed
unit two-torus kernel and an expanding two-torus quotient.
\end{proposition}

\begin{proof}
Direct calculation gives
\[
 \det Y_T=(6m_T)(m_T)-L_T(-q_T)
 =L_Tq_T+6m_T^2=D_T.
\]
The period decomposition
\cref{cusp-spec-primal-lattice-decomposition} gives
\[
 \rho(P_T\mathbb Z^4)=Y_T\mathbb Z^2,
\]
because \(\rho(e_1)=\rho(e_3)=0\) and the projected long generators
are the two columns of \(Y_T\).  Hence
\[
 \pi_T([x]_{P_T\mathbb Z^4})
 :=[\rho(x)]_{Y_T\mathbb Z^2}
\]
is well defined and surjective.

It remains to determine its kernel without suppressing the integral
lattices.  If \(\pi_T([x])=0\), there are \(a,b\in\mathbb Z\) such that
\[
 \rho(x)=Y_T(a,b)^{\mathsf t}
 =\rho(av_{\gamma,T}+bv_{u,T}).
\]
Thus
\[
 x-av_{\gamma,T}-bv_{u,T}
 \in\mathbb Re_1\oplus\mathbb Re_3,
\]
so \([x]\) lies in the displayed copy of
\(\mathbb T^2_{\mathrm{re}}\).  Conversely that copy is plainly killed by
\(\pi_T\).  Finally, if a vector of \(P_T\mathbb Z^4\) lies in
\(\mathbb Re_1\oplus\mathbb Re_3\), its long coefficients \(a,b\) satisfy
\(Y_T(a,b)^{\mathsf t}=0\).  Since \(\det Y_T=D_T>0\), one has
\(a=b=0\); the intersection is therefore exactly
\(\mathbb Ze_1\oplus\mathbb Ze_3\).  This proves injectivity at the left
and exactness.  The systolic bound is
\cref{cusp-spec-long-plane-lower-bound}.
\end{proof}

\begin{theorem}[The maximal regular-locus
\texorpdfstring{\(\mathbb R^4\)}{R4} quotient and its compactification
obstruction]
\label{cusp-spec-globalization-obstruction}
\mbox{\Verified}\ 
Let
\[
 \Lambda=\mathbb Z\langle\widehat\gamma,\widehat u,
                   \widehat w,\widehat\delta\rangle,\qquad
 K_\Lambda=\mathbb Z\langle\widehat w,\widehat\delta\rangle,
\]
and define
\[
 \varepsilon_\gamma
 (a\widehat\gamma+b\widehat u+c\widehat w+d\widehat\delta)=a,\qquad
 L_\Lambda=\ker\varepsilon_\gamma
 =\mathbb Z\langle\widehat u,\widehat w,\widehat\delta\rangle.
\tag{G.1}
\]
Then the following statements hold.

\smallskip
\noindent
\emph{(i) Regular group quotient.}
Over \(B^\circ\), the period family \(J=X^\circ=f^{-1}(B^\circ)\) is a
holomorphic family of complex Lie groups, not merely an affine torus
bundle.  The formula
\begin{equation}
 \chi_\gamma\bigl([z,\Pi(z)x]\bigr)
 =\bigl([z],\varepsilon_\gamma(x)\bmod\mathbb Z\bigr)
 \label{cusp-spec-global-circle-map}
\end{equation}
defines a global smooth real-analytic fibrewise group epimorphism
\begin{equation}
 \chi_\gamma:J\longrightarrow B^\circ\times(\mathbb R/\mathbb Z).
 \label{cusp-spec-global-circle-quotient}
\end{equation}
Write
\begin{equation}
 \theta_\gamma:=\operatorname{pr}_{\mathbb R/\mathbb Z}\circ\chi_\gamma:
 J\longrightarrow\mathbb R/\mathbb Z
 \label{cusp-spec-global-circle-coordinate}
\end{equation}
for its circle coordinate.  Thus
\(\chi_\gamma=(f,\theta_\gamma)\); the two symbols have different
codomains and will not be interchanged below.
Its kernel is the connected real three-torus family
\begin{equation}
 \mathcal H:=
 \frac{\mathfrak h_z^\circ\times
       ((L_\Lambda\otimes\mathbb R)/L_\Lambda)}{\Delta},
 \label{cusp-spec-global-kernel-family}
\end{equation}
where \(\Delta=\pi_1(B^\circ)\) acts diagonally through the deck action
on \(\mathfrak h_z^\circ\) and the restricted monodromy on
\(L_\Lambda\).  Thus every fibre of \(\mathcal H\) is
\((L_\Lambda\otimes\mathbb R)/L_\Lambda\); the family need not be a
product.
The quotient is real and is not a holomorphic quotient of complex tori.

\smallskip
\noindent
\emph{(ii) Exact radial quotient metric.}
On the positive radial cusp, the restriction of
\cref{cusp-spec-global-circle-map} to
\(F_T=\mathbb R^4/P_T\mathbb Z^4\) is
\begin{equation}
 \chi_{\gamma,T}([y])
 =\left\langle P_T^{-\mathsf T}e_\gamma,y\right\rangle
   \bmod\mathbb Z,\qquad e_\gamma=(1,0,0,0)^{\mathsf T}.
 \label{cusp-spec-radial-circle-character}
\end{equation}
For the Riemannian quotient of the Euclidean metric, the target circle has
circumference
\begin{equation}
 \ell_T
 =\frac1{\lVert P_T^{-\mathsf T}e_\gamma\rVert}
 =\frac{D_T}{\sqrt{L_T^2+m_T^2}}
 =T+H-B+O(T^{-1}).
 \label{cusp-spec-global-circle-circumference}
\end{equation}
In particular, \(\ell_T\to\infty\).

\smallskip
\noindent
\emph{(iii) Four copies and minimality.}
The source
\begin{equation}
 J^{[4]}=J\times_{B^\circ}J\times_{B^\circ}J\times_{B^\circ}J
 \label{cusp-spec-fourfold-fibre-product}
\end{equation}
is the actual fourfold fibre product over all of \(B^\circ\).  The product
map
\begin{equation}
 \chi_\gamma^{[4]}=(\chi_\gamma,\chi_\gamma,\chi_\gamma,\chi_\gamma):
 J^{[4]}\longrightarrow B^\circ\times(\mathbb R/\mathbb Z)^4
 \label{cusp-spec-fourfold-circle-map}
\end{equation}
has kernel
\(\mathcal H^{[4]}=\mathcal H\times_{B^\circ}\cdots
\times_{B^\circ}\mathcal H\), whose fibres are
\(\mathbb T^{12}\).  Along the synchronized positive
radial cusp, its quotient fibre, with product quotient metric, is
\begin{equation}
 G_T=(\mathbb R/\ell_T\mathbb Z)^4,\qquad
 \operatorname{sys}(G_T)=\ell_T.
 \label{cusp-spec-four-circle-fibre}
\end{equation}
Consequently
\begin{equation}
 (G_T,0)\longrightarrow(\mathbb R^4,0)
 \label{cusp-spec-four-circle-limit}
\end{equation}
in pointed local-isometric, pointed smooth, and pointed
Gromov--Hausdorff topology.  If \(\ell_T>4R\), the radius-\(R\) ball
about zero is exactly Euclidean.

Equivalently, the external Cartesian product gives a quotient
\[
 J^4\longrightarrow(B^\circ)^4\times(\mathbb R/\mathbb Z)^4,
\]
and \(J^{[4]}\) is its pullback along the diagonal
\(B^\circ\to(B^\circ)^4\).  At four independent cusp parameters its target
lattice is
\(\operatorname{diag}(\ell_{T_1},\ldots,\ell_{T_4})\mathbb Z^4\), so the
same pointed limit holds whenever
\(\min_i\ell_{T_i}\to\infty\).

Among connected torus quotients obtained from \(r\) copies by a primitive
full-monodromy-invariant kernel and whose quotient systole tends to infinity
on the radial cusp, the quotient rank is at most \(r\).  Thus four copies
are minimal for a four-dimensional all-expanding quotient, and
\cref{cusp-spec-fourfold-circle-map} attains the bound.

The cusp-local two-copy map
\begin{equation}
 \pi_T\times\pi_T:F_T\times F_T\longrightarrow Q_T\times Q_T
 \label{cusp-spec-doubled-local-quotient}
\end{equation}
from \cref{cusp-spec-native-quotient} also has pointed limit
\(\mathbb R^4\).  It is not a quotient over the full \(B^\circ\), however,
and no finite base cover makes it one.

\smallskip
\noindent
\emph{(iv) Finite fillings.}
At \(p_j\), use the notation of
\cref{prop:candidate-finite-fillings} and put
\[
 (m_1,v_1,\kappa_1)
 =\bigl(3,\widehat\gamma+2\widehat u-4\widehat w,1\bigr),
\]
\[
 (m_2,v_2,\kappa_2)
 =\bigl(4,-\widehat\gamma-3\widehat u+3\widehat w,-1\bigr),\qquad
 \kappa_j=\varepsilon_\gamma(v_j).
\tag{G.2}
\]
In the un-conjugated flat coordinates, the finite filling generator is
\begin{equation}
 \Phi_j(s,x)=
 \left(e^{-2\pi i/m_j}s,A_jx+\frac{v_j}{m_j}\right).
 \label{cusp-spec-finite-affine-action}
\end{equation}
The raw character is not invariant:
\begin{equation}
 \varepsilon_\gamma\left(A_jx+\frac{v_j}{m_j}\right)
 =\varepsilon_\gamma(x)+\frac{\kappa_j}{m_j}\pmod{\mathbb Z}.
 \label{cusp-spec-finite-character-shift}
\end{equation}
It instead descends equivariantly to the Seifert circle family
\begin{equation}
 C_j=
 \frac{\Delta_j\times(\mathbb R/\mathbb Z)}
 {(s,\theta)\sim
  (e^{-2\pi i/m_j}s,\theta+\kappa_j/m_j)}.
 \label{cusp-spec-finite-Seifert-target}
\end{equation}
The total space \(C_j\) is a smooth solid torus, the map
\(C_j\to D_j\) has a multiple central circle of order \(m_j\), and the
induced map \(N_j\to C_j\) is a smooth proper \(\mathbb T^3\)-bundle.
The target translation in \cref{cusp-spec-finite-Seifert-target} is
exactly compatible with the logarithmic gluing map \(h_j\).

For four synchronized copies, the smooth simultaneous-root extension is
\begin{equation}
 \widetilde N_j^{[4]}
 =\frac{(J'_j)^{\times_{\Delta_j}4}}
 {\langle(\widetilde g_j^{\log})^{\times4}\rangle},
 \label{cusp-spec-finite-fourfold-normalization}
\end{equation}
with target
\begin{equation}
 C_j^{[4]}=
 \frac{\Delta_j\times(\mathbb R/\mathbb Z)^4}
 {(s,\theta_1,\ldots,\theta_4)\sim
  (e^{-2\pi i/m_j}s,
   \theta_1+\kappa_j/m_j,\ldots,\theta_4+\kappa_j/m_j)}.
 \label{cusp-spec-finite-fourfold-target}
\end{equation}
The map \(\widetilde N_j^{[4]}\to C_j^{[4]}\) is a smooth proper
\(\mathbb T^{12}\)-bundle.  There is a canonical finite map
\begin{equation}
 \nu_j:\widetilde N_j^{[4]}
 \longrightarrow N_j\times_{D_j}N_j\times_{D_j}N_j\times_{D_j}N_j
 \label{cusp-spec-finite-normalization-map}
\end{equation}
which is an isomorphism over \(D_j^*\), is the normalization of the
reduced ordinary filled fibre product, and has degree \(m_j^3\) on the
reduced central fibre.  These two degrees are \(27\) and \(64\).  Gluing
these two local normalizations to \(J^{[4]}\) by
\((h_j)^{\times4}\), and gluing their targets to the product
\(B^\circ\times(\mathbb R/\mathbb Z)^4\) by
\((h_j^{S^1})^{\times4}\), gives the corresponding normalization and
Seifert target over \(B\setminus\{p_0\}\).

\smallskip
\noindent
\emph{(v) Cusp nonextension.}
Neither the literal map \(\chi_\gamma\) nor
\(\chi_\gamma^{[4]}\) extends continuously through the toric cusp fibre
\(W\).  Thus the four-copy quotient is global over \(B^\circ\), and its
finite-filling normalization is global over \(B\setminus\{p_0\}\), but
the displayed map \(\chi_\gamma^{[4]}\) has no continuous extension to
the compact fourfold fibre product or to \(X^4\) with the stated
regular-fibre restriction, and it does not factor through
\(X^{\wedge4}\cong S^{24}\).  The cohomology class of
\(\chi_\gamma\) does extend over
\(N_0\); this weaker homotopy statement does not extend the displayed
quotient map or its metric.

\smallskip
\noindent
\emph{(vi) Products, smashes, and sphere constructions.}
If the recognition \(X\cong_{\mathrm{diff}}S^6\) is used, then
\begin{equation}
 X^2\cong_{\mathrm{diff}}S^6\times S^6\not\simeq S^{12},\qquad
 H_6(X^2;\mathbb Z)\cong\mathbb Z^2,
 \label{cusp-spec-product-not-sphere}
\end{equation}
and
\begin{equation}
 \begin{aligned}
 X^4&\cong_{\mathrm{diff}}(S^6)^4\not\simeq S^{24},\\
 H_6(X^4;\mathbb Z)&\cong\mathbb Z^4,&
 H_{12}(X^4;\mathbb Z)&\cong\mathbb Z^6,&
 H_{18}(X^4;\mathbb Z)&\cong\mathbb Z^4.
 \end{aligned}
 \label{cusp-spec-four-product-not-sphere}
\end{equation}
The synchronized regular fourfold \(J^{[4]}\) has real dimension \(18\),
not \(24\).  Choosing a base point gives based homeomorphisms
\begin{equation}
X\wedge X\cong S^{12},\qquad X^{\wedge4}\cong S^{24},
\label{cusp-spec-smash-spheres}
\end{equation}
but these smash spheres do not carry either \emph{literal fixed-target}
quotient above.  The later cone constructions
\cref{cusp-cone-S12-route,cusp-cone-S24-local-route,cusp-cone-S24-route}
change the codomain by retaining the radial parameter and collapsing the
obstructing boundary cross-section; they therefore do not contradict this
nonextension theorem.  Moreover,
neither \(S^{12}\) nor \(S^{24}\), with any differentiable structure,
admits an almost-complex structure
\cite[\S15, Proposition~15.1 and the following remark,
p.~434]{BorelSerre1953}; hence neither can be promoted to a compact
complex manifold.  This is an obstruction to a proposed complex
\(S^{12}\)- or \(S^{24}\)-globalization, not to the real quotient or its
pointed metric limit.  Precisely,
\(\pi_T\times\pi_T\) is a map on \(F_T^2\), and it fails to factor
through \(F_T\wedge F_T\); \(\chi_{\gamma,T}^{[4]}\) is a map on
\(F_T^4\), and it fails to factor through \(F_T^{\wedge4}\).
Neither map is a map on \(X^2\) or \(X^4\), and no extension with the
stated fibre restriction exists.

There is one exact but different sphere construction.  If
\(C_T=\mathbb R/\ell_T\mathbb Z\) is based at zero, postcomposition with
the target fat-wedge collapse gives
\begin{equation}
 F_T^4\xrightarrow{\chi_{\gamma,T}^{[4]}}C_T^4
 \longrightarrow C_T^{\wedge4}\cong S^4.
 \label{cusp-spec-target-smash-map}
\end{equation}
This map does factor through \(F_T^{\wedge4}\).  With the quotient length
metric and base point
\([\ell_T/2,\ell_T/2,\ell_T/2,\ell_T/2]\), its pointed limit is
\(\mathbb R^4\), with every radius-\(R\) ball exactly Euclidean once
\(\ell_T>4R\).  The metric is singular at the collapsed fat wedge, and
the map still does not extend to \(X^4\) or \(X^{\wedge4}\).
\end{theorem}

\begin{proof}
\medskip\noindent\emph{Invariant lattices and the regular group law.}\par
The matrices in \cref{prop:monodromy-certificate} give
\[
 (M_0-I)(a,b,c,d)^{\mathsf T}=(0,0,b,-a)^{\mathsf T},
\]
so
\begin{equation}
 K_\Lambda=\ker(M_0-I)=\operatorname{im}(M_0-I).
 \label{cusp-spec-cusp-kernel}
\end{equation}
They also give
\begin{equation}
 \begin{aligned}
 A_1\widehat u&=-\widehat w+\widehat\delta,&
 A_1\widehat w&=\widehat u-\widehat w,&
 A_1\widehat\delta&=\widehat\delta,\\
 A_2\widehat u&=\widehat w,&
 A_2\widehat w&=-\widehat u+\widehat\delta,&
 A_2\widehat\delta&=\widehat\delta.
 \end{aligned}
 \label{cusp-spec-invariant-flag-calculation}
\end{equation}
Thus \(L_\Lambda\) is invariant.  It is the smallest invariant lattice
containing \(K_\Lambda\), because
\(\widehat u=A_1\widehat w+\widehat w\).  Both lattices are primitive.

For a row covector \((a,b,c,d)\), invariance under \(A_1\) gives
\(b=2c\) and \(d=3c\), whereas invariance under \(A_2\) gives
\(b=c\) and \(d=2c\).  Hence \(b=c=d=0\), and
\begin{equation}
 \operatorname{Hom}(\Lambda,\mathbb Z)^\Delta
 =\mathbb Z\varepsilon_\gamma.
 \label{cusp-spec-invariant-character-lattice}
\end{equation}

The family over the cover is
\((\mathfrak h_z^\circ\times\mathbb C^2)/\Lambda\), and the action used
to form \(J\) is
\[
 [z,\zeta]\longmapsto[gz,R_g(z)\zeta].
\]
It is complex linear in every fibre and fixes the zero section.
Fibrewise addition, inversion, and zero are therefore equivariant and
descend to \(J\).  In the flat real-analytic marking
\[
 \Xi(z,x+\Lambda)=[z,\Pi(z)x],
\]
equation (2.3d) makes the action
\[
 (z,x+\Lambda)\longmapsto(gz,A_gx+\Lambda).
\]
Equation \cref{cusp-spec-invariant-character-lattice} proves that
\cref{cusp-spec-global-circle-map} is well defined.  Primitivity of
\(\varepsilon_\gamma\) proves surjectivity and connectedness of the
kernel.  The induced action on \(\Lambda/L_\Lambda\simeq\mathbb Z\) is
trivial, so the target is the product circle bundle.  Since the kernel has
odd real dimension, this is not a holomorphic quotient.

\medskip\noindent\emph{The quotient metric and the four-copy limit.}\par
The first column of
\cref{cusp-spec-full-inverse-transpose} is
\begin{equation}
 P_T^{-\mathsf T}e_\gamma
 =\left(0,\frac{m_T}{D_T},0,-\frac{L_T}{D_T}\right)^{\mathsf T}.
 \label{cusp-spec-circle-dual-vector}
\end{equation}
The quotient metric on \(\mathbb R/\mathbb Z\) is therefore
\(\ell_T^2d\theta^2\).  Equivalently,
\[
 P_T(L_\Lambda\otimes\mathbb R)
 =\operatorname{span}_{\mathbb R}\{e_1,e_3,v_{u,T}\}
\]
has unit normal
\[
 \nu_T=\frac{(0,m_T,0,-L_T)}{\sqrt{L_T^2+m_T^2}},
\]
and
\[
 \operatorname{dist}
 \bigl(v_{\gamma,T},P_T(L_\Lambda\otimes\mathbb R)\bigr)
 =\left|\langle v_{\gamma,T},\nu_T\rangle\right|
 =\frac{L_Tq_T+6m_T^2}{\sqrt{L_T^2+m_T^2}}
 =\ell_T.
\]
The expansions in
\cref{cusp-spec-radial-notation,cusp-spec-holomorphic-errors} give
\[
 D_T=T^2+(2H-B)T+O(1),\qquad
 \sqrt{L_T^2+m_T^2}=T+H+O(T^{-1}),
\]
which proves \cref{cusp-spec-global-circle-circumference}.

After rescaling the circle coordinate by \(\ell_T\), the product target is
\(\mathbb R^4/\ell_T\mathbb Z^4\).  If
\(x,y\in B_R(0)\) and \(0\ne\lambda\in\ell_T\mathbb Z^4\), then
\[
 \lVert x-y+\lambda\rVert
 \ge\lVert\lambda\rVert-\lVert x-y\rVert
 >\ell_T-2R>2R\ge\lVert x-y\rVert
\]
when \(\ell_T>4R\).  Hence the covering projection is an exact isometry
on \(B_R(0)\), proving all three convergences in
\cref{cusp-spec-four-circle-limit}.  The quotient map is global over
\(B^\circ\); only the metric limit is evaluated along the radial cusp.
The separately chosen Euclidean fibre metric is not claimed to be a
canonical global metric away from that cusp.

\medskip\noindent\emph{Minimality and the finite-base-cover obstruction.}\par
Let \(N\subset\Lambda^{\oplus r}\) be a primitive invariant kernel whose
connected quotient has radial systole tending to infinity.  Every element
of \(K_\Lambda^{\oplus r}\) has uniformly bounded physical length, since
\[
 P_T\widehat w=e_1,\qquad P_T\widehat\delta=e_3.
\]
If such an element survived in the quotient lattice, its quotient norm
would be no greater than its source norm and would give a nonzero uniformly
bounded period.  Therefore \(K_\Lambda^{\oplus r}\subset N\).
Invariance and
\cref{cusp-spec-invariant-flag-calculation} then force
\(L_\Lambda^{\oplus r}\subset N\), and
\[
 \operatorname{rank}(\Lambda^{\oplus r}/N)\le4r-3r=r.
\]
This proves the lower bound \(r\ge4\), while
\(N=L_\Lambda^{\oplus4}\) proves sharpness.

For the finite-cover assertion, put
\begin{equation}
 H_*=A_2A_1
 =\begin{pmatrix}
 1&0&0&0\\
 6&1&1&0\\
 0&0&1&0\\
 -5&0&-1&1
 \end{pmatrix}=I+N_*.
 \label{cusp-spec-no-finite-cover-matrix}
\end{equation}
Direct multiplication gives \(N_*^2=0\), and hence
\begin{equation}
 H_*^n\widehat w
 =\widehat w+n(\widehat u-\widehat\delta),\qquad
 H_*^n\widehat\delta=\widehat\delta.
 \label{cusp-spec-no-finite-cover-orbit}
\end{equation}
Thus \(H_*^nK_\Lambda\ne K_\Lambda\) for every \(n\ne0\).
Every finite-index subgroup of the base fundamental group contains a
positive power of the loop with monodromy \(H_*\).  Such a subgroup cannot
preserve \(K_\Lambda\), proving that no finite base cover globalizes the
native two-torus quotient.

\medskip\noindent\emph{The affine finite fillings and the actual overlap
maps.}\par
Equation \cref{cusp-spec-invariant-character-lattice} gives
\(\varepsilon_\gamma A_j=\varepsilon_\gamma\), so
\cref{cusp-spec-finite-character-shift} follows.  Define
\[
 \overline\Phi_j(s,\theta)
 =\left(e^{-2\pi i/m_j}s,\theta+\frac{\kappa_j}{m_j}\right).
\]
Then \((s,x)\mapsto(s,\varepsilon_\gamma(x))\) is
\((\Phi_j,\overline\Phi_j)\)-equivariant; denote the descended map by
\(\overline\chi_j:N_j\to C_j\).  A nonidentity power which fixes
the base must lie over \(s=0\), where it translates the circle by
\(k\kappa_j/m_j\).  Since \(\gcd(\kappa_j,m_j)=1\), this is nonintegral
for \(0<k<m_j\).  The target action is free; the source action is free by
\cref{prop:candidate-finite-fillings}.

For an explicit solid-torus chart, put
\(\lambda_j=e^{2\pi i/m_j}\), \(u=e^{2\pi i\theta}\), and choose
\(r_j\in\mathbb Z\) with \(r_j\kappa_j\equiv1\pmod {m_j}\).  Then
\[
 w=su^{r_j},\qquad U=u^{m_j}
\]
are invariant and identify \(C_j\) real-analytically with a solid torus.
In these coordinates
\[
 t_j=s^{m_j}=w^{m_j}U^{-r_j},
\]
which proves the central multiplicity.  Equivariance and the primitive
kernel give the proper \(\mathbb T^3\)-bundle \(N_j\to C_j\).

It remains to check the actual logarithmic gluing.  Let
\(\Pi_{\mathbb R}(z):\Lambda\otimes\mathbb R\to\mathbb C^2\) be the
underlying real isomorphism.  On a logarithm branch \(\ell\), define
\(b_j^\ell(z)\) uniquely by
\begin{equation}
 \Pi_{\mathbb R}(z)b_j^\ell(z)
 =\frac{\ell(z)}{2\pi i}\Pi(z)v_j,\qquad
 \beta_j(z)=\varepsilon_\gamma(b_j^\ell(z))\bmod\mathbb Z.
 \label{cusp-spec-finite-target-conjugacy}
\end{equation}
Changing \(\ell\) by \(2\pi in\) adds \(nv_j\) to \(b_j^\ell\), hence
adds \(n\kappa_j\in\mathbb Z\) to its
\(\varepsilon_\gamma\)-value.  Thus \(\beta_j\) is branch independent.
In flat coordinates the source conjugacy (2.9e) is translation by
\(b_j^\ell\).  Applying \(\varepsilon_\gamma\) to that conjugacy gives
\begin{equation}
 \beta_j(g_jz)=\beta_j(z)-\frac{\kappa_j}{m_j}\pmod{\mathbb Z}.
 \label{cusp-spec-finite-target-cocycle}
\end{equation}
Therefore
\[
 h_j^{S^1}(z,\theta)=(z,\theta+\beta_j(z))
\]
conjugates \(\overline\Phi_j\) to the linear target action
\((z,\theta)\mapsto(g_jz,\theta)\), and the overlap identity is
\begin{equation}
 \chi_\gamma\circ h_j=h_j^{S^1}\circ\overline\chi_j.
 \label{cusp-spec-finite-overlap-identity}
\end{equation}

The diagonal action on
\((J'_j)^{\times_{\Delta_j}4}\) is free, and the same coprimality argument
proves that \(C_j^{[4]}\) is smooth.  Every fibre of the resulting map is
a \(\mathbb T^{12}\)-torsor; over \(D_j^*\) its structure torus is the
fourfold fibre product of \(\mathcal H\).  No group zero is asserted at
the affine filling.
In local etale coordinates along the central fibre, the ordinary filled
fibre product has equations
\begin{equation}
 s_i^{m_j}=s_1^{m_j},\qquad i=2,3,4.
 \label{cusp-spec-finite-fibre-product-branches}
\end{equation}
It has \(m_j^3\) smooth branches along the reduced central fibre.  The
simultaneous-root model separates them, is smooth and hence normal, and is
an isomorphism away from the central fibre.  The map \(\nu_j\) is proper:
both spaces are proper over \(D_j\), and the inverse image of a compact
subset of the target is closed in the inverse image of its compact image
in \(D_j\).  It is quasi-finite, since its fibres are singletons away from
the central fibre and have at most \(m_j^3\) points over the central
fibre.  Hence \(\nu_j\) is finite.  The universal property of
normalization now proves that it is the normalization of the reduced
ordinary fibre product.  On the central fibre the relative root choices form
\((\mathbb Z/m_j)^4/\operatorname{diag}(\mathbb Z/m_j)\), of order
\(m_j^3\), which proves the asserted degrees.

On \(D_j^*\), the source identification is \((h_j)^{\times4}\), and
\cref{cusp-spec-finite-overlap-identity} identifies the target overlap
with \((h_j^{S^1})^{\times4}\).  These are precisely the fourfolds of
the one-copy overlap maps and therefore satisfy the cocycle identities.
They glue the two local normalizations and Seifert targets to
\(J^{[4]}\) and to
\(B^\circ\times(\mathbb R/\mathbb Z)^4\), respectively.  Equivalently,
the source gluing follows from uniqueness of normalization.

The fixed product-target trivialization does not extend over \(X\), even if
base winding is allowed.  Indeed, in the presentation of
\cref{prop:geometric-pi1}, let a homomorphism to \(\mathbb Z\) restrict to
\(n\varepsilon_\gamma\) on \(\Lambda\) and take values \(a,b\) on
\(\widehat\rho_1,\widehat\rho_2\).  The three attachment relations give
\[
 3a=n,\qquad4b=-n,\qquad a+b=0.
\]
The first two give \(a=n/3\), \(b=-n/4\), while the third gives
\(n/12=0\).  Hence \(n=0\).  This is the exact global obstruction behind
the fact that the two Seifert targets do not combine into a fixed-circle
map on \(X\).

\medskip\noindent\emph{The explicit toric cusp obstruction.}\par
On the positive radial ray, write
\[
 Z_T=iTB_0+C(t_c),\qquad
 B_0=\begin{pmatrix}0&1\\-1&0\end{pmatrix},\qquad
 t_c=e^{-2\pi T}.
\]
For \(a=(a_1,a_2)^{\mathsf T}\in\mathbb R^2\), take the point of the
regular fibre represented by \(Z_Ta\), with zero short-period coordinates,
and exponentiate it to the dense torus of the cusp model:
\[
 x_T(a)=\exp(2\pi iZ_Ta)\in(\mathbb C^*)^2.
\]
Its logarithmic position is
\begin{equation}
 y_T(a)=
 \frac{(\log|x_{T,1}(a)|,\log|x_{T,2}(a)|)}{\log|t_c|}
 =B_0a+\frac{\operatorname{Im}C(t_c)a}{T}.
 \label{cusp-spec-cusp-nonextension-position}
\end{equation}
The circle character has the constant value
\begin{equation}
 \chi_{\gamma,T}(x_T(a))=a_1\bmod\mathbb Z.
 \label{cusp-spec-cusp-nonextension-character}
\end{equation}

Use the affine toric chart of
\(\sigma=\operatorname{conv}(0,e_1,e_2)\).  Its coordinates are
\begin{equation}
 z_0=\frac{t_c}{x_1x_2},\qquad z_1=x_1,\qquad z_2=x_2.
 \label{cusp-spec-cusp-nonextension-chart}
\end{equation}
If \(B_0a=(y_1,y_2)\) is in the interior of \(\sigma\), then
\(y_1>0\), \(y_2>0\), and \(1-y_1-y_2>0\).  For all sufficiently large
\(T\), the three absolute values in
\cref{cusp-spec-cusp-nonextension-chart} are
\[
 |t_c|^{1-y_{T,1}-y_{T,2}},\qquad
 |t_c|^{y_{T,1}},\qquad |t_c|^{y_{T,2}},
\]
and all tend to zero.  Every such path converges to the same toric fixed
point and hence to the same triple point of \(W\).

Take
\[
 a=\left(-\frac14,\frac14\right),\qquad
 a'=\left(-\frac12,\frac14\right).
\]
Then
\[
 B_0a=\left(\frac14,\frac14\right),\qquad
 B_0a'=\left(\frac14,\frac12\right)
\]
are both in the interior of \(\sigma\).  The paths converge to the same
point of \(W\), whereas
\[
 \chi_{\gamma,T}(x_T(a))=\frac34,\qquad
 \chi_{\gamma,T}(x_T(a'))=\frac12.
\]
No continuous extension of the literal map \(\chi_\gamma\) exists.
Choosing identical marked radial paths in the other three factors and
using these two paths in one factor proves the same assertion for
\(\chi_\gamma^{[4]}\).

The radial family of translations by \(\widehat u/2\) in \(\mathcal H\)
fails to extend as well.  Take two paths
whose limiting logarithmic positions are
\[
 y=\left(\frac18,\frac18\right),\qquad
 y'=\left(\frac38,\frac38\right).
\]
Both approach the fixed point of
\(\operatorname{conv}(0,e_1,e_2)\).  Translation by
\(\widehat u/2\in\mathcal H_T\) adds the exact displacement
\[
 \frac12B_0e_2+
 \frac1{2T}\operatorname{Im}C(t_c)e_2
 \longrightarrow\left(\frac12,0\right)
\]
to logarithmic position.  The translated limits are
\[
 \left(\frac58,\frac18\right)
 \quad\hbox{and}\quad
 \left(\frac78,\frac38\right).
\]
The first lies in the interior of
\(\operatorname{conv}(0,e_1,e_2)\), while the second lies in the interior
of \(\operatorname{conv}(e_1,e_2,e_1+e_2)\).  These two triangle types are
the two translation orbits in the \(A_2\) triangulation and give the two
distinct triple points of \(W\), as proved in
\cref{thm:candidate-cusp-geometry}.  Thus one regular-fibre translation
sends two paths with the same limit to paths with different limits; it has
no continuous extension to \(N_0\).

This pointwise obstruction must not be confused with the cohomology class.
By \cref{prop:geometric-pi1},
\[
 \pi_1(N_0)=\Lambda/K_\Lambda,
\]
and \(\varepsilon_\gamma\) kills \(K_\Lambda\).  It therefore defines an
integral \(H^1\)-class on \(N_0\), hence a map to the Eilenberg--Mac Lane
space \(S^1\) whose restriction is homotopic to \(\chi_\gamma\).  The
explicit paths prove that no such representative agrees with
\(\chi_\gamma\) on the whole punctured cusp.  The phase torus
\((K_\Lambda\otimes\mathbb R)/K_\Lambda\) extends as the compact real torus
action of the toric model, generally with boundary stabilizers; the full
radial \(\widehat u/2\)-translation and the full quotient do not.

\medskip\noindent\emph{Products and smash quotients.}\par
The integral Kunneth theorem gives
\[
 H_k(X^n;\mathbb Z)\cong
 \begin{cases}
  \mathbb Z^{\binom nj},&k=6j,\quad0\le j\le n,\\
  0,&\text{otherwise},
 \end{cases}
 \qquad \chi(X^n)=2^n,
\]
which proves
\cref{cusp-spec-product-not-sphere,cusp-spec-four-product-not-sphere}.
The synchronized regular \(n\)-fold fibre product has dimension
\[
 6n-2(n-1)=4n+2,
\]
which gives \(18\) for four copies.

A product map factors through a smash quotient if and only if it is
constant on the corresponding fat wedge.  On
\(F_T\times\{0\}\), the local two-copy quotient is
\[
 x\longmapsto(\pi_T(x),0),
\]
which is nonconstant.  Thus the literal conclusion is that
\(\pi_T\times\pi_T\) does not factor through \(F_T\wedge F_T\); it would
be a type error to call it a map on \(X\wedge X\).  The restriction of
\(\chi_{\gamma,T}^{[4]}\) to \(F_T\times\{0\}^3\) is similarly
\[
 x\longmapsto(\chi_{\gamma,T}(x),0,0,0),
\]
so it does not factor through \(F_T^{\wedge4}\).

There is no global extension which repairs the domain mismatch:
\(H^1(X^n;\mathbb Z)=0\), whereas both local quotient maps pull nonzero
target \(H^1\)-classes back to nonzero classes on the cusp fibre.  Hence no
map on \(X^2\) or \(X^4\) has the required restriction, and precomposition
with the smash collapses in \cref{cusp-spec-smash-spheres} cannot produce
the local quotient maps.

Identify the metric quotient circle with
\(C_T=\mathbb R/\ell_T\mathbb Z\) by
\(\theta\mapsto\ell_T\theta\).  For
\cref{cusp-spec-target-smash-map}, target-smash collapse is constant
whenever one source coordinate is the identity, so the modified map does
factor through \(F_T^{\wedge4}\).  Its chosen base point has distance
\(\ell_T/2\) from the collapsed target fat wedge.  If two points lie
within distance \(R\) of it, a route through the collapsed point has
length at least \(\ell_T-2R\), while their direct Euclidean distance is at
most \(2R\).  When \(\ell_T>4R\), neither the collapse nor circle
wrap-around shortens distances in that ball.

Let \(u_i\in H^1(F_T^4;\mathbb Z)\) be the pullback of the generator from
the \(i\)-th quotient circle.  The orientation class of
\(S^4=C_T^{\wedge4}\) pulls back under
\cref{cusp-spec-target-smash-map} to
\[
 \pm u_1\smile u_2\smile u_3\smile u_4\ne0.
\]
But \(H^4(X^4;\mathbb Z)=0\).  Thus even this modified \(S^4\)-target map
does not extend to \(X^4\), and a fortiori does not arise from a map on
\(X^{\wedge4}\cong S^{24}\).

For completeness, Borel--Serre prove that \(S^{2n}\) admits no
almost-complex structure for \(n\ge4\).  Their argument applies to every
differentiable structure on the sphere: on \(S^{2n}\), all lower Chern
classes vanish, while their reduced-power calculation forces the top
Chern class to vanish modulo a suitable odd prime \(p<n\); on the other
hand that class is the Euler class, namely \(2h\) for a generator
\(h\in H^{2n}(S^{2n};\mathbb Z)\), a contradiction.  Taking \(n=6\) and
\(n=12\) gives the asserted obstructions for \(S^{12}\) and \(S^{24}\).
This topological obstruction does not itself prohibit the real
four-circle target \(G_T\), the target-smash \(S^4\), or their metric
limits.

Finally,
\(X*X\cong S^{13}\) and \(X^{*4}\cong S^{27}\), so joins do not produce
the proposed \(S^{12}\) or \(S^{24}\) dimensions.  Suspensions can produce
spheres of those dimensions only by discarding the torus-fibration maps.
\end{proof}


\subsection{The Kummer cover extends across the toric cusp}
\label{cusp-cover-section}

This section determines exactly what is true in the assertion that the
degree-$N^2$ cover of the regular fibres extends by a toric power map.  The
answer is positive for the completed \emph{local} cusp family, but the source
action has to be changed, and two different degrees must not be conflated.

Retain the cusp data
\[
 B_0=\begin{pmatrix}0&1\\-1&0\end{pmatrix},\qquad
 Z(s)=sB_0+C(t),\qquad t=e^{2\pi i s},
\]
where $C$ is holomorphic on a disc $\Delta_t$ about zero.  Put
\[
 c_{\bar\lambda}(t)
 :=\exp\!\bigl(2\pi iC(t)\bar\lambda\bigr)
 \in(\mathbb C^*)^2,
 \qquad \bar\lambda\in\bar\Lambda\simeq\mathbb Z^2.
\]
Let $\Sigma$ be the fan in
$N'=\mathbb Z^2\oplus\mathbb Z$ consisting of the cones over the cells of
the translation-invariant $A_2$ triangulation at height one, and write
$Y=Y_\Sigma$.  On the dense torus, with coordinates $(x_1,x_2,t)$, the
target action is
\begin{equation}
 \Psi_{\bar\lambda}(x,t)
 =\bigl(c_{\bar\lambda}(t)t^{B_0\bar\lambda}x,t\bigr).
 \label{cusp-cover-target-action}
\end{equation}
Its quotient is the completed cusp family
$f_0\colon N_0=Y/\bar\Lambda\to\Delta_t$.

Fix an integer $N\geq1$.  Let $v$ be a coordinate on a second disc and
make the base change
\begin{equation}
 \rho_N\colon\Delta_v\longrightarrow\Delta_t,\qquad
 \rho_N(v)=v^N,\qquad t=v^N.
 \label{cusp-cover-base-change}
\end{equation}
On a copy $Y^{(N)}$ of the same toric variety, with dense-torus coordinates
$(\xi_1,\xi_2,v)$, define
\begin{equation}
 C_N(v):=\frac1N C(v^N),\qquad
 c^{(N)}_{\bar\lambda}(v)
 :=\exp\!\bigl(2\pi iC_N(v)\bar\lambda\bigr)
 \label{cusp-cover-source-factor}
\end{equation}
and
\begin{equation}
 \Psi^{(N)}_{\bar\lambda}(\xi,v)
 :=\bigl(c^{(N)}_{\bar\lambda}(v)
          v^{B_0\bar\lambda}\xi,v\bigr).
 \label{cusp-cover-source-action}
\end{equation}

\begin{theorem}[Exact Kummer extension through the cusp]
\label{cusp-cover-extension-theorem}
\mbox{\Verified}\ 
After shrinking the $v$-disc, the following statements hold.

\begin{enumerate}
\item The maps in \eqref{cusp-cover-source-action} extend to a free,
properly discontinuous action on $Y^{(N)}$.  The quotient
\[
 f_N\colon N_0^{\langle N\rangle}
 :=Y^{(N)}/\bar\Lambda\longrightarrow\Delta_v
\]
is a proper holomorphic map from a smooth threefold.  Its central fibre
$W_N$ is reduced, irreducible, and nonnormal, and its normalization is a
toric del Pezzo surface of degree six.

\item Multiplication by $N$ on the cocharacter lattice defines a finite
flat toric morphism
\begin{equation}
 P_N\colon Y^{(N)}\longrightarrow Y,\qquad
 P_N(\xi_1,\xi_2,v)=(\xi_1^N,\xi_2^N,v^N),
 \label{cusp-cover-toric-power-map}
\end{equation}
of degree $N^3$.  It is equivariant for
\eqref{cusp-cover-source-action} and
\eqref{cusp-cover-target-action}, hence descends to a finite flat map
\begin{equation}
 \kappa_N\colon N_0^{\langle N\rangle}\longrightarrow N_0,
 \qquad f_0\circ\kappa_N=\rho_N\circ f_N,
 \label{cusp-cover-completed-map}
\end{equation}
of degree $N^3$.

\item The map $\kappa_N$ is etale over $v\ne0$, and
\begin{equation}
 \kappa_N^*W=NW_N,
 \qquad
 \operatorname{Ram}(\kappa_N)=(N-1)W_N.
 \label{cusp-cover-direct-ramification}
\end{equation}
In particular, the branch support of the map to the original cusp family
is the whole central fibre.

\item Let
\begin{equation}
 N_0[N]:=N_0\times_{\Delta_t}\Delta_v,
 \qquad t=v^N.
 \label{cusp-cover-base-changed-target}
\end{equation}
Then $N_0[N]$ is normal and $\kappa_N$ factors uniquely as
\begin{equation}
 N_0^{\langle N\rangle}
 \xrightarrow{\ Q_N\ }N_0[N]
 \longrightarrow N_0.
 \label{cusp-cover-relative-factorization}
\end{equation}
The map $Q_N$ has generic degree $N^2$.  For $N>1$ it is quasi-etale but
not flat: it is etale over the smooth locus of $N_0[N]$, and its non-etale
and nonflat loci are exactly the inverse images of the double and triple
strata of the base-changed central fibre.  Thus $Q_N$ is not a
normalization map.

\item On the reduced central fibres, the induced map
$W_N\to W$ has generic degree $N^2$.  It lifts to a commutative diagram
\begin{equation}
\begin{CD}
 E_N\simeq\mathrm{dP}_6 @>{[N]}>> E\simeq\mathrm{dP}_6\\
 @V{\nu_N}VV @VV{\nu}V\\
 W_N @>>> W,
\end{CD}
 \label{cusp-cover-normalization-diagram}
\end{equation}
where $[N]$ is toric multiplication by $N$, of degree $N^2$, ramified
with index $N$ along each of the six toric boundary curves.  The diagram
respects the paired conductor curves and the two triple-point classes.

\item Over $v\ne0$, $Q_N$ is exactly the cover defined by
\[
 \mathbb Z\langle\widehat\gamma,\widehat u,
 N\widehat w,N\widehat\delta\rangle
 \subset
 \mathbb Z\langle\widehat\gamma,\widehat u,
 \widehat w,\widehat\delta\rangle.
\]
Its deck group is $(\mathbb Z/N\mathbb Z)^2$.  The pulled-back cusp
monodromy is $M_0^N$, and in the displayed cover basis it is $M_0$:
\begin{equation}
 S_N^{-1}M_0^NS_N=M_0,\qquad
 S_N=\operatorname{diag}(1,1,N,N).
 \label{cusp-cover-monodromy-conjugacy}
\end{equation}
\end{enumerate}
\end{theorem}

\begin{proof}
We prove the assertions in the order in which the constructions occur.

\medskip\noindent
\emph{The source action.}
For $\bar\lambda\in\bar\Lambda$, the integral linear map
\[
 \varphi_{\bar\lambda}(y,y_3)
 =(y+y_3B_0\bar\lambda,y_3)
\]
translates the height-one slice by the integral vector
$B_0\bar\lambda$.  It preserves the $A_2$ triangulation and hence induces
a toric automorphism $\Phi_{\bar\lambda}$ of $Y^{(N)}$.  On the dense
torus this automorphism is
\[
 \Phi_{\bar\lambda}(\xi,v)
 =(v^{B_0\bar\lambda}\xi,v).
\]
Multiplication by the nowhere-zero holomorphic function
$c^{(N)}_{\bar\lambda}(v)$ extends across every toric chart.  Since
$C_N(v)\bar\lambda$ is additive in $\bar\lambda$, the resulting maps
satisfy
$\Psi^{(N)}_{\bar\lambda}\Psi^{(N)}_{\bar\mu}
=\Psi^{(N)}_{\bar\lambda+\bar\mu}$.

For completeness, freeness and proper discontinuity can be checked
directly.  Put
\[
 R_N(v):=-2\pi\operatorname{Im}C_N(v).
\]
It is bounded at $v=0$.  At a dense-torus point fixed by
$\bar\lambda$, logarithmic absolute values give
\[
 \log|v|B_0\bar\lambda+R_N(v)\bar\lambda=0.
\]
Choose the disc so small that
$|\log|v||>2\|R_N\|_\infty$.  Because $B_0$ is orthogonal, the last
equality is impossible for $\bar\lambda\ne0$.  On the central fibre,
$\Psi^{(N)}_{\bar\lambda}$ translates every bounded cell by the nonzero
vector $B_0\bar\lambda$, so it fixes no toric orbit.

For a dense-torus point set
\[
 y(\xi,v)=\frac{(\log|\xi_1|,\log|\xi_2|)}{\log|v|},
 \qquad
 B_{N,v}=B_0+\frac{R_N(v)}{\log|v|}.
\]
Then
\[
 y(\Psi^{(N)}_{\bar\lambda}(\xi,v))
 =y(\xi,v)+B_{N,v}\bar\lambda,
 \qquad
 |B_{N,v}\bar\lambda|\ge\tfrac12|\bar\lambda|.
\]
The logarithmic-position regions associated with any two fixed toric
charts are bounded.  Hence only finitely many $\bar\lambda$ can translate
one of them to meet the other.  Local finiteness of the fan reduces the
same assertion for compact sets to finitely many chart pairs, proving
proper discontinuity.  Finally, reducing
$B_{N,v}^{-1}y(\xi,v)$ modulo $\mathbb Z^2$ places every orbit meeting a
closed subdisc in finitely many closed unit toric charts.  Their union is
compact.  This proves properness of $f_N$.  The central-fibre and
normalization assertions follow exactly from the height-one fan: locally
$v$ is the product of the coordinates belonging to the rays through a
cell, every component occurs with multiplicity one, the lattice action is
transitive on vertices, and the star of a vertex is the complete smooth
hexagonal fan.

\medskip\noindent
\emph{The toric map and equivariance.}
Multiplication by $N$ on $N'$ sends every cone $\sigma\in\Sigma$ onto
itself as a real cone.  It therefore defines the toric morphism $P_N$.
For a maximal cone, let $z_0,z_1,z_2$ be the target characters dual to its
three primitive height-one ray generators, and let $w_0,w_1,w_2$ be the
corresponding source characters.  The map on convergent local rings is
\begin{equation}
 \mathbb C\{z_0,z_1,z_2\}\longrightarrow
\mathbb C\{w_0,w_1,w_2\},\qquad z_i\longmapsto w_i^N.
 \label{cusp-cover-maximal-chart-map}
\end{equation}
The target-chart preimage is the corresponding source chart, and the
right-hand ring is free of rank $N^3$ over the left-hand ring, with basis
\[
 \{w_0^{a_0}w_1^{a_1}w_2^{a_2}:0\le a_i<N\}.
\]
Thus $P_N$ is finite flat of degree $N^3$ on every maximal chart, hence
globally.

The choice \eqref{cusp-cover-source-factor} is forced by equivariance.
Indeed, writing powers coordinatewise,
\begin{align}
 P_N\Psi^{(N)}_{\bar\lambda}(\xi,v)
 &=\bigl((c^{(N)}_{\bar\lambda}(v))^N
          v^{NB_0\bar\lambda}\xi^N,v^N\bigr)\\
 &=\bigl(c_{\bar\lambda}(v^N)
          (v^N)^{B_0\bar\lambda}\xi^N,v^N\bigr)\\
 &=\Psi_{\bar\lambda}P_N(\xi,v).
 \label{cusp-cover-equivariance-computation}
\end{align}
Consequently $P_N$ descends to \eqref{cusp-cover-completed-map}.  To see
finiteness and flatness after the infinite-lattice quotient without any
implicit descent theorem, take an evenly covered open set in $N_0$ and
one of its lifts $V\subset Y$.  Equivariance gives a unique representative
of every source orbit above that open set whose image lies in $V$; two
such representatives cannot differ by a nontrivial lattice translate.
The descended map is therefore locally identified with
$P_N^{-1}(V)\to V$.  It is finite flat of degree $N^3$ by
\eqref{cusp-cover-maximal-chart-map}.

In the same chart the base characters are
\[
 t=z_0z_1z_2,\qquad v=w_0w_1w_2,\qquad P_N^*t=v^N.
\]
The Jacobian determinant of \eqref{cusp-cover-maximal-chart-map} is
\[
 N^3(w_0w_1w_2)^{N-1}=N^3v^{N-1}.
\]
This proves \eqref{cusp-cover-direct-ramification}; the local quotient
maps are biholomorphisms, so the calculation is unchanged downstairs.

\medskip\noindent
\emph{The relative degree-$N^2$ factor.}
The equality
\((f_0\circ\kappa_N)([\xi,v])=v^N=(\rho_N\circ f_N)([\xi,v])\)
gives the factorization
\eqref{cusp-cover-relative-factorization}.  At a triple point the exact
local-ring map is
\begin{equation}
 \begin{gathered}
 A_N:=
 \frac{\mathbb C\{z_0,z_1,z_2,v\}}
      {(z_0z_1z_2-v^N)},
 \qquad
 B:=\mathbb C\{w_0,w_1,w_2\},\\
 A_N\longrightarrow B,\qquad
 z_i\longmapsto w_i^N,\qquad
 v\longmapsto w_0w_1w_2.
 \end{gathered}
 \label{cusp-cover-basechange-local-ring}
\end{equation}
Let
\[
 G_N=\{(\varepsilon_0,\varepsilon_1,\varepsilon_2)\in\mu_N^3:
       \varepsilon_0\varepsilon_1\varepsilon_2=1\}
 \simeq\mu_N^2
\]
act diagonally on $B$.  A monomial
$w_0^{a_0}w_1^{a_1}w_2^{a_2}$ is $G_N$-invariant if and only if
\[
 a_0\equiv a_1\equiv a_2\pmod N.
\]
Writing their common residue as $r$, with $0\le r<N$, expresses that
monomial uniquely in the form
\[
 (w_0w_1w_2)^r
 (w_0^N)^{b_0}(w_1^N)^{b_1}(w_2^N)^{b_2}.
\]
The same coefficientwise argument applies to convergent power series.
Hence
\begin{equation}
 A_N=B^{G_N}.
 \label{cusp-cover-invariant-ring}
\end{equation}
It follows that $A_N$ is normal, that the map has generic degree
$|G_N|=N^2$, and that it is a finite quotient rather than a normalization.
The corresponding face-chart calculations prove normality along the
double and smooth strata as well, so $N_0[N]$ is normal everywhere.

The stabilizer in $G_N$ is trivial when none or exactly one of the $w_i$
vanishes, is isomorphic to $\mu_N$ when exactly two vanish, and is all of
$G_N$ when all three vanish.  Thus $Q_N$ is etale in codimension one and
has no ramification divisor, but it is not etale above the double and
triple strata.  It is also nonflat at every such point when $N>1$.  At a
triple point the closed fibre of \eqref{cusp-cover-basechange-local-ring}
has length
\begin{align}
 \dim_{\mathbb C}
 \frac{\mathbb C\{w_0,w_1,w_2\}}
      {(w_0^N,w_1^N,w_2^N,w_0w_1w_2)}
 &=N^3-(N-1)^3\\
 &=3N^2-3N+1,
 \label{cusp-cover-triple-fibre-length}
\end{align}
which exceeds the generic rank $N^2$.  At a point of a double stratum,
localizing at the nonzero third coordinate and choosing its holomorphic
$N$th root reduces the ramified factor to
\[
 \frac{\mathbb C\{z_0,z_1,v\}}{(z_0z_1-v^N)}
\longrightarrow\mathbb C\{w_0,w_1\},\qquad
 z_i\longmapsto w_i^N,\quad v\longmapsto w_0w_1.
\]
Its local generic rank is $N$, whereas its closed fibre has length
\begin{equation}
 \dim_{\mathbb C}
 \frac{\mathbb C\{w_0,w_1\}}
      {(w_0^N,w_1^N,w_0w_1)}=2N-1>N.
 \label{cusp-cover-double-fibre-length}
\end{equation}
The remaining transverse nonzero torus coordinate is etale locally.
Equations \eqref{cusp-cover-triple-fibre-length} and
\eqref{cusp-cover-double-fibre-length}, together with etaleness elsewhere,
identify the nonflat locus exactly.

\medskip\noindent
\emph{The regular fibres and monodromy.}
Choose a logarithm $r$ with $v=e^{2\pi ir}$.  The target period matrix at
$t=v^N$ is
\[
 Z(Nr)=NrB_0+C(v^N),
\]
whereas the source residual action is determined by
\[
 Z_N(r)=rB_0+\frac1NC(v^N)=\frac1N Z(Nr).
\]
The source fibre is therefore
\[
 \frac{\mathbb C^2}{\mathbb Z^2+Z_N(r)\mathbb Z^2},
\]
and multiplication by $N$ on $\mathbb C^2$ maps it to the target fibre.
Equivalently, in the unscaled additive coordinate it is the identity map
\begin{equation}
 \frac{\mathbb C^2}
      {N\mathbb Z^2+Z(Nr)\mathbb Z^2}
 \longrightarrow
 \frac{\mathbb C^2}
      {\mathbb Z^2+Z(Nr)\mathbb Z^2}.
 \label{cusp-cover-additive-fibre-map}
\end{equation}
The imaginary part of $Z(Nr)$ is invertible on a regular fibre, so its
period lattice meets $\mathbb Z^2$ trivially in the required real sense;
the index in \eqref{cusp-cover-additive-fibre-map} is $N^2$.  Its deck
group is $(\mathbb Z/N\mathbb Z)^2$.  In the ordered lattice basis
$(\widehat\gamma,\widehat u,\widehat w,\widehat\delta)$, this is exactly
the sublattice obtained by $S_N$.  Since $M_0=I+R$ with $R^2=0$ and
$R(a,b,c,d)=(0,0,b,-a)$, one has
\[
 M_0^N=I+NR,\qquad S_N^{-1}M_0^NS_N=M_0,
\]
proving \eqref{cusp-cover-monodromy-conjugacy}.

\medskip\noindent
\emph{The central normalization and conductor.}
Multiplication by $N$ fixes every ray of $\Sigma$.  It therefore maps a
chosen source component of the toric central fibre to the corresponding
target component.  The star of that ray is the hexagonal fan, so the
restriction is the toric power map
$[N]\colon\mathrm{dP}_6\to\mathrm{dP}_6$.  On its dense two-torus it has
degree $N^2$; in each toric boundary chart it is an $N$th-power map in the
normal coordinate, hence has ramification index $N$ along all six boundary
curves.  Equation \eqref{cusp-cover-equivariance-computation} says that
the source and target lattice identifications commute with this map.
Consequently opposite boundary curves, their conductor identifications,
and the two translation-orbits of triangle vertices descend compatibly.
This proves \eqref{cusp-cover-normalization-diagram} and completes the
proof.
\end{proof}

\begin{remark}[Why the unmodified action does not work]
\label{cusp-cover-naive-action-failure}
If one reuses the target action with $v$ in place of $t$ at the source,
then the two sides of equivariance have multiplicative factors
\[
 c_{\bar\lambda}(v)^N v^{NB_0\bar\lambda}
 \quad\hbox{and}\quad
 c_{\bar\lambda}(v^N)v^{NB_0\bar\lambda},
\]
respectively.  In general
$c_{\bar\lambda}(v)^N\ne c_{\bar\lambda}(v^N)$.  Thus the bare formula
$(\xi_1,\xi_2,v)\mapsto(\xi_1^N,\xi_2^N,v^N)$ extends on the ambient toric
variety but does not descend through the manuscript's quotient unless the
source factor \eqref{cusp-cover-source-factor} is inserted.  The explicit
logarithm $2\pi iC(v^N)\bar\lambda/N$ makes this root holomorphic; no
unjustified choice of an $N$th root is involved.
\end{remark}

\begin{remark}[Exact scope]
\label{cusp-cover-global-scope}
The theorem constructs a finite morphism of completed families over the
local cusp discs.  It is a degree-$N^3$ map to the original family and a
degree-$N^2$ map only after the target is base changed.  It is neither a
degree-$N^2$ self-map over the unchanged disc nor the normalization of the
base change.

Nothing here supplies a cover over the entire base $B$.  In fact the
displayed sublattice is not preserved by the original full monodromy when
$N>1$: for example, the printed matrix $A_1$ sends
$\widehat u$ to a vector whose $\widehat w$- and
$\widehat\delta$-coefficients are $-1$ and $1$.  A further global base
cover and compatible extensions through both finite fillings would have
to be constructed separately.  The local Kummer morphism also supplies
no map to a product sphere, smash sphere, $S^{12}$, or $S^{24}$.
\end{remark}


\subsection{Cone-compactified sphere carriers for the native cusp quotient}
\label{cusp-cone-section}

This subsection answers a topological question which is strictly different
from the finite-cover construction.  No period is enlarged by passing to a
cover.  Instead, the native expanding quotient is retained on the cusp, the
parts on which that quotient is unavailable are collapsed, and the resulting
based map is passed through a smash product.  The outcome is an exact global
map with source \(S^{12}\), and there is an analogous \(S^{24}\) map from the
full-monodromy-invariant circle quotient.  The targets are conically
stratified spaces.  The Euclidean limits below use the intrinsic metrics
inherited from the flat quotient tori on the regular slices; they are not
ambient metric limits in one fixed metric
on the compact targets.

Fix a closed coordinate disc
\begin{equation}
 D_0=\{t\in\mathbb C:|t|\leq\epsilon\}\subset B,\qquad
 p_0=\{t=0\},\qquad N_0=f^{-1}(D_0),\qquad W=f^{-1}(0),
 \label{cusp-cone-disc}
\end{equation}
small enough that it contains no singular value other than \(p_0\).  Put
\begin{equation}
 K_\Lambda=\mathbb Z\langle\widehat w,\widehat\delta\rangle,\qquad
 \overline\Lambda=\Lambda/K_\Lambda,\qquad
 \overline T^2=
 \frac{\overline\Lambda\otimes\mathbb R}{\overline\Lambda}.
 \label{cusp-cone-quotient-lattice}
\end{equation}
The ordered basis of \(\overline\Lambda\) is the image of
\((\widehat\gamma,\widehat u)\).

\begin{lemma}[The native quotient over the whole punctured cusp]
\label{cusp-cone-native-descent}
There is a well-defined real-analytic fibrewise group epimorphism
\begin{equation}
 \overline\pi:N_0\setminus W
 \longrightarrow D_0^*\times\overline T^2
 \label{cusp-cone-native-family-map}
\end{equation}
whose kernel is the connected two-torus family defined by \(K_\Lambda\).
The target local system is trivial.  On the positive radial fibre \(F_T\),
\(t=e^{-2\pi T}\), the isomorphism
\begin{equation}
 \iota_T:\overline T^2\longrightarrow Q_T,\qquad
 [(a,b)]\longmapsto[Y_T(a,b)^{\mathsf T}]
 \label{cusp-cone-radial-isometry}
\end{equation}
satisfies
\begin{equation}
 \pi_T=\iota_T\circ\overline\pi|_{F_T}.
 \label{cusp-cone-native-restriction}
\end{equation}
If \(\overline T^2\) is given the pullback by \(\iota_T\) of the
Euclidean quotient metric on \(Q_T=\mathbb R^2/Y_T\mathbb Z^2\), then
\(\iota_T\) is an isometry and
\begin{equation}
 \operatorname{sys}(\overline T^2,g_T)
 =\operatorname{sys}(Q_T)\geq T-C.
 \label{cusp-cone-native-systole}
\end{equation}
\end{lemma}

\begin{proof}
On the universal cover of \(D_0^*\), use the real-analytic flat marking
\[
 (z,x+\Lambda)\longmapsto[z,\Pi(z)x].
\]
The deck generator acts by
\[
 (z,x+\Lambda)\longmapsto(g_0z,M_0x+\Lambda).
\]
The displayed matrix \(M_0\) gives
\begin{equation}
 (M_0-I)(a,b,c,d)^{\mathsf T}=(0,0,b,-a)^{\mathsf T},\qquad
 \operatorname{im}(M_0-I)=K_\Lambda.
 \label{cusp-cone-monodromy-image}
\end{equation}
Consequently \(M_0\) preserves \(K_\Lambda\) and induces the identity on
\(\overline\Lambda\).  The formula
\[
 [z,\Pi(z)x]\longmapsto
 \bigl([z],[x\bmod K_\Lambda]\bigr)
\]
is invariant under the deck action and proves
\cref{cusp-cone-native-family-map}.  The lattice \(K_\Lambda\) is primitive,
so the kernel is connected.  On the radial fibre, projection to the
\((e_2,e_4)\)-coordinates sends
\(a\widehat\gamma+b\widehat u\) to \(Y_T(a,b)^{\mathsf T}\) and kills
\(K_\Lambda\).  This proves
\cref{cusp-cone-radial-isometry,cusp-cone-native-restriction}.  The metric
and systolic assertion is the quotient calculation in
\cref{cusp-spec-native-quotient}.
\end{proof}

Define
\begin{equation}
 \mathcal A=
 \frac{[0,\epsilon]\times S^1\times\overline T^2}
 {\{0\}\times S^1\times\overline T^2\sim v_0,\quad
  \{\epsilon\}\times S^1\times\overline T^2\sim v_1}.
 \label{cusp-cone-A-definition}
\end{equation}
Thus
\begin{equation}
 \mathcal A\cong
 \operatorname{Susp}(S^1\times\overline T^2)
 \cong\operatorname{Susp}(T^3),
 \label{cusp-cone-A-suspension}
\end{equation}
where \(\operatorname{Susp}\) denotes the unreduced suspension.  The extra
circle is the angular coordinate of the punctured base.  If one deliberately
forgets that angle, the evident map
\(\mathcal A\to\operatorname{Susp}(\overline T^2)\) gives a smaller target
with the same restriction to every fixed radial fibre.  We retain
\(\mathcal A\) because it records the entire punctured-cusp quotient.

\begin{theorem}[Exact localized \(S^{12}\) route]
\label{cusp-cone-S12-route}
Choose \(x_0\in W\), and base \(\mathcal A\) at \(v_0\).  There is a
continuous surjective based quotient map
\begin{equation}
 q_2:(X,x_0)\longrightarrow(\mathcal A,v_0)
 \label{cusp-cone-q2}
\end{equation}
with the following exact description.  For \(x\in N_0\setminus W\), write
\[
 t(f(x))=re^{i\phi},\qquad 0<r<\epsilon,
\]
and let \(\overline\pi_f(x)\in\overline T^2\) be the fibre component of
\(\overline\pi(x)\).  Then
\begin{equation}
 q_2(x)=[r,e^{i\phi},\overline\pi_f(x)],\qquad
 q_2(W)=v_0,\qquad
 q_2\bigl(X\setminus\operatorname{int}N_0\bigr)=v_1.
 \label{cusp-cone-q2-formula}
\end{equation}
On \(f^{-1}(|t|=\epsilon)\), the value is \(v_1\).

The smash square is a global based map
\begin{equation}
 q_2^{\wedge2}:X\wedge X\longrightarrow
 \mathcal A\wedge\mathcal A.
 \label{cusp-cone-S12-map}
\end{equation}
Using \(X\cong_{\mathrm{diff}}S^6\), a based choice of the diffeomorphism
identifies its source with
\begin{equation}
 X\wedge X\cong S^6\wedge S^6\cong S^{12}.
 \label{cusp-cone-S12-source}
\end{equation}
For \(0<e^{-2\pi T}<\epsilon\), the canonical product-to-smash quotient
maps restrict to embeddings on \(F_T^2\) and on the corresponding interior
target-product cross-section.  Under those embeddings and \(\iota_T^2\),
the induced smash map restricts exactly to
\begin{equation}
 \pi_T\times\pi_T:F_T\times F_T\longrightarrow Q_T\times Q_T.
 \label{cusp-cone-S12-fibre-map}
\end{equation}
With the Euclidean \(\ell^2\)-product flat quotient metric,
\begin{equation}
 \operatorname{sys}(Q_T\times Q_T)\geq T-C,\qquad
 (Q_T\times Q_T,(0,0))\longrightarrow(\mathbb R^4,0)
 \label{cusp-cone-S12-R4-limit}
\end{equation}
in pointed local-isometric, pointed smooth, and pointed
Gromov--Hausdorff topology.  Once \(T-C>4R\), the radius-\(R\) ball is
exactly Euclidean.

The integral homology of the targets is
\begin{equation}
 \widetilde H_k(\mathcal A;\mathbb Z)\cong
 \begin{cases}
  \mathbb Z^3,&k=2,3,\\
  \mathbb Z,&k=4,\\
  0,&\text{otherwise},
 \end{cases}
 \label{cusp-cone-A-homology}
\end{equation}
and
\begin{equation}
 \widetilde H_k(\mathcal A\wedge\mathcal A;\mathbb Z)\cong
 \begin{cases}
  \mathbb Z^9,&k=4,\\
  \mathbb Z^{18},&k=5,\\
  \mathbb Z^{15},&k=6,\\
  \mathbb Z^6,&k=7,\\
  \mathbb Z,&k=8,\\
  0,&\text{otherwise}.
 \end{cases}
 \label{cusp-cone-A-smash-homology}
\end{equation}
In particular, the target is not \(S^4\), \(S^{12}\), or a compact
four-manifold.
\end{theorem}

\begin{proof}
The equivalence relation in \cref{cusp-cone-A-definition} is closed: it is
the union of the diagonal with the squares of the two closed end
cross-sections.  A compact Hausdorff space modulo a closed equivalence
relation is compact Hausdorff.  Rescaling \(r\) to the unit interval gives
exactly the unreduced suspension in \cref{cusp-cone-A-suspension}.

Well-definedness of the interior formula follows from
\cref{cusp-cone-native-descent}.  A full turn in \(\phi\) changes the
marking by \(M_0\), but \(M_0\) is the identity modulo \(K_\Lambda\).

We check both collapsed strata.  If \(U\) is an open neighborhood of
\(v_0\), its inverse image in the cylinder is open and contains the compact
set \(\{0\}\times S^1\times\overline T^2\).  Hence, for some \(\delta>0\),
it contains
\([0,\delta)\times S^1\times\overline T^2\).  Therefore
\(q_2^{-1}(U)\) contains \(f^{-1}(|t|<\delta)\), a neighborhood of \(W\).
The same compactness argument at the other end shows that every neighborhood
of \(v_1\) contains the image of
\((\epsilon-\delta,\epsilon]\times S^1\times\overline T^2\).  Its preimage
together with \(X\setminus\operatorname{int}N_0\) is a neighborhood of the
outer boundary.  Continuity on the interior is
\cref{cusp-cone-native-family-map}; the closed pasting lemma gives
continuity on \(X\).

Every interior point has a preimage because \(\overline\pi\) is fibrewise
surjective.  The nonempty sets \(W\) and
\(X\setminus\operatorname{int}N_0\) map onto \(v_0\) and \(v_1\),
respectively.  Thus \(q_2\) is surjective.  A continuous surjection from
compact \(X\) to Hausdorff \(\mathcal A\) is a quotient map.

The map is based, so functoriality gives \cref{cusp-cone-S12-map}.  The
standard CW quotient gives \(S^m\wedge S^n\cong S^{m+n}\)
\cite[Chapter~0, Smash Product]{Hatcher2002}, proving
\cref{cusp-cone-S12-source}.

The base point lies in \(W\), so \(F_T^2\) is disjoint from the source fat
wedge.  Its image lies in an interior cross-section of
\(\mathcal A^2\), disjoint from the target fat wedge.  The source and target
fat wedges are closed in the compact Hausdorff products \(X^2\) and
\(\mathcal A^2\), so both smash quotients are Hausdorff.  Their restrictions
to the relevant compact subspaces are continuous and one-to-one, and are
therefore embeddings.
Equations \cref{cusp-cone-native-restriction,cusp-cone-q2-formula} give
\cref{cusp-cone-S12-fibre-map}.

The period lattice of the product target is
\(\operatorname{diag}(Y_T,Y_T)\mathbb Z^4\).  Its systole is the systole of
\(Y_T\mathbb Z^2\), proving the first assertion in
\cref{cusp-cone-S12-R4-limit}.  If \(x,y\in B_R(0)\subset\mathbb R^4\)
and \(0\ne\lambda\) is a product period, then
\[
 \|x-y+\lambda\|\geq\|\lambda\|-\|x-y\|
 \geq(T-C)-2R>2R\geq\|x-y\|
\]
when \(T-C>4R\).  The quotient projection is an exact isometry on the
Euclidean radius-\(R\) ball.  Conversely, any point of quotient distance
less than \(R\) has a shortest lattice representative of norm less than
\(R\), so this embedded ball is the entire quotient radius-\(R\) ball.
This proves all three pointed convergences.

Reduced suspension shifts reduced homology by one.  Since
\[
 H_1(T^3)=\mathbb Z^3,\qquad
 H_2(T^3)=\mathbb Z^3,\qquad H_3(T^3)=\mathbb Z,
\]
\cref{cusp-cone-A-homology} follows.  All groups are free, so the Kunneth
theorem for smash products has no Tor terms.  The ranks in degrees
\(4,5,6,7,8\) are
\[
 3\cdot3,\quad 2(3\cdot3),\quad
 2(3\cdot1)+3\cdot3,\quad 2(3\cdot1),\quad1,
\]
which proves \cref{cusp-cone-A-smash-homology}.
\end{proof}

\paragraph{The localized one-circle construction.}
The full-monodromy-invariant character also gives a smaller cusp-supported
target.  Define
\begin{equation}
 \mathcal B=
 \frac{[0,\epsilon]\times S^1\times S^1}
 {\{0\}\times S^1\times S^1\sim b_0,\quad
  \{\epsilon\}\times S^1\times S^1\sim b_1}
 \cong\operatorname{Susp}(T^2),
 \label{cusp-cone-B-definition}
\end{equation}
where the two circle coordinates in an interior cross-section are the base
angle and the value of \(\theta_\gamma\).

\begin{proposition}[Localized \(S^{24}\) route with target
\(\operatorname{Susp}(T^2)^{\wedge4}\)]
\label{cusp-cone-S24-local-route}
There is a continuous surjective based quotient map
\begin{equation}
 q_{1,\mathrm{cusp}}:(X,x_0)\longrightarrow(\mathcal B,b_0)
 \label{cusp-cone-q1-local}
\end{equation}
given by
\begin{equation}
 q_{1,\mathrm{cusp}}(x)
 =[r,e^{i\phi},\theta_\gamma(x)]
 \quad(0<|t|=r<\epsilon),\qquad
 q_{1,\mathrm{cusp}}(W)=b_0,\qquad
 q_{1,\mathrm{cusp}}
 \bigl(X\setminus\operatorname{int}N_0\bigr)=b_1.
 \label{cusp-cone-q1-local-formula}
\end{equation}
Its fourth smash power has source \(X^{\wedge4}\cong S^{24}\), and on
four synchronized radial fibres it is exactly
\begin{equation}
 \chi_{\gamma,T}^{\times4}:F_T^4\longrightarrow C_T^4.
 \label{cusp-cone-S24-local-fibre-map}
\end{equation}
Consequently its intrinsic target slices satisfy the Euclidean limit
\cref{cusp-cone-S24-R4-limit}.  The target homology is
\begin{equation}
 \widetilde H_k(\mathcal B;\mathbb Z)\cong
 \begin{cases}
  \mathbb Z^2,&k=2,\\
  \mathbb Z,&k=3,\\
  0,&\text{otherwise},
 \end{cases}
 \label{cusp-cone-B-homology}
\end{equation}
and
\begin{equation}
 \widetilde H_k(\mathcal B^{\wedge4};\mathbb Z)\cong
 \begin{cases}
  \mathbb Z^{16},&k=8,\\
  \mathbb Z^{32},&k=9,\\
  \mathbb Z^{24},&k=10,\\
  \mathbb Z^8,&k=11,\\
  \mathbb Z,&k=12,\\
  0,&\text{otherwise}.
 \end{cases}
 \label{cusp-cone-B-smash-homology}
\end{equation}
\end{proposition}

\begin{proof}
The monodromy invariance
\(\varepsilon_\gamma M_0=\varepsilon_\gamma\) makes the interior formula
single-valued after a turn in the base angle.  The two endpoint-continuity
arguments are identical to those in the proof of
\cref{cusp-cone-S12-route}, with compact cross-section \(T^2\) in place of
\(T^3\).  Fibrewise surjectivity of \(\theta_\gamma\), together with the
nonempty inner fibre and exterior, proves surjectivity.  Compact-to-Hausdorff
makes the map a quotient map.

Since \(x_0\in W\), the fourth smash power exists.  The radial regular
fibres and their circle targets avoid the two closed fat wedges.  The
restrictions of the Hausdorff smash quotients to these compact subspaces are
continuous and one-to-one, hence embeddings, and the fibre map is exactly
\cref{cusp-cone-S24-local-fibre-map}.  The intrinsic metric statement is
therefore the four-circle calculation
\cref{cusp-cone-S24-R4-limit}.  Finally, suspension shifts the free homology
of \(T^2\), giving \cref{cusp-cone-B-homology}.  In the fourth smash power,
choosing \(j\) of the four factors in degree \(3\) and the other
\(4-j\) factors in degree \(2\) gives degree \(8+j\) and rank
\(\binom4j2^{4-j}\).  This is \cref{cusp-cone-B-smash-homology}.
\end{proof}

\paragraph{The full-base circle target before coning.}
Choose pairwise disjoint closed coordinate discs \(D_0,D_1,D_2\) about the
three singular values and put
\begin{equation}
 P=B\setminus\bigcup_{j=0}^2\operatorname{int}D_j.
 \label{cusp-cone-pair-of-pants}
\end{equation}
For \(j=1,2\), let \(C_j\) be the solid-torus Seifert target in
\cref{cusp-spec-finite-Seifert-target}, and let \(h_j^{S^1}\) be the
target overlap map in \cref{cusp-spec-finite-overlap-identity}.  Glue
\begin{equation}
 \mathcal M=(P\times S^1)
 \underset{h_1^{S^1}}{\cup}C_1
 \underset{h_2^{S^1}}{\cup}C_2.
 \label{cusp-cone-M-definition}
\end{equation}
The gluing is along boundary tori over \(\partial D_1,\partial D_2\).
The remaining boundary is
\begin{equation}
 \partial\mathcal M=\partial D_0\times S^1\cong T^2.
 \label{cusp-cone-M-boundary}
\end{equation}
Define
\begin{equation}
 \mathcal C=\mathcal M\cup_{\partial\mathcal M}
 \operatorname{Cone}(\partial\mathcal M),
 \label{cusp-cone-C-definition}
\end{equation}
and write \(c_0\) for the cone vertex.
This changes the codomain at the cusp: the whole boundary two-torus,
including the target circle, is collapsed to \(c_0\).  It therefore does
not contradict the proved nonextension of the literal map
\(\chi_\gamma\) to the fixed product target \(B\times S^1\).

\begin{theorem}[Exact cone-compactified \(S^{24}\) route]
\label{cusp-cone-S24-route}
There is a continuous surjective based quotient map
\begin{equation}
 q_1:(X,x_0)\longrightarrow(\mathcal C,c_0),\qquad x_0\in W,
 \label{cusp-cone-q1}
\end{equation}
defined as follows:
\begin{enumerate}
\item over \(P\), it is
\((f,\theta_\gamma):f^{-1}(P)\to P\times S^1\), equivalently the
restriction of \(\chi_\gamma\) to \(f^{-1}(P)\);
\item over \(D_j\), \(j=1,2\), it is the finite-filling map
\(N_j\to C_j\) from \cref{cusp-spec-finite-Seifert-target};
\item on \(N_0\setminus W\), writing \(t=re^{i\phi}\), it is
\begin{equation}
 x\longmapsto[r,(e^{i\phi},\theta_\gamma(x))]
 \in\operatorname{Cone}(\partial D_0\times S^1),
 \label{cusp-cone-q1-cusp-formula}
\end{equation}
and every point of \(W\) maps to \(c_0\).
\end{enumerate}
The fourth smash power is a global based map
\begin{equation}
 q_1^{\wedge4}:X^{\wedge4}\longrightarrow\mathcal C^{\wedge4}.
 \label{cusp-cone-S24-map}
\end{equation}
Using \(X\cong_{\mathrm{diff}}S^6\), its source is
\begin{equation}
 X^{\wedge4}\cong(S^6)^{\wedge4}\cong S^{24}.
 \label{cusp-cone-S24-source}
\end{equation}
On four synchronized radial fibres, the canonical product-to-smash quotient
maps restrict to embeddings on the source fibre product and the interior
target-product cross-section.  Under those embeddings,
\cref{cusp-cone-S24-map} restricts exactly to
\begin{equation}
 \chi_{\gamma,T}^{\times4}:F_T^4\longrightarrow C_T^4,\qquad
 C_T=\mathbb R/\ell_T\mathbb Z.
 \label{cusp-cone-S24-fibre-map}
\end{equation}
With its Euclidean \(\ell^2\)-product flat quotient metric,
\begin{equation}
 \operatorname{sys}(C_T^4)=\ell_T
 =T+H-B+O(T^{-1}),\qquad
 (C_T^4,0)\longrightarrow(\mathbb R^4,0)
 \label{cusp-cone-S24-R4-limit}
\end{equation}
in pointed local-isometric, pointed smooth, and pointed
Gromov--Hausdorff topology.

The target homology is
\begin{equation}
 \widetilde H_k(\mathcal C;\mathbb Z)\cong
 \begin{cases}
  \mathbb Z,&k=2,3,\\
  0,&\text{otherwise},
 \end{cases}
 \label{cusp-cone-C-homology}
\end{equation}
and
\begin{equation}
 \widetilde H_k(\mathcal C^{\wedge4};\mathbb Z)\cong
 \begin{cases}
  \mathbb Z,&k=8,12,\\
  \mathbb Z^4,&k=9,11,\\
  \mathbb Z^6,&k=10,\\
  0,&\text{otherwise}.
 \end{cases}
 \label{cusp-cone-C-smash-homology}
\end{equation}
\end{theorem}

\begin{proof}
The compact spaces \(P\times S^1,C_1,C_2\) are glued along closed boundary
tori by homeomorphisms, producing a compact orientable three-manifold
\(\mathcal M\) with boundary \(\partial D_0\times S^1\).  Attaching the cone
along this closed boundary gives a compact Hausdorff finite CW space
\(\mathcal C\).

The overlap identity
\[
 \chi_\gamma\circ h_j=h_j^{S^1}\circ\overline\chi_j
\]
proves that the first two pieces of \(q_1\) agree.  On \(\partial D_0\), the
cone base is glued to the remaining boundary of \(P\times S^1\), so the
first and third pieces agree.  If \(U\) is a neighborhood of \(c_0\), its
inverse image in the cone contains \([0,\delta)\times T^2\) for some
\(\delta>0\), by compactness of the cone base.  Therefore \(q_1^{-1}(U)\)
contains \(f^{-1}(|t|<\delta)\), a neighborhood of \(W\).  This proves
continuity at the new stratum; continuity elsewhere follows from the input
maps and the closed pasting lemma.

Every regular circle quotient and both finite Seifert targets are
surjective, and \(W\) maps onto the cone vertex.  Thus \(q_1\) is a
continuous surjection from compact \(X\) to Hausdorff \(\mathcal C\), hence
a quotient map.  It is based, so its fourth smash power exists.  Iterating
\(S^m\wedge S^n\cong S^{m+n}\) gives \cref{cusp-cone-S24-source}.

The base point lies in \(W\), while each radial regular fibre is disjoint
from \(W\).  Hence \(F_T^4\) and its product-circle image avoid the closed
source and target fat wedges.  The canonical quotient maps from the source
product and the indicated interior target-product cross-section to their
smash products restrict to embeddings.  Under those embeddings, the induced
smash map restricts to \(\chi_{\gamma,T}^{\times4}\), which itself retains
its displayed nontrivial torus kernel.
The quotient-metric calculation in
\cref{cusp-spec-global-circle-circumference} gives \(\ell_T\).
The product lattice is \(\ell_T\mathbb Z^4\), so the radius-\(R\) argument,
including the shortest-representative proof that the embedded Euclidean ball
is the whole quotient ball, in the proof of \cref{cusp-cone-S12-route} proves
\cref{cusp-cone-S24-R4-limit}.

For homology, let \(\rho_0,\rho_1,\rho_2\) be the boundary loops of the pair
of pants, with \(\rho_0+\rho_1+\rho_2=0\), and let \(\eta\) be the regular
target-circle class.  With the boundary-loop and circle orientations used
in the finite-filling presentation, the meridian relations are
\begin{equation}
 3\rho_1-\eta=0,\qquad4\rho_2+\eta=0.
 \label{cusp-cone-Seifert-meridians}
\end{equation}
Reversing the circle orientation reverses both displayed \(\eta\)-signs
and changes none of the following groups.  Thus
\begin{equation}
 H_1(\mathcal M;\mathbb Z)
 =\frac{\mathbb Z\langle\eta,\rho_1,\rho_2\rangle}
 {\langle3\rho_1-\eta,\;4\rho_2+\eta\rangle}
 \cong\mathbb Z\langle g\rangle,
 \label{cusp-cone-M-H1}
\end{equation}
where
\begin{equation}
 \rho_0=-g,\qquad\rho_1=4g,\qquad
 \rho_2=-3g,\qquad\eta=12g.
 \label{cusp-cone-M-boundary-map}
\end{equation}
In particular,
\(H_1(\partial\mathcal M)=\mathbb Z\langle\rho_0,\eta\rangle
\to H_1(\mathcal M)\) is surjective.  The orientable three-manifold
\(\mathcal M\) has one torus boundary, \(H_3(\mathcal M)=0\), and
\(\chi(\mathcal M)=\frac12\chi(\partial\mathcal M)=0\); hence
\(H_2(\mathcal M;\mathbb Q)=0\).  To obtain the integral statement without
suppressing torsion, Poincare--Lefschetz duality and the universal
coefficient theorem give
\[
 H_2(\mathcal M;\mathbb Z)
 \cong H^1(\mathcal M,\partial\mathcal M;\mathbb Z)
 \cong \operatorname{Hom}
 \bigl(H_1(\mathcal M,\partial\mathcal M;\mathbb Z),\mathbb Z\bigr),
\]
because \(H_0(\mathcal M,\partial\mathcal M;\mathbb Z)=0\).  Thus
\(H_2(\mathcal M;\mathbb Z)\) is free; its rank is zero, so it vanishes.

The cone attachment is the homotopy cofiber of
\(\partial\mathcal M\hookrightarrow\mathcal M\).  Since the boundary
inclusion is a cofibration,
\[
 \widetilde H_k(\mathcal C;\mathbb Z)
 \cong H_k(\mathcal M,\partial\mathcal M;\mathbb Z).
\]
Poincare--Lefschetz duality and
\(H^0(\mathcal M)=H^1(\mathcal M)=\mathbb Z\),
\(H^2(\mathcal M)=H^3(\mathcal M)=0\), give
\cref{cusp-cone-C-homology}.  The homology is free, so the smash Kunneth
theorem says that four factors of degree \(2\) or \(3\) produce ranks
\(\binom4j\) in degree \(8+j\), proving
\cref{cusp-cone-C-smash-homology}.
\end{proof}

\begin{definition}[The three radial target-slice embeddings]
\label{cusp-cone-target-slice-embeddings}
Put \(r_T=e^{-2\pi T}\) and let
\[
 \vartheta_T:\mathbb R/\mathbb Z\longrightarrow
 C_T=\mathbb R/\ell_T\mathbb Z,\qquad [a]\longmapsto[\ell_Ta].
\]
For \(y_1,y_2\in Q_T\), define
\[
 j_T(y_1,y_2)=
 \bigl[[r_T,1,\iota_T^{-1}(y_1)],
       [r_T,1,\iota_T^{-1}(y_2)]\bigr]_{\wedge}
 \in\mathcal A^{\wedge2}.
\]
For \(c_1,\ldots,c_4\in C_T\), define
\[
 k_T^{\mathcal B}(c_1,\ldots,c_4)=
 \bigl[[r_T,1,\vartheta_T^{-1}(c_1)],\ldots,
       [r_T,1,\vartheta_T^{-1}(c_4)]\bigr]_{\wedge}
 \in\mathcal B^{\wedge4},
\]
and, using the cusp-cone piece of \(\mathcal C\),
\[
 k_T^{\mathcal C}(c_1,\ldots,c_4)=
 \bigwedge_{j=1}^4[r_T,(1,\vartheta_T^{-1}(c_j))]
 \in\mathcal C^{\wedge4}.
\]
These are embeddings because every displayed factor lies away from the smash
basepoint, so equality in the smash quotient is equality in every factor;
compact-to-Hausdorff injectivity then gives a topological embedding.
\end{definition}

\begin{proposition}[Intrinsic expansion versus ambient collapse]
\label{cusp-cone-metric-scope}
Let \(D_{12}\), \(D_{24}^{\mathcal B}\), and
\(D_{24}^{\mathcal C}\) be arbitrary fixed metrics compatible with the
topologies of \(\mathcal A^{\wedge2}\), \(\mathcal B^{\wedge4}\), and
\(\mathcal C^{\wedge4}\), respectively.  In their corresponding ambient
metrics, the radial target
slices satisfy
\begin{equation}
 \operatorname{diam}_{\mathrm{amb}}
 \bigl(j_T(Q_T^2)\bigr)\longrightarrow0,\qquad
 \operatorname{diam}_{\mathrm{amb}}
 \bigl(k_T^{\mathcal B}(C_T^4)\bigr)\longrightarrow0,\qquad
 \operatorname{diam}_{\mathrm{amb}}
 \bigl(k_T^{\mathcal C}(C_T^4)\bigr)\longrightarrow0.
 \label{cusp-cone-ambient-collapse}
\end{equation}
By contrast, in the intrinsic quotient metrics,
\begin{equation}
 \operatorname{diam}(Q_T^2)\geq\tfrac12(T-C)\longrightarrow\infty,\qquad
 \operatorname{diam}(C_T^4)\geq\tfrac12\ell_T\longrightarrow\infty.
 \label{cusp-cone-intrinsic-expansion}
\end{equation}
Thus \cref{cusp-cone-S12-R4-limit,cusp-cone-S24-R4-limit} assert
intrinsic, flat pointed limits of varying slices, not Sormani--Wenger
intrinsic-flat convergence and not an ambient pointed \(\mathbb R^4\) limit
inside one fixed compact metric target.
\end{proposition}

\begin{proof}
Every neighborhood of the relevant cone vertex contains every sufficiently
small radial cross-section, uniformly in all compact angular and torus
coordinates.  In a compatible metric, every such cross-section eventually
lies in an arbitrarily small ball about the vertex.  Products followed by
the smash quotient preserve this uniform convergence, proving
\cref{cusp-cone-ambient-collapse}.

For a flat torus with systole \(s\), choose a shortest primitive period
\(\lambda\).  The class of \(\lambda/2\) has distance \(s/2\) from zero:
for every period \(\mu\), the nonzero period \(\lambda-2\mu\) has length at
least \(s\).  Hence the diameter is at least \(s/2\).
Equations \cref{cusp-cone-native-systole,cusp-cone-S24-R4-limit} prove
\cref{cusp-cone-intrinsic-expansion}.
\end{proof}

\begin{proposition}[Exact almost-complex obstruction on \(S^{12}\)]
\label{cusp-cone-S12-almost-complex-obstruction}
Let \(M\) be a smooth manifold homeomorphic to \(S^{12}\).  There is no
bundle endomorphism \(J:TM\to TM\) satisfying \(J^2=-\operatorname{Id}_{TM}\).
This conclusion has two independent proofs which use only the homotopy type
of \(M\): the original Borel--Serre reduced-power calculation at the single
prime \(5\), and Bott integrality.  For the standard embedded sphere there is
also an independent stable-tangent Pontryagin-class contradiction.

These statements concern an unstabilized complex structure on the actual
rank-twelve tangent bundle.  A stable complex structure on
\(TS^{12}\oplus\varepsilon^r_{\mathbb R}\), a complex structure confined to
an open subset or to fibres of a map, or an endomorphism which ceases to be a
bundle endomorphism along a singular set is not an almost-complex structure
on \(TS^{12}\) and therefore lies outside the hypothesis.
\end{proposition}

\begin{proof}
Assume that \(J\) exists and put \(E=(TM,J)\), a complex vector bundle of
rank six.  Orient \(M\) by the interleaved complex frame
\[
 (v_1,Jv_1,\ldots,v_6,Jv_6),
\]
write \([M]_J\) for the resulting fundamental class, and choose
\(u_J\in H^{12}(M;\mathbb Z)\) by
\(\langle u_J,[M]_J\rangle=1\).  Since \(M\) has the integral cohomology of
\(S^{12}\),
\begin{equation}
 c_1(E)=c_2(E)=c_3(E)=c_4(E)=c_5(E)=0.
 \label{cusp-cone-S12-lower-Chern-zero}
\end{equation}
The top-Chern/Euler identity uses exactly the orientation just fixed, so
Poincare--Hopf gives
\begin{equation}
 c_6(E)=e(E_{\mathbb R})=e(TM)=2u_J.
 \label{cusp-cone-S12-top-Chern-two}
\end{equation}
Thus reversing a previously chosen orientation can change the displayed
generator but cannot change the coefficient two.

We first specialize the original calculation of Borel and Serre
\cite[\S12, Proposition~12.4, Lemma~12.5, and Proposition~12.6,
p.~430; \S15, Proposition~15.1, p.~434]{BorelSerre1953}.
They write \(C_{2j}\) for the universal class now denoted \(c_j\).
For \(j=6\), the admissible prime is exactly
\[
 p=5<6,\qquad 5\nmid6,
\]
and their identity in \(H^*(BU(6);\mathbb F_5)\) reads
\begin{equation}
 \overline c_6
 =6^{-1}\mathcal P_5^1(\overline c_2)
   +Q(\overline c_1,\ldots,\overline c_5).
 \label{cusp-cone-S12-Borel-Serre-specialization}
\end{equation}
Here \(6^{-1}=1\) in \(\mathbb F_5\).  To retain the exact coefficient rather
than suppressing the lower-class polynomial, pass to the splitting space,
write \(x_1,\ldots,x_6\) for the degree-two Chern roots, put
\(s_r=\sum_i x_i^r\), and use
\(\mathcal P_5^1(x_i)=x_i^5\) and the Cartan formula.  The resulting universal
identity is
\begin{align}
 \mathcal P_5^1(\overline c_2)
 ={}&\sum_{i<j}(x_i^5x_j+x_ix_j^5)=s_5s_1-s_6 \notag\\
 ={}&6\overline c_6+\overline c_1^4\overline c_2
 -\overline c_1^3\overline c_3
 -4\overline c_1^2\overline c_2^2
 +\overline c_1^2\overline c_4
 +7\overline c_1\overline c_2\overline c_3
 -\overline c_1\overline c_5 \notag\\
 &\quad+2\overline c_2^3-6\overline c_2\overline c_4
 -3\overline c_3^2.
 \label{cusp-cone-S12-P1-exact}
\end{align}
Consequently, for the lower-Chern ideal
\(I=(\overline c_1,\ldots,\overline c_5)\),
\begin{equation}
 \mathcal P_5^1(\overline c_2)
 \equiv6\overline c_6=\overline c_6\pmod I,
 \label{cusp-cone-S12-P1-lower-ideal}
\end{equation}
which verifies the exact top-class coefficient in
\cref{cusp-cone-S12-Borel-Serre-specialization}.  Pulling back along the
classifying map of \(E\), using \cref{cusp-cone-S12-lower-Chern-zero} and
naturality of \(\mathcal P_5^1\), gives
\[
 \overline{c_6(E)}=\mathcal P_5^1(\overline{c_2(E)})
 =\mathcal P_5^1(0)=0.
\]
But \cref{cusp-cone-S12-top-Chern-two} gives
\(\overline{c_6(E)}=2\overline u_J\ne0\) in
\(H^{12}(M;\mathbb F_5)\), the first contradiction.

For the second contradiction, choose an orientation-preserving homotopy
equivalence \(f:S^{12}\to M\).  Bott integrality says that the sixth Chern
class of every complex bundle on \(S^{12}\) is divisible by
\((6-1)!=120\)
\cite[Chapter~20, \S9, Theorem~9.6 and Corollary~9.8,
p.~308]{Husemoller1994}; see also
\cite[\S6]{KonstantisParton2018}.  The coefficient can be checked without
normalization ambiguity: Newton's sixth identity and
\cref{cusp-cone-S12-lower-Chern-zero} give
\[
 \operatorname{ch}_6(f^*E)
 =\frac{s_6(f^*E)}{6!}
 =-\frac{c_6(f^*E)}{5!}.
\]
The integral Bott generator therefore forces
\(c_6(f^*E)=120k f^*u_J\) for some \(k\in\mathbb Z\), whereas
\cref{cusp-cone-S12-top-Chern-two} gives
\(c_6(f^*E)=2f^*u_J\).  Thus \(120k=2\), which has no integral solution.

Finally take \(M=S^{12}\) with its standard smooth embedding in
\(\mathbb R^{13}\).  There is an explicit stable trivialization
\begin{equation}
 TS^{12}\oplus\varepsilon^1_{\mathbb R}
 \xrightarrow{\ \cong\ }\varepsilon^{13}_{\mathbb R},
 \qquad (x;v,a)\longmapsto(x;v+ax),
 \label{cusp-cone-S12-stable-trivialization}
\end{equation}
so \(p_3(TS^{12})=0\).  On the other hand,
\[
 E_{\mathbb R}\otimes_{\mathbb R}\mathbb C
 \cong E\oplus\overline E,
 \qquad
 p_3(E_{\mathbb R})=-c_6(E\oplus\overline E).
\]
Equation \cref{cusp-cone-S12-lower-Chern-zero} and
\(c_6(\overline E)=c_6(E)\) give the exact relation
\begin{equation}
 p_3(TS^{12})=-2c_6(E)=-4u_J,
 \label{cusp-cone-S12-Pontryagin-contradiction}
\end{equation}
contradicting stable triviality in the torsion-free group
\(H^{12}(S^{12};\mathbb Z)\).  This is the Pontryagin--Euler proof recorded
independently in
\cite[Lemma~4.1 and Theorem~4.2]{KonstantisParton2018}.

Every contradiction above requires a globally defined \(J\) on the actual
bundle \(TM\).  Stabilizing that bundle, deleting points, restricting to
regular fibres, or allowing \(J\) to fail at a singular locus removes that
hypothesis rather than producing an exception to it.
\end{proof}

\begin{theorem}[Sharp obstructions to stronger interpretations]
\label{cusp-cone-obstructions}
The sphere routes above have the following exact limitations.
\begin{enumerate}
\item Let \(r\geq1\), and let \(N\subset\Lambda^{\oplus r}\) be a
primitive kernel defining a connected fibrewise torus quotient over the
full regular base, or over any finite cover of that base.  Equip the quotient
with the norm induced by the unrescaled marked Euclidean metric on the
original fibres.  If its quotient systole tends to infinity on the radial
cusp, then
\begin{equation}
 \operatorname{rank}(\Lambda^{\oplus r}/N)\leq r.
 \label{cusp-cone-rank-bound}
\end{equation}
Thus two copies cannot give a full-base four-torus quotient.  The
\(S^{12}\) map in \cref{cusp-cone-S12-route} evades the obstruction only by
localizing to the cusp and collapsing the exterior.

\item No nontrivial finite base cover, including one extending with
branching over the original singular values, makes \(K_\Lambda\) a
sub-local system over the full regular base.  The local cyclic change
\(t=s^M\) used for the auxiliary degree-\(M^2\) cover cannot extend as a
connected cover of \(B\cong\mathbb P^1\) branched only at \(p_0\).

\item There is no map from \(X^2\), \(X\wedge X\), \(X^4\), or
\(X^{\wedge4}\) to a fixed \(T^4\) whose restriction retains all four
target character classes of the displayed product quotient.  There is no
map to a fixed \(S^4\) whose restriction is a degree-one collapse of
\(Q_T^2\), or of the four-circle target, for even one regular fibre.  The
conical targets evade these obstructions because the fibre character and
orientation classes do not extend as global target classes.

\item Cartesian products do not produce the named spheres:
\begin{equation}
 X^2\simeq S^6\times S^6\not\simeq S^{12},\qquad
 X^4\simeq(S^6)^4\not\simeq S^{24}.
 \label{cusp-cone-products-not-spheres}
\end{equation}
The two- and four-fold synchronized fibre products have real dimensions
\(10\) and \(18\), respectively, and generally have singular completed
fibres.  Joins give
\begin{equation}
 X*X\cong S^{13},\qquad X^{*4}\cong S^{27}.
 \label{cusp-cone-join-dimensions}
\end{equation}

\item Since \(X\cong S^6\) is simply connected, every connected unbranched
finite cover of \(X\), \(X^r\), \(X\wedge X\cong S^{12}\), or
\(X^{\wedge4}\cong S^{24}\) is trivial.  The local degree-\(M^2\)
unspooling cover is not the restriction of a nontrivial connected
unbranched cover of any of these compact sources.  Branched or normalized
compactifications are different objects and are not constructed here.

\item By \cref{cusp-cone-S12-almost-complex-obstruction}, no smooth manifold
homeomorphic to \(S^{12}\) admits an almost-complex structure.  Nor does
\(S^{24}\), with any differentiable structure, by the Borel--Serre theorem
\cite[\S15, Proposition~15.1 and the following remark,
p.~434]{BorelSerre1953}.  Thus the smash sources cannot be promoted to
compact complex spheres.
\end{enumerate}
\end{theorem}

\begin{proof}
Every element of \(K_\Lambda^{\oplus r}\) has uniformly bounded physical
length on the radial cusp.  If one survived in the quotient lattice, it
would be a nonzero period of bounded quotient norm.  Hence
\begin{equation}
 K_\Lambda^{\oplus r}\subset N.
 \label{cusp-cone-kernel-contained}
\end{equation}
For the full monodromy group,
\(A_1\widehat w=\widehat u-\widehat w\) and invariance force every copy of
\(\widehat u\) into \(N\).  Therefore
\[
 L_\Lambda^{\oplus r}
 =\mathbb Z\langle\widehat u,\widehat w,\widehat\delta\rangle^{\oplus r}
 \subset N,
\]
which proves \cref{cusp-cone-rank-bound}.

The same conclusion holds after a finite base cover.  Let \(H\) be its
finite-index subgroup and put \(H_*=A_2A_1=I+N_*\), \(N_*^2=0\).
Every finite-index subgroup contains \(H_*^n\) for some \(n>0\), while
\begin{equation}
 H_*^n\widehat w
 =\widehat w+n(\widehat u-\widehat\delta).
 \label{cusp-cone-finite-index-orbit}
\end{equation}
Equations
\cref{cusp-cone-kernel-contained,cusp-cone-finite-index-orbit} give
\(n\widehat u\in N\).  Primitivity gives \(\widehat u\in N\), proving the
same rank bound.  The orbit also gives
\(H_*^nK_\Lambda\ne K_\Lambda\) for every \(n>0\), so no finite cover of
the regular base preserves the native kernel.  A compact cover branched
over the singular values restricts to such a cover.  Finally, a connected
finite cover of \(\mathbb P^1\) branched only at \(p_0\) would restrict to
a connected unbranched finite cover of
\(\mathbb P^1\setminus\{p_0\}\cong\mathbb C\), hence has degree one.

For a fixed \(T^4\) target, \(H^1(X^r;\mathbb Z)=0\), as does the first
cohomology of the smash spheres.  A global map therefore pulls every target
character class back to zero.  On a regular fibre the quotient homomorphisms
inject the target character lattice into the character lattice of the
source torus, giving four independent nonzero pullbacks, a contradiction.

For the fixed \(S^4\) obstruction, let
\(\omega_T\in H^2(Q_T;\mathbb Z)\) be the orientation class.  Since
\(\pi_T^*:H^1(Q_T)\to H^1(F_T)\) injects two independent classes and torus
cohomology is their exterior algebra,
\begin{equation}
 (\pi_T\times\pi_T)^*(\omega_T\times\omega_T)
 \ne0\quad\text{in }H^4(F_T^2;\mathbb Z).
 \label{cusp-cone-fibre-H4-nonzero}
\end{equation}
A degree-one collapse \(Q_T^2\to S^4\) pulls the orientation generator
back to \(\pm(\omega_T\times\omega_T)\).  But
\[
 H^4(X^2;\mathbb Z)=H^4(S^{12};\mathbb Z)=0,
\]
so no global extension has this restriction.  The four-circle case is the
same: its orientation class pulls back to the exterior product of four
independent circle characters, while
\(H^4(X^4)=H^4(S^{24})=0\).

The Kunneth theorem gives \(H_6(X^2)=\mathbb Z^2\), whereas
\(H_6(S^{12})=0\), and gives nonzero groups of ranks \(4,6,4\) in degrees
\(6,12,18\) for \(X^4\), whereas \(S^{24}\) has none.  The real dimension
of an \(r\)-fold synchronized fibre product is
\(6r-2(r-1)=4r+2\), giving \(10\) and \(18\).  The join formula
\(S^m*S^n\cong S^{m+n+1}\) proves
\cref{cusp-cone-join-dimensions}.  The remaining assertions follow from
simple connectivity, \cref{cusp-cone-S12-almost-complex-obstruction}, and
the cited Borel--Serre theorem for \(S^{24}\).
\end{proof}

\paragraph{Canonicity and exact scope.}
The lattice \(K_\Lambda=\ker(M_0-I)=\operatorname{im}(M_0-I)\) and the
punctured-cusp quotient are intrinsic to the marked cusp monodromy.
The compact map \(q_2\) additionally depends on the chosen cusp coordinate,
closed cusp disc, and the decision to collapse its outer boundary and the
exterior.  The identifications
\(X^{\wedge2}\cong S^{12}\) and \(X^{\wedge4}\cong S^{24}\) depend on a
based choice of \(X\cong S^6\).  None of these choices is claimed to be
holomorphic or canonical in the algebraic category.

The map \(q_2^{\wedge2}\) is the direct two-copy construction from the
native \(T^2\times\mathbb R^2\) cusp: it is globally defined on the smash
sphere, and its radial regular-fibre quotients are \(Q_T\times Q_T\), which
converge pointed local-isometrically, smoothly, and Gromov--Hausdorff to
\(\mathbb R^4\).  It is not a torus quotient over the full
regular base.  The map \(q_{1,\mathrm{cusp}}^{\wedge4}\) is the parallel
cusp-supported construction with target
\(\operatorname{Susp}(T^2)^{\wedge4}\).  The stronger map
\(q_1^{\wedge4}\) uses the quotient over every regular and finite-filled
fibre before the cusp cone is attached; four copies are minimal for that
full-base all-expanding group quotient.
Neither construction gives a compact complex \(S^{12}\) or \(S^{24}\), an
ambient Euclidean limit in a fixed compact metric target, or a claim about
the interacting Yang--Mills spectrum.

The argument does not rule out a separately constructed branched or
toroidal compactification of the auxiliary finite covers after additional
branching away from \(p_0\).  Such a construction would have to specify its
global base cover, normalize and resolve every pulled-back singular fibre,
track the lifted metric, and determine the topology of the resulting total
space.  No such smooth compact sphere construction follows from the present
data.


% Canonical continuation of cusp_spectral_geometry.tex.  It uses the notation
% and labels of cusp_globalization_replacement.tex and
% cusp_cone_sphere_routes.tex, which are included immediately before it.

\subsubsection{Optimal sphere exponents and the exact Euclidean boundary}
\label{cusp-cone-optimal-strengthening}

Within the class of native connected fibrewise torus quotients equipped with
the induced, unscaled quotient of the marked Euclidean metric, the two sphere
dimensions in the preceding constructions are forced by two different
integral quotient problems.  Cusp-locally, each copy has an expanding
quotient of rank two, so two copies are necessary and sufficient within that
class.  Over the full regular base, monodromy leaves at most one expanding
quotient direction per copy, so four copies are necessary and sufficient
within the same class.  We now prove this scoped optimality, identify the
compact conical carriers, and construct a noncompact comparison manifold in
which the limiting \(\mathbb R^4\) is an actual boundary fibre.  There is only
a continuous comparison map from this manifold to the compact conical
carrier: it set-theoretically collapses the new boundary, is a homeomorphism
away from that boundary, and is nonproper and nonquotient.

Throughout this subsubsection,
\[
 B_R^Z(z):=\{w\in Z:d_Z(z,w)\leq R\}
\]
denotes the closed radius-\(R\) ball; when the ambient space and center are
clear, superscripts and centers are suppressed.

\begin{theorem}[Optimality within the native connected torus-quotient class]
\label{cusp-cone-optimal-exponents}
Consider only native connected fibrewise torus quotients of copies of the
original marked fibres, and equip every such quotient with the induced,
unscaled Euclidean quotient metric.

\begin{enumerate}
\item For a native connected torus quotient of \(r\) copies over the
punctured cusp, if the quotient systole tends to infinity on the positive
radial ray, then its lattice rank is at most \(2r\).  A rank-four quotient
therefore requires \(r\geq2\).  Equality is attained for \(r=2\) by
\[
 (\Lambda/K_\Lambda)^{\oplus2},
 \qquad K_\Lambda=
 \mathbb Z\langle\widehat w,\widehat\delta\rangle,
\]
whose physical period lattice is
\[
 L_T^{(12)}=\operatorname{diag}(Y_T,Y_T)\mathbb Z^4.
\]
The corresponding smash source is
\(X^{\wedge2}\cong S^{12}\).

\item For a native connected torus quotient of \(r\) copies over the full
regular base, or over a finite cover of that base, radial divergence of the
quotient systole forces its lattice rank to be at most \(r\).  A rank-four
quotient therefore requires \(r\geq4\).  Equality is attained by the four
copies of the full-monodromy-invariant character \(\chi_\gamma\); on the
radial fibre its period lattice is
\[
 L_T^{(24)}=\ell_T\mathbb Z^4,
 \qquad
 \ell_T=T+H-B+O(T^{-1}),
\]
and the corresponding smash source is
\(X^{\wedge4}\cong S^{24}\).
\end{enumerate}

Thus, among precisely these native connected torus quotients with their
induced unscaled marked Euclidean metrics, \(S^{12}\) is the smallest smash
sphere produced by the cusp-local rank-two quotient, and \(S^{24}\) is the
smallest smash sphere produced by a quotient which survives the full
regular-base monodromy.  No optimality is asserted here among arbitrary
continuous quotient maps or after a \(T\)-dependent rescaling of the metric.
\end{theorem}

\begin{proof}
For \(r\) copies, every element of
\[
 K_\Lambda^{\oplus r}
 =\mathbb Z\langle\widehat w,\widehat\delta\rangle^{\oplus r}
\]
has uniformly bounded physical length on the radial cusp.  Let
\(N\subset\Lambda^{\oplus r}\) be a primitive sublattice defining a
connected torus quotient.  If some \(v\in K_\Lambda^{\oplus r}\) had
nonzero class in \(\Lambda^{\oplus r}/N\), that nonzero quotient period
would have quotient norm at most the original norm of \(v\), uniformly in
\(T\).  The quotient systole could not diverge.  Hence
\[
 K_\Lambda^{\oplus r}\subset N,
\]
and
\[
 \operatorname{rank}(\Lambda^{\oplus r}/N)
 \leq 4r-2r=2r.
\]
For \(r=2\), the primitive kernel \(K_\Lambda^{\oplus2}\) attains this
bound.  The cusp monodromy is the identity modulo \(K_\Lambda\), so the
quotient descends over the punctured cusp.  Its radial lattice is exactly
\(\operatorname{diag}(Y_T,Y_T)\mathbb Z^4\), and
\[
 \operatorname{sys}(L_T^{(12)})
 =\operatorname{sys}(Y_T\mathbb Z^2)\geq T-C\longrightarrow\infty.
\]

For a quotient over the full regular base, monodromy invariance gives a
stronger bound.  Since \(\widehat w\in N\) and
\[
 A_1\widehat w=\widehat u-\widehat w,
\]
one has \(\widehat u\in N\).  Together with
\(\widehat w,\widehat\delta\in N\), this places the rank-three lattice
\[
 \mathbb Z\langle\widehat u,\widehat w,
 \widehat\delta\rangle^{\oplus r}
\]
inside \(N\), and hence
\[
 \operatorname{rank}(\Lambda^{\oplus r}/N)\leq r.
\]
The same conclusion holds after a finite base cover.  Indeed, if
\(\Gamma_{\mathrm{mon}}\) is the corresponding finite-index monodromy
subgroup and \(G_*=A_2A_1=I+N_*\), then \(N_*^2=0\), and
\(\Gamma_{\mathrm{mon}}\) contains \(G_*^m\) for
some \(m>0\).  The exact orbit formula
\[
 G_*^m\widehat w
 =\widehat w+m(\widehat u-\widehat\delta)
\]
implies \(m\widehat u\in N\); primitivity of \(N\) implies
\(\widehat u\in N\).

The invariant character \(\chi_\gamma\) has precisely the complementary
rank-one quotient in each copy.  The full-base construction, including the
two finite Seifert fillings and the cone at the unipotent cusp, proves that
four copies attain rank four.  Finally, after choosing the based
diffeomorphism \(X\cong_{\mathrm{diff}}S^6\) already reconstructed in this
workbench,
\(X^{\wedge r}\cong(S^6)^{\wedge r}\cong S^{6r}\), which gives the two
claimed sphere dimensions.
\end{proof}

\begin{theorem}[Exact sphere-to-sphere carriers with moving four-spheres]
\label{cusp-cone-angle-forgotten-sphere-carriers}
There are continuous surjective based quotient maps
\[
 \widehat q_3:(X,x_0)\longrightarrow(S^3,*),
 \qquad
 \widehat q_\gamma:(X,x_0)\longrightarrow(S^2,*),
\]
whose smash powers are continuous surjective quotient maps
\begin{equation}
 \widehat q_3^{\wedge2}:S^{12}\twoheadrightarrow S^6,
 \qquad
 \widehat q_\gamma^{\wedge4}:S^{24}\twoheadrightarrow S^8.
 \label{cusp-cone-sphere-to-sphere-maps}
\end{equation}
Use the marked presentation \(\overline T^2=\mathbb R^2/\mathbb Z^2\) and
put
\[
 K^{(1)}
 =((\mathbb R/\mathbb Z)\times\{0\})
  \cup(\{0\}\times(\mathbb R/\mathbb Z)),
 \qquad
 c:\overline T^2\longrightarrow\overline T^2/K^{(1)}\cong S^2.
\]
Here \(c\) is the quotient collapsing the coordinate one-skeleton.  For
every sufficiently large \(T\), put
\begin{align*}
 r_T&=e^{-2\pi T},\\
 K_T^{(1)}&:=\iota_T(K^{(1)})\subset Q_T,\\
 S_T^2&:=\overline T^2/K^{(1)}\cong Q_T/K_T^{(1)},\\
 c_T&:=c\circ\iota_T^{-1}:Q_T\longrightarrow S_T^2,
\end{align*}
where the identification of the two displayed quotients is induced by
\(\iota_T\).  Within this theorem and its proof, \(c_T\) denotes the
displayed collapse map; the earlier scalar \(c_T=\xi_T-\eta_T\) does not
enter.
At reduced-suspension cross-section height \(r_T\), the first target contains
a based topologically embedded four-sphere, equipped with its intrinsic
quotient length metric,
\[
 \mathbb S_{12,T}^4=S_T^2\wedge S_T^2\hookrightarrow S^6,
\]
and the second target contains, at the same cross-section height, a based
topologically embedded four-sphere, equipped with its intrinsic quotient
length metric,
\[
 \mathbb S_{24,T}^4=C_T^{\wedge4}\hookrightarrow S^8.
\]
On synchronized radial fibres, the exact restrictions factor as
\begin{align}
 F_T^2&\xrightarrow{\ \pi_T\times\pi_T\ }
 Q_T^2\xrightarrow{\ c_T\times c_T\ }(S_T^2)^2
 \xrightarrow{\ \rho_{12,T}\ }S_T^2\wedge S_T^2
 =\mathbb S_{12,T}^4,
 \label{cusp-cone-S12-S4-radial-map}\\
 F_T^4&\xrightarrow{\ \chi_{\gamma,T}^{\times4}\ }
 C_T^4\xrightarrow{\ \rho_{24,T}\ }C_T^{\wedge4}
 =\mathbb S_{24,T}^4,
 \label{cusp-cone-S24-S4-radial-map}
\end{align}
Here \(c_T\) collapses \(K_T^{(1)}\), and \(\rho_{12,T}\) and
\(\rho_{24,T}\) are the smash quotient maps.  Thus the
full collapsed locus in \(Q_T^2\) is exactly
\begin{equation}
 W_{12,T}
 =\bigl(K_T^{(1)}\times Q_T\bigr)
  \cup\bigl(Q_T\times K_T^{(1)}\bigr),
 \label{cusp-cone-S12-exact-collapsed-locus}
\end{equation}
and the full collapsed locus in \(C_T^4\) is exactly
\begin{equation}
 W_{24,T}
 =\{(x_1,x_2,x_3,x_4)\in C_T^4:
       x_j=0\pmod{\ell_T}\ \text{for some }j\}.
 \label{cusp-cone-S24-exact-collapsed-locus}
\end{equation}
The maps
\(\rho_{12,T}\circ(c_T\times c_T):Q_T^2\to
S_T^2\wedge S_T^2\) and
\(\rho_{24,T}:C_T^4\to C_T^{\wedge4}\) send exactly \(W_{12,T}\) and
\(W_{24,T}\), respectively, to the smash basepoint and are injective on
their complements.  The full composites from \(F_T^2\) and \(F_T^4\)
still identify every pair of points identified by
\(\pi_T\times\pi_T\) and \(\chi_{\gamma,T}^{\times4}\), respectively; in
particular, their first arrows retain their nontrivial torus kernels.
Their exact smash-basepoint preimages in the full radial sources are
\[
 \bigl(\pi_T^{-1}(K_T^{(1)})\times F_T\bigr)
 \cup\bigl(F_T\times\pi_T^{-1}(K_T^{(1)})\bigr)
\]
and
\[
 \{(x_1,x_2,x_3,x_4)\in F_T^4:
   \chi_{\gamma,T}(x_j)=0\ \text{in }C_T\ \text{for some }j\},
\]
respectively.  Outside these basepoint preimages, two source points have the
same final image if and only if their images under the corresponding first
arrow in \cref{cusp-cone-S12-S4-radial-map,cusp-cone-S24-S4-radial-map}
are equal.

Choose the Euclidean Riemannian \(\ell^2\)-product length metrics on
\(Q_T^2\) and \(C_T^4\) before taking the smash quotients, and use the
successive quotient length metrics specified in the proof.  With the
specified midpoint basepoints,
\begin{equation}
 (\mathbb S_{12,T}^4,\zeta_{12,T})
 \longrightarrow(\mathbb R^4,0),
 \qquad
 (\mathbb S_{24,T}^4,\zeta_{24,T})
 \longrightarrow(\mathbb R^4,0)
 \label{cusp-cone-moving-S4-R4-limits}
\end{equation}
in pointed local-isometric and pointed Gromov--Hausdorff topology.  Moreover,
the convergence is smooth on the eventually Euclidean basepoint
neighborhoods.  For every \(R>0\), once \(T-C>4R\), respectively once
\(\ell_T>4R\), the corresponding quotient map restricts to an isometry from
the radius-\(R\) ball in \(Q_T^2\), respectively \(C_T^4\), \emph{onto} the
entire radius-\(R\) ball in the indicated smash sphere.  The domain ball, and
hence the quotient ball, is exactly a Euclidean four-ball.
\end{theorem}

\begin{proof}
Write \(\Sigma_r(Z,z_0)\) for the reduced suspension
\[
 \frac{[0,\epsilon]\times Z}
 {(\{0,\epsilon\}\times Z)\cup([0,\epsilon]\times\{z_0\})},
\]
where the displayed union is collapsed to one basepoint.

Use the marked presentation
\(\overline T^2=\mathbb R^2/\mathbb Z^2\), let
\[
 K^{(1)}=
 ((\mathbb R/\mathbb Z)\times\{0\})
 \cup(\{0\}\times(\mathbb R/\mathbb Z)),
\]
and let
\[
 c:\overline T^2\longrightarrow
 \overline T^2/K^{(1)}\cong S^2
\]
be the degree-one quotient obtained by collapsing the one-skeleton.  The
formula
\[
 \widehat p_2:\mathcal A\longrightarrow
 \Sigma_r(\overline T^2/K^{(1)},c(K^{(1)}))\cong S^3,
 \qquad
 [r,e^{i\phi},u]\longmapsto[r,c(u)]
\]
is well defined: the two suspension ends of \(\mathcal A\) and the entire
cylinder over \(K^{(1)}\) map to the reduced-suspension basepoint.  It is a
continuous surjection.  Set
\(\widehat q_3=\widehat p_2\circ q_2\).

For every sufficiently large \(T\), define the transported one-skeleton and
its collapse by
\begin{equation}
 \begin{aligned}
 K_T^{(1)}&:=\iota_T(K^{(1)})\subset Q_T,\\
 S_T^2&:=\overline T^2/K^{(1)}\cong Q_T/K_T^{(1)},\\
 c_T&:=c\circ\iota_T^{-1}:Q_T\longrightarrow S_T^2.
 \end{aligned}
 \label{cusp-cone-transported-skeleton-collapse}
\end{equation}
The displayed target identification is the one induced by \(\iota_T\), so
\(c_T^{-1}(*)=K_T^{(1)}\), where
\(*=c(K^{(1)})=c_T(K_T^{(1)})\).

Give \(Q_T\) its Euclidean quotient length metric and \(S_T^2\) the quotient
length metric under \(c_T\).  Put
\[
 y_*=[(1/2,1/2)]\in\overline T^2,
 \qquad y_T:=\iota_T(y_*)\in Q_T,
 \qquad z_T:=c_T(y_T)\in S_T^2.
\]
Equivalently, the metric on \(\overline T^2\) used below is
\(g_T:=\iota_T^*g_{Q_T}\).  For every \(v\in\mathbb R^2\), the exact
decomposition
\[
 Y_T=T\begin{pmatrix}0&1\\-1&0\end{pmatrix}+E_T,
 \qquad\|E_T\|_{\mathrm{op}}\leq C,
\]
gives
\[
 \|Y_Tv\|\geq(T-C)\|v\|.
\]
The standard flat distance from \(y_*\) to \(K^{(1)}\) is \(1/2\);
therefore
\begin{equation}
 d_T:=\operatorname{dist}_{Q_T}(y_T,K_T^{(1)})
 =\operatorname{dist}_{g_T}(y_*,K^{(1)})
 \geq\frac{T-C}{2}.
 \label{cusp-cone-distance-to-skeleton}
\end{equation}
For a length-space quotient which collapses a closed subset \(K\) to one
point, the distance between two noncollapsed classes represented by \(x,y\)
is
\[
 \min\{d(x,y),d(x,K)+d(y,K)\}.
\]
The two displayed terms are upper bounds, obtained respectively from a
direct path and from paths from \(x\) and \(y\) to \(K\).  Conversely, a
rectifiable quotient path either avoids the collapsed point, in which case
its compact image lifts to \(X\setminus K\) and has length at least
\(d(x,y)\), or meets that point, in which case its two pieces have total
length at least \(d(x,K)+d(y,K)\).  Taking infima proves the formula.
If \(T-C>4R\), then \(d_T>2R\).  The displayed formula and the systolic
bound show that
\[
 c_T:B_R^{Q_T}(y_T)\longrightarrow B_R^{S_T^2}(z_T)
\]
is an isometry onto the entire target ball, and both balls are exactly
Euclidean two-balls.  Indeed, the second term in the minimum cannot realize
a distance at most \(2R\), and any point of the target ball has a unique
noncollapsed representative; the same formula puts that representative in
the source ball.

At reduced-suspension coordinate \(r_T=e^{-2\pi T}\), \(S_T^2\) is a based
embedded subspace of \(S^3\): its point \(*\) is the global
reduced-suspension basepoint.  Hence
\[
 \mathbb S_{12,T}^4=S_T^2\wedge S_T^2
\]
embeds in \(S^3\wedge S^3=S^6\).  Equip \(Q_T^2\) first with the Euclidean
Riemannian \(\ell^2\)-product length metric, equip \((S_T^2)^2\) with the
\(\ell^2\)-product of the two quotient length metrics, and let
\[
 \rho_{12,T}:(S_T^2)^2\longrightarrow S_T^2\wedge S_T^2
\]
carry the quotient length metric obtained by collapsing the fat wedge
\[
 \mathcal W_{12,T}
 =(S_T^2\times\{*\})\cup(\{*\}\times S_T^2).
\]
Since \(c_T^{-1}(*)=K_T^{(1)}\), the inverse image of
\(\mathcal W_{12,T}\) under \(c_T\times c_T\) is exactly \(W_{12,T}\) in
\cref{cusp-cone-S12-exact-collapsed-locus}.  On
\(Q_T^2\setminus W_{12,T}\), the map \(c_T\times c_T\) is injective and its
image avoids \(\mathcal W_{12,T}\); the map \(\rho_{12,T}\) is injective on
the complement of \(\mathcal W_{12,T}\).  Hence their composite identifies
all of \(W_{12,T}\) to the smash basepoint and makes no other
identifications.

The point \(\zeta_{12,T}=[z_T,z_T]\) has distance \(d_T\) from this fat
wedge.  The one-factor onto isometry above gives an onto isometry from the
radius-\(R\) ball about \((y_T,y_T)\) in \(Q_T^2\) to the corresponding
ball about \((z_T,z_T)\) in \((S_T^2)^2\).  Applying the same collapsed-set
distance formula to \(\rho_{12,T}\), and using \(d_T>2R\), shows that
\begin{equation}
 \rho_{12,T}\circ(c_T\times c_T):
 B_R^{Q_T^2}((y_T,y_T))
 \xrightarrow{\ \cong_{\mathrm{isom}}\ }
 B_R^{\mathbb S_{12,T}^4}(\zeta_{12,T})
 \label{cusp-cone-S12-onto-exact-ball}
\end{equation}
is onto as well as isometrically embedded.  The systolic bound
\(\operatorname{sys}(Q_T^2)\geq T-C>4R\) makes its domain exactly a
Euclidean four-ball.  On a radial fibre,
\[
 (\widehat p_2\circ q_2)|_{F_T}
 =[r_T,c(\iota_T^{-1}(\pi_T(\,\cdot\,)))]
 =[r_T,c_T(\pi_T(\,\cdot\,))],
\]
and taking two factors followed by \(\rho_{12,T}\) proves the entire
factorization in \cref{cusp-cone-S12-S4-radial-map}.

For the four-copy construction, choose \(1\in S^1\) as character basepoint
and define
\[
 \widehat p_1:\mathcal B\longrightarrow\Sigma_r(S^1,1)\cong S^2,
 \qquad
 [r,e^{i\phi},\theta]\longmapsto[r,\theta].
\]
The formula forgets the base angle and collapses the two suspension ends
and the entire cylinder with \(\theta=1\) to the reduced-suspension
basepoint.  It is a continuous surjection.  Put
\(\widehat q_\gamma=\widehat p_1\circ q_{1,\mathrm{cusp}}\).

At the interior suspension coordinate \(r_T=e^{-2\pi T}\), identify the
character circle \(C_T=\mathbb R/\ell_T\mathbb Z\) with
\(\mathbb R/\mathbb Z\subset S^2\) by \([x]\mapsto[x/\ell_T]\).  It is a
based embedded circle, so its fourth smash is an embedded
\[
 \mathbb S_{24,T}^4=C_T^{\wedge4}\subset(S^2)^{\wedge4}=S^8.
\]
Equip \(C_T^4\) first with the Euclidean Riemannian \(\ell^2\)-product
length metric, and let
\[
 \rho_{24,T}:C_T^4\longrightarrow C_T^{\wedge4}
\]
carry the quotient length metric obtained by collapsing exactly the fat
wedge \(W_{24,T}\) of
\cref{cusp-cone-S24-exact-collapsed-locus}.  Away from that fat wedge the
smash quotient is injective, so no further points are identified.  The
midpoint
\[
 \zeta_{24,T}
 =[(\ell_T/2,\ell_T/2,\ell_T/2,\ell_T/2)]
\]
has distance \(\ell_T/2\) from that fat wedge.  If \(\ell_T>4R\), neither a
nonzero period nor a path through the collapsed fat wedge can shorten a
distance between points in the radius-\(R\) ball of \(C_T^4\).  Conversely,
the collapsed-set distance formula shows that every point of the
radius-\(R\) quotient ball has its unique noncollapsed representative in
that product ball.  Consequently
\begin{equation}
 \rho_{24,T}:
 B_R^{C_T^4}((\ell_T/2,\ell_T/2,\ell_T/2,\ell_T/2))
 \xrightarrow{\ \cong_{\mathrm{isom}}\ }
 B_R^{\mathbb S_{24,T}^4}(\zeta_{24,T})
 \label{cusp-cone-S24-onto-exact-ball}
\end{equation}
is onto as well as isometrically embedded, and its domain is exactly a
Euclidean four-ball because \(\operatorname{sys}(C_T^4)=\ell_T>4R\).  On
the radial source the formula discards only the base angle and then performs
this smash quotient, giving
\cref{cusp-cone-S24-S4-radial-map}.

Both one-copy maps are continuous surjections from compact \(X\) to
Hausdorff spheres, hence quotient maps, and they are based.  Their smash
powers are surjective; compact-to-Hausdorff again makes them quotient maps.
Using the same chosen based diffeomorphism \(X\cong_{\mathrm{diff}}S^6\),
the homeomorphisms
\[
 X^{\wedge2}\cong S^{12},\quad(S^3)^{\wedge2}\cong S^6,
 \qquad
 X^{\wedge4}\cong S^{24},\quad(S^2)^{\wedge4}\cong S^8
\]
prove \cref{cusp-cone-sphere-to-sphere-maps}.  The exact Euclidean-ball
statements prove all three convergences in
\cref{cusp-cone-moving-S4-R4-limits}.
\end{proof}

\begin{corollary}[The two-copy rank-two sphere carrier is null-homotopic]
\label{cusp-cone-S12-sphere-map-null}
After the based identifications in
\cref{cusp-cone-sphere-to-sphere-maps}, the continuous surjective quotient
\[
 \widehat q_3^{\wedge2}:S^{12}\twoheadrightarrow S^6
\]
is based null-homotopic.  In fact, \(g\wedge g:S^{12}\to S^6\) is based
null-homotopic for every based map \(g:S^6\to S^3\).
\end{corollary}

\begin{proof}
Apply the full external-smash calculation in
\cref{cusp-s12-full-carrier-null} to
\(g=\widehat q_3\).  That theorem proves the stronger assertion that
\(x\wedge y:S^{12}\to S^6\) is null-homotopic for every pair of based
maps \(x,y:S^6\to S^3\).
\end{proof}

\begin{corollary}[Every repeated smash of the circle-quotient carrier is null]
\label{cusp-cone-S24-sphere-map-null}
Put \(f=\widehat q_\gamma:S^6\to S^2\) after the chosen based
identification \(X\cong S^6\).  For every \(r\geq2\),
\begin{equation}
 \bigl[f^{\wedge r}\bigr]=0
 \quad\text{in}\quad
 \pi_{6r}(S^{2r}).
 \label{cusp-cone-S24-null-maps}
\end{equation}
In particular,
\begin{equation}
 \bigl[\widehat q_\gamma^{\wedge2}\bigr]
 =0\in\pi_{12}(S^4),
 \qquad
 \bigl[\widehat q_\gamma^{\wedge4}\bigr]
 =0\in\pi_{24}(S^8).
 \label{cusp-cone-S24-null-special-cases}
\end{equation}
Thus the displayed surjective quotient
\(\widehat q_\gamma^{\wedge4}:S^{24}\twoheadrightarrow S^8\) is globally
null-homotopic.  This does not alter either its point-set surjectivity or the
exact intrinsic \(\mathbb R^4\) limit on the moving four-spheres
\(\mathbb S_{24,T}^4\): those are properties of the specified representative
and its radial restrictions, whereas null-homotopy classifies that
representative only up to based homotopy.
\end{corollary}

\begin{proof}
Associativity of the smash product gives the
exact factorization
\begin{equation}
 \begin{tikzcd}[column sep=large]
 S^6\wedge S^6
 \arrow[r,"f\wedge\operatorname{id}_{S^6}"]&
 S^2\wedge S^6
 \arrow[r,"\operatorname{id}_{S^2}\wedge f"]&
 S^2\wedge S^2 .
 \end{tikzcd}
 \label{cusp-cone-S24-smash-square-factorization}
\end{equation}
Under the standard based sphere identifications, the first arrow is the
sixfold suspension of \(f\), up to the harmless twist homeomorphisms dictated
by the left-versus-right suspension convention:
\begin{equation}
 f\wedge\operatorname{id}_{S^6}
 =\Sigma^6f:S^{12}\longrightarrow S^8.
 \label{cusp-cone-S24-sixfold-suspension}
\end{equation}
Freudenthal's suspension theorem is already in its isomorphism range here,
because \(12\leq 2\cdot8-2\), and remains in that range after every further
suspension.  It therefore identifies \(\pi_{12}(S^8)\) with the fourth
stable stem.  The low-stem calculation gives
\begin{equation}
 \pi_{12}(S^8)\cong\pi_4^{\mathrm S}=0
 \quad
 \text{\cite[Corollary~4.24, p.~360]{Hatcher2002};
 \cite[Table~1, p.~113]{IsaksenWangXu2023}}.
 \label{cusp-cone-fourth-stable-stem}
\end{equation}
Hence \(\Sigma^6f\) is based null-homotopic, so the composite in
\cref{cusp-cone-S24-smash-square-factorization}, namely
\(f\wedge f\), is based null-homotopic.  For \(r\geq2\), the canonical
associativity homeomorphisms identify
\begin{equation}
 f^{\wedge r}=(f\wedge f)\wedge f^{\wedge(r-2)},
\end{equation}
with the convention \(f^{\wedge0}=\operatorname{id}_{S^0}\).  Smashing a
based null-homotopy with any based map is a based null-homotopy.  This proves
\cref{cusp-cone-S24-null-maps,cusp-cone-S24-null-special-cases}.
\end{proof}

\begin{corollary}[Literal \(\mathbb R^4\) open cells in the sphere targets]
\label{cusp-cone-literal-R4-open-cells}
Let \(*\) denote the smash basepoint and set
\[
 U_{12,T}:=\mathbb S_{12,T}^4\setminus\{*\},
 \qquad
 U_{24,T}:=\mathbb S_{24,T}^4\setminus\{*\}.
\]
The embedded moving four-spheres restrict to topological embeddings
\[
 U_{12,T}\cong\mathbb R^4\hookrightarrow S^6,
 \qquad
 U_{24,T}\cong\mathbb R^4\hookrightarrow S^8.
\]
More exactly, the quotient-stage maps restrict to homeomorphisms
\begin{align}
 (Q_T\setminus K_T^{(1)})^2
 &\xrightarrow{\ \cong\ }U_{12,T},
 \label{cusp-cone-U12-homeomorphism}\\
 (C_T\setminus\{0\})^4
 &\xrightarrow{\ \cong\ }U_{24,T}.
 \label{cusp-cone-U24-homeomorphism}
\end{align}
These copies of \(\mathbb R^4\) are nonclosed topological open cells: they
are open in their moving four-spheres, not open in the ambient \(S^6\) or
\(S^8\), and their closures in those ambient spheres are
\(\mathbb S_{12,T}^4\) and \(\mathbb S_{24,T}^4\), respectively.  They are
not globally isometric to Euclidean \(\mathbb R^4\) for either the assigned
compact quotient metrics or the restriction of any fixed compatible metric
on the compact sphere targets.  The exact midpoint balls in
\cref{cusp-cone-S12-onto-exact-ball,cusp-cone-S24-onto-exact-ball} give the
precise metric statement: as \(T\to\infty\), every fixed Euclidean
radius-\(R\) ball occurs isometrically and onto the corresponding pointed
ball.
\end{corollary}

\begin{proof}
For a quotient which collapses a closed subset \(K\) to a point, restriction
to the complement gives a homeomorphism
\[
 X\setminus K\xrightarrow{\ \cong\ }(X/K)\setminus\{*\}.
\]
Indeed, the complement of the quotient point is open, its inverse image is
the saturated open set \(X\setminus K\), and restriction of a quotient map
over an open subset of its target is again a quotient map.  The restricted
map is bijective, and a bijective quotient map is a homeomorphism.
Consequently
\[
 Q_T\setminus K_T^{(1)}
 \cong S_T^2\setminus\{*\}.
\]
The marked coordinate one-skeleton \(K^{(1)}\) is the \(1\)-skeleton of the
standard one-vertex CW structure on \(\overline T^2\), with one \(0\)-cell,
two \(1\)-cells, and one \(2\)-cell.  Its complement is therefore the open
\(2\)-cell \((0,1)^2\cong\mathbb R^2\).  Transport by
\(\iota_T\) gives
\[
 Q_T\setminus K_T^{(1)}\cong\mathbb R^2.
\]
Likewise \(C_T\setminus\{0\}\) is an open interval and hence is homeomorphic
to \(\mathbb R\).

The complement of the smash basepoint is the product of the complements of
the factor basepoints.  Therefore
\[
 (S_T^2\wedge S_T^2)\setminus\{*\}
 \cong(S_T^2\setminus\{*\})^2\cong\mathbb R^4,
\]
and
\[
 C_T^{\wedge4}\setminus\{*\}
 \cong(C_T\setminus\{0\})^4\cong\mathbb R^4.
\]
The full moving four-spheres are topologically embedded in their compact
sphere targets, so these restrictions are topological embeddings.  In each
case the omitted smash basepoint lies in the closure of the open cell, which
proves the closure assertions.  Finally, every compatible metric on a
compact target is bounded, as is each compact quotient metric assigned to a
moving four-sphere, whereas Euclidean \(\mathbb R^4\) is unbounded.  Thus no
global isometry is possible; the exact metric content is precisely the
onto-ball calculation already proved.
\end{proof}

\begin{remark}[The moving four-spheres are null-homotopic in the targets]
\label{cusp-cone-moving-S4-null-homotopy}
For a fixed interior reduced-suspension coordinate, the inclusion
\(S_T^2\hookrightarrow S^3\) contracts through based maps by sliding that
coordinate to a suspension endpoint; throughout the homotopy the collapsed
meridian remains at the global basepoint.  Hence
\(S_T^2\wedge S_T^2\hookrightarrow S^6\) is null-homotopic.  The identical
sliding argument contracts \(C_T\hookrightarrow S^2\), and therefore also
contracts \(C_T^{\wedge4}\hookrightarrow S^8\).  Thus the moving metric
four-spheres do not contradict the fixed-target orientation obstruction:
their fundamental classes die in the ambient targets.
\end{remark}

\begin{remark}[Homotopical status of the two sphere-to-sphere maps]
\label{cusp-cone-sphere-map-homotopy-status}
The maps \(S^{12}\to S^6\) and \(S^{24}\to S^8\) induce zero on all
positive-degree integral homology and cohomology groups and have zero
ordinary Hurewicz image.  They have no classical degree because source and
target dimensions differ.  The full four-copy map
\(\widehat q_\gamma^{\wedge4}:S^{24}\to S^8\) is based null-homotopic by
\cref{cusp-cone-S24-sphere-map-null}.  The distinct two-copy map
\(\widehat q_3^{\wedge2}:S^{12}\to S^6\) is also based null-homotopic, by
\cref{cusp-cone-S12-sphere-map-null}.  In both cases the vanishing is a
theorem about the full specified sphere-valued map, not merely its Hurewicz
image or the inclusion of a moving four-sphere. The one-copy class
\([\widehat q_3]\in\pi_6(S^3)\) is now zero by
Theorem~\ref{cusp-null:thm}; the one-copy class
\([\widehat q_\gamma]\in\pi_6(S^2)\) remains open.
\end{remark}

\begin{proposition}[The compact cone-on-\(T^4\) carriers]
\label{cusp-cone-explicit-carriers}
Fix \(T_0\) sufficiently large and put
\(r_0=e^{-2\pi T_0}<\epsilon\).  Define
\[
 \mathcal D_{12}
 =\frac{[0,r_0]\times(\overline T^2)^2}
 {\{0\}\times(\overline T^2)^2\sim d_{12}},
 \qquad
 \mathcal D_{24}
 =\frac{[0,r_0]\times(S^1)^4}
 {\{0\}\times(S^1)^4\sim d_{24}}.
\]
Both spaces are homeomorphic to \(\operatorname{Cone}(T^4)\).  There are
topological embeddings
\begin{align}
 \delta_{12}:\mathcal D_{12}&\hookrightarrow
 \mathcal A\wedge\mathcal A,
 &
 [r,u_1,u_2]&\longmapsto
 [\,[r,1,u_1],[r,1,u_2],],
 \label{cusp-cone-D12-embedding}\\
 \delta_{24}:\mathcal D_{24}&\hookrightarrow
 \mathcal C^{\wedge4},
 &
 [r,\theta_1,\ldots,\theta_4]&\longmapsto
 \bigwedge_{j=1}^4[r,(1,\theta_j)].
 \label{cusp-cone-D24-embedding}
\end{align}
Here \(1\in S^1\) fixes the positive radial base angle.  For every
\(T\geq T_0\), put \(r_T=e^{-2\pi T}\).  The restrictions of the sphere
quotient maps to the synchronized radial fibres are exactly the
cross-sections at height \(r_T\):
\begin{align}
 F_T^2&\xrightarrow{\ \pi_T\times\pi_T\ }
 Q_T^2\cong\{r_T\}\times T^4
 \subset\mathcal D_{12},
 \label{cusp-cone-D12-slice}\\
 F_T^4&\xrightarrow{\ \chi_{\gamma,T}^{\times4}\ }
 C_T^4\cong\{r_T\}\times T^4
 \subset\mathcal D_{24}.
 \label{cusp-cone-D24-slice}
\end{align}
Every one of these embedded four-torus slices is null-homotopic in its
carrier, and hence in the full smash target.

At the cone vertex \(d_\nu\), \(\nu\in\{12,24\}\), the local integral
homology is
\begin{equation}
 H_k(\mathcal D_\nu,
       \mathcal D_\nu\setminus\{d_\nu\};\mathbb Z)
 \cong
 \begin{cases}
  \mathbb Z^4,&k=2,\\
  \mathbb Z^6,&k=3,\\
  \mathbb Z^4,&k=4,\\
  \mathbb Z,&k=5,\\
  0,&\text{otherwise}.
 \end{cases}
 \label{cusp-cone-carrier-local-homology}
\end{equation}
In particular, \(d_\nu\) is not an integral homology-manifold point, with or
without boundary, of any dimension.  Consequently it is not a topological
manifold point, with or without boundary, of any dimension.
\end{proposition}

\begin{proof}
The two displayed domains are cones on
\((\overline T^2)^2\cong T^4\) and \((S^1)^4\cong T^4\).  At \(r=0\),
every factor on the right of
\cref{cusp-cone-D12-embedding,cusp-cone-D24-embedding} is the relevant
smash basepoint, so the formulae descend from the cylinders to the cones.
For \(0<r\leq r_0\), none of the factors is the smash basepoint.  Equality
in the smash quotient is then equality in every factor, and the coordinate
descriptions of the suspension and cusp cone give equality of \(r\) and of
all torus coordinates.  Thus each \(\delta_\nu\) is injective.  Its domain
is compact and its target is Hausdorff, so it is an embedding.

The already proved one-factor restriction formulae, with
\(r_T=e^{-2\pi T}\), are
\[
 q_2(x)=[r_T,1,\iota_T^{-1}(\pi_T(x))]
 \quad(x\in F_T)
\]
and, in the cusp-cone piece of \(\mathcal C\),
\[
 q_1(x)
 =[r_T,(1,\vartheta_T^{-1}(\chi_{\gamma,T}(x)))]
 \quad(x\in F_T).
\]
Applying these identities in two and four factors gives
\cref{cusp-cone-D12-slice,cusp-cone-D24-slice}.  Sliding the cone coordinate
from \(r_T\) to zero contracts each slice to the cone vertex.  Finally,
excision and the long exact sequence of the pair give, for \(L=T^4\),
\[
 H_k(\operatorname{Cone}(L),
     \operatorname{Cone}(L)\setminus\{v\};\mathbb Z)
 \cong\widetilde H_{k-1}(L;\mathbb Z).
\]
The ranks \(4,6,4,1\) of \(H_1,H_2,H_3,H_4\) of \(T^4\) prove
\cref{cusp-cone-carrier-local-homology}.  At an interior point of an
integral homology \(n\)-manifold the local integral homology is \(\mathbb Z\)
in degree \(n\) and zero in every other degree, whereas at a boundary point
of an integral homology manifold with boundary all these local groups
vanish.  The simultaneous nonzero groups in degrees \(2,3,4,5\), including
ranks greater than one, agree with neither pattern for any \(n\).  Every
topological manifold is locally an integral homology manifold, so the final
assertion follows as well.
\end{proof}

\begin{theorem}[Point-set quotient maps and exact homotopy types]
\label{cusp-cone-target-homotopy-types}
The smash maps are continuous surjective quotient maps
\[
 q_2^{\wedge2}:S^{12}\twoheadrightarrow\mathcal A^{\wedge2},
 \qquad
 q_1^{\wedge4}:S^{24}\twoheadrightarrow\mathcal C^{\wedge4}.
\]
Their targets have homotopy types
\begin{align}
 \mathcal A^{\wedge2}
 &\simeq
 \bigvee^9S^4\vee\bigvee^{18}S^5\vee
 \bigvee^{15}S^6\vee\bigvee^6S^7\vee S^8,
 \label{cusp-cone-A-smash-wedge}\\
 \mathcal C^{\wedge4}
 &\simeq
 S^8\vee\bigvee^4S^9\vee\bigvee^6S^{10}
 \vee\bigvee^4S^{11}\vee S^{12}.
 \label{cusp-cone-C-smash-wedge}
\end{align}
Both maps induce zero on every positive-degree integral homology and
cohomology group and have zero ordinary Hurewicz image.  Neither map has a
continuous homotopy section.  These conclusions do not determine whether
the unstable homotopy class of either map is zero.
\end{theorem}

\begin{proof}
Let \(f:(K,k_0)\to(L,\ell_0)\) be a surjective based map between compact
Hausdorff spaces.  Every non-basepoint class of \(L^{\wedge r}\) has a
representative \((y_1,\ldots,y_r)\) with each \(y_i\ne\ell_0\).  Choosing
\(x_i\in f^{-1}(y_i)\) proves surjectivity of \(f^{\wedge r}\); the smash
basepoint is also in its image.  The source smash is compact.  The target
fat wedge is closed, hence the target smash is Hausdorff.  A continuous
surjection from a compact space to a Hausdorff space is a quotient map.
This proves the first assertion for \(q_2\) and \(q_1\).

The suspension splitting
\[
 \Sigma(X_1\times\cdots\times X_n)
 \simeq
 \bigvee_{\varnothing\ne I\subseteq\{1,\ldots,n\}}
 \Sigma\!\left(\bigwedge_{i\in I}X_i\right)
\]
gives
\[
 \mathcal A=\Sigma T^3
 \simeq\bigvee^3S^2\vee\bigvee^3S^3\vee S^4.
\]
Finite smash products distribute over finite wedges.  The ordered degree
sums give multiplicities \(9,18,15,6,1\) in dimensions
\(4,5,6,7,8\), proving \cref{cusp-cone-A-smash-wedge}.

For \(\mathcal C\), write \(x=\rho_1\), \(y=\rho_2\), and
\(h=\eta\).  The pair-of-pants product and the two exact Seifert fillings
give
\[
 \pi_1(\mathcal M)=
 \langle x,y,h\mid[x,h]=[y,h]=1,
 x^3h^{-1}=1,\ y^4h=1\rangle.
\]
The remaining boundary torus is generated by \(h\) and
\(\rho_0=(xy)^{-1}\).  Coning this torus kills their normal closure, and
therefore
\[
 \pi_1(\mathcal C)
 \cong\langle x,y\mid x^3=1,\ y^4=1,\ xy=1\rangle=1,
\]
because \(y=x^{-1}\) and \(\gcd(3,4)=1\).  The previously computed
relative homology gives
\[
 \widetilde H_2(\mathcal C;\mathbb Z)=\mathbb Z,
 \qquad
 \widetilde H_3(\mathcal C;\mathbb Z)=\mathbb Z,
\]
and no other reduced homology.  Since \(\mathcal C\) is a simply connected
finite three-complex, Hurewicz supplies a sphere representing the
\(H_2\)-generator, and the Whitehead exact sequence makes
\(\pi_3(\mathcal C)\to H_3(\mathcal C)\) surjective.  A wedge of
representing maps
\(S^2\vee S^3\to\mathcal C\) is a homology equivalence between simply
connected CW complexes, hence a homotopy equivalence.  Thus
\[
 \mathcal C\simeq S^2\vee S^3.
\]
Taking the fourth smash power yields \cref{cusp-cone-C-smash-wedge} with
the binomial multiplicities \(1,4,6,4,1\).

The first target has no reduced homology above degree eight and the second
has none above degree twelve.  The source spheres have no positive homology
below degrees twelve and twenty-four.  Hence all positive-degree homology
and cohomology homomorphisms induced by the two maps vanish.  Their
Hurewicz images lie in \(H_{12}(\mathcal A^{\wedge2})=0\) and
\(H_{24}(\mathcal C^{\wedge4})=0\), respectively.  A section would split a
nonzero lower homology group of the target into the zero lower homology of
the source sphere, which is impossible.  Higher homotopy groups of the
displayed wedges need not vanish, so none of these arguments decides the
unstable null-homotopy question.
\end{proof}

\begin{theorem}[Exact decompactification spaces with boundary
\(\mathbb R^4\)]
\label{cusp-cone-exact-decompactification-space}
For \(\nu\in\{12,24\}\), set
\[
 A_T^{(12)}=\operatorname{diag}(Y_T,Y_T),
 \qquad
 A_T^{(24)}=\ell_T I_4,
 \qquad s_0=T_0^{-1}.
\]
On \([0,s_0]\times\mathbb R^4\), define an equivalence relation
\(\sim_\nu\) fibrewise over \(s\) by
\begin{align*}
 (0,x)\sim_\nu(0,y)&\iff x=y,\\
 (s,x)\sim_\nu(s,y)&\iff
 x-y\in A_{1/s}^{(\nu)}\mathbb Z^4
 \quad(0<s\leq s_0).
\end{align*}
Let
\[
 \mathfrak E_\nu
 =([0,s_0]\times\mathbb R^4)/\!\sim_\nu,
 \qquad p_\nu:\mathfrak E_\nu\longrightarrow[0,s_0].
\]
Then \(\mathfrak E_\nu\) is a Hausdorff, second-countable topological
five-manifold with boundary, and
\begin{align}
 p_{12}^{-1}(s)&\cong Q_{1/s}^2 &(s>0),&
 &p_{12}^{-1}(0)&\cong\mathbb R^4,
 \label{cusp-cone-E12-fibres}\\
 p_{24}^{-1}(s)&\cong C_{1/s}^4 &(s>0),&
 &p_{24}^{-1}(0)&\cong\mathbb R^4.
 \label{cusp-cone-E24-fibres}
\end{align}
The other boundary component is the compact four-torus over \(s=s_0\).

Give every positive fibre its Euclidean quotient metric and the zero fibre
its Euclidean metric.  For every \(R>0\), there is \(s_R>0\) such that,
for \(0\leq s<s_R\), writing
\(q_{\nu,s}(x)=[s,x]\), the quotient map restricts to an exact isometry
\[
 q_{\nu,s}:B_R^{\mathbb R^4}(0)
 \xrightarrow{\ \cong_{\mathrm{isom}}\ }
 B_R^{p_\nu^{-1}(s)}([s,0])
\]
onto the entire quotient radius-\(R\) ball, not merely onto an embedded
subset of that ball.
Consequently the positive fibres converge to the zero fibre in pointed
local-isometric, pointed smooth, and pointed Gromov--Hausdorff topology.

There are continuous surjective comparison maps which set-theoretically
collapse the zero boundary fibres,
\[
 \begin{tikzcd}[column sep=large]
 \mathbb R^4=p_{12}^{-1}(0)
 \arrow[r,hook]&
 \mathfrak E_{12}
 \arrow[r,"\beta_{12}"]&
 \mathcal D_{12}
 \arrow[r,hook,"\delta_{12}"]&
 \mathcal A^{\wedge2},
 \end{tikzcd}
\]
and
\[
 \begin{tikzcd}[column sep=large]
 \mathbb R^4=p_{24}^{-1}(0)
 \arrow[r,hook]&
 \mathfrak E_{24}
 \arrow[r,"\beta_{24}"]&
 \mathcal D_{24}
 \arrow[r,hook,"\delta_{24}"]&
 \mathcal C^{\wedge4}.
 \end{tikzcd}
\]
For \(s>0\), \(\beta_\nu\) is the specified marked homeomorphism from
\(p_\nu^{-1}(s)\) to the radial torus slice at \(r=e^{-2\pi/s}\).  Its
restriction
\[
 \beta_\nu:\mathfrak E_\nu\setminus p_\nu^{-1}(0)
 \xrightarrow{\ \cong\ }\mathcal D_\nu\setminus\{d_\nu\}
\]
is a homeomorphism, while
\[
 \beta_\nu^{-1}(d_\nu)=p_\nu^{-1}(0)\cong\mathbb R^4.
\]
Neither \(\beta_\nu\) is proper or a quotient map.  It factors through a
continuous bijection
\[
 \overline\beta_\nu:
 \mathfrak E_\nu/p_\nu^{-1}(0)\longrightarrow\mathcal D_\nu,
\]
but this bijection is not a homeomorphism; equivalently, the topology of
\(\mathcal D_\nu\) is not the quotient topology induced by \(\beta_\nu\).
Thus the limiting \(\mathbb R^4\) is an actual boundary fibre of
\(\mathfrak E_\nu\), whereas the continuous comparison map sends that entire
fibre to the cone vertex.
\end{theorem}

\begin{proof}
We give a common proof.  Write
\[
 L_s=A_{1/s}^{(\nu)}\mathbb Z^4\quad(s>0),
 \qquad L_0=\{0\}.
\]
The matrices vary continuously for \(s>0\), and
\[
 \operatorname{sys}(L_s)\longrightarrow\infty
 \qquad(s\downarrow0)
\]
by \(\operatorname{sys}(Y_T\mathbb Z^2)\geq T-C\) in the
\(S^{12}\) case and \(\operatorname{sys}(\ell_T\mathbb Z^4)=\ell_T\)
in the \(S^{24}\) case.

The equivalence relation is closed.  Indeed, suppose
\((s_j,x_j)\sim_\nu(s_j,y_j)\) and
\((s_j,x_j,y_j)\to(s,x,y)\).  If \(s>0\), write
\(y_j-x_j=A_{1/s_j}^{(\nu)}n_j\).  On a compact neighborhood of \(s\),
the inverse matrices are uniformly bounded, so boundedness of
\(y_j-x_j\) forces \(n_j\) to lie in a finite set.  Passing to a
subsequence makes \(n_j\) constant and proves \(x\sim_\nu y\).  If
\(s=0\), any nonzero \(n_j\) would give
\[
 \|y_j-x_j\|\geq\operatorname{sys}(L_{s_j})\longrightarrow\infty,
\]
contrary to convergence.  Hence \(x=y\).

The quotient projection \(q_\nu:[0,s_0]\times\mathbb R^4\to
\mathfrak E_\nu\) is open.  If \(U\) is open, its saturation is the union
of \(U\) with the images, for \(n\in\mathbb Z^4\), of
\[
 U\cap((0,s_0]\times\mathbb R^4)
 \quad\text{under}\quad
 (s,x)\longmapsto(s,x+A_{1/s}^{(\nu)}n),
\]
and every one of these images is open.  A closed equivalence relation with
an open quotient map has Hausdorff quotient: two inequivalent points admit
a product neighborhood disjoint from the relation, and the saturations of
the two factor neighborhoods are disjoint open sets.  Since the source is
second countable and the quotient map is open, images of a countable base
form a countable base of the quotient.

Fix \(x_0\in\mathbb R^4\).  Choose \(a>0\) and \(s_1>0\) such that
\(\operatorname{sys}(L_s)>2a\) for \(0<s<s_1\).  The quotient projection
is one-to-one on
\([0,s_1)\times B_a^\circ(x_0)\), where
\(B_a^\circ(x_0)=\{x:\|x-x_0\|<a\}\), and its restriction is an open continuous
bijection onto an open subset of \(\mathfrak E_\nu\); it is therefore a
homeomorphism.  These are half-space charts at \(s=0\).  The usual evenly
covered small balls give product charts for \(s>0\), and the same argument
at \(s=s_0\) gives the second boundary component.  This proves the
manifold assertion and the fibre identifications.

For the metric assertion, choose \(s_R\) so that
\(\operatorname{sys}(L_s)>4R\) for \(s<s_R\).  If
\(x,y\in B_R(0)\) and \(0\ne\lambda\in L_s\), then
\[
 \|x-y+\lambda\|
 \geq\|\lambda\|-\|x-y\|
 >4R-2R\geq\|x-y\|.
\]
Thus no translated representative shortens the Euclidean distance on that
ball, so \(q_{\nu,s}\) is an exact isometric embedding there.  It is also
onto the full quotient ball.  Indeed, if \(\xi=[s,x]\) satisfies
\(d([s,0],\xi)\leq R\), discreteness of \(L_s\) makes the infimum in
\[
 d([s,0],[s,x])=\min_{\lambda\in L_s}\|x+\lambda\|
\]
attained.  For a minimizing \(\lambda_0\), the representative
\(x+\lambda_0\) lies in \(B_R^{\mathbb R^4}(0)\) and maps to \(\xi\).  At
\(s=0\), \(q_{\nu,0}\) is the identity.  This proves the asserted onto
isometry, and since it holds for every fixed \(R\), it proves all three
pointed convergences.

It remains to define these comparison maps, which set-theoretically collapse
the zero fibres.  For \(s>0\), put
\(T=1/s\) and \(r=e^{-2\pi T}\).  In the \(S^{12}\) case, write
\(x=(x_1,x_2)\in\mathbb R^2\oplus\mathbb R^2\) and set
\[
 \beta_{12}([s,x])
 =\bigl[r,
 \iota_T^{-1}([x_1]_{Y_T\mathbb Z^2}),
 \iota_T^{-1}([x_2]_{Y_T\mathbb Z^2})\bigr].
\]
In the \(S^{24}\) case, write
\(x=(x_1,\ldots,x_4)\) and set
\[
 \beta_{24}([s,x])
 =\bigl[r,
 [x_1/\ell_T],\ldots,[x_4/\ell_T]\bigr].
\]
These formulae are well defined on the positive fibres and are continuous
there.  Define \(\beta_\nu([0,x])=d_\nu\) for every \(x\in\mathbb R^4\).
Every neighborhood of \(d_\nu\) in the cone contains all cross-sections
with sufficiently small \(r\).  Since \(e^{-2\pi/s}\to0\) as
\(s\downarrow0\), continuity at the zero fibre is uniform in \(x\).
Surjectivity follows because each positive fibre maps onto its complete
cross-section and the zero fibre maps onto the vertex.  The restriction away
from \(p_\nu^{-1}(0)\) is a homeomorphism because
\(s\mapsto e^{-2\pi/s}\) is a homeomorphism from \((0,s_0]\) to
\((0,r_0]\), and the displayed marked-torus identifications are fibrewise
homeomorphisms varying continuously with \(s\).  The vertex fibre is exactly
\(p_\nu^{-1}(0)\), so \(\beta_\nu\) is not proper.

It is also not a quotient map.  Choose
\(a\in\mathbb R^4\setminus\mathbb Z^4\), for example
\(a=(1/2,1/2,1/2,1/2)\), choose \(s_j\downarrow0\), and put
\[
 z_j=[s_j,A_{1/s_j}^{(\nu)}a]\in\mathfrak E_\nu.
\]
For \(\nu=12\),
\[
 \inf_{n\in\mathbb Z^4}
 \|A_{1/s_j}^{(12)}(a+n)\|
 \geq(1/s_j-C)\operatorname{dist}(a,\mathbb Z^4)
 \longrightarrow\infty,
\]
and for \(\nu=24\) the same infimum equals
\[
 \ell_{1/s_j}\operatorname{dist}(a,\mathbb Z^4)
 \longrightarrow\infty.
\]
More explicitly, for every fixed \(x_0\in\mathbb R^4\),
\[
 \inf_{n\in\mathbb Z^4}
 \|A_{1/s_j}^{(\nu)}(a+n)-x_0\|
 \geq
 \inf_{n\in\mathbb Z^4}\|A_{1/s_j}^{(\nu)}(a+n)\|-\|x_0\|
 \longrightarrow\infty.
\]
Hence \(\{z_j:j\geq1\}\) is closed in \(\mathfrak E_\nu\): it has no
accumulation in a positive fibre and the last estimate shows that it
eventually misses every half-space chart about every point of the zero
fibre.  But
\[
 \beta_\nu(z_j)=[e^{-2\pi/s_j},a]\longrightarrow d_\nu,
\]
so \(\{\beta_\nu(z_j):j\geq1\}\) is not closed in \(\mathcal D_\nu\).
Since \(\beta_\nu\) is injective off the zero fibre and the latter maps to
\(d_\nu\), the inverse image of this nonclosed set is the closed set
\(\{z_j:j\geq1\}\).  Thus \(\beta_\nu\) is not a quotient map.  This proves
the two diagrams and their exact topological scope.  Since \(\beta_\nu\) is
constant exactly on \(p_\nu^{-1}(0)\) and otherwise injective, it induces the
continuous bijection \(\overline\beta_\nu\) displayed in the theorem.  Were
that bijection a homeomorphism, \(\beta_\nu\) would be the composite of the
quotient projection with a homeomorphism and hence a quotient map, contrary
to the counterexample.  Therefore \(\overline\beta_\nu\) is not a
homeomorphism.
\end{proof}

\begin{corollary}[Sharp global meaning of the Euclidean limit]
\label{cusp-cone-sharp-global-R4-scope}
The native routes constructed above exhibit the underlying space
\(\mathbb R^4\) in the following three precise senses:
\begin{enumerate}
\item as each of the nonclosed topologically embedded open cells
\(U_{12,T}\subset S^6\) and \(U_{24,T}\subset S^8\);
\item as the pointed intrinsic Euclidean limit of the four-torus slices in the
global quotient carriers of \(S^{12}\) and \(S^{24}\); and
\item as the Euclidean decompactification boundary fibre of
\(\mathfrak E_{12}\) and \(\mathfrak E_{24}\).
\end{enumerate}
It does not produce \(\mathbb R^4\) as the whole compact sphere, as the
whole compact conical target, or as a continuous quotient target of either
compact space.  Nor does the unrescaled restriction of any fixed compatible
ambient target metric have \(\mathbb R^4\) as its slice limit, and the open
cells in item \(1\) are not globally Euclidean-isometric in such a metric.
\end{corollary}

\begin{proof}
The three positive assertions are
\cref{cusp-cone-literal-R4-open-cells,cusp-cone-explicit-carriers,%
cusp-cone-exact-decompactification-space}.
For the negative assertions, a continuous image of a compact sphere or
compact conical target is compact, so it cannot be all of \(\mathbb R^4\).
The local homology calculation
\cref{cusp-cone-carrier-local-homology} shows that the carrier vertex is not
an integral homology-manifold point, and hence not a manifold point, of any
dimension, with or without boundary.  Finally, in every fixed compatible
metric on a compact cone target, all sufficiently deep radial slices lie
uniformly in an arbitrarily small neighborhood of the vertex, so their ambient
diameters tend to zero.  Their intrinsic diameters tend to infinity because
their intrinsic systoles do.  Hence without a \(T\)-dependent rescaling no
fixed ambient metric realizes the intrinsic Euclidean limit on those slices.
The fixed cusp-adapted blow-ups are treated separately in
\cref{sphere-metric-fixed-realization}.  The comparison diagrams record the
unrescaled distinction exactly: they collapse the new \(\mathbb R^4\)
boundary rather than identify it with the compact target.
\end{proof}

\paragraph{Scope.}
The constructions above use only the original, unscaled fibre metrics and
the quotient maps derived from the native \(S^6\) cusp.  They do not use the
auxiliary finite-cover unspooling.  The conical carriers depend on the
chosen cusp coordinate, positive radial ray, and based identification
\(X\cong S^6\).  The spaces \(\mathfrak E_\nu\) additionally remember the
marked Euclidean quotient metrics.  None of these data defines a compact
complex structure on \(S^{12}\) or \(S^{24}\), a four-dimensional physical
spacetime, or an interacting Yang--Mills Hamiltonian.


\subsection{The full \texorpdfstring{\(S^{12}\to S^6\)}{S12-to-S6}
carrier is null-homotopic}
\label{cusp-s12-carrier-homotopy}

This subsection determines the based homotopy class of the full two-copy
sphere carrier.  It is important that the calculation concerns the full
quotient map, not merely the inclusion of one of its moving four-spheres.

Recall the exact maps from
\cref{cusp-cone-S12-route,cusp-cone-angle-forgotten-sphere-carriers}.  The
cusp quotient is
\[
 q_2:(X,x_0)\longrightarrow
 \mathcal A=\operatorname{Susp}(S^1\times\overline T^2),
\]
and, for the coordinate one-skeleton
\[
 K^{(1)}=
 ((\mathbb R/\mathbb Z)\times\{0\})
 \cup(\{0\}\times(\mathbb R/\mathbb Z))\subset\overline T^2,
\]
the degree-one quotient
\(c:\overline T^2\to\overline T^2/K^{(1)}\cong S^2\) defines
\[
 \widehat p_2([r,e^{i\phi},u])=[r,c(u)]
 \in\Sigma_r S^2\cong S^3.
\]
Thus
\[
 g:=\widehat q_3=\widehat p_2\circ q_2:(X,x_0)\longrightarrow(S^3,*)
\]
is the specified continuous surjective based quotient.  After the fixed
based identification \(X\cong_{\mathrm{diff}}S^6\), its smash square is the
specified continuous surjective based quotient
\[
 g^{\wedge2}:S^{12}=S^6\wedge S^6
 \longrightarrow S^3\wedge S^3=S^6.
\]

\begin{theorem}[The full external smash pairing is zero]
\label{cusp-s12-full-carrier-null}
For every pair of based maps \(x,y:S^6\to S^3\),
\[
 [x\wedge y]=0\qquad\text{in}\qquad\pi_{12}(S^6).
\]
In particular,
\(\widehat q_3^{\wedge2}:S^{12}\twoheadrightarrow S^6\) is simultaneously
surjective as a point-set quotient and based null-homotopic.
\end{theorem}

\begin{proof}
Write \(E\) for reduced suspension, using the ordered smash coordinates on
spheres.  The external smash product has the exact factorization
\begin{equation}
 x\wedge y
 =
 (x\wedge\operatorname{id}_{S^3})
 \circ
 (\operatorname{id}_{S^6}\wedge y):
 S^6\wedge S^6\longrightarrow S^6\wedge S^3
 \longrightarrow S^3\wedge S^3.
 \label{cusp-s12-smash-factorization}
\end{equation}
Under the displayed ordered sphere identifications, the two factors, in
order of application, represent \(E^6y\in\pi_{12}(S^9)\) and
\(E^3x\in\pi_9(S^6)\).  Equivalently, Toda's
Barratt--Hilton reduced-join formula gives, with
\(p=q=k=h=3\),
\begin{equation}
 [x\wedge y]=(E^3[x])\circ(E^6[y]).
 \label{cusp-s12-toda-smash-formula}
\end{equation}
Here Toda's sign is \((-1)^{(p+k)h}=(-1)^{18}=+1\).  Reversing the
two displayed factors introduces the separate twist sign
\((-1)^{ph}=-1\); no such reversal is used here.

The three exact groups needed below are
\begin{equation}
 \pi_6(S^3)\cong\mathbb Z/12,
 \qquad
 \pi_{12}(S^9)\cong\mathbb Z/24,
 \qquad
 \pi_{12}(S^6)\cong\mathbb Z/2.
 \label{cusp-s12-three-stem-groups}
\end{equation}
The first and third are the integral low unstable groups displayed in
Hatcher's table extracted from Toda; the middle group is in the stable
three-stem
\cite[opening table, p.~339]{Hatcher2002};
\cite[Table~1, p.~113]{IsaksenWangXu2023}.

Put \(\beta=E^6[y]\in\pi_{12}(S^9)\).  Since
\(12[y]=0\), functoriality of suspension gives \(12\beta=0\).  In the
cyclic group \(\mathbb Z/24\), the elements annihilated by \(12\) are
exactly the even elements.  Hence there is
\(\gamma\in\pi_{12}(S^9)\) such that \(\beta=2\gamma\).
Postcomposition with \(E^3[x]:S^9\to S^6\) induces a group homomorphism
\[
 (E^3[x])_*:\pi_{12}(S^9)\longrightarrow\pi_{12}(S^6).
\]
The target has exponent two, so \cref{cusp-s12-toda-smash-formula} gives
\[
 [x\wedge y]
 =(E^3[x])_*(\beta)
 =2(E^3[x])_*(\gamma)
 =0.
\]
This proves the vanishing of the entire external smash pairing and hence
the theorem.
\end{proof}

\begin{corollary}[All higher smash powers]
\label{cusp-s12-all-higher-smashes-null}
For every based \(g:S^6\to S^3\) and every \(r\geq2\),
\[
 g^{\wedge r}:S^{6r}\longrightarrow S^{3r}
\]
is based null-homotopic.
\end{corollary}

\begin{proof}
With the fixed associativity homeomorphisms for smash products,
\[
 g^{\wedge r}=g^{\wedge2}\wedge g^{\wedge(r-2)}.
\]
Smashing a based null-homotopy with a based map gives a based
null-homotopy.
\end{proof}

\begin{remark}[The one-copy invariant and its computed value]
\label{cusp-s12-one-copy-invariant}
The proof does not need to determine the one-copy class
\([g]\in\pi_6(S^3)\cong\mathbb Z/12\).  Toda's suspension calculation
identifies \(\pi_6(S^3)\) with the even subgroup of the stable three-stem;
Hatcher gives an integral Serre-spectral-sequence derivation of the same
injectivity and identifies twice the stable Hopf--\(\nu\) generator with a
generator of the suspended \(\mathbb Z/12\)-summand
\cite[Theorem~5.41(b), p.~573, and pp.~577--578]{HatcherSpectralSequences}.
Thus, after choosing a stable generator \(\nu_9\), there is a unique
\(a\in\mathbb Z/12\) such that
\begin{equation}
 [E^6g]=2a\nu_9\in\pi_{12}(S^9)\cong\mathbb Z/24.
 \label{cusp-s12-sixfold-suspension-class}
\end{equation}
Because \(a\mapsto2a\) embeds \(\mathbb Z/12\) as the even subgroup of
\(\mathbb Z/24\), \(E^6g\) is null-homotopic if and only if \(g\) is
null-homotopic.  Thus the one-copy class and its sixfold suspension are the
same integer modulo \(12\). It is now zero by Theorem~\ref{cusp-null:thm}.

The exact geometric invariant which computes \(a\) is also identifiable.
Choose \(0<r_0<\epsilon\) and
\(u_0\in\overline T^2\setminus K^{(1)}\), and let
\(y_0=[r_0,c(u_0)]\) lie in the open three-cell of
\(\Sigma_rS^2\).  From the literal formula for \(q_2\),
\begin{equation}
 g^{-1}(y_0)
 =\{x\in f^{-1}(|t|=r_0):\overline\pi_f(x)=u_0\}.
 \label{cusp-s12-regular-inverse-image}
\end{equation}
The fibre of \(\overline\pi\) is the two-torus defined by
\(K_\Lambda=\mathbb Z\langle\widehat w,\widehat\delta\rangle\), and the
cusp monodromy fixes both \(\widehat w\) and \(\widehat\delta\).  Therefore
\eqref{cusp-s12-regular-inverse-image} is the mapping torus of an affine translation of the kernel torus,
which is diffeomorphic to \(S^1\times T^2\cong T^3\): if the translation is
\(k\mapsto k+v\), choose a real lift \(\widetilde v\); the map
\([\theta,k]\mapsto(e^{2\pi i\theta},k+\theta\widetilde v)\) respects
\((1,k)\sim(0,k+v)\) and has the inverse obtained by subtracting
\(\theta\widetilde v\). After choosing a smooth chart in the open
target cell, the differential of \(g\) supplies its normal three-framing.
Pontryagin--Thom identifies the stabilized class of \(g\) with the framed
bordism class of this specifically framed \(T^3\).  Since stabilization on
\(\pi_6(S^3)\) is injective, computing that framing class computes \(a\)
exactly.  The unframed diffeomorphism type \(T^3\), point-set surjectivity,
and the zero ordinary homology map do not compute \(a\).
\end{remark}

\begin{remark}[Scope]
\label{cusp-s12-nullity-scope}
The theorem changes no point-set or metric statement.  A surjective quotient
map may be null-homotopic, and the embedded moving four-spheres and their
intrinsic pointed \(\mathbb R^4\)-limits remain exactly as constructed.  The
theorem concerns the sphere-valued postcomposition
\(\widehat q_3^{\wedge2}\); it does not imply null-homotopy of the earlier
conical-target map
\(q_2^{\wedge2}:S^{12}\to\mathcal A^{\wedge2}\).
\end{remark}

\subsection{The original cusp coefficient: an explicit framed nullbordism}


Let \(X\) be the smooth compact source of the original construction,
\(f:X\to B\) its map, and \(t\) the specified cusp coordinate.
Put \(W=f^{-1}(0)\), choose the original \(x_0\in W\), and retain the
ordered quotient lattice basis induced by
\((\widehat\gamma,\widehat u)\), with kernel generated by
\((\widehat w,\widehat\delta)\). Thus the quotient torus is
\(\overline T^2=\mathbb R^2/\mathbb Z^2\). The original periods are
\begin{equation}
 Z(s)=sB_0+C(t),\qquad e^{2\pi i s}=t,\qquad
 B_0=\begin{pmatrix}0&1\\-1&0\end{pmatrix},\qquad
 C(t)=\begin{pmatrix}6\mu(t)&h(t)\\b(t)-h(t)&\mu(t)\end{pmatrix}.
 \label{cusp-null:eq:periods}
\end{equation}
Every entry of \(C\) is holomorphic at \(t=0\).
The dense torus of the original toric cusp has coordinates
\((x_1,x_2,t)\), with
\(x=\exp(2\pi i\zeta)\). The residual lattice action is exactly
\begin{equation}
 \Psi_\lambda(x,t)=
 \bigl(t^{B_0\lambda}\exp(2\pi iC(t)\lambda)x,t\bigr),
 \qquad\lambda\in\mathbb Z^2.
 \label{cusp-null:eq:action}
\end{equation}
It extends freely and properly discontinuously across the original
unimodular toric filling. In particular the quotient projection is a local
biholomorphism at every toric fixed point. These are the reconstructed
data and proved local properties of the original source
(the reconstructed original manuscript, \S4, source pp.~24--29; its complete cusp data are displayed here).

Let
\[
 K^{(1)}=(\mathbb R/\mathbb Z\times\{0\})\cup(\{0\}\times\mathbb R/\mathbb Z),\qquad
 c:\overline T^2\longrightarrow \overline T^2/K^{(1)}\cong S^2
\]
be the specified degree-one cell quotient. Its restriction to the open
two-cell is taken with the smooth coordinates used in the original
regular-fibre description. Both radial endpoints and the complete
basepoint cylinder are collapsed in
\[
 \Sigma_r(S^2,*)=
 \frac{[0,\epsilon]\times S^2}
 {(\{0,\epsilon\}\times S^2)\cup([0,\epsilon]\times\{*\})}
 \cong S^3.
\]
The original based map is
\begin{equation}
 g(x)=[|t|,c(\overline\pi_f(x))]\quad(0<|t|<\epsilon),\qquad
 g(W)=g(X\setminus\operatorname{int}N_0)=*.
 \label{cusp-null:eq:g}
\end{equation}

\begin{theorem}[The original one-copy cusp map is null-homotopic]
\label{cusp-null:thm}
The based map \(g:(X,x_0)\to(S^3,*)\) in \eqref{cusp-null:eq:g} is based
null-homotopic. More precisely, a regular inverse image with its actual
ordered normal framing induced by \((r,c_1,c_2)\) is framed nullbordant
inside one original cusp chart times \([0,1]\), disjoint from
\(\{x_0\}\times[0,1]\).
\end{theorem}

\begin{proof}
The proof constructs the complete inverse image, compares its actual
framing with an explicitly extendible framing, and gives the extension.

\paragraph{1. The exact complete inverse image.}
Retain the lift
\[
 u^0=\begin{pmatrix}-1/4\\1/4\end{pmatrix},\qquad
 \overline u^0=(3/4,1/4)\pmod{\mathbb Z^2},\qquad
 B_0u^0=\begin{pmatrix}1/4\\1/4\end{pmatrix}.
\]
Thus \(\overline u^0\notin K^{(1)}\), and \(B_0u^0\) lies in the
interior of the original triangle \(\operatorname{conv}(0,e_1,e_2)\).
Use its original unimodular chart
\begin{equation}
 z_0=t/(x_1x_2),\qquad z_1=x_1,\qquad z_2=x_2,
 \qquad t=z_0z_1z_2.
 \label{cusp-null:eq:chart}
\end{equation}
Set
\[
 \ell=\log|t|<0,\qquad R(t)=-2\pi\operatorname{Im}C(t),\qquad
 A(t)=\ell B_0+R(t).
\]
For a representative \(\zeta=Z(s)u+k\), where \(u,k\in\mathbb R^2\),
\(\operatorname{Im}s=-\ell/(2\pi)\) gives the exact identity
\begin{equation}
 \begin{pmatrix}\log|x_1|\\\log|x_2|\end{pmatrix}
 =A(t)u.
 \label{cusp-null:eq:moduli}
\end{equation}
The quotient value is \(u\bmod\mathbb Z^2\).
Since \(B_0\) is orthogonal, for all \(v\in\mathbb R^2\),
\begin{equation}
 \|A(t)v\|\geq (|\ell|-\|R(t)\|)\|v\|.
 \label{cusp-null:eq:invertibility}
\end{equation}
The boundedness of \(R\) therefore makes \(A(t)\) invertible on a
sufficiently small punctured disc.

For any point of the fibre quotient over \(\overline u^0\), subtracting
the appropriate original period, equivalently applying
\(\Psi_{-\lambda}\), gives a representative with the exact lift
\(u=u^0\). This representative is unique: another residual translate
has logarithmic moduli differing by \(A(t)\lambda\), which vanishes
only for \(\lambda=0\). Conversely every dense-torus point satisfying
\(\log|x|=A(t)u^0\) has this quotient value, since the remaining real
parts of \(\zeta\) are exactly the kernel coordinates \(k\).
For each fixed \(t\), varying \(k\bmod\mathbb Z^2\) gives all independent
phases of \(x_1,x_2\).

The following change of coordinates uses the full original period units:
\begin{align}
 h_j^0(t)&=\exp\bigl(-2\pi i(C(t)u^0)_j\bigr),\qquad j=1,2,\\
 \mathcal F(z_0,z_1,z_2)
 &=(w_0,w_1,w_2)
 =\left(\frac{z_0}{h_1^0(t)h_2^0(t)},\,
 h_1^0(t)z_1,\,h_2^0(t)z_2\right).
 \label{cusp-null:eq:straightening}
\end{align}
The functions \(h_j^0\) are nowhere-zero holomorphic units, including
at \(t=0\). The product \(w_0w_1w_2\) is exactly \(t\).
The inverse of \(\mathcal F\) is
\[
 (w_0,w_1,w_2)\longmapsto
 \left(w_0h_1^0(t)h_2^0(t),\,
 \frac{w_1}{h_1^0(t)},\,\frac{w_2}{h_2^0(t)}\right),
 \qquad t=w_0w_1w_2.
\]
Thus this is a biholomorphism of the original chart restricted to the
cusp disc, not a replacement of \(C(t)\).

For \(r_0\in(0,\epsilon)\), the entire inverse image of
\(y_0=[r_0,c(\overline u^0)]\) becomes exactly
\begin{equation}
 T_{r_0}=
 \{|w_0|=\rho_0,\ |w_1|=\rho_1,\ |w_2|=\rho_2\},\qquad
 (\rho_0,\rho_1,\rho_2)=
 (r_0^{1/2},r_0^{1/4},r_0^{1/4}).
 \label{cusp-null:eq:torus}
\end{equation}
Indeed, \(\log|h^0(t)|=2\pi\operatorname{Im}C(t)u^0=-R(t)u^0\),
so \eqref{cusp-null:eq:moduli} gives \(\log|w_j|=\ell(B_0u^0)_j\) for
\(j=1,2\); the product identity gives the exponent \(1/2\) for
\(w_0\). Conversely these moduli recover \(u=u^0\) by invertibility
of \(A(t)\). The angular map is explicitly
\[
 (\arg t,\arg w_1,\arg w_2)\longleftrightarrow
 (\arg w_0,\arg w_1,\arg w_2),\qquad
 \arg w_0=\arg t-\arg w_1-\arg w_2\pmod{2\pi},
\]
with inverse \(\arg t=\arg w_0+\arg w_1+\arg w_2\).
Its integral matrix has determinant one. This includes the complete
angular mapping torus and all kernel phases, with their original signs.

Let \(o\) be the origin of \eqref{cusp-null:eq:chart}. The original free properly
discontinuous action supplies a neighborhood \(V\) of \(o\) on which
the quotient projection is injective. To see the local injectivity
directly, first choose a relatively compact neighborhood; proper
discontinuity leaves only finitely many translates that can intersect it.
None of those nonidentity translates fixes \(o\). Shrinking finitely
many times separates the neighborhood from all of them. The quotient
map is a local biholomorphism, so its restriction to \(V\) is a chart
of the original source.
Choose a polydisc about \(0\) whose closure lies in \(\mathcal F(V)\).
For sufficiently small \(r_0\), both \(T_{r_0}\) and
\begin{equation}
 \{|w_0|\leq\rho_0,\ |w_1|=\rho_1,\ |w_2|=\rho_2\}
 \label{cusp-null:eq:spatial-bound}
\end{equation}
lie in this polydisc. This proves that the whole inverse image embeds
in one chart and supplies space for the bounding manifold below.

\paragraph{2. The actual normal framing, including its angle dependence.}
Put \(L_j=\log|w_j|\), and use the unique real lift of the quotient
coordinate near \(u^0\). From the exact change of coordinates,
\begin{equation}
 \begin{pmatrix}L_1\\L_2\end{pmatrix}
 =\ell B_0u^0+A(t)(u-u^0),\qquad
 L_0=\ell-L_1-L_2.
 \label{cusp-null:eq:nearby-moduli}
\end{equation}
Differentiating on \(u=u^0\) annihilates the term
\((dA(t))(u-u^0)\). Hence no derivative of \(C(t)\) has been omitted:
the exact differential there is
\begin{equation}
 \begin{pmatrix}dL_0\\dL_1\\dL_2\end{pmatrix}
 =J(t)\begin{pmatrix}d\ell\\du_1\\du_2\end{pmatrix},\quad
 J(t)=
 \begin{pmatrix}
 1/2&-A_{11}-A_{21}&-A_{12}-A_{22}\\
 1/4&A_{11}&A_{12}\\
 1/4&A_{21}&A_{22}
 \end{pmatrix}.
 \label{cusp-null:eq:J}
\end{equation}
Adding the second and third rows to the first gives \((1,0,0)\),
and expansion in that row proves \(\det J(t)=\det A(t)\).

Let \(D=Dc_{\overline u^0}\) be the actual two-by-two derivative in
the ordered open-cell target coordinates \((c_1,c_2)\). It is a fixed
invertible matrix; all its entries and signs are retained. Let
\[
 H=\begin{pmatrix}r_0&0&0\\0&D_{11}&D_{12}\\0&D_{21}&D_{22}\end{pmatrix}.
\]
Let \(n_j\) be the real radial vector \(w_j\) in the \(j\)-th complex
coordinate, with other components zero. Then \(dL_i(n_j)=\delta_{ij}\).
Writing \(N=(n_0,n_1,n_2)\) as the matrix with these vector columns, the actual
\(dg\)-dual normal framing is, in the normal quotient bundle,
\begin{equation}
 \mathfrak n(t)=N J(t)H^{-1}.
 \label{cusp-null:eq:actual-frame}
\end{equation}
Indeed, \(dg(N)=HJ(t)^{-1}\), so multiplication by \(J(t)H^{-1}\)
gives the identity. In the Euclidean metric of this chart the radial
vectors are perpendicular to the product torus; a different source
metric changes their representatives by the canonical projection to its
normal bundle and preserves this normal-quotient framing.

Put \(\ell_0=\log r_0\). Decrease \(r_0\) further, if necessary, so
that \(|\ell_0|>2\sup_{|t|\leq r_0}\|R(t)\|\). For
\(0\leq\sigma\leq1\), define only a homotopy of frames by
\[
 A_\sigma(t)=\ell_0B_0+\sigma R(t),
\]
and let \(J_\sigma(t)\) be \eqref{cusp-null:eq:J} with \(A_\sigma\) in place
of \(A\). For every \(v\in\mathbb R^2\),
\[
 \|A_\sigma(t)v\|\geq
 (|\ell_0|-\|R(t)\|)\|v\|>0\quad(v\ne0).
\]
Thus all matrices are invertible. At \(\sigma=0\),
\begin{equation}
 J_0=
 \begin{pmatrix}
 1/2&\ell_0&-\ell_0\\
 1/4&0&\ell_0\\
 1/4&-\ell_0&0
 \end{pmatrix},\qquad \det J_0=\ell_0^2>0.
 \label{cusp-null:eq:J0}
\end{equation}
Continuity also proves \(\det J_\sigma(t)>0\) for every \(\sigma,t\).
The homotopy \(N J_\sigma(t)H^{-1}\), run from \(\sigma=1\) to
\(\sigma=0\), compares the actual frame with the exact constant-matrix
frame \(N P\), where
\begin{equation}
 P=J_0H^{-1}.
 \label{cusp-null:eq:P}
\end{equation}
The original embedded torus and original functions \(C(t)\) have not
changed. A framing homotopy is a framed cylinder: choose a smooth
parameter function constant near its endpoints and use the displayed
frames on \(T_{r_0}\times[0,1]\).

\paragraph{3. A collared framed nullbordism with the exact constant matrix.}
Use a new time coordinate \(v\in[0,1]\). Define
\begin{equation}
 q(v)=
 \begin{cases}
 0,&0\leq v\leq1/4,\\
 \rho_0^2\exp\bigl(4-1/(v-1/4)\bigr),&1/4<v\leq1.
 \end{cases}
 \label{cusp-null:eq:q}
\end{equation}
This function is smooth: all right derivatives of
\(\exp(-1/(v-1/4))\) at \(1/4\) vanish. It is strictly increasing
for \(v>1/4\), satisfies \(q(1/2)=\rho_0^2\), and
\(q'(1/2)=16\rho_0^2>0\).
In the chart times \([0,1]\), put
\begin{equation}
 \mathcal W=
 \{|w_0|^2+q(v)=\rho_0^2,\quad
       |w_1|^2=\rho_1^2,\quad|w_2|^2=\rho_2^2\}.
 \label{cusp-null:eq:bordism}
\end{equation}
No point has \(v>1/2\). The three defining differentials are independent:
the latter two have nonzero gradients in their separate complex factors;
the first has nonzero \(w_0\)-gradient whenever \(w_0\ne0\), and at
\(w_0=0\) it has the nonzero \(v\)-derivative \(q'(1/2)\).
Thus \(\mathcal W\) is a smooth compact four-manifold. Its only
boundary is \(T_{r_0}\times\{0\}\), and on \(0\leq v\leq1/4\)
it is the exact product collar \(T_{r_0}\times[0,1/4]\).
The north locus \(w_0=0,v=1/2\) is an interior cap: the implicit
function theorem expresses \(v\) smoothly in terms of
\((\operatorname{Re}w_0,\operatorname{Im}w_0)\) there. The first
factor is a disk obtained by rotating its radial profile with this smooth
cap and one boundary circle. Therefore
\(\mathcal W\cong D^2\times T^2\).
Its spatial coordinates lie in \eqref{cusp-null:eq:spatial-bound}, so the complete
manifold lies in the injective original chart.

In the Euclidean chart product metric define
\[
 V_0=(w_0,0,0,q'(v)/2),\qquad
 V_1=(0,w_1,0,0),\qquad V_2=(0,0,w_2,0).
\]
These are one half of the gradients of the three defining equations
in \eqref{cusp-null:eq:bordism}. They are mutually orthogonal and never zero.
For example, if \(F_0=|w_0|^2+q(v)-\rho_0^2\), then
\[
 dF_0(V_0)=2|w_0|^2+\frac12q'(v)^2>0
\]
at every point of \(\mathcal W\). On the entire collar \(q'=0\), so
\((V_0,V_1,V_2)=(N,0)\) exactly. Consequently
\begin{equation}
 (V_0,V_1,V_2)P
 \label{cusp-null:eq:extended-frame}
\end{equation}
is a framing of the normal three-plane bundle of \(\mathcal W\),
and its restriction to the collar is exactly the framing \(NP\) from
\eqref{cusp-null:eq:P}, with every original target-coordinate sign retained.
Concatenate the framing cylinder of the preceding paragraph with this
nullbordism, rescaling the two interval lengths. Their framings agree
exactly on the common product collar, so the concatenation is a smooth
framed nullbordism of the actual inverse image.

\paragraph{4. The based map represented by this framed inverse image.}
Choose the injective chart to exclude \(x_0\) if its center is not
\(x_0\). If its center is \(x_0\), the constructed nullbordism still
excludes \(x_0\), since \(|w_1|=\rho_1>0\) and
\(|w_2|=\rho_2>0\) everywhere on it. Compactness therefore gives
tubular neighborhoods for the framed cylinder and cap disjoint from
\(\{x_0\}\times[0,1]\).

The original continuous quotient map is smooth and a submersion near
its entire inverse image of \(y_0\). First make it constant in a
neighborhood of \(x_0\), by a based homotopy that changes nothing near
that inverse image. Explicitly, take a target chart
\(\kappa:O\to\mathbb R^3\) with \(\kappa(*)=0\), whose image is a ball,
and \(y_0\notin O\). By continuity choose a neighborhood \(U\) of
\(x_0\) with \(g(U)\subset O\), and a smooth bump function
\(\beta\) supported in \(U\) and equal to one near \(x_0\).
For \(x\in U\), replace \(g(x)\) during this homotopy by
\(\kappa^{-1}((1-a\beta(x))\kappa(g(x)))\), \(0\leq a\leq1\),
and leave \(g\) fixed off \(U\). Convexity of the target chart image
makes the formula well-defined, the vanishing of \(\beta\) near the
boundary of its support proves continuity of the pasted formula, and
\(\kappa(g(x_0))=0\) makes the homotopy based. It never introduces
an inverse image of \(y_0\). Its endpoint is smooth near \(x_0\) and
agrees with \(g\) near the entire old inverse image.

Denote the endpoint of the preceding homotopy by \(g_b\). Choose one
closed tubular neighborhood \(U_1\) of the inverse image, contained in
the smooth submersion region and with the entire inverse image contained
in \(\operatorname{int}U_1\). The compact set
\(X\setminus\operatorname{int}U_1\) has image under \(g_b\) missing
\(y_0\). In a fixed target metric put
\[
 d_0=\min_{x\in X\setminus\operatorname{int}U_1}
       d(g_b(x),y_0)>0.
\]
Relative smooth approximation, fixed on the whole closed set \(U_1\)
and on a neighborhood of \(x_0\), gives a smooth based representative
\(g_{\mathrm{sm}}\) whose error on
\(X\setminus\operatorname{int}U_1\) is less than \(d_0/2\).
It therefore has the same inverse image and differential on \(U_1\),
and \(d(g_{\mathrm{sm}}(x),y_0)>d_0/2\) on its complement. Thus there
is no uncontrolled intervening region and no additional inverse image.
Make the approximation also closer than a fixed target injectivity
radius; its short geodesics give the based homotopy.

The relative Pontryagin construction now applies to the explicit
framed nullbordism. It produces a based homotopy to a map missing
\(y_0\); stereographic contraction of \(S^3\setminus\{y_0\}\cong\mathbb R^3\)
to the image of \(*\) is based and contracts that last map.
For completeness, the dependence on the framed inverse image in this
argument is precisely the local product-neighborhood construction and
the interpolation lemma of \cite[\S7, Theorem B and Lemma 4,
pp.~43, 48--50]{CuspNullMilnor1965}. Their relative versions fix the source
basepoint because all tubular modifications are away from it and
stereographic interpolation between equal values keeps that value
fixed. A framed cylinder realizes a framing homotopy; the empty
Pontryagin manifold realizes a map missing the chosen regular value.
This proves the claimed based null-homotopy of the exact map \(g\).
\end{proof}

\begin{corollary}
Under the original fixed based identification \(X\cong S^6\), the
one-copy coefficient is
\[
 [g]=0\in\pi_6(S^3)\cong\mathbb Z/12,\qquad a=0\pmod{12},\qquad
 [E^6g]=0\in\pi_{12}(S^9)\cong\mathbb Z/24.
\]
The previously defined homomorphism
\(\Theta_g:\mathbb Z/12\to\pi_6(S^3)\), \(\Theta_g([n])=n[g]\), and
its composite with the marked peripheral character are zero.
\end{corollary}
\begin{proof}
The theorem states that the actual based map is null-homotopic.
Precomposition with the fixed based identification preserves a based
null-homotopy. Suspension preserves it as well. Thus the unique
coefficient of \([g]\) in any generator of \(\mathbb Z/12\) is zero.
For every \([n]\), \(\Theta_g([n])=n\cdot0=0\), and every
composite with this homomorphism is zero.
\end{proof}

\paragraph{Exact scope.}
The proof computes this sphere-valued postcomposition. It does not
compute the homotopy class of the earlier conical-target map
\(q_2:X\to\operatorname{Susp}(S^1\times\overline T^2)\), and it
does not assert a null-homotopy for that different target. Point-set
surjectivity of \(g\) and all of its stated metric restrictions remain
properties of the original map: a homotopy-class computation does not
alter the defining formula \eqref{cusp-null:eq:g}.
The retained-angle companion is proved separately in
Theorem~\ref{h24:ret:main}: $p_3q_2:X\to S^4$ represents
$\eta_4\eta_5\ne0$, so $q_2$ is nonnull. Its Hopf smash
represents $12\nu_6\ne0$ in $\pi_9(S^6)$. These positive
computations use the actual four-framed inverse torus of that different
postcomposition; they do not reverse or contradict the preceding
null-homotopy through $\widehat p_2$.


\begin{thebibliography}{9}
\bibitem{CuspNullMilnor1965}
J.~W. Milnor, \emph{Topology from the Differentiable Viewpoint},
University Press of Virginia, 1965, \S7, Theorem B and Lemma 4,
pp.~42--50. \url{https://www.maths.ed.ac.uk/~aar/papers/milnortop.pdf}.
\end{thebibliography}


\subsection{Fixed-target metric realization and its sharp dimensional scope}
\label{sphere-metric-limit-note}

This note concerns the codomains of the three based maps
\[
 q_2^{\wedge 2}:X^{\wedge 2}\cong S^{12}
   \longrightarrow Z_{12}:=\mathcal A^{\wedge 2},
\]
\[
 q_{1,\mathrm{cusp}}^{\wedge 4}:X^{\wedge 4}\cong S^{24}
   \longrightarrow Z_{24}^{\mathcal B}:=\mathcal B^{\wedge 4},
 \qquad
 q_1^{\wedge 4}:X^{\wedge 4}\cong S^{24}
   \longrightarrow Z_{24}^{\mathcal C}:=\mathcal C^{\wedge 4}.
\]
In particular, \(S^{12}\) and \(S^{24}\) are the sources, not the spaces
in which the metric limits occur.  We construct one fixed compatible
length metric on each codomain which realizes the native
four-dimensional flat limits on the target slices.  The realization is
sharp: after the same rescaling, the full ambient limits have dimensions
\(8\) and \(12\), and the collapsed smash base point goes to infinity.

Put
\[
 r_T=e^{-2\pi T},\qquad
 J=\begin{pmatrix}0&1\\-1&0\end{pmatrix},\qquad
 Y_T=TJ+E_T,
\]
where \(E_T\) is the bounded matrix in the radial cusp expansion.  Choose
and fix \(T_0\) and \(C>0\) so that
\(\|E_T\|_{\mathrm{op}}\leq C\) for every \(T\geq T_0\), in the fixed
Euclidean coordinates used below.  This notation avoids confusing the
error matrix with the quotient circle \(C_T\) used below.  On the
fixed torus
\[
 \overline T^2=(\overline\Lambda\otimes\mathbb R)/\overline\Lambda
              \cong\mathbb R^2/\mathbb Z^2
\]
let
\[
 h_T:=\iota_T^*g_{Q_T}=dx^{\mathsf T}Y_T^{\mathsf T}Y_Tdx.
\]
Thus
\begin{equation}
 T^{-2}h_T\longrightarrow h_\infty:=dx^{\mathsf T}dx
 \label{sphere-metric-normalized-Q-limit}
\end{equation}
in every \(C^k\) norm in the fixed flat coordinates.  Indeed, the cusp
coefficients are holomorphic in \(t=e^{-2\pi T}\), so their \(T\)-derivatives
are exponentially small, while \(T^{-1}Y_T=J+T^{-1}E_T\).

Use normalized angular coordinates
\[
 u,\theta\in\mathbb R/\mathbb Z;
\]
\(u\) is the argument of the base coordinate and \(\theta\) is the value of
the invariant character.  Increase this same \(T_0\) once more, if
necessary, so that all radial cusp formulas used below hold for
\(T\geq T_0\) and \(e^{-2\pi T_0}<\epsilon\), and put
\[
 r_0=e^{-2\pi T_0}.
\]
For \(0<r\leq r_0\), write
\[
 T(r):=-\frac{1}{2\pi}\log r.
\]

\paragraph{Target-slice notation.}
We use the embeddings
\(j_T,k_T^{\mathcal B},k_T^{\mathcal C}\) defined in
\cref{cusp-cone-target-slice-embeddings}; the normalized angle \(u=0\) used
below is the same positive radial base angle denoted by \(1\in S^1\) there.

\begin{theorem}[Topology-forced collapse for every fixed compatible metric]
\label{sphere-metric-uniform-collapse}
Let \(D_{12},D_{24}^{\mathcal B},D_{24}^{\mathcal C}\) be arbitrary fixed
metrics compatible with the respective compact target topologies, and let
\(*\) denote the corresponding smash base point.  Then
\begin{align}
 d_H^{D_{12}}\bigl(j_T(Q_T^2),\{*\}\bigr)&\longrightarrow0,\\
 d_H^{D_{24}^{\mathcal B}}
  \bigl(k_T^{\mathcal B}(C_T^4),\{*\}\bigr)&\longrightarrow0,\\
 d_H^{D_{24}^{\mathcal C}}
  \bigl(k_T^{\mathcal C}(C_T^4),\{*\}\bigr)&\longrightarrow0.
 \label{sphere-metric-Hausdorff-collapse}
\end{align}
Consequently, with the metrics restricted from any one fixed ambient
target, all three target slices converge in Gromov--Hausdorff topology to
a point, not to \(\mathbb R^4\).
\end{theorem}

\begin{proof}
We prove the first assertion; the other two are identical.  Let
\[
 G:[0,r_0]^2\times(T^3)^2\longrightarrow Z_{12}
\]
be the product of the two inner-cone quotient maps followed by the smash
quotient.  It is continuous and
\(G(0,0,z_1,z_2)=*\) for every \((z_1,z_2)\in(T^3)^2\).  Since the domain
is compact, uniform continuity with respect to \(D_{12}\) gives
\[
 \sup_{(z_1,z_2)\in(T^3)^2}
 D_{12}\bigl(G(r,r,z_1,z_2),*\bigr)\longrightarrow0.
\]
The set \(j_T(Q_T^2)\) is contained in such a cross-section, which proves
the first Hausdorff convergence.  Replacing \(T^3\) by \(T^2\) and taking
four factors proves the other two.  Hausdorff convergence to a singleton
implies Gromov--Hausdorff convergence to that singleton.
\end{proof}

The preceding theorem is unrescaled.  It does not prevent a blow-up of one
fixed metric.  The following construction gives the strongest natural
positive statement.

\begin{theorem}[One fixed cusp-adapted metric realizes the vertical
\(\mathbb R^4\) limits]
\label{sphere-metric-fixed-realization}
There are fixed compatible length metrics
\[
 D_{12}\text{ on }Z_{12},\qquad
 D_{24}^{\mathcal B}\text{ on }Z_{24}^{\mathcal B},\qquad
 D_{24}^{\mathcal C}\text{ on }Z_{24}^{\mathcal C}
\]
with the following properties.

On the inner cone of \(\mathcal A\), the metric is induced by
\begin{equation}
 g_{\mathcal A}
 =dr^2+r^2du^2+\frac{r^2}{T(r)^2}h_{T(r)}.
 \label{sphere-metric-A-fixed-metric}
\end{equation}
On the cusp cone of either \(\mathcal B\) or \(\mathcal C\), it is induced
by
\begin{equation}
 g_{\mathcal D}=dr^2+r^2(du^2+d\theta^2),
 \qquad\mathcal D\in\{\mathcal B,\mathcal C\}.
 \label{sphere-metric-BC-fixed-metric}
\end{equation}
The metrics on the smash powers are the metric quotients of the product
length metrics after the closed fat wedge is collapsed.  Write
\(g_{Z_{12}}\) and \(g_{Z_{24}^{\mathcal D}}\) for the resulting smooth
Riemannian tensors on the open regular strata.

Let
\[
 p_T^{12}:=j_T(0,0),\qquad
 p_T^{24,\mathcal D}:=k_T^{\mathcal D}(0,0,0,0).
\]
Put
\[
 \Sigma_T^{12}:=j_T(Q_T^2),\qquad
 \Sigma_T^{24,\mathcal D}:=k_T^{\mathcal D}(C_T^4).
\]
Let \(d_{T,12}^{\mathrm{slice}}\) and
\(d_{T,24,\mathcal D}^{\mathrm{slice}}\) be the intrinsic Riemannian path
metrics of the first fundamental forms on these embedded slices.  These
must be distinguished from the restricted ambient distance functions
\[
 d_{T,12}^{\mathrm{amb}}
 :=D_{12}|_{\Sigma_T^{12}\times\Sigma_T^{12}},\qquad
 d_{T,24,\mathcal D}^{\mathrm{amb}}
 :=D_{24}^{\mathcal D}|_{\Sigma_T^{24,\mathcal D}
                         \times\Sigma_T^{24,\mathcal D}}.
\]
In general the latter can be smaller, because an ambient path is allowed
to leave the slice.

Then the induced first fundamental forms on the target slices satisfy the
exact identities
\begin{equation}
 \left(\frac{T}{r_T}\right)^2
  j_T^*g_{Z_{12}}
 =g_{Q_T}\oplus g_{Q_T},
 \label{sphere-metric-A-slice-exact}
\end{equation}
and
\begin{equation}
 \left(\frac{\ell_T}{r_T}\right)^2
  (k_T^{\mathcal D})^*g_{Z_{24}^{\mathcal D}}
 =g_{C_T}^{\oplus4}.
 \label{sphere-metric-BC-slice-exact}
\end{equation}
Consequently,
\begin{align}
 \left(\Sigma_T^{12},\frac{T}{r_T}d_{T,12}^{\mathrm{slice}},
       p_T^{12}\right)
 &\longrightarrow(\mathbb R^4,0),\\
 \left(\Sigma_T^{24,\mathcal D},
       \frac{\ell_T}{r_T}d_{T,24,\mathcal D}^{\mathrm{slice}},
       p_T^{24,\mathcal D}\right)
 &\longrightarrow(\mathbb R^4,0)
 \label{sphere-metric-slice-fixed-blowups}
\end{align}
in pointed local-isometric and pointed Gromov--Hausdorff topology, and
smoothly on the exhausting Euclidean slice balls.

The full ambient pointed blow-ups are instead
\begin{align}
 \left(Z_{12},\frac{T}{r_T}D_{12},p_T^{12}\right)
 &\longrightarrow(\mathbb R^8,0),\\
 \left(Z_{24}^{\mathcal D},
       \frac{\ell_T}{r_T}D_{24}^{\mathcal D},
       p_T^{24,\mathcal D}\right)
 &\longrightarrow(\mathbb R^{12},0).
 \label{sphere-metric-ambient-fixed-blowups}
\end{align}
as pointed metric spaces in pointed Gromov--Hausdorff topology.  More
precisely, for every fixed radius the corresponding rescaled target ball is
eventually contained in an open regular-stratum coordinate neighborhood,
and the pulled-back Riemannian tensors converge there in every \(C^k\) norm.
Thus smooth convergence is asserted on an exhaustion by regular-stratum
neighborhoods, not globally on the singular compact targets.  Under these
limits, the slice \(\mathbb R^4\) is respectively the coordinate subspace
obtained by fixing the two radial and two base-angle coordinates, or by
fixing the four radial and four base-angle coordinates.
\end{theorem}

\begin{proof}
Equation \eqref{sphere-metric-normalized-Q-limit} shows that the metrics on
the compact cross-sections of \eqref{sphere-metric-A-fixed-metric} remain
uniformly bounded and converge to a positive definite flat metric.
Therefore the cross-section diameter is \(O(r)\).  Moreover the radial
coordinate is \(1\)-Lipschitz and a radial segment from level \(r\) to the
vertex has length \(r\).  The length completion consequently adds exactly
one point at \(r=0\), and its topology is the prescribed cone topology.
The same assertions are immediate for
\eqref{sphere-metric-BC-fixed-metric}.  Interpolation on a closed annulus
away from the vertex extends these metrics to fixed compatible length
metrics on the compact spaces \(\mathcal A,\mathcal B,\mathcal C\).

If a compact metric space \((E,d)\) and a closed subset \(W\subset E\) are
given, the metric quotient which collapses \(W\) to one point is
\[
 d_{E/W}([x],[y])
 =\min\{d(x,y),d(x,W)+d(y,W)\}.
\]
Applying this to the fat wedges in the product length spaces gives the
claimed fixed compatible metrics on the smash targets.

On the \(r_T\)-slice of \(\mathcal A\), with \(u=0\), the first fundamental
form is
\[
 \frac{r_T^2}{T^2}h_T.
\]
The definition \(h_T=\iota_T^*g_{Q_T}\) proves
\eqref{sphere-metric-A-slice-exact}.  On the \(r_T\)-slice of
\(\mathcal B\) or \(\mathcal C\), with \(u=0\), the form is
\(r_T^2d\theta^2\).  Under
\(\vartheta_T:[\theta]\mapsto[\ell_T\theta]\), the native quotient-circle
metric is \(\ell_T^2d\theta^2\).  This proves
\eqref{sphere-metric-BC-slice-exact}.  The native systole estimates
\[
 \operatorname{sys}(Q_T^2)\geq T-C,
 \qquad
 \operatorname{sys}(C_T^4)=\ell_T
\]
then prove \eqref{sphere-metric-slice-fixed-blowups}: every prescribed
radius ball is eventually exactly a Euclidean four-ball.

It remains to compute the ambient limits.  In one \(\mathcal A\)-factor,
put
\[
 \lambda_T=\frac{T}{r_T},\qquad
 s=\lambda_T(r-r_T),\qquad
 \eta=Tu,\qquad z=Y_Tx
\]
on the universal cover of the torus directions.  On every fixed bounded
set in \((s,\eta,z)\), one has
\[
 \frac r{r_T}=1+\frac sT,
 \qquad
 T(r)=T+O(T^{-1}),
 \qquad
 Y_{T(r)}Y_T^{-1}=I+O(T^{-2}),
\]
with the same estimates after any fixed number of derivatives.  Thus
\(\lambda_T^2g_{\mathcal A}\) converges in these coordinates to
\[
 ds^2+d\eta^2+dz_1^2+dz_2^2.
\]
The inward radial boundary is at distance \(T\), the \(u\)-period is \(T\),
and the \(z\)-period lattice is \(Y_T\mathbb Z^2\), whose systole is at
least \(T-C\).  The outward edge of the fixed cusp chart is also at scaled
distance
\[
 \frac{T}{r_T}(r_0-r_T)
 =T\left(\frac{r_0}{r_T}-1\right)\longrightarrow\infty.
\]
Thus the inward vertex, all periods, and the outward edge recede to
infinity.  A single factor converges on exhausting regular-stratum
coordinate neighborhoods to \(\mathbb R^4\), and two factors converge to
\(\mathbb R^8\).

For one \(\mathcal B\)- or \(\mathcal C\)-factor, put
\[
 \lambda_T=\frac{\ell_T}{r_T},\qquad
 s=\lambda_T(r-r_T),\qquad
 \eta=\ell_Tu,\qquad \zeta=\ell_T\theta.
\]
Then
\[
 \lambda_T^2g_{\mathcal D}
 =ds^2+\left(\frac r{r_T}\right)^2
       (d\eta^2+d\zeta^2),
\]
which converges smoothly on bounded coordinate sets to the Euclidean metric
on \(\mathbb R^3\).  The two circle periods and the inward radial boundary
are all at scaled distance \(\ell_T\to\infty\), while the outward edge is
at scaled distance
\[
 \frac{\ell_T}{r_T}(r_0-r_T)
 =\ell_T\left(\frac{r_0}{r_T}-1\right)\longrightarrow\infty.
\]
Four factors therefore converge on exhausting regular-stratum coordinate
neighborhoods to \(\mathbb R^{12}\).

Finally, a level-\(r_T\) factor point has distance exactly \(r_T\) from its
cone vertex.  Indeed, the radial segment has that length.  A competing
curve contained in the inner chart has length at least \(r_T\), while one
which leaves has inner-chart portions of total length at least
\((r_0-r_T)+r_0>r_T\).  The Riemannian-\(\ell^2\)-product distance from a
synchronized point to the fat wedge is the minimum of the factor
distances, and hence is also \(r_T\).

For the chosen scale put
\[
 b_T=\lambda_Tr_T,\qquad a_T=\lambda_T(r_0-r_T),
\]
and let \(\sigma_T\) be the least scaled period in a factor.  Thus
\(\sigma_T\geq T-C\) in the \(\mathcal A\) calculation and
\(\sigma_T=\ell_T\) in the \(\mathcal B,\mathcal C\) calculation.  The
preceding computations give
\[
 a_T\longrightarrow\infty,\qquad
 b_T\longrightarrow\infty,\qquad
 \sigma_T\longrightarrow\infty.
\]
Fix \(R\).  A curve of rescaled length at most \(R\) from \(p_T\) meets
neither the fat wedge nor \(r=r_0\) for all large \(T\), so the
radius-\(R\) target ball lies in the open regular stratum.  If \(x,y\) lie
in that ball, a competitor through the fat wedge has length at least
\[
 (b_T-R)+(b_T-R)>2R,
\]
and a competitor which exits the inner chart and returns has length at
least
\[
 2(a_T-R)>2R.
\]
Every remaining competitor of length at most \(2R\) stays in
\(r/r_T=1+o_R(1)\).  After lifting it to the angular and torus covers, a
nonzero deck translation has length at least
\((1-o_R(1))(\sigma_T-2R)>2R\) for all large \(T\).  Since the path through
\(p_T\) has length at most \(2R\), none of these three kinds of global
shortcut computes the distance between \(x\) and \(y\).  Thus the metric
on every fixed target ball is computed entirely in one lifted
regular-stratum chart.  The \(C^\infty\) convergence of the pulled-back
tensors on these exhausting charts implies the pointed
Gromov--Hausdorff limits in
\eqref{sphere-metric-ambient-fixed-blowups}.
\end{proof}

\begin{theorem}[Complete intrinsic scalar-rescaling trichotomy]
\label{sphere-metric-scaling-trichotomy}
For the fixed metrics of
\cref{sphere-metric-fixed-realization}, let \(\lambda_T>0\) be arbitrary
and put
\[
 \beta_T:=\lambda_Tr_T.
\]
Choose one of the following triples, retaining the indicated ambient
metric:
\[
 \begin{split}
  (\Sigma_T,d_T^{\mathrm{slice}},d_T^{\mathrm{amb}},p_T;D)
  &=(\Sigma_T^{12},d_{T,12}^{\mathrm{slice}},
      d_{T,12}^{\mathrm{amb}},p_T^{12};D_{12}),\\
  (\Sigma_T,d_T^{\mathrm{slice}},d_T^{\mathrm{amb}},p_T;D)
  &=(\Sigma_T^{24,\mathcal D},
      d_{T,24,\mathcal D}^{\mathrm{slice}},
      d_{T,24,\mathcal D}^{\mathrm{amb}},
      p_T^{24,\mathcal D};D_{24}^{\mathcal D}),
      \qquad \mathcal D\in\{\mathcal B,\mathcal C\}.
 \end{split}
\]
After passing to a subsequence so that \(\beta_T\) has a limit in
\([0,\infty]\), the intrinsic pointed slice limit is
\begin{equation}
 (\Sigma_T,\lambda_Td_T^{\mathrm{slice}},p_T)
 \longrightarrow
 \begin{cases}
  \{\mathrm{pt}\},&\beta_T\longrightarrow0,\\
  ((\mathbb R/L\mathbb Z)^4,0),
      &\beta_T\longrightarrow L\in(0,\infty),\\
  (\mathbb R^4,0),&\beta_T\longrightarrow\infty.
 \end{cases}
 \label{sphere-metric-scaling-trichotomy-result}
\end{equation}
The convergence in
\eqref{sphere-metric-scaling-trichotomy-result} is pointed
Gromov--Hausdorff in all three cases.  In the finite positive case the flat
metrics converge smoothly under the fixed marked-torus identifications; in
the unrolling case the convergence is pointed local-isometric and smooth on
the exhausting Euclidean balls.
In the middle line the lattice is orthogonally equivalent to the square
lattice \(L\mathbb Z^4\).  At the same time,
\begin{equation}
 \lambda_TD(p_T,*)=\beta_T
 \label{sphere-metric-vertex-distance}
\end{equation}
for every one of the displayed choices.

Only in the unrolling regime \(\beta_T\to\infty\) do we compare the
intrinsic metric with the restricted ambient distance.  In that regime, for
every \(R\in(0,\infty)\),
\begin{equation}
 \sup_{x,y\in
  \overline B_{\lambda_Td_T^{\mathrm{slice}}}(p_T,R)}
 \left|
  \lambda_Td_T^{\mathrm{amb}}(x,y)
  -\lambda_Td_T^{\mathrm{slice}}(x,y)
 \right|\longrightarrow0.
 \label{sphere-metric-intrinsic-ambient-agreement}
\end{equation}
Moreover, there are \(\varepsilon_{T,R}\to0\) such that
\begin{equation}
 \overline B_{\lambda_Td_T^{\mathrm{slice}}}(p_T,R)
 \subset
 \overline B_{\lambda_Td_T^{\mathrm{amb}}}(p_T,R)
 \subset
 \overline B_{\lambda_Td_T^{\mathrm{slice}}}
             (p_T,(1+\varepsilon_{T,R})R).
 \label{sphere-metric-intrinsic-ambient-ball-comparison}
\end{equation}
Thus the restricted ambient metric has the same pointed
\(\mathbb R^4\) limit on these fixed slice balls.  No agreement is asserted
in the regimes \(\beta_T\to0\) or \(\beta_T\to L\in(0,\infty)\), where
radial or fat-wedge shortcuts need not disappear.  Every rescaling which
unrolls the intrinsic slice sends the cone/fat-wedge base point to infinity,
so no such limit is a tangent cone based at that point.

In the same regime \(\beta_T\to\infty\), the corresponding full ambient
bubbles are
\begin{align}
 (Z_{12},\lambda_TD_{12},p_T^{12})
 &\longrightarrow(\mathbb R^8,0),\\
 (Z_{24}^{\mathcal D},\lambda_TD_{24}^{\mathcal D},
  p_T^{24,\mathcal D})
 &\longrightarrow(\mathbb R^{12},0),
 \qquad\mathcal D\in\{\mathcal B,\mathcal C\}.
 \label{sphere-metric-arbitrary-unrolling-ambient-bubbles}
\end{align}
These are pointed Gromov--Hausdorff limits of the singular metric targets.
The convergence of tensors is \(C^\infty\) only under coordinate
identifications on exhausting open regular-stratum neighborhoods; no
global smooth convergence of the singular targets is claimed.
\end{theorem}

\begin{proof}
For the \(Q_T^2\) route, the rescaled period matrix in each two-dimensional
factor is
\[
 \frac{\lambda_Tr_T}{T}Y_T
 =\beta_T\left(J+\frac{E_T}{T}\right).
\]
If \(\beta_T\to L\in(0,\infty)\), this lattice converges to
\(LJ\mathbb Z^2\); two factors give the square four-torus in
\eqref{sphere-metric-scaling-trichotomy-result}.  If
\(\beta_T\to0\), the lengths of a lattice basis tend to zero, so the
diameter tends to zero.  If \(\beta_T\to\infty\), the reverse triangle
inequality gives systole
\[
 \frac{\beta_T}{T}(T-C)
 =\beta_T\left(1-\frac CT\right)\longrightarrow\infty,
\]
so every fixed-radius ball is eventually Euclidean.

For a \(C_T^4\) slice, the four periods after rescaling have the common
length
\[
 \lambda_Tr_T=\beta_T,
\]
which gives the same three intrinsic cases directly.

Consider first one factor.  A curve from a level-\(r_T\) point to the cone
vertex which remains in the inner chart has length at least \(r_T\), because
the \(dr^2\)-term makes \(r\) \(1\)-Lipschitz there.  If it leaves the
chart, it has a first crossing of \(r=r_0\) and, before ending at the
vertex, a last re-entry through \(r=r_0\); its two inner-chart portions
have total length at least
\((r_0-r_T)+r_0>r_T\).  The radial segment has length \(r_T\), so the
factor distance to its based vertex is exactly \(r_T\).  The
Riemannian-\(\ell^2\)-product distance from every synchronized slice point
to the fat wedge is the minimum of these factor distances and hence is
also \(r_T\).  The metric-quotient formula in the preceding proof therefore
gives \eqref{sphere-metric-vertex-distance}.

It remains to prove the restricted-ambient comparison, so assume
\(\beta_T\to\infty\) and fix \(R\).  Two points in
\(\overline B_{\lambda_Td_T^{\mathrm{slice}}}(p_T,R)\) have intrinsic
distance at most \(2R/\lambda_T\).  The slice path is an ambient competitor,
so
\[
 d_T^{\mathrm{amb}}(x,y)\leq d_T^{\mathrm{slice}}(x,y).
\]
The quotient shortcut through the fat wedge has rescaled length
\(2\lambda_Tr_T=2\beta_T\), and hence is longer than \(2R\) for all large
\(T\).  A path which leaves the fixed cusp chart must reach \(r=r_0\).
Its outward cost is controlled by
\begin{equation}
 \lambda_T(r_0-r_T)
 =\beta_T\left(\frac{r_0}{r_T}-1\right)\longrightarrow\infty,
 \label{sphere-metric-outward-shortcut-exclusion}
\end{equation}
so no path of rescaled length at most \(2R\) can use the interpolating
annulus or the global part of a factor.

Consider now an ambient competitor of length at most
\(2R/\lambda_T+o(\lambda_T^{-1})\) which stays in the cusp chart.  In every
factor its radial coordinate satisfies
\[
 |r-r_T|\leq \frac{2R+o(1)}{\lambda_T},
 \qquad
 \frac r{r_T}=1+O_R(\beta_T^{-1}).
\]
On this tube, the horizontal tensors in the directions tangent to the slice
are uniformly \((1+\varepsilon_{T,R})\)-bilipschitz to their values at
\(r=r_T\), where \(\varepsilon_{T,R}\to0\).  For
\(\mathcal B\) and \(\mathcal C\) this follows directly from
\((r/r_T)^2=1+O_R(\beta_T^{-1})\).  For \(\mathcal A\), the same conclusion
follows, without changing the fixed coordinates or the matrix \(Y_T\), from
\[
 \begin{aligned}
  T(r)-T&=O_R(\beta_T^{-1}),\\
  \frac{T(r)}T&=1+O_R((\beta_TT)^{-1}),\\
  Y_{T(r)}Y_T^{-1}&=I+O_R((\beta_TT)^{-1}).
 \end{aligned}
\]
Projecting the competitor to the torus or character directions, and
discarding its radial and base-angle components, produces a path between
the same slice endpoints.  The bilipschitz estimate therefore gives
\[
 d_T^{\mathrm{slice}}(x,y)
 \leq(1+\varepsilon_{T,R})d_T^{\mathrm{amb}}(x,y).
\]
Together with the opposite inequality and
\(\lambda_Td_T^{\mathrm{slice}}(x,y)\leq2R\), this proves
\eqref{sphere-metric-intrinsic-ambient-agreement}.  The first inclusion in
\eqref{sphere-metric-intrinsic-ambient-ball-comparison} follows from
\(d_T^{\mathrm{amb}}\leq d_T^{\mathrm{slice}}\).  For a point \(x\) in
the middle ball, apply the same fat-wedge and outward exclusions to a
near-minimizing ambient path from \(p_T\) to \(x\), whose rescaled length is
at most \(R+o(1)\), and project that path to the slice directions.  This
gives
\[
 \lambda_Td_T^{\mathrm{slice}}(p_T,x)
 \leq(1+\varepsilon_{T,R})R
\]
after enlarging \(\varepsilon_{T,R}\) by a quantity tending to zero, and
proves the second inclusion.

For completeness, the full ambient bubbles at an arbitrary unrolling scale
are obtained without changing any coordinate or orientation.  In one
\(\mathcal A\)-factor use
\[
 s=\lambda_T(r-r_T),\qquad
 \eta=\beta_Tu,\qquad
 z=\frac{\beta_T}{T}Y_Tx.
\]
Writing \(A_T(r)=Y_{T(r)}Y_T^{-1}\), one obtains the exact expression
\[
 \lambda_T^2g_{\mathcal A}
 =ds^2+\left(\frac r{r_T}\right)^2d\eta^2
  +\left(\frac r{r_T}\right)^2
   \left(\frac{T}{T(r)}\right)^2
   dz^{\mathsf T}A_T(r)^{\mathsf T}A_T(r)dz.
\]
On every fixed bounded coordinate set,
\[
 \frac r{r_T}=1+O_R(\beta_T^{-1}),\qquad
 T(r)-T=O_R(\beta_T^{-1}),\qquad
 A_T(r)=I+O_R((\beta_TT)^{-1}),
\]
with the same estimates after every fixed number of derivatives.  Hence
the tensor converges smoothly to the Euclidean metric on \(\mathbb R^4\).
The \(\eta\)-period is \(\beta_T\), and the \(z\)-period lattice is
\[
 \frac{\beta_T}{T}Y_T\mathbb Z^2,
 \qquad
 \operatorname{sys}\left(
   \frac{\beta_T}{T}Y_T\mathbb Z^2\right)
 \geq\beta_T\left(1-\frac CT\right)\longrightarrow\infty.
\]
In one \(\mathcal B\)- or \(\mathcal C\)-factor use
\[
 s=\lambda_T(r-r_T),\qquad
 \eta=\beta_Tu,\qquad \zeta=\beta_T\theta.
\]
Then
\[
 \lambda_T^2g_{\mathcal D}
 =ds^2+\left(\frac r{r_T}\right)^2
       (d\eta^2+d\zeta^2),
\]
and both angular periods are \(\beta_T\to\infty\).

In all three cases the inward boundary and the fat wedge are at scaled
distance \(\beta_T\), while a lift starting at radius \(r_T\) incurs the
one-way radial cost in
\eqref{sphere-metric-outward-shortcut-exclusion} before its first outward
crossing of \(r=r_0\).  Fix \(R\).  For large \(T\), all these quantities
and all nonzero period lengths exceed \(4R\).
If \(x,y\) lie in the ambient rescaled radius-\(R\) ball, their distances
from the fat wedge are at least \(\beta_T-R\); hence a quotient path through
the smash base point has length at least \(2(\beta_T-R)>2R\).  Every
remaining competitor lifts to the complement of the fat wedge in the
product.  If such a lift exits the inner product chart, then some factor
has a first crossing of \(r=r_0\); the radial variation before that crossing
is at least \(\lambda_T(r_0-r_T)-R\), and returning to the lift of \(y\)
incurs the same lower bound.  Its length is therefore greater than \(2R\).
Finally, a path of length at most \(2R\) remains in the tube
\(r/r_T=1+O_R(\beta_T^{-1})\); there every nonzero deck translation has
length tending to infinity, so such a path cannot wrap around an angular
or torus period.  The displayed coordinate neighborhoods therefore exhaust
the pointed balls without quotient or global shortcuts.
Their smooth tensor convergence proves the pointed Gromov--Hausdorff
limits in
\eqref{sphere-metric-arbitrary-unrolling-ambient-bubbles}, with smooth
convergence confined to this regular-stratum exhaustion.
\end{proof}

\begin{corollary}[Unit-volume normalization does not give
\(\mathbb R^4\)]
\label{sphere-metric-unit-volume}
The four-volume of \(Q_T^2\) is \((\det Y_T)^2\sim T^4\).  Its unit-volume
length factor is therefore \((\det Y_T)^{-1/2}\sim T^{-1}\), and its
unit-volume limit is a square flat four-torus, not \(\mathbb R^4\).
Likewise, the four-volume of \(C_T^4\) is \(\ell_T^4\), and its unit-volume
normalization is exactly the square unit four-torus.
\end{corollary}

\begin{proposition}[What vertex tangents and compact asymptotic cones can
mean]
\label{sphere-metric-tangent-scope}
The compact target topology alone determines no metric tangent cone.
For an exact conic length metric built from fixed piecewise-Riemannian
metrics on the links, however, the vertex links have topological and
Hausdorff dimensions \(7\) for \(Z_{12}\) and \(11\) for both
\(Z_{24}\) targets.  The resulting vertex cones have topological and
Hausdorff dimensions \(8\) and \(12\), and are not \(\mathbb R^4\).

More precisely, writing \(L_3=T^3\) and \(L_2=T^2\), one has the based
topological models
\[
 Z_{12}\cong\operatorname{Susp}(L_3*L_3),
 \qquad
 Z_{24}^{\mathcal B}
 \cong\operatorname{Susp}(L_2*L_2*L_2*L_2).
\]
At the suspension base point their conic links are respectively
\(L_3*L_3\) and \(L_2^{*4}\).  Also
\[
 (\mathcal C,c_0)
 \cong
 (\mathcal M/\partial\mathcal M,[\partial\mathcal M]),
 \qquad
 (Z_{24}^{\mathcal C},*)
 \cong
 (\mathcal M^4/\partial(\mathcal M^4),
  [\partial(\mathcal M^4)]),
\]
so an exact conic metric at the collapsed boundary has link
\(\partial(\mathcal M^4)\).  At the collapsed point of
\(Z_{24}^{\mathcal C}\), the complete integral local-homology calculation is
\begin{equation}
 H_i\bigl(Z_{24}^{\mathcal C},
          Z_{24}^{\mathcal C}\setminus\{*\};\mathbb Z\bigr)
 \cong
 \begin{cases}
  \mathbb Z^{\binom4{i-1}},&2\leq i\leq5,\\
  0,&i=6,7,\\
  \mathbb Z^{\binom4{i-8}},&8\leq i\leq12,\\
  0,&\text{otherwise}.
 \end{cases}
 \label{sphere-metric-C-smash-local-homology}
\end{equation}
In particular, its local homology is not concentrated in top degree.  For
the first two targets the top local homology occurs in degree \(8\) and
\(12\), respectively; for the third target
\eqref{sphere-metric-C-smash-local-homology} includes a top
\(\mathbb Z\) in degree \(12\).  None has the local homology of
\(\mathbb R^4\).

Finally, the standard asymptotic cone obtained by blowing down any fixed
compact metric target is a point.  Indeed, if \(R_n\to\infty\), then
\[
 \operatorname{diam}(Z,R_n^{-1}D)
 =R_n^{-1}\operatorname{diam}(Z,D)\longrightarrow0.
\]
Thus the limits in \cref{sphere-metric-fixed-realization} are
moving-base-point blow-ups, often called bubbles.  They are neither
vertex-based tangent cones nor asymptotic cones of the compact targets.
\end{proposition}

\begin{proof}
For compact spaces \(V_i\) with closed subspaces \(A_i\), there is a
canonical homeomorphism
\[
 (V_1/A_1)\wedge\cdots\wedge(V_m/A_m)
 \cong
 \frac{V_1\times\cdots\times V_m}
      {\bigcup_iV_1\times\cdots\times A_i\times\cdots\times V_m}.
\]
Take \(V_i=\operatorname{Cone}(L_i)\) and \(A_i=L_i\), its outer boundary.
The product of \(m\) cones is homeomorphic to the cone on the iterated join
\(L_1*\cdots*L_m\), and its displayed union of faces is the outer link.
Collapsing that link gives its suspension.  This proves the first two
models and the link dimensions
\[
 \dim(T^3*T^3)=3+3+1=7,
 \qquad
 \dim((T^2)^{*4})=4\cdot2+3=11.
\]

We next prove the based homeomorphism for \(\mathcal C\), rather than
identifying the two presentations without a map.  Put
\(L=\partial\mathcal M\), write the attached cone as
\[
 \operatorname{Cone}(L)
 =L\times[0,1]/(L\times\{0\}),
\]
with base \(L\times\{1\}\), and choose a collar
\(\kappa:L\times[0,1]\hookrightarrow\mathcal M\) satisfying
\(\kappa(x,0)=x\).  There is a homeomorphism
\[
 \Phi:\mathcal C\setminus\{c_0\}
       \longrightarrow\operatorname{int}\mathcal M
\]
given on the punctured cone, the inner collar, and the remaining core by
\[
 \Phi([x,t])=\kappa(x,t/e)\quad(0<t\leq1),\qquad
 \Phi(\kappa(x,s))=\kappa(x,e^{s-1})\quad(0\leq s\leq1),
\]
\[
 \Phi(q)=q
 \quad\text{for}\quad
 q\in\mathcal M\setminus\kappa(L\times[0,1)).
\]
The first two formulas agree at the glued points
\([x,1]=\kappa(x,0)\), and the second and third agree at
\(\kappa(x,1)\).  Their collar-coordinate images are respectively
\(L\times(0,e^{-1}]\) and \(L\times[e^{-1},1]\); hence the displayed
piecewise map and its piecewise inverse are continuous.

The compact Hausdorff spaces \(\mathcal C\) and
\(\mathcal M/L\) are both one-point compactifications of
\(\operatorname{int}\mathcal M\): for \(\mathcal C\) use \(\Phi\), and for
\(\mathcal M/L\) use the quotient map on
\(\mathcal M\setminus L=\operatorname{int}\mathcal M\).  Uniqueness of the
one-point compactification therefore extends \(\Phi\) to a based
homeomorphism
\begin{equation}
 (\mathcal C,c_0)
 \xrightarrow{\ \cong\ }
 (\mathcal M/\partial\mathcal M,[\partial\mathcal M]).
 \label{sphere-metric-C-boundary-quotient-homeomorphism}
\end{equation}
Only after establishing
\eqref{sphere-metric-C-boundary-quotient-homeomorphism} do we apply the
quotient identity at the start of the proof.  It gives
\[
 \mathcal C^{\wedge4}
 \cong
 \frac{\mathcal M^4}
 {\displaystyle\bigcup_{a=1}^4
  \mathcal M^{a-1}\times\partial\mathcal M
  \times\mathcal M^{4-a}}
 =\mathcal M^4/\partial(\mathcal M^4),
\]
and carries the smash base point to the collapsed boundary.

Put \(N=\mathcal M^4\).  Rounding its product corners preserves the
homeomorphism type of the pair \((N,\partial N)\) and makes
\(\partial N\) an eleven-dimensional closed manifold.  A boundary collar
descends in \(N/\partial N\) to a cone neighborhood of the collapsed
point, so the local link is exactly \(\partial N\) and
\begin{equation}
 H_i(N/\partial N,(N/\partial N)\setminus\{*\};\mathbb Z)
 \cong\widetilde H_{i-1}(\partial N;\mathbb Z).
 \label{sphere-metric-boundary-link-local-shift}
\end{equation}

The integral homology calculation in the proof of
\cref{cusp-cone-S24-route} for the specific compact orientable
three-manifold \(\mathcal M\) gives
\[
 H_j(\mathcal M;\mathbb Z)\cong
 \begin{cases}
  \mathbb Z,&j=0,1,\\
  0,&\text{otherwise}.
 \end{cases}
\]
All groups are free, so the integral Kunneth theorem gives
\[
 H_j(N;\mathbb Z)\cong
 \begin{cases}
  \mathbb Z^{\binom4j},&0\leq j\leq4,\\
  0,&\text{otherwise}.
 \end{cases}
\]
Poincare--Lefschetz duality, with the product orientation retained, gives
\[
 H_j(N,\partial N;\mathbb Z)
 \cong H^{12-j}(N;\mathbb Z)
 \cong
 \begin{cases}
  \mathbb Z^{\binom4{j-8}},&8\leq j\leq12,\\
  0,&\text{otherwise}.
 \end{cases}
\]
The long exact sequence of \((N,\partial N)\) now has disjoint absolute
and relative degree ranges.  It gives
\[
 \widetilde H_j(\partial N;\mathbb Z)\cong
 \begin{cases}
  \mathbb Z^{\binom4j},&1\leq j\leq4,\\
  0,&j=5,6,\\
  \mathbb Z^{\binom4{j-7}},&7\leq j\leq11,\\
  0,&\text{otherwise}.
 \end{cases}
\]
Substitution in
\eqref{sphere-metric-boundary-link-local-shift} proves
\eqref{sphere-metric-C-smash-local-homology}.

An exact metric cone is invariant under every radial blow-up at its vertex,
and its Hausdorff dimension is one plus that of its fixed
piecewise-Riemannian link.  For the first target,
\[
 H_8\bigl(\operatorname{Cone}(T^3*T^3),T^3*T^3;\mathbb Z\bigr)
 \cong\mathbb Z;
\]
the analogous top local group for the \(\mathcal B\) target is
\(\mathbb Z\) in degree \(12\), while the full \(\mathcal C\) calculation
is \eqref{sphere-metric-C-smash-local-homology}.  Euclidean
\(\mathbb R^4\) has local homology \(\mathbb Z\) only in degree \(4\), so
none of these conic vertex models is \(\mathbb R^4\).  The diameter
calculation proves the final assertion.
\end{proof}

\paragraph{Consequences for the three sphere routes.}
The exact positive statement is that one fixed metric on each compact
codomain contains the native \(\mathbb R^4\) decompactification as a
vertical four-plane in a moving-point Euclidean bubble.  The whole bubbles
are \(\mathbb R^8\) for the \(S^{12}\) route and \(\mathbb R^{12}\) for the
\(S^{24}\) routes.  The factor \(\beta_T=\lambda_Tr_T\) proves that this
dimension discrepancy cannot be removed by a different scalar rescaling
of the same fixed cusp-adapted metric: the slice unrolls only when the
smash base point has already gone to infinite scaled distance.

No metric on a source sphere is supplied by these quotient constructions.
If an independent smooth Riemannian metric is placed on \(S^{12}\) or
\(S^{24}\), its ordinary tangent spaces are \(\mathbb R^{12}\) or
\(\mathbb R^{24}\).  Pulling a codomain metric back through
\(q_2^{\wedge2}\) or \(q_1^{\wedge4}\), on the regular strata where the
maps are differentiable, gives a degenerate quadratic form along the
positive-dimensional quotient kernels, not a Riemannian source metric.

The moving four-sphere carriers
\[
 \mathbb S_{12,T}^4=S_T^2\wedge S_T^2\subset S^6,\qquad
 \mathbb S_{24,T}^4=C_T^{\wedge4}\subset S^8
\]
of \cref{cusp-cone-angle-forgotten-sphere-carriers} carry separately
assigned metrics.  Namely, \(S_T^2\) first receives the quotient length
metric obtained from
\((\overline T^2,h_T)\cong(Q_T,g_{Q_T})\) by collapsing
\(K^{(1)}\subset\overline T^2\), equivalently by collapsing
\(\iota_T(K^{(1)})\subset Q_T\).  Then
\(\mathbb S_{12,T}^4\) receives the quotient length metric of the
Riemannian \(\ell^2\)-product \((S_T^2)^2\) after its fat wedge is
collapsed.  Likewise, \(\mathbb S_{24,T}^4\) receives the quotient length
metric of the Riemannian \(\ell^2\)-product \(C_T^4\) after its fat wedge is
collapsed.  These metrics are assigned directly to the moving carriers;
they are unrescaled intrinsic quotient metrics, not restrictions of
\(D_{12}\) or \(D_{24}^{\mathcal B}\), whose domains are the conical
targets, and no fixed ambient metric on \(S^6\) or \(S^8\) is used to
define them.  Their Euclidean balls are
based at the midpoint classes
\[
 \begin{aligned}
  y_*&=[(1/2,1/2)]\in\overline T^2,\qquad z_T=c(y_*),\\
  \zeta_{12,T}&=[z_T,z_T],\\
  \zeta_{24,T}
  &=[(\ell_T/2,\ell_T/2,\ell_T/2,\ell_T/2)].
 \end{aligned}
\]
not at the images of \(p_T^{12}\) or \(p_T^{24,\mathcal B}\), which belong
to the subsequently collapsed loci.  In particular, the trichotomy in
\cref{sphere-metric-scaling-trichotomy} is the trichotomy of the
uncollapsed intrinsic flat-torus slices; it is not a finite-\(\beta_T\)
description of the singular quotient-metric moving four-spheres.  Indeed,
the restriction of any one compatible metric on the compact sphere
\(S^6\) or \(S^8\) has diameter bounded by the diameter of that sphere.  By
contrast, in the separately assigned moving metrics the midpoint distances
to the collapsed loci are at least \((T-C)/2\) and exactly
\(\ell_T/2\), respectively, and hence tend to infinity.  Thus the assigned
unrescaled metrics cannot themselves be restrictions of one fixed ambient
metric.  This does not preclude a separately proved realization after a
scalar blow-up of another fixed quotient metric.

Finally, the rank obstruction for a full-base torus quotient must retain
its actual metric hypothesis: \emph{quotient systole tends to infinity}
means the quotient norm induced by the unrescaled Euclidean fibre metric.
Without that qualification the assertion is false, since an arbitrary
\(T\)-dependent scalar multiplication can make the systole of any fixed
nonzero quotient tend to infinity.


\subsection{The exceptional first-coordinate elliptic quotient}

The native real quotient in \cref{cusp-spec-native-exact-sequence} exists
throughout the radial cusp.  A holomorphic quotient obtained from the first
complex coordinate is a different and much more restrictive construction.
We determine its locus without replacing it by an ambient coefficient-space
path.

For \(z\in\mathfrak h_z\), write
\begin{equation}
 \begin{aligned}
  v_{\gamma,z}&=(6\mu(z),\beta(z))^{\mathsf t},&
  v_{u,z}&=(\tau(z),\mu(z))^{\mathsf t},\\
  v_{w,z}&=(1,0)^{\mathsf t},&
  v_{\delta,z}&=(0,1)^{\mathsf t},
 \end{aligned}
 \label{cusp-spec-elliptic-four-periods}
\end{equation}
and put
\begin{equation}
 \mathcal A_z:=\mathbb C^2/
 \mathbb Z\langle v_{\gamma,z},v_{u,z},v_{w,z},v_{\delta,z}\rangle,
 \qquad
 E_{\tau(z)}:=\mathbb C/(\mathbb Z\tau(z)+\mathbb Z).
 \label{cusp-spec-elliptic-source-target}
\end{equation}
These are the marked fibres of \cref{prop:period-input-audit} and the
elliptic curves selected by the same period \(\tau\).

\begin{proposition}[Exact first-coordinate descent and connected kernel]
\label{cusp-spec-first-coordinate-descent}
The linear map
\(\operatorname{pr}_1:\mathbb C^2\to\mathbb C\) induces a surjective
holomorphic homomorphism
\begin{equation}
 \overline{\operatorname{pr}}_{1,z}:\mathcal A_z\longrightarrow
 E_{\tau(z)}
 \label{cusp-spec-first-coordinate-homomorphism}
\end{equation}
if and only if there is a necessarily unique pair
\((a,b)\in\mathbb Z^2\) such that
\begin{equation}
 6\mu(z)=a\tau(z)+b.
 \label{cusp-spec-first-coordinate-condition}
\end{equation}
On this locus define
\begin{equation}
 \tau_f(z):=a\mu(z)-\beta(z).
 \label{cusp-spec-kernel-period}
\end{equation}
Then
\begin{equation}
 \operatorname{Im}\tau_f(z)
 =-D_{\mathrm{can}}(z)>0,
 \qquad
 D_{\mathrm{can}}(z):=\operatorname{Im}\beta(z)
 -\frac{6(\operatorname{Im}\mu(z))^2}{\operatorname{Im}\tau(z)},
 \label{cusp-spec-kernel-period-imaginary-part}
\end{equation}
and the kernel in \cref{cusp-spec-first-coordinate-homomorphism} is connected
and isomorphic to
\begin{equation}
 E_{\tau_f(z)}=\mathbb C/(\mathbb Z+\mathbb Z\tau_f(z)).
 \label{cusp-spec-connected-kernel-elliptic-curve}
\end{equation}
\end{proposition}

\begin{proof}
The four first coordinates in \cref{cusp-spec-elliptic-four-periods} are
\(6\mu(z),\tau(z),1,0\).  Thus the linear projection descends precisely when
\(6\mu(z)\in\mathbb Z\tau(z)+\mathbb Z\), which is
\cref{cusp-spec-first-coordinate-condition}.  The pair is unique because
\(\operatorname{Im}\tau(z)>0\).  The periods \(v_{u,z}\) and \(v_{w,z}\)
map onto the two generators of the target lattice, so the induced map is
surjective.

Replace the marked lattice basis by the unimodular basis whose first vector
is
\begin{equation}
 v_{\gamma,z}-av_{u,z}-bv_{w,z}
 =(0,\beta(z)-a\mu(z))^{\mathsf t}
 =(0,-\tau_f(z))^{\mathsf t}.
 \label{cusp-spec-vertical-unimodular-period}
\end{equation}
Together with \(v_{\delta,z}=(0,1)^{\mathsf t}\), these are exactly the
vertical periods.  Indeed, after subtracting integral multiples of
\(v_{u,z}\) and \(v_{w,z}\), every point in the kernel has first coordinate
zero; and an integral combination of \(v_{u,z},v_{w,z}\) has first
coordinate zero only when both coefficients vanish.  The kernel is therefore
the connected one-dimensional complex torus in
\cref{cusp-spec-connected-kernel-elliptic-curve}.

Taking imaginary parts in
\cref{cusp-spec-first-coordinate-condition} gives
\(6\operatorname{Im}\mu=a\operatorname{Im}\tau\).  Hence
\[
 \operatorname{Im}\tau_f
 =a\operatorname{Im}\mu-\operatorname{Im}\beta
 =\frac{6(\operatorname{Im}\mu)^2}{\operatorname{Im}\tau}
  -\operatorname{Im}\beta
 =-D_{\mathrm{can}}>0
\]
by \cref{prop:period-input-audit}.
\end{proof}

For \((a,b)\in\mathbb Z^2\), define the marked analytic set
\begin{equation}
 \mathcal D_{a,b}:=
 \{z\in\mathfrak h_z:6\mu(z)-a\tau(z)-b=0\},
 \qquad
 \mathcal D:=\bigcup_{(a,b)\in\mathbb Z^2}\mathcal D_{a,b}.
 \label{cusp-spec-elliptic-descent-locus}
\end{equation}

\begin{theorem}[Monodromy, cusp exclusion, and finiteness of the descent locus]
\label{cusp-spec-elliptic-descent-finiteness}
The transformation laws of \cref{prop:period-input-audit} act on the integer
labels in \cref{cusp-spec-elliptic-descent-locus} by
\begin{equation}
 g_0:(a,b)\longmapsto(a,a+b),
 \qquad
 g_2:(a,b)\longmapsto(-b,6+a).
 \label{cusp-spec-descent-label-monodromy}
\end{equation}
Consequently \(\mathcal D\) is invariant under the triangle group and has a
well-defined image in the completed orbifold base \(B\simeq\mathbb P^1\).
That image is finite and is disjoint from some punctured neighborhood of the
distinguished cusp \(p_0\).
\end{theorem}

\begin{proof}
Suppose \(z\in\mathcal D_{a,b}\).  From
\(\tau\circ g_0=\tau-1\) and \(\mu\circ g_0=\mu\),
\[
 6\mu(g_0z)=a(\tau(g_0z)+1)+b
 =a\tau(g_0z)+(a+b),
\]
so \(g_0z\in\mathcal D_{a,a+b}\).  Similarly,
\(\tau\circ g_2=-1/\tau\) and
\(\mu\circ g_2=1+\mu/\tau\) give
\[
 6\mu(g_2z)=6+a+\frac b{\tau(z)}
 =-b\tau(g_2z)+(6+a),
\]
so \(g_2z\in\mathcal D_{-b,6+a}\).  Both label maps are bijections; and
\(g_0,g_2\) generate the triangle group because
\(g_0=(g_1g_2)^{-1}\).  This proves invariance.

We next exclude the cusp with the integer labels allowed to vary.  If no
punctured cusp neighborhood were disjoint from \(\mathcal D\), there would be
points \(z_j\in\mathcal D_{a_j,b_j}\) in the distinguished cusp component
with \(t_c(z_j)\to0\).  Put
\(T_j=\operatorname{Im}\tau(z_j)\) and
\(m_j=\operatorname{Im}\mu(z_j)\).  The cusp construction in
\cref{prop:period-input-audit} and the formulas preceding
\cref{cusp-spec-real-period-matrix} give
\(T_j\to\infty\) while \(m_j\) stays bounded.  Taking imaginary parts yields
\begin{equation}
 a_j=\frac{6m_j}{T_j}\longrightarrow0.
 \label{cusp-spec-cusp-label-a-limit}
\end{equation}
Because \(a_j\) is integral, \(a_j=0\) for all sufficiently large \(j\).
Then \(6\mu(z_j)=b_j\).  The cusp extension of \(\mu\) is bounded, so the
integers \(b_j\) are bounded; after taking a subsequence, they equal one
integer \(b\).  The holomorphic function \(6\mu-b\) extends across
\(t_c=0\) and has zeros accumulating there.  It therefore vanishes on the
cusp disc and, by the identity theorem, on the connected surface
\(\mathfrak h_z\).

More generally, no identity
\(6\mu=a\tau+b\) with fixed integers \(a,b\) can hold on
\(\mathfrak h_z\).  Applying \(g_0\) to such an identity and subtracting it
from the original one gives \(a=0\).  Then \(\mu=b/6\) is constant, whereas
the \(g_2\)-law would give
\begin{equation}
 \frac b6=1+\frac{b}{6\tau}.
 \label{cusp-spec-constant-mu-contradiction}
\end{equation}
This is impossible for the nonconstant map \(\tau\).  The preceding
conclusion \(6\mu-b\equiv0\) is therefore a contradiction, proving cusp
exclusion.

Choose a closed fundamental polygon for the triangle group and remove the
excluded cusp neighborhood.  Its closure \(K\) is compact in
\(\mathfrak h_z\), including the two elliptic fixed points.  The functions
\(\tau,\mu\) are holomorphic there and
\(\inf_K\operatorname{Im}\tau>0\).  At a point of
\(\mathcal D_{a,b}\cap K\),
\begin{equation}
 a=\frac{6\operatorname{Im}\mu}{\operatorname{Im}\tau},
 \qquad
 b=6\mu-a\tau.
 \label{cusp-spec-compact-label-bounds}
\end{equation}
Thus only finitely many integer pairs \((a,b)\) occur on \(K\).  For each of
them, \(6\mu-a\tau-b\) is a nonzero holomorphic function on a neighborhood of
the compact set \(K\), so its zero set meets \(K\) in finitely many points.
The \((3,4,\infty)\) triangle group has one cusp orbit.  Hence cusp exclusion
holds at every lift of the chosen cusp neighborhood.  The set
\(\mathcal D\) is triangle-group invariant by the already proved label
permutation, and it is disjoint from all those lifted cusp neighborhoods.
Consequently every triangle-group orbit contained in \(\mathcal D\) meets
\(K\).  Hence the image of \(\mathcal D\) in the completed orbifold base is
finite.
\end{proof}

\begin{remark}[Exact scope of the elliptic classification]
\label{cusp-spec-elliptic-descent-scope}
The result classifies only the quotient obtained from the specified first
coordinate and the specified target lattice
\(\mathbb Z\tau+\mathbb Z\).  Its cusp exclusion does not prove that
\(\mathcal A_z\) has no other elliptic quotient or elliptic subtorus, that its
N\'eron--Severi group vanishes, or that it carries no smooth integral
magnetic line bundle.  Those are different typed questions.  In particular,
the first-coordinate elliptic construction cannot supply the gap-closing
cusp sequence, whereas the native scalar and classical gauge soft modes
proved above do not use this exceptional locus.
\end{remark}

\begin{theorem}[Base-changed cover monodromy and exact momentum sectors]
\label{cusp-spec-cover-momentum-theorem}
In the ordered integral basis
\((\widehat\gamma,\widehat u,\widehat w,\widehat\delta)\), put
\begin{equation}
 \mathsf S_M:=\operatorname{diag}(1,1,M,M),
 \qquad
 \Lambda^{(M)}:=\mathsf S_M\mathbb Z^4.
 \label{cusp-spec-cover-index-lattice}
\end{equation}
Then \(P_T\Lambda^{(M)}=\Lambda_{T,M}\), and the cover in
\cref{cusp-spec-degree-cover} is induced by the lattice inclusion
\[
 P_T\mathsf S_M\mathbb Z^4\subset P_T\mathbb Z^4.
\]
Writing \(M_0=I+R\), one has
\begin{equation}
 R^2=0,\qquad
 \ker R=\operatorname{im}R
 =\mathbb Z\langle\widehat w,\widehat\delta\rangle,
 \label{cusp-spec-monodromy-kernel-image}
\end{equation}
and
\begin{equation}
 M_0^M=I+MR,\qquad
 \mathsf S_M^{-1}M_0^M\mathsf S_M=M_0.
 \label{cusp-spec-cover-monodromy-conjugacy}
\end{equation}
Thus, after the cyclic base change \(t_c=v^M\), the degree-\(M^2\)
fibre covers form a regular relative cover over \(v\ne0\).  This particular
statement proves only punctured descent.  The distinct completed local
extension, with the required modified source action and the separate
degree-\(M^3\) and degree-\(M^2\) maps, is proved in
\cref{cusp-cover-extension-theorem}.

Equip \(\widetilde F_{T,M}\) with the pullback Euclidean metric.  Its
Riemannian volume is \(M^2D_T\), so normalized Haar measure is exactly
\[
 \mathrm d\widetilde\nu_{T,M}
 :=(M^2D_T)^{-1}\mathrm d\operatorname{vol}_{\widetilde F_{T,M}}.
\]
Let \(\widetilde\Delta_{T,M}\) be the nonnegative self-adjoint scalar
Laplacian on \(L^2(\widetilde F_{T,M},\mathrm d\widetilde\nu_{T,M})\)
obtained from the closed Dirichlet form
\[
 \widetilde{\mathfrak q}_{T,M}[f]
 :=\int_{\widetilde F_{T,M}}|\mathrm df|^2\,
        \mathrm d\widetilde\nu_{T,M},
 \qquad
 \operatorname{Dom}(\widetilde{\mathfrak q}_{T,M})
 =H^1(\widetilde F_{T,M}),
\]
with operator domain
\(\operatorname{Dom}(\widetilde\Delta_{T,M})
=H^2(\widetilde F_{T,M})\).  Its complete orthonormal character basis is
\begin{equation}
 \widetilde\chi_{T,M,k}([x])
 =
 \exp\!\left(
  2\pi i
  \left\langle P_T^{-\mathsf T}\mathsf S_M^{-1}k,x\right\rangle
 \right),
 \qquad k\in\mathbb Z^4,
 \label{cusp-spec-cover-characters}
\end{equation}
and
\begin{equation}
 \operatorname{Spec}(\widetilde\Delta_{T,M})
 =
 \left\{
  4\pi^2
  \left\lvert P_T^{-\mathsf T}\mathsf S_M^{-1}k\right\rvert^2:
  k\in\mathbb Z^4
 \right\}
 \label{cusp-spec-cover-scalar-spectrum}
\end{equation}
with multiplicity.  For this covering map one has
\begin{equation}
 c_{T,M}^*\chi_{T,\ell}
 =\widetilde\chi_{T,M,\mathsf S_M\ell}.
 \label{cusp-spec-cover-character-pullback}
\end{equation}
Consequently the pullback, or equivalently deck-invariant, sector consists
exactly of indices
\begin{equation}
 k\in\mathsf S_M\mathbb Z^4
 \quad\Longleftrightarrow\quad
 k_3\equiv k_4\equiv0\pmod M.
 \label{cusp-spec-cover-invariant-sector}
\end{equation}

There is a constant \(C_{\mathrm{mom}}>0\), independent of sufficiently
large \(T\) and of \(M\ge1\), such that the Euclidean covering radius of
the cover momentum lattice satisfies
\begin{equation}
 \sup_{p\in\mathbb R^4}
 \operatorname{dist}\!\left(
  p,P_T^{-\mathsf T}\mathsf S_M^{-1}\mathbb Z^4
 \right)
 \le C_{\mathrm{mom}}\left(\frac1T+\frac1M\right).
 \label{cusp-spec-cover-momentum-mesh}
\end{equation}
In particular, if \(T_j,M_j\to\infty\), the cover momentum lattices become
locally dense in all of \(\mathbb R^4\), and every \(E\ge0\) is a limit of
scalar eigenvalues in
\cref{cusp-spec-cover-scalar-spectrum}.
\end{theorem}

\begin{proof}
The equality \(P_T\Lambda^{(M)}=\Lambda_{T,M}\) follows by multiplying the
last two period columns \(e_1,e_3\) by \(M\).  The matrix in
\cref{prop:monodromy-certificate} gives
\[
 R(a,b,c,d)^{\mathsf t}=(0,0,b,-a)^{\mathsf t}.
\]
This proves \cref{cusp-spec-monodromy-kernel-image} and \(R^2=0\).
The binomial theorem then gives \(M_0^M=I+MR\).  For every
\((a,b,c,d)^{\mathsf t}\),
\[
 \mathsf S_M^{-1}M_0^M\mathsf S_M
 (a,b,c,d)^{\mathsf t}
 =(a,b,c+b,d-a)^{\mathsf t}
 =M_0(a,b,c,d)^{\mathsf t},
\]
which proves \cref{cusp-spec-cover-monodromy-conjugacy}.  In particular
\(M_0^M\Lambda^{(M)}=\Lambda^{(M)}\), which is precisely the descent
condition after \(t_c=v^M\).

The dual of \(P_T\mathsf S_M\mathbb Z^4\) is
\[
 (P_T\mathsf S_M)^{-\mathsf T}\mathbb Z^4
 =P_T^{-\mathsf T}\mathsf S_M^{-1}\mathbb Z^4.
\]
The standard Fourier theorem for a flat torus proves
\cref{cusp-spec-cover-characters,cusp-spec-cover-scalar-spectrum}.
Substitution of \(k=\mathsf S_M\ell\) proves
\cref{cusp-spec-cover-character-pullback}.  A character on the cover
descends through \(c_{T,M}\) exactly when it is invariant under the deck
group, and the displayed equality shows that this occurs exactly for
\(k\in\mathsf S_M\mathbb Z^4\), proving
\cref{cusp-spec-cover-invariant-sector}.

It remains to prove the quantitative density claim.  Let
\(b_{j,T,M}\) be the four columns of
\(P_T^{-\mathsf T}\mathsf S_M^{-1}\).  The exact matrix
\cref{cusp-spec-full-inverse-transpose}, the boundedness of
\(\alpha_T,\eta_T,c_T,m_T\), and
\[
 L_T,q_T=O(T),\qquad D_T\asymp T^2
\]
give, after increasing a uniform constant \(C_0\),
\[
 \lvert b_{1,T,M}\rvert+\lvert b_{2,T,M}\rvert
 \le\frac{C_0}{T},
 \qquad
 \lvert b_{3,T,M}\rvert+\lvert b_{4,T,M}\rvert
 \le\frac{C_0}{M}.
\]
Given \(p\in\mathbb R^4\), express it uniquely as
\(p=\sum_jt_jb_{j,T,M}\), and choose \(n_j\in\mathbb Z\) with
\(\lvert t_j-n_j\rvert\le\tfrac12\).  Then
\[
 \left\lvert p-\sum_jn_jb_{j,T,M}\right\rvert
 \le\frac12\sum_j\lvert b_{j,T,M}\rvert
 \le\frac{C_0}{2}\left(\frac1T+\frac1M\right).
\]
This proves \cref{cusp-spec-cover-momentum-mesh}.  For \(E>0\), apply it
to any \(p\) with \(4\pi^2|p|^2=E\); continuity gives eigenvalues tending
to \(E\).  For \(E=0\), use \(k=(1,0,0,0)\), whose momentum norm is
\(O(T^{-1})\).
\end{proof}

\paragraph{Scope and residual data.}
The conclusions above use the explicitly declared Euclidean quotient metric.
The canonical period theorem proves the lattice
family but does not itself specify a Riemannian metric, so no assertion about
a different metric follows without an explicit comparison.

The constants \(H\) and \(B\) determine the first relative
\(T^{-1}\)-correction.  When \(B=0\), they do not by themselves determine
which of the two unit pairs is strictly shortest at the next order.  The
canonical functions do determine it through the exact expression
\(B_T(2L_T-B_T)-35m_T^2\) in
\cref{cusp-spec-unit-mode-comparison}; if both \(B=0\) and \(m=0\), its sign
can depend on higher Taylor coefficients of \(b\) and \(\mu\).  This is an
indeterminacy of the two requested zeroth-order constants, not a missing
proof obligation.
The proved construction permits an additive constant in \(\beta\) and chooses
its imaginary part only by the qualitative requirement that the canonical
determinant be negative.  It records no numerical additive normalization.
Accordingly, the formulas retain the resulting value
\(B=\operatorname{Im}b(0)\) and do not silently impose a sign on it.

Finally, the degree-\(M^2\) covers in
\cref{cusp-spec-degree-cover} are local in the marked radial family, exactly
as stated.  The sublattice
\(\mathbb Z\langle\widehat\gamma,\widehat u,
M\widehat w,M\widehat\delta\rangle\) is not invariant under the full cusp
monodromy when \(M>1\): the proved monodromy sends
\(\widehat\gamma\) to \(\widehat\gamma-\widehat\delta\) and
\(\widehat u\) to \(\widehat u+\widehat w\).  Therefore these particular
degree-\(M^2\) covers do not, without additional construction, descend
around the entire punctured cusp.  Its \(M\)-th power does preserve the
sublattice, because it adds \(-M\widehat\delta\) to
\(\widehat\gamma\) and \(M\widehat w\) to \(\widehat u\).  Equivalently,
\(M_0^M=I+M(M_0-I)\).  Thus, after the degree-\(M\) cyclic base change
\(t_c=v^M\), the cover forms a regular relative degree-\(M^2\) cover over
\(v\ne0\).  The pointed-limit proof itself uses no central-fibre extension.
Nevertheless, \cref{cusp-cover-extension-theorem} proves that a corrected
source action extends this cover through \(v=0\) as a finite morphism of
completed \emph{local} cusp families.  It remains true that no cover over the
entire original base, and no cover of the compact threefold \(X\), is used or
constructed here.

\section{Classical gauge fluctuation operators on the positive radial cusp}
\label{cusp-gauge-section}

Retain the positive radial flat tori
\(F_T=\mathbb R^4/P_T\mathbb Z^4\), their Euclidean quotient metrics, and
the real quantities \(D_T,L_T,q_T,m_T,B_T,H,B\) from
\cref{cusp-spec-exact-dual-lattice,cusp-spec-shortest-modes}.  Thus
\[
 D_T=L_Tq_T+6m_T^2>0,\qquad
 L_T=\operatorname{Im}\tau_T,\qquad
 q_T=L_T-B_T=-\operatorname{Im}\beta_T,
\]
and \(D_T\) is the positive covolume.  Let
\[
 \mathbf e_\gamma=(1,0,0,0)^{\mathsf t},\qquad
 \mathbf e_u=(0,1,0,0)^{\mathsf t};
\]
then the already-proved dual covectors are
\[
 \ell_{\gamma,T}=P_T^{-\mathsf T}\mathbf e_\gamma,\qquad
 \ell_{u,T}=P_T^{-\mathsf T}\mathbf e_u,
\]
with exact squared lengths
\begin{equation}
 a_{\gamma,T}:=\lVert\ell_{\gamma,T}\rVert^2
 =\frac{L_T^2+m_T^2}{D_T^2},\qquad
 a_{u,T}:=\lVert\ell_{u,T}\rVert^2
 =\frac{q_T^2+36m_T^2}{D_T^2}.
 \label{cusp-gauge-short-mode-lengths}
\end{equation}
For all sufficiently large \(T\), the shortest nonzero dual vectors are
precisely the minimizing members of
\(\{\pm\ell_{\gamma,T},\pm\ell_{u,T}\}\).  Put
\begin{equation}
 G_T=P_T^{\mathsf T}P_T.
 \label{cusp-gauge-coordinate-gram}
\end{equation}
The matrix \(G_T\) is the tangent Gram matrix in the lattice coordinates
\(y\in\mathbb R^4/\mathbb Z^4\), where \(x=P_Ty\).

\subsection{The continuum Yang--Mills Hessian}
\label{cusp-gauge-continuum-subsection}

Let \(G\) be a compact Lie group with real Lie algebra \(\mathfrak g\).
Fix a positive-definite
\(\operatorname{Ad}(G)\)-invariant inner product on \(\mathfrak g\), put
\[
 d_G=\dim_{\mathbb R}\mathfrak g,\qquad
 \mathfrak g_{\mathbb C}=\mathfrak g\otimes_{\mathbb R}\mathbb C,
\]
and fix a coupling \(g_{\mathrm{YM}}>0\).  On the trivial principal bundle
\[
 \mathcal P_T=F_T\times G\longrightarrow F_T
\]
let \(A_0\) be the product connection.  The Yang--Mills functional is
normalized by
\begin{equation}
 \operatorname{YM}_T(A)
 =\frac1{2g_{\mathrm{YM}}^2}
 \int_{F_T}\lvert F_A\rvert^2\,\mathrm d\operatorname{vol}_{F_T}.
 \label{cusp-gauge-ym-normalization}
\end{equation}

\begin{proposition}[Hessian at the product connection]
\label{cusp-gauge-continuum-hessian}
For smooth
\(a,b\in\Omega^1(F_T;\mathfrak g)\),
\begin{equation}
 \operatorname{Hess}_{A_0}\operatorname{YM}_T(a,b)
 =\frac1{g_{\mathrm{YM}}^2}
 \langle\mathrm da,\mathrm db\rangle_{L^2}.
 \label{cusp-gauge-continuum-hessian-form}
\end{equation}
Relative to the unscaled \(L^2\) inner product on one-forms, the associated
closed Hessian form is
\begin{equation}
 \mathfrak l_T[a]
 =g_{\mathrm{YM}}^{-2}\lVert \mathrm d_1a\rVert_{L^2}^2,
 \qquad
 \operatorname{Dom}(\mathfrak l_T)=\operatorname{Dom}(\mathrm d_1),
 \label{cusp-gauge-continuum-closed-form}
\end{equation}
on \(L^2\Omega^1(F_T;\mathfrak g)\), where \(\mathrm d_1\) is the
maximal closed exterior derivative
\[
 \operatorname{Dom}(\mathrm d_1)
 =\{a\in L^2\Omega^1(F_T;\mathfrak g):
       \mathrm da\in L^2\Omega^2(F_T;\mathfrak g)
       \text{ distributionally}\}.
\]
The associated nonnegative self-adjoint operator is
\begin{equation}
 \mathcal L_T=g_{\mathrm{YM}}^{-2}\mathrm d_1^*\mathrm d_1,
 \qquad
 \operatorname{Dom}(\mathcal L_T)
 =\{a\in\operatorname{Dom}(\mathrm d_1):
       \mathrm d_1a\in\operatorname{Dom}(\mathrm d_1^*)\}.
 \label{cusp-gauge-continuum-hessian-operator}
\end{equation}
Its differential expression on smooth one-forms is
\(g_{\mathrm{YM}}^{-2}\mathrm d^*\mathrm d\).
\end{proposition}

\begin{proof}
With the convention
\(F_{A_0+a}=\mathrm da+\tfrac12[a\wedge a]\), one has
\[
 F_{A_0+sa+tb}
 =s\,\mathrm da+t\,\mathrm db
 +O\bigl((|s|+|t|)^2\bigr).
\]
The coefficient of \(st\) in the squared norm is
\(2\langle\mathrm da,\mathrm db\rangle_{L^2}\).  Taking the mixed
derivative in \cref{cusp-gauge-ym-normalization} proves
\cref{cusp-gauge-continuum-hessian-form}.  The maximal distributional
exterior derivative is closed and densely defined.  Hence its graph-norm
square is a closed nonnegative form, and the first representation theorem
identifies its operator with
\(g_{\mathrm{YM}}^{-2}\mathrm d_1^*\mathrm d_1\) on precisely the domain
displayed in \cref{cusp-gauge-continuum-hessian-operator}.
\end{proof}

\subsection{The finite cubical cochain complex}
\label{cusp-gauge-finite-subsection}

Fix \(N\ge2\) and put
\[
 \Lambda_N=(\mathbb Z/N\mathbb Z)^4.
\]
The vertex \(n\in\Lambda_N\) is realized in \(F_T\) as the class of
\(P_T(n/N)\).  Let \(E=(\mathbb R^4)^*\), with its Euclidean inner product,
and let \(\mathbf e_1,\ldots,\mathbf e_4\) be the standard coordinate
columns in the lattice-coordinate space, also used as the coefficient
columns of its dual coordinate basis.  Define the physical coframe
\begin{equation}
 \vartheta_T^j=P_T^{-\mathsf T}\mathbf e_j\in E,
 \qquad 1\le j\le4.
 \label{cusp-gauge-finite-physical-coframe}
\end{equation}
Thus \(\vartheta_T^j=\mathrm dy^j\) in physical \(x\)-coordinates, and the
Gram matrix of this coframe is \(G_T^{-1}\).

For \(0\le p\le4\), define
\begin{equation}
 C_N^p(T;\mathfrak g_{\mathbb C})
 =\operatorname{Map}\left(
   \Lambda_N,\bigwedge^pE_{\mathbb C}\otimes\mathfrak g_{\mathbb C}
  \right).
 \label{cusp-gauge-finite-cochain-spaces}
\end{equation}
These spaces are finite-dimensional.  Every finite differential, adjoint,
Hodge operator, projection, and reduced Hessian below therefore acts on its
entire displayed finite-dimensional source space; no additional operator
domain is implicit.
The coefficient of
\(\vartheta_T^{i_1}\wedge\cdots\wedge\vartheta_T^{i_p}\) at \(n\) is the
coefficient assigned to the positively oriented coordinate \(p\)-cell
based at \(n\).  Let
\[
 (\mathsf T_ja)(n)=a(n+\mathbf e_j)
\]
and define the forward and backward differences by
\begin{equation}
 D_{N,j}^+=N(\mathsf T_j-\operatorname{id}),\qquad
 (D_{N,j}^-a)(n)
 =N\bigl(a(n)-a(n-\mathbf e_j)\bigr).
 \label{cusp-gauge-finite-differences}
\end{equation}
The finite exterior derivative is
\begin{equation}
 \mathrm d_Na
 =\sum_{j=1}^4
   \varepsilon(\vartheta_T^j)D_{N,j}^+a.
 \label{cusp-gauge-finite-exterior-derivative}
\end{equation}

\begin{lemma}[Cochain property and induced exterior metric]
\label{cusp-gauge-finite-complex-metric}
The maps in
\begin{equation}
 0\longrightarrow C_N^0
 \mathrel{\mathop{\longrightarrow}^{\mathrm d_N}} C_N^1
 \mathrel{\mathop{\longrightarrow}^{\mathrm d_N}} C_N^2
 \mathrel{\mathop{\longrightarrow}^{\mathrm d_N}} C_N^3
 \mathrel{\mathop{\longrightarrow}^{\mathrm d_N}} C_N^4
 \longrightarrow0
 \label{cusp-gauge-finite-cochain-complex}
\end{equation}
form a cochain complex.  Equip \(C_N^p\) with
\begin{equation}
 \langle a,b\rangle_{N,p,T}
 =\frac{D_T}{N^4}\sum_{n\in\Lambda_N}
 \langle a(n),b(n)\rangle_{\wedge^pE\otimes\mathfrak g_{\mathbb C}}.
 \label{cusp-gauge-finite-inner-product}
\end{equation}
For increasing multi-indices
\(I=(i_1,\ldots,i_p)\) and \(J=(j_1,\ldots,j_p)\), the exterior metric is
exactly
\begin{equation}
 \left\langle
  \vartheta_T^{i_1}\wedge\cdots\wedge\vartheta_T^{i_p},
  \vartheta_T^{j_1}\wedge\cdots\wedge\vartheta_T^{j_p}
 \right\rangle
 =\det\left((G_T^{-1})_{i_a j_b}\right)_{a,b=1}^p.
 \label{cusp-gauge-finite-exterior-metric}
\end{equation}
\end{lemma}

\begin{proof}
The differences \(D_{N,j}^+\) commute, whereas the exterior
multiplications \(\varepsilon(\vartheta_T^j)\) anticommute.  Therefore
\[
 \mathrm d_N^2
 =\sum_{j,\ell}
 \varepsilon(\vartheta_T^j)\varepsilon(\vartheta_T^\ell)
 D_{N,j}^+D_{N,\ell}^+=0,
\]
which proves that \cref{cusp-gauge-finite-cochain-complex} is a complex.
The factor \(D_T/N^4\) in
\cref{cusp-gauge-finite-inner-product} is the exact Riemann-sum weight:
\(D_T\) is the covolume of \(P_T\mathbb Z^4\).  The Euclidean inner
product of \(\vartheta_T^i=P_T^{-\mathsf T}\mathbf e_i\) and
\(\vartheta_T^j=P_T^{-\mathsf T}\mathbf e_j\) is
\((G_T^{-1})_{ij}\).  The induced inner product on an exterior power is
the determinant of the matrix of pairings, proving
\cref{cusp-gauge-finite-exterior-metric}.
\end{proof}

\begin{remark}[Integrated cellular cochains give the same operator]
\label{cusp-gauge-finite-integrated-normalization}
The convention above stores form coefficients.  If \(c_I(n)\) instead
stores the integral over a \(p\)-cell, put \(a_I(n)=N^pc_I(n)\).  The
unscaled cellular coboundary \(\mathsf T_j-\operatorname{id}\), together
with the inner product
\[
 \langle c,c'\rangle^{\mathrm{int}}_{N,p,T}
 =D_TN^{2p-4}\sum_{n\in\Lambda_N}
 \langle c(n),c'(n)\rangle_{\wedge^pE\otimes\mathfrak g_{\mathbb C}},
\]
is carried isometrically to
\cref{cusp-gauge-finite-exterior-derivative,cusp-gauge-finite-inner-product}.
Thus the two conventions give identical Hessian symbols and eigenvalues.
\end{remark}

\begin{proposition}[Adjoint and Fourier wedge--contraction symbols]
\label{cusp-gauge-finite-symbols}
The adjoint of \(\mathrm d_N:C_N^p\to C_N^{p+1}\) is
\begin{equation}
 \mathrm d_N^*
 =-\sum_{j=1}^4
 \iota_{(\vartheta_T^j)^\sharp}D_{N,j}^-.
 \label{cusp-gauge-finite-adjoint}
\end{equation}
For \(k\in\Lambda_N\), define
\begin{align}
 \chi_{N,k}(n)&=\exp(2\pi i k\cdot n/N),\notag\\
 \mathfrak d_N(k)_j&=N\left(e^{2\pi i k_j/N}-1\right),
 \label{cusp-gauge-finite-coordinate-symbol}\\
 \zeta_{N,T}(k)&=\sum_{j=1}^4\mathfrak d_N(k)_j\vartheta_T^j
 =P_T^{-\mathsf T}\mathfrak d_N(k)\in E_{\mathbb C}.
 \label{cusp-gauge-finite-physical-symbol}
\end{align}
For every constant exterior amplitude \(\alpha\),
\begin{align}
 \mathrm d_N(\chi_{N,k}\alpha)
 &=\chi_{N,k}\varepsilon(\zeta_{N,T}(k))\alpha,
 \label{cusp-gauge-finite-d-symbol}\\
 \mathrm d_N^*(\chi_{N,k}\alpha)
 &=\chi_{N,k}
   \iota_{\overline{\zeta_{N,T}(k)}^\sharp}\alpha,
 \label{cusp-gauge-finite-dstar-symbol}\\
 (\mathrm d_N\mathrm d_N^*+\mathrm d_N^*\mathrm d_N)
 (\chi_{N,k}\alpha)
 &=\lvert\zeta_{N,T}(k)\rvert^2\chi_{N,k}\alpha.
 \label{cusp-gauge-finite-hodge-symbol}
\end{align}
Consequently the harmonic \(p\)-cochains are exactly the constant
cochains and have complex dimension \(\binom4p d_G\).
\end{proposition}

\begin{proof}
Translation is unitary for
\cref{cusp-gauge-finite-inner-product}; hence
\[
 (D_{N,j}^+)^*
 =N(\mathsf T_j^{-1}-\operatorname{id})=-D_{N,j}^-.
\]
For the induced exterior metric, the adjoint of exterior multiplication
by \(\vartheta_T^j\) is contraction by
\((\vartheta_T^j)^\sharp\).  This proves
\cref{cusp-gauge-finite-adjoint}.  On \(\chi_{N,k}\), the forward
difference has eigenvalue \(\mathfrak d_N(k)_j\), whereas \(-D_{N,j}^-\) has
eigenvalue \(\overline{\mathfrak d_N(k)_j}\).  Summing proves
\cref{cusp-gauge-finite-d-symbol,cusp-gauge-finite-dstar-symbol}.  The
Hermitian exterior-algebra identity
\[
 \iota_{\overline\zeta^\sharp}\varepsilon(\zeta)
 +\varepsilon(\zeta)\iota_{\overline\zeta^\sharp}
 =\lvert\zeta\rvert^2\operatorname{id}
\]
proves \cref{cusp-gauge-finite-hodge-symbol}.  Finally, \(\mathfrak d_N(k)=0\)
exactly when \(k=0\) in \(\Lambda_N\), and \(P_T\) is invertible.  Thus the
Hodge symbol vanishes exactly on the zero Fourier mode, where the entire
constant exterior amplitude is harmonic.
\end{proof}

Define the finite reduced Coulomb slice by
\begin{equation}
 \mathscr C_{N,T}
 =\ker(\mathrm d_N^*:C_N^1\to C_N^0)
  \cap(\ker\Delta_{N,T}^{(1)})^\perp,
 \qquad
 \Delta_{N,T}^{(1)}
 =\mathrm d_N\mathrm d_N^*+\mathrm d_N^*\mathrm d_N.
 \label{cusp-gauge-finite-reduced-coulomb}
\end{equation}
Equivalently, one deletes the zero Fourier mode and, for each \(k\ne0\),
restricts the amplitude to
\begin{equation}
 \mathcal T_{N,T,k}
 =\ker\iota_{\overline{\zeta_{N,T}(k)}^\sharp}
 \subset E_{\mathbb C}.
 \label{cusp-gauge-finite-transverse-space}
\end{equation}
The exact orthogonal projection is
\begin{equation}
 \Pi_{N,T,k}\alpha
 =\alpha-\zeta_{N,T}(k)
 \frac{
  \iota_{\overline{\zeta_{N,T}(k)}^\sharp}\alpha
 }{\lvert\zeta_{N,T}(k)\rvert^2}.
 \label{cusp-gauge-finite-coulomb-projection}
\end{equation}
Thus the projection uses the physical covector
\(P_T^{-\mathsf T}\mathfrak d_N(k)\), including every skew-metric cross term.

The finite quadratic gauge form and its associated operator are
\begin{equation}
 Q_{N,T}(a)
 =\frac1{2g_{\mathrm{YM}}^2}
   \lVert\mathrm d_Na\rVert_{N,2,T}^2,
 \qquad
 \mathcal L_{N,T}
 =g_{\mathrm{YM}}^{-2}\mathrm d_N^*\mathrm d_N.
 \label{cusp-gauge-finite-quadratic-form}
\end{equation}
This is the exact linearized, Lie-algebra-valued cubical discretization of
\cref{cusp-gauge-continuum-hessian-form}; it is not a nonlinear
compact-link Wilson action.

\begin{corollary}[Exact finite Coulomb spectrum]
\label{cusp-gauge-finite-spectrum}
On the \(k\ne0\) Fourier fibre of \(\mathscr C_{N,T}\),
\(\mathcal L_{N,T}\) has the single eigenvalue
\begin{equation}
 \lambda_{N,T}^{\mathrm{YM}}(k)
 =\frac1{g_{\mathrm{YM}}^2}
  \lvert P_T^{-\mathsf T}\mathfrak d_N(k)\rvert^2
 =\frac1{g_{\mathrm{YM}}^2}
  \mathfrak d_N(k)^*G_T^{-1}\mathfrak d_N(k)
 \label{cusp-gauge-finite-mode-eigenvalue}
\end{equation}
with complex multiplicity \(3d_G\).
\end{corollary}

\begin{proof}
On \(\mathcal T_{N,T,k}\), the contraction term in
\[
 \iota_{\overline\zeta^\sharp}\varepsilon(\zeta)\alpha
 =\lvert\zeta\rvert^2\alpha
  -\varepsilon(\zeta)\iota_{\overline\zeta^\sharp}\alpha
\]
vanishes.  This proves the eigenvalue.  Since \(\zeta_{N,T}(k)\ne0\), the
contraction has rank one on \(E_{\mathbb C}\), and its kernel has dimension
three.  Tensoring with \(\mathfrak g_{\mathbb C}\) gives the stated
multiplicity.
\end{proof}

\subsection{Exact coordinate-mode eigenvalues and all mesh regimes}
\label{cusp-gauge-finite-limits-subsection}

Set
\begin{equation}
 \rho_N=4N^2\sin^2\left(\frac\pi N\right).
 \label{cusp-gauge-finite-rho}
\end{equation}
Because \(\mathfrak d_N(\mathbf e_\gamma)\) and \(\mathfrak d_N(\mathbf e_u)\) are the common
scalar \(N(e^{2\pi i/N}-1)\) times the corresponding coordinate vector,
\cref{cusp-gauge-finite-mode-eigenvalue,cusp-gauge-short-mode-lengths}
give the exact eigenvalues
\begin{align}
 \lambda_{N,T}^{\mathrm{YM}}(\mathbf e_\gamma)
 &=\frac{\rho_N}{g_{\mathrm{YM}}^2}
   \frac{L_T^2+m_T^2}{D_T^2},
 \label{cusp-gauge-finite-gamma-eigenvalue}\\
 \lambda_{N,T}^{\mathrm{YM}}(\mathbf e_u)
 &=\frac{\rho_N}{g_{\mathrm{YM}}^2}
   \frac{q_T^2+36m_T^2}{D_T^2}.
 \label{cusp-gauge-finite-u-eigenvalue}
\end{align}
Each individual complex Fourier fibre contributes multiplicity \(3d_G\).
For real cochains, the pair
\(\{\mathbf e_j,-\mathbf e_j\}\) contributes \(6d_G\) dimensions when
\(N>2\).  When \(N=2\), these two labels coincide, so its real contribution
is \(3d_G\).

\begin{theorem}[Finite gap and simultaneous cusp--mesh limits]
\label{cusp-gauge-finite-gap-limits}
There is \(T_1\), independent of \(N\ge2\), such that for \(T\ge T_1\) the
first eigenvalue of \(\mathcal L_{N,T}\) on \(\mathscr C_{N,T}\) is
\begin{equation}
 \gamma_{N,T}
 =\min\left\{
   \lambda_{N,T}^{\mathrm{YM}}(\mathbf e_\gamma),
   \lambda_{N,T}^{\mathrm{YM}}(\mathbf e_u)
 \right\}.
 \label{cusp-gauge-finite-gap-exact}
\end{equation}
Uniformly for \(N\ge2\),
\begin{equation}
 \gamma_{N,T}
 =\frac{\rho_N}{g_{\mathrm{YM}}^2T^2}
 \left(
  1+\frac{-2H+2\min\{B,0\}}T+O(T^{-2})
 \right).
 \label{cusp-gauge-finite-gap-expansion}
\end{equation}
Let
\[
 I_{N,T}=\{j\in\{\gamma,u\}:
 \lambda_{N,T}^{\mathrm{YM}}(\mathbf e_j)=\gamma_{N,T}\},
 \qquad
 \eta_N=\begin{cases}1,&N=2,\\2,&N>2.\end{cases}
\]
The real multiplicity is exactly
\begin{equation}
 3d_G\eta_N\,\#I_{N,T}.
 \label{cusp-gauge-finite-gap-multiplicity}
\end{equation}
The limiting conclusions are:
\begin{enumerate}[label=(\roman*)]
\item For fixed \(N\ge2\),
\begin{equation}
 T^2\gamma_{N,T}
 \longrightarrow\frac{\rho_N}{g_{\mathrm{YM}}^2},
 \qquad \gamma_{N,T}\longrightarrow0.
 \label{cusp-gauge-finite-fixed-N}
\end{equation}
\item For every integer-valued \(N(T)\ge2\), with no growth assumption,
there are constants \(0<c<C<\infty\), independent of \(N(T)\), such that
\begin{equation}
 \frac{c}{T^2}\le\gamma_{N(T),T}\le\frac{C}{T^2}
 \label{cusp-gauge-finite-arbitrary-N}
\end{equation}
for all large \(T\).  Hence \(\gamma_{N(T),T}\to0\), although the scaled
gaps need not converge when \(N(T)\) oscillates.
\item If \(N(T)\to\infty\), then
\begin{equation}
 \gamma_{N(T),T}
 =\frac{4\pi^2}{g_{\mathrm{YM}}^2T^2}
  \left(1+O(T^{-1})+O(N(T)^{-2})\right).
 \label{cusp-gauge-finite-growing-N}
\end{equation}
\item The largest physical elementary edge has length
\begin{equation}
 \max_{1\le j\le4}
 \frac{\lVert P_T\mathbf e_j\rVert}{N(T)}
 =\frac{T}{N(T)}(1+O(T^{-1})).
 \label{cusp-gauge-finite-mesh-size}
\end{equation}
Thus the varying continuum torus is resolved if and only if
\(N(T)/T\to\infty\).  This refinement condition is not needed for the
algebraic soft-mode conclusion in (ii).
\end{enumerate}
\end{theorem}

\begin{proof}
Let
\[
 \sigma_N=2N\sin(\pi/N),\qquad \rho_N=\sigma_N^2.
\]
The least nonzero modulus of a coordinate of \(\mathfrak d_N(k)\) is \(\sigma_N\),
attained precisely at frequency \(\pm1\), with the two signs equal for
\(N=2\).  If \(\mathfrak d_N(k)\) is supported in its first two components but
\(k\notin\{\pm\mathbf e_\gamma,\pm\mathbf e_u\}\), then
\begin{equation}
 |\mathfrak d_N(k)_1|^2+|\mathfrak d_N(k)_2|^2\ge2\sigma_N^2.
 \label{cusp-gauge-finite-next-frequency}
\end{equation}
Indeed, either both coordinates are nonzero, or the single nonzero
frequency is not \(\pm1\).  In the latter case
\[
 \frac{\sin(2\pi/N)}{\sin(\pi/N)}
 =2\cos(\pi/N)\ge\sqrt2
\]
for \(N\ge4\), and the cases \(N=2,3\) are immediate.

For \(z=(z_1,z_2,0,0)^{\mathsf t}\in\mathbb C^4\), the exact inverse
transpose in \cref{cusp-spec-exact-dual-lattice} gives
\begin{equation}
 P_T^{-\mathsf T}z
 =\left(
 0,\frac{m_Tz_1+q_Tz_2}{D_T},
 0,\frac{-L_Tz_1+6m_Tz_2}{D_T}
 \right).
 \label{cusp-gauge-finite-soft-block}
\end{equation}
The scaled \(2\times2\) matrix
\[
 \frac{T}{D_T}
 \begin{pmatrix}m_T&q_T\\-L_T&6m_T\end{pmatrix}
 =
 \begin{pmatrix}0&1\\-1&0\end{pmatrix}
 +O_{\mathrm{op}}(T^{-1})
\]
has all singular values \(1+O(T^{-1})\), also on complex vectors.  Thus
the factor \(\sqrt2\) in
\cref{cusp-gauge-finite-next-frequency} separates every other soft-plane
frequency from the coordinate frequencies uniformly in \(N\) for large
\(T\).

For a general \(z\in\mathbb C^4\), the full inverse transpose in
\cref{cusp-spec-exact-dual-lattice} has first and third components
\(z_3,z_4\).  Hence a finite frequency with nonzero third or fourth
coordinate satisfies
\[
 \lVert P_T^{-\mathsf T}\mathfrak d_N(k)\rVert^2
 \ge |\mathfrak d_N(k)_3|^2+|\mathfrak d_N(k)_4|^2
 \ge\sigma_N^2,
\]
whereas either coordinate candidate is
\(O(\sigma_N^2/T^2)\).  This proves the exact minimizing list and
\cref{cusp-gauge-finite-gap-exact}; the reality discussion before the
theorem proves \cref{cusp-gauge-finite-gap-multiplicity}.

Multiplying the two expansions in
\cref{cusp-spec-shortest-modes} by
\(\rho_N/g_{\mathrm{YM}}^2\) proves
\cref{cusp-gauge-finite-gap-expansion}.  The elementary sine bounds
\(2x/\pi\le\sin x\le x\) for \(0<x\le\pi/2\) give
\begin{equation}
 16\le\rho_N\le4\pi^2.
 \label{cusp-gauge-finite-rho-bounds}
\end{equation}
Together with \cref{cusp-gauge-finite-gap-expansion}, these bounds prove
(i) and (ii).  The Taylor expansion
\(\rho_N=4\pi^2(1+O(N^{-2}))\) proves (iii).

The first two columns of \(P_T\), displayed in
\cref{cusp-spec-exact-dual-lattice}, have Euclidean lengths
\(T(1+O(T^{-1}))\), while its last two columns have length one.  Division
by \(N(T)\) proves \cref{cusp-gauge-finite-mesh-size} and (iv).
\end{proof}

\begin{remark}[Every fixed reduced model remains positively gapped]
\label{cusp-gauge-pointwise-versus-uniform}
For each fixed \(T\), the reduced continuum operator has a positive first
eigenvalue.  For each fixed pair \((T,N)\), the reduced finite operator
also has a positive first eigenvalue.  What fails is a positive lower
bound uniform as \(T\to\infty\).  Before reduction, gauge-longitudinal
directions and constant harmonic one-forms are separate Hessian-zero
directions.
\end{remark}

\subsection{Continuum Coulomb reduction and exact spectrum}
\label{cusp-gauge-continuum-spectrum-subsection}

The infinitesimal gauge directions
\(\mathrm d\Omega^0(F_T;\mathfrak g)\) and the harmonic one-forms are
Hessian-zero directions.  Coulomb gauge removes the former but not the
latter.  The reduced Coulomb space is therefore
\begin{equation}
 \mathscr C_T=
 \left\{
 a\in L^2\Omega^1(F_T;\mathfrak g):
 \mathrm d^*a=0\ \text{distributionally},\
 a\perp\mathcal H^1(F_T;\mathfrak g)
 \right\}.
 \label{cusp-gauge-continuum-reduced-coulomb}
\end{equation}
The subspace \(\mathscr C_T\) reduces \(\mathcal L_T\).  On it the closed
Hessian form has domain
\(H^1\Omega^1(F_T;\mathfrak g)\cap\mathscr C_T\), and the restricted
self-adjoint operator has domain
\(H^2\Omega^1(F_T;\mathfrak g)\cap\mathscr C_T\).  Indeed, on the
distributionally coclosed slice the graph norm of \(\mathrm d_1\), together
with the \(L^2\)-norm, is the Hodge graph norm
\[
 \lVert a\rVert_{L^2}^2+
 \lVert\mathrm da\rVert_{L^2}^2+
 \lVert\mathrm d^*a\rVert_{L^2}^2.
\]
The Gaffney estimate on the compact flat torus identifies its completion
with \(H^1\), and elliptic regularity for the Hodge Laplacian identifies the
operator domain with \(H^2\).  Orthogonality to the finite-dimensional
harmonic space is preserved by the Fourier multiplier.

For \(k\in\mathbb Z^4\), write
\[
 \xi_{T,k}=P_T^{-\mathsf T}k,\qquad
 \chi_{T,k}([x])
 =\exp\bigl(2\pi i\langle\xi_{T,k},x\rangle\bigr).
\]
For a covector \(\xi\), let \(\varepsilon(\xi)\) denote exterior
multiplication by \(\xi\), and use the Euclidean metric to define
\(\xi^\sharp\) and contraction \(\iota_{\xi^\sharp}\).

\begin{proposition}[Continuum Fourier symbol, slice, and multiplicity]
\label{cusp-gauge-continuum-symbol}
For
\(\alpha\in(\mathbb R^4)^*_{\mathbb C}\otimes\mathfrak g_{\mathbb C}\),
\begin{align}
 \mathrm d(\chi_{T,k}\alpha)
 &=2\pi i\,\chi_{T,k}\varepsilon(\xi_{T,k})\alpha,
 \label{cusp-gauge-continuum-d-symbol}\\
 \mathrm d^*(\chi_{T,k}\alpha)
 &=-2\pi i\,\chi_{T,k}\iota_{\xi_{T,k}^\sharp}\alpha,
 \label{cusp-gauge-continuum-dstar-symbol}\\
 \mathrm d^*\mathrm d(\chi_{T,k}\alpha)
 &=4\pi^2\chi_{T,k}
 \left(
  \lvert\xi_{T,k}\rvert^2\alpha
  -\varepsilon(\xi_{T,k})
   \iota_{\xi_{T,k}^\sharp}\alpha
 \right).
 \label{cusp-gauge-continuum-dstard-symbol}
\end{align}
For \(k\ne0\), the Coulomb polarization space is
\begin{equation}
 \mathcal T_{T,k}
 =\ker\iota_{\xi_{T,k}^\sharp}
 \subset(\mathbb R^4)^*_{\mathbb C}.
 \label{cusp-gauge-continuum-transverse-space}
\end{equation}
It has complex dimension three.  On
\(\mathcal T_{T,k}\otimes\mathfrak g_{\mathbb C}\), the Hessian
eigenvalue is
\begin{equation}
 \lambda_T^{\mathrm{YM}}(k)
 =\frac{4\pi^2}{g_{\mathrm{YM}}^2}
 \lvert P_T^{-\mathsf T}k\rvert^2
 \label{cusp-gauge-continuum-mode-eigenvalue}
\end{equation}
with complex multiplicity \(3d_G\) on that signed Fourier fibre.  If
\[
 K_T(\lambda)
 =\{k\in\mathbb Z^4\setminus\{0\}:
       \lambda_T^{\mathrm{YM}}(k)=\lambda\},
\]
the total complex multiplicity, and also the dimension of the real
eigenspace, is
\begin{equation}
 3d_G\,\#K_T(\lambda).
 \label{cusp-gauge-continuum-shell-multiplicity}
\end{equation}
Equivalently, a conjugate pair \(\{k,-k\}\) contributes \(6d_G\) real
dimensions.
\end{proposition}

\begin{proof}
Differentiating the character proves
\cref{cusp-gauge-continuum-d-symbol}.  On the Euclidean cover,
\[
 \mathrm d^*
 =-\sum_{j=1}^4
 \iota_{\partial_{x_j}}\partial_{x_j},
\]
which proves \cref{cusp-gauge-continuum-dstar-symbol}.  The exterior
algebra identity
\[
 \iota_{\xi^\sharp}\varepsilon(\xi)
 +\varepsilon(\xi)\iota_{\xi^\sharp}
 =\lvert\xi\rvert^2\operatorname{id}
\]
proves \cref{cusp-gauge-continuum-dstard-symbol}.
For \(k\ne0\), the invertibility of \(P_T\) gives \(\xi_{T,k}\ne0\).
Contraction by \(\xi_{T,k}^\sharp\) has rank one on one-forms, so its
kernel has dimension three.  On that kernel the second term of
\cref{cusp-gauge-continuum-dstard-symbol} vanishes, proving the eigenvalue
and its per-fibre multiplicity.  The zero Fourier mode consists of the
constant harmonic one-forms and has already been removed in
\cref{cusp-gauge-continuum-reduced-coulomb}.  Fourier completeness gives
all eigenvectors.  Complex conjugation exchanges \(k\) with \(-k\);
the reality condition \(\alpha_{-k}=\overline{\alpha_k}\) turns each
conjugate pair of complex \(3d_G\)-dimensional spaces into a real
\(6d_G\)-dimensional space.  This proves
\cref{cusp-gauge-continuum-shell-multiplicity}.
\end{proof}

\begin{remark}[The Coulomb projection uses the physical covector metric]
\label{cusp-gauge-continuum-projection}
In lattice-coordinate one-form coefficients \(a_y\), the mode \(k\) obeys
\begin{equation}
 k^{\mathsf T}G_T^{-1}a_y=0,\qquad
 \Pi_{T,k}a_y
 =a_y-k\frac{k^{\mathsf T}G_T^{-1}a_y}
 {k^{\mathsf T}G_T^{-1}k}.
 \label{cusp-gauge-continuum-coordinate-projection}
\end{equation}
After conversion to physical \(x\)-coefficients, this is Euclidean
projection perpendicular to \(P_T^{-\mathsf T}k\).  It is generally
incorrect to impose Euclidean coefficient orthogonality to \(k\), or to
delete one coordinate merely because \(k\) is a coordinate lattice vector.
\end{remark}

\begin{theorem}[Exact reduced continuum Hessian gap]
\label{cusp-gauge-continuum-gap}
For all sufficiently large \(T\), the first eigenvalue of
\(\mathcal L_T\) on \(\mathscr C_T\) is
\begin{equation}
 \gamma_T
 =\frac{4\pi^2}{g_{\mathrm{YM}}^2}
 \min\left\{
   \frac{L_T^2+m_T^2}{D_T^2},
   \frac{q_T^2+36m_T^2}{D_T^2}
 \right\}.
 \label{cusp-gauge-continuum-gap-exact}
\end{equation}
It has the radial expansion
\begin{equation}
 \gamma_T
 =\frac{4\pi^2}{g_{\mathrm{YM}}^2T^2}
 \left(
  1+\frac{-2H+2\min\{B,0\}}T+O(T^{-2})
 \right).
 \label{cusp-gauge-continuum-gap-expansion}
\end{equation}
Let
\[
 I_T=\{j\in\{\gamma,u\}:a_{j,T}
             =\min\{a_{\gamma,T},a_{u,T}\}\}.
\]
The real multiplicity of \(\gamma_T\) is exactly
\begin{equation}
 6d_G\,\#I_T.
 \label{cusp-gauge-continuum-gap-multiplicity}
\end{equation}
\end{theorem}

\begin{proof}
By \cref{cusp-gauge-continuum-symbol}, minimizing the reduced Hessian
eigenvalue is exactly the shortest-nonzero-dual-vector problem.  The
shortest-mode theorem
\cref{cusp-spec-shortest-modes} says that, for all large \(T\), the only
possible minimizers are the two signed pairs with exact squared lengths
given in \cref{cusp-gauge-short-mode-lengths}.  This proves
\cref{cusp-gauge-continuum-gap-exact}.  The expansions already proved in
\cref{cusp-spec-shortest-modes} give relative corrections
\(2(B-H)\) and \(-2H\); their minimum is
\(-2H+2\min\{B,0\}\), proving
\cref{cusp-gauge-continuum-gap-expansion}.  Each minimizing signed pair
contributes \(6d_G\) real dimensions by
\cref{cusp-gauge-continuum-symbol}, which proves the multiplicity formula.
\end{proof}

\subsection{The physical boundary and the exact no-gap bridge}
\label{cusp-gauge-physical-boundary-subsection}

The operators proved above are a classical Hessian and a finite quadratic
cochain approximation on the internal flat manifolds \(F_T\).  They do not
construct a quantum gauge theory on standard spacetime \(\mathbb R^4\), and
they do not determine the mass spectrum of such a theory.  The precise
boundaries are as follows.
\begin{enumerate}[label=(\arabic*)]
\item The unscaled cusp retains two compact geometric directions.  The
period parameter \(T\) has not been identified with a spacetime box size,
a lattice spacing, a correlation length, or a renormalization scale.
\item The operator \(g_{\mathrm{YM}}^{-2}\mathrm d^*\mathrm d\) is the
quadratic classical fluctuation operator at the product connection.  It
contains no nonperturbative interacting measure, dimensional
transmutation, or confinement argument.
\item The spaces \(C_N^p(T;\mathfrak g_{\mathbb C})\) are linear
Lie-algebra cochains.  They are not compact-group link variables, and
\cref{cusp-gauge-finite-quadratic-form} is not a Wilson plaquette action.
\item No Euclidean time direction or reflection, reflection-positive
measure, transfer matrix, Osterwalder--Schrader reconstruction, Gauss-law
physical quotient, local observable algebra, thermodynamic limit, or
continuum limit in fixed physical units has been constructed from these
operators.
\item An eigenvalue of a four-dimensional elliptic Hessian has units of
squared inverse length.  Without a time reconstruction and a scale
conversion, it is not an eigenvalue of a physical Hamiltonian.  In
particular, a dimensionless transfer-matrix gap can behave as
\(a_tm\to0\) while the physical mass \(m>0\) remains fixed.
\end{enumerate}

The exact extra statement that would transfer the auxiliary soft sequence
to one already constructed physical Hamiltonian is the following.  It
must use the full zero-energy sector rather than only a chosen vacuum
vector.

\begin{proposition}[Typed state-and-energy no-gap bridge]
\label{cusp-gauge-no-gap-bridge}
Let \(\mathcal H_{\mathrm{phys}}\) be a Hilbert space and let \(H\ge0\) be
self-adjoint on it.  Its closed energy form is
\begin{equation}
 \mathfrak h[\psi]=\lVert H^{1/2}\psi\rVert^2,\qquad
 \mathcal D(\mathfrak h)=\mathcal D(H^{1/2}).
 \label{cusp-gauge-physical-energy-form}
\end{equation}
Let
\begin{equation}
 P_0=\mathbf 1_{\{0\}}(H),\qquad
 \mathcal V=\ker H=P_0\mathcal H_{\mathrm{phys}},\qquad
 \mathcal H_{\mathrm{exc}}
 =(\operatorname{id}-P_0)\mathcal H_{\mathrm{phys}},
 \label{cusp-gauge-vacuum-excitation-spaces}
\end{equation}
and fix a normalized vacuum vector \(\Omega\in\mathcal V\).  Then
\begin{equation}
 \mathcal H_{\mathrm{phys}}
 =\mathcal V\oplus\mathcal H_{\mathrm{exc}}
 \label{cusp-gauge-reducing-decomposition}
\end{equation}
reduces \(H\).  Denote its excitation restriction by
\[
 H_{\mathrm{exc}}
 =H|_{\mathcal D(H)\cap\mathcal H_{\mathrm{exc}}}.
\]

For every \(n\), let \(\mathcal K_n\) be an auxiliary Hilbert space with a
closed nonnegative quadratic form \(\mathfrak q_n\).  Let
\(u_n\in\mathcal D(\mathfrak q_n)\) satisfy
\begin{equation}
 \lVert u_n\rVert_{\mathcal K_n}=1,\qquad
 \mathfrak q_n[u_n]\longrightarrow0.
 \label{cusp-gauge-auxiliary-soft-states}
\end{equation}
Suppose there are linear maps
\begin{equation}
 J_n:\mathcal D(\mathfrak q_n)\longrightarrow\mathcal D(\mathfrak h),
 \qquad
 \psi_n=(\operatorname{id}-P_0)J_nu_n,
 \label{cusp-gauge-bridge-maps}
\end{equation}
and numbers \(c>0\), \(C_n\ge0\), and \(\varepsilon_n\ge0\) such that
\begin{align}
 \lVert\psi_n\rVert&\ge c,
 \label{cusp-gauge-bridge-norm}\\
 \mathfrak h[\psi_n]
 &\le C_n\mathfrak q_n[u_n]+\varepsilon_n,
 \qquad
 C_n\mathfrak q_n[u_n]+\varepsilon_n\longrightarrow0.
 \label{cusp-gauge-bridge-energy}
\end{align}
Then
\begin{equation}
 \inf\operatorname{Spec}(H_{\mathrm{exc}})=0
 \label{cusp-gauge-bridge-zero-infimum}
\end{equation}
and, for every \(\epsilon>0\),
\begin{equation}
 \mathbf 1_{(0,\epsilon)}(H)\ne0.
 \label{cusp-gauge-bridge-spectrum-near-zero}
\end{equation}
Moreover,
\begin{equation}
 0\in\operatorname{Spec}(H_{\mathrm{exc}}),\qquad
 \ker H_{\mathrm{exc}}=\{0\}.
 \label{cusp-gauge-bridge-continuous-threshold}
\end{equation}
Thus zero is a continuous spectral threshold, essential under the usual
Weyl convention, and \(H\) has no positive gap above its full zero-energy
sector.

If the vacuum is unique, so that \(\mathcal V=\mathbb C\Omega\), one may
replace the projection in \cref{cusp-gauge-bridge-maps} by the requirement
\(J_nu_n\perp\Omega\).  Without uniqueness, orthogonality only to
\(\Omega\) is insufficient.
\end{proposition}

\begin{proof}
The operator \(P_0\) is a spectral projection of \(H\), so it commutes
with the resolvent and with the spectral calculus of \(H\).  Hence
\cref{cusp-gauge-reducing-decomposition} is a reducing decomposition and
\[
 \mathcal D(\mathfrak h|_{\mathcal H_{\mathrm{exc}}})
 =\mathcal D(H^{1/2})\cap\mathcal H_{\mathrm{exc}}.
\]
The normalized vectors
\[
 v_n=\frac{\psi_n}{\lVert\psi_n\rVert}
\]
belong to this excitation form domain.  By
\cref{cusp-gauge-bridge-norm,cusp-gauge-bridge-energy},
\[
 \mathfrak h[v_n]
 \le c^{-2}
 \bigl(C_n\mathfrak q_n[u_n]+\varepsilon_n\bigr)
 \longrightarrow0.
\]
The form version of the spectral variational identity is
\[
 \inf\operatorname{Spec}(H_{\mathrm{exc}})
 =\inf_{\substack{
 0\ne v\in\mathcal D(H^{1/2})\cap\mathcal H_{\mathrm{exc}}
 }}
 \frac{\mathfrak h[v]}{\lVert v\rVert^2}.
\]
Nonnegativity and the sequence \(v_n\) prove
\cref{cusp-gauge-bridge-zero-infimum}.

If \(\mathbf 1_{(0,\epsilon)}(H)=0\) for some \(\epsilon>0\), the spectral
measure of every excitation vector would be supported in
\([\epsilon,\infty)\), and therefore
\[
 \mathfrak h[v]\ge\epsilon\lVert v\rVert^2
\]
throughout the excitation form domain.  This contradicts
\(\mathfrak h[v_n]\to0\), proving
\cref{cusp-gauge-bridge-spectrum-near-zero}.  Since the spectrum of a
self-adjoint operator is closed,
\cref{cusp-gauge-bridge-zero-infimum} puts zero in
\(\operatorname{Spec}(H_{\mathrm{exc}})\).  Its kernel is trivial by the
definition of \(\mathcal H_{\mathrm{exc}}\).  A self-adjoint operator has
no residual spectrum, and
\cref{cusp-gauge-bridge-spectrum-near-zero} makes zero an accumulation
point.  This proves
\cref{cusp-gauge-bridge-continuous-threshold} and the essential-threshold
assertion.

If \(\mathcal V=\mathbb C\Omega\), orthogonality to \(\Omega\) is exactly
membership in \(\mathcal H_{\mathrm{exc}}\).  If the kernel is larger, a
zero-energy vector can be orthogonal to the selected \(\Omega\), so that
orthogonality alone does not produce an excitation vector.  This proves
the final assertion.
\end{proof}

\begin{remark}[The present geometry supplies only the auxiliary side]
\label{cusp-gauge-missing-physical-bridge}
Choose any \(T_n\to\infty\), any \(N_n\ge2\), and a unit real transverse
eigenvector \(u_n\in\mathscr C_{N_n,T_n}\) at a minimizing coordinate
mode.  With
\[
 \mathcal K_n=\mathscr C_{N_n,T_n},\qquad
 \mathfrak q_n[u]
 =g_{\mathrm{YM}}^{-2}\lVert\mathrm d_{N_n}u\rVert^2,
\]
\cref{cusp-gauge-finite-gap-limits} gives
\(\mathfrak q_n[u_n]=\gamma_{N_n,T_n}\to0\).  Thus the auxiliary
soft-state data in \cref{cusp-gauge-auxiliary-soft-states} are present.

What is absent is one already constructed, gauge-invariant,
positive-metric physical Hilbert space for four-dimensional quantum
Yang--Mills theory, its Hamiltonian \(H\), and maps \(J_n\) satisfying
\cref{cusp-gauge-bridge-norm,cusp-gauge-bridge-energy} with constants
controlled in physical units after gauge reduction, reconstruction,
renormalization, and all limiting operations.  Consequently
\cref{cusp-gauge-no-gap-bridge} cannot be applied to infer a quantum mass
statement from the cusp geometry.  Even if such maps were supplied, the
proposition would establish only spectrum arbitrarily close to zero for
that specified \(H\); it would not by itself identify a one-particle state,
establish scattering, or construct the required quantum gauge theory.
\end{remark}

\section{Cubical holonomy and finite gauge systems}
\label{cusp-link-section}

This section records three constructions on the regular cusp fibres while
keeping two independent integers separate throughout:
\(M_{\mathrm{cov}}\) is the index used to enlarge two periods of a fibre,
whereas \(N_{\mathrm{mesh}}\) is the number of cubical subdivisions in each
coordinate direction.  In particular, no mesh refinement is identified with
a change of covering space.

Let
\begin{equation}
 P_T=
 \begin{pmatrix}
  6\operatorname{Re}\mu_T&\operatorname{Re}\tau_T&1&0\\
  6\operatorname{Im}\mu_T&\operatorname{Im}\tau_T&0&0\\
  \operatorname{Re}\beta_T&\operatorname{Re}\mu_T&0&1\\
  \operatorname{Im}\beta_T&\operatorname{Im}\mu_T&0&0
 \end{pmatrix},
 \qquad
 F_T=\mathbb R^4/P_T\mathbb Z^4.
 \label{cusp-link-period-matrix}
\end{equation}
Write
\begin{equation}
 L_T=\operatorname{Im}\tau_T,
 \qquad m_T=\operatorname{Im}\mu_T,
 \qquad q_T=-\operatorname{Im}\beta_T,
 \qquad D_T=L_Tq_T+6m_T^2>0.
 \label{cusp-link-radial-data}
\end{equation}
Expanding the determinant along the last two columns gives
\[
 \det P_T=L_T\operatorname{Im}\beta_T-6m_T^2=-D_T.
\]
Along the positive radial cusp the source data give
\begin{equation}
 L_T=T+O(1),\qquad q_T=T+O(1),\qquad m_T=O(1),
 \label{cusp-link-cusp-asymptotics}
\end{equation}
and all three real parts occurring in \eqref{cusp-link-period-matrix} are
bounded.  Equivalently, after ordering physical coordinates as
\((x_1,x_3;x_2,x_4)\), the first two columns have block form
\begin{equation}
 \binom{A_T}{B_T},\qquad
 B_T=
 \begin{pmatrix}6m_T&L_T\\-q_T&m_T\end{pmatrix}
 =T\begin{pmatrix}0&1\\-1&0\end{pmatrix}+C_T,
 \qquad \sup_T\lVert C_T\rVert_{\mathrm{op}}<\infty.
 \label{cusp-link-block-form}
\end{equation}

For \(M_{\mathrm{cov}}\geq1\), put
\begin{equation}
 S_{M_{\mathrm{cov}}}
 =\operatorname{diag}(1,1,M_{\mathrm{cov}},M_{\mathrm{cov}}),
 \qquad
 P_{T,M_{\mathrm{cov}}}=P_TS_{M_{\mathrm{cov}}},
 \label{cusp-link-cover-matrix}
\end{equation}
and define the degree-\(M_{\mathrm{cov}}^2\) cover
\begin{equation}
 \widetilde F_{T,M_{\mathrm{cov}}}
 =\mathbb R^4/P_{T,M_{\mathrm{cov}}}\mathbb Z^4
 \longrightarrow F_T.
 \label{cusp-link-cover-torus}
\end{equation}
Indeed,
\(S_{M_{\mathrm{cov}}}\mathbb Z^4\subset\mathbb Z^4\) and
\[
 \mathbb Z^4/S_{M_{\mathrm{cov}}}\mathbb Z^4
 \cong(\mathbb Z/M_{\mathrm{cov}}\mathbb Z)^2,
\]
so the identity on \(\mathbb R^4\) induces the asserted cover and its deck
group has order \(M_{\mathrm{cov}}^2\).
Its Euclidean quotient metric, written on the coordinate torus
\(\mathbb R^4/\mathbb Z^4\), has Gram matrix and covolume
\begin{equation}
 G_{T,M_{\mathrm{cov}}}
 =P_{T,M_{\mathrm{cov}}}^{\mathsf T}P_{T,M_{\mathrm{cov}}},
 \qquad
 \mathcal D_{T,M_{\mathrm{cov}}}
 =\sqrt{\det G_{T,M_{\mathrm{cov}}}}
 =M_{\mathrm{cov}}^2D_T.
 \label{cusp-link-cover-metric}
\end{equation}

\subsection{The transported cubical subdivision and continuum holonomy}

Fix an integer \(N_{\mathrm{mesh}}\geq2\), and let
\(K_{N_{\mathrm{mesh}}}\) be the standard cubical subdivision of
\(\mathbb R^4/\mathbb Z^4\) obtained by dividing every coordinate circle
into \(N_{\mathrm{mesh}}\) equal intervals.  Its vertex set and positively
oriented coordinate edges are
\begin{equation}
 V_{N_{\mathrm{mesh}}}
 =(\mathbb Z/N_{\mathrm{mesh}}\mathbb Z)^4,
 \qquad
 E_{N_{\mathrm{mesh}}}
 =V_{N_{\mathrm{mesh}}}\times\{1,2,3,4\}.
 \label{cusp-link-cubical-sets}
\end{equation}

\begin{proposition}[Exact cubical realization and the torsion vertices]
\label{cusp-link-cubical-realization}
The linear map \(P_{T,M_{\mathrm{cov}}}\) induces a diffeomorphism
\begin{equation}
 \Phi_{T,M_{\mathrm{cov}}}:\mathbb R^4/\mathbb Z^4
 \longrightarrow\widetilde F_{T,M_{\mathrm{cov}}},
 \qquad [y]\longmapsto[P_{T,M_{\mathrm{cov}}}y].
 \label{cusp-link-coordinate-diffeomorphism}
\end{equation}
Transporting \(K_{N_{\mathrm{mesh}}}\) by this diffeomorphism gives an
honest cubical cell decomposition of the whole torus.  On vertices its
restriction is the bijection
\begin{equation}
 \begin{aligned}
 \iota_{T,M_{\mathrm{cov}},N_{\mathrm{mesh}}}:\quad
 &(\mathbb Z/N_{\mathrm{mesh}}\mathbb Z)^4
 \xrightarrow{\ \cong\ }
 \widetilde F_{T,M_{\mathrm{cov}}}[N_{\mathrm{mesh}}],\\
 &[n]\longmapsto
 \left[\frac1{N_{\mathrm{mesh}}}P_{T,M_{\mathrm{cov}}}n\right].
 \end{aligned}
 \label{cusp-link-torsion-bijection}
\end{equation}
Here the codomain is the finite \(N_{\mathrm{mesh}}\)-torsion subgroup,
not the entire torus.  The coordinate edge \((n,j)\) is the path
\begin{equation}
 \gamma_{n,j}(t)=
 \left[
  \frac1{N_{\mathrm{mesh}}}
  P_{T,M_{\mathrm{cov}}}(n+t\mathbf e_j)
 \right],
 \qquad 0\leq t\leq1.
 \label{cusp-link-coordinate-edge}
\end{equation}
More generally, for an ordered subset
\(I=\{i_1,\ldots,i_p\}\subset\{1,2,3,4\}\), the corresponding closed
\(p\)-cell is parametrized by
\begin{equation}
 c_{n,I}(t_1,\ldots,t_p)=
 \left[
  \frac1{N_{\mathrm{mesh}}}P_{T,M_{\mathrm{cov}}}
  \left(n+\sum_{a=1}^pt_a\mathbf e_{i_a}\right)
 \right],
 \qquad (t_1,\ldots,t_p)\in[0,1]^p.
 \label{cusp-link-general-cell}
\end{equation}
\end{proposition}

\begin{proof}
If \(y-y'\in\mathbb Z^4\), then
\(P_{T,M_{\mathrm{cov}}}(y-y')\) belongs to the quotient lattice, so
\eqref{cusp-link-coordinate-diffeomorphism} is well defined.  The matrix
\(P_{T,M_{\mathrm{cov}}}\) is invertible, and its inverse induces the
inverse map on the two quotients.  Hence it is a diffeomorphism and carries
every open cell of the standard subdivision diffeomorphically onto an open
cell of the target.

Every point displayed in \eqref{cusp-link-torsion-bijection} is killed by
\(N_{\mathrm{mesh}}\).  If it is zero, then
\[
 \frac{P_{T,M_{\mathrm{cov}}}n}{N_{\mathrm{mesh}}}
 \in P_{T,M_{\mathrm{cov}}}\mathbb Z^4.
\]
Invertibility gives
\(n\in N_{\mathrm{mesh}}\mathbb Z^4\), proving injectivity.  Conversely,
if \([x]\) is killed by \(N_{\mathrm{mesh}}\), choose a lift
\(x\in\mathbb R^4\).  There is \(n\in\mathbb Z^4\) such that
\(N_{\mathrm{mesh}}x=P_{T,M_{\mathrm{cov}}}n\), and therefore
\([x]=[P_{T,M_{\mathrm{cov}}}n/N_{\mathrm{mesh}}]\).  This proves
surjectivity.  Formula \eqref{cusp-link-coordinate-edge} is simply the
transport of the corresponding coordinate edge; at a wrap-around endpoint
the extra vector \(P_{T,M_{\mathrm{cov}}}\mathbf e_j\) is a period.
\end{proof}

Let \(G\) be a compact matrix Lie group and use the trivial principal bundle over
\(\widetilde F_{T,M_{\mathrm{cov}}}\).  A connection is represented by
\(A\in\Omega^1(\widetilde F_{T,M_{\mathrm{cov}}};\mathfrak g)\).  For an
oriented path \(\gamma\), define \(U_A(\gamma)\) by
\begin{equation}
 \frac{dU}{dt}=-A_{\gamma(t)}(\dot\gamma(t))U(t),
 \qquad U(0)=1,
 \qquad U_A(\gamma)=U(1).
 \label{cusp-link-parallel-transport-ode}
\end{equation}

\begin{proposition}[Typed continuum-to-link holonomy map]
\label{cusp-link-holonomy-map}
The formula
\begin{equation}
 \operatorname{Hol}_{T,M_{\mathrm{cov}},N_{\mathrm{mesh}}}(A)_{n,j}
 =U_A(\gamma_{n,j})
 \label{cusp-link-holonomy-definition}
\end{equation}
defines a map
\begin{equation}
 \operatorname{Hol}_{T,M_{\mathrm{cov}},N_{\mathrm{mesh}}}:
 \Omega^1(\widetilde F_{T,M_{\mathrm{cov}}};\mathfrak g)
 \longrightarrow G^{E_{N_{\mathrm{mesh}}}}.
 \label{cusp-link-holonomy-typed}
\end{equation}
For the convention
\begin{equation}
 A^g=g^{-1}Ag+g^{-1}dg,
 \qquad
 g\in C^\infty(\widetilde F_{T,M_{\mathrm{cov}}},G),
 \label{cusp-link-gauge-convention}
\end{equation}
it is equivariant for evaluation of \(g\) at the vertices:
\begin{equation}
 \operatorname{Hol}(A^g)_{n,j}
 =g(\gamma_{n,j}(1))^{-1}
  \operatorname{Hol}(A)_{n,j}
  g(\gamma_{n,j}(0)).
 \label{cusp-link-holonomy-equivariance}
\end{equation}
\end{proposition}

\begin{proof}
The smooth matrix ODE \eqref{cusp-link-parallel-transport-ode} has a unique
solution on every edge, which proves that \eqref{cusp-link-holonomy-typed}
is defined.  Let \(U\) solve that ODE for \(A\), set
\(g_t=g(\gamma_{n,j}(t))\), and put
\(V(t)=g_t^{-1}U(t)g_0\).  Direct differentiation gives
\[
 V'(t)=-\bigl(g_t^{-1}A(\dot\gamma_{n,j})g_t
                 +g_t^{-1}g_t'\bigr)V(t),
 \qquad V(0)=1.
\]
This is exactly the parallel-transport equation for \(A^g\).  Uniqueness
and evaluation at \(t=1\) prove \eqref{cusp-link-holonomy-equivariance}.
Notice that the endpoint inverse in this formula is forced by the joint
conventions \eqref{cusp-link-parallel-transport-ode} and
\eqref{cusp-link-gauge-convention}.
\end{proof}

\subsection{An explicit compact-link
\texorpdfstring{\(\mathrm{SU}(2)\)}{SU(2)} flux}

The ordered plaquette convention used below is most naturally written for
the inverse-transport link variable
\begin{equation}
 \mathcal U(A)_{n,j}
 =\operatorname{Hol}_{T,M_{\mathrm{cov}},N_{\mathrm{mesh}}}(A)_{n,j}^{-1}.
 \label{cusp-link-inverse-transport-variable}
\end{equation}
Equation \eqref{cusp-link-holonomy-equivariance} then becomes
\begin{equation}
 \mathcal U(A^g)_{n,j}
 =g(\gamma_{n,j}(0))^{-1}\mathcal U(A)_{n,j}g(\gamma_{n,j}(1)).
 \label{cusp-link-inverse-transport-gauge-law}
\end{equation}
Thus the product in \eqref{cusp-link-plaquette-convention} transforms by
conjugation at its base vertex.  The symbols \(U_j(n)\) in the remainder of
this subsection use this inverse-transport convention.

Use
\begin{equation}
 \mathfrak{su}(2)=\{X\in M_2(\mathbb C):X^*=-X,\ \operatorname{tr}X=0\},
 \qquad
 \langle X,Y\rangle_{\mathfrak{su}(2)}
 =-\frac12\operatorname{tr}(XY),
 \label{cusp-link-su2-inner-product}
\end{equation}
on real elements, and fix
\begin{equation}
 H_{\mathrm{link}}=\begin{pmatrix}i&0\\0&-i\end{pmatrix}.
 \qquad
 \lVert H_{\mathrm{link}}\rVert^2=1,
 \qquad \exp(2\pi H_{\mathrm{link}})=1.
 \label{cusp-link-flux-generator}
\end{equation}
Write \(N=N_{\mathrm{mesh}}\) temporarily inside the next displayed link
formulas, and represent \(n_1,n_2\) by integers in
\(\{0,\ldots,N-1\}\).

\begin{theorem}[A periodic compact-link flux and its exact action]
\label{cusp-link-flux-theorem}
On the oriented edges of \(K_N\), define
\begin{align}
 U_1(n)&=\exp\left(-\frac{2\pi n_2}{N^2}H_{\mathrm{link}}\right),
 \label{cusp-link-flux-u1}\\
 U_2(n)&=
 \begin{cases}
  1,&0\leq n_2\leq N-2,\\[1mm]
  \exp(2\pi n_1H_{\mathrm{link}}/N),&n_2=N-1,
 \end{cases}
 \label{cusp-link-flux-u2}\\
 U_3(n)&=U_4(n)=1.
 \label{cusp-link-flux-u34}
\end{align}
For
\begin{equation}
 U_{ij}(n)=U_i(n)U_j(n+\mathbf e_i)
            U_i(n+\mathbf e_j)^{-1}U_j(n)^{-1},
 \label{cusp-link-plaquette-convention}
\end{equation}
one has
\begin{equation}
 \begin{aligned}
 U_{12}(n)&=\exp(2\pi H_{\mathrm{link}}/N^2)
 &&\text{for every }n,\\
 U_{ij}(n)&=1
 &&\text{for }1\leq i<j\leq4,\quad(i,j)\neq(1,2).
 \end{aligned}
 \label{cusp-link-plaquette-values}
\end{equation}

Let \(\operatorname{Log}\) be the principal logarithm on the displayed
plaquette.  It is unambiguous because \(0<2\pi/N^2\leq\pi/2\).  Define
the logarithmic curvature two-form and its metric-weighted action by
\begin{align}
 \mathcal F_N^{\log}(n)
 &=N^2\operatorname{Log}(U_{12}(n))\,dy^1\wedge dy^2,
 \label{cusp-link-log-curvature}\\
 S^{\log}_{T,M_{\mathrm{cov}},N}(U)
 &=\frac{\mathcal D_{T,M_{\mathrm{cov}}}}{2g_{\mathrm{YM}}^2N^4}
   \sum_{n\in V_N}
   \lVert\mathcal F_N^{\log}(n)\rVert^2_{
       G_{T,M_{\mathrm{cov}}}^{-1}\otimes\mathfrak{su}(2)}.
 \label{cusp-link-log-action}
\end{align}
Then
\begin{equation}
 \mathcal F_N^{\log}(n)=2\pi H_{\mathrm{link}}\,dy^1\wedge dy^2,
 \qquad
 S^{\log}_{T,M_{\mathrm{cov}},N}(U)
 =\frac{2\pi^2M_{\mathrm{cov}}^2}{g_{\mathrm{YM}}^2D_T}.
 \label{cusp-link-log-action-value}
\end{equation}
In particular, the native-fibre value \(M_{\mathrm{cov}}=1\) is
\(2\pi^2/(g_{\mathrm{YM}}^2D_T)\), independently of the mesh.

There is no canonical scalar Wilson action for a general nonorthogonal
metric obtained merely by assigning independent character weights to all
coordinate plaquettes.  For this configuration, however, the exact
continuum-matched \((1,2)\)-sector Wilson functional is unambiguous if it is
defined explicitly by the following formula.  Its ordered plaquette
holonomy and real trace are the Wilson lattice-action construction
\cite{Wilson_1974}; the metric coefficient and normalization here are the
present definition, and their consequences are proved below rather than
imported from that source:
\begin{equation}
 S^{W,12}_{T,M_{\mathrm{cov}},N}(U)
 =\frac{\mathcal D_{T,M_{\mathrm{cov}}}\kappa^{12}_{T,M_{\mathrm{cov}}}}
        {2g_{\mathrm{YM}}^2}
   \sum_{n\in V_N}
   \left(2-\operatorname{Re}\operatorname{tr}U_{12}(n)\right),
 \label{cusp-link-wilson-sector-definition}
\end{equation}
where
\begin{equation}
 \kappa^{12}_{T,M_{\mathrm{cov}}}
 =(G_{T,M_{\mathrm{cov}}}^{-1})_{11}
  (G_{T,M_{\mathrm{cov}}}^{-1})_{22}
 -(G_{T,M_{\mathrm{cov}}}^{-1})_{12}^2
 =D_T^{-2}.
 \label{cusp-link-two-form-minor}
\end{equation}
Its exact value and continuum-mesh limit are
\begin{align}
 S^{W,12}_{T,M_{\mathrm{cov}},N}(U)
 &=\frac{M_{\mathrm{cov}}^2N^4}{2g_{\mathrm{YM}}^2D_T}
   \left(2-2\cos\frac{2\pi}{N^2}\right),
 \label{cusp-link-wilson-sector-value}\\
 &=\frac{M_{\mathrm{cov}}^2}{g_{\mathrm{YM}}^2D_T}
   \left(
    2\pi^2-\frac{2\pi^4}{3N^4}
    +\frac{4\pi^6}{45N^8}+O(N^{-12})
   \right),
 \label{cusp-link-wilson-sector-expansion}\\
 \lim_{N\to\infty}S^{W,12}_{T,M_{\mathrm{cov}},N}(U)
 &=\frac{2\pi^2M_{\mathrm{cov}}^2}{g_{\mathrm{YM}}^2D_T}.
 \label{cusp-link-wilson-sector-limit}
\end{align}
\end{theorem}

\begin{proof}
All links lie in the same diagonal circle subgroup, so they commute.  If
\(0\leq n_2\leq N-2\), the two \(U_2\)-terms in
\eqref{cusp-link-plaquette-convention} are trivial and
\[
 U_{12}(n)=
 \exp(-2\pi n_2H_{\mathrm{link}}/N^2)
 \exp(2\pi(n_2+1)H_{\mathrm{link}}/N^2)
 =\exp(2\pi H_{\mathrm{link}}/N^2).
\]
If \(n_2=N-1\), then \(U_1(n+\mathbf e_2)=1\).  Moreover,
\[
 U_2(n+\mathbf e_1)U_2(n)^{-1}=\exp(2\pi H_{\mathrm{link}}/N).
\]
For \(n_1=N-1\), the exponent computed with representatives differs from
\(2\pi H_{\mathrm{link}}/N\) by \(-2\pi H_{\mathrm{link}}\), whose exponential is one; hence the formula
also holds across that wrap boundary.  It follows that
\[
 U_{12}(n)=
 \exp\left(-\frac{2\pi(N-1)}{N^2}H_{\mathrm{link}}\right)
 \exp\left(\frac{2\pi}{N}H_{\mathrm{link}}\right)
 =\exp(2\pi H_{\mathrm{link}}/N^2).
\]
The functions \(U_1,U_2\) are independent of \(n_3,n_4\), and
\(U_3=U_4=1\), proving every remaining assertion in
\eqref{cusp-link-plaquette-values}.  Reversing the orientation gives
\(U_{21}(n)=U_{12}(n)^{-1}\); it is not an additional trivial plaquette.

The lower-right \(2\times2\) principal block of the cover Gram matrix is
\(M_{\mathrm{cov}}^2I_2\).  Jacobi's
complementary-minor identity therefore gives
\[
 \kappa^{12}_{T,M_{\mathrm{cov}}}
 =\frac{\det(M_{\mathrm{cov}}^2I_2)}
        {\det G_{T,M_{\mathrm{cov}}}}
 =\frac{M_{\mathrm{cov}}^4}{M_{\mathrm{cov}}^4D_T^2}
 =D_T^{-2},
\]
which proves \eqref{cusp-link-two-form-minor}.  By
\eqref{cusp-link-flux-generator} and the principal-log choice,
\(N^2\operatorname{Log}U_{12}=2\pi H_{\mathrm{link}}\).  Hence every one of the \(N^4\)
summands in \eqref{cusp-link-log-action} equals
\(4\pi^2D_T^{-2}\).  Substitution of
\(\mathcal D_{T,M_{\mathrm{cov}}}=M_{\mathrm{cov}}^2D_T\) proves
\eqref{cusp-link-log-action-value}.

Finally, \(H_{\mathrm{link}}\) has eigenvalues \(\pm i\), so
\[
 \operatorname{Re}\operatorname{tr}\exp(2\pi H_{\mathrm{link}}/N^2)
 =2\cos(2\pi/N^2).
\]
There are \(N^4\) equal plaquettes.  Substitution proves
\eqref{cusp-link-wilson-sector-value}.  Applying
\[
 1-\cos x=\frac{x^2}{2}-\frac{x^4}{24}
           +\frac{x^6}{720}+O(x^8)
\]
with \(x=2\pi/N^2\) proves
\eqref{cusp-link-wilson-sector-expansion} and
\eqref{cusp-link-wilson-sector-limit}.
\end{proof}

\subsection{Finite Gaussian quantization with independent cover and mesh}

Let
\begin{equation}
 \vartheta^j_{T,M_{\mathrm{cov}}}
 =P_{T,M_{\mathrm{cov}}}^{-\mathsf T}\mathbf e_j
 \in(\mathbb R^4)^*,
 \qquad 1\leq j\leq4.
 \label{cusp-link-physical-coframe}
\end{equation}
For \(0\leq p\leq4\), define the real finite cochain spaces
\begin{equation}
 C^p_{N_{\mathrm{mesh}}}(T,M_{\mathrm{cov}})
 =\operatorname{Map}\left(
 V_{N_{\mathrm{mesh}}},
 \bigwedge^p(\mathbb R^4)^*\otimes\mathfrak{su}(2)
 \right).
 \label{cusp-link-cochain-spaces}
\end{equation}
The exterior powers carry the metric induced by the physical coframe
\eqref{cusp-link-physical-coframe}; explicitly, for increasing multi-indices
\(I=(i_1,\ldots,i_p)\) and \(J=(j_1,\ldots,j_p)\),
\begin{equation}
 \left\langle
  \vartheta^{i_1}\wedge\cdots\wedge\vartheta^{i_p},
  \vartheta^{j_1}\wedge\cdots\wedge\vartheta^{j_p}
 \right\rangle
 =
 \det\left((G_{T,M_{\mathrm{cov}}}^{-1})_{i_a j_b}\right)_{a,b=1}^p.
 \label{cusp-link-exterior-metric}
\end{equation}
Together with \eqref{cusp-link-su2-inner-product}, this defines the
pointwise tensor-product metric.  The cochain inner product is
\begin{equation}
 \langle a,b\rangle_p
 =\frac{\mathcal D_{T,M_{\mathrm{cov}}}}{N_{\mathrm{mesh}}^4}
   \sum_{n\in V_{N_{\mathrm{mesh}}}}
   \langle a(n),b(n)\rangle.
 \label{cusp-link-cochain-inner-product}
\end{equation}
Let \(\mathsf T_j a(n)=a(n+\mathbf e_j)\), set
\begin{equation}
 D^+_{N_{\mathrm{mesh}},j}
 =N_{\mathrm{mesh}}(\mathsf T_j-1),
 \qquad
 D^-_{N_{\mathrm{mesh}},j}a(n)
 =N_{\mathrm{mesh}}(a(n)-a(n-\mathbf e_j)),
 \label{cusp-link-differences}
\end{equation}
and define
\begin{equation}
 d_{N_{\mathrm{mesh}}}a
 =\sum_{j=1}^4
  \varepsilon(\vartheta^j_{T,M_{\mathrm{cov}}})
  D^+_{N_{\mathrm{mesh}},j}a.
 \label{cusp-link-finite-exterior-derivative}
\end{equation}
The commuting differences and anticommuting exterior multiplications give
\(d_{N_{\mathrm{mesh}}}^2=0\).  Its adjoint is
\begin{equation}
 d_{N_{\mathrm{mesh}}}^*
 =-\sum_{j=1}^4
  \iota_{(\vartheta^j_{T,M_{\mathrm{cov}}})^\sharp}
  D^-_{N_{\mathrm{mesh}},j}.
 \label{cusp-link-finite-adjoint}
\end{equation}
Define the real reduced Coulomb space
\begin{align}
 \mathscr C_{T,M_{\mathrm{cov}},N_{\mathrm{mesh}}}
 &=
 \ker(d_{N_{\mathrm{mesh}}}^*:C^1\to C^0)
 \cap
 \left(\ker\Delta^{(1)}_{N_{\mathrm{mesh}}}\right)^\perp,
 \label{cusp-link-reduced-coulomb-space}\\
 \Delta^{(1)}_{N_{\mathrm{mesh}}}
 &=d_{N_{\mathrm{mesh}}}d_{N_{\mathrm{mesh}}}^*
 +d_{N_{\mathrm{mesh}}}^*d_{N_{\mathrm{mesh}}}.
 \label{cusp-link-finite-hodge-laplacian}
\end{align}
On this finite-dimensional space put
\begin{equation}
 \mathcal L_{T,M_{\mathrm{cov}},N_{\mathrm{mesh}}}
 =g_{\mathrm{YM}}^{-2}d_{N_{\mathrm{mesh}}}^*
                         d_{N_{\mathrm{mesh}}}.
 \label{cusp-link-finite-hessian}
\end{equation}

\begin{proposition}[Exact finite Hessian spectrum]
\label{cusp-link-finite-spectrum}
For
\begin{equation}
 \mathfrak d_{N_{\mathrm{mesh}}}(k)_j
 =N_{\mathrm{mesh}}
  \left(e^{2\pi i k_j/N_{\mathrm{mesh}}}-1\right),
 \qquad
 k\in(\mathbb Z/N_{\mathrm{mesh}}\mathbb Z)^4,
 \label{cusp-link-fourier-symbol}
\end{equation}
the positive spectrum of \(\mathcal L\) is obtained from
\begin{equation}
 \lambda_{T,M_{\mathrm{cov}},N_{\mathrm{mesh}}}(k)
 =\frac1{g_{\mathrm{YM}}^2}
  \left\lVert
   P_T^{-\mathsf T}S_{M_{\mathrm{cov}}}^{-1}
   \mathfrak d_{N_{\mathrm{mesh}}}(k)
  \right\rVert^2,
 \qquad k\neq0.
 \label{cusp-link-finite-mode-eigenvalue}
\end{equation}
Each complex Fourier fibre has transverse multiplicity
\(3\dim\mathfrak{su}(2)=9\).  In particular,
\begin{equation}
 \lambda_{\min}(\mathcal L)
 =\min_{k\neq0}\lambda_{T,M_{\mathrm{cov}},N_{\mathrm{mesh}}}(k)>0.
 \label{cusp-link-finite-minimum}
\end{equation}
For \(\mathbf e_1=(1,0,0,0)^{\mathsf T}\), the exact mode value is
\begin{equation}
 \lambda_{T,M_{\mathrm{cov}},N_{\mathrm{mesh}}}(\mathbf e_1)
 =\frac{4N_{\mathrm{mesh}}^2
              \sin^2(\pi/N_{\mathrm{mesh}})}{g_{\mathrm{YM}}^2}
  \frac{L_T^2+m_T^2}{D_T^2};
 \label{cusp-link-cover-soft-mode}
\end{equation}
it is independent of \(M_{\mathrm{cov}}\).
\end{proposition}

\begin{proof}
Translation is unitary for \eqref{cusp-link-cochain-inner-product}, so
\((D_j^+)^*=-D_j^-\).  Exterior multiplication by \(\vartheta^j\) has
adjoint contraction by \((\vartheta^j)^\sharp\), proving
\eqref{cusp-link-finite-adjoint}.  On the Fourier character
\(\chi_k(n)=e^{2\pi i k\cdot n/N_{\mathrm{mesh}}}\), the derivative symbol
is exterior multiplication by
\[
 \zeta(k)=P_{T,M_{\mathrm{cov}}}^{-\mathsf T}
           \mathfrak d_{N_{\mathrm{mesh}}}(k)
 =P_T^{-\mathsf T}S_{M_{\mathrm{cov}}}^{-1}
           \mathfrak d_{N_{\mathrm{mesh}}}(k).
\]
The Hermitian exterior-algebra identity
\[
 \iota_{\overline\zeta^\sharp}\varepsilon(\zeta)
 +\varepsilon(\zeta)\iota_{\overline\zeta^\sharp}
 =\lVert\zeta\rVert^2\operatorname{id}
\]
shows that on the Coulomb polarization
\(\ker\iota_{\overline\zeta^\sharp}\subset(\mathbb C^4)^*\), the operator
\(d^*d\) is multiplication by \(\lVert\zeta\rVert^2\).  For \(k\neq0\),
\(\mathfrak d_{N_{\mathrm{mesh}}}(k)\neq0\), so invertibility of
\(P_{T,M_{\mathrm{cov}}}\) gives
\(\zeta(k)\neq0\).  The contraction has rank one, and its kernel has
dimension three.  The zero Fourier fibre consists exactly of constant
harmonic cochains and was removed in
\eqref{cusp-link-reduced-coulomb-space}.  Fourier completeness now proves
\eqref{cusp-link-finite-mode-eigenvalue}--\eqref{cusp-link-finite-minimum}.

For \(k=\mathbf e_1\), the only nonzero component of
\(\mathfrak d_{N_{\mathrm{mesh}}}(k)\) has modulus
\[
 2N_{\mathrm{mesh}}\sin(\pi/N_{\mathrm{mesh}}),
\]
and
\(S_{M_{\mathrm{cov}}}^{-1}\mathbf e_1=\mathbf e_1\).  Solving
\(P_T^{\mathsf T}\xi=\mathbf e_1\), or directly inverting the lower
\(2\times2\) block in \eqref{cusp-link-period-matrix}, gives
\[
 \lVert P_T^{-\mathsf T}\mathbf e_1\rVert^2
 =\frac{L_T^2+m_T^2}{D_T^2}.
\]
Substitution proves \eqref{cusp-link-cover-soft-mode}.
\end{proof}

Let \(da\) denote Lebesgue measure on
\(\mathscr C_{T,M_{\mathrm{cov}},N_{\mathrm{mesh}}}\) induced by
\eqref{cusp-link-cochain-inner-product} and define
\begin{equation}
 \mathcal H^{\mathrm{Gauss}}_{T,M_{\mathrm{cov}},N_{\mathrm{mesh}}}
 =L^2(\mathscr C_{T,M_{\mathrm{cov}},N_{\mathrm{mesh}}},da).
 \label{cusp-link-gaussian-hilbert-space}
\end{equation}

\begin{theorem}[Finite Gaussian oscillator and a resolved cusp sequence]
\label{cusp-link-gaussian-theorem}
The Friedrichs realization on \eqref{cusp-link-gaussian-hilbert-space} of
\begin{equation}
 \widehat H^{\mathrm{Gauss}}_{T,M_{\mathrm{cov}},N_{\mathrm{mesh}}}
 =-\frac12\Delta_{\mathscr C}
  +\frac12\langle a,
     \mathcal L_{T,M_{\mathrm{cov}},N_{\mathrm{mesh}}}a\rangle
 \label{cusp-link-gaussian-hamiltonian}
\end{equation}
has a unique Gaussian ground state and spectral gap
\begin{equation}
 \operatorname{gap}(\widehat H^{\mathrm{Gauss}})
 =\sqrt{\lambda_{\min}(
   \mathcal L_{T,M_{\mathrm{cov}},N_{\mathrm{mesh}}})}.
 \label{cusp-link-gaussian-gap}
\end{equation}
Consequently,
\begin{equation}
 \operatorname{gap}(\widehat H^{\mathrm{Gauss}})
 \leq
 \frac{2N_{\mathrm{mesh}}\sin(\pi/N_{\mathrm{mesh}})}{g_{\mathrm{YM}}}
 \frac{\sqrt{L_T^2+m_T^2}}{D_T}.
 \label{cusp-link-gaussian-gap-bound}
\end{equation}

There is an explicit diagonal sequence for which the cover geometry tends locally to
\(\mathbb R^4\), the cubical mesh resolves that geometry, and the Gaussian
gap tends to zero.  Explicitly, for integers \(j\to\infty\), take
\begin{equation}
 T_j=j^2,\qquad M_{\mathrm{cov},j}=j,\qquad
 N_{\mathrm{mesh},j}=j^3.
 \label{cusp-link-diagonal-sequence}
\end{equation}
Then
\begin{equation}
 \begin{aligned}
 \operatorname{sys}(\widetilde F_{T_j,M_{\mathrm{cov},j}})
 &\longrightarrow\infty,\\
 \max_r
 \frac{\lVert P_{T_j,M_{\mathrm{cov},j}}\mathbf e_r\rVert}
      {N_{\mathrm{mesh},j}}
 &\longrightarrow0,\\
 \operatorname{gap}(\widehat H^{\mathrm{Gauss}})
 &\longrightarrow0,\\
 S^{\log}_{T_j,M_{\mathrm{cov},j},N_{\mathrm{mesh},j}}(U)
 &\longrightarrow0.
 \end{aligned}
 \label{cusp-link-diagonal-conclusions}
\end{equation}
Thus the pointed cover tori converge locally isometrically to \(\mathbb R^4\)
while their finite linearized Gaussian systems are simultaneously refined.
\end{theorem}

\begin{proof}
The real finite-dimensional operator \(\mathcal L\) is positive definite on
the reduced Coulomb space.  Choose an orthonormal eigenbasis with eigenvalues
\(0<\lambda_1\leq\cdots\leq\lambda_r\).  In the resulting coordinates,
\eqref{cusp-link-gaussian-hamiltonian} is the tensor sum
\[
 \sum_{s=1}^r
 \left(-\frac12\frac{\partial^2}{\partial x_s^2}
       +\frac12\lambda_sx_s^2\right).
\]
The normalized Hermite basis has eigenvalues
\(\sum_s(n_s+\tfrac12)\sqrt{\lambda_s}\), with
\(n_s\in\mathbb Z_{\geq0}\).  The all-zero multi-index is the unique ground
state, and the smallest possible positive excitation is
\(\sqrt{\lambda_1}\).  This proves \eqref{cusp-link-gaussian-gap}; applying
the minimum to the explicit \(\mathbf e_1\)-mode in
\eqref{cusp-link-cover-soft-mode} proves
\eqref{cusp-link-gaussian-gap-bound}.

For completeness, reorder physical coordinates as in
\eqref{cusp-link-block-form}.  Every vector of the cover lattice has the
form
\[
 \binom{A_Tn+M_{\mathrm{cov}}m}{B_Tn},
 \qquad n,m\in\mathbb Z^2.
\]
Let \(C=\sup_T\lVert C_T\rVert_{\mathrm{op}}\).  If \(n\neq0\), then
\(\lVert B_Tn\rVert\geq(T-C)\lVert n\rVert\); if \(n=0\) and the lattice
vector is nonzero, its norm is at least \(M_{\mathrm{cov}}\).  Hence
\begin{equation}
 \operatorname{sys}(\widetilde F_{T,M_{\mathrm{cov}}})
 \geq\min\{T-C,M_{\mathrm{cov}}\}.
 \label{cusp-link-cover-systole-bound}
\end{equation}
This tends to infinity for the first two choices in
\eqref{cusp-link-diagonal-sequence}.  For any fixed \(R>0\), once the
systole exceeds \(4R\), the quotient map from \(\mathbb R^4\) is injective
and distance-preserving on the radius-\(R\) ball: for a nonzero period
\(\lambda\) and \(x,y\) in that ball,
\[
 \lVert x-y+\lambda\rVert
 \geq\lVert\lambda\rVert-\lVert x-y\rVert
 >4R-2R\geq\lVert x-y\rVert.
\]
This proves the pointed local-isometric \(\mathbb R^4\) limit.

The first two columns of \(P_T\) have norm \(O(T)\), while the last two
columns of \(P_{T,M_{\mathrm{cov}}}\) have norm exactly
\(M_{\mathrm{cov}}\).  Therefore the largest physical edge in the
\(N_{\mathrm{mesh}}\)-subdivision has length
\(O(\max\{T,M_{\mathrm{cov}}\}/N_{\mathrm{mesh}})\), which is \(O(j^{-1})\)
for \eqref{cusp-link-diagonal-sequence}.  Finally,
\eqref{cusp-link-cusp-asymptotics} gives
\(\sqrt{L_T^2+m_T^2}/D_T=T^{-1}(1+O(T^{-1}))\), while
\(2N\sin(\pi/N)\leq2\pi\).  The right-hand side of
\eqref{cusp-link-gaussian-gap-bound} is therefore \(O(j^{-2})\).
Moreover, \eqref{cusp-link-log-action-value} is \(O(j^{-2})\), since
\(M_{\mathrm{cov},j}^2/D_{T_j}=j^2/(j^4+O(j^2))\).
This proves every assertion in \eqref{cusp-link-diagonal-conclusions}.
\end{proof}

\subsection{The exponential map and the exact physical boundary}

Let \(C^1=C^1_{N_{\mathrm{mesh}}}(T,M_{\mathrm{cov}})\), regarded only as a
finite-dimensional Lie-algebra cochain space, and define componentwise
\begin{equation}
 \operatorname{Exp}_{N_{\mathrm{mesh}}}:C^1
 \longrightarrow\mathrm{SU}(2)^{E_{N_{\mathrm{mesh}}}},
 \qquad
 a\longmapsto\left(\exp(a_e/N_{\mathrm{mesh}})\right)_e.
 \label{cusp-link-exponential-map}
\end{equation}
Here, if
\(a(n)=\sum_{j=1}^4\vartheta^j_{T,M_{\mathrm{cov}}}\otimes a_j(n)\),
then \(a_e=a_j(n)\) for the oriented edge \(e=(n,j)\).

\begin{proposition}[Local chart, not a global embedding]
\label{cusp-link-exponential-local-only}
The map \eqref{cusp-link-exponential-map} is smooth and a local
diffeomorphism at zero, but it is not globally injective.  Under left
trivialization of every tangent space at the identity,
\begin{equation}
 D\operatorname{Exp}_{N_{\mathrm{mesh}}}(0)\,a
 =N_{\mathrm{mesh}}^{-1}a.
 \label{cusp-link-exponential-derivative}
\end{equation}
If \(S_W\) is any smooth compact-link action for which the identity
configuration is critical, then
\begin{equation}
 D^2(S_W\circ\operatorname{Exp}_{N_{\mathrm{mesh}}})_0[a,b]
 =N_{\mathrm{mesh}}^{-2}
  D^2(S_W)_{\mathbf1}[a,b]
 \label{cusp-link-exponential-hessian-scaling}
\end{equation}
under the same tangent identification.  Thus, for a prescribed cochain
quadratic form \(Q\),
\[
 D^2(S_W\circ\operatorname{Exp}_{N_{\mathrm{mesh}}})_0=Q
\]
holds exactly when
\[
 D^2(S_W)_{\mathbf1}=N_{\mathrm{mesh}}^2Q
\]
under the coefficient identification.  Both this scaling and the tensor
structure must be verified for the chosen compact-link action; neither
follows from the exponential map alone.
\end{proposition}

\begin{proof}
For the Lie-group exponential, \(D\exp_0\) is the identity.  Applying the
chain rule to \(a_e\mapsto a_e/N_{\mathrm{mesh}}\) gives
\eqref{cusp-link-exponential-derivative}, which is invertible, so the inverse
function theorem gives a local diffeomorphism.  It is not globally
injective: a cochain which is zero on every edge except one, where it equals
\(2\pi N_{\mathrm{mesh}}H_{\mathrm{link}}\), has the same image as the zero cochain by
\eqref{cusp-link-flux-generator}.  Finally, the second-derivative chain rule
contains a term involving \(DS_W(\mathbf1)\), which vanishes by criticality;
the two remaining first derivatives of \(\operatorname{Exp}\) each supply
the factor \(N_{\mathrm{mesh}}^{-1}\).  This proves
\eqref{cusp-link-exponential-hessian-scaling}.
\end{proof}

The Gaussian Hilbert space in \eqref{cusp-link-gaussian-hilbert-space} is
not the compact-link Yang--Mills Hilbert space.  At this finite lattice, the
latter kinematic and gauge-invariant spaces would be
\begin{equation}
 \mathcal H^{\mathrm{link}}_{N_{\mathrm{mesh}}}
 =L^2(\mathrm{SU}(2)^{E_{N_{\mathrm{mesh}}}},dU),
 \qquad
 \mathcal H^{\mathrm{inv}}_{N_{\mathrm{mesh}}}
 =\left(\mathcal H^{\mathrm{link}}_{N_{\mathrm{mesh}}}\right)^{
    \mathrm{SU}(2)^{V_{N_{\mathrm{mesh}}}}},
 \label{cusp-link-physical-hilbert-space}
\end{equation}
where \(dU\) is product normalized Haar measure and the superscript denotes
the nonlinear vertex-gauge action.  Here ``gauge-invariant'' refers only to
the link-configuration space of this four-dimensional cell complex.  Without
a chosen time direction or transfer matrix, it is not asserted to be a
canonical Hamiltonian physical Hilbert space.  L\"uscher's construction of
a selfadjoint, strictly positive transfer matrix
\cite{Luscher_1977} and Menotti--Pelissetto's reflection-positivity theorem
for the Wilson action \cite{MenottiPelissetto_1987} concern the standard
Wilson lattice setting after the requisite reflection/time structure has
been fixed; they do not supply those omitted choices for the nonorthogonal
sector functional \eqref{cusp-link-wilson-sector-definition}.  The
Osterwalder--Schrader reconstruction boundary is likewise the one stated in
the primary Euclidean axioms papers
\cite{OsterwalderSchrader1973,OsterwalderSchrader1975}, not a consequence of
a finite Gaussian oscillator calculation.  The differences are
structural:
\begin{enumerate}
\item \(\mathcal H^{\mathrm{Gauss}}\) is an \(L^2\)-space on a noncompact
linear Coulomb slice with Lebesgue measure; \(\mathcal H^{\mathrm{inv}}\)
is a gauge-invariant subspace of an \(L^2\)-space on compact link variables
with Haar measure.
\item The local exponential map is neither a global coordinate chart nor a
measure-preserving bijection.  No unitary map between the two displayed
Hilbert spaces has been constructed.
\item A link configuration \(U\) determines the Dirac distribution
\(\delta_U\), but
\begin{equation}
 \delta_U\in\mathcal D'(\mathrm{SU}(2)^{E_{N_{\mathrm{mesh}}}}),
 \qquad
 \delta_U\notin L^2(\mathrm{SU}(2)^{E_{N_{\mathrm{mesh}}}},dU).
 \label{cusp-link-dirac-not-l2}
\end{equation}
Indeed, product Haar measure is nonatomic, so a point mass is singular with
respect to it and has no square-integrable density.
\item No nonlinear compact-link Hamiltonian, Gauss-law spectral problem,
time direction, reflection-positive transfer matrix, renormalized-scale
limit, or Osterwalder--Schrader reconstruction is supplied here.
\end{enumerate}
Accordingly, \eqref{cusp-link-gaussian-gap} is an exact theorem about the
finite quadratic oscillator defined by the Coulomb-restricted operator
\(g_{\mathrm{YM}}^{-2}d^*d\).  It is not an eigenvalue statement about a
nonperturbative physical Yang--Mills
Hamiltonian, and the small classical action in
\eqref{cusp-link-log-action-value} is not by itself a physical-state Rayleigh
quotient.

\section{Integral magnetic sectors and a reducible
\texorpdfstring{\(SU(2)\)}{SU(2)} critical point on the radial cusp}
\label{cusp-magnetic-section}

Retain the positive radial flat tori
\(F_T=\mathbb R^4/P_T\mathbb Z^4\), their Euclidean quotient metrics, and
the real quantities \(L_T,q_T,m_T,D_T,H,B,m\) from
\cref{cusp-spec-radial-notation,cusp-spec-exact-dual-lattice,cusp-spec-D-expansion}.
In particular,
\[
 L_T=\operatorname{Im}\tau_T,\qquad
 q_T=-\operatorname{Im}\beta_T,\qquad
 D_T=L_Tq_T+6m_T^2=\operatorname{Vol}(F_T)>0.
\]
The real coordinate order is the already fixed order
\[
 x=(x_1,y_1,x_2,y_2)^{\mathsf t}.
\]

\subsection{The integral coframe}

Let \(e_1,\ldots,e_4\) be the standard basis of \(\mathbb R^4\), put
\(u=P_T^{-1}x\), and define
\begin{equation}
 \theta_T^j:=\mathrm du_j
 =(P_T^{-\mathsf T}e_j)^{\mathsf t}
   (\mathrm dx_1,\mathrm dy_1,\mathrm dx_2,\mathrm dy_2)^{\mathsf t}.
 \label{cusp-magnetic-integral-coframe}
\end{equation}

\begin{proposition}[Integral coframe and the distinguished unit-flux plane]
\label{cusp-magnetic-coframe-proposition}
The forms \(\theta_T^j\) are the integral coframe dual to the ordered
marked lattice basis:
\[
 \int_{[0,1]\ni s\mapsto[x+sP_Te_k]}\theta_T^j=\delta_{jk}.
\]
The first two members are exactly
\begin{equation}
 \theta_T^1=\frac{m_T}{D_T}\,\mathrm dy_1
             -\frac{L_T}{D_T}\,\mathrm dy_2,
 \qquad
 \theta_T^2=\frac{q_T}{D_T}\,\mathrm dy_1
             +\frac{6m_T}{D_T}\,\mathrm dy_2,
 \label{cusp-magnetic-first-two-covectors}
\end{equation}
and hence
\begin{equation}
 \theta_T^1\wedge\theta_T^2
 =\frac1{D_T}\,\mathrm dy_1\wedge\mathrm dy_2,
 \qquad
 \lvert\theta_T^1\wedge\theta_T^2\rvert^2
 =\frac1{D_T^2}.
 \label{cusp-magnetic-wedge-norm}
\end{equation}
\end{proposition}

\begin{proof}
For \(n\in\mathbb Z^4\),
\[
 u(x+P_Tn)=u(x)+n,
\]
which proves the integral-period assertion.  The first two columns of the
already-proved matrix \(P_T^{-\mathsf T}\) in
\cref{cusp-spec-full-inverse-transpose} give
\cref{cusp-magnetic-first-two-covectors}.  Their wedge coefficient is
\[
 \frac{6m_T^2+L_Tq_T}{D_T^2}=\frac1{D_T}.
\]
The coordinate forms \(\mathrm dy_1,\mathrm dy_2\) are orthonormal for the fixed quotient
metric, proving \cref{cusp-magnetic-wedge-norm}.
\end{proof}

\subsection{The unit-flux line and the topologically trivial rank-two sum}

The following construction is made in the marked coordinates \(u\) of
\eqref{cusp-magnetic-integral-coframe}; it therefore applies to the actual
skew lattice \(P_T\mathbb Z^4\), not to a rectangular replacement.

\begin{proposition}[Explicit factor of automorphy and connection]
\label{cusp-magnetic-line-bundle-proposition}
For \(n\in\mathbb Z^4\), set
\begin{equation}
 j_n(u)=\exp(-2\pi i n_1u_2).
 \label{cusp-magnetic-unit-factor}
\end{equation}
The quotient
\[
 \mathscr L_T=(\mathbb R^4\times\mathbb C)/\mathbb Z^4,
 \qquad
 n\cdot(u,z)=(u+n,j_n(u)z),
\]
is a Hermitian line bundle over \(F_T\).  The connection on its pullback
to \(\mathbb R^4\),
\begin{equation}
 \widetilde\nabla^{\mathscr L_T}
 =\mathrm d+2\pi i u_1\,\mathrm du_2,
 \label{cusp-magnetic-line-connection}
\end{equation}
descends, and its curvature is
\begin{equation}
 F_{\mathscr L_T}=2\pi i\,\mathrm du_1\wedge\mathrm du_2
        =2\pi i\,\theta_T^1\wedge\theta_T^2.
 \label{cusp-magnetic-line-curvature}
\end{equation}
With the unitary Chern--Weil convention
\(c_1=[\,iF/(2\pi)\,]\), one has
\(c_1(\mathscr L_T)=-[\theta_T^1\wedge\theta_T^2]\).
\end{proposition}

\begin{proof}
For \(m,n\in\mathbb Z^4\),
\[
 j_m(u+n)j_n(u)
 =j_{m+n}(u)\exp(-2\pi i m_1n_2)
 =j_{m+n}(u).
\]
Thus \eqref{cusp-magnetic-unit-factor} is a factor of automorphy.  Its
values are unitary, so the standard fibre metric descends.  If
\(\mathcal A=2\pi i u_1\,\mathrm du_2\), then
\[
 \mathcal A(u+n)=\mathcal A(u)+2\pi i n_1\,\mathrm du_2,
 \qquad
 j_n^{-1}\mathrm dj_n=-2\pi i n_1\,\mathrm du_2.
\]
Hence \(\mathcal A(u+n)+j_n^{-1}\mathrm dj_n=\mathcal A(u)\), which is precisely
the descent identity for \(\mathrm d+\mathcal A\) under the stated quotient
convention.  Taking \(\mathrm d\mathcal A\) gives
\eqref{cusp-magnetic-line-curvature}.  The Chern--Weil statement follows
without a sign suppression:
\[
 \frac{i}{2\pi}F_{\mathscr L_T}=-\theta_T^1\wedge\theta_T^2.
\]
\end{proof}

Let \(\mathscr L_T^{-1}=\mathscr L_T^*\) carry the dual metric and dual
connection, and set
\begin{equation}
 \mathscr E_T=\mathscr L_T\oplus\mathscr L_T^{-1}.
 \label{cusp-magnetic-rank-two-bundle}
\end{equation}

\begin{proposition}[Determinant, \(c_2\), and explicit
\(SU(2)\)-triviality]
\label{cusp-magnetic-su2-triviality}
The factors of automorphy of \(\mathscr E_T\) are
\[
 G_n(u)=\begin{pmatrix}j_n(u)&0\\0&j_n(u)^{-1}\end{pmatrix}\in SU(2).
\]
Evaluation gives a canonical parallel determinant trivialization
\[
 \det\mathscr E_T=\mathscr L_T\otimes\mathscr L_T^{-1}
 \xrightarrow{\ \mathrm{ev}\ }F_T\times\mathbb C.
\]
Moreover
\[
 c_1(\mathscr E_T)=0,
 \qquad
 c_2(\mathscr E_T)=0\quad\text{in }H^4(F_T;\mathbb Z),
\]
and the underlying principal \(SU(2)\)-bundle is smoothly trivial.  The
last assertion follows from the explicit nullhomotopy below and does not
rely only on the vanishing of a de Rham Chern form.
\end{proposition}

\begin{proof}
The determinant assertion follows pointwise from
\(j_nj_n^{-1}=1\).  The induced determinant connection has connection
form
\[
 \operatorname{tr}
 \begin{pmatrix}
  2\pi i u_1\,\mathrm du_2&0\\0&-2\pi i u_1\,\mathrm du_2
 \end{pmatrix}=0,
\]
so the evaluation trivialization is parallel.

Let \(c=c_1(\mathscr L_T)\).  The integral Whitney identity gives
\[
 c(\mathscr E_T)=(1+c)(1-c)=1-c^2.
\]
The factor \(j_n\) depends only on \((u_1,u_2)\), so \(\mathscr L_T\) is the
pullback of a line bundle \(L_0\) on
\(\mathbb R^2/\mathbb Z^2\) under the marked projection
\([u_1,u_2,u_3,u_4]\mapsto[u_1,u_2]\).  Hence
\[
 c^2=\pi_{12}^*\bigl(c_1(L_0)^2\bigr)=0
\]
integrally because \(H^4(\mathbb T^2;\mathbb Z)=0\).  This proves the two
Chern-class identities on the four-torus.

It remains to trivialize the \(SU(2)\)-bundle itself.  On the fundamental
square \(0\le u_1,u_2\le1\) for the first two marked coordinates, the only
nontrivial seam is
\[
 (0,t,v)\sim(1,t,\gamma(t)v),
 \qquad
 \gamma(t)=
 \begin{pmatrix}e^{-2\pi it}&0\\0&e^{2\pi it}\end{pmatrix}.
\]
For \(0\le r\le1\), define
\[
 a(r,t)=(1-r)e^{-2\pi it}+r,
 \qquad
 s(r)=\sin(\pi r),
 \qquad
 N(r,t)=\bigl(|a(r,t)|^2+s(r)^2\bigr)^{1/2},
\]
and
\begin{equation}
 U(r,t)=\frac1{N(r,t)}
 \begin{pmatrix}
  a(r,t)&-s(r)\\
  s(r)&\overline{a(r,t)}
 \end{pmatrix}.
 \label{cusp-magnetic-explicit-nullhomotopy}
\end{equation}
If \(0<r<1\), then \(s(r)>0\); at \(r=0,1\), one has
\(|a(r,t)|=1\).  Thus \(N\) never vanishes.  The two columns in
\eqref{cusp-magnetic-explicit-nullhomotopy} are orthonormal and its
determinant is one, so \(U(r,t)\in SU(2)\).  It is smooth and periodic in
\(t\), and
\[
 U(0,t)=\gamma(t),
 \qquad
 U(1,t)=I_2.
\]
This is an explicit nullhomotopy of the clutching loop.  To see directly
that it produces a smooth global frame, choose a smooth
\(\chi:[0,1]\to[0,1]\) which is zero near \(0\) and one near \(1\), and put
\[
 K(\sigma,t)=U(1-\chi(\sigma),t).
\]
Then \(K=I_2\) near \(\sigma=0\), \(K=\gamma\) near \(\sigma=1\), and
\(K(\sigma,0)=K(\sigma,1)\).  The required trivialization, written in the
direction from the clutched bundle to the product bundle, is
\begin{equation}
 [\sigma,t,v]\longmapsto
 \bigl([\sigma,t],K(\sigma,t)^{-1}v\bigr).
 \label{cusp-magnetic-explicit-global-frame}
\end{equation}
Indeed, the quotient convention gives
\([0,t,v]=[1,t,\gamma(t)v]\), and the two representatives have the same
image because
\[
 K(1,t)^{-1}\gamma(t)v
 =\gamma(t)^{-1}\gamma(t)v
 =v=K(0,t)^{-1}v.
\]
At the other seam,
\(K(\sigma,0)^{-1}=K(\sigma,1)^{-1}\).  The constant collars used in the
definition of \(\chi\) make the local matrix in
\cref{cusp-magnetic-explicit-global-frame} constant after applying the
seam transition, so the map and its inverse are smooth across the seam.
It is fibrewise unitary and determinant one.  Extending it independently
of \(u_3,u_4\) and
pulling it through the marking gives the required global
\(SU(2)\)-frame of \(\mathscr E_T\).
\end{proof}

\subsection{A non-flat reducible Yang--Mills critical connection}

The symbol \(H=\operatorname{Im}h(0)\) is already fixed by
\cref{cusp-spec-limiting-imaginary-parts}.  For the requested Lie-algebra
element, fix instead the unambiguous notation
\begin{equation}
 \langle X,Y\rangle_{\mathfrak{su}(2)}
   =-\frac12\operatorname{Tr}(XY),
 \qquad
 H_{\mathrm{mag}}=\begin{pmatrix}i&0\\0&-i\end{pmatrix}.
 \label{cusp-magnetic-su2-normalization}
\end{equation}
\begin{theorem}[Reducible critical point and exact action]
\label{cusp-magnetic-ym-critical-point}
The direct-sum connection on \(\mathscr E_T\) has, in the lifted splitting,
\begin{equation}
 \nabla^T=\mathrm d+2\pi u_1\,\mathrm du_2\,H_{\mathrm{mag}},
 \qquad
 F_{\nabla^T}
 =2\pi H_{\mathrm{mag}}\,\theta_T^1\wedge\theta_T^2
 =\frac{2\pi}{D_T}H_{\mathrm{mag}}\,\mathrm dy_1\wedge\mathrm dy_2.
 \label{cusp-magnetic-su2-curvature}
\end{equation}
After any of the global \(SU(2)\)-trivializations from
 \cref{cusp-magnetic-su2-triviality}, this is a connection on the
trivial principal \(SU(2)\)-bundle.  It is non-flat and satisfies
\[
 \mathrm d_{\nabla^T}^*F_{\nabla^T}=0.
\]
For the normalization
\begin{equation}
 \operatorname{YM}_T(\nabla)
 =\frac1{2g_{\mathrm{YM}}^2}
  \int_{F_T}\lvert F_\nabla\rvert^2\,
  \mathrm d\operatorname{vol}_{F_T},
 \label{cusp-magnetic-action-normalization}
\end{equation}
its action is exactly
\begin{equation}
 \operatorname{YM}_T(\nabla^T)
 =\frac{2\pi^2}{g_{\mathrm{YM}}^2D_T}.
 \label{cusp-magnetic-exact-action}
\end{equation}
Its first inverse-power correction, in the canonical scalar notation
\(H=\operatorname{Im}h(0)\), \(B=\operatorname{Im}b(0)\), is
\begin{equation}
 \operatorname{YM}_T(\nabla^T)
 =\frac{2\pi^2}{g_{\mathrm{YM}}^2T^2}
  \left(1+\frac{B-2H}{T}
  +O(T^{-2})\right).
 \label{cusp-magnetic-action-first-correction}
\end{equation}
\end{theorem}

\begin{proof}
The connection matrix is
\[
 \begin{pmatrix}
  2\pi i u_1\,\mathrm du_2&0\\
  0&-2\pi i u_1\,\mathrm du_2
 \end{pmatrix}
 =2\pi u_1\,\mathrm du_2\,H_{\mathrm{mag}},
\]
so \cref{cusp-magnetic-su2-curvature} follows from
\cref{cusp-magnetic-line-bundle-proposition,cusp-magnetic-wedge-norm}.
The diagonal endomorphism \(H_{\mathrm{mag}}\) is a
globally defined parallel section of \(\operatorname{ad}\mathscr E_T\),
because every transition \(G_n\) commutes with it.  The forms
\(\theta_T^j\) are parallel constant forms for the descended flat metric.
Thus \(\nabla^T F_{\nabla^T}=0\), which implies
\(\mathrm d_{\nabla^T}^*F_{\nabla^T}=0\).  The curvature is nonzero because
\(D_T>0\).

The chosen Lie-algebra normalization gives
\[
 \lvert H_{\mathrm{mag}}\rvert^2
 =-\frac12\operatorname{Tr}(H_{\mathrm{mag}}^2)=1.
\]
Equations \cref{cusp-magnetic-wedge-norm,cusp-spec-exact-volume} therefore
give
\[
 \int_{F_T}\lvert F_{\nabla^T}\rvert^2\,
   \mathrm d\operatorname{vol}_{F_T}
 =4\pi^2D_T\frac1{D_T^2}=\frac{4\pi^2}{D_T},
\]
which proves \cref{cusp-magnetic-exact-action}.  Finally,
\cref{cusp-spec-D-expansion} and the elementary reciprocal expansion
\[
 \frac1{T^2+aT+O(1)}
 =\frac1{T^2}\left(1-\frac aT+O(T^{-2})\right)
\]
with \(a=2H-B\) prove
\cref{cusp-magnetic-action-first-correction}.
\end{proof}

\subsection{Constant-curvature magnetic Laplacians for arbitrary integral
nondegenerate flux}

We next give the Hilbert space, closed form, self-adjoint realization, and
oscillator spectrum in a form valid for every nondegenerate integral
alternating form on the marked four-torus.

Let
\[
 \mathsf E=\sum_{1\le j<k\le4}E_{jk}\,\mathrm du_j\wedge\mathrm du_k
 \in\bigwedge^2(\mathbb Z^4)^*,
 \qquad E_{jk}\in\mathbb Z,
\]
and let the same symbol denote its skew matrix.  For \(n\in\mathbb Z^4\),
put
\begin{equation}
 J_n^{\mathsf E}(u)
 =\exp\!\left(-2\pi i\sum_{j<k}E_{jk}n_ju_k\right).
 \label{cusp-magnetic-general-factor}
\end{equation}
The integer in
\[
 \sum_{j<k}E_{jk}m_jn_k
\]
shows directly that
\(J_m^{\mathsf E}(u+n)J_n^{\mathsf E}(u)=J_{m+n}^{\mathsf E}(u)\).
Thus \eqref{cusp-magnetic-general-factor} defines a Hermitian line bundle
\(\mathscr L_{\mathsf E,T}\).  On its pullback, set
\begin{equation}
 \nabla^{\mathsf E}
 =\mathrm d+2\pi i\sum_{j<k}E_{jk}u_j\,\mathrm du_k.
 \label{cusp-magnetic-general-connection}
\end{equation}
Exactly as in \cref{cusp-magnetic-line-bundle-proposition}, the
connection descends and
\begin{equation}
 F_{\nabla^{\mathsf E}}=2\pi i\,\mathsf E.
 \label{cusp-magnetic-general-curvature}
\end{equation}

In \(u\)-coordinates the flat metric matrix is
\begin{equation}
 G_T=P_T^{\mathsf T}P_T,
 \qquad
 \sqrt{\det G_T}=D_T.
 \label{cusp-magnetic-flat-Gram}
\end{equation}
Define the Hilbert space concretely by
\begin{align}
 \mathcal H_{\mathsf E,T}
 =\{\,&f\in L^2_{\mathrm{loc}}(\mathbb R^4):
 f(u+n)=J_n^{\mathsf E}(u)f(u)\text{ for every }n\in\mathbb Z^4\},
 \label{cusp-magnetic-Hilbert-space}\\
 \lVert f\rVert^2
 ={}&\int_{[0,1]^4}\lvert f(u)\rvert^2\sqrt{\det G_T}\,\mathrm d^4u.
 \nonumber
\end{align}
Let \(\mathcal C_{\mathsf E,T}\) be the smooth functions with the same
automorphy.  On this dense subspace define
\begin{equation}
 q_{\mathsf E,T}[f]
 =\int_{[0,1]^4}
  (G_T^{-1})^{jk}
  \overline{\nabla^{\mathsf E}_jf}
  \nabla^{\mathsf E}_kf
  \sqrt{\det G_T}\,\mathrm d^4u.
 \label{cusp-magnetic-quadratic-form}
\end{equation}

\begin{theorem}[Friedrichs realization and oscillator spectrum]
\label{cusp-magnetic-oscillator-theorem}
The form \eqref{cusp-magnetic-quadratic-form} is nonnegative and closable.
Its closed form domain is
\[
 \operatorname{Dom}\overline q_{\mathsf E,T}
 =H^1(F_T;\mathscr L_{\mathsf E,T}).
\]
Its Friedrichs operator is the
magnetic Laplacian
\[
 \Delta_{\mathsf E,T}=(\nabla^{\mathsf E})^*\nabla^{\mathsf E},
 \qquad
 \operatorname{Dom}(\Delta_{\mathsf E,T})
 =H^2(F_T;\mathscr L_{\mathsf E,T}),
\]
and it has compact resolvent.

Assume now that \(\mathsf E\) is nondegenerate.  Let \(b_1,b_2>0\) be
defined by the eigenvalues
\begin{equation}
 \operatorname{Spec}\!\left(
  G_T^{-1/2}(2\pi\mathsf E)G_T^{-1/2}\right)
 =\{ib_1,-ib_1,ib_2,-ib_2\}.
 \label{cusp-magnetic-frequency-definition}
\end{equation}
Then, as a set (with finite multiplicities),
\begin{equation}
 \operatorname{Spec}(\Delta_{\mathsf E,T})
 =\left\{(2n_1+1)b_1+(2n_2+1)b_2:
 n_1,n_2\in\mathbb Z_{\ge0}\right\}.
 \label{cusp-magnetic-oscillator-spectrum}
\end{equation}
In particular, the ground-to-first-distinct-level spacing is
\begin{equation}
 \operatorname{gap}(\Delta_{\mathsf E,T})=2\min(b_1,b_2).
 \label{cusp-magnetic-level-gap}
\end{equation}
No multiplicity formula is needed for, or asserted in, these identities.
\end{theorem}

\begin{proof}
In local unitary trivializations, the integrand in
\eqref{cusp-magnetic-quadratic-form} is the squared norm of a first-order
unitary covariant derivative.  A finite trivializing cover and a partition
of unity show that the form norm
\(\lVert f\rVert^2+q_{\mathsf E,T}[f]\) is equivalent to the usual local \(H^1\)
norm.  Its completion is therefore
\(H^1(F_T;\mathscr L_{\mathsf E,T})\), proving closability and identifying
the closed domain.  The representation theorem for closed nonnegative
forms gives the Friedrichs operator.  Its local principal symbol is
\(G_T^{-1}(\xi,\xi)\), so it is uniformly elliptic.  Standard interior
elliptic regularity on the compact manifold without boundary identifies
the operator domain with \(H^2\).  The compact inclusion
\(H^1\hookrightarrow L^2\) implies compact resolvent.

The real skew-normal-form theorem supplies a constant orthonormal frame
\(v_1,\ldots,v_4\) on the Euclidean cover in which the real two-form
\(2\pi\mathsf E\) is
\[
 b_1v^1\wedge v^2+b_2v^3\wedge v^4.
\]
Changing the sign of \(v^2\) or \(v^4\), if necessary, makes both
frequencies positive.  Let \(\nabla_a\) denote covariant differentiation
along \(v_a\).  Each \(\nabla_a\) is formally skew-adjoint on the common
invariant core \(\mathcal C_{\mathsf E,T}\); the closures used below are the
ones determined by the closed magnetic form.  On that smooth core one has
\[
 [\nabla_1,\nabla_2]=ib_1,
 \qquad
 [\nabla_3,\nabla_4]=ib_2,
\]
and all other commutators vanish.  Put
\[
 C_1=\nabla_1-i\nabla_2,
 \qquad
 C_2=\nabla_3-i\nabla_4,
 \qquad
 \Delta_j=-(\nabla_{2j-1}^2+\nabla_{2j}^2).
\]
Direct multiplication, retaining the commutator signs, gives
\begin{equation}
 C_j^*C_j=\Delta_j-b_j,
 \qquad
 C_jC_j^*=\Delta_j+b_j,
 \qquad
 [\Delta_{\mathsf E,T},C_j]=-2b_jC_j.
 \label{cusp-magnetic-ladder-identities}
\end{equation}

Because the resolvent is compact, the lowest eigenvalue has a smooth
eigenfunction \(\psi_0\neq0\).  If \(C_j\psi_0\neq0\), the last identity
in \eqref{cusp-magnetic-ladder-identities} would make it an eigenfunction
with a strictly smaller eigenvalue.  Hence \(C_1\psi_0=C_2\psi_0=0\), and
the first identity gives
\[
 \Delta_{\mathsf E,T}\psi_0=(b_1+b_2)\psi_0.
\]
Conversely, repeated application of \(C_1^*,C_2^*\) produces nonzero
eigenfunctions with eigenvalues
\((2n_1+1)b_1+(2n_2+1)b_2\).  Nonvanishing follows from
\(C_jC_j^*=C_j^*C_j+2b_j\).

It remains to exclude other eigenvalues.  For an eigenfunction \(\psi\)
with eigenvalue \(\lambda>b_1+b_2\), the operator identity
\[
 (\lambda-b_1-b_2)\psi
 =C_1^*C_1\psi+C_2^*C_2\psi
\]
expresses \(\psi\) as a sum of raising operators applied to the lower
eigenfunctions \(C_j\psi\).  Induction over the finite set of eigenvalues
below \(\lambda\) reduces every eigenfunction to raised ground states.
This proves \eqref{cusp-magnetic-oscillator-spectrum}.  Raising once in a
least-frequency plane proves \eqref{cusp-magnetic-level-gap}.
\end{proof}

\begin{corollary}[Frequency product and cusp spacing bound]
\label{cusp-magnetic-frequency-product}
Let
\[
 d=|\operatorname{Pf}(\mathsf E)|
 =\lvert E_{12}E_{34}-E_{13}E_{24}+E_{14}E_{23}\rvert>0.
\]
Then
\begin{equation}
 b_1b_2=\frac{4\pi^2d}{D_T},
 \label{cusp-magnetic-product-identity}
\end{equation}
and
\begin{align}
 \operatorname{gap}(\Delta_{\mathsf E,T})
 &\le 4\pi\sqrt{\frac d{D_T}}
 \label{cusp-magnetic-gap-bound}\\
 &=\frac{4\pi\sqrt d}{T}
 \left(1+\frac{B-2H}{2T}
 +O(T^{-2})\right).
 \nonumber
\end{align}
\end{corollary}

\begin{proof}
The determinant of the skew matrix in
\eqref{cusp-magnetic-frequency-definition} is both
\(b_1^2b_2^2\) and
\[
 \frac{\det(2\pi\mathsf E)}{\det G_T}
 =\frac{(2\pi)^4\operatorname{Pf}(\mathsf E)^2}{D_T^2}.
\]
Taking the positive square root proves
\eqref{cusp-magnetic-product-identity}.  Since
\(\min(b_1,b_2)\le\sqrt{b_1b_2}\),
\Cref{cusp-magnetic-oscillator-theorem} gives
\eqref{cusp-magnetic-gap-bound}.  Expanding the square root of
\cref{cusp-spec-D-expansion} gives the displayed first correction.
\end{proof}

\subsection{A background-gauge negative mode}

We finish with the classical background-gauge Jacobi operator at the
reducible critical connection.  Every sign below uses exactly the inner
product and the element \(H_{\mathrm{mag}}\) fixed in
\cref{cusp-magnetic-su2-normalization}.

Put
\[
 \mathscr H_T^{\mathrm{ad},1}
 :=L^2\Omega^1(F_T;\operatorname{ad}\mathscr E_T).
\]
Put also
\[
 \mathscr H_T^{\mathrm{ad},0}
 :=L^2\Omega^0(F_T;\operatorname{ad}\mathscr E_T),
 \qquad
 D_T^{\max}:=\mathrm d_{\nabla^T}^{*,\max}.
\]
For the resolvent argument these real Hilbert spaces are complexified;
all operators commute with the real structure, so their restrictions to
real sections are unchanged.
Here the maximal distributional adjoint is the closed, densely defined
operator whose domain and value are specified without an implicit
regularity assumption by
\begin{align}
 \operatorname{Dom}(D_T^{\max})
 :=\{a\in\mathscr H_T^{\mathrm{ad},1}:{}&\text{ there is a }u\in
 \mathscr H_T^{\mathrm{ad},0}\text{ such that}\notag\\[-1mm]
 &\langle a,\mathrm d_{\nabla^T}\phi\rangle_{L^2}
 =\langle u,\phi\rangle_{L^2}
 \text{ for every }\phi\in
 C^\infty(F_T;\operatorname{ad}\mathscr E_T)\},
 \label{cusp-magnetic-maximal-codifferential-domain}\\
 D_T^{\max}a:={}&u.\notag
\end{align}
The test space is dense, so \(u\) is unique.  The definition also proves
closedness directly: if \(a_n\to a\) and \(D_T^{\max}a_n\to u\) in
\(L^2\), passage to the limit in the displayed testing identity gives
\(a\in\operatorname{Dom}(D_T^{\max})\) and \(D_T^{\max}a=u\).
Consequently the background slice
\begin{equation}
 \mathscr S_T
 :=\ker D_T^{\max}
 \subset\mathscr H_T^{\mathrm{ad},1}.
 \label{cusp-magnetic-closed-background-slice}
\end{equation}
is a closed subspace.

The scalar covariant Laplacian used below is fixed as follows.  On
\(C^\infty(F_T;\operatorname{ad}\mathscr E_T)\) start with the
nonnegative form
\[
 q_T^{(0)}[\phi,\psi]
 :=\langle\mathrm d_{\nabla^T}\phi,
             \mathrm d_{\nabla^T}\psi\rangle_{L^2}.
\]
Its closure has domain
\(H^1\Omega^0(F_T;\operatorname{ad}\mathscr E_T)\), and its associated
Friedrichs operator is
\begin{equation}
 \Delta_{\nabla^T}^{(0)}
 :=\mathrm d_{\nabla^T}^*\mathrm d_{\nabla^T},
 \qquad
 \operatorname{Dom}(\Delta_{\nabla^T}^{(0)})
 =H^2\Omega^0(F_T;\operatorname{ad}\mathscr E_T).
 \label{cusp-magnetic-scalar-covariant-laplacian-domain}
\end{equation}
The \(H^2\) domain follows from elliptic regularity on the compact flat
torus; in particular, this notation always denotes the nonnegative
self-adjoint realization, not merely its differential expression.

For an adjoint-valued one-form \(a\), define the smooth bundle
endomorphism
\[
 (V_Ta)_\nu:=-2\sum_{\mu=1}^4[F_{\nu\mu},a_\mu].
\]
It is symmetric: after interchanging \(\mu,\nu\), the identities
\(F_{\mu\nu}=-F_{\nu\mu}\) and skew-adjointness of
\(\operatorname{ad}(F_{\nu\mu})\) give
\(\langle V_Ta,b\rangle=\langle a,V_Tb\rangle\).
On the full one-form Hilbert space define
\begin{equation}
 \mathfrak j_T^{\mathrm{gf}}[a,b]
 :=\langle\nabla^Ta,\nabla^Tb\rangle_{L^2}
   +\langle V_Ta,b\rangle_{L^2},
 \qquad
 \operatorname{Dom}(\mathfrak j_T^{\mathrm{gf}})
 =H^1\Omega^1(F_T;\operatorname{ad}\mathscr E_T).
 \label{cusp-magnetic-full-jacobi-form}
\end{equation}
Because \(V_T\) is a bounded symmetric endomorphism on the compact
torus, this form is closed and bounded below.  Its associated
self-adjoint operator \(\mathcal J_T^{\mathrm{gf}}\) has
\begin{equation}
 \operatorname{Dom}(\mathcal J_T^{\mathrm{gf}})
 =H^2\Omega^1(F_T;\operatorname{ad}\mathscr E_T)
 \label{cusp-magnetic-full-jacobi-domain}
\end{equation}
and differential expression
\[
 (\mathcal J_T^{\mathrm{gf}}a)_\nu
 =-\sum_{\mu=1}^4\nabla_{e_\mu}\nabla_{e_\mu}a_\nu
  -2\sum_{\mu=1}^4[F_{\nu\mu},a_\mu].
\]

\begin{proposition}[Intertwining and the reducing Coulomb slice]
\label{cusp-magnetic-jacobi-reduction-proposition}
At the Yang--Mills connection \(\nabla^T\), the following formal
identities hold on smooth sections:
\begin{equation}
 \mathcal J_T^{\mathrm{gf}}\mathrm d_{\nabla^T}
 =\mathrm d_{\nabla^T}\Delta_{\nabla^T}^{(0)},
 \qquad
 \mathrm d_{\nabla^T}^*\mathcal J_T^{\mathrm{gf}}
 =\Delta_{\nabla^T}^{(0)}\mathrm d_{\nabla^T}^*.
 \label{cusp-magnetic-jacobi-intertwining}
\end{equation}
For every \(z\in\mathbb C\setminus\mathbb R\), the correctly typed
closed-operator resolvent identity is
\begin{equation}
 D_T^{\max}(\mathcal J_T^{\mathrm{gf}}-z)^{-1}f
 =(\Delta_{\nabla^T}^{(0)}-z)^{-1}D_T^{\max}f,
 \qquad f\in\operatorname{Dom}(D_T^{\max}).
 \label{cusp-magnetic-jacobi-resolvent-intertwining}
\end{equation}
In particular, \(\mathscr S_T\) reduces
\(\mathcal J_T^{\mathrm{gf}}\).  The part of the full operator in this
subspace and its closed form are therefore
\begin{equation}
 \mathcal J_T
 :=\mathcal J_T^{\mathrm{gf}}\big|_{\mathscr S_T},
 \quad
 \operatorname{Dom}(\mathcal J_T)
 =H^2\Omega^1(F_T;\operatorname{ad}\mathscr E_T)\cap\mathscr S_T,
 \quad
 \operatorname{Dom}(\mathfrak j_T)
 =H^1\Omega^1(F_T;\operatorname{ad}\mathscr E_T)\cap\mathscr S_T
 \label{cusp-magnetic-reduced-jacobi-domain}
\end{equation}
with
\(\mathfrak j_T=\mathfrak j_T^{\mathrm{gf}}|_{\operatorname{Dom}(\mathfrak j_T)}\).
The operator \(\mathcal J_T\) is self-adjoint on \(\mathscr S_T\), and
\(\mathfrak j_T\) is its closed, lower-bounded form.
\end{proposition}

\begin{proof}
Write \(\nabla_\mu=\nabla^T_{e_\mu}\) in the constant orthonormal frame,
so
\([\nabla_\mu,\nabla_\nu]\phi=[F_{\mu\nu},\phi]\).
For a smooth adjoint-valued scalar \(\phi\), commuting the derivatives
without suppressing either curvature contribution gives
\begin{align*}
 &\bigl(\mathcal J_T^{\mathrm{gf}}\mathrm d_{\nabla^T}\phi
       -\mathrm d_{\nabla^T}\Delta_{\nabla^T}^{(0)}\phi\bigr)_\nu\\
 &\quad=-\nabla_\mu\nabla_\mu\nabla_\nu\phi
         +\nabla_\nu\nabla_\mu\nabla_\mu\phi
         -2[F_{\nu\mu},\nabla_\mu\phi]\\
 &\quad=-[\nabla_\mu F_{\mu\nu},\phi]
   =[ (\mathrm d_{\nabla^T}^*F_{\nabla^T})_\nu,\phi].
\end{align*}
The last member vanishes by
\cref{cusp-magnetic-ym-critical-point}.  This proves the first identity
in \eqref{cusp-magnetic-jacobi-intertwining}.  Since both Laplacians and
\(\mathcal J_T^{\mathrm{gf}}\) are formally symmetric, pairing it against
arbitrary smooth zero- and one-forms proves the second identity.

Fix \(z\notin\mathbb R\), take
\(f\in\operatorname{Dom}(D_T^{\max})\), and put
\(u=(\mathcal J_T^{\mathrm{gf}}-z)^{-1}f\).
The full operator domain gives \(u\in H^2\), hence
\(D_T^{\max}u\in H^1\).  Moreover
\((\mathcal J_T^{\mathrm{gf}}-z)u=f\) and the distributional extension of
the second formal identity imply
\begin{equation}
 (\Delta_{\nabla^T}^{(0)}-z)D_T^{\max}u=D_T^{\max}f
 \quad\text{in }\mathcal D'(F_T;\operatorname{ad}\mathscr E_T).
 \label{cusp-magnetic-resolvent-distributional-equation}
\end{equation}
The right side is in \(L^2\).  Elliptic regularity for the scalar
Laplacian upgrades \(D_T^{\max}u\) from \(H^1\) to \(H^2\), so it belongs
to \(\operatorname{Dom}(\Delta_{\nabla^T}^{(0)})\).  Applying the bounded
scalar resolvent to
\eqref{cusp-magnetic-resolvent-distributional-equation} proves
\eqref{cusp-magnetic-jacobi-resolvent-intertwining}; no extension of
\(D_T^{\max}\) to all of \(L^2\) is asserted or needed.

If \(f\in\mathscr S_T\), the right side of
\eqref{cusp-magnetic-jacobi-resolvent-intertwining} is zero, hence both
resolvents at \(z\) and at \(\overline z\) preserve \(\mathscr S_T\).
For \(g\in\mathscr S_T^\perp\) and \(s\in\mathscr S_T\), self-adjointness
of the full operator gives
\[
 \langle(\mathcal J_T^{\mathrm{gf}}-z)^{-1}g,s\rangle
 =\langle g,(\mathcal J_T^{\mathrm{gf}}-\overline z)^{-1}s\rangle=0.
\]
Thus the resolvent preserves both \(\mathscr S_T\) and its orthogonal
complement, which is precisely the reducing property.  The standard
part of a self-adjoint operator on a reducing subspace is self-adjoint,
and its domain is the intersection in
\eqref{cusp-magnetic-reduced-jacobi-domain}.  The reducing orthogonal
projection commutes with the spectral calculus of the lower-bounded full
operator, including its form square root; this proves the asserted form
domain and completes the proof.
\end{proof}

\begin{theorem}[Exact charged Landau negative mode]
\label{cusp-magnetic-negative-mode}
For the self-adjoint Yang--Mills Jacobi operator \(\mathcal J_T\) on the
closed background slice \(\mathscr S_T\) defined above, the differential
expression in a constant orthonormal frame \((e_\mu)_{\mu=1}^4\) is
\begin{equation}
 (\mathcal J_Ta)_\nu
 =-\sum_{\mu=1}^4\nabla_{e_\mu}\nabla_{e_\mu}a_\nu
  -2\sum_{\mu=1}^4[F_{\nu\mu},a_\mu].
 \label{cusp-magnetic-jacobi-operator}
\end{equation}
There exists a nonzero smooth real
\(\mathfrak{su}(2)\)-valued one-form \(a\) such that
\begin{equation}
 \lVert a\rVert_{L^2}=1,
 \qquad
 \mathrm d_{\nabla^T}^*a=0,
 \qquad
 \mathcal J_Ta=-\frac{4\pi}{D_T}a.
 \label{cusp-magnetic-negative-eigen-equation}
\end{equation}
Equivalently, for the action normalization
\cref{cusp-magnetic-action-normalization},
\begin{equation}
 \operatorname{YM}_T''\big|_{\nabla^T}(a,a)
 =-\frac{4\pi}{g_{\mathrm{YM}}^2D_T}.
 \label{cusp-magnetic-negative-hessian}
\end{equation}
Thus the reducible non-flat critical connection has a strictly negative
classical second-variation direction for every finite \(T\).
\end{theorem}

\begin{proof}
First isolate the exact two-dimensional magnetic problem.  Let
\[
 W=\{x_1=x_2=0\}\subset\mathbb R^4,
\]
so that \((y_1,y_2)\) are oriented orthonormal coordinates on \(W\).
By \cref{cusp-spec-real-period-matrix}, orthogonal projection to \(W\)
sends the first two marked lattice generators to
\[
 v_1=(6m_T,-q_T),
 \qquad
 v_2=(L_T,m_T),
\]
and sends the last two to zero.  Since
\(\det(v_1,v_2)=D_T>0\), projection descends to
\[
 \varpi_T:F_T\longrightarrow
 \Sigma_T:=W/(\mathbb Zv_1+\mathbb Zv_2).
\]
The factors and connection of \(\mathscr L_T^2\) depend only on
\((u_1,u_2)\).  They therefore define a Hermitian line bundle with unitary
connection \((\mathscr M_T,\nabla^{\mathscr M_T})\) on \(\Sigma_T\) such
that
\begin{equation}
 (\mathscr L_T^2,\nabla^{\mathscr L_T^2})
 =\varpi_T^*(\mathscr M_T,\nabla^{\mathscr M_T}).
 \label{cusp-magnetic-descended-root-line}
\end{equation}
Let \(\Delta_+=(\nabla^{\mathscr M_T})^*\nabla^{\mathscr M_T}\) be the
nonnegative self-adjoint Bochner Laplacian on
\(L^2(\Sigma_T;\mathscr M_T)\), defined by its closed \(H^1\) form and
with operator domain \(H^2(\Sigma_T;\mathscr M_T)\).  More precisely, if
\(c_T>0\) is the constant Riemannian volume of a fibre of \(\varpi_T\),
then
\[
 U_T:L^2(\Sigma_T;\mathscr M_T)\longrightarrow
       L^2(F_T;\mathscr L_T^2),
 \qquad U_Ts=c_T^{-1/2}\varpi_T^*s,
\]
is an isometry.  Its image
\[
 \mathscr B_T:=U_TL^2(\Sigma_T;\mathscr M_T)
\]
is the closed fibre-basic subspace.  The pulled-back realization is,
by definition,
\[
 \Delta_{+,\mathrm{bas}}:=U_T\Delta_+U_T^{-1}
 \quad\text{on }\mathscr B_T,
 \qquad
 \operatorname{Dom}(\Delta_{+,\mathrm{bas}})
 =U_TH^2(\Sigma_T;\mathscr M_T),
\]
and its closed form domain is \(U_TH^1(\Sigma_T;\mathscr M_T)\).
Only this operator on \(\mathscr B_T\) is self-adjoint.  In the formulas
below we abbreviate \(\Delta_{+,\mathrm{bas}}\) again to \(\Delta_+\);
no self-adjoint operator on the full
\(L^2(F_T;\mathscr L_T^2)\) is thereby defined or asserted.  On this
root line the curvature is
\begin{equation}
 F_{\mathscr L_T^2}
 =4\pi i\,\theta_T^1\wedge\theta_T^2
 =i\kappa_T\,\mathrm dy_1\wedge\mathrm dy_2,
 \qquad
 \kappa_T=\frac{4\pi}{D_T}>0.
 \label{cusp-magnetic-root-field-strength}
\end{equation}

We now construct an exact lowest Landau section on this actual projected
lattice.  In the marked coordinates put
\[
 \mathcal D_1:=\partial_{u_1},
 \qquad
 \mathcal D_2:=\partial_{u_2}+4\pi i u_1,
\]
and define
\begin{equation}
 \zeta_T:=\frac{q_T-6im_T}{m_T+iL_T},
 \qquad
 v_T:=-\operatorname{Im}\zeta_T
     =\frac{D_T}{m_T^2+L_T^2}>0.
 \label{cusp-magnetic-theta-parameter}
\end{equation}
Indeed,
\[
 \operatorname{Im}\bigl((q_T-6im_T)(m_T-iL_T)\bigr)
 =-(L_Tq_T+6m_T^2)=-D_T,
\]
which also checks the sign \(v_T>0\).
The series
\begin{equation}
 \Theta_T(u_1,u_2)
 =\sum_{\ell\in\mathbb Z}
 \exp\!\left(
  -2\pi i\zeta_T(\ell+u_1)^2+4\pi i\ell u_2
 \right)
 \label{cusp-magnetic-theta-series}
\end{equation}
converges normally with all derivatives, because the absolute value of
its \(\ell\)-th summand is
\(\exp(-2\pi v_T(\ell+u_1)^2)\).  Reindexing the series and using
\(\ell\in\mathbb Z\) give
\begin{equation}
 \Theta_T(u_1+1,u_2)=e^{-4\pi i u_2}\Theta_T(u_1,u_2),
 \qquad
 \Theta_T(u_1,u_2+1)=\Theta_T(u_1,u_2).
 \label{cusp-magnetic-theta-automorphy}
\end{equation}
Thus \(\Theta_T\) is a smooth section of \(\mathscr L_T^2\).  Termwise
differentiation gives
\begin{equation}
 (\mathcal D_1+\zeta_T\mathcal D_2)\Theta_T=0.
 \label{cusp-magnetic-theta-marked-equation}
\end{equation}
The exact period block on the \((y_1,y_2)\)-plane is
\[
 \begin{pmatrix}6m_T&L_T\\-q_T&m_T\end{pmatrix}.
\]
Its inverse transpose gives the physical covariant derivatives
\begin{equation}
 \nabla_{y_1}^+
  =\frac{m_T\mathcal D_1+q_T\mathcal D_2}{D_T},
 \qquad
 \nabla_{y_2}^+
  =\frac{-L_T\mathcal D_1+6m_T\mathcal D_2}{D_T}.
 \label{cusp-magnetic-physical-root-derivatives}
\end{equation}
Consequently
\begin{equation}
 \nabla_{y_1}^+-i\nabla_{y_2}^+
 =\frac{m_T+iL_T}{D_T}
   (\mathcal D_1+\zeta_T\mathcal D_2).
 \label{cusp-magnetic-annihilator-conversion}
\end{equation}
Orthogonality of the \(u_2\)-Fourier modes also gives the nonzero norm
\[
 \int_0^1\!\!\int_0^1\lvert\Theta_T\rvert^2
 \,\mathrm du_1\,\mathrm du_2
 =\sum_{\ell\in\mathbb Z}\int_0^1
   e^{-4\pi v_T(\ell+u_1)^2}\,\mathrm du_1
 =\frac1{2\sqrt{v_T}}.
\]
Set
\[
 \phi_T:=\left(\frac{2\sqrt{v_T}}{D_T}\right)^{1/2}\Theta_T
\]
and pull it back by \(\varpi_T\).  Since the marked-coordinate volume
element is \(D_T\,\mathrm d^4u\), it has unit \(L^2(F_T)\)-norm.  Its
covariant derivatives in the \(x_1,x_2\) directions vanish, and
\begin{equation}
 (\nabla_{y_1}^+-i\nabla_{y_2}^+)\phi_T=0,
 \qquad
 \Delta_+\phi_T=\kappa_T\phi_T,
 \qquad
 \Delta_+\phi_T
 =-\bigl((\nabla_{y_1}^+)^2+(\nabla_{y_2}^+)^2\bigr)\phi_T.
 \label{cusp-magnetic-ground-equation}
\end{equation}
Indeed, \cref{cusp-magnetic-physical-root-derivatives} and
\([\mathcal D_1,\mathcal D_2]=4\pi i\) give
\[
 [\nabla_{y_1}^+,\nabla_{y_2}^+]
 =\frac{m_T(6m_T)-q_T(-L_T)}{D_T^2}\,4\pi i
 =i\kappa_T.
\]
Direct multiplication then gives
\[
 (\nabla_{y_1}^+-i\nabla_{y_2}^+)^*
 (\nabla_{y_1}^+-i\nabla_{y_2}^+)=\Delta_+-\kappa_T,
\]
and the second equation in \eqref{cusp-magnetic-ground-equation} follows
from the first.

In \(\mathfrak{sl}_2(\mathbb C)\), put
\[
 X_+=\begin{pmatrix}0&1\\0&0\end{pmatrix},
 \qquad
 X_-=\begin{pmatrix}0&0\\1&0\end{pmatrix}.
\]
Then
\[
 [H_{\mathrm{mag}},X_+]=2iX_+,
 \qquad
 [H_{\mathrm{mag}},X_-]=-2iX_-.
\]
The \(X_+\)-root line in
\(\operatorname{ad}(\mathscr L_T\oplus\mathscr L_T^{-1})
 \otimes\mathbb C\) is \(\mathscr L_T^2\),
so
\[
 a^+=\phi_TX_+(\mathrm dy_1-i\,\mathrm dy_2)
\]
is a globally defined complexified adjoint-valued one-form.  The first
equation in \eqref{cusp-magnetic-ground-equation} gives
\[
 \mathrm d_{\nabla^T}^*a^+
 =-(\nabla_{y_1}^+-i\nabla_{y_2}^+)\phi_T\,X_+=0.
\]

Put \(\mathcal B_T:=2\pi/D_T\), so that
\(\kappa_T=2\mathcal B_T\) and
\(F_{y_1y_2}=\mathcal B_TH_{\mathrm{mag}}\).  The two nonzero components
of \(a^+\) are \(a_{y_1}^+=\phi_TX_+\) and
\(a_{y_2}^+=-i\phi_TX_+\).  Substitution into
\eqref{cusp-magnetic-jacobi-operator}, with no suppressed charge factor,
gives
\begin{align*}
 (\mathcal J_Ta^+)_{y_1}
 &=\kappa_T\phi_TX_+
   -2[\mathcal B_TH_{\mathrm{mag}},-i\phi_TX_+]
  =-\kappa_T\phi_TX_+,\\
 (\mathcal J_Ta^+)_{y_2}
 &=-i\kappa_T\phi_TX_+
   -2[-\mathcal B_TH_{\mathrm{mag}},\phi_TX_+]
  =i\kappa_T\phi_TX_+.
\end{align*}
Thus \(\mathcal J_Ta^+=-\kappa_Ta^+\).

For completeness, formula \eqref{cusp-magnetic-jacobi-operator} follows by
expanding
\(F_{\nabla^T+ta}=F_{\nabla^T}+t\,\mathrm d_{\nabla^T}a+
\tfrac12t^2[a\wedge a]\), integrating the quadratic form by parts on
the background slice, and using the flat-base Weitzenbock identity.  The
Hodge Laplacian contributes
\(-[F_{\nu\mu},a_\mu]\), and the curvature term in the second variation
contributes the identical term, giving the coefficient \(-2\).

Finally take the anti-Hermitian real form
\begin{equation}
 a:=\frac1{\sqrt2}\bigl(a^+-(a^+)^\dagger\bigr)
  =\frac1{\sqrt2}\left(
   \phi_TX_+(\mathrm dy_1-i\,\mathrm dy_2)
   -\overline{\phi_T}X_-(\mathrm dy_1+i\,\mathrm dy_2)\right).
 \label{cusp-magnetic-real-negative-mode}
\end{equation}
It is nonzero, traceless, and anti-Hermitian.  For
\(X(z):=zX_+-\overline zX_-\), the fixed inner product gives
\(-\tfrac12\operatorname{Tr}(X(z)^2)=\lvert z\rvert^2\).  The two
orthonormal one-form components in
\cref{cusp-magnetic-real-negative-mode} therefore have total pointwise
norm \(\lvert a\rvert^2=\lvert\phi_T\rvert^2\), so
\(\lVert a\rVert_{L^2}=1\).  Complex conjugation and
Hermitian adjunction preserve both the gauge equation and the real Jacobi
operator, so \eqref{cusp-magnetic-negative-eigen-equation} follows from
the calculation for \(a^+\).  The Hessian identity follows from
\(\operatorname{YM}_T''(a,a)
=g_{\mathrm{YM}}^{-2}\langle a,\mathcal J_Ta\rangle_{L^2}\)
on the background slice.
\end{proof}

\section{The native full-flux magnetic sector}
\label{cusp-native-full-flux-section}

This section retains the canonical radial parameter, period matrix, Euclidean
quotient metric, marked lattice, and every sign already fixed in
\cref{cusp-spec-section}.  It treats the monodromy-invariant integral class
that uses both elementary alternating blocks.  It is therefore distinct from
the unit-flux plane treated in \cref{cusp-magnetic-section}.  No metric
rescaling, finite covering, rectangular replacement of the lattice, or
orientation suppression occurs below.

For \(T>T_0\), retain
\begin{equation}
 \begin{gathered}
 L_T=T+H_T=\operatorname{Im}\tau_T,
 \qquad q_T=L_T-B_T=-\operatorname{Im}\beta_T,\\
 m_T=\operatorname{Im}\mu_T,
 \qquad D_T=L_Tq_T+6m_T^2>0,
 \end{gathered}
 \label{cusp-native-full-flux-canonical-scalars}
\end{equation}
and
write
\[
 \eta_T^{\mathrm{re}}:=\operatorname{Re}h_T.
\]
This scalar is denoted by \(\eta_T\) in
\cref{cusp-spec-section}; the superscript is retained throughout this
section so that it cannot be confused with the native integral two-form
introduced below.  The complete period matrix is
\begin{equation}
 P_T=
 \begin{pmatrix}
  6\alpha_T&\eta_T^{\mathrm{re}}&1&0\\
  6m_T&L_T&0&0\\
  c_T&\alpha_T&0&1\\
  -q_T&m_T&0&0
 \end{pmatrix},
 \qquad \det P_T=-D_T.
 \label{cusp-native-full-flux-period-matrix}
\end{equation}
Thus \(F_T=\mathbb R^4/P_T\mathbb Z^4\) has the unrescaled Euclidean quotient
metric and volume \(D_T\).  In the marking coordinates \(y\in\mathbb R^4\),
the map \(y\mapsto x=P_Ty\) pulls the Euclidean metric back to
\(G_T=P_T^{\mathsf T}P_T\).

The limiting cusp constants remain
\begin{equation}
 m=\operatorname{Im}\mu(0),\qquad
 H=\operatorname{Im}h(0),\qquad
 B=\operatorname{Im}b(0),
 \label{cusp-native-full-flux-limiting-constants}
\end{equation}
so holomorphy at \(e^{-2\pi T}=0\) gives
\begin{equation}
 \begin{aligned}
 m_T&=m+O(e^{-2\pi T}),&
 H_T&=H+O(e^{-2\pi T}),\\
 B_T&=B+O(e^{-2\pi T}),&
 D_T&=T^2+(2H-B)T\\
 &&&\quad +(H^2-BH+6m^2)+O(Te^{-2\pi T}).
 \end{aligned}
 \label{cusp-native-full-flux-cusp-expansions}
\end{equation}

\subsection{The invariant class, its cocycle, and both orientations}

In the ordered marking
\(\mathbb Z\langle\gamma,u,w,\delta\rangle\), put
\begin{equation}
 \eta=u\wedge w+6\gamma\wedge\delta.
 \label{cusp-native-full-flux-class}
\end{equation}
The two finite monodromy generators act by
\begin{align*}
T_1\gamma&=\gamma,&T_1u&=-u-w,&
T_1w&=-6\gamma+u,&T_1\delta&=2\gamma+u+w+\delta,\\
T_2\gamma&=\gamma,&T_2u&=6\gamma+w,&
T_2w&=-u,&T_2\delta&=-3\gamma+u+\delta.
\end{align*}
Directly, without passing to a normalized representative,
\begin{align*}
(\bigwedge^2T_1)\eta
&=(-u-w)\wedge(-6\gamma+u)
 +6\gamma\wedge(2\gamma+u+w+\delta)\\
&=u\wedge w+6\gamma\wedge\delta,\\
(\bigwedge^2T_2)\eta
&=(6\gamma+w)\wedge(-u)
 +6\gamma\wedge(-3\gamma+u+\delta)\\
&=u\wedge w+6\gamma\wedge\delta.
\end{align*}
Hence \(\eta\) is invariant under both generators.

On \(Y=\mathbb R^4/\mathbb Z^4\), with ordered coordinates
\(y^1,y^2,y^3,y^4\), its constant representative and alternating matrix are
\begin{equation}
 \eta_T=\mathrm dy^2\wedge\mathrm dy^3
       +6\,\mathrm dy^1\wedge\mathrm dy^4,
 \qquad
 E_\eta=
 \begin{pmatrix}
 0&0&0&6\\
 0&0&1&0\\
 0&-1&0&0\\
 -6&0&0&0
 \end{pmatrix}.
 \label{cusp-native-full-flux-form-matrix}
\end{equation}
In the paired basis \((e_2,e_3;e_1,e_4)\), the alternating elementary
divisors are \((1,6)\); consequently
\begin{equation}
 \operatorname{SNF}(E_\eta)=\operatorname{diag}(1,1,6,6),
 \qquad \operatorname{Pf}(E_\eta)=6,
 \qquad \det E_\eta=36.
 \label{cusp-native-full-flux-smith-pfaffian}
\end{equation}

The actual Hermitian line bundle is the quotient of
\(\mathbb R^4\times\mathbb C\) by
\begin{align}
 (y,z)&\longmapsto(y+e_1,e^{12\pi i y^4}z),&
 (y,z)&\longmapsto(y+e_2,e^{2\pi i y^3}z),
 \label{cusp-native-full-flux-cocycle}\\
 (y,z)&\longmapsto(y+e_3,z),&
 (y,z)&\longmapsto(y+e_4,z).
 \notag
\end{align}
The only possibly nontrivial generator commutators acquire the factors
\(e^{12\pi i}=1\) and \(e^{2\pi i}=1\), so these maps define a
\(\mathbb Z^4\)-action.  Denote the quotient by
\(\mathscr L_{\eta,T}\to F_T\), using the marking
\([y]\mapsto[P_Ty]\).  The connection
\begin{equation}
 \nabla_{\eta,T}
 =\mathrm d-2\pi i(6y^1\,\mathrm dy^4+y^2\,\mathrm dy^3)
 \label{cusp-native-full-flux-connection}
\end{equation}
descends through the displayed cocycle.  For instance translation by \(e_1\)
changes its connection form by \(-12\pi i\,\mathrm dy^4\), which equals the
negative logarithmic derivative of \(e^{12\pi i y^4}\); the \(e_2\) check is
identical.  Therefore
\begin{equation}
 F_{\nabla_{\eta,T}}=-2\pi i\eta_T,
 \qquad c_1(\mathscr L_{\eta,T})
 =\frac{i}{2\pi}F_{\nabla_{\eta,T}}=\eta_T.
 \label{cusp-native-full-flux-curvature}
\end{equation}

There are two opposite orientations which must not be conflated.  Write
\begin{equation}
 \operatorname{or}_{\mathrm{mark}}
 =\mathrm dy^1\wedge\mathrm dy^2\wedge
  \mathrm dy^3\wedge\mathrm dy^4
 \label{cusp-native-full-flux-marking-orientation}
\end{equation}
and let \(\operatorname{or}_{x}\) be the native complex orientation of
\((z_1,z_2)=(x_1+iy_1,x_2+iy_2)\).  Then
\begin{equation}
 P_T^*\operatorname{or}_{x}
 =-D_T\operatorname{or}_{\mathrm{mark}},
 \label{cusp-native-full-flux-orientation-reversal}
\end{equation}
whereas
\begin{equation}
 \frac12\eta_T^2=6\operatorname{or}_{\mathrm{mark}}.
 \label{cusp-native-full-flux-square}
\end{equation}
Consequently
\begin{equation}
 \frac12\langle\eta_T^2,[F_T]_{\mathrm{mark}}\rangle=6,
 \qquad
 \frac12\langle\eta_T^2,[F_T]_{x\text{-}\mathrm{complex}}\rangle=-6.
 \label{cusp-native-full-flux-two-orientation-pairings}
\end{equation}

\subsection{The exact physical frequencies}

In the physical Euclidean coordinates the two-form matrix is
\begin{equation}
 \Omega_T=P_T^{-\mathsf T}E_\eta P_T^{-1}
 =\frac1{D_T}
 \begin{pmatrix}
 0&-q_T&0&-6m_T\\
 q_T&0&6m_T&0\\
 0&-6m_T&0&6L_T\\
 6m_T&0&-6L_T&0
 \end{pmatrix}.
 \label{cusp-native-full-flux-physical-matrix}
\end{equation}
After the exact permutation
\((x^1,x^2,x^3,x^4)\mapsto(x^1,x^3;x^2,x^4)\), its off-diagonal block is
\begin{equation}
 K_T=\begin{pmatrix}-q_T&-6m_T\\-6m_T&6L_T\end{pmatrix},
 \qquad \det K_T=-6D_T.
 \label{cusp-native-full-flux-K}
\end{equation}
Define the magnetic scalars
\begin{equation}
 S_T:=6L_T-q_T,
 \qquad
 \Xi_T:=\sqrt{(6L_T+q_T)^2+144m_T^2}.
 \label{cusp-native-full-flux-S-Xi}
\end{equation}
Here \(S_T\) has no argument and is distinct from the affine residual
\(S_T(k)\) used in the dual-lattice calculation of
\cref{cusp-spec-exact-dual-lattice}.  The signed eigenvalues of \(K_T\) are
\begin{equation}
 \kappa_{T,\pm}=\frac{S_T\pm\Xi_T}{2}.
 \label{cusp-native-full-flux-kappas}
\end{equation}
Direct expansion, with no asymptotic replacement, gives
\begin{equation}
 \Xi_T^2-S_T^2
 =24(L_Tq_T+6m_T^2)=24D_T>0.
 \label{cusp-native-full-flux-Xi-identity}
\end{equation}
Thus \(\kappa_{T,+}>0>\kappa_{T,-}\).  The two positive charge-one
magnetic frequencies are
\begin{equation}
 \boxed{
 \mathcal B_{T,+}=\frac{\pi(S_T+\Xi_T)}{D_T},
 \qquad
 \mathcal B_{T,-}=\frac{\pi(\Xi_T-S_T)}{D_T}.}
 \label{cusp-native-full-flux-frequencies}
\end{equation}
They obey the exact identities
\begin{align}
 \mathcal B_{T,+}\mathcal B_{T,-}
 &=\frac{24\pi^2}{D_T},
 \label{cusp-native-full-flux-frequency-product}\\
 \mathcal B_{T,+}+\mathcal B_{T,-}
 &=\frac{2\pi\Xi_T}{D_T},&
 \mathcal B_{T,+}-\mathcal B_{T,-}
 &=\frac{2\pi S_T}{D_T},
 \label{cusp-native-full-flux-frequency-sum-difference}\\
 \mathcal B_{T,+}&=\frac{24\pi}{\Xi_T-S_T},&
 \mathcal B_{T,-}&=\frac{24\pi}{\Xi_T+S_T}.
 \label{cusp-native-full-flux-rationalized-frequencies}
\end{align}
The last line follows by multiplying by the conjugate factor and using
\cref{cusp-native-full-flux-Xi-identity}.

Keeping the three independent limiting coefficients \(H,B,m\), one obtains
\begin{align}
 \mathcal B_{T,-}
 &=\frac{2\pi}{T}-\frac{2\pi H}{T^2}
 +\frac{2\pi}{T^3}\left(H^2-\frac67m^2\right)+O(T^{-4}),
 \label{cusp-native-full-flux-Bminus-expansion}\\
 \mathcal B_{T,+}
 &=\frac{12\pi}{T}+\frac{12\pi(B-H)}{T^2}
 +\frac{12\pi}{T^3}
 \left((B-H)^2-\frac{36}{7}m^2\right)+O(T^{-4}),
 \label{cusp-native-full-flux-Bplus-expansion}\\
 \rho_T:=\frac{\mathcal B_{T,+}}{\mathcal B_{T,-}}
 &=6+\frac{6B}{T}
 +\frac{6B^2-6BH-180m^2/7}{T^2}+O(T^{-3}).
 \label{cusp-native-full-flux-ratio-expansion}
\end{align}
These follow by substituting \cref{cusp-native-full-flux-cusp-expansions}
into \cref{cusp-native-full-flux-rationalized-frequencies} and expanding the
reciprocals; no sign of \(H,B,m\) is assumed.  In particular, on a sufficiently
deep cusp tail,
\begin{equation}
 S_T>0,
 \qquad \mathcal B_{T,+}>\mathcal B_{T,-}>0,
 \qquad 5<\rho_T<7.
 \label{cusp-native-full-flux-deep-tail-window}
\end{equation}

\subsection{Joint Landau levels and their exact multiplicities}

Let \(\Delta_{\eta,T}=\nabla_{\eta,T}^*\nabla_{\eta,T}\) be the Friedrichs
magnetic Laplacian on \(L^2(F_T,\mathscr L_{\eta,T})\).  Its form domain is
\(H^1(F_T;\mathscr L_{\eta,T})\), its operator domain is
\(H^2(F_T;\mathscr L_{\eta,T})\), and compactness of the torus gives compact
resolvent.  We now prove the multiplicity rather than infer it from the
Pfaffian alone.

Define \(A_T^{\mathrm{pol}}\) on the marked tangent space by
\begin{equation}
 G_T(A_T^{\mathrm{pol}}v,w)=2\pi E_\eta(v,w),
 \qquad
 Q_T^{\mathrm{pol}}=\bigl(-(A_T^{\mathrm{pol}})^2\bigr)^{1/2},
 \qquad
 J_{\eta,T}=A_T^{\mathrm{pol}}(Q_T^{\mathrm{pol}})^{-1}.
 \label{cusp-native-full-flux-polar-complex-structure}
\end{equation}
The endomorphism \(A_T^{\mathrm{pol}}\) is \(G_T\)-skew-adjoint and invertible,
so \(Q_T^{\mathrm{pol}}\) is positive, commutes with
\(A_T^{\mathrm{pol}}\), and
\begin{equation}
 J_{\eta,T}^2=-1,
 \qquad
 E_\eta(J_{\eta,T}v,J_{\eta,T}w)=E_\eta(v,w),
 \qquad
 2\pi E_\eta(v,J_{\eta,T}v)
 =G_T(Q_T^{\mathrm{pol}}v,v)>0.
 \label{cusp-native-full-flux-polar-properties}
\end{equation}
For every signed charge \(q\in\mathbb Z\setminus\{0\}\), use the convention
\begin{equation}
 \mathscr L_{\eta,T}^{\otimes q}=
 \begin{cases}
  (\mathscr L_{\eta,T})^{\otimes q},&q>0,\\
  (\mathscr L_{\eta,T}^{\vee})^{\otimes(-q)},&q<0,
 \end{cases}
 \qquad
 A_{q,T}^{\mathrm{pol}}=qA_T^{\mathrm{pol}},
 \qquad
 Q_{q,T}^{\mathrm{pol}}
 =\bigl(-(A_{q,T}^{\mathrm{pol}})^2\bigr)^{1/2}
 =|q|Q_T^{\mathrm{pol}},
 \label{cusp-native-full-flux-signed-charge-bundle}
\end{equation}
and define
\begin{equation}
 J_{q,T}=A_{q,T}^{\mathrm{pol}}(Q_{q,T}^{\mathrm{pol}})^{-1}
 =\operatorname{sgn}(q)J_{\eta,T}.
 \label{cusp-native-full-flux-signed-charge-complex-structure}
\end{equation}
All these coefficients are constant, hence the Nijenhuis tensor of
\(J_{q,T}\) vanishes and the marking transports it to a complex structure on
\(F_T\).  The induced connection \(\nabla_{\eta,T}^{(q)}\) has
\begin{equation}
 F_{\nabla_{\eta,T}^{(q)}}=-2\pi i q\eta_T,
 \qquad
 2\pi qE_\eta(v,J_{q,T}v)
 =|q|G_T(Q_T^{\mathrm{pol}}v,v)>0.
 \label{cusp-native-full-flux-signed-charge-positivity}
\end{equation}
Thus its curvature is of type \((1,1)\), and
\((\nabla_{\eta,T}^{(q)})^{0,1}\) makes
\(\mathscr L_{\eta,T}^{\otimes q}\) a positive Appell--Humbert line bundle on
the auxiliary complex torus \((F_T,J_{q,T})\).

The auxiliary orientation is fixed, rather than repaired by an absolute value.
For \(q>0\), compatibility identifies the \(J_{q,T}\)-complex orientation
with the \(E_\eta\)-symplectic orientation; for \(q<0\), replacing \(J\) by
\(-J\) preserves complex orientation in complex dimension two.  Hence
\begin{equation}
 [F_T]_{J_{q,T}}=[F_T]_{\mathrm{mark}}
 =-[F_T]_{x\text{-}\mathrm{complex}}.
 \label{cusp-native-full-flux-auxiliary-native-orientations}
\end{equation}
This auxiliary positive structure is not the native complex structure inherited
from the \(S^6\) construction.  Indeed, with \(E_\eta=\operatorname{Im}H_{\eta,T}\),
the native Appell--Humbert calculation is
\begin{equation}
 G_{\eta,T}=
 \begin{pmatrix}
  2L_T&2m_T\\
  2m_T&-q_T/3
 \end{pmatrix},
 \qquad
 \det G_{\eta,T}=-\frac23D_T<0,
 \qquad H_{\eta,T}=2G_{\eta,T}^{-1}.
 \label{cusp-native-full-flux-native-hermitian-signature}
\end{equation}
Consequently \(H_{\eta,T}\), and also \(qH_{\eta,T}\) for every nonzero
\(q\), has signature \((1,1)\).  The native line bundle is holomorphic because
\(q\eta_T\) is of type \((1,1)\), but it is indefinite and is not the positive
bundle just constructed.

The tangent and canonical bundles of the auxiliary complex torus are trivial.
Kodaira vanishing and Hirzebruch--Riemann--Roch therefore give
\begin{equation}
 \begin{aligned}
 \dim_{\mathbb C}H^0\bigl((F_T,J_{q,T}),
                   \mathscr L_{\eta,T}^{\otimes q}\bigr)
 &=\int_{[F_T]_{J_{q,T}}}\frac{(q\eta_T)^2}{2}
  =\int_{[F_T]_{\mathrm{mark}}}\frac{(q\eta_T)^2}{2}
  =6q^2,\\
 \dim_{\mathbb C}H^0\bigl((F_T,J_{\eta,T}),\mathscr L_{\eta,T}\bigr)
 &=6,\\
 \dim_{\mathbb C}H^0\bigl((F_T,J_{2,T}),
                   \mathscr L_{\eta,T}^{\otimes2}\bigr)
 &=24.
 \end{aligned}
 \label{cusp-native-full-flux-vacuum-dimensions}
\end{equation}
No absolute value occurs in this auxiliary Hirzebruch--Riemann--Roch integral.
Equivalently, the positive polarization has elementary divisors
\((|q|,6|q|)\), and the Appell--Humbert theta theorem gives their product
\(6q^2\); see
\href{https://math.univ-cotedazur.fr/~beauvill/pubs/thetaon.pdf}%
{Beauville, Theorems~2.6 (pp.~107--108) and~3.5 (p.~112)}.
On the native complex torus the same signed Hirzebruch--Riemann--Roch
calculation instead gives
\begin{equation}
 \chi\bigl((F_T,J_{\mathrm{native},T}),
           \mathscr L_{\eta,T}^{\otimes q}\bigr)
 =\int_{[F_T]_{x\text{-}\mathrm{complex}}}
       \frac{(q\eta_T)^2}{2}
 =-6q^2,
 \label{cusp-native-full-flux-native-signed-hrr}
\end{equation}
and positivity vanishing is unavailable there.  Thus the positive auxiliary
section count and the native signed Euler characteristic concern distinct
Dolbeault problems and distinct complex orientations.

Choose a \(G_T\)-orthonormal frame
\(e_+,f_+=J_{\eta,T}e_+,e_-,f_-=J_{\eta,T}e_-\) which diagonalizes
\(Q_T^{\mathrm{pol}}\), labelled so its positive eigenvalues are
\(\mathcal B_{T,+},\mathcal B_{T,-}\).  Put
\begin{equation}
 a_\sigma=
 \frac{\nabla_{e_\sigma}+i\nabla_{f_\sigma}}
      {\sqrt{2\mathcal B_{T,\sigma}}},
 \qquad
 a_\sigma^*=
 \frac{-\nabla_{e_\sigma}+i\nabla_{f_\sigma}}
      {\sqrt{2\mathcal B_{T,\sigma}}},
 \qquad \sigma\in\{+,-\}.
 \label{cusp-native-full-flux-ladders}
\end{equation}
The curvature commutators give
\begin{equation}
 [a_\sigma,a_\tau^*]=\delta_{\sigma\tau},
 \qquad [a_\sigma,a_\tau]=[a_\sigma^*,a_\tau^*]=0,
 \label{cusp-native-full-flux-CCR}
\end{equation}
and, for \(N_\sigma=a_\sigma^*a_\sigma\),
\begin{equation}
 \Delta_{\eta,T}
 =\mathcal B_{T,+}(2N_++1)+\mathcal B_{T,-}(2N_-+1).
 \label{cusp-native-full-flux-number-operator}
\end{equation}
The simultaneous vacuum
\(\ker a_+\cap\ker a_-=H^0((F_T,J_{\eta,T}),\mathscr L_{\eta,T})\)
has complex dimension six by
\cref{cusp-native-full-flux-vacuum-dimensions}.

\begin{theorem}[Complete native joint Landau decomposition]
\label{cusp-native-full-flux-landau-theorem}
For every \(n_+,n_-\in\mathbb Z_{\ge0}\), the joint eigenspace
\begin{equation}
 \mathcal H_{T;n_+,n_-}
 =\ker(N_+-n_+)\cap\ker(N_--n_-)
 \label{cusp-native-full-flux-joint-space}
\end{equation}
has complex dimension six and energy
\begin{equation}
 E_{T;n_+,n_-}
 =(2n_++1)\mathcal B_{T,+}
 +(2n_-+1)\mathcal B_{T,-}.
 \label{cusp-native-full-flux-joint-energy}
\end{equation}
Moreover
\begin{equation}
 L^2(F_T,\mathscr L_{\eta,T})
 =\widehat{\bigoplus}_{n_+,n_-\ge0}\mathcal H_{T;n_+,n_-},
 \label{cusp-native-full-flux-complete-decomposition}
\end{equation}
and the map from the vacuum to the joint level is the unitary
\begin{equation}
 \frac{(a_+^*)^{n_+}(a_-^*)^{n_-}}{\sqrt{n_+!n_-!}}.
 \label{cusp-native-full-flux-ladder-isomorphism}
\end{equation}
For a numerical eigenvalue, the multiplicity is six times the number of
joint pairs in \cref{cusp-native-full-flux-joint-energy} producing that
number.
\end{theorem}

\begin{proof}
The elliptic operator has compact resolvent, and the two commuting number
operators preserve its finite-dimensional eigenspaces; they can therefore be
diagonalized simultaneously.  If \(N_\sigma\psi=\lambda_\sigma\psi\), then
\(\lambda_\sigma\ge0\).  Repeated application of \(a_\sigma\) lowers the
Laplacian energy by \(2\mathcal B_{T,\sigma}>0\), so it terminates.  If
\(a_\sigma^r\psi\ne0\) and \(a_\sigma^{r+1}\psi=0\), then
\[
 0=\|a_\sigma^{r+1}\psi\|^2
  =(\lambda_\sigma-r)\|a_\sigma^r\psi\|^2,
\]
and hence \(\lambda_\sigma=r\in\mathbb Z_{\ge0}\).  Lowering both indices
reaches the simultaneous vacuum.

The commutator identities give, for a vacuum vector \(v\),
\[
 \|(a_\sigma^*)^nv\|^2=n!\|v\|^2,
 \qquad a_\sigma^n(a_\sigma^*)^nv=n!v.
\]
On the \(n\)-th level they also give
\((a_\sigma^*)^na_\sigma^n=n!\).  Operators in the two planes commute, so
\cref{cusp-native-full-flux-ladder-isomorphism} is unitary with inverse
\(a_+^{n_+}a_-^{n_-}/\sqrt{n_+!n_-!}\).  This proves the dimension, energy,
and completeness assertions.  Orthogonal joint eigenspaces with coincident
numerical energies must be added, which proves the final sentence.
\end{proof}

The compact-resolvent lowering argument proves that no additional joint
eigenspaces occur.  Independently, this is exactly the case \(m=2\), \(k=1\)
of
\href{https://www.math.hu-berlin.de/~schueth/quantumtori.pdf}%
{Gordon--Kirwin--Schueth--Webb, Proposition~4.5, article pp.~9--10
(volume pp.~175--176)},
\href{https://doi.org/10.1090/pspum/084/1355}%
{doi:10.1090/pspum/084/1355}.  Their signed tensor-power statement gives
every oscillator tuple multiplicity
\(|k|^mV_{E_\eta}=|k|^2\operatorname{Pf}(E_\eta)=6k^2\), so it also supplies
the external completeness check for every \(k\ne0\).

The exact ground energy and first distinct spacing are therefore
\begin{align}
 E_0(T)&=\mathcal B_{T,+}+\mathcal B_{T,-}
 =\frac{2\pi\Xi_T}{D_T}
 =\frac{14\pi}{T}+\frac{\pi(12B-14H)}{T^2}+O(T^{-3}),
 \label{cusp-native-full-flux-ground-energy}\\
 g_0(T)&=2\min\{\mathcal B_{T,+},\mathcal B_{T,-}\}
 =\frac{2\pi(\Xi_T-|S_T|)}{D_T}
 =\frac{48\pi}{\Xi_T+|S_T|}.
 \label{cusp-native-full-flux-global-gap}
\end{align}
On the deep tail \cref{cusp-native-full-flux-deep-tail-window}, this is
\(g_0(T)=2\mathcal B_{T,-}\), with expansion
\begin{equation}
 g_0(T)=\frac{4\pi}{T}-\frac{4\pi H}{T^2}
 +\frac{4\pi}{T^3}\left(H^2-\frac67m^2\right)+O(T^{-4}).
 \label{cusp-native-full-flux-gap-expansion}
\end{equation}

\begin{corollary}[Heat trace and exact fixed-energy count]
\label{cusp-native-full-flux-heat-count}
For \(t>0\),
\begin{equation}
 \operatorname{Tr}(e^{-t\Delta_{\eta,T}})
 =\frac{3}{2\sinh(t\mathcal B_{T,+})
                 \sinh(t\mathcal B_{T,-})}
 =\frac{D_T}{16\pi^2t^2}\bigl(1+O_t(T^{-2})\bigr).
 \label{cusp-native-full-flux-heat-trace}
\end{equation}
For \(E>0\), let \(N_{\eta,T}(E)\) count complex eigenvalues not exceeding
\(E\), with multiplicity.  Then
\begin{equation}
 N_{\eta,T}(E)=
 \begin{cases}
 0,&E<E_0(T),\\[1mm]
 6\displaystyle\sum_{n_+=0}^{\left\lfloor
 (E-E_0(T))/(2\mathcal B_{T,+})\right\rfloor}
 \left(1+\left\lfloor
 \dfrac{E-E_0(T)-2\mathcal B_{T,+}n_+}
       {2\mathcal B_{T,-}}
 \right\rfloor\right),&E\ge E_0(T).
 \end{cases}
 \label{cusp-native-full-flux-counting-function}
\end{equation}
For fixed \(E>0\), as \(T\to\infty\),
\begin{equation}
 N_{\eta,T}(E)
 =\frac{D_T}{32\pi^2}(E-E_0(T))^2+O_E(T)
 =\frac{E^2D_T}{32\pi^2}+O_E(T).
 \label{cusp-native-full-flux-counting-asymptotic}
\end{equation}
\end{corollary}

\begin{proof}
The heat trace is the product of the two geometric series in
\cref{cusp-native-full-flux-joint-energy}, multiplied by six.  Expanding both
hyperbolic sines and using
\cref{cusp-native-full-flux-frequency-product} proves its asymptotic.
The counting formula enumerates the lattice points
\[
 (n_+,n_-)\in\mathbb R_{\ge0}^2,
 \qquad
 2\mathcal B_{T,+}n_++2\mathcal B_{T,-}n_-
 \le E-E_0(T).
\]
Its triangular area is
\((E-E_0(T))^2/(8\mathcal B_{T,+}\mathcal B_{T,-})\).  The difference
between the lattice count and the area is bounded by a constant times the
perimeter plus one, hence is \(O_E(T)\).  Multiplying by six and using the
exact frequency product gives the first asymptotic; the second follows from
\(E_0(T)=O(T^{-1})\) and \(D_T=O(T^2)\).
\end{proof}

\subsection{The split \texorpdfstring{\(SU(2)\)}{SU(2)} critical connection}

Put
\begin{equation}
 \mathscr E_{\eta,T}=\mathscr L_{\eta,T}\oplus
 \mathscr L_{\eta,T}^{-1},
 \qquad
 A_{\eta,T}=\nabla_{\eta,T}\oplus\nabla_{\eta,T}^{-1}.
 \label{cusp-native-full-flux-split-bundle-connection}
\end{equation}
The determinant is canonically trivial and
\begin{equation}
 c(\mathscr E_{\eta,T})=(1+\eta_T)(1-\eta_T)=1-\eta_T^2,
 \qquad c_2(\mathscr E_{\eta,T})=-\eta_T^2.
 \label{cusp-native-full-flux-Chern-class}
\end{equation}
Thus the orientation-sensitive Chern numbers are
\begin{equation}
 \langle c_2(\mathscr E_{\eta,T}),[F_T]_{\mathrm{mark}}\rangle=-12,
 \qquad
 \langle c_2(\mathscr E_{\eta,T}),
 [F_T]_{x\text{-}\mathrm{complex}}\rangle=+12.
 \label{cusp-native-full-flux-Chern-numbers}
\end{equation}
The cohomology class \(c_2=-\eta_T^2\) is nonzero, so this split
\(SU(2)\)-bundle is not smoothly trivial.  This is another exact distinction
from the unit-flux split bundle of \cref{cusp-magnetic-su2-triviality}.

With
\begin{equation}
 H_{\mathrm{mag}}=\begin{pmatrix}i&0\\0&-i\end{pmatrix},
 \qquad
 \langle X,Y\rangle_{\mathfrak{su}(2)}
 =-\frac12\operatorname{Tr}(XY),
 \label{cusp-native-full-flux-Lie-normalization}
\end{equation}
one has \(\lvert H_{\mathrm{mag}}\rvert=1\) and
\begin{equation}
 F_{A_{\eta,T}}=-2\pi\eta_T H_{\mathrm{mag}}.
 \label{cusp-native-full-flux-SU2-curvature}
\end{equation}
This curvature is parallel; hence
\(d_{A_{\eta,T}}^*F_{A_{\eta,T}}=0\), and \(A_{\eta,T}\) is Yang--Mills.
For
\begin{equation}
 \operatorname{YM}_T(A)=\frac1{2g_{\mathrm{YM}}^2}
 \int_{F_T}|F_A|^2\,\mathrm d\operatorname{vol}_{g_T},
 \label{cusp-native-full-flux-action-normalization}
\end{equation}
diagonalizing \cref{cusp-native-full-flux-K} gives the exact action
\begin{align}
 \operatorname{YM}_T(A_{\eta,T})
 &=\frac{D_T}{2g_{\mathrm{YM}}^2}
   (\mathcal B_{T,+}^2+\mathcal B_{T,-}^2)
 \label{cusp-native-full-flux-action-frequency}\\
 &=\frac{2\pi^2}{g_{\mathrm{YM}}^2D_T}
   (q_T^2+36L_T^2+72m_T^2)
 \label{cusp-native-full-flux-action-coordinate}\\
 &=\frac{24\pi^2}{g_{\mathrm{YM}}^2}
 +\frac{2\pi^2S_T^2}{g_{\mathrm{YM}}^2D_T}.
 \label{cusp-native-full-flux-action-excess}
\end{align}
Indeed
\(q_T^2+36L_T^2+72m_T^2=12D_T+(6L_T-q_T)^2\).
Consequently
\begin{equation}
 \operatorname{YM}_T(A_{\eta,T})
 =\frac{74\pi^2}{g_{\mathrm{YM}}^2}
 +\frac{70\pi^2B}{g_{\mathrm{YM}}^2T}+O(T^{-2}).
 \label{cusp-native-full-flux-action-asymptotic}
\end{equation}

\subsection{The closed Jacobi operator and exact Coulomb reduction}

We next define the operator before computing its spectrum.  Let
\begin{equation}
 \mathscr H_T^1=L^2\Omega^1(F_T;\operatorname{ad}\mathscr E_{\eta,T}),
 \qquad
 \mathscr H_T^0=L^2\Omega^0(F_T;\operatorname{ad}\mathscr E_{\eta,T})
 \label{cusp-native-full-flux-real-Hilbert-spaces}
\end{equation}
be real Hilbert spaces.  Their complexifications carry the root
decomposition used below, and every operator constructed here commutes with
the anti-linear real structure.

Define the maximal distributional codifferential by
\begin{align}
 \operatorname{Dom}D_T^{\max}
 =\{a\in\mathscr H_T^1:\;&\text{there is }u\in\mathscr H_T^0\text{ such that}
 \notag\\[-1mm]
 &\langle a,d_{A_{\eta,T}}\phi\rangle_{L^2}
 =\langle u,\phi\rangle_{L^2}
 \text{ for every smooth }\phi\},
 \label{cusp-native-full-flux-max-codiff-domain}\\
 D_T^{\max}a={}&u,
 \qquad \mathscr S_T:=\ker D_T^{\max}.
 \label{cusp-native-full-flux-Coulomb-slice}
\end{align}
Uniqueness follows from density of the test space.  Passing to the limit in
the defining identity proves that \(D_T^{\max}\) is closed, hence the Coulomb
slice \(\mathscr S_T\) is a closed subspace.

On adjoint-valued scalars, let \(\Delta_{A_{\eta,T}}^{(0)}\) be the
nonnegative self-adjoint operator of the closed form
\begin{equation}
 \mathfrak q_T^{(0)}[\phi,\psi]
 =\langle d_{A_{\eta,T}}\phi,d_{A_{\eta,T}}\psi\rangle_{L^2},
 \qquad
 \operatorname{Dom}\mathfrak q_T^{(0)}=H^1\Omega^0.
 \label{cusp-native-full-flux-scalar-form}
\end{equation}
Elliptic regularity gives
\begin{equation}
 \Delta_{A_{\eta,T}}^{(0)}=d_{A_{\eta,T}}^*d_{A_{\eta,T}},
 \qquad
 \operatorname{Dom}\Delta_{A_{\eta,T}}^{(0)}=H^2\Omega^0.
 \label{cusp-native-full-flux-scalar-domain}
\end{equation}

In a constant orthonormal frame, define the symmetric zeroth-order
endomorphism
\begin{equation}
 (V_Ta)_\nu=-2\sum_{\mu=1}^4
 [F_{A_{\eta,T},\nu\mu},a_\mu].
 \label{cusp-native-full-flux-spin-endomorphism}
\end{equation}
Symmetry follows from invariance of the compact Lie-algebra inner product,
antisymmetry of \(F_{\nu\mu}\), and an exchange of \(\nu,\mu\).  Consider
the gauge-fixed form
\begin{equation}
 \mathfrak j_T^{\mathrm{gf}}[a,b]
 =\langle\nabla_{A_{\eta,T}}a,\nabla_{A_{\eta,T}}b\rangle_{L^2}
  +\langle V_Ta,b\rangle_{L^2},
 \qquad
 \operatorname{Dom}\mathfrak j_T^{\mathrm{gf}}=H^1\Omega^1.
 \label{cusp-native-full-flux-Jacobi-form}
\end{equation}

\begin{proposition}[Closed gauge-fixed realization]
\label{cusp-native-full-flux-closed-Jacobi}
The form \cref{cusp-native-full-flux-Jacobi-form} is densely defined, closed,
symmetric, and bounded below.  Its self-adjoint operator
\(\mathcal J_T^{\mathrm{gf}}\) has compact resolvent, domain \(H^2\Omega^1\),
and expression
\begin{equation}
 (\mathcal J_T^{\mathrm{gf}}a)_\nu
 =-\sum_{\mu=1}^4
   \nabla_{A_{\eta,T},e_\mu}\nabla_{A_{\eta,T},e_\mu}a_\nu
  -2\sum_{\mu=1}^4[F_{A_{\eta,T},\nu\mu},a_\mu].
 \label{cusp-native-full-flux-Jacobi-expression}
\end{equation}
It is exactly the Hessian operator with the background-gauge term, before
the overall factor \(g_{\mathrm{YM}}^{-2}\).
\end{proposition}

\begin{proof}
The first summand in \cref{cusp-native-full-flux-Jacobi-form}, together with
the \(L^2\)-norm, is an equivalent \(H^1\)-norm.  The smooth endomorphism
\(V_T\) is bounded and symmetric.  The bounded perturbation theorem for
closed forms gives closedness and a lower bound, and the representation
theorem gives a self-adjoint operator.  Its principal symbol is
\(|\xi|^2\) times the identity.  Elliptic regularity identifies its maximal
\(L^2\)-domain with \(H^2\Omega^1\); direct substitution supplies the reverse
inclusion.  The compact embedding \(H^1\hookrightarrow L^2\) gives compact
resolvent.

For the final assertion, define
\([a\wedge b]_{\mu\nu}=[a_\mu,b_\nu]-[a_\nu,b_\mu]\).  The exact
two-parameter curvature expansion is
\[
 F_{A_{\eta,T}+sa+tb}=F_{A_{\eta,T}}+s d_Aa+t d_Ab
 +\frac{s^2}{2}[a\wedge a]+st[a\wedge b]
 +\frac{t^2}{2}[b\wedge b].
\]
It yields
\begin{equation}
 \operatorname{YM}_T''\big|_{A_{\eta,T}}(a,b)
 =\frac1{g_{\mathrm{YM}}^2}
 \left(\langle d_Aa,d_Ab\rangle
 +\langle F_A,[a\wedge b]\rangle\right).
 \label{cusp-native-full-flux-raw-Hessian}
\end{equation}
On the flat torus, commuting the covariant derivatives with the retained
curvature convention gives
\begin{align*}
 \langle d_Aa,d_Ab\rangle+\langle d_A^*a,d_A^*b\rangle
 &=\langle\nabla_Aa,\nabla_Ab\rangle
 -\sum_{\nu,\mu}\langle a_\nu,[F_{\nu\mu},b_\mu]\rangle,\\
 \langle F_A,[a\wedge b]\rangle
 &=-\sum_{\nu,\mu}\langle a_\nu,[F_{\nu\mu},b_\mu]\rangle.
\end{align*}
The factor \(1/2\) in the two-form inner product is included in the second
identity.  Adding
\(g_{\mathrm{YM}}^{-2}\langle d_A^*a,d_A^*b\rangle\) to
\cref{cusp-native-full-flux-raw-Hessian} gives exactly
\(g_{\mathrm{YM}}^{-2}\mathfrak j_T^{\mathrm{gf}}[a,b]\), including the
coefficient \(-2\) in \cref{cusp-native-full-flux-spin-endomorphism}.
\end{proof}

\begin{proposition}[Resolvent intertwining and the reducing Coulomb slice]
\label{cusp-native-full-flux-Coulomb-reduction}
On smooth sections,
\begin{equation}
 \mathcal J_T^{\mathrm{gf}}d_{A_{\eta,T}}
 =d_{A_{\eta,T}}\Delta_{A_{\eta,T}}^{(0)},
 \qquad
 d_{A_{\eta,T}}^*\mathcal J_T^{\mathrm{gf}}
 =\Delta_{A_{\eta,T}}^{(0)}d_{A_{\eta,T}}^*.
 \label{cusp-native-full-flux-formal-intertwining}
\end{equation}
For \(z\in\mathbb C\setminus\mathbb R\) and
\(f\in\operatorname{Dom}D_T^{\max}\),
\begin{equation}
 D_T^{\max}(\mathcal J_T^{\mathrm{gf}}-z)^{-1}f
 =(\Delta_{A_{\eta,T}}^{(0)}-z)^{-1}D_T^{\max}f.
 \label{cusp-native-full-flux-resolvent-intertwining}
\end{equation}
Thus \(\mathscr S_T\) reduces \(\mathcal J_T^{\mathrm{gf}}\), and its part
\begin{equation}
 \mathcal J_T=\mathcal J_T^{\mathrm{gf}}\big|_{\mathscr S_T},
 \qquad
 \operatorname{Dom}\mathcal J_T=H^2\Omega^1\cap\mathscr S_T
 \label{cusp-native-full-flux-Coulomb-operator}
\end{equation}
is self-adjoint.  Its closed form is the restriction of
\(\mathfrak j_T^{\mathrm{gf}}\) to \(H^1\Omega^1\cap\mathscr S_T\), and the
Yang--Mills Hessian on this slice is
\begin{equation}
 \operatorname{YM}_T''\big|_{A_{\eta,T}}
 =g_{\mathrm{YM}}^{-2}\mathfrak j_T.
 \label{cusp-native-full-flux-Coulomb-Hessian}
\end{equation}
\end{proposition}

\begin{proof}
For a smooth adjoint-valued scalar \(\phi\), write
\(\nabla_\mu=\nabla_{A_{\eta,T},e_\mu}\).  Retaining both curvature terms
when commuting three derivatives gives
\begin{align*}
 &(\mathcal J_T^{\mathrm{gf}}d_A\phi
       -d_A\Delta_A^{(0)}\phi)_\nu\\
 &\quad=-\nabla_\mu\nabla_\mu\nabla_\nu\phi
       +\nabla_\nu\nabla_\mu\nabla_\mu\phi
       -2[F_{\nu\mu},\nabla_\mu\phi]\\
 &\quad=-[\nabla_\mu F_{\mu\nu},\phi]
       =[(d_A^*F_A)_\nu,\phi]=0.
\end{align*}
This proves the first identity in
\cref{cusp-native-full-flux-formal-intertwining}; formal adjunction gives the
second.

Let \(u=(\mathcal J_T^{\mathrm{gf}}-z)^{-1}f\).  The operator domain gives
\(u\in H^2\), hence \(D_T^{\max}u\in H^1\).  Distributionally applying the
second intertwining identity to
\((\mathcal J_T^{\mathrm{gf}}-z)u=f\) yields
\[
 (\Delta_A^{(0)}-z)D_T^{\max}u=D_T^{\max}f.
\]
The right side is in \(L^2\); scalar elliptic regularity upgrades
\(D_T^{\max}u\) to \(H^2\), and the scalar resolvent proves
\cref{cusp-native-full-flux-resolvent-intertwining}.  Both resolvents of the
self-adjoint gauge-fixed operator preserve \(\mathscr S_T\); adjunction shows
that they also preserve its orthogonal complement.  Therefore the slice is
reducing and its operator part is self-adjoint with the displayed domain.
The reducing projection commutes with the spectral calculus, so the form
domain is exactly \(H^1\Omega^1\cap\mathscr S_T\).  The Hessian identity is
the restriction of the last assertion of
\cref{cusp-native-full-flux-closed-Jacobi}.
\end{proof}

\subsection{Charge two, all polarizations, and the Morse index}

Choose root vectors with
\([H_{\mathrm{mag}},E_\pm]=\pm2iE_\pm\).  Then
\begin{equation}
 (\operatorname{ad}\mathscr E_{\eta,T})_{\mathbb C}
 =\underline{\mathbb C}H_{\mathrm{mag}}
 \oplus\mathscr L_{\eta,T}^{\otimes2}E_+
 \oplus\mathscr L_{\eta,T}^{\otimes(-2)}E_-.
 \label{cusp-native-full-flux-root-decomposition}
\end{equation}
The \(E_+\) root connection has curvature \(-4\pi i\eta_T\).  Replacing
\(E_\eta\) by \(2E_\eta\) gives \(J_{2,T}=J_{\eta,T}\), doubles both
frequencies, and changes the vacuum dimension from six to twenty-four by
\cref{cusp-native-full-flux-vacuum-dimensions}.  The \(E_-\) root has
\(q=-2\), uses the conjugate auxiliary structure
\(J_{-2,T}=-J_{\eta,T}\), and has the same positive frequencies and the same
dimension \(6(-2)^2=24\).  Repeating the complete ladder proof of
\cref{cusp-native-full-flux-landau-theorem} therefore gives, for either root,
the joint charge-two scalar levels
\begin{equation}
 E^{(2)}_{T;p,q}
 =2(2p+1)\mathcal B_{T,+}+2(2q+1)\mathcal B_{T,-},
 \qquad p,q\in\mathbb Z_{\ge0},
 \label{cusp-native-full-flux-charge-two-levels}
\end{equation}
each of complex dimension
\begin{equation}
 |\operatorname{Pf}(2E_\eta)|=24.
 \label{cusp-native-full-flux-charge-two-multiplicity}
\end{equation}
Coincident joint pairs must again be summed.

Choose \(O_T\in SO(2)\) with
\(O_T^{\mathsf T}K_TO_T=\operatorname{diag}(\kappa_{T,+},\kappa_{T,-})\),
and define the constant orthonormal coframe
\begin{equation}
 \binom{\mathrm dX_+}{\mathrm dX_-}
 =O_T^{\mathsf T}\binom{\mathrm dx^1}{\mathrm dx^3},
 \qquad
 \binom{\mathrm dY_+}{\mathrm dY_-}
 =O_T^{\mathsf T}\binom{\mathrm dx^2}{\mathrm dx^4}.
 \label{cusp-native-full-flux-orthogonal-coframe}
\end{equation}
Then
\begin{equation}
 \eta_T=\frac{\mathcal B_{T,+}}{2\pi}
          \mathrm dX_+\wedge\mathrm dY_+
         -\frac{\mathcal B_{T,-}}{2\pi}
          \mathrm dX_-\wedge\mathrm dY_-.
 \label{cusp-native-full-flux-normal-form}
\end{equation}
The odd grouping permutation in
\cref{cusp-native-full-flux-orthogonal-coframe} reverses orientation; it is
used only to expose the two blocks and never to infer a Chern sign.

On the \(E_+\) root line, the four constant circular one-forms diagonalize
the spin endomorphism.  Direct substitution into
\cref{cusp-native-full-flux-spin-endomorphism}, with the charge-two root
normalization retained, gives
\begin{center}
\begin{tabular}{>{$}c<{$} >{$}c<{$} >{$}c<{$}}
\toprule
\text{polarization}&V_T\text{-shift}
 &\mathcal J_T^{\mathrm{gf}}\text{-eigenvalue on }(p,q)\\
\midrule
\mathrm dX_+ + i\mathrm dY_+&-4\mathcal B_{T,+}
 &2[(2p-1)\mathcal B_{T,+}+(2q+1)\mathcal B_{T,-}]\\
\mathrm dX_+ - i\mathrm dY_+&+4\mathcal B_{T,+}
 &2[(2p+3)\mathcal B_{T,+}+(2q+1)\mathcal B_{T,-}]\\
\mathrm dX_- - i\mathrm dY_-&-4\mathcal B_{T,-}
 &2[(2p+1)\mathcal B_{T,+}+(2q-1)\mathcal B_{T,-}]\\
\mathrm dX_- + i\mathrm dY_-&+4\mathcal B_{T,-}
 &2[(2p+1)\mathcal B_{T,+}+(2q+3)\mathcal B_{T,-}]\\
\bottomrule
\end{tabular}
\end{center}
The \(E_-\) table is its anti-linear conjugate.  The Cartan sector is the
ordinary nonnegative coefficientwise Laplacian on one-forms.

\begin{theorem}[Exact negative spectrum, walls, and real Morse index]
\label{cusp-native-full-flux-Morse-theorem}
Put
\begin{equation}
 \mathcal B_{T,\max}=\max\{\mathcal B_{T,+},\mathcal B_{T,-}\},
 \quad
 \mathcal B_{T,\min}=\min\{\mathcal B_{T,+},\mathcal B_{T,-}\},
 \quad
 \rho_T^*=\frac{\mathcal B_{T,\max}}{\mathcal B_{T,\min}}
 =\frac{\Xi_T+|S_T|}{\Xi_T-|S_T|}.
 \label{cusp-native-full-flux-ordered-ratio}
\end{equation}
For every regular \(T\), the negative spectrum of the self-adjoint Coulomb
operator is exactly
\begin{equation}
 \left\{
 \nu_{T,n}=2[(2n+1)\mathcal B_{T,\min}-\mathcal B_{T,\max}]
 =\frac{4\pi}{D_T}[n\Xi_T-(n+1)|S_T|]:
 2n+1<\rho_T^*
 \right\}.
 \label{cusp-native-full-flux-general-negative-spectrum}
\end{equation}
Every indicated band has real multiplicity \(48\), and
\begin{equation}
 \operatorname{ind}_{\mathbb R}(\mathcal J_T)
 =48\left\lceil\frac{\rho_T^*-1}{2}\right\rceil.
 \label{cusp-native-full-flux-general-index}
\end{equation}
At the \(n\)-th wall the strict inequality becomes equality.  Equivalently,
\begin{equation}
 \rho_T^*=2n+1
 \iff n\Xi_T=(n+1)|S_T|
 \iff 24n^2D_T=(2n+1)S_T^2.
 \label{cusp-native-full-flux-wall-equations}
\end{equation}
The corresponding branch is then a zero branch and is not counted in the
Morse index.  If one removes the absolute value by choosing the
\(\mathcal B_{T,+}>\mathcal B_{T,-}\) side, the polynomial equation in
\cref{cusp-native-full-flux-wall-equations} must be accompanied by \(S_T>0\);
otherwise squaring introduces the opposite side as an extraneous solution.

On the sufficiently deep tail of
\cref{cusp-native-full-flux-deep-tail-window}, the negative branches are
exactly \(n=0,1,2\),
\begin{equation}
 \nu_{T,n}=\frac{4\pi}{D_T}[n\Xi_T-(n+1)S_T],
 \qquad n=0,1,2,
 \label{cusp-native-full-flux-deep-negative-branches}
\end{equation}
each of real multiplicity \(48\).  Therefore
\begin{equation}
 \operatorname{ind}_{\mathbb R}(\mathcal J_T)=144,
 \label{cusp-native-full-flux-deep-index}
\end{equation}
and
\begin{align}
 \nu_{T,0}&=-\frac{20\pi}{T}+O(T^{-2}),&
 \nu_{T,1}&=-\frac{12\pi}{T}+O(T^{-2}),\notag\\
 \nu_{T,2}&=-\frac{4\pi}{T}+O(T^{-2}),&
 \nu_{T,3}&=+\frac{4\pi}{T}+O(T^{-2}).
 \label{cusp-native-full-flux-negative-asymptotics}
\end{align}
The Hessian operator for
\cref{cusp-native-full-flux-action-normalization} has eigenvalues
\(\nu_{T,n}/g_{\mathrm{YM}}^2\) and the same index.
\end{theorem}

\begin{proof}
Assume first \(\mathcal B_{T,+}>\mathcal B_{T,-}\).  In the first row of
the polarization table, \(p\ge1\) makes the eigenvalue positive.  At \(p=0\)
it is
\[
 2[-\mathcal B_{T,+}+(2q+1)\mathcal B_{T,-}],
\]
which is negative exactly when \(2q+1<\rho_T^*\).  The second and fourth
rows are positive, and the third has minimum
\(2(\mathcal B_{T,+}-\mathcal B_{T,-})>0\).  When
\(\mathcal B_{T,-}>\mathcal B_{T,+}\), the identical argument with the two
planes exchanged shows that only the \(q=0\) part of the third row can be
negative and gives the same formula with the ordered frequencies.  If the
frequencies agree, both candidate lowest branches are zero and none is
negative.  This proves the exhaustive list for the gauge-fixed operator.

Every negative gauge-fixed eigenform is automatically Coulomb.  Elliptic
regularity makes it smooth; if
\(\mathcal J_T^{\mathrm{gf}}a=\lambda a\) with \(\lambda<0\), then
\cref{cusp-native-full-flux-formal-intertwining} gives
\[
 \Delta_{A_{\eta,T}}^{(0)}(d_{A_{\eta,T}}^*a)
 =\lambda d_{A_{\eta,T}}^*a.
\]
Nonnegativity of the scalar operator forces \(d_A^*a=0\).  Hence the full and
Coulomb negative spectral spaces coincide.

One joint charge-two scalar space contributes complex dimension \(24\) in
the \(E_+\) sector.  Its conjugate contributes complex dimension \(24\) in
the complexification of the real adjoint problem.  A complexified real
eigenspace has complex dimension equal to the original real dimension, so
the real multiplicity is \(24+24=48\), not \(24\) or \(96\).

The count of integers \(n\ge0\) satisfying \(2n+1<\rho_T^*\) is
\(\lceil(\rho_T^*-1)/2\rceil\), including the correct exclusion at an odd
integer wall.  Substitution of the ordered frequency formulas gives
\cref{cusp-native-full-flux-general-negative-spectrum}.  The wall equation
is
\(n\Xi_T=(n+1)|S_T|\); squaring and using
\(\Xi_T^2-S_T^2=24D_T\) gives the final polynomial form, and nonnegativity of
both sides gives the converse.  Finally,
\cref{cusp-native-full-flux-deep-tail-window} leaves precisely
\(n=0,1,2\), and the leading frequency terms give
\(\nu_{T,n}=(8n-20)\pi/T+O(T^{-2})\).
\end{proof}

The lowest branch and the action excess are tied by the exact identities
\begin{equation}
 \nu_{T,0}=-2|\mathcal B_{T,+}-\mathcal B_{T,-}|
 =-\frac{4\pi|S_T|}{D_T},
 \qquad
 \operatorname{YM}_T(A_{\eta,T})-\frac{24\pi^2}{g_{\mathrm{YM}}^2}
 =\frac{D_T}{8g_{\mathrm{YM}}^2}\nu_{T,0}^2.
 \label{cusp-native-full-flux-action-instability}
\end{equation}

% Canonical integration of agent_logs/native_direct_integral_complete_canonical_20260903.tex.
% The fixed and varying fields remain separately typed; T is a multiplication
% parameter and never a hidden radial derivative.

\subsection{The fixed and varying radial direct integrals}
\label{cusp-native-full-flux-radial-section}

Choose \(T_1>T_0\) so that the complete deep-tail inequalities in
\cref{cusp-native-full-flux-deep-tail-window} hold on
\[
 I:=(T_1,\infty).
\]
Let \(Y=\mathbb R^4/\mathbb Z^4\).  Denote by
\(\mathscr L_\eta\to Y\) the fixed marked Hermitian line bundle defined
by the factor of automorphy in
\cref{cusp-native-full-flux-cocycle}.  The marking diffeomorphism and its
transported bundle isomorphism are
\[
 \Phi_T:Y\longrightarrow F_T,\qquad [y]\longmapsto[P_Ty],
 \qquad
 \widehat\Phi_T:\mathscr L_\eta\longrightarrow\mathscr L_{\eta,T}.
\]
They are fixed by requiring
\[
 \widehat\Phi_T^*\nabla_{\eta,T}
 =\mathrm d-2\pi i(6y^1\,\mathrm dy^4+y^2\,\mathrm dy^3),
\]
so no gauge, metric, lattice, or orientation normalization is hidden in the
identification.  Put
\[
 \mathscr K_T
 :=L^2(F_T,\mathscr L_{\eta,T};
       \mathrm d\operatorname{vol}_{g_T}).
\]

The pullback of the physical volume is
\begin{equation}
 \Phi_T^*(\mathrm d\operatorname{vol}_{g_T})
 =\sqrt{\det(P_T^{\mathsf T}P_T)}\,\mathrm d^4y
 =D_T\,\mathrm d^4y.
\label{native-di-volume-pullback}
\end{equation}
Put
\[
 \mathscr H_*:=L^2(Y,\mathscr L_\eta;\mathrm d^4y),
 \qquad
 \mathscr D_*:=H^2(Y;\mathscr L_\eta),
\]
where the Sobolev space is defined using any one fixed smooth reference
metric and connection on the fixed compact bundle.  Define
\begin{equation}
 U_T:\mathscr K_T\longrightarrow\mathscr H_*,
 \qquad
 U_Ts:=D_T^{1/2}\widehat\Phi_T^*s.
\label{native-di-marking-unitary}
\end{equation}
Then
\[
 \|U_Ts\|_{\mathscr H_*}^2
 =D_T\int_Y|\widehat\Phi_T^*s|^2\,\mathrm d^4y
 =\|s\|_{\mathscr K_T}^2,
\]
and
\[
 U_T^{-1}v
 =D_T^{-1/2}(\widehat\Phi_T^{-1})^*v.
\]
Thus \(U_T\) is an onto unitary, not merely an isometric injection.

Transport the fibre operator to the fixed space:
\begin{equation}
 A_T:=U_T\Delta_{\eta,T}U_T^{-1},
 \qquad
 \operatorname{Dom}(A_T)=\mathscr D_*.
\label{native-di-fixed-fibre}
\end{equation}
Its closed quadratic form is
\begin{equation}
 \widetilde{\mathfrak a}_T[v,w]
 =\int_Y(G_T^{-1})^{jk}
  \langle\nabla_{\eta,j}v,\nabla_{\eta,k}w\rangle\,\mathrm d^4y,
 \qquad G_T=P_T^{\mathsf T}P_T.
\label{native-di-fixed-form}
\end{equation}
Here \(\nabla_{\eta,j}\) is the \(\mathrm dy^j\)-component of the fixed
marked connection displayed above.  The factor \(D_T^{1/2}\) in \(U_T\) is
constant in \(y\), and it cancels
the factor \(D_T\) in \eqref{native-di-volume-pullback}; hence no unrecorded
factor appears in \eqref{native-di-fixed-form}.

\begin{lemma}[Common domain and measurable resolvents]
\label{native-di-common-domain-lemma}
For every compact interval \(K\Subset I\), the family
\[
 T\longmapsto A_T\in\mathcal B(\mathscr D_*,\mathscr H_*)
\]
is norm-continuous, in fact real analytic.  For every
\(z\in\mathbb C\setminus\mathbb R\), the resolvents
\((A_T-z)^{-1}\) are locally norm-continuous as operators on
\(\mathscr H_*\), and hence form a measurable operator field.
\end{lemma}

\begin{proof}
On \(K\), the matrices \(G_T^{-1}\) are uniformly positive definite, and
their entries are real analytic in \(T\).  In fixed local bundle
trivializations, \eqref{native-di-fixed-form} therefore gives a uniformly
elliptic second-order expression whose coefficients depend real
analytically on \(T\).  Elliptic regularity identifies every operator
domain with the same set \(\mathscr D_*\), and the standard local
coordinate estimate for a differential operator of order two gives
norm-continuity in
\(\mathcal B(\mathscr D_*,\mathscr H_*)\).

Uniform ellipticity on \(K\) gives
\[
 \|v\|_{\mathscr D_*}
 \le C_K\bigl(\|A_Tv\|_{\mathscr H_*}
                  +\|v\|_{\mathscr H_*}\bigr),
 \qquad T\in K.
\]
For \(v=(A_T-z)^{-1}f\), self-adjointness gives
\(\|v\|\le|\operatorname{Im}z|^{-1}\|f\|\), while
\(A_Tv=f+zv\).  Consequently
\begin{equation}
 \sup_{T\in K}
 \|(A_T-z)^{-1}\|_{\mathscr H_*\to\mathscr D_*}<\infty.
\label{native-di-elliptic-resolvent-bound}
\end{equation}
The resolvent identity, with its exact sign,
\[
 (A_T-z)^{-1}-(A_R-z)^{-1}
 =-(A_T-z)^{-1}(A_T-A_R)(A_R-z)^{-1},
\]
together with \eqref{native-di-elliptic-resolvent-bound}, proves local
norm-resolvent continuity.  A countable compact exhaustion of \(I\) gives
measurability on all of \(I\).
\end{proof}

We now define the varying measurable Hilbert field rather than presuppose
one.  A section \(T\mapsto s_T\in\mathscr K_T\) is declared measurable
precisely when
\[
 T\longmapsto U_Ts_T\in\mathscr H_*
\]
is strongly measurable.  This is the measurable field pulled back from the
constant separable field \(I\times\mathscr H_*\).  Its direct integral is
\[
 \mathscr H_{\mathrm{var}}
 :=\int_I^\oplus\mathscr K_T\,\mathrm d T.
\]
The fixed radial Hilbert space is
\[
 \mathscr H_{\mathrm{rad}}
 :=L^2(I,\mathrm d T;\mathscr H_*).
\]
By construction,
\begin{equation}
 \mathscr U:=\int_I^\oplus U_T\,\mathrm d T:
 \mathscr H_{\mathrm{var}}\longrightarrow\mathscr H_{\mathrm{rad}},
 \qquad
 (\mathscr Us)(T)=U_Ts_T,
\label{native-di-direct-integral-unitary}
\end{equation}
is an isometry.  It is onto: if
\(\Psi\in\mathscr H_{\mathrm{rad}}\), then
\(s_T=U_T^{-1}\Psi(T)\) is measurable by the very definition of the
varying field and has the same squared norm integral.

On the fixed direct integral, define the maximal multiplication operator
\begin{equation}
 (\widetilde{\mathscr M}_\eta\Psi)(T):=A_T\Psi(T)
\label{native-di-fixed-multiplication}
\end{equation}
with the exact domain
\begin{equation}
\begin{split}
 \operatorname{Dom}(\widetilde{\mathscr M}_\eta)
 :=\{\Psi\in\mathscr H_{\mathrm{rad}}:\;&
 \Psi(T)\in\mathscr D_*\text{ for almost every }T,\\
 &T\mapsto A_T\Psi(T)\text{ is strongly measurable},\\
 &\int_I\|A_T\Psi(T)\|_{\mathscr H_*}^2\,\mathrm d T<\infty\}.
\end{split}
\label{native-di-fixed-maximal-domain}
\end{equation}
On the varying direct integral, define
\begin{equation}
\begin{split}
 \operatorname{Dom}(\mathscr M_\eta)
 :=\{s\in\mathscr H_{\mathrm{var}}:\;&
 s_T\in\operatorname{Dom}(\Delta_{\eta,T})\text{ for almost every }T,\\
 &T\mapsto U_T\Delta_{\eta,T}s_T
       \text{ is strongly measurable},\\
 &\int_I\|\Delta_{\eta,T}s_T\|_{\mathscr K_T}^2\,\mathrm d T<\infty\},
\end{split}
\label{native-di-varying-maximal-domain}
\end{equation}
and
\begin{equation}
 (\mathscr M_\eta s)_T:=\Delta_{\eta,T}s_T
 \quad\text{for almost every }T.
\label{native-di-varying-multiplication}
\end{equation}
Every expression in \eqref{native-di-fixed-multiplication} and
\eqref{native-di-varying-multiplication} is now typed on its actual fibre.
Equations \eqref{native-di-fixed-fibre},
\eqref{native-di-direct-integral-unitary},
\eqref{native-di-fixed-maximal-domain}, and
\eqref{native-di-varying-maximal-domain} give both the domain equality and
the exact intertwining identity
\begin{equation}
 \mathscr U\operatorname{Dom}(\mathscr M_\eta)
 =\operatorname{Dom}(\widetilde{\mathscr M}_\eta),
 \qquad
 \mathscr U\mathscr M_\eta\mathscr U^{-1}
 =\widetilde{\mathscr M}_\eta.
\label{native-di-intertwining}
\end{equation}

\subsubsection{Density and self-adjointness on the maximal domain}

\begin{proposition}[Maximal-domain realization]
\label{native-di-self-adjoint-proposition}
The operators \(\widetilde{\mathscr M}_\eta\) and
\(\mathscr M_\eta\) are densely defined, nonnegative, and self-adjoint on
the domains \eqref{native-di-fixed-maximal-domain} and
\eqref{native-di-varying-maximal-domain}, respectively.
\end{proposition}

\begin{proof}
For \(\chi\in C_c(I)\) and \(v\in\mathscr D_*\), local continuity of
\(T\mapsto A_Tv\) makes it bounded on \(\operatorname{supp}\chi\).
Thus \(\chi v\in\operatorname{Dom}(\widetilde{\mathscr M}_\eta)\).  Since
\(\mathscr D_*\) is dense in \(\mathscr H_*\), the algebraic tensor product
\[
 C_c(I)\mathbin{\odot}\mathscr D_*
\]
is dense in \(L^2(I;\mathscr H_*)\).  Hence the fixed operator is densely
defined.  Fibrewise self-adjointness and Cauchy--Schwarz give, for
\(\Psi,\Theta\) in its domain,
\[
 \langle\widetilde{\mathscr M}_\eta\Psi,\Theta\rangle
 =\int_I\langle A_T\Psi(T),\Theta(T)\rangle\,\mathrm d T
 =\langle\Psi,\widetilde{\mathscr M}_\eta\Theta\rangle,
\]
so it is symmetric.

Fix \(z\in\mathbb C\setminus\mathbb R\).  For
\(F\in\mathscr H_{\mathrm{rad}}\), define pointwise
\begin{equation}
 (\mathcal R_zF)(T):=(A_T-z)^{-1}F(T).
\label{native-di-radial-resolvent}
\end{equation}
The locally norm-continuous resolvent field, approximation of \(F\) by
simple functions, and
\[
 \|(A_T-z)^{-1}\|_{\mathscr H_*\to\mathscr H_*}
 \le|\operatorname{Im}z|^{-1}
\]
show that \(\mathcal R_zF\) is strongly measurable and
\begin{equation}
 \|\mathcal R_zF\|_{\mathscr H_{\mathrm{rad}}}
 \le|\operatorname{Im}z|^{-1}
    \|F\|_{\mathscr H_{\mathrm{rad}}}.
\label{native-di-radial-resolvent-L2}
\end{equation}
Pointwise, \((A_T-z)^{-1}F(T)\in\mathscr D_*\).  The exact identity
\begin{equation}
 A_T(A_T-z)^{-1}=I+z(A_T-z)^{-1}
\label{native-di-maximal-domain-identity}
\end{equation}
gives the global uniform bound
\[
 \|A_T(A_T-z)^{-1}\|
 \le1+\frac{|z|}{|\operatorname{Im}z|}.
\]
The right-hand side of \eqref{native-di-maximal-domain-identity}, applied
to \(F(T)\), is strongly measurable and square-integrable.  Therefore
\(\mathcal R_zF\) belongs to the full maximal domain
\eqref{native-di-fixed-maximal-domain}; this is the domain step that
pointwise inversion alone does not prove.

Fibrewise identities now yield
\[
 (\widetilde{\mathscr M}_\eta-z)\mathcal R_z=I
 \quad\text{on }\mathscr H_{\mathrm{rad}},
 \qquad
 \mathcal R_z(\widetilde{\mathscr M}_\eta-z)=I
 \quad\text{on }\operatorname{Dom}(\widetilde{\mathscr M}_\eta).
\]
In particular, the ranges of
\(\widetilde{\mathscr M}_\eta-i\) and
\(\widetilde{\mathscr M}_\eta+i\) are both all of
\(\mathscr H_{\mathrm{rad}}\).  The range criterion for a densely defined
symmetric operator proves self-adjointness.  Moreover,
\[
 \langle\widetilde{\mathscr M}_\eta\Psi,\Psi\rangle
 =\int_I\langle A_T\Psi(T),\Psi(T)\rangle\,\mathrm d T\ge0,
\]
so it is nonnegative.  The unitary equivalence
\eqref{native-di-intertwining} transfers density, self-adjointness, and
nonnegativity to the varying-fibre operator.
\end{proof}

\subsubsection{Measurable eigenvectors, including at band crossings}

\begin{lemma}[Explicit measurable eigenvector selection]
\label{native-di-measurable-eigenvector-lemma}
For every fixed
\(\mathbf n=(n_+,n_-)\in\mathbb Z_{\ge0}^2\), there is a strongly
measurable map
\[
 I\longrightarrow\mathscr D_*,
 \qquad T\longmapsto\phi_{\mathbf n}(T),
\]
such that
\begin{equation}
 \|\phi_{\mathbf n}(T)\|_{\mathscr H_*}=1,
 \qquad
 A_T\phi_{\mathbf n}(T)
 =E_{\mathbf n}(T)\phi_{\mathbf n}(T)
\label{native-di-measurable-eigenvector}
\end{equation}
for every \(T\in I\).  The construction remains valid when
\(E_{\mathbf n}(T)\) crosses one or more other bands.
\end{lemma}

\begin{proof}
Set
\[
 C_{\mathbf n,T}:=A_T-E_{\mathbf n}(T),
 \qquad \operatorname{Dom}(C_{\mathbf n,T})=\mathscr D_*.
\]
The band \(E_{\mathbf n}\) is continuous.  The proof of
Lemma~\ref{native-di-common-domain-lemma}, applied also to the scalar
shift, shows that
\((C_{\mathbf n,T}-z)^{-1}\) is a measurable resolvent field for every
nonreal \(z\).

Let
\[
 \mathsf P_{\mathbf n}(T)
 :=\mathbf 1_{\{0\}}(C_{\mathbf n,T})
 =\mathbf 1_{\{E_{\mathbf n}(T)\}}(A_T).
\]
This singleton projection is not assumed continuous at a crossing.
Instead, the spectral theorem gives the pointwise strong limit
\begin{equation}
 \mathsf P_{\mathbf n}(T)
 =\operatorname*{s-lim}_{k\to\infty}
  (I+kC_{\mathbf n,T}^2)^{-1},
\label{native-di-kernel-projection}
\end{equation}
and the exact resolvent formula
\begin{equation}
 (I+kC_{\mathbf n,T}^2)^{-1}
 =\frac1{2i\sqrt{k}}
 \left[
 (C_{\mathbf n,T}-i/\sqrt{k})^{-1}
 -(C_{\mathbf n,T}+i/\sqrt{k})^{-1}
 \right]
\label{native-di-kernel-projection-resolvent}
\end{equation}
shows that every approximating field is strongly measurable.  Therefore
\(T\mapsto\mathsf P_{\mathbf n}(T)v\) is strongly measurable for every
\(v\in\mathscr H_*\).  Proposition
\ref{cusp-native-full-flux-landau-theorem} says that
\(\mathsf P_{\mathbf n}(T)\ne0\) for every \(T\), including crossings.

Choose a countable dense sequence \((v_k)_{k\ge1}\) in the separable
space \(\mathscr H_*\), and define the measurable partition
\[
 X_k=
 \{T:\|\mathsf P_{\mathbf n}(T)v_k\|>0\}
 \cap
 \bigcap_{\ell<k}
 \{T:\|\mathsf P_{\mathbf n}(T)v_\ell\|=0\}.
\]
The sets \(X_k\) are pairwise disjoint and cover \(I\): a nonzero bounded
projection cannot annihilate a dense set.  On \(X_k\), set
\[
 \phi_{\mathbf n}(T)
 =\frac{\mathsf P_{\mathbf n}(T)v_k}
 {\|\mathsf P_{\mathbf n}(T)v_k\|}.
\]
Countable measurable patching gives a strongly measurable unit vector.
The range of the singleton spectral projection lies in
\(\operatorname{Dom}(A_T)\) and satisfies
\((A_T-E_{\mathbf n}(T))\phi_{\mathbf n}(T)=0\), which proves
\eqref{native-di-measurable-eigenvector}.
\end{proof}

\subsubsection{Full spectrum and spectral type}

\begin{theorem}[Purely continuous full half-line spectrum]
\label{native-di-main-theorem}
For both the fixed-fibre and varying-fibre maximal radial multiplication
operators,
\begin{equation}
\begin{aligned}
 \sigma(\widetilde{\mathscr M}_\eta)
 &=\sigma_{\mathrm{ess}}(\widetilde{\mathscr M}_\eta)
 =\sigma_{\mathrm c}(\widetilde{\mathscr M}_\eta)
 =[0,\infty),\\
 \sigma_{\mathrm p}(\widetilde{\mathscr M}_\eta)
 &=\varnothing,
 \qquad
 \ker\widetilde{\mathscr M}_\eta=\{0\},
\end{aligned}
\label{native-di-fixed-spectrum-result}
\end{equation}
and
\begin{equation}
\begin{aligned}
 \sigma(\mathscr M_\eta)
 &=\sigma_{\mathrm{ess}}(\mathscr M_\eta)
 =\sigma_{\mathrm c}(\mathscr M_\eta)
 =[0,\infty),\\
 \sigma_{\mathrm p}(\mathscr M_\eta)
 &=\varnothing,
 \qquad
 \ker\mathscr M_\eta=\{0\}.
\end{aligned}
\label{native-di-varying-spectrum-result}
\end{equation}
Here \(\sigma_{\mathrm{ess}}\) is the Weyl essential spectrum.
\end{theorem}

\begin{proof}
It is enough to prove \eqref{native-di-fixed-spectrum-result}, because
\eqref{native-di-intertwining} is a unitary equivalence.
Proposition~\ref{native-di-self-adjoint-proposition} and nonnegativity give
\[
 \sigma(\widetilde{\mathscr M}_\eta)\subset[0,\infty).
\]

\paragraph{The zero-energy singular sequence.}
Let \(a_j=T_1+2j+1\).  Using the measurable normalized ground eigenvector
from Lemma~\ref{native-di-measurable-eigenvector-lemma}, define
\[
 \Psi_j(T)
 :=\mathbf 1_{[a_j,a_j+1]}(T)\phi_{(0,0)}(T).
\]
These vectors have norm one and have pairwise disjoint supports, hence are
orthonormal and converge weakly to zero.  They belong to the maximal
domain because
\[
 A_T\Psi_j(T)=
 \mathbf 1_{[a_j,a_j+1]}(T)E_0(T)\phi_{(0,0)}(T)
\]
is strongly measurable and square-integrable.  From
\eqref{cusp-native-full-flux-ground-energy},
\[
 \|\widetilde{\mathscr M}_\eta\Psi_j\|^2
 =\int_{a_j}^{a_j+1}E_0(T)^2\,\mathrm d T
 =\frac{196\pi^2}{a_j^2}+O(a_j^{-3})
 \longrightarrow0.
\]
Thus \((\Psi_j)\) is a singular Weyl sequence at zero, and
\[
 0\in\sigma_{\mathrm{ess}}(\widetilde{\mathscr M}_\eta).
\]

\paragraph{Positive energy: escaping crossings before disjoint supports.}
Fix \(\lambda>0\), and fix once and for all
\(T_\bullet=T_1+1\in I\).  For \(n\in\mathbb Z_{\ge0}\), retain the
specific canonical band
\begin{equation}
 e_n(T):=E_{(0,n)}(T)
 =\mathcal B_{T,+}+(2n+1)\mathcal B_{T,-}.
\label{native-di-crossing-band}
\end{equation}
Since \(\mathcal B_{T_\bullet,-}>0\),
\[
 e_n(T_\bullet)\longrightarrow\infty
 \quad(n\to\infty).
\]
For each fixed \(n\), equations
\eqref{cusp-native-full-flux-Bminus-expansion} and
\eqref{cusp-native-full-flux-Bplus-expansion} give
\[
 e_n(T)\longrightarrow0
 \quad(T\to\infty).
\]
Hence, for every sufficiently large \(n\), the intermediate value theorem
gives at least one crossing parameter \(r_n>T_\bullet\) satisfying
\begin{equation}
 e_n(r_n)=\lambda.
\label{native-di-crossing}
\end{equation}

We prove explicitly that every such choice of crossing parameters escapes
to infinity.  Fix \(M>T_\bullet\).  Positivity and continuity give
\[
 c_M:=\min_{T\in[T_\bullet,M]}\mathcal B_{T,-}>0.
\]
For every \(T\in[T_\bullet,M]\),
\[
 e_n(T)\ge(2n+1)c_M.
\]
Thus, once \((2n+1)c_M>\lambda\), equation
\eqref{native-di-crossing} has no solution in
\([T_\bullet,M]\), and the chosen \(r_n\) must satisfy \(r_n>M\).
Since \(M\) was arbitrary,
\begin{equation}
 r_n\longrightarrow\infty.
\label{native-di-crossings-escape}
\end{equation}

Only now, before choosing any support, pass to a subsequence
\(n_j\to\infty\) such that, with \(r_j:=r_{n_j}\),
\begin{equation}
 r_{j+1}>r_j+2.
\label{native-di-separated-crossings}
\end{equation}
Continuity of \(e_{n_j}\) at \(r_j\) permits the subsequent choice of
\(\varepsilon_j\in(0,\tfrac14)\) such that
\[
 J_j:=(r_j-\varepsilon_j,r_j+\varepsilon_j)\Subset I,
 \qquad
 \sup_{T\in J_j}|e_{n_j}(T)-\lambda|<\frac1j.
\]
The order of construction matters: the separation
\eqref{native-di-separated-crossings} makes the intervals \(J_j\)
pairwise disjoint.

Take the measurable normalized eigenvector
\(\phi_{(0,n_j)}(T)\) from
Lemma~\ref{native-di-measurable-eigenvector-lemma}; the singleton
projection used there remains valid if the selected band crosses another
one inside \(J_j\).  Define
\begin{equation}
 \Theta_j(T)
 :=|J_j|^{-1/2}\mathbf 1_{J_j}(T)\phi_{(0,n_j)}(T).
\label{native-di-positive-Weyl-sequence}
\end{equation}
The sections \(\Theta_j\) have norm one and are orthogonal by disjointness
of their supports.  They lie in the maximal domain because
\[
 A_T\Theta_j(T)
 =|J_j|^{-1/2}\mathbf 1_{J_j}(T)
   e_{n_j}(T)\phi_{(0,n_j)}(T)
\]
is strongly measurable and square-integrable.  Moreover,
\[
 \|(\widetilde{\mathscr M}_\eta-\lambda)\Theta_j\|^2
 =\frac1{|J_j|}\int_{J_j}
   |e_{n_j}(T)-\lambda|^2\,\mathrm d T
 <\frac1{j^2}.
\]
Every orthonormal sequence is weakly null, so
\((\Theta_j)\) is a singular Weyl sequence at \(\lambda\).
Thus every \(\lambda>0\) belongs to the Weyl essential spectrum.  Together
with the zero sequence and nonnegativity, this proves
\begin{equation}
 \sigma(\widetilde{\mathscr M}_\eta)
 =\sigma_{\mathrm{ess}}(\widetilde{\mathscr M}_\eta)
 =[0,\infty).
\label{native-di-spectrum-essential}
\end{equation}

\paragraph{Real analyticity, absence of flat bands, and no eigenvalues.}
Equation \eqref{cusp-native-full-flux-Xi-identity} writes
\[
 \Xi_T=\sqrt{S_T^2+24D_T}
\]
with a strictly positive real-analytic radicand.  The positive square-root
function is real analytic on \((0,\infty)\).  Hence
\(\Xi_T\), \(\mathcal B_{T,+}\), \(\mathcal B_{T,-}\), and every
\(E_{\mathbf n}(T)\) are real analytic on the connected interval \(I\).
For each fixed \(\mathbf n\),
\[
 E_{\mathbf n}(T)>0\quad(T<\infty),
 \qquad
 \lim_{T\to\infty}E_{\mathbf n}(T)=0.
\]
It follows that no band is constant: a constant band would have constant
value equal to its limit zero, contradicting strict positivity at every
finite \(T\).  This proves the absence of flat bands, rather than assuming
it.

Fix \(\lambda\in\mathbb R\).  For every \(\mathbf n\), the real-analytic
function \(E_{\mathbf n}-\lambda\) is not identically zero, so its zero set
in \(I\) is discrete and has Lebesgue measure zero.  Since
\(\mathbb Z_{\ge0}^2\) is countable,
\begin{equation}
 Z_\lambda
 :=\{T\in I:\lambda\in\sigma_{\mathrm p}(A_T)\}
 =\bigcup_{\mathbf n\in\mathbb Z_{\ge0}^2}
   \{T:E_{\mathbf n}(T)=\lambda\}
\label{native-di-null-fibre-level-set}
\end{equation}
is a null set.

If
\(\widetilde{\mathscr M}_\eta\Psi=\lambda\Psi\), then
\((A_T-\lambda)\Psi(T)=0\) for almost every \(T\).  At every point where
\(\Psi(T)\ne0\), Proposition
\ref{cusp-native-full-flux-landau-theorem} forces \(T\in Z_\lambda\).
Thus the essential support of \(\Psi\) is contained in the null set
\eqref{native-di-null-fibre-level-set}, and \(\Psi=0\).  A self-adjoint
operator has only real eigenvalues, so
\[
 \sigma_{\mathrm p}(\widetilde{\mathscr M}_\eta)=\varnothing,
 \qquad
 \ker\widetilde{\mathscr M}_\eta=\{0\}.
\]
Self-adjoint operators have empty residual spectrum.  Combining this with
\eqref{native-di-spectrum-essential} proves that the entire spectrum is
continuous and gives \eqref{native-di-fixed-spectrum-result}.  Unitary
equivalence gives \eqref{native-di-varying-spectrum-result}.
\end{proof}

\subsubsection{Exact interpretive boundary}

\begin{remark}[No radial derivative and no interacting Yang--Mills claim]
The variable \(T\) in
\eqref{native-di-fixed-multiplication} and
\eqref{native-di-varying-multiplication} is a direct-integral parameter.
Neither operator contains a \(T\)-derivative.  In particular, there is no
term such as \(-\partial_T^2\), no radial Sobolev condition, no cusp
boundary condition, and no interaction between distinct \(T\)-fibres.
This is why characteristic functions of radial intervals legitimately
occur in the maximal-domain Weyl sequences above.

The theorem is therefore a complete spectral statement for the specified
auxiliary decomposable family of native full-flux magnetic Laplacians.  It
does not construct, identify, or determine the spectrum or mass gap of an
interacting quantum Yang--Mills Hamiltonian, and it makes no
Osterwalder--Schrader reconstruction claim.
\end{remark}

% Canonical split-Cartan SU(r) extension module, 2026-09-03.
% This is a delta to cusp_native_full_flux_magnetic.tex.  It imports that
% module's native torus, full-flux line bundle, orientations, exact magnetic
% frequencies, and SU(2) calculation.  It deliberately does not duplicate
% the radial direct-integral operator or the charge-two SU(2) proof.

\section{The native full-flux split-Cartan \texorpdfstring{\(SU(r)\)}{SU(r)} extension}
\label{cusp-native-full-flux-suN-section}

This section fixes an integer \(r\geq2\) and extends the full-flux
construction of \cref{cusp-native-full-flux-section} from its already proved
split \(SU(2)\) case to arbitrary split-Cartan \(SU(r)\) backgrounds.  Every
object remains on the same regular flat four-torus
\(F_T=\mathbb R^4/P_T\mathbb Z^4\), with the unrescaled Euclidean quotient
metric.  In particular, the integral two-form \(\eta_T\), the line bundle
\(\mathscr L_{\eta,T}\), the two orientations, and the exact quantities
\[
 \begin{gathered}
  L_T,\quad q_T,\quad m_T,\quad D_T=L_T q_T+6m_T^2>0,\\
  S_T=6L_T-q_T,\qquad
  \Xi_T=\sqrt{(6L_T+q_T)^2+144m_T^2}.
 \end{gathered}
\]
are precisely those of
\cref{cusp-native-full-flux-canonical-scalars,%
cusp-native-full-flux-S-Xi}.  The charge-one frequencies are
\begin{equation}
 \mathcal B_{T,+}=\frac{\pi(S_T+\Xi_T)}{D_T},
 \qquad
 \mathcal B_{T,-}=\frac{\pi(\Xi_T-S_T)}{D_T},
 \label{cusp-native-full-flux-suN-signed-frequencies}
\end{equation}
and we shall use the globally ordered notation
\begin{equation}
 \begin{aligned}
  \mathcal B_{T,\max}
  &=\max\{\mathcal B_{T,+},\mathcal B_{T,-}\},&
  \mathcal B_{T,\min}
  &=\min\{\mathcal B_{T,+},\mathcal B_{T,-}\},\\
  \rho_T^*&=\frac{\mathcal B_{T,\max}}{\mathcal B_{T,\min}}.
 \end{aligned}
 \label{cusp-native-full-flux-suN-ordered-frequencies}
\end{equation}
Thus every assertion below remains valid on either side of \(S_T=0\).

\subsection{The split bundle, determinant, and Chern number}

Fix a nonzero trace-zero integral weight vector
\begin{equation}
 \mathbf m=(m_1,\ldots,m_r)\in\mathbb Z^r,
 \qquad \sum_{j=1}^r m_j=0,
 \qquad s_2(\mathbf m):=\sum_{j=1}^r m_j^2.
 \label{cusp-native-full-flux-suN-weight-vector}
\end{equation}
Define
\begin{equation}
 \mathscr E_{\mathbf m,T}:=
 \bigoplus_{j=1}^r\mathscr L_{\eta,T}^{\otimes m_j},
 \qquad
 H_{\mathbf m}:=i\operatorname{diag}(m_1,\ldots,m_r)
 \in\mathfrak{su}(r).
 \label{cusp-native-full-flux-suN-split-bundle}
\end{equation}

\begin{proposition}[Parallel determinant trivialization and exact \(c_2\)]
\label{cusp-native-full-flux-suN-topology}
The ordered splitting in
\cref{cusp-native-full-flux-suN-split-bundle} has a canonical parallel
determinant trivialization.  With \(c=[\eta_T]\),
\begin{equation}
 c_1(\mathscr E_{\mathbf m,T})=0,
 \qquad
 c_2(\mathscr E_{\mathbf m,T})
 =-\frac{s_2(\mathbf m)}2c^2.
 \label{cusp-native-full-flux-suN-Chern-classes}
\end{equation}
The marked and physical-complex Chern numbers are respectively
\begin{equation}
 \left\langle c_2,[F_T]_{\mathrm{mark}}\right\rangle=-6s_2(\mathbf m),
 \qquad
 k_{\mathbf m}:=
 \left\langle c_2,[F_T]_{x\text{-}\mathrm{complex}}\right\rangle
 =6s_2(\mathbf m)>0.
 \label{cusp-native-full-flux-suN-Chern-numbers}
\end{equation}
Moreover \(s_2(\mathbf m)\) is even.  Hence this direct-sum family realizes
only \(k_{\mathbf m}\in12\mathbb Z_{>0}\), not all topological \(SU(r)\)
sectors.
\end{proposition}

\begin{proof}
Let \(j_n^\eta(y)\) denote the factor of automorphy in
\cref{cusp-native-full-flux-cocycle}.  The transition matrix of the ordered
sum is
\[
 \operatorname{diag}\bigl((j_n^\eta)^{m_1},\ldots,
 (j_n^\eta)^{m_r}\bigr).
\]
Its determinant is \((j_n^\eta)^{\sum_j m_j}=1\) identically.  The tensor
product of the ordered unit frames therefore defines a global determinant
section, and the determinant connection has connection form proportional to
\(\sum_j m_j=0\); the section is parallel.

The splitting principle gives
\[
 c(\mathscr E_{\mathbf m,T})=\prod_{j=1}^r(1+m_j c),
 \qquad
 c_2=\sum_{i<j}m_i m_j c^2.
\]
Since
\[
 2\sum_{i<j}m_i m_j
 =\left(\sum_j m_j\right)^2-\sum_j m_j^2=-s_2(\mathbf m),
\]
this proves \cref{cusp-native-full-flux-suN-Chern-classes}.  The two
evaluations follow from
\cref{cusp-native-full-flux-two-orientation-pairings}; retaining its
orientation reversal is essential for the sign.  Finally
\[
 s_2(\mathbf m)=\sum_j m_j^2\equiv\sum_j m_j\equiv0\pmod2,
\]
because \(a^2\equiv a\pmod2\).  This proves the divisibility statement.
\end{proof}

The diagonal connection induced by the summands is
\begin{equation}
 \nabla^{\mathbf m,T}
 =\mathrm d-2\pi(6y^1\,\mathrm dy^4+y^2\,\mathrm dy^3)H_{\mathbf m},
 \qquad
 F_{\mathbf m,T}=-2\pi H_{\mathbf m}\eta_T.
 \label{cusp-native-full-flux-suN-connection-curvature}
\end{equation}
We retain the invariant inner product and functional normalization
\begin{equation}
 \begin{gathered}
  \langle X,Y\rangle_{\mathfrak{su}(r)}
  =-\frac12\operatorname{Tr}(XY),
  \qquad
  \lVert H_{\mathbf m}\rVert^2=\frac12s_2(\mathbf m),\\
  \operatorname{YM}_T(A)=\frac1{2g_{\mathrm{YM}}^2}
  \int_{F_T}|F_A|^2\,\mathrm d\operatorname{vol}_T.
 \end{gathered}
 \label{cusp-native-full-flux-suN-normalizations}
\end{equation}

\begin{theorem}[Yang--Mills criticality and exact action]
\label{cusp-native-full-flux-suN-action-theorem}
The connection \(\nabla^{\mathbf m,T}\) is non-flat, reducible, and
Yang--Mills.  Its action is
\begin{align}
 \operatorname{YM}_T(\nabla^{\mathbf m,T})
 &=\frac{\pi^2s_2(\mathbf m)}{g_{\mathrm{YM}}^2D_T}
   \bigl(q_T^2+36L_T^2+72m_T^2\bigr)
 \label{cusp-native-full-flux-suN-action-coordinate}\\
 &=\frac{2\pi^2k_{\mathbf m}}{g_{\mathrm{YM}}^2}
   +\frac{\pi^2s_2(\mathbf m)S_T^2}
          {g_{\mathrm{YM}}^2D_T}.
 \label{cusp-native-full-flux-suN-action-excess}
\end{align}
On the canonical cusp tail,
\begin{equation}
 \operatorname{YM}_T(\nabla^{\mathbf m,T})
 =\frac{\pi^2s_2(\mathbf m)}{g_{\mathrm{YM}}^2}
 \left(37+\frac{35B}{T}+O(T^{-2})\right).
 \label{cusp-native-full-flux-suN-action-asymptotic}
\end{equation}
\end{theorem}

\begin{proof}
Both \(H_{\mathbf m}\) and \(\eta_T\) are parallel for the diagonal
connection and flat quotient metric, and their brackets vanish.  Hence
\(d_{\nabla^{\mathbf m,T}}^*F_{\mathbf m,T}=0\).  Non-flatness follows
from \(\mathbf m\ne0\) and \(\eta_T\ne0\).

In the physical orthonormal frame used in
\cref{cusp-native-full-flux-normal-form},
\[
 |\eta_T|^2
 =\frac{\mathcal B_{T,+}^2+\mathcal B_{T,-}^2}{4\pi^2}.
\]
Using \cref{cusp-native-full-flux-suN-normalizations},
\(\operatorname{vol}(F_T)=D_T\), and the exact frequency square sum from
the native module yields
\cref{cusp-native-full-flux-suN-action-coordinate}.  The identity
\[
 q_T^2+36L_T^2+72m_T^2
 =12(L_T q_T+6m_T^2)+(6L_T-q_T)^2
 =12D_T+S_T^2
\]
and \(k_{\mathbf m}=6s_2(\mathbf m)\) give
\cref{cusp-native-full-flux-suN-action-excess}.  Substituting the retained
expansions of \(L_T,q_T,m_T,D_T\), without suppressing \(H,B,m\), gives
\[
 \frac{q_T^2+36L_T^2+72m_T^2}{D_T}
 =37+\frac{35B}{T}+O(T^{-2}),
\]
which proves the last formula.
\end{proof}

\subsection{Every nonzero integral root charge}

The native module proves the charge-one ladder decomposition and the
charge-two specialization used by its \(SU(2)\) theorem.  The following
lemma records, and proves, the exact integer-charge extension needed here.

\begin{lemma}[Charge-\(d\) Landau decomposition]
\label{cusp-native-full-flux-suN-charge-d-Landau}
Let \(d\in\mathbb Z\setminus\{0\}\), put \(q=|d|\), and equip
\(\mathscr L_{\eta,T}^{\otimes d}\) with its tensor-power connection.  Its
scalar covariant Laplacian has joint levels
\begin{equation}
 q\bigl((2a+1)\mathcal B_{T,+}+(2b+1)\mathcal B_{T,-}\bigr),
 \qquad a,b\in\mathbb Z_{\ge0},
 \label{cusp-native-full-flux-suN-charge-d-levels}
\end{equation}
and every joint space has complex dimension \(6q^2\).  If different pairs
\((a,b)\) give the same numerical value, their orthogonal joint spaces are
to be summed.
\end{lemma}

\begin{proof}
The curvature is \(-2\pi i d\eta_T\).  In the polar construction of
\cref{cusp-native-full-flux-landau-theorem}, replace the auxiliary complex
structure by
\(J_{d,T}=\operatorname{sgn}(d)J_{\eta,T}\).  Then \(d\eta_T\) is positive
of type \((1,1)\), so the vacuum is the holomorphic-section space of the
positive bundle \(\mathscr L_{\eta,T}^{\otimes d}\) on that auxiliary
complex torus.  Its tangent and canonical bundles are trivial.  Kodaira
vanishing and Hirzebruch--Riemann--Roch give
\[
 \dim_{\mathbb C}H^0
 =\frac12\int_{[F_T]_{d\eta\text{-}\mathrm{symp}}}(d\eta_T)^2
 =\frac12q^2\int_{[F_T]_{\mathrm{mark}}}\eta_T^2
 =6q^2.
\]
Here changing \(d\) to \(-d\) does not change the four-dimensional
symplectic orientation.  Scaling the curvature by \(d\) scales both
positive oscillator frequencies by \(q\), while the two canonical
commutation relations and the lowering argument in
\cref{cusp-native-full-flux-landau-theorem} are unchanged after this
rescaling.  Thus every joint level is unitarily obtained from the vacuum and
has the asserted energy and dimension.  Orthogonal joint spaces remain a
direct sum when their numerical energies coincide.
\end{proof}

\subsection{Closed gauge-fixed and Coulomb operators}

Write \(A=\nabla^{\mathbf m,T}\) in this subsection.  Let
\begin{equation}
 \mathscr H_{\mathbf m,T}^1
 =L^2\Omega^1(F_T;\operatorname{ad}\mathscr E_{\mathbf m,T}),
 \qquad
 \mathscr H_{\mathbf m,T}^0
 =L^2\Omega^0(F_T;\operatorname{ad}\mathscr E_{\mathbf m,T})
 \label{cusp-native-full-flux-suN-real-Hilbert-spaces}
\end{equation}
as real Hilbert spaces.  Define the maximal distributional codifferential
by
\begin{align}
 \operatorname{Dom}D_{\mathbf m,T}^{\max}
 =\{a\in\mathscr H_{\mathbf m,T}^1:\;&
 \exists u\in\mathscr H_{\mathbf m,T}^0\text{ such that}\notag\\[-1mm]
 &\langle a,d_A\phi\rangle_{L^2}
 =\langle u,\phi\rangle_{L^2}
 \quad\text{for every }\notag\\[-1mm]
 &\phi\in C^\infty(F_T;\operatorname{ad}\mathscr E_{\mathbf m,T})\},
 \label{cusp-native-full-flux-suN-max-codiff-domain}\\
 D_{\mathbf m,T}^{\max}a={}&u,
 \qquad
 \mathscr S_{\mathbf m,T}:=\ker D_{\mathbf m,T}^{\max}.
 \label{cusp-native-full-flux-suN-Coulomb-slice}
\end{align}
Density of smooth test scalars makes \(u\) unique.  If
\(a_j\to a\) and \(D_{\mathbf m,T}^{\max}a_j\to u\) in \(L^2\), passage to
the limit in the defining identity gives
\(D_{\mathbf m,T}^{\max}a=u\).  Thus \(D_{\mathbf m,T}^{\max}\) is closed
and \(\mathscr S_{\mathbf m,T}\) is a closed subspace.

Let \(\Delta_A^{(0)}\) be the nonnegative self-adjoint scalar operator of
the closed form
\begin{equation}
 \mathfrak q_{\mathbf m,T}^{(0)}[\phi,\psi]
 =\langle d_A\phi,d_A\psi\rangle_{L^2},
 \qquad
 \operatorname{Dom}\mathfrak q_{\mathbf m,T}^{(0)}=H^1\Omega^0.
 \label{cusp-native-full-flux-suN-scalar-form}
\end{equation}
It has domain \(H^2\Omega^0\) and equals \(d_A^*d_A\).  In a constant
orthonormal frame set
\begin{equation}
 (V_{\mathbf m,T}a)_\nu
 =-2\sum_{\mu=1}^4[F_{\mathbf m,T,\nu\mu},a_\mu].
 \label{cusp-native-full-flux-suN-spin-endomorphism}
\end{equation}
Invariance of the compact Lie-algebra inner product and
\(F_{\nu\mu}=-F_{\mu\nu}\) show directly that \(V_{\mathbf m,T}\) is a
bounded symmetric endomorphism.  Define
\begin{equation}
 \mathfrak j_{\mathbf m,T}^{\mathrm{gf}}[a,b]
 =\langle\nabla_A a,\nabla_A b\rangle_{L^2}
  +\langle V_{\mathbf m,T}a,b\rangle_{L^2},
 \qquad
 \operatorname{Dom}\mathfrak j_{\mathbf m,T}^{\mathrm{gf}}
 =H^1\Omega^1.
 \label{cusp-native-full-flux-suN-Jacobi-form}
\end{equation}

\begin{proposition}[Closed realization and exact Coulomb reduction]
\label{cusp-native-full-flux-suN-Coulomb-reduction}
The form in \cref{cusp-native-full-flux-suN-Jacobi-form} is densely defined,
closed, symmetric, and bounded below.  Its self-adjoint operator has compact
resolvent, domain \(H^2\Omega^1\), and expression
\begin{equation}
 (\mathcal J_{\mathbf m,T}^{\mathrm{gf}}a)_\nu
 =-\sum_{\mu=1}^4\nabla_{A,e_\mu}\nabla_{A,e_\mu}a_\nu
  -2\sum_{\mu=1}^4[F_{\mathbf m,T,\nu\mu},a_\mu].
 \label{cusp-native-full-flux-suN-Jacobi-expression}
\end{equation}
For every \(z\in\mathbb C\setminus\mathbb R\) and
\(f\in\operatorname{Dom}D_{\mathbf m,T}^{\max}\),
\begin{equation}
 D_{\mathbf m,T}^{\max}
 (\mathcal J_{\mathbf m,T}^{\mathrm{gf}}-z)^{-1}f
 =(\Delta_A^{(0)}-z)^{-1}D_{\mathbf m,T}^{\max}f.
 \label{cusp-native-full-flux-suN-resolvent-intertwining}
\end{equation}
Consequently \(\mathscr S_{\mathbf m,T}\) reduces
\(\mathcal J_{\mathbf m,T}^{\mathrm{gf}}\), and
\begin{equation}
 \mathcal J_{\mathbf m,T}
 :=\mathcal J_{\mathbf m,T}^{\mathrm{gf}}
   \big|_{\mathscr S_{\mathbf m,T}},
 \qquad
 \operatorname{Dom}\mathcal J_{\mathbf m,T}
 =H^2\Omega^1\cap\mathscr S_{\mathbf m,T}
 \label{cusp-native-full-flux-suN-Coulomb-operator}
\end{equation}
is self-adjoint, with compact resolvent and closed form
\(\mathfrak j_{\mathbf m,T}\) equal to the restriction of
\(\mathfrak j_{\mathbf m,T}^{\mathrm{gf}}\) to
\(H^1\Omega^1\cap\mathscr S_{\mathbf m,T}\).  The physical Yang--Mills
Hessian on this slice is
\begin{equation}
 \operatorname{YM}_T''\big|_A
 =g_{\mathrm{YM}}^{-2}\mathfrak j_{\mathbf m,T}.
 \label{cusp-native-full-flux-suN-physical-Hessian}
\end{equation}
Thus an eigenvalue \(\nu\) of \(\mathcal J_{\mathbf m,T}\) is the physical
Hessian eigenvalue \(\nu/g_{\mathrm{YM}}^2\), while its sign and Morse-index
contribution are unchanged.
\end{proposition}

\begin{proof}
The first term of
\cref{cusp-native-full-flux-suN-Jacobi-form}, after addition of the
\(L^2\)-norm, is an equivalent \(H^1\)-norm.  The second term is a bounded
symmetric form.  The bounded perturbation theorem for closed forms and the
representation theorem therefore give a lower-bounded self-adjoint
operator.  Its principal symbol is \(|\xi|^2\) times the identity, so
elliptic regularity gives the \(H^2\) domain.  Rellich compactness on the
compact torus gives compact resolvent.

The two-parameter curvature expansion and the flat Weitzenb\"ock identity,
with the inner product and functional normalization in
\cref{cusp-native-full-flux-suN-normalizations}, give
\[
 \operatorname{YM}_T''\big|_A(a,b)
 +g_{\mathrm{YM}}^{-2}\langle d_A^*a,d_A^*b\rangle
 =g_{\mathrm{YM}}^{-2}
   \mathfrak j_{\mathbf m,T}^{\mathrm{gf}}[a,b].
\]
This proves the Hessian normalization after restriction to Coulomb gauge.

It remains to prove that the restriction is an operator rather than a formal
constraint.  For a smooth adjoint scalar \(\phi\), write
\(\nabla_\mu=\nabla_{A,e_\mu}\).  Commuting all three derivatives and using
the Yang--Mills equation gives
\begin{align*}
 &(\mathcal J_{\mathbf m,T}^{\mathrm{gf}}d_A\phi
       -d_A\Delta_A^{(0)}\phi)_\nu\\
 &\quad=-\nabla_\mu\nabla_\mu\nabla_\nu\phi
       +\nabla_\nu\nabla_\mu\nabla_\mu\phi
       -2[F_{\nu\mu},\nabla_\mu\phi]\\
 &\quad=-[\nabla_\mu F_{\mu\nu},\phi]
       =[(d_A^*F_A)_\nu,\phi]=0.
\end{align*}
Formal adjunction yields, on smooth one-forms,
\begin{equation}
 \mathcal J_{\mathbf m,T}^{\mathrm{gf}}d_A
 =d_A\Delta_A^{(0)},
 \qquad
 d_A^*\mathcal J_{\mathbf m,T}^{\mathrm{gf}}
 =\Delta_A^{(0)}d_A^*.
 \label{cusp-native-full-flux-suN-formal-intertwining}
\end{equation}
For
\(u=(\mathcal J_{\mathbf m,T}^{\mathrm{gf}}-z)^{-1}f\), the operator domain
gives \(u\in H^2\), hence \(D_{\mathbf m,T}^{\max}u\in H^1\).  Applying the
second identity distributionally to
\((\mathcal J_{\mathbf m,T}^{\mathrm{gf}}-z)u=f\) gives
\[
 (\Delta_A^{(0)}-z)D_{\mathbf m,T}^{\max}u
 =D_{\mathbf m,T}^{\max}f.
\]
Scalar elliptic regularity and the scalar resolvent give
\cref{cusp-native-full-flux-suN-resolvent-intertwining}.  The two conjugate
resolvents preserve \(\mathscr S_{\mathbf m,T}\); adjunction then makes its
orthogonal complement invariant as well.  The slice is reducing, so the
self-adjoint part has the operator and form domains displayed above.
\end{proof}

\subsection{Rootwise spectrum, coincidences, index, and nullity}

For \(i\ne j\), the matrix unit \(E_{ij}\) obeys
\[
 [H_{\mathbf m},E_{ij}]=i(m_i-m_j)E_{ij},
\]
so its complex root line is
\begin{equation}
 \operatorname{Hom}(\mathscr L_{\eta,T}^{m_j},
                     \mathscr L_{\eta,T}^{m_i})
 =\mathscr L_{\eta,T}^{m_i-m_j},
 \qquad q_{ij}:=|m_i-m_j|.
 \label{cusp-native-full-flux-suN-root-charge}
\end{equation}
The factor \(q_{ij}\) is the complete charge normalization; for the existing
\(SU(2)\) weights \((1,-1)\), it gives \(q_{12}=2\).

Let the distinct numerical values among \(m_1,\ldots,m_r\) have
multiplicities \(r_1,\ldots,r_\ell\).  The neutral centralizer is
\begin{equation}
 \mathfrak z_{\mathbf m}
 =\{X\in\mathfrak{su}(r):[H_{\mathbf m},X]=0\}
 \cong\mathfrak s\!\left(\bigoplus_{a=1}^{\ell}\mathfrak u(r_a)\right),
 \qquad
 \dim_{\mathbb R}\mathfrak z_{\mathbf m}
 =\sum_{a=1}^{\ell}r_a^2-1.
 \label{cusp-native-full-flux-suN-centralizer}
\end{equation}

\begin{theorem}[Complete negative spectrum and numerical multiplicities]
\label{cusp-native-full-flux-suN-Morse-theorem}
Put
\begin{equation}
 N_-(T):=\#\{n\in\mathbb Z_{\ge0}:2n+1<\rho_T^*\}
 =\left\lceil\frac{\rho_T^*-1}{2}\right\rceil.
 \label{cusp-native-full-flux-suN-N-minus}
\end{equation}
For each positive root \(i<j\) with \(q_{ij}>0\), the complete negative
spectrum of the normalized Coulomb operator is
\begin{equation}
 \nu_{ij,n}(T)
 =q_{ij}\bigl((2n+1)\mathcal B_{T,\min}
                  -\mathcal B_{T,\max}\bigr),
 \qquad 0\le n<N_-(T).
 \label{cusp-native-full-flux-suN-negative-bands}
\end{equation}
The contribution of that positive-root pair to each displayed band has real
dimension \(12q_{ij}^2\).  For an actual numerical eigenvalue
\(\lambda<0\), its full real multiplicity is therefore
\begin{equation}
 \operatorname{mult}_{\mathbb R}(\lambda)
 =\sum_{\substack{i<j,\ q_{ij}>0,\ 0\le n<N_-(T)\\
                   \nu_{ij,n}(T)=\lambda}}
   12q_{ij}^2.
 \label{cusp-native-full-flux-suN-coincident-multiplicity}
\end{equation}
In particular, coincident root charges or accidental coincidences between
different \(q_{ij},n\) add rather than being counted as one band.  The real
Morse index is
\begin{equation}
 \boxed{
 \operatorname{ind}_{\mathbb R}\mathcal J_{\mathbf m,T}
 =12N_-(T)\sum_{i<j}(m_i-m_j)^2
 =12rN_-(T)s_2(\mathbf m).}
 \label{cusp-native-full-flux-suN-index}
\end{equation}
Every eigenvalue in
\cref{cusp-native-full-flux-suN-negative-bands} is divided by
\(g_{\mathrm{YM}}^2\) for the physical Hessian.
\end{theorem}

\begin{proof}
On the \(E_{ij}\) line, the scalar oscillator is the charge-\(q_{ij}\)
operator of \cref{cusp-native-full-flux-suN-charge-d-Landau}.  In a
skew-normal orthonormal coframe, the four circular one-form polarizations
diagonalize \cref{cusp-native-full-flux-suN-spin-endomorphism}; the two
polarizations in a plane of charge-\(q_{ij}\) frequency
\(q_{ij}\mathcal B\) shift the scalar energy by
\(\pm2q_{ij}\mathcal B\).  The unfavorable polarizations are positive.
For the favorable polarization in the high-frequency plane, high-plane
oscillator number zero leaves
\[
 q_{ij}\bigl((2n+1)\mathcal B_{T,\min}
             -\mathcal B_{T,\max}\bigr).
\]
Increasing its high-plane oscillator number makes it positive.  The
favorable low-frequency polarization has bottom
\(q_{ij}(\mathcal B_{T,\max}-\mathcal B_{T,\min})\ge0\), and increasing
either number cannot lower it.  This proves that the list is exhaustive.

Every joint scalar space on one complex root line has dimension
\(6q_{ij}^2\).  The opposite root is its anti-linear conjugate.  The pair is
the complexification of a real root subspace, so its real dimension is
\(12q_{ij}^2\).  Root subspaces and joint oscillator spaces are orthogonal;
hence equality of numerical eigenvalues produces precisely the direct sum
in \cref{cusp-native-full-flux-suN-coincident-multiplicity}.

The strict inequality in
\cref{cusp-native-full-flux-suN-N-minus} counts the integers correctly even
when \(\rho_T^*\) is an odd integer.  Summing dimensions gives the first
index formula.  Finally,
\[
 \sum_{i<j}(m_i-m_j)^2
 =r\sum_i m_i^2-\left(\sum_i m_i\right)^2
 =rs_2(\mathbf m),
\]
which gives the second formula.

It remains to justify Coulomb gauge.  If
\(\mathcal J_{\mathbf m,T}^{\mathrm{gf}}a=\lambda a\) with
\(\lambda<0\), elliptic regularity and
\cref{cusp-native-full-flux-suN-formal-intertwining} give
\[
 \Delta_A^{(0)}d_A^*a=\lambda d_A^*a.
\]
The scalar operator is nonnegative, so \(d_A^*a=0\).  Thus the negative
gauge-fixed and Coulomb eigenspaces coincide, completing the proof.
\end{proof}

\begin{theorem}[Exact neutral and wall nullities]
\label{cusp-native-full-flux-suN-nullity-theorem}
Away from an odd wall, the neutral Coulomb kernel has real dimension
\begin{equation}
 \operatorname{null}_{\mathrm{neutral}}\mathcal J_{\mathbf m,T}
 =4\left(\sum_{a=1}^{\ell}r_a^2-1\right).
 \label{cusp-native-full-flux-suN-neutral-nullity}
\end{equation}
At a wall \(\rho_T^*=2n+1>1\), the nonzero-root part adds
\begin{equation}
 12\sum_{i<j}q_{ij}^2=12r s_2(\mathbf m)
 \label{cusp-native-full-flux-suN-odd-wall-nullity}
\end{equation}
real zero modes.  At the equality wall \(S_T=0\), equivalently
\(\rho_T^*=1\), both favorable plane polarizations vanish and the nonzero
root contribution doubles to
\begin{equation}
 24\sum_{i<j}q_{ij}^2=24r s_2(\mathbf m).
 \label{cusp-native-full-flux-suN-equality-wall-nullity}
\end{equation}
The total nullity at a wall is the sum of the neutral term and the stated
root term.
\end{theorem}

\begin{proof}
On \(\mathfrak z_{\mathbf m}\), both the induced connection commutator and
the spin endomorphism vanish.  The operator is therefore the ordinary
coefficientwise Laplacian on one-forms.  Its kernel on a flat four-torus is
the space of constant one-forms, all of which are Coulomb.  Tensoring the
four-dimensional real harmonic-one-form space with
\(\mathfrak z_{\mathbf m}\) proves
\cref{cusp-native-full-flux-suN-neutral-nullity}.

If \(\rho_T^*=2n+1>1\), the high-plane favorable polarization at high
oscillator number zero and low oscillator number \(n\) is the unique zero
branch for each nonzero root.  Its positive-root-pair real dimension is
\(12q_{ij}^2\).  If \(S_T=0\), the two frequencies agree and the lowest
favorable polarization in each of the two planes is zero, so this dimension
doubles.  These charged zero modes are Coulomb: the intertwining identity
shows that their divergences lie in the kernel of the charged scalar
Laplacian, but that scalar Laplacian has strictly positive ground energy by
\cref{cusp-native-full-flux-suN-charge-d-levels}.  Summing over positive
roots and using
\(\sum_{i<j}q_{ij}^2=rs_2(\mathbf m)\) proves the result.
\end{proof}

For \(n\ge1\), the \(n\)-th odd wall is equivalently
\begin{equation}
 \rho_T^*=2n+1
 \iff n\Xi_T=(n+1)|S_T|
 \iff 24n^2D_T=(2n+1)S_T^2.
 \label{cusp-native-full-flux-suN-wall-equations}
\end{equation}
The final polynomial equation must be accompanied by the unsigned ordering
in the middle expression if one needs to distinguish the two sides
\(S_T>0\) and \(S_T<0\).  The \(n=0\) wall is exactly \(S_T=0\).

On the sufficiently deep cusp tail,
\(S_T>0\), \(5<\rho_T^*<7\), and \(N_-(T)=3\).  Hence
\begin{equation}
 \operatorname{ind}_{\mathbb R}\mathcal J_{\mathbf m,T}
 =36r s_2(\mathbf m)=6r k_{\mathbf m},
 \label{cusp-native-full-flux-suN-deep-index}
\end{equation}
and the three negative branches have the retained asymptotics
\begin{equation}
 \nu_{ij,n}(T)
 =\frac{q_{ij}\pi}{T}(4n-10)+O(T^{-2}),
 \qquad n=0,1,2.
 \label{cusp-native-full-flux-suN-negative-asymptotics}
\end{equation}
For example, the genuinely \(SU(3)\) data are
\begin{center}
\renewcommand{\arraystretch}{1.2}
\begin{tabular}{@{}lcccc@{}}
\toprule
weights & nonzero \(q_{ij}\) & \(k_{\mathbf m}\) & deep index
& neutral nullity\\
\midrule
\((1,0,-1)\) & \(1,1,2\) & \(12\) & \(216\) & \(8\)\\
\((1,1,-2)\) & \(3,3\) & \(36\) & \(648\) & \(16\)\\
\((2,1,-3)\) & \(1,4,5\) & \(84\) & \(1512\) & \(8\)\\
\bottomrule
\end{tabular}
\end{center}
The middle row exhibits both required edge phenomena: one neutral root from
the repeated weight and coincident numerical bands from the two charge-three
positive roots.  The specialization \(r=2\), \(\mathbf m=(1,-1)\) is
exactly the already proved native \(SU(2)\) theorem and is not reproduced
here.

\subsection{Stability boundary, possible wall bifurcation, and exact scope}

\begin{proposition}[Where the split background is, and is not, unstable]
\label{cusp-native-full-flux-suN-stability-boundary}
If \(S_T\ne0\), then \(\rho_T^*>1\), \(N_-(T)\ge1\), and every nonzero
positive root contributes a negative Coulomb direction.  Since a nonzero
trace-zero weight vector has at least one nonzero root,
\(\nabla^{\mathbf m,T}\) is then not a local minimum modulo gauge.  This
applies in particular throughout the sufficiently deep cusp tail.

If \(S_T=0\), the negative root spectrum is empty.  In that case
\cref{cusp-native-full-flux-suN-action-excess} attains the topological value
\(2\pi^2k_{\mathbf m}/g_{\mathrm{YM}}^2\), while
\cref{cusp-native-full-flux-suN-equality-wall-nullity} supplies root zero
modes.  No assertion of strict stability follows because the Hessian is
degenerate.
\end{proposition}

\begin{proof}
The first paragraph follows immediately from
\cref{cusp-native-full-flux-suN-ordered-frequencies,%
cusp-native-full-flux-suN-Morse-theorem}.  For the second,
\(S_T=0\) makes the two frequencies equal and hence \(\rho_T^*=1\); the
strict negative-band inequality has no solution.  The action and nullity
statements are exactly
\cref{cusp-native-full-flux-suN-action-excess,%
cusp-native-full-flux-suN-nullity-theorem}.
\end{proof}

A negative eigenspace is a space of descent directions, not a
Lyapunov--Schmidt bifurcation kernel.  A local bifurcation analysis must
instead be posed at one of the zero-mode parameters in
\cref{cusp-native-full-flux-suN-nullity-theorem}.  It must first quotient or
slice the ever-present centralizer-valued harmonic kernel, then identify the
remaining kernel and cokernel of the gauge-fixed nonlinear Yang--Mills map,
and finally compute the parameter transversality and nonlinear obstruction
coefficients.  The present linear calculation supplies the exact candidate
kernel at each wall; it does not assert existence of a nearby irreducible
branch.

The theorem proved here is strictly lower-rung and fibrewise.  It constructs
explicit reducible split-Cartan \(SU(r)\) Yang--Mills critical points on the
already existing regular flat four-torus fibres and determines their
topology, action, normalized Jacobi spectrum, physical Hessian scaling,
Morse index, and nullity.  It constructs no new base, spacetime, cone,
\(S^{12}\), \(S^{24}\), transfer matrix, or interacting quantum Hamiltonian.
The radial multiplication operator already treated in the canonical native
module is not repeated here.

There is also an exact obstruction to interpreting the nonzero-flux bundle
as the restriction of a bundle on a proposed higher carrier.  If
\(i:F_T\to Y\) is any map and \(H^4(Y;\mathbb Z)=0\), then every complex
bundle \(\mathscr V\to Y\) has \(c_2(\mathscr V)=0\), so naturality gives
\[
 c_2(i^*\mathscr V)=i^*c_2(\mathscr V)=0.
\]
This contradicts
\cref{cusp-native-full-flux-suN-Chern-classes,%
cusp-native-full-flux-suN-Chern-numbers} for
\(\mathscr E_{\mathbf m,T}\).  Therefore this nonzero-\(c_2\) bundle cannot
extend over any carrier with vanishing fourth integral cohomology, including
a space diffeomorphic to \(S^6\), \(S^{12}\), or \(S^{24}\).  Changing the
gauge group changes the root system on the fibre; it does not by itself
produce the proposed higher rung or any conclusion about a physical mass
gap.

\section{The higher-rung \texorpdfstring{\(\mathrm{SU}(3)\)}{SU(3)} Gram cone}
\label{sec:higher-rung-su3-gram-cone}

This section isolates the part of the proposed higher-rung construction that
is independent of the candidate complex threefold and proves it from the
underlying linear action.  The compact object used here is the representation
sphere
\[
 \operatorname{Mat}_{3\times2}(\mathbb C)^+,
\]
the one-point compactification of a twelve-real-dimensional vector space.
It must not be identified with a cusp-induced \(S^{12}\) construction unless
an additional based identification has first been chosen.  The proofs below
also distinguish three structures that otherwise look deceptively similar:
the Lorentzian determinant on \(\operatorname{Herm}_2(\mathbb C)\), the
positive-semidefinite Gram cone inside that vector space, and a separately
chosen Euclidean metric on the positive-definite cone.

\subsection{Conventions and the Gram morphism}

Put
\begin{equation}
 \mathscr V=\operatorname{Mat}_{3\times2}(\mathbb C),
 \qquad
 \Gamma:\mathscr V\longrightarrow\operatorname{Herm}_2(\mathbb C),
 \qquad
 \Gamma(A)=A^\dagger A.
 \label{eq:hrgc-gram-map}
\end{equation}
Write \(A=(a\ b)\), with \(a,b\in\mathbb C^3\), and let
\(\mathrm{SU}(3)\) act on the left:
\[
 U\cdot A=UA.
\]
For \(r=0,1,2\), let
\[
 \mathscr V_{(r)}=\{A\in\mathscr V:\operatorname{rank}_{\mathbb C}A=r\}.
\]
The Hermitian products are linear in the second argument.  We orient
\(\mathscr V\) by the complex orientation, in the real coordinate order
\begin{equation}
 \begin{split}
 (&\operatorname{Re}A_{11},\operatorname{Im}A_{11},
   \operatorname{Re}A_{12},\operatorname{Im}A_{12},
   \operatorname{Re}A_{21},\operatorname{Im}A_{21},\\
  &\operatorname{Re}A_{22},\operatorname{Im}A_{22},
   \operatorname{Re}A_{31},\operatorname{Im}A_{31},
   \operatorname{Re}A_{32},\operatorname{Im}A_{32}).
 \end{split}
 \label{eq:hrgc-v-orientation}
\end{equation}
On \(\operatorname{Herm}_2(\mathbb C)\), use the oriented real-linear
coordinates
\begin{equation}
 \Phi(t,x,y,z)=
 \begin{pmatrix}
  t+z&x-iy\\
  x+iy&t-z
 \end{pmatrix},
 \qquad
 \operatorname{or}_{\operatorname{Herm}}=dt\wedge dx\wedge dy\wedge dz.
 \label{eq:hrgc-pauli-coordinates}
\end{equation}
Thus, if \(\Gamma(A)_{12}=\langle a,b\rangle\), then
\begin{equation}
 \begin{aligned}
 t&=\frac12(\lVert a\rVert^2+\lVert b\rVert^2),&
 z&=\frac12(\lVert a\rVert^2-\lVert b\rVert^2),\\
 x&=\operatorname{Re}\langle a,b\rangle,&
 y&=-\operatorname{Im}\langle a,b\rangle.
 \end{aligned}
 \label{eq:hrgc-column-pauli-coordinates}
\end{equation}

\subsection{Exact orbit classification}

\begin{theorem}[The Gram quotient]
\label{thm:hrgc-gram-quotient}
The morphism \(\Gamma\) is \(\mathrm{SU}(3)\)-invariant, has image exactly
\(\operatorname{Herm}_2(\mathbb C)_{\geq0}\), and separates
\(\mathrm{SU}(3)\)-orbits.  It induces a proper homeomorphism
\begin{equation}
 \overline\Gamma:
 \mathscr V/\mathrm{SU}(3)
 \xrightarrow{\ \cong\ }
 \operatorname{Herm}_2(\mathbb C)_{\geq0}.
 \label{eq:hrgc-orbit-homeomorphism}
\end{equation}
\end{theorem}

\begin{proof}
For \(U\in\mathrm{SU}(3)\),
\[
 \Gamma(UA)=A^\dagger U^\dagger UA=A^\dagger A,
\]
so \(\Gamma\) is invariant.  Every matrix \(A^\dagger A\) is positive
semidefinite.  Conversely, for
\(H\in\operatorname{Herm}_2(\mathbb C)_{\geq0}\), let \(H^{1/2}\) be its
unique positive-semidefinite square root and define
\begin{equation}
 s(H)=
 \begin{pmatrix}
  H^{1/2}\\
  0\quad0
 \end{pmatrix}\in\operatorname{Mat}_{3\times2}(\mathbb C).
 \label{eq:hrgc-square-root-section}
\end{equation}
Then \(s(H)^\dagger s(H)=H\), proving surjectivity and giving a continuous
section.

It remains to prove exact orbit separation.  Suppose
\(A^\dagger A=B^\dagger B\).  Define
\[
 T:\operatorname{im}A\longrightarrow\operatorname{im}B,
 \qquad T(Au)=Bu.
\]
If \(Au=0\), then
\[
 0=\lVert Au\rVert^2
  =u^\dagger A^\dagger Au
  =u^\dagger B^\dagger Bu
  =\lVert Bu\rVert^2,
\]
so \(Bu=0\), and \(T\) is well defined.  For \(u,v\in\mathbb C^2\),
\[
 \langle T(Au),T(Av)\rangle
 =u^\dagger B^\dagger Bv
 =u^\dagger A^\dagger Av
 =\langle Au,Av\rangle.
\]
Thus \(T\) is a unitary isomorphism of the image subspaces.  Extend it to
some \(U_0\in\mathrm U(3)\) satisfying \(U_0A=B\).  Since
\(\operatorname{rank}B\leq2\), the orthogonal complement of
\(\operatorname{im}B\) has positive complex dimension.  Choose a unitary
map \(V\) which fixes \(\operatorname{im}B\) pointwise and whose determinant
is \((\det U_0)^{-1}\).  Then
\[
 U=VU_0\in\mathrm{SU}(3),
 \qquad UA=VB=B.
\]
Hence equal Gram matrices are exactly equal \(\mathrm{SU}(3)\)-orbits.

The action is proper because \(\mathrm{SU}(3)\) is compact; hence the orbit
space is Hausdorff.  The continuous map
\(H\mapsto[s(H)]\) is inverse to the continuous induced map
\(\overline\Gamma\), so \eqref{eq:hrgc-orbit-homeomorphism} is a
homeomorphism.  Finally,
\begin{equation}
 \lVert A\rVert_F^2=\operatorname{tr}(A^\dagger A)
 =\operatorname{tr}\Gamma(A).
 \label{eq:hrgc-properness-identity}
\end{equation}
If \(K\) is compact in the positive-semidefinite cone, then
\(\Gamma^{-1}(K)\) is closed and bounded by
\eqref{eq:hrgc-properness-identity}, hence compact.  Thus \(\Gamma\) is
proper.  For compact \(K\),
\(\overline\Gamma^{-1}(K)\) is the continuous image in the orbit space of
the compact set \(\Gamma^{-1}(K)\), so \(\overline\Gamma\) is proper as
well.
\end{proof}

\begin{proposition}[Stabilizers and orbit-type strata]
\label{prop:hrgc-stabilizers}
For \(A\in\mathscr V_{(r)}\),
\begin{equation}
 \operatorname{Stab}_{\mathrm{SU}(3)}(A)\cong\mathrm{SU}(3-r).
 \label{eq:hrgc-stabilizer-formula}
\end{equation}
The exact dimensions are
\begin{equation}
 \begin{array}{c|c|c|c|c}
 r&\dim_{\mathbb R}\mathscr V_{(r)}&\operatorname{Stab}(A)&
 \dim_{\mathbb R}(\mathrm{SU}(3)\cdot A)&
 \dim_{\mathbb R}(\mathscr V_{(r)}/\mathrm{SU}(3))\\ \hline
 2&12&\{1\}&8&4\\
 1&8&\mathrm{SU}(2)&5&3\\
 0&0&\mathrm{SU}(3)&0&0.
 \end{array}
 \label{eq:hrgc-stratum-table}
\end{equation}
Consequently the quotient has topological dimension four, but the action is
free only on the rank-two stratum.
\end{proposition}

\begin{proof}
The equality \(UA=A\) is equivalent to the assertion that \(U\) fixes
\(\operatorname{im}A\) pointwise.  Restriction to its orthogonal complement
identifies the stabilizer with \(\mathrm{SU}(3-r)\), proving
\eqref{eq:hrgc-stabilizer-formula}.  The rank-two locus is open in the
twelve-real-dimensional vector space.  The complex rank-one locus in the
space of \(3\times2\) matrices has complex dimension
\[
 1(3+2-1)=4,
\]
hence real dimension eight away from the origin.  Subtracting
\(8-\dim_{\mathbb R}\operatorname{Stab}(A)\) gives the displayed orbit and
quotient dimensions.  They agree with the target strata: positive-definite
Hermitian matrices form a four-dimensional open cone, while every nonzero
rank-one positive-semidefinite matrix has a unique expression
\(\lambda vv^\dagger\) with \(\lambda>0\) and \(v\in S^3\), modulo
\(v\sim e^{i\theta}v\).  Hence that target stratum is
\[
 (0,\infty)\times\mathbb{CP}^1
 \cong(0,\infty)\times S^2,
\]
of real dimension three.  The rank-zero stratum is the origin.
\end{proof}

\subsection{The free locus, its explicit trivialization, and orientation}

\begin{theorem}[Determinant-one completion]
\label{thm:hrgc-free-trivialization}
There is an \(\mathrm{SU}(3)\)-equivariant diffeomorphism
\begin{equation}
 \mathscr V_{(2)}
 \xrightarrow{\ \cong\ }
 \mathrm{SU}(3)\times\operatorname{Herm}_2(\mathbb C)_{>0}
 \label{eq:hrgc-free-product}
\end{equation}
under which \(\Gamma\) is projection to the second factor and the group acts
by left multiplication on the first.  Thus the principal bundle over the
free quotient has a global section and is trivial.
\end{theorem}

\begin{proof}
For \(A\in\mathscr V_{(2)}\), set
\[
 H=A^\dagger A>0,
 \qquad Q=AH^{-1/2}=(q_1\ q_2).
\]
Then \(Q^\dagger Q=I_2\), so \((q_1,q_2)\) is an ordered orthonormal
two-frame in \(\mathbb C^3\).  With \(\varepsilon_{123}=1\), define
\begin{equation}
 (q_3)_i=\sum_{j,k=1}^3\varepsilon_{ijk}
 \overline{(q_1)_j}\,\overline{(q_2)_k}.
 \label{eq:hrgc-cross-product}
\end{equation}
Antisymmetry gives
\(q_1^\dagger q_3=q_2^\dagger q_3=0\).  The Gram--Pl\"ucker identity, proved
independently in \cref{prop:hrgc-pluecker}, gives \(\lVert q_3\rVert=1\).
Finally, contraction of the two alternating tensors gives
\begin{align*}
 \det[q_1\ q_2\ q_3]
 &=\sum_{i,j,k}\varepsilon_{ijk}(q_1)_i(q_2)_j(q_3)_k\\
 &=\lVert q_1\rVert^2\lVert q_2\rVert^2
   -|q_1^\dagger q_2|^2=1.
\end{align*}
Thus \(U_A=[q_1\ q_2\ q_3]\in\mathrm{SU}(3)\), and
\[
 A=U_A s(H),
\]
where \(s\) is \eqref{eq:hrgc-square-root-section}.  The positive square
root, inverse square root, and the polynomial-conjugate formula
\eqref{eq:hrgc-cross-product} are smooth on the positive-definite locus.
Consequently
\[
 A\longmapsto(U_A,A^\dagger A),
 \qquad
 (U,H)\longmapsto Us(H)
\]
are smooth inverse maps.  Their equivariance is immediate.
\end{proof}

The order of factors in \eqref{eq:hrgc-free-product} does not by itself fix
an orientation.  We now fix and calculate it.  Give
\(\mathfrak{su}(3)\) the ordered real basis
\begin{equation}
 \begin{gathered}
 X_{12},Y_{12},X_{13},Y_{13},X_{23},Y_{23},D_1,D_2,\\
 X_{jk}=E_{jk}-E_{kj},\qquad
 Y_{jk}=i(E_{jk}+E_{kj}),\\
 D_1=i\operatorname{diag}(1,-1,0),\qquad
 D_2=i\operatorname{diag}(0,1,-1).
 \end{gathered}
 \label{eq:hrgc-su3-oriented-basis}
\end{equation}
At
\[
 A_0=\begin{pmatrix}1&0\\0&1\\0&0\end{pmatrix},
\]
use the oriented Hermitian basis
\begin{equation}
 K_t=I_2,\quad
 K_x=\begin{pmatrix}0&1\\1&0\end{pmatrix},\quad
 K_y=\begin{pmatrix}0&-i\\i&0\end{pmatrix},\quad
 K_z=\begin{pmatrix}1&0\\0&-1\end{pmatrix}
 \label{eq:hrgc-hermitian-oriented-basis}
\end{equation}
and horizontal lifts \(B_\alpha=\tfrac12A_0K_\alpha\).  Directly,
\(d\Gamma_{A_0}(B_\alpha)=K_\alpha\).  Relative to
\eqref{eq:hrgc-v-orientation}, the twelve columns
\[
 (X_{12}A_0,Y_{12}A_0,X_{13}A_0,Y_{13}A_0,
 X_{23}A_0,Y_{23}A_0,D_1A_0,D_2A_0,B_t,B_x,B_y,B_z)
\]
form the exact real coordinate matrix
\begingroup
\scriptsize
\setlength{\arraycolsep}{2.4pt}
\begin{equation}
 M_{\mathrm{or}}:=\left(
 \begin{array}{rrrrrrrrrrrr}
0&0&0&0&0&0&0&0&\frac12&0&0&\frac12\\
0&0&0&0&0&0&1&0&0&0&0&0\\
1&0&0&0&0&0&0&0&0&\frac12&0&0\\
0&1&0&0&0&0&0&0&0&0&-\frac12&0\\
-1&0&0&0&0&0&0&0&0&\frac12&0&0\\
0&1&0&0&0&0&0&0&0&0&\frac12&0\\
0&0&0&0&0&0&0&0&\frac12&0&0&-\frac12\\
0&0&0&0&0&0&-1&1&0&0&0&0\\
0&0&-1&0&0&0&0&0&0&0&0&0\\
0&0&0&1&0&0&0&0&0&0&0&0\\
0&0&0&0&-1&0&0&0&0&0&0&0\\
0&0&0&0&0&1&0&0&0&0&0&0
 \end{array}
 \right).
 \label{eq:hrgc-orientation-matrix}
\end{equation}
\endgroup
Laplace expansion gives
\begin{equation}
 \det M_{\mathrm{or}}=\frac12>0.
 \label{eq:hrgc-orientation-determinant}
\end{equation}
Since both factors in \eqref{eq:hrgc-free-product} are connected, the sign is
constant on the free locus.  Therefore the quotient orientation determined
by
\[
 \operatorname{or}_{\mathscr V}
 =\operatorname{or}_{\mathrm{SU}(3)}\wedge
  \operatorname{or}_{\mathscr V_{(2)}/\mathrm{SU}(3)}
\]
agrees under \(\Gamma\) with
\(dt\wedge dx\wedge dy\wedge dz\).

\subsection{Lorentz determinant and Pl\"ucker norm}

For \(H=\Phi(t,x,y,z)\), direct expansion gives
\begin{equation}
 \operatorname{tr}H=2t,
 \qquad
 \det H=t^2-x^2-y^2-z^2,
 \label{eq:hrgc-lorentz-determinant}
\end{equation}
and the eigenvalues are
\(t\pm\sqrt{x^2+y^2+z^2}\).  Hence
\begin{equation}
 \begin{aligned}
 H>0&\Longleftrightarrow t>\sqrt{x^2+y^2+z^2},\\
 H\geq0&\Longleftrightarrow t\geq\sqrt{x^2+y^2+z^2},\\
 \operatorname{rank}H=1&\Longleftrightarrow
 t=\sqrt{x^2+y^2+z^2}>0.
 \end{aligned}
 \label{eq:hrgc-positive-cone-inequalities}
\end{equation}
Thus \(\det\) has signature \((1,3)\) when positive directions are listed
first, while \(-\det\) has signature \((3,1)\).  The quotient image is not
all of this Lorentz vector space: it is precisely the closed future causal
cone, and its free locus is the open future timelike cone.

\begin{proposition}[Gram--Pl\"ucker identity]
\label{prop:hrgc-pluecker}
For \(A=(a\ b)\),
\begin{equation}
 \det\Gamma(A)
 =\lVert a\rVert^2\lVert b\rVert^2-|\langle a,b\rangle|^2
 =\sum_{1\leq i<j\leq3}|a_ib_j-a_jb_i|^2
 =\lVert a\wedge b\rVert^2.
 \label{eq:hrgc-pluecker-identity}
\end{equation}
\end{proposition}

\begin{proof}
Expanding gives
\begin{align*}
 \lVert a\rVert^2\lVert b\rVert^2-|\langle a,b\rangle|^2
 &=\sum_{i,j}|a_i|^2|b_j|^2
   -\sum_{i,j}\overline{a_i}b_i a_j\overline{b_j}\\
 &=\sum_{i<j}
 \bigl(|a_i|^2|b_j|^2+|a_j|^2|b_i|^2
 -\overline{a_i}b_i a_j\overline{b_j}
 -\overline{a_j}b_j a_i\overline{b_i}\bigr)\\
 &=\sum_{i<j}|a_ib_j-a_jb_i|^2.
\end{align*}
The last expression is the squared Hermitian norm of \(a\wedge b\) in the
ordered basis
\((e_1\wedge e_2,e_1\wedge e_3,e_2\wedge e_3)\).
\end{proof}

The determinant is only one invariant and does not classify orbits.  For
example,
\[
 A_1=\begin{pmatrix}1&0\\0&1\\0&0\end{pmatrix},
 \qquad
 A_2=\begin{pmatrix}\sqrt2&0\\0&1/\sqrt2\\0&0\end{pmatrix}
\]
satisfy \(\det\Gamma(A_1)=\det\Gamma(A_2)=1\), but their Gram matrices are
\(I_2\) and \(\operatorname{diag}(2,1/2)\), so
\cref{thm:hrgc-gram-quotient} shows that they lie in distinct orbits.

\subsection{The four-coefficient form is additional input}

Let \(r=(r_1,r_2,r_3,r_4)^{\mathsf T}\in\mathbb R^4\), and define the
orthogonal Hadamard coordinates
\begin{equation}
 \begin{pmatrix}h\\x\\y\\z\end{pmatrix}
 =\frac12
 \begin{pmatrix}
 1&1&1&1\\
 1&1&-1&-1\\
 1&-1&1&-1\\
 1&-1&-1&1
 \end{pmatrix}
 \begin{pmatrix}r_1\\r_2\\r_3\\r_4\end{pmatrix}.
 \label{eq:hrgc-hadamard-transform}
\end{equation}
The proposed coefficient-space form
\begin{equation}
 Q_4(r)=\frac12\left(\sum_{i=1}^4r_i\right)^2-
 \sum_{i=1}^4r_i^2
 \label{eq:hrgc-q4-definition}
\end{equation}
satisfies the exact identity
\begin{equation}
 Q_4(r)=h^2-x^2-y^2-z^2.
 \label{eq:hrgc-q4-lorentz}
\end{equation}
Indeed, orthogonality gives
\(\sum r_i^2=h^2+x^2+y^2+z^2\) and
\(\sum r_i=2h\).

The same Hermitian matrix can be written in the alternative convention
\begin{equation}
 \Psi(h,x,y,z)=
 \begin{pmatrix}h+y&x+iz\\x-iz&h-y\end{pmatrix}.
 \label{eq:hrgc-pro-hermitian-convention}
\end{equation}
It agrees exactly with \(\Phi\) under
\[
 (t,x,y_{\Phi},z_{\Phi})=(h,x,-z,y).
\]
The transformation \((y,z)\mapsto(-z,y)\) has determinant \(+1\), so the
two conventions give the same matrix, Lorentz form, and orientation.

Nothing in the \(\mathrm{SU}(3)\) orbit calculation selects the coefficient
\(1/2\).  More generally,
\begin{equation}
 Q_{4,\lambda}(r)=\lambda\left(\sum_ir_i\right)^2-\sum_ir_i^2
 =(4\lambda-1)h^2-x^2-y^2-z^2.
 \label{eq:hrgc-q4-lambda-family}
\end{equation}
Every \(\lambda>1/4\) therefore has Lorentz signature \((1,3)\), after a
rescaling of \(h\).  The value \(\lambda=1/2\) is the unique member for which
the fixed orthogonal coordinates in \eqref{eq:hrgc-hadamard-transform} have
equal unit coefficients.  This is a normalization choice in coefficient
space, not a theorem forced by the Gram quotient.

Nor does a nonnegative coefficient tuple automatically determine a point of
the positive Gram cone.  The tuple \((1,0,0,0)\) has
\(Q_4=-1/2\).  A coefficient tuple enters the Gram quotient through
\eqref{eq:hrgc-pro-hermitian-convention} only when that Hermitian matrix is
positive semidefinite, equivalently when \(h\geq0\) and \(Q_4\geq0\).
The hypersurface \(z=0\) is a real-symmetric coefficient slice, but that
fact alone does not supply a coordinatewise nonnegative embedding of a
three-component mass domain.

\subsection{The compact representation quotient}

\begin{theorem}[The \(S^{12}\) representation quotient]
\label{thm:hrgc-s12-ball-quotient}
Let
\[
 S^{12}_{\mathrm{col}}=\mathscr V^+=\mathscr V\cup\{\infty\}
\]
be the one-point compactification, and extend the \(\mathrm{SU}(3)\) action
by fixing \(\infty\).  Then
\begin{equation}
 S^{12}_{\mathrm{col}}/\mathrm{SU}(3)
 \cong
 \bigl(\operatorname{Herm}_2(\mathbb C)_{\geq0}\bigr)^+
 \cong B^4.
 \label{eq:hrgc-s12-ball}
\end{equation}
The orbit-type strata are
\begin{equation}
 \begin{array}{c|c|c}
 \text{source stratum}&\text{quotient stratum}&\text{stabilizer}\\ \hline
 \operatorname{rank}A=2&\operatorname{int}B^4&\{1\}\\
 \operatorname{rank}A=1&\partial B^4\setminus\{p_0,p_\infty\}&
   \mathrm{SU}(2)\\
 A=0&\{p_0\}&\mathrm{SU}(3)\\
 A=\infty&\{p_\infty\}&\mathrm{SU}(3).
 \end{array}
 \label{eq:hrgc-compact-strata}
\end{equation}
The ball identification is topological on the full stratified quotient.  On
its smooth interior, it may be chosen to carry the quotient orientation to
the orientation fixed in \eqref{eq:hrgc-pauli-coordinates}.
\end{theorem}

\begin{proof}
As a real vector space, \(\mathscr V\cong\mathbb R^{12}\), so its one-point
compactification is \(S^{12}\).  On
\(\mathscr V\cong\mathbb C^3\oplus\mathbb C^3\), the action of
\(U\in\mathrm{SU}(3)\) has complex determinant \((\det U)^2=1\), hence
real determinant one; it preserves the standard orientation before and after
one-point compactification.

The properness proved from \eqref{eq:hrgc-properness-identity} implies that
\(\Gamma\) extends continuously to
\[
 \Gamma^+:\mathscr V^+\longrightarrow
 \bigl(\operatorname{Herm}_2(\mathbb C)_{\geq0}\bigr)^+,
 \qquad \Gamma^+(\infty)=\infty.
\]
It separates the finite orbits by \cref{thm:hrgc-gram-quotient}, and the
added point is a distinct fixed orbit.  Hence it induces a continuous
bijection
\[
 \mathscr V^+/\mathrm{SU}(3)
 \longrightarrow
 \bigl(\operatorname{Herm}_2(\mathbb C)_{\geq0}\bigr)^+.
\]
The source is compact Hausdorff because a compact group acts on the compact
Hausdorff space \(\mathscr V^+\).  The target is Hausdorff because the
positive-semidefinite cone is locally compact Hausdorff.  The continuous
bijection is therefore a homeomorphism.

By \eqref{eq:hrgc-positive-cone-inequalities}, the map
\[
 (t,\mathbf x)\longmapsto(t-\lVert\mathbf x\rVert,\mathbf x)
\]
is a homeomorphism from the closed future cone to the closed half-space
\(\mathbb R_{\geq0}\times\mathbb R^3\).  For \(u=(s,\mathbf x)\) in that half-space,
inverse stereographic projection is
\begin{equation}
 u\longmapsto
 \frac{1}{1+\lVert u\rVert^2}
 \bigl(2s,2\mathbf x,\lVert u\rVert^2-1\bigr)
 \in S^4\subset\mathbb R^5.
 \label{eq:hrgc-halfspace-stereographic}
\end{equation}
Its image is the closed hemisphere with first coordinate nonnegative, minus
the north pole.  The point at infinity supplies that pole.  A closed
hemisphere is a closed four-ball, proving \eqref{eq:hrgc-s12-ball}.

Positive-definite matrices correspond to the open half-space and hence the
ball interior.  Nonzero rank-one matrices correspond to
\(\{0\}\times(\mathbb R^3\setminus\{0\})\).  The zero matrix and the point at
infinity give two distinct endpoints in the compactification of that
three-dimensional boundary.  Combining this with
\cref{prop:hrgc-stabilizers} proves \eqref{eq:hrgc-compact-strata}.  The
orientation statement on the connected interior is the calculation following
\eqref{eq:hrgc-orientation-determinant}.
\end{proof}

The notation \(S^{12}_{\mathrm{col}}\) records the construction.  There is a
canonical pointed homeomorphism of representation spheres
\[
 (\mathbb C^3)^+\wedge(\mathbb C^3)^+
 \cong(\mathbb C^3\oplus\mathbb C^3)^+
 =S^{12}_{\mathrm{col}}.
\]
If a different based six-sphere \(X\) is used, a based homeomorphism
\(X\cong(\mathbb C^3)^+\) transports the action to \(X\wedge X\).  The
transported action depends on that choice.  Without an independently proved
equivariance statement, it is not induced by a torus fibration, cusp
monodromy, period matrix, or fibrewise flux on \(X\).

\subsection{Exact gauge-theoretic boundary}

\begin{theorem}[No native nonzero \(S^4\) sector from the Gram quotient]
\label{thm:hrgc-no-native-s4-sector}
The construction above supplies no native principal \(\mathrm{SU}(3)\)
bundle of nonzero second Chern number on \(S^4\).  More precisely:
\begin{enumerate}
\item the orbit projection over the rank-two locus is the product bundle in
\eqref{eq:hrgc-free-product}, so its second Chern class is zero;
\item on the rank-one boundary the orbit projection is not a principal
\(\mathrm{SU}(3)\)-bundle, because its stabilizer is \(\mathrm{SU}(2)\);
\item the full compact quotient is \(B^4\), not \(S^4\); and
\item replacing \(B^4\) by the one-point compactification of its interior is
a different construction whose clutching class is additional data.
\end{enumerate}
\end{theorem}

\begin{proof}
The first two assertions are
\cref{thm:hrgc-free-trivialization,prop:hrgc-stabilizers}; the third is
\cref{thm:hrgc-s12-ball-quotient}.  For the fourth, the matrix logarithm is a
diffeomorphism
\[
 \log:\operatorname{Herm}_2(\mathbb C)_{>0}
 \xrightarrow{\ \cong\ }
 \operatorname{Herm}_2(\mathbb C)\cong\mathbb R^4.
\]
Its one-point compactification is therefore \(S^4\), but this operation
collapses the entire finite rank-one boundary of the full Gram quotient into
the new point at infinity.  Principal \(\mathrm{SU}(3)\)-bundles on this
new \(S^4\) are obtained by clutching two trivial bundles over hemispheres
with a map
\[
 \kappa:S^3\longrightarrow\mathrm{SU}(3).
\]
The fibration
\(\mathrm{SU}(2)\to\mathrm{SU}(3)\to S^5\), together with
\(\pi_4(S^5)=\pi_3(S^5)=0\), shows that
\[
 \pi_3(\mathrm{SU}(3))\cong\pi_3(\mathrm{SU}(2))
 =\pi_3(S^3)\cong\mathbb Z.
\]
Choose the generator so that its associated bundle has second Chern number
\(+1\).  A clutching class \(k\in\mathbb Z\) then has second Chern number
\(k\).  Extending the explicit product framing in
\cref{thm:hrgc-free-trivialization} by the same framing gives the constant
clutching map and \(k=0\).  Neither \(A^\dagger A\) nor
\(\lVert a\wedge b\rVert^2\) selects a nonconstant \(\kappa\).  Thus
\(k\ne0\) is valid additional gauge data, but it is not produced by the Gram
construction.
\end{proof}

There are two further cohomological checks.  If a candidate total space is
diffeomorphic to \(S^6\), then
\[
 H^2(S^6;\mathbb Z)=H^4(S^6;\mathbb Z)=0.
\]
Consequently no nonzero fibrewise degree-two class can be the restriction of
a degree-two class on that total space, and no \(c_2\) on the total space can
restrict to a nonzero degree-four class on a fibre.  This does not forbid
independently chosen higher characteristic data; it proves exactly that the
degree-two and degree-four classes in question do not extend.

Likewise, for every continuous map \(f:S^{12}\to S^4\) and generator
\(u_4\in H^4(S^4;\mathbb Z)\),
\[
 f^*u_4=0
\]
because \(H^4(S^{12};\mathbb Z)=0\).  Therefore the pullback along any such
map of an \(\mathrm{SU}(3)\) bundle with nonzero \(c_2=ku_4\) has
\(c_2=0\).  This statement does not assert that every intrinsic
\(\mathrm{SU}(3)\) bundle on \(S^{12}\) is trivial; it says that an
\(S^4\) second-Chern charge is not inherited by pullback to the
representation sphere.

\subsection{The log-Euclidean metric and Yang--Mills theory are extra choices}

The logarithm identifies the open positive cone with a real vector space.
One may therefore declare
\begin{equation}
 X=\log H,
 \qquad
 g_{\log}=\frac12\operatorname{tr}(dX\,dX),
 \label{eq:hrgc-log-metric}
\end{equation}
which is a complete flat Euclidean metric in the \(X\)-coordinates.  This
metric is not induced by the Frobenius metric on \(\mathscr V\).  To see the
difference without discarding any scale factor, set
\[
 H(s)=sI_2,
 \qquad
 A(s)=\sqrt{s}\begin{pmatrix}1&0\\0&1\\0&0\end{pmatrix},
 \qquad s>0.
\]
For \(\xi\in\mathfrak{su}(3)\),
\[
 \operatorname{Re}\operatorname{tr}
 \bigl(\dot A^\dagger\xi A\bigr)
 =\frac{\dot s}{2}
 \operatorname{Re}(\xi_{11}+\xi_{22})=0,
\]
so \(\dot A\) is horizontal for the Frobenius quotient.  Hence
\begin{equation}
 \lVert\dot H\rVert_{\mathrm{quot}}^2
 =\lVert\dot A\rVert_F^2=\frac{\dot s^2}{2s}.
 \label{eq:hrgc-frobenius-ray}
\end{equation}
By contrast, \(X(s)=(\log s)I_2\), and
\begin{equation}
 \lVert\dot H\rVert_{\log}^2
 =\frac12\operatorname{tr}(\dot X^2)
 =\frac{\dot s^2}{s^2}.
 \label{eq:hrgc-log-ray}
\end{equation}
The ratio of \eqref{eq:hrgc-frobenius-ray} and
\eqref{eq:hrgc-log-ray} depends on \(s\), so no global constant rescaling
identifies the metrics.  The Lorentz form
\eqref{eq:hrgc-lorentz-determinant} is a third object and is indefinite.

After choosing \eqref{eq:hrgc-log-metric}, one may define classical
Yang--Mills theory on the trivial product bundle supplied by
\cref{thm:hrgc-free-trivialization}.  The additional choice of the zero
connection has \(F_0=0\), so it solves the nonlinear classical equation
\(d_0^*F_0=0\); its Coulomb Hessian is the transverse free Laplacian.  These
are the standard zero-charge vacuum statements on a chosen Euclidean
\(\mathbb R^4\).  The Gram morphism itself supplies no nonzero clutching
class, no nonflat connection, no interacting Hamiltonian, no quantum
Hilbert-space identification, and no mass-gap estimate.  Thus the exact
logical implication is
\[
 \boxed{
 \begin{array}{c}
 \text{four-dimensional Gram/Lorentz quotient: proved;}\\
 S^{12}_{\mathrm{col}}/\mathrm{SU}(3)\cong B^4: \text{ proved;}\\
 \text{nonzero native }\mathrm{SU}(3)\text{ sector on }S^4:
 \text{ not produced;}\\
 \text{interacting or quantum Yang--Mills inheritance: not produced.}
 \end{array}}
\]

\section{The three-column \texorpdfstring{\(\mathrm{SU}(3)\)}{SU(3)}
Gram cone and its rank-loss \texorpdfstring{\(S^8\)}{S8}}
\label{sec:hr3c-main}

This section proves the topology forced by the left actions of
\(\mathrm U(3)\) and \(\mathrm{SU}(3)\) on three complex columns.  The
calculation produces a canonical topological \(S^8\), but it does not
equip that sphere with a complex structure, a fibration, a gauge bundle,
or a physical interpretation.

\subsection{Conventions}
\label{subsec:hr3c-conventions}

Put
\begin{equation}
 \mathscr V_3=\operatorname{Mat}_{3\times3}(\mathbb C),
 \qquad
 \lVert A\rVert_F^2=\operatorname{tr}(A^\dagger A).
 \label{eq:hr3c-source}
\end{equation}
This is a real vector space of dimension \(18\).  Both
\(\mathrm U(3)\) and \(\mathrm{SU}(3)\) act on the left:
\begin{equation}
 U\cdot A=UA.
 \label{eq:hr3c-left-action}
\end{equation}
Let
\begin{equation}
 \mathcal P_3=
 \{H\in\operatorname{Herm}_3(\mathbb C):H\succeq0\},
 \qquad
 \mathcal K_3=
 \{H\in\mathcal P_3:\operatorname{tr}H=1\}.
 \label{eq:hr3c-gram-cone-slice}
\end{equation}
Write \(S^{17}_F\subset\mathscr V_3\) for the Frobenius unit sphere and
\(\mathscr V_3^+=\mathscr V_3\cup\{\infty\}\cong S^{18}\) for the
one-point compactification, with \(\infty\) fixed by each group.

\subsection{The \texorpdfstring{\(\mathrm U(3)\)}{U(3)} orbit space}
\label{subsec:hr3c-u3}

\begin{theorem}[Exact Gram classification]
\label{thm:hr3c-u3-gram-classification}
The map
\begin{equation}
 \Gamma_U:\mathscr V_3\longrightarrow\mathcal P_3,
 \qquad
 \Gamma_U(A)=A^\dagger A
 \label{eq:hr3c-u3-gram-map}
\end{equation}
is \(\mathrm U(3)\)-invariant, is surjective, and separates
\(\mathrm U(3)\)-orbits.  It induces a homeomorphism
\begin{equation}
 \overline\Gamma_U:
 \mathscr V_3/\mathrm U(3)\xrightarrow{\ \sim\ }\mathcal P_3.
 \label{eq:hr3c-u3-quotient-homeomorphism}
\end{equation}
\end{theorem}

\begin{proof}
Invariance is the identity
\[
 (UA)^\dagger(UA)=A^\dagger U^\dagger UA=A^\dagger A.
\]
For \(H\succeq0\), its positive-semidefinite square root satisfies
\[
 (H^{1/2})^\dagger H^{1/2}=H,
\]
so the map is surjective.

Suppose \(A^\dagger A=B^\dagger B\).  Define
\[
 T:\operatorname{im}A\longrightarrow\operatorname{im}B,
 \qquad
 T(Ax)=Bx.
\]
If \(Ax=0\), then
\[
 \lVert Bx\rVert^2
 =x^\dagger B^\dagger Bx
 =x^\dagger A^\dagger Ax
 =\lVert Ax\rVert^2
 =0,
\]
so \(T\) is well-defined.  For \(x,y\in\mathbb C^3\),
\[
 \langle T(Ax),T(Ay)\rangle
 =x^\dagger B^\dagger By
 =x^\dagger A^\dagger Ay
 =\langle Ax,Ay\rangle.
\]
Thus \(T\) is a complex-linear isometry.  Extend it to a unitary
\(U\in\mathrm U(3)\).  Then \(UA=B\), so the Gram matrix separates
orbits.

Finally,
\[
 \lVert A\rVert_F^2=\operatorname{tr}(A^\dagger A).
\]
Hence \(\Gamma_U\) is proper.  The induced map
\(\overline\Gamma_U\) is a proper continuous bijection between locally
compact Hausdorff spaces and is therefore a homeomorphism.
\end{proof}

\begin{lemma}[Explicit ball chart on the trace-one Gram body]
\label{lem:hr3c-density-ball}
Let
\begin{equation}
 E=\{X\in\operatorname{Herm}_3(\mathbb C):\operatorname{tr}X=0\}
 \label{eq:hr3c-traceless-space}
\end{equation}
with Frobenius inner product.  Then
\(\dim_{\mathbb R}E=8\), and there is an explicit homeomorphism
\begin{equation}
 \chi:\mathcal K_3\xrightarrow{\ \sim\ }B(E)\cong B^8
 \label{eq:hr3c-chi-homeomorphism}
\end{equation}
which sends the singular locus in \(\mathcal K_3\) onto \(S(E)\).
\end{lemma}

\begin{proof}
Put \(c=I_3/3\).  If \(u\in S(E)\), then \(u\ne0\) and
\(\operatorname{tr}u=0\), so
\(\lambda_{\min}(u)<0\).  Define
\begin{equation}
 R(u)=-\frac{1}{3\lambda_{\min}(u)}>0.
 \label{eq:hr3c-radial-length}
\end{equation}
The eigenvalues of \(c+s u\) are
\(1/3+s\lambda_j(u)\).  Therefore, for \(s\geq0\),
\begin{equation}
 c+s u\succeq0
 \quad\Longleftrightarrow\quad
 0\leq s\leq R(u),
 \label{eq:hr3c-positive-radial-interval}
\end{equation}
and \(c+R(u)u\) is singular.

Every \(H\in\mathcal K_3\setminus\{c\}\) has the unique expression
\[
 H=c+s u,
 \qquad
 s=\lVert H-c\rVert_F,
 \qquad
 u=\frac{H-c}{\lVert H-c\rVert_F}.
\]
Define
\begin{equation}
 \chi(c)=0,
 \qquad
 \chi(c+s u)=\frac{s}{R(u)}u,
 \label{eq:hr3c-chi-formula}
\end{equation}
and
\begin{equation}
 \chi^{-1}(0)=c,
 \qquad
 \chi^{-1}(r u)=c+rR(u)u
 \quad(0<r\leq1).
 \label{eq:hr3c-chi-inverse}
\end{equation}
The minimum-eigenvalue function is continuous.  Since it is strictly
negative on the compact sphere \(S(E)\), \(R\) is continuous, finite,
positive, and bounded above and below.  Equations
\eqref{eq:hr3c-chi-formula} and \eqref{eq:hr3c-chi-inverse} are therefore
continuous also at the origin and are inverse homeomorphisms.
Equation \eqref{eq:hr3c-positive-radial-interval} shows that
\[
 \chi\bigl(\{H\in\mathcal K_3:\operatorname{rank}H\leq2\}\bigr)
 =S(E).
\]
\end{proof}

\begin{corollary}[Unit-sphere \(\mathrm U(3)\) quotient]
\label{cor:hr3c-u3-unit-quotient}
There is a homeomorphism of pairs
\begin{equation}
 \left(
 S^{17}_F/\mathrm U(3),
 \{[A]:\operatorname{rank}A\leq2\}
 \right)
 \cong
 (B^8,S^7).
 \label{eq:hr3c-u3-unit-pair}
\end{equation}
\end{corollary}

\begin{proof}
On \(S^{17}_F\),
\[
 \operatorname{tr}(A^\dagger A)=\lVert A\rVert_F^2=1.
\]
Apply Theorem~\ref{thm:hr3c-u3-gram-classification} and
Lemma~\ref{lem:hr3c-density-ball}.
\end{proof}

\begin{theorem}[Full \(\mathrm U(3)\) cone and compactification]
\label{thm:hr3c-u3-compactification}
There are homeomorphisms
\begin{equation}
 \mathscr V_3/\mathrm U(3)
 \cong\mathcal P_3
 \cong\mathbb R^8\times\mathbb R_{\geq0}
 \label{eq:hr3c-u3-half-space}
\end{equation}
and
\begin{equation}
 \mathscr V_3^+/\mathrm U(3)
 \cong
 \mathcal P_3^+
 \cong B^9.
 \label{eq:hr3c-u3-ball-nine}
\end{equation}
Under the second homeomorphism,
\begin{equation}
 \partial B^9
 =
 \{[A]:\operatorname{rank}A\leq2\}
 \cup\{[\infty]\}
 \cong S^8.
 \label{eq:hr3c-u3-boundary-s8}
\end{equation}
\end{theorem}

\begin{proof}
For \(H\ne0\), write
\[
 H=r\rho,
 \qquad
 r=\operatorname{tr}H>0,
 \qquad
 \rho=\frac{H}{\operatorname{tr}H}\in\mathcal K_3.
\]
Thus \(\mathcal P_3\) is the infinite cone on
\(\mathcal K_3\cong B^8\).

We give the compactification map explicitly.  Put
\[
 t=\frac{r-1}{r+1}\in(-1,1)
\]
and define
\begin{equation}
 \mathcal D^9=
 \{(y,t)\in E\oplus\mathbb R:
 \lVert y\rVert_F+|t|\leq1\}.
 \label{eq:hr3c-double-cone}
\end{equation}
Define
\begin{equation}
 \begin{split}
 \Xi_U(r\rho)
  &=\bigl((1-|t|)\chi(\rho),t\bigr),\\
 \Xi_U(0)&=(0,-1),\\
 \Xi_U(\infty)&=(0,1).
 \end{split}
 \label{eq:hr3c-xi-u}
\end{equation}
For \(|t|<1\), the inverse is
\begin{equation}
 x=\frac{y}{1-|t|},
 \qquad
 \rho=\chi^{-1}(x),
 \qquad
 r=\frac{1+t}{1-t},
 \qquad
 H=r\rho.
 \label{eq:hr3c-xi-u-inverse}
\end{equation}
As \(r\downarrow0\) or \(r\to\infty\), the factor \(1-|t|\) tends
uniformly to zero.  A subset of \(\mathcal P_3\) is unbounded if and only
if its trace is unbounded.  Thus \(\Xi_U\) is a homeomorphism
\(\mathcal P_3^+\to\mathcal D^9\).

For \(p=(y,t)\ne0\), put
\[
 \nu(p)=\sqrt{\lVert y\rVert_F^2+t^2},
 \qquad
 \gamma(p)=\lVert y\rVert_F+|t|
\]
and define
\begin{equation}
 \Theta(p)=\frac{\gamma(p)}{\nu(p)}p,
 \qquad
 \Theta(0)=0.
 \label{eq:hr3c-double-cone-ball-map}
\end{equation}
Then \(\lVert\Theta(p)\rVert_2=\gamma(p)\leq1\).  For nonzero
\(q\in B(E\oplus\mathbb R)\), the inverse is
\begin{equation}
 \Theta^{-1}(q)=
 \frac{\lVert q\rVert_2}
 {\lVert q_E\rVert_F+|q_{\mathbb R}|}\,q.
 \label{eq:hr3c-double-cone-ball-inverse}
\end{equation}
Hence \(\Theta:\mathcal D^9\to B^9\) is a homeomorphism and takes
\(\partial\mathcal D^9\) onto \(S^8\).

For finite \(H=r\rho\), Equation~\eqref{eq:hr3c-xi-u} lies on
\(\partial\mathcal D^9\) exactly when
\(\lVert\chi(\rho)\rVert_F=1\), equivalently when \(H\) is singular.
The points \(0\) and \(\infty\) are the two boundary tips.  This proves
\eqref{eq:hr3c-u3-boundary-s8}.

The canonical map
\[
 \mathscr V_3^+/\mathrm U(3)
 \longrightarrow
 (\mathscr V_3/\mathrm U(3))^+
\]
is a continuous bijection: the finite orbits agree and the added orbit is
\(\{\infty\}\).  Its source is compact and its target is Hausdorff, hence
it is a homeomorphism.  This proves
\eqref{eq:hr3c-u3-ball-nine}.

Finally, \(\Theta\circ\Xi_U\) sends \(\infty\) to
\(N=(0,\ldots,0,1)\in\partial B^9\).  The Cayley map
\begin{equation}
 (x',x_9)\longmapsto
 \left(
 \frac{2x'}{\lVert x'\rVert^2+(1-x_9)^2},
 \frac{1-\lVert x\rVert^2}
 {\lVert x'\rVert^2+(1-x_9)^2}
 \right)
 \label{eq:hr3c-cayley-map}
\end{equation}
is a homeomorphism
\[
 B^9\setminus\{N\}\longrightarrow
 \mathbb R^8\times\mathbb R_{\geq0}.
\]
Its inverse is
\begin{equation}
 (v,s)\longmapsto
 \left(
 \frac{2v}{\lVert v\rVert^2+(s+1)^2},
 \frac{\lVert v\rVert^2+s^2-1}
 {\lVert v\rVert^2+(s+1)^2}
 \right).
 \label{eq:hr3c-cayley-inverse}
\end{equation}
This proves \eqref{eq:hr3c-u3-half-space}.
\end{proof}

\subsection{The \texorpdfstring{\(\mathrm{SU}(3)\)}{SU(3)} orbit space}
\label{subsec:hr3c-su3}

\begin{theorem}[Gram matrix plus alternating determinant]
\label{thm:hr3c-su3-invariants}
Define
\begin{equation}
 \Gamma_{SU}:\mathscr V_3\longrightarrow
 \operatorname{Herm}_3(\mathbb C)\times\mathbb C,
 \qquad
 \Gamma_{SU}(A)=(A^\dagger A,\det A).
 \label{eq:hr3c-su-invariant-map}
\end{equation}
Its image is exactly
\begin{equation}
 \mathcal Y_3=
 \{(H,z)\in\mathcal P_3\times\mathbb C:
 |z|^2=\det H\}.
 \label{eq:hr3c-su-invariant-image}
\end{equation}
It separates \(\mathrm{SU}(3)\)-orbits and induces a homeomorphism
\begin{equation}
 \mathscr V_3/\mathrm{SU}(3)\xrightarrow{\ \sim\ }\mathcal Y_3.
 \label{eq:hr3c-su-quotient}
\end{equation}
\end{theorem}

\begin{proof}
The map is \(\mathrm{SU}(3)\)-invariant, and
\[
 \det(A^\dagger A)
 =\overline{\det A}\det A
 =|\det A|^2.
\]
Thus its image is contained in \(\mathcal Y_3\).

Conversely, take \((H,z)\in\mathcal Y_3\).  If \(H\) is singular, then
\(z=0\), and \(A=H^{1/2}\) realizes the pair.  If \(H\succ0\), put
\[
 \omega=\frac{z}{\sqrt{\det H}}\in S^1,
 \qquad
 Q_\omega=\operatorname{diag}(\omega,1,1)\in\mathrm U(3),
 \qquad
 A=Q_\omega H^{1/2}.
\]
Then \(A^\dagger A=H\) and \(\det A=z\).  This proves surjectivity.

Suppose
\[
 A^\dagger A=B^\dagger B,
 \qquad
 \det A=\det B.
\]
Theorem~\ref{thm:hr3c-u3-gram-classification} gives
\(U_0\in\mathrm U(3)\) with \(U_0A=B\).  If \(A\) is invertible, then
\[
 \det U_0=\frac{\det B}{\det A}=1,
\]
so \(U_0\in\mathrm{SU}(3)\).

If \(\operatorname{rank}A\leq2\), choose a unitary \(V\) which is the
identity on \(\operatorname{im}B\) and has determinant
\((\det U_0)^{-1}\) on \((\operatorname{im}B)^\perp\).  This is possible
because the orthogonal complement has positive complex dimension.  Then
\[
 VB=B,
 \qquad
 \det(VU_0)=1,
 \qquad
 (VU_0)A=B.
\]
Thus the invariant separates \(\mathrm{SU}(3)\)-orbits.

Since \(\operatorname{tr}H=\lVert A\rVert_F^2\), the invariant map is
proper.  The induced continuous bijection onto the closed Hausdorff set
\(\mathcal Y_3\) is therefore a homeomorphism.
\end{proof}

\begin{theorem}[The unit-sphere \(\mathrm{SU}(3)\) quotient]
\label{thm:hr3c-su-unit-sphere}
Let
\begin{equation}
 \mathcal Y_{3,1}=
 \{(H,z)\in\mathcal Y_3:\operatorname{tr}H=1\}.
 \label{eq:hr3c-su-unit-image}
\end{equation}
There is an explicit homeomorphism
\begin{equation}
 S^{17}_F/\mathrm{SU}(3)
 \cong
 \mathcal Y_{3,1}
 \xrightarrow{\ \sim\ }
 S(E\oplus\mathbb C)
 \cong S^9.
 \label{eq:hr3c-su-unit-sphere-result}
\end{equation}
\end{theorem}

\begin{proof}
The first homeomorphism follows from
Theorem~\ref{thm:hr3c-su3-invariants}.  To construct the second, let
\[
 x=\chi(H)\in B(E),
 \qquad
 r=\lVert x\rVert_F.
\]
If \(r<1\), then \(H\succ0\) and \(z\ne0\).  Define
\begin{equation}
 \Omega(H,z)=
 \left(
 x,\sqrt{1-r^2}\frac{z}{|z|}
 \right)
 \quad(r<1),
 \qquad
 \Omega(H,0)=(x,0)
 \quad(r=1).
 \label{eq:hr3c-omega}
\end{equation}
Its norm is exactly one:
\[
 \lVert x\rVert_F^2+
 \left|\sqrt{1-r^2}\frac{z}{|z|}\right|^2
 =r^2+(1-r^2)=1.
\]
The map is continuous on the full-rank locus.  As \(H\) approaches the
singular locus, \(r\to1\), so its complex coordinate tends to zero
regardless of phase.  Hence it is continuous everywhere.

Given \((x,w)\in S(E\oplus\mathbb C)\), set
\(H=\chi^{-1}(x)\).  If \(\lVert x\rVert_F<1\), then \(w\ne0\), and the
unique inverse value is
\begin{equation}
 z=\sqrt{\det H}\frac{w}{|w|}.
 \label{eq:hr3c-omega-inverse-z}
\end{equation}
If \(\lVert x\rVert_F=1\), then \(w=0\), \(H\) is singular, and \(z=0\).
Thus \(\Omega\) is a continuous bijection from a compact space to a
Hausdorff space and hence is a homeomorphism.
\end{proof}

\begin{theorem}[The weighted cone and its compactification]
\label{thm:hr3c-su-full-cone}
There is a homeomorphism
\begin{equation}
 \Phi_{SU}:\mathcal Y_3\xrightarrow{\ \sim\ }
 E\oplus\mathbb C\cong\mathbb R^{10}
 \label{eq:hr3c-su-r10}
\end{equation}
and therefore
\begin{equation}
 \mathscr V_3^+/\mathrm{SU}(3)\cong S^{10}.
 \label{eq:hr3c-su-s10}
\end{equation}
\end{theorem}

\begin{proof}
For \((H,z)\ne(0,0)\), set
\begin{equation}
 s=\sqrt{\operatorname{tr}H},
 \qquad
 h=\frac{H}{s^2},
 \qquad
 \zeta=\frac{z}{s^3}.
 \label{eq:hr3c-weighted-normalization}
\end{equation}
The exponents are forced by scalar multiplication in the original source:
for every real \(\lambda>0\),
\begin{equation}
 (\lambda A)^\dagger(\lambda A)=\lambda^2A^\dagger A,
 \qquad
 \det(\lambda A)=\lambda^3\det A,
 \qquad
 s(\lambda A)=\lambda s(A).
 \label{eq:hr3c-scaling-degrees}
\end{equation}
Moreover,
\[
 \operatorname{tr}h=1,
 \qquad
 |\zeta|^2
 =\frac{|z|^2}{s^6}
 =\frac{\det H}{s^6}
 =\det h.
\]
Thus \((h,\zeta)\in\mathcal Y_{3,1}\).  Define
\begin{equation}
 \Phi_{SU}(0,0)=0,
 \qquad
 \Phi_{SU}(H,z)=s\,\Omega(h,\zeta)
 \quad(s>0).
 \label{eq:hr3c-su-radial-homeomorphism}
\end{equation}
Since \(\Omega(h,\zeta)\) has norm one,
\[
 \lVert\Phi_{SU}(H,z)\rVert=s.
\]
For nonzero \(v\in E\oplus\mathbb C\), the inverse is obtained from
\[
 s=\lVert v\rVert,
 \qquad
 (h,\zeta)=
 \Omega^{-1}\left(\frac{v}{\lVert v\rVert}\right),
 \qquad
 H=s^2h,
 \qquad
 z=s^3\zeta.
\]
The norm identity proves continuity at the origin in both directions, so
\(\Phi_{SU}\) is a homeomorphism.

As in the proof of Theorem~\ref{thm:hr3c-u3-compactification}, compactness
of \(\mathrm{SU}(3)\) gives
\[
 \mathscr V_3^+/\mathrm{SU}(3)
 \cong
 (\mathscr V_3/\mathrm{SU}(3))^+.
\]
The one-point compactification of \(\mathbb R^{10}\) is \(S^{10}\).
\end{proof}

\subsection{Rank strata, the standard \texorpdfstring{\(S^8\)}{S8}, and
determinant winding}
\label{subsec:hr3c-rank-complement}

\begin{proposition}[Stabilizers and dimensions]
\label{prop:hr3c-rank-strata}
Let
\[
 \mathscr V_{3,(r)}
 =\{A\in\mathscr V_3:\operatorname{rank}_{\mathbb C}A=r\}.
\]
For \(A\in\mathscr V_{3,(r)}\),
\begin{equation}
 \operatorname{Stab}_{\mathrm U(3)}(A)\cong\mathrm U(3-r),
 \qquad
 \operatorname{Stab}_{\mathrm{SU}(3)}(A)\cong\mathrm{SU}(3-r),
 \label{eq:hr3c-stabilizers}
\end{equation}
with the rank-three and \(\mathrm{SU}(1)\) stabilizers understood to be
trivial.  The quotient-stratum dimensions are
\begin{center}
\begin{tabular}{c|c|c|c|c|c}
\(r\)&
\(\dim_{\mathbb R}\mathscr V_{3,(r)}\)&
\(\dim U(3)\!\cdot\!A\)&
\(\dim \mathscr V_{3,(r)}/U(3)\)&
\(\dim SU(3)\!\cdot\!A\)&
\(\dim \mathscr V_{3,(r)}/SU(3)\)\\ \hline
3&18&9&9&8&10\\
2&16&8&8&8&8\\
1&10&5&5&5&5\\
0&0&0&0&0&0
\end{tabular}
\end{center}
On the unit sphere, the nonzero quotient-stratum dimensions are
\[
 \begin{array}{c|ccc}
 r&3&2&1\\ \hline
 U(3)&8&7&4\\
 SU(3)&9&7&4.
 \end{array}
\]
The trace-one rank-one locus is \(\mathbb CP^2\).
\end{proposition}

\begin{proof}
A left stabilizer fixes \(\operatorname{im}A\) pointwise.  Its unrestricted
unitary action on \((\operatorname{im}A)^\perp\) gives
\(\mathrm U(3-r)\).  Imposing determinant one gives
\(\mathrm{SU}(3-r)\).  The real dimension of the complex rank-\(r\)
matrix stratum is
\[
 2r(6-r).
\]
Subtracting the orbit dimensions gives the table; intersecting a nonzero
rank stratum with the unit sphere subtracts one more dimension.

A trace-one positive-semidefinite matrix has rank one exactly when it has
the form \(vv^\dagger\) with \(\lVert v\rVert=1\).  The phase of \(v\)
does not change \(vv^\dagger\), so this locus is
\[
 S^5/S^1\cong\mathbb CP^2.
\]
\end{proof}

\begin{theorem}[The exact rank-loss pair and complement]
\label{thm:hr3c-rank-loss-pair}
Let
\begin{equation}
 \mathcal Y_{3,\mathrm{sing}}
 =\{(H,0)\in\mathcal Y_3:\operatorname{rank}H\leq2\}.
 \label{eq:hr3c-singular-invariant-locus}
\end{equation}
The homeomorphism \(\Phi_{SU}\) gives homeomorphisms of pairs
\begin{equation}
 (\mathcal Y_3,\mathcal Y_{3,\mathrm{sing}})
 \cong
 (E\oplus\mathbb C,E\oplus\{0\})
 \label{eq:hr3c-affine-pair}
\end{equation}
and
\begin{equation}
 (\mathcal Y_3^+,\mathcal Y_{3,\mathrm{sing}}^+)
 \cong
 (S^{10},S^8_{\mathrm{lin}}).
 \label{eq:hr3c-compact-pair}
\end{equation}
Consequently the compactified rank-loss locus is a standard,
topologically locally flat codimension-two \(S^8\subset S^{10}\), and
\begin{equation}
 S^{10}\setminus S^8
 \cong
 E\times(\mathbb C\setminus\{0\})
 \cong
 \mathbb R^8\times(\mathbb C\setminus\{0\})
 \simeq S^1.
 \label{eq:hr3c-complement}
\end{equation}
The map
\begin{equation}
 \vartheta(H,z)=\frac{z}{|z|}
 =\frac{\det A}{|\det A|}
 \label{eq:hr3c-phase-map}
\end{equation}
is the angular deformation-retraction coordinate and represents a
generator of
\[
 H^1(S^{10}\setminus S^8;\mathbb Z)
 \cong
 \pi_1(S^{10}\setminus S^8)
 \cong\mathbb Z
\]
after choosing the standard positive orientation on \(S^1\).
\end{theorem}

\begin{proof}
In Equation~\eqref{eq:hr3c-omega}, \(H\) is singular exactly when
\(\lVert\chi(H)\rVert_F=1\), in which case the complex coordinate is zero.
Conversely, a zero complex coordinate on the unit \(S^9\) has its
\(E\)-coordinate on \(S(E)\), and its inverse Gram matrix is singular.
Radial extension through Equation~\eqref{eq:hr3c-su-radial-homeomorphism}
therefore gives exactly \eqref{eq:hr3c-affine-pair}.

The homeomorphism \(\Phi_{SU}\) is proper and extends to one-point
compactifications, proving \eqref{eq:hr3c-compact-pair}.  The second member
is the one-point compactification of the linear pair
\[
 (E\oplus\mathbb C,E\oplus\{0\}).
\]
It is the standard topological pair \((S^{10},S^8_{\mathrm{lin}})\);
therefore the \(S^8\) is topologically locally flat at every point.

This local-flatness statement is not a smooth-stratification assertion.
The unit rank-deficient \(S^7\) contains the rank-one orbit stratum
\(\mathbb CP^2\), with rank-two complement
\(S^7\setminus\mathbb CP^2\).  The homeomorphism \(\chi\) uses the minimum
eigenvalue and need not be smooth across eigenvalue-collision strata.

Removing the linear subspace from the affine pair gives
\eqref{eq:hr3c-complement}.  Under \(\Phi_{SU}\), the phase of its complex
coordinate is \(z/|z|\), so \(\vartheta\) is the usual angular projection
\[
 E\times(\mathbb C\setminus\{0\})\longrightarrow S^1.
\]
This is a deformation retraction.  The explicit quotient loop
\begin{equation}
 A_\theta=
 \frac1{\sqrt3}\operatorname{diag}(e^{i\theta},1,1),
 \qquad
 0\leq\theta\leq2\pi,
 \label{eq:hr3c-generator-loop}
\end{equation}
satisfies
\[
 A_\theta^\dagger A_\theta=\frac13I_3,
 \qquad
 \det A_\theta=\frac{e^{i\theta}}{3\sqrt3}.
\]
Thus \(\vartheta([A_\theta])=e^{i\theta}\) winds once and is a generator.
\end{proof}

\begin{corollary}[Residual determinant-phase circle]
\label{cor:hr3c-residual-circle}
The residual group
\[
 \mathrm U(3)/\mathrm{SU}(3)\cong S^1
\]
acts on \(E\oplus\mathbb C\) by
\begin{equation}
 e^{i\theta}\cdot(x,w)=(x,e^{i\theta}w).
 \label{eq:hr3c-residual-action}
\end{equation}
Its fixed set after compactification is the standard
\(S^8\subset S^{10}\), and
\begin{equation}
 S^{10}/S^1\cong B^9,
 \qquad
 \operatorname{Fix}(S^1)=S^8=\partial B^9.
 \label{eq:hr3c-residual-quotient}
\end{equation}
\end{corollary}

\begin{proof}
If \(U\in\mathrm U(3)\), left multiplication fixes \(H=A^\dagger A\) and
sends \(z=\det A\) to \((\det U)z\).  Every phase occurs as a determinant.
Equation~\eqref{eq:hr3c-omega} therefore converts the residual action into
\eqref{eq:hr3c-residual-action}.  On the unit sphere,
\[
 S(E\oplus\mathbb C)/S^1
 \cong
 \{(x,|w|):\lVert x\rVert_F^2+|w|^2=1,\ |w|\geq0\}
 \cong B(E).
\]
On the affine cone,
\[
 (E\oplus\mathbb C)/S^1
 \cong E\times\mathbb R_{\geq0}.
\]
One-point compactification gives \(B^9\), and the fixed set is
\((E\oplus\{0\})^+\cong S^8\), which maps to the boundary.
\end{proof}

\subsection{The exact rank-two/rank-three distinction}
\label{subsec:hr3c-two-ladders}

\begin{theorem}[Quadratic Lorentz cone versus cubic rank-three cone]
\label{thm:hr3c-two-ladders}
For
\(\mathscr V_2=\operatorname{Mat}_{3\times2}(\mathbb C)\), the left
\(\mathrm{SU}(3)\)- and \(\mathrm U(3)\)-orbit relations coincide, and
\begin{equation}
 \mathscr V_2/\mathrm{SU}(3)
 \cong
 \operatorname{Herm}_2(\mathbb C)^{\succeq0}.
 \label{eq:hr3c-two-column-quotient}
\end{equation}
In the coordinates
\begin{equation}
 H_2(t,x_1,x_2,x_3)=
 \begin{pmatrix}
  t+x_3&x_1-i x_2\\
  x_1+i x_2&t-x_3
 \end{pmatrix},
 \label{eq:hr3c-lorentz-coordinates}
\end{equation}
\begin{equation}
 \operatorname{tr}H_2=2t,
 \qquad
 \det H_2=t^2-x_1^2-x_2^2-x_3^2,
 \label{eq:hr3c-lorentz-determinant}
\end{equation}
and
\begin{equation}
 H_2\succeq0
 \quad\Longleftrightarrow\quad
 t\geq0,\qquad
 t^2\geq x_1^2+x_2^2+x_3^2.
 \label{eq:hr3c-lorentz-positivity}
\end{equation}
Thus the two-column quotient is exactly the closed future
\(3+1\)-dimensional Lorentz cone; its one-point compactification is
\(B^4\), with rank-one-plus-infinity boundary \(S^3\).

By contrast, the three-column \(\mathrm U(3)\) quotient is the
nine-real-dimensional, Jordan-rank-three cone
\(\operatorname{Herm}_3(\mathbb C)^{\succeq0}\).  Its determinant is cubic,
its trace-one section is \(B^8\), and its compactified
rank-deficient-plus-infinity boundary is \(S^8\).  It is not a Lorentz
cone.
\end{theorem}

\begin{proof}
Every \(3\times2\) matrix has rank at most two, so the determinant-adjustment
step in the proof of Theorem~\ref{thm:hr3c-su3-invariants} converts every
left \(\mathrm U(3)\) equivalence into a left \(\mathrm{SU}(3)\)
equivalence.  The Gram proof of
Theorem~\ref{thm:hr3c-u3-gram-classification}, with two columns, gives
\eqref{eq:hr3c-two-column-quotient}.

Direct expansion of \eqref{eq:hr3c-lorentz-coordinates} gives
\eqref{eq:hr3c-lorentz-determinant}.  A \(2\times2\) Hermitian matrix is
positive semidefinite exactly when its trace and determinant are
nonnegative, which gives \eqref{eq:hr3c-lorentz-positivity}.  Its trace-one
section is
\[
 t=\frac12,
 \qquad
 x_1^2+x_2^2+x_3^2\leq\frac14,
\]
a closed \(3\)-ball.  The compactification argument of
Theorem~\ref{thm:hr3c-u3-compactification}, with \(E\) of dimension three,
gives \(B^4\) with boundary \(S^3\).

For three columns, Theorem~\ref{thm:hr3c-u3-compactification} gives the
\(B^9/S^8\) statements.  The determinant of a \(3\times3\) Hermitian
matrix is homogeneous of degree three, so this rank-three cone is not the
rank-two quadratic Lorentz cone.
\end{proof}

\subsection{The alternating invariant and the precise nonclaim}
\label{subsec:hr3c-nonclaim}

The extra invariant retained by the square-matrix
\(\mathrm{SU}(3)\) quotient is
\begin{equation}
 z=\det A
 =\varepsilon_{ijk}A_{i1}A_{j2}A_{k3}.
 \label{eq:hr3c-alternating-invariant}
\end{equation}
It is exactly the complex alternating volume of the three columns:
\[
 \det(UA)=\det(U)\det(A).
\]
It is invariant for \(U\in\mathrm{SU}(3)\), its phase is rotated by the
residual \(\mathrm U(3)/\mathrm{SU}(3)\), and it vanishes exactly when the
three columns are linearly dependent.

This invariant theory supplies exact carrier and orbit data for a physical
translation, but does not by itself complete that translation.  The next
construction must specify a quark field, colour-charge and baryon
observables, a principal \(\mathrm{SU}(3)\) bundle with connection and
characteristic classes, and an interacting Yang--Mills sector, together with
maps relating each object to Equation~\eqref{eq:hr3c-alternating-invariant}.
Likewise, Theorem~\ref{thm:hr3c-u3-compactification} supplies a distinguished
\(S^8\), while an almost-complex or integrable-complex structure, a
three-fibre structure, and analogues of the candidate \(S^6\) special-fibre
attaching maps remain exact construction problems rather than consequences
already proved by the compactification theorem.

The exact conclusion is the topological carrier theorem
\begin{equation}
 \boxed{
 \mathscr V_3^+/\mathrm U(3)\cong B^9,\quad
 \partial B^9\cong S^8,\quad
 \mathscr V_3^+/\mathrm{SU}(3)\cong S^{10},\quad
 S^{10}\setminus S^8\simeq S^1.
 }
 \label{eq:hr3c-final-carrier}
\end{equation}
Any gauge-theoretic or geometric use of this carrier requires additional
explicit descent data.

% Canonical workbench input module.
% No document preamble, bibliography, or machine-local path belongs here.

\section{A four-branch toroidal lift and the obstruction to a complex
\texorpdfstring{$S^8$}{S8} total space}
\label{hrftd:section}

\subsection{The exact meaning of the three-way cusp}
\label{hrftd:baseline}

The cusp in the audited threefold is controlled by the
\(\mathbb Z^2\)-periodic \(A_2\)-triangulation of \(\mathbb R^2\).  Its
quotient has one vertex, three edges, and two triangles.  The central fibre
is globally irreducible; its normalization is the degree-six del Pezzo
surface, whose six toric boundary curves are paired into three double curves.
At either of the two deepest points the smooth total threefold maps to the
disc by
\begin{equation}
 (z_0,z_1,z_2)\longmapsto t=z_0z_1z_2.
 \label{hrftd:eq:three-branch}
\end{equation}
Thus the three relevant statements are distinct: there are three local
branches at a deepest point, three global double curves, and only one global
irreducible component.  The generalized dual intersection complex is
\begin{equation}
 \mathcal T_{A_2}/\mathbb Z^2
 \cong \mathbb R^2/\mathbb Z^2=T^2.
 \label{hrftd:eq:dual-torus-two}
\end{equation}
These baseline data are transcribed from the audited canonical-source
snapshot of \(620{,}782\) bytes and \(6{,}315\) lines, with SHA--256 digest
\[
 \text{\scriptsize\ttfamily
 0696B6B7FAB83804DFEC947D45607601E5D7F57D64C6148487969FCCB6C67D2D}.
\]
line \(320\) records the periodic \(A_2\) triangulation and its
\((1,3,2)\) quotient cell counts; lines \(372\) and \(377\) record the
normalization, three double curves, and two triple points; lines
\(1431\)--\(1461\) construct the fan and its unimodular toric charts;
lines \(1647\)--\(1678\) give the quotient local models and orbit strata;
and line \(3966\) gives the deepest central-fibre equation
\(\{z_0z_1z_2=0\}\).  The line references and digest fix the exact source
edition without introducing a machine-local dependency.

The higher-dimensional construction below raises the rank of this periodic
triangulation from two to three.  It does not replace three line bundles by
four line bundles.

\subsection{The rank-three periodic fan}
\label{hrftd:periodic-fan}

Let
\[
 L=\mathbb Z^3,\qquad L_{\mathbb R}=\mathbb R^3,\qquad
 \mathbf 1=e_1+e_2+e_3.
\]
For \(m\in L\) and
\(\pi=(\pi_1,\pi_2,\pi_3)\in S_3\), put
\begin{equation}
 \sigma_{m,\pi}
 =
 \operatorname{conv}\left\{
 m,\ m+e_{\pi_1},\ m+e_{\pi_1}+e_{\pi_2},\
 m+\mathbf 1
 \right\}.
 \label{hrftd:eq:kuhn-simplex}
\end{equation}
Inside the cube \(m+[0,1]^3\), this is the region in which the coordinates,
in the order \(\pi_1,\pi_2,\pi_3\), are weakly decreasing.  The six possible
orders partition the cube, and the induced partitions agree on common
faces.  Their integral translates therefore give a locally finite
\(L\)-periodic triangulation \(\mathcal T_3\) of \(\mathbb R^3\).

Put
\[
 N'=L\oplus\mathbb Z e_4
\]
and, for every cell \(P\) of \(\mathcal T_3\), define
\begin{equation}
 \widehat P
 =
 \operatorname{cone}\{(v,1):v\in P\}
 \subset N'_{\mathbb R}.
 \label{hrftd:eq:cone-over-cell}
\end{equation}
Because all cells are bounded,
\(\widehat P\cap\widehat Q=\widehat{P\cap Q}\).  Hence the cones
\(\widehat P\), together with the zero cone, form a fan
\(\Sigma_3\).

\begin{lemma}[Unimodularity]
\label{hrftd:lemma:unimodular}
Every maximal cone of \(\Sigma_3\) is unimodular.
\end{lemma}

\begin{proof}
For the four lifted vertices of \(\sigma_{m,\pi}\), subtract the first
column \((m,1)\) from the other three columns.  Up to the orientation sign,
the determinant is
\[
 \det\left(
 e_{\pi_1},\
 e_{\pi_1}+e_{\pi_2},\
 \mathbf 1
 \right)=\pm1.
\]
The four primitive ray generators therefore form a basis of \(N'\).
Every face is unimodular as well.
\end{proof}

Let \(Y_3=X_{\Sigma_3}\) be the associated, generally non-quasi-compact,
toric fourfold.  It is smooth by
\cref{hrftd:lemma:unimodular}.  Let
\[
 M'=\operatorname{Hom}(N',\mathbb Z),\qquad
 h=(0,0,0,1)\in M',
\]
and define the toric character
\begin{equation}
 t=\chi^h:Y_3\longrightarrow\mathbb C.
 \label{hrftd:eq:toric-parameter}
\end{equation}
The character \(h\) has value one on every primitive ray generator
\((v,1)\).  Consequently
\begin{equation}
 \operatorname{div}(t)=\sum_{v\in L}E_v
 \label{hrftd:eq:reduced-central-divisor}
\end{equation}
with multiplicity one on every toric divisor \(E_v\).

\begin{proposition}[Four-branch local equation]
\label{hrftd:prop:four-branch}
If a tetrahedron \(\sigma\) has vertices \(v_0,v_1,v_2,v_3\), its toric
chart is \(\mathbb C^4\), and in its dual-character coordinates
\((z_0,z_1,z_2,z_3)\) one has
\begin{equation}
 t=z_0z_1z_2z_3.
 \label{hrftd:eq:four-branch}
\end{equation}
In particular, four distinct components of the divisor upstairs meet
transversely at the torus fixed point of \(\widehat\sigma\).
\end{proposition}

\begin{proof}
Let \(m_0,m_1,m_2,m_3\) be the basis of \(M'\) dual to
\((v_0,1),\ldots,(v_3,1)\), and set \(z_i=\chi^{m_i}\).
Unimodularity identifies the affine chart with \(\mathbb C^4\).
Since \(h\) takes value one on every ray generator,
\[
 h=m_0+m_1+m_2+m_3.
\]
Taking characters gives \eqref{hrftd:eq:four-branch}.
\end{proof}

\subsection{The analytic quotient and the abelian-threefold fibres}
\label{hrftd:quotient}

For \(\lambda\in L\), define the integral shear
\begin{equation}
 \varphi_\lambda:N'\longrightarrow N',
 \qquad
 \varphi_\lambda(y,r)=(y+r\lambda,r).
 \label{hrftd:eq:lattice-shear}
\end{equation}
It sends \(\widehat P\) to \(\widehat{P+\lambda}\), so it induces a toric
automorphism \(\Phi_\lambda\) of \(Y_3\).  These automorphisms preserve \(t\)
and satisfy
\(\Phi_\lambda\Phi_\mu=\Phi_{\lambda+\mu}\).
On the dense torus and over \(t\ne0\), their formula is
\begin{equation}
 \Phi_\lambda(x_1,x_2,x_3,t)
 =
 (t^{\lambda_1}x_1,t^{\lambda_2}x_2,t^{\lambda_3}x_3,t).
 \label{hrftd:eq:dense-action}
\end{equation}
The lattice map \eqref{hrftd:eq:lattice-shear}, rather than the expression
with possibly negative powers in \eqref{hrftd:eq:dense-action}, proves that
the action extends over the toric boundary.

Fix \(0<\epsilon<1\) and set
\[
 Y_{3,\epsilon}=\{p\in Y_3:|t(p)|<\epsilon\}.
\]
Define
\begin{equation}
 N_3=Y_{3,\epsilon}/L,\qquad
 \pi_3:N_3\longrightarrow\Delta_\epsilon.
 \label{hrftd:eq:analytic-quotient}
\end{equation}

\begin{theorem}[The rank-three toroidal cusp]
\label{hrftd:thm:toroidal-cusp}
The \(L\)-action on \(Y_{3,\epsilon}\) is free and properly discontinuous.
The quotient \(N_3\) is a smooth complex fourfold and \(\pi_3\) is proper.
For \(q\ne0\), its fibre is the abelian threefold
\[
 N_{3,q}\cong
 \mathbb C^3/(\mathbb Z^3+\tau\mathbb Z^3)
 \cong E_\tau^3,
 \qquad q=e^{2\pi i\tau},\quad \tau\in\mathbb H.
\]
\end{theorem}

\begin{proof}
A boundary point belongs to a unique torus orbit
\(O_{\widehat P}\).  If \(\Phi_\lambda\) fixes it, then
\(O_{\widehat{P+\lambda}}=O_{\widehat P}\), whence
\(P+\lambda=P\).  Since \(P\) is bounded, \(\lambda=0\).
On a fibre \(t=q\ne0\), \eqref{hrftd:eq:dense-action} is free because
\(|q|<1\) and the equalities \(q^{\lambda_i}=1\) for all \(i\) force
\(\lambda=0\).

For proper discontinuity and properness, use on the dense fibre the
logarithmic coordinate
\begin{equation}
 \operatorname{trop}_q(x)
 =
 \left(
 \frac{\log|x_1|}{\log|q|},
 \frac{\log|x_2|}{\log|q|},
 \frac{\log|x_3|}{\log|q|}
 \right).
 \label{hrftd:eq:tropical-coordinate}
\end{equation}
Equation \eqref{hrftd:eq:dense-action} translates this coordinate by
\(\lambda\).  Every orbit can therefore be moved into
\([0,1]^3\).  If
\[
 r=\sum_{i=0}^3\alpha_i v_i
\]
belongs to
\(\sigma=\operatorname{conv}(v_0,v_1,v_2,v_3)\), then
\(\alpha_i\ge0\), \(\sum_i\alpha_i=1\), and the dual-character
coordinates of \cref{hrftd:prop:four-branch} obey
\[
 \frac{\log|z_i|}{\log|q|}=\alpha_i,
 \qquad |z_i|=|q|^{\alpha_i}\le1.
\]
For \(\epsilon'<\epsilon\), let \(F_{\epsilon'}\) be the closure of the
points with \(0<|q|\le\epsilon'\) and
\(\operatorname{trop}_q(x)\in[0,1]^3\).  For each of the six tetrahedra in
that cube, the preceding identities put these points in the closed unit
polydisc \(\{|z_0|,\ldots,|z_3|\le1\}\) of its
\(\mathbb C^4\)-chart.  The finite union of those polydiscs is compact and
closed in the Hausdorff toric variety; \(F_{\epsilon'}\), being closed in
that union, is compact.  Its translates cover the dense locus, and the same
holds on \(t=0\) because every cell has an \(L\)-translate meeting the
fundamental cube.

It remains to verify proper discontinuity at the boundary, where
\(\operatorname{trop}_q\) itself is undefined.  Let \(K\subset
Y_{3,\epsilon}\) be compact.  Local compactness and the toric affine cover
give finitely many maximal-cone charts \(U_{\sigma_j}\) and compact subsets
\(K_j\subset U_{\sigma_j}\) whose interiors cover \(K\).  On the dense part
of one such compact set, write
\[
 r=\sum_{i=0}^3\alpha_iv_i,\qquad
 \alpha_i=\frac{\log|z_i|}{\log|t|},\qquad
 \sum_i\alpha_i=1.
\]
The coordinate moduli are uniformly bounded above.  When \(|t|\) is small,
this gives a uniform lower bound for every \(\alpha_i\), and the identity
\(\sum_i\alpha_i=1\) gives a uniform upper bound.  When \(|t|\) is bounded
away from zero, the product \(t=z_0z_1z_2z_3\) and the upper bounds give
uniform positive lower bounds for the \(|z_i|\), hence again uniform bounds
for every \(\alpha_i\).  Since only finitely many vertex sets occur, all
dense points of \(K\) consequently have tropical coordinates in one bounded
set \(C_K\subset\mathbb R^3\).

If \(p,\Phi_\lambda(p)\in K\) are dense, then
\(\lambda\in(C_K-C_K)\cap L\), a finite set.  If \(p\) lies on the
boundary, its orbit cell and that of \(\Phi_\lambda(p)\) lie among the
finitely many faces of the \(\sigma_j\), and they satisfy
\(P+\lambda=Q\); again only finitely many \(\lambda\) occur.  Thus a compact
set meets only finitely many of its translates, proving proper
discontinuity.  Finally, the image of \(F_{\epsilon'}\) is the entire inverse
image of \(\{|q|\le\epsilon'\}\), which proves that \(\pi_3\) is proper.
Freeness on the smooth \(Y_{3,\epsilon}\) makes \(N_3\) smooth.

For \(q\ne0\), all non-dense toric orbits lie in \(t=0\); hence
\(Y_{3,q}=(\mathbb C^*)^3\), and
\[
 N_{3,q}=(\mathbb C^*)^3/q^{\mathbb Z^3}.
\]
The map
\[
 [w_1,w_2,w_3]\longmapsto
 [e^{2\pi i w_1},e^{2\pi i w_2},e^{2\pi i w_3}]
\]
induces the asserted isomorphism from
\(\mathbb C^3/(\mathbb Z^3+\tau\mathbb Z^3)\).
\end{proof}

The preceding construction is the untwisted
\(B_0=I\), \(C(t)=0\) instance of the Tate--Mumford construction.
The general analytic construction is due to Mumford,
\emph{An analytic construction of degenerating abelian varieties over
complete rings}, Compositio Math. \textbf{24} (1972), 239--272.
No appeal to that generality is needed for the displayed quotient.

\subsection{Monodromy and the cusp inclusion on fundamental groups}
\label{hrftd:monodromy}

Use the integral basis
\[
 (a_1,a_2,a_3,b_1,b_2,b_3)
\]
of \(H_1(N_{3,q};\mathbb Z)\), where \(a_i\) is the unit period and \(b_i\)
is the \(\tau\)-period in the \(i\)-th factor.  A positive loop about
\(q=0\) sends \(\tau\) to \(\tau+1\).

\begin{proposition}[Cusp monodromy and filling map]
\label{hrftd:prop:cusp-monodromy}
The monodromy is
\begin{equation}
 M_0=
 \begin{pmatrix}I_3&I_3\\0&I_3\end{pmatrix},
 \qquad
 R=M_0-I=
 \begin{pmatrix}0&I_3\\0&0\end{pmatrix}.
 \label{hrftd:eq:cusp-monodromy}
\end{equation}
It satisfies
\[
 R^2=0,\qquad
 \operatorname{rank}R=3,\qquad
 \ker R=\operatorname{im}R
 =\mathbb Z\langle a_1,a_2,a_3\rangle.
\]
Moreover
\[
 \pi_1(N_3)\cong L\cong\mathbb Z^3,
\]
and the inclusion of a smooth fibre induces
\begin{equation}
 a_i\longmapsto0,\qquad b_i\longmapsto e_i\in L.
 \label{hrftd:eq:cusp-pi-one-map}
\end{equation}
\end{proposition}

\begin{proof}
The period change \(b_i\mapsto b_i+a_i\), \(a_i\mapsto a_i\), gives
\eqref{hrftd:eq:cusp-monodromy}; the assertions about \(R\) follow by direct
block multiplication.

The primitive rays
\[
 (0,1),\quad(e_1,1),\quad(e_2,1),\quad(e_3,1)
\]
generate \(N'\).  We record the toric fundamental-group calculation needed
here.  If \(\Sigma_0\) is a finite subfan of this smooth fan and
\(N_{\Sigma_0}\) is the sublattice generated by its primitive ray vectors,
then
\[
 \pi_1(X_{\Sigma_0})
 \cong N'/N_{\Sigma_0}.
\]
Indeed, the dense torus has fundamental group \(N'\).  Adjoining a
neighbourhood of the generic point of the divisor belonging to a primitive
ray \(v\) fills precisely the circle \(v\), so van Kampen quotients by
\(\langle v\rangle\).  The remaining orbit strata have complex codimension
at least two; general position for loops and disc homotopies in the smooth
real manifold shows that adding them changes no fundamental group.  Applying
van Kampen over all rays gives the displayed quotient.

Every loop in \(Y_3\) has compact image and is therefore contained in the
union of finitely many toric affine charts.  Enlarge their finite subfan by
the four rays displayed above.  Those rays generate \(N'\), so the finite
subfan formula makes this enlarged toric open set simply connected.  The
loop bounds there, and hence \(\pi_1(Y_3)=0\).
For \(0<u\le1\), the real one-parameter torus element
\(\vartheta_u=(1,1,1,u)\) satisfies
\[
 t(\vartheta_up)=u\,t(p).
\]
Any loop in \(Y_{3,\epsilon}\) bounds a disc in \(Y_3\).  Choosing \(u\)
small moves that entire disc into \(Y_{3,\epsilon}\), while
\(\vartheta_{u'}\), \(u\le u'\le1\), gives a homotopy back to the original
loop.  Hence \(\pi_1(Y_{3,\epsilon})=0\).  The quotient map in
\eqref{hrftd:eq:analytic-quotient} is therefore the universal cover with
deck group \(L\), so \(\pi_1(N_3)=L\).

The loop \(a_i\) is the angular loop in the \(i\)-th
\(\mathbb C^*\)-coordinate and is killed by the corresponding toric divisor.
The path representing \(b_i\) closes precisely by the deck transformation
\(\Phi_{e_i}\).  This proves \eqref{hrftd:eq:cusp-pi-one-map}.
\end{proof}

\subsection{The central fibre, its normalization, and all strata}
\label{hrftd:central-fibre}

Let
\begin{equation}
 W_3=\pi_3^{-1}(0)
 =
 \left(\bigcup_{v\in L}E_v\right)/L.
 \label{hrftd:eq:central-fibre}
\end{equation}
The toric orbit--cone correspondence identifies its open strata with
\(L\)-orbits of cells of \(\mathcal T_3\).  Its generalized dual
intersection complex is therefore
\begin{equation}
 \mathcal T_3/L\cong\mathbb R^3/L=T^3.
 \label{hrftd:eq:dual-torus-three}
\end{equation}

\begin{theorem}[Forced stratum counts]
\label{hrftd:thm:strata}
The quotient cell structure and the open torus stratification of \(W_3\)
are
\[
\begin{array}{c|c|c|c}
\text{cell}&\text{number}&\text{open stratum}&
\text{local branches}\\ \hline
\text{vertex}&1&(\mathbb C^*)^3&1\\
\text{edge}&7&(\mathbb C^*)^2&2\\
\text{triangle}&12&\mathbb C^*&3\\
\text{tetrahedron}&6&\mathrm{pt}&4.
\end{array}
\]
In particular,
\begin{equation}
 \chi(W_3)=6.
 \label{hrftd:eq:central-euler}
\end{equation}
\end{theorem}

\begin{proof}
There is one vertex orbit.  A fundamental parallelepiped of \(L\) has
normalized three-volume \(3!=6\), whereas every unimodular tetrahedron has
normalized volume one.  Hence there are six tetrahedron orbits.  Every
triangular face is incident to two tetrahedra, so
\[
 f_2=\frac{4f_3}{2}=12.
\]
Since the quotient is \(T^3\),
\[
 0=\chi(T^3)=f_0-f_1+f_2-f_3=1-f_1+12-6,
\]
and therefore \(f_1=7\).  Explicit edge representatives are
\[
 [0,e_1],\ [0,e_2],\ [0,e_3],\
 [0,e_1+e_2],\ [0,e_1+e_3],\
 [0,e_2+e_3],\ [0,\mathbf1].
\]
A \(k\)-cell gives a torus orbit of complex dimension \(3-k\) and an
intersection of \(k+1\) local branches.  Finally every
positive-dimensional \((\mathbb C^*)^j\) has Euler characteristic zero.
Additivity over the finite stratification leaves the six point strata and
proves \eqref{hrftd:eq:central-euler}.
\end{proof}

\begin{theorem}[Normalization]
\label{hrftd:thm:normalization}
\(W_3\) is globally irreducible.  The quotient map induces its normalization
\[
 \nu:E_0\longrightarrow W_3.
\]
The variety \(E_0\) is the smooth three-dimensional permutohedral toric
variety.  It has fourteen toric boundary surfaces, thirty-six invariant
curves, and twenty-four fixed points.  Under \(\nu\), these map respectively
two-to-one, three-to-one, and four-to-one onto seven double surfaces, twelve
triple curves, and six quadruple points.
\end{theorem}

\begin{proof}
The lattice \(L\) acts simply transitively on the vertices and therefore
transitively on the components \(E_v\).  The quotient \(W_3\) is irreducible.
Every cell has a vertex and can be translated to one containing zero, so the
map \(E_0\to W_3\) is surjective.  Only finitely many cells contain zero,
hence the map is finite, and it is injective on the dense torus.  The star
fan at zero is complete and unimodular, so \(E_0\) is smooth, compact, and
normal.  The finite birational map is consequently the normalization.

The rays of the star fan are
\[
 \pm\sum_{i\in I}e_i,
 \qquad
 \varnothing\ne I\subseteq\{1,2,3\}.
\]
These are the fourteen rays of the \(A_3\) braid fan.  Its maximal cones are
indexed by the twenty-four permutations of four symbols; hence the fan is the
permutohedral fan.  There are \(24\) maximal cones.  Each has three
two-dimensional faces and each such face is shared by two maximal cones, so
there are \(24\cdot3/2=36\) two-dimensional cones.  Together with the
fourteen rays this gives the asserted orbit counts.

An edge has two vertices, a triangle has three, and a tetrahedron has four.
The normalization preimage of the corresponding generic quotient stratum
therefore has respectively two, three, or four points.  Division of the
normalization counts by these degrees gives
\[
 14/2=7,\qquad 36/3=12,\qquad 24/4=6.
\]
\end{proof}

\begin{remark}[Local branches are not global components]
\label{hrftd:rem:local-not-global}
At each quadruple point the four local branches are four germs of the same
global irreducible central fibre.  The conductor in the normalization is the
union of fourteen toric boundary surfaces, paired into seven quotient
surfaces.  No canonical four-term decomposition of a vector bundle or line
bundle follows from this local branch number.
\end{remark}

More generally, let \(\bar L\cong\mathbb Z^3\) and let
\(B_0:\bar L\hookrightarrow L\) have finite index \(d\).  Replacing
\(\lambda\) by \(B_0\lambda\) in \eqref{hrftd:eq:lattice-shear} gives
\(d\) vertex orbits and hence \(d\) global central components.  The stratum
counts are
\begin{equation}
 (f_0,f_1,f_2,f_3)=d(1,7,12,6).
 \label{hrftd:eq:index-d-counts}
\end{equation}

\begin{corollary}[Euler characteristic as the determinant of the shear]
\label{hrftd:cor:euler-shear-determinant}
Choose oriented integral bases of \(\bar L\) and \(L\), and let
\([B_0]\in\operatorname{Mat}_{3\times3}(\mathbb Z)\) be the matrix of the
injective shear.  The associated central fibre satisfies
\begin{equation}
 \chi(W_{3,B_0})
 =3!\,[L:B_0\bar L]
 =3!\,|\det[B_0]|.
 \label{hrftd:eq:euler-shear-determinant}
\end{equation}
The right side is unchanged by independent unimodular changes of the two
marked bases.
\end{corollary}

\begin{proof}
The index of the image of a full-rank integral matrix is
\([L:B_0\bar L]=|\det[B_0]|\).  By
\eqref{hrftd:eq:index-d-counts}, the quotient has \(6d\) tetrahedron orbits.
Their zero-dimensional torus strata are the only strata with nonzero Euler
characteristic, so \(\chi(W_{3,B_0})=6d=3!d\).  Replacing the source and
target bases multiplies the determinant by units in \(\mathbb Z\), and hence
does not change its absolute value.
\end{proof}

If a polarization additionally identifies \([B_0]\) with a symmetric
positive-definite matrix, then its determinant is the cubic Jordan norm on
that real symmetric subcone of
\(\operatorname{Herm}_3(\mathbb C)_{\geq0}\).  Without those additional
hypotheses \([B_0]\) is only an integral lattice map, not a positive
Hermitian matrix.  In either case the scalar in
\eqref{hrftd:eq:euler-shear-determinant} does not determine the fan, the
period correction, the singular filling, or the total space.

Thus four global components occur only after the independent choice \(d=4\);
four local branches already occur for \(d=1\).  A nontrivial Fable-style
twist would additionally require holomorphic translations
\(c_\lambda(t)\) satisfying the exact group cocycle, extension, freeness,
and properness identities.

\subsection{Two independent obstructions to a complex
\texorpdfstring{$S^8$}{S8} globalization}
\label{hrftd:s8-obstructions}

\begin{theorem}[Euler obstruction to the literal three-special-fibre pattern]
\label{hrftd:thm:euler-obstruction}
Let \(g:M^8\to S^2\) be a proper triangulable map which is a locally trivial
\(T^6\)-bundle away from finitely many critical values.  If one singular
fibre is \(W_3\) and every other singular fibre is a finite free quotient of
\(T^6\), then
\[
 \chi(M)=6.
\]
In particular, \(M\) is not homeomorphic to \(S^8\).
\end{theorem}

\begin{proof}
Choose a triangulation for \(g\) in which every critical value is a vertex.
Additivity of compactly supported Euler characteristic over the open
regular stratum and the critical singleton strata, proved cell by cell in
this finite triangulation, gives
\[
 \chi(M)
 =
 \chi_c(S^2\setminus\Sigma)\chi(T^6)
 +\sum_{p\in\Sigma}\chi(g^{-1}(p)).
\]
Here multiplicativity on the regular stratum follows by subdividing its
locally trivial bundle over base cells.  Since \(M\) is compact,
\(\chi_c(M)=\chi(M)\), and since \(\chi(T^6)=0\), the first term vanishes.
By \eqref{hrftd:eq:central-euler}, the cusp contributes \(6\).
A finite free quotient by a group \(G\) has Euler characteristic
\(\chi(T^6)/|G|=0\).  Therefore
\(\chi(M)=6\), whereas \(\chi(S^8)=2\).
\end{proof}

Consequently, retaining the cusp \(W_3\) in a real \(S^8\) candidate would
require additional genuinely singular fibres whose Euler characteristics sum
to \(-4\).  Smooth multiple torus fibres cannot supply them.

\begin{theorem}[No smooth complex \(S^8\) total space]
\label{hrftd:thm:no-almost-complex-s8}
A smooth manifold homotopy equivalent to \(S^8\) admits no almost-complex
structure.  Hence a complex structure on the regular locus of a putative
\(T^6\)-fibration on \(S^8\) cannot extend through all its singular fillings.
\end{theorem}

\begin{proof}
Suppose \(E=(TM,J)\) is a complex rank-four bundle whose underlying real
bundle is \(TM\).  Orient \(M\) by \(J\), and let
\(u\in H^8(M;\mathbb Z)\) evaluate to one.  Since
\[
 H^2(M;\mathbb Z)=H^4(M;\mathbb Z)=H^6(M;\mathbb Z)=0,
\]
one has
\[
 c_1(E)=c_2(E)=c_3(E)=0,\qquad
 c_4(E)=e(TM)=\chi(M)u=2u.
\]
For the underlying real bundle,
\begin{equation}
 p_2(TM)
 =c_2(E)^2-2c_1(E)c_3(E)+2c_4(E)
 =4u.
 \label{hrftd:eq:p-two-contradiction}
\end{equation}
On the other hand \(H^4(M)=0\) gives \(p_1(TM)=0\), and the middle
cohomology vanishes, so \(\operatorname{sign}(M)=0\).  Hirzebruch's
eight-dimensional signature formula gives
\[
 0
 =
 \left\langle
 \frac{7p_2(TM)-p_1(TM)^2}{45},[M]
 \right\rangle,
\]
which forces \(p_2(TM)=0\).  This contradicts
\eqref{hrftd:eq:p-two-contradiction}.  For the standard sphere the same
vanishing follows immediately from
\(TS^8\oplus\varepsilon^1_{\mathbb R}\cong\varepsilon^9_{\mathbb R}\).
\end{proof}

If every local filling were a smooth complex toroidal fourfold and every
overlap map were biholomorphic, these local complex tangent bundles would glue
to an almost-complex structure on all of \(S^8\), contradicting
\cref{hrftd:thm:no-almost-complex-s8}.  At least one filling, transition, or
total-space singularity must therefore leave the smooth complex category.

\subsection{An explicit integral monodromy candidate and its failure}
\label{hrftd:finite-monodromy}

Put
\[
 S=
 \begin{pmatrix}0&-1\\1&0\end{pmatrix},
 \qquad
 T=
 \begin{pmatrix}1&1\\0&1\end{pmatrix},
 \qquad
 Q=S^{-1}T^{-1}
 =
 \begin{pmatrix}0&1\\-1&1\end{pmatrix}.
\]
Then
\begin{equation}
 S^4=I,\qquad Q^6=I,\qquad SQ=T^{-1}.
 \label{hrftd:eq:two-by-two-relations}
\end{equation}
On
\(\Lambda=\mathbb Z^3\otimes\mathbb Z^2\cong\mathbb Z^6\), define
\begin{equation}
 M_1=I_3\otimes S,\qquad
 M_2=I_3\otimes Q,\qquad
 M_0=I_3\otimes T.
 \label{hrftd:eq:finite-monodromies}
\end{equation}
They are integral and symplectic, and
\[
 M_1M_2M_0=I.
\]
After reordering the tensor-product basis as
\((a_1,a_2,a_3,b_1,b_2,b_3)\), \(M_0\) is exactly the cusp matrix
\eqref{hrftd:eq:cusp-monodromy}.  Thus
\eqref{hrftd:eq:finite-monodromies} is an explicit local system over a
thrice-punctured sphere with finite orders \(4,6\) and a rank-three
unipotent cusp.

\begin{proposition}[Fixed-point obstruction]
\label{hrftd:prop:finite-fixed-points}
No translation makes either finite linear action \(M_1\) or \(M_2\) free on
\(T^6\).
\end{proposition}

\begin{proof}
Neither \(S\) nor \(Q\) has eigenvalue one.  More explicitly,
\[
 \det(I-S)=2,\qquad \det(I-Q)=1.
\]
Consequently \(I-M_j:T^6\to T^6\) is a surjective isogeny for \(j=1,2\).
For any translation parameter \(c\in T^6\), the fixed-point equation for
\[
 x\longmapsto M_jx+c
\]
is
\[
 (I-M_j)x=c,
\]
and it always has a solution.  Therefore every affine lift has a fixed point.
\end{proof}

This is an integral symplectic homology local system compatible with a
diagonal rank-three period ansatz.  The three matrices alone do not define a
polarized variation of Hodge structure or a holomorphic family
\(E_\tau^3\); those require an equivariant period map in addition.  Any
isotrivial affine quotient filling built from either displayed finite
generator is singular by the fixed-point calculation.  A successful
alternative would need finite
monodromies with invariant torus directions that admit fixed-point-free
affine lifts, while retaining the product relation and the rank-three cusp.
It would still have to overcome both obstructions in
\cref{hrftd:s8-obstructions}.

\subsection{Exact replication requirements for any real singular
\texorpdfstring{$T^6\to S^8\to S^2$}{T6 to S8 to S2} candidate}
\label{hrftd:replication}

The preceding results exclude the direct smooth holomorphic construction, not
an arbitrary real stratified map.  Let
\[
 \Sigma=\{p_0,\ldots,p_s\}\subset S^2,\qquad
 B^\circ=S^2\setminus\Sigma.
\]
A candidate must first specify an affine integral monodromy
\begin{equation}
 \widetilde\rho:\pi_1(B^\circ)\longrightarrow
 \operatorname{GL}(6,\mathbb Z)\ltimes T^6
 \label{hrftd:eq:affine-monodromy}
\end{equation}
with based peripheral matrices satisfying
\[
 M_0M_1\cdots M_s=I.
\]
The corresponding affine peripheral transformations
\(\widetilde M_i=(M_i,c_i)\) must satisfy the stronger identity
\begin{equation}
 \widetilde M_0\widetilde M_1\cdots\widetilde M_s=(I,0),
 \qquad
 c_0+M_0c_1+\cdots+M_0\cdots M_{s-1}c_s=0
 \quad\text{in }T^6.
 \label{hrftd:eq:affine-product}
\end{equation}
Choose pairwise disjoint small closed discs \(D_i\) around the \(p_i\), and
put
\[
 B=S^2\setminus\bigcup_i\operatorname{int}D_i,\qquad
 J=g^{-1}(B).
\]
It must then give each compact filling \(N_i=g^{-1}(D_i)\), its boundary
mapping torus
\[
 A_i=J\cap N_i,
\]
and the actual based inclusion maps
\[
 A_i\longrightarrow J,\qquad A_i\longrightarrow N_i.
\]
Here one fixes a base point in \(J\), a base point in each \(A_i\), and a
whisker in \(J\) joining them; changing a whisker only conjugates the
corresponding homomorphism and does not change the pushout below.

\begin{proposition}[Fundamental-group replication map]
\label{hrftd:prop:van-kampen}
The fundamental group of the glued total space is the iterated pushout
\begin{equation}
 \pi_1(M)
 \cong
 \pi_1(J)
 *_{\pi_1(A_0)}\pi_1(N_0)
 *_{\pi_1(A_1)}\cdots
 *_{\pi_1(A_s)}\pi_1(N_s),
 \label{hrftd:eq:van-kampen}
\end{equation}
where both homomorphisms out of each \(\pi_1(A_i)\) are the literal based
inclusion homomorphisms.  If, in addition, every
\(\pi_1(A_i)\to\pi_1(N_i)\) is surjective, then
\begin{equation}
 \pi_1(M)
 \cong
 \pi_1(J)\Big/
 \Big\langle\!\Big\langle
 (j_{J,i})_*\ker\{\pi_1(A_i)\to\pi_1(N_i)\}:0\le i\le s
 \Big\rangle\!\Big\rangle .
 \label{hrftd:eq:van-kampen-surjective}
\end{equation}
Consequently a proposed \(S^8\) construction must either compute the full
pushout in \eqref{hrftd:eq:van-kampen}, or prove the stated surjectivities
and prove that the displayed normal generators kill \(\pi_1(J)\).
\end{proposition}

\begin{proof}
Apply van Kampen successively to
\[
 M=J\cup N_0\cup\cdots\cup N_s
\]
with connected collar intersections \(A_i\).  Each attachment quotients by
the relations identifying the two images of every element of
\(\pi_1(A_i)\), which is precisely the iterated amalgamated product
\eqref{hrftd:eq:van-kampen}.  If the map to \(\pi_1(N_i)\) is surjective,
then every generator coming from \(\pi_1(N_i)\) can be represented by an
element of \(\pi_1(A_i)\).  Eliminating those generators leaves exactly the
normal closure of the image in \(\pi_1(J)\) of the kernel of the map to
\(\pi_1(N_i)\).  Iterating proves
\eqref{hrftd:eq:van-kampen-surjective}.
\end{proof}

For homology, define the literal Mayer--Vietoris map
\begin{equation}
 \alpha_q
 =
 \bigoplus_i\bigl((j_{J,i})_*,-(j_{N_i})_*\bigr):
 \bigoplus_iH_q(A_i;\mathbb Z)
 \longrightarrow
 H_q(J;\mathbb Z)\oplus\bigoplus_iH_q(N_i;\mathbb Z).
 \label{hrftd:eq:mv-map}
\end{equation}

\begin{theorem}[Integral homology-sphere criterion]
\label{hrftd:thm:mv-criterion}
Assume that \(J\) and the \(N_i\) are compact, connected,
collar-compatible eight-manifolds with nonempty boundary, so their top
homology vanishes.  The glued space has the integral homology of \(S^8\) if
and only if
\[
\begin{array}{c|c}
q&\text{condition on }\alpha_q\\ \hline
0&\alpha_0\text{ is injective and }
   \operatorname{coker}\alpha_0\cong\mathbb Z,\\
1\le q\le6&\alpha_q\text{ is an isomorphism},\\
7&\alpha_7\text{ is surjective and }
   \ker\alpha_7\cong\mathbb Z.
\end{array}
\]
\end{theorem}

\begin{proof}
Insert \eqref{hrftd:eq:mv-map} into the Mayer--Vietoris long exact sequence.
For every \(1\le q\le7\), exactness gives the short exact sequence
\[
 0\longrightarrow\operatorname{coker}\alpha_q
 \longrightarrow H_q(M;\mathbb Z)
 \longrightarrow\ker\alpha_{q-1}
 \longrightarrow0.
\]
Thus simultaneous vanishing in degrees \(1,\ldots,7\) is equivalent to
\(\alpha_0\) being injective, \(\alpha_q\) being an isomorphism for
\(1\le q\le6\), and \(\alpha_7\) being surjective.  In degree zero,
\(H_0(M)\cong\operatorname{coker}\alpha_0\), which gives the remaining
degree-zero condition.  Since the pieces have zero \(H_8\), exactness at
\[
 0\longrightarrow H_8(M)\longrightarrow
 \bigoplus_iH_7(A_i)\xrightarrow{\alpha_7}
 H_7(J)\oplus\bigoplus_iH_7(N_i)\longrightarrow H_7(M)
 \longrightarrow0
\]
identifies \(H_8(M)\) with \(\ker\alpha_7\).  Hence
\(H_8(M)\cong\mathbb Z\) is exactly the last kernel condition.
\end{proof}

If \(A_i\) is the mapping torus of a torus automorphism with integral
monodromy \(M_i\), its integral homology is fixed by the Wang exact sequence:
\begin{equation}
 0\longrightarrow
 \operatorname{coker}\left(I-\bigwedge^qM_i\right)
 \longrightarrow H_q(A_i;\mathbb Z)
 \longrightarrow
 \ker\left(I-\bigwedge^{q-1}M_i\right)
 \longrightarrow0.
 \label{hrftd:eq:wang}
\end{equation}
The right-hand group is free, so this sequence splits as a sequence of
abelian groups, though not canonically.  Equations
\eqref{hrftd:eq:affine-monodromy},
\eqref{hrftd:eq:van-kampen}, \eqref{hrftd:eq:mv-map}, and
\eqref{hrftd:eq:wang}, together with the actual tangent clutching maps, are
the minimum replication data.  For an \(S^8\) total space those clutching
maps must give
\[
 p_1=p_2=0,\qquad e=2u,
\]
while preventing the regular-locus complex structure from extending, as
required by \cref{hrftd:thm:no-almost-complex-s8}.

\subsection{The Hermitian determinant boundary is not the four-branch fibre}
\label{hrftd:hermitian}

Let
\[
 \mathcal C_3=\operatorname{Herm}_3(\mathbb C)_{\ge0}.
\]
The real dimension of the positive-semidefinite rank-\(r\) stratum is
\[
 2\cdot3r-r^2.
\]
Indeed, write \(H=BB^*\) with \(B\) a full-rank complex
\(3\times r\) matrix; the entire redundancy is the right action of
\(U(r)\).  Thus the rank \(3,2,1,0\) strata have real dimensions
\begin{equation}
 9,\quad8,\quad5,\quad0.
 \label{hrftd:eq:rank-dimensions}
\end{equation}
The rank-two stratum is smooth because the derivative of the determinant is
the trace pairing with \(\operatorname{adj}(H)\), which is nonzero exactly
in rank two.  The singular locus is the rank-at-most-one locus.

\begin{proposition}[The real Hermitian boundary compactifies to \(S^8\)]
\label{hrftd:prop:hermitian-compactification}
The determinant-zero boundary
\[
 \partial\mathcal C_3
 =\{H\in\operatorname{Herm}_3(\mathbb C):H\ge0,\ \det H=0\}
\]
is homeomorphic to \(\mathbb R^8\).  Its one-point compactification is
therefore \(S^8\), with the rank strata in
\eqref{hrftd:eq:rank-dimensions} giving a non-smooth stratification.
\end{proposition}

\begin{proof}
Intersect \(\mathcal C_3\) with the affine hyperplane \(\operatorname{tr}H=1\):
\[
 K=\{H\ge0:\operatorname{tr}H=1\}.
\]
This is a compact convex body with nonempty interior in the
eight-dimensional affine space
\(\{H=H^*:\operatorname{tr}H=1\}\).  Radial projection from the interior
point \(I_3/3\) gives a homeomorphism from \(K\) to the closed
eight-ball.  Explicitly, after choosing a Euclidean norm on the trace-zero
translation space, every ray
\[
 I_3/3+sU,\qquad \|U\|=1,\quad s\ge0,
\]
meets \(\partial K\) once, at a radius \(r(U)>0\); compactness and convexity
make \(r\) continuous, and
\((\rho,U)\mapsto I_3/3+\rho r(U)U\), \(0\le\rho\le1\), is the required
radial homeomorphism after collapsing \(\{0\}\times S^7\) to the centre.
Thus \(\partial K\cong S^7\).

Every nonzero \(H\in\partial\mathcal C_3\) has the unique expression
\[
 H=rP,\qquad r=\operatorname{tr}H>0,\qquad P\in\partial K.
\]
Adding \(H=0\) identifies
\[
 \partial\mathcal C_3
 \cong
 \bigl([0,\infty)\times S^7\bigr)/(\{0\}\times S^7)
 \cong\mathbb R^8.
\]
The Alexandroff one-point compactification of \(\mathbb R^8\) is \(S^8\).
\end{proof}

\begin{theorem}[Non-SNC determinant boundary]
\label{hrftd:thm:hermitian-not-snc}
At a rank-one point, the semialgebraic boundary of \(\mathcal C_3\) has a
quadratic-cone normal slice rather than a union of smooth crossing faces.
The ambient real determinant hypersurface, and its complexification, are not
normal-crossings hypersurfaces there.  Neither that rank-one germ nor the
germ at the zero matrix is the four-branch germ in
\eqref{hrftd:eq:four-branch}.
\end{theorem}

\begin{proof}
Unitary conjugation and positive rescaling reduce any rank-one point to
\[
 H_0=\operatorname{diag}(1,0,0).
\]
Near \(H_0\), write
\[
 H=
 \begin{pmatrix}a&b^*\\b&C\end{pmatrix},
 \qquad
 a\in\mathbb R,\quad b\in\mathbb C^2,\quad
 C\in\operatorname{Herm}_2(\mathbb C).
\]
On \(a>0\), the real-analytic coordinate change
\[
 (a,b,C)\longmapsto(a,b,S),\qquad
 S=C-a^{-1}bb^*
\]
is a local diffeomorphism.  The Schur-complement identities are exact:
\begin{equation}
 \det H=a\det S,\qquad
 H\ge0\quad\Longleftrightarrow\quad S\ge0.
 \label{hrftd:eq:schur}
\end{equation}
Writing
\[
 S=\begin{pmatrix}u&x+iy\\x-iy&v\end{pmatrix}
\]
gives the exact semialgebraic normal boundary
\begin{equation}
 uv-x^2-y^2=0,\qquad u\ge0,\quad v\ge0.
 \label{hrftd:eq:quadratic-cone}
\end{equation}
The defining quadratic form has rank four.  A product of two real linear
forms has quadratic rank at most two, so the ambient real-algebraic
hypersurface is not a union of smooth real hyperplane branches; imposing the
displayed positive-semidefinite inequalities cannot turn its quadratic cone
into transverse smooth faces.

After complexification, the normal slice is
\[
 s_{11}s_{22}-s_{12}s_{21}=0.
\]
All four first derivatives vanish only at the origin.  A product of two
linear forms would be singular along their codimension-two intersection, so
the quadratic is irreducible and is not a coordinate crossing.  At the zero
matrix the \(3\times3\) determinant is homogeneous of degree three and
therefore has multiplicity three.  The germ
\(z_0z_1z_2z_3=0\) has multiplicity four.  Neither germ is the other.
\end{proof}

There is also a dimension obstruction.  The toroidal central fibre \(W_3\)
has complex dimension three and real dimension six, while the dense rank-two
boundary of \(\mathcal C_3\) has real dimension eight.  The one-point
compactification in \cref{hrftd:prop:hermitian-compactification} is a
stratified \(S^8\), but it is not the central divisor of the complex
fourfold \(N_3\).

The two occurrences of the number three are independent:
\(\det_{\operatorname{Herm}_3}\) has three eigenvalue factors, whereas a
semistable degeneration of complex three-tori has \(3+1=4\) local branches
at a deepest point because one includes the base parameter.  Moving to
\(\operatorname{Herm}_4\) makes the determinant quartic but changes the real
ambient dimension to \(16\) and the natural compactified boundary dimension
to \(15\); it does not produce an \(S^8\).

\subsection{Exact conclusion and permissible carriers}
\label{hrftd:conclusion}

The local higher-rung toroidal object is exact:
\[
 \boxed{
 \text{rank-three periodic fan}
 \quad\Longrightarrow\quad
 t=z_0z_1z_2z_3,\quad
 \mathcal D(W_3)\cong T^3,\quad
 (f_0,f_1,f_2,f_3)=(1,7,12,6).
 }
\]
It lives on the smooth noncompact complex fourfold \(N_3\to\Delta\), not on
a complex \(S^8\).  A compact complex globalization must have different
topology.  A real \(S^8\) may retain the complex-torus structure only on its
regular locus and requires explicitly non-complex or singular fillings
satisfying every map in \cref{hrftd:replication}.  The Hermitian Gram cone
provides a separate real stratified \(S^8\) boundary, not an SNC toroidal
fibre.  Stable complex data are also permissible because
\[
 TS^8\oplus\varepsilon^2_{\mathbb R}
 \cong\varepsilon^{10}_{\mathbb R},
\]
but this stable structure cannot destabilize to an almost-complex structure
on \(TS^8\).

Finally, the six tetrahedron orbits are naturally indexed by the six
permutations in \(S_3\); neither this integer six nor the rank-six lattice
\(H_1(T^6;\mathbb Z)\) is by itself a cyclic group
\(\mathbb Z/6\mathbb Z\).  Any cyclic-order-six or gauge-centre
interpretation requires an additional exact homomorphism.

\section{An exact toroidal elliptic bridge to the candidate cusp}
\label{sec:teb-main}

This section asks for the smallest typed morphism between the rank-three
toroidal construction and the reconstructed Fabel--Alpöge threefold.  It also
tests whether that morphism is supplied by the nine-dimensional Hermitian Gram
cone.  The conclusions are distinct:
\begin{enumerate}
\item a quotient of integral lattices extends to a morphism of toric cusp
families and commutes with monodromy and the cusp filling maps;
\item the quotient has an explicit unimodular splitting as the candidate cusp
plus one elliptic cusp;
\item the order-three and order-four affine peripheral actions admit exact rank-six
lattice lifts;
\item none of these maps is induced by the determinant-rank stratification or
the rank-one \(\mathbb{CP}^{2}\) normal cones of the Gram construction.
\end{enumerate}
The last assertion is an obstruction to a structure-preserving
identification.  It is not the false assertion that no continuous map between
the underlying spheres exists.

\subsection{The candidate lattice and filling data}
\label{subsec:teb-candidate-data}

Retain the candidate map
\[
 f:X\longrightarrow B\simeq\mathbb P^{1},
 \qquad B^{\circ}=B\setminus\{p_{0},p_{1},p_{2}\}.
\]
Its smooth fibre has marked first homology
\[
 \Lambda_F=
 \mathbb Z\langle\widehat\gamma,\widehat u,
                       \widehat w,\widehat\delta\rangle.
\]
Put
\[
 A_F=\mathbb Z\langle\widehat w,\widehat\delta\rangle,
 \qquad
 B_F=\mathbb Z\langle\widehat\gamma,\widehat u\rangle.
\]
Thus \(\Lambda_F=A_F\oplus B_F\), with the decomposed order
\((\widehat w,\widehat\delta;\widehat\gamma,\widehat u)\).  Define
\[
 B_0=\begin{pmatrix}0&1\\-1&0\end{pmatrix}:B_F\longrightarrow A_F.
 \tag{TEB.1}
\]
The cusp monodromy is
\[
 M_F(a,b)=(a+B_0b,b),
 \tag{TEB.2}
\]
or, in the manuscript order
\((\widehat\gamma,\widehat u,\widehat w,\widehat\delta)\),
\[
 [M_F]=
 \begin{pmatrix}
 1&0&0&0\\
 0&1&0&0\\
 0&1&1&0\\
 -1&0&0&1
 \end{pmatrix}.
 \tag{TEB.3}
\]
In particular,
\[
 (M_F-I)^2=0,
 \qquad
 \ker(M_F-I)=\operatorname{im}(M_F-I)=A_F.
 \tag{TEB.4}
\]

The toric cusp neighbourhood retracts onto its central fibre \(W\).  In the
marking above, the induced filling homomorphism is
\[
 c_F:H_1(T^4;\mathbb Z)=A_F\oplus B_F\longrightarrow H_1(W;\mathbb Z)=B_F,
 \qquad c_F(a,b)=b.
 \tag{TEB.5}
\]
Thus the two toric angular cycles are killed and the two deck cycles survive.
The exact homology of the central fibre is
\[
 H_*(W;\mathbb Z)=
 (\mathbb Z,\mathbb Z^2,\mathbb Z^4,\mathbb Z^2,\mathbb Z).
 \tag{TEB.6}
\]
At the other two special points, the finite monodromies are
\[
 A_1=
 \begin{pmatrix}
 1&0&0&0\\6&0&1&0\\-6&-1&-1&0\\-2&1&0&1
 \end{pmatrix},
 \quad
 A_2=
 \begin{pmatrix}
 1&0&0&0\\0&0&-1&0\\-6&1&0&0\\3&0&1&1
 \end{pmatrix},
 \tag{TEB.7}
\]
in the manuscript order.  They satisfy
\[
 A_1^3=I,\qquad A_2^4=I,\qquad A_1A_2M_F=I.
 \tag{TEB.8}
\]
The invariant primitive vectors
\[
 v_1=\widehat\gamma+2\widehat u-4\widehat w,
 \qquad
 v_2=-\widehat\gamma-3\widehat u+3\widehat w
 \tag{TEB.9}
\]
define free affine actions
\[
 x\longmapsto A_jx+v_j/m_j,
 \qquad (m_1,m_2)=(3,4).
 \tag{TEB.10}
\]

\subsection{The rank-three cusp}
\label{subsec:teb-rank-three}

Let
\[
 L_3=\mathbb Z\langle e_1,e_2,e_3\rangle.
\]
For \(m\in L_3\) and \(\pi=(\pi_1,\pi_2,\pi_3)\in S_3\), put
\[
 \sigma_{m,\pi}=\operatorname{conv}\{m,
 m+e_{\pi_1},m+e_{\pi_1}+e_{\pi_2},m+e_1+e_2+e_3\}.
 \tag{TEB.11}
\]
These tetrahedra give the periodic Kuhn triangulation \(\mathcal T_3\) of
\(L_{3,\mathbb R}\).  Coning all cells at height one gives a smooth toric
fourfold \(Y_3\), because every maximal cone is unimodular.  The height
character gives
\[
 t:Y_3\longrightarrow\mathbb C,
\]
and a maximal affine chart has the exact equation
\[
 t=z_0z_1z_2z_3.
 \tag{TEB.12}
\]

For the untwisted quotient, a noncentral fibre is
\[
 \mathbb C^3/(s\mathbb Z^3+\mathbb Z^3),
 \qquad t=e^{2\pi i s},
\]
and its monodromy, in angular--deck order, is
\[
 M_3=\begin{pmatrix}I_3&I_3\\0&I_3\end{pmatrix},
 \qquad M_3(a,b)=(a+b,b).
 \tag{TEB.13}
\]
The toric filling map is
\[
 c_3:H_1(T^6;\mathbb Z)=L_3\oplus\overline L_3
 \longrightarrow\overline L_3,
 \qquad c_3(a,b)=b.                                    \tag{TEB.14}
\]
Indeed, each angular generator bounds a toric disc, while a deck generator
closes only by the corresponding lattice transformation and survives in the
filling group.

\subsection{The exact \texorpdfstring{\(A_2\)}{A2} quotient of the Kuhn fan}
\label{subsec:teb-fan-map}

In \(A_F\) define
\[
 n_1=\widehat w,\qquad
 n_2=\widehat\delta-\widehat w,\qquad
 n_3=-\widehat\delta.
 \tag{TEB.15}
\]
They are primitive and satisfy \(n_1+n_2+n_3=0\).  Define
\[
 P:L_3\longrightarrow A_F,
 \qquad P(e_i)=n_i.
 \tag{TEB.16}
\]
In the bases \((e_1,e_2,e_3)\) and
\((\widehat w,\widehat\delta)\),
\[
 P=\begin{pmatrix}1&-1&0\\0&1&-1\end{pmatrix},
 \qquad \ker P=\mathbb Z(1,1,1).
 \tag{TEB.17}
\]
Define the deck-lattice quotient
\[
 Q=B_0^{-1}P=
 \begin{pmatrix}0&-1&1\\1&-1&0\end{pmatrix}:
 \overline L_3\longrightarrow B_F.
 \tag{TEB.18}
\]
Then
\[
 B_0Q=P,\qquad \ker Q=\mathbb Z(1,1,1).
 \tag{TEB.19}
\]

\begin{proposition}[The Kuhn-to-\(A_2\) fan morphism]
\label{prop:teb-fan-morphism}
The map \(P\) sends every tetrahedron of \(\mathcal T_3\) onto a triangle of
the candidate's \(A_2\) triangulation.  Consequently
\[
 \widehat P:L_3\oplus\mathbb Z\longrightarrow A_F\oplus\mathbb Z,
 \qquad \widehat P(x,r)=(Px,r),
 \tag{TEB.20}
\]
is a map of height-one fans and induces a toric morphism
\[
 F:Y_3\longrightarrow Y_2,
 \qquad t\circ F=t.
 \tag{TEB.21}
\]
\end{proposition}

\begin{proof}
The four vertices of \(\sigma_{m,\pi}\) map to
\[
 Pm,\quad Pm+n_{\pi_1},\quad
 Pm+n_{\pi_1}+n_{\pi_2}=Pm-n_{\pi_3},\quad Pm.
 \tag{TEB.22}
\]
Write \((n_1,n_2,n_3)=(e_1,e_2-e_1,-e_2)\) in the target basis.  For the six
permutations, the nonrepeated images are
\[
\begin{array}{c|c}
(1,2,3)&\{v,v+e_1,v+e_2\}\\
(2,1,3)&\{v,v-e_1+e_2,v+e_2\}\\
(1,3,2)&\{v,v+e_1,v+e_1-e_2\}\\
(3,1,2)&\{v,v-e_2,v+e_1-e_2\}\\
(2,3,1)&\{v,v-e_1+e_2,v-e_1\}\\
(3,2,1)&\{v,v-e_2,v-e_1\}.
\end{array}                                             \tag{TEB.23}
\]
These are precisely the six triangles incident to \(v=Pm\) in the periodic
triangulation by
\[
 \operatorname{conv}\{u,u+e_1,u+e_2\},
 \quad
 \operatorname{conv}\{u+e_1,u+e_2,u+e_1+e_2\}.
\]
Thus every source cone maps into a target cone.  The toric fan criterion gives
\(F\); the last coordinate in (TEB.20) is unchanged, so the height character
is unchanged.
\end{proof}

The kernel \(\mathbb Z(1,1,1)\) is invariant under permutations of the three
coordinates.  Hence the quotient lattice is natural up to the Weyl action on
the \(A_2\) target.  The ordered basis and orientation used in (TEB.1) choose
one representative of that Weyl orbit.

\subsection{Monodromy and filling compatibility}
\label{subsec:teb-mono-filling}

\begin{theorem}[Exact local lattice bridge]
\label{thm:teb-local-lattice-bridge}
The homomorphism
\[
 \Phi=P\oplus Q:L_3\oplus\overline L_3
 \longrightarrow A_F\oplus B_F,
 \qquad \Phi(a,b)=(Pa,Qb),
 \tag{TEB.24}
\]
is surjective and obeys
\[
 \Phi M_3=M_F\Phi,
 \qquad c_F\Phi=Qc_3.                                  \tag{TEB.25}
\]
Thus the diagram
\[
\begin{CD}
H_1(T^6;\mathbb Z) @>{\Phi}>> H_1(T^4;\mathbb Z)\\
@V{c_3}VV @VV{c_F}V\\
\overline L_3 @>{Q}>> B_F
\end{CD}
\tag{TEB.26}
\]
commutes.
\end{theorem}

\begin{proof}
Both \(P\) and \(Q\) are surjective.  Using (TEB.13),
\[
 \Phi M_3(a,b)=(P(a+b),Qb).
\]
Using (TEB.2) and \(B_0Q=P\),
\[
 M_F\Phi(a,b)=(Pa+B_0Qb,Qb)=(Pa+Pb,Qb).
\]
This proves monodromy equivariance.  The filling identity is the literal
calculation
\[
 c_F\Phi(a,b)=Qb=Qc_3(a,b).
\]
\end{proof}

\subsection{The actual period correction and the cusp-family morphism}
\label{subsec:teb-period-lift}

The candidate period matrix at the cusp has the exact form
\[
 Z_2(t)=sB_0+C(t),
 \qquad t=e^{2\pi i s},
 \tag{TEB.27}
\]
where \(C(t)\) is holomorphic and bounded at \(t=0\).  Fix the integral right
inverse
\[
 S=\begin{pmatrix}1&1\\0&1\\0&0\end{pmatrix},
 \qquad PS=I_2,
 \tag{TEB.28}
\]
and define
\[
 C_3(t)=S C(t)Q,
 \qquad Z_3(t)=sI_3+C_3(t).
 \tag{TEB.29}
\]

\begin{theorem}[Holomorphic toroidal elliptic bridge]
\label{thm:teb-holomorphic-cusp-bridge}
After shrinking the cusp disc, \(\operatorname{Im}Z_3(t)\) is invertible and
\[
 PZ_3(t)=Z_2(t)Q.                                      \tag{TEB.30}
\]
Consequently \(P:\mathbb C^3\to\mathbb C^2\) descends on every punctured
fibre to a surjective holomorphic homomorphism
\[
 \frac{\mathbb C^3}{Z_3(t)\mathbb Z^3+\mathbb Z^3}
 \longrightarrow
 \frac{\mathbb C^2}{Z_2(t)\mathbb Z^2+\mathbb Z^2},    \tag{TEB.31}
\]
with connected elliptic kernel and integral homology map \(\Phi\).

Define the source and target toric actions by
\[
 \Psi^{(3)}_\lambda(x,t)=
 \bigl(\exp(2\pi iC_3(t)\lambda)t^\lambda x,t\bigr),
 \qquad \lambda\in\mathbb Z^3,                         \tag{TEB.32}
\]
\[
 \Psi^{(2)}_\mu(y,t)=
 \bigl(\exp(2\pi iC(t)\mu)t^{B_0\mu}y,t\bigr),
 \qquad \mu\in\mathbb Z^2.                            \tag{TEB.33}
\]
Then
\[
 F\Psi^{(3)}_\lambda=\Psi^{(2)}_{Q\lambda}F.          \tag{TEB.34}
\]
The actions are free and properly discontinuous on sufficiently small toric
cusp neighbourhoods, and (TEB.21) descends to a holomorphic map
\[
 \overline F:N_3^C\longrightarrow N_0                 \tag{TEB.35}
\]
over the cusp disc.  Its smooth-fibre and filling maps are (TEB.24) and
(TEB.26).
\end{theorem}

\begin{proof}
Equations (TEB.19), (TEB.28), and (TEB.29) give
\[
 PZ_3=sP+PSCQ=sB_0Q+CQ=Z_2Q.
\]
Because \(\operatorname{Im}s\to+\infty\) at the cusp while \(C_3(t)\) is
bounded, the determinant of \(\operatorname{Im}Z_3(t)\) has leading term
\((\operatorname{Im}s)^3\); it is nonzero on a sufficiently small punctured
disc.  Hence (TEB.31) is a well-defined homomorphism of complex tori.

Both kernel lattices in (TEB.17) and (TEB.19) are generated by
\(k=(1,1,1)\), and \(C_3(t)k=0\).  The kernel is therefore
\[
 \mathbb C k/(\mathbb Z k+s\mathbb Z k),
\]
an elliptic curve.  Saturation of the primitive kernel proves connectedness.

On the dense torus, (TEB.30) gives
\[
\begin{aligned}
 F\Psi^{(3)}_\lambda(x,t)
 &=\bigl(\exp(2\pi iPC_3(t)\lambda)t^{P\lambda}F(x),t\bigr)\\
 &=\bigl(\exp(2\pi iC(t)Q\lambda)t^{B_0Q\lambda}F(x),t\bigr),
\end{aligned}
\]
which is (TEB.34).  Proposition~\ref{prop:teb-fan-morphism} extends the first
map over every toric stratum.

For completeness, freeness at a boundary point follows from orbit labels.  A
nonzero \(\lambda\) translates the bounded cell labelling its toric orbit to
that cell plus \(\lambda\); equality would force \(\lambda=0\).  On a smooth
fibre a fixed point would give
\(Z_3(t)\lambda\in\mathbb Z^3\), and taking imaginary parts gives
\(\operatorname{Im}Z_3(t)\lambda=0\), again forcing \(\lambda=0\).
Here is the required uniform estimate for proper discontinuity and
properness.  Put \(y=\operatorname{Im}s=-\log|t|/(2\pi)\), and, on a
punctured fibre, use
\[
 \operatorname{trop}_t(x)
 =\frac{1}{\log|t|}
   \bigl(\log|x_1|,\log|x_2|,\log|x_3|\bigr).
\]
The exact transformation law is
\[
 \operatorname{trop}_t(\Psi^{(3)}_\lambda x)
 =\operatorname{trop}_t(x)+A(t)\lambda,\qquad
 A(t)=I_3+\frac{\operatorname{Im}C_3(t)}{y}.           \tag{TEB.35a}
\]
Let
\[
 K=\sup_{|t|\leq\epsilon}
   \|\operatorname{Im}C_3(t)\|_{\mathrm{op}}<\infty.
\]
After shrinking \(\epsilon\), \(K/y\leq1/2\), and hence
\[
 \frac12\|\lambda\|
 \leq\|A(t)\lambda\|
 \leq\frac32\|\lambda\|                               \tag{TEB.35b}
\]
uniformly.  If a compact set meets its \(\lambda\)-translate in the dense
locus, its tropical coordinates lie in one bounded set \(C\), so
\(\|A(t)\lambda\|\leq\operatorname{diam}(C)\) and (TEB.35b) leaves only
finitely many integral \(\lambda\).  On the toric boundary the cell label is
translated by \(\lambda\), because the nonzero factor
\(\exp(2\pi iC_3(t)\lambda)\) does not change an orbit label; finitely many
cells again allow only finitely many \(\lambda\).

For properness, given a tropical point \(r\), choose
\(\lambda\in\mathbb Z^3\) by rounding the three coordinates of
\(A(t)^{-1}r\).  Then \(r-A(t)\lambda\) lies in
\(A(t)[0,1]^3\), a uniformly bounded set by (TEB.35b).  Over a smaller closed
disc, the closures of these representatives lie in the finite union of the
closed toric polydiscs belonging to the six Kuhn tetrahedra, exactly as in
the untwisted rank-three quotient.  That union is compact.  Thus the action
is properly discontinuous, the inverse image of every smaller closed disc
has compact quotient, and the quotient and descended map in (TEB.35) exist.
The homology assertions follow from
Theorem~\ref{thm:teb-local-lattice-bridge}.
\end{proof}

\subsection{A simultaneous unimodular splitting}
\label{subsec:teb-conjugacy}

\begin{proposition}[Candidate cusp plus one elliptic cusp]
\label{prop:teb-explicit-conjugacy}
Let
\[
 r=\begin{pmatrix}0&0&1\end{pmatrix},
\]
\[
 A=\begin{pmatrix}P\\r\end{pmatrix}
 =\begin{pmatrix}1&-1&0\\0&1&-1\\0&0&1\end{pmatrix},
 \qquad
 B=\begin{pmatrix}Q\\r\end{pmatrix}
 =\begin{pmatrix}0&-1&1\\1&-1&0\\0&0&1\end{pmatrix},
 \tag{TEB.36}
\]
and
\[
 D=B_0\oplus[1]
 =\begin{pmatrix}0&1&0\\-1&0&0\\0&0&1\end{pmatrix}.
 \tag{TEB.37}
\]
Then
\[
 \det A=\det B=1,\qquad A=DB.                           \tag{TEB.38}
\]
For \(J=\operatorname{diag}(A,B)\in SL(6,\mathbb Z)\),
\[
 JM_3J^{-1}=
 \begin{pmatrix}I_3&D\\0&I_3\end{pmatrix}.            \tag{TEB.39}
\]
In the order \((A_F,k_A;B_F,k_B)\), this is the literal matrix
\[
\begin{pmatrix}
1&0&0&0&1&0\\
0&1&0&-1&0&0\\
0&0&1&0&0&1\\
0&0&0&1&0&0\\
0&0&0&0&1&0\\
0&0&0&0&0&1
\end{pmatrix}.
\]
In the coordinates
\[
 (A_F,k_A;B_F,k_B)=(Pa,ra;Qb,rb),
\]
and after reordering them as
\[
 (A_F,B_F;k_A,k_B),
\]
the right side of (TEB.39) is
 \[
 M_F\oplus C_0,
 \qquad C_0=\begin{pmatrix}1&1\\0&1\end{pmatrix}.      \tag{TEB.40}
 \]
Explicitly, in the reordered basis it is
\[
\begin{pmatrix}
1&0&0&1&0&0\\
0&1&-1&0&0&0\\
0&0&1&0&0&0\\
0&0&0&1&0&0\\
0&0&0&0&1&1\\
0&0&0&0&0&1
\end{pmatrix}
=
\begin{pmatrix}M_F&0\\0&C_0\end{pmatrix},
\]
where \(M_F\) on the right is in the decomposed order
\((A_F;B_F)\).
Projection to the candidate block is exactly \(\Phi=P\oplus Q\).
\end{proposition}

\begin{proof}
The matrix \(A\) is upper triangular with unit diagonal.  Expanding \(B\)
along its last row gives
\[
 \det B=\det\begin{pmatrix}0&-1\\1&-1\end{pmatrix}=1.
\]
Direct multiplication gives \(DB=A\).  Hence
\[
 JM_3J^{-1}
 =\begin{pmatrix}A&0\\0&B\end{pmatrix}
  \begin{pmatrix}I&I\\0&I\end{pmatrix}
  \begin{pmatrix}A^{-1}&0\\0&B^{-1}\end{pmatrix}
 =\begin{pmatrix}I&AB^{-1}\\0&I\end{pmatrix}
 =\begin{pmatrix}I&D\\0&I\end{pmatrix}.
\]
The first two rows of \(A\) and \(B\) are \(P\) and \(Q\); their last row
records the primitive diagonal kernel coordinate.  Since \(D=B_0\oplus1\),
the asserted permutation separates the candidate and kernel blocks exactly.
\end{proof}

This proposition proves integrally, rather than only rationally, that the
rank-three cusp module is the candidate cusp module plus one standard rank-one
cusp.  It also identifies both kernel cycles with the primitive diagonal
vector \((1,1,1)\).

\subsection{The complete finite-order lattice lift}
\label{subsec:teb-finite-lift}

Define
\[
 C_1=\begin{pmatrix}-1&-1\\1&0\end{pmatrix},
 \qquad
 C_2=\begin{pmatrix}0&1\\-1&0\end{pmatrix},
 \qquad
 C_0=\begin{pmatrix}1&1\\0&1\end{pmatrix}.
 \tag{TEB.41}
\]
Direct multiplication gives
\[
 C_1^3=I,\qquad C_2^4=I,\qquad
 C_1C_2=C_0^{-1},\qquad
 (C_0-I)^2=0,\quad\operatorname{rank}(C_0-I)=1.
 \tag{TEB.42}
\]

\begin{proposition}[Rank-six lift of the \((3,4,\infty)\) peripheral actions]
\label{prop:teb-finite-monodromy-lift}
On \(\widetilde\Lambda=\Lambda_F\oplus\mathbb Z^2\), put
\[
 \widetilde A_1=A_1\oplus C_1,
 \qquad
 \widetilde A_2=A_2\oplus C_2,
 \qquad
 \widetilde M_0=M_F\oplus C_0.
 \tag{TEB.43}
\]
Then
\[
 \widetilde A_1^3=I,\qquad
 \widetilde A_2^4=I,\qquad
 \widetilde A_1\widetilde A_2\widetilde M_0=I,          \tag{TEB.44}
\]
and \(\widetilde M_0-I\) has square zero, rank three, and equal kernel and
image of rank three.  Projection
\[
 \operatorname{pr}:\widetilde\Lambda\longrightarrow\Lambda_F
\]
intertwines all three monodromies.

The affine maps
\[
 \widetilde g_j(x,y)=
 \bigl(A_jx+v_j/m_j,C_jy\bigr)
 \quad\text{on }T^4\times T^2                          \tag{TEB.45}
\]
 have orders \(m_1=3,m_2=4\) and are free.  Their quotient maps project to the
candidate finite-fibre quotient models.
\end{proposition}

\begin{proof}
Equations (TEB.8) and (TEB.42) prove (TEB.44) blockwise.  The nilpotent part
of \(\widetilde M_0\) is the direct sum of the candidate rank-two nilpotent
and the rank-one nilpotent \(C_0-I\), which proves all cusp assertions.
Block projection gives equivariance.

The vectors \((v_j,0)\) are invariant under \(\widetilde A_j\), so the
\(m_j\)-th iterate of (TEB.45) is translation by the integral vector
\((v_j,0)\), hence the identity on the torus.  If a nonidentity iterate had a
fixed point, its projection to \(T^4\) would be a fixed point of the
corresponding nonidentity iterate of (TEB.10).  The latter actions are free.
Thus (TEB.45) is free and has exact order \(m_j\).  Projection descends because
it is equivariant.
\end{proof}

Proposition~\ref{prop:teb-finite-monodromy-lift} proves that the orders
\(3,4\), affine freeness, and a rank-three cusp are mutually compatible at
the integral topological level.  It does not yet identify these block lifts
with one holomorphic period family on the full triangle-group cover.  That is
part of the remaining problem below.

\subsection{The six toric chambers and the diagonal Gram boundary}
\label{subsec:teb-weyl-chambers}

The numerical agreement between the six deepest toric strata and the six
\(A_2\) Weyl chambers has the following exact, but purely
fan-theoretic, form.  Let
\[
 \mathfrak a=
 \left\{\operatorname{diag}(\lambda_1,\lambda_2,\lambda_3):
 \lambda_1+\lambda_2+\lambda_3=0\right\}
 \subset\operatorname{Herm}_3^0(\mathbb C)
\]
and put
\[
 \nu_i=E_{ii}-\frac13I_3,\qquad
 \nu_1+\nu_2+\nu_3=0.                                  \tag{TEB.45a}
\]
Because \(n_1+n_2+n_3=0\), there is a unique real-linear isomorphism
\[
 \jmath:A_{F,\mathbb R}\xrightarrow{\ \sim\ }\mathfrak a,
 \qquad \jmath(n_i)=\nu_i.                             \tag{TEB.45b}
\]
For \(H\in\operatorname{Herm}_3^0(\mathbb C)\), let
\[
 \Phi_{\mathrm G}(H)=H-\lambda_{\min}(H)I_3.
\]
This is the \(SU(3)\)-equivariant homeomorphism onto the singular
positive-semidefinite Gram cone used in the preceding higher-rung analysis.

\begin{proposition}[The Kuhn strata are the marked \(A_2\) chambers]
\label{prop:teb-weyl-chambers}
For every \(\pi=(\pi_1,\pi_2,\pi_3)\in S_3\),
\[
 \jmath\!\left(P(\sigma_{0,\pi})\right)
 =
 \operatorname{conv}
 \{0,\nu_{\pi_1},-\nu_{\pi_3}\},                       \tag{TEB.45c}
\]
and the cone on this triangle is the closed Weyl chamber
\[
 \boxed{
 \mathbb R_{\geq0}\nu_{\pi_1}
 +\mathbb R_{\geq0}(-\nu_{\pi_3})
 =
 \{\lambda_{\pi_1}\geq\lambda_{\pi_2}
                  \geq\lambda_{\pi_3}\}.}              \tag{TEB.45d}
\]
Consequently the six tetrahedron orbits, hence the six quadruple-point
strata of the untwisted toroidal central fibre, are identified
\(S_3\)-equivariantly with the six marked diagonal Weyl chambers.

On this chamber the Gram-boundary homeomorphism is the explicit
piecewise-linear map
\[
 \Phi_{\mathrm G}(H)=H-\lambda_{\pi_3}(H)I_3.          \tag{TEB.45e}
\]
In particular,
\[
 \Phi_{\mathrm G}(\nu_{\pi_1})=E_{\pi_1\pi_1},
 \qquad
 \Phi_{\mathrm G}(-\nu_{\pi_3})=I_3-E_{\pi_3\pi_3}.    \tag{TEB.45f}
\]
\end{proposition}

\begin{proof}
Equation (TEB.22) with \(m=0\) gives the three distinct image vertices
\(0,n_{\pi_1},-n_{\pi_3}\), proving (TEB.45c).  For
\(a,b\geq0\), the diagonal entries of
\(a\nu_{\pi_1}-b\nu_{\pi_3}\) satisfy
\[
 \lambda_{\pi_1}-\lambda_{\pi_2}=a,\qquad
 \lambda_{\pi_2}-\lambda_{\pi_3}=b.
\]
Conversely these two nonnegative differences recover \(a,b\), since the
trace is zero.  This proves (TEB.45d).  Both the permutation action on the
Kuhn labels and the Weyl action on diagonal entries are the same action of
\(S_3\), so the correspondence is equivariant.  On the displayed chamber
the least eigenvalue is \(\lambda_{\pi_3}\), which proves (TEB.45e).
Substitution of the two extreme rays gives (TEB.45f).
\end{proof}

This is stronger than the bare equality \(6=|S_3|\), but it is not an
identification of singular spaces.  After quotienting the target fan by
\(A_F\), the six triangles incident to a marked vertex form two global
triangle orbits.  Moreover \(\Phi_{\mathrm G}\circ\jmath\) is a map on the
real cocharacter space, not a holomorphic map of toroidal total spaces.
It neither supplies the period correction nor changes the analytic
factorization comparison below.

\subsection{The exact Jordan determinant and Lorentz sections}
\label{subsec:teb-jordan-lorentz-sections}

The determinant--index identity in
Corollary~\ref{hrftd:cor:euler-shear-determinant} has a precise extension
inside the nine-dimensional Hermitian cone.  It is an algebraic extension of
the scalar invariant; it is not a map from the toroidal degeneration to the
Hermitian cone.

Write
\[
 \mathcal J_3=\operatorname{Herm}_3(\mathbb C),\qquad
 \mathcal J_3^{\geq0}=\{H\in\mathcal J_3:H\succeq0\},
 \qquad \Delta(H)=\det H.
\]
On the ambient vector spaces, the inclusion
\[
 \iota:\operatorname{Sym}_3(\mathbb R)\hookrightarrow\mathcal J_3,
 \qquad \iota(S)=S+i0                                      \tag{TEB.45g}
\]
is an injective unital homomorphism of real Jordan algebras for the product
\(A\circ B=(AB+BA)/2\).  Its restriction sends the positive-definite cone
into the positive-definite Hermitian cone.  The ambient formulation matters:
a positive cone is not itself a vector space or a Jordan algebra.  The map
also uses the retained marking of the integral lattice; it is not a
coordinate-free map from an unmarked congruence class of polarizations to a
fixed matrix space.

For \(a=(a_1,a_2,a_3)^{\mathsf T}\in\mathbb R^3\), fix the oriented
cross-product matrix
\[
 [a]_{\times}=
 \begin{pmatrix}
  0&-a_3&a_2\\
  a_3&0&-a_1\\
  -a_2&a_1&0
 \end{pmatrix}.                                           \tag{TEB.45h}
\]
Every \(H\in\mathcal J_3\) has one and only one decomposition
\[
 H=S+i[a]_{\times},\qquad
 S\in\operatorname{Sym}_3(\mathbb R),\quad a\in\mathbb R^3. \tag{TEB.45i}
\]

\begin{theorem}[Hermitian determinant splitting]
\label{thm:teb-hermitian-determinant-splitting}
For every real symmetric \(S\) and every \(a\in\mathbb R^3\),
\[
 \boxed{\det\bigl(S+i[a]_{\times}\bigr)
        =\det S-a^{\mathsf T}Sa.}                         \tag{TEB.45j}
\]
If \(S\) is positive definite, then
\[
 S+i[a]_{\times}\succeq0
 \quad\Longleftrightarrow\quad
 a^{\mathsf T}Sa\leq\det S,                              \tag{TEB.45k}
\]
and positive semidefiniteness and the weak scalar inequality may both be
replaced by strict positivity and a strict scalar inequality.
\end{theorem}

\begin{proof}
Direct expansion in the fixed coordinates (TEB.45h) gives
\[
 \det\bigl(S+i[a]_{\times}\bigr)=\det S-a^{\mathsf T}Sa.
\]
For the order statement let \(M=S^{-1/2}\), the positive-definite symmetric
inverse square root.  The identity
\[
 M[a]_{\times}M^{\mathsf T}
 =\det(M)[M^{-\mathsf T}a]_{\times}                       \tag{TEB.45l}
\]
follows by applying both sides to two vectors and evaluating against a third
with the fixed oriented volume form.  Hence
\[
 S^{-1/2}\bigl(S+i[a]_{\times}\bigr)S^{-1/2}
 =I_3+i[b]_{\times},
 \qquad b=\det(S)^{-1/2}S^{1/2}a,
\]
and
\[
 \lVert b\rVert^2=\frac{a^{\mathsf T}Sa}{\det S}.
\]
The three eigenvalues of \(I_3+i[b]_{\times}\) are
\(1-\lVert b\rVert,1,1+\lVert b\rVert\).  Congruence preserves inertia,
so (TEB.45k) and its strict version follow.
\end{proof}

For a fixed marked positive-definite symmetric monodromy form \(B\), define
the injective real-linear map
\[
 \ell_B:\mathbb R\oplus\mathbb R^3\longrightarrow\mathcal J_3,
 \qquad \ell_B(\lambda,a)=\lambda B+i[a]_{\times}.        \tag{TEB.45m}
\]
Theorem~\ref{thm:teb-hermitian-determinant-splitting}, applied first to
\(\lambda>0\) and then directly at \(\lambda=0\), gives
\[
 \det\ell_B(\lambda,a)
 =\lambda\bigl(\lambda^2\det B-a^{\mathsf T}Ba\bigr),    \tag{TEB.45n}
\]
and the exact inverse image
\[
 \boxed{
\ell_B^{-1}(\mathcal J_3^{\geq0})
=\{(\lambda,a):\lambda\geq0,
       \ a^{\mathsf T}Ba\leq\lambda^2\det B\}.}         \tag{TEB.45o}
\]
Indeed, if \(\lambda<0\), then
\(\operatorname{tr}\ell_B(\lambda,a)=\lambda\operatorname{tr}B<0\),
which is impossible for a positive-semidefinite matrix.
Thus every marked positive monodromy ray determines a literal future
\((3+1)\)-Lorentz cone with quadratic form
\[
 q_B(\lambda,a)=\lambda^2\det B-a^{\mathsf T}Ba.          \tag{TEB.45p}
\]
At \(B=I_3\) this is
\(\lambda^2-a_1^2-a_2^2-a_3^2\).  The equality
\[
 \chi(W_{3,I_3})=6=6q_{I_3}(1,0)                         \tag{TEB.45q}
\]
is therefore exact, while remaining an equality of scalar invariants rather
than a morphism of degenerations.

There is also an exact order-cone lift of the rank-two Lorentz cone from
Theorem~\ref{thm:hr3c-two-ladders}.  In its coordinates
\[
 H_2(t,x_1,x_2,x_3)=
 \begin{pmatrix}
  t+x_3&x_1-i x_2\\
  x_1+i x_2&t-x_3
 \end{pmatrix},
\]
put
\[
 \Theta(H_2(t,x_1,x_2,x_3))=
 \begin{pmatrix}
  t&-i x_3&-i x_2\\
  i x_3&t&-i x_1\\
  i x_2&i x_1&t
 \end{pmatrix}.                                         \tag{TEB.45r}
\]

\begin{proposition}[Order-cone lift and its complete rank behaviour]
\label{prop:teb-theta-order-lift}
The map \(\Theta\) is a real-linear bijection from
\(\operatorname{Herm}_2(\mathbb C)\) onto
\(\mathbb RI_3\oplus i\mathfrak{so}(3)\), and
\[
 H_2\succeq0\quad\Longleftrightarrow\quad
 \Theta(H_2)\succeq0.                                  \tag{TEB.45s}
\]
Moreover
\[
 \boxed{\det_3\Theta(H_2)
 =\frac{\operatorname{tr}H_2}{2}\det_2H_2.}             \tag{TEB.45t}
\]
On the positive-semidefinite cone, \(\Theta\) sends ranks
\(2,1,0\) respectively to ranks \(3,2,0\).
\end{proposition}

\begin{proof}
Put \(r=(x_1^2+x_2^2+x_3^2)^{1/2}\).  The eigenvalues of \(H_2\) are
\(t-r,t+r\), while those of \(\Theta(H_2)\) are
\(t-r,t,t+r\).  Both matrices are positive semidefinite exactly when
\(t\geq r\).  Multiplication of the displayed eigenvalues gives
\[
 \det_3\Theta(H_2)=t(t^2-r^2)
 =\frac{\operatorname{tr}H_2}{2}\det_2H_2.
\]
The same eigenvalues give all three rank assertions.  The displayed matrix
also makes real-linearity, injectivity, and the stated image immediate.
\end{proof}

If \(P\) is a rank-one projection in
\(\operatorname{Herm}_2(\mathbb C)\), then its Lorentz coordinates satisfy
\(t=r=1/2\), so \(\Theta(P)\) has eigenvalues \(0,1/2,1\) and is not an
idempotent.  Thus \(\Theta\) is not a Jordan-algebra homomorphism.  Globally,
the pullback of the cubic determinant hypersurface is
\[
 \Theta^{-1}\{\det_3=0\}
 =\{\operatorname{tr}H_2=0\}\cup\{\det_2H_2=0\};         \tag{TEB.45u}
\]
only on the positive cone does the first set add no points beyond the origin.
Neither \(\ell_B\) nor \(\Theta\) is a holomorphic toric map, a map of
integral period lattices, or a singular-fibre filling map.  For the actual
one-parameter toric degeneration the monodromy form is fixed at \(B=I_3\)
while the local parameter in \(z_0z_1z_2z_3=t\) tends to zero.  Consequently
the induced Jordan datum is the constant point \(I_3\), not a path to
\(\det B=0\).  Degeneration at fixed monodromy and rank loss of the monodromy
form are different operations.

Finally, the cyclic group \(C_3\subset S_3\) does have two orbits on the six
Kuhn labels.  Under \(P\), the even labels
\((123),(231),(312)\) give translates of the lower target triangle
\(\operatorname{conv}\{u,u+e_1,u+e_2\}\), while the odd labels
\((132),(213),(321)\) give translates of the upper target triangle
\(\operatorname{conv}\{u+e_1,u+e_2,u+e_1+e_2\}\).  This action is by
coordinate permutations and the exact bridge lowers lattice rank through
\(P\); it is not translation by an index-three superlattice.  No such
translation-lattice realization is constructed by the packet supplying
(TEB.45j)--(TEB.45t).  More decisively,
for any rank-three lattice \(\Lambda'\), a fundamental domain has normalized
volume \(3!=6\), so every \(\Lambda'\)-unimodular triangulation of that domain
has six maximal tetrahedra.  The canonical six-to-two construction is instead
the rank-lowering fan map \(P:L_3\to A_F\) in (TEB.16)--(TEB.23).  It must not
be replaced by an unproved semistable translation quotient.

\begin{auditrecord}[Completed Jordan--toric packet, 3 September 2026]
The source archive has SHA-256
\Source{8BA0BD7FD7BADEAAF2C44E2F2B4EFD42C25DA9119DA3180017FC488CA364E96F}.
Its supplied symbolic identities were replayed independently, and separate
Euler/stratification, morphism/type, fan/Weyl, and theorem-map audits were
performed.  Equations (TEB.45j)--(TEB.45t) are the surviving exact content.
The positive-cone Jordan typing, the normalization stratum attribution, and
the index-three compression were corrected rather than copied.  The audit
reports and exact replay script are retained under \Source{agent_logs/} and
\Source{agent_scratch/jordan_toric_symbolic_replay_20260903/} in the local
research record.
\end{auditrecord}

\subsection{The Gram determinant does not supply this bridge}
\label{subsec:teb-gram-obstructions}

Let
\[
 V=\operatorname{Mat}_3(\mathbb C),
 \qquad
 \mathcal G=V^+/SU(3),
 \qquad
 \mathcal D=\{\det=0\}^+/SU(3).
\]
Here \(V^+=V\cup\{\infty\}\) is the one-point compactification and
\(\{\det=0\}^+=\{\det=0\}\cup\{\infty\}\).
The exact Gram classification gives
\[
 (\mathcal G,\mathcal D)\cong(S^{10},S^8),
 \qquad
 \mathcal G\setminus\mathcal D\cong\mathbb R^9\times S^1.
 \tag{TEB.46}
\]
The real fixed form gives
\[
 \left(\operatorname{Mat}_3(\mathbb R)^+/SO(3),
       \{\det=0\}^+/SO(3)\right)
 \cong(S^6,S^5).                                       \tag{TEB.47}
\]

\begin{theorem}[No determinant-stratified identification]
\label{thm:teb-no-gram-identification}
There is no homeomorphism of pairs from the real Gram pair (TEB.47) to
\[
 \left(X,f^{-1}\{p_0,p_1,p_2\}\right).
 \tag{TEB.48}
\]
There is no homotopy equivalence from the complex Gram complement in
(TEB.46) to the candidate regular part \(f^{-1}(B^\circ)\).  The rank-one
trace-one Gram stratum \(\mathbb{CP}^2\) is not homotopy equivalent to any of
the three candidate special fibres.  Finally, no equivalence of function
germs identifies its normal determinant with any of the three candidate local
filling maps.
\end{theorem}

\begin{proof}
The determinant-zero set in (TEB.47) is a connected codimension-one
\(S^5\), and its complement has two components.  The set in (TEB.48) is the
disjoint union of three connected real four-dimensional fibres, hence has
codimension two and three components.  Its complement is connected because
it is a torus bundle over the connected pair of pants \(B^\circ\).  This
proves the first assertion.

The complex Gram complement has fundamental group \(\mathbb Z\).  The
homotopy sequence of the candidate regular torus bundle gives
\[
 1\longrightarrow\mathbb Z^4\longrightarrow
 \pi_1\bigl(f^{-1}(B^\circ)\bigr)
 \longrightarrow\pi_1(B^\circ)=F_2\longrightarrow1.    \tag{TEB.49}
\]
Its fundamental group surjects onto the nonabelian group \(F_2\), so it is
not \(\mathbb Z\).

The rank-one trace-one stratum has
\[
 H_*(\mathbb{CP}^2;\mathbb Z)=(\mathbb Z,0,\mathbb Z,0,\mathbb Z),
 \qquad \chi(\mathbb{CP}^2)=3.
\]
The cusp fibre has (TEB.6), so \(b_1=2,b_2=4,\chi=2\).  Each finite reduced
fibre is a free finite quotient of \(T^4\), hence has \(\chi=0\), and its
audited first homology has free rank two.  These invariants rule out all three
homotopy equivalences.

At a rank-one Gram matrix, Schur-complement coordinates put the normal
determinant in the exact form
\[
 uv-x^2-y^2,
 \tag{TEB.50}
\]
a nondegenerate quadratic form of signature \((1,3)\).  Its complexification
is the irreducible quadratic determinant
\[
 s_{11}s_{22}-s_{12}s_{21}.
 \tag{TEB.51}
\]
At a deepest candidate cusp point, the map is instead
\[
 (z_0,z_1,z_2)\longmapsto z_0z_1z_2,                   \tag{TEB.52}
\]
with three smooth complex branches.  The source dimensions, multiplicities,
and number of irreducible branches in (TEB.50)--(TEB.52) disagree.  The two
finite transverse maps have multiplicities three and four, while the Gram
determinant has multiplicity two along the rank-one stratum and multiplicity
three at the zero matrix.  It has no order-four stratum.  Orders of vanishing
and analytic factorization are invariant under equivalence of function germs,
which proves the final assertion.
\end{proof}

There are two further typed obstructions.  First, the determinant phase map
\[
 \arg\det:\mathcal G\setminus\mathcal D\longrightarrow S^1
\]
has trivial page monodromy after the \(SU(3)\) quotient.  Before quotienting,
its page is diffeomorphic to
\(\mathbb R_{>0}\times SL(3,\mathbb C)\), hence is simply connected, and its
geometric monodromy is the scalar central element of order three.
It supplies no \(\mathbb Z^4\)- or \(\mathbb Z^6\)-valued first-homology local
system and no order-four peripheral loop.

Second, the trace-one rank-one locus is the homogeneous space
\[
 \mathbb{CP}^2=SU(3)/S(U(1)\times U(2)).
\]
Any \(SU(3)\)-invariant map from it to the candidate base, on which no
\(SU(3)\)-action is part of the construction, is constant.  Choosing a line
\(\ell\in\mathbb{CP}^2\) does define the exact positive compression
\[
 \operatorname{Herm}_3(\mathbb C)_{\ge0}\longrightarrow
 \operatorname{Herm}(\ell^\perp)_{\ge0},
 \qquad H\longmapsto P_{\ell^\perp}H|_{\ell^\perp},     \tag{TEB.53}
\]
but its target is a four-real-dimensional Lorentz cone, not
\(\mathbb P^1\), and it contains no integral torus lattice.  An integral
quotient such as \(P\) requires the additional choice of maximal-torus and
weight-lattice data.  Irreducibility of the fundamental \(SU(3)\)
representation forbids an \(SU(3)\)-equivariant rank-two linear quotient of
\(\mathbb C^3\).

In particular, the homomorphism (TEB.14) belongs to the toroidal degeneration,
not to the Gram cone.  A generic conjugacy orbit in the trace-normalized
Hermitian cone is \(SU(3)/T^2\), which is simply connected, and the determinant
page is simply connected as well.  Neither has first homology
\(\mathbb Z^6\).  Calling any map from it a \(T^6\) filling map before
constructing a torus bundle would be a type error.

\subsection{The remaining global problem}
\label{subsec:teb-remaining}

The local bridge and the finite lattice lift isolate the remaining question.
One must construct a proper complex, or explicitly stratified-real, map
\[
 \widetilde f:\widetilde X^8\longrightarrow\mathbb P^1
\]
and a map over the same base
\[
 \Pi:\widetilde X\longrightarrow X                              \tag{TEB.54}
\]
with all of the following literal properties:
\begin{enumerate}
\item on every regular fibre, \(\Pi\) is the elliptic extension (TEB.31) and
its marked integral map is \(\Phi\);
\item over the cusp disc, \(\Pi\) is the toric morphism (TEB.35), including
the central strata and the filling square (TEB.26);
\item over each finite-order disc, the affine monodromy and translation are
those in (TEB.43)--(TEB.45), and its punctured overlap agrees with the regular
period family;
\item the actual based van Kampen maps and integral Mayer--Vietoris maps are
computed from these inclusions, rather than inferred from the displayed
monodromy matrices.
\end{enumerate}

Equations (TEB.24)--(TEB.45) prove that there is no integral monodromy or
affine-freeness obstruction to these requirements.  They do not supply the
global period and overlap construction in item 3.

There is also an unavoidable topological conclusion.  The rank-three
toroidal central fibre has Euler characteristic six.  The two finite fibres
in (TEB.45) are free finite quotients of \(T^6\), so each has Euler
characteristic zero.  If these are the only three special fibres, Euler
integration gives
\[
 \chi(\widetilde X)=6.                                  \tag{TEB.55}
\]
Therefore this direct complex globalization cannot be \(S^8\), whose Euler
characteristic is two.  Independently, \(S^8\) admits no almost-complex
structure.  A real \(S^8\) construction would have to add non-complex
singular filling data with total Euler contribution \(-4\), then pass the
full fundamental-group, integral-homology, and tangent-clutching tests.

The nine-dimensional Gram cone remains a valid independent
representation-theoretic carrier.  Relating it to (TEB.54) requires a new,
explicit morphism from determinant-rank data to period, lattice, and filling
data.  Neither the determinant, the Cartan cubic, nor the
\(\mathbb{CP}^2\)-family of normal Lorentz cones supplies such a morphism.

\section{The rank-six cusp lattice inside the nine-dimensional cone program}
\label{rsl:sec:rank-six}

The rank-three toroidal construction contains a free rank-six lattice before
any finite quotient is taken.  This section identifies that lattice, its
integral symmetric-group action, its monodromy, and the exact part of it seen
by the Hermitian Gram boundary.  It also separates it from the different
rank-six group which classifies bundle defects on the regular locus of the
reconstructed \(S^6\).

\subsection{The two rank-six groups have different types}
\label{rsl:subsec:two-lattices}

Retain the notation of \cref{subsec:teb-rank-three}.  The lattice native to
the rank-three toroidal cusp is
\begin{equation}
 \Lambda_{\rm tor}
 :=H_1(T^6;\mathbb Z)
 =L_3\oplus\overline L_3
 =\mathbb Z\langle e_1,e_2,e_3\rangle
  \oplus\mathbb Z\langle\bar e_1,\bar e_2,\bar e_3\rangle.
 \label{rsl:eq:tor-lattice}
\end{equation}
Its first summand consists of angular cycles and its second summand consists
of deck cycles.  The cusp monodromy and filling map are
\begin{equation}
 M_3(a,b)=(a+b,b),
 \qquad
 c_3(a,b)=b.
 \label{rsl:eq:monodromy-filling}
\end{equation}

The bundle-defect lattice on the original threefold has a different source
and a different cohomological degree.  With
\(D=W\sqcup S_1\sqcup S_2\) and \(U=X\setminus D\), it is
\begin{equation}
 L_{\rm def}:=H^4(U;\mathbb Z)
 \cong H_1(W;\mathbb Z)\oplus H_1(S_1;\mathbb Z)
                    \oplus H_1(S_2;\mathbb Z)
 \cong(\mathbb Z^2)^3.
 \label{rsl:eq:defect-lattice}
\end{equation}
The map \(\Phi=P\oplus Q\) below is defined on
\(\Lambda_{\rm tor}\), whereas the fibre-degree character of
\cref{zdc:thm:fibre-degree} is defined on \(L_{\rm def}\).  Equality of their
ranks supplies no map between them.  In particular, neither group is the
cyclic group \(\mathbb Z/6\mathbb Z\).

There is also no component-permuting symmetric-group action on
\(L_{\rm def}\) supplied by the marked candidate fibration.  The three
deleted fibres have, respectively,
unipotent, order-three, and order-four local monodromy.  An automorphism of
the marked fibration which permutes the punctures must preserve these
conjugacy-order types and hence fixes all three punctures.  One may impose an
abstract permutation of the three displayed summands in
\eqref{rsl:eq:defect-lattice} after choosing three identifications with
\(\mathbb Z^2\), but that is a re-marking, not an automorphism furnished by
the geometry.  This argument does not classify automorphisms fixing the base
pointwise or internal actions on the individual summands; in particular it
does not assert that the abstract group admits no \(S_3\)-action.

\subsection{The integral \texorpdfstring{\(S_3\)}{S3}-module of the cusp}
\label{rsl:subsec:s3-module}

Let \(R_\sigma\) permute the three standard basis vectors by
\(R_\sigma e_i=e_{\sigma(i)}\), and let \(S_3\) act diagonally on
\(\Lambda_{\rm tor}=L_3\oplus\overline L_3\).  Put
\begin{equation}
 k=e_1+e_2+e_3,
 \qquad
 K=\mathbb Z(k,0)\oplus\mathbb Z(0,\bar k),
 \quad \bar k=\bar e_1+\bar e_2+\bar e_3.
 \label{rsl:eq:kernel-k}
\end{equation}
The group \(S_3\) fixes \(K\) pointwise, commutes with \(M_3\), and makes
\(c_3\) equivariant for the permutation action on \(\overline L_3\):
\(c_3(R_\sigma a,R_\sigma b)=R_\sigma c_3(a,b)\).

Recall
\begin{equation}
 P=\begin{pmatrix}1&-1&0\\0&1&-1\end{pmatrix},
 \qquad
 Q=\begin{pmatrix}0&-1&1\\1&-1&0\end{pmatrix},
 \qquad
 \Phi=P\oplus Q.
 \label{rsl:eq:PQ}
\end{equation}
For \(\sigma\in S_3\), define integral actions on \(A_F\) and \(B_F\) by
\begin{equation}
 \rho_A(\sigma)(Px)=P(R_\sigma x),
 \qquad
 \rho_B(\sigma)(Qx)=Q(R_\sigma x).
 \label{rsl:eq:target-actions}
\end{equation}
These definitions are well posed because
\(\ker P=\ker Q=\mathbb Zk\).

\begin{theorem}[The exact rank-six \(S_3\)-module sequence]
\label{rsl:thm:s3-sequence}
The map \(\Phi\) fits into an exact sequence of integral
\(S_3\)-modules
\begin{equation}
 0\longrightarrow K\longrightarrow\Lambda_{\rm tor}
 \xrightarrow{\ \Phi\ }A_F\oplus B_F\longrightarrow0.
 \label{rsl:eq:s3-sequence}
\end{equation}
It is simultaneously equivariant for the cusp monodromies:
\begin{equation}
 \Phi M_3=M_F\Phi,
 \qquad
 M_3|_K=\begin{pmatrix}1&1\\0&1\end{pmatrix}
 \quad\text{in the basis }((k,0),(0,\bar k)).
 \label{rsl:eq:kernel-monodromy}
\end{equation}
As a sequence of free abelian groups with monodromy it splits integrally, as
shown by \cref{prop:teb-explicit-conjugacy}.  As a sequence of integral
\(S_3\)-modules it does not split.  After tensoring with \(\mathbb Q\),
\begin{equation}
 \Lambda_{\rm tor,\mathbb Q}\cong
 2\mathbf 1\oplus2V_{\rm std},
 \qquad
 (A_F\oplus B_F)_{\mathbb Q}\cong2V_{\rm std},
 \qquad K_{\mathbb Q}\cong2\mathbf1.
 \label{rsl:eq:rational-decomposition}
\end{equation}
\end{theorem}

\begin{proof}
The kernels of \(P\) and \(Q\) are both the primitive diagonal lattice
\(\mathbb Zk\), and both maps are surjective.  This proves exactness.
Equations \eqref{rsl:eq:target-actions} make \(P\), \(Q\), and hence
\(\Phi\), equivariant.  Since \(R_\sigma(a+b)=R_\sigma a+R_\sigma b\), the
source action commutes with \(M_3\).  Moreover
\[
 M_3(xk,y\bar k)=((x+y)k,y\bar k),
\]
which gives \eqref{rsl:eq:kernel-monodromy}.  The target equivariance with
\(M_F\) is equivalently the already verified identity
\(\Phi M_3=M_F\Phi\).

It remains to prove the integral non-splitting.  Any \(S_3\)-equivariant
homomorphism \(r:\mathbb Z^3\to\mathbb Z\), where the target is trivial, has
the same value \(a\in\mathbb Z\) on all three standard basis vectors.  Thus
\(r(k)=3a\), so no such map retracts the inclusion
\(\mathbb Zk\hookrightarrow\mathbb Z^3\).  For two copies, every entry of
the matrix obtained by restricting an equivariant map
\((\mathbb Z^3)^2\to\mathbb Z^2\) to \((\mathbb Zk)^2\) is divisible by
three.  That restriction cannot be the identity matrix.  Hence
\eqref{rsl:eq:s3-sequence} does not split over \(\mathbb Z[S_3]\).

Over \(\mathbb Q\), the invariant diagonal line has the equivariant
complement \(\{x_1+x_2+x_3=0\}\), the standard representation.  Applying
this to both copies gives \eqref{rsl:eq:rational-decomposition}.
\end{proof}

For reference, if \(s_{12}\) and \(s_{23}\) are the adjacent
transpositions, the induced target matrices are
\begin{equation}
 \rho_A(s_{12})=\begin{pmatrix}-1&0\\1&1\end{pmatrix},
 \quad
 \rho_A(s_{23})=\begin{pmatrix}1&1\\0&-1\end{pmatrix},
 \label{rsl:eq:rho-A}
\end{equation}
and
\begin{equation}
 \rho_B(s_{12})=\begin{pmatrix}1&-1\\0&-1\end{pmatrix},
 \quad
 \rho_B(s_{23})=\begin{pmatrix}-1&0\\-1&1\end{pmatrix}.
 \label{rsl:eq:rho-B}
\end{equation}
Direct multiplication gives
\begin{equation}
 \rho_A(\sigma)P=PR_\sigma,
 \qquad
 \rho_B(\sigma)Q=QR_\sigma,
 \qquad
 \rho_A(\sigma)B_0=B_0\rho_B(\sigma).
 \label{rsl:eq:rho-intertwiners}
\end{equation}

The quotient \(\mathbb Z^3/\mathbb Zk\) is the \(A_2\) weight lattice, not
the \(A_2\) root lattice.  Indeed, the augmentation-zero lattice is generated
by \(e_1-e_2\) and \(e_2-e_3\), whose images under \(P\) are the columns of
\begin{equation}
 \begin{pmatrix}2&-1\\-1&2\end{pmatrix}.
 \label{rsl:eq:root-index}
\end{equation}
Its determinant is three.  Thus the root lattice is the index-three
sublattice of \(A_F\) generated by the differences of the three weights
\(n_1,n_2,n_3\).  The integral non-splitting in
\cref{rsl:thm:s3-sequence} is the same index-three phenomenon viewed before
passing to the quotient.

\subsection{The six homology generators are not the six chambers}
\label{rsl:subsec:three-sixes}

There are three related but distinct appearances of the integer six.

\begin{theorem}[The three sixes]
\label{rsl:thm:three-sixes}
For the rank-three toroidal cusp and every injective integral shear
\(\mathcal B:\overline L_3\hookrightarrow L_3\),
\begin{align}
 \operatorname{rk}\Lambda_{\rm tor}&=2\operatorname{rk}L_3=6,
 \label{rsl:eq:six-rank}\\
 \#\{\text{Kuhn maximal cells modulo }L_3\}&=|S_3|=3!=6,
 \label{rsl:eq:six-cells}\\
 \chi(W_{3,\mathcal B})
 &=|S_3|[L_3:\mathcal B\overline L_3]
   =3!\,|\det\mathcal B|.
 \label{rsl:eq:six-euler}
\end{align}
The equalities \eqref{rsl:eq:six-cells} and \eqref{rsl:eq:six-euler} have the
same toric source: the maximal cells are indexed by orderings of three angular
coordinates.  Equation \eqref{rsl:eq:six-rank} has a different source: first
homology contains one angular and one deck copy of the rank-three lattice.
For a general shear \(\mathcal B\), this cell count does not assert that the
diagonal \(S_3\)-action descends to the sheared quotient; that requires the
additional commutation equations
\(R_\sigma\mathcal B=\mathcal B R_\sigma\).

In particular, the six standard basis vectors of \(\Lambda_{\rm tor}\) are
not \(S_3\)-equivariantly identifiable with the six Kuhn cells.  The basis
vectors form two orbits of size three; the maximal cells form one free orbit
of size six.
\end{theorem}

\begin{proof}
The rank formula follows from \eqref{rsl:eq:tor-lattice}.  A Kuhn cell in one
lattice cube is labelled by the order in which the three coordinate
increments are taken, hence by a unique element of \(S_3\).  Left
multiplication is free and transitive on these labels, proving
\eqref{rsl:eq:six-cells}.  The normalization and orbit calculation in
\cref{hrftd:cor:euler-shear-determinant} gives
\eqref{rsl:eq:six-euler}.

Under the diagonal permutation action, the angular basis
\(\{e_1,e_2,e_3\}\) and the deck basis
\(\{\bar e_1,\bar e_2,\bar e_3\}\) are two distinct transitive three-element
sets.  The cell labels form the regular six-element \(S_3\)-set.  No
equivariant bijection exists.  Equivalently, over \(\mathbb Q\) the character
of the basis permutation module on the conjugacy classes
\((1),(12),(123)\) is
\[
 (6,2,0),
\]
whereas the regular representation has character \((6,0,0)\).  The
representations are not isomorphic even over \(\mathbb Q\).
\end{proof}

This distinction also locates the chamber construction: the Kuhn fan uses the
angular summand \(L_3\), while the deck summand is needed for periods,
monodromy, and filling.  The map \(P\) produces the chamber fan; the different
map \(Q\) makes the filling square commute.  Replacing both by a single
untyped ``six-direction map'' loses precisely the information that makes the
cusp morphism valid.

\subsection{The exact map into the Gram boundary}
\label{rsl:subsec:gram-map}

Let
\[
 \mathfrak a
 =\{\operatorname{diag}(\lambda_1,\lambda_2,\lambda_3):
       \lambda_1+\lambda_2+\lambda_3=0\}
 \subset\operatorname{Herm}_3^0(\mathbb C),
\]
let \(\jmath:A_{F,\mathbb R}\to\mathfrak a\) be the marked isomorphism
\(\jmath(n_i)=E_{ii}-I_3/3\), and let
\[
 \Phi_{\rm G}(H)=H-\lambda_{\min}(H)I_3.
\]
Write \(\operatorname{pr}_{\rm ang}(a,b)=a\).  The maps which actually exist
are
\begin{equation}
 \Psi=\jmath\circ P\circ\operatorname{pr}_{\rm ang}:
 \Lambda_{\rm tor,\mathbb R}\longrightarrow\mathfrak a,
 \label{rsl:eq:linear-gram-map}
\end{equation}
and
\begin{equation}
 \Gamma=\Phi_{\rm G}\circ\Psi:
 \Lambda_{\rm tor,\mathbb R}\longrightarrow
 \{H\in\operatorname{Herm}_3(\mathbb C):H\succeq0,\ \det H=0\}.
 \label{rsl:eq:piecewise-gram-map}
\end{equation}

\begin{proposition}[Exact scope of the Gram-boundary map]
\label{rsl:prop:gram-scope}
The map \(\Psi\) is linear and \(S_3\)-equivariant, with
\begin{equation}
 \ker\Psi=\mathbb R(k,0)\oplus(0,\overline L_{3,\mathbb R}).
 \label{rsl:eq:psi-kernel}
\end{equation}
The map \(\Gamma\) is continuous, positively homogeneous,
\(S_3\)-equivariant, and piecewise linear.  On the six projected Kuhn cones it
is exactly the six diagonal Weyl-chamber map of
\cref{prop:teb-weyl-chambers}.

The map \(\Gamma\) is not additive, is not proper, ignores the complete deck
summand, and does not descend to a map from a smooth torus fibre.  Thus it is
an exact continuous piecewise-linear chamber map on a real cocharacter space
into the diagonal part of the nine-dimensional Gram carrier's
eight-dimensional boundary, not a toric fan morphism and not a morphism from
the toroidal total space to the Hermitian Gram cone.
\end{proposition}

\begin{proof}
Linearity and equivariance of \(\Psi\) follow from
\eqref{rsl:eq:rho-intertwiners} and the Weyl permutation of diagonal entries.
Since \(\jmath\) is injective and \(\ker P=\mathbb Rk\), its kernel is exactly
\eqref{rsl:eq:psi-kernel}.  The least-eigenvalue function is continuous,
positively homogeneous, and linear on each closed Weyl chamber, which proves
the corresponding assertions for \(\Gamma\).  The chamber formula is
\cref{prop:teb-weyl-chambers}.

The function \(H\mapsto H-\lambda_{\min}(H)I_3\) is not additive across
different Weyl chambers.  The kernel in \eqref{rsl:eq:psi-kernel} has positive
dimension, so \(\Gamma\) is not proper.  It ignores the deck summand by
definition.  Finally, translation by the angular lattice does not leave it
invariant: for example \(\Gamma(e_1,0)\neq\Gamma(0,0)\).  Hence it cannot
descend through the quotient of the universal cover by
\(\Lambda_{\rm tor}\).
\end{proof}

There is one further integral check on what the elliptic quotient does not
preserve.  On \(\Lambda_{\rm tor}\), define the standard alternating form
\begin{equation}
 E_3((a,b),(a',b'))=a^{\mathsf T}b'-b^{\mathsf T}a'.
 \label{rsl:eq:standard-alternating}
\end{equation}
It is invariant under \(S_3\) and \(M_3\), but on the kernel it satisfies
\begin{equation}
 E_3((k,0),(0,\bar k))=3.
 \label{rsl:eq:kernel-degree-three}
\end{equation}
Every alternating form pulled back along \(\Phi\) vanishes whenever one
argument lies in \(K\).  Therefore \(E_3\) is not the pullback of an
alternating form on \(A_F\oplus B_F\).  In particular, the holomorphic
elliptic quotient in \cref{thm:teb-holomorphic-cusp-bridge} is not a morphism
preserving this chosen principal lattice form.  This is an integral statement
about \eqref{rsl:eq:standard-alternating}; it does not assert that either
nonalgebraic torus family carries a polarization.

\begin{auditrecord}[Rank-six conclusion]
The native higher-rung object is the monodromy lattice
\(\Lambda_{\rm tor}=H_1(T^6;\mathbb Z)\), not a cyclic group.  Its exact
relation to the nine-dimensional Hermitian construction factors through the
angular quotient to the \(A_2\) weight lattice and then through the marked
Cartan/Weyl fan.  The deck half is indispensable for the period and filling
squares but invisible to the Gram-boundary chamber map.  The separate
six-coordinate group \(H^4(U;\mathbb Z)\) remains the correct classifier for
bundle defects on the original \(S^6\) regular locus; no identification of
these two rank-six groups is used.
\end{auditrecord}

\section{The chamber torsor inside the nine-dimensional matrix carrier}
\label{rnmt:sec:matrix-torsor}

The preceding section separates the two three-element basis orbits of
\(\Lambda_{\rm tor}\) from the regular six-element set of Kuhn chambers.
There is nevertheless a natural place where the regular six-set occurs.  It
is not the lattice \(\Lambda_{\rm tor}\) itself, but the off-diagonal part of
the rank-nine bilinear carrier
\[
 \mathcal M_{\mathbb Z}
 :=\operatorname{Hom}(\overline L_3,L_3)
 \cong\operatorname{Mat}_{3\times3}(\mathbb Z).
\]
This is the integral real form underlying the nine-real-dimensional
Hermitian carrier.  The distinction between the full carrier and its
positive-semidefinite cone will be essential.

Retain the marked bases \(e_1,e_2,e_3\) of \(L_3\) and
\(\bar e_1,\bar e_2,\bar e_3\) of \(\overline L_3\).  Let \(E_{ij}\) denote
the homomorphism which sends \(\bar e_j\) to \(e_i\) and annihilates the
other two basis vectors.  The diagonal \(S_3\)-action on the two rank-three
lattices induces
\begin{equation}
 \sigma\cdot X=R_\sigma X R_\sigma^{-1},
 \qquad \sigma\cdot E_{ij}=E_{\sigma(i),\sigma(j)}.
 \label{rnmt:eq:conjugation-action}
\end{equation}
Put
\begin{equation}
 \Delta=\bigoplus_{i=1}^3\mathbb ZE_{ii},
 \qquad
 \mathcal O=\bigoplus_{i\ne j}\mathbb ZE_{ij}.
 \label{rnmt:eq:diagonal-offdiagonal}
\end{equation}

\begin{theorem}[The regular chamber torsor in rank nine]
\label{rnmt:thm:regular-torsor}
There is an integral \(S_3\)-module decomposition
\begin{equation}
 \mathcal M_{\mathbb Z}=\Delta\oplus\mathcal O,
 \qquad
 \Delta\cong\mathbb Z[S_3/S_2],
 \qquad
 \mathcal O\cong\mathbb Z[S_3].
 \label{rnmt:eq:module-decomposition}
\end{equation}
If a chamber is indexed in one-line notation by
\(\pi=(\pi_1\pi_2\pi_3)\in S_3\), then
\begin{equation}
 \vartheta(c_\pi)=E_{\pi_1,\pi_3}
 \label{rnmt:eq:chamber-matrix-bijection}
\end{equation}
is an \(S_3\)-equivariant bijection from the six Kuhn chambers at a marked
vertex to the six-element basis of \(\mathcal O\).  Equivalently, these are
the six ordered mixed pairs
\begin{equation}
 \bigl((e_i,0),(0,\bar e_j)\bigr),\qquad i\ne j.
 \label{rnmt:eq:offdiagonal-pairs}
\end{equation}
They are exactly the mixed basis pairs which are orthogonal for the integral
alternating form \(E_3\), since
\begin{equation}
 E_3\bigl((e_i,0),(0,\bar e_j)\bigr)=\delta_{ij}.
 \label{rnmt:eq:mixed-pairing}
\end{equation}
The three excluded diagonal pairs \(i=j\) comprise the basis of \(\Delta\)
and are the symplectically dual mixed pairs.
\end{theorem}

\begin{proof}
Equation \eqref{rnmt:eq:diagonal-offdiagonal} partitions the nine matrix
units and is preserved by \eqref{rnmt:eq:conjugation-action}, proving the
direct-sum assertion.  The action is transitive on the three diagonal units,
and the stabilizer of \(E_{11}\) is the subgroup permuting \(2\) and \(3\);
hence \(\Delta\cong\mathbb Z[S_3/S_2]\).

The action on ordered pairs \((i,j)\) with \(i\ne j\) is free: a permutation
fixing two distinct letters fixes the remaining letter and is therefore the
identity.  It is transitive because any ordered pair of distinct letters can
be carried to any other.  Thus \(\mathcal O\) is the regular permutation
module.  For each ordered pair \((i,j)\), there is a unique remaining letter
\(k\), and hence a unique one-line permutation \((i,k,j)\).  This proves
that \eqref{rnmt:eq:chamber-matrix-bijection} is a bijection.  Moreover,
\[
 \vartheta(c_{\sigma\pi})
 =E_{\sigma(\pi_1),\sigma(\pi_3)}
 =\sigma\cdot E_{\pi_1,\pi_3},
\]
so it is equivariant.  Finally, \eqref{rnmt:eq:mixed-pairing} is immediate
from the defining formula
\(E_3((a,b),(a',b'))=a^{\mathsf T}b'-b^{\mathsf T}a'\).
\end{proof}

\subsection{The signed quotient and the exact chamber incidence}
\label{rnmt:subsec:signed-quotient}

The repeated indices on the angular and deck bases determine a marking
\(\iota:\overline L_3\to L_3\), \(\iota(\bar e_i)=e_i\).  Define
\begin{equation}
 s:A_F\oplus B_F\longrightarrow A_F,\qquad
 s(x,y)=x-B_0y,
 \qquad
 D=s\circ(P\oplus Q).
 \label{rnmt:eq:D-factorization}
\end{equation}
Since \(B_0Q=P\), this means, in the retained bases,
\begin{equation}
 D(a,b)=Pa-B_0Qb=P(a-\iota b),\qquad [D]=[P\mid-P].
 \label{rnmt:eq:D-formula}
\end{equation}
The formula is typed: \(Q\) is applied to the deck lattice and \(B_0\)
then carries its value from \(B_F\) to \(A_F\).

\begin{theorem}[Unimodular pair--Kuhn--Weyl incidence]
\label{rnmt:thm:pair-kuhn-weyl}
The map \(D\) is surjective, with saturated kernel
\begin{equation}
 \ker D
 =\mathbb Z\langle
 (e_1,\bar e_1),(e_2,\bar e_2),(e_3,\bar e_3),(k,0)\rangle,
 \qquad k=e_1+e_2+e_3.
 \label{rnmt:eq:D-kernel}
\end{equation}
For \(i\ne j\), put
\[
 U_{ij}=\mathbb Z(e_i,0)\oplus\mathbb Z(0,\bar e_j).
\]
Then \(U_{ij}\) is primitive and isotropic for \(E_3\), and
\begin{equation}
 D|_{U_{ij}}:U_{ij}\xrightarrow{\ \sim\ }A_F,\qquad
 (e_i,0)\longmapsto n_i,\quad(0,\bar e_j)\longmapsto-n_j
 \label{rnmt:eq:unimodular-plane}
\end{equation}
is an integral isomorphism.

Let \(k_{ij}\) be the remaining index and
\(\pi_{ij}=(i,k_{ij},j)\).  Under the Cartan marking
\(\jmath(n_i)=\nu_i=E_{ii}-I_3/3\), the positive quadrant in \(U_{ij,\mathbb
R}\) maps isomorphically to the Weyl chamber
\begin{equation}
 \mathcal C_{ij}
 =\mathbb R_{\geq0}\nu_i+\mathbb R_{\geq0}(-\nu_j)
 =\{\operatorname{diag}(\lambda):
 \lambda_i\geq\lambda_{k_{ij}}\geq\lambda_j,\
 \textstyle\sum_r\lambda_r=0\}.
 \label{rnmt:eq:chamber-inequalities}
\end{equation}
The assignments
\begin{equation}
 ((e_i,0),(0,\bar e_j))
 \longmapsto \pi_{ij}
 \longmapsto[\sigma_{0,\pi_{ij}}]
 \longmapsto\mathcal C_{ij}
 \label{rnmt:eq:torsor-chain}
\end{equation}
are equivariant bijections of \(S_3\)-torsors, and
\begin{equation}
 P(\sigma_{0,\pi_{ij}})
 =\operatorname{conv}\{0,n_i,-n_j\}.
 \label{rnmt:eq:projected-kuhn}
\end{equation}
The vector
\begin{equation}
 D(e_i,\bar e_j)=n_i-n_j
 \label{rnmt:eq:interior-root}
\end{equation}
is the interior \(A_2\)-root bisector of this chamber.  With the source and
target orientations fixed by the displayed bases,
\begin{equation}
 \det[e_i\mid e_i+e_{k_{ij}}\mid k]
 =\operatorname{sgn}(\pi_{ij})
 =\det[n_i\mid -n_j].
 \label{rnmt:eq:orientation-sign}
\end{equation}
\end{theorem}

\begin{proof}
Since \(P\) is surjective, so is \(D\).  Equation
\eqref{rnmt:eq:D-formula} gives \(D(a,b)=0\) precisely when
\(a-\iota b=tk\) for a unique \(t\in\mathbb Z\).  Hence
\[
 (a,b)=\sum_{r=1}^3 b_r(e_r,\bar e_r)+t(k,0),
\]
and uniqueness gives \eqref{rnmt:eq:D-kernel}.  Surjectivity implies that
the kernel is saturated.

The plane \(U_{ij}\) is spanned by two members of a basis of
\(\Lambda_{\rm tor}\), hence is primitive.  Equation
\eqref{rnmt:eq:mixed-pairing} makes it isotropic exactly when \(i\ne j\).
In its displayed basis the restriction of \(D\) has columns
\([n_i\mid-n_j]\).  Directly from
\[
 n_1=\binom10,\qquad n_2=\binom{-1}{1},\qquad
 n_3=\binom0{-1},
\]
one obtains
\(\det[n_i\mid-n_j]=\pm1\) for \(i\ne j\), proving
\eqref{rnmt:eq:unimodular-plane}.

For \(x,y\geq0\), the diagonal matrix
\(H=x\nu_i-y\nu_j\) satisfies
\[
 \lambda_i-\lambda_{k_{ij}}=x,\qquad
 \lambda_{k_{ij}}-\lambda_j=y.
\]
Conversely these two differences, together with trace zero, recover \(H\),
which proves \eqref{rnmt:eq:chamber-inequalities} and the cone
isomorphism.  The four vertices of
\(\sigma_{0,(i,k_{ij},j)}\) have \(P\)-images
\[
 0,\quad n_i,\quad n_i+n_{k_{ij}}=-n_j,\quad
 n_1+n_2+n_3=0,
\]
proving \eqref{rnmt:eq:projected-kuhn}.  Ordered distinct pairs, one-line
permutations, Kuhn-cell translation classes, and Weyl chambers are all free
transitive \(S_3\)-sets, and the displayed assignments commute with the
coordinate action.  At \(x=y=1\),
\(\jmath(n_i-n_j)=E_{ii}-E_{jj}\), whose ordered diagonal entries in
positions \((i,k_{ij},j)\) are \((1,0,-1)\); this is the interior bisector.

Finally, relative to the ordered basis
\((e_i,e_{k_{ij}},e_j)\), the columns
\((e_i,e_i+e_{k_{ij}},k)\) form an upper unitriangular matrix.  Their
determinant in the fixed source basis is therefore
\(\operatorname{sgn}(\pi_{ij})\).  For \((i,k_{ij},j)=(1,2,3)\),
\([n_i\mid-n_j]=I_2\); equivariance and
\(\det\rho_A(\sigma)=\operatorname{sgn}(\sigma)\) prove the second equality
in \eqref{rnmt:eq:orientation-sign}.
\end{proof}

\begin{proposition}[Exact marking and monodromy boundary]
\label{rnmt:prop:marking-boundary}
There is no endomorphism \(T:A_F\to A_F\) satisfying
\begin{equation}
 DM_3=TD.
 \label{rnmt:eq:no-monodromy-intertwiner}
\end{equation}
Moreover, for every \(c\in\mathbb Z\), the shear
\begin{equation}
 S_c(a,b)=(a+c\iota b,b)
 \label{rnmt:eq:marking-shear}
\end{equation}
is an integral automorphism which fixes the angular lattice and its Kuhn
fan, preserves \(c_3\) and \(E_3\), and commutes with \(M_3\) and the
diagonal \(S_3\)-action, but
\begin{equation}
 DS_c(a,b)=D(a,b)+cP(\iota b).
 \label{rnmt:eq:D-shear}
\end{equation}
Thus the pair--chamber correspondence is canonical relative to the retained
angular/deck bases, their common indexing, \(P,Q,B_0\), and the sign in
\eqref{rnmt:eq:D-factorization}; it is not an invariant of the unmarked data
\((M_3,E_3,c_3)\) together with the angular fan.
\end{proposition}

\begin{proof}
The vector \((e_i,\bar e_i)\) lies in \(\ker D\), whereas
\[
 M_3(e_i,\bar e_i)=(2e_i,\bar e_i),\qquad
 D(2e_i,\bar e_i)=n_i\ne0.
\]
Thus \(\ker D\) is not \(M_3\)-invariant, which is necessary for
\eqref{rnmt:eq:no-monodromy-intertwiner}.

The inverse of \(S_c\) is \(S_{-c}\).  Direct substitution gives
\[
 S_cM_3(a,b)=(a+(c+1)\iota b,b)=M_3S_c(a,b),
\]
and proves preservation of \(c_3\) and commutation with the diagonal
coordinate action.  Symmetry of the dot product gives
\[
 E_3(S_c(a,b),S_c(a',b'))
 =(a+c\iota b)^{\mathsf T}b'
   -b^{\mathsf T}(a'+c\iota b')
 =E_3((a,b),(a',b')).
\]
Equation \eqref{rnmt:eq:D-shear} follows from
\eqref{rnmt:eq:D-formula}.  For \(c\ne0\), the deck complement is moved, so
the retained marked pair set and \(D\) are changed.
\end{proof}

The parity and opposition operations also retain exact information.  The
even and odd chamber triples are
\begin{align}
 \vartheta\{c_{123},c_{231},c_{312}\}
   &=\{E_{13},E_{21},E_{32}\},
 \label{rnmt:eq:even-cycle}\\
 \vartheta\{c_{132},c_{213},c_{321}\}
   &=\{E_{12},E_{23},E_{31}\}.
 \label{rnmt:eq:odd-cycle}
\end{align}
Thus they are the two orientations of the directed three-cycle.  Reversal of
one-line notation, \((i,k,j)\mapsto(j,k,i)\), is carried by \(\vartheta\) to
matrix transposition \(E_{ij}\mapsto E_{ji}\).  This is the exact chamber
opposition involution; it exchanges the two parity triples.

\subsection{The real Hermitian representation and its integral defect}

Let \(\operatorname{Herm}_3(\mathbb C)\) carry the \(S_3\)-action by
conjugation with the real permutation matrices.  Define
\begin{equation}
 \Xi:\operatorname{Mat}_{3\times3}(\mathbb R)
       \longrightarrow\operatorname{Herm}_3(\mathbb C),
 \qquad
 \Xi(X)=\frac{X+X^{\mathsf T}}2
       +\frac{i(X-X^{\mathsf T})}2.
 \label{rnmt:eq:Xi}
\end{equation}

\begin{proposition}[Exact representation bridge and the positivity boundary]
\label{rnmt:prop:representation-bridge}
The map \(\Xi\) is an \(S_3\)-equivariant real-linear isomorphism.  Its
inverse sends a Hermitian matrix \(H=S+iA\), with \(S\) real symmetric and
\(A\) real skew-symmetric, to \(S+A\).  Consequently
\begin{equation}
 \operatorname{Herm}_3(\mathbb C)
 \cong_{\mathbb R[S_3]}
 2\mathbin{\cdot}\mathbf 1\oplus
 3\mathbin{\cdot}V_{\rm std}\oplus\mathbf{sgn}.
 \label{rnmt:eq:hermitian-character}
\end{equation}

For each unordered pair \(\{i,j\}\), the symmetric and skew generators
\(E_{ij}+E_{ji}\) and \(E_{ij}-E_{ji}\) span an index-two sublattice of
\(\mathbb ZE_{ij}\oplus\mathbb ZE_{ji}\).  Hence the symmetric--skew
off-diagonal splitting has total integral index
\begin{equation}
 2^3=8.
 \label{rnmt:eq:index-eight}
\end{equation}
In particular, \(\Xi\) is a real representation isomorphism, not an
identification of the two preferred integral lattices.

Moreover, for \(i\ne j\), \(\Xi(E_{ij})\) has eigenvalues
\(2^{-1/2},-2^{-1/2},0\).  Thus the regular chamber torsor is located in the
nine-dimensional Hermitian representation, but its matrix-unit image is not
contained in the positive-semidefinite Gram cone or its boundary.
\end{proposition}

\begin{proof}
Every real matrix has the unique decomposition
\(X=S+A\), where
\(S=(X+X^{\mathsf T})/2\) is symmetric and
\(A=(X-X^{\mathsf T})/2\) is skew-symmetric.  Equation
\eqref{rnmt:eq:Xi} sends this pair to \(S+iA\), which is Hermitian, and the
displayed inverse proves bijectivity.  Permutation conjugation commutes with
transpose and with these two projections, proving equivariance.

The real permutation representation on three letters is
\(\mathbf1\oplus V_{\rm std}\).  Therefore
\[
 \operatorname{Mat}_{3\times3}(\mathbb R)
 \cong(\mathbf1\oplus V_{\rm std})^{\otimes2}
 \cong2\mathbin{\cdot}\mathbf 1\oplus
       3\mathbin{\cdot}V_{\rm std}\oplus\mathbf{sgn},
\]
where
\(V_{\rm std}\otimes V_{\rm std}
 \cong\mathbf1\oplus V_{\rm std}\oplus\mathbf{sgn}\).
For one unordered pair, the change-of-basis matrix from
\((E_{ij},E_{ji})\) to
\((E_{ij}+E_{ji},E_{ij}-E_{ji})\) has determinant \(-2\).  Multiplication
over the three unordered pairs proves \eqref{rnmt:eq:index-eight}.

Finally, on the \(i,j\) coordinate plane, \(\Xi(E_{ij})\) has block
\[
 \begin{pmatrix}0&(1+i)/2\\(1-i)/2&0\end{pmatrix},
\]
whose eigenvalues are \(\pm |1+i|/2=\pm2^{-1/2}\); the remaining coordinate
has eigenvalue zero.  Hence it is indefinite.
\end{proof}

\subsection{The positive six-set is a different equivariant realization}

Although the matrix units themselves are indefinite, the spectral-floor map
of the preceding section gives a second realization of the same chamber
torsor in the rank-two positive-semidefinite boundary.  If
\(\pi=(i,k,j)\), then
\begin{equation}
 G_\pi=2E_{ii}+E_{kk}
 \label{rnmt:eq:positive-six-set}
\end{equation}
has diagonal entries equal to a permutation of \((2,1,0)\).  The assignment
\begin{equation}
 E_{ij}\longmapsto 2E_{ii}+E_{kk},
 \qquad \{i,j,k\}=\{1,2,3\},
 \label{rnmt:eq:offdiagonal-positive-map}
\end{equation}
is therefore an equivariant bijection of the two marked six-element sets.
It extends uniquely by linearity to a homomorphism
\(\mathcal O\to\Delta\), but that extension is neither the map \(\Xi\) nor
the nonlinear spectral-floor map on all of
\(\operatorname{Herm}_3^0(\mathbb C)\).  Thus
\eqref{rnmt:eq:offdiagonal-positive-map} records exact chamber incidence; it
does not manufacture a positive monodromy-compatible morphism from the
toroidal family.

\begin{auditrecord}[What the native \(\mathbb Z^6\) result becomes]
The six Kuhn chambers are not the six generators of
\(H_1(T^6;\mathbb Z)\).  They are canonically, relative to the retained
marking, the six symplectically orthogonal ordered angular--deck pairs and
hence the six off-diagonal basis elements in the associated rank-nine
bilinear carrier.  This places the regular \(S_3\)-torsor exactly inside the
nine-dimensional program while preserving the different roles of angular
cycles, deck cycles, monodromy, and the positive-semidefinite Gram boundary.
The construction is a representation and incidence theorem; positivity and
global toroidal descent remain separate questions and are not inferred from
the equality of dimensions or orbit cardinalities.
\end{auditrecord}

\section{The Lorentz observation morphism and the full nine-dimensional fibre}
\label{hlob:sec:observation-morphism}

The rank-three Hermitian carrier and the four-real-dimensional Lorentz cone
are not identical objects.  The required next step is therefore to compute
their exact relation.  This section constructs a split positive observation
morphism, determines every one of its fibres, identifies the five retained
directions beyond the observed four-vector, and proves the complete cubic
coupling among the resulting Lorentz-vector, spinor, and scalar coordinates.
No direction is removed from the source calculation.

\subsection{The marked angular epimorphism and its exact splitting}
\label{hlob:subsec:marked-splitting}

Retain the angular lattice matrix and its primitive kernel vector:
\begin{equation}
 P=
 \begin{pmatrix}
  1&-1&0\\
  0& 1&-1
 \end{pmatrix},
 \qquad
 u=\begin{pmatrix}1\\1\\1\end{pmatrix},
 \qquad Pu=0.
 \label{hlob:eq:P-u}
\end{equation}
The standard Hermitian products in these displayed marked bases give
\begin{equation}
 C:=PP^*
 =\begin{pmatrix}2&-1\\-1&2\end{pmatrix},
 \qquad
 C^{-1}=\frac13\begin{pmatrix}2&1\\1&2\end{pmatrix},
 \qquad \det C=3.
 \label{hlob:eq:cartan}
\end{equation}
Define the corresponding orthogonal right inverse and the projection onto
the displayed kernel line by
\begin{equation}
 R:=P^*C^{-1}
 =\frac13
 \begin{pmatrix}
  2& 1\\
 -1& 1\\
 -1&-2
 \end{pmatrix},
 \qquad
 \Pi:=\frac{uu^*}{3},
 \qquad Q_u:=I_3-\Pi.
 \label{hlob:eq:R-Pi}
\end{equation}
The choice of \(R\) uses the displayed marking and its standard Hermitian
products.  All formulas below retain that choice explicitly.

\begin{lemma}[Exact splitting identities]
\label{hlob:lem:splitting-identities}
The matrices in \eqref{hlob:eq:P-u}--\eqref{hlob:eq:R-Pi} satisfy
\begin{equation}
 PR=I_2,\qquad u^*R=0,\qquad R^*R=C^{-1},\qquad
 RP=Q_u,\qquad \Pi R=0,\qquad P\Pi=0.
 \label{hlob:eq:splitting-identities}
\end{equation}
Moreover, with
\begin{equation}
 X=(R\;u)
 =\begin{pmatrix}
  2/3& 1/3&1\\
 -1/3& 1/3&1\\
 -1/3&-2/3&1
 \end{pmatrix},
 \label{hlob:eq:X}
\end{equation}
one has
\begin{equation}
 \det X=1,\qquad
 X^{-1}=\begin{pmatrix}P\\u^*/3\end{pmatrix}.
 \label{hlob:eq:X-inverse}
\end{equation}
\end{lemma}

\begin{proof}
Direct multiplication, with every denominator retained, gives
\[
 PR
 =\frac13
 \begin{pmatrix}1&-1&0\\0&1&-1\end{pmatrix}
 \begin{pmatrix}2&1\\-1&1\\-1&-2\end{pmatrix}
 =I_2
\]
and
\[
 RP
 =\frac13
 \begin{pmatrix}
  2&-1&-1\\
 -1& 2&-1\\
 -1&-1& 2
 \end{pmatrix}
 =I_3-\frac13
 \begin{pmatrix}1&1&1\\1&1&1\\1&1&1\end{pmatrix}
 =Q_u.
\]
The identities involving \(u\), \(R^*R\), and \(\Pi\) follow from the same
displayed matrices and from \(R=P^*(PP^*)^{-1}\).  Finally,
\[
 \det(3X)
 =\det\begin{pmatrix}2&1&3\\-1&1&3\\-1&-2&3\end{pmatrix}
 =27,
\]
so \(\det X=27/3^3=1\).  Both products
\[
 \begin{pmatrix}P\\u^*/3\end{pmatrix}(R\;u)
 \quad\text{and}\quad
 (R\;u)\begin{pmatrix}P\\u^*/3\end{pmatrix}
\]
are the identity by \eqref{hlob:eq:splitting-identities}; this proves
\eqref{hlob:eq:X-inverse}.
\end{proof}

\subsection{A positive observation morphism with all fibres retained}
\label{hlob:subsec:observation}

Put
\[
 \mathcal J_3=\operatorname{Herm}_3(\mathbb C),\qquad
 \mathcal J_2=\operatorname{Herm}_2(\mathbb C),\qquad
 \mathcal P_r=\{H\in\mathcal J_r:H\succeq0\}.
\]
Define real-linear maps
\begin{equation}
 \mathcal O_P:\mathcal J_3\longrightarrow\mathcal J_2,
 \qquad \mathcal O_P(H)=PHP^*,
 \label{hlob:eq:observation}
\end{equation}
and
\begin{equation}
 \jmath_P:\mathcal J_2\longrightarrow\mathcal J_3,
 \qquad \jmath_P(G)=RGR^*.
 \label{hlob:eq:section}
\end{equation}

\begin{theorem}[Split observation and the exact five-dimensional kernel]
\label{hlob:thm:split-observation}
The observation morphism is surjective and positive, and its displayed
section is a positive linear embedding:
\begin{equation}
 \mathcal O_P\jmath_P=\operatorname{id}_{\mathcal J_2},
 \qquad
 \jmath_P\mathcal O_P(H)=Q_uHQ_u.
 \label{hlob:eq:split-observation}
\end{equation}
Consequently
\begin{equation}
 \mathcal J_3
 =\operatorname{im}\jmath_P\oplus\ker\mathcal O_P,
 \qquad
 \dim_{\mathbb R}\operatorname{im}\jmath_P=4,
 \qquad
 \dim_{\mathbb R}\ker\mathcal O_P=5.
 \label{hlob:eq:dimension-split}
\end{equation}
The kernel is explicitly
\begin{equation}
 \ker\mathcal O_P
 =\{a\,uu^*+uv^*+vu^*:a\in\mathbb R,\ v\in u^\perp\}.
 \label{hlob:eq:linear-kernel}
\end{equation}
It is exactly the tangent space at \(H_0=uu^*\) to the nonzero rank-one
positive-semidefinite stratum
\begin{equation}
 \mathcal R_1=\{ww^*:w\in\mathbb C^3\setminus\{0\}\}
 \subset\mathcal P_3.
 \label{hlob:eq:rank-one-stratum}
\end{equation}
On the positive cone, the complete zero fibre is the retained ray
\begin{equation}
 \mathcal O_P^{-1}(0)\cap\mathcal P_3
 =\{\lambda uu^*: \lambda\geq0\}.
 \label{hlob:eq:positive-zero-fibre}
\end{equation}
\end{theorem}

\begin{proof}
Equations \eqref{hlob:eq:splitting-identities} give
\[
 \mathcal O_P(\jmath_P(G))=PRGR^*P^*=G
\]
and
\[
 \jmath_P(\mathcal O_P(H))=RPH P^*R^*=Q_uHQ_u.
\]
Both maps preserve positive semidefiniteness because they are congruences.
The first identity proves surjectivity, injectivity of \(\jmath_P\), and the
direct sum in \eqref{hlob:eq:dimension-split}.  Since
\(\dim_{\mathbb R}\mathcal J_3=9\) and
\(\dim_{\mathbb R}\mathcal J_2=4\), the dimensions follow.

The condition \(PHP^*=0\) is equivalent to \(Q_uHQ_u=0\).  Relative to
\(\mathbb C^3=u^\perp\oplus\mathbb Cu\), such a Hermitian operator has a
zero \(u^\perp\)-to-\(u^\perp\) block and therefore has the unique form
\eqref{hlob:eq:linear-kernel}.  The differential at \(u\) of
\(F(w)=ww^*\) is
\[
 dF_u(\dot w)=u\dot w^*+\dot w u^*.
\]
Writing \(\dot w=(a/2)u+v+i b u\), with
\(a,b\in\mathbb R\) and \(v\in u^\perp\), shows that its image is exactly
\eqref{hlob:eq:linear-kernel}; the imaginary radial direction \(ibu\) is
the one-real-dimensional kernel of \(dF_u\).  Thus the image has dimension
five and is the stated tangent space.

Finally let \(H\succeq0\) and \(PHP^*=0\).  For every \(\xi\in\mathbb C^2\),
\[
 0=\xi^*PHP^*\xi
   =\lVert H^{1/2}P^*\xi\rVert^2.
\]
Hence \(H^{1/2}\) annihilates
\(\operatorname{im}P^*=u^\perp\).  The range of \(H\) is contained in
\(\mathbb Cu\), and Hermitian positivity then gives
\(H=\lambda uu^*\) for one \(\lambda\geq0\).  The converse follows from
\(Pu=0\).
\end{proof}

The observation morphism is not a dimension count.  It factors through the
already proved two-column Gram carrier.  For \(H\in\mathcal P_3\), put
\begin{equation}
 \mathcal Z_P(H)=H^{1/2}P^*
 \in\operatorname{Mat}_{3\times2}(\mathbb C).
 \label{hlob:eq:gram-factor}
\end{equation}
Then, with no omitted scalar,
\begin{equation}
 \mathcal Z_P(H)^*\mathcal Z_P(H)=PHP^*=\mathcal O_P(H).
 \label{hlob:eq:gram-factor-identity}
\end{equation}

\subsection{Complete block coordinates and the cubic coupling}
\label{hlob:subsec:block-coordinates}

\begin{theorem}[Lorentz-vector--spinor--scalar coordinates]
\label{hlob:thm:block-coordinates}
Every \(H\in\mathcal J_3\) has a unique expression
\begin{equation}
 H=X
 \begin{pmatrix}
  G&\psi\\
  \psi^*&\lambda
 \end{pmatrix}X^*
 =RGR^*+R\psi u^*+u\psi^*R^*+\lambda uu^*,
 \label{hlob:eq:block-expression}
\end{equation}
where
\begin{equation}
 G=PHP^*,\qquad
 \psi=\frac13PHu,\qquad
 \lambda=\frac19u^*Hu,
 \label{hlob:eq:block-inverse}
\end{equation}
with \(G\in\mathcal J_2\), \(\psi\in\mathbb C^2\), and
\(\lambda\in\mathbb R\).  Thus the displayed marking gives the exact
real-vector-space decomposition
\begin{equation}
 \mathcal J_3\xrightarrow{\ \sim\ }
 \mathcal J_2\oplus\mathbb C^2\oplus\mathbb R,
 \qquad 9=4+4+1.
 \label{hlob:eq:four-four-one}
\end{equation}
For
\(G=\begin{psmallmatrix}a&b\\\bar b&d\end{psmallmatrix}\), retain
\begin{equation}
 \operatorname{adj}(G)
 =\begin{pmatrix}d&-b\\-\bar b&a\end{pmatrix}.
 \label{hlob:eq:adjugate}
\end{equation}
Then the complete determinant is
\begin{equation}
 \boxed{\det H
 =\lambda\det G-\psi^*\operatorname{adj}(G)\psi.}
 \label{hlob:eq:cubic-coupling}
\end{equation}

For \(G\succeq0\), let \(G^+\) be its Moore--Penrose inverse on
\(\operatorname{im}G\), extended by zero on \(\ker G\).  The complete
positive fibre is
\begin{equation}
 \begin{split}
 \mathcal O_P^{-1}(G)\cap\mathcal P_3
 =\biggl\{X
 \begin{pmatrix}G&\psi\\\psi^*&\lambda\end{pmatrix}X^*:
 &\ \psi\in\operatorname{im}G,\\[-1mm]
 &\ \lambda\geq\psi^*G^+\psi\biggr\}.
 \end{split}
 \label{hlob:eq:complete-positive-fibre}
\end{equation}
In particular,
\begin{equation}
 H\succ0
 \quad\Longleftrightarrow\quad
 G\succ0\ \text{ and }\ 
 \lambda>\psi^*G^{-1}\psi.
 \label{hlob:eq:positive-definite-condition}
\end{equation}
If \(G\) is singular and \(H\succeq0\), then \(H\) is singular.  If
\(G\succ0\), the boundary of its fibre is the graph
\(\lambda=\psi^*G^{-1}\psi\).
\end{theorem}

\begin{proof}
Equation \eqref{hlob:eq:X-inverse} gives
\[
 X^{-1}HX^{-*}
 =\begin{pmatrix}
   PHP^*&PHu/3\\
   u^*HP^*/3&u^*Hu/9
  \end{pmatrix},
\]
which proves uniqueness, \eqref{hlob:eq:block-expression}, and
\eqref{hlob:eq:block-inverse}.  Because \(\det X=1\), direct expansion of
the \(3\times3\) block determinant gives
\[
 \det H
 =\lambda(ad-|b|^2)
  -\begin{pmatrix}\bar\psi_1&\bar\psi_2\end{pmatrix}
   \begin{pmatrix}d&-b\\-\bar b&a\end{pmatrix}
   \begin{pmatrix}\psi_1\\\psi_2\end{pmatrix},
\]
which is \eqref{hlob:eq:cubic-coupling}.

Invertible congruence by \(X\) preserves positivity.  Suppose the block
matrix is positive semidefinite.  Its upper-left principal block gives
\(G\succeq0\).  If \(v\in\ker G\), positivity of the restriction to
\(\mathbb C(v,0)\oplus\mathbb C(0,1)\) forces \(v^*\psi=0\); hence
\(\psi\perp\ker G\), or equivalently
\(\psi\in\operatorname{im}G\).  Write \(\psi=G^{1/2}w\) with
\(w\in\operatorname{im}G\).  Completing the square gives
\[
 \begin{pmatrix}x\\s\end{pmatrix}^*
 \begin{pmatrix}G&\psi\\\psi^*&\lambda\end{pmatrix}
 \begin{pmatrix}x\\s\end{pmatrix}
 =\lVert G^{1/2}x+s w\rVert^2
  +|s|^2(\lambda-\lVert w\rVert^2).
\]
Here \(\lVert w\rVert^2=\psi^*G^+\psi\).  This proves necessity and
sufficiency of \eqref{hlob:eq:complete-positive-fibre}.  The strict case is
the ordinary Schur-complement criterion
\eqref{hlob:eq:positive-definite-condition}.

If \(G\) is singular, then \(\psi\in\operatorname{im}G\) and
\(\operatorname{adj}(G)\) has image in \(\ker G\).  Thus
\(\psi^*\operatorname{adj}(G)\psi=0=\det G\), so
\eqref{hlob:eq:cubic-coupling} gives \(\det H=0\).  If \(G\succ0\), the
same formula becomes
\[
 \det H=\det G\,(\lambda-\psi^*G^{-1}\psi),
\]
which proves the last assertion.
\end{proof}

\subsection{The embedded Lorentz group and the meaning of
\texorpdfstring{\(9=4+4+1\)}{9=4+4+1}}
\label{hlob:subsec:lorentz-action}

For \(g\in GL_2(\mathbb C)\), define
\begin{equation}
 T_P(g)=RgP+\Pi
 =X\begin{pmatrix}g&0\\0&1\end{pmatrix}X^{-1}.
 \label{hlob:eq:Tg}
\end{equation}

\begin{theorem}[Exact Lorentz-vector, Weyl-spinor, and scalar action]
\label{hlob:thm:lorentz-action}
The map \(T_P:GL_2(\mathbb C)\to GL_3(\mathbb C)\) is an injective group
homomorphism satisfying
\begin{equation}
 T_P(g)u=u,\qquad
 PT_P(g)=gP,\qquad
 T_P(g)R=Rg,\qquad
 \det T_P(g)=\det g.
 \label{hlob:eq:Tg-identities}
\end{equation}
Under congruence by \(T_P(g)\), the block coordinates transform exactly as
\begin{equation}
 (G,\psi,\lambda)\longmapsto
 (gGg^*,g\psi,\lambda).
 \label{hlob:eq:block-action}
\end{equation}
Consequently
\begin{equation}
 \mathcal O_P(T_P(g)HT_P(g)^*)
 =g\mathcal O_P(H)g^*,
 \qquad
 T_P(g)\jmath_P(G)T_P(g)^*=\jmath_P(gGg^*).
 \label{hlob:eq:equivariance}
\end{equation}

For \(g\in SL_2(\mathbb C)\), congruence on \(\mathcal J_2\) is the standard
double-cover action on Minkowski space.  In the retained coordinates
\begin{equation}
 G(t,x_1,x_2,x_3)=
 \begin{pmatrix}
  t+x_3&x_1-i x_2\\
  x_1+i x_2&t-x_3
 \end{pmatrix},
 \label{hlob:eq:pauli-coordinates}
\end{equation}
one has
\begin{equation}
 \det G=t^2-x_1^2-x_2^2-x_3^2,
 \qquad
 G\succeq0\Longleftrightarrow
 t\geq0,\quad t^2\geq x_1^2+x_2^2+x_3^2.
 \label{hlob:eq:lorentz-cone}
\end{equation}
The induced Lorentz-vector homomorphism
\begin{equation}
 \rho_{\mathrm{vec}}:SL_2(\mathbb C)\longrightarrow
 GL(\mathcal J_2),\qquad
 \rho_{\mathrm{vec}}(g)(G)=gGg^*,
 \label{hlob:eq:rho-vector}
\end{equation}
has kernel \(\{I_2,-I_2\}\) and image \(SO^+(1,3)\).  The embedding
\(T_P\), and hence the action on the full \(4+4+1\) carrier, is faithful.
Indeed, on the full nine-dimensional carrier
\(-I_2\) fixes \(G\) and \(\lambda\) but sends
\(\psi\mapsto-\psi\).  Thus the full carrier retains the spin double cover,
and \eqref{hlob:eq:four-four-one} is exactly the real representation sum
\begin{equation}
 \boxed{
 \text{Lorentz four-vector}\ \oplus\
 \text{defining complex spin representation }\mathbb C^2\ \oplus\
 \text{real scalar}.}
 \label{hlob:eq:representation-sum}
\end{equation}
With the convention that \(g\) acts on undotted spinor indices, the middle
summand is called a left Weyl spinor; reversing that convention reverses the
chirality name but not the displayed action \(\psi\mapsto g\psi\).
\end{theorem}

\begin{proof}
The second expression in \eqref{hlob:eq:Tg} and
\eqref{hlob:eq:X-inverse} prove the homomorphism, determinant, and
injectivity assertions.  Equations \eqref{hlob:eq:splitting-identities}
give the other identities directly.  Moreover,
\[
 T_P(g)X=X(g\oplus1),
\]
so congruence of \eqref{hlob:eq:block-expression} gives
\eqref{hlob:eq:block-action}.  The two equivariance identities follow.

For \(g\in SL_2(\mathbb C)\),
\[
 \det(gGg^*)=|\det g|^2\det G=\det G,
\]
and positive semidefiniteness is preserved.  Hence congruence defines a
homomorphism to the identity-component (proper orthochronous) Lorentz group of
\eqref{hlob:eq:lorentz-cone}.  If \(gGg^*=G\) for every Hermitian \(G\),
then \(gg^*=I_2\) by taking \(G=I_2\); the remaining equations say that \(g\)
commutes with every Hermitian matrix.  Thus \(g\) is scalar, and
\(\det g=1\) gives \(g=\pm I_2\).

The differential has zero kernel: if
\(AG+GA^*=0\) for all Hermitian \(G\), then \(A+A^*=0\) at \(G=I_2\), and
the remaining equation is \([A,G]=0\) for every Hermitian \(G\).  Hence \(A\)
is scalar; \(\operatorname{tr}A=0\) gives \(A=0\).  Both
\(SL_2(\mathbb C)\) and \(SO^+(1,3)\) have real dimension six, so the image
is an open subgroup of the connected group \(SO^+(1,3)\), hence the whole
group.  Finally, \eqref{hlob:eq:block-action} gives the action of the
central element and identifies the three standard real representations.
\end{proof}

The cubic in \eqref{hlob:eq:cubic-coupling} is invariant under the full
embedded \(SL_2(\mathbb C)\) action.  Indeed,
\begin{equation}
 \operatorname{adj}(gGg^*)=g^{-*}\operatorname{adj}(G)g^{-1}
 \qquad(g\in SL_2(\mathbb C)),
 \label{hlob:eq:adjugate-covariance}
\end{equation}
so
\((g\psi)^*\operatorname{adj}(gGg^*)(g\psi)
=\psi^*\operatorname{adj}(G)\psi\).
For the larger \(GL_2(\mathbb C)\) action it is instead a relative invariant:
\begin{equation}
 \det\!\left(T_P(g)HT_P(g)^*\right)
 =|\det g|^2\det H.
 \label{hlob:eq:GL2-relative-invariant}
\end{equation}
Thus no \(GL_2\)-invariance is being inferred from the \(SL_2\) formula.

\subsection{The Lorentz cone inside the rank-loss
\texorpdfstring{\(S^8\)}{S8} and its normal opening}
\label{hlob:subsec:S8}

Let
\begin{equation}
 e_2(K)=\frac12\bigl((\operatorname{tr}K)^2
                         -\operatorname{tr}(K^2)\bigr)
 \qquad(K\in\mathcal J_3).
 \label{hlob:eq:e2}
\end{equation}

\begin{theorem}[Proper Lorentz-cone section and determinant opening]
\label{hlob:thm:cone-section}
For every \(G\in\mathcal J_2\),
\begin{equation}
 \operatorname{rank}\jmath_P(G)=\operatorname{rank}G,
 \qquad
 \det\jmath_P(G)=0,
 \qquad
 e_2(\jmath_P(G))=\frac13\det G.
 \label{hlob:eq:e2-section}
\end{equation}
Hence \(\jmath_P\) maps the entire future Lorentz cone properly and
injectively into
\begin{equation}
 \partial\mathcal P_3
 =\{H\succeq0:\det H=0\}
 =\mathcal P_{3,\leq2}.
 \label{hlob:eq:rank-loss-boundary}
\end{equation}
It extends to an embedding of one-point compactifications
\begin{equation}
 \jmath_P^+:\mathcal P_2^+\cong B^4
 \lhook\joinrel\longrightarrow
 (\partial\mathcal P_3)^+\cong S^8.
 \label{hlob:eq:compactified-section}
\end{equation}
The timelike, nonzero null, and zero strata map respectively to rank two,
rank one, and rank zero.

For every real \(\varepsilon\), the affine slice through \(H_0=uu^*\)
satisfies the exact identity
\begin{equation}
 \boxed{
 \det\bigl(uu^*+\varepsilon\jmath_P(G)\bigr)
 =\varepsilon^2\det G.}
 \label{hlob:eq:determinant-opening}
\end{equation}
For \(\varepsilon>0\) and \(G\succeq0\), timelike \(G\) maps to the interior
of \(\mathcal P_3\), nonzero null \(G\) maps to its rank-two boundary, and
\(G=0\) maps to \(H_0\).  The restricted determinant Hessian is
\begin{equation}
 D^2(\det)_{H_0}
 \bigl(\jmath_P(G),\jmath_P(K)\bigr)
 =\operatorname{tr}G\operatorname{tr}K-\operatorname{tr}(GK),
 \label{hlob:eq:restricted-hessian}
\end{equation}
so
\begin{equation}
 \frac12D^2(\det)_{H_0}
 \bigl(\jmath_P(G),\jmath_P(G)\bigr)
 =\det G
 =t^2-x_1^2-x_2^2-x_3^2.
 \label{hlob:eq:hessian-lorentz}
\end{equation}

The surjection
\(\mathcal O_P:\partial\mathcal P_3\to\mathcal P_2\) is not proper and
therefore does not extend by sending infinity to infinity.  Thus
\eqref{hlob:eq:compactified-section} is an actual compactified embedding,
whereas the reverse compactified arrow is not supplied by
\(\mathcal O_P\).
\end{theorem}

\begin{proof}
Because \(R\) has complex rank two, congruence by \(R\) preserves the rank
of a \(2\times2\) Hermitian matrix.  Its image has rank at most two, so its
\(3\times3\) determinant vanishes.  When \(G\) is invertible, the product of
the two nonzero eigenvalues of \(RGR^*\) equals
\[
 \det(GR^*R)=\det G\det(C^{-1})=\frac13\det G.
\]
That product is \(e_2(RGR^*)\).  Both sides are polynomial in the four real
entries of \(G\), so the identity extends to singular \(G\), proving
\eqref{hlob:eq:e2-section}.  The map is a linear embedding between
finite-dimensional spaces, so its restriction to the closed cone is proper.
The compactification statement now follows from the previously proved
identifications
\(\mathcal P_2^+\cong B^4\) and
\((\partial\mathcal P_3)^+\cong S^8\).

Using \eqref{hlob:eq:X},
\[
 uu^*+\varepsilon\jmath_P(G)
 =X\begin{pmatrix}\varepsilon G&0\\0&1\end{pmatrix}X^*.
\]
Since \(\det X=1\), taking determinants proves
\eqref{hlob:eq:determinant-opening}.  The rank and positivity statements
follow from this same block matrix.  Polarizing
\(\det(sG+tK)\) gives
\[
 \det(sG+tK)
 =s^2\det G
  +st\bigl(\operatorname{tr}G\operatorname{tr}K
           -\operatorname{tr}(GK)\bigr)
  +t^2\det K,
\]
which proves \eqref{hlob:eq:restricted-hessian} and
\eqref{hlob:eq:hessian-lorentz}.

For nonproperness choose \(G_0=\operatorname{diag}(1,0)\) and
\[
 H_n=\jmath_P(G_0)+n uu^*,\qquad n\in\mathbb N.
\]
The two summands have orthogonal ranges, so \(H_n\succeq0\) has rank two
and lies in \(\partial\mathcal P_3\).  Its norm tends to infinity, while
\(\mathcal O_P(H_n)=G_0\) for every \(n\).  A continuous extension taking
the source point at infinity to the target point at infinity would therefore
send this convergent compactified sequence both to \(G_0\) and to infinity,
which is impossible.
\end{proof}

\subsection{The ten-dimensional determinant-phase carrier in the same coordinates}
\label{hlob:subsec:ten-dimensional}

The square-matrix \(SU(3)\)-orbit theorem already identifies its affine
quotient with
\begin{equation}
 \mathcal Y_3
 =\{(H,\omega):H\in\mathcal P_3,\ |\omega|^2=\det H\}.
 \label{hlob:eq:Y3}
\end{equation}
Substitution of \eqref{hlob:eq:block-expression} gives an exact
reformulation.

\begin{corollary}[Lorentz-spinor-scalar-phase form of the \(SU(3)\) quotient]
\label{hlob:cor:Y3-block-form}
There is a homeomorphism
\begin{equation}
 \mathcal Y_3\cong
 \left\{(G,\psi,\lambda,\omega):
 \begin{pmatrix}G&\psi\\\psi^*&\lambda\end{pmatrix}\succeq0,\quad
 |\omega|^2
 =\lambda\det G-\psi^*\operatorname{adj}(G)\psi
 \right\}.
 \label{hlob:eq:Y3-block-homeomorphism}
\end{equation}
The \(SL_2(\mathbb C)\)-action is
\begin{equation}
 (G,\psi,\lambda,\omega)\longmapsto
 (gGg^*,g\psi,\lambda,\omega).
 \label{hlob:eq:Y3-action}
\end{equation}
The complex determinant coordinate is invariant, and its zero locus is
exactly the rank-deficient cone:
\begin{equation}
 \omega=0\quad\Longleftrightarrow\quad\det H=0.
 \label{hlob:eq:omega-zero}
\end{equation}
After the proved one-point compactifications, this is the distinguished
\(S^8\) inside the ten-real-dimensional \(S^{10}\) carrier.  The ray
\(\{(G,\psi,\omega)=(0,0,0),\lambda\geq0\}\) is retained inside that
zero locus.
\end{corollary}

\begin{proof}
The block-coordinate map is an invertible real-linear congruence, so it
preserves positivity and is a homeomorphism.  Equation
\eqref{hlob:eq:cubic-coupling} gives the displayed scalar relation.
Equations \eqref{hlob:eq:block-action} and
\eqref{hlob:eq:adjugate-covariance} prove invariance.  The equivalence
\eqref{hlob:eq:omega-zero} follows directly from
\(|\omega|^2=\det H\).  The compactification and the identification of its
zero-determinant locus are the previously proved \(SU(3)\)-orbit and
rank-loss compactification theorems.
\end{proof}

This corollary makes the relation among the dimensions explicit.  The
nine-real-dimensional Hermitian carrier has the \(4+4+1\) decomposition in
\eqref{hlob:eq:four-four-one}.  Adding the complex determinant coordinate
and retaining its one-real-modulus equation gives \(9+2-1=10\) real
dimensions.  The determinant-zero locus removes the phase and imposes the
cubic boundary equation, leaving the eight-real-dimensional rank-loss cone
whose one-point compactification is \(S^8\).

\subsection{The holomorphic toric seed and its Lorentz-orbit completion}
\label{hlob:subsec:toric-seed}

The local four-branch chart has the already typed monomial map
\[
 q(z_0,z_1,z_2,z_3)=(z_0z_3,z_1,z_2).
\]
Define, without identifying angular and filling lattices,
\begin{equation}
 \kappa(w_0,w_1,w_2)=\operatorname{diag}(w_0,w_1w_2),
 \qquad
 K=\kappa q.
 \label{hlob:eq:kappa}
\end{equation}
Then
\begin{equation}
 K(z)=\operatorname{diag}(z_0z_3,z_1z_2),
 \qquad
 \boxed{\det K(z)=z_0z_1z_2z_3.}
 \label{hlob:eq:toric-seed}
\end{equation}
This is a holomorphic toric morphism into the diagonal complex matrices.
On the retained positive real locus put
\[
 a=z_0z_3,\qquad b=z_1z_2,\qquad
 p_0=\frac{a+b}{2},\qquad p_3=\frac{a-b}{2},\qquad p_1=p_2=0.
\]
Then \(K(z)=G(p)\succeq0\) and
\[
 p_0^2-p_3^2=ab=z_0z_1z_2z_3.
\]
Thus the toric products are exactly the two light-cone coordinates
\(p_0+p_3\) and \(p_0-p_3\).

For a retained \(\mu>0\), define the affine section
\begin{equation}
 \Psi_\mu(G)=\mu uu^*+\jmath_P(G).
 \label{hlob:eq:affine-section}
\end{equation}
Equation \eqref{hlob:eq:determinant-opening} gives
\begin{equation}
 \det\Psi_\mu(G)=\mu\det G,\qquad
 \Psi_\mu(G)\succeq0\Longleftrightarrow G\succeq0,\qquad
 \operatorname{rank}\Psi_\mu(G)=1+\operatorname{rank}G.
 \label{hlob:eq:affine-identities}
\end{equation}
Every \(G\succeq0\) has a spectral decomposition
\(G=v\operatorname{diag}(a,b)v^*\) with \(a,b\geq0\) and
\(v\in SU(2)\): a scalar phase changes the determinant of a chosen
unitary diagonalizer without changing its congruence action.  Hence
\begin{equation}
 \Theta_\mu:SL_2(\mathbb C)\times\mathbb R_{\geq0}^4
 \longrightarrow
 \Psi_\mu(\mathcal P_2),\qquad
 \Theta_\mu(g,z)=\Psi_\mu(gK(z)g^*)
 \label{hlob:eq:orbit-completion}
\end{equation}
is surjective.  It obeys
\begin{equation}
 \det\Theta_\mu(g,z)=\mu z_0z_1z_2z_3.
 \label{hlob:eq:orbit-completion-det}
\end{equation}
On \((\mathbb C^*)^4\), the complete auxiliary fibre of \(K\) is
\begin{equation}
 (\rho,\sigma)\cdot(z_0,z_1,z_2,z_3)
 =(\rho z_0,\sigma z_1,\sigma^{-1}z_2,\rho^{-1}z_3),
 \qquad(\rho,\sigma)\in(\mathbb C^*)^2.
 \label{hlob:eq:toric-fibre}
\end{equation}
Thus the local toric chart supplies a diagonal \(1+1\) seed, its
\(SL_2(\mathbb C)\)-orbit supplies all future causal directions, and the
two-complex-dimensional fibre in \eqref{hlob:eq:toric-fibre} is retained
as the exact redundancy of this parametrization.

\subsection{An exact colour-invariant Lorentz momentum and constrained dynamics}
\label{hlob:subsec:physical-translation}

The preceding results supply a pointwise carrier with commuting colour and
spin actions and a complete Lorentz-covariant orbit coordinate.  Let
\begin{equation}
 \mathscr V=\operatorname{Mat}_{3\times2}(\mathbb C),\qquad
 U\cdot Z=UZ\quad(U\in SU(3)),\qquad
 g\triangleright Z=Zg^*\quad(g\in SL_2(\mathbb C)).
 \label{hlob:eq:commuting-actions}
\end{equation}
These two left actions commute.  Retain positive conversion constants
\(\alpha>0\) and \(c>0\), with \(c\) carrying speed dimension, and define
\begin{equation}
 p_\alpha(Z)=\alpha Z^*Z\in\mathcal P_2,\qquad
 m_{\alpha,c}^2(Z)=\frac{\det p_\alpha(Z)}{c^2}.
 \label{hlob:eq:momentum-map}
\end{equation}
If \(p_\alpha\) is assigned momentum dimension, then
\([\alpha]=[p]/[Z]^2\) and
\([m_{\alpha,c}^2]=[p]^2/[c]^2=[\mathrm{mass}]^2\); no value of
\(\alpha\) or \(c\) is selected or suppressed.  Equivalently,
\(\det p_\alpha=m_{\alpha,c}^2c^2\), as required when the time component
of the Hermitian momentum matrix is \(E/c\).

\begin{theorem}[Complete colour-orbit momentum and its mass invariant]
\label{hlob:thm:colour-momentum}
The map \(p_\alpha\) is \(SU(3)\)-invariant, is covariant for the spin action,
and separates \(SU(3)\)-orbits:
\begin{equation}
 p_\alpha(UZ)=p_\alpha(Z),\qquad
 p_\alpha(g\triangleright Z)=g p_\alpha(Z)g^*.
 \label{hlob:eq:momentum-covariance}
\end{equation}
It induces a homeomorphism
\begin{equation}
 \mathscr V/SU(3)\xrightarrow{\ \sim\ }\mathcal P_2,
 \qquad [Z]\longmapsto p_\alpha(Z).
 \label{hlob:eq:momentum-quotient}
\end{equation}
Writing the rows of \(Z\) as three two-component covectors, one has
\begin{equation}
 \boxed{
 m_{\alpha,c}^2(Z)
 =\frac{\alpha^2}{c^2}\sum_{1\leq a<b\leq3}
 |Z_{a1}Z_{b2}-Z_{a2}Z_{b1}|^2.}
 \label{hlob:eq:mass-pluecker}
\end{equation}
It is invariant under both \(SU(3)\) and \(SL_2(\mathbb C)\), and
\begin{equation}
 m_{\alpha,c}^2(Z)>0\Longleftrightarrow\operatorname{rank}_{\mathbb C}Z=2,
 \qquad
 m_{\alpha,c}^2(Z)=0\Longleftrightarrow\operatorname{rank}_{\mathbb C}Z\leq1.
 \label{hlob:eq:mass-rank}
\end{equation}
For every rank-two \(Z_0\),
\begin{equation}
 m_{\alpha,c}^2(rZ_0)=r^4m_{\alpha,c}^2(Z_0)
 \qquad(r>0),
 \label{hlob:eq:mass-scaling}
\end{equation}
so the unrestricted kinematic carrier contains positive masses arbitrarily
close to zero.
\end{theorem}

\begin{proof}
The covariance equations follow directly from
\eqref{hlob:eq:commuting-actions}.  The orbit-separation and homeomorphism
statement is the exact two-column \(SU(3)\) Gram quotient, with multiplication
by the retained positive scalar \(\alpha\).  The \(2\times2\) Gram
determinant is the Cauchy--Binet identity
\[
 \det(Z^*Z)
 =\sum_{a<b}|Z_{a1}Z_{b2}-Z_{a2}Z_{b1}|^2,
\]
which proves \eqref{hlob:eq:mass-pluecker}.  Congruence by an
\(SL_2(\mathbb C)\) matrix preserves the determinant, while left \(SU(3)\)
preserves \(Z^*Z\).  The sum vanishes exactly when all \(2\times2\) minors
vanish, which is equivalent to rank at most one.  Finally
\(p_\alpha(rZ)=r^2p_\alpha(Z)\), and taking the \(2\times2\) determinant
gives \eqref{hlob:eq:mass-scaling}.
\end{proof}

For a retained mass \(m_0\geq0\), the constraint
\begin{equation}
 \mathscr C_{m_0,\alpha,c}
 =\{Z\in\mathscr V:\det p_\alpha(Z)=m_0^2c^2\}
 \label{hlob:eq:mass-shell-preimage}
\end{equation}
is invariant under both actions, and for \(m_0>0\) its \(SU(3)\)-orbit
space is exactly the future mass hyperboloid
\begin{equation}
 \{p\in\mathcal P_2:\det p=m_0^2c^2\}.
 \label{hlob:eq:mass-hyperboloid}
\end{equation}

Retain the Lorentz bilinear pairing
\begin{equation}
 \langle A,B\rangle_L
 =\frac12\bigl(\operatorname{tr}A\operatorname{tr}B
                    -\operatorname{tr}(AB)\bigr),
 \qquad
 \langle A,A\rangle_L=\det A.
 \label{hlob:eq:lorentz-pairing}
\end{equation}
For paths \(x:I\to\mathcal J_2\), \(Z:I\to\mathscr V\), and a real
multiplier \(e:I\to\mathbb R\), define the first-order constrained action
\begin{equation}
 \mathcal S_{m_0,\alpha,c}[x,Z,e]
 =\int_I
 \left(
  \langle p_\alpha(Z),\dot x\rangle_L
  -\frac{e}{2}\bigl(\det p_\alpha(Z)-m_0^2c^2\bigr)
 \right)d\tau.
 \label{hlob:eq:worldline-action}
\end{equation}
It is invariant under arbitrary \(\tau\)-dependent left \(SU(3)\) rotations
of \(Z\) and under the constant \(SL_2(\mathbb C)\) action
\begin{equation}
 x\longmapsto gxg^*,\qquad Z\longmapsto Zg^*.
 \label{hlob:eq:worldline-Lorentz-action}
\end{equation}
If \(x\) has length dimension and \(\tau\) retains its
chosen parameter dimension, then
\([e]=[x]/([\tau][p])\), so both summands in the integrand have dimension
\([p][x]/[\tau]\) and the action has dimension \([p][x]\).
Variation of \(e\) gives
\eqref{hlob:eq:mass-shell-preimage}; compactly supported variation of \(x\)
gives
\begin{equation}
 \frac{d}{d\tau}p_\alpha(Z)=0.
 \label{hlob:eq:momentum-conservation}
\end{equation}
For a variation \(\delta Z:I\to\mathscr V\), retain the exact differential
\begin{equation}
 \delta p_\alpha
 =\alpha\bigl(\delta Z^*Z+Z^*\delta Z\bigr).
 \label{hlob:eq:momentum-differential}
\end{equation}
The \(Z\)-Euler--Lagrange equation is the weak identity
\begin{equation}
 \left\langle
 \alpha\bigl(\delta Z^*Z+Z^*\delta Z\bigr),
 \dot x-e p_\alpha(Z)
 \right\rangle_L=0
 \quad\text{for every }\delta Z\in\mathscr V.
 \label{hlob:eq:Z-weak-equation}
\end{equation}
On the rank-two stratum this is equivalent to
\begin{equation}
 \dot x=e p_\alpha(Z).
 \label{hlob:eq:worldline-velocity}
\end{equation}
Indeed, for arbitrary \(A=A^*\in\mathcal J_2\), the variation
\begin{equation}
 \delta Z_A
 =\frac{1}{2\alpha}Z(Z^*Z)^{-1}A
 \label{hlob:eq:surjective-gram-variation}
\end{equation}
satisfies \(\delta p_\alpha=A\).  Hence the differential of the Gram map is
surjective there, and nondegeneracy of \(\langle\ ,\ \rangle_L\) proves
\eqref{hlob:eq:worldline-velocity}.  On the rank-at-most-one stratum the
differential need not be surjective, so the complete equation is
\eqref{hlob:eq:Z-weak-equation}; it is not replaced by the rank-two formula.
The invariance of \eqref{hlob:eq:lorentz-pairing} follows by polarizing
\(\det(gAg^*)=\det A\).  The multiplier and position equations follow by
direct variation and integration by parts, while
\eqref{hlob:eq:Z-weak-equation} follows from
\(D(\det)_p[A]=2\langle p,A\rangle_L\).

Equation \eqref{hlob:eq:worldline-action} is a complete constrained
kinematic dynamics with an independently retained mass parameter.  It does
not dynamically select \(m_0>0\): choosing \(m_0=0\) is equally compatible,
and \eqref{hlob:eq:mass-scaling} proves that the unrestrained carrier has no
positive kinematic threshold.  The next physical construction is thereby
specified: introduce a field-theoretic kinetic operator and interaction on
the full \((G,\psi,\lambda,\omega)\) carrier, derive rather than insert the
selected spectral sector, and compare its gauge-invariant spectrum with the
measured regime.  The exact morphisms and representations proved above are
the data that such a construction must preserve.

\section{The common marked slice and its dynamics}
\label{sec:q444-reconciliation}

Two Pro-mode packets dated 4 September 2026 describe a thirteen-dimensional
continuation of the marked nine-dimensional Hermitian carrier. One uses
the condition $u^*b=0$ in a complex-plus-triplet decomposition; the other
requires the observed quaternionic block $PHP^*$ to be complex. We prove
that these are the same embedded space, not merely isomorphic quadratic
models. We retain the original nonunitary marking throughout. We then
compute the omitted-coordinate force in the packet's scalar mechanics,
give a constrained alternative, and correct two action conventions.
The results below concern these specified matrices and dynamics; they do
not certify every theorem in either supplied packet.

The subsequent two-cross continuation computes the entire scalar--Gauss
domain in Section~\ref{q444:sgd:section}. Its classified parameters form
$\mathcal D\times\{1,-1\}^2\times U(1)^2$; retaining the frame phase
$\zeta$ gives the literal domain after the sole identification
$(\zeta,p_+,p_-)\sim(-\zeta,-p_+,-p_-)$.
Corollary~\ref{sgd:every-coupling} constructs tuples at every original
$t=g^2>0$. The two spatial residuals cannot vanish simultaneously on this class by
Theorem~\ref{q444:twocross:thm:exclusion}; existence on the remaining
unrestricted exceptional set is open.

Section~\ref{sgd:smooth-addon:section} proves smooth embeddedness and
the exact product
$\mathcal S_{\rm sg}\cong\{1,-1\}^2\times(0,\infty)
\times\mathbb R^2\times\mathbb T^3$, with $t=g^2$ the displayed
positive coordinate. The four literal components have dimension six;
every fixed positive coupling has four embedded five-dimensional
components $\mathbb R^2\times\mathbb T^3$. The inverse charts retain
all phases and both signs, without a physical gauge quotient or an
additional spatial equation.

\subsection{Conventions, determinant, and the retained marking}

Write $q=a+b\mathbf j$, with $a,b\in\mathbb C$,
$\mathbf j^2=-1$, $\mathbf ja=\bar a\mathbf j$. Thus
\[
 (a+b\mathbf j)(c+d\mathbf j)
 =ac-b\bar d+(ad+b\bar c)\mathbf j,
 \qquad \overline{a+b\mathbf j}=\bar a-b\mathbf j.
\]
For complex matrices of the appropriate sizes this gives
\[
 (A+B\mathbf j)(C+D\mathbf j)
 =AC-B\bar D+(AD+B\bar C)\mathbf j,
 \qquad (A+B\mathbf j)^*=A^*-B^{\mathsf T}\mathbf j.
\]
In particular a quaternionic Hermitian matrix is uniquely $A+B\mathbf j$
with $A=A^*$ and $B=-B^{\mathsf T}$. For $b\in\mathbb C^3$ define
\[
 B(b)=\begin{pmatrix}
 0&\bar b_3&-\bar b_2\\
 -\bar b_3&0&\bar b_1\\
 \bar b_2&-\bar b_1&0
 \end{pmatrix},\qquad \Theta(A,b)=A+B(b)\mathbf j.
\]
The map $\Theta:\operatorname{Herm}_3(\mathbb C)\oplus\mathbb C^3
\to\operatorname{Herm}_3(\mathbb H)$ is consequently a real-linear bijection.

We use the Moore determinant with value one on $I_3$. Its Hermitian
rank-three formula is
\begin{equation}\label{q444:moore}
 N\begin{pmatrix}
 a&q_3&\bar q_2\\ \bar q_3&d&q_1\\q_2&\bar q_1&f
 \end{pmatrix}
 =adf-a|q_1|^2-d|q_2|^2-f|q_3|^2
       +2\operatorname{Re}(q_1q_2q_3).
\end{equation}
The quaternionic spectral theorem and determinant congruence law used below
are recorded in S. Alesker, \emph{Non-commutative linear algebra and
plurisubharmonic functions of quaternionic variables}, Theorem 1.1.9,
\href{https://arxiv.org/abs/math/0104209}{arXiv:math/0104209}; see also
W. Wang, \emph{The linear algebra in the quaternionic pluripotential theory},
Propositions 2.1--2.3,
\href{https://arxiv.org/abs/1805.07759}{arXiv:1805.07759}.
In our signs the faithful complex representation is
\[
 \mathcal R(A+B\mathbf j)=
 \begin{pmatrix}A&B\\-\bar B&\bar A\end{pmatrix}.
\]
Wang's representation has the opposite signs on both off-diagonal blocks;
the exact conversion is conjugation by $\operatorname{diag}(I,-I)$.

\begin{lemma}[The full cubic]\label{q444:pair-cubic}
For every $A=A^*\in\operatorname{Mat}_3(\mathbb C)$ and $b\in\mathbb C^3$,
\begin{equation}\label{q444:pair}
 N(\Theta(A,b))=\det A-b^*Ab.
\end{equation}
\end{lemma}
\begin{proof}
Write the cyclic off-diagonal entries of $A$ as
$A_{23}=z_1$, $A_{31}=z_2$, $A_{12}=z_3$. Then
$q_i=z_i+\bar b_i\mathbf j$ in \eqref{q444:moore}, and
$|q_i|^2=|z_i|^2+|b_i|^2$. The complex part of their product is
\[
 z_1z_2z_3-\bar b_1b_2z_3-z_1\bar b_2b_3-\bar b_1\bar z_2b_3.
\]
The three additional cross terms in twice its real part are the negatives
of the off-diagonal terms of $b^*Ab$. The additional diagonal terms in
\eqref{q444:moore} are
$-a|b_1|^2-d|b_2|^2-f|b_3|^2$. All remaining terms are precisely
the ordinary determinant of $A$. This proves the identity with all signs.
\end{proof}

Retain the marked matrices
\begin{equation}\label{q444:marking}
 \begin{gathered}
 P=\begin{pmatrix}1&-1&0\\0&1&-1\end{pmatrix},\quad
 u=\begin{pmatrix}1\\1\\1\end{pmatrix},\quad
 C=PP^*=\begin{pmatrix}2&-1\\-1&2\end{pmatrix},\\
 R=P^*C^{-1}=\frac13\begin{pmatrix}2&1\\-1&1\\-1&-2\end{pmatrix},
 \quad X=(R\;u),\quad
 X^{-1}=\begin{pmatrix}P\\u^*/3\end{pmatrix},\quad \det X=1.
 \end{gathered}
\end{equation}
Direct multiplication gives $PR=I_2$, $u^*R=0$,
$RP=I_3-uu^*/3$. In particular the metric retained by this marking is
\begin{equation}\label{q444:D}
 D=X^*X=\operatorname{diag}(C^{-1},3),\qquad
 C^{-1}=\frac13\begin{pmatrix}2&1\\1&2\end{pmatrix}.
\end{equation}
It is not the identity matrix.

\subsection{The two constructions are the same embedded space}

Define, in the original quaternionic matrix space,
\[
 \mathcal J_P=\{H\in\operatorname{Herm}_3(\mathbb H):
                       PHP^*\in\operatorname{Herm}_2(\mathbb C)\},
 \qquad
 \mathcal L_u=\{\Theta(A,b):u^*b=0\}.
\]
For any $(A,b)$ put
\begin{equation}\label{q444:full-coordinates}
 X^{-1}AX^{-*}=\begin{pmatrix}G&z\\z^*&\lambda\end{pmatrix},
 \qquad X^*b=\binom{\chi}{\rho},\qquad
 \varepsilon=\begin{pmatrix}0&1\\-1&0\end{pmatrix}.
\end{equation}
Here $G$ is complex Hermitian, $z,\chi\in\mathbb C^2$,
$\lambda\in\mathbb R$, $\rho\in\mathbb C$, and $\rho=u^*b$.

\begin{theorem}[Embedded reconciliation]\label{q444:identification}
The two subspaces satisfy $\mathcal J_P=\mathcal L_u$. Their complete
coordinate relation, with no change of the ambient matrix $H$, is
\begin{equation}\label{q444:dictionary}
 H=X\begin{pmatrix}G&z+w\mathbf j\\(z+w\mathbf j)^*&\lambda\end{pmatrix}X^*,
 \qquad w=-\varepsilon\bar\chi,\qquad \chi=\varepsilon\bar w.
\end{equation}
The inverse in the original complex-plus-triplet coordinates is
\[
 A=X\begin{pmatrix}G&z\\z^*&\lambda\end{pmatrix}X^*,
 \qquad b=X^{-*}\binom{\varepsilon\bar w}{0}.
\]
Thus both thirteen-dimensional constructions have the same positive cone,
the same restricted determinant, and the same embedded old complex space,
which is $w=0$, equivalently $\chi=0$.
\end{theorem}
\begin{proof}
The skew matrices obey the following two identities in the fixed basis:
\begin{equation}\label{q444:skew}
 PB(b)P^{\mathsf T}=\overline{u^*b}\,\varepsilon,
 \qquad X^{-1}B(b)X^{-\mathsf T}=B(X^*b).
\end{equation}
For the first identity the upper-right entry is
$\bar b_3+\bar b_2+\bar b_1$; skew symmetry gives the other entries.
For the second, write $B(b)_{ij}=\sum_k\epsilon_{ijk}\bar b_k$.
For any real invertible $T$, expansion of the alternating determinant gives
\[
 \sum_{a,b}T_{ia}T_{jb}\epsilon_{abk}
 =\det(T)\sum_l\epsilon_{ijl}(T^{-1})_{kl}.
\]
Multiplication by $\bar b_k$ therefore gives
$TB(b)T^{\mathsf T}=\det(T)B(T^{-\mathsf T}b)$.
Set $T=X^{-1}$ and use $\det X=1$. This proves \eqref{q444:skew}.
Its first identity proves the asserted equality of subspaces.
At $\rho=0$ the second identity is the explicit matrix
\[
 B\binom{\chi}{0}
 =\begin{pmatrix}0&0&-\bar\chi_2\\
                  0&0&\bar\chi_1\\
                  \bar\chi_2&-\bar\chi_1&0\end{pmatrix}.
\]
Its upper-right column is $-\varepsilon\bar\chi$. This is exactly
\eqref{q444:dictionary}; $\varepsilon^2=-I_2$ gives the inverse. Every
claim about the embedded matrix and its positive cone follows from this
identity, rather than from matching dimensions or determinant levels.
\end{proof}

\begin{proposition}[The restricted and unrestricted cubics]\label{q444:cubics}
In \eqref{q444:full-coordinates} the unrestricted cubic is
\begin{align}\label{q444:N15}
 N_{15}={}&\lambda\det G-z^*\operatorname{adj}(G)z-\chi^*G\chi\notag\\
          &-2\operatorname{Re}(\rho\chi^*z)-\lambda|\rho|^2.
\end{align}
On the common thirteen-space this becomes
\begin{equation}\label{q444:N13}
 N_{13}=\lambda\det G-z^*\operatorname{adj}(G)z
                     -w^*\operatorname{adj}(G)w.
\end{equation}
\end{proposition}
\begin{proof}
Since $\det X=1$, the ordinary block determinant gives
$\det A=\lambda\det G-z^*\operatorname{adj}(G)z$ first for invertible
$G$ and then, as a polynomial identity, for all $G$. Also
\[
 b^*Ab=(\chi^*,\bar\rho)
       \begin{pmatrix}G&z\\z^*&\lambda\end{pmatrix}\binom{\chi}{\rho}.
\]
Expansion and Lemma~\ref{q444:pair-cubic} prove \eqref{q444:N15}.
For $G=\left(\begin{smallmatrix}a&c\\\bar c&d\end{smallmatrix}\right)$,
direct multiplication gives
$\varepsilon^*\operatorname{adj}(G)\varepsilon=G^{\mathsf T}$.
Substituting $w=-\varepsilon\bar\chi$ shows
$w^*\operatorname{adj}(G)w=\chi^*G\chi$, proving \eqref{q444:N13}.
In the second packet's notation $w=\varepsilon\bar\psi_R$, the exact
dictionary is consequently $\chi=-\psi_R$.
\end{proof}

\subsection{Positive fibres and the embedded sphere pair}

\begin{proposition}[All positive fibres]\label{q444:positive}
The matrix in \eqref{q444:dictionary} is positive semidefinite exactly when
\[
 G\succeq0,\qquad z,w\in\operatorname{Ran}_{\mathbb C}G,\qquad
 s:=\lambda-z^*G^\dagger z-w^*G^\dagger w\geq0.
\]
Its quaternionic rank is $\operatorname{rank}_{\mathbb C}G+\mathbf{1}_{s>0}$.
Here $G^\dagger$ is the inverse on $\operatorname{Ran}G$, extended by zero
on $\ker G$. In particular the fibre over $G=0$ is exactly
$\{\lambda uu^*: \lambda\geq0\}$.
\end{proposition}
\begin{proof}
Congruence by the invertible real matrix $X$ preserves positivity and rank.
Write $\psi=z+w\mathbf j$. For any complex Hermitian $T$,
quaternionic multiplication gives
\[
 \psi^*T\psi=z^*Tz+w^*Tw;
\]
the $\mathbf j$ coefficient is
$z^*Tw-w^{\mathsf T}\bar T\bar z=0$, while the scalar $w$ term is
the conjugate of the real number $w^*Tw$.
Testing the full Hermitian form on $(v,0)$ gives $G\succeq0$.
If $v\in\ker G$ and $v^*\psi\neq0$, testing on $(tv,q)$ with
$q$ chosen to make $\operatorname{Re}(v^*\psi q)\neq0$ makes the
form negative for one sufficiently large signed real $t$.
Thus $\psi$ is in the quaternionic range of $G$, exactly the two complex
range conditions above. With $a=G^\dagger\psi$, direct multiplication gives
\[
 \begin{pmatrix}I&0\\-a^*&1\end{pmatrix}
 \begin{pmatrix}G&\psi\\\psi^*&\lambda\end{pmatrix}
 \begin{pmatrix}I&-a\\0&1\end{pmatrix}
 =\begin{pmatrix}G&0\\0&s\end{pmatrix}.
\]
Both triangular factors are invertible, proving positivity and rank in
both directions. At $G=0$ the two spinors vanish, and return through $X$
gives the asserted ray.
\end{proof}

\begin{proposition}[The old boundary inside the new boundary]\label{q444:pair-topology}
Let $\mathcal C=\mathcal J_P\cap\{H\succeq0\}$. Its determinant-zero
locus is its boundary, and its one-point compactification satisfies
\[
 (\mathcal C^+,\partial\mathcal C^+)\cong(B^{13},S^{12}),\qquad
 ((\partial\mathcal C)^+,
   (\partial\operatorname{Herm}_3(\mathbb C)_{\succeq0})^+)
 \cong(S^{12},S^8).
\]
The sphere pair is the standard linear pair and
$S^{12}\setminus S^8\cong S^3\times\mathbb R^9$.
\end{proposition}
\begin{proof}
The identity belongs to $\mathcal J_P$, since $PI_3P^*=C$ is complex.
Set $\mathcal J_P^0=\{K\in\mathcal J_P:\operatorname{tr}K=0\}$.
If $\lambda_{\min}(K)$ is the least quaternionic Hermitian eigenvalue,
the mutually inverse maps are
\[
 K\longmapsto K-\lambda_{\min}(K)I_3,
 \qquad H\longmapsto H-\frac{\operatorname{tr}H}{3}I_3.
\]
They identify $\mathcal J_P^0$ with the positive singular locus. Indeed,
the first produces least eigenvalue zero; for singular positive $H$, the
least eigenvalue after the second map is $-\operatorname{tr}(H)/3$.
Eigenvalue continuity proves both maps continuous. More generally
\[
 (K,t)\longmapsto K+(t-\lambda_{\min}(K))I_3,
 \quad (K,t)\in\mathcal J_P^0\times[0,\infty)
\]
is a homeomorphism onto $\mathcal C$, taking $t=0$ to its boundary.
This also proves that no other boundary points occur. Its inverse sends
$H$ to $(H-(\operatorname{tr}H/3)I_3,\lambda_{\min}(H))$.
The one-point compactification of a closed half-space is a closed ball,
with the boundary compactifying to a sphere; for example inverse
stereographic projection from a point on the equator gives this pair.
On complex matrices the quaternionic eigenvalues are the same ordinary
Hermitian eigenvalues. The trace-free pair is therefore the linear pair
$(\mathbb R^{12},\mathbb R^8)$ under the displayed homeomorphisms.
Both compactifications have their infinity point in the smaller sphere,
so its complement is
$\mathbb R^8\times(\mathbb R^4\setminus\{0\})
\cong\mathbb R^8\times S^3\times\mathbb R$.
This proves the full pair assertion, not just the two dimensions.
\end{proof}

\subsection{The nonunitary Jordan-product correction}

\begin{theorem}[Completion with the retained metric]\label{q444:jordan}
The Jordan subalgebra generated by $\mathcal J_P$ under
$H\circ K=(HK+KH)/2$ is all of $\operatorname{Herm}_3(\mathbb H)$.
In the marked block coordinates the transported product is
\begin{equation}\label{q444:jordan-product}
 M\circ_D N=\frac12(MDN+NDM),\qquad D=\operatorname{diag}(C^{-1},3),
\end{equation}
not the ordinary block product.
\end{theorem}
\begin{proof}
For $F_X(M)=XMX^*$, multiplication gives
$F_X(M)\circ F_X(N)=F_X((MDN+NDM)/2)$ exactly.
Take
\[
 M=\begin{pmatrix}0&\psi\\\psi^*&0\end{pmatrix},\quad
 N=\begin{pmatrix}0&\phi\\\phi^*&0\end{pmatrix},\quad
 \psi=\binom10,\quad \phi=\binom0{\mathbf j}.
\]
The upper-left block of $M\circ_DN$ is
\[
 \frac32\begin{pmatrix}0&-\mathbf j\\\mathbf j&0\end{pmatrix}.
\]
The bottom scalar is zero: its two terms are
$\mathbf j/3$ and $-\mathbf j/3$. The other blocks vanish. Replacing
$\mathbf j$ by $\mathbf k=\mathbf i\mathbf j$ gives the independent
second missing direction. The common slice already contains every other
coordinate in $\operatorname{Herm}_2(\mathbb H)\oplus\mathbb H^2\oplus
\mathbb R$. Hence the linear span of the slice and these two products
is the whole fifteen-dimensional algebra. Any Jordan subalgebra containing
the slice contains that span. This proves the closure.
The $1/2$ product printed in the quaternionic packet is correct for the
ordinary product in the block frame, but the coefficient for the stated
embedded slice under the original product is $3/2$.
\end{proof}

\subsection{The corrected group-action convention}

Put $Z=(z\;w)$. The ordinary direct product has the following left action:
\begin{equation}\label{q444:action}
 (g,k)\cdot(G,Z,\lambda)=(gGg^*,gZk^{-1},\lambda),
 \quad g\in SL_2(\mathbb C),\quad k\in SU(2).
\end{equation}
Composition gives $g_1g_2Zk_2^{-1}k_1^{-1}
=(g_1g_2)Z(k_1k_2)^{-1}$, as required. The packet's $gZk$ describes
commuting left and right actions, or a left action of
$SL_2(\mathbb C)\times SU(2)^{\mathrm{op}}$; it is not a left action
of the usual direct product with those parameters. For instance
$k_1=\varepsilon$, $k_2=\operatorname{diag}(i,-i)$ do not commute.

Both conventions give the same transformations after $k\mapsto k^{-1}$.
To prove invariance, use
\[
 \operatorname{adj}(gGg^*)=(g^*)^{-1}\operatorname{adj}(G)g^{-1},
 \qquad \det g=1,
\]
first on invertible $G$ and then polynomially on all $G$.
The two spinor terms in \eqref{q444:N13} are
$\operatorname{tr}(Z^*\operatorname{adj}(G)Z)$; conjugation by $k$
does not change this trace. Positivity under the left factor follows by
quaternionic congruence with $\operatorname{diag}(g,1)$ in the block frame.
For the right factor, a unit quaternion $a+b\mathbf j$ acts on $Z$ by
\[
 Z\begin{pmatrix}a&b\\-\bar b&\bar a\end{pmatrix}.
\]
This follows by multiplying $(z+w\mathbf j)(a+b\mathbf j)$ explicitly;
it preserves $\psi^*G^\dagger\psi$, a real scalar, and the range conditions
in Proposition~\ref{q444:positive}.
Finally, if an element fixes every $G$, then $g$ is unitary and commutes
with every complex Hermitian matrix; hence $g=\pm I_2$.
Fixing every $Z$ then forces $k=g$. The effective group is therefore
$(SL_2(\mathbb C)\times SU(2))/\{(I,I),(-I,-I)\}$.

\subsection{The full carrier still has an exact four-dimensional observation}

Requiring the observed quaternionic block to be complex defines the
thirteen-slice above. A different operation retains every coordinate of
the full fifteen-dimensional carrier and applies a positive observation
map. These two operations have different domains, with an exact relation.

\begin{proposition}[Full positive observation and the hidden-coordinate bound]
\label{q444:full-observation}
Fix the same complex subfield $\mathbb C=\mathbb R+\mathbb R\mathbf i$
and real matrix $P$. For a quaternionic Hermitian $2\times2$ matrix $A$
put
\[
 \pi_{\mathbb C}(A)=\tfrac12\bigl(A+(\mathbf iI_2)^*A(\mathbf iI_2)\bigr).
\]
Then the real-linear map
\[
 \mathscr O:\operatorname{Herm}_3(\mathbb H)\longrightarrow
       \operatorname{Herm}_2(\mathbb C),\qquad
 \mathscr O(H)=\pi_{\mathbb C}(PHP^*)
\]
is positive, onto, and has positive section $G\mapsto RGR^*$.
Its real kernel has dimension eleven. On writing $H=\Theta(A_0,b)$,
its observed value is $G=PA_0P^*$ and $\rho=u^*b$ obeys
\begin{equation}\label{q444:rho-bound}
 H\succeq0\quad\Longrightarrow\quad
 G\succeq0,\qquad |\rho|^2\leq\det_{\mathbb C}G.
\end{equation}
Every pair $(G,\rho)$ satisfying these inequalities occurs. Thus the
possible hidden $\rho$ at a fixed $G$ are exactly the closed complex disc
of radius $\sqrt{\det G}$. At a null observed block, $\det G=0$ forces
$\rho=0$.
\end{proposition}
\begin{proof}
For $q=a+b\mathbf j$, $a,b\in\mathbb C$, direct multiplication gives
$(q-\mathbf i q\mathbf i)/2=a$.
Thus $\pi_{\mathbb C}$ extracts the complex matrix part. It is positive
as the average of the identity congruence and unitary congruence by
$\mathbf iI_2$. Compression by $P$ is positive as well. Since $PR=I_2$
and $G$ is complex, $\mathscr O(RGR^*)=G$; positivity of this section is
again congruence. Rank--nullity gives $15-4=11$.

The original skew identity \eqref{q444:skew} gives
\[
 PHP^*=G+\bar\rho\varepsilon\mathbf j
 =\begin{pmatrix}
 g_{11}&g_{12}+\bar\rho\mathbf j\\
 \bar g_{12}-\bar\rho\mathbf j&g_{22}
 \end{pmatrix}.
\]
Its Moore determinant is exactly
$g_{11}g_{22}-|g_{12}|^2-|\rho|^2=\det_{\mathbb C}G-|\rho|^2$.
A Hermitian quaternionic $2\times2$ matrix with real diagonal $a,d$
and upper entry $q$ is positive semidefinite precisely when
$a,d\geq0$ and $ad\geq|q|^2$: for $a>0$ its quadratic form is
$a|x+a^{-1}qy|^2+(d-|q|^2/a)|y|^2$; for $a=0$, positivity forces
$q=0$ by varying $x$, and then requires $d\geq0$.
Applying this criterion proves \eqref{q444:rho-bound} and its converse
for the displayed compressed block.

To realize every permitted pair in the original coordinates, take any
$\lambda\geq0$ and set
\[
 A_0=RGR^*+\lambda uu^*,\qquad b=(\rho/3)u,
 \qquad H=\Theta(A_0,b).
\]
The identities for $X=(R\;u)$ give, with $\chi=R^*b=0$,
\[
 X^{-1}HX^{-*}
   =\begin{pmatrix}G+\bar\rho\varepsilon\mathbf j&0\\0&\lambda\end{pmatrix}.
\]
The preceding criterion and congruence prove positivity of $H$, while
its observation and contraction are exactly $G$ and $\rho$.
This proves the claimed complete image and every disc fibre for $(G,\rho)$.
It does not discard the further $z,w,\lambda$ coordinates in a general $H$.
\end{proof}

In explicit Lorentz coordinates, retain
\[
 G=\begin{pmatrix}p_0+p_3&p_1-ip_2\\p_1+ip_2&p_0-p_3\end{pmatrix},
 \qquad \det G=p_0^2-p_1^2-p_2^2-p_3^2.
\]
The bound is therefore an exact inequality between this Lorentz quadratic
form and the omitted complex coordinate. It does not supply a positive
uniform gap: taking $G=\epsilon I_2$, $\rho=\epsilon/2$,
$\lambda=\epsilon$, $\epsilon>0$, in the realization above gives positive
definite $H$, nonzero $\rho$, and $\det G=\epsilon^2\to0$.
An interpretation as physical mass or as a sector with a nonzero threshold
requires a state-space, units, dynamics, and sector dictionary. The formula
identifies the exact quantity that such a dictionary must retain.
It also avoids imposing $\rho=0$ on unrestricted motion merely to obtain
a four-dimensional observation. Positivity of a chosen dynamical trajectory
itself must still be proved. In fact a positive semidefinite trace-free
Hermitian matrix has all three eigenvalues nonnegative with sum zero, so
it is zero. Thus the nonzero scalar configurations below cannot themselves
be positive matrices.

There is an exact bridge that retains this domain distinction. Put
$\mathcal J_{14}=\{Q\in\operatorname{Herm}_3(\mathbb H):\operatorname{tr}Q=0\}$.
The restriction $\mathscr O|_{\mathcal J_{14}}$ is onto all of
$\operatorname{Herm}_2(\mathbb C)$, not only its positive cone: its section is
\[
 j_0(G)=RGR^*-\frac{\operatorname{tr}(C^{-1}G)}3uu^*.
\]
Indeed $R^*R=C^{-1}$, $\operatorname{tr}(uu^*)=3$, and $Pu=0$,
so $\operatorname{tr}j_0(G)=0$ and $\mathscr O(j_0(G))=G$.
Its kernel therefore has dimension ten.

For a positive boundary observation of every scalar configuration retain
the separate, nonlinear map
\begin{equation}\label{q444:spectral-floor}
 \mathcal F(Q)=Q-\lambda_{\min}(Q)I_3,\qquad
 \widehat G(Q)=\mathscr O(\mathcal F(Q))
             =P\Phi P^*-\lambda_{\min}(Q)C,
 \quad Q=\Theta(\Phi,b).
\end{equation}
It maps $\mathcal J_{14}$ bijectively and continuously to the positive
singular boundary of $\operatorname{Herm}_3(\mathbb H)_{\succeq0}$.
The inverse is $H\mapsto H-(\operatorname{tr}H/3)I_3$: the least
eigenvalue of this trace-free matrix is $-\operatorname{tr}H/3$.
Eigenvalue continuity proves continuity in both directions. This is the
full-space version of the earlier sphere-pair construction. The map leaves
$b$, and hence $\rho$, unchanged, and therefore
\[
 \widehat G(Q)\succeq0,\qquad
 \det\widehat G(Q)\geq|u^*b|^2.
\]
The retained matrix $C=\left(\begin{smallmatrix}2&-1\\-1&2\end{smallmatrix}\right)$
is essential in \eqref{q444:spectral-floor}. This exact observation of a
scalar trajectory is not an assertion that the trajectory preserves a
positive linear slice; at repeated least eigenvalues the observation need
not even be differentiable. No Hamiltonian intertwining is asserted.

The fixed marked observation is not automatically a gauge-invariant
four-momentum. There is an exact covariance rule: under
$H'=UHU^*$, $U\in SU(3)$, transform the observer as $P'=PU^*$;
then $P'H'P'^*=PHP^*$ and hence its complex observation agrees.
Keeping $P$ fixed gives a chosen observer; specifying how that observer
and the integral angular lattice participate in a physical model is
additional data, not supplied by an equality of dimensions.

\subsection{The actual period-to-observation bridge}
The common marking also occurs in the rank-three toroidal extension of
the native rank-two cusp. To specify an actual map, retain its data
\[
 B_0=\begin{pmatrix}0&1\\-1&0\end{pmatrix},\quad
 Q_{\rm latt}=B_0^{-1}P=\begin{pmatrix}0&-1&1\\1&-1&0\end{pmatrix},
 \quad S=\begin{pmatrix}1&1\\0&1\\0&0\end{pmatrix}.
\]
Here $PS=I_2$. Let $C_2(t)$ be the retained bounded holomorphic period
correction at $t=0$. On the logarithmic cover $t=e^{2\pi i s}$ put
\[
 Z_2=sB_0+C_2(t),\quad C_3(t)=SC_2(t)Q_{\rm latt},\quad
 Z_3=sI_3+C_3(t),\quad \Phi_{\rm latt}=P\oplus Q_{\rm latt}.
\]
These are the local corrected periods of the toroidal elliptic bridge;
the underlying toric and compact-fibre morphisms are recorded in the
workbench's \emph{Higher-rung toroidal elliptic bridge}, equations
(TEB.16)--(TEB.35). The new calculation below only needs the displayed
matrices and retains the entire correction $C_2(t)$.

\begin{proposition}[Metric-dependent period observation]\label{q444:period-gram}
Set $\Omega_3=(I_3\ Z_3)$ and $\Omega_2=(I_2\ Z_2)$.
For every real symmetric positive-definite $6\times6$ matrix $W_3$,
\[
 H_3=\Omega_3W_3\Omega_3^*,\quad
 W_2=\Phi_{\rm latt}W_3\Phi_{\rm latt}^{\mathsf T}
 \quad\Longrightarrow\quad
 PH_3P^*=\Omega_2W_2\Omega_2^*.
\]
Both sides are positive definite. No fixed positive $W_3$ makes $H_3$
periodic under $s\mapsto s+1$. Nevertheless every retained positive
matrix $W_*$ gives an explicit descending choice
\[
 M(x)=\begin{pmatrix}I_3&xI_3\\0&I_3\end{pmatrix},\quad
 x=\operatorname{Re}s,\quad
 W_3(s)=M(-x)W_*M(-x)^{\mathsf T}.
\]
For this choice, with $y=\operatorname{Im}s$, the descended matrix is
\[
 H_3(t)=(I_3\;iyI_3+C_3(t))W_*(I_3\;iyI_3+C_3(t))^*.
\]
If $W_*=\left(\begin{smallmatrix}E&F\\F^{\mathsf T}&K\end{smallmatrix}\right)$,
then $G(t)=PH_3(t)P^*$ has leading term $y^2PKP^*$, with $PKP^*\succ0$,
and $\det G(t)=y^4\det(PKP^*)+O(y^3)$ at the cusp.
\end{proposition}
\begin{proof}
Direct multiplication gives $PZ_3=Z_2Q_{\rm latt}$, hence
$P\Omega_3=\Omega_2\Phi_{\rm latt}$ and the claimed Gram identity.
Both $P$ and $Q_{\rm latt}$ have rank two, so $\Phi_{\rm latt}$ is
surjective and $W_2$ is positive definite. Each period matrix has an
identity block, proving positivity of both Gram matrices.

Write a fixed $W_3$ in $3+3$ blocks. Since $Z_3(s+n)=Z_3(s)+nI_3$,
the coefficient of $n^2$ in $H_3(s+n)$ is its positive-definite lower
diagonal block. Therefore $H_3(s+n)$ cannot be constant in $n\in\mathbb Z$.
In particular the period-derived observation cannot be made single-valued
by silently retaining a constant lattice metric.

Put $M_3=M(1)$. The proposed varying metric is positive by congruence
and satisfies $W_3(s+1)=M_3^{-1}W_3(s)M_3^{-\mathsf T}$.
Since $\Omega_3(s+1)=\Omega_3(s)M_3$, its Gram matrix is invariant.
Also $\Omega_3(s)M(-x)=(I_3\;iyI_3+C_3(t))$, which proves the exact
descended formula. The complete expansion, without removing its lower
terms, is
\begin{align*}
 H_3={}&y^2K+iy\bigl(F^{\mathsf T}-F+KC_3^*-C_3K\bigr)\\
       &+E+FC_3^*+C_3F^{\mathsf T}+C_3KC_3^*.
\end{align*}
Boundedness of $C_3$ gives the asserted leading term and determinant
estimate. The principal block $K$ is positive definite and $P^*$ is
injective, so $PKP^*$ is positive definite. This proves every assertion.
\end{proof}

For the explicitly specified example $W_*=I_6$ the exact expression is
$H_3=y^2I_3+iy(C_3^*-C_3)+I_3+C_3C_3^*$ and
$\det G=3y^4+O(y^3)$, retaining $C=PP^*$ and $\det C=3$.
Thus this period-derived observation grows in a timelike, full-rank
direction, not toward a null determinant. It is an actual positive
matrix-valued map on the punctured cusp, not an inverse recovery of
periods or of individual torus points. Its extension through the
successor's finite-order overlaps and special fillings, and a dynamical
interpretation of the chosen lattice metric, are further constructions.

% Insertion-ready fragment for the existing quaternionic module.
% It uses the already defined original P,R,u,C,X,B_0,Q_latt,S,C_2,C_3,
% Z_3,Omega_3 and the monodromy M_3 from q444:period-gram.
% No canonical file was edited by the independent checking agent.

\subsection*{All positive period-metric fibres and the cost of rank loss}

The following calculation keeps the periods in
Proposition~\ref{q444:period-gram} fixed.  In particular
$C_3(t)=SC_2(t)Q_{\rm latt}$ still uses the displayed $S$, not the
different right inverse $R=P^*C^{-1}$.  Choose a punctured cusp on which
the retained correction is bounded, and fix a finite constant $K\geq0$
such that $\|C_3(t)\|\leq K$ there.  Such a constant exists because
$\|C_3(t)\|\leq\|S\|\|C_2(t)\|\|Q_{\rm latt}\|$.
All norms below are the Euclidean operator norms.  Write
\[
 t=e^{2\pi i s},\qquad s=x+iy,\qquad
 y=-\frac{\log|t|}{2\pi},
\]
and restrict to the nonempty smaller cusp $y>\max\{K,1\}$.
This changes no period or coefficient.  It identifies a domain where
the inverse used below is proved to exist.

\begin{theorem}[The complete positive metric fibre]
\label{q444:metric-fibre}
Retain
\[
 A_{\rm per}(s)=xI_3+\operatorname{Re}C_3(t),\qquad
 Y_{\rm per}(t)=yI_3+\operatorname{Im}C_3(t),
\]
and define the real period map and its inverse by
\[
 \mathcal N(s)=
 \begin{pmatrix}I_3&A_{\rm per}(s)\\0&Y_{\rm per}(t)\end{pmatrix},
 \qquad
 \mathcal N(s)^{-1}=
 \begin{pmatrix}
 I_3&-A_{\rm per}(s)Y_{\rm per}(t)^{-1}\\
 0&Y_{\rm per}(t)^{-1}
 \end{pmatrix}.                                             \tag{MF.1}
\]
For every prescribed Hermitian matrix $H=E+iF\succ0$, where
$E=E^{\mathsf T}$ and $F=-F^{\mathsf T}$ are real, all real symmetric
positive-definite lattice matrices giving the unchanged value
$\Omega_3W_3\Omega_3^*=H$ are exactly
\[
 W_3=\mathcal N^{-1}
 \left[
  \frac12\begin{pmatrix}E&-F\\F&E\end{pmatrix}
  +\begin{pmatrix}U&V\\V&-U\end{pmatrix}
 \right]\mathcal N^{-\mathsf T},                             \tag{MF.2}
\]
where $U,V$ are real symmetric $3\times3$ matrices and the bracket is
positive definite.  This is a nonempty open convex twelve-real-dimensional
fibre.  An explicit description of its entire parameter domain is
\[
 E_+=E/2+U\succ0,\qquad
 E/2-U-(-F/2+V)^{\mathsf T}E_+^{-1}(-F/2+V)\succ0.           \tag{MF.3}
\]
In particular $U=V=0$ always gives a positive lift.

Every smooth map $H_3(t)\succ0$ on the smaller punctured cusp has a
smooth positive lift (MF.2) with
\[
 W_3(s+1)=M_3^{-1}W_3(s)M_3^{-\mathsf T}.
                                                               \tag{MF.4}
\]
The full set of positive observations $G(t)=PH_3(t)P^*$ obtainable
from such metrics is exactly the set of all smooth
$\operatorname{Herm}_2(\mathbb C)_{\succ0}$-valued maps on that cusp.
For fixed $G\succ0$, all the intermediate positive $H$ are
\[
 H=X\begin{pmatrix}
       G&z\\ z^*&z^*G^{-1}z+\tau
     \end{pmatrix}X^*,\qquad
 z\in\mathbb C^2,\quad \tau>0.                                \tag{MF.5}
\]
The positive real lattice-metric fibre of $G$ has real dimension
seventeen.
\end{theorem}

\begin{proof}
For every real vector $v$,
\[
 \|Y_{\rm per}v\|
 \geq y\|v\|-\|\operatorname{Im}C_3\,v\|
 \geq(y-K)\|v\|.
\]
The inequality $\|\operatorname{Im}C_3\|\leq\|C_3\|$ follows from
$\operatorname{Im}C_3=(C_3-\bar C_3)/(2i)$ and
$\|\bar C_3\|=\|C_3\|$.  Since $y>K$, $Y_{\rm per}$ has zero kernel
and is an invertible square real matrix.  Block multiplication now proves
both inverse identities in (MF.1).  This is the original real period map:
if $a,b\in\mathbb R^3$, then
\[
 a+Z_3b=(a+A_{\rm per}b)+iY_{\rm per}b.
\]
Thus the inverse takes $v+iw$ to
$(v-A_{\rm per}Y_{\rm per}^{-1}w,\;Y_{\rm per}^{-1}w)$.
No correction in $C_3$ has been removed.

Put $\mathcal J=(I_3\;\;iI_3)$.  The exact factorization is
$\Omega_3=\mathcal J\mathcal N$.  Congruence gives a bijection between
real positive $W_3$ and real positive
$T=\mathcal NW_3\mathcal N^{\mathsf T}$.  For
\[
 T=\begin{pmatrix}E_0&F_0\\F_0^{\mathsf T}&K_0\end{pmatrix}
\]
one obtains by multiplication
\[
 \mathcal JT\mathcal J^*
   =E_0+K_0+i(F_0^{\mathsf T}-F_0).                            \tag{MF.6}
\]
Consequently this equals $E+iF$ precisely when
$E_0=E/2+U$, $K_0=E/2-U$, and $F_0=-F/2+V$ with $U,V$ symmetric.
Indeed, the first equality in (MF.6) fixes $K_0=E-E_0$,
and $F_0^{\mathsf T}-F_0=F$ is equivalent to symmetry of $F_0+F/2$.
These recover $U,V$ uniquely, proving the full fibre formula.

For real $v,w$, with $q=v+iw$,
\[
 \begin{pmatrix}v\\w\end{pmatrix}^{\mathsf T}
 \begin{pmatrix}E&-F\\F&E\end{pmatrix}
 \begin{pmatrix}v\\w\end{pmatrix}=q^*Hq.
\]
The right side is positive for every nonzero $(v,w)$, so the first
bracket term in (MF.2) is positive definite.  This proves nonemptiness.
The kernel parameters have dimension $6+6=12$.  Positive definiteness
is an open convex condition on real symmetric matrices, so its intersection
with this affine twelve-dimensional fibre is open and convex and has
dimension twelve.  Finally, for $B_+=-F/2+V$,
\[
 \begin{pmatrix}E_+&B_+\\B_+^{\mathsf T}&E/2-U\end{pmatrix}
 =
 \begin{pmatrix}I&0\\B_+^{\mathsf T}E_+^{-1}&I\end{pmatrix}
 \begin{pmatrix}E_+&0\\0&E/2-U-B_+^{\mathsf T}E_+^{-1}B_+\end{pmatrix}
 \begin{pmatrix}I&E_+^{-1}B_+\\0&I\end{pmatrix}.
\]
The two triangular factors are transposes of each other and are invertible;
the identity proves the necessity and sufficiency of
(MF.3), including every positivity restriction on the kernel parameters.

The original correction is single-valued in $t$, so
\[
 \mathcal N(s+1)=\mathcal N(s)M_3,\qquad
 \mathcal N(s+1)^{-1}=M_3^{-1}\mathcal N(s)^{-1}.               \tag{MF.7}
\]
For any smooth prescribed $H_3(t)$, its real and imaginary parts are smooth
and periodic on the cover.  The inverse of $Y_{\rm per}$ is smooth,
since its entries are its adjugate divided by its nonzero determinant.
Formula (MF.2) with $U=V=0$ is therefore smooth and satisfies (MF.4).
More generally, smooth periodic $U,V$ satisfying (MF.3) give all such
smooth lifts: conversely (MF.7) and (MF.4) make
$\mathcal NW_3\mathcal N^{\mathsf T}$ periodic, and its blocks recover
periodic $U,V$ by (MF.6).  No holomorphic condition is being imposed on
the real metric or on the prescribed Hermitian map.

If $W_3\succ0$ is real, it is also positive on nonzero complex columns:
for $p+iq$, its quadratic form is $p^{\mathsf T}W_3p+
q^{\mathsf T}W_3q$.  Since $\Omega_3^*$ is injective by the identity
block of $\Omega_3$, one has $H_3\succ0$.  Since $P^*$ is injective,
also $G=PH_3P^*\succ0$.  This proves necessity of strict positivity;
no singular $G$ occurs at an interior point for a positive metric.

For the converse, the marked inverse $X^{-1}=(P;\,u^*/3)$ makes the
upper-left block of $X^{-1}HX^{-*}$ equal to $PH P^*$.  With that block
fixed at $G\succ0$, the complete positive block matrices are
\[
 \begin{pmatrix}G&z\\z^*&\ell\end{pmatrix},\qquad
 \ell-z^*G^{-1}z=\tau>0.
\]
To prove this exact positivity statement, multiply
\[
 \begin{pmatrix}I_2&0\\z^*G^{-1}&1\end{pmatrix}
 \begin{pmatrix}G&0\\0&\tau\end{pmatrix}
 \begin{pmatrix}I_2&G^{-1}z\\0&1\end{pmatrix}.
\]
The result is the displayed block with $\ell=z^*G^{-1}z+\tau$.
The invertible triangular factors prove both directions, and congruence
by $X$ gives (MF.5).  Taking $z=0$ and any fixed $\tau=\lambda>0$
already gives the smooth lift $H_3(t)=RG(t)R^*+\lambda uu^*$.
Applying (MF.2) proves every asserted positive observation is realized.
The intermediate $H$ fibre has four real coordinates in $z$ and one in
$\tau$.  Equivalently, the real-linear metric-to-observation map has rank
four: the constructions above give a right inverse on all Hermitian
matrices before imposing positivity.  Its real symmetric domain has
dimension $6\cdot7/2=21$, so its affine kernel has dimension $21-4=17$.
The positive part is nonempty and relatively open, proving the stated
seventeen-dimensional fibre without a dimension-only surjectivity argument.
\end{proof}

The freedom above does not replace the retained periods by an identity
period matrix.  It computes the metric which compensates their full real
period map.  The correction itself has not lost information in passing
from rank two to rank three: the exact inverse on its prescribed image is
\[
 C_2=P C_3 S B_0,
\]
because $PS=I_2$ and $Q_{\rm latt}S=B_0^{-1}$.
To identify exactly what an unrestricted metric observation cannot recover,
take two retained corrections with real period maps $\mathcal N$ and
$\widetilde{\mathcal N}$ on a common sufficiently small cusp.  The
invertible real transformation
\[
 T_{\rm per}=\widetilde{\mathcal N}^{-1}\mathcal N,\qquad
 \widetilde W_3=T_{\rm per}W_3T_{\rm per}^{\mathsf T}
                                                               \tag{MF.8}
\]
preserves positivity and gives
$\widetilde\Omega_3\widetilde W_3\widetilde\Omega_3^*
=\Omega_3W_3\Omega_3^*$ by the two exact factorizations
$\Omega_3=\mathcal J\mathcal N$.
Furthermore $T_{\rm per}(s+1)=M_3^{-1}T_{\rm per}(s)M_3$, so (MF.8)
preserves the prescribed metric monodromy.  This is an explicit
correspondence between period--metric pairs with the same $H_3$ and $G$.
Thus these observations alone cannot recover the period correction when
the metric is allowed this freedom.  The constants remain in the metric
formula; they have not been set to zero.

\begin{proposition}[Relative-metric bounds obstruct projective rank loss]
\label{q444:metric-condition}
For an equivariant positive metric put
\[
 W_*(t)=M(x)W_3(s)M(x)^{\mathsf T},\qquad
 \widetilde\Omega_3(t)=(I_3\;\;iyI_3+C_3(t)).
\]
These are single-valued on the punctured cusp, and
$H_3=\widetilde\Omega_3W_*\widetilde\Omega_3^*$.
If fixed constants $0<a\leq b<\infty$ obey
$aI_6\preceq W_*(t)\preceq bI_6$, then
\[
 a\bigl(1+(y-K)^2\bigr)C
 \preceq G(t)\preceq
 b\bigl(1+(y+K)^2\bigr)C.                                    \tag{MC.1}
\]
In particular the unchanged eigenvalues satisfy
\[
 \lambda_{\min}(G)\geq a\bigl(1+(y-K)^2\bigr)\lambda_{\min}(C),
 \qquad
 \lambda_{\max}(G)\leq b\bigl(1+(y+K)^2\bigr)\lambda_{\max}(C).
                                                               \tag{MC.2}
\]
With the original $C$, $\lambda_{\min}(C)=1$ and
$\lambda_{\max}(C)=3$.  More generally, writing
$\kappa_*(t)=\lambda_{\max}(W_*)/\lambda_{\min}(W_*)$, one always has
\[
 \frac{\lambda_{\min}(G)}{\lambda_{\max}(G)}
 \geq\frac{1+(y-K)^2}{3\kappa_*(t)(1+(y+K)^2)}.               \tag{MC.3}
\]
Thus if the observed eigenvalue ratio tends to zero as $t\to0$, the
relative metric condition number tends to infinity.  Bounded relative
condition number excludes this projective rank loss even when the common
scale of the metric is not bounded above or below.
\end{proposition}

\begin{proof}
The identity $M(x+1)=M(x)M_3$, together with (MF.4), proves that $W_*$
is periodic.  Multiplication gives
$\Omega_3M(-x)=\widetilde\Omega_3$, which proves the stated Gram formula.
For every complex vector $v$,
\[
 (y-K)\|v\|\leq\|(iyI_3+C_3)^*v\|\leq(y+K)\|v\|.
\]
Both inequalities are the triangle inequality, its reverse, and
$\|C_3^*\|\leq K$.  Since
$\widetilde\Omega_3\widetilde\Omega_3^*
=I_3+(iyI_3+C_3)(iyI_3+C_3)^*$, they give
\[
 \bigl(1+(y-K)^2\bigr)I_3
 \preceq\widetilde\Omega_3\widetilde\Omega_3^*
 \preceq\bigl(1+(y+K)^2\bigr)I_3.
\]
Congruencing $aI_6\preceq W_*\preceq bI_6$ by
$\widetilde\Omega_3$, and then by $P$, proves (MC.1).
The minimum and maximum Rayleigh quotients prove (MC.2).
The characteristic polynomial of the original $C$ is
$(z-1)(z-3)$, retaining the factors $1$ and $3$.

At each point take
$a(t)=\lambda_{\min}(W_*)$ and $b(t)=\lambda_{\max}(W_*)$ in the
same inequalities.  Dividing the resulting positive scalar bounds gives
(MC.3); this division is a comparison of eigenvalues, not a rescaling or
replacement of $G$ or $W_*$.  Its period-dependent factor tends to $1$.
For fixed $a,b$, the lower bound tends to $a/(3b)>0$; for example it is
at least $a/(27b)$ once $y\geq2K$.
If the ratio on the left tends to zero, (MC.3) implies
\[
 \kappa_*(t)\geq
 \frac{1+(y-K)^2}
 {3(1+(y+K)^2)}
 \frac{\lambda_{\max}(G)}{\lambda_{\min}(G)}
 \longrightarrow+\infty .
\]
This proves the full asserted obstruction, including metrics whose
overall scale varies.
\end{proof}

\begin{proposition}[A retained-period rank-loss family and its exact cost]
\label{q444:rank-loss-metric}
Fix, without rescaling any matrix, four real constants
\[
 \alpha_{\rm obs}>0,\qquad \beta_{\rm obs}>0,\qquad
 \lambda>0,\qquad p>0,
\]
and set
\[
 \delta(t)=\beta_{\rm obs}y^{-p},\qquad
 G(t)=\begin{pmatrix}\alpha_{\rm obs}&0\\0&\delta(t)\end{pmatrix},
 \qquad H_3(t)=RG(t)R^*+\lambda uu^*.                         \tag{RL.1}
\]
Take the lift (MF.2) with $U=V=0$.  It is a smooth equivariant positive
real lattice metric at every point of the smaller punctured cusp, with
the original periods and correction unchanged.  Its observation has
\[
 \det G(t)=\alpha_{\rm obs}\beta_{\rm obs}y^{-p}
 =\alpha_{\rm obs}\beta_{\rm obs}
       \left(-\frac{\log|t|}{2\pi}\right)^{-p}\longrightarrow0,
 \qquad
 G(t)\longrightarrow\operatorname{diag}(\alpha_{\rm obs},0).
                                                               \tag{RL.2}
\]
Put $D_2(t)=\operatorname{Im}C_2(t)$ and
$Y=Y_{\rm per}(t)=yI_3+\operatorname{Im}C_3(t)$.  The relative metric
has the exact determinant
\[
 \det W_*(t)=
 \frac{\alpha_{\rm obs}^2\lambda^2\beta_{\rm obs}^2\,y^{-2p}}
      {64\,\det(Y)^2},\qquad
 \det(Y)=y\left[
 y^2+y\,\operatorname{tr}(D_2B_0^{-1})+\det D_2\right].        \tag{RL.3}
\]
In particular
\[
 \det W_*(t)\sim
 \frac{\alpha_{\rm obs}^2\lambda^2\beta_{\rm obs}^2}{64}
 y^{-2p-6}.                                                   \tag{RL.4}
\]
If $L_{\max}=\max\{2\alpha_{\rm obs}/3,\,3\lambda\}$, then
\[
 \lambda_{\min}(W_*)\sim\frac{\beta_{\rm obs}}4\,y^{-p-2},
 \qquad
 \lambda_{\max}(W_*)\longrightarrow\frac{L_{\max}}2,
 \qquad
 \kappa_*(t)\sim\frac{2L_{\max}}{\beta_{\rm obs}}\,y^{p+2}.
                                                               \tag{RL.5}
\]
Every constant and the complete correction in the exact formulas
(RL.1)--(RL.3) is retained.
\end{proposition}

\begin{proof}
Formula (MF.5) with $z=0$ and $\tau=\lambda$ gives $H_3\succ0$ and
$PH_3P^*=G$.  Formula (MF.2) proves the metric, smoothness, and monodromy
assertions.  Equation (RL.2) is the direct determinant of the displayed
matrix and uses the original logarithmic-cover factor $2\pi$.
The ratio of its least and greatest eigenvalues is eventually
$\delta/\alpha_{\rm obs}$, so this is projective rank loss in the
precise sense used in Proposition~\ref{q444:metric-condition}.

Write $A_c=\operatorname{Re}C_3$ and
\[
 \widetilde{\mathcal N}=
 \begin{pmatrix}I_3&A_c\\0&Y\end{pmatrix},\qquad
 \mathcal N=\widetilde{\mathcal N}M(x).
\]
Here $H_3$ is real symmetric.  The relative metric is exactly
\[
 W_*=\frac12\widetilde{\mathcal N}^{-1}
       \begin{pmatrix}H_3&0\\0&H_3\end{pmatrix}
       \widetilde{\mathcal N}^{-\mathsf T}
 =\frac12\begin{pmatrix}
 H_3+A_cY^{-1}H_3Y^{-\mathsf T}A_c^{\mathsf T}
       &-A_cY^{-1}H_3Y^{-\mathsf T}\\
 -Y^{-1}H_3Y^{-\mathsf T}A_c^{\mathsf T}
       &Y^{-1}H_3Y^{-\mathsf T}
 \end{pmatrix}.                                               \tag{RL.6}
\]
Since $H_3=X\operatorname{diag}(G,\lambda)X^*$ and $\det X=1$,
$\det H_3=\lambda\alpha_{\rm obs}\delta$.
Taking the real $6\times6$ determinant in (RL.6) gives
$\det W_*=(\det H_3)^2/(2^6\det Y^2)$, proving the first identity
in (RL.3).

The second identity retains the correction rather than replacing $Y$
by $yI_3$.  Indeed
$Y=yI_3+SD_2Q_{\rm latt}$ and $Q_{\rm latt}S=B_0^{-1}$.
For rectangular $U,V$, the identity
$\det(I+UV)=\det(I+VU)$ follows by eliminating either off-diagonal
block of $\left(\begin{smallmatrix}I&U\\-V&I\end{smallmatrix}\right)$
with triangular determinant-one row operations.  Applying it here gives
\[
 \det Y
 =y^3\det\left(I_2+y^{-1}D_2B_0^{-1}\right)
 =y\left[y^2+y\operatorname{tr}(D_2B_0^{-1})+\det D_2\right],
\]
because $\det B_0=1$.  If
$D_2=\left(\begin{smallmatrix}d_{11}&d_{12}\\d_{21}&d_{22}
\end{smallmatrix}\right)$, this is the complete polynomial
\[
 \det Y=y\left[y^2+y(d_{12}-d_{21})
                     +d_{11}d_{22}-d_{12}d_{21}\right].
\]
The entries of $D_2$ remain the prescribed bounded functions of $t$.
Thus $\det Y/y^3\to1$, proving (RL.4) with the displayed coefficient.

For the eigenvalue assertions, $H_3u=3\lambda u$ because $R^*u=0$.
On the two-dimensional invariant plane $u^\perp$, its trace is
\[
 \operatorname{tr}(C^{-1}G)
      =\frac23(\alpha_{\rm obs}+\delta)
\]
and the product of the two eigenvalues is
$\det H_3/(3\lambda)=\alpha_{\rm obs}\delta/3$.
The invariance of the plane follows also from Hermiticity and the
displayed $u$ eigenvector.  Its two eigenvalues are therefore exactly
\[
 h_\pm=
 \frac{\alpha_{\rm obs}+\delta
       \pm\sqrt{\alpha_{\rm obs}^2-\alpha_{\rm obs}\delta+\delta^2}}3.
                                                               \tag{RL.7}
\]
They are positive: their sum and product are positive.  Rationalizing
the numerator of $h_-$ gives
\[
 \frac{h_-}{\delta}
 =\frac{\alpha_{\rm obs}}
 {\alpha_{\rm obs}+\delta+
       \sqrt{\alpha_{\rm obs}^2-\alpha_{\rm obs}\delta+\delta^2}}
 \longrightarrow\frac12.
\]
Also $h_+\to2\alpha_{\rm obs}/3$.
Consequently $\lambda_{\min}(H_3)\sim\delta/2$ and
$\lambda_{\max}(H_3)\to L_{\max}$, with $H_3$ bounded.

The explicit blocks (RL.6), the bound
$\|Y^{-1}\|\leq(y-K)^{-1}$, and bounded $A_c,H_3$ imply
\[
 \left\|W_*-\tfrac12\operatorname{diag}(H_3,0)\right\|
 =O(y^{-2}).
\]
For real symmetric matrices the maximum eigenvalue changes by at most
the operator norm of a perturbation: apply the Rayleigh-quotient maximum
to each matrix and subtract.  Hence
$\lambda_{\max}(W_*)\to L_{\max}/2$.

To control the least eigenvalue without assuming a bounded condition
number of $H_3$, use the exact inverse
\[
 W_*^{-1}
 =2\widetilde{\mathcal N}^{\mathsf T}
      \operatorname{diag}(H_3^{-1},H_3^{-1})
      \widetilde{\mathcal N}.
\]
For $y>1$, the decomposition
$\widetilde{\mathcal N}=\operatorname{diag}(I_3,yI_3)
+\left(\begin{smallmatrix}0&A_c\\0&\operatorname{Im}C_3\end{smallmatrix}\right)$
gives $\|\widetilde{\mathcal N}\|\leq y+\sqrt2K$.
Thus
\[
 \lambda_{\max}(W_*^{-1})
 \leq\frac{2(y+\sqrt2K)^2}{\lambda_{\min}(H_3)}.
\]
Its lower-right principal block includes the positive term
$2Y^{\mathsf T}H_3^{-1}Y$.  If $v$ is a unit eigenvector of
$H_3^{-1}$ for its greatest eigenvalue, testing that block on
$Y^{-1}v/\|Y^{-1}v\|$ gives
\[
 \lambda_{\max}(W_*^{-1})
 \geq\frac{2(y-K)^2}{\lambda_{\min}(H_3)}.
\]
All denominators are positive.  Taking reciprocals yields
\[
 \frac{\lambda_{\min}(H_3)}{2(y+\sqrt2K)^2}
 \leq\lambda_{\min}(W_*)
 \leq\frac{\lambda_{\min}(H_3)}{2(y-K)^2}.                     \tag{RL.8}
\]
Both bounds are asymptotic to $\delta/(4y^2)$, proving the first
assertion of (RL.5).  Combining it with the maximum-eigenvalue limit
proves the condition-number asymptotic with its full coefficient.
\end{proof}

This is metric-choice freedom, not a selection of a physical momentum
or mass.  For the same unchanged periods the allowed choice $W_*=I_6$
instead gives exactly
\[
 H_3=I_3+(iyI_3+C_3)(iyI_3+C_3)^*,\qquad
 G=y^2C+O(y),\qquad \det G=3y^4+O(y^3),
\]
whereas (RL.1)--(RL.2) give a determinant tending to zero and retain the
same $C_2,C_3$.  These two explicit admissible choices, the complete
twelve-dimensional metric fibres, and the seventeen-dimensional observed
fibres prove that the specified period data and positivity alone do not
select an observed trajectory.  The eigenvalue ratios above only diagnose
rank loss; no field, period, metric, or parameter was replaced by a
normalized one.  Units, a metric-selection law, dynamics, and an
interpretation of the marked observation would be additional specified
data.  No extension across other finite-order overlaps or special
fillings, and no physical or quantum gap conclusion, is asserted here.

% Insertion-ready; uses the original P,R,u,C,X,B_0,S,Q_latt,C_2,C_3.
% Independent lane: no canonical file was edited.
\subsection*{Hermitian flat metrics, the original period lattice, and its dual spectrum}

Here a flat metric means a translation-invariant metric on each fixed
torus fibre. Its coefficients may depend smoothly on the cusp parameter
$t$, but not on the point of that fibre. The word \emph{physical} below
distinguishes the original real coordinates of $\mathbb C^3$ from integral
lattice coordinates; it does not assert a physical field theory.
Retain, without changing any basis or period,
\[
 B_0=\begin{pmatrix}0&1\\-1&0\end{pmatrix},\quad
 P=\begin{pmatrix}1&-1&0\\0&1&-1\end{pmatrix},\quad
 Q_{\rm latt}=B_0^{-1}P,\quad
 S=\begin{pmatrix}1&1\\0&1\\0&0\end{pmatrix},
\]
\[
 C_3=SC_2Q_{\rm latt},\qquad
 Z_3=sI_3+C_3,\qquad Z_2=sB_0+C_2,\qquad
 s=x+iy,\quad t=e^{2\pi i s},\quad
 y=-\frac{\log|t|}{2\pi}.
\]
Choose finite bounds $K_2,K\geq0$ on $\|C_2(t)\|,\|C_3(t)\|$,
respectively, and restrict to $y>\max\{K_2,K\}+1$. Norms are the
original Euclidean operator norms. Write
\[
 A=xI_3+\operatorname{Re}C_3,\quad
 Y=yI_3+\operatorname{Im}C_3,\quad
 \mathcal N=\begin{pmatrix}I_3&A\\0&Y\end{pmatrix},\quad
 \mathcal J=(I_3\;\;iI_3),\quad
 J_0=\begin{pmatrix}0&-I_3\\I_3&0\end{pmatrix}.
\]
Thus $\Omega_3=(I_3\;\;Z_3)=\mathcal J\mathcal N$ and
$\mathcal N^{-1}=\left(\begin{smallmatrix}
I_3&-AY^{-1}\\0&Y^{-1}\end{smallmatrix}\right)$.
Indeed $\|Yv\|\geq(y-K)\|v\|$ proves invertibility. The same argument,
using $\|B_0v\|=\|v\|$, proves invertibility of
$Y_2=\operatorname{Im}Z_2=yB_0+\operatorname{Im}C_2$.
For $r=2,3$ put
\[
 \Lambda_r=\Omega_r\mathbb Z^{2r},\qquad
 T_r=\mathbb C^r/\Lambda_r,\qquad
 \mathcal N_2=
 \begin{pmatrix}I_2&\operatorname{Re}Z_2\\0&Y_2\end{pmatrix}.
\]
The real coordinates of $z\in\mathbb C^3$ are
$q=(\operatorname{Re}z,\operatorname{Im}z)^{\mathsf T}
=\mathcal N\xi$, where $\xi\in\mathbb R^6/\mathbb Z^6$.
This is an exact change of coordinates, not a replacement of the lattice.
For any complex matrix $D$ define its realification by
\[
 \mathfrak r(D)=
 \begin{pmatrix}\operatorname{Re}D&-\operatorname{Im}D\\
                 \operatorname{Im}D&\operatorname{Re}D\end{pmatrix}.
\]
Block multiplication gives
$\mathfrak r(D_1D_2)=\mathfrak r(D_1)\mathfrak r(D_2)$ and
$\mathfrak r(D^*)=\mathfrak r(D)^{\mathsf T}$.

\begin{theorem}[The precise Hermitian metric in a period fibre]
\label{q444:hermitian-cometric}
Let $g_{\rm phys}\succ0$ be a real symmetric metric matrix and put
\[
 T=g_{\rm phys}^{-1},\qquad
 g_{\rm latt}=\mathcal N^{\mathsf T}g_{\rm phys}\mathcal N,\qquad
 W=g_{\rm latt}^{-1}=\mathcal N^{-1}T\mathcal N^{-\mathsf T}.
                                                               \tag{HS.1}
\]
Then the earlier period weight $W$ has the exact geometric meaning of a
\emph{lattice cometric}; the covariant lattice metric is its inverse.
The observation is $H=\Omega_3W\Omega_3^*=\mathcal JT\mathcal J^*$.
For a fixed positive Hermitian $H=E+iF$, its full real cometric fibre is
\[
 T=\frac12\mathfrak r(H)+
       \mathcal A(U,V),\qquad
 \mathcal A(U,V)=\begin{pmatrix}U&V\\V&-U\end{pmatrix},\qquad
 U=U^{\mathsf T},\quad V=V^{\mathsf T},\quad T\succ0.              \tag{HS.2}
\]
The real metric is compatible with the given complex structure, meaning
$J_0^{\mathsf T}g_{\rm phys}J_0=g_{\rm phys}$, if and only if
$U=V=0$. With the explicit Hermitian metric convention
\[
 g_{\rm phys}(v,w)=\operatorname{Re}(v^*h_3w),
       \qquad v,w\in\mathbb C^3,
\]
the unique compatible member over $H$ is
\[
 h_3=2H^{-1},\qquad h_3^{-1}=H/2,\qquad
 g_{\rm phys}=2\mathfrak r(H^{-1}),\qquad
 T=\tfrac12\mathfrak r(H),                                    \tag{HS.3}
\]
\[
 g_{\rm latt}=2\mathcal N^{\mathsf T}
                    \mathfrak r(H^{-1})\mathcal N,\qquad
 W=\tfrac12\mathcal N^{-1}\mathfrak r(H)\mathcal N^{-\mathsf T}.
                                                               \tag{HS.4}
\]
For an arbitrary positive real cometric, its observed Hermitian part is
precisely the positive $J_0$-invariant average
\[
 \frac{T+J_0TJ_0^{\mathsf T}}2=\frac12\mathfrak r(H).             \tag{HS.5}
\]
\end{theorem}
\begin{proof}
Substitution of $q=\mathcal N\xi$ into the covariant quadratic form
gives the middle formula in (HS.1); multiplication verifies its inverse.
The observation is positive: $\mathcal J^*$ is injective, and a real
positive matrix $T$ is also positive on complex vectors, since for
$a,b$ real, $(a+ib)^*T(a+ib)=a^{\mathsf T}Ta+b^{\mathsf T}Tb$.
For a real symmetric block matrix
$T=\left(\begin{smallmatrix}D&B\\B^{\mathsf T}&L\end{smallmatrix}\right)$,
\[
 \mathcal JT\mathcal J^*=D+L+i(B^{\mathsf T}-B).
\]
Thus equality to $E+iF$ is equivalent to
$D=E/2+U$, $L=E/2-U$, $B=-F/2+V$ with $U,V$ symmetric.
This proves both directions and uniqueness of the parameters in (HS.2).

Because $J_0^{\mathsf T}=-J_0$ and $J_0^2=-I_6$, metric compatibility
is equivalent to $g_{\rm phys}J_0=J_0g_{\rm phys}$, which is equivalent,
by multiplying by the inverse on both sides, to $TJ_0=J_0T$.
Direct multiplication shows that $\mathfrak r(H)$ commutes with $J_0$,
whereas $\mathcal A J_0=-J_0\mathcal A$. Hence compatibility is equivalent
to $J_0\mathcal A=0$, and the invertibility of $J_0$ makes this equivalent
to $\mathcal A=0$, or $U=V=0$.

For $z=v+iw$ one has
$(v,w)^{\mathsf T}\mathfrak r(H)(v,w)=z^*Hz>0$ unless $z=0$.
Thus $\mathfrak r(H)$ is positive. Realification respects products,
so inversion of $T=\mathfrak r(H)/2$ gives (HS.3), and (HS.1) gives
(HS.4). Finally
$J_0\mathfrak r(H)J_0^{\mathsf T}=\mathfrak r(H)$ and
$J_0\mathcal AJ_0^{\mathsf T}=-\mathcal A$, proving (HS.5).
This average is an average of cometrics, not an average of covariant
metrics. The metric associated to it is the inverse of the average.
\end{proof}

Thus the twelve real fibre coordinates in (HS.2) describe genuine
non-Hermitian translation-invariant real metrics, and compatibility
removes exactly those coordinates. For example, $H=2I_3$ and
$U=\varepsilon\operatorname{diag}(1,0,0)$, $V=0$ give
\[
 T_\varepsilon=\operatorname{diag}
 (1+\varepsilon,1,1,1-\varepsilon,1,1),\qquad -1<\varepsilon<1.
\]
All have the same $H$; only $\varepsilon=0$ is Hermitian-compatible.
They have different marked cometrics, and the next theorem proves
that some marked character eigenvalue therefore changes.

\begin{theorem}[The lattice cometric is the marked character spectrum]
\label{q444:torus-laplace}
For any $T\succ0$ as in (HS.1), define the nonnegative Laplacian by
$\Delta=-\operatorname{div}\operatorname{grad}$. Its complete character
eigenbasis and eigenvalues on $T_3$ are
\[
 \chi_k([\xi])=e^{2\pi i k^{\mathsf T}\xi},\qquad k\in\mathbb Z^6,
 \qquad
 \lambda_k=4\pi^2 k^{\mathsf T}Wk
          =4\pi^2 p_k^{\mathsf T}Tp_k,\quad
 p_k=\mathcal N^{-\mathsf T}k.                               \tag{HS.6}
\]
Here $p_k$ is in the integer-valued dual of
$\mathcal N\mathbb Z^6$: its pairing with $\mathcal Nn$ is
$k^{\mathsf T}n\in\mathbb Z$. Conversely every integer-valued dual
covector is of this form. The eigenspace multiplicity is exactly
the number of integer $k$ with the specified value of (HS.6).
The volume, in the original scales, is
\[
 \operatorname{Vol}(T_3,g_{\rm phys})
 =|\det Y|\sqrt{\det g_{\rm phys}}
 =|\det Y|\det h_3
 =\frac{8|\det Y|}{\det H}                                  \tag{HS.7}
\]
in the compatible case. Without compatibility the first equality
still holds.

The marked spectrum recovers every entry of $W$:
\[
 W_{aa}=\frac{\lambda_{e_a}}{4\pi^2},\qquad
 W_{ab}=
 \frac{\lambda_{e_a+e_b}-\lambda_{e_a}-\lambda_{e_b}}{8\pi^2}
 \quad(a\ne b).                                             \tag{HS.8}
\]
If $H(t)$ is single-valued and (HS.4) is used, all these metrics
descend with the original monodromy
\[
 \mathcal N(s+1)=\mathcal N(s)M_3,\quad
 M_3=\begin{pmatrix}I_3&I_3\\0&I_3\end{pmatrix},\quad
 W(s+1)=M_3^{-1}W(s)M_3^{-\mathsf T}.                         \tag{HS.9}
\]
The same character has labels $k$ and $k'=M_3^{\mathsf T}k$
in the two markings, and $\lambda'_{k'}=\lambda_k$.
\end{theorem}
\begin{proof}
The covariant metric and its determinant are constant in the fibre
coordinate $\xi$, so the defining divergence formula becomes
$\Delta=-\sum_{a,b=1}^6W_{ab}\partial_{\xi_a}\partial_{\xi_b}$.
Two derivatives of $\chi_k$ prove (HS.6). The dual-lattice assertion
follows from $p^{\mathsf T}\mathcal Nn\in\mathbb Z$ for all integer
$n$: taking the six unit vectors is equivalent to
$\mathcal N^{\mathsf T}p\in\mathbb Z^6$.

For completeness, no unlisted eigenfunctions are needed. Integration
over $[0,1]^6$ gives zero for $\chi_k\overline{\chi_l}$ when $k\ne l$;
the common squared norm in Riemannian measure is the volume.
To prove density, the one-dimensional Fejer kernels
\[
 F_M(\theta)=\frac1M\left|\sum_{r=0}^{M-1}e^{2\pi i r\theta}\right|^2
 =\sum_{|j|<M}\left(1-\frac{|j|}{M}\right)e^{2\pi i j\theta}
\]
are nonnegative and have integral one. For
$\operatorname{dist}(\theta,\mathbb Z)\geq\eta>0$ they are bounded by
$1/(M\sin^2(\pi\eta))$, using the finite geometric sum. The product
of six such kernels also has integral one; its mass outside a
coordinate cube of radius $\eta$ tends to zero by the union bound
over its six coordinates. Uniform continuity then proves that
convolution of any continuous periodic function with these product
kernels converges uniformly to that function. Each convolution is a
finite linear combination of characters. Continuous functions are dense
in $L^2([0,1]^6)$: step functions on finite unions of boxes are dense
by regular approximation in Lebesgue measure, and a box indicator is
approximated in $L^2$ by a continuous function that differs only in a
boundary strip of arbitrarily small measure. Multiplication by the
constant Riemannian density preserves this argument.

Consequently the characters, divided by the square root of the volume,
are a complete orthonormal basis. On Fourier coefficients the closure of
$\Delta$ is the real diagonal operator with domain
$\{(a_k):\sum_k\lambda_k^2|a_k|^2<\infty\}$; its adjoint has that same
domain, as testing against each coordinate vector forces the displayed
weighted square summability. It is therefore self-adjoint.
Finite Fourier sums are a core by truncation of both the ordinary
and weighted square sums. For a smooth function, integration by parts
identifies the Fourier coefficients of its smooth Laplacian with
$\lambda_k a_k$; their square sum is finite. Thus the operator on all
smooth functions lies in this diagonal operator and has the same
closure as its restriction to finite Fourier sums. Since
$\lambda_k\geq4\pi^2\lambda_{\min}(W)\|k\|^2$, bounded eigenvalue
intervals contain only finitely many integer indices. This proves
the completeness, discreteness and multiplicity assertions, including
a positive smallest eigenvalue away from $k=0$.

The parallelepiped $\mathcal N[0,1]^6$ has coordinate volume
$|\det\mathcal N|=|\det Y|$, proving the first equality in (HS.7).
For Hermitian $h_3$, the determinant identity
$\det\mathfrak r(h_3)=|\det h_3|^2=(\det h_3)^2$
follows by conjugating the realification over $\mathbb C$ to
$\operatorname{diag}(h_3,\overline{h_3})$ with the invertible change
$(v,w)\mapsto(v+iw,v-iw)$. Positive $h_3$ has positive determinant.
Now $\det(2H^{-1})=8/\det H$, proving (HS.7).
Expansion of the quadratic form at $e_a+e_b$ proves (HS.8).

The correction $C_3(t)$ is unchanged by $s\mapsto s+1$, giving the
first formula in (HS.9). The cometric formula gives the last.
Since $q=\mathcal N\xi=\mathcal N M_3\xi'$, one has
$\xi=M_3\xi'$. Hence $k'=M_3^{\mathsf T}k$ and substitution in
(HS.6) proves the final assertion. All matrix inversions depend
smoothly on the entries where their determinants are nonzero, so a
smooth $H(t)$ gives the asserted smooth descending family.
\end{proof}

\begin{theorem}[The original torus quotient is a metric and spectral map]
\label{q444:metric-submersion}
Let $H_3\succ0$, $G=PH_3P^*$, and equip $T_3$ with
$h_3=2H_3^{-1}$. Retain the original integral map
\[
 L=P\oplus Q_{\rm latt}:\mathbb Z^6\longrightarrow\mathbb Z^4,
 \qquad P_{\mathbb R}=\operatorname{diag}(P,P).
\]
Then $f_P:T_3\to T_2$, $[z]\mapsto[Pz]$, is a surjective holomorphic
homomorphism with connected kernel
\[
 \mathbb C u/(\mathbb Z u+s\mathbb Z u),\qquad
 u=(1,1,1)^{\mathsf T}.                                    \tag{HS.10}
\]
There is a unique target translation-invariant metric making this map
a Riemannian submersion. Its Hermitian metric, complex cometric, real
cometric and lattice cometric are, respectively,
\[
 h_2=2G^{-1},\qquad h_2^{-1}=G/2=Ph_3^{-1}P^*,\qquad
 \mathsf T_2=\tfrac12\mathfrak r(G)
       =P_{\mathbb R}\mathsf T_3P_{\mathbb R}^{\mathsf T},
                                                               \tag{HS.11}
\]
\[
 W_2=\mathcal N_2^{-1}\mathsf T_2\mathcal N_2^{-\mathsf T}
     =LW_3L^{\mathsf T}.                                    \tag{HS.12}
\]
Here $\mathsf T_3=\mathfrak r(H_3)/2$ is the real source cometric,
distinguished from the torus $T_3$. The horizontal lift is the exact complex map
\[
 \eta:\mathbb C^2\longrightarrow\mathbb C^3,\qquad
 \eta=H_3P^*G^{-1}.                                         \tag{HS.13}
\]
For every smooth function $f$ on the target,
\[
 \Delta_3(f\circ f_P)=(\Delta_2f)\circ f_P.                    \tag{HS.14}
\]
For a target character $\ell\in\mathbb Z^4$, its pullback is the source
character labelled by $L^{\mathsf T}\ell$, and its eigenvalue is unchanged.
These are exactly the source characters constant on the fibres. Thus
the target spectrum, with each character's multiplicity counted, is
a subcollection of the source character spectrum, not an identification
with all source modes.
\end{theorem}
\begin{proof}
Direct multiplication gives $PS=I_2$, $B_0Q_{\rm latt}=P$ and
$PZ_3=Z_2Q_{\rm latt}$. Thus
\[
 P\Omega_3=\Omega_2L,\qquad
 P_{\mathbb R}\mathcal N=\mathcal N_2L.                        \tag{HS.15}
\]
Moreover $\operatorname{diag}(S,SB_0)$ is an integral right inverse
of $L$, since $Q_{\rm latt}SB_0=I_2$. Therefore
$P\Lambda_3=\Lambda_2$, proving that the torus homomorphism is
well-defined and surjective.

If $Pz\in\Lambda_2$, surjectivity on lattices supplies
$\gamma\in\Lambda_3$ with $P\gamma=Pz$, so $z-\gamma\in\ker_{\mathbb C}P
=\mathbb C u$. To find its periods, if $\Omega_3(a,b)$ is in
$\mathbb C u$, then $\Omega_2(Pa,Q_{\rm latt}b)=0$.
The invertibility of $\mathcal N_2$ forces $Pa=Q_{\rm latt}b=0$.
Solving the two displayed difference equations gives
$a=ru$, $b=vu$ with $r,v\in\mathbb Z$. Since $C_3u=0$,
$\Omega_3(a,b)=(r+sv)u$. This proves (HS.10), including connectedness.
Its lattice has real rank two because $y>0$.

Put $K_3=h_3^{-1}=H_3/2$, $K_2=PK_3P^*=G/2$.
Surjectivity of $P$ makes $P^*$ injective, so $K_2\succ0$.
The lift in (HS.13) is $K_3P^*K_2^{-1}$ and satisfies
\[
 P\eta=I_2,\qquad
 \eta^*h_3 v=K_2^{-1}Pv=0\quad(v\in\ker P),\qquad
 \eta^*h_3\eta=K_2^{-1}=h_2.                                \tag{HS.16}
\]
Every $w\in\mathbb C^3$ has the orthogonal decomposition
$w=\eta(Pw)+(w-\eta(Pw))$, whose second term lies in $\ker P$.
Hence $P$ restricts to an isometry on the orthogonal complement of
its kernel; this is precisely the definition of a Riemannian
submersion. It also proves uniqueness of $h_2$, since each target
vector has a unique horizontal lift. Taking realifications proves
(HS.11), and multiplying (HS.15) by the inverses on both sides proves
(HS.12).

In lattice coordinates $f_P([\xi])=[L\xi]$. The chain rule gives
$\partial_{\xi_a}\partial_{\xi_b}f(L\xi)
=\sum_{i,j}L_{ia}L_{jb}(\partial_i\partial_jf)(L\xi)$.
Contraction with $-W_{3,ab}$ and (HS.12) proves (HS.14).
Also $e^{2\pi i\ell^{\mathsf T}L\xi}=\chi_{L^{\mathsf T}\ell}(\xi)$;
inserting this index in (HS.6) proves equality of the eigenvalues.

The lattice kernel is $\mathbb Z u\oplus\mathbb Z u$.
The characters constant on the connected kernel annihilate its real
tangent space, so their indices $(m,n)$ obey
$u^{\mathsf T}m=u^{\mathsf T}n=0$. Conversely these equalities make
the character constant there. Every integral zero-sum triple
$m=(m_1,m_2,m_3)$ is
$P^{\mathsf T}(m_1,-m_3)^{\mathsf T}$, and
$\operatorname{im}_{\mathbb Z}Q_{\rm latt}^{\mathsf T}
=\operatorname{im}_{\mathbb Z}P^{\mathsf T}$ because $B_0$ is
unimodular. Thus the stated indices are exactly
$L^{\mathsf T}\mathbb Z^4$. Injectivity of $L^{\mathsf T}$
preserves the number of target characters at each eigenvalue.
\end{proof}

The compatibility restriction still leaves genuine source-metric
freedom over a fixed target metric. Indeed, using the original
$R=P^*C^{-1}$, $C=PP^*$ and $X=(R\;\;u)$, all positive $H_3$
with $PH_3P^*=G$ are
\[
 H_3=X\begin{pmatrix}G&z\\z^*&z^*G^{-1}z+\tau\end{pmatrix}X^*,
 \qquad z\in\mathbb C^2,\quad \tau>0.                         \tag{HS.17}
\]
To prove the assertion, $PX=(I_2\;\;0)$ and
$X^{-1}=\left(\begin{smallmatrix}P\\u^*/3\end{smallmatrix}\right)$
show that the top-left block of $X^{-1}H_3X^{-*}$ is $G$.
For any remaining blocks $z,d$, its quadratic form is
\[
 (v+G^{-1}zw)^*G(v+G^{-1}zw)
       +(d-z^*G^{-1}z)|w|^2.
\]
It is positive exactly when $d=z^*G^{-1}z+\tau$ with $\tau>0$.
This proves (HS.17) in both directions and its five real parameters.
The injectivity of inversion makes distinct $H_3$ distinct compatible
source metrics; (HS.11)--(HS.14) prove that the quotient metric and
quotient spectrum stay fixed. This specifies exactly what the
remaining observation fibre means geometrically.

\begin{proposition}[Native Euclidean observation and expanding dual spectrum]
\label{q444:native-period-spectrum}
Equip the source with its native Euclidean metric in the retained
coordinates: $h_3=I_3$, $g_{\rm phys}=I_6$. Then
\[
 W_3=\mathcal N^{-1}\mathcal N^{-\mathsf T},\qquad
 H_3=2I_3,\qquad G=2C,\qquad
 C=\begin{pmatrix}2&-1\\-1&2\end{pmatrix}.                     \tag{HS.18}
\]
The induced target metric is $h_2=C^{-1}$, not $I_2$.
The source eigenvalues, with no correction suppressed, are
\[
 \lambda_{m,n}
 =4\pi^2\left(\|m\|^2+
       \|Y^{-\mathsf T}(n-A^{\mathsf T}m)\|^2\right),
       \qquad m,n\in\mathbb Z^3.                            \tag{HS.19}
\]
For $y>K+1$, its smallest positive eigenvalue is exactly
\[
 \lambda_1^{(3)}
 =4\pi^2\min_{0\ne n\in\mathbb Z^3}\|Y^{-\mathsf T}n\|^2,
 \qquad
 \frac{4\pi^2}{(y+K)^2}\leq\lambda_1^{(3)}
                  \leq\frac{4\pi^2}{(y-K)^2}.                \tag{HS.20}
\]
For the induced target metric, its smallest positive eigenvalue is
\[
 \lambda_1^{(2)}
 =4\pi^2\min_{\substack{0\ne n\in\mathbb Z^3\\u^{\mathsf T}n=0}}
                \|Y^{-\mathsf T}n\|^2,\qquad
 \frac{8\pi^2}{(y+K)^2}\leq\lambda_1^{(2)}
                  \leq\frac{8\pi^2}{(y-K)^2}.                \tag{HS.21}
\]
In particular the exact leading constants are
\[
 \lambda_1^{(3)}\sim\frac{4\pi^2}{y^2}
   =\frac{16\pi^4}{(\log|t|)^2},\qquad
 \lambda_1^{(2)}\sim\frac{8\pi^2}{y^2}
   =\frac{32\pi^4}{(\log|t|)^2}.                              \tag{HS.22}
\]
The source volume is the retained polynomial
\[
 \operatorname{Vol}(T_3,I)
 =\det Y
 =y\left[y^2+
 y\operatorname{tr}\bigl((\operatorname{Im}C_2)B_0^{-1}\bigr)
 +\det(\operatorname{Im}C_2)\right]\sim y^3.                  \tag{HS.23}
\]
\end{proposition}
\begin{proof}
In (HS.1) the physical cometric is $T=I_6$, so
$H_3=\mathcal J\mathcal J^*=2I_3$. Multiplication by $P,P^*$
and (HS.11) prove all of (HS.18). The exact inverse transpose is
\[
 \mathcal N^{-\mathsf T}
 =\begin{pmatrix}I_3&0\\-Y^{-\mathsf T}A^{\mathsf T}
                              &Y^{-\mathsf T}\end{pmatrix},
\]
which gives (HS.19). The singular values of $Y$ lie between
$y-K$ and $y+K$, since
$(y-K)\|v\|\leq\|Yv\|\leq(y+K)\|v\|$.
The same singular-value bounds for the inverse transpose give
\[
 \frac{\|n\|}{y+K}\leq\|Y^{-\mathsf T}n\|
                             \leq\frac{\|n\|}{y-K}.           \tag{HS.24}
\]
For source modes with $m\ne0$, (HS.19) is at least $4\pi^2$.
The mode $m=0,n=e_1$ has eigenvalue at most
$4\pi^2/(y-K)^2<4\pi^2$. Thus a minimizer has $m=0$.
The minimum exists because the inverse-transpose quadratic form is
positive definite and only finitely many integer vectors have bounded
value. Nonzero integer vectors have norm at least one, and a unit
vector exists, giving (HS.20).

Target modes have $m=P^{\mathsf T}a$, $n=Q_{\rm latt}^{\mathsf T}b$
for $a,b\in\mathbb Z^2$. A nonzero integer zero-sum triple has
squared norm at least two: it has at least one positive and at least
one negative integer coordinate. The value two is attained by
$(1,-1,0)^{\mathsf T}$. Therefore $a\ne0$ costs at least $8\pi^2$,
whereas a mode with $a=0$ and a shortest nonzero zero-sum $n$ has
eigenvalue at most $8\pi^2/(y-K)^2<8\pi^2$.
The complete index characterization in the preceding theorem now
proves (HS.21). Multiplying the bounds by $y^2$ and taking the limit
proves (HS.22), with the original definition of $y$.

Finally $D_2=\operatorname{Im}C_2$ is real and
$Y=yI_3+SD_2Q_{\rm latt}$. Since $Q_{\rm latt}S=B_0^{-1}$,
the identity
$\det(I+UV)=\det(I+VU)$ gives
\[
 \det Y=y^3\det(I_2+y^{-1}D_2B_0^{-1}).
\]
For completeness that determinant identity follows by taking the
determinant of
$\left(\begin{smallmatrix}I&U\\-V&I\end{smallmatrix}\right)$
and eliminating first the upper-right block or first the lower-left
block using block triangular matrices with determinant one.
Expansion of the two-by-two determinant and $\det B_0=1$ give
the displayed polynomial in (HS.23). Its sign is positive:
$yI_3+rSD_2Q_{\rm latt}$ is invertible for $0\leq r\leq1$ by
the same bound $y-rK>0$, so its determinant has the positive sign
of $y^3$. Formula (HS.7) and boundedness of $D_2$ finish the proof.
\end{proof}

Equations (HS.18)--(HS.23) exhibit the exact relationship rather than
a nonidentity alone: the constant observation is twice the native
physical cometric and twice its quotient cometric, while the unchanged
period map enters the dual lattice in (HS.19). The periods expand,
the volume grows, and both displayed positive spectral gaps tend to
zero. No period coefficient was replaced and no metric scale was
divided out. At fixed periods, $H_3$ uniquely specifies the
Hermitian-compatible source metric; at fixed $G$ the five-parameter
family (HS.17) specifies its precise residual freedom. When periods
are also allowed to vary, the constant $H_3=2I_3$ alone cannot recover
their marked cometric; (HS.8) proves that the full marked spectrum
does recover it. These are geometric and spectral assertions, not
a selection of a preferred dynamical metric or a field-theoretic mass.

\paragraph{Primary reference conventions.}
The nonnegative Laplacian and integer-valued dual-lattice convention
agree with A.~Gogolev, \emph{Spectra of flat tori}, Section~2
(27 November 2017), available at
\url{https://andreygogolev.wordpress.com/2017/11/27/spectra-of-flat-tori/}.
V.~Ivrii, \emph{Partial Differential Equations}, Appendix~4.B,
\url{https://www.math.utoronto.ca/courses/apm346h1/20181/PDE-textbook/Chapter4/S4.B.html},
instead places $2\pi$ in the definition of the dual lattice;
his dual is exactly $2\pi\mathcal N^{-\mathsf T}\mathbb Z^6$ here.
The horizontal-isometry definition and uniqueness of quotient metrics
are those in U.~Lang, \emph{Lecture Notes on Riemannian Geometry},
16 June 2020, Definition~4.1 and Proposition~4.4, printed pages~45
and~47,
\url{https://metaphor.ethz.ch/x/2020/fs/401-3532-08L/sc/DG2_16June2020.pdf}.
All specialized matrix identities and period/spectral asymptotics
used here have been proved above in the retained coordinates.

\subsection{The unrestricted and constrained scalar equations}

Equation~\eqref{q444:N15} gives the exact normal covector at $\rho=0$:
\begin{equation}\label{q444:normal-covector}
 d_{\rho}N_{15}(\delta\rho)
 =-2\operatorname{Re}(\delta\rho\,\chi^*z).
\end{equation}
It need not vanish. For $G=I_2$, $z=\chi=(1,0)^{\mathsf T}$,
$\lambda=3$, $\rho=t\in\mathbb R$,
\[
 N_{15}(t)=1-2t-3t^2.
\]
At $t=0$ this is a positive-definite quaternionic matrix:
$w=(0,1)^{\mathsf T}$ and the Schur excess is $3-1-1=1$.
Thus the nonzero normal covector occurs even in the interior of the
positive cone. This alone is a statement about the cubic; to determine
motion one must also retain the specified kinetic form.

\begin{theorem}[Normal force in the supplied scalar model]\label{q444:force}
Use the triality packet's fourteen-real-dimensional scalar variables and
potential
\begin{equation}\label{q444:potential}
 \begin{gathered}
 \Phi=\Phi^*,\quad\operatorname{tr}\Phi=0,\quad b\in\mathbb C^3,
 \quad s=\operatorname{tr}\Phi^2+2b^*b,\\
 V(\Phi,b)=\frac{m^2}{2}s+\frac{\Gamma}{4}s^2
                       -\kappa(\det\Phi-b^*\Phi b),\\
 L=\tfrac12\operatorname{tr}(\dot\Phi^2)+\dot b^*\dot b-V,
 \quad m^2\geq0,\quad\Gamma>0,\quad\kappa\in\mathbb R.
 \end{gathered}
\end{equation}
This is the explicitly specified spatially constant scalar mechanics with
gauge fields omitted, not a reduction claimed for arbitrary gauge fields.
For its solutions, with $\rho=u^*b$,
\begin{equation}\label{q444:rho-eom}
 \ddot b=-[(m^2+\Gamma s)I_3+\kappa\Phi]b,
 \qquad
 \ddot\rho=-(m^2+\Gamma s)\rho-\kappa u^*\Phi b.
\end{equation}
The phase-space constraint $\rho=\dot\rho=0$ is not invariant for
$\kappa\neq0$. At initial data
\[
 \Phi_0=\operatorname{diag}(1,-1,0),\quad
 b_0=(1,-1,0)^{\mathsf T},\quad \dot\Phi_0=0,\quad\dot b_0=0,
\]
one has $\ddot\rho(0)=-2\kappa$.
\end{theorem}
\begin{proof}
For a variation $\delta b$,
\[
 \delta_bV=2\operatorname{Re}
       \bigl(\delta b^*[(m^2+\Gamma s)I_3+\kappa\Phi]b\bigr).
\]
The kinetic variation has the same factor two, giving the first equation
in \eqref{q444:rho-eom}; multiply by $u^*$ for the second.
At the displayed initial data $u^*b_0=0$ but
$\Phi_0b_0=(1,1,0)^{\mathsf T}$, so the acceleration is $-2\kappa$.
The equations have polynomial right-hand sides and hence a unique local
solution for these initial data. Taylor's theorem gives
$\rho(t)=-\kappa t^2+o(t^2)$, which leaves the constraint.
For comparison with the marked cubic, if
$X^{-1}\Phi X^{-*}=\left(\begin{smallmatrix}G&z\\z^*&\lambda\end{smallmatrix}\right)$,
then
\[
 \operatorname{tr}(C^{-1}G)+3\lambda=0,\qquad
 u^*\Phi b=3(z^*\chi+\lambda\rho).
\]
Thus at $\rho=0$ the normal acceleration is $-3\kappa z^*\chi$.
This factor three is fixed by the marking and the kinetic form; it is not
removed by rewriting the cubic.
\end{proof}

\begin{proposition}[Exact constrained alternative]\label{q444:constrained}
Let $\Pi_u=uu^*/3$, $Q_u=I_3-\Pi_u$. Restrict the Lagrangian
\eqref{q444:potential} to $b\in u^\perp$ and velocities in $u^\perp$.
Its Euler--Lagrange equations are
\begin{align}\label{q444:projected}
 \ddot b&=-(m^2+\Gamma s)b-\kappa Q_u\Phi b,\notag\\
 \ddot\Phi&=-(m^2+\Gamma s)\Phi
       +\kappa\bigl(\operatorname{adj}\Phi-bb^*\bigr)_0,
 \qquad T_0=T-\tfrac13(\operatorname{tr}T)I_3.
\end{align}
The reaction force added to the unrestricted $b$ equation is exactly
$\kappa\Pi_u\Phi b$. The constrained solution is also an unrestricted
solution exactly when this force vanishes throughout the solution
(or automatically when $\kappa=0$).
\end{proposition}
\begin{proof}
Allowed $b$ variations are projected by the orthogonal matrix $Q_u$, so
projecting \eqref{q444:rho-eom} gives the first equation. In particular
$u^*\ddot b=0$, and initial data tangent to the constraint remain there.
For the $\Phi$ equation, variations are trace-free Hermitian and
\[
 \delta_\Phi V=
 \operatorname{tr}\bigl([(m^2+\Gamma s)\Phi
              -\kappa\operatorname{adj}\Phi+\kappa bb^*]\delta\Phi\bigr).
\]
Projection onto trace-free matrices gives the second equation. The
unrestricted $\Phi$ equation is unchanged. Subtracting the two $b$
equations gives the stated reaction, proving the last equivalence.
\end{proof}

The unrestricted model retains all six real coordinates of $b$; the
constrained model retains four and supplies the displayed reaction force.
Neither is excluded by the failure of the other identification. The old
complex sector $b=\dot b=0$ \emph{is} invariant by \eqref{q444:rho-eom}.
By contrast, a nonzero fixed complex subspace of $\mathbb C^3$ invariant
under all possible Hermitian $\Phi$ must be the whole space: traceless
Hermitian matrices together with $I_3$ complex-linearly span every matrix
unit, and the matrix units applied to a nonzero vector span $\mathbb C^3$.
Thus a four-real-dimensional complex-linear transverse sector cannot retain all
these unrestricted couplings without the explicit constraint or an
additional restriction on $\Phi$.

\subsection{Preserved interacting sectors of the unrestricted equations}

The failure of the whole slice to be invariant does not exclude invariant
motions within it. We now construct such motions using the original
equations, not by adding the reaction in Proposition~\ref{q444:constrained}.
Write $a(t)=m^2+\Gamma s(t)$, and retain the trace-free Hermitian variable
$\Phi$, the standard metric on $b\in\mathbb C^3$, and $u^*u=3$.

\begin{lemma}[Global existence in the scalar mechanics]\label{q444:global}
Every initial position and velocity in \eqref{q444:potential} has a unique
solution for all real time. Its conserved energy is
\[
 E=\tfrac12\operatorname{tr}\dot\Phi^2+\dot b^*\dot b+V(\Phi,b).
\]
\end{lemma}
\begin{proof}
The unrestricted equations for both variables are
\begin{equation}\label{q444:full-eom}
 \ddot\Phi=-a\Phi+\kappa(\operatorname{adj}\Phi-bb^*)_0,
 \qquad \ddot b=-ab-\kappa\Phi b.
\end{equation}
They form a polynomial first-order system on positions and velocities,
so local existence and uniqueness hold. Put $t_\Phi=\operatorname{tr}\Phi^2$
and $n_b=b^*b$. Direct differentiation of the kinetic energy gives
\[
 -a\operatorname{tr}(\dot\Phi\Phi)
 +\kappa\operatorname{tr}(\dot\Phi\operatorname{adj}\Phi)
 -\kappa b^*\dot\Phi b
 -2a\operatorname{Re}(\dot b^*b)
 -2\kappa\operatorname{Re}(\dot b^*\Phi b).
\]
The trace projection drops out because $\operatorname{tr}\dot\Phi=0$.
Differentiating the displayed potential gives precisely the negative of
this expression, including the two terms from differentiating $b^*\Phi b$.
Thus $E$ is constant.

For completeness, a bound in the unchanged scalar coordinates suffices
to rule out escape. The three real eigenvalues of $\Phi$ give
$|\det\Phi|\leq(t_\Phi/3)^{3/2}$ by the arithmetic--geometric mean
inequality applied to their squares. Cauchy--Schwarz and the Frobenius
norm give $|b^*\Phi b|\leq n_b\sqrt{t_\Phi}$. Since
$s=t_\Phi+2n_b$, these are the explicit inequalities
\[
 |\det\Phi-b^*\Phi b|
 \leq (t_\Phi/3)^{3/2}+n_b\sqrt{t_\Phi}
 \leq \left(\frac1{3\sqrt3}+\frac12\right)s^{3/2}.
\]
This bound is not asserted to be sharp and does not change $V$.
Consequently
\[
 V\geq \frac{\Gamma}{4}s^2+\frac{m^2}{2}s
       -|\kappa|\left(\frac1{3\sqrt3}+\frac12\right)s^{3/2}
       \longrightarrow+\infty\quad(s\longrightarrow\infty).
\]
The energy identity bounds positions; the global lower bound for $V$
then bounds velocities. A polynomial vector field and its derivative are
bounded on this compact phase-space set. A finite endpoint therefore has
a limiting initial state to which local existence applies, extending the
solution. Neither a finite forward nor a finite backward endpoint occurs.
\end{proof}

\begin{theorem}[Fixed-subspace invariant sectors]\label{q444:invariant-W}
Fix a complex subspace $W\subset u^\perp$ of dimension $r=1$ or $2$, and
let $E_W$ be its orthogonal projection in the original metric. Define
\[
 \mathcal T_W=\{(\Phi,b):\Phi=\Phi^*,\ \operatorname{tr}\Phi=0,
                       \ [\Phi,E_W]=0,\ b\in W\}.
\]
With tangent initial velocities, its tangent bundle $T\mathcal T_W$ is
invariant under the full equations \eqref{q444:full-eom} for all real time.
For each fixed $W$, the configuration dimensions are respectively $6$
and $8$, and the phase-space dimensions are $12$ and $16$. Every such
solution satisfies
\[
 \rho=0,\qquad u^*\Phi b=0,\qquad \kappa(uu^*/3)\Phi b=0
\]
identically, without any constraint force.
\end{theorem}
\begin{proof}
For Hermitian $\Phi$, preservation of $W$ also preserves $W^\perp$:
$\langle\Phi x,w\rangle=\langle x,\Phi w\rangle=0$ for
$x\in W^\perp$, $w\in W$. This is equivalent to $[\Phi,E_W]=0$.
The trace-free adjugate identity is
\begin{equation}\label{q444:adj-tracefree}
 \operatorname{adj}\Phi=\Phi^2-\tfrac12(\operatorname{tr}\Phi^2)I_3.
\end{equation}
Indeed the characteristic polynomial is
$z^3-\tfrac12(\operatorname{tr}\Phi^2)z-\det\Phi$;
Cayley--Hamilton and the adjugate identity prove
\eqref{q444:adj-tracefree} when $\Phi$ is invertible, and polynomial
continuity proves it at singular $\Phi$ as well. Thus the adjugate
commutes with $E_W$. So does $bb^*$: it takes $W$ to $W$ and vanishes
on $W^\perp$. Scalar trace subtraction preserves this property.
Both accelerations in \eqref{q444:full-eom} consequently remain in the
specified linear spaces. The position derivatives do so when the
velocities satisfy $[\dot\Phi,E_W]=0$, $\dot b\in W$.
The restricted polynomial vector field has a local solution; its inclusion
is a solution of the full system. Uniqueness and Lemma~\ref{q444:global}
give the claimed global invariance.

The Hermitian blocks on $W$ and $W^\perp$ have real dimensions $r^2$
and $(3-r)^2$. The trace imposes one real equation, and $b$ contributes
$2r$ real coordinates. Hence
$\dim_{\mathbb R}\mathcal T_W=r^2+(3-r)^2-1+2r$.
Velocities double this dimension. Finally $b,\Phi b\in W\subset u^\perp$
proves each asserted zero contraction. The parameter $W$ is fixed, not
an additional moving coordinate: for $r=1$ the choices form
$\mathbb P(u^\perp)\cong\mathbb{CP}^1$, and for $r=2$ the only choice is
$W=u^\perp$.
\end{proof}

Here is an explicit interacting member that also retains a nonzero marked
cross-coordinate $z$. Keep
\[
 v=(1,-1,0)^{\mathsf T},\quad w_0=(1,1,-2)^{\mathsf T},\quad
 E=vv^*/2,\quad F=w_0w_0^*/6,
\]
and take
\begin{equation}\label{q444:nonzero-sector}
 W=\mathbb Cv,\quad \Phi_0=E-F+w_0u^*+uw_0^*,\quad b_0=v,
 \quad \dot\Phi_0=0,\quad\dot b_0=0.
\end{equation}
The vectors $u,v,w_0$ are orthogonal with squared lengths $3,2,6$.
Thus $\operatorname{tr}\Phi_0=0$, $\Phi_0v=v$,
$\operatorname{tr}\Phi_0^2=38$, $s_0=42$, and
\[
 b_0^*\Phi_0b_0=2,\qquad
 \ddot b_0=-(m^2+42\Gamma+\kappa)v.
\]
For example $\kappa>0$ gives nonstationary interacting motion. In the
original marked coordinates, not a newly chosen basis,
\[
 \chi_0=R^*b_0=(1,0)^{\mathsf T},\qquad
 z_0=P\Phi_0u/3=(0,3)^{\mathsf T},\qquad\lambda_0=0.
\]
Hence $z_0\ne0$ while $z_0^*\chi_0=0$. Preservation of the scalar slice
does not require $z=0$ for every solution. The theorem constructs exact
families; it does not assert they exhaust all invariant initial data.

\begin{proposition}[Tangency through third order is insufficient]\label{q444:fourth-jet}
For any unrestricted solution put $c=u^*\Phi b$ and
$h=u^*\dot\Phi\dot b$. The exact identities are
\begin{align}\label{q444:jets}
 \ddot c&=-2ac-\frac{\kappa s}{3}\rho+2h,\notag\\
 \rho^{(4)}&=\left(-\ddot a+a^2+\frac{\kappa^2s}{3}\right)\rho
             -2\dot a\dot\rho+3a\kappa c-2\kappa h.
\end{align}
There are initial data with $\rho^{(j)}(0)=0$ for $0\leq j\leq3$
but $\rho^{(4)}(0)=-12\kappa$.
\end{proposition}
\begin{proof}
Equation~\eqref{q444:adj-tracefree} gives, with
$t_\Phi=\operatorname{tr}\Phi^2$ and $n_b=b^*b$,
\[
 (\operatorname{adj}\Phi-bb^*)_0
   =\Phi^2-bb^*+\frac{n_b-t_\Phi}{3}I_3,
 \qquad
 (\operatorname{adj}\Phi-bb^*)_0b=\Phi^2b-\frac{s}{3}b.
\]
Differentiate $c$ twice and substitute \eqref{q444:full-eom}:
\[
 \ddot c=u^*(\ddot\Phi b+2\dot\Phi\dot b+\Phi\ddot b)
         =-2ac-\frac{\kappa s}{3}\rho+2h.
\]
The two terms $\kappa u^*\Phi^2b$ cancel exactly. Twice differentiating
$\ddot\rho=-a\rho-\kappa c$ and substituting this identity proves the
second formula. In particular, at $\rho=\dot\rho=0$, vanishing second
and third derivatives requires $c=\dot c=0$ when $\kappa\ne0$;
then the fourth derivative is $-2\kappa h$.

For the promised initial data use the same $u,v,w_0$ as above, but now set
\[
 \Phi_0=0,\quad b_0=v,\quad\dot b_0=w_0,\quad
 \dot\Phi_0=(uw_0^*+w_0u^*)/3.
\]
This matrix is trace-free Hermitian, $\dot\Phi_0v=0$ and
$\dot\Phi_0w_0=2u$. Consequently
$\rho_0=\dot\rho_0=c_0=\dot c_0=0$, whereas $h_0=6$.
Formula~\eqref{q444:jets} gives $\rho^{(4)}(0)=-12\kappa$, and Taylor's
theorem gives $\rho(t)=-(\kappa/2)t^4+o(t^4)$.
For $\kappa\ne0$ this solution leaves the slice despite the lower-order
tangencies. For $\kappa=0$, by contrast, the full locus
$\rho=\dot\rho=0$ is invariant because its equation is the homogeneous
linear equation $\ddot\rho=-a(t)\rho$ with zero initial data.
\end{proof}

\subsection{Periodic motions outside every fixed-subspace family}

The families in Theorem~\ref{q444:invariant-W} do not exhaust the motions
inside $\rho=0$. The following exact relative equilibria rotate the
instantaneous kernel of $\Phi$ while retaining all the original couplings.

\begin{theorem}[A moving-kernel family]\label{q444:moving-family}
Choose $\Gamma>0$, $\kappa\ne0$, and real $x,a$ satisfying
\[
 0<\frac{x}{\kappa}\leq\frac1{6\Gamma},\qquad
 a>\frac{\kappa x}{3}.
\]
Choose $B,d>0$ with
\begin{equation}\label{q444:moving-parameters}
 B^2+d^2=\frac{3ax}{\kappa}-x^2,
 \qquad m^2=a\left(1-\frac{6\Gamma x}{\kappa}\right).
\end{equation}
The right-hand sides give a positive sum and $m^2\geq0$.
Retain $u,v,w_0$ from \eqref{q444:nonzero-sector}, and put
\[
 J_0=\begin{pmatrix}0&-1&1\\1&0&-1\\-1&1&0\end{pmatrix},
 \quad \Pi=uu^*/3,\quad Q_u=I_3-\Pi,\quad
 \Omega=\sqrt{a/3}\,J_0,
\]
\begin{equation}\label{q444:rotation}
 O(t)=\Pi+\cos(\sqrt a\,t)Q_u
                 +\frac{\sin(\sqrt a\,t)}{\sqrt3}J_0.
\end{equation}
Define the original-coordinate initial data and their motion by
\begin{align}\label{q444:moving-data}
 b_0&=\frac B{\sqrt2}v,\qquad
 \Phi_0=-\frac x6w_0w_0^*+\frac x3uu^*
          +\frac d{\sqrt{18}}(w_0u^*+uw_0^*),\notag\\
 b(t)&=O(t)b_0,\qquad \Phi(t)=O(t)\Phi_0O(t)^*.
\end{align}
These are global periodic solutions of \eqref{q444:full-eom}. They satisfy
\[
 s=\frac{6ax}{\kappa},\qquad m^2+\Gamma s=a,\qquad
 \Phi(t)b(t)=0,\qquad\rho(t)=0.
\]
No fixed complex $W\subset u^\perp$ contains this trajectory in
$T\mathcal T_W$.
\end{theorem}
\begin{proof}
Direct multiplication gives $J_0^*=-J_0$, $J_0u=0$,
$J_0^2=-3Q_u$, and $\Pi J_0=J_0\Pi=0$.
Consequently $O(t)^*O(t)=I_3$, $O(t)u=u$, and
$\dot O=\Omega O=O\Omega$. Its determinant is one by continuity from
$O(0)=I_3$. This proves that the displayed rotation is in the original
Euclidean metric, with its frequency $\sqrt a$ unchanged.

For the following verification only, use the explicit orthogonal matrix
\[
 K=\left(\frac v{\sqrt2},\frac{w_0}{\sqrt6},\frac u{\sqrt3}\right),
 \quad K^*K=I_3,
\]
and retain its inverse $K^*$ and all its factors. It gives
\[
 K^*\Phi_0K=
 H=\begin{pmatrix}0&0&0\\0&-x&d\\0&d&x\end{pmatrix},
 \quad K^*b_0=Be_1,\quad
 K^*\Omega K=A=\begin{pmatrix}0&-\sqrt a&0\\\sqrt a&0&0\\0&0&0\end{pmatrix}.
\]
Thus $s=2(x^2+d^2+B^2)=6ax/\kappa$, as asserted.
In these exact coordinates $HBe_1=0$ and
\[
 (\operatorname{adj}H-B^2e_1e_1^*)_0
    =\operatorname{diag}(-s/3,s/6,s/6).
\]
The required accelerations are explicitly
\[
 [A,[A,H]]=
 \begin{pmatrix}-2ax&0&0\\0&2ax&-ad\\0&-ad&0\end{pmatrix}
 =-aH+\kappa\operatorname{diag}(-s/3,s/6,s/6),
 \qquad A^2Be_1=-aBe_1.
\]
Conjugation by $K$ therefore verifies both equations at $t=0$.
The full force commutes with orthogonal conjugation on $\Phi$ and the
same orthogonal action on $b$: the adjugate, trace subtraction, norms,
and product $\Phi b$ all transform in this way. Since $O$ commutes with
$\Omega$, conjugation by $O(t)$ verifies both equations at every time,
not merely to a finite Taylor order. This also proves the stated kernel
and normal-coordinate identities.

At times $0$ and $\pi/(2\sqrt a)$, the two values of $b$ are nonzero
multiples of $v$ and $w_0$. They span $u^\perp$ over $\mathbb C$.
Any fixed $W$ containing the whole trajectory would consequently equal
$u^\perp$. But
$u^*\Phi_0w_0=\sqrt{18}\,d\ne0$, so $\Phi_0$ does not preserve this
plane. This rules out every fixed-$W$ family in the earlier theorem.
\end{proof}

A particularly small exact instance is
\begin{equation}\label{q444:periodic-instance}
 \begin{gathered}
 \Gamma=m^2=1,\quad\kappa=6,\quad x=\tfrac12,\quad a=2,
 \quad B=d=\frac1{2\sqrt2},\\
 b_0=\frac14\begin{pmatrix}1\\-1\\0\end{pmatrix},\quad
 \Phi_0=\frac14\begin{pmatrix}1&1&1\\1&1&1\\1&1&-2\end{pmatrix}.
 \end{gathered}
\end{equation}
Here $s=1$ and the period of $b$ is $2\pi/\sqrt2$. Its original marked
coordinates at zero are $\chi_0=(1/4,0)^{\mathsf T}$,
$z_0=(0,1/4)^{\mathsf T}$, and $\lambda_0=1/6$.
Although $\det\Phi-b^*\Phi b=0$ along this solution, its cubic force
does not vanish: at zero it is
$6(\operatorname{adj}\Phi_0-b_0b_0^*)_0\ne0$.
It is the coupled equations, not removal of the cubic, that sustain the
rotating-kernel solution.

\begin{proposition}[Conserved charge and the precise gauge comparison]\label{q444:charge}
In the scalar mechanics the anti-Hermitian matrix
\begin{equation}\label{q444:moment}
 \mathcal J=[\Phi,\dot\Phi]+b\dot b^*-\dot b b^*
\end{equation}
is constant. Its trace-free part is the $SU(3)$ moment matrix, using the
pairing $\operatorname{tr}(\mathcal J\xi)$ with
$\xi\in\mathfrak{su}(3)$. On \eqref{q444:moving-data},
\[
 \mathcal J=-\frac s2\Omega,
 \qquad T=\frac{as}{2},\qquad
 E=m^2s+\frac{3\Gamma}{4}s^2.
\]
In particular \eqref{q444:periodic-instance} has
$\mathcal J=-\Omega/2\ne0$, $T=1$, $V=3/4$, and $E=7/4$.
This scalar solution is not a source-free homogeneous gauge solution
with all $\mathcal A_\mu=0$: it fails the $\mathcal A_0$ field equation.
\end{proposition}
\begin{proof}
Using \eqref{q444:full-eom}, polynomial commutation of $\Phi$ with its
adjugate, and commutation with the scalar trace correction, gives
\[
 \frac{d}{dt}[\Phi,\dot\Phi]=-\kappa[\Phi,bb^*],\qquad
 \frac{d}{dt}(b\dot b^*-\dot b b^*)=\kappa[\Phi,bb^*].
\]
Their sum vanishes. Under the infinitesimal simultaneous symmetry
$\delta\Phi=[\xi,\Phi]$, $\delta b=\xi b$, the kinetic momentum paired
with this variation is
\[
 \operatorname{tr}(\dot\Phi[\xi,\Phi])
       +2\operatorname{Re}(\dot b^*\xi b)
       =\operatorname{tr}(\mathcal J\xi).
\]
This proves the pairing, including its coefficient and sign.
In the explicit frame $K$, compute
\[
 [H,[A,H]]+Be_1(ABe_1)^*-(ABe_1)(Be_1)^*
       =-(x^2+d^2+B^2)A=-\frac s2 A.
\]
It commutes with $A$, so conjugation by $O(t)$ retains the same constant
matrix. Similarly
$\tfrac12\operatorname{tr}[A,H]^2+|ABe_1|^2=as/2$.
The determinant and $b^*\Phi b$ both vanish on the trajectory; substitution
of $a=m^2+\Gamma s$ gives the displayed energy.

For the gauge comparison retain the packet's covariant derivatives
$D_0\Phi=\dot\Phi+[\mathcal A_0,\Phi]$,
$D_0b=\dot b+\mathcal A_0b$. Varying $\mathcal A_0$ by
$\xi\in\mathfrak{su}(3)$ in the scalar kinetic density gives
$\operatorname{tr}(\mathcal J\xi)$ at $\mathcal A_0=0$.
With homogeneous fields and every $\mathcal A_i=0$, the electric field
vanishes and the gauge kinetic variation contributes zero. Thus its
Euler--Lagrange equation requires the trace-free part of $\mathcal J$ to
vanish. The exhibited matrix is already trace-free and nonzero, so this
particular zero-connection lift fails. This is the Gauss equation obtained
by variation before imposing temporal gauge; compare D.~Tong,
\emph{Lectures on Gauge Theory}, Section~2.2.1, equation~(2.27),
\url{https://www.damtp.cam.ac.uk/user/tong/gaugetheory/gt.pdf}.
Our matter contribution has been derived here in the retained
anti-Hermitian convention, not imported from a different convention.
\end{proof}

There is an exact moving-frame map, not an assertion that these theories
are unrelated. With $U(t)=O(t)^*$ the gauge transformation rule is
\[
 \Phi'=U\Phi U^{-1}=\Phi_0,\quad b'=Ub=b_0,\quad
 \mathcal A'_0=U\mathcal A_0U^{-1}-\dot U U^{-1}=\Omega.
\]
The covariant velocities remain
$D'_0\Phi_0=[\Omega,\Phi_0]$ and $D'_0b_0=\Omega b_0$.
The same nonzero charge and Gauss residual therefore remain; making the
ordinary coordinates static does not remove them. A gauge-field lift
would have to supply the electric contribution to Gauss' equation and
satisfy the remaining coupled equations. Neither its existence nor its
nonexistence follows from failure of the zero-connection lift, and no
quantum spectral-gap conclusion is inferred from this periodic motion.

\subsection{The complete scalar vacuum calculation with its kinetic metric}

This subsection validates the homogeneous scalar vacuum theorem in the
triality packet. It does not assert a new discovery of that phase diagram,
and does not include the packet's subsequent gauge-orbit claims.
Set $Q=\Theta(\Phi,b)$ with $\operatorname{tr}Q=0$, and retain
\[
 g(Y,Z)=\operatorname{Re}\operatorname{tr}(YZ),\qquad
 s(Q)=g(Q,Q)=\operatorname{tr}\Phi^2+2b^*b,\qquad
 n(Q)=N(Q)=\det\Phi-b^*\Phi b.
\]
The kinetic term is $g(\dot Q,\dot Q)/2$, not the unweighted sum of
squares of the independent entries. Explicitly,
$g(Y,Y)=\sum_iY_{ii}^2+2\sum_{i<j}|Y_{ij}|^2$ for Hermitian $Y$.

\begin{lemma}[Sharp bound and all equality cases]\label{q444:sharp}
For every trace-free quaternionic Hermitian $Q$,
\[
 54n(Q)^2\leq s(Q)^3.
\]
For $Q\ne0$, equality holds exactly at
$Q=x(3P-I_3)$, $x\in\mathbb R\setminus\{0\}$, where
$P=P^*=P^2$ has quaternionic rank one. There
$s=6x^2$ and $n=2x^3$.
\end{lemma}
\begin{proof}
Here are the spectral and determinant details in the stated conventions.
The real-valued function $v^*Qv$ on the unit sphere of $\mathbb H^3$
attains a maximum. Its Lagrange-multiplier equation gives
$\operatorname{Re}(w^*(Qv-v\lambda))=0$ for every $w$, with real
$\lambda=v^*Qv$. Taking $w=Qv-v\lambda$ proves $Qv=v\lambda$.
The quaternionic orthogonal complement of $v$ is invariant because
$v^*Qw=\lambda v^*w$. Iteration gives a unitary eigenbasis with real
eigenvalues $\lambda_1,\lambda_2,\lambda_3$.

For quaternionic matrices, the real trace is cyclic by
$\operatorname{Re}(pq)=\operatorname{Re}(qp)$, entry by entry. Moreover,
the original cubic \eqref{q444:moore} obeys
\begin{equation}\label{q444:newton-cubic}
 6n(Q)=(\operatorname{tr}Q)^3
       -3\operatorname{tr}(Q)\operatorname{tr}(Q^2)
       +2\operatorname{tr}(Q^3).
\end{equation}
To check it without changing the quaternionic order, in the ordered
entries of \eqref{q444:moore} the cubic trace is the sum of the diagonal
cubes, plus $3(d+f)|q_1|^2+3(a+f)|q_2|^2+3(a+d)|q_3|^2$,
plus $6\operatorname{Re}(q_1q_2q_3)$. Substitution gives each term of
six times \eqref{q444:moore}. Unitary conjugation therefore preserves
$s,n$, which are respectively $\sum_i\lambda_i^2$ and
$\lambda_1\lambda_2\lambda_3$ in the eigenbasis. Their sum is zero.
The exact polynomial identity
\[
 s^3-54n^2=2(\lambda_1-\lambda_2)^2
                (\lambda_1-\lambda_3)^2(\lambda_2-\lambda_3)^2
\]
follows by substituting $\lambda_3=-\lambda_1-\lambda_2$ and expanding.
It proves the bound. For nonzero $Q$, equality means a repeated
eigenvalue, hence spectrum $(2x,-x,-x)$ for some $x\ne0$.
Letting $P$ project onto the simple eigenline gives $Q=x(3P-I_3)$.
The converse follows from this same spectrum.
\end{proof}

\begin{theorem}[All scalar global minima]\label{q444:vacua}
For $\Gamma>0$, $m^2\geq0$, and $\kappa\ne0$, put
$m_{\mathrm c}^2=\kappa^2/(27\Gamma)$.
The global minima of $V=m^2s/2+\Gamma s^2/4-\kappa n$ are exactly:
\begin{itemize}
\item $Q=0$ if $m^2>m_{\mathrm c}^2$;
\item $Q=0$ and $Q=\frac{\kappa}{9\Gamma}(3P-I_3)$ for every
rank-one quaternionic Hermitian projector $P$, if $m^2=m_{\mathrm c}^2$;
\item $Q=x_*(3P-I_3)$ for every such $P$, if $0\leq m^2<m_{\mathrm c}^2$,
where
\[
 x_* =\frac{\kappa+\operatorname{sgn}(\kappa)
              \sqrt{\kappa^2-24\Gamma m^2}}{12\Gamma}.
\]
\end{itemize}
Each nonzero component is diffeomorphic to $\mathbb{HP}^2$.
For $\kappa=0$, the origin is the unique global minimum.
\end{theorem}
\begin{proof}
For $\kappa=0$, the positive quartic proves the assertion, including
$m^2=0$. Otherwise at each fixed original norm $r=\sqrt s$ the sharp
bound is attained with
$x=\operatorname{sgn}(\kappa)r/\sqrt6$ and $Q=x(3P-I_3)$.
For $r>0$, these are exactly its equality and sign cases. Thus the full
minimization is equivalent, with all fibres of the norm map accounted for,
to minimizing
\[
 f(r)=\frac{m^2}{2}r^2+\frac\Gamma4r^4
                   -\frac{|\kappa|}{3\sqrt6}r^3,\qquad r\geq0.
\]
Set $r_{\mathrm c}=2|\kappa|/(3\sqrt6\Gamma)$. Exact factorization gives
\[
 f(r)=\frac\Gamma4r^2(r-r_{\mathrm c})^2
                    +\frac{m^2-m_{\mathrm c}^2}{2}r^2.
\]
This proves the first two cases, including every zero. Below the transition,
$f(r_{\mathrm c})<0$, so a global minimum is nonzero. Put
$a_0=|\kappa|/\sqrt6$. The derivative is
$f'(r)=r(\Gamma r^2-a_0r+m^2)$, with roots
\[
 r_\pm=\frac{a_0\pm\sqrt{a_0^2-4\Gamma m^2}}{2\Gamma}.
\]
Below the transition its discriminant is strictly positive. If $m^2>0$,
$f'$ has signs $+,-,+$ across these two positive roots, so $r_+$ is the
only positive local minimum. If $m^2=0$, $r_-=0$, and the signs on
$(0,r_+)$ and $(r_+,\infty)$ are $-,+$. Coercivity and the negative
value exclude the origin in both cases. Converting $r_+$ to $x$ proves
the third case and the completeness of its matrix equality locus.

For the manifold statement, a right quaternionic line represented by
$v\ne0$ has orthogonal projector $P_v=v(v^*v)^{-1}v^*$.
This is unchanged under $v\mapsto vh$, $h\in\mathbb H^\times$,
and every rank-one Hermitian projector is the unique projection onto
its image line. The formula and its inverse line charts are smooth.
For a fixed nonzero $x$, the maps $P\mapsto x(3P-I_3)$ and
$Q\mapsto(Q/x+I_3)/3$ are mutually inverse. This proves the claimed
$\mathbb{HP}^2$ identification without any gauge quotient.
\end{proof}

\begin{theorem}[Full metric Hessian]\label{q444:vacuum-hessian}
At every nonzero minimum $Q_*=x(3P-I_3)$, the Hessian relative to $g$
has eigenvalues
\[
 \begin{array}{c|c}
 \text{eigenvalue}&\text{real multiplicity}\\\hline
 x(12\Gamma x-\kappa)&1\\
 0&8\\
 3\kappa x&5
 \end{array}
\]
Both nonzero eigenvalues are positive. The entire kernel is the tangent
space of the nonzero minimum manifold. At the origin the metric Hessian
is $m^2I$ in all fourteen directions.
\end{theorem}
\begin{proof}
The preceding spectral proof and cubic trace identity show that
quaternionic unitary conjugation preserves both $g$ and $V$.
Transport the Hessian isometrically to $Q_*=xD$, $D=\operatorname{diag}(2,-1,-1)$.
Every trace-free Hermitian variation has the unique, $g$-orthogonal decomposition
\[
 Y=yD+\begin{pmatrix}0&v^*\\v&0_2\end{pmatrix}
       +\begin{pmatrix}0&0\\0&B_0\end{pmatrix},
 \qquad y\in\mathbb R,\quad v\in\mathbb H^2,
 \quad B_0=B_0^*,\quad\operatorname{tr}B_0=0.
\]
Write $B_0=\left(\begin{smallmatrix}\beta&\omega\\\bar\omega&-\beta\end{smallmatrix}\right)$,
$\beta\in\mathbb R$, $\omega\in\mathbb H$. The respective dimensions are
$1,8,5$, and the retained norm is
\[
 g(Y,Y)=6y^2+2v^*v+\|B_0\|^2,
 \qquad\|B_0\|^2=2\beta^2+2|\omega|^2.
\]
For $h=x+\varepsilon y$, the full cubic, including every mixed direction,
is
\begin{equation}\label{q444:hessian-cubic}
 n(Q_*+\varepsilon Y)
 =2h^3+\varepsilon^2h(v^*v-\|B_0\|^2)
                                      +\varepsilon^3v^*B_0v.
\end{equation}
Indeed the diagonal and lower off-diagonal contributions are
$2h(h^2-\varepsilon^2\beta^2)-2h\varepsilon^2|\omega|^2$.
The upper entries contribute
$h\varepsilon^2(|v_1|^2+|v_2|^2)
+\varepsilon^3\beta(|v_1|^2-|v_2|^2)$, and the triangular term is
$2\varepsilon^3\operatorname{Re}(\bar v_1\omega v_2)$.
These sum to \eqref{q444:hessian-cubic}, preserving the quaternionic order.
Also
\[
 s(Q_*+\varepsilon Y)=6x^2+12\varepsilon xy
            +\varepsilon^2(6y^2+2v^*v+\|B_0\|^2).
\]
The second derivative of the cubic is therefore
$12xy^2+2xv^*v-2x\|B_0\|^2$. Differentiating the potential gives the
complete quadratic form before using stationarity:
\begin{align*}
 D^2V_{xD}[Y,Y]
 ={}&6(m^2+18\Gamma x^2-2\kappa x)y^2\\
 &+2(m^2+6\Gamma x^2-\kappa x)v^*v\\
 &+(m^2+6\Gamma x^2+2\kappa x)\|B_0\|^2.
\end{align*}
Every nonzero minimum satisfies $m^2-\kappa x+6\Gamma x^2=0$.
Thus this becomes
\[
 6x(12\Gamma x-\kappa)y^2+3\kappa x\|B_0\|^2.
\]
There are no mixed quadratic terms; polarization therefore gives no mixed
bilinear blocks. Comparison with $g(Y,Y)$ proves every listed eigenvalue
and multiplicity, with no rescaling of the kinetic metric.

At coexistence the positive values are
$\kappa^2/(27\Gamma)$ and $\kappa^2/(3\Gamma)$.
Below it, $12\Gamma x-\kappa=
\operatorname{sgn}(\kappa)\sqrt{\kappa^2-24\Gamma m^2}$ and $x$ has
the same sign as $\kappa$, proving positivity. In the line chart
\[
 P(z)=\frac1{1+z^*z}\begin{pmatrix}1&z^*\\z&zz^*\end{pmatrix},
 \qquad DP_0[\zeta]=\begin{pmatrix}0&\zeta^*\\\zeta&0_2\end{pmatrix},
\]
the derivative of $x(3P(z)-I_3)$ spans exactly the eight-dimensional
zero block. Isometric conjugation transports this to every $P$.
At the origin, the cubic and quartic contribute zero Hessian, giving
$m^2g$ directly. In particular, at $m^2=\kappa=0$ the origin is still the
unique minimum, although its Hessian vanishes in all directions.
\end{proof}

These are scalar Hessian values of a specified finite-dimensional model.
Their exact multiplicities do not assert a gauge-reduced physical spectrum
or an interacting quantum mass gap. The current calculation includes the
full tangent and normal blocks, not solely the radial restriction.

\subsection{Gauge kinetic sign and precise scope of this audit}

In the packet's gauge lift, $\mathcal A_\mu\in\mathfrak{su}(3)$ and
$D_\mu b=\partial_\mu b+\mathcal A_\mu b$ fix anti-Hermitian gauge
generators. Retain Minkowski signature $(+---)$, consistent with the
positive time-derivative scalar term in that lift. Write
$\|T\|^2=-\operatorname{tr}(T^2)$ for anti-Hermitian $T$, and
$E_i=\mathcal F_{0i}$, $B_i=\tfrac12\epsilon_{ijk}\mathcal F_{jk}$.
Raising the time-space index pair changes its sign and summing both
orders of each antisymmetric pair gives
\[
 \operatorname{tr}(\mathcal F_{\mu\nu}\mathcal F^{\mu\nu})
 =2\sum_i(\|E_i\|^2-\|B_i\|^2).
\]
Consequently the packet's coefficient $-1/(2g_c^2)$ gives negative electric
kinetic energy in these conventions. Keeping its anti-Hermitian fields
and scalar kinetic terms, the positive-energy gauge density is
\[
 \mathcal L_{\mathrm{gauge}}
 =+\frac1{2g_c^2}\operatorname{tr}
       (\mathcal F_{\mu\nu}\mathcal F^{\mu\nu})
 =\frac1{g_c^2}\sum_i(\|E_i\|^2-\|B_i\|^2).
\]
For an immediate check take $\mathcal F_{01}=\operatorname{diag}(i,-i,0)$
and all other independent components zero: the trace contraction is $4$,
so the two signs give $-2/g_c^2$ and $+2/g_c^2$, respectively. This sign
correction changes neither the scalar potential nor its derivative
calculated above. It does matter before interpreting the full gauge lift
as a positive-energy physical model.

The equal-diagonal thirteen-space in the triality packet remains a
different embedded slice: its intersection with
$\operatorname{Herm}_3(\mathbb C)$ has dimension seven, whereas the
common Lorentz slice contains all nine complex-Hermitian coordinates.
Their precise common ambient space is $\operatorname{Herm}_3(\mathbb H)$;
Theorem~\ref{q444:identification} identifies the \emph{Lorentz} constructions,
not the equal-diagonal one. No result here identifies these cone-boundary
spheres with the claimed integrable complex structure on $S^6$, assigns
new spinor periods and toroidal fillings, or proves a continuum interacting
Yang--Mills spectral gap. The exact next data are the transition maps and
lattice extension for the additional spinor, followed by a physical state
space and Hamiltonian comparison retaining the couplings above.

% Ready-to-input fragment; no canonical files are modified by this artifact.
% Uses the existing q444:potential and q444:moving-family data and standard
% theorem/proposition/lemma environments. All new calculations are proved here.
\subsection{The full homogeneous gauge equations and the literal periodic orbit}

Retain signature $(+---)$, $\mathcal A_\mu\in\mathfrak{su}(3)$,
$\Phi=\Phi^*$, $\operatorname{tr}\Phi=0$, and $b\in\mathbb C^3$.
Write $g^2>0$ for the gauge coupling in this subsection and keep
\[
 s=\operatorname{tr}\Phi^2+2b^*b,\qquad
 a=m^2+\Gamma s,\qquad
 V=\frac{m^2}{2}s+\frac{\Gamma}{4}s^2-\kappa(\det\Phi-b^*\Phi b).
\]
Thus $m^2\geq0$, $\Gamma>0$, and $\kappa\in\mathbb R$ are unchanged.
For a matrix $Z$ write
$Z_{\mathrm{tf}}=Z-\operatorname{tr}(Z)I_3/3$. The full density is
\begin{equation}\label{q444:full-gauge-density}
 \mathcal L=\frac1{2g^2}\operatorname{tr}
 (\mathcal F_{\mu\nu}\mathcal F^{\mu\nu})
 +\frac12\operatorname{tr}(D_\mu\Phi D^\mu\Phi)
 +(D_\mu b)^*D^\mu b-V,
\end{equation}
where $\mathcal F_{\mu\nu}=\partial_\mu\mathcal A_\nu-
\partial_\nu\mathcal A_\mu+[\mathcal A_\mu,\mathcal A_\nu]$,
$D_\mu\Phi=\partial_\mu\Phi+[\mathcal A_\mu,\Phi]$, and
$D_\mu b=\partial_\mu b+\mathcal A_\mu b$.
We consider spatially homogeneous fields on a finite-volume flat
$\mathbb T^3$, or locally on $\mathbb R^3$. A nonzero constant energy
density is not asserted to have finite total energy on $\mathbb R^3$.

\begin{proposition}[Every homogeneous field equation]\label{q444:gauge-eom}
Set
\[
 E_i=\dot{\mathcal A}_i+[\mathcal A_0,\mathcal A_i],\qquad
 F_{ij}=[\mathcal A_i,\mathcal A_j],\qquad
 X=\dot\Phi+[\mathcal A_0,\Phi],\qquad v=\dot b+\mathcal A_0b,
\]
and define the lower-index matter currents
\[
 J_0=[\Phi,X]+bv^*-vb^*,\qquad
 J_i=[\Phi,[\mathcal A_i,\Phi]]
                -bb^*\mathcal A_i-\mathcal A_i bb^*.
\]
The complete homogeneous Euler--Lagrange system is
\begin{align}
 D_0X-\sum_i[\mathcal A_i,[\mathcal A_i,\Phi]]
       +a\Phi-\kappa(\operatorname{adj}\Phi-bb^*)_{\mathrm{tf}}&=0,
       \label{q444:gauge-phi}\\
 D_0v-\sum_i\mathcal A_i^2b+(aI_3+\kappa\Phi)b&=0,
       \label{q444:gauge-b}\\
 \frac2{g^2}\sum_i[\mathcal A_i,E_i]-(J_0)_{\mathrm{tf}}&=0,
       \label{q444:full-gauss}\\
 D_0E_i-\sum_j[\mathcal A_j,[\mathcal A_j,\mathcal A_i]]
       -\frac{g^2}{2}(J_i)_{\mathrm{tf}}&=0
       \quad(i=1,2,3). \label{q444:gauge-spatial}
\end{align}
Here $D_0X=\dot X+[\mathcal A_0,X]$,
$D_0v=\dot v+\mathcal A_0v$, and similarly for $E_i$.
\end{proposition}
\begin{proof}
First retain the homogeneous density without deleting $\mathcal A_0$:
\[
 \mathcal L_{\mathrm{hom}}
 =-\frac1{g^2}\sum_i\operatorname{tr}E_i^2
   +\frac1{g^2}\sum_{i<j}\operatorname{tr}F_{ij}^2
   +\frac12\operatorname{tr}X^2+v^*v
   -\frac12\sum_i\operatorname{tr}[\mathcal A_i,\Phi]^2
   -\sum_i|\mathcal A_i b|^2-V.
\]
Both signs in the gauge terms follow by raising indices in $(+---)$
and counting the two orders of each antisymmetric pair.
For trace-free Hermitian $\delta\Phi$ and arbitrary $\delta b$,
\[
 \delta V=
 \operatorname{tr}\!\left[
 \bigl(a\Phi-\kappa(\operatorname{adj}\Phi-bb^*)_{\mathrm{tf}}\bigr)
 \delta\Phi\right]
 +2\operatorname{Re}\bigl(\delta b^*(aI_3+\kappa\Phi)b\bigr).
\]
Covariant integration by parts therefore gives
$D_\mu D^\mu\Phi+a\Phi-
\kappa(\operatorname{adj}\Phi-bb^*)_{\mathrm{tf}}=0$ and
$D_\mu D^\mu b+(aI_3+\kappa\Phi)b=0$. Homogeneity gives the first
two displayed equations, including the minus signs in their spatial sums.

For the gauge variation, antisymmetry gives
\[
 \delta\mathcal L_{\mathrm{gauge}}
 =\frac2{g^2}\operatorname{tr}
       ((D_\mu\delta\mathcal A_\nu)\mathcal F^{\mu\nu}).
\]
After integration by parts its coefficient is
$-2D_\mu\mathcal F^{\mu\nu}/g^2$.
The scalar coefficient of $\delta\mathcal A_\nu$ is
\[
 J^\nu=[\Phi,D^\nu\Phi]
           +b(D^\nu b)^*-(D^\nu b)b^*.
\]
Indeed cyclicity and anti-Hermiticity respectively give
\begin{align*}
 \operatorname{tr}(D^\nu\Phi[\delta\mathcal A_\nu,\Phi])
 &=\operatorname{tr}([\Phi,D^\nu\Phi]\delta\mathcal A_\nu),\\
 2\operatorname{Re}((D^\nu b)^*\delta\mathcal A_\nu b)
 &=\operatorname{tr}\bigl((b(D^\nu b)^*-(D^\nu b)b^*)
                        \delta\mathcal A_\nu\bigr).
\end{align*}
Variation over $\mathfrak{su}(3)$ consequently yields
\[
 D_\mu\mathcal F^{\mu\nu}=\frac{g^2}{2}(J^\nu)_{\mathrm{tf}}.
\]
For $\nu=0$, $\mathcal F^{i0}=E_i$ and $J^0=J_0$.
For $\nu=i$, $\mathcal F^{0i}=-E_i$,
$\mathcal F^{ji}=[\mathcal A_j,\mathcal A_i]$, and $J^i=-J_i$.
These substitutions give exactly \eqref{q444:full-gauss} and
\eqref{q444:gauge-spatial}. In particular the $\mathcal A_0$ equation
survives after choosing temporal gauge. This is the same variational
principle emphasized by D.~Tong, \emph{Lectures on Gauge Theory},
Section~2.2.1, equation~(2.27); the matter currents, signs, and coefficient
$g^2/2$ above have been derived in the present conventions.
\end{proof}

\begin{theorem}[No temporal-gauge lift retaining the literal orbit]
\label{q444:temporal-no-lift}
Keep exactly the moving-kernel scalar functions in
Theorem~\ref{q444:moving-family}, with $B,d>0$ and $x\neq0$.
There are no homogeneous spatial fields $\mathcal A_i(t)$ for which
these same scalar functions, together with $\mathcal A_0=0$, solve
the full system \eqref{q444:gauge-phi}--\eqref{q444:gauge-spatial}.
\end{theorem}
\begin{proof}
The scalar functions already solve the unrestricted scalar equations.
Subtracting them from \eqref{q444:gauge-phi}--\eqref{q444:gauge-b},
with $\mathcal A_0=0$, gives
\[
 \sum_i[\mathcal A_i,[\mathcal A_i,\Phi]]=0,\qquad
 \sum_i\mathcal A_i^2b=0.
\]
Taking the trace pairing with $\Phi$ and the Hermitian pairing with $b$
gives respectively
\[
 -\sum_i\operatorname{tr}[\mathcal A_i,\Phi]^2=0,\qquad
 -\sum_i|\mathcal A_i b|^2=0.
\]
Each $[\mathcal A_i,\Phi]$ is Hermitian, so both are sums of
nonnegative squared norms with the displayed common minus sign. Thus
$[\mathcal A_i,\Phi]=0$ and $\mathcal A_i b=0$ separately for every $i$.

To classify this stabilizer use the retained orthogonal matrix
$K=(v/\sqrt2,w_0/\sqrt6,u/\sqrt3)$ from the moving-family proof, and
put $r^2=x^2+d^2>0$. At time zero the data are
\[
 H=K^*\Phi_0K=\begin{pmatrix}0&0&0\\0&-x&d\\0&d&x\end{pmatrix},
 \qquad K^*b_0=Be_1,\qquad
 A=K^*\Omega K=
 \begin{pmatrix}0&-\sqrt a&0\\\sqrt a&0&0\\0&0&0\end{pmatrix}.
\]
An anti-Hermitian traceless matrix killing $Be_1$ has zero first row
and column and a lower $\mathfrak{su}(2)$ block. Commutation with
$\left(\begin{smallmatrix}-x&d\\d&x\end{smallmatrix}\right)$, whose
distinct eigenvalues are $\pm r$, forces that block to be its real
multiple times $i$. Therefore the stabilizer Lie algebra is
$\mathbb R\,i\Phi(t)$. The full stabilizer group is the circle
\[
 \left\{P_b+\cos\theta(I_3-P_b)
          +i\frac{\sin\theta}{r}\Phi(t):\theta\in\mathbb R\right\},
 \qquad P_b=\frac{b(t)b(t)^*}{B^2}.
\]
Indeed it fixes $b$, has eigenvalues $1,e^{i\theta},e^{-i\theta}$,
and exhausts the unitary commutant fixing the zero eigenspace pointwise.
Consequently
\[
 \mathcal A_i=i c_i(t)\Phi(t),\qquad c_i(t)\in\mathbb R.
\]
The original orbit has
\[
 [\Phi,\dot\Phi]=-r^2\Omega,\qquad
 b\dot b^*-\dot b b^*=-B^2\Omega,\qquad
 J_0=-(r^2+B^2)\Omega.
\]
For verification, $H^2=r^2\operatorname{diag}(0,1,1)$,
$HAH=0$, and hence $[H,[A,H]]=-r^2A$; the vector term
is $-B^2A$. Conjugation back by $O(t)K$, with the rotation commuting
with $\Omega$, gives the stated original-coordinate identities.
Now
\[
 F_{ij}=0,\qquad
 E_i=i\dot c_i\Phi+i c_i\dot\Phi,\qquad
 \sum_i[\mathcal A_i,E_i]
          =r^2\left(\sum_i c_i^2\right)\Omega.
\]
Thus the full Gauss residual, with precisely the convention in
\eqref{q444:full-gauss}, is
\[
 \boxed{\left(\frac{2r^2}{g^2}\sum_i c_i^2+r^2+B^2\right)\Omega.}
\]
Its coefficient is strictly positive and $\Omega\neq0$, proving the
obstruction even before the spatial gauge equations are used.

For completeness those equations force all $c_i$ to vanish as well:
here $J_i=F_{ij}=0$, so they require
$i(\ddot c_i\Phi+2\dot c_i\dot\Phi+c_i\ddot\Phi)=0$.
At a fixed time conjugate by the retained moving frame. The $(1,1)$
entries of $H$ and $[A,H]$ vanish, while that of $[A,[A,H]]$ is
$-2ax\neq0$. Hence $c_i=0$. The remaining Gauss residual is then
$(r^2+B^2)\Omega$, still nonzero.
\end{proof}

\begin{proposition}[Arbitrary temporal connection and exact Gauss inertia]
\label{q444:gauss-inertia}
For arbitrary homogeneous scalar and spatial-gauge positions and
velocities, define
\[
 J_{\mathrm{sc}}=
 [\Phi,\dot\Phi]+b\dot b^*-\dot b b^*,\qquad
 \mathscr I_{\Phi,b}Z=
 \bigl([\Phi,[\Phi,Z]]+bb^*Z+Zbb^*\bigr)_{\mathrm{tf}},
\]
\[
 \mathscr K_{\mathcal A}Z
       =-\sum_i[\mathcal A_i,[\mathcal A_i,Z]].
\]
Then Gauss' equation is exactly
\begin{equation}\label{q444:inertia-equation}
 \left(\mathscr I_{\Phi,b}+\frac2{g^2}\mathscr K_{\mathcal A}\right)
 \mathcal A_0
 =(J_{\mathrm{sc}})_{\mathrm{tf}}
       -\frac2{g^2}\sum_i[\mathcal A_i,\dot{\mathcal A}_i].
\end{equation}
The operator on the left is self-adjoint and positive semidefinite for
$\langle Z,W\rangle=-\operatorname{tr}(ZW)$ on $\mathfrak{su}(3)$.
Its kernel is exactly
\[
 \{Z:[Z,\Phi]=0,\ Zb=0,\ [Z,\mathcal A_i]=0\text{ for every }i\}.
\]
The right-hand side always annihilates that kernel. Thus Gauss' equation
has a pointwise solution for $\mathcal A_0$, unique modulo this kernel.
This pointwise statement asserts neither a smooth selection across
kernel-rank changes nor a solution of the remaining field equations.
\end{proposition}
\begin{proof}
Expansion gives
$J_0=J_{\mathrm{sc}}-\mathscr I_{\Phi,b}\mathcal A_0$
after trace-free projection, and
\[
 \sum_i[\mathcal A_i,E_i]
 =\sum_i[\mathcal A_i,\dot{\mathcal A}_i]
          +\mathscr K_{\mathcal A}\mathcal A_0.
\]
Substitution proves \eqref{q444:inertia-equation}. Cyclicity gives
\[
 \langle Z,\mathscr I_{\Phi,b}Z\rangle
       =\operatorname{tr}[Z,\Phi]^2+2|Zb|^2,\qquad
 \langle Z,\mathscr K_{\mathcal A}Z\rangle
       =\sum_i\|[\mathcal A_i,Z]\|_{\mathfrak{su}}^2,
\]
where $\|T\|_{\mathfrak{su}}^2=-\operatorname{tr}T^2$.
More explicitly, cyclicity gives the bilinear identities
\[
 \langle Z,\mathscr I_{\Phi,b}W\rangle
 =\operatorname{tr}([Z,\Phi][W,\Phi])
                  +2\operatorname{Re}((Zb)^*Wb),\qquad
 \langle Z,\mathscr K_{\mathcal A}W\rangle
 =\sum_i\langle[\mathcal A_i,Z],[\mathcal A_i,W]\rangle.
\]
These are symmetric in $Z,W$ and prove self-adjointness; the sums of
squares prove the stated kernel. If $Z$ is in the kernel, the matter moment
paired with $Z$ vanishes because $[Z,\Phi]=Zb=0$, and
$\operatorname{tr}(Z[\mathcal A_i,\dot{\mathcal A}_i])
=\operatorname{tr}([Z,\mathcal A_i]\dot{\mathcal A}_i)=0$.
Thus the right-hand side is orthogonal to the kernel. In a finite
dimensional real inner-product space a self-adjoint operator has image
equal to the orthogonal complement of its kernel: this follows by
diagonalizing its real symmetric matrix. It gives existence and the
claimed exact ambiguity.
\end{proof}

For the literal moving orbit, subtracting the original scalar equations
while retaining arbitrary $\mathcal A_0$ gives the exact additional
equations
\begin{align*}
 [\dot{\mathcal A}_0,\Phi]+2[\mathcal A_0,\dot\Phi]
 +[\mathcal A_0,[\mathcal A_0,\Phi]]
 -\sum_i[\mathcal A_i,[\mathcal A_i,\Phi]]&=0,\\
 \dot{\mathcal A}_0b+2\mathcal A_0\dot b+\mathcal A_0^2b
                   -\sum_i\mathcal A_i^2b&=0.
\end{align*}
Their norm consequences are
\[
 \sum_i\operatorname{tr}[\mathcal A_i,\Phi]^2
       =\operatorname{tr}X^2-\operatorname{tr}\dot\Phi^2,\qquad
 \sum_i|\mathcal A_i b|^2=|v|^2-|\dot b|^2.
\]
For example the trace pairing of the first equation uses
$\operatorname{tr}(\Phi[\dot{\mathcal A}_0,\Phi])=0$ and
$\operatorname{tr}(\Phi[\mathcal A_0,\dot\Phi])
=-\operatorname{tr}([\mathcal A_0,\Phi]\dot\Phi)$; expansion of $X^2$
then gives the first identity. The real pairing of the second equation
with $b$ gives the second by anti-Hermiticity. These right-hand sides
need not be zero; temporal gauge cannot be imposed while silently keeping
the same literal scalar functions.

Here is an exact counterexample to extending the temporal-gauge
stabilizer deduction to nonzero $\mathcal A_0$. For real constants
$\alpha,\beta_1,\beta_2,\beta_3$ satisfying
$\alpha^2-\sum_i\beta_i^2=1$, keep the same scalar orbit and set
\[
 \mathcal A_0=(\alpha-1)\Omega,\qquad
 \mathcal A_i=\beta_i\Omega.
\]
All curvatures vanish, while
$D_0\Phi=\alpha[\Omega,\Phi]$, $D_0b=\alpha\Omega b$,
$D_i\Phi=\beta_i[\Omega,\Phi]$, and $D_i b=\beta_i\Omega b$.
Their covariant wave operators therefore equal the original second
derivatives, proving both scalar equations exactly. Taking
$\alpha=5/3$, $\beta_1=4/3$, $\beta_2=\beta_3=0$ gives
$\mathcal A_1b\neq0$: the spatial field does not stabilize the scalars.
Nevertheless $J_0=\alpha J_{\mathrm{sc}}\neq0$, since $|\alpha|\ge1$,
and the zero electric field cannot satisfy Gauss' equation. These are
scalar-compatible flat connections, not full gauge solutions.

There is also no lift of the moving orbit with all spatial
$\mathcal A_i=0$, even when $\mathcal A_0$ is arbitrary.
Indeed $J_{\mathrm{sc}}=-\mathscr I_{\Phi,b}\Omega$, so
\eqref{q444:inertia-equation} gives
$\mathscr I_{\Phi,b}(\mathcal A_0+\Omega)=0$.
Its kernel description implies $D_0\Phi=D_0b=0$.
Equation~\eqref{q444:gauge-b} would then require
$(aI_3+\kappa\Phi)b=ab=0$, impossible because $a>0$ and $B>0$.
This argument, and Theorem~\ref{q444:temporal-no-lift}, do not exclude
the general case with both nonzero $\mathcal A_0$ and nonzero spatial
fields satisfying all coupled equations.

% Insertion-ready continuation; no global macros or document wrappers.
% Requires the existing q444:moving-family and q444:gauge-eom sections.
\subsection{Constant-matter coordinates and an all-time real-connection obstruction}

Retain exactly the parameters and literal scalar functions of
Theorem~\ref{q444:moving-family}. In particular $a>0$, $B,d>0$, and
$x\ne0$. Keep the same original-coordinate $O(t)$, $\Omega$, and
the exact orthogonal map and inverse
\[
 K=\left(v/\sqrt2,w_0/\sqrt6,u/\sqrt3\right),\qquad K^{-1}=K^*.
\]
The squared lengths of $v,w_0,u$ are $2,6,3$, their pairwise inner
products vanish, and $\det(v,w_0,u)=6$; hence $K\in SO(3)$.
Put
\[
 H=K^*\Phi_0K=\begin{pmatrix}0&0&0\\0&-x&d\\0&d&x\end{pmatrix},
 \qquad q=K^*b_0=Be_1,\qquad
 A=K^*\Omega K=\begin{pmatrix}0&-\sqrt a&0\\\sqrt a&0&0\\0&0&0\end{pmatrix}.
\]
Thus $Hq=0$. Write $r^2=x^2+d^2$ and $\mu=r^2+B^2$.
The retained force, including the cubic term, is
\[
 F_H=-aH+\kappa(\operatorname{adj}H-qq^*)_{\mathrm{tf}}
 =\begin{pmatrix}-2ax&0&0\\0&2ax&-ad\\0&-ad&0\end{pmatrix}
 =[A,[A,H]],\qquad A^2q=-aq.
\]
Indeed $s=2(r^2+B^2)=6ax/\kappa$,
$\operatorname{adj}H=\operatorname{diag}(-r^2,0,0)$, and
$(\operatorname{adj}H-qq^*)_{\mathrm{tf}}
=\operatorname{diag}(-s/3,s/6,s/6)$, which give these identities
by direct multiplication.

\begin{proposition}[An exact constant-matter reformulation]
\label{q444:static-lift-system}
Set
\[
 C=K^*O^*\mathcal A_0OK+A,\qquad
 Y_i=K^*O^*\mathcal A_iOK,\qquad
 \mathcal E_i=\dot Y_i+[C,Y_i],
\]
\[
 \mathscr I_{H,q}Z=([H,[H,Z]]+qq^*Z+Zqq^*)_{\mathrm{tf}},\qquad
 \mathscr K_YZ=-\sum_i[Y_i,[Y_i,Z]].
\]
The prescribed literal scalar functions have a homogeneous gauge lift
if and only if the following complete system holds:
\begin{align}
 [\dot C,H]+[C,[C,H]]-\sum_i[Y_i,[Y_i,H]]&=F_H,
 \label{q444:static-lift-phi}\\
 (\dot C+C^2-\sum_iY_i^2)q&=-aq,
 \label{q444:static-lift-b}\\
 (\mathscr I_{H,q}+2\mathscr K_Y/g^2)C
   &=-\frac2{g^2}\sum_i[Y_i,\dot Y_i],
 \label{q444:static-lift-gauss}\\
 \dot{\mathcal E}_i+[C,\mathcal E_i]
 -\sum_j[Y_j,[Y_j,Y_i]]+\frac{g^2}{2}\mathscr I_{H,q}Y_i&=0
 \quad(i=1,2,3).
 \label{q444:static-lift-spatial}
\end{align}
The inverse connection maps are exactly
\[
 \mathcal A_0=OKCK^*O^*-\Omega,\qquad
 \mathcal A_i=OKY_iK^*O^*.
\]
Neither the literal functions nor their time scale has been changed.
\end{proposition}
\begin{proof}
For a spatially constant $U(t)\in SU(3)$, the transformations are
$\Phi'=U\Phi U^*$, $b'=Ub$,
$\mathcal A_0'=U\mathcal A_0U^*-\dot UU^*$, and
$\mathcal A_i'=U\mathcal A_iU^*$.
The product rule gives
$D'_\nu(Ub)=UD_\nu b$ and
$D'_\nu(U\Phi U^*)=U(D_\nu\Phi)U^*$.
Their commutators give $\mathcal F'_{\nu\rho}
=U\mathcal F_{\nu\rho}U^*$. Traces, trace projection, and the potential
are invariant under this unitary action; adjugate covariance follows
from the inverse formula on invertible matrices and polynomial
continuity. Thus all the equations of Proposition~\ref{q444:gauge-eom},
including Gauss' equation, transform covariantly.

Taking $U=K^*O^*$ gives the constant matter fields $H,q$ and the displayed
connection maps, because $-\dot UU^*=K^*\Omega K=A$.
The scalar covariant derivatives are $[C,H]$ and $Cq$; their second
derivatives are $[\dot C,H]+[C,[C,H]]$ and $(\dot C+C^2)q$.
The trace-free currents are
$(J_0)_{\mathrm{tf}}=-\mathscr I_{H,q}C$ and
$(J_i)_{\mathrm{tf}}=-\mathscr I_{H,q}Y_i$.
Finally
$\sum_i[Y_i,\mathcal E_i]
=\sum_i[Y_i,\dot Y_i]+\mathscr K_YC$.
Substitution gives every displayed equation with its stated sign and
factor. The explicitly invertible maps prove both directions.
\end{proof}

\begin{theorem}[No real-connection lift with any spatial time dependence]
\label{q444:real-no-lift}
The literal moving-kernel scalar functions admit no homogeneous lift
with $\mathcal A_\nu(t)\in\mathfrak{so}(3)$ for all four indices, on
any nonempty open time interval. The temporal connection is arbitrary;
the spatial connections may be arbitrary twice differentiable functions,
including rotating, multifrequency, and nonperiodic functions. No
commutativity or spatial-stationarity assumption is imposed.
\end{theorem}
\begin{proof}
Since $K,O,\Omega$ are real, the exact frame maps preserve real
skew-symmetry in both directions. Define
\[
 L_1=\begin{pmatrix}0&0&0\\0&0&-1\\0&1&0\end{pmatrix},\quad
 L_2=\begin{pmatrix}0&0&1\\0&0&0\\-1&0&0\end{pmatrix},\quad
 L_3=\begin{pmatrix}0&-1&0\\1&0&0\\0&0&0\end{pmatrix},\quad
 L(z)=\sum_jz_jL_j\quad(z\in\mathbb R^3).
\]
Direct multiplication gives
\[
 L(z)w=z\times w,\quad
 L(z)L(w)=wz^{\mathsf T}-(z\cdot w)I_3,\quad
 [L(z),L(w)]=L(z\times w),\quad
 -\operatorname{tr}(L(z)L(w))=2z\cdot w.
\]
These formulas specify the commutator orientation and every norm factor.
Write
\[
 C=L(c),\quad Y_i=L(y_i),\quad h=\dot c,\quad
 P=\sum_i y_iy_i^{\mathsf T},\quad
 Q=cc^{\mathsf T}-P=(q_{jk})_{j,k=1}^3.
\]
Here the scalar entries $q_{jk}$ do not replace the earlier fixed matter
vector $q=Be_1$. Since $A=\sqrt a L_3$, the residual in
\eqref{q444:static-lift-b}, divided by $B$, is
\[
 R_b=\begin{pmatrix}a-q_{22}-q_{33}\\h_3+q_{12}\\-h_2+q_{13}\end{pmatrix}.
\]
The residual in \eqref{q444:static-lift-phi} is a symmetric matrix
$R_H$ whose six independent entries are
\begin{align*}
 (R_H)_{11}&=2ax-4dq_{23}+2xq_{22}-2xq_{33},\\
 (R_H)_{12}&=dh_2+3dq_{13}+xh_3-3xq_{12},\\
 (R_H)_{13}&=-dh_3+3dq_{12}+xh_2+3xq_{13},\\
 (R_H)_{22}&=-2ax-2dh_1+2dq_{23}+4xq_{11}+2xq_{33},\\
 (R_H)_{23}&=ad-4dq_{11}-dq_{22}-dq_{33}-2xh_1,\\
 (R_H)_{33}&=2dh_1+2dq_{23}-4xq_{11}-2xq_{22}.
\end{align*}
These entries follow by multiplying the explicitly given matrices in
\[
 R_H=[L(h),H]+\sum_{j,k}q_{jk}[L_j,[L_k,H]]
                  -a[L_3,[L_3,H]],\qquad
 R_b=L(h)e_1+\sum_{j,k}q_{jk}L_jL_ke_1+ae_1.
\]
Thus no spatial time-dependence ansatz has entered this calculation.

The equation $R_b=0$ gives $q_{22}+q_{33}=a$, $h_3=-q_{12}$,
$h_2=q_{13}$. The two equations $(R_H)_{12}=(R_H)_{13}=0$ then give
$dq_{13}-xq_{12}=0$ and $dq_{12}+xq_{13}=0$.
Since $x^2+d^2>0$, they force
$q_{12}=q_{13}=h_2=h_3=0$.
Next $(R_H)_{11}=0$ gives $xq_{22}=dq_{23}$.
After these substitutions, $(R_H)_{23}=(R_H)_{33}=0$ give
$xh_1=-2dq_{11}$ and $dh_1=2xq_{11}$.
They imply $2(x^2+d^2)q_{11}=0$, and hence $q_{11}=h_1=0$.
Conversely the following values make every residual entry zero:
\begin{equation}\label{q444:real-scalar-covariance}
 \dot c=0,\qquad Q=
 \begin{pmatrix}0&0&0\\0&d\tau/x&\tau\\0&\tau&a-d\tau/x\end{pmatrix},
 \qquad \tau=q_{23}(t)\in\mathbb R.
\end{equation}
This is the exact necessary-and-sufficient scalar equation within the
real algebra. In particular $Qe_1=0$ and $\operatorname{tr}Q=a>0$.

The covariance $P=\sum_i y_iy_i^{\mathsf T}$ is positive semidefinite.
Suppose $c_1\ne0$ at any time. For $v\in e_1^\perp$ define
$z=v-(c\cdot v)e_1/c_1$. Then $c\cdot z=0$ and $Qe_1=0$, so
\[
 v^{\mathsf T}Qv=z^{\mathsf T}Qz
 =(c\cdot z)^2-z^{\mathsf T}Pz
 =-\sum_i(y_i\cdot z)^2\leq0.
\]
The lower $2\times2$ block of $Q$ would therefore be negative
semidefinite, contradicting its positive trace $a$.
Equivalently, $P_{11}=c_1^2$ and the Schur complement of $P_{11}$ in
$P=cc^{\mathsf T}-Q$ is exactly minus that lower block of $Q$.
Thus $c_1=0$ at every time. Since
$Q_{11}=c_1^2-\sum_i y_{i1}^2=0$, every $y_{i1}=0$ as well.
Differentiation on the open time interval now gives
\[
 c,y_i,\dot y_i\in e_1^\perp,\qquad
 \sum_i y_i\times\dot y_i\in\mathbb R e_1.
\]

The exact inertia on this real algebra is
\[
 \mathscr I_{H,q}L_1=4r^2L_1,\qquad
 \mathscr I_{H,q}L_2=\mu L_2,\qquad
 \mathscr I_{H,q}L_3=\mu L_3.
\]
To prove these formulas, $H^2=r^2\operatorname{diag}(0,1,1)$ and
$HL_1H=-r^2L_1$ give
$[H,[H,L_1]]=4r^2L_1$, while the $qq^*$ terms vanish.
For $j=2,3$, multiplication gives $HL_jH=0$,
$H^2L_j+L_jH^2=r^2L_j$, and
$qq^*L_j+L_jqq^*=B^2L_j$.
Every resulting real-skew matrix is already traceless. This proves the
three identities and gives $\mathscr I_{H,q}C=\mu C$.

For $Z,W\in\mathfrak{su}(3)$ set
$\langle Z,W\rangle=-\operatorname{tr}(ZW)$ and
$\|Z\|^2=\operatorname{tr}(Z^*Z)=-\operatorname{tr}Z^2$.
Cyclicity gives
$\langle Z,[Y,W]\rangle=-\langle[Y,Z],W\rangle$, so
\[
 \langle C,\mathscr K_YC\rangle
 =\sum_i\|[Y_i,C]\|^2
 =2\sum_i|y_i\times c|^2.
\]
Pair the full Gauss equation \eqref{q444:static-lift-gauss} with $C$.
Its right side vanishes because $c\perp\sum_i y_i\times\dot y_i$.
The left side, with the exact factor $2$ from the trace inner product, is
\[
 2\mu|c|^2+\frac4{g^2}\sum_i|y_i\times c|^2=0.
\]
Since $\mu>0$ and $g^2>0$, this forces $c=0$.
Then $Q=-P$ has
$\operatorname{tr}Q=-\sum_i|y_i|^2\leq0$, contradicting
$\operatorname{tr}Q=a>0$.
Every full solution must obey the two scalar equations and Gauss'
equation used here. Their contradiction already excludes a full lift,
without imposing or deleting any spatial gauge equation.
\end{proof}

\begin{corollary}[Constant spatial fields with arbitrary complex color]
\label{q444:static-complex-no-lift}
There is no lift with all $Y_i$ constant in the constant-matter frame,
even when $C(t)$ is arbitrary and the fields take values in the full
complex algebra $\mathfrak{su}(3)$.
\end{corollary}
\begin{proof}
Cyclicity and anti-Hermiticity give
\[
 \langle Z,\mathscr I_{H,q}Z\rangle
 =\operatorname{tr}[Z,H]^2+2|Zq|^2,\qquad
 \langle Z,\mathscr K_YZ\rangle=\sum_i\|[Y_i,Z]\|^2.
\]
For completeness, trace projection pairs trivially with traceless $Z$;
cyclicity gives
$-\operatorname{tr}(Z[H,[H,Z]])=\operatorname{tr}[Z,H]^2$, and
$-\operatorname{tr}(Z(qq^*Z+Zqq^*))=-2q^*Z^2q=2|Zq|^2$.
The second formula follows from the same commutator trace identity used
in the theorem. The matrix $[Z,H]$ is Hermitian, so all displayed
terms are nonnegative.

When $\dot Y_i=0$, Gauss' equation is
$(\mathscr I_{H,q}+2\mathscr K_Y/g^2)C=0$.
Its pairing with $C$ therefore forces $Cq=0$, $[C,H]=0$, and
$[C,Y_i]=0$ for every $i$ at every time.
Since $q$ is constant, $Cq=0$ implies $\dot Cq=0$ and $C^2q=0$.
The vector equation becomes
$(-\sum_iY_i^2+aI_3)q=0$. Pairing with $q$ gives
$\sum_i|Y_iq|^2+aB^2=0$, which contradicts $a>0$, $B>0$.
\end{proof}

\begin{corollary}[Matter-stabilizer gauge saturation]
\label{q444:real-stabilizer-no-lift}
There is no lift for which a differentiable real $\theta(t)$ makes
\[
 C'=SCS^*-i\dot\theta H,\qquad Y_i'=SY_iS^*,\qquad
 S(t)=\exp(i\theta(t)H)
\]
all real skew-symmetric. Nor is there a lift when such a transformation
makes all $Y_i'$ constant, without a reality restriction.
\end{corollary}
\begin{proof}
$H$ is Hermitian and traceless, so $S\in SU(3)$.
It commutes with $H$, and $Hq=0$ gives $Sq=q$.
Differentiation gives $\dot SS^*=i\dot\theta H$.
Thus these are exactly the gauge transformations proved in
Proposition~\ref{q444:static-lift-system}, retaining the same constant
matter fields. They would transform any claimed lift into one excluded
by Theorem~\ref{q444:real-no-lift} or
Corollary~\ref{q444:static-complex-no-lift}, respectively.
\end{proof}

These are exact exclusions for the stated classes, not an assertion
that every $\mathfrak{su}(3)$ connection is real or belongs to their
matter-stabilizer gauge saturation. Arbitrary time-dependent complex
$\mathfrak{su}(3)$ lifts outside those classes remain open here.

\subsection{Constructive homogeneous gauge solutions and their observations}

\begin{theorem}[A real radial family of full gauge solutions]
\label{q444:radial-gauge}
Let $D=\operatorname{diag}(-2,1,1)$ and retain
\[
 \Phi(t)=x(t)D,\qquad b(t)=\beta(t)e_1,\qquad
 x(t),\beta(t)\in\mathbb R,\qquad \mathcal A_\mu=0.
\]
This is a full homogeneous gauge--scalar solution exactly when
\begin{align*}
 \ddot x&=-a x-\kappa x^2+\frac{\kappa}{3}\beta^2,\\
 \ddot\beta&=-(a-2\kappa x)\beta,\qquad
 a=m^2+\Gamma(6x^2+2\beta^2).
\end{align*}
For every real initial position and velocity these equations have a
global solution. They give finite-energy homogeneous solutions on a
finite-volume flat $\mathbb T^3$, and local solutions on $\mathbb R^3$,
without asserting finite total energy on $\mathbb R^3$.
\end{theorem}
\begin{proof}
Here $[\Phi,\dot\Phi]=0$ and
$b\dot b^*-\dot b b^*=0$ because $x,\beta$ are real.
Thus every $J_\mu$ vanishes identically and so do all curvatures.
The gauge equations hold exactly, not after discarding Gauss' equation.
Directly
\[
 (\operatorname{adj}(xD)-\beta^2e_1e_1^*)_{\mathrm{tf}}
          =\left(-x^2+\frac{\beta^2}{3}\right)D,\qquad
 (aI_3+\kappa xD)\beta e_1=(a-2\kappa x)\beta e_1.
\]
The two scalar equations are consequently the displayed real system.
Its conserved density is
\[
 \mathcal E=3\dot x^2+\dot\beta^2+
 V(xD,\beta e_1),\qquad
 V=\frac{m^2}{2}s+\frac{\Gamma}{4}s^2
               -2\kappa x(\beta^2-x^2),\quad s=6x^2+2\beta^2.
\]
The stated conservation follows by differentiating and substituting the
two real equations; the $x$ kinetic coefficient is $3$, not $1/2$.
Moreover $|x|\le\sqrt{s/6}$ and $\beta^2+x^2\le s/2$, so
$|2x(\beta^2-x^2)|\le s^{3/2}/\sqrt6$ and
\[
 V\ge\frac{\Gamma}{4}s^2-\frac{|\kappa|}{\sqrt6}s^{3/2}.
\]
This is coercive and globally bounded below. Conservation bounds the
positions and velocities. The polynomial first-order vector field is
locally Lipschitz and bounded with bounded derivative on their compact
range; hence no finite-time escape is possible, proving global existence.
Multiplying the finite density by the finite volume of $\mathbb T^3$
gives finite total energy there.
\end{proof}

\begin{proposition}[A full-gauge vacuum and its marked Lorentz observation]
\label{q444:vacuum-observation}
At $\Gamma=m^2=1$, $\kappa=6$, let
\[
 x_0=\frac{3+\sqrt3}{12},\qquad
 \Phi=x_0\operatorname{diag}(-2,1,1),\qquad b=3x_0e_1.
\]
Together with $\mathcal A_\mu=0$ this is an equilibrium full gauge
solution at a global minimum of the scalar potential.
For $Q=\Theta(\Phi,b)$ its spectrum is $(4x_0,-2x_0,-2x_0)$.
Its spectral-floor matrix $H=Q+2x_0I_3$ has rank one and is positive.
The complex observation and hidden coordinate are
\[
 G=3x_0\begin{pmatrix}1&-1\\-1&2\end{pmatrix},
 \qquad \rho=3x_0,\qquad \det G=9x_0^2=|\rho|^2.
\]
More generally its $SU(3)$ orbit, parameterized by unit vectors
$e=ge_1$, has
\[
 \Phi=x_0(I_3-3ee^*),\quad b=3x_0e,\quad
 G=3x_0P(I_3-ee^*)P^*,\quad
 \det G=9x_0^2|u^*e|^2=|\rho|^2.
\]
For fixed $P$, this one gauge orbit contains both null and timelike
observations. The fixed-$P$ determinant is therefore not a
gauge-invariant mass observable.
\end{proposition}
\begin{proof}
For a nonzero real equilibrium with $\beta=3x>0$, the radial equations
reduce to
$24\Gamma x^2-2\kappa x+m^2=0$.
The displayed $x_0$ is the larger root at $(1,1,6)$.
Its quaternionic matrix is
\[
 Q=\begin{pmatrix}-2x_0&0&0\\0&x_0&3x_0\mathbf j\\
                       0&-3x_0\mathbf j&x_0\end{pmatrix}.
\]
The lower block has eigenvalues $x_0\pm3x_0$, proving the spectrum.
For an independent global-minimum check, every trace-free quaternionic
Hermitian matrix satisfies $54N(Q)^2\le s(Q)^3$: for eigenvalues summing
to zero, the exact identity is
$s^3-54N^2=2\prod_{i<j}(\lambda_i-\lambda_j)^2$.
Thus, with $r=\sqrt{s}\ge0$,
$V\ge r^2/2+r^4/4-2r^3/\sqrt6$, with equality at eigenvalues
$(2y,-y,-y)$, $y\ge0$, $r=\sqrt6y$.
On this equality family
$V=3y^2+9y^4-12y^3$ and
$dV/dy=6y(1-6y+6y^2)$.
The larger positive root is
$y=(3+\sqrt3)/6=2x_0$; the derivative is positive, negative,
then positive on the three successive intervals, and this minimum has
negative value, since at $y=2/3$ the potential is $-4/9$.
It is therefore the global minimum, attained by the displayed $Q$.

The complex part of $H$ is
$3x_0\operatorname{diag}(0,1,1)$, and compression by the retained $P$
gives the displayed $G$. The $SU(3)$ action on pairs is
$(\Phi,b)\mapsto(g\Phi g^*,gb)$; it is exactly quaternionic congruence
under $\Theta$, since $B(gb)=gB(b)g^{\mathsf T}$ for $g\in SU(3)$.
Consequently the spectrum and spectral-floor shift are unchanged.
Writing $v=Pe$ and using the determinant lemma in dimension two gives
\[
 \det(C-vv^*)=\det C\,(1-v^*C^{-1}v)
 =3\left(1-e^*\left(I_3-\frac{uu^*}{3}\right)e\right)
 =|u^*e|^2.
\]
This proves the general observation formula. A unit $e\in u^\perp$
gives determinant zero, while $e=e_1$ gives a positive determinant.
These vectors are related by $SU(3)$: extend either to a unitary basis
and adjust the phase of one of the remaining columns to make determinant
one. Thus the examples are genuinely on one gauge orbit.

The exact covariance map retains the observer as part of the data:
under $H'=gHg^*$, $P'=Pg^*$ one has
$P'H'P'^*=PHP^*$. Complex-component projection therefore has the same
value. Simultaneously $u'=gu$, $b'=gb$ gives $u'^*b'=u^*b$.
This is covariance of the marked observer, not invariance of a
fixed observer under an active transformation of the state.
\end{proof}

\begin{proposition}[Gauge averaging and distinct minimum orbits]
\label{q444:vacuum-average}
Let $d\mu$ be Haar measure on $SU(3)$ with its retained positive total
mass $V_\mu$, not set to one. For every complex Hermitian $A$,
\[
 \det(PAP^*)=u^*\operatorname{adj}(A)u,\qquad
 \int_{SU(3)}\det(PgAg^*P^*)\,d\mu(g)
                    =V_\mu\operatorname{tr}\operatorname{adj}(A).
\]
For a quaternionic positive rank-one matrix $H=\Theta(A,b)$,
$\operatorname{adj}(A)=bb^*$. Hence its averaged observation is
$V_\mu\|b\|^2$.

Put $\chi=2x_0>0$ for the preceding minimum and, for
$0\le\theta\le\pi/4$, define
\[
 v_\theta=(\cos\theta,\ \sin\theta\,\mathbf j,\ 0)^{\mathsf T},
 \qquad Q_\theta=\chi(3v_\theta v_\theta^*-I_3).
\]
Every $Q_\theta$ is a scalar global minimum, and its complex-plus-triplet
coordinate satisfies
\[
 b_\theta=(0,0,-3\chi\cos\theta\sin\theta)^{\mathsf T},\qquad
 \|b_\theta\|^2=\frac{9\chi^2}{4}\sin^2(2\theta).
\]
Consequently this minimum family contains pairwise $SU(3)$-inequivalent
states with invariant $\|b\|^2$ ranging from $0$ to $9\chi^2/4$.
The coefficient agrees with the supplied packet's invariant formula;
this calculation validates that coefficient rather than claiming it
as a newly discovered one.
\end{proposition}
\begin{proof}
The signed $2$-by-$2$ minors of $P$ form precisely $u=(1,1,1)^{\mathsf T}$.
The exterior-square form of the Cauchy--Binet expansion therefore gives
$\det(PAP^*)=u^*\operatorname{adj}(A)u$; equivalently this follows by
expanding the six entries of $PAP^*$ and collecting the cofactors.
The identity $\operatorname{adj}(gAg^*)=g\operatorname{adj}(A)g^*$
holds for invertible $A$ and then for every $A$ by polynomial continuity.
The Haar average of these conjugates commutes with every element of
$SU(3)$. Commutation with diagonal determinant-one phase matrices kills
its off-diagonal entries, and commutation with determinant-one cyclic
permutations makes its diagonal entries equal. Its trace fixes the
average at $V_\mu\operatorname{tr}(\operatorname{adj}A)I_3/3$.
Since $u^*u=3$, the stated integral follows with all factors retained.

Write $r=b^*b$. The skew convention gives directly
$B(b)\overline{B(b)}=bb^*-rI_3$.
The complex part of $H^2$ is therefore $A^2-bb^*+rI_3$,
and $e_2(H)=e_2(A)-r$.
For positive rank-one $H$, its spectral identity is
$H^2=(\operatorname{tr}H)H$ and $e_2(H)=0$.
Taking complex parts and using $\operatorname{tr}H=\operatorname{tr}A$
gives
$A^2-(\operatorname{tr}A)A+e_2(A)I_3=bb^*$.
The left side is $\operatorname{adj}(A)$, proving the assertion.

Finally $v_\theta^*v_\theta=1$, so $Q_\theta$ has fixed spectrum
$(2\chi,-\chi,-\chi)$ and hence the same minimum value.
Its $(1,2)$ quaternionic entry is
$-3\chi\cos\theta\sin\theta\,\mathbf j$, with the other quaternionic
off-diagonal coefficients zero. The retained definition of $B(b)$
therefore gives the displayed $b_\theta$ with its minus sign.
Taking its norm proves the formula. It is strictly increasing on the
given interval; since $b\mapsto gb$ preserves its norm, distinct
parameters there cannot lie on one $SU(3)$ orbit.
The positive spectral-floor matrix is $3\chi v_\theta v_\theta^*$,
so the averaging identity applies to the same triplet coordinate.
\end{proof}

The classical $SU(3)$ cohomogeneity-one geometry of $\mathbb{HP}^2$
is discussed by A.~Gambioli, \emph{Eight-dimensional $SU(3)$-manifolds
of cohomogeneity one}, arXiv:math/0611796, Section~4, the
$\mathbb{HP}^2=M(\mathbb{CP}^2,S^5)$ example,
\url{https://arxiv.org/abs/math/0611796}.
The explicit family and its separating invariant above are proved
directly; no new classification of this classical orbit space is claimed.
The averaging formula supplies the exact relation of that invariant to
the marked Lorentz determinant. A continuum of classical potential
minima is not, by itself, a calculation of a quantum Hamiltonian spectrum.

% Ready-to-input proof fragment. All edits are confined to agent_work.
% Requires amsmath, amssymb and the standard theorem/proposition/lemma
% environments. Labels beginning q444:vk- are private to this fragment.
\subsection{Vacuum kinetics and exact current-free continuations}

Retain the real scalar state space
\[
 \mathcal E=\operatorname{Herm}_3(\mathbb C)_0\oplus\mathbb C^3,
 \qquad
 \langle(\phi,c),(\psi,d)\rangle
 =\operatorname{tr}(\phi\psi)+2\operatorname{Re}(c^*d).
\]
The subscript \(0\) means trace zero, not a quotient. The scalar kinetic
density is exactly one half of this metric:
\[
 \mathcal L_{\rm sc}=\tfrac12\operatorname{tr}
       (D_\mu\Phi D^\mu\Phi)+(D_\mu b)^*D^\mu b-V,\qquad
 V=\tfrac{m^2}{2}s+\tfrac{\Gamma}{4}s^2
                   -\kappa(\det\Phi-b^*\Phi b),\quad
 s=\operatorname{tr}\Phi^2+2b^*b .
\]
The spacetime metric is \((+---)\), \(m^2\geq0\), \(\Gamma>0\);
the nonzero vacuum results below use \(\kappa>0\).
The gauge action is
\((\Phi,b)\mapsto(g\Phi g^*,gb)\), \(g\in SU(3)\).
Its anti-Hermitian connection is \(\mathcal A_\mu\), and the retained
gauge density is
\(\operatorname{tr}(\mathcal F_{\mu\nu}\mathcal F^{\mu\nu})/(2g^2)\),
with \(g^2>0\). All constructed connections below are identically zero.
We make no claim about connections not considered here.

For \(Z_{\rm tf}=Z-\operatorname{tr}(Z)I_3/3\), variation gives, at
\(\mathcal A_\mu=0\),
\begin{align}
 \Box\Phi&=-a\Phi+\kappa(\operatorname{adj}\Phi-bb^*)_{\rm tf},
       &a&=m^2+\Gamma s,\label{q444:vk-pde}\\
 \Box b&=-(aI_3+\kappa\Phi)b,\qquad&
 J_\mu&=[\Phi,\partial_\mu\Phi]
       +b(\partial_\mu b)^*-(\partial_\mu b)b^* .\notag
\end{align}
Indeed
\[
 dV(\delta\Phi,\delta b)
 =\operatorname{tr}\!\left(
 [a\Phi-\kappa(\operatorname{adj}\Phi-bb^*)_{\rm tf}]\delta\Phi\right)
 +2\operatorname{Re}\!\left[
 ((aI_3+\kappa\Phi)b)^*\delta b\right].
\]
This proves both scalar equations with their original metric factors.
Variation in \(\mathcal A_\mu\) pairs the matter variation with
\(\operatorname{tr}(J^\mu\delta\mathcal A_\mu)\).
Consequently \(J_\mu=0\) and \(\mathcal F_{\mu\nu}=0\) satisfy every
gauge equation, including Gauss, before any restriction of variations.

We retain the quaternionic convention
\[
 \Theta(\Phi,b)=\Phi+B(b)\mathbf j,\qquad
 B(b)=\begin{pmatrix}
 0&\bar b_3&-\bar b_2\\-\bar b_3&0&\bar b_1\\
 \bar b_2&-\bar b_1&0
 \end{pmatrix},\qquad \mathbf j z=\bar z\mathbf j .
\]
Its real trace metric is exactly the displayed metric on \(\mathcal E\):
the diagonal sum of \(B(b)\overline{B(c)}\) has real part
\(-2\operatorname{Re}(b^*c)\). Its cubic is
\(N(\Theta(\Phi,b))=\det\Phi-b^*\Phi b\). To verify the latter identity,
write \(\Phi_{23}=z_1,\Phi_{31}=z_2,\Phi_{12}=z_3\) and use the
Hermitian Moore formula
\[
 N\begin{pmatrix}a&q_3&\bar q_2\\\bar q_3&d&q_1\\q_2&\bar q_1&f\end{pmatrix}
 =adf-a|q_1|^2-d|q_2|^2-f|q_3|^2
                 +2\operatorname{Re}(q_1q_2q_3).
\]
Here \(q_i=z_i+\bar b_i\mathbf j\). The complex part of the triple
product is
\(z_1z_2z_3-\bar b_1b_2z_3-z_1\bar b_2b_3-\bar b_1\bar z_2b_3\).
Its extra real terms and the three extra diagonal terms are precisely
\(-b^*\Phi b\).

\begin{lemma}[Retained nonzero minimum and its curve]\label{q444:vk-vacuum}
Suppose \(0\leq m^2\leq\kappa^2/(27\Gamma)\), \(\kappa>0\), and put
\[
 \chi=\frac{\kappa+\sqrt{\kappa^2-24\Gamma m^2}}{12\Gamma}>0 .
\]
For a unit \(v\in\mathbb H^3\), \(Q=\chi(3vv^*-I_3)\) is a global
minimum of \(V\) on the trace-free quaternionic Hermitian state space.
The origin is an additional minimum only at
\(m^2=\kappa^2/(27\Gamma)\).
For \(0\leq\theta\leq\pi/4\), let
\[
 v_\theta=(\cos\theta,\sin\theta\,\mathbf j,0)^{\mathsf T},
 \qquad Q_\theta=\chi(3v_\theta v_\theta^*-I_3).
\]
Its pair coordinates and potential value are
\begin{align}
 \Phi_\theta&=\chi\operatorname{diag}
     (3\cos^2\theta-1,\,3\sin^2\theta-1,\,-1),\notag\\
 b_\theta&=-3\chi\cos\theta\sin\theta\,e_3,\label{q444:vk-curve}\\
 V_0&=3m^2\chi^2+9\Gamma\chi^4-2\kappa\chi^3
      =\kappa\chi^3-9\Gamma\chi^4 .\notag
\end{align}
\end{lemma}
\begin{proof}
The quaternionic spectral theorem gives real eigenvalues
\(\lambda_1+\lambda_2+\lambda_3=0\) and
\(s=\sum\lambda_i^2\). On each sphere \(s>0\), extrema of the product
have at most two distinct eigenvalues: subtracting two Lagrange equations
for the product gives
\((\lambda_i-\lambda_j)(-\lambda_k-2\beta)=0\).
If all three were distinct they would all equal \(-2\beta\), a
contradiction. The two-eigenvalue trace-free possibilities are
\((2t,-t,-t)\) and their negatives and permutations, where
\(t=\sqrt{s/6}>0\). Thus \(N\leq2t^3\), sharply, and
\[
 V\geq F(t)=3m^2t^2+9\Gamma t^4-2\kappa t^3,\qquad
 F'(t)=6t(m^2-\kappa t+6\Gamma t^2).
\]
For \(m^2>0\) the larger positive root is \(\chi\), the only positive
local minimum; the smaller root is a local maximum. For \(m^2=0\),
\(F'(t)<0\) on \(0<t<\chi\) and \(>0\) on \(t>\chi\).
Using \(m^2=\kappa\chi-6\Gamma\chi^2\) yields
\(F(\chi)=\chi^3(\kappa-9\Gamma\chi)\), negative below the displayed
threshold and zero at it. Comparing with \(F(0)=0\) proves the global
claim and all equality cases used here. The spectrum of \(Q_\theta\)
is \((2\chi,-\chi,-\chi)\). Its \((1,2)\) entry is
\(-3\chi\cos\theta\sin\theta\,\mathbf j\), giving the stated sign of
\(b_\theta\). Direct substitution gives the potential value.
\end{proof}

\begin{proposition}[Horizontal metric and normal acceleration]
\label{q444:vk-horizontal}
Write \(q_\theta=(\Phi_\theta,b_\theta)\),
\(T=\partial_\theta q_\theta\), \(N=\partial_\theta^2q_\theta\).
Then
\[
 \langle T,T\rangle=18\chi^2,\qquad
 \langle T,N\rangle=0,\qquad
 \langle N,N\rangle=72\chi^2 .
\]
The curve is \(SU(3)\)-horizontal and its induced kinetic density is
\(9\chi^2\dot\theta^2\). Every homogeneous unconstrained solution
confined to this fixed-\(\chi\) curve has constant \(\theta\).
The vector \(N\) is normal to the entire fixed-\(\chi\) minimum
manifold. Its radial projection is \(-3q_\theta\), and
\(\|N+3q_\theta\|^2=18\chi^2\); hence there is a nonzero normal
component in addition to the overall radial one.
\end{proposition}
\begin{proof}
With \(D=\operatorname{diag}(-1,1,0)\), differentiation gives
\[
 T_\Phi=3\chi\sin(2\theta)D,\quad T_b=-3\chi\cos(2\theta)e_3,
 \quad
 N_\Phi=6\chi\cos(2\theta)D,\quad N_b=6\chi\sin(2\theta)e_3.
\]
Since \(\operatorname{tr}D^2=2\), the three metric identities follow
without changing the coefficient of the triplet metric.
All commutators of the diagonal matrices vanish, and
\(b_\theta T_b^*-T_b b_\theta^*=0\) because both coefficients are real.
For every \(\xi\in\mathfrak{su}(3)\),
\[
 \langle T,([\xi,\Phi_\theta],\xi b_\theta)\rangle
 =\operatorname{tr}\bigl(
   ([\Phi_\theta,T_\Phi]+b_\theta T_b^*-T_b b_\theta^*)\xi\bigr)=0.
\]
Thus this nonzero tangent is horizontal and cannot itself be a
nonzero vertical tangent in this positive metric.
At a vacuum the potential gradient vanishes, so the homogeneous
equation is \(T\ddot\theta+N\dot\theta^2=0\).
Pairing with \(T\) first gives \(\ddot\theta=0\), and the nonzero
norm of \(N\) then gives \(\dot\theta=0\).
The exact vector that would have to be supplied as an extra force
for a constrained constant-speed motion is \(N\dot\theta^2\).
Its force covector is
\[
 \operatorname{tr}(N_\Phi\dot\theta^2\,\delta\Phi)
       +2\operatorname{Re}((N_b\dot\theta^2)^*\delta b).
\]

For the assertion about the whole minimum manifold,
\(v_\theta^*v'_\theta=0\), \(v_\theta'^*v'_\theta=1\), and
\(v_\theta''=-v_\theta\). Therefore
\[
 Q_\theta''=6\chi(v_\theta'v_\theta'^*-v_\theta v_\theta^*).
\]
Every tangent to the rank-one projector manifold has the form
\(3\chi(zv_\theta^*+v_\theta z^*)\), \(v_\theta^*z=0\).
The real trace pairing with both projectors in \(Q_\theta''\)
is zero, by \(v_\theta^*v_\theta'=0\) and cyclicity of real trace.
This proves normality. Finally
\(\|q_\theta\|^2=6\chi^2\) and
\(\langle N,q_\theta\rangle=-18\chi^2\); subtracting the radial
projection gives the asserted remaining norm.
Thus the constant-speed curve is an intrinsic geodesic of the
minimum manifold but not an unconstrained homogeneous solution.
\end{proof}

\begin{proposition}[Exact three-real-coordinate full-gauge sector]
\label{q444:vk-sector}
Let \(h,k,\beta\) be real fields and set
\[
 \Phi=\operatorname{diag}(h+k,h-k,-2h),\qquad b=\beta e_3,
 \qquad\mathcal A_\mu=0.
\]
This ansatz satisfies all full field equations if and only if
\begin{align}
 \Box h&=-ah-\kappa h^2+\tfrac{\kappa}{3}(k^2+\beta^2),\notag\\
 \Box k&=-(a-2\kappa h)k,\qquad
 \Box\beta=-(a-2\kappa h)\beta,\label{q444:vk-sector-eom}\\
 a&=m^2+\Gamma(6h^2+2k^2+2\beta^2).\notag
\end{align}
In particular it is an unconstrained invariant sector for homogeneous
initial data with tangent velocities. Its homogeneous density is
\[
 \mathcal L_{\rm hom}=3\dot h^2+\dot k^2+\dot\beta^2-V(h,k,\beta),
 \quad
 V=\tfrac{m^2}{2}s+\tfrac{\Gamma}{4}s^2
                    -2\kappa h(k^2+\beta^2-h^2),
 \quad s=6h^2+2k^2+2\beta^2 .
\]
Every homogeneous initial condition in this sector has a global
solution for \(m^2\geq0,\Gamma>0,\kappa\in\mathbb R\).
\end{proposition}
\begin{proof}
Direct multiplication gives
\[
 \det\Phi-b^*\Phi b=2h(k^2+\beta^2-h^2),\qquad
 \operatorname{adj}\Phi-bb^*
 =\operatorname{diag}(-2h(h-k),-2h(h+k),h^2-k^2-\beta^2).
\]
Its trace is \(-3h^2-k^2-\beta^2\). Removing one third of this trace
and substituting in \eqref{q444:vk-pde} gives exactly
\eqref{q444:vk-sector-eom}, including the factor \(1/3\) in the
\(h\) equation. The last component of the triplet force is
\(-(a-2\kappa h)\beta\); all other entries remain zero.
The currents vanish for each spacetime index: the matrices and
their derivatives are diagonal, while \(b\) and its derivative
are real multiples of \(e_3\). The zero connection therefore
satisfies all gauge equations. This also proves consistency
under unrestricted variations, not merely stationary points
of a restricted Lagrangian. As a check, that Lagrangian has
Euler--Lagrange weights \(6,2,2\) on \(h,k,\beta\), respectively.

For homogeneous fields the smooth polynomial vector field is tangent
to the displayed real linear sector, so local uniqueness proves
invariance. Its conserved energy density is
\[
 E=3\dot h^2+\dot k^2+\dot\beta^2+V.
\]
If \(r^2=k^2+\beta^2\), then
\[
 |2h(r^2-h^2)|\leq2|h|(h^2+r^2)
 \leq \frac{s^{3/2}}{\sqrt6},\qquad
 V\geq\frac{\Gamma}{4}s^2-\frac{|\kappa|}{\sqrt6}s^{3/2}.
\]
Here \(|h|\leq\sqrt{s/6}\) and \(h^2+r^2\leq s/2\).
This lower bound is coercive, so a fixed energy bounds all positions
and velocities. The polynomial first-order system remains in a
compact set on every finite interval and extends globally.
This is an ODE existence proof, not a general nonlinear PDE
global-existence claim.
\end{proof}

On \(r>0\), introduce \(k=r\cos(2\theta)\),
\(\beta=-r\sin(2\theta)\). The original kinetic density and homogeneous
equations become
\begin{align}
 \mathcal L_{\rm kin}
   &=3\dot h^2+\dot r^2+4r^2\dot\theta^2,\notag\\
 \ddot h&=-ah-\kappa h^2+\kappa r^2/3,\qquad
 \ddot r=4r\dot\theta^2-(a-2\kappa h)r,\label{q444:vk-polar}\\
 (r^2\dot\theta)^{\boldsymbol\cdot}&=0,\qquad
 a=m^2+\Gamma(6h^2+2r^2).\notag
\end{align}
These follow by differentiating \(k,\beta\), resolving the two
equations into radial and angular components, and using
\(k\dot\beta-\beta\dot k=-2r^2\dot\theta\).
The Cartesian sector includes \(r=0\); the polar chart does not.
Even permitting a varying positive amplitude
\(Q=\chi(t)(3v_\theta v_\theta^*-I_3)\) gives \(r=3h>0\).
Then the \(h\) equation implies
 \(3\ddot h=-(a-2\kappa h)r\), whereas the \(r\) equation
has the extra term \(4r\dot\theta^2\). Since \(\ddot r=3\ddot h\),
one obtains \(4r\dot\theta^2=0\). Thus changing this single amplitude
does not by itself make angular homogeneous motion consistent.

\begin{theorem}[Exact gauge-invariant periodic continuation]
\label{q444:vk-periodic}
Suppose \(0<m^2<\kappa^2/(27\Gamma)\), \(\Gamma>0,\kappa>0\).
Define
\[
 d_0=\sqrt{\kappa^2-24\Gamma m^2},\quad
 h_0=\frac{\kappa+d_0}{24\Gamma}=\frac{\chi}{2},\quad
 h_0<h<\frac{\kappa}{6\Gamma},\quad
 \Delta(h)=1-\frac{6\Gamma h}{\kappa}.
\]
Put
\[
 a(h)=\frac{m^2+12\Gamma h^2}{\Delta(h)},\qquad
 R^2(h)=\frac{3a(h)h}{\kappa}+3h^2,\qquad
 \omega^2(h)=a(h)-2\kappa h
   =\frac{m^2-2\kappa h+24\Gamma h^2}{\Delta(h)}.
\]
Take the positive roots \(R,\omega\). Then the sector fields
\[
 h(t)=h,\quad k(t)=R\cos(\omega t),\quad
 \beta(t)=R\sin(\omega t),\quad \mathcal A_\mu=0
\]
are nonconstant, spatially homogeneous, exact full gauge--scalar
solutions. The gauge-invariant observable
\(\|b(t)\|^2=R^2\sin^2(\omega t)\) is nonconstant.
Their conserved energy density is exactly
\begin{equation}\label{q444:vk-periodic-energy}
 E(h)=R^2\omega^2+
 m^2(3h^2+R^2)+\Gamma(3h^2+R^2)^2-2\kappa h(R^2-h^2).
\end{equation}
Let
\[
 \Delta_0=1-\frac{6\Gamma h_0}{\kappa}
          =\frac{3\kappa-d_0}{4\kappa}>0,\qquad
 V_0=12m^2h_0^2+144\Gamma h_0^4-16\kappa h_0^3 .
\]
As \(h\downarrow h_0\),
\begin{align}
 R&\longrightarrow3h_0,\qquad \omega\longrightarrow0,\notag\\
 \omega^2&=\frac{2d_0}{\Delta_0}(h-h_0)
                            +O((h-h_0)^2),\notag\\
 E(h)-V_0&=\frac{18h_0^2d_0}{\Delta_0}(h-h_0)
                            +O((h-h_0)^2)
          =9h_0^2\omega^2+O(\omega^4).\label{q444:vk-excess}
\end{align}
These are finite-volume classical statements. On a flat torus with
retained volume \(\mathcal V_3>0\), the total energy and its excess are
\(\mathcal V_3 E(h)\) and \(\mathcal V_3(E(h)-V_0)\).
The amplitude tends to the nonzero value \(3h_0\): the family approaches
the minimum manifold while traversing a finite range of its invariant
\(\|b\|^2\), rather than becoming a small oscillation about one fixed
vacuum. Its period tends to infinity.
\end{theorem}
\begin{proof}
The denominator is positive on the stated interval; \(a>0\) and
\(R^2>0\). The numerator of \(\omega^2\) has roots
\((\kappa-d_0)/(24\Gamma)\) and \(h_0\); with its positive leading
coefficient it is positive for \(h>h_0\). Moreover
\(d_0<\kappa\), so \(h_0<\kappa/(12\Gamma)<\kappa/(6\Gamma)\);
the interval is nonempty. The identities
\[
 a=m^2+\Gamma(6h^2+2R^2),\quad
 -ah-\kappa h^2+\kappa R^2/3=0,\quad
 \omega^2=a-2\kappa h
\]
follow by substitution. They are precisely the constant-\(h\)
and the two oscillatory equations in \eqref{q444:vk-sector-eom}.
The previous proposition proves all gauge equations.
The triplet norm is preserved by the full \(SU(3)\) action, so its
nonconstant value rules out a purely gauge motion.

Substitution in the conserved energy gives
\eqref{q444:vk-periodic-energy}. At \(h=h_0\),
\(m^2=2\kappa h_0-24\Gamma h_0^2\), so
\(a=2\kappa h_0\), \(R^2=9h_0^2\) and \(\omega^2=0\).
The limiting field at any phase is a global minimum: the
\(2\)-by-\(2\) quaternionic block has eigenvalues \(h_0\pm3h_0\),
and the last eigenvalue is \(-2h_0\); hence the full spectrum is
\((4h_0,-2h_0,-2h_0)\).
Writing the frequency numerator as \(n(h)\),
\(n(h_0)=0\) and \(n'(h_0)=2d_0\). Taylor expansion proves the
frequency estimate. The curve \(h\mapsto(h,R(h),0)\) is smooth
near \(h_0\), and its endpoint is a critical point of \(V\).
Consequently its potential differs from \(V_0\) by
\(O((h-h_0)^2)\). Its kinetic density is
\[
 R^2\omega^2
 =\left(9h_0^2+O(h-h_0)\right)
   \left(\frac{2d_0}{\Delta_0}(h-h_0)+O((h-h_0)^2)\right).
 \]
This proves both energy expansions. Since the leading coefficient
of \(\omega^2\) is strictly positive, replacing \(h-h_0\) by its
smooth inverse as a function of \(\omega^2\) gives the last equality.
For every \(h>h_0\), \(E(h)>V_0\) because the kinetic term is
positive and the potential is bounded below by \(V_0\).
The torus assertion is integration of a constant spatial density.
Nonzero homogeneous energy density on all of \(\mathbb R^3\) is
not asserted to have finite total energy.
\end{proof}

\begin{proposition}[Linearized physical tangent and exact null waves]
\label{q444:vk-wave}
At each fixed vacuum \(q_{\theta_0}\), the perturbation
\[
 \delta\Phi=T_\Phi(\theta_0)\eta,\qquad
 \delta b=T_b(\theta_0)\eta,\qquad
 \delta\mathcal A_\mu=0
\]
solves the full linearized equations precisely when \(\Box\eta=0\).
Its quadratic density is
\[
 9\chi^2\bigl((\partial_t\eta)^2-|\nabla\eta|^2\bigr).
\]
For a homogeneous mode on a torus of volume \(\mathcal V_3\), its
momentum and quadratic Hamiltonian are
\[
 \pi_\eta=18\chi^2\mathcal V_3\dot\eta,\qquad
 H_\eta=\frac{\pi_\eta^2}{36\chi^2\mathcal V_3}.
\]
There is also an exact nonlinear inhomogeneous continuation at fixed
\(\chi\): for any smooth real \(f\),
\[
 \theta(t,x)=f(t-x_1),\qquad q(t,x)=q_{\theta(t,x)},\qquad
 \mathcal A_\mu=0
\]
solves the full nonlinear field equations wherever these fields are
defined. On the rectangular flat torus with sides \(L_1,L_2,L_3\),
choose \(f\) \(L_1\)-periodic. Its excess over the vacuum energy is
\[
 E_{\rm wave}-\mathcal V_3V_0
 =18\chi^2L_2L_3\int_0^{L_1}(f'(u))^2\,du,\qquad
 \mathcal V_3=L_1L_2L_3.
\]
\end{proposition}
\begin{proof}
Differentiate \(\nabla V(q_\theta)=0\) in \(\theta\). The metric
Hessian annihilates \(T\). The linearized currents are
\[
 \bigl([\Phi_{\theta_0},T_\Phi]+b_{\theta_0}T_b^*
                         -T_b b_{\theta_0}^*\bigr)\partial_\mu\eta=0.
\]
Thus the linearized scalar equation is \(T\Box\eta=0\), and \(T\ne0\).
The gauge equations have zero current and zero field strength.
The kinetic norm \(18\chi^2\) proves the quadratic density and its
Legendre transform with the stated volume.

For the nonlinear assertion, the chain rule in signature \((+---)\)
gives
\[
 \Box q_\theta=T(\theta)\Box\theta
                 +N(\theta)\partial_\mu\theta\,\partial^\mu\theta .
\]
Both coefficients vanish for \(\theta=f(t-x_1)\). Every value remains
a critical point of \(V\); the trigonometric formulas define it for
all real \(\theta\), or one can restrict the range of \(f\) to the
displayed minimum section. The exact currents vanish for all
spacetime derivatives as in the sector proposition. This verifies
every nonlinear equation. The scalar energy density is the sum of
one temporal and one spatial term, each \(9\chi^2(f')^2\), plus \(V_0\).
Periodicity makes its integral independent of time and gives the
formula. For
\(f(u)=\theta_0+\varepsilon\cos(2\pi u/L_1)\) with range inside the
section, the excess is
\(9\chi^2\mathcal V_3\varepsilon^2(2\pi/L_1)^2\).
\end{proof}

\paragraph{Gauge-invariant coordinate and endpoint scope.}
The invariant
\[
 I=\|b_\theta\|^2=\frac{9\chi^2}{4}\sin^2(2\theta),\qquad
 \frac{dI}{d\theta}=\frac{9\chi^2}{2}\sin(4\theta)
\]
has nonzero derivative on \(0<\theta<\pi/4\). Therefore the displayed
linearized mode is visible through the genuine gauge-invariant
linear observable
\(\delta I=(9\chi^2/2)\sin(4\theta_0)\eta\) in the interior.
With \(M=9\chi^2/4\), its interior coordinate expression is
\[
 9\chi^2(\partial\theta)^2
       =\frac{9\chi^2}{16I(M-I)}(\partial I)^2 .
\]
The denominator vanishes at the endpoints because \(I\) ceases to be
a regular linear coordinate there, not because the norm of \(T\)
diverges.

For completeness the stabilizers on this section can be checked
without importing an orbit classification. In the interior the
three diagonal entries of \(\Phi_\theta\) are distinct, so a
stabilizer must be diagonal. Fixing the nonzero \(b_3\) and its
determinant then gives
\(\operatorname{diag}(e^{it},e^{-it},1)\), of dimension one.
At \(\theta=0\), \(b=0\) and the eigenvalues are \(2\chi,-\chi,-\chi\),
giving \(S(U(1)\times U(2))\), of dimension four.
At \(\theta=\pi/4\), fixing \(b_3\ne0\) and the two-dimensional
first eigenspace gives \(\operatorname{diag}(U,1)\), \(U\in SU(2)\),
of dimension three. The gauge-orbit dimensions are consequently
\(7,4,5\), respectively. The fixed-\(\chi\) minimum manifold has
dimension eight: the unit sphere in \(\mathbb H^3\) has dimension
eleven and the free unit-quaternion action \(v\mapsto va\) has
dimension three, with \(vv^*\) precisely its quotient.
Thus the horizontal tangent dimensions within that manifold are
\(1,4,3\), respectively. In particular the interior has exactly one
horizontal vacuum direction. At the endpoints the invariant
linear function \(I\) vanishes to first order; no assertion of a
smooth one-dimensional linear quotient there is made.
This agrees with the classical cohomogeneity-one description in
A.~Gambioli, \emph{Eight-dimensional \(SU(3)\)-manifolds of
cohomogeneity one}, arXiv:math/0611796, Section~4, Examples:
\(\mathbb{HP}^2\) has singular orbit types \(\mathbb{CP}^2\) and
\(S^5\). The kinetic and current statements above are independently
proved here, not imported from that classification.

\paragraph{Fixed observation and classical scope.}
Retain
\[
 P=\begin{pmatrix}1&-1&0\\0&1&-1\end{pmatrix},\qquad
 u=(1,1,1)^{\mathsf T}.
\]
The positive spectral-floor matrix \(H_\theta=Q_\theta+\chi I_3\)
has complex part \(A_\theta=3\chi
\operatorname{diag}(\cos^2\theta,\sin^2\theta,0)\). Consequently
\[
 G_\theta=PA_\theta P^*
 =3\chi\begin{pmatrix}1&-\sin^2\theta\\
                  -\sin^2\theta&\sin^2\theta\end{pmatrix},\qquad
 \rho=u^*b_\theta=-3\chi\cos\theta\sin\theta,\qquad
 \det G_\theta=|\rho|^2=I.
\]
These are direct matrix products with the original \(P,u,\chi\).
They do not assert gauge invariance of a fixed-\(P\) observation away
from this chosen section: under a general gauge transformation its
complex part is \(gA_\theta g^*\), so the observation is
\(PgA_\theta g^*P^*\), not generally \(PA_\theta P^*\).

On a rectangular torus the linearized wave equation has classical
Fourier frequencies
\(\omega_{\mathbf n}^2=4\pi^2\sum_i n_i^2/L_i^2\);
the spatial zero mode has zero linearized frequency. The exact
homogeneous periodic family additionally has arbitrarily small
positive classical frequency and energy excess at fixed volume.
The null waves give exact nonlinear inhomogeneous solutions with
arbitrarily small amplitude. None of these statements quantizes the
model, identifies a Hilbert-space spectrum, proves a quantum mass
gap or its absence, or assigns physical units that the retained
parameters have not been supplied. No finite-total-energy assertion
on all of \(\mathbb R^3\) is inferred from a finite-volume solution.

\subsection{The gauge-vector spectrum and an invariant vector observable}

We now compute masses in the linearized classical field equations of
\eqref{q444:full-gauge-density}, not by assigning a physical meaning to
a determinant in isolation. Retain a positive nonzero minimum amplitude
$\chi$ from Theorem~\ref{q444:vacua}; thus $\kappa>0$,
$0\leq m^2\leq\kappa^2/(27\Gamma)$ and
\[
 \chi=\frac{\kappa+\sqrt{\kappa^2-24\Gamma m^2}}{12\Gamma}.
\]
The nonzero branch is meant at coexistence. Keep $0\leq\theta\leq\pi/4$,
and introduce the notation $q=\cos(2\theta)$, not a change in the field
or its kinetic metric. The same minimum representatives are
\[
 \Phi_\theta=\chi\operatorname{diag}\left(
       \frac{1+3q}{2},\frac{1-3q}{2},-1\right),\qquad
 b_\theta=-\frac{3\chi}{2}\sin(2\theta)e_3,\qquad
 \beta^2:=b_\theta^*b_\theta=\frac{9\chi^2}{4}(1-q^2).
\]
Every constant representative together with $\mathcal A_\mu=0$
solves the full classical field equations.

\begin{theorem}[All eight gauge-inertia eigenvalues]
\label{q444:gauge-spectrum}
Use the positive inner product
$\langle Z,W\rangle=-\operatorname{tr}(ZW)$ on $\mathfrak{su}(3)$.
At the displayed minimum, the exact spectrum of
$\mathscr I_{\Phi_\theta,b_\theta}$ is
\[
\begin{array}{c|c|c}
\text{generator space}&\text{inertia eigenvalue}&
                         \text{real multiplicity}\\ \hline
i\operatorname{diag}(1,-1,0)&0&1\\
i\operatorname{diag}(1,1,-2)&3\chi^2(1-q^2)&1\\
\langle E_{12}-E_{21},\,i(E_{12}+E_{21})\rangle
                         &9\chi^2q^2&2\\
\langle E_{13}-E_{31},\,i(E_{13}+E_{31})\rangle
                         &\frac92\chi^2(1+q)&2\\
\langle E_{23}-E_{32},\,i(E_{23}+E_{32})\rangle
                         &\frac92\chi^2(1-q)&2
\end{array}                                                    \tag{GS.1}
\]
Here $E_{ij}$ are the actual matrix units. The unscaled generators in
the row order with the six off-diagonal generators first have kinetic
Gram matrix $\operatorname{diag}(2,2,2,2,2,2,2,6)$.
The squared masses of transverse solutions of the full linearized
field equations are exactly $g^2/2$ times the eigenvalues in (GS.1).
In particular the diagonal nonzero branch has
\[
 \mu_{\mathrm d}^{\,2}
   =\frac32 g^2\chi^2(1-q^2)=\frac{2g^2}{3}\|b_\theta\|^2.
                                                               \tag{GS.2}
\]
There is one zero eigenvalue for $0<q<1$, four at $q=1$, and three
at $q=0$. These are the Lie-algebra dimensions of the stabilizers
$U(1)_{1,-1,0}$, $S(U(1)\times U(2))$, and $SU(2)$, respectively.
\end{theorem}
\begin{proof}
For a diagonal Hermitian $\Phi=\operatorname{diag}(d_1,d_2,d_3)$,
the operator $[\Phi,[\Phi,Z]]$ multiplies the $ij$ off-diagonal entry
by $(d_i-d_j)^2$. Here $bb^*=\operatorname{diag}(0,0,\beta^2)$.
Consequently $\{bb^*,Z\}$ contributes zero on the $12$ root space
and $\beta^2$ on each of the $13$ and $23$ root spaces; these entries
have zero trace, so no trace-free correction changes them. Retaining
the displayed differences gives
\[
 (d_1-d_2)^2=9\chi^2q^2,\quad
 (d_1-d_3)^2+\beta^2=\frac92\chi^2(1+q),\quad
 (d_2-d_3)^2+\beta^2=\frac92\chi^2(1-q).
\]
For $Z=i\operatorname{diag}(z_1,z_2,z_3)$, $\sum z_i=0$, the
diagonal contribution is
\[
 \mathscr I Z=2i\beta^2z_3
       \left(E_{33}-\frac13I_3\right).
\]
It annihilates $i\operatorname{diag}(1,-1,0)$ and multiplies
$i\operatorname{diag}(1,1,-2)$ by $4\beta^2/3
=3\chi^2(1-q^2)$. The eight displayed generators are linearly
independent and span $\mathfrak{su}(3)$; multiplication and trace give
their stated diagonal Gram matrix. This proves the entire generalized
Hessian spectrum without replacing that metric by an identity matrix.

For the dynamical assertion, perturb only the connection at first order:
$\mathcal A_\nu=\varepsilon Z a_\nu$, with both scalar perturbations zero.
At the constant scalar background the first-order scalar equations
reduce to
$[Z,\Phi_\theta]\partial_\nu a^\nu=0$ and
$Zb_\theta\,\partial_\nu a^\nu=0$.
The first-order current is
$(J^\nu)_{\rm tf}=-\varepsilon\mathscr I(Z)a^\nu$.
For an inertia eigenvector with eigenvalue $\lambda$, the full gauge
equation is therefore
\[
 (\Box+g^2\lambda/2)a^\nu-\partial^\nu(\partial_\mu a^\mu)=0.
                                                               \tag{GS.3}
\]
Choose $a_0=0$ and
$a_j=e_j\cos(k\cdot x-\omega t)$ with $k\cdot e=0$.
Then $\partial_\mu a^\mu=0$ and all scalar equations hold. The temporal
gauge equation is $\sum_j\partial_j\dot a_j=0$, so Gauss is satisfied
as well, not omitted. The remaining equations give
$\omega^2=|k|^2+g^2\lambda/2$. At nonzero frequency or spatial momentum,
the corresponding field strength is nonzero. Thus these are actual
transverse solutions, not pure-gauge variations.
For a specified rectangular flat spatial torus,
$k_j=2\pi n_j/L_j$, with every length $L_j$ retained.
Their positive kinetic sign follows from
$\langle Z,Z\rangle/g^2>0$ in the original gauge density.

For $0<q<1$, every displayed eigenvalue except the first is positive.
At $q=1$ the second and the $23$ pair additionally vanish; at $q=0$
the $12$ pair additionally vanishes. A stabilizer satisfies
$[Z,\Phi_\theta]=0$ and $Zb_\theta=0$ by the proved inertia kernel.
For $0<q<1$ distinct diagonal entries and nonzero $b_\theta$ give
$\operatorname{diag}(e^{i\alpha},e^{-i\alpha},1)$.
At $q=1$, $b_\theta=0$ and the last two entries of $\Phi_\theta$
coincide, giving the stated $U(2)$ subgroup. At $q=0$ the first two
entries coincide and $b_\theta\ne0$ fixes $e_3$, giving the upper-left
$SU(2)$. This proves the full group statements, including the endpoints.
\end{proof}

\begin{proposition}[An exact invariant current and its mass branch]
\label{q444:current-mass}
The full nonlinear gauge--scalar theory has the real, gauge-invariant
conserved vector current
\[
 j_\nu=-i\left(b^*D_\nu b-(D_\nu b)^*b\right),\qquad
 \partial_\nu j^\nu=0.                                      \tag{GS.4}
\]
At any displayed vacuum with $\beta^2>0$, the transverse perturbation
$\mathcal A_\nu=\varepsilon i\operatorname{diag}(1,1,-2)a_\nu$,
with zero scalar perturbations, has
\[
 j_\nu=-4\varepsilon\beta^2 a_\nu+O(\varepsilon^2).
                                                               \tag{GS.5}
\]
It therefore detects the branch (GS.2) in the exact linearized system.
For the positive rank-one spectral-floor matrix
$H_\theta=Q_\theta+\chi I_3=\Theta(A_\theta,b_\theta)$,
\[
 \boxed{\displaystyle
 \mu_{\mathrm d}^{\,2}
 =\frac{2g^2}{3V_\mu}
       \int_{SU(3)}\det(PgA_\theta g^*P^*)\,d\mu(g).}          \tag{GS.6}
\]
The Haar volume $V_\mu>0$ and the original marking $P$ are retained.
At $\beta=0$ the linear coefficient in (GS.5) vanishes; that endpoint
does not supply a nonzero linear interpolator through this current.
\end{proposition}
\begin{proof}
Under $b\mapsto Ub$ and $D_\nu b\mapsto U D_\nu b$, each scalar product
in (GS.4) is unchanged. Their difference is purely imaginary, so $j_\nu$
is real. It is the Noether current for the constant phase action
$b\mapsto e^{i\alpha}b$: allowing $\alpha$ to depend on spacetime
produces $j^\nu\partial_\nu\alpha$ in the variation of the kinetic term.
The potential is phase invariant. Directly, the covariant product rule
gives
\[
 \partial_\nu j^\nu
 =-i\left(b^*D_\nu D^\nu b-(D_\nu D^\nu b)^*b\right).
\]
The exact equation is $D_\nu D^\nu b=-(aI_3+\kappa\Phi)b$;
the matrix in parentheses is Hermitian. The two scalar products
are therefore equal, proving conservation in the full interacting theory.

At the specified constant background,
$b_\theta^*\mathcal A_\nu b_\theta
=-2i\varepsilon\beta^2a_\nu$. Subtracting its conjugate and multiplying
by $-i$ gives $-4\varepsilon\beta^2a_\nu$, with no rescaling of the
Cartan generator. Every off-diagonal generator, and
$i\operatorname{diag}(1,-1,0)$, has zero $33$ entry, so none contributes
to this linear current when scalar perturbations are zero.
For an arbitrary scalar perturbation $\delta b$, its additional linear
term is $2\operatorname{Im}(b_\theta^*\partial_\nu\delta b)$,
a pure gradient at the constant background; the transverse sector
used in (GS.3) sets it to zero. Thus the nonzero transverse mode is
seen by a gauge-invariant observable.
Finally Proposition~\ref{q444:vacuum-average} proves the integral in
(GS.6) equals $V_\mu\beta^2$. Substitution into (GS.2) proves the
exact relation with all coefficients and domains fixed.
\end{proof}

This follows the classical Higgs-phase use of gauge-invariant vector
operators discussed by D.~Tong, \emph{Lectures on Gauge Theory},
Section~2.7.3, especially the fundamental-scalar example on printed
pages~107--108,
\url{https://www.damtp.cam.ac.uk/user/tong/gaugetheory/gt.pdf}.
The eight-eigenvalue calculation, current coefficient, and marked
determinant identity here are derived for this specified adjoint-plus-
fundamental model; they are not transferred from that example.
Equation (GS.6) is a relation to a linearized gauge-vector mass in this
model, not an identification of all entries of $G$ with spacetime momentum,
and not a spectral theorem for pure interacting quantum Yang--Mills.

\begin{corollary}[Variation of the massive branches along the minimum family]
\label{q444:massive-branch-bound}
For $0<q<1$, the smallest strictly positive inertia eigenvalue is
\[
 \lambda_{\rm pos}(q)=
 \begin{cases}
 9\chi^2q^2,&0<q\leq1/2,\\
 \frac92\chi^2(1-q),&1/2\leq q<1 .
 \end{cases}                                                 \tag{GS.7}
\]
It tends to zero at both ends and has unique maximum $9\chi^2/4$
at $q=1/2$. At that point five eigenvalues coincide at $9\chi^2/4$,
two equal $27\chi^2/4$, and one is zero. This positive degeneracy
does not enlarge the stabilizer. The retained spectral traces are
\[
 \operatorname{tr}\mathscr I=\chi^2(36-15\delta),\qquad
 \operatorname{tr}(\mathscr I^2)
       =\chi^4(324-405\delta+171\delta^2),\qquad \delta=1-q^2.
                                                               \tag{GS.8}
\]
\end{corollary}
\begin{proof}
The $13$ eigenvalue is at least the $23$ eigenvalue. Subtracting the
$12$ eigenvalue from the diagonal and $23$ eigenvalues gives
$3\chi^2(1-4q^2)$ and
$\frac92\chi^2(1-q-2q^2)$, respectively; both are nonnegative
exactly for $0\leq q\leq1/2$ on the stated interval.
For $q\geq1/2$, the diagonal minus the $23$ eigenvalue is
$\frac32\chi^2(1-q)(2q-1)\geq0$.
These comparisons prove (GS.7). The first expression increases strictly
and the second decreases strictly, proving the maximum and endpoint
limits. Substitution of $q=1/2$ in (GS.1) gives the multiplicities,
while its kernel remains one-dimensional. Adding the eight eigenvalues
and their squares gives $\chi^2(21+15q^2)$ and
$\chi^4(90+63q^2+171q^4)$, respectively. Replacing the notation
$q^2$ by $1-\delta$ proves (GS.8).
\end{proof}

There is always a zero transverse branch in this linearized system.
The strictly positive branches alone also have no positive lower bound
uniform over the minimum family. These statements concern the proved
quadratic spectrum and its exact parameters. They do not determine
radiative selection of a minimum, a renormalized continuum theory,
or the interacting Hamiltonian spectrum.

\subsection{The retained system and the scope of the result}

The literal constant-matter data are
\[
 H=\begin{pmatrix}0&0&0\\0&-1/2&1/(2\sqrt2)\\
 0&1/(2\sqrt2)&1/2\end{pmatrix},\qquad
 q=\frac{e_1}{2\sqrt2},\qquad
 A=\sqrt2\begin{pmatrix}0&-1&0\\1&0&0\\0&0&0\end{pmatrix}.
\]
The original scalar couplings are $m^2=\Gamma=1$, $\kappa=6$;
$s=\operatorname{tr} H^2+2q^*q=1$ and $a=m^2+\Gamma s=2$.
Write $g^2>0$ for the gauge coupling and keep $\omega\ne0$ real.
Every unknown $C,Z_1,Z_2$ belongs to $\mathfrak{su}(3)$, the real vector space
of traceless anti-Hermitian complex $3\times3$ matrices.
For any matrix $M$, set $M_{\mathrm{tf}}=M-(\operatorname{tr} M)I_3/3$, and define
\[
 \mathcal I X=([H,[H,X]]+qq^*X+Xqq^*)_{\mathrm{tf}},\qquad
 \mathcal KX=-[Z_1,[Z_1,X]]-[Z_2,[Z_2,X]],
\]
\[
 F=[A,[A,H]]=-2H+\operatorname{diag}(-2,1,1).
\]
The last identity follows from $Hq=0$,
$H^2=(3/8)\operatorname{diag}(0,1,1)$ and
$6(\operatorname{adj}H-qq^*)_{\mathrm{tf}}=\operatorname{diag}(-2,1,1)$.
The five equations to be retained are exactly
\begin{align}
 [C,[C,H]]-[Z_1,[Z_1,H]]-[Z_2,[Z_2,H]]&=F,\label{q444:cx:eq:phi}\\
 (C^2-Z_1^2-Z_2^2)q&=-2q,\label{q444:cx:eq:b}\\
 (\mathcal I+2\mathcal K/g^2)C&=(4\omega/g^2)[Z_1,Z_2],\label{q444:cx:eq:gauss}\\
 -\omega^2Z_1-2\omega[C,Z_2]+[C,[C,Z_1]]
             +\mathcal KZ_1+(g^2/2)\mathcal IZ_1&=0,\label{q444:cx:eq:sp1}\\
 -\omega^2Z_2+2\omega[C,Z_1]+[C,[C,Z_2]]
             +\mathcal KZ_2+(g^2/2)\mathcal IZ_2&=0.\label{q444:cx:eq:sp2}
\end{align}

This section does not solve or exclude all triples in $\mathfrak{su}(3)^3$.
It proves necessary identities for every such triple and a no-solution
theorem for the precisely stated complex block class below. No remaining
general equation is assumed true, and no numerical candidate is a result.

\subsubsection*{Exact map to the original literal fields}
For clarity, keep
\[
 u=(1,1,1)^{\mathsf T},\quad v_0=(1,-1,0)^{\mathsf T},\quad
 w_0=(1,1,-2)^{\mathsf T},\quad
 K=\left(v_0/\sqrt2,w_0/\sqrt6,u/\sqrt3\right),
\]
\[
 \Omega=\sqrt{2/3}\begin{pmatrix}0&-1&1\\1&0&-1\\-1&1&0\end{pmatrix},
 \quad O(t)=e^{t\Omega},\quad
 \Phi_0=\frac14\begin{pmatrix}1&1&1\\1&1&1\\1&1&-2\end{pmatrix},
 \quad b_0=\frac14(1,-1,0)^{\mathsf T}.
\]
The dot products and determinant give $K\in SO(3)$, and multiplication
gives $K^*\Phi_0K=H$, $K^*b_0=q$, $K^*\Omega K=A$.
The prescribed scalar fields remain
$\Phi(t)=O(t)\Phi_0O(t)^*$, $b(t)=O(t)b_0$; their frequency is $\sqrt2$.
For a solution of \eqref{q444:cx:eq:phi}--\eqref{q444:cx:eq:sp2}, set
\[
 Y_1=\cos(\omega t)Z_1-\sin(\omega t)Z_2,\quad
 Y_2=\sin(\omega t)Z_1+\cos(\omega t)Z_2,\quad Y_3=0,
\]
\begin{equation}\label{q444:cx:eq:original-inverse}
 \mathcal A_0=OKCK^*O^*-\Omega,\qquad
 \mathcal A_i=OKY_iK^*O^*.
\end{equation}
The inverse is $C=K^*O^*\mathcal A_0OK+A$ and
$Y_i=K^*O^*\mathcal A_iOK$. These are gauge transformations by
$K^*O^*$, including the inhomogeneous temporal term.

Here is a direct verification of the complete finite reduction. In the
constant-matter frame put $E_i=\dot Y_i+[C,Y_i]$. The two scalar
equations are
$[C,[C,H]]-\sum_i[Y_i,[Y_i,H]]=F$ and
$(C^2-\sum_iY_i^2)q=-2q$.
The projected currents are $J_0=-\mathcal IC$ and $J_i=-\mathcal IY_i$.
Gauss is $(2/g^2)\sum_i[Y_i,E_i]+\mathcal IC=0$, and the spatial
equations are
\[
 \dot E_i+[C,E_i]-\sum_j[Y_j,[Y_j,Y_i]]+(g^2/2)\mathcal IY_i=0.
\]
These follow by homogeneous substitution in the covariant equations with
$D_\mu=\partial_\mu+\mathcal A_\mu$, signature $(+---)$ and gauge
density $\operatorname{tr}(\mathcal F_{\mu\nu}\mathcal F^{\mu\nu})/(2g^2)$.
Specifically the adjoint derivative is a commutator, and
$J_\mu=[\Phi,D_\mu\Phi]+b(D_\mu b)^*-(D_\mu b)b^*$; at constant
matter it is minus the displayed inertia. Thus no temporal constraint
has been removed.
The rotation preserves $\sum_iY_i^2$ and every quadratic commutator sum,
and
$\sum_i[Y_i,\dot Y_i]=-2\omega[Z_1,Z_2]$.
Applying $(\partial_t+\operatorname{ad}_C)^2$ to the rotating pair
gives the terms $-\omega^2Z_i$, the two displayed terms with coefficient
$2\omega$, and $[C,[C,Z_i]]$. The third spatial equation is identically
zero. This proves equivalence with all five equations above at every time.
The product rule proves covariance under \eqref{q444:cx:eq:original-inverse},
so this equivalence also holds in the original coordinates.

\subsection{Universal necessary identities in the unrestricted complex class}

Use $\langle X,Y\rangle=-\operatorname{tr}(XY)$ and $\|X\|^2=-\operatorname{tr} X^2$ on
$\mathfrak{su}(3)$. This is a real positive inner product because $X^*=-X$.
Define the nonnegative real quantities
\[
 c=\langle C,\mathcal IC\rangle,\quad
 z=\sum_{j=1}^2\langle Z_j,\mathcal IZ_j\rangle,\quad
 k=\sum_{j=1}^2\|[C,Z_j]\|^2,\quad
 d=\|[Z_1,Z_2]\|^2,\quad z_0=\|Z_1\|^2+\|Z_2\|^2,
\]
and the signed real number $b=\langle C,[Z_1,Z_2]\rangle$.

\begin{proposition}\label{q444:cx:prop:universal}
Every solution of \eqref{q444:cx:eq:phi}--\eqref{q444:cx:eq:sp2} satisfies
\begin{equation}\label{q444:cx:eq:universal}
 c-z=2,\qquad (g^2/2)c+k=2\omega b,\qquad
 \omega^2z_0=k+2d+g^2c+(g^2/2)z.
\end{equation}
In particular $c,k,d,z_0$ are strictly positive, $\omega b>0$, and
\begin{equation}\label{q444:cx:eq:necessary-inequality}
 8kd+2g^2kz\geq g^4c^2.
\end{equation}
\end{proposition}
\begin{proof}
Cyclicity and anti-Hermiticity give
\[
 \langle X,\mathcal IX\rangle=\operatorname{tr}[X,H]^2+2|Xq|^2,
 \qquad
 \langle X,\mathcal KX\rangle=\sum_j\|[Z_j,X]\|^2.
\]
Trace projection pairs trivially with traceless $X$. The first
commutator term follows by cyclically moving $H$ twice, and the vector
term is $-2q^*X^2q=2|Xq|^2$. The second identity follows by applying
$\langle X,[Z,Y]\rangle=-\langle[Z,X],Y\rangle$ twice.
Because $[X,H]$ is Hermitian, all terms are nonnegative.

Pair \eqref{q444:cx:eq:phi} with $H$ using ordinary trace and pair
\eqref{q444:cx:eq:b} with $2q^*$, taking the real part. Their left sides add
to $-c+z$. Their right sides add to
$\operatorname{tr}(HF)-4q^*q=-2\operatorname{tr} H^2-4q^*q=-2$, proving the first identity.
Pair \eqref{q444:cx:eq:gauss} with $C$ to obtain the second.
Pair \eqref{q444:cx:eq:sp1} with $Z_1$ and \eqref{q444:cx:eq:sp2} with $Z_2$.
The gyroscopic terms sum to $4\omega b$, the double $C$ terms sum to
$-k$, and the spatial double commutators sum to $2d$.
The result is
$-\omega^2z_0+4\omega b-k+2d+(g^2/2)z=0$.
Substitute the second identity to obtain the third.

The first identity gives $c\geq2$. The second gives $\omega b>0$.
Thus $b\ne0$, which forces $d>0$, $z_0>0$, and $k>0$, since
$b=\langle[C,Z_1],Z_2\rangle$.
Apply Cauchy--Schwarz to the ordered pairs
$([C,Z_1],[C,Z_2])$ and $(Z_2,-Z_1)$ to obtain
$4b^2\leq kz_0$.
The second identity gives
$\omega^2z_0\geq(k+g^2c/2)^2/k$.
Subtract this lower bound from the third identity and multiply by $4k$;
the result is exactly \eqref{q444:cx:eq:necessary-inequality}.
\end{proof}

These identities rule out any attempted solution with zero spatial
commutator or zero electric coupling to $C$. They do not rule out a
general noncommuting complex triple.

\subsection{A new genuinely complex block exclusion}

Let $P=q q^*/(q^*q)=\operatorname{diag}(1,0,0)$ and put
\[
 \mathfrak p=\{X\in\mathfrak{su}(3):PXP=(I-P)X(I-P)=0\},\qquad
 \mathfrak k=\{X\in\mathfrak{su}(3):[X,P]=0\}.
\]
Thus $\mathfrak p$ consists of arbitrary complex matrices exchanging
the $q$-line with its orthogonal complement; $\mathfrak k$ consists of
the block-diagonal matrices preserving that line. Neither definition
imposes reality of the entries.

\begin{theorem}[Cross--cross--block exclusion]\label{q444:cx:thm:block}
There is no solution of all five equations with
$C,Z_1\in\mathfrak p$ and $Z_2\in\mathfrak k$.
There is likewise no solution for which a constant real rotation of
the ordered pair $(Z_1,Z_2)$ puts it in that class while
$C\in\mathfrak p$.
The assertion allows all complex phases and every $g^2>0$, $\omega\ne0$.
\end{theorem}

\subsubsection*{Exact coordinates used only for the proof}
Set $x=1/2$, $d_0=1/(2\sqrt2)$, $r=\sqrt{x^2+d_0^2}=\sqrt6/4$, and
\[
 S=\frac1{\sqrt{2r(r+x)}}
       \begin{pmatrix}d_0&-(r+x)\\r+x&d_0\end{pmatrix},
 \qquad U=\operatorname{diag}(1,S).
\]
Multiplication gives $S^{\mathsf T}S=I_2$, $\det S=1$,
$U^*HU=\operatorname{diag}(0,r,-r)$, and $U^*q=q$.
All connections and $A$ are conjugated by $U^*$, with inverse
conjugation by $U$. Brackets, trace projection, the force, and the
inertia transform by these exact invertible maps; hence all five
equations are equivalent. The subspaces $\mathfrak p,\mathfrak k$
are preserved because $U$ commutes with $P$.
No field, period, kinetic coefficient, or coupling is divided out.

For $v\in\mathbb C^2$ define
\[
 X(v)=\begin{pmatrix}0&-v^*\\v&0_{2\times2}\end{pmatrix}.
\]
In either of the proof frames below, $H$ has zero first row and column,
and its lower block squares to $r^2I_2$. Since $q^*q=1/8$,
\begin{equation}\label{q444:cx:eq:cross-inertia}
 \mathcal IX(v)=(r^2+1/8)X(v)=\tfrac12X(v).
\end{equation}
Indeed $HX(v)H=0$, $H^2X(v)+X(v)H^2=r^2X(v)$, and the two
$qq^*$ terms sum to $X(v)/8$; all these matrices are traceless.

\begin{lemma}[Scalar classification of the block generator]\label{q444:cx:lem:scalar}
Under the hypotheses of Theorem~\ref{q444:cx:thm:block}, the two scalar
equations put the block generator into exactly one of the following
two forms by an explicitly invertible constant unitary map fixing $q$:
\begin{align*}
 \textnormal{(I)}\quad&H=\operatorname{diag}(0,r,-r),
      &&Z_2=i\operatorname{diag}(0,\sigma,-\sigma),\quad\sigma\in\mathbb R;\\
 \textnormal{(II)}\quad&H=\begin{pmatrix}0&0&0\\0&0&r\\0&r&0\end{pmatrix},
      &&Z_2=i\operatorname{diag}(e,0,-e),\quad e\ne0.
\end{align*}
In both cases write $C=X(v)$, $Z_1=X(w)$ and
$U_v=vv^*-ww^*=(u_{jk})$. In case (I) scalar compatibility is exactly
\[
 u_{11}=1-\delta>0,\quad u_{22}=1+\delta>0,
 \qquad\delta=\frac1{2r}=\sqrt{2/3}.
\]
In case (II) it is exactly
\[
 \operatorname{tr} U_v=2+e^2,\qquad \operatorname{Re}u_{12}=-\delta\ne0.
\]
\end{lemma}
\begin{proof}
Start in the displayed diagonal-$H$ frame. An arbitrary element of
$\mathfrak k$ is $Z_2=i\operatorname{diag}(\eta,B)$, where $B$ is
Hermitian, $\operatorname{tr} B=-\eta$, and $\eta$ is real. Write it without deleting
any coefficient as
\[
 B=-\frac\eta2I_2+
 \begin{pmatrix}z&x_1-i x_2\\x_1+i x_2&-z\end{pmatrix},
 \qquad x_1,x_2,z\in\mathbb R,
 \qquad D=\operatorname{diag}(r,-r).
\]
The lower block of the adjoint scalar equation is
\[
 -\{U_v,D\}+[B,[B,D]]=I_2-2D.
\]
Matrix multiplication gives
\[
 [B,[B,D]]=4r
 \begin{pmatrix}
 x_1^2+x_2^2&-z(x_1-i x_2)\\
 -z(x_1+i x_2)&-(x_1^2+x_2^2)
 \end{pmatrix}.
\]
Its off-diagonal equation is $z(x_1-i x_2)=0$.
The two diagonal equations give
$u_{11}-u_{22}=-1/r$ and
$\operatorname{tr} U_v=2+4(x_1^2+x_2^2)$.
The vector equation gives $\operatorname{tr} U_v=2+\eta^2$.
Thus $\eta^2=4(x_1^2+x_2^2)$.

If $x_1=x_2=0$, then $\eta=0$ and $B=\operatorname{diag}(z,-z)$.
This is case (I) with $\sigma=z$, and its diagonal scalar equations
give the two values stated. They are positive because $\delta^2=2/3<1$.

Otherwise $z=0$ and $e:=\eta\ne0$.
Let $\gamma=2(x_1-i x_2)/\eta$, so $|\gamma|=1$.
Choose either complex square root $\zeta$ of $-\bar\gamma$ and set
\[
 W=\frac\zeta{\sqrt2}\begin{pmatrix}\gamma&\gamma\\1&-1\end{pmatrix}.
\]
Then $|\zeta|=1$, $W^*W=I_2$, $\det W=1$,
$W^*BW=\operatorname{diag}(0,-e)$, and $W^*DW=r\left(
\begin{smallmatrix}0&1\\1&0\end{smallmatrix}\right)$.
These identities follow by multiplying the two columns, which are the
zero and $-e$ eigenvectors of $B$, with the stated phases.
Conjugation by $\operatorname{diag}(1,W)^*$, with inverse its adjoint,
therefore gives case (II), fixes $q$, and preserves both block spaces.

With $\sigma_x=\left(\begin{smallmatrix}0&1\\1&0\end{smallmatrix}\right)$,
in this frame the lower scalar equation is
$-\{U_v,r\sigma_x\}+e^2r\sigma_x=I_2-2r\sigma_x$.
Its diagonal entries give $\operatorname{Re}u_{12}=-1/(2r)$ and
its off-diagonal entries give $\operatorname{tr} U_v=2+e^2$.
The vector equation gives the same trace. Conversely direct substitution
sets every scalar residual entry to zero in each of the two cases.
\end{proof}

\subsubsection*{All gauge equations in the two frames}
In either case write
\[
 Z_2=i\operatorname{diag}(\eta,d_1,d_2),\quad
 n_j=\eta-d_j,\quad V_j=|v_j|^2,\quad W_j=|w_j|^2,
\]
\[
 V=V_1+V_2,\quad W_0=W_1+W_2,\quad N=V+W_0,
\]
\[
 p=v^*w=\alpha+i b,\quad b_j=\operatorname{Im}(\bar v_jw_j),
\]
\[
 A_j=\frac{g^2}{4}+n_j^2+W_0>0,\quad
 P_j=3b+2\omega n_j,\quad B_j=A_j-(\omega^2+N).
\]
Here $b$ is only the imaginary inner product in this section, not the
triplet or the signed pairing in Proposition~\ref{q444:cx:prop:universal}.
Direct block multiplication gives
\begin{equation}\label{q444:cx:eq:block-double}
 -[X(w),[X(w),X(v)]]
       =X\bigl(W_0v+w(\bar p-2p)\bigr),\qquad
 [X(v),Z_2]=X\bigl(i\operatorname{diag}(n_1,n_2)v\bigr).
\end{equation}
Using \eqref{q444:cx:eq:cross-inertia}, the lower first-column equations of
Gauss and the first spatial equation are exactly
\begin{equation}\label{q444:cx:eq:cross-gauge}
 A_jv_j=(\alpha+iP_j)w_j,\qquad
 B_jw_j=-(\alpha-iP_j)v_j\quad(j=1,2).
\end{equation}
There are no other nonzero blocks in these two equations.

The second spatial equation has only a block-diagonal residual.
Write $\mathcal IZ_2=i\operatorname{diag}(\ell_0,\ell_1,\ell_2)$;
it is diagonal in both frames. Its lower diagonal entries are
\begin{equation}\label{q444:cx:eq:spatial-diagonal}
 -\omega^2d_j+2n_j(V_j-W_j)+4\omega b_j+(g^2/2)\ell_j=0.
\end{equation}
Its upper diagonal entry is
\begin{equation}\label{q444:cx:eq:spatial-top}
 -\omega^2\eta-2\sum_jn_j(V_j-W_j)-4\omega b
                                    +(g^2/2)\ell_0=0.
\end{equation}
Its lower off-diagonal entry is
\begin{equation}\label{q444:cx:eq:spatial-off}
 i(n_1+n_2)(v_1\bar v_2-w_1\bar w_2)
                    +2\omega(-v_1\bar w_2+w_1\bar v_2)=0.
\end{equation}
For verification, the lower block of $[C,Z_1]$ is $-vw^*+wv^*$;
the lower block contributed by $[C,[C,Z_2]]+\mathcal KZ_2$ has
entry $i(n_j+n_k)(v_j\bar v_k-w_j\bar w_k)$.
Its upper entry is minus its lower trace. These multiplications give
\eqref{q444:cx:eq:spatial-diagonal}--\eqref{q444:cx:eq:spatial-off} with every sign and
coefficient. They retain the entire second spatial equation.

\begin{lemma}\label{q444:cx:lem:phase}
Every solution of these gauge equations has $\alpha=0$ and
\begin{equation}\label{q444:cx:eq:energy-block}
 \omega^2W_0=\frac{g^2}{4}N+\sum_jn_j^2(V_j+W_j),
\end{equation}
\begin{equation}\label{q444:cx:eq:uj}
 V_j-W_j=\frac{V_j+W_j}{\omega^2+N}
     \left(\omega^2+V-W_0-\frac{g^2}{2}-2n_j^2\right).
\end{equation}
\end{lemma}
\begin{proof}
Conjugate the first equation in \eqref{q444:cx:eq:cross-gauge} and multiply by
$w_j$. Summing gives
$p=\sum_j(\alpha-iP_j)W_j/A_j$.
Taking real parts gives $\alpha=\alpha\sum_jW_j/A_j$.
If $W_0=0$ the sum is zero. Otherwise $A_j>W_0$ gives
$\sum_jW_j/A_j<\sum_jW_j/W_0=1$. Hence $\alpha=0$ in all cases.
The two equations are now $A_jv_j=iP_jw_j$ and $B_jw_j=iP_jv_j$.
Their real pairings with $v_j$ and $w_j$ give
$A_jV_j=-P_jb_j$ and $B_jW_j=P_jb_j$.
Adding them gives
$A_j(V_j+W_j)=(\omega^2+N)W_j$.
Summing and inserting $A_j$ proves \eqref{q444:cx:eq:energy-block}.
Subtracting twice this last formula divided by $\omega^2+N$ from
$V_j+W_j$ proves \eqref{q444:cx:eq:uj}. The denominator is positive because
$\omega\ne0$ and $N\geq0$.
\end{proof}

\begin{proof}[Proof of Theorem~\ref{q444:cx:thm:block}]
Use Lemma~\ref{q444:cx:lem:scalar} and retain every gauge equation as just derived.

In case (I), $\eta=0$, $n_1=-\sigma$, $n_2=\sigma$,
$\mathcal IZ_2=0$, and $V_j-W_j=1\mp\delta>0$.
If $\sigma=0$, then $Z_2=0$ and Gauss paired with $C$ gives
$\langle C,\mathcal IC\rangle+(2/g^2)\langle C,\mathcal KC\rangle=0$.
Both terms are nonnegative and \eqref{q444:cx:eq:cross-inertia} forces $C=0$,
contradicting the two positive scalar differences. Thus $\sigma\ne0$.
Both $v_j$ are nonzero. Gauss then forces $w_j\ne0$ and $P_j\ne0$.
Since $\alpha=0$,
\[
 w_j=i t_jv_j,\qquad t_j=-A_j/P_j\in\mathbb R\setminus\{0\}.
\]
Here $A_1=A_2=g^2/4+\sigma^2+W_0$.
Because $n_1+n_2=0$, \eqref{q444:cx:eq:spatial-off} becomes
$2i\omega(t_1+t_2)v_1\bar v_2=0$.
Thus $t_1+t_2=0$, which implies $P_1+P_2=0$, or $6b=0$.
But \eqref{q444:cx:eq:spatial-top} gives
\[
 b=\frac{\sigma((1-\delta)-(1+\delta))}{2\omega}
   =-\frac{\sigma\delta}{\omega}\ne0.
\]
This contradiction excludes case (I).

In case (II), $\eta=e$, $d_1=0$, $d_2=-e$,
$n_1=e$, $n_2=2e$, and scalar compatibility gives
$V-W_0=2+e^2$ and
$\operatorname{Re}(v_1\bar v_2-w_1\bar w_2)=-\delta\ne0$.
In particular $V_j+W_j>0$ for both $j$.
The transformations $(Z_1,Z_2)\mapsto(-Z_1,-Z_2)$ and
$(C,\omega)\mapsto(-C,-\omega)$ each preserve all five equations,
as substitution verifies. They allow $e>0$ and $\omega>0$ without
changing the given matter fields or their original frequency.

Direct multiplication in this frame gives
\[
 \mathcal IZ_2=ie\operatorname{diag}(1/6,2/3,-5/6).
\]
The adjoint and trace-projected vector parts, respectively, are
\[
 ie\operatorname{diag}(0,2r^2,-2r^2),\qquad
 ie\operatorname{diag}(1/6,-1/12,-1/12).
\]
Their sum proves the formula with $r^2=3/8$.
The right side of \eqref{q444:cx:eq:energy-block} is positive, so $W_0>0$.
Since $N=2W_0+2+e^2>2W_0$ and $n_j^2\geq e^2$, that identity gives
\[
 \omega^2W_0\geq(g^2/4+e^2)N
             >(g^2/2+2e^2)W_0,
 \qquad \omega^2>g^2/2+2e^2.
\]
Formula \eqref{q444:cx:eq:uj} consequently implies
\[
 V_1-W_1=
 \frac{V_1+W_1}{\omega^2+N}
       (\omega^2+2-e^2-g^2/2)>0.
\]
The $j=1$ equation in \eqref{q444:cx:eq:spatial-diagonal} is exactly
\[
 2e(V_1-W_1)+4\omega b_1+eg^2/3=0,
\]
so $b_1<0$. Gauss gives $b_1=-P_1W_1/A_1$.
Since $V_1>0$, it also forces $W_1>0$ and hence $P_1>0$.
Now $P_2=P_1+2e\omega>0$; because $V_2+W_2>0$, Gauss forces
both $V_2,W_2>0$ as well.

Taking imaginary parts in the proof of Lemma~\ref{q444:cx:lem:phase} gives
\[
 b\left(1+3\sum_j\frac{W_j}{A_j}\right)
       =-2\omega\sum_j\frac{n_jW_j}{A_j}<0.
\]
Therefore $b<0$ and $0<P_j=3b+2\omega n_j<2\omega n_j$.
Put $\rho_j=A_j/P_j>0$. Since $A_j>n_j^2$, one has the strict lower bound
\begin{equation}\label{q444:cx:eq:lower-rho}
 \rho_1+\rho_2>
 \frac{n_1+n_2}{2\omega}=\frac{3e}{2\omega}.
\end{equation}
Gauss also gives $w_j=-i\rho_jv_j$.
Substitution into the full off-diagonal spatial equation
\eqref{q444:cx:eq:spatial-off}, followed by division by the nonzero number
$iv_1\bar v_2$, gives
\[
 3e(1-\rho_1\rho_2)-2\omega(\rho_1+\rho_2)=0.
\]
As $\rho_1\rho_2>0$, this proves the contradictory strict upper bound
\[
 \rho_1+\rho_2=
       \frac{3e}{2\omega}(1-\rho_1\rho_2)<\frac{3e}{2\omega}.
\]
This excludes case (II) with arbitrary complex phases of $v,w$.

Finally a constant $SO(2)$ rotation of the ordered pair $(Z_1,Z_2)$
preserves their sum of squares, every quadratic commutator sum, and
$[Z_1,Z_2]$. It commutes with the spatial rotation generator
$\left(\begin{smallmatrix}0&-1\\1&0\end{smallmatrix}\right)$,
so the two spatial equations transform by the same invertible rotation.
It leaves $C,H,q$ unchanged. Thus a solution that could be put in the
excluded class by such a rotation would produce a solution already
excluded, proving the final assertion.
\end{proof}

\subsection{The precise unresolved part}

Theorem~\ref{q444:cx:thm:block} is a complex restriction, not the real-algebra
obstruction restated. In case (II) one has
$q^*Z_2q=ie/8\ne0$. A real-skew matrix paired with the real vector $q$
has $q^*Zq=0$, and any unitary transformation fixing $q$ preserves this
pairing. Thus these nonzero block generators cannot be turned into
real-skew generators by a matter-stabilizer transformation.

Nevertheless the theorem does not put arbitrary $C,Z_1,Z_2$ into its
block class. In particular it assumes $C\in\mathfrak p$ and the stated
decomposition of the spatial pair; neither follows here from all five
unrestricted equations. The general finite system remains unresolved.
The universal identities in Proposition~\ref{q444:cx:prop:universal} continue to
apply outside this class. There is no asserted full-complex lift and no
asserted full-complex no-lift theorem in this section.

The companion \path{check_complex_rotation.py} verifies the constant-frame
matrix reductions with arbitrary complex block coordinates, the two branch
inertia formulas, and the ratio identity by exact arithmetic. It also checks
the universal pairing formulas on one exact full-algebra test triple.
The universal proofs, branch inequalities, symmetry arguments, and original
time-dependent gauge-map equivalence are proved in the text. The checker
does not solve the unrestricted finite system.

\subsection{Exact unrestricted Gauss and frequency elimination}
\label{q444:complex-elimination-section}

The word unrestricted in this subsection refers to all complex colour
coordinates in the specified constant-$C$, two-spatial-rotation ansatz,
not to all time-dependent homogeneous connections. Retain the literal
$H,q,F$, the positive coupling $g^2$, and all five equations
\eqref{q444:cx:eq:phi}--\eqref{q444:cx:eq:sp2}. No block-class restriction
from Theorem~\ref{q444:cx:thm:block} is imposed here.
For arbitrary $Z_1,Z_2\in\mathfrak{su}(3)$ put
\[
 \mathscr B=[Z_1,Z_2],\qquad L=g^2\mathcal I+2\mathcal K.
\]

\begin{lemma}[Positive Gauss inverse on every possible solution]
\label{q444:complex-gauss-inverse}
If $\mathscr B\ne0$, then $L$ is a positive-definite self-adjoint
operator on the original real eight-dimensional metric space
$(\mathfrak{su}(3),-\operatorname{tr})$. Every solution of the five
equations has $\mathscr B\ne0$. Thus its Gauss equation is exactly
\[
 C=\omega T,\qquad T=4L^{-1}\mathscr B.
\]
\end{lemma}
\begin{proof}
Cyclicity gives the full bilinear identities
\begin{align*}
 \langle X,\mathcal IY\rangle
 &=\operatorname{tr}([X,H][Y,H])
                    +2\operatorname{Re}((Xq)^*Yq),\\
 \langle X,\mathcal KY\rangle
 &=\sum_{j=1}^2\langle[Z_j,X],[Z_j,Y]\rangle.
\end{align*}
They are real and symmetric. Their values at $Y=X$ are sums of
nonnegative squared norms, since $[X,H]$ is Hermitian. Consequently
\[
 \ker L=\{X:[X,H]=0,\ Xq=0,\ [X,Z_1]=[X,Z_2]=0\}.
\]
For the converse direction in this equality, each displayed vanishing
condition makes the corresponding double-commutator or vector term
zero directly. For the forward direction, vanishing of their positive
weighted sum forces every one of the norms to vanish.

The eigenvalues of $H$ are the three distinct real numbers $0,r,-r$,
$r=\sqrt6/4$, and $q$ spans its zero eigenline. Conjugate by the exact
matrix $U$ already displayed in the block proof. A matrix commuting with
this diagonal $H$ has no off-diagonal entries. Anti-Hermiticity makes
its diagonal entries imaginary; annihilation of $q$ sets the first
entry to zero; trace zero makes the last two opposite. Therefore
\[
 \{X:[X,H]=0,\ Xq=0\}=\mathbb R\,iH.
\]
If a nonzero element of this line also commuted with both $Z_j$,
both $Z_j$ would commute with $H$. The same distinct-eigenvalue
argument would make them diagonal in its eigenbasis, hence mutually
commuting. This contradicts $\mathscr B\ne0$. Thus $\ker L=0$,
and positivity and finite dimension give the asserted inverse.

For any solution, Proposition~\ref{q444:cx:prop:universal} proves
$\|[Z_1,Z_2]\|^2>0$. Equivalently, if that bracket were zero,
Gauss paired with $C$ would give $c=0$, whereas the scalar contraction
gives $c-z=2$ and $z\geq0$. Finally multiplying Gauss by $g^2$
gives $LC=4\omega\mathscr B$, exactly the stated inverse relation.
No temporal equation was omitted.
\end{proof}

\begin{theorem}[Full residual zero set and a two-sign inverse]
\label{q444:complex-exact-zero-set}
On the explicit open domain
\[
 \mathcal D=\{(Z_1,Z_2,g^2):Z_j\in\mathfrak{su}(3),\ g^2>0,
       \ [Z_1,Z_2]\ne0,\ Tq\ne0\},
\]
retain the just-defined $T$, and set
\begin{equation}\label{q444:complex-frequency}
 \tau=\frac{1/4+\|Z_1q\|^2+\|Z_2q\|^2}{\|Tq\|^2}>0.
\end{equation}
Define the four full residuals, with their original codomains,
\begin{align}
 R_H&=\tau[T,[T,H]]-\sum_{j=1}^2[Z_j,[Z_j,H]]-F
             \in\operatorname{Herm}_3(\mathbb C)_0,\notag\\
 R_q&=(\tau T^2-Z_1^2-Z_2^2)q+2q\in\mathbb C^3,\notag\\
 R_1&=\tau\bigl(-Z_1-2[T,Z_2]+[T,[T,Z_1]]\bigr)
                 +\mathcal KZ_1+(g^2/2)\mathcal IZ_1
                  \in\mathfrak{su}(3),\notag\\
 R_2&=\tau\bigl(-Z_2+2[T,Z_1]+[T,[T,Z_2]]\bigr)
                 +\mathcal KZ_2+(g^2/2)\mathcal IZ_2
                  \in\mathfrak{su}(3).
 \label{q444:complex-full-residuals}
\end{align}
Let $\mathcal Z\subset\mathcal D$ be their common zero set. The map
\begin{equation}\label{q444:complex-two-sign-map}
 (Z_1,Z_2,g^2,\varepsilon)\longmapsto
 \left(C=\varepsilon\sqrt\tau\,T,\ Z_1,Z_2,
       \omega=\varepsilon\sqrt\tau,\ g^2\right),
 \qquad\varepsilon\in\{+1,-1\},
\end{equation}
is a bijection from $\mathcal Z\times\{+1,-1\}$ onto all solutions
of the five retained equations. Its inverse forgets $C,\omega$ and
retains $\varepsilon=\operatorname{sign}(\omega)$.
\end{theorem}
\begin{proof}
For a solution, the preceding lemma gives $C=\omega T$ and a
nonzero spatial bracket. Pair its full vector equation with $q^*$.
Anti-Hermiticity of the three matrices gives exactly
\[
 \omega^2\|Tq\|^2
       =2q^*q+\|Z_1q\|^2+\|Z_2q\|^2
       =\tfrac14+\|Z_1q\|^2+\|Z_2q\|^2>0.
\]
Thus $Tq\ne0$ and $\omega^2=\tau$. The constant $1/4$ is the
original $2q^*q$; no scalar or frequency has been scaled to one.
This contraction determines the frequency but does not replace
the vector equation. Substitution of $C=\omega T$ into all the
scalar and spatial equations gives every entry of
\eqref{q444:complex-full-residuals}. Hence the original solution
determines a point of $\mathcal Z$ and its unique sign.

Conversely, start at any point of $\mathcal Z$ and choose either
sign in \eqref{q444:complex-two-sign-map}. By definition
$LT=4\mathscr B$, so $LC=4\omega\mathscr B$ and Gauss holds.
The substitutions $C=\omega T$, $\omega^2=\tau$ make each of the
other four original residuals identical to its displayed $R$.
They therefore vanish as full matrices or vectors, not only after
contraction. The constructions plainly recover all original
coordinates and signs, proving the bijection.
\end{proof}

This is an exact seventeen-real-variable formulation (sixteen spatial
matrix coordinates and the retained positive $g^2$), not a proof that
its zero set is nonempty or empty. The inverse returns all original
fields through \eqref{q444:cx:eq:original-inverse}. In particular the
nonzero denominator and bracket domain must be retained in any
polynomial elimination. The scalar contraction alone is insufficient:
the accompanying exact test point has that contraction zero while
its full $R_q$ is nonzero. No such point is reported as a solution.

% NEW insertion-ready module. Insert after the already-integrated canonical
% q444:complex-exact-zero-set theorem. Do not insert standalone_context.tex.
\subsection{Complete scalar-coordinate identities and an unrestricted signature constraint}
\label{q444:complex-scalar-coordinate-section}

The following calculation retains all eight real coordinates of each
of $C,Z_1,Z_2$. It is not a cross--cross--block restriction.
Let $x=1/2$, $d_0=1/(2\sqrt2)$, $r=\sqrt6/4$, and
\[
S=\frac1{\sqrt{2r(r+x)}}
\begin{pmatrix}d_0&-(r+x)\\r+x&d_0\end{pmatrix},\qquad
U=\operatorname{diag}(1,S),\qquad\delta=\frac1{2r}=\sqrt{2/3}.
\]
Both columns of the numerator of $S$ have squared length $2r(r+x)$,
their inner product is zero, and its determinant is $2r(r+x)>0$.
Multiplication by the lower block of the literal $H$ shows that the
columns are eigenvectors for $r,-r$ in that order. Thus
\[
U\in SO(3),\qquad U^*HU=H_d=\operatorname{diag}(0,r,-r),\qquad U^*q=q.
\]
Conjugate every connection and force by $U^*$, with inverse conjugation
by $U$. Products, commutators, the adjoint operation, trace projection,
and the five equations all transform by these exact inverse maps.
In this frame $F_d=\operatorname{diag}(-2,1-2r,1+2r)$.

For each original connection $X$ define its eight real coordinates by
\begin{equation}
\widetilde X=U^*XU=
\begin{pmatrix}
ia_X&-\overline{v_{X1}}&-\overline{v_{X2}}\\
v_{X1}&ib_X&z_X\\
v_{X2}&-\overline{z_X}&-i(a_X+b_X)
\end{pmatrix},\qquad
a_X,b_X\in\mathbb R,\quad v_{X1},v_{X2},z_X\in\mathbb C,
\label{q444:complex-eight-coordinates}
\end{equation}
and put $h_X=a_X+2b_X$. This representation and $X=U\widetilde XU^*$
are mutually inverse. In particular no entry or phase has been removed.
The stabilizer coefficient has the original-coordinate formula
\begin{equation}
h_X=\frac{\langle iH,X\rangle}{r}
=\frac{-i\operatorname{tr}(HX)}r.
\label{q444:complex-stabilizer-coefficient}
\end{equation}
Indeed $\operatorname{tr}(H_d\widetilde X)=ir(a_X+2b_X)$.

For a scalar expression $f(X)$ define its signed sum
\[
\mathscr S(f)=f(C)-f(Z_1)-f(Z_2).
\]
This notation retains the temporal plus sign and both spatial minus
signs; it is not a positive covariance. Define
\[
D_a=\mathscr S(a^2),\qquad D_z=\mathscr S(|z|^2),\qquad
D_j=\mathscr S(|v_j|^2)\quad(j=1,2),
\]
\[
m_j=\mathscr S(hv_j),\quad m_z=\mathscr S(hz),\qquad
n_1=\mathscr S(zv_2)+\frac i2\mathscr S(av_1),\quad
n_2=\mathscr S(\bar z v_1)-\frac i2\mathscr S(av_2).
\]

\begin{proposition}[All scalar equations in unrestricted coordinates]
\label{q444:complex-full-scalar-equivalence}
For arbitrary $C,Z_1,Z_2\in\mathfrak{su}(3)$, the two scalar equations
\eqref{q444:cx:eq:phi}--\eqref{q444:cx:eq:b}
hold if and only if the following three real and five complex identities
hold:
\begin{equation}
D_a=4D_z,\qquad
D_1+2D_z=1-\delta,\qquad D_2+2D_z=1+\delta,
\label{q444:complex-scalar-diagonal-identities}
\end{equation}
\begin{equation}
m_1=m_2=m_z=n_1=n_2=0.
\label{q444:complex-scalar-cross-identities}
\end{equation}
\end{proposition}
\begin{proof}
For a diagonal Hermitian matrix with entries $\lambda_j$,
\[
[X,[X,H_d]]_{jk}
=\sum_\ell X_{j\ell}X_{\ell k}
(\lambda_j+\lambda_k-2\lambda_\ell).
\]
Use this identity with $(\lambda_1,\lambda_2,\lambda_3)=(0,r,-r)$
and the full matrix \eqref{q444:complex-eight-coordinates}.
Let $Q_H$ and $Q_q$ denote respectively the left sides minus the right
sides of the two scalar equations in this diagonal frame. Their
diagonal and matter-line entries are
\begin{align*}
(Q_H)_{11}&=2r(D_1-D_2)+2,\\
(Q_H)_{22}&=-2rD_1-4rD_z-(1-2r),\\
(Q_H)_{33}&=2rD_2+4rD_z-(1+2r),\\
(Q_q)_1&=q_1(2-D_a-D_1-D_2),\qquad q_1=1/(2\sqrt2).
\end{align*}
The first of these three matrix entries is minus the sum of the next
two, since the full residual is traceless. The latter two vanish
exactly when $D_1+2D_z=1-\delta$ and $D_2+2D_z=1+\delta$.
Their sum gives $D_1+D_2=2-4D_z$; the vector's first entry vanishes
exactly when $D_a=4D_z$.

For clarity the individual lower matter-column products are
\[
(\widetilde X^2)_{21}=i(a_X+b_X)v_{X1}+z_Xv_{X2},\qquad
(\widetilde X^2)_{31}=-ib_Xv_{X2}-\bar z_Xv_{X1},
\]
and the three independent lower/upper matrix products are
\begin{align*}
[\widetilde X,[\widetilde X,H_d]]_{21}
&=r\bigl(i(a_X-b_X)v_{X1}+3z_Xv_{X2}\bigr),\\
[\widetilde X,[\widetilde X,H_d]]_{31}
&=r\bigl(-i(2a_X+b_X)v_{X2}+3\bar z_Xv_{X1}\bigr),\\
[\widetilde X,[\widetilde X,H_d]]_{23}
&=-2ir h_Xz_X.
\end{align*}
Since $b_X=(h_X-a_X)/2$, the complete remaining residual entries are
\begin{align*}
(Q_q)_2&=q_1(n_1+im_1/2),&
(Q_H)_{21}&=r(3n_1-im_1/2),\\
(Q_q)_3&=q_1(-n_2-im_2/2),&
(Q_H)_{31}&=r(3n_2-im_2/2),\\
&& (Q_H)_{23}&=-2ir m_z.
\end{align*}
The remaining off-diagonal entries of $Q_H$ are the conjugates, since
it is Hermitian. Both $q_1$ and $r$ are nonzero. For either $j$, the
two relevant equations first give $n_j=-im_j/2$ and then
$3n_j-im_j/2=-2im_j=0$. Thus $m_j=n_j=0$; the last matrix entry gives
$m_z=0$. Conversely these five complex equalities make every displayed
off-diagonal and vector entry vanish. Together with the three real
relations already proved, they set the entire scalar residuals to
zero. This proves both directions without a rank or phase restriction.
\end{proof}

\begin{theorem}[Quantitative spacelike stabilizer constraint]
\label{q444:complex-stabilizer-spacelike}
Every scalar-compatible unrestricted triple satisfies the following
inequalities. Put
\[
H_s=h_{Z_1}^2+h_{Z_2}^2,\qquad
S_j=\sum_{\ell=1}^2\bigl(|v_{Z_\ell j}|^2+2|z_{Z_\ell}|^2\bigr),
\qquad c_1=1-\delta>0,\quad c_2=1+\delta>0.
\]
Then
\begin{equation}
h_C^2(S_j+c_j)\le H_s S_j\quad(j=1,2).
\label{q444:complex-stabilizer-margin}
\end{equation}
In particular exactly one of the following possibilities holds:
\[
h_C=h_{Z_1}=h_{Z_2}=0,\qquad\text{or}\qquad
h_C^2<h_{Z_1}^2+h_{Z_2}^2.
\]
Equivalently, the three original-coordinate numbers
$\langle iH,C\rangle,\langle iH,Z_1\rangle,\langle iH,Z_2\rangle$
are all zero or have strictly negative squared norm for the retained
signature $(+--)$.
Every zero of the unrestricted reduced residuals consequently satisfies
\begin{equation}
\tau h_T^2\le H_s
\min\left\{\frac{S_1}{S_1+1-\delta},
\frac{S_2}{S_2+1+\delta}\right\}.
\label{q444:complex-reduced-stabilizer-filter}
\end{equation}
All denominators in this formula are strictly positive.
\end{theorem}
\begin{proof}
For each $j=1,2$ and each $X\in\{C,Z_1,Z_2\}$ put
$V_{Xj}=(v_{Xj},\sqrt2 z_X)\in\mathbb C^2$.
The identities $m_j=m_z=0$ give
\[
h_C V_{Cj}=h_{Z_1}V_{Z_1j}+h_{Z_2}V_{Z_2j}.
\]
Their squared norms and Cauchy--Schwarz on the ordered pair in
$\mathbb C^2\oplus\mathbb C^2$ imply
\[
h_C^2\|V_{Cj}\|^2
\le(h_{Z_1}^2+h_{Z_2}^2)
(\|V_{Z_1j}\|^2+\|V_{Z_2j}\|^2)=H_sS_j.
\]
The diagonal scalar identities give
$\|V_{Cj}\|^2-S_j=D_j+2D_z=c_j$, proving
\eqref{q444:complex-stabilizer-margin}.
Since $0<\delta<1$, each $c_j$ is strictly positive. If $H_s=0$,
the inequality forces $h_C=0$. If $H_s>0$, division by $S_j+c_j>0$
gives $h_C^2\le H_sS_j/(S_j+c_j)<H_s$.
This covers $S_j=0$ as well. The original-coordinate version follows
by multiplication by $r^2>0$ in
\eqref{q444:complex-stabilizer-coefficient}.
Finally $C=\varepsilon\sqrt\tau\,T$ gives $h_C^2=\tau h_T^2$,
so both inequalities give the last formula, independently of the
frequency sign.
\end{proof}

The quantitative inequalities are attained by scalar-compatible
triples; equality in them does not certify the gauge equations.
For an exact witness in the original coordinates retain $A,H,q$ and set
\[
C=\frac53A+\frac43iH,\qquad Z_1=0,\qquad
Z_2=\frac43A+\frac53iH.
\]
The equality $(5/3)^2-(4/3)^2=1$ makes the signed quadratic expressions
equal to those of $(A,0,iH)$: both cross terms cancel in
$C^2-Z_2^2$ and in $[C,[C,H]]-[Z_2,[Z_2,H]]$.
Since $Hq=0$ and $[iH,H]=0$, these are exactly the original two scalar
equations. The $A$ matrix has only matter-cross entries in the diagonal
frame. Their squared moduli are $1-\delta,1+\delta$, either by the
displayed $U$ or by applying the scalar diagonal identities to
$(A,0,0)$. Also $h_A=0$ and $h_{iH}=2r$ by
\eqref{q444:complex-stabilizer-coefficient}. Thus
\[
h_C=8r/3,\qquad h_{Z_2}=10r/3,\qquad
S_j=\frac{16}{9}c_j,
\]
and both sides of \eqref{q444:complex-stabilizer-margin} are equal.
This witness is not a full gauge solution: $[Z_1,Z_2]=0$ contradicts
the necessary noncommutation proved in
Lemma~\ref{q444:complex-gauss-inverse}.

There is a second unrestricted test using both spatial equations, not
only their trace pairing. We give all its coefficients and its exact
nonnegative discriminant below.

\subsection{A two-by-two stress obstruction before evaluating the full residuals}
\label{q444:complex-stress-section}

In the coordinates \eqref{q444:complex-eight-coordinates}, direct
multiplication gives the full inertia matrix
\begin{equation}
U^*(\mathcal IX)U=
\begin{pmatrix}
ia_X/6&-\overline{v_{X1}}/2&-\overline{v_{X2}}/2\\
v_{X1}/2&-ia_X/12&3z_X/2\\
v_{X2}/2&-3\bar z_X/2&-ia_X/12
\end{pmatrix}.
\label{q444:complex-full-inertia-coordinates}
\end{equation}
For off-diagonal entries the double commutator multiplies by the square
of the eigenvalue difference; these are $r^2=3/8$ and $4r^2=3/2$.
The $qq^*$ terms add $1/8$ to the two matter-cross entries.
For the diagonal entries those terms are $i\operatorname{diag}(a_X/4,0,0)$,
whose trace-free part is $i\operatorname{diag}(a_X/6,-a_X/12,-a_X/12)$.
This proves every entry and also
\[
\langle X,\mathcal IX\rangle
=a_X^2/4+|v_{X1}|^2+|v_{X2}|^2+3|z_X|^2.
\]
The eigenvalues of $\mathcal I$ are $0,1/6,1/2,3/2$ with respective
real multiplicities $1,1,4,2$. Indeed the $h$ diagonal direction, the
$a$ diagonal direction with $b=-a/2$, the four cross-entry directions,
and the two lower off-diagonal directions are mutually orthogonal and
are multiplied by precisely these scalars in the displayed formula.

For an arbitrary point of $\mathcal D$ from
Theorem~\ref{q444:complex-exact-zero-set}, define the real
symmetric two-by-two matrices and real scalars
\[
G_{jk}=\langle Z_j,Z_k\rangle,\qquad
D_{jk}=\langle[T,Z_j],[T,Z_k]\rangle,\qquad
M_{jk}=\langle Z_j,\mathcal IZ_k\rangle,
\]
\[
d=\|[Z_1,Z_2]\|^2>0,\qquad
b_T=\langle T,[Z_1,Z_2]\rangle,
\qquad N=G+D-2b_TI_2,\qquad B=dI_2+(g^2/2)M.
\]
All these objects are functions of the unrestricted spatial pair and
$g^2$, with no frequency sign selected. $M$ is positive semidefinite:
for $x\in\mathbb R^2$, $x^TMx=\langle\sum_jx_jZ_j,
\mathcal I\sum_jx_jZ_j\rangle\ge0$. Hence $B$ is positive definite.

\begin{proposition}[Exact stress matrix and nonnegative discriminant]
\label{q444:complex-stress-obstruction}
Every zero of the four full reduced residuals satisfies
\begin{equation}
\tau N=B.
\label{q444:complex-stress-matrix-equation}
\end{equation}
In particular $N$ must be positive definite. On the open domain
$N_{11}>0$, $\det N>0$, define
\begin{equation}
\Delta=(B_{11}N_{22}+B_{22}N_{11}-2B_{12}N_{12})^2
-4\det B\det N.
\label{q444:complex-stress-discriminant}
\end{equation}
Then $\Delta\ge0$, and the stress equation is equivalent to the two
scalar conditions
\begin{equation}
\Delta=0,\qquad \tau=B_{11}/N_{11}.
\label{q444:complex-stress-filter}
\end{equation}
These are necessary tests on the unrestricted reduced zero set, not
replacements for $R_H,R_q,R_1,R_2$.
\end{proposition}
\begin{proof}
For arbitrary $C,Z_1,Z_2,\omega,g^2$, let $Q_1,Q_2$ be the left sides
of the two original spatial equations. Write
$\beta=\langle C,[Z_1,Z_2]\rangle$ and
$D^C_{jk}=\langle[C,Z_j],[C,Z_k]\rangle$. Cyclicity gives the exact
matrix of pairings
\begin{equation}
\langle Z_j,Q_k\rangle
=-\omega^2G_{jk}-D^C_{jk}+(g^2/2)M_{jk}
+(d+2\omega\beta)\delta_{jk}.
\label{q444:complex-spatial-pairing-matrix}
\end{equation}
For the double $C$ commutator the term is $-D^C_{jk}$ by invariance
of the trace inner product. The spatial double commutators have
diagonal pairing $d$ and off-diagonal pairing zero: expand
$\langle Z_j,\mathcal KZ_k\rangle
=\sum_\ell\langle[Z_\ell,Z_j],[Z_\ell,Z_k]\rangle$.
The gyroscopic diagonal pairing in either equation is $2\omega\beta$;
the off-diagonal pairings vanish because $\langle X,[C,X]\rangle=0$.
These facts prove every entry of \eqref{q444:complex-spatial-pairing-matrix},
including its symmetry.
On a reduced zero, $C=\omega T$, $\omega^2=\tau$,
$D^C=\tau D$, and $\omega\beta=\tau b_T$. Setting the pairing matrix
to zero gives $\tau N=B$. Since $B$ is positive definite and $\tau>0$,
$N$ must be positive definite.

Put $n=N_{11}>0$, $b=N_{12}$, $c=N_{22}$, and $h=nc-b^2>0$.
Expansion of the following sum of squares is an exact polynomial identity:
\begin{align}
n^2\Delta={}&
\bigl(B_{11}(nc-2b^2)+2bnB_{12}-n^2B_{22}\bigr)^2\notag\\
&\quad+4h(nB_{12}-bB_{11})^2.
\label{q444:complex-stress-squares}
\end{align}
Indeed the coefficients of $B_{11}^2,B_{22}^2,B_{12}^2,
B_{11}B_{22},B_{11}B_{12},B_{22}B_{12}$ on either side are respectively
$n^2c^2,n^4,4n^3c,-2n^3c+4n^2b^2,-4n^2bc,-4n^3b$.
Thus $\Delta\ge0$. Equality holds exactly when
$nB_{12}=bB_{11}$ and
$B_{11}(nc-2b^2)+2bnB_{12}-n^2B_{22}=0$.
Substitution of the first equality in the second gives
$n(cB_{11}-nB_{22})=0$. Since $n>0$, the two equalities are exactly
$B=(B_{11}/n)N$. Conversely this proportionality makes the sum of
squares zero. Finally $B_{11}>0$, so this proportionality factor is
positive, and $\tau N=B$ holds exactly when $\tau=B_{11}/n$.
This proves the claimed equivalence and every denominator condition.
\end{proof}

The scalar-coordinate equivalence, the quantitative stabilizer
inequality, and the stress discriminant apply to the full complex
finite system. They neither force a cross--cross--block ansatz nor
establish a zero or emptiness of the unrestricted reduced zero set.
Their exact inverse coordinate maps and the reconstruction
\eqref{q444:complex-two-sign-map} preserve all original fields,
couplings, signs, and both frequency branches.

% Insertion-ready after q444:complex-stress-obstruction.
% This module proves a new bounded finite-fibre elimination theorem.
% It re-proves, but does not claim novelty for, the signed norm identity
% already present in coordinate_constraints/orthogonality_normal_form.tex.
\subsection{An exact bounded finite coupling fibre for every unrestricted spatial pair}
\label{q444:finite-coupling-section}

Retain the literal matrices $H,q,F$, all eight real coordinates of each
$Z_j\in\mathfrak{su}(3)$, and the original positive coupling $g^2$.
The inner product remains $\langle X,Y\rangle=-\operatorname{tr}(XY)$.
The letter $t$ below is a polynomial indeterminate whose evaluated value is
the original $g^2$; no physical constant or amplitude is set to one.
This subsection constructs the complete fibre over a fixed ordered pair
$(Z_1,Z_2)$ of the four-residual zero set
\eqref{q444:complex-full-residuals}. The ordered pair is not constrained
to a real subalgebra, a block class, or a prescribed phase relation.

\begin{proposition}[An explicit upper bound for the original coupling]
\label{q444:finite-coupling-bound}
For an arbitrary fixed ordered pair with $[Z_1,Z_2]\ne0$, define
\[
 z_0=\|Z_1\|^2+\|Z_2\|^2,\qquad
 z=\sum_{j=1}^2\langle Z_j,\mathcal IZ_j\rangle,\qquad
 d=\|[Z_1,Z_2]\|^2,
\]
\[
 c_* =2+z,\qquad k_*=2z_0(z_0+4),\qquad
 t_*=
 \frac{k_*z+\sqrt{k_*^2z^2+8k_*d c_*^2}}{c_*^2}.
\]
All these numbers are real, with $z\ge0$ and
$z_0,d,c_*,k_*,t_*>0$.
Every full solution over this ordered pair satisfies
\begin{equation}
 \|C\|^2\le z_0+4,\qquad
 0<g^2\le t_*,\qquad
 0<\omega^2\le
 \frac{k_*+2d+t_*(c_*+z/2)}{z_0}.
 \label{q444:finite-coupling-bounds}
\end{equation}
The same bounds apply to both signs of the original frequency.
\end{proposition}
\begin{proof}
We include the signed norm calculation so that the amplitude estimate has
no unproved input from a proposed coordinate restriction. In the full
coordinates \eqref{q444:complex-eight-coordinates}, put $h_X=a_X+2b_X$.
Multiplication of the diagonal and off-diagonal entries gives
\[
 \|X\|^2=\tfrac32a_X^2+\tfrac12h_X^2
       +2(|v_{X1}|^2+|v_{X2}|^2+|z_X|^2).
\]
The scalar identities
\eqref{q444:complex-scalar-diagonal-identities} give
$D_a=4D_z$ and $D_1+D_2=2-4D_z$. Consequently
\begin{align*}
 \|C\|^2-z_0
 &=\tfrac32D_a+\tfrac12(h_C^2-h_{Z_1}^2-h_{Z_2}^2)
                    +2(D_1+D_2+D_z)\\
 &=4+\tfrac12(h_C^2-h_{Z_1}^2-h_{Z_2}^2).
\end{align*}
For completeness, the sign of the last correction follows directly from
the full scalar cross identities. Write
$V_X=(v_{X1},\sqrt2 z_X)$,
$S_1=\|V_{Z_1}\|^2+\|V_{Z_2}\|^2$, and
$H_s=h_{Z_1}^2+h_{Z_2}^2$. Those identities give
$h_C V_C=h_{Z_1}V_{Z_1}+h_{Z_2}V_{Z_2}$, while the diagonal identities
give $\|V_C\|^2=S_1+1-\sqrt{2/3}$.
Cauchy--Schwarz therefore proves
\[
 h_C^2(S_1+1-\sqrt{2/3})\le H_s S_1.
\]
Since $1-\sqrt{2/3}>0$, this implies $h_C^2\le H_s$, also when
$S_1=0$ or $H_s=0$. This proves $\|C\|^2\le z_0+4$.

For arbitrary anti-Hermitian matrices $X,Y$, diagonalize $X$ unitarily,
writing its eigenvalues as $i\lambda_a$ with real $\lambda_a$.
In that same basis,
\[
 \|[X,Y]\|^2
 =\sum_{a,b}(\lambda_a-\lambda_b)^2|Y_{ab}|^2
 \le2\Bigl(\sum_a\lambda_a^2\Bigr)
               \Bigl(\sum_{a,b}|Y_{ab}|^2\Bigr)
 =2\|X\|^2\|Y\|^2.
\]
Here $(\lambda_a-\lambda_b)^2\le2(\lambda_a^2+\lambda_b^2)
\le2\sum_a\lambda_a^2$ proves the inequality term by term.
Thus the retained $k=\sum_j\|[C,Z_j]\|^2$ satisfies $k\le k_*$.
The universal identities \eqref{q444:cx:eq:universal} give
$c=c_*$ and
\[
 (g^2)^2c_*^2\le8kd+2g^2kz
                    \le8k_*d+2g^2k_*z.
\]
The quadratic on moving the right side to the left has positive leading
coefficient and roots
$\bigl(k_*z\pm\sqrt{k_*^2z^2+8k_*d c_*^2}\bigr)/c_*^2$.
Its negative root is strictly negative, so its nonpositive region at
$g^2>0$ is exactly $g^2\le t_*$. Finally the universal spatial trace
identity is
\[
 \omega^2z_0=k+2d+g^2(c_*+z/2).
\]
Substitution of the two upper bounds proves the last assertion.
Nonzero bracket implies $z_0,d>0$; positivity of $\mathcal I$ and the
definitions prove all remaining strict and weak sign assertions.
\end{proof}

We now eliminate the temporal stabilizer component exactly before
clearing any frequency denominator. This preserves the original
stabilizer vector $J=iH$, whose squared norm is $3/4$.
Set
\[
 W=\{X\in\mathfrak{su}(3):\langle J,X\rangle=0\},\qquad
 \Pi X=X-\frac{\langle J,X\rangle}{3/4}J,
 \qquad \mathscr B=[Z_1,Z_2].
\]
For the fixed spatial pair define
\[
 \eta=\langle J,\mathcal KJ\rangle,\qquad
 w=\Pi\mathcal KJ,\qquad \beta=\langle J,\mathscr B\rangle,
\]
and operators and a vector on the original metric subspace $W$ by
\[
 \mathcal I_7=\mathcal I|_W,\qquad
 \mathcal S_7=\Pi\mathcal K|_W-\frac{w\otimes w}{\eta},\qquad
 f_7=\Pi\mathscr B-\frac{\beta}{\eta}w.
\]
Here $(w\otimes w)x=w\langle w,x\rangle$; no choice of basis
alters any of these maps.

\begin{lemma}[Full stabilizer elimination and polynomial Gauss numerator]
\label{q444:finite-coupling-schur}
If $\mathscr B\ne0$, then $\eta>0$, $\mathcal S_7\succeq0$, and
$\mathcal I_7\succ0$. Define the polynomial endomorphism and its
determinant and adjugate by
\[
 \mathscr M(t)=t\mathcal I_7+2\mathcal S_7,\qquad
 D_7(t)=\det_W\mathscr M(t),\qquad
 V_7(t)=4\operatorname{adj}_W(\mathscr M(t))f_7,
\]
\begin{equation}
 N_7(t)=V_7(t)+
       \frac{2\beta D_7(t)-\langle w,V_7(t)\rangle}{\eta}J.
 \label{q444:finite-coupling-N}
\end{equation}
Then $D_7(t)>0$ for $t>0$, $D_7$ has degree exactly seven and leading
coefficient $3/128$, $\deg V_7\le6$, and $\deg N_7\le7$.
For every original $g^2>0$,
\begin{equation}
 T=4(g^2\mathcal I+2\mathcal K)^{-1}\mathscr B
   =\frac{N_7(g^2)}{D_7(g^2)},\qquad
 Tq=\frac{V_7(g^2)q}{D_7(g^2)}.
 \label{q444:finite-coupling-T}
\end{equation}
In particular all temporal coordinates, including the entire stabilizer
coefficient, are reconstructed by the displayed formula.
\end{lemma}
\begin{proof}
The positive form of $\mathcal K$ gives
$\eta=\sum_j\|[Z_j,J]\|^2\ge0$. If it were zero, both $Z_j$ would
commute with $H$. The three distinct eigenvalues $0,\sqrt6/4,-\sqrt6/4$
of $H$ would make both matrices diagonal in its eigenbasis; hence
$[Z_1,Z_2]=0$, a contradiction. Thus $\eta>0$.
The kernel calculation in Lemma~\ref{q444:complex-gauss-inverse} gives
$\ker\mathcal I=\mathbb RJ$. Self-adjointness makes $W$ invariant, and
positivity then makes $\mathcal I_7$ positive definite.

For $x\in W$ and $s\in\mathbb R$, use
$\langle J,\mathcal Kx\rangle=\langle w,x\rangle$ to complete the square:
\[
 \langle x+sJ,\mathcal K(x+sJ)\rangle
 =\langle x,\mathcal S_7x\rangle
   +\eta\left(s+\frac{\langle w,x\rangle}{\eta}\right)^2.
\]
Taking $s=-\langle w,x\rangle/\eta$ and using positivity of
$\mathcal K$ proves $\mathcal S_7\succeq0$.
Thus $\mathscr M(t)$ is positive definite for every $t>0$.

Write the full Gauss unknown as $T=x+sJ$ with $x\in W$, with no
restriction on $s$. Pair
$(t\mathcal I+2\mathcal K)T=4\mathscr B$ with $J$ and then apply $\Pi$.
These operations give respectively
\[
 s=\frac{2\beta-\langle w,x\rangle}{\eta},\qquad
 (t\mathcal I_7+2\mathcal S_7)x=4f_7.
\]
Conversely these two equations recover the $J$ component and the $W$
projection of Gauss; the orthogonal decomposition
$\mathfrak{su}(3)=W\oplus\mathbb RJ$ therefore recovers the entire
equation. Inverting the positive operator on $W$ yields
$x=V_7(t)/D_7(t)$, and the formula for $s$ yields
\eqref{q444:finite-coupling-N}. Since $Jq=0$, it also yields the second
formula in \eqref{q444:finite-coupling-T}.

Determinants and adjugates may be calculated in any real basis of $W$;
similarity preserves their resulting scalar and endomorphism. The
inertia computation \eqref{q444:complex-full-inertia-coordinates} gives
the seven eigenvalues $1/6,1/2,1/2,1/2,1/2,3/2,3/2$ on $W$.
Thus $\det_W\mathcal I_7=(1/6)(1/2)^4(3/2)^2=3/128$.
The determinant consequently has degree seven with this leading
coefficient; its positivity for $t>0$ follows from positive definiteness.
The adjugate of a seven-by-seven matrix with affine entries has entries
of degree at most six, and \eqref{q444:finite-coupling-N} proves the
remaining degree bound.
\end{proof}

Put
\[
 \alpha=\frac14+\|Z_1q\|^2+\|Z_2q\|^2>0,\qquad
 Q_7(t)=\|V_7(t)q\|^2.
\]
The real polynomial $Q_7$ is nonnegative at every real $t$ and has degree
at most twelve. Its identically-zero case has the following finite test,
which does not require computing a polynomial determinant.

\begin{proposition}[Three exact vector tests for an empty Gauss domain]
\label{q444:finite-coupling-domain-test}
For every fixed noncommuting spatial pair, $Q_7$ is identically zero if
and only if the following three vector equalities all hold:
\begin{equation}
 f_7q=0,\qquad (\mathcal S_7f_7)q=0,\qquad
                 (\mathcal S_7^2f_7)q=0.
 \label{q444:finite-coupling-three-tests}
\end{equation}
When they hold, no positive original coupling over that spatial pair
belongs to the original four-residual domain.
\end{proposition}
\begin{proof}
Put $\mathcal N=\{X\in W:Xq=0\}$. In the complete coordinates
\eqref{q444:complex-eight-coordinates}, membership of $W$ sets $h_X=0$;
the first column condition $Xq=0$ sets $a_X=v_{X1}=v_{X2}=0$.
The only remaining entry is the arbitrary complex $z_X$. Thus
$\dim_{\mathbb R}\mathcal N=2$, and the full inertia formula gives
$\mathcal I_7|_{\mathcal N}=(3/2)\operatorname{id}_{\mathcal N}$.
Self-adjointness also makes $\mathcal N^\perp\cap W$ invariant.
Consequently $\mathcal I_7^{-1}y\in\mathcal N$ if and only if
$y\in\mathcal N$, and on this subspace the inverse is multiplication
by $2/3$.

Suppose first that $Q_7$ is identically zero. For every $t>0$ the vector
$x(t)=4(t\mathcal I_7+2\mathcal S_7)^{-1}f_7$ then belongs to
$\mathcal N$. Set $B_7=\mathcal I_7^{-1}$.
Factoring the pencil and using the exact finite inverse identity
$(I+A)^{-1}=I-A+A^2-A^3(I+A)^{-1}$ gives
\[
 x(t)=\frac4tB_7f_7-\frac8{t^2}B_7\mathcal S_7B_7f_7
       +\frac{16}{t^3}B_7\mathcal S_7B_7\mathcal S_7B_7f_7
       +O(t^{-4}).
\]
Here the remainder has bounded norm after multiplication by $t^4$:
for $A=2B_7\mathcal S_7/t$ and $t>4\|B_7\mathcal S_7\|$, the inverse
has norm at most two, since
$\|(I+A)y\|\ge(1-\|A\|)\|y\|\ge\|y\|/2$.
The finite-dimensional subspace $\mathcal N$ is closed. Multiplying the
expansion by $t$ and taking the limit shows $B_7f_7\in\mathcal N$, hence
$f_7\in\mathcal N$. Subtracting the first term, multiplying by $t^2$
and taking the limit shows $B_7\mathcal S_7B_7f_7\in\mathcal N$,
hence $\mathcal S_7f_7\in\mathcal N$. Subtracting the next term and
multiplying by $t^3$ similarly gives
$\mathcal S_7^2f_7\in\mathcal N$. In the latter two steps every occurrence
of $B_7$ on the already established vectors in $\mathcal N$ is exactly
multiplication by $2/3$. These are precisely the three displayed tests.

Conversely suppose the three tests hold. If $f_7=0$, then $x(t)=0$ for
all $t>0$ and the conclusion follows. If $f_7\ne0$ and
$\mathcal S_7f_7$ is dependent on $f_7$, their one-dimensional span is
invariant under $\mathcal S_7$. Otherwise those two vectors span all of
$\mathcal N$, and the third test makes that two-dimensional span
invariant as well. In either case let $V$ be the resulting invariant
span. It is contained in $\mathcal N$, so it is invariant under
$\mathcal I_7$ and therefore under the positive pencil
$t\mathcal I_7+2\mathcal S_7$.
Its restriction to $V$ is injective and, in finite dimension, onto $V$;
the full inverse thus carries $f_7\in V$ back into $V$. Hence
$x(t)q=0$ for all $t>0$, and the polynomial $Q_7$ vanishes on the
entire positive half-line. A polynomial with infinitely many roots is
identically zero. Formula \eqref{q444:finite-coupling-T} then proves
the asserted exclusion from the original domain for every $g^2>0$.
\end{proof}

Define four polynomial residuals with the original
matrix and vector codomains:
\begin{align}
 P_H(t)&=\alpha[N_7,[N_7,H]]
     -Q_7\left(\sum_j[Z_j,[Z_j,H]]+F\right),\notag\\
 P_q(t)&=\alpha N_7^2q+Q_7(2I_3-Z_1^2-Z_2^2)q,\notag\\
 P_1(t)&=\alpha\left(-D_7^2Z_1-2D_7[N_7,Z_2]
                       +[N_7,[N_7,Z_1]]\right)
                  +Q_7\left(\mathcal KZ_1+\frac t2\mathcal IZ_1\right),\notag\\
 P_2(t)&=\alpha\left(-D_7^2Z_2+2D_7[N_7,Z_1]
                       +[N_7,[N_7,Z_2]]\right)
                  +Q_7\left(\mathcal KZ_2+\frac t2\mathcal IZ_2\right).
 \label{q444:finite-coupling-polynomial-residuals}
\end{align}
All polynomials on the right are evaluated at the same $t$.
Their codomains are respectively
$\operatorname{Herm}_3(\mathbb C)_0$, $\mathbb C^3$,
$\mathfrak{su}(3)$, and $\mathfrak{su}(3)$ for real $t$.
Define the single real polynomial
\begin{equation}
 p_Z(t)=\operatorname{tr}(P_H(t)^2)+P_q(t)^*P_q(t)
                  -\operatorname{tr}(P_1(t)^2)-\operatorname{tr}(P_2(t)^2).
 \label{q444:finite-coupling-one-polynomial}
\end{equation}

% Independent strengthening of finite_coupling_fibres.tex.
% Uses only the definitions stated explicitly in that insertion.
% Canonical and checkpoint files were not edited by the auditor.
\begin{proposition}[The exact degree and positive leading coefficient]
\label{q444:finite-coupling-exact-degree}
Retain an arbitrary fixed ordered pair
$Z_1,Z_2\in\mathfrak{su}(3)$ with $\mathscr B=[Z_1,Z_2]\ne0$,
and all original definitions above: $J,\eta,\beta,\mathcal I_7$,
$D_7,V_7,N_7,Q_7$, $\alpha$, and $P_H,P_q,P_1,P_2,p_Z$.
Put
\[
 \delta_7=\det_W\mathcal I_7=\frac3{128},\qquad
 a_J=\frac{2\beta}{\eta},\qquad A_JX=[J,X],
\]
and define an endomorphism of the original real product metric space
$\mathfrak{su}(3)\oplus\mathfrak{su}(3)$ by
\[
 \mathcal L_{a_J}(X_1,X_2)
    =(X_1+a_J A_JX_2,\ X_2-a_J A_JX_1).
\]
Here $a_J$ is an explicitly defined scalar belonging to the fixed spatial
pair; it is not a new coupling or a rescaling of the original fields.
The coefficient of $t^{28}$ in the original polynomial $p_Z$ is
\begin{equation}
 [t^{28}]p_Z
 =\alpha^2\delta_7^4
   \bigl\|\mathcal L_{a_J}^2(Z_1,Z_2)\bigr\|^2.
 \label{q444:audit-exact-leading-coefficient}
\end{equation}
If $Q_7$ is not identically zero, this coefficient is strictly positive.
Consequently $\deg p_Z=28$ on every such spatial pair.
\end{proposition}
\begin{proof}

First write $N_7=V_7+s_7J$. The previously proved bounds are
$\deg V_7\le6$, $\deg s_7\le7$, and $\deg Q_7\le12$.
Since $[J,H]=0$ and $Jq=0$, the exact expressions are
\[
 [N_7,[N_7,H]]=[V_7,[V_7,H]]+s_7[J,[V_7,H]],\qquad
 N_7^2q=V_7^2q+s_7JV_7q.
\]
Hence $\deg P_H,\deg P_q\le13$. The full spatial formulas give
$\deg P_1,\deg P_2\le14$, including the $tQ_7\mathcal I/2$ terms.
Taking their displayed squared norms yields $\deg p_Z\le28$.
The established degree statements give
\[
 D_7(t)=\delta_7t^7+\text{terms of degree at most six},\qquad
 V_7(t)=\text{terms of degree at most six}.
\]
The defining full stabilizer numerator is
\[
 N_7(t)=V_7(t)+
    \frac{2\beta D_7(t)-\langle w,V_7(t)\rangle}{\eta}J.
\]
Its coefficient of $t^7$ is therefore exactly $a_J\delta_7J$,
also when $\beta=0$ and that coefficient is zero. Since
$\deg Q_7\le12$, its contributions
$Q_7(\mathcal KZ_j+t\mathcal IZ_j/2)$ have degree at most thirteen.
It follows directly from the full spatial numerator formulas that
\begin{align*}
 [t^{14}]P_1
 &=\alpha\delta_7^2
       \bigl(-Z_1-2a_J A_JZ_2+a_J^2A_J^2Z_1\bigr),\\
 [t^{14}]P_2
 &=\alpha\delta_7^2
       \bigl(-Z_2+2a_J A_JZ_1+a_J^2A_J^2Z_2\bigr).
\end{align*}
Direct composition, with both signs retained, gives
\[
 \mathcal L_{a_J}^2(Z_1,Z_2)
 =(Z_1+2a_J A_JZ_2-a_J^2A_J^2Z_1,
   Z_2-2a_J A_JZ_1-a_J^2A_J^2Z_2).
\]
The pair of leading spatial coefficients is thus
$-\alpha\delta_7^2\mathcal L_{a_J}^2(Z_1,Z_2)$.
The scalar numerator degrees satisfy $\deg P_H,\deg P_q\le13$,
so their squared norms contribute no term of degree twenty-eight.
The two spatial squared norms therefore give exactly
\eqref{q444:audit-exact-leading-coefficient}.

It remains to prove strict positivity under the stated domain-polynomial
hypothesis. Invariance of $\langle X,Y\rangle=-\operatorname{tr}(XY)$
gives
\[
 \langle X,A_JY\rangle=-\langle A_JX,Y\rangle.
\]
For arbitrary $X=(X_1,X_2)$ and $Y=(Y_1,Y_2)$, this implies
\begin{align*}
 \langle X,\mathcal L_{a_J}Y\rangle
 &=\langle X_1,Y_1\rangle+\langle X_2,Y_2\rangle
   +a_J\langle X_1,A_JY_2\rangle
   -a_J\langle X_2,A_JY_1\rangle\\
 &=\langle X_1+a_J A_JX_2,Y_1\rangle
   +\langle X_2-a_J A_JX_1,Y_2\rangle
 =\langle\mathcal L_{a_J}X,Y\rangle.
\end{align*}
Thus $\mathcal L_{a_J}$ is self-adjoint on this positive metric space.
If the coefficient in \eqref{q444:audit-exact-leading-coefficient}
were zero, then $\mathcal L_{a_J}^2Z=0$ for $Z=(Z_1,Z_2)$.
Consequently
\[
 \|\mathcal L_{a_J}Z\|^2
 =\langle Z,\mathcal L_{a_J}^2Z\rangle=0,
\]
so $\mathcal L_{a_J}Z=0$. If $a_J=0$, this gives $Z_1=Z_2=0$,
contradicting $\mathscr B\ne0$. Hence $a_J\ne0$, and the two component
equations are
\[
 a_J[J,Z_1]=Z_2,\qquad a_J[J,Z_2]=-Z_1.
\]
The full spatial operator now satisfies the exact identity
\begin{align*}
 a_J\mathcal KJ
 &=-a_J[Z_1,[Z_1,J]]-a_J[Z_2,[Z_2,J]]\\
 &=a_J[Z_1,[J,Z_1]]+a_J[Z_2,[J,Z_2]]\\
 &=[Z_1,Z_2]+[Z_2,-Z_1]=2\mathscr B.
\end{align*}
Apply the original orthogonal projection $\Pi$ to this equality and
use $w=\Pi\mathcal KJ$ and $a_J/2=\beta/\eta$. It yields
\[
 f_7=\Pi\mathscr B-\frac{\beta}{\eta}w
    =\frac{a_J}2w-\frac{a_J}2w=0.
\]
The exact definition $V_7(t)=4\operatorname{adj}_W(\mathscr M(t))f_7$
then makes $V_7$ identically zero, and hence makes
$Q_7(t)=\|V_7(t)q\|^2$ identically zero. This contradicts the
hypothesis. The coefficient is nonnegative by its squared-norm formula,
so it must be strictly positive. Since the global degree bound is
twenty-eight, the degree is exactly twenty-eight.
\end{proof}

\begin{theorem}[Complete finite coupling fibre and its exact inverse]
\label{q444:finite-coupling-fibre}
For every ordered pair $Z_1,Z_2\in\mathfrak{su}(3)$ the complete fibre of
the four-residual zero set is as follows.
If $[Z_1,Z_2]=0$, it is empty. Otherwise form the explicitly defined
polynomials above. If $Q_7$ is identically zero, the fibre is empty.
If $Q_7$ is not identically zero, then $p_Z$ is not identically zero,
$\deg p_Z=28$, and the complete set of admissible couplings is exactly
\begin{equation}
 \left\{t\in(0,t_*]:Q_7(t)>0,\ p_Z(t)=0\right\}.
 \label{q444:finite-coupling-root-set}
\end{equation}
This set has at most fourteen elements. Equivalently, the projection
$\mathcal Z\longrightarrow\mathfrak{su}(3)^2$ that forgets $g^2$ has
fibres of cardinality at most fourteen, including empty fibres.

For each root in \eqref{q444:finite-coupling-root-set} and each
$\varepsilon\in\{+1,-1\}$, the complete original finite-system solution is
\begin{equation}
 g^2=t,\qquad
 \omega=\varepsilon\frac{\sqrt\alpha\,D_7(t)}{\sqrt{Q_7(t)}},\qquad
 C=\varepsilon\frac{\sqrt\alpha\,N_7(t)}{\sqrt{Q_7(t)}},\qquad
 (Z_1,Z_2)=(Z_1,Z_2).
 \label{q444:finite-coupling-reconstruction}
\end{equation}
Thus every fibre of the five-equation solution set over the fixed ordered
spatial pair has at most twenty-eight elements, counting both frequency
signs. No root with $Q_7(t)=0$ is included.
\end{theorem}
\begin{proof}
A zero spatial bracket is impossible by
Lemma~\ref{q444:complex-gauss-inverse}. For a nonzero bracket,
\eqref{q444:finite-coupling-T} gives
$\|Tq\|^2=Q_7(t)/D_7(t)^2$ for every $t>0$.
Thus the exact original domain condition $Tq\ne0$ is precisely
$Q_7(t)>0$. If $Q_7$ is identically zero, the domain has no point over
this pair. On the remaining domain the unchanged frequency is
\[
 \tau=\frac{\alpha D_7(t)^2}{Q_7(t)}.
\]
Substitution in all four full residuals gives, entry by entry,
\begin{equation}
 (R_H,R_q,R_1,R_2)
   =\frac1{Q_7(t)}(P_H(t),P_q(t),P_1(t),P_2(t)).
 \label{q444:finite-coupling-exact-residual-map}
\end{equation}
Every denominator here is strictly positive on the asserted domain.
For real $t$, each summand in \eqref{q444:finite-coupling-one-polynomial}
is a squared norm in its stated real inner-product space. The Hermitian
and anti-Hermitian assertions follow directly from the commutators in
\eqref{q444:finite-coupling-polynomial-residuals}; they do not require a
residual to vanish. Therefore $p_Z(t)\ge0$ for every real $t$, and its
vanishing is equivalent to vanishing of every full polynomial residual.
This proves the exact correspondence of roots and full residual zeros,
before the bound on the number of roots is established.

To prove the degree bounds, write $N_7=V_7+s_7J$, with
$\deg V_7\le6$ and $\deg s_7\le7$. Since $[J,H]=0$ and $Jq=0$,
\[
 [N_7,[N_7,H]]=[V_7,[V_7,H]]+s_7[J,[V_7,H]],
 \qquad
 N_7^2q=V_7^2q+s_7JV_7q.
\]
Their degrees are at most thirteen. Hence $\deg P_H,\deg P_q\le13$.
In the two spatial residuals the terms $D_7^2Z_j$,
$D_7[N_7,Z_k]$, and $[N_7,[N_7,Z_j]]$ have degree at most fourteen;
the terms involving $Q_7\mathcal K$ and $tQ_7\mathcal I$ have degree
at most twelve and thirteen. Thus $\deg P_1,\deg P_2\le14$ and
$\deg p_Z\le28$.

Suppose now that $Q_7$ is not identically zero but $p_Z$ is identically
zero. A nonzero univariate polynomial has only finitely many real
roots, so there exists $t>t_*$ with $Q_7(t)\ne0$. Nonnegativity of
$Q_7$ gives $Q_7(t)>0$. The established exact residual correspondence
then makes this $t$ a full zero; the two-sign inverse of
Theorem~\ref{q444:complex-exact-zero-set} reconstructs a solution of all
five original equations. This contradicts
Proposition~\ref{q444:finite-coupling-bound}. Consequently $p_Z$ is a
nonzero polynomial. Every real root of a real polynomial that is
nonnegative on all of $\mathbb R$ has even multiplicity: at a root
$t_0$, factor it as $(t-t_0)^m u(t)$ with $u(t_0)\ne0$; continuity
makes $u$ have a constant nonzero sign near $t_0$, and an odd $m$ would
give opposite signs on the two sides. Thus a nonzero polynomial of
degree at most twenty-eight of this type has at most fourteen distinct
real roots. Imposing the explicit interval and denominator conditions
can only decrease that number.

The stronger equality $\deg p_Z=28$ follows from the independent
leading-coefficient proof in Proposition~\ref{q444:finite-coupling-exact-degree}.

Finally $D_7(t)>0$ identifies the sign of the displayed $\omega$ with
$\varepsilon$ and gives $C=\omega T$ exactly. The frequency square is
the unchanged $\tau$. Formula
\eqref{q444:finite-coupling-exact-residual-map} proves that all remaining
equations vanish, while the Schur inverse proves the full Gauss equation.
Conversely every original solution determines its fixed pair, its root
$t=g^2$, and its unique sign, and these formulas recover $C,\omega$.
This proves the stated bijection and both cardinality bounds.
\end{proof}

The original spacetime reconstruction is retained literally: with the
already specified $O(t_{\rm phys})$, $K$, and $\Omega$, put
\[
 \begin{aligned}
 Y_1(t_{\rm phys})&=\cos(\omega t_{\rm phys})Z_1
                         -\sin(\omega t_{\rm phys})Z_2,\\
 Y_2(t_{\rm phys})&=\sin(\omega t_{\rm phys})Z_1
                         +\cos(\omega t_{\rm phys})Z_2,\qquad Y_3=0.
 \end{aligned}
\]
\[
 \mathcal A_0=OKCK^*O^*-\Omega,\qquad
 \mathcal A_i=OKY_iK^*O^*.
\]
The prescribed scalar fields remain
$\Phi(t_{\rm phys})=O\Phi_0O^*$ and $b(t_{\rm phys})=Ob_0$, with their
original frequency $\sqrt2$. The polynomial indeterminate $t=g^2$ is
distinct from the physical time $t_{\rm phys}$.

This theorem solves the coupling fibre exactly as a bounded finite root
set while keeping every unrestricted spatial coordinate. It does not
establish that any ordered spatial pair has a nonempty fibre, or that
all such fibres are empty. Existence or emptiness of the complete
sixteen-coordinate spatial projection remains unresolved by this result.

% Insert after the full finite coupling fibre theorem and reconstruction.
\subsection{A proper algebraic exceptional set for the full spatial projection}
\label{q444:generic-projection-section}

We now retain all sixteen spatial coordinates simultaneously. The new
statement is about the entire spatial projection, not only a selected
block class or one value of the coupling. The five original equations and
their prescribed rotating ansatz remain unchanged.

Use the exact fixed matrix $U$ of
\eqref{q444:complex-eight-coordinates}. For each $j=1,2$, write
$\widetilde Z_j=U^*Z_jU$ using the eight real coordinates
$(a_j,h_j,\operatorname{Re}v_{j1},\operatorname{Im}v_{j1},
\operatorname{Re}v_{j2},\operatorname{Im}v_{j2},
\operatorname{Re}z_j,\operatorname{Im}z_j)$, with $b_j=(h_j-a_j)/2$.
This is an invertible linear coordinate map on the original spatial
carrier. Denote the resulting sixteen-coordinate vector by $\xi$.
No coordinate is held at zero in the theorem below.

\begin{lemma}[Polynomial diagonal residuals with retained denominators]
\label{q444:generic-polynomial-diagonals}
Let $\eta=\langle J,\mathcal KJ\rangle$ be the original nonnegative
polynomial in $\xi$. On $\eta>0$ define
\[
 F_\xi(t)=\eta^{16}(U^*P_H(t)U)_{11},\qquad
 G_\xi(t)=\eta^{16}(U^*P_H(t)U)_{22}.
\]
Both extend as polynomials in $(\xi,t)$ with real algebraic
coefficients, and their degree in $t$ is at most twelve.
On the stated domain the inverse formulas are exactly
$(U^*P_HU)_{11}=F_\xi/\eta^{16}$ and
$(U^*P_HU)_{22}=G_\xi/\eta^{16}$.
\end{lemma}
\begin{proof}
All real coordinate components of $\mathcal K,w,\eta,\beta,\Pi\mathscr B$
are polynomials in $\xi$ with fixed real algebraic coefficients; the
projection denominator $\langle J,J\rangle=3/4$ is a fixed nonzero
constant. Choose once and for all a fixed basis of $W$ with real
algebraic coordinates in the displayed diagonal-frame carrier.
The entries of $\mathscr M=t\mathcal I_7+2\mathcal S_7$ have
denominator dividing $\eta$, and those of $f_7$ have the same property.
The determinant expansion shows that $\eta^7D_7$ is polynomial.
Each adjugate entry is a six-by-six minor, so $\eta^7V_7$ is
polynomial as well. The full formula
\eqref{q444:finite-coupling-N} therefore makes $\eta^8N_7$
polynomial. Complex conjugation here conjugates fixed coefficients and
the displayed complex coordinate combinations, while every component of
$\xi$ and $t$ is real. Thus $\eta^{14}Q_7$ is polynomial.
The full formula for $P_H$ now proves polynomial extension after the
displayed multiplication by $\eta^{16}$. No value at $\eta=0$
is treated as an inverse Gauss solution.

To prove the sharper degree bound, work through the exact conjugation
by $U^*$, so $H_d=\operatorname{diag}(0,r,-r)$ and
$J_d=iH_d$ are both diagonal. Write $N_7=V_7+s_7J$ as before.
Since $[J,H]=0$,
\[
 [N_7,[N_7,H]]=[V_7,[V_7,H]]+s_7[J,[V_7,H]].
\]
For every matrix $Y$, $([J_d,Y])_{jj}=0$ follows by subtracting
$(J_d)_{jj}Y_{jj}$ from itself. Consequently the second term contributes
zero to each diagonal entry after conjugation. The first term has
degree at most twelve since $\deg_t V_7\le6$. The other term in
$P_H$ is $Q_7$ times a matrix independent of $t$, and
$\deg_t Q_7\le12$. Multiplication by $\eta^{16}$ does not
change the degree in $t$. This proves every assertion.
\end{proof}

Define a fixed-size polynomial matrix $\mathscr S(\xi)$ by the
linear map
\begin{equation}
 \mathbb R[t]_{<12}\oplus\mathbb R[t]_{<12}
 \longrightarrow\mathbb R[t]_{<24},\qquad
 (A,B)\longmapsto AF_\xi+BG_\xi,
 \qquad \mathscr R(\xi)=\det\mathscr S(\xi).
 \label{q444:generic-sylvester-map}
\end{equation}
Use the increasing-power bases $1,t,\ldots$ and put the $A$ summand
before the $B$ summand. This specifies the determinant sign and its
twenty-four columns, including when either polynomial loses degree.
This is the classical Sylvester resultant map; see the Stacks Project,
\href{https://stacks.math.columbia.edu/tag/00UA}{Tag 00UA}.
The argument below proves the properties needed here directly, without
replacing either polynomial by a monic one.

\begin{theorem}[Generic exclusion in all sixteen spatial coordinates]
\label{q444:generic-projection-obstruction}
The explicitly defined polynomial $\mathscr R$ is not identically
zero. The spatial projection of every full solution of
\eqref{q444:cx:eq:phi}--\eqref{q444:cx:eq:sp2} is contained in
\[
 \{\xi\in\mathbb R^{16}:\mathscr R(\xi)=0\}.
\]
This exceptional set is closed, has empty interior, and has
sixteen-dimensional Lebesgue measure zero. Thus the complement is an
open dense set of unrestricted spatial pairs, each admitting no positive
coupling and no frequency sign solving the five retained equations.
This theorem does not decide whether any point of the exceptional set
admits a solution.
\end{theorem}
\begin{proof}
Any full solution has $[Z_1,Z_2]\ne0$, hence $\eta>0$ by
Lemma~\ref{q444:finite-coupling-schur}. It also has $Q_7(t)>0$ at
the original $t=g^2$, and the exact residual identity
\eqref{q444:finite-coupling-exact-residual-map} makes the entire
$P_H(t)$ zero. Therefore $F_\xi(t)=G_\xi(t)=0$. Evaluation at
that real $t$ is a nonzero linear functional on
$\mathbb R[t]_{<24}$ (it sends $1$ to $1$), and it annihilates
the image of \eqref{q444:generic-sylvester-map}. That image cannot
be all of the codomain. The determinant is zero. This proof includes
every degree-drop case and does not use a converse resultant criterion
outside the original denominator domain.

Here is an exact witness that the determinant polynomial is nonzero.
Write
\[
 M(a,h,v_1,v_2,z)=
 \begin{pmatrix}
 ia&-\bar v_1&-\bar v_2\\
 v_1&i(h-a)/2&z\\
 v_2&-\bar z&-i(a+h)/2
 \end{pmatrix}
\]
and set, in the displayed diagonal coordinates,
\[
 \widetilde Z_1=M(1,2,1+i,-1+2i,2-i),\qquad
 \widetilde Z_2=M(-2,1,2+3i,1-i,-1+2i).
\]
The original pair is $Z_j=U\widetilde Z_jU^*$, not a replacement
of the prescribed $H,q,F$. Every one of its sixteen displayed real
spatial coordinates is nonzero. Substitution in the original definitions
of $D_7,V_7,N_7,Q_7,P_H$ gives two diagonal polynomials of degree
twelve in $\mathbb Q(\sqrt6)[t]$. Their coprimality can be verified
by the following small certificate, with all coefficient lists in
descending powers. Reduce their coefficients by the explicit map
$\sqrt6\mapsto39$ to $\mathbb F_{101}$, where $39^2=6$.
Each rational coefficient denominator is nonzero modulo $101$;
the exact coefficient file and replay test check them individually.
The two reduced polynomials $f,g$ and the Bezout coefficients $a,b$ are
\begin{align*}
 f &: (39,85,79,93,67,20,38,98,26,56,8,25,28),\\
 g &: (1,87,96,7,76,88,75,93,53,73,64,41,16),\\
 a &: (4,60,31,59,74,31,34,18,59,97,95,96),\\
 b &: (46,85,47,5,14,36,38,21,35,48,34,53).
\end{align*}
Multiplication gives $af+bg=1$ in $\mathbb F_{101}[t]$.
Both $f$ and $g$ have degree exactly twelve. If a pair
$(a_0,b_0)$ of degrees less than twelve were in the kernel of their
Sylvester map, coprimality would imply that $f$ divides $b_0$ and
$g$ divides $a_0$. The degree bounds force both to be zero. Thus that
finite-field map, between spaces of equal dimension twenty-four, has
nonzero determinant. Reduction commutes with the determinant polynomial.
It follows that the Sylvester determinant of the two original diagonal
polynomials is nonzero in $\mathbb Q(\sqrt6)$ as well. This uses
the localized coefficient ring with all its denominators inverted,
not a homomorphism from the whole field $\mathbb Q$ to a finite field.

For explicit verification at the witness, multiplication gives
\[
 [\widetilde Z_1,\widetilde Z_2]_{21}=9-\tfrac72 i\ne0,\qquad
 \eta=\tfrac34(7+15)+3(5+5)=\tfrac{93}2>0.
\]
Multiplying both diagonal polynomials by the same nonzero scalar
$\eta^{16}$ multiplies each of the twenty-four columns by
$\eta^{16}$, and hence multiplies their determinant by
$\eta^{384}\ne0$. Consequently $\mathscr R$ is nonzero at
this exact witness and cannot be the zero polynomial. The finite-field
calculation is solely a certificate of this real algebraic statement;
it does not change the physical carrier or any coupling.

As a separate certificate, write $j_0=\langle J,J\rangle=3/4$.
In $W\oplus\mathbb RJ$, without changing the length of $J$, the
full Gauss endomorphism has block matrix
\[
 L=\begin{pmatrix}
 t\mathcal I_7+2\Pi\mathcal K|_W&2w\\
 2\langle w,\cdot\rangle/j_0&2\eta/j_0
 \end{pmatrix}.
\]
Subtracting $j_0w/\eta$ times its last row from the upper block rows
preserves the determinant and makes the upper-right block zero; its
upper-left block becomes $\mathscr M(t)$. Thus
$\det L=(2\eta/j_0)D_7$ exactly.
The independent replay reconstructs both original diagonal polynomials
from that matrix modulo $101$. It obtains Sylvester determinant $40$;
after the retained $\eta^{16}$ multiplication it obtains
$40\eta^{384}=60\ne0$ in $\mathbb F_{101}$. These residues refer
to the exact basis order in \eqref{q444:generic-sylvester-map}.

For completeness, a nonzero real polynomial on $\mathbb R^n$ has a
measure-zero zero set by induction on $n$. For $n=1$ it has finitely
many roots. For $n>1$, write it as a polynomial in the last variable
whose coefficients are polynomials in the first $n-1$ variables.
Choose one nonzero coefficient. The set where all coefficients vanish
is contained in the measure-zero zero set of that coefficient, by
induction. At every other base point the last-variable polynomial has
only finitely many roots. Applying Fubini on each bounded box shows
that the zero set in that box has measure zero; countably many boxes
cover $\mathbb R^n$. Its complement is open because the polynomial
is continuous. Every nonempty open ball has positive measure, so no
ball is contained in the zero set, and the complement is dense.
This proves the assertions at $n=16$.
\end{proof}

The exact large-coefficient identity over $\mathbb Q(\sqrt6)$ and
the small finite-field identity are retained separately in
\path{PAIR_BEZOUT_CERTIFICATE.json} and
\path{PAIR_MODULAR_CERTIFICATE.json}. The replay
\path{check_pair_bezout.py} constructs the original full pair and every
Gauss and scalar numerator, then checks the identities; the integer-only
certificate replay checks the displayed finite-field products separately.
Neither a finite failed search nor two failed test couplings proves this
theorem: the proof uses an explicit nonzero polynomial vanishing on
every possible spatial projection, with its nonzeroness certified at
the literal witness.


% Shared canonical-ready fragment derived from two_cross_exceptional_exclusion.tex.
% Input immediately after q444:generic-projection-section and before the
% existing Sources, independent checks, and reproducibility subsection.
% No preamble, document environment, or global macro definitions.
% Stage the companion proof/checker/receipt at the relative evidence path below.
\subsection{An eight-dimensional exceptional spatial class with no full solution}
\label{q444:twocross:section}

\subsubsection{Literal system and the explicitly selected class}
\label{q444:twocross:literal-system}

All vector spaces in this note retain their real structures. In particular,
$\mathfrak{su}(3)$ is the eight-real-dimensional space of traceless anti-Hermitian
complex $3\times3$ matrices, with positive inner product
$\langle X,Y\rangle=-\operatorname{tr}(XY)$. Retain the original matrices
\[
H=\begin{pmatrix}0&0&0\\0&-1/2&1/(2\sqrt2)\\
0&1/(2\sqrt2)&1/2\end{pmatrix},\quad
q=\frac{e_1}{2\sqrt2},\quad
A=\sqrt2\begin{pmatrix}0&-1&0\\1&0&0\\0&0&0\end{pmatrix},
\]
and the original scalar constants $m^2=\Gamma=1$, $\kappa=6$.
For $Y_{\rm tf}=Y-\operatorname{tr}(Y)I_3/3$, define
\[
\mathcal IX=([H,[H,X]]+qq^*X+Xqq^*)_{\rm tf},\qquad
\mathcal KX=-[Z_1,[Z_1,X]]-[Z_2,[Z_2,X]],
\]
\[
F=[A,[A,H]]=-2H+\operatorname{diag}(-2,1,1).
\]
The last identity is verified by the two displayed matrix products; also
$H^2=(3/8)\operatorname{diag}(0,1,1)$ and $q^*q=1/8$.
The indeterminate/value $t=g^2$ is the original positive gauge coupling,
not physical time. The original frequency satisfies $\omega\in\mathbb R
\setminus\{0\}$. The five retained equations are
\begin{align}
[C,[C,H]]-[Z_1,[Z_1,H]]-[Z_2,[Z_2,H]]&=F,\label{q444:twocross:eq:H}\\
(C^2-Z_1^2-Z_2^2)q&=-2q,\label{q444:twocross:eq:q}\\
(t\mathcal I+2\mathcal K)C&=4\omega[Z_1,Z_2],\label{q444:twocross:eq:G}\\
-\omega^2Z_1-2\omega[C,Z_2]+[C,[C,Z_1]]
+\mathcal KZ_1+(t/2)\mathcal IZ_1&=0,\label{q444:twocross:eq:S1}\\
-\omega^2Z_2+2\omega[C,Z_1]+[C,[C,Z_2]]
+\mathcal KZ_2+(t/2)\mathcal IZ_2&=0.\label{q444:twocross:eq:S2}
\end{align}
Multiplying the original Gauss equation by its retained positive $t$
gives \eqref{q444:twocross:eq:G}; division by that same $t$ is its exact inverse.

Put $P=qq^*/(q^*q)=\operatorname{diag}(1,0,0)$ and define
\[
\mathfrak p=\{X\in\mathfrak{su}(3):PXP=(I_3-P)X(I_3-P)=0\},\qquad
\mathfrak k=\{X\in\mathfrak{su}(3):[X,P]=0\}.
\]
Thus $\mathfrak p$ exchanges the matter line and its orthogonal complement, while
$\mathfrak k$ preserves them. Every $X\in\mathfrak p$ has the unique form
\[
X(v)=\begin{pmatrix}0&-v^*\\v&0_{2\times2}\end{pmatrix},
\qquad v\in\mathbb C^2,
\]
whose inverse takes the lower first column. This is a real linear
isomorphism $\mathbb C^2\to\mathfrak p$, so $\mathfrak p\times\mathfrak p$ has real dimension eight.
Block decomposition gives the orthogonal direct sum $\mathfrak{su}(3)=\mathfrak k\oplus\mathfrak p$,
with both summands of real dimension four.

\begin{theorem}[Two-cross spatial exclusion]\label{q444:twocross:thm:exclusion}
There is no solution of \eqref{q444:twocross:eq:H}--\eqref{q444:twocross:eq:S2} with both
$Z_1,Z_2\in\mathfrak p$. No hypothesis on $C$ is imposed in this statement.
Every complex coordinate of each of the two cross vectors is allowed,
and both signs of $\omega\ne0$ and every original $t>0$ are retained.
\end{theorem}

This is an explicitly restricted spatial class, not a claim that arbitrary
spatial pairs can be put into this class. The proof below also shows that
the entire class lies in the previously defined real algebraic exceptional
zero set. Consequently the result excludes an eight-dimensional linear
subspace inside that zero set, but does not decide existence in its
remaining part.

\subsubsection{Gauss and the complete scalar equations force the temporal form}
\label{q444:twocross:temporal-form}

Block multiplication gives
$[\mathfrak p,\mathfrak p]\subset\mathfrak k$, $[\mathfrak k,\mathfrak p]\subset\mathfrak p$.
Because $H$ and $qq^*$ are block diagonal, $\mathcal I$ preserves
both summands. Direct multiplication using $H^2$ and $q^*q$ gives
\begin{equation}
\mathcal IX(v)=\left(\frac38+\frac18\right)X(v)=\frac12X(v).
\label{q444:twocross:eq:Ip}
\end{equation}
In this computation $HX(v)H=0$, the two $H^2$ products sum to
$(3/8)X(v)$, the two $qq^*$ products sum to $X(v)/8$, and every
term is traceless. If $Z_1,Z_2\in\mathfrak p$, their double commutators
therefore make $\mathcal K$ preserve both summands as well.
Cyclicity gives, for every $Y\in\mathfrak{su}(3)$,
\[
\langle Y,\mathcal KY\rangle=\sum_{j=1}^2\|[Z_j,Y]\|^2\ge0.
\]
Project \eqref{q444:twocross:eq:G} onto $\mathfrak p$. Its right side lies in $\mathfrak k$, hence
\[
\left(\frac t2\operatorname{id}_{\mathfrak p}+2\mathcal K|_{\mathfrak p}\right)C_{\mathfrak p}=0.
\]
Pairing with $C_{\mathfrak p}$ gives a sum whose first term is
$(t/2)\|C_{\mathfrak p}\|^2$ and whose other terms are nonnegative.
Since $t>0$, $C_{\mathfrak p}=0$. Thus Gauss forces $C\in\mathfrak k$; this was
not an additional restriction on the temporal unknown.

Here are all frame changes, together with their inverses. Put
\[
x=\frac12,\quad d_0=\frac1{2\sqrt2},\quad r=\frac{\sqrt6}{4},
\quad \delta=\frac1{2r}=\sqrt{\frac23},
\]
\[
S=\frac1{\sqrt{2r(r+x)}}
\begin{pmatrix}d_0&-(r+x)\\r+x&d_0\end{pmatrix},\qquad
U=\operatorname{diag}(1,S).
\]
The relation $r^2=x^2+d_0^2$ shows that the numerator columns have
squared length $2r(r+x)$, are orthogonal, and have determinant
$2r(r+x)>0$. Multiplication by the lower block of $H$ gives their
eigenvalues $r,-r$ in that order. Hence
\[
U\in SO(3),\qquad U^*HU=\operatorname{diag}(0,r,-r),\qquad U^*q=q.
\]
For every matrix $Y$ the change is $Y\mapsto U^*YU$, with inverse
$Y'\mapsto UY'U^*$. It is applied to $H,F,C,Z_1,Z_2$ simultaneously.
It commutes with products, brackets, adjoints and trace projection,
and fixes $q$. Thus it intertwines $\mathcal I,\mathcal K$ and each
of the five full residuals exactly. It preserves $\mathfrak p,\mathfrak k$ because
it commutes with $P$. No physical constant is rescaled.

In this first frame write the full temporal block matrix as
\[
C_d=\begin{pmatrix}
ia&0&0\\0&i(h-a)/2&z\\0&-\bar z&-i(a+h)/2
\end{pmatrix},\qquad a,h\in\mathbb R,\quad z\in\mathbb C.
\]
Write the spatial matrices as $X(v),X(w)$ and put
$M=vv^*+ww^*$ in this frame only for the following calculation.
The vector scalar residual is exactly
\[
(C_d^2-X(v)^2-X(w)^2)q+2q
=(-a^2+\operatorname{tr}M+2)q.
\]
The complete nonzero block of the adjoint scalar residual, including
its upper entry, is
\[
\begin{pmatrix}
2-2r(M_{11}-M_{22})&0&0\\
0&-4r|z|^2+2rM_{11}-1+2r&-2irhz\\
0&2irh\bar z&4r|z|^2-2rM_{22}-1-2r
\end{pmatrix}.
\]
For example, if the lower Hermitian block $B=-i(C_d)_{\rm lower}$
and $D=\operatorname{diag}(r,-r)$ are used, its lower block is
$-[B,[B,D]]+\{M,D\}-(I_2-2D)$; multiplying these two-by-two
matrices gives every entry just displayed.
The two scalar equations are therefore equivalent to
\begin{equation}
a^2=2+\operatorname{tr}M,\qquad
2|z|^2-M_{11}=1-\delta,\qquad
2|z|^2-M_{22}=1+\delta,\qquad hz=0.
\label{q444:twocross:eq:first-scalar}
\end{equation}
The upper entry is minus the sum of the two lower diagonal entries,
so no scalar entry is left unchecked. Adding the middle two equations
and comparing with the first gives $a^2=4|z|^2$. Since
$M_{22}\ge0$ and $1+\delta>0$, $z\ne0$. Consequently
\begin{equation}
h=0,\qquad a\ne0,\qquad |z|=|a|/2.
\label{q444:twocross:eq:temporal-forced}
\end{equation}

Set $e=a$, retaining either sign, and $\gamma=-2iz/e$, so
$|\gamma|=1$. Choose either square root $\zeta$ of $-\bar\gamma$ and put
\[
W=\frac{\zeta}{\sqrt2}
\begin{pmatrix}\gamma&\gamma\\1&-1\end{pmatrix},\qquad
V=\operatorname{diag}(1,W).
\]
Then $|\zeta|=1$, $W^*W=I_2$ and
$\det W=-\zeta^2\gamma=\bar\gamma\gamma=1$.
The lower Hermitian block of $C_d$ is
\[
\frac e2\begin{pmatrix}-1&\gamma\\\bar\gamma&-1\end{pmatrix}.
\]
Its products with $(\gamma,1)^{\mathsf T}$ and
$(\gamma,-1)^{\mathsf T}$ have eigenvalues $0,-e$, respectively.
Those same columns give
$W^*\operatorname{diag}(r,-r)W=r\left(\begin{smallmatrix}0&1\\1&0\end{smallmatrix}\right)$.
Thus the second exact change $Y_d\mapsto V^*Y_dV$ gives
\begin{equation}
C'=ie\operatorname{diag}(1,0,-1),\quad
H'=\begin{pmatrix}0&0&0\\0&0&r\\0&r&0\end{pmatrix},\quad q'=q,
\quad F'=-2H'+\operatorname{diag}(-2,1,1).
\label{q444:twocross:eq:second-frame}
\end{equation}
The inverse for every original matrix is, explicitly,
\begin{equation}
Y=UVY'V^*U^*.
\label{q444:twocross:eq:original-inverse}
\end{equation}
For spatial vectors the forward change is the corresponding lower
unitary action $v\mapsto W^*v$ and its inverse is $v'\mapsto Wv'$.
The unitary parameters are retained; the assertion is not that $H$ can
be changed while the other data are left fixed. As with $U$, this
simultaneous conjugation intertwines all five equations. In the remainder
of the proof primes are omitted only to ease typography; all matrices
are the exact objects in \eqref{q444:twocross:eq:second-frame} and have inverse
\eqref{q444:twocross:eq:original-inverse}.

\subsubsection{Every gauge entry in the classified frame}
\label{q444:twocross:full-gauge-entries}

Write $Z_1=X(v)$, $Z_2=X(w)$ with arbitrary $v,w\in\mathbb C^2$,
and define, afresh in the second frame,
\[
V_0=v^*v,\quad W_0=w^*w,\quad p_0=v^*w=\alpha_0+ib,
\quad M=vv^*+ww^*,\quad N=i(vw^*-wv^*).
\]
Both $M$ and $N$ are Hermitian; $M\succeq0$. Put
\[
u=M_{11},\quad v_0=M_{22},\quad m=M_{12},\quad
d_M=\det M=|v_1w_2-v_2w_1|^2\ge0.
\]
The determinant identity follows by expanding the four products in
$M_{11}M_{22}-M_{12}M_{21}$ and collecting the squared modulus.
The symbol $v_0$ here is a real diagonal entry, not the vector $v$.
All eight real components of the two vectors remain independent.

Multiplication gives the full scalar residuals in this frame:
\[
Q_H=\begin{pmatrix}
2-4r\operatorname{Re}m&0&0\\
0&2r\operatorname{Re}m-1&r(\operatorname{tr}M-e^2+2)\\
0&r(\operatorname{tr}M-e^2+2)&2r\operatorname{Re}m-1
\end{pmatrix},
\quad Q_q=(\operatorname{tr}M-e^2+2)q.
\]
Consequently scalar compatibility is exactly
\begin{equation}
e^2=2+\operatorname{tr}M,\qquad \operatorname{Re}m=\delta>0.
\label{q444:twocross:eq:second-scalar}
\end{equation}
In particular $m\ne0$. No phase of $m,v,w$ has been set to a chosen value.

Set
\[
Q=\operatorname{diag}(1,2),\qquad
E=\operatorname{diag}(2/3,-5/6).
\]
Direct multiplication, keeping $r^2=3/8$ and $q^*q=1/8$, gives
\begin{equation}
\mathcal IC=ie\operatorname{diag}(1/6,2/3,-5/6),\qquad
\mathcal IX(v)=\tfrac12X(v).
\label{q444:twocross:eq:inertia-classified}
\end{equation}
For the first formula the two contributions, respectively from $H$
and from the trace-free $q$ term, are
\[
 ie\operatorname{diag}(0,3/4,-3/4),\qquad
 ie\operatorname{diag}(1/6,-1/12,-1/12).
\]
The second formula is \eqref{q444:twocross:eq:Ip}, or follows from
$(H')^2=(3/8)\operatorname{diag}(0,1,1)$ directly.
The entire spatial bracket and the entire $\mathcal KC$ are
\[
[Z_1,Z_2]=\begin{pmatrix}-2ib&0\\0&iN\end{pmatrix},\qquad
\mathcal KC=\begin{pmatrix}2ie\operatorname{tr}(QM)&0\\
0&-ie(QM+MQ)\end{pmatrix}.
\]
For example $[C,X(v)]=X(-ieQv)$, and the two double cross products
give the displayed anticommutator. Thus the full Gauss equation is
equivalent to the upper entry and lower Hermitian block
\begin{equation}
\frac{et}{6}+4e\operatorname{tr}(QM)+8\omega b=0,\qquad
etE-2e(QM+MQ)-4\omega N=0.
\label{q444:twocross:eq:gauss-block}
\end{equation}
The upper entry is the negative lower trace, since
$\operatorname{tr}N=2b$. No off-block entries occur. The entries needed
below are therefore exact consequences of the full Gauss equation:
\begin{equation}
b=-\frac e{2\omega}\left(u+2v_0+\frac t{24}\right),\qquad
N_{12}=-\frac{3e}{2\omega}m,\qquad
N_{22}=-\frac e\omega\left(2v_0+\frac{5t}{24}\right).
\label{q444:twocross:eq:gauss-needed}
\end{equation}
All divisions use the original $\omega\ne0$.

Put
\[
D=\left(\frac t4-\omega^2\right)I_2-e^2Q^2.
\]
The complete spatial residuals are $X(s_1),X(s_2)$, respectively, where
\begin{align}
s_1&=(D+W_0I_2)v+
 \left(-\alpha_0I_2+i(2\omega eQ-3bI_2)\right)w,\label{q444:twocross:eq:s1}\\
s_2&=(D+V_0I_2)w+
 \left(-\alpha_0I_2-i(2\omega eQ-3bI_2)\right)v.\label{q444:twocross:eq:s2}
\end{align}
To verify every term, $[C,X(v)]=X(-ieQv)$ gives the two gyroscopic
terms and the term $-e^2Q^2$. The inertia gives $tI_2/4$.
For the remaining double commutators, block multiplication gives
\[
\mathcal KX(v)=X\bigl(W_0v+(\bar p_0-2p_0)w\bigr),\qquad
\mathcal KX(w)=X\bigl(V_0w+(p_0-2\bar p_0)v\bigr).
\]
Since $\bar p_0-2p_0=-\alpha_0-3ib$ and
$p_0-2\bar p_0=-\alpha_0+3ib$, these are exactly
\eqref{q444:twocross:eq:s1}--\eqref{q444:twocross:eq:s2}. Their upper rows are the negative
adjoints of their lower first columns, and all other entries are zero.
Thus $s_1=s_2=0$ is equivalent to the two full spatial equations.

We now use a consequence of those equations; no converse is needed.
Multiplying \eqref{q444:twocross:eq:s1} by $v^*$, \eqref{q444:twocross:eq:s2} by $w^*$ and adding
gives the exact matrix identity
\begin{equation}
s_1v^*+s_2w^*
=DM+d_MI_2+(4bI_2-2\omega eQ)N.
\label{q444:twocross:eq:outer-product}
\end{equation}
Indeed $\operatorname{tr}(M)M-M^2=d_MI_2$ follows by multiplying
the two-by-two matrix $M$. Its off-diagonal entries are zero, and its
diagonal entries are both $M_{11}M_{22}-M_{12}M_{21}$.
The cubic terms on the left of \eqref{q444:twocross:eq:outer-product} are
\[
W_0vv^*+V_0ww^*-\alpha_0(vw^*+wv^*)+3ib(vw^*-wv^*).
\]
Those in $\operatorname{tr}(M)M-M^2$ have $-ib$ instead of $3ib$;
their difference is $4bN$. The gyroscopic terms sum to
$-2\omega eQN$. These expansions prove \eqref{q444:twocross:eq:outer-product}
entry by entry, without a reality or orthogonality restriction on $v,w$.

\begin{proof}[Proof of Theorem~\ref{q444:twocross:thm:exclusion}]
Gauss and the scalar equations have already forced the exact temporal
form and \eqref{q444:twocross:eq:second-scalar}. In particular $e\ne0$, $m\ne0$.
For a full solution, \eqref{q444:twocross:eq:outer-product} is zero. Set
\[
\chi=\frac{3e}{2\omega},\qquad B_2=4b-4\omega e,
\quad D_2=\frac t4-\omega^2-4e^2.
\]
Its $(2,1)$ entry is
$D_2\bar m+B_2N_{21}=\bar m(D_2-\chi B_2)=0$ by Gauss.
Since $m\ne0$, it follows that $D_2=\chi B_2$.
Its $(2,2)$ entry is therefore
\[
0=d_M+B_2(\chi v_0+N_{22}).
\]
Using exactly \eqref{q444:twocross:eq:gauss-needed} gives
\[
\chi v_0+N_{22}=-\frac e{2\omega}\left(v_0+\frac{5t}{12}\right),
\qquad
B_2=-\frac{2e}{\omega}
 \left(u+2v_0+\frac t{24}+2\omega^2\right).
\]
Consequently every proposed full solution in this class would satisfy
\begin{equation}
0=d_M+\frac{e^2}{\omega^2}
\left(u+2v_0+\frac t{24}+2\omega^2\right)
\left(v_0+\frac{5t}{12}\right).
\label{q444:twocross:eq:positive-contradiction}
\end{equation}
But $u,v_0\ge0$, $d_M\ge0$, $t>0$, $\omega^2>0$ and $e^2>0$.
Both parentheses are strictly positive. The second term on the right
is strictly positive for either sign of $e$ and either sign of $\omega$,
contradicting \eqref{q444:twocross:eq:positive-contradiction}. This proves the theorem
without selecting a frequency or deleting any cross-vector phase.
\end{proof}

\begin{corollary}[A further necessary condition on every full solution]
Every solution of the unrestricted five equations satisfies
\[
\sum_{j=1}^2\left(
\|PZ_jP\|_{\rm HS}^2+
\|(I_3-P)Z_j(I_3-P)\|_{\rm HS}^2\right)>0.
\]
\end{corollary}
\begin{proof}
Each summand is a squared Hilbert--Schmidt norm and is nonnegative.
The sum vanishes if and only if both spatial matrices lie in $\mathfrak p$,
which is excluded by Theorem~\ref{q444:twocross:thm:exclusion}.
\end{proof}

The scalar-compatible part of the class is nonempty even with a nonzero
spatial bracket. This is not a full solution: in the frame
\eqref{q444:twocross:eq:second-frame}, choose the exact values
\[
v=\sqrt\delta(1,1)^{\mathsf T},\qquad w=(i,0)^{\mathsf T},
\qquad e=\sqrt{3+2\delta}.
\]
Then $\operatorname{tr}M=1+2\delta$, $M_{12}=\delta$, and
$[X(v),X(w)]_{11}=-2i\sqrt\delta\ne0$.
Thus both full scalar equations hold by their displayed residuals.
To embed this example into the literal frame, use the fixed
$W=i\left(\begin{smallmatrix}1&1\\1&-1\end{smallmatrix}\right)/\sqrt2$
and \eqref{q444:twocross:eq:original-inverse}. The choice of the component $i$ is a
specified example, not a normalization of the general class. The theorem
proves that no retained $t>0,\omega\ne0$ can complete this example to all
five equations.

\subsubsection{Exact inclusion of the selected class in the exceptional locus}
\label{q444:twocross:exceptional-inclusion-section}

We give the determinant map explicitly to specify the comparison with the
full sixteen-coordinate projection, including its degree conventions.
In this section the symbols $H,q,F$ again denote their literal original
matrices in \eqref{q444:twocross:eq:H}--\eqref{q444:twocross:eq:S2},
identical to the earlier retained data of
\eqref{q444:cx:eq:phi}--\eqref{q444:cx:eq:sp2}. For any such matrix write
$\widetilde Y=U^*YU$ when using its first, $H$-diagonal-frame value.
The inverse remains $Y=U\widetilde YU^*$. Put
\[
J=iH,\quad j_0=\langle J,J\rangle=\frac34,\quad
\mathscr W=J^\perp,\quad
\Pi X=X-\frac{\langle J,X\rangle}{j_0}J,\quad \mathscr B=[Z_1,Z_2],
\]
\[
\eta=\langle J,\mathcal KJ\rangle,\quad w_J=\Pi\mathcal KJ,
\quad\beta=\langle J,\mathscr B\rangle,\quad
\mathcal I_7=\mathcal I|_{\mathscr W},
\]
\[
\mathcal S_7=\Pi\mathcal K|_{\mathscr W}
-\frac{w_J\otimes w_J}{\eta},\qquad
f_7=\Pi\mathscr B-\frac\beta\eta w_J.
\]
Here $(w_J\otimes w_J)X=w_J\langle w_J,X\rangle$ and the original length of
$J$ is retained. On $\eta>0$ define
\[
\mathscr M_7(t)=t\mathcal I_7+2\mathcal S_7,\qquad
D_7=\det_{\mathscr W}\mathscr M_7,\qquad
V_7=4\operatorname{adj}_{\mathscr W}(\mathscr M_7)f_7,
\]
\[
N_7=V_7+\frac{2\beta D_7-\langle w_J,V_7\rangle}{\eta}J,
\qquad Q_7=\|V_7q\|^2,
\quad\alpha=\frac14+\|Z_1q\|^2+\|Z_2q\|^2,
\]
\[
P_H(t)=\alpha[N_7,[N_7,H]]
-Q_7\left(\sum_{j=1}^2[Z_j,[Z_j,H]]+F\right).
\]
Let $F_\xi=\eta^{16}(\widetilde P_H)_{11}$ and
$G_\xi=\eta^{16}(\widetilde P_H)_{22}$ denote the exact polynomials used for
the original sixteen-coordinate exceptional set, with inverse division
by $\eta^{16}$ on $\eta>0$. The matrix $P_H$ in its definition is
computed with the literal original $H,q,F$; $\widetilde P_H=U^*P_HU$
is then its exact conjugate, not a second change of already-conjugated data.
They extend polynomially in the sixteen real spatial coordinates $\xi$.
For clarity, their polynomial extension follows from the definitions:
$\eta^7D_7$, $\eta^7V_7$, $\eta^8N_7$ and $\eta^{14}Q_7$ are polynomial,
by determinant and minor expansions. Hence $\eta^{16}P_H$ is polynomial.
All other denominators are fixed nonzero constants. Their two selected
diagonal entries are real because $P_H$ is Hermitian for real $t$.
The universal degree bounds are $\deg V_7\le6$ and
$N_7=V_7+s_7J$ with $\deg s_7\le7$. Since
$[\widetilde J,\widetilde H]=0$ and every diagonal entry of
$[\widetilde J,Y]$ is zero, the diagonal entries of
$[\widetilde N_7,[\widetilde N_7,\widetilde H]]$ have degree at most
twelve; also $\deg Q_7\le12$. Thus $F_\xi,G_\xi$ have degree at most
twelve in $t$, including every specialization where their degrees drop.

The exceptional polynomial is the determinant $\mathscr R(\xi)$ of the
fixed map
\begin{equation}
\mathbb R[t]_{<12}\oplus\mathbb R[t]_{<12}\longrightarrow
\mathbb R[t]_{<24},\qquad (A_0,B_0)\longmapsto A_0F_\xi+B_0G_\xi,
\label{q444:twocross:eq:fixed-map}
\end{equation}
in increasing-power bases, with the first summand before the second.
Theorem~\ref{q444:generic-projection-obstruction} proves this polynomial
is nonzero on the ambient sixteen-dimensional carrier and that all full
solutions project to its zero set. The following inclusion is proved directly, independently
of whether any point of that zero set admits a solution.

\begin{proposition}\label{q444:twocross:prop:exceptional-inclusion}
The exact linear embedding $\mathfrak p\times\mathfrak p\hookrightarrow\mathfrak{su}(3)^2$ has image
inside $\{\mathscr R=0\}$. On every nonzero pair in this class the two
diagonal Gauss-eliminated polynomials have the common factor $D_4(t)^2$,
where
\[
D_4(t)=\det_{\mathfrak p}\left(\frac t2\operatorname{id}_{\mathfrak p}+2\mathcal K|_{\mathfrak p}\right)
\]
has degree exactly four and leading coefficient $1/16$.
\end{proposition}
\begin{proof}
The first-frame $\widetilde J=i\operatorname{diag}(0,r,-r)$ satisfies
$[\widetilde J,X(v)]=X(i\operatorname{diag}(r,-r)v)$.
The change by $U$ preserves $\mathfrak p$ and the inner product.
Thus, for the selected class,
\[
\eta=\sum_j\|[Z_j,J]\|^2
=r^2(\|Z_1\|^2+\|Z_2\|^2).
\]
It is positive except at the zero pair. This also shows why this class is
compatible with the stabilizer decomposition: $-\operatorname{ad}_J^2$
is exactly $r^2\operatorname{id}$ on $\mathfrak p$.

Put $\mathfrak k_0=\mathfrak k\cap J^\perp$. The orthogonal decomposition is
$\mathscr W=\mathfrak k_0\oplus\mathfrak p$, with dimensions three and four.
The already-proved block preservation of $\mathcal I,\mathcal K$,
and $J\in\mathfrak k$, give $w_J,f_7\in\mathfrak k_0$.
Consequently $\mathcal S_7$ preserves both summands and equals
$\mathcal K$ on $\mathfrak p$. Thus
\[
\mathscr M_7=\begin{pmatrix}\mathscr M_3&0\\0&\mathscr M_4\end{pmatrix},
\qquad \mathscr M_4=\frac t2\operatorname{id}_{\mathfrak p}+2\mathcal K|_{\mathfrak p},
\quad D_7=D_3D_4.
\]
Its four-dimensional determinant has degree four and leading coefficient
$(1/2)^4=1/16$. Moreover $\mathscr M_4$ is positive for $t>0$.
For completeness, $\mathcal I$ has kernel $\mathbb RJ$: its quadratic
form is $\operatorname{tr}[X,H]^2+2\|Xq\|^2$; a zero commutes with the
three-distinct-eigenvalue $H$ and annihilates $q$, so has diagonal
entries $0,is,-is$. The same formula proves positivity on $\mathscr W$.
Completing the square in $\langle X+sJ,\mathcal K(X+sJ)\rangle$ gives
\[
\langle X,\mathcal S_7X\rangle
+\eta\left(s+\frac{\langle w_J,X\rangle}{\eta}\right)^2\ge0,
\]
so choosing the displayed square to be zero proves
$\mathcal S_7\succeq0$. Thus $\mathscr M_7$ and its three-dimensional
block are positive for every $t>0$.

Define $V_3=4\operatorname{adj}(\mathscr M_3)f_7\in\mathfrak k_0$.
The block inverse for $t>0$ gives
\[
V_7=D_4V_3,\qquad
N_7=D_4N_3,\qquad
N_3=V_3+\frac{2\beta D_3-\langle w_J,V_3\rangle}\eta J.
\]
These are polynomial identities in $t$: the two sides agree for every
positive $t$, while a nonzero polynomial has only finitely many roots,
as follows by repeated division by its linear root factors.
Therefore $Q_7=D_4^2Q_3$, where $Q_3=\|V_3q\|^2$, and
\begin{equation}
P_H=D_4^2P_{H,3},\qquad
P_{H,3}=\alpha[N_3,[N_3,H]]
-Q_3\left(\sum_j[Z_j,[Z_j,H]]+F\right).
\label{q444:twocross:eq:factor}
\end{equation}
The three-dimensional minors give $\deg V_3\le2$, $\deg D_3\le3$.
Writing $N_3=V_3+s_3J$, one has $[J,H]=0$. In the diagonal frame,
$[J,Y]_{jj}=0$ for every matrix $Y$, by subtracting the identical
diagonal products. Hence the selected diagonal entries of
$[N_3,[N_3,H]]=[V_3,[V_3,H]]+s_3[J,[V_3,H]]$ have degree at most four.
Also $\deg Q_3\le4$. Thus at each fixed pair with $\eta>0$,
\[
F_\xi=D_4^2F_3,\qquad G_\xi=D_4^2G_3,\qquad
\deg F_3,\deg G_3\le4,
\]
where $F_3,G_3$ retain the scalar multiplier $\eta^{16}$.
If they are not both zero, $(A_0,B_0)=(G_3,-F_3)$ is a nonzero
kernel element of \eqref{q444:twocross:eq:fixed-map}, with both degrees less than
twelve. If both are zero, $(1,0)$ is such a kernel element instead.
The determinant therefore vanishes for every nonzero pair in the class,
including every degree-drop case. At the zero pair it vanishes by
continuity of the original polynomial $\mathscr R$ along this linear
subspace. This proves the entire asserted inclusion without using a
converse resultant test or interpreting $\eta=0$ as an inverse solution.
\end{proof}

\subsubsection{An exact certificate retaining nonzero residuals}
\label{q444:twocross:quantitative-section}

The exclusion has a quantitative extension. It also gives a polynomial
identity before the Gauss or spatial residuals are set to zero. Use the
entire second-frame variables above, with every entry of $v,w$ retained.
More precisely, return here to the matrices and inverse in
\eqref{q444:twocross:eq:second-frame}--\eqref{q444:twocross:eq:original-inverse},
not the literal-frame $H$ used in the preceding determinant comparison.
All of $M,N,m,u,v_0,b,s_1,s_2$ in this subsection are computed in
that same second frame as in \eqref{q444:twocross:eq:s1}--\eqref{q444:twocross:eq:s2}.
Define the full spatial residual matrix and the Hermitian Gauss components
\[
 \mathcal A=s_1v^*+s_2w^*,\qquad
 \gamma_0=et/6+4e(u+2v_0)+8\omega b,
 \qquad \mathcal G=etE-2e(QM+MQ)-4\omega N.
\]
Here $\mathcal A$ is not the original generator $A$: its calligraphic
symbol denotes exactly this two-by-two residual matrix. The full Gauss
residual is $i\operatorname{diag}(\gamma_0,\mathcal G)$, with the
second block of size two. Set $B_2=4b-4\omega e$ and
\[
 \mathcal P=d_M+\frac{e^2}{\omega^2}
       (u+2v_0+t/24+2\omega^2)(v_0+5t/12).
\]
This expression is defined for the original nonzero $\omega$; the
identity below has no such denominator after multiplication.

\begin{proposition}[Polynomial residual certificate and strict error bound]
\label{q444:twocross:quantitative}
For every $v,w\in\mathbb C^2$ and real $e,t,\omega$ in the displayed
second-frame matrices, the following polynomial identity holds:
\begin{align}
 &4\omega^2(\bar m\mathcal A_{22}-v_0\mathcal A_{21})
   +\omega B_2(\bar m\mathcal G_{22}-v_0\mathcal G_{21})\notag\\
 &\hspace{8mm}+e(v_0+5t/12)\bar m\gamma_0\notag\\
 &=4\bar m\bigl[\omega^2d_M+
      e^2(u+2v_0+t/24+2\omega^2)(v_0+5t/12)\bigr].
 \label{q444:twocross:polynomial-certificate}
\end{align}
No scalar, Gauss or spatial equation is assumed in this identity.

Now restrict to the exactly specified set of original tuples with
$Z_1,Z_2\in\mathfrak p$, $t>0$, $\omega\ne0$, for which both
scalar equations and Gauss hold. Let $R_1,R_2\in\mathfrak{su}(3)$
be their two full spatial residuals. Then
\begin{equation}
 \bigl(\|R_1\|_{\mathrm{HS}}^2+
             \|R_2\|_{\mathrm{HS}}^2\bigr)^{1/2}
 \geq
 \frac{\sqrt2\,\mathcal P}
 {\sqrt{\operatorname{tr}M\,(1+v_0^2/|m|^2)}}
 >\frac{5\sqrt2\,e^2t}
 {6\sqrt{\operatorname{tr}M\,(1+v_0^2/|m|^2)}}.
 \label{q444:twocross:residual-lower-bound}
\end{equation}
The full matrix norms are preserved by the two exact unitary frame maps.
In particular, on any such set with $\operatorname{tr}M\leq L$ and
$t\geq t_0>0$, where $L>0$, the strictly positive uniform bound is
\begin{equation}
 \bigl(\|R_1\|_{\mathrm{HS}}^2+
             \|R_2\|_{\mathrm{HS}}^2\bigr)^{1/2}
 >\frac{5\sqrt2\,t_0}
 {3\sqrt{L(1+L^2/\delta^2)}},\qquad \delta=1/(2r).
 \label{q444:twocross:uniform-bound}
\end{equation}
The complete scalar--Gauss solution set is parameterized in
Theorems~\ref{sgd:parametrization} and~\ref{sgd:literal-addon:quotient},
and Corollary~\ref{sgd:every-coupling} proves it is nonempty at every
original $t=g^2>0$. No sequence in its displayed bounded portion can
make both spatial residuals tend to zero.
\end{proposition}

\begin{proof}
The complete residual formulas above give
\[
 \mathcal A_{21}=D_2\bar m+B_2N_{21},\qquad
 \mathcal A_{22}=D_2v_0+d_M+B_2N_{22},\quad
 D_2=t/4-\omega^2-4e^2,
\]
and, without any equation imposed,
\[
 \mathcal G_{21}=-6e\bar m-4\omega N_{21},\qquad
 \mathcal G_{22}=-5et/6-8ev_0-4\omega N_{22}.
\]
Subtracting the two indicated spatial entries cancels $D_2$ exactly:
\[
 \bar m\mathcal A_{22}-v_0\mathcal A_{21}
 =\bar m d_M+B_2(\bar mN_{22}-v_0N_{21}).
\]
Multiplication of the indicated Gauss entries gives
\[
 \bar m\mathcal G_{22}-v_0\mathcal G_{21}
 =-e(2v_0+5t/6)\bar m
      -4\omega(\bar mN_{22}-v_0N_{21}).
\]
These cancel the two $N$ entries on the left of
\eqref{q444:twocross:polynomial-certificate}. The remaining factor is
\[
 \gamma_0-2\omega B_2
 =4e(u+2v_0+t/24+2\omega^2).
\]
It yields exactly its stated right side. Each displayed step is a
polynomial identity, including at $m=0$ or $\omega=0$; no inverse or
solution is assigned at those values.

On the stated scalar--Gauss zero set the preceding proof forces the
second-frame form, $e^2=2+\operatorname{tr}M$ and
$\operatorname{Re}m=\delta>0$. Hence $m\ne0$ and
$\operatorname{tr}M>0$. Setting $\gamma_0=\mathcal G=0$ in the
polynomial identity and dividing by its retained nonzero factors gives
\[
 \mathcal A_{22}-\frac{v_0}{\bar m}\mathcal A_{21}
 =\mathcal P>0.
\]
Cauchy--Schwarz on these two entries implies
\[
 \mathcal P\leq
  \sqrt{1+v_0^2/|m|^2}\,\|\mathcal A\|_{\mathrm{HS}}.
\]
Let $S=(s_1\ s_2)$ and $V=(v\ w)$ be the two column matrices.
Then $\mathcal A=SV^*$. Applying Cauchy--Schwarz to each entry and
summing proves $\|SV^*\|_{\mathrm{HS}}^2
\leq\|S\|_{\mathrm{HS}}^2\|V\|_{\mathrm{HS}}^2$.
The right factors are $\|s_1\|^2+\|s_2\|^2$ and
$\operatorname{tr}M$, respectively. Directly from the two off-diagonal
blocks of $X(s)$, $\|X(s)\|_{\mathrm{HS}}^2=2\|s\|^2$.
Thus $R_j=X(s_j)$ proves the first inequality in
\eqref{q444:twocross:residual-lower-bound}.

The exact expansion
\begin{align*}
 &(u+2v_0+t/24+2\omega^2)(v_0+5t/12)\\
 &=5t\omega^2/6
   +v_0(u+2v_0+t/24+2\omega^2)
   +(5t/12)(u+2v_0+t/24)
\end{align*}
has nonnegative remainder terms, the last strictly positive for $t>0$.
Together with $d_M\geq0$ and $e^2>0$ it gives
$\mathcal P>5e^2t/6$, proving the strict second inequality.
Finally $e^2=2+\operatorname{tr}M\geq2$,
$0\leq v_0\leq\operatorname{tr}M\leq L$, and $|m|\geq\delta$
give \eqref{q444:twocross:uniform-bound}. Since unitary conjugation
preserves the Hilbert--Schmidt norm, the same inequalities hold for
the complete spatial residuals in the original literal matter frame.
\end{proof}


\subsubsection{Scope, provenance, and exact replay}
\label{q444:twocross:scope-replay}

Theorem~\ref{q444:twocross:thm:exclusion} and Proposition~\ref{q444:twocross:prop:exceptional-inclusion}
together prove an empty admissible fibre at every point of the explicit
eight-dimensional two-cross linear subspace inside the full exceptional set.
The result does not prove global emptiness or produce a full solution.
It does not assert a result for arbitrary time-dependent gauge connections,
and it makes no quantum mass-gap assertion. The earlier cross--cross--block exclusion in
Theorem~\ref{q444:cx:thm:block} instead assumes $C,Z_1\in\mathfrak p$, $Z_2\in\mathfrak k$;
here only $Z_1,Z_2\in\mathfrak p$ were assumed and $C\in\mathfrak k$ was proved.

The retained literal system is
\eqref{q444:cx:eq:phi}--\eqref{q444:cx:eq:sp2}; its complete scalar
coordinates are given in \eqref{q444:complex-eight-coordinates} and
Proposition~\ref{q444:complex-full-scalar-equivalence}.
The ambient projection statement and determinant convention are
Theorem~\ref{q444:generic-projection-obstruction} and
\eqref{q444:generic-sylvester-map}. The complete source proof, exact
checker, and replay receipt for this additional exclusion are retained
under \path{supporting_materials/workbench/evidence/two_cross_exceptional_20260905/}.
All matrix identities needed for the theorem are derived above; no
external classification theorem, numerical root search, or unproved
literature identity is used.

The companion \path{check_two_cross_exceptional.py} checks 26 exact groups:
the original force and constants, both frame maps (the second over the
entire coordinate ring $\gamma\bar\gamma=1$), every scalar residual before
and after temporal classification, both inertia maps, the full bracket and
Gauss residual, both complete spatial residuals, the nonnegative determinant
identity, the full outer-product identity and the final positivity formula.
Every real component of $v,w$ is an independent symbol. It uses no floating
point acceptance criterion and writes no files. One serial run passed all
26 groups in $0.416$ seconds of measured checker time.
The determinant block-factor proof and all positivity/domain arguments
are analytic proofs in the preceding sections, not claims inferred from
sampled determinant values.

The additional denominator-free residual identity and the full spatial
error bound are proved in Proposition~\ref{q444:twocross:quantitative}.
The companion \path{check_residual_certificate.py} verifies its polynomial
identity and exact positive remainder. The Cauchy--Schwarz estimates and
their scalar--Gauss domain are proved in the text, not inferred from tests.

% Shared complete scalar/Gauss proof, with literal-domain quotient addendum.
% Original sources and exact transformation receipt are in
% supporting_materials/workbench/evidence/scalar_gauss_domain_20260905/.
% No preamble, global macros, or hidden mathematical input.
\subsection{The complete scalar--Gauss domain at every positive coupling}
\label{q444:sgd:section}

\subsubsection{Literal data, residuals, and coordinate maps}

All complex vector spaces below retain their underlying real structures.
The adjoint is denoted by $*$, and $[X,Y]=XY-YX$. Keep exactly
\[
 H_0=\begin{pmatrix}0&0&0\\0&-1/2&1/(2\sqrt2)\\
 0&1/(2\sqrt2)&1/2\end{pmatrix},\qquad
 q=\frac1{2\sqrt2}\begin{pmatrix}1\\0\\0\end{pmatrix},\qquad
 A_0=\sqrt2\begin{pmatrix}0&-1&0\\1&0&0\\0&0&0\end{pmatrix}.
\]
The original constants are $m^2=\Gamma=1$ and $\kappa=6$.
For $Y_{\rm tf}=Y-(\operatorname{tr} Y)I_3/3$, define
\[
 \mathcal I_HY=([H,[H,Y]]+qq^*Y+Yqq^*)_{\rm tf},\qquad
 \mathcal K_{Z_1,Z_2}Y=-\sum_{j=1}^2[Z_j,[Z_j,Y]].
\]
Direct multiplication gives
\[
 F_0=[A_0,[A_0,H_0]]=-2H_0+\operatorname{diag}(-2,1,1),\quad
 H_0^2=\frac38\operatorname{diag}(0,1,1),\quad q^*q=\frac18.
\]
The scalar--Gauss domain in question consists of
\[
 (t,\omega,C_0,Z_{1,0},Z_{2,0}),\quad t=g^2>0,\quad
 \omega\in\mathbb R\setminus\{0\},\quad C_0,Z_{j,0}\in\mathfrak{su}(3),
\]
where $P=qq^*/(q^*q)=\operatorname{diag}(1,0,0)$ and
\[
 PZ_{j,0}P=(I_3-P)Z_{j,0}(I_3-P)=0,
\]
satisfying exactly the following three full residual equations:
\begin{align}
 [C_0,[C_0,H_0]]-\sum_j[Z_{j,0},[Z_{j,0},H_0]]-F_0&=0,
 \label{sgd:literal-H}\\
 (C_0^2-Z_{1,0}^2-Z_{2,0}^2)q+2q&=0,
 \label{sgd:literal-q}\\
 (t\mathcal I_{H_0}+2\mathcal K_{Z_{1,0},Z_{2,0}})C_0
             -4\omega[Z_{1,0},Z_{2,0}]&=0.
 \label{sgd:literal-G}
\end{align}
Equation~\eqref{sgd:literal-G} is the original Gauss equation multiplied
by its retained positive $t$; division by the same $t$ is its inverse.
Neither spatial equation is imposed anywhere in the construction.

The exact maps to and from the classified frame are as follows. Retain
\[
 \begin{gathered}
 r=\frac{\sqrt6}{4},\quad \delta=\frac1{2r}=\sqrt{\frac23},\quad
 x=\frac12,\quad d_0=\frac1{2\sqrt2},\\
 S=\frac1{\sqrt{2r(r+x)}}
       \begin{pmatrix}d_0&-(r+x)\\r+x&d_0\end{pmatrix},\quad
 U=\operatorname{diag}(1,S).
 \end{gathered}
\]
The identity $r^2=x^2+d_0^2$ gives $S^*S=I_2$, $\det S=1$;
multiplication of $H_0$ by its two displayed lower columns gives
\[
 U^*H_0U=\operatorname{diag}(0,r,-r),\qquad U^*q=q.
\]
Choose any $\zeta\in\mathbb C$ with $|\zeta|=1$, and set
\[
 \gamma=-\zeta^{-2},\qquad
 W=\frac\zeta{\sqrt2}\begin{pmatrix}\gamma&\gamma\\1&-1\end{pmatrix}
  =\frac1{\sqrt2}\begin{pmatrix}-\zeta^{-1}&-\zeta^{-1}\\
                         \zeta&-\zeta\end{pmatrix},\qquad
 V=\operatorname{diag}(1,W),\quad B=UV.
\]
Here $|\gamma|=1$ and $\zeta^2=-\overline\gamma$, retaining exactly
the original phase relation. Every unit $\gamma$ occurs, with its two
possible square roots $\zeta$. The displayed entries give
\[
 W^*W=I_2,\quad \det W=1,\quad
 W^*\operatorname{diag}(r,-r)W=\begin{pmatrix}0&r\\r&0\end{pmatrix}.
\]
Thus, simultaneously for all data,
\begin{equation}
 Y'=B^*Y_0B,\qquad Y_0=BY'B^*,\qquad
 H'=\begin{pmatrix}0&0&0\\0&0&r\\0&r&0\end{pmatrix},\quad
 q'=q,\quad F'=-2H'+\operatorname{diag}(-2,1,1).
 \label{sgd:frame}
\end{equation}
These are inverse maps, with no constant or unknown rescaled. In fact
$B^*(XY)B=(B^*XB)(B^*YB)$, the same equality holds for brackets,
trace and adjoints, and $Bq=q$. Therefore $\mathcal I$ and $\mathcal K$
and every full residual in \eqref{sgd:literal-H}--\eqref{sgd:literal-G}
are conjugated exactly (the vector residual is multiplied by $B^*$).
This proves that zero residuals in the classified frame give zero residuals
in the literal frame, with the displayed inverse recovering every datum.
The block-diagonal $B$ commutes with $P$, preserving the two-cross class.

\subsubsection{Every scalar and Gauss entry in the classified frame}

For $v,w\in\mathbb C^2$ write
\[
 X(v)=\begin{pmatrix}0&-v^*\\v&0_{2\times2}\end{pmatrix},\qquad
 Z_1=X(v),\quad Z_2=X(w),\quad
 C=ie\operatorname{diag}(1,0,-1),\quad e\in\mathbb R\setminus\{0\}.
\]
The inverse of $v\mapsto X(v)$ is the lower first column; every
$X(v)$ is traceless anti-Hermitian. In this section all unprimed matrices
are these classified-frame matrices, with literal inverse~\eqref{sgd:frame}.
Define
\[
 \begin{gathered}
 M=vv^*+ww^*=\begin{pmatrix}u&m\\\bar m&v_0\end{pmatrix},\quad
 N=i(vw^*-wv^*),\\
 b=\operatorname{Im}(v^*w),\quad
 Q=\operatorname{diag}(1,2),\quad E=\operatorname{diag}(2/3,-5/6).
 \end{gathered}
\]
The subscript on the real scalar $v_0$ distinguishes it from the vector $v$.

Here are the complete residuals, rather than only contractions of them:
\begin{align}
 Q_H&=\begin{pmatrix}
 2-4r\operatorname{Re}m&0&0\\
 0&2r\operatorname{Re}m-1&r(\operatorname{tr} M-e^2+2)\\
 0&r(\operatorname{tr} M-e^2+2)&2r\operatorname{Re}m-1
 \end{pmatrix},\label{sgd:scalar-H}\\
 Q_q&=(\operatorname{tr} M-e^2+2)q,\label{sgd:scalar-q}\\
 Q_G&=i\operatorname{diag}(\gamma_0,\mathcal G),\label{sgd:Gauss-full}\\
 \gamma_0&=et/6+4e(u+2v_0)+8\omega b,\qquad
 \mathcal G=etE-2e(QM+MQ)-4\omega N.\label{sgd:Gauss-entries}
\end{align}
In \eqref{sgd:Gauss-full} the second diagonal block has size two.
To derive every displayed entry, multiplication gives
\[
 X(v)^2=\operatorname{diag}(-v^*v,-vv^*),\quad
 [C,X(v)]=X(-ieQv),\quad
 [X(v),X(w)]=\operatorname{diag}(-2ib,iN).
\]
The first two identities and $H'_{23}=H'_{32}=r$ give
\eqref{sgd:scalar-H} and \eqref{sgd:scalar-q}; all their omitted entries
are zero. The two contributions to $\mathcal I_{H'}C$ are
\[
 ie\operatorname{diag}(0,3/4,-3/4),\qquad
 ie\operatorname{diag}(1/6,-1/12,-1/12),
\]
so $\mathcal I_{H'}C=ie\operatorname{diag}(1/6,2/3,-5/6)$. Also
\[
 \mathcal KC=\operatorname{diag}\bigl(2ie\operatorname{tr}(QM),-ie(QM+MQ)\bigr).
\]
Substitution of these full blocks gives \eqref{sgd:Gauss-full}.
Finally $\operatorname{tr} N=2b$, since $\operatorname{tr}(vw^*)=w^*v=\overline{v^*w}$.
Thus $\operatorname{tr}\mathcal G=-\gamma_0$; the upper Gauss entry vanishes
whenever the full lower block does, and this implication is an identity,
not an omitted equation. The two scalar equations are equivalent to
\begin{equation}
 e^2=2+u+v_0,\qquad \operatorname{Re}m=\delta.
 \label{sgd:scalar-conditions}
\end{equation}

\subsubsection{The exact realization of the two Hermitian matrices}

\begin{lemma}\label{sgd:realization-necessity}
For every pair $v,w$, let $D=(v\ w)$ and
$J=\left(\begin{smallmatrix}0&i\\-i&0\end{smallmatrix}\right)$.
Then $M=DD^*$, $N=DJD^*$ and
\[
 M+N=(v-iw)(v-iw)^*,\qquad
 M-N=(v+iw)(v+iw)^*.
\]
If $M$ is positive definite, then
\begin{equation}
 \operatorname{tr}(M^{-1}N)=0,\qquad \det N=-\det M.
 \label{sgd:trace-det}
\end{equation}
\end{lemma}
\begin{proof}
Multiplication of the displayed two-by-two blocks gives all three
factorizations. If $M>0$, $D$ is invertible because
$z^*Mz=\|D^*z\|^2$ has zero kernel. Consequently
$M^{-1}N=D^{-*}JD^*$ is similar to $J$, whose trace is zero
and determinant is $-1$. Taking the trace and determinant proves
\eqref{sgd:trace-det}, with every factor retained.
\end{proof}

\begin{lemma}\label{sgd:rank-one-impossible}
Every classified scalar--Gauss solution has $M>0$.
\end{lemma}
\begin{proof}
The scalar equation gives $m=\delta+iy\ne0$, so $u,v_0>0$
and $M\ne0$. It is already positive semidefinite as $DD^*$.
If $M$ had rank one, the range of $N=DJD^*$ would be contained
in its one-dimensional range. A Hermitian matrix with that range
and its orthogonal complement in its kernel is a real scalar multiple
of $M$, so $N=\nu M$ for some real $\nu$.
Set $k=e/\omega\ne0$. The full Gauss equation reads
\begin{equation}
 N=kT,\qquad
 T=\frac t4 E-\frac12(QM+MQ)
  =\begin{pmatrix}t/6-u&-3m/2\\-3\bar m/2&-5t/24-2v_0\end{pmatrix}.
 \label{sgd:T}
\end{equation}
Its off-diagonal entry gives $\nu=-3k/2$ because $m\ne0$.
The $(1,1)$ entry then gives $t/6-u=-3u/2$, or $t/6=-u/2$,
which contradicts $t>0$ and $u>0$. Thus $M$ has rank two and is
positive definite.
\end{proof}

\begin{theorem}[Complete classified-frame parameterization]
\label{sgd:parametrization}
Define the open real parameter set
\[
 \mathcal D=\{(u,v_0,y)\in\mathbb R^3:
      u>0,\quad d=uv_0-\delta^2-y^2>0,\quad c=4v_0-5u>0\}.
\]
For $(u,v_0,y)\in\mathcal D$, retain
\begin{align}
 m&=\delta+iy,&
 M&=\begin{pmatrix}u&m\\\bar m&v_0\end{pmatrix},&
 t&=\frac{72d}{c},\label{sgd:parameters1}\\
 T&=\frac t4E-\frac12(QM+MQ),&
 k&=\sigma\sqrt{\frac d{-\det T}},&
 e&=\epsilon\sqrt{2+u+v_0},\label{sgd:parameters2}\\
 \omega&=e/k,&N&=kT,& \epsilon,\sigma&\in\{1,-1\}.
 \label{sgd:parameters3}
\end{align}
The denominator $-\det T$ is strictly positive. The two matrices
$A_+=M+N$ and $A_-=M-N$ are positive semidefinite of rank one,
and their $(1,1)$ entries $h_+,h_-$ are strictly positive. Define
\begin{equation}
 a_\pm=\begin{pmatrix}\sqrt{h_\pm}\\(A_\pm)_{21}/\sqrt{h_\pm}\end{pmatrix},
 \quad p_\pm\in\mathbb C,\quad |p_\pm|=1,\quad
 v=\frac{p_+a_++p_-a_-}{2},\quad
 w=\frac{i(p_+a_+-p_-a_-)}{2}.
 \label{sgd:vectors}
\end{equation}
These formulas give exactly all classified-frame solutions of
\eqref{sgd:scalar-H}--\eqref{sgd:Gauss-entries}. In fact they are a
bijection from $\mathcal D\times\{1,-1\}^2\times U(1)^2$ to
that solution set with coordinates $(t,\omega,e,v,w)$.
\end{theorem}
\begin{proof}
First consider any solution. Lemma~\ref{sgd:rank-one-impossible}
gives $M>0$ and $d>0$. Inverting its two-by-two matrix gives
\[
 M^{-1}=\frac1d\begin{pmatrix}v_0&-m\\-\bar m&u\end{pmatrix}.
\]
The trace condition in \eqref{sgd:trace-det}, divided by the retained
nonzero $k$, now gives exactly
\begin{align}
 0=d\operatorname{tr}(M^{-1}T)
 &=v_0(t/6-u)+u(-5t/24-2v_0)+3|m|^2\notag\\
 &=\frac{t(4v_0-5u)}{24}-3d.\label{sgd:trace-condition}
\end{align}
Hence $c>0$ and $t=72d/c$. The determinant condition gives
$k^2\det T=-d$, so $\det T<0$ and $k$ is precisely one of
the two choices in \eqref{sgd:parameters2}. The scalar conditions
give precisely the two signs of $e$ stated there. The original
$\omega=e/k$ is unchanged, with sign $\epsilon\sigma$.

Conversely, start with parameters in $\mathcal D$. Positivity of $M$
follows from the exact factorization
\[
 M=LL^*,\qquad
 L=\begin{pmatrix}\sqrt u&0\\\bar m/\sqrt u&\sqrt d/\sqrt u\end{pmatrix},
 \qquad \det L=\sqrt d>0.
\]
The specified positive $t$ makes \eqref{sgd:trace-condition} zero.
Put $R=L^{-1}TL^{-*}$. This is Hermitian and
$\operatorname{tr} R=\operatorname{tr}(M^{-1}T)=0$, so its full form is
\[
 R=\begin{pmatrix}a&z\\\bar z&-a\end{pmatrix},\qquad a\in\mathbb R.
\]
The nonzero $m$ makes $T_{12}=-3m/2\ne0$, hence $R\ne0$.
Thus $\det R=-a^2-|z|^2<0$ and $\det T=d\det R<0$.
This proves that both signs of $k$ are well defined. Let $J_0=kR$.
It is Hermitian, has trace zero and determinant $-1$, and multiplying
the displayed matrix shows $J_0^2=I_2$.
For either sign and every column $z_0$, one has
\[
 z_0^*(I_2\pm J_0)z_0=\frac12\|(I_2\pm J_0)z_0\|^2\ge0.
\]
Also $\det(I_2\pm J_0)=0$ and $\operatorname{tr}(I_2\pm J_0)=2$,
so these matrices have rank one. Therefore
\[
 A_\pm=M\pm N=L(I_2\pm J_0)L^*
\]
are positive semidefinite of rank one, as claimed.

Their first diagonal entries are nonnegative. Suppose $h_s=0$ for
a sign $s\in\{1,-1\}$, where $A_s=M+sN$. Its determinant is
zero, so $|(A_s)_{12}|^2=0$. But
\[
 (A_s)_{12}=m(1-3sk/2),\quad
 h_s=u+sk(t/6-u).
\]
Since $m\ne0$, this forces $sk=2/3$ and then
$h_s=u/3+t/9>0$, a contradiction. Thus $h_\pm>0$.
Using $\det A_\pm=0$ gives
$(A_\pm)_{22}=|(A_\pm)_{21}|^2/h_\pm$; consequently
$A_\pm=a_\pm a_\pm^*$ with exactly the vectors in
\eqref{sgd:vectors}.

For those vectors, $v-iw=p_+a_+$ and $v+iw=p_-a_-$.
Their two outer products are $A_+,A_-$ because the phases have unit
modulus. Adding and subtracting the identities in
Lemma~\ref{sgd:realization-necessity} proves
\[
 vv^*+ww^*=M,\qquad i(vw^*-wv^*)=N.
\]
Equation~\eqref{sgd:scalar-conditions} holds by construction, so both
full scalar residuals vanish. Substitution of $N=kT$ and
$\omega k=e$ gives, with all coefficients retained,
\[
 \mathcal G=etE-2e(QM+MQ)
        -4e\left(\frac t4E-\frac12(QM+MQ)\right)=0.
\]
The exact identity $\gamma_0=-\operatorname{tr}\mathcal G$ then checks the
remaining upper Gauss entry. Thus every entry of all three original
residuals vanishes after applying \eqref{sgd:frame}.

For bijectivity, $u,v_0,y$ are recovered from $M$, $\epsilon$ is
the sign of $e$, and $\sigma$ the sign of $e/\omega$. The two
phase factors are recovered without any zero denominator by
\[
 p_+=\frac{(v-iw)_1}{\sqrt{h_+}},\qquad
 p_-=\frac{(v+iw)_1}{\sqrt{h_-}}.
\]
Indeed a column $c$ with $cc^*=a_\pm a_\pm^*$ and $h_\pm>0$
has $|c_1|=\sqrt{h_\pm}$, and the $(2,1)$ entry forces its second
component to be the same unit multiple of $(a_\pm)_2$.
These formulas prove both surjectivity and the exact inverse, without
setting any spatial phase to a selected value.
\end{proof}

\subsubsection{The complete literal scalar--Gauss domain and its phase quotient}
\label{sgd:literal-addon:section}

The classified-frame bijection of Theorem~\ref{sgd:parametrization}
has an exact extension to the entire literal two-cross domain.
We give the forced temporal classification, the full parameter map,
and the complete fibres of that map. All spaces retain their underlying
real structures, $[X,Y]=XY-YX$, and $*$ denotes the adjoint.

\paragraph{The literal domain and the retained constants}

Retain exactly
\[
 H_0=\begin{pmatrix}0&0&0\\0&-1/2&1/(2\sqrt2)\\
 0&1/(2\sqrt2)&1/2\end{pmatrix},\qquad
 q=\frac1{2\sqrt2}\begin{pmatrix}1\\0\\0\end{pmatrix},\qquad
 A_0=\sqrt2\begin{pmatrix}0&-1&0\\1&0&0\\0&0&0\end{pmatrix},
\]
\[
 m^2=\Gamma=1,\qquad \kappa=6,\qquad
 F_0=[A_0,[A_0,H_0]]=-2H_0+\operatorname{diag}(-2,1,1).
\]
Here $\mathfrak{su}(3)$ consists of the traceless anti-Hermitian complex
$3\times3$ matrices, with positive real inner product
$\langle X,Y\rangle=-\operatorname{tr}(XY)$.
For $Y_{\rm tf}=Y-(\operatorname{tr}Y)I_3/3$, define
\[
 \mathcal I_{H_0}Y
 =([H_0,[H_0,Y]]+qq^*Y+Yqq^*)_{\rm tf},\qquad
 \mathcal K_{Z_1,Z_2}Y=-\sum_{j=1}^2[Z_j,[Z_j,Y]].
\]
Multiplication of the retained matrices gives
$H_0q=0$, $H_0^2=(3/8)\operatorname{diag}(0,1,1)$ and $q^*q=1/8$.
Consequently $\operatorname{tr}H_0^2+2q^*q=1$, so the original
scalar coefficient $m^2+\Gamma(\operatorname{tr}H_0^2+2q^*q)$ is $2$.
Also $\operatorname{adj}H_0=\operatorname{diag}(-3/8,0,0)$ gives
$6(\operatorname{adj}H_0-qq^*)_{\rm tf}
=\operatorname{diag}(-2,1,1)$, agreeing with the displayed force.

Put
\[
 P=\frac{qq^*}{q^*q}=\operatorname{diag}(1,0,0),\qquad
 \mathfrak p=\{X\in\mathfrak{su}(3):
 PXP=(I_3-P)X(I_3-P)=0\}.
\]
Let $\mathcal S_{\rm sg}$ be exactly the set of tuples
\[
 (t,\omega,C_0,Z_{1,0},Z_{2,0}),\qquad
 t=g^2>0,\quad \omega\in\mathbb R\setminus\{0\},\quad
 C_0\in\mathfrak{su}(3),\quad Z_{j,0}\in\mathfrak p,
\]
for which the following three full residuals vanish:
\begin{align}
 [C_0,[C_0,H_0]]
 -\sum_j[Z_{j,0},[Z_{j,0},H_0]]-F_0&=0,
 \label{sgd:literal-addon:H}\\
 (C_0^2-Z_{1,0}^2-Z_{2,0}^2)q+2q&=0,
 \label{sgd:literal-addon:q}\\
 (t\mathcal I_{H_0}+2\mathcal K_{Z_{1,0},Z_{2,0}})C_0
 -4\omega[Z_{1,0},Z_{2,0}]&=0.
 \label{sgd:literal-addon:G}
\end{align}
Their codomains are, respectively,
$\operatorname{Herm}_3(\mathbb C)_0$, $\mathbb C^3$, and
$\mathfrak{su}(3)$. The original Gauss residual is
\[
 (\mathcal I_{H_0}+2\mathcal K_{Z_{1,0},Z_{2,0}}/g^2)C_0
 -(4\omega/g^2)[Z_{1,0},Z_{2,0}].
\]
Multiplication by the same retained positive $t=g^2$ gives
\eqref{sgd:literal-addon:G}, and division by $t$ is its inverse.
No spatial residual is required to vanish in the definition of
$\mathcal S_{\rm sg}$.

\paragraph{Temporal classification and all inverse frame maps}

For $v\in\mathbb C^2$, define
\[
 X(v)=\begin{pmatrix}0&-v^*\\v&0_2\end{pmatrix},\qquad
 \mathfrak k=\{X\in\mathfrak{su}(3):[X,P]=0\}.
\]
Every element of $\mathfrak p$ has exactly this form, whose inverse
is the lower first column. Define the fixed constants and matrix
\[
 r=\frac{\sqrt6}{4},\qquad \delta=\frac1{2r}=\sqrt{\frac23},
 \qquad x=\frac12,\qquad d_0=\frac1{2\sqrt2},
\]
\[
 S=\frac1{\sqrt{2r(r+x)}}
 \begin{pmatrix}d_0&-(r+x)\\r+x&d_0\end{pmatrix},
 \qquad U=\operatorname{diag}(1,S).
\]
The identity $r^2=x^2+d_0^2$ shows that the numerator columns
are orthogonal, have squared length $2r(r+x)$, and have determinant
$2r(r+x)>0$. Multiplication by the lower block of $H_0$ gives
their respective eigenvalues $r,-r$. Thus
\[
 U^*U=I_3,\qquad \det U=1,\qquad
 U^*H_0U=\operatorname{diag}(0,r,-r),\qquad Uq=q.
\]

\begin{lemma}[The temporal form is forced by scalars and Gauss]
\label{sgd:literal-addon:temporal}
Every tuple in $\mathcal S_{\rm sg}$ has $C_0\in\mathfrak k$.
In the fixed first frame $C_d=U^*C_0U$, there are unique
$e\in\mathbb R\setminus\{0\}$ and $z\in\mathbb C$ with
$|z|=|e|/2$ such that
\[
 C_d=\begin{pmatrix}
 ie&0&0\\0&-ie/2&z\\0&-\bar z&-ie/2
 \end{pmatrix}.
\]
Set $\gamma=-2iz/e$. Then $|\gamma|=1$. For either square root
$\zeta^2=-\bar\gamma$, define
\[
 W(\zeta)=\frac{\zeta}{\sqrt2}
 \begin{pmatrix}\gamma&\gamma\\1&-1\end{pmatrix}
 =\frac1{\sqrt2}
 \begin{pmatrix}-\zeta^{-1}&-\zeta^{-1}\\\zeta&-\zeta\end{pmatrix},
 \qquad B(\zeta)=U\operatorname{diag}(1,W(\zeta)).
\]
The simultaneous change and its inverse are
\begin{equation}
 Y'=B(\zeta)^*Y_0B(\zeta),\qquad
 Y_0=B(\zeta)Y'B(\zeta)^*.
 \label{sgd:literal-addon:frame}
\end{equation}
They give
\[
 C'=ie\operatorname{diag}(1,0,-1),\qquad
 H'=\begin{pmatrix}0&0&0\\0&0&r\\0&r&0\end{pmatrix},\qquad
 q'=q,\qquad F'=-2H'+\operatorname{diag}(-2,1,1).
\]
These maps preserve the spatial class and intertwine all three full
residuals in \eqref{sgd:literal-addon:H}--\eqref{sgd:literal-addon:G}.
\end{lemma}

\begin{proof}
Block multiplication gives the orthogonal direct sum
$\mathfrak{su}(3)=\mathfrak k\oplus\mathfrak p$, together with
$[\mathfrak p,\mathfrak p]\subset\mathfrak k$ and
$[\mathfrak k,\mathfrak p]\subset\mathfrak p$.
Since $H_0,qq^*$ are block diagonal, $\mathcal I_{H_0}$ preserves
both summands. More explicitly, $H_0X(v)H_0=0$; the two $H_0^2$
products sum to $(3/8)X(v)$, and the two $qq^*$ products sum
to $X(v)/8$. Each term is traceless. Thus
\[
 \mathcal I_{H_0}|_{\mathfrak p}
 =\tfrac12\operatorname{id}_{\mathfrak p}.
\]
The two bracket inclusions make $\mathcal K=\mathcal K_{Z_{1,0},Z_{2,0}}$
preserve both summands as well. Cyclicity gives, for every
$Y\in\mathfrak{su}(3)$,
\[
 \langle Y,\mathcal KY\rangle
 =-\sum_j\operatorname{tr}([Z_{j,0},Y]^2)
 =\sum_j\|[Z_{j,0},Y]\|^2\ge0.
\]
The $\mathfrak p$ projection of Gauss is therefore
\[
 (\tfrac t2\operatorname{id}_{\mathfrak p}
       +2\mathcal K|_{\mathfrak p})(C_0)_{\mathfrak p}=0.
\]
Pairing with $(C_0)_{\mathfrak p}$ gives nonnegative summands,
including $(t/2)\|(C_0)_{\mathfrak p}\|^2$. Since $t>0$,
$(C_0)_{\mathfrak p}=0$, proving $C_0\in\mathfrak k$.

The arbitrary first-frame temporal block consequently has the form
\[
 C_d=\begin{pmatrix}
 ia&0&0\\0&i(h-a)/2&z\\0&-\bar z&-i(a+h)/2
 \end{pmatrix},\qquad a,h\in\mathbb R,\quad z\in\mathbb C.
\]
Write the spatial matrices in this first frame as
$X(\widetilde v),X(\widetilde w)$, and define
$\widetilde M=\widetilde v\widetilde v^*
                +\widetilde w\widetilde w^*$.
The whole vector residual is
\[
 (-a^2+\operatorname{tr}\widetilde M+2)q,
\]
and the whole adjoint residual is
\[
 \begin{pmatrix}
 2-2r(\widetilde M_{11}-\widetilde M_{22})&0&0\\
 0&-4r|z|^2+2r\widetilde M_{11}-1+2r&-2irhz\\
 0&2irh\bar z&4r|z|^2-2r\widetilde M_{22}-1-2r
 \end{pmatrix}.
\]
For completeness, if $D=\operatorname{diag}(r,-r)$, the lower
spatial contribution is $\widetilde MD+D\widetilde M$, whose
diagonal entries are $2r\widetilde M_{11},-2r\widetilde M_{22}$
and whose off-diagonal entries vanish. The temporal double
commutator has diagonal entries $-4r|z|^2,4r|z|^2$ and
off-diagonal entries $-2irhz,2irh\bar z$.
Subtracting the lower force $I_2-2D$ gives the displayed lower block.
The upper spatial contribution and force give
$2-2r(\widetilde M_{11}-\widetilde M_{22})$.
Finally $X(v)^2=\operatorname{diag}(-v^*v,-vv^*)$ gives the
displayed vector residual, with all its other components zero.

Thus both scalar equations are equivalent to
\[
 a^2=2+\operatorname{tr}\widetilde M,\qquad
 2|z|^2-\widetilde M_{11}=1-\delta,\qquad
 2|z|^2-\widetilde M_{22}=1+\delta,\qquad hz=0.
\]
The upper adjoint entry is minus the lower trace.
Adding the middle equations gives
$4|z|^2=2+\operatorname{tr}\widetilde M=a^2$.
Since $\widetilde M_{22}\ge0$ and $1+\delta>0$, $z\ne0$.
Hence $h=0$ and $a\ne0$. This proves the claimed temporal block
with $e=a$, where its entries recover $e,z$ uniquely.

Now $|\gamma|=1$ and each root $\zeta^2=-\bar\gamma$ has
$|\zeta|=1$, with $\gamma=-\zeta^{-2}$. The displayed entries give
$W^*W=I_2$ and $\det W=-\zeta^2\gamma=1$.
The lower Hermitian block of $-iC_d$ is
\[
 \frac e2\begin{pmatrix}-1&\gamma\\\bar\gamma&-1\end{pmatrix}.
\]
Its products with $(\gamma,1)^{\mathsf T}$ and
$(\gamma,-1)^{\mathsf T}$ have eigenvalues $0,-e$,
respectively. The same columns give
$W^*DW=\left(\begin{smallmatrix}0&r\\r&0\end{smallmatrix}\right)$.
This proves the stated classified matrices.

The block-diagonal unitary $B(\zeta)$ fixes $q$ and commutes with
$P$. Its conjugation therefore preserves $\mathfrak p$ and
$\mathfrak k$. Unitarity gives
$B^*(XY)B=(B^*XB)(B^*YB)$; adjoints, commutators and trace
projection satisfy the corresponding equalities.
It follows that $\mathcal I,\mathcal K$ and every full matrix
residual are intertwined by conjugation. The vector residual is
multiplied by $B^*$. Equation~\eqref{sgd:literal-addon:frame} gives
the exact inverse on every datum, with no rescaling.
\end{proof}

\paragraph{The full parameter map and the quotient bijection}

We restate the classified parameters so that the literal map is explicit.
Put
\[
 Q=\operatorname{diag}(1,2),\qquad E=\operatorname{diag}(2/3,-5/6),
\]
\[
 \mathcal D=\{(u,v_0,y)\in\mathbb R^3:
 u>0,\ d=uv_0-\delta^2-y^2>0,\ c=4v_0-5u>0\}.
\]
For $(u,v_0,y)\in\mathcal D$, independent
$\epsilon,\sigma\in\{1,-1\}$, and $p_+,p_-\in U(1)$, define
\begin{align}
 m&=\delta+iy,&
 M&=\begin{pmatrix}u&m\\\bar m&v_0\end{pmatrix},&
 t&=\frac{72d}{c},\notag\\
 T&=\frac t4E-\frac12(QM+MQ),&
 k&=\sigma\sqrt{\frac d{-\det T}},&
 e&=\epsilon\sqrt{2+u+v_0},\notag\\
 \omega&=e/k,&N&=kT,&A_\pm&=M\pm N.
 \label{sgd:literal-addon:parameters}
\end{align}
Theorem~\ref{sgd:parametrization} proves that $\det T<0$,
that $A_\pm$ are positive semidefinite of rank one, and that
$h_\pm=(A_\pm)_{11}>0$. Thus the following vectors are defined
on the whole displayed domain:
\begin{equation}
 a_\pm=\begin{pmatrix}
 \sqrt{h_\pm}\\(A_\pm)_{21}/\sqrt{h_\pm}
 \end{pmatrix},\qquad
 v=\frac{p_+a_++p_-a_-}{2},\qquad
 w=\frac{i(p_+a_+-p_-a_-)}{2}.
 \label{sgd:literal-addon:vectors}
\end{equation}
In particular $vv^*+ww^*=M$ and $i(vw^*-wv^*)=N$.
The cited theorem proves that these parameters are in bijection
with all solutions of the two full scalar equations and full Gauss
in the frame $C'=ie\operatorname{diag}(1,0,-1)$,
$H'_{23}=H'_{32}=r$, $q'=q$, and
$F'=-2H'+\operatorname{diag}(-2,1,1)$.
It also proves the exact inverse phases
\begin{equation}
 p_+=\frac{(v-iw)_1}{\sqrt{h_+}},\qquad
 p_-=\frac{(v+iw)_1}{\sqrt{h_-}}.
 \label{sgd:literal-addon:phase-inverse}
\end{equation}

Define the full raw parameter space
\[
 \mathcal P_{\rm lit}
 =\mathcal D\times\{1,-1\}^2\times U(1)^2\times U(1)
\]
with coordinates $(u,v_0,y,\epsilon,\sigma,p_+,p_-,\zeta)$.
For the last unit phase set $\gamma=-\zeta^{-2}$ and use the
matrix $B(\zeta)$ of Lemma~\ref{sgd:literal-addon:temporal}.
Define the map $\mathfrak F:\mathcal P_{\rm lit}\to\mathcal S_{\rm sg}$
by \eqref{sgd:literal-addon:parameters}--\eqref{sgd:literal-addon:vectors}
and the following entire literal tuple:
\begin{equation}
 \begin{gathered}
 \mathfrak F(u,v_0,y,\epsilon,\sigma,p_+,p_-,\zeta)
 =(t,\omega,C_0,Z_{1,0},Z_{2,0}),\\
 \begin{aligned}
 C_0&=B(\zeta)\bigl(ie\operatorname{diag}(1,0,-1)\bigr)B(\zeta)^*,\\
 Z_{1,0}&=B(\zeta)X(v)B(\zeta)^*,\\
 Z_{2,0}&=B(\zeta)X(w)B(\zeta)^*.
 \end{aligned}
 \end{gathered}
 \label{sgd:literal-addon:forward}
\end{equation}
The codomain assertion follows from the classified theorem and the
exact residual intertwining in Lemma~\ref{sgd:literal-addon:temporal}.

\begin{theorem}[Exact quotient parameterization of the entire literal domain]
\label{sgd:literal-addon:quotient}
The map $\mathfrak F$ is onto $\mathcal S_{\rm sg}$.
The map
\begin{equation}
 \iota(u,v_0,y,\epsilon,\sigma,p_+,p_-,\zeta)
 =(u,v_0,y,\epsilon,\sigma,-p_+,-p_-,-\zeta)
 \label{sgd:literal-addon:involution}
\end{equation}
is a free involution of $\mathcal P_{\rm lit}$.
Each fibre of $\mathfrak F$ is exactly one two-element
$\iota$-orbit. Consequently \eqref{sgd:literal-addon:forward}
induces a bijection of sets
\begin{equation}
 \overline{\mathfrak F}:
 \mathcal P_{\rm lit}/\langle\iota\rangle
 \ \longrightarrow\ \mathcal S_{\rm sg}.
 \label{sgd:literal-addon:bijection}
\end{equation}
The inverse of this bijection is recovered from the literal matrices
by the formulas in the proof, with both square-root representatives
retained.
\end{theorem}

\begin{proof}
We first construct a preimage of every literal tuple, recording all
inverse coordinates. Given
$(t,\omega,C_0,Z_{1,0},Z_{2,0})\in\mathcal S_{\rm sg}$, compute
\[
 C_d=U^*C_0U,\qquad
 e=-i(C_d)_{11},\qquad z=(C_d)_{23},\qquad
 \gamma=-2iz/e.
\]
Lemma~\ref{sgd:literal-addon:temporal} proves $e\in\mathbb R\setminus\{0\}$
and $|\gamma|=1$. Choose either square root
$\zeta^2=-\bar\gamma$. Define
\[
 Z'_j=B(\zeta)^*Z_{j,0}B(\zeta),\qquad
 v=\begin{pmatrix}(Z'_1)_{21}\\(Z'_1)_{31}\end{pmatrix},\qquad
 w=\begin{pmatrix}(Z'_2)_{21}\\(Z'_2)_{31}\end{pmatrix}.
\]
The lemma proves that these are classified solutions with the original
$t,\omega,e$. Compute
\[
 M=vv^*+ww^*,\qquad N=i(vw^*-wv^*),\qquad
 u=M_{11},\quad v_0=M_{22},\quad y=\operatorname{Im}M_{12},
\]
\[
 \epsilon=\operatorname{sign}e,\qquad
 \sigma=\operatorname{sign}(e/\omega).
\]
Theorem~\ref{sgd:parametrization} proves
$(u,v_0,y)\in\mathcal D$, $M_{12}=\delta+iy$, and all equations
in \eqref{sgd:literal-addon:parameters} with this unchanged
$t,\omega,e$. It proves $h_\pm=(M\pm N)_{11}>0$, so
\eqref{sgd:literal-addon:phase-inverse} recovers unique unit
$p_+,p_-$. The vector formulas recover the displayed $v,w$
and the inverse frame map then recovers all three literal matrices.
Thus \eqref{sgd:literal-addon:forward} sends the recovered raw
parameter to the original tuple. This proves surjectivity.

Next we prove that \eqref{sgd:literal-addon:involution} preserves
each fibre. It preserves the parameter domain, and applying it
twice is the identity. It has no fixed point because
$\zeta=-\zeta$ would imply $\zeta=0$, contrary to $|\zeta|=1$.
Replacing $\zeta$ by $-\zeta$ preserves $\gamma=-\zeta^{-2}$ and
gives $W(-\zeta)=-W(\zeta)$. Put
$D_0=\operatorname{diag}(1,-I_2)$. Then
 \begin{align*}
 B(-\zeta)&=B(\zeta)D_0,\\
 D_0\bigl(ie\operatorname{diag}(1,0,-1)\bigr)D_0
 &=ie\operatorname{diag}(1,0,-1),\\
 D_0X(v)D_0&=X(-v).
 \end{align*}
Replacing $(p_+,p_-)$ by $(-p_+,-p_-)$ replaces $(v,w)$ by
$(-v,-w)$ in \eqref{sgd:literal-addon:vectors}, while leaving
$M,N,t,e,k,\omega$ fixed. The preceding matrix equalities
therefore give $\mathfrak F\circ\iota=\mathfrak F$ on the entire
raw parameter domain. This proves that $\overline{\mathfrak F}$
is well defined.

To determine every fibre, suppose two raw parameters have the same
literal image. The common temporal matrix $C_d$ recovers the same
$e,z,\gamma$ by the inverse formulas above. Their frame roots
therefore satisfy either $\zeta_2=\zeta_1$ or
$\zeta_2=-\zeta_1$, since they square to the same nonzero complex
number. If $\zeta_2=\zeta_1$, the classified temporal and spatial
matrices obtained by that common inverse frame map coincide.
The classified bijection of Theorem~\ref{sgd:parametrization}
then forces equality of all other raw parameters.

If $\zeta_2=-\zeta_1$, the identity
$B(-\zeta)=B(\zeta)D_0$ implies
\[
 Z'_{j,2}=D_0Z'_{j,1}D_0,\qquad
 v_2=-v_1,\qquad w_2=-w_1.
\]
Hence $M_2=M_1$ and $N_2=N_1$, so the inverse coordinates
$u,v_0,y$ are equal. The original $e,\omega$ are equal, which
also gives equal $\epsilon,\sigma$. Consequently $h_\pm$ agree,
and \eqref{sgd:literal-addon:phase-inverse} gives
\[
 p_{+,2}=-p_{+,1},\qquad p_{-,2}=-p_{-,1}.
\]
Thus the two raw parameters differ by exactly $\iota$.
No further case is possible. Each fibre is therefore precisely
one two-element orbit. This proves injectivity of
$\overline{\mathfrak F}$ as well as its already-proved surjectivity.

Finally the inverse formulas starting from a literal tuple give
the same orbit for either root choice, by the second fibre case.
Starting from a raw parameter, applying $\mathfrak F$ and then those
inverse formulas gives its orbit, by the two fibre cases.
Starting from a literal tuple, applying the inverse formulas and
then $\overline{\mathfrak F}$ gives that literal tuple by the first
paragraph. These prove both compositions are identities on their
stated domains, completing the quotient bijection.
\end{proof}

The construction retains
$\operatorname{sign}e=\epsilon$ and
$\operatorname{sign}\omega=\epsilon\sigma$.
The involution changes neither sign. The two independent vector phases
and the frame phase remain coordinates, subject only to the proved
simultaneous-sign identification.

For explicit scope, the additional full spatial residuals would be
\begin{align*}
 R_1={}&-\omega^2Z_{1,0}-2\omega[C_0,Z_{2,0}]
 +[C_0,[C_0,Z_{1,0}]]
 +\mathcal K_{Z_{1,0},Z_{2,0}}Z_{1,0}
 +(t/2)\mathcal I_{H_0}Z_{1,0},\\
 R_2={}&-\omega^2Z_{2,0}+2\omega[C_0,Z_{1,0}]
 +[C_0,[C_0,Z_{2,0}]]
 +\mathcal K_{Z_{1,0},Z_{2,0}}Z_{2,0}
 +(t/2)\mathcal I_{H_0}Z_{2,0}.
\end{align*}
Neither $R_1=0$ nor $R_2=0$ is part of
$\mathcal S_{\rm sg}$ or used in either proof above.
The bijection therefore classifies exactly the two scalar equations
and Gauss; it does not assert a solution of these two additional
spatial equations.


\subsubsection{A short exact family with every equation checked}

\begin{proposition}\label{sgd:explicit-point}
Choose arbitrary $p_+,p_-\in U(1)$ and independent
$\epsilon,\sigma\in\{1,-1\}$. Put
\[
 M=\delta\begin{pmatrix}1&1\\1&2\end{pmatrix},\quad
 t=24\delta,\quad
 e=\epsilon\sqrt{2+3\delta},\quad
 k=\frac{2\sigma}{3\sqrt{13}},\quad
 \omega=\frac{3\sigma\sqrt{13}}2e,
\]
\[
 N=\frac{\sigma\delta}{\sqrt{13}}
        \begin{pmatrix}2&-1\\-1&-6\end{pmatrix},\qquad
 a_\pm=\sqrt\delta\begin{pmatrix}
 \sqrt{1\pm2\sigma/\sqrt{13}}\\[1mm]
 (1\mp\sigma/\sqrt{13})/\sqrt{1\pm2\sigma/\sqrt{13}}
 \end{pmatrix}.
\]
Define $v,w$ by \eqref{sgd:vectors}. For any retained unit frame
phase $\zeta$, these give original tuples by
\[
 C_0=B\bigl(ie\operatorname{diag}(1,0,-1)\bigr)B^*,\qquad
 Z_{1,0}=BX(v)B^*,\qquad Z_{2,0}=BX(w)B^*.
\]
Every such tuple satisfies \eqref{sgd:literal-H}--\eqref{sgd:literal-G}
at the positive original coupling $g^2=24\delta$.
\end{proposition}
\begin{proof}
The first diagonal factors are positive because $2/\sqrt{13}<1$.
For a sign $s$, the exact identity
\[
 (1-s\sigma/\sqrt{13})^2
  =(1+2s\sigma/\sqrt{13})(2-6s\sigma/\sqrt{13})
\]
follows by multiplying out and using $\sigma^2=1$. Hence
$a_\pm a_\pm^*=M\pm N$. The construction therefore has exactly
the displayed $M,N$ for every pair of phases. In particular
\[
 u=\delta,\quad v_0=2\delta,\quad m=\delta,\quad
 \operatorname{tr} M=3\delta,\quad d=\delta^2,\quad
 b=\tfrac12\operatorname{tr} N=-\frac{2\sigma\delta}{\sqrt{13}}.
\]
The scalar identities are $e^2=2+3\delta$ and
$2r\delta=1$; substitution in \eqref{sgd:scalar-H} and
\eqref{sgd:scalar-q} makes every entry zero.
For the full Gauss block, direct multiplication gives
\[
 QM+MQ=\delta\begin{pmatrix}2&3\\3&8\end{pmatrix},\qquad
 tE-2(QM+MQ)
   =\delta\begin{pmatrix}12&-6\\-6&-36\end{pmatrix}.
\]
The same matrix is $(4\omega/e)N$, since
$\omega/e=3\sigma\sqrt{13}/2$. Thus $\mathcal G=0$ entry
by entry, and the upper entry is explicitly
\[
 \gamma_0=4e\delta+20e\delta
       +8\left(\frac{3\sigma\sqrt{13}}2e\right)
             \left(-\frac{2\sigma\delta}{\sqrt{13}}\right)=0.
\]
Finally \eqref{sgd:frame} proves all three literal equations exactly.
As an additional check on the construction, its determinant data are
\[
 T=\frac{3\delta}{2}\begin{pmatrix}2&-1\\-1&-6\end{pmatrix},\quad
 \det T=-\frac{117}{4}\delta^2,\quad k^2=\frac4{117},
\]
agreeing with the general parameterization. Both signs of the original
frequency are present because $\operatorname{sign}\omega=\epsilon\sigma$.
\end{proof}

\begin{corollary}[Every positive original coupling occurs]
\label{sgd:every-coupling}
For every specified original $\tau=g^2>0$, the literal two-cross
scalar--Gauss domain contains tuples at $t=\tau$, with either original
frequency sign and both temporal signs. It contains continuous phase
families at each such coupling.
\end{corollary}
\begin{proof}
For $\rho>1/\sqrt2$ choose
\[
 M_\rho=\delta\begin{pmatrix}\rho&1\\1&2\rho\end{pmatrix}.
\]
Its parameters obey $u=\delta\rho>0$,
$d=\delta^2(2\rho^2-1)>0$, $c=3\delta\rho>0$ and $m=\delta$.
Thus Theorem~\ref{sgd:parametrization} gives the complete solution
family at
\[
 t=\frac{72\delta^2(2\rho^2-1)}{3\delta\rho}
   =\frac{24\delta(2\rho^2-1)}\rho.
\]
For a given $\tau>0$, set
\[
 \rho=\frac{\tau+\sqrt{\tau^2+4608\delta^2}}{96\delta}.
\]
Because $\sqrt{4608}\delta=48\sqrt2\delta$, this is strictly
greater than $1/\sqrt2$. Squaring the displayed defining equality
gives $48\delta\rho^2-\tau\rho-24\delta=0$, which is precisely
$t=\tau$. No limiting passage or change of coupling is used.
The independent choices of $\epsilon,\sigma,p_+,p_-$ supply the
stated signs and phases; the inverse in Theorem~\ref{sgd:parametrization}
ensures different phase pairs yield different classified tuples.
\end{proof}

\subsubsection{Exact scope and reproducibility}

The results concern the simultaneous zero set of both scalar equations
and Gauss. The construction never imposes either spatial equation and
does not claim a full rotating-ansatz solution. It proves that the
previously specified scalar--Gauss domain of the two-cross quantitative
residual certificate is nonempty, in fact at every original positive
coupling. This is an algebraic statement about the retained classical
residuals and makes no quantum assertion.

The companion \texttt{check\_scalar\_gauss\_domain.py} is a finite exact
symbolic replay. It checks the full scalar and Gauss matrix formulas with
eight independent real spatial coordinates, the determinant/trace
identities, the explicit factors and both signs, and the original
unitary frame maps. The rank, positivity, parameter domains, completeness
and inverse-map proofs are given in this document; they are not inferred
from sampled values or from a numerical residual tolerance.

The exact full literal parameter map and its inverse are proved in
Theorem~\ref{sgd:literal-addon:quotient}; its only identification changes
the signs of the frame phase and both vector phases together. Combining
Corollary~\ref{sgd:every-coupling} with
Proposition~\ref{q444:twocross:quantitative} makes the stated spatial-error
bound non-vacuous at every positive original coupling. This does not alter
the full-solution exclusion in Theorem~\ref{q444:twocross:thm:exclusion}.

The complete source versions, the independently checked literal-domain
addendum and the exact symbolic replay are retained under
\path{supporting_materials/workbench/evidence/scalar_gauss_domain_20260905/}.
The retained equations are those derived from the parent model in
\eqref{q444:cx:eq:phi}--\eqref{q444:cx:eq:sp2}; this proof neither changes
that model nor replaces its two spatial equations by the Gauss equation.

% Complete smooth literal-domain refinement; exact staged source and review
% are retained in evidence/s6_global_smooth_reaudit_20260905/.
\subsection{Smooth literal inverse charts and all positive-coupling fibres}
\label{sgd:smooth-addon:section}

The exact algebraic bijections of
Theorems~\ref{sgd:parametrization} and
\ref{sgd:literal-addon:quotient} admit a smooth refinement.
We prove the parameter domain's convexity, give the phase quotient
its explicit smooth coordinates, and construct smooth inverse maps
from the literal matrix entries. We then give a global product
description in which the first continuous coordinate is the unchanged
original coupling $t=g^2$. The topology of every fixed-coupling fibre
will follow from explicit inverse formulas.

\subsubsection{Retained objects and the exact parameter domain}

Keep the literal matrices and constants
\[
 H_0=\begin{pmatrix}0&0&0\\0&-1/2&1/(2\sqrt2)\\
 0&1/(2\sqrt2)&1/2\end{pmatrix},\quad
 q=\frac1{2\sqrt2}\begin{pmatrix}1\\0\\0\end{pmatrix},\quad
 A_0=\sqrt2\begin{pmatrix}0&-1&0\\1&0&0\\0&0&0\end{pmatrix},
\]
\[
 m^2=\Gamma=1,\qquad \kappa=6,\qquad
 F_0=[A_0,[A_0,H_0]]=-2H_0+\operatorname{diag}(-2,1,1),
\]
and set
\[
 r=\frac{\sqrt6}{4},\quad\delta=\frac1{2r}=\sqrt{\frac23},
 \quad x=\frac12,\quad d_0=\frac1{2\sqrt2},\quad
 U=\operatorname{diag}\left(
 1,\frac1{\sqrt{2r(r+x)}}
 \begin{pmatrix}d_0&-(r+x)\\r+x&d_0\end{pmatrix}\right).
\]
Thus $U^*U=I_3$, $\det U=1$, $Uq=q$, and
$U^*H_0U=\operatorname{diag}(0,r,-r)$, as proved in
Lemma~\ref{sgd:literal-addon:temporal}.
Write
\[
 P=\operatorname{diag}(1,0,0),\qquad
 X(v)=\begin{pmatrix}0&-v^*\\v&0_2\end{pmatrix},\qquad
 \mathfrak p=X(\mathbb C^2)\subset\mathfrak{su}(3).
\]
The original real ambient manifold for this selected spatial class is
\[
 \mathcal E=(0,\infty)\times(\mathbb R\setminus\{0\})
             \times\mathfrak{su}(3)\times\mathfrak p\times\mathfrak p.
\]
It has real dimension $1+1+8+4+4=18$.
For $Y_{\rm tf}=Y-(\operatorname{tr}Y)I_3/3$, put
\[
 \mathcal I_{H_0}Y=([H_0,[H_0,Y]]+qq^*Y+Yqq^*)_{\rm tf},
 \qquad \mathcal K_{Z_1,Z_2}Y=-\sum_{j=1}^2[Z_j,[Z_j,Y]].
\]
The subset $\mathcal S_{\rm sg}\subset\mathcal E$ is exactly the
simultaneous zero set
\begin{align*}
 [C_0,[C_0,H_0]]
     -\sum_j[Z_{j,0},[Z_{j,0},H_0]]-F_0&=0,\\
 (C_0^2-Z_{1,0}^2-Z_{2,0}^2)q+2q&=0,\\
 (t\mathcal I_{H_0}+2\mathcal K_{Z_{1,0},Z_{2,0}})C_0
                   -4\omega[Z_{1,0},Z_{2,0}]&=0.
\end{align*}
Neither additional spatial residual is included in this definition.
The topology of $\mathcal S_{\rm sg}$ initially means its subspace
topology in $\mathcal E$; its embedded smooth structure will be proved.

For real coordinates $(u,v_0,y)$ define
\[
 h=\delta^2+y^2,\qquad m=\delta+iy,\qquad
 M=\begin{pmatrix}u&m\\\bar m&v_0\end{pmatrix},\qquad
 d=uv_0-h,\qquad c=4v_0-5u,
\]
\[
 \mathcal D=\{(u,v_0,y)\in\mathbb R^3:u>0,\ d>0,\ c>0\}.
\]
Here $h$ is a scalar, distinct from the two factor entries
$h_\pm$ used below.

\begin{lemma}[Convexity of the entire continuous parameter domain]
\label{sgd:smooth-addon:convex}
The set $\mathcal D$ is a nonempty open convex subset of
$\mathbb R^3$. In particular it is connected and has real dimension three.
\end{lemma}
\begin{proof}
The inequalities are strict inequalities for continuous functions,
so $\mathcal D$ is open. For $u>0$ and $(a,b)\in\mathbb C^2$,
the exact quadratic-form identity is
\[
 \begin{pmatrix}\bar a&\bar b\end{pmatrix}
 M\begin{pmatrix}a\\b\end{pmatrix}
 =u\left|a+\frac m u b\right|^2+\frac d u|b|^2.
\]
It proves $M>0$ when $u>0,d>0$. Conversely, $M>0$ gives
$u>0$ by the first coordinate vector, and evaluation at
$(a,b)=(-m/u,1)$ gives $d/u>0$.
Thus the conditions $u>0,d>0$ are exactly positive definiteness
of the Hermitian matrix in the affine slice $\operatorname{Re}M_{12}=\delta$.

If $M_0,M_1$ are positive definite and $0\le\lambda\le1$,
then for every nonzero $z\in\mathbb C^2$,
\[
 z^*((1-\lambda)M_0+\lambda M_1)z
 =(1-\lambda)z^*M_0z+\lambda z^*M_1z>0.
\]
The affine slice is preserved by this combination, and the additional
linear inequality $4v_0-5u>0$ is preserved as well.
The coordinate map $(u,v_0,y)\mapsto M$ is affine and injective,
so these facts prove convexity in the original real coordinates.
Finally $(u,v_0,y)=(\delta,2\delta,0)$ has
$u=\delta>0$, $d=\delta^2>0$, and $c=3\delta>0$.
It supplies a point and, by openness, a three-dimensional neighbourhood.
Straight segments inside a convex set join any two points,
proving connectedness.
\end{proof}

\subsubsection{The phase quotient with every original phase retained}

Write $\mathbb T=U(1)=\{z\in\mathbb C:|z|=1\}$.
On the original three phases, in the order $(\zeta,p_+,p_-)$,
define
\begin{equation}
 \Psi:\mathbb T^3\longrightarrow\mathbb T^3,\qquad
 \Psi(\zeta,p_+,p_-)
       =\left(\zeta^2,\frac{p_+}{\zeta},\frac{p_-}{\zeta}\right)
       =(\eta,\alpha,\beta).
 \label{sgd:smooth-addon:phase-map}
\end{equation}

\begin{lemma}[Exact smooth phase quotient]
\label{sgd:smooth-addon:phase-quotient}
The map $\Psi$ is a surjective smooth group homomorphism with kernel
$\{(1,1,1),(-1,-1,-1)\}$.
It induces a diffeomorphism
\[
 \mathbb T^3/\{(1,1,1),(-1,-1,-1)\}\ \longrightarrow\ \mathbb T^3.
\]
The quotient smooth charts have explicit local inverses retaining
both original square-root representatives.
\end{lemma}
\begin{proof}
Products, inverses on $\mathbb T$, and the square map are smooth;
the displayed formulas commute with coordinatewise multiplication.
An element is in the kernel precisely when $\zeta^2=1$ and
$p_+=p_-=\zeta$, giving exactly the two listed elements.
Given any $(\eta,\alpha,\beta)\in\mathbb T^3$, choose either root
$\zeta^2=\eta$ and set $p_+=\zeta\alpha$, $p_-=\zeta\beta$.
This gives all preimages, so $\Psi$ is onto and each fibre is
exactly one simultaneous-sign pair.

Here are the smooth local charts, rather than only a set quotient.
Fix $\eta_0\in\mathbb T$ and one root $\zeta_0^2=\eta_0$.
On the open arc $\eta=\eta_0e^{i\theta}$ with $-\pi<\theta<\pi$,
the two lifts are exactly
\[
 \zeta_\pm=\pm\zeta_0e^{i\theta/2},\qquad
 p_{+,\pm}=\zeta_\pm\alpha,\qquad
 p_{-,\pm}=\zeta_\pm\beta.
\]
Each formula is smooth. The two root arcs are disjoint and open
in the preimage, and each maps diffeomorphically to the displayed
arc times $\mathbb T^2$, using \eqref{sgd:smooth-addon:phase-map}
as its inverse. Their simultaneous-sign interchange identifies
exactly these two charts.
The quotient topology agrees with these charts: the quotient
projection is open because the inverse image of the image of any
open set is that set united with its simultaneous-sign translate.
On overlapping arcs the quotient coordinates are the same
$(\eta,\alpha,\beta)$, so their transitions are the ordinary smooth
torus transitions. These charts prove the asserted diffeomorphism.
They do not select a global square root or discard either representative.
\end{proof}

\subsubsection{Global formulas for literal matrices and their smooth inverses}

We record all the algebraic parameters used by the smooth map.
For $(u,v_0,y)\in\mathcal D$, set
\[
 Q=\operatorname{diag}(1,2),\qquad E=\operatorname{diag}(2/3,-5/6),
 \qquad t=\frac{72d}{c},\qquad
 T=\frac t4E-\frac12(QM+MQ).
\]
For independent signs $\epsilon,\sigma\in\{1,-1\}$ put
\[
 e=\epsilon\sqrt{2+u+v_0},\qquad
 k=\sigma\sqrt{\frac d{-\det T}},\qquad
 \omega=e/k,\qquad N=kT,\qquad A_\pm=M\pm N.
\]
Theorem~\ref{sgd:parametrization} proves on this entire domain that
$-\det T>0$, that $A_\pm$ are positive semidefinite of rank one,
and that $h_\pm=(A_\pm)_{11}>0$. Consequently
\[
 a_\pm=\begin{pmatrix}\sqrt{h_\pm}\\
                 (A_\pm)_{21}/\sqrt{h_\pm}\end{pmatrix}
\]
are defined everywhere in $\mathcal D$ for both signs.

For $(\eta,\alpha,\beta)\in\mathbb T^3$ put
\[
v^\flat=\frac{\alpha a_++\beta a_-}{2},\qquad
w^\flat=\frac{i(\alpha a_+-\beta a_-)}{2}.
\]
\[
L(\eta)=\frac1{\sqrt2}\begin{pmatrix}-1&-1\\\eta&-\eta\end{pmatrix},
 \qquad B^\flat(\eta)=U\operatorname{diag}(1,L(\eta)).
\]
The superscript $\flat$ names these exact intermediate coordinates;
it does not denote a metric dual or a change of magnitude.
Multiplication gives
\[
 L(\eta)^*L(\eta)=I_2,\qquad
 \det L(\eta)=\eta,\qquad \det B^\flat(\eta)=\eta.
\]
In particular these auxiliary matrices are unitary with their
displayed determinant; they are not asserted to have determinant one.
The precise relation to the original determinant-one frame is,
for either root $\zeta^2=\eta$,
\[
 \gamma=-\zeta^{-2},\qquad
 W(\zeta)=\frac1{\sqrt2}
 \begin{pmatrix}-\zeta^{-1}&-\zeta^{-1}\\\zeta&-\zeta\end{pmatrix},
 \qquad L(\eta)=\zeta W(\zeta).
\]
If $p_+=\zeta\alpha$, $p_-=\zeta\beta$, the original vectors satisfy
$v=\zeta v^\flat$ and $w=\zeta w^\flat$.
Thus every phase and every determinant factor is retained.

Let
\[
 \mathcal X=\{1,-1\}^2\times\mathcal D\times\mathbb T^3
\]
have coordinates $(\epsilon,\sigma,u,v_0,y,\eta,\alpha,\beta)$.
Define $\mathcal F^\flat:\mathcal X\to\mathcal E$ by
\begin{equation}
 \begin{gathered}
 \mathcal F^\flat(\epsilon,\sigma,u,v_0,y,\eta,\alpha,\beta)
       =(t,\omega,C_0,Z_{1,0},Z_{2,0}),\\
 C_0=B^\flat(\eta)\bigl(ie\operatorname{diag}(1,0,-1)\bigr)
                                    B^\flat(\eta)^*,\\
 \begin{aligned}
 Z_{1,0}&=B^\flat(\eta)X(v^\flat)B^\flat(\eta)^*,\\
 Z_{2,0}&=B^\flat(\eta)X(w^\flat)B^\flat(\eta)^*.
 \end{aligned}
 \end{gathered}
 \label{sgd:smooth-addon:literal-forward}
\end{equation}

\begin{theorem}[Four smooth connected components in literal coordinates]
\label{sgd:smooth-addon:literal-diffeomorphism}
The image of $\mathcal F^\flat$ is exactly $\mathcal S_{\rm sg}$,
and $\mathcal F^\flat$ is a diffeomorphism onto an embedded
six-real-dimensional submanifold of $\mathcal E$.
The submanifold has exactly four connected components, indexed by
\[
 \epsilon=\operatorname{sign}e,\qquad
 \sigma=\operatorname{sign}(e/\omega).
\]
Each component is diffeomorphic to $\mathcal D\times\mathbb T^3$.
The smooth structure and inverse charts refer to the literal
matrix subset; no quotient of literal configurations by a physical
gauge action is being taken.
\end{theorem}
\begin{proof}
First verify exactly which literal matrices the displayed formulas give.
Let $D_\zeta=\operatorname{diag}(1,\zeta I_2)$ and let
$B(\zeta)=U\operatorname{diag}(1,W(\zeta))$ be the original frame.
Then
\begin{align*}
 B^\flat(\eta)&=B(\zeta)D_\zeta,\\
 D_\zeta\bigl(ie\operatorname{diag}(1,0,-1)\bigr)D_\zeta^*
       &=ie\operatorname{diag}(1,0,-1).
\end{align*}
\[
 D_\zeta X(v^\flat)D_\zeta^*=X(\zeta v^\flat)=X(v),
 \qquad D_\zeta X(w^\flat)D_\zeta^*=X(w).
\]
Thus \eqref{sgd:smooth-addon:literal-forward} gives exactly the same
literal tuple as the original map for each of its two phase lifts.
Lemma~\ref{sgd:smooth-addon:phase-quotient} and the already-proved
literal quotient bijection, Theorem~\ref{sgd:literal-addon:quotient},
therefore show that its image is $\mathcal S_{\rm sg}$ and that
$\mathcal F^\flat$ is bijective.

Here is its inverse directly in literal matrix entries.
Given a tuple in $\mathcal S_{\rm sg}$, compute
\begin{equation}
 C_d=U^*C_0U,\qquad
 e=-i(C_d)_{11},\qquad z=(C_d)_{23},\qquad
 \eta=-\frac{2i\bar z}{e}.
 \label{sgd:smooth-addon:literal-temporal-inverse}
\end{equation}
The temporal classification proves that $e$ is real and nonzero,
$|z|=|e|/2$, and hence $|\eta|=1$.
The sign is checked by multiplying the lower temporal block in
\eqref{sgd:smooth-addon:literal-forward}:
\[
 (C_d)_{\rm lower}
 =-\frac{ie}{2}\begin{pmatrix}1&\bar\eta\\\eta&1\end{pmatrix}.
\]
In particular $z=-ie\bar\eta/2$, so
$-2i\bar z/e=\eta$ exactly.

Define the two matrices and their entire lower first columns by
\[
 Z_j^\flat=B^\flat(\eta)^*Z_{j,0}B^\flat(\eta),\qquad
 v^\flat=\begin{pmatrix}(Z_1^\flat)_{21}\\(Z_1^\flat)_{31}\end{pmatrix},
 \qquad
 w^\flat=\begin{pmatrix}(Z_2^\flat)_{21}\\(Z_2^\flat)_{31}\end{pmatrix}.
\]
Then recover
\[
 M=v^\flat(v^\flat)^*+w^\flat(w^\flat)^*,\qquad
 N=i\bigl(v^\flat(w^\flat)^*-w^\flat(v^\flat)^*\bigr),
\]
\begin{equation}
 \begin{gathered}
 u=M_{11},\quad v_0=M_{22},\quad y=\operatorname{Im}M_{12},
 \qquad \epsilon=\operatorname{sign}e,\quad
 \sigma=\operatorname{sign}(e/\omega),\\
 h_\pm=(M\pm N)_{11},\qquad
 \alpha=\frac{(v^\flat-iw^\flat)_1}{\sqrt{h_+}},\qquad
 \beta=\frac{(v^\flat+iw^\flat)_1}{\sqrt{h_-}}.
 \end{gathered}
 \label{sgd:smooth-addon:literal-spatial-inverse}
\end{equation}
Indeed multiplication by the common unit $\zeta$ preserves both
$M$ and $N$. The original phase inverse therefore gives
$\alpha=p_+/\zeta$, $\beta=p_-/\zeta$ in these formulas, with
$h_\pm>0$, $|\alpha|=|\beta|=1$.
The algebraic inverse proves that all recovered coordinates lie in
their exact domains. The displayed formulas recover the original
parameter from \eqref{sgd:smooth-addon:literal-forward}, and the
forward map recovers the original literal tuple from them.

All forward functions are real smooth on each sign sector:
$c>0$, $d>0$, $-\det T>0$, $2+u+v_0>0$, and $h_\pm>0$
are precisely the denominators and square-root arguments appearing
above, while $k\ne0$. Matrix products and conjugates are real smooth.
In the inverse, $e\ne0$ and $h_\pm>0$ ensure smoothness of every
non-discrete formula. The signs are locally constant.
The inverse formulas in particular give a continuous inverse for
the subspace topology on $\mathcal S_{\rm sg}$.

We check embeddedness, rather than just transferring an abstract
smooth structure through a bijection.
At any one image point choose open torus arcs around
$\eta,\alpha,\beta$, with their real angle coordinates.
The inverse formulas extend as real smooth functions to an open
neighbourhood in $\mathcal E$ as follows. Keep $e\ne0$ and the
computed $h_\pm>0$ by taking a sufficiently small neighbourhood.
The same matrix formulas define the complex quantities
$\eta,\alpha,\beta$ there, even when the tuple is off the zero set.
For a unit base point $\eta_0$, on the complex neighbourhood
$\operatorname{Re}(\eta/\eta_0)>0$, its local angle is the smooth
real function
\[
 \arctan\!\left(
 \frac{\operatorname{Im}(\eta/\eta_0)}
      {\operatorname{Re}(\eta/\eta_0)}\right).
\]
Use the corresponding expressions about $\alpha_0,\beta_0$.
On the zero set these are exactly its three local torus coordinates.
Together with $u,v_0,y$, they give a real smooth six-coordinate
map from an ambient neighbourhood whose composition with
$\mathcal F^\flat$ on the parameter chart is the identity.
The derivative of that composition is the identity, so the
derivative of $\mathcal F^\flat$ has rank six at every point.

For clarity, the real inverse function theorem
\cite[\S1, pp.~4--5]{sgd:Milnor1965} now applies with
its full-rank hypothesis verified: choose six ambient linear
coordinate projections whose composite with this rank-six
parametrization has invertible derivative. Locally the image is a
smooth graph over those six coordinates.
The continuous inverse already proved prevents other parameter
patches from contributing additional image points arbitrarily near
this point outside its chart: the inverse image of the chart is
open in the subspace topology. Thus the graph is a neighbourhood
of the point in the whole subset $\mathcal S_{\rm sg}$.
This proves it is an embedded six-dimensional submanifold and that
the displayed inverse is smooth for its embedded structure.

It remains to count components exactly. By
Lemma~\ref{sgd:smooth-addon:convex}, $\mathcal D$ is path connected.
Each circle is path connected: for chosen real arguments of two
unit endpoints, interpolating those arguments gives an explicit
path on the circle. Products of these paths prove $\mathbb T^3$
is path connected. Hence each of the four sign sectors in
$\mathcal X$ is nonempty and connected.
On the literal domain the two sign functions in
\eqref{sgd:smooth-addon:literal-spatial-inverse} are continuous
with values in $\{1,-1\}$ because $e,\omega$ never vanish.
They cannot change on a connected subset, so distinct sectors
cannot belong to the same component. This proves there are
exactly four connected components.

Finally the only quotient used to construct the coordinates was
the two-element kernel of the parameter map.
Its two representatives give identical literal matrices, as the
first displayed equalities proved. Distinct literal tuples remain
distinct under the bijection. No further gauge equivalence or
spatial-pair equivalence has been imposed.
\end{proof}

\subsubsection{The original coupling is a surjective submersion}

On $\mathcal D$ the original coupling function is
\[
 t(u,v_0,y)=\frac{72(uv_0-\delta^2-y^2)}{4v_0-5u}
          =\frac{72d}{c}.
\]

\begin{theorem}[Exact differential and global coupling coordinates]
\label{sgd:smooth-addon:coupling}
The function $t:\mathcal D\to(0,\infty)$ is a smooth surjective
submersion. Its exact partial derivatives are
\begin{equation}
 \begin{gathered}
 \frac{\partial t}{\partial u}
 =\frac{72(4v_0^2-5h)}{c^2},\qquad
 \frac{\partial t}{\partial v_0}
 =\frac{72(4h-5u^2)}{c^2},\\
 \frac{\partial t}{\partial y}=-\frac{144y}{c},
 \qquad h=\delta^2+y^2.
 \end{gathered}
 \label{sgd:smooth-addon:derivatives}
\end{equation}
In fact $\partial t/\partial u>0$ throughout $\mathcal D$.

Define the fixed open set
\[
 \mathcal B=\{(v_0,y)\in\mathbb R^2:
        v_0>\tfrac{\sqrt5}{2}\sqrt{\delta^2+y^2}\}.
\]
The map
\begin{equation}
 \Theta:\mathcal D\longrightarrow(0,\infty)\times\mathcal B,
 \qquad \Theta(u,v_0,y)=(t(u,v_0,y),v_0,y)
 \label{sgd:smooth-addon:trivialization}
\end{equation}
is a diffeomorphism, with exact inverse
\begin{equation}
 u=\frac{4\tau v_0+72(\delta^2+y^2)}{72v_0+5\tau},
 \qquad (\tau,v_0,y)\in(0,\infty)\times\mathcal B.
 \label{sgd:smooth-addon:coupling-inverse}
\end{equation}
Moreover $\mathcal B$ is diffeomorphic to $\mathbb R^2$ by the
explicit inverse pair
\begin{equation}
 \begin{aligned}
 (y,s)&\longmapsto
 \left(\tfrac{\sqrt5}{2}\sqrt{\delta^2+y^2}+e^s,\ y\right),\\
 (v_0,y)&\longmapsto
 \left(y,\log\!\left(v_0-\tfrac{\sqrt5}{2}\sqrt{\delta^2+y^2}\right)\right).
 \end{aligned}
 \label{sgd:smooth-addon:B-chart}
\end{equation}
Here $\tau=t=g^2$ is the unchanged positive original coupling.
\end{theorem}

\begin{proof}
Since $c>0$, the quotient defining $t$ is smooth.
The derivatives $d_u=v_0$, $d_{v_0}=u$, $d_y=-2y$,
$c_u=-5$, $c_{v_0}=4$, $c_y=0$ give
\[
 \frac{\partial t}{\partial u}
 =\frac{72(v_0c+5d)}{c^2}
 =\frac{72(4v_0^2-5h)}{c^2},
 \qquad
 \frac{\partial t}{\partial v_0}
 =\frac{72(uc-4d)}{c^2}
 =\frac{72(4h-5u^2)}{c^2},
\]
and the last derivative in \eqref{sgd:smooth-addon:derivatives}
is $72(-2y)/c$.
On $\mathcal D$, $v_0=(d+h)/u>0$, $c>0$, and $d>0$.
Hence $v_0c+5d>0$, proving $\partial t/\partial u>0$.
The differential therefore has rank one everywhere, proving the
submersion assertion independently of any fibre-topology argument.
In particular its three partial derivatives cannot simultaneously
vanish. The proposed algebraic simultaneous-zero test also agrees:
the first two zero equations would give
$4v_0^2=5h$ and $4h=5u^2$, so positivity of $u,v_0$ would imply
$v_0=5u/4$, and then both $c=0$ and $d=0$, excluded from $\mathcal D$.

We now prove the exact global inverse and surjectivity.
For a point of $\mathcal D$, the identity
\[
 4v_0^2-5h=v_0c+5d>0
\]
and $v_0>0$ imply $(v_0,y)\in\mathcal B$.
Writing $\tau=72d/c>0$, multiply by the retained $c$ to obtain
\[
 \tau(4v_0-5u)=72(uv_0-h),
 \qquad
 (72v_0+5\tau)u=4\tau v_0+72h.
\]
Its denominator $72v_0+5\tau$ is strictly positive, giving
\eqref{sgd:smooth-addon:coupling-inverse}.

Conversely fix any $\tau>0$ and $(v_0,y)\in\mathcal B$, and use
\eqref{sgd:smooth-addon:coupling-inverse}. Here $h>0$, $v_0>0$,
and the numerator and denominator of $u$ are positive.
Exact multiplication, with no sign or factor removed, gives
\begin{align}
 c=4v_0-5u
 &=\frac{72(4v_0^2-5h)}{72v_0+5\tau}>0,\notag\\
 d=uv_0-h
 &=\frac{\tau(4v_0^2-5h)}{72v_0+5\tau}>0.
 \label{sgd:smooth-addon:inverse-domain}
\end{align}
Thus the recovered point is in $\mathcal D$, and
$72d/c=\tau$ exactly.
These calculations prove both compositions of
\eqref{sgd:smooth-addon:trivialization} and
\eqref{sgd:smooth-addon:coupling-inverse} are identities.
Both maps are smooth on their stated open domains.

The set $\mathcal B$ is nonempty; for example
$(v_0,y)=(\sqrt5\delta/2+1,0)$ belongs to it.
Thus every $\tau>0$ has a preimage, proving surjectivity.
Finally $\delta>0$ makes $\sqrt{\delta^2+y^2}$ smooth for all real
$y$, and the defining inequality of $\mathcal B$ makes the
logarithm argument in \eqref{sgd:smooth-addon:B-chart} strictly
positive. Exponential and logarithm are inverse smooth functions
on their stated real domains. Substitution into
\eqref{sgd:smooth-addon:B-chart} proves both compositions are
identities, establishing the diffeomorphism with $\mathbb R^2$.
\end{proof}

\subsubsection{Exact scope and topology of the positive-coupling level sets}

\begin{corollary}[All positive-coupling levels and their literal inverse charts]
\label{sgd:smooth-addon:levels}
In the literal embedded manifold $\mathcal S_{\rm sg}$, the
original coupling projection
\[
 \pi:\mathcal S_{\rm sg}\longrightarrow(0,\infty),\qquad
 (t,\omega,C_0,Z_{1,0},Z_{2,0})\longmapsto t
\]
is a smooth surjective submersion. There is an explicit diffeomorphism
\begin{equation}
 \mathcal S_{\rm sg}
 \ \cong\
 \{1,-1\}^2\times(0,\infty)\times\mathbb R^2\times\mathbb T^3
 \label{sgd:smooth-addon:global-product}
\end{equation}
under which $\pi$ is projection onto the displayed unchanged
$(0,\infty)$ coordinate.
For every specified original $\tau>0$, the complete fixed-coupling
level $\pi^{-1}(\tau)$ is an embedded five-real-dimensional
submanifold, with exactly four connected components.
Each component is diffeomorphic to $\mathbb R^2\times\mathbb T^3$.
These statements concern precisely the scalar--Gauss domain with
both original signs and all three phase coordinates retained.
\end{corollary}
\begin{proof}
Compose the literal diffeomorphism of
Theorem~\ref{sgd:smooth-addon:literal-diffeomorphism} with
the diffeomorphisms of
Theorem~\ref{sgd:smooth-addon:coupling}. More explicitly, given
\[
 (\epsilon,\sigma,\tau,y,s,\eta,\alpha,\beta)
 \in\{1,-1\}^2\times(0,\infty)\times\mathbb R^2\times\mathbb T^3,
\]
first put
\[
 v_0=\frac{\sqrt5}{2}\sqrt{\delta^2+y^2}+e^s,\qquad
 u=\frac{4\tau v_0+72(\delta^2+y^2)}{72v_0+5\tau}.
\]
Then apply \eqref{sgd:smooth-addon:literal-forward} with these
$u,v_0,y$ and the retained signs and phases.
Equation~\eqref{sgd:smooth-addon:inverse-domain} proves that its
coupling is exactly $\tau$ and every denominator is positive.
The inverse begins with the literal matrix formulas
\eqref{sgd:smooth-addon:literal-temporal-inverse} and
\eqref{sgd:smooth-addon:literal-spatial-inverse}, keeps
$\tau=t$ from the original tuple, and recovers
\[
 s=\log\!\left(v_0-\frac{\sqrt5}{2}\sqrt{\delta^2+y^2}\right).
\]
The preceding theorems prove that these maps are smooth, mutually
inverse, and defined on every point of their stated domains.
This proves \eqref{sgd:smooth-addon:global-product}.

In these exact product coordinates $\pi$ is the projection onto
$\tau$. Its derivative is onto the one-dimensional tangent space
of $(0,\infty)$, and every $\tau>0$ occurs.
The fixed-coordinate slice is a smooth embedded codimension-one
submanifold. Transport through the diffeomorphism makes
$\pi^{-1}(\tau)$ an embedded five-dimensional submanifold of
$\mathcal S_{\rm sg}$, and hence of the original ambient manifold
$\mathcal E$. Its inclusion into the fixed-coupling ambient
hyperplane has the same smooth structure, since that coordinate is
identically $\tau$.

The formula identifies the entire fibre, without a missing
parameter portion, with
$\{1,-1\}^2\times\mathbb R^2\times\mathbb T^3$.
The real plane is path connected by line segments; the three
circles are path connected by the explicit angle paths used
earlier. Their product is therefore connected.
The two continuous nonvanishing sign invariants prevent any
connection between distinct sign sectors. This proves both
the exact count of four components and their stated diffeomorphism
type. The fibre topology follows from these explicit inverse maps,
not merely from the fact that $\pi$ is a submersion.
\end{proof}

These manifolds are subsets of the space of ordered literal tuples.
The only phase identification was equality of two parameter
representatives producing the same tuple. No spatial equation has
been imposed, and no physical gauge quotient, frequency-sign
identification, or exchange of $Z_{1,0}$ and $Z_{2,0}$ has been made.
The construction and all level-set statements use the open domain
$t>0$; they assign no boundary fibre at $t=0$.


\subsection{Sources, independent checks, and reproducibility}

The evidence is under \path{supporting_materials/workbench/evidence/},
in subdirectory \path{quaternionic_4441_audit_20260904/}.
The supplied files are preserved unchanged in its \path{source/} directory.
The source claims used here are the complex-plus-triplet and Lorentz-slice
sections of \emph{Triality 4-4-4-1 Dynamic Completion}, and the coordinates,
symmetry, positive fibres, and Jordan-completion sections of
\emph{Quaternionic 4441 Continuation}, both dated 4 September 2026.
Their source TeX and original verification results remain distinguishable
from this audit and from the independent checker \path{verify_exact.py}
in that evidence directory.
It implements quaternionic matrix multiplication independently, expands
the full fifteen-real-variable determinant identity, checks the marked
skew transport and charge conjugation, and verifies the positive-cone
counterexample, retained Jordan coefficient, group composition, actual
scalar acceleration, and electric sign. These checks supplement the proofs
above. The separate \path{verify_invariant_sectors.py} expands the
trace-free adjugate and the full complex tangency identity, and checks the
retained nonzero-sector coordinates and fourth-order escape coefficient.
The invariant-family theorem is proved by tangency and uniqueness, not
inferred from finite tests. These checks do not substitute for checking
other packet claims such as the gauge-orbit part of the vacuum theorem
or the finite-lattice quantization theorem. The separate moving-sector and
vacuum verifiers test the exact periodic equations, conserved charge,
fourteen-coordinate Hessian cubic, and metric eigenvalues. The scalar
global-minimum theorem is proved above, independently of those gauge and
quantum claims.
The portable \path{verify_gauge_lift.py} checks all fourteen potential
derivatives symbolically and all forty-six real homogeneous variational
components at a retained exact two-jet, as well as the orbit, radial,
and observation identities. This finite variational test supplements,
and does not replace, the full variational proof above.
The portable \path{exact_period_metric_checker.py} verifies the unchanged
periods, inverse real period map, full affine fibres, monodromy, rank-loss
constants, and a nonreal rational positive-definite example.
The smoothness and operator inequalities are proved above rather than
inferred from finite tests.
The further portable checks
\path{exact_hermitian_metric_spectrum_checker.py},
\path{verify_vacuum_kinetics.py}, \path{verify_mass_spectrum.py}, and
\path{verify_arbitrary_gauge_lift.py}
test the retained metric/cometric factors, integral quotient and spectral
identities, all current-free continuation equations, gauge-generator
Gram matrix and eigenvalues, and the complete real-connection residual
elimination. Their exact assertion counts and replay logs are in
\path{CLASSICAL_SPECTRUM_EXACT_REPLAYS.json}. Full proofs of existence,
spectral completeness, limits and exclusions are given above; finite
symbolic tests do not replace those arguments.
The further portable \path{check_complex_rotation.py} verifies the
constant-frame blocks, retained original-coordinate map, weighted
Gauss inverse, and both sign reconstructions at an exact unrestricted
test pair. \path{check_complex_constraints.py} expands the full
unrestricted scalar entries, inertia, spatial stress and polynomial
discriminant. Neither checker certifies a zero of the remaining system.
The strict inequalities, equivalences, domains and exclusion proofs are
the analytic arguments above.

The 5 September finite-fibre and projection evidence is in
\path{supporting_materials/workbench/evidence/finite_coupling_20260905/}.
The two original symbolic checkers retain the full eight-dimensional
Gauss equation and all residual codomains. The distinct
\path{check_projection_independent.py} uses only standard-library
rational-pair and finite-field arithmetic, independently reconstructs
the full eight-dimensional Gauss matrix, and reproduces the nonzero
Sylvester determinant. Its 101 evaluations verify degree-twelve
polynomial identities, not a sampled real coupling exclusion.
Complete definitions and proofs are in
Theorems~\ref{q444:finite-coupling-fibre} and
\ref{q444:generic-projection-obstruction}; no existence inside the
exceptional algebraic set, or general complex time-dependent lift,
is inferred. The original scripts and certificates are retained in
the evidence directory's \path{checkers/} subdirectories.

The separately dated scalar--Gauss continuation and independent review are
retained in
\path{supporting_materials/workbench/evidence/scalar_gauss_domain_20260905/}.
The maintained proof is \path{scalar_gauss_domain_shared.tex}; it includes
the full temporal-classification and literal-phase-quotient addendum.
The portable \path{check_scalar_gauss_domain.py} replay passes all 25
exact groups, including every scalar and Gauss entry, both sign choices,
the continuous phases, inverse frames and the prescribed-coupling formula.
Positivity, rank exclusion, completeness and the two-element fibres are
proved analytically in Theorems~\ref{sgd:parametrization}
and~\ref{sgd:literal-addon:quotient}. This continuation supersedes the
earlier absence of a scalar--Gauss nonemptiness assertion; it does not
alter the historical proof snapshot: the two spatial residuals still
cannot vanish simultaneously.
The unrestricted exceptional-set existence problem remains open, and no
quantum spectral conclusion follows from these classical residuals.

The subsequent shared source \path{scalar_gauss_smooth_shared.tex}
proves the smooth phase quotient, embedded literal inverse charts, and
the global trivialization over the unchanged positive coupling axis
(Theorems~\ref{sgd:smooth-addon:literal-diffeomorphism}
and~\ref{sgd:smooth-addon:coupling},
Corollary~\ref{sgd:smooth-addon:levels}). These analytic proofs establish
the component and fibre topology beyond the earlier set bijections;
the earlier exact replay receipt retains its original scope.

\section{Quantum confinement in the homogeneous gauge--scalar model}
\label{qh:chapter}

The object quantized here is a precisely specified interacting
matrix mechanics: three anti-Hermitian traceless gauge matrices,
one Hermitian traceless adjoint scalar, and one complex triplet.
It is obtained by spatially constant substitution in the gauge--scalar
action used in the quaternionic continuation. All of these fields
remain variables. In particular the triplet is not constrained to a
four-real-dimensional slice or to a scalar minimum curve.

The conclusions are a compact-resolvent Hamiltonian, a unique
spatially constant $SU(3)$-invariant ground state, and a strictly
positive first excitation gap at fixed positive volume and couplings.
The explicit form inequality \eqref{qh:full-form-bound} is the
quantitative step. It combines the exact $SU(3)$ root frequencies
with the full quartic scalar potential. Bosonic commutator-matrix
discreteness is a classical result; the primary attribution and
the complete proof in the retained metric are given below.
The present use is to answer a concrete quantization question for
the already specified adjoint-plus-triplet action, not to claim
a new general discreteness mechanism.

For the classical comparison the spacetime signature is $(+---)$
and the original action density is
\[
 \mathcal L=\frac1{2g^2}\operatorname{tr}F_{\mu\nu}F^{\mu\nu}
       +\frac12\operatorname{tr}(D_\mu\Phi D^\mu\Phi)
       +(D_\mu b)^*D^\mu b-V(\Phi,b),
\]
where
\[
 F_{\mu\nu}=\partial_\mu A_\nu-\partial_\nu A_\mu+[A_\mu,A_\nu],
 \quad D_\mu\Phi=\partial_\mu\Phi+[A_\mu,\Phi],
 \quad D_\mu b=\partial_\mu b+A_\mu b.
\]
The potential is written in full below. The finite-model theorem
does not require the construction of a compact complex structure
on $S^6$, and it neither verifies nor refutes that separate global
construction. The precise classical map to the torus field theory,
its noninjective map on gauge orbits, and an explicit obstruction
to one proposed quantum-subsector map are proved at the end.

% Insertion-ready proof; no canonical source is edited by this artifact.
% Required packages/environments: amsmath, amssymb, standard theorem/proof.
\subsection{The retained configuration space and action}

Fix a flat spatial torus of volume \(\mathcal V_3>0\), a gauge coupling
\(g^2>0\), and \(\hbar>0\). Retain \(m^2\geq0\), \(\Gamma>0\) and
\(\kappa\in\mathbb R\). We quantize the finite-dimensional homogeneous
mechanics, not a claimed invariant quantum sector of the full field theory.
The configuration space, with no quotient already performed, is
\[
 \mathcal Q=\mathfrak{su}(3)^3\oplus
       \operatorname{Herm}_3(\mathbb C)_0\oplus\mathbb C^3,
 \qquad \dim_{\mathbb R}\mathcal Q=24+8+6=38.
\]
Its elements are \(q=(A_1,A_2,A_3,\Phi,b)\). On \(\mathfrak{su}(3)\)
use \((X,Y)_{\mathfrak g}=-\operatorname{tr}(XY)\); the scalar metric is
exactly \(\operatorname{tr}(d\Phi^2)+2\,db^*db\).
The trace-free condition on \(\Phi\) is part of the linear state space.

For a homogeneous temporal connection \(A_0\), set
\[
 E_i=\dot A_i+[A_0,A_i],\quad
 X=\dot\Phi+[A_0,\Phi],\quad v=\dot b+A_0b,\quad
 s=\operatorname{tr}\Phi^2+2b^*b.
\]
The spatial integral of the retained \((+---)\), anti-Hermitian density is
\begin{align}
 L_{\rm hom}&=\mathcal V_3\left[
   \frac1{g^2}\sum_i\|E_i\|_{\mathfrak g}^2
      +\frac12\operatorname{tr}X^2+v^*v-U(q)\right],
       \label{qh:hq-L}\\
 U(q)&=\frac1{g^2}\sum_{i<j}\|[A_i,A_j]\|_{\mathfrak g}^2
       +\frac12\sum_i\operatorname{tr}[A_i,\Phi]^2
       +\sum_i\|A_i b\|^2+V(\Phi,b),\notag\\
 V(\Phi,b)&=\frac{m^2}{2}s+\frac{\Gamma}{4}s^2
                      -\kappa(\det\Phi-b^*\Phi b).\notag
\end{align}
Here \([A_i,\Phi]\) is Hermitian, so its displayed squared trace is
nonnegative. There is no velocity for \(A_0\).

\begin{proposition}[Canonical momenta, Hamiltonian and exact Gauss map]
\label{qh:hq-canonical}
Use the canonical one-form
\[
 \vartheta=\sum_i-\operatorname{tr}(\Pi_i\,dA_i)
       +\operatorname{tr}(\Pi_\Phi\,d\Phi)
       +2\operatorname{Re}(\pi_b^*db).
\]
The momenta and Hamiltonian are
\begin{align}
 \Pi_i&=\frac{2\mathcal V_3}{g^2}E_i,\qquad
 \Pi_\Phi=\mathcal V_3X,\qquad \pi_b=\mathcal V_3v,\notag\\
 H_{A_0}&=H_0-\mathcal M_{A_0},\label{qh:hq-H}\\
 H_0&=\frac{g^2}{4\mathcal V_3}\sum_i\|\Pi_i\|_{\mathfrak g}^2
       +\frac1{2\mathcal V_3}\operatorname{tr}\Pi_\Phi^2
       +\frac1{\mathcal V_3}\pi_b^*\pi_b+\mathcal V_3U(q).\notag
\end{align}
For \(\xi\in\mathfrak{su}(3)\), write
\[
 K_\xi(q)=([\xi,A_i]_{i=1}^3,[\xi,\Phi],\xi b),\qquad
 \mathcal M_\xi=\vartheta(K_\xi)=(\mathcal G,\xi)_{\mathfrak g}.
\]
The exact matrix-valued momentum map is
\begin{equation}\label{qh:hq-Gauss}
 \mathcal G=\sum_i[A_i,\Pi_i]-[\Phi,\Pi_\Phi]
          -(b\pi_b^*-\pi_b b^*)_{\rm tf}
 =\mathcal V_3\left(\frac2{g^2}\sum_i[A_i,E_i]-(J_0)_{\rm tf}\right),
\end{equation}
where \(J_0=[\Phi,X]+bv^*-vb^*\).
Thus \(\mathcal G=0\) is precisely the full homogeneous Gauss equation.
\end{proposition}
\begin{proof}
Differentiate \eqref{qh:hq-L} using the stated duality pairings.
In particular, differentiation of \(v^*v\) gives
\(2\mathcal V_3\operatorname{Re}(v^*\delta\dot b)\), explaining
\(\pi_b=\mathcal V_3v\) rather than \(2\mathcal V_3v\) in this complex
duality convention. The velocity is the covariant velocity minus
\(K_{A_0}\); its Legendre transform is consequently \(H_0-\vartheta(K_{A_0})\).
The ordinary quadratic Legendre transform gives every coefficient of
\eqref{qh:hq-H}. Cyclicity gives
\[
 -\operatorname{tr}(\Pi_i[\xi,A_i])
       =-\operatorname{tr}([A_i,\Pi_i]\xi),\quad
 \operatorname{tr}(\Pi_\Phi[\xi,\Phi])
       =\operatorname{tr}([\Phi,\Pi_\Phi]\xi),
\]
while anti-Hermiticity gives
\[
 2\operatorname{Re}(\pi_b^*\xi b)
       =\operatorname{tr}((b\pi_b^*-\pi_b b^*)\xi).
\]
These identities prove the first expression for \(\mathcal G\);
substitution of the momenta proves the second. Variation of the
independent \(A_0\) multiplier imposes \(\mathcal G=0\), even after a
temporal gauge is selected for the evolution. The primary momentum
conjugate to \(A_0\) is zero; \(A_0\) is not an additional coordinate
of \(\mathcal Q\). Since \(\mathcal M_\xi=p\cdot K_\xi\), canonical
Poisson brackets give
\(\{\mathcal M_\xi,\mathcal M_\eta\}=\mathcal M_{[\xi,\eta]}\).
Gauge invariance of every trace and norm in \(H_0\) also gives
\(\{H_0,\mathcal M_\xi\}=0\).
\end{proof}

\paragraph{Coordinates and measure without changing the metrics.}
Choose any fixed orthonormal basis \(T_a\), \(1\leq a\leq8\), for
\(-\operatorname{tr}\) on \(\mathfrak{su}(3)\), and put \(S_a=iT_a\).
Then \(\operatorname{tr}(S_aS_b)=\delta_{ab}\). Write
\[
 A_i=\sum_a a_{ia}T_a,\quad
 \Phi=\sum_a\phi_aS_a,\quad b_r=x_r+iy_r.
\]
These are coordinates on the original spaces, not a rescaling of their
metrics. The mass metric read from \eqref{qh:hq-L} is
\[
 M=\operatorname{diag}\left(
   \frac{2\mathcal V_3}{g^2}I_{24},\
   \mathcal V_3I_8,\
   2\mathcal V_3I_6\right).
\]
We use its constant volume measure
\[
 d\mu_M=C_M\,d^{24}a\,d^8\phi\,d^3x\,d^3y,\qquad
 C_M=\sqrt{\det M}=\frac{2^{15}\mathcal V_3^{19}}{(g^2)^{12}},
\]
retaining this constant rather than setting it to one. The auxiliary
Hilbert space is \(\mathscr H=L^2(\mathcal Q,d\mu_M;\mathbb C)\).
In these coordinates canonical momentum is \(-i\hbar\partial\).
Thus the exact kinetic operator is
\begin{equation}\label{qh:hq-kinetic}
 T=-\frac{\hbar^2g^2}{4\mathcal V_3}\Delta_a
   -\frac{\hbar^2}{2\mathcal V_3}\Delta_\phi
   -\frac{\hbar^2}{4\mathcal V_3}(\Delta_x+\Delta_y).
\end{equation}
Equivalently the complex momentum is
\(\widehat\pi_{b_r}=-i\hbar\partial_{\bar b_r}
=-i\hbar(\partial_{x_r}+i\partial_{y_r})/2\);
its adjoint product is
\(-\hbar^2(\partial_{x_r}^2+\partial_{y_r}^2)/4\).
The Schr\"odinger equation is \(i\hbar\partial_t\psi=H\psi\).

% Root proof: exact scalar bound and combination with the matrix estimate.
\subsection{The retained cubic and a global scalar bound}
\label{qh:scalar-bound-section}

\begin{lemma}[Sharp cubic bound in the original complex coordinates]
\label{qh:scalar-bound}
For every traceless Hermitian $\Phi\in\operatorname{Herm}_3(\mathbb C)$
and $b\in\mathbb C^3$, put
\[
 s=\operatorname{tr}\Phi^2+2b^*b,\qquad
 N=\det\Phi-b^*\Phi b.
\]
Then
\begin{equation}\label{qh:cubic-bound}
 |N|\leq\frac{s^{3/2}}{3\sqrt6}.
\end{equation}
Consequently, with all original coefficients retained,
\begin{equation}\label{qh:scalar-coercivity}
 V(\Phi,b)=\frac{m^2s}{2}+\frac{\Gamma s^2}{4}-\kappa N
 \ \geq\ \frac{\Gamma s^2}{8}-C_\kappa,
 \qquad C_\kappa=\frac{\kappa^4}{54\Gamma^3},
\end{equation}
whenever $m^2\geq0$, $\Gamma>0$, and $\kappa\in\mathbb R$.
This is an inequality for the full potential, not its replacement by
the right-hand side.
\end{lemma}
\begin{proof}
Choose a unitary diagonalization of $\Phi$ and apply the same unitary
to $b$. The inverse conjugation restores both original fields; it
preserves $s,N$ and introduces no restriction on either field.
Write the three real eigenvalues as $\lambda_1,\lambda_2,\lambda_3$,
with sum zero, and let $z_i=|b_i|^2$, $u=\sum_i z_i$.
If $u>0$, the original cubic is the convex combination
\[
 N=\sum_{i=1}^3\frac{z_i}{u}
       \bigl(\lambda_1\lambda_2\lambda_3-u\lambda_i\bigr).
\]
For one index $i$ and the other two indices $j,k$, put
$d_i=\lambda_j-\lambda_k$. Trace zero gives
\[
 \lambda_j\lambda_k=(\lambda_i^2-d_i^2)/4,\qquad
 s=(3\lambda_i^2+d_i^2)/2+2u.
\]
Hence each term in parentheses is exactly
\[
 \lambda_i(\lambda_j\lambda_k-u)
       =\lambda_i^3-\frac{s\lambda_i}{2},\qquad
 |\lambda_i|\leq\sqrt{2s/3}.
\]
If $u=0$, the same identity equals $N$ for each $i$.
If $s=0$, both fields vanish and the assertion is immediate.
For $s>0$, the critical points of $x^3-sx/2$ on
$[-\sqrt{2s/3},\sqrt{2s/3}]$ are $x=\pm\sqrt{s/6}$.
At these two points and the two endpoints, its absolute value is
$s^{3/2}/(3\sqrt6)$. The elementary extreme-value theorem proves
the bound for every term and hence for their convex combination.
Equality is attained, for example, at
$\Phi=\operatorname{diag}(2r,-r,-r)$, $b=0$, $r\ne0$,
since then $s=6r^2$ and $N=2r^3$.

Put $\alpha=|\kappa|/(3\sqrt6)$ and $t=\sqrt{s}\geq0$.
Keep $m^2s/2$ as a nonnegative term when bounding the original
potential. The remaining estimate is
\[
 V\geq\frac{\Gamma t^4}{8}
       +\left(\frac{\Gamma t^4}{8}-\alpha t^3\right).
\]
For $\alpha>0$ the expression in parentheses has its global minimum
at $t=6\alpha/\Gamma$, with value
$-54\alpha^4/\Gamma^3=-\kappa^4/(54\Gamma^3)$.
For $\alpha=0$ it is nonnegative. This proves
\eqref{qh:scalar-coercivity} including both signs of $\kappa$.
\end{proof}


\begin{theorem}[Friedrichs realization with an explicit domain]
\label{qh:hq-Friedrichs}
Let \(W(q)=\mathcal V_3U(q)\) and
\[
 D=\operatorname{diag}\left(
   \frac{\hbar^2g^2}{4\mathcal V_3}I_{24},\
   \frac{\hbar^2}{2\mathcal V_3}I_8,\
   \frac{\hbar^2}{4\mathcal V_3}I_6\right),\qquad
 C=\frac{\mathcal V_3\kappa^4}{54\Gamma^3}.
\]
The differential expression \(T+W=-\nabla\cdot D\nabla+W\) on
\(C_c^\infty(\mathcal Q)\) has the following well-defined semibounded
Friedrichs realization \(H\). Its closed form and form domain are
\begin{align}
 \mathfrak h[\psi]&=\int\left((\nabla\psi)^*D\nabla\psi
                         +W|\psi|^2\right)d\mu_M,\notag\\
 \mathscr D&=H^1(\mathcal Q,d\mu_M)
                \cap L^2(\mathcal Q,(W+C)d\mu_M),\label{qh:hq-form}\\
 D(H)&=\{u\in\mathscr D:
          -\nabla\cdot D\nabla u+Wu\in L^2(\mathcal Q,d\mu_M)
                         \text{ in distributions}\}.\notag
\end{align}
The operator is self-adjoint, bounded below by \(-C\), and
\(C_c^\infty\) is a form core. No compactness assertion is needed
for this construction.
\end{theorem}
\begin{proof}
Lemma~\ref{qh:scalar-bound} proves directly in the original complex
coordinates that $V\geq\Gamma s^2/8-\kappa^4/(54\Gamma^3)$.
All the other terms of \(U\) are nonnegative. In particular \(W+C\)
is a continuous nonnegative polynomial, locally bounded everywhere.

The norm
\[
 \|u\|_{\mathscr D}^2
  =\int\left((\nabla u)^*D\nabla u+(W+C+1)|u|^2\right)d\mu_M
\]
is complete. Indeed a Cauchy sequence converges in \(H^1\) and in the
weighted \(L^2\) space, and their limits agree on every compact set.
This proves closedness. Smooth radial cutoffs converging to one
approximate any \(u\in\mathscr D\) in this norm: the potential and
gradient tails tend to zero, and the cutoff-gradient term is bounded
by a constant times \(R^{-2}\|u\|^2\). On a fixed compact set ordinary
mollification converges in \(H^1\), and the bounded weight there gives
weighted convergence. Thus \(C_c^\infty\) is a form core and the
form is densely defined.

For clarity one may construct its self-adjoint operator directly.
The positive form \(\mathfrak h+(C+1)\langle\cdot,\cdot\rangle\)
is an inner product on \(\mathscr D\). Riesz representation assigns
to every \(f\in\mathscr H\) a unique \(u=Rf\in\mathscr D\) such that
\[
 \mathfrak h(u,v)+(C+1)\langle u,v\rangle=\langle f,v\rangle
                  \quad(v\in\mathscr D).
\]
It follows that \(R\), viewed on \(\mathscr H\), is bounded,
self-adjoint, positive and injective; its range is dense because
its kernel is zero. The spectral theorem for bounded self-adjoint
operators gives the self-adjoint inverse \(R^{-1}\) on
\(\operatorname{Ran}R\). Set \(H=R^{-1}-(C+1)\).
This is precisely the Friedrichs operator of the original form core.
Integration by parts on compact test functions identifies its
distributional expression. Conversely, for \(u\in\mathscr D\) whose
displayed distribution is \(f\in L^2\), that integration-by-parts
identity extends from the form core to all of \(\mathscr D\);
the Riesz characterization gives \(u\in D(H)\), \(Hu=f\).
This proves exactly the asserted domain. It does not replace the
domain by an unjustified separate requirement \(Wu,\Delta u\in L^2\).
\end{proof}

\begin{theorem}[Spatially constant Gauss constraints and averaging]
\label{qh:hq-physical}
For \(g\in SU(3)\), let
\[
 R_gq=(gA_ig^{-1},g\Phi g^*,gb),\qquad
 (\mathcal U(g)\psi)(q)=\psi(R_{g^{-1}}q).
\]
This is a strongly continuous unitary representation on \(\mathscr H\).
For each \(\xi\in\mathfrak{su}(3)\), the operator
\[
 \widehat{\mathcal M}_\xi=-i\hbar K_\xi\cdot\nabla
                          \quad\hbox{on }C_c^\infty(\mathcal Q)
\]
is essentially self-adjoint. Its closure has domain
\[
 D(\widehat{\mathcal M}_\xi)
  =\{\psi\in\mathscr H:K_\xi\cdot\nabla\psi\in\mathscr H
                                      \text{ distributionally}\},
 \qquad
 \mathcal U(e^{t\xi})=e^{-it\widehat{\mathcal M}_\xi/\hbar}.
\]
The spatially constant physical Hilbert space and projection are
\begin{align}
 \mathscr H_{\rm phys}
  &=\{\psi:\mathcal U(g)\psi=\psi\ \forall g\in SU(3)\}
    =\bigcap_{a=1}^8\ker\widehat{\mathcal M}_{T_a},\notag\\
 \mathsf P\psi&=\frac1{\mathcal V_\mu}
       \int_{SU(3)}\mathcal U(g)\psi\,d\mu(g),\qquad
 \mathcal V_\mu=\int_{SU(3)}d\mu(g)>0 .\label{qh:hq-average}
\end{align}
Here the Haar mass is retained explicitly. The space is infinite
dimensional. The set \(C_c^\infty(\mathcal Q)^{SU(3)}
=\mathsf P C_c^\infty(\mathcal Q)\) is dense in it and is a form
core for its physical Hamiltonian
\[
 H_{\rm phys}=H|_{\mathscr H_{\rm phys}},\qquad
 D(H_{\rm phys})=D(H)\cap\mathscr H_{\rm phys}.
\]
Thus the physical Hamiltonian is self-adjoint and is exactly the
Friedrichs operator of the restricted form; no quotient-coordinate
measure or singular-orbit boundary condition has been inserted.
\end{theorem}
\begin{proof}
Each \(R_g\) preserves the three original inner products, hence the
mass metric, its measure and the kinetic form. It preserves \(U\)
by conjugation invariance of the trace/determinant and unitarity of
the triplet action. Continuity holds first on compact smooth
functions and then on \(L^2\) by density and unitarity.
The vector field \(K_\xi\) is linear with zero divergence: its flow
is an orthogonal linear map for each block. Its action on compact
smooth functions is therefore skew-adjoint under integration, so
\(-i\hbar K_\xi\cdot\nabla\) is symmetric. Differentiating the
displayed unitary action gives the stated generator sign.
Stone's theorem gives a self-adjoint generator. Its domain is
the displayed distributional domain: one direction follows by
differentiation against test functions; for the converse,
distributional integration along the flow gives
\[
 \mathcal U(e^{t\xi})\psi-\psi
   =-\int_0^t\mathcal U(e^{s\xi})
                   (K_\xi\cdot\nabla\psi)\,ds
\]
as an \(L^2\) identity, and strong continuity gives differentiability.

To verify that the original compact smooth domain is a core, use
smooth cutoffs depending on the blockwise invariant Euclidean radius.
They commute with \(K_\xi\). Convolution with a radial smooth kernel
also commutes with its orthogonal flow. Cutoff followed by convolution
therefore approximates both \(\psi\) and \(K_\xi\cdot\nabla\psi\)
in \(L^2\). This proves graph-core density and essential
self-adjointness. On that common smooth core,
\[
 [K_\xi\cdot\nabla,K_\eta\cdot\nabla]
   =-K_{[\xi,\eta]}\cdot\nabla,\qquad
 [\widehat{\mathcal M}_\xi,\widehat{\mathcal M}_\eta]
   =i\hbar\widehat{\mathcal M}_{[\xi,\eta]}.
\]
There is no divergence correction or ordering anomaly.

An invariant vector lies in every generator domain with derivative
zero. Conversely the eight vanishing distributional derivatives
give the vanishing derivative for every \(\xi\), hence invariance
under every one-parameter subgroup. These generate connected \(SU(3)\).
For example each unitary determinant-one matrix is unitarily
diagonalizable with eigenphases whose sum can be chosen to be zero,
so it is the exponential of a trace-free anti-Hermitian matrix.
This proves the kernel characterization with its domain included.

The Bochner integral in \eqref{qh:hq-average} is bounded in norm
by one. Haar invariance and inversion imply
\(\mathsf P^2=\mathsf P=\mathsf P^*\), with range precisely the
invariant vectors. Averaging a compact smooth function remains
smooth and has compact support, because the union of its supports
under a compact group is compact. Applying \(\mathsf P\) to any
compact-smooth approximation proves \(L^2\) density.
Nonzero smooth functions supported in pairwise disjoint annuli of
the invariant squared radius are invariant and orthogonal.
This proves infinite dimensionality by explicit test states.

The form and its domain are \(\mathcal U(g)\)-invariant.
The Riesz characterization above therefore gives
\(\mathcal U(g)R=R\mathcal U(g)\); averaging gives
\(\mathsf PR=R\mathsf P\). Functional calculus proves that
\(\mathsf P\) reduces \(H\), so its restriction has the claimed
self-adjoint domain. Jensen's inequality for the positive form
\(\mathfrak h+(C+1)\|\cdot\|^2\) makes \(\mathsf P\) a contraction
in the form norm. Averaging a form-core approximation of an
invariant vector proves the final form-core statement.
\end{proof}

% Integration copy with qh: labels and a_g,b_g; original proof preserved.
% Independent bounded pure-gauge matrix audit; no canonical edits.
\subsection*{An explicit quantum bound for the three retained gauge matrices}

Let $\mathfrak g=\mathfrak{su}(3)$ as a real eight-dimensional vector
space, with the retained metric
\[
 \langle X,Y\rangle=-\operatorname{tr}(XY),\qquad
 \|X\|^2=-\operatorname{tr}(X^2).
\]
For anti-Hermitian $X,Y$, the trace of $XY$ is real, by taking its
conjugate and using cyclicity. The metric is positive definite because
$-\operatorname{tr}(X^2)=\operatorname{tr}(X^*X)$.
Lebesgue measure and all Laplacians below are those of this metric,
not of the negative Killing form. Write $A=(A_1,A_2,A_3)\in\mathfrak g^3$.
For two arbitrary retained numbers $a_g,b_g>0$, consider on
$L^2(\mathfrak g^3,dA)$ the nonnegative quadratic form
\[
 q_{a_g,b_g}[\psi]=
 a_g\sum_{i=1}^3\int\|\nabla_{A_i}\psi\|^2\,dA+
 b_g\sum_{1\leq i<j\leq3}
       \int\|[A_i,A_j]\|^2|\psi|^2\,dA,                       \tag{MQ.1}
\]
initially on $C_c^\infty(\mathfrak g^3;\mathbb C)$.
The coefficients $a_g,b_g$ are independent retained positive numbers;
their subscripts distinguish them from the scalar data. No assertion about the scalar
coupling or a Gauss-reduced operator is included in the theorem below.

\begin{lemma}[The exact adjoint frequency sum, including root walls]
\label{qh:matrix-root-sum}
Let $X\in\mathfrak g$, and let $\theta_1\leq\theta_2\leq\theta_3$
be the real eigenvalues of $-iX$, with $\sum_p\theta_p=0$.
The positive semidefinite real operator
$K_X=-\operatorname{ad}_X^2$ has its eight eigenvalues, with real
multiplicity, equal to
\[
 0,0,\quad
 (\theta_1-\theta_2)^2,(\theta_1-\theta_2)^2,\quad
 (\theta_1-\theta_3)^2,(\theta_1-\theta_3)^2,\quad
 (\theta_2-\theta_3)^2,(\theta_2-\theta_3)^2.                  \tag{MQ.2}
\]
In particular, zeros in this list remain zeros at repeated eigenvalues.
Define the frequency sum with real multiplicities by
$\mathcal S(X)=\operatorname{tr}_{\mathbb R\mathfrak g}\sqrt{K_X}$.
Then
\[
 \mathcal S(X)
 =2\sum_{p<q}|\theta_p-\theta_q|
 =4(\theta_3-\theta_1),\qquad
 2\sqrt6\,\|X\|\leq\mathcal S(X)\leq4\sqrt2\,\|X\|.           \tag{MQ.3}
\]
Both constants are sharp. The lower equality at nonzero $X$ occurs
exactly when two eigenvalues coincide; the upper equality occurs exactly
when the eigenvalues are $(-r,0,r)$ for some $r\geq0$.
If each positive root is counted only once instead of using real
multiplicity, the sum is $\mathcal S(X)/2$ and its sharp lower
constant is $\sqrt6$.
\end{lemma}
\begin{proof}
Cyclicity gives
$\langle[X,Y],Z\rangle=-\langle Y,[X,Z]\rangle$.
Thus $\operatorname{ad}_X$ is skew-adjoint and
$\langle Y,K_XY\rangle=\|[X,Y]\|^2$.
Unitary diagonalization of $X$ is an isometry for the retained metric.
Multiplication of that unitary by a scalar phase adjusts its determinant
to one without changing its conjugation, so this diagonalization can
also be performed by $SU(3)$.

For diagonal $X=i\operatorname{diag}(\theta_1,\theta_2,\theta_3)$,
use the orthonormal Cartan vectors
\[
 D_1=\frac{i}{\sqrt2}\operatorname{diag}(1,-1,0),\qquad
 D_2=\frac{i}{\sqrt6}\operatorname{diag}(1,1,-2),
\]
and, for $p<q$, the orthonormal pair
\[
 U_{pq}=\frac{E_{pq}-E_{qp}}{\sqrt2},\qquad
 V_{pq}=\frac{i(E_{pq}+E_{qp})}{\sqrt2}.
\]
Direct multiplication gives
$[X,U_{pq}]=(\theta_p-\theta_q)V_{pq}$ and
$[X,V_{pq}]=-(\theta_p-\theta_q)U_{pq}$, while $D_1,D_2$
commute with $X$. These eight orthonormal vectors span $\mathfrak g$,
proving (MQ.2) and the exact sum in (MQ.3).

Put $r=\theta_2-\theta_1\geq0$ and $s=\theta_3-\theta_2\geq0$.
The original eigenvalues are exactly
\[
 \theta_1=-(2r+s)/3,\quad \theta_2=(r-s)/3,\quad
 \theta_3=(r+2s)/3,\quad
 \|X\|^2=\frac23(r^2+rs+s^2),\quad \mathcal S(X)=4(r+s).
\]
Consequently
\[
 \mathcal S(X)^2-24\|X\|^2=16rs,\qquad
 32\|X\|^2-\mathcal S(X)^2=\frac{16}{3}(r-s)^2.
\]
Their nonnegativity proves the bounds and their vanishing proves every
equality case. In particular the lower bound does not fail at a repeated
eigenvalue: at a nonzero wall exactly four real frequencies remain
positive. At $X=0$ both sides vanish.
\end{proof}

\begin{lemma}[Slice oscillators with all coefficients]
\label{qh:matrix-slice}
For fixed $X\in\mathfrak g$ and $\alpha,\beta>0$,
\[
 -\alpha\Delta_Y+\beta\|[X,Y]\|^2
 \ \geq\ \sqrt{\alpha\beta}\,\mathcal S(X)
 \quad\hbox{on }L^2(\mathfrak g,dY)                           \tag{MQ.4}
\]
as a quadratic-form inequality. No positive frequency is assigned to
the commuting coordinates.
\end{lemma}
\begin{proof}
For $\omega\geq0$ and $f\in C_c^\infty(\mathbb R;\mathbb C)$,
integration by parts proves the exact identity
\[
 \int|\sqrt\alpha f'+\sqrt\beta\,\omega y f|^2\,dy
 =\alpha\int|f'|^2\,dy+
   \beta\omega^2\int y^2|f|^2\,dy
   -\sqrt{\alpha\beta}\,\omega\int|f|^2\,dy .
\]
Indeed $2\operatorname{Re}(f'\bar f)=(|f|^2)'$ and the boundary
term vanishes. Thus each real oscillator contributes
$\sqrt{\alpha\beta}\,\omega$, including the correct zero contribution
when $\omega=0$.

Diagonalize the real symmetric $K_X$ in an orthonormal basis. This
orthogonal coordinate change preserves measure and $\Delta_Y$, and
the potential becomes the sum $\sum_{\nu=1}^8\omega_\nu^2 y_\nu^2$
with the frequencies in (MQ.2). Apply the preceding identity in
each coordinate, integrate the other coordinates, and sum to obtain
(MQ.4). This proof is for each fixed $X$. It requires no differentiable
choice of the diagonalizing basis as $X$ varies, and introduces no
derivatives in $X$.
\end{proof}

\begin{theorem}[An explicit coercive form bound and compact resolvent]
\label{qh:matrix-compact}
The closure of (MQ.1) has domain
\[
 \mathcal D(q_{a_g,b_g})=
 \left\{\psi\in H^1(\mathfrak g^3):
  \sum_{i<j}\int\|[A_i,A_j]\|^2|\psi|^2\,dA<\infty\right\}.    \tag{MQ.5}
\]
Let $H_A$ be its nonnegative Friedrichs operator. Its formal expression
on compactly supported smooth functions is exactly
\[
 H_A=-a_g\sum_i\Delta_{A_i}
                 +b_g\sum_{i<j}\|[A_i,A_j]\|^2.                \tag{MQ.6}
\]
For every $0\leq\vartheta<1$ the following form bound holds:
\[
 q_{a_g,b_g}[\psi]\geq
 \vartheta a_g\int\|\nabla\psi\|^2\,dA+
 2\sqrt{6(1-\vartheta)a_gb_g}
       \int\left(\sum_i\|A_i\|\right)|\psi|^2\,dA.             \tag{MQ.7}
\]
In particular, retaining half of the original kinetic energy gives
\[
 H_A\geq-\frac{a_g}{2}\Delta_{\mathfrak g^3}
                 +2\sqrt{3a_gb_g}\sum_i\|A_i\|.             \tag{MQ.8}
\]
The form-domain inclusion into $L^2(\mathfrak g^3)$ is compact.
Consequently $(H_A+\mu I)^{-1}$ is compact for every $\mu>0$,
and $H_A$ has only isolated eigenvalues of finite multiplicity,
with no finite accumulation point.
\end{theorem}
\begin{proof}
The potential in (MQ.6) is a continuous nonnegative polynomial.
The right side of (MQ.5), equipped with
$\|\psi\|_2^2+q_{a_g,b_g}[\psi]$, is complete: an associated Cauchy
sequence converges in $H^1$ and also after multiplication by the
square root of the potential; an almost-everywhere convergent
subsequence identifies the second limit as that product.
It is densely defined since it contains $C_c^\infty$.
Smooth cutoffs equal to one on a ball of radius $R$, zero outside
the ball of radius $2R$, with derivative bounded by a constant
times $R^{-1}$, approximate each member in this norm as $R\to\infty$.
The derivative error tends to zero by the $H^1$ norm and the
$R^{-1}$ bound; the weighted error tends to zero by integrability.
For each cutoff, convolution with a compactly supported smooth
approximate identity converges in $H^1$, and also in weighted
$L^2$ because the potential is bounded on the common compact
support. This proves the stated core and closure.

For the lower bound start on this core and split the quadratic form
exactly as
\[
 q_{a_g,b_g}[\psi]=
 \vartheta a_g\|\nabla\psi\|_2^2+
 \sum_{\substack{i,j\in\{1,2,3\}\\i\ne j}}
 \left\{\frac{(1-\vartheta)a_g}{2}\|\nabla_{A_j}\psi\|_2^2
       +\frac{b_g}{2}\int\|[A_i,A_j]\|^2|\psi|^2\,dA\right\}.
                                                               \tag{MQ.9}
\]
Every derivative appears in two directed pairs, and every unordered
potential pair also appears twice. Hence no kinetic or potential
term has been counted more than its original coefficient.
Apply (MQ.4) inside each directed summand with $X=A_i$,
$Y=A_j$, $\alpha=(1-\vartheta)a_g/2$, $\beta=b_g/2$.
Fubini integration in $A_i$ and in the third matrix gives the lower
bound
\[
 \frac{\sqrt{(1-\vartheta)a_gb_g}}2
       \int\mathcal S(A_i)|\psi|^2\,dA
\]
for that summand. There are two partners $j$ for each fixed $i$.
Their sum and (MQ.3) prove (MQ.7), and $\vartheta=1/2$ gives
(MQ.8). Extension to (MQ.5) follows from core convergence and
Fatou's lemma for the nonnegative position weight.

Here is a direct compactness proof. Fix $0<\vartheta<1$ and put
$c_\vartheta=2\sqrt{6(1-\vartheta)a_gb_g}>0$.
In any one retained orthonormal basis of $\mathfrak g^3=\mathbb R^{24}$,
let $Q_R=[-R,R]^{24}$. Outside this cube,
$\sum_i\|A_i\|\geq(\sum_i\|A_i\|^2)^{1/2}\geq R$. Thus
\[
 \int_{\mathbb R^{24}\setminus Q_R}|\psi|^2
 \leq\frac{q_{a_g,b_g}[\psi]}{c_\vartheta R},\qquad
 \int|\nabla\psi|^2\leq\frac{q_{a_g,b_g}[\psi]}{\vartheta a_g}.
                                                               \tag{MQ.10}
\]
Partition $Q_R$ into finitely many cubes of side $h=2R/N$,
and let $\Pi_{R,N}\psi$ be its mean on each small cube and zero
outside $Q_R$. This has finite-dimensional range. On any cube
$Q$ of side $h$,
\[
 \int_Q|\psi-\psi_Q|^2\leq h^2\int_Q|\nabla\psi|^2.            \tag{MQ.11}
\]
For completeness, in one dimension the variance identity gives
$\int_0^h|f-f_{\rm av}|^2
=(2h)^{-1}\int_0^h\int_0^h|f(x)-f(y)|^2\,dx\,dy$.
Cauchy--Schwarz in the fundamental theorem of calculus bounds
this by $h^2\int_0^h|f'|^2$. In a product cube, let $P_j$
be averaging in the $j$th coordinate. The commuting orthogonal
projections give the orthogonal decomposition
\[
 I-\prod_{j=1}^{24}P_j
 =\sum_{j=1}^{24}\left(\prod_{k<j}P_k\right)(I-P_j).
 \]
Applying the one-dimensional estimate to each summand and using
that prior-coordinate averaging contracts the derivative norm
proves (MQ.11), first for smooth functions and then by $H^1$
density.

Equations (MQ.10)--(MQ.11) give the explicit approximation bound
\[
 \|\psi-\Pi_{R,N}\psi\|_2^2
 \leq q_{a_g,b_g}[\psi]
       \left(\frac{h^2}{\vartheta a_g}
                    +\frac1{c_\vartheta R}\right).           \tag{MQ.12}
\]
On a bounded form-domain set choose first $R$ large and then
$N$ large. The error is uniformly as small as prescribed,
while $\Pi_{R,N}$ maps that set to a bounded subset of a fixed
finite-dimensional space. Such a set has a finite net at any
positive radius, by subdividing a bounding coordinate box.
Thus the original set is totally bounded in $L^2$, proving
compactness of the form-domain inclusion without a matrix-model
discreteness theorem.

The closed form defines its Friedrichs operator by the usual
Riesz construction: for each $f\in L^2$ and $\mu>0$ there is a
unique $u$ with
$q_{a_g,b_g}(u,v)+\mu\langle u,v\rangle=\langle f,v\rangle$
for every form-domain $v$. This follows by representing the
bounded functional on the Hilbert space with that form inner
product; coercivity is supplied by $\mu\|v\|_2^2$.
The solution is $(H_A+\mu I)^{-1}f$ and obeys
$q_{a_g,b_g}[u]+\mu\|u\|_2^2\leq\|f\|_2^2/\mu$.
Hence this resolvent factors through a bounded set in the form
domain and its compact inclusion, so it is compact.
The spectral theorem for the resulting compact self-adjoint
resolvent gives finite multiplicities and eigenvalues tending
to zero for the resolvent, equivalently eigenvalues of $H_A$
tending to infinity. Since the resolvent is injective and its
range is dense in an infinite-dimensional space, it has
infinitely many such eigenvalues. This proves all assertions.
\end{proof}

The inequality is a quadratic-form inequality, not a pointwise
lower bound on the quartic potential. That potential still vanishes
on every commuting triple. The exact mechanism is (MQ.9):
localizing the other matrices transversely to a large fixed matrix
costs the nonzero oscillator ground energies in (MQ.4); their
sum remains coercive even at a root wall.

\begin{proposition}[An additional explicit lower bound on the energy origin]
\label{qh:matrix-bottom}
With precisely the energy origin in (MQ.6),
\[
 H_A\geq 9(3/2)^{1/3}a_g^{2/3}b_g^{1/3}\,I.                    \tag{MQ.13}
\]
This is a lower bound on the bottom of the operator, not a claimed
lower bound of the same size on the difference between its first two
energy levels.
\end{proposition}
\begin{proof}
First, for any $\alpha,\gamma>0$,
\[
 -\alpha\frac{d^2}{dx^2}+\gamma|x|
 \geq \frac3{2^{8/3}}(\alpha\gamma^2)^{1/3}.                 \tag{MQ.14}
\]
To prove it, take $f\in C_c^\infty(\mathbb R)$ with $\|f\|_2=1$,
and write $t=\|f'\|_2^2$ and $M=\int|x||f(x)|^2\,dx>0$.
Integration of $(|f|^2)'$ from both ends of the line gives
$\|f\|_\infty^2\leq\|f\|_2\|f'\|_2=\sqrt t$.
Markov's inequality puts at least half of the mass inside
$[-2M,2M]$, so
$1/2\leq4M\|f\|_\infty^2\leq4M\sqrt t$.
Thus $tM^2\geq1/64$. Minimizing
$\alpha/(64M^2)+\gamma M$ over $M>0$, whose unique minimum
occurs at $M^3=\alpha/(32\gamma)$, proves (MQ.14).
Homogeneity in $f$ removes the unit-norm restriction.

Write $x_{i\nu}$ for the original metric-orthonormal real coordinates,
$1\leq i\leq3$, $1\leq\nu\leq8$. Cauchy--Schwarz gives
$\|A_i\|\geq8^{-1/2}\sum_{\nu=1}^8|x_{i\nu}|$.
By (MQ.8),
\[
 H_A\geq
 \sum_{i,\nu}\left\{
    -\frac{a_g}{2}\partial_{x_{i\nu}}^2
           +\sqrt{\frac{3a_gb_g}{2}}\,|x_{i\nu}|\right\}.
\]
Apply (MQ.14) in each of these twenty-four coordinates and
integrate the others. Its total constant is exactly
\[
 24\frac3{2^{8/3}}\left(\frac{a_g}{2}\,\frac{3a_gb_g}{2}\right)^{1/3}
 =9(3/2)^{1/3}a_g^{2/3}b_g^{1/3}.
\]
Core closure gives the operator-form bound stated.
\end{proof}

\paragraph{Exact relation to the retained homogeneous gauge density.}
The canonical density uses signature $(+---)$ and anti-Hermitian
gauge fields, and its pure gauge part is
\[
 \mathcal L_{\rm gauge,hom}
 =\frac1{g^2}\sum_i\|E_i\|^2
       -\frac1{g^2}\sum_{i<j}\|[A_i,A_j]\|^2,\qquad
 E_i=\dot A_i+[A_0,A_i].
\]
The full scalar potential retained there is
$V=m^2s/2+\Gamma s^2/4-\kappa(\det\Phi-b^*\Phi b)$,
with $s=\operatorname{tr}\Phi^2+2b^*b$.
It is not part of the pure matrix operator (MQ.6).
For the pure gauge homogeneous kinetic form on a finite torus of
retained volume $\mathcal V_3>0$, fixing temporal gauge at the
unreduced coordinate level gives
$L=\mathcal V_3\sum_i\|\dot A_i\|^2/g^2
-\mathcal V_3\sum_{i<j}\|[A_i,A_j]\|^2/g^2$.
Its canonical momenta are exactly
$p_i=2\mathcal V_3\dot A_i/g^2$; the Legendre transform is
\[
 H_{\rm cl}=\frac{g^2}{4\mathcal V_3}\sum_i\|p_i\|^2
              +\frac{\mathcal V_3}{g^2}
                       \sum_{i<j}\|[A_i,A_j]\|^2.
\]
The specified canonical substitution $p_i=-i\hbar\nabla_{A_i}$
with $\hbar>0$ therefore gives the exact coefficient map
\[
 a_g=\frac{\hbar^2g^2}{4\mathcal V_3},\qquad
 b_g=\frac{\mathcal V_3}{g^2},\qquad
 a_gb_g=\frac{\hbar^2}{4}.                                  \tag{MQ.15}
\]
No coupling or volume is set to one. Under this map, (MQ.8) reads
\[
 H_A\geq-\frac{\hbar^2g^2}{8\mathcal V_3}\Delta
                   +\sqrt3\,\hbar\sum_i\|A_i\|,
\]
and (MQ.13) reads
$H_A\geq9(3\hbar^4g^2/(32\mathcal V_3))^{1/3}I$.
These coefficient identities concern the displayed unreduced
homogeneous quantum model. They impose neither the Gauss constraint
nor the omitted scalar coupling, and assert no continuum or
infinite-volume field-theoretic gap.

\paragraph{Primary literature and scope.}
The discreteness phenomenon for bosonic commutator matrix potentials
is classical. B.~Simon, \emph{Some quantum operators with discrete
spectrum but classically continuous spectrum}, Annals of Physics
\textbf{146} (1983), 209--220, DOI
\url{https://doi.org/10.1016/0003-4916(83)90057-X}, treats a
compact semisimple Lie-algebra matrix operator in equation~(3)
and Section~7, Corollary~4. Section~2 gives the elementary
zero-point-oscillator argument for the two-variable model.
The present proof combines a direct slice argument with the
explicit $\mathfrak{su}(3)$ root sum in the metric
$-\operatorname{tr}(XY)$ and proves compactness directly by
(MQ.12). It does not use the Fefferman--Phong discreteness theorem
employed in that paper's Section~7, nor claim the underlying
discreteness result as new.

The related finite-volume Yang--Mills analysis is M.~L\"uscher,
\emph{Some analytic results concerning the mass spectrum of Yang--Mills
gauge theories on a torus}, Nuclear Physics B \textbf{219} (1983),
233--261, DOI \url{https://doi.org/10.1016/0550-3213(83)90436-4}.
That field-theoretic small-volume setting is not substituted for the
precisely specified matrix operator here; none of its perturbative
conclusions or coefficients is assumed in the foregoing proof.


\subsection{Compactness for the entire coupled Hamiltonian}
\label{qh:combined-section}

\begin{theorem}[All thirty-eight coordinates are quantum confined]
\label{qh:full-compact}
Retain the configuration metric
\[
 \gamma_E=\operatorname{tr}(d\Phi^2)+2\,db^*db,
 \qquad E=\operatorname{Herm}_3(\mathbb C)_0\oplus\mathbb C^3,
\]
and the metric $-\operatorname{tr}$ on each of the three
$\mathfrak{su}(3)$ factors. Let $T_A=-\sum_i\Delta_{A_i}$ and
$T_E=-\Delta_{\gamma_E}$; these are nonnegative Laplacians.
For every $0<\eta<1$, the full Hamiltonian $H_{\rm hom}=H$ obeys
the quadratic-form inequality
\begin{align}
 H_{\rm hom}\ \geq\;&
 \frac{\eta\hbar^2g^2}{4\mathcal V_3}T_A
 +\frac{\hbar^2}{2\mathcal V_3}T_E
 +\hbar\sqrt{6(1-\eta)}\sum_{i=1}^3\|A_i\|
 \notag\\*[-1mm]
 &+\frac{\mathcal V_3\Gamma}{8}s^2
 -\mathcal V_3 C_\kappa.
 \label{qh:full-form-bound}
\end{align}
Its Friedrichs realization on the entire thirty-eight-dimensional
configuration space has compact resolvent. No commuting or scalar
directions have been deleted in this assertion.
\end{theorem}
\begin{proof}
Apply the matrix estimate with
$a_g=\hbar^2g^2/(4\mathcal V_3)$ and
$b_g=\mathcal V_3/g^2$. Their product is exactly $\hbar^2/4$.
The scalar-gradient and gauge--scalar potentials
\[
 \frac{\mathcal V_3}{2}\sum_i\operatorname{tr}[A_i,\Phi]^2,
 \qquad \mathcal V_3\sum_i\|A_i b\|^2
\]
are nonnegative: $[A_i,\Phi]$ is Hermitian. Retain these terms
in the definition of the operator and use their signs only in
this lower bound. Lemma~\ref{qh:scalar-bound} now proves
\eqref{qh:full-form-bound} on compactly supported smooth
functions, and core closure extends it to the form domain.

Write $q$ for the closed form and
$\widetilde q[\psi]=q[\psi]+\mathcal V_3C_\kappa\|\psi\|_2^2$.
The preceding bound makes $\widetilde q$ nonnegative.
In the original metric the squared radial coordinate is
$r^2=\sum_i\|A_i\|^2+s$.
The retained positive weight
\[
 W_\eta=\hbar\sqrt{6(1-\eta)}\sum_i\|A_i\|
                       +\mathcal V_3\Gamma s^2/8
\]
tends to infinity as $r\to\infty$. More explicitly, outside
the ball $r<R$ either $\sum_i\|A_i\|^2\geq R^2/2$ or
$s\geq R^2/2$, so
\[
 W_\eta\geq
 w_R:=\min\left\{\hbar\sqrt{3(1-\eta)}R,
                        \mathcal V_3\Gamma R^4/32\right\}.
\]
Both coefficients are strictly positive for the stated parameters.
Thus
\[
 \int_{r\geq R}|\psi|^2\leq\widetilde q[\psi]/w_R.
\]
The derivative bound from \eqref{qh:full-form-bound} is
\[
 c_0\int|\nabla\psi|^2\leq\widetilde q[\psi],\qquad
 c_0=\min\{\eta\hbar^2g^2/(4\mathcal V_3),
                          \hbar^2/(2\mathcal V_3)\}>0,
\]
where the gradient uses the product metric just stated.
Metric-orthonormal coordinates are an isometric representation,
not a change to any coupling, original field, or volume.
On the cube $[-R,R]^{38}$, replace a function by its averages on
finitely many cubes of side $h=2R/N$ and by zero outside the
large cube. The coordinate-averaging argument in the matrix proof
works in thirty-eight coordinates without changing its constant:
the squared local error is at most $h^2\int|\nabla\psi|^2$.
The cube's complement lies outside $r<R$, so the total error is at most
\[
 \widetilde q[\psi]\left(h^2/c_0+1/w_R\right).
\]
On every form-norm bounded set this tends uniformly to zero by
first increasing $R$ and then increasing $N$. The approximations
have finite-dimensional range, and their norms are bounded by
$\|\psi\|_2$. Therefore that set is totally bounded in $L^2$.
The complete Hilbert form domain has compact inclusion in $L^2$.
The Riesz construction with the positive form $\widetilde q+\|\cdot\|_2^2$ maps bounded
$L^2$ sets to bounded form-domain sets. Its composition with this
inclusion is the resolvent, proving compactness for the full operator.
\end{proof}

The original classical potential remains zero in its pure-gauge
part for every commuting triple. Inequality~\eqref{qh:full-form-bound}
uses the derivative cost in directions transverse to such triples.
The strictly positive coefficient in front of $\sum_i\|A_i\|$
does not assert a pointwise classical potential that was absent
from the original density.

\begin{theorem}[Strict positivity and the bottom spectral projection]
\label{qh:hq-positive}
For the actual polynomial Hamiltonian above and every \(\tau>0\),
\(e^{-\tau H}\) is positivity improving: if \(0\ne f\geq0\), then
\(e^{-\tau H}f>0\) almost everywhere. The same holds on the invariant
Hilbert space. Moreover
\[
 \inf\sigma(H_{\rm phys})=\inf\sigma(H)=:E_0,\qquad
 \operatorname{rank}\mathbf1_{\{E_0\}}(H)\leq1.
\]
Every nonzero vector in that bottom eigenspace is a scalar multiple
of one strictly positive normalized, \(SU(3)\)-invariant real
function. These conclusions hold without assuming compactness.
\end{theorem}
\begin{proof}
We give the heat-kernel argument with its constants. In the above real
coordinates the free kernel relative to Lebesgue measure is
\[
 p_\tau(q,r)=(4\pi\tau)^{-19}(\det D)^{-1/2}
        \exp\!\left(-\frac{(q-r)^{\mathsf T}D^{-1}(q-r)}{4\tau}\right).
\]
Let \(B_s\), \(0\leq s\leq\tau\), be a standard 38-dimensional
Brownian motion and form the bridge
\[
 Z_s=(1-s/\tau)q+(s/\tau)r
                   +\sqrt{2D}\,(B_s-(s/\tau)B_\tau).
\]
The Friedrichs kernel, relative to the retained measure \(d\mu_M\), is
\begin{equation}\label{qh:hq-FK}
 K_\tau(q,r)=\frac{p_\tau(q,r)}{C_M}\,
       \mathbb E\exp\!\left(-\int_0^\tau W(Z_s)\,ds\right).
\end{equation}
Here \(\tau\) is the semigroup parameter with inverse-energy units;
Euclidean physical time is \(\hbar\tau\). No value of \(\hbar\)
has been set to one in \(D\).

For completeness, the formula can be justified within the present
bounded-below polynomial setting. Put
\(W_n=\min(W+C,n)\). For bounded \(W_n\), the Duhamel iteration of
the free semigroup converges in norm, since its \(j\)-th ordered
term is bounded by \((n\tau)^j/j!\). The Markov Gaussian densities
of the bridge identify exactly these ordered integrals with the
power-series expansion of
\(\mathbb E\exp(-\int_0^\tau W_n(Z_s)\,ds)\).
Thus the kernel formula holds for the bounded potentials.
The forms with potentials \(W_n\) increase to the form with
potential \(W+C\). Their resolvents converge strongly: for
\(u_n=(T+W_n+\lambda)^{-1}f\), \(\lambda>0\), the variational
identities bound \(u_n\) in \(H^1\); testing each fixed truncated
form and passing weakly shows that any weak limit belongs to
the limiting form domain. Compactly supported tests identify
its unique limiting resolvent equation. The identities
\[
 \mathfrak h_n[u_n]+\lambda\|u_n\|^2
       =\operatorname{Re}\langle f,u_n\rangle
\]
and lower semicontinuity then force convergence of the \(L^2\)
norms, hence strong convergence. The spectral functional calculus
gives strong convergence of their semigroups.
Every bridge path is continuous on a compact interval and has
bounded image. The continuous polynomial \(W\) is therefore
bounded on that path, and monotone convergence of its truncated
integrals and dominated convergence of their exponentials yield
\eqref{qh:hq-FK}, restoring the factor \(e^{C\tau}\).
This is the bounded-below Brownian-bridge Feynman--Kac formula,
also stated in B.~Simon, \emph{A Feynman--Kac Formula for Unbounded
Semigroups}, arXiv:math-ph/9907022, Theorem~1.1,
\url{https://arxiv.org/abs/math-ph/9907022}.

In particular the exponential in \eqref{qh:hq-FK} is strictly
positive almost surely and bounded above by \(e^{C\tau}\).
Since \(p_\tau(q,r)>0\), the kernel is strictly positive for all
\(q,r\). Its Gaussian upper bound makes it a valid \(L^2\)
kernel action, and integration against any nonzero nonnegative
function proves positivity improvement. The invariant subspace
is reducing and the action is the same positive kernel, giving
its assertion as well.

To prove equality of the two spectral infima without presupposing
an eigenfunction, take any \(f\in\mathscr D\) and define
\[
 F(q)=\left(\frac1{\mathcal V_\mu}
       \int_{SU(3)}|f(R_{g^{-1}}q)|^2\,d\mu(g)\right)^{1/2}.
\]
This is an invariant nonnegative function with \(\|F\|=\|f\|\).
Apply the norm inequality to the Hilbert-valued Sobolev function
\(q\mapsto(g\mapsto f(R_{g^{-1}}q))\). More explicitly, differentiating
\((\int|f(R_{g^{-1}}q)|^2+\varepsilon^2)^{1/2}\), using
Cauchy--Schwarz, and then taking \(\varepsilon\downarrow0\), gives
\[
 \int(\nabla F)^*D\nabla F\,d\mu_M
 \leq \frac1{\mathcal V_\mu}\int_{SU(3)}
       \int(\nabla\mathcal U(g)f)^*D
                           \nabla\mathcal U(g)f\,d\mu_M\,d\mu(g).
\]
One first argues locally for smooth functions and then uses form
approximation and lower semicontinuity. Invariance of \(W\) gives
equality of the potential integrals. Thus \(F\in\mathscr D\) and
\(\mathfrak h[F]\leq\mathfrak h[f]\). The variational
characterization of the spectral infimum proves
\(\inf\sigma(H_{\rm phys})\leq\inf\sigma(H)\); the reverse inequality
follows from restriction to a subspace.

Finally suppose the bottom spectral projection is nonzero, and take
a real nonzero bottom eigenfunction \(u\); real and imaginary parts
are eigenfunctions because \(H\) is real. The weak-gradient inequality
\(|\nabla|u||\leq|\nabla u|\) gives
\(\mathfrak h[|u|]\leq E_0\|u\|^2\). The lower spectral bound forces
equality, and the spectral theorem then gives
\(|u|\in\ker(H-E_0)\). Positivity improvement makes \(|u|>0\)
almost everywhere. A sign-changing real bottom eigenfunction is
impossible: its positive and negative parts are both sent to
strictly positive functions, so
\[
 |e^{-\tau H}u|<e^{-\tau H}|u|
\]
almost everywhere, contradicting the eigenfunction identities.
Two independent real bottom eigenfunctions would yield, by
orthogonalization against the positive one, such a sign-changing
eigenfunction. Thus the complex eigenspace has dimension one.
Every \(\mathcal U(g)|u|\) is another positive normalized bottom
eigenfunction, so its one-dimensional proportionality constant
is positive and of modulus one, hence one. This proves gauge
invariance and all assertions of the theorem.
\end{proof}

\begin{corollary}[Unique physical ground state and a finite-model gap]
\label{qh:hq-ground}
The actual homogeneous quantum Hamiltonian has one normalized positive
ground state, unique up to phase and invariant under spatially constant
\(SU(3)\). The physical eigenvalues, counted with multiplicity, satisfy
\[
 E_0=E^{\rm phys}_0<E^{\rm phys}_1\leq E^{\rm phys}_2\leq\cdots,
 \qquad E^{\rm phys}_n\longrightarrow+\infty .
\]
In particular \(E^{\rm phys}_1-E_0>0\), without a numerical value being
claimed.
\end{corollary}
\begin{proof}
The positive compact resolvent $(H+C+1)^{-1}$ proved above has a nonzero top eigenvalue,
whose inverse supplies an eigenfunction at the bottom of \(H\).
The strict-positivity theorem makes it unique and invariant.
The physical space is infinite dimensional, and its compact
resolvent has finite-multiplicity eigenvalues accumulating only
at zero. Applying the inverse affine spectral relation proves
the displayed sequence and its strict first inequality.
\end{proof}
This is a gap of the specified finite-dimensional quantum mechanics.
No numerical value follows from positivity, and classical slow
solutions do not compute its level spacings. The proof retains fixed
positive \(\hbar,g^2,\mathcal V_3,\Gamma\); it asserts no uniform
gap in limits of these parameters or in a continuum limit.

\begin{proposition}[Exact classical embedding and its quotient limitations]
\label{qh:hq-embedding}
Let \(\iota\) send each homogeneous configuration and velocity to
the corresponding spatially constant fields on the trivial
\(SU(3)\) bundle over \(\mathbb T^3\).
The full action pulls back to \eqref{qh:hq-L}. Its homogeneous
Euler--Lagrange and Gauss solutions are exact full classical solutions.
On phase space use constant momentum densities
\(\Pi_i/\mathcal V_3,\Pi_\Phi/\mathcal V_3,\pi_b/\mathcal V_3\);
then the full canonical one-form pulls back to \(\vartheta\)
exactly. The map is equivariant for spatially constant \(SU(3)\).
The induced map from its constant-group orbit space to the full
torus gauge orbit space is not injective.
\end{proposition}
\begin{proof}
Spatial derivatives of the retained coefficients vanish. Substitution
in the density gives \eqref{qh:hq-L}, and the full local
Euler--Lagrange residuals are then constant in space.
Vanishing of their homogeneous variations for every time-dependent
matrix variation therefore makes those residuals zero pointwise,
including the \(A_0\) residual. They consequently vanish against
arbitrary, not necessarily homogeneous, spacetime test variations.
This proves the exact classical implication without a symmetry
reduction heuristic. Integrating constant momentum densities against
constant configuration variations multiplies them by \(\mathcal V_3\),
proving the asserted one-form pullback and thus its symplectic
pullback. Constant-group equivariance is direct.

For the limitation, let \(L_1\) be the first torus period and set
\[
 g(x_1)=\operatorname{diag}
   (e^{2\pi ix_1/L_1},e^{-2\pi ix_1/L_1},1).
\]
With \(D=\partial+A\), gauge transformation is
\[
 A_i'=gA_ig^{-1}-(\partial_i g)g^{-1},\quad
 \Phi'=g\Phi g^*,\quad b'=gb.
\]
The homogeneous state \(A_i=0,\Phi=b=0\) is thereby equivalent in
the full torus theory to the homogeneous state
\[
 A'_1=-\frac{2\pi i}{L_1}\operatorname{diag}(1,-1,0),\qquad
 A'_2=A'_3=0,\qquad \Phi'=b'=0 .
\]
No constant conjugation sends the zero \(A_1\) to this nonzero
matrix, proving noninjectivity of the orbit-space map.
On the other hand the same spatial gauge function sends the constant
\(\Phi=E_{12}+E_{21}\) to
\[
 e^{4\pi ix_1/L_1}E_{12}+e^{-4\pi ix_1/L_1}E_{21},
\]
which is not homogeneous. Thus spatial gauge maps do not act on
the whole homogeneous ansatz. More generally, a gauge map preserving
every homogeneous triplet \(b=e_j\), \(j=1,2,3\), must have all
three columns constant, and hence must itself be spatially constant.
This argument applies in particular to any genuinely large
transformation not homotopic to a constant. The explicit diagonal
example above is not being mislabeled as topologically large in
\(SU(3)\); its periodicity alone already proves the required
quotient obstruction.
\end{proof}

\paragraph{Why the classical embedding is not a quantum-sector theorem.}
There is a direct interaction obstruction already in a finite Fourier
cutoff. Choose the trace-free Hermitian
\(S=\operatorname{diag}(1,-1,0)/\sqrt2\), so
\(\operatorname{tr}S^2=1\), \(\det S=0\), and put
\[
 \Phi(x)=\bigl(q+y\cos(2\pi x_1/L_1)\bigr)S,\qquad A_i=b=0 .
\]
The quartic part of the spatially integrated potential is exactly
\[
 \frac{\Gamma\mathcal V_3}{4}
       \left(q^4+3q^2y^2+\frac38y^4\right).
\]
Indeed the averages of the cosine, its square, its cube and its
fourth power are \(0,1/2,0,3/8\). The mass and gradient terms
separate in \(q,y\), the cubic determinant vanishes, but the mixed
term \(3\Gamma\mathcal V_3q^2y^2/4\) remains.
The kinetic density is
\(\mathcal V_3\dot q^2/2+\mathcal V_3\dot y^2/4\);
thus this example also retains the unequal Fourier-mode masses.

A putative product sector consisting of all \(q\)-wavefunctions
times one fixed nonzero normalizable \(y\)-wavefunction \(\Omega\)
is not invariant under this interacting cutoff Hamiltonian.
Here invariance means a reducing quantum subspace, equivalently that
its orthogonal projection commutes with the unitary evolution.
This includes the necessary domain issue: if that projection
commuted with the Hamiltonian, applying it to
\(f(q)v(y)\in C_c^\infty\), with \(\langle\Omega,v\rangle\ne0\),
would put \(f(q)\Omega(y)\) in the operator domain for every
\(f\in C_c^\infty(\mathbb R)\). Its distributional Hamiltonian is
the separated \(q\)-term plus
\(f(q)H_y\Omega+cq^2f(q)y^2\Omega\),
\(c=3\Gamma\mathcal V_3/4\). Integrating against two compact
\(q\)-test functions with distinct averaged \(q^2\) values
isolates \(H_y\Omega\) and \(y^2\Omega\) as \(L^2\) functions.
Thus the following projections are legitimate even if no regularity
of \(\Omega\) was assumed initially.
Projecting its action orthogonally to \(\Omega\), for arbitrary
compact smooth \(q\)-wavefunctions, would require
\[
 (H_y\Omega)_\perp+
       \frac{3\Gamma\mathcal V_3}{4}q^2(y^2\Omega)_\perp=0
                         \quad\hbox{for almost every }q .
\]
Its two coefficients must vanish, implying
\(y^2\Omega=\lambda_\Omega\Omega\). Multiplication by \(y^2\) has no nonzero
\(L^2(\mathbb R)\) eigenfunction: its level sets have Lebesgue
measure zero. Thus no fixed discarded-mode vacuum yields that
product invariant sector. The distribution \(\delta(y)\), which
would impose a literally vanishing discarded coordinate, is not
an \(L^2\) state either.
This does not rule out every possible correspondence with a full
quantum theory; it disproves the specific automatic-subsector
inference and states exactly which finite model has been quantized.
No continuum quantum Hilbert space, full local-gauge quotient,
renormalized Hamiltonian, or continuum mass-gap value is constructed
by the finite-dimensional argument.

\paragraph{Gauss-law reference.}
D.~Tong, \emph{Lectures on Gauge Theory}, Section 2.2.1,
equations (2.27)--\allowbreak(2.31),
\url{https://www.damtp.cam.ac.uk/user/tong/gaugetheory.html}, retains the Gauss equation after temporal
gauge and imposes it on physical states. His Hermitian-field
normalization is not substituted for the retained anti-Hermitian
normalization here; the exact factors and the sign of
\eqref{qh:hq-Gauss} were derived above from the given density.
For this compact spatial torus, spatially constant \(SU(3)\) is
the finite gauge subgroup being imposed, rather than an asymptotic
charge group at spatial infinity.
\par
\section{An exact moving-ground-state comparison for the retained Fourier pair}
\label{fq:section}

The result below concerns precisely the two-coordinate variational
truncation
\[
 \Phi(x)=\bigl(q+y\cos(kx_1)\bigr)S,\qquad
 S=\operatorname{diag}(1,-1,0)/\sqrt2,\qquad
 k=2\pi/L_1,\qquad A_i=b=0.
\]
It does not assume that this two-dimensional plane is an invariant
subspace of the full classical field equations or a reducing subspace
of an independently constructed full quantum field theory.
Throughout, retain
\[
 \hbar>0,\quad\mathcal V_3>0,\quad\Gamma>0,\quad L_1>0,
 \qquad m^2\geq0,
\]
\[
 \qquad \omega_*:=\sqrt{m^2+k^2}>0,
 \qquad \omega(q):=\sqrt{\omega_*^2+3\Gamma q^2}.
\]
All bounds will refer to these original parameters and coordinates.
The quartic discarded-mode interaction is never removed from the operator.

\subsection{Exact coefficients, measures, and self-adjoint operators}

Since $\operatorname{tr}S^2=1$ and $\det S=0$, the pullback of the
cubic determinant is zero. The averages of $\cos(kx_1)$ and its first
four powers are $0,1/2,0,3/8$ for powers one through four. The exact
integrated Lagrangian is therefore
\begin{align}
 L_2={}&\frac{\mathcal V_3}{2}\dot q^2
       +\frac{\mathcal V_3}{4}\dot y^2
       -V_q(q)-\frac{\mathcal V_3\omega(q)^2}{4}y^2-cy^4,
       \label{fq:L}\\
 V_q(q):={}&\mathcal V_3\left(\frac{m^2q^2}{2}
                         +\frac{\Gamma q^4}{4}\right),\qquad
 c:=\frac{3\Gamma\mathcal V_3}{32}>0.\notag
\end{align}
In particular the gradient term is $\mathcal V_3k^2y^2/4$ and the
mixed quartic term is $3\Gamma\mathcal V_3q^2y^2/4$.
The momenta are $p_q=\mathcal V_3\dot q$ and
$p_y=\mathcal V_3\dot y/2$. Thus the mass metric and its measures are
\[
 M_2=\operatorname{diag}(\mathcal V_3,\mathcal V_3/2),\quad
 d\mu_q=\sqrt{\mathcal V_3}\,dq,\quad
 d\mu_y=\sqrt{\mathcal V_3/2}\,dy,\quad
 d\mu_2=\frac{\mathcal V_3}{\sqrt2}\,dq\,dy=d\mu_qd\mu_y.
\]
We retain these constants. Put
\[
 \mathscr K=L^2(\mathbb R,d\mu_y),\qquad
 \mathscr L=L^2(\mathbb R,d\mu_q),\qquad
 \mathscr H_2=L^2(\mathbb R,d\mu_q;\mathscr K).
\]
\[
 a:=\frac{\hbar^2}{2\mathcal V_3},\quad
 a_y:=\frac{\hbar^2}{\mathcal V_3},\quad
 b_q:=\frac{\mathcal V_3\omega(q)^2}{4}.
\]
The Hamiltonian is exactly
\begin{equation}\label{fq:H}
 H_2=-a\partial_q^2+V_q(q)+h(q),\qquad
 h(q)=-a_y\partial_y^2+b_qy^2+cy^4.
\end{equation}
The symbols $a,a_y,b_q,c$ denote these retained coefficients, not a
rescaling of either variable.

Every displayed nonnegative Schr\"odinger expression is realized by
its Friedrichs form. For example
\[
 \mathcal D_y=H^1(\mathbb R,d\mu_y)
                  \cap\{u:y^2u\in\mathscr K\},\qquad
 h_q[u]=a_y\|u'\|^2+b_q\|yu\|^2+c\|y^2u\|^2.
\]
Its form domain is independent of $q$, since $b_q,c>0$.
The full form domain consists of $\Psi\in H^1(\mathbb R^2,d\mu_2)$
for which the original nonnegative multiplication potential in
\eqref{fq:H} has finite integral against $|\Psi|^2$.
Completeness follows by convergence in $H^1$ and after multiplication
by the square root of that potential; an almost-everywhere subsequence
identifies the limits. Smooth cutoffs and then compactly supported
mollification prove that $C_c^\infty$ is a form core. The operator
domain is the form-domain vectors whose displayed differential
expression belongs to $L^2$ distributionally, as follows by testing
against the core and using the Riesz representation of the shifted form.
Since $b_qy^2+cy^4\to\infty$ in $y$, the fibre form-domain inclusion
is compact: the potential controls the $L^2$ tail and the derivative
controls finite-cell averaging errors on bounded intervals.
The same argument proves compactness for $H_2$, using
$V_q(q)+cy^4\to\infty$ as $q^2+y^2\to\infty$.

We shall also use the following elementary operator-core fact.
For a smooth nonnegative real potential on $\mathbb R^n$ and a
constant positive kinetic matrix, $C_c^\infty$ is an operator core
for the Friedrichs operator. Indeed an $L^2$ distributional solution
of $(H^*+1)u=0$ is locally $H^2$ by local elliptic regularity.
Testing against $\chi_R^2u$, with a smooth compact cutoff whose
gradient is $O(R^{-1})$, gives
\[
 \int\bigl(|\nabla(\chi_Ru)|_D^2+(V+1)|\chi_Ru|^2\bigr)
 =\int|\nabla\chi_R|_D^2|u|^2=O(R^{-2})\|u\|^2.
\]
Letting $R\to\infty$ proves $u=0$. Thus the range of $H+1$ on
the minimal domain is dense; its lower bound makes the range of
its closure closed. That closure is surjective at $-1$ and symmetric,
hence self-adjoint. It is therefore the Friedrichs realization and
the stated domain is an operator core. This argument applies to
the fibres, the full operator, and the effective operator below.

\subsection{The actual normalized positive fibre ground state}

\begin{lemma}[Fibre energies and a proved quantitative separation]
\label{fq:gap}
For every $q$, $h(q)$ has a simple lowest eigenvalue $e(q)$ and a
unique real strictly positive eigenfunction $\Omega_q$ with
$\|\Omega_q\|_{\mathscr K}=1$. Write $e_1(q)$ for its second
eigenvalue, counting multiplicity. Then
\begin{equation}\label{fq:energy-bounds}
 \frac{\hbar\omega(q)}2\leq e(q)
 \leq\frac{\hbar\omega(q)}2
       +\frac{9\Gamma\hbar^2}{32\mathcal V_3\omega(q)^2},
 \qquad e_1(q)\geq\frac{3\hbar\omega(q)}2.
\end{equation}
In the explicit parameter regime
\begin{equation}\label{fq:regime}
 \frac{9\Gamma\hbar}{32\mathcal V_3\omega_*^3}\leq\frac12,
 \quad\hbox{equivalently}\quad
 \mathcal V_3\geq\frac{9\Gamma\hbar}{16\omega_*^3},
\end{equation}
these imply, simultaneously for every $q$,
\begin{equation}\label{fq:separation}
 e(q)\leq\hbar\omega(q),\qquad
 e_1(q)-e(q)\geq\lambda(q):=\frac{\hbar\omega(q)}2>0.
\end{equation}
\end{lemma}
\begin{proof}
Compactness gives existence of a bottom eigenfunction. The scalar
kinetic form does not increase on taking its modulus, so a
nonnegative real eigenfunction can be selected. Alternatively, the
continuous nonnegative potential has the Brownian-bridge kernel
\[
 \frac{\exp[-(y-z)^2/(4a_y\tau)]}
 {\sqrt{\mathcal V_3/2}\sqrt{4\pi a_y\tau}}
 \mathbb E\exp\left[-\int_0^\tau
    (b_qZ_s^2+cZ_s^4)\,ds\right]>0\qquad(\tau>0).
\]
Here $Z$ is the bridge from $y$ to $z$ with diffusion covariance
$2a_y$. Each path is continuous with bounded image, so its
exponential is strictly positive and bounded by one. The formula
follows first for bounded truncations by the convergent Duhamel
series and then by increasing-form resolvent convergence and
dominated convergence of the bridge exponential. Consequently the
semigroup is positivity improving. The modulus of any bottom
eigenfunction is again a bottom eigenfunction; if a real bottom
eigenfunction changed sign, positivity improvement on both its
positive and negative parts would imply
$|e^{-\tau h(q)}u|<e^{-\tau h(q)}|u|$, contradicting its eigenvalue
equations. Orthogonalization against a positive bottom eigenfunction
therefore rules out a second independent real one. Reality gives
complex dimension one. The eigenfunction is smooth by the
one-dimensional eigenvalue equation. A zero of a nonnegative smooth
solution would have derivative zero, and uniqueness for that
second-order ODE would force the solution to vanish. Thus the
selected eigenfunction is strictly positive everywhere.

Put $\beta(q)=\mathcal V_3\omega(q)/(2\hbar)$ and use the exactly
normalized harmonic comparison vector
\[
 \chi_q(y)=\bigl(\mathcal V_3/2\bigr)^{-1/4}
       \bigl(\beta(q)/\pi\bigr)^{1/4}
       \exp[-\beta(q)y^2/2],\qquad
 \|\chi_q\|_{\mathscr K}=1.
\]
The harmonic part $h_{\mathrm{osc}}(q)=-a_y\partial_y^2+b_qy^2$
factorizes as
\[
 h_{\mathrm{osc}}(q)
 =a_y(-\partial_y+\beta y)(\partial_y+\beta y)
       +\hbar\omega/2.
\]
Its ground state is $\chi_q$ and its first excitation has energy
$3\hbar\omega/2$. To check the latter without presuming a missing
spectral level, the harmonic operator has compact resolvent; its
lowering operator $\partial_y+\beta y$ maps any eigenfunction above
the ground to an $L^2$ distributional eigenfunction with energy
lowered by $\hbar\omega$. Its norm is nonzero by the factorization.
No energy strictly between $\hbar\omega/2$ and $3\hbar\omega/2$
is possible. The function $y\chi_q$ attains the latter energy.
The operator-core fact above identifies the distributional
eigenfunctions with the self-adjoint operator domain.

Since $cy^4\geq0$, the variational principle and its two-dimensional
min--max version give the lower bounds for $e$ and $e_1$.
Gaussian integration gives, in the retained measure,
\[
 \langle\chi_q,y^2\chi_q\rangle=\frac{\hbar}{\mathcal V_3\omega},
 \qquad
 \langle\chi_q,y^4\chi_q\rangle
       =\frac{3\hbar^2}{\mathcal V_3^2\omega^2}.
\]
Its Rayleigh quotient is the upper bound in
\eqref{fq:energy-bounds}. This is only a trial-vector bound:
$h(q)\chi_q=(\hbar\omega/2+cy^4)\chi_q$ is not proportional
to $\chi_q$ because $c>0$. In particular the true state
$\Omega_q$ is not replaced by this Gaussian.
Finally $\omega(q)\geq\omega_*$ and \eqref{fq:regime} make the
upper correction at most $\hbar\omega(q)/2$, proving
\eqref{fq:separation}.
\end{proof}

\begin{lemma}[An exact smooth eigenline, with its phase fixed]
\label{fq:bundle}
The functions $q\mapsto e(q)$ and $q\mapsto\Omega_q\in\mathscr K$
are locally real analytic, the latter also in the common local
operator graph norm. Thus they define a global normalized positive
ground-state line over the entire original $q$-axis. In this phase,
\[
 \langle\Omega_q,\partial_q\Omega_q\rangle=0,
 \qquad \mathcal A(q):=i\langle\Omega_q,\partial_q\Omega_q\rangle=0.
\]
Put $P_q=|\Omega_q\rangle\langle\Omega_q|$, $Q_q=I-P_q$, and
\[
 R_q=(h(q)-e(q))^{-1}Q_q,\qquad
 d:=\frac{3\Gamma\mathcal V_3}{2}.
\]
The inverse is on the orthogonal complement of the actual ground
state. The exact derivative identities are
\begin{align}
 e'(q)&=dq\,\langle\Omega_q,y^2\Omega_q\rangle,\notag\\
 \Omega'_q&=-dq\,R_qQ_qy^2\Omega_q,\label{fq:first-derivative}\\
 Q_q\Omega''_q
 &=-R_qQ_q\bigl[d y^2\Omega_q
                     +2(dq y^2-e'(q))\Omega'_q\bigr],\notag\\
 \langle\Omega_q,\Omega''_q\rangle&=-\|\Omega'_q\|^2.\notag
\end{align}
These identities hold without imposing \eqref{fq:regime}.
\end{lemma}
\begin{proof}
For $u$ in the operator domain of a fixed fibre $h(q_0)$,
\[
 c\|y^2u\|^2\leq\langle u,h(q_0)u\rangle
                  \leq\|u\|\,\|h(q_0)u\|.
\]
The arithmetic--geometric mean inequality implies that multiplication
by $y^2$ is $h(q_0)$-bounded with arbitrarily small relative bound.
Since $h(q)-h(q_0)=(b_q-b_{q_0})y^2$, nearby real fibres have the
same operator domain. One can see the corresponding resolvent
analyticity directly: at a fixed resolvent point $z$,
\[
 (h(q)-z)^{-1}=(h(q_0)-z)^{-1}
 \bigl[I+(b_q-b_{q_0})y^2(h(q_0)-z)^{-1}\bigr]^{-1},
\]
where the last inverse is a locally convergent Neumann series.
For real $q$, the closed symmetric perturbation is self-adjoint:
the same formula gives surjectivity at two conjugate nonreal
points chosen far enough from the real axis. It agrees with the
Friedrichs realization on the compact smooth operator core.
The displayed formula is analytic also as a map into the original
operator domain equipped with its graph norm.

A small contour around the isolated simple eigenvalue $e(q_0)$
defines the spectral projection by the resolvent integral.
It is locally analytic in that graph norm, remains of rank one
because close projections have isomorphic ranges, and gives the
locally analytic eigenvalue. For a fully fixed positive reference
vector use $\chi_*(y)=\chi_0(y)$, the preceding Gaussian at
$\omega_*$. Then the actual state is exactly
\begin{equation}\label{fq:exact-state}
 \Omega_q=\frac{P_q\chi_*}{\|P_q\chi_*\|_{\mathscr K}}.
\end{equation}
The denominator is nonzero because the ground state and $\chi_*$
are both strictly positive. This formula selects the same real
positive vector on overlaps of the local contour constructions,
so the local analytic sections give a global section.
It can equivalently be obtained by the convergent normalized
heat iteration $e^{-\tau h(q)}\chi_*/\|e^{-\tau h(q)}\chi_*\|$
as $\tau\to\infty$, by the discrete spectral expansion and the
nonzero ground-state overlap.

The real normalization identity gives
$\langle\Omega_q,\Omega'_q\rangle=0$. Differentiating the
eigenvalue equation in its common graph domain gives
$(h-e)\Omega'=-(dq y^2-e')\Omega$. Taking its ground-state
component proves the first identity, and taking its complementary
component proves the second. A second differentiation gives
$(h-e)\Omega''=-(dy^2-e'')\Omega-2(dq y^2-e')\Omega'$.
Its complementary part gives the asserted reduced-resolvent
formula; differentiating the normalization twice gives its
ground-state component. All products in these identities are
in $\mathscr K$, since $y^2$ is bounded in the local graph norm.
\end{proof}

\subsection{Explicit global derivative bounds for the true quartic state}

\begin{lemma}[Moment and reduced-resolvent estimates]
\label{fq:derivative-bounds}
Assume the proved parameter regime \eqref{fq:regime}. Then, globally
in $q$,
\begin{align}
 \|\Omega'_q\|&\leq M_1:=\frac{\sqrt{15\Gamma}}{\omega_*},
       \label{fq:M1}\\
 \|Q_q\Omega''_q\|&\leq M_2:=
 \frac{18\sqrt5\,\Gamma}{\omega_*^2}
 +\frac{192}{5^{3/4}}
        \sqrt{\frac{\Gamma\mathcal V_3}{\hbar\omega_*}}.
       \label{fq:M2}
\end{align}
\end{lemma}
\begin{proof}
In this proof all fibre inner products use the retained measure
$d\mu_y$. Let $\mu_j=\langle\Omega_q,y^j\Omega_q\rangle$ for
even $j$. The energy identity gives $b_q\mu_2\leq e(q)$.
To get the necessary sharper fourth moment, take the bounded
Lipschitz multiplier $F_R(y)=\max(-R,\min(y,R))$.
Testing the weak eigenvalue equation with $F_R^2\Omega_q$ gives
the exact identity
\[
 a_y\|(F_R\Omega_q)'\|^2
 +\int(b_qy^2+cy^4)F_R^2|\Omega_q|^2\,d\mu_y
 =e(q)\|F_R\Omega_q\|^2
          +a_y\int|F_R'|^2|\Omega_q|^2\,d\mu_y.
\]
The test lies in the form domain because $F_R$ and its derivative
are bounded. Since $|F_R|\leq|y|$, $|F_R'|\leq1$, dropping
nonnegative terms and letting $R\to\infty$ by Fatou yields
\[
 b_q\mu_4\leq e(q)\mu_2+a_y,
 \qquad \mu_4\leq\frac{e(q)^2}{b_q^2}+\frac{a_y}{b_q}.
\]
By \eqref{fq:separation}, this gives the explicit estimates
\begin{equation}\label{fq:moments}
 \mu_2\leq\frac{4\hbar}{\mathcal V_3\omega(q)},\qquad
 \|y^2\Omega_q\|=\sqrt{\mu_4}
       \leq\frac{2\sqrt5\,\hbar}{\mathcal V_3\omega(q)}.
\end{equation}
No Gaussian moment has been assigned to the quartic ground state.

From the proved gap, $\|R_q\|\leq1/\lambda(q)$. If $v\in\mathscr K$
and $w=R_qQ_qv$, the exact equation $(h-e)w=Q_qv$ and $w\perp\Omega_q$
give
\[
 c\|y^2w\|^2\leq\langle w,h(q)w\rangle
 =e(q)\|w\|^2+\operatorname{Re}\langle w,v\rangle
 \leq\frac{e(q)+\lambda(q)}{\lambda(q)^2}\|v\|^2.
\]
Thus \eqref{fq:first-derivative} implies the pointwise estimates
\begin{align}
 \|\Omega'_q\|&\leq
       L(q):=\frac{6\sqrt5\,\Gamma|q|}{\omega(q)^2},\notag\\
 \|y^2\Omega'_q\|&\leq
       L(q)\sqrt{\frac{e(q)+\lambda(q)}c}
       \leq4L(q)\sqrt{\frac{\hbar\omega(q)}{\Gamma\mathcal V_3}},
       \label{fq:pointwise-first}\\
 |e'(q)|&\leq\frac{6\Gamma\hbar|q|}{\omega(q)}.\notag
\end{align}
Substitution in the exact second-derivative identity gives
\begin{equation}\label{fq:pointwise-second}
 \|Q_q\Omega''_q\|\leq
 \frac{6\sqrt5\,\Gamma}{\omega(q)^2}
 +144\sqrt5\,\Gamma^{3/2}\sqrt{\frac{\mathcal V_3}{\hbar}}
                      \frac{q^2}{\omega(q)^{5/2}}
 +144\sqrt5\,\Gamma^2\frac{q^2}{\omega(q)^4}.
\end{equation}
For completeness the exact scalar maxima used to make these
bounds global are
\begin{align*}
 \sup_q\frac{|q|}{\omega_*^2+3\Gamma q^2}
       &=\frac1{2\sqrt{3\Gamma}\,\omega_*},\\
 \sup_q\frac{q^2}{(\omega_*^2+3\Gamma q^2)^2}
       &=\frac1{12\Gamma\omega_*^2},\\
 \sup_q\frac{q^2}{(\omega_*^2+3\Gamma q^2)^{5/4}}
       &=\frac4{3\Gamma\,5^{5/4}\sqrt{\omega_*}}.
\end{align*}
Differentiation in $|q|$ or $q^2$ shows that the maxima occur,
respectively, at $q^2=\omega_*^2/(3\Gamma)$,
$q^2=\omega_*^2/(3\Gamma)$, and $q^2=4\omega_*^2/(3\Gamma)$;
each function decreases to zero at infinity. Applying these
maxima to \eqref{fq:pointwise-first} and
\eqref{fq:pointwise-second} gives \eqref{fq:M1} and
\eqref{fq:M2}, including the constants $18\sqrt5$ and
$192/5^{3/4}$.
\end{proof}

\subsection{The exact projected Hamiltonian and its off-diagonal part}

\begin{theorem}[Exact isometric correspondence, not an invariant product]
\label{fq:projection}
For all of the retained positive parameters, without the extra
restriction \eqref{fq:regime}, define the isometry
\[
 J:\mathscr L\longrightarrow\mathscr H_2,\qquad
 (Jf)(q,y)=f(q)\Omega_q(y),\qquad
 (J^*\Psi)(q)=\langle\Omega_q,\Psi(q,\cdot)\rangle_{\mathscr K}.
\]
Its range projection is the fibrewise projection
$(\mathsf P\Psi)(q)=P_q\Psi(q)$, and $\mathsf Q=I-\mathsf P$.
The exact pullback of the full quadratic form is the Friedrichs
form of
\begin{equation}\label{fq:effective}
 H_{\mathrm{eff}}=-\frac{\hbar^2}{2\mathcal V_3}\partial_q^2
       +V_q(q)+e(q)+B(q),\qquad
 B(q)=\frac{\hbar^2}{2\mathcal V_3}\|\Omega'_q\|_{\mathscr K}^2.
\end{equation}
Its domain is precisely
\[
 \mathcal D_{\mathrm{eff}}=
 \left\{f\in H^1(\mathbb R,d\mu_q):
  \int(V_q+e+B)|f|^2\,d\mu_q<\infty\right\},
\]
and $J^{-1}(\mathcal D(H_2^{1/2})\cap\operatorname{Ran}J)
=\mathcal D_{\mathrm{eff}}$.
The Berry connection is zero in the fixed positive real phase,
but $B(0)=0$ and $B(q)>0$ for every $q\ne0$.

For every $f\in C_c^\infty(\mathbb R)$, $Jf\in D(H_2)$ and
the exact operator identity is
\begin{align}
 H_2Jf&=JH_{\mathrm{eff}}f+\mathcal Rf,\label{fq:intertwining}\\
 (\mathcal Rf)(q)&=-a\bigl(2f'(q)\Omega'_q
                                  +f(q)Q_q\Omega''_q\bigr),
 \qquad \mathsf P\mathcal Rf=0.\notag
\end{align}
Thus $J^*H_2J=H_{\mathrm{eff}}$ on this core, while
$\mathsf QH_2J=\mathcal R$ is the retained off-diagonal residual.
The range of $J$ is not a reducing subspace.

Under \eqref{fq:regime}, these identities extend to all
$f\in D(H_{\mathrm{eff}})$, and
\begin{equation}\label{fq:residual}
 \|\mathcal Rf\|\leq
  2\sqrt a\,M_1\|H_{\mathrm{eff}}^{1/2}f\|
        +aM_2\|f\|,
 \qquad 0\leq B(q)\leq\frac{15\Gamma\hbar^2}
                                 {2\mathcal V_3\omega_*^2}.
\end{equation}
\end{theorem}
\begin{proof}
Normalization of the actual fibre state and Fubini give
$\|Jf\|=\|f\|$ and the displayed adjoint and projection.
For compact smooth $f$,
\[
 \partial_q(Jf)=f'\Omega_q+f\Omega'_q,\qquad
 \|\partial_q(Jf)(q)\|^2=|f'|^2+|f|^2\|\Omega'_q\|^2,
\]
because $\langle\Omega_q,\Omega'_q\rangle=0$. The fibre energy
is exactly $e(q)|f(q)|^2$, not a harmonic replacement. These
identities prove equality of the pulled-back forms on the core.
Core approximation gives the forward domain inclusion. Conversely,
if $Jf$ lies in the full form domain, local Hilbert-valued Sobolev
differentiation gives
$f'=\langle\Omega_q,\partial_q(Jf)\rangle$ distributionally;
the term with $\Omega'_q$ vanishes by normalization. Hence $f\in H^1$,
and the displayed nonnegative derivative and fibre-energy identities
give every required integral in $\mathcal D_{\mathrm{eff}}$.
This proves equality of the domains and closed forms. The positive
smooth effective potential defines its self-adjoint Friedrichs
operator with the operator core proved above. It has compact
resolvent because it is bounded below by $V_q(q)$, which tends
to infinity, and its kinetic coefficient $a$ is strictly positive.

By local graph regularity of $\Omega_q$, $Jf$ is in the full form
domain and its distributional Hamiltonian is in $\mathscr H_2$.
It therefore lies in $D(H_2)$. Expanding $\partial_q^2(f\Omega_q)$
and using $\langle\Omega_q,\Omega''_q\rangle=-\|\Omega'_q\|^2$
proves \eqref{fq:intertwining} and the orthogonality of its residual.

At $q=0$, \eqref{fq:first-derivative} gives $\Omega'_0=0$.
If $q\ne0$ and $\Omega'_q=0$, that identity and injectivity of
the reduced inverse on $Q_q\mathscr K$ would give
$Q_qy^2\Omega_q=0$. It would follow that $y^2\Omega_q$ is a
scalar multiple of $\Omega_q$. Multiplication by $y^2$ has no
nonzero $L^2$ eigenfunction, because every level set has measure
zero. Thus $\Omega'_q\ne0$ for $q\ne0$ and $B(q)>0$ there.
If the range of $J$ were reducing, its residual would vanish for
every compact smooth $f$. Since its coefficients are continuous
in $q$, choose at any $q_0\ne0$ a test function with
$f(q_0)=0$, $f'(q_0)=1$. The residual at $q_0$ is then
$-2a\Omega'_{q_0}\ne0$, a contradiction.

Under \eqref{fq:regime}, the global derivative estimates and the
triangle inequality give
\[
 \|\mathcal Rf\|\leq a(2M_1\|f'\|+M_2\|f\|).
\]
All terms in the effective potential are nonnegative, so
$\sqrt a\|f'\|\leq\|H_{\mathrm{eff}}^{1/2}f\|$.
This proves \eqref{fq:residual} and defines $\mathcal R$ as a
bounded map from the effective form domain with its form norm
to $\mathscr H_2$. For $f\in D(H_{\mathrm{eff}})$, choose an
operator-core sequence. It converges also in form norm by the
spectral inequality
$\|H_{\mathrm{eff}}^{1/2}u\|^2\leq\|u\|\|H_{\mathrm{eff}}u\|$.
The right side of \eqref{fq:intertwining} then converges in
$\mathscr H_2$. Closedness of $H_2$ shows that
$Jf\in D(H_2)$ and proves the identity on the full effective
operator domain. The bound on $B$ follows directly from $aM_1^2$.
\end{proof}

\subsection{Controlled physical-time and spectral comparisons}

\begin{theorem}[Explicit finite-energy error in a proved parameter regime]
\label{fq:dynamics}
Assume \eqref{fq:regime}. Let $f\in\mathcal D_{\mathrm{eff}}$
have $\|f\|=1$ and effective energy expectation
$\|H_{\mathrm{eff}}^{1/2}f\|^2\leq E$. For every physical time
$t\in\mathbb R$,
\begin{align}
 &\left\|e^{-itH_2/\hbar}Jf
                      -Je^{-itH_{\mathrm{eff}}/\hbar}f\right\|
       \leq |t|\,\mathcal E(E),\label{fq:time-error}\\
 &\mathcal E(E):=
 \frac{\sqrt{30\Gamma E}}{\omega_*\sqrt{\mathcal V_3}}
 +\frac{9\sqrt5\,\Gamma\hbar}{\mathcal V_3\omega_*^2}
 +\frac{96}{5^{3/4}}
             \sqrt{\frac{\Gamma\hbar}{\mathcal V_3\omega_*}}.
       \notag
\end{align}
The same bound holds for the leakage norm
$\|\mathsf Qe^{-itH_2/\hbar}Jf\|$.
In particular, for any prescribed $T>0$, $0<\rho<1$, the entirely
explicit sufficient volume regime
\begin{equation}\label{fq:target-volume}
 \mathcal V_3\geq\max\left\{
 \frac{9\Gamma\hbar}{16\omega_*^3},\
 \frac{270\Gamma E T^2}{\rho^2\omega_*^2},\
 \frac{27\sqrt5\,\Gamma\hbar T}{\rho\omega_*^2},\
 \frac{288^2\Gamma\hbar T^2}{5^{3/2}\rho^2\omega_*}
 \right\}
\end{equation}
makes both norms at most $\rho$ simultaneously for $|t|\leq T$.
All compared Hamiltonians still contain the original quartic
interaction, and no coupling has been set to one.
\end{theorem}
\begin{proof}
First take $f\in D(H_{\mathrm{eff}})$. The established operator-domain
map permits differentiation of
$e^{-i(t-s)H_2/\hbar}Je^{-isH_{\mathrm{eff}}/\hbar}f$.
Integrating its derivative gives the exact Duhamel identity
\[
 e^{-itH_2/\hbar}Jf-Je^{-itH_{\mathrm{eff}}/\hbar}f
 =-\frac{i}{\hbar}\int_0^t
    e^{-i(t-s)H_2/\hbar}\mathcal R
                       e^{-isH_{\mathrm{eff}}/\hbar}f\,ds.
\]
Unitarity and preservation of the effective form norm give the bound
$|t|(2\sqrt aM_1\sqrt E+aM_2)/\hbar$. Substitution of the exact
$a,M_1,M_2$ yields \eqref{fq:time-error}. For a general form-domain
vector, effective spectral cutoffs converge both in norm and form
norm and give the same estimate by continuity. The comparison
vector on the right lies in $\operatorname{Ran}J$, so applying
the contraction $\mathsf Q$ proves the leakage bound.
The last three lower bounds on volume in \eqref{fq:target-volume}
make respectively the three terms of $T\mathcal E(E)$ at most
$\rho/3$; the first is the separately proved fibre-gap regime.
Squaring the first and third nonnegative inequalities and directly
solving the second gives exactly their displayed constants.
\end{proof}

\begin{corollary}[An explicit approximate-eigenstate relation]
\label{fq:spectral-error}
In \eqref{fq:regime}, every normalized effective eigenvector
$H_{\mathrm{eff}}f=\Lambda f$ satisfies
\[
 \|(H_2-\Lambda)Jf\|\leq
       2\sqrt aM_1\sqrt\Lambda+aM_2
       =\hbar\mathcal E(\Lambda),\qquad
 \operatorname{dist}(\Lambda,\sigma(H_2))
       \leq\hbar\mathcal E(\Lambda).
\]
\end{corollary}
\begin{proof}
The effective operator has compact resolvent as proved above, so
its eigenvectors are in its operator domain, and $\Lambda>0$
because $e(q)\geq\hbar\omega_*/2$. The exact intertwining identity
and residual estimate prove the norm bound. The spectral theorem
for $H_2$ gives
$\|(H_2-\Lambda)u\|\geq\operatorname{dist}(\Lambda,\sigma(H_2))\|u\|$
for every operator-domain vector. Apply it to the normalized $Jf$.
\end{proof}

The preceding comparison is an actual isometry, an exact compressed
Hamiltonian, and an explicit finite-time error estimate for this
specific interacting two-coordinate mechanics. At fixed
$\hbar,\Gamma,L_1,m^2,E,T$, its stated error tends to zero as
$\mathcal V_3\to\infty$, for any family of initial states satisfying
the stated energy bound. This is not a volume-uniform mass-gap
claim and not a continuum limit. The estimate is conservative;
in particular the second-derivative bound deliberately uses the
positive quartic energy to control $y^2R_q$, which is sufficient
for the proved $\mathcal V_3^{-1/2}$ error without an unproved
adiabatic expansion or omission of $y^4$.

\paragraph{Context and attribution.}
Moving spectral subspaces, the Berry connection, diagonal
Born--Huang energy, and errors in time-dependent Born--Oppenheimer
comparisons are established concepts; see G.~Panati, H.~Spohn,
and S.~Teufel, \emph{The time-dependent Born--Oppenheimer approximation},
ESAIM: M2AN \textbf{41} (2007), 297--314,
\url{https://arxiv.org/abs/0712.4369},
\url{https://doi.org/10.1051/m2an:2007023}.
No new general approximation theorem is claimed here. The spectral
bundle, every retained coefficient, and the conservative residual
and time estimates used above were proved directly for
\eqref{fq:H}. The bounded-below bridge formula used in the positivity
argument is also stated in B.~Simon, \emph{A Feynman--Kac Formula for
Unbounded Semigroups}, Theorem~1.1,
\url{https://arxiv.org/abs/math-ph/9907022}; its general unbounded-below
extension is not needed here.

% Input-safe canonical proof module.
% Host requirements: amsmath, amssymb, amsthm, mathtools, mathrsfs,
% array, booktabs, enumitem, hyperref, and theorem environments
% theorem, proposition, lemma, and counterexample.
\begingroup
\def\J{\mathcal J_2}
\def\V{\mathscr V}
\def\Psd{\mathcal P_2}
\def\Ppos{\mathcal P_2^{++}}
\def\Nfuture{\mathcal N_+^\times}
\def\rankC{\operatorname{rank}_{\mathbb C}}
\def\dd{\mathrm d}
\def\SpanR{\operatorname{span}_{\mathbb R}}

\section[Exact presymplectic reduction of the SU(3)-Gram worldline model]
{Exact presymplectic reduction of the\\
\texorpdfstring{$SU(3)$}{SU(3)}-Gram worldline model}
\label{wpr:sec:worldline-reduction}

The constrained Hamiltonian framework used below goes back to Dirac's
analysis of singular Lagrangians and constrained Hamiltonian systems
\cite{Dirac1950}.  The separation of constraint surface, characteristic
distribution, and reduced phase space is the finite-dimensional
presymplectic constraint problem treated systematically by
Gotay--Nester--Hinds \cite{GotayNesterHinds1978}; the regular quotient
step is the symplectic reduction mechanism of Marsden--Weinstein
\cite{MarsdenWeinstein1974}.  These references supply the general
framework only.  Every constraint rank, orbit, stabilizer, kernel,
one-form, two-form, quotient map, sign, equation of motion, and singular
counterexample used for this particular (SU(3))-Gram model is computed
explicitly below.

\subsection{Objects, constants, and equations}
Let
\[
 \J=\{A\in\operatorname{Mat}_{2\times2}(\mathbb C):A=A^*\},
\qquad
 \langle A,B\rangle_L=
 \frac12\bigl(\operatorname{tr}A\operatorname{tr}B-\operatorname{tr}(AB)\bigr).
\]
Then \(\langle A,A\rangle_L=\det A\), and the pairing is nondegenerate.
Write
\[
 \Psd=\{p\in\J:p\succeq0\},\quad
 \Ppos=\{p\in\J:p\succ0\},\quad
 \Nfuture=\{p\in\Psd:\rankC p=1\}.
\]
Retain
\[
 \V=\operatorname{Mat}_{3\times2}(\mathbb C),\quad
 \alpha>0,\quad c>0,\quad m_0\geq0,
\]
and define
\begin{align}
 P_\alpha(Z)&=\alpha Z^*Z,\label{wpr:eq:P}\\
 D P_{\alpha,Z}(\delta Z)&=
 \alpha(\delta Z^*Z+Z^*\delta Z),\label{wpr:eq:DP}\\
 \Phi_{m_0,\alpha,c}(Z)&=
 \frac12\left(\det P_\alpha(Z)-m_0^2c^2\right).\label{wpr:eq:Phi}
\end{align}
The history configuration space on a retained parameter interval \(I\) is
\[
 \mathfrak C_I=
 C^\infty(I,\J)\times C^\infty(I,\V)\times C^\infty(I,\mathbb R),
\]
with coordinates \((x,Z,e)\).  Because the Lagrangian is first order
in \(x\), contains no \(\dot Z\), and contains no \(\dot e\), the
corresponding instantaneous presymplectic phase carrier is
\[
 \mathcal E={\J}_x\times{\V}_Z\times\mathbb R_e.
\]
Define
\begin{equation}
 \theta=\langle P_\alpha(Z),\dd x\rangle_L,\qquad
 \Omega=-\dd\theta,\qquad
 \mathcal H=e\Phi_{m_0,\alpha,c}.
 \label{wpr:eq:forms}
\end{equation}
For tangent vectors \(v=(v_x,v_Z,v_e)\) and \(w=(w_x,w_Z,w_e)\),
\begin{equation}
 \Omega(v,w)=
 \langle v_x,D P_{\alpha,Z}(w_Z)\rangle_L
 -\langle w_x,D P_{\alpha,Z}(v_Z)\rangle_L.
 \label{wpr:eq:Omega}
\end{equation}

\begin{proposition}[Exact presymplectic Euler--Lagrange equation]
\label{wpr:prop:EL}
A path \((x,Z,e)\) is stationary under compactly supported variations
if and only if
\begin{align}
 \det P_\alpha(Z)&=m_0^2c^2,\label{wpr:eq:EL1}\\
 \frac{\dd}{\dd\tau}P_\alpha(Z)&=0,\label{wpr:eq:EL2}\\
 \left\langle D P_{\alpha,Z}(\delta Z),
 \dot x-eP_\alpha(Z)\right\rangle_L&=0
 \quad\text{for every }\delta Z\in\V.\label{wpr:eq:EL3}
\end{align}
Equivalently, \(\iota_{\dot\gamma}\Omega=\dd\mathcal H\).
\end{proposition}
\begin{proof}
Since
\[
 D(\det)_p(A)=2\langle p,A\rangle_L,
\]
one has
\[
 \dd\mathcal H(w)=
 \Phi_{m_0,\alpha,c}w_e+
 e\langle P_\alpha(Z),D P_{\alpha,Z}(w_Z)\rangle_L.
\]
Insert \eqref{wpr:eq:Omega} into
\(\iota_{\dot\gamma}\Omega=\dd\mathcal H\) and test successively
\(w_x,w_e,w_Z\).  This gives \eqref{wpr:eq:EL2},
\eqref{wpr:eq:EL1}, and \eqref{wpr:eq:EL3}.
\end{proof}

\subsection{Colour quotient and orbit types}
The left action is \(U\cdot(x,Z,e)=(x,UZ,e)\).

\begin{lemma}[Local colour is a zero-moment null symmetry]
\label{wpr:lem:null}
For \(\xi\in\mathfrak{su}(3)\), let \(\xi^\#=(0,\xi Z,0)\).  Then
\[
 D P_{\alpha,Z}(\xi Z)=0,\qquad
 \iota_{\xi^\#}\theta=0,\qquad
 \iota_{\xi^\#}\Omega=0.
\]
Thus the zero map to \(\mathfrak{su}(3)^*\) is a presymplectic moment
map.  The action is invariant under every \(\tau\)-dependent
\(U:I\to SU(3)\), with no connection term.
\end{lemma}
\begin{proof}
Because \(\xi^*=-\xi\),
\[
 D P_{\alpha,Z}(\xi Z)=
 \alpha(Z^*\xi^*Z+Z^*\xi Z)=0.
\]
The form identities follow from \eqref{wpr:eq:forms} and
\eqref{wpr:eq:Omega}.  Pointwise,
\((UZ)^*(UZ)=Z^*Z\), even when \(U\) depends on \(\tau\).
\end{proof}

\begin{theorem}[Gram quotient, fibres, and stabilizers]
\label{wpr:thm:quotient}
The map \(P_\alpha\) induces a homeomorphism
\[
 \V/SU(3)\xrightarrow{\ \sim\ }\Psd,\qquad [Z]\longmapsto P_\alpha(Z).
\]
Its orbit types are
\begin{center}
\begin{tabular}{c c c}
\toprule
\(\rankC Z\)&\text{fibre}&\(\operatorname{Stab}_{SU(3)}(Z)\)\\
\midrule
\(2\)&\(SU(3)\)&\(\{I_3\}\)\\
\(1\)&\(SU(3)/SU(2)\cong S^5\)&\(SU(2)\)\\
\(0\)&a point&\(SU(3)\)\\
\bottomrule
\end{tabular}
\end{center}
The rank-two quotient is smooth and free.  Rank one is a smooth
nonprincipal orbit-type stratum with nontrivial stabilizer; the zero
vertex is the singular apex of the total stratified quotient.
\end{theorem}
\begin{proof}
If \(Z_1^*Z_1=Z_2^*Z_2\), the map \(Z_1v\mapsto Z_2v\) is an isometry
between their images and extends to a unitary map of \(\mathbb C^3\).
Its action on the orthogonal complement can be adjusted to make its
determinant one.  In rank two the complement is one-dimensional and
the adjustment is unique, hence the stabilizer is trivial.  In rank
one the stabilizer is \(SU(2)\) on the two-dimensional complement.
At zero it is all of \(SU(3)\).  Surjectivity has the retained section
\[
 Z_p=\alpha^{-1/2}
 \begin{pmatrix}p^{1/2}\\0\quad0\end{pmatrix}.
\]
Compactness of \(SU(3)\) and continuity of the positive square root
give the quotient homeomorphism.  On \(\Ppos\) the section is smooth.
On \(\Nfuture\), the rank-one identity
\[
 p^{1/2}=\frac{p}{\sqrt{\operatorname{tr}p}}
\]
shows directly that the same section is smooth; together with the
constant orbit type \(SU(3)/SU(2)\), this gives the stated smooth
nonprincipal stratum.
\end{proof}

\begin{lemma}[Exact ranks of the Gram differential]
\label{wpr:lem:ranks}
\begin{enumerate}[label=\textup{(\roman*)}]
\item If \(\rankC Z=2\), then
\[
 \operatorname{im}D P_{\alpha,Z}=\J,\qquad
 \ker D P_{\alpha,Z}=T_Z(SU(3)\cdot Z),
\]
of real dimensions \(4\) and \(8\).
\item If \(\rankC Z=1\) and \(p=P_\alpha(Z)\), then
\[
 \operatorname{im}D P_{\alpha,Z}=p^{\perp_L}=T_p\Nfuture
\]
has real dimension \(3\).  The full ambient kernel has dimension \(9\),
but
\[
 \ker(D P_{\alpha,Z}|_{T_Z{\V}_1})=T_Z(SU(3)\cdot Z)
\]
has dimension \(5\), where \({\V}_1=\{Z:\rankC Z=1\}\).  The remaining
four real kernel directions are
\[
 \operatorname{Hom}_{\mathbb C}
 \bigl(\ker Z,(\operatorname{im}Z)^\perp\bigr);
\]
they raise rank and are not colour gauge.
\item At \(Z=0\), \(D P_{\alpha,0}=0\).
\end{enumerate}
\end{lemma}
\begin{proof}
For rank two and arbitrary \(A=A^*\), retain
\[
 \delta Z_A=\frac1{2\alpha}Z(Z^*Z)^{-1}A.
\]
Substitution into \eqref{wpr:eq:DP} gives
\(D P_{\alpha,Z}(\delta Z_A)=A\).  Hence the kernel has dimension
\(12-4=8\), equal to the free colour-orbit tangent.

For rank one write, without fixing either norm,
\[
 Z=ab^*,\qquad a\in\mathbb C^3\setminus\{0\},\quad
 b\in\mathbb C^2\setminus\{0\},\quad r=a^*a>0.
\]
Then \(p=\alpha rbb^*\).  If \(w=\delta Z^*a\), every
\(w\in\mathbb C^2\) occurs by taking \(\delta Z=aw^*/r\), and
\[
 \operatorname{im}D P_{\alpha,Z}
 =\{\alpha(wb^*+bw^*):w\in\mathbb C^2\}.
\]
The kernel of \(w\mapsto wb^*+bw^*\) is
\(\{itb:t\in\mathbb R\}\), so the image has dimension three.  It is
contained in \(p^{\perp_L}\), and equality follows by dimension.

Use the exact orthogonal splittings
\[
 \mathbb C^2=(\ker Z)^\perp\oplus\ker Z,\qquad
 \mathbb C^3=\operatorname{im}Z\oplus(\operatorname{im}Z)^\perp.
\]
Retaining the nonzero singular value \(\sigma>0\), adapted unit vectors
give
\[
 Z=\begin{pmatrix}\sigma&0\\0&0\\0&0\end{pmatrix},\qquad
 \delta Z=\begin{pmatrix}a_0&b_0\\c_1&d_1\\c_2&d_2\end{pmatrix},
\]
and
\[
 \alpha^{-1}D P_{\alpha,Z}(\delta Z)=
 \begin{pmatrix}
 2\sigma\operatorname{Re}a_0&\sigma b_0\\
 \sigma\overline{b_0}&0
 \end{pmatrix}.
\]
The ambient kernel has \(\operatorname{Re}a_0=b_0=0\) and dimension
nine.  The linearized rank-one equations also impose \(d_1=d_2=0\).
The five remaining directions are the \(SU(3)/SU(2)\)-orbit tangent;
the \(d_1,d_2\) block is the displayed complex Hom-space.  The zero
case follows from \eqref{wpr:eq:DP}.
\end{proof}

\subsection{Forms after colour reduction}
Define
\[
 \overline{\mathcal E}
 ={\J}_x\times{\Psd}_p\times\mathbb R_e
 \cong\mathcal E/SU(3),\qquad
 q(x,Z,e)=(x,P_\alpha(Z),e).
\]

\begin{proposition}[Reduced one-form and two-form]
\label{wpr:prop:reduced-forms}
The forms are basic for the colour action and satisfy
\[
 \theta=q^*\overline\theta,\qquad
 \Omega=q^*\overline\Omega,
\]
where
\begin{align}
 \overline\theta&=\langle p,\dd x\rangle_L,\label{wpr:eq:bartheta}\\
 \overline\Omega&=-\dd\overline\theta,\label{wpr:eq:barOmega}\\
 \overline\Omega((v_x,v_p),(w_x,w_p))
 &=\langle v_x,w_p\rangle_L-\langle w_x,v_p\rangle_L.
 \label{wpr:eq:barOmega-explicit}
\end{align}
After omitting the multiplier coordinate, this is symplectic on
\(\J\times\Ppos\).  On \(\J\times\Nfuture\), its kernel is exactly
\[
 \ker\overline\Omega_{(x,p)}=\SpanR\{(p,0)\}.
\]
On \(\J\times\{0\}\), the stratum-restricted two-form is zero.
\end{proposition}
\begin{proof}
Basicness follows from Lemma~\ref{wpr:lem:null}; the formulas follow by
differentiation.  The rank-two assertion uses that \(\Ppos\) is open in
\(\J\).  At nonzero rank-one \(p\), one has
\[
 T_p\Nfuture=p^{\perp_L}.
\]
If \((v_x,v_p)\) annihilates \eqref{wpr:eq:barOmega-explicit}, testing
arbitrary \(w_x\) gives \(v_p=0\), and testing
\(w_p\in p^{\perp_L}\) gives
\(v_x\in(p^{\perp_L})^{\perp_L}=\SpanR\{p\}\).
At \(p=0\), the zero stratum has no momentum tangent direction.
\end{proof}

\subsection{Positive mass}
For \(m_0>0\), define the retained future shell
\[
 \mathcal H^+_{m_0,c}
 =\{p\in\Ppos:\det p=m_0^2c^2\}.
\]
The parameter \(\alpha\) remains present in the exact quotient map
\(p=\alpha Z^*Z\); the displayed subset of momentum space itself
depends on \(m_0,c\).

\begin{theorem}[Exact massive reduction]
\label{wpr:thm:massive}
Let \(m_0>0\).  Then:
\begin{enumerate}[label=\textup{(\alph*)}]
\item The constraint surface
\[
 \Sigma_{m_0,\alpha,c}
 =\J\times\{Z:\det P_\alpha(Z)=m_0^2c^2\}\times\mathbb R_e
\]
is smooth and the colour action on it is free.
\item Its colour quotient is
\[
 \overline\Sigma_{m_0,c}
 =\J\times\mathcal H^+_{m_0,c}\times\mathbb R_e.
\]
Before colour reduction the kernel of the restricted two-form is the
direct sum of the eight-dimensional colour-orbit tangent,
\(\SpanR\{(p,0)\}\), and \(\SpanR\{\partial_e\}\).  After colour
reduction and omission of \(e\), only \(\SpanR\{(p,0)\}\) remains.
\item The equations are
\[
 \dot p=0,\qquad \dot x=e\,p,\qquad
 \det p=m_0^2c^2,
\]
while \(\dot e\) and motion in the colour fibre are arbitrary.
\item The physical phase space
\[
 \mathcal R_{m_0,c}
 =
 \frac{\J\times\mathcal H^+_{m_0,c}}
 {(x,p)\sim(x+s p,p),\ s\in\mathbb R}
\]
is canonically symplectomorphic to \(T^*\mathcal H^+_{m_0,c}\).
\item The colour-reduced positive-mass model is exactly the standard
future-directed spinless first-order relativistic particle of mass
\(m_0\) and speed constant \(c\).
\end{enumerate}
\end{theorem}
\begin{proof}
The constraint implies \(\rankC Z=2\).  By
Lemma~\ref{wpr:lem:ranks}, \(P_\alpha\) is a submersion there, and
\[
 D(\det)_p=2\langle p,\mathord{\cdot}\rangle_L\neq0.
\]
This proves smoothness.  Theorem~\ref{wpr:thm:quotient} gives freeness
and the colour quotient.  Since
\[
 T_p\mathcal H^+_{m_0,c}=p^{\perp_L},
\]
the proof of Proposition~\ref{wpr:prop:reduced-forms} gives the
characteristic line.  Pullback through the Gram submersion adds exactly
the colour-orbit kernel, and the multiplier adds \(\partial_e\).

Surjectivity of \(D P_{\alpha,Z}\) turns \eqref{wpr:eq:EL3} into
\(\dot x=e p\); the other equations are \eqref{wpr:eq:EL1} and
\eqref{wpr:eq:EL2}.  A path with constant \(p\) lies in one free Gram
fibre and differs locally from a constant lift by a
\(\tau\)-dependent colour transformation.

For the physical identification, if
\(A\in T_p\mathcal H^+_{m_0,c}\), define
\[
 \xi_{[x],p}(A)=-\langle x,A\rangle_L.
\]
This is unchanged by \(x\mapsto x+s p\).  Conversely, the annihilator
of \(p^{\perp_L}\) is \(\SpanR\{p\}\), so every cotangent covector has
one and only one class \([x]\).  Hence
\[
 ([x],p)\longmapsto(p,\xi_{[x],p})
\]
is a fibrewise linear bijection to \(T^*\mathcal H^+_{m_0,c}\).
The characteristic \(\mathbb R\)-action is free because every shell
momentum is nonzero.  It is proper: choose any Euclidean norm on
\(\J\).  For compact sets \(K_1,K_2\subset
\J\times\mathcal H^+_{m_0,c}\), the momenta occurring in them have a
positive minimum Euclidean norm \(\delta\); if
\((x,p)\in K_1\) and \((x+s p,p)\in K_2\), boundedness of the two
position projections gives \(|s|\delta\leq C(K_1,K_2)\).
Thus the orbit space is smooth.  The smooth vector-bundle map
\[
 (x,p)\longmapsto
 \bigl(p,\;A\mapsto-\langle x,A\rangle_L\bigr)
\]
is fibrewise onto with kernel \(\SpanR\{p\}\), and therefore descends
with a smooth inverse.  The displayed bijection is consequently a
diffeomorphism.

The basic one-form on the shell is
\begin{equation}
 \theta_{\mathrm{basic}}=-\langle x,\dd p\rangle_L.
 \label{wpr:eq:basic}
\end{equation}
It is the pullback of the tautological cotangent one-form, and
\begin{equation}
 \overline\theta
 =\dd\langle p,x\rangle_L+\theta_{\mathrm{basic}}.
 \label{wpr:eq:boundary}
\end{equation}
Therefore
\(\overline\Omega=-\dd\theta_{\mathrm{basic}}\), the canonical
cotangent symplectic form under the convention
\(\omega_{\mathrm{can}}=-\dd\lambda_{\mathrm{can}}\).
The original \(\overline\theta\) is not horizontal on the massive
characteristic because \(\overline\theta(p,0)=m_0^2c^2\); the exact
boundary term in \eqref{wpr:eq:boundary} produces the reduced
one-form.

Finally the reduced action is
\[
 \int_I\left(\langle p,\dot x\rangle_L
 -\frac e2(\det p-m_0^2c^2)\right)\dd\tau.
\]
Because \(\Ppos\) is open in \(\J\), \(p\) has every infinitesimal
Hermitian variation.  This is the standard covariant first-order
spinless particle action on the future-timelike sector, with no
residual internal symplectic factor.
\end{proof}

\begin{proposition}[Exact Hamiltonian gauge quasi-symmetry]
\label{wpr:prop:gauge}
For a smooth real \(\varepsilon(\tau)\), set
\[
 \delta_\varepsilon x=\varepsilon P_\alpha(Z),\qquad
 \delta_\varepsilon Z=0,\qquad
 \delta_\varepsilon e=\dot\varepsilon.
\]
Then
\begin{equation}
 \delta_\varepsilon L
 =
 \frac{\dd}{\dd\tau}
 \left[
 \frac{\varepsilon}{2}
 \bigl(\det P_\alpha(Z)+m_0^2c^2\bigr)
 \right].
 \label{wpr:eq:gauge-boundary}
\end{equation}
For compactly supported \(\varepsilon\) this is an exact gauge
quasi-symmetry, whose phase-space part generates the shell
characteristic where that characteristic is regular.
\end{proposition}
\begin{proof}
Put \(p=P_\alpha(Z)\) and \(q=\det p=\langle p,p\rangle_L\).  Since
\(\dot q=2\langle p,\dot p\rangle_L\),
\begin{align*}
 \delta_\varepsilon L
 &=
 \langle p,\dot\varepsilon p+\varepsilon\dot p\rangle_L
 -\frac{\dot\varepsilon}{2}(q-m_0^2c^2)\\
 &=
 \frac{\dot\varepsilon}{2}(q+m_0^2c^2)
 +\frac{\varepsilon}{2}\dot q,
\end{align*}
which is \eqref{wpr:eq:gauge-boundary}.
\end{proof}


\subsection{Massless rank one}
Set
\[
 \mathcal C_{0,\alpha,c}
 =\{Z:\det P_\alpha(Z)=0\}
 =\{Z:\rankC Z\leq1\}.
\]
The constant \(c>0\) is retained in the model, although it cancels
from the set-theoretic \(m_0=0\) equation.

\begin{lemma}[The pulled-back massless constraint is irregular]
\label{wpr:lem:irregular}
Writing the rows of \(Z\) as \(Z_{a\bullet}\),
\begin{equation}
 \det P_\alpha(Z)=
 \alpha^2\sum_{1\leq a<b\leq3}
 |Z_{a1}Z_{b2}-Z_{a2}Z_{b1}|^2.
 \label{wpr:eq:CB}
\end{equation}
Consequently,
\[
 \dd(\det P_\alpha)_Z=0
 \quad\text{for every }Z\in\mathcal C_{0,\alpha,c}.
\]
The single real multiplier equation describes the rank-at-most-one
set, but its differential supplies none of the four real normal
equations on the smooth rank-one stratum.
\end{lemma}
\begin{proof}
Equation \eqref{wpr:eq:CB} is Cauchy--Binet with the factor
\(\alpha^2\) retained.  On its zero set every complex minor vanishes,
so differentiation of the sum of squared absolute values gives zero.
Equivalently, at rank one,
\[
 \dd(\det P_\alpha)_Z(\delta Z)
 =2\langle p,D P_{\alpha,Z}(\delta Z)\rangle_L=0
\]
by Lemma~\ref{wpr:lem:ranks}; at \(Z=0\) use
\(D P_{\alpha,0}=0\).
\end{proof}

\begin{theorem}[Nonzero massless equations and physical quotient]
\label{wpr:thm:massless}
Suppose \(m_0=0\) and \(\rankC Z=1\), and put
\(p=P_\alpha(Z)\in\Nfuture\).  Then:
\begin{enumerate}[label=\textup{(\alph*)}]
\item The colour fibre is \(SU(3)/SU(2)\), with stabilizer \(SU(2)\).
\item The complete equations are
\begin{equation}
 \dot p=0,\qquad
 \dot x=\lambda(\tau)p
 \quad\text{for some real }\lambda(\tau),\qquad
 e(\tau)\text{ arbitrary}.
 \label{wpr:eq:null-EL}
\end{equation}
No equation identifies \(\lambda\) with the written multiplier \(e\).
\item On the intrinsic smooth rank-one shell, the kernel before
colour reduction is the direct sum of
\[
 T_Z(SU(3)\cdot Z),\qquad
 \SpanR\{(p,0)\},\qquad
 \SpanR\{\partial_e\}.
\]
After colour reduction and omission of \(e\), only the characteristic
\(\SpanR\{(p,0)\}\) remains.
\item The physical phase space
\[
 \mathcal R_0^\times
 =
 \frac{\J\times\Nfuture}
 {(x,p)\sim(x+s p,p)}
\]
is canonically symplectomorphic to \(T^*\Nfuture\), with reduced
one-form \(-\langle x,\dd p\rangle_L\) and its canonical two-form.
No quotient by positive rescaling of \(p\) occurs.
\item Thus the unparameterized spacetime curves and physical
symplectic quotient agree with the standard future-directed nonzero
spinless massless particle.  The original triples \((x,Z,e)\), their
off-shell neighbourhood, and their multiplier equations are not
exactly the standard massless system.
\end{enumerate}
\end{theorem}
\begin{proof}
The fibre statement is Theorem~\ref{wpr:thm:quotient}.  By
Lemma~\ref{wpr:lem:ranks}, equation \eqref{wpr:eq:EL3} says
\[
 \dot x-e p\in(p^{\perp_L})^{\perp_L}=\SpanR\{p\}.
\]
Since \(ep\) lies in this line, this is equivalent to
\(\dot x=\lambda p\), with no relation between \(\lambda\) and \(e\).
Equation \eqref{wpr:eq:EL2} gives \(\dot p=0\).  A path in a fixed
rank-one Gram fibre differs locally from a constant representative by
a \(\tau\)-dependent \(SU(3)\) transformation; the \(SU(2)\)
stabilizer makes this gauge description reducible.

On the intrinsic rank-one manifold,
Lemma~\ref{wpr:lem:ranks} identifies the vertical kernel with the
five colour directions.  Proposition~\ref{wpr:prop:reduced-forms}
adds the one spacetime characteristic, and \(e\) adds
\(\partial_e\).

Since \(T_p\Nfuture=p^{\perp_L}\), the massive cotangent construction
applies without division by \(\det p\):
\[
 ([x],p)\longmapsto
 \left(p,\;A\longmapsto-\langle x,A\rangle_L\right).
\]
Every \(p\in\Nfuture\) is nonzero.  The freeness, properness, and smooth
vector-bundle argument in the massive proof applies to compact subsets
of \(\J\times\Nfuture\) verbatim, and proves that this map is a
diffeomorphism to \(T^*\Nfuture\).  Moreover,
\eqref{wpr:eq:boundary} identifies the canonical form.  Reparametrizing
the worldline changes \(x\) along \(p\); it does not identify distinct
positive magnitudes of \(p\).  Lemma~\ref{wpr:lem:irregular} and the
absence of \(\lambda=e\) establish the final inequivalence statement.
\end{proof}

\begin{counterexample}[Rank-one multiplier mismatch]
\label{wpr:ce:rank1}
Retain arbitrary \(s>0\), and set
\[
 Z_s=\begin{pmatrix}s&0\\0&0\\0&0\end{pmatrix},\qquad
 p_s=\alpha\begin{pmatrix}s^2&0\\0&0\end{pmatrix}.
\]
The path
\[
 Z(\tau)=Z_s,\qquad x(\tau)=\tau p_s,\qquad e(\tau)=0
\]
satisfies all equations in Proposition~\ref{wpr:prop:EL} for
\(m_0=0\).  The standard action with momentum independently varied in
\(\J\) and the same \(e=0\) would require \(\dot x=e p_s=0\).

Under Proposition~\ref{wpr:prop:gauge}, a solution
\(\dot x=\lambda p\), \(\dot p=0\), transforms by
\[
 \lambda\longmapsto\lambda+\dot\varepsilon,\qquad
 e\longmapsto e+\dot\varepsilon.
\]
Thus \(\lambda-e\) is invariant under the stated Hamiltonian gauge;
the discrepancy remains when the original multiplier is retained.
\end{counterexample}

\begin{counterexample}[Four non-colour first-order null directions]
\label{wpr:ce:transverse}
Retain arbitrary \(s>0\) and \(t\in\mathbb C\setminus\{0\}\), and put
\[
 Z_s=\begin{pmatrix}s&0\\0&0\\0&0\end{pmatrix},\qquad
 \Delta_t=\begin{pmatrix}0&0\\0&t\\0&0\end{pmatrix}.
\]
Then
\begin{equation}
 D P_{\alpha,Z_s}(\Delta_t)=0,\qquad
 \det P_\alpha(Z_s+\varepsilon\Delta_t)
 =\alpha^2s^2|t|^2\varepsilon^2
 \quad(\varepsilon\in\mathbb R).
 \label{wpr:eq:transverse}
\end{equation}
The vector \(\Delta_t\) is not tangent to a colour orbit, since every
\(\xi Z_s\) has zero second column.  It is not tangent to the rank-one
variety, because the perturbed matrix has rank two for
\(\varepsilon\neq0\).  Allowing both complex coordinates in
\((\operatorname{im}Z_s)^\perp\) gives the four real transverse
directions of Lemma~\ref{wpr:lem:ranks}.  They are second-order
shell-opening directions, not gauge.
\end{counterexample}

\subsection{The zero-momentum vertex}
\begin{theorem}[Additional singular branch]
\label{wpr:thm:zero}
Let \(m_0=0\) and \(Z=0\).  Then \(p=0\), the stabilizer is all of
\(SU(3)\), and both \(\theta\) and \(\Omega\) vanish on
\(\J\times\{0\}\times\mathbb R_e\).  Every pair of paths
\[
 x:I\to\J,\qquad e:I\to\mathbb R,
\]
with \(Z(\tau)\equiv0\), satisfies the Euler--Lagrange equations.

The phase part of the Hamiltonian gauge vanishes at \(p=0\), while the
presymplectic kernel contains every \(x\)-direction.  Quotienting by
the actual colour and Hamiltonian-gauge orbits therefore leaves the
whole \(\J\) branch, whereas quotienting by all characteristic leaves
would collapse it to one point.  This is an exact failure of regular
Dirac reduction at the vertex.
\end{theorem}
\begin{proof}
At \(Z=0\), one has \(p=0\) and \(D P_{\alpha,0}=0\).  Every equation
in Proposition~\ref{wpr:prop:EL} is identically satisfied, independent
of \(\dot x,e,\dot e\).  The form and gauge statements follow from
\eqref{wpr:eq:forms} and
\(\delta_\varepsilon x=\varepsilon p=0\).
\end{proof}

\begin{counterexample}[Stationary zero momentum with nonconstant
position]
\label{wpr:ce:zero}
The retained paths
\[
 x(\tau)=\begin{pmatrix}\tau^2&0\\0&0\end{pmatrix},\qquad
 Z(\tau)=0,\qquad e(\tau)=\sin\tau
\]
are stationary for the present \(m_0=0\) action.  The standard
independent-momentum action gives \(\dot x=e p=0\) at \(p=0\), and
therefore excludes this path.  Removing \(p=0\), as is conventional
for physical massless phase space, removes this vertex counterexample
but does not regularize the off-shell rank-one constraint.
\end{counterexample}

\subsection{Kernel dimensions and exact verdict}
\begin{proposition}[Every kernel jump]
\label{wpr:prop:kernels}
Omit the always-null multiplier direction for the table.  On the
ambient \(16\)-real-dimensional \({\J}_x\times{\V}_Z\),
\begin{center}
\begin{tabular}{c c c c}
\toprule
\(\rankC Z\)&\(\operatorname{rank}_{\mathbb R}D P_{\alpha,Z}\)&
\(\operatorname{rank}_{\mathbb R}\Omega\)&
\(\dim_{\mathbb R}\ker\Omega\)\\
\midrule
\(2\)&\(4\)&\(8\)&\(8\)\\
\(1\)&\(3\)&\(6\)&\(10\)\\
\(0\)&\(0\)&\(0\)&\(16\)\\
\bottomrule
\end{tabular}
\end{center}
On the \(12\)-real-dimensional intrinsic rank-one stratum
\({\J}_x\times{\V}_1\), the pulled-back form has rank \(6\) and nullity
\(6\): five colour directions and one spacetime characteristic.
\end{proposition}
\begin{proof}
Using the Lorentz pairing to identify \(\J\cong\J^*\), the matrix of
\(\Omega\) is
\[
 \begin{pmatrix}
 0&G_L D P_{\alpha,Z}\\
 -(D P_{\alpha,Z})^{\mathsf T}G_L&0
 \end{pmatrix}.
\]
Its rank is twice the rank of \(D P_{\alpha,Z}\).  Apply
Lemma~\ref{wpr:lem:ranks} in the ambient space and on \(T_Z{\V}_1\).
\end{proof}

\begin{theorem}[Final comparison]
\label{wpr:thm:verdict}
The strongest correct result is:
\begin{enumerate}[label=\textup{(\arabic*)}]
\item for \(m_0>0\), colour reduction gives exactly the standard
future-timelike first-order spinless particle, including its equations,
constraint characteristic, and physical cotangent symplectic space;
\item for \(m_0=0\) and \(p\neq0\), colour and characteristic reduction
gives exactly the standard future-null physical phase space and
unparameterized null trajectories;
\item the full \(m_0=0\) pulled-back multiplier action is not the
standard massless variational system: its constraint has zero
differential, its rank-one equation is \(\dot x=\lambda p\) rather
than \(\dot x=e p\), its ambient kernel contains four non-colour
transverse directions, and its vertex admits arbitrary \(x(\tau)\).
\end{enumerate}
The exact relating morphism is \(P_\alpha(Z)=\alpha Z^*Z\); the exact
failure locus is the singular differential of
\(\det\circ P_\alpha\) on the rank-at-most-one variety.
\end{theorem}
\begin{proof}
Combine Theorem~\ref{wpr:thm:massive},
Theorem~\ref{wpr:thm:massless},
Lemma~\ref{wpr:lem:irregular}, and
Counterexamples~\ref{wpr:ce:rank1}--\ref{wpr:ce:zero}.
\end{proof}

\subsection{Deterministic verification companion}
The companion program
\texttt{verify\_worldline\_reduction.py} checks the determinant
polarization, the rank-two Gram right inverse, all Gram and
presymplectic ranks, all colour-orbit ranks, the intrinsic rank-one
tangent calculation, the zero gradient and quadratic transverse
opening of the massless constraint, the cotangent-form sign, and
\eqref{wpr:eq:gauge-boundary}, using exact symbolic arithmetic.

\subsection{Comparison with the extended exact Lorentz packet}
\label{wpr:app:packet}

This subsection records a bounded line-specific audit of source lines
875--1654.  The audited source file is
\begin{center}
\small\path{exact_lorentz_observation_bridge.tex}.
\end{center}
It is stored beneath
\path{supporting_materials/workbench/evidence} in the directory
\begin{center}
\small\path{worldline_presymplectic_reduction_20260904}.
\end{center}
Its file SHA-256 is
\[
\texttt{2BE16E1DC71B2E5DC04238F0FBE440138258D0825A4841AB2F29662587681599}.
\]
That packet adds a third column \(\phi\in\mathbb C^3\), writes
\(B=[Z\ \phi]\), and uses
\[
 G=Z^*Z,\quad z=Z^*\phi,\quad
 \lambda=\phi^*\phi,\quad \omega=\det B.
\]
Its mass convention is \(m^2_{\mathrm{packet}}=\det p\), with physical
mass \(M_{\mathrm{phys}}=\sqrt{\det p}/c\).  The convention in the
audited \(S^6\) module is
\(m_{\alpha,c}^2=\det p/c^2\).  The exact dictionary is therefore
\[
 m^2_{\mathrm{packet}}=c^2m_{\alpha,c}^2,\qquad
 M_{\mathrm{phys}}=m_{\alpha,c};
\]
there is no mathematical disagreement.

\subsubsection{Claim ledger}

\begin{center}
\begin{tabular}{>{\raggedright\arraybackslash}p{0.15\textwidth}
                >{\raggedright\arraybackslash}p{0.20\textwidth}
                >{\raggedright\arraybackslash}p{0.55\textwidth}}
\toprule
Source lines&Verdict&Exact scope\\
\midrule
875--950&validated&
\(p=\alpha Z^*Z\) is a proper surjection, its fibres are the
\(SU(3)\)-orbits, and
\(\det p=\alpha^2\sum_{a<b}|\Delta_{ab}(Z)|^2\) has range
\([0,\infty)\).  The zero fibre is rank at most one.\\
\addlinespace
952--1027&validated as a multiplication-operator model&
With the stated pushforward measure, the observer-energy and
mass-squared multiplication operators have essential ranges
\([0,\infty)\), no point spectrum, and no positive gap after a vacuum
is adjoined.  This is a kinematic/free spectral statement, not a
statement about an interacting field Hamiltonian.\\
\addlinespace
1029--1476&validated&
For \(\lambda_0>0\), \(\omega_0\neq0\), the constrained quotient,
associated-ball description, momentum fibres, and sharp lower bound
\(\det p\geq\alpha^2|\omega_0|^2/\lambda_0\) are algebraic and
differential-topological statements.  Since \(\omega_0\neq0\) forces
\(B\) invertible and \(Z\) rank two, the massless rank-one defect
proved above does not invalidate these claims.\\
\addlinespace
1482--1531&valid invariance and constraints, incomplete dynamics&
The displayed action is invariant and its multiplier and \(x\)
variations give the stated constraints and \(\dot p=0\).  Its omitted
\((Z,\phi)\) equations materially change the dynamical interpretation;
Theorem~\ref{wpr:thm:packet-action} below computes them intrinsically.\\
\addlinespace
1533--1603&validated with scope qualification&
The multiplication spectra and the direct-sum/Fock vacuum gap equal
the displayed threshold.  The threshold is inserted by fixing
\(\lambda=\lambda_0\) and \(\omega=\omega_0\neq0\); it is not a
dynamically generated interacting mass gap and is not a Yang--Mills
mass-gap theorem.\\
\addlinespace
1605--1633&validated and scope-sharpening&
The two explicit degenerating sequences prove that controlling only
one of \(\lambda\) and \(\omega\) gives no positive threshold.  Thus
the positive lower endpoint is exactly constraint-selected.\\
\addlinespace
1635--1653&validated&
The diagonal boost excludes a positive-definite quadratic form on the
spinor carrier invariant under all of \(SL_2(\mathbb C)\).\\
\bottomrule
\end{tabular}
\end{center}

The independent replay and line-by-line claim ledger for this packet are
retained under the project-relative path
\path{supporting_materials/workbench/evidence/}\allowbreak
\path{exact_lorentz_observation_bridge_audit_20260904}.
The packet checker passes \(33/33\) assertions and the separately written
checker passes \(30/30\).  The audit-report SHA--256 is
\texttt{c8d01af28836340798c2fa7e27f8e097\allowbreak%
d67e0b8d7c13f75271af473958d391df}; the promotion manifest, which verifies
all retained audit artifacts, has SHA--256
\par\noindent
\texttt{52f53746558fa219260b224abdb71c86d2cd6c13192774d0d2fb1a1ed71bcf35}.

\subsubsection{The full dynamics of the packet's constrained action}

Retain the packet's constants
\[
 \lambda_0>0,\qquad\omega_0\in\mathbb C^*,\qquad\alpha>0
\]
and its smooth colour quotient
\[
 \mathcal Q_{\lambda_0,\omega_0}
 \cong
 \operatorname{Herm}_2^0(\mathbb C)
 \times\{w\in\mathbb C^2:\|w\|^2<\lambda_0\}.
\]
Put
\begin{equation}
 \mu(w)=\frac{|\omega_0|}
 {\sqrt{\lambda_0-\|w\|^2}},\qquad
 G(A,w)=\mu(w)e^A,\qquad
 p(A,w)=\alpha G(A,w).
 \label{wpr:eq:packet-p}
\end{equation}
Eliminating the three regular multiplier constraints from source lines
1484--1495 gives the intrinsic action
\begin{equation}
 S_{\mathrm{packet}}[x,A,w]
 =\int_I\langle p(A,w),\dot x\rangle_L\,\dd\tau
 \label{wpr:eq:packet-action}
\end{equation}
on \({\J}_x\times\mathcal Q_{\lambda_0,\omega_0}\).  The use of the
intrinsic constrained action loses no Euler--Lagrange equation:
the constraint map is a submersion on \(\omega_0\neq0\), so the
Lagrange-multiplier theorem identifies ambient stationarity with
stationarity against every tangent vector of the constrained
manifold.  This is a statement about classical stationary paths; it
does not assert equality of quantum measures or path integrals after
formally integrating out the multipliers.

\begin{theorem}[Complete reduced dynamics of the constrained packet]
\label{wpr:thm:packet-action}
For the momentum projection in \eqref{wpr:eq:packet-p}:
\begin{enumerate}[label=\textup{(\roman*)}]
\item If \(w\neq0\), then \(D p_{(A,w)}\) has real rank \(4\), and its
kernel has dimension \(3\), tangent to the \(S^3\) momentum fibre.
The Euler--Lagrange equations of \eqref{wpr:eq:packet-action} are
\[
 \dot p=0,\qquad \dot x=0,
\]
with motion along the fixed \(S^3\) fibre presymplectically null.
\item If \(w=0\), then
\[
 \det p=\frac{\alpha^2|\omega_0|^2}{\lambda_0},
\]
\(D p_{(A,0)}\) has rank \(3\), its image is
\(p^{\perp_L}\), and its kernel consists of all four real
\(w\)-directions.  The pointwise equations are
\[
 \dot p=0,\qquad \dot x=\lambda(\tau)p.
\]
Every actual solution remaining on the threshold has constant
\((A,w)\), because the momentum fibre over a threshold momentum is a
single point.
\item Consequently, the source action is not a family of ordinary
massive particle actions over all masses above the lower endpoint.
Only its threshold stratum has the ordinary worldline characteristic;
strict-interior configurations have \(\dot x=0\).  The observer
multiplication operator of source lines 1533--1603 is a separate
positive-energy representation and is not the Hamiltonian generating
\(\tau\)-evolution in \eqref{wpr:eq:packet-action}.
\end{enumerate}
\end{theorem}

\begin{proof}
For fixed \(w\), varying \(A\) varies \(p\) through the
three-dimensional fixed-determinant future hyperboloid, so the image
of the \(A\)-differential is \(p^{\perp_L}\).  Direct differentiation
without suppressing any constant gives
\begin{equation}
 D\mu_w(\delta w)
 =
 \frac{|\omega_0|\operatorname{Re}\langle w,\delta w\rangle}
 {(\lambda_0-\|w\|^2)^{3/2}}.
 \label{wpr:eq:Dmu}
\end{equation}
When \(w\neq0\), a radial \(\delta w\) makes
\eqref{wpr:eq:Dmu} nonzero and changes \(\det p=\alpha^2\mu(w)^2\).
This supplies a fourth direction transverse to \(p^{\perp_L}\).
The remaining \(w\)-directions satisfy
\(\operatorname{Re}\langle w,\delta w\rangle=0\), form a
three-dimensional space, and are exactly tangent to the \(S^3\)
momentum fibre.

At \(w=0\), equation \eqref{wpr:eq:Dmu} vanishes in every direction,
so the four \(w\)-directions lie in \(\ker Dp\) and the image is only
\(p^{\perp_L}\).  These directions open the mass above its minimum at
second order:
\[
 \det p(A,\varepsilon v)
 =
 \frac{\alpha^2|\omega_0|^2}
 {\lambda_0-\varepsilon^2\|v\|^2}.
\]

For an action \(\int\langle p(q),\dot x\rangle_L\dd\tau\), variation
of \(x\) gives \(\dot p=0\), while variation of \(q\) gives
\[
 \langle Dp_q(\delta q),\dot x\rangle_L=0
 \quad\text{for all }\delta q.
\]
If \(w\neq0\), the image is all of \(\J\), hence \(\dot x=0\).
If \(w=0\), the image is \(p^{\perp_L}\), hence
\(\dot x\in\SpanR\{p\}\).  The packet's own momentum-fibre theorem
states that the fibre over a threshold momentum is a singleton, which
completes the solution statement.
\end{proof}

\subsubsection{Relation to the rank-one irregularity}

The packet's fixed-\(\omega_0\neq0\) constrained manifold never meets
\(\rankC Z\leq1\), so the zero differential of
\(\det\circ P_\alpha\) does not invalidate its topology or positive
lower bound.  The exact relation is instead:
\begin{enumerate}
\item on the unrestricted carrier, the rank-one and zero strata give
the massless irregularities proved in
Theorems~\ref{wpr:thm:massless} and \ref{wpr:thm:zero};
\item on the packet's constrained carrier, \(Z\) remains rank two, but
the momentum projection develops a different four-real-dimensional
critical kernel at its positive-mass threshold \(w=0\);
\item the algebraic inequality is a sharp, externally
constraint-selected threshold; a dynamically generated gap would
require a specified interacting Hamiltonian or field operator whose
spectrum produces the scale rather than receiving
\(\lambda_0,\omega_0\) as fixed input.
\end{enumerate}

\endgroup

\section{The Hermitian determinant normal slice as a Lorentzian host}
\label{hnsp-section}

This independently extractable section restates one already verified
finite-dimensional identity and constructs exact physical models from it.
When it is placed after the nine-dimensional observation construction, the
notation agrees by
\(C_{A_2}=C\), \(R_P=R\), \(H_0=uu^*\),
\(Q_\perp=Q_u\), and the same map
\(\jmath_P(G)=R_PGR_P^*\).  The positive conclusion is not
merely that the identity resembles Minkowski space: the affine normal slice
has a canonically polarized Lorentz metric, the displayed \(SL_2(\mathbb C)\)
congruence subgroup induces the spin double cover of the proper orthochronous
Lorentz group, and after one explicit choice of length scale it supports a completely
specified Klein--Gordon theory and finite Gaussian regulators.  The same host
also supports classical Yang--Mills theory for any separately chosen compact
gauge group.  In addition, the standard metrics entering the Moore--Penrose
section select one exact compact \(SU(2)\)-type representation internal to the
normal data; promoting its pointwise action to a local gauge symmetry gives
the gauge--scalar model below and is an additional construction.  The proofs
identify exactly which structures come from the determinant and which are
additional physical choices.

\subsection{Fixed matrices and the exact normal metric}

Put
\begin{equation}
 u=\begin{pmatrix}1\\1\\1\end{pmatrix},\qquad
 P=\begin{pmatrix}1&-1&0\\0&1&-1\end{pmatrix},\qquad
 C_{A_2}=PP^*=\begin{pmatrix}2&-1\\-1&2\end{pmatrix},
 \label{hnsp-fixed-data}
\end{equation}
and use the standard Hermitian products in these displayed bases.  The
standard-metric Moore--Penrose right inverse is
\begin{equation}
 R_P=P^*(PP^*)^{-1}
 =\frac13\begin{pmatrix}2&1\\-1&1\\-1&-2\end{pmatrix}.
 \label{hnsp-moore-penrose}
\end{equation}
The denominators in \eqref{hnsp-moore-penrose} are essential: this is a
metric-selected complex-linear section, not a homomorphism splitting the
displayed integral lattices.
Set
\begin{equation}
 H_0=uu^*,\qquad
 \mathcal N_P=\{X\in\operatorname{Herm}_3(\mathbb C):Xu=0\},\qquad
 \jmath_P(G)=R_P G R_P^*.
 \label{hnsp-normal-space}
\end{equation}

\begin{proposition}[Exact identification of the normal slice]
\label{hnsp-normal-identification}
The following identities hold without a change of normalization:
\begin{equation}
 PR_P=I_2,\qquad u^*R_P=0,\qquad
 R_P^*R_P=C_{A_2}^{-1},\qquad
 R_P P=I_3-\frac13uu^*.
 \label{hnsp-right-inverse-identities}
\end{equation}
The real-linear map
\(
 \jmath_P:\operatorname{Herm}_2(\mathbb C)\to\mathcal N_P
\)
is an isomorphism.  Its inverse is
\begin{equation}
 \jmath_P^{-1}(X)=PXP^*.
 \label{hnsp-normal-inverse}
\end{equation}
For all \(G\in\operatorname{Herm}_2(\mathbb C)\) and
\(\varepsilon\in\mathbb R\),
\begin{equation}
 \det\bigl(H_0+\varepsilon\jmath_P(G)\bigr)
 =\varepsilon^2\det G.
 \label{hnsp-determinant-identity}
\end{equation}
\end{proposition}

\begin{proof}
The first three equalities in
\eqref{hnsp-right-inverse-identities} follow by direct multiplication.  The
last member is the orthogonal projection onto the row space of \(P\), which is
\(u^\perp\); direct multiplication gives the displayed matrix as well.  Hence
\(P\jmath_P(G)P^*=G\).  Conversely, if \(X=X^*\) and \(Xu=0\), then also
\(u^*X=0\), and therefore
\[
 R_P(PXP^*)R_P^*=(R_P P)X(R_P P)=X.
\]
This proves the isomorphism and its inverse.

For the determinant, put \(v_0=u/\sqrt3\), choose an orthonormal complement
\(U_\perp\), and let \(U=(v_0\ U_\perp)\) and
\(L=U_\perp^*R_P\).  Then
\[
 U^*H_0U=\operatorname{diag}(3,0,0),\qquad
 U^*\jmath_P(G)U=0\oplus LGL^*,
\]
while
\[
 |\det L|^2=\det(L^*L)=\det(R_P^*R_P)
 =\det(C_{A_2}^{-1})=\frac13.
\]
Taking the determinant in this orthonormal frame gives
\[
 \det(H_0+\varepsilon\jmath_P(G))
 =3\varepsilon^2|\det L|^2\det G
 =\varepsilon^2\det G.
\]
\end{proof}

Put \(X_P=(R_P\;u)\).  The same block computation also records the exact
position of this module inside the full nine-real-dimensional carrier:
\[
 H_0+\varepsilon\jmath_P(G)
 =X_P\begin{pmatrix}\varepsilon G&0\\0&1\end{pmatrix}X_P^*.
\]
Thus, in the full \((G,\psi,\lambda)\) block coordinates, this section uses
the strict sublocus \((\psi,\lambda)=(0,1)\).  Its field models live on the
observed four-vector section; they do not yet define dynamics for the retained
spinor \(\psi\) and scalar \(\lambda\) directions.

For \(G,K\in\operatorname{Herm}_2(\mathbb C)\), define
\begin{equation}
 \eta_2(G,K)=\frac12\bigl(\det(G+K)-\det G-\det K\bigr).
 \label{hnsp-polarized-determinant}
\end{equation}
Transport this form to \(\mathcal N_P\) by
\begin{equation}
 \eta_P(\jmath_P(G),\jmath_P(K))=\eta_2(G,K).
 \label{hnsp-transported-metric}
\end{equation}

\begin{theorem}[The determinant Hessian is exactly Minkowski]
\label{hnsp-lorentz-theorem}
For tangent vectors \(X=\jmath_P(G)\) and \(Y=\jmath_P(K)\),
\begin{equation}
 \eta_P(X,Y)=\frac12D^2(\det)_{H_0}[X,Y].
 \label{hnsp-hessian-metric}
\end{equation}
Writing
\begin{equation}
 G(t,x,y,z)=
 \begin{pmatrix}t+y&x+iz\\x-iz&t-y\end{pmatrix},
 \label{hnsp-pauli-coordinates}
\end{equation}
Relative to the convention
\(
 \left(\begin{smallmatrix}
 t+x_3&x_1-i x_2\\ x_1+i x_2&t-x_3
 \end{smallmatrix}\right)
\)
used in the observation construction, the exact coordinate dictionary is
\((x_1,x_2,x_3)=(x,-z,y)\).  Thus the two displayed Pauli conventions differ
only by this stated spatial relabelling and sign, not by a change of Lorentz
form.
one has
\begin{equation}
 \eta_2(G(t,x,y,z),G(t',x',y',z'))
 =tt'-xx'-yy'-zz'.
 \label{hnsp-minkowski-form}
\end{equation}
Thus \((\mathcal N_P,\eta_P)\) has signature \((1,3)\).  Its distinguished
future cone is
\begin{equation}
 \mathcal C_P^+=\jmath_P\bigl(\operatorname{Herm}_2(\mathbb C)_{\geq0}\bigr)
 =\{\jmath_P(G(t,\mathbf x)):t\geq|\mathbf x|\}.
 \label{hnsp-future-cone}
\end{equation}
Its nonzero null boundary consists exactly of the images of rank-one positive
semidefinite matrices.
\end{theorem}

\begin{proof}
Proposition~\ref{hnsp-normal-identification}, applied to \(sG+rK\), gives
\[
 \det(H_0+s\jmath_P(G)+r\jmath_P(K))=\det(sG+rK).
\]
The coefficient of \(sr\) on the right is \(2\eta_2(G,K)\).  Taking one
derivative in each variable proves \eqref{hnsp-hessian-metric}.  Direct
expansion of \eqref{hnsp-pauli-coordinates} gives
\[
 \det G(t,x,y,z)=t^2-x^2-y^2-z^2,
\]
and polarization gives \eqref{hnsp-minkowski-form}.  A two-by-two Hermitian
matrix is positive semidefinite precisely when its trace is nonnegative and
its determinant is nonnegative, which here is
\(t\geq0\) and \(t^2\geq|\mathbf x|^2\).  On the nonzero boundary its
determinant vanishes and its rank is one.
\end{proof}

The corresponding tetrad inside \(\operatorname{Herm}_3(\mathbb C)\) is
completely explicit.  If
\(E_0=\jmath_P(G(1,0,0,0))\),
\(E_1=\jmath_P(G(0,1,0,0))\),
\(E_2=\jmath_P(G(0,0,1,0))\), and
\(E_3=\jmath_P(G(0,0,0,1))\), then
\begin{align}
 E_0&=\frac19\begin{pmatrix}5&-1&-4\\-1&2&-1\\-4&-1&5\end{pmatrix},&
 E_1&=\frac19\begin{pmatrix}4&1&-5\\1&-2&1\\-5&1&4\end{pmatrix},
 \label{hnsp-explicit-tetrad-real}\\
 E_2&=\frac13\begin{pmatrix}1&-1&0\\-1&0&1\\0&1&-1\end{pmatrix},&
 E_3&=\frac13\begin{pmatrix}0&i&-i\\-i&0&i\\i&-i&0\end{pmatrix}.
 \label{hnsp-explicit-tetrad-imaginary}
\end{align}
Every row and column of each \(E_\mu\) sums to zero, and
\begin{equation}
 \eta_P(E_\mu,E_\nu)=\operatorname{diag}(1,-1,-1,-1)_{\mu\nu}.
 \label{hnsp-tetrad-metric}
\end{equation}
These identities follow by multiplying the matrices in
\eqref{hnsp-moore-penrose} and then applying
\eqref{hnsp-minkowski-form}; no orthonormalization of the marked lattice is
being suppressed.

The metric is a genuine normal Hessian.  It is not a positive energy form:
it takes both signs.  Energy appears only after choosing a time function and
passing to a Hamiltonian formulation below.

Because
\(
 \det(H_0+\jmath_P(G))=\det G
\)
on the entire affine slice, not just to second order, its restriction there is
quadratic.  Hence the polarized Hessian \(\eta_P\) is constant at every point
of the affine slice; affine translations really are isometries of this metric.

\subsection{Exact Lorentz and affine symmetry}

For \(A\in SL_2(\mathbb C)\), define
\begin{equation}
 \widehat A_P=\frac13uu^*+R_P A P\in SL_3(\mathbb C).
 \label{hnsp-lifted-spin-action}
\end{equation}

\begin{theorem}[The spin action inside the normal slice]
\label{hnsp-spin-theorem}
The matrix \(\widehat A_P\) is invertible, with
\begin{equation}
 \widehat A_P^{-1}=\frac13uu^*+R_P A^{-1}P,
 \qquad \det\widehat A_P=\det A=1.
 \label{hnsp-spin-inverse}
\end{equation}
It fixes \(u\) and satisfies
\begin{equation}
 \widehat A_P H_0\widehat A_P^*=H_0,
 \qquad
 \widehat A_P\jmath_P(G)\widehat A_P^*
 =\jmath_P(AGA^*).
 \label{hnsp-spin-congruence}
\end{equation}
Consequently \(SL_2(\mathbb C)\) acts on \(\mathcal N_P\) by future-cone
preserving \(\eta_P\)-isometries.  The kernel is \(\{I_2,-I_2\}\), and the
induced map
\begin{equation}
 PSL_2(\mathbb C)\xrightarrow{\ \sim\ }SO^+(1,3)
 \label{hnsp-spin-cover}
\end{equation}
is an isomorphism of real Lie groups.
\end{theorem}

The kernel \(\{\pm I_2\}\) is a statement about the normal-vector action.
On the full \((G,\psi,\lambda)\) carrier the same embedded
\(SL_2(\mathbb C)\)-action is faithful, because \(-I_2\) fixes \(G\) and
\(\lambda\) but sends \(\psi\mapsto-\psi\).

\begin{proof}
Let \(Q_0=uu^*/3\) and \(Q_\perp=R_P P=I_3-Q_0\).  The relations
\(Q_0R_P=0\), \(PQ_0=0\), and \(PR_P=I_2\) show that the two matrices in
\eqref{hnsp-spin-inverse} multiply to \(Q_0+Q_\perp=I_3\).  In the ordered
decomposition \(\mathbb Cu\oplus R_P\mathbb C^2\), the map is
\(1\oplus A\), proving its determinant.  The same relations prove
\eqref{hnsp-spin-congruence}.  Since
\(\det(AGA^*)=|\det A|^2\det G=\det G\), the action preserves
\(\eta_P\); congruence preserves positive semidefiniteness, hence the chosen
future cone.

If \(AGA^*=G\) for every Hermitian \(G\), taking \(G=I_2\) makes \(A\)
unitary, after which \(A\) commutes with every Hermitian matrix.  Hence
\(A\) is scalar, and determinant one gives \(A=\pm I_2\).  The differential
of the induced map has zero kernel: if
\(XG+GX^*=0\) for every Hermitian \(G\), then \(G=I_2\) makes \(X\)
skew-Hermitian and the remaining equations make it scalar; trace zero then
gives \(X=0\).  Domain and target Lie algebras both have real dimension six,
so the image contains a neighborhood of the identity.  It is therefore an
open subgroup of the connected group \(SO^+(1,3)\), hence the whole group.
\end{proof}

Translations are also intrinsic to the affine slice.  If
\(B\in\operatorname{Herm}_2(\mathbb C)\), define
\begin{equation}
 (B,A)\cdot\bigl(H_0+\jmath_P(G)\bigr)
 =H_0+\jmath_P(B+AGA^*).
 \label{hnsp-poincare-action}
\end{equation}
With multiplication
\(
 (B,A)(C,D)=(B+ACA^*,AD),
\)
this is the spin double cover of the proper orthochronous Poincar\'e action.
Its kernel is \(\{(0,I_2),(0,-I_2)\}\).
Under the dimensional chart \eqref{hnsp-dimensional-chart} introduced below,
a physical translation by \(b\in\mathbb R^{1,3}\) is the element
\(B=G(b/\ell)\).

\begin{proposition}[Ambient \(SU(3)\) is not an internal color action]
\label{hnsp-su3-distinction}
The subgroup of \(SU(3)\) preserving \(H_0\) is
\begin{equation}
 \operatorname{Stab}_{SU(3)}(H_0)
 =\{g\in SU(3):gu=e^{i\theta}u\}
 \cong S(U(1)\times U(2)).
 \label{hnsp-su3-stabilizer}
\end{equation}
It preserves \(\mathcal N_P\).  A general element of \(SU(3)\) moves
\(H_0\) and hence moves the chosen affine normal slice to a different one.
Furthermore, \(\widehat A_P\) in \eqref{hnsp-lifted-spin-action} is unitary
exactly when
\begin{equation}
 A^*C_{A_2}^{-1}A=C_{A_2}^{-1};
 \label{hnsp-compact-spin-subgroup}
\end{equation}
this compact subgroup is conjugate to \(SU(2)\), whereas generic Lorentz
boosts are not unitary.  Thus the Lorentz spin action supplied by the normal
slice is not the ambient \(SU(3)\) orbit action and does not itself define an
internal \(SU(3)\) gauge symmetry.
\end{proposition}

\begin{proof}
The equality \(gH_0g^*=H_0\) is equivalent to equality of the rank-one
lines and norms, hence to \(gu=e^{i\theta}u\).  The orthogonal complement is
then invariant, giving the displayed block stabilizer and preservation of
Hermitian operators supported on \(u^\perp\).  On
\(R_P\mathbb C^2=u^\perp\), the standard norm pulls back to the form
\(R_P^*R_P=C_{A_2}^{-1}\).  Therefore \(\widehat A_P\) is unitary precisely
when \eqref{hnsp-compact-spin-subgroup} holds.  Positive-definite congruence
identifies this group with \(SU(2)\).
\end{proof}

The proposition concerns the congruence action of the ambient \(SU(3)\) on
\(\operatorname{Herm}_3(\mathbb C)\).  It is compatible with, but distinct
from, the left \(SU(3)\)-action on
\(\operatorname{Mat}_{3\times2}(\mathbb C)\) whose invariant Gram map is
\(Z\mapsto Z^*Z\): the source spaces and actions are different, and no
identification of them is being made here.

\subsection{The compact gauge--scalar carrier selected by the standard metrics}

The Lorentz congruence group is noncompact, so it cannot preserve a positive
field-space norm.  The standard metrics already used in
\eqref{hnsp-moore-penrose} do, however, select the positive matrix and compact
subgroup
\begin{equation}
 W_P:=R_P^*R_P=C_{A_2}^{-1}
 =\frac13\begin{pmatrix}2&1\\1&2\end{pmatrix},\qquad
 K_P:=\{a\in SL_2(\mathbb C):a^* W_P a=W_P\}.
 \label{hnsp-intrinsic-compact-group}
\end{equation}
Because \(W_P=C_{A_2}^{-1}\), inversion of the defining equality gives
\[
 a^*W_Pa=W_P
 \quad\Longleftrightarrow\quad
 aC_{A_2}a^*=C_{A_2}.
\]
Accordingly, \(K_P\) is exactly the Lorentz little group of the future
timelike vector \(C_{A_2}\), equivalently the observer stabilizer selected by
the standard metrics.  Its promotion to a local internal gauge symmetry below
is mathematically valid but is an additional physical interpretation.

\begin{proposition}[Exact compact representation on the normal slice]
\label{hnsp-intrinsic-compact-representation}
Conjugation \(a\mapsto W_P^{1/2}aW_P^{-1/2}\) is an isomorphism
\(K_P\cong SU(2)\).  The real bilinear form
\begin{equation}
 E_P(G,K):=\operatorname{tr}\bigl(\jmath_P(G)\jmath_P(K)\bigr)
 =\operatorname{tr}(G W_P K W_P)
 \label{hnsp-positive-normal-metric}
\end{equation}
is a positive-definite \(K_P\)-invariant inner product on
\(\operatorname{Herm}_2(\mathbb C)\).  For \(a\in K_P\), the lift
\begin{equation}
 \widehat a_P=\frac13uu^*+R_P a P
 \label{hnsp-compact-lift}
\end{equation}
lies in \(SU(3)\), fixes \(H_0\), and obeys
\begin{equation}
 \widehat a_P\jmath_P(G)\widehat a_P^*=\jmath_P(aGa^*).
 \label{hnsp-compact-equivariance}
\end{equation}
After whitening, \(\widetilde G=W_P^{1/2}GW_P^{1/2}\), the action is
ordinary unitary congruence.  Therefore, as a real \(K_P\)-module,
\begin{equation}
 \operatorname{Herm}_2(\mathbb C)
 =\mathbb R W_P^{-1}
  \oplus\{G\in\operatorname{Herm}_2(\mathbb C):
              \operatorname{tr}(W_PG)=0\}
 \cong\boldsymbol 1\oplus\boldsymbol 3.
 \label{hnsp-singlet-adjoint-decomposition}
\end{equation}
The centre \(\{\pm I_2\}\) acts trivially.  Thus the metric-selected normal
field is a singlet plus an adjoint triplet; it is neither a fundamental weak
doublet nor an \(SU(3)\) color triplet.
\end{proposition}

\begin{proof}
The defining identity in \eqref{hnsp-intrinsic-compact-group} is equivalent to
\((W_P^{1/2}aW_P^{-1/2})^*(W_P^{1/2}aW_P^{-1/2})=I_2\); determinants are one,
so this identifies \(K_P\) with \(SU(2)\).  Cyclicity of the trace and
\(R_P^*R_P=W_P\) give the second equality in
\eqref{hnsp-positive-normal-metric}.  Moreover
\(E_P(G,G)=\operatorname{tr}(\jmath_P(G)^2)\), which is positive for nonzero
Hermitian \(\jmath_P(G)\).  Equation
\eqref{hnsp-spin-congruence} and unitarity of \(\widehat a_P\) prove
invariance.  Proposition~\ref{hnsp-su3-distinction} gives
\(\widehat a_P\in U(3)\), and the block action on
\(u^\perp\oplus\mathbb Cu\) gives
\(\det\widehat a_P=\det a=1\).  The whitened formula is direct
multiplication.  Unitary congruence fixes the scalar line and acts through the
adjoint rotation representation on traceless Hermitian matrices; the central
elements act by \(G\mapsto(\pm I_2)G(\pm I_2)^*=G\).
\end{proof}

\begin{theorem}[A fully typed positive gauge--scalar realization]
\label{hnsp-intrinsic-gauge-scalar-theorem}
Let \((M,g)\) be an oriented Riemannian \(d\)-manifold and use the trivial
principal \(K_P\)-bundle.  Take
\begin{equation}
 \mathcal A\in\Omega^1(M;\mathfrak k_P),\qquad
 G\in C^\infty(M;\operatorname{Herm}_2(\mathbb C))
 \label{hnsp-intrinsic-fields}
\end{equation}
as the connection and normal scalar.  For
\(a\in C^\infty(M,K_P)\), define
\begin{align}
 G^a&=aGa^*,&
 \mathcal A^a&=a\mathcal A a^{-1}-(da)a^{-1},
 \label{hnsp-intrinsic-gauge-action}\\
 D_{\mathcal A}G&=dG+\mathcal A G+G\mathcal A^*,&
 F_{\mathcal A}&=d\mathcal A+\mathcal A\wedge\mathcal A.
 \label{hnsp-intrinsic-covariant-objects}
\end{align}
Choose on \(\mathfrak k_P\) the positive
\(\operatorname{Ad}(K_P)\)-invariant form transported from
\(-\operatorname{Re}\operatorname{tr}(XY)\) on \(\mathfrak{su}(2)\).  For
\(g_{\rm YM},Z_G>0\) and \(m_G,\lambda_4\geq0\), the Euclidean action
\begin{equation}
 \begin{split}
 S^{\rm cpt}_{g_{\rm YM},Z_G,m_G,\lambda_4}[\mathcal A,G]
 =\int_M\bigg(&\frac1{2g_{\rm YM}^2}|F_{\mathcal A}|^2
 +\frac {Z_G}2|D_{\mathcal A}G|_{E_P}^2
 +\frac{m_G^2}{2}|G|_{E_P}^2\\
 &+\frac{\lambda_4}4(\det G)^2\bigg)\,d\operatorname{vol}_g
 \end{split}
 \label{hnsp-intrinsic-positive-action}
\end{equation}
is nonnegative and gauge invariant.  The associated normal-matrix field
\begin{equation}
 H_\varepsilon=H_0+\varepsilon\jmath_P(G)
 \label{hnsp-intrinsic-normal-field}
\end{equation}
is equivariant under the lifted action
\(H_\varepsilon^a=\widehat a_P H_\varepsilon\widehat a_P^*\).  If
\(\widehat{\mathcal A}=R_P\mathcal A P\), then
\begin{equation}
 D_{\widehat{\mathcal A}}\jmath_P(G)
 =\jmath_P(D_{\mathcal A}G),\qquad
 D_{\widehat{\mathcal A}}H_0=0.
 \label{hnsp-intrinsic-connection-intertwining}
\end{equation}
On \(M=\mathbb R^d\), the scalar Hessian at
\((\mathcal A,G)=(0,0)\) is the nonnegative self-adjoint operator
\begin{equation}
 \begin{aligned}
  \mathcal K_{Z_G,m_G}
   &=(-Z_G\Delta+m_G^2)\otimes I_{\operatorname{Herm}_2},\\
  \operatorname{Dom}(\mathcal K_{Z_G,m_G})
   &=H^2(\mathbb R^d;\operatorname{Herm}_2(\mathbb C)),\\
  \operatorname{Spec}(\mathcal K_{Z_G,m_G})
   &=[m_G^2,\infty).
 \end{aligned}
 \label{hnsp-intrinsic-scalar-hessian}
\end{equation}
Thus the fixed normal data plus the standard metrics select the compact
representation \(K_P\curvearrowright\operatorname{Herm}_2(\mathbb C)\) and
its positive fibre metric.  After the additional choice to gauge that
pointwise action, the formulas above give an exact compact gauge--scalar
model, but they still do not select \(m_G,Z_G,g_{\rm YM},\lambda_4\), a quantum
state, or an interacting gap.
\end{theorem}

\begin{proof}
Direct substitution gives
\(D_{\mathcal A^a}G^a=a(D_{\mathcal A}G)a^*\) and
\(F_{\mathcal A^a}=aF_{\mathcal A}a^{-1}\).  Invariance of the two positive
inner products and
\(\det(aGa^*)=|\det a|^2\det G=\det G\) prove gauge invariance; every term in
\eqref{hnsp-intrinsic-positive-action} is nonnegative.  Equation
\eqref{hnsp-compact-equivariance} gives equivariance of
\eqref{hnsp-intrinsic-normal-field}.  Using \(PR_P=I_2\), \(Pu=0\), and
\(u^*R_P=0\), direct multiplication proves
\eqref{hnsp-intrinsic-connection-intertwining}.  Fourier transformation gives
the self-adjoint operator, its full domain, and its spectrum in
\eqref{hnsp-intrinsic-scalar-hessian}.  Finally, \((\det G)^2\) is homogeneous
of degree four, so it has zero Hessian at \(G=0\).
\end{proof}

\begin{corollary}[The selected normal field cannot screen a fundamental centre charge]
\label{hnsp-centre-screening-corollary}
Every tensor power and every polynomial composite of the
\(\boldsymbol 1\oplus\boldsymbol 3\) normal representation has trivial
action by \(-I_2\in K_P\).  The fundamental two-dimensional representation
has action \(-I_2\).  Therefore the field \(G\) in
Theorem~\ref{hnsp-intrinsic-gauge-scalar-theorem} cannot by itself screen a
fundamental centre charge; such a model requires an additional
centre-nontrivial matter field and its independently specified kinetic,
mass, and interaction terms.
This non-screening statement is confined to the \(G\)-only normal section.
The full \((G,\psi,\lambda)\) carrier already contains the centre-nontrivial
spinor \(\psi\); what remains missing is a specified kinetic and interaction
theory for that retained field.
\end{corollary}

\begin{proof}
The centre acts trivially on \(G\), hence trivially on every tensor power and
subquotient generated from it.  It acts as minus the identity in the
fundamental representation, so that representation cannot occur in this
tensor category.
\end{proof}

\subsection{The candidate cusp matrix is an exact spacelike vector}

The preceding construction applies to every Hermitian two-by-two matrix.  For
the particular inverse-dual matrix retained in the cusp calculation, let
\begin{equation}
 G_{\eta,T}=\begin{pmatrix}2L_T&2m_T\\2m_T&-q_T/3\end{pmatrix},
 \qquad D_T=L_T q_T+6m_T^2>0.
 \label{hnsp-candidate-matrix}
\end{equation}

\begin{proposition}[Exact candidate-specific Lorentz coordinates]
\label{hnsp-candidate-spacelike}
In the coordinates \eqref{hnsp-pauli-coordinates},
\begin{equation}
 (t_T,x_T,y_T,z_T)
 =\left(L_T-\frac{q_T}{6},\ 2m_T,\ L_T+\frac{q_T}{6},\ 0\right),
 \label{hnsp-candidate-coordinates}
\end{equation}
and
\begin{equation}
 \eta_2(G_{\eta,T},G_{\eta,T})
 =\det G_{\eta,T}=-\frac23D_T<0.
 \label{hnsp-candidate-spacelike-norm}
\end{equation}
If \(L_T/T\to1\), \(q_T/T\to1\), and \(m_T/T\to0\), then
\begin{equation}
 \frac1T(t_T,x_T,y_T,z_T)
 \longrightarrow\left(\frac56,0,\frac76,0\right).
 \label{hnsp-candidate-direction-limit}
\end{equation}
Thus this particular candidate vector is spacelike for every \(T\).  It is
not a positive-mass four-momentum.  This statement concerns the indefinite
inverse-dual matrix \(G_{\eta,T}\); it does not concern the different positive
Gram momentum \(p_\alpha(Z)=\alpha Z^*Z\) used in the retained-\(c\) worldline
construction.  The latter is positive semidefinite by definition, and its
image is all of \(\operatorname{Herm}_2(\mathbb C)_{\geq0}\): for every
\(G\succeq0\), the block matrix
\(
 Z=\alpha^{-1/2}
 \left(\begin{smallmatrix}G^{1/2}\\0_{1\times2}\end{smallmatrix}\right)
 \in\operatorname{Mat}_{3\times2}(\mathbb C)
\)
satisfies \(p_\alpha(Z)=G\).  Hence
\[
 G_{\eta,T}\notin\operatorname{im}p_\alpha
 \qquad(T>0,\ \alpha>0).
\]
Positive rescaling preserves the determinant sign and \(GL_2(\mathbb C)\)
congruence preserves inertia, so neither operation repairs this mismatch.
The Gram momentum lies on \(\det p_\alpha=m_0^2c^2\) only after that
mass-shell constraint and the parameters \(\alpha,c,m_0\) are supplied.
Interpreting \(G_{\eta,T}\) as a field value, a normal deformation, or a
spacelike wave-vector requires an additional map; the determinant calculation
alone does not choose among those interpretations.
\end{proposition}

\begin{proof}
Solving \(t+y=2L_T\), \(t-y=-q_T/3\), and
\(x+iz=2m_T\) gives \eqref{hnsp-candidate-coordinates}.  Substitution gives
\[
 t_T^2-x_T^2-y_T^2-z_T^2
 =(2L_T)(-q_T/3)-4m_T^2=-\frac23D_T.
\]
The limit is componentwise.
\end{proof}

\begin{theorem}[Exact determinant-level orbit geometry]
\label{hnsp-determinant-orbit-geometry}
Fix \(\varepsilon>0\) and write
\[
 \mathbb M_\varepsilon
 :=H_0+\varepsilon\jmath_P\bigl(\operatorname{Herm}_2(\mathbb C)\bigr).
\]
For every \(r>0\), the future timelike determinant shell
\begin{equation}
 \mathscr P^+_{r,\varepsilon}
 :=\{H_0+\varepsilon\jmath_P(G):G\succ0,\ \det G=r^2\}
 \cong SL_2(\mathbb C)/SU(2)
 \label{hnsp-future-orbit}
\end{equation}
is one orbit of the lifted congruence action.  It equals
\[
 \{H\in\mathbb M_\varepsilon:H\succ0,\
   \det H=\varepsilon^2r^2\}.
\]
The nonzero future-null locus
\begin{equation}
 \mathscr N^+_\varepsilon
 :=\{H_0+\varepsilon\jmath_P(G):
       G\succeq0,\ \operatorname{rank}G=1\}
 \label{hnsp-null-orbit}
\end{equation}
is also one \(SL_2(\mathbb C)\)-orbit; each of its matrices is positive
semidefinite of ambient rank two.  Finally, with
\(r_T=\sqrt{2D_T/3}\), the candidate-specific spacelike shell is
\begin{equation}
 \mathscr S_{T,\varepsilon}
 :=\{H_0+\varepsilon\jmath_P(G):\det G=-r_T^2\}
 \cong SL_2(\mathbb C)/SU(1,1),
 \label{hnsp-spacelike-orbit}
\end{equation}
and every matrix in it has inertia \((2,1)\).  Consequently the normal slice
contains every timelike mass shell and the future null cone after a unit
dictionary is supplied, but it selects neither \(r\) nor a physical mass.
\end{theorem}

\begin{proof}
If \(G\succ0\) and \(\det G=r^2\), then
\(A=(G/r)^{1/2}\) lies in \(SL_2(\mathbb C)\) and
\(A(rI_2)A^*=G\).  The stabilizer of \(rI_2\) is \(SU(2)\), proving
\eqref{hnsp-future-orbit}; positivity and the ambient determinant follow from
the block decomposition in the proof of
Proposition~\ref{hnsp-normal-identification}.  Every nonzero positive
rank-one \(G\) is \(vv^*\).  Choose \(w\in\mathbb C^2\) with
\(\det(v\ w)=1\); then \(A=(v\ w)\in SL_2(\mathbb C)\) and
\(A\operatorname{diag}(1,0)A^*=vv^*\).  This proves null transitivity, and the
orthogonal block decomposition gives ambient rank \(1+1=2\).

Let \(J=\operatorname{diag}(1,-1)\).  A Hermitian \(G\) with
\(\det G=-r_T^2\) has inertia \((1,1)\), so Sylvester congruence gives
\(G=B(r_T J)B^*\) for some \(B\in GL_2(\mathbb C)\).  Determinants give
\(|\det B|=1\).  If \(\det B=e^{i\theta}\), then
\(D=\operatorname{diag}(e^{-i\theta},1)\) stabilizes \(J\), and
\(A=BD\) has determinant one.  Thus the action is transitive, and the
stabilizer is
\[
 \{A\in SL_2(\mathbb C):AJA^*=J\}=SU(1,1).
\]
Equation~\eqref{hnsp-candidate-spacelike-norm} puts the retained candidate on
this shell.  The block decomposition gives one positive eigenvalue from
\(H_0\) and one of each sign from \(G\), proving the inertia claim.
\end{proof}

\subsection{Units and the minimal continuum field model}

All entries in \eqref{hnsp-fixed-data}--\eqref{hnsp-normal-space} are
dimensionless.  A physical spacetime requires a scale.  For \(\ell>0\), set
\begin{equation}
 \kappa_\ell(x^0,x^1,x^2,x^3)
 =H_0+\jmath_P\!\left(
 G\!\left(\frac{x^0}{\ell},\frac{x^1}{\ell},
           \frac{x^2}{\ell},\frac{x^3}{\ell}\right)\right).
 \label{hnsp-dimensional-chart}
\end{equation}
Equip the affine slice with the constant metric \(g_{P,\ell}=\ell^2\eta_P\).
Then
\begin{equation}
 \kappa_\ell^*g_{P,\ell}
 =(dx^0)^2-(dx^1)^2-(dx^2)^2-(dx^3)^2,
 \qquad
 \det\kappa_\ell(x)=\ell^{-2}\eta_{\mu\nu}x^\mu x^\nu.
 \label{hnsp-dimensional-isometry}
\end{equation}
For the scaled normal displacement
\(H_0+\varepsilon\jmath_P(G(x/\ell))\), the determinant is multiplied by
\(\varepsilon^2\).  For \(\varepsilon\ne0\), the corresponding isometric
metric rescaling is \((\ell/|\varepsilon|)^2\eta_P\).  Hence
\(\varepsilon\) fixes no absolute energy or length by itself.

The exact unit dictionary is:
\begin{center}
\begin{tabular}{@{}>{\raggedright\arraybackslash}p{0.18\linewidth}
                    >{\raggedright\arraybackslash}p{0.34\linewidth}
                    >{\raggedright\arraybackslash}p{0.38\linewidth}@{}}
\toprule
symbol & role & status from the determinant data\\
\midrule
\(P,u,R_P,\jmath_P\) & marked linear algebra & fixed and dimensionless\\
\(\varepsilon\) & normal-displacement amplitude & fixed only when supplied\\
\(\ell\) & conversion from matrix coordinates to length & independent choice\\
\(x^0=ct\) & time coordinate & requires the independent constant \(c\)\\
\(m_{\mathrm{KG}}\) & inverse Compton length in the scalar model & independent choice\\
\(m_{\rm phys}=\hbar m_{\mathrm{KG}}/c\) & physical scalar mass & requires \(\hbar,c\)\\
\(L=Na\) & infrared box size and ultraviolet mesh & regulator choices\\
\(K,g_{\rm YM}\) & compact gauge group and coupling & \(K\) arbitrary; \(K_P\cong SU(2)\) is the metric-selected little group and its gauging is optional\\
\bottomrule
\end{tabular}
\end{center}

The retained worldline mass \(m_0\), the Klein--Gordon inverse length
\(m_{\mathrm{KG}}\), and the dimensionless determinant-shell radius \(r\) are
independent inputs.  Only after declaring that the worldline and field theory
describe the same physical species may one impose the unit conversion
\[
 m_{\mathrm{KG}}=\frac{m_0c}{\hbar},
 \qquad
 \Delta E_{\mathrm{free}}=\hbar c\,m_{\mathrm{KG}}=m_0c^2.
\]
No determinant identity derives this identification.

Use natural units \(c=\hbar=1\) in the remaining formulas.  The minimal
field placed on the Lorentzian host is a real scalar
\(\phi\in C_c^\infty(\mathbb R^{1,3};\mathbb R)\) with action
\begin{equation}
 S_{m_{\mathrm{KG}}}[\phi]=\frac12\int_{\mathbb R^{1,3}}
 \left(\partial_\mu\phi\,\partial^\mu\phi-m_{\mathrm{KG}}^2\phi^2\right)d^4x,
 \qquad m_{\mathrm{KG}}\geq0.
 \label{hnsp-klein-gordon-action}
\end{equation}
Its Euler--Lagrange operator is \(\Box+m_{\mathrm{KG}}^2\).  The Poincar\'e action
\eqref{hnsp-poincare-action} acts on fields by pullback.  The action, the
stress-energy tensor
\begin{equation}
 T_{\mu\nu}=\partial_\mu\phi\partial_\nu\phi
 -\frac12\eta_{\mu\nu}
  (\partial_\alpha\phi\partial^\alpha\phi-m_{\mathrm{KG}}^2\phi^2),
 \label{hnsp-stress-energy}
\end{equation}
and the mass shell \(p_\mu p^\mu=m_{\mathrm{KG}}^2\) are invariant.  Smeared fields
\(\phi(f)=\int\phi f\,d^4x\) are the basic linear observables.

On \(x^0=0\), smooth compactly supported initial data
\((\phi_0,\pi_0)\) carry the symplectic form
\begin{equation}
 \sigma((\phi,\pi),(\psi,\rho))
 =\int_{\mathbb R^3}(\phi\rho-\pi\psi)\,d^3x
 \label{hnsp-scalar-symplectic-form}
\end{equation}
and energy
\begin{equation}
 E_{m_{\mathrm{KG}}}(\phi,\pi)=\frac12\int_{\mathbb R^3}
 (|\pi|^2+|\nabla\phi|^2+m_{\mathrm{KG}}^2|\phi|^2)\,d^3x.
 \label{hnsp-scalar-energy}
\end{equation}
For \(m_{\mathrm{KG}}>0\), the energy completion is
\(H^1(\mathbb R^3)\oplus L^2(\mathbb R^3)\); for \(m_{\mathrm{KG}}=0\), it is
\(\dot H^1(\mathbb R^3)\oplus L^2(\mathbb R^3)\).  The symplectic pairing
\eqref{hnsp-scalar-symplectic-form} is exact on the dense smooth,
compactly supported data.  In the massless homogeneous completion it need not
extend continuously to every pair without an additional dual-space choice.

\begin{theorem}[Exact free-field spectral alternative]
\label{hnsp-free-field-gap}
Let
\begin{equation}
 K_{m_{\mathrm{KG}}}=-\Delta+m_{\mathrm{KG}}^2,
 \quad \operatorname{Dom}(K_{m_{\mathrm{KG}}})=H^2(\mathbb R^3),
 \quad \omega_{m_{\mathrm{KG}}}=K_{m_{\mathrm{KG}}}^{1/2}
 \label{hnsp-one-particle-operator}
\end{equation}
on \(L^2(\mathbb R^3)\).  On the bosonic Fock space
\(\mathscr F_s(L^2(\mathbb R^3))\), define the vacuum-subtracted Hamiltonian
\begin{equation}
 \mathbb H_{m_{\mathrm{KG}}}=d\Gamma(\omega_{m_{\mathrm{KG}}}).
 \label{hnsp-free-fock-hamiltonian}
\end{equation}
Explicitly, on the orthogonal decomposition
\(\mathscr F_s=\bigoplus_{n=0}^\infty L^2(\mathbb R^3)^{\otimes_s n}\), its
domain consists of the vectors \(\Psi=(\Psi_n)\) for which each \(\Psi_n\)
belongs to the domain of
\(
 \sum_{j=1}^n I\otimes\cdots\otimes\omega_{m_{\mathrm{KG}}}\otimes\cdots\otimes I
\)
and the squared norms of those images are summable.
For every \(m_{\mathrm{KG}}\geq0\), the Fock vacuum is the unique zero-energy vector.  Its
excitation gap is exactly
\begin{equation}
 \inf\operatorname{Spec}(\mathbb H_{m_{\mathrm{KG}}}|_{\Omega^\perp})=m_{\mathrm{KG}}.
 \label{hnsp-free-gap-value}
\end{equation}
Thus the same determinant normal geometry admits a positive gap of any chosen
value \(m_{\mathrm{KG}}>0\) and a gapless model at \(m_{\mathrm{KG}}=0\).
\end{theorem}

\begin{proof}
Fourier transform identifies \(K_{m_{\mathrm{KG}}}\) with multiplication by
\(|p|^2+m_{\mathrm{KG}}^2\) and \(\omega_{m_{\mathrm{KG}}}\) with multiplication by
\(\sqrt{|p|^2+m_{\mathrm{KG}}^2}\).  The latter has zero kernel even when
\(m_{\mathrm{KG}}=0\), since
the singleton \(p=0\) has Lebesgue measure zero.  Hence the only Fock vector
annihilated by \(d\Gamma(\omega_{m_{\mathrm{KG}}})\) is the vacuum.  On the
\(n\)-particle sector, the Hamiltonian is the sum of \(n\) copies of
\(\omega_{m_{\mathrm{KG}}}\), so the nonvacuum spectral bottom is at least
\(m_{\mathrm{KG}}\).  The one-particle sector is reducing and has spectral
bottom \(m_{\mathrm{KG}}\), proving equality.  For \(m_{\mathrm{KG}}=0\),
normalized Fourier-space functions supported in balls \(|p|<1/n\) have
energy expectation at most \(1/n\), which makes the absence of a positive gap
explicit.
\end{proof}

There is no contradiction between the absence of a positive kinematic
threshold on the unrestricted Gram carrier and the free Fock gap
\(m_{\mathrm{KG}}\): the former allows the rescaling \(Z\mapsto rZ\), whereas
the latter first fixes an external mass coefficient and quantizes the chosen
free Hamiltonian.

\subsection{A complete finite regulator and its two limits}

Let \(L>0\), \(N\geq2\), \(a=L/N\), and
\(V_N=(\mathbb Z/N\mathbb Z)^3\).  Give
\(Q_{L,N}=\mathbb R^{V_N}\) the inner product
\begin{equation}
 \langle q,r\rangle_a=a^3\sum_{v\in V_N}q(v)r(v).
 \label{hnsp-lattice-inner-product}
\end{equation}
The positive periodic lattice Laplacian is
\begin{equation}
 (-\Delta_a q)(v)=a^{-2}\sum_{j=1}^3
 \bigl(2q(v)-q(v+e_j)-q(v-e_j)\bigr).
 \label{hnsp-lattice-laplacian}
\end{equation}
Let \(d\lambda_a\) be Lebesgue measure associated with the Euclidean space
\((Q_{L,N},\langle\cdot,\cdot\rangle_a)\), and let
\(\Delta_{Q_{L,N}}\) denote its coordinate Laplacian, so that
\(-\Delta_{Q_{L,N}}\) is nonnegative.  For
\(m_{\mathrm{KG}}>0\), define on \(L^2(Q_{L,N},d\lambda_a)\)
\begin{equation}
 \mathbb H_{m_{\mathrm{KG}},L,N}
 =\frac12\left(-\Delta_{Q_{L,N}}
 +\langle q,(-\Delta_a+m_{\mathrm{KG}}^2)q\rangle_a\right)
 -E_{0,m_{\mathrm{KG}},L,N},
 \label{hnsp-lattice-hamiltonian}
\end{equation}
where \(E_{0,m_{\mathrm{KG}},L,N}\) is its finite oscillator ground energy.
More precisely, before subtracting \(E_{0,m_{\mathrm{KG}},L,N}\), this is the self-adjoint
operator associated with the closed form
\begin{equation}
 \mathfrak h_{m_{\mathrm{KG}},L,N}[\Psi]
 =\frac12\int_{Q_{L,N}}
 \left(|\nabla\Psi(q)|_a^2
 +\langle q,(-\Delta_a+m_{\mathrm{KG}}^2)q\rangle_a
             |\Psi(q)|^2\right)d\lambda_a(q),
 \label{hnsp-lattice-closed-form}
\end{equation}
whose domain is
\begin{equation}
 H^1(Q_{L,N})\cap
 L^2\!\left(Q_{L,N},
  \langle q,(-\Delta_a+m_{\mathrm{KG}}^2)q\rangle_a\,
  d\lambda_a(q)\right).
 \label{hnsp-lattice-form-domain}
\end{equation}

\begin{theorem}[Exact regulated spectrum and continuum limits]
\label{hnsp-regulator-theorem}
For \(n=(n_1,n_2,n_3)\in V_N\), the squared normal-mode frequencies are
\begin{equation}
 \omega_{m_{\mathrm{KG}},L,N}(n)^2
 =m_{\mathrm{KG}}^2+\frac4{a^2}\sum_{j=1}^3
 \sin^2\!\left(\frac{\pi n_j}{N}\right).
 \label{hnsp-lattice-dispersion}
\end{equation}
The operator \(\mathbb H_{m_{\mathrm{KG}},L,N}\) is self-adjoint on the tensor-product
oscillator domain, has a unique vacuum, and has exact gap \(m_{\mathrm{KG}}\).  The
subtracted constant is
\begin{equation}
 E_{0,m_{\mathrm{KG}},L,N}
 =\frac12\sum_{n\in V_N}\omega_{m_{\mathrm{KG}},L,N}(n).
 \label{hnsp-regulated-vacuum-energy}
\end{equation}

For \(m_{\mathrm{KG}}=0\), the constant mode is a free particle and the unreduced operator
has no normalized ground state.  On the mean-zero configuration space
\begin{equation}
 Q^0_{L,N}=\{q\in Q_{L,N}:\textstyle\sum_v q(v)=0\},
 \label{hnsp-mean-zero-space}
\end{equation}
the vacuum-subtracted oscillator has a unique ground state and gap
\begin{equation}
 \Delta^0_{L,N}=\frac2a\sin\!\left(\frac\pi N\right).
 \label{hnsp-massless-regulated-gap}
\end{equation}
At fixed \(L\), for every fixed \(r\in\mathbb Z^3\),
\begin{equation}
 \omega_{m_{\mathrm{KG}},L,N}(r)^2\longrightarrow
 m_{\mathrm{KG}}^2+\frac{4\pi^2|r|^2}{L^2}\qquad(N\to\infty),
 \label{hnsp-fixed-box-limit}
\end{equation}
and
\begin{equation}
 \Delta^0_{L,N}\longrightarrow\frac{2\pi}{L}.
 \label{hnsp-fixed-box-gap-limit}
\end{equation}
After the thermodynamic limit \(L\to\infty\), the massive gap remains
\(m_{\mathrm{KG}}\), whereas the massless positive-mode gap tends to zero.
\end{theorem}

\begin{proof}
The characters
\(
 v\mapsto L^{-3/2}e^{2\pi i n\cdot v/N}
\)
diagonalize \(-\Delta_a\).  Substitution gives
\eqref{hnsp-lattice-dispersion}.  Orthogonal diagonalization of the real
configuration space turns \eqref{hnsp-lattice-hamiltonian} into a tensor sum
of one-dimensional harmonic oscillators with frequencies
\(\omega_{m_{\mathrm{KG}},L,N}(n)\).  For \(m_{\mathrm{KG}}>0\), their
minimum is attained at \(n=0\) and equals \(m_{\mathrm{KG}}\), proving the
ground-state and gap claims.

When \(m_{\mathrm{KG}}=0\), the constant coordinate has zero restoring force,
hence gives
the free-particle operator and no \(L^2\) ground state.  Removing that
coordinate leaves the smallest character with one component congruent to
\(\pm1\) and the others zero, giving
\eqref{hnsp-massless-regulated-gap}.  With \(a=L/N\), the elementary limit
\(N\sin(\pi r_j/N)\to\pi r_j\) proves
\eqref{hnsp-fixed-box-limit} and \eqref{hnsp-fixed-box-gap-limit}.  The last
claims follow by sending \(L\) to infinity.
\end{proof}

This regulator supplies a literal configuration space, state Hilbert space,
quadratic Hamiltonian, observables generated by the coordinates and momenta,
and exact fixed-mode ultraviolet and lowest-frequency infrared limits.  The
displayed limits do not assert strong- or norm-resolvent convergence,
convergence of Gaussian measures, or convergence of Fock representations.  The mass
parameter is not hidden in \(P\), \(R_P\), \(\varepsilon\), or the
determinant: it is the independent coefficient in
\eqref{hnsp-klein-gordon-action}.

\subsection{An exact classical gauge-theory completion}

Let \(K\) be any compact matrix Lie group with Lie algebra \(\mathfrak k\)
and an \(\operatorname{Ad}K\)-invariant positive inner product.  On the
trivial bundle over the Lorentzian host, take a smooth compactly supported
gauge potential \(A\in\Omega_c^1(\mathbb R^{1,3};\mathfrak k)\), with
curvature
\begin{equation}
 F_A=dA+A\wedge A.
 \label{hnsp-gauge-curvature}
\end{equation}
The gauge group of smooth maps \(\mathbb R^{1,3}\to K\) that equal the
identity off a compact set acts by
\begin{equation}
 A^g=gAg^{-1}-(dg)g^{-1},\qquad F_{A^g}=gF_A g^{-1}.
 \label{hnsp-gauge-action}
\end{equation}
For an independently chosen coupling \(g_{\mathrm{YM}}>0\), the classical
action is
\begin{equation}
 S_{K,g_{\mathrm{YM}}}[A]
 =-\frac1{4g_{\mathrm{YM}}^2}
 \int\langle F_{\mu\nu},F^{\mu\nu}\rangle\,d^4x.
 \label{hnsp-yang-mills-action}
\end{equation}
It is invariant under \eqref{hnsp-gauge-action} and the Poincar\'e action.
For a closed curve \(\gamma\) and a finite-dimensional representation
\(\rho\) of \(K\),
\begin{equation}
 W_{\gamma,\rho}(A)
 =\operatorname{tr}\rho(\operatorname{Hol}_A(\gamma))
 \label{hnsp-wilson-observable}
\end{equation}
is gauge invariant.

On the spatial regulator \(\Sigma_L=(\mathbb R/L\mathbb Z)^3\), the smooth
Hamiltonian phase space before reduction consists of
\((A,E)\in\Omega^1(\Sigma_L;\mathfrak k)^2\), with gauge action
\begin{equation}
 (A,E)^g=\bigl(gAg^{-1}-(dg)g^{-1},\,gEg^{-1}\bigr).
 \label{hnsp-hamiltonian-gauge-action}
\end{equation}
Its Gauss constraint and
energy are
\begin{equation}
 d_A^*E=0,\qquad
 \mathcal E_{K,g}(A,E)=\frac1{2g_{\mathrm{YM}}^2}
 \int_{\Sigma_L}(|E|^2+|F_A|^2)\,d^3x,
 \label{hnsp-gauge-phase-space}
\end{equation}
and the reduced classical state space is the constraint locus modulo
\(C^\infty(\Sigma_L,K)\).  At this smooth level it is a set-theoretic orbit
space, generally stratified by stabilizer type; making it a smooth symplectic
space requires Sobolev completions and hypotheses on the action.  Equations
\eqref{hnsp-gauge-action} and \eqref{hnsp-gauge-phase-space} show directly
that the constraint locus, energy, and Wilson observables descend to this
quotient: in particular,
\(d_{A^g}^*E^g=g(d_A^*E)g^{-1}\), because the Hodge star acts only on the
base and the adjoint action is orthogonal.  In temporal gauge the canonical momentum is
\(\Pi_E=g_{\mathrm{YM}}^{-2}E\).  Equivalently, the symplectic form on
tangent vectors is
\begin{equation}
 \Omega_{(A,E)}\bigl((a,e),(a',e')\bigr)
 =\frac1{g_{\mathrm{YM}}^2}\int_{\Sigma_L}
 \bigl(\langle e,a'\rangle-\langle e',a\rangle\bigr)\,d\operatorname{vol}.
 \label{hnsp-gauge-symplectic-form}
\end{equation}

\begin{proposition}[The exact nonzero-mode Gaussian gauge gap]
\label{hnsp-gauge-gaussian-gap}
At \((A,E)=(0,0)\), impose the Coulomb conditions
\(d^*A=d^*E=0\) and remove harmonic one-forms.  Relative to the unweighted
\(L^2\) configuration metric, the energy Hessian is
\(g_{\mathrm{YM}}^{-2}\Delta_1\), where the generalized squared-frequency
operator is
\begin{equation}
 \Delta_1=d^*d,
 \quad
 \operatorname{Dom}(\Delta_1)
 =H^2\Omega^1(\Sigma_L;\mathfrak k)
  \cap\ker d^*\cap(\ker\Delta_1)^\perp.
 \label{hnsp-gauge-hessian-domain}
\end{equation}
on the closed Hilbert space
\begin{equation}
 \mathscr K_{L}^{\mathrm{tr}}
 =\bigl(L^2\Omega^1(\Sigma_L;\mathfrak k)\cap\ker d^*\bigr)
  \cap(\ker\Delta_1)^\perp.
 \label{hnsp-transverse-gauge-hilbert-space}
\end{equation}
Its spectrum is
\begin{equation}
 \left\{\frac{4\pi^2|n|^2}{L^2}:n\in\mathbb Z^3\setminus\{0\}\right\},
 \label{hnsp-gauge-hessian-spectrum}
\end{equation}
with transverse-polarization and \(\dim\mathfrak k\) multiplicities.  On the
unreduced transverse Gaussian Fock space after this explicit harmonic-mode
truncation, the gap is \(2\pi/L\), which tends to zero as \(L\to\infty\).
Constant gauge transformations remain after Coulomb gauge, so this statement
does not assert that its one-particle level survives restriction to a
gauge-invariant Fock subspace.
\end{proposition}

\begin{proof}
The curvature linearization at zero is \(A\mapsto dA\).  On the Coulomb
slice, the configuration part of the energy has quadratic form
\(g_{\mathrm{YM}}^{-2}\|dA\|^2/2\), whose Hessian relative to unweighted
\(L^2\) is \(g_{\mathrm{YM}}^{-2}d^*d\) with the displayed elliptic domain.
Dividing by the kinetic metric \(g_{\mathrm{YM}}^{-2}I\) gives the
squared-frequency operator \(\Delta_1=d^*d\).  Fourier characters on the
three-torus give \eqref{hnsp-gauge-hessian-spectrum}; for each nonzero
wave-vector, the condition \(d^*A=0\) leaves its two orthogonal polarizations.
Gaussian quantization replaces each Hessian eigenvalue \(\lambda\) by the
oscillator frequency \(\sqrt\lambda\): in canonical variables a single mode
has Hamiltonian
\(
 \tfrac12(g_{\mathrm{YM}}^2P_Q^2+
 \lambda g_{\mathrm{YM}}^{-2}Q^2)
\), whose frequency is \(\sqrt\lambda\).  The smallest nonzero wave-vector
has length \(2\pi/L\), proving the result.
\end{proof}

Thus the determinant slice is an exact four-dimensional Lorentzian host for
classical gauge theory, not merely an analogy.  It does not select \(K\): the
same construction works for \(U(1)\), \(SU(2)\), \(SU(3)\), or any other
compact matrix group.  In particular, Proposition~\ref{hnsp-su3-distinction}
prevents the ambient \(SU(3)\) orbit label from being silently reused as a
local color gauge group.  Nor does the free Hessian establish the nonlinear
four-dimensional Yang--Mills mass gap; that problem concerns the physical
gauge-invariant quantum Hamiltonian after a continuum construction
\cite{JaffeWitten2000}.

A cubical mesh gives the exact holonomy regulator of Wilson
\cite{Wilson_1974}.  To type the state dependence, let \(V\) be the vertices,
choose one orientation \(E^+\) of the edges, and put
\begin{equation}
 \mathcal C_E=K^{E^+},\qquad \mathcal G_V=K^V,
 \qquad (h\cdot U)_e=h_{s(e)}U_e h_{t(e)}^{-1}.
 \label{hnsp-link-configuration}
\end{equation}
Reverse orientation means \(U_{\bar e}=U_e^{-1}\).  For a closed edge path
\(\gamma=e_1\cdots e_r\) and a nontrivial irreducible representation \(\rho\),
use the normalized Wilson observable
\begin{equation}
 w_{\gamma,\rho}(U)
 =\frac1{\dim\rho}\operatorname{tr}
   \rho(U_{e_1}\cdots U_{e_r}).
 \label{hnsp-normalized-link-wilson}
\end{equation}

\begin{proposition}[Link kinematics selects no confinement state]
\label{hnsp-link-state-countermodels}
There are two \(\mathcal G_V\)-invariant probability measures on the same
configuration space \(\mathcal C_E\) such that
\begin{equation}
 \mathbb E_{\mathrm{pg}}[w_{\gamma,\rho}]=1
 \quad\text{for every closed }\gamma,
 \qquad
 \mathbb E_{\mathrm{Haar}}[w_{\gamma,\rho}]=0
 \label{hnsp-link-state-values}
\end{equation}
whenever \(\gamma\) contains an unoriented edge exactly once in the second
formula.  Their convex mixtures give every expectation between zero and one.
Thus the link configuration and gauge action do not determine an area law or
string tension.
\end{proposition}

\begin{proof}
For the first measure, give \(K^V\) product Haar measure and push it forward
under \(h\mapsto U(h)\), where
\(U(h)_e=h_{s(e)}h_{t(e)}^{-1}\).  Left invariance of vertex Haar measure
makes the push-forward gauge invariant, and every closed holonomy telescopes
to the identity, proving the first value.  For the second, give the independent
variables \(U_e\), \(e\in E^+\), product Haar measure.  Bi-invariance makes
this measure gauge invariant.  Integrating first over an edge that occurs
exactly once inserts
\(\int_K\rho(k)\,dk=0\), because a nontrivial irreducible representation has
no invariant vector.  This proves the second value.  Convexity proves the
remaining assertion.
\end{proof}

The link system is a valid ultraviolet regulator only after a time direction,
Euclidean rotation, coupling normalization, action, and probability measure
have been chosen.  Proposition~\ref{hnsp-link-state-countermodels} shows why
the normal-slice geometry and link space alone do not supply confinement.

\subsection{Exact obstructions and the sufficient spectral bridge}

\begin{proposition}[No Lorentz-invariant positive quadratic energy]
\label{hnsp-no-positive-lorentz-energy}
There is no positive-definite quadratic form on \(\mathcal N_P\) invariant
under the full \(PSL_2(\mathbb C)\) action of
Theorem~\ref{hnsp-spin-theorem}.  Consequently the determinant Hessian cannot
itself be a nonnegative Hamiltonian.
\end{proposition}

\begin{proof}
Suppose a positive-definite norm were invariant.  In the coordinates
\eqref{hnsp-pauli-coordinates}, Lorentz boosts send
\((1,0,0,0)\) to \((\cosh s,\sinh s,0,0)\).  Invariance would keep the
positive-definite norm of this orbit constant.  Every positive-definite norm
on a finite-dimensional space is equivalent to the Euclidean norm, whereas
the Euclidean norms of these vectors diverge with \(|s|\).  This is a
contradiction.
\end{proof}

The next two models make the role of entanglement exact.  They use bipartite
tensor-product Hilbert kinematics but have opposite gap behavior.

\begin{proposition}[Entanglement determines neither positivity nor magnitude of a gap]
\label{hnsp-entanglement-countermodels}
Let
\(\Omega=(e_0\otimes e_0+e_1\otimes e_1)/\sqrt2\) in
\(\mathbb C^2\otimes\mathbb C^2\).  Then
\begin{equation}
 H_{\mathrm{gap}}=I-|\Omega\rangle\langle\Omega|
 \label{hnsp-entangled-gapped}
\end{equation}
is bounded, nonnegative, and self-adjoint, has the entangled vector
\(\Omega\) as its unique vacuum, and has gap one.

On \(\ell^2(\mathbb N_0)\otimes\ell^2(\mathbb N_0)\), put
\begin{equation}
 \Omega_0=\frac{e_0\otimes e_0+e_1\otimes e_1}{\sqrt2},\qquad
 \psi_n=\frac{e_{2n}\otimes e_{2n}+e_{2n+1}\otimes e_{2n+1}}{\sqrt2}
 \quad(n\geq1).
 \label{hnsp-entangled-basis}
\end{equation}
Define a bounded self-adjoint operator by
\begin{equation}
 H_{\mathrm{ent}}\Omega_0=0,\qquad
 H_{\mathrm{ent}}\psi_n=n^{-1}\psi_n,
 \qquad H_{\mathrm{ent}}=I\ \text{on }
 \overline{\operatorname{span}}\{\Omega_0,\psi_n:n\geq1\}^\perp.
 \label{hnsp-entangled-gapless}
\end{equation}
Put
\[
 \mathcal K_{\mathrm{ent}}
 =\overline{\operatorname{span}}\{\psi_n:n\geq1\}.
\]
Every nonzero vector in \(\mathcal K_{\mathrm{ent}}\) is entangled,
\(\mathcal K_{\mathrm{ent}}\) reduces \(H_{\mathrm{ent}}\), and
\[
 \inf\operatorname{Spec}
  (H_{\mathrm{ent}}|_{\mathcal K_{\mathrm{ent}}})=0.
\]
Thus even a closed, completely entangled reducing sector need not have a
positive energy floor.
\end{proposition}

\begin{proof}
Equation \eqref{hnsp-entangled-gapped} is an orthogonal projection with
kernel \(\mathbb C\Omega\).  The vectors in
\eqref{hnsp-entangled-basis} are mutually orthonormal and each has Schmidt
rank two.  Equation \eqref{hnsp-entangled-gapless} is a bounded real diagonal
operator in the indicated orthogonal decomposition.  Its kernel is
\(\mathbb C\Omega_0\), while \(1/n\to0\).  Under the standard
coefficient-operator identification of a bipartite tensor, a nonzero vector
\(\sum_{n\geq1}c_n\psi_n\) has, for every \(n\) with \(c_n\ne0\), two equal
nonzero diagonal entries on the coordinates \(2n,2n+1\).  Its Schmidt rank is
therefore at least two.  The displayed diagonal action also proves reduction
of \(\mathcal K_{\mathrm{ent}}\) and spectral bottom zero.
\end{proof}

Finally, let \(A_{\mathrm{phys}}\) be a lower-semibounded self-adjoint
physical Hamiltonian on \(\mathcal H_{\mathrm{phys}}\), put
\(E_{\mathrm{vac}}=\inf\operatorname{Spec}(A_{\mathrm{phys}})\), and assume
that \(E_{\mathrm{vac}}\) is an eigenvalue.  Let
\(P_0=\boldsymbol 1_{\{E_{\mathrm{vac}}\}}(A_{\mathrm{phys}})\), let
\(\mathfrak a_{\mathrm{phys}}\) be its closed quadratic form, and put
\(\mathcal H_{\mathrm{exc}}=(I-P_0)\mathcal H_{\mathrm{phys}}\).  This is a
reducing subspace.  Let \(\mathcal H_{\mathrm{aux}}\) be another Hilbert
space and let \(B\geq0\) be self-adjoint on it, with closed form
\(\mathfrak b\) satisfying
\begin{equation}
 \mathfrak b[v]\geq\gamma\|v\|_{\mathrm{aux}}^2
 \quad(v\in\operatorname{Dom}\mathfrak b)
 \label{hnsp-auxiliary-gap}
\end{equation}
for some \(\gamma>0\).

\begin{theorem}[Full-domain coercive transfer]
\label{hnsp-coercive-transfer}
 Suppose \(\mathcal H_{\mathrm{exc}}\ne\{0\}\) and a linear map
\begin{equation}
 R:\operatorname{Dom}(\mathfrak a_{\mathrm{phys}})
   \cap\mathcal H_{\mathrm{exc}}
 \longrightarrow\operatorname{Dom}(\mathfrak b)
 \label{hnsp-comparison-map}
\end{equation}
is defined on the entire physical excitation form domain and, for fixed
\(c_R,C_R>0\), satisfies
\begin{equation}
 \|Ru\|_{\mathrm{aux}}\geq c_R\|u\|_{\mathrm{phys}},
 \qquad
 \mathfrak b[Ru]\leq
 C_R\bigl(\mathfrak a_{\mathrm{phys}}[u]
             -E_{\mathrm{vac}}\|u\|^2\bigr).
 \label{hnsp-coercive-comparison}
\end{equation}
Then
\begin{equation}
 \inf\operatorname{Spec}
 \bigl(A_{\mathrm{phys}}|_{\mathcal H_{\mathrm{exc}}}\bigr)-E_{\mathrm{vac}}
 \geq\frac{\gamma c_R^2}{C_R}.
 \label{hnsp-transferred-gap}
\end{equation}
An intertwiner from only one gapped subspace cannot replace
\eqref{hnsp-comparison-map}.
\end{theorem}

\begin{proof}
For every vector in the domain of \eqref{hnsp-comparison-map},
\[
 C_R(\mathfrak a_{\mathrm{phys}}[u]-E_{\mathrm{vac}}\|u\|^2)
 \geq\mathfrak b[Ru]\geq\gamma\|Ru\|_{\mathrm{aux}}^2
 \geq\gamma c_R^2\|u\|_{\mathrm{phys}}^2.
\]
The variational characterization of the spectral bottom gives
\eqref{hnsp-transferred-gap}.  For the final assertion, take \(B=I\) on
\(\mathbb C\) and
\[
 A=0\oplus I\oplus M_x
 \quad\text{on}\quad
 \mathbb C\Omega\oplus\mathbb C\oplus L^2((0,1),dx),
\]
where \((M_x f)(x)=x f(x)\).  The isometry \(Jz=0\oplus z\oplus0\) obeys
\(AJ=JB\), yet \(A\) has unique vacuum \(\Omega\) and excitation spectral
bottom zero because of the multiplication summand.
\end{proof}

\subsection{Exact conclusion}

The determinant identity therefore supplies all of the following, exactly:
a canonical dimensionless affine Minkowski space, its Lorentz metric and
future cone, the spin and Poincar\'e actions, a natural time orientation, and
a dimensionless normal coordinate system.  After the independent choice of
one length/time dictionary this carrier becomes a spacetime model supporting a
standard continuum Klein--Gordon model, an exact finite periodic regulator,
and any chosen classical compact-group gauge theory.  The standard metrics
also select the compact \(K_P\cong SU(2)\) singlet--adjoint representation;
Theorem~\ref{hnsp-intrinsic-gauge-scalar-theorem} gives one exact gauging of
that representation.  Those are substantive physical constructions with
typed configuration spaces, states, actions, observables, spectra, and
limits.

The determinant normal-slice data alone do not select a preferred inertial
observer, an internal gauge interpretation, a coupling, a mass, a probability
measure, a quantum reconstruction, or a nonlinear spectral gap.  A separate
two-column Gram construction supplies a left \(SU(3)\)-action and the
positive momentum \(p_\alpha(Z)=\alpha Z^*Z\), but it retains the independent
parameters \(\alpha,c,m_0\) and does not dynamically select \(m_0>0\).
These limitations are exact rather than semantic: the inverse-dual candidate
cusp vector \eqref{hnsp-candidate-coordinates} is spacelike; no
positive quadratic energy is Lorentz invariant; the same host has the gapped
and gapless scalar realizations of Theorem~\ref{hnsp-free-field-gap}; the
massless gauge Gaussian gap closes as \(L\to\infty\); and entangled states
occur in both gapped and gapless Hamiltonians.  A claim about a physical mass
gap must therefore be proved from the physical spectrum or from a global
comparison controlling its entire excitation domain;
Theorem~\ref{hnsp-coercive-transfer} is one sufficient such criterion.

% Workbench module: exact S^8 constraints for the higher-rung search.
% Date: 2026-09-03.  No standalone preamble is included.
% All labels in this module use the prefix hrs8:.

\section{Exact constraints on an \texorpdfstring{$S^8$}{S8} higher-rung analogue}
\label{hrs8:sec:constraints}

The phrase ``an $S^8$ analogue'' can denote three mathematically different
proposals, which must not be conflated.  A literal smooth complex-fourfold
analogue is impossible because a homotopy $8$-sphere is not almost complex.
A merely real map to $S^2$ with generic fibre $T^6$ is not excluded if it has
genuine singular fibres, but its local fillings must satisfy the exact Euler,
fundamental-group, Wang, and Mayer--Vietoris tests below.  Finally, explicit
real stratified quotient models with boundary or total space $S^8$ do exist;
they do not thereby acquire a complex-fourfold structure.

There is also a different, well-typed alternative: retain the four-real-
dimensional torus fibres of the reconstructed $S^6$ candidate and place
$SU(3)$ bundle data on its regular locus.  Subsection
\ref{hrs8:subsec:su3-regular-locus} computes the exact extension obstruction.
No gauge-theoretic or physical conclusion is asserted without a specified
bundle, connection, trace convention, boundary condition, and analytic
operator statement.

\subsection{The literal complex-fourfold proposal}
\label{hrs8:subsec:no-almost-complex}

\subsubsection{Characteristic-class obstruction}
\label{hrs8:subsubsec:characteristic-class}

Let $M$ be a smooth manifold homotopy equivalent to $S^8$, and suppose that
$J\colon TM\to TM$ satisfies $J^2=-1$.  Then
\[
 E=(TM,J)
\]
is a complex vector bundle of rank four whose underlying oriented real bundle
is $TM$.  Since
\[
 H^2(M;\mathbb Z)=H^4(M;\mathbb Z)=H^6(M;\mathbb Z)=0,
\]
one has
\[
 c_1(E)=c_2(E)=c_3(E)=0.
\]
Orient $M$ by $J$, and choose $u\in H^8(M;\mathbb Z)$ with
$\langle u,[M]\rangle=1$.  The top-Chern--Euler identity and the
Poincar\'e--Hopf theorem give
\begin{equation}
 c_4(E)=e(TM)=\chi(M)u=2u.
 \label{hrs8:eq:top-chern-euler}
\end{equation}
For a complex bundle, the second Pontryagin class of the underlying real
bundle is
\[
 p_2(E_{\mathbb R})
 =c_2(E)^2-2c_1(E)c_3(E)+2c_4(E).
\]
Consequently,
\begin{equation}
 p_2(TM)=4u.
 \label{hrs8:eq:p2-complex}
\end{equation}
On the other hand, $H^4(M;\mathbb Z)=0$ gives $p_1(TM)=0$, and the middle
cohomology vanishes, so $\operatorname{sign}(M)=0$.  Hirzebruch's
eight-dimensional signature formula is
\[
 \operatorname{sign}(M)
 =\left\langle\frac{7p_2(TM)-p_1(TM)^2}{45},[M]\right\rangle.
\]
It forces $\langle p_2(TM),[M]\rangle=0$, contradicting
\eqref{hrs8:eq:p2-complex}.  For the standard embedded sphere, the signature
step may instead be replaced by the explicit stable trivialisation
\[
 TS^8\oplus\varepsilon^1_{\mathbb R}\cong\varepsilon^9_{\mathbb R},
\]
which gives $p_2(TS^8)=0$ directly.  Hence
\begin{equation}
 \boxed{S^8\text{ admits no almost-complex structure}.}
 \label{hrs8:eq:no-almost-complex}
\end{equation}
Integrability is stronger than being almost complex, so a complex-fourfold
structure is excluded a fortiori.

\subsubsection{Independent arithmetic checks}
\label{hrs8:subsubsec:arithmetic-checks}

Borel and Serre prove that $S^{2n}$ admits no almost-complex structure for
$n\geq4$, including arbitrary differentiable structures on the sphere.  In
the $S^8$ specialization, their mod-$3$ reduced-power identity is
\[
 \mathcal P^1(c_2)
 =c_1^2c_2-2c_2^2-c_1c_3+4c_4.
\]
All lower Chern classes vanish on $S^8$, so it would follow that
$c_4=0\pmod 3$, whereas \eqref{hrs8:eq:top-chern-euler} gives
$c_4=\pm2u\neq0\pmod3$.

Bott integrality gives a second independent check.  For a complex bundle
over $S^{2n}$, the top Chern number is divisible by $(n-1)!$.  At $n=4$,
\[
 6\mid\langle c_4(E),[S^8]\rangle,
\]
which is incompatible with the Euler number $\pm2$.

\subsubsection{Exact scope of the obstruction}
\label{hrs8:subsubsec:scope}

The obstruction rules out a smooth complex fourfold diffeomorphic to $S^8$;
a proper holomorphic family of complex three-tori over $\mathbb P^1$ whose
smooth total space is diffeomorphic to $S^8$; and an almost-complex structure
on all of $TS^8$, whether or not integrable.  It also rules out a symplectic
or K\"ahler structure: if $\omega$ were symplectic, then
$H^2(S^8;\mathbb R)=0$ would give $\omega=d\alpha$, while
\[
 \int_{S^8}\omega^4
 =\int_{S^8}d(\alpha\wedge\omega^3)=0,
\]
contradicting symplectic volume.

It does not rule out a complex structure on an open subset, a stable complex
structure, a singular complex analytic space whose tangent sheaf is not a
vector bundle, or a real stratified quotient whose underlying boundary is
$S^8$.  The standard sphere is stably complex because
\[
 TS^8\oplus\varepsilon^2_{\mathbb R}\cong\varepsilon^{10}_{\mathbb R},
\]
but this is not a complex structure on the actual rank-eight tangent bundle.
Moreover, the usual link of an isolated complex singularity of complex
dimension $d$ has real dimension $2d-1$.  Thus an $S^8$ obtained as a quotient
boundary is not, merely for that reason, an ordinary complex-analytic link.

\subsection{Requirements for a real Fabel-style \texorpdfstring{$S^8$}{S8} fibration}
\label{hrs8:subsec:real-fibration}

The reconstructed $S^6$ candidate has a map
\[
 f\colon X^6\longrightarrow B\cong\mathbb P^1
\]
whose regular fibre is a complex two-torus, hence a real $T^4$, and whose
three special fibres comprise one non-normal toric degeneration and two
multiple bielliptic fibres.  The direct dimension analogue on a real
eight-manifold would be a map
\[
 g\colon M^8\longrightarrow S^2
\]
with regular fibre $T^6$.  If $M$ is required to be a smooth complex
fourfold and the regular fibres are complex three-tori, the obstruction in
Subsection \ref{hrs8:subsec:no-almost-complex} ends the proposal.  If $g$ is
only a real smooth or stratified map, the proposal remains topologically
meaningful.

The adjective ``singular'' is essential.  There is no fibre bundle
\[
 T^6\longrightarrow S^8\longrightarrow S^2.
\]
Indeed, its long exact homotopy sequence would contain
\[
 \pi_2(S^8)=0\longrightarrow\pi_2(S^2)=\mathbb Z
 \xrightarrow{\partial}\pi_1(T^6)=\mathbb Z^6
 \longrightarrow\pi_1(S^8)=0.
\]
Exactness would require $\partial\colon\mathbb Z\to\mathbb Z^6$ to be
surjective, which is impossible.  The Euler identity for a fibre bundle gives
the same conclusion:
\[
 \chi(S^8)=2\neq\chi(S^2)\chi(T^6)=2\cdot0=0.
\]

\subsection{Euler constraint for a finite multipole model}
\label{hrs8:subsec:euler}

Let $g\colon M^8\to S^2$ be a proper triangulable map with finite
critical-value set $\Sigma=\{p_1,\ldots,p_s\}$, and suppose that every fibre
over $S^2\setminus\Sigma$ is $T^6$.  Euler integration for a constructible
map gives
\begin{align}
 \chi(M)
 &=\chi(S^2\setminus\Sigma)\chi(T^6)
   +\sum_{i=1}^s\chi\bigl(g^{-1}(p_i)\bigr)\notag\\
 &=\sum_{i=1}^s\chi(F_i).
 \label{hrs8:eq:euler-integration}
\end{align}
Equivalently, one may cut out small discs, use multiplicativity on the regular
bundle, and note that each boundary mapping torus has Euler characteristic
zero.  For $M\cong S^8$ one therefore needs
\begin{equation}
 \boxed{\sum_{i=1}^s\chi(F_i)=2.}
 \label{hrs8:eq:euler-fibres}
\end{equation}
Multiple fibres that are finite free quotients of tori have Euler
characteristic zero and cannot account for $S^8$.  At least one special fibre
must contain a genuine collapse or fixed stratum with nonzero Euler
characteristic.  The numerical pattern consisting of two Euler-zero multiple
fibres and one fibre of Euler characteristic two is compatible with
\eqref{hrs8:eq:euler-fibres}, but that numerical compatibility is not a
construction.

\subsection{Fundamental group and monodromy}
\label{hrs8:subsec:pi1}

Put $B^\circ=S^2\setminus\Sigma$, $E=g^{-1}(B^\circ)$, and
$\Lambda=H_1(T^6;\mathbb Z)\cong\mathbb Z^6$.  The regular bundle has
monodromy
\[
 \rho\colon\pi_1(B^\circ)\longrightarrow\operatorname{GL}(\Lambda).
\]
If the fibres and total space are oriented compatibly, its image lies in
$\operatorname{SL}(6,\mathbb Z)$.  For oriented peripheral loops
$\gamma_1,\ldots,\gamma_s$ satisfying $\gamma_1\cdots\gamma_s=1$, the
matrices $M_i=\rho(\gamma_i)$ satisfy
\begin{equation}
 M_1M_2\cdots M_s=I.
 \label{hrs8:eq:monodromy-product}
\end{equation}

When $s\geq1$, $B^\circ$ has $\pi_2(B^\circ)=0$.  The homotopy sequence of
the regular bundle therefore gives the exact extension
\begin{equation}
 1\longrightarrow\Lambda\longrightarrow\pi_1(E)
 \longrightarrow\pi_1(B^\circ)\longrightarrow1.
 \label{hrs8:eq:regular-pi1-extension}
\end{equation}
After choosing a section and a compatible basepoint system,
\eqref{hrs8:eq:regular-pi1-extension} becomes an affine semidirect
presentation; without those choices it remains the displayed exact
extension.

Let $N_i=g^{-1}(D_i)$ be the filling over a small closed disc $D_i$ about
$p_i$, and let $A_i=E\cap N_i$, a $T^6$ mapping torus.  Van Kampen gives
\begin{equation}
 \pi_1(M)\cong
 \pi_1(E)\Big/
 \Big\langle\!\Big\langle
  \ker\bigl(\pi_1(A_i)\to\pi_1(N_i)\bigr),\ 1\leq i\leq s
 \Big\rangle\!\Big\rangle.
 \label{hrs8:eq:van-kampen-fillings}
\end{equation}
Thus an $S^8$ construction must derive the actual local inclusion maps and
prove that their kernels normally generate all of
\eqref{hrs8:eq:regular-pi1-extension}.  Monodromy matrices or Smith forms
alone do not prove that \eqref{hrs8:eq:van-kampen-fillings} is trivial.

\subsection{Integral homology: the exact gluing tests}
\label{hrs8:subsec:homology}

Replace $B^\circ$ by a compact sphere with $s$ open discs removed and replace
$E,N_i$ by compact collar-compatible pieces.  Define the actual
Mayer--Vietoris maps
\begin{equation}
 \alpha_q=
 \bigoplus_{i=1}^s\bigl((j_{E,i})_*,-(j_{N,i})_*\bigr)\colon
 \bigoplus_iH_q(A_i;\mathbb Z)\longrightarrow
 H_q(E;\mathbb Z)\oplus\bigoplus_iH_q(N_i;\mathbb Z).
 \label{hrs8:eq:mv-map}
\end{equation}
Because the pieces have nonempty boundary, their top homology vanishes.  The
Mayer--Vietoris sequence and
$H_q(S^8;\mathbb Z)=0$ for $1\leq q\leq7$ imply the following necessary and
sufficient conditions:
\begin{equation}
\begin{array}{c|l}
q&\text{condition on }\alpha_q\\ \hline
0&\alpha_0\text{ is injective and }\operatorname{coker}\alpha_0\cong\mathbb Z,\\
1\leq q\leq6&\alpha_q\text{ is an isomorphism},\\
7&\alpha_7\text{ is surjective and }\ker\alpha_7\cong\mathbb Z.
\end{array}
\label{hrs8:eq:mv-tests}
\end{equation}
The generator of $\ker\alpha_7$ is the boundary of the fundamental class of
the closed eight-manifold obtained by gluing the pieces.  Equations
\eqref{hrs8:eq:mv-map}--\eqref{hrs8:eq:mv-tests}, using the actual inclusion
maps, are the integral replication test for a proposed $S^8$ multipole
construction.

Each boundary group is constrained by monodromy.  If $A_i$ is the mapping
torus of $\phi_i\colon T^6\to T^6$ and
$M_i=(\phi_i)_*|_{H_1(T^6;\mathbb Z)}$, the Wang sequence gives
\begin{equation}
 0\longrightarrow
 \operatorname{coker}\bigl(I-\textstyle\bigwedge^qM_i\bigr)
 \longrightarrow H_q(A_i;\mathbb Z)
 \longrightarrow
 \ker\bigl(I-\textstyle\bigwedge^{q-1}M_i\bigr)
 \longrightarrow0.
 \label{hrs8:eq:wang}
\end{equation}
The right group is free, so \eqref{hrs8:eq:wang} splits as an abelian group,
although not canonically.  Since the regular base retracts to a wedge of
$s-1$ circles, the homological Serre spectral sequence gives the filtration
\begin{equation}
 0\longrightarrow
 H_0\!\left(B^\circ;\textstyle\bigwedge^q\Lambda\right)
 \xrightarrow{\ \iota_q\ } H_q(E;\mathbb Z)
 \xrightarrow{\ \pi_q\ }
 H_1\!\left(B^\circ;\textstyle\bigwedge^{q-1}\Lambda\right)
 \longrightarrow0.
 \label{hrs8:eq:serre-regular-part}
\end{equation}
To type both edge terms explicitly, put
\(V_q=\bigwedge^q\Lambda\), choose oriented free generators
\(\gamma_1,\ldots,\gamma_{s-1}\) of \(\pi_1(B^\circ)\), and write
\(M_j=\rho(\gamma_j)\).  In the Wang sign convention the cellular differential is
\begin{equation}
 d_q\colon\bigoplus_{j=1}^{s-1}V_q\longrightarrow V_q,\qquad
 (v_j)_j\longmapsto
 \sum_{j=1}^{s-1}\left(I-\textstyle\bigwedge^qM_j\right)v_j,
 \label{hrs8:eq:serre-cellular-differential}
\end{equation}
so the two homological edge groups are
\begin{equation}
 H_0(B^\circ;V_q)
 =\operatorname{coker}d_q
 =(V_q)_\rho
 =\frac{V_q}{
   \sum_{j=1}^{s-1}(I-\bigwedge^qM_j)V_q},
 \qquad
 H_1(B^\circ;V_{q-1})=\ker d_{q-1}.
 \label{hrs8:eq:serre-coinvariants-kernel}
\end{equation}
Thus the coinvariants are the Serre filtration subobject and the
base-one-cycle group is the quotient.  Since \(\ker d_{q-1}\) is a subgroup
of a free abelian group, it is free, and
\eqref{hrs8:eq:serre-regular-part} splits noncanonically as an abstract
abelian group.

Dually, set \(V_q^\vee=\bigwedge^q\Lambda^\vee\); its monodromy along
\(\gamma_j\) is \(\bigwedge^q(M_j^{-\mathsf T})\).  The cohomological
filtration is
\begin{equation}
 0\longrightarrow H^1(B^\circ;V_{q-1}^\vee)
 \longrightarrow H^q(E;\mathbb Z)
 \longrightarrow H^0(B^\circ;V_q^\vee)\longrightarrow0,
 \qquad
 \begin{aligned}
 H^0(B^\circ;V_q^\vee)
   &=(V_q^\vee)^\rho\\
   &=\bigcap_{j=1}^{s-1}
     \ker\!\left(I-\textstyle\bigwedge^q(M_j^{-\mathsf T})\right).
 \end{aligned}
 \label{hrs8:eq:serre-cohomology-invariants}
\end{equation}
Accordingly, each boundary inclusion in \(\alpha_q\) must send its Wang
coinvariant filtration term through \(\iota_q\), and its base-one-cycle
term must be tested after \(\pi_q\); neither exact sequence supplies a
canonical splitting.  The corrected sequence identifies \(H_q(E;\mathbb Z)\)
noncanonically with the direct sum of its two edge groups, so a rank or
torsion count made from that direct sum survives.  The natural edge maps do
not: every map-level Mayer--Vietoris test must use
\eqref{hrs8:eq:serre-regular-part} in the displayed order.  Equations
\eqref{hrs8:eq:mv-map}--\eqref{hrs8:eq:serre-cohomology-invariants} retain
every lattice, monodromy, kernel, coinvariant, invariant, and inclusion
that a correct proof must compute.

\subsection{The torus-action bound and a maximal real orbit model}
\label{hrs8:subsec:torus-action}

For a smooth torus action on a compact manifold,
\[
 \chi(M)=\chi(M^T).
\]
Since $\chi(S^8)=2$, every torus action on $S^8$ has a fixed point.  Equip
$S^8$ with a $T$-invariant Riemannian metric.  At a fixed point the isotropy
representation is a sum of four real two-dimensional weight planes and
possibly trivial summands.  If the global action is effective, this tangent
representation is faithful: an isometry fixing a point with identity
derivative is the identity on the connected manifold.  Four weights span a
character lattice of rank at most four.  Hence
\begin{equation}
 \boxed{\operatorname{rank}T\leq4.}
 \label{hrs8:eq:torus-rank}
\end{equation}
In particular, a generic $T^6$ fibration of $S^8$ cannot be the orbit map of
a global effective torus action.  This does not rule out a monodromic
$T^6$-bundle over the regular part of $S^2$, because such a bundle need not
supply a global $T^6$ action.

The bound is sharp.  Write
\[
 S^8=\left\{(z_1,z_2,z_3,z_4,y)\in\mathbb C^4\oplus\mathbb R:
       \sum_{i=1}^4|z_i|^2+y^2=1\right\},
\]
and let $T^4$ rotate the four complex coordinates.  The orbit map is
\[
 q(z_1,z_2,z_3,z_4,y)
 =(|z_1|,|z_2|,|z_3|,|z_4|,y),
\]
with image
\[
 Q=\left\{(r_1,r_2,r_3,r_4,y):r_i\geq0,
       \sum r_i^2+y^2=1\right\}.
\]
Two points have the same image exactly when they lie in the same orbit, so
compact--Hausdorff descent gives $S^8/T^4\cong Q$.  Away from the two poles,
write $r_i=\sqrt{1-y^2}\,u_i$.  Then
\[
 (u_1,u_2,u_3,u_4)\in
 \{u_i\geq0,\ \sum u_i^2=1\}\cong\Delta^3,
\]
and the values $y=\pm1$ collapse the two ends.  Therefore
\begin{equation}
 \boxed{S^8/T^4\cong\operatorname{Susp}(\Delta^3)\cong B^4.}
 \label{hrs8:eq:t4-quotient}
\end{equation}
The locus on which all $r_i>0$ has free $T^4$ orbit.  Exactly $j$ zero
complex coordinates give isotropy $T^j$, and the two poles are fixed by the
whole torus.  This is a rigorous two-pole singular torus compactification of
$S^8$, but it has four-dimensional quotient and generic $T^4$ orbit.  It is
not a three-special-fibre map to $S^2$.

\subsection{A three-column Gram quotient whose quotient boundary is \texorpdfstring{$S^8$}{S8}}
\label{hrs8:subsec:gram}

Let
\[
 V=\operatorname{Mat}_{3\times3}(\mathbb C),
 \qquad \dim_{\mathbb R}V=18,
\]
and let $U(3)$ act on the left.  Define
\begin{equation}
 \Gamma\colon V\longrightarrow\operatorname{Herm}_3(\mathbb C)_{\geq0},
 \qquad \Gamma(A)=A^*A.
 \label{hrs8:eq:gram-map}
\end{equation}
The map is surjective because every $Q\geq0$ is the Gram matrix of
$Q^{1/2}$.  It separates $U(3)$-orbits.  Indeed, if $A^*A=B^*B$, the rule
$Ax\mapsto Bx$ is a well-defined isometry from $\operatorname{im}A$ to
$\operatorname{im}B$: equality of the Gram matrices implies both equality of
kernels and preservation of the Hermitian inner product.  Extend this
isometry to a unitary $U$ of $\mathbb C^3$; then $B=UA$.  Hence
\begin{equation}
 V/U(3)\cong\operatorname{Herm}_3(\mathbb C)_{\geq0}.
 \label{hrs8:eq:gram-orbit-space}
\end{equation}

For the Frobenius unit ball
$\mathbb B(V)=\{A:\operatorname{tr}(A^*A)\leq1\}$,
\eqref{hrs8:eq:gram-map} gives
\begin{equation}
 \mathbb B(V)/U(3)
 \cong K:=\{Q\in\operatorname{Herm}_3(\mathbb C):Q\geq0,
             \operatorname{tr}Q\leq1\}.
 \label{hrs8:eq:gram-ball-quotient}
\end{equation}
The real vector space of Hermitian $3\times3$ matrices has dimension nine.
The set $K$ is compact, convex, and has nonempty interior
$\{Q>0,\operatorname{tr}Q<1\}$.  Radial projection from an interior point
gives a homeomorphism of pairs $(K,\partial K)\cong(B^9,S^8)$, and therefore
\begin{equation}
 \boxed{\partial\bigl(\mathbb B^{18}/U(3)\bigr)\cong S^8.}
 \label{hrs8:eq:gram-boundary}
\end{equation}

Two qualifications are essential.  First, the quotient of the source sphere
is not the boundary in \eqref{hrs8:eq:gram-boundary}:
\[
 S^{17}/U(3)\cong\{Q\geq0:\operatorname{tr}Q=1\}.
\]
The space on the right is an eight-dimensional compact convex body and hence
is homeomorphic to $B^8$.  The rest of $\partial K$ consists of
rank-deficient Gram matrices with $\operatorname{tr}Q<1$.  Thus taking the
boundary and taking the orbit quotient do not commute in this construction.

Second, strict $SU(3)$ is different.  For invertible matrices, $A^*A$ does
not record the phase of $\det A$, while left multiplication by $SU(3)$
preserves $\det A$.  The strict quotient is described by
\begin{equation}
 (Q,d)=(A^*A,\det A),
 \qquad Q\geq0,
 \qquad |d|^2=\det Q.
 \label{hrs8:eq:strict-su3-quotient}
\end{equation}
It has generic real dimension ten, not nine.  The construction in
\eqref{hrs8:eq:gram-boundary} therefore uses
\[
 U(3)\cong(SU(3)\times U(1))/\mu_3,
\]
not colour $SU(3)$ alone.  It is a precise $S^8$ real quotient boundary, but
no morphism has been produced from its Gram or rank strata to the period
matrix, torus bundle, logarithmic twists, or three special fibres of the
Fabel--Alp\"oge construction.

\subsection{Retaining the \texorpdfstring{$T^4/\mathbb R^4$}{T4/R4} fibre geometry and placing \texorpdfstring{$SU(3)$}{SU(3)} on the regular locus}
\label{hrs8:subsec:su3-regular-locus}

This subsection classifies the required topological bundle data and its
obstruction to crossing the singular fibres.  It makes no physical
identification.

\subsubsection{The regular locus and its fourth cohomology}
\label{hrs8:subsubsec:regular-cohomology}

Let $X\cong_{\mathrm{diff}}S^6$ be the reconstructed candidate and put
\[
 D=W\sqcup S_1\sqcup S_2,
 \qquad U=X\setminus D=f^{-1}(B^\circ),
 \qquad B^\circ=\mathbb P^1\setminus\{p_0,p_1,p_2\}.
\]
Here $W$ is the non-normal central fibre and $S_1,S_2$ are the underlying
smooth bielliptic surfaces of the two multiple fibres.  The reconstructed
integral calculations give
\[
 H_1(W;\mathbb Z)\cong\mathbb Z^2,
 \qquad H_1(S_1;\mathbb Z)\cong H_1(S_2;\mathbb Z)\cong\mathbb Z^2.
\]
The compact analytic set $D$ is triangulable, so Alexander duality in
$X\cong S^6$ applies:
\[
 \widetilde H^i(U;\mathbb Z)
 \cong\widetilde H_{5-i}(D;\mathbb Z).
\]
At $i=4$, it gives
\begin{equation}
 \boxed{H^4(U;\mathbb Z)\cong H_1(D;\mathbb Z)\cong\mathbb Z^6.}
 \label{hrs8:eq:h4-regular-locus}
\end{equation}

There is an independent fibrewise description.  The pair of pants
$B^\circ$ retracts to a wedge of two circles, and the $T^4$ bundle retracts
to its restriction over that graph.  Hence $U$ has the homotopy type of a CW
complex of dimension at most five.  If $\mathcal V^q=R^qf_*\mathbb Z$ is the
monodromy local system, the Serre spectral sequence has only columns zero and
one and yields
\begin{equation}
 0\longrightarrow H^1(B^\circ;\mathcal V^3)
 \longrightarrow H^4(U;\mathbb Z)
 \xrightarrow{\operatorname{res}_F}
 H^0(B^\circ;\mathcal V^4)
 \longrightarrow0.
 \label{hrs8:eq:h4-serre}
\end{equation}
Holomorphic fibre monodromy preserves the real orientation, so
$H^0(B^\circ;\mathcal V^4)\cong\mathbb Z$.  Thus one coordinate of
\eqref{hrs8:eq:h4-regular-locus} is the constant evaluation on a smooth
$T^4$ fibre; the remaining rank-five kernel in \eqref{hrs8:eq:h4-serre}
consists of mixed base--fibre classes.

\subsubsection{Exact classification of \texorpdfstring{$SU(3)$}{SU(3)} bundles on \texorpdfstring{$U$}{U}}
\label{hrs8:subsubsec:su3-classification}

Let $P\to U$ be a principal $SU(3)$ bundle and
$E_P=P\times_{SU(3)}\mathbb C^3$ its determinant-trivial rank-three complex
bundle.  Through degree five,
\[
 \pi_i(BSU(3))=0\quad(i=1,2,3,5),
 \qquad \pi_4(BSU(3))\cong\mathbb Z,
\]
with universal generator $c_2$.  Because $U$ has CW dimension at most five,
obstruction theory gives a bijection
\begin{equation}
 \boxed{
 \{\text{principal }SU(3)\text{ bundles on }U\}/\!\cong
 \ \xrightarrow{\ c_2\ }\ H^4(U;\mathbb Z)\cong\mathbb Z^6.}
 \label{hrs8:eq:su3-classification}
\end{equation}
The fibre integer
\begin{equation}
 k_F=\left\langle c_2(E_P)|_{F_b},[F_b]\right\rangle,
 \qquad b\in B^\circ,
 \label{hrs8:eq:fibre-instanton-number}
\end{equation}
is independent of $b$.  A nonzero $k_F$ is only one possible obstruction: a
class with $k_F=0$ may still have a nonzero mixed component in the kernel of
\eqref{hrs8:eq:h4-serre}.

\subsubsection{Extension across the three special fibres}
\label{hrs8:subsubsec:extension}

The cohomology sequence of the pair $(X,U)$ contains
\[
 H^4(X;\mathbb Z)\longrightarrow H^4(U;\mathbb Z)
 \xrightarrow{\delta}H^5(X,U;\mathbb Z)
 \longrightarrow H^5(X;\mathbb Z).
\]
Both outer groups vanish because $X\cong S^6$.  Therefore
\begin{equation}
 \boxed{\delta\colon H^4(U;\mathbb Z)
 \xrightarrow{\ \cong\ }H^5(X,U;\mathbb Z).}
 \label{hrs8:eq:relative-delta}
\end{equation}
Every $SU(3)$ bundle on all of $X$ has $c_2=0$, so its restriction to $U$ is
trivial by \eqref{hrs8:eq:su3-classification}.  Conversely, the trivial
bundle on $U$ extends as the trivial bundle on $X$.  Thus, for unframed
topological bundles,
\begin{equation}
 \boxed{P\to U\text{ extends across }D
 \iff c_2(P)=0
 \iff\delta c_2(P)=0.}
 \label{hrs8:eq:extension-criterion}
\end{equation}

Under Alexander duality, \eqref{hrs8:eq:relative-delta} identifies the
obstruction with a six-integer class supported on the one-cycles of
$W\sqcup S_1\sqcup S_2$.  Because the three components are disjoint and
$H_1(D)$ is their direct sum, nonzero components do not vanish merely by
declaring that charges at distinct special fibres cancel.  They cancel as an
extension obstruction only when the complete class in $\mathbb Z^6$ is zero.

This is a statement about the bundle.  Even when $c_2=0$, a specified
connection on the trivial bundle may fail to extend smoothly because of its
holonomy or growth near $D$; this is an analytic boundary-value problem not
decided by \eqref{hrs8:eq:extension-criterion}.

\subsubsection{The distinct six-dimensional clutching charge}
\label{hrs8:subsubsec:c3-clutching}

Principal $SU(3)$ bundles on the complete sphere are classified by
\[
 [S^6,BSU(3)]\cong\pi_5(SU(3))\cong\mathbb Z.
\]
If $k$ is the clutching integer, Bott integrality fixes the top Chern class as
\begin{equation}
 c_3(E_P)=\pm2k\,u\in H^6(S^6;\mathbb Z).
 \label{hrs8:eq:c3-clutching}
\end{equation}
Every such global bundle restricts to the trivial bundle on $U$, because its
restricted $c_2$ vanishes and \eqref{hrs8:eq:su3-classification} is complete.
Thus a global six-dimensional $SU(3)$ clutching charge can be concentrated in
the gluing across the deleted fibres, but it is a
$c_3/\pi_5(SU(3))$ class, not a four-dimensional $c_2$ instanton class on the
regular fibres.

Alexander duality also gives
$H^5(U;\mathbb Z)\cong\widetilde H_0(D;\mathbb Z)\cong\mathbb Z^2$.
The pair sequence therefore yields
\[
 0\longrightarrow\mathbb Z^2
 \longrightarrow H^6(X,U;\mathbb Z)
 \longrightarrow\mathbb Z
 \longrightarrow0,
\]
and hence
\begin{equation}
 H^6(X,U;\mathbb Z)\cong\mathbb Z^3.
 \label{hrs8:eq:h6-relative}
\end{equation}
After choosing the complex orientations of the three fibre components, the
three relative top classes map to the global fundamental class by the
augmentation
\begin{equation}
 (n_0,n_1,n_2)\longmapsto n_0+n_1+n_2.
 \label{hrs8:eq:augmentation}
\end{equation}
The rank-two kernel records redistribution among the three supports with the
same global total.  To promote
\eqref{hrs8:eq:h6-relative}--\eqref{hrs8:eq:augmentation} from exact
cohomological data to a relative $SU(3)$ bundle theorem with prescribed framings,
one must give the three clutching maps and compute their homotopy classes;
arbitrary relative cohomology integers are not silently declared realised.

\subsubsection{Data required before a gauge-theory conclusion}
\label{hrs8:subsubsec:gauge-data}

A mathematical colour or gauge proposal on $U$ must specify:
\begin{enumerate}
 \item a principal $SU(3)$ bundle $P\to U$, equivalently its complete class
       $c_2(P)\in\mathbb Z^6$;
 \item a connection $A$ and a fixed trace convention for its curvature
       Chern--Weil form;
 \item tubular boundary conditions or framings at each of $W,S_1,S_2$;
 \item whether the invariant under study is the fibre number
       \eqref{hrs8:eq:fibre-instanton-number}, the extension defect
       \eqref{hrs8:eq:relative-delta}, or the six-dimensional clutching
       number \eqref{hrs8:eq:c3-clutching}; and
 \item an explicit analytic statement about a Yang--Mills operator, action,
       or spectrum.
\end{enumerate}
The topology alone proves no mass gap, confinement statement, particle
identification, or Yang--Mills dynamics.  It supplies a carrier and an exact
obstruction ledger on which such a calculation can be posed.

\subsection{Consequences for the higher-rung search}
\label{hrs8:subsec:consequences}

Replacing $S^6$ by $S^8$ does not permit a literal smooth complex-fourfold
version of the Fabel--Alp\"oge construction.  A real $T^6$-over-$S^2$
singular fibration could still exist, but it cannot arise as a global
effective torus action and must pass the map-level tests
\eqref{hrs8:eq:euler-fibres}, \eqref{hrs8:eq:van-kampen-fillings}, and
\eqref{hrs8:eq:mv-tests}.  The standard $T^4$ action gives an exact and
economical real $S^8$ multipole model, but its base is $B^4$, not $S^2$, and
its singular locus is a face stratification rather than three isolated
fibres.

The $3\times3$ $U(3)$ Gram construction gives an exact quotient boundary
$S^8$ and is the strongest precise realisation presently obtained for the
``one rung higher'' intuition.  Strict $SU(3)$ retains the determinant phase
and does not give the same quotient.  Keeping the original $T^4$ fibres and
placing $SU(3)$ on $X\setminus D$ is instead mathematically well typed: its
unframed bundle classes form $\mathbb Z^6$, every nonzero $c_2$ is an exact
nonextendable divisor defect, and global $SU(3)$ bundles on $S^6$ carry the
even-valued $c_3$ charge in \eqref{hrs8:eq:c3-clutching} while restricting
trivially to the regular locus.

None of these statements identifies a lepton, quark, or four-dimensional
Yang--Mills theory.

\paragraph{Sources.}
The almost-complex obstruction is cross-checked against A.~Borel and
J.-P.~Serre, \emph{Groupes de Lie et puissances r\'eduites de Steenrod},
Amer. J. Math. \textbf{75} (1953), 409--448, Proposition~15.1,
\texttt{doi:10.2307/2372495}, and P.~Konstantis and M.~Parton,
\emph{Almost complex structures on spheres}, Differential Geom. Appl.
\textbf{57} (2018), 10--22,
\texttt{doi:10.1016/j.difgeo.2017.10.010}.  The fixed-point identity used in
\eqref{hrs8:eq:torus-rank} is from S.~Kobayashi, \emph{Fixed points of
isometries}, Nagoya Math. J. \textbf{13} (1958), 63--68,
\texttt{doi:10.1017/S0027763000023497}.  For singular torus fibrations and
the coordinate sphere action, see respectively Y.~Matsumoto, \emph{Torus
fibrations over the 2-sphere with the simplest singular fibers}, J. Math.
Soc. Japan \textbf{37} (1985), 605--636,
\texttt{doi:10.2969/jmsj/03740605}, and H.~Ishida, Y.~Fukukawa, and M.~Masuda,
\emph{Topological toric manifolds}, Moscow Math. J. \textbf{13} (2013),
57--98, \texttt{doi:10.17323/1609-4514-2013-13-1-57-98}.  The theorem-level
calculations above are given in full and do not rely on these citations as
substitutes for omitted morphisms.

% Mechanically generated from s6_higher_rung_24d_preprint.tex.
% Complete proofs retained; headings and labels adapted for inclusion.
% The root's shared lattice and bibliography inputs are supplied by the host.
\section{The ternary triality structure and its exact carrier morphisms}
\label{sec:h24-continuation}
The complete proof source is also available in the root draft
\path{s6_higher_rung_24d_preprint.tex}. The higher target is a
three-directional structure encoded by the ordered triality modules and
their invariant trilinear tensor. The proofs retain the Peirce coefficient
$1/2$, determinant coefficient $2$, spectral graph, flag-cell inverse and
degree-$+1$ collapse, and the full retained cusp-angle calculation.
The complex-toroidal calculations concern an auxiliary comparison and
do not define the higher target.

The single shared source, beginning at Section~\ref{esn:sec:shared},
retains the original modulo-$840$ lattice, full Weyl quotient, exact
index-six Leech neighbour and cyclic-triality isometry. Its continuation
now computes the fixed-Cayley-frame value groups for
$F=\tau\circ\mathcal L$, with $\tau(q)=\operatorname{Re}((q_1q_2)q_3)$, whose
inclusion has quotient $\mathbb Z/12$; the Albert determinant contains
$2F$, not $F$. It also gives the common Leech intersection $J$ of
determinant $81$ and index $9$, and its fixed fourth $A_2^{12}$ extension
with $72$ roots. Wilson's construction \cite{esn:Wilson2009} supplies
the cyclic fixed Gram matrix $3G_{E_8}$, the $A_2\otimes E_8$
complement, the quotient $\mathbb F_3^8$, and the exact rank-four
literal intersections with the two marked lattice images. All metrics,
cycle directions and the distinction between the cubic image and its
value-generated subgroup remain explicit in the shared proofs.

\subsection{The source construction and the toroidal comparison}

The higher-rung investigation asks how a structure with three interacting
directions can extend the source geometry. The present algebraic datum is
the vector and two half-spin sectors of triality, together with their
invariant trilinear tensor. The complex-torus extension studied first below
is one specified comparison with the source construction. Its complex
structure is a hypothesis of those toroidal theorems; it is not a
requirement on the higher-rung target. The retained cusp-angle map and the
marked $A_5^4D_4$ lattice are retained as inputs to the subsequent comparison
maps.

We use the following exact input from the reconstructed candidate
construction \cite{h24:AlpogeS6_2026}.  Let
\[
 B=\mathbb P^1,\qquad \Sigma=\{p_0,p_1,p_2\},\qquad
 B^\circ=B\setminus\Sigma,
\]
where $p_0$ is a cusp and $p_1,p_2$ have orbifold orders $3,4$.
There is a claimed compact connected complex threefold and proper
holomorphic map
\[
 f:X\longrightarrow B,
 \qquad X\cong_{\mathrm{diff}}S^6,
\]
whose regular fibres are complex two-tori.  In the basis
$(\gamma,u,w,\delta)$ of $V=\mathbb Z^4$ the integral monodromies are generated by
\begin{equation}
 T_1=
 \begin{pmatrix}1&0&-6&2\\0&-1&1&1\\0&-1&0&1\\0&0&0&1\end{pmatrix},
 \qquad
 T_2=
 \begin{pmatrix}1&6&0&-3\\0&0&-1&1\\0&1&0&0\\0&0&0&1\end{pmatrix},
 \label{h24:eq:original-T}
\end{equation}
with $T_1^3=T_2^4=I$ and
$T_0=(T_1T_2)^{-1}=I+N$, $N^2=0$.  On the dual lattice
$\Lambda=V^*$ the monodromies are
$A_j=(T_j^{-1})^{\mathsf T}$.  The period matrix is
\begin{equation}
 \Pi(z)=
 \begin{pmatrix}6\mu(z)&\tau(z)&1&0\\
                 \beta(z)&\mu(z)&0&1\end{pmatrix},
 \qquad F_z=\mathbb C^2/\Pi(z)\Lambda .
 \label{h24:eq:original-period}
\end{equation}
At $p_0$ the filling is a Mumford toric quotient with reduced irreducible
normal-crossing central fibre $W$: its normalization is the degree-six
del Pezzo surface, opposite sides of its anticanonical hexagon are
identified, and $\chi(W)=2$.  At $p_j$, $j=1,2$, the filling is a
logarithmic transform with
\[
 f^*p_1=3S_1,\qquad f^*p_2=4S_2,
\]
where the $S_j$ are smooth bielliptic surfaces and hence have Euler
characteristic zero.  The primitive affine translations have
\begin{equation}
 v_1=\widehat\gamma+2\widehat u-4\widehat w,\qquad
 v_2=-\widehat\gamma-3\widehat u+3\widehat w,\qquad
 (\ell_0,\ell_1,\ell_2)=(0,1,-1).
 \label{h24:eq:original-twists}
\end{equation}
They give
$\pi_1(X)=\mathbb Z/|12\ell_0-4\ell_1-3\ell_2|=0$ and the integral homology of
$S^6$.

The later marked observation uses the integral difference matrix and its
fixed right inverse
\begin{equation}
 P=\begin{pmatrix}1&-1&0\\0&1&-1\end{pmatrix},\quad
 u=\begin{pmatrix}1\\1\\1\end{pmatrix},\quad
 C=PP^*=\begin{pmatrix}2&-1\\-1&2\end{pmatrix},\quad
 R=P^*C^{-1}=\frac13
 \begin{pmatrix}2&1\\-1&1\\-1&-2\end{pmatrix}.
 \label{h24:eq:marked-P}
\end{equation}
Thus $PR=I_2$, $Pu=0$, $u^*R=0$, and for $X_P=(R\;u)$,
\begin{equation}
 X_P^{-1}=\begin{pmatrix}P\\u^*/3\end{pmatrix},
 \qquad \det X_P=1,
 \qquad X_P^*X_P=\operatorname{diag}(C^{-1},3).
 \label{h24:eq:marked-inverse}
\end{equation}

\begin{definition}[Data for the complex-toroidal comparison]
\label[definition]{h24:def:feature-ledger}
A \emph{full higher toroidal wrap} consists of all the following data and
proofs.  Neither its dimension nor its homeomorphism type is prescribed.
\begin{enumerate}[label=\textup{(F\arabic*)},leftmargin=3.2em]
\item A compact connected smooth carrier obtained from the stated pieces,
      with its smooth and homotopy type computed from the attaching maps.
\item A geometric tangent structure on the whole carrier.  For a literal
      toroidal continuation this is an integrable complex structure; a
      different category must state its replacement rather than inherit the
      word ``complex''.
\item A proper holomorphic map to the same one-dimensional base, with a
      complex torus of the dimension forced by total dimension as regular
      fibre.
\item Specified local spaces over the cusp and the two finite points,
      including their multiplicities, reduced fibres, attaching maps, and
      Euler characteristics.
\item An integral period lattice, monodromy representation, affine
      translations, and actual specialization maps at every filling.
\item The van Kampen and integral Mayer--Vietoris maps proving the claimed
      smooth sphere topology.
\item A marked observation whose domain, codomain, lattice map, and
      compatibility with periods and fillings are specified.
\item Rigidity: after all markings above are fixed, no auxiliary metric,
      phase section, fan identification, or gluing choice is left
      unrecorded, and existence and uniqueness up to the stated notion of
      isomorphism are proved.
\end{enumerate}
This definition specifies the complex-toroidal comparison studied in the
following sections. The three-sector triality construction is developed
through its own explicitly defined products, bundles, boundary, and lattice
maps below; the list above does not constrain it to be another complex
structure.
\end{definition}

For a rank-three successor over the regular base, require a complex-linear
surjection $P_z:\mathbb C^3\to\mathbb C^2$ with $K_{\mathbb C,z}=\ker P_z$, period embeddings
$\widetilde\Pi_z:\widetilde\Lambda\hookrightarrow\mathbb C^3$ and
$\Pi_z:\Lambda\hookrightarrow\mathbb C^2$, and
\[
 P_z\widetilde\Pi_z=\Pi_z\Phi,\qquad
 K_{\mathbb C,z}\cap\widetilde\Pi_z(\widetilde\Lambda)
 =\widetilde\Pi_z\kappa(K_\mathbb Z).
\]
Put
\[
 E_z=K_{\mathbb C,z}/\widetilde\Pi_z\kappa(K_\mathbb Z),\qquad
 \widetilde F_z=\mathbb C^3/\widetilde\Pi_z(\widetilde\Lambda),\qquad
 F_z=\mathbb C^2/\Pi_z(\Lambda).
\]
The desired portion of the map diagram is then
\begin{equation}
\begin{CD}
0@>>>K_\mathbb Z@>{\kappa}>>\widetilde\Lambda@>{\Phi}>>\Lambda@>>>0\\
@.@V{\widetilde\Pi_z\kappa}VV@V{\widetilde\Pi_z}VV@V{\Pi_z}VV@.\\
0@>>>K_{\mathbb C,z}@>>>\mathbb C^3@>{P_z}>>\mathbb C^2@>>>0\\
@.@V{q_K}VV@V{q_3}VV@V{q_2}VV@.\\
0@>>>E_z@>>>\widetilde F_z@>{\overline P_z}>>F_z@>>>0.
\end{CD}
\label{h24:eq:required-diagram}
\end{equation}
All three rows are exact.  The maps $\overline P_z$ must assemble to a map
$\widetilde X^\circ\to X^\circ$ over $B^\circ$, together with commuting
specialization squares at all three deleted points.
Sections \ref{h24:sec:torus-successor} and \ref{h24:sec:sphere-ladder} prove exactly
which parts of this diagram exist.

\subsection{Sphere and Euler obstructions}

\begin{theorem}[The $S^8$ tangent obstruction with all hypotheses]
\label[theorem]{h24:thm:s8-obstruction}
Let $M$ be a closed smooth manifold homotopy equivalent to $S^8$.  There is
no bundle endomorphism $J:TM\to TM$ satisfying $J^2=-\operatorname{id}$.  Consequently
no smooth manifold in that homotopy type is a complex fourfold, and in
particular no proper holomorphic complex-three-torus fibration over
$\mathbb P^1$ can have such a total space.
\end{theorem}

\begin{proof}
Suppose $J$ exists and let $E=(TM,J)$, a complex rank-four bundle.  Orient
$M$ by $J$ and choose $a\in H^8(M;\mathbb Z)$ with
$\langle a,[M]\rangle=1$.  Since the cohomology below the top degree
vanishes,
\[
 c_1(E)=c_2(E)=c_3(E)=0.
\]
The top Chern class of a complex bundle is the Euler class of its underlying
oriented real bundle.  Poincar\'e--Hopf and $\chi(M)=2$ therefore give
\[
 c_4(E)=e(TM)=2a.
\]
For the underlying real bundle,
\[
 p_2(E_{\mathbb R})=c_2(E)^2-2c_1(E)c_3(E)+2c_4(E)=4a,
\]
whereas $p_1(TM)=0$ because $H^4(M;\mathbb Z)=0$.  The middle cohomology of $M$
vanishes, so $\operatorname{sign}(M)=0$.  Hirzebruch's degree-eight
signature formula now gives
\[
 0=\operatorname{sign}(M)
 =\left\langle\frac{7p_2(TM)-p_1(TM)^2}{45},[M]\right\rangle
 =\frac{28}{45},
\]
a contradiction.  The last displayed rational number is not being asserted
to be a possible signature: its nonintegrality is precisely the
contradiction forced by the incompatible characteristic classes.
\end{proof}


\begin{theorem}[Parity and the complete higher-sphere exclusion]
\label[theorem]{h24:thm:all-higher-spheres}
Let $M$ be a smooth homotopy $m$-sphere with $m>6$.
Then $M$ admits no almost-complex structure and hence no integrable complex
structure.  In particular $S^9$ cannot be a complex carrier, and there is
no different higher sphere dimension to which the literal complex part of
the candidate can be moved.
\end{theorem}

\begin{proof}
If $m$ is odd, an endomorphism $J$ with $J^2=-I$ would make each real tangent
space a complex vector space, whose real dimension is even.  Hence $m$ must
be even.  Write $m=2n$; then $m>6$ gives $n\ge4$.

If $E=(TM,J)$ existed, $c_i(E)=0$ for $i<n$ and, after using the complex
orientation, $c_n(E)=2a$ for the integral top generator $a$.  Pull $E$ back
along a homotopy equivalence $S^{2n}\to M$.  Bott integrality says that the
top Chern class of every complex bundle on $S^{2n}$ is divisible by
 $(n-1)!$ \cite[Chapter~17, Section~7, Theorem~7.2,
 pp.~274--275]{h24:Husemoller1994}.  Since $(n-1)!\ge6$ does not divide $2$,
 this contradicts $c_n(E)=2a$.  Borel and Serre prove the same conclusion for every
differentiable structure on $S^{2n}$ in \cite[Section~15,
Proposition~15.1 and the following remark, p.~434]{h24:BorelSerre1953}; their
universal reduced-power input is \cite[Section~12, Propositions~12.4 and
12.6 and Lemma~12.5, p.~430]{h24:BorelSerre1953}.
\end{proof}

\begin{proposition}[The direct rank-three Euler mismatch]
\label[proposition]{h24:prop:euler-mismatch}
Consider the direct rank-three toroidal local model obtained by coning the
periodic Kuhn triangulation of $\mathbb R^3$ and dividing by $\mathbb Z^3$.  Its central
fibre $W_3$ has
\[
 \chi(W_3)=6.
\]
If a proper triangulable map $\widetilde f:\widetilde X^8\to S^2$ has this
as one special fibre, has precisely two other special fibres that are free
finite quotients of $T^6$, and is a $T^6$-bundle elsewhere, then
$\chi(\widetilde X)=6$.  Therefore the direct three-fibre pattern cannot
have total space $S^8$ even after forgetting complex structures.
\end{proposition}

\begin{proof}
A fundamental cube of the Kuhn triangulation is divided into the six
tetrahedra indexed by permutations of its three coordinate directions.
The central toric fibre is stratified by algebraic tori; every
positive-dimensional torus stratum has Euler characteristic zero, while
the zero-dimensional strata are indexed by the maximal tetrahedra modulo
$\mathbb Z^3$.  There are exactly six, so additivity over the finite torus
stratification gives $\chi(W_3)=6$.

Let $\Sigma_f$ be the three critical values.  Euler integration, or a cut
into the regular bundle and three filling pieces, gives
\[
 \chi(\widetilde X)
 =\chi(S^2\setminus\Sigma_f)\chi(T^6)
  +\sum_{p\in\Sigma_f}\chi(\widetilde f^{-1}(p)).
\]
The regular term vanishes since $\chi(T^6)=0$.  A free finite quotient of a
torus has Euler characteristic zero: the covering formula multiplies its
Euler characteristic by the positive degree and the covering torus has
Euler characteristic zero.  Only $W_3$ remains, giving six.  Since
$\chi(S^8)=2$, the total spaces cannot be homeomorphic.
\end{proof}

\begin{corollary}[The exact singular-fibre repair equation]
\label[corollary]{h24:cor:euler-repair}
Let $g:M\to S^2$ be a proper triangulable map which is locally trivial
with fibre $T^6$ off a finite set $\Delta$, and suppose the direct Kuhn
cusp fibre $W_3$ occurs over $p_0\in\Delta$.  Then
\begin{equation}
 \chi(M)=6+\sum_{p\in\Delta\setminus\{p_0\}}
                    \chi\bigl(g^{-1}(p)\bigr).
 \label{h24:eq:euler-repair-general}
\end{equation}
Consequently an $S^8$ carrier requires, exactly at the Euler-characteristic
level,
\begin{equation}
 \sum_{p\in\Delta\setminus\{p_0\}}
                    \chi\bigl(g^{-1}(p)\bigr)=-4.
 \label{h24:eq:euler-repair-minus-four}
\end{equation}
If there are only the original two finite values, their replacement fibres
must have Euler characteristics summing to $-4$.  If those two fibres remain
free finite quotients of $T^6$, they still contribute zero, so the cusp
fibre itself must instead be replaced by a filling with Euler
characteristic $2$.  Without either change, topology is an output and every
closed carrier of the unchanged wrap necessarily has Euler characteristic
$6$.
\end{corollary}

\begin{proof}
Euler integration gives
$\chi(M)=\chi(S^2\setminus\Delta)\chi(T^6)
+\sum_{p\in\Delta}\chi(g^{-1}(p))$.  The first term is zero and
$\chi(W_3)=6$ by \cref{h24:prop:euler-mismatch}, giving
\eqref{h24:eq:euler-repair-general}.  Substitution of $\chi(S^8)=2$ gives
\eqref{h24:eq:euler-repair-minus-four}.  The remaining assertions are its two
specializations, using the covering calculation in the proof of
\cref{h24:prop:euler-mismatch}.
\end{proof}

\begin{remark}[What remains possible topologically]
A merely real singular map $S^8\to S^2$ with regular fibre $T^6$ is not
excluded by \cref{h24:thm:s8-obstruction}; it must have special-fibre Euler sum
$2$.  It cannot be a fibre bundle: the homotopy sequence would require the
connecting map $\pi_2(S^2)=\mathbb Z\to\pi_1(T^6)=\mathbb Z^6$ to be surjective.  A real
$S^9\to S^2$ model with regular fibre $T^7$ would likewise require genuine
singularities and special-fibre Euler sum $0$, but neither its carrier nor
its odd-dimensional regular fibre can be complex in the required sense.
\end{remark}

\subsection{The exact toroidal successor over the cusp}
\label{h24:sec:torus-successor}

We now retain, rather than discard, every exact map that survives the
obstructions.  Put
\[
 A_F=\mathbb Z\langle\widehat w,\widehat\delta\rangle,\qquad
 B_F=\mathbb Z\langle\widehat\gamma,\widehat u\rangle,\qquad
 B_0=\begin{pmatrix}0&1\\-1&0\end{pmatrix}:B_F\to A_F.
\]
In angular--deck order the candidate cusp monodromy and filling map are
\begin{equation}
 M_F(a,b)=(a+B_0b,b),\qquad
 c_F:A_F\oplus B_F\longrightarrow B_F,\quad c_F(a,b)=b.
 \label{h24:eq:candidate-cusp-module}
\end{equation}
Let $L_3=\mathbb Z^3$, let $k=(1,1,1)^{\mathsf T}$, and put
\begin{equation}
 P=\begin{pmatrix}1&-1&0\\0&1&-1\end{pmatrix},\qquad
 Q=B_0^{-1}P=\begin{pmatrix}0&-1&1\\1&-1&0\end{pmatrix}.
 \label{h24:eq:PQ}
\end{equation}
Both maps are onto, $\ker P=\ker Q=\mathbb Z k$, and $B_0Q=P$.
The rank-three cusp module and its specialization are
\begin{equation}
 M_3(a,b)=(a+b,b),\qquad
 c_3:L_3\oplus L_3\longrightarrow L_3,\quad c_3(a,b)=b.
 \label{h24:eq:rank3-cusp-module}
\end{equation}

\begin{theorem}[Primitive lattice, monodromy, and filling quotient]
\label[theorem]{h24:thm:lattice-quotient}
The homomorphism
\[
 \Phi=P\oplus Q:L_3\oplus L_3\longrightarrow A_F\oplus B_F,
 \qquad \Phi(a,b)=(Pa,Qb),
\]
fits into the exact sequence
\begin{equation}
 0\longrightarrow\mathbb Z^2
 \xrightarrow{\ \kappa\ }\mathbb Z^6
 \xrightarrow{\ \Phi\ }\mathbb Z^4\longrightarrow0,
 \qquad \kappa(m,n)=(mk,nk),
 \label{h24:eq:lattice-exact}
\end{equation}
and satisfies
\begin{equation}
 \Phi M_3=M_F\Phi,\qquad c_F\Phi=Qc_3.
 \label{h24:eq:mono-fill-compatibility}
\end{equation}
The kernel in \eqref{h24:eq:lattice-exact} is primitive.
\end{theorem}

\begin{proof}
The integer matrix
\[
 S=\begin{pmatrix}1&1\\0&1\\0&0\end{pmatrix}
\]
satisfies $PS=I_2$ and $QSB_0=I_2$.  Thus $P$ and $Q$, and hence $\Phi$,
are integrally surjective.  Their displayed kernels give
$\ker\Phi=\mathbb Z(k,0)\oplus\mathbb Z(0,k)$, proving exactness.  A kernel of a
surjection between free abelian groups is primitive.

Directly,
\[
 \Phi M_3(a,b)=(P(a+b),Qb),
\]
while
\[
 M_F\Phi(a,b)=(Pa+B_0Qb,Qb)=(Pa+Pb,Qb).
\]
The filling identity is equally literal:
$c_F\Phi(a,b)=Qb=Qc_3(a,b)$.
\end{proof}

Let $t=e^{2\pi i s}$ on a logarithmic cusp cover.  Retain the complete
bounded holomorphic correction $C_2(t)$ of the candidate and define
\begin{equation}
 Z_2=sB_0+C_2(t),\qquad
 C_3(t)=S C_2(t)Q,\qquad
 Z_3=sI_3+C_3(t).
 \label{h24:eq:period-extension}
\end{equation}
We use the angular--deck period convention
$\Omega_r=(I_r\;Z_r)$, and write $\Omega_r(t)$ when the dependence through
$s$ and $C_r(t)$ must be displayed.  The manuscript convention
$(Z_r\;I_r)$ is related by the displayed integral block permutation
$\left(\begin{smallmatrix}0&I_r\\I_r&0\end{smallmatrix}\right)$; no
period, sign, or lattice coordinate is silently changed.

\begin{theorem}[Holomorphic elliptic extension and its exact fibre]
\label[theorem]{h24:thm:elliptic-extension}
After shrinking the cusp disc, $\operatorname{Im}Z_3$ is invertible and
\begin{equation}
 PZ_3=Z_2Q,\qquad P\Omega_3=\Omega_2\Phi.
 \label{h24:eq:period-intertwining}
\end{equation}
Consequently $P:\mathbb C^3\to\mathbb C^2$ descends to a short exact sequence of complex
tori
\begin{equation}
 0\longrightarrow E_s\longrightarrow
 \frac{\mathbb C^3}{\mathbb Z^3+Z_3\mathbb Z^3}
 \xrightarrow{\ \overline P_s\ }
 \frac{\mathbb C^2}{\mathbb Z^2+Z_2\mathbb Z^2}
 \longrightarrow0,
 \qquad E_s\cong\frac{\mathbb C}{\mathbb Z+s\mathbb Z}.
 \label{h24:eq:torus-short-exact}
\end{equation}
The correction $C_2$ is recoverable from its prescribed lift by the exact
left inverse
\begin{equation}
 C_2=P C_3 S B_0.
 \label{h24:eq:correction-inverse}
\end{equation}
Thus the extension neither deletes the original periods nor replaces them
by an identity period matrix.
\end{theorem}

\begin{proof}
Since $C_3$ is bounded while $\operatorname{Im}s\to+\infty$,
$\operatorname{Im}Z_3=(\operatorname{Im}s)I_3+\operatorname{Im}C_3$ is
invertible on a sufficiently small cusp.  Using $PS=I_2$ and $B_0Q=P$,
\[
 PZ_3=sP+PSC_2Q=sB_0Q+C_2Q=Z_2Q,
\]
which is exactly \eqref{h24:eq:period-intertwining}.  It makes the complex
linear map $P$ a homomorphism of the displayed quotient tori.  It is
surjective on vector spaces and its lattice image is the whole target
lattice because $\Phi$ is onto; therefore its kernel is connected.

The vector-space kernel is $\mathbb C k$.  Since $Qk=0$, one has $C_3k=0$ and
$Z_3k=sk$.  If a source period lies in $\mathbb C k$, its image under $P$ is the
zero target period.  The real period map $\Omega_2:\mathbb R^4\to\mathbb C^2$ is an
isomorphism, so \eqref{h24:eq:period-intertwining} and
$\ker\Phi=\mathbb Z(k,0)\oplus\mathbb Z(0,k)$ show that all such periods are precisely
$\mathbb Z k+s\mathbb Z k$.  This proves the fibre formula.

Finally $QS=B_0^{-1}$, whence
\[
 PC_3SB_0=PSC_2QSB_0=C_2.
\]
\end{proof}

\subsubsection{The cusp filling map, not merely its boundary lattice}

For $m\in\mathbb Z^3$ and a permutation
$\pi=(\pi_1,\pi_2,\pi_3)\in S_3$, let
\begin{equation}
 \sigma_{m,\pi}=\operatorname{conv}
 \{m,m+e_{\pi_1},m+e_{\pi_1}+e_{\pi_2},m+e_1+e_2+e_3\}.
 \label{h24:eq:kuhn-tetrahedron}
\end{equation}
These are the tetrahedra of the periodic Kuhn triangulation; denote that
triangulation by $\mathcal T_3$.  Let $\mathcal T_2$ be the periodic
triangulation of $\mathbb R^2$ whose triangles are
\[
 \operatorname{conv}\{u,u+e_1,u+e_2\},\qquad
 \operatorname{conv}\{u+e_1,u+e_2,u+e_1+e_2\},
 \qquad u\in\mathbb Z^2.
\]
For $r=2,3$, let $\mathfrak S_r$ be the fan in $\mathbb Z^r\oplus\mathbb Z$ consisting
of all faces of the cones over $\tau\times\{1\}$ for
$\tau\in\mathcal T_r$, and put $Y_r=X_{\mathfrak S_r}$.  In the target
lattice write
\[
 n_1=\widehat w,\qquad n_2=\widehat\delta-\widehat w,\qquad
 n_3=-\widehat\delta;\qquad n_1+n_2+n_3=0.
\]
The map $P$ in \eqref{h24:eq:PQ} is exactly $P(e_i)=n_i$.

\begin{proposition}[Exact fan and quotient-filling morphism]
\label[proposition]{h24:prop:toric-filling-map}
Coning the Kuhn triangulation at height one gives a smooth toric fourfold
$Y_3$ with height map locally equal to
\[
 t=z_0z_1z_2z_3.
\]
The integral map
\[
 \widehat P:\mathbb Z^3\oplus\mathbb Z\longrightarrow A_F\oplus\mathbb Z,\qquad
 \widehat P(x,r)=(Px,r),
\]
is a map from this fan to the candidate's periodic $A_2$ fan and hence
induces a toric morphism $F:Y_3\to Y_2$ over the $t$-line.

On the dense torus define
\begin{align}
 \Psi^{(3)}_\lambda(x,t)
  &=\bigl(\exp(2\pi iC_3(t)\lambda)t^\lambda x,t\bigr),
       &&\lambda\in\mathbb Z^3,\label{h24:eq:source-toric-action}\\
 \Psi^{(2)}_\mu(y,t)
  &=\bigl(\exp(2\pi iC_2(t)\mu)t^{B_0\mu}y,t\bigr),
       &&\mu\in\mathbb Z^2.\label{h24:eq:target-toric-action}
\end{align}
Then
\begin{equation}
 F\Psi^{(3)}_\lambda=\Psi^{(2)}_{Q\lambda}F.
 \label{h24:eq:toric-action-intertwining}
\end{equation}
For a sufficiently small open cusp disc
$\Delta_\epsilon=\{t\in\mathbb C:|t|<\epsilon\}$, set
\begin{equation}
 Y_{r,\epsilon}=t^{-1}(\Delta_\epsilon),\qquad
 N_3^C=Y_{3,\epsilon}/\{\Psi^{(3)}_\lambda:\lambda\in\mathbb Z^3\},\qquad
 N_0=Y_{2,\epsilon}/\{\Psi^{(2)}_\mu:\mu\in\mathbb Z^2\}.
 \label{h24:eq:analytic-filling-quotients}
\end{equation}
After shrinking $\Delta_\epsilon$, both actions are free and properly
discontinuous, and $F$ descends to a holomorphic filling morphism
\begin{equation}
 \overline F:N_3^C\longrightarrow N_0
 \label{h24:eq:filling-morphism}
\end{equation}
whose map on smooth-fibre $H_1$ is $\Phi$ and whose monodromy compatibility
is \eqref{h24:eq:mono-fill-compatibility}.  Under
$H_1(T^6)=L_3\oplus L_3$ and $H_1(T^4)=A_F\oplus B_F$, the angular
generators bound the toric coordinate discs while the deck generators
survive.  Hence the inclusion-induced specialization maps are exactly
$c_3(a,b)=b$ and $c_F(a,b)=b$, and they form the commutative square
\begin{equation}
\begin{CD}
H_1(T^6;\mathbb Z)@>{\Phi}>>H_1(T^4;\mathbb Z)\\
@V{c_3}VV@VV{c_F}V\\
\mathbb Z^3@>{Q}>>\mathbb Z^2.
\end{CD}
\label{h24:eq:filling-square}
\end{equation}
\end{proposition}

\begin{proof}
The four vertices of \eqref{h24:eq:kuhn-tetrahedron} map to
\[
 Pm,\quad Pm+n_{\pi_1},\quad Pm-n_{\pi_3},\quad Pm.
\]
In target coordinates
$(n_1,n_2,n_3)=(e_1,e_2-e_1,-e_2)$.  For the six permutations the three
distinct images are the six $A_2$ triangles incident to $Pm$:
\[
\begin{array}{c|c@{\qquad}c|c}
(1,2,3)&\{v,v+e_1,v+e_2\}&(2,1,3)&\{v,v-e_1+e_2,v+e_2\}\\
(1,3,2)&\{v,v+e_1,v+e_1-e_2\}&(3,1,2)&\{v,v-e_2,v+e_1-e_2\}\\
(2,3,1)&\{v,v-e_1+e_2,v-e_1\}&(3,2,1)&\{v,v-e_2,v-e_1\}.
\end{array}
\]
Thus every source cone maps into a target cone.  The toric fan criterion
gives $F$, and preservation of the height coordinate gives $t\circ F=t$.
The maximal Kuhn cones are unimodular, giving the displayed local monomial
equation and smoothness of $Y_3$.

On the dense torus, $PC_3=C_2Q$ and $P=B_0Q$ give
\[
 F\Psi^{(3)}_\lambda(x,t)
 =\bigl(\exp(2\pi iC_2Q\lambda)t^{B_0Q\lambda}F(x),t\bigr),
\]
which proves \eqref{h24:eq:toric-action-intertwining}; the fan map extends it
over every toric orbit.

For the analytic quotient, put
$y=-\log|t|/(2\pi)$ and
\[
 \operatorname{trop}_t(x)=\frac{1}{\log|t|}
       (\log|x_1|,\log|x_2|,\log|x_3|).
\]
The exact transformation law is
\[
 \operatorname{trop}_t(\Psi^{(3)}_\lambda x)
 =\operatorname{trop}_t(x)+A(t)\lambda,\qquad
 A(t)=I_3+\frac{\operatorname{Im}C_3(t)}{y}.
\]
Boundedness of $C_3$ permits the cusp to be shrunk until
\[
\tfrac12\|\lambda\|\le\|A(t)\lambda\|
 \le\tfrac32\|\lambda\|.
\]
For the target use
\[
 \operatorname{trop}^{(2)}_t(y)
 =\frac{1}{\log|t|}(\log|y_1|,\log|y_2|).
\]
Its exact transformation law is
\[
 \operatorname{trop}^{(2)}_t(\Psi^{(2)}_\mu y)
 =\operatorname{trop}^{(2)}_t(y)+A_2(t)\mu,
 \qquad A_2(t)=B_0+\frac{\operatorname{Im}C_2(t)}{y}.
\]
Since $B_0$ is orthogonal, after one further shrinking the same bounds hold
with $A_2(t),\mu$ in place of $A(t),\lambda$.  Therefore a compact set can
meet only finitely many of its translates for either action.
On a boundary orbit the bounded cell label is translated by $\lambda$, so
the same conclusion holds there; the target argument uses the invertible
lattice map $B_0$ in exactly the same way.  A fixed point in a smooth fibre would
give $\operatorname{Im}Z_3\lambda=0$, hence $\lambda=0$; equality of a
boundary cell with its translate likewise forces $\lambda=0$.  The target
equalities give $\operatorname{Im}Z_2\mu=0$ in the smooth fibre and
$B_0\mu=0$ on the boundary, hence $\mu=0$.  Rounding
the coordinates of $A(t)^{-1}\operatorname{trop}_t(x)$ gives a uniformly
 bounded set of representatives.  The closures of these representatives lie
 in the finite union of closed toric polydiscs belonging to the six Kuhn
 tetrahedra; that finite union is compact.  This proves proper
discontinuity.  The same construction with $A_2(t)$ supplies target
representatives.  Over every smaller closed disc
$\overline\Delta_{\epsilon'}$, $0<\epsilon'<\epsilon$, the displayed
finite unions are compact, proving properness there.  The equivariant $F$
therefore descends, and
\eqref{h24:eq:filling-square} is exactly the second identity in
\eqref{h24:eq:mono-fill-compatibility}.
\end{proof}

\subsubsection{Finite peripheral data and the precise globalization gap}

In the manuscript lattice order the candidate matrices are
\begin{equation}
 A_1=\begin{pmatrix}1&0&0&0\\6&0&1&0\\-6&-1&-1&0\\-2&1&0&1\end{pmatrix},
 \qquad
 A_2=\begin{pmatrix}1&0&0&0\\0&0&-1&0\\-6&1&0&0\\3&0&1&1\end{pmatrix}.
 \label{h24:eq:finite-candidate-matrices}
\end{equation}
The matrices $A_1,A_2$ are written in the manuscript order
$(\widehat\gamma,\widehat u,\widehat w,\widehat\delta)=B_F\oplus A_F$,
whereas $M_F$ in \eqref{h24:eq:candidate-cusp-module} is written in the
angular--deck order $A_F\oplus B_F$.  With the explicit block swap
\begin{equation}
 J=\begin{pmatrix}0&I_2\\ I_2&0\end{pmatrix},\qquad
 M_F^{\rm man}=JM_FJ=
 \begin{pmatrix}I_2&0\\ B_0&I_2\end{pmatrix},
 \label{h24:eq:manuscript-cusp-order}
\end{equation}
one has
$A_1^3=A_2^4=I$ and $A_1A_2M_F^{\rm man}=I$.  Define
\begin{equation}
 C_1=\begin{pmatrix}-1&-1\\1&0\end{pmatrix},\quad
 C_2'=\begin{pmatrix}0&1\\-1&0\end{pmatrix},\quad
 C_0=\begin{pmatrix}1&1\\0&1\end{pmatrix}.
\label{h24:eq:elliptic-peripheral-matrices}
\end{equation}
For uniform notation below set $D_1=C_1$ and $D_2=C_2'$.
Then $(C_1)^3=(C_2')^4=I$, $C_1C_2'=C_0^{-1}$, and
$(C_0-I)^2=0$ with rank one.

The kernel block in this presentation is not an unrelated direct summand of
the rank-three cusp lattice.  Put
\begin{equation}
 \begin{gathered}
 r=\begin{pmatrix}0&0&1\end{pmatrix},\qquad
 A=\begin{pmatrix}P\\r\end{pmatrix}
 =\begin{pmatrix}1&-1&0\\0&1&-1\\0&0&1\end{pmatrix},\qquad
 B=\begin{pmatrix}Q\\r\end{pmatrix}
 =\begin{pmatrix}0&-1&1\\1&-1&0\\0&0&1\end{pmatrix},\\
 D=B_0\oplus[1]
 =\begin{pmatrix}0&1&0\\-1&0&0\\0&0&1\end{pmatrix},\qquad
 U_6=\operatorname{diag}(A,B)\in SL(6,\mathbb Z).
 \end{gathered}
 \label{h24:eq:rank-six-basis-change}
\end{equation}

\begin{lemma}[The identical cusp lattice in quotient--kernel coordinates]
\label[lemma]{h24:lem:cusp-basis-identification}
Let
\[
 \rho(Pa,ra;Qb,rb)=(Pa,Qb;ra,rb),\qquad
 \rho_{\rm man}=(J\oplus I_2)\rho.
\]
Then
\begin{align}
 U_6M_3U_6^{-1}
 &=\begin{pmatrix}I_3&D\\0&I_3\end{pmatrix},
 \label{h24:eq:rank-six-conjugacy}\\
 \rho U_6M_3U_6^{-1}\rho^{-1}
 &=M_F\oplus C_0,
 \label{h24:eq:angular-kernel-splitting}\\
 \rho_{\rm man}U_6M_3U_6^{-1}\rho_{\rm man}^{-1}
 &=M_F^{\rm man}\oplus C_0.
 \label{h24:eq:manuscript-kernel-splitting}
\end{align}
Projection onto the candidate block is exactly
$\operatorname{pr}_{1,2}\rho U_6(a,b)=(Pa,Qb)=\Phi(a,b)$ in
angular--deck coordinates and $J\Phi(a,b)$ in manuscript coordinates.
\end{lemma}

\begin{proof}
Both $A$ and $B$ have determinant one, and direct multiplication gives
$A=DB$, so $AB^{-1}=D$.  Conjugating
$M_3=\left(\begin{smallmatrix}I_3&I_3\\0&I_3\end{smallmatrix}\right)$
by $U_6$ proves \eqref{h24:eq:rank-six-conjugacy}.  The first two rows of
$A,B$ are exactly $P,Q$, while $r(1,1,1)^{\mathsf T}=1$; regrouping the
first two and last coordinates therefore gives
\eqref{h24:eq:angular-kernel-splitting}.  Applying the already displayed block
swap $J$ proves \eqref{h24:eq:manuscript-kernel-splitting} and both projection
formulas.
\end{proof}

\begin{proposition}[Exact finite-order affine lift]
\label[proposition]{h24:prop:finite-lift}
On $\widetilde\Lambda=\Lambda\oplus\mathbb Z^2$, identified with the original
$\mathbb Z^3\oplus\mathbb Z^3$ by \cref{h24:lem:cusp-basis-identification}, set
\[
 \widetilde A_1=A_1\oplus C_1,\qquad
 \widetilde A_2=A_2\oplus C_2',\qquad
 \widetilde M_0=M_F^{\rm man}\oplus C_0.
\]
Then
\[
 \widetilde A_1^3=I,\qquad \widetilde A_2^4=I,\qquad
 \widetilde A_1\widetilde A_2\widetilde M_0=I.
\]
The affine transformations
\begin{equation}
 \widetilde g_j(x,y)=(A_jx+v_j/m_j,D_jy),\qquad
 (m_1,m_2)=(3,4),
 \label{h24:eq:affine-lifts}
\end{equation}
of $T^4\times T^2$ have exact orders $m_j$, act freely, and project to
the candidate finite-fibre actions.
\end{proposition}

\begin{proof}
All three matrix identities follow blockwise from the identities just
displayed and the candidate identities.  The vector $(v_j,0)$ is fixed by
$\widetilde A_j$.  The $m_j$-th power of \eqref{h24:eq:affine-lifts} is
translation by the integral vector $(v_j,0)$ and hence is the identity on
the torus.  If a smaller positive power fixed a point, its projection to
$T^4$ would fix a point under the corresponding smaller power of the
candidate affine transformation.  The logarithmic-transform action in the
candidate is free, so this is impossible.  Projection to the first block is
equivariant by construction.
\end{proof}

\begin{remark}[Exact endpoint]
Theorems \ref{h24:thm:lattice-quotient}--\ref{h24:thm:elliptic-extension} and
Propositions \ref{h24:prop:toric-filling-map}--\ref{h24:prop:finite-lift} prove the
regular cusp, the entire toric cusp filling, and compatible finite
\emph{integral affine} data.  They do not construct one holomorphic
rank-three period family on the full $(3,4,\infty)$ cover whose two finite
analytic fillings glue to the cusp family.  Even if that analytic gluing is
constructed, \cref{h24:prop:euler-mismatch} proves that the unchanged
three-fibre pattern has Euler characteristic six and hence is not $S^8$.
This is a proved morphism with a proved endpoint, not an assertion that the
rank-three and rank-two objects are unrelated.
\end{remark}

\subsection[The forced S9, S10, and rank-loss S8]
{The forced $S^9$, $S^{10}$, and rank-loss $S^8$}
\label{h24:sec:sphere-ladder}

Put
\[
 \mathscr V=\operatorname{Mat}_{3\times3}(\mathbb C),\qquad
 \mathcal P_3=\{H\in\operatorname{Herm}_3(\mathbb C):H\succeq0\},\qquad
 \mathcal K_3=\{H\in\mathcal P_3:\operatorname{tr} H=1\}.
\]
The superscript $+$ denotes one-point compactification.

\begin{theorem}[Complete $U(3)$ and $SU(3)$ orbit invariants]
\label[theorem]{h24:thm:orbit-invariants}
The maps
\begin{align}
 \Gamma_U(A)&=A^*A,\label{h24:eq:U-gram}\\
 \Gamma_{SU}(A)&=(A^*A,\det A)\label{h24:eq:SU-gram}
\end{align}
induce homeomorphisms
\begin{align}
 \mathscr V/U(3)&\xrightarrow{\sim}\mathcal P_3,\label{h24:eq:U-quotient}\\
 \mathscr V/SU(3)&\xrightarrow{\sim}
 \mathcal Y_3:=\{(H,z)\in\mathcal P_3\times\mathbb C:|z|^2=\det H\}.
 \label{h24:eq:SU-quotient}
\end{align}
For a rank-$r$ matrix $A$, the two orbit fibres and stabilizers are
\[
 U(3)\cdot A\cong U(3)/U(3-r),\qquad
 SU(3)\cdot A\cong SU(3)/SU(3-r),
\]
with $SU(1)$ interpreted as the trivial group.
\end{theorem}

\begin{proof}
Both maps are invariant under the indicated left actions, and
$|\det A|^2=\det(A^*A)$.  Every $H\succeq0$ is the Gram matrix of
$H^{1/2}$.  If $A^*A=B^*B$, the rule $Ax\mapsto Bx$ is well-defined because
the two matrices have the same kernel, and is isometric because
\[
 \langle Bx,By\rangle=x^*B^*By=x^*A^*Ay=\langle Ax,Ay\rangle.
\]
Extending it orthogonally gives $U\in U(3)$ with $UA=B$.  Thus
\eqref{h24:eq:U-gram} separates $U(3)$-orbits.

Given $(H,z)\in\mathcal Y_3$, if $H$ is singular then $z=0$ and
$H^{1/2}$ realizes it.  If $H$ is positive definite, put
\[
 \omega=z/\sqrt{\det H},\qquad
 A=\operatorname{diag}(\omega,1,1)H^{1/2}.
\]
Then $(A^*A,\det A)=(H,z)$.  If two matrices have both invariants equal,
the unitary $U$ above has determinant one when they are invertible.  In
rank at most two, multiply $U$ on the left by a unitary that fixes
$\operatorname{im} B$ pointwise and corrects the determinant on its nonzero orthogonal
complement.  Hence \eqref{h24:eq:SU-gram} separates $SU(3)$-orbits.

Both invariant maps are proper because $\operatorname{tr}(A^*A)=\|A\|_F^2$.  The induced
continuous bijections to the displayed closed Hausdorff sets are therefore
homeomorphisms.  A left stabilizer fixes $\operatorname{im} A$ pointwise and is arbitrary
unitary, respectively special unitary, on $(\operatorname{im} A)^\perp$, proving the
fibre statements.
\end{proof}

Let
\[
 E=\{X\in\operatorname{Herm}_3(\mathbb C):\operatorname{tr} X=0\},\qquad \dim_{\mathbb R}E=8,\qquad c=I_3/3.
\]
Equip $E$ with the Frobenius inner product
$\langle X,Y\rangle_F=\operatorname{tr}(XY)$ and its norm; $B(E)$ and $S(E)$ below
denote its closed unit ball and unit sphere.
For $v\in S(E)$ define
\begin{equation}
 R_*(v)=-\frac{1}{3\lambda_{\min}(v)}.
 \label{h24:eq:radial-length}
\end{equation}
Since a nonzero traceless Hermitian matrix has a negative eigenvalue,
$R_*(v)$ is finite and positive.

\begin{lemma}[Explicit eight-ball chart]
\label[lemma]{h24:lem:K3-ball}
The maps
\begin{align}
 \chi(c)&=0,&
 \chi(c+s v)&=\frac{s}{R_*(v)}v,\quad0<s\le R_*(v),
 \label{h24:eq:chi}\\
 \chi^{-1}(rv)&=c+rR_*(v)v,\quad0<r\le1
 \label{h24:eq:chi-inverse}
\end{align}
are inverse homeomorphisms $\mathcal K_3\cong B(E)\cong B^8$.  They send
the rank-deficient locus onto $S(E)$.
\end{lemma}

\begin{proof}
The eigenvalues of $c+sv$ are $1/3+s\lambda_j(v)$.  Hence this matrix is
positive semidefinite exactly for $0\le s\le R_*(v)$ and is singular
exactly at the upper endpoint.  Every $H\ne c$ in $\mathcal K_3$ lies on
one unique such ray.  Continuity of the least eigenvalue makes
$v\mapsto-\lambda_{\min}(v)$ a positive continuous function on the compact
space $S(E)$, hence bounded above and bounded away from zero.  Thus
$R_*(v)$ is continuous and uniformly bounded above and away from zero.
The formulas \eqref{h24:eq:chi}--\eqref{h24:eq:chi-inverse} are consequently
continuous, including at $c$ and at the origin, and prove all assertions.
\end{proof}

\begin{theorem}[The forced sphere ladder and explicit inverses]
\label[theorem]{h24:thm:sphere-ladder}
Let $S_F^{17}=\{A\in\mathscr V:\|A\|_F=1\}$.  Then
\begin{equation}
 S_F^{17}/SU(3)\cong
 \mathcal Y_{3,1}:=\{(H,z)\in\mathcal Y_3:\operatorname{tr} H=1\}\cong S^9.
 \label{h24:eq:S9}
\end{equation}
An explicit second homeomorphism is
\begin{equation}
 \Omega(H,z)=
 \begin{cases}
 \left(\chi(H),\sqrt{1-\|\chi(H)\|^2}\,z/|z|\right),&H\succ0,\\
 (\chi(H),0),&H\text{ singular},
 \end{cases}
 \label{h24:eq:Omega-S9}
\end{equation}
from $\mathcal Y_{3,1}$ to $S(E\oplus\mathbb C)$.

Furthermore the weighted radial map
\begin{equation}
 \Phi_{SU}(H,z)=
 \begin{cases}
 0,&(H,z)=(0,0),\\
 s\,\Omega(H/s^2,z/s^3),&s=\sqrt{\operatorname{tr} H}>0
 \end{cases}
 \label{h24:eq:weighted-radial-map}
\end{equation}
is a homeomorphism $\mathcal Y_3\cong E\oplus\mathbb C\cong\mathbb R^{10}$.  It extends
to one-point compactifications and gives a homeomorphism of pairs
\begin{equation}
 \left(\mathscr V^+/SU(3),\{\det A=0\}^+/SU(3)\right)
 \cong\left(S^{10},S^8_{\mathrm{lin}}\right).
 \label{h24:eq:S10-S8-pair}
\end{equation}
The complement is
\begin{equation}
 S^{10}\setminus S^8\cong E\times(\mathbb C\setminus\{0\})
 \simeq S^1,
 \label{h24:eq:rank-complement}
\end{equation}
and $\arg\det A$ is the angular deformation-retraction coordinate.
\end{theorem}

\begin{proof}
The first identification in \eqref{h24:eq:S9} follows from
\cref{h24:thm:orbit-invariants}.  If $H\succ0$, then $z\ne0$; if $H$ is
singular, then $z=0$ and \cref{h24:lem:K3-ball} gives
$\|\chi(H)\|=1$.  Thus \eqref{h24:eq:Omega-S9} lands on the unit sphere and is
continuous even at rank loss because the coefficient of $z/|z|$ tends to
zero.  Given $(x,w)\in S(E\oplus\mathbb C)$, put $H=\chi^{-1}(x)$.  If
$\|x\|<1$, then $w\ne0$ and the unique inverse determinant is
\begin{equation}
 z=\sqrt{\det H}\,w/|w|.
 \label{h24:eq:Omega-inverse}
\end{equation}
If $\|x\|=1$, then $w=0$, $H$ is singular, and $z=0$.  This proves the
$S^9$ homeomorphism with its inverse.

For $(H,z)\ne0$, the normalized pair
$(H/s^2,z/s^3)$ belongs to $\mathcal Y_{3,1}$ because
$|z/s^3|^2=\det(H/s^2)$.  Conversely, for $v\ne0$ in $E\oplus\mathbb C$, set
\[
 s=\|v\|,\qquad (h,\zeta)=\Omega^{-1}(v/\|v\|),\qquad
 H=s^2h,\qquad z=s^3\zeta.
\]
This is the inverse of \eqref{h24:eq:weighted-radial-map}.  The norm identity
$\|\Phi_{SU}(H,z)\|=s$ proves continuity of both maps at zero and properness,
so the map extends to compactifications.

Rank loss is exactly $z=0$.  Under $\Omega$, the normalized rank-loss set
maps to $S(E\oplus0)$, and radialization maps the entire affine rank-loss
set to $E\oplus0$.  Its one-point compactification is the standard linear
$S^8\subset S^{10}$.  Removing it gives
$E\times(\mathbb C\setminus0)$; projection of the second factor to its phase is a
deformation retraction.  By construction this phase is
$z/|z|=\det A/|\det A|$.
\end{proof}

Put
\[
 \Sigma_3=\{H\in\operatorname{Herm}_3(\mathbb C):H\succeq0,\ \det H=0\}.
\]

\begin{proposition}[The $U(3)$ ball and why its boundary is the $S^8$]
\label[proposition]{h24:prop:U3-ball}
The map
\begin{equation}
 E\times\mathbb R_{\ge0}\longrightarrow\mathcal P_3,\qquad
 (K,t)\longmapsto K+(t-\lambda_{\min}(K))I_3
 \label{h24:eq:halfspace-map}
\end{equation}
is a homeomorphism with inverse
\[
 H\longmapsto
 \left(H-\frac{\operatorname{tr} H}{3}I_3,\lambda_{\min}(H)\right).
\]
Consequently
\begin{equation}
 \mathscr V^+/U(3)\cong\mathcal P_3^+\cong B^9,\qquad
 \partial B^9=\Sigma_3^+\cong S^8.
 \label{h24:eq:U3-ball-boundary}
\end{equation}
By contrast, the quotient of the source unit sphere is
\begin{equation}
 S_F^{17}/U(3)\cong\mathcal K_3\cong B^8,
 \label{h24:eq:U3-unit-ball}
\end{equation}
whose boundary is the trace-one rank-loss $S^7$.  Thus ``take the source
sphere quotient'' and ``take the boundary of the compactified full quotient''
are different operations.
\end{proposition}

\begin{proof}
The least eigenvalue of the image of \eqref{h24:eq:halfspace-map} is $t$.
Conversely, subtracting one third of the trace from $H$ produces a traceless
matrix $K$ whose least eigenvalue is
$\lambda_{\min}(H)-\operatorname{tr}(H)/3$; substitution recovers $H$.  Eigenvalue
continuity proves that the formulas are inverse homeomorphisms.  The
one-point compactification of a closed half-space is a closed ball, and its
boundary is the one-point compactification of the boundary hyperplane.
Now use \cref{h24:thm:orbit-invariants}.  Equation \eqref{h24:eq:U3-unit-ball}
follows from $\|A\|_F^2=\operatorname{tr}(A^*A)$ and \cref{h24:lem:K3-ball}.
\end{proof}

\begin{corollary}[Residual determinant circle and its fibres]
\label[corollary]{h24:cor:residual-circle}
The residual group $U(3)/SU(3)\cong S^1$ acts on the pair in
\eqref{h24:eq:S10-S8-pair} by
\[
 e^{i\theta}\cdot(x,w)=(x,e^{i\theta}w).
\]
Its fixed set is precisely the rank-loss $S^8$, and the orbit map has
\begin{equation}
 S^{10}/S^1\cong B^9,\qquad
 q^{-1}(q(p))\cong
 \begin{cases}S^1,&p\notin S^8,\\\{p\},&p\in S^8.\end{cases}
 \label{h24:eq:circle-fibres}
\end{equation}
On the normalized section, $S^9/S^1\cong B^8$ and the fixed set is the
rank-deficient $S^7=\partial B^8$.
\end{corollary}

\begin{proof}
Left multiplication by $U\in U(3)$ fixes $H=A^*A$ and sends
$z=\det A$ to $(\det U)z$.  Every phase is a determinant.  Equations
\eqref{h24:eq:Omega-S9} and \eqref{h24:eq:weighted-radial-map} therefore turn the
residual action into scalar rotation of the complex coordinate.  Its orbit
space on $E\oplus\mathbb C$ is $E\times\mathbb R_{\ge0}$, whose one-point compactification
is a closed $9$-ball; the zero complex coordinate is its boundary.  The
stabilizer is trivial for $w\ne0$ and all of $S^1$ for $w=0$, giving the
fibres.  On the unit sphere, replacing $w$ by $|w|$ gives the closed unit
ball in $E$ and the last statement.
\end{proof}

The $S^8$ in \eqref{h24:eq:S10-S8-pair} can also be defined without first
forming an orbit quotient.

\begin{proposition}[The exact residual-circle degeneration and join]
\label[proposition]{h24:prop:S9-join}
Under \eqref{h24:eq:S9}, forgetting determinant phase is exactly
\begin{equation}
 q_1:S(E\oplus\mathbb C)=S^9\longrightarrow B(E)=B^8,\qquad
 q_1(x,w)=x.
 \label{h24:eq:S9-to-B8}
\end{equation}
Its fibre is $S^1$ for $\|x\|<1$ and a singleton for
$\|x\|=1$.  More precisely,
\begin{equation}
 S^9=S(E\oplus\mathbb C)\cong S(E)*S(\mathbb C)\cong S^7*S^1,
 \label{h24:eq:S9-join}
\end{equation}
and $q_1$ collapses the $S^1$ join factor at the $S^7$ endpoint.  On the
full compactification the analogous map is
\begin{equation}
 q_+:S^{10}\longrightarrow B^9,
 \label{h24:eq:S10-to-B9}
\end{equation}
with circle fibre off the rank-loss $S^8$ and point fibre on that $S^8$.
The integral class of the circle on either complement is generated by
$\arg\det A$.
\end{proposition}

\begin{proof}
For fixed $H\succ0$ of trace one, the equation
$|z|^2=\det H>0$ leaves exactly its phase, an $S^1$; for singular $H$ it
forces $z=0$.  Formula \eqref{h24:eq:Omega-S9} converts $H$ into $x=\chi(H)$
and the determinant phase into the phase of $w$, so the orbit map is
literally \eqref{h24:eq:S9-to-B8}.  Every unit pair has the unique form
\[
 (x,w)=(\cos\alpha\,u,\sin\alpha\,v),\qquad
 u\in S(E),\ v\in S(\mathbb C),\ 0\le\alpha\le\pi/2,
\]
away from the endpoints, with the usual endpoint identifications.  This is
the defining quotient of the join and proves \eqref{h24:eq:S9-join}.  The full
statement is \cref{h24:cor:residual-circle,h24:prop:U3-ball}.  The generator claim
follows from the loop
\[
 A_\theta=3^{-1/2}\operatorname{diag}(e^{i\theta},1,1),\qquad
 0\le\theta\le2\pi,
\]
whose Gram matrix is $I_3/3$ and whose determinant phase winds once.
\end{proof}

\begin{proposition}[Spectral-floor chart and the marked fibres]
\label[proposition]{h24:prop:spectral-floor-observation}
The maps
\begin{equation}
 \mathcal F:E\longrightarrow\Sigma_3,\qquad
 \mathcal F(K)=K-\lambda_{\min}(K)I_3,\qquad
 \mathcal F^{-1}(H)=H-\frac{\operatorname{tr} H}{3}I_3
 \label{h24:eq:spectral-floor}
\end{equation}
are inverse homeomorphisms.  Hence
$\Sigma_3^+\cong S^8$, and it is exactly the $S^8$ of
\eqref{h24:eq:S10-S8-pair}.

The fixed marked compression
\begin{equation}
 \mathscr O:\Sigma_3\longrightarrow\operatorname{Herm}_2(\mathbb C)_{\succeq0},\qquad
 \mathscr O(H)=PHP^*
 \label{h24:eq:marked-compression}
\end{equation}
is onto.  For a fixed $G\succeq0$, put
\[
 X_P^{-1}HX_P^{-*}=\begin{pmatrix}G&a\\a^*&\lambda\end{pmatrix}.
\]
Its complete fibre in $\Sigma_3$ is
\begin{equation}
 \mathscr O^{-1}(G)=
 \left\{
 X_P\begin{pmatrix}G&a\\a^*&a^*G^\dagger a+\tau\end{pmatrix}X_P^*:
 a\in\operatorname{Ran}G,\ \tau\ge0,\
 \bigl(\operatorname{rank} G<2\ \text{or}\ \tau=0\bigr)
 \right\}.
 \label{h24:eq:marked-complete-fibre}
\end{equation}
Consequently the fibres are homeomorphic to
\begin{equation}
 \begin{cases}
 \mathbb C^2,&G\succ0,\\
 \mathbb C\times\mathbb R_{\ge0},&\operatorname{rank} G=1,\\
 \mathbb R_{\ge0},&G=0.
 \end{cases}
 \label{h24:eq:marked-fibre-types}
\end{equation}
The positive section is $G\mapsto RGR^*$.  The affine map
\eqref{h24:eq:marked-compression} is not proper, so these formulas do not assert
an unproved extension through the one-point compactification point.
\end{proposition}

\begin{proof}
For $K\in E$, the matrix $\mathcal F(K)$ is positive semidefinite and has
least eigenvalue zero.  If $H\in\Sigma_3$, then
$K=H-(\operatorname{tr} H/3)I_3$ is traceless and has least eigenvalue $-\operatorname{tr} H/3$;
therefore the two formulas in \eqref{h24:eq:spectral-floor} are inverse.
Continuity of the least eigenvalue proves the homeomorphism.  Its one-point
compactification is $E^+\cong S^8$.  The proof of
\cref{h24:thm:sphere-ladder} identifies the same rank-loss affine space with
$E\oplus0$, so the two $S^8$ descriptions agree.

Equations \eqref{h24:eq:marked-P}--\eqref{h24:eq:marked-inverse} show that the
upper-left block of $X_P^{-1}HX_P^{-*}$ is $PHP^*=G$.  Congruence by the
invertible $X_P$ reduces positivity to positivity of the displayed block
matrix.  The exact generalized Schur-complement criterion is
\[
 G\succeq0,\qquad a\in\operatorname{Ran}G,\qquad
 \tau:=\lambda-a^*G^\dagger a\ge0.
\]
Indeed, multiplying on both sides by the invertible triangular matrices
with off-diagonal entry $-G^\dagger a$ reduces it to
$\operatorname{diag}(G,\tau)$; if $a$ had a component on $\ker G$, testing the quadratic
form on a large signed multiple of a kernel vector would violate
positivity.  The same congruence gives
\[
 \operatorname{rank} H=\operatorname{rank} G+\mathbf 1_{\tau>0}.
\]
Thus $H$ is singular precisely when $\operatorname{rank} G<2$ or $\tau=0$, proving
\eqref{h24:eq:marked-complete-fibre}.  The three fibre types follow by using
$\dim_\mathbb C\operatorname{Ran}G=2,1,0$.  Setting $a=0$ and $\tau=0$ gives
$H=RGR^*$ and proves surjectivity and the section.  Finally, for fixed
$G\succ0$ the parameter $a\in\mathbb C^2$ is unbounded, so the inverse image of
the compact singleton $\{G\}$ is not compact; $\mathscr O$ is not proper.
\end{proof}

\subsubsection[A global marked observation on the S9 quotient]
{A global marked observation on the $S^9$ quotient}

Although compression on the affine rank-loss $S^8$ is nonproper, the
trace-one $S^9$ has a global compact observation with exactly computable
fibres.  Define
\begin{equation}
 \mathscr O_9:\mathcal Y_{3,1}\longrightarrow
 \mathcal B_P:=\{G\succeq0:\operatorname{tr}(C^{-1}G)\le1\},\qquad
 \mathscr O_9(H,z)=PHP^*.
 \label{h24:eq:O9}
\end{equation}
The set $\mathcal B_P$ is a compact convex body with nonempty interior in
$\operatorname{Herm}_2(\mathbb C)$, hence is homeomorphic to $B^4$ by radial projection from an
interior point.

\begin{theorem}[Every fibre of the $S^9$ marked observation]
\label[theorem]{h24:thm:O9-fibres}
The map \eqref{h24:eq:O9} is a continuous surjection.  For
$G\in\mathcal B_P$ set
\begin{equation}
 L(G)=\frac{1-\operatorname{tr}(C^{-1}G)}{3}.
 \label{h24:eq:L-G}
\end{equation}
Its fibre is exactly the set of pairs
\begin{equation}
 a\in\operatorname{Ran}G,\qquad
 a^*G^\dagger a\le L(G),\qquad
 z\in\mathbb C,\qquad
 |z|^2=\det G\bigl(L(G)-a^*G^\dagger a\bigr),
 \label{h24:eq:O9-complete-fibre}
\end{equation}
where $(H,z)$ is recovered uniquely from
\begin{equation}
 H=X_P\begin{pmatrix}G&a\\a^*&L(G)\end{pmatrix}X_P^*.
 \label{h24:eq:O9-H-recovery}
\end{equation}
The fibre is $S^5$ when $G\succ0$ and $L(G)>0$; it is a point when
$G\succ0$ and $L(G)=0$; it is a closed two-disc when $\operatorname{rank} G=1$ and
$L(G)>0$, degenerating to a point when $L(G)=0$; and it is a point at
$G=0$.  These are orbit-space fibres, not torus fibres.
\end{theorem}

\begin{proof}
Since $\operatorname{tr} H=\operatorname{tr}(X_P^*X_PK)$ for
$K=X_P^{-1}HX_P^{-*}$, \eqref{h24:eq:marked-inverse} gives
\[
 1=\operatorname{tr} H=\operatorname{tr}(C^{-1}G)+3\lambda.
\]
Thus $\lambda=L(G)$.  The Schur-complement argument in
\cref{h24:prop:spectral-floor-observation} gives exactly the first two
conditions in \eqref{h24:eq:O9-complete-fibre}.  Since $\det X_P=1$,
\[
 \det H=\det G\bigl(L(G)-a^*G^\dagger a\bigr)
\]
when $G$ is invertible, and the same identity follows polynomially, or from
rank, when $G$ is singular.  The defining equation
$|z|^2=\det H$ of $\mathcal Y_{3,1}$ gives the last condition.  Conversely
these conditions construct a unique point by \eqref{h24:eq:O9-H-recovery};
taking $a=0$ proves surjectivity for every $G\in\mathcal B_P$.

If $G\succ0$, put $b=G^{-1/2}a$ and
$w=z/\sqrt{\det G}$.  Equation \eqref{h24:eq:O9-complete-fibre} becomes
$\|b\|^2+|w|^2=L(G)$ in $\mathbb C^2\oplus\mathbb C$.  It is $S^5$ for $L(G)>0$ and a
point for $L(G)=0$.  If $\operatorname{rank} G=1$, then $z=0$ and the allowed $a$ form
the closed Hermitian norm disc in the one-dimensional complex space
$\operatorname{Ran}G$.  If $G=0$, then $a=z=0$ and
$H=uu^*/3$ is the unique point.
\end{proof}

\subsubsection[The exact period-to-S9 and period-to-S8 maps]
{The exact period-to-$S^9$ and period-to-$S^8$ maps}

For $t$ in the punctured cusp and
$G_6\in\operatorname{Sym}^+_6(\mathbb R)$, retain both arguments and put
\begin{equation}
 G_4=\Phi G_6\Phi^{\mathsf T},\qquad
 H_3(t;G_6)=\Omega_3(t)G_6\Omega_3(t)^*,\qquad
 H_2(t;G_6)=\Omega_2(t)G_4\Omega_2(t)^*.
 \label{h24:eq:period-Grams}
\end{equation}

\begin{proposition}[Commuting marked Gram square]
\label[proposition]{h24:prop:Gram-square}
The matrices $G_4,H_3(t;G_6)$, and $H_2(t;G_6)$ are positive definite, and
\begin{equation}
 PH_3(t;G_6)P^*=H_2(t;G_6).
 \label{h24:eq:Gram-observation}
\end{equation}
Equivalently, for every fixed $t$ the square
\begin{equation}
\begin{CD}
\operatorname{Sym}^+_6(\mathbb R)@>{G\mapsto\Omega_3(t)G\Omega_3(t)^*}>>
\operatorname{Herm}_3(\mathbb C)_{\succ0}\\
@V{G\mapsto\Phi G\Phi^{\mathsf T}}VV
@VV{H\mapsto PHP^*}V\\
\operatorname{Sym}^+_4(\mathbb R)@>{G\mapsto\Omega_2(t)G\Omega_2(t)^*}>>
\operatorname{Herm}_2(\mathbb C)_{\succ0}
\end{CD}
\label{h24:eq:Gram-square}
\end{equation}
commutes.  Normalize without suppressing either argument:
\[
 h_3(t;G_6)=\frac{H_3(t;G_6)}{\operatorname{tr} H_3(t;G_6)},\qquad
 z_3(t;G_6)=+\sqrt{\det h_3(t;G_6)}.
\]
Then the explicit point
\begin{equation}
 \jmath(t;G_6)=\bigl(h_3(t;G_6),z_3(t;G_6)\bigr)
 \in\mathcal Y_{3,1}\cong S^9,
 \qquad
 \mathscr O_9\bigl(\jmath(t;G_6)\bigr)
 =\frac{H_2(t;G_6)}{\operatorname{tr} H_3(t;G_6)}.
 \label{h24:eq:period-to-S9}
\end{equation}
The positive square root is a specified phase-zero section; it does not
erase the existence of the other determinant phases in $S^9$.

The spectral-floor map gives a second exact, generally noncompact-chart map
\begin{equation}
 \begin{aligned}
 \sigma(t;G_6)
 &=H_3(t;G_6)-\lambda_{\min}\bigl(H_3(t;G_6)\bigr)I_3
   \in\Sigma_3\hookrightarrow\Sigma_3^+\cong S^8,\\
 \mathscr O\bigl(\sigma(t;G_6)\bigr)
 &=H_2(t;G_6)-\lambda_{\min}\bigl(H_3(t;G_6)\bigr)C.
 \end{aligned}
 \label{h24:eq:period-to-S8}
\end{equation}
\end{proposition}

\begin{proof}
Surjectivity of $\Phi$ makes $\Phi^{\mathsf T}$ injective; hence
$G_4=\Phi G_6\Phi^{\mathsf T}$ is positive definite.  The identity blocks
in $\Omega_r(t)$ make $\Omega_r(t)^*$ injective, so both $H_r(t;G_6)$ are
positive definite.  Using $P\Omega_3(t)=\Omega_2(t)\Phi$ from
\eqref{h24:eq:period-intertwining},
\[
 PH_3(t;G_6)P^*=P\Omega_3(t)G_6\Omega_3(t)^*P^*
 =\Omega_2(t)\Phi G_6\Phi^{\mathsf T}\Omega_2(t)^*=H_2(t;G_6).
\]
This proves the square and \eqref{h24:eq:period-to-S9}.  Subtracting the least
eigenvalue makes a positive matrix singular without altering its
eigenspaces; applying $P$ and $PI_3P^*=C$ proves
\eqref{h24:eq:period-to-S8}.
\end{proof}

There is one useful cusp endpoint for a fixed, explicitly chosen metric.
Take $G_6=I_6$ in the real period frame, choose the standard logarithmic
strip $0\le x<1$ for $s=x+iy$, and abbreviate
$H_{3,I}(t)=H_3(t;I_6)$.  The exact formula is
\begin{equation}
 \begin{split}
 H_{3,I}(t)
 &=I_3+Z_3(t)Z_3(t)^*\\
 &=(1+|s|^2)I_3+sC_3(t)^*+\overline{s}C_3(t)
   +C_3(t)C_3(t)^*.
 \end{split}
 \label{h24:eq:identity-period-metric}
\end{equation}
Put
$h_{3,I}(t)=H_{3,I}(t)/\operatorname{tr} H_{3,I}(t)$,
$z_{3,I}(t)=+\sqrt{\det h_{3,I}(t)}$, and
$\jmath_I(t)=\jmath(t;I_6)$.

\begin{corollary}[A proved cusp limit in the $S^9$ carrier]
\label[corollary]{h24:cor:S9-cusp-limit}
For \eqref{h24:eq:identity-period-metric}, boundedness of $C_3(t)$ gives
\begin{equation}
 h_{3,I}(t)\longrightarrow I_3/3,\qquad
 z_{3,I}(t)\longrightarrow1/(3\sqrt3),\qquad
 \mathscr O_9\bigl(\jmath_I(t)\bigr)\longrightarrow C/3
 \quad(t\to0).
 \label{h24:eq:S9-limit}
\end{equation}
Thus the period metric has an exact continuous marked limit in the interior
of the $S^9$ orbit carrier.  This is a map of period--metric data, not a
compactification of every point of the torus family and not a fourth
singular filling.
\end{corollary}

\begin{proof}
Equation \eqref{h24:eq:identity-period-metric} and boundedness of $C_3$ give
$y^{-2}H_{3,I}\to I_3$ as $y\to\infty$: indeed $0\le x<1$ gives
$|s|^2/y^2\to1$, while the two terms linear in $s$ are $O(y)$ and the
remaining correction is $O(1)$.  Taking traces and determinants yields
$h_{3,I}\to I_3/3$ and the positive square-root limit for $z_{3,I}$.
Compressing by $P$ and using \eqref{h24:eq:period-to-S9} gives the last limit.
\end{proof}

\begin{theorem}[Which Gram carrier is strongest]
\label[theorem]{h24:thm:strongest-gram}
Within the specified left-action invariant theory, the information order is
strict:
\[
 \mathcal Y_3\xrightarrow{\text{forget }\arg z}\mathcal P_3,
 \qquad
 \mathcal Y_{3,1}\xrightarrow{\text{forget }\arg z}\mathcal K_3.
\]
The maps have circle fibres on full-rank matrices and point fibres at rank
loss.  Hence:
\begin{enumerate}[label=\textup{(\roman*)}]
\item $S^9=S_F^{17}/SU(3)$ is the strongest \emph{normalized single-sphere}
carrier: it retains the complete Gram matrix, determinant phase, its
integral winding, the exact circle degeneration \eqref{h24:eq:S9-to-B8}, and
the global marked observation \eqref{h24:eq:O9}.
\item $(S^{10},S^8)$ is the strongest \emph{unnormalized compact pair}: it
additionally retains the radial scale, and its $S^8$ is exactly the fixed
rank-loss locus.
\item The standalone $S^8=\partial(\mathscr V^+/U(3))$ is the strongest
rank-loss sphere, but it has forgotten determinant phase and does not carry
the global compact extension of the affine compression
\eqref{h24:eq:marked-compression}.
\end{enumerate}
None of this ordering supplies a complex tangent structure, a $T^6$ regular
fibre, the base $\mathbb P^1$, the three special values, or their integral
specialization maps.
\end{theorem}

\begin{proof}
The orbit-separation theorem shows that $(H,z)$ is the complete $SU(3)$
invariant and $H$ the complete $U(3)$ invariant.  For $H\succ0$, changing
the phase of $z$ changes the $SU(3)$ orbit but not the $U(3)$ orbit; at rank
loss $z=0$, so no phase remains.  This proves strict information loss and
the fibre claim.  Propositions \ref{h24:prop:U3-ball} and \ref{h24:prop:S9-join},
Theorem \ref{h24:thm:O9-fibres}, and Equation \eqref{h24:eq:S10-S8-pair} prove
(i)--(iii).  The absent structures cannot be inferred from the quotient
map: its base has dimensions eight or nine, its generic fibre is $S^1$, its
singular values form the full boundary $S^7$ or $S^8$, and its complement
has fundamental group $\mathbb Z$.  In contrast, since
$f|_{X^\circ}:X^\circ\to B^\circ$ is a $T^4$-bundle and
$\pi_2(B^\circ)=0$, its homotopy exact sequence gives
\[
 1\longrightarrow\mathbb Z^4\longrightarrow\pi_1(X^\circ)
 \longrightarrow\pi_1(B^\circ)=F_2\longrightarrow1,
\]
which surjects onto the nonabelian free group $F_2$.  Thus even the regular
complements cannot be homotopy equivalent.  The tangent obstructions in
\cref{h24:thm:s8-obstruction,h24:thm:all-higher-spheres} finish the claimed scope.
\end{proof}

\subsection{The 24-dimensional octonionic flag}
\label{h24:sec:octonionic-flag}

We now construct the exceptional Jordan object instead of selecting it by
dimension.  Let
\[
 \mathfrak h_3(\mathbb O)=
 \left\{
 \begin{pmatrix}
 x_1&p&q\\ \overline p&x_2&r\\ \overline q&\overline r&x_3
 \end{pmatrix}:
 x_1,x_2,x_3\in\mathbb R,\ p,q,r\in\mathbb O
 \right\},
 \qquad
 a\circ b=\frac12(ab+ba).
\]
Define
\[
 \mathbb O\mathbb P^2=\{a\in\mathfrak h_3(\mathbb O):a^2=a,\ \operatorname{tr} a=1\}
\]
and, without replacing the incidence equation by a mnemonic,
\begin{equation}
 \operatorname{Fl}(\mathbb O)=
 \{(a,b)\in\mathbb O\mathbb P^2\times\mathbb O\mathbb P^2:
                 \operatorname{Re}\operatorname{tr}(ab)=0\}.
 \label{h24:eq:octonionic-flag-definition}
\end{equation}
Let
\[
 F_4=\operatorname{Aut}_{\mathbb R}\bigl(\mathfrak h_3(\mathbb O),\circ\bigr),\qquad
 d_1=\operatorname{diag}(1,0,0),\qquad d_2=\operatorname{diag}(0,1,0),\qquad
 x_{\mathbb O,0}=(d_1,d_2).
\]
Here \(F_4\) denotes the compact real form.  These definitions and all
claims in the next theorem are recorded in the primary calculation of
 Mare and Willems \cite[Definition~2.2, p.~488;
 Proposition~2.4, p.~489, and Proposition~2.7, p.~491, and the
 root-space discussion on pp.~491--492]{h24:MareWillems2013}.

\begin{theorem}[The exact triality tangent module]
\label[theorem]{h24:thm:octonionic-triality}
The diagonal \(F_4\)-action on \(\operatorname{Fl}(\mathbb O)\) is transitive,
the stabilizer of \(x_{\mathbb O,0}\) is \(\operatorname{Spin}(8)\), and hence
\[
 \operatorname{Fl}(\mathbb O)\cong F_4/\operatorname{Spin}(8).
\]
Fix pairwise distinct real numbers \(\lambda_1,\lambda_2,\lambda_3\)
with \(\lambda_1+\lambda_2+\lambda_3=0\), and set
\(d_3=I_3-d_1-d_2\) and \(d_\lambda=\operatorname{diag}(\lambda_1,\lambda_2,\lambda_3)\).
The exact comparison with the diagonal-orbit model is the equivariant
diffeomorphism
\begin{equation}
 \begin{aligned}
 \Xi_\lambda:\operatorname{Fl}(\mathbb O)&\longrightarrow
                 F_4d_\lambda\subset\mathfrak h_3^0(\mathbb O),\\
 (a,b)&\longmapsto
       \lambda_1a+\lambda_2b+\lambda_3(I_3-a-b).
 \end{aligned}
 \label{h24:eq:flag-diagonal-orbit-map}
\end{equation}
For \(z\in F_4d_\lambda\), its inverse is \((p_1(z),p_2(z))\), where
\begin{equation}
 p_i(z)=\frac{(z-\lambda_j I_3)\circ(z-\lambda_k I_3)}
                  {(\lambda_i-\lambda_j)(\lambda_i-\lambda_k)},
 \qquad\{i,j,k\}=\{1,2,3\}.
 \label{h24:eq:flag-spectral-idempotents}
\end{equation}
Define the three coordinate subspaces of \(\mathfrak h_3^0(\mathbb O)\) by
\begin{align}
 \mathfrak h_{\gamma_1}
 &=\left\{
 \begin{pmatrix}0&0&0\\0&0&r\\0&\overline r&0\end{pmatrix}:r\in\mathbb O
 \right\},\nonumber\\
 \mathfrak h_{\gamma_2}
 &=\left\{
 \begin{pmatrix}0&0&q\\0&0&0\\\overline q&0&0\end{pmatrix}:q\in\mathbb O
 \right\},\label{h24:eq:octonionic-root-spaces}\\
 \mathfrak h_{\gamma_3}
 &=\left\{
 \begin{pmatrix}0&p&0\\\overline p&0&0\\0&0&0\end{pmatrix}:p\in\mathbb O
 \right\}.\nonumber
\end{align}
As real \(\operatorname{Spin}(8)\)-modules,
\begin{equation}
 T_{x_{\mathbb O,0}}\operatorname{Fl}(\mathbb O)
 \xrightarrow[\cong]{\ d\Xi_\lambda\ }
 T_{d_\lambda}(F_4d_\lambda)
 =\mathfrak h_{\gamma_1}\oplus
  \mathfrak h_{\gamma_2}\oplus
  \mathfrak h_{\gamma_3}
 \cong V_8\oplus S_8^+\oplus S_8^-.
 \label{h24:eq:triality-tangent}
\end{equation}
In particular,
\[
 \dim_{\mathbb R}\operatorname{Fl}(\mathbb O)=8+8+8=24.
\]
\end{theorem}

\begin{proof}
Mare and Willems prove that the stabilizer of \((d_1,d_2)\) consists
exactly of the automorphisms fixing the diagonal Jordan subalgebra
pointwise, which is \(\operatorname{Spin}(8)\), and that the action is
transitive.  Their restricted-root decomposition of
\(\mathfrak h_3^0(\mathbb O)\) has the three spaces in
\eqref{h24:eq:octonionic-root-spaces}; the tangent space to the orbit is their
direct sum.  To verify the model comparison, write
\((a,b)=(h d_1,h d_2)\) with \(h\in F_4\).  Since \(hI_3=I_3\),
\(I_3-a-b=hd_3\), so \(\Xi_\lambda(a,b)=hd_\lambda\).
The three orthogonal diagonal idempotents satisfy the polynomial
identities \(p_i(d_\lambda)=d_i\).  Jordan automorphisms preserve these
identities, whence \(p_i(hd_\lambda)=hd_i\).  This proves both inverse
compositions in \eqref{h24:eq:flag-diagonal-orbit-map}--
\eqref{h24:eq:flag-spectral-idempotents}.  All formulas are polynomial on
the ambient real vector spaces with nonzero constant denominators, so the
maps are smooth.  Their differentials are inverse, which gives the
displayed tangent identification at the correctly typed basepoints.
Under the stabilizer, the three summands are respectively
the vector representation and the two real half-spin representations
\cite[pp.~491--492, especially the discussion around
Equation~(2.10)]{h24:MareWillems2013}.  Each is the displayed copy of the
eight-dimensional real vector space \(\mathbb O\), which proves the dimension
calculation without suppressing any summand.
\end{proof}

For \(k=1,2,3\), let
\[
 \mathcal E_k=F_4\times_{\operatorname{Spin}(8)}
                         \mathfrak h_{\gamma_k}.
\]
These are rank-eight real subbundles, and
\begin{equation}
 T\operatorname{Fl}(\mathbb O)=
 \mathcal E_1\oplus\mathcal E_2\oplus\mathcal E_3.
 \label{h24:eq:octonionic-tangent-bundles}
\end{equation}
They are associated bundles; the notation does not assert that any
\(\mathcal E_k\) is trivial.

\begin{theorem}[Integral ring, cells, and Euler class]
\label[theorem]{h24:thm:octonionic-cohomology}
There are orientations and a basis
\(\beta_1,\beta_2\in H^8(\operatorname{Fl}(\mathbb O);\mathbb Z)\) for which
\begin{equation}
 \begin{aligned}
 e_1:=e(\mathcal E_1)&=2\beta_1-\beta_2,\\
 e_2:=e(\mathcal E_2)&=-\beta_1+2\beta_2,\\
 e_3:=e(\mathcal E_3)&=e_1+e_2=\beta_1+\beta_2,
 \end{aligned}
 \label{h24:eq:octonionic-Euler-generators}
\end{equation}
and
\begin{equation}
 H^*(\operatorname{Fl}(\mathbb O);\mathbb Z)
 \cong
 \frac{\mathbb Z[\beta_1,\beta_2]}
 {\bigl(\beta_1^2-\beta_1\beta_2+\beta_2^2,\,
         \beta_1\beta_2(\beta_1-\beta_2)\bigr)},
 \qquad |\beta_1|=|\beta_2|=8.
 \label{h24:eq:octonionic-cohomology-ring}
\end{equation}
Writing
\[
 u_{\mathbb O}=\beta_1^2\beta_2=\beta_1\beta_2^2
 \in H^{24}(\operatorname{Fl}(\mathbb O);\mathbb Z)
\]
with the induced orientation, \(u_{\mathbb O}\) is a primitive top generator and
\begin{equation}
 e\bigl(T\operatorname{Fl}(\mathbb O)\bigr)
 =e_1e_2e_3=6u_{\mathbb O},\qquad
 \left\langle e(T\operatorname{Fl}(\mathbb O)),
          [\operatorname{Fl}(\mathbb O)]\right\rangle=6.
 \label{h24:eq:octonionic-Euler-six}
\end{equation}
The integral cohomology is free and has Poincar\'e polynomial
\begin{equation}
 1+2t^8+2t^{16}+t^{24}.
 \label{h24:eq:octonionic-Poincare}
\end{equation}

More precisely, on
\(\mathfrak d^0=\{\operatorname{diag}(x_1,x_2,x_3):x_1+x_2+x_3=0\}\), put
\[
 \gamma_1(x)=x_3-x_2,\qquad
 \gamma_2(x)=x_1-x_3,\qquad
 \gamma_3(x)=x_1-x_2=\gamma_1(x)+\gamma_2(x).
\]
In the symmetric-pair normalization the actual restricted roots are
\(\{\pm\gamma_1/2,\pm\gamma_2/2,\pm\gamma_3/2\}\), of type \(A_2\).
The doubled labels \(\gamma_i\) name the same root spaces and reflection
hyperplanes; this factor of two is the one stated in the source's
root-normalization footnote \cite[footnote~3, p.~491]{h24:MareWillems2013}.
There is a cell decomposition
\begin{equation}
 \operatorname{Fl}(\mathbb O)=\bigsqcup_{\pi\in S_3}C_\pi,\qquad
 C_\pi\cong\mathbb R^{8\ell(\pi)}.
 \label{h24:eq:octonionic-six-cells}
\end{equation}
For the stabilizer action,
\begin{equation}
 \operatorname{Fl}(\mathbb O)^{\operatorname{Spin}(8)}
 =\{\pi x_{\mathbb O,0}:\pi\in S_3\}.
 \label{h24:eq:octonionic-six-fixed-points}
\end{equation}
Thus both the cell calculation and the Euler-class calculation give
\(\chi(\operatorname{Fl}(\mathbb O))=6\).
\end{theorem}

\begin{proof}
Mare and Willems' integral theorem gives a free ring generated by
\(\beta_1,\beta_2\), with relations
\[
 S_i(\beta_1,\beta_2-\beta_1,-\beta_2)=0,\qquad i=2,3,
\]
and gives \eqref{h24:eq:octonionic-Euler-generators}
\cite[Theorem~1.1, p.~484, and Equation~(3.2),
p.~496]{h24:MareWillems2013}.  Expanding without changing the generators,
\[
 S_2=-\beta_1^2+\beta_1\beta_2-\beta_2^2,\qquad
 S_3=\beta_1\beta_2(\beta_1-\beta_2),
\]
which is exactly \eqref{h24:eq:octonionic-cohomology-ring}.  The second
relation gives
\(\beta_1^2\beta_2=\beta_1\beta_2^2\).  Multiplying the first relation by
\(\beta_1\) and by \(\beta_2\) then gives
\(\beta_1^3=\beta_2^3=0\).  Consequently
\[
 (2\beta_1-\beta_2)(-\beta_1+2\beta_2)(\beta_1+\beta_2)
 =3\beta_1^2\beta_2+3\beta_1\beta_2^2=6u_{\mathbb O}.
\]
The primary calculation states that one sixth of this Euler product is a
basis of the top integral cohomology
\cite[Equation~(3.3) and the following paragraph,
p.~497]{h24:MareWillems2013}; hence \(u_{\mathbb O}\) is primitive and
\eqref{h24:eq:octonionic-Euler-six} follows.

The same source constructs the \(S_3\)-indexed cells, with one
eight-dimensional real root space for each inversion of \(\pi\), and
identifies the fixed-point set
\cite[Equations~(2.11)--(2.12) and Proposition~2.10,
pp.~493--494]{h24:MareWillems2013}.  Their real dimensions are
\[
 0,\quad8,8,\quad16,16,\quad24.
\]
All cells are even-dimensional.  The cellular chain groups therefore have
zero differentials, yielding \eqref{h24:eq:octonionic-Poincare} and a second
calculation of the Euler characteristic as the number of cells, namely
six.
Equivalently, the cellwise complex coordinates have dimensions
\(0,4,4,8,8,12\), since \(C_\pi\cong\mathbb C^{4\ell(\pi)}\) as real
cells.  This cell notation supplies no global complex tangent structure.
\end{proof}

\begin{proposition}[The exact relation between the two Weyl data]
\label[proposition]{h24:prop:two-Coxeter-systems}
The ordinary Coxeter number of the ambient root system \(F_4\) is
\[
 h(F_4)=12.
\]
The restricted root system that governs
\eqref{h24:eq:octonionic-six-cells}, however, is \(A_2\); its Weyl group is
\(S_3\) and its Coxeter number is \(3\).  For a common maximal torus of
\(\operatorname{Spin}(8)\) and \(F_4\), there is a split exact sequence
\[
 1\longrightarrow W(D_4)\longrightarrow W(F_4)
   \longrightarrow S_3\longrightarrow1,
 \qquad [W(F_4):W(D_4)]=1152/192=6.
\]
The quotient permutes the three triality representations.  It gives an
exact Weyl-group relation behind the number six; the scalar
\(h(F_4)=12\) alone does not determine the cells or Euler characteristic.
\end{proposition}

\begin{proof}
For a finite root system the invariant degrees satisfy \(d_i=e_i+1\),
and the largest degree is the Coxeter number.  The authoritative \(F_4\)
table gives
\[
 (d_1,d_2,d_3,d_4)=(2,6,8,12),
 \qquad(e_1,e_2,e_3,e_4)=(1,5,7,11),
\]
so \(h(F_4)=12\)
\cite[Section~10.1, p.~213, and the table on p.~216]
{h24:KomoriMatsumotoTsumura2023}.  By contrast, the roots
\(\gamma_1,\gamma_2,\gamma_1+\gamma_2\) in the proof of
\cref{h24:thm:octonionic-cohomology} form \(A_2\).  Its two simple reflections
act as the adjacent transpositions of three coordinates, and their product
is a three-cycle; hence its Coxeter element has order \(3\).
Mare and Willems explicitly identify this restricted Weyl group with
\(S_3\) \cite[pp.~492--493]{h24:MareWillems2013}.

For completeness the extension can be checked in root coordinates.
Take an orthonormal basis \(\varepsilon_1,\ldots,\varepsilon_4\).
The long roots of \(F_4\) are
\(\{\pm\varepsilon_i\pm\varepsilon_j:i<j\}\), the \(D_4\) roots.
Partition the short roots into
\[
 \mathcal V=\{\pm\varepsilon_i\},\qquad
 \mathcal S_\pm=
 \left\{\tfrac12\sum_{i=1}^4\delta_i\varepsilon_i:
           \delta_i\in\{1,-1\},\ \prod_i\delta_i=\pm1\right\}.
\]
These are exactly the three connected components of the graph joining
orthogonal short roots: distinct components have inner products
\(\pm1/2\), and within each component every non-antipodal distinct pair
is orthogonal.  Thus \(W(F_4)\) permutes the three sets.  An element in
the kernel preserves \(\mathcal V\), so is a signed permutation matrix;
preserving both half-spin sets forces an even number of sign changes.
These are exactly \(W(D_4)\), of order \(2^3\cdot4!=192\).
The reflections in \(\varepsilon_1\) and
\(h=(\varepsilon_1+\varepsilon_2+\varepsilon_3+\varepsilon_4)/2\)
induce two distinct transpositions: the first exchanges
\(\mathcal S_+,\mathcal S_-\), and the second exchanges
\(\mathcal V,\mathcal S_-\).  Their product has order three, since
\(\langle\varepsilon_1,h\rangle=1/2\).  They generate a copy of
\(S_3\) mapping isomorphically onto the permutation quotient.  This proves
surjectivity and splitting, and gives \(|W(F_4)|=6\cdot192=1152\).
The three displayed short-root sets are the weights of the three triality
modules, so this is the asserted permutation of those modules.
\end{proof}

\subsubsection{The exact six-chamber comparison}

Let \(\mathscr K_3^{\max}/\mathbb Z^3\) be the translation classes of the maximal
Kuhn tetrahedra \(\sigma_{m,\pi}\) in
\eqref{h24:eq:kuhn-tetrahedron}.  Translation by \(-m\) gives the exact
identification
\[
 \mathscr K_3^{\max}/\mathbb Z^3
 =\{[\sigma_{0,\pi}]:\pi\in S_3\}.
\]
The zero-dimensional torus strata of \(W_3\) are indexed by this set; call
the stratum corresponding to \([\sigma_{0,\pi}]\) by \(z_\pi\).

\begin{theorem}[The exact marked map of finite indexing sets]
\label[theorem]{h24:thm:Kuhn-octonionic-comparison}
After fixing the coordinate order \((e_1,e_2,e_3)\) and the base flag
\(x_{\mathbb O,0}\), the formula
\begin{equation}
 \Theta:W_3^{(0)}\longrightarrow
 \operatorname{Fl}(\mathbb O)^{\operatorname{Spin}(8)},\qquad
 \Theta(z_\pi)=\pi x_{\mathbb O,0}
 \label{h24:eq:Kuhn-flag-bijection}
\end{equation}
is an \(S_3\)-equivariant bijection.  Here
\(W_3^{(0)}=\{z_\pi:\pi\in S_3\}\) is a finite set of torus zero-strata,
and \(S_3\) acts on the source by
permuting coordinate directions and on the target as the restricted Weyl
group.  It identifies the six summands contributing to the two Euler
characteristics, but it does not extend to an asserted identification of
the spaces.
\end{theorem}

\begin{proof}
Every translation class of a maximal Kuhn tetrahedron has one and only one
representative \(\sigma_{0,\pi}\), so the source of \(\Theta\) has exactly
the six displayed elements.  Equation
\eqref{h24:eq:octonionic-six-fixed-points} gives exactly the six displayed
target elements.  Thus \(\Theta\) is bijective.  For
\(\tau\in S_3\), permutation of coordinate directions gives
\[
 \tau\cdot[\sigma_{0,\pi}]=[\sigma_{0,\tau\pi}],
\]
while the restricted Weyl action gives
\(\tau\cdot(\pi x_{\mathbb O,0})=(\tau\pi)x_{\mathbb O,0}\).  Hence \(\Theta\) is
equivariant.

For \(W_3\), every positive-dimensional algebraic-torus stratum has Euler
characteristic zero, so its six zero-strata produce \(\chi(W_3)=6\).
For the flag, the six cells in
\eqref{h24:eq:octonionic-six-cells} are all even-dimensional, so each
contributes \(+1\).  The bijection matches the index sets; the two
stratifications and their positive-dimensional pieces are not identified
by \(\Theta\).
More formally, let \(I_{\mathrm{cell}}=\{\operatorname{label}(C_\pi):
\pi\in S_3\}\) be a set of six labels and put
\(\tau\operatorname{label}(C_\pi)=\operatorname{label}(C_{\tau\pi})\).
Then \(z_\pi\mapsto\pi x_{\mathbb O,0}\mapsto
\operatorname{label}(C_\pi)\) is a chain of equivariant bijections of
finite sets.  The last action is defined on labels only: cells of different
dimensions occur, so no action carrying each cell homeomorphically to
the cell with its translated label is asserted.
\end{proof}

\begin{proposition}[A genuine embedded lower flag, and the degree boundary]
\label[proposition]{h24:prop:complex-flag-embedding}
Choose \(\iota\in\mathbb O\) with \(\iota^2=-1\) and put
\(\mathbb C_\iota=\operatorname{span}_{\mathbb R}\{1,\iota\}\).  Entrywise inclusion in
\(\mathfrak h_3(\mathbb O)\) gives a topological embedding
\begin{equation}
 \jmath_\iota:\operatorname{Fl}_3(\mathbb C_\iota)
 \hookrightarrow\operatorname{Fl}(\mathbb O)
 \label{h24:eq:complex-flag-embedding}
\end{equation}
which sends the six coordinate flags to the six points in
\eqref{h24:eq:octonionic-six-fixed-points}.

There is also an ungraded-ring isomorphism
\begin{equation}
 \rho:
 H^*(\operatorname{Fl}_3(\mathbb C);\mathbb Z)
 \xrightarrow{\ \cong\ }
 H^*(\operatorname{Fl}(\mathbb O);\mathbb Z),\qquad
 x_i\longmapsto\beta_i,\quad i=1,2,
 \label{h24:eq:degree-regraded-ring-map}
\end{equation}
where the source has the presentation in
\eqref{h24:eq:octonionic-cohomology-ring} with \(|x_i|=2\), while
\(|\beta_i|=8\).  No continuous map in either direction can induce this
degree regrading.
\end{proposition}

\begin{proof}
The subalgebra \(\mathbb C_\iota\) is associative and isomorphic to \(\mathbb C\).
Rank-one Hermitian projectors in
\(\mathfrak h_3(\mathbb C_\iota)\) remain idempotent trace-one elements after
entrywise inclusion, and orthogonality is the unchanged equation
\(\operatorname{Re}\operatorname{tr}(ab)=0\).  The resulting map is injective.
Its domain is compact and its target Hausdorff, so it is a topological
embedding.  The diagonal projectors are unchanged, which proves the
fixed-point assertion.

The standard complete-flag presentation is
\[
 H^*(\operatorname{Fl}_3(\mathbb C);\mathbb Z)
 \cong
 \frac{\mathbb Z[x_1,x_2]}
 {(x_1^2-x_1x_2+x_2^2,\,
   x_1x_2(x_1-x_2))},\qquad |x_i|=2;
\]
Mare and Willems derive the octonionic presentation precisely by comparison
with this \(A_2\) flag ring
\cite[Theorem~1.1, p.~484, and the proof on
pp.~496--497]{h24:MareWillems2013}.  Thus \(\rho\) is an isomorphism after
forgetting degrees.  If \(f:\operatorname{Fl}(\mathbb O)\to
\operatorname{Fl}_3(\mathbb C)\) were to induce \(\rho\), then
\(f^*(x_i)\in H^2(\operatorname{Fl}(\mathbb O);\mathbb Z)=0\).  If a map in the other
direction were to induce \(\rho^{-1}\), then the image of each
\(\beta_i\) would lie in
\(H^8(\operatorname{Fl}_3(\mathbb C);\mathbb Z)=0\).  Either conclusion contradicts the
nonzero generators in \eqref{h24:eq:degree-regraded-ring-map}.
\end{proof}

\subsubsection{The octonionic determinant-boundary sphere}

For \(\lambda\in\mathbb R\) and \(q_1,q_2,q_3\in\mathbb O\), retain the ordered
coordinates
\begin{equation}
 Q_{\mathbb O}(\lambda;q_1,q_2,q_3)=
 \begin{pmatrix}
 \lambda&q_3&\overline q_2\\
 \overline q_3&\lambda&q_1\\
 q_2&\overline q_1&\lambda
 \end{pmatrix}
 \in\mathfrak h_3(\mathbb O).
 \label{h24:eq:octonionic-equal-diagonal}
\end{equation}
The equal-diagonal slice and its positive cone are
\begin{equation}
 \mathcal T(\mathbb O)=
 \{Q_{\mathbb O}(\lambda;q_1,q_2,q_3)\}
 \cong\mathbb R\oplus\mathbb O^3,\qquad
 \mathcal C(\mathbb O)=\mathcal T(\mathbb O)\cap
                 \mathfrak h_3(\mathbb O)_{\succeq0}.
 \label{h24:eq:octonionic-equal-diagonal-cone}
\end{equation}
Thus \(\dim_{\mathbb R}\mathcal T(\mathbb O)=1+3\cdot8=25\).  The Jordan determinant
on this slice is the exact cubic
\begin{equation}
 \det_{\mathbb O}Q_{\mathbb O}
 =\lambda^3-\lambda\bigl(
 |q_1|^2+|q_2|^2+|q_3|^2\bigr)
 +2\operatorname{Re}\bigl((q_1q_2)q_3\bigr).
 \label{h24:eq:octonionic-equal-diagonal-determinant}
\end{equation}
The real part is unchanged if \(q_1(q_2q_3)\) is used, because the
octonionic associator has zero real part.

\begin{theorem}[The exact octonionic \(S^{24}\) determinant carrier]
\label[theorem]{h24:thm:octonionic-determinant-S24}
The relative boundary of \(\mathcal C(\mathbb O)\) inside
\(\mathcal T(\mathbb O)\) is
\[
 \partial\mathcal C(\mathbb O)
 =\{Q\in\mathcal C(\mathbb O):\det_{\mathbb O}Q=0\}.
\]
There are homeomorphisms of pairs
\begin{equation}
 \bigl(\mathcal C(\mathbb O)^+,
       (\partial\mathcal C(\mathbb O))^+\bigr)
 \cong(B^{25},S^{24}).
 \label{h24:eq:octonionic-determinant-pair}
\end{equation}
In particular
\[
 S_{\det,\mathbb O}^{24}:=(\partial\mathcal C(\mathbb O))^+\cong S^{24}
\]
is an exact determinant-boundary sphere.  It is not the
\(24\)-dimensional flag \(F_4/\operatorname{Spin}(8)\).
\end{theorem}

\begin{proof}
The nonnegative cone of the Euclidean Jordan algebra
\(\mathfrak h_3(\mathbb O)\) is closed and convex.  Its intersection
\[
 K_{\mathbb O}=\{(q_1,q_2,q_3)\in\mathbb O^3:
              Q_{\mathbb O}(1;q_1,q_2,q_3)\succeq0\}
\]
with the affine hyperplane \(\lambda=1\) is therefore closed and convex.
Every positive semidefinite \(2\times2\) principal corner gives
\(|q_i|\le1\), so \(K_{\mathbb O}\) is bounded and hence compact.
The origin of \(\mathbb O^3\) corresponds to \(I_3\), which is positive
definite, so it is an interior point of \(K_{\mathbb O}\).

For \(v\in S(\mathbb O^3)\), let
\[
 R_{\mathbb O}(v)=\max\{r\ge0:rv\in K_{\mathbb O}\}.
\]
Convexity, compactness, and the interior-point property make
\(R_{\mathbb O}:S(\mathbb O^3)\to\mathbb R_{>0}\) continuous.  Indeed, choose
\(0<\epsilon<M\) with
\(\epsilon B^{24}\subset K_{\mathbb O}\subset M B^{24}\).  The Minkowski
functional
\[
 p_K(x)=\inf\{t>0:x\in tK_{\mathbb O}\}
\]
is positive for \(x\ne0\), positively homogeneous, and subadditive:
if \(x\in aK_{\mathbb O}\) and \(y\in bK_{\mathbb O}\), convexity gives
\(x+y\in(a+b)K_{\mathbb O}\).  Taking infima proves subadditivity.
As \(p_K(w)\le\|w\|/\epsilon\), subadditivity in both directions gives
\(|p_K(x)-p_K(y)|\le\|x-y\|/\epsilon\).  Hence
\(R_{\mathbb O}(v)=1/p_K(v)\) is continuous.  Radial rescaling
\[
 0\longmapsto0,\qquad rv\longmapsto
 \frac{r}{R_{\mathbb O}(v)}v
\]
is consequently a homeomorphism \(K_{\mathbb O}\cong B^{24}\) carrying
\(\partial K_{\mathbb O}\) onto \(S^{23}\).  Denote it by \(F_K\); its
inverse sends \(sv\), \(0<s\le1\), to \(sR_{\mathbb O}(v)v\), and sends
zero to zero.  Both maps are continuous at zero because the radial
function is bounded above and below by positive constants.

A nonzero element of \(\mathcal C(\mathbb O)\) has \(\lambda>0\): if
\(\lambda=0\), the same principal-corner inequalities force
\(q_1=q_2=q_3=0\).  Positive scaling therefore reconstructs the whole
cone as
\[
 \mathcal C(\mathbb O)=\{0\}\cup
 \{\lambda Q_{\mathbb O}(1;q_1,q_2,q_3):
       \lambda>0,\ (q_1,q_2,q_3)\in K_{\mathbb O}\}.
\]
Here are explicit maps through the Lorentz cone, with
\(q=(q_1,q_2,q_3)\):
\[
 \begin{aligned}
 \mathcal C(\mathbb O)&\longrightarrow
       L_{25}:=\{(s,x)\in\mathbb R\oplus\mathbb O^3:s\ge\|x\|\},\\
 Q_{\mathbb O}(\lambda;q)&\longmapsto
       (\lambda,\lambda F_K(q/\lambda))\quad(\lambda>0),\\
 0&\longmapsto(0,0),\\
 L_{25}&\longrightarrow[0,\infty)\times\mathbb O^3,\\
 (s,x)&\longmapsto(s-\|x\|,x).
 \end{aligned}
\]
The inverses are \((s,x)\mapsto Q_{\mathbb O}(s;sF_K^{-1}(x/s))\) for
\(s>0\), with zero sent to zero, and
\((t,x)\mapsto(t+\|x\|,x)\), respectively.  The boundedness of
\(F_K\) and \(F_K^{-1}\) proves continuity at the cone vertices.
The boundary maps to \(s=\|x\|\), then to \(t=0\), proving an explicit
pair homeomorphism with a closed half-space and its boundary.

To identify this boundary with determinant zero in the relative topology,
use the spectral decomposition in the Euclidean Jordan algebra.
A positive definite element has three positive eigenvalues and is an
interior point even in the full algebra, hence in the slice.  Conversely,
if a positive semidefinite \(Q\) has a zero eigenvalue, then
\(Q-\delta I_3\in\mathcal T(\mathbb O)\), for every \(\delta>0\), has a
negative eigenvalue.  Such points lie outside \(\mathcal C(\mathbb O)\) and
converge to \(Q\); thus \(Q\) is not a relative interior point.
Since the cone is closed, these are exactly its relative boundary points.
The determinant is the product of the three nonnegative eigenvalues, so
its zero set in the cone is exactly that boundary.

Fixing an orthonormal identification \(\mathbb O^3\cong\mathbb R^{24}\), the map
\[
 (t,x)\longmapsto
 \left(\frac{2x}{D},\frac{t^2+\|x\|^2-1}{D}\right)\in B^{25},
 \qquad D=1+t^2+\|x\|^2,
\]
extends at infinity to \((0,1)\in\partial B^{25}\).
It is inverse stereographic projection to the closed hemisphere in
\(S^{25}\) with first coordinate nonnegative, followed by deletion of
that first coordinate.  The hemisphere projection is bijective, with
inverse adjoining \(\sqrt{1-\|y\|^2}\); hence the displayed map is a
homeomorphism after one-point compactification.  At \(t=0\) it is
\[
 x\longmapsto
 \left(\frac{2x}{1+\|x\|^2},
       \frac{\|x\|^2-1}{1+\|x\|^2}\right)\in S^{24},
\]
with infinity sent to \((0,1)\).  This proves
\eqref{h24:eq:octonionic-determinant-pair}, including its specified boundary.
\end{proof}

% Staged inclusion. No canonical file was edited.
\subsubsection{The ordered Albert triality tensor and its global descent}
\label{h24:att:sec:ordered-tensor}

\paragraph{Original coordinates and octonion identities}

Let $\mathbb O$ be the real octonion division algebra, with its original
conjugation, norm $N(a)=a\bar a=\bar a a$, and real inner product
$h(a,b)=\operatorname{Re}(a\bar b)$. Put $J=\mathfrak h_3(\mathbb O)$,
with its Albert Jordan product $X\circ Y=(XY+YX)/2$. Matrix multiplication
here involves only products of two octonions in each entry; no product
of three octonionic matrices is left unparenthesized. Set
$e_i=\operatorname{diag}(\delta_{1i},\delta_{2i},\delta_{3i})$,
$I=e_1+e_2+e_3$, and
\[
 E_1(a)=\begin{pmatrix}0&0&0\\0&0&a\\0&\bar a&0\end{pmatrix},\quad
 E_2(b)=\begin{pmatrix}0&0&\bar b\\0&0&0\\b&0&0\end{pmatrix},\quad
 E_3(c)=\begin{pmatrix}0&c&0\\\bar c&0&0\\0&0&0\end{pmatrix}.
\]
The specified equal-diagonal object remains exactly
\begin{equation}
 Q_{\mathbb O}(\lambda;q_1,q_2,q_3)
 =\lambda I+E_1(q_1)+E_2(q_2)+E_3(q_3)
 =\begin{pmatrix}\lambda&q_3&\bar q_2\\
 \bar q_3&\lambda&q_1\\q_2&\bar q_1&\lambda\end{pmatrix}.
 \label{h24:att:eq:Q}
\end{equation}
Let $V_i$ be three labelled copies of $\mathbb O$ and $J_i=E_i(V_i)$.
Labelling does not identify their group actions. The original trace
metric on $J$ is
\begin{equation}
\begin{split}
 G(X,Y)&=\operatorname{tr}(X\circ Y),\\
 G\left(\sum_i\xi_i e_i+\sum_iE_i(a_i),
        \sum_i\eta_i e_i+\sum_iE_i(b_i)\right)
 &=\sum_i\xi_i\eta_i+2\sum_i h(a_i,b_i).
\end{split}
\label{h24:att:eq:metrics}
\end{equation}
In particular $E_i^*G=2h$, not $h$.

We use alternativity, $N(ab)=N(a)N(b)$, and conjugation reversing
products. The parenthesis changes used below follow as follows.
Because $\bar a=2\operatorname{Re}(a)-a$, alternativity gives
\begin{equation}
 \bar a(ab)=N(a)b,\quad a(\bar a b)=N(a)b,\quad
 (ba)\bar a=N(a)b,\quad (b\bar a)a=N(a)b.
 \label{h24:att:eq:inverse-property}
\end{equation}
Polarizing norm composition gives
$h(ab,ac)=N(a)h(b,c)$ and $h(ba,ca)=N(a)h(b,c)$.
For $a\ne0$, these equations and \eqref{h24:att:eq:inverse-property}
imply $L_a^*=L_{\bar a}$ and $R_a^*=R_{\bar a}$; for $a=0$
the identities are immediate. Consequently
\begin{align}
 h(ab,c)&=h(b,\bar a c)=h(a,c\bar b),\nonumber\\
 \operatorname{Re}((ab)c)&=\operatorname{Re}(a(bc))
                         =\operatorname{Re}((bc)a).
 \label{h24:att:eq:real-associativity}
\end{align}
For the first real-part equality, both expressions equal
$h(c,\bar b\bar a)$ by the adjoint identities. For the second, use
$\operatorname{Re}(xy)=h(x,\bar y)=h(y,\bar x)=\operatorname{Re}(yx)$.
These are identities of real parts; they do not equate the octonions
$(ab)c$ and $a(bc)$.

\paragraph{The complete Peirce product table}

A cyclic triple means exactly $(i,j,k)=(1,2,3),(2,3,1),(3,1,2)$.

\begin{proposition}\label{h24:att:prop:peirce}
The decomposition $J=\bigoplus_i\mathbb R e_i\oplus\bigoplus_iJ_i$
has the complete product table
\begin{align}
 e_i\circ e_j&=\delta_{ij}e_i,\nonumber\\
 e_i\circ E_i(a)&=0,&
 e_j\circ E_i(a)&=e_k\circ E_i(a)=\tfrac12E_i(a),\nonumber\\
 E_i(a)\circ E_i(b)&=h(a,b)(e_j+e_k),\label{h24:att:eq:peirce-table}\\
 E_i(a)\circ E_j(b)&=\tfrac12E_k(\overline{ab}).\nonumber
\end{align}
Reversing the Jordan arguments uses commutativity of $\circ$,
without reversing the octonion order in the cyclic formula.
\end{proposition}

\begin{proof}
The first three formulas follow entry by entry. For example the two
nonzero diagonal entries of $E_1(a)\circ E_1(b)$ are
$(a\bar b+b\bar a)/2$ and $(\bar a b+\bar b a)/2$, both $h(a,b)$.
The mixed products have precisely the following nonzero entries:
\[
\begin{array}{c|c|c}
\text{product}&\text{first entry}&\text{second entry}\\\hline
E_1(a)\circ E_2(b)&(1,2)=\bar b\bar a/2&(2,1)=ab/2\\
E_2(b)\circ E_3(c)&(2,3)=\bar c\bar b/2&(3,2)=bc/2\\
E_3(c)\circ E_1(a)&(3,1)=\bar a\bar c/2&(1,3)=ca/2
\end{array}
\]
These are respectively $E_3(\overline{ab})/2$,
$E_1(\overline{bc})/2$, and $E_2(\overline{ca})/2$, by the original
definitions of $E_1,E_2,E_3$.
\end{proof}

\paragraph{Three typed products, adjoints, and inverses}

Define the trilinear form and all three bilinear maps by
\begin{align}
 \tau:V_1\times V_2\times V_3&\longrightarrow\mathbb R,&
 \tau(a,b,c)&=\operatorname{Re}((ab)c),\label{h24:att:eq:tau}\\
 \mu_{12}:V_1\times V_2&\longrightarrow V_3,&
 \mu_{12}(a,b)&=\overline{ab}=\bar b\bar a,\nonumber\\
 \mu_{23}:V_2\times V_3&\longrightarrow V_1,&
 \mu_{23}(b,c)&=\overline{bc}=\bar c\bar b,\nonumber\\
 \mu_{31}:V_3\times V_1&\longrightarrow V_2,&
 \mu_{31}(c,a)&=\overline{ca}=\bar a\bar c.\nonumber
\end{align}
Every $V_i$ carries $h$, with its label retained. Equation
\eqref{h24:att:eq:real-associativity} and conjugation give
\begin{equation}
 \tau(a,b,c)=h(\mu_{12}(a,b),c)
 =h(\mu_{23}(b,c),a)=h(\mu_{31}(c,a),b).
 \label{h24:att:eq:riesz}
\end{equation}
For example $h(\overline{ab},c)=
\operatorname{Re}(\overline{ab}\bar c)=
\operatorname{Re}(\overline{c(ab)})=\operatorname{Re}((ab)c)$.
Thus these are exactly the three Riesz-curried maps for $h$.
In the original matrix metric,
\begin{equation}
 \tau(a,b,c)=G(E_1(a)\circ E_2(b),E_3(c))
 =\operatorname{tr}((E_1(a)\circ E_2(b))\circ E_3(c)).
 \label{h24:att:eq:tau-jordan}
\end{equation}

\begin{proposition}\label{h24:att:prop:norm-inverse}
For cyclic $(i,j,k)$, $a\in V_i$, and $b\in V_j$,
\begin{align}
 N(\mu_{ij}(a,b))&=N(a)N(b),\label{h24:att:eq:composition}\\
 \mu_{ki}(\mu_{ij}(a,b),a)&=N(a)b,\label{h24:att:eq:inverse-one}\\
 \mu_{jk}(b,\mu_{ij}(a,b))&=N(b)a.\label{h24:att:eq:inverse-two}
\end{align}
For fixed $a\ne0$, the real-linear isomorphism
$M_a:V_j\to V_k$, $M_a(b)=\mu_{ij}(a,b)$, has
\[
 M_a^*(c)=\mu_{ki}(c,a),\quad
 M_a^*M_a=N(a)\operatorname{id}_{V_j},\quad
 M_aM_a^*=N(a)\operatorname{id}_{V_k},\quad
 M_a^{-1}(c)=\frac{\mu_{ki}(c,a)}{N(a)}.
\]
For fixed $b\ne0$, the inverse of $a\mapsto\mu_{ij}(a,b)$ is
$c\mapsto\mu_{jk}(b,c)/N(b)$, and its numerator is also the adjoint.
\end{proposition}

\begin{proof}
Norm composition proves the first formula. The inverse numerators are,
with their parentheses preserved,
\[
 \overline{\overline{ab}\,a}=\bar a(ab)=N(a)b,\qquad
 \overline{b\,\overline{ab}}=(ab)\bar b=N(b)a.
\]
Their final equalities are \eqref{h24:att:eq:inverse-property}. Relabelling
by the three stated cyclic triples proves all six inverse identities.
The adjoint statements are \eqref{h24:att:eq:riesz}. Both adjoint-composition
identities follow from the two inverse identities, including with the
other argument nonzero. Division by $N(a)$ or $N(b)$ gives the inverses.
No unit-length hypothesis has been imposed.
\end{proof}

Retaining $G$ gives the following different numerical identities.
For $X=E_i(a)\in J_i,Y=E_j(b)\in J_j$ with $i\ne j$, put
$B_{ij}(X,Y)=2X\circ Y=E_k(\mu_{ij}(a,b))$, using cyclic argument order.
Then
\begin{equation}
 G(B_{ij}(X,Y),B_{ij}(X,Y))=\tfrac12G(X,X)G(Y,Y),\qquad
 G(X\circ Y,X\circ Y)=\tfrac18G(X,X)G(Y,Y).
 \label{h24:att:eq:trace-norm-factors}
\end{equation}
Indeed the left sides are $2N(a)N(b)$ and $N(a)N(b)/2$,
whereas $G(X,X)G(Y,Y)=4N(a)N(b)$. The coefficient-one composition
identity uses the explicitly stated $h$, not the trace metric $G$.

\paragraph{The determinant and the factor six in polarization}

\begin{proposition}\label{h24:att:prop:determinant}
For arbitrary real $\xi_i$ and the original octonionic coordinates
$q_i$, let $X=\sum_i\xi_i e_i+\sum_iE_i(q_i)$. Its Jordan determinant is
\begin{equation}
 d(X)=\xi_1\xi_2\xi_3-\sum_{i=1}^3\xi_iN(q_i)
                    +2\operatorname{Re}((q_1q_2)q_3).
 \label{h24:att:eq:determinant-general}
\end{equation}
In particular,
\[
 d(Q_{\mathbb O}(\lambda;q_1,q_2,q_3))
 =\lambda^3-\lambda\bigl(N(q_1)+N(q_2)+N(q_3)\bigr)
                  +2\operatorname{Re}((q_1q_2)q_3).
\]
\end{proposition}

\begin{proof}
The Albert cubic norm is given by the Jordan trace formula
\[
 d(X)=\tfrac13\operatorname{tr}(X\circ(X\circ X))
 -\tfrac12\operatorname{tr}(X)\operatorname{tr}(X\circ X)
 +\tfrac16\operatorname{tr}(X)^3.
\]
The product table computes the diagonal entries and three coordinates
of $X^2$:
\[
 (X^2)_{ii}=\xi_i^2+N(q_j)+N(q_k),\qquad
 (X^2)_i=(\xi_j+\xi_k)q_i+\overline{q_jq_k},
 \quad(i,j,k)\text{ cyclic}.
\]
With $s=\xi_1+\xi_2+\xi_3$, it follows that
\begin{align*}
 \operatorname{tr}(X)&=s,\\
 \operatorname{tr}(X^2)&=\sum_i\xi_i^2+2\sum_iN(q_i),\\
 \operatorname{tr}(X\circ X^2)
 &=\sum_i\xi_i^3+3\sum_i(\xi_j+\xi_k)N(q_i)
                     +6\operatorname{Re}((q_1q_2)q_3).
\end{align*}
Each scalar product of $q_i$ with $\overline{q_jq_k}$ equals
$\tau(q_1,q_2,q_3)$ by \eqref{h24:att:eq:riesz}; the trace metric supplies
the factor $2$. Substitution into the cubic gives coefficient
$(\xi_j+\xi_k)-s=-\xi_i$ for $N(q_i)$ and coefficient $6/3=2$ for
$\tau$. Direct expansion gives the scalar part
\[
 \tfrac13\sum_i\xi_i^3-\tfrac12s\sum_i\xi_i^2+\tfrac16s^3
 =\xi_1\xi_2\xi_3.
\]
This proves \eqref{h24:att:eq:determinant-general}. Setting all three
$\xi_i=\lambda$ retains exactly the original equal-diagonal slice.
\end{proof}

Define the symmetric polarization $\mathcal D$ by
$\mathcal D(X,X,X)=d(X)$, explicitly
\begin{align}
 6\mathcal D(X,Y,Z)={}&d(X+Y+Z)-d(X+Y)-d(X+Z)-d(Y+Z)
 \nonumber\\&+d(X)+d(Y)+d(Z).
 \label{h24:att:eq:polarization}
\end{align}
Expansion of homogeneous cubic monomials proves symmetry and
trilinearity: inclusion--exclusion removes every term missing an
argument, and degree three forces each remaining argument to occur
once. On the diagonal the multiplier is $27-3\cdot8+3=6$.
The same expansion proves
\[
 \left.\frac{\partial^3}{\partial r\partial s\partial t}
 d(rX+sY+tZ)\right|_{r=s=t=0}=6\mathcal D(X,Y,Z).
\]
For the three distinct sectors one obtains
\begin{equation}
\begin{split}
 \left.\partial_r\partial_s\partial_t
 d(rE_1(a)+sE_2(b)+tE_3(c))\right|_0&=2\tau(a,b,c),\\
 \mathcal D(E_1(a),E_2(b),E_3(c))&=\tfrac13\tau(a,b,c).
\end{split}
 \label{h24:att:eq:factor-six}
\end{equation}
Indeed that determinant is $2rst\tau(a,b,c)$. Relative to the original
$G$, the three Riesz representatives for fixed $E_1(a),E_2(b)$ are
\[
\begin{array}{c|c}
\text{scalar form on }J_1\times J_2\times J_3&
             \text{representative in }J_3\\\hline
\tau&E_3(\overline{ab})/2=E_1(a)\circ E_2(b)\\
\partial_r\partial_s\partial_t d&
                  E_3(\overline{ab})=2E_1(a)\circ E_2(b)\\
\mathcal D&E_3(\overline{ab})/6=(E_1(a)\circ E_2(b))/3
\end{array}
\]
Thus the determinant coefficient $2$ and polarization factor $6$
are separate, fully specified operations.

\paragraph{The exact stabilizer and its three actions}

Let $\kappa(a)=\bar a$ and define
\begin{equation}
\begin{split}
H_\tau=\{(g_1,g_2,g_3)\in O(V_1,h)\times O(V_2,h)\times O(V_3,h):\quad&\\
\tau(g_1a,g_2b,g_3c)=\tau(a,b,c)\text{ for every }a,b,c&\}.
\end{split}
\label{h24:att:eq:metric-stabilizer}
\end{equation}
By \eqref{h24:att:eq:riesz}, orthogonality and nondegeneracy of $h$,
this is equivalent to each of the equations
\begin{align}
 \mu_{12}(g_1a,g_2b)&=g_3\mu_{12}(a,b),\nonumber\\
 \mu_{23}(g_2b,g_3c)&=g_1\mu_{23}(b,c),\label{h24:att:eq:equivariance}\\
 \mu_{31}(g_3c,g_1a)&=g_2\mu_{31}(c,a).\nonumber
\end{align}
With every conjugation retained, the first equation reads
\begin{equation}
 (g_1a)(g_2b)=\kappa g_3\kappa(ab).
 \label{h24:att:eq:related-triple}
\end{equation}
Consequently $(g_1,g_2,g_3)\mapsto(g_1,g_2,\kappa g_3\kappa)$
is a group isomorphism onto the ordinary product-related triples
$\mathcal R=\{(A,B,C)\in O(8)^3:(Ax)(By)=C(xy)\text{ for all }x,y\}$.

\begin{theorem}\label{h24:att:thm:stabilizer}
Every element of $H_\tau$ has $g_i\in SO(8)$, and $H_\tau$ is
isomorphic to the compact simply connected $\operatorname{Spin}(8)$.
The map
\[
\begin{split}
 \Phi:H_\tau&\longrightarrow
 \{g\in\operatorname{Aut}(J,\circ):g(e_i)=e_i\ (i=1,2,3)\},\\
 \Phi(g)\left(\sum_i\xi_i e_i+\sum_iE_i(a_i)\right)
 &=\sum_i\xi_i e_i+\sum_iE_i(g_i a_i)
\end{split}
\]
is an exact group isomorphism. Thus \eqref{h24:att:eq:equivariance} is
$\operatorname{Spin}(8)$ equivariance for the actions on the original
Albert coordinates.
\end{theorem}

\begin{proof}
For unit $u,v\in\mathbb O$, reflection in $u^\perp$ is
\[
 D_u(x)=x-2h(x,u)u=-u\bar x u.
\]
To verify the second expression, multiply
$x\bar u+u\bar x=2h(x,u)$ on the right by $u$ and use
$(x\bar u)u=x$. Define
\[
 A_{u,v}(x)=v(\bar u x),\quad
 B_{u,v}(y)=(y\bar u)v,\quad
 C_{u,v}(z)=v(\bar u z\bar u)v=D_vD_u(z).
\]
Repeated outside letters are parenthesized by flexibility. The
Moufang identity $(vx)(yv)=v(xy)v$, first for $v$ and then for $\bar u$,
proves
\[
 (v(\bar u x))((y\bar u)v)=v(\bar u(xy)\bar u)v.
\]
This is a product-related triple. All its coordinates lie in $SO(8)$:
multiplication by a unit octonion is orthogonal, and its determinant
is continuously $+1$ on the connected unit sphere because it is $+1$
at $1$. Every element of $SO(8)$ is a product of an even number of
reflections. To obtain such a decomposition, reflect the image of the
first vector of an orthonormal basis onto that vector, then repeat
inside its orthogonal complement. After finitely many steps the map is
the identity; the determinant records the parity of the reflections.
Projection of $\mathcal R\cap SO(8)^3$ to $C$ is therefore surjective.

Its kernel contains exactly two elements. If $(Ax)(By)=xy$, put
$p=A1$, $q=B1$. Then $pq=1$, $By=\bar p y$, and $Ax=xp$.
The equation becomes $(xp)(\bar p y)=xy$. Replacing $y$ by $py$
gives $(xp)y=x(py)$ for every $x,y$. Hence $p$ lies in the middle
nucleus of $\mathbb O$, which is $\mathbb R$. The latter fact can be
checked without an implicit associativity assumption. In the
Cayley--Dickson presentation $\mathbb O=\mathbb H\oplus\mathbb H\ell$
with
\[
 (a+b\ell)(c+d\ell)=(ac-\bar d b)+(da+b\bar c)\ell,
\]
write $p=a+b\ell$ and use quaternion units
$\mathsf i\mathsf j=\mathsf k$. The $\ell$ coefficient of
$(\mathsf i p)\mathsf j-\mathsf i(p\mathsf j)$ is $-2b\mathsf k$,
so $b=0$. The $\ell$ coefficient of
$(\ell a)x-\ell(ax)$ is $\bar a\bar x-\bar x\bar a$.
Its vanishing for $x=\mathsf i,\mathsf j$ forces $a$ real.
Since $|p|=1$, the kernel is
$\{(I,I,I),(-I,-I,I)\}$.

There are no orientation-reversing related triples. First, the cyclic
invariance of $\tau$ gives surjectivity onto $SO(8)$ in each of its
three coordinates, and hence the same holds for the three projections
of $\mathcal R\cap SO(8)^3$. If a product-related triple had
$\det A=-1$, compose it with an orientation-preserving product-related
triple whose first coordinate is $\kappa A^{-1}$. The result has first
coordinate $\kappa$, so satisfies $\bar x(By)=C(xy)$.
Setting $x=1$ yields $B=C$. Setting $y=1$, and writing $p=B1$, yields
$Bx=\bar x p$. Hence
\[
 \bar x(\bar y p)=\overline{xy}\,p\qquad\text{for all }x,y.
\]
Set $y=p$ and multiply on the right by $\bar p$. The inverse property
gives $\bar x\bar p=\bar p\bar x$ for all $x$. The center of
$\mathbb O$ is $\mathbb R$: in the displayed Cayley--Dickson
presentation, commuting $a+b\ell$ with $\mathsf i$ forces $b=0$,
and commuting $a$ with $\mathsf i,\mathsf j$ forces $a$ real.
Thus $p=\pm1$. The preceding equation would imply
$\bar x\bar y=\overline{xy}=\bar y\bar x$ for all $x,y$,
contradicting $\mathsf i\mathsf j=-\mathsf j\mathsf i$.
Thus $\det A=1$. Cyclic permutation of $H_\tau$ excludes negative
determinants in the other two coordinates too.

The group is connected. The reflection-pair triples are continuously
parametrized by $S^7\times S^7$, contain the identity at $u=v=1$,
and contain $(-I,-I,I)$ at $u=1,v=-1$. Their third coordinates
generate $SO(8)$, so every related triple is a product of these
triples and a kernel element. Every such element therefore lies in
the identity component. Projection to $C$ is a two-sheeted covering:
it is a surjective homomorphism of compact Lie groups with finite
kernel, hence a local diffeomorphism with those two-element fibres.
Here $\pi_1(SO(8))=\mathbb Z/2$. One proof starts with the unit
quaternion conjugation action on imaginary quaternions, which
identifies $SO(3)$ with $S^3/\{\pm1\}$, and uses the fibrations
$SO(n-1)\to SO(n)\to S^{n-1}$ for $n\ge4$; their bases have zero
$\pi_1$ and $\pi_2$, so the fundamental group remains $\mathbb Z/2$.
A connected two-sheeted cover of this group is simply connected.
It is $\operatorname{Spin}(8)$, the simply connected cover of $SO(8)$.

The map $\Phi$ preserves every product in
\eqref{h24:att:eq:peirce-table}, by orthogonality and
\eqref{h24:att:eq:equivariance}; bilinearity proves preservation of all
Jordan products. Its inverse is $\Phi(g^{-1})$. Conversely, a Jordan
automorphism fixing all $e_i$ preserves
\[
 J_i=\{X\in J:e_i\circ X=0,\quad e_j\circ X=e_k\circ X=X/2\}.
\]
Its restrictions define linear maps $g_i$ in the original $E_i$
coordinates. Applying the automorphism to the same-sector Peirce
product proves $h(g_i a,g_i b)=h(a,b)$; applying it to the mixed
product proves \eqref{h24:att:eq:related-triple}. Thus the restrictions
belong to $H_\tau$. The direct-sum decomposition of $J$ proves
injectivity and surjectivity of $\Phi$ with no further kernel.
\end{proof}

\begin{proposition}\label{h24:att:prop:linear-stabilizer}
Without the metrics, the stabilizer of the same typed tensor in
$GL(V_1)\times GL(V_2)\times GL(V_3)$ is exactly
\[
 \{(s_1g_1,s_2g_2,s_3g_3):s_i>0,\quad s_1s_2s_3=1,\quad
                          (g_1,g_2,g_3)\in H_\tau\}.
\]
Thus the metric-preserving stabilizer is $\operatorname{Spin}(8)$;
the unrestricted typed linear stabilizer has two additional
independent positive scale parameters.
\end{proposition}

\begin{proof}
Tensor invariance for $(A,B,C)$ and the $h$ adjoint imply
$\mu_{12}(Aa,Bb)=C^{-T}\mu_{12}(a,b)$.
For nonzero $a$, Proposition~\ref{h24:att:prop:norm-inverse} gives
$|\det M_a|=N(a)^4$ on dimension eight. Taking determinants yields
$N(Aa)^4=|\det B\det C|^{-1}N(a)^4$. Positivity therefore gives
$A=s_1g_1$, where $s_1>0$ and $g_1$ is orthogonal. The cyclic equations
give $B=s_2g_2$, $C=s_3g_3$ likewise. For unit $a,b$, the two sides
of the adjoint equation have norms $s_1s_2$ and $s_3^{-1}$, so
$s_1s_2s_3=1$. Dividing by those explicit scalars proves
$(g_1,g_2,g_3)\in H_\tau$. The converse follows by trilinearity.
\end{proof}

The three actions can also be read directly as the vector and two
real half-spin actions. On $S=V_2\oplus V_3$ define
\begin{equation}
 c(a)(b,c_0)=(-\mu_{31}(c_0,a),\mu_{12}(a,b)),\qquad a\in V_1.
 \label{h24:att:eq:clifford}
\end{equation}
The adjoint and inverse identities prove
$c(a)^*=-c(a)$ and $c(a)^2=-N(a)I$. Polarization gives
\[
 c(a)c(a')+c(a')c(a)=-2h(a,a')I.
\]
This is the real Clifford action with negative norm convention.
Even products preserve each of the two summands.
Equation \eqref{h24:att:eq:equivariance} gives the exact intertwining map
\[
 \operatorname{diag}(g_2,g_3)c(a)
 =c(g_1a)\operatorname{diag}(g_2,g_3).
\]
Conversely, an even product of unit Clifford operators covers the
corresponding product of reflections on $V_1$, is orthogonal on $S$,
and satisfies this equation. Its two diagonal blocks, together with
its vector action, therefore form a triple in $H_\tau$. This map
from the spin double cover is onto: it covers $SO(V_1)$ and includes
both elements over the identity, because the Clifford scalar $-1$
acts as $-I$ on both summands. The two groups are connected double
covers and the map is an isomorphism. This identifies $V_1$ as the
vector representation and $V_2,V_3$ as the two half-spin modules in
the displayed Clifford realization, without conjugating any $q_i$.
Their names $S^+,S^-$ can be fixed by the parity summands; when the
names are instead defined by an oriented volume element, that
orientation must be specified as additional notation.

There is no nonzero $\operatorname{Spin}(8)$ equivariant linear map
$V_i\to V_j$ for $i\ne j$. The four triples
\[
 (I,I,I),\quad(I,-I,-I),\quad(-I,I,-I),\quad(-I,-I,I)
\]
belong to $H_\tau$ by the product equation. For any two different
sectors, one of these acts as $+I$ on the source and $-I$ on the
target. An equivariant linear map then equals its own negative and
vanishes. Their exact positive relation is the three nonzero
equivariant bilinear maps \eqref{h24:att:eq:tau}; fixing any nonzero
argument gives the intersector isomorphisms and inverses already
proved in Proposition~\ref{h24:att:prop:norm-inverse}.

\paragraph{An explicit tangent comparison and global bundle descent}

Put $F_4=\operatorname{Aut}(J,\circ)$ and use
Theorem~\ref{h24:att:thm:stabilizer} to identify $H_\tau$ with the ordered
frame stabilizer. Then $M=F_4/H_\tau$ parametrizes the orbit of the
ordered frame $(e_1,e_2,e_3)$.

Every Jordan automorphism preserves $\operatorname{tr}$ and $G$.
Indeed the Peirce table gives
$\operatorname{Tr}_{\operatorname{End}(J)}L_{e_i}=1+16/2=9$ and
$\operatorname{Tr}_{\operatorname{End}(J)}L_{E_i(a)}=0$.
The latter operator has zero diagonal blocks in the Peirce
decomposition, as each product in its table changes summand.
Thus $\operatorname{Tr}L_X=9\operatorname{tr}X$.
Conjugating $L_X$ by a Jordan automorphism gives $L_{gX}$, proving
$\operatorname{tr}(gX)=\operatorname{tr}X$. Product preservation now
gives $G(gX,gY)=G(X,Y)$.

Write $L_X(Y)=X\circ Y$. In any Jordan algebra $[L_X,L_Y]$ is a
derivation, with the following algebraic proof. The coefficient of
three distinct parameters in the Jordan identity gives
\begin{align*}
x\circ((z\circ w)\circ y)+z\circ((w\circ x)\circ y)
 +w\circ((x\circ z)\circ y)
={}&(z\circ w)\circ(y\circ x)+(w\circ x)\circ(y\circ z)\\
&+(x\circ z)\circ(y\circ w).
\end{align*}
Subtracting this identity with $x,y$ interchanged cancels the right
sides and gives exactly
\[
 [L_x,L_y](z\circ w)
 =([L_x,L_y]z)\circ w+z\circ([L_x,L_y]w).
\]
For cyclic $(i,j,k)$ define
\begin{equation}
 D_i(a)=4[L_{E_i(a)},L_{e_j}].
 \label{h24:att:eq:derivations}
\end{equation}
Its complete values follow from the Peirce product table:
\begin{align}
 D_i(a)e_i&=0,& D_i(a)e_j&=E_i(a),&
 D_i(a)e_k&=-E_i(a),\nonumber\\
 D_i(a)E_i(b)&=2h(a,b)(e_k-e_j),\label{h24:att:eq:D-table}\\
 D_i(a)E_j(b)&=-E_k(\overline{ab}),&
 D_i(a)E_k(c)&=E_j(\overline{ca}).\nonumber
\end{align}
For example on $e_j$ the commutator before multiplication by $4$ is
$E_i(a)/2-E_i(a)/4=E_i(a)/4$, and on $E_j(b)$ it is
$0-E_k(\overline{ab})/4$. The same product table supplies the other
entries. Exponentiating a derivation gives Jordan automorphisms,
since the derivation rule and the exponential power series give
$e^{tD}(X\circ Y)=(e^{tD}X)\circ(e^{tD}Y)$.

Let $\mathfrak f_4$ be the derivation algebra and $\mathfrak h$
its subalgebra fixing each $e_i$. Put
$\mathfrak m_i=\{D_i(a):a\in V_i\}$ and
$\mathfrak m=\bigoplus_i\mathfrak m_i$. Then
\begin{equation}
 \mathfrak f_4=\mathfrak h\oplus\mathfrak m,\qquad
 \mathfrak m_i\cong V_i,\qquad
 \operatorname{Ad}(h)D_i(a)=D_i(h_i a).
 \label{h24:att:eq:reductive}
\end{equation}
To check that no tangent directions are missing, differentiate
$e_i^2=e_i$ to get $De_i=2e_i\circ De_i$. Thus $De_i$ has only the
components $J_j\oplus J_k$. Differentiating $I^2=I$ gives $DI=0$,
so the two possible occurrences of any $J_i$ component in the triple
$(De_1,De_2,De_3)$ are opposite. Equation \eqref{h24:att:eq:D-table}
therefore supplies unique $a_i$ such that $D-\sum_iD_i(a_i)$ fixes
all three $e_i$. That equation also proves the sum direct.
Conjugation of the commutator defining $D_i$ by an automorphism
fixing $e_j$ proves the displayed equivariance.

Retain arbitrary distinct original spectral parameters
$\lambda_1,\lambda_2,\lambda_3\in\mathbb R$, with
$\lambda_1+\lambda_2+\lambda_3=0$, and set
\[
 d_\lambda=\sum_i\lambda_i e_i,\qquad
 \Delta_1=\lambda_2-\lambda_3,\quad
 \Delta_2=\lambda_3-\lambda_1,\quad
 \Delta_3=\lambda_1-\lambda_2.
\]
Every $\Delta_i$ is nonzero and
\begin{equation}
 D_i(a)d_\lambda=\Delta_i E_i(a).
 \label{h24:att:eq:spectral-differential}
\end{equation}
The equivariant map
$\Xi_\lambda:M\to F_4d_\lambda$, $[g]\mapsto gd_\lambda$,
is a diffeomorphism. Its inverse recovers the ordered frame from the
Jordan polynomials
\[
 p_i(z)=\frac{(z-\lambda_jI)\circ(z-\lambda_kI)}
              {(\lambda_i-\lambda_j)(\lambda_i-\lambda_k)}.
\]
For $z=gd_\lambda$ these equal $ge_i$: evaluate on $d_\lambda$ first,
then apply the Jordan automorphism $g$. In particular the stabilizer
of $d_\lambda$ is exactly $H_\tau$, and both maps are smooth.
Equation \eqref{h24:att:eq:spectral-differential} is the differential of
this comparison with all scale factors retained.

Use the associated-bundle convention
$(g,v)\sim(gh,h^{-1}v)$ for $F_4\times_{H_\tau}V_i$, and denote this
rank-eight real vector bundle by $\mathcal V_i$. Define
\begin{align}
 K_i:\mathcal V_i&\longrightarrow TM,&
 K_i[g,a]&=\left.\frac{d}{dt}\right|_0
                   g\exp(tD_i(a))H_\tau,\label{h24:att:eq:Kmap}\\
 P_i:\mathcal V_i&\longrightarrow M\times J,&
 P_i[g,a]&=(gH_\tau,gE_i(a)).\nonumber
\end{align}
Both maps are well defined by the proved equivariance.
Equation \eqref{h24:att:eq:reductive} shows that
$\bigoplus_iK_i:\bigoplus_i\mathcal V_i\to TM$ is an isomorphism.
The image of $P_i$ above the frame $(p_1,p_2,p_3)=(ge_1,ge_2,ge_3)$
is precisely the intrinsic Peirce space
\[
 \{X\in J:p_i\circ X=0,\quad p_j\circ X=p_k\circ X=X/2\}.
\]
Suppressing the recorded base-point component of $P_i$ when viewing
it as a tangent vector to the orbit, the exact comparison is
\begin{equation}
 d\Xi_\lambda\circ K_i=\Delta_iP_i.
 \label{h24:att:eq:bundle-comparison}
\end{equation}
Thus the original Albert slots and frame tangent slots are related
by an explicit map, rather than an unrecorded change of coordinates.

The tensor and products descend by the formulas
\begin{align}
 \boldsymbol\tau([g,a],[g,b],[g,c])&=\tau(a,b,c),\nonumber\\
 \boldsymbol\mu_{ij}([g,a],[g,b])&=[g,\mu_{ij}(a,b)]
                                  \in\mathcal V_k.
 \label{h24:att:eq:descent}
\end{align}
Representatives over one point can always be chosen with common $g$.
Replacing it by $gh$ replaces the arguments by
$h_i^{-1}a,h_j^{-1}b,h_k^{-1}c$.
Tensor invariance and \eqref{h24:att:eq:equivariance} give exactly the
same scalar and the same associated class. This proves well-defined
global descent, beyond invariance at a single fibre. The associated
metrics induced by $h$ carry all norm, adjoint, and inverse formulas
of Proposition~\ref{h24:att:prop:norm-inverse} to identities of bundle maps.

Transport by $K_i$ gives the three subbundles
$\mathcal T_i=K_i(\mathcal V_i)$, the direct sum
$TM=\mathcal T_1\oplus\mathcal T_2\oplus\mathcal T_3$, and global
bilinear maps $\mathcal T_i\otimes\mathcal T_j\to\mathcal T_k$.
Define a second, orbit-coordinate realization explicitly by
$S_i=K_i/\Delta_i$. Then $d\Xi_\lambda S_i=P_i$.
Transport of $\boldsymbol\mu$ by $S_i$ is exactly the original
coordinate map $E_i(a),E_j(b)\mapsto E_k(\overline{ab})$.
These transports obey the following exact comparison. For
$u_i\in J_i,u_j\in J_j$ at the reference orbit point, the
$K$-transported product has orbit-coordinate value
\begin{equation}
 \frac{\Delta_k}{\Delta_i\Delta_j}\,B_{ij}(u_i,u_j),
 \label{h24:att:eq:tangent-scale-products}
\end{equation}
whereas the $S$-transported product has value $B_{ij}(u_i,u_j)$.
Indeed the $K$-coordinate arguments are
$E_i^{-1}u_i/\Delta_i$ and $E_j^{-1}u_j/\Delta_j$; bilinearity and
the output multiplier $\Delta_k$ give the displayed coefficient.
At every other point the same identity is transported by $g\in F_4$.
The $K$-transported scalar in orbit coordinates is
\[
 \frac{\tau(E_1^{-1}u_1,E_2^{-1}u_2,E_3^{-1}u_3)}
             {\Delta_1\Delta_2\Delta_3},
\]
and the $S$-transported scalar is its numerator.
For the pullback $G_\lambda=\Xi_\lambda^*G$ of the original trace
metric on the orbit, trace invariance and \eqref{h24:att:eq:metrics} give
\[
 G_\lambda(K_i[g,a],K_i[g,b])=2\Delta_i^2h(a,b),\qquad
 G_\lambda(S_i[g,a],S_i[g,b])=2h(a,b).
\]
No $\Delta_i$ is set to $1$ and no coordinate is conjugated to absorb
a sign or a metric factor.

There is a specified connection for which these descended tensors
are parallel. On $F_4\to F_4/H_\tau$ take
\[
 \omega=\operatorname{pr}_{\mathfrak h}(g^{-1}dg)
 \quad\text{for the decomposition }\eqref{h24:att:eq:reductive}.
\]
Because the projection commutes with $\operatorname{Ad}(H_\tau)$,
this form is right equivariant and sends each fundamental vertical
vector to its generating element of $\mathfrak h$. These are exactly
the principal-connection conditions. Upstairs, the descended tensor
and products are represented by the constant $H_\tau$-invariant
tensors $\tau,\mu_{ij}$. Their covariant derivatives are the
ordinary derivatives of those constant coefficients plus the
infinitesimal $\mathfrak h$ action. Both terms vanish.
They are therefore parallel for this explicitly defined homogeneous
connection and for the connections transported to $TM$ by $K$ or $S$.
No assertion about Levi--Civita parallelism follows from this argument.

\paragraph{Source conventions and scope}

The matrix convention, related-triple equation, and Jordan frame
action agree entry for entry with Ichiro Yokota,
\emph{Exceptional Lie groups}, arXiv:0902.0431,
Theorem~1.14.2, Lemmas~1.14.3--1.14.4, Theorem~1.16.2, and
Theorem~2.7.1 \cite{h24:Yokota2009}. The present proof supplies the coordinate calculations,
norm and inverse maps, determinant coefficients, and global bundle
maps in that convention. Its algebraic starting objects are the
real alternative composition algebra $\mathbb O$ and its Albert
Jordan algebra $J$; those defining structures are retained.

The proved result is an ordered three-sector tangent structure on
$F_4/\operatorname{Spin}(8)$ with a normed triality tensor and three
exact bilinear morphisms. It does not assert toroidal filling,
an integrable complex structure, or physical dynamics.
The obstruction to equivariant unary identifications and the
positive bilinear relations between sectors were both proved.
The dependence on the diagonal-orbit parameters is the exact
morphism \eqref{h24:att:eq:bundle-comparison}, with its product and metric
consequences; it is not a missing comparison.


\subsubsection{A coordinate-preserving triality map onto the determinant boundary}
\label{h24:tri:sec:spectral-graph}

Retain the ordered matrix in \eqref{h24:eq:octonionic-equal-diagonal} and write
\[
 E(q)=Q_{\mathbb O}(0;q_1,q_2,q_3),\qquad
 s(q)=|q_1|^2+|q_2|^2+|q_3|^2,\qquad
 \tau(q)=\operatorname{Re}((q_1q_2)q_3).
\]
Here the eigenvalues are the three real Jordan eigenvalues, with their
multiplicities. They are the roots of
\(z^3-s(q)z-2\tau(q)\). Define
\[
 a(q)=-\lambda_{\min}(E(q)).
\]
Thus \(a(q)\) is exactly the largest real root of
\(\lambda^3-s(q)\lambda+2\tau(q)=0\). This prescription keeps the
complete cubic, including its mixed term and its factor two.

\begin{theorem}[The spectral graph and its exact inverse]
\label{h24:tri:thm:spectral-graph}
The map
\begin{equation}
 \mathcal B:\mathbb O^3\longrightarrow\partial\mathcal C(\mathbb O),
 \qquad q\longmapsto Q_{\mathbb O}(a(q);q)
 \label{h24:tri:eq:spectral-graph}
\end{equation}
is a proper homeomorphism. Its inverse is the coordinate projection
\[
 \mathcal P\left(
 \begin{pmatrix}
 \lambda&q_3&\overline q_2\\
 \overline q_3&\lambda&q_1\\
 q_2&\overline q_1&\lambda
 \end{pmatrix}\right)=(q_1,q_2,q_3).
\]
It extends to a based homeomorphism
\[
 \mathcal B^+:(\mathbb O^3)^+
       \xrightarrow{\ \cong\ }S_{\det,\mathbb O}^{24}.
\]
Let \(K=\operatorname{Spin}(8)\subset F_4\) be the pointwise stabilizer
of the ordered diagonal Jordan frame. On \(\mathbb O^3\), use the
three representations defined by
\(kE_i(q_i)=E_i(\rho_i(k)q_i)\), with \(E_i\) the three ordered
off-diagonal inclusions. Then \(\mathcal B\) and \(\mathcal P\) are
\(K\)-equivariant. The norm and the cubic interaction are recovered
on the boundary without discarding any coordinate:
\begin{equation}
 2\tau(q)=a(q)s(q)-a(q)^3.
 \label{h24:tri:eq:boundary-triality}
\end{equation}
\end{theorem}

\begin{proof}
Use the Euclidean Jordan trace inner product
\(\langle u,v\rangle_J=\operatorname{tr}(u\circ v)\).
For \(q,p\in\mathbb O^3\), direct multiplication gives
\[
 \|E(q)-E(p)\|_J^2=2\sum_{i=1}^3|q_i-p_i|^2.
\]
Let \(\mathcal I\) be the compact set of primitive idempotents of
trace one. The spectral theorem for the Euclidean Albert algebra gives
\[
 \lambda_{\min}(u)=\min_{e\in\mathcal I}\langle u,e\rangle_J.
\]
For completeness, in a spectral Jordan frame write
\(u=\sum_i\mu_i e_i\). For every \(e\in\mathcal I\), the numbers
\(\langle e_i,e\rangle_J\) are nonnegative and have sum one: positivity
uses self-duality of the cone of squares, and the sum is
\(\langle I_3,e\rangle_J=1\). Thus
\(\langle u,e\rangle_J\geq\min_i\mu_i\), with equality at an
idempotent for a least eigenvalue. This proves the displayed formula.
Every such \(e\) has \(\|e\|_J^2=\operatorname{tr}(e^2)=1\).
Cauchy--Schwarz and the minimum formula now give
\begin{equation}
 |a(q)-a(p)|\leq\|E(q)-E(p)\|_J
       =\sqrt{2\sum_{i=1}^3|q_i-p_i|^2}.
 \label{h24:tri:eq:spectral-Lipschitz}
\end{equation}
In particular \(a\) is continuous.

Since \(E(q)\) has trace zero, its least eigenvalue is nonpositive.
If \(q\ne0\), all eigenvalues cannot be zero, since the spectral theorem
would then give \(E(q)=0\). Their zero sum consequently forces a
strictly negative least eigenvalue. Hence \(a(0)=0\) and
\(a(q)>0\) for \(q\ne0\).
Adding \(\lambda I_3\) adds \(\lambda\) to each of the three Jordan
eigenvalues. Therefore
\[
 Q_{\mathbb O}(\lambda;q)\succeq0
       \quad\Longleftrightarrow\quad\lambda\geq a(q),
\]
and within this positive cone the determinant vanishes exactly at
\(\lambda=a(q)\). This is the relative boundary established in
\cref{h24:thm:octonionic-determinant-S24}. It proves both inverse
compositions in \eqref{h24:tri:eq:spectral-graph}.
Continuity of the graph and coordinate projection proves the
homeomorphism. Moreover
\[
 \|Q_{\mathbb O}(a(q);q)\|_J^2=3a(q)^2+2s(q)\geq2s(q),
\]
so unbounded \(q\) leaves every compact subset of the boundary.
This proves properness directly and the extension at infinity.

A member of \(K\) fixes \(I_3\), preserves the Jordan product and trace,
and sends each \(E_i(q_i)\) to \(E_i(\rho_i(k)q_i)\). It preserves
the characteristic polynomial and the ordered eigenvalues. Thus
\(a(\rho_1(k)q_1,\rho_2(k)q_2,\rho_3(k)q_3)=a(q)\), which proves
both equivariance identities. Substitution of \(\lambda=a(q)\) into
the unchanged determinant formula gives
\eqref{h24:tri:eq:boundary-triality}.
\end{proof}

\paragraph{The exact tangent map from the ordered flag.}
For pairwise different real \(\lambda_1,\lambda_2,\lambda_3\) with
zero sum, retain \(\Xi_\lambda\) of
\eqref{h24:eq:flag-diagonal-orbit-map}. At \(x_0=(d_1,d_2)\), the inverse
of its differential, written in the original off-diagonal coordinates, is
\begin{equation}
 \begin{aligned}
 A(q)&=E_2\left(\frac{q_2}{\lambda_1-\lambda_3}\right)
       +E_3\left(\frac{q_3}{\lambda_1-\lambda_2}\right),\\
 B(q)&=E_1\left(\frac{q_1}{\lambda_2-\lambda_3}\right)
       -E_3\left(\frac{q_3}{\lambda_1-\lambda_2}\right),\\
 (d\Xi_\lambda)_{x_0}(A(q),B(q))&=E(q).
 \end{aligned}
 \label{h24:tri:eq:flag-tangent-inverse}
\end{equation}
Indeed, differentiate the three idempotents in a Jordan frame. The
idempotence equations and the identity that their sum is \(I_3\)
force \(A\) to have only sectors two and three, \(B\) to have only
sectors one and three, and their sector-three parts to be negatives.
To check that these are actual tangent vectors, the differential of the
already proved diffeomorphism \(\Xi_\lambda\) is an isomorphism onto
\(E_1(\mathbb O)\oplus E_2(\mathbb O)\oplus E_3(\mathbb O)\).
Its formula is
\((\lambda_1-\lambda_3)A+(\lambda_2-\lambda_3)B\).
Solving its three sector equations gives precisely
\eqref{h24:tri:eq:flag-tangent-inverse}, so no extra tangent solutions or
denominators are suppressed. Composing this tangent coordinate
isomorphism with \(\mathcal B\) gives an explicit equivariant
homeomorphism from the complete tangent space at \(x_0\) onto
\(\partial\mathcal C(\mathbb O)\).


% Staged proof fragment. No canonical file has been edited.
% Parent integration prefix: fc:
\subsubsection{An explicit equivariant flag cell and its degree-one collapse}
\label{h24:fc:section}

Write $J=\mathfrak h_3(\mathbb O)$, $J_0=\ker\operatorname{tr}$,
$e=I_3$, and $d_i=\operatorname{diag}(\delta_{i1},\delta_{i2},\delta_{i3})$.
The inner product on $J$ is
\[
 \langle x,y\rangle_J=\operatorname{tr}(x\circ y),\qquad
 L_x(y)=x\circ y.
\]
All operator products and brackets below belong to
$\operatorname{End}_{\mathbb R}(J)$; thus $[A,B]=AB-BA$ is a commutator
of real linear operators. An adjoint $A^*$ in these formulas is the
adjoint for $\langle\ ,\ \rangle_J$.

We preserve the ordered coordinates of the equal-diagonal slice by putting
\begin{equation}
 \begin{aligned}
 E_1(r)&=\begin{pmatrix}0&0&0\\0&0&r\\0&\bar r&0\end{pmatrix},\qquad
 E_2(s)=\begin{pmatrix}0&0&\bar s\\0&0&0\\s&0&0\end{pmatrix},\\
 E_3(t)&=\begin{pmatrix}0&t&0\\\bar t&0&0\\0&0&0\end{pmatrix},\qquad
 W=\mathbb O\oplus\mathbb O\oplus\mathbb O,\\
 E(q)&=E_1(q_1)+E_2(q_2)+E_3(q_3)
      =Q_{\mathbb O}(0;q_1,q_2,q_3).
 \end{aligned}
 \label{h24:fc:coordinate-operators}
\end{equation}
In particular the second coordinate is the lower $(3,1)$ entry, and
conjugating it in the upper $(1,3)$ entry remains part of every formula.
Let $K=\operatorname{Aut}(J)=F_4$ and
$M=\{m\in K:m d_i=d_i\ (i=1,2,3)\}=\operatorname{Spin}(8)$.
The action of $M$ on $W$ is defined, without choosing three unrelated
orthogonal actions, by
\begin{equation}
 E_i((m\cdot q)_i)=mE_i(q_i),\qquad i=1,2,3.
 \label{h24:fc:triality-action}
\end{equation}
The three summands are the vector and the two real half-spin modules in
this specified coordinate convention.

We use the real symmetric pair
\[
 G=E_{6(-26)},\qquad
 \mathfrak g=\operatorname{Der}(J)\oplus L_{J_0},\qquad
 \mathfrak k=\operatorname{Der}(J),\qquad
 \mathfrak s=L_{J_0}.
\]
Here $G$ is the connected reduced structure group in its real linear
representation on $J$. Its Cartan involution on the Lie algebra is
$\theta(A)=-A^*$; hence
$\theta(D+L_x)=D-L_x$ for $D\in\operatorname{Der}(J)$ and $x\in J_0$.
The symmetric-pair realization and restricted-root spaces are as in
\cite[Section~2.6, pp.~491--493]{h24:MareWillems2013}.
The real-flag theorem we use is precisely
\cite[Appendix~B, Theorem~B.2, Proposition~B.3 and Corollary~B.4,
p.~519]{h24:MareWillems2013}: for $P=MAN$ the space $G/P=K/M$
has cells $Nw$, and $X\mapsto\exp(X)wP$ is a diffeomorphism from
the sum of the positive root spaces made negative by $w^{-1}$ onto
that cell. Its Cartan comparison is $\Theta(X)=X-\theta X$.
We now determine that comparison, the exponential, its flag coordinates,
its inverse, and the collapse with every coefficient specified.

\begin{lemma}[The Cartan inverse with its exact coefficient]
\label{h24:fc:cartan-inverse}
For $h=\operatorname{diag}(h_1,h_2,h_3)\in J_0$, set
\[
 \rho_1(h)=\frac{h_3-h_2}{2},\qquad
 \rho_2(h)=\frac{h_1-h_3}{2},\qquad
 \rho_3(h)=\frac{h_1-h_2}{2}=\rho_1(h)+\rho_2(h).
\]
Fix any
$h_*=\operatorname{diag}(h_{*,1},h_{*,2},h_{*,3})\in J_0$
with $h_{*,1}>h_{*,3}>h_{*,2}$. For $u\in E_i(\mathbb O)$ define
\begin{equation}
 \nu_i(u)=\frac12\left(L_u+
             \frac{[L_{h_*},L_u]}{\rho_i(h_*)}\right).
 \label{h24:fc:nu}
\end{equation}
Then $\nu_i(u)$ is independent of the permitted choice of $h_*$,
belongs to $\mathfrak g_{\rho_i}$, and satisfies
\begin{equation}
 [L_h,\nu_i(u)]=\rho_i(h)\nu_i(u),\qquad
 \nu_i(u)-\theta\nu_i(u)=L_u,\qquad
 \nu_i(u)e=\frac12u.
 \label{h24:fc:nu-identities}
\end{equation}
Thus $\nu_i$ is exactly the inverse of $\Theta$ on the $i$th summand;
the factor $1/2$ in \eqref{h24:fc:nu} is necessary.
The isomorphism
\[
 Z:W\longrightarrow\mathfrak n
       =\mathfrak g_{\rho_1}\oplus\mathfrak g_{\rho_2}
                               \oplus\mathfrak g_{\rho_3},\qquad
 Z(q)=\sum_{i=1}^3\nu_i(E_i(q_i))
\]
is $M$-equivariant.
\end{lemma}

\begin{proof}
We first record the Jordan identity needed to ensure that the operators
are in the correct Lie algebra. For any $a,b\in J$,
$D_{a,b}=[L_a,L_b]$ is a derivation, and therefore
\begin{equation}
 [D_{a,b},L_c]=L_{D_{a,b}c}.
 \label{h24:fc:inner-derivation}
\end{equation}
Here is the polarization proving this assertion. The coefficient of the
product of three distinct real parameters in
$(x^2\circ y)\circ x=x^2\circ(y\circ x)$ gives
\[
 \begin{split}
 &d\circ\bigl(b\circ(a\circ c)\bigr)
 +a\circ\bigl(b\circ(c\circ d)\bigr)
 +c\circ\bigl(b\circ(d\circ a)\bigr)\\
 &\hspace{1em}=(a\circ c)\circ(b\circ d)
 +(c\circ d)\circ(b\circ a)
 +(d\circ a)\circ(b\circ c).
 \end{split}
\]
Subtract the same identity with $a$ and $b$ exchanged. The three
terms on the right cancel, by commutativity of the Jordan product.
The result, rearranged, is
\[
 D_{a,b}(c\circ d)=(D_{a,b}c)\circ d+c\circ(D_{a,b}d),
\]
which proves the assertion and \eqref{h24:fc:inner-derivation}.

Let $u$ have its two off-diagonal entries in positions $(j,k)$ and
$(k,j)$, and write $D_h=[L_h,L_u]$ for real diagonal $h$.
If $h'$ is another real diagonal element, direct use of the binary
Jordan product gives
\begin{equation}
 \begin{split}
 D_hh'
 &=\left(\frac{(h_j+h_k)(h'_j+h'_k)}4
                 -\frac{h_jh'_j+h_kh'_k}2\right)u\\
 &=-\frac{(h_j-h_k)(h'_j-h'_k)}4u
  =-\rho_i(h)\rho_i(h')u.
 \end{split}
 \label{h24:fc:diagonal-calculation}
\end{equation}
For completeness, the proportionality of $D_h$ in $h$ can also be
checked on every $z\in J$ with no associativity assumption on three
octonions. Expanding the two Jordan products, the fact that every
entry of $h$ is real gives the restricted identity
\[
 4[L_h,L_u]z=(hu-uh)z-z(hu-uh).
\]
In this displayed identity only, the products are binary matrix
products; the expansion is valid because every reassociation it uses
contains a real entry from $h$. Since $u$ has only the indicated two
off-diagonal entries, $hu-uh$ depends linearly on the single real
scalar $h_j-h_k$. Consequently
\begin{equation}
 [L_h,L_u]=\frac{\rho_i(h)}{\rho_i(h_*)}[L_{h_*},L_u].
 \label{h24:fc:proportionality}
\end{equation}
This assertion includes $\rho_i(h)=0$. We use no corresponding
identity with arbitrary octonionic $h$, and no octonionic matrix
conjugation is used to define a group action.

Put $D_*=[L_{h_*},L_u]$. Equations
\eqref{h24:fc:inner-derivation} and
\eqref{h24:fc:diagonal-calculation} give
\[
 [L_h,D_*]=-L_{D_*h}=\rho_i(h_*)\rho_i(h)L_u.
\]
Combining this with \eqref{h24:fc:proportionality} proves
$[L_h,\nu_i(u)]=\rho_i(h)\nu_i(u)$ and also proves independence of
$h_*$. Since $\theta L_u=-L_u$ and $\theta D_*=D_*$, subtracting
$\theta\nu_i(u)$ from \eqref{h24:fc:nu} gives $L_u$ exactly.
Every derivation kills the unit, so $\nu_i(u)e=u/2$.
The map $u\mapsto L_u$ is injective because $L_ue=u$.
The positive restricted-root spaces have the three stated
eight-dimensional Cartan images; thus these injective maps exhaust
them. Finally, an $m\in M$ fixes $h_*$ and satisfies
$mL_um^{-1}=L_{mu}$ and
$m[L_{h_*},L_u]m^{-1}=[L_{h_*},L_{mu}]$.
These formulas prove the asserted equivariance.
\end{proof}

\begin{lemma}[Finite exponential and the coefficient in the nilpotent law]
\label{h24:fc:finite-exponential}
For every $q\in W$,
\begin{equation}
 Z(q)^5=0,\qquad
 \exp Z(q)=I+Z(q)+\frac{Z(q)^2}{2}
                   +\frac{Z(q)^3}{6}+\frac{Z(q)^4}{24}.
 \label{h24:fc:finite-exp}
\end{equation}
The same formula holds with $Z(q)$ replaced by $-Z(q)^*$.
Writing $\nu_i(r)$ temporarily for $\nu_i(E_i(r))$, the brackets are
\begin{equation}
 [\nu_1(r),\nu_2(s)]=\nu_3\left(-\frac12\overline{rs}\right),
 \qquad [\nu_i(r),\nu_i(s)]=0,
 \qquad[\nu_3(t),\nu_j(r)]=0.
 \label{h24:fc:nil-bracket}
\end{equation}
The last equality holds for $j=1,2,3$. In exponential coordinates the
group law on $N$ is
\begin{equation}
 \begin{split}
 (q_1,q_2,q_3)\star(q'_1,q'_2,q'_3)
 =\biggl(&q_1+q'_1,\ q_2+q'_2,\\
         &q_3+q'_3-\frac14\overline{q_1q'_2}
                         +\frac14\overline{q'_1q_2}\biggr).
 \end{split}
 \label{h24:fc:nil-group-law}
\end{equation}
\end{lemma}

\begin{proof}
The auxiliary operator $H_0=L_{\operatorname{diag}(2,-2,0)}$ has
the following eigenspaces in its representation on $J$:
\[
 \begin{array}{c|ccccc}
 \text{eigenvalue}&-2&-1&0&1&2\\ \hline
 \text{eigenspace}&\mathbb R d_2&E_1(\mathbb O)&
 \mathbb R d_3\oplus E_3(\mathbb O)&E_2(\mathbb O)&\mathbb R d_1.
 \end{array}
\]
This table follows from
$L_h d_j=h_jd_j$ and
$L_hu=(h_j+h_k)u/2$ for an off-diagonal $(j,k)$ element.
The actual positive roots evaluated on $\operatorname{diag}(2,-2,0)$
are $1,1,2$. Therefore every term of $Z(q)$ strictly raises these
weights by at least one. Five successive applications give zero.
Adjoints strictly lower the weights, proving the analogous assertion
for $-Z(q)^*$. This proves both finite exponential formulas.

Brackets add restricted roots. There is no root $2\rho_i$,
$\rho_1+\rho_3$, or $\rho_2+\rho_3$, so the brackets claimed to
vanish in \eqref{h24:fc:nil-bracket} do vanish. To compute the
remaining coefficient, put $u=E_1(r)$, $v=E_2(s)$. Binary matrix
multiplication gives, in the retained coordinate order,
\[
 u\circ v=\frac12E_3(\overline{rs})=:w.
\]
The weight of $u$, $v$, and $w$ under $L_h$ is respectively
$-h_1/2$, $-h_2/2$, and $-h_3/2$. Hence
\[
 [L_h,L_u]v=-\rho_1(h)w,\qquad
 [L_h,L_v]u=\rho_2(h)w,
\]
and \eqref{h24:fc:nu} gives
$\nu_1(r)v=0$ and $\nu_2(s)u=w$. Applying the commutator to $e$,
\[
 [\nu_1(r),\nu_2(s)]e
       =\frac12(\nu_1(r)v-\nu_2(s)u)
       =-\frac14E_3(\overline{rs}).
\]
On the third root space the map $A\mapsto2Ae$ is inverse to $\nu_3$.
This proves the first formula of \eqref{h24:fc:nil-bracket}.
All iterated brackets of length three vanish. The exact
Baker--Campbell--Hausdorff formula is therefore
$\log(\exp X\exp Y)=X+Y+[X,Y]/2$; substituting the just-computed
bracket gives \eqref{h24:fc:nil-group-law}, including both signs and
the coefficient $1/4$.
\end{proof}

\begin{theorem}[A finite operator formula for the entire open cell]
\label{h24:fc:open-cell}
For $q\in W$ define
\begin{equation}
 \begin{aligned}
 A(q)&=\left(I+Z+\frac{Z^2}{2}+\frac{Z^3}{6}
                       +\frac{Z^4}{24}\right)d_2,\\
 B(q)&=\left(I-Z^*+\frac{(Z^*)^2}{2}-\frac{(Z^*)^3}{6}
                       +\frac{(Z^*)^4}{24}\right)d_1,
       \qquad Z=Z(q),\\
 \Psi(q)&=\left(\frac{A(q)}{\operatorname{tr}A(q)},
               \frac{B(q)}{\operatorname{tr}B(q)}\right).
 \end{aligned}
 \label{h24:fc:explicit-chart}
\end{equation}
Both traces in these formulas are strictly positive. The map $\Psi$
is an $M$-equivariant diffeomorphism
\begin{equation}
 W\xrightarrow{\ \Psi\ }
 U=\{(a,b)\in\operatorname{Fl}(\mathbb O):a_{22}>0, b_{11}>0\}
   =C_{w_0},\qquad w_0=(12).
 \label{h24:fc:open-cell-locus}
\end{equation}
The inverse $q=\Psi^{-1}(a,b)$ is the exact formula
\begin{equation}
 \boxed{\quad
 q_1=\frac{2a_{23}}{a_{22}},\qquad
 q_2=-\frac{2b_{31}}{b_{11}},\qquad
 q_3=\frac{2a_{12}}{a_{22}}-\frac14\overline{q_1q_2}.
 \quad}
 \label{h24:fc:explicit-inverse}
\end{equation}
Every denominator in \eqref{h24:fc:explicit-inverse} is a positive
real scalar. In particular $\Psi(0)=(d_2,d_1)$.
\end{theorem}

\begin{proof}
Let $\widehat w_0$ be the Jordan automorphism that interchanges
the first and second real diagonal positions, obtained by conjugation
with the real permutation matrix for $(12)$. For this real matrix
the conjugation is a well-defined Jordan automorphism. Its action on
the retained positive roots is
\[
 w_0\rho_1=-\rho_2,\qquad
 w_0\rho_2=-\rho_1,\qquad
 w_0\rho_3=-\rho_3.
\]
It is consequently the longest Weyl element, and the real-flag
cell theorem gives a diffeomorphism
\begin{equation}
 W\longrightarrow N\widehat w_0P/P,\qquad
 q\longmapsto\exp Z(q)\widehat w_0P.
 \label{h24:fc:group-chart}
\end{equation}

Here are the exact Iwasawa factors that identify the flag pair in
\eqref{h24:fc:group-chart}. Write $g=kan$ uniquely with
$k\in K$, $a=\exp L_h\in A$, $n\in N$. Each positive root
operator kills $d_1$. This follows either by applying
\eqref{h24:fc:nu} to $d_1$, or from the $H_0$ weight table,
whose highest weight is $\mathbb R d_1$. Each operator
$-\nu_i^*$ kills $d_2$ by the lowest-weight statement. Thus
\[
 nd_1=d_1,\qquad (n^*)^{-1}d_2=d_2,
\]
and consequently
\begin{equation}
 gd_1=e^{h_1}kd_1,\qquad
 (g^*)^{-1}d_2=e^{-h_2}kd_2.
 \label{h24:fc:iwasawa-flag-factors}
\end{equation}
Since Jordan automorphisms preserve the trace,
the traces on the right are exactly $e^{h_1}>0$ and $e^{-h_2}>0$.
It follows that the flag corresponding to $gP$ is
\[
 \left(\frac{gd_1}{\operatorname{tr}(gd_1)},
       \frac{(g^*)^{-1}d_2}{\operatorname{tr}((g^*)^{-1}d_2)}\right)
       =(kd_1,kd_2).
\]
For $g=\exp Z(q)\widehat w_0$, its two numerators are
$\exp Z(q)d_2$ and $\exp(-Z(q)^*)d_1$, respectively.
The finite exponential identities prove
\eqref{h24:fc:explicit-chart}, all its positivity assertions,
and the identification with \eqref{h24:fc:group-chart}.

We next verify the image description, so that no membership criterion
is left implicit. For any $n\in N$, the $H_0$ weight table shows
that the coefficient of $d_2$ in $nd_j$ is $1$ if $j=2$ and $0$
otherwise. Indeed $n$ strictly raises weights apart from its identity
term. Similarly the coefficient of $d_1$ in $(n^*)^{-1}d_j$ is
$1$ if $j=1$ and $0$ otherwise. A flag in $N\widehat wP/P$ is
represented by the positive-trace multiples of
$nd_{w(1)}$ and $(n^*)^{-1}d_{w(2)}$. It satisfies
$a_{22}>0$ precisely when $w(1)=2$, and $b_{11}>0$ precisely
when $w(2)=1$. Both conditions hold precisely for $w=w_0=(12)$.
The six-cell decomposition covers the flag, proving
\eqref{h24:fc:open-cell-locus}.

For the inverse we record the required coefficients of the finite
exponentials. The same weights show
\begin{equation}
 \begin{gathered}
 A(q)_{22}=1,\qquad B(q)_{11}=1,\qquad
 A(q)_{23}=\frac12q_1,\qquad
 B(q)_{31}=-\frac12q_2,\\
 A(q)_{12}=\frac12q_3+\frac18\overline{q_1q_2}.
 \label{h24:fc:coefficient-extraction}
 \end{gathered}
\end{equation}
To see also the last coefficient directly, the relevant first-order
terms are
\[
 Zd_2=\frac12E_1(q_1)+\frac12E_3(q_3),\qquad
 -Z^*d_1=-\frac12E_2(q_2)-\frac12E_3(q_3).
\]
An $E_3$ contribution to $Z^2d_2$ can arise only from
$\nu_2(q_2)\nu_1(q_1)d_2$ or
$\nu_1(q_1)\nu_2(q_2)d_2$; the latter is zero, and the former
is $\frac14E_3(\overline{q_1q_2})$ by the calculation in
Lemma~\ref{h24:fc:finite-exponential}. The other weight-two
contribution $\nu_1(q_1)^2d_2$ lies in $\mathbb R d_3$,
and $\nu_2(q_2)^2d_2=0$. The coefficient $1/2$ in the
exponential gives $1/8$. Terms of degree three or four raise
weight too far to have an $E_3$ component. The other four
identities in \eqref{h24:fc:coefficient-extraction} follow from
the unique lowest, highest, and adjacent weights.

For $(a,b)=\Psi(q)$, the first two identities imply
\[
 A(q)=\frac{a}{a_{22}},\qquad
 B(q)=\frac{b}{b_{11}}.
\]
The remaining identities solve for $q_1,q_2,q_3$ in precisely
the order \eqref{h24:fc:explicit-inverse}. Since
\eqref{h24:fc:group-chart} is onto $U$, this proves both inverse
compositions everywhere on $U$, not only a formal local inverse.
The displayed formulas show their smoothness.

Finally, $Z(mq)=mZ(q)m^{-1}$ and $m$ is orthogonal for the
trace inner product. It fixes $d_1,d_2$ and preserves the trace.
Each expression in \eqref{h24:fc:explicit-chart} therefore
commutes with the $M$ action. This proves equivariance of $\Psi$
and, by its proved bijectivity, of its inverse.
\end{proof}

\begin{proposition}[The orientation sign, including the conjugated coordinate]
\label{h24:fc:orientation}
Fix an orientation of $\mathbb O$, and give $W$ the ordered product
orientation $(q_1,q_2,q_3)$. Give the flag the orientation specified
in the integral Euler-class calculation: orient its first root
summand by the chosen orientation of $E_1(\mathbb O)$, its third
by $n_2E_1(\mathbb O)$, and its second by
$n_1^{-1}n_2E_1(\mathbb O)$, where $n_1=(23)$ and $n_2=(13)$
are real-permutation lifts. The total orientation is the ordered
sum of root summands $(1,2,3)$, transported by $K$.
Then $\Psi:W\to U$ is orientation preserving.
\end{proposition}

\begin{proof}
These are the orientation transport conventions of
\cite[Section~3, pp.~495--496]{h24:MareWillems2013}.
They are independent of the lifts within their cosets because
$M$ is connected and each of its orthogonal representations has
positive real determinant.
We compute the sign explicitly. For this computation only, call
the product orientation of the three upper entries $(r,q,p)$ in
positions $(23),(13),(12)$ the upper-entry orientation. Octonionic
conjugation fixes the real line and acts by $-1$ on the seven
imaginary directions, so its real determinant is $(-1)^7=-1$.
The lift $n_2=(13)$ takes an upper $(23)$ entry $r$ to the upper
$(12)$ entry $\bar r$. The lift $n_1^{-1}=(23)$ then takes that
upper $(12)$ entry to the upper $(13)$ entry with the same value.
The source orientations of summands $1,2,3$ thus have signs
$+1,-1,-1$ relative to the upper-entry coordinates. Their product
is $+1$, so the flag tangent orientation at $d_\lambda$ agrees
with the upper-entry orientation. Here $d_\lambda$ has the original
pairwise distinct parameters $\lambda_i$ and the trace-zero relation;
no ordering or rescaling of these parameters is needed.

The lift $\widehat w_0=(12)$ interchanges the upper-entry
$(23)$ and $(13)$ blocks and conjugates the $(12)$ block.
Interchanging two eight-dimensional blocks has determinant
$(-1)^{8\cdot8}=+1$, and conjugation has determinant $-1$.
Thus the transported flag orientation at
$\widehat w_0d_\lambda$ is the negative of its upper-entry
orientation. This transport preserves the flag orientation because
it is the action of an element of connected $K$.

In the retained $q$ coordinates, $E$ conjugates exactly the second
upper-entry block; its determinant relative to the upper-entry
orientation is therefore $-1$. Differentiate
\eqref{h24:fc:explicit-chart} and the original orbit comparison
$\Xi_\lambda(a,b)=\lambda_1a+\lambda_2b+\lambda_3(e-a-b)$
at $q=0$. The traces have zero first derivative, and the result is
\begin{equation}
 \begin{split}
 d(\Xi_\lambda\Psi)_0(q)
 ={}&\frac{\lambda_1-\lambda_3}{2}E_1(q_1)
    +\frac{\lambda_3-\lambda_2}{2}E_2(q_2)
    +\frac{\lambda_1-\lambda_2}{2}E_3(q_3)\\
 ={}&-\sum_{i=1}^3\rho_i(\widehat w_0d_\lambda)E_i(q_i).
 \end{split}
 \label{h24:fc:chart-differential}
\end{equation}
Its three real scalar factors are nonzero; each acts on eight
real coordinates, so their determinant contribution is positive.
The total differential has sign $-1$ relative to upper-entry
coordinates, exactly the sign of the transported flag orientation
at this point. Therefore its sign as an oriented map is $+1$.
The domain is connected and the map is a diffeomorphism, so this
sign is $+1$ everywhere.
\end{proof}

\begin{theorem}[The exact equivariant quotient onto the representation sphere]
\label{h24:fc:collapse}
Let $F=\operatorname{Fl}(\mathbb O)$ and put
\[
 A=F\setminus U=\{(a,b)\in F:a_{22}=0\ \text{or}\ b_{11}=0\}.
\]
It is the union of the five cells other than $C_{w_0}$, and is a
closed $16$-dimensional CW subcomplex. Define
\begin{equation}
 c:F\longrightarrow W^+,
 \qquad
 c(a,b)=
 \begin{cases}
   \Psi^{-1}(a,b),&(a,b)\in U,\\
   \infty,&(a,b)\in A.
 \end{cases}
 \label{h24:fc:collapse-map}
\end{equation}
Then $c$ is continuous, surjective, and $M$-equivariant and induces
a based $M$-equivariant homeomorphism
\begin{equation}
 F/A\xrightarrow{\ \cong\ }W^+.
 \label{h24:fc:quotient-homeomorphism}
\end{equation}
With the orientations in Proposition~\ref{h24:fc:orientation},
$c$ has degree $+1$. If $\omega_W\in\widetilde H^{24}(W^+;\mathbb Z)$
is its positive generator, then
\begin{equation}
 c_*[F]=[W^+],\qquad c^*\omega_W=u_{\mathbb O},
 \qquad u_{\mathbb O}=\beta_1^2\beta_2=\beta_1\beta_2^2.
 \label{h24:fc:degree-one}
\end{equation}
On integral homology $c_*$ is the identity in degree zero, an
isomorphism in degree $24$, and the zero map on the two copies
of $\mathbb Z$ in each of degrees $8$ and $16$.
On fixed points its exact formula is
\[
 c(w_0x_{\mathbb O,0})=0,\qquad
 c(wx_{\mathbb O,0})=\infty\quad(w\ne w_0).
\]
These are the six points of $F^M$ mapping onto the two points
$(W^+)^M=\{0,\infty\}$.
\end{theorem}

\begin{proof}
All diagonal entries of primitive positive idempotents are
nonnegative. The open-cell locus of
Theorem~\ref{h24:fc:open-cell} consequently gives the stated
closed-set formula for $A$. The restricted roots have multiplicity
eight. The six cell dimensions are $0,8,8,16,16,24$, with the
unique top cell $C_{w_0}$. The CW property of the real-flag
decomposition implies that their union below degree $24$ is the
closed $16$-skeleton, which proves the subcomplex assertion.

Continuity of $c$ on $U$ is already proved. Let $K_0\subset W$
be compact. Its image $\Psi(K_0)$ is compact in the Hausdorff
compact manifold $F$, and hence closed in $F$. Thus
\[
 c^{-1}\bigl(W^+\setminus K_0\bigr)=F\setminus\Psi(K_0)
\]
is open and contains $A$. Complements of compact subsets of $W$
form a neighborhood basis of $\infty$. This proves continuity
at every point of $A$, including all its lower-dimensional strata.
It also proves the exact escape statement: a net in $U$ approaching
$A$ eventually leaves $\Psi(K_0)$ for every compact $K_0\subset W$.
Surjectivity and the description of the fibers are immediate from
the two inverse compositions of the chart: every finite point has
one preimage, and $c^{-1}(\infty)=A$. Therefore the induced map
$F/A\to W^+$ is a continuous bijection from a compact space to a
Hausdorff space and hence a homeomorphism.
The open cell is $M$-invariant by the proved chart equivariance;
its complement is invariant too. The formula for $c$ is therefore
$M$-equivariant, with $M$ fixing $\infty$.

Since $A$ has dimension $16$, its integral homology in degrees
$23$ and $24$ is zero. The exact sequence of the pair gives an
isomorphism
\[
 H_{24}(F;\mathbb Z)\xrightarrow{\ \cong\ }H_{24}(F,A;\mathbb Z).
\]
Because $A$ is a subcomplex, quotienting identifies this relative
group with $\widetilde H_{24}(F/A;\mathbb Z)$.
At any point of $U$, the local map is the inverse of the
orientation-preserving chart. It has local degree $+1$.
There is exactly one preimage of that finite regular value.
This fixes the sign of the preceding homology isomorphism and
proves $c_*[F]=[W^+]$. Pairing with cohomology proves
$c^*\omega_W=u_{\mathbb O}$ for the already fixed positive
top generator of the flag. The remaining homology assertions
follow because $W^+$ is a $24$-sphere and $F$ is connected with
the stated six-cell homology.
Each of the three nontrivial irreducible eight-dimensional
$M$-modules has zero invariant subspace, so $W^M=\{0\}$.
The point $w_0x_{\mathbb O,0}$ is $\Psi(0)$ and belongs to the
open cell; the other five Weyl points lie in $A$. Substitution
in the defining formula for $c$ proves the fixed-point assertions.
\end{proof}

\begin{corollary}[Composition with the retained spectral boundary graph]
\label{h24:fc:boundary-composite}
For $q\in W$, let
\[
 \alpha(q)=-\lambda_{\min}(E(q)),\qquad
 \Gamma(q)=Q_{\mathbb O}(\alpha(q);q_1,q_2,q_3).
\]
The exact boundary graph has inverse the off-diagonal coordinate
projection $Q_{\mathbb O}(\lambda;q)\mapsto q$. It extends to an
$M$-equivariant based homeomorphism
\[
 \Gamma^+:W^+\longrightarrow
 (\partial\mathcal C(\mathbb O))^+,
 \qquad\Gamma^+(\infty)=\infty.
\]
Consequently the fully specified map
\begin{equation}
 C_{\det}=\Gamma^+c:F\longrightarrow
           S_{\det,\mathbb O}^{24}
 \label{h24:fc:determinant-collapse}
\end{equation}
is continuous, $M$-equivariant, and has degree $+1$ when the
boundary sphere carries the orientation transported from its
retained ordered coordinates by $\Gamma^+$. Its exact fibers are
\[
 C_{\det}^{-1}\bigl(\Gamma(q)\bigr)=\{\Psi(q)\}
 \quad(q\in W),\qquad C_{\det}^{-1}(\infty)=A.
\]
\end{corollary}

\begin{proof}
The eigenvalues of $\lambda e+E(q)$ are
$\lambda+\lambda_j(E(q))$, with their multiplicities retained.
It is positive semidefinite precisely when
$\lambda\ge-\lambda_{\min}(E(q))$, and lies on the relative
boundary precisely when equality holds. This proves both
compositions with the stated coordinate projection. The minimum
eigenvalue is continuous, so $\Gamma$ is continuous, and the
projection is continuous. A homeomorphism is proper, because the
inverse maps compact sets to compact sets; therefore it extends
to the indicated one-point compactifications.
An element of $M$ fixes $e$, preserves eigenvalues as a Jordan
automorphism, and sends $E(q)$ to $E(mq)$. Thus
$\alpha(mq)=\alpha(q)$ and $\Gamma(mq)=m\Gamma(q)$.
This proves equivariance also at infinity. The asserted
properties, degree, and fibers of $C_{\det}$ now follow from
Theorem~\ref{h24:fc:collapse} by composition with the explicitly
oriented homeomorphism $\Gamma^+$.
\end{proof}

\begin{corollary}[The exact limit on a reverse sphere map]
\label{h24:fc:no-reverse-degree}
Every continuous map $f:S^{24}\to F$ has degree zero.
In particular neither $c$ nor $C_{\det}$ admits a continuous
section or a homotopy section. These assertions coexist with
the degree-one equivariant maps just constructed.
\end{corollary}

\begin{proof}
The cohomology groups $H^8(S^{24};\mathbb Z)$ vanish, so
$f^*\beta_1=f^*\beta_2=0$. Hence
\[
 f^*u_{\mathbb O}=f^*(\beta_1^2\beta_2)
                 =(f^*\beta_1)^2f^*\beta_2=0.
\]
This is exactly the statement that the degree is zero.
A homotopy section of $c$ would satisfy
$f^*c^*\omega_W=\omega_W$, contradicting the equality just
proved and \eqref{h24:fc:degree-one}. The same argument applies
after the homeomorphism $\Gamma^+$.
\end{proof}

\begin{proposition}[Exact maps among the three marked 24-dimensional constructions]
\label[proposition]{h24:prop:three-24d-objects}
The following constructions are related by the explicit spectral graph,
flag-cell collapse, and smash-coordinate maps proved here.
\begin{enumerate}[label=\textup{(\roman*)}]
\item The octonionic flag is the homogeneous \(24\)-manifold
      \(F_4/\operatorname{Spin}(8)\), with
      \(H^8\cong H^{16}\cong\mathbb Z^2\) and \(\chi=6\).
\item After a based diffeomorphism \(X\cong S^6\), the fourth smash power
      \(X^{\wedge4}\) is a based sphere \(S_{\wedge}^{24}\); it has
      \(\widetilde H^k=0\) for \(k\ne24\) and \(\chi=2\), and its quotient
      maps depend on the chosen based candidate and cusp collapses.
\item The octonionic equal-diagonal cone has the determinant-boundary
      compactification
      \(S_{\det,\mathbb O}^{24}=(\partial\mathcal C(\mathbb O))^+\).  It is produced by
      the explicit cubic \eqref{h24:eq:octonionic-equal-diagonal-determinant},
      not by a smash product.
\end{enumerate}
The first package is not homotopy equivalent to either sphere.  The two
sphere carriers in (ii) and (iii) are related by the based homeomorphism
displayed below. The flag maps to the determinant sphere by the explicit
degree-one map of Corollary~\ref{h24:fc:boundary-composite}; its entire
infinite fibre is the stated $16$-skeleton. The obstruction to a reverse
homotopy section is Corollary~\ref{h24:fc:no-reverse-degree}.

The lower complex determinant orbit pair \((S^{10},S^8)\) from
\cref{h24:thm:sphere-ladder} is a lower rung: its \(S^8\) is the compactified
complex rank-loss locus and its complement has determinant-phase
fundamental group \(\mathbb Z\).  It is not a third \(24\)-dimensional object.
\end{proposition}

\begin{proof}
Item (i) is \cref{h24:thm:octonionic-triality,h24:thm:octonionic-cohomology}.
The based smash identity
\((S^6)^{\wedge4}\cong S^{24}\) proves the sphere and cohomology assertions
in (ii).  Item (iii) is
\cref{h24:thm:octonionic-determinant-S24}.  The nonzero degree-eight and
degree-sixteen cohomology of the flag precludes a homotopy equivalence with
either sphere; the Euler characteristics give an independent obstruction.
For an explicit choice-dependent comparison, let
\(\eta:(X,x_{\mathrm{base}})\to(\mathbb R^6)^+\) be a based homeomorphism,
write the four six-coordinate vectors consecutively as
$(x_0,\ldots,x_{23})$, and fix the explicit linear isometry
$L:(\mathbb R^6)^4\to\mathbb O^3$ given by
$q_j=\sum_{a=0}^7x_{8(j-1)+a}u_a$ in the fixed ordered orthonormal
octonion basis. Let $b=\mathcal P$ be the off-diagonal coordinate inverse
of the spectral graph in Theorem~\ref{h24:tri:thm:spectral-graph}.
The smash compactification map
\[
 \kappa:((\mathbb R^6)^+)^{\wedge4}\longrightarrow((\mathbb R^6)^4)^+
\]
sends a tuple of finite vectors to their ordered tuple and sends the
collapsed wedge to infinity.  It is a continuous bijection from a compact
space to a Hausdorff space, hence a homeomorphism.  Thus
\[
 (b^+)^{-1}\circ L^+\circ\kappa\circ\eta^{\wedge4}:
 X^{\wedge4}\xrightarrow{\ \cong\ }
       (\partial\mathcal C(\mathbb O))^+
\]
is a based homeomorphism with inverse obtained by reversing these four
maps.  It depends explicitly on \(\eta\) and \(L\); neither marked
construction specifies them.  The final
lower-rung assertions are \eqref{h24:eq:S10-S8-pair} and
\eqref{h24:eq:rank-complement}.
\end{proof}

\subsection{What the mod--12 and mod--24 arithmetic actually is}
\label{h24:sec:mod-12-24}

\subsubsection{The integral peripheral representation}

The matrices \(C_1,C_2',C_0\) in
\eqref{h24:eq:elliptic-peripheral-matrices} are not merely of orders three,
four, and infinite.  Since \(C_2'\) and \(C_0\) generate
\(SL_2(\mathbb Z)\), their abelianization detects one exact modulus.

\begin{theorem}[Peripheral modulus and the mod--24 boundary]
\label[theorem]{h24:thm:peripheral-moduli}
Let
\[
 \operatorname{ab}:SL_2(\mathbb Z)\twoheadrightarrow SL_2(\mathbb Z)_{\mathrm{ab}}
\]
be the abelianization map.  There is a unique isomorphism
\[
 \overline\nu:SL_2(\mathbb Z)_{\mathrm{ab}}
       \xrightarrow{\ \cong\ }\mathbb Z/12
\]
for which the marked character
\begin{equation}
 \chi=\overline\nu\circ\operatorname{ab}:
 SL_2(\mathbb Z)\longrightarrow\mathbb Z/12
 \label{h24:eq:marked-peripheral-character}
\end{equation}
sends \(C_0\) to \(-1\).  Its values on the three peripheral matrices are
\begin{equation}
 \chi(C_1)=4,\qquad \chi(C_2')=-3,\qquad \chi(C_0)=-1
 \pmod {12}.
 \label{h24:eq:peripheral-mod12-values}
\end{equation}
The residues \(4\) and \(-3=9\) are complementary idempotents in
\(\mathbb Z/12\); additively they have exact orders \(3\) and \(4\).

Let
\[
 j:\mathbb Z/12\hookrightarrow\mathbb Z/24,\qquad j(a)=2a.
\]
Every group homomorphism \(\psi:SL_2(\mathbb Z)\to\mathbb Z/24\) has image in the even
subgroup and factors uniquely through \(j\).  In particular
\begin{equation}
 \chi_{24}=j\chi,\qquad
 \bigl(\chi_{24}(C_1),\chi_{24}(C_2'),\chi_{24}(C_0)\bigr)
 =(8,-6,-2)\pmod {24}
 \label{h24:eq:peripheral-mod24-values}
\end{equation}
is the unique mod--\(24\) re-encoding with
\[
 \exp\!\left(\frac{2\pi i\,\chi_{24}(M)}{24}\right)
 =
 \exp\!\left(\frac{2\pi i\,\chi(M)}{12}\right)
 \qquad(M\in SL_2(\mathbb Z)).
\]
Its image has order \(12\), not \(24\).  If
\(r:\mathbb Z/24\to\mathbb Z/12\) is ordinary reduction, then
\(r\circ j=[2]\), and there is no group homomorphism
\(\widetilde\chi:SL_2(\mathbb Z)\to\mathbb Z/24\) with
\(r\circ\widetilde\chi=\chi\).
\end{theorem}

\begin{proof}
Put
\[
 S_0=\begin{pmatrix}0&-1\\1&0\end{pmatrix},\qquad T=C_0.
\]
Then \(C_2'=S_0^3\), \(C_1=T^{-1}S_0\), and the standard integral modular
presentation is
\[
 SL_2(\mathbb Z)=
 \langle S_0,T\mid S_0^4=1,\ S_0^2=(S_0T)^3\rangle.
\]
Here is why these relations suffice.  With \(R=S_0T\), the presented
group is \(G=\langle S_0,R:S_0^4=1,S_0^2=R^3\rangle\).
The element \(c=S_0^2=R^3\) is central and has order at most two;
\(G/\langle c\rangle=C_2*C_3\).  The modular normal-form theorem
identifies \(PSL_2(\mathbb Z)\) with this same free product on the images of
\(S_0,R\) \cite[Theorem~C.1, pp.~18--19]{h24:ConradSL2Z}.
The matrices \(S_0,T\) generate \(SL_2(\mathbb Z)\)
\cite[Theorem~1.1, pp.~1--2]{h24:ConradSL2Z}; hence the map from \(G\)
onto \(SL_2(\mathbb Z)\) is surjective and its kernel lies in
\(\langle c\rangle\).  Since \(c\) maps to \(-I_2\ne I_2\), the
kernel is trivial.  This proves the presentation (equivalently
\(C_4*_{C_2}C_6\)) with the stated matrix generators.
After abelianization, with additive generators \(s,t\), its two relations
are
\[
 4s=0,\qquad 2s=3(s+t),
\]
or \(4s=0\) and \(s+3t=0\).  The relation matrix
\(\left(\begin{smallmatrix}4&0\\1&3\end{smallmatrix}\right)\) has Smith
normal form \(\operatorname{diag}(1,12)\), so the abelianization is cyclic of exact order
\(12\), generated by \(t\).  Choosing \(t\mapsto-1\) gives
\[
 s\mapsto3,\qquad
 C_2'=S_0^3\mapsto9=-3,\qquad
 C_1=T^{-1}S_0\mapsto1+3=4,
\]
which proves \eqref{h24:eq:peripheral-mod12-values}.  Direct calculation in
\(\mathbb Z/12\) gives
\[
 4^2=4,\qquad9^2=9,\qquad4+9=1,\qquad4\cdot9=0.
\]

Every group homomorphism from \(SL_2(\mathbb Z)\) to an abelian group factors through its
abelianization.  An element \(x\in\mathbb Z/24\) satisfies \(12x=0\) exactly when
\(x\) is even, proving the unique factorization through \(j\) and
excluding a surjective character.  Also \(rj(a)=2a\pmod {12}\).
Finally a reduction-compatible lift
would send \(C_0\) to either \(11\) or \(23\) modulo \(24\), both of which
have order \(24\); an image of the order-\(12\) abelianized generator
cannot have that order.  This proves the final assertion.
\end{proof}

\begin{proposition}[The primitive cusp quotient singles out mod--3, not
mod--12 or mod--24]
\label[proposition]{h24:prop:cusp-diagonal-modulus}
For \(N\ge1\), let
\[
 d_N:\mathbb Z^3\longrightarrow\mathbb Z/N,\qquad
 d_N(a_1,a_2,a_3)=a_1+a_2+a_3\pmod N.
\]
There is a homomorphism \(\overline d_N:\mathbb Z^2\to\mathbb Z/N\) satisfying
\(d_N=\overline d_NP\) if and only if \(N\mid3\).  For \(N=3\), it is
explicitly
\[
 \overline d_3(y_1,y_2)=y_1-y_2\pmod3.
\]
Similarly, \(d_N\) factors through \(Q\) if and only if \(N\mid3\); for
\(N=3\) the exact formula is
\[
 d_3=(1,1)Q.
\]
Thus neither primitive lattice quotient in \cref{h24:thm:lattice-quotient}
produces a diagonal-sum character modulo \(12\) or modulo \(24\).
\end{proposition}

\begin{proof}
Both \(\ker P\) and \(\ker Q\) are generated by
\(k=(1,1,1)^{\mathsf T}\).  A homomorphism out of \(\mathbb Z^3\) factors
through either quotient precisely when it kills that kernel.  Since
\(d_N(k)=3\pmod N\), this occurs precisely when \(N\mid3\).  At \(N=3\),
direct multiplication with the matrices in \eqref{h24:eq:PQ} gives
\[
 (1,-1)P=(1,-2,1)\equiv(1,1,1)\pmod3,\qquad
 (1,1)Q=(1,-2,1)\equiv(1,1,1)\pmod3,
\]
which proves both formulas and the claimed boundary.
\end{proof}

\begin{proposition}[Coefficient reduction and determinant winding do not
select a modulus]
\label[proposition]{h24:prop:no-modulus-from-cohomology-or-determinant}
For every \(N\ge1\), coefficient reduction gives a graded-ring
isomorphism
\[
 H^*(\operatorname{Fl}(\mathbb O);\mathbb Z/N)
 \cong H^*(\operatorname{Fl}(\mathbb O);\mathbb Z)\otimes\mathbb Z/N,
\]
and its reduced Euler class
\[
 e_N=6u_{\mathbb O}\bmod N\in H^{24}(\operatorname{Fl}(\mathbb O);\mathbb Z/N)
\]
has additive order \(N/\gcd(N,6)\).  In particular \(e_{12}\) has order
two and \(e_{24}\) has order four.  Ordinary coefficient reduction is
defined for every \(N\); this operation alone makes no choice of
\(N=12\) or \(N=24\).

Likewise, in the lower complex orbit pair \((S^{10},S^8)\), the
determinant phase identifies the fundamental group of the rank-full
complement in \eqref{h24:eq:rank-complement} with \(\mathbb Z\).  Reducing its winding
number modulo \(N\) gives a surjection to \(\mathbb Z/N\) for every \(N\), so
this plain winding-number reduction makes no choice of \(12\) or \(24\).  This
phase statement is not being transferred to the real octonionic
determinant in \eqref{h24:eq:octonionic-equal-diagonal-determinant}.
\end{proposition}

\begin{proof}
All integral cohomology groups in
\eqref{h24:eq:octonionic-Poincare} are free.  The universal coefficient
sequence therefore has no \(\operatorname{Tor}\) term, and multiplication
is preserved by coefficient change.  Equation
\eqref{h24:eq:octonionic-Euler-six} gives the stated reduced Euler class.
Since the reduction of the primitive \(u_{\mathbb O}\) generates
\(H^{24}(\operatorname{Fl}(\mathbb O);\mathbb Z/N)\cong\mathbb Z/N\), the least positive
integer \(m\) annihilating \(e_N\) is the least \(m\) with \(N\mid6m\),
namely \(N/\gcd(N,6)\).
For the determinant complement,
\eqref{h24:eq:rank-complement} is a deformation retraction onto the phase
circle, and \(\arg\det\) is that retraction coordinate.  Its winding is an
isomorphism on \(\pi_1\), whose reduction modulo any \(N\) is surjective.
\end{proof}

\subsubsection[The distinct stable mod--12-to-mod--24 injection]
{The distinct stable \(\mathbb Z/12\hookrightarrow\mathbb Z/24\)}

The injection \(j(a)=2a\) is realized, after compatible generator choices,
by the universal sixfold suspension homomorphism below.  The cusp quotient
supplies one particular element in its domain; it does not supply the
injection itself.  Neither construction comes from the octonionic flag or
from a \(24\)-dimensional smash source.  We state all data needed to type
these assertions.

Let \(v_0\) be the suspension vertex to which the central cusp fibre is
collapsed, write \(W=q_2^{-1}(v_0)\), and choose explicitly
\(x_{\mathrm{base}}\in W\).  The exact cusp collapse in the canonical audit
then gives a based quotient
\[
 q_2:(X,x_{\mathrm{base}})\longrightarrow
 (\mathcal A,v_0),\qquad
 \mathcal A=\operatorname{Susp}(S^1\times\overline T^2).
\]
Let
\[
 K^{(1)}=(S^1\times\{0\})\cup(\{0\}\times S^1)
 \subset\overline T^2,
\]
and let \(c:\overline T^2\to\overline T^2/K^{(1)}\cong S^2\) be the
degree-one cell quotient, with \(*=c(K^{(1)})\).  Write
\[
 \Sigma_r(S^2,*)=
 \frac{[0,\epsilon]\times S^2}
 { (\{0,\epsilon\}\times S^2)\cup([0,\epsilon]\times\{*\})}
\]
for reduced suspension, the entire denominator being collapsed to one
basepoint.  The formula
\[
 \widehat p_2([r,e^{i\phi},u])=[r,c(u)]
 \in\Sigma_r(S^2,*)\cong S^3
\]
defines a based map: it sends both suspension vertices of \(\mathcal A\)
and the entire cylinder over \(K^{(1)}\) to that one basepoint.  Thus
\begin{equation}
 g=\widehat p_2q_2:(X,x_{\mathrm{base}})\longrightarrow(S^3,*).
 \label{h24:eq:one-copy-g}
\end{equation}
After the fixed based identification \(X\cong S^6\), this determines
\([g]\in\pi_6(S^3)\).  The construction and its endpoint continuity are
part of the cited candidate audit \cite{h24:AlpogeS6_2026}.
The exact research-file locators are the reduced-suspension construction in
\begin{center}\small
\texttt{cusp\_sphere\_route\_strengthening.tex}
\end{center}
and the one-copy map in
\begin{center}\small
\texttt{cusp\_s12\_carrier\_homotopy.tex}.
\end{center}

\begin{theorem}[The stable even-subgroup morphism and its exact scope]
\label[theorem]{h24:thm:stable-12-to-24}
Sixfold reduced suspension is an injection
\begin{equation}
 E^6:\pi_6(S^3)\cong\mathbb Z/12
 \hookrightarrow
 \pi_{12}(S^9)\cong\mathbb Z/24
 \label{h24:eq:stable-even-injection}
\end{equation}
whose image is the even subgroup.  With compatible generators \(\alpha\) and
\(\nu_9\),
\[
 E^6(a\alpha)=2a\nu_9.
\]
Consequently there is a unique \(a\in\mathbb Z/12\) for which the particular
map in \eqref{h24:eq:one-copy-g} satisfies
\begin{equation}
 [E^6g]=2a\nu_9.
 \label{h24:eq:one-copy-stable-class}
\end{equation}
The coefficient \(a\) is not determined by point-set surjectivity or by
the ordinary homology map.

For every pair of based maps \(x,y:S^6\to S^3\),
\begin{equation}
 [x\wedge y]=0\in\pi_{12}(S^6).
 \label{h24:eq:external-smash-zero}
\end{equation}
In particular,
\[
 g^{\wedge4}:S^{24}\longrightarrow S^{12}
\]
is based null-homotopic.  Thus the \(\mathbb Z/24\) in
\eqref{h24:eq:stable-even-injection} is the ambient target group of the
sixfold-suspended one-copy map \(S^{12}\to S^9\).  The separate
\(24\)-dimensional fourth-smash map is null-homotopic.
\end{theorem}

\begin{proof}
The relevant initial suspension chain is
\[
 \pi_6(S^3)\xrightarrow{E}\pi_7(S^4)
                   \xrightarrow{E}\pi_8(S^5).
\]
Hatcher computes these groups as \(\mathbb Z/12\), \(\mathbb Z\oplus\mathbb Z/12\),
and \(\mathbb Z/24\), respectively.  The Hopf fibration identifies the torsion
summand of the middle group with the suspended first group.  In the
quotient defining the second suspension, twice the Hopf generator of
the free summand is identified with a generator of this torsion summand
\cite[Theorem~5.41(b), p.~573, and pp.~577--578]
{h24:HatcherSpectralSequences}.
Choose compatible signs and generators \(\alpha\in\pi_6(S^3)\) and
\(\nu_5=E\nu_4\in\pi_8(S^5)\).  The displayed quotient relation gives
\(E^2\alpha=2\nu_5\).  Multiplication by two from \(\mathbb Z/12\) to
\(\mathbb Z/24\) is injective: \(24\mid2a\) implies \(12\mid a\).
For \(n\ge5\), Freudenthal gives an isomorphism
\[
 E:\pi_{n+3}(S^n)\longrightarrow\pi_{n+4}(S^{n+1}),
 \qquad n+3\le2n-2.
\]
Applying it for \(n=5,6,7,8\) and putting \(\nu_9=E^4\nu_5\)
gives \(E^6\alpha=2\nu_9\) and
\(\pi_{12}(S^9)\cong\pi_3^{\mathrm S}\cong\mathbb Z/24\).
The stable group is independently tabulated in
\cite[Table~1, p.~113]{h24:IsaksenWangXu2023}.  This proves
\eqref{h24:eq:stable-even-injection} and the unique formula
\eqref{h24:eq:one-copy-stable-class}.

For \eqref{h24:eq:external-smash-zero}, the ordered external smash has the
exact factorization
\[
 [x\wedge y]=(E^3[x])\circ(E^6[y])
\]
by the reduced-join formula
\cite[Proposition~3.1]{h24:Toda1962}; its ordered-coordinate sign is
\((-1)^{3\cdot6}=+1\).  The integral low-dimensional table gives
\(\pi_{12}(S^6)\cong\mathbb Z/2\)
\cite[opening table, p.~339]{h24:Hatcher2002}.  The class \(E^6[y]\) is even in
\(\pi_{12}(S^9)\cong\mathbb Z/24\), and postcomposition by \(E^3[x]\) is a group
homomorphism into the exponent-two group \(\pi_{12}(S^6)\).  Its value on
an even element is therefore zero.  Hence \(g^{\wedge2}\) is
null-homotopic, and smashing its based null-homotopy with
\(g^{\wedge2}\) proves the assertion for \(g^{\wedge4}\).
\end{proof}

\begin{proposition}[The generator-free bridge between the two
\(\mathbb Z/12\)-objects]
\label[proposition]{h24:prop:generator-free-bridge}
The marked one-copy class defines a homomorphism
\begin{equation}
 \Theta_g:\mathbb Z/12\longrightarrow\pi_6(S^3),\qquad
 \Theta_g([n])=n[g].
 \label{h24:eq:Theta-g}
\end{equation}
Therefore
\[
 \Theta_g\circ\chi:SL_2(\mathbb Z)\longrightarrow\pi_6(S^3)
\]
is a well-defined homomorphism from the marked peripheral data to the one-copy
homotopy group.  If \(\alpha\) and \(\nu_9\) are compatible generators as
in \cref{h24:thm:stable-12-to-24} and \([g]=a\alpha\), then
\begin{equation}
 E^6\Theta_g\chi(M)=2a\,\chi(M)\nu_9
 \in\pi_{12}(S^9)\cong\mathbb Z/24.
 \label{h24:eq:generator-free-stable-bridge}
\end{equation}
For an integer lift \(\widetilde a\) put \(d=\gcd(\widetilde a,12)\).
Then
\[
 \ker\Theta_g=\langle[12/d]\rangle\subset\mathbb Z/12,
 \qquad \operatorname{im}\Theta_g=\langle[g]\rangle=\langle d\alpha\rangle,
 \qquad |\operatorname{im}\Theta_g|=12/d.
\]
The first group has order \(d\).  Since \(\chi\) is surjective and
\(E^6\) injective, both composites have image order \(12/d\), and
\(\ker(E^6\Theta_g\chi)=\chi^{-1}(\langle[12/d]\rangle)\).
This bridge is canonical once the marked character \(\chi\), the based
map \(g\), and the based identification \(S^6\cong X\) are fixed.
It does not prove \(d=1\), equivalently that \([g]\) is a
generator or that \(\Theta_g\) is an isomorphism, and it does not define an
equivariant map of degenerations.
\end{proposition}

\begin{proof}
The source of \(\Theta_g\) is well-defined because
\(\pi_6(S^3)\cong\mathbb Z/12\), so \(12[g]=0\).  It is a homomorphism by the
abelian group law on \(\pi_6(S^3)\).  Composition with the already defined
\(\chi\) is therefore a well-defined homomorphism.  Writing
\([g]=a\alpha\) and applying the homomorphism \(E^6\) gives
\[
 E^6\Theta_g\chi(M)
 =\chi(M)E^6(a\alpha)=2a\,\chi(M)\nu_9,
\]
which is \eqref{h24:eq:generator-free-stable-bridge} for every \(M\in SL_2(\mathbb Z)\).
The congruence \(n\widetilde a\equiv0\pmod{12}\) is equivalent to
\(12/d\mid n\), giving the kernel.  B\'ezout's identity gives
\(\langle a\alpha\rangle=\langle d\alpha\rangle\), of order \(12/d\).
Surjectivity of \(\chi\) preserves that image, and injectivity of
\(E^6\) preserves its order and the composite kernel.  This group calculation by itself does not compute \(a\).
The explicit framed calculation in \cref{h24:thm:cusp-g-zero} below does:
\(a=0\), so \(d=12\), \(\ker\Theta_g=\mathbb Z/12\), and every
image displayed here is zero.
\end{proof}

\subsubsection{The original framed cusp coefficient is zero}
Let \(X\) be the smooth compact source of the original construction,
\(f:X\to B\) its map, and \(t\) the specified cusp coordinate.
Put \(W=f^{-1}(0)\), choose the original \(x_0\in W\), and retain the
ordered quotient lattice basis induced by
\((\widehat\gamma,\widehat u)\), with kernel generated by
\((\widehat w,\widehat\delta)\). Thus the quotient torus is
\(\overline T^2=\mathbb R^2/\mathbb Z^2\). The original periods are
\begin{equation}
 Z(s)=sB_0+C(t),\qquad e^{2\pi i s}=t,\qquad
 B_0=\begin{pmatrix}0&1\\-1&0\end{pmatrix},\qquad
 C(t)=\begin{pmatrix}6\mu(t)&h(t)\\b(t)-h(t)&\mu(t)\end{pmatrix}.
 \label{h24:gzero:eq:periods}
\end{equation}
Every entry of \(C\) is holomorphic at \(t=0\).
The dense torus of the original toric cusp has coordinates
\((x_1,x_2,t)\), with
\(x=\exp(2\pi i\zeta)\). The residual lattice action is exactly
\begin{equation}
 \Psi_\lambda(x,t)=
 \bigl(t^{B_0\lambda}\exp(2\pi iC(t)\lambda)x,t\bigr),
 \qquad\lambda\in\mathbb Z^2.
 \label{h24:gzero:eq:action}
\end{equation}
It extends freely and properly discontinuously across the original
unimodular toric filling. In particular the quotient projection is a local
biholomorphism at every toric fixed point. These are the reconstructed
data and proved local properties of the original source
\cite[\S4, source pp.~24--29]{h24:AlpogeS6_2026}.

Let
\[
 K^{(1)}=(\mathbb R/\mathbb Z\times\{0\})\cup(\{0\}\times\mathbb R/\mathbb Z),\qquad
 c:\overline T^2\longrightarrow \overline T^2/K^{(1)}\cong S^2
\]
be the specified degree-one cell quotient. Its restriction to the open
two-cell is taken with the smooth coordinates used in the original
regular-fibre description. Both radial endpoints and the complete
basepoint cylinder are collapsed in
\[
 \Sigma_r(S^2,*)=
 \frac{[0,\epsilon]\times S^2}
 {(\{0,\epsilon\}\times S^2)\cup([0,\epsilon]\times\{*\})}
 \cong S^3.
\]
The original based map is
\begin{equation}
 g(x)=[|t|,c(\overline\pi_f(x))]\quad(0<|t|<\epsilon),\qquad
 g(W)=g(X\setminus\operatorname{int}N_0)=*.
 \label{h24:gzero:eq:g}
\end{equation}

\begin{theorem}[The original cusp map is null-homotopic]
\label[theorem]{h24:thm:cusp-g-zero}
The based map \(g:(X,x_0)\to(S^3,*)\) in \eqref{h24:gzero:eq:g} is based
null-homotopic. More precisely, a regular inverse image with its actual
ordered normal framing induced by \((r,c_1,c_2)\) is framed nullbordant
inside one original cusp chart times \([0,1]\), disjoint from
\(\{x_0\}\times[0,1]\).
\end{theorem}

\begin{proof}
The proof constructs the complete inverse image, compares its actual
framing with an explicitly extendible framing, and gives the extension.

\paragraph{1. The exact complete inverse image.}
Retain the lift
\[
 u^0=\begin{pmatrix}-1/4\\1/4\end{pmatrix},\qquad
 \overline u^0=(3/4,1/4)\pmod{\mathbb Z^2},\qquad
 B_0u^0=\begin{pmatrix}1/4\\1/4\end{pmatrix}.
\]
Thus \(\overline u^0\notin K^{(1)}\), and \(B_0u^0\) lies in the
interior of the original triangle \(\operatorname{conv}(0,e_1,e_2)\).
Use its original unimodular chart
\begin{equation}
 z_0=t/(x_1x_2),\qquad z_1=x_1,\qquad z_2=x_2,
 \qquad t=z_0z_1z_2.
 \label{h24:gzero:eq:chart}
\end{equation}
Set
\[
 \ell=\log|t|<0,\qquad R(t)=-2\pi\operatorname{Im}C(t),\qquad
 A(t)=\ell B_0+R(t).
\]
For a representative \(\zeta=Z(s)u+k\), where \(u,k\in\mathbb R^2\),
\(\operatorname{Im}s=-\ell/(2\pi)\) gives the exact identity
\begin{equation}
 \begin{pmatrix}\log|x_1|\\\log|x_2|\end{pmatrix}
 =A(t)u.
 \label{h24:gzero:eq:moduli}
\end{equation}
The quotient value is \(u\bmod\mathbb Z^2\).
Since \(B_0\) is orthogonal, for all \(v\in\mathbb R^2\),
\begin{equation}
 \|A(t)v\|\geq (|\ell|-\|R(t)\|)\|v\|.
 \label{h24:gzero:eq:invertibility}
\end{equation}
The boundedness of \(R\) therefore makes \(A(t)\) invertible on a
sufficiently small punctured disc.

For any point of the fibre quotient over \(\overline u^0\), subtracting
the appropriate original period, equivalently applying
\(\Psi_{-\lambda}\), gives a representative with the exact lift
\(u=u^0\). This representative is unique: another residual translate
has logarithmic moduli differing by \(A(t)\lambda\), which vanishes
only for \(\lambda=0\). Conversely every dense-torus point satisfying
\(\log|x|=A(t)u^0\) has this quotient value, since the remaining real
parts of \(\zeta\) are exactly the kernel coordinates \(k\).
For each fixed \(t\), varying \(k\bmod\mathbb Z^2\) gives all independent
phases of \(x_1,x_2\).

The following change of coordinates uses the full original period units:
\begin{align}
 h_j^0(t)&=\exp\bigl(-2\pi i(C(t)u^0)_j\bigr),\qquad j=1,2,\\
 \mathcal F(z_0,z_1,z_2)
 &=(w_0,w_1,w_2)
 =\left(\frac{z_0}{h_1^0(t)h_2^0(t)},\,
 h_1^0(t)z_1,\,h_2^0(t)z_2\right).
 \label{h24:gzero:eq:straightening}
\end{align}
The functions \(h_j^0\) are nowhere-zero holomorphic units, including
at \(t=0\). The product \(w_0w_1w_2\) is exactly \(t\).
The inverse of \(\mathcal F\) is
\[
 (w_0,w_1,w_2)\longmapsto
 \left(w_0h_1^0(t)h_2^0(t),\,
 \frac{w_1}{h_1^0(t)},\,\frac{w_2}{h_2^0(t)}\right),
 \qquad t=w_0w_1w_2.
\]
Thus this is a biholomorphism of the original chart restricted to the
cusp disc, not a replacement of \(C(t)\).

For \(r_0\in(0,\epsilon)\), the entire inverse image of
\(y_0=[r_0,c(\overline u^0)]\) becomes exactly
\begin{equation}
 T_{r_0}=
 \{|w_0|=\rho_0,\ |w_1|=\rho_1,\ |w_2|=\rho_2\},\qquad
 (\rho_0,\rho_1,\rho_2)=
 (r_0^{1/2},r_0^{1/4},r_0^{1/4}).
 \label{h24:gzero:eq:torus}
\end{equation}
Indeed, \(\log|h^0(t)|=2\pi\operatorname{Im}C(t)u^0=-R(t)u^0\),
so \eqref{h24:gzero:eq:moduli} gives \(\log|w_j|=\ell(B_0u^0)_j\) for
\(j=1,2\); the product identity gives the exponent \(1/2\) for
\(w_0\). Conversely these moduli recover \(u=u^0\) by invertibility
of \(A(t)\). The angular map is explicitly
\[
 (\arg t,\arg w_1,\arg w_2)\longleftrightarrow
 (\arg w_0,\arg w_1,\arg w_2),\qquad
 \arg w_0=\arg t-\arg w_1-\arg w_2\pmod{2\pi},
\]
with inverse \(\arg t=\arg w_0+\arg w_1+\arg w_2\).
Its integral matrix has determinant one. This includes the complete
angular mapping torus and all kernel phases, with their original signs.

Let \(o\) be the origin of \eqref{h24:gzero:eq:chart}. The original free properly
discontinuous action supplies a neighborhood \(V\) of \(o\) on which
the quotient projection is injective. To see the local injectivity
directly, first choose a relatively compact neighborhood; proper
discontinuity leaves only finitely many translates that can intersect it.
None of those nonidentity translates fixes \(o\). Shrinking finitely
many times separates the neighborhood from all of them. The quotient
map is a local biholomorphism, so its restriction to \(V\) is a chart
of the original source.
Choose a polydisc about \(0\) whose closure lies in \(\mathcal F(V)\).
For sufficiently small \(r_0\), both \(T_{r_0}\) and
\begin{equation}
 \{|w_0|\leq\rho_0,\ |w_1|=\rho_1,\ |w_2|=\rho_2\}
 \label{h24:gzero:eq:spatial-bound}
\end{equation}
lie in this polydisc. This proves that the whole inverse image embeds
in one chart and supplies space for the bounding manifold below.

\paragraph{2. The actual normal framing, including its angle dependence.}
Put \(L_j=\log|w_j|\), and use the unique real lift of the quotient
coordinate near \(u^0\). From the exact change of coordinates,
\begin{equation}
 \begin{pmatrix}L_1\\L_2\end{pmatrix}
 =\ell B_0u^0+A(t)(u-u^0),\qquad
 L_0=\ell-L_1-L_2.
 \label{h24:gzero:eq:nearby-moduli}
\end{equation}
Differentiating on \(u=u^0\) annihilates the term
\((dA(t))(u-u^0)\). Hence no derivative of \(C(t)\) has been omitted:
the exact differential there is
\begin{equation}
 \begin{pmatrix}dL_0\\dL_1\\dL_2\end{pmatrix}
 =J(t)\begin{pmatrix}d\ell\\du_1\\du_2\end{pmatrix},\quad
 J(t)=
 \begin{pmatrix}
 1/2&-A_{11}-A_{21}&-A_{12}-A_{22}\\
 1/4&A_{11}&A_{12}\\
 1/4&A_{21}&A_{22}
 \end{pmatrix}.
 \label{h24:gzero:eq:J}
\end{equation}
Adding the second and third rows to the first gives \((1,0,0)\),
and expansion in that row proves \(\det J(t)=\det A(t)\).

Let \(D=Dc_{\overline u^0}\) be the actual two-by-two derivative in
the ordered open-cell target coordinates \((c_1,c_2)\). It is a fixed
invertible matrix; all its entries and signs are retained. Let
\[
 H=\begin{pmatrix}r_0&0&0\\0&D_{11}&D_{12}\\0&D_{21}&D_{22}\end{pmatrix}.
\]
Let \(n_j\) be the real radial vector \(w_j\) in the \(j\)-th complex
coordinate, with other components zero. Then \(dL_i(n_j)=\delta_{ij}\).
Writing \(N=(n_0,n_1,n_2)\) as the matrix with these vector columns, the actual
\(dg\)-dual normal framing is, in the normal quotient bundle,
\begin{equation}
 \mathfrak n(t)=N J(t)H^{-1}.
 \label{h24:gzero:eq:actual-frame}
\end{equation}
Indeed, \(dg(N)=HJ(t)^{-1}\), so multiplication by \(J(t)H^{-1}\)
gives the identity. In the Euclidean metric of this chart the radial
vectors are perpendicular to the product torus; a different source
metric changes their representatives by the canonical projection to its
normal bundle and preserves this normal-quotient framing.

Put \(\ell_0=\log r_0\). Decrease \(r_0\) further, if necessary, so
that \(|\ell_0|>2\sup_{|t|\leq r_0}\|R(t)\|\). For
\(0\leq\sigma\leq1\), define only a homotopy of frames by
\[
 A_\sigma(t)=\ell_0B_0+\sigma R(t),
\]
and let \(J_\sigma(t)\) be \eqref{h24:gzero:eq:J} with \(A_\sigma\) in place
of \(A\). For every \(v\in\mathbb R^2\),
\[
 \|A_\sigma(t)v\|\geq
 (|\ell_0|-\|R(t)\|)\|v\|>0\quad(v\ne0).
\]
Thus all matrices are invertible. At \(\sigma=0\),
\begin{equation}
 J_0=
 \begin{pmatrix}
 1/2&\ell_0&-\ell_0\\
 1/4&0&\ell_0\\
 1/4&-\ell_0&0
 \end{pmatrix},\qquad \det J_0=\ell_0^2>0.
 \label{h24:gzero:eq:J0}
\end{equation}
Continuity also proves \(\det J_\sigma(t)>0\) for every \(\sigma,t\).
The homotopy \(N J_\sigma(t)H^{-1}\), run from \(\sigma=1\) to
\(\sigma=0\), compares the actual frame with the exact constant-matrix
frame \(N P\), where
\begin{equation}
 P=J_0H^{-1}.
 \label{h24:gzero:eq:P}
\end{equation}
The original embedded torus and original functions \(C(t)\) have not
changed. A framing homotopy is a framed cylinder: choose a smooth
parameter function constant near its endpoints and use the displayed
frames on \(T_{r_0}\times[0,1]\).

\paragraph{3. A collared framed nullbordism with the exact constant matrix.}
Use a new time coordinate \(v\in[0,1]\). Define
\begin{equation}
 q(v)=
 \begin{cases}
 0,&0\leq v\leq1/4,\\
 \rho_0^2\exp\bigl(4-1/(v-1/4)\bigr),&1/4<v\leq1.
 \end{cases}
 \label{h24:gzero:eq:q}
\end{equation}
This function is smooth: all right derivatives of
\(\exp(-1/(v-1/4))\) at \(1/4\) vanish. It is strictly increasing
for \(v>1/4\), satisfies \(q(1/2)=\rho_0^2\), and
\(q'(1/2)=16\rho_0^2>0\).
In the chart times \([0,1]\), put
\begin{equation}
 \mathcal W=
 \{|w_0|^2+q(v)=\rho_0^2,\quad
       |w_1|^2=\rho_1^2,\quad|w_2|^2=\rho_2^2\}.
 \label{h24:gzero:eq:bordism}
\end{equation}
No point has \(v>1/2\). The three defining differentials are independent:
the latter two have nonzero gradients in their separate complex factors;
the first has nonzero \(w_0\)-gradient whenever \(w_0\ne0\), and at
\(w_0=0\) it has the nonzero \(v\)-derivative \(q'(1/2)\).
Thus \(\mathcal W\) is a smooth compact four-manifold. Its only
boundary is \(T_{r_0}\times\{0\}\), and on \(0\leq v\leq1/4\)
it is the exact product collar \(T_{r_0}\times[0,1/4]\).
The north locus \(w_0=0,v=1/2\) is an interior cap: the implicit
function theorem expresses \(v\) smoothly in terms of
\((\operatorname{Re}w_0,\operatorname{Im}w_0)\) there. The first
factor is a disk obtained by rotating its radial profile with this smooth
cap and one boundary circle. Therefore
\(\mathcal W\cong D^2\times T^2\).
Its spatial coordinates lie in \eqref{h24:gzero:eq:spatial-bound}, so the complete
manifold lies in the injective original chart.

In the Euclidean chart product metric define
\[
 V_0=(w_0,0,0,q'(v)/2),\qquad
 V_1=(0,w_1,0,0),\qquad V_2=(0,0,w_2,0).
\]
These are one half of the gradients of the three defining equations
in \eqref{h24:gzero:eq:bordism}. They are mutually orthogonal and never zero.
For example, if \(F_0=|w_0|^2+q(v)-\rho_0^2\), then
\[
 dF_0(V_0)=2|w_0|^2+\frac12q'(v)^2>0
\]
at every point of \(\mathcal W\). On the entire collar \(q'=0\), so
\((V_0,V_1,V_2)=(N,0)\) exactly. Consequently
\begin{equation}
 (V_0,V_1,V_2)P
 \label{h24:gzero:eq:extended-frame}
\end{equation}
is a framing of the normal three-plane bundle of \(\mathcal W\),
and its restriction to the collar is exactly the framing \(NP\) from
\eqref{h24:gzero:eq:P}, with every original target-coordinate sign retained.
Concatenate the framing cylinder of the preceding paragraph with this
nullbordism, rescaling the two interval lengths. Their framings agree
exactly on the common product collar, so the concatenation is a smooth
framed nullbordism of the actual inverse image.

\paragraph{4. The based map represented by this framed inverse image.}
Choose the injective chart to exclude \(x_0\) if its center is not
\(x_0\). If its center is \(x_0\), the constructed nullbordism still
excludes \(x_0\), since \(|w_1|=\rho_1>0\) and
\(|w_2|=\rho_2>0\) everywhere on it. Compactness therefore gives
tubular neighborhoods for the framed cylinder and cap disjoint from
\(\{x_0\}\times[0,1]\).

The original continuous quotient map is smooth and a submersion near
its entire inverse image of \(y_0\). First make it constant in a
neighborhood of \(x_0\), by a based homotopy that changes nothing near
that inverse image. Explicitly, take a target chart
\(\kappa:O\to\mathbb R^3\) with \(\kappa(*)=0\), whose image is a ball,
and \(y_0\notin O\). By continuity choose a neighborhood \(U\) of
\(x_0\) with \(g(U)\subset O\), and a smooth bump function
\(\beta\) supported in \(U\) and equal to one near \(x_0\).
For \(x\in U\), replace \(g(x)\) during this homotopy by
\(\kappa^{-1}((1-a\beta(x))\kappa(g(x)))\), \(0\leq a\leq1\),
and leave \(g\) fixed off \(U\). Convexity of the target chart image
makes the formula well-defined, the vanishing of \(\beta\) near the
boundary of its support proves continuity of the pasted formula, and
\(\kappa(g(x_0))=0\) makes the homotopy based. It never introduces
an inverse image of \(y_0\). Its endpoint is smooth near \(x_0\) and
agrees with \(g\) near the entire old inverse image.

Denote the endpoint of the preceding homotopy by \(g_b\). Choose one
closed tubular neighborhood \(U_1\) of the inverse image, contained in
the smooth submersion region and with the entire inverse image contained
in \(\operatorname{int}U_1\). The compact set
\(X\setminus\operatorname{int}U_1\) has image under \(g_b\) missing
\(y_0\). In a fixed target metric put
\[
 d_0=\min_{x\in X\setminus\operatorname{int}U_1}
       d(g_b(x),y_0)>0.
\]
Relative smooth approximation, fixed on the whole closed set \(U_1\)
and on a neighborhood of \(x_0\), gives a smooth based representative
\(g_{\mathrm{sm}}\) whose error on
\(X\setminus\operatorname{int}U_1\) is less than \(d_0/2\).
It therefore has the same inverse image and differential on \(U_1\),
and \(d(g_{\mathrm{sm}}(x),y_0)>d_0/2\) on its complement. Thus there
is no uncontrolled intervening region and no additional inverse image.
Make the approximation also closer than a fixed target injectivity
radius; its short geodesics give the based homotopy.

The relative Pontryagin construction now applies to the explicit
framed nullbordism. It produces a based homotopy to a map missing
\(y_0\); stereographic contraction of \(S^3\setminus\{y_0\}\cong\mathbb R^3\)
to the image of \(*\) is based and contracts that last map.
For completeness, the dependence on the framed inverse image in this
argument is precisely the local product-neighborhood construction and
the interpolation lemma of \cite[\S7, Theorem B and Lemma 4,
pp.~43, 48--50]{h24:Milnor1965}. Their relative versions fix the source
basepoint because all tubular modifications are away from it and
stereographic interpolation between equal values keeps that value
fixed. A framed cylinder realizes a framing homotopy; the empty
Pontryagin manifold realizes a map missing the chosen regular value.
This proves the claimed based null-homotopy of the exact map \(g\).
\end{proof}

\begin{corollary}
Under the original fixed based identification \(X\cong S^6\), the
one-copy coefficient is
\[
 [g]=0\in\pi_6(S^3)\cong\mathbb Z/12,\qquad a=0\pmod{12},\qquad
 [E^6g]=0\in\pi_{12}(S^9)\cong\mathbb Z/24.
\]
The previously defined homomorphism
\(\Theta_g:\mathbb Z/12\to\pi_6(S^3)\), \(\Theta_g([n])=n[g]\), and
its composite with the marked peripheral character are zero.
\end{corollary}
\begin{proof}
The theorem states that the actual based map is null-homotopic.
Precomposition with the fixed based identification preserves a based
null-homotopy. Suspension preserves it as well. Thus the unique
coefficient of \([g]\) in any generator of \(\mathbb Z/12\) is zero.
For every \([n]\), \(\Theta_g([n])=n\cdot0=0\), and every
composite with this homomorphism is zero.
\end{proof}

\paragraph{Exact scope.}
The proof computes this sphere-valued postcomposition. It does not
compute the homotopy class of the earlier conical-target map
\(q_2:X\to\operatorname{Susp}(S^1\times\overline T^2)\), and it
does not assert a null-homotopy for that different target. Point-set
surjectivity of \(g\) and all of its stated metric restrictions remain
properties of the original map: a homotopy-class computation does not
alter the defining formula \eqref{h24:gzero:eq:g}.

% BEGIN RETAINED_ANGLE_AND_HOPF_BRIDGE_20260905

\subsubsection{The retained-angle cusp quotient and its nonzero top-cell class}
\label{h24:ret:section}

\paragraph{The exact object.}
The notation refers to the specified threefold \(X\), its fixed based
smooth identification with \(S^6\), and its cusp filling from
\cite[\S\S3.7,4.1--4.5]{h24:retbib:Canonical}.  The native quotient and its compact
extension are the maps in \cite{h24:retbib:ConeRoutes}; their full formulas are
recalled below.  This calculation uses the original period matrix and
native period lattices.  It computes a new postcomposition of that
particular quotient; it does not use the homotopy class of the
angle-forgetting map \(g:X\to S^3\).

\paragraph{Original coordinates and the specified sphere-valued map}

The ordered quotient lattice is
\[
 \overline\Lambda=\mathbb Z\langle\overline{\widehat\gamma},
              \overline{\widehat u}\rangle,
 \qquad
 K_\Lambda=\mathbb Z\langle\widehat w,\widehat\delta\rangle,
 \qquad
 \overline T^2=\mathbb R^2/\mathbb Z^2.
\]
In these ordered bases the original period data are
\begin{equation}
 B_0=\begin{pmatrix}0&1\\-1&0\end{pmatrix},\qquad
 Z(s)=sB_0+C(t),\qquad t=e^{2\pi i s},\qquad
 C(t)=\begin{pmatrix}6\mu(t)&h_{\rm per}(t)\\
 b(t)-h_{\rm per}(t)&\mu(t)\end{pmatrix}.
 \label{h24:ret:eq:original-periods}
\end{equation}
Here \(b=\beta+\tau\), \(s=\tau-h_{\rm per}(t)\), and every entry of
\(C\) is holomorphic at \(t=0\), exactly as in the original source.
The subscript in \(h_{\rm per}\) distinguishes that period function
from the sphere-valued map below.  Write
\begin{equation}
 R(t)=-2\pi\operatorname{Im}C(t),\qquad
 \ell=\log|t|,\qquad A(t)=\ell B_0+R(t).
 \label{h24:ret:eq:A}
\end{equation}
Choose a closed cusp disc \(|t|\leq\epsilon\) inside the original
coordinate disc and small enough for the quotient construction.  On a
slightly larger closed disc put \(M=\sup\|R(t)\|<\infty\).  We will
choose a regular value with \(0<r_0<\epsilon\) and
\(|\log r_0|>M\).  This chooses a point of the already specified map;
it does not change that map or the cutoff radius \(\epsilon\).

The original covering model on its dense torus is
\begin{equation}
 x_j=e^{2\pi i\zeta_j},\qquad
 \Psi_n(x,t)=
 \bigl(e^{2\pi i C(t)n}t^{B_0n}x,t\bigr),\qquad n\in\mathbb Z^2.
 \label{h24:ret:eq:deck}
\end{equation}
For a branch of \(s\), writing
\(\zeta=Z(s)u+k\) with \(u,k\in\mathbb R^2\) gives
\begin{equation}
 \log|x|=A(t)u.
 \label{h24:ret:eq:logx}
\end{equation}
Indeed, \(\operatorname{Im}s=-\ell/(2\pi)\), so
\(-2\pi\operatorname{Im}(Z(s)u+k)=\ell B_0u+R(t)u\).
Since \(B_0\) is orthogonal,
\begin{equation}
 \bigl|\bigl(\ell B_0+vR(t)\bigr)\xi\bigr|
 \geq (|\ell|-vM)|\xi|>0
 \quad(0\leq v\leq1,\ \xi\ne0)
 \label{h24:ret:eq:invertibility}
\end{equation}
whenever \(|\ell|>M\).  Thus \(A(t)\) is invertible there.  Equation
\eqref{h24:ret:eq:deck} changes \(A(t)^{-1}\log|x|\) by \(n\), so the native
quotient is exactly
\begin{equation}
 \overline\pi_f([x,t])=[A(t)^{-1}\log|x|]\in\mathbb R^2/\mathbb Z^2.
 \label{h24:ret:eq:native}
\end{equation}
This also follows from the flat marking in \cite{h24:retbib:ConeRoutes}.

Let
\[
 T^3=S^1_\phi\times\overline T^2,
 \qquad S^1_\phi=\mathbb R/(2\pi\mathbb Z),
\]
with coordinate basepoints \(\phi=0\), \(u_1=0\), \(u_2=0\).
Its coordinate two-skeleton is the following specified closed subset:
\begin{equation}
 K^{(2)}=
 \bigl(\{0\}\times\overline T^2\bigr)
 \cup\bigl(S^1_\phi\times\{u_1=0\}\bigr)
 \cup\bigl(S^1_\phi\times\{u_2=0\}\bigr).
 \label{h24:ret:eq:two-skeleton}
\end{equation}
The product CW structure has one open three-cell, with coordinates
\((\phi,u_1,u_2)\in(0,2\pi)\times(0,1)^2\).  Consequently
\[
 c_3:T^3\longrightarrow T^3/K^{(2)}\cong S^3
\]
is the degree-one quotient, where the orientation of \(S^3\) is fixed
by the displayed coordinate order.  One explicit identification of its
open cell with \(\mathbb R^3\) is
\[
 (\phi,u_1,u_2)\longmapsto
 \bigl(-\cot(\phi/2),-\cot(\pi u_1),-\cot(\pi u_2)\bigr).
\]
Each derivative in its own variable is strictly positive, and approach
to any boundary face leaves every compact subset of \(\mathbb R^3\).
Extending by its point at infinity gives the stated based sphere
identification and verifies its continuity on the collapsed skeleton.

The exact conical target and map are
\begin{align}
 \mathcal A&=
 \frac{[0,\epsilon]\times T^3}
 {\{0\}\times T^3\sim v_0,\quad
  \{\epsilon\}\times T^3\sim v_1},
 \label{h24:ret:eq:A-target}\\
 q_2(x)&=[r,e^{i\phi},\overline\pi_f(x)]
       &&(t(f(x))=re^{i\phi},\ 0<r<\epsilon),\notag\\
 q_2(W)&=v_0,\qquad
 q_2\bigl(X\setminus\operatorname{int}N_0\bigr)=v_1.
 \label{h24:ret:eq:q2}
\end{align}
Fix the same source basepoint \(x_{\rm base}\in W\).  Define
\begin{equation}
 p_3:\mathcal A\longrightarrow\Sigma_r(S^3,*),\qquad
 p_3([r,z])=[r,c_3(z)],\qquad
 H=p_3q_2:(X,x_{\rm base})\longrightarrow(S^4,*).
 \label{h24:ret:eq:H}
\end{equation}
The reduced suspension here collapses both endpoint copies of \(S^3\)
and the entire interval over \(*\) to one point.  Thus \(p_3\) sends
both \(v_0\) and \(v_1\) to \(*\) and is well-defined and based.
It is continuous by the defining quotient topology.  An explicit
identification of its open four-cell with \(\mathbb R^4\) is
\begin{equation}
 (r,\phi,u_1,u_2)\longmapsto
 \bigl(-\cot(\pi r/\epsilon),-\cot(\phi/2),
       -\cot(\pi u_1),-\cot(\pi u_2)\bigr).
 \label{h24:ret:eq:four-cell}
\end{equation}
This fixes the target orientation.  Replacing \(r\) locally by
\(\ell=\log r\) preserves this orientation, since \(dr/d\ell=r>0\).
The ordered coordinates in \eqref{h24:ret:eq:four-cell} specify the sphere
identification used throughout this calculation.

\begin{theorem}
\label{h24:ret:main}
For the exact map \eqref{h24:ret:eq:H},
\begin{equation}
 [H]=\eta_4\circ\eta_5\ne0
       \quad\hbox{in}\quad\pi_6(S^4)\cong\mathbb Z/2,
 \label{h24:ret:eq:main-class}
\end{equation}
where \(\eta_n:S^{n+1}\to S^n\) is the \((n-2)\)-fold reduced
suspension of the complex Hopf map.  In particular \(q_2\) is not
based null-homotopic.  Its stabilized top-cell postcomposition is the
element \(\eta^2\in\pi_2^{\rm S}\).
\end{theorem}

\paragraph{The entire regular inverse image lies in one original chart}

Choose the actual lift
\begin{equation}
 u_0=(-1/4,1/4)^{\mathsf T},\qquad
 p_0=B_0u_0=(1/4,1/4)^{\mathsf T},\qquad \phi_0=\pi.
 \label{h24:ret:eq:chosen-value}
\end{equation}
Modulo \(\mathbb Z^2\), \(u_0\) is \((3/4,1/4)\) in the open cell.  Put
\(t_0=r_0e^{i\phi_0}\) and
\(y_0=[r_0,c_3(e^{i\phi_0},[u_0])]\).  The numbers \(r_0\) and
\(t_0\) will remain nonzero throughout the framing calculation.

The original triangle \(\{0,e_1,e_2\}\) has toric chart
\begin{equation}
 \mathcal U_\sigma\cong\mathbb C^3,\qquad
 (z_0,z_1,z_2)=(t/(x_1x_2),x_1,x_2),\qquad
 t=z_0z_1z_2.
 \label{h24:ret:eq:z-chart}
\end{equation}
These characters are dual to \((0,0,1),(1,0,1),(0,1,1)\), respectively,
as direct pairing verifies.  This is the chart of
\cite[Lemma 4.2(iii)]{h24:retbib:Canonical}, with its original coordinates.

For \(j=1,2\), introduce the nowhere vanishing holomorphic functions
\begin{equation}
 h_j(t)=\exp\bigl(-2\pi i(C(t)u_0)_j\bigr).
 \label{h24:ret:eq:units}
\end{equation}
The following is an exact coordinate map, with the exact inverse
displayed alongside it:
\begin{align}
 F(z)&=(w_0,w_1,w_2)
 =\bigl(z_0/(h_1(t)h_2(t)),\ h_1(t)z_1,\ h_2(t)z_2\bigr),
 \label{h24:ret:eq:F}\\
 F^{-1}(w)&=
 \bigl(w_0h_1(t)h_2(t),\ w_1/h_1(t),\ w_2/h_2(t)\bigr),
 \quad t=w_0w_1w_2.
 \label{h24:ret:eq:Finv}
\end{align}
Both maps are holomorphic wherever \(|t|\) is in the period disc,
both preserve \(t\), and substitution proves that both composites are
the identity.  In particular the maps extend across the origin and
across every coordinate divisor.  No period matrix is replaced by a
different matrix in this coordinate change.

\begin{lemma}
For sufficiently small \(r_0>0\), the whole inverse image
\(M=H^{-1}(y_0)\) is represented injectively inside one contractible
coordinate neighborhood in \(X\).  In its \(w\)-coordinates it is
exactly the embedded two-torus
\begin{equation}
 w(a,b)=
 \bigl(\rho_0e^{i(\phi_0-a-b)},\rho_1e^{ia},\rho_2e^{ib}\bigr),
 \qquad
 \rho_0=r_0^{1/2},\quad \rho_1=\rho_2=r_0^{1/4},
 \quad (a,b)\in\bigl(\mathbb R/(2\pi\mathbb Z)\bigr)^2.
 \label{h24:ret:eq:M}
\end{equation}
\end{lemma}
\begin{proof}
The point \(y_0\) lies off the collapsed subset of the target, so
\eqref{h24:ret:eq:H} and \eqref{h24:ret:eq:q2} imply
\[
 M=\{x:t(f(x))=t_0,\ \overline\pi_f(x)=[u_0]\}.
\]
Any lift of such a point to \((\mathbb C^*)^2\times\{t_0\}\) has
\(A(t_0)^{-1}\log|x|=u_0+n\) for one \(n\in\mathbb Z^2\).
Applying the actual transformation \(\Psi_{-n}\) from
\eqref{h24:ret:eq:deck} produces a lift with
\(\log|x|=A(t_0)u_0\).  If two such lifts represent the same point,
\(\Psi_n\) relates them and changes their lifted quotient coordinate
by \(n\); hence \(n=0\), and the lifts are identical.  This proves
both completeness and injectivity of this chosen representative.

Since
\(\log|h_j(t)|=2\pi\operatorname{Im}(C(t)u_0)_j
=-(R(t)u_0)_j\), its \(w\)-coordinates satisfy
\[
 (\log|w_1|,\log|w_2|)^{\mathsf T}
 =A(t_0)u_0-R(t_0)u_0=\ell_0p_0,
 \qquad \ell_0=\log r_0.
\]
Thus \(|w_1|=|w_2|=r_0^{1/4}\), and the unchanged product
\(w_0w_1w_2=t_0\) forces \(|w_0|=r_0^{1/2}\) and the displayed
phase.  Conversely every pair \((a,b)\) in \eqref{h24:ret:eq:M} gives one
such lift.  The original kernel coordinates \(k_j\) had period one
in \(x_j=e^{2\pi i\zeta_j}\); the phase coordinates \(a,b\) have
period \(2\pi\), with no cover or change of lattice.

The covering \(Y_\epsilon\to N_0\) is injective on a neighborhood of
the origin in \eqref{h24:ret:eq:z-chart}, because the original action is free
and properly discontinuous.  Choose a ball \(V\) centered at that
origin with closure inside such a neighborhood and the period disc.
Choose a ball \(V_w\) about the origin with closure contained in
\(F(V)\).  Every coordinate in \eqref{h24:ret:eq:M} tends uniformly to zero
as \(r_0\downarrow0\).  Therefore for sufficiently small \(r_0\)
the entire torus lies in \(V_w\).  The inverse image of \(V_w\)
under \(F\), followed by the injective covering projection, is the
required contractible chart in \(X\).
\end{proof}

\paragraph{The full differential and its ordered normal framing}

We calculate in a neighborhood of the compact torus \eqref{h24:ret:eq:M}.
The quotient coordinates admit a unique local lift near \(u_0\).
Set
\[
 v=(\log|w_1|,\log|w_2|)^{\mathsf T},\quad
 \ell=\log|w_0w_1w_2|,\quad
 \phi=\arg(w_0w_1w_2)
\]
using the branch of \(\phi\) near \(\phi_0=\pi\).
Equations \eqref{h24:ret:eq:logx} and \eqref{h24:ret:eq:F} give the exact identity
\begin{equation}
 v=\ell B_0u+R(t)(u-u_0),\qquad
 u=u_0+A(t)^{-1}(v-\ell p_0).
 \label{h24:ret:eq:exact-u}
\end{equation}
Differentiate this identity before restricting to the inverse image:
\begin{equation}
 du=-A(t)^{-1}(dA(t))A(t)^{-1}(v-\ell p_0)
        +A(t)^{-1}(dv-p_0\,d\ell).
 \label{h24:ret:eq:full-derivative}
\end{equation}
On \(M\) the vector \(v-\ell p_0\) is exactly zero.  Therefore
\begin{equation}
 du|_M=A_0^{-1}(dv-p_0\,d\ell),\qquad A_0=A(t_0).
 \label{h24:ret:eq:derivative-on-M}
\end{equation}
This explicitly accounts for all derivatives of the original bounded
period terms; none are omitted or estimated away.

Identify \(T_w\mathbb C^3\) with \(\mathbb C^3\) in the given coordinate chart,
and define real tangent vectors
\begin{equation}
 L_j=(0,\ldots,w_j,\ldots,0),\qquad
 Q_j=(0,\ldots,iw_j,\ldots,0)\qquad(0\leq j\leq2).
 \label{h24:ret:eq:LQ}
\end{equation}
Thus \(d\log|w_i|(L_j)=\delta_{ij}\),
\(d\arg w_i(Q_j)=\delta_{ij}\), and the crossed evaluations vanish.
The ordered tangent vectors of the parametrization \eqref{h24:ret:eq:M} are
\begin{equation}
 T_a=-Q_0+Q_1,\qquad T_b=-Q_0+Q_2.
 \label{h24:ret:eq:tangents}
\end{equation}
They are independent because their \(w_1\) and \(w_2\) components,
respectively, are nonzero.  For the target order
\((\ell,\phi,u_1,u_2)\), define the following four lifts of its
coordinate vectors:
\begin{align}
 N_\ell&=\tfrac12L_0+\tfrac14L_1+\tfrac14L_2,
       &N_\phi&=Q_0,\label{h24:ret:eq:normal-first}\\
 N_i&=-\bigl((A_0)_{1i}+(A_0)_{2i}\bigr)L_0
           +(A_0)_{1i}L_1+(A_0)_{2i}L_2
       &&(i=1,2).\label{h24:ret:eq:normal-u}
\end{align}
Let \(P=(\ell,\phi,u_1,u_2)\).  Evaluation using
\eqref{h24:ret:eq:derivative-on-M} gives the entire differential table
\begin{equation}
\begin{array}{c|rrrrrr}
 &N_\ell&N_\phi&N_1&N_2&T_a&T_b\\\hline
 d\ell&1&0&0&0&0&0\\
 d\phi&0&1&0&0&0&0\\
 du_1&0&0&1&0&0&0\\
 du_2&0&0&0&1&0&0
\end{array}.
 \label{h24:ret:eq:differential-table}
\end{equation}
In particular \(P\), and therefore \(H\) in its open target cell,
is a submersion along \(M\).  The classes of
\((N_\ell,N_\phi,N_1,N_2)\) in \(TX|_M/TM\) give the specified
normal four-framing.  Orthogonal representatives are not required for
this quotient-bundle definition; projecting the displayed lifts along
\(TM\) to any chosen orthogonal complement preserves the quotient
frame.

If the normal frame is specified using the literal four real coordinates
in \eqref{h24:ret:eq:four-cell}, their differential at the chosen value is
\[
 dH_{\rm cell}=K\,dP,\qquad
 K=\operatorname{diag}\!\left(
 \frac{\pi r_0}{\epsilon}\csc^2\frac{\pi r_0}{\epsilon},
 \frac12,2\pi,2\pi\right).
\]
The first factor includes \(dr=r_0\,d\ell\), the second uses
\(\phi_0=\pi\), and the final two use the actual representatives
\((3/4,1/4)\) of \(u_0\). Thus, writing
\(\mathcal N=(N_\ell,N_\phi,N_1,N_2)\), the literal dual frame is
\(\mathcal N K^{-1}\). All four entries of \(K\) are positive
constants on \(M\). The explicitly positive diagonal matrices
\((1-v)K^{-1}+vI_4\), \(0\leq v\leq1\), give a homotopy
from that literal frame to \(\mathcal N\). This retains the original
target coordinates and establishes that either coordinate description
gives the same framed bordism class.

For completeness, the original period matrix admits the following
explicit homotopy of this normal framing.  Put
\begin{equation}
 A_{0,v}=\ell_0B_0+vR(t_0),\qquad 0\leq v\leq1,
 \label{h24:ret:eq:A-homotopy}
\end{equation}
and replace \(A_0\) by \(A_{0,v}\) only in
\eqref{h24:ret:eq:normal-u}.  Equation \eqref{h24:ret:eq:invertibility} proves
invertibility for every \(v\).  Their images under the actual
\((du_1,du_2)\) are the columns of \(A_0^{-1}A_{0,v}\), so these
are normal frames throughout the homotopy.  This is a homotopy of
the computed frame on the same embedded torus, not a replacement of
the original degeneration by a different one.

\begin{lemma}[The exact orientation sign]
Use the complex real-coordinate orientation, whose ordered coordinates are
\[
 (\operatorname{Re}w_0,\operatorname{Im}w_0,
  \operatorname{Re}w_1,\operatorname{Im}w_1,
  \operatorname{Re}w_2,\operatorname{Im}w_2).
\]
Then one has
\begin{equation}
 \det(N_\ell,N_\phi,N_1,N_2,T_a,T_b)
       =-\det(A_0)|w_0w_1w_2|^2<0.
 \label{h24:ret:eq:sign}
\end{equation}
Hence the tangent orientation induced by the ordered normal frame is
the orientation of \((-T_a,T_b)\).
\end{lemma}
\begin{proof}
The ordered frame \((L_0,Q_0,L_1,Q_1,L_2,Q_2)\) has determinant
\(|w_0|^2|w_1|^2|w_2|^2>0\) in real coordinates.  Regroup its rows
as \((L_0,L_1,L_2,Q_0,Q_1,Q_2)\); this permutation has three
inversions and sign \(-1\).  Regroup the columns of the frame in
\eqref{h24:ret:eq:sign} as
\((N_\ell,N_1,N_2,N_\phi,T_a,T_b)\); this has two inversions and
sign \(+1\).  The angular block is
\(\left(\begin{smallmatrix}1&-1&-1\\0&1&0\\0&0&1\end{smallmatrix}\right)\),
of determinant one.  The radial block has first column
\((1/2,1/4,1/4)^{\mathsf T}\) and the other columns given by
\eqref{h24:ret:eq:normal-u}.  Adding its last two rows to its first row gives
first row \((1,0,0)\), with remaining lower-right block \(A_0\).
Its determinant is \(\det A_0\).  This proves the equality.
Equation \eqref{h24:ret:eq:A-homotopy} joins \(A_0\) to \(\ell_0B_0\)
through invertible matrices, so
\(\operatorname{sign}\det A_0=\operatorname{sign}\det(\ell_0B_0)=+1\).
The inequality and the tangent orientation follow.
\end{proof}

\paragraph{An explicit contraction of the full frame}

Define a complex-linear matrix, also regarded as a real-linear matrix,
\begin{equation}
 D(a,b)=\operatorname{diag}(e^{-i(a+b)},e^{ia},e^{ib})\in\operatorname{SU}(3).
 \label{h24:ret:eq:D}
\end{equation}
By \eqref{h24:ret:eq:M}, every \(L_j\) and \(Q_j\) at \((a,b)\) is
\(D(a,b)\) times its value at \((0,0)\).  The coefficients in
\eqref{h24:ret:eq:normal-first}, \eqref{h24:ret:eq:normal-u}, and
\eqref{h24:ret:eq:tangents} are constant on \(M\), since \(t=t_0\).
It follows that the positively oriented full frame is exactly
\begin{equation}
 (N_\ell,N_\phi,N_1,N_2,-T_a,T_b)_{(a,b)}
       =D(a,b)G_0,
 \label{h24:ret:eq:full-frame}
\end{equation}
where \(G_0\in\operatorname{GL}^+(6,\mathbb R)\) is its single value at \((0,0)\).
Thus neither an unspecified phase nor a position-dependent
real-linear matrix remains.

Here is a periodic contraction of \(D\), with all parameters
specified.  For \(v\in[0,1]\) put
\[
 s_v=\sin(\pi v/2),\qquad c_v=\cos(\pi v/2),\qquad
 Q_v(\theta)=
 \begin{pmatrix}s_ve^{-i\theta}&-c_v\\c_v&s_ve^{i\theta}\end{pmatrix},
 \qquad K_v(\theta)=Q_v(\theta)Q_v(0)^*.
\]
The two columns of \(Q_v\) have norm one and zero Hermitian inner
product, and \(\det Q_v=s_v^2+c_v^2=1\).  Thus
\(Q_v,K_v\in\operatorname{SU}(2)\).  Direct substitution gives
\[
 K_0(\theta)=I,\quad
 K_1(\theta)=\operatorname{diag}(e^{-i\theta},e^{i\theta}),\quad
 K_v(0)=I.
\]
Let \(\iota_{01},\iota_{02}:\operatorname{SU}(2)\to\operatorname{SU}(3)\) act in the
ordered coordinate planes \((0,1)\) and \((0,2)\), respectively,
and as the identity in the unused coordinate.  Then
\begin{equation}
 D_v(a,b)=\iota_{01}(K_v(a))\,\iota_{02}(K_v(b))
 \label{h24:ret:eq:D-homotopy}
\end{equation}
is a well-defined continuous map on
\([0,1]\times(\mathbb R/(2\pi\mathbb Z))^2\), with
\(D_0=I\) and \(D_1=D\).  This proves the needed null-homotopy
directly, including both phase seams.  In particular
\(D_v(a,b)G_0\) contracts the comparison matrix in
\eqref{h24:ret:eq:full-frame} to one constant positive matrix.

\begin{proposition}
The stable tangential framing induced by the actual normal four-frame
of \(H^{-1}(y_0)\) agrees, up to homotopy and the orientation sign
\eqref{h24:ret:eq:sign}, with the Lie framing
\((\partial_a,\partial_b)\) on the native two-torus.
\end{proposition}
\begin{proof}
In the coordinate chart, the stable tangent trivialization is the
isomorphism
\[
 \mathbb R^4\oplus TM\longrightarrow T\mathbb C^3|_M,
 \qquad (v,V)\longmapsto
 v_1N_\ell+v_2N_\phi+v_3N_1+v_4N_2+V.
\]
Using the intrinsic frame \((-\partial_a,\partial_b)\) on \(TM\),
its matrix is precisely \eqref{h24:ret:eq:full-frame}.  The homotopy
\eqref{h24:ret:eq:D-homotopy} therefore makes this comparison constant.
A constant positive matrix changes no homotopy class of stable
framing: Gram--Schmidt deforms it to an element of \(SO(6)\), and
an orthogonal matrix of determinant one can be reduced to the identity
by finitely many oriented plane rotations, each joined to the
identity by varying its angle.  This proves the asserted comparison
within the chart.

It remains to compare the chart with the ambient sphere framing.
Pull back the usual stable tangent trivialization
\(TS^6\oplus\mathbb R\cong\mathbb R^7\) under the fixed based diffeomorphism
\(X\cong S^6\).  On the contractible \(w\)-ball containing \(M\),
it differs from the coordinate trivialization plus one real line by
a continuous map into \(\operatorname{GL}(7,\mathbb R)\).  Contracting the ball to its
center contracts this comparison map to a constant.  Thus its
restriction to \(M\) contributes no varying stable framing class.
If the chosen sphere orientation disagrees with the complex
orientation, the remaining constant has negative determinant and
reverses the resulting framed bordism class.  Both that sign and the
explicit sign \eqref{h24:ret:eq:sign} disappear in the order-two class
computed below.  This argument compares actual ambient
trivializations on a contractible chart; it does not suppose that
\(\mathbb C^3\) globally trivializes \(TX\).
\end{proof}

\paragraph{Nonzero framed bordism and the exact unstable class}

We give the nonvanishing test for the resulting torus explicitly.
A stable framing gives a spin structure by lifting its oriented frame
to the trivial \(\operatorname{Spin}\)-bundle.  On the Lie-framed
circle this is the disconnected double cover, whereas the spin
structure induced from the boundary of a framed disc is the
connected double cover: as the boundary tangent turns once, its
\(2\pi\)-rotation lifts from \(1\) to \(-1\) in
\(\operatorname{Spin}(2)\).  We call the former nonbounding and the
latter bounding.

For a closed oriented spin surface \(S\), the spin quadratic
refinement
\[
 q_S:H_1(S;\mathbb F_2)\longrightarrow\mathbb F_2
\]
takes a class represented by an embedded circle to zero or one
according to whether the induced spin structure, using its normal
line in the surface, is bounding or nonbounding.  It satisfies
\begin{equation}
 q_S(x+y)=q_S(x)+q_S(y)+(x\mathbin{\cdot}y).
 \label{h24:ret:eq:quadratic}
\end{equation}
This is the usual spin quadratic form: resolving transverse
intersections adds one tangent turn per intersection, and the
\(2\pi\)-spin-lift calculation above gives the last term.  The
well-defined correspondence between spin structures and such
quadratic refinements is given in
\cite[Theorem 3.2 and Remark 3.3]{h24:retbib:DebrayGunningham}; its original
surface proof is \cite{h24:retbib:Johnson}.

For a symplectic basis \((e_i,f_i)_{i=1}^g\), put
\begin{equation}
 \operatorname{Arf}(S)=\sum_{i=1}^gq_S(e_i)q_S(f_i)\in\mathbb F_2.
 \label{h24:ret:eq:Arf}
\end{equation}
The basis independence can be seen without selecting a preferred
basis: factor the finite sum over the symplectic planes to obtain
\[
 \sum_{x\in H_1(S;\mathbb F_2)}(-1)^{q_S(x)}
   =\prod_{i=1}^g
     \sum_{a,b\in\mathbb F_2}(-1)^{a q_S(e_i)+b q_S(f_i)+ab}
   =2^g(-1)^{\sum_i q_S(e_i)q_S(f_i)}.
\]
The left side has no choice of basis, proving the assertion.

Here is also the bordism-invariance argument needed for the present
obstruction.  A spin bordism has a handle decomposition.  An index
zero or three handle creates or removes a spin sphere, whose first
homology and Arf invariant are zero.  An index one handle either
connects two components, in which case the first-homology quadratic
forms add orthogonally, or increases the genus of one component.
In the latter case the new symplectic pair contains the meridian
of the handle.  That meridian bounds the cocore disc with its
induced spin structure, so its quadratic value is zero; the pair
contributes zero to \eqref{h24:ret:eq:Arf}, regardless of the value on its
dual.  Existing cycles can be represented outside the attaching
discs, so their quadratic values are preserved.  An index two
handle is the reverse of an index one handle.  Therefore each
handle preserves the total Arf invariant of the boundary, and a
spin boundary has Arf invariant zero.  This also proves that a
stably framed boundary has Arf invariant zero, because forgetting
the framing gives a spin bordism.  The full spin bordism statement
is recorded in \cite[Theorem 3.6]{h24:retbib:DebrayGunningham}.

For the intrinsic Lie framing on
\(M=(\mathbb R/(2\pi\mathbb Z))^2\), take
\[
 e=[a\mapsto(a,0)],\qquad f=[b\mapsto(0,b)].
\]
Both the tangent and normal vectors in the surface are constant in
the Lie frame along either circle.  Thus both circles have the
nonbounding spin structure, and
\begin{equation}
 e\mathbin{\cdot}f=1,\qquad q_M(e)=q_M(f)=1,
 \qquad q_M(xe+yf)=x+xy+y,
 \qquad \operatorname{Arf}(M)=1.
 \label{h24:ret:eq:actual-Arf}
\end{equation}
Reversing the first Lie tangent vector changes the orientation but
does not change these mod-two values.  In particular the
specifically framed inverse image computed above is not a framed
boundary.

To identify the element, let \(\eta\in\pi_1^{\rm S}\) denote the
stable class of the complex Hopf map.  Its Pontryagin--Thom inverse
image is the circle with its Lie framing.  This can be verified
directly on the Hopf map \(S^3\to\mathbb{CP}^1\): near the fibre
\((z_1,z_2)=(e^{ia},0)\), the target coordinate is \(z_2/z_1\), so
the two horizontal normal vectors are \((0,e^{ia})\) and
\((0,ie^{ia})\).  Add the tangent vector \((ie^{ia},0)\) and the
outward radial vector \((e^{ia},0)\).  In \(\mathbb C^2\) the resulting
full frame makes one complete rotation in each of its two real
coordinate planes.  This loop is contractible explicitly: for the
fixed real-linear involution \(J(z_1,z_2)=(z_1,\overline z_2)\),
\[
 J\operatorname{diag}(e^{ia},e^{ia})J=\operatorname{diag}(e^{ia},e^{-ia})=K_1(-a),
\]
and \(JK_v(-a)J\) contracts the original loop to the identity using
the periodic homotopy already constructed.  Thus the induced stable
tangent frame on the circle is the Lie frame.  This is the basic
framed-circle identification in \cite{h24:retbib:Pontryagin1955}.

The product of two normally framed inverse images represents the
product in stable homotopy: its tubular collapse is the smash of
the two tubular collapse maps, as follows immediately from the
product tubular neighborhood \(\nu_1\times\nu_2\) and its ordered
normal coordinates.  Consequently the product of the two Lie
circles represents \(\eta^2\).  Equation \eqref{h24:ret:eq:actual-Arf}
proves \(\eta^2\ne0\) by an explicit bordism invariant.  The
stable group \(\pi_2^{\rm S}=\mathbb Z/2\) is Pontryagin's computation
\cite{h24:retbib:Pontryagin1950,h24:retbib:Pontryagin1955}; the stable table in
\cite[pp.~384--385]{h24:retbib:Hatcher} gives the same group.  Thus the sign from
\eqref{h24:ret:eq:sign} satisfies \(-\eta^2=\eta^2\).

\begin{proof}[Proof of the theorem]
The globally continuous quotient \(H\) is smooth on a neighborhood
of its compact inverse image \(M\), by
\eqref{h24:ret:eq:exact-u} and \eqref{h24:ret:eq:four-cell}. First make it constant
near the source basepoint, before choosing any approximation bound.
Choose a target chart \(\kappa:O\to\mathbb R^4\) about \(*\) with
\(\kappa(*)=0\), convex ball image, and \(y_0\notin O\).
Choose a source neighborhood \(U\) of \(x_{\rm base}\), disjoint
from \(M\), with \(H(U)\subset O\). Let \(\beta\) be a smooth
function supported in \(U\), equal to one near \(x_{\rm base}\),
with \(0\leq\beta\leq1\). On \(U\) use the based homotopy
\[
 H_s(x)=\kappa^{-1}\bigl((1-s\beta(x))\kappa(H(x))\bigr),
 \qquad 0\leq s\leq1,
\]
and leave \(H\) unchanged off \(U\). Its endpoint \(H_b\) is
smooth and constant near the basepoint; it agrees with \(H\)
near \(M\) and has no new inverse image of \(y_0\).

Choose one closed tubular neighborhood \(U_1\) of \(M\), with
\(M\subset\operatorname{int}U_1\), contained in the unchanged
smooth region and disjoint from the basepoint neighborhood.
Compactness gives the positive bound
\[
 d_0=\min_{x\in X\setminus\operatorname{int}U_1}
             d\bigl(H_b(x),y_0\bigr)>0.
\]
Relative smooth approximation, fixing the same closed set \(U_1\)
and a neighborhood of the basepoint, produces a smooth based map
\(H_{\rm sm}\) whose error from \(H_b\) is less than \(d_0/2\)
on its complement. Choose the approximation closer also than a
fixed target injectivity radius; the short target geodesics give a
based homotopy. Thus the exact inverse image and its differential
are unchanged on \(U_1\), and the complement acquires no new
inverse image. The relative Pontryagin construction applies as in
\cite[\S7, Theorem B and Lemma 4, pp.~43, 48--50]{h24:retbib:Milnor1965}.

Pontryagin--Thom therefore identifies the stabilization of
\([H]\) with the bordism class of \(M\) carrying the four normal
vectors \eqref{h24:ret:eq:normal-first}--\eqref{h24:ret:eq:normal-u}.  The comparison
\eqref{h24:ret:eq:full-frame} and its explicit contraction identify that
class with the Lie torus, with the harmless order-two orientation
sign already computed.  Hence
\[
 E^\infty[H]=\eta^2\ne0.
\]
Freudenthal suspension is an isomorphism
\[
 E:\pi_{n+2}(S^n)\longrightarrow\pi_{n+3}(S^{n+1})
 \qquad(n\geq4),
\]
because \(n+2\leq2n-2\).  Therefore stabilization from
\(\pi_6(S^4)\) is an isomorphism onto \(\pi_2^{\rm S}\).
The composite \(\eta_4\circ\eta_5:S^6\to S^4\) stabilizes to
\(\eta^2\), so injectivity gives the exact equality
\eqref{h24:ret:eq:main-class}.  Finally, a based null-homotopy of \(q_2\)
would postcompose with \(p_3\) to a based null-homotopy of \(H\),
contradicting \eqref{h24:ret:eq:main-class}.
\end{proof}

\paragraph{Exact scope of the comparison with the angle-forgetting map.}
The map here retains the angular coordinate before collapsing the
coordinate two-skeleton.  Its regular inverse image is the specific
normally four-framed \(T^2\) above.  The angle-forgetting
postcomposition has a normally three-framed \(T^3\) inverse image.
The calculation above is independent of its class and supplies an
actual obstruction to null-homotopy of the original \(q_2\), even
when that other postcomposition is null-homotopic.  The exact
morphism carrying this obstruction is
\[
 (p_3)_*:[X,\mathcal A]_*\longrightarrow[X,S^4]_*,\qquad
 [q_2]\longmapsto[H]=\eta_4\eta_5.
\]
No conclusion about an arbitrary earlier quotient follows merely
from the diffeomorphism type of one inverse image; the specified
four-framing is what proves this value.

\begin{corollary}[The exact order-two quotient detected by the original cusp map]
\label{h24:ret:q2-detected-quotient}
Let \(d:(S^6,*)\to(X,x_{\mathrm{base}})\) be the fixed based
diffeomorphism, and put
\[
\alpha=[q_2\circ d]\in\pi_6(\mathcal A,v_0).
\]
The original top-cell postcomposition induces a surjective homomorphism
\[
\varepsilon=(p_3)_*:\pi_6(\mathcal A,v_0)
  \longrightarrow\pi_6(S^4,*)\cong\mathbb Z/2,
\qquad \varepsilon(\alpha)=\eta_4\circ\eta_5.
\]
It gives the exact quotient isomorphism
\[
\pi_6(\mathcal A,v_0)/\ker\varepsilon
 \xrightarrow{\ [\beta]\mapsto\varepsilon(\beta)\ }
\mathbb Z/2.
\]
Every odd multiple of \(\alpha\) is nonzero. Its order is either
infinite or an even positive integer; this calculation does not determine
that order or assert a splitting of the displayed quotient.
\end{corollary}
\begin{proof}
Postcomposition distributes over the pinch sum on the source sphere,
so \((p_3)_*\) is a group homomorphism. Theorem~\ref{h24:ret:main}
gives \((p_3)_*(\alpha)=[p_3q_2d]=\eta_4\eta_5\ne0\).
This is the nonzero element of the cyclic group of order two, proving
surjectivity. The first isomorphism theorem gives the stated quotient
with exactly the displayed map. For every integer \(n\),
\[
\varepsilon(n\alpha)=(n\bmod2)\eta_4\eta_5.
\]
An odd multiple therefore has nonzero image and cannot vanish.
If \(m\alpha=0\), applying \(\varepsilon\) forces \(m\)
to be even. These arguments prove all the assertions without inferring
the exact order of \(\alpha\).
\end{proof}

% Staged proof only. No canonical source has been edited by this worker.
% Uses the already proved retained-angle theorem ret:main.

\subsubsection{The retained cusp class as a three-dimensional bundle inside the triality carrier}
\label{h24:rettri:section}

Retain the exact based map $H=p_3q_2:X\to S^4$, the fixed based
diffeomorphism $d:S^6\to X$, the original cutoff $\epsilon$, the ordered
native coordinates $(r,\phi,u_1,u_2)$, and every period appearing in
\eqref{h24:ret:eq:original-periods}. Theorem~\ref{h24:ret:main} has proved
$[H d]=\eta_4\eta_5\ne0$ in $\pi_6(S^4)\cong\mathbb Z/2$.
We construct an exact bundle-valued image of this specific map and then
its action on the original three Peirce sectors. A quaternion subalgebra
and its ordered basis are explicit choices in this construction.

\paragraph{The literal Hopf pullback and its transition.}
Fix $\mathbb H=\mathbb R\langle1,i,j,k\rangle$ with
$i^2=j^2=k^2=ijk=-1$. Identify the ordered open four-cell of the target
with $\mathbb H$ by
\begin{equation}
 z=-\cot(\pi r/\epsilon)-\cot(\phi/2)i
       -\cot(\pi u_1)j-\cot(\pi u_2)k.
 \label{h24:rettri:z}
\end{equation}
This is exactly the ordered coordinate map
\eqref{h24:ret:eq:four-cell}, followed by the real linear map carrying
the four standard basis vectors to $(1,i,j,k)$. Each of the four
coordinate derivatives in the original open cell is positive.
The based homeomorphism
\begin{equation}
 \chi:S^4=\mathbb H\cup\{\infty\}\longrightarrow\mathbb H\mathbb P^1,
 \qquad \chi(z)=[1:z],\qquad \chi(*)=[0:1]
 \label{h24:rettri:chi}
\end{equation}
uses right quaternionic lines. The orientation on this presentation of
$\mathbb H\mathbb P^1$ is the one induced by the ordered affine basis
$(1,i,j,k)$. Thus no orientation convention on another presentation of
$\mathbb H\mathbb P^1$ is silently substituted.

Let
\[
 S^7=\{(v_0,v_1)\in\mathbb H^2:|v_0|^2+|v_1|^2=1\},\qquad
 \pi(v_0,v_1)=[v_0:v_1].
\]
The right action is $(v_0,v_1)u=(v_0u,v_1u)$ for
$u\in\operatorname{Sp}(1)=\{u\in\mathbb H:|u|=1\}$. It is free:
at least one $v_j$ is nonzero, and $v_ju=v_j$ implies $u=1$.
Two unit vectors on the same right line differ by a unique unit
quaternion. Hence $\pi$ is the quaternionic Hopf principal bundle.
Its local sections in the affine coordinates $z=v_1v_0^{-1}$ and
$w=v_0v_1^{-1}$ are
\begin{equation}
 s_0(z)=\frac{(1,z)}{\sqrt{1+|z|^2}},\qquad
 s_\infty(w)=\frac{(w,1)}{\sqrt{1+|w|^2}}.
 \label{h24:rettri:sections}
\end{equation}
On $z\ne0$, the coordinate change is $w=z^{-1}$ and
\begin{equation}
 s_\infty(z^{-1})=s_0(z)\,g(z),\qquad
 g(z)=\frac{\overline z}{|z|}.
 \label{h24:rettri:transition}
\end{equation}
Indeed, multiplying the right side gives
$(\overline z/|z|,|z|)/\sqrt{1+|z|^2}$, and the same expression
results from the left side because
$\sqrt{1+|z^{-1}|^2}=\sqrt{1+|z|^2}/|z|$.
This calculation retains both the inverse and the conjugation.

Define
\begin{equation}
 P_H=\{(x,v)\in X\times S^7:\pi(v)=\chi(H(x))\}
       \longrightarrow X,
 \qquad (x,v)u=(x,vu).
 \label{h24:rettri:P}
\end{equation}
The preimages under $\chi H$ of the two affine charts cover $X$.
Equations~\eqref{h24:rettri:sections}--\eqref{h24:rettri:transition}
are local sections and transitions for \eqref{h24:rettri:P} after
composition with $\chi H$. In particular, wherever both charts apply,
the transition is the exact $g(z)$ from \eqref{h24:rettri:transition},
with $z$ given by the original four-coordinate expression
\eqref{h24:rettri:z}. The angle $\phi$ has not been dropped.

\begin{theorem}
\label{h24:rettri:bundleclass}
The principal bundle $P_H$ is nontrivial. Its clutching class is the
nonzero element
\begin{equation}
 c(P_H)=\partial[\chi H d]\in\pi_5(\operatorname{Sp}(1))
                  \cong\mathbb Z/2,
 \label{h24:rettri:clutch}
\end{equation}
where $\partial$ is the boundary of the specified right Hopf fibration.
The associated real three-plane bundle
\begin{equation}
 \mathcal A_H=P_H\times_{\operatorname{Ad}}\operatorname{Im}\mathbb H,
 \qquad \operatorname{Ad}_u(v)=uv\overline u,
 \label{h24:rettri:rankthree}
\end{equation}
is nontrivial, has nonzero order-two clutching class in
$\pi_5(SO(3))$, and admits no nowhere-zero continuous section.
Its oriented orthonormal frame bundle has $P_H$ as its unique spin lift
up to isomorphism over $X$. The maps from the retained homotopy group
to these two clutching groups are isomorphisms:
\begin{equation}
 \pi_6(S^4)\xrightarrow{\chi_*}\pi_6(\mathbb H\mathbb P^1)
 \xrightarrow{\partial}\pi_5(\operatorname{Sp}(1))
 \xrightarrow{\operatorname{Ad}_*}\pi_5(SO(3)).
 \label{h24:rettri:exact-detection}
\end{equation}
\end{theorem}

\begin{proof}
The first map is an isomorphism because $\chi$ is a based
homeomorphism. In the long exact sequence of the displayed Hopf
bundle, the segment is
\[
 \pi_6(S^7)\longrightarrow\pi_6(\mathbb H\mathbb P^1)
 \xrightarrow{\partial}\pi_5(S^3)\longrightarrow\pi_5(S^7).
\]
Both outside groups vanish: cellular approximation homotopes a map
from a sphere of dimension less than seven into the zero-cell of the
CW structure on $S^7$. Exactness therefore makes $\partial$ an
isomorphism. To fix its clutching convention, write the oriented
source $S^6=D^6_+\cup D^6_-$, give the equator the boundary orientation
of $D^6_+$, and choose local lifts $\sigma_\pm$ of the pulled-back
bundle. On the equator define $g_{+-}$ by
$\sigma_- = \sigma_+ g_{+-}$. Choose the values of the two lifts
equal at its fixed basepoint. We use the convention
$\partial[\chi H d]=[g_{+-}]$. Changing disk lifts changes this map
by null-homotopic factors and does not change its homotopy class.
This is the defining boundary construction of the fibration exact
sequence. The already proved nonzero class $[H d]$ therefore gives
the nonzero order-two class in \eqref{h24:rettri:clutch}.

One can also see nontriviality directly. A section of $P_H$ gives a
map $\widetilde h:S^6\to S^7$ with
$\pi\widetilde h=\chi H d$. Since $\pi_6(S^7)=0$, this would make
$\chi H d$ null-homotopic, contrary to Theorem~\ref{h24:ret:main}.
All spaces at issue are simply connected, so based and free
null-homotopy agree here.

Conjugation by a unit quaternion preserves $\operatorname{Im}\mathbb H$
and its Euclidean norm. Its determinant is $+1$ by continuity in
$u$ and its value at $u=1$. The kernel consists of the real unit
quaternions $\{1,-1\}$: a quaternion commuting with $i,j,k$ is real.
For a unit imaginary quaternion $n$ and $v\in\operatorname{Im}\mathbb H$,
write $v=v_\parallel+v_\perp$ relative to the line $\mathbb R n$.
Direct multiplication with $u=\cos\alpha+n\sin\alpha$ gives
\begin{equation}
 u v\overline u=v_\parallel+\cos(2\alpha)v_\perp
                     +\sin(2\alpha)(n\times v_\perp),
 \label{h24:rettri:Rodrigues}
\end{equation}
where $n\times v=(nv-vn)/2$. Every orientation-preserving
orthogonal map of a real three-space fixes an axis and rotates its
orthogonal plane: its real odd-degree characteristic polynomial has
a real eigenvalue, norm preservation gives $\pm1$, and determinant
$+1$ forces an eigenvalue $+1$. Thus
\eqref{h24:rettri:Rodrigues} proves that
$\operatorname{Ad}:S^3\to SO(3)$ is onto with kernel $\{\pm1\}$.
It is the two-sheeted covering and induces an isomorphism on
$\pi_5$. Its value on \eqref{h24:rettri:clutch} is the clutching
class of \eqref{h24:rettri:rankthree}, proving nontriviality and
order two.

The bundle $P_H\to P_H/\{\pm1\}$ is the spin double covering of
the oriented orthonormal frames of $\mathcal A_H$. If that frame
bundle had a section, its inverse image would be a double covering
of the simply connected space $X\cong S^6$. This covering is
trivial, so one of its two sections would give a section of $P_H$,
again a contradiction. Any two lifts of an $SO(3)$ bundle differ
by a principal $\mathbb Z/2$ bundle: locally the ratio between two
lifts of each transition lies in the central kernel $\{\pm1\}$,
and the cocycle identity gives a $\mathbb Z/2$ cocycle. Every such
bundle on the simply connected $X$ is trivial. This proves the
claimed uniqueness of the spin lift.

If $\mathcal A_H$ had a nowhere-zero section, divide that section
by its positive length to obtain a unit section. Orthogonal
complement gives $\mathcal A_H\cong\mathbb R\oplus\mathcal B$
for an oriented two-plane bundle $\mathcal B$. A two-plane
bundle on $S^6$ has clutching map $S^5\to SO(2)=S^1$.
Every such map lifts to $\mathbb R$ and is null-homotopic, so
$\mathcal B$ is trivial. This would trivialize $\mathcal A_H$,
contradicting its nonzero clutching class.
\end{proof}

The fibrewise scalar product, orientation and cross product of
$\operatorname{Im}\mathbb H$ descend to $\mathcal A_H$.
For example,
\[
 (uv\overline u)(uw\overline u)=u(vw)\overline u
\]
by associativity in $\mathbb H$, so conjugation preserves
$v\times w=(vw-wv)/2$. The induced fibrewise Lie bracket is
exactly $[v,w]=2(v\times w)$. No factor of two is discarded.

\paragraph{An explicit subgroup of the three-sector triality group.}
Fix the quaternion subalgebra above inside
$\mathbb O=\mathbb H\oplus\mathbb H\ell$. Its ordered real basis is
\[
 (1,i,j,k,\ell,i\ell,j\ell,k\ell),
\]
and its multiplication is
\begin{equation}
 (a,b)(c,d)=(ac-\overline d b,\;da+b\overline c),\qquad
 K(a,b)=\overline{(a,b)}=(\overline a,-b).
 \label{h24:rettri:CD}
\end{equation}
This is the literal convention of Yokota, \S1.10. In his cyclic
related-triple presentation the group is
\begin{equation}
 \operatorname{RT}(\mathbb O)=
 \{(T_1,T_2,T_3)\in SO(\mathbb O)^3:
       (T_1x)(T_2y)=K T_3 K(xy)\text{ for all }x,y\in\mathbb O\}.
 \label{h24:rettri:RT}
\end{equation}
Yokota, Theorem 1.16.2, identifies this group with
$\operatorname{Spin}(8)$; Theorem 2.7.1 \cite{h24:Yokota2009} gives its diagonal-fixing
Jordan action. In particular the three projection representations
in \eqref{h24:rettri:RT} are the vector and two half-spin sectors
with that fixed marking.

For each $u\in\operatorname{Sp}(1)$ define
\begin{equation}
 \begin{aligned}
 B_u(a,b)&=(u a\overline u,b),\\
 A_u(a,b)&=(u a,b\overline u),\\
 D_u(a,b)&=(a\overline u,b\overline u)=K A_u K(a,b).
 \end{aligned}
 \label{h24:rettri:threeactions}
\end{equation}

\begin{theorem}
\label{h24:rettri:subgroup}
The formula
\begin{equation}
 \rho:\operatorname{Sp}(1)\longrightarrow\operatorname{Spin}(8)
      =\operatorname{RT}(\mathbb O),\qquad
 \rho(u)=(B_u,A_u,D_u)
 \label{h24:rettri:rho}
\end{equation}
is an injective continuous group homomorphism. It satisfies
\begin{equation}
 (B_u x)(A_u y)=A_u(xy),\qquad
 \operatorname{Re}(((B_u x)(A_u y))(D_u z))
           =\operatorname{Re}((xy)z)
 \label{h24:rettri:product}
\end{equation}
for all octonions $x,y,z$ and every unit quaternion $u$.
\end{theorem}

\begin{proof}
All three maps are real linear norm isometries by the quaternionic
norm identity. Their determinants are $+1$: they vary continuously
with $u$ on the connected $S^3$ and equal the identity at $u=1$.
For $u,v\in\operatorname{Sp}(1)$, associativity in $\mathbb H$
and $\overline{uv}=\overline v\,\overline u$ give
\[
 B_uB_v=B_{uv},\qquad A_uA_v=A_{uv},\qquad D_uD_v=D_{uv}.
\]
For example, the second coordinate of $A_uA_v(a,b)$ is
$(b\overline v)\overline u=b\overline{uv}$, retaining the order.

For $x=(a,b)$ and $y=(c,d)$, expansion of
\eqref{h24:rettri:CD} gives
\begin{align*}
 (B_u x)(A_u y)
 &= (u a\overline u,b)(u c,d\overline u)\\
 &=\bigl((u a\overline u)(u c)
          -\overline{d\overline u}\,b,
          (d\overline u)(u a\overline u)
             +b\overline{u c}\bigr)\\
 &=\bigl(u(ac-\overline d b),
                    (da+b\overline c)\overline u\bigr)
 =A_u(xy).
\end{align*}
This uses only quaternionic associativity in the displayed
components. Since $A_u=K D_u K$, it proves membership in
\eqref{h24:rettri:RT}. If $\rho(u)$ is the identity, evaluate
$A_u$ at $(1,0)$ to get $(u,0)=(1,0)$, so $u=1$.
Thus $\rho$ is injective.

The other two cyclic pair identities, with every conjugation
retained, are
\begin{align*}
 (A_u x)(D_u y)
 &=(u a,b\overline u)(c\overline u,d\overline u)\\
 &=\bigl(u(ac-\overline d b)\overline u,\;da+b\overline c\bigr)
 =B_u(xy)=K B_u K(xy),\\
 (D_u x)(B_u y)
 &=(a\overline u,b\overline u)(u c\overline u,d)\\
 &=\bigl((ac-\overline d b)\overline u,
                    (da+b\overline c)\overline u\bigr)
 =D_u(xy)=K A_u K(xy).
\end{align*}

Put $\langle v,w\rangle=\operatorname{Re}(v\overline w)$ on
$\mathbb O$. Then $\operatorname{Re}(vz)=\langle v,Kz\rangle$.
The already proved product relation and $K D_u=A_u K$ imply
\begin{align*}
 \operatorname{Re}(((B_u x)(A_u y))(D_u z))
 &=\langle A_u(xy),K D_u z\rangle\\
 &=\langle A_u(xy),A_u Kz\rangle
 =\langle xy,Kz\rangle
 =\operatorname{Re}((xy)z).
\end{align*}
No reassociation of a non-real octonionic triple product was used.
\end{proof}

\paragraph{The exact original Jordan coordinates and all three bundles.}
Write the original Hermitian matrix with upper entries $(p,q,r)$ as
\begin{equation}
 J(\lambda;p,q,r)=
 \begin{pmatrix}
 \lambda_1&p&q\\
 \overline p&\lambda_2&r\\
 \overline q&\overline r&\lambda_3
 \end{pmatrix},\qquad
 (q_1,q_2,q_3)=(r,\overline q,p).
 \label{h24:rettri:J}
\end{equation}
Thus $q_2$ here is the cyclic Jordan coordinate, not the previously
defined cusp quotient map; only the explicitly displayed ordered
triple in this paragraph uses this coordinate notation.
Under \eqref{h24:rettri:rho}, the cyclic coordinates transform by
\[
 (q_1,q_2,q_3)\longmapsto(B_u q_1,A_u q_2,D_u q_3).
\]
Since $K A_u K=D_u$, the action in the original upper entries is
\begin{equation}
 (r,q,p)\longmapsto(B_u r,D_u q,D_u p),\qquad
 J(\lambda;p,q,r)\longmapsto
       J(\lambda;D_u p,D_u q,B_u r).
 \label{h24:rettri:original-action}
\end{equation}
In particular this is an action on the three specific spaces
$\mathfrak h_{\gamma_1},\mathfrak h_{\gamma_2},
\mathfrak h_{\gamma_3}$ of
\eqref{h24:eq:octonionic-root-spaces}, in their original order.

The determinant and trace form are
\begin{align}
 \det J&=\lambda_1\lambda_2\lambda_3
       -\lambda_1|r|^2-\lambda_2|q|^2-\lambda_3|p|^2
       +2\operatorname{Re}((r\overline q)p),
 \label{h24:rettri:det}\\
 \operatorname{tr}(J^2)&=\lambda_1^2+\lambda_2^2+\lambda_3^2
                       +2(|p|^2+|q|^2+|r|^2).
 \label{h24:rettri:metric}
\end{align}
Theorem~\ref{h24:rettri:subgroup} and norm preservation prove that
\eqref{h24:rettri:original-action} preserves every term, with its
displayed sign and coefficient. Its being a Jordan automorphism
also follows directly from Yokota's diagonal-fixing action in
Theorem 2.7.1, with exactly the identification
\eqref{h24:rettri:J}. Equivalently, the Peirce multiplication rules
$F_i(x)\circ F_i(y)=\langle x,y\rangle(E_{i+1}+E_{i+2})$ and
$F_i(x)\circ F_{i+1}(y)=\tfrac12 F_{i+2}(\overline{xy})$,
where indices are cyclic modulo three, are preserved by the norm
identity and the three cyclic pair identities proved above.
Here $F_1(x)$ has upper entry $x$ at $(2,3)$,
$F_2(x)$ has upper entry $\overline x$ at $(1,3)$, and
$F_3(x)$ has upper entry $x$ at $(1,2)$; all other entries are
zero apart from their Hermitian conjugates. Direct matrix
multiplication gives the displayed rules, with coefficient
$1/2$ from the Jordan product. Finally
$E_i\circ E_j=\delta_{ij}E_i$, $E_i\circ F_i(x)=0$,
and $E_i\circ F_j(x)=F_j(x)/2$ for $i\ne j$, and these rules
are preserved because all $E_i$ are fixed. The displayed
identities specify the Jordan product on every pair by real
bilinearity.

For a right principal bundle our associated-bundle convention is
$(p u,v)\sim(p,\sigma(u)v)$. Let
\[
 \mathcal L_H=P_H\times_{u:a\mapsto u a}\mathbb H.
\]
This is the tautological right quaternionic line pulled back by
$\chi H$: the fibre map $[(x,v),a]\mapsto (x,v a)$ is
well-defined, right quaternionic linear and invertible onto the
line $\chi H(x)$. Its unit sphere bundle is $P_H$, by
$p\mapsto[p,1]$. Consequently it is nontrivial as a quaternionic
line. It is also nontrivial as a real four-plane bundle: if it had
a nowhere-zero real section, dividing by length would give a
section of its unit sphere bundle and hence of $P_H$.

The original ordered Peirce representation produces the real
rank-24 bundle
\begin{equation}
 \mathcal W_H=\mathcal V_{r,H}\oplus\mathcal V_{q,H}
                              \oplus\mathcal V_{p,H},
 \quad
 \mathcal V_{r,H}=P_H\times_B\mathbb O,
 \quad
 \mathcal V_{q,H}=\mathcal V_{p,H}=P_H\times_D\mathbb O.
 \label{h24:rettri:W}
\end{equation}
The equality between the last two names specifies isomorphic
representations with two separately retained labelled copies.
Their explicit decompositions are
\begin{equation}
 \begin{aligned}
 \mathcal V_{r,H}&\cong
       \mathbb R\oplus\mathcal A_H\oplus\mathbb R^4,
       &&(\operatorname{Re}a,\operatorname{Im}a,b),\\
 \mathcal V_{q,H}&\cong\mathcal L_H\oplus\mathcal L_H,
       &&(a,b)\longmapsto(\overline a,\overline b),\\
 \mathcal V_{p,H}&\cong\mathcal L_H\oplus\mathcal L_H,
       &&(a,b)\longmapsto(\overline a,\overline b).
 \end{aligned}
 \label{h24:rettri:splittings}
\end{equation}
For the last two maps,
$\overline{a\overline u}=u\overline a$ proves equivariance,
and similarly for $b$. Each quaternionic conjugation has real
determinant $-1$ in $(1,i,j,k)$, so their ordered pair has
determinant $+1$. The first line keeps the actual coordinate order
$(1,\operatorname{Im}\mathbb H,\mathbb H\ell)$; moving the
one-dimensional first factor across the imaginary three-space
would contribute the sign $(-1)^3=-1$, and no such reorder is
being made. In cyclic coordinates the middle action is $A$;
the equivariant map to two left standard representations is
$(a,b)\mapsto(a,\overline b)$, of determinant $-1$. This is
why the original $\overline q$ is retained explicitly.

All fibrewise maps in \eqref{h24:rettri:product} and
\eqref{h24:rettri:det} descend to the associated bundles: replacing
a local frame $p$ by $p u$ changes inputs by precisely the three
displayed actions, whose product and scalar identities have been
proved. In particular
\[
 \mathbb R^3\oplus\mathcal W_H
       =P_H\times_\rho J_3(\mathbb O)
\]
is a bundle of Jordan algebras with its three fixed diagonal
idempotents, its full product, trace form and determinant. The
three-sector structure on \eqref{h24:rettri:W} includes the
explicit nonzero invariant trilinear map
$(r,q,p)\mapsto2\operatorname{Re}((r\overline q)p)$;
it is nonzero because its value at $r=q=p=1$ is exactly $2$.

\paragraph{What extension to the full triality group forgets.}
Form the principal bundle
\begin{equation}
 Q_H=P_H\times_\rho\operatorname{Spin}(8)\longrightarrow X,
 \quad (p u,s)\sim(p,\rho(u)s).
 \label{h24:rettri:Q}
\end{equation}
It contains the reduction $P_H$ by $p\mapsto[p,1]$. This is
equivariant because $[p u,1]=[p,\rho(u)]$.

\begin{proposition}
\label{h24:rettri:forgetting}
The principal bundle $Q_H$ is trivial. All three rank-eight
bundles in \eqref{h24:rettri:W} are therefore trivial as real
vector bundles, while their specified subbundles
$\mathcal A_H$ and $\mathcal L_H$ remain nontrivial.
After any chosen trivialization of $Q_H$, its reduction determines
a continuous map
\begin{equation}
 f_H:X\longrightarrow\operatorname{Spin}(8)/\rho(\operatorname{Sp}(1))
 \label{h24:rettri:reductionmap}
\end{equation}
whose pullback of the principal bundle
$\operatorname{Spin}(8)\to\operatorname{Spin}(8)/\rho(\operatorname{Sp}(1))$
is $P_H$. This map is non-null-homotopic.
\end{proposition}

\begin{proof}
Let $g:S^5\to\operatorname{Sp}(1)$ be a based clutching
representative for $P_H$ in the convention used in
\eqref{h24:rettri:clutch}. Its homotopy class has order two.
The transition $D_g$ on an eight-dimensional sector is conjugate
by $C(a,b)=(\overline a,\overline b)$ to
$\operatorname{diag}(L_g,L_g)$ on $\mathbb H^2$.
For $0\le t\le\pi/2$ put
\[
 R_t=\begin{pmatrix}\cos t&-\sin t\\
                       \sin t&\cos t\end{pmatrix}
       \quad\text{on }\mathbb H^2.
\]
This is an orthogonal path commuting with quaternionic scalar
right multiplication. Its conjugation carries
$\operatorname{diag}(L_g,I)$ at $t=0$ to
$\operatorname{diag}(I,L_g)$ at $t=\pi/2$. The maps into $SO(8)$
represented by these two matrices are therefore based homotopic.
Pointwise product of maps from a based sphere into a topological
group represents addition in its homotopy group: concatenation
in one loop coordinate and pointwise multiplication have the
same unit and satisfy interchange, yielding the same operation.
Consequently
\[
 [\operatorname{diag}(L_g,L_g)]
 =2[\operatorname{diag}(L_g,I)]
 =\iota_*(2[g])=0\quad\text{in }\pi_5(SO(8)),
\]
where $\iota(u)=\operatorname{diag}(L_u,I)$. Conjugating back by
$C$ gives $[D_g]=0$.

The third projection of \eqref{h24:rettri:RT} is a double covering
$\operatorname{Spin}(8)\to SO(8)$: this is the triality covering
in the third marking, with the same conjugation $K$ as in
\eqref{h24:rettri:RT}. Coverings induce isomorphisms on
$\pi_n$ for $n\ge2$ by unique lifting of based sphere maps and
their homotopies. The projection sends $[\rho g]$ to $[D_g]=0$,
so $[\rho g]=0$. A bundle obtained by gluing two trivial bundles
over disks with a null-homotopic clutching map is trivial: extend
that clutching map across a disk and change its local frame by
the extension. This proves triviality of $Q_H$ directly.
Every associated vector bundle of a trivial principal bundle is
trivial; applying the three projections gives the asserted
rank-eight trivializations. The earlier proofs of nontriviality
of $\mathcal A_H$ and $\mathcal L_H$ still apply to these actual
subbundles, so forgetting their inclusions loses information.

Choose an isomorphism $\tau:Q_H\to X\times G$, where
$G=\operatorname{Spin}(8)$, and put $K_0=\rho(\operatorname{Sp}(1))$.
For $p\in(P_H)_x$ write $\tau([p,1])=(x,a_p)$. Then
$a_{pu}=a_p\rho(u)$, so
$f_H(x)=a_p K_0$ is well-defined, independent of the choice of
$p$, and continuous in local sections. The explicit principal
bundle isomorphism to the pullback is
\[
 P_H\longrightarrow f_H^*(G\to G/K_0),\qquad
 p\longmapsto(x,a_p),
\]
using the identification of $\operatorname{Sp}(1)$ with $K_0$
by the proved injective map $\rho$. The formula is fibrewise
bijective and has inverse in each local trivialization.
If $f_H$ were null-homotopic, its pullback bundle would be
isomorphic to the pullback along a constant map and hence
trivial, contradicting Theorem~\ref{h24:rettri:bundleclass}.
Equivalently, in the fibration exact sequence its clutching
boundary is precisely $c(P_H)\ne0$.
\end{proof}

There is also an explicit map into the space of oriented
three-planes in the first eight-dimensional sector. If
$T_1:G\to SO(\mathbb O)$ is the first projection of
\eqref{h24:rettri:RT}, define
\begin{equation}
 \begin{aligned}
 \Gamma:G/K_0&\longrightarrow\operatorname{Gr}_3^+(\mathbb O),\\
 gK_0&\longmapsto
 \operatorname{span}_{\mathbb R}^{+}
       \bigl(T_1(g)i,T_1(g)j,T_1(g)k\bigr).
 \end{aligned}
 \label{h24:rettri:Grassmann}
\end{equation}
Here the superscript records the orientation of the displayed
ordered basis. This formula is well-defined: multiplication of
$g$ on the right by $\rho(u)$ changes that ordered basis by
the orientation-preserving map $\operatorname{Ad}_u$ of
$\operatorname{Im}\mathbb H$, because $T_1\rho(u)=B_u$.
Continuity follows from local sections of $G\to G/K_0$.
Let $\mathcal T_3$ be the tautological oriented three-plane
bundle on this Grassmannian. The explicit fibre map
\[
 G\times_{K_0}\operatorname{Im}\mathbb H
       \longrightarrow\Gamma^*\mathcal T_3,\qquad
 [g,v]\longmapsto(gK_0,T_1(g)v)
\]
is well-defined since
$T_1(g\rho(u))v=T_1(g)\operatorname{Ad}_u v$.
It is a linear isometry on every fibre, hence a bundle
isomorphism. Pulling back by \eqref{h24:rettri:reductionmap}
therefore gives
\[
 (\Gamma f_H)^*\mathcal T_3\cong\mathcal A_H.
\]
In particular $\Gamma f_H:X\to\operatorname{Gr}_3^+(\mathbb O)$
is non-null-homotopic, since pullback along a null-homotopic
map would make the tautological bundle trivial. This locates
the three-plane in the actual first triality sector, with its
orientation and its nonzero clutching class retained.

Thus the exact bridge is the sequence of explicit constructions
\[
 H\ \longmapsto\ P_H\ \longmapsto\ \mathcal A_H
 \ \longmapsto\ (\mathcal W_H,\rho\text{-reduction},
                   2\operatorname{Re}((r\overline q)p)).
\]
The first two steps retain the entire nonzero order-two class
by \eqref{h24:rettri:exact-detection}. In the last step the
specified reduction, or equivalently the distinguished
three-plane with its spin frame data, retains that class.
The triple of unrestricted rank-eight bundles alone does not.
The space $G/K_0$ parametrizes reductions acting on this very
same rank-24 carrier; no claim that $\dim(G/K_0)=24$ is made.

\paragraph{The precise half-angle sign law and its relation to the cusp.}
For the fixed imaginary unit $i$ define the path
\begin{equation}
 u(\theta)=\cos(\theta/2)+i\sin(\theta/2),
                   \qquad 0\le\theta\le2\pi.
 \label{h24:rettri:halfangle}
\end{equation}
Equation~\eqref{h24:rettri:Rodrigues} gives
\[
 \operatorname{Ad}_{u(\theta)}i=i,\quad
 \operatorname{Ad}_{u(\theta)}j=\cos\theta\,j+\sin\theta\,k,\quad
 \operatorname{Ad}_{u(\theta)}k=-\sin\theta\,j+\cos\theta\,k.
\]
This is one full $2\pi$ rotation in the oriented $(j,k)$ plane.
Its lift starts at $1$ and ends at $-1$. It is therefore the
nontrivial loop in $SO(3)$: a null-homotopy of the loop would
lift over a disk and force its lifted endpoints to agree.
Substitution in \eqref{h24:rettri:threeactions} gives exactly
\begin{equation}
 \rho(-1)=(I,-I,-I).
 \label{h24:rettri:central-sign}
\end{equation}
It is central in \eqref{h24:rettri:RT} because each component
is scalar. Thus a full rotation in the three-dimensional
adjoint sector acts by minus the identity on both associated
spinor sectors while acting by the identity on the vector
sector. This conclusion follows from actual related-triple
maps and their product identity.

The parameter $\theta$ in \eqref{h24:rettri:halfangle} is a
structure-group parameter; it has not been equated with the
original cusp phase $\phi$. The actual dependence on that cusp
phase is already specified, with all four coordinates, by
\eqref{h24:rettri:z} and \eqref{h24:rettri:transition}.
No connection or holonomy is part of the definition of $P_H$.
In particular the order-two class proved here lies in
$\pi_5(\operatorname{Sp}(1))$ through degree-six clutching,
and is not being identified with a degree-one line-bundle
class. Every real line bundle on $X\cong S^6$ is trivial:
its unit sphere bundle is a double covering of a simply
connected space, so it has a section. The exact relation
established here is instead the faithful order-two detection
\eqref{h24:rettri:exact-detection}, the nontrivial real
three-plane \eqref{h24:rettri:rankthree}, its explicit
three-sector triality action, and the structural sign law
\eqref{h24:rettri:central-sign}.

% Primary convention source for integration bibliography:
% I. Yokota, Exceptional Lie Groups, arXiv:0902.0431,
% sections 1.10, 1.14, 1.16 and Theorem 2.7.1.

\subsubsection{A nonzero Hopf--cusp morphism to the stable mod--24 group}
\label{h24:hbridge:section}
The retained-angle input used below is the complete Theorem~\ref{h24:ret:main} above.

\paragraph{The exact source maps and conventions}

Retain the original marked source \((X,x_0)\), its fixed based
diffeomorphism \(d:(S^6,*)\to(X,x_0)\), the cusp coordinate \(t\),
the fixed cutoff \(\epsilon\), and the native quotient
\(\overline\pi_f\) of the cusp fibres by
\(\mathbb Z\langle\widehat w,\widehat\delta\rangle\).
The quotient coordinates are ordered by
\((\overline{\widehat\gamma},\overline{\widehat u})\).
Thus the fixed continuous based map is
\[
q_2:X\longrightarrow
\mathcal A=\operatorname{Susp}(S^1\times\overline T^2),
\qquad
q_2(x)=[|t|,t/|t|,\overline\pi_f(x)]
\quad(0<|t|<\epsilon).
\]
It sends the entire central fibre \(W\) to the first suspension
vertex and \(X\setminus\operatorname{int}N_0\) to the second.
The unreduced suspension \(\mathcal A\) is based at its first
vertex. In the original product torus set
\[
K^{(2)}=(\{1\}\times\overline T^2)
\cup(S^1\times\{u_1=0\})
\cup(S^1\times\{u_2=0\}),
\qquad c_3:T^3\longrightarrow T^3/K^{(2)}\cong S^3.
\]
The retained-angle sphere map and its sphere-source representative are
\begin{equation}
p_3([r,z])=[r,c_3(z)],\qquad
H=p_3q_2:X\longrightarrow\Sigma_r(S^3,*)\cong S^4,
\qquad h=H\circ d:S^6\longrightarrow S^4.
\label{h24:hbridge:eq:retained-map}
\end{equation}
Both suspension vertices and the whole cylinder over the basepoint
are identified in \(\Sigma_r(S^3,*)\). The preceding complete
retained-angle framing theorem proves
\begin{equation}
[h]=\eta_4\circ\eta_5\ne0
\quad\hbox{in}\quad\pi_6(S^4)\cong\mathbb Z/2.
\label{h24:hbridge:eq:proved-H}
\end{equation}
Its proof identifies the actual derivative-framed inverse torus
with the Lie-framed torus, with Arf invariant one. We use this
already proved class, retaining the particular map
\eqref{h24:hbridge:eq:retained-map}; see \cite{h24:hbridgebib:HBRetained}.

Fix the complex Hopf map
\[
\eta_2:(S^3,(1,0))\longrightarrow
(\mathbb{CP}^1,[1:0])\cong(S^2,*),\qquad
(z_1,z_2)\longmapsto[z_1:z_2].
\]
Use ordered smash coordinates on spheres. In this note
\(E^k f=f\wedge\operatorname{id}_{S^k}\) denotes right reduced
suspension; the usual ordered sphere homeomorphisms are understood.
Put \(\eta_n=E^{n-2}\eta_2\), and write \(\eta\) for the
stable Hopf class. The order of each factor is retained below.

\paragraph{The Hopf--cusp homomorphism and its nonzero value}

\begin{theorem}
\label{h24:hbridge:mu-theorem}
The assignment
\begin{equation}
\mu_\eta:\pi_6(S^4)\longrightarrow\pi_9(S^6),\qquad
\mu_\eta([x])=[\eta_2\wedge x]
\label{h24:hbridge:eq:mu}
\end{equation}
is an injective group homomorphism.\par
Its literal representative has
the complete domain and codomain
\[
\eta_2\wedge x:S^3\wedge S^6\cong S^9
\longrightarrow S^2\wedge S^4\cong S^6.
\]
Choose any generator \(\nu_6\) of
\(\pi_9(S^6)\cong\mathbb Z/24\). Then the particular original
cusp map supplies
\begin{equation}
\kappa_H=[\eta_2\wedge h]
=\mu_\eta([h])=12\nu_6\ne0.
\label{h24:hbridge:eq:kappa}
\end{equation}
It has exact order two. Its value is independent of orientation
reversals in the sphere identifications and of the choice of the
generator \(\nu_6\).
\end{theorem}

\begin{proof}
A based homotopy of \(x\) smashes with the fixed based Hopf map,
so \eqref{h24:hbridge:eq:mu} is well-defined. The exact ordered
point-set factorization is
\begin{equation}
\eta_2\wedge x
=(\eta_2\wedge\operatorname{id}_{S^4})
 \circ(\operatorname{id}_{S^3}\wedge x):
S^3\wedge S^6\longrightarrow S^3\wedge S^4
\longrightarrow S^2\wedge S^4.
\label{h24:hbridge:eq:factorization}
\end{equation}
The first applied map has source dimension nine and target dimension
seven; the second has source dimension seven and target dimension six.
The second map is \(E^4\eta_2\). To compare the first with
\(E^3x=x\wedge\operatorname{id}_{S^3}\), interchange its sphere
factors at source and target. The interchange degrees are respectively
\((-1)^{3\cdot6}=+1\) and \((-1)^{3\cdot4}=+1\). A based map
of a sphere of degree one is based homotopic to the identity, so
\begin{equation}
\mu_\eta([x])=(E^4\eta_2)_*(E^3[x]).
\label{h24:hbridge:eq:hom-formula}
\end{equation}
Suspension is a group homomorphism. Postcomposition is also a group
homomorphism on homotopy groups, since it distributes over the
source pinch map defining the sum. Thus
\eqref{h24:hbridge:eq:hom-formula} proves additivity, including zero
and inverses, before passage to stable homotopy.

Freudenthal makes
\[
E:\pi_{n+3}(S^n)\longrightarrow\pi_{n+4}(S^{n+1})
\]
an isomorphism for \(n\ge5\), because \(n+3<2n-1\).
In particular stabilization
\(\pi_9(S^6)\to\pi_3^{\mathrm S}\) is an isomorphism.
Stabilization from \(\pi_6(S^4)\) is likewise an isomorphism
onto \(\pi_2^{\mathrm S}\), because \(n+2<2n-1\) for
\(n\ge4\). The exact factorization
\eqref{h24:hbridge:eq:factorization} shows that its stable class is
the stable composition product of \(\eta\) with \([x]\).
Using the proved original class \eqref{h24:hbridge:eq:proved-H},
\begin{equation}
E^\infty\kappa_H
=\eta\,E^\infty[h]
=\eta\,\eta^2
=\eta^3
=12\nu\ne0\quad\hbox{in}\quad
\pi_3^{\mathrm S}\cong\mathbb Z/24.
\label{h24:hbridge:eq:stable-product}
\end{equation}
The integral relation here is the one explicitly stated in
\cite[p.~385]{h24:hbridgebib:HBHatcher}, not the coefficient \(4\) in the
two-primary quotient. More precisely, the two-primary relation
is \(\eta^3=4\nu\) modulo eight. Bilinearity and \(2\eta=0\)
give \(2\eta^3=0\) integrally. In \(\mathbb Z/24\) the only
elements killed by two are \(0\) and \(12\nu\), and only the
second has nonzero image \(4\nu\) modulo eight. This recovers
the exact integral coefficient and its nonvanishing.

The inverse of the stabilization isomorphism therefore gives
\eqref{h24:hbridge:eq:kappa}. Since the source of \(\mu_\eta\) is
cyclic of order two, generated by \([h]\), its nonzero value
proves injectivity and identifies its image as
\(\{0,12\nu_6\}\).

An orientation reversal at a sphere source or target multiplies
the corresponding stable homotopy class by \(-1\); here
\(-12\nu_6=12\nu_6\). A sign change in the chosen Hopf
representative also leaves the stable product unchanged since
\(-\eta=\eta\). Finally, replacing the generator by
\(u\nu_6\), with \(u\) a unit modulo twenty-four, does not
change the element: every such \(u\) is odd, and
\(12u\equiv12\pmod{24}\). All orientation and generator
assertions consequently hold for the actual order-two element.
\end{proof}

\begin{corollary}[The same nonzero class in the twelve-dimensional receiver]
Put \(\nu_9=E^3\nu_6\). The particular based map
\[
E^3(\eta_2\wedge h):S^{12}\longrightarrow S^9
\]
has class \(12\nu_9\ne0\) in
\(\pi_{12}(S^9)\cong\mathbb Z/24\).
\end{corollary}
\begin{proof}
Each of these three suspensions is an isomorphism in the preceding
stable range. Apply their composite to
\eqref{h24:hbridge:eq:kappa}; it preserves the exact order two.
\end{proof}

\paragraph{The actual marked peripheral character}

Retain the original three elliptic-kernel matrices, with their
original signs and ordering:
\begin{equation}
C_0=\begin{pmatrix}1&1\\0&1\end{pmatrix},\qquad
C_1=\begin{pmatrix}-1&-1\\1&0\end{pmatrix},\qquad
C_2'=\begin{pmatrix}0&1\\-1&0\end{pmatrix}.
\label{h24:hbridge:eq:matrices}
\end{equation}
The existing peripheral theorem \cite{h24:hbridgebib:HBHigher} defines the
surjective marked abelianization character
\begin{equation}
\chi:SL_2(\mathbb Z)\longrightarrow\mathbb Z/12,
\qquad
(\chi(C_0),\chi(C_1),\chi(C_2'))=(-1,4,-3).
\label{h24:hbridge:eq:chi}
\end{equation}
For an explicit check of this normalization, let
\[
S_0=\begin{pmatrix}0&-1\\1&0\end{pmatrix},\qquad T=C_0.
\]
Then \(C_2'=S_0^3\), \(C_1=T^{-1}S_0\), and direct matrix
multiplication gives \(C_1C_2'=C_0^{-1}\). In the modular
presentation \(S_0^4=1\), \(S_0^2=(S_0T)^3\), the
abelianized generators satisfy \(4s=0\) and \(s+3t=0\).
Thus \(s=-3t\) and \(12t=0\). The assignment
\(t\mapsto-1\), \(s\mapsto3\) satisfies both relations and
maps onto \(\mathbb Z/12\); consequently the order of the
abelianized generator is exactly twelve. It gives, without a
change of peripheral matrices, precisely
\(\chi(C_0)=-1\), \(\chi(C_1)=1+3=4\), and
\(\chi(C_2')=3\cdot3=-3\pmod{12}\).

\begin{theorem}[An exact nonzero peripheral homotopy morphism]
\label{h24:hbridge:peripheral-theorem}
Define
\begin{equation}
\Theta_H:\mathbb Z/12\longrightarrow\pi_6(S^4),
\qquad \Theta_H([n])=n[h],
\qquad
\Psi_H=\mu_\eta\circ\Theta_H\circ\chi:
SL_2(\mathbb Z)\longrightarrow\pi_9(S^6).
\label{h24:hbridge:eq:Psi}
\end{equation}
These are group homomorphisms. The complete coefficient formula is
\begin{equation}
\Psi_H(M)=12\,\widetilde{\chi(M)}\,\nu_6,
\label{h24:hbridge:eq:coefficient}
\end{equation}
where \(\widetilde{\chi(M)}\) is any integer lift of the
marked residue. In the original matrix order of
\eqref{h24:hbridge:eq:matrices}, its exact values are
\begin{equation}
(\Psi_H(C_0),\Psi_H(C_1),\Psi_H(C_2'))
=(12\nu_6,0,12\nu_6).
\label{h24:hbridge:eq:values}
\end{equation}
Its image and kernel are exactly
\begin{equation}
\operatorname{im}\Psi_H=\{0,12\nu_6\},\qquad
\ker\Psi_H=\chi^{-1}(2\mathbb Z/12).
\label{h24:hbridge:eq:kernel}
\end{equation}
In particular its kernel has index two in \(SL_2(\mathbb Z)\).
\end{theorem}

\begin{proof}
The element \([h]\) has exact order two, so replacing \(n\)
by \(n+12k\) leaves \(n[h]\) unchanged. Addition in its
abelian homotopy group proves that \(\Theta_H\) is a group
homomorphism. Moreover
\[
\ker\Theta_H=2\mathbb Z/12,
\qquad \operatorname{im}\Theta_H=\pi_6(S^4).
\]
The composite in \eqref{h24:hbridge:eq:Psi} is consequently a
group homomorphism, and \eqref{h24:hbridge:eq:kappa} gives
\(\mu_\eta(n[h])=12n\nu_6\). Changing the integer lift by
\(12k\) changes its coefficient by \(144k\), which is a
multiple of twenty-four. This proves
\eqref{h24:hbridge:eq:coefficient} with its independence of lifts.
The three coefficients are respectively
\(-12\equiv12\), \(48\equiv0\), and
\(-36\equiv12\pmod{24}\), proving
\eqref{h24:hbridge:eq:values}. Since \(12n\equiv0\pmod{24}\)
holds exactly when \(n\) is even, the asserted kernel is exact.
Surjectivity of \(\chi\) and nonvanishing at \(C_0\) prove
that the image is the full order-two subgroup, so the first
isomorphism theorem gives the index assertion.
\end{proof}

Here is the complete parity factorization. Let
\(\operatorname{red}_2:\mathbb Z/12\to\mathbb Z/2\) be
ordinary parity reduction, let
\(\iota_2([e])=12e\pmod{24}\), and let
\(c_{\nu_6}:\pi_9(S^6)\to\mathbb Z/24\) send \(\nu_6\)
to one. Every arrow and its codomain in the diagram are fixed:
\begin{equation}
\begin{CD}
SL_2(\mathbb Z) @>{\Psi_H}>> \pi_9(S^6)\\
@V{\operatorname{red}_2\circ\chi}VV @VV{c_{\nu_6}}V\\
\mathbb Z/2 @>{\iota_2}>> \mathbb Z/24.
\end{CD}
\label{h24:hbridge:eq:parity-square}
\end{equation}
Commutativity is exactly the coefficient calculation above; the
bottom arrow is injective because \(12e\) is zero modulo
twenty-four exactly when \(e\) is zero modulo two.

\paragraph{Comparison with the earlier maps, with their exact types}

Write \(\psi_H=c_{\nu_6}\circ\Psi_H\). Ordinary reduction
\(r:\mathbb Z/24\to\mathbb Z/12\) gives
\begin{equation}
r\circ\psi_H=0,
\qquad (r\circ\psi_H)(C_0)=0\ne-1=\chi(C_0).
\label{h24:hbridge:eq:not-lift}
\end{equation}
Thus this nonzero character is not a reduction-compatible lift of
\(\chi\). Its natural commuting square is
\eqref{h24:hbridge:eq:parity-square}. The existing nonexistence
theorem for a reduction-compatible lift remains valid: every
homomorphism from \(SL_2(\mathbb Z)\) to an abelian group
factors through its abelianization \(\mathbb Z/12\), whereas
a lift of \(\chi(C_0)=-1\) would require \(C_0\) to map
to \(11\) or \(23\) modulo twenty-four, each of order
twenty-four. The abelianized generator has order twelve and
cannot have either image.

The universal even-subgroup injection is the separate map
\[
j:\mathbb Z/12\hookrightarrow\mathbb Z/24,
\qquad j([a])=[2a].
\]
The new coefficient morphism factors through it by the exact
identity
\begin{equation}
\psi_H=j\circ[6]\circ\chi,
\qquad [6]:\mathbb Z/12\longrightarrow\mathbb Z/12,
\quad [a]\longmapsto[6a].
\label{h24:hbridge:eq:j-factor}
\end{equation}
Indeed \(2(6a)=12a\pmod{24}\). The map \(j\) is
injective with image of order twelve; the composite
\eqref{h24:hbridge:eq:j-factor} has image of order two because
\([6]\) retains precisely parity. The map \(j\) is realized
by the already proved universal suspension
\(E^6:\pi_6(S^3)\to\pi_{12}(S^9)\) after compatible
generator choices. Equation \eqref{h24:hbridge:eq:j-factor} is a
coefficient identity; the actual new geometric representative is
\(\eta_2\wedge h\), followed, when desired, by three
suspensions as above.

For the original angle-forgetting map
\(g=\widehat p_2q_2:X\to S^3\), the separate framed
calculation has proved \([g\circ d]=0\). Therefore the
previous morphism
\[
\Theta_g:\mathbb Z/12\longrightarrow\pi_6(S^3),
\qquad [n]\longmapsto n[g\circ d]
\]
and \(E^6\Theta_g\chi\) are zero. In contrast,
\(E^3\Psi_H(C_0)=12\nu_9\ne0\) in this same stable
receiver \(\pi_{12}(S^9)\). These statements agree because
the exact postcompositions from \(\mathcal A\) are different:
\(p_3\) retains the angular coordinate in its top-cell map,
while \(\widehat p_2\) forgets it. The new homomorphism
uses the already proved value of \(p_3q_2\) and the specified
Hopf map, with all maps displayed in
\eqref{h24:hbridge:eq:retained-map}--\eqref{h24:hbridge:eq:Psi}.

The constructed representative has source \(S^3\wedge X\cong
S^9\); its threefold suspension has source \(S^{12}\).
The number twenty-four is the order of the receiving stable
homotopy group. This calculation does not identify either source
with a twenty-four-dimensional carrier, does not supply new
periods or analytic overlap maps, and does not complete the
global higher torus family. It supplies the exact nonzero
homotopy and peripheral morphisms just proved.

\paragraph{Source checks.}
The three matrices and the marked values were checked directly
against equations \emph{elliptic-peripheral-matrices} and
\emph{peripheral-mod12-values} in the existing higher-rung draft
\cite{h24:hbridgebib:HBHigher}. The calculation above reproduces their signs
using the same \(S_0,T\). The retained-angle input is the
complete derivative-framing theorem in \cite{h24:hbridgebib:HBRetained}.
For the stable facts, the author-hosted text
\cite{h24:hbridgebib:HBHatcher} was read at Corollary~4.24, printed p.~360,
and at printed pp.~384--385. Page~385 explicitly distinguishes
the two-primary coefficient four from the integral coefficient
twelve. Proposition~4.56 on that page, with its proof on
p.~388, gives the stable composition product and its smash
description. The fixed-factor proof
\eqref{h24:hbridge:eq:factorization}--\eqref{h24:hbridge:eq:hom-formula}
establishes the needed homomorphism here directly.

\subsubsection{Boundary of the three tested numerical mechanisms}


% END RETAINED_ANGLE_AND_HOPF_BRIDGE_20260905

\begin{theorem}[Boundary of the three tested numerical mechanisms]
\label[theorem]{h24:thm:exact-numerical-boundary}
For the chamber indexing, peripheral abelianization, stable suspension,
and the three additional mechanisms tested below, the verified relations
are the following.
\begin{enumerate}[label=\textup{(\roman*)}]
\item The six Kuhn zero-strata and the six octonionic fixed flags have
      the marked \(S_3\)-equivariant finite-set bijection
      \eqref{h24:eq:Kuhn-flag-bijection}.  Their common labels also index
      the flag cells, and both Euler characteristics are \(6\).
\item The elliptic-kernel peripheral matrices produce the surjective
      character \(\chi:SL_2(\mathbb Z)\to\mathbb Z/12\), with no
      reduction-compatible or surjective mod--\(24\) lift.
\item The separate stable suspension map embeds \(\mathbb Z/12\) as
      \(2(\mathbb Z/24)\), and \cref{h24:prop:generator-free-bridge} gives the exact
      composite determined by \(g\). The actual coefficient is \(a=0\)
      by \cref{h24:thm:cusp-g-zero}; the composite is zero, and the associated \(S^{24}\)
      fourth-smash map is null-homotopic.
\item The ambient \(F_4\) root system has \(h=12\), and
      \(W(F_4)/W(D_4)\cong S_3\).  The three specific tests give:
      diagonal-sum factorization through \(P\) or \(Q\) exactly for
      \(N\mid3\); ordinary cohomology reduction for every \(N\), with
      \(e_N\) of order \(N/\gcd(N,6)\); and plain determinant-winding
      reduction onto \(\mathbb Z/N\) for every \(N\).  None of these specified
      operations selects \(N=12\) or \(N=24\) from the dimension or
      Coxeter number alone.
\end{enumerate}
No arrow between these four lines is defined merely by equality of their
integers.
\end{theorem}

\begin{proof}
Item (i) is \cref{h24:thm:Kuhn-octonionic-comparison}; item (ii) is
\cref{h24:thm:peripheral-moduli}; item (iii) is
\cref{h24:thm:stable-12-to-24,h24:prop:generator-free-bridge}; and the Coxeter
assertion in (iv) is \cref{h24:prop:two-Coxeter-systems}.  The three stated
tests for an additional arithmetic arrow are evaluated by
\cref{h24:prop:cusp-diagonal-modulus,h24:prop:no-modulus-from-cohomology-or-determinant}:
the primitive lattice
diagonal-sum factorization allows only \(N\mid3\), cohomology reduces for every
modulus, and determinant winding reduces onto every modulus.  These
computations prove the stated boundary for those mechanisms without
asserting that arbitrary future structures cannot be added.
\end{proof}

% BEGIN ES-NIEMEIER HOST-ONLY INPUT: exclude from generated h24: fragment.
% Proof-preserving conversion of ORBIT_CARTAN_BRIDGE.md, Sections 1--7.
% Staged fragment only; uses the host's original Albert coordinates.
\subsection{The full Albert orbit quotient and the Cartan family}
\label{h24:ocb:sec:full-family}

The fixed-spectrum comparison of
\cref{h24:thm:octonionic-triality,h24:att:sec:ordered-tensor} extends to the
whole Albert algebra with all three eigenvalue variables retained.
We give the invariant quotient, its smooth regular part, the exact
normal framing, and the two focal strata. The diagonalization and flag
geometry are established in \cite[Proposition~2.8.2 and
Theorems~2.7.4, 2.8.3]{h24:Yokota2009} and
\cite[Definition~2.2 and Sections~2.3--2.6]{h24:MareWillems2013}; the
calculations below specify their maps in the retained coordinates.

\subsubsection{Original space, invariants, and direct diagonalization}
\label{h24:ocb:sec:diagonalization}

Let $\mathbb O$ be the real octonion division algebra with its retained
conjugation, norm $N$, and inner product
$h(a,b)=\operatorname{Re}(a\bar b)$. Set
$J=\mathfrak h_3(\mathbb O)$, with Jordan product
$X\circ Y=(XY+YX)/2$, and retain the cyclic coordinates of
\eqref{h24:att:eq:Q}:
\begin{equation}
 X=\begin{pmatrix}\xi_1&q_3&\bar q_2\\
 \bar q_3&\xi_2&q_1\\q_2&\bar q_1&\xi_3\end{pmatrix}
 =\sum_{i=1}^3\xi_i e_i+\sum_{i=1}^3E_i(q_i).
 \label{h24:ocb:eq:coordinates}\tag{OCB1}
\end{equation}
Here $e_i$ are the original diagonal idempotents, $E_i$ are the original
off-diagonal inclusions, and $I=e_1+e_2+e_3$. Matrix products in the
Jordan product involve only products of two octonions. Powers of a single
Jordan element are unambiguous by the Jordan identity; all subsequent
polynomial formulas use these powers.

The positive trace metric and determinant are
\begin{equation}
 \begin{split}
 G(X,Y)&=\operatorname{tr}(X\circ Y),\\
 G(X,X)&=\xi_1^2+\xi_2^2+\xi_3^2+2\sum_iN(q_i),\\
 d(X)&=\xi_1\xi_2\xi_3-\sum_i\xi_iN(q_i)
              +2\operatorname{Re}((q_1q_2)q_3).
 \end{split}
 \label{h24:ocb:eq:metric-determinant}\tag{OCB2}
\end{equation}
Thus $\dim_{\mathbb R}J=27$, with every diagonal coordinate, metric
factor, and mixed cubic retained. Define
\begin{equation}
 t(X)=\operatorname{tr}X,\qquad u(X)=G(X,X),\qquad
 e_2(X)=\frac{t(X)^2-u(X)}2,\qquad
 \Phi(X)=(t(X),u(X),d(X)).
 \label{h24:ocb:eq:invariants}\tag{OCB3}
\end{equation}
Let $F_4=\operatorname{Aut}_{\mathbb R}(J,\circ)$. An automorphism
fixes $I$, preserves trace and $G$, and therefore preserves the
determinant by the trace identity
\begin{equation}
 d(X)=\tfrac13\operatorname{tr}(X^3)
 -\tfrac12t(X)u(X)+\tfrac16t(X)^3.
 \label{h24:ocb:eq:trace-determinant}\tag{OCB4}
\end{equation}
For completeness, trace preservation follows without choosing a
preferred matrix realization: the Peirce table gives
$\operatorname{Tr}_{\operatorname{End}(J)}(L_X)=9\operatorname{tr}X$;
conjugation of $L_X$ by an automorphism proves the assertion. Product
preservation is a closed polynomial condition, so $F_4$ is closed in the
orthogonal group of $(J,G)$ and hence compact.

\begin{proposition}[Diagonalization and the exact invariant fibres]
\label{h24:ocb:prop:invariant-fibres}
Every $X\in J$ is $F_4$-conjugate to a real diagonal matrix. The fibres
of $\Phi$ are exactly the $F_4$ orbits, including repeated-eigenvalue
and scalar orbits. The map $\Phi:J\to\Phi(J)$ is proper.
\end{proposition}

\begin{proof}
On the compact orbit of $X$, maximize
$\xi_1^2+\xi_2^2+\xi_3^2$. If $q_i\ne0$ at a maximizing point,
put $a=q_i/|q_i|$, let $(i,j,k)$ be cyclic, and use the already proved
derivation
\begin{equation}
 D_i(a)=4[L_{E_i(a)},L_{e_j}].
 \label{h24:ocb:eq:derivation}\tag{OCB5}
\end{equation}
Its action satisfies
$D_i(a)(e_j-e_k)=2E_i(a)$ and
$D_i(a)E_i(a)=-2(e_j-e_k)$. The remaining off-diagonal slots do not
contribute diagonal entries under this action. Write the relevant
component as $v(e_j-e_k)+wE_i(a)$, with
$v=(\xi_j-\xi_k)/2$ and $w=|q_i|>0$. Exponentiation rotates $(v,w)$
through angle $2s$. Some rotation gives squared diagonal-difference
coordinate $v^2+w^2$, increasing the diagonal-square sum by $2w^2$.
This contradicts maximality. Hence every off-diagonal coordinate
vanishes at a maximum, proving diagonalization with the diagonal mean
retained.

If $X=g\sum_i\lambda_i e_i$, then
\eqref{h24:ocb:eq:invariants}--\eqref{h24:ocb:eq:trace-determinant} show that
the multiset of real eigenvalues is exactly the multiset of roots of
\begin{equation}
 p_X(z)=z^3-t(X)z^2+e_2(X)z-d(X).
 \label{h24:ocb:eq:characteristic-polynomial}\tag{OCB6}
\end{equation}
It is independent of the diagonalizing automorphism. Every permutation
of a real diagonal is realized by conjugation with a real orthogonal
permutation matrix. Such conjugation preserves the Jordan product;
a commuting sign adjustment can be used if determinant $+1$ is
desired. Equal invariant triples thus give the same eigenvalue multiset
and the same orbit. Conversely, the invariants are constant on orbits,
as proved above. This includes repeated roots and scalar matrices.
Finally, a compact set of invariant triples bounds $u=G(X,X)$, and
its inverse image is closed and bounded in $J$, hence compact.
\end{proof}

This is the compact-orbit diagonalization argument of
\cite[Proposition~2.8.2, pp.~56--57]{h24:Yokota2009}, reproduced in the
retained derivation coordinates. It also proves the polynomial identities
used below by transporting their diagonal versions.

\subsubsection{Exact projectors, smooth orbit inverse, and regular family}
\label{h24:ocb:sec:regular-family}

Fix labelled, pairwise distinct real numbers
$\lambda_1,\lambda_2,\lambda_3$, with arbitrary sum. Put
$d_\lambda=\sum_i\lambda_i e_i$ and
$\mathcal O_\lambda=F_4d_\lambda$. The ordered frame space is
\[
 \mathcal F=\{(p_1,p_2,p_3):p_i^2=p_i,\
 \operatorname{tr}p_i=1,\ p_i\circ p_j=0\ (i\ne j),\
 \textstyle\sum_i p_i=I\}.
\]
The correspondence $(p_1,p_2,p_3)\mapsto(p_1,p_2)$ and its inverse
$(p_1,p_2)\mapsto(p_1,p_2,I-p_1-p_2)$ identify this space with the
flag space of \eqref{h24:eq:octonionic-flag-definition}; in particular
$\mathcal F=F_4/\operatorname{Spin}(8)$ by
\cref{h24:thm:octonionic-triality}. Define
\begin{equation}
 \Xi_\lambda:\mathcal F\longrightarrow\mathcal O_\lambda,
 \qquad (p_i)\longmapsto\sum_i\lambda_i p_i,
 \label{h24:ocb:eq:orbit-map}\tag{OCB7}
\end{equation}
and, for $\{i,j,k\}=\{1,2,3\}$, define
\begin{equation}
 P_i^\lambda(X)=
 \frac{(X-\lambda_jI)\circ(X-\lambda_kI)}
      {(\lambda_i-\lambda_j)(\lambda_i-\lambda_k)}.
 \label{h24:ocb:eq:projectors}\tag{OCB8}
\end{equation}

\begin{proposition}[The ordered spectral diffeomorphism]
\label{h24:ocb:prop:spectral-diffeomorphism}
The map $\Xi_\lambda$ is an $F_4$-equivariant diffeomorphism, with
inverse $(P_1^\lambda,P_2^\lambda,P_3^\lambda)$. In particular
$\dim\mathcal O_\lambda=24$. If
$\mathcal C=\{\lambda_1>\lambda_2>\lambda_3\}$ and
$J_{\mathrm{reg}}$ is the set of elements with distinct eigenvalues,
then $\mathcal F\times\mathcal C\to J_{\mathrm{reg}}$, given by the
same spectral sum, is a diffeomorphism.
\end{proposition}

\begin{proof}
For $X=\sum_l\lambda_l p_l$, orthogonality and idempotency give
\[
 \begin{aligned}
 (X-\lambda_jI)\circ(X-\lambda_kI)
 &=\sum_l(\lambda_l-\lambda_j)(\lambda_l-\lambda_k)p_l\\
 &=(\lambda_i-\lambda_j)(\lambda_i-\lambda_k)p_i.
 \end{aligned}
\]
Thus \eqref{h24:ocb:eq:projectors} recovers every ordered projector and
proves both inverse compositions. It also proves that the stabilizer
of $d_\lambda$ is the pointwise diagonal-frame stabilizer
$\operatorname{Spin}(8)$. The quotient orbit map is an injective
immersion of the compact homogeneous space
$F_4/\operatorname{Spin}(8)$ into the Hausdorff vector space $J$,
hence an embedding. Its inverse is polynomial in $X$ divided by
nonzero fixed real constants, proving smoothness of the inverse and
the asserted equivariant diffeomorphism.

The differential at fixed $\lambda$ is explicitly
\begin{equation}
 dP_i^\lambda|_X(H)=
 \frac{2X\circ H-(\lambda_j+\lambda_k)H}
      {(\lambda_i-\lambda_j)(\lambda_i-\lambda_k)}.
 \label{h24:ocb:eq:projector-differential}\tag{OCB9}
\end{equation}
At $d_\lambda$, the sector $E_i(\mathbb O)$ joins eigenspaces $j,k$.
On this sector $dP_j(H)=H/(\lambda_j-\lambda_k)$,
$dP_k(H)=-H/(\lambda_j-\lambda_k)$, and $dP_i(H)=0$.
Substitution into $\sum_l\lambda_l dP_l(H)$ returns $H$, verifying
the differential inverse with its signs.

The original derivations satisfy
\begin{equation}
 \begin{gathered}
 D_i(a)d_\lambda=(\lambda_j-\lambda_k)E_i(a)=\Delta_i E_i(a),\\
 (\Delta_1,\Delta_2,\Delta_3)=
 (\lambda_2-\lambda_3,\lambda_3-\lambda_1,\lambda_1-\lambda_2),
 \qquad(i,j,k)\text{ cyclic}.
 \end{gathered}
 \label{h24:ocb:eq:tangent-multipliers}\tag{OCB10}
\end{equation}
All three multipliers are nonzero, so the three eight-dimensional
off-diagonal sectors constitute the tangent space. Their transported
Peirce subbundles at every orbit point have metric $2\Delta_i^2h$
in the homogeneous tangent coordinates of
\eqref{h24:att:eq:bundle-comparison}.

For the full map
\begin{equation}
 \mathcal F\times\mathcal C\longrightarrow J_{\mathrm{reg}},
 \qquad ((p_i),(\lambda_i))\longmapsto\sum_i\lambda_i p_i,
 \label{h24:ocb:eq:full-regular-map}\tag{OCB11}
\end{equation}
surjectivity follows from diagonalization and sorting. Injectivity
follows from the unique ordered roots and \eqref{h24:ocb:eq:projectors}.
At a simple root,
$\partial p_X/\partial z=\prod_{j\ne i}(\lambda_i-\lambda_j)\ne0$;
the implicit function theorem therefore gives a smooth local root
function. Ordering glues these functions globally.
Equation~\eqref{h24:ocb:eq:projectors}, with the smoothly varying roots,
provides the smooth inverse of \eqref{h24:ocb:eq:full-regular-map}.
\end{proof}

\subsubsection{Retained mean, radius, and angular parameter}
\label{h24:ocb:sec:mean-radius-angle}

For every original $X$, put
\begin{equation}
 \mu=\frac{t(X)}3,\qquad Y=X-\mu I,\qquad
 s=G(Y,Y)=u(X)-\frac{t(X)^2}3,\qquad d_0=d(Y).
 \label{h24:ocb:eq:mean-radius}\tag{OCB12}
\end{equation}
Write $J_0=\mathfrak h_3^0(\mathbb O)=\{Y\in J:\operatorname{tr}Y=0\}$.
The displayed change of variables has inverse $X=\mu I+Y$;
$\mu$ and $s$ retain their actual values. The eigenvalues
$a_i=\lambda_i-\mu$ of $Y$ satisfy $\sum_i a_i=0$, so expansion gives
\begin{equation}
 d(X)=\mu^3-\frac{\mu s}{2}+d_0.
 \label{h24:ocb:eq:mean-determinant}\tag{OCB13}
\end{equation}
For $s=r^2>0$, the ordered eigenvalues have the exact parametrization
\begin{equation}
 \begin{aligned}
 \lambda_1&=\mu+\sqrt{2/3}\,r\cos\theta,\\
 \lambda_2&=\mu+\sqrt{2/3}\,r\cos(\theta-2\pi/3),\\
 \lambda_3&=\mu+\sqrt{2/3}\,r\cos(\theta+2\pi/3),
 \end{aligned}
 \qquad 0\le\theta\le\pi/3.
 \label{h24:ocb:eq:angular-eigenvalues}\tag{OCB14}
\end{equation}
The cosines sum to zero, their squared sum is $3/2$, and their
product is $\tfrac14\cos3\theta$. Consequently
\begin{equation}
 \begin{aligned}
 t(X)&=3\mu,&u(X)&=3\mu^2+r^2,\\
 d_0&=\frac{r^3}{3\sqrt6}\cos3\theta,&
 d(X)&=\mu^3-\frac{\mu r^2}{2}
                    +\frac{r^3}{3\sqrt6}\cos3\theta.
 \end{aligned}
 \label{h24:ocb:eq:angular-invariants}\tag{OCB15}
\end{equation}
For $0<\theta<\pi/3$ the ordered inequalities are strict. The inverse is
\begin{equation}
 \begin{gathered}
 \mu=t/3,\qquad r=\sqrt{u-t^2/3},\\
 \theta=\tfrac13\arccos\left(
 \frac{3\sqrt6\,[d-\mu^3+\mu r^2/2]}{r^3}\right).
 \end{gathered}
 \label{h24:ocb:eq:angular-inverse}\tag{OCB16}
\end{equation}
To prove that every ordered triple is listed, put $a_3=-a_1-a_2$
and expand the discriminant:
\begin{equation}
 \mathcal D=\prod_{i<j}(a_i-a_j)^2
            =\frac{s^3}{2}-27d_0^2\ge0.
 \label{h24:ocb:eq:discriminant}\tag{OCB17}
\end{equation}
Thus the arccos argument lies in $[-1,1]$. The triple in
\eqref{h24:ocb:eq:angular-eigenvalues} has the same characteristic
polynomial and the correct ordering, hence equals the original
triple. Equality in \eqref{h24:ocb:eq:discriminant} means repeated
eigenvalues. If $r=0$, positivity of $G$ gives $Y=0$ and $X=\mu I$;
this scalar stratum requires no angle.

The full invariant image is therefore exactly
\begin{equation}
 \left\{(t,u,d):s=u-t^2/3\ge0,\quad
 27(d-\mu^3+\mu s/2)^2\le s^3/2,\quad\mu=t/3\right\}.
 \label{h24:ocb:eq:invariant-image}\tag{OCB18}
\end{equation}
Every triple in this set is attained by
\eqref{h24:ocb:eq:angular-eigenvalues}, including the scalar case just
described. This gives the whole three-dimensional orbit quotient;
specified mean and radius select its corresponding fibre.

The full regular metric also records these parameters. Let $h_i$ be
the metric $h$ on the $i$th associated bundle, expressed by the
homogeneous tangent coordinate $K_i[g,a]$ of \eqref{h24:att:eq:Kmap}.
The pullback of $G$ by
\eqref{h24:ocb:eq:full-regular-map}--\eqref{h24:ocb:eq:angular-eigenvalues} is
\begin{equation}
 3\,d\mu^2+dr^2+r^2d\theta^2
             +\sum_i2\Delta_i(\mu,r,\theta)^2h_i.
 \label{h24:ocb:eq:full-metric}\tag{OCB19}
\end{equation}
Indeed the diagonal differential has coefficient vectors
$\mathbf1$, $a/r$, and $\partial_\theta a$. Their pairwise inner
products vanish, and their squared norms are $3,1,r^2$. Every
diagonal differential is orthogonal to every off-diagonal Peirce
sector. The sector terms follow from
\eqref{h24:ocb:eq:tangent-multipliers} and $E_i^*G=2h$, with every
eigenvalue difference and radial factor retained.

\subsubsection{Exact codimension and global normal framing}
\label{h24:ocb:sec:normal-framing}

The same $24$-dimensional regular orbit has codimension three in $J$,
two in the affine fixed-trace space $\mu I+J_0$, and one in the
fixed-trace, fixed-radius sphere $\mu I+S_r^{25}$. The change of
variables \eqref{h24:ocb:eq:mean-radius} and the radius-level inclusion
give the actual maps between these ambient descriptions.

For a fixed distinct triple, its normal space in $J$ is
\begin{equation}
 \nu_X\mathcal O_\lambda
 =\operatorname{span}_{\mathbb R}
     \{P_1^\lambda(X),P_2^\lambda(X),P_3^\lambda(X)\}.
 \label{h24:ocb:eq:projector-normal-space}\tag{OCB20}
\end{equation}
At a diagonal point this is the orthogonal complement of the three
off-diagonal tangent sectors; transport by $F_4$ proves the equality
everywhere. The projectors are pairwise $G$-orthogonal with squared
norm one.

The differential of the invariant map has gradients
\begin{equation}
 \nabla t=I,\qquad\nabla u=2X,\qquad
 \nabla d=X^\#=X^2-tX+e_2I.
 \label{h24:ocb:eq:invariant-gradients}\tag{OCB21}
\end{equation}
To check the determinant gradient, diagonalize $X$. By
\eqref{h24:ocb:eq:metric-determinant}, its derivative on a diagonal
variation is $\sum_i\lambda_j\lambda_k H_{ii}$, where
$\{i,j,k\}=\{1,2,3\}$, and its derivative on an off-diagonal
variation is zero. The right side of
\eqref{h24:ocb:eq:invariant-gradients} is the diagonal vector of these
coefficients; automorphism invariance proves the formula for every
$X,H$. In the projector basis the coefficient determinant is
\begin{equation}
 \det\begin{pmatrix}
 1&1&1\\2\lambda_1&2\lambda_2&2\lambda_3\\
 \lambda_2\lambda_3&\lambda_1\lambda_3&\lambda_1\lambda_2
 \end{pmatrix}
 =-2(\lambda_1-\lambda_2)(\lambda_1-\lambda_3)
                                      (\lambda_2-\lambda_3).
 \label{h24:ocb:eq:gradient-determinant}\tag{OCB22}
\end{equation}
Hence $\Phi$ has rank three exactly on the regular stratum, where
its tangent kernel is the orbit tangent space.

In the fixed-trace hyperplane define the two normal vectors
\begin{equation}
 Y=X-\mu I,\qquad Z=Y^2-\frac{s}{3}I.
 \label{h24:ocb:eq:two-normal-fields}\tag{OCB23}
\end{equation}
They have trace zero and are normal by
\eqref{h24:ocb:eq:projector-normal-space}. The trace-zero eigenvalue
identities give
\begin{equation}
 \operatorname{tr}Y^3=3d_0,\qquad
 \operatorname{tr}Y^4=s^2/2,\qquad
 \operatorname{Gram}(Y,Z)=
 \begin{pmatrix}s&3d_0\\3d_0&s^2/6\end{pmatrix}.
 \label{h24:ocb:eq:normal-gram}\tag{OCB24}
\end{equation}
Specifically, $\sum_i a_i^3=3a_1a_2a_3$ follows from
$\sum_i a_i=0$. Squaring $\sum_{i<j}a_ia_j=-s/2$ gives
$\sum_{i<j}a_i^2a_j^2=s^2/4$, since the remaining cross term is
$2a_1a_2a_3\sum_i a_i=0$. Expanding $s^2$ now gives
$\sum_i a_i^4=s^2/2$. These identities and
$Z=Y^2-sI/3$ prove every Gram entry. Therefore
\begin{equation}
 \det\operatorname{Gram}(Y,Z)=s^3/6-9d_0^2
                             =\mathcal D/3>0
 \label{h24:ocb:eq:normal-gram-determinant}\tag{OCB25}
\end{equation}
for distinct eigenvalues. Thus $(Y,Z)$ is a smooth ordered global
two-framing of the orbit normal bundle in $\mu I+J_0$.

The corresponding explicitly specified orthonormal frame is
\begin{equation}
 n_1=\frac{Y}{\sqrt s},\qquad
 n_2=\frac{Z-(3d_0/s)Y}{\sqrt{s^2/6-9d_0^2/s}}.
 \label{h24:ocb:eq:orthonormal-frame}\tag{OCB26}
\end{equation}
Every square root here is the positive real root; its argument is
positive by \eqref{h24:ocb:eq:normal-gram-determinant}. In $J$, prepend
$n_0=I/\sqrt3$ to obtain $(n_0,n_1,n_2)$. In the radius sphere the
normal direction is $n_2$. These formulas identify the three normal
bundles explicitly.

These normal fields are parallel for the normal connection along the
fixed orbit. Each is a linear combination of the $P_i^\lambda$ with
coefficients constant on that orbit. Along $g(s)X$, its ambient
derivative is $g'(s)$ applied to the corresponding diagonal normal
vector, and lies in the transported off-diagonal tangent sectors by
\eqref{h24:ocb:eq:tangent-multipliers}. Its normal component is zero.
Thus the normal bundle is flat with the displayed global parallel
orthonormal frame.

For any transported normal $\eta=\sum_i\nu_iP_i^\lambda$ with fixed
real $\nu_i$, use the convention $A_\eta H=-(D_H\eta)^T$. Then
\begin{equation}
 A_\eta\big|_{E_i(\mathbb O)}
 =-\frac{\nu_j-\nu_k}{\lambda_j-\lambda_k}\,\operatorname{id},
 \qquad(i,j,k)\text{ cyclic},
 \label{h24:ocb:eq:shape-operator}\tag{OCB27}
\end{equation}
with the sectors transported by $g$ at other points. Indeed a path
generated by $D_i(a)$ has tangent
$(\lambda_j-\lambda_k)E_i(a)$, whereas the derivative of its
normal field is $(\nu_j-\nu_k)E_i(a)$. This proves the formula
and its minus sign. The shape operator for $n_0$ is zero, and the
one for the radial normal $n_1$ is $-\operatorname{id}/r$.

\subsubsection{Cartan equations and spherical principal curvatures}
\label{h24:ocb:sec:cartan-equations}

On the whole trace-zero space $J_0$, with its original trace metric,
define the explicitly scaled cubic
\begin{equation}
 C(Y)=3\sqrt6\,d(Y).
 \label{h24:ocb:eq:cartan-polynomial}\tag{OCB28}
\end{equation}
The inverse relation is $d(Y)=C(Y)/(3\sqrt6)$, and
\eqref{h24:ocb:eq:mean-determinant} recovers the full determinant with
its mean and radius. Orthogonal projection of
\eqref{h24:ocb:eq:invariant-gradients} to $J_0$ gives
$\nabla_{J_0}d=Y^2-sI/3=Z$. Hence
\eqref{h24:ocb:eq:normal-gram} gives
\begin{equation}
 \|\nabla C\|_G^2=54\|Z\|_G^2=9s^2,
 \qquad\Delta_{J_0}C=0.
 \label{h24:ocb:eq:cartan-euclidean}\tag{OCB29}
\end{equation}
For the second equality, the full $J$ Laplacian of
\eqref{h24:ocb:eq:metric-determinant} is $-8\operatorname{tr}X$.
Each eight-dimensional octonion slot has metric $2h$ and contributes
$-8\xi_i$; each pure diagonal second derivative is zero, and the
mixed cubic is separately linear in each real coordinate. Restricting
to $J_0$ subtracts the second derivative in the unit scalar direction
$I/\sqrt3$. The identity
\[
 d(Y+zI/\sqrt3)=d(Y)-\frac{s z}{2\sqrt3}
                              +\frac{z^3}{3\sqrt3}
\]
shows that this second derivative is zero at $z=0$. The full
Laplacian also vanishes at trace zero, proving
\eqref{h24:ocb:eq:cartan-euclidean} with the original metric factor two.

On each specified radius sphere $S_r^{25}\subset J_0$, $r>0$,
set $f_r=C/r^3$. Then \eqref{h24:ocb:eq:angular-invariants} gives
$f_r=\cos3\theta$, and homogeneity yields
\begin{equation}
 \|\nabla_{S_r}f_r\|^2=\frac9{r^2}(1-f_r^2),\qquad
 \Delta_{S_r}f_r=-\frac{81}{r^2}f_r.
 \label{h24:ocb:eq:cartan-spherical}\tag{OCB30}
\end{equation}
Indeed the radial component of $\nabla C$ is $(3C/r)n_1$.
The polar Laplacian on the $26$-dimensional Euclidean space,
applied to the degree-three homogeneous harmonic function, is
$r[3(3+24)f_1+\Delta_{S_1}f_1]$. These calculations prove both
constants in \eqref{h24:ocb:eq:cartan-spherical}. Its regular level
hypersurfaces are therefore isoparametric: both the squared gradient
and the Laplacian depend only on the function value. By
\eqref{h24:ocb:eq:angular-invariants}--\eqref{h24:ocb:eq:invariant-image},
the regular levels are precisely the fixed-mean, fixed-radius regular
$F_4/\operatorname{Spin}(8)$ orbits.

The principal curvatures admit a direct calculation. On the diagonal
section let $N_\theta=(1/r)\partial_\theta Y$, a unit vector
perpendicular to $Y$, and transport it by the projectors. The exact
eigenvalue differences in \eqref{h24:ocb:eq:angular-eigenvalues} are
\begin{equation}
 \Delta_1=\sqrt2r\sin\theta,\qquad
 \Delta_2=-\sqrt2r\sin(\theta+\pi/3),\qquad
 \Delta_3=-\sqrt2r\sin(\theta-\pi/3).
 \label{h24:ocb:eq:angular-differences}\tag{OCB31}
\end{equation}
Apply \eqref{h24:ocb:eq:shape-operator} with
$\nu_i=(\partial_\theta\lambda_i)/r$. For the declared normal
$N_\theta$, the three principal curvatures are
\begin{equation}
 -\frac1r\cot\theta,\qquad
 -\frac1r\cot(\theta+\pi/3),\qquad
 -\frac1r\cot(\theta-\pi/3),
 \label{h24:ocb:eq:principal-curvatures}\tag{OCB32}
\end{equation}
each on its original eight-dimensional Peirce sector. They are
distinct for $0<\theta<\pi/3$. In the invariant frame
\eqref{h24:ocb:eq:orthonormal-frame}, one has $n_2=-N_\theta$:
differentiating $C=r^3\cos3\theta$ gives
$\nabla_{S_r}C=-3r^2\sin3\theta\,N_\theta$, whereas $n_2$
is the positive-length version of the projected determinant
gradient. The principal curvatures for $n_2$ are consequently the
opposites of \eqref{h24:ocb:eq:principal-curvatures}. This proves the
three-curvature, multiplicity-eight Cartan family with both normal
orientations, mean, and radius specified.

\subsubsection{Focal projective planes and their endpoint maps}
\label{h24:ocb:sec:focal-strata}

At $\theta=0$ one has
$\lambda_2=\lambda_3=\mu-r/\sqrt6$ and
$\lambda_1=\mu+2r/\sqrt6$; at $\theta=\pi/3$ one has
$\lambda_1=\lambda_2=\mu+r/\sqrt6$ and
$\lambda_3=\mu-2r/\sqrt6$. The two singular orbits in
$\mu I+S_r^{25}$ are exactly
\begin{equation}
 \begin{aligned}
 \mathcal P^+_{\mu,r}
 &=\{\mu I+\tfrac r{\sqrt6}(3p-I):p^2=p,\
                                       \operatorname{tr}p=1\},\\
 \mathcal P^-_{\mu,r}
 &=\{\mu I-\tfrac r{\sqrt6}(3p-I):p^2=p,\
                                       \operatorname{tr}p=1\}.
 \end{aligned}
 \label{h24:ocb:eq:focal-planes}\tag{OCB33}
\end{equation}
Their inverse maps are, respectively,
\begin{equation}
 p=\tfrac13\left(I+\tfrac{\sqrt6}{r}(X-\mu I)\right),\qquad
 p=\tfrac13\left(I-\tfrac{\sqrt6}{r}(X-\mu I)\right).
 \label{h24:ocb:eq:focal-inverses}\tag{OCB34}
\end{equation}
These affine maps are smooth, have the displayed two-sided inverses
on the specified sets, and are $F_4$-equivariant. The primitive
idempotent space is the sixteen-dimensional Cayley projective plane
$\mathbb O\mathbb P^2=F_4/\operatorname{Spin}(9)$
\cite[Theorems~2.7.4 and 2.8.3]{h24:Yokota2009}
\cite[Section~2.3]{h24:MareWillems2013}. Hence the formulas give embedded
projective-plane strata.

The endpoint maps of the regular family are the explicit projections
\begin{equation}
 \begin{aligned}
 (p_1,p_2,p_3)&\longmapsto
             \mu I+\tfrac r{\sqrt6}(3p_1-I),\\
 (p_1,p_2,p_3)&\longmapsto
             \mu I-\tfrac r{\sqrt6}(3p_3-I).
 \end{aligned}
 \label{h24:ocb:eq:endpoint-maps}\tag{OCB35}
\end{equation}
Each fibre is $\operatorname{Spin}(9)/\operatorname{Spin}(8)\cong S^8$:
fixing one primitive idempotent leaves two ordered complementary
idempotents, the octonionic projective line. These frame bundles
are identified in \cite[Sections~2.3--2.6]{h24:MareWillems2013}.
They are also the focal maps. Indeed the diagonal path $Y(\theta)$
is a great circle because $\partial_\theta^2Y=-Y$ and
$|\partial_\theta Y|=r$. Transport by $F_4$ therefore gives the
normal geodesics to the regular hypersurfaces. At each endpoint one
of the tangent multipliers in \eqref{h24:ocb:eq:angular-differences}
vanishes with its eight-dimensional sector, giving exactly the
collapse in \eqref{h24:ocb:eq:endpoint-maps}. Each focal stratum has
codimension nine in the ambient radius sphere and a normal sphere
of dimension eight.

For $r=0$, the construction gives the scalar line $X=\mu I$,
each point fixed by $F_4$. For $r>0$ and endpoint angle, the full
repeated-eigenvalue stratum retains the two parameters $(\mu,r)$
in addition to its projective-plane coordinate, and has dimension
eighteen in $J$. These descriptions retain every scalar, radial,
and mean stratum.

\subsubsection{Exact map from the equal-diagonal triality slice}
\label{h24:ocb:sec:triality-slice}

For the original matrix
$Q_{\mathbb O}(\lambda;q)=\lambda I+E(q)$, with
$E(q)=\sum_iE_i(q_i)$ and
$\tau(q)=\operatorname{Re}((q_1q_2)q_3)$, the maps above give
\begin{equation}
 \mu=\lambda,\qquad Y=E(q),\qquad
 r^2=2\sum_iN(q_i),\qquad d_0=2\tau(q),\qquad
 C(Y)=6\sqrt6\,\tau(q).
 \label{h24:ocb:eq:triality-parameters}\tag{OCB36}
\end{equation}
Thus its invariant triple is
\begin{equation}
 \begin{aligned}
 \Phi(Q_{\mathbb O}(\lambda;q))
 =\bigl(&3\lambda,\ 3\lambda^2+2\sum_iN(q_i),\\
        &\lambda^3-\lambda\sum_iN(q_i)+2\tau(q)\bigr).
 \end{aligned}
 \label{h24:ocb:eq:triality-invariant-map}\tag{OCB37}
\end{equation}
This is the explicit invariant morphism from the original triality
coordinates to the full orbit quotient. At every simple spectrum,
its ordered-projector lift is \eqref{h24:ocb:eq:projectors} applied to
$Q_{\mathbb O}(\lambda;q)$. The positive boundary graph of
\cref{h24:tri:thm:spectral-graph} chooses the original scalar
$\lambda=-\lambda_{\min}(E(q))$. In the full orbit coordinates,
this says that the least eigenvalue is zero and the other two are
nonnegative. Equations~\eqref{h24:ocb:eq:triality-parameters}--
\eqref{h24:ocb:eq:triality-invariant-map} prove the relation to the whole
Albert orbit family with all coefficients and this scalar choice
retained.

\paragraph{Attribution and scope.}
Yokota's diagonalization is Proposition~2.8.2, pp.~56--57;
his primitive-idempotent stabilizer and Cayley-plane identification
are Theorems~2.7.4 and 2.8.3, pp.~54 and 57
\cite{h24:Yokota2009}. Mare and Willems give the primitive idempotents,
full real diagonal space, frame stabilizer, projective plane,
$S^8$ frame projections, and the polar isotropy representation
in Definition~2.2 and Sections~2.3--2.6, pp.~488--493
\cite{h24:MareWillems2013}. Their description identifies $J_0$ with
the isotropy representation of $E_{6(-26)}/F_4$ and displays its
diagonal two-plane and three eight-dimensional root spaces. The
Cartan-family equations and focal maps used here are proved in
\eqref{h24:ocb:eq:cartan-euclidean}--\eqref{h24:ocb:eq:endpoint-maps}.
The orbit and flag geometry is established literature; these
calculations assert neither complex integrability nor physical
dynamics.


% Proof-preserving conversion of NIEMEIER_CARTAN_SPECTRAL_BRIDGE.md.
% Staged fragment only; depends on the common-lattice and orbit fragments.
\subsection{The exact spectral map of the common marked lattice}
\label{h24:ncs:sec:spectral-bridge}

The literal common-lattice embedding of
\cref{wcl:thm:literal-intersection} has a matrix-valued and
frame-valued lift into the Albert orbit family of
\cref{h24:ocb:sec:full-family}. We calculate this map with the three
octonionic slots separately labelled, the trace metric, the arbitrary
scalar diagonal, and the original determinant coefficient two.

\subsubsection{Domain and the original matrix}
\label{h24:ncs:sec:domain}

Use the ordered quaternion subalgebra
\[
 \mathbb H=\mathbb R\langle1,i,j,k\rangle\subset\mathbb O,
 \qquad i^2=j^2=k^2=ijk=-1.
\]
For $a=(a_0,a_1,a_2,a_3)$, write
\begin{equation}
 \begin{gathered}
 v=a_1i+a_2j+a_3k,\qquad
 t=\sqrt{a_1^2+a_2^2+a_3^2},\qquad a=a_0+v,\\
 Q=a_0^2+t^2,\qquad P=a_0(a_0^2-3t^2).
 \end{gathered}
 \label{h24:ncs:eq:quaternion-coordinates}\tag{NCS1}
\end{equation}
Every square root in this section means the nonnegative real root.
Here $t$ denotes the length of the imaginary quaternion, while the
trace invariant will always be written $\operatorname{tr}X$.
The integral domain and its retained embedding are
\begin{equation}
 D_4=\{a\in\mathbb Z^4:a_0+a_1+a_2+a_3\in2\mathbb Z\},
 \qquad j(a)=\sqrt2(a,a,a)\in\mathbb O^3.
 \label{h24:ncs:eq:lattice-embedding}\tag{NCS2}
\end{equation}
The literal intersection theorem
\cref{wcl:thm:literal-intersection} identifies $j(D_4)$ with the
common lattice of the two specified marked Niemeier/Leech embeddings
and the specified Wilson embedding. Its original Niemeier coordinates
are
\begin{equation}
 x_{i,m}=(-1)^m a_{i-1}\quad(1\le i\le4,\ 0\le m\le5),
 \qquad d=0.
 \label{h24:ncs:eq:niemeier-coordinates}\tag{NCS3}
\end{equation}
The six coordinates in each $A_5$ block sum to zero. Pairing with
the retained Weyl vector gives $3\sum_i a_i$, which is divisible
by six for $a\in D_4$; hence these vectors belong to the retained
neighbor sublattice $K$. The full Wilson membership certificate and
both intersection equalities are proved in
\cref{wcl:thm:literal-intersection}.

For arbitrary $\mu\in\mathbb R$, define
\begin{equation}
 Y(a)=\begin{pmatrix}
 0&\sqrt2a&\sqrt2\bar a\\
 \sqrt2\bar a&0&\sqrt2a\\
 \sqrt2a&\sqrt2\bar a&0
 \end{pmatrix},\qquad
 X(\mu,a)=\mu I+Y(a).
 \label{h24:ncs:eq:original-matrix}\tag{NCS4}
\end{equation}
This is the original cyclic convention
$(q_1,q_2,q_3)=(\sqrt2a,\sqrt2a,\sqrt2a)$, including the
conjugated upper-right entry. On $J=\mathfrak h_3(\mathbb O)$,
the product is $A\circ B=(AB+BA)/2$ and the metric is
$G(A,B)=\operatorname{tr}(A\circ B)$. Consequently
\begin{equation}
 \|j(a)\|^2=6Q,\qquad r^2:=G(Y,Y)=12Q.
 \label{h24:ncs:eq:two-metrics}\tag{NCS5}
\end{equation}
This retains the factor two between the two metrics. The map
$(\mu,a)\mapsto X(\mu,a)$ from $\mathbb R\times D_4$ to $J$
is injective: its inverse on its image is
$\mu=\operatorname{tr}X/3$ and $a=X_{12}/\sqrt2$.

\subsubsection{Complete determinant and invariant map}
\label{h24:ncs:sec:invariant-map}

The quaternion identity $v^2=-t^2$ and associativity give
\[
 \operatorname{Re}(a^3)=a_0^3-3a_0t^2=P.
\]
Thus the triality cubic and determinant of $Y$ are exactly
\begin{equation}
 \begin{aligned}
 \tau(j(a))
 &=\operatorname{Re}((\sqrt2a\sqrt2a)\sqrt2a)=2\sqrt2P,\\
 \det Y&=2\tau(j(a))=4\sqrt2P.
 \end{aligned}
 \label{h24:ncs:eq:cubic-determinant}\tag{NCS6}
\end{equation}
The complete invariant map is therefore
\begin{equation}
 \boxed{
 (\operatorname{tr}X,\operatorname{tr}X^2,\det X)
   =\bigl(3\mu,\ 3\mu^2+12Q,\
                     \mu^3-6\mu Q+4\sqrt2P\bigr).
 }
 \label{h24:ncs:eq:full-invariants}\tag{NCS7}
\end{equation}
Indeed each of the three squared off-diagonal norms is $2Q$.
The negative terms in the original determinant formula thus sum
to $-6\mu Q$, while the mixed term is precisely
\eqref{h24:ncs:eq:cubic-determinant}. This proves the displayed
invariant map with every diagonal variable and original lattice
scale retained. The characteristic polynomial of $Y$ is
\begin{equation}
 p_Y(z)=z^3-6Qz-4\sqrt2P.
 \label{h24:ncs:eq:characteristic-polynomial}\tag{NCS8}
\end{equation}

\subsubsection{Exact eigenvalues, projectors, and directional inverse}
\label{h24:ncs:sec:projectors}

The three labelled eigenvalues of $Y$ are
\begin{equation}
 \nu_0=2\sqrt2a_0,\qquad
 \nu_1=-\sqrt2a_0-\sqrt6t,\qquad
 \nu_2=-\sqrt2a_0+\sqrt6t.
 \label{h24:ncs:eq:labelled-eigenvalues}\tag{NCS9}
\end{equation}
Their sum is zero, their squared sum is $12Q$, and their product
is $4\sqrt2P$. Direct multiplication of
$\prod_{j=0}^2(z-\nu_j)$ gives
\eqref{h24:ncs:eq:characteristic-polynomial}, proving that these are
the eigenvalues. The eigenvalues of $X$ are $\mu+\nu_j$, with
$\mu$ arbitrary.

For $t>0$, let $n=v/t$, so $n^2=-1$ and $v=tn$.
The plane $\mathbb C_n=\mathbb R\langle1,n\rangle$ is a
commutative associative subalgebra. Set
\begin{equation}
 \begin{gathered}
 T=\begin{pmatrix}0&1&0\\0&0&1\\1&0&0\end{pmatrix},
 \qquad T^3=I,\\
 \omega_n=-\frac12+\frac{\sqrt3}{2}n,\qquad
 \omega_n^3=1,\qquad1+\omega_n+\omega_n^2=0.
 \end{gathered}
 \label{h24:ncs:eq:cyclic-matrix}\tag{NCS10}
\end{equation}
Every entry of \eqref{h24:ncs:eq:original-matrix} lies in
$\mathbb C_n$, so associative matrix multiplication is legitimate
throughout this calculation. Define
\begin{equation}
 \Pi_j(n)=\frac13\left(I+\omega_n^{-j}T
                               +\omega_n^{-2j}T^2\right),
 \qquad j=0,1,2.
 \label{h24:ncs:eq:fourier-projectors}\tag{NCS11}
\end{equation}
These matrices are Hermitian: $T^*=T^2$,
$\bar\omega_n=\omega_n^{-1}$, and the two coefficients are
conjugates. Each trace is one. To verify all projector identities,
use
\[
 \sum_{k=0}^2\omega_n^{mk}=
 \begin{cases}3,&3\mid m,\\0,&3\nmid m.\end{cases}
\]
Expanding \eqref{h24:ncs:eq:fourier-projectors} gives
\begin{equation}
 \Pi_i\Pi_j=\delta_{ij}\Pi_j,\qquad
 \sum_{j=0}^2\Pi_j=I,\qquad T\Pi_j=\omega_n^j\Pi_j.
 \label{h24:ncs:eq:projector-identities}\tag{NCS12}
\end{equation}
Furthermore $Y=\sqrt2(aT+\bar aT^2)$, whence
\begin{equation}
 \begin{aligned}
 Y\Pi_j&=\sqrt2(a\omega_n^j+\bar a\omega_n^{2j})\Pi_j
            =\nu_j\Pi_j,\\
 X&=\sum_{j=0}^2(\mu+\nu_j)\Pi_j.
 \end{aligned}
 \label{h24:ncs:eq:spectral-decomposition}\tag{NCS13}
\end{equation}
Multiplication in $\mathbb C_n$ gives exactly the three signs in
\eqref{h24:ncs:eq:labelled-eigenvalues}. Thus
\eqref{h24:ncs:eq:fourier-projectors} gives an actual map into the
ordered Jordan-frame space $F_4/\operatorname{Spin}(8)$.

When the eigenvalues are pairwise distinct, these are also the
intrinsic spectral projectors:
\begin{equation}
 \Pi_j=
 \frac{(Y-\nu_kI)\circ(Y-\nu_lI)}
      {(\nu_j-\nu_k)(\nu_j-\nu_l)},
 \qquad\{j,k,l\}=\{0,1,2\}.
 \label{h24:ncs:eq:intrinsic-projectors}\tag{NCS14}
\end{equation}
Substitution of \eqref{h24:ncs:eq:spectral-decomposition}, together
with \eqref{h24:ncs:eq:projector-identities}, proves this formula.
Its denominators are retained; their exceptional loci are
computed below.

The labelled frame retains the direction $n$. Specifically,
\begin{equation}
 (\Pi_1)_{12}=-\frac16-\frac{\sqrt3}{6}n,
 \qquad n=-2\sqrt3\,\operatorname{Im}(\Pi_1)_{12}.
 \label{h24:ncs:eq:directional-inverse}\tag{NCS15}
\end{equation}
Together with
$a_0=\nu_0/(2\sqrt2)$,
$t=(\nu_2-\nu_1)/(2\sqrt6)$, and $v=tn$, this recovers the
original quaternion from the labelled spectrum and frame. The
scalar parameter is recovered from $\operatorname{tr}X/3$.
Thus the matrix and frame lift has an exact inverse on its
labelled image. The invariant triple
\eqref{h24:ncs:eq:full-invariants} alone loses $n$; the lift
\eqref{h24:ncs:eq:fourier-projectors} records it explicitly.

At $t=0$, one has $Y=\sqrt2a_0(T+T^2)$.
The rank-one projector $\Pi_0=(I+T+T^2)/3$ remains fixed, and
$I-\Pi_0$ is the rank-two projector of the repeated eigenvalue
$-\sqrt2a_0$ when $a_0\ne0$. No direction $n$ is assigned to a
zero imaginary part. At $a=0$, the matrix is the scalar $\mu I$.

\subsubsection{Singular loci and the exact Cartan angle}
\label{h24:ncs:sec:singular-angle}

The discriminant of the original polynomial
\eqref{h24:ncs:eq:characteristic-polynomial}, with its full scalar
factor, is
\begin{equation}
 \boxed{
 \prod_{i<j}(\nu_i-\nu_j)^2
       =864(Q^3-P^2)
       =864t^2(3a_0^2-t^2)^2.
 }
 \label{h24:ncs:eq:discriminant}\tag{NCS16}
\end{equation}
Indeed expansion of
$(a_0^2+t^2)^3-a_0^2(a_0^2-3t^2)^2$ gives
$t^2(3a_0^2-t^2)^2$. The first equality also follows by
multiplying the three differences in
\eqref{h24:ncs:eq:labelled-eigenvalues}. The spectrum is simple
exactly when $t>0$ and $t^2\ne3a_0^2$.

For nonzero $a$, use the Cartan polynomial of
\eqref{h24:ocb:eq:cartan-polynomial} in the original trace metric.
Equations~\eqref{h24:ncs:eq:two-metrics}--\eqref{h24:ncs:eq:cubic-determinant}
give
\begin{equation}
 \begin{gathered}
 C(Y)=3\sqrt6\det Y=24\sqrt3P,\qquad r^2=12Q,\\
 \boxed{C(Y)=r^3\cos3\theta,\qquad
 \theta=\frac13\arccos\left(\frac{P}{Q^{3/2}}\right)
                              \in[0,\pi/3].}
 \end{gathered}
 \label{h24:ncs:eq:cartan-angle}\tag{NCS17}
\end{equation}
The discriminant identity proves that the arccos argument belongs
to $[-1,1]$. This is an identity for $C$ with the whole amplitude
retained: $Y$ still has radius $r=\sqrt{12Q}$, and its scalar-shifted
determinant is still \eqref{h24:ncs:eq:full-invariants}.

To link this to the quaternionic argument, define
$\varphi\in[0,\pi]$ by
\[
 a_0=\sqrt Q\cos\varphi,\qquad t=\sqrt Q\sin\varphi.
\]
When $t>0$, the inverse is
$a=\sqrt Q(\cos\varphi+n\sin\varphi)$, retaining
$Q,\varphi,n$. The identity
$4\cos^3\varphi-3\cos\varphi=\cos3\varphi$ gives
\begin{equation}
 \cos3\theta=\cos3\varphi.
 \label{h24:ncs:eq:quaternion-angle}\tag{NCS18}
\end{equation}
The ordered Cartan chamber thus folds the quaternionic angular
variable by this explicit triple-angle map. The labelled frame,
\eqref{h24:ncs:eq:labelled-eigenvalues}, and
\eqref{h24:ncs:eq:directional-inverse} retain the inverse information.
This quaternionic angle is defined by
\eqref{h24:ncs:eq:quaternion-coordinates}; its definition makes no
identification with a cusp angular coordinate.

For $t>0$, the decreasing orders of the labelled spectrum are
\begin{equation}
 \begin{array}{c|c}
 a_0>t/\sqrt3&\nu_0>\nu_2>\nu_1\\
 -t/\sqrt3<a_0<t/\sqrt3&\nu_2>\nu_0>\nu_1\\
 a_0<-t/\sqrt3&\nu_2>\nu_1>\nu_0.
 \end{array}
 \label{h24:ncs:eq:ordered-chambers}\tag{NCS19}
\end{equation}
The differences
$\nu_0-\nu_2=\sqrt2(3a_0-\sqrt3t)$ and
$\nu_0-\nu_1=\sqrt2(3a_0+\sqrt3t)$ prove each row.
At $a_0=t/\sqrt3$, precisely $\nu_0=\nu_2$; at
$a_0=-t/\sqrt3$, precisely $\nu_0=\nu_1$. Thus these
rows specify how to order the projectors before applying
\eqref{h24:ocb:eq:orbit-map}. At the two double-eigenvalue loci,
the formulas \eqref{h24:ncs:eq:fourier-projectors} remain defined,
and \eqref{h24:ncs:eq:projector-identities}--
\eqref{h24:ncs:eq:spectral-decomposition} give repeated-eigenspace
projectors $\Pi_0+\Pi_2$ and $\Pi_0+\Pi_1$, respectively.
Their complementary simple-eigenvalue projectors are $\Pi_1$
and $\Pi_2$. This follows directly by grouping the equal
coefficients in \eqref{h24:ncs:eq:spectral-decomposition}; no
vanishing denominator in \eqref{h24:ncs:eq:intrinsic-projectors}
is used.

\subsubsection{Actual integral witnesses in the orbit family}
\label{h24:ncs:sec:integral-witnesses}

The image proof of \cref{d4i:thm:complete-image} supplies three
original-coordinate lattice lines:
\begin{equation}
 \begin{gathered}
 L_\varepsilon(h)=(-3h-\varepsilon,h,h,h+\varepsilon),
       \qquad\varepsilon=\pm1,\\
 L_0(h)=(-3h,h+1,h-1,h),\qquad h\in\mathbb Z.
 \end{gathered}
 \label{h24:ncs:eq:integral-lines}\tag{NCS20}
\end{equation}
Their coordinate sums are zero. Direct expansion gives
\begin{equation}
 \begin{array}{c|c|c}
 a&P(a)&Q(a)\\\hline
 L_\varepsilon(h)&6h+2\varepsilon&12h^2+8\varepsilon h+2\\
 L_0(h)&18h&12h^2+2.
 \end{array}
 \label{h24:ncs:eq:line-values}\tag{NCS21}
\end{equation}
The eigenvalues, projectors, full radius, and scalar-shifted
determinant of each actual common-lattice vector are therefore
given by \eqref{h24:ncs:eq:full-invariants}--
\eqref{h24:ncs:eq:cartan-angle}, evaluated at
\eqref{h24:ncs:eq:line-values} and its original imaginary vector $v$.
The Cartan angles satisfy
\begin{equation}
 \begin{aligned}
 \cos3\theta_\varepsilon(h)
    &=\frac{6h+2\varepsilon}
                  {(12h^2+8\varepsilon h+2)^{3/2}},\\
 \cos3\theta_0(h)&=\frac{18h}{(12h^2+2)^{3/2}}.
 \end{aligned}
 \label{h24:ncs:eq:line-angles}\tag{NCS22}
\end{equation}
As $|h|\to\infty$, both ratios tend to zero: the numerators
have degree one, while their positive denominators grow as
$|h|^3$. Hence the angles tend to $\pi/6$, and simultaneously
$r^2=12Q\to\infty$. The original matrices and their radii
remain those in \eqref{h24:ncs:eq:original-matrix}--
\eqref{h24:ncs:eq:two-metrics}.

On the shared direction $a=h(-3,1,1,1)$ one has $P=0$ and
$Q=12h^2$. For $h\ne0$, the decreasing spectrum is exactly
\begin{equation}
 (6\sqrt2|h|,\ 0,\ -6\sqrt2|h|),\qquad
 r=12|h|,\qquad\theta=\pi/6.
 \label{h24:ncs:eq:zero-cubic-direction}\tag{NCS23}
\end{equation}
The discriminant \eqref{h24:ncs:eq:discriminant} is strictly positive,
so this zero-cubic direction lies on a regular Cartan orbit.
The spectrum of $X(\mu,a)$ is obtained by adding $\mu$ to each
displayed eigenvalue.

\paragraph{Sources and scope of the bridge.}
The literal-intersection theorem \cref{wcl:thm:literal-intersection}
supplies the marked Niemeier/Leech and Wilson membership proofs for
\eqref{h24:ncs:eq:lattice-embedding}--\eqref{h24:ncs:eq:niemeier-coordinates};
the mixed cubic is \eqref{wcl:eq:common-cubic}. Wilson's construction
is \cite{esn:Wilson2009}. The full invariant quotient, retained
mean and radius, ordered spectral diffeomorphism, and focal maps are
proved in \cref{h24:ocb:sec:full-family}, using the flag geometry of
\cite[Sections~2.3--2.6]{h24:MareWillems2013} and
\cite[Proposition~2.8.2 and Theorems~2.7.4, 2.8.3]{h24:Yokota2009}.
The complete image theorem \cref{d4i:thm:complete-image} proves
\[
 P(D_4)=18\mathbb Z\cup(6\mathbb Z+2)\cup(6\mathbb Z-2)
\]
using exactly \eqref{h24:ncs:eq:integral-lines}. This section provides
its matrix-valued and frame-valued lift with an inverse on the
labelled image; it makes no assertion that every Cartan orbit meets
this rank-four lattice. The construction is independent of the
assertion that the original completed complex threefold is
diffeomorphic to $S^6$. No interacting quantum Hamiltonian is
specified here. The retained cusp bundle has its separate explicit
triality-subgroup map; its source angular coordinate and the
quaternionic angular coordinate above retain their respective
definitions.



% END ES-NIEMEIER HOST-ONLY INPUT.

\begingroup
\sloppy
\raggedright
\begin{thebibliography}{99}

\bibitem{h24:Milnor1965}
J.~W. Milnor, \emph{Topology from the Differentiable Viewpoint},
University Press of Virginia, 1965, \S7, Theorem B and Lemma 4,
pp.~42--50.
\url{https://www.maths.ed.ac.uk/~aar/papers/milnortop.pdf}.



\bibitem{h24:AlpogeS6_2026}
L.~Alp\"oge, with the disclosed Fable/Claude contribution,
\emph{The \((3,4,\infty)\) modular family of \(2\)-tori, completed at its
three special points, is a complex structure on \(S^6\)}, together with
the citable canonical audit and source bundle, Zenodo record 22238779
(2026), \href{https://doi.org/10.5281/zenodo.22238779}
{doi:10.5281/zenodo.22238779}.  Local audit locators are below
\texttt{supporting\_materials/workbench/research/}; the one-copy quotient
and framing discussion are in
\texttt{cusp\_sphere\_route\_strengthening.tex} and
\texttt{cusp\_s12\_carrier\_homotopy.tex}.

\bibitem{h24:ConradSL2Z}
K.~Conrad, \emph{SL\(_2(\mathbb Z)\)}, author-hosted expository notes,
Theorem~1.1, pp.~1--2; Theorem~C.1, pp.~18--19,
\href{https://kconrad.math.uconn.edu/blurbs/grouptheory/SL(2,Z).pdf}
{author's PDF}, accessed 5 September 2026.

\bibitem{h24:BorelSerre1953}
A.~Borel and J.-P.~Serre,
\emph{Groupes de Lie et puissances r\'eduites de Steenrod},
Amer. J. Math. \textbf{75} (1953), no.~3, 409--448;
Sections~12 and~15,
\href{https://doi.org/10.2307/2372495}{doi:10.2307/2372495}.

\bibitem{h24:HatcherSpectralSequences}
A.~Hatcher,
\emph{Spectral Sequences in Algebraic Topology}, Chapter~5 draft,
Section~5.1; Theorem~5.41(b), p.~573, and the suspension calculation,
pp.~577--578,
\href{https://pi.math.cornell.edu/~hatcher/AT/ATch5.pdf}
{author-hosted chapter PDF}.

\bibitem{h24:Hatcher2002}
A.~Hatcher,
\emph{Algebraic Topology},
Cambridge University Press, Cambridge, 2002; the opening
low-dimensional homotopy-group table, p.~339,
\href{https://pi.math.cornell.edu/~hatcher/AT/ATpage.html}
{author's electronic edition}.

\bibitem{h24:Husemoller1994}
D.~Husemoller,
\emph{Fibre Bundles}, third edition,
Graduate Texts in Mathematics, vol.~20, Springer, New York, 1994;
Chapter~17, Section~7, Theorem~7.2, pp.~274--275,
\href{https://doi.org/10.1007/978-1-4757-2261-1}
{doi:10.1007/978-1-4757-2261-1}.

\bibitem{h24:IsaksenWangXu2023}
D.~C.~Isaksen, G.~Wang, and Z.~Xu,
\emph{Stable homotopy groups of spheres: from dimension 0 to 90},
Publ. Math. IH\'ES \textbf{137} (2023), 107--243; Table~1, p.~113,
\href{https://doi.org/10.1007/s10240-023-00139-1}
{doi:10.1007/s10240-023-00139-1}.

\bibitem{h24:KomoriMatsumotoTsumura2023}
Y.~Komori, K.~Matsumoto, and H.~Tsumura,
\emph{The Theory of Zeta-Functions of Root Systems},
Springer Monographs in Mathematics, Springer, Singapore, 2023;
Section~10.1, pp.~213 and~216,
\href{https://doi.org/10.1007/978-981-99-0910-0}
{doi:10.1007/978-981-99-0910-0}.

\bibitem{h24:MareWillems2013}
A.-L.~Mare and M.~Willems,
\emph{Topology of the octonionic flag manifold},
M\"unster J. Math. \textbf{6} (2013), 483--523,
\href{https://www.uni-muenster.de/FB10/mjm/vol_6/mjm_vol_6_14.pdf}
{published article};
\href{https://arxiv.org/abs/0809.4318}{arXiv:0809.4318}.

\bibitem{h24:Toda1962}
H.~Toda,
\emph{Composition Methods in Homotopy Groups of Spheres},
Annals of Mathematics Studies \textbf{49},
Princeton University Press, Princeton, NJ, 1962;
Proposition~3.1,
\href{https://doi.org/10.1515/9781400882625}
{doi:10.1515/9781400882625}.

\bibitem{h24:retbib:Canonical}
L.~Alp\"oge, with the disclosed Fable/Claude contribution,
\emph{The \((3,4,\infty)\) modular family of \(2\)-tori, completed at
its three special points, is a complex structure on \(S^6\)},
canonical audit and reconstruction source bundle, Zenodo record 22238779
(2026), original manuscript \S3.7 and \S\S4.1--4.5.
\url{https://doi.org/10.5281/zenodo.22238779}.

\bibitem{h24:retbib:ConeRoutes}
\emph{Real \(S^{12}\) and \(S^{24}\) Routes from Alp\"oge's
Candidate \((3,4,\infty)\) Two-Torus Fibration},
companion research draft in the preceding canonical source bundle;
native quotient, localized cusp map, and framed-cusp calculation.
Current proof-bearing edition:
\url{https://www.overleaf.com/read/rtmyqxyrzprn}.

\bibitem{h24:retbib:Milnor1965}
J.~W.~Milnor, \emph{Topology from the Differentiable Viewpoint},
University Press of Virginia, 1965, \S7, Theorem B and Lemma 4.
\url{https://www.maths.ed.ac.uk/~aar/papers/milnortop.pdf}.

\bibitem{h24:retbib:Pontryagin1950}
L.~S.~Pontryagin,
\emph{Homotopy classification of mappings of an \((n+2)\)-dimensional
sphere on an \(n\)-dimensional one},
Dokl. Akad. Nauk SSSR \textbf{70} (1950), 957--959.
\url{https://www.sas.rochester.edu/mth/sites/doug-ravenel/otherpapers/pont3.pdf}.

\bibitem{h24:retbib:Pontryagin1955}
L.~S.~Pontryagin,
\emph{Smooth manifolds and their applications in homotopy theory},
Trudy Mat. Inst. Steklov \textbf{45} (1955);
English translation, Amer. Math. Soc. Transl. (2) \textbf{11} (1959).
\url{https://www.sas.rochester.edu/mth/sites/doug-ravenel/otherpapers/pont4.pdf}.

\bibitem{h24:retbib:Johnson}
D.~Johnson, \emph{Spin structures and quadratic forms on surfaces},
J. London Math. Soc. (2) \textbf{22} (1980), 365--373.

\bibitem{h24:retbib:DebrayGunningham}
A.~Debray and S.~Gunningham,
\emph{The Arf--Brown TQFT of pin-minus surfaces},
arXiv:1803.11183, especially \S3.1.1 and Proposition 5.1.
\url{https://adebray.github.io/papers/arf_brown.pdf}.

\bibitem{h24:retbib:Hatcher}
A.~Hatcher, \emph{Algebraic Topology}, Cambridge University Press,
2002; group table p.~339, Corollary 4.24, p.~360, and pp.~384--385.
\url{https://pi.math.cornell.edu/~hatcher/AT/AT.pdf}.

\bibitem{h24:hbridgebib:HBRetained}
\emph{The retained-angle cusp quotient has nonzero top-cell class},
complete derivative-framing theorem and its based
Pontryagin--Thom proof, 5 September 2026;
companion to the canonical cusp/sphere-route draft.
Current cumulative source collection:
\url{https://www.overleaf.com/read/rtmyqxyrzprn}.

\bibitem{h24:hbridgebib:HBHigher}
\emph{The higher torus wrap: obstruction, exact local lift, and the
Jordan--triality rung},
canonical higher-rung research draft, section on peripheral
modulus and the mod--24 boundary, 5 September 2026.
Current cumulative source collection:
\url{https://www.overleaf.com/read/rtmyqxyrzprn}.

\bibitem{h24:hbridgebib:HBHatcher}
A.~Hatcher, \emph{Algebraic Topology}, Cambridge University Press,
2002, Corollary~4.24, p.~360, stable homotopy groups and the
integral relation \(\eta^3=12\nu\), pp.~384--385,
and Proposition~4.56 with proof on p.~388.
\url{https://pi.math.cornell.edu/~hatcher/AT/AT.pdf}.

% Host-only helper: the workbench supplies it in its central bibliography.
\bibitem{h24:Yokota2009}
I.~Yokota, \emph{Exceptional Lie groups}, 2009;
Sections~1.10, 1.14, 1.16 and Theorem~2.7.1,
\url{https://arxiv.org/abs/0902.0431}.


\end{thebibliography}
\endgroup

\subsection{The Koide cone and the canonical period discriminant}
\label{koide-s6-subsection-ambient-discriminant}

For \(r=(r_1,r_2,r_3)\in\mathbb R^3\), set
\begin{equation}
 Q_K(r)
 :=\frac{2}{3}(r_1+r_2+r_3)^2-(r_1^2+r_2^2+r_3^2)
 \label{koide-s6-eq-koide-form}
\end{equation}
and introduce the orthogonal coordinates
\begin{equation}
 \begin{pmatrix}h\\x\\y\end{pmatrix}
 =\mathcal F\begin{pmatrix}r_1\\r_2\\r_3\end{pmatrix},
 \qquad
 \mathcal F=
 \begin{pmatrix}
  1/\sqrt3&1/\sqrt3&1/\sqrt3\\
  1/\sqrt2&0&-1/\sqrt2\\
  1/\sqrt6&-2/\sqrt6&1/\sqrt6
 \end{pmatrix}.
 \label{koide-s6-eq-orthogonal-coordinates}
\end{equation}

We use exactly the canonical period convention
\begin{equation}
 Z=\begin{pmatrix}6\mu&\tau\\ \beta&\mu\end{pmatrix},
 \qquad \Pi=[Z\mid I_2].
 \label{koide-s6-eq-canonical-complex-period}
\end{equation}
Write
\[
 t=\operatorname{Im}\tau,\qquad
 m=\operatorname{Im}\mu,\qquad
 \varrho=\operatorname{Im}\beta.
\]
Relative to the period-lattice basis
\[
 (\widehat\gamma,\widehat u,\widehat w,\widehat\delta)
\]
and the interleaved real codomain basis
\[
 (\operatorname{Re}z_1,\operatorname{Im}z_1,
  \operatorname{Re}z_2,\operatorname{Im}z_2),
\]
the realification of \(\Pi\) is
\begin{equation}
 P(\tau,\mu,\beta)=
 \begin{pmatrix}
  6\operatorname{Re}\mu&\operatorname{Re}\tau&1&0\\
  6m&t&0&0\\
  \operatorname{Re}\beta&\operatorname{Re}\mu&0&1\\
  \varrho&m&0&0
 \end{pmatrix}.
 \label{koide-s6-eq-canonical-real-period}
\end{equation}
For comparison with the Lorentz cone, define the auxiliary symmetric
coefficient matrix
\begin{equation}
 S(t,m,\varrho):=
 \begin{pmatrix}t&\sqrt6\,m\\ \sqrt6\,m&\varrho\end{pmatrix}.
 \label{koide-s6-eq-symmetric-coefficient-matrix}
\end{equation}
Thus \(S\) represents the same coefficient discriminant, but it is not the
canonical \(4\times4\) real period matrix \(P\).

\begin{theorem}[Ambient coefficient-space discriminant isomorphism]
\label{koide-s6-thm-ambient-discriminant}
The matrix \(\mathcal F\) is orthogonal, and
\begin{equation}
 Q_K(r)=h^2-x^2-y^2.
 \label{koide-s6-eq-lorentz-identity}
\end{equation}
The real-linear map
\begin{equation}
 \mathcal K(h,x,y)
 =\bigl(\operatorname{Im}\tau,
         \operatorname{Im}\mu,
         \operatorname{Im}\beta\bigr)
 =\left(h+y,\frac{x}{\sqrt6},h-y\right)
 \label{koide-s6-eq-coefficient-map}
\end{equation}
is an isomorphism from \(\mathbb R^3\) onto the ambient real coefficient
space, with inverse
\begin{equation}
 \mathcal K^{-1}(t,m,\varrho)
 =\left(\frac{t+\varrho}{2},\sqrt6\,m,
        \frac{t-\varrho}{2}\right).
 \label{koide-s6-eq-coefficient-map-inverse}
\end{equation}
For every choice of the real parts of \(\tau,\mu,\beta\),
\begin{equation}
 \det P(\tau,\mu,\beta)
 =\det S(t,m,\varrho)
 =t\varrho-6m^2.
 \label{koide-s6-eq-canonical-discriminant}
\end{equation}
Consequently, if
\((t,m,\varrho)=\mathcal K(h,x,y)\), then
\begin{equation}
 \det P(\tau,\mu,\beta)
 =\det S\bigl(\mathcal K(h,x,y)\bigr)
 =Q_K(r).
 \label{koide-s6-eq-period-koide-identity}
\end{equation}
Moreover, for \(r\in\mathbb R_{\geq0}^3\setminus\{0\}\), the equation
\(Q_K(r)=0\) is equivalent to the assertion that
\begin{equation}
 S\bigl(\mathcal K(\mathcal F(r))\bigr)
 =\begin{pmatrix}h+y&x\\x&h-y\end{pmatrix}
 \label{koide-s6-eq-positive-sector-matrix}
\end{equation}
is positive semidefinite of rank one.
\end{theorem}

\begin{proof}
The rows of \(\mathcal F\) have norm one and pairwise scalar product zero,
so \(\mathcal F\mathcal F^{\mathsf T}=I_3\).  Hence
\[
 r_1^2+r_2^2+r_3^2=h^2+x^2+y^2,
 \qquad r_1+r_2+r_3=\sqrt3\,h.
\]
Substitution in \cref{koide-s6-eq-koide-form} proves
\cref{koide-s6-eq-lorentz-identity}.

Adding and subtracting the first and third entries of
\cref{koide-s6-eq-coefficient-map}, and multiplying the second by
\(\sqrt6\), proves \cref{koide-s6-eq-coefficient-map-inverse}.  Expanding
\cref{koide-s6-eq-canonical-real-period} along its last two columns gives
\[
 \det P
 =-\det\begin{pmatrix}6m&t\\\varrho&m\end{pmatrix}
 =t\varrho-6m^2.
\]
Equivalently, in the exact canonical convention,
\(\det P=-\det\operatorname{Im}Z\).  Taking the determinant in
\cref{koide-s6-eq-symmetric-coefficient-matrix} proves the other equality
in \cref{koide-s6-eq-canonical-discriminant}, while
\[
 (h+y)(h-y)-6\left(\frac{x}{\sqrt6}\right)^2
 =h^2-x^2-y^2
\]
proves \cref{koide-s6-eq-period-koide-identity}.

Finally, let \(r\in\mathbb R_{\geq0}^3\setminus\{0\}\).  Then
\(h=(r_1+r_2+r_3)/\sqrt3>0\).  If \(Q_K(r)=0\), the symmetric matrix in
\cref{koide-s6-eq-positive-sector-matrix} has determinant zero and trace
\(2h>0\); its eigenvalues are therefore \(0\) and \(2h\), so it is positive
semidefinite of rank one.  Conversely, rank one forces its determinant,
and hence \(Q_K(r)\), to vanish.
\end{proof}

\begin{corollary}[The Koide cone is not a family of regular tori]
\label{koide-s6-cor-ambient-only}
The identification in
\cref{koide-s6-thm-ambient-discriminant} is only a real-linear isomorphism
of ambient coefficient spaces together with their quadratic discriminants.
It neither supplies the real parts of the periods nor imposes the
holomorphic and equivariance conditions of the canonical period family.
At every point of the Koide cone, every matrix of the canonical form
\cref{koide-s6-eq-canonical-real-period} with coefficient triple
\cref{koide-s6-eq-coefficient-map} satisfies \(\det P=0\).  Its columns
therefore do not span \(\mathbb R^4\), do not generate a full lattice, and
do not define a regular compact complex two-torus.  Thus the ambient
discriminant identity is not a map from Koide triples to regular fibres of
the canonical \(S^6\) family.
\end{corollary}

\begin{corollary}[Separation from the actual canonical cusp]
\label{koide-s6-cor-actual-cusp}
In the notation
\[
 L_T=T+H_T,\qquad q_T=L_T-B_T,\qquad
 D_T=L_Tq_T+6m_T^2
\]
of the positive radial cusp, the already-proved determinant formula
\cref{cusp-spec-determinant} gives
\begin{equation}
 \det P_T=-D_T<0.
 \label{koide-s6-eq-actual-cusp-determinant}
\end{equation}
Moreover, with
\(H=\operatorname{Im}h(0)\) and \(B=\operatorname{Im}b(0)\),
\cref{cusp-spec-D-expansion} gives
\begin{equation}
 D_T=T^2+(2H-B)T+O(1),
 \qquad
 \det P_T\sim-T^2.
 \label{koide-s6-eq-actual-cusp-asymptotic}
\end{equation}
Thus every actual radial cusp fibre is regular and has nonzero oriented
period determinant, whereas every point in the Koide image has determinant
zero.
\end{corollary}

\begin{proof}
The first assertion is exactly \cref{cusp-spec-determinant}; the stated
asymptotic follows from \cref{cusp-spec-D-expansion}.  Since \(D_T>0\), the
matrix \(P_T\) is invertible and its columns form the full period lattice of
the corresponding regular fibre.
\end{proof}

% Shared mathematical module. Input exactly once in each host document.
% No document class, package loads, macros, or theorem environments are defined.
% Requirements and bibliography entries: INTEGRATION_MAP.md.
% Source locators and snapshot provenance: ES_NIEMEIER_PROVENANCE.md.
% All labels and citation keys introduced here use the esn: namespace.

\section{The modulo-840 lattice, triality, and exact shell maps}
\label{esn:sec:shared}

This section constructs the original modulo-$840$ lattice $N$ with
root system $A_5^4D_4$, its full integral isometry group, and the exact
maps from denominator shells and finite prime samples to the retained
residue coordinates. The original data and source labels are recorded in
\cite{esn:Provenance}. The shared continuation proves the index-six Leech
neighbour $\Lambda$ and the marked isometry $\mathcal L$ into the three
octonion sectors. In the specified Cayley frame, put
$\tau(q)=\operatorname{Re}((q_1q_2)q_3)$ and $F=\tau\circ\mathcal L$;
the Albert determinant is $\lambda^3-\lambda(y,y)+2F(y)$.
The value-generated groups have quotient
$\mathcal V_F(\Lambda)/\mathcal V_F(N)\cong\mathbb Z/12$
(Theorem~\ref{nca:thm:value-modules}). The three residue-cycled neighbours
meet in $J$, with determinant $81$ and index $9$ in each; the fixed
fourth extension has exactly $72$ roots of type $A_2^{12}$
(Theorems~\ref{esrf:thm:common-J} and~\ref{esrf:thm:fourth-extension}).
The common lattice and this fourth extension have the exact same
value-generated group
$\frac1{27}\mathbb Z\oplus\frac{\sqrt2}{216}\mathbb Z
\oplus\frac{\sqrt3}{18}\mathbb Z$.
The inclusions from $\mathcal V_F(N)$ and into $\mathcal V_F(\Lambda)$
have matrices $\operatorname{diag}(1,4,1)$ and
$\operatorname{diag}(1,1,3)$ and quotients $\mathbb Z/4$, $\mathbb Z/3$
(Theorems~\ref{jmv:thm:JM-values} and~\ref{jmv:thm:chain}).
Wilson's construction \cite{esn:Wilson2009} is realized with both metrics
and the direction $c_{\rm W}=c^{-1}$ retained. Its fixed Gram matrix is
$3G_{E_8}$, its complement is explicitly isometric to $A_2\otimes E_8$,
and the index is $3^8$, with quotient $\mathbb F_3^8$. Its literal
intersections with $\mathcal L(N)$ and $\mathcal L(\Lambda)$ are the
same rank-four lattice of Gram matrix $6G_{D_4}$, minimum squared norm
$12$ and value-generated subgroup $4\sqrt2\mathbb Z$ for $\tau$
($8\sqrt2\mathbb Z$ for $2\tau$). This is not the actual image:
$12\sqrt2$ is not attained by $\tau$ on that intersection.
All coordinates, multiplication conventions, domain lattices and maps
are specified in the complete proofs below.

\subsection{The finite multiplicative carrier and its four fibres}
\label{esn:subsec:residues}

For an odd prime \(q\) and a unit \(r\) modulo \(q\), write
\(\left(\frac rq\right)=1\) when \(r\) is a square modulo \(q\),
and \(-1\) otherwise.  Let
\begin{equation}
 \begin{split}
 V_{840}&=\{r\in(\mathbb Z/840\mathbb Z)^\times:
                                      r\equiv1\pmod {24}\},\\
 \beta(r)&=\left(\frac{1-\left(\frac r5\right)}2,
                  \frac{1-\left(\frac r7\right)}2\right)
                          \in\mathbb F_2^2,\\
 S_{840}&=\ker\beta,\qquad F_\eta=\beta^{-1}(\eta)
                                      \quad(\eta\in\mathbb F_2^2).
 \end{split}
 \label{esn:eq:residue-carrier}
\end{equation}
Multiplication is used in \(V_{840}\) and addition in \(\mathbb F_2^2\).

\begin{theorem}[The marked cyclic fibre]
\label{esn:thm:cyclic-fibre}
The map \(\beta\) is an onto group homomorphism, each \(F_\eta\) has
six elements, and
\begin{equation}
 \begin{split}
 S_{840}&=\{1,121,169,289,361,529\}=\langle289\rangle,\\
 (289^j)_{j=0}^5&=(1,289,361,169,121,529)\pmod {840}.
 \end{split}
 \label{esn:eq:ordered-fibre}
\end{equation}
Each fibre is a torsor under multiplication by \(S_{840}\).  In
particular the marking
\begin{equation}
 \kappa:S_{840}\xrightarrow{\ \cong\ }\mathbb Z/6\mathbb Z,
 \qquad r=289^{\kappa(r)}\pmod {840},
 \label{esn:eq:kappa}
\end{equation}
is an isomorphism of groups.
\end{theorem}

\begin{proof}
The Chinese remainder map modulo \(24,5,7\), whose moduli are pairwise
coprime, identifies \(V_{840}\) with
\((\mathbb Z/5)^\times\times(\mathbb Z/7)^\times\).  The powers of
\(2\) modulo \(5\) are \(1,2,4,3\), and the powers of \(3\) modulo
\(7\) are \(1,3,2,6,4,5\); these prove the respective cyclic orders
four and six.  In a cyclic group of either even order, the squares are
exactly the even powers.  Therefore each Legendre character is the
homomorphism recording exponent parity and is onto.  Their product is
onto \(\mathbb F_2^2\).  Its kernel has
\((4/2)(6/2)=6\) elements and every fibre has that same cardinality.

Successive multiplication of the six displayed residues by \(289\)
gives the next entry and then returns from \(529\) to \(1\).
They are distinct.  Moreover \(289\equiv1\pmod {24}\),
\(289\equiv4\pmod5\), and \(289\equiv2\pmod7\), so it belongs to
the simultaneous square kernel.  Its six powers therefore exhaust that
kernel.  If \(r,s\in F_\eta\), then \(sr^{-1}\in S_{840}\) and
\(s=(sr^{-1})r\), uniquely.  This proves the torsor assertion and all
claims about \(\kappa\).
\end{proof}

Fix the coordinate order
\[
 (\eta_1,\eta_2,\eta_3,\eta_4)=(00,10,01,11)
\]
and choose one origin \(f_i\in F_{\eta_i}\) in each fibre, with
\(f_1=1\).  The six ordered elements of its fibre are
\(f_i289^j\), \(0\le j\le5\).  These origins are part of the
coordinate identifications used below; no selection of them is inferred
from an individual denominator certificate.

\subsection{The root lattice and its exact discriminant form}
\label{esn:subsec:root-and-discriminant}

For a finite set \(F\), give the free abelian group
\(\mathbb Z[F]=\bigoplus_{r\in F}\mathbb Ze_r\) the metric
\((e_r,e_s)=1\) for \(r=s\) and zero otherwise.  Put
\[
 A(F)=\left\{\sum_{r\in F}n_re_r:\sum_{r\in F}n_r=0\right\}.
\]
The ordered fibre coordinates identify each \(A(F_{\eta_i})\)
isometrically with
\[
 A=\{(n_0,\ldots,n_5)\in\mathbb Z^6:\textstyle\sum_j n_j=0\}.
\]
In a separate Euclidean space with orthonormal basis
\(e_1,e_2,e_3,e_4\), define
\begin{equation}
 D=\{x\in\mathbb Z^4:x_1+x_2+x_3+x_4\equiv0\pmod2\},
 \qquad R_{\mathrm{ES}}=
       \left(\bigoplus_{i=1}^4 A(F_{\eta_i})\right)\mathbin{\perp}D.
 \label{esn:eq:root-lattice}
\end{equation}
All five summands are orthogonal with precisely their displayed metrics.
The symbol \(R_{\mathrm{ES}}\) denotes this lattice; the unadorned
\(R\) later used for an integer shell modulus has a different declared
domain.

For a full integral lattice \(L\) in its real span, let
\[
 L^\vee=\{x:(x,L)\subset\mathbb Z\},\qquad D_L=L^\vee/L.
\]
If \(L\) is even integral, its discriminant quadratic form is
\(q_L(x+L)=(x,x)\pmod {2\mathbb Z}\).  This is well-defined because
replacing \(x\) by \(x+\ell\) changes its squared norm by
\(2(x,\ell)+(\ell,\ell)\in2\mathbb Z\).

\begin{theorem}[Root metric and marked discriminant]
\label{esn:thm:root-discriminant}
The lattice \(R_{\mathrm{ES}}\) is even integral of rank 24 and has
Gram matrix
\[
 G_{R_{\mathrm{ES}}}=G_A^{\oplus4}\oplus G_D,
\]
where
\[
 G_A=\begin{pmatrix}
 2&-1&0&0&0\\-1&2&-1&0&0\\0&-1&2&-1&0\\
 0&0&-1&2&-1\\0&0&0&-1&2
 \end{pmatrix},\qquad
 G_D=\begin{pmatrix}
 2&-1&0&0\\-1&2&-1&-1\\0&-1&2&0\\0&-1&0&2
 \end{pmatrix}.
\]
Their determinants are 6 and 4, respectively, and
\(\det R_{\mathrm{ES}}=6^4\cdot4=5184\).

In each identified \(A\) summand use
\(\lambda=(5,-1,-1,-1,-1,-1)/6\) as the discriminant generator.
For \(D\), use the representatives
\begin{equation}
 v=(1,0,0,0),\qquad s=\tfrac12(1,1,1,1),\qquad
 c=\tfrac12(-1,1,1,1)
 \label{esn:eq:d4-representatives}
\end{equation}
and the exact marking \(10\mapsto v+D\), \(01\mapsto s+D\),
\(11\mapsto c+D\).  Then
\begin{equation}
 \begin{split}
 D_{R_{\mathrm{ES}}}&=(\mathbb Z/6\mathbb Z)^4\oplus\mathbb F_2^2,\\
 q_{R_{\mathrm{ES}}}(x_1,x_2,x_3,x_4;y)
   &=\frac56\sum_{i=1}^4x_i^2+\mathbf1_{y\ne0}
                                              \pmod {2\mathbb Z}.
 \end{split}
 \label{esn:eq:full-discriminant}
\end{equation}
\end{theorem}

\begin{proof}
For a six-element ordered fibre the five differences
\(e_j-e_{j+1}\), \(0\le j\le4\), form a basis of its augmentation
lattice: the coefficients of a zero-sum vector in this basis are its
first five successive partial sums.  Dot products give \(G_A\).
The determinant \(d_n\) of the analogous tridiagonal \(n\)-by-\(n\)
matrix satisfies \(d_n=2d_{n-1}-d_{n-2}\), \(d_0=1,d_1=2\), by
first-row expansion; induction gives \(d_n=n+1\).
The four vectors
\(e_1-e_2,e_2-e_3,e_3-e_4,e_3+e_4\) lie in \(D\) and their
coordinate determinant has absolute value 2.  Their integer span
therefore has index 2 in \(\mathbb Z^4\).  The parity homomorphism
\(\mathbb Z^4\to\mathbb F_2\), \(x\mapsto\sum x_i\pmod2\), is
onto and has kernel \(D\), so this span equals \(D\).
Their dot products give \(G_D\), whose determinant is the square of
the coordinate determinant, namely 4.  Integer coordinates make all
pairings integral.  In either lattice the squared norm is congruent
modulo two to the coordinate sum, hence is even.  Rank and determinant
now follow from the orthogonal direct sum.

The vector \(\lambda\) has sum zero and integral pairing with each
augmentation root.  Its norm is \(5/6\), and \(k\lambda\in A\)
holds exactly when \(6\) divides \(k\).  Thus it has exact order six
in \(D_A\).  In an integral basis of an integral lattice with Gram
matrix \(G\), the dual coordinates are \(G^{-1}\mathbb Z^n\);
therefore \(|L^\vee/L|=|\det G|\).  Since \(|D_A|=6\), this
element generates the full discriminant group and gives its stated form.

A vector pairing integrally with all \(e_i-e_j\) has common
fractional parts in its four coordinates.  Pairing also with
\(e_i+e_j\) forces this common fractional part to be either zero or
one half.  Conversely every entirely integral or entirely half-integral
vector pairs integrally with every root of \(D\).  There are two
integer classes according to sum parity, and two half-integer classes
according to the parity of their difference from \(s\).  The four
representatives \(0,v,s,c\) are therefore all the classes, and their
squared norms are \(0,1,1,1\).  Their addition is as marked because
\(v+s-c=(2,0,0,0)\in D\).  The dual and quadratic form of an
orthogonal sum are the direct sums of those of its summands, which proves
\eqref{esn:eq:full-discriminant}.
\end{proof}

\subsection{The index-72 glue and all glue-class minima}
\label{esn:subsec:glue}

Define the binary and ternary codes
\begin{equation}
 \begin{split}
 C_2&=\operatorname{span}_{\mathbb F_2}
       \{(1111;00),(1100;10),(1010;01)\},\\
 C_3&=\operatorname{span}_{\mathbb F_3}
       \{z_1=(1,1,1,0),\ z_2=(1,2,0,1)\}.
 \end{split}
 \label{esn:eq:codes}
\end{equation}
The semicolon distinguishes the four binary A5 coordinates from the
two D4 discriminant coordinates.  Let
\begin{equation}
 \begin{split}
 \mathcal H&=\{(3u+2z;y):(u;y)\in C_2, z\in C_3\}
                                  \subset D_{R_{\mathrm{ES}}},\\
 N&=\{x\in R_{\mathrm{ES}}^\vee:
                                x+R_{\mathrm{ES}}\in\mathcal H\}.
 \end{split}
 \label{esn:eq:glue-and-overlattice}
\end{equation}
Here \(3u+2z\) is computed coordinatewise modulo six using the
integer representatives of \(u\) and \(z\).

\begin{theorem}[Even unimodular completion with exactly 144 roots]
\label{esn:thm:niemeier}
The subgroup \(\mathcal H\) is isotropic of order 72.  Its inverse
image \(N\) is a positive-definite even unimodular lattice of rank 24,
with \([N:R_{\mathrm{ES}}]=72\).  Among the 71 nonzero glue classes,
46 have minimum squared norm 4 and 25 have minimum squared norm 6.
The roots of \(N\) are exactly those of \(R_{\mathrm{ES}}\), with
root system \(A_5^4D_4\): 30 roots in each A5 summand and 24 in D4.
\end{theorem}

\begin{proof}
The binary generators are independent, hence give eight words.  Listing
their sums shows that the seven nonzero words comprise one word
\((1111;00)\) and six words whose first part has weight two and whose
second part is nonzero.  A ternary word is
\[
 az_1+bz_2=(a+b,a+2b,a,b).
\]
Its last two coordinates prove injectivity of the parameter pair
\((a,b)\), so there are nine words.  If \(a=0\ne b\) the third
coordinate alone vanishes; if \(b=0\ne a\) the fourth alone vanishes;
and if both are nonzero exactly one of \(a+b,a+2b\) vanishes.
Thus every nonzero word has weight three, and each of the four
coordinates is the unique vanishing coordinate for two words.

The map \((u_i,z_i)\mapsto3u_i+2z_i\pmod6\) is bijective:
reduction modulo two recovers \(u_i\), and reduction modulo three
recovers \(2z_i\), from which multiplication by 2 recovers \(z_i\).
It is a homomorphism of additive groups, since changing binary or
ternary representatives by their moduli changes \(3u_i+2z_i\) by a
multiple of six.  Its product with the retained binary \(y\)-coordinate
therefore maps \(C_2\times C_3\) injectively onto the subgroup
\(\mathcal H\).  It follows that \(|\mathcal H|=8\cdot9=72\).  For
\(u_i\in\{0,1\}\) and \(z_i\in\{0,1,2\}\),
\[
 \frac56(3u_i+2z_i)^2
 =\frac{15}{2}u_i+10u_iz_i+\frac{10}{3}z_i^2
 \equiv\frac32u_i+\frac43z_i^2\pmod {2\mathbb Z}.
\]
The cross-term is even; a nonzero \(z_i\) has square 1 or 4, whose
contributions in the last expression differ by 4.  The binary
contribution including D4 is therefore zero, 6, or \(3+1=4\), and
the ternary contribution is zero or \((4/3)3=4\) modulo
\(2\mathbb Z\).  This proves isotropy.

The inverse image is a lattice because it is a subgroup between the
full lattices \(R_{\mathrm{ES}}\) and \(R_{\mathrm{ES}}^\vee\).
Isotropy gives \((x,x)\in2\mathbb Z\) for every \(x\in N\), and
\[
 (x,y)=\frac{(x+y,x+y)-(x,x)-(y,y)}2\in\mathbb Z
                                    \qquad(x,y\in N).
\]
Thus it is even integral.  Its index is \(|\mathcal H|=72\).
For a finite-index inclusion of equal-rank lattices, the covolume
changes by that index, and the Gram determinant is the square of the
covolume in orthonormal coordinates.  Hence
\[
 \det N=\frac{6^4\cdot4}{72^2}=1.
\]
The dual-index formula proved above implies \(N=N^\vee\).
Its rank and positive definiteness are inherited from the unchanged real
span and metric.

For \(0\le k\le5\), every vector in the A5 class \(k\lambda+A\)
has the form
\[
 m-\frac{k}{6}(1,1,1,1,1,1),\qquad
 m\in\mathbb Z^6,\quad\sum m_i=k.
\]
Its squared norm is \(\sum m_i^2-k^2/6\).  Since
\(m_i^2\ge m_i\) for every integer \(m_i\), the minimum is at least
\(k-k^2/6\).  It is attained by choosing exactly \(k\) of the six
entries equal to one and the others zero.  Thus
\begin{equation}
 \mu_A(k)=\frac{k(6-k)}6.
 \label{esn:eq:a5-coset-minimum}
\end{equation}
Every nonzero integral D4 class has minimum one, attained by \(v\);
every half-integral vector has norm at least
\(4(1/2)^2=1\), attained by \(s,c\) in their respective classes.
The four A5 minima and the D4 minimum add and are simultaneously attained
because the metric is an orthogonal sum.

For a coordinate with \(k_i=3u_i+2z_i\pmod6\), its minimum
contribution is 0 for \((u_i,z_i)=(0,0)\), \(3/2\) for
\((1,0)\), \(4/3\) for \((0,z_i\ne0)\), and \(5/6\) for
\((1,z_i\ne0)\).  The complete nonzero case table is
\begin{equation}
\begin{array}{c|c|c|r|l}
 \operatorname{wt}(u)&\operatorname{wt}(z)&
 |\operatorname{supp}(u)\cap\operatorname{supp}(z)|&
 \text{classes}&\text{minimum squared norm}\\\hline
 2&0&0&6&2(3/2)+1=4\\
 4&0&0&1&4(3/2)=6\\
 0&3&0&8&3(4/3)=4\\
 2&3&1&24&(3/2)+2(4/3)+(5/6)+1=6\\
 2&3&2&24&(4/3)+2(5/6)+1=4\\
 4&3&3&8&(3/2)+3(5/6)=4.
\end{array}
 \label{esn:eq:glue-minimum-table}
\end{equation}
For each of the six weight-two binary words, four nonzero ternary
words omit a coordinate inside its support, and four omit a coordinate
outside.  This proves the two counts \(6\cdot4=24\); all other counts
follow from the enumerated binary and ternary words.  Thus the total
number at minimum 4 is \(6+8+24+8=46\), and the total at minimum 6 is
\(1+24=25\).  Their sum is 71, leaving exactly the zero glue class.

No nonzero glue class can contain a norm-two vector.  In the zero
class, an A5 norm-two vector has exactly one coordinate 1 and another
\(-1\), so is \(e_i-e_j\), with \(6\cdot5=30\) possibilities.
A D4 norm-two vector has exactly two coordinates of absolute value one,
so is \(\pm e_i\pm e_j\), with \(4\binom42=24\) possibilities.
An orthogonal sum of two nonzero root-lattice vectors has norm at least
4 because each summand is positive even.  These are therefore all
\(4\cdot30+24=144\) roots of \(N\).
\end{proof}

A Niemeier lattice means a positive-definite even unimodular lattice of
rank 24.  The theorem constructs one with root system \(A_5^4D_4\)
directly.  Its root type appears with lambency 6 in the primary umbral
table \cite[Table~3, p.~21]{esn:CDH2014}.  This table reference is not
needed to prove the inverse-image construction or any minimum above.

For comparison with the earlier labeled glue presentation in the source
\cite{esn:ESMain2026,esn:Provenance}, put
\[
 h_1=(1,3,1,1;00),\qquad h_2=(1,4,2,3;11),\qquad
 h_3=(3,0,3,0;01),
\]
and let \(\mathcal H_{\mathrm{old}}\) be their subgroup.  The exact
transport is
\begin{equation}
 P(x_1,x_2,x_3,x_4;y)=(x_3,x_4,x_1,x_2;y),\qquad
 P(\mathcal H_{\mathrm{old}})=\mathcal H.
 \label{esn:eq:old-code-transporter}
\end{equation}
Indeed, if \(rh_1+th_2+wh_3=0\) with \(r,t\in\mathbb Z/6\)
and \(w\in\mathbb Z/2\), the final two coordinates imply that \(t\)
is even and \(w=0\).  The first and third coordinates then give
\(r+t=0\) and \(r+2t=0\) modulo six, whence \(t=r=0\).
Thus \(|\mathcal H_{\mathrm{old}}|=6\cdot6\cdot2=72\).  Directly,
\[
\begin{split}
 P(h_1)&=3(1111;00)+2(2,2,2,0;00),\\
 P(h_2)&=3(0110;11)+2(1,0,2,2;00),\\
 P(h_3)&=3(1010;01).
\end{split}
\]
The displayed ternary parts are \(2z_1\) and \(2z_1+2z_2\).
Consequently \(P(\mathcal H_{\mathrm{old}})\subset\mathcal H\),
and the equal orders prove equality.  This transporter is realized by
the corresponding isometry permuting the full identified A5 summands.

\subsection{The common D4 root lattice, shortest weights, and triality}
\label{esn:subsec:d4-flag}

Use orthonormal coordinates \(\varepsilon_1,\ldots,\varepsilon_4\)
for the ordinary F4 root realization.  Its long roots and the three
parts of its short roots are
\begin{equation}
\begin{split}
 \Phi_D&=\{\pm\varepsilon_i\pm\varepsilon_j:1\le i<j\le4\},\\
 \mathcal V&=\{\pm\varepsilon_i:1\le i\le4\},\\
 \mathcal S_\pm&=
 \left\{\frac12\sum_{i=1}^4\delta_i\varepsilon_i:
          \delta_i\in\{1,-1\},\ \prod_i\delta_i=\pm1\right\},\\
 \Phi_F&=\Phi_D\cup\mathcal V\cup\mathcal S_+\cup\mathcal S_-.
\end{split}
 \label{esn:eq:root-weight-sets}
\end{equation}
For a nonzero vector \(a\), reflection means the exact map
\(s_a(x)=x-2(x,a)a/(a,a)\).  Write \(W_D\) and \(W_F\) for the
groups generated by the reflections in \(\Phi_D\) and \(\Phi_F\),
respectively.

\begin{theorem}[Integral D4 and the marked triality quotient]
\label{esn:thm:d4-f4-triality}
The map
\begin{equation}
 \iota_D:D\longrightarrow\mathbb Z\langle\Phi_D\rangle,
 \qquad(x_1,x_2,x_3,x_4)\longmapsto
                         \sum_{i=1}^4x_i\varepsilon_i
 \label{esn:eq:d4-isometry}
\end{equation}
is an integral isometry onto the long-root lattice.  It extends to duals
and discriminant forms, sending the shortest-vector sets of the marked
classes \(v,s,c\) respectively to
\(\mathcal V,\mathcal S_+,\mathcal S_-\).

The two matrices
\begin{equation}
 R_D=\operatorname{diag}(-1,1,1,1),\qquad
 H_D=\frac12\begin{pmatrix}
 1&1&1&1\\1&1&-1&-1\\1&-1&1&-1\\1&-1&-1&1
 \end{pmatrix}
 \label{esn:eq:d4-triality-matrices}
\end{equation}
are the short-root reflections \(s_{\varepsilon_1}\) and \(s_c\),
where \(c=(-1,1,1,1)/2\) is read in the \(\varepsilon\) coordinates.
They induce \((s\ c)\) and \((v\ s)\) on the three nonzero
discriminant classes.  There is a split exact sequence
\begin{equation}
 1\longrightarrow W_D\longrightarrow W_F
   \longrightarrow S_3\longrightarrow1,
 \qquad W_D\cong\{\text{even signed permutations of four coordinates}\},
 \label{esn:eq:weyl-triality-sequence}
\end{equation}
whose displayed section is \(\langle R_D,H_D\rangle\cong S_3\).
In particular \(|W_D|=192\) and \(|W_F|=1152\).
\end{theorem}

\begin{proof}
The four D4 basis vectors in Theorem~\ref{esn:thm:root-discriminant}
are among \(\Phi_D\), and every vector of \(\Phi_D\) has even
coordinate sum.  Thus the displayed identity of coordinates has exactly
the stated image, preserves the metric, and extends to duals by the
integer-pairing definition.  In a nonzero integral dual class, norm one
forces one coordinate \(\pm1\) and all others zero.  In a
half-integral class, norm one forces all four coordinates to be
\(\pm1/2\).  The difference from \(s\) has even coordinate sum
exactly when the number of negative signs is even.  This proves the
three shortest-vector identifications with the stated marking.

The reflection identity for \(R_D\) is immediate, and
\(H_D=I-2cc^T\) follows by multiplying the four coordinates of \(c\).
Both are orthogonal involutions.  The first preserves the parity of an
integer coordinate sum.  For \(x\in D\), the four coordinates of
\(H_Dx\) are integers because each numerator is congruent to
\(\sum x_i\) modulo two, and their sum is \(2x_1\).  Thus both
preserve \(D\).  Their effects on the representatives show that
\(R_D\) fixes the class \(v\) and exchanges \(s,c\), whereas
\(H_Dv=s\), \(H_Ds=v\), and \(H_Dc=-c\equiv c\pmod D\).
They therefore preserve all of \(\Phi_F\) by the shortest-vector
description.  On the plane with ordered basis \(\varepsilon_1,c\),
their product has matrix
\[
 \begin{pmatrix}0&-1\\1&-1\end{pmatrix},
\]
whose cube is the identity, while both reflections fix the orthogonal
complement.  Consequently \(R_D^2=H_D^2=(R_DH_D)^3=1\).
Every word in two such involutions reduces to
\((R_DH_D)^k\) or \((R_DH_D)^kR_D\), \(0\le k\le2\).
Their discriminant permutations are the six elements of \(S_3\),
so these six words are distinct and give the stated section.

Reflection in \(\varepsilon_i-\varepsilon_j\) swaps coordinates
\(i,j\).  Reflection in \(\varepsilon_i+\varepsilon_j\) swaps them
and changes both signs.  Hence the long-root reflections generate every
permutation and every even sign change, and each long-root reflection
is of this form.  This proves the assertion about \(W_D\) and its
order \(2^3\cdot4!=192\).  Such a reflection acts trivially on
\(D^\vee/D\): for a root \(a\) of squared norm two,
\(s_a(x)-x=-(x,a)a\in D\) for all \(x\in D^\vee\).

Let \(K=\langle W_D,R_D,H_D\rangle\).  Each of its generators
preserves the long roots, so \(W_D\) is normal in \(K\) by
\(g s_a g^{-1}=s_{ga}\).  The even signed permutations act
transitively on each of \(\mathcal V,\mathcal S_+,\mathcal S_-\),
and the displayed section permutes these three sets transitively.
Thus every short root is \(g\varepsilon_1\) for some \(g\in K\),
and its reflection \(gR_Dg^{-1}\) belongs to \(K\).  All generators
of \(W_F\) therefore belong to \(K\); the reverse inclusion follows
from their definitions, proving \(W_F=K\).
An element in the kernel of its permutation action on the three short
sets preserves \(\mathcal V\), so it is a signed permutation matrix.
Preservation of \(\mathcal S_+\) forces an even number of sign
changes, so the kernel is \(W_D\).  This proves exactness, splitting,
and the order \(6\cdot192=1152\).
\end{proof}

\paragraph{The exact geometric datum and weight marking.}
Let \(\mathbb O\) denote the real normed octonion algebra and let
\(J=\mathfrak h_3(\mathbb O)\) be the real vector space of Hermitian
three-by-three matrices, with Jordan product
\(a\circ b=(ab+ba)/2\).  Define
\[
 \mathbb O\mathbb P^2=\{a\in J:a^2=a,\ \operatorname{tr}a=1\},
 \qquad
 \operatorname{Fl}(\mathbb O)=
 \{(a,b)\in(\mathbb O\mathbb P^2)^2:
                                 \operatorname{Re}\operatorname{tr}(ab)=0\}.
\]
For the compact group \(F_4=\operatorname{Aut}_{\mathbb R}(J,\circ)\),
Mare and Willems prove the homogeneous identification
\(\operatorname{Fl}(\mathbb O)=F_4/\operatorname{Spin}(8)\) and
identify its tangent representation as
\(V_8\oplus S_8^+\oplus S_8^-\)
\cite[Definition~2.2, Proposition~2.4, Proposition~2.7, and
pp.~491--492]{esn:MareWillems2013}.  These are the prior geometric
results used here.  With their torus coordinate forms \(L_i\), their
equations (5.4)--(5.6) assign the vector, positive half-spin, and negative
half-spin weights respectively the patterns
\(\pm L_i\), even half-sign sums, and odd half-sign sums
\cite[pp.~510--511]{esn:MareWillems2013}.
The coordinate map \(L_i\mapsto\varepsilon_i\) retains that sign
marking and identifies these three weight sets with those of
\eqref{esn:eq:root-weight-sets}.

For completeness the half-spin weight patterns admit the following
direct calculation.  In the complex exterior algebra
\(\Lambda^*\mathbb C^4\), let \(a_i\) be exterior multiplication
by the \(i\)-th basis vector and \(b_i\) contraction by its dual.
On a basis wedge, insertion and deletion give
\[
 a_ia_j+a_ja_i=b_ib_j+b_jb_i=0,\qquad
 a_ib_j+b_ja_i=\delta_{ij}I.
\]
Thus \(C_i=a_i+b_i\) and \(D_i=\mathrm i(b_i-a_i)\) square to
one and mutually anticommute for distinct Clifford generators.
Writing \(N_i=a_ib_i\), direct multiplication gives
\[
 \tfrac12 C_iD_i=\mathrm i(N_i-\tfrac12 I).
\]
The four commuting infinitesimal torus operators therefore have weights
\(\sum_i(\mathbf1_{i\in I}-1/2)\varepsilon_i\) on the basis
wedge indexed by \(I\subset\{1,2,3,4\}\).  The even exterior
part has an even number of negative signs, and the odd part has an odd
number.  There are eight one-dimensional weight spaces in each part.
For the vector representation the four real coordinate planes
complexify into their two rotation eigenspaces, with weights
\(\varepsilon_i,-\varepsilon_i\).  This proves the displayed
weight formulas directly in the stated marking.

The map proved here is consequently the actual chain
\[
 D\xrightarrow{\ \iota_D\ }\mathbb Z\langle\Phi_D\rangle,
 \qquad
 \{v,s,c\}\longleftrightarrow
       \{\mathcal V,\mathcal S_+,\mathcal S_-\}
       =\{\operatorname{wt}(V_8),
          \operatorname{wt}(S_8^+),\operatorname{wt}(S_8^-)\}.
\]
The same \(S_3\) acts through the explicitly proved permutations of
these three sets.  The companion higher-rung treatment records the
geometric datum and root-coordinate comparison at the respective labels
\begin{center}
\texttt{h24:thm:octonionic-triality}\\
\texttt{h24:prop:two-Coxeter-systems}.
\end{center}
The exact source pointers and passages are included in
\cite{esn:Provenance}.  The integral metric map
in this section is the map on D4 and its dual.  The definitions make no
map from the whole rank-24 lattice \(N\) into the flag tangent space.

\subsection{The full indicated glue stabilizer and its triality quotient}
\label{esn:subsec:glue-gl2-triality}

All vector spaces over finite fields in this subsection retain their displayed
bases.  Write
\[
 j:\mathbb F_3^2\longrightarrow\mathbb F_3^4,
 \qquad j(a,b)=(a+b,a+2b,a,b),
 \qquad C_3=j(\mathbb F_3^2).
\]
The four row forms defining this injection are
\[
 \ell_1=(1,1),\quad \ell_2=(1,2),\quad
 \ell_3=(1,0),\quad \ell_4=(0,1).
\]
They represent all four one-dimensional subspaces of the dual of
\(\mathbb F_3^2\).  Indeed each nonzero row has either first coordinate zero,
in which case its line is represented by \(\ell_4\), or a unique scalar
multiple with first coordinate one, namely one of
\((1,0),(1,1),(1,2)\).

Let \(E\subset\mathbb F_2^4\) be the even-weight subspace and define
\[
 \tau:E\longrightarrow\mathbb F_2^2,
 \qquad \tau(u_1,u_2,u_3,u_4)
       =(u_2+u_4,u_3+u_4).
\]
The binary code in the index-72 glue is exactly
\[
 C_2=\{(u,\tau(u)):u\in E\}.
\]
To verify this, the independent vectors
\(1111,1100,1010\) form a basis of the three-dimensional space \(E\),
and their images under \(\tau\) are respectively
\(00,10,01\).  These are precisely the three stated binary code generators.
In particular \(\tau\) is onto and
\[
 \ker\tau=\langle1111\rangle.
\]
Explicitly, \(\tau(u)=0\) gives \(u_2=u_3=u_4\), and the even-weight
equation then gives \(u_1=u_4\).

For a permutation \(p\) of \(\{1,2,3,4\}\), use the coordinate convention
\[
 (P_pu)_i=u_{p(i)}.
\]
The permutation preserves \(E\) and \(\langle1111\rangle\), so it induces
a unique invertible map \(B_p\) satisfying
\begin{equation}
 B_p\tau(u)=\tau(P_pu)\qquad(u\in E).
 \label{esn:eq:binary-permutation-map}
\end{equation}
For an entirely explicit formula, the two columns of \(B_p\) are
\(\tau(P_p(1100))\) and \(\tau(P_p(1010))\).
The convention implies \(P_pP_q=P_{q\circ p}\).  Accordingly
\(B_pB_q=B_{q\circ p}\); no reversal is suppressed in what follows.

The subgroup of \eqref{esn:eq:glue-and-overlattice} can therefore be written
\[
 \mathcal H=
 \{(3u+2z;\tau(u)):u\in E,\ z\in C_3\}
 \subset(\mathbb Z/6)^4\oplus\mathbb F_2^2.
\]
Let \(G\) be its stabilizer among the transformations
\[
 g=(\varepsilon,p,B),\qquad
 (x;y)\longmapsto
 \bigl((\varepsilon_i x_{p(i)})_{i=1}^4;By\bigr),
\]
where \(\varepsilon_i\in\{1,-1\}\), \(p\in S_4\), and
\(B\in GL_2(\mathbb F_2)\).  Denote the signed permutation matrix on the
first four coordinates by \(M_g\).

\begin{theorem}[An explicit isomorphism to \(GL_2(\mathbb F_3)\)]
\label{esn:thm:glue-gl2}
Restriction to the ternary code gives an isomorphism
\[
 \Theta:G\xrightarrow{\ \cong\ }GL_2(\mathbb F_3),
 \qquad M_gj=j\Theta(g).
\]
For every \(A\in GL_2(\mathbb F_3)\), its inverse image is specified
uniquely by
\begin{equation}
 \ell_iA=\varepsilon_i\ell_{p(i)}\quad(1\le i\le4),
 \qquad B=B_p.
 \label{esn:eq:inverse-gl2-glue-map}
\end{equation}
In particular \(|G|=48\).
\end{theorem}

\begin{proof}
The elements of \(\mathcal H\) of order dividing two are precisely
\((3u;\tau(u))\), and those of order dividing three are precisely
\((2z;0)\).  Every indicated ambient transformation preserves these primary
parts, so an element of \(G\) satisfies \(M_gC_3=C_3\).  The injection
\(j\) therefore yields one and only one \(\Theta(g)\in GL_2(\mathbb F_3)\).
The binary part is preserved exactly when
\(B\tau(u)=\tau(P_pu)\), because all signs become one modulo two.
Surjectivity of \(\tau\) forces \(B=B_p\).

If \(\Theta(g)=I\), the equality \(M_gj=j\) implies
\(\varepsilon_i\ell_{p(i)}=\ell_i\) for every \(i\).
The four row lines are distinct, so \(p(i)=i\).  Since each \(\ell_i\)
is nonzero and \(1\ne-1\) in \(\mathbb F_3\), also
\(\varepsilon_i=1\) for every \(i\).  It follows that \(B_p=I\) and
\(g=1\).  Thus \(\Theta\) is injective.

Conversely fix \(A\in GL_2(\mathbb F_3)\).  The rows \(\ell_iA\) are
nonzero and belong to four distinct row lines.  Since the \(\ell_i\)
represent every row line, each \(\ell_iA\) has a unique expression
\(\varepsilon_i\ell_{p(i)}\), where the nonzero scalar in \(\mathbb F_3\)
is represented by the exact integer sign \(\varepsilon_i=1\) or \(-1\).
Distinctness of the row lines makes \(p\) a permutation.
The resulting signed permutation satisfies \(M_gj=jA\), hence preserves
\(C_3\).  Choosing \(B=B_p\) preserves \(C_2\) as proved above, hence
preserves \(\mathcal H\).  This constructs an element of \(G\) mapping
to \(A\), proving surjectivity and the stated inverse formula.

Finally
\[
 M_{gh}j=M_gM_hj=M_gj\Theta(h)
 =j\Theta(g)\Theta(h),
\]
so \(\Theta(gh)=\Theta(g)\Theta(h)\).
An invertible \(2\)-by-\(2\) matrix over \(\mathbb F_3\) has eight
choices for its first nonzero column and six choices for its second column
outside the span of the first.  Thus its group order is \(8\cdot6=48\).
\end{proof}

The map to the permutation of the four component labels is more precisely
\(g\mapsto p^{-1}\), because a vector supported on component \(k\) is
sent to component \(p^{-1}(k)\).  Its kernel under \(\Theta\) is
\(\{I,-I\}\).  In fact a matrix preserving all four row lines is diagonal
by preservation of \(\ell_3\) and \(\ell_4\), and its diagonal entries
are equal by preservation of \(\ell_1\).  Conversely scalar matrices
preserve all four lines.  Therefore the component permutation image has
order \(48/2=24\), so is all \(S_4\), with
\[
 G/\Theta^{-1}(\{\pm I\})
       \xrightarrow{\ \overline\Theta\ }PGL_2(\mathbb F_3)\cong S_4.
\]
The center of \(GL_2(\mathbb F_3)\) is exactly \(\{\pm I\}\): a
central matrix commutes with \(\operatorname{diag}(1,-1)\), forcing its
off-diagonal entries to vanish, and commuting with
\(\left(\begin{smallmatrix}0&1\\1&0\end{smallmatrix}\right)\) forces
the diagonal entries to agree.

\begin{theorem}[The exact triality quotient and its quaternion kernel]
\label{esn:thm:triality-q8}
The projection \(\pi:G\to GL_2(\mathbb F_2)\),
\((\varepsilon,p,B)\mapsto B\), is onto, and under \(\Theta\) its
kernel is the quaternion subgroup generated by
\[
 I_q=\begin{pmatrix}0&1\\2&0\end{pmatrix},
 \qquad J_q=\begin{pmatrix}1&1\\1&2\end{pmatrix}.
\]
There is a split exact sequence
\[
 1\longrightarrow Q_8\longrightarrow GL_2(\mathbb F_3)
 \xrightarrow{\ \pi\Theta^{-1}\ }GL_2(\mathbb F_2)
 \longrightarrow1.
\]
The splitting is the subgroup generated by
\[
 A_R=\begin{pmatrix}1&0\\0&2\end{pmatrix},
 \qquad A_H=\begin{pmatrix}1&2\\0&2\end{pmatrix}.
\]
It identifies \(GL_2(\mathbb F_3)\) with \(Q_8\rtimes S_3\) for the
conjugation action of this displayed subgroup on this displayed \(Q_8\).
\end{theorem}

\begin{proof}
The nonzero elements of \(E/\langle1111\rangle\) are represented by
the complementary pairs of two-element subsets
\[
 \{12\mid34\},\qquad\{13\mid24\},\qquad\{14\mid23\},
\]
marked respectively by \(10,01,11\) via \(\tau\).
Thus \(B_p\) is exactly the permutation action on these three
partitions, expressed in the displayed binary marking.

Using \eqref{esn:eq:inverse-gl2-glue-map} gives
\[
 M(A_R)x=(x_2,x_1,x_3,-x_4),\quad
 B_R=\begin{pmatrix}1&1\\0&1\end{pmatrix},
\]
\[
 M(A_H)x=(x_1,x_3,x_2,-x_4),\quad
 B_H=\begin{pmatrix}0&1\\1&0\end{pmatrix}.
\]
These binary matrices act respectively as the transpositions
\((01\ 11)\) and \((10\ 01)\), so they generate all permutations of
the three nonzero binary vectors.  There are
\((4-1)(4-2)=6\) invertible binary matrices; a matrix fixing all nonzero
vectors fixes the basis.  Hence this permutation action identifies
\(GL_2(\mathbb F_2)\) with \(S_3\) and proves surjectivity of \(\pi\).
Its kernel therefore has order eight.

For \(I_q\) the row rule gives
\[
 p=(2,1,4,3),\qquad\varepsilon=(-1,1,1,-1),
\]
and for \(J_q\) it gives
\[
 p=(3,4,1,2),\qquad\varepsilon=(-1,-1,1,1).
\]
Both permutations fix each of the three displayed complementary
partitions, so both matrices lie in the kernel.  Direct matrix
multiplication over \(\mathbb F_3\) gives
\[
 I_q^2=J_q^2=-I,\qquad I_qJ_q=-J_qI_q,
 \qquad I_qJ_q=\begin{pmatrix}1&2\\2&2\end{pmatrix}.
\]
The eight matrices
\(I,-I,\pm I_q,\pm J_q,\pm I_qJ_q\) are pairwise distinct by their
displayed entries.  The relations show this set is a group: powers reduce
using \(I_q^2=J_q^2=-I\), and factors can be reordered using
\(J_qI_q=-I_qJ_q\).  These relations identify it with the quaternion
group \(Q_8\).  Since the kernel has order eight, this is the whole kernel.

The matrices \(A_R,A_H\) satisfy
\[
 A_R^2=A_H^2=I,\qquad
 A_RA_H=\begin{pmatrix}1&2\\0&1\end{pmatrix},\qquad
 (A_RA_H)^3=I.
\]
Every word in the two involutions reduces to one of
\((A_RA_H)^k\) or \((A_RA_H)^kA_R\), \(k=0,1,2\), by alternation
and the cubic relation.  Their images under \(\pi\) are the six distinct
permutations generated by \(B_R,B_H\).  Hence the subgroup has exactly
six elements and maps isomorphically onto \(GL_2(\mathbb F_2)\).
This gives the splitting.  For any \(g\), the unique splitting element
with image \(\pi(g)\) leaves \(g\) times its inverse in the kernel;
the intersection of kernel and splitting is trivial.  Multiplication
therefore realizes the asserted semidirect product with the indicated
conjugation action.
\end{proof}

The matrices in the kernel have an additional marked interpretation:
their projective component permutations are the four permutations
\(1,(12)(34),(13)(24),(14)(23)\).  Each fixes all three partitions, and
these are exactly the kernel of the onto action \(S_4\to S_3\), whose
kernel has order \(24/6=4\).  Thus the preceding exact sequence is the
precise map from the order-48 ternary-code symmetry to the three
\(D_4\) discriminant classes.

\begin{theorem}[A faithful action on the actual index-72 lattice]
\label{esn:thm:gl2-full-lattice-lift}
Fix the four identifications of the \(A_5\) coordinate lattices used in
the glue presentation, and the discriminant marking
\(10\mapsto v,01\mapsto s,11\mapsto c\) of \(D_4\).
The group \(G\cong GL_2(\mathbb F_3)\) has a faithful action by integral
isometries of the index-72 overlattice.  Write
\(\varphi_i:A\longrightarrow A(F_{\eta_i})\) for the restriction of
the fixed free-coordinate isometry \(e_j\mapsto e_{f_i289^j}\).
For \(X_i\in A(F_{\eta_i})\), the \(i\)-th output component is exactly
\(\varepsilon_i\varphi_i\varphi_{p(i)}^{-1}X_{p(i)}\).
In the identified \(A\) coordinates this is the action
\[
 L_g(X_1,X_2,X_3,X_4;d)
 =\bigl(\varepsilon_1X_{p(1)},\ldots,
        \varepsilon_4X_{p(4)};\delta(B)d\bigr),
\]
where \(\delta:GL_2(\mathbb F_2)\to O(D_4)\) is the section given by
\[
 \delta(B_R)=R_D=\operatorname{diag}(-1,1,1,1),\qquad
 \delta(B_H)=H_D=\frac12
 \begin{pmatrix}
 1&1&1&1\\1&1&-1&-1\\1&-1&1&-1\\1&-1&-1&1
 \end{pmatrix}.
\]
Under \(\Theta\), the two displayed full triality lifts are exactly
\(A_R\) and \(A_H\).
\end{theorem}

\begin{proof}
The four identified \(A_5\) summands carry the same Euclidean metric,
so the displayed signed permutation is a root-lattice isometry.  The
matrices \(R_D,H_D\) preserve \(D_4\): \(R_D\) negates one integral
coordinate and preserves parity of the sum; for \(x\in D_4\), each
coordinate of \(H_Dx\) is integral because its numerator has the same
parity as \(\sum x_i\), and the coordinate sum of \(H_Dx\) is \(2x_1\).
Both matrices are orthogonal involutions by multiplication.
Furthermore \(R_D=s_{e_1}\), \(H_D=s_c\), with
\(c=(-1,1,1,1)/2\), and \((e_1,c)=-1/2\).
On the plane with basis \(e_1,c\), their product has matrix
\(\left(\begin{smallmatrix}0&-1\\1&-1\end{smallmatrix}\right)\)
and has cube the identity; both reflections fix the orthogonal
complement pointwise.  The discriminant actions are the two distinct
transpositions \(B_R,B_H\).  The same word reduction as in the preceding
proof therefore shows that \(\langle R_D,H_D\rangle\) has exactly six
elements and that its discriminant action is an isomorphism onto
\(GL_2(\mathbb F_2)\).  Its inverse is the homomorphism \(\delta\).

Thus \(g\mapsto L_g\) is a homomorphism of root-lattice isometries.
It extends to the dual lattice because an isometry preserves the
integer-pairing condition defining the dual.  Its discriminant action is
the specified \((\varepsilon,p,B)\), which preserves \(\mathcal H\).
Consequently it preserves the inverse image
\[
 N=\{x\in R_{\mathrm{ES}}^\vee:
                  x+R_{\mathrm{ES}}\in\mathcal H\}.
\]
It is therefore an integral isometry of \(N\).
If \(L_g\) is the identity on \(N\), it is the identity on the full
real span, because \(N\) contains the full-rank lattice
\(R_{\mathrm{ES}}\).  It is therefore also the identity on
\(R_{\mathrm{ES}}^\vee/R_{\mathrm{ES}}\), in particular on the ternary
code.  Hence \(\Theta(g)=I\), and injectivity of \(\Theta\) gives \(g=1\).
This proves faithfulness.
The exact formulas for \(A_R,A_H\) were established above, so their
lifts are the stated full triality transformations.
\end{proof}

This constructs an explicit finite subgroup of \(O(N)\), including its
action on all five root summands and on their glue, and its exact
\(D_4\)-triality quotient.  The group and its projective quotient agree
with the \(A_5^4D_4\) row of \cite[Table~3, p.~21]{esn:CDH2014};
the isomorphisms here have been proved by their actual code actions.

For clarity, the two full triality elements, denoted
\(\widehat R=L_{\Theta^{-1}(A_R)}\) and
\(\widehat H=L_{\Theta^{-1}(A_H)}\), act in the fixed coordinates by
\[
\begin{split}
 \widehat R(X_1,X_2,X_3,X_4;d)
     &=(X_2,X_1,X_3,-X_4;R_Dd),\\
 \widehat H(X_1,X_2,X_3,X_4;d)
     &=(X_1,X_3,X_2,-X_4;H_Dd).
\end{split}
\]
Their squares and the cube of their product are the identity on the
entire lattice, by the proved faithful group homomorphism.
The product on the first three A5 summands is the indicated
three-cycle; its two negations on the fourth summand cancel.

\subsection{The exact reduced-shell divisor bijection}
\label{esn:subsec:shells}

Let \(p\) be a positive prime with \(p\equiv1\pmod {24}\), and
let the positive integer shell \(R\) satisfy \(R\equiv3\pmod4\).
Retain
\begin{equation}
 a=a_R(p)=\frac{p+R}{4},\qquad
 N_{\mathrm{sh}}=pa,\qquad g=\gcd(R,N_{\mathrm{sh}}),\qquad
 R_0=\frac Rg,\quad N_0=\frac{N_{\mathrm{sh}}}{g}.
 \label{esn:eq:shell-data}
\end{equation}
The integer \(N_{\mathrm{sh}}\) is the source's scalar \(N=pa\);
the subscript distinguishes it from the lattice \(N\) without changing
either object.  In particular \(a,R_0,N_0\) are positive integers,
\(\gcd(R_0,N_0)=1\), and the prescribed first denominator equation is
\begin{equation}
 \frac4p-\frac1a=\frac{4a-p}{pa}
       =\frac{R}{N_{\mathrm{sh}}}=\frac{R_0}{N_0}
       =\frac1b+\frac1c.
 \label{esn:eq:fixed-shell-equation}
\end{equation}
Here \(b,c\) are positive integer denominators, and their order is
initially retained.  Define the finite set
\[
 \mathcal D_{R,p}=
 \{d\in\mathbb Z_{>0}:d\mid N_0^2,
                             \ d\equiv-N_0\pmod {R_0}\},
 \qquad A_R(p)=|\mathcal D_{R,p}|.
\]

\begin{theorem}[Certificates, ordered pairs, and complementary pairs]
\label{esn:thm:shell-bijection}
There is an exact bijection
\begin{equation}
 \begin{split}
 \mathcal D_{R,p}&\longrightarrow
  \{(b,c)\in\mathbb Z_{>0}^2:
                    4/p=1/a+1/b+1/c\},\\
 d&\longmapsto
 \left(\frac{N_0+d}{R_0},
       \frac{N_0+N_0^2/d}{R_0}\right),
 \end{split}
 \label{esn:eq:shell-bijection}
\end{equation}
with inverse \(d=R_0b-N_0\).  The involution \(d\mapsto N_0^2/d\)
corresponds exactly to \((b,c)\mapsto(c,b)\).  Consequently the
number of unordered pairs of remaining denominators is
\begin{equation}
 U_R(p)=\frac{A_R(p)+\delta_{R,p}}2,
 \qquad
 \delta_{R,p}=\begin{cases}1,&R_0\mid2,\\0,&R_0\nmid2.\end{cases}
 \label{esn:eq:unordered-count}
\end{equation}
The unique possible fixed certificate is \(d=N_0\), giving
\(b=c=2N_0/R_0\).
\end{theorem}

\begin{proof}
For a positive solution of \eqref{esn:eq:fixed-shell-equation},
\(R_0b/N_0=1+b/c>1\), so \(d=R_0b-N_0>0\); likewise
\(e=R_0c-N_0>0\).  Multiplying the fraction equation by
\(N_0bc\) gives \(R_0bc=N_0(b+c)\), and therefore
\[
 (R_0b-N_0)(R_0c-N_0)=N_0^2.
\]
Thus \(d\mid N_0^2\), \(d\equiv-N_0\pmod {R_0}\), and
\(e=N_0^2/d\).

Conversely, let \(d\in\mathcal D_{R,p}\) and \(e=N_0^2/d\).
For \(R_0>1\), the coprimality \(\gcd(R_0,N_0)=1\) makes
\(d\equiv-N_0\) a unit, and hence
\[
 e\equiv N_0^2(-N_0)^{-1}\equiv-N_0\pmod {R_0}.
\]
For \(R_0=1\), this congruence is vacuous and holds automatically.
Thus both expressions in \eqref{esn:eq:shell-bijection} are positive
integers.  Substitution into \((R_0b-N_0)(R_0c-N_0)=N_0^2\)
and expansion give the original fraction equation.  Each construction
recovers exactly \(d=R_0b-N_0\) and \(e=R_0c-N_0\), so they
are inverse, proving the bijection and the involution statement.

For a positive divisor, \(d=N_0^2/d\) holds exactly when
\(d=N_0\).  This divisor is a certificate exactly when
\(N_0\equiv-N_0\pmod {R_0}\), that is, \(R_0\mid2N_0\).
Coprimality makes this equivalent to \(R_0\mid2\).
Every other orbit contains two certificates, so counting the fixed
orbits and the two-element orbits gives
\(U_R(p)=\delta_{R,p}+(A_R(p)-\delta_{R,p})/2\), exactly
\eqref{esn:eq:unordered-count}.
\end{proof}

Factor the exact integer \(N_0\) as
\(N_0=\prod_{j=1}^{m}q_j^{e_j}\), with distinct positive primes
\(q_j\) and positive exponents \(e_j\), and define
\begin{equation}
 \begin{split}
 B(N_0)&=\prod_{j=1}^{m}\{-e_j,-e_j+1,\ldots,e_j\},\\
 \lambda_{R_0,N_0}:B(N_0)&\longrightarrow(\mathbb Z/R_0\mathbb Z)^\times,
 \qquad
 \lambda_{R_0,N_0}(\beta)=\prod_{j=1}^{m}q_j^{\beta_j}\pmod {R_0}.
 \end{split}
 \label{esn:eq:exponent-box}
\end{equation}
The group for modulus one is taken to be the one-element group, so that
the same notation also expresses the vacuous congruence in that case.
For \(R_0>1\), every \(q_j\) is a unit by coprimality, so all
negative powers in this formula are defined.

\begin{theorem}[The exact exponent-box map and the two raw channels]
\label{esn:thm:raw-channels}
The map
\begin{equation}
 \beta\longmapsto d=\prod_{j=1}^{m}q_j^{e_j+\beta_j}
 \label{esn:eq:exponent-to-divisor}
\end{equation}
bijects \(\lambda_{R_0,N_0}^{-1}(-1)\) with
\(\mathcal D_{R,p}\).  It retains each valuation and both original
integers \(R_0,N_0\).

In addition, suppose \(\gcd(R,pa)=1\) and \(p\nmid a\).
Write
\[
 a=\prod_{j=1}^{s}q_j^{e_j},\qquad
 B(a)=\prod_{j=1}^{s}\{-e_j,\ldots,e_j\},
\]
and define, for each unit
\(t\) modulo \(R\),
\[
 c_R(a;t)=\#\left\{\beta\in B(a):
                                 \prod_j q_j^{\beta_j}\equiv t\pmod R\right\}.
\]
Then the number of raw divisor certificates, equivalently ordered
denominator pairs, is
\begin{equation}
 A_R(p)=c_R(a;-1)+2c_R(a;-p^{-1}).
 \label{esn:eq:two-raw-channels}
\end{equation}
Its complementary-divisor quotient is the count in
\eqref{esn:eq:unordered-count}, with \(R_0=R\) under these additional
coprimality hypotheses.
\end{theorem}

\begin{proof}
A divisor of \(N_0^2\) has one and only one exponent vector
\(\alpha_j\in\{0,\ldots,2e_j\}\).  The exact change
\(\beta_j=\alpha_j-e_j\) gives the entire displayed box and
\(dN_0^{-1}\equiv\prod_jq_j^{\beta_j}\pmod {R_0}\).
Multiplication of \(d\equiv-N_0\) by the unit \(N_0^{-1}\)
gives precisely the target \(-1\); modulo one both statements are
vacuous.  This proves the first bijection and its inverse.

Under the additional hypotheses, \(R_0=R\), \(N_0=pa\), and
the exponent of \(p\) in \(pa\) is exactly one.  Every positive integer
divisor \(d\) of \(p^2a^2\) therefore has exactly one valuation-index
representation
\[
 d=p^{t+1}\prod_jq_j^{e_j+\beta_j},\qquad
 (t,\beta)\in\{-1,0,1\}\times B(a).
\]
After division by \(pa\), its unit-group residue is
\[
 d(pa)^{-1}\equiv p^t\prod_jq_j^{\beta_j}\pmod R.
\]
Distinct valuation-index pairs may have the same unit-group residue;
they remain distinct integer divisors in the count.
The three conditions for this to equal \(-1\) are
\[
\begin{array}{c|c}
 t&\prod_jq_j^{\beta_j}\pmod R\\\hline
 -1&-p\\0&-1\\1&-p^{-1}.
\end{array}
\]
The negation \(\beta\mapsto-\beta\) is a bijection of \(B(a)\)
and inverts the unit-group value, so it bijects the fibres at
\(-p\) and \(-p^{-1}\).  Summing the sizes of these three disjoint
raw exponent classes proves \eqref{esn:eq:two-raw-channels}.  The
complementary divisor sends \((t,\beta)\) to \((-t,-\beta)\),
so identifying it changes the raw count as proved in
Theorem~\ref{esn:thm:shell-bijection}; it does not leave the factor
two in \eqref{esn:eq:two-raw-channels} unchanged.
\end{proof}

This counting convention corrects the wording attached to the source's
two-channel formula: the displayed sum is a raw count, whereas a count
after identifying complementary divisors is \eqref{esn:eq:unordered-count}.
The exact source locator and formula are retained in
\cite{esn:Residual2026,esn:Provenance}.

\subsection{A fixed-shell count does not descend to the residue fibre}
\label{esn:subsec:non-descent}

\begin{theorem}[Two primes in the same residue with different shell counts]
\label{esn:thm:non-descent}
The two primes \(2521\) and \(3361\) both reduce to \(1\in S_{840}\),
but for the same shell \(R=3\) their raw and unordered counts are
\begin{equation}
 A_3(2521)=0,\qquad A_3(3361)=6,
 \qquad U_3(2521)=0,\quad U_3(3361)=3.
 \label{esn:eq:non-descent-counts}
\end{equation}
One actual denominator certificate in the second case is
\begin{equation}
 \frac4{3361}=\frac1{841}+\frac1{974690}+\frac1{28266010}.
 \label{esn:eq:explicit-denominator-example}
\end{equation}
Thus the raw count, its unordered quotient, and the indicator of
completion of this fixed shell each fail to factor through prime
reduction modulo 840.
\end{theorem}

\begin{proof}
The residue equations are the literal equalities
\(2521=3\cdot840+1\) and \(3361=4\cdot840+1\).
For an explicit finite primality verification, the primes at most 57,
in increasing order, are
\[
 \mathcal P=(2,3,5,7,11,13,17,19,23,29,31,37,41,43,47,53).
\]
The remainders against its first 15 entries for \(2521\), all 16
entries for \(3361\), and its first nine entries for \(631\), are
respectively
\[
\begin{split}
 &(1,1,1,1,2,12,5,13,14,27,10,5,20,27,30),\\
 &(1,1,1,1,6,7,12,17,3,26,13,31,40,7,24,22),\\
 &(1,1,1,1,4,7,2,4,10).
\end{split}
\]
All are nonzero, and the respective integer square roots are 50, 57,
and 25.  The prime 29 has remainders \(1,2,4\) against \(2,3,5\),
the primes at most its integer square root 5.
Every composite positive integer greater than one has a prime divisor
at most its square root: in a factorization into two factors greater
than one, the smaller factor is at most the square root and has a prime
divisor no larger than itself.  These nonzero remainder lists therefore
prove each stated primality.

For \(R=3\),
\[
 a_3(2521)=631,\qquad a_3(3361)=841=29^2.
\]
Both primes are 1 modulo three; also \(631\equiv1\pmod3\) and
\(841\equiv1\pmod3\).  Thus \(\gcd(3,pa)=1\) in both cases.
Both \(a\) values are smaller than their respective primes, so
\(p\nmid a\).  Every additional hypothesis of
Theorem~\ref{esn:thm:raw-channels} is now verified.
The two targets \(-1\) and \(-p^{-1}\) coincide modulo three.

For \(a=631\), all exponents \(-1,0,1\) give residue one, so
the target \(-1\) has count zero.  For \(a=29^2\), among the five
exponents \(-2,-1,0,1,2\), exactly \(-1,1\) give \(-1\) modulo
three because \(29\equiv-1\pmod3\).  Both target counts are
therefore 2 and the raw count is \(2+2\cdot2=6\).  Here
\(R_0=3\nmid2\), so both unordered counts are half the corresponding
raw counts.  This proves \eqref{esn:eq:non-descent-counts}.

For the second prime, retain
\[
 N_{\mathrm{sh}}=3361\cdot841=2826601,
 \qquad d=3361\cdot29=97469,
 \qquad \frac{N_{\mathrm{sh}}^2}{d}=3361\cdot29^3=81971429.
\]
Then
\[
 \frac{N_{\mathrm{sh}}+d}{3}=974690,
 \qquad
 \frac{N_{\mathrm{sh}}+N_{\mathrm{sh}}^2/d}{3}=28266010.
\]
The bijection of Theorem~\ref{esn:thm:shell-bijection} now proves
\eqref{esn:eq:explicit-denominator-example}.  The first prime has
no completion with this prescribed first denominator by its zero
certificate count.  Two arguments with the same residue and unequal
function values cannot be the pullback of a single residue function:
if \(f(p)=h(p\bmod840)\), equality of the residues forces equality
of the values.  Applying this observation to the three displayed
functions proves the claimed failures of factorization.
\end{proof}

\subsection{An actual residue cocycle and a retained-sample Fourier map}
\label{esn:subsec:cocycle-fourier}

For \(r\in S_{840}\), multiplication by \(r\) permutes each fibre.
Let \(\rho(r)\) act by
\begin{equation}
 \rho(r)e_s=e_{rs}\quad(s\in V_{840}),\qquad
 \rho(r)|_D=I,
 \qquad b(r)=e_r-e_1\in A(S_{840})\subset N.
 \label{esn:eq:residue-action-cocycle}
\end{equation}
The two free coordinates in \(b(r)\) belong to the principal fibre.

\begin{theorem}[The full-lattice residue action and its exact cocycle]
\label{esn:thm:residue-cocycle}
The permutations \(\rho(r)\) extend to a faithful homomorphism
\(\rho:S_{840}\to O(N)\).  The map \(b:S_{840}\to N\) satisfies
\begin{equation}
 b(rs)=b(r)+\rho(r)b(s)\qquad(r,s\in S_{840}),
 \qquad b(1)=0,\qquad (b(r),b(r))=2\quad(r\ne1).
 \label{esn:eq:cocycle-identity}
\end{equation}
Thus reduction of any prime into \(S_{840}\), followed by \(b\),
is an explicit lattice-valued map on that prime domain.
\end{theorem}

\begin{proof}
Multiplication permutes the orthonormal free coordinates, preserves
their augmentation kernels, and obeys \(\rho(r)\rho(s)=\rho(rs)\).
On the \(i\)-th fibre its dual discriminant generator is
\[
 \lambda_i=e_{f_i}-\frac16\sum_{t\in F_{\eta_i}}e_t.
\]
Its image differs from it by
\(e_{rf_i}-e_{f_i}\in A(F_{\eta_i})\).  Hence \(\rho(r)\)
acts trivially on the entire A5 discriminant group, since
\(\lambda_i\) generates it.  It fixes D and therefore the whole
discriminant group, in particular \(\mathcal H\).
It follows from the inverse-image definition that it preserves \(N\).
If \(\rho(r)\) were the identity, applying it to
\(e_1-e_{289}\) would give
\(e_r-e_{289r}=e_1-e_{289}\).  Since \(289\ne1\), the
positive coordinate on the left is exactly \(r\), so comparison of
free-coordinate coefficients forces \(r=1\).  This proves faithfulness.

Finally
\[
 b(r)+\rho(r)b(s)
   =(e_r-e_1)+(e_{rs}-e_r)=e_{rs}-e_1=b(rs).
\]
For \(r\ne1\), the two basis vectors in the difference are
orthogonal unit vectors, giving squared norm two.  The stated prime map
is the composition \(p\mapsto p\bmod840\mapsto b(p\bmod840)\),
whose domains have all been specified.
\end{proof}

Fix a positive shell \(R\equiv3\pmod4\), and let \(P\) be any
finite set of positive primes whose reduction lies in \(S_{840}\).
For every \(p\in P\), use its actual count \(A_R(p)\) defined by
the reduced divisor set, including the full data of
\eqref{esn:eq:shell-data}.  Define
\begin{equation}
 \begin{split}
 w_{R,P}(r)&=\sum_{\substack{p\in P\\p\equiv r\ (840)}}A_R(p),\\
 \widehat w_{R,P}(k)&=
       \sum_{r\in S_{840}}w_{R,P}(r)
                     \exp\left(-\frac{2\pi\mathrm i k\kappa(r)}6\right),
             \qquad k\in\mathbb Z/6\mathbb Z.
 \end{split}
 \label{esn:eq:sample-pushforward}
\end{equation}
This is the pushforward along prime reduction, followed by the finite
Fourier transform with the already fixed generator 289.

\begin{theorem}[Inversion and an integral augmentation-valued count map]
\label{esn:thm:sample-fourier}
The function \(w_{R,P}\) is recovered by
\begin{equation}
 w_{R,P}(r)=\frac16\sum_{k=0}^5\widehat w_{R,P}(k)
                      \exp\left(\frac{2\pi\mathrm i k\kappa(r)}6\right).
 \label{esn:eq:fourier-inverse}
\end{equation}
Let \(u=\sum_{r\in S_{840}}e_r\),
\(W_{R,P}=\sum_rw_{R,P}(r)e_r\), and
\(M_{R,P}=\sum_rw_{R,P}(r)=\widehat w_{R,P}(0)\).
Then
\begin{equation}
 Q_{R,P}=6W_{R,P}-M_{R,P}u\in A(S_{840})\subset N
 \label{esn:eq:integral-count-augmentation}
\end{equation}
is a well-defined integral lattice-valued count map.  Put
\(\zeta_6=e^{2\pi\mathrm i/6}\) and
\[
 f_k=\sum_{j=0}^5\zeta_6^{kj}e_{289^j}\qquad(0\le k\le5).
\]
The five vectors \(f_1,\ldots,f_5\) form a basis of
\(A(S_{840})\otimes\mathbb C\), with
\begin{equation}
 \rho(289)f_k=\zeta_6^{-k}f_k,\qquad
 Q_{R,P}=\sum_{k=1}^5\widehat w_{R,P}(k)f_k.
 \label{esn:eq:fourier-augmentation-eigenvectors}
\end{equation}
Both the prime sample \(P\) and the shell \(R\) are retained in the
domain of these maps.
\end{theorem}

\begin{proof}
For an integer \(m\), the sum
\(\sum_{k=0}^5\zeta_6^{km}\) equals 6 if \(6\mid m\).
Otherwise it is the geometric sum
\((1-\zeta_6^{6m})/(1-\zeta_6^m)=0\).
Substituting \eqref{esn:eq:sample-pushforward} into the right side
of \eqref{esn:eq:fourier-inverse} therefore gives
\[
 \sum_{s\in S_{840}}w_{R,P}(s)
       \frac16\sum_{k=0}^5
              \zeta_6^{k(\kappa(r)-\kappa(s))}=w_{R,P}(r),
\]
because \(\kappa\) is bijective.  This proves inversion.

The coefficient sum of \(6W_{R,P}-M_{R,P}u\) is
\(6M_{R,P}-6M_{R,P}=0\), and its coefficients are integers, so
it belongs to the exact augmentation lattice.  The same geometric sum
shows that \(f_k\) has coefficient sum zero for \(1\le k\le5\)
and \(f_0=u\).  Moreover the six \(f_k\) are independent: the
coefficient Fourier transform of \(f_k\) is six in coordinate \(k\)
and zero in all other coordinates, by that same identity.  Since the
augmentation space has dimension five, its five nonzero-frequency
vectors form a basis.  Replacing \(j\) by \(j+1\) in
\(\rho(289)f_k\) gives the multiplier \(\zeta_6^{-k}\), including
its displayed sign.

Inversion in each free coordinate gives
\(W_{R,P}=(1/6)\sum_{k=0}^5\widehat w_{R,P}(k)f_k\).
Multiplication by 6 and subtraction of
\(M_{R,P}u=\widehat w_{R,P}(0)f_0\) prove the second identity
in \eqref{esn:eq:fourier-augmentation-eigenvectors}.
All sums were taken over the specified finite set \(P\) and its
residue image.  The transform therefore recovers these aggregate counts;
it makes no identification between the separate counts of distinct
primes with the same residue.  Their possible inequality is proved in
Theorem~\ref{esn:thm:non-descent}.
\end{proof}

% Input after esn:thm:gl2-full-lattice-lift and before the shell subsection.
% This fragment introduces no packages, macros, or theorem environments.
% New labels have prefix esn:fa:.  Its residue corollary uses the formula
% defining rho, which is restated here so placement before that formula is safe.

\subsection{The entire lattice automorphism group and its Weyl quotient}
\label{esn:fa:subsec:full-automorphisms}

Cheng, Duncan and Harvey define the umbral group as the quotient of
the Niemeier lattice automorphism group by its root-reflection subgroup;
their Table~3 identifies this quotient as \(GL_2(3)\) for
\(A_5^4D_4\) \cite[Section~2.4 and Table~3]{esn:CDH2014}.
The argument here proves that identification in the present residue,
metric, discriminant and glue coordinates, and locates the arithmetic
six-cycle in the full group. It does not claim the abstract group as new.
The displayed complement is the particular lift already constructed here;
no equality of its representation with a simple-root-preserving section
is assumed.

Retain the exact lattice \(N\), the root lattice \(R_{\mathrm{ES}}\),
the four fibre identifications \(\varphi_i\), the discriminant marking,
the glue subgroup \(\mathcal H\), and the lift \(g\mapsto L_g\)
constructed above.  For a lattice \(K\) in a Euclidean space, \(O(K)\)
means the group of real linear isometries of its real span that map
\(K\) onto itself.  For the present root lattice define
\[
 W(R_{\mathrm{ES}})
 =\langle s_\alpha:\alpha\in R_{\mathrm{ES}},\ (\alpha,\alpha)=2\rangle,
 \qquad s_\alpha(x)=x-(x,\alpha)\alpha.
\]
All actions below retain the previously fixed six-coordinate order in
each \(A_5\) summand and the marking \(10\mapsto v\),
\(01\mapsto s\), \(11\mapsto c\) in \(D^\vee/D\).

\begin{theorem}[All isometries of the displayed \(A_5\) lattice]
\label{esn:fa:thm:a5-automorphisms}
Let \(A=\{x\in\mathbb Z^6:\sum_{j=0}^5x_j=0\}\) with its displayed
Euclidean metric, and for a permutation \(\sigma\) of
\(\{0,1,2,3,4,5\}\) let \(Q_\sigma e_j=e_{\sigma(j)}\).
Every element of \(O(A)\) has exactly one expression
\(\epsilon Q_\sigma|_{\mathbb R A}\), where
\(\epsilon\in\{1,-1\}\).  Thus
\[
 O(A)=\langle Q_\sigma|_{\mathbb R A}:\sigma\in S_6\rangle
                  \times\{I,-I\},
 \qquad W(A)=\langle Q_\sigma|_{\mathbb R A}:\sigma\in S_6\rangle
                  \cong S_6.
\]
On the marked discriminant group
\(D_A=(\mathbb Z/6)\lambda\),
\(\lambda=e_0-\frac16\sum_{j=0}^5e_j\), this isometry acts by
\(k\lambda+A\mapsto\epsilon k\lambda+A\).  The kernel of this
discriminant action is exactly \(W(A)\).
\end{theorem}

\begin{proof}
The norm-two vectors are precisely
\(\alpha_{ij}=e_i-e_j\) for \(i\ne j\).  For two distinct such roots,
\[
 (\alpha_{ij},\alpha_{kl})
       =\delta_{ik}-\delta_{il}-\delta_{jk}+\delta_{jl}.
\]
Their inner product is one exactly when they have the same first index
or the same second index.  Indeed equality of first indices, with
distinct second indices, makes the displayed expression one; equality
of second indices, with distinct first indices, does the same.  If both
positive Kronecker terms vanish the expression is nonpositive, and if
both are one then the roots are equal, the excluded case.

Put
\[
 C_i=\{\alpha_{ij}:j\ne i\},\qquad
 -C_i=\{\alpha_{ji}:j\ne i\}\qquad(0\le i\le5).
\]
These twelve sets are exactly the maximal sets of roots in which every
two distinct members have inner product one.  To prove this, a singleton
can be enlarged by keeping its first index and choosing a third index,
so a maximal such set contains a pair.  If the pair is
\(\alpha_{ij},\alpha_{ik}\), with \(j\ne k\), a further root
\(\alpha_{lm}\) has inner product one with both only if
\[
 (l=i\ \text{or}\ m=j)\quad\text{and}\quad
 (l=i\ \text{or}\ m=k).
\]
Since \(j\ne k\), these conditions force \(l=i\).  Thus the whole
set is contained in \(C_i\), and maximality makes it equal to \(C_i\).
For a pair with common second index the same argument, after negating
both roots and every possible further member, gives \(-C_i\).
Each \(C_i\) and \(-C_i\) has five roots and is maximal by this argument.

Form a graph whose vertices are these twelve maximal sets, with an edge
when the two sets have nonempty intersection.  Its intersections are
exactly
\[
 C_i\cap C_j=(-C_i)\cap(-C_j)=\varnothing\quad(i\ne j),
 \qquad
 C_i\cap(-C_j)=
 \begin{cases}\{\alpha_{ij}\},&i\ne j,\\
               \varnothing,&i=j.
 \end{cases}
\]
Consequently it is the bipartite graph on the two displayed six-element
families with every cross-edge except the six pairs \((C_i,-C_i)\).
It is connected: two distinct vertices on the same side have a common
neighbor indexed by any index different from both; vertices on opposite
sides with different indices are adjacent; and \(C_i\) and \(-C_i\)
are joined by the path
\(C_i,-C_j,C_k,-C_i\) for distinct \(i,j,k\).
A connected bipartite graph has only two choices for its ordered
bipartition, because the side of a fixed vertex determines the side of
each other vertex along any connecting path.  Hence an isometry of
\(A\), which permutes roots and preserves their inner products and set
intersections, either preserves both displayed families or exchanges them.

In the first case write
\(C_i\mapsto C_{\sigma(i)}\) and
\(-C_i\mapsto-C_{\tau(i)}\).
The unique opposite-family vertex not adjacent to \(C_i\) is \(-C_i\).
Preservation of that missing edge forces \(\tau(i)=\sigma(i)\)
for each \(i\).  The singleton intersection formula then gives
\(\alpha_{ij}\mapsto\alpha_{\sigma(i)\sigma(j)}\).
In the case exchanging the families, composition with \(-I\) reduces
to the first case, so the root action is
\(\alpha_{ij}\mapsto-\alpha_{\sigma(i)\sigma(j)}\).
The roots \(e_j-e_5\), \(0\le j\le4\), span \(\mathbb R A\),
so the root action determines the whole linear isometry.  These cases
give precisely the asserted expressions.

The permutation \(\sigma\) is determined by the permutation of the
six vertices \(C_i\), after the sign is determined by whether the two
families are preserved or exchanged.  Thus the expression is unique.
Every \(Q_\sigma\) and \(-I\) preserves \(A\), and \(-I\)
commutes with every linear map; uniqueness proves the direct product.
Reflection in \(e_i-e_j\) interchanges free coordinates \(i,j\)
and fixes the others, as follows on substituting this root in
\(x\mapsto x-(x,e_i-e_j)(e_i-e_j)\).
Every permutation is a product of transpositions: each cycle
\((i_1\ i_2\ \cdots\ i_m)\) is
\((i_1\ i_m)(i_1\ i_{m-1})\cdots(i_1\ i_2)\).
These reflections therefore generate precisely the displayed \(S_6\).
Finally
\[
 Q_\sigma\lambda-\lambda=e_{\sigma(0)}-e_0\in A.
\]
Thus the discriminant action is the stated sign.  Since \(\lambda+A\)
has exact order six, its negative is distinct from itself; the action
is trivial exactly when \(\epsilon=1\).
\end{proof}

\begin{theorem}[All isometries of the displayed \(D_4\) lattice]
\label{esn:fa:thm:d4-automorphisms}
For the exact lattice
\(D=\{x\in\mathbb Z^4:\sum_i x_i\equiv0\pmod2\}\),
the action on its marked discriminant group gives the split exact sequence
\[
 1\longrightarrow W_D\longrightarrow O(D)
   \xrightarrow{\ d_D\ }GL_2(\mathbb F_2)\longrightarrow1,
 \qquad
 O(D)=W_D\rtimes\delta(GL_2(\mathbb F_2)).
\]
Here \(W_D\) consists of all signed coordinate permutations with an even
number of minus signs, and \(\delta\) is exactly the section generated
by the earlier matrices \(R_D,H_D\).  In particular \(O(D)=W_F\)
in the root and weight coordinates of
Theorem~\ref{esn:thm:d4-f4-triality}.
\end{theorem}

\begin{proof}
Every isometry of \(D\) preserves its dual: if \(T\in O(D)\),
\(x\in D^\vee\), and \(a\in D\), then
\((Tx,a)=(x,T^{-1}a)\in\mathbb Z\).
It therefore induces an additive automorphism of
\(D^\vee/D\cong\mathbb F_2^2\), giving the homomorphism \(d_D\).
The three nonzero classes have the exact norm-one representative sets
\[
 \begin{aligned}
 \mathcal V&=\{\pm e_i:1\le i\le4\},\\
 \mathcal S_+&=\{\tfrac12\textstyle\sum_i\delta_i e_i:
                                      \prod_i\delta_i=1\},\\
 \mathcal S_-&=\{\tfrac12\textstyle\sum_i\delta_i e_i:
                                      \prod_i\delta_i=-1\},
 \qquad \delta_i\in\{1,-1\}.
 \end{aligned}
\]
in the respective marked classes \(v,s,c\).
For the integral nonzero class, norm one forces exactly one coordinate
\(\pm1\).  A half-integral vector has norm at least
\(4(1/2)^2=1\), with equality exactly when every coordinate is
\(\pm1/2\).  The difference of such a vector from
\(s=(1,1,1,1)/2\) has one coordinate \(-1\) per minus sign,
so it belongs to \(D\) exactly when their number is even.
This proves the three sets and their marking directly.

If \(d_D(T)=I\), then \(T\) preserves \(\mathcal V\).
The vectors \(Te_i\) are orthonormal members of \(\mathcal V\),
so \(Te_i=\delta_i e_{\sigma(i)}\) for a permutation \(\sigma\)
of four indices and signs \(\delta_i\).  Preservation of
\(\mathcal S_+\) forces \(Ts\in\mathcal S_+\), hence
\(\prod_i\delta_i=1\).  Thus the kernel consists only of even
signed permutations.

Conversely, reflection in \(e_i-e_j\) interchanges coordinates
\(i,j\), and reflection in \(e_i+e_j\) interchanges them and
negates both.  Their composition produces a sign change in exactly
those two coordinates.  Any even set of sign changes is a product of
such disjoint pairs, and transpositions generate all permutations.
The root reflections therefore generate every even signed permutation,
and each root reflection is itself such a matrix.  If
\(\alpha\in D\) has norm two and \(x\in D^\vee\), then
\(s_\alpha x-x=-(x,\alpha)\alpha\in D\).  Every root reflection
acts trivially on the discriminant group.  This proves
\(\ker d_D=W_D\) in both directions.

For completeness the section has the original exact matrices
\[
 R_D=\operatorname{diag}(-1,1,1,1),\qquad
 H_D=\frac12
 \begin{pmatrix}
 1&1&1&1\\1&1&-1&-1\\1&-1&1&-1\\1&-1&-1&1
 \end{pmatrix}.
\]
They preserve \(D\), are orthogonal involutions, and act on
\((v,s,c)\) by \((s\ c)\) and \((v\ s)\), respectively,
as calculated in Theorem~\ref{esn:thm:d4-f4-triality}.
Their binary matrices in the retained marking are
\[
 B_R=\begin{pmatrix}1&1\\0&1\end{pmatrix},\qquad
 B_H=\begin{pmatrix}0&1\\1&0\end{pmatrix}.
\]
The same earlier coordinate calculation gives
\(R_D^2=H_D^2=(R_DH_D)^3=I\); thus their subgroup has at most six
elements, by reducing words to \((R_DH_D)^k\) or
\((R_DH_D)^kR_D\), \(k=0,1,2\).
The induced two transpositions generate the six permutations of the
three nonzero binary vectors.  There are exactly six invertible binary
matrices, by the \(3\cdot2\) choices of their independent columns,
and their action on the three nonzero vectors is faithful.  Hence the
six words have distinct discriminant actions, and the restriction of
\(d_D\) to \(\langle R_D,H_D\rangle\) is an isomorphism onto
\(GL_2(\mathbb F_2)\).  Its inverse is the specified \(\delta\).
For every \(T\in O(D)\),
\(T\delta(d_D(T))^{-1}\in\ker d_D=W_D\).
The two factors intersect trivially because \(d_D\delta=I\).
This proves the asserted splitting and exhausts \(O(D)\).
The earlier theorem proved
\(W_F=\langle W_D,R_D,H_D\rangle\), so it also proves the last equality.
\end{proof}

\begin{theorem}[The exact full automorphism quotient of \(N\)]
\label{esn:fa:thm:full-quotient}
Every isometry of \(N\) has a unique factorization
\[
 T=wL_g,\qquad w\in W(R_{\mathrm{ES}}),\quad g\in G.
\]
The Weyl group is normal, and the induced discriminant action gives
the split exact sequence and the indicated isomorphisms
\begin{equation}
 \begin{split}
 1&\longrightarrow W(R_{\mathrm{ES}})
   \longrightarrow O(N)\xrightarrow{\ d_{\mathrm{ES}}\ }G
   \longrightarrow1,\\
 O(N)&=W(R_{\mathrm{ES}})\rtimes L(G),\\
 L(G)&\xrightarrow{\ \cong\ }O(N)/W(R_{\mathrm{ES}})
               \cong G\xrightarrow{\ \Theta\ }GL_2(\mathbb F_3),\\
 W(R_{\mathrm{ES}})&\cong S_6^4\times W_D,
 \qquad |O(N)|=48\cdot(6!)^4\cdot(2^3\cdot4!).
 \end{split}
 \label{esn:fa:eq:full-group}
\end{equation}
In identified coordinates, if
\(w=(Q_{\sigma_1},Q_{\sigma_2},Q_{\sigma_3},Q_{\sigma_4};w_D)\)
and \(g=(\varepsilon,p,B)\), the exact semidirect-product action is
\begin{equation}
 L_gwL_g^{-1}
   =\bigl(Q_{\sigma_{p(1)}},Q_{\sigma_{p(2)}},
          Q_{\sigma_{p(3)}},Q_{\sigma_{p(4)}};
          \delta(B)w_D\delta(B)^{-1}\bigr).
 \label{esn:fa:eq:weyl-conjugation}
\end{equation}
In the original fibre coordinates the \(i\)-th permutation factor
on the right means
\(\varphi_iQ_{\sigma_{p(i)}}\varphi_i^{-1}\).
\end{theorem}

\begin{proof}
Theorem~\ref{esn:thm:niemeier} proves that the norm-two vectors of \(N\)
are exactly the 144 roots of its five original orthogonal summands.
Join two distinct roots by an edge when their inner product is nonzero.
The five summands are exactly the connected components of this graph.
There are no edges between different summands by orthogonality.
Within an \(A_5\) summand, the simple roots
\(e_0-e_1,e_1-e_2,e_2-e_3,e_3-e_4,e_4-e_5\)
have consecutive inner product \(-1\), forming a connected graph,
and span its real space.  Every root not already among them has nonzero
inner product with at least one of them; otherwise it is orthogonal
to its entire real space and has norm zero, contrary to norm two.
Thus its whole root graph is connected.  In \(D\), the four roots
\[
 e_1-e_2,\quad e_2-e_3,\quad e_3-e_4,\quad e_3+e_4
\]
span \(\mathbb R D\), and the second has inner product \(-1\)
with each of the other three.  The same spanning argument proves
connectedness of the \(D_4\) root graph.

An element \(T\in O(N)\) permutes the norm-two vectors and preserves
these edges, hence permutes their connected components.  Their
cardinalities are \(30,30,30,30,24\), so \(T\) preserves the
\(D_4\) component and permutes the four \(A_5\) components.
The roots generate each original lattice summand: the five and four
displayed simple roots are their integral bases by
Theorem~\ref{esn:thm:root-discriminant}.
Consequently \(T\) preserves their integer span
\(R_{\mathrm{ES}}\).  The same component argument applies to every
element of \(O(R_{\mathrm{ES}})\), since its norm-two vectors are
the same 144 roots.  It follows that every such element has, in the
fixed identified coordinates, the form
\begin{equation}
 T(X_1,X_2,X_3,X_4;d)
    =\bigl(\varepsilon_iQ_{\sigma_i}X_{p(i)}\bigr)_{i=1}^4
                                  \mathbin{;}T_Dd,
 \label{esn:fa:eq:all-root-isometries}
\end{equation}
where \(p\in S_4\), \(\varepsilon_i\in\{1,-1\}\),
\(\sigma_i\in S_6\), and \(T_D\in O(D)\).
This follows by transporting each isometry from its input \(A_5\)
summand to its output summand using the fixed \(\varphi_i\), then
applying Theorem~\ref{esn:fa:thm:a5-automorphisms}.
All choices in \eqref{esn:fa:eq:all-root-isometries} are unique.
The convention is precisely the earlier convention: a vector supported
on component \(k\) is sent to component \(p^{-1}(k)\).

Let \(\mathcal A_{\mathrm{ES}}\) be the group of transformations
of the exact discriminant group
\((\mathbb Z/6)^4\oplus\mathbb F_2^2\) given by
\[
 (x;y)\longmapsto
        \bigl((\varepsilon_i x_{p(i)})_{i=1}^4;By\bigr),
 \qquad \varepsilon_i\in\{1,-1\},\quad p\in S_4,
                 \quad B\in GL_2(\mathbb F_2).
\]
Each tuple gives a distinct transformation: the image of the order-six
generator in each individual summand determines its output summand
and its sign, and the images of the two binary basis vectors determine
\(B\).  The discriminant action of
\eqref{esn:fa:eq:all-root-isometries}, calculated by the two preceding
theorems, is exactly this tuple with \(B=d_D(T_D)\).
This defines a homomorphism
\[
 d_{\mathrm{ES}}:O(R_{\mathrm{ES}})\longrightarrow
                                      \mathcal A_{\mathrm{ES}}.
\]
It is a homomorphism because extending an isometry to the dual and then
passing to the quotient commutes with composition.  In particular the
component permutation rule remains
\(P_pP_q=P_{q\circ p}\); there is no change of the earlier convention.

This homomorphism is onto.  Given any tuple
\(a=(\varepsilon,p,B)\in\mathcal A_{\mathrm{ES}}\), define
\[
 \widetilde L_a(X_1,X_2,X_3,X_4;d)
   =\bigl(\varepsilon_iX_{p(i)}\bigr)_{i=1}^4
                                     \mathbin{;}\delta(B)d.
\]
It is a root-lattice isometry and has discriminant action \(a\).
The map \(a\mapsto\widetilde L_a\) is a homomorphism: if
\(a=(\varepsilon,p,B)\) and \(a'=(\varepsilon',q,B')\),
its composite has \(i\)-th output
\(\varepsilon_i\varepsilon'_{p(i)}X_{q(p(i))}\)
and \(D\)-component \(\delta(B)\delta(B')d=\delta(BB')d\),
which are exactly the tuple product in the stated convention.

The kernel of \(d_{\mathrm{ES}}\) consists exactly of the tuples
in \eqref{esn:fa:eq:all-root-isometries} with
\(p=1\), every \(\varepsilon_i=1\), and \(T_D\in W_D\).
It is therefore the direct product \(S_6^4\times W_D\).
Every root reflection acts on its own orthogonal summand and fixes
the other four pointwise, and the preceding two theorems identify the
groups generated inside those summands.  Hence this direct product
is exactly \(W(R_{\mathrm{ES}})\).

Every root reflection also preserves \(N\).  Indeed \(\alpha\in N\)
and integrality of \(N\) give
\(s_\alpha x=x-(x,\alpha)\alpha\in N\) for every \(x\in N\).
The reflection is an involution, so inclusion is equality.
Thus \(W(R_{\mathrm{ES}})\subset O(N)\).
It is normal there, since for \(T\in O(N)\)
\[
 Ts_\alpha T^{-1}(x)
   =x-(x,T\alpha)T\alpha=s_{T\alpha}(x),
\]
and \(T\alpha\) is again a norm-two root.

It remains to impose the actual glue, which completely determines the
remaining isometries.  For \(T\in O(R_{\mathrm{ES}})\) and
\(x\in R_{\mathrm{ES}}^\vee\), the inverse-image definition gives
\[
 Tx\in N
 \quad\Longleftrightarrow\quad
 d_{\mathrm{ES}}(T)(x+R_{\mathrm{ES}})\in\mathcal H.
\]
Consequently \(T(N)=N\) if and only if
\(d_{\mathrm{ES}}(T)(\mathcal H)=\mathcal H\).
For the forward implication, represent each class of \(\mathcal H\)
by a vector of \(N\) and apply both \(T\) and \(T^{-1}\).
For the reverse implication, the displayed equivalence applies to
\(T\) and \(T^{-1}\) because the action on \(\mathcal H\) is
bijective.  The group preserving \(\mathcal H\) inside precisely
\(\mathcal A_{\mathrm{ES}}\) is the group \(G\) of
Theorem~\ref{esn:thm:glue-gl2}, proved there to be exactly
\(GL_2(\mathbb F_3)\) through the code map \(\Theta\).
We have therefore proved the equality of subgroups
\[
 O(N)=d_{\mathrm{ES}}^{-1}(G)
                          \quad\text{inside }O(R_{\mathrm{ES}}).
\]
The section \(\widetilde L\) restricts on \(G\) to the already
specified lift \(L\).  This proves surjectivity of the restricted
map \(d_{\mathrm{ES}}:O(N)\to G\), with exactly the asserted kernel.
For \(T\in O(N)\), set \(g=d_{\mathrm{ES}}(T)\).
Then \(w=TL_g^{-1}\) lies in the kernel and \(T=wL_g\).
If \(wL_g=w'L_{g'}\), applying \(d_{\mathrm{ES}}\) gives
\(g=g'\), then cancellation gives \(w=w'\).
In particular \(L(G)\cap W(R_{\mathrm{ES}})=\{I\}\).
Normality and this unique factorization prove the semidirect product
and the quotient isomorphism, including its surjectivity.
The four \(S_6\) factors have order \((6!)^4\);
\(W_D\) has \(2^3\) sign choices and \(4!\) permutation choices,
and \(|G|=48\).  Unique factorization proves the displayed order.

Finally put \(Z=L_g^{-1}X\).  From the \(i\)-th output equation
\(X_i=\varepsilon_iZ_{p(i)}\) follows
\(Z_{p(i)}=\varepsilon_iX_i\), with each vector read in the fixed
identified \(A\) coordinates.  Therefore the \(i\)-th output of
\(L_gwL_g^{-1}X\) is
\[
 \varepsilon_iQ_{\sigma_{p(i)}}Z_{p(i)}
   =\varepsilon_iQ_{\sigma_{p(i)}}\varepsilon_iX_i
   =Q_{\sigma_{p(i)}}X_i.
\]
The \(D\)-component is
\(\delta(B)w_D\delta(B)^{-1}d\).
It belongs to \(W_D\) by normality of \(\ker d_D\).
These are exactly the components of
\eqref{esn:fa:eq:weyl-conjugation}, proving the asserted action with
all signs and component indices retained.
\end{proof}

\begin{theorem}[The residue six-cycle in the full automorphism group]
\label{esn:fa:thm:residue-weyl}
For \(r\in S_{840}=\langle289\rangle\), let
\(\rho(r)e_s=e_{rs}\) on each of the four fibres and
\(\rho(r)|_D=I\), with the exact fibre action used in
\eqref{esn:eq:residue-action-cocycle}.
Then
\begin{equation}
 \rho(S_{840})\subset W(R_{\mathrm{ES}}),\qquad
 d_{\mathrm{ES}}(\rho(r))=1,\qquad
 L_g\rho(r)L_g^{-1}=\rho(r)
       \quad(g\in G,\ r\in S_{840}).
 \label{esn:fa:eq:residue-position}
\end{equation}
In particular the map
\[
 S_{840}\times G\longrightarrow O(N),\qquad(r,g)\longmapsto\rho(r)L_g
\]
is an injective homomorphism with image of order \(6\cdot48=288\),
isomorphic to \(C_6\times GL_2(\mathbb F_3)\).
\end{theorem}

\begin{proof}
Put \(C e_j=e_{j+1\bmod6}\) on the free six-coordinate space and
restrict it to \(\mathbb R A\).  If
\(r=289^m\), with \(m\in\{0,1,2,3,4,5\}\), the exact fibre
identification gives
\[
 \rho(r)e_{f_i289^j}=e_{f_i289^{j+m}},\qquad
 \rho(r)|_{A(F_{\eta_i})}=\varphi_iC^m\varphi_i^{-1}.
\]
Thus in all four identified summands \(\rho(r)\) is the same
coordinate permutation \(C^m\), and it is the identity on \(D\).
Theorem~\ref{esn:fa:thm:a5-automorphisms} puts it in the corresponding
product of Weyl groups, giving the first two assertions.
Each sign \(\varepsilon_i I\) commutes with \(C^m\), and the
same operator \(C^m\) occurs in every component.  Directly the
\(i\)-th outputs of \(L_g\rho(r)\) and \(\rho(r)L_g\) are
respectively
\(\varepsilon_iC^mX_{p(i)}\) and
\(C^m\varepsilon_iX_{p(i)}\), which agree; the \(D\)-outputs
also agree since \(\rho(r)|_D=I\).
This proves the conjugation identity without changing the chosen origins.

The six restrictions \(C^m|_{\mathbb R A}\) are distinct: if
\(C^m(e_0-e_1)=e_0-e_1\), comparison of the unique coordinate with
coefficient \(+1\) gives \(m=0\pmod6\).
Thus \(\rho\) is faithful.  Commutation makes the displayed product
map a homomorphism.  If \(\rho(r)L_g=I\), then
\(L_g=\rho(r)^{-1}\) lies in
\(L(G)\cap W(R_{\mathrm{ES}})=\{I\}\).
Faithfulness of both factors gives \(r=1\) and \(g=1\).
This proves injectivity and the exact image order.
\end{proof}

% Exact proof module, 6 September 2026. Input after es_niemeier_shared.tex.
% All original coordinates and Euclidean root lengths are retained.
\section{The exact Leech neighbor of the retained index-72 lattice}
\label{esnl:sec:neighbor}

Retain the lattice, metric, discriminant marking, and glue literally:
\[
 \begin{split}
 A&=\{(x_0,\ldots,x_5)\in\mathbb Z^6:\sum x_j=0\},\\
 D&=\{(d_1,\ldots,d_4)\in\mathbb Z^4:\sum d_i\equiv0\pmod2\},
 \qquad R=A^{\perp4}\perp D,\\
 C_2&=\langle b_0=(1111;00),b_1=(1100;10),b_2=(1010;01)\rangle_{\mathbb F_2},\\
 C_3&=\langle z_1=(1,1,1,0),z_2=(1,2,0,1)\rangle_{\mathbb F_3},\\
 H&=\{(3u+2z;y):(u;y)\in C_2,\ z\in C_3\},\\
 N&=\{x\in R^\vee:x+R\in H\}.
 \end{split}
\]
The coordinates in each $A$ remain the six ordered elements $f_i289^j$
of its original residue fibre. The $A$ discriminant generator is
$\lambda=(5,-1,-1,-1,-1,-1)/6$. The $D$ markings $10,01,11$
remain represented by $v_D=(1,0,0,0)$,
$s_D=(1,1,1,1)/2$, $c_D=(-1,1,1,1)/2$.
The preceding proof establishes that $N$ is even unimodular,
$[N:R]=72$, $\det R=5184$, and its roots are exactly those of
$A_5^4D_4$. Every root has squared norm two. No metric rescaling
occurs anywhere below.

\subsection{The Weyl vector and the exact neighbor}

Use the original simple roots $a_j=e_j-e_{j+1}$, $0\le j\le4$,
in each $A$ component and
\[
 d_1=e_1-e_2,\quad d_2=e_2-e_3,\quad
 d_3=e_3-e_4,\quad d_4=e_3+e_4
\]
in $D$. Their highest roots are
$\theta_A=e_0-e_5=\sum_ja_j$ and
$\theta_D=e_1+e_2=d_1+2d_2+d_3+d_4$.
Both highest-root heights are five, so the common Coxeter number
is six. Define
\[
 \rho_A=(5,3,1,-1,-3,-5)/2,\qquad
 \rho_D=(3,2,1,0),\qquad \rho=\rho_A^{\oplus4}\perp\rho_D.
 \label{esnl:eq:rho}
\]
Each simple root pairs to one with $\rho$.
Moreover $\rho_A-3\lambda=(0,2,1,0,-1,-2)\in A$ and
$\rho_D\in D$. Thus the exact glue class of $\rho$ is
$(3,3,3,3;00)$, the image of $b_0$, and $\rho\in N$.
Directly
\[
 \rho_A^2=35/2,\quad \rho_D^2=14,\quad
 \rho^2=4(35/2)+14=84=2\cdot6\cdot7.
\]
Consequently $(\rho,N)\subset\mathbb Z$.
Fix $\alpha=a_0$ in the first $A$ component and put
\[
 \ell_\rho:N\to\mathbb Z/6,\quad x\mapsto(\rho,x)\bmod6,\qquad
 K=\ker\ell_\rho,\quad R_0=R\cap K,\quad
 v=\rho/6-\alpha,\quad \Lambda=K+\mathbb Z v.
 \tag{NL1}
\]
The pairing $(\rho,\alpha)=1$ proves
$[N:K]=[R:R_0]=6$. For $x\in K$, the number
$(v,x)=(\rho,x)/6-(\alpha,x)$ is integral, and
\[
 v^2=84/36-2/6+2=4,\qquad
 6v=\rho-6\alpha\in K,\qquad(\rho,6v)=78.
\]
Hence $\Lambda$ is positive definite even integral of rank 24.
If $mv\in N$, its pairing after adding $m\alpha$ with $\alpha$
is $m/6\in\mathbb Z$, so $6\mid m$.
It follows that
\[
 N\cap\Lambda=K,\qquad [\Lambda:K]=6,\qquad
 \det K=36,\qquad \det\Lambda=36/6^2=1.
 \tag{NL2}
\]
Changing $\alpha$ to another root of $\rho$-pairing one changes
$v$ by an element of $K$, and so does not change $\Lambda$.
We shall use the equivalent exact congruence description
\[
 \Lambda=\{x+k\rho/6:x\in N,\ k\in\mathbb Z,\
                                (\rho,x)+k\equiv0\pmod6\}.
 \tag{NL3}
\]
Indeed write $x+k\rho/6=(x+k\alpha)+kv$; the first summand
is in $K$ exactly under this congruence. Changing $k$ by six
can be absorbed in $x$ using $\rho\in N$, and preserves the
congruence since $\rho^2=84$.

\subsection{A complete rootlessness proof}

\begin{theorem}[The Leech neighbor in the retained coordinates]
\label{esnl:thm:leech-neighbor}
The lattice $\Lambda=K+\mathbb Z(\rho/6-\alpha)$ is positive
definite even unimodular of rank 24 and has minimum squared norm
exactly four. It is the holy-construction Leech lattice of the
retained $A_5^4D_4$ glue, and
$N\cap\Lambda=K$ with $[N:K]=[\Lambda:K]=6$.
\end{theorem}

\begin{proof}
The rank, parity, unimodularity and indices were proved in
(NL1)--(NL2). The holy-presentation identification is proved in
(NL8)--(NL11) below. We now prove the required minimum directly.
For $k=0$ in (NL3), a possible root is an original root of $N$.
An $A$ root pairs with $\rho$ to one of $\pm1,\ldots,\pm5$.
For $D$, all roots $\pm e_i\pm e_j$ have nonzero $\rho_D$-pairing
of absolute value at most five. Thus none belongs to $K$.
Negation and the reduction of $k$ modulo six leave $k=1,2,3$.

Define the exact shifted class minima by
\[
 m_A(k,t)=\min_{x\in t\lambda+A}|x+k\rho_A/6|^2,\qquad
 m_D(k,y)=\min_{x+D=y}|x+k\rho_D/6|^2.
\]
Their values are
\[
 \begin{array}{c|rrrrrr}
 k\backslash t&0&1&2&3&4&5\\\hline
 1&35/72&35/72&35/72&35/72&35/72&35/72\\
 2&11/18&4/9&11/18&4/9&11/18&4/9\\
 3&3/8&17/24&17/24&3/8&17/24&17/24
 \end{array}
 \quad
 \begin{array}{c|rrrr}
 k\backslash y&00&10&01&11\\\hline
 1&7/18&7/18&7/18&7/18\\
 2&2/9&5/9&5/9&5/9\\
 3&1/2&1/2&1/2&1/2
 \end{array}.
 \tag{NL4}
\]
We prove the table and its needed equality information.
The shifted $A$ coordinates are precisely
\[
 \left(n_j-\frac t6+\frac{k(5-2j)}{12}\right)_{j=0}^5,
               \qquad n_j\in\mathbb Z,\quad\sum n_j=t.
 \tag{NL5}
\]
At $k=1$, coordinatewise nearest representatives are the six
numbers $\pm1/12,\pm3/12,\pm5/12$ in an order depending on $t$.
They sum to zero and their squares sum to $35/72$.
At $k=2$ and even $t$, the nearest coordinates consist of
two copies of each sign of $1/6$ and one copy of each sign of
$1/2$; the half-integer ties admit exactly the needed zero-sum
choice and the squared sum is $11/18$.
At odd $t$, there are two each of $0,1/3,-1/3$, giving $4/9$.
For $t=1,3,5$ their actual ordered shifted vectors are
\[
 (-1,1,0,-1,1,0)/3,\quad
 (1,0,-1,1,0,-1)/3,\quad
 (0,-1,1,0,-1,1)/3.
 \tag{NL6}
\]
Subtracting $\rho_A/3$ gives original squared norms
$17/6,3/2,17/6$, respectively.
At $k=3,t=0,3$, the nearest coordinates alternate $\pm1/4$,
giving $3/8$ with zero sum.
At each other $t$, the coordinatewise nearest representatives
have three absolute values $1/12$ and three $5/12$, square sum
$13/24$, and coordinate sum $1$ or $-1$.
Meeting the zero-sum constraint requires an integer change.
Every nonzero coordinate change increases the squared sum by
at least $(7/12)^2-(5/12)^2=1/6$.
Changing one suitable $5/12$ coordinate attains this increase
and restores the zero sum. Thus the constrained minimum is
$17/24$. At $t=0,3$, the original minimizing vectors are
\[
 (-1,-1,0,0,1,1),\qquad (-3,-1,-1,1,1,3)/2,
 \tag{NL7}
\]
of squared norms $4,11/2$.

For $D$ and $k=1$ the shift is $(1/2,1/3,1/6,0)$.
The nearest integer coordinate distances are
$(1/2,1/3,1/6,0)$; the first-coordinate tie meets either integer
parity. The nearest half-integer distances are
$(0,1/6,1/3,1/2)$; the last-coordinate tie meets either
half-integer class. Both squared sums are $7/18$.
At $k=2$ the shift is $(1,2/3,1/3,0)$.
The unique nearest integer vector before shifting is
$(-1,-1,0,0)\in D$, giving $2/9$.
Changing parity costs at least $1/3$, attained by changing the
second or third coordinate, giving $5/9$.
The half-integer coordinate distances are
$(1/2,1/6,1/6,1/2)$, giving $5/9$; ties meet either marking.
At $k=3$ the shift is $(3/2,1,1/2,0)$.
Integer-class minimum distances are $(1/2,0,1/2,0)$,
and half-integer-class distances $(0,1/2,0,1/2)$.
Ties meet every marking and all square sums are $1/2$.
In class $00$, the two original minimizing vectors are
\[
 (-1,-1,0,0),\qquad(-2,-1,-1,0),
\]
of squared norms $2,6$. In class $10$, they are
\[
 (-1,-1,-1,0),\qquad(-2,-1,0,0),
\]
of squared norms $3,5$.
In either half-integer class they have the shapes
$(-3/2,-1/2,-1/2,\pm1/2)$ and
$(-3/2,-3/2,-1/2,\pm1/2)$.
Exactly one last-coordinate sign in each shape meets each
original marking; their squared norms are $3,5$.
This proves every entry and the relevant equality cases in (NL4).

For $k=1$, every full shifted vector has squared norm at least
$4(35/72)+7/18=7/3>2$.
For $k=2$, the lower bound from (NL4) is two.
Equality forces all $t_i$ odd and $y=00$.
In the actual glue this forces $u=1111$.
If $z=0$, all $t_i=3$ and (NL6) gives
$x^2=4(3/2)+2=8$.
If $z\ne0$, its weight is three, so exactly one $t_i=3$
and three $t_i\in\{1,5\}$; then
$x^2=3/2+3(17/6)+2=12$.
The norm-two equation is
\[
 2=x^2+\tfrac23(\rho,x)+\tfrac{28}{3},
 \qquad(\rho,x)=-11-\tfrac32x^2.
\]
Its pairing is $-23$ or $-29$, always $1$ modulo six.
Condition (NL3) requires residue $4$, so all nine equality
candidates, one per ternary word, are excluded.

For $k=3$, the lower bound two requires $z=0$:
every nonzero ternary word raises the minimum by
$3(17/24-3/8)=1$.
If $u=0$, the possible original squared norms are $18,22$.
If $u=1111$, they are $24,28$.
If $u$ has weight two, its $D$ class is nonzero and they are
$2(4)+2(11/2)+3=22$ or $2(4)+2(11/2)+5=24$.
These are all sixteen equality candidates.
But
\[
 2=x^2+(\rho,x)+21,\qquad(\rho,x)=-19-x^2,
\]
and (NL3) would require $x^2\equiv2\pmod6$.
None of $18,22,24,28$ has that residue.
Thus $\Lambda$ has no roots. Together with (NL2), this proves
that it is positive definite even unimodular of rank 24,
with minimum squared norm exactly four, attained by $v$.
\end{proof}

\subsection{The exact holy-construction presentation}

For $0\le t\le5$ set
\[
 \omega_t=(\underbrace{1,\ldots,1}_{t},
             \underbrace{0,\ldots,0}_{6-t})-\tfrac t6(1,\ldots,1).
\]
Then $\omega_t+A=t\lambda+A$ and
$(\rho_A,\omega_t)=t(6-t)/2=3\omega_t^2$.
For the $D$ markings set
\[
 \omega_{00}=0,\quad\omega_{10}=e_1,\quad
 \omega_{01}=(1,1,1,1)/2,\quad
 \omega_{11}=(1,1,1,-1)/2.
\]
The last differs from the original $c_D$ by $(1,0,0,-1)\in D$,
so retains its exact marking.
For $a=(t_1,t_2,t_3,t_4;y)\in H$ define
\[
 u_a=(\omega_{t_1},\omega_{t_2},\omega_{t_3},\omega_{t_4};\omega_y),
 \qquad g_a=u_a-\rho/6.
 \tag{NL8}
\]
These give all 72 representatives of $N/R$. Each nonzero $D$
representative has norm one and pairing three; hence
\[
 (\rho,u_a)=3u_a^2\in6\mathbb Z,\quad
 u_a\in K,\quad g_a^2=u_a^2-(\rho,u_a)/3+84/36=7/3.
 \tag{NL9}
\]
Let $\mathcal F$ be the 29 vectors consisting, in each original
component, of the negative original simple roots and the positive
highest root. Their pairings with $\rho$ are $-1$ or $5$.
There are positive integral zero relations of coefficient sum six:
\[
 \theta_A+\sum_{j=0}^4(-a_j)=0,\qquad
 \theta_D+(-d_1)+2(-d_2)+(-d_3)+(-d_4)=0.
 \tag{NL10}
\]
Every component of $g_a$ pairs to $1/6$ with the corresponding
vertices except one tip, where the pairing is $-5/6$.
For $A$ this follows because $\omega_t$ has simple-root pairing
one only at $a_{t-1}$ when $t\ne0$, and has
$(\omega_t,\theta_A)=1$ for $t\ne0$.
For $D$ the three nonzero representatives have simple-root
pairing one only at $d_1,d_4,d_3$, respectively; all have
highest-root pairing one. The zero representative points to
the highest root. These computations fix all orientations in
the holy generator diagram.

The following coefficient presentations are exact equalities:
\[
 \begin{split}
 N&=\left\{\sum_fa_ff+\sum_ab_ag_a:
                    a_f,b_a\in\mathbb Z,\ \sum_ab_a=0\right\},\\
 \Lambda&=\left\{\sum_fa_ff+\sum_ab_ag_a:
            a_f,b_a\in\mathbb Z,\ \sum_fa_f+\sum_ab_a=0\right\}.
 \end{split}
 \tag{NL11}
\]
For the first, the $f$ span $R$ and $g_a-g_0=u_a$ generate its
given extension $N$; the zero coefficient sum cancels $\rho/6$.
For the second, write $A=\sum a_f$, $B=\sum b_a$,
$r=\sum a_ff$, and $x=r+\sum b_au_a$.
The vector is $x-B\rho/6$, and by (NL9)
$(\rho,x)\equiv-A\equiv B$, exactly (NL3) with $k=-B$.
Conversely, for any $x+k\rho/6$ in (NL3), write
$x=r+\sum b_au_a$. The coefficient of $u_0=0$ can be set so
$B=-k$. Write $r$ with the negative simple-root vertices.
The congruence gives $A+B\equiv0\pmod6$.
Adding a suitable multiple of a zero relation (NL10) makes
$A+B=0$ exactly. Thus both directions of (NL11) are proved.

These are exactly the holy coefficient rules of Conway--Sloane
\cite[pp.~277--280]{esn:ConwaySloane1982}.
Borcherds gives the uniform proof of the holy-construction
theorem \cite[Section~7, especially the coefficient rules
(i)--(ii) and Lemma~7.3]{esn:Borcherds1985}.
Thus (NL11) is an explicit realization of that Leech
construction from the actual 72 marked glue words; all its
metric, determinant, integrality and rootlessness assertions
have also been independently proved here.

\subsection{The indices, discriminant morphism, and cyclic order twelve}

Reduction modulo $R$ gives an isomorphism $K/R_0\cong H$:
its kernel is $R_0$, and (NL9) supplies representatives in $K$
for every $H$ class. Consequently
\[
 [K:R_0]=72,\quad [N:R_0]=[\Lambda:R_0]=432,\quad
 \det R_0=6^2\cdot5184=186624.
 \tag{NL12}
\]
Thus the original and new inclusions, with every index, are
\[
 R_0\subset R\subset N,\quad [R:R_0]=6,\ [N:R]=72;
 \qquad
 R_0\subset K\subset\Lambda,\quad[K:R_0]=72,\ [\Lambda:K]=6.
\]
In particular $R\cap\Lambda=R_0$, an exact common sublattice.

\begin{theorem}[The exact order-twelve quotient generator]
\label{esnl:thm:cyclic-twelve}
The quotient of $\Lambda$ by the original height-zero root
sublattice $R_0$ is
$\Lambda/R_0\cong\mathbb Z/12\oplus(\mathbb Z/6)^2$.
The class of $v=\rho/6-\alpha$ has exact order twelve.
\end{theorem}

\begin{proof}
Let $\bar b_i$ denote the image in $H$ of the binary generator
$b_i$, namely its first four coordinates multiplied by three
and its last two coordinates retained.
Since $6v=\rho-6\alpha$ and $6\alpha\in R_0$,
\[
 6[v]=\bar b_0\quad\text{in }\Lambda/R_0.
 \tag{NL13}
\]
This is a nonzero order-two element. The classes
$\bar b_0,\bar b_1,\bar b_2,2z_1,2z_2$ form the displayed
independent binary and ternary generators of $H$.
Eliminate $\bar b_0$ using (NL13). The resulting surjective
presentation is
\[
 \Lambda/R_0\cong
 \mathbb Z/12\oplus(\mathbb Z/2)^2\oplus(\mathbb Z/3)^2
 \cong\mathbb Z/12\oplus(\mathbb Z/6)^2.
 \tag{NL14}
\]
There are no further relations: this presentation has order
$12\cdot4\cdot9=432$, equal to (NL12).
Thus $[v]$ has exact order twelve. This particular occurrence
of twelve is proved from the Coxeter denominator six and the
actual order-two Weyl-vector glue class.
\end{proof}

Set $M=N+\Lambda=N+\mathbb Z(\rho/6)$.
Then
\[
 M=K^\vee,\qquad K^\vee/K=
    \langle[\alpha],[v]\rangle\cong(\mathbb Z/6)^2.
 \tag{NL15}
\]
Indeed $M\subset K^\vee$, its index over $N$ is six by pairing
with $\alpha$, and $[K^\vee:N]=[N:K]=6$.
The cyclic images of $N/K$ and $\Lambda/K$ have order six and
trivial intersection, proving the stated generators.
Their exact discriminant data are
\[
 q_K([\alpha])=q_K([v])=0,\quad
 b_K([\alpha],[v])=1/6\pmod{\mathbb Z},\quad
 q_K(a[\alpha]+b[v])=ab/3\pmod{2\mathbb Z}.
 \tag{NL16}
\]
These follow from $\alpha^2=2$, $v^2=4$ and
$(\alpha,v)=1/6-2=-11/6$.
The two original lattices are therefore the inverse images
of the two indicated order-six isotropic subgroups in this
specified discriminant form.

\subsection{An exact Lorentzian quotient and symmetry transport}

Take $\mathcal I=N\perp U$ in coordinates $(x,m,n)$ with
\[
 \langle(x,m,n),(x',m',n')\rangle
 =(x,x')-mn'-m'n .
\]
Thus $U$ has Gram matrix
$\left(\begin{smallmatrix}0&-1\\-1&0\end{smallmatrix}\right)$.
The vectors $z=(0,0,1)$ and $w=(\rho,6,7)$ are primitive
isotropic: $w^2=84-2\cdot6\cdot7=0$, and six and seven are
coprime. There are integral quotient isometries
\[
 \begin{split}
 z^\perp/\mathbb Zz&\longrightarrow N,\quad[(x,0,n)]\mapsto x,\\
 w^\perp/\mathbb Zw&\longrightarrow\Lambda,\quad
                                  [(x,m,n)]\mapsto x-m\rho/6.
 \end{split}
 \tag{NL17}
\]
For the second map, orthogonality means $(\rho,x)=7m+6n$,
giving (NL3) with $k=-m$. Conversely take $m=-k$ and
$n=((\rho,x)+7k)/6$ for a vector in (NL3).
If the image vanishes, $x=m\rho/6$; pairing with $\alpha$
forces $m=6j$, and then $n=7j$, so the kernel is exactly
$\mathbb Zw$. Finally
\[
 |x-m\rho/6|^2
 =x^2-\tfrac m3(\rho,x)+\tfrac73m^2=x^2-2mn.
\]
Polarization gives the full pairing, proving the isometry.

For every $T\in O(N)$ the exact transport laws are
\[
 T(K_\rho)=K_{T\rho},\qquad
 T(\Lambda_\rho)=\Lambda_{T\rho},\qquad
 (x,m,n)\mapsto(Tx,m,n):w_\rho\mapsto w_{T\rho}.
 \tag{NL18}
\]
They follow by substitution in (NL3).
In particular the stabilizer of $\rho$ acts on the fixed
$\Lambda_\rho$. General isometries act on this specified
family of neighbors; no assertion that an arbitrary chosen
automorphism complement fixes this Weyl vector is required.

\subsection{The exact $F_4$ Coxeter action of order twelve}

The order twelve in (NL14) is a lattice quotient invariant of
this explicitly marked construction. There is also the ambient
$F_4$ Coxeter action on the same original $D_4$ real root space.
Use the precise $F_4$ simple roots
\[
 \beta_1=e_2-e_3,\quad\beta_2=e_3-e_4,\quad
 \beta_3=e_4,\quad\beta_4=(e_1-e_2-e_3-e_4)/2,
 \qquad C=s_{\beta_1}s_{\beta_2}s_{\beta_3}s_{\beta_4},
\]
where $s_\beta(x)=x-2(x,\beta)\beta/\beta^2$ and the rightmost
factor acts first. The first two roots have squared norm two,
and the last two squared norm one. No rescaling of the $D_4$
roots has occurred.

\begin{theorem}[Coxeter twelve and its marked triality quotient]
\label{esnl:thm:f4-coxeter-twelve}
In the retained four-coordinate metric,
\[
 C=\frac12
 \begin{pmatrix}
 1&1&1&1\\-1&1&1&-1\\1&1&-1&-1\\1&-1&1&-1
 \end{pmatrix},\qquad
 C^3=
 \begin{pmatrix}
 0&1&0&0\\-1&0&0&0\\0&0&0&-1\\0&0&1&0
 \end{pmatrix}.
 \tag{NL19}
\]
It satisfies $C^4-C^2+I=0$, $C^6=-I$ and has exact order
twelve. Its marked discriminant action is
$v_D+D\mapsto c_D+D\mapsto s_D+D\mapsto v_D+D$.
The restriction of the original triality quotient is the exact
sequence
\[
 1\longrightarrow\langle C^3\rangle\cong\mathbb Z/4
 \longrightarrow\langle C\rangle\cong\mathbb Z/12
 \longrightarrow\langle(v_D\,c_D\,s_D)\rangle\cong\mathbb Z/3
 \longrightarrow1,
 \tag{NL20}
\]
with $\langle C^3\rangle\subset W(D_4)$.
\end{theorem}

\begin{proof}
The first reflection exchanges coordinates two and three; the
second exchanges coordinates three and four; the third changes
the sign of coordinate four; the last is
$I-2\beta_4\beta_4^T$. Their product gives the first matrix in
(NL19). Multiplying it three times gives the second, whose
square is $-I$. Multiplying the first matrix also gives
$C^4-C^2+I=0$, and hence $C^{12}=I$.
Its first column gives
$Cv_D=c_D+(1,-1,0,0)$, with the displayed correction in $D$.
Directly $Cs_D=v_D$ and $Cc_D=(0,0,-1,-1)+s_D$.
The displayed integer correction lies in $D$ and proves the
stated three-cycle with the original $v_D,s_D,c_D$ markings.
The first two reflections preserve $D$ as root reflections.
The last two preserve $D$ because a sign change preserves
integer sum parity, and reflection in $\beta_4$ sends an
integer even-sum vector to an integer vector of even sum.
The inverse reflections have the same property, so these
maps are isometries of $D$.
Consequently their induced action on $D^\vee/D$ is defined
and has the exact order-three image calculated above.
The matrix $C^3$ is an even signed permutation of exact order
four, and thus lies in the already identified group $W(D_4)$.
If $C^m=I$, its discriminant image forces $3\mid m$,
and the exact order of $C^3$ then forces $4\mid m/3$.
Thus $C$ has exact order twelve and the discriminant kernel
on its cyclic subgroup is exactly $\langle C^3\rangle$.
This proves every map and group order in (NL20).
\end{proof}

In particular $C^2$ still has nontrivial order-three
discriminant action, while $C^3$ is in the long-root Weyl
group. The Coxeter denominator six in (NL1), the cyclic order
twelve in (NL14), and the $F_4$ Coxeter twelve in (NL20)
therefore have their stated, distinct domains and explicit
relationships through the original $D_4$ marking.

% Staged only. No canonical file has been edited.
% Requires amsmath, amssymb, amsthm, and the existing esn theorem labels.
% All new labels have namespace n24:.

\subsection{An exact cyclic-triality isometry of the entire Niemeier space}
\label{n24:sec:cyclic-isometry}

Retain the ordered fibres
\(f_i289^j\), \(1\leq i\leq4\), \(0\leq j\leq5\),
the Euclidean augmentation coordinates of \(A(F_{\eta_i})\), the
ordered orthonormal coordinates of \(D\subset\mathbb R^4\), and the
72-element glue \(\mathcal H\) of
\eqref{esn:eq:glue-and-overlattice}.  Thus the actual real span is
\[
 V_N=N\otimes_{\mathbb Z}\mathbb R
 =\left\{(x_{i,j};d_k)\in\mathbb R^{24}\oplus\mathbb R^4:
               \sum_{j=0}^5x_{i,j}=0\ (1\leq i\leq4)\right\},
\]
with squared norm \(\sum_{i,j}x_{i,j}^2+\sum_kd_k^2\).
The four displayed linear constraints, rather than an implicit change of
metric, give its dimension 24.  Choose an ordered orthonormal octonion
basis \((u_0,\ldots,u_7)\) with \(u_0=1\).  Every octonion product
below is the product of the chosen real octonion algebra, without a
change of multiplication convention.

For \(p\in\{0,1\}\) and \(j\in\{0,1,2\}\), set
\begin{align}
 m_{i,p}&=\frac13\sum_{j=0}^2x_{i,p+2j},&
 t_{i,p,j}&=x_{i,p+2j}-m_{i,p},\label{n24:eq:triplets}\\
 \alpha_i&=\frac{\sum_{j=0}^2x_{i,2j}
                    -\sum_{j=0}^2x_{i,1+2j}}{\sqrt6},&
 (b_0,\ldots,b_7)&=(\alpha_1,\alpha_2,\alpha_3,\alpha_4,
                           d_1,d_2,d_3,d_4).
 \label{n24:eq:fixed-coordinates}
\end{align}
In particular \(m_{i,0}=\alpha_i/\sqrt6\),
\(m_{i,1}=-\alpha_i/\sqrt6\), and
\(\sum_jt_{i,p,j}=0\).  The ordered map is
\begin{equation}
 \begin{split}
 \mathcal L:V_N&\longrightarrow\mathbb O^3,\\
 (x;d)&\longmapsto(q_1,q_2,q_3),\\
 q_{j+1}&=\sum_{i=1}^4\sum_{p=0}^1
       \left(t_{i,p,j}+\frac{b_{2(i-1)+p}}{\sqrt3}\right)
                                          u_{2(i-1)+p}.
 \end{split}
 \label{n24:eq:full-isometry}
\end{equation}

\begin{theorem}[Actual rank-24 map, inverse, pairing, and residue action]
\label{n24:thm:full-isometry}
The map \(\mathcal L\) is a real linear isometry onto
\(\mathbb O^3\) with its product norm.  Write
\(q_{j+1}=\sum_{a=0}^7q_{j,a}u_a\), and define
\begin{equation}
 b_a=\frac1{\sqrt3}\sum_{j=0}^2q_{j,a},\qquad
 t_{i,p,j}=q_{j,2(i-1)+p}
                    -\frac13\sum_{h=0}^2q_{h,2(i-1)+p}.
 \label{n24:eq:inverse-first}
\end{equation}
Its inverse is exactly
\begin{equation}
 x_{i,p+2j}=t_{i,p,j}+
                 \frac{(-1)^p b_{i-1}}{\sqrt6},
 \qquad d_k=b_{k+3}.
 \label{n24:eq:inverse-second}
\end{equation}
If \(r=\rho(289^2)=\rho(361)\), with the residue action
of Theorem~\ref{esn:fa:thm:residue-weyl}, and
\(c(q_1,q_2,q_3)=(q_3,q_1,q_2)\), then
\begin{equation}
 \mathcal Lr=c\mathcal L,\qquad r^3=c^3=I,
 \qquad rN=N.
 \label{n24:eq:cyclic-intertwining}
\end{equation}
Thus \(\mathcal L|_N\) embeds the original even unimodular lattice,
with every pairing and glue class retained, into the actual three
octonion sectors and intertwines an actual order-three lattice
isometry with cyclic triality.
\end{theorem}

\begin{proof}
The two means sum to zero because each six-coordinate fibre has total
sum zero.  The equations for \(m_{i,p}\) follow from
\(\alpha_i=6m_{i,0}/\sqrt6\).  Orthogonality to constant triplets gives
\[
 \sum_{j=0}^5x_{i,j}^2
 =\sum_{p=0}^1\sum_{j=0}^2t_{i,p,j}^2
                    +3m_{i,0}^2+3m_{i,1}^2
 =\sum_{p,j}t_{i,p,j}^2+\alpha_i^2.
\]
For each \(a=2(i-1)+p\), the same zero-sum identity gives
\[
 \sum_{j=0}^2\left(t_{i,p,j}+\frac{b_a}{\sqrt3}\right)^2
                      =\sum_{j=0}^2t_{i,p,j}^2+b_a^2.
\]
Summing over all eight values of \(a\) proves
\(\sum_{h=1}^3|q_h|^2=\sum_{i,j}x_{i,j}^2+\sum_kd_k^2\).
Polarization proves equality of all pairings.  Equations
\eqref{n24:eq:inverse-first} recover the means and residuals of each
output triplet.  Equation~\eqref{n24:eq:inverse-second} then recovers
the original means, the original fibre coordinates, and \(d\).
Conversely those formulas applied to arbitrary \(q\) yield
\(\sum_jx_{i,j}=0\) and return the prescribed \(q\).  This proves
bijectivity without a dimension-only argument.

The residue action sends the coordinate vector \(e_{f_i289^h}\)
to \(e_{f_i289^{h+2}}\), so the coefficient in output position
\(p+2j\) is the old coefficient in position \(p+2(j-1)\), with
\(j\) read modulo three.  Both means, all \(\alpha_i\), and all
\(d_k\) are fixed.  Hence the three output octonions become
\((q_3,q_1,q_2)\), proving the intertwining formula.  Each coordinate
permutation acts trivially on \(A^\vee/A\): for
\(\lambda=e_0-\frac16\sum_he_h\), its difference from \(\lambda\)
is an augmentation root.  The \(D\)-coordinates are fixed.  Thus
\(r\) acts trivially on the full discriminant group, preserves
\(\mathcal H\), and preserves its exact inverse image \(N\).
Its order is three, as seen on \(e_0-e_1\) in any fibre.
\end{proof}

\paragraph{The exact octonion-coordinate glue test.}
This embedding has the following fully explicit image, denoted
\(\Lambda_{\mathbb O}=\mathcal L(N)\).  Given \(q\), first recover
\((x;d)\) by \eqref{n24:eq:inverse-first}--\eqref{n24:eq:inverse-second}.
For each fibre require
\(x_{i,h}-x_{i,k}\in\mathbb Z\) for every \(h,k\), and put
\[
 k_i=-6x_{i,0}\pmod6\in\mathbb Z/6\mathbb Z.
\]
This is defined because
\(6x_{i,0}=\sum_h(x_{i,0}-x_{i,h})\in\mathbb Z\).
Require also that the four \(d_k\) are all integral or all
half-integral.  Let \(y(d)\in\mathbb F_2^2\) be the unique class with
\[
 d+D=0+D,\ v+D,\ s+D,\ c+D
 \quad\hbox{marked respectively by }00,10,01,11.
\]
Then the final and exact test is
\begin{equation}
 (k_1,k_2,k_3,k_4;y(d))
   =\bigl(3u+2(a+b,a+2b,a,b);\tau(u)\bigr)
 \label{n24:eq:image-glue}
\end{equation}
for \(u\in\mathbb F_2^4\) of even weight,
\((a,b)\in\mathbb F_3^2\), with the first four coordinates read
modulo six, and \(\tau(u)=(u_2+u_4,u_3+u_4)\).
Indeed the integral-difference conditions are precisely membership in
\(A^\vee\), and the common fractional part equals
\(-k_i/6\pmod{\mathbb Z}\), which is the fractional part of
\(k_i\lambda\).  The half-integral test is precisely membership in
\(D^\vee\), and \eqref{n24:eq:image-glue} is the original
\(\mathcal H\).  This proves both directions of the image
description.  No new lattice or rescaled glue is substituted.

\paragraph{The determinant and the factor two in the Albert metric.}
The matrix convention, also specifying this construction independently
of its host document, is
\[
 Q_{\mathbb O}(\lambda;q)=
 \begin{pmatrix}
 \lambda&q_3&\overline q_2\\
 \overline q_3&\lambda&q_1\\
 q_2&\overline q_1&\lambda
 \end{pmatrix}\in\mathfrak h_3(\mathbb O),\qquad
 X\circ Y=(XY+YX)/2.
\]
Put
\[
 E(q)=Q_{\mathbb O}(0;q_1,q_2,q_3),\qquad
 \mathfrak t(q_1,q_2,q_3)=\operatorname{Re}((q_1q_2)q_3),
 \qquad \mathfrak t_N(X)=\mathfrak t(\mathcal LX).
\]
The original ordered Albert determinant becomes the exact polynomial
\begin{equation}
 \det_{\mathbb O}\bigl(\lambda I+E(\mathcal LX)\bigr)
   =\lambda^3-\lambda(X,X)+2\mathfrak t_N(X).
 \label{n24:eq:niemeier-albert-cubic}
\end{equation}
Moreover
\begin{equation}
 \operatorname{tr}\bigl(E(\mathcal LX)\circ E(\mathcal LY)\bigr)
                  =2(X,Y).
 \label{n24:eq:albert-metric-two}
\end{equation}
For completeness, direct binary multiplication gives
$\operatorname{tr}E=0$, $\operatorname{tr}(E\circ E)=2\sum_i|q_i|^2$,
and $\operatorname{tr}(E\circ(E\circ E))=6\mathfrak t(q)$.
For the last equality, the three cyclic off-diagonal coordinates of
$E\circ E$ are $\overline{q_2q_3},\overline{q_3q_1},
\overline{q_1q_2}$; their trace pairings with the corresponding
$q_i$ each give $2\mathfrak t(q)$.
The Jordan trace formula
$\det X=\operatorname{tr}(X\circ(X\circ X))/3
-\operatorname{tr}X\operatorname{tr}(X\circ X)/2
+(\operatorname{tr}X)^3/6$
therefore gives the first identity after expanding $X=\lambda I+E$
and substituting the proved norm equality. For the second,
each off-diagonal octonion
occurs in both transposed matrix entries, giving exactly
\(2\sum_i\operatorname{Re}(q_i\overline{p_i})\).
Real associativity and cyclicity of octonion real parts give
\(\mathfrak t(q_3,q_1,q_2)=\mathfrak t(q_1,q_2,q_3)\):
\[
 \operatorname{Re}((q_3q_1)q_2)
 =\operatorname{Re}(q_3(q_1q_2))
 =\operatorname{Re}((q_1q_2)q_3).
\]
Consequently \(\mathfrak t_N(rX)=\mathfrak t_N(X)\), and the
whole determinant in \eqref{n24:eq:niemeier-albert-cubic} is
\(r\)-invariant with \(\lambda\) fixed.  Equivalently, for the real
permutation matrix \(P e_i=e_{i+1}\), indices modulo three,
\[
 P\bigl(\lambda I+E(q)\bigr)P^T
             =\lambda I+E(q_3,q_1,q_2).
\]
This is also a Jordan automorphism: multiplication by a real
permutation matrix merely reindexes the two-factor matrix products,
so it preserves \((AB+BA)/2\).  It therefore carries the positive
cone and its determinant boundary to themselves.  The cyclic bridge
is simultaneously linear, isometric on \(\mathbb O^3\), integral
on the stated lattice image, and compatible with the actual Albert
cubic and positive cone.

\paragraph{Dependence and the four-dimensional lattice summand.}
The map uses precisely the four specified origins \(f_i\), the
generator \(289\), the even/odd triplet order, the displayed order
of \(b\), the ordered \(D\)-coordinates, and the chosen orthonormal
octonion basis.  Changing one of these changes the displayed map;
the construction makes no claim of independence from these data.
Its restriction to the actual \(D\)-summand is
\begin{equation}
 d\longmapsto\frac1{\sqrt3}(z(d),z(d),z(d)),\qquad
 z(d)=d_1u_4+d_2u_5+d_3u_6+d_4u_7.
 \label{n24:eq:d-diagonal-embedding}
\end{equation}
This is an isometry on \(D\), and on \(D^\vee\) in its real span.
In particular its 24 shortest nonzero dual vectors become 24
specified points in this four-dimensional diagonal subspace.
The separate Cartan-weight map \(d\mapsto\sum_kd_k\varepsilon_k\)
is related to it by the exact isometry
\(\sum_kd_k\varepsilon_k\mapsto(z(d),z(d),z(d))/\sqrt3\).

\subsection{The Leech neighbour in the same three octonion coordinates}
\label{n24:sec:leech-coordinates}

The preceding map also transports the entire neighbour construction,
not only the real span of $N$. Put $N_{\mathbb O}=\mathcal L(N)$ and
retain the actual $\rho,\alpha,K,\Lambda$ of
Theorem~\ref{esnl:thm:leech-neighbor}. In the already fixed basis let
\[
 U=\sum_{a=0}^7u_a,\qquad
 V=\frac{u_0+u_1+u_2+u_3}{\sqrt2}
       +\sqrt3u_4+\frac{2u_5+u_6}{\sqrt3}.
\]
Substitution in the original isometry gives the explicit vectors
\begin{equation}
 \begin{aligned}
 \varrho=\mathcal L\rho&=(V+2U,V,V-2U),\\
 \delta=\mathcal L\alpha
 &=\left(\frac{(2+\sqrt2)u_0-2u_1}{3},
          \frac{(\sqrt2-1)u_0+u_1}{3},
          \frac{(\sqrt2-1)u_0+u_1}{3}\right).
 \end{aligned}
 \label{n24:eq:weyl-in-octonions}
\end{equation}
Indeed each even triplet of $\rho_A$ has mean $1/2$, each odd
triplet has mean $-1/2$, and both residual triplets are $(2,0,-2)$.
The four mean-difference coordinates are $\sqrt6/2$, while the
four $D$ coordinates are $(3,2,1,0)$. For $\alpha$, the two residual
triplets in the first fibre are $(2,-1,-1)/3$ and $(-2,1,1)/3$;
its only nonzero mean-difference coordinate is $2/\sqrt6$.
These facts prove every coefficient in
\eqref{n24:eq:weyl-in-octonions}. In particular
$|\varrho|^2=84$, $|\delta|^2=2$, and
$\langle\varrho,\delta\rangle=1$ in the product octonion metric.

Define
\begin{equation}
 \begin{aligned}
 K_{\mathbb O}&=\{q\in N_{\mathbb O}:
                          \langle\varrho,q\rangle\equiv0\pmod6\},\\
 \Lambda_{\mathrm{Leech},\mathbb O}
       &=K_{\mathbb O}+\mathbb Z(\varrho/6-\delta).
 \end{aligned}
 \label{n24:eq:leech-in-octonions}
\end{equation}
The isometry and its inverse prove, in both directions,
$\mathcal L(K)=K_{\mathbb O}$ and
$\mathcal L(\Lambda)=\Lambda_{\mathrm{Leech},\mathbb O}$.
Thus \eqref{n24:eq:leech-in-octonions} is a full coordinate
description of the Leech lattice inside the very same three sectors.
The glue test above, followed by its displayed height congruence,
is an exact membership test for $K_{\mathbb O}$; adjoining the six
cosets $k(\varrho/6-\delta)$, $0\le k<6$, gives all of the Leech
lattice. The common-sublattice indices, discriminant form, and
order-twelve quotient are transported without alteration.

The cyclic action also acts on the marked family of neighbours.
Retain $r=\rho(361)$ and its corresponding map
$c(q_1,q_2,q_3)=(q_3,q_1,q_2)$.
For any transported Weyl vector $\rho'=T\rho$, $T\in O(N)$,
the neighbour is defined using $T\alpha$, so that
$\langle\rho',T\alpha\rangle=1$. Then
\[
 c\,\mathcal L(K_{\rho'})=\mathcal L(K_{r\rho'}),\qquad
 c\,\mathcal L(\Lambda_{\rho'})=\mathcal L(\Lambda_{r\rho'}).
\]
This follows directly from $\mathcal Lr=c\mathcal L$ and
(NL18). The inverse maps use $r^{-1}$ and $c^{-1}$.
It specifies exactly which marked neighbour is the target; it does
not silently assume that the chosen $\rho$ is fixed by $r$.
Finally the common cubic evaluates on every one of these lattice
points as
\[
 \det Q_{\mathbb O}(\lambda;\mathcal LX)
    =\lambda^3-\lambda(X,X)
                   +2\operatorname{Re}(((\mathcal LX)_1
                                         (\mathcal LX)_2)(\mathcal LX)_3).
\]
This identity holds for all $X\in N_{\mathbb R}$, hence also
on the explicitly embedded $N$, $K$ and $\Lambda$.

\subsection{Exactly which root isometries preserve this glue}
\label{n24:sec:glue-symmetry}

Retain the coordinate convention \((P_px)_i=x_{p(i)}\).
An arbitrary root-lattice isometry has the original form
\[
 T(X_1,X_2,X_3,X_4;d)
   =\bigl((\varepsilon_iQ_{\sigma_i}X_{p(i)})_{i=1}^4;\ T_Dd\bigr),
\]
where \(\varepsilon_i\in\{\pm1\}\), \(\sigma_i\in S_6\),
\(p\in S_4\), and \(T_D\in O(D)\).  If \(B=d_D(T_D)\),
the necessary and sufficient condition for \(T(N)=N\) is
\begin{equation}
 B\tau(u)=\tau(P_pu)\ (u\in E),\qquad
 \bigl(\varepsilon_i\ell_{p(i)}\bigr)_{i=1}^4
                  =\bigl(\ell_iA\bigr)_{i=1}^4
 \quad\hbox{for some }A\in GL_2(\mathbb F_3),
 \label{n24:eq:all-glue-preservers}
\end{equation}
where
\(\ell_1=(1,1),\ell_2=(1,2),\ell_3=(1,0),\ell_4=(0,1)\).
The \(A\), and hence \(B=B_p\), are unique.  Indeed the
two-primary part of the glue is the graph of \(\tau\), and signs
are trivial modulo two.  Its three-primary part is the image of the
injective map \(j(a,b)=(a+b,a+2b,a,b)\), so preserving it is
equivalent to \(M_\varepsilon P_pj=jA\).  These are exactly
the two conditions in \eqref{n24:eq:all-glue-preservers}; the
primary decomposition proves their sufficiency as well as necessity.
Since the four \(\ell_i\) represent the four distinct projective
row lines over \(\mathbb F_3\), each \(A\) gives exactly one
signed permutation by this equation.  This recovers the precise
48-element glue quotient, not merely its action on three labels.

The canonical full-automorphism theorem proves, with these same
coordinates and section, that
\[
 O(N)=\bigl(S_6^4\times W(D_4)\bigr)\rtimes L(G),
 \qquad G\cong GL_2(\mathbb F_3).
\]
Conjugation by the actual map \(\mathcal L\) now gives the entire
orthogonal stabilizer of the embedded lattice, with explicit matrices
obtained from \eqref{n24:eq:full-isometry} and its inverse:
\begin{equation}
 O(\Lambda_{\mathbb O})
       =\mathcal L O(N)\mathcal L^{-1},\qquad
 \Gamma_G(g)=\mathcal LL_g\mathcal L^{-1},\qquad
 \mathcal LL_g=\Gamma_G(g)\mathcal L.
 \label{n24:eq:full-transported-action}
\end{equation}
Both inclusions in the first equality follow by conjugating an
isometry preserving the specified lattice.  This is a faithful
action on the whole 24-dimensional space. Each \(\Gamma_G(g)\)
commutes with \(c\),
because the canonical equation
\(L_g\rho(289^2)=\rho(289^2)L_g\) conjugates to that equation.

For completeness, the previously displayed \(D_4\)-triality
section and geometric frame permutations have a precise comparison
even though that particular proposed linear conjugacy fails.  On
\(V_N\), each of \(\widehat R,\widehat H\) has trace
\(5-5+2=2\): the transposed \(A_5\) pair contributes zero,
the other two components contribute \(5,-5\), and the
\(D_4\) reflection contributes 2.  Their three-cycle has trace
\(5+1=6\).  A geometric Albert-frame transposition acts on
\(\mathbb O^3\) by swapping two sectors with conjugation and by
conjugating the fixed sector; its trace is
\(1-7=-6\).  Its three-cycle has trace zero.  Therefore an
invertible map cannot conjugate these two particular \(S_3\)
representations, since trace is invariant under conjugation.
The complete real character decompositions are, respectively,
\[
 V_N|_{\langle\widehat R,\widehat H\rangle}
       \cong7\mathbf1\oplus5\operatorname{sgn}\oplus6\operatorname{Std},
 \qquad
 \mathbb O^3|_{S_3^{\rm frame}}
       \cong\mathbf1\oplus7\operatorname{sgn}\oplus8\operatorname{Std}.
\]
These multiplicities follow by taking inner products with the
characters \((1,1,1),(1,-1,1),(2,0,-1)\), on classes of sizes
\(1,3,2\).  The exact whole-space relation for the original section
is \eqref{n24:eq:full-transported-action}; the exact
determinant-compatible cyclic relation is
\eqref{n24:eq:cyclic-intertwining}.  In the latter relation the
residue element has trivial discriminant action while its octonion
action cycles the sectors.  Those are two explicitly calculated
actions of the same operator on two explicitly different constructions;
no identification between their quotient homomorphisms is assumed.

Finally, the connected group \(\operatorname{Spin}(8)\) acts on
the ambient triality module, and its subgroup preserving
\(\Lambda_{\mathbb O}\) is exactly its intersection with
\eqref{n24:eq:full-transported-action}.  It is finite: the columns
of a lattice isometry in a fixed lattice basis have prescribed
bounded lengths, so each column has finitely many possible lattice
values.  A connected group preserving a discrete full lattice
pointwise as a set acts trivially, since each orbit of a lattice
vector is connected and discrete and the lattice spans.  The
nontrivial connected triality action therefore does not preserve
the whole lattice.  The proved cyclic subgroup \(\langle c\rangle\)
does preserve both the lattice and its Albert cubic.  Equations
\eqref{n24:eq:all-glue-preservers} and
\eqref{n24:eq:full-transported-action} specify the full glue
stabilizer; they do not assert that every element of that finite
orthogonal group also preserves the cubic.

% Proof-bearing staged module. Namespace nca:. No canonical edits.
\section{The arithmetic of the retained Albert cubic on both lattices}
\label{nca:sec:cubic-arithmetic}

\subsection{Original data and the exact polynomial}

Retain the ordered coordinates \(X=(x_{i,j};d_k)\),
\(1\leq i\leq4\), \(0\leq j\leq5\), \(1\leq k\leq4\), with
\(\sum_{j=0}^5x_{i,j}=0\), and their original Euclidean metric.
Let
\[
 A_5=\{x\in\mathbb Z^6:\sum_jx_j=0\},\qquad
 D_4=\{d\in\mathbb Z^4:\sum_kd_k\equiv0\pmod2\},\qquad
 R=A_5^{\perp4}\perp D_4.
\]
Write \(\lambda=(5,-1,-1,-1,-1,-1)/6\),
\(v_D=(1,0,0,0)\), and \(s_D=(1,1,1,1)/2\).
The original discriminant glue is generated by
\[
 \begin{split}
 &(3,3,3,3;00),\quad(3,3,0,0;10),\quad(3,0,3,0;01),\\
 &(2,2,2,0;00),\quad(2,4,0,2;00)
 \end{split}
\]
in \((\mathbb Z/6)^4\oplus\mathbb F_2^2\), with the same marking
\(10=v_D+D_4\), \(01=s_D+D_4\). Its order is 72: the first
three generators have independent two-primary reductions, and the
last two independent three-primary reductions.  The lattice \(N\)
is its inverse image in \(R^\vee\).

Retain
\[
 \rho_A=(5,3,1,-1,-3,-5)/2,\qquad
 \rho_D=(3,2,1,0),\qquad \rho=\rho_A^{\oplus4}\perp\rho_D,
 \qquad \alpha=(e_0-e_1,0,0,0;0),
\]
and the exact neighbour
\[
 K=\{X\in N:(\rho,X)\equiv0\pmod6\},\qquad
 v=\rho/6-\alpha,\qquad \Lambda=K+\mathbb Zv.
\]
Here \(\rho\in N\), \(\rho^2=84\), and \((\rho,\alpha)=1\),
as follows immediately from these coordinates and the first glue
generator.  The name \(\Lambda\) always denotes this neighbour.

Fix an ordered orthonormal octonion basis \(u_0=1,u_1,\ldots,u_7\).
For the first part of the calculation no multiplication table is
imposed beyond the actual multiplication of these basis elements.
Define, in the original coordinates,
\begin{align}
 s_i&=\sum_{j=0}^5(-1)^jx_{i,j},\label{nca:eq:original-s}\\
 A_{2(i-1)+p}
  &=x_{i,p}-\frac{x_{i,p}+x_{i,p+2}+x_{i,p+4}}3,
       &&p=0,1,\label{nca:eq:original-A}\\
 B_{2(i-1)+p}
  &=x_{i,p+2}-\frac{x_{i,p}+x_{i,p+2}+x_{i,p+4}}3,
       &&p=0,1.\label{nca:eq:original-B}
\end{align}
Write \(A=A_0+\mathbf a\), \(B=B_0+\mathbf b\), where
\(\mathbf a=\sum_{h=1}^7A_hu_h\) and
\(\mathbf b=\sum_{h=1}^7B_hu_h\).  The common mean octonion is
\begin{equation}
 M=\frac{\sqrt2}{6}\sum_{a=0}^3s_{a+1}u_a
       +\frac{\sqrt3}{3}\sum_{k=1}^4d_ku_{k+3}.
 \label{nca:eq:mean}
\end{equation}
The retained isometry therefore has precisely the ordered form
\[
 \mathcal LX=(A+M,B+M,-A-B+M).
\]
This is an identity with its original three triplets: the residuals
sum to zero, the first four common mean coefficients are
\(s_i/\sqrt{18}=\sqrt2s_i/6\), and the last four are
\(d_k/\sqrt3=\sqrt3d_k/3\).
For completeness the original squared norm and the output norm are both
\[
 |A|^2+|B|^2+|A+B|^2+\frac16\sum_i s_i^2+\sum_kd_k^2.
\]
In the original fibres this follows by separating the means
\(s_i/6,-s_i/6\) from their zero-sum triplets; in the output it
follows from \(A+B+(-A-B)=0\) and
\(3|M|^2=\sum_i s_i^2/6+\sum_kd_k^2\). The inverse recovers
\(M=(q_1+q_2+q_3)/3\), then \(s_i=3\sqrt2M_{i-1}\),
\(d_k=\sqrt3M_{k+3}\), and \(A=q_1-M,B=q_2-M\).
Adding \((-1)^ps_i/6\) to the \(p\)-th residual triplet in the
\(i\)-th fibre recovers each original \(x_{i,p+2j}\). Thus the
map is onto and preserves the original pairing by polarization.

For imaginary octonions put
\[
 \mathbf a\times\mathbf b=\frac{\mathbf a\mathbf b-
                                    \mathbf b\mathbf a}{2},\qquad
 \phi(\mathbf a,\mathbf b,\mathbf c)
                     =\langle\mathbf a\times\mathbf b,\mathbf c\rangle.
\]
Let \(\phi_{ijk}=\phi(u_i,u_j,u_k)\).  Thus \(\phi\) is the
actual alternating octonion three-form in the chosen frame, including
all signs. Define
\begin{align}
 D&=A_0^2+A_0B_0+B_0^2-\langle\mathbf a,\mathbf a\rangle
             -\langle\mathbf a,\mathbf b\rangle
             -\langle\mathbf b,\mathbf b\rangle,\label{nca:eq:D}\\
 Z&=(2A_0+B_0)\mathbf a+(A_0+2B_0)\mathbf b
                               -3\mathbf a\times\mathbf b
       =\sum_{h=1}^7Z_hu_h.\label{nca:eq:Z}
\end{align}
Equivalently, the component in the final formula is the exact
quadratic polynomial
\[
 Z_h=(2A_0+B_0)A_h+(A_0+2B_0)B_h
                  -3\sum_{i,j=1}^7\phi_{ijh}A_iB_j.
\]

\begin{theorem}[Exact pullback with every radical and coefficient]
\label{nca:thm:polynomial}
The original cubic
\(F(X)=\operatorname{Re}(((\mathcal LX)_1(\mathcal LX)_2)
                                           (\mathcal LX)_3)\)
has the complete polynomial expression
\begin{equation}
 F(X)=P_0(X)+\sqrt2P_2(X)+\sqrt3P_3(X),
 \label{nca:eq:polynomial}
\end{equation}
where every variable is the displayed linear function of the original
\(x_{i,j},d_k\), and
\begin{align}
 P_0={}&-A_0B_0(A_0+B_0)+B_0\langle\mathbf a,\mathbf a\rangle
              +A_0\langle\mathbf b,\mathbf b\rangle
              +2(A_0+B_0)\langle\mathbf a,\mathbf b\rangle,
                \label{nca:eq:P0}\\
 P_2={}&\frac{-s_1D+s_2Z_1+s_3Z_2+s_4Z_3}{6}
       +\frac{s_1(4s_1^2-3\sum_{i=1}^4s_i^2)}{108}
       -\frac{s_1\sum_{k=1}^4d_k^2}{6},\label{nca:eq:P2}\\
 P_3={}&\frac{d_1Z_4+d_2Z_5+d_3Z_6+d_4Z_7}{3}.
                  \label{nca:eq:P3}
\end{align}
These identities hold throughout the original real vector space, in
particular on \(N\), \(K\), and \(\Lambda\). They retain the ordered
Albert determinant as
\[
 \det Q_{\mathbb O}(\lambda;\mathcal LX)
                =\lambda^3-\lambda(X,X)+2F(X).
\]
\end{theorem}

\begin{proof}
For \(a=a_0+\mathbf a\), \(b=b_0+\mathbf b\),
\(c=c_0+\mathbf c\), binary octonion multiplication gives
\[
 \operatorname{Re}((ab)c)
  =a_0b_0c_0-a_0\langle\mathbf b,\mathbf c\rangle
   -b_0\langle\mathbf a,\mathbf c\rangle
   -c_0\langle\mathbf a,\mathbf b\rangle
   -\phi(\mathbf a,\mathbf b,\mathbf c).
\]
Call the first four, symmetric, terms \(S(a,b,c)\).  Substitute
\((a,b,c)=(A+M,B+M,-A-B+M)\).  Terms with exactly two copies
of \(M\) in \(S\) sum to \(S(M,M,A+B-A-B)=0\). Terms
with exactly one copy sum to
\(-S(M,A,A)-S(M,A,B)-S(M,B,B)\).  The term with no copy is
\(S(A,B,-A-B)=P_0\), by multiplying its four displayed terms.
The alternating term with no copy vanishes; terms with two or
three copies vanish. Its three terms with one copy all equal
\(\phi(\mathbf a,\mathbf b,\operatorname{Im}M)\), since
\[
 \phi(\mathbf a,\operatorname{Im}M,-\mathbf a-\mathbf b)
 =\phi(\mathbf a,\mathbf b,\operatorname{Im}M),\quad
 \phi(\operatorname{Im}M,\mathbf b,-\mathbf a-\mathbf b)
 =\phi(\mathbf a,\mathbf b,\operatorname{Im}M).
\]
Thus all one-copy terms together are
\(-M_0D+\langle\operatorname{Im}M,Z\rangle\).
Finally the term with three copies is
\[
 \operatorname{Re}(M^3)=4M_0^3-3M_0|M|^2,
 \qquad |M|^2=\frac{\sum_i s_i^2}{18}
                              +\frac{\sum_kd_k^2}{3},
 \qquad M_0=\frac{\sqrt2s_1}{6}.
\]
Substitution proves \eqref{nca:eq:P0}--\eqref{nca:eq:P3} with
their displayed factors. In particular it proves the absence of a
\(\sqrt6\) term without discarding any mixed term. The determinant
formula uses the original isometry and its unchanged squared norm.
\end{proof}

\subsection{Finite generators of both original lattices}

Order the 24 original simple roots as
\[
 g_{5(i-1)+j+1}=a_{i,j}=e_{i,j}-e_{i,j+1}
 \quad(1\leq i\leq4,\ 0\leq j\leq4),
\]
followed by
\[
 g_{21}=e_1-e_2,\quad g_{22}=e_2-e_3,\quad
 g_{23}=e_3-e_4,\quad g_{24}=e_3+e_4
\]
in the original \(D_4\) coordinates.  All unspecified components
are zero.  Adjoin the five actual glue lifts
\begin{align*}
 g_{25}&=(3\lambda,3\lambda,3\lambda,3\lambda;0),&
 g_{26}&=(3\lambda,3\lambda,0,0;v_D),\\
 g_{27}&=(3\lambda,0,3\lambda,0;s_D),&
 g_{28}&=(2\lambda,2\lambda,2\lambda,0;0),\\
 g_{29}&=(2\lambda,4\lambda,0,2\lambda;0).
\end{align*}
Then
\begin{equation}
 N=\sum_{i=1}^{29}\mathbb Zg_i,\qquad
 H_i=(\rho,g_i)=
 \begin{cases}
 1,&1\leq i\leq24,\\
 30,18,18,15,20,&i=25,26,27,28,29\text{ respectively}.
 \end{cases}
 \label{nca:eq:N-generators}
\end{equation}
Indeed the first 24 vectors span \(R\), the last five generate
the original glue, and subtraction of the corresponding glue lift
leaves every element of \(N\) in \(R\).  The height calculation
uses \((\rho_A,\lambda)=5/2\),
\((\rho_D,v_D)=(\rho_D,s_D)=3\), and the simple-root heights one.

Define the following 30, fully specified neighbour generators:
\begin{equation}
 h_1=6\alpha,\qquad h_i=g_i-H_i\alpha\quad(2\leq i\leq29),
 \qquad h_{30}=v=\rho/6-\alpha.
 \label{nca:eq:Lambda-generators}
\end{equation}
Then \(K=\sum_{i=1}^{29}\mathbb Zh_i\) and
\(\Lambda=\sum_{i=1}^{30}\mathbb Zh_i\).
For if \(X=\sum_i n_ig_i\in K\), then
\[
 X=\sum_{i=2}^{29}n_i(g_i-H_i\alpha)
                         +\left(\sum_{i=1}^{29}n_iH_i\right)\alpha,
\]
and the last coefficient is divisible by six. Conversely every
displayed generator before \(h_{30}\) has height divisible by six.
Adjoining \(h_{30}\) is exactly the original neighbour definition.

\subsection{A complete finite calculation of the cubic value modules}

For a lattice \(L\) and a homogeneous real cubic \(F\), define
its additive value module by
\[
 \mathcal V_F(L)=\sum_{X\in L}\mathbb ZF(X)\subset\mathbb R.
\]
This definition retains every value and takes its generated additive
group. It does not claim that every element of the group is one
individual value of \(F\).

\begin{theorem}[Finite cubic-value generators, including negative integers]
\label{nca:thm:finite-values}
If \(L=\sum_{i=1}^m\mathbb Zb_i\), allowing dependent generators,
then \(\mathcal V_F(L)\) is generated by the following finite list:
\begin{align}
 A_i={}&F(b_i),&&1\leq i\leq m,\label{nca:eq:single}\\
 D^+_{ij}={}&F(b_i+b_j)-F(b_i)-F(b_j),&&i<j,\label{nca:eq:Dplus}\\
 D^-_{ij}={}&F(b_i-b_j)-F(b_i)+F(b_j),&&i<j,\label{nca:eq:Dminus}\\
 D_{ijk}={}&F(b_i+b_j+b_k)-F(b_i+b_j)-F(b_i+b_k)\notag\\
          &{}-F(b_j+b_k)+F(b_i)+F(b_j)+F(b_k),&&i<j<k.
 \label{nca:eq:Dthree}
\end{align}
Consequently \eqref{nca:eq:N-generators} gives exactly 4495
specified generators of \(\mathcal V_F(N)\), and
\eqref{nca:eq:Lambda-generators} gives exactly 4960 specified
generators of \(\mathcal V_F(\Lambda)\), for every fixed
orthonormal octonion frame and its actual \(\phi_{ijk}\).
\end{theorem}

\begin{proof}
Let \(B_F\) be the symmetric trilinear polarization, with
\(B_F(X,X,X)=F(X)\); explicitly its value is
\(1/6\) times the seven-term expression in
\eqref{nca:eq:Dthree}, with its three inputs replacing the three
generators. For every integer tuple \(n\), expansion gives
\[
 F\left(\sum_i n_ib_i\right)
 =\sum_i n_i^3A_i
 +\sum_{i<j}\bigl(n_i^2n_j C_{ij}+n_in_j^2 E_{ij}\bigr)
 +\sum_{i<j<k}n_in_jn_kD_{ijk},
\]
where \(C_{ij}=3B_F(b_i,b_i,b_j)\) and
\(E_{ij}=3B_F(b_i,b_j,b_j)\).  The two pair generators are
\(D^+_{ij}=C_{ij}+E_{ij}\) and
\(D^-_{ij}=-C_{ij}+E_{ij}\). Therefore their contribution is
\[
 \frac{n_in_j(n_i+n_j)}2D^+_{ij}
       +\frac{n_in_j(n_j-n_i)}2D^-_{ij}.
\]
Both coefficients are integers, for all positive, zero, and negative
integers \(n_i,n_j\): if either integer is even the claim is immediate,
and if both are odd their sum and difference are even. This proves
that the finite list generates all values. Conversely each entry
of the list is an integer linear combination of actual values on
\(L\). The proof does not use linear independence of the generators.
Finally \(29+2\binom{29}{2}+\binom{29}{3}=4495\) and
\(30+2\binom{30}{2}+\binom{30}{3}=4960\).
\end{proof}

For numerical arithmetic now specify the additional octonion frame
choice completely.  Use \(\mathbb O=\mathbb H\oplus\mathbb H\ell\)
in ordered coordinates \((a,b)\), with
\begin{equation}
 \begin{gathered}
 (a,b)(c,d)=(ac-\overline d\,b,\ da+b\overline c),\\
 u_0=(1,0),\quad u_1=(i,0),\quad u_2=(j,0),\quad u_3=(k,0),\\
 u_4=(0,1),\quad u_5=(0,i),\quad u_6=(0,j),\quad u_7=(0,k).
 \end{gathered}
 \label{nca:eq:Cayley-frame}
\end{equation}
Here \(ij=k\), \(jk=i\), \(ki=j\), and the reversed products
are negative.  The positive oriented three-form triples are exactly
\begin{equation}
 (123),\ (145),\ (176),\ (246),\ (257),\ (347),\ (365).
 \label{nca:eq:seven-triples}
\end{equation}
Every \(\phi_{ijk}\) is obtained from these by alternation, or is
zero.  This is a specified frame for the retained, previously
frame-dependent isometry, without a change to any lattice coordinates
or to the Albert matrix convention.

\begin{theorem}[Exact value modules and their quotient of order twelve]
\label{nca:thm:value-modules}
For the frame \eqref{nca:eq:Cayley-frame}, the exact additive modules are
\begin{align}
 \mathcal V_F(N)
    &=\frac1{27}\mathbb Z\ \oplus\
       \frac{\sqrt2}{54}\mathbb Z\ \oplus\
       \frac{\sqrt3}{18}\mathbb Z,\label{nca:eq:N-values}\\
 \mathcal V_F(\Lambda)
    &=\frac1{27}\mathbb Z\ \oplus\
       \frac{\sqrt2}{216}\mathbb Z\ \oplus\
       \frac{\sqrt3}{54}\mathbb Z.\label{nca:eq:Lambda-values}
\end{align}
In particular the inclusion of these two additive subgroups of
\(\mathbb R\) has the exact quotient
\begin{equation}
 \mathcal V_F(\Lambda)/\mathcal V_F(N)
                  \cong\mathbb Z/4\oplus\mathbb Z/3
                  \cong\mathbb Z/12.
 \label{nca:eq:value-quotient}
\end{equation}
This is an inclusion of cubic value modules; it does not assert the
false lattice inclusion \(N\subset\Lambda\).
\end{theorem}

\begin{proof}
The polynomial coefficients in
\eqref{nca:eq:P0}--\eqref{nca:eq:P3} are rational on both lattices
for this exact multiplication table.  The numbers
\(1,\sqrt2,\sqrt3\) are linearly independent over \(\mathbb Q\).
For example, isolating one radical in a rational relation and
squaring either gives a rational multiple of another irrational
radical, or a rational square equal to two, three, or \(3/2\),
which is impossible by parity of prime exponents. Thus a value has
a unique rational coefficient triple \((P_0,P_2,P_3)\).

Apply the complete finite list in
Theorem~\ref{nca:thm:finite-values} to the displayed original
generators. Exact substitution gives the following least common
multiples of coefficient denominators. The denominators in each
entry are ordered as \((P_0,P_2,P_3)\), and a zero coefficient has
denominator one:
\begin{equation}
\begin{array}{c|c|c|c|c}
 &A_i&D^+_{ij}&D^-_{ij}&D_{ijk}\\\hline
 N&(27,27,1)&(9,18,18)&(9,18,18)&(9,18,18)\\
 \Lambda&(27,216,27)&(27,108,54)&(27,108,54)&(9,36,54)
\end{array}.
\label{nca:eq:denominator-certificate}
\end{equation}
This table is a finite rational calculation, with no bound on
lattice points or heuristic sampling.  The complete arithmetic
certificate is the accompanying \texttt{verify\_cubic.py}: its
\texttt{residual\_data} and \texttt{cubic} functions are exactly
\eqref{nca:eq:original-s}--\eqref{nca:eq:P3}; it forms every
\(A_i,D^+_{ij},D^-_{ij},D_{ijk}\) for the 29 or 30 explicitly
defined vectors, using rational arithmetic.  The \texttt{cross}
function executes the three cyclic positive and three negative
entries for each of the seven triples
\eqref{nca:eq:seven-triples}.  Hence the full finite calculation
is specified by these equations, the finite index bounds in
Theorem~\ref{nca:thm:finite-values}, and integer arithmetic.
Its 4495 and 4960 entries verify the upper inclusions in
\eqref{nca:eq:N-values}--\eqref{nca:eq:Lambda-values}.

Here are short certificates for the reverse inclusions. Coordinates
in the next two tables are relative to the respective three
proposed generators on the right sides of
\eqref{nca:eq:N-values} and \eqref{nca:eq:Lambda-values}.
In the first table all pair symbols use the \(g_i\):
\begin{equation}
\begin{array}{c|r r r}
 F(g_1)&-4&4&0\\
 F(g_2)&-5&-4&0\\
 F(g_7)&0&18&0\\
 F(g_{28})&-80&52&0\\
 D^+_{21,28}&0&-36&16\\
 D^+_{2,27}&-27&27&3\\
 D^+_{2,25}&-81&45&0
\end{array}.
\label{nca:eq:N-lower-certificate}
\end{equation}
The determinants of the three columns obtained by transposing
rows \((1,4,5)\), \((2,3,6)\), and \((2,7,6)\) are,
respectively, \(1792,-270,-1647\). Their greatest common divisor
is one. In the second table all pair symbols use the \(h_i\):
\begin{equation}
\begin{array}{c|r r r}
 F(h_2)&14&-272&0\\
 F(h_4)&5&-128&0\\
 F(h_{28})&13780&-74672&0\\
 F(h_{30})&14&-63&-2\\
 D^+_{2,27}&7425&-60624&-9
\end{array}.
\label{nca:eq:Lambda-lower-certificate}
\end{equation}
Rows \((1,3,4)\), \((1,4,5)\), and \((2,4,5)\) give
determinants \(-5405504,2315394,1281267\), whose greatest
common divisor is also one.

For clarity, the elementary integer argument using these
determinants does not require a normal-form theorem.  Let
\(W\subset\mathbb Z^3\) be the group generated by one displayed
table.  For any three of its columns forming a matrix \(C\),
the adjugate identity \(C\operatorname{adj}(C)=\det(C)I\)
implies \(\det(C)\mathbb Z^3\subset W\).  Bezout's identity
for the three displayed determinants, whose gcd is one, therefore
gives \(\mathbb Z^3\subset W\). The opposite inclusion is
immediate from the integer entries. Every table entry is an
integer combination of actual values by its definition, proving
the two reverse inclusions.  Finally the coordinates of the
inclusion of the first module in the second are
\((a,b,c)\mapsto(a,4b,3c)\). Its cokernel is the displayed
product of cyclic groups, with order twelve. This also proves
the asserted exact quotient map.
\end{proof}

\subsection{Frame dependence and exact symmetry transport}

The finite-list theorem is valid for the actual coefficients of
every orthonormal frame. The numerical modules
\eqref{nca:eq:N-values}--\eqref{nca:eq:Lambda-values} refer to
the frame \eqref{nca:eq:Cayley-frame}.  This dependence has an
exact comparison map. If \(u_a'=Tu_a\) for a real isometry
\(T\) fixing \(1\), then
\[
 \mathcal L_{u'}=(T,T,T)\mathcal L_u.
\]
The symmetric part \(S\) of the trilinear tensor is unchanged,
because \(T\) preserves the scalar coordinate and inner product.
Writing \(\mathbf m=\operatorname{Im}M\), its exact change is
\begin{equation}
 F_{u'}(X)-F_u(X)
  =-3\bigl(\phi(T\mathbf a,T\mathbf b,T\mathbf m)
                  -\phi(\mathbf a,\mathbf b,\mathbf m)\bigr).
 \label{nca:eq:frame-transport}
\end{equation}
This follows from the complete expansion in the polynomial proof.
If \(T\) is an octonion automorphism it preserves \(\phi\), so
the cubic and its value modules are unchanged.

The dependence on a general orthonormal frame is substantive. Take
the actual lattice vector with the only nonzero fibre coordinates
\[
 x_{1,1}=1,\ x_{1,5}=-1,\quad
 x_{2,2}=1,\ x_{2,4}=-1,\qquad d=(1,1,0,0).
\]
It lies in \(R\subset N\), has \(A=u_1\), \(B=u_2\), and
\(M=(u_4+u_5)/\sqrt3\). In the Cayley frame its value is zero.
Replace only
\[
 u'_3=(2u_3-\sqrt5u_4)/3,\qquad
 u'_4=(\sqrt5u_3+2u_4)/3.
\]
This is an exact orthonormal frame change; the same original
lattice vector now has value \(-\sqrt{15}/3\), because
\(\phi(u_1,u_2,u'_4)=\sqrt5/3\) and
\(\phi(u_1,u_2,u_5)=0\). This value is outside
\(\mathbb Q+\sqrt2\mathbb Q+\sqrt3\mathbb Q\). Indeed the
independent sign changes of \(\sqrt2,\sqrt3\) are automorphisms
of \(\mathbb Q(\sqrt2,\sqrt3)\). If an element \(z\) in that
field had \(z^2=15\), each sign change would send it to
\(z\) or \(-z\). In the basis
\(1,\sqrt2,\sqrt3,\sqrt6\) this makes \(z\) a rational
multiple of a single basis vector. Its square would force one
of \(15,15/2,5,5/2\) to be a rational square, impossible by
parity of its prime exponents.
Thus the numerical specialization is not silently extended to
an unspecified orthonormal frame.

Return to the explicit Cayley frame for the following numerical
witnesses. Let \(\sigma=\rho(289)\) act simultaneously by
\(e_{i,j}\mapsto e_{i,j+1\bmod6}\), and trivially on \(D_4\).
The exact isometry gives
\(\mathcal L\sigma^2=c\mathcal L\) for
\(c(q_1,q_2,q_3)=(q_3,q_1,q_2)\). Cyclic real associativity
therefore proves \(F(\sigma^{2k}X)=F(X)\) for every real \(X\).
These maps preserve \(N\) by their original discriminant action.
The stabilizer of this cubic within the six-element residue group
is exactly its three-element subgroup: direct substitution at
\(\alpha=g_1\) gives
\begin{equation}
 F(\alpha)=\frac{-4+2\sqrt2}{27},\qquad
 F(\sigma\alpha)=\frac{-5-2\sqrt2}{27},
 \label{nca:eq:residue-nonpreserver}
\end{equation}
which differ. All odd powers have the same effect on the cubic
as \(\sigma\), by its proved invariance under even powers.

The original reflection exchanging \(e_{1,0},e_{1,1}\) sends
\(\alpha\) to \(-\alpha\), hence changes its nonzero cubic
value to its negative. The original full glue-preserving lifts
\[
 \begin{aligned}
 \widehat R(X_1,X_2,X_3,X_4;d)
       &=(X_2,X_1,X_3,-X_4;R_Dd),\\
 \widehat H(X_1,X_2,X_3,X_4;d)
       &=(X_1,X_3,X_2,-X_4;H_Dd).
 \end{aligned}
\]
also have explicit witnesses:
\[
 F(\widehat R\alpha)=0,\qquad
 F(g_7)=\frac{\sqrt2}{3},\qquad F(\widehat Hg_7)=0.
\]
Their original matrices remain
\[
 R_D=\operatorname{diag}(-1,1,1,1),\qquad
 H_D=\frac12\begin{pmatrix}1&1&1&1\\1&1&-1&-1\\
                           1&-1&1&-1\\1&-1&-1&1\end{pmatrix}.
\]
These maps preserve the original lattice, as can be checked on its
displayed generators.  Both \(R_D\) and \(H_D\) preserve \(D_4\):
the former changes one sign, while the latter has integral coordinates
on every even-sum integer vector and has output sum twice its first
input coordinate. They are orthogonal involutions. Their marked
discriminant actions fix \(v_D\) and exchange \(s_D,c_D\), and
exchange \(v_D,s_D\) while fixing \(c_D\), respectively.
On the binary glue generators, \(\widehat R\) fixes the first
two and sends the third to their second-plus-third sum;
\(\widehat H\) fixes the first and exchanges the second and third.
On the ternary code generators \(z_1=(1,1,1,0)\),
\(z_2=(1,2,0,1)\), the two signed component permutations act as
\((z_1,z_2)\mapsto(z_1,2z_2)\) and
\((z_1,z_2)\mapsto(z_1,2z_1+2z_2)\), respectively.
Thus both preserve the full original glue and its inverse image.
Every within-fibre coordinate permutation also preserves that inverse
image, since its action on \(\lambda+A_5\) is trivial:
\(Q\lambda-\lambda=Qe_0-e_0\in A_5\).
These are concrete failures to preserve this cubic, without
discarding their actual lattice actions or glue.

Negation preserves each lattice \(N\) and \(\Lambda\), and its
exact polynomial law is \(F(-X)=-F(X)\). In the neighbour the
nonzero witness is its displayed norm-four generator:
\begin{equation}
 F(\rho)=15\sqrt2+16\sqrt3,\qquad
 F(v)=\frac{14}{27}-\frac{7\sqrt2}{24}-\frac{\sqrt3}{27}.
 \label{nca:eq:special-values}
\end{equation}
Both equalities follow by substituting the original coordinates
of \(\rho\) and \(v\) into the proved polynomial.

Every original lattice isometry \(T\in O(N)\), whether or not
it preserves the cubic, has the exact replacement transport
\begin{equation}
 \Gamma_T=\mathcal LT\mathcal L^{-1},\qquad
 F^T(Y)=F(T^{-1}Y),\qquad F^T(TX)=F(X),
 \label{nca:eq:all-transport}
\end{equation}
and
\[
 T\Lambda_\rho=\Lambda_{T\rho},\qquad
 \mathcal V_{F^T}(\Lambda_{T\rho})=\mathcal V_F(\Lambda_\rho),
 \qquad \mathcal V_{F^T}(N)=\mathcal V_F(N).
\]
For the neighbour transport, write its exact elements as
\(x+k\rho/6\), where \(x\in N\) and
\((\rho,x)+k\equiv0\pmod6\); application of \(T\) gives
the identical congruence with \((T\rho,Tx)\). The inverse is
\(T^{-1}\). The value-module equalities then follow by the
bijections and the displayed identity \(F^T(TX)=F(X)\).
In particular for \(T=\sigma^{2k}\), one has \(F^T=F\), so
the same retained cubic has exactly the same value module on
every neighbour \(\Lambda_{\sigma^{2k}\rho}\) in this cyclic
orbit. This statement names its target neighbour and does not
assume the chosen \(\rho\) is fixed.

% Exact follow-on module. Input after the original Niemeier and Leech proofs.
% Original Euclidean coordinates, residue action, scales and glue are retained.
\section{The residue cubic action and the common integral Leech correspondence}
\label{esrf:sec:fixed-neighbor}

Retain the exact lattices $R=A^{\perp4}\perp D\subset N$ and the
binary and ternary codes from the index-72 construction.
The four $A$ coordinates remain the ordered six-element residue
fibres. To avoid confusing a vector with a representation, write
$\varrho$ for the Weyl vector
\[
 \varrho_A=(5,3,1,-1,-3,-5)/2,\qquad
 \varrho_D=(3,2,1,0),\qquad
 \varrho=\varrho_A^{\oplus4}\perp\varrho_D.
\]
Thus $\varrho\in N$ and $\varrho^2=84$.
The original residue representation is still denoted by $\rho$.
The specific element studied here is
\[
 r=\rho(361)=\rho(289)^2,\qquad
 r|_{A_i}=C^2,\quad Ce_j=e_{j+1\bmod6},\qquad r|_D=I.
 \tag{RF1}
\]
In particular $r^3=I$. This is the actual residue coordinate
permutation, with no conjugation or replacement.

For each $\sigma\in\Omega=O(N)\varrho$, choose an original root
$\alpha_\sigma$ with $(\sigma,\alpha_\sigma)=1$ and set
\[
 \begin{split}
 \ell_\sigma(x)&=(\sigma,x)\bmod6,\qquad
 K_\sigma=\ker(\ell_\sigma:N\to\mathbb Z/6),\\
 \Lambda_\sigma&=K_\sigma+
                      \mathbb Z(\sigma/6-\alpha_\sigma)\\
 &=\{x+k\sigma/6:x\in N,\ k\in\mathbb Z,\
                                    (\sigma,x)+k\equiv0\pmod6\}.
 \end{split}
 \tag{RF2}
\]
The preceding neighbor proof shows that the definition is independent
of $\alpha_\sigma$, every $\Lambda_\sigma$ is the specified rootless
even unimodular lattice, and
$N\cap\Lambda_\sigma=K_\sigma$ with both indices six.
The family considered in this section is precisely
$\mathscr F=\{\Lambda_\sigma:\sigma\in\Omega\}$.

\subsection{A necessary and sufficient equality criterion}

\begin{theorem}[Equality of the marked Weyl neighbors]
\label{esrf:thm:equality}
For $\sigma,\tau\in\Omega$, the following three statements are equivalent:
\[
 \Lambda_\sigma=\Lambda_\tau,\qquad
 K_\sigma=K_\tau,\qquad
 \tau-\epsilon\sigma\in6N
                 \text{ for some }\epsilon\in\{1,-1\}.
 \tag{RF3}
\]
\end{theorem}

\begin{proof}
Equality of the two neighbors gives equality of the kernels by
intersection with $N$.
Each $\ell_\sigma$ is onto because
$(\sigma,\alpha_\sigma)=1$.
If the kernels are equal, their induced isomorphisms
$N/K_\sigma\to\mathbb Z/6$ differ by an automorphism of
$\mathbb Z/6$. Its two units are $1,-1$.
Therefore $(\tau-\epsilon\sigma,N)\subset6\mathbb Z$.
Unimodularity $N^\vee=N$ gives
$(\tau-\epsilon\sigma)/6\in N$.
Conversely this congruence makes the two kernels equal.

It remains to show that this kernel criterion actually gives
equality of the chosen unimodular extensions.
The exact formula (RF2), after replacing $k$ by $-k$, gives
$\Lambda_{-\sigma}=\Lambda_\sigma$.
We may consequently suppose $\tau=\sigma+6t$ with $t\in N$.
Both Weyl vectors have squared norm 84, so
\[
 0=\tau^2-\sigma^2=12(\sigma,t)+36t^2,\qquad
                       (\sigma,t)=-3t^2\in6\mathbb Z.
 \tag{RF4}
\]
The last inclusion uses the actual evenness of $N$.
For $x+k\tau/6$ in (RF2), put $x'=x+kt\in N$.
Then $x+k\tau/6=x'+k\sigma/6$, and
\[
 (\sigma,x')+k
 \equiv(\sigma,x)+k
 \equiv(\tau,x)+k\equiv0\pmod6.
\]
This proves $\Lambda_\tau\subset\Lambda_\sigma$.
Interchanging $\sigma,\tau$ proves the reverse inclusion.
Thus no additional unverified choice of an extension is being
identified merely from equality of its common kernel.
\end{proof}

It follows in particular that the map
\[
 \Omega/{\sim}\ \longrightarrow\ \mathscr F,\qquad
 [\sigma]\longmapsto\Lambda_\sigma,\qquad
 \sigma\sim\tau\ \Longleftrightarrow\
                          \tau\equiv\pm\sigma\pmod{6N},
 \tag{RF5}
\]
is a bijection. It is $O(N)$-equivariant: the original formula
$T\Lambda_\sigma=\Lambda_{T\sigma}$ follows by substitution in
(RF2), and $T$ preserves the congruence relation since $T(N)=N$.

\subsection{The exact tetracode obstruction to a fixed neighbor}

\begin{theorem}[No fixed member of this specified family]
\label{esrf:thm:no-fixed-neighbor}
The actual residue element $r=\rho(361)$ fixes no lattice of
$\mathscr F$. Every orbit of its action on $\mathscr F$ has
exactly three members.
\end{theorem}

\begin{proof}
Suppose first that $r\sigma\equiv-\sigma\pmod{6N}$.
Applying $r$ successively, using $r(N)=N$ and $r^3=I$, gives
$\sigma\equiv-\sigma\pmod{6N}$.
Then $\sigma/3\in N$, contradicting its pairing
$1/3$ with the root $\alpha_\sigma\in N$.
Thus (RF3) permits fixedness only if
$t=(r\sigma-\sigma)/6\in N$.

Every isometry of $N$ preserves the original root lattice, permutes
its four $A_5$ components, and acts in each of them by a signed
coordinate permutation, as proved in the full automorphism theorem.
Negation permutes the six entries of $\varrho_A$.
Consequently each $A$ block $s=(s_0,\ldots,s_5)$ of every
$\sigma\in\Omega$ is a permutation of the six distinct numbers
$5/2,3/2,1/2,-1/2,-3/2,-5/2$.
In this block put $\Delta_j=s_{j-2\bmod6}-s_j$.
If the corresponding block $\Delta/6$ of $t$ belongs to $A^\vee$,
its coordinate differences are integral, so
\[
 \Delta_j\equiv c\pmod6\qquad\text{for every }j
 \tag{RF6}
\]
for one integer $c$ modulo six.
On either three-cycle of the index permutation, the three
$\Delta_j$ sum to zero. Hence $3c\equiv0\pmod6$ and
$c\in\{0,2,4\}$.
Each $\Delta_j$ is an integer between $-5$ and $5$.
If $c=0$, (RF6) forces all $\Delta_j=0$, contrary to the
distinctness of the entries on each three-cycle.
Therefore $c=2$ or $4$.

In the original discriminant marking, a vector of the $A$ class
$t_i\lambda+A$ has common fractional part $-t_i/6$.
The corresponding class of $\Delta/6$ is therefore
$t_i=-c\pmod6$, also in $\{2,4\}$.
The difference on the $D$ block is exactly zero because $r|_D=I$.
Thus membership of the full vector $(r\sigma-\sigma)/6$ in $N$
would require a glue word
\[
 (t_1,t_2,t_3,t_4;00)\in H,\qquad t_i\in\{2,4\}
                                                \text{ for every }i.
 \tag{RF7}
\]
In $H=\{(3u+2z;y)\}$ this forces $u=0$ and all four $z_i$
nonzero. But the exact ternary code is
\[
 C_3=\{(a+b,a+2b,a,b):a,b\in\mathbb F_3\}.
 \tag{RF8}
\]
If exactly one parameter is zero, exactly one coordinate is zero.
If both are nonzero, precisely one of $a+b,a+2b$ is zero.
Thus every nonzero word has weight three, and the zero word
has weight zero. No word in (RF7) belongs to $H$.
This contradicts $t\in N$.

Neither sign in (RF3) can occur. Since $r^3=I$, every orbit
therefore has cardinality three.
\end{proof}

The conclusion concerns exactly $\mathscr F$ and the action (RF1).
The integral correspondence for these three-member orbits is
constructed next; no assertion about other embeddings or
constructions of a Leech lattice follows from this obstruction.

\subsection{The canonical orbit and its common integral lattice}

Put $\sigma_j=r^j\varrho$, $\Lambda_j=\Lambda_{\sigma_j}$ and
$K_j=K_{\sigma_j}$, with $j$ read modulo three.
Keep the same original root
$\alpha=e_0-e_1$ in the first $A$ block.
Directly $(\sigma_j,\alpha)=1$ for all three $j$.
Define
\[
 v_j=\sigma_j/6-\alpha,\quad v=v_0,\qquad
 q_1=(\sigma_1-\sigma_0)/6,\quad
 q_2=(\sigma_2-\sigma_0)/6.
 \tag{RF9}
\]
The complete original coordinates are
\[
 \begin{split}
 q_1&=\bigl((-2,-2,1,1,1,1)/3\bigr)^{\oplus4}\perp0,\\
 q_2&=\bigl((-1,-1,-1,-1,2,2)/3\bigr)^{\oplus4}\perp0.
 \end{split}
 \tag{RF10}
\]
Their classes in $R^\vee/R$ are respectively
$(4,4,4,4;00)$ and $(2,2,2,2;00)$.
Neither is in $H$, while $3q_1,3q_2\in R$.
In particular each $q_i+N$ has exact order three.
Direct pairings give
\[
 q_1^2=q_2^2=16/3,\quad(q_1,q_2)=8/3,\quad
                    (\varrho,q_1)=(\varrho,q_2)=-16,
 \qquad q_1+q_2\in R.
 \tag{RF11}
\]

For $x\in N$, write its original glue class as $(3u+2z;y)$,
with $z=(a+b,a+2b,a,b)\in C_3$, and define the homomorphism
\[
 \beta:N\longrightarrow\mathbb F_3,\qquad\beta(x)=b=z_4.
 \tag{RF12}
\]
This map is defined through the original quotient $N/R$.
This definition uses the retained residue marking without changing it.
The sum of the four $z$ coordinates is also $b$ modulo three.
Write $\ell_j(x)=(\sigma_j,x)\bmod6$.
Then the exact three character formulas are
\[
 \ell_1(x)=\ell_0(x)+4\beta(x),\qquad
 \ell_2(x)=\ell_0(x)+2\beta(x)\quad\text{in }\mathbb Z/6.
 \tag{RF13}
\]
The maps $b\mapsto2b,4b$ from $\mathbb F_3$ to $\mathbb Z/6$
are well-defined homomorphisms.
To prove (RF13), the discriminant bilinear form gives
\[
 (q_1,x)\equiv\tfrac56\,4\sum t_i
       \equiv\tfrac13\sum t_i
       \equiv\tfrac23\beta(x)\pmod{\mathbb Z},
 \qquad t_i=3u_i+2z_i\pmod6.
\]
For $q_2$ the same calculation gives
$(q_2,x)\equiv\beta(x)/3\pmod{\mathbb Z}$.
Multiplication by six and (RF9) prove both formulas.

Using the original dominant representatives
$\omega_t=(1^t,0^{6-t})-\frac t6(1,\ldots,1)$, put
\[
 u_*=(\omega_2,\omega_4,0,\omega_2;0)\in N.
 \tag{RF14}
\]
Its ternary class is exactly the original generator
$z_2=(1,2,0,1)$.
It has $u_*^2=4$, $(\varrho,u_*)=12$ and $\beta(u_*)=1$.
The map
\[
 (\ell_0,\beta):N\longrightarrow\mathbb Z/6\oplus\mathbb F_3
 \tag{RF15}
\]
is onto: $\alpha$ maps to $(1,0)$ and $u_*$ to $(0,1)$.
Consequently
\[
 \begin{gathered}
 I=\ker(\ell_0,\beta)=K_0\cap K_1=K_1\cap K_2
     =K_2\cap K_0=\bigcap_jK_j,\\
 [N:I]=18,\qquad[K_j:I]=3,\qquad\det I=324.
 \end{gathered}
 \tag{RF16}
\]
The kernel equalities follow from (RF13), because either
$2b=0$ or $4b=0$ in $\mathbb Z/6$ forces $b=0$ in
$\mathbb F_3$.
The action $r$ permutes the $K_j$, so it preserves $I$.
Also $3q_i\in I$: these are in $R$, hence have $\beta=0$,
and $(\varrho,3q_i)=-48$ is divisible by six.

Define the full-rank lattice
\[
 J=I+\mathbb Z(3v).
 \tag{RF17}
\]
It is contained in $\Lambda_0$.
The element $6v=\varrho-6\alpha$ lies in $I$: its
$\ell_0$ value is zero, and its glue has zero ternary part.
The element $3v$ is not in $N$, because
$(3v,\alpha)=-11/2$ is not an integer.
Thus
\[
 N\cap J=I,\qquad[J:I]=2,\qquad
                    \det J=324/4=81,\qquad[\Lambda_0:J]=9.
 \tag{RF18}
\]
This lattice is positive definite, even, and rootless, because
it is a sublattice of $\Lambda_0$.

All exact root corrections in the cyclic action are retained.
Put $\kappa_j=\alpha-r^j\alpha\in R$.
The vectors $r\alpha=e_2-e_3$ and $r^2\alpha=e_4-e_5$
are original simple roots with $\varrho$-pairing one.
Hence $\kappa_j\in I$.
The precise formulas, rather than literal identification of
different generators, are
\[
 r^jv=v+q_j+\kappa_j\quad(j=1,2),\qquad
 r(3v)=3v+3q_1+3\kappa_1.
 \tag{RF19}
\]
Since $rI=I$, $3q_1\in I$ and $\kappa_1\in I$, this proves
$rJ\subset J$. Applying $r^3=I$ gives equality $rJ=J$.

\begin{theorem}[All three pairwise Leech intersections]
\label{esrf:thm:common-J}
The actual three-member orbit satisfies
\[
 \Lambda_0\cap\Lambda_1
 =\Lambda_1\cap\Lambda_2
 =\Lambda_2\cap\Lambda_0
 =\bigcap_j\Lambda_j=J,\qquad[\Lambda_j:J]=9.
 \tag{RF20}
\]
The element $r$ is an integral isometry of $J$, and induces
isometries $\Lambda_j\to\Lambda_{j+1}$ with this common restriction.
\end{theorem}

\begin{proof}
First $J\subset\Lambda_0\cap\Lambda_1$: $I$ lies in both,
and $3v=3v_1-3q_1$ with $3q_1\in I$.
Suppose $y\in\Lambda_0\cap\Lambda_1$.
Using the same $\alpha$, write
$y=x_0+k_0v_0=x_1+k_1v_1$, $x_j\in K_j$.
The fractional part of $(v_j,\alpha)$ is $1/6$ for both
$j$, so pairing with $\alpha\in N$ gives
$k_0\equiv k_1\pmod6$.
Absorbing multiples of $6v_j\in K_j$, take the same
$k\in\{0,\ldots,5\}$ in both expressions.
Then $x_0-x_1=kq_1\in N$.
The exact order three of $q_1+N$ forces $k=0$ or $3$.
For these values $kq_1\in I\subset K_1$, so
$x_0=x_1+kq_1\in K_0\cap K_1=I$.
Thus $y=x_0+kv\in J$, proving the first equality.
Apply $r$ and $r^2$, using $rJ=J$, to prove the other pairwise
equalities. Their common index is nine by (RF18) and the
isometries induced by $r$.
\end{proof}

For comparison with the original index-72 glue, let
$R_0=R\cap K_0$. Since $\beta|_R=0$, (RF13) gives
$R\cap K_j=R_0$ for every $j$, and $R\cap I=R_0$.
The reduction map $K_0/R_0\cong H$ identifies $I/R_0$
with the kernel of the original ternary coefficient $b$.
Thus
\[
 [I:R_0]=24,\qquad[J:R_0]=48,\qquad
 [K_0:R_0]=72,\qquad[\Lambda_j:R_0]=432.
 \tag{RF21}
\]
All these are inclusions of the displayed original lattices.

\subsection{The exact discriminant form and its cyclic three planes}

\begin{theorem}[The integral correspondence over $\mathbb F_3^4$]
\label{esrf:thm:discriminant-correspondence}
One has
\[
 J^\vee=\Lambda_i+\Lambda_j\quad(i\ne j),\qquad
 J^\vee/J\cong\mathbb F_3^4.
 \tag{RF22}
\]
In the ordered basis
\[
 e_1=u_*+J,\quad e_2=v+J,\quad
 e_3=ru_*+J,\quad e_4=rv+J,
 \tag{RF23}
\]
the discriminant quadratic form is
\[
 q_J(a,b,c,d)=\frac23(2ad+bc+bd)\pmod{2\mathbb Z},
 \tag{RF24}
\]
and three times its discriminant bilinear form has matrix
\[
 B=
 \begin{pmatrix}
 0&0&0&2\\0&0&1&1\\0&1&0&0\\2&1&0&0
 \end{pmatrix}\quad\text{over }\mathbb F_3.
 \tag{RF25}
\]
The three original Leech extensions are exactly the inverse
images in $J^\vee$ of the isotropic planes
\[
 \begin{split}
 P_0&=\{(a,b,0,0):a,b\in\mathbb F_3\},\\
 P_1&=\{(0,0,c,d):c,d\in\mathbb F_3\},\\
 P_2&=\{(a,b,a-b,b):a,b\in\mathbb F_3\}.
 \end{split}
 \tag{RF26}
\]
The induced action of the actual $r$ is the matrix
\[
 \mathcal R=
 \begin{pmatrix}
 0&0&2&1\\
 0&0&0&2\\
 1&0&2&2\\
 0&1&0&2
 \end{pmatrix},\qquad
 \mathcal R^3=I,\quad \mathcal R^TB\mathcal R=B,
 \quad\mathcal R P_j=P_{j+1}.
 \tag{RF27}
\]
All entries in (RF23)--(RF27) use the original glue generator
$z_2$, root $\alpha$, pairing, and residue permutation.
\end{theorem}

\begin{proof}
Since $J\subset\Lambda_i=\Lambda_i^\vee$, each $\Lambda_i$
lies in $J^\vee$.
By (RF20), the two subgroups $\Lambda_i/J$ and $\Lambda_j/J$
have order nine and intersect trivially for $i\ne j$.
Their sum has order 81, equal to
$|J^\vee/J|=\det J=81$, proving the first equality in (RF22).

The homomorphism $\beta:K_0/I\to\mathbb F_3$ is an isomorphism,
with generator $u_*+I$. Hence $\Lambda_0/J$ is generated by
$u_*+J,v+J$.
Both have order dividing three because $3u_*\in I$ and
$3v\in J$.
They are independent: if $a u_*+b v\in J$, reduce first modulo
$N$; the subgroup $J+N$ has the two cosets represented by
$0,3v$, whereas $v+N$ has order six.
For $a,b\in\{0,1,2\}$ this forces $b=0$.
Then membership of $a u_*$ in $I=N\cap J$ forces $a=0$.
Thus $\Lambda_0/J\cong\mathbb F_3^2$ with the first two
generators in (RF23).
Applying $r$ proves the corresponding statement for
$\Lambda_1/J$ and the last two generators.
Their trivial intersection and total order 81 prove that
(RF23) is indeed a basis of the full discriminant group.

Here are the full Euclidean Gram entries before any passage
to a quotient:
\[
 \operatorname{Gram}(u_*,v,ru_*,rv)=
 \begin{pmatrix}
 4&2&-2&-4/3\\
 2&4&-2/3&-2/3\\
 -2&-2/3&4&2\\
 -4/3&-2/3&2&4
 \end{pmatrix}.
 \tag{RF28}
\]
For an explicit verification, each of the three nonzero
$A$ blocks of $u_*$ pairs to $-2/3$ with its shifted block,
giving $(u_*,ru_*)=-2$.
One has
$(u_*,\alpha)=(u_*,r\alpha)=(ru_*,\alpha)=0$,
$(\varrho,u_*)=12$,
$(\sigma_1,u_*)=-8$, and
$(\sigma_2,u_*)=-4$, directly from (RF10) and (RF14).
These give respectively
$(u_*,v)=2$, $(u_*,rv)=-4/3$, and
$(v,ru_*)=-2/3$.
Also $(\varrho,\sigma_1)=-12$,
$(\varrho,r\alpha)=(\sigma_1,\alpha)=1$,
and $(\alpha,r\alpha)=0$. Thus
\[
 (v,rv)=-12/36-1/6-1/6=-2/3.
\]
The four diagonal entries equal four, and
$(ru_*,rv)=(u_*,v)=2$ by the exact isometry $r$.
This proves (RF28).
Reducing its entries modulo $\mathbb Z$ gives (RF25).
The diagonal and integral cross terms in the squared norm
are even. Reducing the remaining cross terms modulo
$2\mathbb Z$ gives exactly (RF24).
Changing any coordinate by three changes (RF24) by an even
integer, so it is well-defined on the asserted field basis.

It remains to calculate the action without identifying a
family permutation with a fixed Leech automorphism.
By definition $r e_1=e_3$, $r e_2=e_4$.
The vector $u_*+ru_*+r^2u_*$ lies in $I$: its $\beta$ value
is $1+1+1=0$, because $r$ acts trivially on $N/R$,
and its $\ell_0$ value is $12-4-8=0$.
Consequently
\[
 r^2u_*\equiv-u_*-ru_*\pmod J.
 \tag{RF29}
\]
By (RF19),
\[
 v+rv+r^2v
       =3v+q_1+q_2+\kappa_1+\kappa_2.
\]
The $\kappa_j$ belong to $I$, and $3v\in J$.
Furthermore $q_1+q_2-u_*+ru_*\in I$.
Indeed $q_1+q_2\in R$, so its $\beta$ value is zero;
subtracting $u_*$ and adding $ru_*$ keeps that value zero.
Its pairing with $\varrho$ is
$-32-12-4=-48$, divisible by six.
Thus
\[
 r^2v\equiv u_*-v-ru_*-rv\pmod J.
 \tag{RF30}
\]
The four images in (RF29)--(RF30) give the four columns of
$\mathcal R$ in (RF27).
The matrix has cube the identity and preserves $B$ either
by multiplication, or by the already proved facts $r^3=I$
and $rJ=J$ with its original isometric pairing.

The first two planes in (RF26) are respectively
$\Lambda_0/J$ and $\Lambda_1/J$ by their proved bases.
Their next image is the span of
$(-1,0,-1,0)$ and $(1,-1,-1,-1)$, from (RF29)--(RF30).
Solving their two linear parameters gives exactly
$(a,b,a-b,b)$, proving the stated $P_2=\Lambda_2/J$.
Each plane is isotropic: the first two vanish in (RF24),
and the third has
$2ab+b(a-b)+b^2=3ab=0$ in $\mathbb F_3$.
They have nine elements and pairwise intersection zero,
also following directly from (RF26).
Finally each $\Lambda_j$ is exactly the inverse image of
its quotient plane because it contains $J$.
This proves the claimed integral correspondence and all
its maps, indices, pairings, and cyclic actions.
\end{proof}

\subsection{The fixed fourth plane and its exact \texorpdfstring{$A_2^{12}$}{A2 to the twelfth power} extension}

The same finite discriminant form has the additional plane
\[
 P_\infty=\{(a,b,-a+b,-b):a,b\in\mathbb F_3\}.
 \tag{RF31}
\]
Put
\[
 p=u_*-ru_*,\qquad q=v+ru_*-rv,\qquad
                       M=J+\mathbb Zp+\mathbb Zq.
 \tag{RF32}
\]
These are literal vectors in the original Euclidean space.

\begin{theorem}[The fixed fourth extension]
\label{esrf:thm:fourth-extension}
The plane $P_\infty$ is isotropic and is fixed by the induced
residue action. Its inverse image $M$ is a positive definite
even unimodular lattice with exactly 72 roots, of root system
$A_2^{12}$ and minimum squared norm two.
The original residue element $r$ is an integral isometry of $M$.
It fixes four of these $A_2$ components pointwise and acts as
a coordinate three-cycle on each of the other eight.

The exact intersection and index relations are
\[
 \begin{gathered}
 M\cap\Lambda_j=J,\qquad [M:J]=[\Lambda_j:J]=9,\\
 M\cap N=L_0:=\{x\in N:\beta(x)=0,\
                                (\varrho,x)\equiv0\pmod2\},\\
 [M:L_0]=[N:L_0]=6.
 \end{gathered}
 \tag{RF33}
\]
If $R_M$ denotes the lattice generated by all its roots, then
\[
 R_M\cong A_2^{\perp12},\qquad
 \det R_M=3^{12}=531441,\qquad[M:R_M]=3^6=729.
 \tag{RF34}
\]
\end{theorem}

\begin{proof}
Substitution of (RF31) in (RF24) gives the numerator
$2a(-b)+b(-a+b)+b(-b)=-3ab=0$ in $\mathbb F_3$.
Thus $P_\infty$ is isotropic.
The matrix (RF27) maps its vector with parameters $(a,b)$ to
\[
 (a+b,b,-a,-b),
 \tag{RF35}
\]
which is again in $P_\infty$, now with parameters $(a+b,b)$.
The images of $p,q$ in (RF23) are
$(1,0,-1,0)$ and $(0,1,1,-1)$, a basis of $P_\infty$.
Consequently (RF32) is its exact inverse image under
$J^\vee\to J^\vee/J$, and $[M:J]=9$.
Isotropy gives even norms on $M$; polarization gives integral
pairings. Its determinant is $\det J/9^2=1$, and its
positive definiteness and rank are inherited from the unchanged
ambient real space. Equation (RF35), together with $rJ=J$,
proves $rM=M$.
Each of $P_0,P_1,P_2$ intersects $P_\infty$ in zero, as follows
by equating their coordinates in (RF26) and (RF31).
Taking inverse images proves $M\cap\Lambda_j=J$.

For its relation with $N$, (RF19) gives the exact expression
\[
 q=ru_*-q_1-\kappa_1 .
 \tag{RF36}
\]
Therefore
\[
 N+M=N+\mathbb Z(\varrho/2)+\mathbb Zq_1.
 \tag{RF37}
\]
The two indicated cosets modulo $N$ have orders two and three.
For the first, pair with $\alpha$; for the second, use (RF10).
Subgroups of coprime orders intersect trivially, so the quotient
in (RF37) has order six.
Also $p\in N$ has $\beta(p)=0$ and
$(\varrho,p)=12-(-4)=16$, hence image $(4,0)$ in (RF15).
Thus $p+I$ has order three.
The vector $3q$ lies in $J$ because $q+J$ has order three,
and lies in $N$ by (RF36) and $3q_1\in N$.
It is therefore in $N\cap J=I$.
We also have $6v\in I$.
Write an element of $M$ using generators $I,3v,p,q$ and reduce
the coefficient of $3v$ modulo two and of $q$ modulo three.
If this element is in $N$, their two coprime-order cosets in
(RF37) force both reduced coefficients to vanish.
Hence
\[
 M\cap N=I+\mathbb Zp.
 \tag{RF38}
\]
Its image in (RF15) is precisely $2(\mathbb Z/6)\oplus0$,
which proves the displayed formula for $L_0$ and
$[N:L_0]=6$. Since $\det N=\det M=1$, the equal covolumes
give $[M:L_0]=6$, proving (RF33).

We next prove that the following root count is exhaustive.
Every vector in $N+M$ has one of the six forms
\[
 x+a\varrho/2-bq_1,\qquad
                   x\in N,\ a\in\{0,1\},\ b\in\{0,1,2\}.
 \tag{RF39}
\]
When $a=0$, the vector is in $R^\vee$.
Every norm-two vector of this particular $R^\vee$ is already
a root of $R$.
Indeed, the $A$ class minima are
$0,5/6,4/3,3/2,4/3,5/6$, and the three nonzero $D$ classes
have minimum one. In a nonzero class of $R^\vee/R$ with
minimum at most two, the possible minima are exactly among
\[
 \{5/6,1,4/3,3/2,5/3,11/6\}.
 \tag{RF40}
\]
This exhausts them because at most two $A$ classes can be
nonzero under the bound; if two are nonzero, each has minimum
$5/6$; with a nonzero $D$ class there is at most one nonzero
$A$ class and its minimum must be $5/6$.
No number in (RF40) is congruent to two modulo $2\mathbb Z$.
All norms within a fixed discriminant class differ by an even
integer, so no such class has a norm-two vector.
This proves the claim about $R^\vee$.

All original roots have $\beta=0$, so (RF33) says that those
lying in $M$ are exactly the ones with even $\varrho$ height.
In each original $A$ block they are
$e_i-e_j$ with $i\ne j$ and $i\equiv j\pmod2$.
These give two orthogonal $A_2$ systems per block, on the
coordinate triples $\{0,2,4\}$ and $\{1,3,5\}$, and 48 roots
in total. In $D$ they are the eight roots
$\pm e_1\pm e_3,\ \pm e_2\pm e_4$.
They form four pairwise orthogonal $A_1$ systems.
This accounts for exactly 56 roots with $a=0$ in (RF39).

For $a=1$, absorb the $A$ dual vector $-bq_1$ into the
marked $A$ discriminant classes.
If the original class of $x$ is $(t_1,t_2,t_3,t_4;y)\in H$,
the shifted class in each $A$ is $t_i-4b$.
The exact half-Weyl-shift minima proved in the neighbor
calculation are
\[
 \min_{w\in t\lambda+A}|w+\varrho_A/2|^2
 =\begin{cases}3/8& t\equiv0\pmod3,\\17/24&t\not\equiv0\pmod3,\end{cases}
 \qquad
 \min_{w+D=y}|w+\varrho_D/2|^2=1/2.
 \tag{RF41}
\]
For completeness the first line follows by choosing centered
quarter coordinates when $t=0,3$; for other $t$, the nearest
twelfth coordinates have square sum $13/24$ and coordinate
sum $\pm1$, and restoring zero sum costs $1/6$.
The second line follows from the integer coordinate distances
$(1/2,0,1/2,0)$ and half-integer distances
$(0,1/2,0,1/2)$, with their ties meeting every $D$ marking.
Thus the total minimum in (RF39) with $a=1$ is at least two.
Equality requires $t_i-4b\equiv0\pmod3$ in every component.
Since $t_i\equiv2z_i\pmod3$, this says
$z_i=2b$ for all four coordinates.
For $b=1,2$ this is excluded by the exact code (RF8).
For $b=0$ it requires $z=0$.
The eight binary codewords and the two $D$ ties then give
exactly sixteen equality vectors $y=x+\varrho/2$.
Every one belongs to $M$: the norm equation gives
\[
 2=x^2+(\varrho,x)+21,\qquad
                         (\varrho,x)=-19-x^2,
 \tag{RF42}
\]
which is odd because $N$ is even.
Also $\beta(x)=0$ because $z=0$.
Therefore $x+3\alpha\in L_0$ and
$y=(x+3\alpha)+3v\in M$.
This proves exhaustively that $M$ has exactly $56+16=72$ roots.

Here are the full root coordinates and the resulting exact
orthogonal decomposition.
Let $\zeta_i$ be supported in the $i$th original $A$ block
with entries
\[
 \zeta=(1,-1,1,-1,1,-1)/4,\qquad \zeta^2=3/8.
\]
In the following table every $Z_y$ is an actual sum of those
four mutually orthogonal vectors, and every $\delta_y$ is
an original $D$ root:
\[
 \begin{array}{c|l|l}
 y&Z_y&\delta_y\\\hline
 00&\zeta_1+\zeta_2+\zeta_3+\zeta_4&e_1+e_3\\
 10&-\zeta_1-\zeta_2+\zeta_3+\zeta_4&e_1-e_3\\
 01&-\zeta_1+\zeta_2-\zeta_3+\zeta_4&e_2+e_4\\
 11&\zeta_1-\zeta_2-\zeta_3+\zeta_4&e_2-e_4
 \end{array}.
 \tag{RF43}
\]
The sixteen new roots are exactly
\[
 \{\epsilon Z_y+\eta\delta_y/2:
                   y\in\mathbb F_2^2,\ \epsilon,\eta\in\{1,-1\}\}.
 \tag{RF44}
\]
To verify the signs from the original glue, the $A$ half-shift
minimizer for binary coordinate $u_i=0$ is $\zeta_i$ and
for $u_i=1$ it is $-\zeta_i$.
For $y=00$ the two codewords are $0000,1111$; for $10$ they
are $1100,0011$; for $01$ they are $1010,0101$; and for $11$
they are $0110,1001$.
These give respectively the four opposite pairs $\pm Z_y$
in (RF43).
The two shifted $D$ minimizers in these exact original
markings are respectively
$\pm(e_1+e_3)/2,\pm(e_1-e_3)/2,
 \pm(e_2+e_4)/2,\pm(e_2-e_4)/2$.
This proves (RF44) with every sign and marking retained.

The four $Z_y$ are mutually orthogonal and have squared norm
$3/2$, as follows by dotting their four displayed sign rows.
The $\delta_y$ are mutually orthogonal of squared norm two.
The $Z_y$ are orthogonal to the entire original $D$ space,
and to the eight old $A_2$ systems because their coordinates
are constant on each parity triple.
For each $y$, the six roots
\[
 \{\pm\delta_y\}
       \ \cup\ \{\epsilon Z_y+\eta\delta_y/2:\epsilon,\eta=\pm1\}
 \tag{RF45}
\]
are an $A_2$ system: take simple roots
$\delta_y$ and $Z_y-\delta_y/2$, whose Gram matrix is
$\left(\begin{smallmatrix}2&-1\\-1&2\end{smallmatrix}\right)$;
their sum is $Z_y+\delta_y/2$.
The four such systems are mutually orthogonal and orthogonal
to the eight old systems. They exhaust the 72 proved roots.
Thus their root lattice is $A_2^{\perp12}$ of determinant
$3^{12}$, and unimodularity of $M$ gives the index $3^6$.
This proves (RF34) and the Coxeter number three of this
rooted Niemeier lattice directly from its twelve components.

Finally $r$ fixes each $\zeta_i$ because it permutes coordinates
within their parity triples; it fixes each $\delta_y$ because
it is the identity on $D$. Hence it fixes the four systems
(RF45) pointwise. On each of the eight old $A_2$ systems it
cyclically permutes its three coordinate positions, giving its
exact order-three Coxeter action.
This completes the asserted classification of the fixed
fourth extension and its precise relation to the three Leech
lattices through $J$.
\end{proof}

% New J/M extension; prerequisite definitions are in es_niemeier_cubic_arithmetic.tex.
\section{The cubic on the common lattice and the fourth Niemeier extension}
\label{jmv:sec:JM-extension}

Retain the original coordinates, Cayley multiplication, $F$, the generators $g_i$
and their heights $H_i$ from Section~\ref{nca:sec:cubic-arithmetic}, and the
exact common lattice and fourth extension of Section~\ref{esrf:sec:fixed-neighbor}.
The symbol $u$ below is the unchanged vector $u_*$ in (RF32).
\subsection{Integral generators of the common lattice and the fixed extension}

Define the unchanged residue action by
\[
 (rX)_{i,j}=x_{i,j-2\pmod6},\qquad (rX)_{D,k}=d_k.
\]
The two primary parts of the original glue have the exact descriptions
\begin{align*}
 H_2={}&\{(\varepsilon_1+\varepsilon_2+\varepsilon_3,
 \varepsilon_1+\varepsilon_2,\varepsilon_1+\varepsilon_3,
 \varepsilon_1;\varepsilon_2,\varepsilon_3):\varepsilon_i\in\mathbb F_2\},\\
 H_3={}&\{(t+s,t+2s,t,s;00):t,s\in\mathbb F_3\}.
\end{align*}
In the original \((\mathbb Z/6)^4\) coordinates this means
\(k=3w+2z\), with \(w\) the four binary entries of \(H_2\)
and \(z\) the four ternary entries of \(H_3\). The binary
\(D_4\) marking is retained. Chinese remainders recover \(w,z\)
uniquely, and the displayed parameterizations are injective.
Consequently
\[
 \beta:N\longrightarrow\mathbb F_3,\qquad \beta(X)=z_4=s
\]
is a well-defined additive homomorphism. The actual lifts
\eqref{nca:eq:N-generators} satisfy
\[
 \beta(g_i)=0\quad(1\leq i\leq28),\qquad \beta(g_{29})=1.
\]
Indeed roots have zero glue, the first three lifts are two-primary,
and the last two have \((t,s)=(1,0),(0,1)\), respectively.

Set, with no change of representatives,
\begin{align}
 I&=\{X\in N:(\rho,X)\equiv0\pmod6,\ \beta(X)=0\},
       &J&=I+\mathbb Z(3v),\label{jmv:eq:I-J}\\
 \omega_t&=(\underbrace{1,\ldots,1}_{t},0,\ldots,0)
                    -\frac t6(1,1,1,1,1,1),
       &u&=(\omega_2,\omega_4,0,\omega_2;0),\notag\\
 p&=u-r(u),&q&=v+r(u)-r(v),\notag\\
 M&=J+\mathbb Zp+\mathbb Zq.&&\label{jmv:eq:M}
\end{align}
Thus \(J\) and \(M\) are precisely the common-lattice and
residue-fixed extension presentations being evaluated here; the
following value-group proof needs no root count, lattice determinant,
or classification label for \(M\). For explicitness, the two extra
vectors in the original five blocks are
\begin{align*}
 p={}&((1,1,-1,-1,0,0),(1,1,0,0,-1,-1),0,
                         (1,1,-1,-1,0,0);0),\\
 q={}&\tfrac13((-2,4,4,-2,-2,-2),0,(2,2,-1,-1,-1,-1),
                         (1,1,1,1,-2,-2);0).
\end{align*}
The zero blocks have their original sizes six and four. These
coordinates follow by substituting \(v=\rho/6-\alpha\) into
\eqref{jmv:eq:M}; in particular the shifted \(\alpha\) term in
\(r(v)\) is retained. Also \(u\in N\),
\(\beta(u)=1\), \((\rho,u)=12\), and \((u,u)=4\):
its class is \((2,4,0,2;00)\),
\((\rho_A,\omega_t)=t(6-t)/2\), and
\((\omega_t,\omega_t)=t(6-t)/6\).

\begin{theorem}[Integral generating lists with both kernel conditions]
\label{jmv:thm:IJM-generators}
Define the following 32 ordered vectors:
\begin{align}
 b_1&=6\alpha,&
 b_i&=g_i-H_i\alpha\quad(2\leq i\leq28),\notag\\
 b_{29}&=3g_{29}-60\alpha,&
 b_{30}&=3v=\rho/2-3\alpha,\notag\\
 b_{31}&=p,&b_{32}&=q.\label{jmv:eq:IJM-generators}
\end{align}
Then the integral, not merely rational, equalities are
\[
 I=\sum_{i=1}^{29}\mathbb Zb_i,\qquad
 J=\sum_{i=1}^{30}\mathbb Zb_i,\qquad
 M=\sum_{i=1}^{32}\mathbb Zb_i.
\]
\end{theorem}
\begin{proof}
Every one of the first 29 vectors has \(\beta=0\) and height
divisible by six, by the exact heights and characters already
calculated. Conversely take any integral expression
\(X=\sum_{i=1}^{29}n_ig_i\) for \(X\in I\). Applying
\(\beta\) shows \(n_{29}=3a\) for an integer \(a\),
regardless of which expression was chosen. The exact equality
\[
 X=\sum_{i=2}^{28}n_i(g_i-H_i\alpha)
        +a(3g_{29}-60\alpha)
        +\left(n_1+\sum_{i=2}^{28}n_iH_i+60a\right)\alpha
\]
has last coefficient \((\rho,X)\), which is divisible by six.
It therefore expresses \(X\) as an integer combination of the
first 29 \(b_i\). The reverse containment was proved for each
generator. The last two equalities follow by adjoining exactly
the vectors in \eqref{jmv:eq:I-J} and \eqref{jmv:eq:M}.
No assertion of linear independence is used.
\end{proof}

\subsection{Complete finite certificates for the two new value groups}

\begin{theorem}[Equality of the common-lattice and fixed-extension value-generated groups]
\label{jmv:thm:JM-values}
For the retained Cayley multiplication and the original ordered
coordinates, the two additive cubic-value groups are exactly
\begin{equation}
 \mathcal V_F(J)=\mathcal V_F(M)
 =\frac1{27}\mathbb Z\ \oplus\frac{\sqrt2}{216}\mathbb Z
                         \ \oplus\frac{\sqrt3}{18}\mathbb Z.
 \label{jmv:eq:JM-values}
\end{equation}
The equality concerns generated additive groups, not equality
of individual-value sets \(F(J)\) and \(F(M)\).
\end{theorem}
\begin{proof}
Apply Theorem~\ref{nca:thm:finite-values} to the first 30 and
all 32 vectors in \eqref{jmv:eq:IJM-generators}. The resulting
list lengths are exactly
\[
 30+2\binom{30}{2}+\binom{30}{3}=4960,\qquad
 32+2\binom{32}{2}+\binom{32}{3}=5984.
\]
Write all pair and triple symbols in the next tables using these
\(b_i\), not the \(g_i\) or \(h_i\). Substitution in the full
polynomial gives the complete denominator certificate
\begin{equation}
\begin{array}{c|c|c|c|c}
 &A_i&D^+_{ij}&D^-_{ij}&D_{ijk}\\\hline
 J&(27,216,1)&(9,36,18)&(9,36,18)&(9,36,18)\\
 M&(27,216,1)&(9,36,18)&(9,36,18)&(9,36,18)
\end{array}.
\label{jmv:eq:JM-denominators}
\end{equation}
As before each triple is the least common multiple of all reduced
rational denominators in its full category, in the order
\((P_0,P_2,P_3)\); a zero has denominator one. The finite
arithmetic is completely specified by
\eqref{nca:eq:original-s}--\eqref{nca:eq:P3},
\eqref{jmv:eq:IJM-generators}, and all indices in
\eqref{nca:eq:single}--\eqref{nca:eq:Dthree}.
The accompanying standard-library checker
\texttt{check\_jm\_values.py}, using the included
\texttt{cubic\_engine.py}, executes every entry with exact
fractions and asserts this table. It reconstructs all 72 marked
glue classes, checks all 29 \(\beta(g_i)\), and checks both
kernel conditions on all the displayed \(I\) generators.
These finite computations establish the upper containment in
\eqref{jmv:eq:JM-values}, including all mixed polarization terms
involving \(p,q\), not just their individual cubic values.

Here is a short certificate for the reverse containment, using
only \(J\). Coordinates of each row are taken in the ordered
proposed value-group basis
\((1/27,\sqrt2/216,\sqrt3/18)\):
\begin{equation}
\begin{array}{c|r r r}
 F(b_2)&14&-272&0\\
 F(b_4)&5&-128&0\\
 F(b_{27})&23220&-111672&-90\\
 F(b_{28})&13780&-74672&0\\
 F(b_{30})&378&-1701&-18\\
 D^+_{2,27}&7425&-60624&-3\\
 D^+_{11,21}&36&48&-16
\end{array}.
\label{jmv:eq:JM-lower-certificate}
\end{equation}
Let \(C_1,C_2,C_3\) have, respectively, columns obtained by
transposing rows \((2,3,5)\), \((2,5,6)\), and \((1,4,7)\)
of this table. Direct three-by-three determinants are
\begin{align*}
 d_1&=-39859290,&d_2&=11531403,&d_3&=-43244032.
\end{align*}
The following explicit integer identity certifies their gcd one:
\begin{equation}
 (-6882196252)d_1+(-23788908901)d_2+(-907)d_3=1.
 \label{jmv:eq:JM-bezout}
\end{equation}
Put \((c_1,c_2,c_3)=(-6882196252,-23788908901,-907)\).
For each standard basis column \(e_t\in\mathbb Z^3\) the exact
integer reconstruction is
\begin{equation}
 e_t=\sum_{j=1}^3 C_j\bigl(c_j\operatorname{adj}(C_j)e_t\bigr).
 \label{jmv:eq:JM-axis-reconstruction}
\end{equation}
All columns of every \(C_j\) represent integer combinations of
actual cubic values on \(J\); their integer coefficient vectors
in \eqref{jmv:eq:JM-axis-reconstruction} therefore give each
of the three claimed generators as such a combination. This is
an explicit reverse-containment morphism, not an index inference.
The checker also expands every adjugate and verifies all three
axis reconstructions, saving their integer coefficients in
\texttt{EXACT\_RESULTS.json}. Thus \(\mathcal V_F(J)\) contains
the right-hand side of \eqref{jmv:eq:JM-values}. Since \(J\subset M\),
the same certificate proves its containment in \(\mathcal V_F(M)\).
The upper bounds already established finish both equalities.
\end{proof}

For additional exact checks on the unchanged representatives, the
three individual values are
\[
 F(3v)=14-\frac{63}{8}\sqrt2-\sqrt3,\qquad
 F(p)=-2,\qquad F(q)=\frac{64}{27}.
\]
They agree with the corresponding rows of the complete calculation;
they are not substituted for the mixed-term certificate.

\subsection{Exact inclusions, quotient maps, and the factor of the Albert determinant}

\begin{theorem}[The exact four-to-three chain of value groups]
\label{jmv:thm:chain}
There is a strict chain of additive subgroups of the original real
coefficient field
\[
 \mathcal V_F(N)\ \subsetneq\ \mathcal V_F(J)=\mathcal V_F(M)
                       \ \subsetneq\ \mathcal V_F(\Lambda).
\]
The inclusion matrices in the ordered bases already displayed are
\[
 \mathcal V_F(N)\longrightarrow\mathcal V_F(J)=\mathcal V_F(M):
       (a,b,c)\longmapsto(a,4b,c),
\]
\[
 \mathcal V_F(J)=\mathcal V_F(M)\longrightarrow\mathcal V_F(\Lambda):
       (a,b,c)\longmapsto(a,b,3c).
\]
Their exact quotient maps are
\begin{align*}
 \frac a{27}+\frac{b\sqrt2}{216}+\frac{c\sqrt3}{18}
       &\longmapsto b\pmod4,\\
 \frac a{27}+\frac{b\sqrt2}{216}+\frac{c\sqrt3}{54}
       &\longmapsto c\pmod3.
\end{align*}
In particular
\begin{align*}
 \mathcal V_F(J)/\mathcal V_F(N)&\cong\mathbb Z/4,&
 \mathcal V_F(M)/\mathcal V_F(N)&\cong\mathbb Z/4,\\
 \mathcal V_F(\Lambda)/\mathcal V_F(J)&\cong\mathbb Z/3,&
 \mathcal V_F(\Lambda)/\mathcal V_F(M)&\cong\mathbb Z/3,\\
 \mathcal V_F(M)/\mathcal V_F(J)&=0.&&
\end{align*}
\end{theorem}
\begin{proof}
The first basis comparison uses
\(\sqrt2/54=4(\sqrt2/216)\), with the other basis entries equal.
The second uses \(\sqrt3/18=3(\sqrt3/54)\), with its other
entries equal. Uniqueness of rational radical coefficients proves
that the displayed kernels are exactly the corresponding subgroups.
Surjectivity is witnessed by \(\sqrt2/216\) and \(\sqrt3/54\),
respectively, both actual elements of the proved additive groups.
Their orders modulo the smaller groups are exactly four and three:
the coefficient divisibility conditions are necessary and sufficient.
This also proves strictness. The last quotient is zero by the equality
of the two proved groups. In particular the composite inclusion is
\((a,b,c)\mapsto(a,4b,3c)\), recovering the full quotient
\(\mathbb Z/4\oplus\mathbb Z/3\), with explicit map
\[
 \frac a{27}+\frac{b\sqrt2}{216}+\frac{c\sqrt3}{54}
       \longmapsto (b\bmod4,c\bmod3).
\]
Its identification with \(\mathbb Z/12\) is
\((\bar b,\bar c)\mapsto9b+4c\pmod{12}\), whose inverse is
reduction modulo four and modulo three. Every assertion follows
from the proved cubic coefficient groups, not from an index of
the original lattices.
\end{proof}

The factor two in the original Albert determinant is unchanged:
\[
 \det Q_{\mathbb O}(\lambda;\mathcal LX)
       =\lambda^3-\lambda(X,X)+2F(X).
\]
Thus, if one instead asks for the additive group generated by the
actual cubic term \(2F\), it is exactly \(2\mathcal V_F(L)\).
For \(J,M\) this is
\[
 \frac2{27}\mathbb Z\ \oplus\frac{\sqrt2}{108}\mathbb Z
                           \ \oplus\frac{\sqrt3}{9}\mathbb Z.
\]
Multiplication by two is an injective additive map of the real
coefficient field and restricts to the exact isomorphism
$\mathcal V_F(L)\to\mathcal V_{2F}(L)=2\mathcal V_F(L)$,
$x\mapsto2x$, with inverse $y\mapsto y/2$.
It intertwines the same two inclusion matrices and
the same quotients; no factor has been absorbed into \(F\).

\subsection{Reproduction and scope of the exact certificate}

The accompanying evidence folder is
\begin{flushleft}\small
\texttt{supporting\_materials/workbench/evidence/}\\
\texttt{period\_lattice\_continuation\_20260906/jm\_cubic/}
\end{flushleft}
It contains both full source files \texttt{cubic\_engine.py} and
\texttt{check\_jm\_values.py}. Running the latter with Python's
standard library performs the following finite proof certificate:
all \(20399\) polarization entries for \(N,\Lambda,J,M\), all
displayed denominator and determinant assertions, the integer
Bezout and adjugate axis reconstructions, and the exact inclusion
coefficient comparisons. The output \texttt{EXACT\_RESULTS.json}
retains each selected column, determinant, Bezout coefficient,
and every expanded integer linear combination reconstructing a
value-group axis.

There is additionally an independent coefficient-identity check:
the original 24 simple roots form a real basis of the original
24-dimensional space. The direct Cayley--Dickson implementation
and the rational polynomial implementation are compared on the
24 basis vectors, their 552 ordered-by-index plus/minus pairs,
and their 2024 three-element sums, a total of 2600 exact inputs.
The finite polarization proof shows these evaluations determine
every coefficient of a homogeneous cubic in those 24 variables.
Consequently this is a full polynomial-identity certificate on
the real ambient space, not sampling of lattice points. The
field implementation retains all four rational coefficients in
\(\mathbb Q(\sqrt2,\sqrt3)\), and verifies that the \(\sqrt6\)
coefficient is identically zero. A further 65 direct checks cover
the displayed new generating lists and \(u,p,q\).

This computation proves the stated additive value groups and
their morphisms. It asserts neither equality of individual-value
sets nor that the cubic group reconstructs a lattice or classifies
all cubic-preserving lattice isometries.

% Complete proof module: Wilson construction and retained-coordinate comparison.
\section{A cyclic Wilson Leech lattice in the original triality space}
\label{wcl:sec:construction}

\subsection{The two retained metrics and the octonion convention}

Let $V=\mathbb O_1\oplus\mathbb O_2\oplus\mathbb O_3$ be the original
ordered octonion space. Retain its metric and cubic
\[
\begin{aligned}
 h_0(x,x)&=N(x_1)+N(x_2)+N(x_3),\\
 \tau(x_1,x_2,x_3)&=\operatorname{Re}((x_1x_2)x_3),\qquad
 d(Q_{\mathbb O}(0;x))=2\tau(x_1,x_2,x_3),
\end{aligned}
\]
where
\[
 Q_{\mathbb O}(\lambda;x)=
 \begin{pmatrix}\lambda&x_3&\bar x_2\\
 \bar x_3&\lambda&x_1\\x_2&\bar x_1&\lambda\end{pmatrix}.
\]
Wilson's norm on the same underlying ordered space is
\[
 q_{\rm W}(x)=\tfrac12 h_0(x,x).
\]
The map
\begin{equation}
 \iota:(V,q_{\rm W})\longrightarrow(V,h_0),\qquad
 \iota(x_1,x_2,x_3)=(x_1/\sqrt2,x_2/\sqrt2,x_3/\sqrt2)
 \label{wcl:eq:metric-isometry}
\end{equation}
is an explicit isometry. Both metrics remain specified. In particular,
for the original $\lambda$, which is not rescaled,
\begin{equation}
 d(Q_{\mathbb O}(\lambda;\iota x))
 =\lambda^3-\lambda q_{\rm W}(x)+\frac1{\sqrt2}\tau(x_1,x_2,x_3).
 \label{wcl:eq:cubic-scaling}
\end{equation}
Indeed the quadratic coordinates scale by $1/2$, while the trilinear
term scales by $(1/\sqrt2)^3=1/(2\sqrt2)$ and its original coefficient
is $2$. Thus the off-diagonal cubic after $\iota$ is
$1/(2\sqrt2)$ times its value before $\iota$.

For an explicit realization inside the same $\mathbb O$, choose a
Cayley--Dickson basis
$(1,\mathsf i,\mathsf j,\mathsf k,\ell,\mathsf i\ell,
\mathsf j\ell,\mathsf k\ell)$, with
\[
 (a+b\ell)(c+d\ell)=(ac-\bar d b)+(da+b\bar c)\ell,\qquad
 \mathsf i\mathsf j=\mathsf k .
\]
Wilson's ordered orthonormal basis is realized by
\begin{equation}
 (i_\infty,i_0,i_1,i_2,i_3,i_4,i_5,i_6)
 =(1,\mathsf i,\mathsf j,\ell,\mathsf k,
                  \mathsf j\ell,-\mathsf k\ell,\mathsf i\ell).
 \label{wcl:eq:actual-basis-map}
\end{equation}
The seven oriented Fano triples are $(t,t+1,t+3)$ modulo $7$:
$i_ti_{t+1}=i_{t+3}$, with cyclic products within each triple and
opposite signs on reversal. To verify \eqref{wcl:eq:actual-basis-map},
its seven corresponding products, in order $t=0,\ldots,6$, are
\[
 \mathsf i\mathsf j=\mathsf k,\quad
 \mathsf j\ell=\mathsf j\ell,\quad
 \ell\mathsf k=-\mathsf k\ell,\quad
 \mathsf k(\mathsf j\ell)=\mathsf i\ell,\quad
 (\mathsf j\ell)(-\mathsf k\ell)=\mathsf i,\quad
 (-\mathsf k\ell)(\mathsf i\ell)=\mathsf j,\quad
 (\mathsf i\ell)\mathsf i=\ell.
\]
Alternativity supplies cyclic products on each quaternionic line.
The seven lines cover all pairs of different imaginary basis units;
squares are $-1$ and $1$ is the unit. This proves the algebra
isomorphism on the full multiplication table, hence by bilinearity on
all octonions. It also preserves conjugation and the norm. No
conjugation of a coordinate of $Q_{\mathbb O}$ is performed.

\subsection{The explicit eight-dimensional lattice and two sublattices}

In the displayed Wilson basis write $e_0=i_\infty,e_{t+1}=i_t$ and put
\[
\begin{aligned}
 D_8&=\{(z_0,\ldots,z_7)\in\mathbb Z^8:\textstyle\sum z_i\equiv0\pmod2\},\\
 s&=\tfrac12(-1+i_0+i_1+i_2+i_3+i_4+i_5+i_6),\qquad
 L_{\rm W}=D_8+\mathbb Zs.
\end{aligned}
\]
This is Wilson's lattice $L$, renamed $L_{\rm W}$ to distinguish it
from the project's original rank-24 Niemeier lattice. Since $2s\in
D_8$ and $s\notin D_8$, it is exactly $D_8\cup(s+D_8)$.

\begin{lemma}\label{wcl:lem:e8}
The lattice $L_{\rm W}$ is even and unimodular, has minimum norm $2$,
and has the explicit basis
\[
 b_1=s,\quad b_j=e_{j-1}-e_j\ (2\le j\le7),\quad b_8=e_6+e_7.
\]
Its Gram matrix for the unchanged octonion norm is
\begin{equation}
 G_8=
 \begin{pmatrix}
 2&0&0&0&0&0&0&1\\
 0&2&-1&0&0&0&0&0\\
 0&-1&2&-1&0&0&0&0\\
 0&0&-1&2&-1&0&0&0\\
 0&0&0&-1&2&-1&0&0\\
 0&0&0&0&-1&2&-1&-1\\
 0&0&0&0&0&-1&2&0\\
 1&0&0&0&0&-1&0&2
 \end{pmatrix},\qquad \det G_8=1 .
 \label{wcl:eq:e8-gram}
\end{equation}
Its norm-two vectors are precisely the $112$ vectors $\pm e_i\pm e_j$
with $i<j$ and the $128$ vectors
$\frac12(\pm e_0\pm\cdots\pm e_7)$ with an odd number of minus signs.
\end{lemma}

\begin{proof}
For $z\in D_8$, $N(z)\equiv\sum z_i=0\pmod2$ and
$h(s,z)=\frac12\sum z_i-z_0\in\mathbb Z$.
As $N(s)=2$, every norm in $D_8+\mathbb Zs$ is even.
The covolume of $D_8$ is $2$, and its displayed extension has index
$2$, so $L_{\rm W}$ has covolume $1$. The integral bilinear form follows
by polarizing the even norm, and covolume one proves self-duality.
The columns $b_i$ all belong to $L_{\rm W}$. Their determinant has
absolute value one: expansion along the first coordinate gives
$1/2$ times the determinant $2$ of the indicated $D_7$ basis.
They therefore form a basis. Their scalar products give the displayed
matrix. Positivity and evenness give minimum at least $2$, and
$N(b_i)=2$ attains it. Integer vectors of norm two have exactly two
nonzero entries $\pm1$. Half-integer vectors have norm at least
$8/4=2$, with equality exactly when all eight entries are $\pm1/2$.
Their membership in $s+D_8$ is exactly the stated odd-minus parity.
\end{proof}

Define, with the bars in the positions printed in Wilson,
\[
 \mathsf M=L_{\rm W}\bar s,\qquad \mathsf N=L_{\rm W}s .
\]
All products are right products. Let $B$ be the basis matrix with
columns $b_i$. Right multiplication by $s$ in this basis is the
following integer matrix:
\begin{equation}
 S=B^{-1}R_s B=
 \begin{pmatrix}
 3&0&0&0&0&0&0&2\\
 -2&-1&0&1&-1&1&0&-2\\
 -4&-1&-1&1&0&0&1&-3\\
 -6&-1&-1&0&0&1&0&-4\\
 -8&0&-1&0&-1&1&1&-5\\
 -10&-1&0&0&-1&0&1&-5\\
 -5&0&-1&1&-1&0&0&-2\\
 -7&0&0&0&0&0&0&-4
 \end{pmatrix}.
 \label{wcl:eq:S}
\end{equation}
This is a direct multiplication-table calculation using
\eqref{wcl:eq:actual-basis-map}; it can be checked by multiplying each
column $b_j$ by the displayed $s$. Set $T=-I_8-S$. Since
$s+\bar s=-1$, $s\bar s=2$, and alternativity allows repeated $s$
and $\bar s= -1-s$, one has
\begin{equation}
 T=B^{-1}R_{\bar s}B,\qquad S+T=-I_8,\qquad ST=TS=2I_8,
 \qquad S^{\mathsf T}G_8S=T^{\mathsf T}G_8T=2G_8 .
 \label{wcl:eq:ST-identities}
\end{equation}

\begin{lemma}\label{wcl:lem:complementary}
Both $\mathsf M,\mathsf N$ are sublattices of $L_{\rm W}$ of index
$16$ and minimum norm $4$, and
\[
 2L_{\rm W}\subset\mathsf M,\mathsf N\subset L_{\rm W},\qquad
 \mathsf M+\mathsf N=L_{\rm W},\qquad
 \mathsf M\cap\mathsf N=2L_{\rm W}.
\]
Their images $\overline{\mathsf M},\overline{\mathsf N}$ in
$\mathcal E=L_{\rm W}/2L_{\rm W}$ are complementary
four-dimensional $\mathbb F_2$ subspaces.
\end{lemma}

\begin{proof}
The integer matrices $S,T$ give the inclusions in $L_{\rm W}$;
$ST=2I$ gives the inclusions of $2L_{\rm W}$.
Norm composition shows that each right multiplication scales all
squared lengths by $N(s)=N(\bar s)=2$; thus its determinant has
absolute value $(\sqrt2)^8=16$, proving the indices and minimum norms.
The identity $S+T=-I$ gives $\mathsf M+\mathsf N=L_{\rm W}$.
Modulo $2L_{\rm W}$, each image has size $2^8/16=2^4$.
Their sum is the eight-dimensional space $\mathcal E$, so their
intersection is zero. This proves the asserted intersection upstairs.
\end{proof}

\subsection{The binary glue code and its exact generators}

Wilson's original lattice is exactly
\begin{equation}
\begin{split}
 \Lambda_{\rm W}=\{(x,y,z)\in L_{\rm W}^3:\ &
 x+y,x+z,y+z\in\mathsf M,\quad x+y+z\in\mathsf N\}.
\end{split}
 \label{wcl:eq:Wilson-definition}
\end{equation}
It carries Wilson's norm $q_{\rm W}$, as printed. Define the lattice
inside the original triality metric by
\[
 \Lambda_{\rm cyc}=\iota(\Lambda_{\rm W})\subset(V,h_0).
\]
The simultaneous use of $\mathsf M=L_{\rm W}\bar s$ for pair sums
and $\mathsf N=L_{\rm W}s$ for the total sum is essential.

Write any class in $\mathcal E$ uniquely as $m+n$, with
$m\in\overline{\mathsf M}$ and $n\in\overline{\mathsf N}$.
Then the image of \eqref{wcl:eq:Wilson-definition} in $\mathcal E^3$
is the explicit binary code
\begin{equation}
 \mathcal C=
 \{(m_1+n,m_2+n,m_1+m_2+n):
       m_1,m_2\in\overline{\mathsf M},\
       n\in\overline{\mathsf N}\}.
 \label{wcl:eq:code}
\end{equation}
Indeed the pair conditions say that the three $\mathsf N$ components
are equal; the total-sum condition says that the sum of the three
$\mathsf M$ components is zero. This proves both directions of the
code description. Its dimension is $4+4+4=12$.

For a concrete lift, the first four columns of $S$ and of $T$ are
bases of their images modulo $2$. The determinant of the eight-column
matrix $(T_1,T_2,T_3,T_4,S_1,S_2,S_3,S_4)$ is $-7$, so these columns
are independent modulo $2$; each of the two images already has
dimension four. Put $m_j=b_j\bar s$, $n_j=b_js$ for $1\le j\le4$.
A complete set of $36$ integer generators of $\Lambda_{\rm W}$ is
\begin{equation}
\begin{gathered}
 (2b_i,0,0),\ (0,2b_i,0),\ (0,0,2b_i)\quad(1\le i\le8),\\
 (m_j,m_j,0),\ (m_j,0,m_j),\ (n_j,n_j,n_j)\quad(1\le j\le4).
\end{gathered}
 \label{wcl:eq:generators}
\end{equation}
These vectors are all in \eqref{wcl:eq:Wilson-definition}.
Their reductions span precisely \eqref{wcl:eq:code}; the first
$24$ generate its full kernel $(2L_{\rm W})^3$. This proves that
they generate the entire lattice, not only a sublattice.

\begin{theorem}\label{wcl:thm:Leech}
The explicitly generated lattice $\Lambda_{\rm W}$ with $q_{\rm W}$,
and therefore $\Lambda_{\rm cyc}$ with the unaltered $h_0$, is an
even unimodular rank-24 lattice of minimum norm $4$. In particular
it is rootless. This is the actual Wilson Leech construction, through
the explicit isometry \eqref{wcl:eq:metric-isometry}.
\end{theorem}

\begin{proof}
The code proves
$[L_{\rm W}^3:\Lambda_{\rm W}]=2^{24-12}=2^{12}$.
Since $L_{\rm W}$ has covolume one for its original octonion metric,
$\Lambda_{\rm W}$ has covolume $2^{12}$ for $h_0$.
Changing to the separately recorded metric $q_{\rm W}$ multiplies
volume in dimension $24$ by $(1/\sqrt2)^{24}=2^{-12}$,
so its covolume is exactly one.

For every triple of vectors in a real inner product space,
\[
 N(x)+N(y)+N(z)
 =N(x+y)+N(x+z)+N(y+z)-N(x+y+z).
\]
For a vector of $\Lambda_{\rm W}$ the four terms on the right
are norms in $\mathsf M$ or $\mathsf N$, so each is divisible by
$4$: these are twice norms in the even lattice $L_{\rm W}$.
It follows that $q_{\rm W}(x,y,z)$ is even. Polarization of these
even norms gives an integral bilinear form. An integral lattice of
covolume one equals its dual, proving unimodularity.

If a nonzero vector had $q_{\rm W}=2$, its three coordinate norms
would sum to $4$. Each nonzero coordinate in $L_{\rm W}$ has norm
at least $2$, so one coordinate is zero. The other two then lie in
$\mathsf M$, each of whose nonzero vectors has norm at least $4$.
At most one can be nonzero. That coordinate also lies in $\mathsf N$
by the total-sum condition, so it lies in
$\mathsf M\cap\mathsf N=2L_{\rm W}$ and has norm at least $8$,
a contradiction. Hence all nonzero norms are at least $4$.
The vectors $(2b_i,0,0)$ have $q_{\rm W}=4$, so the minimum is
exactly $4$. The isometry $\iota$ transfers every statement to the
original metric.
\end{proof}

\subsection{Cyclic symmetry, the fixed lattice, and exact three-primary glue}

Retain the canonical cycle $c(x,y,z)=(z,x,y)$. Denote the opposite
cycle used in Wilson coordinates by $c_{\rm W}(x,y,z)=(y,z,x)$.
Direct substitution gives
\[
 c_{\rm W}=c^{-1}=c^2,\qquad c_{\rm W}^2=c,\qquad
 \langle c_{\rm W}\rangle=\langle c\rangle.
\]
Both actions are retained with these directions; no action is renamed
by silently reversing it.
The defining three coordinate memberships, three pair conditions,
and total-sum condition are all unchanged by $c_{\rm W}$; hence both cycles preserve
$\Lambda_{\rm W}$ and $\Lambda_{\rm cyc}$. They preserve their metrics
and the original cubic:
\[
 \operatorname{Re}((xy)z)=\operatorname{Re}((yz)x).
\]
This follows from the octonion adjoint identities
$h(ab,d)=h(b,\bar a d)=h(a,d\bar b)$, which imply real associativity,
and $\operatorname{Re}(ab)=\operatorname{Re}(ba)$.
In the Albert algebra $c_{\rm W}$ is realized by conjugation with the
\emph{real} permutation matrix for $(1\,3\,2)$.
It sends the coordinates $(x_1,x_2,x_3)$ to $(x_2,x_3,x_1)$ with no
conjugations, and permutes the ordered diagonal frame accordingly.
The inverse real permutation matrix, for $(1\,2\,3)$, realizes the
canonical $c$ and sends $(x_1,x_2,x_3)$ to $(x_3,x_1,x_2)$.
Thus the fixed lattices of the two cycles are identical. For the
retained residue isometry $\mathcal Lr=c\mathcal L$, the opposite
action satisfies $\mathcal Lr^{-1}=c_{\rm W}\mathcal L$.

\begin{theorem}\label{wcl:thm:fixed-coinvariant}
Before applying $\iota$, the cyclic fixed lattice and its orthogonal
complement in $\Lambda_{\rm W}$ are exactly
\[
 F=\{(n,n,n):n\in\mathsf N\},\qquad
 C=\{(m_1,m_2,-m_1-m_2):m_1,m_2\in\mathsf M\}.
\]
Their ranks are $8$ and $16$. For the retained Wilson metric, and
therefore after $\iota$ for the original triality metric, their Gram
matrices in the displayed bases are
\begin{equation}
 G_F=3G_8,\qquad
 G_C=\begin{pmatrix}2G_8&G_8\\G_8&2G_8\end{pmatrix}.
 \label{wcl:eq:fixed-coinvariant-grams}
\end{equation}
The fixed-lattice basis is $(b_js,b_js,b_js)$, $1\le j\le8$;
the coinvariant bases are $(b_j\bar s,0,-b_j\bar s)$ and
$(0,b_j\bar s,-b_j\bar s)$. In particular the fixed lattice has
minimum norm $6$ and determinant $3^8$.
The coinvariant lattice has determinant $3^8$, and the map which
negates its second $L_{\rm W}$ parameter identifies its displayed
Gram matrix with the usual $A_2\otimes E_8$ Gram matrix.
\end{theorem}

\begin{proof}
A fixed triple has the form $(x,x,x)$. Its pair condition holds for
every $x\in L_{\rm W}$ because $2L_{\rm W}\subset\mathsf M$.
Its total condition is $3x\in\mathsf N$. As $L_{\rm W}/\mathsf N$
has exponent two, $3x\in\mathsf N$ is equivalent to $x\in\mathsf N$.
This proves the first description. Orthogonality to the whole
diagonal real eight-space says $x+y+z=0$. The pair conditions are
then exactly $x,y,z\in\mathsf M$, and the total condition is automatic.
This proves the description of $C$.

Right multiplication by either $s$ or $\bar s$ multiplies the
octonion metric by $2$. Thus
\[
 q_{\rm W}(as,as,as)=3N(a),
\]
and
\[
 q_{\rm W}(a\bar s,b\bar s,-(a+b)\bar s)
 =N(a)+N(b)+N(a+b)=2N(a)+2N(b)+2h(a,b).
\]
Polarization gives both full Gram matrices. Their determinants are
\[
 \det G_F=3^8\det G_8=3^8,\qquad
 \det G_C=\det\begin{pmatrix}2&1\\1&2\end{pmatrix}^{8}
 (\det G_8)^2=3^8.
\]
Minimum norm $6$ for $F$ follows from the exact
factor $3$ and the minimum $2$ in $L_{\rm W}$.
The sign map $b\mapsto-b$ changes both off-diagonal blocks from
$G_8$ to $-G_8$, proving the stated exact comparison to
$A_2\otimes E_8$.
\end{proof}

There is a complete integral comparison, including the glue:
\begin{equation}
 0\longrightarrow F\oplus C\longrightarrow\Lambda_{\rm W}
 \xrightarrow{\ \Sigma\ }\mathsf N/3\mathsf N
 \longrightarrow0,\qquad
 \Sigma(x,y,z)=x+y+z\pmod{3\mathsf N}.
 \label{wcl:eq:three-primary-exact}
\end{equation}
To prove surjectivity, for every $a\in L_{\rm W}$ the explicit vector
\[
 v(a)=(as,a,-a)
\]
lies in $\Lambda_{\rm W}$ and has sum $as$. Its pair sums are
$a(s+1)=-a\bar s\in\mathsf M$,
$a(s-1)=-a\bar s-2a\in\mathsf M$, and $0$;
its total lies in $\mathsf N$. If a lattice vector has total $3n$,
subtracting $(n,n,n)\in F$ leaves a vector of $C$. This proves the
kernel assertion. Thus
\[
 \Lambda_{\rm W}/(F\oplus C)\cong\mathsf N/3\mathsf N
 \cong L_{\rm W}/3L_{\rm W}\cong(\mathbb Z/3)^8,\qquad
 [\Lambda_{\rm W}:F\oplus C]=3^8.
\]
All maps commute with $c_{\rm W}$ and with its inverse $c$; both
cycles act trivially on the quotient.

The discriminant anti-isometry is also explicit. Use the quadratic
discriminant convention $q_D(u+D)=h_D(u,u)\pmod{2\mathbb Z}$
for an even lattice $D$ with its retained metric. The orthogonal
projections of the lift $v(a)$ are
\begin{align}
 f(a)&=\tfrac13(as,as,as),\nonumber\\
 \gamma(a)&=\left(\tfrac23as,\ a-\tfrac13as,\ -a-\tfrac13as\right).
 \label{wcl:eq:glue-projections}
\end{align}
They belong to $F^*$ and $C^*$, respectively: their pairings with
$F$ or $C$ are the pairings of the integral vector $v(a)$ with the
same lattice, because the other projection is orthogonal.
The classes $f(a)+F$ identify $L_{\rm W}/3L_{\rm W}$ with $F^*/F$.
If $\gamma(a)\in C$, then $f(a)=v(a)-\gamma(a)\in
\Lambda_{\rm W}\cap F_{\mathbb R}=F$, so $a\in3L_{\rm W}$.
Since $C^*/C$ has order $\det G_C=3^8$, the map
\[
 f(a)+F\longmapsto\gamma(a)+C
\]
is an isomorphism of the two discriminant groups. Its quadratic
sign is determined without an unspecified convention:
\[
 q_F(f(a))=\tfrac13N(a),\qquad
 q_C(\gamma(a))=q_{\rm W}(v(a))-\tfrac13N(a)
              =2N(a)-\tfrac13N(a).
\]
As $N(a)$ is even, their sum $2N(a)$ is zero modulo $2\mathbb Z$.
This is the required anti-isometry. Equations
\eqref{wcl:eq:three-primary-exact} and
\eqref{wcl:eq:glue-projections} specify the overlattice itself,
including its actual coset representatives.

\subsection{Finite lattice symmetries and exact cubic tests}

Independent sign changes of the three sectors preserve
$\Lambda_{\rm W}$. Changing the sign of $x$ changes a pair sum by
$-2x\in2L_{\rm W}\subset\mathsf M$ and changes the total by
$-2x\in2L_{\rm W}\subset\mathsf N$. The same argument applies to
each sector. All six sector permutations preserve the defining
conditions too. Under signs $(\epsilon_1,\epsilon_2,\epsilon_3)$
the cubic $\tau$ is multiplied by $\epsilon_1\epsilon_2\epsilon_3$;
hence exactly the even sign changes in this subgroup preserve it.
Cyclic permutations preserve $\tau$ as already proved.

The odd permutation $(x,y,z)\mapsto(y,x,z)$ need not preserve it,
even on the lattice itself. In Wilson's basis put
\[
 a=2(1+i_0),\qquad b=2(1+i_1),\qquad d=2(i_2+i_3).
\]
Every coordinate belongs to $2L_{\rm W}$, so $(a,b,d)\in
\Lambda_{\rm W}$. The real part of
$((1+i_0)(1+i_1))(i_2+i_3)$ is $-1$, because
$i_0i_1=i_3$ and the $i_3^2$ term is the only real contribution.
Reversing the first two factors changes this contribution to $+1$.
Consequently
\begin{equation}
 \tau(a,b,d)=-8,\qquad \tau(b,a,d)=8.
 \label{wcl:eq:transposition-cubic}
\end{equation}
After the specified scalar isometry the original determinant cubic
has values $-4\sqrt2$ and $4\sqrt2$.

The simultaneous Fano-index operations
\[
 A:i_t\longmapsto i_{t+1},\qquad
 B_0:i_t\longmapsto i_{2t},\qquad i_\infty\longmapsto i_\infty
\]
are octonion algebra automorphisms, as the defining oriented Fano
products show. They are orthogonal permutations preserving $D_8$ and
the odd-minus half coset, and they fix the displayed $s$ and $\bar s$.
They consequently preserve $\mathsf M,\mathsf N,\Lambda_{\rm W}$.
Applied simultaneously to all three sectors, they preserve $\tau$.
They satisfy $A^7=B_0^3=1$ and $B_0AB_0^{-1}=A^2$; their permutations
are the distinct maps $t\mapsto 2^j t+k$ with $j=0,1,2$ and
$k\in\mathbb F_7$. Thus they form an explicitly realized
$C_7\rtimes C_3$ of order $21$.

The even sign changes and the cycle $c_{\rm W}$ form
$(\mathbb Z/2)^2\rtimes C_3\cong A_4$. They commute with the
simultaneous Fano operations, and their intersection is trivial:
a simultaneous Fano operation fixes the real unit in every sector
and does not permute sectors. Hence an explicit finite group
\[
 A_4\times(C_7\rtimes C_3)
\]
of order $252$ preserves both the literal lattice and the original
cubic. This is a proved subgroup; no assertion that it is the full
cubic stabilizer is needed.

Wilson also gives the lattice symmetries
\[
 r_t(x,y,z)=(x,yi_t,zi_t)\qquad(t=0,\ldots,6).
\]
For completeness their lattice action can be checked solely from
the explicit finite matrices above. Let
$U_t=B^{-1}R_{i_t}B$. The signed permutation multiplication table
and \eqref{wcl:eq:S} give
\begin{equation}
 U_t\in GL_8(\mathbb Z),\qquad
 T^{-1}(U_t-I_8)\in M_8(\mathbb Z),\qquad
 \tfrac12(U_t-I_8)T\in M_8(\mathbb Z)
 \quad(0\le t\le6).
 \label{wcl:eq:unit-matrix-certificates}
\end{equation}
These are finite exact matrix assertions, not assumptions about
octonionic matrix associativity. The full multiplication rule and
the exact rational algorithm checking every entry for all seven
values are included in the accompanying \texttt{check\_wilson.py}.
Concretely, start with the $64$ products
$1i_t=i_t1=i_t$, $i_t^2=-1$, and the seven oriented Fano triples,
form the eight columns $R_{i_t}(e_j)$, conjugate by the explicitly
displayed $B$, and substitute the integer $T=-I_8-S$.
For each of the three matrices in \eqref{wcl:eq:unit-matrix-certificates}
the reduced denominators of every entry are $1$, for each
$t=0,1,2,3,4,5,6$. Also $\det U_t=1$.
This is a reproducible $7\cdot3\cdot64$-entry certificate.
The first condition is also immediate from the signed permutation:
right multiplication by $i_t$ flips four coordinate signs, preserving
the odd-minus parity of $L_{\rm W}$.

The second matrix assertion says
$L_{\rm W}(i_t-1)\subset\mathsf M$, and the third says
$\mathsf M(i_t-1)\subset2L_{\rm W}$. Thus $x+yi_t$ and $x+zi_t$
differ from the original pair sums by elements of $\mathsf M$.
The pair $(y+z)i_t$ remains in $\mathsf M$, and the transformed total
differs from $x+y+z$ by
$(y+z)(i_t-1)\in2L_{\rm W}\subset\mathsf N$.
All defining conditions are preserved. The inverse follows from
$r_t^2(x,y,z)=(x,-y,-z)$ and the proved sign symmetries.
Norm composition proves that $r_t$ is an isometry.

Nevertheless none of these seven $r_t$ preserves the original cubic.
For $a_t=2(1+i_t)$, the triple $(a_t,a_t,a_t)$ lies in
$(2L_{\rm W})^3\subset\Lambda_{\rm W}$. All products below occur in
the associative real subalgebra generated by $i_t$, and give
\[
 \tau(a_t,a_t,a_t)=8\operatorname{Re}((1+i_t)^3)=-16,
\]
whereas $(1+i_t)i_t=-1+i_t$ and
\[
 \tau(a_t,a_ti_t,a_ti_t)
 =8\operatorname{Re}((1+i_t)(-1+i_t)^2)=16.
\]
Therefore the original determinant cubic on the isometric lattice
$\Lambda_{\rm cyc}$ changes from $-8\sqrt2$ to $8\sqrt2$.
This supplies actual lattice vectors as counterexamples, not only
ambient real vectors outside the lattice.

\subsection{Exact comparison with coordinate lattices and the marked lane}

The construction has the explicit integral sandwich
\begin{equation}
 (\sqrt2L_{\rm W})^3
 \ \subset\ \Lambda_{\rm cyc}\
 \subset\ (L_{\rm W}/\sqrt2)^3,
 \qquad
 [\Lambda_{\rm cyc}:(\sqrt2L_{\rm W})^3]
 =[(L_{\rm W}/\sqrt2)^3:\Lambda_{\rm cyc}]=2^{12}.
 \label{wcl:eq:coordinate-sandwich}
\end{equation}
This is \eqref{wcl:eq:code} transported by the exact isometry
\eqref{wcl:eq:metric-isometry}. The inclusions are literal inclusions
in the fixed three octonion sectors, and each commutes with both
$c_{\rm W}$ and the canonical $c=c_{\rm W}^{-1}$.
Together with the explicit bases, the code, and the three-primary
glue \eqref{wcl:eq:glue-projections}, it determines actual
commensurability maps and cosets rather than only an abstract
isometry class.

Under \eqref{wcl:eq:actual-basis-map}, the single negative column
takes the Wilson odd-minus half roots to even-minus half roots in
the current ordered Cayley--Dickson basis. The element $s$ remains
exactly
\[
 s=\tfrac12(-1+\mathsf i+\mathsf j+\mathsf k+\ell
                     +\mathsf i\ell+\mathsf j\ell-\mathsf k\ell),
\]
with coordinate vector $(-1,1,1,1,1,1,1,-1)/2$.
The signed permutation has determinant $+1$ and all $64$ basis
products agree, as independently checked in the accompanying code.
Thus the lattice, its cyclic action, and the triality cubic live
in one specified multiplication convention.

\subsection{Literal intersections with both retained lattice embeddings}
\label{wcl:sec:literal-intersections}

Let $U=\operatorname{span}_{\mathbb R}(1,\mathsf i,\mathsf j,\mathsf k)$
be the first four coordinates in the \emph{current} ordered
Cayley--Dickson basis, and let
\[
 D_4=\{a=(a_0,a_1,a_2,a_3)\in\mathbb Z^4:
                          a_0+a_1+a_2+a_3\in2\mathbb Z\}\subset U.
\]
This quaternionic $D_4$ denotes the displayed four-coordinate lattice;
the original Niemeier $D$ summand below remains in its original place.

Let $J_{\rm W}$ denote the signed permutation matrix in
\eqref{wcl:eq:actual-basis-map}. If $U_0:\mathbb R^4\to\mathbb R^8$
is the coordinate inclusion in the current basis, then
\begin{equation}
 S^{-1}B^{-1}J_{\rm W}^{-1}U_0=
 \begin{pmatrix}
 1/2&-1/2&-1/2&-1/2\\
 -1/2&0&0&1/2\\
 -1&1/2&0&1/2\\
 -3/2&1&1/2&1/2\\
 -2&1&1&1/2\\
 -5/2&3/2&1&1\\
 -5/4&3/4&3/4&3/4\\
 -7/4&3/4&3/4&3/4
 \end{pmatrix}.
 \label{wcl:eq:quaternion-membership-matrix}
\end{equation}
All matrices on the left have already been explicitly specified.
Thus the equation is a direct rational multiplication certificate.

\begin{lemma}\label{wcl:lem:quaternion-intersection}
In this fixed coordinate frame,
\[
 \mathsf N\cap U=2D_4 .
\]
\end{lemma}
\begin{proof}
Since $\mathsf N\subset L_{\rm W}$, an element of $\mathsf N\cap U$
has integral first four coordinates: the half-integer coset of the
current $E_8$ description has no zero coordinate among its last four.
Write those integral coordinates as $k_0,k_1,k_2,k_3$. Membership in
$\mathsf N$ is exactly integrality of the matrix
\eqref{wcl:eq:quaternion-membership-matrix} applied to $k$.
Subtracting its last two rows gives $k_0/2$, so $k_0$ is even.
The second row then forces $k_3$ even, the third forces $k_1$ even,
and the fourth forces $k_2$ even. Write $k=2a$.
The first six rows are now integral. The seventh and eighth rows
are integral precisely when
$-5a_0+3a_1+3a_2+3a_3$ is even, equivalently
$a_0+a_1+a_2+a_3$ is even. Conversely these conditions make all eight
rows integral. This proves both inclusions with no unstated lattice
identification.
\end{proof}

We now retain the original lattices and their actual map $\mathcal L$.
The original coordinates are $X=(x_{i,j};d_k)$,
$1\le i\le4$, $0\le j\le5$, $1\le k\le4$, with
$\sum_jx_{i,j}=0$. The root lattice is
$R=A_5^{\perp4}\perp D$, with its unchanged Euclidean metric.
Its index-$72$ overlattice $N$ is the inverse image of the original
glue
\[
 (k_1,k_2,k_3,k_4;y)
   =\bigl(3u+2(a+b,a+2b,a,b);\,(u_2+u_4,u_3+u_4)\bigr),
\]
where $u\in\mathbb F_2^4$ has even weight and $a,b\in\mathbb F_3$;
the first four entries are modulo $6$. Here $k_i=-6x_{i,0}\pmod6$
and $y$ is the original $D^\vee/D$ marking, so $d=0$ has $y=00$.
Retain the original neighbour
\[
 \rho_A=(5,3,1,-1,-3,-5)/2,\quad
 \rho_D=(3,2,1,0),\quad \rho=\rho_A^{\oplus4}\perp\rho_D,
\]
\[
 K=\{X\in N:(\rho,X)\in6\mathbb Z\},\qquad
 \alpha=(e_0-e_1,0,0,0;0),\qquad
 \Lambda=K+\mathbb Z(\rho/6-\alpha).
\]
Thus $\rho\in N$, $\rho^2=84$, $(\rho,\alpha)=1$.
Every vector of $N$ or $\Lambda$ has rational original coordinates.

For exact use of the retained isometry, set
\[
 \sigma_i=\sum_{j=0}^5(-1)^jx_{i,j},\qquad
 t_{i,p,j}=x_{i,p+2j}-\frac13\sum_{h=0}^2x_{i,p+2h}.
\]
Write $u_0=1,u_1=\mathsf i,u_2=\mathsf j,u_3=\mathsf k$ and retain
the remaining current Cayley--Dickson basis vectors $u_4,\ldots,u_7$.
The already specified map has the exact coordinate form
\[
 (\mathcal LX)_{j+1,a}=
 \begin{cases}
 t_{i,p,j}+\sqrt2\,\sigma_{a+1}/6,&0\le a\le3,\\
 t_{i,p,j}+\sqrt3\,d_{a-3}/3,&4\le a\le7,
 \end{cases}
 \qquad a=2(i-1)+p.
\]
The placement of $\sigma_{a+1}$ is retained, including its difference
from the index $i$ of the residual. No coordinate is reassigned.

\begin{theorem}\label{wcl:thm:literal-intersection}
The literal intersections in the original metric space are exactly
\begin{equation}
 \Lambda_{\rm cyc}\cap\mathcal L(N)
 =\Lambda_{\rm cyc}\cap\mathcal L(\Lambda)
 =\mathcal I
 :=\{\sqrt2(a,a,a):a\in D_4\subset U\}.
 \label{wcl:eq:literal-common-lattice}
\end{equation}
In original Niemeier coordinates the inverse image of this common
lattice is
\[
 \mathcal L^{-1}(\mathcal I)
 =\{x_{i,j}=(-1)^j a_{i-1},\ d=0:\ a\in D_4\}\subset K.
\]
It has rank four, minimum norm $12$, and Gram matrix
\begin{equation}
 G_{\mathcal I}=6
 \begin{pmatrix}
 2&-1&0&0\\
 -1&2&-1&-1\\
 0&-1&2&0\\
 0&-1&0&2
 \end{pmatrix},
 \qquad\det G_{\mathcal I}=4\cdot6^4=5184.
 \label{wcl:eq:common-gram}
\end{equation}
An explicit basis is $\sqrt2(\delta_j,\delta_j,\delta_j)$,
where
$\delta_1=(1,-1,0,0)$, $\delta_2=(0,1,-1,0)$,
$\delta_3=(0,0,1,-1)$, $\delta_4=(0,0,1,1)$.
\end{theorem}

\begin{proof}
Each coordinate of $\Lambda_{\rm cyc}$ belongs to
$\mathbb Q/\sqrt2=\sqrt2\mathbb Q$ in the current frame.
For a rational original vector $X$, the first four coordinate
formulas above therefore force $t_{i,p,j}=0$ at the corresponding
positions, by linear independence of $1,\sqrt2$ over $\mathbb Q$.
The last four force both $t_{i,p,j}=0$ and $d_{a-3}=0$, by linear
independence of $1,\sqrt2,\sqrt3$. For the latter independence, a
relation with nonzero $\sqrt3$ coefficient would place $\sqrt3$ in
$\mathbb Q(\sqrt2)$; squaring $A+B\sqrt2$ forces $AB=0$ and then
would make $3$ or $3/2$ a rational square, a contradiction.
All residuals therefore vanish and $d=0$. The zero fibre sums give
\[
 x_{i,j}=(-1)^j\sigma_i/6.
\]

First let $X\in N$. Its $A_5^\vee$ integral-difference conditions
are exactly $k_i:=\sigma_i/3\in\mathbb Z$, yielding
$x_{i,j}=(-1)^jk_i/2$. Their discriminant labels are
$-3k_i=3k_i\pmod6$. Their three-primary parts are zero, so the
injectivity of $(a,b)\mapsto(a+b,a+2b,a,b)$ forces $a=b=0$.
Since $d=0$, the binary condition is
$u_2+u_4=u_3+u_4=0$ and $\sum u_i=0$. These equations give
$u_1=u_2=u_3=u_4$. Thus all four $k_i$ have the same parity.
Conversely these conditions satisfy every original glue requirement.
The image is the diagonal triple
\[
 \mathcal LX=(k/\sqrt2,k/\sqrt2,k/\sqrt2),\qquad k\in U.
\]
By the fixed-lattice description of Wilson's construction it
belongs to $\Lambda_{\rm cyc}$ exactly when $k\in\mathsf N$.
The previous lemma makes this equivalent to $k=2a$, $a\in D_4$,
which proves the first intersection equality.

For the second equality, a vector of $\Lambda$ can be written
$X=n+m\rho/6$ with $n\in N$, $m\in\mathbb Z$.
The already proved necessary condition $d_X=0$ gives
$d_n=-m(3,2,1,0)/6\in D^\vee$.
All coordinates of a vector of $D^\vee$ are integral or all
half-integral; the fourth coordinate is zero, so all are integral.
Its third coordinate forces $6\mid m$, and hence $X\in N$.
In fact $\Lambda\cap N=K$: in
$k+m(\rho/6-\alpha)\in N$ the pairing with the original root
$\alpha$ forces $6\mid m$, and
$(\rho,\rho-6\alpha)=84-6=78\in6\mathbb Z$.
Thus an intersection vector belongs to $K$.
For the diagonal-fibre vectors just found,
\[
 (\rho,X)=\frac32\sum_{i=1}^4 k_i.
\]
Consequently the condition $X\in K$ is $\sum k_i\equiv0\pmod4$.
Every $k=2a$ with $a\in D_4$ satisfies it, proving equality of
the two intersections and the displayed inverse image.

Finally
$h_0(\sqrt2(a,a,a),\sqrt2(b,b,b))=6h(a,b)$.
The four $\delta_j$ form a basis of $D_4$, since their determinant
has absolute value two and $D_4$ has index two in $\mathbb Z^4$.
Their scalar products give \eqref{wcl:eq:common-gram}.
The least nonzero norm in $D_4$ is two: a norm-one integer vector
has odd coordinate sum, while $(1,-1,0,0)$ has norm two.
This proves the exact minimum and determinant.
\end{proof}

The two slice indices are also literal. Define the common
four-dimensional real subspace
$\mathcal U=\{(u,u,u):u\in U\}$. The proof gives
\[
 \mathcal L(N)\cap\mathcal U
 =\{(k,k,k)/\sqrt2:k\in P\},\quad
 P=2\mathbb Z^4+\mathbb Z(1,1,1,1),
\]
\[
 \mathcal L(\Lambda)\cap\mathcal U
 =\{(k,k,k)/\sqrt2:k\in P_0\},\quad
 P_0=\{k\in P:\textstyle\sum k_i\equiv0\pmod4\}
     =2D_4+\mathbb Z(1,1,1,1).
\]
The first proof does not need Wilson membership to obtain these
slice equations. The lattice $P$ has basis
$2e_0,2e_1,2e_2,(1,1,1,1)$ and covolume $8$.
The lattice $2D_4$ has covolume $2^4\cdot2=32$, so
$[P:2D_4]=4$. The class of $(1,1,1,1)$ has exact order two
modulo $2D_4$, yielding $[P_0:2D_4]=2$. Therefore
\[
 [\mathcal L(N)\cap\mathcal U:\mathcal I]=4,\qquad
 [\mathcal L(\Lambda)\cap\mathcal U:\mathcal I]=2.
\]
The rank-four intersection, strictly below rank $24$, also proves
that these full lattices are not commensurable \emph{in the retained
embedding}: a common finite-index subgroup would have full rank.
Their exact existing relation is the common lattice with the
explicit embeddings, inverse map, metrics, and indices just proved.

\subsection{The cubic value group on the literal common lattice}
\label{wcl:sec:common-cubic}

Retain $F(y)=\tau(y_1,y_2,y_3)$, so the Albert off-diagonal
determinant is $2F$. For $y=\sqrt2(a,a,a)\in\mathcal I$ write
$a=a_0+\mathbf a$ in the same quaternionic $U$. Since
$\mathbf a^2=-N(\mathbf a)$ and $a_0$ is real, direct multiplication
inside this associative subalgebra gives
\[
 \operatorname{Re}(a^3)
 =a_0^3-3a_0N(\mathbf a)=4a_0^3-3a_0N(a).
\]
Every scalar and factor is therefore
\begin{equation}
 F(\sqrt2(a,a,a))
 =2\sqrt2\bigl(4a_0^3-3a_0N(a)\bigr).
 \label{wcl:eq:common-cubic}
\end{equation}
For $a\in D_4$, $N(a)$ is even, because
$N(a)\equiv\sum a_i=0\pmod2$. The integer in parentheses is
therefore even, so every value belongs to $4\sqrt2\mathbb Z$.
For the actual lattice vector $a=(1,1,0,0)$ it equals
$4-3\cdot2=-2$, giving $F=-4\sqrt2$.
Consequently the \emph{additive subgroup generated by all values}
is exactly
\[
 \bigl\langle F(\mathcal I)\bigr\rangle_{\mathbb Z}
       =4\sqrt2\mathbb Z,\qquad
 \bigl\langle d(Q_{\mathbb O}(0;\mathcal I))\bigr\rangle_{\mathbb Z}
       =8\sqrt2\mathbb Z.
\]
This assertion concerns the generated additive groups; it does not
claim that every element of either group is an individually attained
cubic value.
For an explicit distinction, put
$P(a)=a_0(a_0^2-3(a_1^2+a_2^2+a_3^2))$.
If $3\mid P(a)$, reduction modulo three forces $3\mid a_0$,
and substitution $a_0=3b$ gives
$P(a)=9b(3b^2-a_1^2-a_2^2-a_3^2)$.
Thus $3\mid P(a)$ implies $9\mid P(a)$.
The value $P(a)=6$ is impossible, so $F=12\sqrt2$ is not
attained, although it belongs to the proved additive value group.


\subsection{Primary source and exact reproducibility data}

Wilson's octonionic construction is the source of the pair and total
conditions \cite[Sections~2--4, pp.~2--6]{esn:Wilson2009}. His page~2
uses the odd-minus lattice $L$ and $R=\bar L$; page~3 gives pair sums
in $L\bar s$, the total in $Ls$, and the metric equal to one half of
the sum of octonion norms. The displayed basis map and
\eqref{wcl:eq:metric-isometry} retain these conjugations and the exact
scale in the present coordinates. The distinct coordinate change on
Wilson's page~6 is not substituted for this isometry.

The complete integer-arithmetic program and its output are supplied in the
following directory (the two displayed lines form one path):
\begin{quote}\small
\path{supporting_materials/workbench/evidence/}\\
\path{lattice_foundation_20260906/wilson/}
\end{quote}
The files are \path{check_wilson.py} and \path{EXACT_CHECKS.json}.
They record the generating-lattice index $4096$, Gram determinant $1$,
even diagonal, integral cyclic action, the fixed and coinvariant Gram
matrices, and the index $6561$ of their sum. The same program checks
all $64$ multiplication-table products for
\eqref{wcl:eq:actual-basis-map} and the stated exact cubic tests.
The proofs above establish rootlessness and the global descriptions;
these conclusions are not inferred from a finite search.


% Complete TeX conversion of THREE_LINE_IMAGE_PROOF.md, Sections 1--5.
% Retains the original F=tau from es_niemeier_wilson_cyclic.tex.
\section{The exact cubic image of the original common lattice}
\label{d4i:sec:complete-image}

Retain the coordinates $a_0,a_1,a_2,a_3$ and the original lattice
\[
 D_4=\{a\in\mathbb Z^4:a_0+a_1+a_2+a_3\in2\mathbb Z\}.
 \label{d4i:eq:original-domain}\tag{D4I1}
\]
The polynomial already used in \cref{wcl:sec:common-cubic} is
\[
 P(a)=a_0\bigl(a_0^2-3(a_1^2+a_2^2+a_3^2)\bigr).
 \label{d4i:eq:original-polynomial}\tag{D4I2}
\]
Here and throughout this section $F$ is the original triality cubic
$F=\tau$, with its ordered octonionic arguments as in
\cref{wcl:sec:construction,wcl:sec:common-cubic}. In particular, on the
literal common lattice of \cref{wcl:thm:literal-intersection},
\[
 \mathcal I=\{\sqrt2(a,a,a):a\in D_4\},\qquad
 F(\sqrt2(a,a,a))=2\sqrt2\,P(a).
 \label{d4i:eq:retained-cubic}\tag{D4I3}
\]
No coordinate permutation, scaling of $P$, or alteration of the lattice
is made. The witnesses below are actual points of this original lattice.
The image proof is elementary and requires no external arithmetic
existence theorem.

\subsection{Necessary restrictions in the original coordinates}
\label{d4i:sec:necessary}

For every integer $b$, one has $b^2\equiv b\pmod2$. Consequently,
for $a\in D_4$,
\[
 a_0^2+a_1^2+a_2^2+a_3^2
 \equiv a_0+a_1+a_2+a_3\equiv0\pmod2.
\]
Writing $s=a_1^2+a_2^2+a_3^2$, the factor $a_0^2-3s$ is congruent
to $a_0^2+s\equiv0\pmod2$. Thus $P(a)$ is even.

Modulo three the original polynomial satisfies $P(a)\equiv a_0^3$.
If $3\mid P(a)$, the field $\mathbb Z/3\mathbb Z$ has no nonzero
element with zero cube, so $3\mid a_0$. Substitution of $a_0=3b$
into the original formula gives
\[
 P(a)=3b(9b^2-3s)=9b(3b^2-s).
 \label{d4i:eq:three-divisibility}\tag{D4I4}
\]
Therefore
\[
 P(a)\in2\mathbb Z,\qquad
                 3\mid P(a)\ \Longrightarrow\ 9\mid P(a).
 \label{d4i:eq:necessary-restrictions}\tag{D4I5}
\]

\subsection{Three explicit affine lines and their exact values}
\label{d4i:sec:three-lines}

Define the common direction
\[
 v=(-3,1,1,1).
 \label{d4i:eq:common-direction}\tag{D4I6}
\]
Its coordinate sum is zero, and direct substitution gives
\[
                   P(v)=-3\bigl(9-3(1+1+1)\bigr)=0.
\]
For every $r\in\mathbb Z$ and $\varepsilon\in\{-1,1\}$, put
\[
 L_\varepsilon(r)=(-3r-\varepsilon,r,r,r+\varepsilon)
                =rv+(-\varepsilon,0,0,\varepsilon).
 \label{d4i:eq:two-lines}\tag{D4I7}
\]
Its coordinate sum is
$-3r-\varepsilon+r+r+r+\varepsilon=0$, so it belongs to $D_4$.
The original last-three-coordinate square sum and first-coordinate
square are respectively
\[
 s=r^2+r^2+(r+\varepsilon)^2=3r^2+2r\varepsilon+1,
 \qquad (-3r-\varepsilon)^2=9r^2+6r\varepsilon+1.
\]
Thus the factor in the original cubic is exactly
\[
 a_0^2-3s=(9r^2+6r\varepsilon+1)
                 -(9r^2+6r\varepsilon+3)=-2,
\]
and therefore
\[
 P(L_\varepsilon(r))=(-3r-\varepsilon)(-2)=6r+2\varepsilon.
 \label{d4i:eq:two-line-values}\tag{D4I8}
\]

For every $h\in\mathbb Z$, define the third line
\[
 L_0(h)=(-3h,h+1,h-1,h)=hv+(0,1,-1,0).
 \label{d4i:eq:third-line}\tag{D4I9}
\]
Its coordinate sum is zero, so it also belongs to $D_4$. Here
\[
 s=(h+1)^2+(h-1)^2+h^2=3h^2+2,
\]
and consequently
\[
 P(L_0(h))=(-3h)\bigl(9h^2-3(3h^2+2)\bigr)
                         =(-3h)(-6)=18h.
 \label{d4i:eq:third-line-values}\tag{D4I10}
\]
These identities hold for all signed integer parameters. In particular,
$L_0(0)=(0,1,-1,0)$ has cubic value zero, and the zero vector also
has cubic value zero.

\subsection{The global image and a proved reconstruction map}
\label{d4i:sec:global-image}

Set
\[
 S=\{n\in\mathbb Z:n\in2\mathbb Z\text{ and }
                                  (3\nmid n\text{ or }9\mid n)\}.
 \label{d4i:eq:image-domain}\tag{D4I11}
\]

\begin{theorem}[Complete cubic image and explicit section]
\label{d4i:thm:complete-image}
The original polynomial has the exact image
\[
 P(D_4)=18\mathbb Z\ \cup\ (6\mathbb Z+2)\ \cup\ (6\mathbb Z-2)=S.
 \label{d4i:eq:complete-image}\tag{D4I12}
\]
Equivalently,
\[
 P(D_4)=\{n\in\mathbb Z:
                 n\bmod18\in\{0,2,4,8,10,14,16\}\}.
 \label{d4i:eq:residue-image}\tag{D4I13}
\]
The following map is an actual section $\sigma:S\to D_4$:
\[
 \sigma(n)=
 \begin{cases}
 (-3h,h+1,h-1,h),&n=18h,\\
 (-3r-\varepsilon,r,r,r+\varepsilon),
        &n=2(3r+\varepsilon),\quad\varepsilon\in\{-1,1\}.
 \end{cases}
 \label{d4i:eq:global-section}\tag{D4I14}
\]
Its displayed parameters are uniquely determined in their respective
cases, and $P\circ\sigma=\operatorname{id}_S$.
\end{theorem}

\begin{proof}
Equation~\eqref{d4i:eq:necessary-restrictions} proves
$P(D_4)\subseteq S$. For the reverse inclusion let $n\in S$.
If $3\mid n$, then $9\mid n$, and evenness gives $n=18h$ for
a unique $h\in\mathbb Z$. Equation~\eqref{d4i:eq:third-line-values}
gives $P(L_0(h))=n$.

If $3\nmid n$, write $n=2q$. Then $3\nmid q$, so there is a
unique $\varepsilon\in\{-1,1\}$ with
$q\equiv\varepsilon\pmod3$. The number
$r=(q-\varepsilon)/3$ is consequently a unique integer.
Equation~\eqref{d4i:eq:two-line-values} gives
\[
                    P(L_\varepsilon(r))=6r+2\varepsilon=2q=n.
\]
The two displayed cases are disjoint and cover $S$, and both
witnesses lie in $D_4$ by \cref{d4i:sec:three-lines}. These
calculations prove the asserted section identity and both inclusions
in \eqref{d4i:eq:complete-image}. The first arithmetic progression
contributes residue zero modulo eighteen; the second contributes
$2,8,14$; and the third contributes $16,4,10$. This proves
\eqref{d4i:eq:residue-image}.
\end{proof}

The classification includes every positive, zero, and negative value.
The full excluded set is the union of the odd integers with
$18\mathbb Z+6$ and $18\mathbb Z+12$.
Let
\[
 A=\{a\in\mathbb Z^4:a_0+a_1+a_2+a_3=0\}
 \label{d4i:eq:sum-zero-slice}\tag{D4I15}
\]
be the literal sum-zero sublattice in the original coordinates. The
three witnesses lie in $A\subseteq D_4$, so the two inclusions
$S\subseteq P(A)\subseteq P(D_4)=S$ prove
\[
                              P(A)=P(D_4)=S.
 \label{d4i:eq:slice-image}\tag{D4I16}
\]
The image of the original cubic $F=\tau$ on the common lattice is
therefore exactly
\[
\begin{aligned}
 F(\mathcal I)&=2\sqrt2\,S\\
 &=4\sqrt2\bigl((9\mathbb Z)\cup(3\mathbb Z+1)
                                     \cup(3\mathbb Z-1)\bigr).
\end{aligned}
 \label{d4i:eq:original-cubic-image}\tag{D4I17}
\]
Because the original off-diagonal Albert determinant is $2F$, its
corresponding image is $4\sqrt2\,S$. The additive group generated
by that determinant image is $8\sqrt2\mathbb Z$, while the
additive group generated by $F(\mathcal I)$ is $4\sqrt2\mathbb Z$.
Indeed every $P$ value is even, and the first two lines at parameter
zero attain $P=2$ and $P=-2$. Thus the generated additive group of
the $P$ values is exactly $2\mathbb Z$; applying the two original
scalar factors proves both group assertions without identifying an
image set with its generated group.

\subsection{Exact local images with explicit local sections}
\label{d4i:sec:local-images}

For a prime $p$, let
\[
 D_{4,p}=\{a\in\mathbb Z_p^4:
                            a_0+a_1+a_2+a_3\in2\mathbb Z_p\}.
 \label{d4i:eq:local-domain}\tag{D4I18}
\]
This is precisely $D_4\otimes_{\mathbb Z}\mathbb Z_p$ in the original
coordinates. To verify the equality directly, use the integer basis
\[
\begin{gathered}
 b_0=(1,-1,0,0),\quad b_1=(0,1,-1,0),\\
 b_2=(0,0,1,-1),\quad b_3=(0,0,1,1).
\end{gathered}
 \label{d4i:eq:integral-basis}\tag{D4I19}
\]
For any $a$ in the displayed definition of $D_{4,p}$, the unique
coefficients of these basis vectors are
\[
\begin{aligned}
 c_0&=a_0,&c_1&=a_0+a_1,\\
 c_3&=(a_0+a_1+a_2+a_3)/2,&c_2&=c_3-a_3.
\end{aligned}
 \label{d4i:eq:local-basis-inverse}\tag{D4I20}
\]
All belong to $\mathbb Z_p$, and substitution gives the original
four coordinates. Conversely a $\mathbb Z_p$-linear combination
of the four vectors has coordinate sum $2c_3$. These two direct
calculations prove the tensor-product identification, including
at $p=2$. Taking integral coefficients gives the same basis
description of the original $D_4$.

The three line identities of \cref{d4i:sec:three-lines} hold over
$\mathbb Z_p$, since they are integer polynomial identities. Each
line stays in the sum-zero submodule of $D_{4,p}$.

\begin{theorem}[All local images and the exact local-to-global statement]
\label{d4i:thm:local-images}
For the unchanged polynomial $P$,
\[
\begin{aligned}
 P(D_{4,2})&=2\mathbb Z_2,\\
 P(D_{4,3})&=\mathbb Z_3^{\times}\cup9\mathbb Z_3,\\
 P(D_{4,p})&=\mathbb Z_p\qquad(p\ge5).
\end{aligned}
 \label{d4i:eq:local-images}\tag{D4I21}
\]
An integer $n$ belongs to every local image if and only if
$n\in S$, and hence if and only if $n\in P(D_4)$. The global
section is exactly \eqref{d4i:eq:global-section}; explicit local
sections are given in the proof.
\end{theorem}

\begin{proof}
For $p=2$, the parity calculation of \cref{d4i:sec:necessary}, reduced
modulo two, gives $P(D_{4,2})\subseteq2\mathbb Z_2$. Conversely,
given an arbitrary $n=2q\in2\mathbb Z_2$, put
\[
 r=(q-1)/3\in\mathbb Z_2.
\]
Division by three is valid since three is a unit in $\mathbb Z_2$.
Then $L_1(r)\in D_{4,2}$ and
\[
                         P(L_1(r))=6r+2=2q=n.
\]
Thus $n\mapsto L_1((n/2-1)/3)$ is an explicit section on
$2\mathbb Z_2$, and proves the first equality.

For $p=3$, reduction of the original polynomial modulo three and
the substitution $a_0=3b$ in
\eqref{d4i:eq:three-divisibility} prove that any value divisible by
three belongs to $9\mathbb Z_3$. Every unit
$n\in\mathbb Z_3^{\times}$ is congruent modulo three to exactly
one of $2\varepsilon$, with $\varepsilon\in\{-1,1\}$. Set
\[
 r=(n-2\varepsilon)/6\in\mathbb Z_3.
\]
This quotient is integral because the numerator is divisible by
three and two is a unit. The line $L_\varepsilon(r)$ then has
$P$ value $6r+2\varepsilon=n$.
Every $n\in9\mathbb Z_3$ can be written uniquely as $18h$ with
$h=n/18\in\mathbb Z_3$, again because two is a unit.
The third line has $P(L_0(h))=18h=n$. These two disjoint cases
give the explicit local section
\[
 n\longmapsto
 \begin{cases}
 L_\varepsilon((n-2\varepsilon)/6),
        &n\in\mathbb Z_3^{\times},\quad n\equiv2\varepsilon\pmod3,\\
 L_0(n/18),&n\in9\mathbb Z_3.
 \end{cases}
 \label{d4i:eq:three-adic-section}\tag{D4I22}
\]
They prove the second equality in \eqref{d4i:eq:local-images}.

For every prime $p\ge5$, six is a unit. Given any
$n\in\mathbb Z_p$, choose $r=(n-2)/6\in\mathbb Z_p$.
The point $L_1(r)$ belongs to $D_{4,p}$ and has
$P(L_1(r))=6r+2=n$. Thus $n\mapsto L_1((n-2)/6)$ is a
section on all of $\mathbb Z_p$, proving the third equality.

For integers $n$, the first local condition is precisely evenness;
the second is precisely $3\nmid n$ or $9\mid n$; and the other
primes impose no condition. Membership in every local image is
therefore exactly $n\in S$. The explicit integral section in
\cref{d4i:thm:complete-image} proves that those conditions suffice
for global integral representability. This proves the local-to-global
assertion without approximation or inference from a finite modular
calculation.
\end{proof}

\subsection{The exact factorization and quotient metric on the sum-zero slice}
\label{d4i:sec:factorization}

Define the integral map into the original sublattice $A$ of
\eqref{d4i:eq:sum-zero-slice} by
\[
 \Phi:\mathbb Z^3\longrightarrow A,\qquad
 \Phi(w,u,v)=(-3w-u-v,w+u,w+v,w).
 \label{d4i:eq:slice-map}\tag{D4I23}
\]
Its coordinate sum is zero. For $a\in A$, its proposed inverse is
\[
                 \Phi^{-1}(a)=(a_3,a_1-a_3,a_2-a_3).
 \label{d4i:eq:slice-inverse}\tag{D4I24}
\]
Substitution returns the original second, third, and fourth
coordinates immediately; the first becomes $-a_1-a_2-a_3=a_0$
by the defining equation of $A$. Substituting $\Phi$ into the
inverse returns $(w,u,v)$. These two equalities prove that $\Phi$
is an exact integral-module isomorphism.

For $a=(a_0,y,z,w)\in A$, put $T=y+z+w=-a_0$. Direct expansion
of the original polynomial gives
\[
\begin{aligned}
 P(a)&=-T\bigl(T^2-3(y^2+z^2+w^2)\bigr)\\
     &=2T(y^2+z^2+w^2-yz-yw-zw)\\
     &=T\bigl((y-z)^2+(y-w)^2+(z-w)^2\bigr).
\end{aligned}
 \label{d4i:eq:slice-expansion}\tag{D4I25}
\]
Under $\Phi$, one has $T=3w+u+v$, and the three differences are
$u-v,u,v$. Consequently
\[
                 P(\Phi(w,u,v))=2(3w+u+v)(u^2-uv+v^2).
 \label{d4i:eq:factorization}\tag{D4I26}
\]
This identity explains the three original lines without changing
their coordinates. On $L_\varepsilon(r)$, the quotient differences
are $(u,v)=(-\varepsilon,-\varepsilon)$, so
$u^2-uv+v^2=1$; on $L_0(h)$, they are $(u,v)=(1,-1)$, so
$u^2-uv+v^2=3$. The original first coordinates remain
$-3r-\varepsilon$ and $-3h$, respectively.

The exact quotient map is
\[
 \chi:A\longrightarrow\mathbb Z^2,\qquad
                  \chi(a)=(a_1-a_3,a_2-a_3).
 \label{d4i:eq:quotient-map}\tag{D4I27}
\]
It is onto because $\chi\Phi(0,u,v)=(u,v)$. An element lies in
its kernel exactly when its last three coordinates all equal a
common integer $w$; its first coordinate is then $-3w$.
Thus $\ker\chi=\mathbb Z(-3,1,1,1)$. The section
$(u,v)\mapsto\Phi(0,u,v)$ proves the split exact sequence
\[
 0\longrightarrow\mathbb Z(-3,1,1,1)\longrightarrow A
       \xrightarrow{\ \chi\ }\mathbb Z^2\longrightarrow0.
 \label{d4i:eq:split-exact-sequence}\tag{D4I28}
\]

For $\omega=(-1+\sqrt{-3})/2$, one has
$\omega+\overline\omega=-1$ and $\omega\overline\omega=1$.
Direct multiplication therefore proves
\[
                 (u+v\omega)(u+v\overline\omega)=u^2-uv+v^2.
 \label{d4i:eq:eisenstein-norm}\tag{D4I29}
\]
The quadratic factor in the original cubic is exactly this
Eisenstein norm, with the preceding coefficient two unchanged.

The original Euclidean metric on these same coordinates is,
by expansion,
\[
 \|\Phi(w,u,v)\|^2
            =12w^2+8w(u+v)+2(u^2+uv+v^2).
 \label{d4i:eq:full-slice-metric}\tag{D4I30}
\]
Write $d=(-3,1,1,1)$ for the direction when used in the metric
calculation. Its squared norm is twelve, and the exact pairing is
\[
          \langle\Phi(w,u,v),d\rangle=12w+4(u+v).
 \label{d4i:eq:direction-pairing}\tag{D4I31}
\]
Hence orthogonal projection away from the real line $\mathbb R d$
is given in the original coordinates by
\[
\begin{aligned}
 \pi\Phi(w,u,v)
  &=\Phi(w,u,v)-\left(w+\frac{u+v}{3}\right)d\\
  &=\left(0,\frac{2u-v}{3},\frac{2v-u}{3},-\frac{u+v}{3}\right).
\end{aligned}
 \label{d4i:eq:orthogonal-projection}\tag{D4I32}
\]
This depends only on the quotient coordinate $\chi(a)=(u,v)$.
Its squared norm is
\[
\begin{aligned}
 \|\pi\Phi(w,u,v)\|^2
  &=\frac{(2u-v)^2+(2v-u)^2+(u+v)^2}{9}\\
  &=\frac23(u^2-uv+v^2).
\end{aligned}
 \label{d4i:eq:projected-norm}\tag{D4I33}
\]
The quotient metric on $A/\mathbb Z d$, realized by this exact
projection, thus has Gram matrix
\[
              \frac13\begin{pmatrix}2&-1\\-1&2\end{pmatrix}
 \label{d4i:eq:quotient-gram}\tag{D4I34}
\]
in the displayed quotient coordinates $(u,v)$.
The retained common-lattice map $a\mapsto\sqrt2(a,a,a)$ multiplies
every squared norm by six, as proved in
\cref{wcl:thm:literal-intersection}. Consequently the corresponding
quotient Gram matrix on the image of $A$ in $\mathcal I$ is
\[
                         \begin{pmatrix}4&-2\\-2&4\end{pmatrix}.
 \label{d4i:eq:common-lattice-quotient-gram}\tag{D4I35}
\]
This is the exact quotient-metric calculation on the indicated
rank-three slice. It neither identifies that slice with the full
$D_4$ nor changes the original common-lattice metric.

% Integration-ready proof module; no canonical file has been edited.
% Requires amsmath, amssymb, amsthm, and hyperref in a standalone host.
% All coordinates and quadratic-form conventions are retained explicitly.

\section{The marked Niemeier lattice, its twelve theta indices, and the umbral shadow}
\label{nus:sec:shadow}

The original objects are
\[
 A=\{(x_0,\ldots,x_5)\in\mathbb Z^6:\sum_jx_j=0\},\qquad
 D=\{d\in\mathbb Z^4:\sum_id_i\equiv0\pmod2\},\qquad
 R=A^{\perp4}\perp D.
\]
The metrics are the displayed Euclidean metrics.  The four copies of $A$
retain their original fibre order $f_i289^j$ for $0\le j<6$, with fibre
labels $(00,10,01,11)$.  Their discriminant generator is
$\lambda=(5,-1,-1,-1,-1,-1)/6$.  The discriminant of $D$ retains
$10\mapsto v=(1,0,0,0)$, $01\mapsto s=(1,1,1,1)/2$, and
$11\mapsto c=(-1,1,1,1)/2$.  Consequently the original norm-valued
quadratic refinement, with codomain $\mathbb Q/2\mathbb Z$, is
\[
 q_R(k_1,k_2,k_3,k_4;y)
 =\frac56\sum_i k_i^2+\mathbf1_{y\ne0}\pmod{2\mathbb Z}.
\]
The original glue and Niemeier lattice are
\[
 \begin{split}
 E&=\{u\in\mathbb F_2^4:\sum_i u_i=0\},\qquad
 \tau(u)=(u_2+u_4,u_3+u_4),\\
 j(a,b)&=(a+b,a+2b,a,b),\qquad C_3=j(\mathbb F_3^2),\\
 H&=\{(3u+2j(a,b);\tau(u)):u\in E,\ (a,b)\in\mathbb F_3^2\},\\
 N&=\{x\in R^\vee:x+R\in H\}.
 \end{split}
\]
Thus $|H|=8\cdot9=72$ and $N/R=H$, with the root system
$X=A_5^4D_4$ proved in the retained lattice calculation.  In particular,
no norm-two vectors are introduced by the nonzero glue classes: their
original minima are four or six.  We use this proved literal root system,
rather than a rank-only identification, throughout the construction below.

\subsection{The primary prescription and its complete coefficient matrix}

Write $e(t)=\exp(2\pi\mathrm i t)$ and $q=e(\tau)$ for
$\tau\in\mathbb H$.  Cheng--Duncan--Harvey, arXiv:1307.5793v2,
Section~3.2 \cite[Sections~3.2--3.3]{esn:CDH2014}, source labels
\texttt{thetexp}, \texttt{transf\_theta},
and \texttt{transf\_theta\_2}, use
\[
 \theta_{m,r}(\tau,z)
 =\sum_{k\equiv r\ (2m)}q^{k^2/(4m)}e(kz),
 \qquad r\in\mathbb Z/(2m).
\]
Their Section~3.3, source label \texttt{def:OmegaMatrices}, defines
\[
 \Omega_m(n)_{r,s}
 =\begin{cases}
 1,&r+s\equiv0\pmod{2n},\quad r-s\equiv0\pmod{2m/n},\\
 0,&\text{otherwise},
 \end{cases}
 \qquad n\mid m.
\]
Their ADE table, source label \texttt{ADE1}, assigns
\[
 \Omega^{A_{m-1}}=\Omega_m(1),\qquad
 \Omega^{D_{m/2+1}}=\Omega_m(1)+\Omega_m(m/2),
\]
and extends by addition over components with the same Coxeter number.
These are the primary authors' prescribed maps.  The following calculation
specializes and proves their values in the present marked coordinates; it
does not assert a new ADE classification.

The six-cycle $e_j\mapsto e_{j+1\bmod6}$ on an $A$ fibre has
characteristic polynomial $(x^6-1)/(x-1)$ on its augmentation space.
Its order is six.  In $D$, take the original simple roots
$\alpha_1=e_1-e_2$, $\alpha_2=e_2-e_3$, $\alpha_3=e_3-e_4$,
$\alpha_4=e_3+e_4$.  Their reflection product
$s_{\alpha_1}s_{\alpha_2}s_{\alpha_3}s_{\alpha_4}$ sends
\[
 e_1\mapsto e_2,\quad e_2\mapsto e_3,\quad
 e_3\mapsto-e_1,\quad e_4\mapsto-e_4.
\]
It has order six and characteristic polynomial
\[
 (x^3+1)(x+1)
 =(x^2-x+1)(x+1)^2
 =\frac{(x^2-1)(x^6-1)}{(x-1)(x^3-1)}.
\]
Thus the original five components have Coxeter number $h=m=6$, and the
primary theta prescription has precisely $2h=12$ indices and denominator
$4h=24$ in its exponents.  The root rank remains $4\cdot5+4=24$;
the glue index remains $72$.

Put $J=\Omega_6(3)$.  Its two congruences are
\[
 r+s\equiv0\pmod6,\qquad r-s\equiv0\pmod4.
\]
The simultaneous solutions modulo twelve consist of the single value
$s\equiv5r\pmod{12}$: this value satisfies both congruences, and the
difference of any two solutions is divisible by both six and four, hence
by twelve.  Therefore
\[
 \boxed{\Omega^X=5I_{12}+J,\qquad J_{r,s}=\mathbf1_{s\equiv5r\ (12)}.}
\]
In the full original index order $0,1,\ldots,11$ this is
\begingroup
\csname c@MaxMatrixCols\endcsname=12\relax
\[
\Omega^X=\begin{pmatrix}
6&0&0&0&0&0&0&0&0&0&0&0\\
0&5&0&0&0&1&0&0&0&0&0&0\\
0&0&5&0&0&0&0&0&0&0&1&0\\
0&0&0&6&0&0&0&0&0&0&0&0\\
0&0&0&0&5&0&0&0&1&0&0&0\\
0&1&0&0&0&5&0&0&0&0&0&0\\
0&0&0&0&0&0&6&0&0&0&0&0\\
0&0&0&0&0&0&0&5&0&0&0&1\\
0&0&0&0&1&0&0&0&5&0&0&0\\
0&0&0&0&0&0&0&0&0&6&0&0\\
0&0&1&0&0&0&0&0&0&0&5&0\\
0&0&0&0&0&0&0&1&0&0&0&5
\end{pmatrix}.
\]
\endgroup
The unary series and the actual shadow are
\[
 S_{6,r}(\tau)=\frac1{2\pi\mathrm i}
  \left.\frac{\partial\theta_{6,r}}{\partial z}\right|_{z=0}
 =\sum_{k\equiv r\ (12)}kq^{k^2/24},
 \qquad S^X=\Omega^X S_6.
\]
All series and their derivatives converge normally on compact subsets of
their stated domains, by Gaussian decay.  Substitution $k\mapsto-k$
proves $S_{6,-r}=-S_{6,r}$, so $S_{6,0}=S_{6,6}=0$.
Let $\mathcal A:\mathbb C^5\to\mathbb C^{12}$ insert zero in positions
zero and six and put $(\mathcal Au)_{12-r}=-u_r$ for $1\le r\le5$.
This retains, and explicitly reconstructs, all twelve components.  Then
\[
 \Omega^X\mathcal A=\mathcal A M^X,\qquad
 \boxed{M^X=\begin{pmatrix}
 5&0&0&0&1\\0&4&0&0&0\\0&0&6&0&0\\
 0&0&0&4&0\\1&0&0&0&5
 \end{pmatrix}.}
\]
Indeed its $(r,s)$ entry is
$\Omega^X_{r,s}-\Omega^X_{r,12-s}$; direct substitution gives the matrix.
In particular the fourth A-type multiplicity and the full D-type
contribution have both been retained.  The diagonal entries
$(5,4,6,4,5)$ are exactly the Coxeter exponent multiplicities:
four copies of $(1,2,3,4,5)$ and one copy of $(1,3,3,5)$.

\subsection{Quadratic refinements and the finite theta transformations}

Let $L=\mathbb Z\ell$ with $(\ell,\ell)=12$.  Its dual is
$L^\vee=\mathbb Z\ell/12$.  Hence
\[
 A_L=L^\vee/L=\mathbb Z/12,\qquad
 q_L(r)=r^2/12\pmod{2\mathbb Z},\qquad
 Q_L(r)=q_L(r)/2=r^2/24\pmod{\mathbb Z},
\]
and the associated bilinear form is $B_L(r,s)=rs/12\pmod{\mathbb Z}$.
The symbols $q_L$ and $Q_L$ deliberately have different codomains and
retain the factor two.

Define the unitary Fourier matrices and reflection by
\[
 (F_\pm)_{r,s}=12^{-1/2}e(\pm rs/12),\qquad
 R_{r,s}=\mathbf1_{s\equiv-r\ (12)}.
\]
Finite geometric-series summation gives
\[
 F_\pm^2=R,\qquad F_-=RF_+=F_+R,\qquad F_+F_-=I.
\]
The substitution $k\mapsto k+12$ shows that $e(k^2/24)$ depends only on
$r=k\bmod12$.  Thus, with
$D_{r,s}=e(r^2/24)\delta_{r,s}$,
\[
 \theta_6(\tau+1,z)=D\theta_6(\tau,z).
\]
Poisson summation of the Gaussian gives
\begin{align*}
 \theta_6(-1/\tau,z/\tau)
 &=\sqrt{-\mathrm i\tau}\,e(6z^2/\tau)F_-\theta_6(\tau,z),\\
 \theta_6(-1/\tau,-z/\tau)
 &=\sqrt{-\mathrm i\tau}\,e(6z^2/\tau)F_+\theta_6(\tau,z).
\end{align*}
Here $\sqrt{-\mathrm i\tau}$ is the holomorphic square root positive
when $\tau$ is positive imaginary.  For completeness, apply the Fourier
transform convention $\widehat f(u)=\int_{\mathbb R}f(x)e(-ux)\,dx$
to $f(x)=\exp(\pi\mathrm i\tau x^2/12+2\pi\mathrm i zx)$.
Completing the square evaluates its transform; the real positive Gaussian
integral is obtained by squaring the integral and using polar coordinates,
and holomorphic continuation gives the expression for $\Im\tau>0$.
Periodizing $f(r+12x)$, its $n$-th Fourier coefficient is
$12^{-1}e(nr/12)\widehat f(n/12)$.  Summing these coefficients at zero
and grouping $n$ by its residue modulo twelve gives the first displayed
transformation.  Absolute Gaussian convergence justifies the integration,
grouping, and Fourier evaluation.  Replacing $z$ by $-z$ and using
$\theta_{6,-r}(\tau,z)=\theta_{6,r}(\tau,-z)$ gives the second.
The second, with its plus-sign Fourier matrix, is precisely CDH's
convention; the first uses the usual generator
$S=\left(\begin{smallmatrix}0&-1\\1&0\end{smallmatrix}\right)$
with elliptic argument $z/\tau$.

The matrices carrying the finite multiplier for this first convention are
\[
 U(T)=D,\qquad U(S)=e(-1/8)F_-.
\]
The square-root weight factor has not been included in $U$.
Differentiation of the theta transformation at $z=0$ gives
\[
 S_6(-1/\tau)=\tau\sqrt{-\mathrm i\tau}F_-S_6(\tau),\qquad
 S_6(\tau+1)=D S_6(\tau).
\]
Equivalently, CDH's reflected convention gives
$S_6(-1/\tau)=-\tau\sqrt{-\mathrm i\tau}F_+S_6(\tau)$.
The equality follows from $RS_6=-S_6$.  Thus the unary series has weight
$3/2$ with the same finite multiplier; theta has weight $1/2$.

The permutation $J$ preserves $Q_L$ because
$((5r)^2-r^2)/24=r^2\in\mathbb Z$.  Since $5^{-1}=5$ modulo twelve,
\[
 (JF_\pm)_{r,s}=12^{-1/2}e(\pm5rs/12)
 =(F_\pm J)_{r,s}.
\]
Therefore $\Omega^X$ commutes with both finite multiplier matrices.
The shadow transformation follows by applying this proved intertwiner.
Its component series vanish at infinity since every nonzero summand has
strictly positive exponent; the modular transformation shows the same
cuspidal behavior at every cusp of the full modular group.

There is also an exact comparison with the original A5 discriminant,
without replacing its form.  Reduction
\[
 p:\mathbb Z/12\longrightarrow\mathbb Z/6,\qquad p(r)=r\bmod6
\]
has kernel $\{0,6\}$ and satisfies
\[
 q_A(p(r))=\frac56r^2=10q_L(r)\pmod{2\mathbb Z}.
\]
Thus it is a surjective quadratic similitude with precisely multiplier
ten, and its bilinear forms satisfy the same multiplier relation.
It is not asserted to be an isometry.  Applying $p$ to four coordinates
and the identity to the retained $D$ discriminant gives
\[
 \widetilde D=(\mathbb Z/12)^4\oplus\mathbb F_2^2
 \longrightarrow D_R,
 \qquad (r;y)\longmapsto(p(r_1),\ldots,p(r_4);y).
\]
The pullback quadratic refinement is exactly
$10\sum_iq_L(r_i)+\mathbf1_{y\ne0}$, its radical is
$\{0,6\}^4\oplus0$, and its quotient by that radical is the original
$(D_R,q_R)$.  To prove the radical claim, pairing with each coordinate
generator forces $5r_i/6\in\mathbb Z$, hence $6\mid r_i$; the
nondegenerate $D$ form forces $y=0$.  The vectors so obtained have
quadratic value $30$ in each nonzero coordinate, hence zero modulo two.
The inverse image of the original $H$ has order $16\cdot72=1152$ and
is isotropic for this pullback form.  This gives the exact relation of the
theta indices to all original glue coordinates, with the form and
multiplier specified on both sides.

\subsection{The full Fourier determinant and the retained order-twelve character}
\label{nus:subsec:determinant}

Let $\chi:\mathrm{SL}_2(\mathbb Z)\to\mathbb Z/12$ be the retained
marking
\[
 \chi(T)=-1,\qquad \chi(S)=3.
\]
These values respect the presentation
$S^4=I$, $(ST)^3=S^2$: in the abelian target the required relations
are $4\cdot3=0$ and $3(3-1)=2\cdot3$ modulo twelve.
They therefore define a character, and its value at $T$ makes it onto.
We now calculate the determinant of the entire twelve-dimensional
finite theta representation, rather than that of a reduced component
list.

Put $\zeta=e(1/12)$.  The Vandermonde determinant gives
\[
 \det F_+=12^{-6}\prod_{0\le j<k\le11}(\zeta^k-\zeta^j).
\]
For every such pair,
\[
 \zeta^k-\zeta^j
 =2\mathrm i\sin\!\left(\frac{\pi(k-j)}{12}\right)
       \exp\!\left(\frac{\pi\mathrm i(j+k)}{12}\right),
\]
whose sine is strictly positive.  There are $\binom{12}{2}=66$ pairs,
and each index occurs in eleven of them, so
\[
 \sum_{j<k}(j+k)=11\sum_{j=0}^{11}j=11\cdot66=726.
\]
The phase of the determinant is consequently
\[
 \mathrm i^{66}\exp(726\pi\mathrm i/12)=(-1)\mathrm i=-\mathrm i.
\]
Its remaining factor is positive.  Unitarity, proved by the finite
geometric sum above, makes its absolute value one.  It follows exactly
that
\[
 \boxed{\det F_+=-\mathrm i,\qquad\det F_-=\mathrm i.}
\]
In particular retaining CDH's reflection changes the Fourier determinant
by $\det R=(-1)^5=-1$; it cannot be dropped.

The sum of the squares of all twelve indices is
\[
 \sum_{r=0}^{11}r^2=\frac{11\cdot12\cdot23}{6}=506.
\]
Consequently
\[
 \boxed{\det U(T)=e(506/24)=e(1/12),\qquad
 \det U(S)=e(-12/8)\det F_-=-\mathrm i=e(-3/12).}
\]
For an explicit check that the finite matrices have the relevant
metaplectic relations, let
$G=\sum_{r=0}^{11}e(r^2/24)$.  Listing the squares modulo twenty-four
gives
\[
 G=1+4e(1/24)+2e(1/6)+2e(3/8)+2e(2/3)+e(1/2).
\]
The first and last terms cancel, as do the terms with arguments $1/6$
and $2/3$.  The elementary sine and cosine values at $\pi/12$ and
$3\pi/4$ give
\[
 4e(1/24)+2e(3/8)=\sqrt6(1+\mathrm i)
 =\sqrt{12}\,e(1/8).
\]
Completing the square in the finite sum proves
\[
 (F_-D)^2=e(1/8)D^{-1}F_-,\qquad
 (F_-D)^3=e(1/8)R.
\]
For the first equality, its $(r,s)$ entry before evaluation is
$12^{-1}e(s^2/24)\sum_te(t^2/24-t(r+s)/12)$; shifting $t$ by
$r+s$ produces exactly the stated Gauss sum and matrix entry.
The second follows by multiplication and $F_-^2=R$.
It follows that
\[
 (U(S)U(T))^3=U(S)^2=-\mathrm iR,\qquad U(S)^4=-I_{12}.
\]
Taking determinants kills this metaplectic central sign because the
dimension twelve is even.  The determinant therefore descends to an
ordinary character of $\mathrm{SL}_2(\mathbb Z)$.  Its two computed
generator values, together with generation by $S,T$, prove
\[
 \boxed{\det U(\gamma)=e(-\chi(\gamma)/12).}
\]
The dual determinant is its inverse and is exactly
$e(\chi(\gamma)/12)$, with the retained signs $T\mapsto e(-1/12)$,
$S\mapsto\mathrm i$.

This is a statement about the determinant multiplier, with its domain
and codomain specified.  It is not a statement that a vector of weight
$1/2$ has scalar weight one after taking a determinant.  The separate
Dedekind-eta identities
\[
 \eta(\tau+1)=e(1/24)\eta(\tau),\qquad
 \eta(-1/\tau)=\sqrt{-\mathrm i\tau}\,\eta(\tau)
\]
show that $\eta^2$ has weight one and multiplier $\det U$, whereas
$\eta^{-2}$ has weight minus one and multiplier $(\det U)^{-1}$.
Indeed $\eta^2(S\tau)=(-\mathrm i)\tau\eta^2(\tau)$ and
$\eta^{-2}(S\tau)=\mathrm i\tau^{-1}\eta^{-2}(\tau)$.
The integer-weight factors $(c\tau+d)^{\pm1}$ obey the cocycle identity
by matrix multiplication, so the two generator formulas prove the
statements for every $\gamma$.  At $-I$ the character and odd integer
weight each contribute a minus sign, whose product is one.

\subsection{The residue Fourier map and what the Weyl quotient retains}

The original marking is
$\kappa(289^j)=j\pmod6$, with the full ordered residue list
\[
 (289^j)_{j=0}^5=(1,289,361,169,121,529)\pmod{840}.
\]
On the six-dimensional complex coordinate carrier define
\[
 E_2:\mathbb C[S_{840}]\longrightarrow\mathbb C[\mathbb Z/12],
 \qquad E_2(e_{289^j})=e_{2j},
\]
and on a separately indexed $\mathbb C^6$ define
\[
 B(e_k)=\frac{e_k+e_{k+6}}{\sqrt2},\qquad0\le k<6.
\]
Both are injective and preserve the indicated coordinate Hermitian
products.  Retain the original negative-sign normalized Fourier matrix
$F_6^-=(6^{-1/2}e(-kj/6))_{k,j}$.  Entry-by-entry calculation gives
\[
 \boxed{F_-E_2=B F_6^-.}
\]
Indeed the $(r,j)$ entry of the left side is
$12^{-1/2}e(-rj/6)$, and the right side has the same value, using
$r\bmod6$.  Thus the original unnormalized sample transform is carried
into this formula by its exact factor $\sqrt6$.
Translation $j\mapsto j+1$ intertwines under $E_2$ with translation
$r\mapsto r+2$ modulo twelve.  This gives a literal Fourier and
translation map, not an identification of different quadratic forms.
In particular $Q_L(6)=3/2=1/2\pmod{\mathbb Z}$, so the quotient
$\mathbb Z/12\to\mathbb Z/6$ does not carry $Q_L$ itself to a
quadratic form on the quotient.  Its exact multiplier-ten comparison
with the original A5 form was proved above.

For the original lattice action, $\rho(289)$ is the same six-cycle in
each of the four $A$ coordinate sets and is the identity on $D$.
Each coordinate transposition of $A$ is reflection in the corresponding
root $e_i-e_j$, so the full residue action is in
$W(R)=W(A)^4\times W(D)$.  It acts trivially on every A discriminant
because
$Q_\sigma\lambda-\lambda=e_{\sigma(0)}-e_0\in A$, and it fixes
the D discriminant.  Consequently
\[
 S_{840}\xrightarrow{\rho}O(N)\longrightarrow O(N)/W(R)
\]
is the trivial homomorphism.  This proves exactly which residue-group
action is killed, without discarding its already constructed Fourier map.

The original marked quotient is
\[
 G=O(N)/W(R)\cong\mathrm{GL}_2(\mathbb F_3).
\]
Here its actual coordinate map uses
$\ell_1=(1,1)$, $\ell_2=(1,2)$, $\ell_3=(1,0)$,
$\ell_4=(0,1)$.  For $A\in\mathrm{GL}_2(\mathbb F_3)$ determine
the unique signs and permutation by
\[
 \ell_iA=\varepsilon_i\ell_{p(i)},\qquad
 (P_pu)_i=u_{p(i)},\qquad B_p\tau(u)=\tau(P_pu).
\]
The four row lines exhaust the projective line, so the signs and
permutation exist uniquely.  The binary equation defines $B_p$ uniquely
since $\tau$ is onto.  Then $(k;y)\mapsto(\varepsilon_i k_{p(i)};B_py)$
preserves precisely the original two glue codes, and its ternary
restriction is $M_gj=jA$.  Conversely a marked glue stabilizer gives this
same $A$.  Multiplication of these identities proves the group
isomorphism, with the original row/column convention retained.

To compare its chosen signed lifts with CDH's simple-root permutation
section, let $Q_{\rm rev}e_j=e_{5-j}$ in an $A$ fibre.  When
$\varepsilon_i=-1$, replace the representative $-I$ by
$-Q_{\rm rev}$; their quotient is $Q_{\rm rev}\in W(A)$ and
$-Q_{\rm rev}(e_j-e_{j+1})=e_{4-j}-e_{5-j}$.
Thus the resulting map permutes the original simple roots by diagram
reversal.  When $\varepsilon_i=1$ no change is needed.
For D, put $\Delta=(\alpha_1\ \alpha_2\ \alpha_3\ \alpha_4)$.
Associate the original classes $v,s,c$ with outer nodes $1,4,3$,
respectively; their fundamental weights are
$e_1$, $(1,1,1,1)/2$, and $(1,1,1,-1)/2$, the last of which differs
from the retained $c$ by $(1,0,0,-1)\in D$.
For $B_p$, let $P_{B_p}$ permute these outer nodes exactly as $B_p$
permutes $v,s,c$ and fix node two.  Then
\[
 d_{B_p}=\Delta P_{B_p}\Delta^{-1}
\]
is the exact original-coordinate diagram representative.  Its
orthogonality follows because permuting the outer nodes preserves the
original Gram matrix.  It preserves $D$, and its action on fundamental
weights proves its discriminant action is exactly $B_p$.
Its quotient with the originally chosen D lift has trivial discriminant
action, hence belongs to the proved Weyl kernel $W(D)$.
This specifies the actual change of section on every summand.

Define the following four class functions using those retained maps:
\[
 a_g=\#\{i:p(i)=i\},\quad
 b_g=\sum_{p(i)=i}\varepsilon_i,\quad
 s_g=\operatorname{sgn}(B_p|_{\{v,s,c\}}),\quad
 c_g=\#\operatorname{Fix}(B_p|_{\{v,s,c\}})-1.
\]
The first two are the traces on the four-component permutation and
signed permutation representations.  The last is the trace on the
two-dimensional sum-zero subspace of the three outer-node permutation
representation, and $s_g$ is its sign character.  Thus these are precisely
CDH's $(\bar\chi^{X_A},\chi^{X_A},\chi^{X_D},\check\chi^{X_D})$;
the remaining $\bar\chi^{X_D}$ is one because the single D component
is fixed.  The comparison of sections above proves this identification
for the actual marked quotient.

CDH Section~5.1 \cite[Section~5.1]{esn:CDH2014} defines the full
twisted coefficient matrix by
\[
 (\Omega_g^{X_A})_{r,s}
 =\bigl(b_g\mathbf1_{r\ \text{even}}+
          a_g\mathbf1_{r\ \text{odd}}\bigr)\delta_{r,s},
\]
and, for its exceptional $m=6$ D4 case, by
\[
 (\Omega_g^{X_D})_{r,s}=
 \begin{cases}
 c_g\delta_{r,s},&r\equiv0,3\pmod6,\\
 \delta_{r,s}+s_gJ_{r,s},&r\equiv1,5\pmod6,\\
 s_g\delta_{r,s}+s_gJ_{r,s},&r\equiv2,4\pmod6.
 \end{cases}
\]
All deltas in this formula have indices modulo twelve.  The exact maps
from the marked glue quotient to the prescribed twisted shadows are
\[
 g\longmapsto\Omega_g^X=\Omega_g^{X_A}+\Omega_g^{X_D},\qquad
 S_g^X=\Omega_g^XS_6.
\]
They are class-function constructions from traces of representations;
the assignment $g\mapsto\Omega_g^X$ is not asserted to be a group
representation.  On the odd five-component presentation the explicit
matrix is
\[
 \boxed{M_g^X=\begin{pmatrix}
 a_g+1&0&0&0&s_g\\
 0&b_g&0&0&0\\
 0&0&a_g+c_g&0&0\\
 0&0&0&b_g&0\\
 s_g&0&0&0&a_g+1
 \end{pmatrix}.}
\]
This follows by the same signed component reconstruction $\mathcal A$:
in rows two and four the two D terms cancel, in row three the contribution
is $c_g$, and rows one and five pair with one another.
For reference, all possible character tuples are
\[
\begin{array}{c|rrrr|r}
 [g]&a_g&b_g&s_g&c_g&\text{number of elements}\\\hline
 1A&4&4&1&2&1\\
 2A&4&-4&1&2&1\\
 2B&2&0&-1&0&12\\
 4A&0&0&1&2&6\\
 3A&1&1&1&-1&8\\
 6A&1&-1&1&-1&8\\
 8A,8B&0&0&-1&0&12
\end{array}
\]
with the CDH class names of their table
\texttt{tab:chars:eul:6}.  The formulas apply to each actual $A$, so
their validity does not depend on a choice of conjugacy-class names.
The finite checker obtains the table from all $48$ invertible matrices
over $\mathbb F_3$ and the original row forms.

The distinguished residue action has quotient $g=1$ and therefore
produces the untwisted shadow, while a product $\rho(r)L_g$ produces
the same $S_g^X$ for every $r$.  The four-component signed action and
the D4 triality quotient survive and enter the displayed matrix.  The
central element $-I\in\mathrm{GL}_2(\mathbb F_3)$, for example,
changes $b_g$ from four to minus four while fixing $a_g,s_g,c_g$;
it changes the even-index odd-space entries with exactly that sign.

The already marked composite
\[
 \mathrm{SL}_2(\mathbb Z)\xrightarrow{\chi}\mathbb Z/12
 \longrightarrow\mathbb Z/6\xrightarrow{j\mapsto289^j}S_{840}
\]
is onto, has kernel $\chi^{-1}(\{0,6\})$, and sends $T$ to
$289^5=529$.  Following it by the Weyl quotient kills it, by the
proved residue calculation.  Following $\chi$ into the dual determinant
multiplier retains its complete order twelve, by the determinant theorem.
These are explicit different outgoing morphisms from the same marking,
including their kernels and signs.

\subsection{The shadow-to-Fricke map with the original scalar seventy-two}

The Coxeter calculation above gives the original formal frame exponent
\[
 \pi_X=-5[1]+[2]-[3]+5[6].
\]
Let $\mathcal E(\sum a_d[d])=\prod_d\eta(d\tau)^{a_d}$.
It is a group homomorphism from addition of formal exponents to
multiplication of nonzero holomorphic functions, by the exponent laws.
The retained scalar function is
\[
 \boxed{T_6(\tau)=\mathcal E(-\pi_X)
 =\frac{\eta(\tau)^5\eta(3\tau)}
        {\eta(2\tau)\eta(6\tau)^5}.}
\]
The minus sign on $\pi_X$ is essential: $\mathcal E(\pi_X)=T_6^{-1}$.

Here is a direct coefficient map from the actual shadow matrix to this
eta function.  Its diagonal difference vector is
$\alpha=(5,4,6,4,5)$, and set $\alpha_0=0$ modulo six.
Define the Lambert-series map at this marked matrix by
\[
 \mathscr L(\Omega^X)(\tau)
 =1+\sum_{n\ge1}\alpha_{n\bmod6}
                    \frac{nq^n}{1-q^n}.
\]
Convergence is absolute and locally uniform since $|q|<1$ and
$\alpha$ is bounded.  The coefficient of $q^N$ is exactly
$\sum_{d\mid N}d\alpha_{d\bmod6}$; thus this specifies every
coefficient, not only its initial terms.
For all positive integers $n$ the retained multiplicities satisfy
\[
 \alpha_{n\bmod6}
 =5-\mathbf1_{2\mid n}+\mathbf1_{3\mid n}-5\mathbf1_{6\mid n}.
\]
Checking the six residues proves this identity.  The product definition
$\eta(\tau)=q^{1/24}\prod_{n\ge1}(1-q^n)$ and termwise logarithmic
differentiation then give
\[
 \boxed{f^X:=\mathscr L(\Omega^X)
 =-q\frac{d}{dq}\log T_6
 =5\lambda_6+\lambda_2-\lambda_3,}
\]
where the original functions on the right are
\[
 \lambda_d=q\frac{d}{dq}\log\frac{\eta(d\tau)}{\eta(\tau)}
 =\frac{d-1}{24}+\sum_{n\ge1}\sigma_1(n)(q^n-dq^{dn}).
\]
The constant is
$[5(6-1)+(2-1)-(3-1)]/24=1$, and the other coefficients agree by
the six-residue identity.  This proves the shadow-to-eta calculation
directly for every coefficient.

The precise primary connecting map is CDH
Section~4.2 \cite[Section~4.2]{esn:CDH2014}:
\[
 \sigma^X(\tau,z)=\sum_{r\ (12)}\overline{S_r^X(\tau)}
                                \theta_{6,r}(\tau,z),\qquad
 \mathscr S_{1,1}(\sigma^X)=f^X.
\]
Its source labels are \texttt{eqn:sigmaX}, \texttt{def:weight2},
\texttt{weight2:Aseries}, and \texttt{weight2\_EZ}.
The displayed Lambert map is its explicit specialization at this
root-system shadow.  It verifies the resulting function independently
of the general Shimura-lift theorem.  In particular the shadow and the
retained Fricke function are connected by this complete coefficient map.

The original modular subgroup here is $\Gamma_0(6)$, and
$[\mathrm{SL}_2(\mathbb Z):\Gamma_0(6)]=12$.  To compute the index,
reduction modulo six gives
$\mathrm{SL}_2(\mathbb Z/6)\cong
\mathrm{SL}_2(\mathbb F_2)\times\mathrm{SL}_2(\mathbb F_3)$ of
order $6\cdot24=144$.  Reduction is onto: elementary upper and lower
unipotent matrices generate over each field, and CRT choices of their
entries realize each factor while acting trivially on the other.
The upper triangular subgroup modulo six has $\varphi(6)\cdot6=12$
elements (choose its first unit diagonal entry and its upper-right entry;
the other diagonal entry is its inverse).  Its inverse image is
$\Gamma_0(6)$, proving index $144/12=12$.
This modular-group index and the twelve theta indices have their
specified constructions; neither replaces the lattice index seventy-two.

Under the exact Fricke transformation $W_6\tau=-1/(6\tau)$,
the eta inversion formula sends each $\eta(dW_6\tau)$ to
$\sqrt{-\mathrm i(6/d)\tau}\,\eta((6/d)\tau)$.
The retained exponent vector at $d=1,2,3,6$ is $(5,-1,1,-5)$;
its reversed vector is its negative and its exponent sum is zero.
The $\tau$ factors therefore cancel.  Since each $6/d$ is positive,
the consistent square-root branch gives the positive scalar
\[
 \prod_{d\mid6}(6/d)^{a_d/2}
 =6^{5/2}3^{-1/2}2^{1/2}=72.
\]
Consequently
\[
 \boxed{T_6(W_6\tau)=72/T_6(\tau).}
\]
For the weight-two slash action of the determinant-six representative,
\[
 (f|_2W_6)(\tau)=(6\tau^2)^{-1}f(-1/(6\tau)).
\]
Differentiating the last identity and using
$d(W_6\tau)/d\tau=(6\tau^2)^{-1}$ gives
\[
 \boxed{f^X|_2W_6=-f^X.}
\]
On the Jacobi side, the primary Eichler--Zagier operator
$\mathcal W_6(6)$ acts on theta by $\Omega_6(6)=R$.
This is verified by its defining congruences, $r+s=0\pmod{12}$,
$r-s=0\pmod2$.  On the odd shadow and on the corresponding
skew-Jacobi form its eigenvalue is minus one.  The specialized coefficient
map just proved carries that negative eigenline to the negative Fricke
eigenline of $f^X$.  This proves an exact equivariant linear comparison
on these lines; it does not evaluate an $\mathrm{SL}_2(\mathbb Z)$
character on the determinant-six matrix.

Finally, the original sheet coordinates remain
\[
 J_6=T_6+5+72/T_6,\qquad Y_6=T_6-72/T_6.
\]
Direct substitution and differentiation prove
\[
 J_6(W_6\tau)=J_6(\tau),\quad
 Y_6(W_6\tau)=-Y_6(\tau),\quad
 Y_6^2=(J_6-5)^2-288,\quad
 dJ_6=Y_6\,dT_6/T_6=-2\pi\mathrm iY_6 f^X\,d\tau.
\]
Thus the original scalar $72$, the additive shift $5$, the branch
constant $288$, and the differential factor $-2\pi\mathrm i$ all
survive the complete shadow-to-sheet calculation.

\subsection{Established scope and reproducibility}

The constructed maps are the ADE coefficient specialization, the exact
finite-theta intertwiner, the multiplier determinant and its dual, the
quadratic similitude to the original discriminants, the marked glue-group
trace prescription, the original six-to-twelve Fourier map, and the
coefficient-level shadow-to-Fricke map.  Their kernels, weights, signs,
and constants were specified and proved above.  They establish concrete
relations even where a direct isometry or an identical group action is
absent.  They do not construct a new graded moonshine module or a new
moonshine theorem.  The name ``umbral shadow'' here refers to precisely
the CDH prescription, with its source provenance retained.

The standard-library program \texttt{check\_shadow.py} uses exact
arithmetic in $\mathbb Q[z]/(z^8-z^4+1)$ with $z=e(1/24)$.
It verifies both full and odd matrices, the Fourier determinant by the
Vandermonde product, the Gauss sum, the character signs, the Fourier
embedding, all $48$ original-code character records, the first $240$
Lambert coefficients, and the squared positive Fricke scalar $72^2$.
The proofs above establish the unbounded identities independently of
those finite checks.


% Workbench-ready, self-contained mathematical section.
% Required bibliography keys: Borovoi_2001, GorodnikOhShah_2009,
% Voight_2021, BenoistOh_2012, Good_1983, and ChatzakosPetridis_2016.
% The exact parametrization, census, quaternion bridge, and primary-source
% theorem audit are included as namespaced modules at the end of this file.

\section{Exact arithmetic of the Koide defect-one shell}
\label{koide-count-sec-main}

This section isolates an arithmetic result that does not use the existence of
the proposed complex structure on \(S^6\), any period-map assertion, or any
Yang--Mills interpretation.  All coordinates, chambers, local factors, and
normalizations used in the final constant are fixed explicitly below.

For \(r=(r_1,r_2,r_3)\), define the Koide quadratic form
\begin{equation}
 Q_K(r_1,r_2,r_3)
 :=\frac23(r_1+r_2+r_3)^2-(r_1^2+r_2^2+r_3^2).
 \label{koide-count-eq-qk}
\end{equation}
For integral coordinates put
\begin{equation}
 \Delta(x,y,z):=4(xy+xz+yz)-(x^2+y^2+z^2).
 \label{koide-count-eq-defect}
\end{equation}
Expanding \eqref{koide-count-eq-qk} gives the exact identity
\begin{equation}
 \Delta(x,y,z)=3Q_K(x,y,z).
 \label{koide-count-eq-defect-qk}
\end{equation}
Thus the equation studied here,
\(\Delta=-1\), is the nonzero level \(Q_K=-1/3\); it is not an affine
translate of the null cone \(Q_K=0\).

\subsection{Vieta reduction and exact ternary coordinates}
\label{koide-count-subsec-vieta}

Regarded as a monic quadratic equation in \(z\), the defect equation is
\begin{equation}
 z^2-4(x+y)z+x^2+y^2-4xy-1=0.
 \label{koide-count-eq-monic-z}
\end{equation}
Consequently the other root is
\begin{equation}
 z'=4(x+y)-z.
 \label{koide-count-eq-vieta-root}
\end{equation}
The analogous substitutions in \(x\) and \(y\) are integral involutions
preserving \(\Delta\).

Let $e_1,e_2,e_3$ be the standard basis and let $B_K$ be the integral
bilinear form with Gram matrix
\begin{equation}
 A_K=
 \begin{pmatrix}
 -1&2&2\\
 2&-1&2\\
 2&2&-1
 \end{pmatrix},
 \qquad
 B_K(u,v)=u^{\mathsf T}A_Kv.
 \label{koide-count-eq-vieta-gram}
\end{equation}
Thus $\Delta(u)=B_K(u,u)$.  The eigenvalue of $A_K$ on
$\mathbb R(1,1,1)$ is $3$, and its eigenvalue on the orthogonal complement is
$-3$, so this integral presentation has signature $(1,2)$.

\begin{proposition}[Vieta moves are Lorentzian reflections]
\label{koide-count-prop-vieta-reflections}
For $i=1,2,3$, define
\begin{equation}
 s_i(u)=u-
 \frac{2B_K(u,e_i)}{B_K(e_i,e_i)}e_i.
 \label{koide-count-eq-reflection-definition}
\end{equation}
Then $s_i\in\operatorname O(B_K,\mathbb Z)$,
$s_i^2=1$, and $\det s_i=-1$.  Explicitly,
\begin{align}
 s_1(x,y,z)&=(-x+4y+4z,y,z),\label{koide-count-eq-vieta-s1}\\
 s_2(x,y,z)&=(x,-y+4x+4z,z),\label{koide-count-eq-vieta-s2}\\
 s_3(x,y,z)&=(x,y,-z+4x+4y).\label{koide-count-eq-vieta-s3}
\end{align}
Consequently these are exactly the three Vieta root substitutions for
$\Delta=-1$, now identified as reflections of the unchanged Lorentzian
lattice $(\mathbb Z^3,B_K)$.
\end{proposition}

\begin{proof}
Since $B_K(e_i,e_i)=-1$, formula
\eqref{koide-count-eq-reflection-definition} is integral and direct
evaluation gives \eqref{koide-count-eq-vieta-s1}--
\eqref{koide-count-eq-vieta-s3}.  The standard reflection calculation
gives
\[
 B_K(s_i(u),s_i(v))=B_K(u,v),\qquad s_i^2=1.
\]
The map fixes the hyperplane $e_i^\perp$ pointwise and sends $e_i$ to
$-e_i$, hence has determinant $-1$.  Finally, when the equation is viewed
as a monic quadratic in one coordinate, the displayed replacement is the
sum of its two roots minus the original root, proving the last assertion.
\end{proof}

\begin{proposition}[Strict Vieta reduction]
\label{koide-count-prop-vieta-reduction}
Let \(0<x\leq y\leq z\) be integers satisfying
\(\Delta(x,y,z)=-1\).  If \(z<4(x+y)\), then the other root in
\eqref{koide-count-eq-vieta-root} satisfies \(0<z'<y\).  If
\(z\geq4(x+y)\), then \(z'\leq0\).
\end{proposition}

\begin{proof}
In the first case \(z'>0\) follows immediately from
\eqref{koide-count-eq-vieta-root}.  Let \(P(Z)\) be the polynomial on the
left of \eqref{koide-count-eq-monic-z}.  Since \(x\leq y\),
\[
 P(y)=x^2-8xy-2y^2-1
 \leq xy-8xy-2y^2-1<0.
\]
But \(P(y)=(y-z)(y-z')\), and \(y-z\leq0\); hence \(y-z'>0\), so
\(z'<y\).  The second assertion again follows directly from
\eqref{koide-count-eq-vieta-root}.
\end{proof}

Define the reduced-root set and its height count by
\begin{align}
 \mathcal R
 &:={}
 \{(x,y,z)\in\mathbb Z_{>0}^3:
 x\leq y\leq z,\ \Delta(x,y,z)=-1,\ z\geq4(x+y)\},
 \label{koide-count-eq-reduced-set}\\
 R(X)&:=\#\{(x,y,z)\in\mathcal R:z\leq X\}.
 \label{koide-count-eq-rx-definition}
\end{align}
Proposition~\ref{koide-count-prop-vieta-reduction} shows that the wall
\(z\geq4(x+y)\) is exactly the terminal side of the decreasing Vieta move
after the coordinates have been ordered.

\begin{proposition}[The positive Vieta graph is a rooted forest]
\label{koide-count-prop-vieta-forest}
Let $\mathcal S^+$ be the set of unordered triples of positive integers
satisfying $\Delta=-1$.  Join two vertices in a simple graph when one is
obtained from the other by a Vieta reflection $s_i$ and the reflected
coordinate is positive.  Every connected component of this graph is a
tree with exactly one root, and the roots are precisely the unordered
triples whose ordered representatives lie in $\mathcal R$.  Equivalently,
every positive integral solution has a unique reduced representative, and
it is obtained by repeatedly ordering the coordinates and reflecting the
largest coordinate whenever it lies below the terminal wall.
\end{proposition}

\begin{proof}
First, no Vieta reflection fixes a point of $\Delta=-1$.  Indeed, a fixed
reflection in a coordinate $w$, with the other two coordinates denoted by
$u,v$, would give $w=2(u+v)$, and substitution in the unchanged form gives
\[
 \Delta(u,v,2(u+v))=3(u^2+4uv+v^2)>0.
\]
Thus the sum of the three coordinates changes strictly across every edge.

Now order a vertex as $0<x\leq y\leq z$.  In fact $y<z$, because $y=z$
would give
\[
 \Delta(x,y,y)=-x^2+8xy+2y^2>0.
\]
Reflecting either of the two smaller coordinates produces a new coordinate
strictly larger than $z$:
\[
 4(y+z)-x>z,\qquad 4(x+z)-y>z.
\]
If $z<4(x+y)$, Proposition~\ref{koide-count-prop-vieta-reduction} says
that reflection in the largest coordinate produces $0<z'<y$ and strictly
decreases the coordinate sum.  If $z\geq4(x+y)$, that reflection is
nonpositive and hence is not an edge of $\mathcal S^+$.  Therefore every
nonterminal vertex has exactly one neighbour of smaller coordinate sum,
whereas a terminal vertex has none.

Iteration of the unique decreasing move terminates because the coordinate
sum is a positive integer.  The endpoint is represented by an element of
$\mathcal R$.  The graph has no cycle: on a putative simple cycle choose a
vertex of largest coordinate sum; its two neighbours on the cycle would be
two distinct smaller neighbours, contrary to uniqueness.  Thus each
component is a tree.  A tree in which following the unique decreasing edge
always terminates has exactly one terminal vertex, proving both uniqueness
and the characterization of the roots.
\end{proof}

For an ordered triple define
\begin{equation}
 n=z-4(x+y),\qquad q=y-x,\qquad a=x+y-4n,
 \label{koide-count-eq-forward-coordinates}
\end{equation}
and let
\begin{equation}
 \mathcal T:=
 \left\{(a,q,n)\in\mathbb Z^3:
 \begin{array}{c}
 a>0,\ q\geq0,\ n\geq0,\ a+4n>q,\\
 a^2-3q^2-18n^2=-2
 \end{array}
 \right\}.
 \label{koide-count-eq-ternary-chamber}
\end{equation}

\begin{theorem}[Integral ternary-coordinate bijection]
\label{koide-count-thm-coordinate-bijection}
The map in \eqref{koide-count-eq-forward-coordinates} restricts to a
bijection \(\Phi:\mathcal R\to\mathcal T\).  Its inverse is
\begin{equation}
 (a,q,n)\longmapsto
 \left(
 \frac{a+4n-q}{2},
 \frac{a+4n+q}{2},
 4a+17n
 \right).
 \label{koide-count-eq-inverse-coordinates}
\end{equation}
For arbitrary integral \((x,y,z)\) and the associated \((a,q,n)\),
\begin{equation}
 2\Delta(x,y,z)=a^2-3q^2-18n^2,
 \qquad z=4a+17n.
 \label{koide-count-eq-polynomial-identity}
\end{equation}
In particular,
\begin{equation}
 R(X)=\#\{(a,q,n)\in\mathcal T:4a+17n\leq X\}.
 \label{koide-count-eq-exact-count}
\end{equation}
\end{theorem}

\begin{proof}
Put \(p=x+y\).  Since \(q=y-x\), one has
\(2x=p-q\) and \(2y=p+q\), and direct substitution into
\eqref{koide-count-eq-defect} gives
\[
 2\Delta=p^2-3q^2+8pz-2z^2.
\]
The definition \(n=z-4p\) gives \(z=4p+n\), whence
\[
 2\Delta=p^2-3q^2-8pn-2n^2.
\]
Finally \(a=p-4n\), equivalently \(p=a+4n\), gives both identities in
\eqref{koide-count-eq-polynomial-identity}.

Suppose first that \((x,y,z)\in\mathcal R\).  Then \(q,n\geq0\), while
\(a+4n=p>q\) because \(p-q=2x>0\).  The defect equation and
\eqref{koide-count-eq-polynomial-identity} give
\(a^2-3q^2-18n^2=-2\).  If \(n=0\), then \(a=p>0\).  If \(n\geq1\),
then
\[
 a^2=3q^2+18n^2-2\geq18n^2-2\geq16n^2.
\]
Were \(a\leq0\), this inequality would force \(a\leq-4n\), contrary to
\(a+4n>q\geq0\).  Hence \(a>0\), and the forward map lands in
\(\mathcal T\).

Conversely, reducing the ternary equation modulo \(2\) gives
\(a\equiv q\pmod2\), so all three coordinates in
\eqref{koide-count-eq-inverse-coordinates} are integral.  Put
\(p=a+4n\).  The chamber inequality gives
\[
 x=\frac{p-q}{2}>0,\qquad
 y=\frac{p+q}{2}\geq x.
\]
Moreover
\[
 z=4a+17n=4p+n\geq4p>p>y.
\]
The identities \eqref{koide-count-eq-polynomial-identity} now give
\(\Delta=-1\) and \(z-4(x+y)=n\geq0\).  Thus the inverse map lands in
\(\mathcal R\).  The displayed formulae show directly that the two
compositions are the identity.  The last assertion follows from the height
identity in \eqref{koide-count-eq-polynomial-identity}.
\end{proof}

\begin{proposition}[Complete Pell classification of the chamber faces]
\label{koide-count-prop-pell-faces}
The integral points of $\mathcal T$ on its two coordinate faces are the
following two disjoint families.

On the terminal Vieta wall $n=0$, define integers $a_k,q_k$ by
\begin{equation}
 a_k+q_k\sqrt3=(1+\sqrt3)(2+\sqrt3)^k
 \qquad(k\geq0).
 \label{koide-count-eq-wall-pell-sequence}
\end{equation}
Then the points of $\mathcal T$ with $n=0$ are exactly
$(a_k,q_k,0)$ for $k\geq1$, and their reduced roots are
\begin{equation}
 \left(
  \frac{a_k-q_k}{2},
  \frac{a_k+q_k}{2},
  4a_k
 \right).
 \label{koide-count-eq-wall-pell-roots}
\end{equation}
The first is $(1,4,20)$.

On the symmetry face $q=0$, define integers $A_k,N_k$ by
\begin{equation}
 A_k+N_k\sqrt{18}=(4+\sqrt{18})(17+4\sqrt{18})^k
 \qquad(k\geq0).
 \label{koide-count-eq-symmetric-pell-sequence}
\end{equation}
Then the points of $\mathcal T$ with $q=0$ are exactly
$(A_k,0,N_k)$ for $k\geq0$, and their reduced roots are
\begin{equation}
 \left(
  \frac{A_k+4N_k}{2},
  \frac{A_k+4N_k}{2},
  4A_k+17N_k
 \right).
 \label{koide-count-eq-symmetric-pell-roots}
\end{equation}
The first is $(4,4,33)$.

In the projectivized real quadric these families tend, respectively, to
the two chamber-boundary rays
\begin{equation}
 (a,q,n)=\mathbb R_{>0}(\sqrt3,1,0),
 \qquad
 (a,q,n)=\mathbb R_{>0}(3\sqrt2,0,1).
 \label{koide-count-eq-pell-limit-rays}
\end{equation}
Each family contributes only $O(\log X)$ points to $R(X)$.
\end{proposition}

\begin{proof}
On $n=0$, the equation of $\mathcal T$ is
$a^2-3q^2=-2$.  Here $q=0$ is impossible.  If $q=1$, positivity gives
$a=1$, which fails the strict chamber inequality $a>q$.  For any positive
solution with $q\geq2$, multiplication by the norm-one unit
$2-\sqrt3$ gives
\begin{equation}
 a'+q'\sqrt3=(a+q\sqrt3)(2-\sqrt3),\qquad
 a'=2a-3q,\quad q'=2q-a.
 \label{koide-count-eq-wall-pell-descent}
\end{equation}
The unchanged equation implies
\[
 a<2q,\qquad
 a>\frac32q,\qquad
 a>q
 \quad(q\geq2),
\]
because its squares differ from $4q^2$, $9q^2/4$, and $q^2$ by,
respectively, $q^2+2$, $3q^2/4-2>0$, and $2q^2-2>0$.
Thus $a',q'>0$ and $q'<q$.  Repeated descent reaches $(1,1)$, and reversing
the descent proves \eqref{koide-count-eq-wall-pell-sequence} is the complete
positive solution set.  The recurrence
\[
 a_{k+1}=2a_k+3q_k,\qquad q_{k+1}=a_k+2q_k
\]
shows inductively that $a_k,q_k$ are both odd.  The chamber excludes only
$k=0$, and Theorem~\ref{koide-count-thm-coordinate-bijection} now gives
\eqref{koide-count-eq-wall-pell-roots}.

On $q=0$, one has $a^2-18n^2=-2$.  Necessarily $n>0$.  Direct substitution
for $1\leq n\leq5$ gives the sole solution $(a,n)=(4,1)$.  If $n\geq6$,
multiplication by the norm-one unit $17-4\sqrt{18}$ gives
\begin{equation}
 a'+n'\sqrt{18}=(a+n\sqrt{18})(17-4\sqrt{18}),\qquad
 a'=17a-72n,\quad n'=17n-4a.
 \label{koide-count-eq-symmetric-pell-descent}
\end{equation}
Both new coordinates are positive.  Indeed,
\[
 \frac an=\sqrt{18-\frac2{n^2}}
 >\frac{72}{17}\quad(n\geq6),
 \qquad
 \frac an<\sqrt{18}<\frac{17}{4}.
\]
Moreover $a>4n$ and hence $0<n'<n$.  Descent therefore reaches $(4,1)$,
and reversal proves \eqref{koide-count-eq-symmetric-pell-sequence} is the
complete positive solution set.  Its recurrence,
\[
 A_{k+1}=17A_k+72N_k,\qquad
 N_{k+1}=4A_k+17N_k,
\]
keeps $A_k$ even and $N_k$ odd, so
\eqref{koide-count-eq-symmetric-pell-roots} follows from the inverse
coordinate formula.

Conjugating the two identities defining the sequences shows
\[
 \frac{a_k}{q_k}\longrightarrow\sqrt3,\qquad
 \frac{A_k}{N_k}\longrightarrow\sqrt{18}=3\sqrt2,
\]
which proves \eqref{koide-count-eq-pell-limit-rays}.  Finally, the dominant
factors $2+\sqrt3>1$ and $17+4\sqrt{18}>1$ give exponential growth in $k$,
whereas the heights in \eqref{koide-count-eq-wall-pell-roots} and
\eqref{koide-count-eq-symmetric-pell-roots} are positive linear forms in
the displayed coordinates.  Hence either face contains only $O(\log X)$
points of height at most $X$.
\end{proof}

\subsection{The affine quadric and its local densities}
\label{koide-count-subsec-local}

Put
\begin{equation}
 F(a,q,n):=a^2-3q^2-18n^2,
 \qquad
 \mathscr X:=\{F=-2\}\subset\mathbb A^3.
 \label{koide-count-eq-affine-quadric}
\end{equation}
The matrix of \(F\) is
\(A=\operatorname{diag}(1,-3,-18)\), so
\(\det A=54\), and
\(-(-2)\det A=108\) is not a square.  The integral point
\(x_0=(1,1,0)\) lies on \(\mathscr X\).

For every prime \(p\), define
\begin{equation}
 N_{p,k}:=
 \#\{(a,q,n)\bmod p^k:F(a,q,n)\equiv-2\pmod{p^k}\},
 \qquad
 \sigma_p:=\lim_{k\to\infty}p^{-2k}N_{p,k}.
 \label{koide-count-eq-local-density-definition}
\end{equation}

\begin{theorem}[Exact local densities]
\label{koide-count-thm-local-densities}
The limits in \eqref{koide-count-eq-local-density-definition} exist and
are
\begin{equation}
 \sigma_2=\frac32,
 \qquad
 \sigma_3=2,
 \qquad
 \sigma_p=
 \frac{1-p^{-2}}{1-\left(\frac3p\right)p^{-1}}
 \quad(p\geq5).
 \label{koide-count-eq-all-local-densities}
\end{equation}
Writing \(\chi_{12}\) for the primitive real character modulo \(12\), the
singular series, with the primes taken in increasing order (equivalently,
with Borovoi's convergence-factor prescription), is therefore
\begin{equation}
 \mathfrak S(\mathscr X):=\prod_p\sigma_p
 =\frac92\frac{L(1,\chi_{12})}{\zeta(2)}
 =\frac92\frac{\log(2+\sqrt3)}{\sqrt3\,\zeta(2)}.
 \label{koide-count-eq-singular-series}
\end{equation}
\end{theorem}

\begin{proof}
We compute the two bad primes first.

For \(p=3\), once \(q,n\bmod3^k\) have been chosen, the equation for
\(a\) is
\[
 a^2\equiv-2+3q^2+18n^2\equiv1\pmod3.
\]
The two roots modulo \(3\) are nonsingular, and each has a unique Hensel
lift modulo every \(3^k\).  Thus
\begin{equation}
 N_{3,k}=2\cdot3^{2k},
 \label{koide-count-eq-three-adic-count}
\end{equation}
which gives \(\sigma_3=2\).

For \(p=2\), reduction modulo \(4\) leaves exactly two parity branches:
\begin{equation}
 (a,q,n)\equiv(0,0,1)\pmod2,
 \qquad
 (a,q,n)\equiv(1,1,0)\pmod2.
 \label{koide-count-eq-two-adic-branches}
\end{equation}
Fix \(k\geq4\) and put \(j=k-1\).  On the first branch write
\(a=2A\), \(q=2Q\), and put \(r=3n\).  Then
\begin{equation}
 F(a,q,n)+2
 =-2\bigl(r^2-2A^2+6Q^2-1\bigr).
 \label{koide-count-eq-two-adic-even-branch}
\end{equation}
Hence we must solve
\[
 r^2\equiv1+2A^2-6Q^2\pmod{2^j}.
\]
An odd unit is a square modulo \(2^j\), \(j\geq3\), precisely when it is
\(1\pmod8\), in which case it has four roots.  Here that condition is
equivalent to \(A,Q\) both being even.  There are \(2^{2j-2}\) such pairs
\((A,Q)\bmod2^j\), four choices of \(r\bmod2^j\), and, because \(3\) is
a unit, two odd choices of \(n\bmod2^k\) above each such \(r\).  This
branch therefore contributes
\begin{equation}
 2^{2j-2}\cdot4\cdot2=2^{2k-1}
 \label{koide-count-eq-two-adic-first-branch-count}
\end{equation}
solutions modulo \(2^k\).

On the second branch write \(n=2m\) and set
\begin{equation}
 v=\frac{a-q}{2},
 \qquad
 r=\frac{-a+3q}{2}.
 \label{koide-count-eq-two-adic-divided-coordinates}
\end{equation}
Then \(r,v\) have opposite parity, the inverse formula is
\(a=r+3v\), \(q=r+v\), and
\begin{equation}
 F(a,q,n)+2
 =-2\bigl(r^2-3v^2+36m^2-1\bigr).
 \label{koide-count-eq-two-adic-odd-branch}
\end{equation}
First suppose \(v=2V\) and \(r\) is odd.  For fixed
\(V\bmod2^{j-1}\) and \(m\bmod2^j\), the equation for \(r\) is
\[
 r^2\equiv1+12V^2-36m^2\pmod{2^j}.
\]
Its right side is \(1\pmod8\) exactly when \(V,m\) have the same parity.
There are
\(2\cdot2^{j-2}2^{j-1}=2^{2j-2}\) such pairs, and four roots \(r\),
so this subbranch gives \(2^{2j}\) triples \((r,v,m)\bmod2^j\).

Now suppose \(r=2R\) and \(v\) is odd.  The congruence in
\eqref{koide-count-eq-two-adic-odd-branch}, solved for \(v^2\), has an
odd square on the left; its right side is \(1\pmod8\) exactly when
\(R,m\) have opposite parity.  Again there are \(2^{2j-2}\) admissible
pairs and four square roots, hence \(2^{2j}\) transformed triples.

It remains to account for the divided-coordinate map.  Each class
\((r,v)\bmod2^j\) of opposite parity has exactly two inverse classes
\((a,q)\bmod2^k\): shifting either \(r\) or \(v\) by \(2^j\) changes both
\(a=r+3v\) and \(q=r+v\) by \(2^j\), while shifting both changes neither
modulo \(2^k\).  The class \(m\bmod2^j\) gives one even class
\(n=2m\bmod2^k\).  Thus the second branch contributes
\begin{equation}
 2\bigl(2^{2j}+2^{2j}\bigr)=2^{2k}.
 \label{koide-count-eq-two-adic-second-branch-count}
\end{equation}
Combining \eqref{koide-count-eq-two-adic-first-branch-count} and
\eqref{koide-count-eq-two-adic-second-branch-count} gives
\begin{equation}
 N_{2,k}=3\cdot2^{2k-1}\qquad(k\geq4),
 \label{koide-count-eq-two-adic-total-count}
\end{equation}
and hence \(\sigma_2=3/2\).

For \(p\geq5\), one has \(p\nmid108\).  Borovoi's good-prime formula
\cite[Section~1.3]{Borovoi_2001} gives
\[
 \sigma_p=\frac{1-p^{-2}}{1-\chi(p)p^{-1}},
 \qquad
 \chi(p)=\left(\frac{-(-2)\det A}{p}\right)
 =\left(\frac{108}{p}\right)=\left(\frac3p\right),
\]
which proves the remaining formula in
\eqref{koide-count-eq-all-local-densities}.  Since
\(\chi(p)=\chi_{12}(p)\) for \(p\geq5\),
\[
 \prod_{p\geq5}\sigma_p
 =L(1,\chi_{12})\prod_{p\geq5}(1-p^{-2})
 =\frac32\frac{L(1,\chi_{12})}{\zeta(2)}.
\]
Multiplication by \(\sigma_2\sigma_3=3\) gives the first equality in
\eqref{koide-count-eq-singular-series}.  Finally, the residue classes on
which \(\chi_{12}\) is nonzero give
\begin{align*}
 L(1,\chi_{12})
 &=\sum_{m\geq0}
 \left(
 \frac1{12m+1}-\frac1{12m+5}
 -\frac1{12m+7}+\frac1{12m+11}
 \right)\\
 &=\int_0^1\frac{1-x^2}{x^4-x^2+1}\,\mathrm dx
 =\frac{\log(2+\sqrt3)}{\sqrt3},
\end{align*}
where the middle equality follows by summing the geometric series and the
last by an elementary partial-fraction integration.  This completes the
local computation.
\end{proof}

\subsection{The archimedean chamber factor}
\label{koide-count-subsec-archimedean}

The Poincar\'e-residue gauge form on the whole smooth affine quadric is the
unique regular \(2\)-form \(\omega_\infty\) whose local representatives
satisfy
\begin{equation}
 \omega_\infty\wedge\mathrm dF
 =\mathrm da\wedge\mathrm dq\wedge\mathrm dn,
 \label{koide-count-eq-gauge-form}
\end{equation}
Explicitly, on the three standard charts it is
\begin{equation}
 \omega_\infty
 =\frac{\mathrm dq\wedge\mathrm dn}{2a}\quad(a\neq0),
 \qquad
 \omega_\infty
 =\frac{\mathrm da\wedge\mathrm dn}{6q}\quad(q\neq0),
 \qquad
 \omega_\infty
 =-\frac{\mathrm da\wedge\mathrm dq}{36n}\quad(n\neq0).
 \label{koide-count-eq-global-gauge-charts}
\end{equation}
These expressions agree on overlaps after restriction to the tangent bundle
of \(F=-2\).  Since the gradient
\((2a,-6q,-36n)\) cannot vanish on that level, the three charts cover
\(\mathscr X\), so \(\omega_\infty\) is globally defined and nowhere
vanishing.  Let \(m_\infty:=|\omega_\infty|\) be its associated positive
measure.  On the sheet \(a>0\), used for the chamber count, the first
expression in \eqref{koide-count-eq-global-gauge-charts} applies.
Let
\begin{equation}
 \mathcal C:=
 \{(a,q,n)\in\mathbb R^3:a>0,\ q\geq0,\ n\geq0,\ a+4n>q\}.
 \label{koide-count-eq-real-chamber}
\end{equation}

\begin{theorem}[Exact archimedean factor]
\label{koide-count-thm-archimedean-factor}
As \(X\to\infty\),
\begin{equation}
 m_\infty\{u\in\mathscr X(\mathbb R)\cap\mathcal C:
              4a+17n\leq X\}
 \sim I_\infty X,
 \qquad
 I_\infty=\frac1{2\sqrt3}
 \log\!\left(\frac3{2\sqrt2}\right).
 \label{koide-count-eq-archimedean-asymptotic}
\end{equation}
\end{theorem}

\begin{proof}
First, the wall \(a+4n>q\) removes only a bounded part of the positive
orthant.  Indeed, if \(q\geq a+4n\), then \(a\leq q-4n\); using
\(a^2=3q^2+18n^2-2\) gives
\[
 3q^2+18n^2-2\leq(q-4n)^2,
 \qquad\text{hence}\qquad
 q^2+4qn+n^2\leq1.
\]
Thus the removed real region has bounded \(m_\infty\)-measure.  Its only
integral point on \(\mathscr X\) with \(a>0,q,n\geq0\) is
\((a,q,n)=(1,1,0)\).

Introduce polar coordinates in the negative plane by
\begin{equation}
 \sqrt3\,q=r\cos\theta,
 \qquad
 3\sqrt2\,n=r\sin\theta,
 \qquad
 0\leq\theta\leq\frac\pi2.
 \label{koide-count-eq-polar-coordinates}
\end{equation}
On \(F=-2\), one has \(a=\sqrt{r^2-2}\) and \(r\,\mathrm dr=a\,\mathrm da\).
The Jacobian in \eqref{koide-count-eq-polar-coordinates} therefore gives
\begin{equation}
 \mathrm dm_\infty
 =\frac1{6\sqrt6}\,\mathrm da\,\mathrm d\theta.
 \label{koide-count-eq-polar-measure}
\end{equation}
At fixed \(\theta\), the height boundary is
\begin{equation}
 4a+\frac{17}{3\sqrt2}\sqrt{a^2+2}\sin\theta=X.
 \label{koide-count-eq-height-boundary}
\end{equation}
Let \(A_X(\theta)\) be its nonnegative solution when one exists, and put
it equal to zero otherwise.  Then \(0\leq A_X(\theta)/X\leq1/4\), and
for every fixed \(\theta\),
\[
 \frac{A_X(\theta)}{X}
 \longrightarrow
 \frac1{4+\frac{17}{3\sqrt2}\sin\theta}.
\]
Dominated convergence and \eqref{koide-count-eq-polar-measure} give
\begin{equation}
 I_\infty
 =\frac1{6\sqrt6}\int_0^{\pi/2}
 \frac{\mathrm d\theta}
 {4+\frac{17}{3\sqrt2}\sin\theta}.
 \label{koide-count-eq-archimedean-integral}
\end{equation}
For completeness, set \(t=\tan(\theta/2)\).  If
\(\alpha=2\sqrt2/3\) and \(\beta=3/(2\sqrt2)\), then
\(\alpha\beta=1\), and
\[
 \int_0^{\pi/2}
 \frac{\mathrm d\theta}
 {4+\frac{17}{3\sqrt2}\sin\theta}
 =\frac12\int_0^1\frac{\mathrm dt}{(t+\alpha)(t+\beta)}
 =3\sqrt2\log\!\left(\frac3{2\sqrt2}\right).
\]
Substitution in \eqref{koide-count-eq-archimedean-integral} proves
\eqref{koide-count-eq-archimedean-asymptotic}.
\end{proof}

\subsection{A genuine norm and the homogeneous counting theorem}
\label{koide-count-subsec-gos}

\begin{lemma}[A genuine norm realizing the height]
\label{koide-count-lem-genuine-norm}
Define
\begin{equation}
 \lVert(a,q,n)\rVert_*
 :=\max\{|4a+17n|,|q|,|a|\}.
 \label{koide-count-eq-star-norm}
\end{equation}
Then \(\lVert\cdot\rVert_*\) is a norm on \(\mathbb R^3\), and on the
chamber \(\mathcal C\) it satisfies
\begin{equation}
 \lVert(a,q,n)\rVert_*=4a+17n.
 \label{koide-count-eq-norm-is-height}
\end{equation}
For every integer \(X\geq0\),
\begin{equation}
 R(X)=\#\{u\in\mathscr X(\mathbb Z)\cap\mathcal C:
                    \lVert u\rVert_*<X+1\}.
 \label{koide-count-eq-r-as-sector-count}
\end{equation}
\end{lemma}

\begin{proof}
The linear map
\[
 (a,q,n)\longmapsto(4a+17n,q,a)
\]
has determinant \(-17\).  Thus \eqref{koide-count-eq-star-norm} is the
pullback of the supremum norm by an invertible linear map and is a genuine
norm.  If \((a,q,n)\in\mathcal C\), then
\[
 0<a<4a+17n,
 \qquad
 0\leq q<a+4n<4a+17n,
\]
which proves \eqref{koide-count-eq-norm-is-height}.  Since \(4a+17n\)
is integral on \(\mathscr X(\mathbb Z)\), equation
\eqref{koide-count-eq-r-as-sector-count} follows from
\eqref{koide-count-eq-exact-count}.
\end{proof}

\begin{lemma}[Verification of the Gorodnik--Oh--Shah hypotheses]
\label{koide-count-lem-gos-hypotheses}
Let \(\mathbf G=\operatorname{Spin}(F)\), acting through its standard
three-dimensional orthogonal representation on \(W=\mathbb R^3\), and let
\(x_0=(1,1,0)\).  Then
\[
 \mathscr X(\mathbb R)=\mathbf G(\mathbb R)^\circ x_0.
\]
This orbit is a closed affine symmetric variety; \(\mathbf G\) is connected
and absolutely almost simple over \(\mathbb Q\), is isotropic over
\(\mathbb R\), and acts irreducibly on \(W\) over \(\mathbb R\); and the
stabilizer \(\mathbf H=\operatorname{Stab}_{\mathbf G}(x_0)\) has no
nontrivial \(\mathbb Q\)-characters.  The cone \(\mathcal C\) in
\eqref{koide-count-eq-real-chamber} is admissible and generic for
\cite[Theorem~1.6]{GorodnikOhShah_2009}.
\end{lemma}

\begin{proof}
The form \(F\) has signature \((1,2)\).  Hence \(\mathbf G\) is a
connected absolutely almost simple group of type \(A_1\), it is
\(\mathbb R\)-isotropic, and its standard three-dimensional representation
is \(\mathbb R\)-irreducible.  The nonzero level \(\mathscr X\) is
Zariski closed.  With
\[
 Q=\sqrt3\,q,
 \qquad
 N=3\sqrt2\,n,
\]
its real equation becomes
\[
 Q^2+N^2-a^2=2.
\]
The parametrization
\[
 (a,\theta)\longmapsto
 \bigl(a,\sqrt{a^2+2}\cos\theta,
          \sqrt{a^2+2}\sin\theta\bigr),
 \qquad (a,\theta)\in\mathbb R\times(\mathbb R/2\pi\mathbb Z),
\]
shows that \(\mathscr X(\mathbb R)\) is connected.  Witt extension first
produces an element \(t\in\operatorname O(F)(\mathbb R)\) taking any
prescribed point of the level set to any other.  If \(\det t=-1\), choose
an anisotropic vector in the orthogonal complement of the target point and
compose \(t\) on the left with its reflection.  That reflection fixes the
target point and has determinant \(-1\); hence the composite lies in
\(\operatorname{SO}(F)(\mathbb R)\).  Thus
\(\operatorname{SO}(F)(\mathbb R)\) is transitive on the level set.  The
identity-component orbits have the full dimension of the level surface and
are open; because their complements are unions of open orbits, each is also
closed.  Connectedness therefore leaves only one such orbit, proving
\(\mathscr X(\mathbb R)=\mathbf G(\mathbb R)^\circ x_0\).  The stabilizer
of a nonisotropic vector is the fixed subgroup of the involution induced by
reflection in that vector, so the orbit is affine symmetric.

The orthogonal complement of \(x_0\) is
\begin{equation}
 x_0^\perp=\{(3s,s,t):s,t\in\mathbb Q\},
 \qquad
 F(3s,s,t)=6(s^2-3t^2).
 \label{koide-count-eq-rational-complement}
\end{equation}
Thus \(\mathbf H\) is the norm-one torus of
\(K=\mathbb Q(\sqrt3)\).  Its character lattice is \(\mathbb Z\), with
the nontrivial element of \(\operatorname{Gal}(K/\mathbb Q)\) acting by
\(-1\); consequently it has no nonzero Galois-invariant character and no
nontrivial \(\mathbb Q\)-character.  The integral point \(x_0\) supplies
the required nonempty integral locus.

It remains to check the sector rather than merely the affine orbit.  For the
given representation and any Euclidean sphere in \(W\), the Satake boundary
is
\begin{equation}
 \mathscr X^\infty=\{F=0\}\cap S(W)
 =\mathcal O_+\sqcup\mathcal O_-,
 \qquad
 \mathcal O_\pm:=\{F=0,\ \pm a>0\}\cap S(W),
 \label{koide-count-eq-satake-boundary}
\end{equation}
and \(\mathbf G(\mathbb R)^\circ\) is transitive on each displayed circle.
For a nonzero null vector with \(a,q,n\geq0\),
\[
 a=\sqrt{3q^2+18n^2}>q-4n.
\]
Hence both the closure and the relative interior of \(\mathcal C\) meet
exactly \(\mathcal O_+\), not \(\mathcal O_-\).  The intersection of
\(\partial\mathcal C\) with the Satake boundary consists of the two radial
directions
\begin{equation}
 (a,q,n)=(\sqrt3,1,0),
 \qquad
 (a,q,n)=(3\sqrt2,0,1).
 \label{koide-count-eq-sector-boundary-rays}
\end{equation}
These are exactly the two projective limits of the integral Pell families
in Proposition~\ref{koide-count-prop-pell-faces}.
This two-point intersection has measure zero for the smooth measure class
on the boundary orbit.
This is precisely the admissibility and genericity condition required by
the cited sector theorem.
\end{proof}

\subsection{The 2-adic orbit sign and Borovoi's coefficient}
\label{koide-count-subsec-borovoi}

Write $\widehat{\mathbb Z}=\prod_p\mathbb Z_p$.  For the nonsquare
ternary quadric under consideration, Borovoi associates to every local
point an orbit sign $\nu_v\in\{\pm1\}$, constant on
$\operatorname{Spin}(F)(\mathbb Q_v)$-orbits and equal to $+1$ at almost
every finite place.  Thus
$\nu_f((x_p)_p)=\prod_p\nu_p(x_p)$ is defined on
$\mathscr X(\widehat{\mathbb Z})$.  Because $\mathscr X(\mathbb R)$ is
connected, his normalization reads
\begin{equation}
 c_{\mathscr X}
 =\frac{2m_f\{x\in\mathscr X(\widehat{\mathbb Z}):\nu_f(x)=+1\}}
 {m_f(\mathscr X(\widehat{\mathbb Z}))},
 \label{koide-count-eq-borovoi-sign-definition}
\end{equation}
where $m_f$ is the finite Tamagawa measure attached to the gauge form
\eqref{koide-count-eq-gauge-form} and Borovoi's convergence factors.

\begin{proposition}[Finite adelic sign balance]
\label{koide-count-prop-finite-sign-balance}
For the affine quadric \(\mathscr X=\{F=-2\}\), the coefficient
\(c_{\mathscr X}\) in Borovoi's counting formula
\cite[Theorem~1.4.1 and Proposition~1.5.1]{Borovoi_2001} is exactly
\begin{equation}
 c_{\mathscr X}=1.
 \label{koide-count-eq-borovoi-coefficient}
\end{equation}
\end{proposition}

\begin{proof}
Continue to write \(A=\operatorname{diag}(1,-3,-18)\).  Put
\(e=(1,0,0)\), \(v=(1,0,1)\), and for every anisotropic vector \(w\) let
\[
 r_w(u)=u-\frac{2u^{\mathsf T}Aw}{w^{\mathsf T}Aw}\,w
\]
be its orthogonal reflection.  The element \(g=r_vr_e\) is
\begin{equation}
 g=
 \begin{pmatrix}
  -19/17&0&-36/17\\
  0&1&0\\
  -2/17&0&-19/17
 \end{pmatrix}.
 \label{koide-count-eq-sign-toggle-matrix}
\end{equation}
Direct multiplication gives
\begin{equation}
 g^{\mathsf T}Ag=A,
 \qquad
 \det g=1.
 \label{koide-count-eq-sign-toggle-orthogonal}
\end{equation}
Because \(17\in\mathbb Z_2^\times\), both \(g\) and \(g^{-1}\) have entries
in \(\mathbb Z_2\), so \(g\in\operatorname{SO}(F,\mathbb Z_2)\).  Its
spinor norm is
\begin{equation}
 \vartheta(g)=F(v)F(e)=-17
 \quad\text{in }\mathbb Q_2^\times/\mathbb Q_2^{\times2}.
 \label{koide-count-eq-toggle-spinor-norm}
\end{equation}

We identify the stabilizer's spinor-norm image.  By
\eqref{koide-count-eq-rational-complement}, over \(\mathbb Q_2\) the
orthogonal complement of \(x_0\) is
\begin{equation}
 F|_{x_0^\perp}
 =6(s^2-3t^2)
 =6N_{K_2/\mathbb Q_2}(s+t\sqrt3),
 \qquad K_2=\mathbb Q_2(\sqrt3).
 \label{koide-count-eq-two-adic-complement}
\end{equation}
The odd unit \(3\) is not a square in \(\mathbb Q_2\), so \(K_2\) is a
quadratic field and this binary plane is anisotropic.  Put
\(u_0=(3,1,0)\), so \(F(u_0)=6\).  If
\(h\in\operatorname{SO}(F|_{x_0^\perp})(\mathbb Q_2)\), then
\(hr_{u_0}\) is a two-dimensional orthogonal transformation of determinant
\(-1\), hence is a reflection \(r_w\) in some nonisotropic vector
\(w=(3s,s,t)\).  Therefore \(h=r_wr_{u_0}\) and
\[
 \vartheta(h)=F(w)F(u_0)
 =36N_{K_2/\mathbb Q_2}(s+t\sqrt3)
 \quad\bmod\mathbb Q_2^{\times2}.
\]
Conversely these products realize every norm square class as \(w\) ranges
through the nonzero vectors of the anisotropic plane.  Hence
\begin{equation}
 \vartheta\!\left(
 \operatorname{Stab}_{\operatorname{SO}(F)}(x_0)(\mathbb Q_2)
 \right)
 =N_{K_2/\mathbb Q_2}(K_2^\times)
 \quad\bmod\mathbb Q_2^{\times2}.
 \label{koide-count-eq-stabilizer-spinor-image}
\end{equation}

For a quadratic extension, its norm group is the kernel of the Hilbert
symbol with the defining square class
\cite[Corollary~13.4.6]{Voight_2021}.  The odd-unit formula
\cite[Equation~12.4.14]{Voight_2021} gives
\begin{equation}
 (-17,3)_2
 =(-1)^{((-17-1)/2)((3-1)/2)}=-1.
 \label{koide-count-eq-minus17-hilbert-symbol}
\end{equation}
Thus \(-17\notin N_{K_2/\mathbb Q_2}(K_2^\times)\).

By Witt extension, every \(x\in\mathscr X(\mathbb Q_2)\) is \(tx_0\) for
some \(t\in\operatorname O(F)(\mathbb Q_2)\).  If \(\det t=-1\), choose
an anisotropic vector in \(x^\perp\) and compose on the left with its
reflection; this fixes \(x\) and changes the determinant.  Hence every
such \(x\) has the form \(x=sx_0\) with
\(s\in\operatorname{SO}(F)(\mathbb Q_2)\).  The image of
\(\operatorname{Spin}(F)(\mathbb Q_2)\) in
\(\operatorname{SO}(F)(\mathbb Q_2)\) is the kernel of the spinor norm,
as used in \cite[Lemma~3.1]{Borovoi_2001}.  Consequently the
\(\operatorname{Spin}(F)(\mathbb Q_2)\)-orbit of \(x=sx_0\) is determined
by
\begin{equation}
 \vartheta(s)\in
 \mathbb Q_2^\times/
 \bigl(\mathbb Q_2^{\times2}
 N_{K_2/\mathbb Q_2}(K_2^\times)\bigr),
 \label{koide-count-eq-two-adic-orbit-quotient}
\end{equation}
a quotient of order two.  Equations
\eqref{koide-count-eq-toggle-spinor-norm} and
\eqref{koide-count-eq-minus17-hilbert-symbol} show that multiplication by
\(g\) exchanges these two orbits.  In Borovoi's convention, the orbit of
\(x_0\) has sign \(\nu_2=+1\) and the other orbit has sign
\(\nu_2=-1\).  Hence
\begin{equation}
 \nu_2(gx)=-\nu_2(x)
 \qquad(x\in\mathscr X(\mathbb Q_2)).
 \label{koide-count-eq-local-sign-toggle}
\end{equation}

Apply \(g\) at the prime \(2\) and the identity at every other finite
prime.  This gives a bijection of \(\mathscr X(\widehat{\mathbb Z})\).  It
preserves the finite Tamagawa measure: at \(2\) this follows from invariance
of the local measure, equivalently from \(g^*F=F\) and \(\det g=1\), and at
all other primes the map is the identity.  Equation
\eqref{koide-count-eq-local-sign-toggle} shows that it exchanges the finite
sign loci \(\nu_f=+1\) and \(\nu_f=-1\).  They therefore have equal
measure.  Since \(\mathscr X(\mathbb R)\) is connected, Borovoi's formula
becomes
\[
 c_{\mathscr X}
 =\frac{2m_f(\mathscr X(\widehat{\mathbb Z})_+)}
        {m_f(\mathscr X(\widehat{\mathbb Z}))}=1.
\]
The argument uses a measure-preserving, sign-reversing bijection; it does
not require \(g\) to be an involution (indeed \(g^2\neq I\)).
\end{proof}

\subsection{Calibration and the counting constant}
\label{koide-count-subsec-calibration}

\begin{theorem}[Calibrated reduced-root asymptotic]
\label{koide-count-thm-linear-asymptotic}
The number of reduced positive defect-one roots satisfies
\begin{equation}
 R(X)\sim C_0X\qquad(X\to\infty),
 \label{koide-count-eq-final-asymptotic}
\end{equation}
where
\begin{equation}
 \begin{aligned}
 C_0
 &=\mathfrak S(\mathscr X)I_\infty\\
 &=\frac9{2\pi^2}\log(2+\sqrt3)
   \log\!\left(\frac3{2\sqrt2}\right)\\
 &=0.035362047799731853\ldots.
 \end{aligned}
 \label{koide-count-eq-final-constant}
\end{equation}
\end{theorem}

\begin{proof}
Let \(\mu_{\mathrm{GOS}}\) denote the invariant measure on
\(\mathscr X(\mathbb R)\) with the normalization appearing in the
Gorodnik--Oh--Shah sector theorem
\cite[Theorem~1.6]{GorodnikOhShah_2009}.  The gauge measure \(m_\infty\)
in \eqref{koide-count-eq-gauge-form} is also a nonzero
\(\mathbf G(\mathbb R)^\circ\)-invariant Radon measure.  Because
\(\mathscr X(\mathbb R)\) is a single homogeneous orbit, uniqueness of
invariant measure gives
\begin{equation}
 \mu_{\mathrm{GOS}}=\kappa m_\infty
 \label{koide-count-eq-measure-scalar}
\end{equation}
for some \(\kappa>0\).

The scalar is fixed by applying both counting theorems to the same Euclidean
balls and the full cone \(W\).  Gorodnik--Oh--Shah give
\[
 \#\{x\in\mathscr X(\mathbb Z):|x|_2<T\}
 \sim\kappa\,m_\infty
 \{x\in\mathscr X(\mathbb R):|x|_2<T\}.
\]
Borovoi's Theorem~1.4.1 and Proposition~1.5.1, together with
Proposition~\ref{koide-count-prop-finite-sign-balance}, give for the same
balls
\[
 \#\{x\in\mathscr X(\mathbb Z):|x|_2<T\}
 \sim\mathfrak S(\mathscr X)\,m_\infty
 \{x\in\mathscr X(\mathbb R):|x|_2<T\}.
\]
The real volume here is asymptotic to a positive constant times \(T\).
Comparison therefore gives the invariant-measure calibration
\begin{equation}
 \kappa=\mathfrak S(\mathscr X).
 \label{koide-count-eq-measure-calibration}
\end{equation}

Lemma~\ref{koide-count-lem-genuine-norm} and
Lemma~\ref{koide-count-lem-gos-hypotheses} now allow the sector theorem to
be applied with the actual norm and chamber of the reduced-root problem:
\begin{align*}
 R(X)
 &\sim\mu_{\mathrm{GOS}}
 \{u\in\mathscr X(\mathbb R)\cap\mathcal C:
       \lVert u\rVert_*\leq X\}\\
 &=\mathfrak S(\mathscr X)\,
 m_\infty\{u\in\mathscr X(\mathbb R)\cap\mathcal C:
       4a+17n\leq X\}\\
 &\sim\mathfrak S(\mathscr X)I_\infty X.
\end{align*}
The switch between strict and weak inequalities is harmless by a
\((1\pm\varepsilon)X\) squeeze; on integral points it also follows directly
from \eqref{koide-count-eq-r-as-sector-count}.  Since
\(R(X)=R(\lfloor X\rfloor)\), the integer-height asymptotic extends to all
real \(X\to\infty\).  Finally, substituting
\eqref{koide-count-eq-singular-series} and
\eqref{koide-count-eq-archimedean-asymptotic}, and using
\(\zeta(2)=\pi^2/6\), gives exactly
\eqref{koide-count-eq-final-constant}.
\end{proof}

\subsection{Exact computation and the remaining open problem}
\label{koide-count-subsec-computation}

The first five reduced roots are
\begin{equation}
 (1,4,20),\quad(4,4,33),\quad(4,8,49),\quad
 (4,15,76),\quad(8,24,129).
 \label{koide-count-eq-first-roots}
\end{equation}
The first root is the first admissible term of the $n=0$ Pell family and
the second is the first term of the $q=0$ Pell family in
Proposition~\ref{koide-count-prop-pell-faces}.
Exact integer enumeration in the ternary coordinates gives the following
values.  The residual column uses the full constant in
\eqref{koide-count-eq-final-constant}, not a fitted constant.
\begin{center}
\begin{tabular}{r@{\qquad}r@{\qquad}r}
\toprule
\(X\)&\(R(X)\)&\(R(X)-C_0X\)\\
\midrule
10,000&353&\(-0.62048\)\\
20,000&715&\(7.75904\)\\
50,000&1,755&\(-13.10239\)\\
75,000&2,654&\(1.84642\)\\
100,000&3,535&\(-1.20478\)\\
125,000&4,460&\(39.74403\)\\
150,000&5,294&\(-10.30717\)\\
175,000&6,195&\(6.64164\)\\
200,000&7,109&\(36.59044\)\\
\bottomrule
\end{tabular}
\end{center}
The enumeration loops over \(n\geq0\) and \(a>0\) with
\(4a+17n\leq X\), tests whether
\[
 q^2=\frac{a^2-18n^2+2}{3}
\]
is the square of a nonnegative integer satisfying \(q<a+4n\), and then
reconstructs \((x,y,z)\) with
\eqref{koide-count-eq-inverse-coordinates}.  Every retained triple is
rechecked in the original polynomial \eqref{koide-count-eq-defect} and
against all ordering and wall inequalities.  Separate exhaustive congruence
enumerations agree with the formulae above at
\(p=2,3,5,7,11,13\) through the computed powers.  These finite checks
verify the implementation; the asymptotic theorem rests on the proofs above,
not on numerical fitting.

The reproducibility script is
\begin{center}
 \path{supporting_materials/workbench/reproduction/t1_koide_exact/verify_koide_narain_root_density.py}.
\end{center}
Running it with \texttt{--max-z 200000} emits the coordinate, local-density,
$2$-adic isometry, and bounded-enumeration certificates recorded in
\begin{center}
 \path{supporting_materials/workbench/reproduction/t1_koide_exact/koide_verification_200000.json}.
\end{center}
The script uses vectorized square detection only to propose candidates; it
rechecks every retained candidate with integer arithmetic in both coordinate
systems and in the original polynomial $\Delta$.

\begin{test}[Square-root discrepancy]
\label{koide-count-open-square-root-discrepancy}
\Open Determine whether, for every \(\varepsilon>0\),
\begin{equation}
 R(X)-C_0X=O_\varepsilon(X^{1/2+\varepsilon}).
 \label{koide-count-eq-square-root-discrepancy}
\end{equation}
The homogeneous-dynamics argument above and the exact
effective-well-roundedness verification in
\cref{koide-param-thm:power-saving} below prove, for some fixed
\(\eta>0\),
\[
 R(X)=C_0X+O(X^{1-\eta}),
\]
but do not prove
\eqref{koide-count-eq-square-root-discrepancy}.  No unproved spectral gap,
GRH assumption, or implicit ``conditional theorem'' is being used here:
the sharper error term is an open problem in this workbench; see the
primary-source comparison in
\cref{koide-lit-determination}.
\end{test}

% Integrated 4 September 2026 from the isolated exact-parametrization audit.
% Every label is namespaced koide-param-*; the original variables, chamber,
% strict walls, lattice, and height are unchanged.
\providecommand{\OO}{\mathcal O}
\providecommand{\NN}{\mathbb N}
\providecommand{\ZZ}{\mathbb Z}
\providecommand{\RR}{\mathbb R}
\providecommand{\QQ}{\mathbb Q}
\providecommand{\Norm}{\operatorname N}
\providecommand{\Tr}{\operatorname {Tr}}
\providecommand{\e}{\varepsilon}

\section{Exact ideal and unit-orbit parametrization of the unchanged sharp Koide count}

Throughout this section the variables, equation, chamber, and height are
exactly
\begin{equation}
 a^2-3q^2-18n^2=-2,
 \qquad
 a>0,\quad q\geq0,\quad n\geq0,\quad a+4n>q,
 \qquad
 4a+17n\leq X.
 \label{koide-param-eq:original-problem}
\end{equation}
Every auxiliary symbol introduced below is connected to these same integer
points by an explicitly proved map and inverse.  No count for a normalized
form, a rescaled lattice, or a substituted cutoff is used in place of
\eqref{koide-param-eq:original-problem}.
Thus
\begin{equation}
 R(X)=\#\left\{(a,q,n)\in\ZZ^3:\text{all conditions in
 \eqref{koide-param-eq:original-problem} hold}\right\}.
 \label{koide-param-eq:R-definition}
\end{equation}
The already calibrated leading constant is
\begin{equation}
 C_0=
 \frac9{2\pi^2}\log(2+\sqrt3)
 \log\!\left(\frac3{2\sqrt2}\right).
 \label{koide-param-eq:C0}
\end{equation}

The exact parametrization proved below has four pieces of data for each
fixed $n\geq1$: a positive integer $C$, a positive integer $D$, a square
root $E$ of $3$ modulo $D$, and a nonnegative Pell exponent $k$.  There is
no quotient left implicit: equations \eqref{koide-param-eq:bounded-seed-search},
\eqref{koide-param-eq:HNF-to-seed}, and \eqref{koide-param-eq:unit-orbit-coordinates} construct
$(a,q,n)$ by finite integer operations, while
\eqref{koide-param-eq:forward-HNF} and the descent \eqref{koide-param-eq:inverse-descent} recover
$(C,D,E,k)$ from $(a,q,n)$.  The face $n=0$ is recorded separately and
exactly.

This parametrization gives the exact sharp-cutoff formula
\eqref{koide-param-eq:exact-floor-formula}.  It also separates the discrepancy into
the radial shifted-divisor sum and the unit-phase discrepancy in
\eqref{koide-param-eq:exact-discrepancy-splitting}.  The square-root estimate
\begin{equation}
 R(X)-C_0X=O_\e(X^{1/2+\e})
 \label{koide-param-eq:square-root-target}
\end{equation}
is not proved or refuted here.  What is proved unconditionally, by applying
the effective symmetric-variety theorem to the exact chamber, is that
there is a constant $\eta>0$ for which
\begin{equation}
 R(X)=C_0X+O(X^{1-\eta}).
 \label{koide-param-eq:unconditional-power-saving}
\end{equation}
The cited general theorem does not supply $\eta\geq1/2$, and its proof does
not justify replacing \eqref{koide-param-eq:unconditional-power-saving} by
\eqref{koide-param-eq:square-root-target}.

\section{Arithmetic of $\ZZ[\sqrt3]$}

Put
\begin{equation}
 K=\QQ(\sqrt3),\qquad \OO=\ZZ[\sqrt3],\qquad
 \epsilon=2+\sqrt3.
 \label{koide-param-eq:field-data}
\end{equation}
For $\alpha=u+v\sqrt3\in K$, write
\begin{equation}
 \alpha'=u-v\sqrt3,
 \qquad
 \Norm(\alpha)=\alpha\alpha'=u^2-3v^2.
 \label{koide-param-eq:conjugate-norm}
\end{equation}

\begin{lemma}[Ring, class number, and totally positive units]
\label{koide-param-lem:ring-units}
The ring of integers of $K$ is $\OO$.  Every nonzero integral ideal of
$\OO$ is principal.  Every unit has norm $+1$, and the totally positive
unit group is
\begin{equation}
 \OO^{\times,+}=\{\epsilon^j:j\in\ZZ\}.
 \label{koide-param-eq:positive-units}
\end{equation}
\end{lemma}

\begin{proof}
Since $3\equiv3\pmod4$, the standard integral-basis calculation for a
quadratic field gives $\OO_K=\ZZ[\sqrt3]$ and discriminant $12$.
Minkowski's bound says that every ideal class in a real quadratic field of
discriminant $12$ has an integral representative of norm at most
\(\frac{2!}{2^2}\sqrt{12}=\sqrt3<2\).  Its norm must therefore be $1$,
so it is $\OO$ and the ordinary class number is one.

The norm of a unit is $\pm1$.  The equation $u^2-3v^2=-1$ is impossible
modulo $3$, because it would give $u^2\equiv2\pmod3$.  Hence every unit
has norm $+1$.  Dirichlet's unit theorem makes the positive units an
infinite cyclic group.  To identify its least generator, suppose
$\gamma=u+v\sqrt3>1$ and $\gamma'>0$.  Replacing $\gamma$ by its inverse
if necessary makes $v\geq0$.  If $v=0$, then $\gamma=1$.  If $v\geq1$,
then $u^2-3v^2=1$ forces $u\geq2$, and consequently
$\gamma\geq2+\sqrt3=\epsilon$.  Since $\epsilon$ itself has norm one,
it is the least positive unit greater than one and generates the group.
\end{proof}

For $m\geq1$ define
\begin{equation}
 \chi_{12}(d)=\left(\frac{12}{d}\right),
 \qquad
 \rho(m)=\sum_{d\mid m}\chi_{12}(d).
 \label{koide-param-eq:rho-definition}
\end{equation}
Here the symbol is Kronecker's symbol; in particular
\begin{equation}
 \chi_{12}(d)=
 \begin{cases}
  0,&2\mid d\text{ or }3\mid d,\\
  1,&d\equiv1,11\pmod {12},\\
  -1,&d\equiv5,7\pmod {12}.
 \end{cases}
 \label{koide-param-eq:chi12-table}
\end{equation}

\begin{lemma}[Ideal number]
\label{koide-param-lem:ideal-number}
The number of integral ideals of $\OO$ having absolute norm $m$ is
$\rho(m)$.
\end{lemma}

\begin{proof}
The Dedekind zeta function factors as
$\zeta_K(s)=\zeta(s)L(s,\chi_{12})$.  Comparing Dirichlet coefficients
gives the assertion.  Equivalently, the equality follows prime by prime.
At $2$ and $3$ there is one ramified prime ideal, so there is one ideal of
norm $p^e$ for every $e\geq0$.  If $\chi_{12}(p)=1$, then $p$ splits and
there are $e+1$ ideals of norm $p^e$.  If $\chi_{12}(p)=-1$, then $p$ is
inert, so there is one such ideal when $e$ is even and none when $e$ is
odd.  These are respectively the local values of
$\sum_{j=0}^e\chi_{12}(p)^j$.
Consequently, if
$m=2^{e_2}3^{e_3}\prod_{p\geq5}p^{e_p}$, then
\begin{equation}
 \rho(m)=
 \begin{cases}
  \displaystyle\prod_{\substack{p\mid m,\ p\geq5\\p\equiv1,11\ (12)}}(e_p+1),
   &e_p\text{ is even for every }p\equiv5,7\pmod {12},\\[6pt]
  0,&\text{otherwise}.
 \end{cases}
 \label{koide-param-eq:rho-Euler-product}
\end{equation}
\end{proof}

\section{Hermite-normal-form parameters for every ideal}

For $m\geq1$ let
\begin{equation}
 \mathcal P(m)=
 \left\{(C,D,E):
 \begin{array}{l}
 C,D\in\ZZ_{>0},\quad C^2D=m,\\
 0\leq E<D,\quad E^2\equiv3\pmod D
 \end{array}
 \right\}.
 \label{koide-param-eq:P-m}
\end{equation}

\begin{proposition}[Exact HNF parametrization of ideals]
\label{koide-param-prop:HNF-ideals}
The map
\begin{equation}
 (C,D,E)\longmapsto
 \mathfrak a(C,D,E)
 :=C\bigl(D\ZZ+(E+\sqrt3)\ZZ\bigr)
 \label{koide-param-eq:param-ideal}
\end{equation}
is a bijection from $\mathcal P(m)$ to the integral ideals of $\OO$ of
norm $m$.  Consequently
\begin{equation}
 \#\mathcal P(m)
 =\sum_{C^2D=m}\#\{E\bmod D:E^2\equiv3\pmod D\}
 =\rho(m).
 \label{koide-param-eq:root-divisor-identity}
\end{equation}
\end{proposition}

\begin{proof}
Relative to the basis $1,\sqrt3$ of $\OO$, the lattice in
\eqref{koide-param-eq:param-ideal} has column basis
\[
 \begin{pmatrix}CD&CE\\0&C\end{pmatrix}.
\]
Its index is $C^2D=m$.  It is stable under multiplication by $\sqrt3$
because
\begin{align*}
 \sqrt3(CD)&=D\,C(E+\sqrt3)-E(CD),\\
 \sqrt3\,C(E+\sqrt3)
 &=E\,C(E+\sqrt3)+\frac{3-E^2}{D}(CD),
\end{align*}
and the last coefficient is integral.  It is therefore an integral ideal
of norm $m$.

Conversely, write the unique column Hermite normal form of an integral
ideal $\mathfrak a$ as
\[
 \mathfrak a=A\ZZ+(B+C\sqrt3)\ZZ,
 \qquad A,C>0,\quad0\leq B<A.
\]
Stability of $A$ under multiplication by $\sqrt3$ implies $C\mid A$ and
$C\mid B$.  Write $A=CD$ and $B=CE$.  Stability of $B+C\sqrt3$ then
implies $D\mid3-E^2$.  The HNF inequality gives $0\leq E<D$, and the
index is $AC=C^2D$.  HNF uniqueness proves injectivity and surjectivity.

For completeness, the last equality in \eqref{koide-param-eq:root-divisor-identity}
can be checked without zeta functions.  Let
$r_3(D)=\#\{E\bmod D:E^2\equiv3\pmod D\}$.  For $p\geq5$, Hensel's
lemma gives $r_3(p^j)=2$ for every $j\geq1$ if
$\chi_{12}(p)=1$, and $r_3(p^j)=0$ if $\chi_{12}(p)=-1$.  Also
$r_3(2)=r_3(3)=1$ and $r_3(2^j)=r_3(3^j)=0$ for $j\geq2$.
Thus, for every prime power $p^e$,
\[
 \sum_{2c+j=e}r_3(p^j)=\sum_{j=0}^e\chi_{12}(p)^j.
\]
Multiplicativity proves \eqref{koide-param-eq:root-divisor-identity}.
\end{proof}

The HNF parameters can be recovered directly from a generator.  Suppose
$\alpha=a+q\sqrt3$ has positive norm $m$, with $a>0$ and $q\geq0$.  Put
\begin{equation}
 C=\gcd(a,q),\qquad A=a/C,\qquad B=q/C,\qquad D=m/C^2.
 \label{koide-param-eq:forward-HNF-pre}
\end{equation}
If $D=1$, put $E=0$.  If $D>1$, then
\begin{equation}
 \gcd(B,D)=1,\qquad E\equiv A B^{-1}\pmod D,\qquad0\leq E<D.
 \label{koide-param-eq:forward-HNF}
\end{equation}
Indeed,
$\gcd(B,D)=\gcd(B,A^2-3B^2)=\gcd(B,A^2)=1$, and
$E^2\equiv3\pmod D$.  The principal ideal $(\alpha)$ is exactly
$\mathfrak a(C,D,E)$: the element $A+B\sqrt3$ lies in
$D\ZZ+(E+\sqrt3)\ZZ$, so the principal ideal is contained in the HNF
ideal, while both ideals have norm $D$ after division by $C$.

\section{Complete integral parametrization of the original equation}

For $n\geq1$ retain the exact integer
\begin{equation}
 m_n=18n^2-2.
 \label{koide-param-eq:mn}
\end{equation}
It satisfies $m_n>0$ and $m_n\equiv1\pmod3$.

\begin{lemma}[Every ideal of norm $m_n$ has positive-norm generators]
\label{koide-param-lem:positive-norm-generator}
If $\mathfrak a$ is an integral ideal of $\OO$ of norm $m_n$, then every
generator has norm $+m_n$, rather than $-m_n$.
\end{lemma}

\begin{proof}
Lemma~\ref{koide-param-lem:ring-units} makes $\mathfrak a$ principal.  A generator
$u+v\sqrt3$ has norm $\pm m_n$.  Negative norm would give
\[
 u^2-3v^2=-m_n,
\]
and reduction modulo $3$ would give
$u^2\equiv-m_n\equiv2\pmod3$, which is impossible.  Thus the norm is
$+m_n$.
\end{proof}

For each $\mathbf p=(C,D,E)\in\mathcal P(m_n)$, choose a generator of
$\mathfrak a(C,D,E)$.  By Lemma~\ref{koide-param-lem:positive-norm-generator} its
norm is $m_n$.  Change its sign so that both real embeddings are positive,
and multiply by a unique power of $\epsilon$ so as to obtain the unique
generator
\begin{equation}
 \alpha_{n,\mathbf p}=a_{n,\mathbf p,0}
   +q_{n,\mathbf p,0}\sqrt3,
 \qquad
 \sqrt{m_n}\leq\alpha_{n,\mathbf p}
       <\epsilon\sqrt{m_n}.
 \label{koide-param-eq:canonical-generator}
\end{equation}
Existence and uniqueness follow from
\eqref{koide-param-eq:positive-units}, because the half-open intervals
$[\epsilon^j\sqrt{m_n},\epsilon^{j+1}\sqrt{m_n})$ partition
$\RR_{>0}$.  Since
$\alpha_{n,\mathbf p}'=m_n/\alpha_{n,\mathbf p}$, the lower inequality
in \eqref{koide-param-eq:canonical-generator} implies
$q_{n,\mathbf p,0}\geq0$.

\begin{lemma}[Purely integral reduced seed and descent]
\label{koide-param-lem:integral-seed-descent}
For an integral pair $(a,q)$ satisfying
\[
 a^2-3q^2=m_n,\qquad a>0,\qquad q\geq0,
\]
the half-open inequalities
\[
 \sqrt{m_n}\leq a+q\sqrt3<\epsilon\sqrt{m_n}
\]
are equivalent to the purely integral inequalities
\begin{equation}
 a>2q\geq0.
 \label{koide-param-eq:integral-seed-condition}
\end{equation}
If $a\leq2q$, multiplication by $\epsilon^{-1}=2-\sqrt3$ is the
integral inverse step
\begin{equation}
 (a,q)\longmapsto(a^-,q^-)=(2a-3q,2q-a).
 \label{koide-param-eq:inverse-descent}
\end{equation}
It preserves the norm, has $a^->0$ and $0\leq q^-<q$, and can therefore
be iterated until the unique pair satisfying \eqref{koide-param-eq:integral-seed-condition}
is reached.
\end{lemma}

\begin{proof}
Put $u=(a+q\sqrt3)/\sqrt{m_n}$.  Norm positivity and the sign conditions
give $u\geq1$.  Moreover
\[
 \frac a{\sqrt{m_n}}=\frac{u+u^{-1}}2,
 \qquad
 \frac{q\sqrt3}{\sqrt{m_n}}=\frac{u-u^{-1}}2.
\]
For $u\geq1$, the ratio
\[
 \frac aq=\sqrt3\,\frac{u^2+1}{u^2-1}
\]
(with value $+\infty$ when $q=0$) is strictly decreasing.  At
$u=\epsilon$ it equals $2$, since
$\epsilon+\epsilon^{-1}=4$.  Thus $u<\epsilon$ is equivalent to $a>2q$;
the strict inequality exactly matches the half-open upper endpoint.

If $a\leq2q$, norm positivity gives $a>\sqrt3q$, and hence
\[
 a^-=2a-3q>(2\sqrt3-3)q>0,
 \qquad q^-=2q-a\geq0.
\]
Also $a>\sqrt3q>q$ gives $q^-=2q-a<q$.  Thus the positive integer $q$
strictly decreases at every inverse step, so the process terminates.  A
terminal pair has $a>2q$.  Conversely, for such a pair the second
coordinate after one inverse step is negative, so no further inverse step
stays in the nonnegative-$q$ region.  This proves uniqueness as well as
termination.
\end{proof}

The preceding seed can be reconstructed from the HNF data by a bounded
integer search; thus the parametrization does not require an unspecified
choice of a generator.  More precisely, for
$\mathbf p=(C,D,E)\in\mathcal P(m_n)$ there is a unique pair $(A,B)$ with
\begin{equation}
 \begin{gathered}
  A,B\in\ZZ,\qquad A>0,\qquad 0\leq B,\qquad B^2<D,\\
  A^2-3B^2=D,\qquad A\equiv EB\pmod D.
 \end{gathered}
 \label{koide-param-eq:bounded-seed-search}
\end{equation}
It gives the generator in \eqref{koide-param-eq:canonical-generator} by the exact
integer formula
\begin{equation}
 (a_{n,\mathbf p,0},q_{n,\mathbf p,0})=(CA,CB).
 \label{koide-param-eq:HNF-to-seed}
\end{equation}
Indeed, existence follows by writing the reduced generator as
$C(A+B\sqrt3)$.  The seed inequality $A>2B$ gives
$D=A^2-3B^2>B^2$.  Conversely, the conditions in
\eqref{koide-param-eq:bounded-seed-search} imply $A>2B$ and put
$A+B\sqrt3$ in $D\ZZ+(E+\sqrt3)\ZZ$.  Its principal ideal is contained
in that HNF ideal, and both have norm $D$, so they are equal.  Lemma
\ref{koide-param-lem:integral-seed-descent} and HNF uniqueness give uniqueness.  In
particular one needs to test only the integers
$0\leq B\leq\lfloor\sqrt{D-1}\rfloor$ and the positive square
$A=\sqrt{D+3B^2}$.
The phase introduced below is therefore also explicit in the same integer
data:
\begin{equation}
 \tau_{n,(C,D,E)}
 =\frac{\log((A+B\sqrt3)/\sqrt D)}{\log\epsilon}.
 \label{koide-param-eq:HNF-phase-explicit}
\end{equation}
Here the factor $C$ cancels because $m_n=C^2D$; no ideal generator remains
to be chosen.

For $k\geq0$ define integers $a_{n,\mathbf p,k}$ and
$q_{n,\mathbf p,k}$ by
\begin{equation}
 a_{n,\mathbf p,k}+q_{n,\mathbf p,k}\sqrt3
 =\alpha_{n,\mathbf p}\epsilon^k.
 \label{koide-param-eq:unit-orbit-coordinates}
\end{equation}
Equivalently,
\begin{equation}
 \begin{pmatrix}a_{n,\mathbf p,k+1}\\q_{n,\mathbf p,k+1}\end{pmatrix}
 =
 \begin{pmatrix}2&3\\1&2\end{pmatrix}
 \begin{pmatrix}a_{n,\mathbf p,k}\\q_{n,\mathbf p,k}\end{pmatrix}.
 \label{koide-param-eq:unit-recurrence}
\end{equation}

\begin{theorem}[Integral parametrization, with inverse and multiplicity]
\label{koide-param-thm:integral-parametrization}
The map
\begin{equation}
 (n,\mathbf p,k)
 \longmapsto
 (a_{n,\mathbf p,k},q_{n,\mathbf p,k},n),
 \qquad
 n\geq1,\quad\mathbf p\in\mathcal P(m_n),\quad k\geq0,
 \label{koide-param-eq:parameter-map}
\end{equation}
is a bijection onto all integral solutions of
\begin{equation}
 a^2-3q^2-18n^2=-2,\qquad a>0,\ q\geq0,\ n\geq1.
 \label{koide-param-eq:n-positive-solutions}
\end{equation}
Every such solution automatically satisfies $a+4n>q$.

The inverse is as follows.  From $(a,q,n)$ form
$\alpha=a+q\sqrt3$ and recover $(C,D,E)$ by
\eqref{koide-param-eq:forward-HNF-pre}--\eqref{koide-param-eq:forward-HNF}.  There is a unique
$k\geq0$ such that
\begin{equation}
 \alpha=\alpha_{n,(C,D,E)}\epsilon^k.
 \label{koide-param-eq:recover-k}
\end{equation}
Thus every point occurs with multiplicity exactly one.
\end{theorem}

\begin{proof}
Equation \eqref{koide-param-eq:unit-orbit-coordinates} preserves the norm, so its
image satisfies
\[
 a_{n,\mathbf p,k}^2-3q_{n,\mathbf p,k}^2=m_n,
\]
which is exactly the equation in \eqref{koide-param-eq:n-positive-solutions}.
Both embeddings of the generator are positive, hence $a>0$; and
$\alpha_{n,\mathbf p}\geq\sqrt{m_n}$ and $k\geq0$ imply $q\geq0$.

Conversely, a solution in \eqref{koide-param-eq:n-positive-solutions} has
\(\Norm(a+q\sqrt3)=m_n>0\).  Since $a>0$ and $q\geq0$, the inequality
$a^2-3q^2=m_n>0$ gives $a>\sqrt3q$; hence both embeddings are positive.
Its principal ideal has unique HNF data $(C,D,E)$.  The quotient of
 $a+q\sqrt3$ by the canonical generator of that ideal is a totally positive
 unit, hence is $\epsilon^k$ for a unique $k\in\ZZ$.  Lemma
 \ref{koide-param-lem:integral-seed-descent} supplies the completely integral recovery
 of $k$: iterate \eqref{koide-param-eq:inverse-descent} precisely while $a\leq2q$;
 the number of steps is $k$, and the terminal pair is
 $(a_{n,\mathbf p,0},q_{n,\mathbf p,0})$.  In particular $k\geq0$.  This
 proves surjectivity and the stated inverse.  HNF uniqueness, seed
 uniqueness, and unit-exponent uniqueness prove multiplicity one.

It remains to prove the chamber wall.  If a point with $n\geq1$ had
$q\geq a+4n$, then $a\leq q-4n$, and the unchanged equation would give
\[
 3q^2+18n^2-2=a^2\leq(q-4n)^2.
\]
After moving terms and dividing by $2$, this is
\[
 q^2+4qn+n^2\leq1.
\]
For nonnegative integers with $n\geq1$, equality could occur only at
$(q,n)=(0,1)$, but that pair does not satisfy $q\geq a+4n$ with $a>0$.
Thus the assumed wall inequality is impossible.
\end{proof}

The congruence needed by the original inverse map to the $(x,y,z)$
coordinates is retained automatically: reducing
$a^2-3q^2-18n^2=-2$ modulo $2$ gives
\begin{equation}
 a\equiv q\pmod2.
 \label{koide-param-eq:parity}
\end{equation}
No parity class has been lost or added in
Theorem~\ref{koide-param-thm:integral-parametrization}.

\begin{proposition}[The face $n=0$]
\label{koide-param-prop:n-zero}
All points satisfying the original equation and chamber with $n=0$ are
\begin{equation}
 a_k+q_k\sqrt3=(1+\sqrt3)\epsilon^k,\qquad k\geq1.
 \label{koide-param-eq:n-zero-family}
\end{equation}
They have height $4a_k$ and occur with multiplicity one.
\end{proposition}

\begin{proof}
When $n=0$, the norm is $-2$.  Its parity implies $a\equiv q\pmod2$, and
therefore
\[
 \frac{a+q\sqrt3}{1+\sqrt3}
 =\frac{3q-a}{2}+\frac{a-q}{2}\sqrt3\in\OO.
\]
The quotient has norm one, so it is $\pm\epsilon^k$ by
Lemma~\ref{koide-param-lem:ring-units}.  The conditions $a>0$ and $q\geq0$ select the
positive sign.  At $k=0$ one gets $(a,q)=(1,1)$, which is excluded by the
strict chamber wall $a>q$.  At every $k\geq1$, the recurrence
\eqref{koide-param-eq:unit-recurrence} gives
$a_{k+1}-q_{k+1}=a_k+q_k>0$.  For $k\leq-1$, direct use of
$a=(t-2/t)/2$ with $t=(1+\sqrt3)\epsilon^k\leq\sqrt3-1<\sqrt2$
gives $a<0$.  Hence precisely the exponents $k\geq1$ occur.
\end{proof}

\section{Exact divisor and sharp-height formulas}

The number of unit-orbits at a fixed positive $n$ has an especially simple
divisor formula.

\begin{proposition}[Orbit divisor sum and shifted factorization]
\label{koide-param-prop:orbit-divisor}
For $n\geq1$, the number of orbits under
\((a,q)\mapsto(2a+3q,a+2q)\) among the solutions with $a>0$ and
$q\geq0$ is
\begin{equation}
 \rho(m_n)=\sum_{d\mid18n^2-2}\chi_{12}(d)
 =\rho(3n-1)\rho(3n+1).
 \label{koide-param-eq:orbit-divisor-sum}
\end{equation}
\end{proposition}

\begin{proof}
The first equality is Theorem~\ref{koide-param-thm:integral-parametrization} and
Proposition~\ref{koide-param-prop:HNF-ideals}.  For the second, write
\[
 m_n=2(3n-1)(3n+1).
\]
The two displayed linear factors have greatest common divisor dividing
$2$.  The function $\rho=1*\chi_{12}$ is multiplicative and
$\rho(2^j)=1$ for all $j\geq0$.  Removing all powers of $2$ therefore
makes the two factors coprime without changing any of the three $\rho$
values, proving the product identity.
\end{proof}

For $\mathbf p\in\mathcal P(m_n)$ define its exact unit phase by
\begin{equation}
 \tau_{n,\mathbf p}
 =\frac{\log(\alpha_{n,\mathbf p}/\sqrt{m_n})}{\log\epsilon},
 \qquad0\leq\tau_{n,\mathbf p}<1.
 \label{koide-param-eq:unit-phase}
\end{equation}
For $n\geq1$ put
\begin{equation}
 Y_n(X)=\frac{X-17n}{4}.
 \label{koide-param-eq:Y-n}
\end{equation}
Whenever $Y_n(X)\geq\sqrt{m_n}$, put
\begin{equation}
 L_n(X)=
 \frac1{\log\epsilon}
 \log\!\left(
 \frac{Y_n(X)+\sqrt{Y_n(X)^2-m_n}}{\sqrt{m_n}}
 \right).
 \label{koide-param-eq:L-n}
\end{equation}
Let
\begin{equation}
 \mathcal N(X)=
 \{n\in\ZZ_{\geq1}:Y_n(X)\geq\sqrt{m_n}\}.
 \label{koide-param-eq:N-X}
\end{equation}
This is the exact finite range
\begin{equation}
 X\geq17n+4\sqrt{18n^2-2};
 \label{koide-param-eq:n-boundary}
\end{equation}
no asymptotic replacement of this boundary is made.

For the $n=0$ face define
\begin{equation}
 R_0(X)=
 \max\!\left\{0,
 \left\lfloor
 \frac{
 \log\!\left(
 \dfrac{X/4+\sqrt{(X/4)^2+2}}{1+\sqrt3}
 \right)}{\log\epsilon}
 \right\rfloor\right\}.
 \label{koide-param-eq:R0-exact}
\end{equation}

\begin{theorem}[Exact sharp-height floor formula]
\label{koide-param-thm:exact-floor}
For every real $X\geq0$,
\begin{equation}
 R(X)=R_0(X)+
 \sum_{n\in\mathcal N(X)}
 \sum_{\mathbf p\in\mathcal P(m_n)}
 \left(\left\lfloor L_n(X)-\tau_{n,\mathbf p}\right\rfloor+1\right).
 \label{koide-param-eq:exact-floor-formula}
\end{equation}
Every equality case in $4a+17n\leq X$ is included by the floor in
\eqref{koide-param-eq:exact-floor-formula}.
\end{theorem}

\begin{proof}
For a fixed positive $n$ and parameter $\mathbf p$, put
$t_k=\alpha_{n,\mathbf p}\epsilon^k$.  Since
$t_k\geq\sqrt{m_n}$,
\begin{equation}
 a_{n,\mathbf p,k}=\frac12\left(t_k+\frac{m_n}{t_k}\right)
 \label{koide-param-eq:a-from-t}
\end{equation}
is nondecreasing, and is strictly increasing after the possible initial
case $q=0$.  The unchanged height inequality is
$a_{n,\mathbf p,k}\leq Y_n(X)$.  It has no solution unless
$Y_n(X)\geq\sqrt{m_n}$.  In that range, solving the quadratic inequality
in $t_k$ gives
\begin{equation}
 t_k\leq Y_n(X)+\sqrt{Y_n(X)^2-m_n}.
 \label{koide-param-eq:t-height-bound}
\end{equation}
After division by $\sqrt{m_n}$ and logarithms, this is precisely
\begin{equation}
 k+\tau_{n,\mathbf p}\leq L_n(X).
 \label{koide-param-eq:k-phase-bound}
\end{equation}
Because $L_n(X)\geq0$ and $0\leq\tau_{n,\mathbf p}<1$, the number of
integers $k\geq0$ satisfying \eqref{koide-param-eq:k-phase-bound} is exactly
$\lfloor L_n(X)-\tau_{n,\mathbf p}\rfloor+1$; this expression is zero,
not negative, when $L_n(X)<\tau_{n,\mathbf p}$.  Equality is retained.
Theorem~\ref{koide-param-thm:integral-parametrization} now proves the positive-$n$
sum.

For $n=0$, put $t=(1+\sqrt3)\epsilon^k$.  Since the norm is $-2$,
$a_k=(t-2/t)/2$.  The inequality $4a_k\leq X$ is equivalent to
\[
 t\leq X/4+\sqrt{(X/4)^2+2}.
\]
Counting the integers $k\geq1$ gives \eqref{koide-param-eq:R0-exact} and completes
the proof.
\end{proof}

Write
\begin{equation}
 J_n(X)=\lfloor L_n(X)\rfloor,
 \qquad
 \lambda_n(X)=L_n(X)-J_n(X)\in[0,1),
 \label{koide-param-eq:J-lambda}
\end{equation}
and define the exact phase distribution
\begin{equation}
 A_n(t)=\#\{\mathbf p\in\mathcal P(m_n):
                     \tau_{n,\mathbf p}\leq t\},
 \qquad0\leq t<1,
 \label{koide-param-eq:A-n}
\end{equation}
with discrepancy
\begin{equation}
 \Delta_n(t)=A_n(t)-t\rho(m_n).
 \label{koide-param-eq:Delta-n}
\end{equation}
Since
\(\lfloor J+\lambda-\tau\rfloor+1
=J+\mathbf1_{\{\tau\leq\lambda\}}\), Theorem~\ref{koide-param-thm:exact-floor}
gives the exact identity
\begin{align}
 R(X)-C_0X
 &=R_0(X)
 +\left[
   \sum_{n\in\mathcal N(X)}
     \rho(m_n)L_n(X)-C_0X
  \right]
 +\sum_{n\in\mathcal N(X)}
    \Delta_n(\lambda_n(X)).
 \label{koide-param-eq:exact-discrepancy-splitting}
\end{align}
This identity includes the two unbounded coordinate faces and every sharp
height boundary point.  Formula \eqref{koide-param-eq:R0-exact} also gives
$R_0(X)=O(\log(2+X))$, so the $n=0$ face is below the requested
$X^{1/2+\varepsilon}$ scale; the two positive-$n$ brackets are the genuine
analytic terms.

\section{The phase term as exact Hecke-character data}

The last sum in \eqref{koide-param-eq:exact-discrepancy-splitting} is not determined
by the divisor count $\rho(m_n)$.  It records where the generator of each
ideal lies inside a fundamental interval for the positive units.

For $\ell\in\ZZ$ and a nonzero integral ideal $\mathfrak a=(\gamma)$,
define
\begin{equation}
 \xi_\ell(\mathfrak a)=
 \exp\!\left(
  \frac{\pi i\ell}{\log\epsilon}
  \log\left|\frac{\gamma}{\gamma'}\right|
 \right).
 \label{koide-param-eq:Hecke-character}
\end{equation}
This is independent of the generator: replacing $\gamma$ by
$\pm\epsilon^j\gamma$ changes the logarithm by $2j\log\epsilon$, and
the exponential by $e^{2\pi i\ell j}=1$.  It is multiplicative, hence is
the unramified ideal part of a unitary Hecke character of $K$.  For
$m=m_n$ and the canonical generator,
\begin{equation}
 \xi_\ell(\mathfrak a(C,D,E))
 =e^{2\pi i\ell\tau_{n,(C,D,E)}}.
 \label{koide-param-eq:character-phase}
\end{equation}
Define its norm-$m$ coefficient by
\begin{equation}
 S_\ell(m)=\sum_{\Norm\mathfrak a=m}\xi_\ell(\mathfrak a).
 \label{koide-param-eq:S-ell}
\end{equation}
Then $S_0(m)=\rho(m)$, whereas the nonzero $S_\ell(m_n)$ are the Fourier
coefficients of the phase multiset.

The coefficient $S_\ell(m)$ also has a completely explicit Euler-factor
description.  It is multiplicative in $m$.  If $p\geq5$ is split, choose
one prime ideal $\mathfrak p$ above $p$ and put
$\omega_{\ell,p}=\xi_\ell(\mathfrak p)$.  Since
$\mathfrak p\overline{\mathfrak p}=(p)$ and
$\xi_\ell((p))=1$, one has
\begin{equation}
 S_\ell(p^e)=\sum_{j=0}^{e}\omega_{\ell,p}^{2j-e}
 \qquad(\chi_{12}(p)=1).
 \label{koide-param-eq:split-Hecke-factor}
\end{equation}
If $p\geq5$ is inert, then
\begin{equation}
 S_\ell(p^e)=
 \begin{cases}1,&2\mid e,\\0,&2\nmid e.
 \end{cases}
 \qquad(\chi_{12}(p)=-1).
 \label{koide-param-eq:inert-Hecke-factor}
\end{equation}
Finally the unique primes above $2$ and $3$ are generated by
$1+\sqrt3$ and $\sqrt3$, respectively.  Since
\[
 \left|\frac{1+\sqrt3}{1-\sqrt3}\right|=\epsilon,
 \qquad
 \left|\frac{\sqrt3}{-\sqrt3}\right|=1,
\]
their factors are
\begin{equation}
 S_\ell(2^e)=(-1)^{\ell e},
 \qquad S_\ell(3^e)=1.
 \label{koide-param-eq:ramified-Hecke-factors}
\end{equation}
These formulas enumerate every ideal factor with multiplicity one.  At
$\ell=0$ they reduce prime by prime to the divisor sum $\rho(m)$; at
$\ell\ne0$ they show exactly what angular information the unweighted
divisor sum omits.

They also retain the shifted factorization of $m_n$ at every Fourier mode.
Writing $3n-1=2^r u$ and $3n+1=2^s v$ with $u,v$ odd, one has
$\gcd(u,v)=1$.  Multiplicativity and
\eqref{koide-param-eq:ramified-Hecke-factors} therefore give the exact identity
\begin{equation}
 S_\ell(18n^2-2)
 =(-1)^\ell S_\ell(3n-1)S_\ell(3n+1).
 \label{koide-param-eq:shifted-Hecke-factorization}
\end{equation}
Indeed, the exponent of $2$ on the left is $1+r+s$, while the product on
the right has total exponent $r+s$; all odd prime factors are disjoint.
Thus the angular obstruction is itself a shifted correlation, now of
nontrivial Hecke coefficients rather than merely of their dimensions.

Here is an endpoint-exact Fourier statement.  Suppose $0<\lambda<1$, and
let
\begin{equation}
 b_n(t)=\#\{\mathbf p\in\mathcal P(m_n):\tau_{n,\mathbf p}=t\}.
 \label{koide-param-eq:b-n}
\end{equation}
Fej\'er summation of the Fourier series of
$\mathbf1_{[0,\lambda]}-\lambda$ gives
\begin{align}
 \Delta_n(\lambda)
 &=\frac12\bigl(b_n(0)+b_n(\lambda)\bigr)\notag\\
 &\quad+
 \lim_{M\to\infty}
 \sum_{0<|\ell|\leq M}
 \left(1-\frac{|\ell|}{M+1}\right)
 \frac{1-e^{-2\pi i\ell\lambda}}{2\pi i\ell}
 S_\ell(m_n).
 \label{koide-param-eq:Fejer-exact}
\end{align}
The two half-weights correct the midpoint convention at the circular
endpoint $0$ and at $\lambda$.  When $\lambda=0$, the exact value is
simply $\Delta_n(0)=b_n(0)$.

Combining this with \eqref{koide-param-eq:shifted-Hecke-factorization} makes the
aggregate obstruction a literal shifted correlation.  Define the exact
endpoint contribution
\begin{equation}
 \mathcal E(X)=
 \sum_{\substack{n\in\mathcal N(X)\\\lambda_n(X)=0}}b_n(0)
 +\frac12
 \sum_{\substack{n\in\mathcal N(X)\\\lambda_n(X)>0}}
       \bigl(b_n(0)+b_n(\lambda_n(X))\bigr).
 \label{koide-param-eq:endpoint-aggregate}
\end{equation}
Since $\mathcal N(X)$ is finite, summation over $n$ commutes with the
Fej\'er limit and gives
\begin{align}
 \sum_{n\in\mathcal N(X)}\Delta_n(\lambda_n(X))
 &=\mathcal E(X)\notag\\
 &\quad+\lim_{M\to\infty}
 \sum_{0<|\ell|\leq M}
 \left(1-\frac{|\ell|}{M+1}\right)
 \frac{(-1)^\ell}{2\pi i\ell}\notag\\
 &\qquad\times
 \sum_{\substack{n\in\mathcal N(X)\\\lambda_n(X)>0}}
 \bigl(1-e^{-2\pi i\ell\lambda_n(X)}\bigr)
 S_\ell(3n-1)S_\ell(3n+1).
 \label{koide-param-eq:exact-shifted-Hecke-correlation}
\end{align}
Every coefficient and endpoint in this identity comes from the original
weak height inequality; no smoothing has been inserted.

Equations \eqref{koide-param-eq:exact-discrepancy-splitting} and
\eqref{koide-param-eq:Fejer-exact} identify the analytic problem without discarding
the original height.  Its radial part is the sharply weighted shifted
correlation
\begin{equation}
 \sum_{n\in\mathcal N(X)}
 \rho(3n-1)\rho(3n+1)L_n(X),
 \label{koide-param-eq:radial-shifted-correlation}
\end{equation}
and its angular part requires cancellation in the polynomially sampled
Hecke coefficients
\begin{equation}
 S_\ell(18n^2-2),
 \label{koide-param-eq:polynomial-Hecke-samples}
\end{equation}
uniformly over the Fourier modes needed to approximate the discontinuous
threshold $\tau\leq\lambda_n(X)$.  The endpoint
\eqref{koide-param-eq:n-boundary} gives $L_n(X)=0$ and has a square-root singularity
in its derivative; it is not a fixed smooth weight.  A bound for the
unweighted orbit total $\rho(m_n)$ therefore does not bound the actual
sharp-height count.

\section{An unconditional effective exponent for the exact sector}

This section records the power saving that is presently justified without
assuming the square-root estimate.

Let
\begin{equation}
 F(a,q,n)=a^2-3q^2-18n^2,
 \qquad
 \mathscr X=\{F=-2\},
 \qquad
 \mathbf G=\operatorname{Spin}(F).
 \label{koide-param-eq:group-variety}
\end{equation}
Let $x_0=(1,1,0)$.  Its orthogonal complement is
\begin{equation}
 x_0^\perp=\{(3s,s,t):s,t\in\QQ\},
 \qquad
 F(3s,s,t)=6(s^2-3t^2),
 \label{koide-param-eq:stabilizer-complement}
\end{equation}
so the identity component of its stabilizer is the norm-one torus of
$\QQ(\sqrt3)$.  Its rational character group is zero because Galois
conjugation acts by $-1$ on its geometric character lattice.

The form $F$ is anisotropic over $\QQ$.  Indeed, suppose a primitive
integral triple satisfies $F(a,q,n)=0$.  Modulo $3$ one first gets
$3\mid a$.  Writing $a=3a_1$ and dividing the equation by $3$ gives
$3a_1^2-q^2-6n^2=0$, hence $3\mid q$.  Write $q=3q_1$ and divide by
$3$ again to get
\begin{equation}
 a_1^2-3q_1^2-2n^2=0.
\end{equation}
Modulo $3$, the sum $a_1^2+n^2$ is zero, which forces
$3\mid a_1$ and $3\mid n$.  Thus $3$ divides $a,q,n$, contradicting
primitivity.  In particular the associated arithmetic Fuchsian quotient
is cocompact.

For $X>0$ let $B_X$ be the exact subset of $\mathscr X(\RR)$ cut out by
\begin{equation}
 a>0,\quad q\geq0,\quad n\geq0,\quad a+4n>q,
 \quad4a+17n\leq X.
 \label{koide-param-eq:B-X}
\end{equation}
With the gauge measure
\begin{equation}
 dm_\infty=\frac{dq\,dn}{2a}\qquad(a>0),
 \label{koide-param-eq:gauge-measure}
\end{equation}
put
\begin{equation}
 I_\infty=\frac1{2\sqrt3}
 \log\!\left(\frac3{2\sqrt2}\right).
 \label{koide-param-eq:I-infty}
\end{equation}

\begin{lemma}[Exact-sector volume with a power-saving remainder]
\label{koide-param-lem:volume}
As $X\to\infty$,
\begin{equation}
 m_\infty(B_X)=I_\infty X+O(1).
 \label{koide-param-eq:volume-linear-O1}
\end{equation}
Moreover, the family $B_X$ is effectively well-rounded in the sense of
Benoist--Oh, Definition~13.1~\cite[p.~1931]{BenoistOh_2012}.
\end{lemma}

\begin{proof}
Introduce exact polar coordinates on the negative plane by
\begin{equation}
 \sqrt3q=r\cos\theta,\qquad
 3\sqrt2n=r\sin\theta,\qquad
 0\leq\theta\leq\pi/2.
 \label{koide-param-eq:polar}
\end{equation}
On $F=-2$ one has $r=\sqrt{a^2+2}$ and
\begin{equation}
 dm_\infty=\frac1{6\sqrt6}\,da\,d\theta.
 \label{koide-param-eq:polar-measure}
\end{equation}
Write $b(\theta)=17\sin\theta/(3\sqrt2)$.  In the positive orthant the
outer boundary is the unique nonnegative solution $A_X(\theta)$ of
\begin{equation}
 4A_X(\theta)+b(\theta)\sqrt{A_X(\theta)^2+2}=X.
 \label{koide-param-eq:A-boundary}
\end{equation}
For all sufficiently large $X$, uniformly in $\theta$,
\begin{equation}
 A_X(\theta)=\frac{X}{4+b(\theta)}+O(X^{-1}).
 \label{koide-param-eq:A-asymptotic}
\end{equation}
To see this without a formal expansion, subtract
$(4+b)A_X$ from \eqref{koide-param-eq:A-boundary} and use
\[
 \sqrt{A_X^2+2}-A_X
 =\frac2{\sqrt{A_X^2+2}+A_X}=O(X^{-1}),
\]
where \eqref{koide-param-eq:A-boundary} uniformly gives $A_X\asymp X$.
Integration of \eqref{koide-param-eq:A-asymptotic} in
\eqref{koide-param-eq:polar-measure} gives
\[
 \frac{X}{6\sqrt6}\int_0^{\pi/2}
 \frac{d\theta}{4+\frac{17}{3\sqrt2}\sin\theta}+O(X^{-1})
 =I_\infty X+O(X^{-1}).
\]
For the asserted exact value, set
$c=17/(3\sqrt2)$, $d=\sqrt{c^2-16}=1/(3\sqrt2)$, and
$t=\tan(\theta/2)$.  Then
\begin{align*}
 \int_0^{\pi/2}\frac{d\theta}{4+c\sin\theta}
 &=\int_0^1\frac{dt}{2t^2+ct+2}\\
 &=\frac1d
   \left[\log\!\left(\frac{4t+c-d}{4t+c+d}\right)\right]_{t=0}^{t=1}\\
 &=3\sqrt2\log\!\left(\frac3{2\sqrt2}\right).
\end{align*}
Multiplication by $1/(6\sqrt6)$ is exactly the constant in
\eqref{koide-param-eq:I-infty}.
If the strict chamber wall fails, then the calculation in the proof of
Theorem~\ref{koide-param-thm:integral-parametrization}, now over $\RR$, gives
$q^2+4qn+n^2\leq1$.  The removed real subset is therefore compact.
For all large $X$ it contributes a fixed finite measure, proving
\eqref{koide-param-eq:volume-linear-O1}.

It remains to check effective well-roundedness for the actual sector.
Let $U_\delta$ be the image in the fixed vector representation of the
group ball used by Benoist--Oh~\cite[Definition~13.1]{BenoistOh_2012}.
For every linear functional $\ell$ on
$\RR^3$ there is a constant, independent of $x$ and $\delta$, such that
\begin{equation}
 |\ell(xg)-\ell(x)|\leq C_\ell\delta\|x\|
 \qquad(g\in U_\delta).
 \label{koide-param-eq:linear-functional-motion}
\end{equation}
Apply this to the five literal boundary functionals
\[
 a,\qquad q,\qquad n,\qquad a+4n-q,\qquad4a+17n.
\]
On the unbounded part of the positive sheet, $\|x\|\asymp a+1$, and the
last functional is also comparable to $\|x\|$.  Thus membership in the
outer height inequality can change only in a shell of radial width
$O(\delta X)$.  Membership in $q\geq0$ or $n\geq0$ can change only in
angular strips of width $O(\delta)$ about $\theta=\pi/2$ or $\theta=0$.
The boundary $a=0$ is compact.  The boundary $a+4n-q=0$ is compact by
the calculation preceding \eqref{koide-param-eq:volume-linear-O1}; moreover its
restriction to $\mathscr X$ is regular, because proportionality of
$(1,-1,4)$ to $(2a,-6q,-36n)$ is incompatible with $a,q,n\geq0$.
Consequently its group-action boundary tube has measure $O(\delta)$.
By \eqref{koide-param-eq:polar-measure}, all these regions together have measure
$O(\delta X)$.  The same analysis applies to the inner intersection that
defines $B^-_{X,\delta}$.  Hence
\begin{equation}
 m_\infty(B^+_{X,\delta}\setminus B^-_{X,\delta})
 \ll \delta\,m_\infty(B_X),
 \label{koide-param-eq:BO-boundary-condition}
\end{equation}
which is precisely Definition 13.1(2), with boundary exponent $1$, for
the group-action thickenings rather than merely Euclidean tubes.

For Definition~13.1(3) of \cite{BenoistOh_2012}, take the archimedean height to be the
Euclidean norm in the displayed integral basis.  It is bounded below on
$\mathscr X$.  Uniformly for $0<\delta<1$, the thickening
$B^+_{X,\delta}$ lies in a fixed finite enlargement of the polar sector;
there $H_\infty(x)\asymp1+a$ and the measure in a radial interval is
$O(da)$.  Hence for every $s>0$, uniformly in $\delta$,
\[
 \int_{B_X^+}\|u\|^{-s}\,dm_\infty(u)
 \ll
 \begin{cases}
 X^{1-s},&0<s<1,\\
 \log X,&s=1,\\
 1,&s>1.
 \end{cases}
\]
Each line is $O(m_\infty(B_X)^{1-\delta_s})$, for example with
$\delta_s=\min(s/2,1/2)>0$.  This is the remaining condition.  Since
$S=\{\infty\}$, Definition 13.1(1), concerning a finite-place compact
open subgroup, is void.  The choices of strict
or weak inequality on the finitely many piecewise smooth boundary arcs do
not change these neighborhood estimates.
\end{proof}

\begin{theorem}[Unconditional effective conclusion]
\label{koide-param-thm:power-saving}
There exists $\eta>0$, depending only on the fixed form, lattice, chamber,
and height in \eqref{koide-param-eq:original-problem}, such that
\begin{equation}
 R(X)=C_0X+O(X^{1-\eta}).
 \label{koide-param-eq:power-saving-again}
\end{equation}
\end{theorem}

\begin{proof}
Use the standard vector representation
$\rho:\operatorname{Spin}(F)\to\operatorname{SO}(F)\subset
\operatorname{GL}_3$, whose kernel is finite, and the displayed integral
basis of the exact lattice $\ZZ^3$.  Let
\[
 \mathbf G_{\ZZ}
 =\{g\in\mathbf G(\QQ):\rho(g)\ZZ^3=\ZZ^3\}.
\]
This is the integral group attached to the chosen representation and
lattice in the Benoist--Oh setup~\cite[pp.~1892--1893]{BenoistOh_2012};
take their congruence subgroup
$\Gamma_S$ to be the level-one group $\mathbf G_{\ZZ}$ itself.  The group
$\mathbf G$ is connected,
simply connected, and almost $\QQ$-simple.  It is noncompact over $\RR$
because $F$ has signature $(1,2)$.  Thus $S=\{\infty\}$ contains every
noncompact archimedean place and $G_S$ is noncompact.  The orbit
$x_0\mathbf G$ is the affine symmetric variety $\mathscr X$, and its
stabilizer has no nontrivial rational character by
\eqref{koide-param-eq:stabilizer-complement}.  These verify every standing hypothesis
on pp.~1892--1893 of \cite{BenoistOh_2012}.  Their Theorem~12.5 supplies the required
effective equidistribution of the symmetric-stabilizer translates.

The Borel--Harish-Chandra finiteness theorem quoted in \cite[p.~1892]{BenoistOh_2012}
p.~1892 says that $\mathscr X(\ZZ)$ is a finite union of
$\mathbf G_{\ZZ}$-orbits.  A rational representative of any such orbit
again has a symmetric stabilizer with no rational
character: its orthogonal complement is anisotropic over $\QQ$, since an
isotropic vector in it would be an isotropic vector for $F$.  Apply
Benoist--Oh, Theorem~13.2~\cite[pp.~1931--1932]{BenoistOh_2012}, to each orbit and to the exact effectively
well-rounded family $B_X$ of Lemma~\ref{koide-param-lem:volume}.  For each orbit it
gives its invariant-volume main term with error
$O(m_\infty(B_X)^{1-\eta_j})$ for some $\eta_j>0$.  There are finitely
many orbits, so their sum has the same form with
$\eta=\min_j\eta_j>0$.

All invariant measures on the same real homogeneous variety are
proportional.  Hence the sum of the finitely many orbit main terms is
$K\,m_\infty(B_X)$ for one fixed $K>0$.  Lemma~\ref{koide-param-lem:volume} and the
already established asymptotic $R(X)\sim C_0X$ force
$K I_\infty=C_0$.  Therefore
\[
 K\,m_\infty(B_X)=K(I_\infty X+O(1))=C_0X+O(1).
\]
Since $m_\infty(B_X)\asymp X$, summing the effective errors proves
\eqref{koide-param-eq:power-saving-again}; the $O(1)$ term is absorbed after decreasing
$\eta$ if necessary.  The theorem is applied first to the sequence of
positive integral $X$.  For real $X$, the original height $4a+17n$ is an
integer, so $R(X)=R(\lfloor X\rfloor)$; replacing $X$ by $\lfloor X\rfloor$
changes $C_0X$ by $O(1)$ and proves the same estimate.
\end{proof}

\begin{remark}[What exponent has actually been obtained]
Theorem~\ref{koide-param-thm:power-saving} is an unconditional exponent strictly
smaller than $1$: the error is $O(X^{1-\eta})$ for some $\eta>0$.
Neither the statement nor the proof of \cite[Theorem~13.2]{BenoistOh_2012} identifies an exponent
$1-\eta\leq1/2$ for this sharp sector.  Tracing their Sobolev and mixing
constants would make a positive $\eta$ effective in principle, but that
calculation has not been carried out in this section and cannot be replaced
by the assertion $\eta=1/2$.
\end{remark}

\section{Exact obstruction to the square-root claim}

There are two complementary descriptions of the remaining obstruction;
the first is an exact identity, while the second explains the relevant
spectral mechanisms without asserting an unproved intertwining formula.

First, the elementary ideal description gives the literal identity
\eqref{koide-param-eq:exact-discrepancy-splitting}.  Thus the square-root claim is
equivalent to a bound of that size for the displayed sum of the radial
bracket, the aggregate phase discrepancy, and $R_0(X)$.  Controlling the
radial and phase terms separately is a sufficient route, though the exact
identity also permits cancellation between them.  Formula
\eqref{koide-param-eq:Fejer-exact} shows that a separate estimate of the latter entails
uniform cancellation in the nonzero archimedean Hecke modes
\eqref{koide-param-eq:polynomial-Hecke-samples}.  Merely using
\(\#\mathcal P(m_n)=\rho(m_n)\), or replacing each phase distribution by
Lebesgue measure, deletes a term that is exactly present in $R(X)$.

Second, $F$ is rationally anisotropic, so $\mathbf G$ gives a cocompact
arithmetic Fuchsian quotient.  A negative-norm vector represents an
oriented geodesic, and its stabilizer is the split real torus appearing in
\eqref{koide-param-eq:stabilizer-complement}.  In the standard radial torus-orbit
counting kernel, the spectral expansion has two possible sources larger
than the square-root scale:
\begin{enumerate}
 \item an exceptional eigenvalue $\lambda=s(1-s)<1/4$ can contribute
 $c_sX^s$ with $s>1/2$, unless the exact torus period and exact sector
 transform make $c_s=0$;
 \item after exceptional terms are absent or removed, estimating the
 tempered spectral sum for a discontinuous cutoff at the conjectural
 $X^{1/2+\e}$ scale requires cancellation beyond the general
 effective-mixing argument.
\end{enumerate}
For the present nonradial chamber, one must first construct a
remainder-preserving sector spectral expansion before these candidate
terms can even be asserted with their exact coefficients.  An actual
refutation by a secondary term would then require proving both $s>1/2$
and $c_s\ne0$ for this fixed arithmetic quotient and exact chamber.  No
such construction or nonvanishing has been established here.  Conversely,
absence of exceptional terms alone would not estimate the full tempered
sharp-cutoff sum.  This is why the exact parametrization neither proves
nor refutes \eqref{koide-param-eq:square-root-target}, while it does isolate the data a
future proof must bound.

For comparison only, not as a theorem about the present nonradial sector,
the classical hyperbolic lattice-point problem for a conjugacy class has
an unconditional $O(X^{2/3})$ remainder due to Good
\cite[Theorem~4]{Good_1983}, with the modern proof of
Chatzakos--Petridis~\cite[Theorem~1.1 and Conjecture~5.7]{ChatzakosPetridis_2016},
and a conjectural $O_\e(X^{1/2+\e})$ scale.  That
published $O(X^{2/3})$ theorem counts a different radial kernel.  An exact
intertwining of its kernel with the simultaneous chamber and height in
\eqref{koide-param-eq:original-problem} has not been proved, so $O(X^{2/3})$ is not
claimed here for $R(X)$.

The independent finite replay used to check the HNF multiplicities, seed
reduction, inverse descent, and divisor factorization for \(1\leq n\leq200\)
is retained at
\path{supporting_materials/workbench/evidence/koide_exact_checkpoint_20260904/parametrization/verify_parametrization.py};
its SHA--256 digest is
\texttt{061c5e22258628169fd1ecb8bcc5e7cf\allowbreak%
8b8d19bfdb918650668c3e90f2f9730d}.
That replay is a finite consistency check.  The proofs above, rather than the
finite range, establish the parametrization for every admissible integer.

\section{Boundary ledger}

For clarity, every boundary used above has the following exact treatment.
\begin{itemize}
 \item $a>0$ is strict.  There is no omitted integral solution with
 $a=0$, because the original equation would give a multiple of $3$ equal
 to $2$.
 \item $q=0$ is included.  It corresponds to phase $\tau=0$, is included
 by $\tau\leq\lambda$, and is retained in the endpoint term
 $b_n(0)/2$ in \eqref{koide-param-eq:Fejer-exact}.
 \item $n=0$ is included and is counted exactly by $R_0(X)$; the sole
 solution $(a,q,n)=(1,1,0)$ on $a+4n=q$ is excluded.
 \item $a+4n=q$ is strict.  Its real intersection is bounded and its sole
 integral point in the nonnegative chamber is $(1,1,0)$.
 \item $4a+17n=X$ is included.  Equality is included in
 \eqref{koide-param-eq:t-height-bound}, \eqref{koide-param-eq:k-phase-bound}, and every floor in
 \eqref{koide-param-eq:exact-floor-formula}.
 \item The canonical unit interval is half-open,
 $[\sqrt{m_n},\epsilon\sqrt{m_n})$.  Consequently no generator is counted
 twice at a unit boundary.
 \item Conjugate ideals are distinct unless the ideal is fixed by
 conjugation.  They yield distinct canonical phases and distinct
 $q\geq0$ unit-orbits; no division by $2$ is made.  When $q=0$, the ideal
 is self-conjugate and still occurs once by HNF uniqueness.
\end{itemize}

% Integrated 4 September 2026 from the exact 10^8 census and independent replay.
% This section retains the literal ternary equation, chamber, strict wall, height,
% and original-coordinate inverse.

\section{Exact unit-orbit census through height \(10^8\)}
\label{koide-count-sec-census-1e8}

The census retains exactly the same equation, chamber, strict wall, and
height as the preceding theorems:
\begin{equation}
 (a,q,n)\in\mathbb Z^3,\qquad
 a^2-3q^2-18n^2=-2,\qquad
 a>0,\quad q,n\geq0,\quad a+4n>q,\qquad
 4a+17n\leq X.
 \label{koide-count-eq-census-domain}
\end{equation}
It does not infer an asymptotic from finite data.  Its exact fixed-\(n\)
equation and factorization are
\begin{equation}
 a^2-3q^2=m_n,\qquad
 m_n=18n^2-2=2(3n-1)(3n+1).
 \label{koide-count-eq-census-fixed-n}
\end{equation}
For each \(n\geq1\), every ideal of norm \(m_n\) is enumerated by its
Hermite-normal-form triple
\begin{equation}
 (C,D,E)\in\mathbb Z_{>0}^2\times\mathbb Z,\qquad
 C^2D=m_n,\qquad
 0\leq E<D,\qquad E^2\equiv3\pmod D.
 \label{koide-count-eq-census-hnf}
\end{equation}
The generator is reduced by exact multiplication by the fundamental unit
\(2+\sqrt3\) to the unique seed \(a_0>2q_0\geq0\).  Its forward unit
orbit is
\begin{equation}
 \binom{a_{k+1}}{q_{k+1}}
 =
 \begin{pmatrix}2&3\\1&2\end{pmatrix}
 \binom{a_k}{q_k},\qquad k\geq0.
 \label{koide-count-eq-census-unit-step}
\end{equation}
Each emitted point is reconstructed by
\begin{equation}
 x=\frac{a+4n-q}{2},\qquad
 y=\frac{a+4n+q}{2},\qquad
 z=4a+17n,
 \label{koide-count-eq-census-inverse}
\end{equation}
and is rechecked in both coordinate systems and in the original polynomial
\(\Delta(x,y,z)=-1\).  Thus the accelerated enumeration does not replace
the original objects by a floating-point or rescaled surrogate.

\subsection{Complete result and checkpoint table}

The complete integer-domain audit for
\(1\leq X\leq100{,}000{,}000\) gives
\begin{equation}
 \boxed{R(100{,}000{,}000)=3{,}536{,}002.}
 \label{koide-count-eq-census-final-count}
\end{equation}
The endpoint in \(n\) is \(2{,}943{,}725\); the enumeration contains
\(2{,}327{,}513\) canonical unit orbits.  The following checkpoints are
exact integer counts.  The residuals use the proved constant
\(C_0\) in \eqref{koide-count-eq-final-constant}.
\begin{center}
\begin{tabular}{r@{\qquad}r@{\qquad}r@{\qquad}r}
\toprule
\(X\)&\(R(X)\)&\(R(X)-C_0X\)&\((R(X)-C_0X)/\sqrt X\)\\
\midrule
200,000&7,109&\(36.590440053629\)&\(0.081818711287\)\\
250,000&8,854&\(13.488050067037\)&\(0.026976100134\)\\
500,000&17,695&\(13.976100134073\)&\(0.019765190359\)\\
1,000,000&35,381&\(18.952200268147\)&\(0.018952200268\)\\
2,000,000&70,637&\(-87.095599463706\)&\(-0.061585888992\)\\
5,000,000&176,943&\(132.761001340734\)&\(0.059372524752\)\\
10,000,000&353,598&\(-22.477997318532\)&\(-0.007108166877\)\\
20,000,000&707,458&\(217.044005362937\)&\(0.048532515010\)\\
50,000,000&1,768,230&\(127.610013407342\)&\(0.018046781166\)\\
100,000,000&3,536,002&\(-202.779973185317\)&\(-0.020277997319\)\\
\bottomrule
\end{tabular}
\end{center}

\subsection{Certified extrema over every integer}

Write \(E(X)=R(X)-C_0X\).  Between consecutive occupied heights,
\(R(X)\) is constant and both \(E(X)\) and
\(E(X)/\sqrt X\) are strictly decreasing, since
\begin{equation}
 \frac{d}{dX}\left(\frac{R}{\sqrt X}-C_0\sqrt X\right)
 =-\frac{R}{2X^{3/2}}-\frac{C_0}{2\sqrt X}<0.
 \label{koide-count-eq-census-monotonicity}
\end{equation}
Checking both integer endpoints of every constant-count interval therefore
certifies extrema over every integer in the audited range, rather than at
selected checkpoints.  On \(1\leq X\leq10^8\), the maximum signed residual
is
\begin{equation}
 \begin{aligned}
 E(77{,}078{,}404)
 &=721.7934249571297569656681553538691093379\ldots,\\
 R(77{,}078{,}404)&=2{,}726{,}372,
 \end{aligned}
 \label{koide-count-eq-census-max-positive}
\end{equation}
and the minimum signed, hence largest absolute, residual is
\begin{equation}
 \begin{aligned}
 E(32{,}461{,}294)
 &=-812.8300691488068897091966006797491655933\ldots,\\
 R(32{,}461{,}294)&=1{,}147{,}085.
 \end{aligned}
 \label{koide-count-eq-census-max-negative}
\end{equation}
The maximum of
\(\lvert E(X)\rvert/\sqrt X\) over the full audited range is attained at
\(X=160\), where
\[
 \frac{\lvert E(160)\rvert}{\sqrt{160}}
 =0.185157076934082141384010185515317233816\ldots.
\]
On \(10{,}000\leq X\leq10^8\), the maximum is
\begin{equation}
 \begin{aligned}
 \frac{\lvert E(867{,}240)\rvert}{\sqrt{867{,}240}}
 &=0.171400192808814559728949658032345552502\ldots,\\
 R(867{,}240)&=30{,}827,
 \end{aligned}
 \label{koide-count-eq-census-tail-extremum}
\end{equation}
and the root entering at that height is
\begin{equation}
 (a,q,n)=(199997,115061,3956),\qquad
 (x,y,z)=(50380,165441,867240).
 \label{koide-count-eq-census-tail-witness}
\end{equation}
The bound \(\lvert E(X)\rvert\leq0.171\sqrt X\) first fails on this tail at
\(X=867{,}216\), where \(R(X)=30{,}826\) and
\[
 E(X)=159.4663553077412220330449353234319817199\ldots.
\]
No counterexample to the finite statement with constant \(0.172\) occurs
through \(10^8\).  These are bounded computations, not a claim that either
constant persists.

The audit finds \(12{,}916\) sign changes.  The \(3{,}536{,}002\) roots
occupy \(2{,}472{,}015\) distinct heights; the maximum shell multiplicity
is \(35\), at \(z=92{,}923{,}120\), and the largest gap between consecutive
occupied heights is \(564\), from \(71{,}586{,}276\) to \(71{,}586{,}840\).

\subsection{Independent replay and exact algorithm}

An independent pass decompressed and rechecked all \(3{,}536{,}002\)
six-coordinate rows, recomputed every checkpoint, verified the coordinate
maps, equations, ordering, chamber, height, nonduplication, and gzip CRC,
and returned \texttt{PASS}.  It also compared the literal \((n,a)\) scan
through \(X=20{,}000\), every generalized-Pell seed for \(1\leq n\leq1000\),
and the full six-coordinate point set through \(X=200{,}000\).

The reproducibility sources and receipts are retained project-relatively under
\path{supporting_materials/workbench/evidence/koide_exact_checkpoint_20260904/census}.
The two \(10^8\)-range outputs and the full coordinate receipt are in its
\path{final_100000000} subdirectory.  The principal digests are
\begin{center}
\begin{tabular}{>{\raggedright\arraybackslash}p{0.43\linewidth}
                >{\raggedright\arraybackslash}p{0.51\linewidth}}
\toprule
artifact&SHA-256\\
\midrule
\texttt{koide\_fast\_enum.py}&
\texttt{dddb73b828b63be08e2ff0b5346edb32\allowbreak%
31662005f83d5b2448c06fef6ba31017}\\
\texttt{koide\_discrepancy\_100000000.json}&
\texttt{da92f0866074e530b8174a48915caa0f\allowbreak%
65cd4cd8957e209306adc200739d1961}\\
\texttt{koide\_roots\_through\_\allowbreak 100000000.csv.gz}&
\texttt{dfc313492aedb396dcfa10520bfefc699\allowbreak%
ba601357bc9933d921dfa3884310bd4}\\
\texttt{coordinate\_audit\_receipt.json}&
\texttt{3010363108e7b460b7331a92f3263bbd\allowbreak%
77992aedfc3064973750e16d707da80f}\\
\bottomrule
\end{tabular}
\end{center}

The exact algorithm first finds the largest admissible \(n\) by integer
binary search, then factors the unchanged integers \(3n-1\) and \(3n+1\)
with a smallest-prime-factor sieve.  It enumerates every HNF triple in
\eqref{koide-count-eq-census-hnf}, constructs all roots of \(3\) modulo
\(D\) by prime-power lifting and the Chinese remainder theorem, reduces
each generator to the unique seed, and iterates
\eqref{koide-count-eq-census-unit-step} until the literal height fails.
For each \(n\), the independently checked identity
\[
 \#\mathcal P(m_n)=\rho(m_n)=\rho(3n-1)\rho(3n+1)
\]
certifies the orbit multiplicity.  The work is
\[
 X^{1+o(1)}+O(R(X))
\]
integer operations before the optional height sort.  This complexity and
the finite census use no conjectural discrepancy estimate.

% Integrated 4 September 2026; exact in-project provenance is recorded below.
% The algebra, order, lattice, shell, centralizer, and cutoff are retained exactly.
\section{Exact quaternion bridge: order, shell, and hyperbolic conjugacy classes}\label{koide-quat-exact-quaternion-bridge-order-shell-and-hyperbolic-conjugacy-classes}

The proposed quaternionic realization is correct. The displayed lattice is an order in an indefinite rational quaternion division algebra of discriminant \(6\). The parity-restricted ternary shell is exactly the reduced-trace-\(4\), reduced-norm-\(1\) locus in that order. Every resulting element is primitive hyperbolic, with eigenvalues \(2\pm\sqrt3\) and translation length \(2\log(2+\sqrt3)\).

The centralizer statement is correct when written intrinsically as

\[
 \mathcal O\cap\mathbb Q(\gamma)
 =\mathbb Z[\gamma-2]\cong\mathbb Z[\sqrt3].
\]

The entire integral shell is a finite union of conjugacy classes under \(\Gamma=\mathcal O^1\). This gives an exact reduction of the Koide count to a finite sum of sharp, non-radial sector counts inside primitive hyperbolic conjugacy classes of a cocompact arithmetic Fuchsian group.

It does \textbf{not} prove the open estimate

\[
 R(X)-C_0X=O_\varepsilon(X^{1/2+\varepsilon}).
\]

Rather, it converts the count exactly into a sharp sector problem on a finite union of hyperbolic conjugacy classes. This places it beside, but does not prove it equivalent to, the standard hyperbolic conjugacy-class lattice-point problem. Even for the radial Huber cutoff, the general pointwise theorem gives an \(O(X^{2/3})\) remainder after all exceptional spectral terms are included, while the \(O(X^{1/2+\varepsilon})\) remainder is conjectural. The Koide cutoff is a different sharp polyhedral-sector height, so the standard radial theorem cannot be imported without a new transfer or trace-formula argument.

\subsection{Quaternion algebra and reduced invariants}\label{koide-quat-quaternion-algebra-and-reduced-invariants}

Let

\[
 B=(18,-6)_{\mathbb Q}
 =\mathbb Q\langle I,J\rangle/
 (I^2-18,\ J^2+6,\ IJ+JI),
 \qquad K=IJ.
\]

Then

\[
 K^2=-I^2J^2=108.
\]

For

\[
 x=x_0+x_1I+x_2J+x_3K,
\]

the standard involution, reduced trace, and reduced norm are

\[
 \bar x=x_0-x_1I-x_2J-x_3K,
 \qquad \operatorname{trd}(x)=x+\bar x=2x_0,
\]

and

\begin{equation}\label{koide-quat-eq-reduced-norm}
 \operatorname{nrd}(x)=x\bar x
 =x_0^2-18x_1^2+6x_2^2-108x_3^2.
\end{equation}

All signs in \eqref{koide-quat-eq-reduced-norm} are forced by \(\operatorname{nrd}(I)=-18\), \(\operatorname{nrd}(J)=6\), and \(\operatorname{nrd}(K)=-108\).

\subsection{The proposed lattice is an order}\label{koide-quat-the-proposed-lattice-is-an-order}

Put

\[
 E=\frac{I+J}{2},\qquad
 F_0=\frac{I-J}{2},\qquad
 G=\frac K2,
\]

and let

\begin{equation}\label{koide-quat-eq-order-lattice}
 \mathcal O=\mathbb Z1\oplus\mathbb ZE
             \oplus\mathbb ZF_0\oplus\mathbb ZG.
\end{equation}

The displayed elements are a rational basis of \(B\), because

\[
 I=E+F_0,\qquad J=E-F_0,\qquad K=2G.
\]

Direct multiplication gives the complete nontrivial table

\begin{equation}\label{koide-quat-eq-multiplication-table}
\begin{array}{c|ccc}
   &E&F_0&G\\ \hline
 E&3&6-G&6E-3F_0\\
 F_0&6+G&3&3E-6F_0\\
 G&-6E+3F_0&-3E+6F_0&27.
\end{array}
\end{equation}

Every structure constant in \eqref{koide-quat-eq-multiplication-table} is integral. Thus \eqref{koide-quat-eq-order-lattice} is a unital subring and a full rank-four \(\mathbb Z\)-lattice in \(B\), hence an order. In particular, the unital ring sometimes denoted \(\mathbb Z[1,E,F_0,G]\) is exactly the direct sum in \eqref{koide-quat-eq-order-lattice}.

An arbitrary order element is

\begin{equation}\label{koide-quat-eq-order-coordinates}
 x=r+sE+tF_0+uG
  =r+\frac{s+t}{2}I+\frac{s-t}{2}J+\frac u2K
 \qquad(r,s,t,u\in\mathbb Z),
\end{equation}

and \eqref{koide-quat-eq-reduced-norm} gives

\begin{equation}\label{koide-quat-eq-order-trace-norm}
 \operatorname{trd}(x)=2r,
 \qquad
 \operatorname{nrd}(x)
 =r^2-3s^2-12st-3t^2-27u^2.
\end{equation}

The reduced-trace pairing on the basis \((1,E,F_0,G)\) is

\[
 \bigl(\operatorname{trd}(b_ib_j)\bigr)_{i,j}
 =
 \begin{pmatrix}
  2&0&0&0\\
  0&6&12&0\\
  0&12&6&0\\
  0&0&0&54
 \end{pmatrix},
\]

with determinant

\begin{equation}\label{koide-quat-eq-order-discriminant-determinant}
 -11664=-108^2.
\end{equation}

Thus the reduced discriminant of this order is \(108\), not \(6\). The algebra has discriminant \(6\), proved below, so \(\mathcal O\) is nonmaximal and has index \(18\) in a maximal order containing it. This index is not an Eichler level: the order is nonmaximal at the ramified primes \(2\) and \(3\), so it is not an Eichler order in the standard sense. Any optimal-embedding or trace-formula result that assumes a maximal order must account for this exact nonmaximal order instead of replacing \(\mathcal O\).

\subsection{Trace-zero lattice and exact shell bijection}\label{koide-quat-trace-zero-lattice-and-exact-shell-bijection}

Equation~\eqref{koide-quat-eq-order-trace-norm} gives

\[
 \mathcal O^0
 :=\{x\in\mathcal O:\operatorname{trd}(x)=0\}
 =\mathbb ZE\oplus\mathbb ZF_0\oplus\mathbb ZG.
\]

Writing

\[
 q=s+t,\qquad a=s-t,\qquad n=u
\]

gives the exact description

\begin{equation}\label{koide-quat-eq-trace-zero-lattice}
 \boxed{
 \mathcal O^0=
 \left\{
 \frac{qI+aJ+nK}{2}:
 a,q,n\in\mathbb Z,\ a\equiv q\pmod2
 \right\}.}
\end{equation}

The inverse formulas are

\[
 s=\frac{q+a}{2},\qquad
 t=\frac{q-a}{2},\qquad
 u=n,
\]

which are integral exactly when \(a\equiv q\pmod2\).

Define

\[
 F(a,q,n)=a^2-3q^2-18n^2,
 \qquad
 v(a,q,n)=\frac{qI+aJ+nK}{2},
\]

and

\begin{equation}\label{koide-quat-eq-gamma-parametrization}
 \gamma(a,q,n)=2+v(a,q,n).
\end{equation}

Distinct pure quaternion basis elements anticommute, so

\begin{equation}\label{koide-quat-eq-pure-square}
 v(a,q,n)^2
 =\frac{18q^2-6a^2+108n^2}{4}
 =-\frac32F(a,q,n),
\end{equation}

and

\[
 \operatorname{nrd}(v)=\frac32F(a,q,n).
\]

Since \(v\) has reduced trace zero,

\begin{equation}\label{koide-quat-eq-gamma-trace-norm}
 \operatorname{trd}(\gamma)=4,
 \qquad
 \operatorname{nrd}(\gamma)
 =4+\operatorname{nrd}(v)
 =4+\frac32F(a,q,n).
\end{equation}

Therefore \eqref{koide-quat-eq-gamma-parametrization} defines a bijection

\begin{equation}\label{koide-quat-eq-shell-bijection}
 \boxed{
 \left\{(a,q,n)\in\mathbb Z^3:
 F(a,q,n)=-2,\ a\equiv q\pmod2\right\}
 \xrightarrow{\ \sim\ }
 \left\{\gamma\in\mathcal O:
 \operatorname{trd}(\gamma)=4,\
 \operatorname{nrd}(\gamma)=1\right\}.}
\end{equation}

The inverse sends

\[
 \gamma=2+sE+tF_0+uG
\]

to

\[
 (a,q,n)=(s-t,s+t,u).
\]

There is no hidden congruence loss. On the integral equation \(F=-2\), reduction modulo \(2\) gives

\[
 a^2-q^2\equiv0\pmod2,
\]

so \(a\equiv q\pmod2\) is automatic. The parity condition in \eqref{koide-quat-eq-shell-bijection} records the order lattice exactly but removes no integral shell point.

Finally,

\begin{equation}\label{koide-quat-eq-gamma-quadratic}
 (\gamma-2)^2=3,\qquad
 \gamma^2-4\gamma+1=0,\qquad
 \gamma^{-1}=4-\gamma=2-v\in\mathcal O.
\end{equation}

Thus every element on the right of \eqref{koide-quat-eq-shell-bijection} lies in \(\Gamma=\mathcal O^1\).

\subsection{\texorpdfstring{Division over \(\mathbb Q\) and splitting over \(\mathbb R\)}{Division over \textbackslash mathbb Q and splitting over \textbackslash mathbb R}}\label{koide-quat-division-over-mathbb-q-and-splitting-over-mathbb-r}

Square factors may be removed from a Hilbert-symbol presentation, so

\[
 (18,-6)_{\mathbb Q}\simeq(2,-6)_{\mathbb Q}.
\]

At \(p=3\), write

\[
 18=3^2\cdot2,\qquad -6=3\cdot(-2).
\]

The odd-prime Hilbert-symbol formula gives

\begin{equation}\label{koide-quat-eq-hilbert-three}
 (18,-6)_3
 =(-1)^{2\cdot1\cdot(3-1)/2}
   \left(\frac2{3}\right)
   \left(\frac{-2}{3}\right)^2
 =-1.
\end{equation}

Thus \(B\otimes\mathbb Q_3\) is division, and \(B\) is a division algebra over \(\mathbb Q\).

At \(p=2\), write

\[
 18=2\cdot9,\qquad -6=2\cdot(-3).
\]

The \(2\)-adic formula

\[
 (2^\alpha u,2^\beta w)_2
 =(-1)^{
 ((u-1)/2)((w-1)/2)
 +\alpha(w^2-1)/8+\beta(u^2-1)/8}
\]

gives

\begin{equation}\label{koide-quat-eq-hilbert-two}
 (18,-6)_2=-1.
\end{equation}

At every odd prime other than \(3\), both valuations are zero and the Hilbert symbol is \(+1\). At the real place, \(18>0\), so the real Hilbert symbol is \(+1\). Consequently

\begin{equation}\label{koide-quat-eq-ramification-discriminant}
 \operatorname{Ram}(B)=\{2,3\},\qquad
 \operatorname{disc}(B)=6,
\end{equation}

and

\[
 B\otimes_{\mathbb Q}\mathbb R\simeq M_2(\mathbb R).
\]

An explicit real splitting is

\begin{equation}\label{koide-quat-eq-real-splitting}
 \rho(I)=
 \begin{pmatrix}3\sqrt2&0\\0&-3\sqrt2\end{pmatrix},
 \qquad
 \rho(J)=
 \begin{pmatrix}0&6\\-1&0\end{pmatrix}.
\end{equation}

These matrices square to \(18\) and \(-6\) and anticommute. Therefore

\[
 \rho(K)=
 \begin{pmatrix}0&18\sqrt2\\3\sqrt2&0\end{pmatrix}.
\]

For the element associated with \((a,q,n)\),

\begin{equation}\label{koide-quat-eq-gamma-matrix}
 \rho(\gamma)=
 \begin{pmatrix}
 2+\frac{3\sqrt2}{2}q&3a+9\sqrt2\,n\\
 -\frac12a+\frac{3\sqrt2}{2}n&
 2-\frac{3\sqrt2}{2}q
 \end{pmatrix}.
\end{equation}

Its determinant is the reduced norm in \eqref{koide-quat-eq-gamma-trace-norm}.

Let \(\mathcal M\) be a maximal overorder containing \(\mathcal O\).
Since \([\mathcal M:\mathcal O]=18\), the finite abelian group
\(\mathcal M/\mathcal O\) has order \(18\).  Lagrange's theorem therefore
shows that multiplication by \(18\) annihilates every element of this
additive quotient; equivalently, the ideal \((18)=18\mathbb Z\) annihilates
\(\mathcal M/\mathcal O\), and \(18\mathcal M\subset\mathcal O\).  Hence the
principal congruence subgroup
\[
 \mathcal M^1(18\mathcal M)
 :=\{u\in\mathcal M^1:u-1\in18\mathcal M\}
\]
satisfies the exact inclusions
\[
 \mathcal M^1(18\mathcal M)\subset\mathcal O^1\subset\mathcal M^1.
\]
Thus \(\mathcal O^1\) has finite index in the norm-one group of a maximal
order and contains a principal congruence subgroup.  This is the required
finite-index and congruence bridge; it does not replace the local type of
\(\mathcal O\) by the local type of \(\mathcal M\).

The image

\[
 \bar\Gamma=\rho(\mathcal O^1)/\{\pm1\}
 \subset PSL_2(\mathbb R)
\]

is a discrete congruence arithmetic Fuchsian group. Since the defining
rational quaternion algebra is division,
\(\bar\Gamma\backslash\mathbb H\) is compact; see
\cite[Lemma~38.1.2, Proposition~38.1.3, and Definition~38.7.2]{Voight_2021}.

\subsection{Centralizer, optimal quadratic order, and primitivity}\label{koide-quat-centralizer-optimal-quadratic-order-and-primitivity}

Fix a shell point and set

\[
 v=\gamma-2.
\]

Equation~\eqref{koide-quat-eq-gamma-quadratic} gives \(v^2=3\), hence

\begin{equation}\label{koide-quat-eq-quadratic-field}
 \mathbb Q(\gamma)=\mathbb Q(v)\cong\mathbb Q(\sqrt3).
\end{equation}

Since \(\gamma\) is noncentral, the centralizer theorem for quaternion algebras gives

\begin{equation}\label{koide-quat-eq-algebra-centralizer}
 C_B(\gamma)=\mathbb Q(\gamma).
\end{equation}

Moreover,

\[
 \mathbb Z[v]\subseteq\mathcal O\cap\mathbb Q(v).
\]

Every element of \(\mathcal O\) is integral over \(\mathbb Z\), so every element of the intersection is an algebraic integer in \(\mathbb Q(\sqrt3)\). Because

\[
 \mathcal O_{\mathbb Q(\sqrt3)}=\mathbb Z[\sqrt3]
\]

and \(\mathbb Z[v]\) is already isomorphic to this maximal quadratic order, both inclusions force

\begin{equation}\label{koide-quat-eq-optimal-intersection}
 \boxed{
 \mathcal O\cap\mathbb Q(\gamma)
 =\mathbb Z[v]=\mathbb Z[\gamma-2]
 \cong\mathbb Z[\sqrt3].}
\end{equation}

Thus each shell point gives an optimal embedding

\begin{equation}\label{koide-quat-eq-optimal-embedding}
 \iota_\gamma:\mathbb Z[\sqrt3]\hookrightarrow\mathcal O,
 \qquad \sqrt3\longmapsto\gamma-2.
\end{equation}

Conversely, every embedding \eqref{koide-quat-eq-optimal-embedding} gives a trace-\(4\), norm-\(1\) element \(2+\iota(\sqrt3)\), hence one integral ternary point. The shell is exactly the set of these optimal embeddings before quotienting by conjugation.

The norm-one units of \(\mathbb Z[\sqrt3]\) are

\begin{equation}\label{koide-quat-eq-norm-one-units}
 \{\pm(2+\sqrt3)^k:k\in\mathbb Z\}.
\end{equation}

For completeness, \(2+\sqrt3\) is the smallest positive norm-one unit greater than \(1\). If \(x+y\sqrt3>1\) has norm \(1\) and \(y>0\), then \(y\ge1\), \(x^2=1+3y^2\), and \(x+y\sqrt3\ge2+\sqrt3\). Multiplication by a suitable power of \(2-\sqrt3\) reduces any positive norm-one unit to \([1,2+\sqrt3)\), proving \eqref{koide-quat-eq-norm-one-units}.

Equations~\eqref{koide-quat-eq-algebra-centralizer}--\eqref{koide-quat-eq-norm-one-units} imply

\begin{equation}\label{koide-quat-eq-group-centralizer}
 C_\Gamma(\gamma)=\{\pm\gamma^k:k\in\mathbb Z\}.
\end{equation}

Hence \(\gamma\) is primitive in \(\Gamma\) and its image is primitive in \(\bar\Gamma\). Indeed, a root of \(\gamma\), or a root of its class modulo the center, would commute with \(\gamma\) and therefore be a unit from \eqref{koide-quat-eq-norm-one-units}; the exponent equation \(rk=1\) has no solution with \(k\ge2\).

The reduced characteristic polynomial is

\[
 T^2-4T+1,
\]

so the real eigenvalues are

\begin{equation}\label{koide-quat-eq-eigenvalues}
 2+\sqrt3,\qquad 2-\sqrt3.
\end{equation}

For a hyperbolic element of \(PSL_2(\mathbb R)\),

\[
 2\cosh(\ell/2)=|\operatorname{tr}(\rho(\gamma))|=4.
\]

Thus every shell element has primitive translation length

\begin{equation}\label{koide-quat-eq-translation-length}
 \boxed{
 \ell_0=2\operatorname{arcosh}(2)
       =2\log(2+\sqrt3).}
\end{equation}

\subsection{Integral adjoint action and finitely many classes}\label{koide-quat-integral-adjoint-action-and-finitely-many-classes}

Let

\[
 \mathscr S
 =\{\gamma\in\Gamma:\operatorname{trd}(\gamma)=4\}.
\]

By \eqref{koide-quat-eq-shell-bijection}, \(\mathscr S\) is exactly the entire integral ternary shell. By \eqref{koide-quat-eq-translation-length}, every element of \(\mathscr S\) is hyperbolic with translation length \(\ell_0\).

Put

\[
 L=\{(a,q,n)\in\mathbb Z^3:a\equiv q\pmod2\}.
\]

Conjugation preserves \(\mathcal O\), \(\mathcal O^0\), and reduced norm. Because

\[
 \operatorname{nrd}(v(a,q,n))=\frac32F(a,q,n),
\]

it gives an exact orthogonal representation on the actual parity lattice

\begin{equation}\label{koide-quat-eq-adjoint-representation}
 \operatorname{Ad}:\Gamma/\{\pm1\}
 \hookrightarrow SO^+(L,F).
\end{equation}

The kernel before dividing by the center is \(\Gamma\cap\mathbb Q^\times=\{\pm1\}\). No action on the larger lattice \(\mathbb Z^3\), and no finite-index equality with its full integral orthogonal group, is asserted.

There are only finitely many \(\Gamma\)-conjugacy classes in \(\mathscr S\). Choose a compact fundamental set \(D\subset\mathbb H\) for \(\bar\Gamma\), and fix \(z_0\in\mathbb H\). For each conjugacy class, conjugate a representative so that its invariant axis meets \(D\). If \(z\in D\) lies on that axis, then

\[
 d(z,\gamma z)=\ell_0.
\]

Writing

\[
 R_D=\max_{w\in D}d(z_0,w),
\]

the triangle inequality gives

\[
 d(z_0,\gamma z_0)\le2R_D+\ell_0.
\]

Proper discontinuity leaves only finitely many elements of \(\bar\Gamma\) moving \(z_0\) by at most this amount. Hence only finitely many \(\bar\Gamma\)-classes occur. This also proves finiteness of the corresponding \(\Gamma\)-classes: if two images are conjugate in \(\bar\Gamma\), lifting a conjugator gives

\[
 \delta\gamma_1\delta^{-1}=\pm\gamma_2.
\]

Both sides before the possible sign have reduced trace \(4\); the minus sign would have reduced trace \(-4\), so it is impossible.

Choose representatives \(\gamma_1,\ldots,\gamma_h\). Then

\begin{equation}\label{koide-quat-eq-shell-class-decomposition}
 \boxed{
 \mathscr S=
 \coprod_{j=1}^h
 \{\delta\gamma_j\delta^{-1}:
 [\delta]\in\Gamma/C_\Gamma(\gamma_j)\},
 \qquad
 C_\Gamma(\gamma_j)=\{\pm\gamma_j^k:k\in\mathbb Z\}.}
\end{equation}

This proves finiteness but does not compute \(h\) or a representative list. Those finite arithmetic data are needed by a class-by-class spectral formula whose coefficients contain geodesic periods.

\subsection{Exact reformulation of the Koide chamber count}\label{koide-quat-exact-reformulation-of-the-koide-chamber-count}

For \(\gamma\in\mathscr S\), let

\[
 c(\gamma)=(a(\gamma),q(\gamma),n(\gamma))
\]

be the inverse coordinate map in
\cref{koide-quat-trace-zero-lattice-and-exact-shell-bijection}. Let

\[
 \mathcal C=\{(a,q,n)\in\mathbb R^3:
 a>0,\ q\ge0,\ n\ge0,\ a+4n>q\},
\]

and retain the canonical height

\[
 H(a,q,n)=4a+17n.
\]

The integral bijections in
\cref{koide-count-thm-coordinate-bijection,koide-param-thm:integral-parametrization}
give

\begin{equation}\label{koide-quat-eq-koide-class-count}
 \boxed{
 R(X)=\sum_{j=1}^h
 \#\left\{
 [\delta]\in\Gamma/C_\Gamma(\gamma_j):
 \begin{array}{l}
 c(\delta\gamma_j\delta^{-1})\in\mathcal C,\\
 H(c(\delta\gamma_j\delta^{-1}))\le X
 \end{array}
 \right\}.}
\end{equation}

There is no hidden finite-to-one factor in
\eqref{koide-quat-eq-koide-class-count}: cosets by the centralizer
parametrize the elements of each class exactly once, and
\eqref{koide-quat-eq-shell-bijection} parametrizes shell elements by triples
exactly once.  Put \(C_j=C_\Gamma(\gamma_j)\).  Here
\(\Gamma/C_j\) is the set of left cosets \(\delta C_j\), not the set of
right cosets \(C_j\delta\), and the exact orbit bijection is

\[
 \Phi_j:\Gamma/C_j\longrightarrow
 \{\delta\gamma_j\delta^{-1}:\delta\in\Gamma\},
 \qquad
 \delta C_j\longmapsto\delta\gamma_j\delta^{-1}.
\]

Indeed, right multiplication of \(\delta\) by an element of \(C_j\) does
not change its conjugate, so \(\Phi_j\) is well defined and surjective; and

\[
 \Phi_j(\delta_1C_j)=\Phi_j(\delta_2C_j)
 \iff \delta_2^{-1}\delta_1\in C_j
 \iff \delta_1C_j=\delta_2C_j,
\]

so it is injective.  If \(\pi:\Gamma\to\bar\Gamma\) is the quotient map,
then

\[
 \Gamma/C_j
 \cong
 \bar\Gamma/C_{\bar\Gamma}(\bar\gamma_j)
\]

via \(\delta C_j\mapsto\pi(\delta)C_{\bar\Gamma}(\bar\gamma_j)\).
To prove the required equality of stabilizers, observe first that
\(\pi(C_j)\subseteq C_{\bar\Gamma}(\bar\gamma_j)\).  Conversely, if
\(\pi(\delta)\) centralizes \(\bar\gamma_j\), then
\(\delta\gamma_j\delta^{-1}=\pm\gamma_j\).  Reduced-trace invariance and
\(\operatorname{trd}(\gamma_j)=4\) exclude the minus sign, so
\(\delta\in C_j\).  Thus
\(\pi(C_j)=C_{\bar\Gamma}(\bar\gamma_j)\); together with
\(\ker\pi=\{\pm1\}\subset C_j\), this proves the displayed bijection and
shows that passage to \(PSL_2\) introduces no factor \(2\). Positive trace
selects the lift sign, but the chamber is not invariant under
\(\gamma\mapsto\gamma^{-1}\); inverse classes and their orientations must
still be retained separately.

Under the splitting \eqref{koide-quat-eq-real-splitting},

\begin{equation}\label{koide-quat-eq-frobenius-norm}
 \|\rho(\gamma)\|_{\mathrm F}^2
 =8+9q^2+\frac{37}{4}a^2
   +\frac{105\sqrt2}{2}an+\frac{333}{2}n^2.
\end{equation}

For \(i\in\mathbb H\),

\[
 2\cosh d(i,\rho(\gamma)i)=\|\rho(\gamma)\|_{\mathrm F}^2.
\]

Since \(\sinh(\ell_0/2)=\sqrt3\), the standard Huber parameter is

\begin{equation}\label{koide-quat-eq-huber-parameter}
 X_i(\gamma)^2
 =\frac{\|\rho(\gamma)\|_{\mathrm F}^2-2}{12}.
\end{equation}

The shell equation simplifies \eqref{koide-quat-eq-frobenius-norm} exactly:

\[
 \|\rho(\gamma)\|_{\mathrm F}^2
 =14+\frac{(7a+15\sqrt2\,n)^2}{4},
\]

and hence

\begin{equation}\label{koide-quat-eq-radial-height}
 \boxed{
 X_i(\gamma)^2
 =1+\frac{(7a+15\sqrt2\,n)^2}{48}.}
\end{equation}

Thus, in the positive chamber, the radial cutoff has the exact irrational linear height \(7a+15\sqrt2\,n\). The Koide height \(H=4a+17n\) is a different linear height. The two sizes are comparable:

\[
 a\le\frac H4,\qquad
 n\le\frac H{17},\qquad
 0\le q<a+4n\le\frac{33}{68}H,
\]

and conversely

\[
 H\le\sqrt{305}\sqrt{a^2+q^2+n^2}.
\]

Norm equivalence in the trace-zero space then compares \(H\) coarsely with \(X_i\). Such comparability identifies the growth scale but does not transfer a main coefficient or a remainder between the sharply bounded regions.

\subsection{Consequence for the open square-root discrepancy}\label{koide-quat-consequence-for-the-open-square-root-discrepancy}

The quaternion construction proves:

\begin{enumerate}
\def\labelenumi{\arabic{enumi}.}
\item
  the ternary shell is exactly a finite union of primitive hyperbolic conjugacy classes in one explicit cocompact arithmetic Fuchsian group;
\item
  \(\log(2+\sqrt3)=\ell_0/2\), so its occurrence in the main constant is tied to the primitive translation length;
\item
  the conjectural exponent \(1/2\) is the same exponent proposed for the standard hyperbolic lattice-point problem in conjugacy classes, giving a precise analytic analogy but not an equivalence of the two cutoff problems.
\end{enumerate}

The third item classifies the remaining analytic problem; it does not prove the
error estimate. Chatzakos--Petridis~\cite[Theorem~1.1]{ChatzakosPetridis_2016}
give

\[
 E(\mathcal H,X;z)=O(X^{2/3})
\]

for the standard radial count after the main expression, including every small-eigenvalue term, is subtracted. Their Conjecture 5.7 proposes

\[
 E(\mathcal H,X;z)=O_\varepsilon(X^{1/2+\varepsilon}).
\]

The cutoff in \eqref{koide-quat-eq-koide-class-count} uses both a sharp chamber and \(4a+17n\), not the radial height \(7a+15\sqrt2\,n\) in \eqref{koide-quat-eq-radial-height}. Their symmetric difference can have main-order size, so even the conjectured radial bound would not by itself prove the Koide bound. Therefore the status of

\[
 R(X)-C_0X=O_\varepsilon(X^{1/2+\varepsilon})
\]

remains \textbf{OPEN}. Closing it requires a spectral or automorphic argument for the exact finite sum \eqref{koide-quat-eq-koide-class-count}, with its exact sector and height.

\subsection{Missing mathematics before an error theorem can be claimed}\label{koide-quat-missing-mathematics-before-an-error-theorem-can-be-claimed}

The bridge does not supply the following results.

\begin{enumerate}
\def\labelenumi{\arabic{enumi}.}
\item
  \textbf{Class data.} Determine \(h\), representatives \(\gamma_j\), and orientation/inversion conventions. Finiteness does not compute the periods in a spectral expansion.
\item
  \textbf{Exact nonmaximal order.} The order has reduced discriminant \(108\), is nonmaximal at the ramified primes, and is not an Eichler order in the standard sense. A result for maximal or specified Eichler orders must be checked at this order or transferred through finite-index groups with all coset multiplicities retained. Index and discriminant alone do not determine the precise local order type required by a trace formula.
\item
  \textbf{Central quotient.} Track the passage from \(\Gamma\subset SL_2(\mathbb R)\) to \(\bar\Gamma\subset PSL_2(\mathbb R)\), including \(\pm1\), inverse classes, and possible orbifold torsion.
\item
  \textbf{Exact cutoff.} Prove a trace formula or effective orbital theorem for all four inequalities defining \(\mathcal C\) and for \(4a+17n\le X\). The radial Huber cutoff is different.
\item
  \textbf{Sharp-boundary control.} Remove smoothing with an explicit estimate at the required scale. The two Pell faces contribute only \(O(\log X)\), as already proved canonically, but this alone does not control every smoothing strip adjacent to all chamber walls.
\item
  \textbf{Exceptional spectrum.} Once an appropriate sector spectral expansion is constructed, one must determine whether it contains \(X^{s_j}\) terms for \(\lambda_j=s_j(1-s_j)<1/4\), weighted by closed-geodesic periods. To subtract only \(C_0X\) with a \(O(X^{1/2+\varepsilon})\) remainder, every relevant exceptional term must be shown to vanish or cancel, or it must be displayed and subtracted. No such vanishing is proved here.
\item
  \textbf{Pointwise square-root estimate.} After secondary terms and the exact sector are handled, the needed pointwise \(1/2+\varepsilon\) bound remains at the conjectural strength of the classical problem. Mean-square results do not give it for every \(X\).
\item
  \textbf{Main coefficient.} A spectral derivation must reproduce \[
  C_0=\frac{9}{2\pi^2}\log(2+\sqrt3)
      \log\!\left(\frac3{2\sqrt2}\right)
  \] using the same lattice, chamber, strict walls, height, and Haar measure. The existing local-density proof fixes this constant but does not calibrate a new trace formula automatically.
\item
  \textbf{Transfer between the two exact heights.} Formula~\eqref{koide-quat-eq-radial-height} already gives the exact standard radial parameter. What is missing is a remainder-level transfer between its height \(7a+15\sqrt2\,n\) and the canonical height \(4a+17n\), together with the chamber walls. Coarse norm comparability does not preserve a discrepancy estimate.
\end{enumerate}

These are concrete mathematical tasks. None invalidates the quaternion bridge; together they explain why it does not close the open claim.

\subsection{Reproducibility and references}\label{koide-quat-reproducibility-and-references}

The canonical quaternion replay is
\path{supporting_materials/workbench/evidence/koide_exact_checkpoint_20260904/quaternion/verify_quaternion_bridge.py}
(SHA-256
\path{6675ba2fc4fa8c2b10e14558e412178b6b5b579c9c96afca78e1db9bcdf5c6b4}).
Its recorded successful output is
\path{supporting_materials/workbench/evidence/koide_exact_checkpoint_20260904/quaternion/verification.log}
(SHA-256
\path{66853c6191a000c7c827aacd1e0e33dcf34dbe3dd61d746f66fa4d805f34c91e}),
and its local manifest is
\path{supporting_materials/workbench/evidence/koide_exact_checkpoint_20260904/quaternion/SHA256SUMS.txt}
(SHA-256
\path{f9d044db0d4b003e7aef7649b3394e6b299c09a6b2eb7adca095f8f9100ba623}).
The checkpoint-wide manifest is
\path{supporting_materials/workbench/evidence/koide_exact_checkpoint_20260904/SHA256SUMS}
(SHA-256
\path{8d108eb928e5d07093badde169f443a4a38b6fcf3eb59d5192bdd33ecef8d7a4}).
The replay checks the basis multiplication table; the trace, norm, and
coordinate formulas for
\((r,s,t,u)\in\{-4,-3,\ldots,4\}^4\); the shell identities for
\((a,q,n)\in\{-24,-23,\ldots,24\}^2\times\{-8,-7,\ldots,8\}\), containing
\(48\) shell points; the
inverse identity on those points; the radial-height identity on that box;
and the Hilbert-symbol formulas at \(2\) and \(3\).  These finite ranges are
consistency checks, not proofs outside the stated boxes.  The exact canonical
sources used for the counting statements imported into this bridge are:

\begin{itemize}
\item
  \path{supporting_materials/workbench/research/koide_counting_repair.tex},
  SHA-256
  \path{965b0a3dca8bfe4f22b9e9208368fa04697070fd25bd06bd67f2481d17cb9d9d};
\item
  \path{supporting_materials/workbench/research/koide_exact_parametrization.tex},
  SHA-256
  \path{a457709dcce0fce0774eb9326b84b00d5d8144c2087f480da704974e499bcdd3};
\item
  \path{supporting_materials/workbench/research/koide_exact_census.tex},
  SHA-256
  \path{18ed5968b7e67d9a0b23d8675b40c37ba46170ad29efdbb9ce10ed632c99600a}.
\end{itemize}

The bounded coordinate-bijection replay, whose scope does not include the
quaternion multiplication table, Hilbert symbols, centralizer, conjugacy
finiteness, or spectral conclusions, is
\path{supporting_materials/workbench/reproduction/t1_koide_exact/verify_koide_narain_root_density.py}
(SHA-256
\path{e4bf2982b26b07add6057b5d2d5c771a52ef52c629b5ad5f0aab24e440e41986}).
Its stored height-\(200000\) output is
\path{supporting_materials/workbench/reproduction/t1_koide_exact/koide_verification_200000.json}
(SHA-256
\path{b907f0ee5bdd3ad80622b669722d218b79a11a5baab0d4571e141936100706b7});
the exact invocation and scope are stated in
\path{supporting_materials/workbench/reproduction/t1_koide_exact/README.md}
(SHA-256
\path{5e9cdcb9e67954c37d8fde69e13549172aa03245674b358e4a1d25f8d9d3bff0}).
The quaternion identities and deductions in this section are proved in the
displayed algebraic calculations themselves, not inferred from that bounded
replay.

References:

\begin{itemize}
\item
  John Voight, \emph{Quaternion Algebras}, Graduate Texts in Mathematics 288,
  especially Section 10.1, Lemma 38.1.2, and Proposition 38.1.3;
  \href{https://doi.org/10.1007/978-3-030-56694-4}{DOI}.
\item
  Dimitrios Chatzakos and Yiannis N. Petridis,
  \emph{The hyperbolic lattice point problem in conjugacy classes},
  arXiv:1504.01307v2, Theorem 1.1 and Conjecture 5.7;
  \href{https://arxiv.org/abs/1504.01307}{arXiv}.
\item
  The exact Koide chamber, height, main term, and Pell-face estimates are in
  Sections~\ref{koide-count-sec-main}--\ref{koide-count-sec-census-1e8}
  of this workbench.
\end{itemize}

% Integrated 4 September 2026; exact in-project provenance is recorded below.
% The exact fixed level, lattice, chamber, height, theorem hypotheses, and locators are retained.
\section{Primary-source theorem audit of effective counting for the fixed Koide ternary quadric sector}\label{koide-lit-primary-source-theorem-audit-of-effective-counting-for-the-fixed-koide-ternary-quadric-sector}

\textbf{Date:} 4 September 2026\\
\textbf{Exact counting function:}

\[
 R(X)=\#\left\{(a,q,n)\in\mathbb Z^3:
 \begin{array}{l}
 a^2-3q^2-18n^2=-2,\\
 a>0,\ q\geq0,\ n\geq0,\ a+4n>q,\\
 4a+17n\leq X
 \end{array}\right\}.
\]

The leading constant proved in
\cref{koide-count-thm-linear-asymptotic} is

\[
 C_0=\frac{9}{2\pi^2}\log(2+\sqrt3)
       \log\!\left(\frac{3}{2\sqrt2}\right).
\]

The canonical project-relative mathematical sources used by this audit are:

\begin{itemize}
\item
  \path{supporting_materials/workbench/research/koide_counting_repair.tex},
  which proves the coordinate bijection, local densities, invariant-volume
  factor, Gorodnik--Oh--Shah hypotheses, and calibrated constant, SHA-256
  \path{965b0a3dca8bfe4f22b9e9208368fa04697070fd25bd06bd67f2481d17cb9d9d};
\item
  \path{supporting_materials/workbench/research/koide_exact_parametrization.tex},
  which proves the exact parametrization and the Benoist--Oh power-saving
  application, SHA-256
  \path{a457709dcce0fce0774eb9326b84b00d5d8144c2087f480da704974e499bcdd3};
\item
  \path{supporting_materials/workbench/research/koide_quaternion_bridge.tex},
  which proves the exact order, shell, centralizer, conjugacy-class, and
  radial-height statements used below, SHA-256
  \path{93759f7ba99870c84c9144a373a53d2e9a8bc3ae909518e47991373ace0b47bb};
\item
  \path{supporting_materials/workbench/research/koide_exact_census.tex},
  which records the separately bounded finite census, SHA-256
  \path{18ed5968b7e67d9a0b23d8675b40c37ba46170ad29efdbb9ce10ed632c99600a}.
\end{itemize}

The primary publications themselves are identified by exact versions,
dates where version-sensitive, locators, and DOI or arXiv links in
\cref{koide-lit-primary-references-and-exact-locators}.  They are not
claimed to be stored in a separate canonical evidence manifest.

\subsection{Determination}\label{koide-lit-determination}

Among the source versions enumerated in
\cref{koide-lit-primary-references-and-exact-locators}, no theorem proves,
for this exact fixed level, full integral lattice, sharp chamber, and sharp
polyhedral height, the unconditional estimate

\[
 R(X)-C_0X=O_\varepsilon(X^{1/2+\varepsilon}).
\]

The strongest exact-domain conclusion obtained from the general theorems
enumerated there is

\[
 \boxed{R(X)=C_0X+O(X^{1-\eta})\quad\text{for some unspecified }\eta>0.}
\]

It follows by applying Benoist--Oh, Definition 13.1 and Theorem 13.2, orbit by orbit after directly verifying effective well-roundedness for the present family. Benoist--Oh do not give a numerical lower bound for \(\eta\), and their proof does not imply \(\eta\geq 1/2\).

The superficially stronger exponents in nearby problems do not answer the question:

\begin{itemize}

\item
  Good's sharp hyperbolic conjugacy-class theorem, as formulated and reproved by Chatzakos--Petridis, gives \(O(X^{2/3})\) only for a different radial cutoff and only after every small-eigenvalue term \(X^{s_j}\) has been placed in the main expression. The same paper states \(O_\varepsilon(X^{1/2+\varepsilon})\) as Conjecture 5.7.
\item
  Liu--Sarnak obtain a smooth radial error \(O_\varepsilon(X^{1/2+\theta+\varepsilon})\), hence \(O_\varepsilon(X^{39/64+\varepsilon})\) from \(\theta=7/64\), and square-root strength only under the generalized Selberg/Ramanujan-at-infinity conjecture for the relevant Jacquet--Langlands transfers. Their paper's global affine-sieve setup assumes \(t\det(f)\) square-free, but the analytic Proposition 3.2 itself is formulated for a variable cocompact lattice with a spectral gap. After the exact quaternionic identification it applies class by class to this lattice. What it still does not supply is a sharp cutoff, the coordinate chamber, or the nonradial height \(4a+17n\).
\item
  Gorodnik--Nevo's published sharp polar-sector theorem loses powers through set thickening. For \(\dim G=3\), it gives \(O(X^{19/20+\epsilon})\) even with tempered spectrum; the Kim--Sarnak complementary-series threshold, which does apply to the exact congruence arithmetic lattice by Jacquet--Langlands, gives the still weaker \(O(X^{39/40+\epsilon})\). This remains a comparison theorem because Gorodnik--Nevo's fixed polar-sector family has not been identified with the exact polyhedron, including all integral orbits.
\item
  Square-root-scale statements in Chatzakos--Petridis, Chatzakos, Voskou, and Lekkas--Petridis are mean-square, averaged, Omega, or conjectural statements. They are not pointwise sharp-sector upper bounds.
\item
  Recent split-stabilizer theorems assume \(-m\det(F)\) is a square and have a \(T\log T\)-type main term. Here \(-(-2)\det(F)=108\) is not a square, so they concern the opposite arithmetic regime.
\end{itemize}

The two independent missing inputs are therefore (i) a sharp-boundary spectral or arithmetic estimate for the actual height and chamber, and (ii) control or elimination of any complementary-spectrum secondary terms for the actual sector. The Pell faces \(q=0\) and \(n=0\), which contribute only \(O(\log X)\), do not solve the moving outer-boundary problem.

\subsection{Exact arithmetic and geometric classification}\label{koide-lit-exact-arithmetic-and-geometric-classification}

\subsubsection{The rational group and stabilizer}\label{koide-lit-the-rational-group-and-stabilizer}

Let

\[
 F(a,q,n)=a^2-3q^2-18n^2,
 \qquad \mathscr X=\{F=-2\}.
\]

The matrix of \(F\) is \(\operatorname{diag}(1,-3,-18)\), so \(\det(F)=54\) and

\[
 -(-2)\det(F)=108\notin\mathbb Q^{\times2}.
\]

The homogeneous form \(F\) is anisotropic over \(\mathbb Q\). Indeed, suppose a primitive integral triple satisfies

\[
 a^2-3q^2-18n^2=0.
\]

Modulo \(3\), one gets \(3\mid a\); write \(a=3a_1\). Dividing by \(3\) and reducing modulo \(3\) gives \(3\mid q\); write \(q=3q_1\). The resulting equation is

\[
 a_1^2=3q_1^2+2n^2.
\]

Modulo \(3\), this says \(a_1^2+n^2\equiv0\pmod3\), whence \(3\mid a_1,n\). Thus \(3\mid a,q,n\), contradicting primitivity.

Put \(G=\operatorname{Spin}(F)\) and \(x_0=(1,1,0)\). The orthogonal complement of \(x_0\) is

\[
 x_0^\perp=\{(3s,s,t):s,t\in\mathbb Q\},
 \qquad F(3s,s,t)=6(s^2-3t^2).
\]

Consequently the connected stabilizer \(H=\operatorname{Stab}_G(x_0)^\circ\) is the norm-one torus

\[
 \operatorname{Res}^{1}_{\mathbb Q(\sqrt3)/\mathbb Q}\mathbb G_m.
\]

It is split over \(\mathbb R\), anisotropic over \(\mathbb Q\), and has no
nontrivial \(\mathbb Q\)-character.  Under the quaternionic identification
in \cref{koide-quat-centralizer-optimal-quadratic-order-and-primitivity},
\(x_0\) corresponds to \(v_0=(I+J)/2\), with \(v_0^2=3\), and

\[
 \mathcal O\cap\mathbb Q(v_0)=\mathbb Z[v_0],
 \qquad
 H(\mathbb Q)\cap\Gamma
 =\{\pm(2+v_0)^k:k\in\mathbb Z\}.
\]

Thus \((H(\mathbb Q)\cap\Gamma)/\{\pm1\}\) is infinite cyclic.  Under
\(H(\mathbb R)^\circ\cong\mathbb R_{>0}\), the logarithm sends its positive
part to \(\log(2+\sqrt3)\mathbb Z\subset\mathbb R\); the quotient is compact,
and in particular \((H\cap\Gamma)\backslash H(\mathbb R)\) has finite
volume. Since \(F\) is \(\mathbb Q\)-anisotropic, the arithmetic Fuchsian
quotient attached to \(G\) is cocompact as well. There is no ambient
hyperbolic cusp.

This is the finite-volume stabilizer regime associated with a closed geodesic.
It is not the rationally split regime with divergent stabilizer in which one
obtains \(T\log T\) growth.

\subsubsection{The exact height really is a norm-sector height}\label{koide-lit-the-exact-height-really-is-a-norm-sector-height}

Define

\[
 \|(a,q,n)\|_*=max\{|4a+17n|,|q|,|a|\}.
\]

This is a genuine norm because it is the pullback of the sup norm under the invertible linear map

\[
 (a,q,n)\longmapsto(4a+17n,q,a),
\]

whose determinant is \(-17\). On the target chamber,

\[
 0<a<4a+17n,
 \qquad 0\leq q<a+4n<4a+17n,
\]

and therefore \(\|(a,q,n)\|_*=4a+17n\). Thus the outer cutoff is not merely comparable to a norm: it is exactly the restriction of a fixed ambient norm.

The additional wall \(a+4n>q\) removes only a bounded real set. If \(q\geq a+4n\), then

\[
 a^2+2=3q^2+18n^2
 \geq3(a+4n)^2+18n^2,
\]

so \(2a^2+24an+66n^2\leq2\). Its only integral point in the positive orthant on \(\mathscr X\) is \((1,1,0)\). The unbounded faces \(q=0\) and \(n=0\) are Pell families and contribute \(O(\log X)\) integral points, but their real boundary strips still have to be included in a sharp-domain argument.

\subsubsection{Invariant volume with a strong enough secondary estimate}\label{koide-lit-invariant-volume-with-a-strong-enough-secondary-estimate}

Use

\[
 \sqrt3q=r\cos\theta,
 \qquad 3\sqrt2n=r\sin\theta,
 \qquad 0\leq\theta\leq\frac\pi2.
\]

On \(F=-2\), \(r=\sqrt{a^2+2}\), and the gauge measure is

\[
 dm_\infty=\frac1{6\sqrt6}\,da\,d\theta.
\]

Writing \(c=17/(3\sqrt2)\), the outer boundary is

\[
 4a+c\sqrt{a^2+2}\sin\theta=X.
\]

If \(A_X(\theta)\) is its positive solution, then uniformly in \(\theta\),

\[
 A_X(\theta)=\frac{X}{4+c\sin\theta}+O(X^{-1}).
\]

Indeed \(A_X(\theta)\asymp X\) uniformly and \(\sqrt{A^2+2}=A+O(A^{-1})\). Once \(X\) contains the bounded wall excision, it follows that

\[
 m_\infty(B_X)=I_\infty X+O(1),
 \qquad
 I_\infty=\frac1{2\sqrt3}\log\!\left(\frac3{2\sqrt2}\right).
\]

This \(O(1)\) refinement matters: a theorem whose main term is the exact invariant volume can be converted to the displayed \(C_0X\) main term without spending a power of \(X\).

\subsubsection{The full integral locus is not assumed to be one orbit}\label{koide-lit-the-full-integral-locus-is-not-assumed-to-be-one-orbit}

The exact lattice must not be silently replaced by a convenient single orbit or by a commensurable maximal-order lattice. Borel--Harish-Chandra finiteness, in the form quoted by Benoist--Oh on journal p.~1892, says that \(G(\mathbb Z)\) has finitely many orbits on \(\mathscr X(\mathbb Z)\). Thus

\[
 \mathscr X(\mathbb Z)=\bigsqcup_{j=1}^{h}x_j\Gamma
\]

for the integral spin group \(\Gamma\) attached to the chosen lattice and finitely many integral representatives \(x_j\). An effective theorem for one orbit must be applied to every \(x_j\Gamma\), and the finitely many main-measure scalars must be summed. The full-lattice scalar is calibrated from Borovoi's local/Brauer--Manin density formula in \cref{koide-count-prop-finite-sign-balance,koide-count-thm-linear-asymptotic}; it is not obtained by declaring that all integral points form the orbit of \(x_0\).

\subsubsection{Exact quaternionic lattice and its congruence status}\label{koide-lit-exact-quaternionic-lattice-and-its-congruence-status}

For the spectral comparisons below, the precise arithmetic lattice can be retained. Put

\[
 B=(18,-6)_{\mathbb Q}
   =\mathbb Q\langle I,J\rangle/(I^2-18,\ J^2+6,\ IJ+JI),
 \qquad K=IJ,
\]

and

\[
 \mathcal O=\mathbb Z1\oplus
  \mathbb Z\frac{I+J}{2}\oplus
  \mathbb Z\frac{I-J}{2}\oplus
  \mathbb Z\frac K2.
\]

The complete integral multiplication table is
\eqref{koide-quat-eq-multiplication-table}, and its reduced-trace Gram
determinant is \(-108^2\) by
\eqref{koide-quat-eq-order-discriminant-determinant}; hence
\(\operatorname{discrd}(\mathcal O)=108\).  The Hilbert-symbol calculations
\eqref{koide-quat-eq-hilbert-three} and
\eqref{koide-quat-eq-hilbert-two}, together with the unramified-place and
real-place calculation summarized in
\eqref{koide-quat-eq-ramification-discriminant}, show that the ramification
set is exactly \(\{2,3\}\).  Thus \(B\) is the indefinite quaternion
division algebra of discriminant \(6\). The map

\[
 (a,q,n)\longmapsto
 \gamma(a,q,n)=2+\frac{qI+aJ+nK}{2}
\]

is the bijection \eqref{koide-quat-eq-shell-bijection} from
\(F(a,q,n)=-2\) to the reduced-trace-\(4\), reduced-norm-one locus in
\(\mathcal O\). Indeed, the parity condition \(a\equiv q\pmod2\), which is
exactly what membership in \(\mathcal O\) requires, follows automatically
from \(F=-2\) modulo \(2\), and

\[
 \operatorname{nrd}(\gamma)=4+\frac32F(a,q,n)=1.
\]

Thus the acting Fuchsian lattice is exactly \(\Gamma=\mathcal O^1\), not a substituted maximal-order lattice. If \(\mathcal M\supset\mathcal O\) is maximal, the reduced-discriminant index formula gives

\[
 [\mathcal M:\mathcal O]=\frac{108}{6}=18,
 \qquad 18\mathcal M\subset\mathcal O,
\]

and consequently

\[
 \mathcal M^1(18\mathcal M)
 \subset\mathcal O^1\subset\mathcal M^1.
\]

Here
\(\mathcal M^1(18\mathcal M)=\{u\in\mathcal M^1:u-1\in18\mathcal M\}\).
Therefore \(\mathcal O^1\) is a congruence arithmetic Fuchsian group in the
exact sense used by Chu--Li. The number \(108\) is the order's reduced
discriminant, not an Eichler level: the order is nonmaximal at the ramified
primes \(2,3\). This non-Eichler feature changes the finite compact-open
subgroup and hence which automorphic representations have fixed vectors and
the dimensions of those fixed-vector spaces; it does not change the
Jacquet--Langlands/Kim--Sarnak archimedean bound.

\subsection{Result matrix}\label{koide-lit-result-matrix}

\begin{longtable}[]{@{}
  >{\raggedright\arraybackslash}p{0.31\linewidth}
  >{\raggedright\arraybackslash}p{0.31\linewidth}
  >{\raggedright\arraybackslash}p{0.31\linewidth}@{}}
\toprule\noalign{}
\begin{minipage}[b]{\linewidth}\raggedright
Source and exact locator
\end{minipage} & \begin{minipage}[b]{\linewidth}\raggedright
What is actually proved
\end{minipage} & \begin{minipage}[b]{\linewidth}\raggedright
Status for the target
\end{minipage} \\
\midrule\noalign{}
\endhead
\bottomrule\noalign{}
\endlastfoot
Duke--Rudnick--Sarnak (1993), Theorem 1.2, p.~145; Example 1.5, pp.~146--147 & One arithmetic orbit on an affine symmetric variety: count \(\sim\) invariant volume; ternary anisotropic stabilizer gives linear rather than \(T\log T\) growth & Correct qualitative regime, no rate and no exact sharp-sector theorem \\
Eskin--McMullen (1993), Theorem 1.4, author PDF p.~3 & Every well-rounded sequence has count asymptotic to volume & Qualitative only; no exponent \\
Gorodnik--Oh--Shah (2009), definitions before Theorem 1.6 and Theorem 1.6, p.~5; Theorem 1.18, pp.~8--9 & Any norm and an admissible generic cone: full integral count \(\sim\) volume & Direct qualitative theorem for the exact norm and cone, whose hypotheses are proved in \cref{koide-count-lem-genuine-norm,koide-count-lem-gos-hypotheses}; no error term \\
Gorodnik--Oh--Shah (2010), Theorem 1.13, author PDF pp.~4--5 & Well-rounded polar sectors and their asymptotic & Qualitative only \\
Benoist--Oh (2012), Theorem 1.3, p.~1894 & Standard product/Euclidean height ball: relative \(O(T^{-\delta})\) for some \(\delta>0\) & Theorem 1.3 alone is not the exact polyhedron \\
Benoist--Oh (2012), Definition 13.1, p.~1931, and Theorem 13.2, pp.~1931--1932 & Any independently verified effectively well-rounded family: \(\operatorname{vol}(B)(1+O(\operatorname{vol}(B)^{-\delta_0}))\) & Applies to the exact family after the proof in \cref{koide-lit-definition-13.1-for-the-actual-b_x}; yields only unspecified \(O(X^{1-\eta})\) \\
Gorodnik--Nevo (2012), Definition 3.1, p.~13; Theorem 8.1 and Corollary 8.3, pp.~34--35 & Sharp fixed polar sectors: \(O(\operatorname{vol}^{1-\zeta+\epsilon})\), \(\zeta=(2n_e(p))^{-1}(1+3\dim G)^{-1}\) & Explicit but much weaker; target polyhedron and finite orbit union require an additional reduction \\
Liu--Sarnak (2010), equations (2.3)--(2.8), pp.~4--5; Proposition 3.2, pp.~10--11; equation (3.65), p.~16 & Smooth radial count for a cocompact lattice containing the fixed cyclic stabilizer: \(O_\epsilon(T^{1/2+\theta+\epsilon})\) & Exact lattice and finite class sum are available; unconditional \(39/64\), conditional \(1/2\), but fixed-width smoothing and the radial height do not match the sharp target \\
Good (1983), Theorem 4, p.~116, as explicitly restated and reproved in Chatzakos--Petridis (2016), Theorem 1.1, p.~3 & Sharp radial hyperbolic conjugacy-class count: \(O(X^{2/3})\) after all \(X^{s_j}\) small-spectrum terms & Closest sharp spectral model, not the exact polyhedron; square root is CP Conjecture 5.7 \\
Bourgain--Kontorovich--Sarnak (2010), Theorem 1.1, p.~2 & Fixed angular Fourier modes in a sharp group norm ball: \(O_{n,k}(T^{4/3})\) after \(T^{2s_j}\) terms & With affine height \(X\asymp T^2\), again \(X^{2/3}\); constants are not uniform enough to recover a sharp interval for free \\
Browning--Gorodnik (2017), Theorem 2.1, p.~11; Corollaries 2.5--2.6, pp.~24--25 & Full Euclidean balls, finite adelic/orbit sums, and congruence conditions: unspecified \(O(\operatorname{vol}^{1-\varrho})\) & Correct nonsquare homogeneous regime, wrong height geometry and no explicit exponent \\
Cao--Huang (2022), Theorem 1.5, p.~4; Theorem 3.4, pp.~21--22 & Arithmetic purity for anisotropic nonsquare ternary quadrics; effective input for Euclidean balls has an unspecified power & Qualitative full-lattice theorem, not a square-root sector result \\
Cao--Huang--Zhang (arXiv:2509.23320), Setting 2.1 and Theorem 2.1 & Product Euclidean/max height balls with congruence-uniform \(O(\mathbf N(\mathfrak l)^\kappa\operatorname{vol}^{1-\delta})\) & Current extension, but \(\delta\) is unspecified and the growing real factor is the standard height ball, not this cone \\
Oh--Shah (2014), Theorem 1.2, p.~3; Theorem 6.1, pp.~13--14 & Divergent-geodesic/split-stabilizer \(T\log T\) count; smooth error \(O(T)\), sharp relative logarithmic error & Excluded by the finite-volume anisotropic stabilizer \\
Xu--Zhang (2023), main theorem; Huang (arXiv:2412.03350), Theorem 1.1; Zhang (arXiv:2603.21080), Theorem 1.1 & Split-stabilizer results assuming \(-m\det(F)\) is a square, with smooth or radial \(T\log T\)-type expansions & \(108\) is nonsquare; wrong main-term regime \\
Huang (arXiv:2412.05903), Theorems 1.1--1.2 & Smooth compact real weight while the level varies as \(m_0p_0^{2h}\); error of logarithmic relative strength & Varying level and smooth weight, not the fixed-level sharp count \\
Li--Liu (arXiv:2501.15033v2), Theorems 1.1--1.2, pp.~2--3; Lemma 2.3, pp.~5--6 & Almost-prime results using a smooth radial orbit count with \(O_\epsilon(d^{1+\epsilon}T^{1/2+\theta+\epsilon})\) & Current use of the Liu--Sarnak input; globally assumes \(t\det(f)\) square-free and supplies no sharp-sector count \\
Voskou (arXiv:2407.03134v3), Theorem 3 and Conjecture 1, p.~3 & Four sharp orientation sectors in a closed-geodesic double-coset problem: \(O(X^{2/3})\); square root conjectured & Current sharp comparator, not the target \\
Lekkas--Petridis (arXiv:2509.12902, version dated 24 Aug.~2026), Theorem 1.2, Theorem 1.5, Conjecture 1.6 & Sharp closed-geodesic double-coset \(O(X^{2/3})\), square-root mean square, pointwise square-root conjecture & Version-specific confirmation that mean square does not furnish the requested pointwise bound \\
\end{longtable}

\subsection{Benoist--Oh: exact applicability and exact limitation}\label{koide-lit-benoistoh-exact-applicability-and-exact-limitation}

Among the general effective theorems enumerated in
\cref{koide-lit-result-matrix}, this is the one whose stated hypotheses can
be verified for the actual family without changing its height or chamber.

\subsubsection{What the theorem assumes}\label{koide-lit-what-the-theorem-assumes}

Benoist--Oh work with a connected semisimple simply connected almost \(K\)-simple group \(G\), a finite-kernel rational representation, a symmetric subgroup \(H\), and the hypothesis that \(H^\circ\) has no nontrivial \(K\)-character; see journal pp.~1892--1893. For a congruence subgroup \(\Gamma_S\), Theorem 12.5, journal p.~1929, supplies the effective equidistribution of the translates of \((\Gamma_S\cap H_S)\backslash H_S\) used in Theorem 13.2. Example (A), journal p.~1896, explicitly treats ternary quadratic forms provided the homogeneous form does not represent zero over \(\mathbb Q\).

All of these hypotheses hold here:

\begin{itemize}

\item
  \(G=\operatorname{Spin}(F)\) is simply connected and absolutely almost simple of type \(A_1\), and \(G(\mathbb R)\) is noncompact.
\item
  The standard three-dimensional orthogonal representation has finite kernel.
\item
  The stabilizer of a nonisotropic vector is symmetric.
\item
  The mod-\(3\) descent in
  \cref{koide-lit-the-rational-group-and-stabilizer} proves that \(F\) does
  not represent zero over \(\mathbb Q\).
\item
  \(H^\circ\) is the \(\mathbb Q\)-anisotropic norm-one torus of \(\mathbb Q(\sqrt3)\), so it has no nontrivial \(\mathbb Q\)-character and its arithmetic quotient is compact.
\item
  With \(S=\{\infty\}\), the full integral spin group is itself a congruence subgroup in Benoist--Oh's sense.
\end{itemize}

There is no omitted wavefront or properness condition: symmetry of \(H\) is exactly what invokes their effective wavefront argument in Lemma 12.3 and Theorem 12.5. Cocompactness makes escape in the automorphic quotient simpler, but does not remove Definition 13.1(3), which still has to be checked on the growing affine set.

Here ``congruence'' is an exact property of the chosen integral group, not an inference from mere commensurability. Under the quaternionic realization of \(G\), the finite-integrality conditions define a compact-open subgroup of \(G(\mathbb A_f)\); equivalently, after placing the corresponding order \(\mathcal O\) in a maximal order \(\mathcal O_{\max}\), some principal subgroup \(\{u\in\mathcal O_{\max}^1:u\equiv1\pmod{N\mathcal O_{\max}}\}\) lies in the exact integral spin group. This distinction matters for spectral-gap transfer: an arbitrary finite-index noncongruence subgroup need not inherit the Kim--Sarnak gap.

\subsubsection{\texorpdfstring{Definition 13.1 for the actual \(B_X\)}{Definition 13.1 for the actual B\_X}}\label{koide-lit-definition-13.1-for-the-actual-b_x}

Let \(B_X\subset\mathscr X(\mathbb R)\) be the exact chamber cut out by the five linear conditions and \(4a+17n\leq X\). In Benoist--Oh's notation, \(B_{X,\epsilon}^+=B_XB_\epsilon\) and \(B_{X,\epsilon}^-=\bigcap_{u\in B_\epsilon}B_Xu\).

Condition (1) is void because there is no finite-place factor.

For condition (2), if \(u\) lies in an \(\epsilon\)-ball about the identity, every fixed linear functional \(L\) defining the chamber or height satisfies

\[
 |L(xu)-L(x)|\ll\epsilon\|x\|_2.
\]

On the unbounded portion of this chamber, \(\|x\|_2\asymp 1+a\asymp1+(4a+17n)\). Hence a membership change under \(u\) is confined to:

\begin{enumerate}
\def\labelenumi{\arabic{enumi}.}

\item
  an angular strip of width \(O(\epsilon)\) at \(q=0\) or \(n=0\);
\item
  an \(O(\epsilon X)\)-thick radial shell at \(4a+17n=X\); or
\item
  an \(O(\epsilon)\)-neighborhood of one of the compact boundary pieces \(a=0\) or \(a+4n=q\).
\end{enumerate}

Since \(dm_\infty=(6\sqrt6)^{-1}da\,d\theta\), their total measure is \(O(\epsilon X)\). Thus

\[
 m_\infty(B_{X,\epsilon}^+\setminus B_{X,\epsilon}^-)
 \ll \epsilon\,m_\infty(B_X),
\]

so Definition 13.1(2) holds with \(\kappa=1\).

For condition (3), Benoist--Oh's Euclidean height is comparable to \(1+a\) throughout a fixed group-neighborhood enlargement of the chamber. For every \(k>0\), uniformly in \(0<\epsilon<1\),

\[
 \int_{B_{X,\epsilon}^+}H_\infty(x)^{-k}\,dm_\infty(x)
 \ll_k\int_0^{O(X)}(1+a)^{-k}\,da.
\]

The right side is \(O(X^{1-k})\) for \(0<k<1\), \(O(\log X)\) for \(k=1\), and \(O(1)\) for \(k>1\). As \(m_\infty(B_X)\asymp X\), this is \(O(m_\infty(B_X)^{1-\delta(k)})\) for some \(\delta(k)>0\) in every case. Definition 13.1(3) is therefore satisfied.

Proposition 14.2 of Benoist--Oh verifies only their standard product height balls. It should not be cited as though it automatically covers this polyhedron. The preceding direct verification is the necessary bridge to Theorem 13.2.

\subsubsection{Finite-orbit passage and main-term normalization}\label{koide-lit-finite-orbit-passage-and-main-term-normalization}

Benoist--Oh state on journal p.~1892 that \(G(\mathcal O_S)\) has finitely many orbits on \(\mathscr X(\mathcal O_S)\). Apply Theorem 13.2 to each exact integral orbit. If \(\eta_j>0\) is the exponent for the \(j\)-th orbit, take \(\eta=\min_j\eta_j>0\). The finite sum has the form

\[
 R(X)=\kappa_{\mathrm{orb}}\,m_\infty(B_X)+O(X^{1-\eta}).
\]

Here \(\kappa_{\mathrm{orb}}\) is the sum of the orbit-volume normalization
constants.  The comparison with Borovoi's Theorem 1.4.1 and Proposition
1.5.1 is carried out in
\cref{koide-count-thm-linear-asymptotic}; specifically,
\eqref{koide-count-eq-singular-series},
\eqref{koide-count-eq-archimedean-asymptotic}, and
\eqref{koide-count-eq-final-constant} give

\[
 \kappa_{\mathrm{orb}}=\mathfrak S(\mathscr X)
 =\frac{9\sqrt3}{\pi^2}\log(2+\sqrt3),
 \qquad
 \kappa_{\mathrm{orb}} I_\infty=C_0.
\]

Together with \(m_\infty(B_X)=I_\infty X+O(1)\), this yields

\[
 R(X)=C_0X+O(X^{1-\eta}).
\]

This counts the actual integral lattice. No assertion that \(\mathscr X(\mathbb Z)\) is one orbit, and no replacement by a maximal order, is involved.

\subsubsection{Why the exponent remains unspecified}\label{koide-lit-why-the-exponent-remains-unspecified}

The proof of Benoist--Oh Theorem 13.2 balances an effective-equidistribution/Sobolev error and a boundary-thickening error:

\[
 c_1\epsilon^{-p}v_X^{1-\delta}+c_2\epsilon^\kappa v_X,
 \qquad v_X=\operatorname{vol}(B_X).
\]

The resulting saving is of the form

\[
 \eta=\frac{\delta\kappa}{p+\kappa}>0.
\]

Neither the Sobolev exponent \(p\) nor the effective equidistribution exponent \(\delta\) is made numerically strong enough in the theorem to show \(\eta\geq1/2\). Thus Benoist--Oh supplies a genuine exact-family power saving, but not the requested square-root bound.

\subsection{The sharp hyperbolic frontier}\label{koide-lit-the-sharp-hyperbolic-frontier}

\subsubsection{Good and Chatzakos--Petridis}\label{koide-lit-good-and-chatzakospetridis}

For a cocompact or cofinite Fuchsian group and a fixed hyperbolic conjugacy class \(\mathcal H\), Chatzakos--Petridis define, in equations (1.4)--(1.6),

\[
 M(\mathcal H,X;z)
 =\sum_{1/2<s_j\leq1}
 A(s_j)\widehat u_j u_j(z)X^{s_j},
\]

where \(\lambda_j=s_j(1-s_j)<1/4\) and \(\widehat u_j\) is the period of \(u_j\) along the invariant closed geodesic. Their Theorem 1.1, author manuscript p.~3, restates Good's Theorem 4, p.~116, and proves

\[
 N(\mathcal H,X;z)-M(\mathcal H,X;z)=O(X^{2/3}).
\]

The \(2/3\) arises from the sharp majorant/minorant error

\[
 O(XY^{-1/2}+Y+X^{1/2}\log X)
\]

after choosing \(Y=X^{2/3}\). Theorems 1.2--1.3 are mean-square results, not pointwise improvements. Conjecture 5.7, author manuscript p.~22, explicitly asks for

\[
 O_\varepsilon(X^{1/2+\varepsilon}).
\]

There are two independent failures of direct transfer to \(R(X)\):

\begin{enumerate}
\def\labelenumi{\arabic{enumi}.}

\item
  the Huber cutoff is radial, whereas \(4a+17n\leq X\) plus the chamber is a nonradial sharp polyhedron; norm comparability does not preserve the leading constant or error because the symmetric difference can have volume of order \(X\);
\item
  even in the radial problem, the theorem subtracts every nonzero small-eigenvalue contribution, not only the linear term.
\end{enumerate}

For the exact Koide sector no spectral expansion establishing actual \(X^{s_j}\) terms has yet been produced. They are therefore a necessary obstruction to check, not terms that this audit claims are already known to occur. A proof of \(R(X)-C_0X=O(X^{1/2+\varepsilon})\) would have to show absence or vanishing of all relevant complementary-spectrum coefficients, or explicitly subtract them and then prove they vanish in the target formulation.

\subsubsection{\texorpdfstring{Current sharp refinements remain at \(2/3\)}{Current sharp refinements remain at 2/3}}\label{koide-lit-current-sharp-refinements-remain-at-23}

Voskou, Theorem 3 of arXiv:2407.03134v3, proves \(O(X^{2/3})\) for four prescribed orientation/side sectors in a closed-geodesic-to-closed-geodesic problem, with all small-spectrum period terms included in the main expression. Conjecture 1 on the same page proposes \(O_\varepsilon(X^{1/2+\varepsilon})\), while Theorem 4 is a mean-square result.

Lekkas--Petridis, in the 24 August 2026 version of arXiv:2509.12902, Theorem 1.2, likewise prove

\[
 N(X,l)=\sum_{1/2<s_j\leq1}
 \frac2\pi\gamma_1(s_j)|\widehat u_j|^2X^{s_j}
 +O(X^{2/3})
\]

for a closed-geodesic double-coset count. Their Theorem 1.5 gives

\[
 \frac1X\int_X^{2X}|E(x,l)|^2\,dx\ll X\log^2X,
\]

and Conjecture 1.6 proposes the pointwise \(X^{1/2+\epsilon}\) bound. Their introduction states that Good's \(2/3\) has not been improved for any of the nine point/geodesic/cusp double-coset configurations. This is strong current evidence about the sharp frontier, but neither paper treats the exact Koide height.

\subsubsection{Fixed angular harmonics do not give a sharp interval for free}\label{koide-lit-fixed-angular-harmonics-do-not-give-a-sharp-interval-for-free}

Bourgain--Kontorovich--Sarnak, Theorem 1.1, prove for fixed left and right \(K\)-Fourier modes \(n,k\)

\[
 \sum_{\|\gamma\|<T}e^{2in\theta_1(\gamma)}e^{2ik\theta_2(\gamma)}
 =\mathbf1_{n=0}\mathbf1_{k=0}c_0T^2
  +\sum_jc_j(n,k)T^{2s_j}+O_{n,k}(T^{4/3}).
\]

On the adjoint orbit, affine height is quadratic in the group norm, so \(T^{4/3}\) again corresponds to \(X^{2/3}\). Recovering a sharp angular interval requires summing many Fourier modes, while the theorem's implied constants depend on \(n,k\). It is not a uniform desmoothing theorem for this chamber.

\subsection{Smooth spectral counting}\label{koide-lit-smooth-spectral-counting}

Liu--Sarnak's paper is the closest source with a square-root-plus-gap exponent. Its global arithmetic setup and the hypotheses of its analytic proposition must be separated.

The introduction, author PDF p.~2, fixes an indefinite anisotropic integral ternary form \(f\), a nonzero level \(t\), and assumes \(t\det(f)\) is square-free. Thus the paper's affine-sieve theorem, as globally stated, does not cover \(t\det(f)=-108\). Equations (2.3)--(2.4), author PDF p.~4, fix one arithmetic orbit \(\mathcal O\) and a smooth \(K\)-invariant radial weight \(F_T\) satisfying

\[
 F_T=1\quad(|x|\leq T/c_0),
 \qquad
 F_T=0\quad(|x|\geq c_0T),
\]

with \(c_0>1\) fixed. In Fermi coordinates, Proposition 3.2 uses a smooth even function equal to one for \(|\rho|\leq R-1\), zero for \(|\rho|\geq R+1\), and with uniformly bounded derivatives. Unlike the global sieve statement, Proposition 3.2 is formulated for any variable subgroup \(\Gamma<\mathrm{SL}_2(\mathbb R)\) satisfying its uniform local condition (3.1), the spectral gap (2.7), and containing the fixed primitive cyclic group \(\Gamma_\xi\). It proves

\[
 H_R(z)=\frac{X(R)}{\operatorname{Area}(\Gamma\backslash\mathbb H)}
 +O_\varepsilon(e^{(1/2+\theta+\varepsilon)R}),
 \qquad X(R)\asymp e^R.
\]

Equation (3.60) has \(T=\sinh R\), and equation (3.65) transfers this to the one-sheeted affine orbit. Thus the smooth affine error is

\[
 O_\varepsilon(T^{1/2+\theta+\varepsilon}).
\]

Equation (2.7) assumes \(\lambda_1\geq1/4-\theta^2\). The paper invokes Jacquet--Langlands and Kim--Sarnak to take \(\theta=7/64\), giving

\[
 O_\varepsilon(T^{39/64+\varepsilon})
\]

unconditionally in its congruence-arithmetic setting; \(\theta=0\), and
hence square root, is conditional on the generalized
Selberg/Ramanujan-at-infinity conjecture for the relevant
Jacquet--Langlands transfers.

The numerical input itself is Kim--Sarnak, Appendix 2, Proposition 2, printed p.~176. It applies to an arbitrary cuspidal automorphic representation of \(\mathrm{GL}_2/\mathbb Q\): for a spherical real component, \(|\Re\mu_{j,\infty}|\leq7/64\), independently of finite conductor. Jacquet--Langlands preserves the real component because the quaternion algebra is split at infinity. Chu--Li, proof of Corollary 1.2, equation (5.2), printed p.~5126, records the resulting

\[
 \lambda_1\geq\frac{975}{4096}
\]

for every congruence arithmetic Fuchsian group. Thus the nonmaximal or
non-Eichler structure of the exact finite order does not invalidate
\(s\leq39/64\).  It changes the finite compact-open subgroup and hence which
automorphic representations have fixed vectors and the dimensions of those
fixed-vector spaces; it does not change this archimedean bound. This
validates the spectral-gap input for the exact lattice. It does not provide
sharp-boundary control or force the periods of possible
complementary-spectrum eigenfunctions to vanish.

For the exact order lattice \(\Gamma=\mathcal O^1\), each primitive trace-four conjugacy class supplies the cyclic subgroup in Proposition 3.2, and there are only finitely many such classes. Hence the proposition and equation (3.65) do give an unconditional \(39/64\) estimate after summing the corresponding \textbf{smooth standard-radial} counts. The square-free hypothesis is not an obstruction to this isolated analytic application.

This still does not desmooth to the target. A transition of fixed width in \(R\) is a fixed multiplicative annulus in \(T\), whose invariant volume is of order \(T\), not \(T^{1/2}\). Shrinking the transition makes the derivative/Sobolev constants in Proposition 3.2 grow, and the paper gives no uniform estimate that balances this loss at square-root scale. It also has no sharp coordinate chamber, and its radial height is not \(4a+17n\). These geometric failures remain even after the exact lattice and finite class sum have been supplied.

The same distinction applies to any informal statement that a tempered spectral component has square-root size. Duke--Rudnick--Sarnak's discussion on p.~155 estimates individual tempered frequencies; it does not control the complete spectral sum evaluated at a point with a shrinking approximation to a sharp characteristic function.

\subsection{General effective sector and full-lattice theorems}\label{koide-lit-general-effective-sector-and-full-lattice-theorems}

\subsubsection{Gorodnik--Nevo's explicit generic exponent}\label{koide-lit-gorodniknevos-explicit-generic-exponent}

Gorodnik--Nevo Definition 3.1 defines

\[
 n_e(2)=1,
 \qquad
 n_e(p)=\text{the least even integer at least }p/2\quad(p>2).
\]

Their Section 8 considers

\[
 S_t(\Phi,\Psi)
 =\{g:\log\|gv_0\|\leq t\}\cap\Phi A^+\Psi.
\]

Theorem 8.1, author PDF p.~34, assumes \(L^{p+}\) matrix coefficients and Lipschitz-well-rounded fixed angular sets; Corollary 8.3, p.~35, passes to an orbit when \((\Gamma\cap H)\backslash H\) is compact. The exponent is

\[
 \zeta=\frac{1}{2n_e(p)(1+3\dim G)}.
\]

For \(\dim G=3\),

\[
 \zeta=\frac1{20n_e(p)}.
\]

Even tempered spectrum has \(p=2\), \(n_e(p)=1\), and hence absolute error \(O(X^{19/20+\epsilon})\). For the exact congruence arithmetic lattice, Jacquet--Langlands and Kim--Sarnak give \(s\leq39/64\). The corresponding matrix coefficients are \(L^{p+}\) for \(p=64/25\), so \(n_e(p)=2\) and the theorem gives only \(O(X^{39/40+\epsilon})\).

This numerical calculation is an upper benchmark for the published generic polar-sector theorem, not a claim that Corollary 8.3 literally counts \(R(X)\). A fixed polar sector contains every radial parameter in fixed polar directions. An arbitrary affine linear cone and its finite-radius polyhedral height need an exact polar decomposition or a quantitatively controlled symmetric-difference argument; all integral orbits must then be included. None of those identifications appears in the theorem statement.

Gorodnik--Nevo themselves note after Corollary 5.3, author PDF p.~23, that their sharp-set counting mechanism remains worse than square-root volume error even for tempered spectrum, whereas square-root quality is available there only for a smooth weighted problem.

\subsubsection{Browning--Gorodnik and subsequent full-lattice work}\label{koide-lit-browninggorodnik-and-subsequent-full-lattice-work}

Browning--Gorodnik Theorem 2.1 gives, for a Euclidean ball \(B_H\) in one congruence orbit,

\[
 |\Gamma_\ell y\cap B_H|
 =\frac{m_{Z_{y,\ell}}(Z_{y,\ell})}{m_{X_\ell}(X_\ell)}m_Y(B_H)
 +O\!\left(\ell^{\dim H+\dim G}m_Y(B_H)^{1-\varrho}\right)
\]

for some unspecified \(\varrho>0\). Corollaries 2.5--2.6 sum adelic and integral orbits and retain the Brauer/Kottwitz density, at the cost of \(\ell^{\dim H+2\dim G}\). This is a useful exact-lattice precedent, but its archimedean set is the full norm ball, not the present sector, and no square-root exponent is asserted.

Cao--Huang Theorem 1.5 proves arithmetic purity of the Hardy--Littlewood property for an anisotropic ternary affine quadric with \(-m\det(q)\) nonsquare. Its Theorem 3.4 invokes Nevo--Sarnak, Gorodnik--Nevo, and Browning--Gorodnik to furnish an unspecified positive exponent, denoted \(\beta_X\) there with \(X\) meaning the affine variety; write it here as \(\beta_{\mathscr X}>0\) to distinguish it from the height variable. It is not a sharp-sector remainder theorem.

The current Cao--Huang--Zhang preprint, arXiv:2509.23320, Setting 2.1 and Theorem 2.1, gives a congruence-uniform power saving

\[
 O\!\left(\mathbf N(\mathfrak l)^\kappa
 \operatorname{vol}(\mathbb B_S(N))^{1-\delta}\right)
\]

for standard product Euclidean/max height balls and a compact piecewise-smooth factor. The exponent \(\delta\) is unspecified, and the unbounded real factor in the target is neither compact nor their standard Euclidean ball. This preprint extends the general effective framework but does not improve the conclusion of the direct Benoist--Oh argument in \cref{koide-lit-definition-13.1-for-the-actual-b_x,koide-lit-finite-orbit-passage-and-main-term-normalization}.

\subsection{Results in different arithmetic regimes}\label{koide-lit-results-in-different-arithmetic-regimes}

\subsubsection{Square discriminant and divergent stabilizer}\label{koide-lit-square-discriminant-and-divergent-stabilizer}

Oh--Shah Theorem 1.2 assumes

\[
 \operatorname{vol}((H\cap\Gamma)\backslash H)=\infty
\]

and obtains a smoothed \(cT\log T+O(T)\) count. Its sharp version has only a logarithmic relative error. Theorem 6.1 allows an open cone and arbitrary norm, but retains the logarithmic main-factor characteristic of its divergent-stabilizer regime; that regime is excluded here.

Xu--Zhang, \emph{Quarterly Journal of Mathematics} 74 (2023), 659--685, assumes \(-m\det(f)\) is a square over the number field. Huang, arXiv:2412.03350, Theorem 1.1, assumes the same split condition and a fixed smooth compact weight; its principal error is \(B(\log B)^{1983/1984}\) in a \(B\log B\) count. Zhang, arXiv:2603.21080, Section 1.1 and Theorem 1.1, again assumes \(-m\det(q)\in(k^\times)^2\) and proves a smooth \(K_0\)-invariant expansion with an error \(O(T^{l_1+l_2-\delta})\), \(0<\delta<1/2\); Remark 1.2 obtains an authentic count for \(\Gamma=\mathrm{SL}_2(\mathcal O_k)\). These are important current results, but all lie in the rationally split torus regime excluded by \(108\notin\mathbb Q^{\times2}\), and none uses the target polyhedral height.

\subsubsection{Varying level, isotropic forms, and special rotational quadrics}\label{koide-lit-varying-level-isotropic-forms-and-special-rotational-quadrics}

Huang, arXiv:2412.05903, studies \(F(x)=m_0N\) with \(N=p_0^{2h}\to\infty\), a smooth compactly supported real weight \(w(x/\sqrt N)\), and prescribed local conditions. Theorems 1.1--1.2 have a \(\sqrt N/(\log N)^A\) error. This is a Linnik-type varying-level theorem, not a fixed affine shell.

Li--Liu, arXiv:2501.15033v2 (18 April 2025), Theorems 1.1--1.2, prove new almost-prime conclusions for anisotropic ternary equations. Their global assumption remains \(t\det(f)\) square-free. The quantitative counting input in Lemma 2.3, pp.~5--6, is

\[
 O_\varepsilon(d^{1+\varepsilon}T^{1/2+\theta+\varepsilon})
\]

for the same kind of smooth radial orbit weight inherited from Liu--Sarnak. It is a current application of the \(7/64\) gap, not a new unsmoothed sector theorem.

Lowry-Duda's thesis, arXiv:1704.02376, Theorem 5.5.4, treats the special rotational form \(x_1^2+x_2^2=x_3^2+h\) in a full Euclidean ball. Written in ordinary radius \(T\), its three-dimensional sharp error is \(O(T^{21/22+\epsilon})\). The homogeneous form is \(\mathbb Q\)-isotropic, and neither the lattice nor the height sector is the target one.

Results on values of irrational ternary forms, shrinking sets while the discriminant varies, or positive-definite ternary representation numbers answer different problems. Lindelöf or generalized Lindelöf hypotheses occurring there do not yield the present fixed-level pointwise sector bound.

\subsection{Spectral-gap, GRH, smoothing, and averaging taxonomy}\label{koide-lit-spectral-gap-grh-smoothing-and-averaging-taxonomy}

\begin{longtable}[]{@{}
  >{\raggedright\arraybackslash}p{0.31\linewidth}
  >{\raggedright\arraybackslash}p{0.31\linewidth}
  >{\raggedright\arraybackslash}p{0.31\linewidth}@{}}
\toprule\noalign{}
\begin{minipage}[b]{\linewidth}\raggedright
Classification
\end{minipage} & \begin{minipage}[b]{\linewidth}\raggedright
Source result
\end{minipage} & \begin{minipage}[b]{\linewidth}\raggedright
Why it does not prove the target
\end{minipage} \\
\midrule\noalign{}
\endhead
\bottomrule\noalign{}
\endlastfoot
Unconditional, exact sharp family & Benoist--Oh Theorem 13.2 plus the direct Definition 13.1 verification & Gives only \(O(X^{1-\eta})\) with \(\eta>0\) unspecified \\
Unconditional, exact norm/cone, full lattice & Gorodnik--Oh--Shah Theorem 1.6, with the exact hypotheses proved in \cref{koide-count-lem-genuine-norm,koide-count-lem-gos-hypotheses} & Qualitative \(C_0X+o(X)\) only \\
Unconditional, sharp radial closed-geodesic count & Good / Chatzakos--Petridis & \(O(X^{2/3})\), different height, and main expression includes small-spectrum terms \\
Unconditional, smooth radial count & Liu--Sarnak Proposition 3.2 with \(\theta=7/64\), summed over the exact finite conjugacy-class decomposition & \(39/64\) for the exact lattice, but fixed multiplicative smoothing and a different radial height \\
Selberg/Ramanujan conditional & Liu--Sarnak with \(\theta=0\) & Square-root only for the smooth radial problem; still not the exact count \\
GRH/Lindelöf conditional & None among the source versions in \cref{koide-lit-primary-references-and-exact-locators} proves this fixed sharp-sector estimate & The occurrences listed there concern varying levels/forms, shrinking-set equidistribution, or lower/Omega results \\
Smoothed & Liu--Sarnak; Oh--Shah; Huang I; Zhang 2026 & Smoothing or split arithmetic changes the theorem; no square-root desmoothing supplied \\
Averaged/mean square & Chatzakos--Petridis Theorems 1.2--1.3; Voskou Theorem 4; Lekkas--Petridis Theorem 1.5 & Does not imply a uniform pointwise estimate at every \(X\) \\
Omega/lower-bound & Chatzakos (2019), including a cocompact geodesic-average \(\Omega(X^{1/2}\log\log\log X)\) & Evidence for scale, not an upper bound \\
Isotropic or split & Oh--Shah; Xu--Zhang; Huang I; Zhang 2026; Lowry-Duda & Different stabilizer, quotient, main order, form, or level \\
Varying level & Huang II & Counts \(F=m_0p_0^{2h}\), not \(F=-2\) \\
\end{longtable}

In particular, ``assuming GRH'' is not a substitute for the missing theorem. The conditional hypothesis that occurs naturally in the closest smooth spectral result is the Selberg/Ramanujan eigenvalue conjecture, and even it only reaches square root before the sharp-boundary unsmoothing problem is addressed.

\subsection{Citation-ready conclusion}\label{koide-lit-citation-ready-conclusion}

Subject to the bounded-search qualification in
\cref{koide-lit-audit-limitations}, the following paragraph records the
conclusion supported by the enumerated source versions:

\begin{quote}
Among the source versions enumerated in \cref{koide-lit-primary-references-and-exact-locators}, none implies the unconditional sharp estimate \(R(X)-C_0X=O_\varepsilon(X^{1/2+\varepsilon})\) for the fixed ternary quadric \(a^2-3q^2-18n^2=-2\) in the chamber \(a>0\), \(q,n\geq0\), \(a+4n>q\), with height \(4a+17n\leq X\). Gorodnik--Oh--Shah, Theorem 1.6, applies to the exact admissible generic cone and an arbitrary norm---with those exact hypotheses proved in \cref{koide-count-lem-genuine-norm,koide-count-lem-gos-hypotheses}---but gives only a qualitative asymptotic. Benoist--Oh, Definition 13.1 and Theorem 13.2, applies after the direct effective-well-roundedness verification in \cref{koide-lit-definition-13.1-for-the-actual-b_x} and, after summing the finitely many integral spin-group orbits, yields \(R(X)=C_0X+O(X^{1-\eta})\) for some unspecified \(\eta>0\), not \(\eta=1/2\). In the closest sharp closed-geodesic problem, Good's theorem, as restated and reproved by Chatzakos--Petridis, Theorem 1.1, gives \(O(X^{2/3})\) only after all small-eigenvalue terms \(X^{s_j}\) are included; Chatzakos--Petridis, Conjecture 5.7, states the \(O_\varepsilon(X^{1/2+\varepsilon})\) improvement as a conjecture. Liu--Sarnak, Proposition 3.2, applies class by class to the exact congruence quaternionic lattice and proves the smooth radial bound \(O_\varepsilon(X^{1/2+\theta+\varepsilon})\), unconditionally with \(\theta=7/64\) and conditionally with \(\theta=0\), but its fixed-width smoothing and radial height do not match the exact sharp count. Thus complementary-spectrum terms and sharp nonradial boundary control remain separate unresolved inputs.
\end{quote}

\subsection{Primary references and exact locators}\label{koide-lit-primary-references-and-exact-locators}

\begin{enumerate}
\def\labelenumi{\arabic{enumi}.}

\item
  Yves Benoist and Hee Oh, ``Effective equidistribution of \(S\)-integral points on symmetric varieties,'' \emph{Annales de l'Institut Fourier} \textbf{62} (2012), 1889--1942. \href{https://doi.org/10.5802/aif.2738}{DOI}, \href{https://www.numdam.org/item/AIF_2012__62_5_1889_0/}{official article/PDF}, \href{https://arxiv.org/abs/0706.1621}{arXiv}. Standing setting pp.~1892--1893; finite orbit statement p.~1892; Theorem 1.3 p.~1894; Theorem 12.5 p.~1929; Definition 13.1 p.~1931; Theorem 13.2 pp.~1931--1932; Proposition 14.2 p.~1933; ternary Example (A) p.~1896.
\item
  Alexander Gorodnik, Hee Oh, and Nimish Shah, ``Integral points on symmetric varieties and Satake compactifications,'' \emph{American Journal of Mathematics} \textbf{131} (2009), 1--57. \href{https://doi.org/10.1353/ajm.0.0034}{DOI}, \href{https://arxiv.org/abs/math/0610497}{arXiv}, \href{https://gauss.math.yale.edu/~ho2/doc/satake_ajm.pdf}{author PDF}. Admissible/generic definitions and Theorem 1.6 pp.~4--5; Theorem 1.18 pp.~8--9.
\item
  Alexander Gorodnik, Hee Oh, and Nimish Shah, ``Strong wavefront lemma and counting lattice points in sectors,'' \emph{Israel Journal of Mathematics} \textbf{176} (2010), 419--444. \href{https://doi.org/10.1007/s11856-010-0035-8}{DOI}, \href{https://arxiv.org/abs/0708.3188}{arXiv}. Theorem 1.13, author PDF pp.~4--5.
\item
  Alexander Gorodnik and Amos Nevo, ``Counting lattice points,'' \emph{Journal für die reine und angewandte Mathematik} \textbf{663} (2012), 127--176. \href{https://doi.org/10.1515/CRELLE.2011.096}{DOI}, \href{https://arxiv.org/abs/0903.1515}{arXiv}, \href{https://www.math.uzh.ch/gorodnik/papers/GN2-final.pdf}{author PDF}. Theorem 1.9 pp.~4--5; Definition 3.1 p.~13; Theorem 4.5 pp.~17--18; Theorem 8.1 and Corollary 8.3 pp.~34--35; Proposition 8.4 p.~35.
\item
  William Duke, Zeév Rudnick, and Peter Sarnak, ``Density of integer points on affine homogeneous varieties,'' \emph{Duke Mathematical Journal} \textbf{71} (1993), 143--179. \href{https://doi.org/10.1215/S0012-7094-93-07107-4}{DOI}, \href{https://www.math.ucla.edu/~wdduke/preprints/density.pdf}{author PDF}. Theorem 1.2 p.~145; Example 1.5 pp.~146--147; discussion of individual tempered frequencies p.~155; Theorem 3.5 p.~161.
\item
  Alex Eskin and Curtis McMullen, ``Mixing, counting, and equidistribution in Lie groups,'' \emph{Duke Mathematical Journal} \textbf{71} (1993), 181--209. \href{https://doi.org/10.1215/S0012-7094-93-07108-6}{DOI}, \href{https://people.math.harvard.edu/~ctm/papers/home/text/papers/mixing/mixing.pdf}{author PDF}. Theorem 1.4, author PDF p.~3; Theorems 6.1 and 6.4 pp.~30--31; Open Question 1 p.~35.
\item
  Jianya Liu and Peter Sarnak, ``Integral points on quadrics in three variables whose coordinates have few prime factors,'' \emph{Israel Journal of Mathematics} \textbf{178} (2010), 393--426. \href{https://doi.org/10.1007/s11856-010-0069-y}{DOI}, \href{https://web.math.princeton.edu/facultypapers/sarnak/few-final.pdf}{author PDF}. Arithmetic hypothesis p.~2; equations (2.3)--(2.8) pp.~4--5; Proposition 3.2 pp.~10--11; equations (3.60) and (3.65) pp.~15--16.
\item
  Anton Good, \emph{Local Analysis of Selberg's Trace Formula}, Lecture Notes in Mathematics 1040, Springer, 1983. \href{https://doi.org/10.1007/BFb0073074}{DOI}. Theorem 4, p.~116. The theorem statement and a new proof were checked through Chatzakos--Petridis (2016) because the monograph text was not directly available in this audit.
\item
  Dimitrios Chatzakos and Yiannis N. Petridis, ``The hyperbolic lattice point problem in conjugacy classes,'' \emph{Forum Mathematicum} \textbf{28} (2016), 981--1003. \href{https://doi.org/10.1515/forum-2015-0102}{DOI}, \href{https://arxiv.org/abs/1504.01307}{arXiv}, \href{https://discovery.ucl.ac.uk/1477487/7/Petridis_forum-2015-0102.pdf}{author manuscript}. Equations (1.4)--(1.6) and Theorem 1.1 p.~3; Theorems 1.2--1.3 pp.~3--4; Conjecture 5.7 p.~22.
\item
  Dimitrios Chatzakos, ``Mean value results and \(\Omega\)-results for the hyperbolic lattice point problem in conjugacy classes,'' \emph{Revista Matemática Iberoamericana} \textbf{35} (2019), 1123--1152. \href{https://doi.org/10.4171/RMI/1080}{DOI}, \href{https://arxiv.org/abs/1610.01462}{arXiv}. Mean-value and Omega theorems, not a pointwise upper bound.
\item
  Jean Bourgain, Alex Kontorovich, and Peter Sarnak, ``Sector estimates for hyperbolic isometries,'' \emph{Geometric and Functional Analysis} \textbf{20} (2010), 1175--1200. \href{https://doi.org/10.1007/s00039-010-0092-5}{DOI}, \href{https://arxiv.org/abs/1001.4541}{arXiv}. Theorem 1.1, p.~2.
\item
  Hee Oh and Nimish Shah, ``Limits of translates of divergent geodesics and integral points on one-sheeted hyperboloids,'' \emph{Israel Journal of Mathematics} \textbf{199} (2014), 915--931. \href{https://doi.org/10.1007/s11856-013-0063-2}{DOI}, \href{https://arxiv.org/abs/1104.4988}{arXiv}. Theorem 1.2 p.~3; Corollary 1.4 pp.~3--4; Theorem 6.1 pp.~13--14.
\item
  Tim D. Browning and Alexander Gorodnik, ``Power-free values of polynomials on symmetric varieties,'' \emph{Proceedings of the London Mathematical Society} \textbf{114} (2017), 1044--1080. \href{https://doi.org/10.1112/plms.12030}{DOI}, \href{https://arxiv.org/abs/1606.06342}{arXiv}. Theorem 2.1 p.~11; Corollaries 2.5--2.6 pp.~24--25.
\item
  Yang Cao and Zhizhong Huang, ``Arithmetic purity of the Hardy--Littlewood property and geometric sieve for affine quadrics,'' \emph{Advances in Mathematics} \textbf{398} (2022), 108236. \href{https://doi.org/10.1016/j.aim.2022.108236}{DOI}, \href{https://arxiv.org/abs/2003.07287}{arXiv}. Theorem 1.5 p.~4; Proposition 2.6 pp.~12--13; Theorem 3.4 pp.~21--22.
\item
  Yang Cao, Zhizhong Huang, and Runlin Zhang, ``Effective equidistribution, arithmetic purity of strong approximation, and geometric sieve for affine quadrics,'' \href{https://arxiv.org/abs/2509.23320}{arXiv:2509.23320}. Setting 2.1 and Theorem 2.1; arXiv HTML retrieved 4 September 2026.
\item
  Mikhail Borovoi, ``On representations of integers by indefinite ternary quadratic forms,'' \emph{Journal of Number Theory} \textbf{90} (2001), 281--293. \href{https://doi.org/10.1006/jnth.2001.2662}{DOI}, \href{https://arxiv.org/abs/math/0006141}{arXiv}, \href{https://www.math.tau.ac.il/~borovoi/papers/reprfin.pdf}{author PDF}. Theorem 1.4.1 and Proposition 1.5.1; qualitative local/Brauer--Manin coefficient, no error exponent.
\item
  Fei Xu and Runlin Zhang, ``Counting integral points on indefinite ternary quadratic equations over number fields,'' \emph{Quarterly Journal of Mathematics} \textbf{74} (2023), 659--685. \href{https://doi.org/10.1093/qmath/haac039}{DOI}, \href{https://arxiv.org/abs/2103.10707}{arXiv}. Square-discriminant hypothesis in the main theorem.
\item
  Zhizhong Huang, ``Quantitative strong approximation for ternary quadratic
  forms I.''

  \href{https://arxiv.org/abs/2412.03350}{arXiv:2412.03350}; Theorem 1.1;
  square-discriminant, smooth fixed-weight \(B\log B\) count.
\item
  Zhizhong Huang, ``Quantitative strong approximation for ternary quadratic
  forms II.''

  \href{https://arxiv.org/abs/2412.05903}{arXiv:2412.05903}; Theorems 1.1--1.2;
  varying level \(m_0p_0^{2h}\) and smooth compact weight.
\item
  Runlin Zhang, ``Effective count of integer points on ternary affine quadrics and effective equidistribution,'' \href{https://arxiv.org/abs/2603.21080}{arXiv:2603.21080}. Section 1.1 and Theorem 1.1; split-stabilizer smooth count; Remark 1.2 on the authentic radial count.
\item
  Marios Voskou, ``Refined Counting of Geodesic Segments in the Hyperbolic
  Plane.''

  \href{https://arxiv.org/abs/2407.03134}{arXiv:2407.03134v3}; Theorems 3--4
  and Conjecture 1, pp.~3--4.
\item
  Dimitrios Lekkas and Yiannis N. Petridis, ``A relative trace formula and counting geodesic arcs in the hyperbolic plane,'' \href{https://arxiv.org/abs/2509.12902}{arXiv:2509.12902}, version dated 24 August 2026. Theorem 1.2, Theorem 1.5, and Conjecture 1.6.
\item
  Henry H. Kim, ``Functoriality for the exterior square of \(\mathrm{GL}_4\) and the symmetric fourth of \(\mathrm{GL}_2\),'' with Appendix 2 by Henry H. Kim and Peter Sarnak, ``Refined estimates towards the Ramanujan and Selberg conjectures,'' \emph{Journal of the American Mathematical Society} \textbf{16} (2003), 139--183. \href{https://doi.org/10.1090/S0894-0347-02-00410-1}{DOI}, \href{https://www.ams.org/journals/jams/2003-16-01/S0894-0347-02-00410-1/S0894-0347-02-00410-1.pdf}{official AMS PDF}. Appendix 2, Proposition 2, printed p.~176: \(|\Re\mu_{j,\infty}|\leq7/64\) for a spherical real component of any cuspidal \(\mathrm{GL}_2/\mathbb Q\) representation, and the displayed consequence \(\lambda_1\geq975/4096\).
\item
  Michelle Chu and Han Li, ``Small generators of cocompact arithmetic Fuchsian groups,'' \emph{Proceedings of the American Mathematical Society} \textbf{144} (2016), 5121--5127. \href{https://doi.org/10.1090/proc/13177}{DOI}, \href{https://arxiv.org/abs/1603.00926}{arXiv}. Section 3.1, p.~5122, defines congruence arithmetic Fuchsian groups; proof of Corollary 1.2, equation (5.2), p.~5126, records \(\lambda_1\geq975/4096\) for every such group via Jacquet--Langlands and Kim--Sarnak.
\item
  Hervé Jacquet and Robert P. Langlands, \emph{Automorphic Forms on \(\mathrm{GL}(2)\)}, Lecture Notes in Mathematics 114, Springer, 1970. \href{https://doi.org/10.1007/BFb0058988}{DOI}, \href{https://publications.ias.edu/sites/default/files/automorphic-forms-on-gl2_rpl_9.pdf}{IAS-hosted text}. Chapter III, especially §16, supplies the global quaternionic transfer; at split places, including the real place here, the local component is preserved.
\item
  John Voight, \emph{Quaternion Algebras}, Graduate Texts in Mathematics 288, Springer, 2021. \href{https://doi.org/10.1007/978-3-030-56694-4}{DOI}, \href{https://jvoight.github.io/quat-book.pdf}{open author version}. Proposition 15.5.2 (maximal superorders); §23.1.3 (Eichler orders at ramified primes); equations (27.6.6)--(27.6.7) (compact-open finite adelic order/unit group); Theorem 28.2.10 (strong approximation for indefinite \(B^1\)); Definition 38.7.2 (principal congruence groups).
\item
  Jiamin Li and Jianya Liu, ``Anisotropic quadratic equations in three
  variables.''

  \href{https://arxiv.org/abs/2501.15033}{arXiv:2501.15033v2};
  Theorems 1.1--1.2, pp.~2--3; Lemma 2.3, pp.~5--6. The global form-level
  hypothesis is \(t\det(f)\) square-free; the quantitative input is a smooth
  radial, congruence-orbit estimate with
  \[
   O_\varepsilon\!\left(
     d^{1+\varepsilon}T^{1/2+\theta+\varepsilon}
   \right).
  \]
\item
  David Lowry-Duda, \emph{On Some Variants of the Gauss Circle Problem}, \href{https://arxiv.org/abs/1704.02376}{arXiv:1704.02376}. Theorem 5.5.4, thesis pp.~97--98.
\end{enumerate}

\subsection{Audit limitations}\label{koide-lit-audit-limitations}

The claim-by-claim source versions and locators used in this bounded audit
are enumerated in
\cref{koide-lit-primary-references-and-exact-locators}.  The canonical tree
does not preserve byte-identical snapshots or retrieval hashes for those
external publications, so this file does not assert a reproducible local
source corpus.  Good's 1983 book was the sole exception to direct full-text
access: its theorem locator and statement were checked through the explicit
restatement and independent proof in Chatzakos--Petridis (2016) and the
version-dated Lekkas--Petridis treatment.

The statement that no theorem was located is a result about the source
versions enumerated here, not a proof that no unpublished or differently
indexed result exists.  It does establish that none of the enumerated works
may be cited as a proof of the exact unconditional square-root discrepancy.


\section{The exact gapless auxiliary operator and the physical boundary}
\label{cusp-interpretation-section}

The preceding flat-torus calculation produces eigenvalues tending to zero
along the radial cusp.  This section identifies a single self-adjoint
operator for which those numbers genuinely form a gapless spectrum, proves
that complex and Hodge-theoretic data alone cannot determine such a
spectrum, and states the exact additional morphism needed to transfer an
auxiliary Rayleigh sequence to a physical Hamiltonian.  In particular, the
word \emph{gapless} below is attached to a specified operator and Hilbert
space, not to the complex threefold in isolation.

\subsection{A direct-integral scalar Laplacian with zero essential threshold}
\label{cusp-interpretation-direct-integral-subsection}

Use the angular coordinates \(y\in\mathbb R^4/\mathbb Z^4\) of
\cref{cusp-spec-exact-dual-lattice}.  The map \(y\mapsto P_Ty\) identifies
this fixed torus with \(F_T\), and normalized Haar measure on \(F_T\) pulls
back to \(dy\).  Write
\[
 e_k(y):=e^{2\pi i k\cdot y},
 \qquad
 \lambda_k(T):=4\pi^2\lvert P_T^{-\mathsf T}k\rvert^2,
 \qquad k\in\mathbb Z^4.
\]
By \cref{cusp-spec-exact-dual-lattice}, these are exactly the
scalar-Laplacian eigenfunctions and eigenvalues on \((F_T,g_T)\), normalized
in \(L^2\) for Haar probability measure.

Fix \(T_0\) large enough that all radial formulas of
\cref{cusp-spec-exact-dual-lattice,cusp-spec-shortest-modes} hold, and put
\begin{equation}
 \mathscr H_{\mathrm{aux}}
 :=L^2\!\left([T_0,\infty),dT;
 L^2_0(\mathbb R^4/\mathbb Z^4,dy)\right).
 \label{cusp-interpretation-auxiliary-Hilbert-space}
\end{equation}
Fourier expansion gives the concrete unitary identification
\begin{equation}
 \mathscr H_{\mathrm{aux}}
 \cong
 \bigoplus_{k\in\mathbb Z^4\setminus\{0\}}
 L^2([T_0,\infty),dT).
 \label{cusp-interpretation-Fourier-direct-sum}
\end{equation}
Define the multiplication operator
\begin{equation}
 (\mathscr L_{\mathrm{aux}}u)_k(T)
 :=\lambda_k(T)u_k(T)
 \label{cusp-interpretation-direct-integral-operator}
\end{equation}
on the exact domain
\begin{equation}
 D(\mathscr L_{\mathrm{aux}})
 :=\left\{u\in\mathscr H_{\mathrm{aux}}:
 \sum_{k\ne0}\int_{T_0}^{\infty}
 \lambda_k(T)^2\lvert u_k(T)\rvert^2\,dT<\infty\right\}.
 \label{cusp-interpretation-direct-integral-domain}
\end{equation}
Thus \(\mathscr L_{\mathrm{aux}}\) is, by definition and not by analogy,
the direct integral of the mean-zero scalar Laplacians \(\Delta_{F_T}\).

\begin{theorem}[A proved gapless auxiliary operator]
\label{cusp-interpretation-direct-integral-gapless}
The operator \(\mathscr L_{\mathrm{aux}}\) in
\cref{cusp-interpretation-direct-integral-operator} is nonnegative and
self-adjoint.  It has no zero eigenvector, but
\begin{equation}
 0\in\sigma_{\mathrm{ess}}(\mathscr L_{\mathrm{aux}}).
 \label{cusp-interpretation-zero-essential-spectrum}
\end{equation}
In particular, this specified auxiliary operator has no positive spectral
gap above zero.
\end{theorem}

\begin{proof}
For each \(k\ne0\), \(\lambda_k(T)\) is a positive continuous real function.
Multiplication by it on
\[
 \{v\in L^2([T_0,\infty)):\lambda_kv\in L^2([T_0,\infty))\}
\]
is self-adjoint and nonnegative: its adjoint domain is obtained by testing
against compactly supported bounded functions, which forces
\(\lambda_kv\in L^2\), and the adjoint acts by the same real multiplier.
The orthogonal sum of these multiplication operators has precisely the
domain in \cref{cusp-interpretation-direct-integral-domain}; hence it is
self-adjoint and nonnegative.  Since every \(\lambda_k(T)\) is strictly
positive, \(\mathscr L_{\mathrm{aux}}u=0\) implies \(u_k=0\) almost
everywhere for every \(k\ne0\), so its kernel is zero.

Let \(k_0=(1,0,0,0)\).  The exact formula
\cref{cusp-spec-unit-mode-exact-lengths} gives
\begin{equation}
 \lambda_{k_0}(T)
 =4\pi^2\frac{L_T^2+m_T^2}{D_T^2}
 =\frac{4\pi^2}{T^2}\bigl(1+O(T^{-1})\bigr).
 \label{cusp-interpretation-soft-mode-eigenvalue}
\end{equation}
Choose integers \(n>T_0\), and define \(u^{(n)}\) by
\begin{equation}
 u^{(n)}_{k_0}(T)=\mathbf1_{[n,n+1]}(T),
 \qquad u^{(n)}_k(T)=0\quad(k\ne k_0).
 \label{cusp-interpretation-Weyl-sequence}
\end{equation}
The vectors \(u^{(n)}\) lie in
\(D(\mathscr L_{\mathrm{aux}})\), are orthonormal, and converge weakly to
zero because their \(T\)-supports escape to infinity.  Moreover,
\[
 \lVert\mathscr L_{\mathrm{aux}}u^{(n)}\rVert^2
 =\int_n^{n+1}\lambda_{k_0}(T)^2\,dT
 \le\sup_{T\in[n,n+1]}\lambda_{k_0}(T)^2
 \longrightarrow0
\]
by \cref{cusp-interpretation-soft-mode-eigenvalue}.  This orthonormal
singular sequence proves
\(0\in\sigma_{\mathrm{ess}}(\mathscr L_{\mathrm{aux}})\).  Explicitly, if
zero were outside the essential spectrum, then, because the kernel is zero,
some interval \([0,\varepsilon)\) would have finite-dimensional spectral
projection.  Its complement would satisfy
\(\lVert\mathscr L_{\mathrm{aux}}v\rVert\ge
\varepsilon\lVert v\rVert\); weak convergence would make the components of
\(u^{(n)}\) in the finite-dimensional spectral subspace tend strongly to
zero, contradicting
\(\lVert\mathscr L_{\mathrm{aux}}u^{(n)}\rVert\to0\) and
\(\lVert u^{(n)}\rVert=1\).
\end{proof}

\begin{remark}[The unit-volume normalization is a different operator family]
The theorem uses the unscaled Euclidean quotient metrics \(g_T\).  By
\cref{cusp-spec-unit-volume}, the first positive scalar eigenvalue for
\(\widehat g_T=D_T^{-1/2}g_T\) still tends to zero, but at the different
rate
\[
 \lambda_1(F_T,\widehat g_T)
 =\frac{4\pi^2}{T}\bigl(1+O(T^{-1})\bigr),
\]
rather than the unscaled \(T^{-2}\) rate.  Rescaling a metric changes its
Laplacian, so the two operator families and their asymptotic scales must not
be identified.
\end{remark}

\subsection{Complex data do not select a spectral gap}
\label{cusp-interpretation-metric-indeterminacy-subsection}

\begin{proposition}[Fixed complex torus, fixed volume, arbitrarily small gap]
\label{cusp-interpretation-fixed-complex-counterexample}
Let
\[
 A=(\mathbb C/(\mathbb Z+i\mathbb Z))^2
\]
with fixed complex coordinates
\(z_j=x_j+iy_j\).  For \(R\ge1\), the metric
\begin{equation}
 g_R=R^2(dx_1^2+dy_1^2)+R^{-2}(dx_2^2+dy_2^2)
 \label{cusp-interpretation-fixed-complex-metric}
\end{equation}
is a flat K\"ahler metric of volume one on the same complex torus, while
\begin{equation}
 \lambda_1(\Delta_{g_R})=\frac{4\pi^2}{R^2}\longrightarrow0.
 \label{cusp-interpretation-fixed-complex-gap}
\end{equation}
Consequently the complex structure, Hodge decomposition, integral
cohomology, and coherent direct-image data of a trivial family do not, even
together with total volume, determine a positive lower bound for the scalar
Laplacian gap.
\end{proposition}

\begin{proof}
The associated two-form
\[
 \omega_R=R^2dx_1\wedge dy_1+R^{-2}dx_2\wedge dy_2
\]
is closed, positive, and \(J\)-invariant, so \(g_R\) is K\"ahler for the fixed
complex structure.  In the ordered real coordinates
\((x_1,y_1,x_2,y_2)\), its matrix is
\(\operatorname{diag}(R^2,R^2,R^{-2},R^{-2})\), whose determinant is one;
therefore the volume is one.  The Fourier character
\[
 e_{a,b,c,d}
 =\exp(2\pi i(ax_1+by_1+cx_2+dy_2))
\]
has eigenvalue
\[
 4\pi^2\left(R^{-2}(a^2+b^2)+R^2(c^2+d^2)\right).
\]
For \(R\ge1\), the minimum over nonzero integer quadruples is
\(4\pi^2R^{-2}\), attained at the four vectors
\[
 (a,b,c,d)\in
 \{(1,0,0,0),(-1,0,0,0),(0,1,0,0),(0,-1,0,0)\}.
\]
This proves
\cref{cusp-interpretation-fixed-complex-gap}.  None of the listed complex or
cohomological objects depends on \(R\), and for a trivial holomorphic family
their coherent direct images are likewise unchanged.
\end{proof}

\subsection{The exact bridge to a physical Hamiltonian}
\label{cusp-interpretation-physical-bridge-subsection}

The mathematical mass gap in the Clay formulation is a property of the
Hamiltonian of a constructed four-dimensional quantum Yang--Mills theory,
not of an internal scalar Laplacian or of a classical Hessian
\cite{JaffeWitten2000}.  Euclidean input reaches such a Hamiltonian only
after a positive-metric reconstruction satisfying, in particular, the
Osterwalder--Schrader requirements
\cite{OsterwalderSchrader1973,OsterwalderSchrader1975}.  The following
proposition isolates the exact operator-theoretic statement that any proposed
transfer must prove.

\begin{proposition}[Typed Rayleigh-sequence transfer]
\label{cusp-interpretation-typed-physical-transfer}
Let \(H\ge0\) be self-adjoint on a Hilbert space
\(\mathcal H_{\mathrm{phys}}\).  Define
\begin{equation}
 \mathcal V_0:=\ker H,
 \qquad P_0:=\mathbf1_{\{0\}}(H),
 \qquad
 \mathcal H_{\mathrm{exc}}:=(I-P_0)\mathcal H_{\mathrm{phys}},
 \label{cusp-interpretation-excitation-space}
\end{equation}
and let
\begin{equation}
 \mathfrak h[\psi]:=\lVert H^{1/2}\psi\rVert^2,
 \qquad
 \mathcal Q(H):=D(H^{1/2}).
 \label{cusp-interpretation-Hamiltonian-form}
\end{equation}
For each \(n\), let \(\mathcal K_n\) be a Hilbert space,
\(\mathfrak q_n\) a nonnegative quadratic form on a domain
\(D(\mathfrak q_n)\subseteq\mathcal K_n\), and
\(0\ne u_n\in D(\mathfrak q_n)\).  Suppose there is a linear map
\[
 J_n:\operatorname{span}\{u_n\}\longrightarrow\mathcal Q(H),
\]
and put
\(\psi_n=(I-P_0)J_nu_n\).  If there are \(c>0\), \(C_n\ge0\), and
\(\varepsilon_n\ge0\) such that
\begin{align}
 \lVert\psi_n\rVert
 &\ge c\lVert u_n\rVert,
 \label{cusp-interpretation-norm-noncollapse}\\
 \mathfrak h[\psi_n]
 &\le C_n\mathfrak q_n[u_n]
       +\varepsilon_n\lVert u_n\rVert^2,
 \label{cusp-interpretation-energy-transfer}\\
 C_n\frac{\mathfrak q_n[u_n]}{\lVert u_n\rVert^2}
   +\varepsilon_n&\longrightarrow0.
 \label{cusp-interpretation-physical-smallness}
\end{align}
Then the excitation Hamiltonian
\[
 H_{\mathrm{exc}}
 :=H\big|_{D(H)\cap\mathcal H_{\mathrm{exc}}}
\]
is self-adjoint and satisfies
\begin{equation}
 \inf\sigma(H_{\mathrm{exc}})=0,
 \qquad
 0\in\sigma_{\mathrm{ess}}(H_{\mathrm{exc}}).
 \label{cusp-interpretation-physical-gapless-conclusion}
\end{equation}
\end{proposition}

\begin{proof}
The spectral projection \(P_0\) commutes with \(H\) and \(H^{1/2}\), so
\[
 \mathcal H_{\mathrm{phys}}
 =\ker H\oplus\mathcal H_{\mathrm{exc}}
\]
is a reducing decomposition.  Hence \(H_{\mathrm{exc}}\) is self-adjoint,
nonnegative, has form domain
\(\mathcal Q(H)\cap\mathcal H_{\mathrm{exc}}\), and has zero kernel.
Define \(w_n=\psi_n/\lVert\psi_n\rVert\).  Then
\(\lVert w_n\rVert=1\),
\(w_n\in\mathcal Q(H)\cap\mathcal H_{\mathrm{exc}}\), and
\begin{align*}
 \mathfrak h[w_n]
 &=\frac{\mathfrak h[\psi_n]}{\lVert\psi_n\rVert^2}\\
 &\le c^{-2}\left(
 C_n\frac{\mathfrak q_n[u_n]}{\lVert u_n\rVert^2}
 +\varepsilon_n\right)\longrightarrow0.
\end{align*}
The variational characterization of the bottom of a nonnegative
self-adjoint operator gives
\[
 \inf\sigma(H_{\mathrm{exc}})
 =\inf_{0\ne\psi\in
 \mathcal Q(H)\cap\mathcal H_{\mathrm{exc}}}
 \frac{\mathfrak h[\psi]}{\lVert\psi\rVert^2}=0.
\]
Since the spectrum is closed, zero belongs to it; since
\(\ker H_{\mathrm{exc}}=0\), zero is not an eigenvalue.  It is therefore a
continuous spectral point, and hence belongs to the essential spectrum under
the standard definition that removes isolated eigenvalues of finite
multiplicity.
\end{proof}

\begin{remark}[Why one vacuum vector is not enough]
If the vacuum is unique, choose a normalized vector
\(\Omega\) spanning \(\ker H\).  Then
\(P_0=|\Omega\rangle\langle\Omega|\).  Without uniqueness, merely requiring
\(J_nu_n\perp\Omega\) does not project away every zero-energy state.  For
example,
\[
 H=\operatorname{diag}(0,0,1),\qquad
 \Omega=(1,0,0),\qquad \psi=(0,1,0)
\]
has \(\psi\perp\Omega\) and zero Rayleigh quotient, although the genuine
excitation space \((\ker H)^\perp\) is the third coordinate and has gap one.
Thus \cref{cusp-interpretation-excitation-space} must use the full kernel
unless vacuum uniqueness is independently proved.
\end{remark}

\begin{auditrecord}[What the \(S^6\) calculation proves]
\label{cusp-interpretation-present-boundary}
The radial period matrix, scalar spectra, classical gauge Hessian, finite
cochain operator, magnetic Laplacian, and reducible \(SU(2)\) Jacobi operator
are explicit auxiliary geometric or classical operators.  They prove the
operator statements made in
\cref{cusp-interpretation-direct-integral-gapless,cusp-gauge-continuum-gap,cusp-gauge-finite-gap-limits,cusp-magnetic-oscillator-theorem,cusp-magnetic-negative-mode}.
They do not construct a common gauge-invariant positive-metric physical
Hilbert space \(\mathcal H_{\mathrm{phys}}\), a quantum Hamiltonian \(H\),
or maps \(J_n\) satisfying
\cref{cusp-interpretation-norm-noncollapse,cusp-interpretation-energy-transfer,cusp-interpretation-physical-smallness}
after continuum, infinite-volume, and renormalization limits.  Nor do the
period, Hodge, Leray, or direct-image calculations determine those data, as
\cref{cusp-interpretation-fixed-complex-counterexample} already disproves at
the level of geometric Laplacians.  Therefore the auxiliary gaplessness is
\Verified, while a conclusion about the four-dimensional quantum
Yang--Mills Hamiltonian is \Open.
\end{auditrecord}

\section{Integral topology without the cusp-collapse argument}
\label{sec:integral-leray}

This section follows the independent nearby-cycle/Leray route.  It does not
use the candidate's cellular collapse, its Mayer--Vietoris homology matrix,
or the asserted equality between those two constructions.  All coefficient
groups are integral.

Put
\[
 B=\mathbb P^1,\quad \Sigma=\{p_0,p_1,p_2\},\quad
 B^\circ=B\setminus\Sigma,\quad j:B^\circ\hookrightarrow B,
\]
and let \(F\) be a smooth fibre.  Retain
\[
 V=H^1(F)=\mathbb Z\langle\gamma,u,w,\delta\rangle,
 \quad q=uw+6\gamma\delta,
 \quad \operatorname{vol}=\gamma uw\delta.
\]

\subsection{Integral specialization at the cusp}

\begin{theorem}[nearby-cycle specialization]
\label{thm:p0-integral-specialization}
\Verified\ For the proper holomorphic cusp filling
\(f_0:N_0\to\Delta\) of \cref{thm:candidate-cusp-geometry}, with its
proved local normal-crossings equations and normalization combinatorics,
the specialization map of the central fibre \(W\),
\[
 \operatorname{sp}^r:H^r(W;\mathbb Z)\longrightarrow H^r(F;\mathbb Z)
\]
is injective and
\[
 \operatorname{im}\operatorname{sp}^r
 =H^r(F;\mathbb Z)^{T_0}
\tag{5.1}
\]
for \(0\le r\le4\).  In particular, the image is saturated and there is no
unrecorded finite-index defect at \(p_0\).
\end{theorem}

\begin{proof}
Let \(\widetilde\Delta^*\to\Delta^*\) be the universal cover, put
\(\widetilde N_0^*=N_0\times_\Delta\widetilde\Delta^*\), and let
\[
 k:\widetilde N_0^*\longrightarrow N_0,
 \qquad i:W\hookrightarrow N_0
\]
be the natural maps.  The analytic nearby-cycle complex used here is
\[
 \psi=\psi_{f_0}\mathbb Z
 :=i^*Rk_*k^*\mathbb Z_{N_0}\in D_c^b(W).
\tag{5.0}
\]
Dimca defines analytic nearby cycles in
\cite[Definition~4.2.1]{Dimca2004}; the Milnor-fibre description of their
stalks is \cite[Proposition~4.2.2]{Dimca2004}.  The proper-direct-image
theorem for nearby cycles is
\cite[Proposition~4.2.11]{Dimca2004}; applied to \(f_0\) and the identity
map of the disc, it gives
\[
 \mathbb H^n(W;\psi)\simeq H^n(F;\mathbb Z)
\]
and identifies the morphism induced by the unit
\(\mathbb Z_W\to\psi\) with specialization.  Proper base change itself is
\cite[Theorem~2.3.26]{Dimca2004}.  Dimca's coefficient hypothesis is a
commutative noetherian ring of finite global dimension, which
\(\mathbb Z\) satisfies.  Thus the general integral analytic formalism is
available here.  The candidate-specific existence and local description of
\(f_0\) are proved in \cref{thm:candidate-cusp-geometry,% 
prop:candidate-cusp-overlap}.  Kashiwara--Schapira
\cite[\S8.6]{KashiwaraSchapira1990}
corroborates the analytic formalism, while SGA~7, Expos\'e~XIII
\cite{SGA7II1973}, is its algebraic antecedent.

At a point where \(f=z_0\cdots z_k\), \(k=0,1,2\), the Milnor fibre
deformation retracts onto \((S^1)^k\): in polar coordinates the radial
variables form a contractible open simplex and the phases satisfy one linear
relation.  Hence
\[
 \mathcal H^0\psi=\mathbb Z_W,\qquad
 \mathcal H^1\psi=:\mathscr V,\qquad
 \mathcal H^2\psi=\mathbb Z_P\oplus\mathbb Z_Q,\qquad
 \mathcal H^b\psi=0\ (b\ge3),
\tag{5.2}
\]
where \(\mathscr V\) is supported on the three double curves
\(\bar C_1,\bar C_2,\bar C_3\simeq\mathbb P^1\).

Orient each vanishing circle by the displayed order of its two branches in
the normalization.  At \(P\), the branches of
\((\bar C_1,\bar C_2,\bar C_3)\) are respectively
\[
 (b_0,b_1),\qquad(b_0,b_2),\qquad(b_1,b_2).
\]
The three restrictions are therefore evaluations on
\[
 e_0-e_1,\qquad e_0-e_2,\qquad e_1-e_2,
\]
whose alternating sum is zero.  At \(Q\), the ordered branch pairs are
\[
 (c_2,c_1),\qquad(c_0,c_1),\qquad(c_0,c_2),
\]
so the evaluations are
\[
 e_2-e_1,\qquad e_0-e_1,\qquad e_0-e_2,
\]
and again the first minus the second plus the third is zero.  The global
normalization labels \(C_k\) and \(C_{k+3}\) distinguish the two branches
of each double curve and therefore fix these signs simultaneously.  Thus
\[
0\longrightarrow\mathscr V
\longrightarrow\bigoplus_{k=1}^3 i_{k*}\mathbb Z_{\bar C_k}
\xrightarrow{\delta_1}\mathbb Z_P\oplus\mathbb Z_Q
\longrightarrow0,
\]
\[
 \delta_1(n_1,n_2,n_3)=(s,s),\qquad s=n_1-n_2+n_3.
\tag{5.3}
\]
The image \(\mathbb Z(1,1)\) is primitive.  Taking cohomology of (5.3)
gives free groups with ranks
\[
 \operatorname{rk}H^a(W;\mathscr V)=(2,1,3,0)
 \quad(a=0,1,2,3).
\tag{5.4}
\]

The normalization resolution for the constant sheaf is
\[
0\longrightarrow\mathbb Z_W
\longrightarrow\nu_*\mathbb Z_{\mathrm{dP}_6}
\longrightarrow\bigoplus_{k=1}^3 i_{k*}\mathbb Z_{\bar C_k}
\xrightarrow{\delta_1}\mathbb Z_P\oplus\mathbb Z_Q
\longrightarrow0.
\tag{5.5}
\]
This is exact stalkwise.  At a smooth point it is the identity
\(\mathbb Z\to\mathbb Z\).  At a generic double point it is the augmented
complex
\[
 0\longrightarrow\mathbb Z\longrightarrow\mathbb Z^2
 \xrightarrow{(x_0,x_1)\mapsto x_0-x_1}\mathbb Z\longrightarrow0.
\]
At a triple point, in either of the branch orders just fixed, it is
\[
 0\longrightarrow\mathbb Z\longrightarrow\mathbb Z^3
 \xrightarrow{d_0}\mathbb Z^3
 \xrightarrow{d_1}\mathbb Z\longrightarrow0,
\]
where at \(P\)
\[
 d_0(x_0,x_1,x_2)=(x_0-x_1,x_0-x_2,x_1-x_2),
 \qquad d_1(n_1,n_2,n_3)=n_1-n_2+n_3,
\]
while at \(Q\)
\[
 d_0(x_0,x_1,x_2)=(x_2-x_1,x_0-x_1,x_0-x_2),
 \qquad d_1(n_1,n_2,n_3)=n_1-n_2+n_3.
\]
These augmented simplicial complexes are exact, which proves (5.5),
including its endpoint signs.

On \(H^2\), the only nontrivial map is
\[
 d:H^2(\mathrm{dP}_6;\mathbb Z)\longrightarrow\mathbb Z^3,
 \qquad
 d(L)=(L\cdot C_k-L\cdot C_{k+3})_{k=1}^3.
\]
Write \(\mathrm{dP}_6\) as the blow-up of \(\mathbb P^2\) at three
noncollinear points, with basis \(H,E_1,E_2,E_3\), and order its hexagon as
\[
 E_1, H-E_1-E_2, E_2, H-E_2-E_3, E_3, H-E_1-E_3.
\]
For \(L=aH-\sum b_iE_i\), direct intersection gives
\[
 d(L)=(s,-s,s),\qquad s=b_1+b_2+b_3-a.
\tag{5.6}
\]
Thus \(\operatorname{im}d=\mathbb Z(1,-1,1)\) is primitive,
\(\ker d\simeq\mathbb Z^3\), and
\(\operatorname{coker}d\simeq\mathbb Z^2\).  The resolution (5.5) has
width two and the cohomology of its smooth strata is concentrated in even
degrees.  The only possible higher differential would be
\(d_2:E_2^{0,q}\to E_2^{2,q-1}\), whose target is an odd row and is zero;
differentials with \(r\ge3\) leave the width-two range.  Therefore
\[
 H^a(W;\mathbb Z)=
 (\mathbb Z,\mathbb Z^2,\mathbb Z^4,\mathbb Z^2,\mathbb Z)
 \quad(a=0,1,2,3,4),
\tag{5.7}
\]
with no torsion.

The hypercohomology spectral sequence
\[
 E_2^{a,b}=H^a(W;\mathcal H^b\psi)
 \Longrightarrow H^{a+b}(F;\mathbb Z)
\tag{5.8}
\]
has the following ranks, all on free groups:
\[
\begin{array}{c|ccccc}
b=2&2&0&0&0&0\\
b=1&2&1&3&0&0\\
b=0&1&2&4&2&1\\ \hline
&a=0&a=1&a=2&a=3&a=4.
\end{array}
\]
Let the possible differential ranks be
\[
\alpha=\operatorname{rk}d_2^{0,1},\quad
\beta=\operatorname{rk}d_2^{1,1},\quad
\gamma'=\operatorname{rk}d_2^{2,1},\quad
\delta'=\operatorname{rk}d_2^{0,2},\quad
\epsilon=\operatorname{rk}d_3^{0,2}.
\]
Comparison with the torus Betti numbers \((1,4,6,4,1)\) in total degrees
one, four, and two gives
\[
 \alpha=0,\qquad \gamma'=0,\qquad
 \beta+\delta'+\epsilon=1.
\tag{5.9}
\]

We determine the remaining ambiguity geometrically, keeping the supported
specialization map explicit.  Choose a compact \(T_f\)-orbit
\(K=T_f\cdot w_0\simeq T^2\) in the smooth stratum
\(W^\circ\simeq(\mathbb C^*)^2\).  On a relatively compact neighbourhood
of \(K\), average a Riemannian metric over \(T_f\) and horizontally lift
\(\partial_{\operatorname{Re}t}\) and
\(\partial_{\operatorname{Im}t}\).  If \(\xi_1^h,\xi_2^h\) are those lifts,
put
\[
 \xi_t=(\operatorname{Re}t)\xi_1^h+
       (\operatorname{Im}t)\xi_2^h.
\]
Choose a relatively compact \(T_f\)-invariant open neighbourhood \(K'\)
with \(K\Subset K'\).  For some \(r_0>0\), the time-one flow
\[
 \Theta:K'\times\Delta_{r_0}\longrightarrow N_0,
 \qquad \Theta(w,t)=\Phi^{\xi_t}_1(w),
\tag{5.9a}
\]
is a \(T_f\)-equivariant open embedding over \(\Delta_{r_0}\), with
\(\Theta(w,0)=w\).  Indeed \(df_0(\xi_t)=t\), so
\(f_0(\Theta(w,t))=t\); its differential is invertible along
\(K'\times\{0\}\), and injectivity follows first by applying \(f_0\) and
then from uniqueness for the time-one flow.  Shrinking \(K'\) and \(r_0\)
gives the claimed embedding.

For \(0<r<r_0\), set
\[
 A_r=f_0^{-1}(\Delta_r),\qquad
 Z_r=\Theta(K\times\Delta_r),\qquad
 K_t=\Theta(K\times\{t\}).
\]
The subset \(Z_r\) is closed in \(A_r\).  Indeed, if a sequence in \(Z_r\)
converges in \(A_r\), compactness of \(K\) gives a convergent subsequence of
its \(K\)-coordinates; its base coordinates converge to the image under
\(f_0\), which still lies in \(\Delta_r\), and the embedding \(\Theta\)
then identifies the limit with a point of \(Z_r\).

For \(0<|t|<r\), functoriality of relative cohomology gives the two
commutative squares
\[
\begin{CD}
H^3(A_r,A_r\setminus Z_r) @>>> H^3(A_r)\\
@V{u_0\;(\simeq)}VV                  @VV{\rho_0}V\\
H^3(W,W\setminus K) @>>> H^3(W)
\end{CD}
\qquad
\begin{CD}
H^3(A_r,A_r\setminus Z_r) @>>> H^3(A_r)\\
@V{u_t\;(\simeq)}VV                  @VV{\rho_t}V\\
H^3(F_t,F_t\setminus K_t) @>>> H^3(F_t).
\end{CD}
\tag{5.9b}
\]
The two left vertical arrows are isomorphisms.  By excision they are the
restrictions to the two slices of
\[
 H^3\bigl(K'\times\Delta_r,
          (K'\setminus K)\times\Delta_r\bigr),
\]
and homotopy invariance makes both restrictions isomorphisms.  Taking the
direct limit over shrinking \(r\), the maps \(\rho_0\) induce the
proper-base-change isomorphism
\[
 \varinjlim_{r\downarrow0}H^3(A_r;\mathbb Z)
 \xrightarrow{\ \sim\ }H^3(W;\mathbb Z).
\]
Choose the points \(t\) on one radial ray and identify the \(F_t\) with the
fixed marked fibre \(F\) by parallel transport.  Restriction to \(F_t\),
followed by this identification, is compatible under shrinking \(r\).
The nearby-cycle specialization is precisely the composite
\[
 \operatorname{sp}^3:
 H^3(W;\mathbb Z)
 \xleftarrow[\ \sim\ ]{\rho_0}
 \varinjlim_{r\downarrow0}H^3(A_r;\mathbb Z)
 \xrightarrow{\rho_{\rm near}}H^3(F;\mathbb Z);
\]
this is the neighbourhood realization of the unit
\(\mathbb Z_W\to\psi_{f_0}\mathbb Z\) in the comparison cited above.
Taking the direct limit in the two functorial squares therefore proves,
rather than assumes, that
\[
 H^3(W,W\setminus K)\longrightarrow H^3(W)
 \xrightarrow{\operatorname{sp}^3}H^3(F_t)
\tag{5.9c}
\]
is the transported supported class followed by the Gysin map for
\(K_t\hookrightarrow F_t\).

Under the marking of the torus fibre,
\[
 \Lambda_{\rm tor}=\mathbb Z\widehat w\oplus
                    \mathbb Z\widehat\delta
\]
is a saturated direct summand of \(\Lambda\), and \(K_t\) is a translate of
\((\Lambda_{\rm tor}\otimes\mathbb R)/\Lambda_{\rm tor}\).  Choose a
complementary summand \(\Lambda'\), and write
\(F_t\simeq T_f\times T_B\) as oriented real tori, where
\(T_B=(\Lambda'\otimes\mathbb R)/\Lambda'\), with the product orientation
chosen to agree with that of \(F_t\).  The real rank-two normal bundle of
the translate \(K_t\) is trivial and oriented.  The Thom isomorphism gives
\[
H^3(F_t,F_t\setminus K_t)\simeq H^1(T_f),
\]
and its map to absolute cohomology is
\[
 \iota_!(a)=a\otimes[T_B]^\vee
 \in H^1(T_f)\otimes H^2(T_B)\subset H^3(F_t).
\tag{5.9d}
\]
This is an injective primitive rank-two K\"unneth direct summand.  On the
central smooth stratum,
\(W^\circ\simeq(\mathbb C^*)^2\simeq T_f\times\mathbb R^2\) as an
oriented real manifold; hence the real rank-two normal bundle of \(K\) is
also trivial and oriented.  The same Thom calculation gives
\(H^3(W,W\setminus K)\simeq H^1(K)\simeq\mathbb Z^2\).
Hence (5.9c) has primitive rank-two image.  Since
\(\operatorname{rk}H^3(W)=2\), specialization in degree three is injective
and has precisely that primitive image.  In (5.9), this forces
\(\beta=\epsilon=0\), and consequently \(\delta'=1\).

We make the integral saturation argument explicit.  The unit
\(\mathbb Z_W\to\psi\) is a morphism from the complex concentrated in
cohomological degree zero to the Postnikov filtration of \(\psi\).  It
therefore induces a morphism of hypercohomology spectral sequences from the
single row
\[
 H^a(W;\mathbb Z)=E_2^{a,0}(\mathbb Z_W)
\]
to (5.8).  On the \(b=0\) row the induced \(E_2\)-map is the identity,
because \(\mathcal H^0\psi=\mathbb Z_W\).  Consequently, in total degree
\(r\), specialization factors through the last filtration step
\[
 H^r(W;\mathbb Z)\xrightarrow{\ \sim\ }
 E_\infty^{r,0}=F^rH^r(F;\mathbb Z)
 \lhook\joinrel\longrightarrow H^r(F;\mathbb Z)
 \tag{5.9e}
\]
whenever no differential
\(d_k:E_k^{r-k,k-1}\to E_k^{r,0}\) enters on any page.  For \(r=1\) this
is automatic; for \(r=2\) the only possible entering differential is
\(d_2^{0,1}\), and \(\alpha=0\).  Thus (5.9e) is an isomorphism onto the
specialization image in degrees one and two.  The quotient by \(F^r\) has
the finite filtration whose graded groups are the remaining
\(E_\infty^{a,r-a}\), \(a<r\).  In degree one these are kernels of maps of
free abelian groups.  In degree two they are
\(E_\infty^{1,1}=\mathbb Z\) and
\(E_\infty^{0,2}=\ker d_2^{0,2}\); here the possible subsequent map is
\(d_3^{0,2}\), and it is zero because its already determined rank
\(\epsilon\) is zero.  Both groups are therefore free.  Successive extensions of
free abelian groups split, so \(H^r(F;\mathbb Z)/F^r\) is free for
\(r=1,2\).  Hence the specialization images in those degrees are
primitive.

In degree three, (5.9c)--(5.9d) give a primitive rank-two lattice contained
in \(\operatorname{im}\operatorname{sp}^3\).  Since
\(H^3(W;\mathbb Z)\) itself has rank two, the specialization image has the
same rational span; a primitive lattice has no proper finite-index
overlattice in that span inside \(H^3(F;\mathbb Z)\).  Therefore the two
images coincide, proving both injectivity and primitivity in degree three.
In degree four, \(E_\infty^{4,0}=H^4(F;\mathbb Z)\), so specialization is an
isomorphism.  Thus specialization is injective with primitive image in every
degree.  Monodromy invariance gives
\[
 \operatorname{im}\operatorname{sp}^r
 \subseteq H^r(F;\mathbb Z)^{T_0}.
\]
By \cref{thm:cusp-lattices}, the invariant ranks are
\((1,2,4,2,1)\), exactly those in (5.7), and each invariant lattice is the
kernel of an integral endomorphism, hence primitive.  Two primitive
sublattices of the same rational subspace coincide: both equal the
intersection of that subspace with the ambient lattice.  This proves (5.1).
\end{proof}

\subsection{The global \texorpdfstring{\(E_2\)}{E2}-page}

The Leray spectral sequence used in this subsection is
\cite[Corollary~2.3.4]{Dimca2004}.
Let \(R^r=R^rf_*\mathbb Z\) and
\(\mathcal L^r=j^*R^r\).  The finite-fibre lattices have already been proved
as abstract lattice statements in \cref{prop:finite-indices}.  Their
identification with the stalks of the proposed family requires the explicit
radial-stalk compatibility below.  This compatibility is treated separately
in \cref{lem:finite-filling-stalk-map} and is not silently absorbed here.

\begin{proposition}[all Leray terms used below]
\label{prop:leray-e2}
\Verified\ The specialization theorem
\cref{thm:p0-integral-specialization}, the finite fillings of
\cref{prop:candidate-finite-fillings}, and
\cref{lem:finite-filling-stalk-map} give, for \(j=1,2\) and
\(0\le r\le4\), the finite-filling compatibility
\[
\tag{FF}
 N_j\ \text{strongly deformation retracts onto }S_j,\qquad
 H^r(S_j;\mathbb Z)\longrightarrow H^r(F;\mathbb Z)^{T_j}
 \ \text{is }\pi_j^*
\]
under the marked punctured-fibre identification.  Consequently,
\[
H^0(B,R^1)=\mathbb Z(12\gamma),\quad
H^0(B,R^2)=\mathbb Z(2q),\quad
H^0(B,R^3)=\mathbb Z(2\gamma uw),
\tag{5.10}
\]
\[
 H^1(B,R^1)=H^1(B,R^2)=0,
\tag{5.11}
\]
and, with \(\omega\) a generator of \(H^2(B;\mathbb Z)\),
\[
H^2(B,R^1)=\mathbb Z\omega[\delta],\qquad
H^2(B,R^2)=\mathbb Z\omega[\gamma\delta],\qquad
H^2(B,R^4)=\mathbb Z\omega[\operatorname{vol}].
\tag{5.12}
\]
These identifications are compatible with exterior products in \(V\).
\end{proposition}

\begin{proof}
Under the stated hypotheses, proper base change identifies the stalk of
\(R^r\) at \(p_0\) with \(H^r(W;\mathbb Z)\), while
\((j_*\mathcal L^r)_{p_0}=H^r(F;\mathbb Z)^{T_0}\).  The stalk map is
specialization and is an isomorphism by
\cref{thm:p0-integral-specialization}; hence its cokernel at \(p_0\) is zero.
At \(p_j\), condition \textup{(FF)} identifies the stalk with
\(H^r(S_j;\mathbb Z)\), the target with
\(H^r(F;\mathbb Z)^{T_j}\), and the stalk map with \(\pi_j^*\).  The
finite-fibre calculation \cref{prop:finite-indices} makes each of these maps
injective.  Away from the three special points the map is the identity of
the local system.  Stalkwise kernels and cokernels therefore give the exact
sequence
\[
0\longrightarrow R^r\longrightarrow j_*\mathcal L^r
\longrightarrow Q^r\longrightarrow0,
\tag{5.13}
\]
with
\[
Q^1=(\mathbb Z/3)_{p_1}\oplus(\mathbb Z/4)_{p_2},
\quad Q^2=(\mathbb Z/2)_{p_2},
\quad Q^3=(\mathbb Z/2)_{p_2},
\quad Q^4=(\mathbb Z/3)_{p_1}\oplus(\mathbb Z/4)_{p_2}.
\]
The degree-one quotients have orders three and four and are generated by the
class of \(\gamma\).  In degree two the \(p_1\)-quotient is zero, while
\[
 (\textstyle\bigwedge^2V)^{T_2}/
 \{a\gamma\eta_2+bq:a\equiv b\pmod2\}\simeq\mathbb Z/2
\]
is generated by the class of \(q\).  In degree three the \(p_1\)-quotient
is zero and the \(p_2\)-quotient has order two, generated by the class of
\(\gamma uw\).  This last identification uses the actual twist
\(v_2=-\varepsilon'\), for which the source's admissibility quantity
\(\eta_2(v_2)\) is zero and hence divisible by four.  In degree four,
pullback multiplies the orientation class by the covering
degree \(m_j\), giving the displayed \(Q^4\).  Hence the skyscraper sheaves
have exactly the asserted stalks; no claim for an arbitrary twist is being
imported.
By \cref{thm:coinvariants},
\[
V^G=\mathbb Z\gamma,\quad
(\textstyle\bigwedge^2V)^G=\mathbb Zq,\quad
(\textstyle\bigwedge^3V)^G=\mathbb Z\gamma uw,
\tag{5.14}
\]
and
\[
V_G=\mathbb Z[\delta],\qquad
(\textstyle\bigwedge^2V)_G=\mathbb Z[\gamma\delta],\qquad
[q]=12[\gamma\delta].
\tag{5.15}
\]
Since \(\det T_1=\det T_2=1\), the action on
\(\bigwedge^4V=\mathbb Z\operatorname{vol}\) is trivial and
\[
 (\textstyle\bigwedge^4V)_G=\mathbb Z[\operatorname{vol}].
\tag{5.15a}
\]
The maps from the invariant groups in (5.14) onto the corresponding
skyscraper fibres are surjective, with kernels
\(12\mathbb Z\gamma\), \(2\mathbb Zq\), and
\(2\mathbb Z\gamma uw\).  This proves (5.10).

Because each \(Q^r\) is a skyscraper, it has no positive cohomology; the
long exact sequence of (5.13) therefore identifies
\(H^2(B,R^r)\) with \(H^2(B,j_*\mathcal L^r)\).
For any free local system \(\mathcal L\) on \(B^\circ\), with stalk \(M\),
the quotient \(j_*\mathcal L/j_!\mathcal L\) is a skyscraper.  Hence
\[
 H^2(B,j_*\mathcal L)\simeq H_c^2(B^\circ,\mathcal L)\simeq
 H_0(B^\circ,\mathcal L)=M_G,
\tag{5.16}
\]
where the middle isomorphism is Poincar\'e duality with local coefficients
\cite[\S3.H]{Hatcher2002}; see also
\href{https://doi.org/10.1007/978-1-4612-0647-7}{G.~E.~Bredon,
\emph{Sheaf Theory}, 2nd ed., Chapter~VI}.
Cap product with the locally finite fundamental class makes (5.16) natural
for coefficient pairings, proving (5.12) and the product assertion.

It remains to prove (5.11), for which we spell out the parabolic-cohomology
quotient.  The low-degree Leray sequence for \(j:B^\circ\hookrightarrow B\)
contains
\[
 0\longrightarrow H^1(B,j_*\mathcal L)
 \longrightarrow H^1(B^\circ,\mathcal L)
 \longrightarrow H^0(B,R^1j_*\mathcal L).
\]
If \(I_i=\langle t_i\rangle\) is the local inertia group at \(p_i\), then
\[
 (R^1j_*\mathcal L)_{p_i}
 =H^1(I_i,M)=M/(T_i-I)M.
\]
Consequently the last arrow sends the class of a cocycle \(c\) to the three
classes of \(c(t_i)\), and \(H^1(B,j_*\mathcal L)\) is exactly the kernel
defined by the three parabolic conditions below.  Use left cocycles, so
\(c(gh)=c(g)+g c(h)\).  A cocycle on
\(\pi_1(B^\circ)=F(t_1,t_2)\) is determined by
\(x=c(t_1)\) and \(y=c(t_2)\).  Extension across the three punctures is the
parabolic condition
\[
 x\in(T_1-I)M,\qquad y\in(T_2-I)M,\qquad
 c(t_0)\in(T_0-I)M,
 \qquad t_0=(t_1t_2)^{-1}.
\]
Write \(y=(T_2-I)m\) and subtract the coboundary of \(m\); this normalizes
\(y=0\).  A further coboundary preserves the normalization precisely when
its defining element lies in \(M^{T_2}\), and then changes \(x\) by an
element of \((T_1-I)M^{T_2}\).  With \(y=0\),
\(c(t_0)=-T_0x\).  Since \(T_0\) preserves
\((T_0-I)M\), the last parabolic condition is equivalent to
\(x\in(T_0-I)M\).  Therefore
\[
 H^1(B,j_*\mathcal L)\simeq
 \frac{(T_1-I)M\cap(T_0-I)M}{(T_1-I)M^{T_2}}.
\tag{5.17}
\]
For \(M=V\),
\[
(T_1-I)V=\langle2\gamma+u+w,2\gamma-u\rangle,\qquad
(T_0-I)V=\langle\gamma,u\rangle,
\]
and their intersection is
\(\mathbb Z(2\gamma-u)=(T_1-I)\mathbb Z\eta_2\).  For
\(M=\bigwedge^2V\),
\[
\begin{aligned}
(\textstyle\bigwedge^2T_1-I)\textstyle\bigwedge^2V
=\langle&\gamma u,\gamma w,2u\delta+w\delta,\\
&uw+u\delta-w\delta-6\gamma\delta\rangle,
\end{aligned}
\]
whereas
\[
(\textstyle\bigwedge^2T_0-I)\textstyle\bigwedge^2V
=\langle\gamma u,\gamma w+u\delta\rangle.
\]
Their intersection is \(\mathbb Z\gamma u\), equal to
\((\bigwedge^2T_1-I)\mathbb Z(\gamma\eta_2)\).  Thus both groups in
(5.17) vanish.  The maps
\(H^0(B,j_*\mathcal L^r)\to H^0(B,Q^r)\) were proved surjective above for
\(r=1,2\), and skyscraper sheaves have no positive cohomology.  The long
exact sequence of (5.13) therefore gives (5.11) for \(R^1,R^2\).
\end{proof}

\subsection{The three transgressions}

\begin{lemma}[the primitive first transgression]
\label{lem:first-transgression}
\Verified\ For the constructed map and the based meridians of
\cref{prop:geometric-pi1}, one may orient \(\omega\) so that
\[
 d_2^{0,1}(12\gamma)=\omega.
\tag{5.18}
\]
\end{lemma}

\begin{proof}
Choose disjoint discs \(D_i\) about the special points and a pair-of-pants
neighbourhood \(V_0\) of their complement.  The class
\(\theta=12\gamma\) is a section of \(R^1f_*\mathbb Z\) over each \(D_i\):
at \(p_0\) this is \cref{thm:p0-integral-specialization}, and at \(p_j\)
it is the degree-one finite-fibre lattice.  It is also a section over
\(V_0\), because it is globally monodromy invariant.  For every member
\(U\in\{D_0,D_1,D_2,V_0\}\), the local Leray five-term sequence contains
\[
 H^1(f^{-1}(U);\mathbb Z)\longrightarrow
 H^0(U,R^1f_*\mathbb Z)\longrightarrow
 H^2(U,R^0f_*\mathbb Z)=H^2(U;\mathbb Z)=0.
\]
Thus the section \(\theta\) has a cohomological lift
\(\Theta_U\in H^1(f^{-1}(U);\mathbb Z)\) on every member of the cover.
On an overlap \(D_i\cap V_0\), the two lifts have the same image in
\(H^0(R^1f_*\mathbb Z)\); local five-term exactness therefore identifies
their difference with a class in
\(H^1(D_i\cap V_0,R^0f_*\mathbb Z)=H^1(D_i\cap V_0;\mathbb Z)\).
These three differences form the \v{C}ech connecting cocycle used below.

Use clockwise meridians.  The affine deck relation at \(p_j\) is
\[
 \widehat\rho_j^{m_j}=t_{v_j}.
\tag{5.19}
\]
Put \(\Theta_j=\Theta_{D_j}\).  Since \(\Theta_j\) is integral and restricts
to \(\theta\) on the fibre, (5.19) gives
\(m_j\Theta_j(\widehat\rho_j)=\Theta_j(t_{v_j})=\theta(v_j)\).  Therefore
\[
 \Theta_1(\widehat\rho_1)=12\gamma(v_1)/3=4,
 \qquad
 \Theta_2(\widehat\rho_2)=12\gamma(v_2)/4=-3.
\tag{5.20}
\]
By \cref{prop:geometric-pi1}, the toric meridian
\(\widetilde a_0\) is the
boundary of the holomorphic disc
\(\{(1,1,t):|t|\le\epsilon\}\).  It is therefore null-homotopic and the
local lift evaluates to zero.

Write \(\alpha_i=\Theta_{V_0}(\widehat\rho_i)\).  By the exact based cusp
relation (5.23a), the three annular discrepancies have
meridian values
\[
 4-\alpha_1,\qquad -3-\alpha_2,\qquad \alpha_1+\alpha_2,
\tag{5.21}
\]
whose sum is one.  Mayer--Vietoris on the base identifies
\[
H^1\!\left(\coprod_i(D_i\cap V_0)\right)/\operatorname{im}H^1(V_0)
\simeq H^2(B;\mathbb Z).
\]
Under meridian evaluation the numerator is \(\mathbb Z^3\), the image is
the sum-zero lattice, and the connecting map is the signed sum of the three
coordinates.

We now identify this connecting class with the Leray differential rather
than taking that identification as an extra hypothesis.  Put
\(K=Rf_*\mathbb Z_X\).  The canonical truncation triangle
\cite[Tag~08J5]{StacksProject} is
\[
 R^0f_*\mathbb Z_X\longrightarrow\tau_{\le1}K
 \longrightarrow R^1f_*\mathbb Z_X[-1]
 \xrightarrow{\ \kappa\ }R^0f_*\mathbb Z_X[1].
\tag{5.21a}
\]
After applying \(R\Gamma(B,-)\), the connecting homomorphism
\[
 H^0(B,R^1f_*\mathbb Z)\xrightarrow{\,H^0(\kappa)\,}
 H^2(B,R^0f_*\mathbb Z)
\tag{5.21b}
\]
is precisely \(d_2^{0,1}\) in the hypercohomology, hence Leray, spectral
sequence; this is the canonical-truncation construction of that spectral
sequence \cite[Tag~0FLJ]{StacksProject}.  On the chosen disc--pair-of-pants
cover, functorial Mayer--Vietoris for the triangle
\cite[Tag~0F16]{StacksProject} computes \(H^0(\kappa)\) by choosing the local
degree-one lifts \(\Theta_U\) and taking their differences on the annular
overlaps.  Those differences are exactly the cocycle (5.21).  Thus the
signed sum just computed is \(d_2^{0,1}(\theta)\).
Equation (5.21) therefore represents a primitive generator; orient \(\omega\)
to obtain (5.18).
\end{proof}

\begin{lemma}[the higher transgressions]
\label{lem:higher-transgressions}
\Verified\ In the integral Leray spectral sequence of the constructed map,
\[
 d_2^{0,2}(2q)=\omega[\delta],
 \qquad
 d_2^{0,3}(2\gamma uw)=\omega[\gamma\delta].
\tag{5.22}
\]
\end{lemma}

\begin{proof}
Use the multiplicative Leray spectral sequence obtained from a filtered
multiplicative cochain resolution.  Its \(E_2\)-product is the natural cup
product
\[
 H^p(B,R^q)\otimes H^{p'}(B,R^{q'})
 \longrightarrow H^{p+p'}(B,R^{q+q'}),
\]
and every \(d_r\) is a graded derivation.  This is the standard filtered-DGA
construction; compare the filtered-complex spectral sequence
\cite[Tags~012K and~0FLJ]{StacksProject} with the derived cup product
\cite[\S20.31]{StacksProject}.  The generator identifications and their
products are natural by the last sentence of
\cref{prop:leray-e2}.
Put \(\theta=12\gamma\), \(a=2q\), and \(b=2\gamma uw\).  By
\cref{prop:leray-e2}, write
\[
 d_2a=r\omega[\delta],\qquad d_2b=s\omega[\gamma\delta].
\]
In the \(E_2\)-ring,
\[
 \theta a=12b,\qquad ab=0,\qquad [q]=12[\gamma\delta].
\]
The derivation rule and \(|\theta|=1\) give
\[
12d_2b=d_2(\theta a)=d_2\theta\cdot a-\theta\cdot d_2a,
\]
hence \(s=2-r\).  Since \(|a|=2\),
\(0=d_2(ab)=d_2a\cdot b+a\cdot d_2b\).  The exterior products are
\[
 \delta\wedge2\gamma uw=-2\operatorname{vol},
 \qquad
 2q\wedge\gamma\delta=2\operatorname{vol}.
\]
Therefore \(-2r+2s=0\), so \(r=s\).  Together with \(s=2-r\), this gives
\(r=s=1\), proving (5.22).
\end{proof}

\begin{theorem}[integral cohomology in low degrees]
\label{thm:h2-vanishing}
\Verified\ For the constructed holomorphic map,
\[
 H^1(X;\mathbb Z)=H^2(X;\mathbb Z)=H^3(X;\mathbb Z)=0,
 \qquad H_2(X;\mathbb Z)=0.
\tag{5.23}
\]
\end{theorem}

\begin{proof}
The Leray terms in total degree two are
\[
E_2^{2,0}=\mathbb Z\omega,\qquad E_2^{1,1}=0,\qquad
E_2^{0,2}=\mathbb Z(2q).
\]
By \cref{lem:first-transgression,lem:higher-transgressions}, the incoming
map to the first term and the outgoing map from the last term are
isomorphisms.  Thus all three \(E_\infty\)-pieces vanish and \(H^2(X)=0\).
The first map has zero kernel, giving \(H^1(X)=0\).  In total degree three,
the quotient of \(E_2^{2,1}\) by the image of \(E_2^{0,2}\) is zero,
\(E_2^{1,2}=0\), and the outgoing map from \(E_2^{0,3}\) is injective by
(5.22).  Hence \(H^3(X)=0\).  No later differential can meet these terms
because the base has cohomological dimension two.

The universal coefficient sequences show both
\(\operatorname{Hom}(H_2(X),\mathbb Z)=0\) and
\(\operatorname{Ext}(H_2(X),\mathbb Z)=0\).  Since \(H_2(X)\) is finitely
generated, it has neither a free nor a torsion summand, so it is zero.
\end{proof}

\subsection{Based van Kampen and sphere recognition}

\begin{proposition}[the geometric fundamental-group presentation]
\label{prop:geometric-pi1}
\Verified\ For the punctured-overlap gluings of
\cref{prop:candidate-cusp-overlap,prop:candidate-finite-fillings}, the based
van Kampen presentation is
\[
 \pi_1(X)\simeq
 \langle c,x,y\mid c\text{ central},\ xy=1,\ x^3=c,\ y^4=c^{-1}\rangle
 \simeq1.
\tag{5.24}
\]
\end{proposition}

\begin{proof}
Choose \(b_*\) in a pair-of-pants core of \(B^\circ\), put \(x_*\) on the
holomorphic zero section above it, and choose three nonintersecting base
whiskers from \(b_*\) to clockwise boundary circles.  Lift all whiskers by
the zero section.  With the ordering compatible with \(g_1g_2g_0=1\), the
resulting based loops obey
\[
 \widehat\rho_1\widehat\rho_2\widetilde a_0=1.
\tag{5.23a}
\]
The period-family action is linear and hence preserves the zero section.
It splits the homotopy sequence over the pair of pants, giving
\[
 \pi_1(f^{-1}(B^\circ))=
 \Lambda\rtimes F(\widehat\rho_1,\widehat\rho_2),
\]
with conjugation actions \(A_1,A_2\).  The clockwise meridian at
\(p_j\) lifts to the affine generator in (2.9), with its exact
sign and based translation.  To verify this in the actual
real marking, put
\[
 \alpha(s)=\frac{\log s_j}{2\pi i},\qquad
 J_s=\Pi(s)^{-1}i\Pi(s),\qquad
 a_j(s)=\operatorname{Re}\alpha(s)\,v_j+
        \operatorname{Im}\alpha(s)\,J_s v_j.
\]
The overlap translation by
\(\alpha(s)\Pi(s)v_j\) is translation by \(a_j(s)\) in the real
marking.  Thus the inverse image of the period-family zero section
is \(x(s)=-a_j(s)\).

The covariance of \(J_s\) is a consequence of the proved period
identity (2.8e), not an additional assumption.  That identity is
\[
 R_{g_j}(s)\Pi(s)=\Pi(g_js)A_j.
\]
The map \(R_{g_j}(s)\) is complex linear.  Hence
\[
 \begin{split}
 J_{g_js}A_j
 &=\Pi(g_js)^{-1}i\,\Pi(g_js)A_j\\
 &=\Pi(g_js)^{-1}i\,R_{g_j}(s)\Pi(s)\\
 &=\Pi(g_js)^{-1}R_{g_j}(s)i\,\Pi(s)
 =A_jJ_s.
 \end{split}
\]
The original vectors are
\[
 v_1=(1,2,-4,0)^{\mathsf t},\qquad
 v_2=(-1,-3,3,0)^{\mathsf t},
\]
in the basis
\((\widehat\gamma,\widehat u,\widehat w,\widehat\delta)\).
Multiplication by the printed matrices gives
\[
 A_1v_1=(1,\,6-4,\,-6-2+4,\,-2+2)^{\mathsf t}
        =(1,2,-4,0)^{\mathsf t},
\]
\[
 A_2v_2=(-1,\,-3,\,6-3,\,-3+3)^{\mathsf t}
        =(-1,-3,3,0)^{\mathsf t}.
\]
Thus \(A_jv_j=v_j\) exactly.  Clockwise continuation by \(g_j\)
changes the chosen logarithm to
\(\alpha(g_js)=\alpha(s)-1/m_j\).  Using this real change, the
unchanged imaginary part, and \(J_{g_js}A_j=A_jJ_s\), one obtains
\[
 a_j(g_js)=A_ja_j(s)-v_j/m_j.
\]
Therefore the affine generator sends the initial lift of the
zero section to its actual continued endpoint:
\[
 A_j(-a_j(s))+v_j/m_j=-a_j(g_js).
\]
Changing the logarithm branch by an integer changes the initial
lift by an integer multiple of \(v_j\).  Such a translation
commutes with the affine generator because \(A_jv_j=v_j\), so the
endpoint identification does not introduce a translated conjugate.
The fixed zero-section whiskers retain the based generator.
After \(m_j\) clockwise iterations, the coordinate returns to the
same point and \(\alpha\) has decreased by one.  Thus the endpoint
of \(-a_j\) has increased by \(v_j\), in agreement with (2.8f).
Consequently
\[
 \widehat\rho_j^{m_j}=t_{v_j}.
 \tag{5.23b}
\]
The radial model of \cref{lem:finite-filling-stalk-map} shows that attaching
\(N_j\) imposes exactly this relation.  Since the global and local whiskers
are the same zero-section path under the overlap, conjugating back to
\(x_*\) adds no fibre drift.

For completeness, the cusp covering has no hidden fundamental group.  A loop
in the infinite toric variety has compact image and therefore lies in the
open \(X_{\Sigma_{\rm tor}}\) associated with a finite subfan
\(\Sigma_{\rm tor}\).  Enlarge that subfan by the rays through
\((0,1),(e_1,1),(e_2,1)\).  Write
\[
 N_{\Sigma_{\rm tor}}
 :=\bigl\langle|\Sigma_{\rm tor}|\cap N'\bigr\rangle_{\mathbb Z}.
\]
The three displayed primitive generators span \(N'\), so the classical
formula
\(\pi_1(X_{\Sigma_{\rm tor}})=N'/N_{\Sigma_{\rm tor}}\)
\cite[Theorem~12.1.10]{CoxLittleSchenck2011} makes the resulting finite toric
open simply connected; the original loop therefore has a null-homotopy in
that open.  Thus \(Y\) is simply connected.  The one-parameter subgroup in
the third torus coordinate acts by \(t\mapsto r t\), \(r>0\).  Given the
compact image of a null-homotopy, choose \(r\) so small that its translate
lies in \(|t|<\epsilon\); the path \(r_s=(1-s)+sr\) gives a homotopy from the
original loop to its translate inside \(Y_\epsilon\).  Hence
\(Y_\epsilon\) is also simply connected.  The verified free quotient
\(N_0=Y_\epsilon/\bar\Lambda\) therefore has universal cover
\(Y_\epsilon\), so
\(\pi_1(N_0)=\bar\Lambda\).  Equation (2.8iii) identifies a fibre-lattice
loop \(\lambda\in\Lambda\) with the deck transformation
\(\bar\lambda\); hence the inclusion of a nearby fibre is exactly
\[
 \Lambda\longrightarrow\bar\Lambda,\qquad
 \lambda\longmapsto\bar\lambda,
\tag{5.23c}
\]
with kernel
\(\Lambda_{\rm tor}=\langle\widehat w,\widehat\delta\rangle\).
On the zero section the same overlap sends \((z,0)\) to
\((1,1,t_c)\).  In the toric chart of the triangle
\(\{0,e_1,e_2\}\), whose coordinates may be written
\((t_c/(x_1x_2),x_1,x_2)\), this is the holomorphic disc
\(t_c\mapsto(t_c,1,1)\).  Thus the cusp meridian
\(\widetilde a_0\) is null-homotopic, with no fibre translation.  Combining
this fact with (5.23a), the cusp attachment imposes precisely
\(\widehat\rho_1\widehat\rho_2=1\).

Van Kampen now gives
\[
\left\langle\Lambda,\widehat\rho_1,\widehat\rho_2\ \middle|\
\widehat\rho_j\lambda\widehat\rho_j^{-1}=A_j\lambda,
\widehat\rho_1^3=t_{v_1},\widehat\rho_2^4=t_{v_2},
\widehat\rho_1\widehat\rho_2=1,
\Lambda_{\rm tor}=1\right\rangle.
\]
The normal closure of \(\Lambda_{\rm tor}\) under \(A_1,A_2\) is
\(\ker\gamma=\langle\widehat u,\widehat w,\widehat\delta\rangle\); for
example \(A_1\widehat w=\widehat u-\widehat w\).  The image of \(\Lambda\)
is therefore generated by \(c=[\widehat\gamma]\), which is central because
\(\gamma\) is invariant.  Since \(\gamma(v_1)=1\) and
\(\gamma(v_2)=-1\), the presentation becomes (5.24).  It is the case
\((\ell_0,\ell_1,\ell_2)=(0,1,-1)\) of
\cref{prop:terminal-group}, where \(p=-1\); hence the group is trivial.
\end{proof}

\begin{corollary}[topology and the third Chern number]
\label{cor:candidate-sphere}
\Verified\ The underlying smooth manifold of the constructed threefold is
diffeomorphic to the standard six-sphere, and
\[
 \langle c_3(T_X),[X]\rangle=2.
\tag{5.25}
\]
\end{corollary}

\begin{proof}
Combine \cref{thm:h2-vanishing,prop:geometric-pi1}.  The universal
coefficient theorem and Poincar\'e duality give the integral homology of
\(S^6\), exactly as in \cref{lem:cohomological-closure}.  The complete
recognition argument of \cref{thm:sphere-recognition} then gives
\(X\cong_{\rm diff}S^6\).  Its Euler characteristic is two.  On a complex
threefold the top Chern class of the complex tangent bundle is the Euler
class of the underlying oriented real tangent bundle, proving (5.25).
\end{proof}

\paragraph{Independent role of the collapse route.}
The proof above is independent of the cellular route.  The candidate's cusp
collapse, its integral homology map, and its equality with nearby-cycle
specialization are now independently verified in
\cref{prop:verified-cellular-collapse,prop:verified-collapse-homology,%
thm:collapse-equals-specialization}.  The actual finite-band and cusp
Mayer--Vietoris maps, their primitive blocks, and their commuting comparison
with the differentials proved in this section are now
\cref{prop:actual-finite-band-maps,prop:cusp-mv-blocks,%
thm:mv-leray-comparison}.  This supplies a second complete integral topology
certificate; it remains logically independent of the sphere-recognition
conclusion already obtained in this section.

\section{The direct algebraic-reduction proof}
\label{sec:algebraic-dimension}

The calculation in this section is independent of the integral topology of
\(X\).  It uses the period family of
\cref{prop:period-input-audit} and the proper connected fibration of
\cref{thm:candidate-global-geometry}; it does not use \(H^2(X;\mathbb Z)\),
the sphere-recognition argument, or any theorem of
Campana--Demailly--Peternell.

\begin{lemma}[algebraic reduction of a two-dimensional torus]
\label{lem:torus-algebraic-reduction-audit}
Let \(T\) be a complex torus of dimension two.  Suppose
\[
 \operatorname{NS}(T)=\mathbb Z\eta,\qquad \eta^2\ne0,
\]
and the Appell--Humbert Hermitian form \(H_\eta\) is nondegenerate and
indefinite.  Then \(a(T)=0\).
\end{lemma}

\begin{proof}
The algebraic reduction of a compact complex torus is a quotient
\(q:T\to A=T/K\) by a connected complex subtorus, with \(A\) an abelian
variety and \(\dim A=a(T)\)
\cite[Chapter~2, \S6]{BirkenhakeLange1999}; see also
\cite[\S8, p.~2437]{PareschiPopaSchnell2017}.

Assume \(a(T)>0\).  If \(a(T)=2\), then \(T\) is an abelian surface.  An
ample line bundle has first Chern class \(n\eta\), \(n\ne0\), and its
Appell--Humbert form \(nH_\eta\) is positive definite.  This contradicts
the indefiniteness of \(H_\eta\).

If \(a(T)=1\), then \(A\) is an elliptic curve.  For a point \(p\in A\),
put
\[
 \varphi=q^*c_1(\mathcal O_A(p))\in\operatorname{NS}(T).
\]
The class is nonzero: for any translation-invariant Kahler form \(\omega\)
on \(T\), the projection formula gives
\(\int_T\varphi\wedge\omega=\int_K\omega>0\).  But
\(\varphi^2=0\), since \(A\) is a curve.  Writing
\(\varphi=n\eta\), \(n\ne0\), yields
\(0=n^2\eta^2\ne0\), a contradiction.  Both possible positive algebraic
dimensions are excluded.
\end{proof}

Let
\[
 V=\mathbb Z\langle\gamma,u,w,\delta\rangle,
 \qquad
 \eta=u\wedge w+6\gamma\wedge\delta.
\]
For an integral two-class \(E\in\bigwedge^2V\), write its coordinates in
the ordered basis
\((\gamma u,\gamma w,\gamma\delta,uw,u\delta,w\delta)\) as
\(E_{\gamma u},\ldots,E_{w\delta}\), and set
\[
\begin{aligned}
P_E={}&E_{\gamma u}-E_{\gamma w}\tau-E_{\gamma\delta}\mu
 +6E_{uw}\mu+E_{u\delta}\beta\\
&+E_{w\delta}(6\mu^2-\tau\beta).
\end{aligned}
\tag{3.1}
\]

\begin{proposition}[very general Neron--Severi group]
\label{prop:very-general-torus-audit}
\Verified\ For the descended period family of
\cref{prop:period-input-audit}, outside a countable subset of
\(B^\circ\) the smooth fibre
\(F_b\) satisfies
\[
 \operatorname{NS}(F_b)=\mathbb Z\eta,
 \qquad a(F_b)=0.
\tag{3.2}
\]
\end{proposition}

\begin{proof}
The two holomorphic one-forms defining the Hodge subspace are
\[
 \sigma_1=6\mu\gamma+\tau u+w,
 \qquad
 \sigma_2=\beta\gamma+\mu u+\delta.
\]
The determinant inequality (2.2) gives
\(V_{\mathbb C}=F^1\oplus\overline{F^1}\), where
\(F^1=\langle\sigma_1,\sigma_2\rangle\).  For a real two-class \(E\),
the Hodge decomposition therefore gives
\[
 E\in H^{1,1}(F_b)
 \quad\Longleftrightarrow\quad
 E\wedge\sigma_1\wedge\sigma_2=0.
\]
Expanding the wedge product gives
\[
 E\wedge\sigma_1\wedge\sigma_2
 =P_E\,\gamma\wedge u\wedge w\wedge\delta.
\tag{3.3}
\]

We next determine exactly which \(P_E\) vanish identically.  Suppose
\[
 a+b\tau+c\mu+d\beta+e(6\mu^2-\tau\beta)\equiv0.
\tag{3.4}
\]
Apply the cusp transformation
\((\tau,\mu,\beta)\mapsto(\tau-1,\mu,\beta+1)\) and subtract.  This gives
\[
 -b+d+e(\beta-\tau+1)=0.
\]
Because \(\beta-\tau\) is translated by two, \(e=0\), and then \(d=b\).
Apply the order-four transformation in (2.1), subtract (3.4), and multiply
by \(-\tau\).  One obtains
\[
 6b\mu^2+c(\tau-1)\mu+b(\tau^2+3\tau+1)-c\tau=0.
\]
Subtract its cusp translate:
\[
 c\mu+2b(\tau+1)-c=0.
\tag{3.5}
\]
Translate (3.5) once around the cusp.  Since \(\mu\) is invariant and
\(\tau\mapsto\tau-1\), subtracting the translated identity from (3.5)
gives \(2b=0\).  Thus \(b=0\), and (3.5) becomes
\(c(\mu-1)=0\).  If \(c\ne0\), then \(\mu\equiv1\), whereas the
order-four law would give
\(1=\mu\circ g_2=1+1/\tau\), an impossibility.  Hence \(c=0\), then
\(b=d=0\), and finally \(a=0\).
Thus the five functions
\[
 1,\quad\tau,\quad\mu,\quad\beta,\quad6\mu^2-\tau\beta
\]
are linearly independent over \(\mathbb C\).

Reading the coefficients in (3.1) now gives
\[
 P_E\equiv0\quad\Longleftrightarrow\quad E\in\mathbb Z\eta.
\tag{3.6}
\]
For each other integral class, the zero set of \(P_E\) on the connected
Riemann surface \(\mathfrak h_z\) is discrete, hence countable.  There are
only countably many integral classes.  The class
\(\eta=q=uw+6\gamma\delta\) is monodromy invariant by
\cref{thm:coinvariants}, so it defines a class on every marked fibre after
descent.  The image in \(B^\circ\) of the countable union of the remaining
discrete zero sets is still countable.  A base point outside that image has
no additional integral \((1,1)\)-class on any lift, proving the
Neron--Severi statement on the quotient family.

Finally,
\[
 \eta^2=12\,\gamma\wedge u\wedge w\wedge\delta.
\tag{3.7}
\]
For a nondegenerate alternating form \(E\), write \(E^\vee\) for its
inverse alternating form on the dual real vector space.  If
\(G_\eta=(i\eta^\vee(\sigma_a,\overline{\sigma_b}))_{a,b}\), direct
inversion of the alternating-form matrix and substitution give
\[
G_\eta=
\begin{pmatrix}
2\operatorname{Im}\tau&2\operatorname{Im}\mu\\
2\operatorname{Im}\mu&\frac13\operatorname{Im}\beta
\end{pmatrix},
\]
\[
\det G_\eta=
\frac23\operatorname{Im}\tau
\left(\operatorname{Im}\beta-
\frac{6(\operatorname{Im}\mu)^2}{\operatorname{Im}\tau}\right)<0.
\tag{3.8}
\]
Here is the Appell--Humbert bridge needed to read the inertia of this
matrix.  With the convention \(E=\operatorname{Im}H_E\), one has
\[
 E=\frac1{2i}\sum_{a,b}(H_E)_{ab}\,
       \sigma_a\wedge\overline{\sigma_b}.
\]
Thus, in the ordered basis
\((\sigma_1,\sigma_2,\overline{\sigma_1},\overline{\sigma_2})\), the
upper-right block of the alternating-form matrix of \(E\) is
\(A=-(i/2)H_E\).  Inverting the block matrix gives
\[
 E^\vee(\sigma_a,\overline{\sigma_b})
   =[-(A^t)^{-1}]_{ab}
   =-2i(\overline{H_E}^{\,-1})_{ab},
 \qquad
 G_E=2\overline{H_E}^{\,-1}.
\tag{3.8a}
\]
The displayed \(G_\eta\) is real symmetric, so
\(H_\eta=2G_\eta^{-1}\).  Inversion and multiplication by a positive scalar
preserve inertia.  Hence \(H_\eta\) is nondegenerate and indefinite, while
(3.7) is nonzero.
Apply \cref{lem:torus-algebraic-reduction-audit}.
\end{proof}

\begin{lemma}[meromorphic descent]
\label{lem:meromorphic-descent-audit}
Let \(f:X\to B\) be a proper surjective holomorphic map from a connected
complex threefold to a connected compact Riemann surface, with connected
fibres, and a submersion over
\(B^\circ=B\setminus\Sigma\), where \(\Sigma\) is finite and
\(B^\circ\ne\varnothing\).  If
\(g\in\mathcal M(X)\) is constant on the fibres over an uncountable subset
of \(B^\circ\), then there is a unique \(\varphi\in\mathcal M(B)\) such that
\(g=f^*\varphi\).
\end{lemma}

\begin{proof}
Let \(P\) be the polar divisor of \(g\), and put
\(\Omega=f^{-1}(B^\circ)\setminus|P|\).  A proper analytic subset of a
connected complex manifold has real codimension at least two and does not
disconnect it.  First note that \(f^{-1}(B^\circ)\) is connected.  Indeed,
each of its connected components is open, and its image is open because
\(f\) is a submersion and closed because the base change
\(f^{-1}(B^\circ)\to B^\circ\) is proper.  Since \(B^\circ\) is connected,
every component would meet every fibre; connectedness of the fibres
therefore permits only one component.  Removing \(|P|\) now leaves
\(\Omega\) connected.  Only finitely many fibres are
components of \(|P|\).  For every other \(b\), the complement
\(F_b\cap\Omega\) of a proper analytic subset in the connected manifold
\(F_b\) is connected.

On \(\Omega\), define the analytic set
\[
 Z=\{dg\wedge f^*(dt)=0\},
\]
where \(t\) is any local coordinate on \(B\).  If \(g|_{F_b}\) is
constant, then \(F_b\cap\Omega\subset Z\).  If \(Z\ne\Omega\), then, because
\(\Omega\) is second countable, \(Z\) has at most countably many
irreducible components.  Each irreducible two-dimensional fibre contained
in \(Z\) is one of those components, and distinct fibres give distinct
components.  This contradicts the uncountable hypothesis.  Therefore
\(Z=\Omega\).

It follows that \(dg\) annihilates \(\ker df\).  Connectedness of
\(F_b\cap\Omega\), together with local holomorphic sections of the
submersion, gives \(g=\varphi\circ f\) on \(\Omega\), for a holomorphic
\(\varphi\) away from finitely many points of \(B\).

Fix such a point \(t_0\).  Choose a general holomorphic curve germ
\(\kappa:(\Delta,0)\to(X,x)\), not contained in \(|P|\), with
\[
 f\circ\kappa(s)=t_0+s^ku(s),\qquad u(0)\ne0.
\]
If \(g\circ\kappa\) has pole order at most \(M\), then
\(|g\circ\kappa(s)|\le C|s|^{-M}\).  After shrinking the disc, boundedness
above and below of \(|u|\) gives constants \(c,C_1>0\) such that
\[
 c|t-t_0|^{1/k}\le |s|\le C_1|t-t_0|^{1/k}
\]
on each local inverse branch.  Therefore
\[
 |\varphi(t)|\le C_2|t-t_0|^{-M/k}
\]
away from the image of a discrete exceptional set, and then everywhere on
the punctured disc by holomorphy.  For any integer \(N>M/k\), the function
\((t-t_0)^N\varphi(t)\) is bounded and extends holomorphically.  Thus
\(\varphi\) is meromorphic at \(t_0\).  Repeating at the finitely many
exceptional points produces \(\varphi\in\mathcal M(B)\).  Equality on the
dense open set \(\Omega\) gives equality of meromorphic functions on \(X\),
and surjectivity of \(f\) gives uniqueness.
\end{proof}

\begin{theorem}[algebraic dimension and algebraic reduction]
\label{thm:algebraic-dimension-audit}
\Verified\ For the global holomorphic family of
\cref{thm:candidate-global-geometry},
\[
 \mathcal M(X)=f^*\mathcal M(\mathbb P^1)=f^*\mathbb C(t),
 \qquad a(X)=1,
\tag{3.9}
\]
and \(f\) is the algebraic reduction.
\end{theorem}

\begin{proof}
By \cref{prop:very-general-torus-audit}, all but countably many smooth
fibres have algebraic dimension zero.  Let \(g\in\mathcal M(X)\).  Only
finitely many fibres are components of its polar divisor, so on uncountably
many of the preceding fibres the restriction of \(g\) is a meromorphic
function on a torus of algebraic dimension zero and is therefore constant.
By \cref{lem:meromorphic-descent-audit}, \(g=f^*\varphi\) for a unique
\(\varphi\in\mathcal M(\mathbb P^1)\).  The reverse inclusion is immediate,
and (3.9) follows.  Since \(f\) is proper, surjective, has connected fibres,
and induces the equality of meromorphic function fields, it is the algebraic
reduction.
\end{proof}

\paragraph{Disposition.}
The earlier label ``algebraic reduction open'' was too broad.  The direct
proof above validates the assertion independently of Section~7 and of the
CDP comparison.  Its provenance is the candidate's Lemmas~9.2--9.6 and
Proposition~9.7, reconstructed here with an explicit Appell--Humbert matrix
bridge, a proof that the smooth locus remains connected, and the omitted
root estimate in the final meromorphic-extension step.  This does not by
itself prove the claimed topology of \(X\); that is a separate calculation.

\section{Structure-sheaf direct images and uniform non-torsion}
\label{sec:analytic-picard}

This section reconstructs the analytic inputs used by the candidate's
Proposition~10.8.  The proof differs at one important point from the
candidate's route: extension across the cusp is obtained from the intrinsic
nearby-cycle isomorphism of \cref{thm:p0-integral-specialization}, not from
the cellular collapse.  The two cusp maps are now independently proved equal
in \cref{thm:collapse-equals-specialization}; the argument below nevertheless
retains its logically independent intrinsic formulation.

Put
\[
 e=d\zeta _1\wedge d\zeta _2
\]
on the marked torus family.  Retain the local coordinates and notation of
\cref{sec:candidate-geometry}:
\[
s_j\circ g_j=\zeta_js_j,\qquad
(m_1,m_2)=(3,4),\qquad t_j=s_j^{m_j},\qquad
f^{-1}(p_j)=m_jS_j.
\]

\subsection{The Hodge line and the finite normal characters}

\begin{lemma}[determinant characters and the Hodge line]
\label{lem:hodge-line-audit}
\Verified\ Let
\[
 \mathfrak K=f_*\omega_{X/B}.
\]
Then \(\mathfrak K\) is invertible and
\[
 \mathfrak K\simeq\mathcal O_B(1).
\tag{9.1}
\]
More precisely,
\[
 \det R_{g_1}=-\frac1\tau,\qquad
 \det R_{g_2}=\frac1\tau,\qquad
 \det R_{g_0}=1,
\tag{9.2}
\]
and at the fixed points
\[
 \det R_{g_j}(z_j)=\zeta_j^{a_j},
 \qquad (a_1,a_2)=(2,1).
\tag{9.3}
\]
The form \(e\) generates \(\mathfrak K\) at the cusp, while a unit times
\(s_j^{-a_j}e\) generates it at \(p_j\).
\end{lemma}

\begin{proof}
Equation (2.3d) in \cref{prop:period-input-audit} gives
\[
 R_g(z)\Pi(z)\lambda=\Pi(gz)A_g\lambda.
\]
Using
\[
 A_1\widehat w=\widehat u-\widehat w,\quad
 A_1\widehat\delta=\widehat\delta,\qquad
 A_2\widehat w=-\widehat u+\widehat\delta,\quad
 A_2\widehat\delta=\widehat\delta,
\]
and writing \(\tau^\sharp=\tau(g_jz)\),
\(\mu^\sharp=\mu(g_jz)\), direct multiplication gives
\[
 R_{g_1}=
 \begin{pmatrix}\tau^\sharp-1&0\\ \mu^\sharp&1\end{pmatrix},
 \qquad
 R_{g_2}=
 \begin{pmatrix}-\tau^\sharp&0\\1-\mu^\sharp&1\end{pmatrix}.
\]
The transformation laws
\(\tau\circ g_1=(\tau-1)/\tau\) and
\(\tau\circ g_2=-1/\tau\) prove (9.2) at \(g_1,g_2\); the cusp formulas
give \(R_{g_0}=I\).  Since \(\widehat\delta\) is fixed, one eigenvalue is
\(1\).  At \(\tau(z_1)=\rho=e^{\pi i/3}\), the other is
\(\rho-1=e^{2\pi i/3}=\zeta_1^2\); at \(\tau(z_2)=i\), it is
\(-i=\zeta_2\).  This proves (9.3).

We determine the extension over a finite filling without suppressing the
ramification.  On \(J'_j\), the nowhere-zero form \(ds_j\wedge e\)
transforms by
\[
 (\widetilde g_j^{\log})^*(ds_j\wedge e)
 =\chi_j(s_j)\,ds_j\wedge e,
 \qquad \chi_j=\zeta_j\det R_{g_j}.
\tag{9.4}
\]
Iteration \(m_j\) times shows
\(\prod_{r=0}^{m_j-1}\chi_j(\zeta_j^rs)=1\).
Write \(\chi_j(s)=\chi_j(0)e^{\psi(s)}\), with \(\psi(0)=0\).
The product identity says that the power series of \(\psi\) has no terms
whose degree is divisible by \(m_j\).  Hence
\[
 \varphi(s)-\varphi(\zeta_js)=\psi(s)
\]
has a holomorphic solution, obtained coefficientwise by division by
\(1-\zeta_j^n\).  Therefore
\[
 s^{k_j}e^{\varphi(s)}ds\wedge e,
 \qquad k_j=m_j-1-a_j,
\tag{9.5}
\]
is invariant: \(\chi_j(0)=\zeta_j^{1+a_j}\), so invariance is equivalent
to \(k_j+1+a_j\equiv0\pmod {m_j}\).  We now distinguish the descended
canonical section from the generator of its direct image.  Let
\[
 q_j:J'_j\longrightarrow N_j,\qquad
 \widetilde f_j:J'_j\longrightarrow\Delta_j,\qquad
 \varpi_j:\Delta_j\longrightarrow D_j
\]
denote the finite quotient, the torus-family projection, and the cyclic base
quotient, respectively, so that
\[
 f_j\circ q_j=\varpi_j\circ\widetilde f_j,
 \qquad t_j=s_j^{m_j}.
\]
Put \(\mathcal A_j=\widetilde f_j^{-1}(0)=\{s_j=0\}\) and
\(G_j=\langle\widetilde g_j^{\log}\rangle\).  The action is free by
\cref{prop:candidate-finite-fillings}; hence \(q_j\) is finite \'etale
and \(q_j^{-1}(S_j)=\mathcal A_j\).

The invariant form in (9.5), now denoted
\[
 \sigma_j=s_j^{k_j}e^{\varphi(s_j)}ds_j\wedge e,
 \qquad k_j=m_j-1-a_j,
\tag{9.5a}
\]
descends uniquely to
\(\overline\sigma_j\in H^0(N_j,K_{N_j})\) with
\(q_j^*\overline\sigma_j=\sigma_j\).  Since \(e^{\varphi(s_j)}\) is a
unit and \(ds_j\wedge e\) is nowhere vanishing,
\[
 \operatorname{div}_{J'_j}(\sigma_j)=k_j\mathcal A_j.
\]
Because \(q_j\) is \'etale,
\[
 q_j^*\operatorname{div}_{N_j}(\overline\sigma_j)
 =\operatorname{div}_{J'_j}(q_j^*\overline\sigma_j)
 =k_j\mathcal A_j=q_j^*(k_jS_j).
\]
Faithfulness of finite surjective pullback on Cartier divisors therefore gives
\[
 \operatorname{div}_{N_j}(\overline\sigma_j)=k_jS_j.
\tag{9.5b}
\]
Thus (9.5a) is a canonical section with that divisor, not a nowhere-vanishing
frame when \(k_j>0\); explicitly, \(k_1=0\) and \(k_2=2\).

For the direct-image generator, let
\[
 \vartheta_j=(ds_j\wedge e)\otimes(dt_j)^{-1}
\tag{9.5c}
\]
be the nowhere-vanishing frame of
\[
 q_j^*\omega_{N_j/D_j}
 =K_{J'_j}\otimes(f_j\circ q_j)^*K_{D_j}^{-1}.
\]
Here \((dt_j)^{-1}\) is the pulled-back frame of the inverse canonical line
of \(D_j\), not division by the differential of \(s_j^{m_j}\).  Put
\(\vartheta'_j=e^{\varphi(s_j)}\vartheta_j\).  The defining equation for
\(\varphi\), together with
\(\chi_j(s)=\chi_j(0)e^{\psi(s)}\) and
\(\chi_j(0)=\zeta_j^{1+a_j}\), gives
\[
 (\widetilde g_j^{\log})^*\vartheta'_j
 =\zeta_j^{1+a_j}\vartheta'_j.
\tag{9.5d}
\]
Every holomorphic coefficient of a relative top form is pulled back from
\(\Delta_j\), because the fibres of \(\widetilde f_j\) are connected compact
complex tori.  Hence
\[
 \widetilde f_{j*}q_j^*\omega_{N_j/D_j}
 =\mathcal O_{\Delta_j}\vartheta'_j.
\]
For \(F(s_j)\in\mathcal O_{\Delta_j}\), the section
\(F(s_j)\vartheta'_j\) is \(G_j\)-invariant exactly when
\[
 F(\zeta_js_j)\zeta_j^{1+a_j}=F(s_j).
\tag{9.5e}
\]
Writing \(F(s_j)=\sum_{n\ge0}c_ns_j^n\), this says that \(c_n=0\) unless
\[
 n+1+a_j\equiv0\pmod {m_j},
 \qquad\text{equivalently}\qquad n\equiv k_j\pmod {m_j}.
\]
Consequently,
\[
 \left((\varpi_j)_*\widetilde f_{j*}
 q_j^*\omega_{N_j/D_j}\right)^{G_j}
 =\mathcal O_{D_j}\cdot s_j^{k_j}\vartheta'_j,
\]
and finite descent gives the precise invariant-module identity
\[
 f_{j*}\omega_{N_j/D_j}
 =\mathcal O_{D_j}\cdot s_j^{k_j}\vartheta'_j.
\tag{9.5f}
\]
Finally, over \(\Delta_j^*\), the differential identity
\[
 dt_j=m_js_j^{m_j-1}ds_j
\]
identifies this generator with
\[
 s_j^{k_j}\vartheta'_j
 =\frac{e^{\varphi(s_j)}}{m_j}
 s_j^{k_j-(m_j-1)}e
 =\frac{e^{\varphi(s_j)}}{m_j}s_j^{-a_j}e.
\tag{9.5g}
\]
Thus a unit times \(s_j^{-a_j}e\) generates
\(f_{j*}\omega_{N_j/D_j}=(f_*\omega_{X/B})|_{D_j}\), whereas the
canonical section (9.5a) itself has divisor \(k_jS_j\).

At the cusp, in every smooth toric chart the unimodular monomial change of
coordinates turns
\[
 \Omega_0=\frac{dx_1}{x_1}\wedge
          \frac{dx_2}{x_2}\wedge dt_c
\]
into a unit times the ordinary coordinate volume form.  Thus it generates
the canonical line.  It is invariant under translations and under the
unimodular shears defining the quotient.  Under the punctured-overlap map,
\(\Omega_0=(2\pi i)^2dt_c\wedge e\), so \(e\) generates
\(\mathfrak K\) at \(p_0\).

It remains to compute the degree.  The modular-form calculation already
used in the proof of \cref{prop:period-input-audit} gives
\[
 \Theta=
 \left(\frac{E_4^2E_6^{1/2}}{\Delta_{\rm mod}}\right)\!\circ\tau.
\]
The square root exists because the zeros of \(E_6\circ\tau\) have even
order.  The weights \(8,3,12\) give the same automorphy factor as \(e\);
the sign at \(z_2\) is fixed by its linearizing coordinate.  The orders are
\[
 \operatorname{ord}_{z_1}\Theta=2=a_1,\qquad
 \operatorname{ord}_{z_2}\Theta=1=a_2,\qquad
 \Theta=t_c^{-1}\cdot(\text{unit})
\tag{9.6}
\]
at the cusp.  Consequently \(\mathsf e=\Theta^{-1}e\) is single-valued.
At \(p_j\) it is a unit times \(s_j^{-a_j}e\); at \(p_0\) it is a unit
times \(t_ce\).  Hence it is a global section of \(\mathfrak K\) with the
single simple zero \(p_0\).  This proves
\(\mathfrak K\simeq\mathcal O_B(p_0)\simeq\mathcal O_B(1)\).
Finally, \(\mathfrak K\) is torsion free of generic rank one on the smooth
curve \(B\), hence is invertible.
\end{proof}

\begin{lemma}[exact order of the finite-fibre normal line]
\label{lem:finite-normal-order-audit}
\Verified\ For \(j=1,2\),
\[
 \mathcal O_X(S_j)|_{S_j}=N_{S_j/X}
\]
has exact order \(m_j\) in \(\operatorname{Pic}(S_j)\).
\end{lemma}

\begin{proof}
On the finite cover of \cref{prop:candidate-finite-fillings}, the central
torus \(A_0^{(j)}=\{s_j=0\}\) maps to \(S_j\) as a connected principal
\(G_j\simeq\mathbb Z/m_j\)-cover.  Its normal line is trivial, and a
generator of \(G_j\) acts on the fibre by
\(s_j\mapsto\zeta_js_j\).  Thus \(N_{S_j/X}\) is associated to this
principal cover by the character \(\zeta_j\), and its \(m_j\)-th power is
trivial.

Suppose its \(k\)-th power were trivial for some \(0<k<m_j\).  Pulling a
nowhere-zero trivializing section back to \(A_0^{(j)}\) would give a
nowhere-zero holomorphic function \(h\) satisfying
\[
 h(gx)=\zeta_j^{k}h(x).
\]
Every holomorphic function on the compact connected torus \(A_0^{(j)}\)
is constant.  The displayed equation would force \(\zeta_j^k=1\), contrary
to \(0<k<m_j\).  Hence the order is exactly \(m_j\).
\end{proof}

% Complete finite-filling certificates, retaining the original marked periods.
% Host prerequisites: candidate_geometry.tex, exact_algebra.tex;
% normal-line comparison: analytic_picard.tex. No package or macro definitions.
\subsection{Exact product covers, normalization branches, and peripheral kernels of the finite fillings}
\label{ffcert:sec:finite-fillings}

We retain the actual period functions and finite fillings of
\cref{prop:period-input-audit,prop:monodromy-certificate,%
prop:candidate-finite-fillings}. The following explicit product maps
describe the original central tori and surfaces, not replacements of
their periods. We also identify the complete varying family as the
normalization of the actual ramified base change, compute the entire
peripheral inclusion kernel, and give a smooth trivialization of the
finite normal line. These calculations use no sphere identification,
global integral cohomology, or CDP theorem. The separate period and cusp
constructions retain their own proofs and are not inferred from these
finite-local calculations.

\subsubsection{The original lattice, periods, and their complex summands}

Use the ordered marked lattice and its dual basis
\[
 \Lambda=\mathbb Z\langle\widehat\gamma,\widehat u,
            \widehat w,\widehat\delta\rangle,
 \qquad \Lambda^*=\mathbb Z\langle\gamma,u,w,\delta\rangle.
 \tag{FF1}\label{ffcert:eq:marking}
\]
In these original column coordinates the matrices and translations are
\[
\begin{gathered}
 A_1=\begin{pmatrix}1&0&0&0\\6&0&1&0\\-6&-1&-1&0\\-2&1&0&1\end{pmatrix},
 \qquad
 A_2=\begin{pmatrix}1&0&0&0\\0&0&-1&0\\-6&1&0&0\\3&0&1&1\end{pmatrix},\\
 \varepsilon=(1,2,-4,0)^{\mathsf t},\quad
 \varepsilon'=(1,3,-3,0)^{\mathsf t},\quad
 v_1=\varepsilon,\quad v_2=-\varepsilon',\quad (m_1,m_2)=(3,4).
\end{gathered}
\tag{FF2}\label{ffcert:eq:matrices-translations}
\]
They satisfy $A_j^{m_j}=I$, $A_jv_j=v_j$, and
$\gamma A_j=\gamma$, with $\gamma(v_1)=1$ and $\gamma(v_2)=-1$.
Retain the entire period matrix
\[
 \Pi(z)=\begin{pmatrix}6\mu(z)&\tau(z)&1&0\\
                         \beta(z)&\mu(z)&0&1\end{pmatrix},
 \qquad
 D(z)=\operatorname{Im}\beta(z)
       -\frac{6(\operatorname{Im}\mu(z))^2}{\operatorname{Im}\tau(z)}<0.
 \tag{FF3}\label{ffcert:eq:periods}
\]
At the original elliptic centres let
\[
\begin{gathered}
 \rho=e^{\pi i/3},\quad \tau(z_1)=\rho,\quad \tau(z_2)=i,
 \quad \mu_1:=\mu(z_1)=\frac{2-\rho}{3},\quad
 \mu_2:=\mu(z_2)=\frac{1-i}{2},\\
 \beta_j=\beta(z_j),\qquad
 \kappa_1=\beta_1+2\mu_1,\qquad
 \kappa_2=\beta_2+3\mu_2.
\end{gathered}
\tag{FF4}\label{ffcert:eq:central-parameters}
\]
Neither $\beta_j$ nor either of its real components is adjusted.
Since $\operatorname{Im}\mu_1=-\operatorname{Im}\rho/3$ and
$\operatorname{Im}\mu_2=-1/2$, the exact inequalities are
\[
 \operatorname{Im}\kappa_1
 =\operatorname{Im}\beta_1-\tfrac23\operatorname{Im}\rho=D(z_1)<0,
 \qquad
 \operatorname{Im}\kappa_2
 =\operatorname{Im}\beta_2-\tfrac32=D(z_2)<0.
 \tag{FF5}\label{ffcert:eq:negative-orientation}
\]
Thus $(1,\kappa_j)$ is real-linearly independent with its original
negative imaginary orientation. No conjugation or positive-imaginary
replacement is made.

Put $\delta_*:=\widehat\delta$ and define
\[
\begin{aligned}
 e_1&=(0,1,-2,0)^{\mathsf t},&f_1&=(0,0,-3,1)^{\mathsf t},\\
 e_2&=(0,1,-1,0)^{\mathsf t},&f_2&=(0,0,-2,1)^{\mathsf t}.
\end{aligned}
\tag{FF6}\label{ffcert:eq:moving-generators}
\]

\begin{proposition}[The full fixed and moving integral lattices]
\label{ffcert:prop:integral-summands}
The exact fixed and image lattices are
\[
\begin{aligned}
 \Lambda^{A_1}&=\mathbb Z\varepsilon\oplus\mathbb Z\delta_*,&
 (A_1-I)\Lambda&=\mathbb Ze_1\oplus\mathbb Zf_1,\\
 \Lambda^{A_2}&=\mathbb Z\varepsilon'\oplus\mathbb Z\delta_*,&
 (A_2-I)\Lambda&=\mathbb Ze_2\oplus\mathbb Zf_2.
\end{aligned}
\tag{FF7}\label{ffcert:eq:integral-summands}
\]
The image lattices are saturated, with complete annihilators
\[
\begin{gathered}
 \operatorname{Ann}((A_1-I)\Lambda)
 =\mathbb Z\gamma\oplus\mathbb Z\psi_1,
 \quad \psi_1=2u+w+3\delta,\\
 \operatorname{Ann}((A_2-I)\Lambda)
 =\mathbb Z\gamma\oplus\mathbb Z\psi_2,
 \quad \psi_2=u+w+2\delta,\\
 \det[\varepsilon,\delta_*,e_1,f_1]=-3,
 \qquad \det[\varepsilon',\delta_*,e_2,f_2]=-2.
\end{gathered}
\tag{FF8}\label{ffcert:eq:annihilators-indices}
\]
In particular the sum of the fixed and moving lattices is not being
asserted to be all of $\Lambda$.
\end{proposition}

\begin{proof}
Solving $(A_1-I)(a,b,c,d)^{\mathsf t}=0$ gives $b=2a,c=-4a$;
the second equation gives $b=3a,c=-3a$. In both cases $a,d$ are arbitrary
integers. The columns of the two image matrices are, in order,
\[
 \begin{aligned}
 A_1-I &: \quad 6e_1-2f_1,\ -e_1+f_1,\ e_1,\ 0,\\
 A_2-I &: \quad 3f_2,\ -e_2,\ -e_2+f_2,\ 0.
 \end{aligned}
 \tag{FF9}\label{ffcert:eq:image-columns}
\]
These columns generate precisely the two stated image lattices. The first
is exactly the set of integer columns satisfying $a=0$ and
$2b+c+3d=0$: such a column is $be_1+df_1$. The second is exactly
the set satisfying $a=0$ and $b+c+2d=0$, with the analogous expression
$be_2+df_2$. They are therefore saturated. Their annihilators are the
displayed full integer spans: in the first case an annihilating row has
coefficients $(a,2c,c,3c)$, and in the second $(a,c,c,2c)$.
Expansion of the two four-by-four determinants gives $-3$ and $-2$.
\end{proof}

The complex images under the unchanged periods are
\[
\begin{aligned}
 \Pi(z_1)\varepsilon&=(0,\kappa_1),&
 \Pi(z_2)\varepsilon'&=(0,\kappa_2),&
 \Pi(z_j)\delta_*&=(0,1),\\
 \Pi(z_1)e_1&=(\rho-2,-(\rho-2)/3),&
 \Pi(z_1)f_1&=(-3,1),\\
 \Pi(z_2)e_2&=(i-1,-(i-1)/2),&
 \Pi(z_2)f_2&=(-2,1).
\end{aligned}
\tag{FF10}\label{ffcert:eq:complex-images}
\]
For example the first fixed coordinate vanishes because
$6(2-\rho)/3+2\rho-4=0$ or $6(1-i)/2+3i-3=0$, respectively.
The fixed real plane consequently maps onto the complex line
$\{(0,\xi):\xi\in\mathbb C\}$. The moving real planes map onto
$\{(\eta,-\eta/3)\}$ and $\{(\eta,-\eta/2)\}$, respectively,
because their displayed period pairs are real-linearly independent.
These are complementary complex lines in the original $\mathbb C^2$.

\subsubsection{The exact elliptic-product covers of the central tori and surfaces}

Define the one-dimensional complex tori
\[
\begin{gathered}
 E_j=\mathbb C/(\mathbb Z+\kappa_j\mathbb Z),\\
 F_1=\mathbb C/(\mathbb Z(\rho-2)+\mathbb Z(-3)),\qquad
 F_2=\mathbb C/(\mathbb Z(i-1)+\mathbb Z(-2)).
\end{gathered}
\tag{FF11}\label{ffcert:eq:elliptic-factors}
\]
Each period pair is real-linearly independent. Translation by its rank-two
lattice is free and properly discontinuous; the image of a closed real
parallelotope covers the quotient, proving compactness. Thus these are
smooth elliptic curves with the class of zero distinguished.

\begin{theorem}[Product maps onto the original finite central surfaces]
\label{ffcert:thm:product-covers}
The complex-linear maps
\[
 L_1(\xi,\eta)=(\eta,\xi-\eta/3),\qquad
 L_2(\xi,\eta)=(\eta,\xi-\eta/2)
 \tag{FF12}\label{ffcert:eq:addition-maps}
\]
induce surjective local biholomorphic group homomorphisms
\[
 \ell_j:E_j\times F_j\longrightarrow
 A_0^{(j)}:=\mathbb C^2/\Pi(z_j)\Lambda,
 \quad \ker\ell_1=\langle k_1\rangle,
 \quad \ker\ell_2=\langle k_2\rangle,
 \tag{FF13}\label{ffcert:eq:torus-covers}
\]
where
\[
 k_1=(2/3,\rho),\quad \operatorname{ord}(k_1)=3,
 \qquad k_2=(1/2,i),\quad \operatorname{ord}(k_2)=2.
 \tag{FF14}\label{ffcert:eq:kernel-points}
\]
The following automorphisms lift the original affine generators:
\[
 g^{\mathrm{prod}}_1(\xi,\eta)
   =(\xi+\kappa_1/3,(\rho-1)\eta),\qquad
 g^{\mathrm{prod}}_2(\xi,\eta)
   =(\xi-\kappa_2/4,-i\eta).
 \tag{FF15}\label{ffcert:eq:product-generators}
\]
Writing $k_j$ also for translation by that point, their joint groups
are $\langle k_1,g^{\mathrm{prod}}_1\rangle\cong C_3\times C_3$
and $\langle k_2,g^{\mathrm{prod}}_2\rangle\cong C_2\times C_4$.
They act freely, and the original reduced central surfaces are exactly
\[
 S_1=(E_1\times F_1)/\langle k_1,g^{\mathrm{prod}}_1\rangle,
 \qquad
 S_2=(E_2\times F_2)/\langle k_2,g^{\mathrm{prod}}_2\rangle.
 \tag{FF16}\label{ffcert:eq:surface-product-covers}
\]
The finite \'etale product-cover degrees are $9$ and $8$.
\end{theorem}

\begin{proof}
Each $L_j$ has determinant $-1$ in the displayed product and original
$(\zeta_1,\zeta_2)$ coordinates. Equation~\eqref{ffcert:eq:complex-images}
shows that the product period lattice maps exactly to the sum of the
fixed and moving sublattices inside $\Pi(z_j)\Lambda$.
Thus $L_j$ induces the stated surjective local biholomorphism, and its
kernel is identified with
$\Lambda/(\Lambda^{A_j}+(A_j-I)\Lambda)$, of order $3$ or $2$ by
\eqref{ffcert:eq:annihilators-indices}. Explicitly,
\[
 \widehat u=\tfrac23\delta_*+e_1-\tfrac23f_1,
 \qquad
 \widehat u=\tfrac12\delta_*+e_2-\tfrac12f_2.
 \tag{FF17}\label{ffcert:eq:kernel-lifts}
\]
Consequently $L_j(k_j)=\Pi(z_j)\widehat u$. The first kernel point
has order exactly three, since its fixed coordinate does, and
$3\rho=3(\rho-2)-2(-3)$ is a moving period. The second has order
exactly two, since its fixed coordinate does, and
$2i=2(i-1)-(-2)$ is a moving period. These orders equal the full kernel
cardinalities, proving the complete kernel assertions.

At the centres the original covariance matrices are
\[
 R_1=\begin{pmatrix}\rho-1&0\\(2-\rho)/3&1\end{pmatrix},\qquad
 R_2=\begin{pmatrix}-i&0\\(1+i)/2&1\end{pmatrix}.
 \tag{FF18}\label{ffcert:eq:central-covariance}
\]
They act by $1$ on the fixed lines and by
$\rho-1=e^{2\pi i/3}$ or $-i$ on the corresponding moving line.
Substitution gives the full equality of affine maps on $\mathbb C^2$,
\[
 L_jg_j^{\mathrm{prod}}
   =R_jL_j+\Pi(z_j)v_j/m_j.
 \tag{FF19}\label{ffcert:eq:product-intertwining}
\]
In particular the negative translation $-\kappa_2/4$ is forced by
$v_2=-\varepsilon'$. It has not been replaced by its opposite.

The moving rotations preserve their exact period lattices. For example,
using $\rho^2=\rho-1$, their images of the two period generators are
\[
\begin{aligned}
 (\rho-1)(\rho-2)&=-2(\rho-2)+(-3),&
 (\rho-1)(-3)&=-3(\rho-2)+(-3),\\
 (-i)(i-1)&=(i-1)-(-2),&
 (-i)(-2)&=2(i-1)-(-2).
\end{aligned}
\tag{FF20}\label{ffcert:eq:moving-rotations}
\]
The first matrix in these period bases has columns $(-2,1)$ and
$(-3,1)$, determinant $1$; the second has columns $(1,-1)$ and
$(2,-1)$, determinant $1$.
Equivalently, invariance and invertibility follow by applying the finite
orders of the rotations to these integral period images.
The kernel classes are fixed by the moving rotations because
\[
 (\rho-1)\rho-\rho=-1-\rho=-(\rho-2)+(-3),\qquad
 (-i)i-i=1-i=-(i-1).
 \tag{FF21}\label{ffcert:eq:kernel-rotation}
\]
Hence the product automorphism commutes with translation by $k_j$ on
the product torus. It has exact order $m_j$: its moving rotation has
that order; its $m_j$-th fixed-coordinate translation is
$\kappa_1$ or $-\kappa_2$; and no smaller positive multiple of that
translation is a fixed-coordinate period, since $\kappa_j$ is nonreal.

To prove both the direct-product assertion and freeness of the entire
joint group, put $(d_1,d_2)=(3,2)$, $(c_1,c_2)=(2/3,1/2)$ and
$(\epsilon_1,\epsilon_2)=(1,-1)$. A fixed point for $k_j^a
(g_j^{\mathrm{prod}})^b$, with $0\le a<d_j$ and $0\le b<m_j$,
would require
\[
 ac_j+\epsilon_jb\kappa_j/m_j=n+q\kappa_j,
 \qquad n,q\in\mathbb Z.
 \tag{FF22}\label{ffcert:eq:joint-freeness}
\]
Imaginary parts give $\epsilon_jb/m_j=q$, hence $b=0$. The real
equality then gives $ac_j\in\mathbb Z$, hence $a=0$ in its stated
range. Every nonidentity pair is therefore fixed-point free; in
particular it is not the identity transformation. This proves the
claimed joint groups. Quotienting first by $k_j$ and then by the lifted
affine generator gives exactly the original $S_j$ of
\cref{prop:candidate-finite-fillings}, by
\eqref{ffcert:eq:product-intertwining}. The total covering degrees are
$d_jm_j=9,8$ and the actions are free, so the covers are \'etale.
\end{proof}

\subsubsection{The varying quotient and its full local analytic coordinates}

Retain the original linearizing coordinate $s_j$ on $\Delta_j$, with
its actual round-disc radius $r_j>0$, and the original phase
\[
 s_j(g_jz)=\zeta_js_j(z),\quad
 \zeta_j=e^{-2\pi i/m_j},\quad t_j=s_j^{m_j},
 \qquad |s_j|<r_j,\quad |t_j|<r_j^{m_j}.
 \tag{FF23}\label{ffcert:eq:disc-and-phase}
\]
The radius is not replaced by one. On the universal marked cover the
holomorphic affine map and the exact period covariance are
\[
\begin{gathered}
 \widehat G_j(z,\zeta)
  =\bigl(g_jz,R_{g_j}(z)\zeta+\Pi(g_jz)v_j/m_j\bigr),\\
 R_{g_j}(z)\Pi(z)=\Pi(g_jz)A_j,
 \qquad t_\lambda(z,\zeta)=(z,\zeta+\Pi(z)\lambda).
\end{gathered}
\tag{FF24}\label{ffcert:eq:varying-action}
\]
Let $G_j$ denote the induced map of the original varying torus family
$J'_j=(\Delta_j\times\mathbb C^2)/\Lambda$ and put
$N_j=J'_j/\langle G_j\rangle$, with quotient map $q_j$.
The real-analytic marking $\Xi(z,x)=(z,\Pi(z)x)$ converts its powers
to the exact formulas
\[
 \Xi^{-1}\widehat G_j^k\Xi(z,x)
   =(g_j^kz,A_j^kx+kv_j/m_j),\qquad
 \widehat G_j^{m_j}=t_{v_j}.
 \tag{FF25}\label{ffcert:eq:varying-powers}
\]

\begin{proposition}[Smooth quotient charts and the exact multiplicities]
\label{ffcert:prop:quotient-charts}
The map $q_j:J'_j\to N_j$ is a finite \'etale covering of degree
$m_j$. The quotient is a Hausdorff smooth complex threefold, and its
proper map $f_j:N_j\to D_j$ has exact scheme fibre $m_jS_j$ at the
centre. In its local quotient coordinates the transverse map is
literally $t_j=s_j^{m_j}$.
\end{proposition}

\begin{proof}
At $s_j=0$ the freeness is proved in
\cref{ffcert:thm:product-covers}. Away from zero, a power with
$0<k<m_j$ has base coordinate multiplied by $\zeta_j^k\ne1$, so
has no fixed point. Equation~\eqref{ffcert:eq:varying-powers} shows
that the $m_j$-th power is the original lattice translation, hence is
the identity on $J'_j$. This gives a free action of exact order $m_j$.

For $y\in J'_j$, its finitely many distinct orbit points can be
separated by pairwise disjoint coordinate neighbourhoods. Intersecting
their inverse translates and shrinking gives a coordinate neighbourhood
$U$ of $y$ with the $G_j^aU$ pairwise disjoint. Then
$q_j|_U:U\to q_j(U)$ is a biholomorphism. Two distinct finite orbits
have disjoint saturated neighbourhoods by the same finite separation
argument. These charts prove the Hausdorff and smooth assertions, and
the quotient is a finite \'etale covering. They give directly the
trivial-stabilizer case of the quotient analytic structure in
\cite[conditions (I)--(II), Theorem~1]{Cartan1954Quotient}; smoothness
is not inferred from a theorem about arbitrary finite quotients.

At a point of $S_j$, use the quotient chart to retain $s_j$ as a
transverse coordinate, together with two fibre coordinates. The full
map is $t_j=s_j^{m_j}$. Thus $S_j$ is a smooth reduced Cartier
divisor, and its scheme multiplicity is exactly $m_j$.

For compact $C\subset\Delta_j$, the map induced by $\Xi$ from
$C\times[0,1]^4$ onto the inverse image of $C$ in $J'_j$ is
continuous and surjective, so that inverse image is compact. The actual
power map between the discs of radii $r_j$ and $r_j^{m_j}$ is proper:
the inverse image of a compact target set is closed, bounded, and stays
strictly inside the radius-$r_j$ disc. For compact $K\subset D_j$,
the inverse image in $J'_j$ of all roots of its $t_j$-values is
therefore compact. Its finite quotient image is exactly $f_j^{-1}(K)$,
proving properness of the full restriction.
\end{proof}

\subsubsection{Every branch in the normalization of the actual base change}

Let $\Delta_j^{\mathrm{base}}$ be the same radius-$r_j$ disc with
coordinate $u$, mapped to $D_j$ by $t_j=u^{m_j}$, and form
\[
 P_j=N_j\times_{D_j}\Delta_j^{\mathrm{base}},\qquad
 \nu_j:J'_j\longrightarrow P_j,\quad
 \nu_j(y)=(q_j(y),s_j(y)).
 \tag{FF26}\label{ffcert:eq:normalization-map}
\]

\begin{theorem}[The original varying family is the normalization]
\label{ffcert:thm:normalization}
The map $\nu_j$ is exactly the normalization of $P_j$. It is a
biholomorphism away from $u=0$; at the centre it has all $m_j$ local
branches prescribed by the original cyclic action. Its central fibre
is the connected original torus $A_0^{(j)}$, and
\[
 \nu_j(G_jy)=(q_j(y),\zeta_js_j(y)).
 \tag{FF27}\label{ffcert:eq:normalization-deck}
\]
Thus the original negative-phase action lifts $u\mapsto\zeta_ju$,
permutes the local branches cyclically, and has quotient $N_j$.
\end{theorem}

\begin{proof}
The definition uses no logarithm choice. If $u\ne0$, the $s_j$-values
in a covering orbit are exactly the $m_j$ distinct roots of its
$t_j$-value, so there is exactly one representative with value $u$.
If $u=0$, all $m_j$ covering points have $s_j=0$. This proves
surjectivity and describes all fibres. A compact subset of $P_j$
projects to a compact subset of $N_j$, whose inverse image under
$q_j$ is compact. Its inverse image under $\nu_j$ is a closed subset
of that compact set, hence is compact. Thus $\nu_j$ is proper with
finite fibres. Away from the centre the unique orbit representative
and quotient charts give the holomorphic inverse.

More precisely, around $n\in S_j$ choose a chart $q_j(U)$ as in
\cref{ffcert:prop:quotient-charts}, with inherited transverse coordinate
$s$ and fibre coordinates $z_1,z_2$. The fibre product is locally
\[
 s^{m_j}-u^{m_j}
   =\prod_{a=0}^{m_j-1}(s-e^{2\pi ia/m_j}u)=0.
 \tag{FF28}\label{ffcert:eq:all-branches}
\]
The factors are distinct smooth prime hypersurface factors in the
regular convergent-power-series ring $\mathbb C\{z_1,z_2,s,u\}$.
Their product ideal is radical: in this factorial ring, divisibility
of a power by every distinct prime factor implies divisibility of the
element by their product. Thus the fibre-product germ is reduced.
Its total ring of fractions is the product of the branch fraction
fields. Its integral closure is the product of the regular branch
rings $\mathbb C\{z_1,z_2,u\}$: adjoining the branch idempotents is
integral, since each satisfies $e^2-e=0$, and each component is
integrally closed. This proves the complete local normalization,
including all central branches.

On $G_j^aU$ the second component of $\nu_j$ is $u=\zeta_j^a s$,
or equivalently
\[
 s=\zeta_j^{-a}u=e^{2\pi ia/m_j}u.
 \tag{FF29}\label{ffcert:eq:branch-identification}
\]
Each of the $m_j$ disjoint covering charts therefore maps
isomorphically onto exactly one branch of
\eqref{ffcert:eq:all-branches}; these exhaust the branches. The
local identifications glue by the specified map $\nu_j$, giving the
global normalization. The source is the original varying family,
not an isotrivial product. Restricting its coordinate to zero gives
its connected torus $A_0^{(j)}$. Equation~\eqref{ffcert:eq:normalization-deck}
follows directly from $s_j(G_jy)=\zeta_js_j(y)$ and $q_jG_j=q_j$.
\end{proof}

\subsubsection{Universal covers and the full peripheral inclusion kernel}

\begin{theorem}[Exact finite-filling deck groups and attachment kernel]
\label{ffcert:thm:peripheral-kernel}
The universal cover $\Delta_j\times\mathbb C^2\to N_j$ has the
following complete deck-group presentation:
\[
 \Gamma_j=\left\langle\Lambda,\widehat G_j\ \middle|\
 \begin{array}{l}
 \Lambda\text{ abelian},\\
 \widehat G_jt_\lambda\widehat G_j^{-1}=t_{A_j\lambda},\\
 \widehat G_j^{m_j}=t_{v_j}
 \end{array}\right\rangle.
 \tag{FF30}\label{ffcert:eq:full-deck-group}
\]
This is also the universal-cover group of $S_j$. On the punctured
filling the complete group is $\Lambda\rtimes_{A_j}\mathbb Z$,
with the generator specified below, and the entire inclusion kernel is
\[
 \ker\{\pi_1(N_j\setminus S_j)\to\pi_1(N_j)\}
   =\langle\widetilde G_j^{m_j}t_{-v_j}\rangle
   \cong\mathbb Z.
 \tag{FF31}\label{ffcert:eq:whole-kernel}
\]
It is a central subgroup, not merely a necessary power relation.
\end{theorem}

\begin{proof}
In the real marking the original universal-cover transformations are
\[
 t_\lambda(s,x)=(s,x+\lambda),\qquad
 \widehat G_j(s,x)=(\zeta_js,A_jx+v_j/m_j).
 \tag{FF32}\label{ffcert:eq:cover-real-action}
\]
They satisfy the two relations in
\eqref{ffcert:eq:full-deck-group} by $A_jv_j=v_j$.
Moving translations to the left and using the power relation expresses
every word as $t_\lambda\widehat G_j^k$ with $\lambda\in\Lambda$
and $0\le k<m_j$. These forms are unique: the base action determines
$k$, and the remaining translation determines $\lambda$. Therefore
the abstract presented group acts faithfully and no additional relation
is implicit.

For $0<k<m_j$, a fixed-point equation would imply, on applying the
original integral covector $\gamma$ to the real fibre coordinate,
\[
 \gamma(\lambda)+k\gamma(v_j)/m_j=0.
 \tag{FF33}\label{ffcert:eq:cover-freeness}
\]
This is impossible because $\gamma(\lambda)$ is integral and
$\gamma(v_j)=1$ or $-1$. A nonzero translation with $k=0$ has no
fixed point either. The lattice action is properly discontinuous:
on compact base sets the real inverses $\Pi(z)^{-1}$ are uniformly
bounded, leaving finitely many lattice translations that can meet two
fixed compact sets. The full group is a finite union of lattice cosets,
so the same property holds for it. Thus the simply connected
$\Delta_j\times\mathbb C^2$ is the universal cover. Its restriction
to $s=0$ is the simply connected $\mathbb C^2$ covering $S_j$ with
the same action and group.

On the punctured disc retain exactly
\[
 w=\frac{\log s_j}{2\pi i},\qquad s_j=e^{2\pi iw},\qquad
 \operatorname{Im}w> -\frac{\log r_j}{2\pi},
 \qquad
 \widetilde G_j(w,x)=(w-1/m_j,A_jx+v_j/m_j).
 \tag{FF34}\label{ffcert:eq:log-cover-action}
\]
The displayed half-plane times $\mathbb C^2$, with its pulled-back
holomorphic coordinates, is simply connected. Along with the lattice
translations, $\widetilde G_j$ generates its entire punctured deck
group: the ordinary logarithm deck map is already
\[
 Q(w,x)=(w+1,x)=t_{v_j}\widetilde G_j^{-m_j}(w,x).
 \tag{FF35}\label{ffcert:eq:ordinary-log-deck}
\]
There is no finite-power relation for $\widetilde G_j$, since its
$k$-th power has $w$-displacement $-k/m_j$. Thus the exact normal
forms are $t_\lambda\widetilde G_j^k$ for unrestricted
$k\in\mathbb Z$, and the group is precisely the stated semidirect
product with action $A_j$.

The inclusion sends $t_\lambda$ to $t_\lambda$ and
$\widetilde G_j$ to $\widehat G_j$. Write $k=qm_j+r$ with
$0\le r<m_j$. Its image has normal form
$t_{\lambda+qv_j}\widehat G_j^r$, which is the identity exactly
when $r=0$ and $\lambda=-qv_j$. This proves
\eqref{ffcert:eq:whole-kernel} in both directions. Its generator has
$w$-displacement $-1$, so has infinite order. It commutes with all
translations because $A_j^{m_j}=I$, and with
$\widetilde G_j$ because $A_jv_j=v_j$, proving centrality.
\end{proof}

\subsubsection{The original punctured overlap and integral cohomology pullback}

On the affine punctured model choose a local logarithm and put
\[
 \alpha(z)=\frac{\log s_j(z)}{2\pi i},\qquad
 h_j[z,\zeta]=[z,\zeta+\alpha(z)\Pi(z)v_j].
 \tag{FF36}\label{ffcert:eq:overlap-map}
\]
A branch change adds $n\Pi(z)v_j$, an original lattice period,
so the map is well-defined and holomorphic on the punctured torus
quotient; subtraction of the same term gives its inverse. Under the
clockwise step $s_j\mapsto\zeta_js_j$, the compatible branch is
$\alpha(g_jz)=\alpha(z)-1/m_j$. The covariance
\eqref{ffcert:eq:varying-action} and $A_jv_j=v_j$ therefore give
the exact conjugacy to the original linear lift
$\widehat G_j^0(z,\zeta)=(g_jz,R_{g_j}(z)\zeta)$:
\[
 h_j\widehat G_jh_j^{-1}=\widehat G_j^0.
 \tag{FF37}\label{ffcert:eq:overlap-conjugacy}
\]
Indeed the two translations in $h_j\widehat G_j$ add to
$\Pi(g_jz)v_j(1/m_j+\alpha(g_jz))
=\alpha(z)\Pi(g_jz)v_j$, exactly the translation in
$\widehat G_j^0h_j$.

The corresponding real translation retains the full complex structure:
\[
\begin{gathered}
 J_z=\Pi(z)^{-1}i\Pi(z),\qquad
 a_j(z)=\operatorname{Re}\alpha(z)v_j
               +\operatorname{Im}\alpha(z)J_zv_j,\\
 J_{g_jz}A_j=A_jJ_z,\qquad
 a_j(g_jz)=A_ja_j(z)-v_j/m_j.
\end{gathered}
\tag{FF38}\label{ffcert:eq:real-overlap}
\]
This follows by applying $\Pi(z)^{-1}$ to
$\alpha(z)\Pi(z)v_j$ and then using covariance; complex scalar
multiplication is not being imposed on the marked real lattice.
The inverse image of the middle zero section is consequently
$-a_j(z)$, whose endpoint after one clockwise step is exactly
$A_j(-a_j(z))+v_j/m_j$. Iterating $m_j$ such steps increases
this marked endpoint by precisely $v_j$, not by an unspecified
translated generator. Thus the overlap map realizes the based
peripheral convention in \eqref{ffcert:eq:whole-kernel}.

The map from the fibre over $z$ to its deck-translated fibre over
$g_jz$ has marked linear part $A_j$; its inverse has part
$A_j^{-1}$. The original contragredient convention is
\[
 (A_j^{-1})^{\mathsf t}=T_j,
 \qquad A_j=(T_j^{-1})^{\mathsf t}.
 \tag{FF39}\label{ffcert:eq:contragredient-direction}
\]
These directions keep a based loop distinct from its inverse deck
map. No matrix or generator is reversed to obtain the attachment relation.

\begin{corollary}[The exact integral first-cohomology covering image]
\label{ffcert:cor:covering-h1}
For the original covering $\pi_j:A_0^{(j)}\to S_j$,
\[
 H_1(S_j,\mathbb Z)\cong\mathbb Z^2,\qquad
 \pi_j^*H^1(S_j,\mathbb Z)
   =\mathbb Z(m_j\gamma)\oplus\mathbb Z\psi_j
        \subset H^1(A_0^{(j)},\mathbb Z).
 \tag{FF40}\label{ffcert:eq:integral-h1-image}
\]
The index $m_j$ in the invariant covector lattice is retained.
\end{corollary}

\begin{proof}
By \cref{ffcert:prop:integral-summands}, the map
\[
 \Lambda\longrightarrow\mathbb Z^2,\qquad
 \lambda\longmapsto(\gamma(\lambda),\psi_j(\lambda))
 \tag{FF41}\label{ffcert:eq:coinvariant-coordinates}
\]
has kernel exactly $(A_j-I)\Lambda$. It is onto:
$\widehat\gamma$ gives $(1,0)$, and $\widehat w$ for $j=1$,
or $\widehat u$ for $j=2$, gives $(0,1)$. Thus the full
coinvariant group is $\mathbb Z^2$. In these coordinates $v_j$
is $(1,0)$ or $(-1,0)$. Abelianizing
\eqref{ffcert:eq:full-deck-group} adjoins one generator with relation
$m_j\widehat G_j=v_j$, eliminating a primitive generator and
leaving $\mathbb Z^2$. This computes $H_1(S_j,\mathbb Z)$ without
global topology. A homomorphism from this group to $\mathbb Z$
restricts on $\Lambda$ to $a\gamma+b\psi_j$ precisely when
$m_j\mid a$, since it must assign the integer
$a\gamma(v_j)/m_j$ to $\widehat G_j$. This necessary and sufficient
condition proves the exact covering pullback.
\end{proof}

\subsubsection{The finite-to-global analytic gluing interface}

The global construction of \cref{thm:candidate-global-geometry} glues
complete punctured restrictions of the original middle family and
the three local fillings. The finite identifications are precisely
\eqref{ffcert:eq:overlap-map}. The special discs are pairwise disjoint,
and a middle point over each punctured disc has exactly one
representative in its filling. Thus an equivalence class meets any
given summand at most once. Saturations of open subsets are open,
because the overlap identifications are biholomorphisms, so each summand
embeds openly and injectively.

Two points at one special base value have representatives in its same
filling; two points at one ordinary base value have representatives in
the middle piece. They have disjoint neighbourhoods in that Hausdorff
open summand, and their images remain disjoint because no equivalence
class meets that summand twice. Points at distinct base values are
separated by the base. This proves the gluing's Hausdorff assertion
using the actual full restrictions, not an unrestricted gluing principle.

In particular $f^{-1}(D_j)$ is the entire corresponding filling:
every middle point over $D_j^*$ is identified with its unique filling
representative. Thus the properness in
\cref{ffcert:prop:quotient-charts} is properness of the full inverse
image. For a compact base set $K$, finitely many relatively compact
$V_a$ whose closures lie inside members $U_a$ of the original
four-piece base cover can be chosen to cover $K$. The compact sets
$K\cap\overline V_a$ have compact preimages under the corresponding
full proper restrictions, and these preimages cover exactly
$f^{-1}(K)$. This is the compact-cover argument for global properness.
The separate cusp and middle constructions supply their own full
proper restrictions; the finite-local proof neither replaces those
constructions nor implies smooth sphere recognition.

\subsubsection{The exact normal character and a smooth trivialization}

\begin{theorem}[Holomorphic normal order and the explicit smooth product]
\label{ffcert:thm:normal-trivialization}
Let $\mathcal N_j=N_{S_j/N_j}$, equivalently
$\mathcal O_X(S_j)|_{S_j}$ after the specified open embedding into
the original global analytic construction. This line has exact order
$m_j$ in the holomorphic Picard group, but has the following explicit
smooth trivialization and hence zero integral first Chern class.
\end{theorem}

\begin{proof}
Differentiating the actual transverse coordinate in
\eqref{ffcert:eq:disc-and-phase} identifies the normal line as the
quotient of the trivial line over $A_0^{(j)}$ by
\[
 (x,n)\longmapsto(A_jx+v_j/m_j,\zeta_j n),
 \qquad \zeta_j=e^{-2\pi i/m_j}.
 \tag{FF42}\label{ffcert:eq:normal-action}
\]
Here $x$ is the original real marking of the central complex torus;
the action on the torus and line is holomorphic. A section of its
$k$-th power therefore pulls back to a function satisfying
\[
 h(A_jx+v_j/m_j)=\zeta_j^k h(x).
 \tag{FF43}\label{ffcert:eq:normal-section-character}
\]
The positive exponent follows from the quotient itself: the representative
$(x,h(x))$ is identified at the image point with
$(A_jx+v_j/m_j,\zeta_j^kh(x))$, which must equal the representative
of the section there. The negative exponent would instead describe
the dual normal line under the declared action.

The $m_j$-th power has the trivial character, so the constant function
one descends to a nowhere-zero holomorphic section of that power.
If a power $0<k<m_j$ were holomorphically trivial, its trivializing
section would pull back to a nowhere-zero holomorphic function on the
compact connected torus $A_0^{(j)}$. Such a function is constant,
and \eqref{ffcert:eq:normal-section-character} would force
$\zeta_j^k=1$, impossible in this range. This proves the exact
holomorphic order in \cref{lem:finite-normal-order-audit}, with the
normal tangent character fully specified.

For the smooth trivialization put
\[
 \epsilon_j=\gamma(v_j)\in\{1,-1\},\qquad
 H_j(x)=\exp(-2\pi i\epsilon_j\gamma(x))
 \quad(x\in\Lambda\otimes\mathbb R).
 \tag{FF44}\label{ffcert:eq:smooth-normal-section}
\]
Integrality of $\gamma(\lambda)$ gives $H_j(x+\lambda)=H_j(x)$,
so this is a well-defined nowhere-zero smooth function on the actual
marked central torus. Moreover $\gamma A_j=\gamma$ and
$\epsilon_j^2=1$ give
\[
 \begin{aligned}
 H_j(A_jx+v_j/m_j)
 &=\exp(-2\pi i\epsilon_j\gamma(x))
       \exp(-2\pi i\epsilon_j^2/m_j)\\
 &=\zeta_jH_j(x).
 \end{aligned}
 \tag{FF45}\label{ffcert:eq:smooth-equivariance}
\]
It consequently descends to a nowhere-zero section of the actual
normal line. Explicitly the smooth fibrewise-linear trivialization and
its inverse are
\[
 \begin{aligned}
 \Theta_j:\mathcal N_j&\longrightarrow S_j\times\mathbb C,
 & [x,n]&\longmapsto([x],n/H_j(x)),\\
 \Theta_j^{-1}:S_j\times\mathbb C&\longrightarrow\mathcal N_j,
 & ([x],z)&\longmapsto[x,zH_j(x)].
 \end{aligned}
 \tag{FF46}\label{ffcert:eq:normal-product-isomorphism}
\]
Replacing $x$ by a lattice translate leaves both formulas unchanged.
Replacing $(x,n)$ by
$(A_jx+v_j/m_j,\zeta_jn)$ also leaves $n/H_j(x)$ unchanged by
\eqref{ffcert:eq:smooth-equivariance}; the inverse has the same
equivariance. The two maps are therefore well-defined, smooth, and
inverse, retaining the original signs for each $j$. Thus
$c_1(\mathcal N_j)=0$ in integral cohomology while its holomorphic
order is exactly $m_j$. This conclusion uses only the original affine
action and integral covector, not the topology of a global $X$.
\end{proof}


\subsection{The cusp structure sheaf}

\begin{lemma}[structure-sheaf cohomology of the cusp fibre]
\label{lem:cusp-structure-sheaf-audit}
\Verified\ The reduced normal-crossings fibre \(W\) satisfies
\[
 h^q(W,\mathcal O_W)=(1,2,1)\qquad(q=0,1,2).
\tag{9.7}
\]
The natural map
\[
 H^1(W;\mathbb C)\longrightarrow H^1(W,\mathcal O_W)
\tag{9.8}
\]
is an isomorphism.  In degree two the corresponding map is surjective and
is injective on the kernel of pullback to the normalization
\(\nu:\mathrm{dP}_6\to W\).  Finally,
\[
 \bigwedge^2H^1(W,\mathcal O_W)
 \longrightarrow H^2(W,\mathcal O_W)
\tag{9.9}
\]
is an isomorphism.
\end{lemma}

\begin{proof}
Let \(D_1,D_2,D_3\simeq\mathbb P^1\) be the double curves and \(P,Q\) the
two triple points, with the branch order fixed in the proof of
\cref{thm:p0-integral-specialization}.  The normalization complex is
\[
0\longrightarrow\mathcal O_W\longrightarrow
 \nu_*\mathcal O_{\mathrm{dP}_6}\xrightarrow{\delta_0}
 \bigoplus_{k=1}^3\mathcal O_{D_k}\xrightarrow{\delta_1}
 \mathbb C_P\oplus\mathbb C_Q\longrightarrow0,
\tag{9.10}
\]
where \(\delta_0\) is the difference of the two branch restrictions and
\[
 \delta_1(g_1,g_2,g_3)
 =(g_1(P)-g_2(P)+g_3(P),
   g_1(Q)-g_2(Q)+g_3(Q)).
\]
For completeness, exactness is stalkwise.  At a generic double point it is
the fibre-product sequence
\[
0\to\mathbb C\{x,y,z\}/(xy)
 \to\mathbb C\{y,z\}\oplus\mathbb C\{x,z\}
 \to\mathbb C\{z\}\to0.
\]
At a triple point it is the augmented branch complex for
\(\mathbb C\{x,y,z\}/(xyz)\): three branch rings, three pairwise
intersection rings, and the common residue field, with alternating
differences.  The augmented two-simplex complex is exact; successively
lifting compatible power series proves exactness of the analytic rings.
Away from the singular locus the first arrow is an isomorphism.

The same formulas give an exact complex with \(\mathbb C\) in place of
\(\mathcal O\).  Since \(\mathrm{dP}_6\) is rational and each \(D_k\) is a
rational curve, the positive structure-sheaf cohomology of every term on
the right of (9.10) vanishes.  Hypercohomology therefore reduces to
\[
 \mathbb C\xrightarrow{0}\mathbb C^3
 \xrightarrow{d_1}\mathbb C^2,
 \qquad
 d_1(a_1,a_2,a_3)=(s,s),\quad s=a_1-a_2+a_3.
\tag{9.11}
\]
Its cohomology has dimensions \(1,2,1\), proving (9.7).

For the constant-sheaf complex, the row of degree zero is again (9.11),
the row of degree one vanishes, and the remaining degree-two row begins
with \(H^2(\mathrm{dP}_6;\mathbb C)\).  The comparison induced by
\(\mathbb C_W\hookrightarrow\mathcal O_W\) is the identity on (9.11) and
zero on the positive rows of the target.  It follows that (9.8) is an
isomorphism, while in degree two it maps the dual-complex summand
\(\ker\nu^*\) isomorphically onto \(H^2(W,\mathcal O_W)\).

Choose the basis \(\gamma_W,u_W\in H^1(W;\mathbb Z)\) whose intrinsic
specializations are \(\gamma,u\), respectively; it exists and is unique by
\cref{thm:p0-integral-specialization}.  Their cup product is nonzero,
because its specialization is
\(\gamma\wedge u\ne0\in H^2(F;\mathbb Z)\).  Its pullback to
\(\mathrm{dP}_6\) is zero since
\(H^1(\mathrm{dP}_6;\mathbb C)=0\).  The degree-two comparison therefore
sends it to a nonzero element of \(H^2(W,\mathcal O_W)\).  The images of
\(\gamma_W,u_W\) form a basis by (9.8), and both sides of (9.9) are
one-dimensional.  Hence (9.9) is an isomorphism.
\end{proof}

\subsection{Direct images, base change, and the Picard line}

\begin{theorem}[structure-sheaf direct images]
\label{thm:structure-direct-images-audit}
\Verified\ For the constructed map \(f:X\to B=\mathbb P^1\),
\[
 f_*\mathcal O_X=\mathcal O_B,\qquad
 R^1f_*\mathcal O_X\simeq
 \mathcal O_B\oplus\mathcal O_B(-1),\qquad
 R^2f_*\mathcal O_X\simeq\mathcal O_B(-1).
\tag{9.12}
\]
All three direct images are locally free and every base-change morphism
\[
 (R^qf_*\mathcal O_X)\otimes k(b)
 \longrightarrow H^q(X_b,\mathcal O_{X_b}),
 \qquad 0\le q\le2,
\tag{9.13}
\]
is an isomorphism, including at the two nonreduced fibres.  Moreover the
invariant class \(\gamma\) defines the trivial summand in
\(R^1f_*\mathcal O_X\) and is nowhere zero on \(B\).
\end{theorem}

\begin{proof}
The proper map has connected fibres and \(B\) is normal, so Stein
factorization gives \(f_*\mathcal O_X=\mathcal O_B\).

The map is flat.  A local parameter on \(B\) pulls back to a nonzero
element of the regular local domain \(\mathcal O_{X,x}\); if it vanished
identically near one point, analytic continuation on the connected
manifold would make \(f\) constant.  Thus \(\mathcal O_{X,x}\) is
torsion-free, hence flat, over the discrete valuation ring
\(\mathcal O_{B,f(x)}\).

Over \(B^\circ\), the fibres are two-dimensional complex tori and their
structure-sheaf cohomology dimensions are \(1,2,1\).  At the cusp the same
dimensions hold by \cref{lem:cusp-structure-sheaf-audit}.  Grauert's
theorem therefore gives local freeness of ranks \(1,2,1\) and base change
on a neighbourhood of \(B^\circ\cup\{p_0\}\); the precise analytic statement is
\cite[Chapter~III, \S4, Theorem~4.12]{BanicaStanasila1976}, originally due to
Grauert \cite{Grauert1960}.

At \(p_j\), use the finite cover \(J'_j\to N_j\).  If
\(\widetilde f:J'_j\to\Delta_j\) is the smooth torus family and
\(\pi(s)=s^{m_j}\), exactness of finite-group invariants in characteristic
zero gives
\[
 R^qf_*\mathcal O_{N_j}
 =\bigl(\pi_*R^q\widetilde f_*\mathcal O_{J'_j}\bigr)^{G_j}.
\tag{9.14}
\]
After averaging, a free eigenframe \(v\) may be chosen for each character
\(\nu\).  The invariant sections are exactly
\(s^rv\,\mathcal O_{D_j}\), where \(0\le r<m_j\) is determined by
\(\nu\zeta_j^r=1\).  Hence (9.14) is free of ranks \(1,2,1\).
Higher direct images vanish by the fibre-dimension bound, so these are all
the nonzero direct images.  Since all direct images are locally free, the analytic
cohomology-and-base-change theorem, applied by descending induction from
\(q=2\), proves (9.13) at the finite fibres as well.  In the notation of
\cite[Chapter~III, \S3, Theorem~3.4, Corollaries~3.5, 3.7, and 3.10]
  {BanicaStanasila1976}.  Corollary~3.10 says that, for each fixed \(q\),
  freeness of \(R^p f_*\mathcal O_X\) for every \(p\ge q\) is equivalent to
  cohomological flatness in all dimensions \(p\ge q\).  Applying this
  successively with \(q=2,1,0\) gives the displayed base-change maps in every
  relevant degree.  This is a local statement on the base and so applies
  equally to the nonreduced scheme-theoretic fibres.

Write \(G=R^1f_*\mathcal O_X\).  We construct its distinguished section
without using a cellular collapse.  On \(B^\circ\), the monodromy-invariant
integral class \(\gamma\) maps to a holomorphic section
\(\gamma^{0,1}\) of \(G\).  It is nowhere zero there: if its image in
\(H^{0,1}(F)\) vanished, then
\(\gamma\in H^{1,0}(F)\cap H^1(F;\mathbb R)=0\).  On a finite cover it is
a constant invariant class and remains nonzero at \(s_j=0\), so it
descends across \(p_j\).

At the cusp, let \(j:\Delta^*\hookrightarrow\Delta\).  The natural
specialization morphism
\[
 R^1f_*\mathbb C\longrightarrow
 j_*R^1f_*\mathbb C|_{\Delta^*}
\]
is an isomorphism: away from \(0\) this is tautological, while at \(0\) it
is exactly the isomorphism from \(H^1(W;\mathbb C)\) onto the
\(T_0\)-invariant subspace proved in
\cref{thm:p0-integral-specialization}.  Hence the invariant classes
\(\gamma,u\) extend as sections of \(R^1f_*\mathbb C\) on the whole disc.
The natural morphism induced by \(\mathbb C\hookrightarrow\mathcal O\),
\[
 R^1f_*\mathbb C\otimes\mathcal O_\Delta
 \longrightarrow G|_\Delta,
\]
sends them at the central point to the two classes in
\(H^1(W,\mathcal O_W)\).  These form a basis by
\cref{lem:cusp-structure-sheaf-audit}.  Thus
\(\gamma^{0,1}\) extends across the cusp and is nowhere zero everywhere.
In particular \(\mathcal O_B\gamma^{0,1}\subset G\) is a line subbundle.

Relative duality applies to the proper flat Gorenstein map \(f\): \(X\) is
smooth, the base is a smooth curve, and \(\omega_{X/B}\) is invertible.
Using \cref{lem:hodge-line-audit}, it gives
\[
 R^2f_*\mathcal O_X\simeq\mathfrak K^\vee
 \simeq\mathcal O_B(-1).
\tag{9.15}
\]
The relative theorem is the one of Ramis--Ruget--Verdier
\cite[pp.~261--283]{RamisRugetVerdier1971}; the absolute singular
dualizing-complex comparison is independently recorded in
\cite[Theorems~1--2 and Proposition~2]{RamisRuget1970}.

The relative cup product
\[
 \bigwedge^2G\longrightarrow R^2f_*\mathcal O_X
\tag{9.16}
\]
is an isomorphism.  It is so on every smooth torus fibre.  At the cusp it
is so by (9.13) and \cref{lem:cusp-structure-sheaf-audit}.  At \(p_j\),
choose on the cover an eigenframe \(v_1,v_2\), with
\(v_1=\gamma^{0,1}\) invariant.  Its invariant exponent is \(r_1=0\).
The action on \(v_1\wedge v_2\) is dual to the Hodge character in (9.3),
so its invariant generator is \(s_j^{a_j}(v_1\wedge v_2)\).
The invariant generator for the second line is consequently
\(s_j^{a_j}v_2\), because \(0<a_j<m_j\).  Wedge product sends
\(v_1\wedge(s_j^{a_j}v_2)\) to exactly the invariant generator of the
right side of (9.16), with no wrap modulo \(m_j\).

Consequently
\(\det G\simeq R^2f_*\mathcal O_X\simeq\mathcal O_B(-1)\).  The quotient
of \(G\) by its trivial subbundle is \(\mathcal O_B(-1)\), and the
extension splits because
\[
 \operatorname{Ext}^1(\mathcal O_B(-1),\mathcal O_B)
 =H^1(B,\mathcal O_B(1))=0.
\]
This proves (9.12) and identifies the trivial summand.
\end{proof}

\begin{corollary}[the Picard parameter and fibrewise restriction]
\label{cor:picard-line-audit}
\Verified\ One has
\[
 H^1(X,\mathcal O_X)=\mathbb C,\qquad
 \operatorname{Pic}(X)=\operatorname{Pic}^0(X)\simeq(\mathbb C,+).
\tag{9.17}
\]
For every scheme-theoretic fibre \(X_b\), restriction induces an injective
complex-linear map
\[
 \iota_b:H^1(X,\mathcal O_X)\longrightarrow
 H^1(X_b,\mathcal O_{X_b}).
\tag{9.18}
\]
On a smooth fibre its image is the complex line generated by the
\((0,1)\)-image of \(\gamma\).
\end{corollary}

\begin{proof}
The Leray spectral sequence and \(H^1(B,\mathcal O_B)=0\) identify
\[
 H^1(X,\mathcal O_X)=H^0(B,R^1f_*\mathcal O_X)=\mathbb C
\]
by \cref{thm:structure-direct-images-audit}.  Since the verified integral
topology gives \(H^1(X;\mathbb Z)=H^2(X;\mathbb Z)=0\), the analytic
exponential sequence
\[
0\to\mathbb Z\to\mathcal O_X
 \xrightarrow{u\mapsto e^{2\pi iu}}\mathcal O_X^*\to1
\]
identifies \(H^1(X,\mathcal O_X)\) with \(\operatorname{Pic}(X)\) and
shows that every line bundle is topologically trivial.  Under Leray,
restriction is evaluation of a global section of
\(R^1f_*\mathcal O_X\), followed by base change.  The only global section
is a scalar multiple of the nowhere-zero section \(\gamma^{0,1}\), so
evaluation is injective at every point and has the stated image on a
smooth fibre.
\end{proof}

\subsection{One line bundle non-torsion on every fibre}

\begin{theorem}[uniform non-torsion]
\label{thm:uniform-nontorsion-audit}
\Verified\ There is a line bundle \(L\in\operatorname{Pic}(X)\) such that
\(L|_{X_b}\) is non-torsion for every scheme-theoretic fibre
\(X_b=f^*(b)\), and \(L|_{S_j}\) is non-torsion for \(j=1,2\).
More precisely, identify \(\operatorname{Pic}(X)\) with \(\mathbb C\),
\(z\mapsto L_z\).  Every \(z\) outside the union of \(\mathbb R\) with a
countable set has these properties.
\end{theorem}

\begin{proof}
The pullback \(f^*\mathcal O_B(1)\) is nontrivial: by the projection
formula and \(f_*\mathcal O_X=\mathcal O_B\),
\[
 h^0(X,f^*\mathcal O_B(1))=h^0(B,\mathcal O_B(1))=2,
\]
whereas \(h^0(X,\mathcal O_X)=1\).  Rescale the additive coordinate in
\cref{cor:picard-line-audit} so that
\[
 L_{12}=f^*\mathcal O_B(1).
\tag{9.19}
\]

Fix a smooth fibre \(F_q\).  Put
\[
 \Lambda_q=\iota_q^{-1}\!\left(
 \operatorname{im}[H^1(F_q;\mathbb Z)\to
 H^1(F_q,\mathcal O_{F_q})]\right)\subset\mathbb C.
\tag{9.20}
\]
The integral image is a lattice in the underlying real vector space of
\(H^1(F_q,\mathcal O_{F_q})\), so \(\Lambda_q\) is discrete.  It contains
\(12\), because (9.19) is trivial on every fibre.  The exponential
sequence on \(F_q\) shows
\[
 L_z|_{F_q}\text{ is torsion}
 \quad\Longleftrightarrow\quad z\in\mathbb Q\Lambda_q.
\tag{9.21}
\]
If \(\operatorname{rk}\Lambda_q=1\), its nonzero element \(12\) forces
\(\Lambda_q\subset\mathbb R\); all its bad parameters are real.

Suppose \(\operatorname{rk}\Lambda_q=2\).  By
\cref{cor:picard-line-audit}, the image of \(\iota_q\) is the line
generated by the image of \(\gamma\).  Thus there is an integral class
\[
 \lambda=a\gamma+bu+cw+d\delta\notin\mathbb Q\gamma
\tag{9.22}
\]
whose \((0,1)\)-image lies on that line.  On the covering family, the four
marked lattice columns are
\[
 \pi_{\widehat\gamma}=(6\mu,\beta),\quad
 \pi_{\widehat u}=(\tau,\mu),\quad
 \pi_{\widehat w}=(1,0),\quad
 \pi_{\widehat\delta}=(0,1).
\tag{9.23}
\]
Every real cohomology class has a unique expression
\[
 2\operatorname{Re}(\bar\xi\,d\zeta_1+
                       \bar\upsilon\,d\zeta_2).
\]
Evaluating it on (9.23) gives
\[
\begin{aligned}
2\operatorname{Re}\bar\xi&=c,&
2\operatorname{Re}\bar\upsilon&=d,\\
2\operatorname{Re}(\bar\xi\tau+\bar\upsilon\mu)&=b,&
2\operatorname{Re}(6\bar\xi\mu+\bar\upsilon\beta)&=a.
\end{aligned}
\tag{9.24}
\]
For \(\gamma\), the first two right sides vanish.  The third equation then
shows that its \((0,1)\)-line is represented, up to a nonzero scalar, by
\((\operatorname{Im}\mu,-\operatorname{Im}\tau)\); the fourth equation
only fixes the scalar.  Write
\(\bar\xi=c/2+ir\), \(\bar\upsilon=d/2+is\).  The condition that the class
in (9.22) have \((0,1)\)-image on the \(\gamma\)-line is
\[
 \bar\xi\operatorname{Im}\tau+
 \bar\upsilon\operatorname{Im}\mu=0.
\]
Its real part, together with the third equation of (9.24), is equivalent
to the single holomorphic equation
\[
 c\tau+d\mu-b=0.
\tag{9.25}
\]
Neither \(a\) nor \(\beta\) occurs, as adding a multiple of \(\gamma\)
must not change the condition.

The functions \(1,\tau,\mu\) are linearly independent.  Applying the cusp
transformation \((\tau,\mu)\mapsto(\tau-1,\mu)\) to a constant linear
relation first kills the coefficient of \(\tau\).  If the remaining
relation made \(\mu\) constant, the order-four law
\(\mu\circ g_2=1+\mu/\tau\), with nonconstant \(\tau\), would force an
impossible identity.  Hence all coefficients vanish.

For every integral \(\lambda\notin\mathbb Q\gamma\), equation (9.25) is
therefore a nonzero holomorphic function on the covering Riemann surface.
Its zero set is discrete and countable.  There are countably many
\(\lambda\), so their union, and its image in \(B^\circ\), are countable.
This image contains every \(q\) for which
\(\operatorname{rk}\Lambda_q=2\).  Each such fibre contributes the
countable bad set \(\mathbb Q\Lambda_q\); all rank-one fibres together
contribute only parameters in \(\mathbb R\).  Thus the smooth fibres
exclude only \(\mathbb R\) together with a countable set.

There remain the three scheme-theoretic special fibres.  By
\cref{thm:structure-direct-images-audit,cor:picard-line-audit}, restriction
\[
 \mathbb C\longrightarrow H^1(X_b,\mathcal O_{X_b})
\]
is complex-linear and injective for each of them.  The exponential sequence
and countability of \(H^1(X_b;\mathbb Z)\) show that the parameters for
which \(L_z|_{X_b}\) is torsion form a countable set.

Finally consider the reductions \(S_j\).  The restriction
\(\varphi_j:\mathbb C\to H^1(S_j,\mathcal O_{S_j})\) is complex-linear.
It is nonzero.  Write \(\mathcal O_X(S_j)=L_{z_j}\).  Since
\[
 \mathcal O_X(m_jS_j)=f^*\mathcal O_B(1)=L_{12},
\]
the additive parameter gives \(z_j=12/m_j\), while
\cref{lem:finite-normal-order-audit} says that
\[
 L_{12/m_j}|_{S_j}=\mathcal O_X(S_j)|_{S_j}
\]
is nontrivial.  If \(\varphi_j\) were zero, the exponential sequence and
the vanishing of the restricted first Chern class would make every
\(L_z|_{S_j}\) trivial, a contradiction.  Thus \(\varphi_j\) is
injective, and the same exponential-sequence argument makes its torsion
parameters countable.  Removing these two additional countable sets proves
the theorem.
\end{proof}

\begin{remark}[provenance of the repair]
The period equation \(c\tau+d\mu-b=0\) and the countability strategy are
the candidate's.  The exact normal-character proof above makes the order of
\(\mathcal O_X(S_j)|_{S_j}\) explicit, while the cusp extension in
\cref{thm:structure-direct-images-audit} replaces the candidate's
collapse-dependent argument by the independently verified intrinsic
nearby-cycle map.  Identifying the restriction line with the global
\(\gamma^{0,1}\)-summand also removes the candidate's separate continuity
argument for a moving integral lift.
\end{remark}

\subsection{A local analytic obstruction to a cusp-only Picard argument}
\label{subsec:local-cusp-picard-obstruction}

The relative Picard space of the cusp family cannot by itself force the
normalised central restriction to vanish.  The obstruction is visible
directly in the toric model of \cref{thm:candidate-cusp-geometry}.  We give
the complete factor-of-automorphy calculation, including the descent
cocycle and the sign of the generic Chern class.

Retain the cusp notation of
\cref{prop:period-input-audit,thm:candidate-cusp-geometry}:
\[
 B_0=\begin{pmatrix}0&1\\-1&0\end{pmatrix},
 \qquad
 Z(s)=sB_0+C(t),
 \qquad
 C(t)=\begin{pmatrix}6\mu(t)&h(t)\\ b(t)-h(t)&\mu(t)\end{pmatrix},
 \qquad t=e^{2\pi i s},
\tag{5.17t}
\]
\[
 c_{\lambda}(t)=\exp(2\pi iC(t)\lambda),
 \qquad
 \exp(2\pi iZ(s)\lambda)=t^{B_0\lambda}c_\lambda(t)
 \quad(\lambda\in\bar\Lambda\simeq\mathbb Z^2),
\tag{5.17u}
\]
and
\[
 \Phi_\lambda(x,t)=(t^{B_0\lambda}x,t),
 \qquad
 \Psi_\lambda(p)=(c_\lambda(t(p)),1)\cdot\Phi_\lambda(p),
 \qquad
 \Psi_\lambda\Psi_\mu=\Psi_{\lambda+\mu}.
\tag{5.17v}
\]
Write \(E_v\), \(v\in\mathbb Z^2\), for the toric prime divisors of
\(Y_\epsilon\), so that \(E_0\simeq\mathrm{dP}_6\), and let
\[
 \pi:Y_\epsilon\longrightarrow N_0=Y_\epsilon/\bar\Lambda,
 \qquad
 \nu=\pi|_{E_0}:E_0\longrightarrow W
\]
be the quotient and the normalization of the central fibre.
Put
\[
 N'=\mathbb Z^2\oplus\mathbb Z,
 \qquad
 M'=\operatorname{Hom}(N',\mathbb Z).
\]
For \(m=(\ell,k)\in M'=\mathbb Z^2\oplus\mathbb Z\), fix the character
convention
\[
 \chi^m(x,t)=x^\ell t^k,
 \qquad
 \operatorname{div}(\chi^m)=
 \sum_{v\in\mathbb Z^2}\langle m,(v,1)\rangle E_v.
\tag{5.17w}
\]

\begin{proposition}[an explicit nonzero local cusp restriction]
\label{prop:local-cusp-picard-obstruction}
\Verified\ Let
\[
 H=\begin{pmatrix}1&0\\0&-6\end{pmatrix},
 \qquad
 P=B_0^{\mathsf t}H=\begin{pmatrix}0&6\\1&0\end{pmatrix},
 \qquad
 Q(v)=v^{\mathsf t}Hv=v_1^2-6v_2^2.
\tag{5.17x}
\]
Then
\[
 \mathcal D_Q=\sum_{v\in\mathbb Z^2}Q(v)E_v
\tag{5.17y}
\]
is a locally finite Cartier divisor on \(Y_\epsilon\).
The line bundle \(\mathcal O_{Y_\epsilon}(\mathcal D_Q)\) has an explicit
\((\bar\Lambda,\Psi)\)-equivariant structure and descends to a holomorphic
line bundle \(\mathcal L_Q\) on \(N_0\).  If
\[
 (a_1,a_2,a_3,a_4,a_5,a_6)
 =(e_1,e_2,e_2-e_1,-e_1,-e_2,e_1-e_2),
 \qquad C_i=E_0\cap E_{a_i},
\]
then
\[
 \nu^*(\mathcal L_Q|_W)
 \simeq
 \mathcal O_{E_0}\!\left(\sum_{i=1}^6Q(a_i)C_i\right),
\qquad
 \left(\deg_{C_i}\nu^*(\mathcal L_Q|_W)\right)_{i=1}^6
 =(-12,2,0,-12,2,0).
\tag{5.17z}
\]
In particular \(\nu^*(\mathcal L_Q|_W)\) is nontrivial.

On every nonzero fibre
\[
 F_t=\mathbb C^2/(Z(s)\mathbb Z^2+\mathbb Z^2),
\]
the first Chern class of \(\mathcal L_Q|_{F_t}\), in the ordered lattice
basis
\((\widehat\gamma,\widehat u,\widehat w,\widehat\delta)\) indexing the
columns of \([Z\mid I_2]\), is exactly
\[
 c_1(\mathcal L_Q|_{F_t})=-2q,
 \qquad
 q=u\wedge w+6\gamma\wedge\delta.
\tag{5.17aa}
\]
\end{proposition}

\begin{proof}
\emph{Local finiteness and Cartier status.}
The fan of \(Y\) consists of the cones over the cells of the locally finite
\(A_2\)-triangulation at height one.  If
\(U_\sigma\simeq\mathbb C^3\) is the affine chart belonging to a triangle
\(\sigma=\{v_0,v_1,v_2\}\), then \(U_\sigma\) meets precisely
\(E_{v_0},E_{v_1},E_{v_2}\).  Indeed, by the orbit--cone
correspondence, \(U_\sigma\cap E_v\ne\varnothing\) exactly when
\(\mathbb R_{\geq0}(v,1)\) is a face of the cone over \(\sigma\).
Thus \(\{E_v\}_{v\in\mathbb Z^2}\) is locally finite and
\(\mathcal D_Q\) is a well-defined locally finite Weil divisor.

Every triangle of the fan is unimodular, as proved in
\cref{thm:candidate-cusp-geometry}.  If \(z_0,z_1,z_2\) are the
characters dual to
\((v_0,1),(v_1,1),(v_2,1)\), then
\[
 \mathcal D_Q|_{U_\sigma}
 =\operatorname{div}\!
   \left(z_0^{Q(v_0)}z_1^{Q(v_1)}z_2^{Q(v_2)}\right).
\tag{5.17ab}
\]
Negative exponents merely make the displayed monomial meromorphic, as
permitted for a local equation of a Cartier divisor.  Hence
\(\mathcal D_Q\) is Cartier on every chart and therefore on
\(Y_\epsilon\).

\emph{The signed divisor pullback identity.}
For \(\lambda\in\bar\Lambda=\mathbb Z^2\), set
\[
 b_\lambda=B_0\lambda,
 \qquad
 \ell_\lambda=2Hb_\lambda,
 \qquad
 k_\lambda=Q(b_\lambda),
 \qquad
 m_\lambda=(\ell_\lambda,k_\lambda)\in M'.
\tag{5.17ac}
\]
The fan automorphism sends \(E_v\) to \(E_{v+b_\lambda}\).
Consequently \(\Phi_\lambda^*E_w=E_{w-b_\lambda}\), and
\begin{align*}
 \Phi_\lambda^*\mathcal D_Q-\mathcal D_Q
 &=\sum_{v\in\mathbb Z^2}
   \bigl(Q(v+b_\lambda)-Q(v)\bigr)E_v\\
 &=\sum_{v\in\mathbb Z^2}
   \bigl(2v^{\mathsf t}Hb_\lambda+
         b_\lambda^{\mathsf t}Hb_\lambda\bigr)E_v
 =\operatorname{div}(\chi^{m_\lambda}).
\end{align*}
The \(t\)-dependent horizontal torus translation
\(p\mapsto(c_\lambda(t(p)),1)\cdot p\) preserves every
torus-invariant prime divisor \(E_v\).  Hence
\[
 \boxed{\ 
 \Psi_\lambda^*\mathcal D_Q
 =\mathcal D_Q+\operatorname{div}(\chi^{m_\lambda})\ }.
\tag{5.17ad}
\]

\emph{The complete \(\Psi\)-cocycle.}
The matrices in (5.17t) and (5.17x) give
\[
 PC(t)=
 \begin{pmatrix}
 6(b(t)-h(t))&6\mu(t)\\
 6\mu(t)&h(t)
 \end{pmatrix},
\tag{5.17ae}
\]
which is symmetric.  Thus
\[
 A(t):=2PC(t)
\tag{5.17af}
\]
is a symmetric holomorphic matrix on the full cusp disc.  Define the
holomorphic units \(h_\lambda\) and meromorphic functions \(f_\lambda\) by
\[
 h_\lambda(t)=
 \exp\!\left(\pi i\lambda^{\mathsf t}A(t)\lambda\right),
 \qquad
 f_\lambda=h_\lambda(t)\chi^{m_\lambda}.
\tag{5.17ag}
\]
Equation (5.17ad) gives
\[
 \operatorname{div}(f_\lambda)
 =\Psi_\lambda^*\mathcal D_Q-\mathcal D_Q.
\tag{5.17ah}
\]

The cocharacter map underlying \(\Phi_\mu\) is
\[
 \varphi_\mu(y,y_3)=(y+y_3b_\mu,y_3),
\]
so its pullback on characters is
\[
 \varphi_\mu^*(\ell,k)
 =(\ell,k+\ell^{\mathsf t}b_\mu).
\tag{5.17ai}
\]
The quadratic identity for \(Q\) gives
\[
 m_{\lambda+\mu}=m_\mu+\varphi_\mu^*m_\lambda.
\tag{5.17aj}
\]
Indeed, the first coordinates add, while the last coordinate on the
right is
\[
 Q(b_\mu)+Q(b_\lambda)+2b_\lambda^{\mathsf t}Hb_\mu
 =Q(b_\lambda+b_\mu).
\]
Using (5.17u), (5.17v), and (5.17ai), one obtains
\begin{align*}
 \Psi_\mu^*\chi^{m_\lambda}
 &=\chi^{m_\lambda}(c_\mu(t),1)
   \chi^{\varphi_\mu^*m_\lambda}\\
 &=\exp\!\left(2\pi i\ell_\lambda^{\mathsf t}C(t)\mu\right)
   \chi^{\varphi_\mu^*m_\lambda}\\
 &=\exp\!\left(2\pi i\lambda^{\mathsf t}A(t)\mu\right)
   \chi^{\varphi_\mu^*m_\lambda}.
\tag{5.17ak}
\end{align*}
The last equality is exact because
\[
 \ell_\lambda^{\mathsf t}C\mu
 =(2HB_0\lambda)^{\mathsf t}C\mu
 =2\lambda^{\mathsf t}B_0^{\mathsf t}HC\mu
 =2\lambda^{\mathsf t}PC\mu
 =\lambda^{\mathsf t}A\mu.
\]
Since \(A=A^{\mathsf t}\), (5.17ag), (5.17aj), and (5.17ak) imply
\begin{align*}
 f_\mu\,\Psi_\mu^*f_\lambda
 &=h_\mu h_\lambda
   \exp(2\pi i\lambda^{\mathsf t}A\mu)
   \chi^{m_\mu+\varphi_\mu^*m_\lambda}\\
 &=\exp\!\left(
   \pi i(\lambda+\mu)^{\mathsf t}A(\lambda+\mu)\right)
   \chi^{m_{\lambda+\mu}}
 =f_{\lambda+\mu}.
\tag{5.17al}
\end{align*}

Use the divisor-sheaf convention
\[
 \mathcal O(D)(U)=
 \{g\in\mathcal M_Y(U):
   \operatorname{div}(g)+D|_U\geq0\}.
\]
If \(D'=D+\operatorname{div}(f)\), multiplication by \(f\) maps
\(\mathcal O(D')\) isomorphically to \(\mathcal O(D)\).  Thus
(5.17ad) and (5.17ag) define
\[
 \theta_\lambda:
 \Psi_\lambda^*\mathcal O(\mathcal D_Q)
 \xrightarrow{\ \times f_\lambda\ }
 \mathcal O(\mathcal D_Q).
\tag{5.17am}
\]
Because
\((\Psi_\lambda\Psi_\mu)^*=\Psi_\mu^*\Psi_\lambda^*\),
equation (5.17al) is exactly
\[
 \theta_{\lambda+\mu}
 =\theta_\mu\circ\Psi_\mu^*\theta_\lambda.
\tag{5.17an}
\]
This proves a complete equivariance datum, not merely invariance of the
divisor class.

\emph{Analytic descent.}
By \cref{thm:candidate-cusp-geometry}, the action on \(Y_\epsilon\) is
free and properly discontinuous.  Every point therefore has an open
neighbourhood \(U\) whose distinct translates are disjoint and on which
\(\pi|_U\) is a biholomorphism onto \(\pi(U)\).  Transport
\(\mathcal O(\mathcal D_Q)|_U\) to \(\pi(U)\), and on overlaps arising
from translates glue by (5.17am).  Equation (5.17an) is precisely the
triple-overlap condition.  The resulting holomorphic line bundle
\(\mathcal L_Q\) on \(N_0\) satisfies
\[
 \pi^*\mathcal L_Q
 \simeq\mathcal O_{Y_\epsilon}(\mathcal D_Q).
\tag{5.17ao}
\]

\emph{The normalized central restriction.}
Let \(j_0:E_0\hookrightarrow Y_\epsilon\).  Since
\(\nu=\pi\circ j_0\), (5.17ao) gives
\[
 \nu^*(\mathcal L_Q|_W)
 =j_0^*\pi^*\mathcal L_Q
 \simeq\mathcal O_{E_0}(\mathcal D_Q|_{E_0}).
\tag{5.17ap}
\]
The only \(E_v\), \(v\ne0\), meeting \(E_0\) are those indexed by the
six \(a_i\) in the statement, and the intersections are the six toric
\((-1)\)-curves \(C_i\).  Since \(Q(0)=0\), there is no
\(E_0|_{E_0}\) term, so
\[
 \mathcal D_Q|_{E_0}
 =\sum_{i=1}^6Q(a_i)C_i,
 \qquad
 (Q(a_1),\ldots,Q(a_6))=(1,-6,-5,1,-6,-5).
\tag{5.17aq}
\]
The hexagonal intersection table is
\[
 C_i^2=-1,
 \qquad
 C_i\cdot C_{i\pm1}=1,
 \qquad
 C_i\cdot C_j=0
 \quad(j\notin\{i,i\pm1\}),
\tag{5.17ar}
\]
with indices modulo six.  Therefore
\[
 \deg_{C_i}\mathcal O_{E_0}
 \!\left(\sum_jQ(a_j)C_j\right)
 =-Q(a_i)+Q(a_{i-1})+Q(a_{i+1}),
\]
and direct substitution gives
\[
 \begin{array}{c|rrrrrr}
 i&1&2&3&4&5&6\\ \hline
 -Q(a_i)+Q(a_{i-1})+Q(a_{i+1})
 &-12&2&0&-12&2&0.
 \end{array}
\tag{5.17as}
\]
The degree on \(C_1=C_{e_1}\) is \(-12\), proving nontriviality.

\emph{The generic-fibre factor and its exact Chern class.}
Fix \(t\ne0\), choose \(s\) with \(t=e^{2\pi i s}\), and write
\(x=e^{2\pi i\zeta}\), \(\zeta\in\mathbb C^2\).
Equation (5.17u) identifies the deck transformation indexed by
\(\lambda\in\mathbb Z^2\) with
\[
 \zeta\longmapsto\zeta+Z(s)\lambda
\]
modulo the horizontal translations
\(\zeta\mapsto\zeta+n\), \(n\in\mathbb Z^2\).  The matrix identities
\[
 HB_0=P^{\mathsf t},
 \qquad
 B_0^{\mathsf t}HB_0=PB_0
\tag{5.17at}
\]
give
\[
 \chi^{m_\lambda}(x,t)
 =x^{2P^{\mathsf t}\lambda}
  t^{\lambda^{\mathsf t}PB_0\lambda},
\qquad
 h_\lambda(t)
 =\exp(2\pi i\lambda^{\mathsf t}PC(t)\lambda).
\]
Since \(Z=sB_0+C\), it follows that
\[
 f_\lambda(\zeta)
 =\exp\!\left(
 2\pi i\lambda^{\mathsf t}PZ\lambda+
 4\pi i\lambda^{\mathsf t}P\zeta\right).
\tag{5.17au}
\]
The matrix \(PZ=sPB_0+PC\) is symmetric:
\[
 PB_0=\operatorname{diag}(-6,1),
\]
and \(PC\) is symmetric by (5.17ae).

Under the descent convention (5.17am), a pulled-back local section
written as \(\varphi(\zeta)\) in the canonical frame satisfies
\[
 \varphi(\zeta+Z\lambda)
 =f_\lambda(\zeta)^{-1}\varphi(\zeta).
\]
Thus the factor of automorphy for the lattice translation
\(\zeta\mapsto\zeta+Z\lambda+n\) is
\[
 a_{(\lambda,n)}(\zeta)
 =\exp(2\pi i g_{(\lambda,n)}(\zeta)),
\qquad
 g_{(\lambda,n)}(\zeta)
 =-\lambda^{\mathsf t}PZ\lambda
  -2\lambda^{\mathsf t}P\zeta.
\tag{5.17av}
\]
This is independent of \(n\) as a factor of automorphy: replacing
\(\zeta\) by \(\zeta+n\) changes \(g\) by the integer
\(-2\lambda^{\mathsf t}Pn\).

For \(l=(\lambda,n)\), \(l'=(\mu,m)\), the exponential connecting
cocycle is
\begin{align*}
 c(l,l')
 &=g_l(\zeta+Z\mu+m)+g_{l'}(\zeta)-g_{l+l'}(\zeta)\\
 &=-2\lambda^{\mathsf t}Pm\in\mathbb Z,
\tag{5.17aw}
\end{align*}
where the symmetry of \(PZ\) cancels the two mixed \(Z\)-terms.
The corresponding alternating integral form is
\[
 E(l,l')=c(l,l')-c(l',l)
 =-2\lambda^{\mathsf t}Pm+2\mu^{\mathsf t}Pn.
\tag{5.17ax}
\]
This is \(c_1(\mathcal L_Q|_{F_t})\).  In the ordered column basis of
\([Z\mid I_2]\), the cross-block of the class
\(q=u\wedge w+6\gamma\wedge\delta\) from
\((\widehat\gamma,\widehat u)\) to
\((\widehat w,\widehat\delta)\) is
\[
 \begin{pmatrix}
 q(\widehat\gamma,\widehat w)&
 q(\widehat\gamma,\widehat\delta)\\
 q(\widehat u,\widehat w)&
 q(\widehat u,\widehat\delta)
 \end{pmatrix}
 =\begin{pmatrix}0&6\\1&0\end{pmatrix}=P.
\tag{5.17ay}
\]
Hence (5.17ax) is exactly \(-2q\), as asserted.
\end{proof}

\begin{remark}[the exact obstruction to the purely analytic route]
\label{rem:local-cusp-picard-obstruction}
The proposition proves that the local homomorphism
\[
 \operatorname{Pic}(N_0)\longrightarrow\operatorname{Pic}(E_0),
 \qquad
 \mathcal L\longmapsto\nu^*(\mathcal L|_W),
\]
is not zero.  Thus the toric cusp geometry, conductor descent on \(W\),
or local relative-Picard representability cannot alone prove the desired
global vanishing
\[
 \nu^*i^*:\operatorname{Pic}(X)\longrightarrow
 \operatorname{Pic}(E_0).
\]
There is no contradiction with
\cref{prop:leray-edge-restriction,prop:direct-boundary-picard}: the local
bundle above has nonzero generic class \(-2q\), whereas the two global
proofs use the order-three and order-four fillings and their common
boundary relation to exclude every nonzero multiple of \(q\).

More generally, Bingener proves that if a holomorphic map is proper,
flat, and cohomologically flat in dimension zero, then its relative
Picard functor is represented by a \emph{not necessarily separated}
complex space
\cite[Introduction, p.~318; Satz~6.2, pp.~336--337]{Bingener1980}.
The same introduction explains that the exponential-sequence quotient
works directly under topological local triviality and cohomological
flatness in dimension one, but that passage to the quotient fails in
general.  The cusp map of \cref{thm:candidate-cusp-geometry} is not
topologically locally trivial across \(0\).  Bingener's theorem therefore
supplies representability but neither separatedness nor the
component-specialization assertion needed here.

Consequently a purely analytic completion must add global compatibility
with both finite fillings and must control the boundary support function;
neither relative-Picard representability nor conductor descent supplies
those statements.  The later calculation in
\cref{prop:analytic-ah-boundary,lem:analytic-support-specialization,%
thm:analytic-all-picard} now supplies them: it first forces the generic
Appell--Humbert form to vanish using the two finite fillings, and then
proves that its mixed block is the polarization of the three normalized
central degrees.  Thus this proposition remains the sharp cusp-local
obstruction, while the global purely analytic route is complete.
\end{remark}

\subsection{The canonical divisor and the exact anticanonical ring}
\label{sec:canonical-ring-recovery}

The analytic construction admits a further explicit calculation which does
not use its claimed identification with a sphere. Retain the actual proper
map $f:X\to B=\mathbb P^1$, the cusp $p_0$, the finite points $p_1,p_2$,
the reduced finite fibres $S_1,S_2$, and
$f^*(p_1)=3S_1$, $f^*(p_2)=4S_2$. All line-bundle isomorphisms below
refer to this constructed analytic space. No replacement period matrix or
new compactification is introduced.

\begin{theorem}[Canonical bundle and recovery of the original fibration]
\label{thm:canonical-ring-recovery}
The evaluation map
\[
 f^*(f_*\omega_{X/B})\longrightarrow\omega_{X/B}
\]
has zero divisor exactly $2S_2$. Consequently
\[
 \omega_X\simeq f^*\mathcal O_B(-p_2)\otimes\mathcal O_X(2S_2)
       \simeq\mathcal O_X(-2S_2),\qquad
 \omega_X^{-2}\simeq f^*\mathcal O_B(p_2).
 \tag{CR1}
\]
For every integer $k$ the actual direct image of the vertical divisor is
\[
 f_*\mathcal O_X(kS_2)
       =\mathcal O_B\bigl(\lfloor k/4\rfloor p_2\bigr)
 \tag{CR2}
\]
as subsheaves of the meromorphic-function sheaf, using pullback to identify
their sections. In particular, for each integer $m\geq0$,
\[
 h^0(X,\omega_X^{-m})=\lfloor m/2\rfloor+1,
 \qquad h^0(X,\omega_X^m)=0\quad(m>0).
 \tag{CR3}
\]
There is a graded algebra isomorphism
\[
 \mathbb C[U,V]\xrightarrow{\sim}
       \bigoplus_{m\geq0}H^0(X,\omega_X^{-m}),
       \qquad\deg U=1,\quad\deg V=2.
 \tag{CR4}
\]
The degree-two complete linear system has no base points and its map is
the original $f$, followed by a specified projective-coordinate map of $B$.
The degree-one system consists of the single divisor $2S_2$.
\end{theorem}

\begin{proof}
Write $\mathfrak K=f_*\omega_{X/B}$. The map is the natural adjunction
evaluation, not an identification inferred from the ranks of its two
terms. On a smooth torus fibre the invariant top form is nowhere zero.
At the cusp, the original form
$(dx_1/x_1)\wedge(dx_2/x_2)\wedge dt_c$ is a nowhere-zero canonical
frame in each of the unimodular toric charts. Its relative version is
the local generator of $\mathfrak K$ used in
\cref{lem:hodge-line-audit}, so evaluation has no zero there.

For a finite filling retain $t_j=s_j^{m_j}$,
$(m_1,m_2)=(3,4)$ and $(a_1,a_2)=(2,1)$. On its finite \'etale
total-space cover $q_j:J'_j\to N_j$, the actual direct-image generator
pulls back, by (9.5c)--(9.5f), to
\[
 s_j^{m_j-1-a_j}e^{\varphi_j(s_j)}
       (ds_j\wedge e)\otimes(dt_j)^{-1}.
 \tag{CR5}
\]
Here $(dt_j)^{-1}$ is the pulled-back inverse canonical frame on the
base, not division by $m_js_j^{m_j-1}ds_j$. The coefficient
$e^{\varphi_j}$ is a unit. Thus the zero divisor on the cover is
$(m_j-1-a_j)\{s_j=0\}$. Since $q_j$ is \'etale, faithful pullback
of Cartier divisors makes the downstairs divisor
$(m_j-1-a_j)S_j$. The two coefficients are respectively $0$ and $2$.
These neighbourhoods and the smooth locus cover $X$, proving the stated
global zero divisor with no omitted component.

The evaluation therefore identifies
\[
 \omega_{X/B}\simeq f^*\mathfrak K\otimes\mathcal O_X(2S_2).
\]
By \cref{lem:hodge-line-audit}, $\mathfrak K\simeq\mathcal O_B(p_0)$.
In the original affine base coordinate $t$, with
$t(p_1)=0$, $t(p_2)=1$, $t(p_0)=\infty$, retain the explicit choices
\[
 h=\frac1{t-1},\qquad \eta_B=\frac{dt}{(t-1)^2}.
\]
Their divisors are $p_0-p_2$ and $-2p_2$, respectively: at infinity,
with $w=1/t$, they are $w/(1-w)$ and $-dw/(1-w)^2$.
Thus $\eta_B$ specifies $\omega_B\simeq\mathcal O_B(-2p_2)$.
Tensoring with $f^*\omega_B$ gives
\[
 \omega_X\simeq
 \mathcal O_X\bigl(f^*(p_0)-2f^*(p_2)+2S_2\bigr).
\]
Multiplication by $f^*h$, with its displayed divisor, identifies the
base factor $\mathcal O_B(p_0-2p_2)$ with $\mathcal O_B(-p_2)$.
Using the actual multiplicity $f^*(p_2)=4S_2$ proves (CR1).

To prove (CR2), a section of $\mathcal O_X(kS_2)$ is represented by
a meromorphic function on $X$ with the indicated vertical order bound.
Away from $p_2$ it is holomorphic. Its restriction to each connected
compact fibre there is constant; $f_*\mathcal O_X=\mathcal O_B$
therefore descends it uniquely to a holomorphic function on that base
open set. Near $p_2$, pull back to the original cover $J'_2$, including
its zero section. The descended function becomes a function $a(s_2)$,
and $s_2^k a(s_2)$ is holomorphic at zero (with the evident vanishing
interpretation when $k<0$). Thus $a$ is meromorphic and has a Laurent
series bounded below. Invariance under the actual quotient, whose base
coordinate is $s_2\mapsto\zeta_2s_2$, says
$a(\zeta_2s_2)=a(s_2)$. Only exponents divisible by four occur, so
$a(s_2)=b(s_2^4)$ for a meromorphic $b$ on the original base disc.
The order inequality is exactly
\[
 4\operatorname{ord}_{p_2}(b)+k\geq0
 \quad\Longleftrightarrow\quad
 \operatorname{ord}_{p_2}(b)\geq-\lfloor k/4\rfloor.
\]
Conversely every such $b$ pulls back to an allowed meromorphic function.
This proves equality in both directions, as well as compatibility on
overlaps. No assumption about the field of all meromorphic functions
on $X$ is used. Applying (CR2) with $k=2m$ and $k=-2m$ and using
$h^0(\mathbb P^1,\mathcal O(d))=\max\{d+1,0\}$ proves (CR3).

Here is the exact algebra map in (CR4). Fix the isomorphism
$\omega_X^{-1}\simeq\mathcal O_X(2S_2)$ obtained above. Let $u$ be
the section corresponding to the canonical divisor section of
$\mathcal O_X(2S_2)$, so $\operatorname{div}(u)=2S_2$.
Its square, under the induced degree-two identification in (CR1),
is $f^*\sigma_{p_2}$, where $\sigma_{p_2}$ is the canonical divisor
section of $\mathcal O_B(p_2)$. Choose a second section $\ell$ of this
line bundle with $\ell(p_2)\ne0$, retaining that choice, and let $v$
correspond to $f^*\ell$ under the same degree-two identification.
Define $U\mapsto u$, $V\mapsto v$.

For $m=2r$, the monomials $u^{2r-2j}v^j$, $0\leq j\leq r$,
correspond exactly to
$f^*(\sigma_{p_2}^{r-j}\ell^j)$, a basis of
$H^0(B,\mathcal O_B(rp_2))$. For $m=2r+1$, multiplication by $u$
identifies the same basis with the full space: in the meromorphic
description (CR2), both $4rS_2$ and $(4r+2)S_2$ allow precisely
the base pole orders at most $r$. Thus the displayed map is an
isomorphism in every degree, proving both surjectivity and absence of
polynomial relations. Replacing $\ell$ by $a\ell+b\sigma_{p_2}$,
$a\ne0$, changes the second generator by exactly
$v\mapsto av+bu^2$; no scalar is silently discarded.
Replacing the chosen isomorphism $\omega_X^{-1}\to\mathcal O_X(2S_2)$
by $c$ times itself changes the generators by
$(u,v)\mapsto(c^{-1}u,c^{-2}v)$, with $c\in\mathbb C^*$ retained.

The two degree-two generators have no common zero because
$\sigma_{p_2},\ell$ do not. Their map is
$x\mapsto[\sigma_{p_2}(f(x)):\ell(f(x))]$, which is exactly
$f$ followed by that chosen coordinate map $B\to\mathbb P^1$.
Explicitly, writing $\ell=A/(t-1)+B$ with $A\ne0$, it is
$t\mapsto[t-1:A+B(t-1)]$, including the values $[0:A]$ at $p_2$
and $[1:B]$ at $p_0$.
Degree one has basis $u$ and its divisor is $2S_2$, as claimed.
\end{proof}

In the already specified additive Picard coordinate (9.19), with
$L_{12}=f^*\mathcal O_B(1)$, the identities retain the original constants:
\[
 \mathcal O_X(S_1)=L_4,\qquad
 \mathcal O_X(S_2)=L_3,\qquad
 \omega_X=L_{-6},\qquad \omega_X^{-2}=L_{12}.
\]
This last expression uses the previously proved Picard identification;
the divisor, direct-image and ring proofs above do not depend on it.

\subsection{The retained period constant is an effective deformation parameter}
\label{sec:retained-period-deformation}

Keep the functions, lattice marking, base coordinate and local fillings of
\cref{sec:candidate-geometry}. In particular, fix the actual cusp-regular
solution $\beta^{\rm part}$, not a replacement with a cusp pole, and write
\[
 \beta_c=\beta^{\rm part}+c,\qquad
 \mathcal P=\{c\in\mathbb C:\operatorname{Im}c<-M\},
 \quad M=\max_{B\setminus\{p_0\}}D_{\rm part}.
 \tag{PD1}
\]
The maximum and the exact admissible domain were proved in
\cref{prop:period-exact-admissible-constants}. The following calculation
uses the analytic construction, not its integral cohomology or sphere
recognition. Its purpose is to determine whether the freedom in $c$
actually changes the constructed complex manifold.

\begin{theorem}[Exact period deformation]
\label{thm:retained-period-deformation}
For every open disc $Q$ with $\overline Q\subset\mathcal P$, the retained
construction gives a proper holomorphic submersion
\[
 \pi:\mathcal X_Q\longrightarrow Q,
 \qquad F:\mathcal X_Q\longrightarrow B\times Q,
 \qquad \pi=\operatorname{pr}_Q\circ F,
 \tag{PD2}
\]
whose fibre over $c$ is the analytically constructed threefold $X_c$ and
whose restriction there is the original $f_c:X_c\to B$.
The three distinguished values remain the original
$t(p_1)=0$, $t(p_2)=1$, $t(p_0)=\infty$.

Fix a regular base point and one of its original lifts $z$, and fix
$c_*\in Q$. With all functions below evaluated at that $z$, set
\[
 T=\operatorname{Im}\tau>0,\quad m=\operatorname{Im}\mu,\quad
 D=\operatorname{Im}\beta_{c_*}-\frac{6m^2}{T}<0,\quad a=\frac mT,
 \qquad P=\begin{pmatrix}0&0\\-a&1\end{pmatrix}.
 \tag{PD3}
\]
For $\delta=c-c_*$, the exact marked real-linear map from the universal
cover of the fibre at $c_*$ to that at $c$ is
\[
 w=A_\delta\zeta+B_\delta\overline\zeta,\qquad
 A_\delta=I+\frac{\delta}{2iD}P,\qquad
 B_\delta=-\frac{\delta}{2iD}P.
 \tag{PD4}
\]
It maps the original period lattice onto the new period lattice, and
therefore descends to a real-analytic marked diffeomorphism of the two
tori. In the explicit convention that the transported $(0,1)$ tangent
space is spanned by
$\partial_{\overline\zeta_k}+\sum_j\nu^j_k\partial_{\zeta_j}$,
its exact coefficient matrix and its first derivative are
\[
 \nu_\delta=\frac{\delta}{2iD+\delta}P,\qquad
 \dot\nu_z=\frac{1}{2iD}\,
 \partial_{\zeta_2}\otimes
 (d\overline\zeta_2-a\,d\overline\zeta_1).
 \tag{PD5}
\]
The class of $\dot\nu_z$ in
$H^{0,1}(f_{c_*}^{-1}(t(z)),T^{1,0})$ is nonzero.
With the tangent-exact-sequence convention
$\operatorname{KS}(\partial_c)=[\bar\partial h]$ for a smooth
$(1,0)$ lift $h$ of $\partial_c$, the actual fibre class is
$\operatorname{KS}_z(\partial_c)=-[\dot\nu_z]$.
Moreover the first-order deformation of the entire $X_{c_*}$ obtained
from (PD2) is nontrivial: its Kodaira--Spencer class in
$H^1(X_{c_*},T_{X_{c_*}})$ is nonzero. In particular
$h^1(X_{c_*},T_{X_{c_*}})\geq1$.
\end{theorem}

\begin{proof}
\emph{Construction with the parameter retained.}
On an ordinary lifted base open set use the original period map
\[
 \Pi(z,c)=
 \begin{pmatrix}
 6\mu(z)&\tau(z)&1&0\\
 \beta^{\rm part}(z)+c&\mu(z)&0&1
 \end{pmatrix}:
 \Lambda\otimes\mathbb R\longrightarrow\mathbb C^2.
 \tag{PD6}
\]
Its real determinant is $\operatorname{Im}\tau\,D_c$, which never
vanishes on $Q$. Thus $(z,c,u)\mapsto[z,c,\Pi(z,c)u]$ identifies the
quotient, as a real-analytic family, with the original marked real torus
over this base open set times $Q$. The lattice action is holomorphic in
$(z,c,\zeta)$; it is free and properly discontinuous locally uniformly
on compact base-parameter sets, since the inverse real period matrices
are uniformly bounded there. The quotient is therefore a complex
manifold. Projection to $Q$ is a submersion because it is so before
the covering quotient. Its map to the lifted base times $Q$ is proper:
over any compact subset the image of that subset times a fixed closed
lattice parallelotope is the full inverse image.

The deck matrices are the original integral $M_g$ and the original
matrices $R_g(z)$, with
$\Pi(gz,c)=R_g(z)\Pi(z,c)M_g$. The $R_g$ do not depend on $c$.
Consequently the ordinary base quotient is jointly holomorphic.
Proper discontinuity there follows from proper discontinuity on the
base, exactly as in \cref{lem:middle-period-proper}, with compact
parameter sets carried along.

At a finite filling retain $m_1=3$, $m_2=4$,
$s_j(g_jz)=e^{-2\pi i/m_j}s_j(z)$, and the actual logarithmic lift
\[
 (z,c,\zeta)\longmapsto
 \left(g_jz,c,R_{g_j}(z)\zeta+
                  \Pi(g_jz,c)\frac{v_j}{m_j}\right).
 \tag{PD7}
\]
In the real marking this is $(z,c,u)\mapsto
(g_jz,c,A_ju+v_j/m_j)$, independent of $c$.
Every nonidentity power remains free: the invariant integral covector
$\gamma$ takes the translation of its $k$-th power to $k/3$ or $-k/4$,
which is not an integer for $0<k<m_j$.
The finite quotient is therefore a local biholomorphism and is a
submersion to $Q$. The same compact-parallelotope argument proves
properness over $D_j\times Q$. Its map is still exactly
$t_j=s_j^{m_j}$. On the punctured overlap the translating section is
\[
 \frac{\log s_j(z)}{2\pi i}\Pi(z,c)v_j.
 \tag{PD8}
\]
A change of logarithm adds an actual period. It is therefore a
branch-independent holomorphic torus section, jointly in $z,c$.
The calculation (2.9e) applies with $c$ retained and proves the
punctured-overlap conjugacy, not merely equality of the monodromy orders.

For the cusp retain the same infinite unimodular fan and the full matrix
\[
 C(t_c,c)=
 \begin{pmatrix}6\mu(t_c)&h(t_c)\\
 b^{\rm part}(t_c)-h(t_c)+c&\mu(t_c)
 \end{pmatrix},\qquad
 B_0=\begin{pmatrix}0&1\\-1&0\end{pmatrix}.
 \tag{PD9}
\]
Choose a closed cusp disc of radius $r_1$ on which these functions are
holomorphic. Put
\[
 K_Q=\max_{|t_c|\le r_1,\ c\in\overline Q}
          \|-2\pi\operatorname{Im}C(t_c,c)\|,
 \qquad 0<\epsilon<\min(r_1,1),\quad
          |\log\epsilon|>2K_Q.
 \tag{PD10}
\]
The action on $Y_\epsilon\times Q$ is the original one,
\[
 (x,t_c,c)\longmapsto
 \left(e^{2\pi iC(t_c,c)\lambda}
                  t_c^{B_0\lambda}x,t_c,c\right),
 \quad\lambda\in\mathbb Z^2.
 \tag{PD11}
\]
It is holomorphic across the boundary by the same toric automorphism
and nowhere-zero torus multiplier. The estimates in (2.8) now hold
uniformly on $Q$:
\[
 B_{t_c,c}=B_0+
       \frac{-2\pi\operatorname{Im}C(t_c,c)}{\log|t_c|},\quad
 \|B_{t_c,c}-B_0\|<\tfrac12,\quad
 \|B_{t_c,c}\lambda\|>\tfrac12\|\lambda\|\quad(\lambda\ne0).
 \tag{PD12}
\]
They exclude stabilizers on the dense torus and bound the finite set
of translations that can join each pair of bounded charts, uniformly
in $c$. At $t_c=0$ a nonzero translation moves the bounded cell, so
there is no stabilizer there either. A compact set is contained in finitely
many of the open charts $\mathcal U_\sigma(2,\epsilon)\times Q$ from
\cref{lem:toric-chart-control}.
The chart-pair estimate proves proper discontinuity on the whole
product, including sequences tending to its boundary.

For clarity properness is also uniform, not only pointwise in $c$.
For $\epsilon_1<\epsilon$, let $K_0$ be the finite union of closed
unit polydiscs associated to triangles meeting the closed radius
$\sqrt2$ ball, intersected with $|t_c|\le\epsilon_1$, exactly as in
\cref{thm:candidate-cusp-geometry}. For every $c$ and $t_c\ne0$,
choose $\lambda$ so that
$\|B_{t_c,c}^{-1}y+\lambda\|_\infty\le1/2$.
The translated logarithmic position has norm at most
$3/(2\sqrt2)<\sqrt2$, and hence lies in a chart defining $K_0$.
At $t_c=0$, the six closed bidiscs covering $E_0$ give the same
representatives after the exact translation by $B_0v$, with multiplier
$\exp(2\pi iC(0,c)B_0v)$ retained. Therefore, for compact $K\subset Q$,
the image of the compact $K_0\times K$ is the full inverse image of
$\{|t_c|\le\epsilon_1\}\times K$. Every compact set in the cusp
base times $Q$ is closed in a set of this form. This proves properness.
The quotient is a submersion to $Q$ because its covering charts are
open subsets of $Y_\epsilon\times Q$.

Exponentiating the original logarithmic torus coordinates gives the
jointly holomorphic cusp overlap, with period identity
$Z=sB_0+C(t_c,c)$. Glue it and (PD8) to the middle family.
The special base discs are disjoint. Each resulting inverse image of
a member of this base cover is its entire Hausdorff local piece.
Two points with different $(t,c)$ values are separated on the base;
two with the same value lie in one Hausdorff open piece and can be
separated there. This proves Hausdorffness. Local complex charts give
a complex fourfold and the submersion $\pi$.
Properness of $F$ follows by covering any compact set of $B\times Q$
by finitely many relatively compact sets inside these base pieces and
using the just proved local properness. Since $B$ is compact, projection
$B\times Q\to Q$ is proper, proving (PD2). The fibres and overlaps
are precisely those in the original analytic construction.

\emph{The exact variation on a regular fibre.}
Write $u=(u_1,u_2,u_3,u_4)^t$ in the original ordered lattice basis and
$\zeta=\Pi(z,c_*)u$. Taking imaginary parts gives the two exact equations
\[
 \operatorname{Im}\zeta_1=6m u_1+T u_2,\qquad
 \operatorname{Im}\zeta_2=(D+6m^2/T)u_1+m u_2.
 \tag{PD13}
\]
Subtracting $a$ times the first equation from the second yields
\[
 u_1=\frac{\operatorname{Im}\zeta_2-a\operatorname{Im}\zeta_1}{D}.
 \tag{PD14}
\]
In the changed periods the same $u$ maps to
$w=\Pi(z,c)u=\zeta+\delta e_2u_1$. Substitution of
$\operatorname{Im}\zeta=(\zeta-\overline\zeta)/(2i)$ proves (PD4).
This also proves its asserted lattice map and real invertibility:
it is exactly $\Pi(z,c)\Pi(z,c_*)^{-1}$ as a real-linear map.
Its real determinant is $(D+\operatorname{Im}\delta)/D>0$.
Neither the nonholomorphic part nor the original marking is dropped.

Here $P^2=P$, and
\[
 \det A_\delta=1+\frac\delta{2iD},\qquad
 A_\delta^{-1}=I-\frac{\delta}{2iD+\delta}P.
 \tag{PD15}
\]
The denominator can vanish only for $\delta=-2iD$. At that parameter
$D+\operatorname{Im}\delta=-D>0$, which is outside the admissible
domain. Thus this inverse is defined at every $c\in Q$.
Annihilating $w$ by the declared new $(0,1)$ fields gives
$A_\delta\nu_\delta+B_\delta=0$, which proves (PD5), including
its sign. All its coefficients are constant in the fibre coordinates.
Thus its first derivative is a globally defined $\bar\partial$-closed
tangent-valued $(0,1)$ form on the torus.

It cannot be $\bar\partial V$ for a globally defined smooth vector
field $V$. Indeed the invariant holomorphic frame
$\partial_{\zeta_1},\partial_{\zeta_2}$ writes $V$ with periodic
smooth coefficients. The integral of each of their derivatives
$\partial_{\overline\zeta_k}V^j$ against the translation-invariant
probability volume is zero: real constant vector fields preserve this
volume, and integration of their derivatives on the compact torus is
zero. But the coefficient of
$\partial_{\zeta_2}\otimes d\overline\zeta_2$ in $\dot\nu_z$ is
the nonzero constant $1/(2iD)$. Its integral is that same constant.
This contradiction proves nonzero Dolbeault class without a dimension
count or an appeal to numerical period variation.

The conversion to the tangent-exact-sequence convention retains a
minus sign. In the moving coordinates the real-marking lift has
$(1,0)$ part $h=\partial_c+u_1\partial_{w_2}$ at $c_*$.
Since (PD14) gives
\[
 \bar\partial u_1=-\frac{d\overline\zeta_2-a\,d\overline\zeta_1}{2iD},
\]
the connecting class $[\bar\partial h]$ is exactly $-[\dot\nu_z]$.
Both conventions therefore give a nonzero class; they are not identified
without this sign conversion.

\emph{The intrinsic pencil in the parameter family.}
Let $\mathcal S_2\subset\mathcal X_Q$ be the reduced finite divisor
over $p_2\times Q$, so $F^*(p_2\times Q)=4\mathcal S_2$.
The proof of \cref{thm:canonical-ring-recovery} works for relative
forms over $Q$, as can be checked on each original chart.
On the middle family the frame is $d\zeta_1\wedge d\zeta_2$ relative
to $B\times Q$. At the cusp the relative canonical frame over $Q$ is
$(dx_1/x_1)\wedge(dx_2/x_2)\wedge dt_c$.
Extra $dc$ terms in transition differentials vanish in relative forms.
The finite determinant characters $\chi_j$, their correcting units
$e^{\varphi_j(s_j)}$, and the modular function $\Theta$ used in
(9.4)--(9.6) are independent of $c$, since $R_g,\tau,\mu$ are.
Thus their coefficient identities hold jointly over $Q$.
The evaluation exponents remain $m_j-1-a_j=0,2$.
It follows by the actual same transition functions and divisors that
\[
 \omega_{\mathcal X_Q/Q}\simeq\mathcal O_{\mathcal X_Q}(-2\mathcal S_2),
 \qquad
 \omega_{\mathcal X_Q/Q}^{-2}\simeq
          F^*\operatorname{pr}_B^*\mathcal O_B(p_2).
 \tag{PD16}
\]
Here the base identifications retain $h=1/(t-1)$ and
$\eta_B=dt/(t-1)^2$ from (CR1); these functions also do not depend on $c$.
In particular (PD16) is a relative line-bundle isomorphism, not a
collection of separately chosen isomorphisms of its fibres.

We need its complete degree-two system after an infinitesimal base
change, and prove this explicitly. Put
$A=\mathbb C[\epsilon]/(\epsilon^2)$, base-change by
$c=c_*+\epsilon$, write $X_A$ for the resulting deformation and
$L_A=\omega_{X_A/A}^{-2}$. For $L_*=\omega_{X_{c_*}}^{-2}$,
flatness of the deformation and local freeness of $L_A$ give
\[
 0\longrightarrow L_*\xrightarrow{\ \epsilon\ }L_A
        \longrightarrow L_*\longrightarrow0.
 \tag{PD17}
\]
The two sections of $\mathcal O_B(p_2)$ used in (CR4), namely
$\sigma_{p_2}$ and $\ell=A_0/(t-1)+B_0'$ with $A_0\ne0$,
pull back by (PD16) to sections $s_0,s_1$ of $L_A$ whose reductions
are a basis of $H^0(X_{c_*},L_*)$. This basis assertion is exactly
the proved degree-two case of (CR2)--(CR4).
Any section of $L_A$ reduces to a complex linear combination of these
two reductions; subtract that lifted combination. Exactness of (PD17)
then writes the remainder as $\epsilon$ times a section of $L_*$,
again a combination of that basis. Conversely, if an $A$-linear
combination of $s_0,s_1$ is zero, reduce modulo $\epsilon$ and then
use the injective first arrow of (PD17) to see that all four complex
coefficients vanish. Therefore
\[
 H^0(X_A,L_A)=A s_0\oplus A s_1.
 \tag{PD18}
\]
The system is basepoint-free and its morphism is $f_A$ followed by
the fixed coordinate map $t\mapsto[t-1:A_0+B_0'(t-1)]$.
This proves the required base-change statement without presuming that
pointwise dimensions alone imply it.

\emph{Nontriviality for the complete threefold.}
Suppose its first-order deformation were trivial, with an isomorphism
$X_A\simeq X_{c_*}\times\operatorname{Spec}A$ reducing to the identity.
Intrinsic relative canonical bundles and their complete spaces of
sections are preserved by that isomorphism. By (PD18), it would
therefore identify $f_A$ with $a_A\circ f_{c_*}$ for a projective
linear change $a_A\in\operatorname{PGL}_2(A)$ reducing to the identity.

This change fixes all three distinguished points even to first order.
Here is a scheme-sensitive check. In the finite covering charts
$f_A$ is $t_j=s_j^{m_j}$, so its relative critical ideal is
$(s_j^{m_j-1})$, and its restriction to that critical scheme has
$t_j=0$. The map from the parameter ring $A$ to this critical
scheme is injective; it contains a constant-parameter section given
by $s_j=0$ and the zero section of the covering torus.
At the cusp, in a chart containing a triple point,
$t_c=z_0z_1z_2$ and the critical ideal is
$(z_1z_2,z_0z_2,z_0z_1)$. Again $t_c=0$ in its ring and $A$ injects,
using $z_0=z_1=z_2=0$. At a double chart the same assertion follows
from $t_c=z_0z_1$; a smooth chart has no critical point.
These critical ideals are intrinsic to the differential of the map
over $A$, and the charts are local biholomorphisms onto $X_A$.
Away from the three distinguished base values the map is a
submersion, so it has no critical point.
Consequently the scheme-theoretic critical-value sections are exactly
the three constant sections $p_1,p_2,p_0$. An automorphism reducing
to the identity cannot permute these three distinct closed points.
Its value at each is therefore the same point over $A$, not a
nonzero infinitesimal displacement: such a displacement would be a
nonzero element of $\epsilon\mathbb C$ in an injective copy of $A$
on the just displayed critical chart.

For completeness write $a_A(t)=t+\epsilon(b_0+b_1t+b_2t^2)$.
Fixing $0$ gives $b_0=0$, fixing $1$ gives $b_0+b_1+b_2=0$,
and fixing infinity, in $u=1/t$, gives $b_2=0$.
Hence $a_A$ is the identity. The hypothesized trivialization thus
preserves $f$, and restricting it to any regular base value would
trivialize the torus deformation (PD4). Its nonzero class (PD5)
excludes this.

Finally, the zero Kodaira--Spencer class is precisely the condition for
such a first-order trivialization. One may see the needed direction
directly: choose local product charts for the holomorphic submersion
$\pi$ and differentiate their transition functions at $c_*$. The
resulting tangent-valued overlap functions form a Cech $1$-cocycle.
If its class vanishes, after refinement it is the coboundary of local
holomorphic vector fields. Changing each chart by that vector field
times $\epsilon$ cancels the first-order overlap terms, producing
the identity-reducing trivialization. This is the usual
Kodaira--Spencer connecting class (for the standard deformation-space
interpretation see \cite[Section~67, p.~15]{McMullen1996Complex});
the explicitly recorded conversion
from the plus-graph convention in (PD5) changes its sign, not its
vanishing. The contradiction proves the final assertion.
\end{proof}

The theorem establishes a genuinely varying analytic family with its
original marked fibration and its intrinsic degree-two canonical pencil.
It does not assert that every two distinct parameters give nonisomorphic
unmarked manifolds, or that this family exhausts their deformation space.
The global integral topology and the CDP argument are not inputs to this
deformation calculation.


The global identifications within this retained family are computed in
\cref{pshift:thm:classification}: they are precisely $c'-c\in2\mathbb Z$,
including the finite-filling obstruction and every cusp overlap.


% New proof fragment. No canonical source has been edited by this lane.
% Uses the actual period construction, finite normalization certificate,
% toric fan, and analytically proved canonical-pencil recovery.
\subsection{Exact global identifications of the retained period parameter}
\label{pshift:sec:classification}

Keep the original base $B$, coordinate $t$, values
$t(p_1)=0,t(p_2)=1,t(p_0)=\infty$, and the original functions
$\tau,\mu,\beta^{\mathrm{part}}$. Write
\[
 \beta_c=\beta^{\mathrm{part}}+c,\qquad
 \mathcal P=\{c\in\mathbb C:\operatorname{Im}c<-M\},\qquad
 \Pi_c(z)=\begin{pmatrix}6\mu(z)&\tau(z)&1&0\\
                      \beta_c(z)&\mu(z)&0&1\end{pmatrix}.
 \tag{PS1}\label{pshift:eq:periods}
\]
All assertions below concern the analytic pieces and their actual gluing.
No integral cohomology, sphere recognition, or CDP assertion is an input.
In particular, the global parameter classification is not inferred from
an equality between two fibre lattices.

Put $e=(0,1)^{\mathsf t}\in\mathbb C^2$ and retain the ordered lattice
\[
 \Lambda=\mathbb Z\langle\widehat\gamma,\widehat u,
                        \widehat w,\widehat\delta\rangle,
 \qquad E=\widehat\delta\otimes\gamma.
 \tag{PS2}\label{pshift:eq:rank-one-map}
\]
Thus $E$ is the matrix with its only nonzero entry, equal to one, in
row four and column one; $E^2=0$ and $\Pi_cE=e\gamma$.
The finite data, with every sign retained, are
\[
\begin{gathered}
 A_1=\begin{pmatrix}1&0&0&0\\6&0&1&0\\-6&-1&-1&0\\-2&1&0&1\end{pmatrix},
 \qquad
 A_2=\begin{pmatrix}1&0&0&0\\0&0&-1&0\\-6&1&0&0\\3&0&1&1\end{pmatrix},\\
 v_1=(1,2,-4,0)^{\mathsf t},\quad
 v_2=(-1,-3,3,0)^{\mathsf t},\quad
 (m_1,m_2)=(3,4),\quad
 (\epsilon_1,\epsilon_2)=(1,-1),\quad \epsilon_j=\gamma(v_j).
\end{gathered}
\tag{PS3}\label{pshift:eq:finite-data}
\]
The covariance is $R_g\Pi_c=\Pi_c(gz)A_g$, with
\[
 R_{g_1}=\begin{pmatrix}-1/\tau&0\\(1-\mu)/\tau&1\end{pmatrix},\qquad
 R_{g_2}=\begin{pmatrix}1/\tau&0\\-\mu/\tau&1\end{pmatrix},\qquad
 R_{g_0}=I_2.
 \tag{PS4}\label{pshift:eq:covariance}
\]
In particular $R_ge=e$ for every $g$.

\begin{theorem}[The exact identification subgroup is $2\mathbb Z$]
\label{pshift:thm:classification}
For $c,c'\in\mathcal P$, the retained analytic threefolds are
biholomorphic if and only if
\[
                         c'-c\in2\mathbb Z.
 \tag{PS5}\label{pshift:eq:classification}
\]
Every such biholomorphism preserves the original map to $B$. For every
even integer $n$, the proof below constructs an isomorphism
$X_c\longrightarrow X_{c+n}$ whose complex-linear part on the original
universal fibre coordinates is $I_2$. Its map on the marked fibre lattice
is $I-nE$.

The equality of the middle lattices for all integer shifts loses a
nonzero order-two finite-filling invariant: the shift $n$ changes that
invariant by $n\bmod2$. The cusp contributes no additional obstruction
to the displayed even-shift construction.
\end{theorem}

\subsubsection{The middle family and its complete rational centralizer}

For an integer $n$, put $S_n=I+nE$. Since $S_n^{-1}=S_{-n}$,
\[
 \Pi_{c+n}=\Pi_c S_n,\qquad
 \Pi_{c+n}\Lambda=\Pi_c\Lambda.
 \tag{PS6}\label{pshift:eq:integer-shear}
\]
Both original matrices fix $\widehat\delta$ and preserve $\gamma$.
Consequently $A_jE=EA_j=E$, and the same commutation holds for every
monodromy word. The identity of $\mathbb C^2$ therefore identifies
the two full middle torus families with their actual $R_g$-action.
Its source-to-target lattice map is $S_{-n}$, since
$\Pi_{c+n}S_{-n}=\Pi_c$.

\begin{lemma}[All middle-family linear parts]
\label{pshift:lem:centralizer}
The simultaneous rational centralizer of $A_1,A_2$ is exactly
\[
 \{aI+bE:a,b\in\mathbb Q\}.
 \tag{PS7}\label{pshift:eq:centralizer}
\]
A biholomorphism between the middle families for $c,c'$ over the
identity of $B^\circ$ forces $c'-c\in\mathbb Z$. Its complex-linear
part on every fibre is either $I_2$ or $-I_2$.
\end{lemma}

\begin{proof}
Solving the two fixed-vector systems from (PS3) gives
\[
 (\Lambda\otimes\mathbb Q)^{A_1}
   =\mathbb Q(1,2,-4,0)^{\mathsf t}+\mathbb Q\widehat\delta,
 \quad
 (\Lambda\otimes\mathbb Q)^{A_2}
   =\mathbb Q(1,3,-3,0)^{\mathsf t}+\mathbb Q\widehat\delta.
\]
Their intersection is $\mathbb Q\widehat\delta$. The fixed covectors
are respectively
\[
 \mathbb Q\gamma+\mathbb Q(2u+w+3\delta),\qquad
 \mathbb Q\gamma+\mathbb Q(u+w+2\delta).
\]
Their intersection is $\mathbb Q\gamma$: comparing the $u$ and $w$
coefficients of a common multiple gives $2x=y$ and $x=y$, hence
$x=y=0$. Every commuting endomorphism therefore preserves
$\mathbb Q\widehat\delta$ and $\ker\gamma$.

On $\ker\gamma/\mathbb Q\widehat\delta$, in the classes of
$(\widehat u,\widehat w)$, the two matrices are
\[
 Q_1=\begin{pmatrix}0&1\\-1&-1\end{pmatrix},\qquad
 Q_2=\begin{pmatrix}0&-1\\1&0\end{pmatrix}.
\]
A matrix commuting with $Q_2$ has the form
$\left(\begin{smallmatrix}b&d\\-d&b\end{smallmatrix}\right)$.
Commutation with $Q_1$ gives $d=0$. Thus a commuting endomorphism $H$
has columns
\[
 H\widehat\gamma=(a,r,s,k)^{\mathsf t},\quad
 H\widehat u=b\widehat u+u_0\widehat\delta,\quad
 H\widehat w=b\widehat w+w_0\widehat\delta,\quad
 H\widehat\delta=d_0\widehat\delta.
\]
Use the actual columns $A_1\widehat u=-\widehat w+\widehat\delta$,
$A_1\widehat w=\widehat u-\widehat w$, and
$A_2\widehat u=\widehat w$. Commutation on these columns gives
\[
 d_0-w_0=b+u_0,\qquad u_0=2w_0,\qquad w_0=u_0.
\]
Hence $u_0=w_0=0$ and $d_0=b$. Commutation on
$\widehat\gamma$ with $A_1$ gives
\[
 r=2(a-b),\qquad s=-4(a-b).
\]
Commutation on that column with $A_2$ also gives $r=-s$.
Therefore $a=b$, $r=s=0$, and $H=aI+kE$. Conversely these matrices
commute, by the fixed-vector and fixed-covector identities preceding
(PS6). This proves (PS7).

For the geometric assertion, a holomorphic map between two compact
complex tori lifts to their universal covers. The derivative of the
lift is periodic, so its entries descend to holomorphic functions on
the compact source torus and are constant. The lift is consequently
affine, with a complex-linear part $L$. In a holomorphic family these
linear parts vary holomorphically. Their lattice maps are integral
and locally constant, hence give an isomorphism of the two marked
local systems. Because the map is over the identity of the base,
that integral map $H$ commutes with both $A_j$. Thus
$H=aI+bE$, with $a,b\in\mathbb Z$. Its determinant is $a^4$;
invertibility over $\mathbb Z$ gives $a=1$ or $a=-1$.

The complete period equation is
\[
                   L(z)\Pi_c(z)=\Pi_{c'}(z)(aI+bE).
 \tag{PS8}\label{pshift:eq:period-morphism}
\]
Its third and fourth columns give $L(z)=aI_2$. Its first column
then gives $a(c'-c)+b=0$. Therefore $c'-c=-b/a\in\mathbb Z$,
and the asserted two linear parts exhaust the possibilities.
\end{proof}

\subsubsection{The exact finite obstruction, including normalization}

The two cyclic norm matrices are, by multiplication of (PS3),
\[
 N_1=I+A_1+A_1^2=
 \begin{pmatrix}3&0&0&0\\6&0&0&0\\-12&0&0&0\\0&2&1&3\end{pmatrix},
\]
\[
 N_2=I+A_2+A_2^2+A_2^3=
 \begin{pmatrix}4&0&0&0\\12&0&0&0\\-12&0&0&0\\0&2&2&4\end{pmatrix}.
 \tag{PS9}\label{pshift:eq:norms}
\]
For example these formulas give
$N_1\widehat w=\widehat\delta$,
$N_2\widehat u=N_2\widehat w=2\widehat\delta$,
$N_1\widehat\delta=3\widehat\delta$, and
$N_2\widehat\delta=4\widehat\delta$.

We first justify why a local isomorphism must satisfy the exact affine
equations on the original finite cover. The actual ramified base change
$t_j=u^{m_j}$ has local equation
\[
 s_j^{m_j}-u^{m_j}
   =\prod_{a=0}^{m_j-1}(s_j-e^{2\pi ia/m_j}u)=0.
\]
Its normalization separates all these smooth branches, each with
coordinates $(z_1,z_2,u)$. The original map
$y\mapsto(q_j(y),s_j(y))$ identifies the original torus family $J'_j$
with this normalization, as proved by the explicit branches in
\cref{ffcert:thm:normalization}. A biholomorphism over the identity of
$D_j$ base-changes over the identity of the $u$-disc and lifts uniquely
to the normalization: on local rings it carries integral elements in
the total fraction ring to integral elements, and its inverse does
the same. This lift commutes with the original deck action because
both composites lift the same map $(x,u)\mapsto(F(x),\zeta_j u)$.
In particular one may not replace the generator by an arbitrary
coprime power.

Let $c'=c+n$, where $n$ is an integer, and suppose the linear part of
the lifted map is $aI_2$, $a\in\{1,-1\}$. The lifted map on the torus
family is $[z,\zeta]\mapsto[z,a\zeta+B_j(z)]$. Its translating
section has a holomorphic lift $B_j:\Delta_j\to\mathbb C^2$:
$\Delta_j\times\mathbb C^2\to J'_j$ is a covering map and the disc is
simply connected. Local inverses of this covering show that the lift
is holomorphic. The deck-equivariance identity holds modulo an
actual lattice vector, so, after choosing this lift, it is
\[
 B_j(g_jz)-R_{g_j}(z)B_j(z)
  =\Pi_c(g_jz)\left(
       \frac{(1-a)v_j+n\epsilon_j\widehat\delta}{m_j}
          +\ell_j\right),\qquad \ell_j\in\Lambda.
 \tag{PS10}\label{pshift:eq:finite-equation}
\]
The integer column $\ell_j$ is constant: it is locally constant by
discreteness of the marked lattice and the disc is connected.
Iterating $m_j$ times, with $A_jv_j=v_j$ and
$A_j\widehat\delta=\widehat\delta$, yields the necessary equality
\[
 N_j\ell_j=-(1-a)v_j-n\epsilon_j\widehat\delta.
 \tag{PS11}\label{pshift:eq:finite-norm-obstruction}
\]
Here the left side of the iterated equation is zero because the
base point and the original vector coordinate return after $m_j$
steps; real invertibility of $\Pi_c$ makes the lattice equality exact.

If $a=-1$, applying $\gamma$ to (PS11) gives
\[
                   m_j\gamma(\ell_j)=-2\epsilon_j.
\]
This is impossible already at $j=1$, since $3$ does not divide $2$.
Thus a global isomorphism cannot have negative linear part.

For $a=1$, the first row of (PS11) gives $\gamma(\ell_j)=0$.
Equation (PS9) then proves that the order-three filling imposes no
restriction on $n$, whereas the order-four filling imposes precisely
$n\in2\mathbb Z$. More intrinsically, the finite obstruction is the
class of $-n\epsilon_j\widehat\delta$ in
\[
 \Lambda^{A_j}/N_j\Lambda.
 \tag{PS12}\label{pshift:eq:finite-class}
\]
In the original fixed bases this quotient is
$\mathbb Z/3$ at $j=1$, with $\widehat\delta$ zero, and
$\mathbb Z/4\oplus\mathbb Z/2$ at $j=2$, with
$\widehat\delta$ generating the second summand. Indeed (PS9) gives
\[
 N_1\Lambda=3\mathbb Z v_1\oplus\mathbb Z\widehat\delta,
 \quad
 N_2\Lambda=4\mathbb Z\varepsilon'\oplus
                      2\mathbb Z\widehat\delta,
 \qquad \varepsilon'=(1,3,-3,0)^{\mathsf t}.
\]
This proves the exact parity obstruction rather than merely failure
of the untranslated identity.

\subsubsection{Local maps for every even shift}

Let $n\in2\mathbb Z$. Define the following original rational columns,
integral cocycles, and holomorphic translation vectors:
\[
\begin{aligned}
 r_1&=-\frac n3(\widehat u+\widehat w),&
 \ell_1&=-n\widehat w,&
 b_1(z)&=\Pi_c(z)r_1=-\frac n3(\tau(z)+1,\mu(z))^{\mathsf t},\\
 r_2&= \frac n4(\widehat u+\widehat w),&
 \ell_2&=\frac n2\widehat u,&
 b_2(z)&=\Pi_c(z)r_2=\frac n4(\tau(z)+1,\mu(z))^{\mathsf t}.
\end{aligned}
\tag{PS13}\label{pshift:eq:local-corrections}
\]
Direct multiplication of the complete matrices in (PS3) gives
\[
 (I-A_j)r_j=\frac{n\epsilon_j}{m_j}\widehat\delta+\ell_j.
 \tag{PS14}\label{pshift:eq:local-matrix-identity}
\]
For $j=1$ the four entries are $(0,0,-n,n/3)$, and for $j=2$
they are $(0,n/2,0,-n/4)$. Pairing with the full $\Pi_c(g_jz)$
shows that $b_j$ satisfies (PS10) for $a=1$.
Translation by $b_j$ is therefore already an isomorphism of the
two finite local fillings. The remaining issue is its exact overlap
with the fixed middle identification.

Keep the original linearizing coordinate, radius, and phase
$s_j(g_jz)=e^{-2\pi i/m_j}s_j(z)$. On a logarithm branch put
\[
 w_j(z)=\frac{\log s_j(z)}{2\pi i},\qquad
 a_j(z)=b_j(z)+n\epsilon_j w_j(z)e.
 \tag{PS15}\label{pshift:eq:finite-middle-sections}
\]
Changing the logarithm adds $n\epsilon_j\widehat\delta$ under the
period map, so $a_j$ is branch independent as a torus section.
For the specified negative generator, $w_j(g_jz)=w_j(z)-1/m_j$.
Equations (PS4) and (PS14) consequently give the exact equality
\[
 a_j(g_jz)=R_{g_j}(z)a_j(z)+\Pi_c(g_jz)\ell_j.
 \tag{PS16}\label{pshift:eq:local-affine-cocycle}
\]
This is why the logarithmic term in (PS15) cannot be omitted.
In fact the original finite overlap maps are translations by
$\sigma_{j,c}=w_j\Pi_c v_j$, and
\[
                  \sigma_{j,c+n}-\sigma_{j,c}
                     =n\epsilon_jw_je.
 \tag{PS17}\label{pshift:eq:actual-overlap-difference}
\]
Thus $a_j$ is exactly the middle translation induced by the local
translation $b_j$, with the original gluing maps retained.

At the cusp retain $s=\tau-h$, $e^{2\pi is}=t_c$, $s(g_0z)=s(z)-1$,
and the complete period matrix
\[
 Z_c=sB_0+C_c(t_c),\quad
 B_0=\begin{pmatrix}0&1\\-1&0\end{pmatrix},\quad
 C_c=\begin{pmatrix}6\mu&h\\ b^{\mathrm{part}}-h+c&\mu\end{pmatrix}.
 \tag{PS18}\label{pshift:eq:cusp-periods}
\]
For $\lambda=(\lambda_1,\lambda_2)^{\mathsf t}\in\mathbb Z^2$,
$C_{c+n}\lambda-C_c\lambda=(0,n\lambda_1)^{\mathsf t}$. Therefore
\[
 \exp(2\pi iC_{c+n}\lambda)t_c^{B_0\lambda}
       =\exp(2\pi iC_c\lambda)t_c^{B_0\lambda}.
 \tag{PS19}\label{pshift:eq:cusp-actions-identical}
\]
The entire cusp actions, including their nonconstant holomorphic units,
are identical for an integer shift.

The affine cocycle convention in (PS16) has multiplication law
\[
                    \ell_{gh}=\ell_g+A_g\ell_h.
 \tag{PS20}\label{pshift:eq:cocycle-law}
\]
Since $g_0=(g_1g_2)^{-1}$, it gives
\[
\begin{aligned}
 \ell_1+A_1\ell_2
   &=-\frac{3n}{2}\widehat w+\frac n2\widehat\delta,\\
 \ell_0=-A_{g_0}(\ell_1+A_1\ell_2)
   &=\frac{3n}{2}\widehat w-\frac n2\widehat\delta.
\end{aligned}
\tag{PS21}\label{pshift:eq:cusp-cocycle}
\]
Here $A_{g_0}$ fixes both toric lattice vectors, as is also immediate
from $A_{g_0}=\left(\begin{smallmatrix}I&0\\B_0&I\end{smallmatrix}\right)$.
Put
\[
 q_n=(-3n/2,n/2)^{\mathsf t}\in\mathbb Z^2,
                  \qquad a_0(z)=s(z)q_n.
 \tag{PS22}\label{pshift:eq:cusp-translation}
\]
Then $a_0(g_0z)-a_0(z)=-q_n=\Pi_c(g_0z)\ell_0$, exactly as
required by (PS16). Exponentiating gives the map
$(x,t_c)\mapsto(t_c^{q_n}x,t_c)$ on the dense cusp torus.
It extends holomorphically and invertibly to the original infinite fan:
the integral map
\[
 (y,y_3)\longmapsto(y+y_3q_n,y_3)
\]
translates each height-one triangle by the integer vector $q_n$ and
therefore permutes the actual cones, with inverse given by $-q_n$.
It commutes with the cusp quotient action in (PS19). This proves
extension across every cusp stratum, not only across its dense torus.

\subsubsection{The global additive correction and all overlap identities}

We give the global correction explicitly as an additive gluing
construction and prove its solvability. Let $\widetilde B^\circ$ be
the universal cover of the three-punctured base. Its fundamental group
is free on the original two negative peripheral generators. Pull back
the original $\Pi_c,R$ to this cover. Assign $\ell_1,\ell_2$ from
(PS13) to its free generators and extend by (PS20). Freeness makes
this an actual integral cocycle; its third peripheral value is
precisely (PS21). The maps
\[
 (z,y)\longmapsto
       (gz,R_g(z)y+\Pi_c(gz)\ell_g)
 \tag{PS23}\label{pshift:eq:affine-bundle-action}
\]
form an affine action. To check the action law, covariance gives
$R_g(hz)\Pi_c(hz)=\Pi_c(ghz)A_g$; substituting this equality
in the composition gives exactly (PS20).
The quotient is an affine holomorphic bundle over $B^\circ$, since
the deck action on the base is free and properly discontinuous and
one may use evenly covered base opens as affine charts.

On a finite punctured disc, (PS16) says that $a_j$ is a section of
this affine bundle. On the cusp punctured disc, (PS21)--(PS22) say
the same for $a_0$. Define a sheaf $\mathscr H_n$ on $B$ as follows.
Away from the special points it is the sheaf of sections of (PS23).
Near $p_j$, a section is required to have the form $a_j+\eta_j$,
where $\eta_j:\Delta_j\to\mathbb C^2$ is holomorphic and obeys
$\eta_j(g_jz)=R_{g_j}(z)\eta_j(z)$. Near the cusp it is required
to have the form $a_0+\eta_0$, with $\eta_0(t_c)$ holomorphic on
the full cusp disc. These are compatible definitions on the punctured
overlaps by the equal affine cocycles just proved. They also prove
local nonemptiness, including at all three special points.

The differences of these local sections form a vector sheaf
$\mathscr V$. Its ordinary sections have covariance $\eta(gz)=R_g\eta(z)$;
at the two finite points and the cusp its extensions are precisely
the holomorphic extensions specified above. Projection to the first
coordinate gives an exact sequence on the original compact base
\[
 0\longrightarrow\mathcal O_B e\longrightarrow\mathscr V
        \longrightarrow\mathcal L\longrightarrow0,
                  \qquad \mathcal L\simeq\mathcal O_B(-p_0).
 \tag{PS24}\label{pshift:eq:vector-sheaf}
\]
Here $\mathcal L$ is exactly the homogeneous period sheaf from
\cref{prop:period-input-audit}, not a newly chosen extension. Indeed
the first diagonal entries in (PS4) are its factors $-1/\tau$ and
$1/\tau$, and at the cusp both extensions require ordinary
holomorphicity. The original function
\[
 \Theta=\frac{(E_4\circ\tau)^2(E_6\circ\tau)^{1/2}}
                 {\Delta_{\mathrm{mod}}\circ\tau}
\]
has those exact factors, orders two and one in $s_1,s_2$, and a
simple pole in $t_c$. Dividing by $\Theta$ identifies the ordinary
and finite stalks with $\mathcal O_B$, and the cusp stalk with
$\mathfrak m_{p_0}$. This recovers the displayed $\mathcal L$
with its original cusp condition. These transformation and order
calculations are the ones used to construct the retained $\mu$;
no topological conclusion about $X_c$ is involved.

For completeness, exactness on the finite stalks in (PS24) retains
the off-diagonal entry of $R_g$. Given an equivariant first component
$u(z)$, begin with the vector $(u(z),0)^{\mathsf t}$ and average it
under the finite group action on vector-valued holomorphic functions.
The resulting vector is equivariant, and its first component is
still $u$. Thus projection is surjective there. Its kernel is
exactly the invariant second-component functions, which are
holomorphic functions of $s_j^{m_j}=t_j$, multiplied by $e$.
The ordinary and cusp stalks are immediate. This also proves that
$\mathscr V$ is locally free of rank two: lift a local frame of
$\mathcal L$ and join it to $e$.

Both $H^1(B,\mathcal O_B)$ and
$H^1(B,\mathcal O_B(-p_0))$ vanish. One can compute these two
groups without a degree approximation: on the two standard affine
charts of $\mathbb P^1$, split a holomorphic Laurent series on
their overlap into its nonnegative and negative powers. For
$\mathcal O_B$ these are the two coboundary contributions, and
for $\mathcal O_B(-p_0)$ its transition factor shifts the second
contribution by one, still leaving no missing exponent. The positive
part converges on the full first affine chart and the transformed
negative part on the full second affine chart. This proves the two
vanishings, equivalently the usual additive Cousin calculation on
these charts. The exact sequence (PS24) now gives
\[
                             H^1(B,\mathscr V)=0.
 \tag{PS25}\label{pshift:eq:no-analytic-obstruction}
\]

Here is how to choose the correction entering the maps. Choose
local sections $h_\alpha$ of $\mathscr H_n$, taking the displayed
$a_j,a_0$ on sufficiently small special discs and affine charts
elsewhere. Their differences $d_{\alpha\beta}=h_\beta-h_\alpha$
are a Cech $1$-cocycle with values in $\mathscr V$. Project it to
$\mathcal L$ and solve that additive cocycle, using its just proved
$H^1$ vanishing. Lift the resulting local $\mathcal L$ sections to
$\mathscr V$ by the local lifting construction in (PS24). Subtracting
their coboundary from $d$ leaves a scalar $\mathcal O_Be$ cocycle.
Solve that scalar additive cocycle by its $H^1$ vanishing. This
produces sections $v_\alpha$ with
$d_{\alpha\beta}=v_\beta-v_\alpha$. The local sections
$h_\alpha-v_\alpha$ agree and define a global section
$a\in\mathscr H_n(B)$. This is a two-stage construction of the
required correction, not an unproved additional gluing condition.
It gives holomorphic vectors $\eta_j,\eta_0$ with
\[
 a=a_j+\eta_j\quad(j=1,2),\qquad a=a_0+\eta_0
                              \quad\hbox{at the cusp}.
 \tag{PS26}\label{pshift:eq:global-correction}
\]
Any choices made in the additive solutions give maps with the same
asserted properties. Their differences as sections of this affine
bundle lie in $H^0(B,\mathscr V)=\mathbb C e$: this follows from
(PS24), $H^0(B,\mathcal L)=0$, and $H^0(B,\mathcal O_B)=\mathbb C$.

Projecting $a$ modulo $\Pi_c\Lambda$ gives a holomorphic section of
the middle torus family because every $\ell_g$ in (PS23) is integral.
The promised three kinds of local maps are therefore
\[
\begin{aligned}
 F_{\mathrm{mid}}[z,\zeta]&=[z,\zeta+a(z)],\\
 F_j[z,\zeta]&=[z,\zeta+b_j(z)+\eta_j(z)]\quad(j=1,2),\\
 F_0[x,t_c]&=[\exp(2\pi i\eta_0(t_c))t_c^{q_n}x,t_c].
\end{aligned}
\tag{PS27}\label{pshift:eq:full-isomorphism}
\]
The first descends by (PS23) and equality of the integer-shift
lattices. The finite maps descend by (PS10), (PS14), and the exact
covariance of $\eta_j$. The cusp map descends by (PS19); its
holomorphic-unit factor extends by torus multiplication on every
original chart, and its integer factor extends by the fan map in
(PS22). Each map is a biholomorphism, with inverse given by the
negative translation or the inverse torus multiplier and fan map.

On a finite overlap, write the original gluing as
$h_{j,c}(\zeta)=\zeta+\sigma_{j,c}$. Then (PS17), (PS26), and
(PS27) give the equality of actual torus maps
\[
                     h_{j,c+n}F_j=F_{\mathrm{mid}}h_{j,c}.
 \tag{PS28}\label{pshift:eq:finite-overlap-identity}
\]
Both sides add
$\sigma_{j,c}+a=\sigma_{j,c+n}+b_j+\eta_j$.
On the cusp overlap the original map is exponentiation of $\zeta$;
(PS22) and (PS26) give
$\exp(2\pi ia)=t_c^{q_n}\exp(2\pi i\eta_0)$, exactly the
third line of (PS27). Thus that overlap also agrees, including
all nonconstant period units. The maps and their inverses glue to
a global biholomorphism $X_c\to X_{c+n}$ over $B$.

The choices can be made coherently for every parameter and every even
shift. Both initial cocycle columns in (PS13) have $\gamma=0$.
Because every $A_g$ preserves $\gamma$, (PS20) shows that
$\gamma(\ell_g)=0$ for every word. Therefore the affine term
$\Pi_c(gz)\ell_g$ in (PS23) is independent of $c$. The displayed
$a_j,a_0$, and the vector sheaf $\mathscr V$ are also independent
of $c$. Perform the additive construction once for $n=2$, keeping
one resulting section $a^{(2)}$. For $n=2k$ choose
$a^{(2k)}=k a^{(2)}$ and multiply its local correcting vectors by
$k$. All their affine cocycles and local formulas are then exactly
those for $n=2k$. The resulting maps satisfy the full group law
\[
 F_{2h,c+2k}\circ F_{2k,c}=F_{2(h+k),c},\qquad
                 F_{0,c}=\mathrm{id}_{X_c}.
 \tag{PS28a}\label{pshift:eq:coherent-shift-action}
\]
This follows by adding the actual translation vectors in the first
two lines of (PS27), and by multiplying their units and adding the
integer exponents in its third line. The formulas are jointly
holomorphic in $c$, since their corrections are independent of it.

\subsubsection{Intrinsic necessity and the resulting parameter space}

The analytically computed complete linear system
$|\omega_{X_c}^{-2}|$ is the original map $f_c$, by
\cref{thm:canonical-ring-recovery}. A biholomorphism preserves its
canonical bundle, the complete space of sections, and its evaluation
map. Consequently it induces $\alpha\in\operatorname{PGL}_2(\mathbb C)$
with $f_{c'}F=\alpha f_c$. The three special fibres have distinct
intrinsic types: the cusp is reduced and nonnormal, while the other
two have multiplicities three and four. All remaining fibres are
smooth reduced tori. Thus $\alpha$ fixes $p_0,p_1,p_2$ individually.
In the original coordinate, a fractional-linear map fixing
$\infty,0,1$ has zero lower-left and upper-right entries and equal
diagonal entries; hence $\alpha=\mathrm{id}_B$.

Now \cref{pshift:lem:centralizer} forces $c'-c=n\in\mathbb Z$
and linear part $\pm I_2$. Equations (PS10)--(PS11), obtained on
the original normalized finite covers, exclude $-I_2$ and then
exclude odd $n$ for $+I_2$. Conversely (PS27)--(PS28) construct
the global map for every even $n$. This proves
\cref{pshift:thm:classification} without converting a failed
particular map into a nonisomorphism assertion.

As a complex parameter space for the isomorphism classes represented
by this retained family, the quotient $\mathcal P/(2\mathbb Z)$ has
the explicit coordinate
\[
 w=\exp(-\pi ic),\qquad 0<|w|<\exp(-\pi M).
 \tag{PS29}\label{pshift:eq:parameter-disc}
\]
Indeed $|w|=\exp(\pi\operatorname{Im}c)$, and two complex numbers
have the same value of this exponential exactly when they differ
by an even integer. Every point in the displayed punctured disc
has a logarithm giving a $c\in\mathcal P$. The exponential has
nonzero derivative, so these local inverse logarithms give the
claimed quotient complex structure. The integer-period quotient
seen by the middle family instead has coordinate $w^2$; it is
covered two-to-one by (PS29). Its two lifts are distinguished by
the exact order-two summand in (PS12). This statement classifies
the parameters of the retained family; it does not assert that
the family exhausts all deformations or all complex threefolds
with the same singular-fibre types.


\section{Hodge numbers, cotangent direct images, and the Fr\"olicher page}
\label{sec:hodge-frolicher}

This section first derives the one-parameter Hodge table and the complete
Fr\"olicher spectral sequence from the candidate-specific structure-sheaf
calculation and the verified sphere topology.  It then performs the missing
logarithmic direct-image calculation through all three special fibres.  The
cyclic invariant lattices, logarithmic Kodaira--Spencer divisor, residue
cancellation, and one remaining global scalar give
$h^{1,2}=1$ and $h^{1,1}=2$; thus no Hodge number remains undetermined.

\subsection{Vanishing of holomorphic forms}

We first place the holomorphic-form calculation inside the workbench rather
than using it as an unexpanded assertion from the reconstructed source.

\begin{proposition}[candidate-specific holomorphic-form vanishing]
\label{prop:workbench-holomorphic-form-vanishing}
\Verified\ For the constructed threefold,
\[
 H^0(X,\Omega_X^p)=0\qquad (p=1,2,3).
\]
\end{proposition}

\begin{proof}
Over the marked cover of the smooth part, write the fibre coordinates as
$\zeta=(\zeta_1,\zeta_2)$ and
\[
 \Pi=[Z\mid I_2],\qquad
 Z=\begin{pmatrix}6\mu&\tau\\ \beta&\mu\end{pmatrix}.
\]
Translation by $\lambda\in\Lambda$ sends
$d\zeta$ to $d\zeta+\dot\Pi\lambda\,dz$.  Compactness of each torus fibre
therefore puts every translation-invariant holomorphic form into one of the
following normal forms:
\begin{align*}
 p=1:&\quad a(z)\,dz+c(z)\cdot d\zeta,
       &c(z)\dot Z(z)&=0,\\
 p=2:&\quad dz\wedge\bigl(b(z)\cdot d\zeta\bigr),\\
 p=3:&\quad c(z)\,dz\wedge d\zeta_1\wedge d\zeta_2.
\end{align*}
Indeed, fibrewise periodic coefficients are constant.  A possible affine
part in a coefficient vanishes after restriction to the real span of the
two constant columns of $\Pi$; in degree two this also kills the putative
$d\zeta_1\wedge d\zeta_2$ term because $\dot\tau$ is nonzero on the smooth
base.

The finite fillings are obtained from smooth torus families by free cyclic
quotients, so the displayed coefficients extend holomorphically across the
two elliptic lifts.  At the cusp use the verified toric charts
$t=w_0w_1w_2$.  Substitution of
\[
 d\zeta_i=(2\pi i)^{-1}dx_i/x_i,
 \qquad ds=(2\pi i)^{-1}dt/t
\]
and restriction successively to the coordinate lines and planes shows that
the coefficients of $d\zeta_i$, $ds$, $ds\wedge d\zeta_i$, and
$ds\wedge d\zeta_1\wedge d\zeta_2$ vanish to order at least one in $t$.
Indeed, on every unimodular monomial chart the triple
$(d\log x_1,d\log x_2,d\log t)$ is a constant invertible linear combination
of the $d\log w_j$.  A coefficient $F(t)$ multiplying a nonzero logarithmic
residue is holomorphic only if its constant term vanishes: restricting to
$w_j=0$ gives that residue as $F(0)$, so $F(t)=tG(t)$.  Invertibility of the
monomial change of basis applies this argument to each displayed
coefficient.
Thus every coefficient used below descends with the asserted regularity at
all three points of the base.

The linear part of a deck transformation has the form
\[
 R_g=\begin{pmatrix}r_g&0\\ *&1\end{pmatrix},
 \qquad \jmath_g:=r_g^{-1}=c_g\tau+d_g,
 \qquad g^*(d\tau)=\jmath_g^{-2}d\tau.
\]
For a one-form, the second component $c_2$ is a holomorphic invariant
function on the base and vanishes at the cusp, hence is zero.  The second
column of $c\dot Z=0$ then gives $c_1\dot\tau=0$, so $c_1=0$.  What remains
is a holomorphic one-form on $\mathbb P^1$, and hence vanishes.

For a two-form, $b_2\,dz$ descends to a holomorphic one-form on
$\mathbb P^1$ and is zero.  The remaining form
$\psi=b_1\,dz$ satisfies $g^*\psi=\jmath_g\psi$.  Consequently
\[
 \psi^{\otimes2}\otimes d\tau
\]
is invariant.  At the order-three and order-four points it descends as a
cubic differential with a pole of order at most two, while the toric cusp
estimate gives a pole of order at most one.  It is therefore a section of
\[
 (\Omega^1_{\mathbb P^1})^{\otimes3}
 (2p_1+2p_2+p_0)\simeq\mathcal O_{\mathbb P^1}(-1),
\]
so it is zero.  Since $\dot\tau$ is nonzero on a dense open set,
$\psi=0$.  The coefficient of a three-form obeys the same weight-one law
and the same argument applies.  Hence all three spaces vanish.
\end{proof}

\subsection{The exact one-parameter Hodge diamond}

\begin{theorem}[Hodge diamond determined up to one integer]
\label{thm:parametric-hodge-diamond}
\Verified\ Put $m=h^{1,2}(X)$.  Then $m\in\mathbb Z_{\ge0}$ and the
only nonzero Hodge numbers are
\[
 h^{0,0}=h^{0,1}=h^{3,2}=h^{3,3}=1,
 \qquad
 h^{1,1}=h^{2,2}=m+1,
 \qquad
 h^{1,2}=h^{2,1}=m.
\]
Equivalently,
\[
\begin{array}{c|cccc}
q\backslash p&0&1&2&3\\ \hline
3&0&0&0&1\\
2&0&m&m+1&1\\
1&1&m+1&m&0\\
0&1&0&0&0
\end{array}
\]
where the row index is $q$ and the column index is $p$.
\end{theorem}

\begin{proof}
The independently reconstructed structure-sheaf direct images are
\[
 f_*\mathcal O_X=\mathcal O,
 \qquad R^1f_*\mathcal O_X=\mathcal O\oplus\mathcal O(-1),
 \qquad R^2f_*\mathcal O_X=\mathcal O(-1)
\]
by \cref{thm:structure-direct-images-audit}.  Coherent Leray on
$\mathbb P^1$ therefore gives
\[
 (h^{0,0},h^{0,1},h^{0,2},h^{0,3})=(1,1,0,0).
\]
Together with \cref{prop:workbench-holomorphic-form-vanishing} and Serre
duality \cite[\S13, Theorem~4, p.~22]{Serre1955Duality}, this fixes every
entry except the four middle entries, with
$h^{1,2}=h^{2,1}=m$ and $h^{1,1}=h^{2,2}$.

The verified integral sphere topology gives $c_1(X)=c_2(X)=0$ and
$\int_Xc_3(X)=2$.  In complex dimension three,
\begin{align*}
 \operatorname{td}(X)
 &=1+\frac12c_1+\frac1{12}(c_1^2+c_2)+\frac1{24}c_1c_2,\\
 \operatorname{ch}(\Omega_X^1)
 &=3-c_1+\frac12(c_1^2-2c_2)
   -\frac16(c_1^3-3c_1c_2+3c_3).
\end{align*}
Hirzebruch--Riemann--Roch for arbitrary compact complex manifolds
\cite[Theorem~3, p.~427]{AtiyahSinger1963} hence gives
\[
 \chi(\Omega_X^1)=-\frac12\int_Xc_3=-1.
\]
The already known zero entries reduce this identity to
$-h^{1,1}+h^{1,2}=-1$, proving $h^{1,1}=m+1$.
\end{proof}

This specializes the general constraints of Gray and Ugarte; a convenient
published synthesis is Angella \cite{Angella2018}, while the second-page
symmetry and inequalities originate in \cite{Ugarte2000}.  The argument
above uses the stronger candidate-specific direct images and form
vanishings, rather than assuming a general Hodge pattern for a hypothetical
complex $S^6$.

\begin{theorem}[the complete Fr\"olicher page]
\label{thm:frolicher-e2-candidate}
\Verified\ The Fr\"olicher spectral sequence of $X$ does not degenerate at
$E_1$, but it does degenerate at $E_2$.  Apart from the endpoint copies of
$\mathbb C$, the two nonzero $d_1=\partial$ rows are exact complexes
\[
 0\longrightarrow\mathbb C
 \longrightarrow\mathbb C^{m+1}
 \longrightarrow\mathbb C^m
 \longrightarrow0,
\]
\[
 0\longrightarrow\mathbb C^m
 \longrightarrow\mathbb C^{m+1}
 \longrightarrow\mathbb C
 \longrightarrow0.
\]
Their differential ranks are respectively $(1,m)$ and $(m,1)$.
Consequently
\[
 \dim_{\mathbb C}E_1=6+4m,
 \qquad \dim_{\mathbb C}E_2=2,
 \qquad E_2=E_\infty.
\]
\end{theorem}

\begin{proof}
Fr\"olicher's original construction identifies $E_1$ with Dolbeault
cohomology and $E_\infty$ with the associated graded de Rham cohomology
\cite[p.~643, Lemmas~1--2 and Theorem~3]{Frolicher1955}.  The table in
\cref{thm:parametric-hodge-diamond} leaves only the displayed
$q=1$ and $q=2$ row complexes between the endpoint classes.  Since
$H^k_{\mathrm{dR}}(X;\mathbb C)=0$ for $1\le k\le5$, every middle
$E_\infty^{p,q}$ is zero.  There is only one a priori possible higher
differential between middle positions,
$d_2:E_2^{0,2}\to E_2^{2,1}$, but already
$E_1^{0,2}=H^2(X,\mathcal O_X)=0$.  Every other $d_r$ with $r\ge2$ has
zero source or target by the same Hodge table.  Thus $E_2=E_\infty$, so the
two rows must be exact.  The ranks follow from their dimensions.
\end{proof}

The theorem explains why neither topology nor the Fr\"olicher sequence alone
can select a value of $m$: the two exact row complexes exist numerically for
every $m\ge0$.  The singular-fibre calculation below supplies the missing
analytic information.

\subsection{Generic maximal variation and its elliptic zero}

The most tempting shortcut to the remaining Hodge number is to treat the
Kodaira--Spencer matrix as invertible everywhere.  The following calculation
shows precisely what is true and what is false.

\begin{proposition}[determinant of the printed period derivative]
\label{prop:period-derivative-determinant}
\Verified\ With dots denoting differentiation in the coordinate on
$\mathfrak h_z$, set
\[
 K=\det(\dot Z)=6\dot\mu^{\,2}-\dot\tau\dot\beta.
\]
Then $K$ is not identically zero, so the Kodaira--Spencer map has rank two
on a dense open set.  However $K(z_1)=0$ at the chosen order-three elliptic
point.  In particular, the assertion $\dot Z\ne0$ does not imply that
$\dot Z$ is invertible, and direct images across the finite fibres require
their actual cyclic invariant lattices.
\end{proposition}

\begin{proof}
At the cusp, \cref{prop:period-input-audit} gives
\[
 t=e^{2\pi i s},\qquad Z=sB_0+C(t),\qquad
 B_0=\begin{pmatrix}0&1\\-1&0\end{pmatrix},
\]
with $C$ holomorphic at $t=0$.  Hence
\[
 \frac{dZ}{ds}=B_0+2\pi i\,tC'(t),
 \qquad
 \det\!\left(\frac{dZ}{ds}\right)=1+O(t).
\]
The cusp coordinate $s$ is locally biholomorphic on the marked cover, so
$K$ is nonzero near the punctured cusp and therefore is not identically
zero.

Now put $\rho=e^{\pi i/3}=\tau(z_1)$.  The difference between $\mu$ and
the local affine solution
\[
 \mu_1=\frac{2-\tau}{3}
\]
vanishes to order at least two at $z_1$: the homogeneous line has local
generator of order two there.  Since $\tau-\rho$ has a simple zero,
\[
 \dot\mu(z_1)=-\frac13\dot\tau(z_1).
\]
Write
\[
 \phi_1=2-\frac{6(1-\mu)^2}{\tau},
 \qquad \beta\circ g_1=\beta+\phi_1.
\]
Using $\rho^2-\rho+1=0$ and the preceding derivative gives
\[
 \dot\phi_1(z_1)
 =-\frac23(1+\rho)\dot\tau(z_1).
\]
The derivative of the linearized order-three action at its fixed point is
$g_1'(z_1)=-\rho$.  Differentiating the transformation law for $\beta$
therefore yields
\[
 -(1+\rho)\dot\beta(z_1)=\dot\phi_1(z_1),
 \qquad
 \dot\beta(z_1)=\frac23\dot\tau(z_1).
\]
Substitution gives
\[
 K(z_1)=6\left(-\frac13\dot\tau(z_1)\right)^2
 -\dot\tau(z_1)\left(\frac23\dot\tau(z_1)\right)=0.
\]
\end{proof}

The zero divisor of $K$ cannot be read from the single equality $K(z_1)=0$.
It is determined below only after the invariant lattices at the two multiple
fibres have been inserted.  This distinction is essential: direct image does
not commute with taking the dual lattice at a ramified coarse fibre.

\subsection{Logarithmic cotangent direct images}

Put
\[
 B=\mathbb P^1,\qquad
 \Sigma=p_0+p_1+p_2,\qquad
 D=f^{-1}(\Sigma)_{\mathrm{red}}=W+S_1+S_2,
\]
and set
\[
 \mathcal L=\Omega_B^1(\log\Sigma)\simeq\mathcal O_B(1),\qquad
 \mathcal A=\Omega_X^1(\log D),\qquad
 \mathcal E=\mathcal A/f^*\mathcal L.
\tag{11.1}
\]
At the cusp the local equations $t=z_1\cdots z_r$, $1\le r\le3$, show
directly that $\mathcal A$ and $\mathcal E$ are locally free.  At a finite
fibre the covering family has divisor $s=0$, the cyclic action is free on the
total space, and $dt/t=m\,ds/s$; hence its locally free logarithmic sequence
descends through the finite \`etale quotient.  Thus everywhere one has the
logarithmic conormal sequence
\[
 0\longrightarrow f^*\mathcal L\longrightarrow\mathcal A
 \longrightarrow\mathcal E\longrightarrow0.
\tag{11.2}
\]
The local $d\log$ bases and their relative form, including
$t=\prod z_i^{e_i}$ and analytic base change, are the logarithmic conormal
formalism of \cite[Chapter~II, \S3.2--3.3, pp.~72--75]{Deligne1970}; the
display above also follows directly from the stated local equations.

\begin{proposition}[the three logarithmic relative direct images]
\label{prop:log-relative-direct-images}
\Verified\ One has
\begin{align*}
 f_*\mathcal E&\simeq\mathcal O\oplus\mathcal O(-1),\\
 R^1f_*\mathcal E&\simeq
       \mathcal O^{\oplus2}\oplus\mathcal O(-1)^{\oplus2},\\
 R^2f_*\mathcal E&\simeq\mathcal O\oplus\mathcal O(-1).
\end{align*}
All three sheaves are locally free and commute with base change.
\end{proposition}

\begin{proof}
At the cusp, the toric coordinates give
\[
 \mathcal E|_{N_0}
 =\mathcal O_{N_0}[d\log x_1]\oplus
  \mathcal O_{N_0}[d\log x_2].
\tag{11.3}
\]
Indeed a quotient shear changes each $d\log x_i$ only by a holomorphic
multiple of $d\log t$, which is zero in $\mathcal E$.  The already proved
structure-sheaf calculation therefore gives free local direct images of
ranks $2,4,2$ at $p_0$.

We now keep the ramification at $p_j$.  Write
\[
 A=\mathbb C\{t\},\qquad C=\mathbb C\{s\},\qquad t=s^m,
 \qquad (m,a)=(3,2)\text{ or }(4,1).
\]
Cartan's analytic quotient theorem identifies local holomorphic functions on
the quotient with invariants \cite[\S1, Theorem~1, pp.~1--2]{Cartan1954Quotient}.
The action on the total space is free, so the quotient map is finite \'etale;
the corresponding locally free sheaves and their equivariant structures
descend by finite \'etale descent.  This sheaf statement is not being
attributed to Cartan's theorem on invariant functions.  Finite pushforward
is exact, and Reynolds averaging
$v\mapsto |G_j|^{-1}\sum_{g\in G_j}gv$ makes the invariant functor exact in
characteristic zero.  Hence
\[
 R^qf_*\mathcal E|_{D_j}
 =\bigl(\pi_*R^q\widetilde f_*\widetilde{\mathcal E}\bigr)^{G_j}.
\tag{11.3a}
\]
Choose equivariant frames $e_0,e_1$ of the relative cotangent and
$b_0,b_1$ of $R^1\widetilde f_*\mathcal O$ whose semilinear characters are
respectively
\[
 (1,\zeta^a),\qquad (1,\zeta^{-a});
\]
$b_0$ is the nowhere-zero class $\gamma^{0,1}$.  These may be chosen as
holomorphic constant-character frames, rather than merely as eigenvectors on
the central fibre.  Indeed, if $w$ lifts a $\nu$-eigenvector at $s=0$, then
\[
 P_\nu(w)=\frac1m\sum_{k=0}^{m-1}\nu^{-k}(g^*)^kw
\]
satisfies $g^*P_\nu(w)=\nu P_\nu(w)$ and has the same nonzero central value;
after shrinking the disc, the projected eigenvectors form a frame.  On the
marked torus family the projection formula gives
\[
 R^q\widetilde f_*\widetilde{\mathcal E}
 \simeq
 (\mathcal O e_0\oplus\mathcal O e_1)
 \otimes
 \bigwedge^q(\mathcal O b_0\oplus\mathcal O b_1),
 \qquad 0\le q\le2.
\]
Solving
$\nu\zeta^r=1$ for the invariant exponent gives the following \emph{actual}
$A$-bases:
\begin{align*}
 f_*\mathcal E:\quad& e_0,\ s^{m-a}e_1,\\
 R^1f_*\mathcal E:\quad&
 e_0b_0,\ s^ae_0b_1,\ s^{m-a}e_1b_0,\ e_1b_1,\\
 R^2f_*\mathcal E:\quad&
 s^ae_0(b_0\wedge b_1),\ e_1(b_0\wedge b_1).
\tag{11.4}
\end{align*}
For comparison, the variable tangent lattice is generated by $s^av_1$.
Thus evaluation of the descended variable cotangent on the descended
variable tangent is $s^m=t$, not a unit.  This proves explicitly why
\[
 f_*\mathcal E\longrightarrow(f_*T_{X/B})^\vee
\]
is multiplication by $\operatorname{diag}(1,t)$ locally, although
$\mathcal E^\vee\simeq T_{X/B}$ before direct image.  This last identity is
an identity of locally free relative sheaves, not an assertion that formation
of the kernel commutes with taking a tangent-space fibre.  For example, at a
triple chart
\[
 R=\mathbb C\{x,y,z\},\qquad t=xyz,
\]
one has
\[
 T_{X/B}=\ker\!\left[R^3\xrightarrow{(yz,xz,xy)}R\right].
\]
If $yz\,u+xz\,v+xy\,w=0$, reduction modulo $x$ and $y$ gives
$u=xU$ and $v=yV$, and cancellation in the domain gives
$w=-z(U+V)$.  Hence
\[
 T_{X/B}
 =R(x\partial_x-z\partial_z)\oplus
  R(y\partial_y-z\partial_z)
 =\mathcal E^\vee.
\]
The kernel of the pointwise linear map $df_x$ at the triple point is
three-dimensional, whereas the fibre of this sheaf kernel is
two-dimensional: tensoring the kernel sequence with the residue field is
not left exact because the image ideal $(xy,xz,yz)$ is not flat.  Thus a
section can be nowhere zero in $T_{X/B}=\mathcal E^\vee$ even when its image
in the ambient tangent-space fibre vanishes.

The invariant vertical field $\xi$ gives a trivial subline of
$V:=f_*T_{X/B}$; the quotient has the reciprocal automorphy of the Hodge
line $\mathfrak K\simeq\mathcal O(1)$ from
\cref{lem:hodge-line-audit}.  Since $H^1(\mathcal O(1))=0$, the resulting
extension splits:
\[
 V\simeq\mathcal O\xi\oplus\mathcal O(-1),
 \qquad V^\vee\simeq\mathcal O\oplus\mathcal O(1).
\]
Evaluation gives a generically invertible map
$f_*\mathcal E\to V^\vee$.  Formula (11.4) shows that its elementary
divisors are $(1,t)$ at each of $p_1,p_2$ and $(1,1)$ elsewhere.  Thus its
cokernel has total length two and
\(\deg(f_*\mathcal E)=1-2=-1\).  Evaluation on $\xi$ is stalkwise
surjective: the covector dual to the fixed eigenvector is a unit generator
at the cusp and in both finite invariant bases.  Its kernel is consequently
a line bundle of degree $-1$, and
\[
 0\longrightarrow\mathcal O(-1)\longrightarrow f_*\mathcal E
 \xrightarrow{\operatorname{ev}_\xi}\mathcal O\longrightarrow0.
\]
Because $H^1(\mathcal O(-1))=0$, this extension also splits, proving
\[
 f_*\mathcal E\simeq\mathcal O\oplus\mathcal O(-1).
\]
Also, by 
\cref{thm:structure-direct-images-audit},
\[
 G:=R^1f_*\mathcal O_X=\mathcal O\oplus\mathcal O(-1),
 \qquad R^2f_*\mathcal O_X=\det G=\mathcal O(-1).
\tag{11.5}
\]
Here is the global filtration implicit in those bases.  Choose the
degree-zero summand in the preceding split sequence so that its generator
evaluates to one on $\xi$; locally this is the fixed eigen-covector $e_0$.
Together with the fixed class $b_0=\gamma^{0,1}$, take the corresponding
saturated sublines in $f_*\mathcal E$ and $G$.  Formula (11.3) at the cusp
and the bases (11.4) at the finite points show that these sublines are locally
generated by $e_0,b_0$ and that their quotients are locally generated by
$s^{m-a}e_1,s^ab_1$; hence both inclusions and both quotients are locally
free.  The evaluation splitting and Hodge-line calculation identify the two
filtrations as
\[
 0\subset\mathcal O\subset f_*\mathcal E,\qquad
 0\subset\mathcal O\subset G,\qquad
 \operatorname{gr}(f_*\mathcal E)=
 \operatorname{gr}(G)=\mathcal O\oplus\mathcal O(-1).
\]
Over the smooth locus the cup-product identification is
$R^1f_*\mathcal E=(f_*\mathcal E)\otimes G$.  Formula (11.4) extends its four
character lines as saturated rank-one subsheaves over all of $B$.  Three
associated-graded lines are generated locally by
$e_0b_0,s^ae_0b_1,s^{m-a}e_1b_0$ and have degrees $0,-1,-1$.  On the fourth,
the separately descended factors would
give
\[
 (s^{m-a}e_1)(s^ab_1)=t(e_1b_1),
\]
whereas $e_1b_1$ itself is invariant.  Thus that graded line is the upper
modification of $\mathcal O(-2)$ at $p_1+p_2$, hence is $\mathcal O$.
Similarly, the two graded lines of $R^2f_*\mathcal E$ are
$\mathcal O(-1)$ and
$\mathcal O(-2)(p_1+p_2)=\mathcal O$.  Every extension between the resulting
degree-zero and degree-minus-one lines splits: its $\operatorname{Ext}^1$
is one of $H^1(\mathcal O(-1))$, $H^1(\mathcal O)$, or
$H^1(\mathcal O(1))$, all zero.  This proves the three asserted splittings.
This is also the elementary filtered form of the
Birkhoff--Grothendieck classification
\cite[Theorem~2.1, p.~126]{Grothendieck1957}; the divisor modifications and
$\operatorname{Ext}^1$ identifications may equivalently be read in
\cite[Tags 0C6K, 0B3A, and 01XT]{StacksProject}.

Finally, (11.3a)--(11.4) make every direct image free of ranks $2,4,2$ at
the finite points; (11.3) does the same at the cusp, and the smooth locus is
standard.  The multiple-fibre base change can also be read directly from the
invariant modules.  The Reynolds idempotent is $A$-linear, so invariants
commute with arbitrary finite $A$-base change.  For the scheme-theoretic
special fibre,
\[
 C\otimes_AA/(t)=C/(s^m),
 \qquad
 \bigl((C/(s^m))v\bigr)^G=\mathbb C\,s^rv
\]
on a character line whose invariant exponent satisfies $0\le r<m$.
Every basis vector in (11.4) therefore survives reduction modulo $t$, and
the special-fibre dimensions are exactly $2,4,2$.  No $R^q$ with $q>2$
occurs.  Descending induction in analytic cohomology and base change then
makes every base-change map an isomorphism, including for the nonreduced
finite fibres; precisely the same criterion, with citation, is recorded in
the proof of
\cref{thm:structure-direct-images-audit}; see in particular
\cite[Chapter~III, \S3, Theorem~3.4 and Corollary~3.10]{BanicaStanasila1976}.
\end{proof}

\begin{proposition}[the logarithmic Kodaira--Spencer divisor]
\label{prop:log-ks-divisor}
\Verified\ The connecting maps of (11.2) are
\begin{align*}
 \kappa_0&:f_*\mathcal E\longrightarrow
  \mathcal L\otimes R^1f_*\mathcal O_X,\\
 \kappa_1&:R^1f_*\mathcal E\longrightarrow
  \mathcal L\otimes R^2f_*\mathcal O_X\simeq\mathcal O.
\end{align*}
Then
\[
 \operatorname{div}(\det\kappa_0)=p_1+p_2,
 \qquad \kappa_1\text{ is surjective as a sheaf}.
\tag{11.6}
\]
In particular $\kappa_0$ is injective and
$\operatorname{coker}\kappa_0=\mathbb C_{p_1}\oplus\mathbb C_{p_2}$.
Moreover the raw period determinant has
\[
 \operatorname{ord}_{z_1}K=2,\qquad
 \operatorname{ord}_{z_2}K=0.
\tag{11.7}
\]
\end{proposition}

\begin{proof}
On the smooth locus, the connecting construction and its
exterior-product Leibniz rule are the Gauss--Manin/Kodaira--Spencer
formalism of
\cite[\S2, equations~(5)--(10), pp.~202--204]{KatzOda1968}.
Katz--Oda work with a smooth algebraic morphism; the citation is not being
used as a logarithmic analytic special-fibre theorem.  In the present family,
Deligne's local monomial calculation supplies the locally split logarithmic
sequence (11.2) \cite[Chapter~II, \S3.3]{Deligne1970}.  Taking a local
splitting and a \v Cech representative, the connecting cochain is the
difference of the two splittings.  Since exterior product of cochains obeys
\[
 \delta(\alpha\wedge\beta)
 =\delta\alpha\wedge\beta+(-1)^{\deg\alpha}
   \alpha\wedge\delta\beta,
\]
the two displayed connecting maps extend across every logarithmic chart and
satisfy the same Leibniz rule.  Thus the singular-fibre formula below is a
direct calculation from (11.2), not an extrapolation of the Katz--Oda
hypotheses.
The determinant is a section of
\[
 \det(\mathcal L\otimes G)\otimes\det(f_*\mathcal E)^{-1}
 =\mathcal L^{\otimes2}\otimes\det G
   \otimes\det(f_*\mathcal E)^{-1}
 \simeq\mathcal O(2)\otimes\mathcal O(-1)\otimes\mathcal O(1)
 \simeq\mathcal O(2).
\tag{11.8}
\]
At the cusp, $Z=sB_0+C(t)$ and $dZ/ds=B_0+O(t)$, so it is a unit.  At a
finite point use $\ell=dt/t=m\,ds/s$.  In invariant determinant frames
with no normalization imposed against the raw marked frames, direct
calculation gives
\[
 \det\kappa_0
 =\varepsilon(t)m^{-2}s^{m+2-2a}K(s),
 \qquad \varepsilon(0)\ne0.
\tag{11.9}
\]
Here $\varepsilon$ is the nowhere-zero invariant holomorphic unit recording the two
frame normalizations; compatible eigenframes may be normalized so that
$\varepsilon=1$.  The dot in $K$ means $d/ds$ in the chosen linearizing coordinate.
Replacing $s$ by another such coordinate changes only the unit, so every
divisor and order calculation below is coordinate-independent.
Equivariance gives
\[
 \operatorname{ord}_sK\equiv2a-2\pmod m.
\tag{11.10}
\]
At $(m,a)=(3,2)$, 
\cref{prop:period-derivative-determinant} supplies $K(z_1)=0$; hence the
right side of (11.9) has coarse order at least one.  At $(m,a)=(4,1)$,
the prefactor is $s^4=t$, so the coarse order is again at least one.
The two forced zeros exhaust the degree in (11.8).  They are therefore
simple, there are no others, and (11.7) follows from (11.9)--(11.10).

For completeness, we derive the coarse matrices rather than only recording
their shape.  Let
\[
 D_{\mathrm{KS}}(s)=(d_{rs}(s))_{0\leq r,s\leq1}
\]
be the Kodaira--Spencer matrix on the marked cover after the common factor
$ds$ has been removed.  Since $t=s^m$ and
$ds=(s/m)\ell$, passage from the marked frames to the actual invariant
frames in (11.4) gives
\[
 A=\frac{s}{m}d_{00},\qquad
 B=\frac{s^{1-a}}m d_{01},\qquad
 C=\frac{s^{1-a}}m d_{10},\qquad
 D_0=\frac{s^{m+1-2a}}m d_{11}.
\tag{11.10a}
\]
The character of each expression in (11.10a) is trivial, so each is a
holomorphic function of $t$.  In the bases
$(e_0,s^{m-a}e_1)$ and $(b_0,s^ab_1)$ one consequently has the exact matrix
\[
 [\kappa_0]=\begin{pmatrix}A&tB\\ C&D_0\end{pmatrix}.
\tag{11.10b}
\]
Its determinant is (11.9).  Applying the exterior-product Leibniz rule to
the four ordered generators
\[
 e_0b_0,\quad s^ae_0b_1,\quad
 s^{m-a}e_1b_0,\quad e_1b_1
\]
now gives, without a change of basis,
\[
 [\kappa_1]=(-C,\ A,\ -D_0,\ B)
\tag{11.11}
\]
into the invariant generator $s^a(b_0\wedge b_1)\ell$.  If (11.11) were
not surjective at a finite point, all four entries would lie in $(t)$ and
$\det\kappa_0$ would have order at least two, contrary to (11.6).  Away
from $p_1,p_2$, the determinant divisor (11.6) makes $\kappa_0$
invertible, so the exterior-product row (11.11) contains a unit and is
surjective there as well.  At the cusp this reduction is not an additional
choice.  Use the logarithmic base
frame
\(\ell=ds=(2\pi i)^{-1}dt/t\), the relative cotangent frame
\((e_0,e_1)=([d\log x_1],[d\log x_2])\) of (11.3), and the quotient frame
\((\gamma,u)\) of \(G\), which extends across the cusp by
\cref{thm:structure-direct-images-audit}.  In these frames the smooth-locus Kodaira--Spencer matrix is the
transpose of the derivative of the period matrix: for
\(h=a e_0+b e_1\),
\[
 \kappa_0(h)=
 (6a\dot\mu+b\dot\beta)\gamma+
 (a\dot\tau+b\dot\mu)u .
 \tag{11.11a}
\]
Here the dot is \(d/ds\).  The cusp expansion
\(Z=sB_0+C(t)\), \(t=e^{2\pi i s}\), gives
\[
 \frac{dZ}{ds}=B_0+2\pi i\,tC'(t),
 \qquad
 B_0=\begin{pmatrix}0&1\\-1&0\end{pmatrix}.
\]
These logarithmic frames extend over \(t=0\), so (11.11a) reduces exactly
to \(B_0^{\mathsf t}\).  Write the cusp matrix, separately from the
finite-fibre normalization above, as
\([\kappa_0]_{p_0}=\left(\begin{smallmatrix}A&E\\C&D\end{smallmatrix}\right)\).
In the same ordered exterior basis used for (11.11), its Leibniz row is
\((-C,A,-D,E)\).  Substituting
\((A,E,C,D)=(0,-1,1,0)\) gives the second reduction.  Thus, in the named
frames,
\[
 [\kappa_0]\equiv\begin{pmatrix}0&-1\\1&0\end{pmatrix}\pmod t,
 \qquad [\kappa_1]\equiv(-1,0,0,-1)\pmod t,
\]
exactly.  Under arbitrary holomorphic changes of source and target frames
the matrices change by invertible left and right factors; in particular
invertibility of \(\kappa_0\) and the existence of a unit entry in the row
of \(\kappa_1\) are unchanged.  Thus $\kappa_1$ is surjective everywhere.
\end{proof}

\subsection{Residues and the absolute cotangent direct images}

Let $i_D:D\hookrightarrow X$, and let
\[
 \nu:\widetilde D=\mathrm{dP}_6\sqcup S_1\sqcup S_2\longrightarrow D
\]
be the normalization.  In formulas below, $\nu_*\mathcal O_{\widetilde D}$
abbreviates $(i_D\circ\nu)_*\mathcal O_{\widetilde D}$ as a sheaf on $X$.
The ordinary/logarithmic local calculation for
$z_1\cdots z_r=0$ glues to
\[
 0\longrightarrow\Omega_X^1\longrightarrow\mathcal A
 \xrightarrow{\operatorname{res}}\nu_*\mathcal O_{\widetilde D}
 \longrightarrow0.
\tag{11.12}
\]
The normalization in the target is indispensable at the self-crossing
cusp.  This is the degree-one residue sequence associated with the normalized
strata of a normal-crossings divisor
\cite[Chapter~II, \S3.4--3.7, pp.~76--78]{Deligne1970}; the displayed local
calculation fixes its convention in the present self-crossing model.

\begin{proposition}[absolute cotangent direct images]
\label{prop:absolute-cotangent-direct-images}
\Verified\ Put $\mathcal K=\ker\kappa_1$.  Then
\[
 f_*\Omega_X^1\simeq\mathcal O(-2),\qquad
 R^1f_*\Omega_X^1\simeq\mathcal K,\qquad
 R^2f_*\Omega_X^1\simeq\mathcal O\oplus\mathcal O(-1).
\tag{11.13}
\]
None of these sheaves has torsion.
\end{proposition}

\begin{proof}
The three normalized components have
\[
 R^0f_*\nu_*\mathcal O_{\widetilde D}
  =\mathbb C_{p_0}\oplus\mathbb C_{p_1}\oplus\mathbb C_{p_2},
 \quad
 R^1f_*\nu_*\mathcal O_{\widetilde D}
  =\mathbb C_{p_1}\oplus\mathbb C_{p_2},
 \quad R^2=0.
\tag{11.14}
\]
Indeed $\mathrm{dP}_6$ is rational and $S_1,S_2$ are bielliptic, with
$h^1(\mathcal O_{S_j})=1$ and $h^2(\mathcal O_{S_j})=0$.

The long direct-image sequence of (11.2), together with
\cref{prop:log-ks-divisor}, gives
\[
 f_*\mathcal A=\mathcal O(1),\qquad
 0\longrightarrow T\longrightarrow R^1f_*\mathcal A
 \longrightarrow\mathcal K\longrightarrow0,
 \qquad R^2f_*\mathcal A=\mathcal O\oplus\mathcal O(-1),
\tag{11.15}
\]
where $T=\mathbb C_{p_1}\oplus\mathbb C_{p_2}$.
The residue of the base logarithmic form is weighted evaluation:
\[
 \mathcal O(1)\longrightarrow
 \mathbb C_{p_0}\oplus\mathbb C_{p_1}\oplus\mathbb C_{p_2},
 \qquad q\longmapsto(q(p_0),\,3q(p_1),\,4q(p_2)).
\tag{11.16}
\]
It is surjective as a sheaf and has kernel
$\mathcal O(1)(-p_0-p_1-p_2)=\mathcal O(-2)$.

It remains to identify the torsion in (11.15).  At $p_j$, the residue map on
$G\otimes\mathbb C_{p_j}$ sends
\[
 x b_0+y s^ab_1\longmapsto m_jx[b_0]
 \in H^1(S_j,\mathcal O_{S_j}).
\tag{11.17}
\]
More precisely, after trivializing $\mathcal L$ by
$\ell=dt/t$, the logarithmic conormal and residue maps form the exact
commutative diagram
\[
\begin{array}{ccccccccc}
0&\longrightarrow&f^*\mathcal L&\longrightarrow&\mathcal A
 &\longrightarrow&\mathcal E&\longrightarrow&0\\
&&\big\downarrow\,{\scriptstyle\operatorname{res}_B}
&&\big\downarrow\,{\scriptstyle\operatorname{res}_{S_j}}
&&\big\downarrow\,{\scriptstyle 0}&&\\
0&\longrightarrow&i_{j*}\mathcal O_{S_j}&\xrightarrow{\ =\ }
 i_{j*}\mathcal O_{S_j}&\longrightarrow&0&&
\end{array}
\tag{11.16a}
\]
with
\[
 \operatorname{res}_B(q\ell)=m_jq|_{S_j}.
\tag{11.16b}
\]
Applying $Rf_*$ to (11.16a) gives the typed composite
\[
 f_*\mathcal E|_{p_j}\xrightarrow{\ \kappa_0\ }
 (\mathcal L\otimes G)|_{p_j}
 \xrightarrow{\ \operatorname{res}_{1,j}\ }
 H^1(S_j,\mathcal O_{S_j}),
 \qquad \operatorname{res}_{1,j}\kappa_0=0,
\tag{11.16c}
\]
and (11.16b) gives exactly
\[
 \operatorname{res}_{1,j}(xb_0+ys^ab_1)=m_jx[b_0].
\]
Thus its kernel is the one-dimensional line
$\mathbb C\,s^ab_1$.  The simple zero of $\det\kappa_0$ proved in
\cref{prop:log-ks-divisor} implies that $\kappa_0\bmod t$ has rank one;
the zero composite in (11.16c) therefore forces
\[
 \operatorname{im}(\kappa_0\bmod t)
 =\ker\operatorname{res}_{1,j}
 =\mathbb C\,s^ab_1.
\tag{$11.17a_{\mathrm{fin}}$}
\]
Consequently the induced map
$\operatorname{coker}\kappa_0\to H^1(S_j,\mathcal O_{S_j})$ is an
isomorphism at each finite fibre.  Thus each summand
$T_{p_j}\to H^1(S_j,\mathcal O_{S_j})$ is an isomorphism.  In the long
sequence of (11.12), (11.16) first removes the entire degree-zero residue,
and (11.17) then cancels $T$.  What remains is precisely (11.13).
\end{proof}

\begin{remark}[the cusp is not the torsion locus]
At $p_0$, if
$Q=\nu_*\mathcal O_{\mathrm{dP}_6}/\mathcal O_W$, then
\[
 0\longrightarrow\Omega^1_{N_0/\Delta}\longrightarrow
 \mathcal E|_{N_0}\longrightarrow i_*Q\longrightarrow0,
 \qquad h^\bullet(Q)=(2,1,0).
\tag{11.17a}
\]
This relative sequence is the right column of the following exact
normalization--residue diagram; all sheaves in the bottom row are regarded as
sheaves on $N_0$ by $i_*$:
\[
\begin{array}{ccccccccc}
&&0&&0&&0&&\\
&&\big\downarrow&&\big\downarrow&&\big\downarrow&&\\
0&\longrightarrow&f^*\Omega^1_\Delta&\longrightarrow&\Omega^1_{N_0}
 &\longrightarrow&\Omega^1_{N_0/\Delta}&\longrightarrow&0\\
&&\big\downarrow&&\big\downarrow&&\big\downarrow&&\\
0&\longrightarrow&f^*\Omega^1_\Delta(\log0)&\longrightarrow&
 \Omega^1_{N_0}(\log W)&\longrightarrow&\mathcal E|_{N_0}
 &\longrightarrow&0\\
&&\big\downarrow\,{\scriptstyle\operatorname{res}_B}
&&\big\downarrow\,{\scriptstyle\operatorname{res}_W}
&&\big\downarrow\,{\scriptstyle\rho}&&\\
0&\longrightarrow&i_*\mathcal O_W&\longrightarrow&
 i_*\nu_*\mathcal O_{\mathrm{dP}_6}&\longrightarrow&i_*Q
 &\longrightarrow&0\\
&&\big\downarrow&&\big\downarrow&&\big\downarrow&&\\
&&0&&0&&0&&
\end{array}
\tag{11.17a$'$}
\]
The first two columns are the base and normalized-divisor residue sequences;
the bottom row is the definition of $Q$.  Hence the snake lemma identifies
the right column with (11.17a), including its displayed map $\rho$.

The exactness assertion is also visible stalkwise.  At a double chart $t=xy$, with one
coordinate along the double stratum and logarithmic generator $e$, the image
of ordinary relative differentials in $\mathcal E$ contains the stratum
differential and has transverse part $(x,y)e$; its cokernel is therefore
$R/(x,y)$.  At a triple chart $t=xyz$, write
$\mathcal E=Re_1\oplus Re_2$ with
$d\log x=e_1$, $d\log y=e_2$, and
$d\log z=-e_1-e_2$.  The image is generated by
\[
 xe_1,\qquad ye_2,\qquad z(e_1+e_2),
\]
Writing a normalized branch value as a triple modulo the diagonal, the
quotient map is explicitly
\[
 e_1\longmapsto[(1,0,0)],\qquad
 e_2\longmapsto[(0,1,0)].
\tag{11.17a$''$}
\]
Its kernel is exactly
$\langle xe_1,ye_2,z(e_1+e_2)\rangle$, so this is the standard local
presentation of the normalization quotient $Q$.
Away from the double locus the map is an isomorphism.  This proves (11.17a)
stalkwise.
Here are the maps, rather than only their ranks.  Quotienting the normalization
complex (9.10) by $\mathcal O_W$ gives
\[
 0\longrightarrow Q\longrightarrow
 \bigoplus_{k=1}^3\mathcal O_{D_k}
 \xrightarrow{\delta_1}\mathbb C_P\oplus\mathbb C_Q
 \longrightarrow0.
\]
Because $D_k\simeq\mathbb P^1$, its cohomology is the cohomology of
\[
 \mathbb C^3\longrightarrow\mathbb C^2,\qquad
 (a_1,a_2,a_3)\longmapsto(s,s),\qquad s=a_1-a_2+a_3.
\]
Thus
\[
 H^0(Q)=\{(a_1,a_2,a_3):a_1-a_2+a_3=0\}\simeq\mathbb C^2,
 \qquad H^1(Q)=\mathbb C^2/\mathbb C(1,1)\simeq\mathbb C,
\tag{11.17b}
\]
and $H^2(Q)=0$.  Since $\mathrm{dP}_6$ is rational, the connecting maps for
$0\to\mathcal O_W\to\nu_*\mathcal O_{\mathrm{dP}_6}\to Q\to0$ are
isomorphisms
\[
 \delta_0:H^0(Q)\xrightarrow{\sim}H^1(W,\mathcal O_W),
 \qquad
 \delta_1:H^1(Q)\xrightarrow{\sim}H^2(W,\mathcal O_W).
\tag{11.17c}
\]

At a triple chart $t=z_0z_1z_2$, a relative logarithmic form is represented
by $\eta=\sum a_i\,d\log z_i$, with holomorphic coefficient triple
$(a_0,a_1,a_2)$ taken modulo $\mathcal O(1,1,1)$.  Restricting the branch
differences to the three double strata gives
\[
 \rho(\eta)=(a_0-a_1,\ a_0-a_2,\ a_1-a_2).
\tag{11.17d}
\]
Formula (11.17d) is written in the $P$-ordering of the globally labelled
double curves.  The $Q$-point sign is fixed by a separate calculation, not by
an implicit relabelling.  In the common $e_1,e_2$-plane, use the ordered
branch vertices
\[
 (b_0,b_1,b_2)=(0,e_1,e_2),\qquad
 (c_0,c_1,c_2)=(e_1,e_1+e_2,e_2)
\]
at $P$ and $Q$, respectively.  The coefficient triples of
$d\log x_1,d\log x_2$ are
\[
 (0,1,0),\ (0,0,1)\quad\text{at }P,
 \qquad
 (1,1,0),\ (0,1,1)\quad\text{at }Q.
\]
The globally labelled branch-difference maps are
\[
 d_P(a)=(a_0-a_1,a_0-a_2,a_1-a_2),
 \qquad
 d_Q(a)=(a_2-a_1,a_0-a_1,a_0-a_2).
\tag{11.17d$'$}
\]
Direct substitution gives, at both triple points, the same two global
sections
\[
 q_1=(-1,0,1),\qquad q_2=(0,-1,-1).
\tag{11.17d$''$}
\]
For constant coefficients this induces an isomorphism
$\mathbb C^3/\mathbb C(1,1,1)\to H^0(Q)$.  In the toric chart
$(z_0,z_1,z_2)=(t/(x_1x_2),x_1,x_2)$ it sends
\[
 d\log x_1\longmapsto(-1,0,1),\qquad
 d\log x_2\longmapsto(0,-1,-1),
\]
which is a basis of the hyperplane in (11.17b); unimodular changes of cusp
chart merely change this basis by $\mathrm{GL}_2(\mathbb Z)$.

Write $\eta_1=d\log x_1$, $\eta_2=d\log x_2$, take the $q_i$ of
(11.17d$''$), and put $\alpha_i=\delta_0(q_i)$.  By (11.17c),
$\alpha_1,\alpha_2$ form a basis of $H^1(W,\mathcal O_W)$.  Naturality of
the connecting maps fixes a sign that must be retained.  With the standard
\v Cech differential
$(\delta c)_{i_0\ldots i_{r+1}}
=\sum_k(-1)^kc_{i_0\ldots\widehat{i_k}\ldots i_{r+1}}$, lift $q_i$ to the
normalization complex.  The total differential of a degree-one cocycle
$\alpha$ times that lift is
$(-1)^1\alpha\smile\delta\widetilde q_i$.  The projection formula therefore
gives exactly
\[
 \delta_1\bigl(\rho_*(\alpha\otimes\eta_i)\bigr)
 =-\alpha\smile\delta_0(q_i)=-\alpha\smile\alpha_i.
\]
The cup product in \cref{lem:cusp-structure-sheaf-audit} is an isomorphism.
Define, with this convention,
\[
 \theta=-\delta_1^{-1}(\alpha_1\smile\alpha_2).
\]
Graded commutativity now gives, with no residual sign ambiguity,
\[
 \rho_*=(0,1,-1,0):
 \mathbb C\langle
 \alpha_1\eta_1,\alpha_1\eta_2,\alpha_2\eta_1,\alpha_2\eta_2
 \rangle\longrightarrow\mathbb C\theta.
\tag{11.17e}
\]
In particular the map is surjective.  Applying direct images to (11.17a),
the first map is reduction modulo $t$ followed by the isomorphism (11.17d),
and the second is (11.17e).  Consequently
\[
 f_*\Omega^1_{N_0/\Delta}=t\mathcal O^2,\quad
 R^1f_*\Omega^1_{N_0/\Delta}=\ker(\mathcal O^4\to\mathbb C_{p_0}),
 \quad R^2f_*\Omega^1_{N_0/\Delta}=\mathcal O^2.
\]
For the row in (11.17e), the middle kernel is generated by
$e_1,e_4,e_2+e_3,te_2$.
Thus $\Omega^1_{X/B}$ is torsion-free in a neighbourhood of the cusp, but
non-locally-free along its $xy$ and $xyz$ strata.  Globally its torsion is
supported precisely on the two multiple fibres, as the following local module
calculation proves.  On the smooth cover $J'_j$, let
$R=\mathcal O_{J'_j,x}$ and choose coordinates $(s,y_1,y_2)$ smooth over the
$s$-disc.  Relative to the original base $B$, where $t=s^{m_j}$, one has
\[
 \Omega^1_{J'_j/B,x}
 =\frac{R\,ds\oplus R\,dy_1\oplus R\,dy_2}
        {(m_js^{m_j-1}ds)}
 \simeq R\,dy_1\oplus R\,dy_2
 \oplus\bigl(R/(s^{m_j-1})\bigr)ds.
\tag{11.17f}
\]
Thus the last summand is exactly the torsion submodule.  The quotient map
$q:J'_j\to N_j$ is finite \'{e}tale, so the relative cotangent sequence gives
\[
 q^*\Omega^1_{N_j/B}\simeq\Omega^1_{J'_j/B}.
\]
Faithful flatness descends both the presence of the displayed torsion and its
absence off $s=0$.  Hence ordinary relative torsion is supported exactly on
the two multiple fibres and is absent at the cusp.  The reconstructed
manuscript's assertion that torsion occurs at all three singular fibres is
false.
\end{remark}

\subsection{The scalar and the exact Hodge diamond}

\begin{theorem}[the remaining scalar is nonzero]
\label{thm:h12-scalar}
\Verified\ The induced map
\[
 H^0(B,R^1f_*\mathcal E)\simeq\mathbb C^2
 \xrightarrow{\ H^0(\kappa_1)\ }
 H^0(B,\mathcal O)\simeq\mathbb C
\]
has rank one.
\end{theorem}

\begin{proof}
Every matrix $R_g$ in (2.3d) has second column $(0,1)^{\mathsf t}$.
Consequently the marked vertical field
\[
 \xi=\frac{\partial}{\partial\zeta_2}
\]
is invariant under every deck transformation.  The affine translations in
\cref{prop:candidate-finite-fillings} have identity differential in this
direction, and the corresponding toric logarithmic direction extends across
the cusp charts of \cref{thm:candidate-cusp-geometry}.  Thus $\xi$ is a
global section nowhere zero in the locally free relative sheaf
$T_{X/B}=\mathcal E^\vee$ and is the section used in the evaluation sequence
\[
 0\longrightarrow\mathcal O(-1)\longrightarrow f_*\mathcal E
 \xrightarrow{\operatorname{ev}_\xi}\mathcal O\longrightarrow0.
\]
This sequence splits, and the lift $h$ of $1\in H^0(\mathcal O)$ is unique:
the difference of two lifts is a global section of $\mathcal O(-1)$, while
$H^0(B,\mathcal O(-1))=0$.  On the marked cover write that unique lift as
\[
 h=(a(z),1)\binom{d\zeta_1}{d\zeta_2}.
\tag{11.18}
\]
The function $a$ is holomorphic on each finite cover because $h$ is a
section of the actual invariant lattice, not merely of its generic fibre.
For $g_2$, invariance of (11.18) and
$R_{g_2}=\left(\begin{smallmatrix}1/\tau&0\\-\mu/\tau&1\end{smallmatrix}\right)$
mean exactly
\[
 (a(z),1)=(a(g_2z),1)R_{g_2}(z),
 \qquad\text{hence}\qquad
 a(g_2z)=\tau(z)a(z)+\mu(z).
\tag{11.18a}
\]
At the fixed point $z_2$, the normalized values
$\tau(z_2)=i$ and $\mu(z_2)=(1-i)/2$ therefore give
\[
 a(z_2)=\frac12.
\tag{11.18b}
\]
The cup-product map
$f_*\mathcal E\otimes G\to R^1f_*\mathcal E$ sends
$h\otimes\gamma$ to the unit generator $e_0b_0$ of the degree-zero
character line in (11.4).  This identifies the precise entry of the row
(11.11) that must be evaluated.

Normalize the order-four linearizing coordinate by
\[
 s=(z_2-\bar z_2)s_2,
 \qquad s\circ g_2=-is,
 \qquad \frac{ds}{dz}(z_2)=1.
\]
Write a dot for $d/ds$, and put
$\ell=dt/t=4\,ds/s$.  In the quotient basis $(\gamma,u)$ of $G$, the
marked period derivative and $ds=(s/4)\ell$ give the exact logarithmic map
\[
 \kappa_0(h)=\frac{s}{4}
 \Bigl[(6a\dot\mu+\dot\beta)\gamma
 +(a\dot\tau+\dot\mu)u\Bigr]\otimes\ell.
\tag{11.19}
\]
For $(m,a)=(4,1)$, the actual invariant target generator from (11.4) is
$s(b_0\wedge b_1)\ell$.  Normalize $b_1$ so that
\[
 \gamma\wedge b_1=u_2(s)(\gamma\wedge u),
 \qquad u_2(0)=1;
\]
equivalently, use the invariant frame
\[
 \omega_2=s\,u_2(s)(\gamma\wedge u)\ell.
\tag{11.19f}
\]
The uncorrected expression $s(\gamma\wedge u)\ell$ is only invariant at
$s=0$: indeed $g_2^*(\gamma\wedge u)=\tau(z)(\gamma\wedge u)$, so its
pullback factor is $(-i)\tau(z)$.  Formula (11.11), in the genuine invariant
frame (11.19f), sends the first ordered source generator $e_0b_0$ to
minus the second coefficient of $\kappa_0(e_0)$.  Since $u_2(0)=1$,
(11.19) therefore gives the exact signed coarse value
\[
 \kappa_1(h\otimes\gamma)(p_2)
 =-\frac14\bigl(a(z_2)\dot\tau(z_2)+\dot\mu(z_2)\bigr).
\tag{11.19a}
\]

At the order-four lift $z_2$, the quotient coordinate has order four while
the modular invariant has order two at $\tau=i$; see
\cite[Chapter~VII]{Serre1973}.  Equivalently,
$j(\tau)-1728$ has a double zero at $i$ and
$j(\tau(z))-1728$ has order four at $z_2$.  Hence
\[
 \operatorname{ord}_{z_2}(\tau-i)=2,\qquad \dot\tau(z_2)=0.
\tag{11.20}
\]
By (11.7), $K(z_2)\ne0$; since
$K=6\dot\mu^2-\dot\tau\dot\beta$, (11.20) gives
$K(z_2)=6\dot\mu(z_2)^2\ne0$.  Substitution in (11.19a) therefore gives
\[
 \kappa_1(h\otimes\gamma)(p_2)=-\frac14\dot\mu(z_2)\ne0.
\tag{11.20a}
\]
The image is a global section of
$\mathcal L\otimes R^2f_*\mathcal O_X\simeq\mathcal O_B$; because it is
nonzero at $p_2$, it is a nonzero constant.  Hence
$H^0(\kappa_1)$ has rank one.
\end{proof}

\begin{corollary}[exact Hodge diamond]
\label{cor:exact-hodge-diamond}
\Verified\ The last parameter in
\cref{thm:parametric-hodge-diamond} is
\[
 h^{1,2}(X)=1.
\]
Consequently $h^{1,1}=2$, and the Hodge table is
\[
\begin{array}{c|cccc}
q\backslash p&0&1&2&3\\ \hline
3&0&0&0&1\\
2&0&1&2&1\\
1&1&2&1&0\\
0&1&0&0&0.
\end{array}
\]
The two nontrivial Fr\"olicher $d_1$ rows are both exact complexes
$0\to\mathbb C\to\mathbb C^2\to\mathbb C\to0$, and
$\dim E_1=10$, $\dim E_2=2$.
\end{corollary}

\begin{proof}
Put $\mathcal K=\ker\kappa_1$.  Since
\[
0\longrightarrow\mathcal K\longrightarrow
\mathcal O^{\oplus2}\oplus\mathcal O(-1)^{\oplus2}
\xrightarrow{\kappa_1}\mathcal O\longrightarrow0
\]
and the global map has rank one, its restriction to the two trivial summands
is a nonzero constant row (the negative summands have no global sections).
After a constant change of basis it is the identity on one copy of
$\mathcal O$.  Subtracting the remaining component from that summand makes
the map the projection, so its kernel is
\[
 \mathcal K\simeq\mathcal O\oplus\mathcal O(-1)^{\oplus2}.
\]
By 
\cref{prop:absolute-cotangent-direct-images} and coherent Leray,
\[
 h^{1,2}=h^1(B,\mathcal K)
 +h^0(B,\mathcal O\oplus\mathcal O(-1))=0+1=1.
\]
The remaining assertions follow from
\cref{thm:parametric-hodge-diamond,thm:frolicher-e2-candidate}.
\end{proof}

\subsection{Bott--Chern and Aeppli dimensions}

For a compact complex manifold write
\[
 H^{p,q}_{BC}(X)=
 \frac{\ker\partial\cap\ker\bar\partial\subset A^{p,q}(X)}
      {\operatorname{im}\partial\bar\partial},
 \qquad
 H^{p,q}_{A}(X)=
 \frac{\ker\partial\bar\partial\subset A^{p,q}(X)}
      {\operatorname{im}\partial+\operatorname{im}\bar\partial}.
\]

\begin{theorem}[exact Bott--Chern and Aeppli diamonds]
\label{thm:exact-bott-chern-aeppli}
\Verified\ The Bott--Chern dimensions of the candidate are
\[
\begin{array}{c|cccc}
q\backslash p&0&1&2&3\\ \hline
3&0&0&1&1\\
2&0&1&2&1\\
1&0&2&1&0\\
0&1&0&0&0,
\end{array}
\tag{11.21}
\]
and its Aeppli dimensions are
\[
\begin{array}{c|cccc}
q\backslash p&0&1&2&3\\ \hline
3&0&0&0&1\\
2&0&1&2&0\\
1&1&2&1&0\\
0&1&1&0&0.
\end{array}
\tag{11.22}
\]
In total degree $k=0,\ldots,6$ this is
\[
 (h^k_{BC})=(1,0,2,2,2,2,1),\qquad
 (h^k_A)=(1,2,2,2,2,0,1).
\tag{11.23}
\]
Consequently, set
\[
 \Delta^k=h^k_{BC}+h^k_A-2b_k
\]
for the degree-$k$ non-$\partial\bar\partial$ defects.  The
Angella--Tomassini inequality
\cite[Theorem~A]{AngellaTomassini2013} gives $\Delta^k\ge0$ for every compact
complex manifold; here the exact values are
\[
 (\Delta^k)=(0,2,4,4,4,2,0).
\tag{11.24}
\]
\end{theorem}

\begin{proof}
McHugh associates to a compact complex manifold the finite-dimensional
quotient
\[
 K^{p,q}=
 \frac{\ker\partial\cap\ker\bar\partial\subset A^{p,q}}
 {\bigl(\operatorname{im}\partial\cap\ker\bar\partial\bigr)
  +\bigl(\operatorname{im}\bar\partial\cap\ker\partial\bigr)},
 \qquad k^{p,q}=\dim_{\mathbb C}K^{p,q}.
\]
Put
\[
 Z^{p,q}=\ker\partial\cap\ker\bar\partial\subset A^{p,q},
 \qquad
 B_K^{p,q}=
 \bigl(\operatorname{im}\partial\cap\ker\bar\partial\bigr)
 +\bigl(\operatorname{im}\bar\partial\cap\ker\partial\bigr).
\]
The natural map $Z^{p,q}\to H^{p+q}_{\mathrm{dR}}(X;\mathbb C)$ need not
descend to an injection from $K^{p,q}$; the dimension estimate has the
opposite quotient mechanism.  If a pure form $\alpha\in Z^{p,q}$ is
$d$-exact, the $(p,q)$ component of a primitive gives
\[
 \alpha=\partial\beta^{p-1,q}+\bar\partial\beta^{p,q-1}.
\]
Because $\bar\partial\alpha=0$ and $\partial\alpha=0$, respectively,
$\partial\beta^{p-1,q}$ is $\bar\partial$-closed and
$\bar\partial\beta^{p,q-1}$ is $\partial$-closed.  Thus
$\ker(Z^{p,q}\to H^{p+q}_{\mathrm{dR}})\subset B_K^{p,q}$, so the de Rham
image of $Z^{p,q}$ surjects onto $K^{p,q}$.  Hence
\[
 k^{p,q}\le b_{p+q}.
\tag{11.25}
\]
Because $X$ has the integral cohomology of $S^6$, (11.25) gives
$k^{p,q}=0$ whenever $0<p+q<6$.

For a compact complex threefold, McHugh's Theorem~6.3,
Proposition~6.4, equation~(7), and the ensuing table give
\begin{align*}
 h^{1,1}_{BC}
 &=2h^{0,1}-b_1+\dim_{\mathbb R}H^{1,1}(X;\mathbb R),\\
 h^{1,2}_{BC}
 &=h^{1,2}+\dim_{\mathbb C}E_2^{3,1}-k^{2,0},\\
 h^{2,2}_{BC}
 &=-h^{1,1}_{BC}+h^{0,1}-h^{0,2}-h^{1,0}+h^{1,1}
   +h^{1,2}+h^{2,0}\\
 &\hspace{2em}+2\dim_{\mathbb C}E_2^{3,1}
   +k^{1,1}-k^{1,2}-2k^{2,0},
\end{align*}
and his threefold table gives the remaining entries
\cite[Theorem~6.3, Proposition~6.4, equation~(7), \S7.1]{McHugh2019}; see also the
independent synthesis in \cite[\S4]{Angella2018}.  Here
$H^{1,1}(X;\mathbb R)$ denotes the subspace of
$H^2_{\mathrm{dR}}(X;\mathbb R)$ represented by real $d$-closed forms of
type $(1,1)$, and it is zero because $H^2_{\mathrm{dR}}(X;\mathbb R)=0$,
while \cref{thm:frolicher-e2-candidate} gives $E_2=E_\infty$; since
$b_4=0$, this in particular gives $E_2^{3,1}=0$.  Substitution of
\[
 h^{0,1}=1,\quad h^{0,2}=h^{1,0}=h^{2,0}=0,\quad
 h^{1,1}=2,\quad h^{1,2}=1
\]
therefore gives
\[
 h^{1,1}_{BC}=2,\qquad h^{1,2}_{BC}=h^{2,1}_{BC}=1,\qquad
 h^{2,2}_{BC}=2.
\]
The same table gives
$h^{2,3}_{BC}=h^{3,2}_{BC}=1$, the two endpoint dimensions equal to one,
and every other entry zero.  This proves (11.21).  The Hodge-star duality
of Schweitzer \cite[\S2.c, pp.~9--10]{Schweitzer2007},
$h_A^{p,q}=h_{BC}^{3-p,3-q}$, proves (11.22).  Summing along diagonals and
using $b_0=b_6=1$ and $b_k=0$ for $1\le k\le5$ proves
(11.23)--(11.24).
\end{proof}

The dimension calculation above is completed by the explicit representative
and product theorem in the following section.

\section{Period-coordinate Dolbeault classes and the Bott--Chern algebra}
\label{sec:bcpr-product-structure}

\subsection{Fixed input and conventions}

Let $f\colon X\to B=\mathbb P^1$ be the compact complex threefold of
\cref{thm:candidate-global-geometry}.  We use
$d=\partial+\bar\partial$ and, following
\cite[\S1.a]{Schweitzer2007},
\begin{equation}
 H^{p,q}_{BC}(X)=
 \frac{\ker\partial\cap\ker\bar\partial
       \subset\mathcal A^{p,q}(X)}
      {\operatorname{im}(\partial\bar\partial)}.
\label{eq:bcpr-bc-definition}
\end{equation}
The natural map to Dolbeault cohomology is denoted
\begin{equation}
 \pi_{\bar\partial}^{p,q}\colon H^{p,q}_{BC}(X)
 \longrightarrow H^{p,q}_{\bar\partial}(X).
\label{eq:bcpr-dolbeault-map}
\end{equation}
Exterior multiplication makes the direct sum of the Bott--Chern groups a
bigraded algebra, and complex conjugation interchanges bidegrees $(p,q)$
and $(q,p)$ \cite[\S1.a]{Schweitzer2007}.

By \cref{cor:exact-hodge-diamond}, the two nonzero rows of the
Fr\"olicher differential are the exact complexes
\begin{equation}
 0\longrightarrow H^{0,1}_{\bar\partial}(X)
 \xrightarrow{d_1}H^{1,1}_{\bar\partial}(X)
 \xrightarrow{d_1}H^{2,1}_{\bar\partial}(X)
 \longrightarrow0,
\label{eq:bcpr-hodge-row-one}
\end{equation}
\begin{equation}
 0\longrightarrow H^{1,2}_{\bar\partial}(X)
 \xrightarrow{d_1}H^{2,2}_{\bar\partial}(X)
 \xrightarrow{d_1}H^{3,2}_{\bar\partial}(X)
 \longrightarrow0,
\label{eq:bcpr-hodge-row-two}
\end{equation}
and each row has dimensions $1,2,1$.  The exact Bott--Chern calculation
of \cref{thm:exact-bott-chern-aeppli}, obtained from the compact-threefold
formulae of \cite{McHugh2019}, gives
\begin{equation}
\begin{array}{c|cccccccc}
(p,q)&(0,0)&(1,1)&(2,1)&(1,2)&(2,2)&(3,2)&(2,3)&(3,3)\\ \hline
h^{p,q}_{BC}&1&2&1&1&2&1&1&1.
\end{array}
\label{eq:bcpr-nonzero-dimensions}
\end{equation}
All other Bott--Chern groups vanish.  By
\cref{cor:candidate-sphere}, $b_k(X)=0$ for $1\leq k\leq5$.

On a marked smooth fibre retain the integral basis
\[
 V=\mathbb Z\langle\gamma,u,w,\delta\rangle
\]
without renaming its element $u$, and put
\begin{equation}
 \eta=u\wedge w+6\gamma\wedge\delta,
 \qquad
 \eta^2=12\gamma\wedge u\wedge w\wedge\delta.
\label{eq:bcpr-eta-square}
\end{equation}
The holomorphic Hodge rows supplied by
\cref{prop:period-input-audit,prop:very-general-torus-audit} are
\begin{equation}
 \sigma_1=6\mu\gamma+\tau u+w,
 \qquad
 \sigma_2=\beta\gamma+\mu u+\delta.
\label{eq:bcpr-period-rows}
\end{equation}
Thus no global Dolbeault or Bott--Chern generator below is denoted by
$u$ or $x$; those letters remain reserved for the printed fibre basis and
for local real linear functions.

\subsection{Exact fibre representatives}

Write
\begin{equation}
 y=\operatorname{Im}\tau,
 \qquad m=\operatorname{Im}\mu,
 \qquad r=\operatorname{Im}\beta,
 \qquad \Delta=r-\frac{6m^2}{y}<0.
\label{eq:bcpr-delta}
\end{equation}
The four marked lattice columns in the fibre coordinate
$\zeta=(\zeta_1,\zeta_2)$ are
\begin{equation}
 (6\mu,\beta),\quad(\tau,\mu),\quad(1,0),\quad(0,1),
\label{eq:bcpr-lattice-columns}
\end{equation}
in the order $\widehat\gamma,\widehat u,\widehat w,
\widehat\delta$.  On the universal cover of a marked fibre define the
real linear functions
\begin{equation}
 x_\gamma=
 \frac{y\,\operatorname{Im}\zeta_2-m\,\operatorname{Im}\zeta_1}
      {y\Delta},
 \qquad
 x_u=
 \frac{r\,\operatorname{Im}\zeta_1-6m\,\operatorname{Im}\zeta_2}
      {y\Delta}.
\label{eq:bcpr-real-linear-functions}
\end{equation}
For the first two columns in \cref{eq:bcpr-lattice-columns}, the numerators
of the changes in $x_\gamma$ are respectively
$yr-6m^2=y\Delta$ and $ym-my=0$, while the numerators of the changes in
$x_u$ are $6rm-6mr=0$ and $ry-6m^2=y\Delta$.  The last two columns are
real and change neither imaginary part.  Hence
\[
 \bigl(\Delta_{\widehat\gamma}x_\gamma,
       \Delta_{\widehat u}x_\gamma,
       \Delta_{\widehat w}x_\gamma,
       \Delta_{\widehat\delta}x_\gamma\bigr)=(1,0,0,0),
\]
\begin{equation}
 \bigl(\Delta_{\widehat\gamma}x_u,
       \Delta_{\widehat u}x_u,
       \Delta_{\widehat w}x_u,
       \Delta_{\widehat\delta}x_u\bigr)=(0,1,0,0).
\label{eq:bcpr-linear-function-increments}
\end{equation}
Their differentials, although not the functions themselves, therefore
descend to the fibre.  Their $(0,1)$ parts are
\begin{equation}
 \Gamma:=\bar\partial_Fx_\gamma
 =\frac{i}{2\Delta}
 \left(d\bar\zeta_2-\frac{m}{y}d\bar\zeta_1\right),
\label{eq:bcpr-gamma-form}
\end{equation}
\begin{equation}
 U:=\bar\partial_Fx_u
 =\frac{i}{2y\Delta}
 \left(r\,d\bar\zeta_1-6m\,d\bar\zeta_2\right).
\label{eq:bcpr-u-form}
\end{equation}
They represent respectively the $(0,1)$ images of the integral classes
$\gamma$ and $u$.

The identity
\begin{equation}
 \eta=u\wedge\sigma_1+6\gamma\wedge\sigma_2
     =-\sigma_1\wedge u-6\sigma_2\wedge\gamma
\label{eq:bcpr-eta-decomposition}
\end{equation}
holds in $\bigwedge^2V_{\mathbb C}$.  Since $\eta$ is of type $(1,1)$
on every marked fibre, \cref{eq:bcpr-eta-decomposition} gives its exact
Dolbeault representative
\begin{equation}
 Q_F=-d\zeta_1\wedge U-6d\zeta_2\wedge\Gamma.
\label{eq:bcpr-q-fibre}
\end{equation}
Equivalently, with the Appell--Humbert convention
$\eta=(2i)^{-1}\sum(H_\eta)_{ab}\sigma_a\wedge\bar\sigma_b$, the matrix
calculation in the proof of \cref{prop:very-general-torus-audit} gives
\begin{equation}
 H_\eta=\frac1{y\Delta}
 \begin{pmatrix}r&-6m\\-6m&6y\end{pmatrix},
\label{eq:bcpr-appell-humbert-matrix}
\end{equation}
and therefore
\begin{equation}
 Q_F=\frac1{2iy\Delta}\Bigl(
 r\,d\zeta_1\wedge d\bar\zeta_1
 -6m\,d\zeta_1\wedge d\bar\zeta_2
 -6m\,d\zeta_2\wedge d\bar\zeta_1
 +6y\,d\zeta_2\wedge d\bar\zeta_2\Bigr).
\label{eq:bcpr-q-fibre-expanded}
\end{equation}
Direct exterior multiplication, with holomorphic differentials ordered
before antiholomorphic differentials, gives
\begin{equation}
 Q_F\wedge\Gamma
 =\frac1{4y\Delta}
 d\zeta_1\wedge d\bar\zeta_1\wedge d\bar\zeta_2\neq0,
\label{eq:bcpr-q-gamma-product}
\end{equation}
\begin{equation}
 Q_F^2
 =\frac3{y\Delta}
 d\zeta_1\wedge d\zeta_2\wedge
 d\bar\zeta_1\wedge d\bar\zeta_2\neq0.
\label{eq:bcpr-q-square}
\end{equation}
Indeed,
$\det H_\eta=(6ry-36m^2)/(y^2\Delta^2)=6/(y\Delta)$,
and the coefficient in \cref{eq:bcpr-q-square} is
$\tfrac12\det H_\eta$ in the displayed ordering.  The identities
\Cref{eq:bcpr-q-gamma-product,eq:bcpr-q-square} will be used only in
fibre cohomology.

\subsection{Extension of the two Dolbeault inputs}

\begin{lemma}[the degree-one input]
\label{lem:bcpr-degree-one-input}
\Verified\ There is a distinguished class
$\mathfrak g\in H^{0,1}_{\bar\partial}(X)$ whose restriction to every
smooth marked fibre is represented by $\Gamma$.
\end{lemma}

\begin{proof}
By \cref{thm:structure-direct-images-audit}, the monodromy-invariant
integral class $\gamma$ maps to the nowhere-zero holomorphic section
$\gamma^{0,1}$ of
$R^1f_*\mathcal O_X\simeq\mathcal O\oplus\mathcal O(-1)$.  Coherent
Leray gives
\[
 H^1(X,\mathcal O_X)=H^0(B,R^1f_*\mathcal O_X)
 =\mathbb C\gamma^{0,1}.
\]
Let $\mathfrak g$ be the class mapping to $\gamma^{0,1}$.
\Cref{eq:bcpr-gamma-form} is the exact period-coordinate
representative of its fibre restriction.  This is the distinguished
coordinate of \cref{cor:picard-line-audit} before any analytic rescaling.
\end{proof}

\begin{lemma}[extension of the invariant $(1,1)$ class]
\label{lem:bcpr-eta-extension}
\Verified\ The fibre class $\eta$ extends to a nonzero holomorphic section
\begin{equation}
 \boldsymbol\eta\in
 H^0\bigl(B,R^1f_*\Omega_X^1\bigr).
\label{eq:bcpr-eta-global-section}
\end{equation}
Consequently there is a class
$\mathfrak q\in H^{1,1}_{\bar\partial}(X)$ whose coherent-Leray edge
value is $\boldsymbol\eta$.  For every smooth fibre inclusion
$i_b\colon F_b\to X$,
\begin{equation}
 i_b^*\mathfrak q=[Q_F]_{\bar\partial}
 \in H^{1,1}_{\bar\partial}(F_b).
\label{eq:bcpr-q-restriction}
\end{equation}
\end{lemma}

\begin{proof}
Over the smooth marked locus, $\eta$ is flat, monodromy invariant, and of
type $(1,1)$ identically.  Equation
\Cref{eq:bcpr-eta-decomposition} identifies the corresponding holomorphic
section of the relative group as
\begin{equation}
 -\sigma_1\wedge u^{0,1}-6\sigma_2\wedge\gamma^{0,1}.
\label{eq:bcpr-relative-eta-section}
\end{equation}

At a finite filling write $t=s^M$, with $(M,a)=(3,2)$ or $(4,1)$.
The covering family is smooth at $s=0$, and
\cref{eq:bcpr-eta-decomposition} holds identically there, so
\cref{eq:bcpr-relative-eta-section} remains holomorphic.  It is invariant
under the cyclic action.  In equivariant frames $e_0,e_1$ for the relative
cotangent and $b_0,b_1$ for $R^1\mathcal O$, the actual invariant
$\mathbb C\{t\}$-basis from
\cref{prop:log-relative-direct-images} is
\begin{equation}
 e_0b_0,\quad s^ae_0b_1,\quad
 s^{M-a}e_1b_0,\quad e_1b_1.
\label{eq:bcpr-finite-invariant-basis}
\end{equation}
Exactness of cyclic invariants in characteristic zero therefore makes
\cref{eq:bcpr-relative-eta-section} a holomorphic descended section at
both finite points.

At the cusp, $\gamma$ and $u$ extend as sections of $R^1f_*\mathbb C$,
and their $(0,1)$ images form a basis of $H^1(W,\mathcal O_W)$ by
\cref{thm:structure-direct-images-audit}.  On the dense torus,
\[
 \sigma_i=d\zeta_i=(2\pi i)^{-1}d\log x_i.
\]
The classes of these logarithmic forms extend in the relative logarithmic
sheaf as
\begin{equation}
 \mathcal E|_{N_0}
 =\mathcal O_{N_0}[d\log x_1]\oplus
  \mathcal O_{N_0}[d\log x_2]
\label{eq:bcpr-cusp-log-basis}
\end{equation}
by \cref{prop:log-relative-direct-images}.  A cusp quotient shear changes
$d\log x_i$ only by a holomorphic multiple of $d\log t$, which is zero in
the relative logarithmic sheaf.  Thus
\cref{eq:bcpr-relative-eta-section} also extends holomorphically through
the cusp.  This concerns classes in the relative logarithmic sheaf; it
does not assert that the $d\log x_i$ are global ordinary one-forms on $X$.

The Gauss--Manin/Kodaira--Spencer map
\[
 \kappa_1\colon R^1f_*\mathcal E\longrightarrow
 \Omega_B^1(\log\Sigma)\otimes R^2f_*\mathcal O_X
\]
annihilates \cref{eq:bcpr-relative-eta-section} on the smooth locus: a
flat class which remains of type $(1,1)$ has zero $(0,2)$ Higgs component.
Both the section and $\kappa_1$ extend holomorphically, so the identity
theorem gives $\kappa_1(\boldsymbol\eta)=0$ on $B$.  The ordinary-cotangent
calculation in
\cref{prop:log-ks-divisor,prop:absolute-cotangent-direct-images,thm:h12-scalar}
identifies
\begin{equation}
 R^1f_*\Omega_X^1=\ker\kappa_1=:\mathcal K,
 \qquad
 \mathcal K\simeq\mathcal O\oplus\mathcal O(-1)^{\oplus2}.
\label{eq:bcpr-cotangent-kernel}
\end{equation}
This proves \cref{eq:bcpr-eta-global-section}.  Its value on every smooth
fibre is nonzero in Dolbeault cohomology, not merely as a differential form:
by \cref{eq:bcpr-eta-square} its square is the integral top class
$12\gamma\wedge u\wedge w\wedge\delta$, and therefore
\begin{equation}
 \int_{(F_b,\mathrm{or}_{\rm mark})}Q_F^2=12,
 \qquad
 \int_{(F_b,\mathrm{or}_{\mathbb C})}Q_F^2=-12
\qquad(b\in B^\circ),
\label{eq:bcpr-q-square-integral}
\end{equation}
where $\mathrm{or}_{\rm mark}$ is the orientation determined by
$(\gamma,u,w,\delta)$ and $\mathrm{or}_{\mathbb C}$ is the complex
orientation.  They are opposite because the determinant of the four marked
lattice columns in real coordinates is $y\Delta<0$.  Thus
$[Q_F]_{\bar\partial}\ne0$, so the holomorphic
section $\boldsymbol\eta$ is nonzero.  Explicitly, if
$Q_F=\bar\partial_F\alpha$ for a $(1,0)$-form $\alpha$, then
$Q_F^2=\bar\partial_F(\alpha\wedge Q_F)$; integration and Stokes' theorem
would contradict \cref{eq:bcpr-q-square-integral}.  Here
$\alpha\wedge Q_F$ has bidegree $(2,1)$, so its $\partial_F$ is zero in
complex dimension two and the displayed $\bar\partial_F$-boundary is also a
$d$-boundary.

Finally, coherent Leray over the curve gives the exact edge sequence
\begin{equation}
 0\longrightarrow H^1(B,f_*\Omega_X^1)
 \longrightarrow H^1(X,\Omega_X^1)
 \longrightarrow H^0(B,R^1f_*\Omega_X^1)
 \longrightarrow0,
\label{eq:bcpr-cotangent-edge}
\end{equation}
where $f_*\Omega_X^1\simeq\mathcal O(-2)$ by
\cref{prop:absolute-cotangent-direct-images}.  Choose $\mathfrak q$ over
$\boldsymbol\eta$.  Base change on a smooth fibre gives
\cref{eq:bcpr-q-restriction}.  Any two choices differ by the
one-dimensional left term in \cref{eq:bcpr-cotangent-edge}.  That term
restricts to zero on every smooth fibre.  The nonzero class
$d_1\mathfrak g$ also restricts to zero, because its restriction is
$[\partial_F\Gamma]_{\bar\partial}=0$; since
$H^1(B,\mathcal O(-2))$ is one-dimensional, the left term is exactly
$\operatorname{im}(d_1\colon H^{0,1}_{\bar\partial}\to
H^{1,1}_{\bar\partial})$.
\end{proof}

\begin{remark}[what is and is not explicit]
\label{rem:bcpr-literal-gauge}
Choose global $\bar\partial$-closed forms
\begin{equation}
 G\in\mathcal A^{0,1}(X),\qquad
 Q\in\mathcal A^{1,1}(X)
\label{eq:bcpr-global-representatives}
\end{equation}
representing $\mathfrak g$ and $\mathfrak q$.
\Cref{eq:bcpr-q-restriction} says that
$i_b^*[Q]_{\bar\partial}$ is the Dolbeault class of the explicit
translation-invariant
form displayed in \cref{eq:bcpr-q-fibre-expanded}; it does not say that an arbitrary
representative $Q$ restricts literally to that form.  A literal global
coefficient formula would require additional $C^\infty$ gauge choices and
horizontal correction forms on the disc--pair-of-pants cover, compatible
with the two affine cyclic conjugacies and the toric cusp.  Those choices
are not part of the canonical data.  No such literal gauge is needed for
the class and product calculations below.
\end{remark}

\subsection{The factor \texorpdfstring{$12$}{12} and the Bott--Chern first Chern class}

Let
\[
 \exp\colon H^1(X,\mathcal O_X)\longrightarrow\operatorname{Pic}(X)
\]
be induced by $z\mapsto e^{2\pi iz}$.  By
\cref{prop:leray-e2},
\begin{equation}
 H^0(B,R^1f_*\mathbb Z)=\mathbb Z(12\gamma).
\label{eq:bcpr-global-integral-line}
\end{equation}
The exponential--Leray comparison gives a line bundle $M$ on $B$ such
that
\begin{equation}
 \exp(12\mathfrak g)=f^*M,
 \qquad c_1(M)=d_2(12\gamma).
\label{eq:bcpr-exponential-transgression}
\end{equation}
For completeness, the first assertion follows because
\cref{eq:bcpr-global-integral-line} says that the restriction of
$12\mathfrak g$ to every fibre lies in the image of integral cohomology
and hence exponentiates to the trivial bundle.  The low-degree Leray
sequence for $\mathcal O^*$ then makes the resulting bundle a pullback.
On a disc--pair-of-pants cover of $B$, choose local integral lifts of
$12\gamma$.  Their differences on the annular overlaps form a Cech
one-cochain, and its Cech coboundary is simultaneously the two-cocycle
defining the Leray transgression and the two-cocycle of logarithms defining
$c_1(M)$.  This proves the second assertion with the same sign convention;
it is the exponential-sequence version of the local-lift calculation in
the proof of \cref{lem:first-transgression}.

Orient the generator $\omega\in H^2(B;\mathbb Z)$ by
\begin{equation}
 \omega=c_1(\mathcal O_B(1)).
\label{eq:bcpr-base-orientation}
\end{equation}
Then \cref{lem:first-transgression} and
\cref{eq:bcpr-exponential-transgression} give
\begin{equation}
 \exp(12\mathfrak g)=f^*\mathcal O_B(1).
\label{eq:bcpr-factor-twelve}
\end{equation}
If the opposite orientation is used for the primitive generator in
\cref{lem:first-transgression}, the corresponding statement is
$\exp(-12\mathfrak g)=f^*\mathcal O_B(1)$.  Thus the canonical integral
data fixes the absolute factor $12$, while its sign is the orientation
choice fixed in \cref{eq:bcpr-base-orientation}.

Put
\begin{equation}
 a=\partial G,
 \qquad a^\dagger=\bar\partial\bar G,
 \qquad
 \mathsf A=[a]_{BC},
 \qquad \mathsf B=[a^\dagger]_{BC}.
\label{eq:bcpr-a-b-definitions}
\end{equation}
Both forms are $d$-closed.  We now compute their Chern-class normalization
without choosing local metric weights.  For $z\in\mathbb C$, equip the
smooth trivial line with a global frame $s$ and the holomorphic structure
\begin{equation}
 \bar\partial_zs=-2\pi i zG\otimes s.
\label{eq:bcpr-holomorphic-structure}
\end{equation}
It is integrable because $\bar\partial G=0$.  If
$G=\bar\partial\phi_i$ on a good cover, then
$e_i=e^{2\pi iz\phi_i}s$ is $\bar\partial_z$-holomorphic and
\[
 e_j=e^{2\pi iz(\phi_j-\phi_i)}e_i.
\]
The holomorphic cocycle $\phi_j-\phi_i$ represents $\mathfrak g$ under
the Cech--Dolbeault comparison, so
\cref{eq:bcpr-holomorphic-structure} is precisely the line bundle
$\exp(z\mathfrak g)$ for the displayed exponential convention.

Give $s$ unit length.  The unit-metric Chern connection is
\begin{equation}
 \nabla_z=d-2\pi i\bigl(zG+\bar z\,\bar G\bigr).
\label{eq:bcpr-unit-chern-connection}
\end{equation}
Indeed, its $(0,1)$ part is
\cref{eq:bcpr-holomorphic-structure}, and its connection one-form is
purely imaginary, so it is metric compatible; uniqueness identifies it
with the Chern connection.  Since $\bar\partial G=0$ and
$\partial\bar G=0$, its curvature is the global $(1,1)$ form
\begin{equation}
 F_{\nabla_z}
 =-2\pi i\bigl(z\,\partial G+
                 \bar z\,\bar\partial\bar G\bigr)
 =-2\pi i\bigl(za+\bar z\,a^\dagger\bigr).
\label{eq:bcpr-unit-chern-curvature}
\end{equation}
With the Bott--Chern convention
\begin{equation}
 c_{1,BC}(L,h)=
 \left[\frac{i}{2\pi}F_{\nabla^h}\right]_{BC}
\label{eq:bcpr-chern-convention}
\end{equation}
of \cite[\S8.a]{Schweitzer2007},
\cref{eq:bcpr-unit-chern-curvature} gives, for every $z\in\mathbb C$,
\begin{equation}
 c_{1,BC}\bigl(\exp(z\mathfrak g)\bigr)
 =z\mathsf A+\bar z\,\mathsf B.
\label{eq:bcpr-chern-exp-identity}
\end{equation}
Consequently, if $\varepsilon=1$ for the orientation fixed by
\cref{eq:bcpr-base-orientation} and $\varepsilon=-1$ for its opposite,
then
\begin{equation}
 c_{1,BC}\bigl(f^*\mathcal O_B(1)\bigr)
 =12\varepsilon\bigl(\mathsf A+\mathsf B\bigr).
\label{eq:bcpr-c1-pullback}
\end{equation}
This is the convention-consistent form of the factor-$12$ identity; no
analytic rescaling of the Picard coordinate changes its integral content.

\subsection{All Bott--Chern generators}

Using \cref{eq:bcpr-global-representatives}, define the following
$d$-closed forms and their classes:
\begin{equation}
\begin{array}{c|c|c}
\text{class}&\text{representative}&\text{bidegree}\\ \hline
\mathsf A&a=\partial G&(1,1)\\
\mathsf B&a^\dagger=\bar\partial\bar G&(1,1)\\
\mathsf C&c=\partial Q&(2,1)\\
\overline{\mathsf C}&c^\dagger=\bar\partial\bar Q&(1,2)\\
\mathsf D&d_0=\partial(G\wedge Q)&(2,2)\\
\overline{\mathsf D}&d_0^\dagger=
 \bar\partial(\bar G\wedge\bar Q)&(2,2)\\
\mathsf E&e=\partial(Q^2)&(3,2)\\
\overline{\mathsf E}&e^\dagger=\bar\partial(\bar Q^2)&(2,3)\\
\mathsf T&\nu,\quad\int_X\nu=1&(3,3).
\end{array}
\label{eq:bcpr-generators}
\end{equation}
Here $\nu$ is any smooth positive top-degree form divided by its integral.
The chosen $Q$ cannot be real: its fibre class is real, but the proof below
shows that $\partial Q$ is nonzero globally, whereas a real
$\bar\partial$-closed $Q$ would also be $\partial$-closed.

\begin{proposition}[nonvanishing and bases]
\label{prop:bcpr-generators-are-bases}
\Verified\ Together with the unit $\mathbf1\in H^{0,0}_{BC}(X)$, the entries
in \cref{eq:bcpr-generators} give bases of all the nonzero Bott--Chern groups
in \cref{eq:bcpr-nonzero-dimensions}:
$\mathsf A,\mathsf B$ form the basis in bidegree $(1,1)$ and
$\mathsf D,\overline{\mathsf D}$ form the basis in bidegree $(2,2)$;
each other listed bidegree is one-dimensional and is generated by its
displayed class.
\end{proposition}

\begin{proof}
Because $X$ is compact and connected, a smooth function annihilated by both
$\partial$ and $\bar\partial$ is constant.  There is no
$\partial\bar\partial$ image in bidegree $(0,0)$, so
$\mathbf1:=[1]_{BC}$ generates
$H^{0,0}_{BC}(X)=\mathbb C\mathbf1$.

The Dolbeault image of $\mathsf A$ is
$d_1\mathfrak g\neq0$ by injectivity in
\cref{eq:bcpr-hodge-row-one}.  Thus $\partial G$ is a nonzero
Bott--Chern class.  The Dolbeault image of $\mathsf B$ is zero because
$a^\dagger=\bar\partial\bar G$, while
$\mathsf B=\overline{\mathsf A}\neq0$.  Hence
$\mathsf A,\mathsf B$ are independent, and
\cref{eq:bcpr-nonzero-dimensions} makes them a basis.

The classes $d_1\mathfrak g$ and $\mathfrak q$ are independent in
$H^{1,1}_{\bar\partial}(X)$: the first restricts to zero on every smooth
torus fibre, while the second restricts to the nonzero class
by \cref{eq:bcpr-q-restriction,eq:bcpr-q-square-integral}.  Exactness of
\cref{eq:bcpr-hodge-row-one} therefore gives
\begin{equation}
 d_1\mathfrak q=[\partial Q]_{\bar\partial}\neq0.
\label{eq:bcpr-c-nonzero}
\end{equation}
Thus $\mathsf C\neq0$, and conjugation gives
$\overline{\mathsf C}\neq0$.

Set
\begin{equation}
 v=[G\wedge Q]_{\bar\partial}
 \in H^{1,2}_{\bar\partial}(X).
\label{eq:bcpr-v-class}
\end{equation}
Its fibre restriction is the class represented by
$\Gamma\wedge Q_F$.  This class is nonzero, rather than merely represented
by a nonzero invariant form: wedging with the holomorphic one-form
$d\zeta_2$ and using \cref{eq:bcpr-q-gamma-product} gives
\begin{equation}
 \int_{F_b}Q_F\wedge\Gamma\wedge d\zeta_2
 =\frac1{4y\Delta}
   \int_{F_b}d\zeta_1\wedge d\zeta_2\wedge
                    d\bar\zeta_1\wedge d\bar\zeta_2
 =-1
\label{eq:bcpr-q-gamma-serre-pairing}
\end{equation}
for the complex orientation: the middle integral equals
$-4y\Delta>0$.  Thus the displayed constant top form has nonzero integral.
Dolbeault--Serre duality on $F_b$ therefore gives
$[\Gamma\wedge Q_F]_{\bar\partial}\ne0$.  Equivalently, if
$Q_F\wedge\Gamma=\bar\partial_F\Psi$ with $\Psi$ of bidegree $(1,1)$,
then, since $\bar\partial_Fd\zeta_2=0$,
\[
 Q_F\wedge\Gamma\wedge d\zeta_2
 =\bar\partial_F(\Psi\wedge d\zeta_2),
\]
and Stokes' theorem would make the integral in
\cref{eq:bcpr-q-gamma-serre-pairing} zero.  Here
$Q_F\wedge\Gamma$ has bidegree $(1,2)$,
$d\zeta_2$ has bidegree $(1,0)$, and
$\Psi\wedge d\zeta_2$ has bidegree $(2,1)$; its $\partial_F$ is zero in
complex dimension two, so the displayed $\bar\partial_F$-boundary is a
$d$-boundary.  Thus $v\neq0$, and it spans
$H^{1,2}_{\bar\partial}(X)$.  Injectivity in
\cref{eq:bcpr-hodge-row-two} gives
\begin{equation}
 d_1v=[\partial(G\wedge Q)]_{\bar\partial}\neq0.
\label{eq:bcpr-d-nonzero}
\end{equation}
Hence $\mathsf D\neq0$, and conjugation gives
$\overline{\mathsf D}\neq0$.  The latter has zero Dolbeault image because
its representative is $\bar\partial$-exact, so it is independent of
$\mathsf D$; \cref{eq:bcpr-nonzero-dimensions} gives the asserted
$(2,2)$ basis.

The class $[Q^2]_{\bar\partial}$ has nonzero fibre restriction by
\cref{eq:bcpr-q-square-integral}.  Since $v$ spans
$H^{1,2}_{\bar\partial}(X)$, the image of $d_1$ is spanned by $d_1v$.
On a smooth fibre,
\[
 i_b^*(d_1v)
 =d_1[\Gamma\wedge Q_F]_{\bar\partial}
 =[\partial_F(\Gamma\wedge Q_F)]_{\bar\partial}=0,
\]
because $\Gamma$ and $Q_F$ have constant fibre coefficients.  Thus every
class in $\operatorname{im}(d_1\colon H^{1,2}_{\bar\partial}\to
H^{2,2}_{\bar\partial})$ restricts to zero.
Therefore $[Q^2]_{\bar\partial}$ is independent of $d_1v$.  Exactness of
\cref{eq:bcpr-hodge-row-two} yields
\begin{equation}
 d_1[Q^2]_{\bar\partial}
 =[\partial(Q^2)]_{\bar\partial}\neq0.
\label{eq:bcpr-e-nonzero}
\end{equation}
Hence $\mathsf E\neq0$, and conjugation gives
$\overline{\mathsf E}\neq0$.  Finally, $\int_X\nu=1$ proves
$\mathsf T\neq0$.  The dimensions in
\cref{eq:bcpr-nonzero-dimensions} complete the proof.
\end{proof}

\subsection{Every Bott--Chern product}

\begin{theorem}[square-zero positive ideal]
\label{thm:bcpr-square-zero-algebra}
\Verified\ Let $N\subset H^{*,*}_{BC}(X)$ be the span of the nine positive-degree
generators in \cref{eq:bcpr-generators}.  Then
\begin{equation}
 H^{*,*}_{BC}(X)=\mathbb C\mathbf1\oplus N,
 \qquad N^2=0.
\label{eq:bcpr-algebra}
\end{equation}
Equivalently, every product of two positive-degree Bott--Chern classes is
zero.
\end{theorem}

\begin{proof}
We first calculate the products landing below top degree.  Since
$\bar\partial G=\bar\partial Q=0$ and, after conjugation,
$\partial\bar G=\partial\bar Q=0$, the graded Leibniz rule gives the exact
identities
\begin{equation}
 \partial\bar\partial(G\wedge\bar G)
 =-a\wedge a^\dagger,
\label{eq:bcpr-product-ab}
\end{equation}
\begin{equation}
 \partial\bar\partial(G\wedge\bar Q)
 =-a\wedge c^\dagger,
 \qquad
 \partial\bar\partial(\bar G\wedge Q)
 =a^\dagger\wedge c.
\label{eq:bcpr-product-acbar-bc}
\end{equation}
Consequently
\begin{equation}
 \mathsf A\mathsf B=0,\qquad
 \mathsf A\overline{\mathsf C}=0,\qquad
 \mathsf B\mathsf C=0.
\label{eq:bcpr-first-zero-products}
\end{equation}

Choose the pullback Fubini--Study metric on
$f^*\mathcal O_B(1)$.  Its Chern form is the pullback of a $(1,1)$ form on
the curve $B$, so its square vanishes pointwise.  Squaring
\cref{eq:bcpr-c1-pullback} and using
\cref{eq:bcpr-first-zero-products} gives
\begin{equation}
 \mathsf A^2+\mathsf B^2=0.
\label{eq:bcpr-sum-of-squares}
\end{equation}
\Cref{eq:bcpr-sum-of-squares} alone does not separate the two
squares.  Apply $\pi_{\bar\partial}^{2,2}$.  The class
$\mathsf B^2$ has zero Dolbeault image because
\begin{equation}
 (a^\dagger)^2
 =\bar\partial(\bar G\wedge a^\dagger).
\label{eq:bcpr-b-square-dolbeault-zero}
\end{equation}
It follows from \cref{eq:bcpr-sum-of-squares} that
$[a^2]_{\bar\partial}=0$.  On the other hand,
\begin{equation}
 [a^2]_{\bar\partial}
 =d_1[G\wedge a]_{\bar\partial},
 \qquad
 \partial(G\wedge a)=a^2.
\label{eq:bcpr-a-square-d1}
\end{equation}
The first arrow in \cref{eq:bcpr-hodge-row-two} is injective, so
\cref{eq:bcpr-a-square-d1} implies
$[G\wedge a]_{\bar\partial}=0$.  There is therefore a form
$\Phi\in\mathcal A^{1,1}(X)$ such that
\begin{equation}
 G\wedge a=\bar\partial\Phi.
\label{eq:bcpr-phi-primitive}
\end{equation}
Applying $\partial$ gives
\begin{equation}
 a^2=\partial\bar\partial\Phi,
\label{eq:bcpr-a-square-bc-zero}
\end{equation}
and hence $\mathsf A^2=0$.  Conjugation gives
$\mathsf B^2=0$.

It remains among the products below top bidegree to prove
$\mathsf A\mathsf C=0$; its conjugate will give
$\mathsf B\overline{\mathsf C}=0$.  First, the natural map
\begin{equation}
 \pi_{\bar\partial}^{3,2}\colon H^{3,2}_{BC}(X)
 \longrightarrow H^{3,2}_{\bar\partial}(X)
\label{eq:bcpr-bc-to-dolbeault-32}
\end{equation}
is an isomorphism.  Indeed, both spaces have dimension one by
\cref{eq:bcpr-nonzero-dimensions,cor:exact-hodge-diamond}, and
\cref{eq:bcpr-e-nonzero} says that $\mathsf E$ has nonzero image.
Pair $\pi_{\bar\partial}^{3,2}(\mathsf A\mathsf C)$ with the generator
$[G]_{\bar\partial}\in H^{0,1}_{\bar\partial}(X)$.  By
\cref{eq:bcpr-phi-primitive}, $\bar\partial c=0$, and Stokes' theorem,
\begin{equation}
 \int_Xa\wedge c\wedge G
 =-\int_XG\wedge a\wedge c
 =-\int_X\bar\partial\Phi\wedge c
 =-\int_X\bar\partial(\Phi\wedge c)=0.
\label{eq:bcpr-serre-pairing-zero}
\end{equation}
Dolbeault--Serre duality makes
\begin{equation}
 H^{3,2}_{\bar\partial}(X)\times
 H^{0,1}_{\bar\partial}(X)\longrightarrow\mathbb C,
 \qquad
 ([\alpha],[\rho])\longmapsto\int_X\alpha\wedge\rho,
\label{eq:bcpr-serre-pairing}
\end{equation}
a perfect pairing.  Since
$H^{0,1}_{\bar\partial}(X)=\mathbb C[G]_{\bar\partial}$,
\cref{eq:bcpr-serre-pairing-zero} implies
$\pi_{\bar\partial}^{3,2}(\mathsf A\mathsf C)=0$; injectivity in
\cref{eq:bcpr-bc-to-dolbeault-32} gives
\begin{equation}
 \mathsf A\mathsf C=0,
 \qquad
 \mathsf B\overline{\mathsf C}=0.
\label{eq:bcpr-second-zero-products}
\end{equation}
This step uses Dolbeault--Serre duality together with the independently
proved injectivity of \cref{eq:bcpr-bc-to-dolbeault-32}.  The
Bott--Chern--Aeppli pairing of
\cite[\S2.c]{Schweitzer2007} by itself would supply a perfect pairing but
would not annihilate this product.

The only products not yet considered which can have bidegree at most
$(3,3)$ are
\begin{equation}
 \mathsf A\mathsf D,\quad
 \mathsf A\overline{\mathsf D},\quad
 \mathsf B\mathsf D,\quad
 \mathsf B\overline{\mathsf D},\quad
 \mathsf C\overline{\mathsf C},
\label{eq:bcpr-top-products}
\end{equation}
all of bidegree $(3,3)$.  Every non-top generator in
\cref{eq:bcpr-generators} is de Rham exact:
\begin{equation}
\begin{gathered}
 a=dG,\qquad a^\dagger=d\bar G,\qquad
 c=dQ,\qquad c^\dagger=d\bar Q,\\
 d_0=d(G\wedge Q),\qquad
 d_0^\dagger=d(\bar G\wedge\bar Q),\qquad
 e=d(Q^2),\qquad e^\dagger=d(\bar Q^2).
\end{gathered}
\label{eq:bcpr-de-rham-primitives}
\end{equation}
Each factor paired with an exact factor is $d$-closed, so every product in
\cref{eq:bcpr-top-products} is $d$-exact and has integral zero.  Integration
induces an isomorphism
\begin{equation}
 H^{3,3}_{BC}(X)\xrightarrow{\ \int_X\ }\mathbb C.
\label{eq:bcpr-top-integration}
\end{equation}
It is well defined by Stokes' theorem, the source has dimension one by
\cref{eq:bcpr-nonzero-dimensions}, and it is nonzero on $\mathsf T$.
Therefore all five products in \cref{eq:bcpr-top-products} vanish in
Bott--Chern cohomology.

All other products of positive-degree generators have $p>3$ or $q>3$ and
vanish already as differential forms.
\Cref{eq:bcpr-first-zero-products,eq:bcpr-a-square-bc-zero,eq:bcpr-second-zero-products,eq:bcpr-top-products,eq:bcpr-top-integration}
exhaust every possible positive-degree product
and prove \cref{eq:bcpr-algebra}.
\end{proof}

\begin{remark}[literal global gauge boundary]
\label{rem:bcpr-global-gauge-boundary}
The Bott--Chern algebra \cref{eq:bcpr-algebra} and the fibre-coordinate
restrictions displayed in
\cref{eq:bcpr-gamma-form,eq:bcpr-u-form,eq:bcpr-q-fibre,eq:bcpr-q-fibre-expanded,eq:bcpr-q-gamma-product,eq:bcpr-q-square}
are complete.  The canonical data do not select a preferred literal global
coefficient atlas for the smooth forms $G$ and $Q$.  Such an atlas would
require explicit smooth local trivializations and horizontal correction
functions through both finite affine fillings and the cusp.  After fixing
$\mathfrak g$ and one coherent-Leray lift $\mathfrak q$, changing those
smooth gauges changes $G$ and $Q$ by Dolbeault coboundaries but changes
none of the cohomology classes, nonvanishing arguments, or products proved
above.  This is a limitation only on a preferred global coefficient
presentation, not on the representatives as cohomology classes or on the
Bott--Chern algebra.
\end{remark}


\begin{auditrecord}[H2: independent replication of the cotangent calculation]
\label{record:h12-direct-images}
\Open\ A further context-isolated AI derivation has independently recovered
the cyclic invariant lattices (11.4), the determinant divisor (11.6), the
residue cancellation (11.17), the scalar test (11.19)--(11.20), and the
splittings
\[
 R^0f_*\Omega_X^1=\mathcal O(-2),\quad
 R^1f_*\Omega_X^1=\mathcal O\oplus\mathcal O(-1)^{\oplus2},\quad
 R^2f_*\Omega_X^1=\mathcal O\oplus\mathcal O(-1).
\]
It also independently gives $h^{1,2}=1$ and $h^{1,1}=2$.  Independent human
specialist replication remains outstanding, so the record stays \Open; this
external-replication boundary is not an unresolved premise of the theorem
above.
\end{auditrecord}

\section{The alternate collapse and Mayer--Vietoris route}

The nearby-cycle/Leray proof in \cref{sec:integral-leray} gives one verified
route to the topology of $X$.  This section audits the manuscript's independent
cellular-collapse/Mayer--Vietoris route.  It now closes the entire alternate
certificate: the equivariant cusp retraction, its primitive integral cellular
map, equality with analytic nearby-cycle specialization, both geometrically
realized Mayer--Vietoris stages, and the comparison of their obstruction maps
with Leray.

We first correct a conflation in the earlier audit.  For any two-set
decomposition $Y=A\cup B$, our Mayer--Vietoris convention is
\[
\Delta_q=(i_{A*},-i_{B*}):
H_q(A\cap B;\mathbb Z)\longrightarrow
H_q(A;\mathbb Z)\oplus H_q(B;\mathbb Z).
\tag{4.1}
\]
The manuscript does not identify its printed $\alpha_q$ with a one-shot map
for three collars.  It uses two stages.  First,
\[
X^\circ=X_1\cup X_2,
\qquad X_1\cap X_2\simeq F,
\qquad
\alpha_q=(\pi_{1*},-\pi_{2*}).
\tag{4.2}
\]
Second, it attaches the cusp through
\[
X=N'_0\cup X^\circ,
\qquad N'_0\cap X^\circ=M,
\qquad
K_q=\ker\bigl(H_q(M)\to H_q(W)\bigr).
\tag{4.3}
\]
The proof below realizes both stages by actual oriented inclusions and then
also records the equivalent one-shot degree-one collar matrix.

\subsection{The abstract cusp-retraction lemma}

\begin{proposition}[equivariant topological spine]\label{prop:topological-spine}
\Verified\ Let $A$ be a contractible manifold with boundary and let
$A_0\subset\partial A$ be a contractible boundary component.  Let a
group $\Gamma$ act freely and properly discontinuously on $(A,A_0)$.  Assume
$Q=A/\Gamma$ and $H=A_0/\Gamma$ are manifolds, $Q$ is compact,
$H\hookrightarrow Q$ is collared, and the two quotient maps are connected
coverings with deck group $\Gamma$.  If the inclusion identifies the two
deck groups by the identity automorphism of $\Gamma$, then
$H\hookrightarrow Q$ is a strong deformation retract.  The retraction lifts
uniquely from the identity lift to a $\Gamma$-equivariant strong deformation
retraction of $A$ onto $A_0$.
\end{proposition}

\begin{proof}
Because $A$ and $A_0$ are contractible universal covers,
$\pi_n(Q)=\pi_n(H)=0$ for $n\ge2$, while
$\pi_1(Q)\cong\Gamma\cong\pi_1(H)$.  By hypothesis the inclusion induces the
identity on $\Gamma$, so it is a weak homotopy equivalence.  Manifolds have
the homotopy type of CW complexes; Whitehead's theorem
\cite[Theorem~4.5]{Hatcher2002} therefore makes the inclusion a homotopy
equivalence.

The collar makes $(Q,H)$ an NDR pair, hence the inclusion is a cofibration.
We record the standard cofibration argument.  Choose a homotopy inverse
$r_0:Q\to H$.  Since $r_0|_H\simeq\operatorname{id}_H$, the homotopy
extension property modifies $r_0$ through maps $Q\to H$ to a retraction
$r$ satisfying $r|_H=\operatorname{id}_H$.  The maps
$\operatorname{id}_Q$ and $i\circ r$ are homotopic and agree on $H$.
Applying the homotopy extension property once more makes that homotopy
stationary on $H$.  It is therefore a strong deformation retraction.

Let $P:A\to Q$ be the covering.  Lift the deformation homotopy from the
initial lift $\operatorname{id}_A$.  Since the base homotopy fixes $H$ and
the initial lift fixes $A_0$, uniqueness of path lifting makes the lifted
homotopy fix $A_0$.  For $\gamma\in\Gamma$, both
$\widetilde H_t\circ\gamma$ and $\gamma\circ\widetilde H_t$ lift the same
map and agree at $t=0$; uniqueness gives equality for every $t$.  The lift is
therefore $\Gamma$-equivariant and ends in $A_0$.
\end{proof}

\subsection{The candidate-specific equivariant collapse}

We now verify every hypothesis of \cref{prop:topological-spine} in the
candidate.  Retain the toric threefold $Y$, the compact torus
\[
 T_c=(S^1)^3,\qquad T_f=(S^1)^2\times\{1\}\subset T_c,
\]
and the free properly discontinuous action
$\Gamma=\{\Psi_{\bar\lambda}:\bar\lambda\in\bar\Lambda\}$ from
\cref{thm:candidate-cusp-geometry}.  Put
\[
 Y_{\leq\eta}=\{|t|\leq\eta\},\qquad
 Y_{\geq0}=\overline{(\mathbb R_{>0})^3}\subset Y,\qquad
 A=(Y_{\leq\eta})_{\geq0},\qquad A_0=(Y_0)_{\geq0}.
\tag{C.1}
\]
For a torus point write $y=\log|x|/\log|t|$.  The period matrix on a
punctured cusp fibre is
\[
 Z_t=sB_0+C(t),\qquad s=\frac{\log t}{2\pi i},\qquad
 B_t=B_0+\frac{R(t)}{\log|t|},\quad R(t)=-2\pi\operatorname{Im}C(t).
\tag{C.2}
\]
Here $B_0:\bar\Lambda\to\Lambda_{\rm tor}$ is unimodular.

\begin{lemma}[straightening the twist]\label{lem:collapse-straightening}
\Verified\ For all sufficiently small $\eta>0$ there is a homeomorphism
$G:Y_{\leq\eta}\to Y_{\leq\eta}$, fixing $Y_0$ pointwise and preserving
$t$, which conjugates $\Gamma$ to the action $\Gamma^0$ obtained by replacing
$C(t)$ by $C_0=C(0)$.  It commutes with $T_f$ and, on every smooth quotient
fibre, induces a real-linear isomorphism which is the identity on the marked
lattice $\Lambda$.
\end{lemma}

\begin{proof}
Shrink $\eta$ until $B_t$ and
$B_t^0=B_0+R(0)/\log|t|$ are invertible for $0<|t|\leq\eta$.  In logarithmic
coordinates $\xi=(2\pi i)^{-1}\log x$ define
\[
 L_t(\xi)=\xi-(C(t)-C_0)(\operatorname{Im}Z_t)^{-1}
                         \operatorname{Im}\xi.
\tag{C.3}
\]
Every $\xi$ has the unique expression $\xi=Z_ta+b$, with
$a,b\in\mathbb R^2$, and then $L_t(\xi)=Z_t^0a+b$.  Thus $L_t$ is invertible,
fixes $\mathbb R^2$ and its sublattice $\mathbb Z^2$, and sends each deck
translation $Z_t\bar\lambda$ to $Z_t^0\bar\lambda$.  It therefore descends
on $Y_t=\mathbb C^2/\mathbb Z^2$ and conjugates the two $\bar\Lambda$-actions.
In toric coordinates it is
\[
 G(x,t)=\bigl(\exp(-2\pi i(C(t)-C_0)B_t^{-1}y)x,t\bigr),
 \qquad G|_{Y_0}=\operatorname{id}.
\tag{C.4}
\]
It commutes with $T_f$ because $y$ is $T_f$-invariant.

Only continuity at $Y_0$ is nonformal.  Around any $p\in Y_0$, choose one of
the bounded-position subcharts used in the properness proof of
\cref{thm:candidate-cusp-geometry}.  On that subchart $y$ stays bounded,
$B_t^{-1}$ stays bounded, and $C(t)-C_0\to0$.  Hence the torus factor in
(C.4) tends to $1$.  The same argument applies to the inverse, obtained by
interchanging $C(t)$ and $C_0$.  Thus $G$ is a homeomorphism.  Because it
fixes the real lattice and sends the two period columns $Z_t\bar\lambda$ to
$Z_t^0\bar\lambda$, it induces the identity on the marked $H_1=\Lambda$.
\end{proof}

Write $c^0_{\bar\lambda}=u_{\bar\lambda}a_{\bar\lambda}$ for its polar
decomposition, with $u_{\bar\lambda}\in T_f$ and
$a_{\bar\lambda}=\exp(R(0)\bar\lambda)\in(\mathbb R_{>0})^2$.  Let
$\Gamma^+$ be the positive action
\[
 \Psi^+_{\bar\lambda}=(a_{\bar\lambda},1)\Phi_{\bar\lambda}.
\tag{C.5}
\]

\begin{lemma}[positive quotient and honeycomb]\label{lem:positive-collapse}
\Verified\ The pair $(A,A_0)$ in (C.1) consists of contractible topological
manifolds, $A_0$ being a boundary component of $A$.  The action $\Gamma^+$ is
free and properly discontinuous, and
\[
 Q^+=A/\Gamma^+,\qquad H=A_0/\Gamma^+
\tag{C.6}
\]
are respectively a compact $3$-manifold with boundary and one of its boundary
$2$-tori.  The inclusion $H\hookrightarrow Q^+$ is a strong deformation
retract, and its retraction lifts to a $\Gamma^+$-equivariant strong
deformation retraction $P$ of $A$ onto $A_0$.
\end{lemma}

\begin{proof}
Polar multiplication $m:T_c\times Y_{\geq0}\to Y$ is a closed quotient map:
in every smooth unimodular toric chart it is the usual map
$(S^1)^3\times[0,\infty)^3\to\mathbb C^3$, and $T_c$ is compact.  If $q$ is
in the orbit indexed by a face $C$ of the $A_2$ triangulation, its stabilizer
is the connected compact subtorus on
\[
 N_{\widehat C}=\bigoplus_{v\in V(C)}\mathbb Z(v,1).
\]
Unimodularity makes this lattice saturated, and direct intersection with
$\Lambda_{\rm tor}$ gives
\[
 \operatorname{Stab}_{T_c}(q)\cap T_f=
 \begin{cases}
  1,&C\text{ is a vertex},\\
  S^1_{v'-v},&C=\{v,v'\}\text{ is an edge},\\
  T_f,&C\text{ is a triangle}.
 \end{cases}
\tag{C.7}
\]
There are no hidden finite components.  The lattices
$N_{\widehat C}$ and $\Lambda_{\rm tor}$ are saturated, and every face contains
a vertex $(v,1)$, so
$N_{\widehat C}+\Lambda_{\rm tor}=N'$.  Thus the displayed intersections are
the connected subtori attached to their primitive intersection lattices.

In a toric chart the positive set $A$ is
\[
 \{(r_0,r_1,r_2)\in[0,\infty)^3:r_0r_1r_2\leq\eta\}.
\]
Its corner links are discs, so it is a topological $3$-manifold with boundary
$A_0\sqcup\{|t|=\eta\}_{>0}$.  Its interior is
$\mathbb R_y^2\times(0,\eta)$.  Brown's collaring theorem
\cite{Brown1962} makes the inclusion of the interior a homotopy equivalence;
therefore $A$ is contractible.

The nonnegative part of each component $E_v\simeq\mathrm{dP}_6$ is its moment
hexagon.  These hexagons meet along exactly the arcs prescribed by the dual
honeycomb $\mathcal H$ of the $A_2$ triangulation.  Map the two vertex orbits
of $\mathcal H$ to the corresponding torus fixed points, choose homeomorphisms
on representatives of the three edge orbits, and extend the resulting circle
homeomorphism over one hexagon.  Propagating by $B_0\bar\Lambda$ defines a
$\Gamma^+$-equivariant continuous bijection
\[
 \vartheta:\mathbb R^2\xrightarrow{\sim}A_0,\qquad
 \vartheta(z+B_0\bar\lambda)=
 \Psi^+_{\bar\lambda}\vartheta(z).
\tag{C.8}
\]
Local finiteness of the fan makes this bijection proper, hence a homeomorphism.
Thus $A_0$ is contractible.

The action in (C.5) is the already verified cusp action with constant twist
$C^+=(2\pi i)^{-1}R(0)$, since
$e^{2\pi iC^+\bar\lambda}=a_{\bar\lambda}$; hence it is free and properly
discontinuous.  Properness of its quotient over
the disc, proved by the same bounded fundamental-domain argument as
\cref{thm:candidate-cusp-geometry}, makes $Q^+$ compact.  Its two boundary
components are $H\simeq\mathbb R^2/B_0\bar\Lambda$ and
$\mathbb R^2/B_\eta^+\bar\Lambda$, both tori.  Since a free properly
discontinuous action preserves the boundary charts, $Q^+$ is a topological
$3$-manifold with boundary; Brown's collaring theorem makes
$H\hookrightarrow Q^+$ collared.  The covers $A\to Q^+$ and
$A_0\to H$ have the same deck group, and the equivariant inclusion identifies
that group by the identity.  The result is now
\cref{prop:topological-spine}.
\end{proof}

\begin{lemma}[spreading the positive retraction by phases]
\label{lem:phase-spread-collapse}
\Verified\ The lifted retraction $P$ of \cref{lem:positive-collapse} spreads
to a $\Gamma^0$- and $T_c$-equivariant strong deformation retraction of
$Y_{\leq\eta}$ onto $Y_0$.  Conjugating by $G$ gives a $\Gamma$- and
$T_f$-equivariant strong deformation retraction, and hence a strong
deformation retraction
\[
 r_s:N'_0=Y_{\leq\eta}/\Gamma\longrightarrow W=Y_0/\Gamma.
\tag{C.9}
\]
\end{lemma}

\begin{proof}
For $\phi\in T_c$ and $q\in A$ set
\[
 R_s(\phi q)=\phi P_s(q).
\tag{C.10}
\]
This is well defined.  On the open orbit and the outer boundary the
stabilizer is trivial; on $A_0$ the homotopy fixes $q$, so (C.7) causes no
ambiguity.  It is continuous because its pullback along the closed quotient
map $m\times\operatorname{id}_{[0,1]}$ is continuous.

The frozen and positive actions satisfy
\[
 \Psi^0_{\bar\lambda}(\phi q)=
 u_{\bar\lambda}\varphi_{\bar\lambda}(\phi)
 \Psi^+_{\bar\lambda}(q),
\tag{C.11}
\]
where $\varphi_{\bar\lambda}$ is the compact-torus automorphism induced by
$(z,z_3)\mapsto(z+z_3B_0\bar\lambda,z_3)$.  Equation (C.11) and
$\Gamma^+$-equivariance of $P$ prove $\Gamma^0$-equivariance of $R$.
The remaining deformation-retraction identities follow directly from those
of $P$.  Finally \cref{lem:collapse-straightening} conjugates $R$ to the
original action and permits descent to (C.9).
\end{proof}

\begin{proposition}[fixed-fibre cellular collapse]
\label{prop:verified-cellular-collapse}
\Verified\ Fix $t_0$ with $0<|t_0|\leq\eta$ and put $F=f^{-1}(t_0)$.  The
retraction in \cref{lem:phase-spread-collapse} can be chosen so that its end
map $r:N'_0\to W$ restricts to a cellular collapse $c=r|_F$ with the following
description.  There is a splitting
\[
 F\xrightarrow[\bar G]{\sim}F_{t_0}^0\simeq T_B\times T_f,
\tag{C.12}
\]
where $\bar G$ is the identity on the marking $\Lambda$, and the quotient
honeycomb $\Theta\subset T_B$ has two vertices and three edges $e_k$.  Over a
vertex, $c$ collapses all of $T_f$ to a triple point of $W$; over
$e_k^\circ$ it collapses the primitive circle $S^1_{n_k}$ and maps onto the
double curve $\bar C_k$; over $T_B\setminus\Theta$ it is injective onto
$W\setminus D$.
\end{proposition}

\begin{proof}
Let $S\subset A$ be the positive slice $|t|=|t_0|$ and put
\[
 B_\rho^+=B_0+\frac{R(0)}{\log\rho},\qquad
 \ell_\rho=B_0(B_\rho^+)^{-1}.
\]
The map $y:S\to\mathbb R^2$ identifies the positive action with translations
$y\mapsto y+B_{|t_0|}^+\bar\lambda$, while $\ell_{|t_0|}$ conjugates these to
$z\mapsto z+B_0\bar\lambda$.  Define on the frozen fibre
\[
 \widetilde c((\phi,\arg t_0)q)=
 (\phi,\arg t_0)\,
 \vartheta(\ell_{|t_0|}y(q)),\qquad \phi\in T_f, q\in S.
\tag{C.13}
\]
The stabilizer on $S$ is trivial, so the formula is well defined.  Equations
(C.8) and (C.11) make it $\Gamma^0$-equivariant, hence it descends to
$c_\vartheta:F_{t_0}^0\to W$.

Writing
\[
 K_{t_0}=\operatorname{Re}C_0+\frac{\arg t_0}{2\pi}B_0,
\]
the frozen action on $T_f\times\mathbb R_y^2$ is
\[
 \bar\lambda(\phi,y)=
 (\phi e^{2\pi iK_{t_0}\bar\lambda},
  y+B_{|t_0|}^+\bar\lambda).
\]
The change of variables
\[
 \Xi(\phi,y)=
 (\phi e^{-2\pi iK_{t_0}(B_{|t_0|}^+)^{-1}y},y)
\tag{C.14}
\]
conjugates this to the product action and gives (C.12).  Formula (C.7) now
proves the asserted cell description, including primitivity of the collapsed
circle.

It remains to force the end map of an actual strong deformation retraction to
equal (C.13), rather than merely be homotopic to it.  In the chart
$\vartheta:\mathbb R^2\to A_0$ put
\[
 \widehat a=\vartheta^{-1}P_1.
\]
Put $\rho_0=|t_0|$.  If $\rho_0<\eta$, choose a closed positive band
$\{\rho_-\leq |t|\leq\rho_+\}$, disjoint from $A_0$, with $\rho_0$ in its
interior.  If $\rho_0=\eta$, instead use the one-sided band
$\{\rho_-\leq |t|\leq\eta\}$.  Define
$\widehat b=\ell_{|t|}y$ only on the chosen band.  Wherever they are defined,
both maps satisfy
$\widehat a(\Psi^+_{\bar\lambda}q)=\widehat a(q)+B_0\bar\lambda$ and the same
identity for $\widehat b$.  Choose a continuous tent function $\chi(|t|)$,
extended by zero off the band, with $\chi(\rho_0)=1$ and vanishing near every
boundary of the band across which it must be pasted; in the one-sided case no
vanishing is required at the ambient boundary $|t|=\eta$.  With
$w(q)=\chi(|t(q)|)$ define
\[
 K_u(q)=\vartheta\bigl((1-uw(q))\widehat a(q)
                         +uw(q)\widehat b(q)\bigr)
\tag{C.15}
\]
inside the band and $K_u=P_1$ outside.  Because the transition action is by
translations, affine interpolation preserves $\Gamma^+$-equivariance.  The
weight vanishes near every gluing boundary and near $A_0$, so (C.15) pastes
continuously, fixes $A_0$, and takes values in $A_0$.

Concatenate the first half of $P$ with $K_u$.  The result $P'$ still satisfies
all four properties used in \cref{lem:phase-spread-collapse}, and its end map
on $S$ is exactly $\vartheta\ell_{|t_0|}y$.  Spreading $P'$ by phases and
conjugating by $G$ therefore yields a strong deformation retraction whose
restriction to $F$ is
\[
 r|_F=c_\vartheta\circ\bar G.
\]
This proves the proposition, including continuity across all toric strata.
\end{proof}

\subsection{The integral collapse map}

Let $D=\bar C_1\cup\bar C_2\cup\bar C_3$ be the double locus.  Choose
primitive orientations of its three curve classes and of the honeycomb edges
so that the annulus degrees below have the displayed signs.  These orientations
are compared with the clockwise collar convention in
\cref{prop:actual-finite-band-maps,prop:cusp-mv-blocks}.  In
$\Lambda_{\rm tor}=\mathbb Z\langle\widehat w,\widehat\delta\rangle$ put
\[
 n_1=\widehat w,\qquad n_2=\widehat\delta-\widehat w,\qquad
 n_3=-\widehat\delta,\qquad
 \nu_k(v)=\det(n_k,v).
\tag{C.16}
\]

\begin{proposition}[integral homology of the cellular collapse]
\label{prop:verified-collapse-homology}
\Verified\ The map $c_*$ of \cref{prop:verified-cellular-collapse} is
surjective in every degree, and
\[
 0\longrightarrow\operatorname{im}(\textstyle\bigwedge^qM_0-I)
 \longrightarrow\textstyle\bigwedge^q\Lambda
 \xrightarrow{c_*}H_q(W;\mathbb Z)\longrightarrow0
\tag{C.17}
\]
is an exact sequence of free abelian groups.  In the splitting
$\Lambda=\bar\Lambda\oplus\Lambda_{\rm tor}$,
\[
 c_*:\Lambda\to H_1(W)=\bar\Lambda
\]
is projection, and in degree two
\[
 c_*(\widehat\gamma\widehat u)=[\mathbb T_B],\qquad
 c_*(\widehat w\widehat\delta)=0,
\tag{C.18}
\]
\[
 c_*(\beta\otimes v)=
 \sum_{k=1}^3\nu_k(B_0\beta)\nu_k(v)[\bar C_k].
\tag{C.19}
\]
Equivalently, in the ordered source basis
\[
(\widehat\gamma\widehat u,\widehat\gamma\widehat w,
 \widehat\gamma\widehat\delta,\widehat u\widehat w,
 \widehat u\widehat\delta,\widehat w\widehat\delta)
\]
and target basis
$([\mathbb T_B],[\bar C_1],[\bar C_2],[\bar C_3])$,
\[
[c_*]_2=
\begin{pmatrix}
1&0&0&0&0&0\\
0&0&-1&0&0&0\\
0&-1&-1&1&1&0\\
0&0&0&1&0&0
\end{pmatrix}.
\tag{C.20}
\]
In degree three,
\[
 c_*(\widehat\gamma\widehat u\,v)=v
 \quad(v\in\Lambda_{\rm tor}),
 \qquad
 c_*(\beta\,\widehat w\widehat\delta)=0
 \quad(\beta\in\bar\Lambda),
\tag{C.21}
\]
where the first formula fixes the displayed orientation of $H_3(W)$.
\end{proposition}

\begin{proof}
First identify the target integrally.  The normalization
$\nu:\mathrm{dP}_6\to W$ identifies opposite sides of the anticanonical
hexagon.  A plumbing neighbourhood of that hexagon descends to a regular
neighbourhood of $D$.  If $P_0$ is its complement in $\mathrm{dP}_6$, then
$P_0\simeq(\mathbb C^*)^2\simeq T^2$, and excision followed by Lefschetz
duality gives
\[
 H_k(W,D)\simeq H_k(P_0,\partial P_0)
 \simeq H^{4-k}(T^2).
\tag{C.22}
\]
The class $[\mathbb T_B]=c_*[T_B\times\{*\}]$ pairs once with a compact
$T_f$-orbit in $W\setminus D$, so it lifts the generator of $H_2(W,D)$.
Likewise the cycles $c_*(T_B\times S^1_{\widehat w})$ and
$c_*(T_B\times S^1_{\widehat\delta})$ pair unimodularly with the two
coordinate circles in $T_f$.  Thus the boundary maps of the pair vanish and
\[
 H_*(W)=(\mathbb Z,\mathbb Z^2,\mathbb Z^4,\mathbb Z^2,\mathbb Z),
\tag{C.23}
\]
with
$H_2(W)=\mathbb Z\langle[\bar C_1],[\bar C_2],[\bar C_3],
[\mathbb T_B]\rangle$ and $H_3(W)$ generated by the two cycles in (C.21).

The cell description gives projection in degree one and (C.18).  Give $T_B$
the CW structure with one open hexagonal $2$-cell and $1$-skeleton $\Theta$.
A class $\beta\in H_1(T_B)$ has the unique cellular representative
\[
 \sum_{k=1}^3\nu_k(B_0\beta)e_k.
\]
On $e_k\times S^1_v$, the collapse kills $S^1_{n_k}$ and maps the quotient
circle to $\bar C_k$ with degree $\nu_k(v)$.  This proves (C.19), without a
rank inference.  The three choices
$(B_0\beta,v)=(\widehat w,\widehat w),
(\widehat\delta,\widehat\delta),(\widehat w,\widehat\delta)$ give coefficient
vectors
\[
 (0,1,1),\qquad(1,1,0),\qquad(0,1,0),
\]
whose determinant is a unit.  Hence degree two is surjective.  Formula
(C.21) follows directly by crossing the base $2$-cell with a primitive fibre
circle; it proves degree-three surjectivity and records the orientations.
 In degree four the top-dimensional open stratum
$(T_B\setminus\Theta)\times T_f$ maps homeomorphically to the open
$4$-stratum of $W$, so the map is multiplication by a unit after choosing the
target orientation.  Degree zero is immediate.

Let $h$ be geometric monodromy around the cusp.  Fibre transport in $N'_0$
is a homotopy from the inclusion of $F$ to its composite with $h$.  Composing
with $r$ yields $c\simeq c\circ h$, and therefore
\[
 \operatorname{im}(\textstyle\bigwedge^qM_0-I)\subseteq\ker c_*.
\tag{C.24}
\]
The inclusion is an equality integrally, not only rationally.  In the ordered
basis $(\widehat\gamma,\widehat u,\widehat w,\widehat\delta)$, direct exterior
multiplication gives the primitive image lattices
\[
\begin{array}{c|c}
q&\operatorname{im}(\bigwedge^qM_0-I)\\ \hline
0,4&0\\
1&\langle\widehat w,\widehat\delta\rangle\\
2&\langle\widehat w\widehat\delta,
          \widehat\gamma\widehat w+\widehat u\widehat\delta\rangle\\
3&\langle\widehat\gamma\widehat w\widehat\delta,
          \widehat u\widehat w\widehat\delta\rangle.
\end{array}
\tag{C.25}
\]
Each displayed subgroup is a direct summand.  Its rank is
$\binom4q-b_q(W)$ by (C.23), while surjectivity gives the same rank for
$\ker c_*$.  The quotient $H_q(W)$ is free, so $\ker c_*$ is saturated.
Two saturated sublattices of the same rank and rational span coincide, proving
(C.17).
\end{proof}

\subsection{Comparison with analytic nearby cycles}

\begin{theorem}[collapse equals nearby-cycle specialization]
\label{thm:collapse-equals-specialization}
\Verified\ With the radial marking of the punctured cusp family, for every
$0\leq q\leq4$ one has an equality of integral homomorphisms
\[
 \operatorname{sp}^q=i_F^*r^*=c^*:
 H^q(W;\mathbb Z)\longrightarrow H^q(F;\mathbb Z).
\tag{C.26}
\]
Consequently the cellular calculation gives
\[
 \operatorname{im}c^*=(\textstyle\bigwedge^qV)^{T_0}
\]
as a saturated lattice.  In particular, the degree-three orientation in the
cellular computation is the degree-three orientation of intrinsic
specialization; there is no independent sign or index.
\end{theorem}

\begin{proof}
Let $N_\rho=f^{-1}(|t|\leq\rho)\subset N'_0$ and let
$i_{W,\rho}:W\hookrightarrow N_\rho$ be the central-fibre inclusion.  The
closed tubes are cofinal among neighbourhoods of $W$: if an open set
$\mathcal O\subset N'_0$ contains $W$, properness makes
$f(N'_0\setminus\mathcal O)$ closed and disjoint from $0$, so
$N_\rho\subset\mathcal O$ for all sufficiently small $\rho$.  For each such
$\rho$, choose $t_\rho$ on the radial ray through the fixed point $t_0$ whose
fibre is $F$, with $0<|t_\rho|<\rho$, and write
$i_{t_\rho}:F_{t_\rho}\hookrightarrow N_\rho$.  Identify
$H^q(F_{t_\rho};\mathbb Z)$ with $H^q(F;\mathbb Z)$ by parallel transport
along that ray.  The neighborhood realization of the unit
$\mathbb Z_W\to\psi_f\mathbb Z$, together with proper base change, identifies
specialization with
\[
 H^q(W;\mathbb Z)
 \xleftarrow[\sim]{\ \rho_0\ }
 \varinjlim_{\rho\downarrow0}H^q(N_\rho;\mathbb Z)
 \xrightarrow{\ \rho_{\rm near}\ }H^q(F;\mathbb Z),
\tag{C.27}
\]
where $\rho_0$ is induced by $i_{W,\rho}^*$ and $\rho_{\rm near}$ is induced
by the transported maps $i_{t_\rho}^*$.  These restriction maps are compatible
under shrinking $\rho$.  This is the same analytic integral
construction used in
\cref{thm:p0-integral-specialization}; see
\cite[Definition~4.2.1 and Propositions~4.2.2, 4.2.11]{Dimca2004}.

Applying the construction of \cref{lem:phase-spread-collapse} at radius
$\rho$ shows that $i_{W,\rho}^*$ is an isomorphism; no compatibility of those
radius-dependent deformation homotopies is needed.  Let
$j_\rho:N_\rho\hookrightarrow N'_0$ be inclusion.  The single retraction
$r:N'_0\to W$ satisfies
\[
 i_{W,\rho}^*j_\rho^*r^*
 =(r\circ j_\rho\circ i_{W,\rho})^*
 =\operatorname{id}_{H^q(W;\mathbb Z)}.
\]
Hence $j_\rho^*r^*$ is the unique inverse of $i_{W,\rho}^*$.  These inverses
are compatible under shrinking because they are restrictions of the single
map $r^*:H^q(W;\mathbb Z)\to H^q(N'_0;\mathbb Z)$.  Consequently
$\rho_0^{-1}(\alpha)$ in (C.27) is represented, for every sufficiently small
$\rho$, by $j_\rho^*r^*\alpha$.

The restrictions of one class on $N'_0$ to the fibres along the punctured
radial ray form a section of the Gauss--Manin local system.  Therefore, after
the parallel-transport identifications used to define $\rho_{\rm near}$,
$i_{t_\rho}^*j_\rho^*r^*\alpha$ equals $i_F^*r^*\alpha$.  It follows that the
composite in
(C.27) is $i_F^*r^*$.  Since $c=r|_F$, functoriality gives
$i_F^*r^*=(r\circ i_F)^*=c^*$.

Changing the lift of the radial marking composes the target with monodromy.
By (C.24), $c^*$ lands in the monodromy-invariant lattice, so the equality is
independent of that harmless choice.  Finally
\cref{thm:p0-integral-specialization} identifies the common image with the
full saturated invariant lattice.
\end{proof}

\begin{test}[resolved cusp-collapse falsifiers]
\Verified\ The audit tested the failure modes previously attached to C1--C3:
continuity at every toric stratum, finite stabilizer components, equivariance,
primitive edge circles, degree-two and degree-three orientations, torsion and
saturation, and functorial equality with the nearby-cycle unit.  The proofs
above exclude each one.  This status concerns the cusp collapse only; it does
not infer the global Mayer--Vietoris matrices.
\end{test}

\begin{auditrecord}[V1: based van Kampen inclusions]
\label{record:vankampen}
\Verified\ The actual overlap coordinate maps, with the zero section as
base locus and clockwise meridians, show that the
finite fillings impose $\widehat\rho_j^{m_j}=t_{v_j}$ with the same sign,
the cusp filling kills $\Lambda_{\rm tor}$ and its toric meridian, and the
tautological gluing has zero translation drift.  Thus the abstract terminal
presentation is the actual $\pi_1(X)$.  The deck generators, common clockwise
convention, zero-section disc, normal closure, boundary inclusion
$\Lambda\to\Lambda/\Lambda_{\rm tor}$,
$\widetilde a_0\mapsto1$, and absence of extra whisker drift are proved in
\cref{prop:geometric-pi1}.
\end{auditrecord}

\subsection{The actual finite-band Mayer--Vietoris maps}

Fix the ordered orientation
\[
H_1(F)=\Lambda=
\mathbb Z\langle\widehat\gamma,\widehat u,
                    \widehat w,\widehat\delta\rangle
\]
dual to $(\gamma,u,w,\delta)$, use the complex orientations of
$F,S_1,S_2$, and orient every base meridian clockwise.  Put
\[
 \eta_1=2u+w+3\delta,
 \qquad \eta_2=u+w+2\delta,
 \qquad q=uw+6\gamma\delta.
\tag{MV.1}
\]

\begin{proposition}[actual finite-band maps in primitive bases]
\label{prop:actual-finite-band-maps}
\Verified\ In the decomposition $X^\circ=X_1\cup X_2$, the maps denoted
$\alpha_r$ in the manuscript are the actual Mayer--Vietoris maps
\[
 \alpha_r=(\pi_{1*},-\pi_{2*}):
 H_r(F;\mathbb Z)\longrightarrow
 H_r(S_1;\mathbb Z)\oplus H_r(S_2;\mathbb Z).
\tag{MV.2}
\]
There is no conjugacy, unrecorded index, or independent orientation sign in
(MV.2).

Choose primitive bases
\[
 H_1(S_1)=\mathbb Z\langle g_1,h_1\rangle,
 \quad [\widehat\gamma]_1=3g_1,\quad
 h_1=\pi_{1*}\widehat w,
\]
\[
 H_1(S_2)=\mathbb Z\langle g_2,h_2\rangle,
 \quad [\widehat\gamma]_2=-4g_2,\quad
 h_2=\pi_{2*}\widehat w.
\]
Then, in the source basis
$(\widehat\gamma,\widehat u,\widehat w,\widehat\delta)$ and the target
basis $(g_1,h_1,g_2,h_2)$,
\[
 [\alpha_1]=
 \begin{pmatrix}
 3&0&0&0\\
 0&2&1&3\\
 4&0&0&0\\
 0&-1&-1&-2
 \end{pmatrix}.
\tag{MV.3}
\]
It has Smith form $\operatorname{diag}(1,1,1,0)$ and
\[
 \ker\alpha_1=\mathbb Z\nu_0,
 \qquad
 \nu_0=-\widehat u-\widehat w+\widehat\delta.
\tag{MV.4}
\]

Define cohomology bases $(a_j,b_j)$ by
\[
 \pi_1^*a_1=\gamma\eta_1,
 \qquad \pi_1^*b_1=q-2\gamma\eta_1,
\]
\[
 \pi_2^*a_2=2\gamma\eta_2,
 \qquad \pi_2^*b_2=q-3\gamma\eta_2,
\tag{MV.5}
\]
and let $(a_j^\sharp,b_j^\sharp)$ be the Kronecker-dual homology bases.
In the ordered source basis
\[
(\widehat\gamma\widehat u,
 \widehat\gamma\widehat w,
 \widehat\gamma\widehat\delta,
 \widehat u\widehat w,
 \widehat u\widehat\delta,
 \widehat w\widehat\delta)
\]
and target basis
$(a_1^\sharp,b_1^\sharp,a_2^\sharp,b_2^\sharp)$,
\[
 [\alpha_2]=
 \begin{pmatrix}
  2& 1& 3& 0&0&0\\
 -4&-2& 0& 1&0&0\\
 -2&-2&-4& 0&0&0\\
  3& 3& 0&-1&0&0
 \end{pmatrix}.
\tag{MV.6}
\]
This matrix also has Smith form $\operatorname{diag}(1,1,1,0)$; its
primitive left kernel is
\[
 \Phi=(4,2,3,2).
\tag{MV.7}
\]
The class $\mathcal A\in H^2(X^\circ;\mathbb Z)$ represented by $\Phi$
restricts to $2q$, because
\[
 2q=4\pi_1^*a_1+2\pi_1^*b_1
    =3\pi_2^*a_2+2\pi_2^*b_2.
\tag{MV.8}
\]

If the degree-three bases are Poincar\'e dual to the degree-one cohomology
bases just fixed, the projection formula makes $[\alpha_3]$ the same matrix
as (MV.3).  Finally
\[
 [\alpha_0]=\binom{1}{-1},
 \qquad
 [\alpha_4]=\binom{3}{-4}.
\tag{MV.9}
\]
Every image in (MV.3), (MV.6), and (MV.9) is saturated.
\end{proposition}

\begin{proof}
Choose closed discs $D'_1,D'_2$ whose intersection is a contractible band
disjoint from the three special points.  Over the annular shell between
$D'_j$ and the smaller filling disc, $f$ is a proper submersion.  An
Ehresmann connection lifts a base deformation retraction that is stationary
near the smaller disc; adjoining the identity there gives a deformation
retraction $X_j\to N_j$.  The verified local radial retraction of
\cref{lem:finite-filling-stalk-map} keeps the marked flat coordinate $x$
fixed and sends its endpoint to the orbit $\pi_j(x)\in S_j$.  Since the band
is contractible, choose its marking by the same radial paths.  The composite
of the actual band inclusion with $X_j\to N_j\to S_j$ is therefore literally
$\pi_j$.  The standard chain-level Mayer--Vietoris injection is
$z\mapsto(i_{1*}z,-i_{2*}z)$
\cite[\S2.2]{Hatcher2002}, proving (MV.2) with its sign.

The affine fundamental-group presentations of the finite fibres and the
exact finite pullback lattices of \cref{prop:finite-indices} give the stated
primitive degree-one bases and (MV.3).  Its $3\times3$ minors include $3$
and $4$, so their gcd is one; direct elimination gives (MV.4).
The two bases in (MV.5) are unimodular changes of
\[
 \langle\gamma\eta_1,q\rangle,
 \qquad
 \langle2\gamma\eta_2,\gamma\eta_2+q\rangle,
\]
the exact integral images of $\pi_j^*$ in degree two.  Evaluating them on the
six source tori gives (MV.6).  Its rank-three minors include $2$ and $-3$,
so the image is saturated; multiplication verifies (MV.7) and (MV.8).

For degree three, if $x\in H_3(F)$ and $\theta\in H^1(S_j)$, then
\[
 \langle\theta,\pi_{j*}x\rangle
 =\langle\pi_j^*\theta,x\rangle.
\]
Thus in the Poincar\'e-dual bases the degree-three matrix equals the
degree-one matrix.  In degree four, the holomorphic quotient coverings have
positive degrees $3$ and $4$; the sole minus sign is the Mayer--Vietoris
minus sign.  This proves (MV.9) and completes the integral certificate.
\end{proof}

\subsection{The cusp attachment and its primitive blocks}

\begin{proposition}[actual cusp kernels and global blocks]
\label{prop:cusp-mv-blocks}
\Verified\ For the second decomposition $X=N'_0\cup X^\circ$, the map
$H_r(M)\to H_r(W)$ is surjective and
\[
 K_1=\mathbb Z[\widetilde a_0],
 \qquad
 K_2=\mathbb Z\langle
       \xi_{\widehat w},\xi_{\widehat\delta}\rangle,
\tag{MV.10}
\]
where
\[
 \xi_\nu=a_*([S^1_\nu]\times[\widetilde a_0])
\]
uses the oriented fibrewise translation action and the clockwise toric
meridian.  Both displayed subgroups are primitive.

Let $\zeta$ be the generator of $H_1(X^\circ)$ normalized by
$\Phi_1(\zeta)=1$.  In the source basis
$(\widehat\gamma,\widehat u,\widetilde a_0)$ of $H_1(M)$ and target basis
$(\bar\gamma,\bar u,\zeta)$ of
$H_1(W)\oplus H_1(X^\circ)$, the actual map
$(c_*,-j_*)$ is
\[
 D_1=
 \begin{pmatrix}
 1&0&0\\
 0&1&0\\
 -12&0&-p
 \end{pmatrix},
 \qquad \det D_1=-p.
\tag{MV.11}
\]

Orient the connecting map by
\[
 \partial(j_*\xi_\nu)=(I-A_2)\nu=(I-A_1^{-1})\nu
 \quad(\nu\in\Lambda_{\rm tor}).
\tag{MV.12}
\]
Let $e$ be the primitive generator of
$C_2:=\operatorname{coker}\alpha_2$ with
$\langle\mathcal A,e\rangle=1$, and put
$z=j_*\xi_{\widehat w}$.  Then $(e,z)$ is a basis of
$H_2(X^\circ)$ and
\[
 j_*\xi_{\widehat\delta}=-p\,e,
 \qquad
 [j_*|_{K_2}]_{(\xi_{\widehat\delta},\xi_{\widehat w})}^{(e,z)}
 =\begin{pmatrix}-p&0\\0&1\end{pmatrix}.
\tag{MV.13}
\]
After splitting the surjections onto $H_r(W)$, the full degree-one and
degree-two Mayer--Vietoris maps are integrally equivalent to
\[
 I_2\oplus[p],
 \qquad
 I_4\oplus
 \begin{pmatrix}-p&0\\0&1\end{pmatrix},
\tag{MV.14}
\]
respectively.  They are unimodular for the candidate $p=-1$.
\end{proposition}

\begin{proof}
The composite $H_r(F)\to H_r(M)\to H_r(W)$ is the verified collapse $c_*$,
which is surjective by \cref{prop:verified-collapse-homology}; hence the
second arrow is surjective.  Under the Wang projection, the base class
$[\widetilde a_0]$ is primitive.  It bounds the explicit zero-section disc
in the toric filling, proving the first formula in (MV.10).

For $K_2$, use the toric chart $t=z_1z_2z_3$.  The literal solid tori
\[
 \{|z_1|\leq\rho,\ |z_2|=r,\ z_3=c\}
 \quad\hbox{and}\quad
 \{|z_2|\leq\rho,\ |z_1|=r,\ z_3=c\}
\]
have oriented boundaries $\xi_{\widehat w}$ and
$\xi_{\widehat\delta}$.  Under
\[
 H_2(M)\twoheadrightarrow
 \ker(M_0-I|\Lambda)=\Lambda_{\rm tor},
\]
these classes map to $\widehat w$ and $\widehat\delta$.  Their span is
therefore primitive; the Wang ranks and the verified homology of $W$ show
that it is all of $K_2$.

The actual meridian calculation gives
\[
 j_*\widehat\gamma=12\zeta,
 \qquad j_*\widehat u=0,
 \qquad j_*[\widetilde a_0]=p\zeta,
\]
while $c_*$ is projection onto
$\mathbb Z\langle\bar\gamma,\bar u\rangle$.  This is (MV.11).
Because $A_1A_2=M_0^{-1}$ acts trivially on $\Lambda_{\rm tor}$, the two
expressions in (MV.12) agree.  Since
$A_2\widehat w=-\widehat u+\widehat\delta$ and
$A_2\widehat\delta=\widehat\delta$,
\[
 \partial(j_*\xi_{\widehat w})=-\nu_0,
 \qquad \partial(j_*\xi_{\widehat\delta})=0.
\tag{MV.15}
\]
Thus $z$ is a primitive lift, up to the fixed unit sign, of the generator of
$\ker\alpha_1$.

It remains to determine the coefficient of $e$.  Normalize contraction by
\[
 \iota_v(\theta_1\wedge\cdots\wedge\theta_k)
 =\sum_r(-1)^{r-1}\theta_r(v)
   \theta_1\wedge\cdots\widehat{\theta_r}\cdots\wedge\theta_k.
\]
The globally free $S^1_{\widehat\delta}$-action and the sweep/slant identity
give
\[
 \langle\mathcal A,j_*\xi_{\widehat\delta}\rangle
 =\langle\iota_{\widehat\delta}\mathcal A,
          j_*[\widetilde a_0]\rangle,
 \qquad
 \iota_{\widehat\delta}(2q)=-12\gamma.
\tag{MV.16}
\]
Let $\Theta\in H^1(X^\circ)$ be the primitive class restricting to
$12\gamma$.  Restriction to $F$ is injective, so (MV.16) implies
$\iota_{\widehat\delta}\mathcal A=-\Theta$.  The based meridian relation
gives
\[
 \Theta(j_*[\widetilde a_0])=p.
\]
Hence the left side of (MV.16) is $-p$, which proves (MV.13).  Splitting the
surjective cusp blocks gives (MV.14).
\end{proof}

\begin{remark}[one-shot $H_1$ check]
The preceding two-stage calculation is the manuscript's argument.  There is
also a useful one-shot check.  Let $U$ be the torus bundle over the
pair-of-pants core and
$N=N'_0\sqcup N'_1\sqcup N'_2$.  In the target basis
\[
(c,\rho_1,\rho_2;
 \bar\gamma,\bar u;
 g_1,h_1;g_2,h_2)
\]
and domain basis
\[
(\gamma_0,u_0,a_0;
 \gamma_1,w_1,\rho_1;
 \gamma_2,w_2,\rho_2),
\]
the actual map $(i_{U*},-i_{N*})$ for the candidate is
\[
\Delta_1=
\begin{pmatrix}
1&0&\ell_0&1&0&0&1&0&0\\
0&0&-1&0&0&1&0&0&0\\
0&0&-1&0&0&0&0&0&1\\
-1&0&0&0&0&0&0&0&0\\
0&-1&0&0&0&0&0&0&0\\
0&0&0&-3&0&-1&0&0&0\\
0&0&0&0&-1&0&0&0&0\\
0&0&0&0&0&0&4&0&-1\\
0&0&0&0&0&0&0&-1&0
\end{pmatrix}.
\tag{MV.17}
\]
For example, $a_0$ maps to
$\ell_0c-\rho_1-\rho_2$ on the pair-of-pants side and to zero on the cusp
side; $\gamma_1$ maps to $c$ and $3g_1$, while $\gamma_2$ maps to $c$ and
$-4g_2$.  Thus every column is an actual primitive collar cycle.  Exact
elimination gives
\[
 \det\Delta_1=1-12\ell_0=-p
\]
for $(\ell_1,\ell_2)=(1,-1)$.  At $\ell_0=0$, this determinant is one.
\end{remark}

\subsection{Exact comparison with the Leray transgressions}

\begin{theorem}[Mayer--Vietoris/Leray comparison]
\label{thm:mv-leray-comparison}
\Verified\ With the preceding primitive bases and clockwise meridians, the
actual Mayer--Vietoris obstruction maps commute with the intrinsic Leray
transgressions in degrees one, two, and three.  If
$\omega_{\rm MV}\in H^2(B;\mathbb Z)$ is the generator fixed below, then
\[
 d_2^{0,1}(12\gamma)=-p\,\omega_{\rm MV},
\tag{L.1}
\]
\[
 d_2^{0,2}(2q)=-p\,\omega_{\rm MV}[\delta],
 \qquad
 d_2^{0,3}(2\gamma uw)
   =-p\,\omega_{\rm MV}[\gamma\delta].
\tag{L.2}
\]
For the candidate $p=-1$, all three maps are multiplication by $+1$ in
these bases.
\end{theorem}

\begin{proof}
Let $V$ be the pair-of-pants neighbourhood in the base and let
$D_i$ be the three special discs.  Order the annuli as
$(D_1\cap V,D_2\cap V,D_0\cap V)$ and orient their meridians clockwise.
Define
\[
 \omega_{\rm MV}=\partial_B[(1,0,0)].
\]
The chain-level Mayer--Vietoris sequence on the base gives
\[
 \partial_B(n_1,n_2,n_0)
 =(n_1+n_2+n_0)\omega_{\rm MV},
\tag{L.3}
\]
because the image of $H^1(V)$ is the sum-zero sublattice.

Choose a lift $\Theta_V$ of $12\gamma$ over the pair of pants and write
$b_i=\Theta_V(\widehat\rho_i)$.  The finite local lifts have meridian values
$4\ell_1$ and $3\ell_2$, while the cusp lift vanishes on
$\widetilde a_0$.  For the general affine cusp regluing used in this
comparison, let $t_\mu$ denote translation by the class
$\mu\in\Lambda/\Lambda_{\rm tor}$ and set
$\ell_0=\gamma(\mu)$.  Thus
$\Theta_V(t_\mu)=12\ell_0$.  The tautological gluing of the constructed
candidate has $\mu=0$ and $\ell_0=0$; in that case the following relation
has no vertical drift.  In general,
\[
 \widetilde a_0=t_\mu\widehat\rho_2^{-1}
                         \widehat\rho_1^{-1},
\]
the local-minus-pair-of-pants discrepancy is
\[
 \bigl(4\ell_1-b_1,\ 3\ell_2-b_2,\
 -(12\ell_0-b_1-b_2)\bigr).
\tag{L.4}
\]
Its signed sum is $-p$.  The canonical truncation triangle identifies this
very cocycle, functorially on the induced total-space cover, with Leray
$d_2$; this is the construction proved in
\cref{lem:first-transgression}.  Equations (L.3)--(L.4) prove (L.1).

This is more than equality of the terminal integer.  The transpose of the
oriented collar map applied to the tuple of local lifts has meridian values
\[
 (b_1-4\ell_1,\ b_2-3\ell_2,\
 12\ell_0-b_1-b_2),
\]
the negative of (L.4).  Consequently the square
\[
\begin{CD}
\mathbb Z\langle12\gamma\rangle
 @>{d_2^{0,1}}>>
\mathbb Z\langle\omega_{\rm MV}\rangle\\
@V{\cong}VV @VV{\cong}V\\
H^1(X^\circ;\mathbb Z)
@>{(j_*|_{K_1})^\vee}>>
K_1^\vee
\end{CD}
\tag{L.5}
\]
commutes when the left vertical map sends $12\gamma$ to $\Theta$ and the
right sends $\omega_{\rm MV}$ to $-\widetilde a_0^\vee$.  The lower map is
multiplication by $p$, while the upper map is multiplication by $-p$.

For degree two, quotient the primitive lift in (MV.15):
\[
 \overline K_2=K_2/\mathbb Z\xi_{\widehat w},
 \qquad C_2=\operatorname{coker}\alpha_2.
\]
Equation (MV.13) gives an actual reduced map
\[
 \overline j:\overline K_2\longrightarrow C_2,
 \qquad [\xi_{\widehat\delta}]\longmapsto -p\,e.
\tag{L.6}
\]
Under the primitive identifications
\[
 2q\longmapsto\Phi\in C_2^\vee,
 \qquad
 \omega_{\rm MV}[\delta]\longmapsto
 \xi_{\widehat\delta}^{\,\vee}\in\overline K_2^\vee,
\]
the square
\[
\begin{CD}
\mathbb Z\langle2q\rangle
@>{d_2^{0,2}}>>
\mathbb Z\langle\omega_{\rm MV}[\delta]\rangle\\
@V{\cong}VV @VV{\cong}V\\
C_2^\vee @>{\overline j^{\,\vee}}>> \overline K_2^\vee
\end{CD}
\tag{L.7}
\]
commutes, and both horizontal maps are multiplication by $-p$.  This proves
the first formula in (L.2).

The final formula has no new sign choice.  Write
\[
 d_2(2q)=r\,\omega_{\rm MV}[\delta],
 \qquad
 d_2(2\gamma uw)=s\,\omega_{\rm MV}[\gamma\delta].
\]
Using
\[
 (12\gamma)(2q)=12(2\gamma uw),
 \qquad (2q)(2\gamma uw)=0,
\]
the derivation rule and the already fixed degree-one sign give
\[
 s=2(-p)-r.
\]
The identities
\[
 \delta\wedge2\gamma uw=-2\operatorname{vol},
 \qquad
 2q\wedge\gamma\delta=2\operatorname{vol}
\]
applied to the second product give $r=s$.  Hence $r=s=-p$, proving the
second formula in (L.2).

Changing every meridian orientation reverses $p$ and
$\omega_{\rm MV}$ simultaneously.  Reversing the
$S^1_{\widehat\delta}$ orientation reverses both
$\xi_{\widehat\delta}$ and its dual coefficient generator.  Thus all
remaining convention changes are common units and cannot change an index.
\end{proof}

\begin{test}[resolved global topological falsifiers]
\Verified\ The actual finite-band maps are (MV.2), their degree-one and
degree-two images are saturated, the multiple-fibre top class is
$(3,-4)$, the cusp kernels are literal primitive collar cycles, the reduced
cusp blocks are (MV.11) and (MV.13), and the comparison is the pair of
commuting squares (L.5) and (L.7), not a comparison of absolute determinants.
The former MV1 and L1 failure modes are therefore excluded for the candidate.
\end{test}

\subsection{Complete recognition theorem}

\begin{lemma}[cohomological closure in dimension six]
\label{lem:cohomological-closure}
\Verified\ Let $X$ be a smooth, closed, connected, oriented six-manifold.
If
\[
\pi_1(X)=0,
\qquad H^2(X;\mathbb Z)=H^3(X;\mathbb Z)=0,
\tag{4.4}
\]
then $H_k(X;\mathbb Z)$ is $\mathbb Z$ for $k=0,6$ and zero for
$1\le k\le5$.
\end{lemma}

\begin{proof}
The Hurewicz map in degree one identifies
$H_1(X;\mathbb Z)=\pi_1(X)^{\mathrm{ab}}=0$.  The universal coefficient
sequence
\[
0\longrightarrow\operatorname{Ext}(H_1(X),\mathbb Z)
\longrightarrow H^2(X;\mathbb Z)
\longrightarrow\operatorname{Hom}(H_2(X),\mathbb Z)
\longrightarrow0
\]
shows that $\operatorname{Hom}(H_2(X),\mathbb Z)=0$, so the finitely
generated group $H_2(X)$ is torsion.  The next universal coefficient
sequence contains an injection
\[
\operatorname{Ext}(H_2(X),\mathbb Z)
\hookrightarrow H^3(X;\mathbb Z)=0.
\]
For a finitely generated torsion group
$A=\bigoplus_i\mathbb Z/n_i$, one has
$\operatorname{Ext}(A,\mathbb Z)=\bigoplus_i\mathbb Z/n_i$; hence
$H_2(X)=0$.  Poincar\'e duality gives
\[
H_3(X)\simeq H^3(X)=0,\quad
H_4(X)\simeq H^2(X)=0,\quad
H_5(X)\simeq H^1(X)=0.
\]
Finally connectedness and orientation give $H_0(X)=H_6(X)=\mathbb Z$.
\end{proof}

\begin{theorem}[sphere recognition from the exact inputs]
\label{thm:sphere-recognition}
\Verified\ Let $X$ be a smooth, closed, connected six-manifold.  If
\[
\pi_1(X)=0,\qquad
H_k(X;\mathbb Z)\cong
\begin{cases}
\mathbb Z,&k=0,6,\\
0,&1\le k\le5,
\end{cases}
\tag{4.5}
\]
then $X$ is diffeomorphic to the standard $S^6$.
\end{theorem}

\begin{proof}
A closed smooth manifold has the homotopy type of a finite CW complex.
We first prove successively that $\pi_i(X)=0$ for $1\le i\le5$.
The case $i=1$ is assumed.  Suppose $2\le n\le5$ and
$\pi_i(X)=0$ for $i<n$.  Then $X$ is $(n-1)$-connected, so the Hurewicz
homomorphism
\[
h_n:\pi_n(X)\longrightarrow H_n(X;\mathbb Z)
\]
is an isomorphism by the Hurewicz theorem
\cite[Theorem~4.32]{Hatcher2002}.  The target vanishes by (4.5), hence
$\pi_n(X)=0$.  Induction proves the assertion.

Now $X$ is $5$-connected.  Hurewicz in degree six gives
$\pi_6(X)\cong H_6(X;\mathbb Z)\cong\mathbb Z$.  Choose a map
$g:S^6\to X$ whose homotopy class maps to a chosen orientation generator.
It has degree one.  Hence $g_*$ is an isomorphism in degrees zero and six;
in every other degree both integral homology groups vanish.  Thus $g$ is an
integral homology equivalence.  Both spaces are simply connected CW
complexes, so the homological Whitehead theorem
\cite[Corollary~4.33]{Hatcher2002} makes $g$ a homotopy equivalence.
Therefore $X$ is a smooth homotopy $6$-sphere.

Smale's generalized Poincar\'e theorem in dimensions greater than four
\cite[Theorem~A]{Smale1961} identifies the underlying topological manifold
with $S^6$.  The smooth endpoint is stated directly by Smale:
every homotopy $6$-sphere is diffeomorphic to $S^6$
\cite[Corollary~1.2]{Smale1962}.  Equivalently, smooth homotopy
$6$-spheres up to oriented $h$-cobordism form $\Theta_6$;
Kervaire--Milnor's opening table gives $|\Theta_6|=1$, and their following
remark identifies $h$-cobordism with diffeomorphism in dimension six
\cite[opening table and following Remark, pp.~504--505]{KervaireMilnor1963}.
Hence there is no exotic smooth structure in dimension six, and $X$ is
diffeomorphic to the standard sphere.
\end{proof}

\begin{corollary}[abstract application theorem]\label{cor:abstract-s6}
\Verified\ Let $X$ be a smooth, closed, connected, oriented six-manifold.
Suppose V1 identifies the actual based van Kampen quotient
with $G(\ell_0,\ell_1,\ell_2)$ of
\cref{prop:terminal-group}, and suppose an independently constructed
MV1/L1 comparison proves
\[
H^2(X;\mathbb Z)\cong H^3(X;\mathbb Z)
\cong\mathbb Z/|p|\mathbb Z,
\qquad p=12\ell_0-4\ell_1-3\ell_2.
\tag{4.6}
\]
If $|p|=1$, then $X\cong_{\mathrm{diff}}S^6$.
\end{corollary}

\begin{proof}
For $|p|=1$, 
\cref{prop:terminal-group} and the assumed geometric identification give
$\pi_1(X)=0$, while (4.6) gives $H^2(X)=H^3(X)=0$.
\Cref{lem:cohomological-closure} converts these inputs into the homology
condition (4.5), and
\cref{thm:sphere-recognition} gives the claimed diffeomorphism.
\end{proof}

\begin{corollary}[independent Mayer--Vietoris certificate for the candidate]
\label{cor:candidate-mv-sphere}
\Verified\ For the tautological candidate
$(\ell_0,\ell_1,\ell_2)=(0,1,-1)$, the two-stage geometric
Mayer--Vietoris calculation gives
\[
 H_k(X;\mathbb Z)\cong
 \begin{cases}
 \mathbb Z,&k=0,6,\\
 0,&1\leq k\leq5.
 \end{cases}
\]
Together with the independently verified based van Kampen calculation, this
proves $X\cong_{\mathrm{diff}}S^6$ without using the Leray computation of the
same groups.
\end{corollary}

\begin{proof}
Here $p=-1$.  Thus (MV.11) and the reduced block (MV.13) are unimodular.
The two Mayer--Vietoris sequences give $H_1(X)=H_2(X)=0$ integrally, with no
torsion hidden by a rank calculation.  The verified collapse homology
(C.23) gives $e(W)=1-2+4-2+1=2$.  Away from $W$, the map to the punctured
base is stratified by smooth four-tori and the two bielliptic fibres, all of
Euler characteristic zero.  Additivity therefore gives $e(X)=e(W)=2$.

Poincar\'e duality and universal coefficients give $H_5(X)=H_4(X)=0$.
Moreover
\[
 H_3(X)\cong H^3(X)\cong\operatorname{Hom}(H_3(X),\mathbb Z),
\]
because $H_2(X)=0$; hence $H_3(X)$ is free.  The equality $e(X)=2$ forces
its rank to be zero.  The based van Kampen result
\cref{prop:geometric-pi1} gives $\pi_1(X)=1$, and
\cref{thm:sphere-recognition} completes the proof.
\end{proof}

Thus C1--C3, V1, MV1, and L1 are all verified for the candidate.  The
nearby-cycle/Leray route and the collapse/Mayer--Vietoris route are now two
separate integral certificates for the same sphere-recognition input.

\section{Fibre degree and the cyclic-six centre}
\label{zdc:sec:bridge}

This section supplies the theorem--proof module behind the fibre-degree and
cyclic-centre formulas recorded in the orientation material.  There are three
different objects in the calculation:
\[
 H^4(U;\mathbb Z),\qquad H^4(U;\mathbb Z/6),\qquad
 Z(SU(6))=\mu _6.
\]
The first is a free rank-six characteristic lattice, the second has six
independent residue coordinates, and the third is a single cyclic group.
The candidate determines one particular surjection from the first object to
the third.  It does not identify all three objects, and it does not construct
a gauge connection or a spectral operator.

\subsection{Spaces, markings, and the finite-band presentation}
\label{zdc:subsec:markings}

Retain the reconstructed map
\[
 f:X\longrightarrow\mathbb P^1,
 \qquad X\cong_{\mathrm{diff}}S^6,
\]
and write
\[
 D_0=W,\qquad D_1=S_1,\qquad D_2=S_2,
 \qquad D=D_0\sqcup D_1\sqcup D_2,
 \qquad U=X\setminus D.
\]
Here $W$ is the reduced non-normal fibre and $3S_1,4S_2$ are the
two multiple fibres.  The verified local homology computations give
\[
 H_1(W;\mathbb Z)=\mathbb Z\langle a_0,b_0\rangle,
 \qquad
 H_1(S_j;\mathbb Z)=\mathbb Z\langle a_j,b_j\rangle
 \quad(j=1,2).
\]
The marking is
\[
 a_0=[\widehat\gamma],\quad b_0=[\widehat u],
 \quad a_j=[\widehat g_j],
\]
where $\widehat g_j$ is the clockwise affine deck generator.  The second
generator $b_j$ is the invariant coinvariant coordinate.  With the
marking used in the finite-band calculation one may take
\(b_1=\pi_{1*}[\widehat w]\) and
\(b_2=\pi_{2*}[\widehat w]\); in the second coinvariant quotient
\(\pi_{2*}[\widehat u]=\pi_{2*}[\widehat w]\), so this agrees with the
alternative representative used elsewhere.  The two specializations of
the marked $\widehat\gamma$-cycle are
\begin{equation}
 \pi_{1*}[\widehat\gamma]=3a_1,
 \qquad
 \pi_{2*}[\widehat\gamma]=-4a_2.
 \label{zdc:eq:gamma-specializations}
\end{equation}

Fix Alexander duality using the complex orientation of $X$, outward-first
boundary orientation, and the convention that the Poincare dual of a
positive nearby fibre is the negative covector of its clockwise meridian.
It identifies
\begin{equation}
 L:=H^4(U;\mathbb Z)
 \xrightarrow{\ \sim\ }
 H_1(D;\mathbb Z)
 =\mathbb Z\langle a_0,b_0;a_1,b_1;a_2,b_2\rangle
  \cong\mathbb Z^6.
 \label{zdc:eq:defect-lattice}
\end{equation}
Indeed, $D$ is compact and locally contractible, so
\[
 \widetilde H^q(S^6\setminus D;\mathbb Z)
 \cong\widetilde H_{5-q}(D;\mathbb Z).
\]
Take $q=4$ and use the disjoint-union decomposition of $D$.

Let $X^\circ$ be the space obtained after the two finite fillings and
before the central toric filling.  For logarithmic twist integers
\(\ell_1,\ell_2\), its verified degree-one Mayer--Vietoris calculation is
the presentation
\begin{equation}
 H_1(X^\circ;\mathbb Z)
 =\frac{\mathbb Z\langle c,\bar g_1,\bar g_2\rangle}
 {\mathbb Z\langle3\bar g_1-\ell_1c,
                         4\bar g_2-\ell_2c\rangle}.
 \label{zdc:eq:h1xcirc-presentation}
\end{equation}
Thus $c,\bar g_1,\bar g_2$ are three marked generators, not a basis.
The relation matrix is
\[
 R_1=
 \begin{pmatrix}-\ell_1&3&0\\-\ell_2&0&4\end{pmatrix}.
\]
For admissible twists $3\nmid\ell_1$ and $2\nmid\ell_2$, the gcd of
its rank-two minors
\[
 3\ell_2,\qquad-4\ell_1,\qquad12
\]
is one.  Therefore the presented group is torsion-free of rank one.  The
homomorphism
\begin{equation}
 \begin{aligned}
 \Phi_1:H_1(X^\circ;\mathbb Z)&\longrightarrow\mathbb Z,\\
 \Phi_1(c)&=12,&
 \Phi_1(\bar g_1)&=4\ell_1,&
 \Phi_1(\bar g_2)&=3\ell_2
 \end{aligned}
 \label{zdc:eq:phi1}
\end{equation}
annihilates both relations.  Its three displayed values have gcd one under
the same admissibility hypotheses; hence it is an isomorphism.  For the
candidate
\begin{equation}
 (\ell_0,\ell_1,\ell_2)=(0,1,-1),
 \qquad
 \Phi_1(c)=12,\quad
 \Phi_1(\bar g_1)=4,\quad
 \Phi_1(\bar g_2)=-3.
 \label{zdc:eq:candidate-twists}
\end{equation}
Put
\begin{equation}
 \zeta:=\bar g_1+\bar g_2.
 \label{zdc:eq:zeta-generator}
\end{equation}
Then $\Phi_1(\zeta)=1$.  Since $\Phi_1$ is an isomorphism, $\zeta$ is the
positive basis element of $H_1(X^\circ;\mathbb Z)$ selected by $\Phi_1$, and
\begin{equation}
 c=12\zeta,\qquad \bar g_1=4\zeta,\qquad \bar g_2=-3\zeta.
 \label{zdc:eq:presentation-to-basis}
\end{equation}

Let
\[
 \overline q_{\mathrm{col}}:H_1(D;\mathbb Z)
 \longrightarrow H_1(X^\circ;\mathbb Z)
\]
be the homomorphism sending $a_0,a_1,a_2$ to
$c,\bar g_1,\bar g_2$, respectively, and sending every $b_i$ to zero.
Define $q_{\mathrm{col}}=\overline q_{\mathrm{col}}\circ\operatorname{AD}$,
where $\operatorname{AD}$ is the fixed Alexander isomorphism above.  Thus
the marked collar homomorphism is well typed:
\begin{equation}
 \begin{aligned}
 q_{\mathrm{col}}:L&\longrightarrow H_1(X^\circ;\mathbb Z),\\
 q_{\mathrm{col}}(a_0)&=c,&q_{\mathrm{col}}(b_0)&=0,\\
 q_{\mathrm{col}}(a_1)&=\bar g_1,&q_{\mathrm{col}}(b_1)&=0,\\
 q_{\mathrm{col}}(a_2)&=\bar g_2,&q_{\mathrm{col}}(b_2)&=0.
 \end{aligned}
 \label{zdc:eq:qcol}
\end{equation}
The central-band projection sends $a_0=[\widehat\gamma]$ to $c$ and
$b_0=[\widehat u]$ to zero.  In the two finite-band cokernels, the actual
Mayer--Vietoris map is $(\pi_{1*},-\pi_{2*})$; it sends the affine generators
to $\bar g_1,\bar g_2$ and kills the invariant coordinates $b_1,b_2$.
These statements derive the six displayed values from the marked collar
calculation.  The map is surjective because its images contain the three
presentation generators in \eqref{zdc:eq:h1xcirc-presentation}.  It is not
an inclusion-induced map $H_1(D)\to H_1(X^\circ)$: no inclusion of the
deleted central fibre $W$ into $X^\circ$ exists.

\subsection{The exact fibre-degree row}
\label{zdc:subsec:fibre-degree}

Let $F\cong T^4$ be a complex-oriented smooth fibre and choose
\(\omega_F\in H^4(F;\mathbb Z)\) with
\(\langle\omega_F,[F]\rangle=1\).  If $i_F:F\hookrightarrow U$ is the
inclusion, put $\operatorname{res}_F=i_F^*:H^4(U;\mathbb Z)\to
H^4(F;\mathbb Z)$.  Define
\begin{equation}
 \operatorname{res}_F(\alpha)=\deg_F(\alpha)\omega_F,
 \qquad
 \deg_F:L\longrightarrow\mathbb Z.
 \label{zdc:eq:degree-definition}
\end{equation}

\begin{theorem}[Fibre-degree row]
\label{zdc:thm:fibre-degree}
With the fixed Alexander-duality convention above, one has
\begin{equation}
 \boxed{
 \deg_F(n_0a_0+r_0b_0+n_1a_1+r_1b_1+n_2a_2+r_2b_2)
 =12n_0+4n_1-3n_2.}
 \label{zdc:eq:degree-row}
\end{equation}
Equivalently,
\[
 [\deg_F]_{(a_0,b_0,a_1,b_1,a_2,b_2)}
 =(12,0,4,0,-3,0).
\]
Moreover,
\begin{equation}
 \boxed{\deg_F=\Phi_1\circ q_{\mathrm{col}}}
 \label{zdc:eq:degree-factorization}
\end{equation}
for the candidate twists, and
\begin{equation}
 \boxed{
 \ker\deg_F
 =\mathbb Z\langle
 b_0,b_1,b_2,a_0-3a_1,3a_1+4a_2\rangle.}
 \label{zdc:eq:degree-kernel}
\end{equation}
This kernel is saturated of rank five.
\end{theorem}

\begin{proof}
We retain every sign in the one-shot collar calculation.  In the ordered
boundary-source basis
\[
 (\gamma_0,u_0,\widetilde a_0;
   \gamma_1,w_1,\rho_1;
   \gamma_2,w_2,\rho_2)
\]
and ordered target basis
\[
 (c,\bar\rho_1,\bar\rho_2;
   a_0,b_0;a_1,b_1;a_2,b_2),
\]
the actual map $(i_{U*},-i_{N*})$ at
\(\left(\ell_0,\ell_1,\ell_2\right)=(0,1,-1)\) is
\begin{equation}
 \Delta_1=
 \begin{pmatrix}
 1&0&0&1&0&0&1&0&0\\
 0&0&-1&0&0&1&0&0&0\\
 0&0&-1&0&0&0&0&0&1\\
 -1&0&0&0&0&0&0&0&0\\
 0&-1&0&0&0&0&0&0&0\\
 0&0&0&-3&0&-1&0&0&0\\
 0&0&0&0&-1&0&0&0&0\\
 0&0&0&0&0&0&4&0&-1\\
 0&0&0&0&0&0&0&-1&0
 \end{pmatrix}.
 \label{zdc:eq:collar-matrix}
\end{equation}
For example, the columns say
\[
 \gamma_0\mapsto c-a_0,\quad u_0\mapsto-b_0,\quad
 \widetilde a_0\mapsto-\bar\rho_1-\bar\rho_2,
\]
\[
 \gamma_1\mapsto c-3a_1,\quad w_1\mapsto-b_1,\quad
 \rho_1\mapsto\bar\rho_1-a_1,
\]
\[
 \gamma_2\mapsto c+4a_2,\quad w_2\mapsto-b_2,\quad
 \rho_2\mapsto\bar\rho_2-a_2.
\]
Thus the entries $-3,+4$ are exactly
\eqref{zdc:eq:gamma-specializations}; all other displayed minus signs are
the second-inclusion signs in the Mayer--Vietoris convention.  Direct
integral elimination gives $\det\Delta_1=1$.

Let $\widetilde a_0^\vee$ be the integral covector on the ordered
boundary-source group that is one on $\widetilde a_0$ and zero on the other
eight basis elements.  Pullback of target covectors is represented by
$\Delta_1^{\mathsf T}$.  By the fixed orientation convention, the target
covector extending the Poincare dual of a complex-oriented nearby fibre
must pull back to $-\widetilde a_0^\vee$.  Since $\Delta_1$ is unimodular,
there is a unique integral extension.  Its last six entries are precisely
the Alexander coordinates of the restriction functional on $H_1(D)$.
Consequently they are obtained from
\begin{equation}
 \Delta_1^{\mathsf T}y=-\widetilde a_0^\vee.
 \label{zdc:eq:transpose-system}
\end{equation}
Writing
\[
 y=(A,B,C;X_0,Y_0;X_1,Y_1;X_2,Y_2),
\]
the nine equations are
\[
 A=X_0,\quad Y_0=Y_1=Y_2=0,\quad
 A=3X_1,\quad A=-4X_2,
\]
\[
 B=X_1,\qquad C=X_2,\qquad -B-C=-1.
\]
Their unique solution is
\[
 y=(12,4,-3;12,0;4,0;-3,0).
\]
The last six entries prove \eqref{zdc:eq:degree-row}.  Equations
\eqref{zdc:eq:phi1}, \eqref{zdc:eq:candidate-twists}, and
\eqref{zdc:eq:qcol} give the same row on the six basis elements, proving
the typed equality \eqref{zdc:eq:degree-factorization}.

The row is primitive because $\gcd(12,4,3)=1$; explicitly,
\(\deg_F(a_1+a_2)=1\).  Each of the five displayed elements in
\eqref{zdc:eq:degree-kernel} lies in the kernel.  Conversely, if
\(12n_0+4n_1-3n_2=0\), reduction modulo four gives $n_2=4v$, and then
$n_1=3v-3n_0$.  Therefore
\[
 n_0a_0+n_1a_1+n_2a_2
 =n_0(a_0-3a_1)+v(3a_1+4a_2).
\]
The three $b_i$-coordinates are free, and the five displayed generators
are independent.  They therefore form a basis of the kernel.  Finally, a
kernel of a surjection from a free abelian group to $\mathbb Z$ is
saturated, proving the last assertion.
\end{proof}

Reversing the global Alexander-duality convention negates all six entries
of \eqref{zdc:eq:degree-row} simultaneously.  It cannot change the sign of
only one component.  In particular, with the convention fixed here,
$q_{\mathrm{col}}(a_2)=+\bar g_2$ while the coordinate of $\bar g_2$ in the
positive generator of $H_1(X^\circ;\mathbb Z)$ is $-3$.

For completeness, if $j:M\hookrightarrow X^\circ$ denotes the central
collar inclusion, the toric attaching meridian obeys
\[
 j_*[\widetilde a_0]=\ell_0c-\bar g_1-\bar g_2,
\]
and hence
\[
 \Phi_1(j_*[\widetilde a_0])
 =12\ell_0-4\ell_1-3\ell_2=p.
\]
At the candidate values this is $p=-1$.  Thus the same correctly typed
functional controls the smooth-fibre degree and the terminal meridian, but
the two statements use the maps displayed above rather than an untyped
identification of their source groups.

\subsection{Normal-line holonomy and the two finite reductions}
\label{zdc:subsec:finite-holonomy}

The transverse disc in the order-$m_j$ local quotient is acted on by
\(e^{-2\pi i/m_j}\), with $m_1=3,m_2=4$.  Therefore the normal-line
holonomy characters, in the marked bases, are
\begin{equation}
 \begin{aligned}
 h_1:H_1(S_1;\mathbb Z)&\longrightarrow\mu_3,&
 h_1(a_1)&=e^{-2\pi i/3},& h_1(b_1)&=1,\\
 h_2:H_1(S_2;\mathbb Z)&\longrightarrow\mu_4,&
 h_2(a_2)&=e^{-2\pi i/4},& h_2(b_2)&=1.
 \end{aligned}
 \label{zdc:eq:normal-holonomies}
\end{equation}
Via \eqref{zdc:eq:defect-lattice}, write each $x\in L$ uniquely as
$x=(x_0,x_1,x_2)$ with $x_i\in H_1(D_i;\mathbb Z)$.  Regard
$\mu_3$ and $\mu_4$ as their standard subgroups of $\mu_{12}$.  The inverse
on the second factor below records the candidate twist $\ell_2=-1$, rather
than a discretionary change of sign.  Define
\begin{equation}
 \begin{aligned}
 \chi_{12}:L&\longrightarrow\mu_{12},&
 \chi_{12}(x_0,x_1,x_2)&=h_1(x_1)h_2(x_2)^{-1},\\
 \chi_6:L&\longrightarrow\mu_6,&
 \chi_6(x)&=\chi_{12}(x)^2.
 \end{aligned}
 \label{zdc:eq:finite-characters}
\end{equation}

\begin{proposition}[The order-twelve and order-six reductions]
\label{zdc:prop:finite-reductions}
Put
\[
 \rho_m:\mathbb Z\longrightarrow\mu_m,
 \qquad \rho_m(k)=e^{-2\pi ik/m}.
\]
Then
\begin{equation}
 \boxed{\chi_{12}=\rho_{12}\circ\deg_F,
 \qquad \chi_6=\rho_6\circ\deg_F.}
 \label{zdc:eq:holonomy-degree}
\end{equation}
Both characters are surjective, and
\begin{equation}
 L/\ker\chi_{12}\cong\mu_{12},
 \qquad L/\ker\chi_6\cong\mu_6.
 \label{zdc:eq:finite-quotients}
\end{equation}
\end{proposition}

\begin{proof}
For
\(x=\sum_{i=0}^2(n_i a_i+r_i b_i)\),
\eqref{zdc:eq:normal-holonomies} gives, with
\(\zeta_{12}=e^{-2\pi i/12}\),
\[
 \chi_{12}(x)=\zeta_{12}^{\,4n_1-3n_2}.
\]
Equation \eqref{zdc:eq:degree-row} differs in its exponent only by
\(12n_0\), which vanishes modulo twelve.  This proves the first identity.
Squaring it is reduction of the same integer modulo six and proves the
second.  Since $\deg_F(a_1+a_2)=1$, both reductions are surjective.  The
first isomorphism theorem gives \eqref{zdc:eq:finite-quotients}.
\end{proof}

The kernels can be read without suppressing any coordinate:
\[
 \ker\chi_{12}=
 \mathbb Za_0\oplus\mathbb Zb_0\oplus
 3\mathbb Za_1\oplus\mathbb Zb_1\oplus
 4\mathbb Za_2\oplus\mathbb Zb_2,
\]
\[
 \ker\chi_6=
 \mathbb Za_0\oplus\mathbb Zb_0\oplus
 3\mathbb Za_1\oplus\mathbb Zb_1\oplus
 2\mathbb Za_2\oplus\mathbb Zb_2.
\]
The first follows by reducing $4n_1-3n_2$ modulo three and four; the
second follows by reducing it modulo three and two.  In particular, the
natural combined character has order twelve before its order-six quotient.

\subsection{Coefficient reduction is not the distinguished quotient}
\label{zdc:subsec:coefficients}

Alexander duality also gives
\[
 H^5(U;\mathbb Z)\cong\widetilde H_0(D;\mathbb Z)\cong\mathbb Z^2.
\]
The cohomological universal-coefficient sequence with coefficients
\(\mathbb Z/6\) is therefore
\[
 0\longrightarrow H^4(U;\mathbb Z)\otimes\mathbb Z/6
 \longrightarrow H^4(U;\mathbb Z/6)
 \longrightarrow
 \operatorname{Tor}_1^{\mathbb Z}(H^5(U;\mathbb Z),\mathbb Z/6)
 \longrightarrow0,
\]
whose last term is zero.  Thus the marked coefficient group is
\begin{equation}
 \boxed{H^4(U;\mathbb Z/6)\cong L/6L\cong(\mathbb Z/6)^6.}
 \label{zdc:eq:coefficient-group}
\end{equation}
If $r_6:L\to L/6L$ is coefficient reduction, then $\chi_6$ factors as
\[
 L\xrightarrow{r_6}(\mathbb Z/6)^6
 \xrightarrow{\ \overline d_F\ }\mathbb Z/6
 \xrightarrow{k\mapsto e^{-2\pi ik/6}}\mu_6,
\]
where, in the six marked residue coordinates,
\begin{equation}
 [\overline d_F]=(0,0,4,0,3,0)\pmod6.
 \label{zdc:eq:distinguished-mod-six-row}
\end{equation}
Thus
\[
 \overline d_F(\bar n_0,\bar r_0,\bar n_1,
                    \bar r_1,\bar n_2,\bar r_2)
 =4\bar n_1+3\bar n_2\in\mathbb Z/6.
\]
It is surjective because $\overline d_F(0,0,1,0,1,0)=1$ in
$\mathbb Z/6$.
Consequently
\[
 L/\ker\chi_6
 \cong (L/6L)/\ker\overline d_F
 \cong\mathbb Z/6.
\]
This is one distinguished cyclic quotient of the six-coordinate group in
\eqref{zdc:eq:coefficient-group}; it is not that group itself.

\subsection{The tensor subgroup and the genuine common-centre quotient}
\label{zdc:subsec:groups}

\begin{theorem}[Tensor subgroup and centres]
\label{zdc:thm:tensor-centres}
Tensor product defines a faithful closed homomorphism
\begin{equation}
 \iota:SU(2)\times SU(3)\hookrightarrow SU(6),
 \qquad \iota(A,B)=A\otimes B.
 \label{zdc:eq:tensor-map}
\end{equation}
Its restriction to the centres is the isomorphism
\begin{equation}
 \mu_2\times\mu_3\xrightarrow{\ \sim\ }\mu_6=Z(SU(6)),
 \qquad (\lambda,\mu)\longmapsto\lambda\mu.
 \label{zdc:eq:centre-isomorphism}
\end{equation}
The image of \(\iota\) is a proper subgroup of $SU(6)$, of real
dimension $3+8=11$.
\end{theorem}

\begin{proof}
The matrix $A\otimes B$ is unitary and
\[
 \det(A\otimes B)=(\det A)^3(\det B)^2=1,
\]
so \(\iota\) has the stated codomain.  Suppose more generally that
$A\otimes B=sI_6$.  Choose an eigenvector $w$ of $B$, with eigenvalue
$\beta$.  Then
\[
 (A\otimes B)(v\otimes w)=\beta Av\otimes w=sv\otimes w
\]
for every $v$, so $A=(s/\beta)I_2$; substitution shows that $B$ is scalar
as well.  Applying this scalar--tensor lemma with $s=1$ gives
\[
 A=\lambda I_2,\qquad B=\mu I_3,\qquad \lambda\mu=1.
\]
The determinant conditions give $\lambda^2=1$ and $\mu^3=1$.  Hence
\(\lambda=\mu^{-1}\in\mu_2\cap\mu_3=\{1\}\), proving injectivity.  A
continuous injective homomorphism from a compact group into a Hausdorff
group has closed image.

On centres the tensor matrix is $(\lambda\mu)I_6$.  Multiplication is
injective because $\mu_2\cap\mu_3=\{1\}$, and both its domain and target
have six elements, proving \eqref{zdc:eq:centre-isomorphism}.  Explicitly,
the inverse sends $s\in\mu_6$ to $(s^3,s^{-2})$.  Finally the dimension of
the image is $11<35=\dim SU(6)$, so it is proper.
\end{proof}

The isomorphisms proved above give the explicit chain
\begin{equation}
 L/\ker\chi_6\xrightarrow{\ \overline\chi_6\ }\mu_6
 \xrightarrow{\lambda\mapsto\lambda I_6}Z(SU(6)).
 \label{zdc:eq:lattice-to-centre}
\end{equation}
This is an isomorphism of finite groups.  It does not provide transition
functions for an $SU(6)$-bundle.  Also, \(\iota\) does not factor through
the quotient of $SU(2)\times SU(3)$ by its full centre: that centre maps
isomorphically, rather than trivially, to $Z(SU(6))$.

Put $\zeta_6=e^{2\pi i/6}$.  In the integer $U(1)$-charge
normalization, the common-centre generator used for the Standard-Model
global form is
\[
 (\zeta_6,-I_2,e^{2\pi i/3}I_3).
\]
If $\omega=e^{2\pi i/3}=\zeta_6^2$, then
\[
 \iota(-I_2,\omega I_3)=(-\omega)I_6=\zeta_6^{-1}I_6;
\]
the first coordinate $\zeta_6$ cancels this scalar under the
tensor--scalar map below.
This agrees with the convention in David Tong, \emph{Line Operators in the
Standard Model}, JHEP 07 (2017) 104, arXiv:1705.01853v2.  The following
matrix theorem proves the quotient directly and does not rely on a physical
interpretation of that reference.

\begin{theorem}[Common-centre quotient with $U(1)$]
\label{zdc:thm:common-centre}
Define
\begin{equation}
 \delta:\mu_6\hookrightarrow U(1)\times SU(2)\times SU(3),
 \qquad
 \delta(\xi)=(\xi,\xi^3I_2,\xi^2I_3),
 \label{zdc:eq:delta}
\end{equation}
and
\begin{equation}
 \widehat\iota:U(1)\times SU(2)\times SU(3)\longrightarrow U(6),
 \qquad
 \widehat\iota(z,A,B)=z(A\otimes B).
 \label{zdc:eq:extended-map}
\end{equation}
Then
\begin{equation}
 \ker\widehat\iota=\delta(\mu_6),
 \label{zdc:eq:extended-kernel}
\end{equation}
so the induced map is a faithful closed embedding
\begin{equation}
 \frac{U(1)\times SU(2)\times SU(3)}{\delta(\mu_6)}
 \hookrightarrow U(6).
 \label{zdc:eq:g6-embedding}
\end{equation}
Moreover,
\begin{equation}
 \boxed{\det\widehat\iota(z,A,B)=z^6,}
 \qquad
 \boxed{\operatorname{im}\widehat\iota\cap SU(6)
        =\iota(SU(2)\times SU(3)).}
 \label{zdc:eq:det-intersection}
\end{equation}
\end{theorem}

\begin{proof}
For $\xi\in\mu_6$, one has
\[
 \det(\xi^3I_2)=\xi^6=1,
 \qquad
 \det(\xi^2I_3)=\xi^6=1.
\]
Thus $\delta$ has the stated codomain.  It is central, is a homomorphism,
and is injective because its first coordinate is $\xi$.

For $\xi\in\mu_6$,
\[
 \widehat\iota(\delta(\xi))
 =\xi(\xi^3I_2\otimes\xi^2I_3)
 =\xi^6I_6=I_6,
\]
so \(\delta(\mu_6)\) lies in the kernel.  Conversely, suppose
\(z(A\otimes B)=I_6\).  The eigenvalue argument in the preceding proof,
now applied to $A\otimes B=z^{-1}I_6$, gives
\[
 A=\lambda I_2,\qquad B=\mu I_3,\qquad
 \lambda^2=1,\quad\mu^3=1,\quad z\lambda\mu=1.
\]
The Chinese-remainder isomorphism
\[
 \mu_6\longrightarrow\mu_2\times\mu_3,
 \qquad \xi\longmapsto(\xi^3,\xi^2)
\]
provides a unique $\xi\in\mu_6$ with
\(\lambda=\xi^3\) and $\mu=\xi^2$.  Then
\(z=(\lambda\mu)^{-1}=\xi\), because $\xi^{-5}=\xi$.  Hence the kernel
is exactly \(\delta(\mu_6)\).  Compactness makes the induced faithful image
closed.

The determinant formula is
\[
 \det\bigl(z(A\otimes B)\bigr)
 =z^6(\det A)^3(\det B)^2=z^6.
\]
Every tensor matrix lies in the intersection by taking $z=1$.  Conversely,
if $z(A\otimes B)$ lies in $SU(6)$, then $z^6=1$.  By
\eqref{zdc:eq:centre-isomorphism}, choose
\((\lambda,\mu)\in\mu_2\times\mu_3\) with $\lambda\mu=z$.  Then
\[
 z(A\otimes B)=(\lambda A)\otimes(\mu B),
\]
and $\lambda A\in SU(2)$, $\mu B\in SU(3)$.  This proves the reverse
inclusion in \eqref{zdc:eq:det-intersection}.
\end{proof}

Writing the quotient in \eqref{zdc:eq:g6-embedding} as $G_6$, determinant
therefore gives the exact sequence
\[
 1\longrightarrow SU(2)\times SU(3)
 \xrightarrow{\ \iota\ }G_6
 \xrightarrow{\ \det\circ\widehat\iota\ }U(1)
 \longrightarrow1.
\]
Here the middle arrow means the tensor embedding followed by the quotient
identification with $\operatorname{im}\widehat\iota$; surjectivity on the
right follows because $z\mapsto z^6$ is surjective on $U(1)$, and its kernel
is \eqref{zdc:eq:det-intersection}.

The quotient in \eqref{zdc:eq:g6-embedding} is the genuine common-centre
quotient: the $U(1)$ factor cancels the nontrivial scalar produced by the
two special-unitary centres.  The displayed tensor--scalar image in $U(6)$
is not contained in $SU(6)$; by \eqref{zdc:eq:det-intersection}, its
intersection with $SU(6)$ is exactly the tensor subgroup.  The quotient is
not isomorphic to $SU(6)$: their real dimensions are $12$ and $35$.

\subsection{Exact boundary of the result}
\label{zdc:subsec:boundary}

The construction proves the primitive fibre-degree map, its saturated
kernel, its order-twelve and order-six finite reductions, the distinction
between the full coefficient group and the selected cyclic quotient, and
the two matrix-group embeddings.  It does not produce a degree-changing
morphism from
\[
 H^4(U;\mathbb Z)\longrightarrow H^2(M;\mu_6)
\]
for a four-dimensional spacetime $M$.  In particular,
\(\chi_6\) is not, merely by its codomain, the obstruction to lifting a
common-centre quotient bundle.  No physical gauge bundle, connection,
Yang--Mills Hamiltonian, confinement statement, or mass-gap conclusion is
inferred here.

\section{The non-normal fibre and the CDP obstruction}

Campana--Demailly--Peternell prove that a compact complex threefold with
$H^1(X;\mathbb Z)=H^2(X;\mathbb Z)=0$, algebraic dimension one, and
holomorphic algebraic reduction has $c_3(X)\le0$
\cite[Theorem~2.2, p.~680]{CDP2020}.  Their corollary gives $a(X)=0$ for
a compact complex threefold homeomorphic to $S^6$
\cite[Corollary~2.3, p.~680]{CDP2020}.  The candidate asserts all those
hypotheses together with $c_3(X)=2$ and $a(X)=1$.  This section separates:
\begin{enumerate}[label=(\roman*)]
\item validity of one printed inference on p.~692 of \cite{CDP2020};
\item validity of the published theorem by any other argument;
\item validity of the candidate's proposed non-normal-fibre mechanism; and
\item truth of the candidate's upstream construction and topology.
\end{enumerate}

\subsection{The exact p.~692 inference}

Let $S$ be the reduction of a fibre in a smooth threefold $\mathcal X$,
and let $A$ denote the line-bundle twist after duality.  The two relevant
sequences are
\[
0\longrightarrow N^*_{S/\mathcal X}\otimes A
\longrightarrow\Omega^1_{\mathcal X}|_S\otimes A
\xrightarrow{\rho}\Omega^1_S\otimes A\longrightarrow0
\tag{5.1}
\]
and
\[
0\longrightarrow\mathcal T(\Omega^1_S)\otimes A
\longrightarrow\Omega^1_S\otimes A
\longrightarrow\widetilde\Omega^1_S\otimes A\longrightarrow0,
\tag{5.2}
\]
where $\mathcal T$ denotes sections supported on a nowhere dense analytic
set.  The opening reduction of \cite[\S7, p.~692]{CDP2020} treats the
vanishing of the global sections of the left end of (5.1) and the right end
of (5.2) as sufficient for the vanishing of the middle term of (5.1).
That inference leaves untreated a section whose image under $\rho$ is a
nonzero torsion section.  The omitted sheaf can be identified exactly.  Put
\[
 \mathscr T_S:=\ker\bigl(\Omega_S^1\to
 j_*\Omega^1_{S_{\rm reg}}\bigr),
 \qquad \widetilde\Omega_S^1=\Omega_S^1/\mathscr T_S,
\]
and
\[
 \mathscr K_S:=
 \ker\bigl(\Omega_X^1|_S\longrightarrow\widetilde\Omega_S^1\bigr).
\]

\begin{proposition}[the missing cotangent filtration]
\label{prop:missing-cotangent-filtration}
Let $S$ be a reduced effective Cartier divisor in a smooth complex
threefold $X$.  There are natural short exact sequences
\[
0\longrightarrow N^*_{S/X}\longrightarrow\mathscr K_S
\longrightarrow\mathscr T_S\longrightarrow0,
\]
\[
0\longrightarrow\mathscr K_S\longrightarrow\Omega_X^1|_S
\longrightarrow\widetilde\Omega_S^1\longrightarrow0.
\]
\end{proposition}

\begin{proof}
The conormal sequence is right exact.  It is also left exact here.  Locally
write $S=(g=0)$ with $g$ reduced.  If $h\,dg|_S=0$, then $h$ vanishes at the
generic smooth point of every irreducible component, because $dg$ is nonzero
there.  The ring $\mathcal O_S$ is reduced, so the intersection of its
minimal primes is zero; hence $h=0$.  Thus
\[
0\to N^*_{S/X}\to\Omega_X^1|_S\to\Omega_S^1\to0
\]
is exact.  By definition, $\mathscr K_S$ is the inverse image of
$\mathscr T_S$ under the last surjection.  Restriction to that inverse image
is onto $\mathscr T_S$ and has kernel $N^*_{S/X}$, proving the first
sequence.  Composing with the quotient by $\mathscr T_S$ proves the second.
\end{proof}

\begin{proposition}[the exact global-section criterion]
\label{prop:missing-global-section-criterion}
Let $A$ be a line bundle on $S$.  If
\[
H^0(S,\widetilde\Omega_S^1\otimes A)=0,
\]
then
\[
H^0(S,\Omega_X^1|_S\otimes A)
=H^0(S,\mathscr K_S\otimes A).
\]
If also $H^0(S,N^*_{S/X}\otimes A)=0$, then this group is canonically
\[
\ker\!\left[
H^0(S,\mathscr T_S\otimes A)
\xrightarrow{\delta_A}H^1(S,N^*_{S/X}\otimes A)
\right].
\]
Hence the two vanishings printed in the p.~692 reduction imply the desired
ambient vanishing if and only if $\delta_A$ is injective.
\end{proposition}

\begin{proof}
Tensoring the second sequence of
\cref{prop:missing-cotangent-filtration} by the locally free sheaf $A$
preserves exactness; its global-section sequence gives the first equality.
Tensor the first sequence and take global sections.  If the conormal
$H^0$ vanishes, exactness identifies the middle group with the displayed
kernel.  No condition on the torsion-free quotient controls that kernel.
\end{proof}

\subsection{Local algebra of a reduced normal-crossings fibre}

Let $g:\mathcal X\to\Delta$ be a proper holomorphic map from a smooth complex
threefold to a disc, and put $S=g^{-1}(0)$.  Assume $S$ is reduced and has
normal crossings: at every point there are coordinates $(x,y,z)$ in which
\[
g=ux,\qquad g=uxy,\qquad\text{or}\qquad g=uxyz
\tag{5.3}
\]
for a unit $u$.  Let $\eta:\widetilde S\to S$ be the normalization,
$D=S_{\mathrm{sing,red}}$, and $\widetilde D=\eta^{-1}(D)_{\mathrm{red}}$.

\begin{lemma}[the conormal map is injective]\label{lem:conormal-injective}
\Verified\ In the setting above, $N^*_{S/\mathcal X}\simeq\mathcal O_S$,
and its generator maps to $dg|_S$.  Multiplication by $dg|_S$ is injective,
so
\[
0\longrightarrow\mathcal O_S
\xrightarrow{\,h\mapsto h\,dg\,}
\Omega^1_{\mathcal X}|_S\longrightarrow\Omega^1_S\longrightarrow0
\tag{5.4}
\]
is exact.
\end{lemma}

\begin{proof}
The fibre ideal is the principal ideal $(g)$ and is reduced, so
$(g)/(g^2)\simeq\mathcal O_S$.  Multiplication by a unit does not affect
injectivity, and the term $f\,du$ vanishes after restriction to $f=0$.
It therefore suffices to use $f=x,xy,xyz$.

For $f=xy$, if a lift $\widetilde h$ satisfies
$\widetilde h(y\,dx+x\,dy)=0$ modulo $(xy)$, then
$\widetilde h y,\widetilde h x\in(xy)$.  Hence
$\widetilde h\in(x)\cap(y)=(xy)$ in the regular local ring, so $h=0$ on
$S$.  For $f=xyz$, the three coefficient conditions say
$\widetilde h yz,\widetilde h xz,\widetilde h xy\in(xyz)$; hence
$\widetilde h\in(x)\cap(y)\cap(z)=(xyz)$.  The case $f=x$ is immediate.
Thus the conormal map is injective in every local model, proving (5.4).
\end{proof}

\begin{lemma}[normalization--conductor exact sequence]
\label{lem:conductor-square}
\Verified\ The sequence
\[
0\longrightarrow\mathcal O_S
\xrightarrow{\alpha}
\eta_*\mathcal O_{\widetilde S}\oplus\mathcal O_D
\xrightarrow{\beta}
\eta_*\mathcal O_{\widetilde D}
\longrightarrow0,
\tag{5.5}
\]
\[
\alpha(h)=(\eta^*h,h|_D),\qquad
\beta(\widetilde h,k)=
\widetilde h|_{\widetilde D}-\eta^*k,
\]
is exact.
\end{lemma}

\begin{proof}
The assertion is local.  Away from $D$, $\eta$ is an isomorphism.
At a double point $S=\{xy=0\}$, a member of $\ker\beta$ is
\[
\bigl((h_1(y,z),h_2(x,z)),k(z)\bigr),
\qquad
h_1(0,z)=h_2(0,z)=k(z).
\]
The function
\[
h=h_1(y,z)+h_2(x,z)-h_1(0,z)
\]
on $\{xy=0\}$ restricts to $h_1,h_2$ on the two branches and to $k$ on
their intersection.  This proves middle exactness in the double model.

At a triple point $S=\{xyz=0\}$, write the branch functions as
$h_1(y,z)$, $h_2(x,z)$, $h_3(x,y)$.  Membership in $\ker\beta$ is exactly
\[
h_1(0,z)=h_2(0,z),\quad
h_1(y,0)=h_3(0,y),\quad
h_2(x,0)=h_3(x,0).
\]
Then
\[
\begin{split}
h={}&h_1(y,z)+h_2(x,z)+h_3(x,y)\\
&-h_1(0,z)-h_1(y,0)-h_2(x,0)+h_1(0,0)
\end{split}
\]
restricts to the prescribed function on every branch.  Injectivity of
$\alpha$ follows because $S$ is reduced and $\eta$ is an isomorphism over a
dense open set.  Surjectivity of $\beta$ follows locally by lifting a
function on the branch conductor to the smooth branch and taking $k=0$.
This proves exactness in every local model.
\end{proof}

\begin{theorem}[conductor-supported ambient differential]
\label{prop:conductor}
\Verified\ In the setting (5.3), let $A$ be a line bundle on $S$ with
$\eta^*A\simeq\mathcal O_{\widetilde S}$, and choose a nowhere-vanishing
section $e$ of $\eta^*A$.  There is a unique nonzero
\[
s\in H^0(S,\Omega^1_{\mathcal X}|_S\otimes A)
\tag{5.6}
\]
such that
\[
\eta^*s=\eta^*(dg|_S)\otimes e,\qquad s|_D=0.
\tag{5.7}
\]
Its image $\rho(s)$ in $\Omega^1_S\otimes A$ is supported on $D$ and
therefore torsion; its image in
$\widetilde\Omega^1_S\otimes A$ is zero.  Moreover:
\begin{enumerate}[label=(\alph*)]
\item $\rho(s)=0$ if and only if $e=\eta^*\sigma$ for a nowhere-vanishing
$\sigma\in H^0(S,A)$;
\item if $S$ is connected and $A\not\simeq\mathcal O_S$, then
$H^0(S,A)=0$ and $s$ does not come from the conormal term in (5.1).
\end{enumerate}
\end{theorem}

\begin{proof}
In the double model, $d(xy)=y\,dx+x\,dy$ has coefficients in the reduced
ideal $(x,y)$ of $D$.  In the triple model,
\[
d(xyz)=yz\,dx+xz\,dy+xy\,dz
\]
has coefficients in $(xy,xz,yz)$, the reduced ideal of the union of the
three double curves.  Unit factors in (5.3) do not change this after
restriction to $S$.  Hence $\eta^*(dg|_S)$ restricts to zero on
$\widetilde D$.  Tensor (5.5) by the locally free sheaf
$E=\Omega^1_{\mathcal X}|_S\otimes A$.  The compatible pair
\[
\bigl(\eta^*(dg|_S)\otimes e,0\bigr)
\]
therefore descends uniquely to $s$.  On the smooth locus of $S$, $dg|_S$
is a nonzero conormal vector and $e$ is nowhere zero, so $s\ne0$.

On $S\setminus D$, the normalization is an isomorphism and
$s=dg|_S\otimes e$.  By (5.4) its intrinsic image vanishes there.  Thus
$\rho(s)$ is supported on $D$ and lies in the torsion subsheaf; its image
modulo torsion is zero.

If $\rho(s)=0$, exactness of (5.4) after tensoring by $A$ gives
$s=dg|_S\otimes\sigma$ for a section $\sigma$ of $A$.  Pulling back and
comparing with (5.7) on the dense open set where $dg\ne0$ gives
$\eta^*\sigma=e$ there, hence everywhere by holomorphic continuation.
Since $e$ never vanishes and $\eta$ is surjective, $\sigma$ never vanishes
and trivializes $A$.  The converse is immediate.

Finally suppose $S$ is connected and $\sigma\in H^0(S,A)$.  On each compact
connected component of $\widetilde S$, write
$\eta^*\sigma=c_i e$; the holomorphic function $c_i$ is constant.  If one
$c_i$ is zero and another is nonzero, connectedness of the finite union of
irreducible components of $S$ provides two such components meeting at a
point.  The value of $\sigma$ at that point would be both zero and nonzero,
a contradiction.  Hence either every $c_i=0$, so $\sigma=0$, or every
$c_i\ne0$, so $\sigma$ is nowhere zero and $A$ is trivial.  When $A$ is
nontrivial only the first case occurs.  Since
$N^*_{S/\mathcal X}\simeq\mathcal O_S$, the conormal global-section space is
$H^0(S,A)=0$, while $s\ne0$.
\end{proof}

\subsection{A proper counterexample to the printed inference}

\begin{theorem}[nodal-cubic test family]\label{prop:nodal-counterexample}
\Refuted\ The p.~692 implication described after (5.2) is false as a formal
sheaf-theoretic statement, even for a proper family with smooth total space
and an irreducible reduced normal-crossings central fibre.
\end{theorem}

\begin{proof}
Let $\Delta$ be a sufficiently small disc with coordinate $t$ and put
\[
\mathcal C=
\{([x:y:z],t)\in\mathbb P^2\times\Delta:
y^2z=x^2(x+z)+tz^3\}.
\tag{5.8}
\]
This is proper over $\Delta$.  Its total space is smooth.  On $z\ne0$, the
partial derivative of the defining equation with respect to $t$ is
$-z^3\ne0$.  On $z=0$, the equation forces $x=0$; then $y\ne0$ and the
$z$-derivative is $y^2\ne0$.  The central fibre $C_0$ is the irreducible
nodal cubic $y^2z=x^2(x+z)$, with node $[0:0:1]$.

Set
\[
\mathcal X=\mathcal C\times\mathbb P^1,\qquad
S=C_0\times\mathbb P^1,
\]
and let $g:\mathcal X\to\Delta$ be projection.  Then $\mathcal X$ is a
smooth threefold, $g$ is proper, and $S$ is an irreducible reduced
normal-crossings surface.  Its normalization is
$\widetilde S\simeq\mathbb P^1\times\mathbb P^1$.

Let $p,q$ be the two inverse images of the node in the normalization of
$C_0$.  Fix $\lambda\in\mathbb C^*$ which is not a root of unity (for
example, $\lambda=2$).  Glue the fibres of
the trivial line bundle on the normalization at $p$ and $q$ by the scalar
$\lambda$.  This defines a line bundle
$A_\lambda\in\operatorname{Pic}^0(C_0)$ whose pullback is trivial.  It has
infinite order: its $n$th tensor power has gluing scalar $\lambda^n$, which
is never one.  It is not trivial: a global section pulls back to a constant
$c$ satisfying
$c=\lambda c$, so
\[
H^0(C_0,A_\lambda)=0.
\tag{5.9}
\]
Put $A=\operatorname{pr}_1^*A_\lambda$ on $S$.  Then
$\eta^*A\simeq\mathcal O_{\widetilde S}$ and $H^0(S,A)=0$.
Because $S$ is the full fibre,
$N^*_{S/\mathcal X}\simeq\mathcal O_S$, so the conormal-twist global
sections vanish.

The natural map
$\widetilde\Omega^1_S\to\eta_*\Omega^1_{\widetilde S}$ is injective:
its kernel is supported where normalization is not an isomorphism, while
the source is torsion-free by definition.  Projection formula and
$\eta^*A\simeq\mathcal O_{\widetilde S}$ therefore give
\[
H^0(S,\widetilde\Omega^1_S\otimes A)
\hookrightarrow
H^0(\widetilde S,\Omega^1_{\widetilde S})=0,
\tag{5.10}
\]
the last equality following from
$\widetilde S=\mathbb P^1\times\mathbb P^1$.
Nevertheless \cref{prop:conductor} produces
\[
0\ne s\in H^0(S,\Omega^1_{\mathcal X}|_S\otimes A).
\]
Thus both vanishings invoked in the printed reduction hold while the desired
ambient vanishing fails.  This test chooses the line bundle $A$ directly on
$S$; it does not assert that $A$ extends to a line bundle on the total space
$\mathcal X$.  It therefore refutes the claimed \emph{formal implication
between the three sheaf-cohomology groups} for a non-torsion $A$, but does
not by itself instantiate the narrower global-line-bundle situation on
p.~692.  The family (5.8) also does not satisfy all hypotheses of CDP
Theorem~2.2, so this is not a counterexample to that theorem.
\end{proof}

\subsection{Cartier-divisor duality without a normality hypothesis}

\begin{theorem}[complete Cartier-fibre duality]
\label{thm:cartier-duality}
\Verified\ Let $X$ be a compact smooth complex threefold,
$i:W\hookrightarrow X$ an effective Cartier divisor, and $E$ a holomorphic
vector bundle.  Put $F=E|_W$ and
\[
\omega_W=(\omega_X\otimes\mathcal O_X(W))|_W.
\]
For $q=0,1,2$ there are natural perfect pairings
\[
H^q(W,F)^\vee
\simeq H^{2-q}(W,F^\vee\otimes\omega_W).
\tag{5.11}
\]
Moreover $H^q(W,F)=0$ for $q>2$.  No normality or reducedness hypothesis on
$W$ is used.
\end{theorem}

\begin{proof}
Let \(i:W\hookrightarrow X\) be the given Cartier inclusion.  The
length-one locally free resolution
\[
 0\longrightarrow\mathcal O_X(-W)
 \xrightarrow{\;\cdot g\;}\mathcal O_X
 \longrightarrow i_*\mathcal O_W\longrightarrow0
\]
is locally defined by a nonzerodivisor \(g\); changing \(g\) by a unit
gives the usual transition functions for \(\mathcal O_X(W)\).
Applying sheaf Hom into \(\omega_X\), the degree-zero map is
injective and its cokernel is
\[
 \mathcal E xt^1_X(i_*\mathcal O_W,\omega_X)
 =i_*\bigl((\omega_X\otimes\mathcal O_X(W))|_W\bigr)
 =i_*\omega_W.
\]
The other sheaf Ext groups vanish.  Equivalently, \(W\) is a
two-dimensional Cohen--Macaulay Gorenstein complex space and its
analytic dualizing complex is
\[
 K_W^\bullet\simeq\omega_W[2].
\]
This is the hypersurface calculation in
\cite[Proposition~2 and pp.~90--91]{RamisRuget1970}; no normality or
reducedness of the Cartier divisor is used.

For the locally free sheaf \(F=E|_W\), the analytic duality trace of
\cite[Theorem~2, pp.~82--83]{RamisRuget1970} supplies the actual
functorial pairing, and identifies the dual of \(H_c^q(W,F)\) with
\[
 \operatorname{Ext}^{-q}_W(F,K_W^\bullet)
 =\operatorname{Ext}^{2-q}_W(F,\omega_W)
 =H^{2-q}(W,F^\vee\otimes\omega_W).
\]
The last equality follows because the positive sheaf Ext groups
of a locally free first argument vanish.  Since \(W\) is compact,
\(H_c^q(W,F)=H^q(W,F)\); coherent cohomology is finite-dimensional,
so this pairing is perfect in the ordinary algebraic sense.
It is this trace pairing, rather than an unspecified arrow between
two matching five-term sequences, that constructs the isomorphism
(5.11).  Its contravariance in \(F\) is precomposition in the Ext
argument.

For \(q>2\), the displayed Ext degree \(2-q\) is negative.  The
derived Hom of the locally free sheaf \(F\), concentrated in degree
zero, into the sheaf \(\omega_W\), concentrated in degree zero, has
no negative cohomology.  Hence the duality theorem gives
\(H^q(W,F)=0\).  This also follows directly from coherent
cohomological dimension on the compact surface \(W\).
\end{proof}

\begin{proposition}[absolute duality on a compact Cartier fibre]
\label{prop:cartier-duality-absolute}
\Verified\ Let $X$ be a smooth complex threefold, not necessarily compact,
and let $i:W\hookrightarrow X$ be a compact effective Cartier divisor that
is the full scheme-theoretic fibre of a holomorphic map $f:X\to B$ to a
smooth curve.  For every vector bundle $E$ on $X$ there is a functorial
perfect pairing
\[
 H^2(W,E|_W)\times
 H^0(W,E^\vee|_W\otimes\omega_X|_W)\longrightarrow\mathbb C.
\]
Equivalently,
\[
 H^2(W,E|_W)^\vee\simeq
 H^0(W,E^\vee|_W\otimes\omega_X|_W).
\tag{5.15a}
\]
Neither normality of $W$ nor compactness of the ambient threefold is
required.
\end{proposition}

\begin{proof}
An effective Cartier divisor in a smooth threefold is locally a
two-dimensional hypersurface quotient of a regular local ring.  Hence $W$
is Cohen--Macaulay and Gorenstein.  If $K_W^\bullet$ denotes its analytic
dualizing complex, the complete-intersection calculation of Ramis--Ruget
gives
\[
 K_W^\bullet\simeq\omega_W[2]
\]
with $\omega_W$ invertible
\cite[Proposition~2 and the complete-intersection calculation,
pp.~90--91]{RamisRuget1970}.  Cartier adjunction gives
\[
 \omega_W\simeq
 (\omega_X\otimes\mathcal O_X(W))|_W.
\]
Because $W$ is the full fibre over $b\in B$, one has
$\mathcal O_X(W)\simeq f^*\mathcal O_B(b)$ locally along $W$; its
restriction to $W$ is the pullback of a one-dimensional vector space.
Choosing a nonzero vector identifies
$\mathcal O_X(W)|_W\simeq\mathcal O_W$ and
$\omega_W\simeq\omega_X|_W$.  A different choice rescales the final trace
pairing by a nonzero scalar.

Ramis--Ruget duality pairs $H_c^p(W,\mathcal G)$ with
$\operatorname{Ext}^{-p}_W(\mathcal G,K_W^\bullet)$
\cite[Theorem~2, pp.~82--83]{RamisRuget1970}.  Since $W$ is compact,
$H_c^p=H^p$.  Taking $p=2$ and $\mathcal G=E|_W$, the shift and local
freeness give
\[
 \operatorname{Ext}^{-2}_W(\mathcal G,K_W^\bullet)
 \simeq\operatorname{Hom}_W(\mathcal G,\omega_W)
 =H^0(W,E^\vee|_W\otimes\omega_W).
\]
Grauert finiteness makes both coherent cohomology spaces finite-dimensional
\cite{Grauert1960}; the strong analytic duality is therefore the ordinary
algebraic perfect pairing (5.15a).  Its variance is contravariant in $E|_W$
by precomposition on the Hom term.  This proves the assertion.
\end{proof}

\subsection{Flatness and top-degree analytic base change}

\begin{lemma}[automatic properness, flatness, and Cartier fibres]
\label{lem:curve-flatness}
\Verified\ Let $f:X\to B$ be a nonconstant holomorphic map from a connected
complex manifold to a Riemann surface.  Then $f$ is flat, every holomorphic
vector bundle on $X$ is $f$-flat, and every full scheme-theoretic fibre is an
effective Cartier divisor of pure codimension one.  If $X$ is compact, then
$f$ is proper.
\end{lemma}

\begin{proof}
Fix $x\in X$, $c=f(x)$, and put
$R=\mathcal O_{B,c}$, $S=\mathcal O_{X,x}$.  The ring $R$ is a discrete
valuation ring and $S$ is a regular local domain.  The homomorphism
$R\to S$ is injective.  Indeed, if a nonzero germ $h\in R$ mapped to zero,
$h\circ f$ would vanish on a neighbourhood of $x$.  The zero set of a
nonzero one-variable holomorphic germ is isolated, so $f$ would be constant
on that neighbourhood and therefore, by the identity theorem on connected
$X$, constant everywhere.

Thus every nonzero element of $R$ acts on the domain $S$ as a nonzero
element and hence as a nonzerodivisor.  Therefore $S$ is torsion-free as an
$R$-module.  A module over a discrete valuation ring is flat if and only if
it is torsion-free \cite[Tag~0AUW]{StacksProject}; hence $S$ is $R$-flat.
This holds at every point, so $f$ is flat.  A vector bundle is locally a
finite free $\mathcal O_X$-module, and a finite direct sum of $R$-flat
modules is flat; hence it is $f$-flat.  If $t$ is a local parameter at
$c$, then $f^*t$ is a nonzero element of the domain $S$ and thus a
nonzerodivisor.  The full scheme-theoretic fibre is therefore locally cut
out by one nonzerodivisor, so it is an effective Cartier divisor of pure
codimension one.  Finally, when $X$ is compact, the underlying continuous
map to the Hausdorff surface $B$ is proper.
\end{proof}

\begin{proposition}[the exact analytic base-change input]
\label{prop:analytic-base-change}
\Verified\ Let $f:X\to B$ be proper between complex spaces, let
$\mathcal F$ be coherent and $f$-flat, and fix $c\in B$.  If
\[
H^3(X_c,\mathcal F|_{X_c})=0,
\tag{5.16}
\]
then the canonical fibre map
\[
\beta_c^2:
(R^2f_*\mathcal F)_c\otimes_{\mathcal O_{B,c}}\mathbb C(c)
\longrightarrow H^2(X_c,\mathcal F|_{X_c})
\tag{5.17}
\]
is an isomorphism.
\end{proposition}

\begin{proof}
Apply B\u{a}nic\u{a}--St\u{a}n\u{a}\c{s}il\u{a},
Chapter~III, \S3, Corollary~3.5, with $q=3$
\cite[pp.~120--121]{BanicaStanasila1976}.  Its hypotheses are precisely
properness, coherence, and flatness.  Condition (5.16) first forces the
degree-three direct image to vanish near $c$.  The corollary then applies the
equivalent base-change conditions of its preceding Theorem~3.4 in degree
$q-1=2$ and states that the canonical map (5.17) is bijective.  This is the
top-degree argument used in \cite[Lemma~1.5]{CDP1998}; it does not require
the stronger, false assertion that all direct images are locally free.
\end{proof}

\begin{corollary}[uniform top-degree base change for surface fibres]
\label{cor:surface-top-base-change}
\Verified\ Let $f:X\to B$ be a proper flat holomorphic map to a Riemann
surface whose full scheme-theoretic fibres have pure complex dimension two,
and let $E$ be a holomorphic vector bundle on $X$.  Then
\[
 R^qf_*E=0\quad(q>2)
 \tag{5.17-V}
\]
and, for every $b\in B$, including a point with nonreduced fibre, the
canonical map
\[
 (R^2f_*E)_b\otimes_{\mathcal O_{B,b}}\mathbb C(b)
 \xrightarrow{\ \sim\ }H^2(X_b,E|_{X_b})
 \tag{5.17-T}
\]
is an isomorphism.  No constancy of $h^2$, local freeness of $R^2f_*E$, or
reducedness of $X_b$ is required.
\end{corollary}

\begin{proof}
On a sufficiently small neighbourhood $U$ of $b$, the finite-complex theorem
used in analytic cohomology and base change represents $Rf_*E|_U$ by a
bounded complex $K^\bullet$ of finite free $\mathcal O_U$-modules
\cite[Chapter~III, \S3, Theorem~3.4]{BanicaStanasila1976}.  Flatness of
$E$ over $B$ identifies its derived fibre with the ordinary restriction to
the full scheme-theoretic fibre.  Over the discrete valuation ring
$A=\mathcal O_{B,b}$, the universal-coefficient sequence is
\[
0\longrightarrow H^q(K^\bullet)_b\otimes_A\mathbb C(b)
 \longrightarrow H^q(K^\bullet\otimes_A\mathbb C(b))
 \longrightarrow
 \operatorname{Tor}_1^A(H^{q+1}(K^\bullet)_b,\mathbb C(b))
 \longrightarrow0.
\tag{5.17-U}
\]
If some $R^qf_*E$ with $q>2$ were nonzero, choose the largest such $q$ and a
point of its support.  Nakayama's lemma makes the left term of (5.17-U)
nonzero there, while maximality kills the Tor term.  This would give nonzero
coherent cohomology in degree $q>2$ on a compact two-dimensional complex
space, a contradiction.  Thus (5.17-V) holds.  For $q=2$, the only
obstruction in (5.17-U) is
$\operatorname{Tor}_1^A((R^3f_*E)_b,\mathbb C(b))$, which is zero by
(5.17-V), proving (5.17-T).  This is the explicit obstruction calculation
behind \cref{prop:analytic-base-change}; it was independently supplied, with
eight named local certificates, in Reviewer~B's candidate-specific follow-up
\cite{ReviewerBValidation2026}.
\end{proof}

\subsection{A topology-independent \texorpdfstring{$\operatorname{Pic}^0$}{Pic0} route and three all-Pic routes}

Let $i:W\hookrightarrow X$ denote the cusp-fibre inclusion and
$\eta:\widetilde W\simeq\mathrm{dP}_6\to W$ its normalization.

\begin{proposition}[topologically trivial twists pull back trivially]
\label{prop:pic0-pullback}
\Verified\ Assume that $\omega_X$ is trivial on a neighbourhood of $W$.
For every $L\in\operatorname{Pic}^0(X)$, put
\[
 A_L=(L^\vee\otimes\omega_X)|_W.
\]
Then
\[
 \eta^*A_L\simeq\mathcal O_{\mathrm{dP}_6}.
\tag{5.17a}
\]
No assertion about $H^2(X;\mathbb Z)$ is used.
\end{proposition}

\begin{proof}
Every line bundle in $\operatorname{Pic}^0(X)$ has $c_1(L)=0$.  Naturality
therefore gives $c_1(\eta^*i^*L)=0$.  The degree-six del Pezzo surface is the
blow-up of $\mathbb P^2$ in three noncollinear points, so
$H^1(\mathrm{dP}_6,\mathcal O)=0$.  The exponential sequence consequently
makes
\[
 c_1:\operatorname{Pic}(\mathrm{dP}_6)\longrightarrow
 H^2(\mathrm{dP}_6;\mathbb Z)
\]
injective.  Hence $\eta^*i^*L$ is trivial.  The assumed neighbourhood
trivialization of $\omega_X$ also pulls back trivially, which proves
(5.17a).
\end{proof}

\begin{proposition}[the exact all-line-bundle restriction obstruction]
\label{prop:all-pic-restriction}
\Verified\ Define
\[
 r=\eta^*i^*:\operatorname{Pic}(X)\longrightarrow
 \operatorname{Pic}(\mathrm{dP}_6).
\]
Under the same local canonical trivialization,
\[
 \eta^*A_L\simeq r(L)^\vee
 \qquad(L\in\operatorname{Pic}(X)).
\tag{5.17b}
\]
Thus the trivial-pullback hypothesis used by
\cref{prop:conductor} holds simultaneously for every line bundle exactly
when $r=0$.  It always holds on $\operatorname{Pic}^0(X)$.
\end{proposition}

\begin{proof}
Pullback commutes with tensor product and duals, while the local
trivialization gives $\eta^*i^*\omega_X\simeq\mathcal O$.  This is (5.17b),
and its uniform equivalence with $r=0$ is immediate.  The final assertion is
\cref{prop:pic0-pullback}.
\end{proof}

\begin{proposition}[Leray-edge vanishing of normalized cusp restrictions]
\label{prop:leray-edge-restriction}
\Verified\ For the constructed map $f:X\to B$ and every
$L\in\operatorname{Pic}(X)$, the degree-two Leray edge section satisfies
\[
 e(c_1(L))=0\in H^0(B,R^2f_*\mathbb Z),
 \qquad
 \eta^*i^*L\simeq\mathcal O_{\mathrm{dP}_6}.
\tag{5.17c}
\]
This proof does not invoke the packaged conclusion
$H^2(X;\mathbb Z)=0$.
\end{proposition}

\begin{proof}
First record explicitly the compatibility used here.  For every $b\in B$,
proper base change and functoriality of the filtered-complex spectral
sequence give a commutative square
\[
\begin{CD}
 H^2(X;\mathbb Z) @>{e}>> H^0(B,R^2f_*\mathbb Z)\\
 @V{i_b^*}VV @VV{\operatorname{ev}_b}V\\
 H^2(X_b;\mathbb Z) @= H^2(X_b;\mathbb Z).
\end{CD}
\tag{5.17d}
\]
Indeed, pull the Leray filtered complex back along $\{b\}\to B$; the
spectral sequence over a point has only column zero, so its edge is the
identity, while the right vertical map is the proper-base-change
identification.  This uses the functoriality and convergence statements in
\cite[Tags~012O, 012U and~012W]{StacksProject} and analytic proper base
change \cite[Theorem~2.3.26]{Dimca2004}.  The connecting homomorphism in the
exponential sequence is natural under $i_b^*$, so (5.17d) sends
$c_1(L)$ to $c_1(L|_{X_b})$.

Put $s_L=e(c_1(L))$.  The image of an edge morphism consists of permanent
cycles, hence
\[
 s_L\in E_\infty^{0,2}\subseteq\ker
 \bigl(d_2^{0,2}:E_2^{0,2}\to E_2^{2,1}\bigr);
\tag{5.17e}
\]
see \cite[Tags~011N, 011O and~012U]{StacksProject}.  For the constructed
candidate, \cref{prop:leray-e2,lem:higher-transgressions} give
\[
 E_2^{0,2}=\mathbb Z(2q),\qquad
 E_2^{2,1}=\mathbb Z\,\omega[\delta],\qquad
 d_2^{0,2}(2q)=\omega[\delta].
\]
Thus $d_2^{0,2}$ is an isomorphism and (5.17e) forces $s_L=0$.
Taking $b=p_0$ in (5.17d) gives $c_1(i^*L)=0$, and hence
$c_1(\eta^*i^*L)=0$.  Finally, the injectivity of
$c_1:\operatorname{Pic}(\mathrm{dP}_6)\to H^2(\mathrm{dP}_6;\mathbb Z)$
proved in \cref{prop:pic0-pullback} makes $\eta^*i^*L$ trivial.
\end{proof}

\begin{lemma}[the pair-of-pants defect class]
\label{lem:picard-crossed-defect}
\Verified\ Choose the compact pair-of-pants core $P\subset B^\circ$ and
the zero-section whiskers used in \cref{prop:geometric-pi1}, and put
$J_P=f^{-1}(P)$.  Orient its boundary meridians clockwise.  Write
$\rho_0$ for the projection of the zero-section cusp meridian
$\widetilde a_0$.  The constructed candidate uses the tautological cusp
gluing, whose affine regluing class is $\mu=0$ (equivalently
$\ell_0=\gamma(\mu)=0$); hence the flat-marked cusp lift has no vertical
translation.  The base loops and their chosen lifts therefore satisfy
\[
 \rho_1\rho_2\rho_0=1,
 \qquad
 \widehat\rho_1\widehat\rho_2\widetilde a_0=1.
\tag{5.17g}
\]
For $L\in\operatorname{Pic}(X)$, write
$E=c_1(L|_F)\in\bigwedge^2V$.  Then $E\in(\bigwedge^2V)^G$, and there is a
crossed homomorphism
\[
 z:\pi_1(P)\longrightarrow V_{\mathbb R},\qquad
 z(gh)=z(g)+T_gz(h),
\tag{5.17h}
\]
whose class records
$c_1(L|_{J_P})-[\Omega_E]$.  Here $\Omega_E$ is the closed, translation-
invariant vertical two-form induced by $E$ in the flat marking.  If
$d_i=z(\rho_i)$ and a bar denotes the image in
$V_G\otimes\mathbb R$, then
\[
 \bar d_1+\bar d_2+\bar d_0=0.
\tag{5.17i}
\]
The three barred classes, and the number
$\langle d_0,\widehat\delta\rangle$, do not depend on the chosen crossed-
homomorphism representative.  More precisely, use the base-first convention
for the edge injection in (5.17j).  For any oriented loop $\rho\subset P$
and any $A_\rho$-invariant $\nu\in\Lambda$, its flat-marked sweep satisfies
\[
 \bigl\langle z(\rho),\nu\bigr\rangle
 =-\bigl\langle c_1(L|_{J_P})-[\Omega_E],
 a_*([S^1_\nu]\times[\widetilde\rho])\bigr\rangle.
\tag{5.17j'}
\]
Here $a$ is the fibrewise translation action and $\widetilde\rho$ is the
flat-marked lift of $\rho$.
\end{lemma}

\begin{proof}
The flat marking identifies the pullback of $J_P$ to the universal cover of
$P$ with
$\widetilde P\times(\Lambda_{\mathbb R}/\Lambda)$.  Since $E$ is the
restriction of a class on the total space, it is fixed by every monodromy
$T_g$.  The constant vertical form representing $E$ is therefore invariant
under the deck action and descends to the asserted form $\Omega_E$ on
$J_P$.

The pair of pants retracts onto a graph.  Hence the Serre spectral sequence
of the marked real torus bundle has only columns zero and one and gives the
short exact edge sequence
\[
0\longrightarrow H^1(P;V_{\mathbb R})
 \longrightarrow H^2(J_P;\mathbb R)
 \longrightarrow H^0(P;\textstyle\bigwedge^2V_{\mathbb R})
 \longrightarrow0.
\tag{5.17j}
\]
Both $c_1(L|_{J_P})$ and $[\Omega_E]$ map to $E$, so their difference is
the image of a unique class in $H^1(P;V_{\mathbb R})$.  The standard group-
cohomology model for a graph identifies this group with crossed
homomorphisms modulo principal crossed homomorphisms
\cite[Chapter~I, \S5]{Brown1982}; this gives (5.17h).

For completeness, the slant convention in (5.17j') is part of this
identification, not an additional geometric assertion.  Give an oriented
edge of a graph spine of $P$ a coordinate $s$ increasing with its based
loop and represent the corresponding $V_{\mathbb R}$-valued cellular
one-cochain by a base-first form $ds\wedge\beta_s$.  On the sweep of an
$A_\rho$-invariant fibre circle $S^1_\nu$, Fubini's formula gives
\[
\int_{S^1_s\times S^1_\nu}ds\wedge\beta_s
 =\int_{S^1_s}\langle\beta_s,\nu\rangle,ds
 =\langle z(\rho),\nu\rangle.
\]
The sweep in (5.17j') is written in the fibre-first order
$[S^1_\nu]\times[\widetilde\rho]$, which differs from this base-first
cross-product orientation by $(-1)^{1\cdot1}=-1$; this is the displayed
minus sign.
Changing the cellular representative by a principal cochain changes both
sides by the same boundary term, which vanishes because $A_\rho\nu=\nu$.
This proves (5.17j').

Applying $z$ to (5.17g) gives
\[
 d_1+T_1d_2+T_1T_2d_0=0.
\]
All $T_g$ act trivially on the global coinvariants, proving (5.17i).  A
change of representative has the form
$z(g)\mapsto z(g)+(T_g-I)a$.  Its image in $V_G\otimes\mathbb R$ is zero.
Moreover the displayed matrices in
\cref{prop:monodromy-certificate} fix $\widehat\delta$ under every dual
action $A_g$.  Therefore
\[
 \langle(T_g-I)a,\widehat\delta\rangle
 =\langle a,(A_g^{-1}-I)\widehat\delta\rangle=0,
\]
which proves the final assertion.
\end{proof}

\begin{lemma}[the finite-cover norm identity]
\label{lem:picard-finite-norm}
\Verified\ Retain the notation of
\cref{prop:candidate-finite-fillings,lem:picard-crossed-defect}.  For
$j=1,2$, let
\[
 N_{T_j}=I+T_j+\cdots+T_j^{m_j-1}.
\]
With the clockwise conventions of (5.17g), the pullback to the
$m_j$-fold $s_j$-cover gives
\[
 N_{T_j}d_j=\iota_{v_j}E
 \quad\text{in }V_{\mathbb R},
\qquad
 m_j\bar d_j=[\iota_{v_j}E]
 \quad\text{in }V_G\otimes\mathbb R.
\tag{5.17k}
\]
\end{lemma}

\begin{proof}
Restriction of the crossed homomorphism to the subgroup generated by
$\rho_j^{m_j}$ has value
\[
 z(\rho_j^{m_j})
 =\sum_{r=0}^{m_j-1}T_j^rz(\rho_j)=N_{T_j}d_j.
\tag{5.17l}
\]
This is the cohomological pullback of the annular defect to the
$s_j$-cover; the norm, rather than unqualified multiplication by $m_j$, is
essential before passage to global coinvariants.  A principal change
$d_j\mapsto d_j+(T_j-I)a$ does not alter this value, since
$N_{T_j}(T_j-I)a=(T_j^{m_j}-I)a=0$.

Let $Q_j:J'_j\to N_j$ be the finite quotient map.  The global flat marking
$\Xi$ over the full disc $\Delta_j$ makes the constant vertical form
$\Omega_E$ well defined on $J'_j$.  Radial retraction onto the central
torus makes
\[
 H^2(J'_j;\mathbb R)\longrightarrow H^2(A_0^{(j)};\mathbb R)
\]
an isomorphism.  On the central fibre the quotient is exactly $\pi_j$, and
naturality gives
\[
 (Q_j|_{A_0^{(j)}})^*c_1(L|_{S_j})
 =\pi_j^*c_1(L|_{S_j})=E.
\tag{5.17l'}
\]
The second equality is the actual radial/stalk identification in
\cref{lem:finite-filling-stalk-map}; on a nearby fibre the gluing map $h_j$
is a translation and therefore acts trivially on $H^2$ of that fibre.
Consequently the two classes
$c_1(Q_j^*(L|_{N_j}))$ and $[\Omega_E]$ have the same central restriction,
and the retraction isomorphism proves
\[
 c_1(Q_j^*(L|_{N_j}))-[\Omega_E]=0
 \quad\text{in }H^2(J'_j;\mathbb R).
\tag{5.17l''}
\]

On $\Delta_j^*$, the source identity
$h_j\widetilde g_j^{\log}h_j^{-1}=\widetilde g_j^0$ shows that $h_j$
descends from the logarithmic quotient $N_j^*$ to the linear/global
quotient $J^*$.  Thus
\[
 Q_j^*(L|_{N_j})|_{\Delta_j^*}
 \simeq h_j^*(L|_{J^*}).
\tag{5.17l'''}
\]
The logarithmic overlap is the holomorphic translation
\[
 h_j[s,\zeta]=[s,\zeta+\alpha(s)\Pi(s)v_j],
 \qquad \alpha(s)=\frac{\log s_j}{2\pi i}.
\]
To compute its effect on a real cohomology class, retain the exact
real marking
\[
 \Xi(s,x+\Lambda)=[s,\Pi(s)x],\qquad
 J_s=\Pi(s)^{-1}\circ i\circ\Pi(s)
 \quad\text{on }\Lambda_{\mathbb R}.
\]
The period matrix \(\Pi(s)\) is a real isomorphism by (2.2)--(2.4);
thus \(J_s\) is a well-defined smooth family of real complex
structures.  On a branch of the logarithm the actual real
translation vector is
\[
 a_j(s)=\operatorname{Re}\alpha(s)\,v_j+
        \operatorname{Im}\alpha(s)\,J_s v_j.
 \tag{5.17m}
\]
Indeed, applying \(\Pi(s)\) gives precisely
\(\alpha(s)\Pi(s)v_j\).  A different logarithm branch changes
\(a_j\) by an integral multiple of \(v_j\), so its torus translation
is branch-independent.  Complex multiplication by \(\alpha\) is
not multiplication by a real scalar in this marking.

Write \(E=\frac12\sum_{a,b}E_{ab}\,dx^a\wedge dx^b\) in the
ordered real coordinates dual to the marked lattice; the matrix
\((E_{ab})\) is alternating and is constant over the covering
disc.  Direct substitution in the real map
\(h_j(s,x)=(s,x+a_j(s))\) gives
\[
 \begin{split}
 h_j^*\Omega_E-\Omega_E
  ={}&\sum_{a,b}E_{ab}\,d a_j^a\wedge dx^b\\
    &+\frac12\sum_{a,b}E_{ab}\,d a_j^a\wedge d a_j^b.
 \end{split}
 \tag{5.17m'}
\]
The last term has both arguments in the base directions.
In particular, it restricts to zero when the base is an oriented
circle.

Take one full clockwise circle in the \(s_j\)-cover,
\[
 s_j(\theta)=r e^{-2\pi i\theta},\qquad 0\le\theta\le1.
\]
Choose the logarithm continuously along that circle, so
\(\alpha(1)=\alpha(0)-1\).  The endpoints represent the same
point of the covering disc.  The marking is defined on the full
disc, and therefore \(J_{s_j(1)}=J_{s_j(0)}\).  Consequently
\[
 a_j(1)-a_j(0)=-v_j.
 \tag{5.17m''}
\]
For any \(\lambda\in\Lambda\), restrict (5.17m') to the sweep of
the marked circle \(x(u)=u\lambda\), using the base-first
orientation \(d\theta\wedge du\).  Its integral is exactly
\[
 \int_0^1 E\!\left(\frac{d a_j}{d\theta},\lambda\right)d\theta
 =E(a_j(1)-a_j(0),\lambda)
 =-E(v_j,\lambda).
\]
The horizontal term contributes zero.  Since marked lattice
circles span \(H_1(F;\mathbb R)\), these integrals determine the
mixed component in
\[
 H^1(S^1;\mathbb R)\otimes H^1(F;\mathbb R).
\]
Thus the mixed component of \(h_j^*\Omega_E-\Omega_E\) is
\(-\iota_{v_j}E\), with \((\iota_{v_j}E)(\lambda)=E(v_j,\lambda)\).

The pullback of the defect
\(c_1(L|_{J^*})-[\Omega_E]\) has zero fibre component and mixed
component \(N_{T_j}d_j\) by (5.17l).  Translation by \(h_j\) does
not change this mixed component: on the covering circle such
a fibre-zero class is represented by
\(d\theta\wedge\beta\), with \(\beta\) a constant fibre one-form,
and substituting \(dx\mapsto dx+d a_j\) adds only a base two-form,
which vanishes on the circle.  Equations (5.17l'') and
(5.17l''') state that the total pulled-back class extends over
the covering disc and has zero mixed component.  It follows that
\[
 N_{T_j}d_j-\iota_{v_j}E=0.
\]
In the global coinvariants every \(T_j\) acts as the identity,
so \(m_j\bar d_j=[\iota_{v_j}E]\).  This proves (5.17k), with the
original clockwise convention and no change to either translation
vector.
\end{proof}

\begin{lemma}[the cusp boundary defect and Chern-class specialization]
\label{lem:picard-cusp-defect}
\Verified\ In the notation of
\cref{lem:picard-crossed-defect}, one has
\[
 \bar d_0=0\quad\text{in }V_G\otimes\mathbb R.
\tag{5.17n}
\]
Moreover, for every line bundle $L$ on the cusp neighbourhood $N_0$,
nearby-cycle specialization is compatible with first Chern classes:
\[
 \operatorname{sp}^2\bigl(c_1(L|_W)\bigr)=c_1(L|_F).
\tag{5.17o}
\]
\end{lemma}

\begin{proof}
Let $\widetilde a_0$ be the flat-marked clockwise zero-section meridian in
the lifted relation of (5.17g), and let
\[
 \xi_{\widehat\delta}
 =a_*([S^1_{\widehat\delta}]\times[\widetilde a_0])
\]
be its sweep under fibrewise translation by $\widehat\delta$.  The boundary
$\rho_0$ of $P$ was chosen to be its projection.  This is the zero-section
sweep: there is no hidden vertical drift in $\widetilde a_0$, by the
explicit overlap calculation in \cref{prop:geometric-pi1}.  Hence the
vertical form $\Omega_E$ vanishes on the swept torus, and the slant identity
(5.17j') gives
\[
 \langle d_0,\widehat\delta\rangle
 =-\langle c_1(L|_{J_P})-[\Omega_E],
          \xi_{\widehat\delta}\rangle
 =-\langle c_1(L|_{J_P}),\xi_{\widehat\delta}\rangle.
\tag{5.17p}
\]
Combine the exponential overlap coordinates $(x_1,x_2,t)$ of (2.8ii) with
the unimodular cusp chart
\[
 (u,x_1,x_2)=\left(\frac{t}{x_1x_2},x_1,x_2\right).
\]
The literal solid torus
\[
 \mathbb S_{\widehat\delta}
 =\{|u|=\rho,\ x_1=1,\ |x_2|\le1\}
\]
has boundary parametrization
\[
 (\theta,\phi)\longmapsto
 \bigl(\rho e^{i(\theta-\phi)},1,e^{i\phi}\bigr)
\tag{5.17p0}
\]
and hence $t=\rho e^{i\theta}$; its $\phi$-circle is translation by
$\widehat\delta$.  Direct evaluation in the proved cusp marking therefore
identifies $\partial\mathbb S_{\widehat\delta}$ with
$\xi_{\widehat\delta}$.  Reversing $\theta$ supplies the fixed clockwise
convention and only changes the common orientation sign.  No reduced
Mayer--Vietoris block or Leray comparison is used here.
Since $L$ extends across $N_0$, its Chern class evaluates to zero on that
boundary.  Thus the left side of (5.17p) is zero.  By
\cref{thm:coinvariants}, $V_G=\mathbb Z[\delta]$, and evaluation on the
globally invariant dual vector $\widehat\delta$ identifies
$V_G\otimes\mathbb R$ with $\mathbb R$.  This proves (5.17n).

For (5.17o), let $\ell=c_1(L|_{N_0})$, let
$\iota_0:W\hookrightarrow N_0$, and let
$\iota_t:F_t\hookrightarrow N_0$ be a sufficiently near smooth cusp fibre.
The neighbourhood realization of the nearby-cycle unit
$\mathbb Z_W\to\psi_{f_0}\mathbb Z$ recalled in the proof of
\cref{thm:p0-integral-specialization} gives the commutative square
\[
\begin{CD}
 H^2(N_0;\mathbb Z) @>{\iota_0^*}>> H^2(W;\mathbb Z)\\
 @V{\iota_t^*}VV @VV{\operatorname{sp}^2}V\\
 H^2(F_t;\mathbb Z) @= H^2(F_t;\mathbb Z).
\end{CD}
\tag{5.17p'}
\]
Applying this square to $\ell$ gives
$\operatorname{sp}^2(c_1(L|_W))=c_1(L|_{F_t})$.  Flat parallel transport
in the marked smooth family identifies the latter with $c_1(L|_F)=E$.
Equivalently, $c_1$ is the connecting morphism in the exponential sequence,
and restriction to the shrinking neighbourhoods and to $F_t$ commutes with
that morphism before passage to the direct limit.  This proves the claimed
naturality without using any global Leray differential.
\end{proof}

\begin{proposition}[direct boundary-Picard vanishing]
\label{prop:direct-boundary-picard}
\Verified\ For the constructed threefold and every analytic line bundle
$L\in\operatorname{Pic}(X)$,
\[
 c_1(L|_F)=0,
 \qquad
 \eta^*i^*L\simeq\mathcal O_{\mathrm{dP}_6}.
\tag{5.17q}
\]
This proof is independent of the global Leray transgressions
$d_2^{0,q}$ and of the packaged conclusion $H^2(X;\mathbb Z)=0$; it uses
only the actual finite-filling lattices, their logarithmic overlap maps, the
based boundary relation, the local cusp solid torus, and the local integral
specialization theorem.
\end{proposition}

\begin{proof}
The common invariant calculation in \cref{thm:coinvariants} first gives
\[
 E=c_1(L|_F)=kq.
\]
Because $L$ extends across the order-four filling, its fibre class belongs
to the actual specialization lattice
\[
 \pi_2^*H^2(S_2;\mathbb Z)
 =\{a\gamma\eta_2+bq:a\equiv b\pmod2\}
\]
of \cref{prop:finite-indices,lem:finite-filling-stalk-map}.  Taking
coefficient independence in the displayed invariant basis forces
$a=0$ and $b=k$; hence $k$ is even.  Write $E=2nq$.

Contraction with the two affine translation vectors is an exact integral
calculation:
\[
 \iota_{v_1}(2q)=4\eta_1,
 \qquad
 \iota_{v_2}(2q)=-6\eta_2.
\tag{5.17r}
\]
Indeed $v_1=(1,2,-4,0)^t$, $v_2=(-1,-3,3,0)^t$ in the basis dual to
$(\gamma,u,w,\delta)$, while
$q=uw+6\gamma\delta$.  In the global coinvariants,
\[
 [\eta_1]=3[\delta],\qquad [\eta_2]=2[\delta].
\]
Applying \cref{lem:picard-finite-norm} with $m_1=3$, $m_2=4$ therefore
gives
\[
 \bar d_1=4n[\delta],\qquad
 \bar d_2=-3n[\delta].
\tag{5.17s}
\]
Now combine the balance relation (5.17i) with the cusp vanishing (5.17n):
\[
 0=\bar d_1+\bar d_2+\bar d_0=n[\delta].
\]
The coinvariant group is torsion-free, so $n=0$ and hence $E=0$.

Finally (5.17o) and the injectivity of the local integral specialization in
\cref{thm:p0-integral-specialization} imply
$c_1(L|_W)=0$.  Pulling back to the normalization gives
$c_1(\eta^*i^*L)=0$.  The injectivity of
$c_1:\operatorname{Pic}(\mathrm{dP}_6)\to
H^2(\mathrm{dP}_6;\mathbb Z)$ proved in
\cref{prop:pic0-pullback} yields the second assertion of (5.17q).
\end{proof}

\subsection{An analytic completion of the all-Picard certificate}
\label{subsec:analytic-all-picard}

The local example in \cref{prop:local-cusp-picard-obstruction} shows that
the cusp neighbourhood alone cannot annihilate normalized central
restrictions.  We now close the stronger analytic route by retaining the
order-three and order-four fillings.  The proof has two independent parts.
First, logarithms of the actual \(\mathcal O^*\)-factors give the boundary
cocycle and its two finite-filling norm identities.  Second, a
candidate-specific support-function calculation reads every possible
central restriction directly from the mixed Appell--Humbert block on a
nearby torus.  No global Leray differential, nearby-cycle specialization,
or Chern-class evaluation on a cusp bounding chain occurs.

\begin{proposition}[analytic Appell--Humbert boundary cocycle]
\label{prop:analytic-ah-boundary}
\Verified\ Let \(P\subset B^\circ\) be the compact pair-of-pants core,
orient its boundary meridians clockwise, and retain the zero-section lifts
of \cref{prop:geometric-pi1}.  Thus
\[
 \rho_1\rho_2\rho_0=1,
 \qquad
 \widehat\rho_1\widehat\rho_2\widetilde a_0=1,
\tag{5.17-A1}
\]
and the second relation has no terminal vertical translation.  For
\(L\in\operatorname{Pic}(X)\), let
\[
 E=c_1(L|_F)\in\operatorname{Alt}^2(\Lambda,\mathbb Z)
   =\bigwedge^2V
\]
be its fibrewise Appell--Humbert form.  Put
\(V_{\mathbb R}=\operatorname{Hom}(\Lambda,\mathbb R)\), and use
\[
 A_g=(T_g^{-1})^{\mathsf t},
 \qquad
 \langle T_g\alpha,\lambda\rangle
 =\langle\alpha,A_g^{-1}\lambda\rangle.
\tag{5.17-A2}
\]
There is a crossed homomorphism, unique up to a principal crossed
homomorphism,
\[
 d:\pi_1(P)\longrightarrow V_{\mathbb R},
 \qquad
 d(gh)=d(g)+T_gd(h),
 \qquad
 d(g)\sim d(g)+(T_g-I)a.
\tag{5.17-A3}
\]
If \(d_i=d(\rho_i)\), then
\[
 d_1+T_1d_2+T_1T_2d_0=0,
 \qquad
 \bar d_1+\bar d_2+\bar d_0=0
 \quad\text{in }V_G\otimes\mathbb R.
\tag{5.17-A4}
\]
For the finite filling of order \(m_j\), \(j=1,2\),
\[
 N_{T_j}d_j=\iota_{v_j}E,
 \qquad
 m_j\bar d_j=[\iota_{v_j}E],
 \qquad
 N_{T_j}=I+T_j+\cdots+T_j^{m_j-1},
\tag{5.17-A5}
\]
where \((\iota_vE)(\lambda)=E(v,\lambda)\).  At the cusp,
\[
 \langle d_0,\widehat\delta\rangle=0,
 \qquad
 \bar d_0=0\quad\text{in }V_G\otimes\mathbb R.
\tag{5.17-A6}
\]
These identities are obtained from holomorphic transition factors and the
argument principle.
\end{proposition}

\begin{proof}
Choose an open Riemann-surface neighbourhood \(P^{\rm an}\Subset B^\circ\)
which deformation retracts onto \(P\), and pull the period family to its
universal cover.  Before quotienting by the fibre lattice, the resulting
space is
\[
 \widetilde{\mathcal J}
 =\widetilde P^{\rm an}\times\mathbb C^2.
\]
It is Stein and contractible, so the pullback of \(L\) is holomorphically
trivial.  In a nowhere-zero frame, write the factor belonging to a deck
transformation \(\alpha\) as
\[
 a_\alpha\in\mathcal O^*(\widetilde{\mathcal J}),
 \qquad
 a_{\alpha\beta}(x)=a_\alpha(\beta x)a_\beta(x).
\tag{5.17-A7}
\]
The space is simply connected, hence \(a_\alpha=e^{2\pi i h_\alpha}\) for
a holomorphic logarithm.  Therefore
\[
 c(\alpha,\beta)
 =h_\alpha(\beta x)+h_\beta(x)-h_{\alpha\beta}(x)\in\mathbb Z
\tag{5.17-A8}
\]
is independent of \(x\), and associativity makes \(c\) an integral group
two-cocycle.

Let \(\tau_\lambda\) denote fibre translation by
\(\Pi(z)\lambda\), \(\lambda\in\Lambda\), and let \(\widehat g\) be the
zero-section lift of a base deck transformation.  Period equivariance gives
\[
 \widehat g\,\tau_\lambda\,\widehat g^{-1}
 =\tau_{A_g\lambda}.
\tag{5.17-A9}
\]
On fibre translations, the commutator of (5.17-A8) is
\[
 E(\lambda,\mu)
 =c(\tau_\lambda,\tau_\mu)-c(\tau_\mu,\tau_\lambda).
\tag{5.17-A10}
\]
The four oriented sides of the \((\lambda,\mu)\)-parallelogram identify
this integer with the exponential connecting class of \(L|_F\); this is
the Appell--Humbert convention used in
(5.17av)--(5.17ax), and agrees with the standard factor-of-automorphy
calculation \cite[Chapter~1, \S2, pp.~5--8, especially
Theorem~2.1, p.~7]{BirkenhakeLange1999}%
\cite[equations~(3.3)--(3.10), pp.~3--7]{SchnellLecture3}.
Conjugating the commutator by \(\widehat g\) yields
\[
 E(A_g\lambda,A_g\mu)=E(\lambda,\mu).
\tag{5.17-A11}
\]

Push the logarithmic central extension defined by \(c\) out along
\(\mathbb Z\hookrightarrow\mathbb R\), and denote its central element of
weight \(r\) by \(\mathbf c(r)\).  Choose lifts \(U_\lambda\) of the fibre
translations satisfying the real Appell--Humbert normalization
\[
 U_\lambda U_\mu
 =\mathbf c\!\left(\frac12E(\lambda,\mu)\right)U_{\lambda+\mu}.
\tag{5.17-A12}
\]
To see that this normalization exists, start with arbitrary lifts.  The
difference between their real cocycle and \(E/2\) has zero commutator, so
it defines an abelian extension of the free abelian group \(\Lambda\) by
\(\mathbb R\).  Choose lifts of an integral basis and extend
multiplicatively; this splits that abelian extension and gives (5.17-A12).
If \(U'_\lambda\) is a second normalized system, then
\(U'_\lambda U_\lambda^{-1}\) is central, say
\(\mathbf c(\langle a,\lambda\rangle)\).  Comparing the two copies of
(5.17-A12) proves additivity in \(\lambda\), so
\(a\in V_{\mathbb R}\).  Thus every two normalized systems differ by exactly
the character change used below.

For \(g\in\pi_1(P)\), choose a lift \(G_g\) of \(\widehat g\) and define
\(d(g)\in V_{\mathbb R}\) by
\[
 G_gU_{A_g^{-1}\lambda}G_g^{-1}
 =\mathbf c\bigl(\langle d(g),\lambda\rangle\bigr)U_\lambda.
\tag{5.17-A13}
\]
The exponent is additive in \(\lambda\): conjugate (5.17-A12) and use
(5.17-A11).  The lifts of \(gh\) and \(G_gG_h\) differ centrally, so
\[
\begin{aligned}
 \langle d(gh),\lambda\rangle
 &=\langle d(g),\lambda\rangle
   +\langle d(h),A_g^{-1}\lambda\rangle\\
 &=\langle d(g)+T_gd(h),\lambda\rangle.
\end{aligned}
\]
This proves the crossed-homomorphism law.  Replacing
\(U_\lambda\) by
\(\mathbf c(\langle a,\lambda\rangle)U_\lambda\) changes (5.17-A13) by
\((T_g-I)a\).  Replacing \(G_g\) by a central multiple does not change
(5.17-A13), so this proves the stated independence and uniqueness.  Applying
the crossed law to
(5.17-A1) proves the first identity in (5.17-A4); passage to coinvariants
proves the second.

For the order-\(m_j\) filling, distinguish every covering level.  Put
\[
 \widetilde J_j=\Delta_j\times\mathbb C^2,\qquad
 r_j:\widetilde J_j\longrightarrow J'_j
\]
for the lattice quotient, whose deck translations are \(\tau_\lambda\).
On the vector cover define
\[
 Q_{j,\log}:=\widetilde g_j^{\,v_j/m_j},\qquad
 Q_{j,\mathrm{lin}}:=\widetilde g_j^{\,0}.
\]
They satisfy
\[
\begin{aligned}
 Q_{j,\log}\tau_\lambda Q_{j,\log}^{-1}
  &=\tau_{A_j\lambda},&
 (Q_{j,\log})^{m_j}&=\tau_{v_j},\\
 Q_{j,\mathrm{lin}}\tau_\lambda Q_{j,\mathrm{lin}}^{-1}
  &=\tau_{A_j\lambda},&
 (Q_{j,\mathrm{lin}})^{m_j}&=1.
\end{aligned}
\tag{5.17-A14}
\]
The positive sign is fixed by
\(s_j\circ g_j=e^{-2\pi i/m_j}s_j\) and (2.8f): after exactly
\(m_j\) clockwise iterations, the logarithmic vector-cover lift translates
by \(+v_j\).

Let \(q_{j,\log},q_{j,\mathrm{lin}}\) be the descended automorphisms of \(J'_j\),
and put \(J_j^{\prime *}=J'_j|_{\Delta_j^*}\).  Both descended automorphisms have
order dividing \(m_j\), because \(\tau_{v_j}\) is trivial after the lattice
quotient; each covers \(g_j\), whose order is exactly \(m_j\), so each
descended order is exactly \(m_j\).
The branch-independent torus translation
\[
 h_j:J_j^{\prime *}\longrightarrow J_j^{\prime *}
\]
constructed in (2.9d) satisfies, on this quotient level,
\[
 h_jq_{j,\log}h_j^{-1}=q_{j,\mathrm{lin}}.
\]
Define
\[
\begin{aligned}
 p_{j,\log}:J_j^{\prime *}&\longrightarrow
  N_j^*:=J_j^{\prime *}/\langle q_{j,\log}\rangle,\\
 p_{j,\mathrm{lin}}:J_j^{\prime *}&\longrightarrow
 J_j^*:=J_j^{\prime *}/\langle q_{j,\mathrm{lin}}\rangle .
\end{aligned}
\]
The conjugacy induces a biholomorphism
\[
 \bar h_j:N_j^*\longrightarrow J_j^*,\qquad
 p_{j,\mathrm{lin}}\circ h_j=\bar h_j\circ p_{j,\log}.
\]
The gluing of \(X\) identifies
\(L|_{N_j^*}\simeq\bar h_j^*(L|_{J_j^*})\); hence the exactly typed
line-bundle square gives
\[
 p_{j,\log}^*(L|_{N_j^*})
 \simeq h_j^*p_{j,\mathrm{lin}}^*(L|_{J_j^*}).
\tag{5.17-A14a}
\]

For the factor extension return to the filled vector cover.  Let
\[
 P_{j,\log}:J'_j\longrightarrow N_j,\qquad
 \pi_j^{\mathrm{vec}}:=P_{j,\log}\circ r_j:
 \widetilde J_j\longrightarrow N_j .
\]
Here \(P_{j,\log}\) is the full quotient and
\(P_{j,\log}|_{J_j^{\prime *}}=p_{j,\log}\).
The deck group is
\[
 \Gamma_j^{\log}
 =\langle\tau_\lambda,Q_{j,\log}:\lambda\in\Lambda\rangle.
\]
The pullback \((\pi_j^{\mathrm{vec}})^*(L|_{N_j})\) is trivial on the
contractible Stein manifold \(\widetilde J_j\); its factors define an
integral central extension, which we push out along
\(\mathbb Z\hookrightarrow\mathbb R\).  Over \(\Delta_j^*\), the
line-bundle square (5.17-A14a) identifies each fibre with the global factor
extension by a translation, so its commutator form is \(E\).  That form is
integer-valued and locally constant on the connected disc \(\Delta_j\);
therefore it is \(E\) on the filled cover as well.  Choose vertical lifts
\(U_\lambda^{(j)}\) normalized by
\[
 U_\lambda^{(j)}U_\mu^{(j)}
 =\mathbf c\!\left(\frac12E(\lambda,\mu)\right)
  U_{\lambda+\mu}^{(j)}.
\]
For \(\alpha\in\Gamma_j^{\log}\), define \(A_\alpha,T_\alpha\) by
\[
 \alpha\tau_\lambda\alpha^{-1}=\tau_{A_\alpha\lambda},\qquad
 T_\alpha=(A_\alpha^{-1})^{\mathsf t}.
\]
For a lift \(\widetilde G_\alpha\), define
\(\widetilde d_j(\alpha)\in V_{\mathbb R}\) through
\[
 \widetilde G_\alpha U_{A_\alpha^{-1}\lambda}^{(j)}
 \widetilde G_\alpha^{-1}
 =\mathbf c\!\left(\langle\widetilde d_j(\alpha),\lambda\rangle\right)
  U_\lambda^{(j)}.
\tag{5.17-A14c}
\]
Indeed, before naming \(\widetilde d_j(\alpha)\), the quotient of the two
sides of (5.17-A14c) is central because both cover the same translation.
Conjugating the normalized product and using
\[
 E(A_\alpha\lambda,A_\alpha\mu)=E(\lambda,\mu)
\]
shows that its exponent \(r_\alpha(\lambda)\) satisfies
\[
 r_\alpha(\lambda+\mu)=r_\alpha(\lambda)+r_\alpha(\mu).
\]
Thus it is the pairing with a unique element of \(V_{\mathbb R}\), as
asserted.  Since \(A_{\alpha\beta}=A_\alpha A_\beta\), choose the lift of
\(\alpha\beta\) to be \(\widetilde G_\alpha\widetilde G_\beta\); any other
choice differs centrally.  Then
\[
\begin{aligned}
 &\widetilde G_\alpha\widetilde G_\beta
 U_{A_\beta^{-1}A_\alpha^{-1}\lambda}^{(j)}
 \widetilde G_\beta^{-1}\widetilde G_\alpha^{-1}\\
 &\quad=
 \mathbf c\!\left(
  \langle\widetilde d_j(\beta),A_\alpha^{-1}\lambda\rangle
  +\langle\widetilde d_j(\alpha),\lambda\rangle\right)
 U_\lambda^{(j)}.
\end{aligned}
\]
The defining dual action
\(\langle T_\alpha v,\lambda\rangle
=\langle v,A_\alpha^{-1}\lambda\rangle\) therefore proves the crossed law
\[
 \widetilde d_j(\alpha\beta)
 =\widetilde d_j(\alpha)+T_\alpha\widetilde d_j(\beta).
\tag{5.17-A14d}
\]

Put
\[
 \widetilde J_j^*:=\Delta_j^*\times\mathbb C^2,
 \qquad
 r_j^*:=r_j|_{\widetilde J_j^*}:
 \widetilde J_j^*\longrightarrow J_j^{\prime *}
\]
and define
\[
\begin{aligned}
 \Phi_{j,\log}
  &:=p_{j,\log}\circ r_j^*,\\
 \Phi_{j,\mathrm{glob}}
  &:=p_{j,\mathrm{lin}}\circ h_j\circ r_j^*
   =\bar h_j\circ\Phi_{j,\log}.
\end{aligned}
\]
Both maps are invariant under \(\Gamma_j^{\log}\).  This is immediate for
every \(\tau_\lambda\), while
\[
\begin{aligned}
 \Phi_{j,\log}Q_{j,\log}
 &=p_{j,\log}q_{j,\log}r_j^*
  =\Phi_{j,\log},\\
 \Phi_{j,\mathrm{glob}}Q_{j,\log}
 &=p_{j,\mathrm{lin}}h_jq_{j,\log}r_j^*
  =p_{j,\mathrm{lin}}q_{j,\mathrm{lin}}h_jr_j^*
  =\Phi_{j,\mathrm{glob}}.
\end{aligned}
\]
Thus \(\widetilde J_j^*\) is the same \(\Gamma_j^{\log}\)-cover of the two
biholomorphic quotients.  Pulling the gluing isomorphism to this cover gives
\[
 \Theta_j:
 \Phi_{j,\log}^*(L|_{N_j^*})
 \xrightarrow{\ \sim\ }
 \Phi_{j,\mathrm{glob}}^*(L|_{J_j^*}).
\]
The filled factor extension restricts to this cover, and the global factor
extension pulls back to it through the fixed boundary whisker.  Because
\(\Theta_j\) intertwines descent, it induces an isomorphism of the resulting
pushed-out \(\Gamma_j^{\log}\)-factor extensions which sends
\(\mathbf c(r)\) to \(\mathbf c(r)\) for every \(r\in\mathbb R\).

Let \(\gamma_j^J:[0,1]\to J_j^*\) be the zero-section image of the
unwhiskered clockwise boundary meridian, and let \(\omega_j\) be the fixed
zero-section whisker from the global basepoint to \(\gamma_j^J(0)\).  Thus
the based total-space loop retained in (5.17-A1) is
\[
 \widehat\rho_j=\omega_j*\gamma_j^J*\omega_j^{-1}.
\]
Choose \(x_j\in J_j^{\prime *}\) such that
\(p_{j,\mathrm{lin}}(h_j(x_j))=\gamma_j^J(0)\), put
\(x_{j,\mathrm{lin}}=h_j(x_j)\), and let
\(\widetilde\gamma_{j,\mathrm{lin}}\) be the
\(p_{j,\mathrm{lin}}\)-lift of \(\gamma_j^J\) beginning at
\(x_{j,\mathrm{lin}}\).  For the chosen clockwise generator,
\[
 \widetilde\gamma_{j,\mathrm{lin}}(1)
 =q_{j,\mathrm{lin}}(x_{j,\mathrm{lin}}).
\]
Set
\[
 \gamma_j^N:=\bar h_j^{-1}\circ\gamma_j^J,
 \qquad
 \widetilde\gamma_{j,\log}
 :=h_j^{-1}\circ\widetilde\gamma_{j,\mathrm{lin}}.
\]
Then the quotient square gives the point-valued identities
\[
\begin{aligned}
 p_{j,\log}\circ\widetilde\gamma_{j,\log}
  &=\gamma_j^N,\\
 \widetilde\gamma_{j,\log}(0)&=x_j,\\
 \widetilde\gamma_{j,\log}(1)&=q_{j,\log}(x_j).
\end{aligned}
\]
Choose \(\widetilde x_j\in\widetilde J_j^*\) with
\(r_j^*(\widetilde x_j)=x_j\), and let
\(\widetilde\gamma_{j,\mathrm{vec}}\) be the \(r_j^*\)-lift of
\(\widetilde\gamma_{j,\log}\) beginning at \(\widetilde x_j\).  On
adjacent sheets choose compatible logarithm branches \(\ell,\ell'\) with
\[
 \ell'(w)=\ell(g_j^{-1}w)-\frac{2\pi i}{m_j}.
\]
The calculation in (2.9e) gives the point-valued endpoint formula
\[
 \widetilde\gamma_{j,\mathrm{vec}}(1)
 =Q_{j,\log}(\widetilde x_j),
\]
excluding both \(Q_{j,\log}^{-1}(\widetilde x_j)\) and every
\(\tau_\lambda Q_{j,\log}(\widetilde x_j)\) with \(\lambda\ne0\).
Fibre loops remain \(\tau_\lambda\), since \(h_j\) is a fibre translation.

The factor of \(L\) along the fixed lift of \(\omega_j\) identifies the
target factor over the local boundary segment with the globally based
factor.  Let \(G_{\rho_j}^{(j)}\) be the target lift of the local boundary
deck element whose whiskered image is the chosen global lift
\(G_{\rho_j}\).  Thus it covers \(\widehat\rho_j\) after whiskering and,
by (5.17-A13), satisfies
\[
 G_{\rho_j}^{(j)}U_{A_j^{-1}\lambda}(G_{\rho_j}^{(j)})^{-1}
 =\mathbf c(\langle d_j,\lambda\rangle)U_\lambda,
 \qquad d_j=d(\rho_j).
\]
Consequently, for some
\(a_j\in V_{\mathbb R}\) and \(b_j\in\mathbb R\),
\[
 \Theta_{j,*}(U_\lambda^{(j)})
 =\mathbf c(\langle a_j,\lambda\rangle)U_\lambda,
 \qquad
 \Theta_{j,*}(\widetilde G_{Q_{j,\log}})
 =\mathbf c(b_j)G_{\rho_j}^{(j)}.
\]
Here \(U_\lambda\) is the normalized global system in (5.17-A12), and
\(\mathbf c(b_j)\) disappears under conjugation.
Therefore
\[
\begin{aligned}
 \langle\widetilde d_j(Q_{j,\log}),\lambda\rangle
 +\langle a_j,\lambda\rangle
 &=\langle d_j,\lambda\rangle
   +\langle a_j,A_j^{-1}\lambda\rangle\\
 &=\langle d_j+T_ja_j,\lambda\rangle,
\end{aligned}
\]
so
\[
 \widetilde d_j(Q_{j,\log})=d_j+(T_j-I)a_j.
\tag{5.17-A14e}
\]

For a translation choose
\(\widetilde G_{\tau_v}=U_v^{(j)}\).  The filled normalization above gives
\[
 U_v^{(j)}U_\lambda^{(j)}(U_v^{(j)})^{-1}
 =\mathbf c(E(v,\lambda))U_\lambda^{(j)},
\]
and hence
\[
 \widetilde d_j(\tau_v)=\iota_vE,\qquad
 (\iota_vE)(\lambda)=E(v,\lambda).
\tag{5.17-A14f}
\]
Here
\[
 A_{Q_{j,\log}}=A_j,\qquad T_{Q_{j,\log}}=T_j.
\]
Since \(N_{T_j}(T_j-I)=T_j^{m_j}-I=0\), equations
(5.17-A14), (5.17-A14d), and (5.17-A14e) give
\[
\begin{aligned}
 N_{T_j}d_j
 &=N_{T_j}\widetilde d_j(Q_{j,\log})\\
 &=\widetilde d_j((Q_{j,\log})^{m_j})\\
 &=\widetilde d_j(\tau_{v_j})\\
 &=\iota_{v_j}E.
\end{aligned}
\tag{5.17-A15}
\]
Changing a line-bundle frame gives an isomorphic factor extension; changing
\(\widetilde G_\alpha\) by a central factor does not change conjugation; and
every second normalized vertical system differs by a character, changing
\(\widetilde d_j(\alpha)\) by \((T_\alpha-I)a\).  Finally,
\(\ell\mapsto\ell+2\pi in\) changes a vector-cover lift of \(h_j\) by
\[
 H_j\longmapsto H'_j=\tau_{nv_j}H_j.
\]
At the factor-extension level this means
\[
 \Theta'_{j,*}
 =\operatorname{Ad}(U_{nv_j})\circ\Theta_{j,*}.
\]
Consequently the comparison character changes by
\(n\iota_{v_j}E\), because
\[
 U_{nv_j}
 \bigl(\mathbf c(\langle a_j,\lambda\rangle)U_\lambda\bigr)
 U_{nv_j}^{-1}
 =
 \mathbf c(\langle a_j+n\iota_{v_j}E,\lambda\rangle)U_\lambda.
\]
Since \(A_jv_j=v_j\) and \(E\) is \(A_j\)-invariant,
\[
 \langle T_j(\iota_{v_j}E),\lambda\rangle
 =E(v_j,A_j^{-1}\lambda)
 =E(A_jv_j,\lambda)
 =E(v_j,\lambda),
\]
so this branch change leaves both (5.17-A14e) and (5.17-A15) unchanged.
The translation \(\tau_{nv_j}\) also commutes with
\(Q_{j,\mathrm{lin}}\) on the vector cover because \(A_jv_j=v_j\); on the
lattice quotient it is already the identity.  Hence it changes only the
central lift choice there.
The descended \(h_j,\bar h_j\), and the quotient square are branch
independent.  Thus (5.17-A15) is choice-independent and proves (5.17-A5),
including its sign.

It remains to obtain (5.17-A6) analytically.  Since
\(A_0\widehat\delta=\widehat\delta\), equation (5.17-A13) says that
\[
 [G_0,U_{\widehat\delta}]
 =\mathbf c\bigl(\langle d_0,\widehat\delta\rangle\bigr).
\tag{5.17-A16}
\]
This logarithmic commutator is the clutching winding with the base loop
first.  Thus, if \(b\) and \(f\) denote the oriented base and fibre loops,
respectively,
\[
 \operatorname{wind}_{b\wedge f}(L)
  =\langle d_0,\widehat\delta\rangle,
 \qquad
 \operatorname{wind}_{\xi_{\widehat\delta}}(L)
 =-\langle d_0,\widehat\delta\rangle.
\tag{5.17-A17}
\]
The second equality uses the fibre-first convention
\(\xi_{\widehat\delta}=[S^1_{\widehat\delta}]\times[\rho_0]\), hence the
single sign change.
Use the actual unimodular cusp chart
\[
 (u,x_1,x_2)=\left(\frac{t}{x_1x_2},x_1,x_2\right),
 \qquad t=ux_1x_2.
\tag{5.17-A18}
\]
For \(0<\rho\ll1\), the boundary parametrization
\[
 (\theta,\phi)\longmapsto
 \bigl(\rho e^{i(\theta-\phi)},1,e^{i\phi}\bigr),
 \qquad t=\rho e^{i\theta},
\tag{5.17-A19}
\]
has base coordinate \(\theta\) and fibre translation
\(\widehat\delta\) in the \(\phi\)-direction.  The latter fixed-base loop
does \emph{not} itself bound the \(x_2\)-disc.  To identify the correct
meridian, put \(\alpha=\theta-\phi\), let \(a_u\) be the positively
oriented \(u\)-longitude, and let \(m_{x_2}\) be the positively oriented
\(x_2\)-meridian at fixed \(u\).  In the basis
\((a_u,m_{x_2})\), the base and fibre loops are
\[
 b=a_u,
 \qquad
 f=m_{x_2}-a_u,
 \qquad
 b\wedge f=a_u\wedge m_{x_2}.
\tag{5.17-A19a}
\]
The change-of-basis determinant is \(+1\), so this preserves the
base-first orientation in (5.17-A17).  Now \(m_{x_2}\) does bound the
analytic disc
\[
 z\longmapsto(\rho e^{i\alpha_0},1,z),
 \qquad |z|\le1.
\tag{5.17-A19b}
\]

Choose a thin annulus \(A_u\) about \(|u|=\rho\), a disc \(D_{x_1}\)
about \(1\), and a disc \(D_{x_2}\) containing \(\{|x_2|\le1\}\), small
enough that
\[
 A_u\times D_{x_1}\times D_{x_2}
 \subset\{|ux_1x_2|<\epsilon\}.
\tag{5.17-A19c}
\]
This neighbourhood is Stein and retracts onto \(S^1\), so its second
integral cohomology vanishes.  Cartan's theorem~B and the exponential
sequence therefore give trivial Picard group
\cite[\S2, fundamental theorem, p.~81]{Cartan1957}%
\cite[Proposition~6.28 and equations~(6.35)--(6.36),
pp.~180--182]{Lee2024}.  Pull \(L|_{N_0}\) back to
the cusp cover and choose one holomorphic frame on this entire
neighbourhood.  Compute the clutching cocycle with this extending frame on
both sides of a cut in the \(u\)-longitude.  Its transition unit on
\(D_{x_2}\) is then identically \(F=1\), so by (5.17-A19a) its
base--fibre commutator winding is
\[
 \frac1{2\pi i}\int_{|x_2|=1}\frac{dF}{F}=0.
\tag{5.17-A20}
\]
This computes the same alternating logarithmic cocycle as (5.17-A16).
Indeed, changing from the global Appell--Humbert frame to the extending
chart frame changes \(c\) by a group coboundary.  Put
\[
 g=\widehat g_0,
 \qquad
 h=\tau_{\widehat\delta}
\]
in the underlying deck group.  Equation (5.17-A9) and
\(A_0\widehat\delta=\widehat\delta\) give \(gh=hg\).  If
\(c'=c+\delta\beta\) is the cocycle after a frame change, then
\[
 \bigl[c'(g,h)-c'(h,g)\bigr]
 -\bigl[c(g,h)-c(h,g)\bigr]
 =\beta(hg)-\beta(gh)=0.
\]
Thus the alternating cocycle is frame-invariant without assuming that its
central-extension lifts commute.  Equations (5.17-A17), (5.17-A19a), and
(5.17-A20) therefore give
\(\langle d_0,\widehat\delta\rangle=0\).  This is a calculation in
\(\mathcal O^*\), not an evaluation of a Chern class on a bounding chain.
Finally, the matrices \(A_1,A_2\) fix
\(\widehat\delta\), so the evaluation is unchanged by every principal
change in (5.17-A3).  Since
\(V_G=\mathbb Z[\delta]\) and
\(\delta(\widehat\delta)=1\), its vanishing is equivalent to
\(\bar d_0=0\).  This proves (5.17-A6).
\end{proof}

\begin{lemma}[the toric support function detects central restriction]
\label{lem:analytic-support-specialization}
\Verified\ Let \(M\in\operatorname{Pic}(N_0)\), let
\(\pi:Y_\epsilon\to N_0\) be the cusp covering, and let
\(\nu:E_0\to W\) be the normalization.  If the Appell--Humbert form of
\(M|_{F_t}\), \(t\ne0\), is \(E_M\), then the three independent degrees
of \(\nu^*(M|_W)\) determine an integral symmetric form
\(B:\mathbb Z^2\times\mathbb Z^2\to\mathbb Z\) satisfying
\[
 E_M\bigl((\lambda,0),(0,m)\bigr)
 =-B(B_0\lambda,m)
 \qquad(\lambda,m\in\mathbb Z^2).
\tag{5.17-A21}
\]
Consequently
\[
 E_M=0
 \quad\Longrightarrow\quad
 \nu^*(M|_W)\simeq\mathcal O_{E_0}.
\tag{5.17-A22}
\]
\end{lemma}

\begin{proof}
Put \(\mathcal M=\pi^*M\).  The identity
\(\pi\circ\Psi_\lambda=\pi\) gives \(\mathcal M\) its genuine
\(\Psi\)-equivariant descent structure.  In the notation of
\cref{prop:local-cusp-picard-obstruction}, order the six rays in the star
of the origin as
\[
 (r_1,r_2,r_3,r_4,r_5,r_6)
 =(e_1,e_2,e_2-e_1,-e_1,-e_2,e_1-e_2),
 \qquad C_i=E_0\cap E_{r_i}.
\tag{5.17-A23}
\]
The quotient identifies opposite conductor branches.  Therefore, for
\(A=\mathcal M|_{E_0}=\nu^*(M|_W)\),
\[
 d_i=\deg_{C_i}A,
 \qquad
 (\deg_{C_1}A,\ldots,\deg_{C_6}A)
 =(d_1,d_2,d_3,d_1,d_2,d_3).
\tag{5.17-A24}
\]

Define the integer-valued quadratic function
\[
 a(x,y)
 =(d_2+d_3)\binom{x}{2}+d_3xy+(d_1+d_3)\binom{y}{2}.
\tag{5.17-A25}
\]
Its polarization is
\[
 B(v,w)=a(v+w)-a(v)-a(w),
 \qquad
 [B]=
 \begin{pmatrix}
  d_2+d_3&d_3\\ d_3&d_1+d_3
 \end{pmatrix}.
\tag{5.17-A26}
\]
The locally finite divisor
\[
 D_a=\sum_{v\in\mathbb Z^2}a(v)E_v
\tag{5.17-A27}
\]
is Cartier: on every maximal chart, the three height-one ray generators
form an integral basis, so the three assigned coefficients are the values
of one integral character.  In the standard toric convention the associated
support function has ray values \(-a(v)\); the divisor coefficients here are
the values \(a(v)\), consistent with
\(\operatorname{div}(\chi^m)=\sum_\rho\langle m,u_\rho\rangle D_\rho\)
\cite[\S3.3, pp.~60--63; \S3.4, pp.~63--72]{Fulton1993}%
\cite[Propositions~4.1.1--4.1.2, p.~171; Theorems~4.2.8,
p.~179 and~4.2.12, p.~182]{CoxLittleSchenck2011}.
This is the same chartwise Cartier argument as
in (5.17ab), and it applies to an arbitrary integer-valued \(a\).

Direct substitution gives
\[
 (a(r_1),\ldots,a(r_6))
 =(0,0,d_2,d_2+d_3,d_1+d_3,d_1).
\tag{5.17-A28}
\]
Using \(C_i^2=-1\) and \(C_iC_{i\pm1}=1\), the degree vector of
\(\mathcal O(D_a)|_{E_0}\) is exactly the vector in (5.17-A24).  More
generally, restrict the principal identity
\(\operatorname{div}(t)=\sum_wE_w\) to \(E_v\).  Since \(E_v\) meets
precisely \(E_{v+r_i}\), it gives
\[
 \mathcal O(D_a)|_{E_v}
 \simeq
 \mathcal O_{E_v}\!\left(
   \sum_{i=1}^6(a(v+r_i)-a(v))C_{v,i}
 \right).
\tag{5.17-A29}
\]
The coefficient difference is
\[
 a(v+r_i)-a(v)=a(r_i)+B(v,r_i).
\tag{5.17-A30}
\]
The second summand in (5.17-A30) is the divisor of the toric character
with exponent \(B(v,-)\).  Thus every restriction in (5.17-A29) has the
same six degrees as \(\mathcal O(D_a)|_{E_0}\).

Because \(B_0\bar\Lambda=\mathbb Z^2\), the deck action is transitive on
the components \(E_v\).  The genuine descent structure on \(\mathcal M\)
therefore transports (5.17-A24), with its ordered six boundary directions,
to every \(E_v\).  It follows that
\[
 \mathcal R:=\mathcal M\otimes\mathcal O(-D_a)
 \quad\text{satisfies}\quad
 \mathcal R|_{E_v}\simeq\mathcal O_{E_v}
 \quad(v\in\mathbb Z^2).
\tag{5.17-A31}
\]
Indeed, the six toric curves detect \(\operatorname{Pic}(\mathrm{dP}_6)\):
in the blow-up basis \(H,E_1,E_2,E_3\), vanishing of the degrees on the
three exceptional curves kills the three \(E_i\)-coefficients, and
vanishing on any strict transform \(H-E_i-E_j\) kills the remaining
\(H\)-coefficient.  Since
\(H^1(\mathrm{dP}_6,\mathcal O)=0\), equality of those degrees implies
isomorphism, not merely numerical equivalence.

For \(b_\lambda=B_0\lambda\), the torus-translation part of
\(\Psi_\lambda\) preserves each toric prime divisor, while its fan part
translates the index by \(b_\lambda\).  Hence
\[
 \Psi_\lambda^*D_a-D_a
 =\sum_v(a(v+b_\lambda)-a(v))E_v
 =\operatorname{div}(\chi^{m_\lambda}),
\tag{5.17-A32}
\]
where
\[
 m_\lambda=(\ell_\lambda,k_\lambda),
 \qquad
 \ell_\lambda=B(b_\lambda,-),
 \qquad
 k_\lambda=a(b_\lambda).
\tag{5.17-A33}
\]
Thus multiplication by \(\chi^{m_\lambda}\) gives an individual
isomorphism
\(\Psi_\lambda^*\mathcal O(D_a)\simeq\mathcal O(D_a)\).
No cocycle assertion for these individual isomorphisms is needed: we use
them only to separate the horizontal winding of the genuine descent factor
of \(\mathcal M\).

For a ray \(\mathbb R_{\ge0}(v,1)\), its toric open set in
\(Y_\epsilon\) is
\[
 U_v\simeq\Delta_\epsilon\times(\mathbb C^*)^2,
 \qquad E_v^\circ=\{0\}\times(\mathbb C^*)^2.
\tag{5.17-A34}
\]
Indeed \((e_1,0),(e_2,0),(v,1)\) is an integral basis and \(t\) is the
character dual to the last vector; compare
\cite[\S2.1, pp.~28--29]{Fulton1993} and
\cite[Example~1.2.21, p.~31]{CoxLittleSchenck2011}.
The set \(U_v\) is Stein and retracts
onto \(E_v^\circ\).  By (5.17-A31),
\(c_1(\mathcal R|_{U_v})=0\); the exponential sequence and Cartan's
theorem~B therefore trivialize \(\mathcal R|_{U_v}\).  Choose a global
nowhere-zero frame \(s_v\) of \(\mathcal R|_{E_v}\).  A frame on \(U_v\) can be
multiplied by a unit pulled back from \((\mathbb C^*)^2\) so that its
restriction to \(E_v^\circ\) agrees with \(s_v|_{E_v^\circ}\); denote the
resulting frame by \(e_v\).

We also need to compare these frames on the generic torus.  The direction
of the toric affine-open inclusions is essential here.  If \(v,w\) are
adjacent vertices, choose a triangle \(\sigma=\{v,w,u\}\) containing them,
write \(z_v,z_w,z_u\) for the coordinates dual to its three height-one ray
generators, and put
\[
 S_\sigma:=U_\sigma\cap Y_\epsilon
 =\{(z_v,z_w,z_u)\in\mathbb C^3:
       |z_vz_wz_u|<\epsilon\}.
\tag{5.17-A34a}
\]
Because \(\rho_v<\sigma\), localization of the maximal affine chart gives
\begin{equation}
 U_{\rho_v}=D(z_wz_u)\subset U_\sigma,\qquad
 U_{\rho_w}=D(z_vz_u)\subset U_\sigma,\qquad
 U_{\rho_v}\cap U_{\rho_w}=D(z_vz_wz_u)=U_0.
\label{eq:analytic-toric-open-directions}
\end{equation}
Thus \(U_\sigma\) is not contained in \(U_{\rho_v}\cap U_{\rho_w}\), and
the coordinate discs in \(S_\sigma\) cannot be applied directly to the
ratio \(e_w/e_v\).

The space \(S_\sigma\) is nevertheless the required comparison chart.  It
is the closed graph
\[
 \{(z_v,z_w,z_u,t)\in\mathbb C^3\times\Delta_\epsilon:
   t=z_vz_wz_u\}
\]
in the Stein space \(\mathbb C^3\times\Delta_\epsilon\), and hence is
Stein.  Radial scaling of the three \(z\)-coordinates contracts it to the
origin.  The exponential sequence and Cartan's theorem~B therefore give
\(\operatorname{Pic}(S_\sigma)=0\).  Choose a frame \(e_\sigma\) of
\(\mathcal R|_{S_\sigma}\).

Let \(s_v\) be the global frame of
\(\mathcal R|_{E_v}\simeq\mathcal O_{E_v}\) used to normalize \(e_v\).
On \(U_{\rho_v}\cap Y_\epsilon\) put
\[
 A_{\sigma,v}=e_v/e_\sigma.
\]
At \(t=0\), this is the restriction to \((\mathbb C^*)^2\) of the
holomorphic unit
\[
 G_{\sigma,v}
 =\frac{s_v}{e_\sigma|_{E_v\cap S_\sigma}}
 \in\mathcal O^*(E_v\cap S_\sigma),
 \qquad E_v\cap S_\sigma=\{z_v=0\}\simeq\mathbb C^2.
\]
Consequently its winding on every horizontal cocharacter is zero at
\(t=0\).  Under
\(U_{\rho_v}\cap Y_\epsilon\simeq
\Delta_\epsilon\times(\mathbb C^*)^2\), that winding is an integer locally
constant in \(t\), and is therefore zero on every nearby nonzero fibre.
The same argument applies to \(A_{\sigma,w}\).  On their actual common
domain \(U_0\cap Y_\epsilon\),
\[
 e_w/e_v=A_{\sigma,w}/A_{\sigma,v},
\]
so this ratio has zero horizontal winding.  The one-skeleton of the
\(A_2\)-triangulation is connected, and multiplication along an edge path
now gives
\[
 \operatorname{wind}_m(e_w/e_v)=0
 \qquad(v,w,m\in\mathbb Z^2)
\tag{5.17-A34b}
\]
on every nonzero fibre.

Let
\(\theta^{\mathcal M}_\lambda:\Psi_\lambda^*\mathcal M\to\mathcal M\)
be the genuine descent isomorphism, let
\(\mu^+_\lambda:\Psi_\lambda^*\mathcal O(D_a)\to\mathcal O(D_a)\)
be multiplication by \(\chi^{m_\lambda}\) from (5.17-A32), and define the
inverse-dual isomorphism
\[
 \mu^-_\lambda:=((\mu^+_\lambda)^{-1})^\vee:
 \Psi_\lambda^*\mathcal O(-D_a)\longrightarrow\mathcal O(-D_a),
 \qquad s\longmapsto\chi^{-m_\lambda}s.
\]
The ordinary dual of \(\mu^+_\lambda\) has the opposite arrow; the
inverse-dual above is the map required by the pullback tensor product.  Then
\[
 \theta^{\mathcal R}_\lambda
 :=\theta^{\mathcal M}_\lambda\otimes\mu^-_\lambda:
 \Psi_\lambda^*\mathcal R\longrightarrow\mathcal R.
\]
Since \(\Psi_\lambda(U_v)=U_{v+b_\lambda}\), define
\(r_{\lambda,v}\in\mathcal O^*(U_v)\) by the typed identity
\begin{equation}
 \theta^{\mathcal R}_\lambda
   (\Psi_\lambda^*e_{v+b_\lambda})
 =r_{\lambda,v}e_v.
\label{eq:analytic-residual-descent-factor}
\end{equation}
At \(t=0\) it extends to an isomorphism between two trivial line bundles
on the compact surface \(E_v\).  Its factor in the chosen global frames is
therefore a nonzero constant.  The winding of
\(r_{\lambda,v}(t,-)\) around every horizontal one-parameter subgroup is
an integer locally constant in \(t\); at \(t=0\) it is zero, hence it is
zero for every nearby \(t\ne0\).

To express this isomorphism in one frame on the generic torus, put
\(q_{\lambda,v}=e_{v+b_\lambda}/e_v\).  Equation (5.17-A34b) gives zero
horizontal winding for \(q_{\lambda,v}\); because
\(\Psi_\lambda|_{Y_t}\) is multiplication by a fixed torus point, it
induces the identity on horizontal \(H_1\), and
\(\Psi_\lambda^*q_{\lambda,v}\) also has zero winding.  Relative to the
single frame \(e_v\), the residual multiplier is therefore
\[
 h_{\lambda,v}
 =\frac{r_{\lambda,v}}{\Psi_\lambda^*q_{\lambda,v}},
 \qquad
 \operatorname{wind}_m(h_{\lambda,v})=0.
\tag{5.17-A34c}
\]

On \(U_0\) let \(\mathbf1_a\) be the canonical toric frame of
\(\mathcal O(D_a)\), and put \(F_v=e_v\otimes\mathbf1_a\).  Equations
\cref{eq:analytic-residual-descent-factor} and (5.17-A32) give
\[
 \theta^{\mathcal M}_\lambda(\Psi_\lambda^*F_v)
 =h_{\lambda,v}\chi^{m_\lambda}F_v.
\]
On a nonzero fibre, \(\chi^{m_\lambda}=x^{\ell_\lambda}t^{k_\lambda}\).
With the section-descent convention (5.17am), the factor of automorphy is
therefore
\[
 \alpha_\lambda=(h_{\lambda,v}\chi^{m_\lambda})^{-1}.
\]
The residual factor has zero horizontal winding by (5.17-A34c), and hence
\begin{equation}
 \operatorname{wind}_m(\alpha_\lambda)
 =-\ell_\lambda(m)=-B(B_0\lambda,m).
\label{eq:analytic-factor-horizontal-winding}
\end{equation}

This winding is exactly the mixed Chern-class entry.  Pull \(F_v\) back to
the universal cover
\[
 p_t:\mathbb C^2\longrightarrow Y_t,
 \qquad \zeta\longmapsto(e^{2\pi i\zeta},t).
\]
The period indexed by \((\lambda,n)\in\mathbb Z^2\oplus\mathbb Z^2\)
acts by \(\zeta\mapsto\zeta+Z(s)\lambda+n\).  The horizontal factor is one,
because \(F_v\) is already defined on \(Y_t\).  Choose a logarithm
\(g_\lambda\in\mathcal O(\mathbb C^2)\) with
\[
 p_t^*\alpha_\lambda=e^{2\pi i g_\lambda}.
\]
For every horizontal period \(n\),
\[
 g_\lambda(\zeta+n)-g_\lambda(\zeta)
 =\operatorname{wind}_n(\alpha_\lambda).
\]
In the exponential connecting cocycle
\[
 c(l,l')=g_l(\zeta+l')+g_{l'}(\zeta)-g_{l+l'}(\zeta),
\]
take the horizontal logarithm to be zero and the logarithm for the combined
period \((\lambda,n)\) to be \(g_\lambda\).  Then
\[
 c((\lambda,0),(0,n))
 =\operatorname{wind}_n(\alpha_\lambda),
 \qquad
 c((0,n),(\lambda,0))=0.
\]
Therefore the alternating exponential connecting integer between
the vertical deck
translation indexed by \(\lambda\) and the horizontal loop indexed by
\(m\) is
\[
 E_M\bigl((\lambda,0),(0,m)\bigr)
 =-\ell_\lambda(m)
 =-B(B_0\lambda,m),
\]
which proves (5.17-A21).  This sign also recovers the explicit check in
\cref{prop:local-cusp-picard-obstruction}: for
\(a(v)=v^{\mathsf t}Hv\), one has \(B=2H\), so (5.17-A21) is the
cross-block \(-2B_0^{\mathsf t}H=-2P\), namely \(-2q\).

If \(E_M=0\), equation (5.17-A21) and the unimodularity of \(B_0\) force
\(B=0\).  The matrix in (5.17-A26) then gives
\(d_3=0\), \(d_2=0\), and \(d_1=0\).  Thus all six degrees of \(A\)
vanish, and the Picard detection argument following (5.17-A31) gives
\(A\simeq\mathcal O_{E_0}\).  This is (5.17-A22).
\end{proof}

\begin{theorem}[purely analytic vanishing of the global restriction]
\label{thm:analytic-all-picard}
\Verified\ For the constructed threefold,
\[
 r=\eta^*i^*:\operatorname{Pic}(X)\longrightarrow
 \operatorname{Pic}(\mathrm{dP}_6)
\]
is the zero homomorphism.  This proof uses only holomorphic
\(\mathcal O^*\)-factors, the two finite logarithmic fillings, the cusp
argument principle, and the \(\Psi\)-equivariant toric support function.
It uses neither a global integral Leray transgression nor local integral
nearby-cycle specialization.
\end{theorem}

\begin{proof}
Let \(L\in\operatorname{Pic}(X)\).  By
\cref{thm:coinvariants}, its invariant fibre class is
\[
 E=c_1(L|_F)=kq,
 \qquad k\in\mathbb Z,
 \qquad q=uw+6\gamma\delta.
\tag{5.17-A35}
\]
The exact contractions, now used without an integral specialization
lattice, are
\[
 \iota_{v_1}q=2\eta_1,
 \qquad
 \iota_{v_2}q=-3\eta_2,
 \qquad
 [\eta_1]=3[\delta],
 \qquad
 [\eta_2]=2[\delta].
\tag{5.17-A36}
\]
Apply (5.17-A5) with \((m_1,m_2)=(3,4)\).  In the real global
coinvariants,
\[
 \bar d_1=2k[\delta],
 \qquad
 \bar d_2=-\frac32k[\delta].
\tag{5.17-A37}
\]
Equations (5.17-A4) and (5.17-A6) therefore give
\[
 0=\bar d_1+\bar d_2+\bar d_0
   =\frac{k}{2}[\delta].
\tag{5.17-A38}
\]
Since \(V_G\otimes\mathbb R=\mathbb R[\delta]\), one has \(k=0\), so
\(E=0\).  Apply \cref{lem:analytic-support-specialization} to
\(M=L|_{N_0}\).  It gives
\[
 \eta^*i^*L=\nu^*(L|_W)\simeq\mathcal O_{E_0}.
\]
This holds for every \(L\), proving \(r=0\).
\end{proof}


\begin{corollary}[candidate nonvanishing for \texorpdfstring{$\operatorname{Pic}^0$}{Pic0} without topology]
\label{cor:candidate-pic0-route}
\Verified\ For the constructed threefold and every
$L\in\operatorname{Pic}^0(X)$,
\[
 H^0(W,\Omega_X^1|_W\otimes A_L)\ne0,
 \qquad
 H^2(W,(T_X\otimes L)|_W)\ne0,
\]
and
\[
 R^2f_*(T_X\otimes L)\ne0.
\tag{5.17f}
\]
This proof uses the analytic construction of the cusp fibre but no global
homology, Mayer--Vietoris, or Leray calculation.
\end{corollary}

\begin{proof}
The cusp theorem \cref{thm:candidate-cusp-geometry} supplies the full compact
reduced normal-crossings fibre, its $\mathrm{dP}_6$ normalization, and the
trivialization of $\omega_X$ near $W$.  By
\cref{prop:pic0-pullback}, $\eta^*A_L$ is trivial, so
\cref{prop:conductor} gives the first nonvanishing.  Apply
\cref{prop:cartier-duality-absolute} to $E=T_X\otimes L$; its dual space is
exactly the preceding nonzero section space, giving the fibrewise $H^2$
nonvanishing.  The fibre has complex dimension two, so its coherent
$H^3$ vanishes.  The global construction makes $X$ compact and $f$
nonconstant; \cref{lem:curve-flatness} therefore gives properness and
$f$-flatness of $T_X\otimes L$.  The base-change isomorphism of
\cref{prop:analytic-base-change} now identifies the fibre of the $R^2$
stalk with the nonzero $H^2$ group.  Hence the stalk, and therefore the
direct-image sheaf, is nonzero.
\end{proof}

\begin{remark}[reconciliation of the three restriction routes]
The Pic$^0$ argument proves (5.17f) without global homology,
Mayer--Vietoris, or Leray.  The all-Pic statement
\cref{prop:leray-edge-restriction} avoids citing the packaged theorem
$H^2(X;\mathbb Z)=0$, but it reuses the same verified global integral
Leray and specialization machinery from which that theorem is deduced.
By contrast, \cref{prop:direct-boundary-picard} supplies a second all-Pic
proof without any global Leray differential: it recovers the primitive
$4-3=1$ obstruction directly from the finite logarithmic gluing and the
cusp boundary.  It still uses the local integral cusp-specialization
theorem.  The third proof in \cref{thm:analytic-all-picard} replaces both
that specialization and the cusp Chern-class bounding-chain evaluation by
logarithmic $\mathcal O^*$ factors, the argument principle, and the exact
$\Psi$-equivariant support function.  Thus the purely analytic
strengthening is now complete; the three proofs have genuinely different
dependency sets.
\end{remark}

\subsection{Abstract conductor theorem for the candidate fibre}

\begin{theorem}[the candidate conductor mechanism, with all hypotheses]
\label{thm:candidate-conductor}
\Verified\ Assume:
\begin{enumerate}[label=(H\arabic*)]
\item $X$ is a connected compact smooth complex threefold and
$f:X\to B$ is a nonconstant holomorphic map to a smooth compact curve;
\item $W=f^{-1}(p_0)$ is the full scheme-theoretic fibre, is reduced, and
has the normal-crossings local models (5.3);
\item the normalization $\eta:\widetilde W\to W$ is the degree-six del
Pezzo surface $\mathrm{dP}_6$;
\item $H^2(X;\mathbb Z)=0$.
\end{enumerate}
Then, for every $L\in\operatorname{Pic}(X)$,
\[
H^2(W,(T_X\otimes L)|_W)\ne0
\quad\text{and}\quad
R^2f_*(T_X\otimes L)\ne0.
\tag{5.18}
\]
\end{theorem}

\begin{proof}
Fix $L$ and put
\[
\mathcal F_L=T_X\otimes L,\qquad
A_L=(L^\vee\otimes\omega_X)|_W.
\]
Hypothesis (H4) gives $c_1(L^\vee\otimes\omega_X)=0$.  Pullback therefore
gives $c_1(\eta^*A_L)=0$.  The surface $\mathrm{dP}_6$ is the blow-up of
$\mathbb P^2$ at three noncollinear points; the standard blow-up calculation
gives $H^1(\mathrm{dP}_6,\mathcal O)=0$
\cite[Chapter~V]{Hartshorne1977}.  In the exponential sequence, the kernel
of
\[
c_1:\operatorname{Pic}(\mathrm{dP}_6)\longrightarrow
H^2(\mathrm{dP}_6;\mathbb Z)
\]
is the image of $H^1(\mathrm{dP}_6,\mathcal O)$ and is therefore zero.
Consequently
\[
\eta^*A_L\simeq\mathcal O_{\widetilde W}.
\tag{5.19}
\]
\Cref{prop:conductor} now gives
\[
H^0(W,\Omega^1_X|_W\otimes A_L)\ne0.
\tag{5.20}
\]

Because $W=f^*(p_0)$, one has
$\mathcal O_X(W)=f^*\mathcal O_B(p_0)$, whose restriction to the fibre is
trivial.  Adjunction therefore gives $\omega_W\simeq\omega_X|_W$.
Apply \cref{thm:cartier-duality} to $E=\mathcal F_L$:
\[
\begin{split}
H^2(W,\mathcal F_L|_W)^\vee
&\simeq H^0(W,(\mathcal F_L|_W)^\vee\otimes\omega_W)\\
&=H^0(W,\Omega^1_X|_W\otimes A_L).
\end{split}
\]
The right side is nonzero by (5.20), proving the first assertion of (5.18).
The same duality theorem gives
$H^3(W,\mathcal F_L|_W)=0$.

Since $X$ is compact, $f$ is proper.  By
\cref{lem:curve-flatness}, $f$ is flat and the vector bundle
$\mathcal F_L$ is $f$-flat.  Thus every hypothesis of
\cref{prop:analytic-base-change} is satisfied at $p_0$, and
\[
(R^2f_*\mathcal F_L)_{p_0}\otimes\mathbb C(p_0)
\simeq H^2(W,\mathcal F_L|_W)\ne0.
\]
The stalk, and hence the coherent direct-image sheaf, is nonzero.  This is
the second assertion of (5.18).
\end{proof}

\begin{corollary}[application to the constructed threefold]
\label{cor:candidate-conductor-actual}
\Verified\ For the global construction of
\cref{thm:candidate-global-geometry} and every
$L\in\operatorname{Pic}(X)$,
\[
H^2(W,(T_X\otimes L)|_W)\ne0,
\qquad R^2f_*(T_X\otimes L)\ne0.
\]
\end{corollary}

\begin{proof}
The cusp construction gives the compact full Cartier fibre, its
$\mathrm{dP}_6$ normalization, and the local trivialization of $\omega_X$.
By the local canonical trivialization and
\cref{prop:direct-boundary-picard}, $\eta^*A_L$ is trivial for every $L$.
The conductor proposition then gives
$H^0(W,\Omega_X^1|_W\otimes A_L)\ne0$; absolute Cartier-fibre duality turns
this into the displayed $H^2$ nonvanishing, and the uniform top-degree
base-change isomorphism turns it into a nonzero $R^2$ stalk.  This is exactly
the chain proved in \cref{cor:candidate-pic0-route}, now with the all-Pic
restriction supplied directly by \cref{prop:direct-boundary-picard}.  The
global Leray-edge proof in \cref{prop:leray-edge-restriction} is an
independent certificate of the same restriction statement.
\end{proof}

\subsection{A second printed issue: abelianization of split extensions}

\begin{lemma}[abelianization of a semidirect product]
\label{lem:semidirect-abelianization}
\Verified\ If
$1\to K\to G\to Q\to1$ is split, then
\[
G^{\mathrm{ab}}\simeq
Q^{\mathrm{ab}}\oplus(K^{\mathrm{ab}})_Q,
\tag{5.21}
\]
where the subscript denotes coinvariants for the conjugation action of $Q$.
\end{lemma}

\begin{proof}
Choose a splitting and write $G=K\rtimes Q$.  Let $N\triangleleft K$ be
generated by $[K,K]$ and by all elements
$qkq^{-1}k^{-1}$ with $q\in Q$, $k\in K$.  It is normal in $G$:
$[K,K]$ makes it stable under conjugation by $K$, while conjugation by
$Q$ permutes the displayed cross-commutators.  By construction
$K/N=(K^{\mathrm{ab}})_Q$, and $Q$ acts trivially on this quotient.  Hence
\[
G/N\simeq(K^{\mathrm{ab}})_Q\times Q.
\]
Since $N\subseteq[G,G]$, abelianizing the last display gives (5.21).
\end{proof}

The sentence in \cite[Lemma~3.2, pp.~84--85]{CDP1998} that passes from a
split fundamental-group extension to a direct sum with $H_1(F;\mathbb Z)$
is therefore valid only when monodromy acts trivially after abelianization.
A minimal counterexample retaining a free quotient group is
\[
G=\mathbb Z\rtimes_{-1}\mathbb Z.
\]
The extension splits, but
$G^{\mathrm{ab}}\simeq\mathbb Z\oplus\mathbb Z/2$, whereas the printed
direct-product formula gives $\mathbb Z^2$.
For the candidate's printed monodromy, \cref{thm:coinvariants} gives
\[
H^1(F;\mathbb Q)^{\pi_1(B^\circ)}=\mathbb Q\gamma,
\]
so the rational coinvariant rank is one, not
$b_1(F)=4$.

If $C^\circ$ is a sphere punctured at $s$ points, the correct rank formula is
\[
b_1(X^\circ)=s-1+
\dim_{\mathbb Q}H_1(F;\mathbb Q)_{\pi_1(C^\circ)}.
\tag{5.22}
\]
Thus the proof of CDP98 Theorem~3.1 reaches only the coinvariant rank unless
an additional theorem proves trivial monodromy.  CDP20 Lemma~4.2 repeats the
direct-sum formula, and CDP20 Corollary~5.2 imports CDP98 Theorem~3.1.
Repairing the local-freeness issue discussed there does not repair this
distinct monodromy error.  The direct-product step in the torus branch of the
CDP20 proof is therefore refuted, independently of the singular-fibre defect
above.

\subsection{The corrected theorem that survives both defects}

The original hypotheses do not force the CDP conclusion: the constructed
threefold above is a counterexample.  There is nevertheless a sharp useful
replacement.  It isolates the exact extra global vanishing needed by the
intact part of the argument, and the following fibrewise criterion replaces
the invalid p.~692 reduction.

\begin{lemma}[generic endpoint vanishing]\label{lem:cdp-generic-endpoints}
\Verified\ Let \(Y\) be a connected compact complex manifold, let \(D\) be
a nonzero effective Cartier divisor, and suppose
\(\mathcal O_Y(D)\in\operatorname{Pic}^0(Y)\).  For every holomorphic vector
bundle \(B\), a general \(P\in\operatorname{Pic}^0(Y)\) satisfies
\[
 H^0(Y,B\otimes P)=0,\qquad
 H^{\dim Y}(Y,B\otimes P)=0.
\tag{5.23}
\]
\end{lemma}

\begin{proof}
The inclusions \(\mathcal O_Y(-(k+1)D)\hookrightarrow
\mathcal O_Y(-kD)\) give a descending chain
\[
 V_k=H^0(Y,B\otimes\mathcal O_Y(-kD))
 \subseteq H^0(Y,B).
\]
It stabilizes because the last vector space is finite-dimensional.  A
section in the stable value belongs to every \(V_k\).  In a local equation
\(z=0\) for \(D\), every coefficient of that section is divisible by
\(z^k\) for all \(k\), and hence is zero.  The section vanishes on a
neighbourhood of \(D\), and therefore on connected \(Y\) by analytic
continuation.  Thus \(V_k=0\) for some \(k\).  Since
\(\mathcal O_Y(-kD)\in\operatorname{Pic}^0(Y)\), upper semicontinuity in a
local universal family of topologically trivial line bundles shows that
\[
 \{P:h^0(Y,B\otimes P)>0\}
\]
is a proper analytic subset of \(\operatorname{Pic}^0(Y)\).  This proves
the first vanishing for general \(P\).  Serre duality identifies the dual
of the second group in (5.23) with
\[
 H^0(Y,B^\vee\otimes\omega_Y\otimes P^\vee).
\]
Apply the first part to \(B^\vee\otimes\omega_Y\) and use that inversion is
an automorphism of \(\operatorname{Pic}^0(Y)\).
\end{proof}

\begin{theorem}[corrected CDP global-vanishing criterion]
\label{thm:corrected-cdp-global}
\Verified\ Let \(X\) be a compact connected complex threefold such that
\[
 H^1(X;\mathbb Z)=H^2(X;\mathbb Z)=0,\qquad a(X)=1,
\]
and suppose its algebraic reduction is a holomorphic map
\(f:X\to C\).  If there exists \(L_0\in\operatorname{Pic}(X)\) with
\[
 H^2(X,T_X\otimes L_0)=0,
\tag{5.24}
\]
then, for \(M\) in a dense open subset of \(\operatorname{Pic}(X)\),
\[
 H^i(X,T_X\otimes M)=0\quad(i\ne1),
\tag{5.25}
\]
\[
 \chi(X,T_X\otimes M)\le0,\qquad
 \left\langle c_3(T_X),[X]\right\rangle\le0.
\tag{5.26}
\]
\end{theorem}

\begin{proof}
The exponential sequence and the two integral vanishings give
\[
 \operatorname{Pic}(X)=\operatorname{Pic}^0(X)
 \simeq H^1(X,\mathcal O_X).
\tag{5.27}
\]
Thus this Picard group is a connected irreducible complex vector group.  Any fibre divisor \(D\) of
\(f\) is nonzero and effective, and its line bundle lies in
\(\operatorname{Pic}^0(X)\) by \(H^2(X;\mathbb Z)=0\).  Applying
\cref{lem:cdp-generic-endpoints} to \(B=T_X\) shows that the degree-zero and
degree-three groups in (5.25) vanish away from two proper analytic subsets
of \(\operatorname{Pic}(X)\).

The cohomology-jump locus
\[
 \{M\in\operatorname{Pic}(X):h^2(X,T_X\otimes M)\ge1\}
\]
is analytic by upper semicontinuity.  It is proper because (5.24) gives a
point outside it.  Its complement therefore meets the two endpoint
vanishing loci in a dense open set.  This proves (5.25), and on that set
\[
 \chi(X,T_X\otimes M)=-h^1(X,T_X\otimes M)\le0.
\tag{5.28}
\]

Every \(M\in\operatorname{Pic}(X)\) has rationally trivial first Chern
class.  The hypotheses also give \(c_1(T_X)=0\) in rational cohomology and,
by Poincar\'e duality, \(H^4(X;\mathbb Q)=0\), so
\(c_2(T_X)=0\) rationally.  The degree-six term of
\(\operatorname{ch}(T_X\otimes M)\operatorname{td}(X)\) is consequently
\(\tfrac12c_3(T_X)\).  Hirzebruch--Riemann--Roch turns (5.28) into
\[
 \chi(X,T_X\otimes M)
 =\frac12\left\langle c_3(T_X),[X]\right\rangle\le0,
\]
which proves (5.26).
\end{proof}

\begin{corollary}[intact relative criterion]\label{cor:corrected-cdp-relative}
\Verified\ In the setting of
\cref{thm:corrected-cdp-global}, it is enough that one line bundle
\(L\in\operatorname{Pic}(X)\) satisfy
\[
 R^2f_*(T_X\otimes L)=0.
\tag{5.29}
\]
This is the valid content of the sufficient criterion printed as
\cite[Proposition~2.4, p.~680]{CDP2020}.
\end{corollary}

\begin{proof}
Let \(A\) be a positive-degree line bundle on the compact Riemann surface
\(C\), and put \(E=T_X\otimes L\).  Projection formula gives
\[
 R^2f_*(E\otimes f^*A^m)=R^2f_*E\otimes A^m=0.
\]
For \(m\gg0\), Serre vanishing on \(C\) gives
\[
 H^1(C,R^1f_*E\otimes A^m)=0.
\]
The Leray spectral sequence has only columns zero and one, so its
total-degree-two terms yield
\[
 H^2(X,E\otimes f^*A^m)=0.
\]
Thus \(L_0=L\otimes f^*A^m\) satisfies (5.24), and
\cref{thm:corrected-cdp-global} applies.
\end{proof}

\begin{theorem}[exact reduced-fibre repair of the p.~692 step]
\label{thm:corrected-cdp-fibrewise}
\Verified\ Retain the setting of
\cref{thm:corrected-cdp-global}, assume every full scheme-theoretic fibre
\(W_c=f^{-1}(c)\) is reduced, and fix \(L\in\operatorname{Pic}(X)\).  Put
\[
 A_c=(\omega_X\otimes L^\vee)|_{W_c},\qquad
 \mathscr T_c=\ker\!\left(
 \Omega^1_{W_c}\longrightarrow
 j_*\Omega^1_{(W_c)_{\rm reg}}\right).
\]
For each \(c\), let
\[
 \delta_c:
 H^0(W_c,\mathscr T_c\otimes A_c)
 \longrightarrow
 H^1(W_c,N^*_{W_c/X}\otimes A_c)
\tag{5.30}
\]
be the connecting map from the missing cotangent filtration.  If, for every
\(c\),
\[
 H^0(W_c,\widetilde\Omega^1_{W_c}\otimes A_c)=0,\qquad
 H^0(W_c,N^*_{W_c/X}\otimes A_c)=0,\qquad
 \delta_c\ \text{is injective},
\tag{5.31}
\]
then (5.29), (5.25), and (5.26) hold.
\end{theorem}

\begin{proof}
Apply \cref{prop:missing-global-section-criterion} to the whole reduced
Cartier fibre \(W_c\), not separately to its components.  Conditions (5.31)
give
\[
 H^0(W_c,\Omega_X^1|_{W_c}\otimes A_c)
 =\ker\delta_c=0.
\tag{5.32}
\]
The map \(f\) is flat: a local parameter at \(c\) pulls back to a
non-zero-divisor in every regular local ring of the connected manifold \(X\),
and torsion-freeness over the base DVR is flatness.  Moreover, the normal
line of the full fibre is the pullback of the one-dimensional vector space
\(T_cC\), hence is trivial.  Adjunction therefore identifies
\(\omega_{W_c}\simeq\omega_X|_{W_c}\).  With \(E=T_X\otimes L\), absolute
Cartier-fibre duality now gives
\[
 H^2(W_c,E|_{W_c})^\vee
 \simeq H^0(W_c,\Omega_X^1|_{W_c}\otimes A_c)=0.
\]
The uniform top-degree base-change isomorphism of
\cref{cor:surface-top-base-change} therefore gives
\[
 (R^2f_*E)_c\otimes_{\mathcal O_{C,c}}\mathbb C(c)=0
\]
for every \(c\).  Nakayama's lemma kills every stalk of the coherent sheaf
\(R^2f_*E\), proving (5.29).  The remaining conclusions follow from
\cref{cor:corrected-cdp-relative}.
\end{proof}

\begin{remark}[scope and sharpness of the repair]
Under the two endpoint vanishings in (5.31), injectivity of \(\delta_c\) is
equivalent to the desired ambient vanishing (5.32); it is not a politeness
assumption hiding the original claim.  For a reducible reduced fibre the
map must be computed on the whole Cartier fibre.  A nonreduced fibre requires
additional nilpotent-thickening layers and is not covered by
\cref{thm:corrected-cdp-fibrewise}.  The global criterion
\cref{thm:corrected-cdp-global} and the relative criterion
\cref{cor:corrected-cdp-relative} do allow nonreduced fibres.

The reconstructed candidate violates both added conditions rather than being
silently excluded by terminology.  Since \(c_3(X)=2\), the contrapositive of
\cref{thm:corrected-cdp-global} proves
\[
 H^2(X,T_X\otimes L_0)\ne0
 \quad\text{for every }L_0\in\operatorname{Pic}(X).
\]
Independently, \cref{cor:candidate-conductor-actual} proves
\[
 R^2f_*(T_X\otimes L)\ne0
 \quad\text{for every }L\in\operatorname{Pic}(X),
\]
and its cusp conductor supplies a nonzero element of the relevant
\(\ker\delta_c\).  The first conclusion rules out the global hypothesis; the
second independently rules out the stronger relative sufficient condition.
Thus the corrected theorem is compatible with, and
precisely separates, the verified counterexample to CDP20
Theorem~2.2 as stated.
\end{remark}

\paragraph{The exact monodromy repair.}
The split-extension argument can reach only the invariant integral lattice
\(H^1(F;\mathbb Z)^{\pi_1(C^\circ)}\).  Its final complex-linear inference is
valid precisely when
\[
 \operatorname{Span}_{\mathbb C}\operatorname{im}\!\left(
 H^1(F;\mathbb Z)^{\pi_1(C^\circ)}
 \longrightarrow H^1(F,\mathcal O_F)\right)
 =H^1(F,\mathcal O_F).
\tag{5.33}
\]
Trivial monodromy on \(H_1(F;\mathbb Z)\) is a stronger sufficient
hypothesis.  Neither (5.33) nor trivial monodromy follows from the hypotheses
of CDP20 Theorem~2.2.  The corrected theorem above bypasses this rank route
entirely; it does not relabel the false direct-sum formula as proved.

\subsection{Exact present disposition}

\begin{enumerate}[label=\arabic*.]
\item \Refuted\ The p.~692 reduction in \cite{CDP2020} is not a formal
consequence of the two cotangent sequences.
\Cref{prop:nodal-counterexample} is a proper geometric counterexample to
that exact sheaf-level inference for a non-torsion bundle on the central
fibre; no extension of that bundle to the displayed total space is claimed.
\item \Refuted\ The direct-product homology formula in
\cite[Lemma~3.2]{CDP1998} omits monodromy coinvariants.  This leaves the
torus branch imported by \cite{CDP2020} without its printed rank argument.
\item \Verified\ The normalization--conductor construction,
Cartier-fibre duality, flatness over the base curve, and top-degree analytic
base change are proved above under their stated hypotheses.
\item \Verified\ The period-family descent, toric and finite fillings,
global gluing, and integral analytic nearby-cycle/Leray bridges are proved in
\cref{sec:candidate-geometry,sec:integral-leray}.  Therefore the candidate's
hypotheses (H1)--(H4) hold and
$R^2f_*(T_X\otimes L)\ne0$ for every $L$ by
\cref{cor:candidate-conductor-actual}.
\item \Verified\ Those same construction and topology results, together
with \cref{thm:algebraic-dimension-audit}, hold for the constructed compact
complex threefold.  Hence
$H^1(X;\mathbb Z)=H^2(X;\mathbb Z)=0$, $a(X)=1$, $f$ is a holomorphic
algebraic reduction, and $c_3(X)=2$.  The resulting $X$ satisfies every
stated hypothesis of CDP Theorem~2.2 and contradicts its conclusion; it is
therefore a counterexample to that theorem as stated.  This theorem-level
conclusion uses the complete construction and topology proofs; it is not
inferred from the two printed proof defects alone.
\item \Verified\ The corrected global criterion
\cref{thm:corrected-cdp-global} bypasses both refuted dependencies:
existence of one \(L_0\) with
\(H^2(X,T_X\otimes L_0)=0\) implies generic off-degree-one vanishing and
\(c_3(X)\le0\).  The intact relative sufficient condition is
\(R^2f_*(T_X\otimes L)=0\), and
\cref{thm:corrected-cdp-fibrewise} replaces the invalid reduced-fibre step
by the exact injectivity condition on \(\delta_c\).  The original theorem
without these added hypotheses remains refuted by the constructed
threefold.  Independent specialist replication remains desirable as
external adjudication, but it is neither a mathematical hypothesis nor a
gate on the present proof program.
\end{enumerate}

\begin{test}[decisive theorem-level outcomes]\label{test:cdp-comparison}
The current workbench records two exact defects in the printed CDP chain, a
completed route from the constructed threefold to a theorem-level
counterexample, and the strongest corrected conclusion justified by the
intact argument.  The global vanishing in
\cref{thm:corrected-cdp-global} is false for every \(L_0\) by contraposition
from \(c_3(X)=2\).  Independently, the cusp-conductor calculation forces
\(R^2f_*(T_X\otimes L)\ne0\) for every \(L\), so the stronger relative
sufficient condition fails as well.  Thus the repair does not contradict the
counterexample or recover the original
theorem by renaming its conclusion.  Any further attempt to prove the
original statement must falsify a named step in the candidate reconstruction
or derive one of the new vanishing hypotheses from genuinely stronger
geometry.  Merely invoking the published theorem does neither.
\end{test}

\section{Confirmed local repairs kept separate from the theorem}

These corrections are mathematically definite but do not settle the central construction.

\begin{enumerate}[leftmargin=2em]
\item \Verified\ Proposition~4.7(iv)'s closed-image conclusion is correct,
although its printed justification is not.  If \(s:B\to E\) is the section
of \(f:E\to B\), then
\[
 s(B)=\{x\in E:s(f(x))=x\},
\]
the equalizer of \(\operatorname{id}_E\) and \(s\circ f\).  Since \(E\) is
Hausdorff, its diagonal is closed, so this equalizer is closed.  Thus the
explicit holomorphic section, its cusp disc, and the null-homotopy used in
\cref{lem:first-transgression,prop:geometric-pi1} require no additional
chartwise or local-finiteness hypothesis.
\item \Refuted\ Remark~7.28 is false for a fixed triple. For example,
\[
(\ell_0,\ell_1,\ell_2)=(100,1,-1),\qquad
p=12(100)-4(1)-3(-1)=1199\ne\pm1.
\]
The intended statement must quantify existentially over $\ell_0$ for fixed admissible $(\ell_1,\ell_2)$.
\item \Verified\ Appendix Lemma~A.3 prints cohomological arrows while labelling them as homology. The sequence must be written
\[
H^k(W)\longrightarrow H^k(\widetilde W)\oplus H^k(D)
\longrightarrow H^k(\widetilde D)\longrightarrow H^{k+1}(W),
\]
followed by universal coefficients when homology ranks are required.
\item \Verified\ Appendix Remark~A.6 must read $\pi_1(N_0)=\bar\Lambda=\pi_1(W)$, not $\pi_1(N_0)=\Lambda$.
\item \Verified\ The modular assembler has a page-71 portability defect under the tested MiK\TeX{} route; the canonical monolithic source renders correctly. This is a build defect, not a mathematical counterexample.
\item \Refuted\ The manuscript's assertion that $\Omega^1_{X/B}$ has
torsion at all three singular fibres is false literally.  The explicit
$t=xy$ and $t=xyz$ modules show that it is torsion-free but non-locally-free
at the cusp; torsion occurs only at the two multiple fibres.
\item \Verified\ In the finite-filling Hodge-line computation, the invariant
canonical section
\[
 \sigma_j=s_j^{m_j-1-a_j}e^{\varphi(s_j)}ds_j\wedge e
\]
has divisor \((m_j-1-a_j)S_j\) after descent and therefore is not a
nowhere-vanishing frame when \(j=2\).  The correct direct-image statement is
the invariant-module identity
\[
 f_{j*}\omega_{N_j/D_j}
 =\mathcal O_{D_j}\cdot s_j^{m_j-1-a_j}\vartheta'_j,
 \qquad
 \vartheta'_j=e^{\varphi(s_j)}(ds_j\wedge e)\otimes(dt_j)^{-1}.
\]
On the punctured cover its generator is a unit times \(s_j^{-a_j}e\).
This repairs the local wording while leaving
\(f_*\omega_{X/B}\simeq\mathcal O_{\mathbb P^1}(1)\), the structure-sheaf
direct images, and uniform non-torsion unchanged.
\end{enumerate}

\section{Tool and computational resource disclosure}

The source manuscript available for this audit was reconstructed from a PDF
in ten consecutive source strips---eight strips of eleven pages and two of
ten pages---in ChatGPT Professional-mode sessions and then consolidated into
a single \LaTeX{} source.  Two separate ChatGPT
Professional-mode sessions produced the first full-manuscript proof audits.
The present consolidation, citation audit, exact-arithmetic reproduction, and
validate/repair/disprove workbench were prepared with Codex using ChatGPT
5.6 Sol in Ultra mode.  The exact computations reported here use integer and
rational arithmetic; their mathematical certificates are reproduced in the
body rather than delegated to software output.

These systems are disclosed as tools, not treated as mathematical
authorities or authors.  Every imported claim is checked against a publicly
resolvable human source, and every new assertion remains subject to human
verification.  This disclosure and the point-of-use citation policy follow
the recommendations of the Leiden Declaration on Artificial Intelligence and
Mathematics \cite{Leiden2026}.  Responsibility for any version submitted or
published, including the adequacy of its proofs and references, belongs to
the human author or editor who releases that version.

The explicit candidate and its original theorem claims remain attributed to
the reported Claude/Fabel source.  The purpose of the present work is narrower
but substantive: to make the construction's load-bearing maps and calculations
explicit, reproduce them independently where possible, repair identified proof
gaps when the geometry permits, and isolate the exact boundary of what still
requires outside adjudication.  This is offered as a verification aid that may
save expert reconstruction time; it is not a claim of priority for the
candidate, a transfer of authorship, or a substitute for specialist review.

\begingroup
\hbadness=1200
\begin{thebibliography}{99}

\bibitem{McMullen1996Complex}
C.~T.~McMullen,
\emph{Complex Manifolds}, Math~241 course outline, Berkeley,
Spring~1996, Section~67, p.~15,
\href{https://people.math.harvard.edu/~ctm/home/text/class/berkeley/241/96/course/course.pdf}%
{author-hosted course notes}.

\bibitem{Leiden2026}
J.~Alper, M.~Barany, A.~Chavarri Villarello, S.~Dahmen, W.~Dean,
K.~Ganapathy, M.~Harris, D.~Holmes, M.~Jamnik, S.~Kelk, B.~Kra,
U.~Martin, B.~Naskr\k{e}cki, R.~Ochigame, J.~Portegies, and J.~Schmitt,
\emph{Leiden Declaration on Artificial Intelligence and Mathematics},
2026,
\href{https://doi.org/10.5281/zenodo.20302944}{doi:10.5281/zenodo.20302944}.

\bibitem{CDP2020}
F.~Campana, J.-P.~Demailly, and T.~Peternell,
\emph{The algebraic dimension of compact complex threefolds with vanishing
second Betti number},
Compos. Math. \textbf{156} (2020), 679--696,
\href{https://doi.org/10.1112/S0010437X19007802}{doi:10.1112/S0010437X19007802}.

\bibitem{CDP1998}
F.~Campana, J.-P.~Demailly, and T.~Peternell,
\emph{The algebraic dimension of compact complex threefolds with vanishing
second Betti number},
Compos. Math. \textbf{112} (1998), 77--91,
\href{https://doi.org/10.1023/A:1000313214795}{doi:10.1023/A:1000313214795};
\href{https://arxiv.org/abs/math/9607215}{arXiv:math/9607215}.

\bibitem{RamisRuget1970}
J.-P.~Ramis and G.~Ruget,
\emph{Complexe dualisant et th\'eor\`emes de dualit\'e en g\'eom\'etrie
analytique complexe},
Publ. Math. IH\'ES \textbf{38} (1970), 77--91,
\href{https://doi.org/10.1007/BF02684652}{doi:10.1007/BF02684652}.

\bibitem{RamisRugetVerdier1971}
J.-P.~Ramis, G.~Ruget, and J.-L.~Verdier,
\emph{Dualit\'e relative en g\'eom\'etrie analytique complexe},
Invent. Math. \textbf{13} (1971), 261--283,
\href{https://doi.org/10.1007/BF01406078}{doi:10.1007/BF01406078};
\href{https://eudml.org/doc/142096}{EuDML record}.

\bibitem{BanicaStanasila1976}
C.~B\u{a}nic\u{a} and O.~St\u{a}n\u{a}\c{s}il\u{a},
\emph{Algebraic Methods in the Global Theory of Complex Spaces},
Editura Academiei and John Wiley \& Sons, Bucharest and London, 1976;
  Chapter~III, \S3, Theorem~3.4 and Corollaries~3.5 and~3.10, pp.~116--123;
  Chapter~III, \S4, Theorem~4.12, pp.~134--135,
  \href{https://djvu.online/file/2hQiOnz8WuoQk}{publicly accessible scan}.

\bibitem{Grauert1960}
H.~Grauert, \emph{Ein Theorem der analytischen Garbentheorie und die
Modulr\"aume komplexer Strukturen}, Publ. Math. IH\'ES \textbf{5} (1960),
5--64, \href{https://doi.org/10.1007/BF02684746}{doi:10.1007/BF02684746}.

\bibitem{Lee2024}
J.~M.~Lee,
\emph{Introduction to Complex Manifolds},
Graduate Studies in Mathematics \textbf{244}, American Mathematical Society,
Providence, RI, 2024; Theorem~9.37,
Proposition~6.28, p.~180, equations~(6.35)--(6.36), pp.~181--182,
and Cartan theorem~B, p.~278,
\href{https://bookstore.ams.org/gsm-244}{publisher page}.

\bibitem{Hatcher2002}
A.~Hatcher,
\emph{Algebraic Topology},
Cambridge University Press, Cambridge, 2002; Corollary~4.24, p.~360,
the opening low-dimensional homotopy-group table, p.~339,
Theorems~4.5 and~4.32, and Corollary~4.33,
\href{https://pi.math.cornell.edu/~hatcher/AT/ATpage.html}{author's electronic edition}.

\bibitem{HatcherSpectralSequences}
A.~Hatcher,
\emph{Spectral Sequences in Algebraic Topology}, Chapter~5 draft,
\S5.1; Theorem~5.41(b), p.~573, and the suspension calculation,
pp.~577--578,
\href{https://pi.math.cornell.edu/~hatcher/AT/ATch5.pdf}%
{author-hosted chapter PDF}.

\bibitem{IsaksenWangXu2023}
D.~C.~Isaksen, G.~Wang, and Z.~Xu,
\emph{Stable homotopy groups of spheres: from dimension 0 to 90},
Publ. Math. IH\'ES \textbf{137} (2023), 107--243; Table~1, p.~113,
\href{https://doi.org/10.1007/s10240-023-00139-1}%
{doi:10.1007/s10240-023-00139-1}.

\bibitem{Toda1962}
H.~Toda,
\emph{Composition Methods in Homotopy Groups of Spheres},
Annals of Mathematics Studies \textbf{49}, Princeton University Press,
Princeton, NJ, 1962; Proposition~3.1, Lemma~5.4, and
Propositions~5.6 and~5.11,
\href{https://doi.org/10.1515/9781400882625}%
{doi:10.1515/9781400882625}.

\bibitem{Bredon1997}
G.~E.~Bredon,
\emph{Sheaf Theory}, second edition,
Graduate Texts in Mathematics \textbf{170}, Springer, New York, 1997;
Chapter~VI,
\href{https://doi.org/10.1007/978-1-4612-0647-7}{doi:10.1007/978-1-4612-0647-7}.

\bibitem{Smale1961}
S.~Smale,
\emph{Generalized Poincar\'e's conjecture in dimensions greater than four},
Ann. of Math. (2) \textbf{74} (1961), 391--406; Theorem~A, p.~391,
\href{https://doi.org/10.2307/1970239}{doi:10.2307/1970239}.

\bibitem{Smale1962}
S.~Smale,
\emph{On the structure of manifolds},
Amer. J. Math. \textbf{84} (1962), 387--399;
Theorem~1.1 and Corollary~1.2,
\href{https://doi.org/10.2307/2372978}{doi:10.2307/2372978};
\href{https://www.sas.rochester.edu/mth/sites/doug-ravenel/otherpapers/Novikov-Cobordism.pdf}%
{public English reprint, p.~270}.

\bibitem{KervaireMilnor1963}
M.~A.~Kervaire and J.~W.~Milnor,
\emph{Groups of homotopy spheres: I},
Ann. of Math. (2) \textbf{77} (1963), 504--537; opening order table,
p.~504, and the following Remark, p.~505,
\href{https://doi.org/10.2307/1970128}{doi:10.2307/1970128};
\href{https://www.sas.rochester.edu/mth/sites/doug-ravenel/otherpapers/kervaire-milnor.pdf}%
{public primary scan}.

\bibitem{Hartshorne1977}
R.~Hartshorne,
\emph{Algebraic Geometry},
Graduate Texts in Mathematics \textbf{52}, Springer, New York, 1977.

\bibitem{BirkenhakeLange1999}
C.~Birkenhake and H.~Lange,
\emph{Complex Tori},
Progress in Mathematics \textbf{177}, Birkh\"auser, Boston, 1999;
Chapter~1, \S2, pp.~5--8, especially Appell--Humbert Theorem~2.1,
p.~7, and Chapter~2, \S6,
\href{https://doi.org/10.1007/978-1-4612-1566-0}{doi:10.1007/978-1-4612-1566-0}.

\bibitem{SchnellLecture3}
C.~Schnell,
\emph{MAT 615, Lecture 3: Holomorphic line bundles,
Appell--Humbert theorem},
equations~(3.3)--(3.10), pp.~3--7,
\href{https://www.math.stonybrook.edu/~cschnell/mat615/lectures/lecture3.pdf}%
{author-hosted lecture notes}.

\bibitem{SchnellLecture6}
C.~Schnell,
\emph{MAT 615, Lecture 6: Translations},
equation~(6.2), p.~1,
\href{https://www.math.stonybrook.edu/~cschnell/mat615/lectures/lecture6.pdf}%
{author-hosted lecture notes}.

\bibitem{PareschiPopaSchnell2017}
G.~Pareschi, M.~Popa, and C.~Schnell,
\emph{Hodge modules on complex tori and generic vanishing for compact
K\"ahler manifolds},
Geom. Topol. \textbf{21} (2017), 2419--2460; \S8, p.~2437,
\href{https://doi.org/10.2140/gt.2017.21.2419}{doi:10.2140/gt.2017.21.2419}.

\bibitem{Fulton1993}
W.~Fulton,
\emph{Introduction to Toric Varieties},
Annals of Mathematics Studies \textbf{131}, Princeton University Press,
Princeton, 1993; \S2.1, pp.~28--29; \S3.3, pp.~60--63; and \S3.4,
pp.~63--72,
\href{https://www.jstor.org/stable/j.ctt1b7x7vc}{JSTOR book record}.

\bibitem{CoxLittleSchenck2011}
D.~A.~Cox, J.~B.~Little, and H.~K.~Schenck,
\emph{Toric Varieties},
Graduate Studies in Mathematics \textbf{124}, American Mathematical Society,
Providence, RI, 2011; Theorem~12.1.10,
Example~1.2.21, p.~31; Propositions~4.1.1--4.1.2, p.~171;
Theorems~4.2.8, p.~179 and~4.2.12, p.~182;
\href{https://doi.org/10.1090/gsm/124}{doi:10.1090/gsm/124}.

\bibitem{Danilov1978}
V.~I.~Danilov,
\emph{The geometry of toric varieties},
Russian Math. Surveys \textbf{33} (1978), no.~2, 97--154; \S9,
\href{https://doi.org/10.1070/RM1978v033n02ABEH002305}{doi:10.1070/RM1978v033n02ABEH002305}.

\bibitem{Serre1973}
J.-P.~Serre,
\emph{A Course in Arithmetic},
Graduate Texts in Mathematics \textbf{7}, Springer, New York, 1973;
Chapter~VII,
\href{https://doi.org/10.1007/978-1-4684-9884-4}{doi:10.1007/978-1-4684-9884-4}.

\bibitem{Beardon1983}
A.~F.~Beardon,
\emph{The Geometry of Discrete Groups},
Graduate Texts in Mathematics \textbf{91}, Springer, New York, 1983;
Chapters~7 and~10,
\href{https://doi.org/10.1007/978-1-4612-1146-4}{doi:10.1007/978-1-4612-1146-4}.

\bibitem{Brown1982}
K.~S.~Brown,
\emph{Cohomology of Groups},
Graduate Texts in Mathematics \textbf{87}, Springer, New York, 1982,
\href{https://doi.org/10.1007/978-1-4684-9327-6}%
{doi:10.1007/978-1-4684-9327-6}.

\bibitem{Brown1962}
M.~Brown,
\emph{Locally flat imbeddings of topological manifolds},
Ann. of Math. (2) \textbf{75} (1962), 331--341,
\href{https://doi.org/10.2307/1970177}{doi:10.2307/1970177}.

\bibitem{Clemens1977}
C.~H.~Clemens,
\emph{Degeneration of K\"ahler manifolds},
Duke Math. J. \textbf{44} (1977), 215--290,
\href{https://doi.org/10.1215/S0012-7094-77-04410-6}%
{\nolinkurl{doi:10.1215/S0012-7094-77-04410-6}}.

\bibitem{SGA7I1972}
A.~Grothendieck, M.~Raynaud, and D.~S.~Rim (eds.),
\emph{Groupes de monodromie en g\'eom\'etrie alg\'ebrique. I},
S\'eminaire de G\'eom\'etrie Alg\'ebrique du Bois-Marie 1967--1969,
Lecture Notes in Mathematics \textbf{288}, Springer, Berlin, 1972,
\href{https://doi.org/10.1007/BFb0068688}{doi:10.1007/BFb0068688}.

\bibitem{SGA7II1973}
P.~Deligne and N.~Katz (eds.),
\emph{Groupes de monodromie en g\'eom\'etrie alg\'ebrique. II},
S\'eminaire de G\'eom\'etrie Alg\'ebrique du Bois-Marie 1967--1969,
Lecture Notes in Mathematics \textbf{340}, Springer, Berlin, 1973.

\bibitem{Dimca2004}
A.~Dimca,
\emph{Sheaves in Topology},
Universitext, Springer, Berlin, 2004; Definition~4.2.1,
Propositions~4.2.2 and~4.2.11, Theorem~2.3.26, and Corollary~2.3.4,
\href{https://doi.org/10.1007/978-3-642-18868-8}{doi:10.1007/978-3-642-18868-8}.

\bibitem{KashiwaraSchapira1990}
M.~Kashiwara and P.~Schapira,
\emph{Sheaves on Manifolds},
Grundlehren der mathematischen Wissenschaften \textbf{292}, Springer,
Berlin, 1990; \S8.6,
\href{https://doi.org/10.1007/978-3-662-02661-8}{doi:10.1007/978-3-662-02661-8}.

\bibitem{SGAEnglish2026}
A.~Grothendieck and the SGA contributors,
\emph{SGA in English: maintained working edition}, 2026,
\href{https://doi.org/10.5281/zenodo.21925497}{doi:10.5281/zenodo.21925497};
\href{https://github.com/KokunoYumeto/sga-en}{source repository}.
This working edition is cited only as a searchable reading layer; the
original volumes remain the authority.

\bibitem{ReviewerBValidation2026}
Reviewer~B,
\emph{S6 Candidate-Specific Validation Update}, independent mathematical
follow-up, 26 August 2026; seven-page report, claim ledger, finite test
program, and eight exact certificates included in this project under
\nolinkurl{supporting_materials/peer_reviews/reviewer_B_candidate_validation_2026-08-26/}.
This is an audit/provenance record, not external mathematical literature.

\bibitem{StacksProject}
The Stacks Project Authors,
\emph{The Stacks Project},
\href{https://stacks.math.columbia.edu}{stacks.math.columbia.edu},
especially Tags
\href{https://stacks.math.columbia.edu/tag/0AUW}{0AUW} and
\href{https://stacks.math.columbia.edu/tag/0EX2}{0EX2}, and also
\href{https://stacks.math.columbia.edu/tag/0CDQ}{0CDQ},
\href{https://stacks.math.columbia.edu/tag/0C6K}{0C6K},
\href{https://stacks.math.columbia.edu/tag/0B3A}{0B3A}, and
\href{https://stacks.math.columbia.edu/tag/01XT}{01XT}; the filtered-complex,
Leray, and base-change locators used here are Tags
\href{https://stacks.math.columbia.edu/tag/08J5}{08J5},
\href{https://stacks.math.columbia.edu/tag/0FLJ}{0FLJ},
\href{https://stacks.math.columbia.edu/tag/0F16}{0F16},
\href{https://stacks.math.columbia.edu/tag/012K}{012K},
\href{https://stacks.math.columbia.edu/tag/01F2}{01F2},
\href{https://stacks.math.columbia.edu/tag/011N}{011N},
\href{https://stacks.math.columbia.edu/tag/011O}{011O},
\href{https://stacks.math.columbia.edu/tag/012O}{012O},
\href{https://stacks.math.columbia.edu/tag/012U}{012U}, and
\href{https://stacks.math.columbia.edu/tag/012W}{012W}.

\bibitem{Bingener1980}
J.~Bingener,
\emph{Darstellbarkeitskriterien f\"ur analytische Funktoren},
Ann. Sci. \`Ecole Norm. Sup. (4) \textbf{13} (1980), no.~3, 317--347;
Introduction, p.~318, and Satz~6.2, pp.~336--337,
\href{https://doi.org/10.24033/asens.1385}{doi:10.24033/asens.1385};
\href{https://www.numdam.org/item/ASENS_1980_4_13_3_317_0/}%
{Numdam record and full text}.

\bibitem{Cartan1954Quotient}
H.~Cartan,
\emph{Quotient d'une vari\'et\'e analytique par un groupe discret
d'automorphismes},
S\'eminaire Henri Cartan \textbf{6} (1953--1954), Expos\'e~12, 1--13,
\href{https://www.numdam.org/item/SHC_1953-1954__6__A12_0/}{Numdam record}.

\bibitem{Cartan1957}
H.~Cartan,
\emph{Vari\'et\'es analytiques r\'eelles et vari\'et\'es analytiques complexes},
Bull. Soc. Math. France \textbf{85} (1957), 77--99; \S2,
fundamental theorem, p.~81,
\href{https://doi.org/10.24033/bsmf.1481}{doi:10.24033/bsmf.1481};
\href{https://www.numdam.org/articles/10.24033/bsmf.1481/}%
{Numdam primary text}.

\bibitem{Deligne1970}
P.~Deligne,
\emph{\'Equations diff\'erentielles \`a points singuliers r\'eguliers},
Lecture Notes in Mathematics, vol.~163, Springer, 1970,
\href{https://doi.org/10.1007/BFb0061194}{doi:10.1007/BFb0061194};
\href{https://publications.ias.edu/sites/default/files/Number9.pdf}%
{author-hosted scan}.

\bibitem{Frolicher1955}
A.~Fr\"olicher,
\emph{Relations between the cohomology groups of Dolbeault and topological
invariants},
Proc. Natl. Acad. Sci. USA \textbf{41} (1955), no.~9, 641--644,
\href{https://doi.org/10.1073/pnas.41.9.641}{doi:10.1073/pnas.41.9.641}.

\bibitem{AtiyahSinger1963}
M.~F.~Atiyah and I.~M.~Singer,
\emph{The index of elliptic operators on compact manifolds},
Bull. Amer. Math. Soc. \textbf{69} (1963), 422--433;
Theorem~3, p.~427,
\href{https://doi.org/10.1090/S0002-9904-1963-10957-X}%
{doi:10.1090/S0002-9904-1963-10957-X}.

\bibitem{Serre1955Duality}
J.-P.~Serre,
\emph{Un th\'eor\`eme de dualit\'e},
Comment. Math. Helv. \textbf{29} (1955), 9--26;
\S13, Theorem~4, p.~22,
\href{https://doi.org/10.1007/BF02564268}{doi:10.1007/BF02564268}.

\bibitem{AngellaTomassini2013}
D.~Angella and A.~Tomassini,
\emph{On the $\partial\bar\partial$-Lemma and Bott--Chern cohomology},
Invent. Math. \textbf{192} (2013), no.~1, 71--81,
\href{https://doi.org/10.1007/s00222-012-0406-3}%
{doi:10.1007/s00222-012-0406-3};
\href{https://arxiv.org/abs/1402.1954}{arXiv:1402.1954}.

\bibitem{Grothendieck1957}
A.~Grothendieck,
\emph{Sur la classification des fibr\'es holomorphes sur la sph\`ere de
Riemann},
Amer. J. Math. \textbf{79} (1957), no.~1, 121--138,
\href{https://doi.org/10.2307/2372388}{doi:10.2307/2372388}.

\bibitem{KatzOda1968}
N.~M.~Katz and T.~Oda,
\emph{On the differentiation of de Rham cohomology classes with respect to
parameters},
J. Math. Kyoto Univ. \textbf{8} (1968), no.~2, 199--213,
\href{https://doi.org/10.1215/kjm/1250524135}%
{doi:10.1215/kjm/1250524135}.

\bibitem{Angella2018}
D.~Angella,
\emph{Hodge numbers of a hypothetical complex structure on $S^6$},
Differential Geom. Appl. \textbf{57} (2018), 105--120,
\href{https://doi.org/10.1016/j.difgeo.2017.10.012}%
{doi:10.1016/j.difgeo.2017.10.012};
\href{https://arxiv.org/abs/1705.10518}{arXiv:1705.10518}.

\bibitem{McHugh2019}
A.~McHugh,
\emph{Bott--Chern--Aeppli, Dolbeault, and Fr\"olicher on compact complex
$3$-folds},
Riv. Mat. Univ. Parma (N.S.) \textbf{10} (2019), no.~1, 25--62,
\href{https://www.rivmat.unipr.it/vols/2019-10-1/02-mchugh.html}%
{publisher record};
\href{https://arxiv.org/abs/1708.03251}{arXiv:1708.03251}.

\bibitem{Schweitzer2007}
M.~Schweitzer,
\emph{Autour de la cohomologie de Bott--Chern},
preprint (2007), \S2.c;
\href{https://arxiv.org/abs/0709.3528}{arXiv:0709.3528}.

\bibitem{Ugarte2000}
L.~Ugarte,
\emph{Hodge numbers of a hypothetical complex structure on the six sphere},
Geom. Dedicata \textbf{81} (2000), 173--179,
\href{https://doi.org/10.1023/A:1005236308351}%
{doi:10.1023/A:1005236308351}.

\bibitem{Borovoi_2001}
M.~Borovoi,
\emph{On representations of integers by indefinite ternary quadratic forms},
J. Number Theory \textbf{90} (2001), no.~2, 281--293,
\href{https://doi.org/10.1006/jnth.2001.2662}%
{doi:10.1006/jnth.2001.2662};
\href{https://arxiv.org/abs/math/0006141}{arXiv:math/0006141}.

\bibitem{GorodnikOhShah_2009}
A.~Gorodnik, H.~Oh, and N.~Shah,
\emph{Integral points on symmetric varieties and Satake compactifications},
Amer. J. Math. \textbf{131} (2009), no.~1, 1--57,
\href{https://doi.org/10.1353/ajm.0.0034}{doi:10.1353/ajm.0.0034};
\href{https://arxiv.org/abs/math/0610497}{arXiv:math/0610497}.

\bibitem{BenoistOh_2012}
Y.~Benoist and H.~Oh,
\emph{Effective equidistribution of (S)-integral points on symmetric
varieties},
Ann. Inst. Fourier (Grenoble) \textbf{62} (2012), no.~5, 1889--1942;
Definition~13.1, Theorem~12.5, and Theorem~13.2,
\href{https://doi.org/10.5802/aif.2738}{doi:10.5802/aif.2738};
\href{https://arxiv.org/abs/0706.1621}{arXiv:0706.1621}.

\bibitem{Good_1983}
A.~Good,
\emph{Local Analysis of Selberg's Trace Formula},
Lecture Notes in Mathematics, vol.~1040, Springer, Berlin, 1983;
Theorem~4, p.~116,
\href{https://doi.org/10.1007/BFb0073074}{doi:10.1007/BFb0073074}.

\bibitem{ChatzakosPetridis_2016}
D.~Chatzakos and Y.~N.~Petridis,
\emph{The hyperbolic lattice point problem in conjugacy classes},
Forum Math. \textbf{28} (2016), no.~5, 981--1003;
Theorem~1.1 and Conjecture~5.7,
\href{https://doi.org/10.1515/forum-2015-0102}%
{doi:10.1515/forum-2015-0102};
\href{https://arxiv.org/abs/1504.01307}{arXiv:1504.01307}.

\bibitem{LiuSarnak_2010}
J.~Liu and P.~Sarnak,
\emph{Integral points on quadrics in three variables whose coordinates have
few prime factors},
Israel J. Math. \textbf{178} (2010), 393--426;
Proposition~3.2 and equations~(3.60), (3.65),
\href{https://doi.org/10.1007/s11856-010-0069-y}%
{doi:10.1007/s11856-010-0069-y}.

\bibitem{KimSarnak_2003}
H.~H.~Kim, with Appendix~2 by H.~H.~Kim and P.~Sarnak,
\emph{Functoriality for the exterior square of $\mathrm{GL}_4$ and the
symmetric fourth of $\mathrm{GL}_2$},
J. Amer. Math. Soc. \textbf{16} (2003), 139--183;
Appendix~2, Proposition~2,
\href{https://doi.org/10.1090/S0894-0347-02-00410-1}%
{doi:10.1090/S0894-0347-02-00410-1}.

\bibitem{ChuLi_2016}
M.~Chu and H.~Li,
\emph{Small generators of cocompact arithmetic Fuchsian groups},
Proc. Amer. Math. Soc. \textbf{144} (2016), 5121--5127;
\S3.1 and equation~(5.2),
\href{https://doi.org/10.1090/proc/13177}{doi:10.1090/proc/13177};
\href{https://arxiv.org/abs/1603.00926}{arXiv:1603.00926}.

\bibitem{Voight_2021}
J.~Voight,
\emph{Quaternion Algebras}, Graduate Texts in Mathematics, vol.~288,
Springer, Cham, 2021,
\href{https://doi.org/10.1007/978-3-030-56694-4}%
{doi:10.1007/978-3-030-56694-4}.

\bibitem{Husemoller1994}
D.~Husemoller,
\emph{Fibre Bundles}, third edition,
Graduate Texts in Mathematics, vol.~20, Springer, New York, 1994;
Chapter~17, \S7, Theorem~7.2, pp.~274--275, and
Chapter~20, \S9, Theorem~9.6 and Corollary~9.8, pp.~307--308,
\href{https://doi.org/10.1007/978-1-4757-2261-1}%
{doi:10.1007/978-1-4757-2261-1}.

\bibitem{KonstantisParton2018}
P.~Konstantis and M.~Parton,
\emph{Almost complex structures on spheres},
Differential Geom. Appl. \textbf{57} (2018), 10--22;
Lemma~4.1, Theorem~4.2, and \S6,
\href{https://doi.org/10.1016/j.difgeo.2017.10.010}%
{doi:10.1016/j.difgeo.2017.10.010};
\href{https://arxiv.org/abs/1707.03883}{arXiv:1707.03883}.

\bibitem{BorelSerre1953}
A.~Borel and J.-P.~Serre,
\emph{Groupes de Lie et puissances r\'eduites de Steenrod},
Amer. J. Math. \textbf{75} (1953), no.~3, 409--448;
\S12, Proposition~12.4, Lemma~12.5, and Proposition~12.6, p.~430;
\S15, Proposition~15.1 and the following remark, p.~434,
\href{https://doi.org/10.2307/2372495}{doi:10.2307/2372495};
\href{https://www.maths.ed.ac.uk/~v1ranick/papers/borelserre.pdf}%
{accessible primary scan}.

\bibitem{JaffeWitten2000}
A.~Jaffe and E.~Witten,
\emph{Quantum Yang--Mills theory}, in
\emph{The Millennium Prize Problems}, Clay Mathematics Institute,
Cambridge, MA, 2006, pp.~129--152; official problem description first
released in 2000,
\href{https://www.claymath.org/wp-content/uploads/2022/06/yangmills.pdf}%
{Clay Mathematics Institute PDF}.

\bibitem{Wilson_1974}
\begingroup
\setlength{\emergencystretch}{3em}
\sloppy
K.~G.~Wilson,
\emph{Confinement of quarks},
Phys. Rev. D \textbf{10} (1974), no.~8, 2445--2459,
\url{https://doi.org/10.1103/PhysRevD.10.2445}.
\par\endgroup

\bibitem{Luscher_1977}
M.~L\"uscher,
\emph{Construction of a selfadjoint, strictly positive transfer matrix for
Euclidean lattice gauge theories},
Commun. Math. Phys. \textbf{54} (1977), no.~3, 283--292,
\href{https://doi.org/10.1007/BF01614090}{doi:10.1007/BF01614090}.

\bibitem{MenottiPelissetto_1987}
P.~Menotti and A.~Pelissetto,
\emph{General proof of Osterwalder--Schrader positivity for the Wilson
action},
Commun. Math. Phys. \textbf{113} (1987), 369--373,
\href{https://doi.org/10.1007/BF01221251}{doi:10.1007/BF01221251}.

\bibitem{OsterwalderSchrader1973}
K.~Osterwalder and R.~Schrader,
\emph{Axioms for Euclidean Green's functions},
Commun. Math. Phys. \textbf{31} (1973), 83--112,
\href{https://doi.org/10.1007/BF01645738}{doi:10.1007/BF01645738};
\href{https://projecteuclid.org/euclid.cmp/1103858969}%
{Project Euclid primary text}.

\bibitem{OsterwalderSchrader1975}
K.~Osterwalder and R.~Schrader,
\emph{Axioms for Euclidean Green's functions II},
Commun. Math. Phys. \textbf{42} (1975), 281--305,
\href{https://doi.org/10.1007/BF01608978}{doi:10.1007/BF01608978};
\href{https://projecteuclid.org/euclid.cmp/1103899050}%
{Project Euclid primary text}.

\bibitem{ClankersES2026}
The Clankers,
\emph{Erd\H{o}s--Straus Project Archive: Residual Divisor Shells, Quotient
Surpluses, and Finite Verification}, version 2026.08.31.1, Zenodo (2026),
\href{https://doi.org/10.5281/zenodo.22209908}%
{doi:10.5281/zenodo.22209908}.

\bibitem{Dirac1950}
P.~A.~M.~Dirac,
\emph{Generalized Hamiltonian dynamics},
Canad. J. Math. \textbf{2} (1950), 129--148,
\href{https://doi.org/10.4153/CJM-1950-012-1}%
{doi:10.4153/CJM-1950-012-1}.

\bibitem{GotayNesterHinds1978}
M.~J.~Gotay, J.~M.~Nester, and G.~Hinds,
\emph{Presymplectic manifolds and the Dirac--Bergmann theory of constraints},
J. Math. Phys. \textbf{19} (1978), 2388--2399,
\href{https://doi.org/10.1063/1.523597}{doi:10.1063/1.523597}.

\bibitem{MarsdenWeinstein1974}
J.~Marsden and A.~Weinstein,
\emph{Reduction of symplectic manifolds with symmetry},
Rep. Math. Phys. \textbf{5} (1974), 121--130,
\href{https://doi.org/10.1016/0034-4877(74)90021-4}%
{doi:10.1016/0034-4877(74)90021-4}.

% Bibliography entries for es_niemeier_shared.tex.
% Input once inside the host's existing thebibliography environment.
% This file deliberately does not open or close that environment.

\bibitem{esn:CDH2014}
M.~C.~N.~Cheng, J.~F.~R.~Duncan, and J.~A.~Harvey,
\emph{Umbral Moonshine and Niemeier Lattices},
Research in the Mathematical Sciences \textbf{1} (2014), Article~3;
Section~2.4, definition of the Weyl quotient and Table~3,
\url{https://arxiv.org/abs/1307.5793v2}.

\bibitem{esn:MareWillems2013}
A.-L.~Mare and M.~Willems,
\emph{Topology of the octonionic flag manifold},
M\"unster Journal of Mathematics \textbf{6} (2013), 483--523;
Definition~2.2, Propositions~2.4 and~2.7, pp.~491--492,
and equations (5.4)--(5.6), pp.~510--511,
\url{https://www.uni-muenster.de/FB10/mjm/vol_6/mjm_vol_6_14.pdf};
\url{https://arxiv.org/abs/0809.4318}.

\bibitem{esn:Residual2026}
The Clankers,
\emph{Residual divisor shells, support-homogeneous signs, and the
\(A_5^4D_4\) Niemeier refinement of the modulo-\(840\)
Erd\H{o}s--Straus datum},
April 2026, maintained-source snapshot of 5 August 2026;
theorems \texttt{thm:reduced-shell}, \texttt{thm:isotropic},
\texttt{thm:root-preservation}, and \texttt{thm:two-channel}.
The exact source title, statement excerpts, snapshot digest, and
counting correction are included in the accompanying
\emph{ES--Niemeier source provenance} \cite{esn:Provenance}.

\bibitem{esn:ESMain2026}
\emph{Exact support--glue--Fricke transport},
theorem \nolinkurl{thm:exact-support-glue-Fricke-transport}
in the maintained Erd\H{o}s--Straus Foundation source,
inspected 5 September 2026.
The exact earlier-glue generators and transporter, with their original
equation labels, are reproduced in the accompanying
\emph{ES--Niemeier source provenance} \cite{esn:Provenance}.

\bibitem{esn:Provenance}
\emph{ES--Niemeier source provenance},
companion document to \emph{The modulo-840 lattice, triality, and exact
shell maps}, 5 September 2026;
\nolinkurl{ES_NIEMEIER_PROVENANCE.md}, included with this module.
It supplies exact passages and labels from the maintained residual
note, the support--glue theorem, and the canonical
\emph{Higher torus wrapping and the octonionic 24-dimensional continuation}.


\bibitem{esn:ConwaySloane1982}
J.~H.~Conway and N.~J.~A.~Sloane,
\emph{Twenty-three constructions for the Leech lattice},
Proc. R. Soc. Lond. A \textbf{381} (1982), 275--283, pp.~277--280,
\url{https://doi.org/10.1098/rspa.1982.0071}.

\bibitem{esn:Borcherds1985}
R.~E.~Borcherds, \emph{The Leech lattice},
Proc. R. Soc. Lond. A \textbf{398} (1985), 365--376,
Sections~6--7,
\url{https://math.berkeley.edu/~reb/papers/leech/leech.pdf}.

\bibitem{esn:Wilson2009}
R.~A.~Wilson, \emph{Octonions and the Leech lattice},
Journal of Algebra \textbf{322} (2009), 2186--2190,
\url{https://doi.org/10.1016/j.jalgebra.2009.03.021}.
Author version submitted 18 December 2008, revised 20 March 2009,
Sections~2--4, preprint pp.~2--6,
\url{https://webspace.maths.qmul.ac.uk/r.a.wilson/pubs_files/octoLeech1rev.pdf}.

\bibitem{sgd:Milnor1965}
J.~W.~Milnor, \emph{Topology from the Differentiable Viewpoint},
University Press of Virginia, 1965, \S1, pp.~1, 4--5.
\url{https://www.maths.ed.ac.uk/~aar/papers/milnortop.pdf}.
\end{thebibliography}

\endgroup

\end{document}
